\documentclass[spanish, 11pt]{exam}

%These tell TeX which packages to use.
\usepackage{array,epsfig}
\usepackage{amsmath}
\usepackage{amsfonts}
\usepackage{amssymb}
\usepackage{amsxtra}
\usepackage{amsthm}
\usepackage{mathrsfs}
\usepackage{color}
\usepackage{multicol}
\usepackage{verbatim}

\usepackage[utf8]{inputenc}
\usepackage[spanish]{babel}
\usepackage{eurosym}

\usepackage{graphicx}
\graphicspath{{../img/}}


\printanswers
\nopointsinmargin
\pointformat{}

%Pagination stuff.
%\setlength{\topmargin}{-.3 in}
%\setlength{\oddsidemargin}{0in}
%\setlength{\evensidemargin}{0in}
%\setlength{\textheight}{9.in}
%\setlength{\textwidth}{6.5in}
%\pagestyle{empty}

\renewcommand{\solutiontitle}{\noindent\textbf{Sol:}\enspace}

\newcommand{\class}{4º Académicas}
\newcommand{\examdate}{\today}
\newcommand{\examnum}{Inecuaciones 2}
\newcommand{\tipo}{A}


\newcommand{\timelimit}{50 minutos}



\pagestyle{head}
\firstpageheader{\includegraphics[width=0.2\columnwidth]{header_left}}{\textbf{Departamento de Matemáticas\linebreak \class}\linebreak \examnum}{\includegraphics[width=0.1\columnwidth]{header_right}}
\runningheader{\class}{\examnum}{Página \thepage\ of \numpages}
\runningheadrule

\begin{document}



\begin{questions}

\question Resolver las siguientes ecuaciones polinómicas:
\begin{multicols}{2}
\begin{parts}
\part[]$x^3-5x^2+6x\leqslant 0$  
\begin{solution} $ \left(-\infty, 0\right] \cup	\left[2, 3\right]$ \end{solution}
\part[]$ 2x^3+4x^2+2x \geqslant 0 $  
\begin{solution} $ \left\{-1\right\}\cup\left[0, \infty\right)$ \end{solution}
\begin{comment}
\part[]$ x^3 - 2x^2-5x+6 \geqslant 0$  
\begin{solution} $ \left[-2, 1\right]\cup\left[3, \infty\right)$ \end{solution}
\part[]$x^4+3x^3-3x^2\leqslant11x+6 $  
\begin{solution}$\left[-3, 2\right]$ \end{solution}
\end{comment}
\part[]$x^4+6x^3\leqslant-9x^2+4x+12 $  
\begin{solution} $\left[-3, 1\right]$\end{solution}
\part[]$x^4+2x^3-12x^2+14x-5>0 $  
\begin{solution} $\left(-\infty, -5\right)\cup\left(1, \infty\right) $ \end{solution}

\end{parts}
\end{multicols}

\question Resolver las siguientes inecuaciones racionales
\begin{multicols}{3}
\begin{parts}
\part[]  $\frac{{5x - 4}}{{x + 3}}\, - \,2\, \geqslant \,\frac{{2x}}{{x + 3}}$
%from sympy.solvers.inequalities import reduce_rational_inequalities
%reduce_rational_inequalities([[(5*x-4)/(x+3)-2 >= 2*x/(x+3)]], x,relational=0)
\begin{solution} $ \left(-\infty, -3\right) \cup \left[10, \infty\right)$ \end{solution}
\part[]  $\frac{x}{{4 - 2x}}\, > \,\frac{3}{{4 - 2x}}$
\begin{solution} $ \left(2, 3\right)$ \end{solution}
\part[]  $ \frac{{{x^2} - 9}}{{x - 1}}\, \leqslant \,0$
\begin{solution} $ \left(-\infty, -3\right] \cup \left(1, 3\right]$ \end{solution}
\part[]  $\frac{{{x^2} - 1}}{{{x^2}}}\, \geqslant \,0$
\begin{solution} $\left(-\infty, -1\right] \cup \left[1, \infty\right) $ \end{solution}
\part[]  $ \frac{{{x^2} - 3x + 2}}{{4 - {x^2}}}\, \leqslant \,0$
\begin{solution} $ \left(-\infty, -2\right) \cup \left[1, 2\right) \cup \left(2, \infty\right)$ \end{solution}
\part[]  $\frac{{{x^2} - 3x - 2}}{{{x^2} + 2x + 6}}\, < \,3$
\begin{solution} $ \left(-\infty, \infty\right)$ \end{solution}


\end{parts}
\end{multicols}

\question Resolver las siguientes inecuaciones con valor absoluto:
\begin{multicols}{2}
\begin{parts}

\part[]  $\left| {3x - 1} \right|\, \leqslant \,5$
%reduce_abs_inequalities([(Abs(3*x - 1) - 5, '<=')], x)
\begin{solution} $\left[-\frac{4}{3},2\right]$\end{solution}
\part[]  $\left| {4x + 3} \right|\,\, > \,2$ 
\begin{solution} $\left(-\infty,-\frac{5}{4}\right)\cup\left(-\frac{1}{4},+\infty\right) $ \end{solution}
\part[]  $\left| {3 - 4x} \right|\, \geqslant \,5$
\begin{solution} $\left(-\infty,-\frac{1}{2}\right]\cup\left[2,+\infty\right) $ \end{solution}
\part[]  $\left| {5x - 3} \right|\, \leqslant 2$
\begin{solution} $\left[\frac{1}{5},1\right]$ \end{solution}
\end{parts}
\end{multicols}


\question Resolver los siguientes sistemas de inecuaciones con dos incógnitas:
\begin{multicols}{2}
\begin{parts}
\part[]  $\left\{ {\begin{matrix}
   {y <  - 2x + 4}  \\ 
   {y \geqslant x}  \\ 

 \end{matrix} } \right.$
\begin{solution} %% Creator: Matplotlib, PGF backend
%%
%% To include the figure in your LaTeX document, write
%%   \input{<filename>.pgf}
%%
%% Make sure the required packages are loaded in your preamble
%%   \usepackage{pgf}
%%
%% and, on pdftex
%%   \usepackage[utf8]{inputenc}\DeclareUnicodeCharacter{2212}{-}
%%
%% or, on luatex and xetex
%%   \usepackage{unicode-math}
%%
%% Figures using additional raster images can only be included by \input if
%% they are in the same directory as the main LaTeX file. For loading figures
%% from other directories you can use the `import` package
%%   \usepackage{import}
%%
%% and then include the figures with
%%   \import{<path to file>}{<filename>.pgf}
%%
%% Matplotlib used the following preamble
%%   \usepackage{fontspec}
%%   \setmainfont{DejaVuSerif.ttf}[Path=/home/hp/Mis_aplicaciones/anaconda3/lib/python3.6/site-packages/matplotlib/mpl-data/fonts/ttf/]
%%   \setsansfont{DejaVuSans.ttf}[Path=/home/hp/Mis_aplicaciones/anaconda3/lib/python3.6/site-packages/matplotlib/mpl-data/fonts/ttf/]
%%   \setmonofont{DejaVuSansMono.ttf}[Path=/home/hp/Mis_aplicaciones/anaconda3/lib/python3.6/site-packages/matplotlib/mpl-data/fonts/ttf/]
%%
\begingroup%
\makeatletter%
\begin{pgfpicture}%
\pgfpathrectangle{\pgfpointorigin}{\pgfqpoint{6.000000in}{4.000000in}}%
\pgfusepath{use as bounding box, clip}%
\begin{pgfscope}%
\pgfsetbuttcap%
\pgfsetmiterjoin%
\pgfsetlinewidth{0.000000pt}%
\definecolor{currentstroke}{rgb}{1.000000,1.000000,1.000000}%
\pgfsetstrokecolor{currentstroke}%
\pgfsetstrokeopacity{0.000000}%
\pgfsetdash{}{0pt}%
\pgfpathmoveto{\pgfqpoint{0.000000in}{0.000000in}}%
\pgfpathlineto{\pgfqpoint{6.000000in}{0.000000in}}%
\pgfpathlineto{\pgfqpoint{6.000000in}{4.000000in}}%
\pgfpathlineto{\pgfqpoint{0.000000in}{4.000000in}}%
\pgfpathclose%
\pgfusepath{}%
\end{pgfscope}%
\begin{pgfscope}%
\pgfsetbuttcap%
\pgfsetmiterjoin%
\definecolor{currentfill}{rgb}{1.000000,1.000000,1.000000}%
\pgfsetfillcolor{currentfill}%
\pgfsetlinewidth{0.000000pt}%
\definecolor{currentstroke}{rgb}{0.000000,0.000000,0.000000}%
\pgfsetstrokecolor{currentstroke}%
\pgfsetstrokeopacity{0.000000}%
\pgfsetdash{}{0pt}%
\pgfpathmoveto{\pgfqpoint{0.750000in}{0.500000in}}%
\pgfpathlineto{\pgfqpoint{5.400000in}{0.500000in}}%
\pgfpathlineto{\pgfqpoint{5.400000in}{3.520000in}}%
\pgfpathlineto{\pgfqpoint{0.750000in}{3.520000in}}%
\pgfpathclose%
\pgfusepath{fill}%
\end{pgfscope}%
\begin{pgfscope}%
\pgfpathrectangle{\pgfqpoint{0.750000in}{0.500000in}}{\pgfqpoint{4.650000in}{3.020000in}}%
\pgfusepath{clip}%
\pgfsetbuttcap%
\pgfsetmiterjoin%
\definecolor{currentfill}{rgb}{0.000000,0.000000,1.000000}%
\pgfsetfillcolor{currentfill}%
\pgfsetlinewidth{0.000000pt}%
\definecolor{currentstroke}{rgb}{0.000000,0.000000,0.000000}%
\pgfsetstrokecolor{currentstroke}%
\pgfsetstrokeopacity{0.000000}%
\pgfsetdash{}{0pt}%
\pgfpathmoveto{\pgfqpoint{0.749998in}{0.688750in}}%
\pgfpathlineto{\pgfqpoint{0.749998in}{0.783124in}}%
\pgfpathlineto{\pgfqpoint{0.895312in}{0.783124in}}%
\pgfpathlineto{\pgfqpoint{0.895312in}{0.688750in}}%
\pgfpathmoveto{\pgfqpoint{0.749998in}{0.783124in}}%
\pgfpathlineto{\pgfqpoint{0.749998in}{0.783124in}}%
\pgfpathlineto{\pgfqpoint{0.749998in}{0.877501in}}%
\pgfpathlineto{\pgfqpoint{0.895312in}{0.877501in}}%
\pgfpathlineto{\pgfqpoint{0.895312in}{0.783124in}}%
\pgfpathmoveto{\pgfqpoint{0.749998in}{0.877501in}}%
\pgfpathlineto{\pgfqpoint{0.749998in}{0.877501in}}%
\pgfpathlineto{\pgfqpoint{0.749998in}{0.971872in}}%
\pgfpathlineto{\pgfqpoint{0.895312in}{0.971872in}}%
\pgfpathlineto{\pgfqpoint{0.895312in}{0.877501in}}%
\pgfpathmoveto{\pgfqpoint{0.749998in}{0.971872in}}%
\pgfpathlineto{\pgfqpoint{0.749998in}{0.971872in}}%
\pgfpathlineto{\pgfqpoint{0.749998in}{1.066249in}}%
\pgfpathlineto{\pgfqpoint{0.895312in}{1.066249in}}%
\pgfpathlineto{\pgfqpoint{0.895312in}{0.971872in}}%
\pgfpathmoveto{\pgfqpoint{0.749998in}{1.066249in}}%
\pgfpathlineto{\pgfqpoint{0.749998in}{1.066249in}}%
\pgfpathlineto{\pgfqpoint{0.749998in}{1.160625in}}%
\pgfpathlineto{\pgfqpoint{0.895312in}{1.160625in}}%
\pgfpathlineto{\pgfqpoint{0.895312in}{1.066249in}}%
\pgfpathmoveto{\pgfqpoint{0.749998in}{1.160625in}}%
\pgfpathlineto{\pgfqpoint{0.749998in}{1.160625in}}%
\pgfpathlineto{\pgfqpoint{0.749998in}{1.254998in}}%
\pgfpathlineto{\pgfqpoint{0.895312in}{1.254998in}}%
\pgfpathlineto{\pgfqpoint{0.895312in}{1.160625in}}%
\pgfpathmoveto{\pgfqpoint{0.749998in}{1.254998in}}%
\pgfpathlineto{\pgfqpoint{0.749998in}{1.254998in}}%
\pgfpathlineto{\pgfqpoint{0.749998in}{1.349376in}}%
\pgfpathlineto{\pgfqpoint{0.895312in}{1.349376in}}%
\pgfpathlineto{\pgfqpoint{0.895312in}{1.254998in}}%
\pgfpathmoveto{\pgfqpoint{0.749998in}{1.349376in}}%
\pgfpathlineto{\pgfqpoint{0.749998in}{1.349376in}}%
\pgfpathlineto{\pgfqpoint{0.749998in}{1.443750in}}%
\pgfpathlineto{\pgfqpoint{0.895312in}{1.443750in}}%
\pgfpathlineto{\pgfqpoint{0.895312in}{1.349376in}}%
\pgfpathmoveto{\pgfqpoint{0.749998in}{1.443750in}}%
\pgfpathlineto{\pgfqpoint{0.749998in}{1.443750in}}%
\pgfpathlineto{\pgfqpoint{0.749998in}{1.538127in}}%
\pgfpathlineto{\pgfqpoint{0.895312in}{1.538127in}}%
\pgfpathlineto{\pgfqpoint{0.895312in}{1.443750in}}%
\pgfpathmoveto{\pgfqpoint{0.749998in}{1.538127in}}%
\pgfpathlineto{\pgfqpoint{0.749998in}{1.538127in}}%
\pgfpathlineto{\pgfqpoint{0.749998in}{1.632501in}}%
\pgfpathlineto{\pgfqpoint{0.895312in}{1.632501in}}%
\pgfpathlineto{\pgfqpoint{0.895312in}{1.538127in}}%
\pgfpathmoveto{\pgfqpoint{0.749998in}{1.632501in}}%
\pgfpathlineto{\pgfqpoint{0.749998in}{1.632501in}}%
\pgfpathlineto{\pgfqpoint{0.749998in}{1.726876in}}%
\pgfpathlineto{\pgfqpoint{0.895312in}{1.726876in}}%
\pgfpathlineto{\pgfqpoint{0.895312in}{1.632501in}}%
\pgfpathmoveto{\pgfqpoint{0.749998in}{1.726876in}}%
\pgfpathlineto{\pgfqpoint{0.749998in}{1.726876in}}%
\pgfpathlineto{\pgfqpoint{0.749998in}{1.821249in}}%
\pgfpathlineto{\pgfqpoint{0.895312in}{1.821249in}}%
\pgfpathlineto{\pgfqpoint{0.895312in}{1.726876in}}%
\pgfpathmoveto{\pgfqpoint{0.749998in}{1.821249in}}%
\pgfpathlineto{\pgfqpoint{0.749998in}{1.821249in}}%
\pgfpathlineto{\pgfqpoint{0.749998in}{1.915627in}}%
\pgfpathlineto{\pgfqpoint{0.895312in}{1.915627in}}%
\pgfpathlineto{\pgfqpoint{0.895312in}{1.821249in}}%
\pgfpathmoveto{\pgfqpoint{0.749998in}{1.915627in}}%
\pgfpathlineto{\pgfqpoint{0.749998in}{1.915627in}}%
\pgfpathlineto{\pgfqpoint{0.749998in}{2.010000in}}%
\pgfpathlineto{\pgfqpoint{0.895312in}{2.010000in}}%
\pgfpathlineto{\pgfqpoint{0.895312in}{1.915627in}}%
\pgfpathmoveto{\pgfqpoint{0.749998in}{2.010000in}}%
\pgfpathlineto{\pgfqpoint{0.749998in}{2.010000in}}%
\pgfpathlineto{\pgfqpoint{0.749998in}{2.104377in}}%
\pgfpathlineto{\pgfqpoint{0.895312in}{2.104377in}}%
\pgfpathlineto{\pgfqpoint{0.895312in}{2.010000in}}%
\pgfpathmoveto{\pgfqpoint{0.749998in}{2.104377in}}%
\pgfpathlineto{\pgfqpoint{0.749998in}{2.104377in}}%
\pgfpathlineto{\pgfqpoint{0.749998in}{2.198748in}}%
\pgfpathlineto{\pgfqpoint{0.895312in}{2.198748in}}%
\pgfpathlineto{\pgfqpoint{0.895312in}{2.104377in}}%
\pgfpathmoveto{\pgfqpoint{0.749998in}{2.198748in}}%
\pgfpathlineto{\pgfqpoint{0.749998in}{2.198748in}}%
\pgfpathlineto{\pgfqpoint{0.749998in}{2.293122in}}%
\pgfpathlineto{\pgfqpoint{0.895312in}{2.293122in}}%
\pgfpathlineto{\pgfqpoint{0.895312in}{2.198748in}}%
\pgfpathmoveto{\pgfqpoint{0.749998in}{2.293122in}}%
\pgfpathlineto{\pgfqpoint{0.749998in}{2.293122in}}%
\pgfpathlineto{\pgfqpoint{0.749998in}{2.387499in}}%
\pgfpathlineto{\pgfqpoint{0.895312in}{2.387499in}}%
\pgfpathlineto{\pgfqpoint{0.895312in}{2.293122in}}%
\pgfpathmoveto{\pgfqpoint{0.749998in}{2.387499in}}%
\pgfpathlineto{\pgfqpoint{0.749998in}{2.387499in}}%
\pgfpathlineto{\pgfqpoint{0.749998in}{2.481878in}}%
\pgfpathlineto{\pgfqpoint{0.895312in}{2.481878in}}%
\pgfpathlineto{\pgfqpoint{0.895312in}{2.387499in}}%
\pgfpathmoveto{\pgfqpoint{0.749998in}{2.481878in}}%
\pgfpathlineto{\pgfqpoint{0.749998in}{2.481878in}}%
\pgfpathlineto{\pgfqpoint{0.749998in}{2.576251in}}%
\pgfpathlineto{\pgfqpoint{0.895312in}{2.576251in}}%
\pgfpathlineto{\pgfqpoint{0.895312in}{2.481878in}}%
\pgfpathmoveto{\pgfqpoint{0.749998in}{2.576251in}}%
\pgfpathlineto{\pgfqpoint{0.749998in}{2.576251in}}%
\pgfpathlineto{\pgfqpoint{0.749998in}{2.670624in}}%
\pgfpathlineto{\pgfqpoint{0.895312in}{2.670624in}}%
\pgfpathlineto{\pgfqpoint{0.895312in}{2.576251in}}%
\pgfpathmoveto{\pgfqpoint{0.749998in}{2.670624in}}%
\pgfpathlineto{\pgfqpoint{0.749998in}{2.670624in}}%
\pgfpathlineto{\pgfqpoint{0.749998in}{2.764999in}}%
\pgfpathlineto{\pgfqpoint{0.895312in}{2.764999in}}%
\pgfpathlineto{\pgfqpoint{0.895312in}{2.670624in}}%
\pgfpathmoveto{\pgfqpoint{0.749998in}{2.764999in}}%
\pgfpathlineto{\pgfqpoint{0.749998in}{2.764999in}}%
\pgfpathlineto{\pgfqpoint{0.749998in}{2.859374in}}%
\pgfpathlineto{\pgfqpoint{0.895312in}{2.859374in}}%
\pgfpathlineto{\pgfqpoint{0.895312in}{2.764999in}}%
\pgfpathmoveto{\pgfqpoint{0.749998in}{2.859374in}}%
\pgfpathlineto{\pgfqpoint{0.749998in}{2.859374in}}%
\pgfpathlineto{\pgfqpoint{0.749998in}{2.953751in}}%
\pgfpathlineto{\pgfqpoint{0.895312in}{2.953751in}}%
\pgfpathlineto{\pgfqpoint{0.895312in}{2.859374in}}%
\pgfpathmoveto{\pgfqpoint{0.749998in}{2.953751in}}%
\pgfpathlineto{\pgfqpoint{0.749998in}{2.953751in}}%
\pgfpathlineto{\pgfqpoint{0.749998in}{3.048124in}}%
\pgfpathlineto{\pgfqpoint{0.895312in}{3.048124in}}%
\pgfpathlineto{\pgfqpoint{0.895312in}{2.953751in}}%
\pgfpathmoveto{\pgfqpoint{0.749998in}{3.048124in}}%
\pgfpathlineto{\pgfqpoint{0.749998in}{3.048124in}}%
\pgfpathlineto{\pgfqpoint{0.749998in}{3.142502in}}%
\pgfpathlineto{\pgfqpoint{0.895312in}{3.142502in}}%
\pgfpathlineto{\pgfqpoint{0.895312in}{3.048124in}}%
\pgfpathmoveto{\pgfqpoint{0.749998in}{3.142502in}}%
\pgfpathlineto{\pgfqpoint{0.749998in}{3.142502in}}%
\pgfpathlineto{\pgfqpoint{0.749998in}{3.236874in}}%
\pgfpathlineto{\pgfqpoint{0.895312in}{3.236874in}}%
\pgfpathlineto{\pgfqpoint{0.895312in}{3.142502in}}%
\pgfpathmoveto{\pgfqpoint{0.749998in}{3.236874in}}%
\pgfpathlineto{\pgfqpoint{0.749998in}{3.236874in}}%
\pgfpathlineto{\pgfqpoint{0.749998in}{3.331252in}}%
\pgfpathlineto{\pgfqpoint{0.895312in}{3.331252in}}%
\pgfpathlineto{\pgfqpoint{0.895312in}{3.236874in}}%
\pgfpathmoveto{\pgfqpoint{0.749998in}{3.331252in}}%
\pgfpathlineto{\pgfqpoint{0.749998in}{3.331252in}}%
\pgfpathlineto{\pgfqpoint{0.749998in}{3.425627in}}%
\pgfpathlineto{\pgfqpoint{0.895312in}{3.425627in}}%
\pgfpathlineto{\pgfqpoint{0.895312in}{3.331252in}}%
\pgfpathmoveto{\pgfqpoint{0.749998in}{3.425627in}}%
\pgfpathlineto{\pgfqpoint{0.749998in}{3.425627in}}%
\pgfpathlineto{\pgfqpoint{0.749998in}{3.519998in}}%
\pgfpathlineto{\pgfqpoint{0.895312in}{3.519998in}}%
\pgfpathlineto{\pgfqpoint{0.895312in}{3.425627in}}%
\pgfpathmoveto{\pgfqpoint{0.895312in}{0.688750in}}%
\pgfpathlineto{\pgfqpoint{0.895312in}{0.688750in}}%
\pgfpathlineto{\pgfqpoint{0.895312in}{0.783124in}}%
\pgfpathlineto{\pgfqpoint{1.040622in}{0.783124in}}%
\pgfpathlineto{\pgfqpoint{1.040622in}{0.688750in}}%
\pgfpathmoveto{\pgfqpoint{0.895312in}{0.783124in}}%
\pgfpathlineto{\pgfqpoint{0.895312in}{0.783124in}}%
\pgfpathlineto{\pgfqpoint{0.895312in}{0.877501in}}%
\pgfpathlineto{\pgfqpoint{1.040622in}{0.877501in}}%
\pgfpathlineto{\pgfqpoint{1.040622in}{0.783124in}}%
\pgfpathmoveto{\pgfqpoint{0.895312in}{0.877501in}}%
\pgfpathlineto{\pgfqpoint{0.895312in}{0.877501in}}%
\pgfpathlineto{\pgfqpoint{0.895312in}{0.971872in}}%
\pgfpathlineto{\pgfqpoint{1.040622in}{0.971872in}}%
\pgfpathlineto{\pgfqpoint{1.040622in}{0.877501in}}%
\pgfpathmoveto{\pgfqpoint{0.895312in}{0.971872in}}%
\pgfpathlineto{\pgfqpoint{0.895312in}{0.971872in}}%
\pgfpathlineto{\pgfqpoint{0.895312in}{1.066249in}}%
\pgfpathlineto{\pgfqpoint{1.040622in}{1.066249in}}%
\pgfpathlineto{\pgfqpoint{1.040622in}{0.971872in}}%
\pgfpathmoveto{\pgfqpoint{0.895312in}{1.066249in}}%
\pgfpathlineto{\pgfqpoint{0.895312in}{1.066249in}}%
\pgfpathlineto{\pgfqpoint{0.895312in}{1.160625in}}%
\pgfpathlineto{\pgfqpoint{1.040622in}{1.160625in}}%
\pgfpathlineto{\pgfqpoint{1.040622in}{1.066249in}}%
\pgfpathmoveto{\pgfqpoint{0.895312in}{1.160625in}}%
\pgfpathlineto{\pgfqpoint{0.895312in}{1.160625in}}%
\pgfpathlineto{\pgfqpoint{0.895312in}{1.254998in}}%
\pgfpathlineto{\pgfqpoint{1.040622in}{1.254998in}}%
\pgfpathlineto{\pgfqpoint{1.040622in}{1.160625in}}%
\pgfpathmoveto{\pgfqpoint{0.895312in}{1.254998in}}%
\pgfpathlineto{\pgfqpoint{0.895312in}{1.254998in}}%
\pgfpathlineto{\pgfqpoint{0.895312in}{1.349376in}}%
\pgfpathlineto{\pgfqpoint{1.040622in}{1.349376in}}%
\pgfpathlineto{\pgfqpoint{1.040622in}{1.254998in}}%
\pgfpathmoveto{\pgfqpoint{0.895312in}{1.349376in}}%
\pgfpathlineto{\pgfqpoint{0.895312in}{1.349376in}}%
\pgfpathlineto{\pgfqpoint{0.895312in}{1.443750in}}%
\pgfpathlineto{\pgfqpoint{1.040622in}{1.443750in}}%
\pgfpathlineto{\pgfqpoint{1.040622in}{1.349376in}}%
\pgfpathmoveto{\pgfqpoint{0.895312in}{1.443750in}}%
\pgfpathlineto{\pgfqpoint{0.895312in}{1.443750in}}%
\pgfpathlineto{\pgfqpoint{0.895312in}{1.538127in}}%
\pgfpathlineto{\pgfqpoint{1.040622in}{1.538127in}}%
\pgfpathlineto{\pgfqpoint{1.040622in}{1.443750in}}%
\pgfpathmoveto{\pgfqpoint{0.895312in}{1.538127in}}%
\pgfpathlineto{\pgfqpoint{0.895312in}{1.538127in}}%
\pgfpathlineto{\pgfqpoint{0.895312in}{1.632501in}}%
\pgfpathlineto{\pgfqpoint{1.040622in}{1.632501in}}%
\pgfpathlineto{\pgfqpoint{1.040622in}{1.538127in}}%
\pgfpathmoveto{\pgfqpoint{0.895312in}{1.632501in}}%
\pgfpathlineto{\pgfqpoint{0.895312in}{1.632501in}}%
\pgfpathlineto{\pgfqpoint{0.895312in}{1.726876in}}%
\pgfpathlineto{\pgfqpoint{1.040622in}{1.726876in}}%
\pgfpathlineto{\pgfqpoint{1.040622in}{1.632501in}}%
\pgfpathmoveto{\pgfqpoint{0.895312in}{1.726876in}}%
\pgfpathlineto{\pgfqpoint{0.895312in}{1.726876in}}%
\pgfpathlineto{\pgfqpoint{0.895312in}{1.821249in}}%
\pgfpathlineto{\pgfqpoint{1.040622in}{1.821249in}}%
\pgfpathlineto{\pgfqpoint{1.040622in}{1.726876in}}%
\pgfpathmoveto{\pgfqpoint{0.895312in}{1.821249in}}%
\pgfpathlineto{\pgfqpoint{0.895312in}{1.821249in}}%
\pgfpathlineto{\pgfqpoint{0.895312in}{1.915627in}}%
\pgfpathlineto{\pgfqpoint{1.040622in}{1.915627in}}%
\pgfpathlineto{\pgfqpoint{1.040622in}{1.821249in}}%
\pgfpathmoveto{\pgfqpoint{0.895312in}{1.915627in}}%
\pgfpathlineto{\pgfqpoint{0.895312in}{1.915627in}}%
\pgfpathlineto{\pgfqpoint{0.895312in}{2.010000in}}%
\pgfpathlineto{\pgfqpoint{1.040622in}{2.010000in}}%
\pgfpathlineto{\pgfqpoint{1.040622in}{1.915627in}}%
\pgfpathmoveto{\pgfqpoint{0.895312in}{2.010000in}}%
\pgfpathlineto{\pgfqpoint{0.895312in}{2.010000in}}%
\pgfpathlineto{\pgfqpoint{0.895312in}{2.104377in}}%
\pgfpathlineto{\pgfqpoint{1.040622in}{2.104377in}}%
\pgfpathlineto{\pgfqpoint{1.040622in}{2.010000in}}%
\pgfpathmoveto{\pgfqpoint{0.895312in}{2.104377in}}%
\pgfpathlineto{\pgfqpoint{0.895312in}{2.104377in}}%
\pgfpathlineto{\pgfqpoint{0.895312in}{2.198748in}}%
\pgfpathlineto{\pgfqpoint{1.040622in}{2.198748in}}%
\pgfpathlineto{\pgfqpoint{1.040622in}{2.104377in}}%
\pgfpathmoveto{\pgfqpoint{0.895312in}{2.198748in}}%
\pgfpathlineto{\pgfqpoint{0.895312in}{2.198748in}}%
\pgfpathlineto{\pgfqpoint{0.895312in}{2.293122in}}%
\pgfpathlineto{\pgfqpoint{1.040622in}{2.293122in}}%
\pgfpathlineto{\pgfqpoint{1.040622in}{2.198748in}}%
\pgfpathmoveto{\pgfqpoint{0.895312in}{2.293122in}}%
\pgfpathlineto{\pgfqpoint{0.895312in}{2.293122in}}%
\pgfpathlineto{\pgfqpoint{0.895312in}{2.387499in}}%
\pgfpathlineto{\pgfqpoint{1.040622in}{2.387499in}}%
\pgfpathlineto{\pgfqpoint{1.040622in}{2.293122in}}%
\pgfpathmoveto{\pgfqpoint{0.895312in}{2.387499in}}%
\pgfpathlineto{\pgfqpoint{0.895312in}{2.387499in}}%
\pgfpathlineto{\pgfqpoint{0.895312in}{2.481878in}}%
\pgfpathlineto{\pgfqpoint{1.040622in}{2.481878in}}%
\pgfpathlineto{\pgfqpoint{1.040622in}{2.387499in}}%
\pgfpathmoveto{\pgfqpoint{0.895312in}{2.481878in}}%
\pgfpathlineto{\pgfqpoint{0.895312in}{2.481878in}}%
\pgfpathlineto{\pgfqpoint{0.895312in}{2.576251in}}%
\pgfpathlineto{\pgfqpoint{1.040622in}{2.576251in}}%
\pgfpathlineto{\pgfqpoint{1.040622in}{2.481878in}}%
\pgfpathmoveto{\pgfqpoint{0.895312in}{2.576251in}}%
\pgfpathlineto{\pgfqpoint{0.895312in}{2.576251in}}%
\pgfpathlineto{\pgfqpoint{0.895312in}{2.670624in}}%
\pgfpathlineto{\pgfqpoint{1.040622in}{2.670624in}}%
\pgfpathlineto{\pgfqpoint{1.040622in}{2.576251in}}%
\pgfpathmoveto{\pgfqpoint{0.895312in}{2.670624in}}%
\pgfpathlineto{\pgfqpoint{0.895312in}{2.670624in}}%
\pgfpathlineto{\pgfqpoint{0.895312in}{2.764999in}}%
\pgfpathlineto{\pgfqpoint{1.040622in}{2.764999in}}%
\pgfpathlineto{\pgfqpoint{1.040622in}{2.670624in}}%
\pgfpathmoveto{\pgfqpoint{0.895312in}{2.764999in}}%
\pgfpathlineto{\pgfqpoint{0.895312in}{2.764999in}}%
\pgfpathlineto{\pgfqpoint{0.895312in}{2.859374in}}%
\pgfpathlineto{\pgfqpoint{1.040622in}{2.859374in}}%
\pgfpathlineto{\pgfqpoint{1.040622in}{2.764999in}}%
\pgfpathmoveto{\pgfqpoint{0.895312in}{2.859374in}}%
\pgfpathlineto{\pgfqpoint{0.895312in}{2.859374in}}%
\pgfpathlineto{\pgfqpoint{0.895312in}{2.953751in}}%
\pgfpathlineto{\pgfqpoint{1.040622in}{2.953751in}}%
\pgfpathlineto{\pgfqpoint{1.040622in}{2.859374in}}%
\pgfpathmoveto{\pgfqpoint{0.895312in}{2.953751in}}%
\pgfpathlineto{\pgfqpoint{0.895312in}{2.953751in}}%
\pgfpathlineto{\pgfqpoint{0.895312in}{3.048124in}}%
\pgfpathlineto{\pgfqpoint{1.040622in}{3.048124in}}%
\pgfpathlineto{\pgfqpoint{1.040622in}{2.953751in}}%
\pgfpathmoveto{\pgfqpoint{0.895312in}{3.048124in}}%
\pgfpathlineto{\pgfqpoint{0.895312in}{3.048124in}}%
\pgfpathlineto{\pgfqpoint{0.895312in}{3.142502in}}%
\pgfpathlineto{\pgfqpoint{1.040622in}{3.142502in}}%
\pgfpathlineto{\pgfqpoint{1.040622in}{3.048124in}}%
\pgfpathmoveto{\pgfqpoint{0.895312in}{3.142502in}}%
\pgfpathlineto{\pgfqpoint{0.895312in}{3.142502in}}%
\pgfpathlineto{\pgfqpoint{0.895312in}{3.236874in}}%
\pgfpathlineto{\pgfqpoint{1.040622in}{3.236874in}}%
\pgfpathlineto{\pgfqpoint{1.040622in}{3.142502in}}%
\pgfpathmoveto{\pgfqpoint{0.895312in}{3.236874in}}%
\pgfpathlineto{\pgfqpoint{0.895312in}{3.236874in}}%
\pgfpathlineto{\pgfqpoint{0.895312in}{3.331252in}}%
\pgfpathlineto{\pgfqpoint{1.040622in}{3.331252in}}%
\pgfpathlineto{\pgfqpoint{1.040622in}{3.236874in}}%
\pgfpathmoveto{\pgfqpoint{0.895312in}{3.331252in}}%
\pgfpathlineto{\pgfqpoint{0.895312in}{3.331252in}}%
\pgfpathlineto{\pgfqpoint{0.895312in}{3.425627in}}%
\pgfpathlineto{\pgfqpoint{1.040622in}{3.425627in}}%
\pgfpathlineto{\pgfqpoint{1.040622in}{3.331252in}}%
\pgfpathmoveto{\pgfqpoint{0.895312in}{3.425627in}}%
\pgfpathlineto{\pgfqpoint{0.895312in}{3.425627in}}%
\pgfpathlineto{\pgfqpoint{0.895312in}{3.519998in}}%
\pgfpathlineto{\pgfqpoint{1.040622in}{3.519998in}}%
\pgfpathlineto{\pgfqpoint{1.040622in}{3.425627in}}%
\pgfpathmoveto{\pgfqpoint{1.040622in}{0.877501in}}%
\pgfpathlineto{\pgfqpoint{1.040622in}{0.877501in}}%
\pgfpathlineto{\pgfqpoint{1.040622in}{0.971872in}}%
\pgfpathlineto{\pgfqpoint{1.185941in}{0.971872in}}%
\pgfpathlineto{\pgfqpoint{1.185941in}{0.877501in}}%
\pgfpathmoveto{\pgfqpoint{1.040622in}{0.971872in}}%
\pgfpathlineto{\pgfqpoint{1.040622in}{0.971872in}}%
\pgfpathlineto{\pgfqpoint{1.040622in}{1.066249in}}%
\pgfpathlineto{\pgfqpoint{1.185941in}{1.066249in}}%
\pgfpathlineto{\pgfqpoint{1.185941in}{0.971872in}}%
\pgfpathmoveto{\pgfqpoint{1.040622in}{1.066249in}}%
\pgfpathlineto{\pgfqpoint{1.040622in}{1.066249in}}%
\pgfpathlineto{\pgfqpoint{1.040622in}{1.160625in}}%
\pgfpathlineto{\pgfqpoint{1.185941in}{1.160625in}}%
\pgfpathlineto{\pgfqpoint{1.185941in}{1.066249in}}%
\pgfpathmoveto{\pgfqpoint{1.040622in}{1.160625in}}%
\pgfpathlineto{\pgfqpoint{1.040622in}{1.160625in}}%
\pgfpathlineto{\pgfqpoint{1.040622in}{1.254998in}}%
\pgfpathlineto{\pgfqpoint{1.185941in}{1.254998in}}%
\pgfpathlineto{\pgfqpoint{1.185941in}{1.160625in}}%
\pgfpathmoveto{\pgfqpoint{1.040622in}{1.254998in}}%
\pgfpathlineto{\pgfqpoint{1.040622in}{1.254998in}}%
\pgfpathlineto{\pgfqpoint{1.040622in}{1.349376in}}%
\pgfpathlineto{\pgfqpoint{1.185941in}{1.349376in}}%
\pgfpathlineto{\pgfqpoint{1.185941in}{1.254998in}}%
\pgfpathmoveto{\pgfqpoint{1.040622in}{1.349376in}}%
\pgfpathlineto{\pgfqpoint{1.040622in}{1.349376in}}%
\pgfpathlineto{\pgfqpoint{1.040622in}{1.443750in}}%
\pgfpathlineto{\pgfqpoint{1.185941in}{1.443750in}}%
\pgfpathlineto{\pgfqpoint{1.185941in}{1.349376in}}%
\pgfpathmoveto{\pgfqpoint{1.040622in}{1.443750in}}%
\pgfpathlineto{\pgfqpoint{1.040622in}{1.443750in}}%
\pgfpathlineto{\pgfqpoint{1.040622in}{1.538127in}}%
\pgfpathlineto{\pgfqpoint{1.185941in}{1.538127in}}%
\pgfpathlineto{\pgfqpoint{1.185941in}{1.443750in}}%
\pgfpathmoveto{\pgfqpoint{1.040622in}{1.538127in}}%
\pgfpathlineto{\pgfqpoint{1.040622in}{1.538127in}}%
\pgfpathlineto{\pgfqpoint{1.040622in}{1.632501in}}%
\pgfpathlineto{\pgfqpoint{1.185941in}{1.632501in}}%
\pgfpathlineto{\pgfqpoint{1.185941in}{1.538127in}}%
\pgfpathmoveto{\pgfqpoint{1.040622in}{1.632501in}}%
\pgfpathlineto{\pgfqpoint{1.040622in}{1.632501in}}%
\pgfpathlineto{\pgfqpoint{1.040622in}{1.726876in}}%
\pgfpathlineto{\pgfqpoint{1.185941in}{1.726876in}}%
\pgfpathlineto{\pgfqpoint{1.185941in}{1.632501in}}%
\pgfpathmoveto{\pgfqpoint{1.040622in}{1.726876in}}%
\pgfpathlineto{\pgfqpoint{1.040622in}{1.726876in}}%
\pgfpathlineto{\pgfqpoint{1.040622in}{1.821249in}}%
\pgfpathlineto{\pgfqpoint{1.185941in}{1.821249in}}%
\pgfpathlineto{\pgfqpoint{1.185941in}{1.726876in}}%
\pgfpathmoveto{\pgfqpoint{1.040622in}{1.821249in}}%
\pgfpathlineto{\pgfqpoint{1.040622in}{1.821249in}}%
\pgfpathlineto{\pgfqpoint{1.040622in}{1.915627in}}%
\pgfpathlineto{\pgfqpoint{1.185941in}{1.915627in}}%
\pgfpathlineto{\pgfqpoint{1.185941in}{1.821249in}}%
\pgfpathmoveto{\pgfqpoint{1.040622in}{1.915627in}}%
\pgfpathlineto{\pgfqpoint{1.040622in}{1.915627in}}%
\pgfpathlineto{\pgfqpoint{1.040622in}{2.010000in}}%
\pgfpathlineto{\pgfqpoint{1.185941in}{2.010000in}}%
\pgfpathlineto{\pgfqpoint{1.185941in}{1.915627in}}%
\pgfpathmoveto{\pgfqpoint{1.040622in}{2.010000in}}%
\pgfpathlineto{\pgfqpoint{1.040622in}{2.010000in}}%
\pgfpathlineto{\pgfqpoint{1.040622in}{2.104377in}}%
\pgfpathlineto{\pgfqpoint{1.185941in}{2.104377in}}%
\pgfpathlineto{\pgfqpoint{1.185941in}{2.010000in}}%
\pgfpathmoveto{\pgfqpoint{1.040622in}{2.104377in}}%
\pgfpathlineto{\pgfqpoint{1.040622in}{2.104377in}}%
\pgfpathlineto{\pgfqpoint{1.040622in}{2.198748in}}%
\pgfpathlineto{\pgfqpoint{1.185941in}{2.198748in}}%
\pgfpathlineto{\pgfqpoint{1.185941in}{2.104377in}}%
\pgfpathmoveto{\pgfqpoint{1.040622in}{2.198748in}}%
\pgfpathlineto{\pgfqpoint{1.040622in}{2.198748in}}%
\pgfpathlineto{\pgfqpoint{1.040622in}{2.293122in}}%
\pgfpathlineto{\pgfqpoint{1.185941in}{2.293122in}}%
\pgfpathlineto{\pgfqpoint{1.185941in}{2.198748in}}%
\pgfpathmoveto{\pgfqpoint{1.040622in}{2.293122in}}%
\pgfpathlineto{\pgfqpoint{1.040622in}{2.293122in}}%
\pgfpathlineto{\pgfqpoint{1.040622in}{2.387499in}}%
\pgfpathlineto{\pgfqpoint{1.185941in}{2.387499in}}%
\pgfpathlineto{\pgfqpoint{1.185941in}{2.293122in}}%
\pgfpathmoveto{\pgfqpoint{1.040622in}{2.387499in}}%
\pgfpathlineto{\pgfqpoint{1.040622in}{2.387499in}}%
\pgfpathlineto{\pgfqpoint{1.040622in}{2.481878in}}%
\pgfpathlineto{\pgfqpoint{1.185941in}{2.481878in}}%
\pgfpathlineto{\pgfqpoint{1.185941in}{2.387499in}}%
\pgfpathmoveto{\pgfqpoint{1.040622in}{2.481878in}}%
\pgfpathlineto{\pgfqpoint{1.040622in}{2.481878in}}%
\pgfpathlineto{\pgfqpoint{1.040622in}{2.576251in}}%
\pgfpathlineto{\pgfqpoint{1.185941in}{2.576251in}}%
\pgfpathlineto{\pgfqpoint{1.185941in}{2.481878in}}%
\pgfpathmoveto{\pgfqpoint{1.040622in}{2.576251in}}%
\pgfpathlineto{\pgfqpoint{1.040622in}{2.576251in}}%
\pgfpathlineto{\pgfqpoint{1.040622in}{2.670624in}}%
\pgfpathlineto{\pgfqpoint{1.185941in}{2.670624in}}%
\pgfpathlineto{\pgfqpoint{1.185941in}{2.576251in}}%
\pgfpathmoveto{\pgfqpoint{1.040622in}{2.670624in}}%
\pgfpathlineto{\pgfqpoint{1.040622in}{2.670624in}}%
\pgfpathlineto{\pgfqpoint{1.040622in}{2.764999in}}%
\pgfpathlineto{\pgfqpoint{1.185941in}{2.764999in}}%
\pgfpathlineto{\pgfqpoint{1.185941in}{2.670624in}}%
\pgfpathmoveto{\pgfqpoint{1.040622in}{2.764999in}}%
\pgfpathlineto{\pgfqpoint{1.040622in}{2.764999in}}%
\pgfpathlineto{\pgfqpoint{1.040622in}{2.859374in}}%
\pgfpathlineto{\pgfqpoint{1.185941in}{2.859374in}}%
\pgfpathlineto{\pgfqpoint{1.185941in}{2.764999in}}%
\pgfpathmoveto{\pgfqpoint{1.040622in}{2.859374in}}%
\pgfpathlineto{\pgfqpoint{1.040622in}{2.859374in}}%
\pgfpathlineto{\pgfqpoint{1.040622in}{2.953751in}}%
\pgfpathlineto{\pgfqpoint{1.185941in}{2.953751in}}%
\pgfpathlineto{\pgfqpoint{1.185941in}{2.859374in}}%
\pgfpathmoveto{\pgfqpoint{1.040622in}{2.953751in}}%
\pgfpathlineto{\pgfqpoint{1.040622in}{2.953751in}}%
\pgfpathlineto{\pgfqpoint{1.040622in}{3.048124in}}%
\pgfpathlineto{\pgfqpoint{1.185941in}{3.048124in}}%
\pgfpathlineto{\pgfqpoint{1.185941in}{2.953751in}}%
\pgfpathmoveto{\pgfqpoint{1.040622in}{3.048124in}}%
\pgfpathlineto{\pgfqpoint{1.040622in}{3.048124in}}%
\pgfpathlineto{\pgfqpoint{1.040622in}{3.142502in}}%
\pgfpathlineto{\pgfqpoint{1.185941in}{3.142502in}}%
\pgfpathlineto{\pgfqpoint{1.185941in}{3.048124in}}%
\pgfpathmoveto{\pgfqpoint{1.040622in}{3.142502in}}%
\pgfpathlineto{\pgfqpoint{1.040622in}{3.142502in}}%
\pgfpathlineto{\pgfqpoint{1.040622in}{3.236874in}}%
\pgfpathlineto{\pgfqpoint{1.185941in}{3.236874in}}%
\pgfpathlineto{\pgfqpoint{1.185941in}{3.142502in}}%
\pgfpathmoveto{\pgfqpoint{1.040622in}{3.236874in}}%
\pgfpathlineto{\pgfqpoint{1.040622in}{3.236874in}}%
\pgfpathlineto{\pgfqpoint{1.040622in}{3.331252in}}%
\pgfpathlineto{\pgfqpoint{1.185941in}{3.331252in}}%
\pgfpathlineto{\pgfqpoint{1.185941in}{3.236874in}}%
\pgfpathmoveto{\pgfqpoint{1.040622in}{3.331252in}}%
\pgfpathlineto{\pgfqpoint{1.040622in}{3.331252in}}%
\pgfpathlineto{\pgfqpoint{1.040622in}{3.425627in}}%
\pgfpathlineto{\pgfqpoint{1.185941in}{3.425627in}}%
\pgfpathlineto{\pgfqpoint{1.185941in}{3.331252in}}%
\pgfpathmoveto{\pgfqpoint{1.040622in}{3.425627in}}%
\pgfpathlineto{\pgfqpoint{1.040622in}{3.425627in}}%
\pgfpathlineto{\pgfqpoint{1.040622in}{3.519998in}}%
\pgfpathlineto{\pgfqpoint{1.185941in}{3.519998in}}%
\pgfpathlineto{\pgfqpoint{1.185941in}{3.425627in}}%
\pgfpathmoveto{\pgfqpoint{1.185941in}{0.971872in}}%
\pgfpathlineto{\pgfqpoint{1.185941in}{0.971872in}}%
\pgfpathlineto{\pgfqpoint{1.185941in}{1.066249in}}%
\pgfpathlineto{\pgfqpoint{1.331254in}{1.066249in}}%
\pgfpathlineto{\pgfqpoint{1.331254in}{0.971872in}}%
\pgfpathmoveto{\pgfqpoint{1.185941in}{1.066249in}}%
\pgfpathlineto{\pgfqpoint{1.185941in}{1.066249in}}%
\pgfpathlineto{\pgfqpoint{1.185941in}{1.160625in}}%
\pgfpathlineto{\pgfqpoint{1.331254in}{1.160625in}}%
\pgfpathlineto{\pgfqpoint{1.331254in}{1.066249in}}%
\pgfpathmoveto{\pgfqpoint{1.185941in}{1.160625in}}%
\pgfpathlineto{\pgfqpoint{1.185941in}{1.160625in}}%
\pgfpathlineto{\pgfqpoint{1.185941in}{1.254998in}}%
\pgfpathlineto{\pgfqpoint{1.331254in}{1.254998in}}%
\pgfpathlineto{\pgfqpoint{1.331254in}{1.160625in}}%
\pgfpathmoveto{\pgfqpoint{1.185941in}{1.254998in}}%
\pgfpathlineto{\pgfqpoint{1.185941in}{1.254998in}}%
\pgfpathlineto{\pgfqpoint{1.185941in}{1.349376in}}%
\pgfpathlineto{\pgfqpoint{1.331254in}{1.349376in}}%
\pgfpathlineto{\pgfqpoint{1.331254in}{1.254998in}}%
\pgfpathmoveto{\pgfqpoint{1.185941in}{1.349376in}}%
\pgfpathlineto{\pgfqpoint{1.185941in}{1.349376in}}%
\pgfpathlineto{\pgfqpoint{1.185941in}{1.443750in}}%
\pgfpathlineto{\pgfqpoint{1.331254in}{1.443750in}}%
\pgfpathlineto{\pgfqpoint{1.331254in}{1.349376in}}%
\pgfpathmoveto{\pgfqpoint{1.185941in}{1.443750in}}%
\pgfpathlineto{\pgfqpoint{1.185941in}{1.443750in}}%
\pgfpathlineto{\pgfqpoint{1.185941in}{1.538127in}}%
\pgfpathlineto{\pgfqpoint{1.331254in}{1.538127in}}%
\pgfpathlineto{\pgfqpoint{1.331254in}{1.443750in}}%
\pgfpathmoveto{\pgfqpoint{1.185941in}{1.538127in}}%
\pgfpathlineto{\pgfqpoint{1.185941in}{1.538127in}}%
\pgfpathlineto{\pgfqpoint{1.185941in}{1.632501in}}%
\pgfpathlineto{\pgfqpoint{1.331254in}{1.632501in}}%
\pgfpathlineto{\pgfqpoint{1.331254in}{1.538127in}}%
\pgfpathmoveto{\pgfqpoint{1.185941in}{1.632501in}}%
\pgfpathlineto{\pgfqpoint{1.185941in}{1.632501in}}%
\pgfpathlineto{\pgfqpoint{1.185941in}{1.726876in}}%
\pgfpathlineto{\pgfqpoint{1.331254in}{1.726876in}}%
\pgfpathlineto{\pgfqpoint{1.331254in}{1.632501in}}%
\pgfpathmoveto{\pgfqpoint{1.185941in}{1.726876in}}%
\pgfpathlineto{\pgfqpoint{1.185941in}{1.726876in}}%
\pgfpathlineto{\pgfqpoint{1.185941in}{1.821249in}}%
\pgfpathlineto{\pgfqpoint{1.331254in}{1.821249in}}%
\pgfpathlineto{\pgfqpoint{1.331254in}{1.726876in}}%
\pgfpathmoveto{\pgfqpoint{1.185941in}{1.821249in}}%
\pgfpathlineto{\pgfqpoint{1.185941in}{1.821249in}}%
\pgfpathlineto{\pgfqpoint{1.185941in}{1.915627in}}%
\pgfpathlineto{\pgfqpoint{1.331254in}{1.915627in}}%
\pgfpathlineto{\pgfqpoint{1.331254in}{1.821249in}}%
\pgfpathmoveto{\pgfqpoint{1.185941in}{1.915627in}}%
\pgfpathlineto{\pgfqpoint{1.185941in}{1.915627in}}%
\pgfpathlineto{\pgfqpoint{1.185941in}{2.010000in}}%
\pgfpathlineto{\pgfqpoint{1.331254in}{2.010000in}}%
\pgfpathlineto{\pgfqpoint{1.331254in}{1.915627in}}%
\pgfpathmoveto{\pgfqpoint{1.185941in}{2.010000in}}%
\pgfpathlineto{\pgfqpoint{1.185941in}{2.010000in}}%
\pgfpathlineto{\pgfqpoint{1.185941in}{2.104377in}}%
\pgfpathlineto{\pgfqpoint{1.331254in}{2.104377in}}%
\pgfpathlineto{\pgfqpoint{1.331254in}{2.010000in}}%
\pgfpathmoveto{\pgfqpoint{1.185941in}{2.104377in}}%
\pgfpathlineto{\pgfqpoint{1.185941in}{2.104377in}}%
\pgfpathlineto{\pgfqpoint{1.185941in}{2.198748in}}%
\pgfpathlineto{\pgfqpoint{1.331254in}{2.198748in}}%
\pgfpathlineto{\pgfqpoint{1.331254in}{2.104377in}}%
\pgfpathmoveto{\pgfqpoint{1.185941in}{2.198748in}}%
\pgfpathlineto{\pgfqpoint{1.185941in}{2.198748in}}%
\pgfpathlineto{\pgfqpoint{1.185941in}{2.293122in}}%
\pgfpathlineto{\pgfqpoint{1.331254in}{2.293122in}}%
\pgfpathlineto{\pgfqpoint{1.331254in}{2.198748in}}%
\pgfpathmoveto{\pgfqpoint{1.185941in}{2.293122in}}%
\pgfpathlineto{\pgfqpoint{1.185941in}{2.293122in}}%
\pgfpathlineto{\pgfqpoint{1.185941in}{2.387499in}}%
\pgfpathlineto{\pgfqpoint{1.331254in}{2.387499in}}%
\pgfpathlineto{\pgfqpoint{1.331254in}{2.293122in}}%
\pgfpathmoveto{\pgfqpoint{1.185941in}{2.387499in}}%
\pgfpathlineto{\pgfqpoint{1.185941in}{2.387499in}}%
\pgfpathlineto{\pgfqpoint{1.185941in}{2.481878in}}%
\pgfpathlineto{\pgfqpoint{1.331254in}{2.481878in}}%
\pgfpathlineto{\pgfqpoint{1.331254in}{2.387499in}}%
\pgfpathmoveto{\pgfqpoint{1.185941in}{2.481878in}}%
\pgfpathlineto{\pgfqpoint{1.185941in}{2.481878in}}%
\pgfpathlineto{\pgfqpoint{1.185941in}{2.576251in}}%
\pgfpathlineto{\pgfqpoint{1.331254in}{2.576251in}}%
\pgfpathlineto{\pgfqpoint{1.331254in}{2.481878in}}%
\pgfpathmoveto{\pgfqpoint{1.185941in}{2.576251in}}%
\pgfpathlineto{\pgfqpoint{1.185941in}{2.576251in}}%
\pgfpathlineto{\pgfqpoint{1.185941in}{2.670624in}}%
\pgfpathlineto{\pgfqpoint{1.331254in}{2.670624in}}%
\pgfpathlineto{\pgfqpoint{1.331254in}{2.576251in}}%
\pgfpathmoveto{\pgfqpoint{1.185941in}{2.670624in}}%
\pgfpathlineto{\pgfqpoint{1.185941in}{2.670624in}}%
\pgfpathlineto{\pgfqpoint{1.185941in}{2.764999in}}%
\pgfpathlineto{\pgfqpoint{1.331254in}{2.764999in}}%
\pgfpathlineto{\pgfqpoint{1.331254in}{2.670624in}}%
\pgfpathmoveto{\pgfqpoint{1.185941in}{2.764999in}}%
\pgfpathlineto{\pgfqpoint{1.185941in}{2.764999in}}%
\pgfpathlineto{\pgfqpoint{1.185941in}{2.859374in}}%
\pgfpathlineto{\pgfqpoint{1.331254in}{2.859374in}}%
\pgfpathlineto{\pgfqpoint{1.331254in}{2.764999in}}%
\pgfpathmoveto{\pgfqpoint{1.185941in}{2.859374in}}%
\pgfpathlineto{\pgfqpoint{1.185941in}{2.859374in}}%
\pgfpathlineto{\pgfqpoint{1.185941in}{2.953751in}}%
\pgfpathlineto{\pgfqpoint{1.331254in}{2.953751in}}%
\pgfpathlineto{\pgfqpoint{1.331254in}{2.859374in}}%
\pgfpathmoveto{\pgfqpoint{1.185941in}{2.953751in}}%
\pgfpathlineto{\pgfqpoint{1.185941in}{2.953751in}}%
\pgfpathlineto{\pgfqpoint{1.185941in}{3.048124in}}%
\pgfpathlineto{\pgfqpoint{1.331254in}{3.048124in}}%
\pgfpathlineto{\pgfqpoint{1.331254in}{2.953751in}}%
\pgfpathmoveto{\pgfqpoint{1.185941in}{3.048124in}}%
\pgfpathlineto{\pgfqpoint{1.185941in}{3.048124in}}%
\pgfpathlineto{\pgfqpoint{1.185941in}{3.142502in}}%
\pgfpathlineto{\pgfqpoint{1.331254in}{3.142502in}}%
\pgfpathlineto{\pgfqpoint{1.331254in}{3.048124in}}%
\pgfpathmoveto{\pgfqpoint{1.185941in}{3.142502in}}%
\pgfpathlineto{\pgfqpoint{1.185941in}{3.142502in}}%
\pgfpathlineto{\pgfqpoint{1.185941in}{3.236874in}}%
\pgfpathlineto{\pgfqpoint{1.331254in}{3.236874in}}%
\pgfpathlineto{\pgfqpoint{1.331254in}{3.142502in}}%
\pgfpathmoveto{\pgfqpoint{1.185941in}{3.236874in}}%
\pgfpathlineto{\pgfqpoint{1.185941in}{3.236874in}}%
\pgfpathlineto{\pgfqpoint{1.185941in}{3.331252in}}%
\pgfpathlineto{\pgfqpoint{1.331254in}{3.331252in}}%
\pgfpathlineto{\pgfqpoint{1.331254in}{3.236874in}}%
\pgfpathmoveto{\pgfqpoint{1.185941in}{3.331252in}}%
\pgfpathlineto{\pgfqpoint{1.185941in}{3.331252in}}%
\pgfpathlineto{\pgfqpoint{1.185941in}{3.425627in}}%
\pgfpathlineto{\pgfqpoint{1.331254in}{3.425627in}}%
\pgfpathlineto{\pgfqpoint{1.331254in}{3.331252in}}%
\pgfpathmoveto{\pgfqpoint{1.185941in}{3.425627in}}%
\pgfpathlineto{\pgfqpoint{1.185941in}{3.425627in}}%
\pgfpathlineto{\pgfqpoint{1.185941in}{3.519998in}}%
\pgfpathlineto{\pgfqpoint{1.331254in}{3.519998in}}%
\pgfpathlineto{\pgfqpoint{1.331254in}{3.425627in}}%
\pgfpathmoveto{\pgfqpoint{1.331254in}{1.066249in}}%
\pgfpathlineto{\pgfqpoint{1.331254in}{1.066249in}}%
\pgfpathlineto{\pgfqpoint{1.331254in}{1.160625in}}%
\pgfpathlineto{\pgfqpoint{1.476564in}{1.160625in}}%
\pgfpathlineto{\pgfqpoint{1.476564in}{1.066249in}}%
\pgfpathmoveto{\pgfqpoint{1.331254in}{1.160625in}}%
\pgfpathlineto{\pgfqpoint{1.331254in}{1.160625in}}%
\pgfpathlineto{\pgfqpoint{1.331254in}{1.254998in}}%
\pgfpathlineto{\pgfqpoint{1.476564in}{1.254998in}}%
\pgfpathlineto{\pgfqpoint{1.476564in}{1.160625in}}%
\pgfpathmoveto{\pgfqpoint{1.331254in}{1.254998in}}%
\pgfpathlineto{\pgfqpoint{1.331254in}{1.254998in}}%
\pgfpathlineto{\pgfqpoint{1.331254in}{1.349376in}}%
\pgfpathlineto{\pgfqpoint{1.476564in}{1.349376in}}%
\pgfpathlineto{\pgfqpoint{1.476564in}{1.254998in}}%
\pgfpathmoveto{\pgfqpoint{1.331254in}{1.349376in}}%
\pgfpathlineto{\pgfqpoint{1.331254in}{1.349376in}}%
\pgfpathlineto{\pgfqpoint{1.331254in}{1.443750in}}%
\pgfpathlineto{\pgfqpoint{1.476564in}{1.443750in}}%
\pgfpathlineto{\pgfqpoint{1.476564in}{1.349376in}}%
\pgfpathmoveto{\pgfqpoint{1.331254in}{1.443750in}}%
\pgfpathlineto{\pgfqpoint{1.331254in}{1.443750in}}%
\pgfpathlineto{\pgfqpoint{1.331254in}{1.538127in}}%
\pgfpathlineto{\pgfqpoint{1.476564in}{1.538127in}}%
\pgfpathlineto{\pgfqpoint{1.476564in}{1.443750in}}%
\pgfpathmoveto{\pgfqpoint{1.331254in}{1.538127in}}%
\pgfpathlineto{\pgfqpoint{1.331254in}{1.538127in}}%
\pgfpathlineto{\pgfqpoint{1.331254in}{1.632501in}}%
\pgfpathlineto{\pgfqpoint{1.476564in}{1.632501in}}%
\pgfpathlineto{\pgfqpoint{1.476564in}{1.538127in}}%
\pgfpathmoveto{\pgfqpoint{1.331254in}{1.632501in}}%
\pgfpathlineto{\pgfqpoint{1.331254in}{1.632501in}}%
\pgfpathlineto{\pgfqpoint{1.331254in}{1.726876in}}%
\pgfpathlineto{\pgfqpoint{1.476564in}{1.726876in}}%
\pgfpathlineto{\pgfqpoint{1.476564in}{1.632501in}}%
\pgfpathmoveto{\pgfqpoint{1.331254in}{1.726876in}}%
\pgfpathlineto{\pgfqpoint{1.331254in}{1.726876in}}%
\pgfpathlineto{\pgfqpoint{1.331254in}{1.821249in}}%
\pgfpathlineto{\pgfqpoint{1.476564in}{1.821249in}}%
\pgfpathlineto{\pgfqpoint{1.476564in}{1.726876in}}%
\pgfpathmoveto{\pgfqpoint{1.331254in}{1.821249in}}%
\pgfpathlineto{\pgfqpoint{1.331254in}{1.821249in}}%
\pgfpathlineto{\pgfqpoint{1.331254in}{1.915627in}}%
\pgfpathlineto{\pgfqpoint{1.476564in}{1.915627in}}%
\pgfpathlineto{\pgfqpoint{1.476564in}{1.821249in}}%
\pgfpathmoveto{\pgfqpoint{1.331254in}{1.915627in}}%
\pgfpathlineto{\pgfqpoint{1.331254in}{1.915627in}}%
\pgfpathlineto{\pgfqpoint{1.331254in}{2.010000in}}%
\pgfpathlineto{\pgfqpoint{1.476564in}{2.010000in}}%
\pgfpathlineto{\pgfqpoint{1.476564in}{1.915627in}}%
\pgfpathmoveto{\pgfqpoint{1.331254in}{2.010000in}}%
\pgfpathlineto{\pgfqpoint{1.331254in}{2.010000in}}%
\pgfpathlineto{\pgfqpoint{1.331254in}{2.104377in}}%
\pgfpathlineto{\pgfqpoint{1.476564in}{2.104377in}}%
\pgfpathlineto{\pgfqpoint{1.476564in}{2.010000in}}%
\pgfpathmoveto{\pgfqpoint{1.331254in}{2.104377in}}%
\pgfpathlineto{\pgfqpoint{1.331254in}{2.104377in}}%
\pgfpathlineto{\pgfqpoint{1.331254in}{2.198748in}}%
\pgfpathlineto{\pgfqpoint{1.476564in}{2.198748in}}%
\pgfpathlineto{\pgfqpoint{1.476564in}{2.104377in}}%
\pgfpathmoveto{\pgfqpoint{1.331254in}{2.198748in}}%
\pgfpathlineto{\pgfqpoint{1.331254in}{2.198748in}}%
\pgfpathlineto{\pgfqpoint{1.331254in}{2.293122in}}%
\pgfpathlineto{\pgfqpoint{1.476564in}{2.293122in}}%
\pgfpathlineto{\pgfqpoint{1.476564in}{2.198748in}}%
\pgfpathmoveto{\pgfqpoint{1.331254in}{2.293122in}}%
\pgfpathlineto{\pgfqpoint{1.331254in}{2.293122in}}%
\pgfpathlineto{\pgfqpoint{1.331254in}{2.387499in}}%
\pgfpathlineto{\pgfqpoint{1.476564in}{2.387499in}}%
\pgfpathlineto{\pgfqpoint{1.476564in}{2.293122in}}%
\pgfpathmoveto{\pgfqpoint{1.331254in}{2.387499in}}%
\pgfpathlineto{\pgfqpoint{1.331254in}{2.387499in}}%
\pgfpathlineto{\pgfqpoint{1.331254in}{2.481878in}}%
\pgfpathlineto{\pgfqpoint{1.476564in}{2.481878in}}%
\pgfpathlineto{\pgfqpoint{1.476564in}{2.387499in}}%
\pgfpathmoveto{\pgfqpoint{1.331254in}{2.481878in}}%
\pgfpathlineto{\pgfqpoint{1.331254in}{2.481878in}}%
\pgfpathlineto{\pgfqpoint{1.331254in}{2.576251in}}%
\pgfpathlineto{\pgfqpoint{1.476564in}{2.576251in}}%
\pgfpathlineto{\pgfqpoint{1.476564in}{2.481878in}}%
\pgfpathmoveto{\pgfqpoint{1.331254in}{2.576251in}}%
\pgfpathlineto{\pgfqpoint{1.331254in}{2.576251in}}%
\pgfpathlineto{\pgfqpoint{1.331254in}{2.670624in}}%
\pgfpathlineto{\pgfqpoint{1.476564in}{2.670624in}}%
\pgfpathlineto{\pgfqpoint{1.476564in}{2.576251in}}%
\pgfpathmoveto{\pgfqpoint{1.331254in}{2.670624in}}%
\pgfpathlineto{\pgfqpoint{1.331254in}{2.670624in}}%
\pgfpathlineto{\pgfqpoint{1.331254in}{2.764999in}}%
\pgfpathlineto{\pgfqpoint{1.476564in}{2.764999in}}%
\pgfpathlineto{\pgfqpoint{1.476564in}{2.670624in}}%
\pgfpathmoveto{\pgfqpoint{1.331254in}{2.764999in}}%
\pgfpathlineto{\pgfqpoint{1.331254in}{2.764999in}}%
\pgfpathlineto{\pgfqpoint{1.331254in}{2.859374in}}%
\pgfpathlineto{\pgfqpoint{1.476564in}{2.859374in}}%
\pgfpathlineto{\pgfqpoint{1.476564in}{2.764999in}}%
\pgfpathmoveto{\pgfqpoint{1.331254in}{2.859374in}}%
\pgfpathlineto{\pgfqpoint{1.331254in}{2.859374in}}%
\pgfpathlineto{\pgfqpoint{1.331254in}{2.953751in}}%
\pgfpathlineto{\pgfqpoint{1.476564in}{2.953751in}}%
\pgfpathlineto{\pgfqpoint{1.476564in}{2.859374in}}%
\pgfpathmoveto{\pgfqpoint{1.331254in}{2.953751in}}%
\pgfpathlineto{\pgfqpoint{1.331254in}{2.953751in}}%
\pgfpathlineto{\pgfqpoint{1.331254in}{3.048124in}}%
\pgfpathlineto{\pgfqpoint{1.476564in}{3.048124in}}%
\pgfpathlineto{\pgfqpoint{1.476564in}{2.953751in}}%
\pgfpathmoveto{\pgfqpoint{1.331254in}{3.048124in}}%
\pgfpathlineto{\pgfqpoint{1.331254in}{3.048124in}}%
\pgfpathlineto{\pgfqpoint{1.331254in}{3.142502in}}%
\pgfpathlineto{\pgfqpoint{1.476564in}{3.142502in}}%
\pgfpathlineto{\pgfqpoint{1.476564in}{3.048124in}}%
\pgfpathmoveto{\pgfqpoint{1.331254in}{3.142502in}}%
\pgfpathlineto{\pgfqpoint{1.331254in}{3.142502in}}%
\pgfpathlineto{\pgfqpoint{1.331254in}{3.236874in}}%
\pgfpathlineto{\pgfqpoint{1.476564in}{3.236874in}}%
\pgfpathlineto{\pgfqpoint{1.476564in}{3.142502in}}%
\pgfpathmoveto{\pgfqpoint{1.331254in}{3.236874in}}%
\pgfpathlineto{\pgfqpoint{1.331254in}{3.236874in}}%
\pgfpathlineto{\pgfqpoint{1.331254in}{3.331252in}}%
\pgfpathlineto{\pgfqpoint{1.476564in}{3.331252in}}%
\pgfpathlineto{\pgfqpoint{1.476564in}{3.236874in}}%
\pgfpathmoveto{\pgfqpoint{1.331254in}{3.331252in}}%
\pgfpathlineto{\pgfqpoint{1.331254in}{3.331252in}}%
\pgfpathlineto{\pgfqpoint{1.331254in}{3.425627in}}%
\pgfpathlineto{\pgfqpoint{1.476564in}{3.425627in}}%
\pgfpathlineto{\pgfqpoint{1.476564in}{3.331252in}}%
\pgfpathmoveto{\pgfqpoint{1.331254in}{3.425627in}}%
\pgfpathlineto{\pgfqpoint{1.331254in}{3.425627in}}%
\pgfpathlineto{\pgfqpoint{1.331254in}{3.519998in}}%
\pgfpathlineto{\pgfqpoint{1.476564in}{3.519998in}}%
\pgfpathlineto{\pgfqpoint{1.476564in}{3.425627in}}%
\pgfpathmoveto{\pgfqpoint{1.476564in}{1.160625in}}%
\pgfpathlineto{\pgfqpoint{1.476564in}{1.160625in}}%
\pgfpathlineto{\pgfqpoint{1.476564in}{1.254998in}}%
\pgfpathlineto{\pgfqpoint{1.621877in}{1.254998in}}%
\pgfpathlineto{\pgfqpoint{1.621877in}{1.160625in}}%
\pgfpathmoveto{\pgfqpoint{1.476564in}{1.254998in}}%
\pgfpathlineto{\pgfqpoint{1.476564in}{1.254998in}}%
\pgfpathlineto{\pgfqpoint{1.476564in}{1.349376in}}%
\pgfpathlineto{\pgfqpoint{1.621877in}{1.349376in}}%
\pgfpathlineto{\pgfqpoint{1.621877in}{1.254998in}}%
\pgfpathmoveto{\pgfqpoint{1.476564in}{1.349376in}}%
\pgfpathlineto{\pgfqpoint{1.476564in}{1.349376in}}%
\pgfpathlineto{\pgfqpoint{1.476564in}{1.443750in}}%
\pgfpathlineto{\pgfqpoint{1.621877in}{1.443750in}}%
\pgfpathlineto{\pgfqpoint{1.621877in}{1.349376in}}%
\pgfpathmoveto{\pgfqpoint{1.476564in}{1.443750in}}%
\pgfpathlineto{\pgfqpoint{1.476564in}{1.443750in}}%
\pgfpathlineto{\pgfqpoint{1.476564in}{1.538127in}}%
\pgfpathlineto{\pgfqpoint{1.621877in}{1.538127in}}%
\pgfpathlineto{\pgfqpoint{1.621877in}{1.443750in}}%
\pgfpathmoveto{\pgfqpoint{1.476564in}{1.538127in}}%
\pgfpathlineto{\pgfqpoint{1.476564in}{1.538127in}}%
\pgfpathlineto{\pgfqpoint{1.476564in}{1.632501in}}%
\pgfpathlineto{\pgfqpoint{1.621877in}{1.632501in}}%
\pgfpathlineto{\pgfqpoint{1.621877in}{1.538127in}}%
\pgfpathmoveto{\pgfqpoint{1.476564in}{1.632501in}}%
\pgfpathlineto{\pgfqpoint{1.476564in}{1.632501in}}%
\pgfpathlineto{\pgfqpoint{1.476564in}{1.726876in}}%
\pgfpathlineto{\pgfqpoint{1.621877in}{1.726876in}}%
\pgfpathlineto{\pgfqpoint{1.621877in}{1.632501in}}%
\pgfpathmoveto{\pgfqpoint{1.476564in}{1.726876in}}%
\pgfpathlineto{\pgfqpoint{1.476564in}{1.726876in}}%
\pgfpathlineto{\pgfqpoint{1.476564in}{1.821249in}}%
\pgfpathlineto{\pgfqpoint{1.621877in}{1.821249in}}%
\pgfpathlineto{\pgfqpoint{1.621877in}{1.726876in}}%
\pgfpathmoveto{\pgfqpoint{1.476564in}{1.821249in}}%
\pgfpathlineto{\pgfqpoint{1.476564in}{1.821249in}}%
\pgfpathlineto{\pgfqpoint{1.476564in}{1.915627in}}%
\pgfpathlineto{\pgfqpoint{1.621877in}{1.915627in}}%
\pgfpathlineto{\pgfqpoint{1.621877in}{1.821249in}}%
\pgfpathmoveto{\pgfqpoint{1.476564in}{1.915627in}}%
\pgfpathlineto{\pgfqpoint{1.476564in}{1.915627in}}%
\pgfpathlineto{\pgfqpoint{1.476564in}{2.010000in}}%
\pgfpathlineto{\pgfqpoint{1.621877in}{2.010000in}}%
\pgfpathlineto{\pgfqpoint{1.621877in}{1.915627in}}%
\pgfpathmoveto{\pgfqpoint{1.476564in}{2.010000in}}%
\pgfpathlineto{\pgfqpoint{1.476564in}{2.010000in}}%
\pgfpathlineto{\pgfqpoint{1.476564in}{2.104377in}}%
\pgfpathlineto{\pgfqpoint{1.621877in}{2.104377in}}%
\pgfpathlineto{\pgfqpoint{1.621877in}{2.010000in}}%
\pgfpathmoveto{\pgfqpoint{1.476564in}{2.104377in}}%
\pgfpathlineto{\pgfqpoint{1.476564in}{2.104377in}}%
\pgfpathlineto{\pgfqpoint{1.476564in}{2.198748in}}%
\pgfpathlineto{\pgfqpoint{1.621877in}{2.198748in}}%
\pgfpathlineto{\pgfqpoint{1.621877in}{2.104377in}}%
\pgfpathmoveto{\pgfqpoint{1.476564in}{2.198748in}}%
\pgfpathlineto{\pgfqpoint{1.476564in}{2.198748in}}%
\pgfpathlineto{\pgfqpoint{1.476564in}{2.293122in}}%
\pgfpathlineto{\pgfqpoint{1.621877in}{2.293122in}}%
\pgfpathlineto{\pgfqpoint{1.621877in}{2.198748in}}%
\pgfpathmoveto{\pgfqpoint{1.476564in}{2.293122in}}%
\pgfpathlineto{\pgfqpoint{1.476564in}{2.293122in}}%
\pgfpathlineto{\pgfqpoint{1.476564in}{2.387499in}}%
\pgfpathlineto{\pgfqpoint{1.621877in}{2.387499in}}%
\pgfpathlineto{\pgfqpoint{1.621877in}{2.293122in}}%
\pgfpathmoveto{\pgfqpoint{1.476564in}{2.387499in}}%
\pgfpathlineto{\pgfqpoint{1.476564in}{2.387499in}}%
\pgfpathlineto{\pgfqpoint{1.476564in}{2.481878in}}%
\pgfpathlineto{\pgfqpoint{1.621877in}{2.481878in}}%
\pgfpathlineto{\pgfqpoint{1.621877in}{2.387499in}}%
\pgfpathmoveto{\pgfqpoint{1.476564in}{2.481878in}}%
\pgfpathlineto{\pgfqpoint{1.476564in}{2.481878in}}%
\pgfpathlineto{\pgfqpoint{1.476564in}{2.576251in}}%
\pgfpathlineto{\pgfqpoint{1.621877in}{2.576251in}}%
\pgfpathlineto{\pgfqpoint{1.621877in}{2.481878in}}%
\pgfpathmoveto{\pgfqpoint{1.476564in}{2.576251in}}%
\pgfpathlineto{\pgfqpoint{1.476564in}{2.576251in}}%
\pgfpathlineto{\pgfqpoint{1.476564in}{2.670624in}}%
\pgfpathlineto{\pgfqpoint{1.621877in}{2.670624in}}%
\pgfpathlineto{\pgfqpoint{1.621877in}{2.576251in}}%
\pgfpathmoveto{\pgfqpoint{1.476564in}{2.670624in}}%
\pgfpathlineto{\pgfqpoint{1.476564in}{2.670624in}}%
\pgfpathlineto{\pgfqpoint{1.476564in}{2.764999in}}%
\pgfpathlineto{\pgfqpoint{1.621877in}{2.764999in}}%
\pgfpathlineto{\pgfqpoint{1.621877in}{2.670624in}}%
\pgfpathmoveto{\pgfqpoint{1.476564in}{2.764999in}}%
\pgfpathlineto{\pgfqpoint{1.476564in}{2.764999in}}%
\pgfpathlineto{\pgfqpoint{1.476564in}{2.859374in}}%
\pgfpathlineto{\pgfqpoint{1.621877in}{2.859374in}}%
\pgfpathlineto{\pgfqpoint{1.621877in}{2.764999in}}%
\pgfpathmoveto{\pgfqpoint{1.476564in}{2.859374in}}%
\pgfpathlineto{\pgfqpoint{1.476564in}{2.859374in}}%
\pgfpathlineto{\pgfqpoint{1.476564in}{2.953751in}}%
\pgfpathlineto{\pgfqpoint{1.621877in}{2.953751in}}%
\pgfpathlineto{\pgfqpoint{1.621877in}{2.859374in}}%
\pgfpathmoveto{\pgfqpoint{1.476564in}{2.953751in}}%
\pgfpathlineto{\pgfqpoint{1.476564in}{2.953751in}}%
\pgfpathlineto{\pgfqpoint{1.476564in}{3.048124in}}%
\pgfpathlineto{\pgfqpoint{1.621877in}{3.048124in}}%
\pgfpathlineto{\pgfqpoint{1.621877in}{2.953751in}}%
\pgfpathmoveto{\pgfqpoint{1.476564in}{3.048124in}}%
\pgfpathlineto{\pgfqpoint{1.476564in}{3.048124in}}%
\pgfpathlineto{\pgfqpoint{1.476564in}{3.142502in}}%
\pgfpathlineto{\pgfqpoint{1.621877in}{3.142502in}}%
\pgfpathlineto{\pgfqpoint{1.621877in}{3.048124in}}%
\pgfpathmoveto{\pgfqpoint{1.476564in}{3.142502in}}%
\pgfpathlineto{\pgfqpoint{1.476564in}{3.142502in}}%
\pgfpathlineto{\pgfqpoint{1.476564in}{3.236874in}}%
\pgfpathlineto{\pgfqpoint{1.621877in}{3.236874in}}%
\pgfpathlineto{\pgfqpoint{1.621877in}{3.142502in}}%
\pgfpathmoveto{\pgfqpoint{1.476564in}{3.236874in}}%
\pgfpathlineto{\pgfqpoint{1.476564in}{3.236874in}}%
\pgfpathlineto{\pgfqpoint{1.476564in}{3.331252in}}%
\pgfpathlineto{\pgfqpoint{1.621877in}{3.331252in}}%
\pgfpathlineto{\pgfqpoint{1.621877in}{3.236874in}}%
\pgfpathmoveto{\pgfqpoint{1.476564in}{3.331252in}}%
\pgfpathlineto{\pgfqpoint{1.476564in}{3.331252in}}%
\pgfpathlineto{\pgfqpoint{1.476564in}{3.425627in}}%
\pgfpathlineto{\pgfqpoint{1.621877in}{3.425627in}}%
\pgfpathlineto{\pgfqpoint{1.621877in}{3.331252in}}%
\pgfpathmoveto{\pgfqpoint{1.476564in}{3.425627in}}%
\pgfpathlineto{\pgfqpoint{1.476564in}{3.425627in}}%
\pgfpathlineto{\pgfqpoint{1.476564in}{3.519998in}}%
\pgfpathlineto{\pgfqpoint{1.621877in}{3.519998in}}%
\pgfpathlineto{\pgfqpoint{1.621877in}{3.425627in}}%
\pgfpathmoveto{\pgfqpoint{1.621877in}{1.160625in}}%
\pgfpathlineto{\pgfqpoint{1.621877in}{1.160625in}}%
\pgfpathlineto{\pgfqpoint{1.621877in}{1.254998in}}%
\pgfpathlineto{\pgfqpoint{1.767186in}{1.254998in}}%
\pgfpathlineto{\pgfqpoint{1.767186in}{1.160625in}}%
\pgfpathmoveto{\pgfqpoint{1.621877in}{1.254998in}}%
\pgfpathlineto{\pgfqpoint{1.621877in}{1.254998in}}%
\pgfpathlineto{\pgfqpoint{1.621877in}{1.349376in}}%
\pgfpathlineto{\pgfqpoint{1.767186in}{1.349376in}}%
\pgfpathlineto{\pgfqpoint{1.767186in}{1.254998in}}%
\pgfpathmoveto{\pgfqpoint{1.621877in}{1.349376in}}%
\pgfpathlineto{\pgfqpoint{1.621877in}{1.349376in}}%
\pgfpathlineto{\pgfqpoint{1.621877in}{1.443750in}}%
\pgfpathlineto{\pgfqpoint{1.767186in}{1.443750in}}%
\pgfpathlineto{\pgfqpoint{1.767186in}{1.349376in}}%
\pgfpathmoveto{\pgfqpoint{1.621877in}{1.443750in}}%
\pgfpathlineto{\pgfqpoint{1.621877in}{1.443750in}}%
\pgfpathlineto{\pgfqpoint{1.621877in}{1.538127in}}%
\pgfpathlineto{\pgfqpoint{1.767186in}{1.538127in}}%
\pgfpathlineto{\pgfqpoint{1.767186in}{1.443750in}}%
\pgfpathmoveto{\pgfqpoint{1.621877in}{1.538127in}}%
\pgfpathlineto{\pgfqpoint{1.621877in}{1.538127in}}%
\pgfpathlineto{\pgfqpoint{1.621877in}{1.632501in}}%
\pgfpathlineto{\pgfqpoint{1.767186in}{1.632501in}}%
\pgfpathlineto{\pgfqpoint{1.767186in}{1.538127in}}%
\pgfpathmoveto{\pgfqpoint{1.621877in}{1.632501in}}%
\pgfpathlineto{\pgfqpoint{1.621877in}{1.632501in}}%
\pgfpathlineto{\pgfqpoint{1.621877in}{1.726876in}}%
\pgfpathlineto{\pgfqpoint{1.767186in}{1.726876in}}%
\pgfpathlineto{\pgfqpoint{1.767186in}{1.632501in}}%
\pgfpathmoveto{\pgfqpoint{1.621877in}{1.726876in}}%
\pgfpathlineto{\pgfqpoint{1.621877in}{1.726876in}}%
\pgfpathlineto{\pgfqpoint{1.621877in}{1.821249in}}%
\pgfpathlineto{\pgfqpoint{1.767186in}{1.821249in}}%
\pgfpathlineto{\pgfqpoint{1.767186in}{1.726876in}}%
\pgfpathmoveto{\pgfqpoint{1.621877in}{1.821249in}}%
\pgfpathlineto{\pgfqpoint{1.621877in}{1.821249in}}%
\pgfpathlineto{\pgfqpoint{1.621877in}{1.915627in}}%
\pgfpathlineto{\pgfqpoint{1.767186in}{1.915627in}}%
\pgfpathlineto{\pgfqpoint{1.767186in}{1.821249in}}%
\pgfpathmoveto{\pgfqpoint{1.621877in}{1.915627in}}%
\pgfpathlineto{\pgfqpoint{1.621877in}{1.915627in}}%
\pgfpathlineto{\pgfqpoint{1.621877in}{2.010000in}}%
\pgfpathlineto{\pgfqpoint{1.767186in}{2.010000in}}%
\pgfpathlineto{\pgfqpoint{1.767186in}{1.915627in}}%
\pgfpathmoveto{\pgfqpoint{1.621877in}{2.010000in}}%
\pgfpathlineto{\pgfqpoint{1.621877in}{2.010000in}}%
\pgfpathlineto{\pgfqpoint{1.621877in}{2.104377in}}%
\pgfpathlineto{\pgfqpoint{1.767186in}{2.104377in}}%
\pgfpathlineto{\pgfqpoint{1.767186in}{2.010000in}}%
\pgfpathmoveto{\pgfqpoint{1.621877in}{2.104377in}}%
\pgfpathlineto{\pgfqpoint{1.621877in}{2.104377in}}%
\pgfpathlineto{\pgfqpoint{1.621877in}{2.198748in}}%
\pgfpathlineto{\pgfqpoint{1.767186in}{2.198748in}}%
\pgfpathlineto{\pgfqpoint{1.767186in}{2.104377in}}%
\pgfpathmoveto{\pgfqpoint{1.621877in}{2.198748in}}%
\pgfpathlineto{\pgfqpoint{1.621877in}{2.198748in}}%
\pgfpathlineto{\pgfqpoint{1.621877in}{2.293122in}}%
\pgfpathlineto{\pgfqpoint{1.767186in}{2.293122in}}%
\pgfpathlineto{\pgfqpoint{1.767186in}{2.198748in}}%
\pgfpathmoveto{\pgfqpoint{1.621877in}{2.293122in}}%
\pgfpathlineto{\pgfqpoint{1.621877in}{2.293122in}}%
\pgfpathlineto{\pgfqpoint{1.621877in}{2.387499in}}%
\pgfpathlineto{\pgfqpoint{1.767186in}{2.387499in}}%
\pgfpathlineto{\pgfqpoint{1.767186in}{2.293122in}}%
\pgfpathmoveto{\pgfqpoint{1.621877in}{2.387499in}}%
\pgfpathlineto{\pgfqpoint{1.621877in}{2.387499in}}%
\pgfpathlineto{\pgfqpoint{1.621877in}{2.481878in}}%
\pgfpathlineto{\pgfqpoint{1.767186in}{2.481878in}}%
\pgfpathlineto{\pgfqpoint{1.767186in}{2.387499in}}%
\pgfpathmoveto{\pgfqpoint{1.621877in}{2.481878in}}%
\pgfpathlineto{\pgfqpoint{1.621877in}{2.481878in}}%
\pgfpathlineto{\pgfqpoint{1.621877in}{2.576251in}}%
\pgfpathlineto{\pgfqpoint{1.767186in}{2.576251in}}%
\pgfpathlineto{\pgfqpoint{1.767186in}{2.481878in}}%
\pgfpathmoveto{\pgfqpoint{1.621877in}{2.576251in}}%
\pgfpathlineto{\pgfqpoint{1.621877in}{2.576251in}}%
\pgfpathlineto{\pgfqpoint{1.621877in}{2.670624in}}%
\pgfpathlineto{\pgfqpoint{1.767186in}{2.670624in}}%
\pgfpathlineto{\pgfqpoint{1.767186in}{2.576251in}}%
\pgfpathmoveto{\pgfqpoint{1.621877in}{2.670624in}}%
\pgfpathlineto{\pgfqpoint{1.621877in}{2.670624in}}%
\pgfpathlineto{\pgfqpoint{1.621877in}{2.764999in}}%
\pgfpathlineto{\pgfqpoint{1.767186in}{2.764999in}}%
\pgfpathlineto{\pgfqpoint{1.767186in}{2.670624in}}%
\pgfpathmoveto{\pgfqpoint{1.621877in}{2.764999in}}%
\pgfpathlineto{\pgfqpoint{1.621877in}{2.764999in}}%
\pgfpathlineto{\pgfqpoint{1.621877in}{2.859374in}}%
\pgfpathlineto{\pgfqpoint{1.767186in}{2.859374in}}%
\pgfpathlineto{\pgfqpoint{1.767186in}{2.764999in}}%
\pgfpathmoveto{\pgfqpoint{1.621877in}{2.859374in}}%
\pgfpathlineto{\pgfqpoint{1.621877in}{2.859374in}}%
\pgfpathlineto{\pgfqpoint{1.621877in}{2.953751in}}%
\pgfpathlineto{\pgfqpoint{1.767186in}{2.953751in}}%
\pgfpathlineto{\pgfqpoint{1.767186in}{2.859374in}}%
\pgfpathmoveto{\pgfqpoint{1.621877in}{2.953751in}}%
\pgfpathlineto{\pgfqpoint{1.621877in}{2.953751in}}%
\pgfpathlineto{\pgfqpoint{1.621877in}{3.048124in}}%
\pgfpathlineto{\pgfqpoint{1.767186in}{3.048124in}}%
\pgfpathlineto{\pgfqpoint{1.767186in}{2.953751in}}%
\pgfpathmoveto{\pgfqpoint{1.621877in}{3.048124in}}%
\pgfpathlineto{\pgfqpoint{1.621877in}{3.048124in}}%
\pgfpathlineto{\pgfqpoint{1.621877in}{3.142502in}}%
\pgfpathlineto{\pgfqpoint{1.767186in}{3.142502in}}%
\pgfpathlineto{\pgfqpoint{1.767186in}{3.048124in}}%
\pgfpathmoveto{\pgfqpoint{1.621877in}{3.142502in}}%
\pgfpathlineto{\pgfqpoint{1.621877in}{3.142502in}}%
\pgfpathlineto{\pgfqpoint{1.621877in}{3.236874in}}%
\pgfpathlineto{\pgfqpoint{1.767186in}{3.236874in}}%
\pgfpathlineto{\pgfqpoint{1.767186in}{3.142502in}}%
\pgfpathmoveto{\pgfqpoint{1.621877in}{3.236874in}}%
\pgfpathlineto{\pgfqpoint{1.621877in}{3.236874in}}%
\pgfpathlineto{\pgfqpoint{1.621877in}{3.331252in}}%
\pgfpathlineto{\pgfqpoint{1.767186in}{3.331252in}}%
\pgfpathlineto{\pgfqpoint{1.767186in}{3.236874in}}%
\pgfpathmoveto{\pgfqpoint{1.621877in}{3.331252in}}%
\pgfpathlineto{\pgfqpoint{1.621877in}{3.331252in}}%
\pgfpathlineto{\pgfqpoint{1.621877in}{3.425627in}}%
\pgfpathlineto{\pgfqpoint{1.767186in}{3.425627in}}%
\pgfpathlineto{\pgfqpoint{1.767186in}{3.331252in}}%
\pgfpathmoveto{\pgfqpoint{1.621877in}{3.425627in}}%
\pgfpathlineto{\pgfqpoint{1.621877in}{3.425627in}}%
\pgfpathlineto{\pgfqpoint{1.621877in}{3.519998in}}%
\pgfpathlineto{\pgfqpoint{1.767186in}{3.519998in}}%
\pgfpathlineto{\pgfqpoint{1.767186in}{3.425627in}}%
\pgfpathmoveto{\pgfqpoint{1.767186in}{1.254998in}}%
\pgfpathlineto{\pgfqpoint{1.767186in}{1.254998in}}%
\pgfpathlineto{\pgfqpoint{1.767186in}{1.349376in}}%
\pgfpathlineto{\pgfqpoint{1.912496in}{1.349376in}}%
\pgfpathlineto{\pgfqpoint{1.912496in}{1.254998in}}%
\pgfpathmoveto{\pgfqpoint{1.767186in}{1.349376in}}%
\pgfpathlineto{\pgfqpoint{1.767186in}{1.349376in}}%
\pgfpathlineto{\pgfqpoint{1.767186in}{1.443750in}}%
\pgfpathlineto{\pgfqpoint{1.912496in}{1.443750in}}%
\pgfpathlineto{\pgfqpoint{1.912496in}{1.349376in}}%
\pgfpathmoveto{\pgfqpoint{1.767186in}{1.443750in}}%
\pgfpathlineto{\pgfqpoint{1.767186in}{1.443750in}}%
\pgfpathlineto{\pgfqpoint{1.767186in}{1.538127in}}%
\pgfpathlineto{\pgfqpoint{1.912496in}{1.538127in}}%
\pgfpathlineto{\pgfqpoint{1.912496in}{1.443750in}}%
\pgfpathmoveto{\pgfqpoint{1.767186in}{1.538127in}}%
\pgfpathlineto{\pgfqpoint{1.767186in}{1.538127in}}%
\pgfpathlineto{\pgfqpoint{1.767186in}{1.632501in}}%
\pgfpathlineto{\pgfqpoint{1.912496in}{1.632501in}}%
\pgfpathlineto{\pgfqpoint{1.912496in}{1.538127in}}%
\pgfpathmoveto{\pgfqpoint{1.767186in}{1.632501in}}%
\pgfpathlineto{\pgfqpoint{1.767186in}{1.632501in}}%
\pgfpathlineto{\pgfqpoint{1.767186in}{1.726876in}}%
\pgfpathlineto{\pgfqpoint{1.912496in}{1.726876in}}%
\pgfpathlineto{\pgfqpoint{1.912496in}{1.632501in}}%
\pgfpathmoveto{\pgfqpoint{1.767186in}{1.726876in}}%
\pgfpathlineto{\pgfqpoint{1.767186in}{1.726876in}}%
\pgfpathlineto{\pgfqpoint{1.767186in}{1.821249in}}%
\pgfpathlineto{\pgfqpoint{1.912496in}{1.821249in}}%
\pgfpathlineto{\pgfqpoint{1.912496in}{1.726876in}}%
\pgfpathmoveto{\pgfqpoint{1.767186in}{1.821249in}}%
\pgfpathlineto{\pgfqpoint{1.767186in}{1.821249in}}%
\pgfpathlineto{\pgfqpoint{1.767186in}{1.915627in}}%
\pgfpathlineto{\pgfqpoint{1.912496in}{1.915627in}}%
\pgfpathlineto{\pgfqpoint{1.912496in}{1.821249in}}%
\pgfpathmoveto{\pgfqpoint{1.767186in}{1.915627in}}%
\pgfpathlineto{\pgfqpoint{1.767186in}{1.915627in}}%
\pgfpathlineto{\pgfqpoint{1.767186in}{2.010000in}}%
\pgfpathlineto{\pgfqpoint{1.912496in}{2.010000in}}%
\pgfpathlineto{\pgfqpoint{1.912496in}{1.915627in}}%
\pgfpathmoveto{\pgfqpoint{1.767186in}{2.010000in}}%
\pgfpathlineto{\pgfqpoint{1.767186in}{2.010000in}}%
\pgfpathlineto{\pgfqpoint{1.767186in}{2.104377in}}%
\pgfpathlineto{\pgfqpoint{1.912496in}{2.104377in}}%
\pgfpathlineto{\pgfqpoint{1.912496in}{2.010000in}}%
\pgfpathmoveto{\pgfqpoint{1.767186in}{2.104377in}}%
\pgfpathlineto{\pgfqpoint{1.767186in}{2.104377in}}%
\pgfpathlineto{\pgfqpoint{1.767186in}{2.198748in}}%
\pgfpathlineto{\pgfqpoint{1.912496in}{2.198748in}}%
\pgfpathlineto{\pgfqpoint{1.912496in}{2.104377in}}%
\pgfpathmoveto{\pgfqpoint{1.767186in}{2.198748in}}%
\pgfpathlineto{\pgfqpoint{1.767186in}{2.198748in}}%
\pgfpathlineto{\pgfqpoint{1.767186in}{2.293122in}}%
\pgfpathlineto{\pgfqpoint{1.912496in}{2.293122in}}%
\pgfpathlineto{\pgfqpoint{1.912496in}{2.198748in}}%
\pgfpathmoveto{\pgfqpoint{1.767186in}{2.293122in}}%
\pgfpathlineto{\pgfqpoint{1.767186in}{2.293122in}}%
\pgfpathlineto{\pgfqpoint{1.767186in}{2.387499in}}%
\pgfpathlineto{\pgfqpoint{1.912496in}{2.387499in}}%
\pgfpathlineto{\pgfqpoint{1.912496in}{2.293122in}}%
\pgfpathmoveto{\pgfqpoint{1.767186in}{2.387499in}}%
\pgfpathlineto{\pgfqpoint{1.767186in}{2.387499in}}%
\pgfpathlineto{\pgfqpoint{1.767186in}{2.481878in}}%
\pgfpathlineto{\pgfqpoint{1.912496in}{2.481878in}}%
\pgfpathlineto{\pgfqpoint{1.912496in}{2.387499in}}%
\pgfpathmoveto{\pgfqpoint{1.767186in}{2.481878in}}%
\pgfpathlineto{\pgfqpoint{1.767186in}{2.481878in}}%
\pgfpathlineto{\pgfqpoint{1.767186in}{2.576251in}}%
\pgfpathlineto{\pgfqpoint{1.912496in}{2.576251in}}%
\pgfpathlineto{\pgfqpoint{1.912496in}{2.481878in}}%
\pgfpathmoveto{\pgfqpoint{1.767186in}{2.576251in}}%
\pgfpathlineto{\pgfqpoint{1.767186in}{2.576251in}}%
\pgfpathlineto{\pgfqpoint{1.767186in}{2.670624in}}%
\pgfpathlineto{\pgfqpoint{1.912496in}{2.670624in}}%
\pgfpathlineto{\pgfqpoint{1.912496in}{2.576251in}}%
\pgfpathmoveto{\pgfqpoint{1.767186in}{2.670624in}}%
\pgfpathlineto{\pgfqpoint{1.767186in}{2.670624in}}%
\pgfpathlineto{\pgfqpoint{1.767186in}{2.764999in}}%
\pgfpathlineto{\pgfqpoint{1.912496in}{2.764999in}}%
\pgfpathlineto{\pgfqpoint{1.912496in}{2.670624in}}%
\pgfpathmoveto{\pgfqpoint{1.767186in}{2.764999in}}%
\pgfpathlineto{\pgfqpoint{1.767186in}{2.764999in}}%
\pgfpathlineto{\pgfqpoint{1.767186in}{2.859374in}}%
\pgfpathlineto{\pgfqpoint{1.912496in}{2.859374in}}%
\pgfpathlineto{\pgfqpoint{1.912496in}{2.764999in}}%
\pgfpathmoveto{\pgfqpoint{1.767186in}{2.859374in}}%
\pgfpathlineto{\pgfqpoint{1.767186in}{2.859374in}}%
\pgfpathlineto{\pgfqpoint{1.767186in}{2.953751in}}%
\pgfpathlineto{\pgfqpoint{1.912496in}{2.953751in}}%
\pgfpathlineto{\pgfqpoint{1.912496in}{2.859374in}}%
\pgfpathmoveto{\pgfqpoint{1.767186in}{2.953751in}}%
\pgfpathlineto{\pgfqpoint{1.767186in}{2.953751in}}%
\pgfpathlineto{\pgfqpoint{1.767186in}{3.048124in}}%
\pgfpathlineto{\pgfqpoint{1.912496in}{3.048124in}}%
\pgfpathlineto{\pgfqpoint{1.912496in}{2.953751in}}%
\pgfpathmoveto{\pgfqpoint{1.767186in}{3.048124in}}%
\pgfpathlineto{\pgfqpoint{1.767186in}{3.048124in}}%
\pgfpathlineto{\pgfqpoint{1.767186in}{3.142502in}}%
\pgfpathlineto{\pgfqpoint{1.912496in}{3.142502in}}%
\pgfpathlineto{\pgfqpoint{1.912496in}{3.048124in}}%
\pgfpathmoveto{\pgfqpoint{1.767186in}{3.142502in}}%
\pgfpathlineto{\pgfqpoint{1.767186in}{3.142502in}}%
\pgfpathlineto{\pgfqpoint{1.767186in}{3.236874in}}%
\pgfpathlineto{\pgfqpoint{1.912496in}{3.236874in}}%
\pgfpathlineto{\pgfqpoint{1.912496in}{3.142502in}}%
\pgfpathmoveto{\pgfqpoint{1.767186in}{3.236874in}}%
\pgfpathlineto{\pgfqpoint{1.767186in}{3.236874in}}%
\pgfpathlineto{\pgfqpoint{1.767186in}{3.331252in}}%
\pgfpathlineto{\pgfqpoint{1.912496in}{3.331252in}}%
\pgfpathlineto{\pgfqpoint{1.912496in}{3.236874in}}%
\pgfpathmoveto{\pgfqpoint{1.767186in}{3.331252in}}%
\pgfpathlineto{\pgfqpoint{1.767186in}{3.331252in}}%
\pgfpathlineto{\pgfqpoint{1.767186in}{3.425627in}}%
\pgfpathlineto{\pgfqpoint{1.912496in}{3.425627in}}%
\pgfpathlineto{\pgfqpoint{1.912496in}{3.331252in}}%
\pgfpathmoveto{\pgfqpoint{1.767186in}{3.425627in}}%
\pgfpathlineto{\pgfqpoint{1.767186in}{3.425627in}}%
\pgfpathlineto{\pgfqpoint{1.767186in}{3.519998in}}%
\pgfpathlineto{\pgfqpoint{1.912496in}{3.519998in}}%
\pgfpathlineto{\pgfqpoint{1.912496in}{3.425627in}}%
\pgfpathmoveto{\pgfqpoint{1.912496in}{1.349376in}}%
\pgfpathlineto{\pgfqpoint{1.912496in}{1.349376in}}%
\pgfpathlineto{\pgfqpoint{1.912496in}{1.443750in}}%
\pgfpathlineto{\pgfqpoint{2.057813in}{1.443750in}}%
\pgfpathlineto{\pgfqpoint{2.057813in}{1.349376in}}%
\pgfpathmoveto{\pgfqpoint{1.912496in}{1.443750in}}%
\pgfpathlineto{\pgfqpoint{1.912496in}{1.443750in}}%
\pgfpathlineto{\pgfqpoint{1.912496in}{1.538127in}}%
\pgfpathlineto{\pgfqpoint{2.057813in}{1.538127in}}%
\pgfpathlineto{\pgfqpoint{2.057813in}{1.443750in}}%
\pgfpathmoveto{\pgfqpoint{1.912496in}{1.538127in}}%
\pgfpathlineto{\pgfqpoint{1.912496in}{1.538127in}}%
\pgfpathlineto{\pgfqpoint{1.912496in}{1.632501in}}%
\pgfpathlineto{\pgfqpoint{2.057813in}{1.632501in}}%
\pgfpathlineto{\pgfqpoint{2.057813in}{1.538127in}}%
\pgfpathmoveto{\pgfqpoint{1.912496in}{1.632501in}}%
\pgfpathlineto{\pgfqpoint{1.912496in}{1.632501in}}%
\pgfpathlineto{\pgfqpoint{1.912496in}{1.726876in}}%
\pgfpathlineto{\pgfqpoint{2.057813in}{1.726876in}}%
\pgfpathlineto{\pgfqpoint{2.057813in}{1.632501in}}%
\pgfpathmoveto{\pgfqpoint{1.912496in}{1.726876in}}%
\pgfpathlineto{\pgfqpoint{1.912496in}{1.726876in}}%
\pgfpathlineto{\pgfqpoint{1.912496in}{1.821249in}}%
\pgfpathlineto{\pgfqpoint{2.057813in}{1.821249in}}%
\pgfpathlineto{\pgfqpoint{2.057813in}{1.726876in}}%
\pgfpathmoveto{\pgfqpoint{1.912496in}{1.821249in}}%
\pgfpathlineto{\pgfqpoint{1.912496in}{1.821249in}}%
\pgfpathlineto{\pgfqpoint{1.912496in}{1.915627in}}%
\pgfpathlineto{\pgfqpoint{2.057813in}{1.915627in}}%
\pgfpathlineto{\pgfqpoint{2.057813in}{1.821249in}}%
\pgfpathmoveto{\pgfqpoint{1.912496in}{1.915627in}}%
\pgfpathlineto{\pgfqpoint{1.912496in}{1.915627in}}%
\pgfpathlineto{\pgfqpoint{1.912496in}{2.010000in}}%
\pgfpathlineto{\pgfqpoint{2.057813in}{2.010000in}}%
\pgfpathlineto{\pgfqpoint{2.057813in}{1.915627in}}%
\pgfpathmoveto{\pgfqpoint{1.912496in}{2.010000in}}%
\pgfpathlineto{\pgfqpoint{1.912496in}{2.010000in}}%
\pgfpathlineto{\pgfqpoint{1.912496in}{2.104377in}}%
\pgfpathlineto{\pgfqpoint{2.057813in}{2.104377in}}%
\pgfpathlineto{\pgfqpoint{2.057813in}{2.010000in}}%
\pgfpathmoveto{\pgfqpoint{1.912496in}{2.104377in}}%
\pgfpathlineto{\pgfqpoint{1.912496in}{2.104377in}}%
\pgfpathlineto{\pgfqpoint{1.912496in}{2.198748in}}%
\pgfpathlineto{\pgfqpoint{2.057813in}{2.198748in}}%
\pgfpathlineto{\pgfqpoint{2.057813in}{2.104377in}}%
\pgfpathmoveto{\pgfqpoint{1.912496in}{2.198748in}}%
\pgfpathlineto{\pgfqpoint{1.912496in}{2.198748in}}%
\pgfpathlineto{\pgfqpoint{1.912496in}{2.293122in}}%
\pgfpathlineto{\pgfqpoint{2.057813in}{2.293122in}}%
\pgfpathlineto{\pgfqpoint{2.057813in}{2.198748in}}%
\pgfpathmoveto{\pgfqpoint{1.912496in}{2.293122in}}%
\pgfpathlineto{\pgfqpoint{1.912496in}{2.293122in}}%
\pgfpathlineto{\pgfqpoint{1.912496in}{2.387499in}}%
\pgfpathlineto{\pgfqpoint{2.057813in}{2.387499in}}%
\pgfpathlineto{\pgfqpoint{2.057813in}{2.293122in}}%
\pgfpathmoveto{\pgfqpoint{1.912496in}{2.387499in}}%
\pgfpathlineto{\pgfqpoint{1.912496in}{2.387499in}}%
\pgfpathlineto{\pgfqpoint{1.912496in}{2.481878in}}%
\pgfpathlineto{\pgfqpoint{2.057813in}{2.481878in}}%
\pgfpathlineto{\pgfqpoint{2.057813in}{2.387499in}}%
\pgfpathmoveto{\pgfqpoint{1.912496in}{2.481878in}}%
\pgfpathlineto{\pgfqpoint{1.912496in}{2.481878in}}%
\pgfpathlineto{\pgfqpoint{1.912496in}{2.576251in}}%
\pgfpathlineto{\pgfqpoint{2.057813in}{2.576251in}}%
\pgfpathlineto{\pgfqpoint{2.057813in}{2.481878in}}%
\pgfpathmoveto{\pgfqpoint{1.912496in}{2.576251in}}%
\pgfpathlineto{\pgfqpoint{1.912496in}{2.576251in}}%
\pgfpathlineto{\pgfqpoint{1.912496in}{2.670624in}}%
\pgfpathlineto{\pgfqpoint{2.057813in}{2.670624in}}%
\pgfpathlineto{\pgfqpoint{2.057813in}{2.576251in}}%
\pgfpathmoveto{\pgfqpoint{1.912496in}{2.670624in}}%
\pgfpathlineto{\pgfqpoint{1.912496in}{2.670624in}}%
\pgfpathlineto{\pgfqpoint{1.912496in}{2.764999in}}%
\pgfpathlineto{\pgfqpoint{2.057813in}{2.764999in}}%
\pgfpathlineto{\pgfqpoint{2.057813in}{2.670624in}}%
\pgfpathmoveto{\pgfqpoint{1.912496in}{2.764999in}}%
\pgfpathlineto{\pgfqpoint{1.912496in}{2.764999in}}%
\pgfpathlineto{\pgfqpoint{1.912496in}{2.859374in}}%
\pgfpathlineto{\pgfqpoint{2.057813in}{2.859374in}}%
\pgfpathlineto{\pgfqpoint{2.057813in}{2.764999in}}%
\pgfpathmoveto{\pgfqpoint{1.912496in}{2.859374in}}%
\pgfpathlineto{\pgfqpoint{1.912496in}{2.859374in}}%
\pgfpathlineto{\pgfqpoint{1.912496in}{2.953751in}}%
\pgfpathlineto{\pgfqpoint{2.057813in}{2.953751in}}%
\pgfpathlineto{\pgfqpoint{2.057813in}{2.859374in}}%
\pgfpathmoveto{\pgfqpoint{1.912496in}{2.953751in}}%
\pgfpathlineto{\pgfqpoint{1.912496in}{2.953751in}}%
\pgfpathlineto{\pgfqpoint{1.912496in}{3.048124in}}%
\pgfpathlineto{\pgfqpoint{2.057813in}{3.048124in}}%
\pgfpathlineto{\pgfqpoint{2.057813in}{2.953751in}}%
\pgfpathmoveto{\pgfqpoint{1.912496in}{3.048124in}}%
\pgfpathlineto{\pgfqpoint{1.912496in}{3.048124in}}%
\pgfpathlineto{\pgfqpoint{1.912496in}{3.142502in}}%
\pgfpathlineto{\pgfqpoint{2.057813in}{3.142502in}}%
\pgfpathlineto{\pgfqpoint{2.057813in}{3.048124in}}%
\pgfpathmoveto{\pgfqpoint{1.912496in}{3.142502in}}%
\pgfpathlineto{\pgfqpoint{1.912496in}{3.142502in}}%
\pgfpathlineto{\pgfqpoint{1.912496in}{3.236874in}}%
\pgfpathlineto{\pgfqpoint{2.057813in}{3.236874in}}%
\pgfpathlineto{\pgfqpoint{2.057813in}{3.142502in}}%
\pgfpathmoveto{\pgfqpoint{1.912496in}{3.236874in}}%
\pgfpathlineto{\pgfqpoint{1.912496in}{3.236874in}}%
\pgfpathlineto{\pgfqpoint{1.912496in}{3.331252in}}%
\pgfpathlineto{\pgfqpoint{2.057813in}{3.331252in}}%
\pgfpathlineto{\pgfqpoint{2.057813in}{3.236874in}}%
\pgfpathmoveto{\pgfqpoint{1.912496in}{3.331252in}}%
\pgfpathlineto{\pgfqpoint{1.912496in}{3.331252in}}%
\pgfpathlineto{\pgfqpoint{1.912496in}{3.425627in}}%
\pgfpathlineto{\pgfqpoint{2.057813in}{3.425627in}}%
\pgfpathlineto{\pgfqpoint{2.057813in}{3.331252in}}%
\pgfpathmoveto{\pgfqpoint{1.912496in}{3.425627in}}%
\pgfpathlineto{\pgfqpoint{1.912496in}{3.425627in}}%
\pgfpathlineto{\pgfqpoint{1.912496in}{3.519998in}}%
\pgfpathlineto{\pgfqpoint{2.057813in}{3.519998in}}%
\pgfpathlineto{\pgfqpoint{2.057813in}{3.425627in}}%
\pgfpathmoveto{\pgfqpoint{2.057813in}{1.538127in}}%
\pgfpathlineto{\pgfqpoint{2.057813in}{1.538127in}}%
\pgfpathlineto{\pgfqpoint{2.057813in}{1.632501in}}%
\pgfpathlineto{\pgfqpoint{2.203129in}{1.632501in}}%
\pgfpathlineto{\pgfqpoint{2.203129in}{1.538127in}}%
\pgfpathmoveto{\pgfqpoint{2.057813in}{1.632501in}}%
\pgfpathlineto{\pgfqpoint{2.057813in}{1.632501in}}%
\pgfpathlineto{\pgfqpoint{2.057813in}{1.726876in}}%
\pgfpathlineto{\pgfqpoint{2.203129in}{1.726876in}}%
\pgfpathlineto{\pgfqpoint{2.203129in}{1.632501in}}%
\pgfpathmoveto{\pgfqpoint{2.057813in}{1.726876in}}%
\pgfpathlineto{\pgfqpoint{2.057813in}{1.726876in}}%
\pgfpathlineto{\pgfqpoint{2.057813in}{1.821249in}}%
\pgfpathlineto{\pgfqpoint{2.203129in}{1.821249in}}%
\pgfpathlineto{\pgfqpoint{2.203129in}{1.726876in}}%
\pgfpathmoveto{\pgfqpoint{2.057813in}{1.821249in}}%
\pgfpathlineto{\pgfqpoint{2.057813in}{1.821249in}}%
\pgfpathlineto{\pgfqpoint{2.057813in}{1.915627in}}%
\pgfpathlineto{\pgfqpoint{2.203129in}{1.915627in}}%
\pgfpathlineto{\pgfqpoint{2.203129in}{1.821249in}}%
\pgfpathmoveto{\pgfqpoint{2.057813in}{1.915627in}}%
\pgfpathlineto{\pgfqpoint{2.057813in}{1.915627in}}%
\pgfpathlineto{\pgfqpoint{2.057813in}{2.010000in}}%
\pgfpathlineto{\pgfqpoint{2.203129in}{2.010000in}}%
\pgfpathlineto{\pgfqpoint{2.203129in}{1.915627in}}%
\pgfpathmoveto{\pgfqpoint{2.057813in}{2.010000in}}%
\pgfpathlineto{\pgfqpoint{2.057813in}{2.010000in}}%
\pgfpathlineto{\pgfqpoint{2.057813in}{2.104377in}}%
\pgfpathlineto{\pgfqpoint{2.203129in}{2.104377in}}%
\pgfpathlineto{\pgfqpoint{2.203129in}{2.010000in}}%
\pgfpathmoveto{\pgfqpoint{2.057813in}{2.104377in}}%
\pgfpathlineto{\pgfqpoint{2.057813in}{2.104377in}}%
\pgfpathlineto{\pgfqpoint{2.057813in}{2.198748in}}%
\pgfpathlineto{\pgfqpoint{2.203129in}{2.198748in}}%
\pgfpathlineto{\pgfqpoint{2.203129in}{2.104377in}}%
\pgfpathmoveto{\pgfqpoint{2.057813in}{2.198748in}}%
\pgfpathlineto{\pgfqpoint{2.057813in}{2.198748in}}%
\pgfpathlineto{\pgfqpoint{2.057813in}{2.293122in}}%
\pgfpathlineto{\pgfqpoint{2.203129in}{2.293122in}}%
\pgfpathlineto{\pgfqpoint{2.203129in}{2.198748in}}%
\pgfpathmoveto{\pgfqpoint{2.057813in}{2.293122in}}%
\pgfpathlineto{\pgfqpoint{2.057813in}{2.293122in}}%
\pgfpathlineto{\pgfqpoint{2.057813in}{2.387499in}}%
\pgfpathlineto{\pgfqpoint{2.203129in}{2.387499in}}%
\pgfpathlineto{\pgfqpoint{2.203129in}{2.293122in}}%
\pgfpathmoveto{\pgfqpoint{2.057813in}{2.387499in}}%
\pgfpathlineto{\pgfqpoint{2.057813in}{2.387499in}}%
\pgfpathlineto{\pgfqpoint{2.057813in}{2.481878in}}%
\pgfpathlineto{\pgfqpoint{2.203129in}{2.481878in}}%
\pgfpathlineto{\pgfqpoint{2.203129in}{2.387499in}}%
\pgfpathmoveto{\pgfqpoint{2.057813in}{2.481878in}}%
\pgfpathlineto{\pgfqpoint{2.057813in}{2.481878in}}%
\pgfpathlineto{\pgfqpoint{2.057813in}{2.576251in}}%
\pgfpathlineto{\pgfqpoint{2.203129in}{2.576251in}}%
\pgfpathlineto{\pgfqpoint{2.203129in}{2.481878in}}%
\pgfpathmoveto{\pgfqpoint{2.057813in}{2.576251in}}%
\pgfpathlineto{\pgfqpoint{2.057813in}{2.576251in}}%
\pgfpathlineto{\pgfqpoint{2.057813in}{2.670624in}}%
\pgfpathlineto{\pgfqpoint{2.203129in}{2.670624in}}%
\pgfpathlineto{\pgfqpoint{2.203129in}{2.576251in}}%
\pgfpathmoveto{\pgfqpoint{2.057813in}{2.670624in}}%
\pgfpathlineto{\pgfqpoint{2.057813in}{2.670624in}}%
\pgfpathlineto{\pgfqpoint{2.057813in}{2.764999in}}%
\pgfpathlineto{\pgfqpoint{2.203129in}{2.764999in}}%
\pgfpathlineto{\pgfqpoint{2.203129in}{2.670624in}}%
\pgfpathmoveto{\pgfqpoint{2.057813in}{2.764999in}}%
\pgfpathlineto{\pgfqpoint{2.057813in}{2.764999in}}%
\pgfpathlineto{\pgfqpoint{2.057813in}{2.859374in}}%
\pgfpathlineto{\pgfqpoint{2.203129in}{2.859374in}}%
\pgfpathlineto{\pgfqpoint{2.203129in}{2.764999in}}%
\pgfpathmoveto{\pgfqpoint{2.057813in}{2.859374in}}%
\pgfpathlineto{\pgfqpoint{2.057813in}{2.859374in}}%
\pgfpathlineto{\pgfqpoint{2.057813in}{2.953751in}}%
\pgfpathlineto{\pgfqpoint{2.203129in}{2.953751in}}%
\pgfpathlineto{\pgfqpoint{2.203129in}{2.859374in}}%
\pgfpathmoveto{\pgfqpoint{2.057813in}{2.953751in}}%
\pgfpathlineto{\pgfqpoint{2.057813in}{2.953751in}}%
\pgfpathlineto{\pgfqpoint{2.057813in}{3.048124in}}%
\pgfpathlineto{\pgfqpoint{2.203129in}{3.048124in}}%
\pgfpathlineto{\pgfqpoint{2.203129in}{2.953751in}}%
\pgfpathmoveto{\pgfqpoint{2.057813in}{3.048124in}}%
\pgfpathlineto{\pgfqpoint{2.057813in}{3.048124in}}%
\pgfpathlineto{\pgfqpoint{2.057813in}{3.142502in}}%
\pgfpathlineto{\pgfqpoint{2.203129in}{3.142502in}}%
\pgfpathlineto{\pgfqpoint{2.203129in}{3.048124in}}%
\pgfpathmoveto{\pgfqpoint{2.057813in}{3.142502in}}%
\pgfpathlineto{\pgfqpoint{2.057813in}{3.142502in}}%
\pgfpathlineto{\pgfqpoint{2.057813in}{3.236874in}}%
\pgfpathlineto{\pgfqpoint{2.203129in}{3.236874in}}%
\pgfpathlineto{\pgfqpoint{2.203129in}{3.142502in}}%
\pgfpathmoveto{\pgfqpoint{2.057813in}{3.236874in}}%
\pgfpathlineto{\pgfqpoint{2.057813in}{3.236874in}}%
\pgfpathlineto{\pgfqpoint{2.057813in}{3.331252in}}%
\pgfpathlineto{\pgfqpoint{2.203129in}{3.331252in}}%
\pgfpathlineto{\pgfqpoint{2.203129in}{3.236874in}}%
\pgfpathmoveto{\pgfqpoint{2.057813in}{3.331252in}}%
\pgfpathlineto{\pgfqpoint{2.057813in}{3.331252in}}%
\pgfpathlineto{\pgfqpoint{2.057813in}{3.425627in}}%
\pgfpathlineto{\pgfqpoint{2.203129in}{3.425627in}}%
\pgfpathlineto{\pgfqpoint{2.203129in}{3.331252in}}%
\pgfpathmoveto{\pgfqpoint{2.057813in}{3.425627in}}%
\pgfpathlineto{\pgfqpoint{2.057813in}{3.425627in}}%
\pgfpathlineto{\pgfqpoint{2.057813in}{3.519998in}}%
\pgfpathlineto{\pgfqpoint{2.203129in}{3.519998in}}%
\pgfpathlineto{\pgfqpoint{2.203129in}{3.425627in}}%
\pgfpathmoveto{\pgfqpoint{2.203129in}{1.632501in}}%
\pgfpathlineto{\pgfqpoint{2.203129in}{1.632501in}}%
\pgfpathlineto{\pgfqpoint{2.203129in}{1.726876in}}%
\pgfpathlineto{\pgfqpoint{2.348441in}{1.726876in}}%
\pgfpathlineto{\pgfqpoint{2.348441in}{1.632501in}}%
\pgfpathmoveto{\pgfqpoint{2.203129in}{1.726876in}}%
\pgfpathlineto{\pgfqpoint{2.203129in}{1.726876in}}%
\pgfpathlineto{\pgfqpoint{2.203129in}{1.821249in}}%
\pgfpathlineto{\pgfqpoint{2.348441in}{1.821249in}}%
\pgfpathlineto{\pgfqpoint{2.348441in}{1.726876in}}%
\pgfpathmoveto{\pgfqpoint{2.203129in}{1.821249in}}%
\pgfpathlineto{\pgfqpoint{2.203129in}{1.821249in}}%
\pgfpathlineto{\pgfqpoint{2.203129in}{1.915627in}}%
\pgfpathlineto{\pgfqpoint{2.348441in}{1.915627in}}%
\pgfpathlineto{\pgfqpoint{2.348441in}{1.821249in}}%
\pgfpathmoveto{\pgfqpoint{2.203129in}{1.915627in}}%
\pgfpathlineto{\pgfqpoint{2.203129in}{1.915627in}}%
\pgfpathlineto{\pgfqpoint{2.203129in}{2.010000in}}%
\pgfpathlineto{\pgfqpoint{2.348441in}{2.010000in}}%
\pgfpathlineto{\pgfqpoint{2.348441in}{1.915627in}}%
\pgfpathmoveto{\pgfqpoint{2.203129in}{2.010000in}}%
\pgfpathlineto{\pgfqpoint{2.203129in}{2.010000in}}%
\pgfpathlineto{\pgfqpoint{2.203129in}{2.104377in}}%
\pgfpathlineto{\pgfqpoint{2.348441in}{2.104377in}}%
\pgfpathlineto{\pgfqpoint{2.348441in}{2.010000in}}%
\pgfpathmoveto{\pgfqpoint{2.203129in}{2.104377in}}%
\pgfpathlineto{\pgfqpoint{2.203129in}{2.104377in}}%
\pgfpathlineto{\pgfqpoint{2.203129in}{2.198748in}}%
\pgfpathlineto{\pgfqpoint{2.348441in}{2.198748in}}%
\pgfpathlineto{\pgfqpoint{2.348441in}{2.104377in}}%
\pgfpathmoveto{\pgfqpoint{2.203129in}{2.198748in}}%
\pgfpathlineto{\pgfqpoint{2.203129in}{2.198748in}}%
\pgfpathlineto{\pgfqpoint{2.203129in}{2.293122in}}%
\pgfpathlineto{\pgfqpoint{2.348441in}{2.293122in}}%
\pgfpathlineto{\pgfqpoint{2.348441in}{2.198748in}}%
\pgfpathmoveto{\pgfqpoint{2.203129in}{2.293122in}}%
\pgfpathlineto{\pgfqpoint{2.203129in}{2.293122in}}%
\pgfpathlineto{\pgfqpoint{2.203129in}{2.387499in}}%
\pgfpathlineto{\pgfqpoint{2.348441in}{2.387499in}}%
\pgfpathlineto{\pgfqpoint{2.348441in}{2.293122in}}%
\pgfpathmoveto{\pgfqpoint{2.203129in}{2.387499in}}%
\pgfpathlineto{\pgfqpoint{2.203129in}{2.387499in}}%
\pgfpathlineto{\pgfqpoint{2.203129in}{2.481878in}}%
\pgfpathlineto{\pgfqpoint{2.348441in}{2.481878in}}%
\pgfpathlineto{\pgfqpoint{2.348441in}{2.387499in}}%
\pgfpathmoveto{\pgfqpoint{2.203129in}{2.481878in}}%
\pgfpathlineto{\pgfqpoint{2.203129in}{2.481878in}}%
\pgfpathlineto{\pgfqpoint{2.203129in}{2.576251in}}%
\pgfpathlineto{\pgfqpoint{2.348441in}{2.576251in}}%
\pgfpathlineto{\pgfqpoint{2.348441in}{2.481878in}}%
\pgfpathmoveto{\pgfqpoint{2.203129in}{2.576251in}}%
\pgfpathlineto{\pgfqpoint{2.203129in}{2.576251in}}%
\pgfpathlineto{\pgfqpoint{2.203129in}{2.670624in}}%
\pgfpathlineto{\pgfqpoint{2.348441in}{2.670624in}}%
\pgfpathlineto{\pgfqpoint{2.348441in}{2.576251in}}%
\pgfpathmoveto{\pgfqpoint{2.203129in}{2.670624in}}%
\pgfpathlineto{\pgfqpoint{2.203129in}{2.670624in}}%
\pgfpathlineto{\pgfqpoint{2.203129in}{2.764999in}}%
\pgfpathlineto{\pgfqpoint{2.348441in}{2.764999in}}%
\pgfpathlineto{\pgfqpoint{2.348441in}{2.670624in}}%
\pgfpathmoveto{\pgfqpoint{2.203129in}{2.764999in}}%
\pgfpathlineto{\pgfqpoint{2.203129in}{2.764999in}}%
\pgfpathlineto{\pgfqpoint{2.203129in}{2.859374in}}%
\pgfpathlineto{\pgfqpoint{2.348441in}{2.859374in}}%
\pgfpathlineto{\pgfqpoint{2.348441in}{2.764999in}}%
\pgfpathmoveto{\pgfqpoint{2.203129in}{2.859374in}}%
\pgfpathlineto{\pgfqpoint{2.203129in}{2.859374in}}%
\pgfpathlineto{\pgfqpoint{2.203129in}{2.953751in}}%
\pgfpathlineto{\pgfqpoint{2.348441in}{2.953751in}}%
\pgfpathlineto{\pgfqpoint{2.348441in}{2.859374in}}%
\pgfpathmoveto{\pgfqpoint{2.203129in}{2.953751in}}%
\pgfpathlineto{\pgfqpoint{2.203129in}{2.953751in}}%
\pgfpathlineto{\pgfqpoint{2.203129in}{3.048124in}}%
\pgfpathlineto{\pgfqpoint{2.348441in}{3.048124in}}%
\pgfpathlineto{\pgfqpoint{2.348441in}{2.953751in}}%
\pgfpathmoveto{\pgfqpoint{2.203129in}{3.048124in}}%
\pgfpathlineto{\pgfqpoint{2.203129in}{3.048124in}}%
\pgfpathlineto{\pgfqpoint{2.203129in}{3.142502in}}%
\pgfpathlineto{\pgfqpoint{2.348441in}{3.142502in}}%
\pgfpathlineto{\pgfqpoint{2.348441in}{3.048124in}}%
\pgfpathmoveto{\pgfqpoint{2.203129in}{3.142502in}}%
\pgfpathlineto{\pgfqpoint{2.203129in}{3.142502in}}%
\pgfpathlineto{\pgfqpoint{2.203129in}{3.236874in}}%
\pgfpathlineto{\pgfqpoint{2.348441in}{3.236874in}}%
\pgfpathlineto{\pgfqpoint{2.348441in}{3.142502in}}%
\pgfpathmoveto{\pgfqpoint{2.203129in}{3.236874in}}%
\pgfpathlineto{\pgfqpoint{2.203129in}{3.236874in}}%
\pgfpathlineto{\pgfqpoint{2.203129in}{3.331252in}}%
\pgfpathlineto{\pgfqpoint{2.348441in}{3.331252in}}%
\pgfpathlineto{\pgfqpoint{2.348441in}{3.236874in}}%
\pgfpathmoveto{\pgfqpoint{2.203129in}{3.331252in}}%
\pgfpathlineto{\pgfqpoint{2.203129in}{3.331252in}}%
\pgfpathlineto{\pgfqpoint{2.203129in}{3.425627in}}%
\pgfpathlineto{\pgfqpoint{2.348441in}{3.425627in}}%
\pgfpathlineto{\pgfqpoint{2.348441in}{3.331252in}}%
\pgfpathmoveto{\pgfqpoint{2.203129in}{3.425627in}}%
\pgfpathlineto{\pgfqpoint{2.203129in}{3.425627in}}%
\pgfpathlineto{\pgfqpoint{2.203129in}{3.519998in}}%
\pgfpathlineto{\pgfqpoint{2.348441in}{3.519998in}}%
\pgfpathlineto{\pgfqpoint{2.348441in}{3.425627in}}%
\pgfpathmoveto{\pgfqpoint{2.348441in}{1.632501in}}%
\pgfpathlineto{\pgfqpoint{2.348441in}{1.632501in}}%
\pgfpathlineto{\pgfqpoint{2.348441in}{1.726876in}}%
\pgfpathlineto{\pgfqpoint{2.493750in}{1.726876in}}%
\pgfpathlineto{\pgfqpoint{2.493750in}{1.632501in}}%
\pgfpathmoveto{\pgfqpoint{2.348441in}{1.726876in}}%
\pgfpathlineto{\pgfqpoint{2.348441in}{1.726876in}}%
\pgfpathlineto{\pgfqpoint{2.348441in}{1.821249in}}%
\pgfpathlineto{\pgfqpoint{2.493750in}{1.821249in}}%
\pgfpathlineto{\pgfqpoint{2.493750in}{1.726876in}}%
\pgfpathmoveto{\pgfqpoint{2.348441in}{1.821249in}}%
\pgfpathlineto{\pgfqpoint{2.348441in}{1.821249in}}%
\pgfpathlineto{\pgfqpoint{2.348441in}{1.915627in}}%
\pgfpathlineto{\pgfqpoint{2.493750in}{1.915627in}}%
\pgfpathlineto{\pgfqpoint{2.493750in}{1.821249in}}%
\pgfpathmoveto{\pgfqpoint{2.348441in}{1.915627in}}%
\pgfpathlineto{\pgfqpoint{2.348441in}{1.915627in}}%
\pgfpathlineto{\pgfqpoint{2.348441in}{2.010000in}}%
\pgfpathlineto{\pgfqpoint{2.493750in}{2.010000in}}%
\pgfpathlineto{\pgfqpoint{2.493750in}{1.915627in}}%
\pgfpathmoveto{\pgfqpoint{2.348441in}{2.010000in}}%
\pgfpathlineto{\pgfqpoint{2.348441in}{2.010000in}}%
\pgfpathlineto{\pgfqpoint{2.348441in}{2.104377in}}%
\pgfpathlineto{\pgfqpoint{2.493750in}{2.104377in}}%
\pgfpathlineto{\pgfqpoint{2.493750in}{2.010000in}}%
\pgfpathmoveto{\pgfqpoint{2.348441in}{2.104377in}}%
\pgfpathlineto{\pgfqpoint{2.348441in}{2.104377in}}%
\pgfpathlineto{\pgfqpoint{2.348441in}{2.198748in}}%
\pgfpathlineto{\pgfqpoint{2.493750in}{2.198748in}}%
\pgfpathlineto{\pgfqpoint{2.493750in}{2.104377in}}%
\pgfpathmoveto{\pgfqpoint{2.348441in}{2.198748in}}%
\pgfpathlineto{\pgfqpoint{2.348441in}{2.198748in}}%
\pgfpathlineto{\pgfqpoint{2.348441in}{2.293122in}}%
\pgfpathlineto{\pgfqpoint{2.493750in}{2.293122in}}%
\pgfpathlineto{\pgfqpoint{2.493750in}{2.198748in}}%
\pgfpathmoveto{\pgfqpoint{2.348441in}{2.293122in}}%
\pgfpathlineto{\pgfqpoint{2.348441in}{2.293122in}}%
\pgfpathlineto{\pgfqpoint{2.348441in}{2.387499in}}%
\pgfpathlineto{\pgfqpoint{2.493750in}{2.387499in}}%
\pgfpathlineto{\pgfqpoint{2.493750in}{2.293122in}}%
\pgfpathmoveto{\pgfqpoint{2.348441in}{2.387499in}}%
\pgfpathlineto{\pgfqpoint{2.348441in}{2.387499in}}%
\pgfpathlineto{\pgfqpoint{2.348441in}{2.481878in}}%
\pgfpathlineto{\pgfqpoint{2.493750in}{2.481878in}}%
\pgfpathlineto{\pgfqpoint{2.493750in}{2.387499in}}%
\pgfpathmoveto{\pgfqpoint{2.348441in}{2.481878in}}%
\pgfpathlineto{\pgfqpoint{2.348441in}{2.481878in}}%
\pgfpathlineto{\pgfqpoint{2.348441in}{2.576251in}}%
\pgfpathlineto{\pgfqpoint{2.493750in}{2.576251in}}%
\pgfpathlineto{\pgfqpoint{2.493750in}{2.481878in}}%
\pgfpathmoveto{\pgfqpoint{2.348441in}{2.576251in}}%
\pgfpathlineto{\pgfqpoint{2.348441in}{2.576251in}}%
\pgfpathlineto{\pgfqpoint{2.348441in}{2.670624in}}%
\pgfpathlineto{\pgfqpoint{2.493750in}{2.670624in}}%
\pgfpathlineto{\pgfqpoint{2.493750in}{2.576251in}}%
\pgfpathmoveto{\pgfqpoint{2.348441in}{2.670624in}}%
\pgfpathlineto{\pgfqpoint{2.348441in}{2.670624in}}%
\pgfpathlineto{\pgfqpoint{2.348441in}{2.764999in}}%
\pgfpathlineto{\pgfqpoint{2.493750in}{2.764999in}}%
\pgfpathlineto{\pgfqpoint{2.493750in}{2.670624in}}%
\pgfpathmoveto{\pgfqpoint{2.348441in}{2.764999in}}%
\pgfpathlineto{\pgfqpoint{2.348441in}{2.764999in}}%
\pgfpathlineto{\pgfqpoint{2.348441in}{2.859374in}}%
\pgfpathlineto{\pgfqpoint{2.493750in}{2.859374in}}%
\pgfpathlineto{\pgfqpoint{2.493750in}{2.764999in}}%
\pgfpathmoveto{\pgfqpoint{2.348441in}{2.859374in}}%
\pgfpathlineto{\pgfqpoint{2.348441in}{2.859374in}}%
\pgfpathlineto{\pgfqpoint{2.348441in}{2.953751in}}%
\pgfpathlineto{\pgfqpoint{2.493750in}{2.953751in}}%
\pgfpathlineto{\pgfqpoint{2.493750in}{2.859374in}}%
\pgfpathmoveto{\pgfqpoint{2.348441in}{2.953751in}}%
\pgfpathlineto{\pgfqpoint{2.348441in}{2.953751in}}%
\pgfpathlineto{\pgfqpoint{2.348441in}{3.048124in}}%
\pgfpathlineto{\pgfqpoint{2.493750in}{3.048124in}}%
\pgfpathlineto{\pgfqpoint{2.493750in}{2.953751in}}%
\pgfpathmoveto{\pgfqpoint{2.348441in}{3.048124in}}%
\pgfpathlineto{\pgfqpoint{2.348441in}{3.048124in}}%
\pgfpathlineto{\pgfqpoint{2.348441in}{3.142502in}}%
\pgfpathlineto{\pgfqpoint{2.493750in}{3.142502in}}%
\pgfpathlineto{\pgfqpoint{2.493750in}{3.048124in}}%
\pgfpathmoveto{\pgfqpoint{2.348441in}{3.142502in}}%
\pgfpathlineto{\pgfqpoint{2.348441in}{3.142502in}}%
\pgfpathlineto{\pgfqpoint{2.348441in}{3.236874in}}%
\pgfpathlineto{\pgfqpoint{2.493750in}{3.236874in}}%
\pgfpathlineto{\pgfqpoint{2.493750in}{3.142502in}}%
\pgfpathmoveto{\pgfqpoint{2.348441in}{3.236874in}}%
\pgfpathlineto{\pgfqpoint{2.348441in}{3.236874in}}%
\pgfpathlineto{\pgfqpoint{2.348441in}{3.331252in}}%
\pgfpathlineto{\pgfqpoint{2.493750in}{3.331252in}}%
\pgfpathlineto{\pgfqpoint{2.493750in}{3.236874in}}%
\pgfpathmoveto{\pgfqpoint{2.348441in}{3.331252in}}%
\pgfpathlineto{\pgfqpoint{2.348441in}{3.331252in}}%
\pgfpathlineto{\pgfqpoint{2.348441in}{3.425627in}}%
\pgfpathlineto{\pgfqpoint{2.493750in}{3.425627in}}%
\pgfpathlineto{\pgfqpoint{2.493750in}{3.331252in}}%
\pgfpathmoveto{\pgfqpoint{2.348441in}{3.425627in}}%
\pgfpathlineto{\pgfqpoint{2.348441in}{3.425627in}}%
\pgfpathlineto{\pgfqpoint{2.348441in}{3.519998in}}%
\pgfpathlineto{\pgfqpoint{2.493750in}{3.519998in}}%
\pgfpathlineto{\pgfqpoint{2.493750in}{3.425627in}}%
\pgfpathmoveto{\pgfqpoint{2.493750in}{1.726876in}}%
\pgfpathlineto{\pgfqpoint{2.493750in}{1.726876in}}%
\pgfpathlineto{\pgfqpoint{2.493750in}{1.821249in}}%
\pgfpathlineto{\pgfqpoint{2.639063in}{1.821249in}}%
\pgfpathlineto{\pgfqpoint{2.639063in}{1.726876in}}%
\pgfpathmoveto{\pgfqpoint{2.493750in}{1.821249in}}%
\pgfpathlineto{\pgfqpoint{2.493750in}{1.821249in}}%
\pgfpathlineto{\pgfqpoint{2.493750in}{1.915627in}}%
\pgfpathlineto{\pgfqpoint{2.639063in}{1.915627in}}%
\pgfpathlineto{\pgfqpoint{2.639063in}{1.821249in}}%
\pgfpathmoveto{\pgfqpoint{2.493750in}{1.915627in}}%
\pgfpathlineto{\pgfqpoint{2.493750in}{1.915627in}}%
\pgfpathlineto{\pgfqpoint{2.493750in}{2.010000in}}%
\pgfpathlineto{\pgfqpoint{2.639063in}{2.010000in}}%
\pgfpathlineto{\pgfqpoint{2.639063in}{1.915627in}}%
\pgfpathmoveto{\pgfqpoint{2.493750in}{2.010000in}}%
\pgfpathlineto{\pgfqpoint{2.493750in}{2.010000in}}%
\pgfpathlineto{\pgfqpoint{2.493750in}{2.104377in}}%
\pgfpathlineto{\pgfqpoint{2.639063in}{2.104377in}}%
\pgfpathlineto{\pgfqpoint{2.639063in}{2.010000in}}%
\pgfpathmoveto{\pgfqpoint{2.493750in}{2.104377in}}%
\pgfpathlineto{\pgfqpoint{2.493750in}{2.104377in}}%
\pgfpathlineto{\pgfqpoint{2.493750in}{2.198748in}}%
\pgfpathlineto{\pgfqpoint{2.639063in}{2.198748in}}%
\pgfpathlineto{\pgfqpoint{2.639063in}{2.104377in}}%
\pgfpathmoveto{\pgfqpoint{2.493750in}{2.198748in}}%
\pgfpathlineto{\pgfqpoint{2.493750in}{2.198748in}}%
\pgfpathlineto{\pgfqpoint{2.493750in}{2.293122in}}%
\pgfpathlineto{\pgfqpoint{2.639063in}{2.293122in}}%
\pgfpathlineto{\pgfqpoint{2.639063in}{2.198748in}}%
\pgfpathmoveto{\pgfqpoint{2.493750in}{2.293122in}}%
\pgfpathlineto{\pgfqpoint{2.493750in}{2.293122in}}%
\pgfpathlineto{\pgfqpoint{2.493750in}{2.387499in}}%
\pgfpathlineto{\pgfqpoint{2.639063in}{2.387499in}}%
\pgfpathlineto{\pgfqpoint{2.639063in}{2.293122in}}%
\pgfpathmoveto{\pgfqpoint{2.493750in}{2.387499in}}%
\pgfpathlineto{\pgfqpoint{2.493750in}{2.387499in}}%
\pgfpathlineto{\pgfqpoint{2.493750in}{2.481878in}}%
\pgfpathlineto{\pgfqpoint{2.639063in}{2.481878in}}%
\pgfpathlineto{\pgfqpoint{2.639063in}{2.387499in}}%
\pgfpathmoveto{\pgfqpoint{2.493750in}{2.481878in}}%
\pgfpathlineto{\pgfqpoint{2.493750in}{2.481878in}}%
\pgfpathlineto{\pgfqpoint{2.493750in}{2.576251in}}%
\pgfpathlineto{\pgfqpoint{2.639063in}{2.576251in}}%
\pgfpathlineto{\pgfqpoint{2.639063in}{2.481878in}}%
\pgfpathmoveto{\pgfqpoint{2.493750in}{2.576251in}}%
\pgfpathlineto{\pgfqpoint{2.493750in}{2.576251in}}%
\pgfpathlineto{\pgfqpoint{2.493750in}{2.670624in}}%
\pgfpathlineto{\pgfqpoint{2.639063in}{2.670624in}}%
\pgfpathlineto{\pgfqpoint{2.639063in}{2.576251in}}%
\pgfpathmoveto{\pgfqpoint{2.493750in}{2.670624in}}%
\pgfpathlineto{\pgfqpoint{2.493750in}{2.670624in}}%
\pgfpathlineto{\pgfqpoint{2.493750in}{2.764999in}}%
\pgfpathlineto{\pgfqpoint{2.639063in}{2.764999in}}%
\pgfpathlineto{\pgfqpoint{2.639063in}{2.670624in}}%
\pgfpathmoveto{\pgfqpoint{2.493750in}{2.764999in}}%
\pgfpathlineto{\pgfqpoint{2.493750in}{2.764999in}}%
\pgfpathlineto{\pgfqpoint{2.493750in}{2.859374in}}%
\pgfpathlineto{\pgfqpoint{2.639063in}{2.859374in}}%
\pgfpathlineto{\pgfqpoint{2.639063in}{2.764999in}}%
\pgfpathmoveto{\pgfqpoint{2.493750in}{2.859374in}}%
\pgfpathlineto{\pgfqpoint{2.493750in}{2.859374in}}%
\pgfpathlineto{\pgfqpoint{2.493750in}{2.953751in}}%
\pgfpathlineto{\pgfqpoint{2.639063in}{2.953751in}}%
\pgfpathlineto{\pgfqpoint{2.639063in}{2.859374in}}%
\pgfpathmoveto{\pgfqpoint{2.493750in}{2.953751in}}%
\pgfpathlineto{\pgfqpoint{2.493750in}{2.953751in}}%
\pgfpathlineto{\pgfqpoint{2.493750in}{3.048124in}}%
\pgfpathlineto{\pgfqpoint{2.639063in}{3.048124in}}%
\pgfpathlineto{\pgfqpoint{2.639063in}{2.953751in}}%
\pgfpathmoveto{\pgfqpoint{2.493750in}{3.048124in}}%
\pgfpathlineto{\pgfqpoint{2.493750in}{3.048124in}}%
\pgfpathlineto{\pgfqpoint{2.493750in}{3.142502in}}%
\pgfpathlineto{\pgfqpoint{2.639063in}{3.142502in}}%
\pgfpathlineto{\pgfqpoint{2.639063in}{3.048124in}}%
\pgfpathmoveto{\pgfqpoint{2.493750in}{3.142502in}}%
\pgfpathlineto{\pgfqpoint{2.493750in}{3.142502in}}%
\pgfpathlineto{\pgfqpoint{2.493750in}{3.236874in}}%
\pgfpathlineto{\pgfqpoint{2.639063in}{3.236874in}}%
\pgfpathlineto{\pgfqpoint{2.639063in}{3.142502in}}%
\pgfpathmoveto{\pgfqpoint{2.493750in}{3.236874in}}%
\pgfpathlineto{\pgfqpoint{2.493750in}{3.236874in}}%
\pgfpathlineto{\pgfqpoint{2.493750in}{3.331252in}}%
\pgfpathlineto{\pgfqpoint{2.639063in}{3.331252in}}%
\pgfpathlineto{\pgfqpoint{2.639063in}{3.236874in}}%
\pgfpathmoveto{\pgfqpoint{2.493750in}{3.331252in}}%
\pgfpathlineto{\pgfqpoint{2.493750in}{3.331252in}}%
\pgfpathlineto{\pgfqpoint{2.493750in}{3.425627in}}%
\pgfpathlineto{\pgfqpoint{2.639063in}{3.425627in}}%
\pgfpathlineto{\pgfqpoint{2.639063in}{3.331252in}}%
\pgfpathmoveto{\pgfqpoint{2.493750in}{3.425627in}}%
\pgfpathlineto{\pgfqpoint{2.493750in}{3.425627in}}%
\pgfpathlineto{\pgfqpoint{2.493750in}{3.519998in}}%
\pgfpathlineto{\pgfqpoint{2.639063in}{3.519998in}}%
\pgfpathlineto{\pgfqpoint{2.639063in}{3.425627in}}%
\pgfpathmoveto{\pgfqpoint{2.639063in}{1.915627in}}%
\pgfpathlineto{\pgfqpoint{2.639063in}{1.915627in}}%
\pgfpathlineto{\pgfqpoint{2.639063in}{2.010000in}}%
\pgfpathlineto{\pgfqpoint{2.784375in}{2.010000in}}%
\pgfpathlineto{\pgfqpoint{2.784375in}{1.915627in}}%
\pgfpathmoveto{\pgfqpoint{2.639063in}{2.010000in}}%
\pgfpathlineto{\pgfqpoint{2.639063in}{2.010000in}}%
\pgfpathlineto{\pgfqpoint{2.639063in}{2.104377in}}%
\pgfpathlineto{\pgfqpoint{2.784375in}{2.104377in}}%
\pgfpathlineto{\pgfqpoint{2.784375in}{2.010000in}}%
\pgfpathmoveto{\pgfqpoint{2.639063in}{2.104377in}}%
\pgfpathlineto{\pgfqpoint{2.639063in}{2.104377in}}%
\pgfpathlineto{\pgfqpoint{2.639063in}{2.198748in}}%
\pgfpathlineto{\pgfqpoint{2.784375in}{2.198748in}}%
\pgfpathlineto{\pgfqpoint{2.784375in}{2.104377in}}%
\pgfpathmoveto{\pgfqpoint{2.639063in}{2.198748in}}%
\pgfpathlineto{\pgfqpoint{2.639063in}{2.198748in}}%
\pgfpathlineto{\pgfqpoint{2.639063in}{2.293122in}}%
\pgfpathlineto{\pgfqpoint{2.784375in}{2.293122in}}%
\pgfpathlineto{\pgfqpoint{2.784375in}{2.198748in}}%
\pgfpathmoveto{\pgfqpoint{2.639063in}{2.293122in}}%
\pgfpathlineto{\pgfqpoint{2.639063in}{2.293122in}}%
\pgfpathlineto{\pgfqpoint{2.639063in}{2.387499in}}%
\pgfpathlineto{\pgfqpoint{2.784375in}{2.387499in}}%
\pgfpathlineto{\pgfqpoint{2.784375in}{2.293122in}}%
\pgfpathmoveto{\pgfqpoint{2.639063in}{2.387499in}}%
\pgfpathlineto{\pgfqpoint{2.639063in}{2.387499in}}%
\pgfpathlineto{\pgfqpoint{2.639063in}{2.481878in}}%
\pgfpathlineto{\pgfqpoint{2.784375in}{2.481878in}}%
\pgfpathlineto{\pgfqpoint{2.784375in}{2.387499in}}%
\pgfpathmoveto{\pgfqpoint{2.639063in}{2.481878in}}%
\pgfpathlineto{\pgfqpoint{2.639063in}{2.481878in}}%
\pgfpathlineto{\pgfqpoint{2.639063in}{2.576251in}}%
\pgfpathlineto{\pgfqpoint{2.784375in}{2.576251in}}%
\pgfpathlineto{\pgfqpoint{2.784375in}{2.481878in}}%
\pgfpathmoveto{\pgfqpoint{2.639063in}{2.576251in}}%
\pgfpathlineto{\pgfqpoint{2.639063in}{2.576251in}}%
\pgfpathlineto{\pgfqpoint{2.639063in}{2.670624in}}%
\pgfpathlineto{\pgfqpoint{2.784375in}{2.670624in}}%
\pgfpathlineto{\pgfqpoint{2.784375in}{2.576251in}}%
\pgfpathmoveto{\pgfqpoint{2.639063in}{2.670624in}}%
\pgfpathlineto{\pgfqpoint{2.639063in}{2.670624in}}%
\pgfpathlineto{\pgfqpoint{2.639063in}{2.764999in}}%
\pgfpathlineto{\pgfqpoint{2.784375in}{2.764999in}}%
\pgfpathlineto{\pgfqpoint{2.784375in}{2.670624in}}%
\pgfpathmoveto{\pgfqpoint{2.639063in}{2.764999in}}%
\pgfpathlineto{\pgfqpoint{2.639063in}{2.764999in}}%
\pgfpathlineto{\pgfqpoint{2.639063in}{2.859374in}}%
\pgfpathlineto{\pgfqpoint{2.784375in}{2.859374in}}%
\pgfpathlineto{\pgfqpoint{2.784375in}{2.764999in}}%
\pgfpathmoveto{\pgfqpoint{2.639063in}{2.859374in}}%
\pgfpathlineto{\pgfqpoint{2.639063in}{2.859374in}}%
\pgfpathlineto{\pgfqpoint{2.639063in}{2.953751in}}%
\pgfpathlineto{\pgfqpoint{2.784375in}{2.953751in}}%
\pgfpathlineto{\pgfqpoint{2.784375in}{2.859374in}}%
\pgfpathmoveto{\pgfqpoint{2.639063in}{2.953751in}}%
\pgfpathlineto{\pgfqpoint{2.639063in}{2.953751in}}%
\pgfpathlineto{\pgfqpoint{2.639063in}{3.048124in}}%
\pgfpathlineto{\pgfqpoint{2.784375in}{3.048124in}}%
\pgfpathlineto{\pgfqpoint{2.784375in}{2.953751in}}%
\pgfpathmoveto{\pgfqpoint{2.639063in}{3.048124in}}%
\pgfpathlineto{\pgfqpoint{2.639063in}{3.048124in}}%
\pgfpathlineto{\pgfqpoint{2.639063in}{3.142502in}}%
\pgfpathlineto{\pgfqpoint{2.784375in}{3.142502in}}%
\pgfpathlineto{\pgfqpoint{2.784375in}{3.048124in}}%
\pgfpathmoveto{\pgfqpoint{2.639063in}{3.142502in}}%
\pgfpathlineto{\pgfqpoint{2.639063in}{3.142502in}}%
\pgfpathlineto{\pgfqpoint{2.639063in}{3.236874in}}%
\pgfpathlineto{\pgfqpoint{2.784375in}{3.236874in}}%
\pgfpathlineto{\pgfqpoint{2.784375in}{3.142502in}}%
\pgfpathmoveto{\pgfqpoint{2.639063in}{3.236874in}}%
\pgfpathlineto{\pgfqpoint{2.639063in}{3.236874in}}%
\pgfpathlineto{\pgfqpoint{2.639063in}{3.331252in}}%
\pgfpathlineto{\pgfqpoint{2.784375in}{3.331252in}}%
\pgfpathlineto{\pgfqpoint{2.784375in}{3.236874in}}%
\pgfpathmoveto{\pgfqpoint{2.639063in}{3.331252in}}%
\pgfpathlineto{\pgfqpoint{2.639063in}{3.331252in}}%
\pgfpathlineto{\pgfqpoint{2.639063in}{3.425627in}}%
\pgfpathlineto{\pgfqpoint{2.784375in}{3.425627in}}%
\pgfpathlineto{\pgfqpoint{2.784375in}{3.331252in}}%
\pgfpathmoveto{\pgfqpoint{2.639063in}{3.425627in}}%
\pgfpathlineto{\pgfqpoint{2.639063in}{3.425627in}}%
\pgfpathlineto{\pgfqpoint{2.639063in}{3.519998in}}%
\pgfpathlineto{\pgfqpoint{2.784375in}{3.519998in}}%
\pgfpathlineto{\pgfqpoint{2.784375in}{3.425627in}}%
\pgfpathmoveto{\pgfqpoint{2.784375in}{1.915627in}}%
\pgfpathlineto{\pgfqpoint{2.784375in}{1.915627in}}%
\pgfpathlineto{\pgfqpoint{2.784375in}{2.010000in}}%
\pgfpathlineto{\pgfqpoint{2.929688in}{2.010000in}}%
\pgfpathlineto{\pgfqpoint{2.929688in}{1.915627in}}%
\pgfpathmoveto{\pgfqpoint{2.784375in}{2.010000in}}%
\pgfpathlineto{\pgfqpoint{2.784375in}{2.010000in}}%
\pgfpathlineto{\pgfqpoint{2.784375in}{2.104377in}}%
\pgfpathlineto{\pgfqpoint{2.929688in}{2.104377in}}%
\pgfpathlineto{\pgfqpoint{2.929688in}{2.010000in}}%
\pgfpathmoveto{\pgfqpoint{2.784375in}{2.104377in}}%
\pgfpathlineto{\pgfqpoint{2.784375in}{2.104377in}}%
\pgfpathlineto{\pgfqpoint{2.784375in}{2.198748in}}%
\pgfpathlineto{\pgfqpoint{2.929688in}{2.198748in}}%
\pgfpathlineto{\pgfqpoint{2.929688in}{2.104377in}}%
\pgfpathmoveto{\pgfqpoint{2.784375in}{2.198748in}}%
\pgfpathlineto{\pgfqpoint{2.784375in}{2.198748in}}%
\pgfpathlineto{\pgfqpoint{2.784375in}{2.293122in}}%
\pgfpathlineto{\pgfqpoint{2.929688in}{2.293122in}}%
\pgfpathlineto{\pgfqpoint{2.929688in}{2.198748in}}%
\pgfpathmoveto{\pgfqpoint{2.784375in}{2.293122in}}%
\pgfpathlineto{\pgfqpoint{2.784375in}{2.293122in}}%
\pgfpathlineto{\pgfqpoint{2.784375in}{2.387499in}}%
\pgfpathlineto{\pgfqpoint{2.929688in}{2.387499in}}%
\pgfpathlineto{\pgfqpoint{2.929688in}{2.293122in}}%
\pgfpathmoveto{\pgfqpoint{2.784375in}{2.387499in}}%
\pgfpathlineto{\pgfqpoint{2.784375in}{2.387499in}}%
\pgfpathlineto{\pgfqpoint{2.784375in}{2.481878in}}%
\pgfpathlineto{\pgfqpoint{2.929688in}{2.481878in}}%
\pgfpathlineto{\pgfqpoint{2.929688in}{2.387499in}}%
\pgfpathmoveto{\pgfqpoint{2.784375in}{2.481878in}}%
\pgfpathlineto{\pgfqpoint{2.784375in}{2.481878in}}%
\pgfpathlineto{\pgfqpoint{2.784375in}{2.576251in}}%
\pgfpathlineto{\pgfqpoint{2.929688in}{2.576251in}}%
\pgfpathlineto{\pgfqpoint{2.929688in}{2.481878in}}%
\pgfpathmoveto{\pgfqpoint{2.784375in}{2.576251in}}%
\pgfpathlineto{\pgfqpoint{2.784375in}{2.576251in}}%
\pgfpathlineto{\pgfqpoint{2.784375in}{2.670624in}}%
\pgfpathlineto{\pgfqpoint{2.929688in}{2.670624in}}%
\pgfpathlineto{\pgfqpoint{2.929688in}{2.576251in}}%
\pgfpathmoveto{\pgfqpoint{2.784375in}{2.670624in}}%
\pgfpathlineto{\pgfqpoint{2.784375in}{2.670624in}}%
\pgfpathlineto{\pgfqpoint{2.784375in}{2.764999in}}%
\pgfpathlineto{\pgfqpoint{2.929688in}{2.764999in}}%
\pgfpathlineto{\pgfqpoint{2.929688in}{2.670624in}}%
\pgfpathmoveto{\pgfqpoint{2.784375in}{2.764999in}}%
\pgfpathlineto{\pgfqpoint{2.784375in}{2.764999in}}%
\pgfpathlineto{\pgfqpoint{2.784375in}{2.859374in}}%
\pgfpathlineto{\pgfqpoint{2.929688in}{2.859374in}}%
\pgfpathlineto{\pgfqpoint{2.929688in}{2.764999in}}%
\pgfpathmoveto{\pgfqpoint{2.784375in}{2.859374in}}%
\pgfpathlineto{\pgfqpoint{2.784375in}{2.859374in}}%
\pgfpathlineto{\pgfqpoint{2.784375in}{2.953751in}}%
\pgfpathlineto{\pgfqpoint{2.929688in}{2.953751in}}%
\pgfpathlineto{\pgfqpoint{2.929688in}{2.859374in}}%
\pgfpathmoveto{\pgfqpoint{2.784375in}{2.953751in}}%
\pgfpathlineto{\pgfqpoint{2.784375in}{2.953751in}}%
\pgfpathlineto{\pgfqpoint{2.784375in}{3.048124in}}%
\pgfpathlineto{\pgfqpoint{2.929688in}{3.048124in}}%
\pgfpathlineto{\pgfqpoint{2.929688in}{2.953751in}}%
\pgfpathmoveto{\pgfqpoint{2.784375in}{3.048124in}}%
\pgfpathlineto{\pgfqpoint{2.784375in}{3.048124in}}%
\pgfpathlineto{\pgfqpoint{2.784375in}{3.142502in}}%
\pgfpathlineto{\pgfqpoint{2.929688in}{3.142502in}}%
\pgfpathlineto{\pgfqpoint{2.929688in}{3.048124in}}%
\pgfpathmoveto{\pgfqpoint{2.784375in}{3.142502in}}%
\pgfpathlineto{\pgfqpoint{2.784375in}{3.142502in}}%
\pgfpathlineto{\pgfqpoint{2.784375in}{3.236874in}}%
\pgfpathlineto{\pgfqpoint{2.929688in}{3.236874in}}%
\pgfpathlineto{\pgfqpoint{2.929688in}{3.142502in}}%
\pgfpathmoveto{\pgfqpoint{2.784375in}{3.236874in}}%
\pgfpathlineto{\pgfqpoint{2.784375in}{3.236874in}}%
\pgfpathlineto{\pgfqpoint{2.784375in}{3.331252in}}%
\pgfpathlineto{\pgfqpoint{2.929688in}{3.331252in}}%
\pgfpathlineto{\pgfqpoint{2.929688in}{3.236874in}}%
\pgfpathmoveto{\pgfqpoint{2.929688in}{2.104377in}}%
\pgfpathlineto{\pgfqpoint{2.929688in}{2.104377in}}%
\pgfpathlineto{\pgfqpoint{2.929688in}{2.198748in}}%
\pgfpathlineto{\pgfqpoint{3.075004in}{2.198748in}}%
\pgfpathlineto{\pgfqpoint{3.075004in}{2.104377in}}%
\pgfpathmoveto{\pgfqpoint{2.929688in}{2.198748in}}%
\pgfpathlineto{\pgfqpoint{2.929688in}{2.198748in}}%
\pgfpathlineto{\pgfqpoint{2.929688in}{2.293122in}}%
\pgfpathlineto{\pgfqpoint{3.075004in}{2.293122in}}%
\pgfpathlineto{\pgfqpoint{3.075004in}{2.198748in}}%
\pgfpathmoveto{\pgfqpoint{2.929688in}{2.293122in}}%
\pgfpathlineto{\pgfqpoint{2.929688in}{2.293122in}}%
\pgfpathlineto{\pgfqpoint{2.929688in}{2.387499in}}%
\pgfpathlineto{\pgfqpoint{3.075004in}{2.387499in}}%
\pgfpathlineto{\pgfqpoint{3.075004in}{2.293122in}}%
\pgfpathmoveto{\pgfqpoint{2.929688in}{2.387499in}}%
\pgfpathlineto{\pgfqpoint{2.929688in}{2.387499in}}%
\pgfpathlineto{\pgfqpoint{2.929688in}{2.481878in}}%
\pgfpathlineto{\pgfqpoint{3.075004in}{2.481878in}}%
\pgfpathlineto{\pgfqpoint{3.075004in}{2.387499in}}%
\pgfpathmoveto{\pgfqpoint{2.929688in}{2.481878in}}%
\pgfpathlineto{\pgfqpoint{2.929688in}{2.481878in}}%
\pgfpathlineto{\pgfqpoint{2.929688in}{2.576251in}}%
\pgfpathlineto{\pgfqpoint{3.075004in}{2.576251in}}%
\pgfpathlineto{\pgfqpoint{3.075004in}{2.481878in}}%
\pgfpathmoveto{\pgfqpoint{2.929688in}{2.576251in}}%
\pgfpathlineto{\pgfqpoint{2.929688in}{2.576251in}}%
\pgfpathlineto{\pgfqpoint{2.929688in}{2.670624in}}%
\pgfpathlineto{\pgfqpoint{3.075004in}{2.670624in}}%
\pgfpathlineto{\pgfqpoint{3.075004in}{2.576251in}}%
\pgfpathmoveto{\pgfqpoint{2.929688in}{2.670624in}}%
\pgfpathlineto{\pgfqpoint{2.929688in}{2.670624in}}%
\pgfpathlineto{\pgfqpoint{2.929688in}{2.764999in}}%
\pgfpathlineto{\pgfqpoint{3.075004in}{2.764999in}}%
\pgfpathlineto{\pgfqpoint{3.075004in}{2.670624in}}%
\pgfpathmoveto{\pgfqpoint{2.929688in}{2.764999in}}%
\pgfpathlineto{\pgfqpoint{2.929688in}{2.764999in}}%
\pgfpathlineto{\pgfqpoint{2.929688in}{2.859374in}}%
\pgfpathlineto{\pgfqpoint{3.075004in}{2.859374in}}%
\pgfpathlineto{\pgfqpoint{3.075004in}{2.764999in}}%
\pgfpathmoveto{\pgfqpoint{2.929688in}{2.859374in}}%
\pgfpathlineto{\pgfqpoint{2.929688in}{2.859374in}}%
\pgfpathlineto{\pgfqpoint{2.929688in}{2.953751in}}%
\pgfpathlineto{\pgfqpoint{3.075004in}{2.953751in}}%
\pgfpathlineto{\pgfqpoint{3.075004in}{2.859374in}}%
\pgfpathmoveto{\pgfqpoint{2.929688in}{2.953751in}}%
\pgfpathlineto{\pgfqpoint{2.929688in}{2.953751in}}%
\pgfpathlineto{\pgfqpoint{2.929688in}{3.048124in}}%
\pgfpathlineto{\pgfqpoint{3.075004in}{3.048124in}}%
\pgfpathlineto{\pgfqpoint{3.075004in}{2.953751in}}%
\pgfpathmoveto{\pgfqpoint{2.929688in}{3.048124in}}%
\pgfpathlineto{\pgfqpoint{2.929688in}{3.048124in}}%
\pgfpathlineto{\pgfqpoint{2.929688in}{3.142502in}}%
\pgfpathlineto{\pgfqpoint{3.075004in}{3.142502in}}%
\pgfpathlineto{\pgfqpoint{3.075004in}{3.048124in}}%
\pgfpathmoveto{\pgfqpoint{3.075004in}{2.104377in}}%
\pgfpathlineto{\pgfqpoint{3.075004in}{2.104377in}}%
\pgfpathlineto{\pgfqpoint{3.075004in}{2.198748in}}%
\pgfpathlineto{\pgfqpoint{3.220310in}{2.198748in}}%
\pgfpathlineto{\pgfqpoint{3.220310in}{2.104377in}}%
\pgfpathmoveto{\pgfqpoint{3.075004in}{2.198748in}}%
\pgfpathlineto{\pgfqpoint{3.075004in}{2.198748in}}%
\pgfpathlineto{\pgfqpoint{3.075004in}{2.293122in}}%
\pgfpathlineto{\pgfqpoint{3.220310in}{2.293122in}}%
\pgfpathlineto{\pgfqpoint{3.220310in}{2.198748in}}%
\pgfpathmoveto{\pgfqpoint{3.075004in}{2.293122in}}%
\pgfpathlineto{\pgfqpoint{3.075004in}{2.293122in}}%
\pgfpathlineto{\pgfqpoint{3.075004in}{2.387499in}}%
\pgfpathlineto{\pgfqpoint{3.220310in}{2.387499in}}%
\pgfpathlineto{\pgfqpoint{3.220310in}{2.293122in}}%
\pgfpathmoveto{\pgfqpoint{3.075004in}{2.387499in}}%
\pgfpathlineto{\pgfqpoint{3.075004in}{2.387499in}}%
\pgfpathlineto{\pgfqpoint{3.075004in}{2.481878in}}%
\pgfpathlineto{\pgfqpoint{3.220310in}{2.481878in}}%
\pgfpathlineto{\pgfqpoint{3.220310in}{2.387499in}}%
\pgfpathmoveto{\pgfqpoint{3.075004in}{2.481878in}}%
\pgfpathlineto{\pgfqpoint{3.075004in}{2.481878in}}%
\pgfpathlineto{\pgfqpoint{3.075004in}{2.576251in}}%
\pgfpathlineto{\pgfqpoint{3.220310in}{2.576251in}}%
\pgfpathlineto{\pgfqpoint{3.220310in}{2.481878in}}%
\pgfpathmoveto{\pgfqpoint{3.075004in}{2.576251in}}%
\pgfpathlineto{\pgfqpoint{3.075004in}{2.576251in}}%
\pgfpathlineto{\pgfqpoint{3.075004in}{2.670624in}}%
\pgfpathlineto{\pgfqpoint{3.220310in}{2.670624in}}%
\pgfpathlineto{\pgfqpoint{3.220310in}{2.576251in}}%
\pgfpathmoveto{\pgfqpoint{3.075004in}{2.670624in}}%
\pgfpathlineto{\pgfqpoint{3.075004in}{2.670624in}}%
\pgfpathlineto{\pgfqpoint{3.075004in}{2.764999in}}%
\pgfpathlineto{\pgfqpoint{3.220310in}{2.764999in}}%
\pgfpathlineto{\pgfqpoint{3.220310in}{2.670624in}}%
\pgfpathmoveto{\pgfqpoint{3.075004in}{2.764999in}}%
\pgfpathlineto{\pgfqpoint{3.075004in}{2.764999in}}%
\pgfpathlineto{\pgfqpoint{3.075004in}{2.859374in}}%
\pgfpathlineto{\pgfqpoint{3.220310in}{2.859374in}}%
\pgfpathlineto{\pgfqpoint{3.220310in}{2.764999in}}%
\pgfpathmoveto{\pgfqpoint{3.075004in}{2.859374in}}%
\pgfpathlineto{\pgfqpoint{3.075004in}{2.859374in}}%
\pgfpathlineto{\pgfqpoint{3.075004in}{2.953751in}}%
\pgfpathlineto{\pgfqpoint{3.220310in}{2.953751in}}%
\pgfpathlineto{\pgfqpoint{3.220310in}{2.859374in}}%
\pgfpathmoveto{\pgfqpoint{3.220310in}{2.293122in}}%
\pgfpathlineto{\pgfqpoint{3.220310in}{2.293122in}}%
\pgfpathlineto{\pgfqpoint{3.220310in}{2.387499in}}%
\pgfpathlineto{\pgfqpoint{3.365627in}{2.387499in}}%
\pgfpathlineto{\pgfqpoint{3.365627in}{2.293122in}}%
\pgfpathmoveto{\pgfqpoint{3.220310in}{2.387499in}}%
\pgfpathlineto{\pgfqpoint{3.220310in}{2.387499in}}%
\pgfpathlineto{\pgfqpoint{3.220310in}{2.481878in}}%
\pgfpathlineto{\pgfqpoint{3.365627in}{2.481878in}}%
\pgfpathlineto{\pgfqpoint{3.365627in}{2.387499in}}%
\pgfpathmoveto{\pgfqpoint{3.220310in}{2.481878in}}%
\pgfpathlineto{\pgfqpoint{3.220310in}{2.481878in}}%
\pgfpathlineto{\pgfqpoint{3.220310in}{2.576251in}}%
\pgfpathlineto{\pgfqpoint{3.365627in}{2.576251in}}%
\pgfpathlineto{\pgfqpoint{3.365627in}{2.481878in}}%
\pgfpathmoveto{\pgfqpoint{3.220310in}{2.576251in}}%
\pgfpathlineto{\pgfqpoint{3.220310in}{2.576251in}}%
\pgfpathlineto{\pgfqpoint{3.220310in}{2.670624in}}%
\pgfpathlineto{\pgfqpoint{3.365627in}{2.670624in}}%
\pgfpathlineto{\pgfqpoint{3.365627in}{2.576251in}}%
\pgfpathmoveto{\pgfqpoint{3.220310in}{2.670624in}}%
\pgfpathlineto{\pgfqpoint{3.220310in}{2.670624in}}%
\pgfpathlineto{\pgfqpoint{3.220310in}{2.764999in}}%
\pgfpathlineto{\pgfqpoint{3.365627in}{2.764999in}}%
\pgfpathlineto{\pgfqpoint{3.365627in}{2.670624in}}%
\pgfpathmoveto{\pgfqpoint{3.365627in}{2.387499in}}%
\pgfpathlineto{\pgfqpoint{3.365627in}{2.387499in}}%
\pgfpathlineto{\pgfqpoint{3.365627in}{2.481878in}}%
\pgfpathlineto{\pgfqpoint{3.510939in}{2.481878in}}%
\pgfpathlineto{\pgfqpoint{3.510939in}{2.387499in}}%
\pgfpathmoveto{\pgfqpoint{3.365627in}{2.481878in}}%
\pgfpathlineto{\pgfqpoint{3.365627in}{2.481878in}}%
\pgfpathlineto{\pgfqpoint{3.365627in}{2.576251in}}%
\pgfpathlineto{\pgfqpoint{3.510939in}{2.576251in}}%
\pgfpathlineto{\pgfqpoint{3.510939in}{2.481878in}}%
\pgfpathmoveto{\pgfqpoint{0.749998in}{0.594374in}}%
\pgfpathlineto{\pgfqpoint{0.749998in}{0.594374in}}%
\pgfpathlineto{\pgfqpoint{0.749998in}{0.641562in}}%
\pgfpathlineto{\pgfqpoint{0.822655in}{0.641562in}}%
\pgfpathlineto{\pgfqpoint{0.822655in}{0.594374in}}%
\pgfpathmoveto{\pgfqpoint{0.749998in}{0.641562in}}%
\pgfpathlineto{\pgfqpoint{0.749998in}{0.641562in}}%
\pgfpathlineto{\pgfqpoint{0.749998in}{0.688750in}}%
\pgfpathlineto{\pgfqpoint{0.822655in}{0.688750in}}%
\pgfpathlineto{\pgfqpoint{0.822655in}{0.641562in}}%
\pgfpathmoveto{\pgfqpoint{0.822655in}{0.641562in}}%
\pgfpathlineto{\pgfqpoint{0.822655in}{0.641562in}}%
\pgfpathlineto{\pgfqpoint{0.822655in}{0.688750in}}%
\pgfpathlineto{\pgfqpoint{0.895312in}{0.688750in}}%
\pgfpathlineto{\pgfqpoint{0.895312in}{0.641562in}}%
\pgfpathmoveto{\pgfqpoint{0.895312in}{0.641562in}}%
\pgfpathlineto{\pgfqpoint{0.895312in}{0.641562in}}%
\pgfpathlineto{\pgfqpoint{0.895312in}{0.688750in}}%
\pgfpathlineto{\pgfqpoint{0.967967in}{0.688750in}}%
\pgfpathlineto{\pgfqpoint{0.967967in}{0.641562in}}%
\pgfpathmoveto{\pgfqpoint{1.040622in}{0.783124in}}%
\pgfpathlineto{\pgfqpoint{1.040622in}{0.783124in}}%
\pgfpathlineto{\pgfqpoint{1.040622in}{0.830313in}}%
\pgfpathlineto{\pgfqpoint{1.113282in}{0.830313in}}%
\pgfpathlineto{\pgfqpoint{1.113282in}{0.783124in}}%
\pgfpathmoveto{\pgfqpoint{1.040622in}{0.830313in}}%
\pgfpathlineto{\pgfqpoint{1.040622in}{0.830313in}}%
\pgfpathlineto{\pgfqpoint{1.040622in}{0.877501in}}%
\pgfpathlineto{\pgfqpoint{1.113282in}{0.877501in}}%
\pgfpathlineto{\pgfqpoint{1.113282in}{0.830313in}}%
\pgfpathmoveto{\pgfqpoint{1.113282in}{0.830313in}}%
\pgfpathlineto{\pgfqpoint{1.113282in}{0.830313in}}%
\pgfpathlineto{\pgfqpoint{1.113282in}{0.877501in}}%
\pgfpathlineto{\pgfqpoint{1.185941in}{0.877501in}}%
\pgfpathlineto{\pgfqpoint{1.185941in}{0.830313in}}%
\pgfpathmoveto{\pgfqpoint{1.185941in}{0.877501in}}%
\pgfpathlineto{\pgfqpoint{1.185941in}{0.877501in}}%
\pgfpathlineto{\pgfqpoint{1.185941in}{0.924687in}}%
\pgfpathlineto{\pgfqpoint{1.258598in}{0.924687in}}%
\pgfpathlineto{\pgfqpoint{1.258598in}{0.877501in}}%
\pgfpathmoveto{\pgfqpoint{1.185941in}{0.924687in}}%
\pgfpathlineto{\pgfqpoint{1.185941in}{0.924687in}}%
\pgfpathlineto{\pgfqpoint{1.185941in}{0.971872in}}%
\pgfpathlineto{\pgfqpoint{1.258598in}{0.971872in}}%
\pgfpathlineto{\pgfqpoint{1.258598in}{0.924687in}}%
\pgfpathmoveto{\pgfqpoint{1.258598in}{0.924687in}}%
\pgfpathlineto{\pgfqpoint{1.258598in}{0.924687in}}%
\pgfpathlineto{\pgfqpoint{1.258598in}{0.971872in}}%
\pgfpathlineto{\pgfqpoint{1.331254in}{0.971872in}}%
\pgfpathlineto{\pgfqpoint{1.331254in}{0.924687in}}%
\pgfpathmoveto{\pgfqpoint{1.331254in}{0.971872in}}%
\pgfpathlineto{\pgfqpoint{1.331254in}{0.971872in}}%
\pgfpathlineto{\pgfqpoint{1.331254in}{1.019061in}}%
\pgfpathlineto{\pgfqpoint{1.403909in}{1.019061in}}%
\pgfpathlineto{\pgfqpoint{1.403909in}{0.971872in}}%
\pgfpathmoveto{\pgfqpoint{1.331254in}{1.019061in}}%
\pgfpathlineto{\pgfqpoint{1.331254in}{1.019061in}}%
\pgfpathlineto{\pgfqpoint{1.331254in}{1.066249in}}%
\pgfpathlineto{\pgfqpoint{1.403909in}{1.066249in}}%
\pgfpathlineto{\pgfqpoint{1.403909in}{1.019061in}}%
\pgfpathmoveto{\pgfqpoint{1.403909in}{1.019061in}}%
\pgfpathlineto{\pgfqpoint{1.403909in}{1.019061in}}%
\pgfpathlineto{\pgfqpoint{1.403909in}{1.066249in}}%
\pgfpathlineto{\pgfqpoint{1.476564in}{1.066249in}}%
\pgfpathlineto{\pgfqpoint{1.476564in}{1.019061in}}%
\pgfpathmoveto{\pgfqpoint{1.476564in}{1.066249in}}%
\pgfpathlineto{\pgfqpoint{1.476564in}{1.066249in}}%
\pgfpathlineto{\pgfqpoint{1.476564in}{1.113437in}}%
\pgfpathlineto{\pgfqpoint{1.549220in}{1.113437in}}%
\pgfpathlineto{\pgfqpoint{1.549220in}{1.066249in}}%
\pgfpathmoveto{\pgfqpoint{1.476564in}{1.113437in}}%
\pgfpathlineto{\pgfqpoint{1.476564in}{1.113437in}}%
\pgfpathlineto{\pgfqpoint{1.476564in}{1.160625in}}%
\pgfpathlineto{\pgfqpoint{1.549220in}{1.160625in}}%
\pgfpathlineto{\pgfqpoint{1.549220in}{1.113437in}}%
\pgfpathmoveto{\pgfqpoint{1.549220in}{1.113437in}}%
\pgfpathlineto{\pgfqpoint{1.549220in}{1.113437in}}%
\pgfpathlineto{\pgfqpoint{1.549220in}{1.160625in}}%
\pgfpathlineto{\pgfqpoint{1.621877in}{1.160625in}}%
\pgfpathlineto{\pgfqpoint{1.621877in}{1.113437in}}%
\pgfpathmoveto{\pgfqpoint{1.767186in}{1.207812in}}%
\pgfpathlineto{\pgfqpoint{1.767186in}{1.207812in}}%
\pgfpathlineto{\pgfqpoint{1.767186in}{1.254998in}}%
\pgfpathlineto{\pgfqpoint{1.839841in}{1.254998in}}%
\pgfpathlineto{\pgfqpoint{1.839841in}{1.207812in}}%
\pgfpathmoveto{\pgfqpoint{1.912496in}{1.302187in}}%
\pgfpathlineto{\pgfqpoint{1.912496in}{1.302187in}}%
\pgfpathlineto{\pgfqpoint{1.912496in}{1.349376in}}%
\pgfpathlineto{\pgfqpoint{1.985155in}{1.349376in}}%
\pgfpathlineto{\pgfqpoint{1.985155in}{1.302187in}}%
\pgfpathmoveto{\pgfqpoint{2.057813in}{1.443750in}}%
\pgfpathlineto{\pgfqpoint{2.057813in}{1.443750in}}%
\pgfpathlineto{\pgfqpoint{2.057813in}{1.490938in}}%
\pgfpathlineto{\pgfqpoint{2.130471in}{1.490938in}}%
\pgfpathlineto{\pgfqpoint{2.130471in}{1.443750in}}%
\pgfpathmoveto{\pgfqpoint{2.057813in}{1.490938in}}%
\pgfpathlineto{\pgfqpoint{2.057813in}{1.490938in}}%
\pgfpathlineto{\pgfqpoint{2.057813in}{1.538127in}}%
\pgfpathlineto{\pgfqpoint{2.130471in}{1.538127in}}%
\pgfpathlineto{\pgfqpoint{2.130471in}{1.490938in}}%
\pgfpathmoveto{\pgfqpoint{2.130471in}{1.490938in}}%
\pgfpathlineto{\pgfqpoint{2.130471in}{1.490938in}}%
\pgfpathlineto{\pgfqpoint{2.130471in}{1.538127in}}%
\pgfpathlineto{\pgfqpoint{2.203129in}{1.538127in}}%
\pgfpathlineto{\pgfqpoint{2.203129in}{1.490938in}}%
\pgfpathmoveto{\pgfqpoint{2.203129in}{1.538127in}}%
\pgfpathlineto{\pgfqpoint{2.203129in}{1.538127in}}%
\pgfpathlineto{\pgfqpoint{2.203129in}{1.585314in}}%
\pgfpathlineto{\pgfqpoint{2.275785in}{1.585314in}}%
\pgfpathlineto{\pgfqpoint{2.275785in}{1.538127in}}%
\pgfpathmoveto{\pgfqpoint{2.203129in}{1.585314in}}%
\pgfpathlineto{\pgfqpoint{2.203129in}{1.585314in}}%
\pgfpathlineto{\pgfqpoint{2.203129in}{1.632501in}}%
\pgfpathlineto{\pgfqpoint{2.275785in}{1.632501in}}%
\pgfpathlineto{\pgfqpoint{2.275785in}{1.585314in}}%
\pgfpathmoveto{\pgfqpoint{2.275785in}{1.585314in}}%
\pgfpathlineto{\pgfqpoint{2.275785in}{1.585314in}}%
\pgfpathlineto{\pgfqpoint{2.275785in}{1.632501in}}%
\pgfpathlineto{\pgfqpoint{2.348441in}{1.632501in}}%
\pgfpathlineto{\pgfqpoint{2.348441in}{1.585314in}}%
\pgfpathmoveto{\pgfqpoint{2.348441in}{1.585314in}}%
\pgfpathlineto{\pgfqpoint{2.348441in}{1.585314in}}%
\pgfpathlineto{\pgfqpoint{2.348441in}{1.632501in}}%
\pgfpathlineto{\pgfqpoint{2.421095in}{1.632501in}}%
\pgfpathlineto{\pgfqpoint{2.421095in}{1.585314in}}%
\pgfpathmoveto{\pgfqpoint{2.493750in}{1.679689in}}%
\pgfpathlineto{\pgfqpoint{2.493750in}{1.679689in}}%
\pgfpathlineto{\pgfqpoint{2.493750in}{1.726876in}}%
\pgfpathlineto{\pgfqpoint{2.566406in}{1.726876in}}%
\pgfpathlineto{\pgfqpoint{2.566406in}{1.679689in}}%
\pgfpathmoveto{\pgfqpoint{2.639063in}{1.821249in}}%
\pgfpathlineto{\pgfqpoint{2.639063in}{1.821249in}}%
\pgfpathlineto{\pgfqpoint{2.639063in}{1.868438in}}%
\pgfpathlineto{\pgfqpoint{2.711719in}{1.868438in}}%
\pgfpathlineto{\pgfqpoint{2.711719in}{1.821249in}}%
\pgfpathmoveto{\pgfqpoint{2.639063in}{1.868438in}}%
\pgfpathlineto{\pgfqpoint{2.639063in}{1.868438in}}%
\pgfpathlineto{\pgfqpoint{2.639063in}{1.915627in}}%
\pgfpathlineto{\pgfqpoint{2.711719in}{1.915627in}}%
\pgfpathlineto{\pgfqpoint{2.711719in}{1.868438in}}%
\pgfpathmoveto{\pgfqpoint{2.711719in}{1.868438in}}%
\pgfpathlineto{\pgfqpoint{2.711719in}{1.868438in}}%
\pgfpathlineto{\pgfqpoint{2.711719in}{1.915627in}}%
\pgfpathlineto{\pgfqpoint{2.784375in}{1.915627in}}%
\pgfpathlineto{\pgfqpoint{2.784375in}{1.868438in}}%
\pgfpathmoveto{\pgfqpoint{2.784375in}{1.868438in}}%
\pgfpathlineto{\pgfqpoint{2.784375in}{1.868438in}}%
\pgfpathlineto{\pgfqpoint{2.784375in}{1.915627in}}%
\pgfpathlineto{\pgfqpoint{2.857031in}{1.915627in}}%
\pgfpathlineto{\pgfqpoint{2.857031in}{1.868438in}}%
\pgfpathmoveto{\pgfqpoint{2.784375in}{3.331252in}}%
\pgfpathlineto{\pgfqpoint{2.784375in}{3.331252in}}%
\pgfpathlineto{\pgfqpoint{2.784375in}{3.378439in}}%
\pgfpathlineto{\pgfqpoint{2.857031in}{3.378439in}}%
\pgfpathlineto{\pgfqpoint{2.857031in}{3.331252in}}%
\pgfpathmoveto{\pgfqpoint{2.784375in}{3.378439in}}%
\pgfpathlineto{\pgfqpoint{2.784375in}{3.378439in}}%
\pgfpathlineto{\pgfqpoint{2.784375in}{3.425627in}}%
\pgfpathlineto{\pgfqpoint{2.857031in}{3.425627in}}%
\pgfpathlineto{\pgfqpoint{2.857031in}{3.378439in}}%
\pgfpathmoveto{\pgfqpoint{2.857031in}{3.331252in}}%
\pgfpathlineto{\pgfqpoint{2.857031in}{3.331252in}}%
\pgfpathlineto{\pgfqpoint{2.857031in}{3.378439in}}%
\pgfpathlineto{\pgfqpoint{2.929688in}{3.378439in}}%
\pgfpathlineto{\pgfqpoint{2.929688in}{3.331252in}}%
\pgfpathmoveto{\pgfqpoint{2.784375in}{3.425627in}}%
\pgfpathlineto{\pgfqpoint{2.784375in}{3.425627in}}%
\pgfpathlineto{\pgfqpoint{2.784375in}{3.472812in}}%
\pgfpathlineto{\pgfqpoint{2.857031in}{3.472812in}}%
\pgfpathlineto{\pgfqpoint{2.857031in}{3.425627in}}%
\pgfpathmoveto{\pgfqpoint{2.929688in}{2.010000in}}%
\pgfpathlineto{\pgfqpoint{2.929688in}{2.010000in}}%
\pgfpathlineto{\pgfqpoint{2.929688in}{2.057189in}}%
\pgfpathlineto{\pgfqpoint{3.002346in}{2.057189in}}%
\pgfpathlineto{\pgfqpoint{3.002346in}{2.010000in}}%
\pgfpathmoveto{\pgfqpoint{2.929688in}{2.057189in}}%
\pgfpathlineto{\pgfqpoint{2.929688in}{2.057189in}}%
\pgfpathlineto{\pgfqpoint{2.929688in}{2.104377in}}%
\pgfpathlineto{\pgfqpoint{3.002346in}{2.104377in}}%
\pgfpathlineto{\pgfqpoint{3.002346in}{2.057189in}}%
\pgfpathmoveto{\pgfqpoint{3.002346in}{2.057189in}}%
\pgfpathlineto{\pgfqpoint{3.002346in}{2.057189in}}%
\pgfpathlineto{\pgfqpoint{3.002346in}{2.104377in}}%
\pgfpathlineto{\pgfqpoint{3.075004in}{2.104377in}}%
\pgfpathlineto{\pgfqpoint{3.075004in}{2.057189in}}%
\pgfpathmoveto{\pgfqpoint{2.929688in}{3.142502in}}%
\pgfpathlineto{\pgfqpoint{2.929688in}{3.142502in}}%
\pgfpathlineto{\pgfqpoint{2.929688in}{3.189688in}}%
\pgfpathlineto{\pgfqpoint{3.002346in}{3.189688in}}%
\pgfpathlineto{\pgfqpoint{3.002346in}{3.142502in}}%
\pgfpathmoveto{\pgfqpoint{2.929688in}{3.189688in}}%
\pgfpathlineto{\pgfqpoint{2.929688in}{3.189688in}}%
\pgfpathlineto{\pgfqpoint{2.929688in}{3.236874in}}%
\pgfpathlineto{\pgfqpoint{3.002346in}{3.236874in}}%
\pgfpathlineto{\pgfqpoint{3.002346in}{3.189688in}}%
\pgfpathmoveto{\pgfqpoint{3.002346in}{3.142502in}}%
\pgfpathlineto{\pgfqpoint{3.002346in}{3.142502in}}%
\pgfpathlineto{\pgfqpoint{3.002346in}{3.189688in}}%
\pgfpathlineto{\pgfqpoint{3.075004in}{3.189688in}}%
\pgfpathlineto{\pgfqpoint{3.075004in}{3.142502in}}%
\pgfpathmoveto{\pgfqpoint{2.929688in}{3.236874in}}%
\pgfpathlineto{\pgfqpoint{2.929688in}{3.236874in}}%
\pgfpathlineto{\pgfqpoint{2.929688in}{3.284063in}}%
\pgfpathlineto{\pgfqpoint{3.002346in}{3.284063in}}%
\pgfpathlineto{\pgfqpoint{3.002346in}{3.236874in}}%
\pgfpathmoveto{\pgfqpoint{3.075004in}{2.057189in}}%
\pgfpathlineto{\pgfqpoint{3.075004in}{2.057189in}}%
\pgfpathlineto{\pgfqpoint{3.075004in}{2.104377in}}%
\pgfpathlineto{\pgfqpoint{3.147657in}{2.104377in}}%
\pgfpathlineto{\pgfqpoint{3.147657in}{2.057189in}}%
\pgfpathmoveto{\pgfqpoint{3.075004in}{2.953751in}}%
\pgfpathlineto{\pgfqpoint{3.075004in}{2.953751in}}%
\pgfpathlineto{\pgfqpoint{3.075004in}{3.000938in}}%
\pgfpathlineto{\pgfqpoint{3.147657in}{3.000938in}}%
\pgfpathlineto{\pgfqpoint{3.147657in}{2.953751in}}%
\pgfpathmoveto{\pgfqpoint{3.075004in}{3.000938in}}%
\pgfpathlineto{\pgfqpoint{3.075004in}{3.000938in}}%
\pgfpathlineto{\pgfqpoint{3.075004in}{3.048124in}}%
\pgfpathlineto{\pgfqpoint{3.147657in}{3.048124in}}%
\pgfpathlineto{\pgfqpoint{3.147657in}{3.000938in}}%
\pgfpathmoveto{\pgfqpoint{3.147657in}{2.953751in}}%
\pgfpathlineto{\pgfqpoint{3.147657in}{2.953751in}}%
\pgfpathlineto{\pgfqpoint{3.147657in}{3.000938in}}%
\pgfpathlineto{\pgfqpoint{3.220310in}{3.000938in}}%
\pgfpathlineto{\pgfqpoint{3.220310in}{2.953751in}}%
\pgfpathmoveto{\pgfqpoint{3.075004in}{3.048124in}}%
\pgfpathlineto{\pgfqpoint{3.075004in}{3.048124in}}%
\pgfpathlineto{\pgfqpoint{3.075004in}{3.095313in}}%
\pgfpathlineto{\pgfqpoint{3.147657in}{3.095313in}}%
\pgfpathlineto{\pgfqpoint{3.147657in}{3.048124in}}%
\pgfpathmoveto{\pgfqpoint{3.220310in}{2.151563in}}%
\pgfpathlineto{\pgfqpoint{3.220310in}{2.151563in}}%
\pgfpathlineto{\pgfqpoint{3.220310in}{2.198748in}}%
\pgfpathlineto{\pgfqpoint{3.292968in}{2.198748in}}%
\pgfpathlineto{\pgfqpoint{3.292968in}{2.151563in}}%
\pgfpathmoveto{\pgfqpoint{3.220310in}{2.198748in}}%
\pgfpathlineto{\pgfqpoint{3.220310in}{2.198748in}}%
\pgfpathlineto{\pgfqpoint{3.220310in}{2.245935in}}%
\pgfpathlineto{\pgfqpoint{3.292968in}{2.245935in}}%
\pgfpathlineto{\pgfqpoint{3.292968in}{2.198748in}}%
\pgfpathmoveto{\pgfqpoint{3.220310in}{2.245935in}}%
\pgfpathlineto{\pgfqpoint{3.220310in}{2.245935in}}%
\pgfpathlineto{\pgfqpoint{3.220310in}{2.293122in}}%
\pgfpathlineto{\pgfqpoint{3.292968in}{2.293122in}}%
\pgfpathlineto{\pgfqpoint{3.292968in}{2.245935in}}%
\pgfpathmoveto{\pgfqpoint{3.292968in}{2.245935in}}%
\pgfpathlineto{\pgfqpoint{3.292968in}{2.245935in}}%
\pgfpathlineto{\pgfqpoint{3.292968in}{2.293122in}}%
\pgfpathlineto{\pgfqpoint{3.365627in}{2.293122in}}%
\pgfpathlineto{\pgfqpoint{3.365627in}{2.245935in}}%
\pgfpathmoveto{\pgfqpoint{3.220310in}{2.764999in}}%
\pgfpathlineto{\pgfqpoint{3.220310in}{2.764999in}}%
\pgfpathlineto{\pgfqpoint{3.220310in}{2.812187in}}%
\pgfpathlineto{\pgfqpoint{3.292968in}{2.812187in}}%
\pgfpathlineto{\pgfqpoint{3.292968in}{2.764999in}}%
\pgfpathmoveto{\pgfqpoint{3.220310in}{2.812187in}}%
\pgfpathlineto{\pgfqpoint{3.220310in}{2.812187in}}%
\pgfpathlineto{\pgfqpoint{3.220310in}{2.859374in}}%
\pgfpathlineto{\pgfqpoint{3.292968in}{2.859374in}}%
\pgfpathlineto{\pgfqpoint{3.292968in}{2.812187in}}%
\pgfpathmoveto{\pgfqpoint{3.292968in}{2.764999in}}%
\pgfpathlineto{\pgfqpoint{3.292968in}{2.764999in}}%
\pgfpathlineto{\pgfqpoint{3.292968in}{2.812187in}}%
\pgfpathlineto{\pgfqpoint{3.365627in}{2.812187in}}%
\pgfpathlineto{\pgfqpoint{3.365627in}{2.764999in}}%
\pgfpathmoveto{\pgfqpoint{3.220310in}{2.859374in}}%
\pgfpathlineto{\pgfqpoint{3.220310in}{2.859374in}}%
\pgfpathlineto{\pgfqpoint{3.220310in}{2.906563in}}%
\pgfpathlineto{\pgfqpoint{3.292968in}{2.906563in}}%
\pgfpathlineto{\pgfqpoint{3.292968in}{2.859374in}}%
\pgfpathmoveto{\pgfqpoint{3.365627in}{2.293122in}}%
\pgfpathlineto{\pgfqpoint{3.365627in}{2.293122in}}%
\pgfpathlineto{\pgfqpoint{3.365627in}{2.340311in}}%
\pgfpathlineto{\pgfqpoint{3.438283in}{2.340311in}}%
\pgfpathlineto{\pgfqpoint{3.438283in}{2.293122in}}%
\pgfpathmoveto{\pgfqpoint{3.365627in}{2.340311in}}%
\pgfpathlineto{\pgfqpoint{3.365627in}{2.340311in}}%
\pgfpathlineto{\pgfqpoint{3.365627in}{2.387499in}}%
\pgfpathlineto{\pgfqpoint{3.438283in}{2.387499in}}%
\pgfpathlineto{\pgfqpoint{3.438283in}{2.340311in}}%
\pgfpathmoveto{\pgfqpoint{3.438283in}{2.340311in}}%
\pgfpathlineto{\pgfqpoint{3.438283in}{2.340311in}}%
\pgfpathlineto{\pgfqpoint{3.438283in}{2.387499in}}%
\pgfpathlineto{\pgfqpoint{3.510939in}{2.387499in}}%
\pgfpathlineto{\pgfqpoint{3.510939in}{2.340311in}}%
\pgfpathmoveto{\pgfqpoint{3.365627in}{2.576251in}}%
\pgfpathlineto{\pgfqpoint{3.365627in}{2.576251in}}%
\pgfpathlineto{\pgfqpoint{3.365627in}{2.623437in}}%
\pgfpathlineto{\pgfqpoint{3.438283in}{2.623437in}}%
\pgfpathlineto{\pgfqpoint{3.438283in}{2.576251in}}%
\pgfpathmoveto{\pgfqpoint{3.365627in}{2.623437in}}%
\pgfpathlineto{\pgfqpoint{3.365627in}{2.623437in}}%
\pgfpathlineto{\pgfqpoint{3.365627in}{2.670624in}}%
\pgfpathlineto{\pgfqpoint{3.438283in}{2.670624in}}%
\pgfpathlineto{\pgfqpoint{3.438283in}{2.623437in}}%
\pgfpathmoveto{\pgfqpoint{3.438283in}{2.576251in}}%
\pgfpathlineto{\pgfqpoint{3.438283in}{2.576251in}}%
\pgfpathlineto{\pgfqpoint{3.438283in}{2.623437in}}%
\pgfpathlineto{\pgfqpoint{3.510939in}{2.623437in}}%
\pgfpathlineto{\pgfqpoint{3.510939in}{2.576251in}}%
\pgfpathmoveto{\pgfqpoint{3.365627in}{2.670624in}}%
\pgfpathlineto{\pgfqpoint{3.365627in}{2.670624in}}%
\pgfpathlineto{\pgfqpoint{3.365627in}{2.717812in}}%
\pgfpathlineto{\pgfqpoint{3.438283in}{2.717812in}}%
\pgfpathlineto{\pgfqpoint{3.438283in}{2.670624in}}%
\pgfpathmoveto{\pgfqpoint{3.510939in}{2.387499in}}%
\pgfpathlineto{\pgfqpoint{3.510939in}{2.387499in}}%
\pgfpathlineto{\pgfqpoint{3.510939in}{2.434688in}}%
\pgfpathlineto{\pgfqpoint{3.583592in}{2.434688in}}%
\pgfpathlineto{\pgfqpoint{3.583592in}{2.387499in}}%
\pgfpathmoveto{\pgfqpoint{3.510939in}{2.434688in}}%
\pgfpathlineto{\pgfqpoint{3.510939in}{2.434688in}}%
\pgfpathlineto{\pgfqpoint{3.510939in}{2.481878in}}%
\pgfpathlineto{\pgfqpoint{3.583592in}{2.481878in}}%
\pgfpathlineto{\pgfqpoint{3.583592in}{2.434688in}}%
\pgfpathmoveto{\pgfqpoint{3.583592in}{2.387499in}}%
\pgfpathlineto{\pgfqpoint{3.583592in}{2.387499in}}%
\pgfpathlineto{\pgfqpoint{3.583592in}{2.434688in}}%
\pgfpathlineto{\pgfqpoint{3.656246in}{2.434688in}}%
\pgfpathlineto{\pgfqpoint{3.656246in}{2.387499in}}%
\pgfpathmoveto{\pgfqpoint{3.510939in}{2.481878in}}%
\pgfpathlineto{\pgfqpoint{3.510939in}{2.481878in}}%
\pgfpathlineto{\pgfqpoint{3.510939in}{2.529064in}}%
\pgfpathlineto{\pgfqpoint{3.583592in}{2.529064in}}%
\pgfpathlineto{\pgfqpoint{3.583592in}{2.481878in}}%
\pgfpathmoveto{\pgfqpoint{0.749998in}{0.547186in}}%
\pgfpathlineto{\pgfqpoint{0.749998in}{0.547186in}}%
\pgfpathlineto{\pgfqpoint{0.749998in}{0.570780in}}%
\pgfpathlineto{\pgfqpoint{0.786327in}{0.570780in}}%
\pgfpathlineto{\pgfqpoint{0.786327in}{0.547186in}}%
\pgfpathmoveto{\pgfqpoint{0.749998in}{0.570780in}}%
\pgfpathlineto{\pgfqpoint{0.749998in}{0.570780in}}%
\pgfpathlineto{\pgfqpoint{0.749998in}{0.594374in}}%
\pgfpathlineto{\pgfqpoint{0.786327in}{0.594374in}}%
\pgfpathlineto{\pgfqpoint{0.786327in}{0.570780in}}%
\pgfpathmoveto{\pgfqpoint{0.786327in}{0.570780in}}%
\pgfpathlineto{\pgfqpoint{0.786327in}{0.570780in}}%
\pgfpathlineto{\pgfqpoint{0.786327in}{0.594374in}}%
\pgfpathlineto{\pgfqpoint{0.822655in}{0.594374in}}%
\pgfpathlineto{\pgfqpoint{0.822655in}{0.570780in}}%
\pgfpathmoveto{\pgfqpoint{0.822655in}{0.594374in}}%
\pgfpathlineto{\pgfqpoint{0.822655in}{0.594374in}}%
\pgfpathlineto{\pgfqpoint{0.822655in}{0.617968in}}%
\pgfpathlineto{\pgfqpoint{0.858983in}{0.617968in}}%
\pgfpathlineto{\pgfqpoint{0.858983in}{0.594374in}}%
\pgfpathmoveto{\pgfqpoint{0.822655in}{0.617968in}}%
\pgfpathlineto{\pgfqpoint{0.822655in}{0.617968in}}%
\pgfpathlineto{\pgfqpoint{0.822655in}{0.641562in}}%
\pgfpathlineto{\pgfqpoint{0.858983in}{0.641562in}}%
\pgfpathlineto{\pgfqpoint{0.858983in}{0.617968in}}%
\pgfpathmoveto{\pgfqpoint{0.858983in}{0.617968in}}%
\pgfpathlineto{\pgfqpoint{0.858983in}{0.617968in}}%
\pgfpathlineto{\pgfqpoint{0.858983in}{0.641562in}}%
\pgfpathlineto{\pgfqpoint{0.895312in}{0.641562in}}%
\pgfpathlineto{\pgfqpoint{0.895312in}{0.617968in}}%
\pgfpathmoveto{\pgfqpoint{0.895312in}{0.617968in}}%
\pgfpathlineto{\pgfqpoint{0.895312in}{0.617968in}}%
\pgfpathlineto{\pgfqpoint{0.895312in}{0.641562in}}%
\pgfpathlineto{\pgfqpoint{0.931639in}{0.641562in}}%
\pgfpathlineto{\pgfqpoint{0.931639in}{0.617968in}}%
\pgfpathmoveto{\pgfqpoint{0.967967in}{0.665156in}}%
\pgfpathlineto{\pgfqpoint{0.967967in}{0.665156in}}%
\pgfpathlineto{\pgfqpoint{0.967967in}{0.688750in}}%
\pgfpathlineto{\pgfqpoint{1.004294in}{0.688750in}}%
\pgfpathlineto{\pgfqpoint{1.004294in}{0.665156in}}%
\pgfpathmoveto{\pgfqpoint{1.040622in}{0.712343in}}%
\pgfpathlineto{\pgfqpoint{1.040622in}{0.712343in}}%
\pgfpathlineto{\pgfqpoint{1.040622in}{0.735937in}}%
\pgfpathlineto{\pgfqpoint{1.076952in}{0.735937in}}%
\pgfpathlineto{\pgfqpoint{1.076952in}{0.712343in}}%
\pgfpathmoveto{\pgfqpoint{1.040622in}{0.735937in}}%
\pgfpathlineto{\pgfqpoint{1.040622in}{0.735937in}}%
\pgfpathlineto{\pgfqpoint{1.040622in}{0.759530in}}%
\pgfpathlineto{\pgfqpoint{1.076952in}{0.759530in}}%
\pgfpathlineto{\pgfqpoint{1.076952in}{0.735937in}}%
\pgfpathmoveto{\pgfqpoint{1.040622in}{0.759530in}}%
\pgfpathlineto{\pgfqpoint{1.040622in}{0.759530in}}%
\pgfpathlineto{\pgfqpoint{1.040622in}{0.783124in}}%
\pgfpathlineto{\pgfqpoint{1.076952in}{0.783124in}}%
\pgfpathlineto{\pgfqpoint{1.076952in}{0.759530in}}%
\pgfpathmoveto{\pgfqpoint{1.076952in}{0.759530in}}%
\pgfpathlineto{\pgfqpoint{1.076952in}{0.759530in}}%
\pgfpathlineto{\pgfqpoint{1.076952in}{0.783124in}}%
\pgfpathlineto{\pgfqpoint{1.113282in}{0.783124in}}%
\pgfpathlineto{\pgfqpoint{1.113282in}{0.759530in}}%
\pgfpathmoveto{\pgfqpoint{1.113282in}{0.783124in}}%
\pgfpathlineto{\pgfqpoint{1.113282in}{0.783124in}}%
\pgfpathlineto{\pgfqpoint{1.113282in}{0.806718in}}%
\pgfpathlineto{\pgfqpoint{1.149611in}{0.806718in}}%
\pgfpathlineto{\pgfqpoint{1.149611in}{0.783124in}}%
\pgfpathmoveto{\pgfqpoint{1.113282in}{0.806718in}}%
\pgfpathlineto{\pgfqpoint{1.113282in}{0.806718in}}%
\pgfpathlineto{\pgfqpoint{1.113282in}{0.830313in}}%
\pgfpathlineto{\pgfqpoint{1.149611in}{0.830313in}}%
\pgfpathlineto{\pgfqpoint{1.149611in}{0.806718in}}%
\pgfpathmoveto{\pgfqpoint{1.149611in}{0.806718in}}%
\pgfpathlineto{\pgfqpoint{1.149611in}{0.806718in}}%
\pgfpathlineto{\pgfqpoint{1.149611in}{0.830313in}}%
\pgfpathlineto{\pgfqpoint{1.185941in}{0.830313in}}%
\pgfpathlineto{\pgfqpoint{1.185941in}{0.806718in}}%
\pgfpathmoveto{\pgfqpoint{1.185941in}{0.830313in}}%
\pgfpathlineto{\pgfqpoint{1.185941in}{0.830313in}}%
\pgfpathlineto{\pgfqpoint{1.185941in}{0.853907in}}%
\pgfpathlineto{\pgfqpoint{1.222269in}{0.853907in}}%
\pgfpathlineto{\pgfqpoint{1.222269in}{0.830313in}}%
\pgfpathmoveto{\pgfqpoint{1.185941in}{0.853907in}}%
\pgfpathlineto{\pgfqpoint{1.185941in}{0.853907in}}%
\pgfpathlineto{\pgfqpoint{1.185941in}{0.877501in}}%
\pgfpathlineto{\pgfqpoint{1.222269in}{0.877501in}}%
\pgfpathlineto{\pgfqpoint{1.222269in}{0.853907in}}%
\pgfpathmoveto{\pgfqpoint{1.222269in}{0.853907in}}%
\pgfpathlineto{\pgfqpoint{1.222269in}{0.853907in}}%
\pgfpathlineto{\pgfqpoint{1.222269in}{0.877501in}}%
\pgfpathlineto{\pgfqpoint{1.258598in}{0.877501in}}%
\pgfpathlineto{\pgfqpoint{1.258598in}{0.853907in}}%
\pgfpathmoveto{\pgfqpoint{1.258598in}{0.877501in}}%
\pgfpathlineto{\pgfqpoint{1.258598in}{0.877501in}}%
\pgfpathlineto{\pgfqpoint{1.258598in}{0.901094in}}%
\pgfpathlineto{\pgfqpoint{1.294926in}{0.901094in}}%
\pgfpathlineto{\pgfqpoint{1.294926in}{0.877501in}}%
\pgfpathmoveto{\pgfqpoint{1.258598in}{0.901094in}}%
\pgfpathlineto{\pgfqpoint{1.258598in}{0.901094in}}%
\pgfpathlineto{\pgfqpoint{1.258598in}{0.924687in}}%
\pgfpathlineto{\pgfqpoint{1.294926in}{0.924687in}}%
\pgfpathlineto{\pgfqpoint{1.294926in}{0.901094in}}%
\pgfpathmoveto{\pgfqpoint{1.294926in}{0.901094in}}%
\pgfpathlineto{\pgfqpoint{1.294926in}{0.901094in}}%
\pgfpathlineto{\pgfqpoint{1.294926in}{0.924687in}}%
\pgfpathlineto{\pgfqpoint{1.331254in}{0.924687in}}%
\pgfpathlineto{\pgfqpoint{1.331254in}{0.901094in}}%
\pgfpathmoveto{\pgfqpoint{1.331254in}{0.924687in}}%
\pgfpathlineto{\pgfqpoint{1.331254in}{0.924687in}}%
\pgfpathlineto{\pgfqpoint{1.331254in}{0.948280in}}%
\pgfpathlineto{\pgfqpoint{1.367581in}{0.948280in}}%
\pgfpathlineto{\pgfqpoint{1.367581in}{0.924687in}}%
\pgfpathmoveto{\pgfqpoint{1.331254in}{0.948280in}}%
\pgfpathlineto{\pgfqpoint{1.331254in}{0.948280in}}%
\pgfpathlineto{\pgfqpoint{1.331254in}{0.971872in}}%
\pgfpathlineto{\pgfqpoint{1.367581in}{0.971872in}}%
\pgfpathlineto{\pgfqpoint{1.367581in}{0.948280in}}%
\pgfpathmoveto{\pgfqpoint{1.367581in}{0.948280in}}%
\pgfpathlineto{\pgfqpoint{1.367581in}{0.948280in}}%
\pgfpathlineto{\pgfqpoint{1.367581in}{0.971872in}}%
\pgfpathlineto{\pgfqpoint{1.403909in}{0.971872in}}%
\pgfpathlineto{\pgfqpoint{1.403909in}{0.948280in}}%
\pgfpathmoveto{\pgfqpoint{1.403909in}{0.971872in}}%
\pgfpathlineto{\pgfqpoint{1.403909in}{0.971872in}}%
\pgfpathlineto{\pgfqpoint{1.403909in}{0.995466in}}%
\pgfpathlineto{\pgfqpoint{1.440237in}{0.995466in}}%
\pgfpathlineto{\pgfqpoint{1.440237in}{0.971872in}}%
\pgfpathmoveto{\pgfqpoint{1.403909in}{0.995466in}}%
\pgfpathlineto{\pgfqpoint{1.403909in}{0.995466in}}%
\pgfpathlineto{\pgfqpoint{1.403909in}{1.019061in}}%
\pgfpathlineto{\pgfqpoint{1.440237in}{1.019061in}}%
\pgfpathlineto{\pgfqpoint{1.440237in}{0.995466in}}%
\pgfpathmoveto{\pgfqpoint{1.440237in}{0.995466in}}%
\pgfpathlineto{\pgfqpoint{1.440237in}{0.995466in}}%
\pgfpathlineto{\pgfqpoint{1.440237in}{1.019061in}}%
\pgfpathlineto{\pgfqpoint{1.476564in}{1.019061in}}%
\pgfpathlineto{\pgfqpoint{1.476564in}{0.995466in}}%
\pgfpathmoveto{\pgfqpoint{1.476564in}{1.019061in}}%
\pgfpathlineto{\pgfqpoint{1.476564in}{1.019061in}}%
\pgfpathlineto{\pgfqpoint{1.476564in}{1.042655in}}%
\pgfpathlineto{\pgfqpoint{1.512892in}{1.042655in}}%
\pgfpathlineto{\pgfqpoint{1.512892in}{1.019061in}}%
\pgfpathmoveto{\pgfqpoint{1.476564in}{1.042655in}}%
\pgfpathlineto{\pgfqpoint{1.476564in}{1.042655in}}%
\pgfpathlineto{\pgfqpoint{1.476564in}{1.066249in}}%
\pgfpathlineto{\pgfqpoint{1.512892in}{1.066249in}}%
\pgfpathlineto{\pgfqpoint{1.512892in}{1.042655in}}%
\pgfpathmoveto{\pgfqpoint{1.512892in}{1.042655in}}%
\pgfpathlineto{\pgfqpoint{1.512892in}{1.042655in}}%
\pgfpathlineto{\pgfqpoint{1.512892in}{1.066249in}}%
\pgfpathlineto{\pgfqpoint{1.549220in}{1.066249in}}%
\pgfpathlineto{\pgfqpoint{1.549220in}{1.042655in}}%
\pgfpathmoveto{\pgfqpoint{1.549220in}{1.066249in}}%
\pgfpathlineto{\pgfqpoint{1.549220in}{1.066249in}}%
\pgfpathlineto{\pgfqpoint{1.549220in}{1.089843in}}%
\pgfpathlineto{\pgfqpoint{1.585549in}{1.089843in}}%
\pgfpathlineto{\pgfqpoint{1.585549in}{1.066249in}}%
\pgfpathmoveto{\pgfqpoint{1.549220in}{1.089843in}}%
\pgfpathlineto{\pgfqpoint{1.549220in}{1.089843in}}%
\pgfpathlineto{\pgfqpoint{1.549220in}{1.113437in}}%
\pgfpathlineto{\pgfqpoint{1.585549in}{1.113437in}}%
\pgfpathlineto{\pgfqpoint{1.585549in}{1.089843in}}%
\pgfpathmoveto{\pgfqpoint{1.585549in}{1.089843in}}%
\pgfpathlineto{\pgfqpoint{1.585549in}{1.089843in}}%
\pgfpathlineto{\pgfqpoint{1.585549in}{1.113437in}}%
\pgfpathlineto{\pgfqpoint{1.621877in}{1.113437in}}%
\pgfpathlineto{\pgfqpoint{1.621877in}{1.089843in}}%
\pgfpathmoveto{\pgfqpoint{1.621877in}{1.113437in}}%
\pgfpathlineto{\pgfqpoint{1.621877in}{1.113437in}}%
\pgfpathlineto{\pgfqpoint{1.621877in}{1.137031in}}%
\pgfpathlineto{\pgfqpoint{1.658204in}{1.137031in}}%
\pgfpathlineto{\pgfqpoint{1.658204in}{1.113437in}}%
\pgfpathmoveto{\pgfqpoint{1.621877in}{1.137031in}}%
\pgfpathlineto{\pgfqpoint{1.621877in}{1.137031in}}%
\pgfpathlineto{\pgfqpoint{1.621877in}{1.160625in}}%
\pgfpathlineto{\pgfqpoint{1.658204in}{1.160625in}}%
\pgfpathlineto{\pgfqpoint{1.658204in}{1.137031in}}%
\pgfpathmoveto{\pgfqpoint{1.658204in}{1.137031in}}%
\pgfpathlineto{\pgfqpoint{1.658204in}{1.137031in}}%
\pgfpathlineto{\pgfqpoint{1.658204in}{1.160625in}}%
\pgfpathlineto{\pgfqpoint{1.694531in}{1.160625in}}%
\pgfpathlineto{\pgfqpoint{1.694531in}{1.137031in}}%
\pgfpathmoveto{\pgfqpoint{1.694531in}{1.137031in}}%
\pgfpathlineto{\pgfqpoint{1.694531in}{1.137031in}}%
\pgfpathlineto{\pgfqpoint{1.694531in}{1.160625in}}%
\pgfpathlineto{\pgfqpoint{1.730859in}{1.160625in}}%
\pgfpathlineto{\pgfqpoint{1.730859in}{1.137031in}}%
\pgfpathmoveto{\pgfqpoint{1.767186in}{1.184218in}}%
\pgfpathlineto{\pgfqpoint{1.767186in}{1.184218in}}%
\pgfpathlineto{\pgfqpoint{1.767186in}{1.207812in}}%
\pgfpathlineto{\pgfqpoint{1.803513in}{1.207812in}}%
\pgfpathlineto{\pgfqpoint{1.803513in}{1.184218in}}%
\pgfpathmoveto{\pgfqpoint{1.839841in}{1.231405in}}%
\pgfpathlineto{\pgfqpoint{1.839841in}{1.231405in}}%
\pgfpathlineto{\pgfqpoint{1.839841in}{1.254998in}}%
\pgfpathlineto{\pgfqpoint{1.876168in}{1.254998in}}%
\pgfpathlineto{\pgfqpoint{1.876168in}{1.231405in}}%
\pgfpathmoveto{\pgfqpoint{1.912496in}{1.278592in}}%
\pgfpathlineto{\pgfqpoint{1.912496in}{1.278592in}}%
\pgfpathlineto{\pgfqpoint{1.912496in}{1.302187in}}%
\pgfpathlineto{\pgfqpoint{1.948825in}{1.302187in}}%
\pgfpathlineto{\pgfqpoint{1.948825in}{1.278592in}}%
\pgfpathmoveto{\pgfqpoint{1.985155in}{1.325781in}}%
\pgfpathlineto{\pgfqpoint{1.985155in}{1.325781in}}%
\pgfpathlineto{\pgfqpoint{1.985155in}{1.349376in}}%
\pgfpathlineto{\pgfqpoint{2.021484in}{1.349376in}}%
\pgfpathlineto{\pgfqpoint{2.021484in}{1.325781in}}%
\pgfpathmoveto{\pgfqpoint{2.057813in}{1.396563in}}%
\pgfpathlineto{\pgfqpoint{2.057813in}{1.396563in}}%
\pgfpathlineto{\pgfqpoint{2.057813in}{1.420156in}}%
\pgfpathlineto{\pgfqpoint{2.094142in}{1.420156in}}%
\pgfpathlineto{\pgfqpoint{2.094142in}{1.396563in}}%
\pgfpathmoveto{\pgfqpoint{2.057813in}{1.420156in}}%
\pgfpathlineto{\pgfqpoint{2.057813in}{1.420156in}}%
\pgfpathlineto{\pgfqpoint{2.057813in}{1.443750in}}%
\pgfpathlineto{\pgfqpoint{2.094142in}{1.443750in}}%
\pgfpathlineto{\pgfqpoint{2.094142in}{1.420156in}}%
\pgfpathmoveto{\pgfqpoint{2.094142in}{1.420156in}}%
\pgfpathlineto{\pgfqpoint{2.094142in}{1.420156in}}%
\pgfpathlineto{\pgfqpoint{2.094142in}{1.443750in}}%
\pgfpathlineto{\pgfqpoint{2.130471in}{1.443750in}}%
\pgfpathlineto{\pgfqpoint{2.130471in}{1.420156in}}%
\pgfpathmoveto{\pgfqpoint{2.130471in}{1.443750in}}%
\pgfpathlineto{\pgfqpoint{2.130471in}{1.443750in}}%
\pgfpathlineto{\pgfqpoint{2.130471in}{1.467344in}}%
\pgfpathlineto{\pgfqpoint{2.166800in}{1.467344in}}%
\pgfpathlineto{\pgfqpoint{2.166800in}{1.443750in}}%
\pgfpathmoveto{\pgfqpoint{2.130471in}{1.467344in}}%
\pgfpathlineto{\pgfqpoint{2.130471in}{1.467344in}}%
\pgfpathlineto{\pgfqpoint{2.130471in}{1.490938in}}%
\pgfpathlineto{\pgfqpoint{2.166800in}{1.490938in}}%
\pgfpathlineto{\pgfqpoint{2.166800in}{1.467344in}}%
\pgfpathmoveto{\pgfqpoint{2.166800in}{1.467344in}}%
\pgfpathlineto{\pgfqpoint{2.166800in}{1.467344in}}%
\pgfpathlineto{\pgfqpoint{2.166800in}{1.490938in}}%
\pgfpathlineto{\pgfqpoint{2.203129in}{1.490938in}}%
\pgfpathlineto{\pgfqpoint{2.203129in}{1.467344in}}%
\pgfpathmoveto{\pgfqpoint{2.203129in}{1.490938in}}%
\pgfpathlineto{\pgfqpoint{2.203129in}{1.490938in}}%
\pgfpathlineto{\pgfqpoint{2.203129in}{1.514533in}}%
\pgfpathlineto{\pgfqpoint{2.239457in}{1.514533in}}%
\pgfpathlineto{\pgfqpoint{2.239457in}{1.490938in}}%
\pgfpathmoveto{\pgfqpoint{2.203129in}{1.514533in}}%
\pgfpathlineto{\pgfqpoint{2.203129in}{1.514533in}}%
\pgfpathlineto{\pgfqpoint{2.203129in}{1.538127in}}%
\pgfpathlineto{\pgfqpoint{2.239457in}{1.538127in}}%
\pgfpathlineto{\pgfqpoint{2.239457in}{1.514533in}}%
\pgfpathmoveto{\pgfqpoint{2.239457in}{1.514533in}}%
\pgfpathlineto{\pgfqpoint{2.239457in}{1.514533in}}%
\pgfpathlineto{\pgfqpoint{2.239457in}{1.538127in}}%
\pgfpathlineto{\pgfqpoint{2.275785in}{1.538127in}}%
\pgfpathlineto{\pgfqpoint{2.275785in}{1.514533in}}%
\pgfpathmoveto{\pgfqpoint{2.275785in}{1.538127in}}%
\pgfpathlineto{\pgfqpoint{2.275785in}{1.538127in}}%
\pgfpathlineto{\pgfqpoint{2.275785in}{1.561720in}}%
\pgfpathlineto{\pgfqpoint{2.312113in}{1.561720in}}%
\pgfpathlineto{\pgfqpoint{2.312113in}{1.538127in}}%
\pgfpathmoveto{\pgfqpoint{2.275785in}{1.561720in}}%
\pgfpathlineto{\pgfqpoint{2.275785in}{1.561720in}}%
\pgfpathlineto{\pgfqpoint{2.275785in}{1.585314in}}%
\pgfpathlineto{\pgfqpoint{2.312113in}{1.585314in}}%
\pgfpathlineto{\pgfqpoint{2.312113in}{1.561720in}}%
\pgfpathmoveto{\pgfqpoint{2.312113in}{1.561720in}}%
\pgfpathlineto{\pgfqpoint{2.312113in}{1.561720in}}%
\pgfpathlineto{\pgfqpoint{2.312113in}{1.585314in}}%
\pgfpathlineto{\pgfqpoint{2.348441in}{1.585314in}}%
\pgfpathlineto{\pgfqpoint{2.348441in}{1.561720in}}%
\pgfpathmoveto{\pgfqpoint{2.348441in}{1.561720in}}%
\pgfpathlineto{\pgfqpoint{2.348441in}{1.561720in}}%
\pgfpathlineto{\pgfqpoint{2.348441in}{1.585314in}}%
\pgfpathlineto{\pgfqpoint{2.384768in}{1.585314in}}%
\pgfpathlineto{\pgfqpoint{2.384768in}{1.561720in}}%
\pgfpathmoveto{\pgfqpoint{2.421095in}{1.608908in}}%
\pgfpathlineto{\pgfqpoint{2.421095in}{1.608908in}}%
\pgfpathlineto{\pgfqpoint{2.421095in}{1.632501in}}%
\pgfpathlineto{\pgfqpoint{2.457422in}{1.632501in}}%
\pgfpathlineto{\pgfqpoint{2.457422in}{1.608908in}}%
\pgfpathmoveto{\pgfqpoint{2.493750in}{1.656095in}}%
\pgfpathlineto{\pgfqpoint{2.493750in}{1.656095in}}%
\pgfpathlineto{\pgfqpoint{2.493750in}{1.679689in}}%
\pgfpathlineto{\pgfqpoint{2.530078in}{1.679689in}}%
\pgfpathlineto{\pgfqpoint{2.530078in}{1.656095in}}%
\pgfpathmoveto{\pgfqpoint{2.566406in}{1.703282in}}%
\pgfpathlineto{\pgfqpoint{2.566406in}{1.703282in}}%
\pgfpathlineto{\pgfqpoint{2.566406in}{1.726876in}}%
\pgfpathlineto{\pgfqpoint{2.602735in}{1.726876in}}%
\pgfpathlineto{\pgfqpoint{2.602735in}{1.703282in}}%
\pgfpathmoveto{\pgfqpoint{2.639063in}{1.774062in}}%
\pgfpathlineto{\pgfqpoint{2.639063in}{1.774062in}}%
\pgfpathlineto{\pgfqpoint{2.639063in}{1.797655in}}%
\pgfpathlineto{\pgfqpoint{2.675391in}{1.797655in}}%
\pgfpathlineto{\pgfqpoint{2.675391in}{1.774062in}}%
\pgfpathmoveto{\pgfqpoint{2.639063in}{1.797655in}}%
\pgfpathlineto{\pgfqpoint{2.639063in}{1.797655in}}%
\pgfpathlineto{\pgfqpoint{2.639063in}{1.821249in}}%
\pgfpathlineto{\pgfqpoint{2.675391in}{1.821249in}}%
\pgfpathlineto{\pgfqpoint{2.675391in}{1.797655in}}%
\pgfpathmoveto{\pgfqpoint{2.675391in}{1.797655in}}%
\pgfpathlineto{\pgfqpoint{2.675391in}{1.797655in}}%
\pgfpathlineto{\pgfqpoint{2.675391in}{1.821249in}}%
\pgfpathlineto{\pgfqpoint{2.711719in}{1.821249in}}%
\pgfpathlineto{\pgfqpoint{2.711719in}{1.797655in}}%
\pgfpathmoveto{\pgfqpoint{2.711719in}{1.821249in}}%
\pgfpathlineto{\pgfqpoint{2.711719in}{1.821249in}}%
\pgfpathlineto{\pgfqpoint{2.711719in}{1.844843in}}%
\pgfpathlineto{\pgfqpoint{2.748047in}{1.844843in}}%
\pgfpathlineto{\pgfqpoint{2.748047in}{1.821249in}}%
\pgfpathmoveto{\pgfqpoint{2.711719in}{1.844843in}}%
\pgfpathlineto{\pgfqpoint{2.711719in}{1.844843in}}%
\pgfpathlineto{\pgfqpoint{2.711719in}{1.868438in}}%
\pgfpathlineto{\pgfqpoint{2.748047in}{1.868438in}}%
\pgfpathlineto{\pgfqpoint{2.748047in}{1.844843in}}%
\pgfpathmoveto{\pgfqpoint{2.748047in}{1.844843in}}%
\pgfpathlineto{\pgfqpoint{2.748047in}{1.844843in}}%
\pgfpathlineto{\pgfqpoint{2.748047in}{1.868438in}}%
\pgfpathlineto{\pgfqpoint{2.784375in}{1.868438in}}%
\pgfpathlineto{\pgfqpoint{2.784375in}{1.844843in}}%
\pgfpathmoveto{\pgfqpoint{2.857031in}{1.892032in}}%
\pgfpathlineto{\pgfqpoint{2.857031in}{1.892032in}}%
\pgfpathlineto{\pgfqpoint{2.857031in}{1.915627in}}%
\pgfpathlineto{\pgfqpoint{2.893360in}{1.915627in}}%
\pgfpathlineto{\pgfqpoint{2.893360in}{1.892032in}}%
\pgfpathmoveto{\pgfqpoint{2.857031in}{3.378439in}}%
\pgfpathlineto{\pgfqpoint{2.857031in}{3.378439in}}%
\pgfpathlineto{\pgfqpoint{2.857031in}{3.402033in}}%
\pgfpathlineto{\pgfqpoint{2.893360in}{3.402033in}}%
\pgfpathlineto{\pgfqpoint{2.893360in}{3.378439in}}%
\pgfpathmoveto{\pgfqpoint{2.857031in}{3.402033in}}%
\pgfpathlineto{\pgfqpoint{2.857031in}{3.402033in}}%
\pgfpathlineto{\pgfqpoint{2.857031in}{3.425627in}}%
\pgfpathlineto{\pgfqpoint{2.893360in}{3.425627in}}%
\pgfpathlineto{\pgfqpoint{2.893360in}{3.402033in}}%
\pgfpathmoveto{\pgfqpoint{2.893360in}{3.378439in}}%
\pgfpathlineto{\pgfqpoint{2.893360in}{3.378439in}}%
\pgfpathlineto{\pgfqpoint{2.893360in}{3.402033in}}%
\pgfpathlineto{\pgfqpoint{2.929688in}{3.402033in}}%
\pgfpathlineto{\pgfqpoint{2.929688in}{3.378439in}}%
\pgfpathmoveto{\pgfqpoint{2.784375in}{3.472812in}}%
\pgfpathlineto{\pgfqpoint{2.784375in}{3.472812in}}%
\pgfpathlineto{\pgfqpoint{2.784375in}{3.496405in}}%
\pgfpathlineto{\pgfqpoint{2.820703in}{3.496405in}}%
\pgfpathlineto{\pgfqpoint{2.820703in}{3.472812in}}%
\pgfpathmoveto{\pgfqpoint{2.784375in}{3.496405in}}%
\pgfpathlineto{\pgfqpoint{2.784375in}{3.496405in}}%
\pgfpathlineto{\pgfqpoint{2.784375in}{3.519998in}}%
\pgfpathlineto{\pgfqpoint{2.820703in}{3.519998in}}%
\pgfpathlineto{\pgfqpoint{2.820703in}{3.496405in}}%
\pgfpathmoveto{\pgfqpoint{2.820703in}{3.472812in}}%
\pgfpathlineto{\pgfqpoint{2.820703in}{3.472812in}}%
\pgfpathlineto{\pgfqpoint{2.820703in}{3.496405in}}%
\pgfpathlineto{\pgfqpoint{2.857031in}{3.496405in}}%
\pgfpathlineto{\pgfqpoint{2.857031in}{3.472812in}}%
\pgfpathmoveto{\pgfqpoint{2.857031in}{3.425627in}}%
\pgfpathlineto{\pgfqpoint{2.857031in}{3.425627in}}%
\pgfpathlineto{\pgfqpoint{2.857031in}{3.449220in}}%
\pgfpathlineto{\pgfqpoint{2.893360in}{3.449220in}}%
\pgfpathlineto{\pgfqpoint{2.893360in}{3.425627in}}%
\pgfpathmoveto{\pgfqpoint{2.929688in}{1.939220in}}%
\pgfpathlineto{\pgfqpoint{2.929688in}{1.939220in}}%
\pgfpathlineto{\pgfqpoint{2.929688in}{1.962813in}}%
\pgfpathlineto{\pgfqpoint{2.966017in}{1.962813in}}%
\pgfpathlineto{\pgfqpoint{2.966017in}{1.939220in}}%
\pgfpathmoveto{\pgfqpoint{2.929688in}{1.962813in}}%
\pgfpathlineto{\pgfqpoint{2.929688in}{1.962813in}}%
\pgfpathlineto{\pgfqpoint{2.929688in}{1.986407in}}%
\pgfpathlineto{\pgfqpoint{2.966017in}{1.986407in}}%
\pgfpathlineto{\pgfqpoint{2.966017in}{1.962813in}}%
\pgfpathmoveto{\pgfqpoint{2.929688in}{1.986407in}}%
\pgfpathlineto{\pgfqpoint{2.929688in}{1.986407in}}%
\pgfpathlineto{\pgfqpoint{2.929688in}{2.010000in}}%
\pgfpathlineto{\pgfqpoint{2.966017in}{2.010000in}}%
\pgfpathlineto{\pgfqpoint{2.966017in}{1.986407in}}%
\pgfpathmoveto{\pgfqpoint{2.966017in}{1.986407in}}%
\pgfpathlineto{\pgfqpoint{2.966017in}{1.986407in}}%
\pgfpathlineto{\pgfqpoint{2.966017in}{2.010000in}}%
\pgfpathlineto{\pgfqpoint{3.002346in}{2.010000in}}%
\pgfpathlineto{\pgfqpoint{3.002346in}{1.986407in}}%
\pgfpathmoveto{\pgfqpoint{3.002346in}{2.010000in}}%
\pgfpathlineto{\pgfqpoint{3.002346in}{2.010000in}}%
\pgfpathlineto{\pgfqpoint{3.002346in}{2.033594in}}%
\pgfpathlineto{\pgfqpoint{3.038675in}{2.033594in}}%
\pgfpathlineto{\pgfqpoint{3.038675in}{2.010000in}}%
\pgfpathmoveto{\pgfqpoint{3.002346in}{2.033594in}}%
\pgfpathlineto{\pgfqpoint{3.002346in}{2.033594in}}%
\pgfpathlineto{\pgfqpoint{3.002346in}{2.057189in}}%
\pgfpathlineto{\pgfqpoint{3.038675in}{2.057189in}}%
\pgfpathlineto{\pgfqpoint{3.038675in}{2.033594in}}%
\pgfpathmoveto{\pgfqpoint{3.038675in}{2.033594in}}%
\pgfpathlineto{\pgfqpoint{3.038675in}{2.033594in}}%
\pgfpathlineto{\pgfqpoint{3.038675in}{2.057189in}}%
\pgfpathlineto{\pgfqpoint{3.075004in}{2.057189in}}%
\pgfpathlineto{\pgfqpoint{3.075004in}{2.033594in}}%
\pgfpathmoveto{\pgfqpoint{3.002346in}{3.189688in}}%
\pgfpathlineto{\pgfqpoint{3.002346in}{3.189688in}}%
\pgfpathlineto{\pgfqpoint{3.002346in}{3.213281in}}%
\pgfpathlineto{\pgfqpoint{3.038675in}{3.213281in}}%
\pgfpathlineto{\pgfqpoint{3.038675in}{3.189688in}}%
\pgfpathmoveto{\pgfqpoint{3.002346in}{3.213281in}}%
\pgfpathlineto{\pgfqpoint{3.002346in}{3.213281in}}%
\pgfpathlineto{\pgfqpoint{3.002346in}{3.236874in}}%
\pgfpathlineto{\pgfqpoint{3.038675in}{3.236874in}}%
\pgfpathlineto{\pgfqpoint{3.038675in}{3.213281in}}%
\pgfpathmoveto{\pgfqpoint{3.038675in}{3.189688in}}%
\pgfpathlineto{\pgfqpoint{3.038675in}{3.189688in}}%
\pgfpathlineto{\pgfqpoint{3.038675in}{3.213281in}}%
\pgfpathlineto{\pgfqpoint{3.075004in}{3.213281in}}%
\pgfpathlineto{\pgfqpoint{3.075004in}{3.189688in}}%
\pgfpathmoveto{\pgfqpoint{2.929688in}{3.284063in}}%
\pgfpathlineto{\pgfqpoint{2.929688in}{3.284063in}}%
\pgfpathlineto{\pgfqpoint{2.929688in}{3.307657in}}%
\pgfpathlineto{\pgfqpoint{2.966017in}{3.307657in}}%
\pgfpathlineto{\pgfqpoint{2.966017in}{3.284063in}}%
\pgfpathmoveto{\pgfqpoint{2.929688in}{3.307657in}}%
\pgfpathlineto{\pgfqpoint{2.929688in}{3.307657in}}%
\pgfpathlineto{\pgfqpoint{2.929688in}{3.331252in}}%
\pgfpathlineto{\pgfqpoint{2.966017in}{3.331252in}}%
\pgfpathlineto{\pgfqpoint{2.966017in}{3.307657in}}%
\pgfpathmoveto{\pgfqpoint{2.966017in}{3.284063in}}%
\pgfpathlineto{\pgfqpoint{2.966017in}{3.284063in}}%
\pgfpathlineto{\pgfqpoint{2.966017in}{3.307657in}}%
\pgfpathlineto{\pgfqpoint{3.002346in}{3.307657in}}%
\pgfpathlineto{\pgfqpoint{3.002346in}{3.284063in}}%
\pgfpathmoveto{\pgfqpoint{3.002346in}{3.236874in}}%
\pgfpathlineto{\pgfqpoint{3.002346in}{3.236874in}}%
\pgfpathlineto{\pgfqpoint{3.002346in}{3.260469in}}%
\pgfpathlineto{\pgfqpoint{3.038675in}{3.260469in}}%
\pgfpathlineto{\pgfqpoint{3.038675in}{3.236874in}}%
\pgfpathmoveto{\pgfqpoint{2.929688in}{3.331252in}}%
\pgfpathlineto{\pgfqpoint{2.929688in}{3.331252in}}%
\pgfpathlineto{\pgfqpoint{2.929688in}{3.354845in}}%
\pgfpathlineto{\pgfqpoint{2.966017in}{3.354845in}}%
\pgfpathlineto{\pgfqpoint{2.966017in}{3.331252in}}%
\pgfpathmoveto{\pgfqpoint{3.147657in}{2.080783in}}%
\pgfpathlineto{\pgfqpoint{3.147657in}{2.080783in}}%
\pgfpathlineto{\pgfqpoint{3.147657in}{2.104377in}}%
\pgfpathlineto{\pgfqpoint{3.183983in}{2.104377in}}%
\pgfpathlineto{\pgfqpoint{3.183983in}{2.080783in}}%
\pgfpathmoveto{\pgfqpoint{3.147657in}{3.000938in}}%
\pgfpathlineto{\pgfqpoint{3.147657in}{3.000938in}}%
\pgfpathlineto{\pgfqpoint{3.147657in}{3.024531in}}%
\pgfpathlineto{\pgfqpoint{3.183983in}{3.024531in}}%
\pgfpathlineto{\pgfqpoint{3.183983in}{3.000938in}}%
\pgfpathmoveto{\pgfqpoint{3.147657in}{3.024531in}}%
\pgfpathlineto{\pgfqpoint{3.147657in}{3.024531in}}%
\pgfpathlineto{\pgfqpoint{3.147657in}{3.048124in}}%
\pgfpathlineto{\pgfqpoint{3.183983in}{3.048124in}}%
\pgfpathlineto{\pgfqpoint{3.183983in}{3.024531in}}%
\pgfpathmoveto{\pgfqpoint{3.183983in}{3.000938in}}%
\pgfpathlineto{\pgfqpoint{3.183983in}{3.000938in}}%
\pgfpathlineto{\pgfqpoint{3.183983in}{3.024531in}}%
\pgfpathlineto{\pgfqpoint{3.220310in}{3.024531in}}%
\pgfpathlineto{\pgfqpoint{3.220310in}{3.000938in}}%
\pgfpathmoveto{\pgfqpoint{3.075004in}{3.095313in}}%
\pgfpathlineto{\pgfqpoint{3.075004in}{3.095313in}}%
\pgfpathlineto{\pgfqpoint{3.075004in}{3.118907in}}%
\pgfpathlineto{\pgfqpoint{3.111330in}{3.118907in}}%
\pgfpathlineto{\pgfqpoint{3.111330in}{3.095313in}}%
\pgfpathmoveto{\pgfqpoint{3.075004in}{3.118907in}}%
\pgfpathlineto{\pgfqpoint{3.075004in}{3.118907in}}%
\pgfpathlineto{\pgfqpoint{3.075004in}{3.142502in}}%
\pgfpathlineto{\pgfqpoint{3.111330in}{3.142502in}}%
\pgfpathlineto{\pgfqpoint{3.111330in}{3.118907in}}%
\pgfpathmoveto{\pgfqpoint{3.111330in}{3.095313in}}%
\pgfpathlineto{\pgfqpoint{3.111330in}{3.095313in}}%
\pgfpathlineto{\pgfqpoint{3.111330in}{3.118907in}}%
\pgfpathlineto{\pgfqpoint{3.147657in}{3.118907in}}%
\pgfpathlineto{\pgfqpoint{3.147657in}{3.095313in}}%
\pgfpathmoveto{\pgfqpoint{3.147657in}{3.048124in}}%
\pgfpathlineto{\pgfqpoint{3.147657in}{3.048124in}}%
\pgfpathlineto{\pgfqpoint{3.147657in}{3.071719in}}%
\pgfpathlineto{\pgfqpoint{3.183983in}{3.071719in}}%
\pgfpathlineto{\pgfqpoint{3.183983in}{3.048124in}}%
\pgfpathmoveto{\pgfqpoint{3.075004in}{3.142502in}}%
\pgfpathlineto{\pgfqpoint{3.075004in}{3.142502in}}%
\pgfpathlineto{\pgfqpoint{3.075004in}{3.166095in}}%
\pgfpathlineto{\pgfqpoint{3.111330in}{3.166095in}}%
\pgfpathlineto{\pgfqpoint{3.111330in}{3.142502in}}%
\pgfpathmoveto{\pgfqpoint{3.220310in}{2.127970in}}%
\pgfpathlineto{\pgfqpoint{3.220310in}{2.127970in}}%
\pgfpathlineto{\pgfqpoint{3.220310in}{2.151563in}}%
\pgfpathlineto{\pgfqpoint{3.256639in}{2.151563in}}%
\pgfpathlineto{\pgfqpoint{3.256639in}{2.127970in}}%
\pgfpathmoveto{\pgfqpoint{3.292968in}{2.198748in}}%
\pgfpathlineto{\pgfqpoint{3.292968in}{2.198748in}}%
\pgfpathlineto{\pgfqpoint{3.292968in}{2.222342in}}%
\pgfpathlineto{\pgfqpoint{3.329297in}{2.222342in}}%
\pgfpathlineto{\pgfqpoint{3.329297in}{2.198748in}}%
\pgfpathmoveto{\pgfqpoint{3.292968in}{2.222342in}}%
\pgfpathlineto{\pgfqpoint{3.292968in}{2.222342in}}%
\pgfpathlineto{\pgfqpoint{3.292968in}{2.245935in}}%
\pgfpathlineto{\pgfqpoint{3.329297in}{2.245935in}}%
\pgfpathlineto{\pgfqpoint{3.329297in}{2.222342in}}%
\pgfpathmoveto{\pgfqpoint{3.329297in}{2.222342in}}%
\pgfpathlineto{\pgfqpoint{3.329297in}{2.222342in}}%
\pgfpathlineto{\pgfqpoint{3.329297in}{2.245935in}}%
\pgfpathlineto{\pgfqpoint{3.365627in}{2.245935in}}%
\pgfpathlineto{\pgfqpoint{3.365627in}{2.222342in}}%
\pgfpathmoveto{\pgfqpoint{3.292968in}{2.812187in}}%
\pgfpathlineto{\pgfqpoint{3.292968in}{2.812187in}}%
\pgfpathlineto{\pgfqpoint{3.292968in}{2.835780in}}%
\pgfpathlineto{\pgfqpoint{3.329297in}{2.835780in}}%
\pgfpathlineto{\pgfqpoint{3.329297in}{2.812187in}}%
\pgfpathmoveto{\pgfqpoint{3.292968in}{2.835780in}}%
\pgfpathlineto{\pgfqpoint{3.292968in}{2.835780in}}%
\pgfpathlineto{\pgfqpoint{3.292968in}{2.859374in}}%
\pgfpathlineto{\pgfqpoint{3.329297in}{2.859374in}}%
\pgfpathlineto{\pgfqpoint{3.329297in}{2.835780in}}%
\pgfpathmoveto{\pgfqpoint{3.329297in}{2.812187in}}%
\pgfpathlineto{\pgfqpoint{3.329297in}{2.812187in}}%
\pgfpathlineto{\pgfqpoint{3.329297in}{2.835780in}}%
\pgfpathlineto{\pgfqpoint{3.365627in}{2.835780in}}%
\pgfpathlineto{\pgfqpoint{3.365627in}{2.812187in}}%
\pgfpathmoveto{\pgfqpoint{3.220310in}{2.906563in}}%
\pgfpathlineto{\pgfqpoint{3.220310in}{2.906563in}}%
\pgfpathlineto{\pgfqpoint{3.220310in}{2.930157in}}%
\pgfpathlineto{\pgfqpoint{3.256639in}{2.930157in}}%
\pgfpathlineto{\pgfqpoint{3.256639in}{2.906563in}}%
\pgfpathmoveto{\pgfqpoint{3.220310in}{2.930157in}}%
\pgfpathlineto{\pgfqpoint{3.220310in}{2.930157in}}%
\pgfpathlineto{\pgfqpoint{3.220310in}{2.953751in}}%
\pgfpathlineto{\pgfqpoint{3.256639in}{2.953751in}}%
\pgfpathlineto{\pgfqpoint{3.256639in}{2.930157in}}%
\pgfpathmoveto{\pgfqpoint{3.256639in}{2.906563in}}%
\pgfpathlineto{\pgfqpoint{3.256639in}{2.906563in}}%
\pgfpathlineto{\pgfqpoint{3.256639in}{2.930157in}}%
\pgfpathlineto{\pgfqpoint{3.292968in}{2.930157in}}%
\pgfpathlineto{\pgfqpoint{3.292968in}{2.906563in}}%
\pgfpathmoveto{\pgfqpoint{3.292968in}{2.859374in}}%
\pgfpathlineto{\pgfqpoint{3.292968in}{2.859374in}}%
\pgfpathlineto{\pgfqpoint{3.292968in}{2.882968in}}%
\pgfpathlineto{\pgfqpoint{3.329297in}{2.882968in}}%
\pgfpathlineto{\pgfqpoint{3.329297in}{2.859374in}}%
\pgfpathmoveto{\pgfqpoint{3.220310in}{2.953751in}}%
\pgfpathlineto{\pgfqpoint{3.220310in}{2.953751in}}%
\pgfpathlineto{\pgfqpoint{3.220310in}{2.977345in}}%
\pgfpathlineto{\pgfqpoint{3.256639in}{2.977345in}}%
\pgfpathlineto{\pgfqpoint{3.256639in}{2.953751in}}%
\pgfpathmoveto{\pgfqpoint{3.365627in}{2.245935in}}%
\pgfpathlineto{\pgfqpoint{3.365627in}{2.245935in}}%
\pgfpathlineto{\pgfqpoint{3.365627in}{2.269529in}}%
\pgfpathlineto{\pgfqpoint{3.401955in}{2.269529in}}%
\pgfpathlineto{\pgfqpoint{3.401955in}{2.245935in}}%
\pgfpathmoveto{\pgfqpoint{3.365627in}{2.269529in}}%
\pgfpathlineto{\pgfqpoint{3.365627in}{2.269529in}}%
\pgfpathlineto{\pgfqpoint{3.365627in}{2.293122in}}%
\pgfpathlineto{\pgfqpoint{3.401955in}{2.293122in}}%
\pgfpathlineto{\pgfqpoint{3.401955in}{2.269529in}}%
\pgfpathmoveto{\pgfqpoint{3.401955in}{2.269529in}}%
\pgfpathlineto{\pgfqpoint{3.401955in}{2.269529in}}%
\pgfpathlineto{\pgfqpoint{3.401955in}{2.293122in}}%
\pgfpathlineto{\pgfqpoint{3.438283in}{2.293122in}}%
\pgfpathlineto{\pgfqpoint{3.438283in}{2.269529in}}%
\pgfpathmoveto{\pgfqpoint{3.438283in}{2.293122in}}%
\pgfpathlineto{\pgfqpoint{3.438283in}{2.293122in}}%
\pgfpathlineto{\pgfqpoint{3.438283in}{2.316716in}}%
\pgfpathlineto{\pgfqpoint{3.474611in}{2.316716in}}%
\pgfpathlineto{\pgfqpoint{3.474611in}{2.293122in}}%
\pgfpathmoveto{\pgfqpoint{3.438283in}{2.316716in}}%
\pgfpathlineto{\pgfqpoint{3.438283in}{2.316716in}}%
\pgfpathlineto{\pgfqpoint{3.438283in}{2.340311in}}%
\pgfpathlineto{\pgfqpoint{3.474611in}{2.340311in}}%
\pgfpathlineto{\pgfqpoint{3.474611in}{2.316716in}}%
\pgfpathmoveto{\pgfqpoint{3.474611in}{2.316716in}}%
\pgfpathlineto{\pgfqpoint{3.474611in}{2.316716in}}%
\pgfpathlineto{\pgfqpoint{3.474611in}{2.340311in}}%
\pgfpathlineto{\pgfqpoint{3.510939in}{2.340311in}}%
\pgfpathlineto{\pgfqpoint{3.510939in}{2.316716in}}%
\pgfpathmoveto{\pgfqpoint{3.438283in}{2.623437in}}%
\pgfpathlineto{\pgfqpoint{3.438283in}{2.623437in}}%
\pgfpathlineto{\pgfqpoint{3.438283in}{2.647031in}}%
\pgfpathlineto{\pgfqpoint{3.474611in}{2.647031in}}%
\pgfpathlineto{\pgfqpoint{3.474611in}{2.623437in}}%
\pgfpathmoveto{\pgfqpoint{3.438283in}{2.647031in}}%
\pgfpathlineto{\pgfqpoint{3.438283in}{2.647031in}}%
\pgfpathlineto{\pgfqpoint{3.438283in}{2.670624in}}%
\pgfpathlineto{\pgfqpoint{3.474611in}{2.670624in}}%
\pgfpathlineto{\pgfqpoint{3.474611in}{2.647031in}}%
\pgfpathmoveto{\pgfqpoint{3.474611in}{2.623437in}}%
\pgfpathlineto{\pgfqpoint{3.474611in}{2.623437in}}%
\pgfpathlineto{\pgfqpoint{3.474611in}{2.647031in}}%
\pgfpathlineto{\pgfqpoint{3.510939in}{2.647031in}}%
\pgfpathlineto{\pgfqpoint{3.510939in}{2.623437in}}%
\pgfpathmoveto{\pgfqpoint{3.365627in}{2.717812in}}%
\pgfpathlineto{\pgfqpoint{3.365627in}{2.717812in}}%
\pgfpathlineto{\pgfqpoint{3.365627in}{2.741405in}}%
\pgfpathlineto{\pgfqpoint{3.401955in}{2.741405in}}%
\pgfpathlineto{\pgfqpoint{3.401955in}{2.717812in}}%
\pgfpathmoveto{\pgfqpoint{3.365627in}{2.741405in}}%
\pgfpathlineto{\pgfqpoint{3.365627in}{2.741405in}}%
\pgfpathlineto{\pgfqpoint{3.365627in}{2.764999in}}%
\pgfpathlineto{\pgfqpoint{3.401955in}{2.764999in}}%
\pgfpathlineto{\pgfqpoint{3.401955in}{2.741405in}}%
\pgfpathmoveto{\pgfqpoint{3.401955in}{2.717812in}}%
\pgfpathlineto{\pgfqpoint{3.401955in}{2.717812in}}%
\pgfpathlineto{\pgfqpoint{3.401955in}{2.741405in}}%
\pgfpathlineto{\pgfqpoint{3.438283in}{2.741405in}}%
\pgfpathlineto{\pgfqpoint{3.438283in}{2.717812in}}%
\pgfpathmoveto{\pgfqpoint{3.438283in}{2.670624in}}%
\pgfpathlineto{\pgfqpoint{3.438283in}{2.670624in}}%
\pgfpathlineto{\pgfqpoint{3.438283in}{2.694218in}}%
\pgfpathlineto{\pgfqpoint{3.474611in}{2.694218in}}%
\pgfpathlineto{\pgfqpoint{3.474611in}{2.670624in}}%
\pgfpathmoveto{\pgfqpoint{3.365627in}{2.764999in}}%
\pgfpathlineto{\pgfqpoint{3.365627in}{2.764999in}}%
\pgfpathlineto{\pgfqpoint{3.365627in}{2.788593in}}%
\pgfpathlineto{\pgfqpoint{3.401955in}{2.788593in}}%
\pgfpathlineto{\pgfqpoint{3.401955in}{2.764999in}}%
\pgfpathmoveto{\pgfqpoint{3.510939in}{2.340311in}}%
\pgfpathlineto{\pgfqpoint{3.510939in}{2.340311in}}%
\pgfpathlineto{\pgfqpoint{3.510939in}{2.363905in}}%
\pgfpathlineto{\pgfqpoint{3.547265in}{2.363905in}}%
\pgfpathlineto{\pgfqpoint{3.547265in}{2.340311in}}%
\pgfpathmoveto{\pgfqpoint{3.510939in}{2.363905in}}%
\pgfpathlineto{\pgfqpoint{3.510939in}{2.363905in}}%
\pgfpathlineto{\pgfqpoint{3.510939in}{2.387499in}}%
\pgfpathlineto{\pgfqpoint{3.547265in}{2.387499in}}%
\pgfpathlineto{\pgfqpoint{3.547265in}{2.363905in}}%
\pgfpathmoveto{\pgfqpoint{3.547265in}{2.363905in}}%
\pgfpathlineto{\pgfqpoint{3.547265in}{2.363905in}}%
\pgfpathlineto{\pgfqpoint{3.547265in}{2.387499in}}%
\pgfpathlineto{\pgfqpoint{3.583592in}{2.387499in}}%
\pgfpathlineto{\pgfqpoint{3.583592in}{2.363905in}}%
\pgfpathmoveto{\pgfqpoint{3.583592in}{2.363905in}}%
\pgfpathlineto{\pgfqpoint{3.583592in}{2.363905in}}%
\pgfpathlineto{\pgfqpoint{3.583592in}{2.387499in}}%
\pgfpathlineto{\pgfqpoint{3.619919in}{2.387499in}}%
\pgfpathlineto{\pgfqpoint{3.619919in}{2.363905in}}%
\pgfpathmoveto{\pgfqpoint{3.583592in}{2.434688in}}%
\pgfpathlineto{\pgfqpoint{3.583592in}{2.434688in}}%
\pgfpathlineto{\pgfqpoint{3.583592in}{2.458283in}}%
\pgfpathlineto{\pgfqpoint{3.619919in}{2.458283in}}%
\pgfpathlineto{\pgfqpoint{3.619919in}{2.434688in}}%
\pgfpathmoveto{\pgfqpoint{3.583592in}{2.458283in}}%
\pgfpathlineto{\pgfqpoint{3.583592in}{2.458283in}}%
\pgfpathlineto{\pgfqpoint{3.583592in}{2.481878in}}%
\pgfpathlineto{\pgfqpoint{3.619919in}{2.481878in}}%
\pgfpathlineto{\pgfqpoint{3.619919in}{2.458283in}}%
\pgfpathmoveto{\pgfqpoint{3.619919in}{2.434688in}}%
\pgfpathlineto{\pgfqpoint{3.619919in}{2.434688in}}%
\pgfpathlineto{\pgfqpoint{3.619919in}{2.458283in}}%
\pgfpathlineto{\pgfqpoint{3.656246in}{2.458283in}}%
\pgfpathlineto{\pgfqpoint{3.656246in}{2.434688in}}%
\pgfpathmoveto{\pgfqpoint{3.510939in}{2.529064in}}%
\pgfpathlineto{\pgfqpoint{3.510939in}{2.529064in}}%
\pgfpathlineto{\pgfqpoint{3.510939in}{2.552658in}}%
\pgfpathlineto{\pgfqpoint{3.547265in}{2.552658in}}%
\pgfpathlineto{\pgfqpoint{3.547265in}{2.529064in}}%
\pgfpathmoveto{\pgfqpoint{3.510939in}{2.552658in}}%
\pgfpathlineto{\pgfqpoint{3.510939in}{2.552658in}}%
\pgfpathlineto{\pgfqpoint{3.510939in}{2.576251in}}%
\pgfpathlineto{\pgfqpoint{3.547265in}{2.576251in}}%
\pgfpathlineto{\pgfqpoint{3.547265in}{2.552658in}}%
\pgfpathmoveto{\pgfqpoint{3.547265in}{2.529064in}}%
\pgfpathlineto{\pgfqpoint{3.547265in}{2.529064in}}%
\pgfpathlineto{\pgfqpoint{3.547265in}{2.552658in}}%
\pgfpathlineto{\pgfqpoint{3.583592in}{2.552658in}}%
\pgfpathlineto{\pgfqpoint{3.583592in}{2.529064in}}%
\pgfpathmoveto{\pgfqpoint{3.583592in}{2.481878in}}%
\pgfpathlineto{\pgfqpoint{3.583592in}{2.481878in}}%
\pgfpathlineto{\pgfqpoint{3.583592in}{2.505471in}}%
\pgfpathlineto{\pgfqpoint{3.619919in}{2.505471in}}%
\pgfpathlineto{\pgfqpoint{3.619919in}{2.481878in}}%
\pgfpathmoveto{\pgfqpoint{3.510939in}{2.576251in}}%
\pgfpathlineto{\pgfqpoint{3.510939in}{2.576251in}}%
\pgfpathlineto{\pgfqpoint{3.510939in}{2.599844in}}%
\pgfpathlineto{\pgfqpoint{3.547265in}{2.599844in}}%
\pgfpathlineto{\pgfqpoint{3.547265in}{2.576251in}}%
\pgfpathmoveto{\pgfqpoint{0.749998in}{0.523592in}}%
\pgfpathlineto{\pgfqpoint{0.749998in}{0.523592in}}%
\pgfpathlineto{\pgfqpoint{0.749998in}{0.535389in}}%
\pgfpathlineto{\pgfqpoint{0.768162in}{0.535389in}}%
\pgfpathlineto{\pgfqpoint{0.768162in}{0.523592in}}%
\pgfpathmoveto{\pgfqpoint{0.749998in}{0.535389in}}%
\pgfpathlineto{\pgfqpoint{0.749998in}{0.535389in}}%
\pgfpathlineto{\pgfqpoint{0.749998in}{0.547186in}}%
\pgfpathlineto{\pgfqpoint{0.768162in}{0.547186in}}%
\pgfpathlineto{\pgfqpoint{0.768162in}{0.535389in}}%
\pgfpathmoveto{\pgfqpoint{0.768162in}{0.535389in}}%
\pgfpathlineto{\pgfqpoint{0.768162in}{0.535389in}}%
\pgfpathlineto{\pgfqpoint{0.768162in}{0.547186in}}%
\pgfpathlineto{\pgfqpoint{0.786327in}{0.547186in}}%
\pgfpathlineto{\pgfqpoint{0.786327in}{0.535389in}}%
\pgfpathmoveto{\pgfqpoint{0.786327in}{0.547186in}}%
\pgfpathlineto{\pgfqpoint{0.786327in}{0.547186in}}%
\pgfpathlineto{\pgfqpoint{0.786327in}{0.558983in}}%
\pgfpathlineto{\pgfqpoint{0.804491in}{0.558983in}}%
\pgfpathlineto{\pgfqpoint{0.804491in}{0.547186in}}%
\pgfpathmoveto{\pgfqpoint{0.786327in}{0.558983in}}%
\pgfpathlineto{\pgfqpoint{0.786327in}{0.558983in}}%
\pgfpathlineto{\pgfqpoint{0.786327in}{0.570780in}}%
\pgfpathlineto{\pgfqpoint{0.804491in}{0.570780in}}%
\pgfpathlineto{\pgfqpoint{0.804491in}{0.558983in}}%
\pgfpathmoveto{\pgfqpoint{0.804491in}{0.558983in}}%
\pgfpathlineto{\pgfqpoint{0.804491in}{0.558983in}}%
\pgfpathlineto{\pgfqpoint{0.804491in}{0.570780in}}%
\pgfpathlineto{\pgfqpoint{0.822655in}{0.570780in}}%
\pgfpathlineto{\pgfqpoint{0.822655in}{0.558983in}}%
\pgfpathmoveto{\pgfqpoint{0.822655in}{0.570780in}}%
\pgfpathlineto{\pgfqpoint{0.822655in}{0.570780in}}%
\pgfpathlineto{\pgfqpoint{0.822655in}{0.582577in}}%
\pgfpathlineto{\pgfqpoint{0.840819in}{0.582577in}}%
\pgfpathlineto{\pgfqpoint{0.840819in}{0.570780in}}%
\pgfpathmoveto{\pgfqpoint{0.822655in}{0.582577in}}%
\pgfpathlineto{\pgfqpoint{0.822655in}{0.582577in}}%
\pgfpathlineto{\pgfqpoint{0.822655in}{0.594374in}}%
\pgfpathlineto{\pgfqpoint{0.840819in}{0.594374in}}%
\pgfpathlineto{\pgfqpoint{0.840819in}{0.582577in}}%
\pgfpathmoveto{\pgfqpoint{0.840819in}{0.582577in}}%
\pgfpathlineto{\pgfqpoint{0.840819in}{0.582577in}}%
\pgfpathlineto{\pgfqpoint{0.840819in}{0.594374in}}%
\pgfpathlineto{\pgfqpoint{0.858983in}{0.594374in}}%
\pgfpathlineto{\pgfqpoint{0.858983in}{0.582577in}}%
\pgfpathmoveto{\pgfqpoint{0.858983in}{0.594374in}}%
\pgfpathlineto{\pgfqpoint{0.858983in}{0.594374in}}%
\pgfpathlineto{\pgfqpoint{0.858983in}{0.606171in}}%
\pgfpathlineto{\pgfqpoint{0.877147in}{0.606171in}}%
\pgfpathlineto{\pgfqpoint{0.877147in}{0.594374in}}%
\pgfpathmoveto{\pgfqpoint{0.858983in}{0.606171in}}%
\pgfpathlineto{\pgfqpoint{0.858983in}{0.606171in}}%
\pgfpathlineto{\pgfqpoint{0.858983in}{0.617968in}}%
\pgfpathlineto{\pgfqpoint{0.877147in}{0.617968in}}%
\pgfpathlineto{\pgfqpoint{0.877147in}{0.606171in}}%
\pgfpathmoveto{\pgfqpoint{0.877147in}{0.606171in}}%
\pgfpathlineto{\pgfqpoint{0.877147in}{0.606171in}}%
\pgfpathlineto{\pgfqpoint{0.877147in}{0.617968in}}%
\pgfpathlineto{\pgfqpoint{0.895312in}{0.617968in}}%
\pgfpathlineto{\pgfqpoint{0.895312in}{0.606171in}}%
\pgfpathmoveto{\pgfqpoint{0.931639in}{0.629765in}}%
\pgfpathlineto{\pgfqpoint{0.931639in}{0.629765in}}%
\pgfpathlineto{\pgfqpoint{0.931639in}{0.641562in}}%
\pgfpathlineto{\pgfqpoint{0.949803in}{0.641562in}}%
\pgfpathlineto{\pgfqpoint{0.949803in}{0.629765in}}%
\pgfpathmoveto{\pgfqpoint{0.967967in}{0.653359in}}%
\pgfpathlineto{\pgfqpoint{0.967967in}{0.653359in}}%
\pgfpathlineto{\pgfqpoint{0.967967in}{0.665156in}}%
\pgfpathlineto{\pgfqpoint{0.986131in}{0.665156in}}%
\pgfpathlineto{\pgfqpoint{0.986131in}{0.653359in}}%
\pgfpathmoveto{\pgfqpoint{1.004294in}{0.676953in}}%
\pgfpathlineto{\pgfqpoint{1.004294in}{0.676953in}}%
\pgfpathlineto{\pgfqpoint{1.004294in}{0.688750in}}%
\pgfpathlineto{\pgfqpoint{1.022458in}{0.688750in}}%
\pgfpathlineto{\pgfqpoint{1.022458in}{0.676953in}}%
\pgfpathmoveto{\pgfqpoint{1.040622in}{0.700547in}}%
\pgfpathlineto{\pgfqpoint{1.040622in}{0.700547in}}%
\pgfpathlineto{\pgfqpoint{1.040622in}{0.712343in}}%
\pgfpathlineto{\pgfqpoint{1.058787in}{0.712343in}}%
\pgfpathlineto{\pgfqpoint{1.058787in}{0.700547in}}%
\pgfpathmoveto{\pgfqpoint{1.076952in}{0.735937in}}%
\pgfpathlineto{\pgfqpoint{1.076952in}{0.735937in}}%
\pgfpathlineto{\pgfqpoint{1.076952in}{0.747734in}}%
\pgfpathlineto{\pgfqpoint{1.095117in}{0.747734in}}%
\pgfpathlineto{\pgfqpoint{1.095117in}{0.735937in}}%
\pgfpathmoveto{\pgfqpoint{1.076952in}{0.747734in}}%
\pgfpathlineto{\pgfqpoint{1.076952in}{0.747734in}}%
\pgfpathlineto{\pgfqpoint{1.076952in}{0.759530in}}%
\pgfpathlineto{\pgfqpoint{1.095117in}{0.759530in}}%
\pgfpathlineto{\pgfqpoint{1.095117in}{0.747734in}}%
\pgfpathmoveto{\pgfqpoint{1.095117in}{0.747734in}}%
\pgfpathlineto{\pgfqpoint{1.095117in}{0.747734in}}%
\pgfpathlineto{\pgfqpoint{1.095117in}{0.759530in}}%
\pgfpathlineto{\pgfqpoint{1.113282in}{0.759530in}}%
\pgfpathlineto{\pgfqpoint{1.113282in}{0.747734in}}%
\pgfpathmoveto{\pgfqpoint{1.113282in}{0.759530in}}%
\pgfpathlineto{\pgfqpoint{1.113282in}{0.759530in}}%
\pgfpathlineto{\pgfqpoint{1.113282in}{0.771327in}}%
\pgfpathlineto{\pgfqpoint{1.131447in}{0.771327in}}%
\pgfpathlineto{\pgfqpoint{1.131447in}{0.759530in}}%
\pgfpathmoveto{\pgfqpoint{1.113282in}{0.771327in}}%
\pgfpathlineto{\pgfqpoint{1.113282in}{0.771327in}}%
\pgfpathlineto{\pgfqpoint{1.113282in}{0.783124in}}%
\pgfpathlineto{\pgfqpoint{1.131447in}{0.783124in}}%
\pgfpathlineto{\pgfqpoint{1.131447in}{0.771327in}}%
\pgfpathmoveto{\pgfqpoint{1.131447in}{0.771327in}}%
\pgfpathlineto{\pgfqpoint{1.131447in}{0.771327in}}%
\pgfpathlineto{\pgfqpoint{1.131447in}{0.783124in}}%
\pgfpathlineto{\pgfqpoint{1.149611in}{0.783124in}}%
\pgfpathlineto{\pgfqpoint{1.149611in}{0.771327in}}%
\pgfpathmoveto{\pgfqpoint{1.149611in}{0.783124in}}%
\pgfpathlineto{\pgfqpoint{1.149611in}{0.783124in}}%
\pgfpathlineto{\pgfqpoint{1.149611in}{0.794921in}}%
\pgfpathlineto{\pgfqpoint{1.167776in}{0.794921in}}%
\pgfpathlineto{\pgfqpoint{1.167776in}{0.783124in}}%
\pgfpathmoveto{\pgfqpoint{1.149611in}{0.794921in}}%
\pgfpathlineto{\pgfqpoint{1.149611in}{0.794921in}}%
\pgfpathlineto{\pgfqpoint{1.149611in}{0.806718in}}%
\pgfpathlineto{\pgfqpoint{1.167776in}{0.806718in}}%
\pgfpathlineto{\pgfqpoint{1.167776in}{0.794921in}}%
\pgfpathmoveto{\pgfqpoint{1.167776in}{0.794921in}}%
\pgfpathlineto{\pgfqpoint{1.167776in}{0.794921in}}%
\pgfpathlineto{\pgfqpoint{1.167776in}{0.806718in}}%
\pgfpathlineto{\pgfqpoint{1.185941in}{0.806718in}}%
\pgfpathlineto{\pgfqpoint{1.185941in}{0.794921in}}%
\pgfpathmoveto{\pgfqpoint{1.185941in}{0.806718in}}%
\pgfpathlineto{\pgfqpoint{1.185941in}{0.806718in}}%
\pgfpathlineto{\pgfqpoint{1.185941in}{0.818515in}}%
\pgfpathlineto{\pgfqpoint{1.204105in}{0.818515in}}%
\pgfpathlineto{\pgfqpoint{1.204105in}{0.806718in}}%
\pgfpathmoveto{\pgfqpoint{1.185941in}{0.818515in}}%
\pgfpathlineto{\pgfqpoint{1.185941in}{0.818515in}}%
\pgfpathlineto{\pgfqpoint{1.185941in}{0.830313in}}%
\pgfpathlineto{\pgfqpoint{1.204105in}{0.830313in}}%
\pgfpathlineto{\pgfqpoint{1.204105in}{0.818515in}}%
\pgfpathmoveto{\pgfqpoint{1.204105in}{0.818515in}}%
\pgfpathlineto{\pgfqpoint{1.204105in}{0.818515in}}%
\pgfpathlineto{\pgfqpoint{1.204105in}{0.830313in}}%
\pgfpathlineto{\pgfqpoint{1.222269in}{0.830313in}}%
\pgfpathlineto{\pgfqpoint{1.222269in}{0.818515in}}%
\pgfpathmoveto{\pgfqpoint{1.222269in}{0.830313in}}%
\pgfpathlineto{\pgfqpoint{1.222269in}{0.830313in}}%
\pgfpathlineto{\pgfqpoint{1.222269in}{0.842110in}}%
\pgfpathlineto{\pgfqpoint{1.240433in}{0.842110in}}%
\pgfpathlineto{\pgfqpoint{1.240433in}{0.830313in}}%
\pgfpathmoveto{\pgfqpoint{1.222269in}{0.842110in}}%
\pgfpathlineto{\pgfqpoint{1.222269in}{0.842110in}}%
\pgfpathlineto{\pgfqpoint{1.222269in}{0.853907in}}%
\pgfpathlineto{\pgfqpoint{1.240433in}{0.853907in}}%
\pgfpathlineto{\pgfqpoint{1.240433in}{0.842110in}}%
\pgfpathmoveto{\pgfqpoint{1.240433in}{0.842110in}}%
\pgfpathlineto{\pgfqpoint{1.240433in}{0.842110in}}%
\pgfpathlineto{\pgfqpoint{1.240433in}{0.853907in}}%
\pgfpathlineto{\pgfqpoint{1.258598in}{0.853907in}}%
\pgfpathlineto{\pgfqpoint{1.258598in}{0.842110in}}%
\pgfpathmoveto{\pgfqpoint{1.258598in}{0.853907in}}%
\pgfpathlineto{\pgfqpoint{1.258598in}{0.853907in}}%
\pgfpathlineto{\pgfqpoint{1.258598in}{0.865704in}}%
\pgfpathlineto{\pgfqpoint{1.276762in}{0.865704in}}%
\pgfpathlineto{\pgfqpoint{1.276762in}{0.853907in}}%
\pgfpathmoveto{\pgfqpoint{1.258598in}{0.865704in}}%
\pgfpathlineto{\pgfqpoint{1.258598in}{0.865704in}}%
\pgfpathlineto{\pgfqpoint{1.258598in}{0.877501in}}%
\pgfpathlineto{\pgfqpoint{1.276762in}{0.877501in}}%
\pgfpathlineto{\pgfqpoint{1.276762in}{0.865704in}}%
\pgfpathmoveto{\pgfqpoint{1.276762in}{0.865704in}}%
\pgfpathlineto{\pgfqpoint{1.276762in}{0.865704in}}%
\pgfpathlineto{\pgfqpoint{1.276762in}{0.877501in}}%
\pgfpathlineto{\pgfqpoint{1.294926in}{0.877501in}}%
\pgfpathlineto{\pgfqpoint{1.294926in}{0.865704in}}%
\pgfpathmoveto{\pgfqpoint{1.294926in}{0.877501in}}%
\pgfpathlineto{\pgfqpoint{1.294926in}{0.877501in}}%
\pgfpathlineto{\pgfqpoint{1.294926in}{0.889298in}}%
\pgfpathlineto{\pgfqpoint{1.313090in}{0.889298in}}%
\pgfpathlineto{\pgfqpoint{1.313090in}{0.877501in}}%
\pgfpathmoveto{\pgfqpoint{1.294926in}{0.889298in}}%
\pgfpathlineto{\pgfqpoint{1.294926in}{0.889298in}}%
\pgfpathlineto{\pgfqpoint{1.294926in}{0.901094in}}%
\pgfpathlineto{\pgfqpoint{1.313090in}{0.901094in}}%
\pgfpathlineto{\pgfqpoint{1.313090in}{0.889298in}}%
\pgfpathmoveto{\pgfqpoint{1.313090in}{0.889298in}}%
\pgfpathlineto{\pgfqpoint{1.313090in}{0.889298in}}%
\pgfpathlineto{\pgfqpoint{1.313090in}{0.901094in}}%
\pgfpathlineto{\pgfqpoint{1.331254in}{0.901094in}}%
\pgfpathlineto{\pgfqpoint{1.331254in}{0.889298in}}%
\pgfpathmoveto{\pgfqpoint{1.331254in}{0.901094in}}%
\pgfpathlineto{\pgfqpoint{1.331254in}{0.901094in}}%
\pgfpathlineto{\pgfqpoint{1.331254in}{0.912890in}}%
\pgfpathlineto{\pgfqpoint{1.349418in}{0.912890in}}%
\pgfpathlineto{\pgfqpoint{1.349418in}{0.901094in}}%
\pgfpathmoveto{\pgfqpoint{1.331254in}{0.912890in}}%
\pgfpathlineto{\pgfqpoint{1.331254in}{0.912890in}}%
\pgfpathlineto{\pgfqpoint{1.331254in}{0.924687in}}%
\pgfpathlineto{\pgfqpoint{1.349418in}{0.924687in}}%
\pgfpathlineto{\pgfqpoint{1.349418in}{0.912890in}}%
\pgfpathmoveto{\pgfqpoint{1.349418in}{0.912890in}}%
\pgfpathlineto{\pgfqpoint{1.349418in}{0.912890in}}%
\pgfpathlineto{\pgfqpoint{1.349418in}{0.924687in}}%
\pgfpathlineto{\pgfqpoint{1.367581in}{0.924687in}}%
\pgfpathlineto{\pgfqpoint{1.367581in}{0.912890in}}%
\pgfpathmoveto{\pgfqpoint{1.367581in}{0.924687in}}%
\pgfpathlineto{\pgfqpoint{1.367581in}{0.924687in}}%
\pgfpathlineto{\pgfqpoint{1.367581in}{0.936483in}}%
\pgfpathlineto{\pgfqpoint{1.385745in}{0.936483in}}%
\pgfpathlineto{\pgfqpoint{1.385745in}{0.924687in}}%
\pgfpathmoveto{\pgfqpoint{1.367581in}{0.936483in}}%
\pgfpathlineto{\pgfqpoint{1.367581in}{0.936483in}}%
\pgfpathlineto{\pgfqpoint{1.367581in}{0.948280in}}%
\pgfpathlineto{\pgfqpoint{1.385745in}{0.948280in}}%
\pgfpathlineto{\pgfqpoint{1.385745in}{0.936483in}}%
\pgfpathmoveto{\pgfqpoint{1.385745in}{0.936483in}}%
\pgfpathlineto{\pgfqpoint{1.385745in}{0.936483in}}%
\pgfpathlineto{\pgfqpoint{1.385745in}{0.948280in}}%
\pgfpathlineto{\pgfqpoint{1.403909in}{0.948280in}}%
\pgfpathlineto{\pgfqpoint{1.403909in}{0.936483in}}%
\pgfpathmoveto{\pgfqpoint{1.403909in}{0.948280in}}%
\pgfpathlineto{\pgfqpoint{1.403909in}{0.948280in}}%
\pgfpathlineto{\pgfqpoint{1.403909in}{0.960076in}}%
\pgfpathlineto{\pgfqpoint{1.422073in}{0.960076in}}%
\pgfpathlineto{\pgfqpoint{1.422073in}{0.948280in}}%
\pgfpathmoveto{\pgfqpoint{1.403909in}{0.960076in}}%
\pgfpathlineto{\pgfqpoint{1.403909in}{0.960076in}}%
\pgfpathlineto{\pgfqpoint{1.403909in}{0.971872in}}%
\pgfpathlineto{\pgfqpoint{1.422073in}{0.971872in}}%
\pgfpathlineto{\pgfqpoint{1.422073in}{0.960076in}}%
\pgfpathmoveto{\pgfqpoint{1.422073in}{0.960076in}}%
\pgfpathlineto{\pgfqpoint{1.422073in}{0.960076in}}%
\pgfpathlineto{\pgfqpoint{1.422073in}{0.971872in}}%
\pgfpathlineto{\pgfqpoint{1.440237in}{0.971872in}}%
\pgfpathlineto{\pgfqpoint{1.440237in}{0.960076in}}%
\pgfpathmoveto{\pgfqpoint{1.440237in}{0.971872in}}%
\pgfpathlineto{\pgfqpoint{1.440237in}{0.971872in}}%
\pgfpathlineto{\pgfqpoint{1.440237in}{0.983669in}}%
\pgfpathlineto{\pgfqpoint{1.458400in}{0.983669in}}%
\pgfpathlineto{\pgfqpoint{1.458400in}{0.971872in}}%
\pgfpathmoveto{\pgfqpoint{1.440237in}{0.983669in}}%
\pgfpathlineto{\pgfqpoint{1.440237in}{0.983669in}}%
\pgfpathlineto{\pgfqpoint{1.440237in}{0.995466in}}%
\pgfpathlineto{\pgfqpoint{1.458400in}{0.995466in}}%
\pgfpathlineto{\pgfqpoint{1.458400in}{0.983669in}}%
\pgfpathmoveto{\pgfqpoint{1.458400in}{0.983669in}}%
\pgfpathlineto{\pgfqpoint{1.458400in}{0.983669in}}%
\pgfpathlineto{\pgfqpoint{1.458400in}{0.995466in}}%
\pgfpathlineto{\pgfqpoint{1.476564in}{0.995466in}}%
\pgfpathlineto{\pgfqpoint{1.476564in}{0.983669in}}%
\pgfpathmoveto{\pgfqpoint{1.476564in}{0.995466in}}%
\pgfpathlineto{\pgfqpoint{1.476564in}{0.995466in}}%
\pgfpathlineto{\pgfqpoint{1.476564in}{1.007264in}}%
\pgfpathlineto{\pgfqpoint{1.494728in}{1.007264in}}%
\pgfpathlineto{\pgfqpoint{1.494728in}{0.995466in}}%
\pgfpathmoveto{\pgfqpoint{1.476564in}{1.007264in}}%
\pgfpathlineto{\pgfqpoint{1.476564in}{1.007264in}}%
\pgfpathlineto{\pgfqpoint{1.476564in}{1.019061in}}%
\pgfpathlineto{\pgfqpoint{1.494728in}{1.019061in}}%
\pgfpathlineto{\pgfqpoint{1.494728in}{1.007264in}}%
\pgfpathmoveto{\pgfqpoint{1.494728in}{1.007264in}}%
\pgfpathlineto{\pgfqpoint{1.494728in}{1.007264in}}%
\pgfpathlineto{\pgfqpoint{1.494728in}{1.019061in}}%
\pgfpathlineto{\pgfqpoint{1.512892in}{1.019061in}}%
\pgfpathlineto{\pgfqpoint{1.512892in}{1.007264in}}%
\pgfpathmoveto{\pgfqpoint{1.512892in}{1.019061in}}%
\pgfpathlineto{\pgfqpoint{1.512892in}{1.019061in}}%
\pgfpathlineto{\pgfqpoint{1.512892in}{1.030858in}}%
\pgfpathlineto{\pgfqpoint{1.531056in}{1.030858in}}%
\pgfpathlineto{\pgfqpoint{1.531056in}{1.019061in}}%
\pgfpathmoveto{\pgfqpoint{1.512892in}{1.030858in}}%
\pgfpathlineto{\pgfqpoint{1.512892in}{1.030858in}}%
\pgfpathlineto{\pgfqpoint{1.512892in}{1.042655in}}%
\pgfpathlineto{\pgfqpoint{1.531056in}{1.042655in}}%
\pgfpathlineto{\pgfqpoint{1.531056in}{1.030858in}}%
\pgfpathmoveto{\pgfqpoint{1.531056in}{1.030858in}}%
\pgfpathlineto{\pgfqpoint{1.531056in}{1.030858in}}%
\pgfpathlineto{\pgfqpoint{1.531056in}{1.042655in}}%
\pgfpathlineto{\pgfqpoint{1.549220in}{1.042655in}}%
\pgfpathlineto{\pgfqpoint{1.549220in}{1.030858in}}%
\pgfpathmoveto{\pgfqpoint{1.549220in}{1.042655in}}%
\pgfpathlineto{\pgfqpoint{1.549220in}{1.042655in}}%
\pgfpathlineto{\pgfqpoint{1.549220in}{1.054452in}}%
\pgfpathlineto{\pgfqpoint{1.567384in}{1.054452in}}%
\pgfpathlineto{\pgfqpoint{1.567384in}{1.042655in}}%
\pgfpathmoveto{\pgfqpoint{1.549220in}{1.054452in}}%
\pgfpathlineto{\pgfqpoint{1.549220in}{1.054452in}}%
\pgfpathlineto{\pgfqpoint{1.549220in}{1.066249in}}%
\pgfpathlineto{\pgfqpoint{1.567384in}{1.066249in}}%
\pgfpathlineto{\pgfqpoint{1.567384in}{1.054452in}}%
\pgfpathmoveto{\pgfqpoint{1.567384in}{1.054452in}}%
\pgfpathlineto{\pgfqpoint{1.567384in}{1.054452in}}%
\pgfpathlineto{\pgfqpoint{1.567384in}{1.066249in}}%
\pgfpathlineto{\pgfqpoint{1.585549in}{1.066249in}}%
\pgfpathlineto{\pgfqpoint{1.585549in}{1.054452in}}%
\pgfpathmoveto{\pgfqpoint{1.585549in}{1.066249in}}%
\pgfpathlineto{\pgfqpoint{1.585549in}{1.066249in}}%
\pgfpathlineto{\pgfqpoint{1.585549in}{1.078046in}}%
\pgfpathlineto{\pgfqpoint{1.603713in}{1.078046in}}%
\pgfpathlineto{\pgfqpoint{1.603713in}{1.066249in}}%
\pgfpathmoveto{\pgfqpoint{1.585549in}{1.078046in}}%
\pgfpathlineto{\pgfqpoint{1.585549in}{1.078046in}}%
\pgfpathlineto{\pgfqpoint{1.585549in}{1.089843in}}%
\pgfpathlineto{\pgfqpoint{1.603713in}{1.089843in}}%
\pgfpathlineto{\pgfqpoint{1.603713in}{1.078046in}}%
\pgfpathmoveto{\pgfqpoint{1.603713in}{1.078046in}}%
\pgfpathlineto{\pgfqpoint{1.603713in}{1.078046in}}%
\pgfpathlineto{\pgfqpoint{1.603713in}{1.089843in}}%
\pgfpathlineto{\pgfqpoint{1.621877in}{1.089843in}}%
\pgfpathlineto{\pgfqpoint{1.621877in}{1.078046in}}%
\pgfpathmoveto{\pgfqpoint{1.621877in}{1.089843in}}%
\pgfpathlineto{\pgfqpoint{1.621877in}{1.089843in}}%
\pgfpathlineto{\pgfqpoint{1.621877in}{1.101640in}}%
\pgfpathlineto{\pgfqpoint{1.640040in}{1.101640in}}%
\pgfpathlineto{\pgfqpoint{1.640040in}{1.089843in}}%
\pgfpathmoveto{\pgfqpoint{1.621877in}{1.101640in}}%
\pgfpathlineto{\pgfqpoint{1.621877in}{1.101640in}}%
\pgfpathlineto{\pgfqpoint{1.621877in}{1.113437in}}%
\pgfpathlineto{\pgfqpoint{1.640040in}{1.113437in}}%
\pgfpathlineto{\pgfqpoint{1.640040in}{1.101640in}}%
\pgfpathmoveto{\pgfqpoint{1.640040in}{1.101640in}}%
\pgfpathlineto{\pgfqpoint{1.640040in}{1.101640in}}%
\pgfpathlineto{\pgfqpoint{1.640040in}{1.113437in}}%
\pgfpathlineto{\pgfqpoint{1.658204in}{1.113437in}}%
\pgfpathlineto{\pgfqpoint{1.658204in}{1.101640in}}%
\pgfpathmoveto{\pgfqpoint{1.658204in}{1.113437in}}%
\pgfpathlineto{\pgfqpoint{1.658204in}{1.113437in}}%
\pgfpathlineto{\pgfqpoint{1.658204in}{1.125234in}}%
\pgfpathlineto{\pgfqpoint{1.676368in}{1.125234in}}%
\pgfpathlineto{\pgfqpoint{1.676368in}{1.113437in}}%
\pgfpathmoveto{\pgfqpoint{1.658204in}{1.125234in}}%
\pgfpathlineto{\pgfqpoint{1.658204in}{1.125234in}}%
\pgfpathlineto{\pgfqpoint{1.658204in}{1.137031in}}%
\pgfpathlineto{\pgfqpoint{1.676368in}{1.137031in}}%
\pgfpathlineto{\pgfqpoint{1.676368in}{1.125234in}}%
\pgfpathmoveto{\pgfqpoint{1.676368in}{1.125234in}}%
\pgfpathlineto{\pgfqpoint{1.676368in}{1.125234in}}%
\pgfpathlineto{\pgfqpoint{1.676368in}{1.137031in}}%
\pgfpathlineto{\pgfqpoint{1.694531in}{1.137031in}}%
\pgfpathlineto{\pgfqpoint{1.694531in}{1.125234in}}%
\pgfpathmoveto{\pgfqpoint{1.730859in}{1.148828in}}%
\pgfpathlineto{\pgfqpoint{1.730859in}{1.148828in}}%
\pgfpathlineto{\pgfqpoint{1.730859in}{1.160625in}}%
\pgfpathlineto{\pgfqpoint{1.749022in}{1.160625in}}%
\pgfpathlineto{\pgfqpoint{1.749022in}{1.148828in}}%
\pgfpathmoveto{\pgfqpoint{1.767186in}{1.172422in}}%
\pgfpathlineto{\pgfqpoint{1.767186in}{1.172422in}}%
\pgfpathlineto{\pgfqpoint{1.767186in}{1.184218in}}%
\pgfpathlineto{\pgfqpoint{1.785350in}{1.184218in}}%
\pgfpathlineto{\pgfqpoint{1.785350in}{1.172422in}}%
\pgfpathmoveto{\pgfqpoint{1.803513in}{1.196015in}}%
\pgfpathlineto{\pgfqpoint{1.803513in}{1.196015in}}%
\pgfpathlineto{\pgfqpoint{1.803513in}{1.207812in}}%
\pgfpathlineto{\pgfqpoint{1.821677in}{1.207812in}}%
\pgfpathlineto{\pgfqpoint{1.821677in}{1.196015in}}%
\pgfpathmoveto{\pgfqpoint{1.839841in}{1.219608in}}%
\pgfpathlineto{\pgfqpoint{1.839841in}{1.219608in}}%
\pgfpathlineto{\pgfqpoint{1.839841in}{1.231405in}}%
\pgfpathlineto{\pgfqpoint{1.858005in}{1.231405in}}%
\pgfpathlineto{\pgfqpoint{1.858005in}{1.219608in}}%
\pgfpathmoveto{\pgfqpoint{1.876168in}{1.243201in}}%
\pgfpathlineto{\pgfqpoint{1.876168in}{1.243201in}}%
\pgfpathlineto{\pgfqpoint{1.876168in}{1.254998in}}%
\pgfpathlineto{\pgfqpoint{1.894332in}{1.254998in}}%
\pgfpathlineto{\pgfqpoint{1.894332in}{1.243201in}}%
\pgfpathmoveto{\pgfqpoint{1.912496in}{1.266795in}}%
\pgfpathlineto{\pgfqpoint{1.912496in}{1.266795in}}%
\pgfpathlineto{\pgfqpoint{1.912496in}{1.278592in}}%
\pgfpathlineto{\pgfqpoint{1.930661in}{1.278592in}}%
\pgfpathlineto{\pgfqpoint{1.930661in}{1.266795in}}%
\pgfpathmoveto{\pgfqpoint{1.948825in}{1.290390in}}%
\pgfpathlineto{\pgfqpoint{1.948825in}{1.290390in}}%
\pgfpathlineto{\pgfqpoint{1.948825in}{1.302187in}}%
\pgfpathlineto{\pgfqpoint{1.966990in}{1.302187in}}%
\pgfpathlineto{\pgfqpoint{1.966990in}{1.290390in}}%
\pgfpathmoveto{\pgfqpoint{1.985155in}{1.313984in}}%
\pgfpathlineto{\pgfqpoint{1.985155in}{1.313984in}}%
\pgfpathlineto{\pgfqpoint{1.985155in}{1.325781in}}%
\pgfpathlineto{\pgfqpoint{2.003319in}{1.325781in}}%
\pgfpathlineto{\pgfqpoint{2.003319in}{1.313984in}}%
\pgfpathmoveto{\pgfqpoint{2.021484in}{1.337579in}}%
\pgfpathlineto{\pgfqpoint{2.021484in}{1.337579in}}%
\pgfpathlineto{\pgfqpoint{2.021484in}{1.349376in}}%
\pgfpathlineto{\pgfqpoint{2.039649in}{1.349376in}}%
\pgfpathlineto{\pgfqpoint{2.039649in}{1.337579in}}%
\pgfpathmoveto{\pgfqpoint{2.057813in}{1.372969in}}%
\pgfpathlineto{\pgfqpoint{2.057813in}{1.372969in}}%
\pgfpathlineto{\pgfqpoint{2.057813in}{1.384766in}}%
\pgfpathlineto{\pgfqpoint{2.075978in}{1.384766in}}%
\pgfpathlineto{\pgfqpoint{2.075978in}{1.372969in}}%
\pgfpathmoveto{\pgfqpoint{2.057813in}{1.384766in}}%
\pgfpathlineto{\pgfqpoint{2.057813in}{1.384766in}}%
\pgfpathlineto{\pgfqpoint{2.057813in}{1.396563in}}%
\pgfpathlineto{\pgfqpoint{2.075978in}{1.396563in}}%
\pgfpathlineto{\pgfqpoint{2.075978in}{1.384766in}}%
\pgfpathmoveto{\pgfqpoint{2.075978in}{1.384766in}}%
\pgfpathlineto{\pgfqpoint{2.075978in}{1.384766in}}%
\pgfpathlineto{\pgfqpoint{2.075978in}{1.396563in}}%
\pgfpathlineto{\pgfqpoint{2.094142in}{1.396563in}}%
\pgfpathlineto{\pgfqpoint{2.094142in}{1.384766in}}%
\pgfpathmoveto{\pgfqpoint{2.094142in}{1.396563in}}%
\pgfpathlineto{\pgfqpoint{2.094142in}{1.396563in}}%
\pgfpathlineto{\pgfqpoint{2.094142in}{1.408360in}}%
\pgfpathlineto{\pgfqpoint{2.112307in}{1.408360in}}%
\pgfpathlineto{\pgfqpoint{2.112307in}{1.396563in}}%
\pgfpathmoveto{\pgfqpoint{2.094142in}{1.408360in}}%
\pgfpathlineto{\pgfqpoint{2.094142in}{1.408360in}}%
\pgfpathlineto{\pgfqpoint{2.094142in}{1.420156in}}%
\pgfpathlineto{\pgfqpoint{2.112307in}{1.420156in}}%
\pgfpathlineto{\pgfqpoint{2.112307in}{1.408360in}}%
\pgfpathmoveto{\pgfqpoint{2.112307in}{1.408360in}}%
\pgfpathlineto{\pgfqpoint{2.112307in}{1.408360in}}%
\pgfpathlineto{\pgfqpoint{2.112307in}{1.420156in}}%
\pgfpathlineto{\pgfqpoint{2.130471in}{1.420156in}}%
\pgfpathlineto{\pgfqpoint{2.130471in}{1.408360in}}%
\pgfpathmoveto{\pgfqpoint{2.130471in}{1.420156in}}%
\pgfpathlineto{\pgfqpoint{2.130471in}{1.420156in}}%
\pgfpathlineto{\pgfqpoint{2.130471in}{1.431953in}}%
\pgfpathlineto{\pgfqpoint{2.148636in}{1.431953in}}%
\pgfpathlineto{\pgfqpoint{2.148636in}{1.420156in}}%
\pgfpathmoveto{\pgfqpoint{2.130471in}{1.431953in}}%
\pgfpathlineto{\pgfqpoint{2.130471in}{1.431953in}}%
\pgfpathlineto{\pgfqpoint{2.130471in}{1.443750in}}%
\pgfpathlineto{\pgfqpoint{2.148636in}{1.443750in}}%
\pgfpathlineto{\pgfqpoint{2.148636in}{1.431953in}}%
\pgfpathmoveto{\pgfqpoint{2.148636in}{1.431953in}}%
\pgfpathlineto{\pgfqpoint{2.148636in}{1.431953in}}%
\pgfpathlineto{\pgfqpoint{2.148636in}{1.443750in}}%
\pgfpathlineto{\pgfqpoint{2.166800in}{1.443750in}}%
\pgfpathlineto{\pgfqpoint{2.166800in}{1.431953in}}%
\pgfpathmoveto{\pgfqpoint{2.166800in}{1.443750in}}%
\pgfpathlineto{\pgfqpoint{2.166800in}{1.443750in}}%
\pgfpathlineto{\pgfqpoint{2.166800in}{1.455547in}}%
\pgfpathlineto{\pgfqpoint{2.184965in}{1.455547in}}%
\pgfpathlineto{\pgfqpoint{2.184965in}{1.443750in}}%
\pgfpathmoveto{\pgfqpoint{2.166800in}{1.455547in}}%
\pgfpathlineto{\pgfqpoint{2.166800in}{1.455547in}}%
\pgfpathlineto{\pgfqpoint{2.166800in}{1.467344in}}%
\pgfpathlineto{\pgfqpoint{2.184965in}{1.467344in}}%
\pgfpathlineto{\pgfqpoint{2.184965in}{1.455547in}}%
\pgfpathmoveto{\pgfqpoint{2.184965in}{1.455547in}}%
\pgfpathlineto{\pgfqpoint{2.184965in}{1.455547in}}%
\pgfpathlineto{\pgfqpoint{2.184965in}{1.467344in}}%
\pgfpathlineto{\pgfqpoint{2.203129in}{1.467344in}}%
\pgfpathlineto{\pgfqpoint{2.203129in}{1.455547in}}%
\pgfpathmoveto{\pgfqpoint{2.203129in}{1.467344in}}%
\pgfpathlineto{\pgfqpoint{2.203129in}{1.467344in}}%
\pgfpathlineto{\pgfqpoint{2.203129in}{1.479141in}}%
\pgfpathlineto{\pgfqpoint{2.221293in}{1.479141in}}%
\pgfpathlineto{\pgfqpoint{2.221293in}{1.467344in}}%
\pgfpathmoveto{\pgfqpoint{2.203129in}{1.479141in}}%
\pgfpathlineto{\pgfqpoint{2.203129in}{1.479141in}}%
\pgfpathlineto{\pgfqpoint{2.203129in}{1.490938in}}%
\pgfpathlineto{\pgfqpoint{2.221293in}{1.490938in}}%
\pgfpathlineto{\pgfqpoint{2.221293in}{1.479141in}}%
\pgfpathmoveto{\pgfqpoint{2.221293in}{1.479141in}}%
\pgfpathlineto{\pgfqpoint{2.221293in}{1.479141in}}%
\pgfpathlineto{\pgfqpoint{2.221293in}{1.490938in}}%
\pgfpathlineto{\pgfqpoint{2.239457in}{1.490938in}}%
\pgfpathlineto{\pgfqpoint{2.239457in}{1.479141in}}%
\pgfpathmoveto{\pgfqpoint{2.239457in}{1.490938in}}%
\pgfpathlineto{\pgfqpoint{2.239457in}{1.490938in}}%
\pgfpathlineto{\pgfqpoint{2.239457in}{1.502735in}}%
\pgfpathlineto{\pgfqpoint{2.257621in}{1.502735in}}%
\pgfpathlineto{\pgfqpoint{2.257621in}{1.490938in}}%
\pgfpathmoveto{\pgfqpoint{2.239457in}{1.502735in}}%
\pgfpathlineto{\pgfqpoint{2.239457in}{1.502735in}}%
\pgfpathlineto{\pgfqpoint{2.239457in}{1.514533in}}%
\pgfpathlineto{\pgfqpoint{2.257621in}{1.514533in}}%
\pgfpathlineto{\pgfqpoint{2.257621in}{1.502735in}}%
\pgfpathmoveto{\pgfqpoint{2.257621in}{1.502735in}}%
\pgfpathlineto{\pgfqpoint{2.257621in}{1.502735in}}%
\pgfpathlineto{\pgfqpoint{2.257621in}{1.514533in}}%
\pgfpathlineto{\pgfqpoint{2.275785in}{1.514533in}}%
\pgfpathlineto{\pgfqpoint{2.275785in}{1.502735in}}%
\pgfpathmoveto{\pgfqpoint{2.275785in}{1.514533in}}%
\pgfpathlineto{\pgfqpoint{2.275785in}{1.514533in}}%
\pgfpathlineto{\pgfqpoint{2.275785in}{1.526330in}}%
\pgfpathlineto{\pgfqpoint{2.293949in}{1.526330in}}%
\pgfpathlineto{\pgfqpoint{2.293949in}{1.514533in}}%
\pgfpathmoveto{\pgfqpoint{2.275785in}{1.526330in}}%
\pgfpathlineto{\pgfqpoint{2.275785in}{1.526330in}}%
\pgfpathlineto{\pgfqpoint{2.275785in}{1.538127in}}%
\pgfpathlineto{\pgfqpoint{2.293949in}{1.538127in}}%
\pgfpathlineto{\pgfqpoint{2.293949in}{1.526330in}}%
\pgfpathmoveto{\pgfqpoint{2.293949in}{1.526330in}}%
\pgfpathlineto{\pgfqpoint{2.293949in}{1.526330in}}%
\pgfpathlineto{\pgfqpoint{2.293949in}{1.538127in}}%
\pgfpathlineto{\pgfqpoint{2.312113in}{1.538127in}}%
\pgfpathlineto{\pgfqpoint{2.312113in}{1.526330in}}%
\pgfpathmoveto{\pgfqpoint{2.312113in}{1.538127in}}%
\pgfpathlineto{\pgfqpoint{2.312113in}{1.538127in}}%
\pgfpathlineto{\pgfqpoint{2.312113in}{1.549924in}}%
\pgfpathlineto{\pgfqpoint{2.330277in}{1.549924in}}%
\pgfpathlineto{\pgfqpoint{2.330277in}{1.538127in}}%
\pgfpathmoveto{\pgfqpoint{2.312113in}{1.549924in}}%
\pgfpathlineto{\pgfqpoint{2.312113in}{1.549924in}}%
\pgfpathlineto{\pgfqpoint{2.312113in}{1.561720in}}%
\pgfpathlineto{\pgfqpoint{2.330277in}{1.561720in}}%
\pgfpathlineto{\pgfqpoint{2.330277in}{1.549924in}}%
\pgfpathmoveto{\pgfqpoint{2.330277in}{1.549924in}}%
\pgfpathlineto{\pgfqpoint{2.330277in}{1.549924in}}%
\pgfpathlineto{\pgfqpoint{2.330277in}{1.561720in}}%
\pgfpathlineto{\pgfqpoint{2.348441in}{1.561720in}}%
\pgfpathlineto{\pgfqpoint{2.348441in}{1.549924in}}%
\pgfpathmoveto{\pgfqpoint{2.384768in}{1.573517in}}%
\pgfpathlineto{\pgfqpoint{2.384768in}{1.573517in}}%
\pgfpathlineto{\pgfqpoint{2.384768in}{1.585314in}}%
\pgfpathlineto{\pgfqpoint{2.402931in}{1.585314in}}%
\pgfpathlineto{\pgfqpoint{2.402931in}{1.573517in}}%
\pgfpathmoveto{\pgfqpoint{2.421095in}{1.597111in}}%
\pgfpathlineto{\pgfqpoint{2.421095in}{1.597111in}}%
\pgfpathlineto{\pgfqpoint{2.421095in}{1.608908in}}%
\pgfpathlineto{\pgfqpoint{2.439259in}{1.608908in}}%
\pgfpathlineto{\pgfqpoint{2.439259in}{1.597111in}}%
\pgfpathmoveto{\pgfqpoint{2.457422in}{1.620704in}}%
\pgfpathlineto{\pgfqpoint{2.457422in}{1.620704in}}%
\pgfpathlineto{\pgfqpoint{2.457422in}{1.632501in}}%
\pgfpathlineto{\pgfqpoint{2.475586in}{1.632501in}}%
\pgfpathlineto{\pgfqpoint{2.475586in}{1.620704in}}%
\pgfpathmoveto{\pgfqpoint{2.493750in}{1.644298in}}%
\pgfpathlineto{\pgfqpoint{2.493750in}{1.644298in}}%
\pgfpathlineto{\pgfqpoint{2.493750in}{1.656095in}}%
\pgfpathlineto{\pgfqpoint{2.511914in}{1.656095in}}%
\pgfpathlineto{\pgfqpoint{2.511914in}{1.644298in}}%
\pgfpathmoveto{\pgfqpoint{2.530078in}{1.667892in}}%
\pgfpathlineto{\pgfqpoint{2.530078in}{1.667892in}}%
\pgfpathlineto{\pgfqpoint{2.530078in}{1.679689in}}%
\pgfpathlineto{\pgfqpoint{2.548242in}{1.679689in}}%
\pgfpathlineto{\pgfqpoint{2.548242in}{1.667892in}}%
\pgfpathmoveto{\pgfqpoint{2.566406in}{1.691485in}}%
\pgfpathlineto{\pgfqpoint{2.566406in}{1.691485in}}%
\pgfpathlineto{\pgfqpoint{2.566406in}{1.703282in}}%
\pgfpathlineto{\pgfqpoint{2.584571in}{1.703282in}}%
\pgfpathlineto{\pgfqpoint{2.584571in}{1.691485in}}%
\pgfpathmoveto{\pgfqpoint{2.602735in}{1.715079in}}%
\pgfpathlineto{\pgfqpoint{2.602735in}{1.715079in}}%
\pgfpathlineto{\pgfqpoint{2.602735in}{1.726876in}}%
\pgfpathlineto{\pgfqpoint{2.620899in}{1.726876in}}%
\pgfpathlineto{\pgfqpoint{2.620899in}{1.715079in}}%
\pgfpathmoveto{\pgfqpoint{2.639063in}{1.738672in}}%
\pgfpathlineto{\pgfqpoint{2.639063in}{1.738672in}}%
\pgfpathlineto{\pgfqpoint{2.639063in}{1.750469in}}%
\pgfpathlineto{\pgfqpoint{2.657227in}{1.750469in}}%
\pgfpathlineto{\pgfqpoint{2.657227in}{1.738672in}}%
\pgfpathmoveto{\pgfqpoint{2.639063in}{1.750469in}}%
\pgfpathlineto{\pgfqpoint{2.639063in}{1.750469in}}%
\pgfpathlineto{\pgfqpoint{2.639063in}{1.762266in}}%
\pgfpathlineto{\pgfqpoint{2.657227in}{1.762266in}}%
\pgfpathlineto{\pgfqpoint{2.657227in}{1.750469in}}%
\pgfpathmoveto{\pgfqpoint{2.639063in}{1.762266in}}%
\pgfpathlineto{\pgfqpoint{2.639063in}{1.762266in}}%
\pgfpathlineto{\pgfqpoint{2.639063in}{1.774062in}}%
\pgfpathlineto{\pgfqpoint{2.657227in}{1.774062in}}%
\pgfpathlineto{\pgfqpoint{2.657227in}{1.762266in}}%
\pgfpathmoveto{\pgfqpoint{2.657227in}{1.762266in}}%
\pgfpathlineto{\pgfqpoint{2.657227in}{1.762266in}}%
\pgfpathlineto{\pgfqpoint{2.657227in}{1.774062in}}%
\pgfpathlineto{\pgfqpoint{2.675391in}{1.774062in}}%
\pgfpathlineto{\pgfqpoint{2.675391in}{1.762266in}}%
\pgfpathmoveto{\pgfqpoint{2.675391in}{1.774062in}}%
\pgfpathlineto{\pgfqpoint{2.675391in}{1.774062in}}%
\pgfpathlineto{\pgfqpoint{2.675391in}{1.785859in}}%
\pgfpathlineto{\pgfqpoint{2.693555in}{1.785859in}}%
\pgfpathlineto{\pgfqpoint{2.693555in}{1.774062in}}%
\pgfpathmoveto{\pgfqpoint{2.675391in}{1.785859in}}%
\pgfpathlineto{\pgfqpoint{2.675391in}{1.785859in}}%
\pgfpathlineto{\pgfqpoint{2.675391in}{1.797655in}}%
\pgfpathlineto{\pgfqpoint{2.693555in}{1.797655in}}%
\pgfpathlineto{\pgfqpoint{2.693555in}{1.785859in}}%
\pgfpathmoveto{\pgfqpoint{2.693555in}{1.785859in}}%
\pgfpathlineto{\pgfqpoint{2.693555in}{1.785859in}}%
\pgfpathlineto{\pgfqpoint{2.693555in}{1.797655in}}%
\pgfpathlineto{\pgfqpoint{2.711719in}{1.797655in}}%
\pgfpathlineto{\pgfqpoint{2.711719in}{1.785859in}}%
\pgfpathmoveto{\pgfqpoint{2.711719in}{1.797655in}}%
\pgfpathlineto{\pgfqpoint{2.711719in}{1.797655in}}%
\pgfpathlineto{\pgfqpoint{2.711719in}{1.809452in}}%
\pgfpathlineto{\pgfqpoint{2.729883in}{1.809452in}}%
\pgfpathlineto{\pgfqpoint{2.729883in}{1.797655in}}%
\pgfpathmoveto{\pgfqpoint{2.711719in}{1.809452in}}%
\pgfpathlineto{\pgfqpoint{2.711719in}{1.809452in}}%
\pgfpathlineto{\pgfqpoint{2.711719in}{1.821249in}}%
\pgfpathlineto{\pgfqpoint{2.729883in}{1.821249in}}%
\pgfpathlineto{\pgfqpoint{2.729883in}{1.809452in}}%
\pgfpathmoveto{\pgfqpoint{2.729883in}{1.809452in}}%
\pgfpathlineto{\pgfqpoint{2.729883in}{1.809452in}}%
\pgfpathlineto{\pgfqpoint{2.729883in}{1.821249in}}%
\pgfpathlineto{\pgfqpoint{2.748047in}{1.821249in}}%
\pgfpathlineto{\pgfqpoint{2.748047in}{1.809452in}}%
\pgfpathmoveto{\pgfqpoint{2.748047in}{1.821249in}}%
\pgfpathlineto{\pgfqpoint{2.748047in}{1.821249in}}%
\pgfpathlineto{\pgfqpoint{2.748047in}{1.833046in}}%
\pgfpathlineto{\pgfqpoint{2.766211in}{1.833046in}}%
\pgfpathlineto{\pgfqpoint{2.766211in}{1.821249in}}%
\pgfpathmoveto{\pgfqpoint{2.748047in}{1.833046in}}%
\pgfpathlineto{\pgfqpoint{2.748047in}{1.833046in}}%
\pgfpathlineto{\pgfqpoint{2.748047in}{1.844843in}}%
\pgfpathlineto{\pgfqpoint{2.766211in}{1.844843in}}%
\pgfpathlineto{\pgfqpoint{2.766211in}{1.833046in}}%
\pgfpathmoveto{\pgfqpoint{2.766211in}{1.833046in}}%
\pgfpathlineto{\pgfqpoint{2.766211in}{1.833046in}}%
\pgfpathlineto{\pgfqpoint{2.766211in}{1.844843in}}%
\pgfpathlineto{\pgfqpoint{2.784375in}{1.844843in}}%
\pgfpathlineto{\pgfqpoint{2.784375in}{1.833046in}}%
\pgfpathmoveto{\pgfqpoint{2.784375in}{1.844843in}}%
\pgfpathlineto{\pgfqpoint{2.784375in}{1.844843in}}%
\pgfpathlineto{\pgfqpoint{2.784375in}{1.856640in}}%
\pgfpathlineto{\pgfqpoint{2.802539in}{1.856640in}}%
\pgfpathlineto{\pgfqpoint{2.802539in}{1.844843in}}%
\pgfpathmoveto{\pgfqpoint{2.784375in}{1.856640in}}%
\pgfpathlineto{\pgfqpoint{2.784375in}{1.856640in}}%
\pgfpathlineto{\pgfqpoint{2.784375in}{1.868438in}}%
\pgfpathlineto{\pgfqpoint{2.802539in}{1.868438in}}%
\pgfpathlineto{\pgfqpoint{2.802539in}{1.856640in}}%
\pgfpathmoveto{\pgfqpoint{2.802539in}{1.856640in}}%
\pgfpathlineto{\pgfqpoint{2.802539in}{1.856640in}}%
\pgfpathlineto{\pgfqpoint{2.802539in}{1.868438in}}%
\pgfpathlineto{\pgfqpoint{2.820703in}{1.868438in}}%
\pgfpathlineto{\pgfqpoint{2.820703in}{1.856640in}}%
\pgfpathmoveto{\pgfqpoint{2.857031in}{1.880235in}}%
\pgfpathlineto{\pgfqpoint{2.857031in}{1.880235in}}%
\pgfpathlineto{\pgfqpoint{2.857031in}{1.892032in}}%
\pgfpathlineto{\pgfqpoint{2.875196in}{1.892032in}}%
\pgfpathlineto{\pgfqpoint{2.875196in}{1.880235in}}%
\pgfpathmoveto{\pgfqpoint{2.893360in}{1.903829in}}%
\pgfpathlineto{\pgfqpoint{2.893360in}{1.903829in}}%
\pgfpathlineto{\pgfqpoint{2.893360in}{1.915627in}}%
\pgfpathlineto{\pgfqpoint{2.911524in}{1.915627in}}%
\pgfpathlineto{\pgfqpoint{2.911524in}{1.903829in}}%
\pgfpathmoveto{\pgfqpoint{2.893360in}{3.402033in}}%
\pgfpathlineto{\pgfqpoint{2.893360in}{3.402033in}}%
\pgfpathlineto{\pgfqpoint{2.893360in}{3.413830in}}%
\pgfpathlineto{\pgfqpoint{2.911524in}{3.413830in}}%
\pgfpathlineto{\pgfqpoint{2.911524in}{3.402033in}}%
\pgfpathmoveto{\pgfqpoint{2.893360in}{3.413830in}}%
\pgfpathlineto{\pgfqpoint{2.893360in}{3.413830in}}%
\pgfpathlineto{\pgfqpoint{2.893360in}{3.425627in}}%
\pgfpathlineto{\pgfqpoint{2.911524in}{3.425627in}}%
\pgfpathlineto{\pgfqpoint{2.911524in}{3.413830in}}%
\pgfpathmoveto{\pgfqpoint{2.820703in}{3.496405in}}%
\pgfpathlineto{\pgfqpoint{2.820703in}{3.496405in}}%
\pgfpathlineto{\pgfqpoint{2.820703in}{3.508202in}}%
\pgfpathlineto{\pgfqpoint{2.838867in}{3.508202in}}%
\pgfpathlineto{\pgfqpoint{2.838867in}{3.496405in}}%
\pgfpathmoveto{\pgfqpoint{2.820703in}{3.508202in}}%
\pgfpathlineto{\pgfqpoint{2.820703in}{3.508202in}}%
\pgfpathlineto{\pgfqpoint{2.820703in}{3.519998in}}%
\pgfpathlineto{\pgfqpoint{2.838867in}{3.519998in}}%
\pgfpathlineto{\pgfqpoint{2.838867in}{3.508202in}}%
\pgfpathmoveto{\pgfqpoint{2.857031in}{3.449220in}}%
\pgfpathlineto{\pgfqpoint{2.857031in}{3.449220in}}%
\pgfpathlineto{\pgfqpoint{2.857031in}{3.461016in}}%
\pgfpathlineto{\pgfqpoint{2.875196in}{3.461016in}}%
\pgfpathlineto{\pgfqpoint{2.875196in}{3.449220in}}%
\pgfpathmoveto{\pgfqpoint{2.857031in}{3.461016in}}%
\pgfpathlineto{\pgfqpoint{2.857031in}{3.461016in}}%
\pgfpathlineto{\pgfqpoint{2.857031in}{3.472812in}}%
\pgfpathlineto{\pgfqpoint{2.875196in}{3.472812in}}%
\pgfpathlineto{\pgfqpoint{2.875196in}{3.461016in}}%
\pgfpathmoveto{\pgfqpoint{2.929688in}{1.927423in}}%
\pgfpathlineto{\pgfqpoint{2.929688in}{1.927423in}}%
\pgfpathlineto{\pgfqpoint{2.929688in}{1.939220in}}%
\pgfpathlineto{\pgfqpoint{2.947852in}{1.939220in}}%
\pgfpathlineto{\pgfqpoint{2.947852in}{1.927423in}}%
\pgfpathmoveto{\pgfqpoint{2.966017in}{1.951017in}}%
\pgfpathlineto{\pgfqpoint{2.966017in}{1.951017in}}%
\pgfpathlineto{\pgfqpoint{2.966017in}{1.962813in}}%
\pgfpathlineto{\pgfqpoint{2.984181in}{1.962813in}}%
\pgfpathlineto{\pgfqpoint{2.984181in}{1.951017in}}%
\pgfpathmoveto{\pgfqpoint{2.966017in}{1.962813in}}%
\pgfpathlineto{\pgfqpoint{2.966017in}{1.962813in}}%
\pgfpathlineto{\pgfqpoint{2.966017in}{1.974610in}}%
\pgfpathlineto{\pgfqpoint{2.984181in}{1.974610in}}%
\pgfpathlineto{\pgfqpoint{2.984181in}{1.962813in}}%
\pgfpathmoveto{\pgfqpoint{2.966017in}{1.974610in}}%
\pgfpathlineto{\pgfqpoint{2.966017in}{1.974610in}}%
\pgfpathlineto{\pgfqpoint{2.966017in}{1.986407in}}%
\pgfpathlineto{\pgfqpoint{2.984181in}{1.986407in}}%
\pgfpathlineto{\pgfqpoint{2.984181in}{1.974610in}}%
\pgfpathmoveto{\pgfqpoint{2.984181in}{1.974610in}}%
\pgfpathlineto{\pgfqpoint{2.984181in}{1.974610in}}%
\pgfpathlineto{\pgfqpoint{2.984181in}{1.986407in}}%
\pgfpathlineto{\pgfqpoint{3.002346in}{1.986407in}}%
\pgfpathlineto{\pgfqpoint{3.002346in}{1.974610in}}%
\pgfpathmoveto{\pgfqpoint{3.002346in}{1.986407in}}%
\pgfpathlineto{\pgfqpoint{3.002346in}{1.986407in}}%
\pgfpathlineto{\pgfqpoint{3.002346in}{1.998204in}}%
\pgfpathlineto{\pgfqpoint{3.020510in}{1.998204in}}%
\pgfpathlineto{\pgfqpoint{3.020510in}{1.986407in}}%
\pgfpathmoveto{\pgfqpoint{3.002346in}{1.998204in}}%
\pgfpathlineto{\pgfqpoint{3.002346in}{1.998204in}}%
\pgfpathlineto{\pgfqpoint{3.002346in}{2.010000in}}%
\pgfpathlineto{\pgfqpoint{3.020510in}{2.010000in}}%
\pgfpathlineto{\pgfqpoint{3.020510in}{1.998204in}}%
\pgfpathmoveto{\pgfqpoint{3.020510in}{1.998204in}}%
\pgfpathlineto{\pgfqpoint{3.020510in}{1.998204in}}%
\pgfpathlineto{\pgfqpoint{3.020510in}{2.010000in}}%
\pgfpathlineto{\pgfqpoint{3.038675in}{2.010000in}}%
\pgfpathlineto{\pgfqpoint{3.038675in}{1.998204in}}%
\pgfpathmoveto{\pgfqpoint{3.038675in}{2.010000in}}%
\pgfpathlineto{\pgfqpoint{3.038675in}{2.010000in}}%
\pgfpathlineto{\pgfqpoint{3.038675in}{2.021797in}}%
\pgfpathlineto{\pgfqpoint{3.056839in}{2.021797in}}%
\pgfpathlineto{\pgfqpoint{3.056839in}{2.010000in}}%
\pgfpathmoveto{\pgfqpoint{3.038675in}{2.021797in}}%
\pgfpathlineto{\pgfqpoint{3.038675in}{2.021797in}}%
\pgfpathlineto{\pgfqpoint{3.038675in}{2.033594in}}%
\pgfpathlineto{\pgfqpoint{3.056839in}{2.033594in}}%
\pgfpathlineto{\pgfqpoint{3.056839in}{2.021797in}}%
\pgfpathmoveto{\pgfqpoint{3.056839in}{2.021797in}}%
\pgfpathlineto{\pgfqpoint{3.056839in}{2.021797in}}%
\pgfpathlineto{\pgfqpoint{3.056839in}{2.033594in}}%
\pgfpathlineto{\pgfqpoint{3.075004in}{2.033594in}}%
\pgfpathlineto{\pgfqpoint{3.075004in}{2.021797in}}%
\pgfpathmoveto{\pgfqpoint{3.038675in}{3.213281in}}%
\pgfpathlineto{\pgfqpoint{3.038675in}{3.213281in}}%
\pgfpathlineto{\pgfqpoint{3.038675in}{3.225078in}}%
\pgfpathlineto{\pgfqpoint{3.056839in}{3.225078in}}%
\pgfpathlineto{\pgfqpoint{3.056839in}{3.213281in}}%
\pgfpathmoveto{\pgfqpoint{3.038675in}{3.225078in}}%
\pgfpathlineto{\pgfqpoint{3.038675in}{3.225078in}}%
\pgfpathlineto{\pgfqpoint{3.038675in}{3.236874in}}%
\pgfpathlineto{\pgfqpoint{3.056839in}{3.236874in}}%
\pgfpathlineto{\pgfqpoint{3.056839in}{3.225078in}}%
\pgfpathmoveto{\pgfqpoint{2.966017in}{3.307657in}}%
\pgfpathlineto{\pgfqpoint{2.966017in}{3.307657in}}%
\pgfpathlineto{\pgfqpoint{2.966017in}{3.319455in}}%
\pgfpathlineto{\pgfqpoint{2.984181in}{3.319455in}}%
\pgfpathlineto{\pgfqpoint{2.984181in}{3.307657in}}%
\pgfpathmoveto{\pgfqpoint{2.966017in}{3.319455in}}%
\pgfpathlineto{\pgfqpoint{2.966017in}{3.319455in}}%
\pgfpathlineto{\pgfqpoint{2.966017in}{3.331252in}}%
\pgfpathlineto{\pgfqpoint{2.984181in}{3.331252in}}%
\pgfpathlineto{\pgfqpoint{2.984181in}{3.319455in}}%
\pgfpathmoveto{\pgfqpoint{3.002346in}{3.260469in}}%
\pgfpathlineto{\pgfqpoint{3.002346in}{3.260469in}}%
\pgfpathlineto{\pgfqpoint{3.002346in}{3.272266in}}%
\pgfpathlineto{\pgfqpoint{3.020510in}{3.272266in}}%
\pgfpathlineto{\pgfqpoint{3.020510in}{3.260469in}}%
\pgfpathmoveto{\pgfqpoint{3.002346in}{3.272266in}}%
\pgfpathlineto{\pgfqpoint{3.002346in}{3.272266in}}%
\pgfpathlineto{\pgfqpoint{3.002346in}{3.284063in}}%
\pgfpathlineto{\pgfqpoint{3.020510in}{3.284063in}}%
\pgfpathlineto{\pgfqpoint{3.020510in}{3.272266in}}%
\pgfpathmoveto{\pgfqpoint{2.929688in}{3.354845in}}%
\pgfpathlineto{\pgfqpoint{2.929688in}{3.354845in}}%
\pgfpathlineto{\pgfqpoint{2.929688in}{3.366642in}}%
\pgfpathlineto{\pgfqpoint{2.947852in}{3.366642in}}%
\pgfpathlineto{\pgfqpoint{2.947852in}{3.354845in}}%
\pgfpathmoveto{\pgfqpoint{2.929688in}{3.366642in}}%
\pgfpathlineto{\pgfqpoint{2.929688in}{3.366642in}}%
\pgfpathlineto{\pgfqpoint{2.929688in}{3.378439in}}%
\pgfpathlineto{\pgfqpoint{2.947852in}{3.378439in}}%
\pgfpathlineto{\pgfqpoint{2.947852in}{3.366642in}}%
\pgfpathmoveto{\pgfqpoint{3.075004in}{2.033594in}}%
\pgfpathlineto{\pgfqpoint{3.075004in}{2.033594in}}%
\pgfpathlineto{\pgfqpoint{3.075004in}{2.045392in}}%
\pgfpathlineto{\pgfqpoint{3.093167in}{2.045392in}}%
\pgfpathlineto{\pgfqpoint{3.093167in}{2.033594in}}%
\pgfpathmoveto{\pgfqpoint{3.075004in}{2.045392in}}%
\pgfpathlineto{\pgfqpoint{3.075004in}{2.045392in}}%
\pgfpathlineto{\pgfqpoint{3.075004in}{2.057189in}}%
\pgfpathlineto{\pgfqpoint{3.093167in}{2.057189in}}%
\pgfpathlineto{\pgfqpoint{3.093167in}{2.045392in}}%
\pgfpathmoveto{\pgfqpoint{3.093167in}{2.045392in}}%
\pgfpathlineto{\pgfqpoint{3.093167in}{2.045392in}}%
\pgfpathlineto{\pgfqpoint{3.093167in}{2.057189in}}%
\pgfpathlineto{\pgfqpoint{3.111330in}{2.057189in}}%
\pgfpathlineto{\pgfqpoint{3.111330in}{2.045392in}}%
\pgfpathmoveto{\pgfqpoint{3.111330in}{2.045392in}}%
\pgfpathlineto{\pgfqpoint{3.111330in}{2.045392in}}%
\pgfpathlineto{\pgfqpoint{3.111330in}{2.057189in}}%
\pgfpathlineto{\pgfqpoint{3.129493in}{2.057189in}}%
\pgfpathlineto{\pgfqpoint{3.129493in}{2.045392in}}%
\pgfpathmoveto{\pgfqpoint{3.147657in}{2.068986in}}%
\pgfpathlineto{\pgfqpoint{3.147657in}{2.068986in}}%
\pgfpathlineto{\pgfqpoint{3.147657in}{2.080783in}}%
\pgfpathlineto{\pgfqpoint{3.165820in}{2.080783in}}%
\pgfpathlineto{\pgfqpoint{3.165820in}{2.068986in}}%
\pgfpathmoveto{\pgfqpoint{3.183983in}{2.092580in}}%
\pgfpathlineto{\pgfqpoint{3.183983in}{2.092580in}}%
\pgfpathlineto{\pgfqpoint{3.183983in}{2.104377in}}%
\pgfpathlineto{\pgfqpoint{3.202146in}{2.104377in}}%
\pgfpathlineto{\pgfqpoint{3.202146in}{2.092580in}}%
\pgfpathmoveto{\pgfqpoint{3.183983in}{3.024531in}}%
\pgfpathlineto{\pgfqpoint{3.183983in}{3.024531in}}%
\pgfpathlineto{\pgfqpoint{3.183983in}{3.036328in}}%
\pgfpathlineto{\pgfqpoint{3.202146in}{3.036328in}}%
\pgfpathlineto{\pgfqpoint{3.202146in}{3.024531in}}%
\pgfpathmoveto{\pgfqpoint{3.183983in}{3.036328in}}%
\pgfpathlineto{\pgfqpoint{3.183983in}{3.036328in}}%
\pgfpathlineto{\pgfqpoint{3.183983in}{3.048124in}}%
\pgfpathlineto{\pgfqpoint{3.202146in}{3.048124in}}%
\pgfpathlineto{\pgfqpoint{3.202146in}{3.036328in}}%
\pgfpathmoveto{\pgfqpoint{3.111330in}{3.118907in}}%
\pgfpathlineto{\pgfqpoint{3.111330in}{3.118907in}}%
\pgfpathlineto{\pgfqpoint{3.111330in}{3.130704in}}%
\pgfpathlineto{\pgfqpoint{3.129493in}{3.130704in}}%
\pgfpathlineto{\pgfqpoint{3.129493in}{3.118907in}}%
\pgfpathmoveto{\pgfqpoint{3.111330in}{3.130704in}}%
\pgfpathlineto{\pgfqpoint{3.111330in}{3.130704in}}%
\pgfpathlineto{\pgfqpoint{3.111330in}{3.142502in}}%
\pgfpathlineto{\pgfqpoint{3.129493in}{3.142502in}}%
\pgfpathlineto{\pgfqpoint{3.129493in}{3.130704in}}%
\pgfpathmoveto{\pgfqpoint{3.147657in}{3.071719in}}%
\pgfpathlineto{\pgfqpoint{3.147657in}{3.071719in}}%
\pgfpathlineto{\pgfqpoint{3.147657in}{3.083516in}}%
\pgfpathlineto{\pgfqpoint{3.165820in}{3.083516in}}%
\pgfpathlineto{\pgfqpoint{3.165820in}{3.071719in}}%
\pgfpathmoveto{\pgfqpoint{3.147657in}{3.083516in}}%
\pgfpathlineto{\pgfqpoint{3.147657in}{3.083516in}}%
\pgfpathlineto{\pgfqpoint{3.147657in}{3.095313in}}%
\pgfpathlineto{\pgfqpoint{3.165820in}{3.095313in}}%
\pgfpathlineto{\pgfqpoint{3.165820in}{3.083516in}}%
\pgfpathmoveto{\pgfqpoint{3.075004in}{3.166095in}}%
\pgfpathlineto{\pgfqpoint{3.075004in}{3.166095in}}%
\pgfpathlineto{\pgfqpoint{3.075004in}{3.177891in}}%
\pgfpathlineto{\pgfqpoint{3.093167in}{3.177891in}}%
\pgfpathlineto{\pgfqpoint{3.093167in}{3.166095in}}%
\pgfpathmoveto{\pgfqpoint{3.075004in}{3.177891in}}%
\pgfpathlineto{\pgfqpoint{3.075004in}{3.177891in}}%
\pgfpathlineto{\pgfqpoint{3.075004in}{3.189688in}}%
\pgfpathlineto{\pgfqpoint{3.093167in}{3.189688in}}%
\pgfpathlineto{\pgfqpoint{3.093167in}{3.177891in}}%
\pgfpathmoveto{\pgfqpoint{3.220310in}{2.116174in}}%
\pgfpathlineto{\pgfqpoint{3.220310in}{2.116174in}}%
\pgfpathlineto{\pgfqpoint{3.220310in}{2.127970in}}%
\pgfpathlineto{\pgfqpoint{3.238474in}{2.127970in}}%
\pgfpathlineto{\pgfqpoint{3.238474in}{2.116174in}}%
\pgfpathmoveto{\pgfqpoint{3.256639in}{2.139766in}}%
\pgfpathlineto{\pgfqpoint{3.256639in}{2.139766in}}%
\pgfpathlineto{\pgfqpoint{3.256639in}{2.151563in}}%
\pgfpathlineto{\pgfqpoint{3.274803in}{2.151563in}}%
\pgfpathlineto{\pgfqpoint{3.274803in}{2.139766in}}%
\pgfpathmoveto{\pgfqpoint{3.292968in}{2.175155in}}%
\pgfpathlineto{\pgfqpoint{3.292968in}{2.175155in}}%
\pgfpathlineto{\pgfqpoint{3.292968in}{2.186952in}}%
\pgfpathlineto{\pgfqpoint{3.311133in}{2.186952in}}%
\pgfpathlineto{\pgfqpoint{3.311133in}{2.175155in}}%
\pgfpathmoveto{\pgfqpoint{3.292968in}{2.186952in}}%
\pgfpathlineto{\pgfqpoint{3.292968in}{2.186952in}}%
\pgfpathlineto{\pgfqpoint{3.292968in}{2.198748in}}%
\pgfpathlineto{\pgfqpoint{3.311133in}{2.198748in}}%
\pgfpathlineto{\pgfqpoint{3.311133in}{2.186952in}}%
\pgfpathmoveto{\pgfqpoint{3.311133in}{2.186952in}}%
\pgfpathlineto{\pgfqpoint{3.311133in}{2.186952in}}%
\pgfpathlineto{\pgfqpoint{3.311133in}{2.198748in}}%
\pgfpathlineto{\pgfqpoint{3.329297in}{2.198748in}}%
\pgfpathlineto{\pgfqpoint{3.329297in}{2.186952in}}%
\pgfpathmoveto{\pgfqpoint{3.329297in}{2.198748in}}%
\pgfpathlineto{\pgfqpoint{3.329297in}{2.198748in}}%
\pgfpathlineto{\pgfqpoint{3.329297in}{2.210545in}}%
\pgfpathlineto{\pgfqpoint{3.347462in}{2.210545in}}%
\pgfpathlineto{\pgfqpoint{3.347462in}{2.198748in}}%
\pgfpathmoveto{\pgfqpoint{3.329297in}{2.210545in}}%
\pgfpathlineto{\pgfqpoint{3.329297in}{2.210545in}}%
\pgfpathlineto{\pgfqpoint{3.329297in}{2.222342in}}%
\pgfpathlineto{\pgfqpoint{3.347462in}{2.222342in}}%
\pgfpathlineto{\pgfqpoint{3.347462in}{2.210545in}}%
\pgfpathmoveto{\pgfqpoint{3.347462in}{2.210545in}}%
\pgfpathlineto{\pgfqpoint{3.347462in}{2.210545in}}%
\pgfpathlineto{\pgfqpoint{3.347462in}{2.222342in}}%
\pgfpathlineto{\pgfqpoint{3.365627in}{2.222342in}}%
\pgfpathlineto{\pgfqpoint{3.365627in}{2.210545in}}%
\pgfpathmoveto{\pgfqpoint{3.329297in}{2.835780in}}%
\pgfpathlineto{\pgfqpoint{3.329297in}{2.835780in}}%
\pgfpathlineto{\pgfqpoint{3.329297in}{2.847577in}}%
\pgfpathlineto{\pgfqpoint{3.347462in}{2.847577in}}%
\pgfpathlineto{\pgfqpoint{3.347462in}{2.835780in}}%
\pgfpathmoveto{\pgfqpoint{3.329297in}{2.847577in}}%
\pgfpathlineto{\pgfqpoint{3.329297in}{2.847577in}}%
\pgfpathlineto{\pgfqpoint{3.329297in}{2.859374in}}%
\pgfpathlineto{\pgfqpoint{3.347462in}{2.859374in}}%
\pgfpathlineto{\pgfqpoint{3.347462in}{2.847577in}}%
\pgfpathmoveto{\pgfqpoint{3.256639in}{2.930157in}}%
\pgfpathlineto{\pgfqpoint{3.256639in}{2.930157in}}%
\pgfpathlineto{\pgfqpoint{3.256639in}{2.941954in}}%
\pgfpathlineto{\pgfqpoint{3.274803in}{2.941954in}}%
\pgfpathlineto{\pgfqpoint{3.274803in}{2.930157in}}%
\pgfpathmoveto{\pgfqpoint{3.256639in}{2.941954in}}%
\pgfpathlineto{\pgfqpoint{3.256639in}{2.941954in}}%
\pgfpathlineto{\pgfqpoint{3.256639in}{2.953751in}}%
\pgfpathlineto{\pgfqpoint{3.274803in}{2.953751in}}%
\pgfpathlineto{\pgfqpoint{3.274803in}{2.941954in}}%
\pgfpathmoveto{\pgfqpoint{3.292968in}{2.882968in}}%
\pgfpathlineto{\pgfqpoint{3.292968in}{2.882968in}}%
\pgfpathlineto{\pgfqpoint{3.292968in}{2.894766in}}%
\pgfpathlineto{\pgfqpoint{3.311133in}{2.894766in}}%
\pgfpathlineto{\pgfqpoint{3.311133in}{2.882968in}}%
\pgfpathmoveto{\pgfqpoint{3.292968in}{2.894766in}}%
\pgfpathlineto{\pgfqpoint{3.292968in}{2.894766in}}%
\pgfpathlineto{\pgfqpoint{3.292968in}{2.906563in}}%
\pgfpathlineto{\pgfqpoint{3.311133in}{2.906563in}}%
\pgfpathlineto{\pgfqpoint{3.311133in}{2.894766in}}%
\pgfpathmoveto{\pgfqpoint{3.220310in}{2.977345in}}%
\pgfpathlineto{\pgfqpoint{3.220310in}{2.977345in}}%
\pgfpathlineto{\pgfqpoint{3.220310in}{2.989141in}}%
\pgfpathlineto{\pgfqpoint{3.238474in}{2.989141in}}%
\pgfpathlineto{\pgfqpoint{3.238474in}{2.977345in}}%
\pgfpathmoveto{\pgfqpoint{3.220310in}{2.989141in}}%
\pgfpathlineto{\pgfqpoint{3.220310in}{2.989141in}}%
\pgfpathlineto{\pgfqpoint{3.220310in}{3.000938in}}%
\pgfpathlineto{\pgfqpoint{3.238474in}{3.000938in}}%
\pgfpathlineto{\pgfqpoint{3.238474in}{2.989141in}}%
\pgfpathmoveto{\pgfqpoint{3.365627in}{2.222342in}}%
\pgfpathlineto{\pgfqpoint{3.365627in}{2.222342in}}%
\pgfpathlineto{\pgfqpoint{3.365627in}{2.234139in}}%
\pgfpathlineto{\pgfqpoint{3.383791in}{2.234139in}}%
\pgfpathlineto{\pgfqpoint{3.383791in}{2.222342in}}%
\pgfpathmoveto{\pgfqpoint{3.365627in}{2.234139in}}%
\pgfpathlineto{\pgfqpoint{3.365627in}{2.234139in}}%
\pgfpathlineto{\pgfqpoint{3.365627in}{2.245935in}}%
\pgfpathlineto{\pgfqpoint{3.383791in}{2.245935in}}%
\pgfpathlineto{\pgfqpoint{3.383791in}{2.234139in}}%
\pgfpathmoveto{\pgfqpoint{3.383791in}{2.234139in}}%
\pgfpathlineto{\pgfqpoint{3.383791in}{2.234139in}}%
\pgfpathlineto{\pgfqpoint{3.383791in}{2.245935in}}%
\pgfpathlineto{\pgfqpoint{3.401955in}{2.245935in}}%
\pgfpathlineto{\pgfqpoint{3.401955in}{2.234139in}}%
\pgfpathmoveto{\pgfqpoint{3.401955in}{2.245935in}}%
\pgfpathlineto{\pgfqpoint{3.401955in}{2.245935in}}%
\pgfpathlineto{\pgfqpoint{3.401955in}{2.257732in}}%
\pgfpathlineto{\pgfqpoint{3.420119in}{2.257732in}}%
\pgfpathlineto{\pgfqpoint{3.420119in}{2.245935in}}%
\pgfpathmoveto{\pgfqpoint{3.401955in}{2.257732in}}%
\pgfpathlineto{\pgfqpoint{3.401955in}{2.257732in}}%
\pgfpathlineto{\pgfqpoint{3.401955in}{2.269529in}}%
\pgfpathlineto{\pgfqpoint{3.420119in}{2.269529in}}%
\pgfpathlineto{\pgfqpoint{3.420119in}{2.257732in}}%
\pgfpathmoveto{\pgfqpoint{3.420119in}{2.257732in}}%
\pgfpathlineto{\pgfqpoint{3.420119in}{2.257732in}}%
\pgfpathlineto{\pgfqpoint{3.420119in}{2.269529in}}%
\pgfpathlineto{\pgfqpoint{3.438283in}{2.269529in}}%
\pgfpathlineto{\pgfqpoint{3.438283in}{2.257732in}}%
\pgfpathmoveto{\pgfqpoint{3.438283in}{2.269529in}}%
\pgfpathlineto{\pgfqpoint{3.438283in}{2.269529in}}%
\pgfpathlineto{\pgfqpoint{3.438283in}{2.281326in}}%
\pgfpathlineto{\pgfqpoint{3.456447in}{2.281326in}}%
\pgfpathlineto{\pgfqpoint{3.456447in}{2.269529in}}%
\pgfpathmoveto{\pgfqpoint{3.438283in}{2.281326in}}%
\pgfpathlineto{\pgfqpoint{3.438283in}{2.281326in}}%
\pgfpathlineto{\pgfqpoint{3.438283in}{2.293122in}}%
\pgfpathlineto{\pgfqpoint{3.456447in}{2.293122in}}%
\pgfpathlineto{\pgfqpoint{3.456447in}{2.281326in}}%
\pgfpathmoveto{\pgfqpoint{3.456447in}{2.281326in}}%
\pgfpathlineto{\pgfqpoint{3.456447in}{2.281326in}}%
\pgfpathlineto{\pgfqpoint{3.456447in}{2.293122in}}%
\pgfpathlineto{\pgfqpoint{3.474611in}{2.293122in}}%
\pgfpathlineto{\pgfqpoint{3.474611in}{2.281326in}}%
\pgfpathmoveto{\pgfqpoint{3.474611in}{2.293122in}}%
\pgfpathlineto{\pgfqpoint{3.474611in}{2.293122in}}%
\pgfpathlineto{\pgfqpoint{3.474611in}{2.304919in}}%
\pgfpathlineto{\pgfqpoint{3.492775in}{2.304919in}}%
\pgfpathlineto{\pgfqpoint{3.492775in}{2.293122in}}%
\pgfpathmoveto{\pgfqpoint{3.474611in}{2.304919in}}%
\pgfpathlineto{\pgfqpoint{3.474611in}{2.304919in}}%
\pgfpathlineto{\pgfqpoint{3.474611in}{2.316716in}}%
\pgfpathlineto{\pgfqpoint{3.492775in}{2.316716in}}%
\pgfpathlineto{\pgfqpoint{3.492775in}{2.304919in}}%
\pgfpathmoveto{\pgfqpoint{3.492775in}{2.304919in}}%
\pgfpathlineto{\pgfqpoint{3.492775in}{2.304919in}}%
\pgfpathlineto{\pgfqpoint{3.492775in}{2.316716in}}%
\pgfpathlineto{\pgfqpoint{3.510939in}{2.316716in}}%
\pgfpathlineto{\pgfqpoint{3.510939in}{2.304919in}}%
\pgfpathmoveto{\pgfqpoint{3.474611in}{2.647031in}}%
\pgfpathlineto{\pgfqpoint{3.474611in}{2.647031in}}%
\pgfpathlineto{\pgfqpoint{3.474611in}{2.658827in}}%
\pgfpathlineto{\pgfqpoint{3.492775in}{2.658827in}}%
\pgfpathlineto{\pgfqpoint{3.492775in}{2.647031in}}%
\pgfpathmoveto{\pgfqpoint{3.474611in}{2.658827in}}%
\pgfpathlineto{\pgfqpoint{3.474611in}{2.658827in}}%
\pgfpathlineto{\pgfqpoint{3.474611in}{2.670624in}}%
\pgfpathlineto{\pgfqpoint{3.492775in}{2.670624in}}%
\pgfpathlineto{\pgfqpoint{3.492775in}{2.658827in}}%
\pgfpathmoveto{\pgfqpoint{3.401955in}{2.741405in}}%
\pgfpathlineto{\pgfqpoint{3.401955in}{2.741405in}}%
\pgfpathlineto{\pgfqpoint{3.401955in}{2.753202in}}%
\pgfpathlineto{\pgfqpoint{3.420119in}{2.753202in}}%
\pgfpathlineto{\pgfqpoint{3.420119in}{2.741405in}}%
\pgfpathmoveto{\pgfqpoint{3.401955in}{2.753202in}}%
\pgfpathlineto{\pgfqpoint{3.401955in}{2.753202in}}%
\pgfpathlineto{\pgfqpoint{3.401955in}{2.764999in}}%
\pgfpathlineto{\pgfqpoint{3.420119in}{2.764999in}}%
\pgfpathlineto{\pgfqpoint{3.420119in}{2.753202in}}%
\pgfpathmoveto{\pgfqpoint{3.438283in}{2.694218in}}%
\pgfpathlineto{\pgfqpoint{3.438283in}{2.694218in}}%
\pgfpathlineto{\pgfqpoint{3.438283in}{2.706015in}}%
\pgfpathlineto{\pgfqpoint{3.456447in}{2.706015in}}%
\pgfpathlineto{\pgfqpoint{3.456447in}{2.694218in}}%
\pgfpathmoveto{\pgfqpoint{3.438283in}{2.706015in}}%
\pgfpathlineto{\pgfqpoint{3.438283in}{2.706015in}}%
\pgfpathlineto{\pgfqpoint{3.438283in}{2.717812in}}%
\pgfpathlineto{\pgfqpoint{3.456447in}{2.717812in}}%
\pgfpathlineto{\pgfqpoint{3.456447in}{2.706015in}}%
\pgfpathmoveto{\pgfqpoint{3.365627in}{2.788593in}}%
\pgfpathlineto{\pgfqpoint{3.365627in}{2.788593in}}%
\pgfpathlineto{\pgfqpoint{3.365627in}{2.800390in}}%
\pgfpathlineto{\pgfqpoint{3.383791in}{2.800390in}}%
\pgfpathlineto{\pgfqpoint{3.383791in}{2.788593in}}%
\pgfpathmoveto{\pgfqpoint{3.365627in}{2.800390in}}%
\pgfpathlineto{\pgfqpoint{3.365627in}{2.800390in}}%
\pgfpathlineto{\pgfqpoint{3.365627in}{2.812187in}}%
\pgfpathlineto{\pgfqpoint{3.383791in}{2.812187in}}%
\pgfpathlineto{\pgfqpoint{3.383791in}{2.800390in}}%
\pgfpathmoveto{\pgfqpoint{3.510939in}{2.316716in}}%
\pgfpathlineto{\pgfqpoint{3.510939in}{2.316716in}}%
\pgfpathlineto{\pgfqpoint{3.510939in}{2.328514in}}%
\pgfpathlineto{\pgfqpoint{3.529102in}{2.328514in}}%
\pgfpathlineto{\pgfqpoint{3.529102in}{2.316716in}}%
\pgfpathmoveto{\pgfqpoint{3.510939in}{2.328514in}}%
\pgfpathlineto{\pgfqpoint{3.510939in}{2.328514in}}%
\pgfpathlineto{\pgfqpoint{3.510939in}{2.340311in}}%
\pgfpathlineto{\pgfqpoint{3.529102in}{2.340311in}}%
\pgfpathlineto{\pgfqpoint{3.529102in}{2.328514in}}%
\pgfpathmoveto{\pgfqpoint{3.529102in}{2.328514in}}%
\pgfpathlineto{\pgfqpoint{3.529102in}{2.328514in}}%
\pgfpathlineto{\pgfqpoint{3.529102in}{2.340311in}}%
\pgfpathlineto{\pgfqpoint{3.547265in}{2.340311in}}%
\pgfpathlineto{\pgfqpoint{3.547265in}{2.328514in}}%
\pgfpathmoveto{\pgfqpoint{3.547265in}{2.340311in}}%
\pgfpathlineto{\pgfqpoint{3.547265in}{2.340311in}}%
\pgfpathlineto{\pgfqpoint{3.547265in}{2.352108in}}%
\pgfpathlineto{\pgfqpoint{3.565429in}{2.352108in}}%
\pgfpathlineto{\pgfqpoint{3.565429in}{2.340311in}}%
\pgfpathmoveto{\pgfqpoint{3.547265in}{2.352108in}}%
\pgfpathlineto{\pgfqpoint{3.547265in}{2.352108in}}%
\pgfpathlineto{\pgfqpoint{3.547265in}{2.363905in}}%
\pgfpathlineto{\pgfqpoint{3.565429in}{2.363905in}}%
\pgfpathlineto{\pgfqpoint{3.565429in}{2.352108in}}%
\pgfpathmoveto{\pgfqpoint{3.565429in}{2.352108in}}%
\pgfpathlineto{\pgfqpoint{3.565429in}{2.352108in}}%
\pgfpathlineto{\pgfqpoint{3.565429in}{2.363905in}}%
\pgfpathlineto{\pgfqpoint{3.583592in}{2.363905in}}%
\pgfpathlineto{\pgfqpoint{3.583592in}{2.352108in}}%
\pgfpathmoveto{\pgfqpoint{3.619919in}{2.375702in}}%
\pgfpathlineto{\pgfqpoint{3.619919in}{2.375702in}}%
\pgfpathlineto{\pgfqpoint{3.619919in}{2.387499in}}%
\pgfpathlineto{\pgfqpoint{3.638082in}{2.387499in}}%
\pgfpathlineto{\pgfqpoint{3.638082in}{2.375702in}}%
\pgfpathmoveto{\pgfqpoint{3.619919in}{2.458283in}}%
\pgfpathlineto{\pgfqpoint{3.619919in}{2.458283in}}%
\pgfpathlineto{\pgfqpoint{3.619919in}{2.470080in}}%
\pgfpathlineto{\pgfqpoint{3.638082in}{2.470080in}}%
\pgfpathlineto{\pgfqpoint{3.638082in}{2.458283in}}%
\pgfpathmoveto{\pgfqpoint{3.619919in}{2.470080in}}%
\pgfpathlineto{\pgfqpoint{3.619919in}{2.470080in}}%
\pgfpathlineto{\pgfqpoint{3.619919in}{2.481878in}}%
\pgfpathlineto{\pgfqpoint{3.638082in}{2.481878in}}%
\pgfpathlineto{\pgfqpoint{3.638082in}{2.470080in}}%
\pgfpathmoveto{\pgfqpoint{3.547265in}{2.552658in}}%
\pgfpathlineto{\pgfqpoint{3.547265in}{2.552658in}}%
\pgfpathlineto{\pgfqpoint{3.547265in}{2.564454in}}%
\pgfpathlineto{\pgfqpoint{3.565429in}{2.564454in}}%
\pgfpathlineto{\pgfqpoint{3.565429in}{2.552658in}}%
\pgfpathmoveto{\pgfqpoint{3.547265in}{2.564454in}}%
\pgfpathlineto{\pgfqpoint{3.547265in}{2.564454in}}%
\pgfpathlineto{\pgfqpoint{3.547265in}{2.576251in}}%
\pgfpathlineto{\pgfqpoint{3.565429in}{2.576251in}}%
\pgfpathlineto{\pgfqpoint{3.565429in}{2.564454in}}%
\pgfpathmoveto{\pgfqpoint{3.583592in}{2.505471in}}%
\pgfpathlineto{\pgfqpoint{3.583592in}{2.505471in}}%
\pgfpathlineto{\pgfqpoint{3.583592in}{2.517268in}}%
\pgfpathlineto{\pgfqpoint{3.601755in}{2.517268in}}%
\pgfpathlineto{\pgfqpoint{3.601755in}{2.505471in}}%
\pgfpathmoveto{\pgfqpoint{3.583592in}{2.517268in}}%
\pgfpathlineto{\pgfqpoint{3.583592in}{2.517268in}}%
\pgfpathlineto{\pgfqpoint{3.583592in}{2.529064in}}%
\pgfpathlineto{\pgfqpoint{3.601755in}{2.529064in}}%
\pgfpathlineto{\pgfqpoint{3.601755in}{2.517268in}}%
\pgfpathmoveto{\pgfqpoint{3.510939in}{2.599844in}}%
\pgfpathlineto{\pgfqpoint{3.510939in}{2.599844in}}%
\pgfpathlineto{\pgfqpoint{3.510939in}{2.611641in}}%
\pgfpathlineto{\pgfqpoint{3.529102in}{2.611641in}}%
\pgfpathlineto{\pgfqpoint{3.529102in}{2.599844in}}%
\pgfpathmoveto{\pgfqpoint{3.510939in}{2.611641in}}%
\pgfpathlineto{\pgfqpoint{3.510939in}{2.611641in}}%
\pgfpathlineto{\pgfqpoint{3.510939in}{2.623437in}}%
\pgfpathlineto{\pgfqpoint{3.529102in}{2.623437in}}%
\pgfpathlineto{\pgfqpoint{3.529102in}{2.611641in}}%
\pgfpathmoveto{\pgfqpoint{3.656246in}{2.399296in}}%
\pgfpathlineto{\pgfqpoint{3.656246in}{2.399296in}}%
\pgfpathlineto{\pgfqpoint{3.656246in}{2.411093in}}%
\pgfpathlineto{\pgfqpoint{3.674410in}{2.411093in}}%
\pgfpathlineto{\pgfqpoint{3.674410in}{2.399296in}}%
\pgfpathmoveto{\pgfqpoint{3.656246in}{2.411093in}}%
\pgfpathlineto{\pgfqpoint{3.656246in}{2.411093in}}%
\pgfpathlineto{\pgfqpoint{3.656246in}{2.422891in}}%
\pgfpathlineto{\pgfqpoint{3.674410in}{2.422891in}}%
\pgfpathlineto{\pgfqpoint{3.674410in}{2.411093in}}%
\pgfpathmoveto{\pgfqpoint{3.656246in}{2.422891in}}%
\pgfpathlineto{\pgfqpoint{3.656246in}{2.422891in}}%
\pgfpathlineto{\pgfqpoint{3.656246in}{2.434688in}}%
\pgfpathlineto{\pgfqpoint{3.674410in}{2.434688in}}%
\pgfpathlineto{\pgfqpoint{3.674410in}{2.422891in}}%
\pgfpathmoveto{\pgfqpoint{0.749998in}{0.511795in}}%
\pgfpathlineto{\pgfqpoint{0.749998in}{0.511795in}}%
\pgfpathlineto{\pgfqpoint{0.749998in}{0.517694in}}%
\pgfpathlineto{\pgfqpoint{0.759080in}{0.517694in}}%
\pgfpathlineto{\pgfqpoint{0.759080in}{0.511795in}}%
\pgfpathmoveto{\pgfqpoint{0.749998in}{0.517694in}}%
\pgfpathlineto{\pgfqpoint{0.749998in}{0.517694in}}%
\pgfpathlineto{\pgfqpoint{0.749998in}{0.523592in}}%
\pgfpathlineto{\pgfqpoint{0.759080in}{0.523592in}}%
\pgfpathlineto{\pgfqpoint{0.759080in}{0.517694in}}%
\pgfpathmoveto{\pgfqpoint{0.759080in}{0.517694in}}%
\pgfpathlineto{\pgfqpoint{0.759080in}{0.517694in}}%
\pgfpathlineto{\pgfqpoint{0.759080in}{0.523592in}}%
\pgfpathlineto{\pgfqpoint{0.768162in}{0.523592in}}%
\pgfpathlineto{\pgfqpoint{0.768162in}{0.517694in}}%
\pgfpathmoveto{\pgfqpoint{0.768162in}{0.523592in}}%
\pgfpathlineto{\pgfqpoint{0.768162in}{0.523592in}}%
\pgfpathlineto{\pgfqpoint{0.768162in}{0.529491in}}%
\pgfpathlineto{\pgfqpoint{0.777244in}{0.529491in}}%
\pgfpathlineto{\pgfqpoint{0.777244in}{0.523592in}}%
\pgfpathmoveto{\pgfqpoint{0.768162in}{0.529491in}}%
\pgfpathlineto{\pgfqpoint{0.768162in}{0.529491in}}%
\pgfpathlineto{\pgfqpoint{0.768162in}{0.535389in}}%
\pgfpathlineto{\pgfqpoint{0.777244in}{0.535389in}}%
\pgfpathlineto{\pgfqpoint{0.777244in}{0.529491in}}%
\pgfpathmoveto{\pgfqpoint{0.777244in}{0.529491in}}%
\pgfpathlineto{\pgfqpoint{0.777244in}{0.529491in}}%
\pgfpathlineto{\pgfqpoint{0.777244in}{0.535389in}}%
\pgfpathlineto{\pgfqpoint{0.786327in}{0.535389in}}%
\pgfpathlineto{\pgfqpoint{0.786327in}{0.529491in}}%
\pgfpathmoveto{\pgfqpoint{0.786327in}{0.535389in}}%
\pgfpathlineto{\pgfqpoint{0.786327in}{0.535389in}}%
\pgfpathlineto{\pgfqpoint{0.786327in}{0.541288in}}%
\pgfpathlineto{\pgfqpoint{0.795409in}{0.541288in}}%
\pgfpathlineto{\pgfqpoint{0.795409in}{0.535389in}}%
\pgfpathmoveto{\pgfqpoint{0.786327in}{0.541288in}}%
\pgfpathlineto{\pgfqpoint{0.786327in}{0.541288in}}%
\pgfpathlineto{\pgfqpoint{0.786327in}{0.547186in}}%
\pgfpathlineto{\pgfqpoint{0.795409in}{0.547186in}}%
\pgfpathlineto{\pgfqpoint{0.795409in}{0.541288in}}%
\pgfpathmoveto{\pgfqpoint{0.795409in}{0.541288in}}%
\pgfpathlineto{\pgfqpoint{0.795409in}{0.541288in}}%
\pgfpathlineto{\pgfqpoint{0.795409in}{0.547186in}}%
\pgfpathlineto{\pgfqpoint{0.804491in}{0.547186in}}%
\pgfpathlineto{\pgfqpoint{0.804491in}{0.541288in}}%
\pgfpathmoveto{\pgfqpoint{0.804491in}{0.547186in}}%
\pgfpathlineto{\pgfqpoint{0.804491in}{0.547186in}}%
\pgfpathlineto{\pgfqpoint{0.804491in}{0.553085in}}%
\pgfpathlineto{\pgfqpoint{0.813573in}{0.553085in}}%
\pgfpathlineto{\pgfqpoint{0.813573in}{0.547186in}}%
\pgfpathmoveto{\pgfqpoint{0.804491in}{0.553085in}}%
\pgfpathlineto{\pgfqpoint{0.804491in}{0.553085in}}%
\pgfpathlineto{\pgfqpoint{0.804491in}{0.558983in}}%
\pgfpathlineto{\pgfqpoint{0.813573in}{0.558983in}}%
\pgfpathlineto{\pgfqpoint{0.813573in}{0.553085in}}%
\pgfpathmoveto{\pgfqpoint{0.813573in}{0.553085in}}%
\pgfpathlineto{\pgfqpoint{0.813573in}{0.553085in}}%
\pgfpathlineto{\pgfqpoint{0.813573in}{0.558983in}}%
\pgfpathlineto{\pgfqpoint{0.822655in}{0.558983in}}%
\pgfpathlineto{\pgfqpoint{0.822655in}{0.553085in}}%
\pgfpathmoveto{\pgfqpoint{0.822655in}{0.558983in}}%
\pgfpathlineto{\pgfqpoint{0.822655in}{0.558983in}}%
\pgfpathlineto{\pgfqpoint{0.822655in}{0.564882in}}%
\pgfpathlineto{\pgfqpoint{0.831737in}{0.564882in}}%
\pgfpathlineto{\pgfqpoint{0.831737in}{0.558983in}}%
\pgfpathmoveto{\pgfqpoint{0.822655in}{0.564882in}}%
\pgfpathlineto{\pgfqpoint{0.822655in}{0.564882in}}%
\pgfpathlineto{\pgfqpoint{0.822655in}{0.570780in}}%
\pgfpathlineto{\pgfqpoint{0.831737in}{0.570780in}}%
\pgfpathlineto{\pgfqpoint{0.831737in}{0.564882in}}%
\pgfpathmoveto{\pgfqpoint{0.831737in}{0.564882in}}%
\pgfpathlineto{\pgfqpoint{0.831737in}{0.564882in}}%
\pgfpathlineto{\pgfqpoint{0.831737in}{0.570780in}}%
\pgfpathlineto{\pgfqpoint{0.840819in}{0.570780in}}%
\pgfpathlineto{\pgfqpoint{0.840819in}{0.564882in}}%
\pgfpathmoveto{\pgfqpoint{0.840819in}{0.570780in}}%
\pgfpathlineto{\pgfqpoint{0.840819in}{0.570780in}}%
\pgfpathlineto{\pgfqpoint{0.840819in}{0.576679in}}%
\pgfpathlineto{\pgfqpoint{0.849901in}{0.576679in}}%
\pgfpathlineto{\pgfqpoint{0.849901in}{0.570780in}}%
\pgfpathmoveto{\pgfqpoint{0.840819in}{0.576679in}}%
\pgfpathlineto{\pgfqpoint{0.840819in}{0.576679in}}%
\pgfpathlineto{\pgfqpoint{0.840819in}{0.582577in}}%
\pgfpathlineto{\pgfqpoint{0.849901in}{0.582577in}}%
\pgfpathlineto{\pgfqpoint{0.849901in}{0.576679in}}%
\pgfpathmoveto{\pgfqpoint{0.849901in}{0.576679in}}%
\pgfpathlineto{\pgfqpoint{0.849901in}{0.576679in}}%
\pgfpathlineto{\pgfqpoint{0.849901in}{0.582577in}}%
\pgfpathlineto{\pgfqpoint{0.858983in}{0.582577in}}%
\pgfpathlineto{\pgfqpoint{0.858983in}{0.576679in}}%
\pgfpathmoveto{\pgfqpoint{0.858983in}{0.582577in}}%
\pgfpathlineto{\pgfqpoint{0.858983in}{0.582577in}}%
\pgfpathlineto{\pgfqpoint{0.858983in}{0.588476in}}%
\pgfpathlineto{\pgfqpoint{0.868065in}{0.588476in}}%
\pgfpathlineto{\pgfqpoint{0.868065in}{0.582577in}}%
\pgfpathmoveto{\pgfqpoint{0.858983in}{0.588476in}}%
\pgfpathlineto{\pgfqpoint{0.858983in}{0.588476in}}%
\pgfpathlineto{\pgfqpoint{0.858983in}{0.594374in}}%
\pgfpathlineto{\pgfqpoint{0.868065in}{0.594374in}}%
\pgfpathlineto{\pgfqpoint{0.868065in}{0.588476in}}%
\pgfpathmoveto{\pgfqpoint{0.868065in}{0.588476in}}%
\pgfpathlineto{\pgfqpoint{0.868065in}{0.588476in}}%
\pgfpathlineto{\pgfqpoint{0.868065in}{0.594374in}}%
\pgfpathlineto{\pgfqpoint{0.877147in}{0.594374in}}%
\pgfpathlineto{\pgfqpoint{0.877147in}{0.588476in}}%
\pgfpathmoveto{\pgfqpoint{0.877147in}{0.594374in}}%
\pgfpathlineto{\pgfqpoint{0.877147in}{0.594374in}}%
\pgfpathlineto{\pgfqpoint{0.877147in}{0.600273in}}%
\pgfpathlineto{\pgfqpoint{0.886230in}{0.600273in}}%
\pgfpathlineto{\pgfqpoint{0.886230in}{0.594374in}}%
\pgfpathmoveto{\pgfqpoint{0.877147in}{0.600273in}}%
\pgfpathlineto{\pgfqpoint{0.877147in}{0.600273in}}%
\pgfpathlineto{\pgfqpoint{0.877147in}{0.606171in}}%
\pgfpathlineto{\pgfqpoint{0.886230in}{0.606171in}}%
\pgfpathlineto{\pgfqpoint{0.886230in}{0.600273in}}%
\pgfpathmoveto{\pgfqpoint{0.886230in}{0.600273in}}%
\pgfpathlineto{\pgfqpoint{0.886230in}{0.600273in}}%
\pgfpathlineto{\pgfqpoint{0.886230in}{0.606171in}}%
\pgfpathlineto{\pgfqpoint{0.895312in}{0.606171in}}%
\pgfpathlineto{\pgfqpoint{0.895312in}{0.600273in}}%
\pgfpathmoveto{\pgfqpoint{0.895312in}{0.606171in}}%
\pgfpathlineto{\pgfqpoint{0.895312in}{0.606171in}}%
\pgfpathlineto{\pgfqpoint{0.895312in}{0.612070in}}%
\pgfpathlineto{\pgfqpoint{0.904394in}{0.612070in}}%
\pgfpathlineto{\pgfqpoint{0.904394in}{0.606171in}}%
\pgfpathmoveto{\pgfqpoint{0.895312in}{0.612070in}}%
\pgfpathlineto{\pgfqpoint{0.895312in}{0.612070in}}%
\pgfpathlineto{\pgfqpoint{0.895312in}{0.617968in}}%
\pgfpathlineto{\pgfqpoint{0.904394in}{0.617968in}}%
\pgfpathlineto{\pgfqpoint{0.904394in}{0.612070in}}%
\pgfpathmoveto{\pgfqpoint{0.904394in}{0.612070in}}%
\pgfpathlineto{\pgfqpoint{0.904394in}{0.612070in}}%
\pgfpathlineto{\pgfqpoint{0.904394in}{0.617968in}}%
\pgfpathlineto{\pgfqpoint{0.913475in}{0.617968in}}%
\pgfpathlineto{\pgfqpoint{0.913475in}{0.612070in}}%
\pgfpathmoveto{\pgfqpoint{0.913475in}{0.612070in}}%
\pgfpathlineto{\pgfqpoint{0.913475in}{0.612070in}}%
\pgfpathlineto{\pgfqpoint{0.913475in}{0.617968in}}%
\pgfpathlineto{\pgfqpoint{0.922557in}{0.617968in}}%
\pgfpathlineto{\pgfqpoint{0.922557in}{0.612070in}}%
\pgfpathmoveto{\pgfqpoint{0.931639in}{0.623867in}}%
\pgfpathlineto{\pgfqpoint{0.931639in}{0.623867in}}%
\pgfpathlineto{\pgfqpoint{0.931639in}{0.629765in}}%
\pgfpathlineto{\pgfqpoint{0.940721in}{0.629765in}}%
\pgfpathlineto{\pgfqpoint{0.940721in}{0.623867in}}%
\pgfpathmoveto{\pgfqpoint{0.949803in}{0.635664in}}%
\pgfpathlineto{\pgfqpoint{0.949803in}{0.635664in}}%
\pgfpathlineto{\pgfqpoint{0.949803in}{0.641562in}}%
\pgfpathlineto{\pgfqpoint{0.958885in}{0.641562in}}%
\pgfpathlineto{\pgfqpoint{0.958885in}{0.635664in}}%
\pgfpathmoveto{\pgfqpoint{0.967967in}{0.647461in}}%
\pgfpathlineto{\pgfqpoint{0.967967in}{0.647461in}}%
\pgfpathlineto{\pgfqpoint{0.967967in}{0.653359in}}%
\pgfpathlineto{\pgfqpoint{0.977049in}{0.653359in}}%
\pgfpathlineto{\pgfqpoint{0.977049in}{0.647461in}}%
\pgfpathmoveto{\pgfqpoint{0.986131in}{0.659258in}}%
\pgfpathlineto{\pgfqpoint{0.986131in}{0.659258in}}%
\pgfpathlineto{\pgfqpoint{0.986131in}{0.665156in}}%
\pgfpathlineto{\pgfqpoint{0.995213in}{0.665156in}}%
\pgfpathlineto{\pgfqpoint{0.995213in}{0.659258in}}%
\pgfpathmoveto{\pgfqpoint{1.004294in}{0.671054in}}%
\pgfpathlineto{\pgfqpoint{1.004294in}{0.671054in}}%
\pgfpathlineto{\pgfqpoint{1.004294in}{0.676953in}}%
\pgfpathlineto{\pgfqpoint{1.013376in}{0.676953in}}%
\pgfpathlineto{\pgfqpoint{1.013376in}{0.671054in}}%
\pgfpathmoveto{\pgfqpoint{1.022458in}{0.682851in}}%
\pgfpathlineto{\pgfqpoint{1.022458in}{0.682851in}}%
\pgfpathlineto{\pgfqpoint{1.022458in}{0.688750in}}%
\pgfpathlineto{\pgfqpoint{1.031540in}{0.688750in}}%
\pgfpathlineto{\pgfqpoint{1.031540in}{0.682851in}}%
\pgfpathmoveto{\pgfqpoint{1.040622in}{0.694648in}}%
\pgfpathlineto{\pgfqpoint{1.040622in}{0.694648in}}%
\pgfpathlineto{\pgfqpoint{1.040622in}{0.700547in}}%
\pgfpathlineto{\pgfqpoint{1.049704in}{0.700547in}}%
\pgfpathlineto{\pgfqpoint{1.049704in}{0.694648in}}%
\pgfpathmoveto{\pgfqpoint{1.058787in}{0.706445in}}%
\pgfpathlineto{\pgfqpoint{1.058787in}{0.706445in}}%
\pgfpathlineto{\pgfqpoint{1.058787in}{0.712343in}}%
\pgfpathlineto{\pgfqpoint{1.067869in}{0.712343in}}%
\pgfpathlineto{\pgfqpoint{1.067869in}{0.706445in}}%
\pgfpathmoveto{\pgfqpoint{1.076952in}{0.718242in}}%
\pgfpathlineto{\pgfqpoint{1.076952in}{0.718242in}}%
\pgfpathlineto{\pgfqpoint{1.076952in}{0.724140in}}%
\pgfpathlineto{\pgfqpoint{1.086034in}{0.724140in}}%
\pgfpathlineto{\pgfqpoint{1.086034in}{0.718242in}}%
\pgfpathmoveto{\pgfqpoint{1.076952in}{0.724140in}}%
\pgfpathlineto{\pgfqpoint{1.076952in}{0.724140in}}%
\pgfpathlineto{\pgfqpoint{1.076952in}{0.730039in}}%
\pgfpathlineto{\pgfqpoint{1.086034in}{0.730039in}}%
\pgfpathlineto{\pgfqpoint{1.086034in}{0.724140in}}%
\pgfpathmoveto{\pgfqpoint{1.076952in}{0.730039in}}%
\pgfpathlineto{\pgfqpoint{1.076952in}{0.730039in}}%
\pgfpathlineto{\pgfqpoint{1.076952in}{0.735937in}}%
\pgfpathlineto{\pgfqpoint{1.086034in}{0.735937in}}%
\pgfpathlineto{\pgfqpoint{1.086034in}{0.730039in}}%
\pgfpathmoveto{\pgfqpoint{1.086034in}{0.730039in}}%
\pgfpathlineto{\pgfqpoint{1.086034in}{0.730039in}}%
\pgfpathlineto{\pgfqpoint{1.086034in}{0.735937in}}%
\pgfpathlineto{\pgfqpoint{1.095117in}{0.735937in}}%
\pgfpathlineto{\pgfqpoint{1.095117in}{0.730039in}}%
\pgfpathmoveto{\pgfqpoint{1.095117in}{0.735937in}}%
\pgfpathlineto{\pgfqpoint{1.095117in}{0.735937in}}%
\pgfpathlineto{\pgfqpoint{1.095117in}{0.741835in}}%
\pgfpathlineto{\pgfqpoint{1.104199in}{0.741835in}}%
\pgfpathlineto{\pgfqpoint{1.104199in}{0.735937in}}%
\pgfpathmoveto{\pgfqpoint{1.095117in}{0.741835in}}%
\pgfpathlineto{\pgfqpoint{1.095117in}{0.741835in}}%
\pgfpathlineto{\pgfqpoint{1.095117in}{0.747734in}}%
\pgfpathlineto{\pgfqpoint{1.104199in}{0.747734in}}%
\pgfpathlineto{\pgfqpoint{1.104199in}{0.741835in}}%
\pgfpathmoveto{\pgfqpoint{1.104199in}{0.741835in}}%
\pgfpathlineto{\pgfqpoint{1.104199in}{0.741835in}}%
\pgfpathlineto{\pgfqpoint{1.104199in}{0.747734in}}%
\pgfpathlineto{\pgfqpoint{1.113282in}{0.747734in}}%
\pgfpathlineto{\pgfqpoint{1.113282in}{0.741835in}}%
\pgfpathmoveto{\pgfqpoint{1.113282in}{0.747734in}}%
\pgfpathlineto{\pgfqpoint{1.113282in}{0.747734in}}%
\pgfpathlineto{\pgfqpoint{1.113282in}{0.753632in}}%
\pgfpathlineto{\pgfqpoint{1.122364in}{0.753632in}}%
\pgfpathlineto{\pgfqpoint{1.122364in}{0.747734in}}%
\pgfpathmoveto{\pgfqpoint{1.113282in}{0.753632in}}%
\pgfpathlineto{\pgfqpoint{1.113282in}{0.753632in}}%
\pgfpathlineto{\pgfqpoint{1.113282in}{0.759530in}}%
\pgfpathlineto{\pgfqpoint{1.122364in}{0.759530in}}%
\pgfpathlineto{\pgfqpoint{1.122364in}{0.753632in}}%
\pgfpathmoveto{\pgfqpoint{1.122364in}{0.753632in}}%
\pgfpathlineto{\pgfqpoint{1.122364in}{0.753632in}}%
\pgfpathlineto{\pgfqpoint{1.122364in}{0.759530in}}%
\pgfpathlineto{\pgfqpoint{1.131447in}{0.759530in}}%
\pgfpathlineto{\pgfqpoint{1.131447in}{0.753632in}}%
\pgfpathmoveto{\pgfqpoint{1.131447in}{0.759530in}}%
\pgfpathlineto{\pgfqpoint{1.131447in}{0.759530in}}%
\pgfpathlineto{\pgfqpoint{1.131447in}{0.765429in}}%
\pgfpathlineto{\pgfqpoint{1.140529in}{0.765429in}}%
\pgfpathlineto{\pgfqpoint{1.140529in}{0.759530in}}%
\pgfpathmoveto{\pgfqpoint{1.131447in}{0.765429in}}%
\pgfpathlineto{\pgfqpoint{1.131447in}{0.765429in}}%
\pgfpathlineto{\pgfqpoint{1.131447in}{0.771327in}}%
\pgfpathlineto{\pgfqpoint{1.140529in}{0.771327in}}%
\pgfpathlineto{\pgfqpoint{1.140529in}{0.765429in}}%
\pgfpathmoveto{\pgfqpoint{1.140529in}{0.765429in}}%
\pgfpathlineto{\pgfqpoint{1.140529in}{0.765429in}}%
\pgfpathlineto{\pgfqpoint{1.140529in}{0.771327in}}%
\pgfpathlineto{\pgfqpoint{1.149611in}{0.771327in}}%
\pgfpathlineto{\pgfqpoint{1.149611in}{0.765429in}}%
\pgfpathmoveto{\pgfqpoint{1.149611in}{0.771327in}}%
\pgfpathlineto{\pgfqpoint{1.149611in}{0.771327in}}%
\pgfpathlineto{\pgfqpoint{1.149611in}{0.777225in}}%
\pgfpathlineto{\pgfqpoint{1.158694in}{0.777225in}}%
\pgfpathlineto{\pgfqpoint{1.158694in}{0.771327in}}%
\pgfpathmoveto{\pgfqpoint{1.149611in}{0.777225in}}%
\pgfpathlineto{\pgfqpoint{1.149611in}{0.777225in}}%
\pgfpathlineto{\pgfqpoint{1.149611in}{0.783124in}}%
\pgfpathlineto{\pgfqpoint{1.158694in}{0.783124in}}%
\pgfpathlineto{\pgfqpoint{1.158694in}{0.777225in}}%
\pgfpathmoveto{\pgfqpoint{1.158694in}{0.777225in}}%
\pgfpathlineto{\pgfqpoint{1.158694in}{0.777225in}}%
\pgfpathlineto{\pgfqpoint{1.158694in}{0.783124in}}%
\pgfpathlineto{\pgfqpoint{1.167776in}{0.783124in}}%
\pgfpathlineto{\pgfqpoint{1.167776in}{0.777225in}}%
\pgfpathmoveto{\pgfqpoint{1.167776in}{0.783124in}}%
\pgfpathlineto{\pgfqpoint{1.167776in}{0.783124in}}%
\pgfpathlineto{\pgfqpoint{1.167776in}{0.789022in}}%
\pgfpathlineto{\pgfqpoint{1.176859in}{0.789022in}}%
\pgfpathlineto{\pgfqpoint{1.176859in}{0.783124in}}%
\pgfpathmoveto{\pgfqpoint{1.167776in}{0.789022in}}%
\pgfpathlineto{\pgfqpoint{1.167776in}{0.789022in}}%
\pgfpathlineto{\pgfqpoint{1.167776in}{0.794921in}}%
\pgfpathlineto{\pgfqpoint{1.176859in}{0.794921in}}%
\pgfpathlineto{\pgfqpoint{1.176859in}{0.789022in}}%
\pgfpathmoveto{\pgfqpoint{1.176859in}{0.789022in}}%
\pgfpathlineto{\pgfqpoint{1.176859in}{0.789022in}}%
\pgfpathlineto{\pgfqpoint{1.176859in}{0.794921in}}%
\pgfpathlineto{\pgfqpoint{1.185941in}{0.794921in}}%
\pgfpathlineto{\pgfqpoint{1.185941in}{0.789022in}}%
\pgfpathmoveto{\pgfqpoint{1.185941in}{0.794921in}}%
\pgfpathlineto{\pgfqpoint{1.185941in}{0.794921in}}%
\pgfpathlineto{\pgfqpoint{1.185941in}{0.800820in}}%
\pgfpathlineto{\pgfqpoint{1.195023in}{0.800820in}}%
\pgfpathlineto{\pgfqpoint{1.195023in}{0.794921in}}%
\pgfpathmoveto{\pgfqpoint{1.185941in}{0.800820in}}%
\pgfpathlineto{\pgfqpoint{1.185941in}{0.800820in}}%
\pgfpathlineto{\pgfqpoint{1.185941in}{0.806718in}}%
\pgfpathlineto{\pgfqpoint{1.195023in}{0.806718in}}%
\pgfpathlineto{\pgfqpoint{1.195023in}{0.800820in}}%
\pgfpathmoveto{\pgfqpoint{1.195023in}{0.800820in}}%
\pgfpathlineto{\pgfqpoint{1.195023in}{0.800820in}}%
\pgfpathlineto{\pgfqpoint{1.195023in}{0.806718in}}%
\pgfpathlineto{\pgfqpoint{1.204105in}{0.806718in}}%
\pgfpathlineto{\pgfqpoint{1.204105in}{0.800820in}}%
\pgfpathmoveto{\pgfqpoint{1.204105in}{0.806718in}}%
\pgfpathlineto{\pgfqpoint{1.204105in}{0.806718in}}%
\pgfpathlineto{\pgfqpoint{1.204105in}{0.812617in}}%
\pgfpathlineto{\pgfqpoint{1.213187in}{0.812617in}}%
\pgfpathlineto{\pgfqpoint{1.213187in}{0.806718in}}%
\pgfpathmoveto{\pgfqpoint{1.204105in}{0.812617in}}%
\pgfpathlineto{\pgfqpoint{1.204105in}{0.812617in}}%
\pgfpathlineto{\pgfqpoint{1.204105in}{0.818515in}}%
\pgfpathlineto{\pgfqpoint{1.213187in}{0.818515in}}%
\pgfpathlineto{\pgfqpoint{1.213187in}{0.812617in}}%
\pgfpathmoveto{\pgfqpoint{1.213187in}{0.812617in}}%
\pgfpathlineto{\pgfqpoint{1.213187in}{0.812617in}}%
\pgfpathlineto{\pgfqpoint{1.213187in}{0.818515in}}%
\pgfpathlineto{\pgfqpoint{1.222269in}{0.818515in}}%
\pgfpathlineto{\pgfqpoint{1.222269in}{0.812617in}}%
\pgfpathmoveto{\pgfqpoint{1.222269in}{0.818515in}}%
\pgfpathlineto{\pgfqpoint{1.222269in}{0.818515in}}%
\pgfpathlineto{\pgfqpoint{1.222269in}{0.824414in}}%
\pgfpathlineto{\pgfqpoint{1.231351in}{0.824414in}}%
\pgfpathlineto{\pgfqpoint{1.231351in}{0.818515in}}%
\pgfpathmoveto{\pgfqpoint{1.222269in}{0.824414in}}%
\pgfpathlineto{\pgfqpoint{1.222269in}{0.824414in}}%
\pgfpathlineto{\pgfqpoint{1.222269in}{0.830313in}}%
\pgfpathlineto{\pgfqpoint{1.231351in}{0.830313in}}%
\pgfpathlineto{\pgfqpoint{1.231351in}{0.824414in}}%
\pgfpathmoveto{\pgfqpoint{1.231351in}{0.824414in}}%
\pgfpathlineto{\pgfqpoint{1.231351in}{0.824414in}}%
\pgfpathlineto{\pgfqpoint{1.231351in}{0.830313in}}%
\pgfpathlineto{\pgfqpoint{1.240433in}{0.830313in}}%
\pgfpathlineto{\pgfqpoint{1.240433in}{0.824414in}}%
\pgfpathmoveto{\pgfqpoint{1.240433in}{0.830313in}}%
\pgfpathlineto{\pgfqpoint{1.240433in}{0.830313in}}%
\pgfpathlineto{\pgfqpoint{1.240433in}{0.836211in}}%
\pgfpathlineto{\pgfqpoint{1.249516in}{0.836211in}}%
\pgfpathlineto{\pgfqpoint{1.249516in}{0.830313in}}%
\pgfpathmoveto{\pgfqpoint{1.240433in}{0.836211in}}%
\pgfpathlineto{\pgfqpoint{1.240433in}{0.836211in}}%
\pgfpathlineto{\pgfqpoint{1.240433in}{0.842110in}}%
\pgfpathlineto{\pgfqpoint{1.249516in}{0.842110in}}%
\pgfpathlineto{\pgfqpoint{1.249516in}{0.836211in}}%
\pgfpathmoveto{\pgfqpoint{1.249516in}{0.836211in}}%
\pgfpathlineto{\pgfqpoint{1.249516in}{0.836211in}}%
\pgfpathlineto{\pgfqpoint{1.249516in}{0.842110in}}%
\pgfpathlineto{\pgfqpoint{1.258598in}{0.842110in}}%
\pgfpathlineto{\pgfqpoint{1.258598in}{0.836211in}}%
\pgfpathmoveto{\pgfqpoint{1.258598in}{0.842110in}}%
\pgfpathlineto{\pgfqpoint{1.258598in}{0.842110in}}%
\pgfpathlineto{\pgfqpoint{1.258598in}{0.848008in}}%
\pgfpathlineto{\pgfqpoint{1.267680in}{0.848008in}}%
\pgfpathlineto{\pgfqpoint{1.267680in}{0.842110in}}%
\pgfpathmoveto{\pgfqpoint{1.258598in}{0.848008in}}%
\pgfpathlineto{\pgfqpoint{1.258598in}{0.848008in}}%
\pgfpathlineto{\pgfqpoint{1.258598in}{0.853907in}}%
\pgfpathlineto{\pgfqpoint{1.267680in}{0.853907in}}%
\pgfpathlineto{\pgfqpoint{1.267680in}{0.848008in}}%
\pgfpathmoveto{\pgfqpoint{1.267680in}{0.848008in}}%
\pgfpathlineto{\pgfqpoint{1.267680in}{0.848008in}}%
\pgfpathlineto{\pgfqpoint{1.267680in}{0.853907in}}%
\pgfpathlineto{\pgfqpoint{1.276762in}{0.853907in}}%
\pgfpathlineto{\pgfqpoint{1.276762in}{0.848008in}}%
\pgfpathmoveto{\pgfqpoint{1.276762in}{0.853907in}}%
\pgfpathlineto{\pgfqpoint{1.276762in}{0.853907in}}%
\pgfpathlineto{\pgfqpoint{1.276762in}{0.859806in}}%
\pgfpathlineto{\pgfqpoint{1.285844in}{0.859806in}}%
\pgfpathlineto{\pgfqpoint{1.285844in}{0.853907in}}%
\pgfpathmoveto{\pgfqpoint{1.276762in}{0.859806in}}%
\pgfpathlineto{\pgfqpoint{1.276762in}{0.859806in}}%
\pgfpathlineto{\pgfqpoint{1.276762in}{0.865704in}}%
\pgfpathlineto{\pgfqpoint{1.285844in}{0.865704in}}%
\pgfpathlineto{\pgfqpoint{1.285844in}{0.859806in}}%
\pgfpathmoveto{\pgfqpoint{1.285844in}{0.859806in}}%
\pgfpathlineto{\pgfqpoint{1.285844in}{0.859806in}}%
\pgfpathlineto{\pgfqpoint{1.285844in}{0.865704in}}%
\pgfpathlineto{\pgfqpoint{1.294926in}{0.865704in}}%
\pgfpathlineto{\pgfqpoint{1.294926in}{0.859806in}}%
\pgfpathmoveto{\pgfqpoint{1.294926in}{0.865704in}}%
\pgfpathlineto{\pgfqpoint{1.294926in}{0.865704in}}%
\pgfpathlineto{\pgfqpoint{1.294926in}{0.871603in}}%
\pgfpathlineto{\pgfqpoint{1.304008in}{0.871603in}}%
\pgfpathlineto{\pgfqpoint{1.304008in}{0.865704in}}%
\pgfpathmoveto{\pgfqpoint{1.294926in}{0.871603in}}%
\pgfpathlineto{\pgfqpoint{1.294926in}{0.871603in}}%
\pgfpathlineto{\pgfqpoint{1.294926in}{0.877501in}}%
\pgfpathlineto{\pgfqpoint{1.304008in}{0.877501in}}%
\pgfpathlineto{\pgfqpoint{1.304008in}{0.871603in}}%
\pgfpathmoveto{\pgfqpoint{1.304008in}{0.871603in}}%
\pgfpathlineto{\pgfqpoint{1.304008in}{0.871603in}}%
\pgfpathlineto{\pgfqpoint{1.304008in}{0.877501in}}%
\pgfpathlineto{\pgfqpoint{1.313090in}{0.877501in}}%
\pgfpathlineto{\pgfqpoint{1.313090in}{0.871603in}}%
\pgfpathmoveto{\pgfqpoint{1.313090in}{0.877501in}}%
\pgfpathlineto{\pgfqpoint{1.313090in}{0.877501in}}%
\pgfpathlineto{\pgfqpoint{1.313090in}{0.883400in}}%
\pgfpathlineto{\pgfqpoint{1.322172in}{0.883400in}}%
\pgfpathlineto{\pgfqpoint{1.322172in}{0.877501in}}%
\pgfpathmoveto{\pgfqpoint{1.313090in}{0.883400in}}%
\pgfpathlineto{\pgfqpoint{1.313090in}{0.883400in}}%
\pgfpathlineto{\pgfqpoint{1.313090in}{0.889298in}}%
\pgfpathlineto{\pgfqpoint{1.322172in}{0.889298in}}%
\pgfpathlineto{\pgfqpoint{1.322172in}{0.883400in}}%
\pgfpathmoveto{\pgfqpoint{1.322172in}{0.883400in}}%
\pgfpathlineto{\pgfqpoint{1.322172in}{0.883400in}}%
\pgfpathlineto{\pgfqpoint{1.322172in}{0.889298in}}%
\pgfpathlineto{\pgfqpoint{1.331254in}{0.889298in}}%
\pgfpathlineto{\pgfqpoint{1.331254in}{0.883400in}}%
\pgfpathmoveto{\pgfqpoint{1.331254in}{0.889298in}}%
\pgfpathlineto{\pgfqpoint{1.331254in}{0.889298in}}%
\pgfpathlineto{\pgfqpoint{1.331254in}{0.895196in}}%
\pgfpathlineto{\pgfqpoint{1.340336in}{0.895196in}}%
\pgfpathlineto{\pgfqpoint{1.340336in}{0.889298in}}%
\pgfpathmoveto{\pgfqpoint{1.331254in}{0.895196in}}%
\pgfpathlineto{\pgfqpoint{1.331254in}{0.895196in}}%
\pgfpathlineto{\pgfqpoint{1.331254in}{0.901094in}}%
\pgfpathlineto{\pgfqpoint{1.340336in}{0.901094in}}%
\pgfpathlineto{\pgfqpoint{1.340336in}{0.895196in}}%
\pgfpathmoveto{\pgfqpoint{1.340336in}{0.895196in}}%
\pgfpathlineto{\pgfqpoint{1.340336in}{0.895196in}}%
\pgfpathlineto{\pgfqpoint{1.340336in}{0.901094in}}%
\pgfpathlineto{\pgfqpoint{1.349418in}{0.901094in}}%
\pgfpathlineto{\pgfqpoint{1.349418in}{0.895196in}}%
\pgfpathmoveto{\pgfqpoint{1.349418in}{0.901094in}}%
\pgfpathlineto{\pgfqpoint{1.349418in}{0.901094in}}%
\pgfpathlineto{\pgfqpoint{1.349418in}{0.906992in}}%
\pgfpathlineto{\pgfqpoint{1.358500in}{0.906992in}}%
\pgfpathlineto{\pgfqpoint{1.358500in}{0.901094in}}%
\pgfpathmoveto{\pgfqpoint{1.349418in}{0.906992in}}%
\pgfpathlineto{\pgfqpoint{1.349418in}{0.906992in}}%
\pgfpathlineto{\pgfqpoint{1.349418in}{0.912890in}}%
\pgfpathlineto{\pgfqpoint{1.358500in}{0.912890in}}%
\pgfpathlineto{\pgfqpoint{1.358500in}{0.906992in}}%
\pgfpathmoveto{\pgfqpoint{1.358500in}{0.906992in}}%
\pgfpathlineto{\pgfqpoint{1.358500in}{0.906992in}}%
\pgfpathlineto{\pgfqpoint{1.358500in}{0.912890in}}%
\pgfpathlineto{\pgfqpoint{1.367581in}{0.912890in}}%
\pgfpathlineto{\pgfqpoint{1.367581in}{0.906992in}}%
\pgfpathmoveto{\pgfqpoint{1.367581in}{0.912890in}}%
\pgfpathlineto{\pgfqpoint{1.367581in}{0.912890in}}%
\pgfpathlineto{\pgfqpoint{1.367581in}{0.918789in}}%
\pgfpathlineto{\pgfqpoint{1.376663in}{0.918789in}}%
\pgfpathlineto{\pgfqpoint{1.376663in}{0.912890in}}%
\pgfpathmoveto{\pgfqpoint{1.367581in}{0.918789in}}%
\pgfpathlineto{\pgfqpoint{1.367581in}{0.918789in}}%
\pgfpathlineto{\pgfqpoint{1.367581in}{0.924687in}}%
\pgfpathlineto{\pgfqpoint{1.376663in}{0.924687in}}%
\pgfpathlineto{\pgfqpoint{1.376663in}{0.918789in}}%
\pgfpathmoveto{\pgfqpoint{1.376663in}{0.918789in}}%
\pgfpathlineto{\pgfqpoint{1.376663in}{0.918789in}}%
\pgfpathlineto{\pgfqpoint{1.376663in}{0.924687in}}%
\pgfpathlineto{\pgfqpoint{1.385745in}{0.924687in}}%
\pgfpathlineto{\pgfqpoint{1.385745in}{0.918789in}}%
\pgfpathmoveto{\pgfqpoint{1.385745in}{0.924687in}}%
\pgfpathlineto{\pgfqpoint{1.385745in}{0.924687in}}%
\pgfpathlineto{\pgfqpoint{1.385745in}{0.930585in}}%
\pgfpathlineto{\pgfqpoint{1.394827in}{0.930585in}}%
\pgfpathlineto{\pgfqpoint{1.394827in}{0.924687in}}%
\pgfpathmoveto{\pgfqpoint{1.385745in}{0.930585in}}%
\pgfpathlineto{\pgfqpoint{1.385745in}{0.930585in}}%
\pgfpathlineto{\pgfqpoint{1.385745in}{0.936483in}}%
\pgfpathlineto{\pgfqpoint{1.394827in}{0.936483in}}%
\pgfpathlineto{\pgfqpoint{1.394827in}{0.930585in}}%
\pgfpathmoveto{\pgfqpoint{1.394827in}{0.930585in}}%
\pgfpathlineto{\pgfqpoint{1.394827in}{0.930585in}}%
\pgfpathlineto{\pgfqpoint{1.394827in}{0.936483in}}%
\pgfpathlineto{\pgfqpoint{1.403909in}{0.936483in}}%
\pgfpathlineto{\pgfqpoint{1.403909in}{0.930585in}}%
\pgfpathmoveto{\pgfqpoint{1.403909in}{0.936483in}}%
\pgfpathlineto{\pgfqpoint{1.403909in}{0.936483in}}%
\pgfpathlineto{\pgfqpoint{1.403909in}{0.942381in}}%
\pgfpathlineto{\pgfqpoint{1.412991in}{0.942381in}}%
\pgfpathlineto{\pgfqpoint{1.412991in}{0.936483in}}%
\pgfpathmoveto{\pgfqpoint{1.403909in}{0.942381in}}%
\pgfpathlineto{\pgfqpoint{1.403909in}{0.942381in}}%
\pgfpathlineto{\pgfqpoint{1.403909in}{0.948280in}}%
\pgfpathlineto{\pgfqpoint{1.412991in}{0.948280in}}%
\pgfpathlineto{\pgfqpoint{1.412991in}{0.942381in}}%
\pgfpathmoveto{\pgfqpoint{1.412991in}{0.942381in}}%
\pgfpathlineto{\pgfqpoint{1.412991in}{0.942381in}}%
\pgfpathlineto{\pgfqpoint{1.412991in}{0.948280in}}%
\pgfpathlineto{\pgfqpoint{1.422073in}{0.948280in}}%
\pgfpathlineto{\pgfqpoint{1.422073in}{0.942381in}}%
\pgfpathmoveto{\pgfqpoint{1.422073in}{0.948280in}}%
\pgfpathlineto{\pgfqpoint{1.422073in}{0.948280in}}%
\pgfpathlineto{\pgfqpoint{1.422073in}{0.954178in}}%
\pgfpathlineto{\pgfqpoint{1.431155in}{0.954178in}}%
\pgfpathlineto{\pgfqpoint{1.431155in}{0.948280in}}%
\pgfpathmoveto{\pgfqpoint{1.422073in}{0.954178in}}%
\pgfpathlineto{\pgfqpoint{1.422073in}{0.954178in}}%
\pgfpathlineto{\pgfqpoint{1.422073in}{0.960076in}}%
\pgfpathlineto{\pgfqpoint{1.431155in}{0.960076in}}%
\pgfpathlineto{\pgfqpoint{1.431155in}{0.954178in}}%
\pgfpathmoveto{\pgfqpoint{1.431155in}{0.954178in}}%
\pgfpathlineto{\pgfqpoint{1.431155in}{0.954178in}}%
\pgfpathlineto{\pgfqpoint{1.431155in}{0.960076in}}%
\pgfpathlineto{\pgfqpoint{1.440237in}{0.960076in}}%
\pgfpathlineto{\pgfqpoint{1.440237in}{0.954178in}}%
\pgfpathmoveto{\pgfqpoint{1.440237in}{0.960076in}}%
\pgfpathlineto{\pgfqpoint{1.440237in}{0.960076in}}%
\pgfpathlineto{\pgfqpoint{1.440237in}{0.965974in}}%
\pgfpathlineto{\pgfqpoint{1.449318in}{0.965974in}}%
\pgfpathlineto{\pgfqpoint{1.449318in}{0.960076in}}%
\pgfpathmoveto{\pgfqpoint{1.440237in}{0.965974in}}%
\pgfpathlineto{\pgfqpoint{1.440237in}{0.965974in}}%
\pgfpathlineto{\pgfqpoint{1.440237in}{0.971872in}}%
\pgfpathlineto{\pgfqpoint{1.449318in}{0.971872in}}%
\pgfpathlineto{\pgfqpoint{1.449318in}{0.965974in}}%
\pgfpathmoveto{\pgfqpoint{1.449318in}{0.965974in}}%
\pgfpathlineto{\pgfqpoint{1.449318in}{0.965974in}}%
\pgfpathlineto{\pgfqpoint{1.449318in}{0.971872in}}%
\pgfpathlineto{\pgfqpoint{1.458400in}{0.971872in}}%
\pgfpathlineto{\pgfqpoint{1.458400in}{0.965974in}}%
\pgfpathmoveto{\pgfqpoint{1.458400in}{0.971872in}}%
\pgfpathlineto{\pgfqpoint{1.458400in}{0.971872in}}%
\pgfpathlineto{\pgfqpoint{1.458400in}{0.977771in}}%
\pgfpathlineto{\pgfqpoint{1.467482in}{0.977771in}}%
\pgfpathlineto{\pgfqpoint{1.467482in}{0.971872in}}%
\pgfpathmoveto{\pgfqpoint{1.458400in}{0.977771in}}%
\pgfpathlineto{\pgfqpoint{1.458400in}{0.977771in}}%
\pgfpathlineto{\pgfqpoint{1.458400in}{0.983669in}}%
\pgfpathlineto{\pgfqpoint{1.467482in}{0.983669in}}%
\pgfpathlineto{\pgfqpoint{1.467482in}{0.977771in}}%
\pgfpathmoveto{\pgfqpoint{1.467482in}{0.977771in}}%
\pgfpathlineto{\pgfqpoint{1.467482in}{0.977771in}}%
\pgfpathlineto{\pgfqpoint{1.467482in}{0.983669in}}%
\pgfpathlineto{\pgfqpoint{1.476564in}{0.983669in}}%
\pgfpathlineto{\pgfqpoint{1.476564in}{0.977771in}}%
\pgfpathmoveto{\pgfqpoint{1.476564in}{0.983669in}}%
\pgfpathlineto{\pgfqpoint{1.476564in}{0.983669in}}%
\pgfpathlineto{\pgfqpoint{1.476564in}{0.989568in}}%
\pgfpathlineto{\pgfqpoint{1.485646in}{0.989568in}}%
\pgfpathlineto{\pgfqpoint{1.485646in}{0.983669in}}%
\pgfpathmoveto{\pgfqpoint{1.476564in}{0.989568in}}%
\pgfpathlineto{\pgfqpoint{1.476564in}{0.989568in}}%
\pgfpathlineto{\pgfqpoint{1.476564in}{0.995466in}}%
\pgfpathlineto{\pgfqpoint{1.485646in}{0.995466in}}%
\pgfpathlineto{\pgfqpoint{1.485646in}{0.989568in}}%
\pgfpathmoveto{\pgfqpoint{1.485646in}{0.989568in}}%
\pgfpathlineto{\pgfqpoint{1.485646in}{0.989568in}}%
\pgfpathlineto{\pgfqpoint{1.485646in}{0.995466in}}%
\pgfpathlineto{\pgfqpoint{1.494728in}{0.995466in}}%
\pgfpathlineto{\pgfqpoint{1.494728in}{0.989568in}}%
\pgfpathmoveto{\pgfqpoint{1.494728in}{0.995466in}}%
\pgfpathlineto{\pgfqpoint{1.494728in}{0.995466in}}%
\pgfpathlineto{\pgfqpoint{1.494728in}{1.001365in}}%
\pgfpathlineto{\pgfqpoint{1.503810in}{1.001365in}}%
\pgfpathlineto{\pgfqpoint{1.503810in}{0.995466in}}%
\pgfpathmoveto{\pgfqpoint{1.494728in}{1.001365in}}%
\pgfpathlineto{\pgfqpoint{1.494728in}{1.001365in}}%
\pgfpathlineto{\pgfqpoint{1.494728in}{1.007264in}}%
\pgfpathlineto{\pgfqpoint{1.503810in}{1.007264in}}%
\pgfpathlineto{\pgfqpoint{1.503810in}{1.001365in}}%
\pgfpathmoveto{\pgfqpoint{1.503810in}{1.001365in}}%
\pgfpathlineto{\pgfqpoint{1.503810in}{1.001365in}}%
\pgfpathlineto{\pgfqpoint{1.503810in}{1.007264in}}%
\pgfpathlineto{\pgfqpoint{1.512892in}{1.007264in}}%
\pgfpathlineto{\pgfqpoint{1.512892in}{1.001365in}}%
\pgfpathmoveto{\pgfqpoint{1.512892in}{1.007264in}}%
\pgfpathlineto{\pgfqpoint{1.512892in}{1.007264in}}%
\pgfpathlineto{\pgfqpoint{1.512892in}{1.013162in}}%
\pgfpathlineto{\pgfqpoint{1.521974in}{1.013162in}}%
\pgfpathlineto{\pgfqpoint{1.521974in}{1.007264in}}%
\pgfpathmoveto{\pgfqpoint{1.512892in}{1.013162in}}%
\pgfpathlineto{\pgfqpoint{1.512892in}{1.013162in}}%
\pgfpathlineto{\pgfqpoint{1.512892in}{1.019061in}}%
\pgfpathlineto{\pgfqpoint{1.521974in}{1.019061in}}%
\pgfpathlineto{\pgfqpoint{1.521974in}{1.013162in}}%
\pgfpathmoveto{\pgfqpoint{1.521974in}{1.013162in}}%
\pgfpathlineto{\pgfqpoint{1.521974in}{1.013162in}}%
\pgfpathlineto{\pgfqpoint{1.521974in}{1.019061in}}%
\pgfpathlineto{\pgfqpoint{1.531056in}{1.019061in}}%
\pgfpathlineto{\pgfqpoint{1.531056in}{1.013162in}}%
\pgfpathmoveto{\pgfqpoint{1.531056in}{1.019061in}}%
\pgfpathlineto{\pgfqpoint{1.531056in}{1.019061in}}%
\pgfpathlineto{\pgfqpoint{1.531056in}{1.024959in}}%
\pgfpathlineto{\pgfqpoint{1.540138in}{1.024959in}}%
\pgfpathlineto{\pgfqpoint{1.540138in}{1.019061in}}%
\pgfpathmoveto{\pgfqpoint{1.531056in}{1.024959in}}%
\pgfpathlineto{\pgfqpoint{1.531056in}{1.024959in}}%
\pgfpathlineto{\pgfqpoint{1.531056in}{1.030858in}}%
\pgfpathlineto{\pgfqpoint{1.540138in}{1.030858in}}%
\pgfpathlineto{\pgfqpoint{1.540138in}{1.024959in}}%
\pgfpathmoveto{\pgfqpoint{1.540138in}{1.024959in}}%
\pgfpathlineto{\pgfqpoint{1.540138in}{1.024959in}}%
\pgfpathlineto{\pgfqpoint{1.540138in}{1.030858in}}%
\pgfpathlineto{\pgfqpoint{1.549220in}{1.030858in}}%
\pgfpathlineto{\pgfqpoint{1.549220in}{1.024959in}}%
\pgfpathmoveto{\pgfqpoint{1.549220in}{1.030858in}}%
\pgfpathlineto{\pgfqpoint{1.549220in}{1.030858in}}%
\pgfpathlineto{\pgfqpoint{1.549220in}{1.036756in}}%
\pgfpathlineto{\pgfqpoint{1.558302in}{1.036756in}}%
\pgfpathlineto{\pgfqpoint{1.558302in}{1.030858in}}%
\pgfpathmoveto{\pgfqpoint{1.549220in}{1.036756in}}%
\pgfpathlineto{\pgfqpoint{1.549220in}{1.036756in}}%
\pgfpathlineto{\pgfqpoint{1.549220in}{1.042655in}}%
\pgfpathlineto{\pgfqpoint{1.558302in}{1.042655in}}%
\pgfpathlineto{\pgfqpoint{1.558302in}{1.036756in}}%
\pgfpathmoveto{\pgfqpoint{1.558302in}{1.036756in}}%
\pgfpathlineto{\pgfqpoint{1.558302in}{1.036756in}}%
\pgfpathlineto{\pgfqpoint{1.558302in}{1.042655in}}%
\pgfpathlineto{\pgfqpoint{1.567384in}{1.042655in}}%
\pgfpathlineto{\pgfqpoint{1.567384in}{1.036756in}}%
\pgfpathmoveto{\pgfqpoint{1.567384in}{1.042655in}}%
\pgfpathlineto{\pgfqpoint{1.567384in}{1.042655in}}%
\pgfpathlineto{\pgfqpoint{1.567384in}{1.048553in}}%
\pgfpathlineto{\pgfqpoint{1.576466in}{1.048553in}}%
\pgfpathlineto{\pgfqpoint{1.576466in}{1.042655in}}%
\pgfpathmoveto{\pgfqpoint{1.567384in}{1.048553in}}%
\pgfpathlineto{\pgfqpoint{1.567384in}{1.048553in}}%
\pgfpathlineto{\pgfqpoint{1.567384in}{1.054452in}}%
\pgfpathlineto{\pgfqpoint{1.576466in}{1.054452in}}%
\pgfpathlineto{\pgfqpoint{1.576466in}{1.048553in}}%
\pgfpathmoveto{\pgfqpoint{1.576466in}{1.048553in}}%
\pgfpathlineto{\pgfqpoint{1.576466in}{1.048553in}}%
\pgfpathlineto{\pgfqpoint{1.576466in}{1.054452in}}%
\pgfpathlineto{\pgfqpoint{1.585549in}{1.054452in}}%
\pgfpathlineto{\pgfqpoint{1.585549in}{1.048553in}}%
\pgfpathmoveto{\pgfqpoint{1.585549in}{1.054452in}}%
\pgfpathlineto{\pgfqpoint{1.585549in}{1.054452in}}%
\pgfpathlineto{\pgfqpoint{1.585549in}{1.060350in}}%
\pgfpathlineto{\pgfqpoint{1.594631in}{1.060350in}}%
\pgfpathlineto{\pgfqpoint{1.594631in}{1.054452in}}%
\pgfpathmoveto{\pgfqpoint{1.585549in}{1.060350in}}%
\pgfpathlineto{\pgfqpoint{1.585549in}{1.060350in}}%
\pgfpathlineto{\pgfqpoint{1.585549in}{1.066249in}}%
\pgfpathlineto{\pgfqpoint{1.594631in}{1.066249in}}%
\pgfpathlineto{\pgfqpoint{1.594631in}{1.060350in}}%
\pgfpathmoveto{\pgfqpoint{1.594631in}{1.060350in}}%
\pgfpathlineto{\pgfqpoint{1.594631in}{1.060350in}}%
\pgfpathlineto{\pgfqpoint{1.594631in}{1.066249in}}%
\pgfpathlineto{\pgfqpoint{1.603713in}{1.066249in}}%
\pgfpathlineto{\pgfqpoint{1.603713in}{1.060350in}}%
\pgfpathmoveto{\pgfqpoint{1.603713in}{1.066249in}}%
\pgfpathlineto{\pgfqpoint{1.603713in}{1.066249in}}%
\pgfpathlineto{\pgfqpoint{1.603713in}{1.072147in}}%
\pgfpathlineto{\pgfqpoint{1.612795in}{1.072147in}}%
\pgfpathlineto{\pgfqpoint{1.612795in}{1.066249in}}%
\pgfpathmoveto{\pgfqpoint{1.603713in}{1.072147in}}%
\pgfpathlineto{\pgfqpoint{1.603713in}{1.072147in}}%
\pgfpathlineto{\pgfqpoint{1.603713in}{1.078046in}}%
\pgfpathlineto{\pgfqpoint{1.612795in}{1.078046in}}%
\pgfpathlineto{\pgfqpoint{1.612795in}{1.072147in}}%
\pgfpathmoveto{\pgfqpoint{1.612795in}{1.072147in}}%
\pgfpathlineto{\pgfqpoint{1.612795in}{1.072147in}}%
\pgfpathlineto{\pgfqpoint{1.612795in}{1.078046in}}%
\pgfpathlineto{\pgfqpoint{1.621877in}{1.078046in}}%
\pgfpathlineto{\pgfqpoint{1.621877in}{1.072147in}}%
\pgfpathmoveto{\pgfqpoint{1.621877in}{1.078046in}}%
\pgfpathlineto{\pgfqpoint{1.621877in}{1.078046in}}%
\pgfpathlineto{\pgfqpoint{1.621877in}{1.083944in}}%
\pgfpathlineto{\pgfqpoint{1.630958in}{1.083944in}}%
\pgfpathlineto{\pgfqpoint{1.630958in}{1.078046in}}%
\pgfpathmoveto{\pgfqpoint{1.621877in}{1.083944in}}%
\pgfpathlineto{\pgfqpoint{1.621877in}{1.083944in}}%
\pgfpathlineto{\pgfqpoint{1.621877in}{1.089843in}}%
\pgfpathlineto{\pgfqpoint{1.630958in}{1.089843in}}%
\pgfpathlineto{\pgfqpoint{1.630958in}{1.083944in}}%
\pgfpathmoveto{\pgfqpoint{1.630958in}{1.083944in}}%
\pgfpathlineto{\pgfqpoint{1.630958in}{1.083944in}}%
\pgfpathlineto{\pgfqpoint{1.630958in}{1.089843in}}%
\pgfpathlineto{\pgfqpoint{1.640040in}{1.089843in}}%
\pgfpathlineto{\pgfqpoint{1.640040in}{1.083944in}}%
\pgfpathmoveto{\pgfqpoint{1.640040in}{1.089843in}}%
\pgfpathlineto{\pgfqpoint{1.640040in}{1.089843in}}%
\pgfpathlineto{\pgfqpoint{1.640040in}{1.095741in}}%
\pgfpathlineto{\pgfqpoint{1.649122in}{1.095741in}}%
\pgfpathlineto{\pgfqpoint{1.649122in}{1.089843in}}%
\pgfpathmoveto{\pgfqpoint{1.640040in}{1.095741in}}%
\pgfpathlineto{\pgfqpoint{1.640040in}{1.095741in}}%
\pgfpathlineto{\pgfqpoint{1.640040in}{1.101640in}}%
\pgfpathlineto{\pgfqpoint{1.649122in}{1.101640in}}%
\pgfpathlineto{\pgfqpoint{1.649122in}{1.095741in}}%
\pgfpathmoveto{\pgfqpoint{1.649122in}{1.095741in}}%
\pgfpathlineto{\pgfqpoint{1.649122in}{1.095741in}}%
\pgfpathlineto{\pgfqpoint{1.649122in}{1.101640in}}%
\pgfpathlineto{\pgfqpoint{1.658204in}{1.101640in}}%
\pgfpathlineto{\pgfqpoint{1.658204in}{1.095741in}}%
\pgfpathmoveto{\pgfqpoint{1.658204in}{1.101640in}}%
\pgfpathlineto{\pgfqpoint{1.658204in}{1.101640in}}%
\pgfpathlineto{\pgfqpoint{1.658204in}{1.107538in}}%
\pgfpathlineto{\pgfqpoint{1.667286in}{1.107538in}}%
\pgfpathlineto{\pgfqpoint{1.667286in}{1.101640in}}%
\pgfpathmoveto{\pgfqpoint{1.658204in}{1.107538in}}%
\pgfpathlineto{\pgfqpoint{1.658204in}{1.107538in}}%
\pgfpathlineto{\pgfqpoint{1.658204in}{1.113437in}}%
\pgfpathlineto{\pgfqpoint{1.667286in}{1.113437in}}%
\pgfpathlineto{\pgfqpoint{1.667286in}{1.107538in}}%
\pgfpathmoveto{\pgfqpoint{1.667286in}{1.107538in}}%
\pgfpathlineto{\pgfqpoint{1.667286in}{1.107538in}}%
\pgfpathlineto{\pgfqpoint{1.667286in}{1.113437in}}%
\pgfpathlineto{\pgfqpoint{1.676368in}{1.113437in}}%
\pgfpathlineto{\pgfqpoint{1.676368in}{1.107538in}}%
\pgfpathmoveto{\pgfqpoint{1.676368in}{1.113437in}}%
\pgfpathlineto{\pgfqpoint{1.676368in}{1.113437in}}%
\pgfpathlineto{\pgfqpoint{1.676368in}{1.119335in}}%
\pgfpathlineto{\pgfqpoint{1.685449in}{1.119335in}}%
\pgfpathlineto{\pgfqpoint{1.685449in}{1.113437in}}%
\pgfpathmoveto{\pgfqpoint{1.676368in}{1.119335in}}%
\pgfpathlineto{\pgfqpoint{1.676368in}{1.119335in}}%
\pgfpathlineto{\pgfqpoint{1.676368in}{1.125234in}}%
\pgfpathlineto{\pgfqpoint{1.685449in}{1.125234in}}%
\pgfpathlineto{\pgfqpoint{1.685449in}{1.119335in}}%
\pgfpathmoveto{\pgfqpoint{1.685449in}{1.119335in}}%
\pgfpathlineto{\pgfqpoint{1.685449in}{1.119335in}}%
\pgfpathlineto{\pgfqpoint{1.685449in}{1.125234in}}%
\pgfpathlineto{\pgfqpoint{1.694531in}{1.125234in}}%
\pgfpathlineto{\pgfqpoint{1.694531in}{1.119335in}}%
\pgfpathmoveto{\pgfqpoint{1.694531in}{1.125234in}}%
\pgfpathlineto{\pgfqpoint{1.694531in}{1.125234in}}%
\pgfpathlineto{\pgfqpoint{1.694531in}{1.131132in}}%
\pgfpathlineto{\pgfqpoint{1.703613in}{1.131132in}}%
\pgfpathlineto{\pgfqpoint{1.703613in}{1.125234in}}%
\pgfpathmoveto{\pgfqpoint{1.694531in}{1.131132in}}%
\pgfpathlineto{\pgfqpoint{1.694531in}{1.131132in}}%
\pgfpathlineto{\pgfqpoint{1.694531in}{1.137031in}}%
\pgfpathlineto{\pgfqpoint{1.703613in}{1.137031in}}%
\pgfpathlineto{\pgfqpoint{1.703613in}{1.131132in}}%
\pgfpathmoveto{\pgfqpoint{1.703613in}{1.131132in}}%
\pgfpathlineto{\pgfqpoint{1.703613in}{1.131132in}}%
\pgfpathlineto{\pgfqpoint{1.703613in}{1.137031in}}%
\pgfpathlineto{\pgfqpoint{1.712695in}{1.137031in}}%
\pgfpathlineto{\pgfqpoint{1.712695in}{1.131132in}}%
\pgfpathmoveto{\pgfqpoint{1.712695in}{1.131132in}}%
\pgfpathlineto{\pgfqpoint{1.712695in}{1.131132in}}%
\pgfpathlineto{\pgfqpoint{1.712695in}{1.137031in}}%
\pgfpathlineto{\pgfqpoint{1.721777in}{1.137031in}}%
\pgfpathlineto{\pgfqpoint{1.721777in}{1.131132in}}%
\pgfpathmoveto{\pgfqpoint{1.730859in}{1.142929in}}%
\pgfpathlineto{\pgfqpoint{1.730859in}{1.142929in}}%
\pgfpathlineto{\pgfqpoint{1.730859in}{1.148828in}}%
\pgfpathlineto{\pgfqpoint{1.739940in}{1.148828in}}%
\pgfpathlineto{\pgfqpoint{1.739940in}{1.142929in}}%
\pgfpathmoveto{\pgfqpoint{1.749022in}{1.154727in}}%
\pgfpathlineto{\pgfqpoint{1.749022in}{1.154727in}}%
\pgfpathlineto{\pgfqpoint{1.749022in}{1.160625in}}%
\pgfpathlineto{\pgfqpoint{1.758104in}{1.160625in}}%
\pgfpathlineto{\pgfqpoint{1.758104in}{1.154727in}}%
\pgfpathmoveto{\pgfqpoint{1.767186in}{1.166523in}}%
\pgfpathlineto{\pgfqpoint{1.767186in}{1.166523in}}%
\pgfpathlineto{\pgfqpoint{1.767186in}{1.172422in}}%
\pgfpathlineto{\pgfqpoint{1.776268in}{1.172422in}}%
\pgfpathlineto{\pgfqpoint{1.776268in}{1.166523in}}%
\pgfpathmoveto{\pgfqpoint{1.785350in}{1.178320in}}%
\pgfpathlineto{\pgfqpoint{1.785350in}{1.178320in}}%
\pgfpathlineto{\pgfqpoint{1.785350in}{1.184218in}}%
\pgfpathlineto{\pgfqpoint{1.794432in}{1.184218in}}%
\pgfpathlineto{\pgfqpoint{1.794432in}{1.178320in}}%
\pgfpathmoveto{\pgfqpoint{1.803513in}{1.190117in}}%
\pgfpathlineto{\pgfqpoint{1.803513in}{1.190117in}}%
\pgfpathlineto{\pgfqpoint{1.803513in}{1.196015in}}%
\pgfpathlineto{\pgfqpoint{1.812595in}{1.196015in}}%
\pgfpathlineto{\pgfqpoint{1.812595in}{1.190117in}}%
\pgfpathmoveto{\pgfqpoint{1.821677in}{1.201913in}}%
\pgfpathlineto{\pgfqpoint{1.821677in}{1.201913in}}%
\pgfpathlineto{\pgfqpoint{1.821677in}{1.207812in}}%
\pgfpathlineto{\pgfqpoint{1.830759in}{1.207812in}}%
\pgfpathlineto{\pgfqpoint{1.830759in}{1.201913in}}%
\pgfpathmoveto{\pgfqpoint{1.839841in}{1.213710in}}%
\pgfpathlineto{\pgfqpoint{1.839841in}{1.213710in}}%
\pgfpathlineto{\pgfqpoint{1.839841in}{1.219608in}}%
\pgfpathlineto{\pgfqpoint{1.848923in}{1.219608in}}%
\pgfpathlineto{\pgfqpoint{1.848923in}{1.213710in}}%
\pgfpathmoveto{\pgfqpoint{1.858005in}{1.225506in}}%
\pgfpathlineto{\pgfqpoint{1.858005in}{1.225506in}}%
\pgfpathlineto{\pgfqpoint{1.858005in}{1.231405in}}%
\pgfpathlineto{\pgfqpoint{1.867087in}{1.231405in}}%
\pgfpathlineto{\pgfqpoint{1.867087in}{1.225506in}}%
\pgfpathmoveto{\pgfqpoint{1.876168in}{1.237303in}}%
\pgfpathlineto{\pgfqpoint{1.876168in}{1.237303in}}%
\pgfpathlineto{\pgfqpoint{1.876168in}{1.243201in}}%
\pgfpathlineto{\pgfqpoint{1.885250in}{1.243201in}}%
\pgfpathlineto{\pgfqpoint{1.885250in}{1.237303in}}%
\pgfpathmoveto{\pgfqpoint{1.894332in}{1.249100in}}%
\pgfpathlineto{\pgfqpoint{1.894332in}{1.249100in}}%
\pgfpathlineto{\pgfqpoint{1.894332in}{1.254998in}}%
\pgfpathlineto{\pgfqpoint{1.903414in}{1.254998in}}%
\pgfpathlineto{\pgfqpoint{1.903414in}{1.249100in}}%
\pgfpathmoveto{\pgfqpoint{1.912496in}{1.260897in}}%
\pgfpathlineto{\pgfqpoint{1.912496in}{1.260897in}}%
\pgfpathlineto{\pgfqpoint{1.912496in}{1.266795in}}%
\pgfpathlineto{\pgfqpoint{1.921578in}{1.266795in}}%
\pgfpathlineto{\pgfqpoint{1.921578in}{1.260897in}}%
\pgfpathmoveto{\pgfqpoint{1.930661in}{1.272694in}}%
\pgfpathlineto{\pgfqpoint{1.930661in}{1.272694in}}%
\pgfpathlineto{\pgfqpoint{1.930661in}{1.278592in}}%
\pgfpathlineto{\pgfqpoint{1.939743in}{1.278592in}}%
\pgfpathlineto{\pgfqpoint{1.939743in}{1.272694in}}%
\pgfpathmoveto{\pgfqpoint{1.948825in}{1.284491in}}%
\pgfpathlineto{\pgfqpoint{1.948825in}{1.284491in}}%
\pgfpathlineto{\pgfqpoint{1.948825in}{1.290390in}}%
\pgfpathlineto{\pgfqpoint{1.957908in}{1.290390in}}%
\pgfpathlineto{\pgfqpoint{1.957908in}{1.284491in}}%
\pgfpathmoveto{\pgfqpoint{1.966990in}{1.296288in}}%
\pgfpathlineto{\pgfqpoint{1.966990in}{1.296288in}}%
\pgfpathlineto{\pgfqpoint{1.966990in}{1.302187in}}%
\pgfpathlineto{\pgfqpoint{1.976072in}{1.302187in}}%
\pgfpathlineto{\pgfqpoint{1.976072in}{1.296288in}}%
\pgfpathmoveto{\pgfqpoint{1.985155in}{1.308086in}}%
\pgfpathlineto{\pgfqpoint{1.985155in}{1.308086in}}%
\pgfpathlineto{\pgfqpoint{1.985155in}{1.313984in}}%
\pgfpathlineto{\pgfqpoint{1.994237in}{1.313984in}}%
\pgfpathlineto{\pgfqpoint{1.994237in}{1.308086in}}%
\pgfpathmoveto{\pgfqpoint{2.003319in}{1.319883in}}%
\pgfpathlineto{\pgfqpoint{2.003319in}{1.319883in}}%
\pgfpathlineto{\pgfqpoint{2.003319in}{1.325781in}}%
\pgfpathlineto{\pgfqpoint{2.012402in}{1.325781in}}%
\pgfpathlineto{\pgfqpoint{2.012402in}{1.319883in}}%
\pgfpathmoveto{\pgfqpoint{2.021484in}{1.331680in}}%
\pgfpathlineto{\pgfqpoint{2.021484in}{1.331680in}}%
\pgfpathlineto{\pgfqpoint{2.021484in}{1.337579in}}%
\pgfpathlineto{\pgfqpoint{2.030566in}{1.337579in}}%
\pgfpathlineto{\pgfqpoint{2.030566in}{1.331680in}}%
\pgfpathmoveto{\pgfqpoint{2.039649in}{1.343477in}}%
\pgfpathlineto{\pgfqpoint{2.039649in}{1.343477in}}%
\pgfpathlineto{\pgfqpoint{2.039649in}{1.349376in}}%
\pgfpathlineto{\pgfqpoint{2.048731in}{1.349376in}}%
\pgfpathlineto{\pgfqpoint{2.048731in}{1.343477in}}%
\pgfpathmoveto{\pgfqpoint{2.057813in}{1.355274in}}%
\pgfpathlineto{\pgfqpoint{2.057813in}{1.355274in}}%
\pgfpathlineto{\pgfqpoint{2.057813in}{1.361173in}}%
\pgfpathlineto{\pgfqpoint{2.066896in}{1.361173in}}%
\pgfpathlineto{\pgfqpoint{2.066896in}{1.355274in}}%
\pgfpathmoveto{\pgfqpoint{2.057813in}{1.361173in}}%
\pgfpathlineto{\pgfqpoint{2.057813in}{1.361173in}}%
\pgfpathlineto{\pgfqpoint{2.057813in}{1.367071in}}%
\pgfpathlineto{\pgfqpoint{2.066896in}{1.367071in}}%
\pgfpathlineto{\pgfqpoint{2.066896in}{1.361173in}}%
\pgfpathmoveto{\pgfqpoint{2.057813in}{1.367071in}}%
\pgfpathlineto{\pgfqpoint{2.057813in}{1.367071in}}%
\pgfpathlineto{\pgfqpoint{2.057813in}{1.372969in}}%
\pgfpathlineto{\pgfqpoint{2.066896in}{1.372969in}}%
\pgfpathlineto{\pgfqpoint{2.066896in}{1.367071in}}%
\pgfpathmoveto{\pgfqpoint{2.066896in}{1.367071in}}%
\pgfpathlineto{\pgfqpoint{2.066896in}{1.367071in}}%
\pgfpathlineto{\pgfqpoint{2.066896in}{1.372969in}}%
\pgfpathlineto{\pgfqpoint{2.075978in}{1.372969in}}%
\pgfpathlineto{\pgfqpoint{2.075978in}{1.367071in}}%
\pgfpathmoveto{\pgfqpoint{2.075978in}{1.372969in}}%
\pgfpathlineto{\pgfqpoint{2.075978in}{1.372969in}}%
\pgfpathlineto{\pgfqpoint{2.075978in}{1.378868in}}%
\pgfpathlineto{\pgfqpoint{2.085060in}{1.378868in}}%
\pgfpathlineto{\pgfqpoint{2.085060in}{1.372969in}}%
\pgfpathmoveto{\pgfqpoint{2.075978in}{1.378868in}}%
\pgfpathlineto{\pgfqpoint{2.075978in}{1.378868in}}%
\pgfpathlineto{\pgfqpoint{2.075978in}{1.384766in}}%
\pgfpathlineto{\pgfqpoint{2.085060in}{1.384766in}}%
\pgfpathlineto{\pgfqpoint{2.085060in}{1.378868in}}%
\pgfpathmoveto{\pgfqpoint{2.085060in}{1.378868in}}%
\pgfpathlineto{\pgfqpoint{2.085060in}{1.378868in}}%
\pgfpathlineto{\pgfqpoint{2.085060in}{1.384766in}}%
\pgfpathlineto{\pgfqpoint{2.094142in}{1.384766in}}%
\pgfpathlineto{\pgfqpoint{2.094142in}{1.378868in}}%
\pgfpathmoveto{\pgfqpoint{2.094142in}{1.384766in}}%
\pgfpathlineto{\pgfqpoint{2.094142in}{1.384766in}}%
\pgfpathlineto{\pgfqpoint{2.094142in}{1.390665in}}%
\pgfpathlineto{\pgfqpoint{2.103225in}{1.390665in}}%
\pgfpathlineto{\pgfqpoint{2.103225in}{1.384766in}}%
\pgfpathmoveto{\pgfqpoint{2.094142in}{1.390665in}}%
\pgfpathlineto{\pgfqpoint{2.094142in}{1.390665in}}%
\pgfpathlineto{\pgfqpoint{2.094142in}{1.396563in}}%
\pgfpathlineto{\pgfqpoint{2.103225in}{1.396563in}}%
\pgfpathlineto{\pgfqpoint{2.103225in}{1.390665in}}%
\pgfpathmoveto{\pgfqpoint{2.103225in}{1.390665in}}%
\pgfpathlineto{\pgfqpoint{2.103225in}{1.390665in}}%
\pgfpathlineto{\pgfqpoint{2.103225in}{1.396563in}}%
\pgfpathlineto{\pgfqpoint{2.112307in}{1.396563in}}%
\pgfpathlineto{\pgfqpoint{2.112307in}{1.390665in}}%
\pgfpathmoveto{\pgfqpoint{2.112307in}{1.396563in}}%
\pgfpathlineto{\pgfqpoint{2.112307in}{1.396563in}}%
\pgfpathlineto{\pgfqpoint{2.112307in}{1.402461in}}%
\pgfpathlineto{\pgfqpoint{2.121389in}{1.402461in}}%
\pgfpathlineto{\pgfqpoint{2.121389in}{1.396563in}}%
\pgfpathmoveto{\pgfqpoint{2.112307in}{1.402461in}}%
\pgfpathlineto{\pgfqpoint{2.112307in}{1.402461in}}%
\pgfpathlineto{\pgfqpoint{2.112307in}{1.408360in}}%
\pgfpathlineto{\pgfqpoint{2.121389in}{1.408360in}}%
\pgfpathlineto{\pgfqpoint{2.121389in}{1.402461in}}%
\pgfpathmoveto{\pgfqpoint{2.121389in}{1.402461in}}%
\pgfpathlineto{\pgfqpoint{2.121389in}{1.402461in}}%
\pgfpathlineto{\pgfqpoint{2.121389in}{1.408360in}}%
\pgfpathlineto{\pgfqpoint{2.130471in}{1.408360in}}%
\pgfpathlineto{\pgfqpoint{2.130471in}{1.402461in}}%
\pgfpathmoveto{\pgfqpoint{2.130471in}{1.408360in}}%
\pgfpathlineto{\pgfqpoint{2.130471in}{1.408360in}}%
\pgfpathlineto{\pgfqpoint{2.130471in}{1.414258in}}%
\pgfpathlineto{\pgfqpoint{2.139554in}{1.414258in}}%
\pgfpathlineto{\pgfqpoint{2.139554in}{1.408360in}}%
\pgfpathmoveto{\pgfqpoint{2.130471in}{1.414258in}}%
\pgfpathlineto{\pgfqpoint{2.130471in}{1.414258in}}%
\pgfpathlineto{\pgfqpoint{2.130471in}{1.420156in}}%
\pgfpathlineto{\pgfqpoint{2.139554in}{1.420156in}}%
\pgfpathlineto{\pgfqpoint{2.139554in}{1.414258in}}%
\pgfpathmoveto{\pgfqpoint{2.139554in}{1.414258in}}%
\pgfpathlineto{\pgfqpoint{2.139554in}{1.414258in}}%
\pgfpathlineto{\pgfqpoint{2.139554in}{1.420156in}}%
\pgfpathlineto{\pgfqpoint{2.148636in}{1.420156in}}%
\pgfpathlineto{\pgfqpoint{2.148636in}{1.414258in}}%
\pgfpathmoveto{\pgfqpoint{2.148636in}{1.420156in}}%
\pgfpathlineto{\pgfqpoint{2.148636in}{1.420156in}}%
\pgfpathlineto{\pgfqpoint{2.148636in}{1.426055in}}%
\pgfpathlineto{\pgfqpoint{2.157718in}{1.426055in}}%
\pgfpathlineto{\pgfqpoint{2.157718in}{1.420156in}}%
\pgfpathmoveto{\pgfqpoint{2.148636in}{1.426055in}}%
\pgfpathlineto{\pgfqpoint{2.148636in}{1.426055in}}%
\pgfpathlineto{\pgfqpoint{2.148636in}{1.431953in}}%
\pgfpathlineto{\pgfqpoint{2.157718in}{1.431953in}}%
\pgfpathlineto{\pgfqpoint{2.157718in}{1.426055in}}%
\pgfpathmoveto{\pgfqpoint{2.157718in}{1.426055in}}%
\pgfpathlineto{\pgfqpoint{2.157718in}{1.426055in}}%
\pgfpathlineto{\pgfqpoint{2.157718in}{1.431953in}}%
\pgfpathlineto{\pgfqpoint{2.166800in}{1.431953in}}%
\pgfpathlineto{\pgfqpoint{2.166800in}{1.426055in}}%
\pgfpathmoveto{\pgfqpoint{2.166800in}{1.431953in}}%
\pgfpathlineto{\pgfqpoint{2.166800in}{1.431953in}}%
\pgfpathlineto{\pgfqpoint{2.166800in}{1.437852in}}%
\pgfpathlineto{\pgfqpoint{2.175882in}{1.437852in}}%
\pgfpathlineto{\pgfqpoint{2.175882in}{1.431953in}}%
\pgfpathmoveto{\pgfqpoint{2.166800in}{1.437852in}}%
\pgfpathlineto{\pgfqpoint{2.166800in}{1.437852in}}%
\pgfpathlineto{\pgfqpoint{2.166800in}{1.443750in}}%
\pgfpathlineto{\pgfqpoint{2.175882in}{1.443750in}}%
\pgfpathlineto{\pgfqpoint{2.175882in}{1.437852in}}%
\pgfpathmoveto{\pgfqpoint{2.175882in}{1.437852in}}%
\pgfpathlineto{\pgfqpoint{2.175882in}{1.437852in}}%
\pgfpathlineto{\pgfqpoint{2.175882in}{1.443750in}}%
\pgfpathlineto{\pgfqpoint{2.184965in}{1.443750in}}%
\pgfpathlineto{\pgfqpoint{2.184965in}{1.437852in}}%
\pgfpathmoveto{\pgfqpoint{2.184965in}{1.443750in}}%
\pgfpathlineto{\pgfqpoint{2.184965in}{1.443750in}}%
\pgfpathlineto{\pgfqpoint{2.184965in}{1.449649in}}%
\pgfpathlineto{\pgfqpoint{2.194047in}{1.449649in}}%
\pgfpathlineto{\pgfqpoint{2.194047in}{1.443750in}}%
\pgfpathmoveto{\pgfqpoint{2.184965in}{1.449649in}}%
\pgfpathlineto{\pgfqpoint{2.184965in}{1.449649in}}%
\pgfpathlineto{\pgfqpoint{2.184965in}{1.455547in}}%
\pgfpathlineto{\pgfqpoint{2.194047in}{1.455547in}}%
\pgfpathlineto{\pgfqpoint{2.194047in}{1.449649in}}%
\pgfpathmoveto{\pgfqpoint{2.194047in}{1.449649in}}%
\pgfpathlineto{\pgfqpoint{2.194047in}{1.449649in}}%
\pgfpathlineto{\pgfqpoint{2.194047in}{1.455547in}}%
\pgfpathlineto{\pgfqpoint{2.203129in}{1.455547in}}%
\pgfpathlineto{\pgfqpoint{2.203129in}{1.449649in}}%
\pgfpathmoveto{\pgfqpoint{2.203129in}{1.455547in}}%
\pgfpathlineto{\pgfqpoint{2.203129in}{1.455547in}}%
\pgfpathlineto{\pgfqpoint{2.203129in}{1.461446in}}%
\pgfpathlineto{\pgfqpoint{2.212211in}{1.461446in}}%
\pgfpathlineto{\pgfqpoint{2.212211in}{1.455547in}}%
\pgfpathmoveto{\pgfqpoint{2.203129in}{1.461446in}}%
\pgfpathlineto{\pgfqpoint{2.203129in}{1.461446in}}%
\pgfpathlineto{\pgfqpoint{2.203129in}{1.467344in}}%
\pgfpathlineto{\pgfqpoint{2.212211in}{1.467344in}}%
\pgfpathlineto{\pgfqpoint{2.212211in}{1.461446in}}%
\pgfpathmoveto{\pgfqpoint{2.212211in}{1.461446in}}%
\pgfpathlineto{\pgfqpoint{2.212211in}{1.461446in}}%
\pgfpathlineto{\pgfqpoint{2.212211in}{1.467344in}}%
\pgfpathlineto{\pgfqpoint{2.221293in}{1.467344in}}%
\pgfpathlineto{\pgfqpoint{2.221293in}{1.461446in}}%
\pgfpathmoveto{\pgfqpoint{2.221293in}{1.467344in}}%
\pgfpathlineto{\pgfqpoint{2.221293in}{1.467344in}}%
\pgfpathlineto{\pgfqpoint{2.221293in}{1.473243in}}%
\pgfpathlineto{\pgfqpoint{2.230375in}{1.473243in}}%
\pgfpathlineto{\pgfqpoint{2.230375in}{1.467344in}}%
\pgfpathmoveto{\pgfqpoint{2.221293in}{1.473243in}}%
\pgfpathlineto{\pgfqpoint{2.221293in}{1.473243in}}%
\pgfpathlineto{\pgfqpoint{2.221293in}{1.479141in}}%
\pgfpathlineto{\pgfqpoint{2.230375in}{1.479141in}}%
\pgfpathlineto{\pgfqpoint{2.230375in}{1.473243in}}%
\pgfpathmoveto{\pgfqpoint{2.230375in}{1.473243in}}%
\pgfpathlineto{\pgfqpoint{2.230375in}{1.473243in}}%
\pgfpathlineto{\pgfqpoint{2.230375in}{1.479141in}}%
\pgfpathlineto{\pgfqpoint{2.239457in}{1.479141in}}%
\pgfpathlineto{\pgfqpoint{2.239457in}{1.473243in}}%
\pgfpathmoveto{\pgfqpoint{2.239457in}{1.479141in}}%
\pgfpathlineto{\pgfqpoint{2.239457in}{1.479141in}}%
\pgfpathlineto{\pgfqpoint{2.239457in}{1.485040in}}%
\pgfpathlineto{\pgfqpoint{2.248539in}{1.485040in}}%
\pgfpathlineto{\pgfqpoint{2.248539in}{1.479141in}}%
\pgfpathmoveto{\pgfqpoint{2.239457in}{1.485040in}}%
\pgfpathlineto{\pgfqpoint{2.239457in}{1.485040in}}%
\pgfpathlineto{\pgfqpoint{2.239457in}{1.490938in}}%
\pgfpathlineto{\pgfqpoint{2.248539in}{1.490938in}}%
\pgfpathlineto{\pgfqpoint{2.248539in}{1.485040in}}%
\pgfpathmoveto{\pgfqpoint{2.248539in}{1.485040in}}%
\pgfpathlineto{\pgfqpoint{2.248539in}{1.485040in}}%
\pgfpathlineto{\pgfqpoint{2.248539in}{1.490938in}}%
\pgfpathlineto{\pgfqpoint{2.257621in}{1.490938in}}%
\pgfpathlineto{\pgfqpoint{2.257621in}{1.485040in}}%
\pgfpathmoveto{\pgfqpoint{2.257621in}{1.490938in}}%
\pgfpathlineto{\pgfqpoint{2.257621in}{1.490938in}}%
\pgfpathlineto{\pgfqpoint{2.257621in}{1.496837in}}%
\pgfpathlineto{\pgfqpoint{2.266703in}{1.496837in}}%
\pgfpathlineto{\pgfqpoint{2.266703in}{1.490938in}}%
\pgfpathmoveto{\pgfqpoint{2.257621in}{1.496837in}}%
\pgfpathlineto{\pgfqpoint{2.257621in}{1.496837in}}%
\pgfpathlineto{\pgfqpoint{2.257621in}{1.502735in}}%
\pgfpathlineto{\pgfqpoint{2.266703in}{1.502735in}}%
\pgfpathlineto{\pgfqpoint{2.266703in}{1.496837in}}%
\pgfpathmoveto{\pgfqpoint{2.266703in}{1.496837in}}%
\pgfpathlineto{\pgfqpoint{2.266703in}{1.496837in}}%
\pgfpathlineto{\pgfqpoint{2.266703in}{1.502735in}}%
\pgfpathlineto{\pgfqpoint{2.275785in}{1.502735in}}%
\pgfpathlineto{\pgfqpoint{2.275785in}{1.496837in}}%
\pgfpathmoveto{\pgfqpoint{2.275785in}{1.502735in}}%
\pgfpathlineto{\pgfqpoint{2.275785in}{1.502735in}}%
\pgfpathlineto{\pgfqpoint{2.275785in}{1.508634in}}%
\pgfpathlineto{\pgfqpoint{2.284867in}{1.508634in}}%
\pgfpathlineto{\pgfqpoint{2.284867in}{1.502735in}}%
\pgfpathmoveto{\pgfqpoint{2.275785in}{1.508634in}}%
\pgfpathlineto{\pgfqpoint{2.275785in}{1.508634in}}%
\pgfpathlineto{\pgfqpoint{2.275785in}{1.514533in}}%
\pgfpathlineto{\pgfqpoint{2.284867in}{1.514533in}}%
\pgfpathlineto{\pgfqpoint{2.284867in}{1.508634in}}%
\pgfpathmoveto{\pgfqpoint{2.284867in}{1.508634in}}%
\pgfpathlineto{\pgfqpoint{2.284867in}{1.508634in}}%
\pgfpathlineto{\pgfqpoint{2.284867in}{1.514533in}}%
\pgfpathlineto{\pgfqpoint{2.293949in}{1.514533in}}%
\pgfpathlineto{\pgfqpoint{2.293949in}{1.508634in}}%
\pgfpathmoveto{\pgfqpoint{2.293949in}{1.514533in}}%
\pgfpathlineto{\pgfqpoint{2.293949in}{1.514533in}}%
\pgfpathlineto{\pgfqpoint{2.293949in}{1.520431in}}%
\pgfpathlineto{\pgfqpoint{2.303031in}{1.520431in}}%
\pgfpathlineto{\pgfqpoint{2.303031in}{1.514533in}}%
\pgfpathmoveto{\pgfqpoint{2.293949in}{1.520431in}}%
\pgfpathlineto{\pgfqpoint{2.293949in}{1.520431in}}%
\pgfpathlineto{\pgfqpoint{2.293949in}{1.526330in}}%
\pgfpathlineto{\pgfqpoint{2.303031in}{1.526330in}}%
\pgfpathlineto{\pgfqpoint{2.303031in}{1.520431in}}%
\pgfpathmoveto{\pgfqpoint{2.303031in}{1.520431in}}%
\pgfpathlineto{\pgfqpoint{2.303031in}{1.520431in}}%
\pgfpathlineto{\pgfqpoint{2.303031in}{1.526330in}}%
\pgfpathlineto{\pgfqpoint{2.312113in}{1.526330in}}%
\pgfpathlineto{\pgfqpoint{2.312113in}{1.520431in}}%
\pgfpathmoveto{\pgfqpoint{2.312113in}{1.526330in}}%
\pgfpathlineto{\pgfqpoint{2.312113in}{1.526330in}}%
\pgfpathlineto{\pgfqpoint{2.312113in}{1.532228in}}%
\pgfpathlineto{\pgfqpoint{2.321195in}{1.532228in}}%
\pgfpathlineto{\pgfqpoint{2.321195in}{1.526330in}}%
\pgfpathmoveto{\pgfqpoint{2.312113in}{1.532228in}}%
\pgfpathlineto{\pgfqpoint{2.312113in}{1.532228in}}%
\pgfpathlineto{\pgfqpoint{2.312113in}{1.538127in}}%
\pgfpathlineto{\pgfqpoint{2.321195in}{1.538127in}}%
\pgfpathlineto{\pgfqpoint{2.321195in}{1.532228in}}%
\pgfpathmoveto{\pgfqpoint{2.321195in}{1.532228in}}%
\pgfpathlineto{\pgfqpoint{2.321195in}{1.532228in}}%
\pgfpathlineto{\pgfqpoint{2.321195in}{1.538127in}}%
\pgfpathlineto{\pgfqpoint{2.330277in}{1.538127in}}%
\pgfpathlineto{\pgfqpoint{2.330277in}{1.532228in}}%
\pgfpathmoveto{\pgfqpoint{2.330277in}{1.538127in}}%
\pgfpathlineto{\pgfqpoint{2.330277in}{1.538127in}}%
\pgfpathlineto{\pgfqpoint{2.330277in}{1.544025in}}%
\pgfpathlineto{\pgfqpoint{2.339359in}{1.544025in}}%
\pgfpathlineto{\pgfqpoint{2.339359in}{1.538127in}}%
\pgfpathmoveto{\pgfqpoint{2.330277in}{1.544025in}}%
\pgfpathlineto{\pgfqpoint{2.330277in}{1.544025in}}%
\pgfpathlineto{\pgfqpoint{2.330277in}{1.549924in}}%
\pgfpathlineto{\pgfqpoint{2.339359in}{1.549924in}}%
\pgfpathlineto{\pgfqpoint{2.339359in}{1.544025in}}%
\pgfpathmoveto{\pgfqpoint{2.339359in}{1.544025in}}%
\pgfpathlineto{\pgfqpoint{2.339359in}{1.544025in}}%
\pgfpathlineto{\pgfqpoint{2.339359in}{1.549924in}}%
\pgfpathlineto{\pgfqpoint{2.348441in}{1.549924in}}%
\pgfpathlineto{\pgfqpoint{2.348441in}{1.544025in}}%
\pgfpathmoveto{\pgfqpoint{2.348441in}{1.549924in}}%
\pgfpathlineto{\pgfqpoint{2.348441in}{1.549924in}}%
\pgfpathlineto{\pgfqpoint{2.348441in}{1.555822in}}%
\pgfpathlineto{\pgfqpoint{2.357522in}{1.555822in}}%
\pgfpathlineto{\pgfqpoint{2.357522in}{1.549924in}}%
\pgfpathmoveto{\pgfqpoint{2.348441in}{1.555822in}}%
\pgfpathlineto{\pgfqpoint{2.348441in}{1.555822in}}%
\pgfpathlineto{\pgfqpoint{2.348441in}{1.561720in}}%
\pgfpathlineto{\pgfqpoint{2.357522in}{1.561720in}}%
\pgfpathlineto{\pgfqpoint{2.357522in}{1.555822in}}%
\pgfpathmoveto{\pgfqpoint{2.357522in}{1.555822in}}%
\pgfpathlineto{\pgfqpoint{2.357522in}{1.555822in}}%
\pgfpathlineto{\pgfqpoint{2.357522in}{1.561720in}}%
\pgfpathlineto{\pgfqpoint{2.366604in}{1.561720in}}%
\pgfpathlineto{\pgfqpoint{2.366604in}{1.555822in}}%
\pgfpathmoveto{\pgfqpoint{2.366604in}{1.555822in}}%
\pgfpathlineto{\pgfqpoint{2.366604in}{1.555822in}}%
\pgfpathlineto{\pgfqpoint{2.366604in}{1.561720in}}%
\pgfpathlineto{\pgfqpoint{2.375686in}{1.561720in}}%
\pgfpathlineto{\pgfqpoint{2.375686in}{1.555822in}}%
\pgfpathmoveto{\pgfqpoint{2.384768in}{1.567619in}}%
\pgfpathlineto{\pgfqpoint{2.384768in}{1.567619in}}%
\pgfpathlineto{\pgfqpoint{2.384768in}{1.573517in}}%
\pgfpathlineto{\pgfqpoint{2.393850in}{1.573517in}}%
\pgfpathlineto{\pgfqpoint{2.393850in}{1.567619in}}%
\pgfpathmoveto{\pgfqpoint{2.402931in}{1.579416in}}%
\pgfpathlineto{\pgfqpoint{2.402931in}{1.579416in}}%
\pgfpathlineto{\pgfqpoint{2.402931in}{1.585314in}}%
\pgfpathlineto{\pgfqpoint{2.412013in}{1.585314in}}%
\pgfpathlineto{\pgfqpoint{2.412013in}{1.579416in}}%
\pgfpathmoveto{\pgfqpoint{2.421095in}{1.591212in}}%
\pgfpathlineto{\pgfqpoint{2.421095in}{1.591212in}}%
\pgfpathlineto{\pgfqpoint{2.421095in}{1.597111in}}%
\pgfpathlineto{\pgfqpoint{2.430177in}{1.597111in}}%
\pgfpathlineto{\pgfqpoint{2.430177in}{1.591212in}}%
\pgfpathmoveto{\pgfqpoint{2.439259in}{1.603009in}}%
\pgfpathlineto{\pgfqpoint{2.439259in}{1.603009in}}%
\pgfpathlineto{\pgfqpoint{2.439259in}{1.608908in}}%
\pgfpathlineto{\pgfqpoint{2.448341in}{1.608908in}}%
\pgfpathlineto{\pgfqpoint{2.448341in}{1.603009in}}%
\pgfpathmoveto{\pgfqpoint{2.457422in}{1.614806in}}%
\pgfpathlineto{\pgfqpoint{2.457422in}{1.614806in}}%
\pgfpathlineto{\pgfqpoint{2.457422in}{1.620704in}}%
\pgfpathlineto{\pgfqpoint{2.466504in}{1.620704in}}%
\pgfpathlineto{\pgfqpoint{2.466504in}{1.614806in}}%
\pgfpathmoveto{\pgfqpoint{2.475586in}{1.626603in}}%
\pgfpathlineto{\pgfqpoint{2.475586in}{1.626603in}}%
\pgfpathlineto{\pgfqpoint{2.475586in}{1.632501in}}%
\pgfpathlineto{\pgfqpoint{2.484668in}{1.632501in}}%
\pgfpathlineto{\pgfqpoint{2.484668in}{1.626603in}}%
\pgfpathmoveto{\pgfqpoint{2.493750in}{1.638400in}}%
\pgfpathlineto{\pgfqpoint{2.493750in}{1.638400in}}%
\pgfpathlineto{\pgfqpoint{2.493750in}{1.644298in}}%
\pgfpathlineto{\pgfqpoint{2.502832in}{1.644298in}}%
\pgfpathlineto{\pgfqpoint{2.502832in}{1.638400in}}%
\pgfpathmoveto{\pgfqpoint{2.511914in}{1.650196in}}%
\pgfpathlineto{\pgfqpoint{2.511914in}{1.650196in}}%
\pgfpathlineto{\pgfqpoint{2.511914in}{1.656095in}}%
\pgfpathlineto{\pgfqpoint{2.520996in}{1.656095in}}%
\pgfpathlineto{\pgfqpoint{2.520996in}{1.650196in}}%
\pgfpathmoveto{\pgfqpoint{2.530078in}{1.661993in}}%
\pgfpathlineto{\pgfqpoint{2.530078in}{1.661993in}}%
\pgfpathlineto{\pgfqpoint{2.530078in}{1.667892in}}%
\pgfpathlineto{\pgfqpoint{2.539160in}{1.667892in}}%
\pgfpathlineto{\pgfqpoint{2.539160in}{1.661993in}}%
\pgfpathmoveto{\pgfqpoint{2.548242in}{1.673790in}}%
\pgfpathlineto{\pgfqpoint{2.548242in}{1.673790in}}%
\pgfpathlineto{\pgfqpoint{2.548242in}{1.679689in}}%
\pgfpathlineto{\pgfqpoint{2.557324in}{1.679689in}}%
\pgfpathlineto{\pgfqpoint{2.557324in}{1.673790in}}%
\pgfpathmoveto{\pgfqpoint{2.566406in}{1.685587in}}%
\pgfpathlineto{\pgfqpoint{2.566406in}{1.685587in}}%
\pgfpathlineto{\pgfqpoint{2.566406in}{1.691485in}}%
\pgfpathlineto{\pgfqpoint{2.575489in}{1.691485in}}%
\pgfpathlineto{\pgfqpoint{2.575489in}{1.685587in}}%
\pgfpathmoveto{\pgfqpoint{2.584571in}{1.697384in}}%
\pgfpathlineto{\pgfqpoint{2.584571in}{1.697384in}}%
\pgfpathlineto{\pgfqpoint{2.584571in}{1.703282in}}%
\pgfpathlineto{\pgfqpoint{2.593653in}{1.703282in}}%
\pgfpathlineto{\pgfqpoint{2.593653in}{1.697384in}}%
\pgfpathmoveto{\pgfqpoint{2.602735in}{1.709181in}}%
\pgfpathlineto{\pgfqpoint{2.602735in}{1.709181in}}%
\pgfpathlineto{\pgfqpoint{2.602735in}{1.715079in}}%
\pgfpathlineto{\pgfqpoint{2.611817in}{1.715079in}}%
\pgfpathlineto{\pgfqpoint{2.611817in}{1.709181in}}%
\pgfpathmoveto{\pgfqpoint{2.620899in}{1.720977in}}%
\pgfpathlineto{\pgfqpoint{2.620899in}{1.720977in}}%
\pgfpathlineto{\pgfqpoint{2.620899in}{1.726876in}}%
\pgfpathlineto{\pgfqpoint{2.629981in}{1.726876in}}%
\pgfpathlineto{\pgfqpoint{2.629981in}{1.720977in}}%
\pgfpathmoveto{\pgfqpoint{2.639063in}{1.732774in}}%
\pgfpathlineto{\pgfqpoint{2.639063in}{1.732774in}}%
\pgfpathlineto{\pgfqpoint{2.639063in}{1.738672in}}%
\pgfpathlineto{\pgfqpoint{2.648145in}{1.738672in}}%
\pgfpathlineto{\pgfqpoint{2.648145in}{1.732774in}}%
\pgfpathmoveto{\pgfqpoint{2.657227in}{1.744571in}}%
\pgfpathlineto{\pgfqpoint{2.657227in}{1.744571in}}%
\pgfpathlineto{\pgfqpoint{2.657227in}{1.750469in}}%
\pgfpathlineto{\pgfqpoint{2.666309in}{1.750469in}}%
\pgfpathlineto{\pgfqpoint{2.666309in}{1.744571in}}%
\pgfpathmoveto{\pgfqpoint{2.657227in}{1.750469in}}%
\pgfpathlineto{\pgfqpoint{2.657227in}{1.750469in}}%
\pgfpathlineto{\pgfqpoint{2.657227in}{1.756367in}}%
\pgfpathlineto{\pgfqpoint{2.666309in}{1.756367in}}%
\pgfpathlineto{\pgfqpoint{2.666309in}{1.750469in}}%
\pgfpathmoveto{\pgfqpoint{2.657227in}{1.756367in}}%
\pgfpathlineto{\pgfqpoint{2.657227in}{1.756367in}}%
\pgfpathlineto{\pgfqpoint{2.657227in}{1.762266in}}%
\pgfpathlineto{\pgfqpoint{2.666309in}{1.762266in}}%
\pgfpathlineto{\pgfqpoint{2.666309in}{1.756367in}}%
\pgfpathmoveto{\pgfqpoint{2.666309in}{1.756367in}}%
\pgfpathlineto{\pgfqpoint{2.666309in}{1.756367in}}%
\pgfpathlineto{\pgfqpoint{2.666309in}{1.762266in}}%
\pgfpathlineto{\pgfqpoint{2.675391in}{1.762266in}}%
\pgfpathlineto{\pgfqpoint{2.675391in}{1.756367in}}%
\pgfpathmoveto{\pgfqpoint{2.675391in}{1.762266in}}%
\pgfpathlineto{\pgfqpoint{2.675391in}{1.762266in}}%
\pgfpathlineto{\pgfqpoint{2.675391in}{1.768164in}}%
\pgfpathlineto{\pgfqpoint{2.684473in}{1.768164in}}%
\pgfpathlineto{\pgfqpoint{2.684473in}{1.762266in}}%
\pgfpathmoveto{\pgfqpoint{2.675391in}{1.768164in}}%
\pgfpathlineto{\pgfqpoint{2.675391in}{1.768164in}}%
\pgfpathlineto{\pgfqpoint{2.675391in}{1.774062in}}%
\pgfpathlineto{\pgfqpoint{2.684473in}{1.774062in}}%
\pgfpathlineto{\pgfqpoint{2.684473in}{1.768164in}}%
\pgfpathmoveto{\pgfqpoint{2.684473in}{1.768164in}}%
\pgfpathlineto{\pgfqpoint{2.684473in}{1.768164in}}%
\pgfpathlineto{\pgfqpoint{2.684473in}{1.774062in}}%
\pgfpathlineto{\pgfqpoint{2.693555in}{1.774062in}}%
\pgfpathlineto{\pgfqpoint{2.693555in}{1.768164in}}%
\pgfpathmoveto{\pgfqpoint{2.693555in}{1.774062in}}%
\pgfpathlineto{\pgfqpoint{2.693555in}{1.774062in}}%
\pgfpathlineto{\pgfqpoint{2.693555in}{1.779961in}}%
\pgfpathlineto{\pgfqpoint{2.702637in}{1.779961in}}%
\pgfpathlineto{\pgfqpoint{2.702637in}{1.774062in}}%
\pgfpathmoveto{\pgfqpoint{2.693555in}{1.779961in}}%
\pgfpathlineto{\pgfqpoint{2.693555in}{1.779961in}}%
\pgfpathlineto{\pgfqpoint{2.693555in}{1.785859in}}%
\pgfpathlineto{\pgfqpoint{2.702637in}{1.785859in}}%
\pgfpathlineto{\pgfqpoint{2.702637in}{1.779961in}}%
\pgfpathmoveto{\pgfqpoint{2.702637in}{1.779961in}}%
\pgfpathlineto{\pgfqpoint{2.702637in}{1.779961in}}%
\pgfpathlineto{\pgfqpoint{2.702637in}{1.785859in}}%
\pgfpathlineto{\pgfqpoint{2.711719in}{1.785859in}}%
\pgfpathlineto{\pgfqpoint{2.711719in}{1.779961in}}%
\pgfpathmoveto{\pgfqpoint{2.711719in}{1.785859in}}%
\pgfpathlineto{\pgfqpoint{2.711719in}{1.785859in}}%
\pgfpathlineto{\pgfqpoint{2.711719in}{1.791757in}}%
\pgfpathlineto{\pgfqpoint{2.720801in}{1.791757in}}%
\pgfpathlineto{\pgfqpoint{2.720801in}{1.785859in}}%
\pgfpathmoveto{\pgfqpoint{2.711719in}{1.791757in}}%
\pgfpathlineto{\pgfqpoint{2.711719in}{1.791757in}}%
\pgfpathlineto{\pgfqpoint{2.711719in}{1.797655in}}%
\pgfpathlineto{\pgfqpoint{2.720801in}{1.797655in}}%
\pgfpathlineto{\pgfqpoint{2.720801in}{1.791757in}}%
\pgfpathmoveto{\pgfqpoint{2.720801in}{1.791757in}}%
\pgfpathlineto{\pgfqpoint{2.720801in}{1.791757in}}%
\pgfpathlineto{\pgfqpoint{2.720801in}{1.797655in}}%
\pgfpathlineto{\pgfqpoint{2.729883in}{1.797655in}}%
\pgfpathlineto{\pgfqpoint{2.729883in}{1.791757in}}%
\pgfpathmoveto{\pgfqpoint{2.729883in}{1.797655in}}%
\pgfpathlineto{\pgfqpoint{2.729883in}{1.797655in}}%
\pgfpathlineto{\pgfqpoint{2.729883in}{1.803554in}}%
\pgfpathlineto{\pgfqpoint{2.738965in}{1.803554in}}%
\pgfpathlineto{\pgfqpoint{2.738965in}{1.797655in}}%
\pgfpathmoveto{\pgfqpoint{2.729883in}{1.803554in}}%
\pgfpathlineto{\pgfqpoint{2.729883in}{1.803554in}}%
\pgfpathlineto{\pgfqpoint{2.729883in}{1.809452in}}%
\pgfpathlineto{\pgfqpoint{2.738965in}{1.809452in}}%
\pgfpathlineto{\pgfqpoint{2.738965in}{1.803554in}}%
\pgfpathmoveto{\pgfqpoint{2.738965in}{1.803554in}}%
\pgfpathlineto{\pgfqpoint{2.738965in}{1.803554in}}%
\pgfpathlineto{\pgfqpoint{2.738965in}{1.809452in}}%
\pgfpathlineto{\pgfqpoint{2.748047in}{1.809452in}}%
\pgfpathlineto{\pgfqpoint{2.748047in}{1.803554in}}%
\pgfpathmoveto{\pgfqpoint{2.748047in}{1.809452in}}%
\pgfpathlineto{\pgfqpoint{2.748047in}{1.809452in}}%
\pgfpathlineto{\pgfqpoint{2.748047in}{1.815350in}}%
\pgfpathlineto{\pgfqpoint{2.757129in}{1.815350in}}%
\pgfpathlineto{\pgfqpoint{2.757129in}{1.809452in}}%
\pgfpathmoveto{\pgfqpoint{2.748047in}{1.815350in}}%
\pgfpathlineto{\pgfqpoint{2.748047in}{1.815350in}}%
\pgfpathlineto{\pgfqpoint{2.748047in}{1.821249in}}%
\pgfpathlineto{\pgfqpoint{2.757129in}{1.821249in}}%
\pgfpathlineto{\pgfqpoint{2.757129in}{1.815350in}}%
\pgfpathmoveto{\pgfqpoint{2.757129in}{1.815350in}}%
\pgfpathlineto{\pgfqpoint{2.757129in}{1.815350in}}%
\pgfpathlineto{\pgfqpoint{2.757129in}{1.821249in}}%
\pgfpathlineto{\pgfqpoint{2.766211in}{1.821249in}}%
\pgfpathlineto{\pgfqpoint{2.766211in}{1.815350in}}%
\pgfpathmoveto{\pgfqpoint{2.766211in}{1.821249in}}%
\pgfpathlineto{\pgfqpoint{2.766211in}{1.821249in}}%
\pgfpathlineto{\pgfqpoint{2.766211in}{1.827147in}}%
\pgfpathlineto{\pgfqpoint{2.775293in}{1.827147in}}%
\pgfpathlineto{\pgfqpoint{2.775293in}{1.821249in}}%
\pgfpathmoveto{\pgfqpoint{2.766211in}{1.827147in}}%
\pgfpathlineto{\pgfqpoint{2.766211in}{1.827147in}}%
\pgfpathlineto{\pgfqpoint{2.766211in}{1.833046in}}%
\pgfpathlineto{\pgfqpoint{2.775293in}{1.833046in}}%
\pgfpathlineto{\pgfqpoint{2.775293in}{1.827147in}}%
\pgfpathmoveto{\pgfqpoint{2.775293in}{1.827147in}}%
\pgfpathlineto{\pgfqpoint{2.775293in}{1.827147in}}%
\pgfpathlineto{\pgfqpoint{2.775293in}{1.833046in}}%
\pgfpathlineto{\pgfqpoint{2.784375in}{1.833046in}}%
\pgfpathlineto{\pgfqpoint{2.784375in}{1.827147in}}%
\pgfpathmoveto{\pgfqpoint{2.784375in}{1.833046in}}%
\pgfpathlineto{\pgfqpoint{2.784375in}{1.833046in}}%
\pgfpathlineto{\pgfqpoint{2.784375in}{1.838944in}}%
\pgfpathlineto{\pgfqpoint{2.793457in}{1.838944in}}%
\pgfpathlineto{\pgfqpoint{2.793457in}{1.833046in}}%
\pgfpathmoveto{\pgfqpoint{2.784375in}{1.838944in}}%
\pgfpathlineto{\pgfqpoint{2.784375in}{1.838944in}}%
\pgfpathlineto{\pgfqpoint{2.784375in}{1.844843in}}%
\pgfpathlineto{\pgfqpoint{2.793457in}{1.844843in}}%
\pgfpathlineto{\pgfqpoint{2.793457in}{1.838944in}}%
\pgfpathmoveto{\pgfqpoint{2.793457in}{1.838944in}}%
\pgfpathlineto{\pgfqpoint{2.793457in}{1.838944in}}%
\pgfpathlineto{\pgfqpoint{2.793457in}{1.844843in}}%
\pgfpathlineto{\pgfqpoint{2.802539in}{1.844843in}}%
\pgfpathlineto{\pgfqpoint{2.802539in}{1.838944in}}%
\pgfpathmoveto{\pgfqpoint{2.802539in}{1.844843in}}%
\pgfpathlineto{\pgfqpoint{2.802539in}{1.844843in}}%
\pgfpathlineto{\pgfqpoint{2.802539in}{1.850742in}}%
\pgfpathlineto{\pgfqpoint{2.811621in}{1.850742in}}%
\pgfpathlineto{\pgfqpoint{2.811621in}{1.844843in}}%
\pgfpathmoveto{\pgfqpoint{2.802539in}{1.850742in}}%
\pgfpathlineto{\pgfqpoint{2.802539in}{1.850742in}}%
\pgfpathlineto{\pgfqpoint{2.802539in}{1.856640in}}%
\pgfpathlineto{\pgfqpoint{2.811621in}{1.856640in}}%
\pgfpathlineto{\pgfqpoint{2.811621in}{1.850742in}}%
\pgfpathmoveto{\pgfqpoint{2.811621in}{1.850742in}}%
\pgfpathlineto{\pgfqpoint{2.811621in}{1.850742in}}%
\pgfpathlineto{\pgfqpoint{2.811621in}{1.856640in}}%
\pgfpathlineto{\pgfqpoint{2.820703in}{1.856640in}}%
\pgfpathlineto{\pgfqpoint{2.820703in}{1.850742in}}%
\pgfpathmoveto{\pgfqpoint{2.820703in}{1.856640in}}%
\pgfpathlineto{\pgfqpoint{2.820703in}{1.856640in}}%
\pgfpathlineto{\pgfqpoint{2.820703in}{1.862539in}}%
\pgfpathlineto{\pgfqpoint{2.829785in}{1.862539in}}%
\pgfpathlineto{\pgfqpoint{2.829785in}{1.856640in}}%
\pgfpathmoveto{\pgfqpoint{2.820703in}{1.862539in}}%
\pgfpathlineto{\pgfqpoint{2.820703in}{1.862539in}}%
\pgfpathlineto{\pgfqpoint{2.820703in}{1.868438in}}%
\pgfpathlineto{\pgfqpoint{2.829785in}{1.868438in}}%
\pgfpathlineto{\pgfqpoint{2.829785in}{1.862539in}}%
\pgfpathmoveto{\pgfqpoint{2.829785in}{1.862539in}}%
\pgfpathlineto{\pgfqpoint{2.829785in}{1.862539in}}%
\pgfpathlineto{\pgfqpoint{2.829785in}{1.868438in}}%
\pgfpathlineto{\pgfqpoint{2.838867in}{1.868438in}}%
\pgfpathlineto{\pgfqpoint{2.838867in}{1.862539in}}%
\pgfpathmoveto{\pgfqpoint{2.857031in}{1.874336in}}%
\pgfpathlineto{\pgfqpoint{2.857031in}{1.874336in}}%
\pgfpathlineto{\pgfqpoint{2.857031in}{1.880235in}}%
\pgfpathlineto{\pgfqpoint{2.866113in}{1.880235in}}%
\pgfpathlineto{\pgfqpoint{2.866113in}{1.874336in}}%
\pgfpathmoveto{\pgfqpoint{2.875196in}{1.886134in}}%
\pgfpathlineto{\pgfqpoint{2.875196in}{1.886134in}}%
\pgfpathlineto{\pgfqpoint{2.875196in}{1.892032in}}%
\pgfpathlineto{\pgfqpoint{2.884278in}{1.892032in}}%
\pgfpathlineto{\pgfqpoint{2.884278in}{1.886134in}}%
\pgfpathmoveto{\pgfqpoint{2.893360in}{1.897931in}}%
\pgfpathlineto{\pgfqpoint{2.893360in}{1.897931in}}%
\pgfpathlineto{\pgfqpoint{2.893360in}{1.903829in}}%
\pgfpathlineto{\pgfqpoint{2.902442in}{1.903829in}}%
\pgfpathlineto{\pgfqpoint{2.902442in}{1.897931in}}%
\pgfpathmoveto{\pgfqpoint{2.911524in}{1.909728in}}%
\pgfpathlineto{\pgfqpoint{2.911524in}{1.909728in}}%
\pgfpathlineto{\pgfqpoint{2.911524in}{1.915627in}}%
\pgfpathlineto{\pgfqpoint{2.920606in}{1.915627in}}%
\pgfpathlineto{\pgfqpoint{2.920606in}{1.909728in}}%
\pgfpathmoveto{\pgfqpoint{2.911524in}{3.402033in}}%
\pgfpathlineto{\pgfqpoint{2.911524in}{3.402033in}}%
\pgfpathlineto{\pgfqpoint{2.911524in}{3.407931in}}%
\pgfpathlineto{\pgfqpoint{2.920606in}{3.407931in}}%
\pgfpathlineto{\pgfqpoint{2.920606in}{3.402033in}}%
\pgfpathmoveto{\pgfqpoint{2.911524in}{3.407931in}}%
\pgfpathlineto{\pgfqpoint{2.911524in}{3.407931in}}%
\pgfpathlineto{\pgfqpoint{2.911524in}{3.413830in}}%
\pgfpathlineto{\pgfqpoint{2.920606in}{3.413830in}}%
\pgfpathlineto{\pgfqpoint{2.920606in}{3.407931in}}%
\pgfpathmoveto{\pgfqpoint{2.838867in}{3.496405in}}%
\pgfpathlineto{\pgfqpoint{2.838867in}{3.496405in}}%
\pgfpathlineto{\pgfqpoint{2.838867in}{3.502304in}}%
\pgfpathlineto{\pgfqpoint{2.847949in}{3.502304in}}%
\pgfpathlineto{\pgfqpoint{2.847949in}{3.496405in}}%
\pgfpathmoveto{\pgfqpoint{2.838867in}{3.502304in}}%
\pgfpathlineto{\pgfqpoint{2.838867in}{3.502304in}}%
\pgfpathlineto{\pgfqpoint{2.838867in}{3.508202in}}%
\pgfpathlineto{\pgfqpoint{2.847949in}{3.508202in}}%
\pgfpathlineto{\pgfqpoint{2.847949in}{3.502304in}}%
\pgfpathmoveto{\pgfqpoint{2.875196in}{3.449220in}}%
\pgfpathlineto{\pgfqpoint{2.875196in}{3.449220in}}%
\pgfpathlineto{\pgfqpoint{2.875196in}{3.455118in}}%
\pgfpathlineto{\pgfqpoint{2.884278in}{3.455118in}}%
\pgfpathlineto{\pgfqpoint{2.884278in}{3.449220in}}%
\pgfpathmoveto{\pgfqpoint{2.875196in}{3.455118in}}%
\pgfpathlineto{\pgfqpoint{2.875196in}{3.455118in}}%
\pgfpathlineto{\pgfqpoint{2.875196in}{3.461016in}}%
\pgfpathlineto{\pgfqpoint{2.884278in}{3.461016in}}%
\pgfpathlineto{\pgfqpoint{2.884278in}{3.455118in}}%
\pgfpathmoveto{\pgfqpoint{2.893360in}{3.425627in}}%
\pgfpathlineto{\pgfqpoint{2.893360in}{3.425627in}}%
\pgfpathlineto{\pgfqpoint{2.893360in}{3.431525in}}%
\pgfpathlineto{\pgfqpoint{2.902442in}{3.431525in}}%
\pgfpathlineto{\pgfqpoint{2.902442in}{3.425627in}}%
\pgfpathmoveto{\pgfqpoint{2.893360in}{3.431525in}}%
\pgfpathlineto{\pgfqpoint{2.893360in}{3.431525in}}%
\pgfpathlineto{\pgfqpoint{2.893360in}{3.437423in}}%
\pgfpathlineto{\pgfqpoint{2.902442in}{3.437423in}}%
\pgfpathlineto{\pgfqpoint{2.902442in}{3.431525in}}%
\pgfpathmoveto{\pgfqpoint{2.857031in}{3.472812in}}%
\pgfpathlineto{\pgfqpoint{2.857031in}{3.472812in}}%
\pgfpathlineto{\pgfqpoint{2.857031in}{3.478711in}}%
\pgfpathlineto{\pgfqpoint{2.866113in}{3.478711in}}%
\pgfpathlineto{\pgfqpoint{2.866113in}{3.472812in}}%
\pgfpathmoveto{\pgfqpoint{2.857031in}{3.478711in}}%
\pgfpathlineto{\pgfqpoint{2.857031in}{3.478711in}}%
\pgfpathlineto{\pgfqpoint{2.857031in}{3.484609in}}%
\pgfpathlineto{\pgfqpoint{2.866113in}{3.484609in}}%
\pgfpathlineto{\pgfqpoint{2.866113in}{3.478711in}}%
\pgfpathmoveto{\pgfqpoint{2.929688in}{1.921525in}}%
\pgfpathlineto{\pgfqpoint{2.929688in}{1.921525in}}%
\pgfpathlineto{\pgfqpoint{2.929688in}{1.927423in}}%
\pgfpathlineto{\pgfqpoint{2.938770in}{1.927423in}}%
\pgfpathlineto{\pgfqpoint{2.938770in}{1.921525in}}%
\pgfpathmoveto{\pgfqpoint{2.947852in}{1.933322in}}%
\pgfpathlineto{\pgfqpoint{2.947852in}{1.933322in}}%
\pgfpathlineto{\pgfqpoint{2.947852in}{1.939220in}}%
\pgfpathlineto{\pgfqpoint{2.956935in}{1.939220in}}%
\pgfpathlineto{\pgfqpoint{2.956935in}{1.933322in}}%
\pgfpathmoveto{\pgfqpoint{2.966017in}{1.945118in}}%
\pgfpathlineto{\pgfqpoint{2.966017in}{1.945118in}}%
\pgfpathlineto{\pgfqpoint{2.966017in}{1.951017in}}%
\pgfpathlineto{\pgfqpoint{2.975099in}{1.951017in}}%
\pgfpathlineto{\pgfqpoint{2.975099in}{1.945118in}}%
\pgfpathmoveto{\pgfqpoint{2.984181in}{1.962813in}}%
\pgfpathlineto{\pgfqpoint{2.984181in}{1.962813in}}%
\pgfpathlineto{\pgfqpoint{2.984181in}{1.968712in}}%
\pgfpathlineto{\pgfqpoint{2.993264in}{1.968712in}}%
\pgfpathlineto{\pgfqpoint{2.993264in}{1.962813in}}%
\pgfpathmoveto{\pgfqpoint{2.984181in}{1.968712in}}%
\pgfpathlineto{\pgfqpoint{2.984181in}{1.968712in}}%
\pgfpathlineto{\pgfqpoint{2.984181in}{1.974610in}}%
\pgfpathlineto{\pgfqpoint{2.993264in}{1.974610in}}%
\pgfpathlineto{\pgfqpoint{2.993264in}{1.968712in}}%
\pgfpathmoveto{\pgfqpoint{2.993264in}{1.968712in}}%
\pgfpathlineto{\pgfqpoint{2.993264in}{1.968712in}}%
\pgfpathlineto{\pgfqpoint{2.993264in}{1.974610in}}%
\pgfpathlineto{\pgfqpoint{3.002346in}{1.974610in}}%
\pgfpathlineto{\pgfqpoint{3.002346in}{1.968712in}}%
\pgfpathmoveto{\pgfqpoint{3.002346in}{1.974610in}}%
\pgfpathlineto{\pgfqpoint{3.002346in}{1.974610in}}%
\pgfpathlineto{\pgfqpoint{3.002346in}{1.980508in}}%
\pgfpathlineto{\pgfqpoint{3.011428in}{1.980508in}}%
\pgfpathlineto{\pgfqpoint{3.011428in}{1.974610in}}%
\pgfpathmoveto{\pgfqpoint{3.002346in}{1.980508in}}%
\pgfpathlineto{\pgfqpoint{3.002346in}{1.980508in}}%
\pgfpathlineto{\pgfqpoint{3.002346in}{1.986407in}}%
\pgfpathlineto{\pgfqpoint{3.011428in}{1.986407in}}%
\pgfpathlineto{\pgfqpoint{3.011428in}{1.980508in}}%
\pgfpathmoveto{\pgfqpoint{3.011428in}{1.980508in}}%
\pgfpathlineto{\pgfqpoint{3.011428in}{1.980508in}}%
\pgfpathlineto{\pgfqpoint{3.011428in}{1.986407in}}%
\pgfpathlineto{\pgfqpoint{3.020510in}{1.986407in}}%
\pgfpathlineto{\pgfqpoint{3.020510in}{1.980508in}}%
\pgfpathmoveto{\pgfqpoint{3.020510in}{1.986407in}}%
\pgfpathlineto{\pgfqpoint{3.020510in}{1.986407in}}%
\pgfpathlineto{\pgfqpoint{3.020510in}{1.992305in}}%
\pgfpathlineto{\pgfqpoint{3.029593in}{1.992305in}}%
\pgfpathlineto{\pgfqpoint{3.029593in}{1.986407in}}%
\pgfpathmoveto{\pgfqpoint{3.020510in}{1.992305in}}%
\pgfpathlineto{\pgfqpoint{3.020510in}{1.992305in}}%
\pgfpathlineto{\pgfqpoint{3.020510in}{1.998204in}}%
\pgfpathlineto{\pgfqpoint{3.029593in}{1.998204in}}%
\pgfpathlineto{\pgfqpoint{3.029593in}{1.992305in}}%
\pgfpathmoveto{\pgfqpoint{3.029593in}{1.992305in}}%
\pgfpathlineto{\pgfqpoint{3.029593in}{1.992305in}}%
\pgfpathlineto{\pgfqpoint{3.029593in}{1.998204in}}%
\pgfpathlineto{\pgfqpoint{3.038675in}{1.998204in}}%
\pgfpathlineto{\pgfqpoint{3.038675in}{1.992305in}}%
\pgfpathmoveto{\pgfqpoint{3.038675in}{1.998204in}}%
\pgfpathlineto{\pgfqpoint{3.038675in}{1.998204in}}%
\pgfpathlineto{\pgfqpoint{3.038675in}{2.004102in}}%
\pgfpathlineto{\pgfqpoint{3.047757in}{2.004102in}}%
\pgfpathlineto{\pgfqpoint{3.047757in}{1.998204in}}%
\pgfpathmoveto{\pgfqpoint{3.038675in}{2.004102in}}%
\pgfpathlineto{\pgfqpoint{3.038675in}{2.004102in}}%
\pgfpathlineto{\pgfqpoint{3.038675in}{2.010000in}}%
\pgfpathlineto{\pgfqpoint{3.047757in}{2.010000in}}%
\pgfpathlineto{\pgfqpoint{3.047757in}{2.004102in}}%
\pgfpathmoveto{\pgfqpoint{3.047757in}{2.004102in}}%
\pgfpathlineto{\pgfqpoint{3.047757in}{2.004102in}}%
\pgfpathlineto{\pgfqpoint{3.047757in}{2.010000in}}%
\pgfpathlineto{\pgfqpoint{3.056839in}{2.010000in}}%
\pgfpathlineto{\pgfqpoint{3.056839in}{2.004102in}}%
\pgfpathmoveto{\pgfqpoint{3.056839in}{2.010000in}}%
\pgfpathlineto{\pgfqpoint{3.056839in}{2.010000in}}%
\pgfpathlineto{\pgfqpoint{3.056839in}{2.015899in}}%
\pgfpathlineto{\pgfqpoint{3.065922in}{2.015899in}}%
\pgfpathlineto{\pgfqpoint{3.065922in}{2.010000in}}%
\pgfpathmoveto{\pgfqpoint{3.056839in}{2.015899in}}%
\pgfpathlineto{\pgfqpoint{3.056839in}{2.015899in}}%
\pgfpathlineto{\pgfqpoint{3.056839in}{2.021797in}}%
\pgfpathlineto{\pgfqpoint{3.065922in}{2.021797in}}%
\pgfpathlineto{\pgfqpoint{3.065922in}{2.015899in}}%
\pgfpathmoveto{\pgfqpoint{3.065922in}{2.015899in}}%
\pgfpathlineto{\pgfqpoint{3.065922in}{2.015899in}}%
\pgfpathlineto{\pgfqpoint{3.065922in}{2.021797in}}%
\pgfpathlineto{\pgfqpoint{3.075004in}{2.021797in}}%
\pgfpathlineto{\pgfqpoint{3.075004in}{2.015899in}}%
\pgfpathmoveto{\pgfqpoint{3.056839in}{3.213281in}}%
\pgfpathlineto{\pgfqpoint{3.056839in}{3.213281in}}%
\pgfpathlineto{\pgfqpoint{3.056839in}{3.219180in}}%
\pgfpathlineto{\pgfqpoint{3.065922in}{3.219180in}}%
\pgfpathlineto{\pgfqpoint{3.065922in}{3.213281in}}%
\pgfpathmoveto{\pgfqpoint{3.056839in}{3.219180in}}%
\pgfpathlineto{\pgfqpoint{3.056839in}{3.219180in}}%
\pgfpathlineto{\pgfqpoint{3.056839in}{3.225078in}}%
\pgfpathlineto{\pgfqpoint{3.065922in}{3.225078in}}%
\pgfpathlineto{\pgfqpoint{3.065922in}{3.219180in}}%
\pgfpathmoveto{\pgfqpoint{2.984181in}{3.307657in}}%
\pgfpathlineto{\pgfqpoint{2.984181in}{3.307657in}}%
\pgfpathlineto{\pgfqpoint{2.984181in}{3.313556in}}%
\pgfpathlineto{\pgfqpoint{2.993264in}{3.313556in}}%
\pgfpathlineto{\pgfqpoint{2.993264in}{3.307657in}}%
\pgfpathmoveto{\pgfqpoint{2.984181in}{3.313556in}}%
\pgfpathlineto{\pgfqpoint{2.984181in}{3.313556in}}%
\pgfpathlineto{\pgfqpoint{2.984181in}{3.319455in}}%
\pgfpathlineto{\pgfqpoint{2.993264in}{3.319455in}}%
\pgfpathlineto{\pgfqpoint{2.993264in}{3.313556in}}%
\pgfpathmoveto{\pgfqpoint{3.020510in}{3.260469in}}%
\pgfpathlineto{\pgfqpoint{3.020510in}{3.260469in}}%
\pgfpathlineto{\pgfqpoint{3.020510in}{3.266367in}}%
\pgfpathlineto{\pgfqpoint{3.029593in}{3.266367in}}%
\pgfpathlineto{\pgfqpoint{3.029593in}{3.260469in}}%
\pgfpathmoveto{\pgfqpoint{3.020510in}{3.266367in}}%
\pgfpathlineto{\pgfqpoint{3.020510in}{3.266367in}}%
\pgfpathlineto{\pgfqpoint{3.020510in}{3.272266in}}%
\pgfpathlineto{\pgfqpoint{3.029593in}{3.272266in}}%
\pgfpathlineto{\pgfqpoint{3.029593in}{3.266367in}}%
\pgfpathmoveto{\pgfqpoint{3.038675in}{3.236874in}}%
\pgfpathlineto{\pgfqpoint{3.038675in}{3.236874in}}%
\pgfpathlineto{\pgfqpoint{3.038675in}{3.242773in}}%
\pgfpathlineto{\pgfqpoint{3.047757in}{3.242773in}}%
\pgfpathlineto{\pgfqpoint{3.047757in}{3.236874in}}%
\pgfpathmoveto{\pgfqpoint{3.038675in}{3.242773in}}%
\pgfpathlineto{\pgfqpoint{3.038675in}{3.242773in}}%
\pgfpathlineto{\pgfqpoint{3.038675in}{3.248672in}}%
\pgfpathlineto{\pgfqpoint{3.047757in}{3.248672in}}%
\pgfpathlineto{\pgfqpoint{3.047757in}{3.242773in}}%
\pgfpathmoveto{\pgfqpoint{3.002346in}{3.284063in}}%
\pgfpathlineto{\pgfqpoint{3.002346in}{3.284063in}}%
\pgfpathlineto{\pgfqpoint{3.002346in}{3.289962in}}%
\pgfpathlineto{\pgfqpoint{3.011428in}{3.289962in}}%
\pgfpathlineto{\pgfqpoint{3.011428in}{3.284063in}}%
\pgfpathmoveto{\pgfqpoint{3.002346in}{3.289962in}}%
\pgfpathlineto{\pgfqpoint{3.002346in}{3.289962in}}%
\pgfpathlineto{\pgfqpoint{3.002346in}{3.295860in}}%
\pgfpathlineto{\pgfqpoint{3.011428in}{3.295860in}}%
\pgfpathlineto{\pgfqpoint{3.011428in}{3.289962in}}%
\pgfpathmoveto{\pgfqpoint{2.947852in}{3.354845in}}%
\pgfpathlineto{\pgfqpoint{2.947852in}{3.354845in}}%
\pgfpathlineto{\pgfqpoint{2.947852in}{3.360744in}}%
\pgfpathlineto{\pgfqpoint{2.956935in}{3.360744in}}%
\pgfpathlineto{\pgfqpoint{2.956935in}{3.354845in}}%
\pgfpathmoveto{\pgfqpoint{2.947852in}{3.360744in}}%
\pgfpathlineto{\pgfqpoint{2.947852in}{3.360744in}}%
\pgfpathlineto{\pgfqpoint{2.947852in}{3.366642in}}%
\pgfpathlineto{\pgfqpoint{2.956935in}{3.366642in}}%
\pgfpathlineto{\pgfqpoint{2.956935in}{3.360744in}}%
\pgfpathmoveto{\pgfqpoint{2.966017in}{3.331252in}}%
\pgfpathlineto{\pgfqpoint{2.966017in}{3.331252in}}%
\pgfpathlineto{\pgfqpoint{2.966017in}{3.337150in}}%
\pgfpathlineto{\pgfqpoint{2.975099in}{3.337150in}}%
\pgfpathlineto{\pgfqpoint{2.975099in}{3.331252in}}%
\pgfpathmoveto{\pgfqpoint{2.966017in}{3.337150in}}%
\pgfpathlineto{\pgfqpoint{2.966017in}{3.337150in}}%
\pgfpathlineto{\pgfqpoint{2.966017in}{3.343049in}}%
\pgfpathlineto{\pgfqpoint{2.975099in}{3.343049in}}%
\pgfpathlineto{\pgfqpoint{2.975099in}{3.337150in}}%
\pgfpathmoveto{\pgfqpoint{2.929688in}{3.378439in}}%
\pgfpathlineto{\pgfqpoint{2.929688in}{3.378439in}}%
\pgfpathlineto{\pgfqpoint{2.929688in}{3.384338in}}%
\pgfpathlineto{\pgfqpoint{2.938770in}{3.384338in}}%
\pgfpathlineto{\pgfqpoint{2.938770in}{3.378439in}}%
\pgfpathmoveto{\pgfqpoint{2.929688in}{3.384338in}}%
\pgfpathlineto{\pgfqpoint{2.929688in}{3.384338in}}%
\pgfpathlineto{\pgfqpoint{2.929688in}{3.390236in}}%
\pgfpathlineto{\pgfqpoint{2.938770in}{3.390236in}}%
\pgfpathlineto{\pgfqpoint{2.938770in}{3.384338in}}%
\pgfpathmoveto{\pgfqpoint{3.075004in}{2.021797in}}%
\pgfpathlineto{\pgfqpoint{3.075004in}{2.021797in}}%
\pgfpathlineto{\pgfqpoint{3.075004in}{2.027696in}}%
\pgfpathlineto{\pgfqpoint{3.084085in}{2.027696in}}%
\pgfpathlineto{\pgfqpoint{3.084085in}{2.021797in}}%
\pgfpathmoveto{\pgfqpoint{3.075004in}{2.027696in}}%
\pgfpathlineto{\pgfqpoint{3.075004in}{2.027696in}}%
\pgfpathlineto{\pgfqpoint{3.075004in}{2.033594in}}%
\pgfpathlineto{\pgfqpoint{3.084085in}{2.033594in}}%
\pgfpathlineto{\pgfqpoint{3.084085in}{2.027696in}}%
\pgfpathmoveto{\pgfqpoint{3.084085in}{2.027696in}}%
\pgfpathlineto{\pgfqpoint{3.084085in}{2.027696in}}%
\pgfpathlineto{\pgfqpoint{3.084085in}{2.033594in}}%
\pgfpathlineto{\pgfqpoint{3.093167in}{2.033594in}}%
\pgfpathlineto{\pgfqpoint{3.093167in}{2.027696in}}%
\pgfpathmoveto{\pgfqpoint{3.093167in}{2.033594in}}%
\pgfpathlineto{\pgfqpoint{3.093167in}{2.033594in}}%
\pgfpathlineto{\pgfqpoint{3.093167in}{2.039493in}}%
\pgfpathlineto{\pgfqpoint{3.102249in}{2.039493in}}%
\pgfpathlineto{\pgfqpoint{3.102249in}{2.033594in}}%
\pgfpathmoveto{\pgfqpoint{3.093167in}{2.039493in}}%
\pgfpathlineto{\pgfqpoint{3.093167in}{2.039493in}}%
\pgfpathlineto{\pgfqpoint{3.093167in}{2.045392in}}%
\pgfpathlineto{\pgfqpoint{3.102249in}{2.045392in}}%
\pgfpathlineto{\pgfqpoint{3.102249in}{2.039493in}}%
\pgfpathmoveto{\pgfqpoint{3.102249in}{2.039493in}}%
\pgfpathlineto{\pgfqpoint{3.102249in}{2.039493in}}%
\pgfpathlineto{\pgfqpoint{3.102249in}{2.045392in}}%
\pgfpathlineto{\pgfqpoint{3.111330in}{2.045392in}}%
\pgfpathlineto{\pgfqpoint{3.111330in}{2.039493in}}%
\pgfpathmoveto{\pgfqpoint{3.129493in}{2.051290in}}%
\pgfpathlineto{\pgfqpoint{3.129493in}{2.051290in}}%
\pgfpathlineto{\pgfqpoint{3.129493in}{2.057189in}}%
\pgfpathlineto{\pgfqpoint{3.138575in}{2.057189in}}%
\pgfpathlineto{\pgfqpoint{3.138575in}{2.051290in}}%
\pgfpathmoveto{\pgfqpoint{3.147657in}{2.063087in}}%
\pgfpathlineto{\pgfqpoint{3.147657in}{2.063087in}}%
\pgfpathlineto{\pgfqpoint{3.147657in}{2.068986in}}%
\pgfpathlineto{\pgfqpoint{3.156738in}{2.068986in}}%
\pgfpathlineto{\pgfqpoint{3.156738in}{2.063087in}}%
\pgfpathmoveto{\pgfqpoint{3.165820in}{2.074884in}}%
\pgfpathlineto{\pgfqpoint{3.165820in}{2.074884in}}%
\pgfpathlineto{\pgfqpoint{3.165820in}{2.080783in}}%
\pgfpathlineto{\pgfqpoint{3.174901in}{2.080783in}}%
\pgfpathlineto{\pgfqpoint{3.174901in}{2.074884in}}%
\pgfpathmoveto{\pgfqpoint{3.183983in}{2.086682in}}%
\pgfpathlineto{\pgfqpoint{3.183983in}{2.086682in}}%
\pgfpathlineto{\pgfqpoint{3.183983in}{2.092580in}}%
\pgfpathlineto{\pgfqpoint{3.193065in}{2.092580in}}%
\pgfpathlineto{\pgfqpoint{3.193065in}{2.086682in}}%
\pgfpathmoveto{\pgfqpoint{3.202146in}{2.098479in}}%
\pgfpathlineto{\pgfqpoint{3.202146in}{2.098479in}}%
\pgfpathlineto{\pgfqpoint{3.202146in}{2.104377in}}%
\pgfpathlineto{\pgfqpoint{3.211228in}{2.104377in}}%
\pgfpathlineto{\pgfqpoint{3.211228in}{2.098479in}}%
\pgfpathmoveto{\pgfqpoint{3.202146in}{3.024531in}}%
\pgfpathlineto{\pgfqpoint{3.202146in}{3.024531in}}%
\pgfpathlineto{\pgfqpoint{3.202146in}{3.030430in}}%
\pgfpathlineto{\pgfqpoint{3.211228in}{3.030430in}}%
\pgfpathlineto{\pgfqpoint{3.211228in}{3.024531in}}%
\pgfpathmoveto{\pgfqpoint{3.202146in}{3.030430in}}%
\pgfpathlineto{\pgfqpoint{3.202146in}{3.030430in}}%
\pgfpathlineto{\pgfqpoint{3.202146in}{3.036328in}}%
\pgfpathlineto{\pgfqpoint{3.211228in}{3.036328in}}%
\pgfpathlineto{\pgfqpoint{3.211228in}{3.030430in}}%
\pgfpathmoveto{\pgfqpoint{3.129493in}{3.118907in}}%
\pgfpathlineto{\pgfqpoint{3.129493in}{3.118907in}}%
\pgfpathlineto{\pgfqpoint{3.129493in}{3.124806in}}%
\pgfpathlineto{\pgfqpoint{3.138575in}{3.124806in}}%
\pgfpathlineto{\pgfqpoint{3.138575in}{3.118907in}}%
\pgfpathmoveto{\pgfqpoint{3.129493in}{3.124806in}}%
\pgfpathlineto{\pgfqpoint{3.129493in}{3.124806in}}%
\pgfpathlineto{\pgfqpoint{3.129493in}{3.130704in}}%
\pgfpathlineto{\pgfqpoint{3.138575in}{3.130704in}}%
\pgfpathlineto{\pgfqpoint{3.138575in}{3.124806in}}%
\pgfpathmoveto{\pgfqpoint{3.165820in}{3.071719in}}%
\pgfpathlineto{\pgfqpoint{3.165820in}{3.071719in}}%
\pgfpathlineto{\pgfqpoint{3.165820in}{3.077617in}}%
\pgfpathlineto{\pgfqpoint{3.174901in}{3.077617in}}%
\pgfpathlineto{\pgfqpoint{3.174901in}{3.071719in}}%
\pgfpathmoveto{\pgfqpoint{3.165820in}{3.077617in}}%
\pgfpathlineto{\pgfqpoint{3.165820in}{3.077617in}}%
\pgfpathlineto{\pgfqpoint{3.165820in}{3.083516in}}%
\pgfpathlineto{\pgfqpoint{3.174901in}{3.083516in}}%
\pgfpathlineto{\pgfqpoint{3.174901in}{3.077617in}}%
\pgfpathmoveto{\pgfqpoint{3.183983in}{3.048124in}}%
\pgfpathlineto{\pgfqpoint{3.183983in}{3.048124in}}%
\pgfpathlineto{\pgfqpoint{3.183983in}{3.054023in}}%
\pgfpathlineto{\pgfqpoint{3.193065in}{3.054023in}}%
\pgfpathlineto{\pgfqpoint{3.193065in}{3.048124in}}%
\pgfpathmoveto{\pgfqpoint{3.183983in}{3.054023in}}%
\pgfpathlineto{\pgfqpoint{3.183983in}{3.054023in}}%
\pgfpathlineto{\pgfqpoint{3.183983in}{3.059922in}}%
\pgfpathlineto{\pgfqpoint{3.193065in}{3.059922in}}%
\pgfpathlineto{\pgfqpoint{3.193065in}{3.054023in}}%
\pgfpathmoveto{\pgfqpoint{3.147657in}{3.095313in}}%
\pgfpathlineto{\pgfqpoint{3.147657in}{3.095313in}}%
\pgfpathlineto{\pgfqpoint{3.147657in}{3.101212in}}%
\pgfpathlineto{\pgfqpoint{3.156738in}{3.101212in}}%
\pgfpathlineto{\pgfqpoint{3.156738in}{3.095313in}}%
\pgfpathmoveto{\pgfqpoint{3.147657in}{3.101212in}}%
\pgfpathlineto{\pgfqpoint{3.147657in}{3.101212in}}%
\pgfpathlineto{\pgfqpoint{3.147657in}{3.107110in}}%
\pgfpathlineto{\pgfqpoint{3.156738in}{3.107110in}}%
\pgfpathlineto{\pgfqpoint{3.156738in}{3.101212in}}%
\pgfpathmoveto{\pgfqpoint{3.093167in}{3.166095in}}%
\pgfpathlineto{\pgfqpoint{3.093167in}{3.166095in}}%
\pgfpathlineto{\pgfqpoint{3.093167in}{3.171993in}}%
\pgfpathlineto{\pgfqpoint{3.102249in}{3.171993in}}%
\pgfpathlineto{\pgfqpoint{3.102249in}{3.166095in}}%
\pgfpathmoveto{\pgfqpoint{3.093167in}{3.171993in}}%
\pgfpathlineto{\pgfqpoint{3.093167in}{3.171993in}}%
\pgfpathlineto{\pgfqpoint{3.093167in}{3.177891in}}%
\pgfpathlineto{\pgfqpoint{3.102249in}{3.177891in}}%
\pgfpathlineto{\pgfqpoint{3.102249in}{3.171993in}}%
\pgfpathmoveto{\pgfqpoint{3.111330in}{3.142502in}}%
\pgfpathlineto{\pgfqpoint{3.111330in}{3.142502in}}%
\pgfpathlineto{\pgfqpoint{3.111330in}{3.148400in}}%
\pgfpathlineto{\pgfqpoint{3.120412in}{3.148400in}}%
\pgfpathlineto{\pgfqpoint{3.120412in}{3.142502in}}%
\pgfpathmoveto{\pgfqpoint{3.111330in}{3.148400in}}%
\pgfpathlineto{\pgfqpoint{3.111330in}{3.148400in}}%
\pgfpathlineto{\pgfqpoint{3.111330in}{3.154298in}}%
\pgfpathlineto{\pgfqpoint{3.120412in}{3.154298in}}%
\pgfpathlineto{\pgfqpoint{3.120412in}{3.148400in}}%
\pgfpathmoveto{\pgfqpoint{3.075004in}{3.189688in}}%
\pgfpathlineto{\pgfqpoint{3.075004in}{3.189688in}}%
\pgfpathlineto{\pgfqpoint{3.075004in}{3.195586in}}%
\pgfpathlineto{\pgfqpoint{3.084085in}{3.195586in}}%
\pgfpathlineto{\pgfqpoint{3.084085in}{3.189688in}}%
\pgfpathmoveto{\pgfqpoint{3.075004in}{3.195586in}}%
\pgfpathlineto{\pgfqpoint{3.075004in}{3.195586in}}%
\pgfpathlineto{\pgfqpoint{3.075004in}{3.201485in}}%
\pgfpathlineto{\pgfqpoint{3.084085in}{3.201485in}}%
\pgfpathlineto{\pgfqpoint{3.084085in}{3.195586in}}%
\pgfpathmoveto{\pgfqpoint{3.220310in}{2.110275in}}%
\pgfpathlineto{\pgfqpoint{3.220310in}{2.110275in}}%
\pgfpathlineto{\pgfqpoint{3.220310in}{2.116174in}}%
\pgfpathlineto{\pgfqpoint{3.229392in}{2.116174in}}%
\pgfpathlineto{\pgfqpoint{3.229392in}{2.110275in}}%
\pgfpathmoveto{\pgfqpoint{3.238474in}{2.122072in}}%
\pgfpathlineto{\pgfqpoint{3.238474in}{2.122072in}}%
\pgfpathlineto{\pgfqpoint{3.238474in}{2.127970in}}%
\pgfpathlineto{\pgfqpoint{3.247557in}{2.127970in}}%
\pgfpathlineto{\pgfqpoint{3.247557in}{2.122072in}}%
\pgfpathmoveto{\pgfqpoint{3.256639in}{2.133868in}}%
\pgfpathlineto{\pgfqpoint{3.256639in}{2.133868in}}%
\pgfpathlineto{\pgfqpoint{3.256639in}{2.139766in}}%
\pgfpathlineto{\pgfqpoint{3.265721in}{2.139766in}}%
\pgfpathlineto{\pgfqpoint{3.265721in}{2.133868in}}%
\pgfpathmoveto{\pgfqpoint{3.274803in}{2.145665in}}%
\pgfpathlineto{\pgfqpoint{3.274803in}{2.145665in}}%
\pgfpathlineto{\pgfqpoint{3.274803in}{2.151563in}}%
\pgfpathlineto{\pgfqpoint{3.283886in}{2.151563in}}%
\pgfpathlineto{\pgfqpoint{3.283886in}{2.145665in}}%
\pgfpathmoveto{\pgfqpoint{3.292968in}{2.157461in}}%
\pgfpathlineto{\pgfqpoint{3.292968in}{2.157461in}}%
\pgfpathlineto{\pgfqpoint{3.292968in}{2.163359in}}%
\pgfpathlineto{\pgfqpoint{3.302050in}{2.163359in}}%
\pgfpathlineto{\pgfqpoint{3.302050in}{2.157461in}}%
\pgfpathmoveto{\pgfqpoint{3.292968in}{2.163359in}}%
\pgfpathlineto{\pgfqpoint{3.292968in}{2.163359in}}%
\pgfpathlineto{\pgfqpoint{3.292968in}{2.169257in}}%
\pgfpathlineto{\pgfqpoint{3.302050in}{2.169257in}}%
\pgfpathlineto{\pgfqpoint{3.302050in}{2.163359in}}%
\pgfpathmoveto{\pgfqpoint{3.292968in}{2.169257in}}%
\pgfpathlineto{\pgfqpoint{3.292968in}{2.169257in}}%
\pgfpathlineto{\pgfqpoint{3.292968in}{2.175155in}}%
\pgfpathlineto{\pgfqpoint{3.302050in}{2.175155in}}%
\pgfpathlineto{\pgfqpoint{3.302050in}{2.169257in}}%
\pgfpathmoveto{\pgfqpoint{3.302050in}{2.169257in}}%
\pgfpathlineto{\pgfqpoint{3.302050in}{2.169257in}}%
\pgfpathlineto{\pgfqpoint{3.302050in}{2.175155in}}%
\pgfpathlineto{\pgfqpoint{3.311133in}{2.175155in}}%
\pgfpathlineto{\pgfqpoint{3.311133in}{2.169257in}}%
\pgfpathmoveto{\pgfqpoint{3.311133in}{2.175155in}}%
\pgfpathlineto{\pgfqpoint{3.311133in}{2.175155in}}%
\pgfpathlineto{\pgfqpoint{3.311133in}{2.181054in}}%
\pgfpathlineto{\pgfqpoint{3.320215in}{2.181054in}}%
\pgfpathlineto{\pgfqpoint{3.320215in}{2.175155in}}%
\pgfpathmoveto{\pgfqpoint{3.311133in}{2.181054in}}%
\pgfpathlineto{\pgfqpoint{3.311133in}{2.181054in}}%
\pgfpathlineto{\pgfqpoint{3.311133in}{2.186952in}}%
\pgfpathlineto{\pgfqpoint{3.320215in}{2.186952in}}%
\pgfpathlineto{\pgfqpoint{3.320215in}{2.181054in}}%
\pgfpathmoveto{\pgfqpoint{3.320215in}{2.181054in}}%
\pgfpathlineto{\pgfqpoint{3.320215in}{2.181054in}}%
\pgfpathlineto{\pgfqpoint{3.320215in}{2.186952in}}%
\pgfpathlineto{\pgfqpoint{3.329297in}{2.186952in}}%
\pgfpathlineto{\pgfqpoint{3.329297in}{2.181054in}}%
\pgfpathmoveto{\pgfqpoint{3.329297in}{2.186952in}}%
\pgfpathlineto{\pgfqpoint{3.329297in}{2.186952in}}%
\pgfpathlineto{\pgfqpoint{3.329297in}{2.192850in}}%
\pgfpathlineto{\pgfqpoint{3.338380in}{2.192850in}}%
\pgfpathlineto{\pgfqpoint{3.338380in}{2.186952in}}%
\pgfpathmoveto{\pgfqpoint{3.329297in}{2.192850in}}%
\pgfpathlineto{\pgfqpoint{3.329297in}{2.192850in}}%
\pgfpathlineto{\pgfqpoint{3.329297in}{2.198748in}}%
\pgfpathlineto{\pgfqpoint{3.338380in}{2.198748in}}%
\pgfpathlineto{\pgfqpoint{3.338380in}{2.192850in}}%
\pgfpathmoveto{\pgfqpoint{3.338380in}{2.192850in}}%
\pgfpathlineto{\pgfqpoint{3.338380in}{2.192850in}}%
\pgfpathlineto{\pgfqpoint{3.338380in}{2.198748in}}%
\pgfpathlineto{\pgfqpoint{3.347462in}{2.198748in}}%
\pgfpathlineto{\pgfqpoint{3.347462in}{2.192850in}}%
\pgfpathmoveto{\pgfqpoint{3.347462in}{2.198748in}}%
\pgfpathlineto{\pgfqpoint{3.347462in}{2.198748in}}%
\pgfpathlineto{\pgfqpoint{3.347462in}{2.204647in}}%
\pgfpathlineto{\pgfqpoint{3.356544in}{2.204647in}}%
\pgfpathlineto{\pgfqpoint{3.356544in}{2.198748in}}%
\pgfpathmoveto{\pgfqpoint{3.347462in}{2.204647in}}%
\pgfpathlineto{\pgfqpoint{3.347462in}{2.204647in}}%
\pgfpathlineto{\pgfqpoint{3.347462in}{2.210545in}}%
\pgfpathlineto{\pgfqpoint{3.356544in}{2.210545in}}%
\pgfpathlineto{\pgfqpoint{3.356544in}{2.204647in}}%
\pgfpathmoveto{\pgfqpoint{3.356544in}{2.204647in}}%
\pgfpathlineto{\pgfqpoint{3.356544in}{2.204647in}}%
\pgfpathlineto{\pgfqpoint{3.356544in}{2.210545in}}%
\pgfpathlineto{\pgfqpoint{3.365627in}{2.210545in}}%
\pgfpathlineto{\pgfqpoint{3.365627in}{2.204647in}}%
\pgfpathmoveto{\pgfqpoint{3.347462in}{2.835780in}}%
\pgfpathlineto{\pgfqpoint{3.347462in}{2.835780in}}%
\pgfpathlineto{\pgfqpoint{3.347462in}{2.841679in}}%
\pgfpathlineto{\pgfqpoint{3.356544in}{2.841679in}}%
\pgfpathlineto{\pgfqpoint{3.356544in}{2.835780in}}%
\pgfpathmoveto{\pgfqpoint{3.347462in}{2.841679in}}%
\pgfpathlineto{\pgfqpoint{3.347462in}{2.841679in}}%
\pgfpathlineto{\pgfqpoint{3.347462in}{2.847577in}}%
\pgfpathlineto{\pgfqpoint{3.356544in}{2.847577in}}%
\pgfpathlineto{\pgfqpoint{3.356544in}{2.841679in}}%
\pgfpathmoveto{\pgfqpoint{3.274803in}{2.930157in}}%
\pgfpathlineto{\pgfqpoint{3.274803in}{2.930157in}}%
\pgfpathlineto{\pgfqpoint{3.274803in}{2.936056in}}%
\pgfpathlineto{\pgfqpoint{3.283886in}{2.936056in}}%
\pgfpathlineto{\pgfqpoint{3.283886in}{2.930157in}}%
\pgfpathmoveto{\pgfqpoint{3.274803in}{2.936056in}}%
\pgfpathlineto{\pgfqpoint{3.274803in}{2.936056in}}%
\pgfpathlineto{\pgfqpoint{3.274803in}{2.941954in}}%
\pgfpathlineto{\pgfqpoint{3.283886in}{2.941954in}}%
\pgfpathlineto{\pgfqpoint{3.283886in}{2.936056in}}%
\pgfpathmoveto{\pgfqpoint{3.311133in}{2.882968in}}%
\pgfpathlineto{\pgfqpoint{3.311133in}{2.882968in}}%
\pgfpathlineto{\pgfqpoint{3.311133in}{2.888867in}}%
\pgfpathlineto{\pgfqpoint{3.320215in}{2.888867in}}%
\pgfpathlineto{\pgfqpoint{3.320215in}{2.882968in}}%
\pgfpathmoveto{\pgfqpoint{3.311133in}{2.888867in}}%
\pgfpathlineto{\pgfqpoint{3.311133in}{2.888867in}}%
\pgfpathlineto{\pgfqpoint{3.311133in}{2.894766in}}%
\pgfpathlineto{\pgfqpoint{3.320215in}{2.894766in}}%
\pgfpathlineto{\pgfqpoint{3.320215in}{2.888867in}}%
\pgfpathmoveto{\pgfqpoint{3.329297in}{2.859374in}}%
\pgfpathlineto{\pgfqpoint{3.329297in}{2.859374in}}%
\pgfpathlineto{\pgfqpoint{3.329297in}{2.865273in}}%
\pgfpathlineto{\pgfqpoint{3.338380in}{2.865273in}}%
\pgfpathlineto{\pgfqpoint{3.338380in}{2.859374in}}%
\pgfpathmoveto{\pgfqpoint{3.329297in}{2.865273in}}%
\pgfpathlineto{\pgfqpoint{3.329297in}{2.865273in}}%
\pgfpathlineto{\pgfqpoint{3.329297in}{2.871171in}}%
\pgfpathlineto{\pgfqpoint{3.338380in}{2.871171in}}%
\pgfpathlineto{\pgfqpoint{3.338380in}{2.865273in}}%
\pgfpathmoveto{\pgfqpoint{3.292968in}{2.906563in}}%
\pgfpathlineto{\pgfqpoint{3.292968in}{2.906563in}}%
\pgfpathlineto{\pgfqpoint{3.292968in}{2.912461in}}%
\pgfpathlineto{\pgfqpoint{3.302050in}{2.912461in}}%
\pgfpathlineto{\pgfqpoint{3.302050in}{2.906563in}}%
\pgfpathmoveto{\pgfqpoint{3.292968in}{2.912461in}}%
\pgfpathlineto{\pgfqpoint{3.292968in}{2.912461in}}%
\pgfpathlineto{\pgfqpoint{3.292968in}{2.918360in}}%
\pgfpathlineto{\pgfqpoint{3.302050in}{2.918360in}}%
\pgfpathlineto{\pgfqpoint{3.302050in}{2.912461in}}%
\pgfpathmoveto{\pgfqpoint{3.238474in}{2.977345in}}%
\pgfpathlineto{\pgfqpoint{3.238474in}{2.977345in}}%
\pgfpathlineto{\pgfqpoint{3.238474in}{2.983243in}}%
\pgfpathlineto{\pgfqpoint{3.247557in}{2.983243in}}%
\pgfpathlineto{\pgfqpoint{3.247557in}{2.977345in}}%
\pgfpathmoveto{\pgfqpoint{3.238474in}{2.983243in}}%
\pgfpathlineto{\pgfqpoint{3.238474in}{2.983243in}}%
\pgfpathlineto{\pgfqpoint{3.238474in}{2.989141in}}%
\pgfpathlineto{\pgfqpoint{3.247557in}{2.989141in}}%
\pgfpathlineto{\pgfqpoint{3.247557in}{2.983243in}}%
\pgfpathmoveto{\pgfqpoint{3.256639in}{2.953751in}}%
\pgfpathlineto{\pgfqpoint{3.256639in}{2.953751in}}%
\pgfpathlineto{\pgfqpoint{3.256639in}{2.959650in}}%
\pgfpathlineto{\pgfqpoint{3.265721in}{2.959650in}}%
\pgfpathlineto{\pgfqpoint{3.265721in}{2.953751in}}%
\pgfpathmoveto{\pgfqpoint{3.256639in}{2.959650in}}%
\pgfpathlineto{\pgfqpoint{3.256639in}{2.959650in}}%
\pgfpathlineto{\pgfqpoint{3.256639in}{2.965548in}}%
\pgfpathlineto{\pgfqpoint{3.265721in}{2.965548in}}%
\pgfpathlineto{\pgfqpoint{3.265721in}{2.959650in}}%
\pgfpathmoveto{\pgfqpoint{3.220310in}{3.000938in}}%
\pgfpathlineto{\pgfqpoint{3.220310in}{3.000938in}}%
\pgfpathlineto{\pgfqpoint{3.220310in}{3.006836in}}%
\pgfpathlineto{\pgfqpoint{3.229392in}{3.006836in}}%
\pgfpathlineto{\pgfqpoint{3.229392in}{3.000938in}}%
\pgfpathmoveto{\pgfqpoint{3.220310in}{3.006836in}}%
\pgfpathlineto{\pgfqpoint{3.220310in}{3.006836in}}%
\pgfpathlineto{\pgfqpoint{3.220310in}{3.012735in}}%
\pgfpathlineto{\pgfqpoint{3.229392in}{3.012735in}}%
\pgfpathlineto{\pgfqpoint{3.229392in}{3.006836in}}%
\pgfpathmoveto{\pgfqpoint{3.365627in}{2.210545in}}%
\pgfpathlineto{\pgfqpoint{3.365627in}{2.210545in}}%
\pgfpathlineto{\pgfqpoint{3.365627in}{2.216443in}}%
\pgfpathlineto{\pgfqpoint{3.374709in}{2.216443in}}%
\pgfpathlineto{\pgfqpoint{3.374709in}{2.210545in}}%
\pgfpathmoveto{\pgfqpoint{3.365627in}{2.216443in}}%
\pgfpathlineto{\pgfqpoint{3.365627in}{2.216443in}}%
\pgfpathlineto{\pgfqpoint{3.365627in}{2.222342in}}%
\pgfpathlineto{\pgfqpoint{3.374709in}{2.222342in}}%
\pgfpathlineto{\pgfqpoint{3.374709in}{2.216443in}}%
\pgfpathmoveto{\pgfqpoint{3.374709in}{2.216443in}}%
\pgfpathlineto{\pgfqpoint{3.374709in}{2.216443in}}%
\pgfpathlineto{\pgfqpoint{3.374709in}{2.222342in}}%
\pgfpathlineto{\pgfqpoint{3.383791in}{2.222342in}}%
\pgfpathlineto{\pgfqpoint{3.383791in}{2.216443in}}%
\pgfpathmoveto{\pgfqpoint{3.383791in}{2.222342in}}%
\pgfpathlineto{\pgfqpoint{3.383791in}{2.222342in}}%
\pgfpathlineto{\pgfqpoint{3.383791in}{2.228240in}}%
\pgfpathlineto{\pgfqpoint{3.392873in}{2.228240in}}%
\pgfpathlineto{\pgfqpoint{3.392873in}{2.222342in}}%
\pgfpathmoveto{\pgfqpoint{3.383791in}{2.228240in}}%
\pgfpathlineto{\pgfqpoint{3.383791in}{2.228240in}}%
\pgfpathlineto{\pgfqpoint{3.383791in}{2.234139in}}%
\pgfpathlineto{\pgfqpoint{3.392873in}{2.234139in}}%
\pgfpathlineto{\pgfqpoint{3.392873in}{2.228240in}}%
\pgfpathmoveto{\pgfqpoint{3.392873in}{2.228240in}}%
\pgfpathlineto{\pgfqpoint{3.392873in}{2.228240in}}%
\pgfpathlineto{\pgfqpoint{3.392873in}{2.234139in}}%
\pgfpathlineto{\pgfqpoint{3.401955in}{2.234139in}}%
\pgfpathlineto{\pgfqpoint{3.401955in}{2.228240in}}%
\pgfpathmoveto{\pgfqpoint{3.401955in}{2.234139in}}%
\pgfpathlineto{\pgfqpoint{3.401955in}{2.234139in}}%
\pgfpathlineto{\pgfqpoint{3.401955in}{2.240037in}}%
\pgfpathlineto{\pgfqpoint{3.411037in}{2.240037in}}%
\pgfpathlineto{\pgfqpoint{3.411037in}{2.234139in}}%
\pgfpathmoveto{\pgfqpoint{3.401955in}{2.240037in}}%
\pgfpathlineto{\pgfqpoint{3.401955in}{2.240037in}}%
\pgfpathlineto{\pgfqpoint{3.401955in}{2.245935in}}%
\pgfpathlineto{\pgfqpoint{3.411037in}{2.245935in}}%
\pgfpathlineto{\pgfqpoint{3.411037in}{2.240037in}}%
\pgfpathmoveto{\pgfqpoint{3.411037in}{2.240037in}}%
\pgfpathlineto{\pgfqpoint{3.411037in}{2.240037in}}%
\pgfpathlineto{\pgfqpoint{3.411037in}{2.245935in}}%
\pgfpathlineto{\pgfqpoint{3.420119in}{2.245935in}}%
\pgfpathlineto{\pgfqpoint{3.420119in}{2.240037in}}%
\pgfpathmoveto{\pgfqpoint{3.420119in}{2.245935in}}%
\pgfpathlineto{\pgfqpoint{3.420119in}{2.245935in}}%
\pgfpathlineto{\pgfqpoint{3.420119in}{2.251834in}}%
\pgfpathlineto{\pgfqpoint{3.429201in}{2.251834in}}%
\pgfpathlineto{\pgfqpoint{3.429201in}{2.245935in}}%
\pgfpathmoveto{\pgfqpoint{3.420119in}{2.251834in}}%
\pgfpathlineto{\pgfqpoint{3.420119in}{2.251834in}}%
\pgfpathlineto{\pgfqpoint{3.420119in}{2.257732in}}%
\pgfpathlineto{\pgfqpoint{3.429201in}{2.257732in}}%
\pgfpathlineto{\pgfqpoint{3.429201in}{2.251834in}}%
\pgfpathmoveto{\pgfqpoint{3.429201in}{2.251834in}}%
\pgfpathlineto{\pgfqpoint{3.429201in}{2.251834in}}%
\pgfpathlineto{\pgfqpoint{3.429201in}{2.257732in}}%
\pgfpathlineto{\pgfqpoint{3.438283in}{2.257732in}}%
\pgfpathlineto{\pgfqpoint{3.438283in}{2.251834in}}%
\pgfpathmoveto{\pgfqpoint{3.438283in}{2.257732in}}%
\pgfpathlineto{\pgfqpoint{3.438283in}{2.257732in}}%
\pgfpathlineto{\pgfqpoint{3.438283in}{2.263630in}}%
\pgfpathlineto{\pgfqpoint{3.447365in}{2.263630in}}%
\pgfpathlineto{\pgfqpoint{3.447365in}{2.257732in}}%
\pgfpathmoveto{\pgfqpoint{3.438283in}{2.263630in}}%
\pgfpathlineto{\pgfqpoint{3.438283in}{2.263630in}}%
\pgfpathlineto{\pgfqpoint{3.438283in}{2.269529in}}%
\pgfpathlineto{\pgfqpoint{3.447365in}{2.269529in}}%
\pgfpathlineto{\pgfqpoint{3.447365in}{2.263630in}}%
\pgfpathmoveto{\pgfqpoint{3.447365in}{2.263630in}}%
\pgfpathlineto{\pgfqpoint{3.447365in}{2.263630in}}%
\pgfpathlineto{\pgfqpoint{3.447365in}{2.269529in}}%
\pgfpathlineto{\pgfqpoint{3.456447in}{2.269529in}}%
\pgfpathlineto{\pgfqpoint{3.456447in}{2.263630in}}%
\pgfpathmoveto{\pgfqpoint{3.456447in}{2.269529in}}%
\pgfpathlineto{\pgfqpoint{3.456447in}{2.269529in}}%
\pgfpathlineto{\pgfqpoint{3.456447in}{2.275427in}}%
\pgfpathlineto{\pgfqpoint{3.465529in}{2.275427in}}%
\pgfpathlineto{\pgfqpoint{3.465529in}{2.269529in}}%
\pgfpathmoveto{\pgfqpoint{3.456447in}{2.275427in}}%
\pgfpathlineto{\pgfqpoint{3.456447in}{2.275427in}}%
\pgfpathlineto{\pgfqpoint{3.456447in}{2.281326in}}%
\pgfpathlineto{\pgfqpoint{3.465529in}{2.281326in}}%
\pgfpathlineto{\pgfqpoint{3.465529in}{2.275427in}}%
\pgfpathmoveto{\pgfqpoint{3.465529in}{2.275427in}}%
\pgfpathlineto{\pgfqpoint{3.465529in}{2.275427in}}%
\pgfpathlineto{\pgfqpoint{3.465529in}{2.281326in}}%
\pgfpathlineto{\pgfqpoint{3.474611in}{2.281326in}}%
\pgfpathlineto{\pgfqpoint{3.474611in}{2.275427in}}%
\pgfpathmoveto{\pgfqpoint{3.474611in}{2.281326in}}%
\pgfpathlineto{\pgfqpoint{3.474611in}{2.281326in}}%
\pgfpathlineto{\pgfqpoint{3.474611in}{2.287224in}}%
\pgfpathlineto{\pgfqpoint{3.483693in}{2.287224in}}%
\pgfpathlineto{\pgfqpoint{3.483693in}{2.281326in}}%
\pgfpathmoveto{\pgfqpoint{3.474611in}{2.287224in}}%
\pgfpathlineto{\pgfqpoint{3.474611in}{2.287224in}}%
\pgfpathlineto{\pgfqpoint{3.474611in}{2.293122in}}%
\pgfpathlineto{\pgfqpoint{3.483693in}{2.293122in}}%
\pgfpathlineto{\pgfqpoint{3.483693in}{2.287224in}}%
\pgfpathmoveto{\pgfqpoint{3.483693in}{2.287224in}}%
\pgfpathlineto{\pgfqpoint{3.483693in}{2.287224in}}%
\pgfpathlineto{\pgfqpoint{3.483693in}{2.293122in}}%
\pgfpathlineto{\pgfqpoint{3.492775in}{2.293122in}}%
\pgfpathlineto{\pgfqpoint{3.492775in}{2.287224in}}%
\pgfpathmoveto{\pgfqpoint{3.492775in}{2.293122in}}%
\pgfpathlineto{\pgfqpoint{3.492775in}{2.293122in}}%
\pgfpathlineto{\pgfqpoint{3.492775in}{2.299021in}}%
\pgfpathlineto{\pgfqpoint{3.501857in}{2.299021in}}%
\pgfpathlineto{\pgfqpoint{3.501857in}{2.293122in}}%
\pgfpathmoveto{\pgfqpoint{3.492775in}{2.299021in}}%
\pgfpathlineto{\pgfqpoint{3.492775in}{2.299021in}}%
\pgfpathlineto{\pgfqpoint{3.492775in}{2.304919in}}%
\pgfpathlineto{\pgfqpoint{3.501857in}{2.304919in}}%
\pgfpathlineto{\pgfqpoint{3.501857in}{2.299021in}}%
\pgfpathmoveto{\pgfqpoint{3.501857in}{2.299021in}}%
\pgfpathlineto{\pgfqpoint{3.501857in}{2.299021in}}%
\pgfpathlineto{\pgfqpoint{3.501857in}{2.304919in}}%
\pgfpathlineto{\pgfqpoint{3.510939in}{2.304919in}}%
\pgfpathlineto{\pgfqpoint{3.510939in}{2.299021in}}%
\pgfpathmoveto{\pgfqpoint{3.492775in}{2.647031in}}%
\pgfpathlineto{\pgfqpoint{3.492775in}{2.647031in}}%
\pgfpathlineto{\pgfqpoint{3.492775in}{2.652929in}}%
\pgfpathlineto{\pgfqpoint{3.501857in}{2.652929in}}%
\pgfpathlineto{\pgfqpoint{3.501857in}{2.647031in}}%
\pgfpathmoveto{\pgfqpoint{3.492775in}{2.652929in}}%
\pgfpathlineto{\pgfqpoint{3.492775in}{2.652929in}}%
\pgfpathlineto{\pgfqpoint{3.492775in}{2.658827in}}%
\pgfpathlineto{\pgfqpoint{3.501857in}{2.658827in}}%
\pgfpathlineto{\pgfqpoint{3.501857in}{2.652929in}}%
\pgfpathmoveto{\pgfqpoint{3.420119in}{2.741405in}}%
\pgfpathlineto{\pgfqpoint{3.420119in}{2.741405in}}%
\pgfpathlineto{\pgfqpoint{3.420119in}{2.747304in}}%
\pgfpathlineto{\pgfqpoint{3.429201in}{2.747304in}}%
\pgfpathlineto{\pgfqpoint{3.429201in}{2.741405in}}%
\pgfpathmoveto{\pgfqpoint{3.420119in}{2.747304in}}%
\pgfpathlineto{\pgfqpoint{3.420119in}{2.747304in}}%
\pgfpathlineto{\pgfqpoint{3.420119in}{2.753202in}}%
\pgfpathlineto{\pgfqpoint{3.429201in}{2.753202in}}%
\pgfpathlineto{\pgfqpoint{3.429201in}{2.747304in}}%
\pgfpathmoveto{\pgfqpoint{3.456447in}{2.694218in}}%
\pgfpathlineto{\pgfqpoint{3.456447in}{2.694218in}}%
\pgfpathlineto{\pgfqpoint{3.456447in}{2.700116in}}%
\pgfpathlineto{\pgfqpoint{3.465529in}{2.700116in}}%
\pgfpathlineto{\pgfqpoint{3.465529in}{2.694218in}}%
\pgfpathmoveto{\pgfqpoint{3.456447in}{2.700116in}}%
\pgfpathlineto{\pgfqpoint{3.456447in}{2.700116in}}%
\pgfpathlineto{\pgfqpoint{3.456447in}{2.706015in}}%
\pgfpathlineto{\pgfqpoint{3.465529in}{2.706015in}}%
\pgfpathlineto{\pgfqpoint{3.465529in}{2.700116in}}%
\pgfpathmoveto{\pgfqpoint{3.474611in}{2.670624in}}%
\pgfpathlineto{\pgfqpoint{3.474611in}{2.670624in}}%
\pgfpathlineto{\pgfqpoint{3.474611in}{2.676522in}}%
\pgfpathlineto{\pgfqpoint{3.483693in}{2.676522in}}%
\pgfpathlineto{\pgfqpoint{3.483693in}{2.670624in}}%
\pgfpathmoveto{\pgfqpoint{3.474611in}{2.676522in}}%
\pgfpathlineto{\pgfqpoint{3.474611in}{2.676522in}}%
\pgfpathlineto{\pgfqpoint{3.474611in}{2.682421in}}%
\pgfpathlineto{\pgfqpoint{3.483693in}{2.682421in}}%
\pgfpathlineto{\pgfqpoint{3.483693in}{2.676522in}}%
\pgfpathmoveto{\pgfqpoint{3.438283in}{2.717812in}}%
\pgfpathlineto{\pgfqpoint{3.438283in}{2.717812in}}%
\pgfpathlineto{\pgfqpoint{3.438283in}{2.723710in}}%
\pgfpathlineto{\pgfqpoint{3.447365in}{2.723710in}}%
\pgfpathlineto{\pgfqpoint{3.447365in}{2.717812in}}%
\pgfpathmoveto{\pgfqpoint{3.438283in}{2.723710in}}%
\pgfpathlineto{\pgfqpoint{3.438283in}{2.723710in}}%
\pgfpathlineto{\pgfqpoint{3.438283in}{2.729608in}}%
\pgfpathlineto{\pgfqpoint{3.447365in}{2.729608in}}%
\pgfpathlineto{\pgfqpoint{3.447365in}{2.723710in}}%
\pgfpathmoveto{\pgfqpoint{3.383791in}{2.788593in}}%
\pgfpathlineto{\pgfqpoint{3.383791in}{2.788593in}}%
\pgfpathlineto{\pgfqpoint{3.383791in}{2.794491in}}%
\pgfpathlineto{\pgfqpoint{3.392873in}{2.794491in}}%
\pgfpathlineto{\pgfqpoint{3.392873in}{2.788593in}}%
\pgfpathmoveto{\pgfqpoint{3.383791in}{2.794491in}}%
\pgfpathlineto{\pgfqpoint{3.383791in}{2.794491in}}%
\pgfpathlineto{\pgfqpoint{3.383791in}{2.800390in}}%
\pgfpathlineto{\pgfqpoint{3.392873in}{2.800390in}}%
\pgfpathlineto{\pgfqpoint{3.392873in}{2.794491in}}%
\pgfpathmoveto{\pgfqpoint{3.401955in}{2.764999in}}%
\pgfpathlineto{\pgfqpoint{3.401955in}{2.764999in}}%
\pgfpathlineto{\pgfqpoint{3.401955in}{2.770898in}}%
\pgfpathlineto{\pgfqpoint{3.411037in}{2.770898in}}%
\pgfpathlineto{\pgfqpoint{3.411037in}{2.764999in}}%
\pgfpathmoveto{\pgfqpoint{3.401955in}{2.770898in}}%
\pgfpathlineto{\pgfqpoint{3.401955in}{2.770898in}}%
\pgfpathlineto{\pgfqpoint{3.401955in}{2.776796in}}%
\pgfpathlineto{\pgfqpoint{3.411037in}{2.776796in}}%
\pgfpathlineto{\pgfqpoint{3.411037in}{2.770898in}}%
\pgfpathmoveto{\pgfqpoint{3.365627in}{2.812187in}}%
\pgfpathlineto{\pgfqpoint{3.365627in}{2.812187in}}%
\pgfpathlineto{\pgfqpoint{3.365627in}{2.818085in}}%
\pgfpathlineto{\pgfqpoint{3.374709in}{2.818085in}}%
\pgfpathlineto{\pgfqpoint{3.374709in}{2.812187in}}%
\pgfpathmoveto{\pgfqpoint{3.365627in}{2.818085in}}%
\pgfpathlineto{\pgfqpoint{3.365627in}{2.818085in}}%
\pgfpathlineto{\pgfqpoint{3.365627in}{2.823983in}}%
\pgfpathlineto{\pgfqpoint{3.374709in}{2.823983in}}%
\pgfpathlineto{\pgfqpoint{3.374709in}{2.818085in}}%
\pgfpathmoveto{\pgfqpoint{3.510939in}{2.304919in}}%
\pgfpathlineto{\pgfqpoint{3.510939in}{2.304919in}}%
\pgfpathlineto{\pgfqpoint{3.510939in}{2.310818in}}%
\pgfpathlineto{\pgfqpoint{3.520020in}{2.310818in}}%
\pgfpathlineto{\pgfqpoint{3.520020in}{2.304919in}}%
\pgfpathmoveto{\pgfqpoint{3.510939in}{2.310818in}}%
\pgfpathlineto{\pgfqpoint{3.510939in}{2.310818in}}%
\pgfpathlineto{\pgfqpoint{3.510939in}{2.316716in}}%
\pgfpathlineto{\pgfqpoint{3.520020in}{2.316716in}}%
\pgfpathlineto{\pgfqpoint{3.520020in}{2.310818in}}%
\pgfpathmoveto{\pgfqpoint{3.520020in}{2.310818in}}%
\pgfpathlineto{\pgfqpoint{3.520020in}{2.310818in}}%
\pgfpathlineto{\pgfqpoint{3.520020in}{2.316716in}}%
\pgfpathlineto{\pgfqpoint{3.529102in}{2.316716in}}%
\pgfpathlineto{\pgfqpoint{3.529102in}{2.310818in}}%
\pgfpathmoveto{\pgfqpoint{3.529102in}{2.316716in}}%
\pgfpathlineto{\pgfqpoint{3.529102in}{2.316716in}}%
\pgfpathlineto{\pgfqpoint{3.529102in}{2.322615in}}%
\pgfpathlineto{\pgfqpoint{3.538184in}{2.322615in}}%
\pgfpathlineto{\pgfqpoint{3.538184in}{2.316716in}}%
\pgfpathmoveto{\pgfqpoint{3.529102in}{2.322615in}}%
\pgfpathlineto{\pgfqpoint{3.529102in}{2.322615in}}%
\pgfpathlineto{\pgfqpoint{3.529102in}{2.328514in}}%
\pgfpathlineto{\pgfqpoint{3.538184in}{2.328514in}}%
\pgfpathlineto{\pgfqpoint{3.538184in}{2.322615in}}%
\pgfpathmoveto{\pgfqpoint{3.538184in}{2.322615in}}%
\pgfpathlineto{\pgfqpoint{3.538184in}{2.322615in}}%
\pgfpathlineto{\pgfqpoint{3.538184in}{2.328514in}}%
\pgfpathlineto{\pgfqpoint{3.547265in}{2.328514in}}%
\pgfpathlineto{\pgfqpoint{3.547265in}{2.322615in}}%
\pgfpathmoveto{\pgfqpoint{3.547265in}{2.328514in}}%
\pgfpathlineto{\pgfqpoint{3.547265in}{2.328514in}}%
\pgfpathlineto{\pgfqpoint{3.547265in}{2.334412in}}%
\pgfpathlineto{\pgfqpoint{3.556347in}{2.334412in}}%
\pgfpathlineto{\pgfqpoint{3.556347in}{2.328514in}}%
\pgfpathmoveto{\pgfqpoint{3.547265in}{2.334412in}}%
\pgfpathlineto{\pgfqpoint{3.547265in}{2.334412in}}%
\pgfpathlineto{\pgfqpoint{3.547265in}{2.340311in}}%
\pgfpathlineto{\pgfqpoint{3.556347in}{2.340311in}}%
\pgfpathlineto{\pgfqpoint{3.556347in}{2.334412in}}%
\pgfpathmoveto{\pgfqpoint{3.556347in}{2.334412in}}%
\pgfpathlineto{\pgfqpoint{3.556347in}{2.334412in}}%
\pgfpathlineto{\pgfqpoint{3.556347in}{2.340311in}}%
\pgfpathlineto{\pgfqpoint{3.565429in}{2.340311in}}%
\pgfpathlineto{\pgfqpoint{3.565429in}{2.334412in}}%
\pgfpathmoveto{\pgfqpoint{3.565429in}{2.340311in}}%
\pgfpathlineto{\pgfqpoint{3.565429in}{2.340311in}}%
\pgfpathlineto{\pgfqpoint{3.565429in}{2.346209in}}%
\pgfpathlineto{\pgfqpoint{3.574510in}{2.346209in}}%
\pgfpathlineto{\pgfqpoint{3.574510in}{2.340311in}}%
\pgfpathmoveto{\pgfqpoint{3.565429in}{2.346209in}}%
\pgfpathlineto{\pgfqpoint{3.565429in}{2.346209in}}%
\pgfpathlineto{\pgfqpoint{3.565429in}{2.352108in}}%
\pgfpathlineto{\pgfqpoint{3.574510in}{2.352108in}}%
\pgfpathlineto{\pgfqpoint{3.574510in}{2.346209in}}%
\pgfpathmoveto{\pgfqpoint{3.574510in}{2.346209in}}%
\pgfpathlineto{\pgfqpoint{3.574510in}{2.346209in}}%
\pgfpathlineto{\pgfqpoint{3.574510in}{2.352108in}}%
\pgfpathlineto{\pgfqpoint{3.583592in}{2.352108in}}%
\pgfpathlineto{\pgfqpoint{3.583592in}{2.346209in}}%
\pgfpathmoveto{\pgfqpoint{3.583592in}{2.352108in}}%
\pgfpathlineto{\pgfqpoint{3.583592in}{2.352108in}}%
\pgfpathlineto{\pgfqpoint{3.583592in}{2.358006in}}%
\pgfpathlineto{\pgfqpoint{3.592674in}{2.358006in}}%
\pgfpathlineto{\pgfqpoint{3.592674in}{2.352108in}}%
\pgfpathmoveto{\pgfqpoint{3.583592in}{2.358006in}}%
\pgfpathlineto{\pgfqpoint{3.583592in}{2.358006in}}%
\pgfpathlineto{\pgfqpoint{3.583592in}{2.363905in}}%
\pgfpathlineto{\pgfqpoint{3.592674in}{2.363905in}}%
\pgfpathlineto{\pgfqpoint{3.592674in}{2.358006in}}%
\pgfpathmoveto{\pgfqpoint{3.592674in}{2.358006in}}%
\pgfpathlineto{\pgfqpoint{3.592674in}{2.358006in}}%
\pgfpathlineto{\pgfqpoint{3.592674in}{2.363905in}}%
\pgfpathlineto{\pgfqpoint{3.601755in}{2.363905in}}%
\pgfpathlineto{\pgfqpoint{3.601755in}{2.358006in}}%
\pgfpathmoveto{\pgfqpoint{3.601755in}{2.358006in}}%
\pgfpathlineto{\pgfqpoint{3.601755in}{2.358006in}}%
\pgfpathlineto{\pgfqpoint{3.601755in}{2.363905in}}%
\pgfpathlineto{\pgfqpoint{3.610837in}{2.363905in}}%
\pgfpathlineto{\pgfqpoint{3.610837in}{2.358006in}}%
\pgfpathmoveto{\pgfqpoint{3.619919in}{2.369803in}}%
\pgfpathlineto{\pgfqpoint{3.619919in}{2.369803in}}%
\pgfpathlineto{\pgfqpoint{3.619919in}{2.375702in}}%
\pgfpathlineto{\pgfqpoint{3.629001in}{2.375702in}}%
\pgfpathlineto{\pgfqpoint{3.629001in}{2.369803in}}%
\pgfpathmoveto{\pgfqpoint{3.638082in}{2.381600in}}%
\pgfpathlineto{\pgfqpoint{3.638082in}{2.381600in}}%
\pgfpathlineto{\pgfqpoint{3.638082in}{2.387499in}}%
\pgfpathlineto{\pgfqpoint{3.647164in}{2.387499in}}%
\pgfpathlineto{\pgfqpoint{3.647164in}{2.381600in}}%
\pgfpathmoveto{\pgfqpoint{3.638082in}{2.458283in}}%
\pgfpathlineto{\pgfqpoint{3.638082in}{2.458283in}}%
\pgfpathlineto{\pgfqpoint{3.638082in}{2.464182in}}%
\pgfpathlineto{\pgfqpoint{3.647164in}{2.464182in}}%
\pgfpathlineto{\pgfqpoint{3.647164in}{2.458283in}}%
\pgfpathmoveto{\pgfqpoint{3.638082in}{2.464182in}}%
\pgfpathlineto{\pgfqpoint{3.638082in}{2.464182in}}%
\pgfpathlineto{\pgfqpoint{3.638082in}{2.470080in}}%
\pgfpathlineto{\pgfqpoint{3.647164in}{2.470080in}}%
\pgfpathlineto{\pgfqpoint{3.647164in}{2.464182in}}%
\pgfpathmoveto{\pgfqpoint{3.565429in}{2.552658in}}%
\pgfpathlineto{\pgfqpoint{3.565429in}{2.552658in}}%
\pgfpathlineto{\pgfqpoint{3.565429in}{2.558556in}}%
\pgfpathlineto{\pgfqpoint{3.574510in}{2.558556in}}%
\pgfpathlineto{\pgfqpoint{3.574510in}{2.552658in}}%
\pgfpathmoveto{\pgfqpoint{3.565429in}{2.558556in}}%
\pgfpathlineto{\pgfqpoint{3.565429in}{2.558556in}}%
\pgfpathlineto{\pgfqpoint{3.565429in}{2.564454in}}%
\pgfpathlineto{\pgfqpoint{3.574510in}{2.564454in}}%
\pgfpathlineto{\pgfqpoint{3.574510in}{2.558556in}}%
\pgfpathmoveto{\pgfqpoint{3.601755in}{2.505471in}}%
\pgfpathlineto{\pgfqpoint{3.601755in}{2.505471in}}%
\pgfpathlineto{\pgfqpoint{3.601755in}{2.511369in}}%
\pgfpathlineto{\pgfqpoint{3.610837in}{2.511369in}}%
\pgfpathlineto{\pgfqpoint{3.610837in}{2.505471in}}%
\pgfpathmoveto{\pgfqpoint{3.601755in}{2.511369in}}%
\pgfpathlineto{\pgfqpoint{3.601755in}{2.511369in}}%
\pgfpathlineto{\pgfqpoint{3.601755in}{2.517268in}}%
\pgfpathlineto{\pgfqpoint{3.610837in}{2.517268in}}%
\pgfpathlineto{\pgfqpoint{3.610837in}{2.511369in}}%
\pgfpathmoveto{\pgfqpoint{3.619919in}{2.481878in}}%
\pgfpathlineto{\pgfqpoint{3.619919in}{2.481878in}}%
\pgfpathlineto{\pgfqpoint{3.619919in}{2.487776in}}%
\pgfpathlineto{\pgfqpoint{3.629001in}{2.487776in}}%
\pgfpathlineto{\pgfqpoint{3.629001in}{2.481878in}}%
\pgfpathmoveto{\pgfqpoint{3.619919in}{2.487776in}}%
\pgfpathlineto{\pgfqpoint{3.619919in}{2.487776in}}%
\pgfpathlineto{\pgfqpoint{3.619919in}{2.493674in}}%
\pgfpathlineto{\pgfqpoint{3.629001in}{2.493674in}}%
\pgfpathlineto{\pgfqpoint{3.629001in}{2.487776in}}%
\pgfpathmoveto{\pgfqpoint{3.583592in}{2.529064in}}%
\pgfpathlineto{\pgfqpoint{3.583592in}{2.529064in}}%
\pgfpathlineto{\pgfqpoint{3.583592in}{2.534963in}}%
\pgfpathlineto{\pgfqpoint{3.592674in}{2.534963in}}%
\pgfpathlineto{\pgfqpoint{3.592674in}{2.529064in}}%
\pgfpathmoveto{\pgfqpoint{3.583592in}{2.534963in}}%
\pgfpathlineto{\pgfqpoint{3.583592in}{2.534963in}}%
\pgfpathlineto{\pgfqpoint{3.583592in}{2.540861in}}%
\pgfpathlineto{\pgfqpoint{3.592674in}{2.540861in}}%
\pgfpathlineto{\pgfqpoint{3.592674in}{2.534963in}}%
\pgfpathmoveto{\pgfqpoint{3.529102in}{2.599844in}}%
\pgfpathlineto{\pgfqpoint{3.529102in}{2.599844in}}%
\pgfpathlineto{\pgfqpoint{3.529102in}{2.605742in}}%
\pgfpathlineto{\pgfqpoint{3.538184in}{2.605742in}}%
\pgfpathlineto{\pgfqpoint{3.538184in}{2.599844in}}%
\pgfpathmoveto{\pgfqpoint{3.529102in}{2.605742in}}%
\pgfpathlineto{\pgfqpoint{3.529102in}{2.605742in}}%
\pgfpathlineto{\pgfqpoint{3.529102in}{2.611641in}}%
\pgfpathlineto{\pgfqpoint{3.538184in}{2.611641in}}%
\pgfpathlineto{\pgfqpoint{3.538184in}{2.605742in}}%
\pgfpathmoveto{\pgfqpoint{3.547265in}{2.576251in}}%
\pgfpathlineto{\pgfqpoint{3.547265in}{2.576251in}}%
\pgfpathlineto{\pgfqpoint{3.547265in}{2.582149in}}%
\pgfpathlineto{\pgfqpoint{3.556347in}{2.582149in}}%
\pgfpathlineto{\pgfqpoint{3.556347in}{2.576251in}}%
\pgfpathmoveto{\pgfqpoint{3.547265in}{2.582149in}}%
\pgfpathlineto{\pgfqpoint{3.547265in}{2.582149in}}%
\pgfpathlineto{\pgfqpoint{3.547265in}{2.588047in}}%
\pgfpathlineto{\pgfqpoint{3.556347in}{2.588047in}}%
\pgfpathlineto{\pgfqpoint{3.556347in}{2.582149in}}%
\pgfpathmoveto{\pgfqpoint{3.510939in}{2.623437in}}%
\pgfpathlineto{\pgfqpoint{3.510939in}{2.623437in}}%
\pgfpathlineto{\pgfqpoint{3.510939in}{2.629336in}}%
\pgfpathlineto{\pgfqpoint{3.520020in}{2.629336in}}%
\pgfpathlineto{\pgfqpoint{3.520020in}{2.623437in}}%
\pgfpathmoveto{\pgfqpoint{3.510939in}{2.629336in}}%
\pgfpathlineto{\pgfqpoint{3.510939in}{2.629336in}}%
\pgfpathlineto{\pgfqpoint{3.510939in}{2.635234in}}%
\pgfpathlineto{\pgfqpoint{3.520020in}{2.635234in}}%
\pgfpathlineto{\pgfqpoint{3.520020in}{2.629336in}}%
\pgfpathmoveto{\pgfqpoint{3.656246in}{2.393397in}}%
\pgfpathlineto{\pgfqpoint{3.656246in}{2.393397in}}%
\pgfpathlineto{\pgfqpoint{3.656246in}{2.399296in}}%
\pgfpathlineto{\pgfqpoint{3.665328in}{2.399296in}}%
\pgfpathlineto{\pgfqpoint{3.665328in}{2.393397in}}%
\pgfpathmoveto{\pgfqpoint{3.674410in}{2.405195in}}%
\pgfpathlineto{\pgfqpoint{3.674410in}{2.405195in}}%
\pgfpathlineto{\pgfqpoint{3.674410in}{2.411093in}}%
\pgfpathlineto{\pgfqpoint{3.683492in}{2.411093in}}%
\pgfpathlineto{\pgfqpoint{3.683492in}{2.405195in}}%
\pgfpathmoveto{\pgfqpoint{3.674410in}{2.411093in}}%
\pgfpathlineto{\pgfqpoint{3.674410in}{2.411093in}}%
\pgfpathlineto{\pgfqpoint{3.674410in}{2.416992in}}%
\pgfpathlineto{\pgfqpoint{3.683492in}{2.416992in}}%
\pgfpathlineto{\pgfqpoint{3.683492in}{2.411093in}}%
\pgfpathmoveto{\pgfqpoint{3.674410in}{2.416992in}}%
\pgfpathlineto{\pgfqpoint{3.674410in}{2.416992in}}%
\pgfpathlineto{\pgfqpoint{3.674410in}{2.422891in}}%
\pgfpathlineto{\pgfqpoint{3.683492in}{2.422891in}}%
\pgfpathlineto{\pgfqpoint{3.683492in}{2.416992in}}%
\pgfpathmoveto{\pgfqpoint{3.656246in}{2.434688in}}%
\pgfpathlineto{\pgfqpoint{3.656246in}{2.434688in}}%
\pgfpathlineto{\pgfqpoint{3.656246in}{2.440587in}}%
\pgfpathlineto{\pgfqpoint{3.665328in}{2.440587in}}%
\pgfpathlineto{\pgfqpoint{3.665328in}{2.434688in}}%
\pgfpathmoveto{\pgfqpoint{3.656246in}{2.440587in}}%
\pgfpathlineto{\pgfqpoint{3.656246in}{2.440587in}}%
\pgfpathlineto{\pgfqpoint{3.656246in}{2.446486in}}%
\pgfpathlineto{\pgfqpoint{3.665328in}{2.446486in}}%
\pgfpathlineto{\pgfqpoint{3.665328in}{2.440587in}}%
\pgfpathmoveto{\pgfqpoint{0.749998in}{0.499998in}}%
\pgfpathlineto{\pgfqpoint{0.749998in}{0.499998in}}%
\pgfpathlineto{\pgfqpoint{0.749998in}{0.502947in}}%
\pgfpathlineto{\pgfqpoint{0.754539in}{0.502947in}}%
\pgfpathlineto{\pgfqpoint{0.754539in}{0.499998in}}%
\pgfpathmoveto{\pgfqpoint{0.749998in}{0.502947in}}%
\pgfpathlineto{\pgfqpoint{0.749998in}{0.502947in}}%
\pgfpathlineto{\pgfqpoint{0.749998in}{0.505897in}}%
\pgfpathlineto{\pgfqpoint{0.754539in}{0.505897in}}%
\pgfpathlineto{\pgfqpoint{0.754539in}{0.502947in}}%
\pgfpathmoveto{\pgfqpoint{0.754539in}{0.502947in}}%
\pgfpathlineto{\pgfqpoint{0.754539in}{0.502947in}}%
\pgfpathlineto{\pgfqpoint{0.754539in}{0.505897in}}%
\pgfpathlineto{\pgfqpoint{0.759080in}{0.505897in}}%
\pgfpathlineto{\pgfqpoint{0.759080in}{0.502947in}}%
\pgfpathmoveto{\pgfqpoint{0.749998in}{0.505897in}}%
\pgfpathlineto{\pgfqpoint{0.749998in}{0.505897in}}%
\pgfpathlineto{\pgfqpoint{0.749998in}{0.508846in}}%
\pgfpathlineto{\pgfqpoint{0.754539in}{0.508846in}}%
\pgfpathlineto{\pgfqpoint{0.754539in}{0.505897in}}%
\pgfpathmoveto{\pgfqpoint{0.749998in}{0.508846in}}%
\pgfpathlineto{\pgfqpoint{0.749998in}{0.508846in}}%
\pgfpathlineto{\pgfqpoint{0.749998in}{0.511795in}}%
\pgfpathlineto{\pgfqpoint{0.754539in}{0.511795in}}%
\pgfpathlineto{\pgfqpoint{0.754539in}{0.508846in}}%
\pgfpathmoveto{\pgfqpoint{0.754539in}{0.505897in}}%
\pgfpathlineto{\pgfqpoint{0.754539in}{0.505897in}}%
\pgfpathlineto{\pgfqpoint{0.754539in}{0.508846in}}%
\pgfpathlineto{\pgfqpoint{0.759080in}{0.508846in}}%
\pgfpathlineto{\pgfqpoint{0.759080in}{0.505897in}}%
\pgfpathmoveto{\pgfqpoint{0.754539in}{0.508846in}}%
\pgfpathlineto{\pgfqpoint{0.754539in}{0.508846in}}%
\pgfpathlineto{\pgfqpoint{0.754539in}{0.511795in}}%
\pgfpathlineto{\pgfqpoint{0.759080in}{0.511795in}}%
\pgfpathlineto{\pgfqpoint{0.759080in}{0.508846in}}%
\pgfpathmoveto{\pgfqpoint{0.759080in}{0.505897in}}%
\pgfpathlineto{\pgfqpoint{0.759080in}{0.505897in}}%
\pgfpathlineto{\pgfqpoint{0.759080in}{0.508846in}}%
\pgfpathlineto{\pgfqpoint{0.763621in}{0.508846in}}%
\pgfpathlineto{\pgfqpoint{0.763621in}{0.505897in}}%
\pgfpathmoveto{\pgfqpoint{0.759080in}{0.508846in}}%
\pgfpathlineto{\pgfqpoint{0.759080in}{0.508846in}}%
\pgfpathlineto{\pgfqpoint{0.759080in}{0.511795in}}%
\pgfpathlineto{\pgfqpoint{0.763621in}{0.511795in}}%
\pgfpathlineto{\pgfqpoint{0.763621in}{0.508846in}}%
\pgfpathmoveto{\pgfqpoint{0.763621in}{0.508846in}}%
\pgfpathlineto{\pgfqpoint{0.763621in}{0.508846in}}%
\pgfpathlineto{\pgfqpoint{0.763621in}{0.511795in}}%
\pgfpathlineto{\pgfqpoint{0.768162in}{0.511795in}}%
\pgfpathlineto{\pgfqpoint{0.768162in}{0.508846in}}%
\pgfpathmoveto{\pgfqpoint{0.759080in}{0.511795in}}%
\pgfpathlineto{\pgfqpoint{0.759080in}{0.511795in}}%
\pgfpathlineto{\pgfqpoint{0.759080in}{0.514744in}}%
\pgfpathlineto{\pgfqpoint{0.763621in}{0.514744in}}%
\pgfpathlineto{\pgfqpoint{0.763621in}{0.511795in}}%
\pgfpathmoveto{\pgfqpoint{0.759080in}{0.514744in}}%
\pgfpathlineto{\pgfqpoint{0.759080in}{0.514744in}}%
\pgfpathlineto{\pgfqpoint{0.759080in}{0.517694in}}%
\pgfpathlineto{\pgfqpoint{0.763621in}{0.517694in}}%
\pgfpathlineto{\pgfqpoint{0.763621in}{0.514744in}}%
\pgfpathmoveto{\pgfqpoint{0.763621in}{0.511795in}}%
\pgfpathlineto{\pgfqpoint{0.763621in}{0.511795in}}%
\pgfpathlineto{\pgfqpoint{0.763621in}{0.514744in}}%
\pgfpathlineto{\pgfqpoint{0.768162in}{0.514744in}}%
\pgfpathlineto{\pgfqpoint{0.768162in}{0.511795in}}%
\pgfpathmoveto{\pgfqpoint{0.763621in}{0.514744in}}%
\pgfpathlineto{\pgfqpoint{0.763621in}{0.514744in}}%
\pgfpathlineto{\pgfqpoint{0.763621in}{0.517694in}}%
\pgfpathlineto{\pgfqpoint{0.768162in}{0.517694in}}%
\pgfpathlineto{\pgfqpoint{0.768162in}{0.514744in}}%
\pgfpathmoveto{\pgfqpoint{0.768162in}{0.511795in}}%
\pgfpathlineto{\pgfqpoint{0.768162in}{0.511795in}}%
\pgfpathlineto{\pgfqpoint{0.768162in}{0.514744in}}%
\pgfpathlineto{\pgfqpoint{0.772703in}{0.514744in}}%
\pgfpathlineto{\pgfqpoint{0.772703in}{0.511795in}}%
\pgfpathmoveto{\pgfqpoint{0.768162in}{0.514744in}}%
\pgfpathlineto{\pgfqpoint{0.768162in}{0.514744in}}%
\pgfpathlineto{\pgfqpoint{0.768162in}{0.517694in}}%
\pgfpathlineto{\pgfqpoint{0.772703in}{0.517694in}}%
\pgfpathlineto{\pgfqpoint{0.772703in}{0.514744in}}%
\pgfpathmoveto{\pgfqpoint{0.772703in}{0.514744in}}%
\pgfpathlineto{\pgfqpoint{0.772703in}{0.514744in}}%
\pgfpathlineto{\pgfqpoint{0.772703in}{0.517694in}}%
\pgfpathlineto{\pgfqpoint{0.777244in}{0.517694in}}%
\pgfpathlineto{\pgfqpoint{0.777244in}{0.514744in}}%
\pgfpathmoveto{\pgfqpoint{0.768162in}{0.517694in}}%
\pgfpathlineto{\pgfqpoint{0.768162in}{0.517694in}}%
\pgfpathlineto{\pgfqpoint{0.768162in}{0.520643in}}%
\pgfpathlineto{\pgfqpoint{0.772703in}{0.520643in}}%
\pgfpathlineto{\pgfqpoint{0.772703in}{0.517694in}}%
\pgfpathmoveto{\pgfqpoint{0.768162in}{0.520643in}}%
\pgfpathlineto{\pgfqpoint{0.768162in}{0.520643in}}%
\pgfpathlineto{\pgfqpoint{0.768162in}{0.523592in}}%
\pgfpathlineto{\pgfqpoint{0.772703in}{0.523592in}}%
\pgfpathlineto{\pgfqpoint{0.772703in}{0.520643in}}%
\pgfpathmoveto{\pgfqpoint{0.772703in}{0.517694in}}%
\pgfpathlineto{\pgfqpoint{0.772703in}{0.517694in}}%
\pgfpathlineto{\pgfqpoint{0.772703in}{0.520643in}}%
\pgfpathlineto{\pgfqpoint{0.777244in}{0.520643in}}%
\pgfpathlineto{\pgfqpoint{0.777244in}{0.517694in}}%
\pgfpathmoveto{\pgfqpoint{0.772703in}{0.520643in}}%
\pgfpathlineto{\pgfqpoint{0.772703in}{0.520643in}}%
\pgfpathlineto{\pgfqpoint{0.772703in}{0.523592in}}%
\pgfpathlineto{\pgfqpoint{0.777244in}{0.523592in}}%
\pgfpathlineto{\pgfqpoint{0.777244in}{0.520643in}}%
\pgfpathmoveto{\pgfqpoint{0.777244in}{0.517694in}}%
\pgfpathlineto{\pgfqpoint{0.777244in}{0.517694in}}%
\pgfpathlineto{\pgfqpoint{0.777244in}{0.520643in}}%
\pgfpathlineto{\pgfqpoint{0.781785in}{0.520643in}}%
\pgfpathlineto{\pgfqpoint{0.781785in}{0.517694in}}%
\pgfpathmoveto{\pgfqpoint{0.777244in}{0.520643in}}%
\pgfpathlineto{\pgfqpoint{0.777244in}{0.520643in}}%
\pgfpathlineto{\pgfqpoint{0.777244in}{0.523592in}}%
\pgfpathlineto{\pgfqpoint{0.781785in}{0.523592in}}%
\pgfpathlineto{\pgfqpoint{0.781785in}{0.520643in}}%
\pgfpathmoveto{\pgfqpoint{0.781785in}{0.520643in}}%
\pgfpathlineto{\pgfqpoint{0.781785in}{0.520643in}}%
\pgfpathlineto{\pgfqpoint{0.781785in}{0.523592in}}%
\pgfpathlineto{\pgfqpoint{0.786327in}{0.523592in}}%
\pgfpathlineto{\pgfqpoint{0.786327in}{0.520643in}}%
\pgfpathmoveto{\pgfqpoint{0.777244in}{0.523592in}}%
\pgfpathlineto{\pgfqpoint{0.777244in}{0.523592in}}%
\pgfpathlineto{\pgfqpoint{0.777244in}{0.526541in}}%
\pgfpathlineto{\pgfqpoint{0.781785in}{0.526541in}}%
\pgfpathlineto{\pgfqpoint{0.781785in}{0.523592in}}%
\pgfpathmoveto{\pgfqpoint{0.777244in}{0.526541in}}%
\pgfpathlineto{\pgfqpoint{0.777244in}{0.526541in}}%
\pgfpathlineto{\pgfqpoint{0.777244in}{0.529491in}}%
\pgfpathlineto{\pgfqpoint{0.781785in}{0.529491in}}%
\pgfpathlineto{\pgfqpoint{0.781785in}{0.526541in}}%
\pgfpathmoveto{\pgfqpoint{0.781785in}{0.523592in}}%
\pgfpathlineto{\pgfqpoint{0.781785in}{0.523592in}}%
\pgfpathlineto{\pgfqpoint{0.781785in}{0.526541in}}%
\pgfpathlineto{\pgfqpoint{0.786327in}{0.526541in}}%
\pgfpathlineto{\pgfqpoint{0.786327in}{0.523592in}}%
\pgfpathmoveto{\pgfqpoint{0.781785in}{0.526541in}}%
\pgfpathlineto{\pgfqpoint{0.781785in}{0.526541in}}%
\pgfpathlineto{\pgfqpoint{0.781785in}{0.529491in}}%
\pgfpathlineto{\pgfqpoint{0.786327in}{0.529491in}}%
\pgfpathlineto{\pgfqpoint{0.786327in}{0.526541in}}%
\pgfpathmoveto{\pgfqpoint{0.786327in}{0.523592in}}%
\pgfpathlineto{\pgfqpoint{0.786327in}{0.523592in}}%
\pgfpathlineto{\pgfqpoint{0.786327in}{0.526541in}}%
\pgfpathlineto{\pgfqpoint{0.790868in}{0.526541in}}%
\pgfpathlineto{\pgfqpoint{0.790868in}{0.523592in}}%
\pgfpathmoveto{\pgfqpoint{0.786327in}{0.526541in}}%
\pgfpathlineto{\pgfqpoint{0.786327in}{0.526541in}}%
\pgfpathlineto{\pgfqpoint{0.786327in}{0.529491in}}%
\pgfpathlineto{\pgfqpoint{0.790868in}{0.529491in}}%
\pgfpathlineto{\pgfqpoint{0.790868in}{0.526541in}}%
\pgfpathmoveto{\pgfqpoint{0.790868in}{0.526541in}}%
\pgfpathlineto{\pgfqpoint{0.790868in}{0.526541in}}%
\pgfpathlineto{\pgfqpoint{0.790868in}{0.529491in}}%
\pgfpathlineto{\pgfqpoint{0.795409in}{0.529491in}}%
\pgfpathlineto{\pgfqpoint{0.795409in}{0.526541in}}%
\pgfpathmoveto{\pgfqpoint{0.786327in}{0.529491in}}%
\pgfpathlineto{\pgfqpoint{0.786327in}{0.529491in}}%
\pgfpathlineto{\pgfqpoint{0.786327in}{0.532440in}}%
\pgfpathlineto{\pgfqpoint{0.790868in}{0.532440in}}%
\pgfpathlineto{\pgfqpoint{0.790868in}{0.529491in}}%
\pgfpathmoveto{\pgfqpoint{0.786327in}{0.532440in}}%
\pgfpathlineto{\pgfqpoint{0.786327in}{0.532440in}}%
\pgfpathlineto{\pgfqpoint{0.786327in}{0.535389in}}%
\pgfpathlineto{\pgfqpoint{0.790868in}{0.535389in}}%
\pgfpathlineto{\pgfqpoint{0.790868in}{0.532440in}}%
\pgfpathmoveto{\pgfqpoint{0.790868in}{0.529491in}}%
\pgfpathlineto{\pgfqpoint{0.790868in}{0.529491in}}%
\pgfpathlineto{\pgfqpoint{0.790868in}{0.532440in}}%
\pgfpathlineto{\pgfqpoint{0.795409in}{0.532440in}}%
\pgfpathlineto{\pgfqpoint{0.795409in}{0.529491in}}%
\pgfpathmoveto{\pgfqpoint{0.790868in}{0.532440in}}%
\pgfpathlineto{\pgfqpoint{0.790868in}{0.532440in}}%
\pgfpathlineto{\pgfqpoint{0.790868in}{0.535389in}}%
\pgfpathlineto{\pgfqpoint{0.795409in}{0.535389in}}%
\pgfpathlineto{\pgfqpoint{0.795409in}{0.532440in}}%
\pgfpathmoveto{\pgfqpoint{0.795409in}{0.529491in}}%
\pgfpathlineto{\pgfqpoint{0.795409in}{0.529491in}}%
\pgfpathlineto{\pgfqpoint{0.795409in}{0.532440in}}%
\pgfpathlineto{\pgfqpoint{0.799950in}{0.532440in}}%
\pgfpathlineto{\pgfqpoint{0.799950in}{0.529491in}}%
\pgfpathmoveto{\pgfqpoint{0.795409in}{0.532440in}}%
\pgfpathlineto{\pgfqpoint{0.795409in}{0.532440in}}%
\pgfpathlineto{\pgfqpoint{0.795409in}{0.535389in}}%
\pgfpathlineto{\pgfqpoint{0.799950in}{0.535389in}}%
\pgfpathlineto{\pgfqpoint{0.799950in}{0.532440in}}%
\pgfpathmoveto{\pgfqpoint{0.799950in}{0.532440in}}%
\pgfpathlineto{\pgfqpoint{0.799950in}{0.532440in}}%
\pgfpathlineto{\pgfqpoint{0.799950in}{0.535389in}}%
\pgfpathlineto{\pgfqpoint{0.804491in}{0.535389in}}%
\pgfpathlineto{\pgfqpoint{0.804491in}{0.532440in}}%
\pgfpathmoveto{\pgfqpoint{0.795409in}{0.535389in}}%
\pgfpathlineto{\pgfqpoint{0.795409in}{0.535389in}}%
\pgfpathlineto{\pgfqpoint{0.795409in}{0.538338in}}%
\pgfpathlineto{\pgfqpoint{0.799950in}{0.538338in}}%
\pgfpathlineto{\pgfqpoint{0.799950in}{0.535389in}}%
\pgfpathmoveto{\pgfqpoint{0.795409in}{0.538338in}}%
\pgfpathlineto{\pgfqpoint{0.795409in}{0.538338in}}%
\pgfpathlineto{\pgfqpoint{0.795409in}{0.541288in}}%
\pgfpathlineto{\pgfqpoint{0.799950in}{0.541288in}}%
\pgfpathlineto{\pgfqpoint{0.799950in}{0.538338in}}%
\pgfpathmoveto{\pgfqpoint{0.799950in}{0.535389in}}%
\pgfpathlineto{\pgfqpoint{0.799950in}{0.535389in}}%
\pgfpathlineto{\pgfqpoint{0.799950in}{0.538338in}}%
\pgfpathlineto{\pgfqpoint{0.804491in}{0.538338in}}%
\pgfpathlineto{\pgfqpoint{0.804491in}{0.535389in}}%
\pgfpathmoveto{\pgfqpoint{0.799950in}{0.538338in}}%
\pgfpathlineto{\pgfqpoint{0.799950in}{0.538338in}}%
\pgfpathlineto{\pgfqpoint{0.799950in}{0.541288in}}%
\pgfpathlineto{\pgfqpoint{0.804491in}{0.541288in}}%
\pgfpathlineto{\pgfqpoint{0.804491in}{0.538338in}}%
\pgfpathmoveto{\pgfqpoint{0.804491in}{0.535389in}}%
\pgfpathlineto{\pgfqpoint{0.804491in}{0.535389in}}%
\pgfpathlineto{\pgfqpoint{0.804491in}{0.538338in}}%
\pgfpathlineto{\pgfqpoint{0.809032in}{0.538338in}}%
\pgfpathlineto{\pgfqpoint{0.809032in}{0.535389in}}%
\pgfpathmoveto{\pgfqpoint{0.804491in}{0.538338in}}%
\pgfpathlineto{\pgfqpoint{0.804491in}{0.538338in}}%
\pgfpathlineto{\pgfqpoint{0.804491in}{0.541288in}}%
\pgfpathlineto{\pgfqpoint{0.809032in}{0.541288in}}%
\pgfpathlineto{\pgfqpoint{0.809032in}{0.538338in}}%
\pgfpathmoveto{\pgfqpoint{0.809032in}{0.538338in}}%
\pgfpathlineto{\pgfqpoint{0.809032in}{0.538338in}}%
\pgfpathlineto{\pgfqpoint{0.809032in}{0.541288in}}%
\pgfpathlineto{\pgfqpoint{0.813573in}{0.541288in}}%
\pgfpathlineto{\pgfqpoint{0.813573in}{0.538338in}}%
\pgfpathmoveto{\pgfqpoint{0.804491in}{0.541288in}}%
\pgfpathlineto{\pgfqpoint{0.804491in}{0.541288in}}%
\pgfpathlineto{\pgfqpoint{0.804491in}{0.544237in}}%
\pgfpathlineto{\pgfqpoint{0.809032in}{0.544237in}}%
\pgfpathlineto{\pgfqpoint{0.809032in}{0.541288in}}%
\pgfpathmoveto{\pgfqpoint{0.804491in}{0.544237in}}%
\pgfpathlineto{\pgfqpoint{0.804491in}{0.544237in}}%
\pgfpathlineto{\pgfqpoint{0.804491in}{0.547186in}}%
\pgfpathlineto{\pgfqpoint{0.809032in}{0.547186in}}%
\pgfpathlineto{\pgfqpoint{0.809032in}{0.544237in}}%
\pgfpathmoveto{\pgfqpoint{0.809032in}{0.541288in}}%
\pgfpathlineto{\pgfqpoint{0.809032in}{0.541288in}}%
\pgfpathlineto{\pgfqpoint{0.809032in}{0.544237in}}%
\pgfpathlineto{\pgfqpoint{0.813573in}{0.544237in}}%
\pgfpathlineto{\pgfqpoint{0.813573in}{0.541288in}}%
\pgfpathmoveto{\pgfqpoint{0.809032in}{0.544237in}}%
\pgfpathlineto{\pgfqpoint{0.809032in}{0.544237in}}%
\pgfpathlineto{\pgfqpoint{0.809032in}{0.547186in}}%
\pgfpathlineto{\pgfqpoint{0.813573in}{0.547186in}}%
\pgfpathlineto{\pgfqpoint{0.813573in}{0.544237in}}%
\pgfpathmoveto{\pgfqpoint{0.813573in}{0.541288in}}%
\pgfpathlineto{\pgfqpoint{0.813573in}{0.541288in}}%
\pgfpathlineto{\pgfqpoint{0.813573in}{0.544237in}}%
\pgfpathlineto{\pgfqpoint{0.818114in}{0.544237in}}%
\pgfpathlineto{\pgfqpoint{0.818114in}{0.541288in}}%
\pgfpathmoveto{\pgfqpoint{0.813573in}{0.544237in}}%
\pgfpathlineto{\pgfqpoint{0.813573in}{0.544237in}}%
\pgfpathlineto{\pgfqpoint{0.813573in}{0.547186in}}%
\pgfpathlineto{\pgfqpoint{0.818114in}{0.547186in}}%
\pgfpathlineto{\pgfqpoint{0.818114in}{0.544237in}}%
\pgfpathmoveto{\pgfqpoint{0.818114in}{0.544237in}}%
\pgfpathlineto{\pgfqpoint{0.818114in}{0.544237in}}%
\pgfpathlineto{\pgfqpoint{0.818114in}{0.547186in}}%
\pgfpathlineto{\pgfqpoint{0.822655in}{0.547186in}}%
\pgfpathlineto{\pgfqpoint{0.822655in}{0.544237in}}%
\pgfpathmoveto{\pgfqpoint{0.813573in}{0.547186in}}%
\pgfpathlineto{\pgfqpoint{0.813573in}{0.547186in}}%
\pgfpathlineto{\pgfqpoint{0.813573in}{0.550135in}}%
\pgfpathlineto{\pgfqpoint{0.818114in}{0.550135in}}%
\pgfpathlineto{\pgfqpoint{0.818114in}{0.547186in}}%
\pgfpathmoveto{\pgfqpoint{0.813573in}{0.550135in}}%
\pgfpathlineto{\pgfqpoint{0.813573in}{0.550135in}}%
\pgfpathlineto{\pgfqpoint{0.813573in}{0.553085in}}%
\pgfpathlineto{\pgfqpoint{0.818114in}{0.553085in}}%
\pgfpathlineto{\pgfqpoint{0.818114in}{0.550135in}}%
\pgfpathmoveto{\pgfqpoint{0.818114in}{0.547186in}}%
\pgfpathlineto{\pgfqpoint{0.818114in}{0.547186in}}%
\pgfpathlineto{\pgfqpoint{0.818114in}{0.550135in}}%
\pgfpathlineto{\pgfqpoint{0.822655in}{0.550135in}}%
\pgfpathlineto{\pgfqpoint{0.822655in}{0.547186in}}%
\pgfpathmoveto{\pgfqpoint{0.818114in}{0.550135in}}%
\pgfpathlineto{\pgfqpoint{0.818114in}{0.550135in}}%
\pgfpathlineto{\pgfqpoint{0.818114in}{0.553085in}}%
\pgfpathlineto{\pgfqpoint{0.822655in}{0.553085in}}%
\pgfpathlineto{\pgfqpoint{0.822655in}{0.550135in}}%
\pgfpathmoveto{\pgfqpoint{0.822655in}{0.547186in}}%
\pgfpathlineto{\pgfqpoint{0.822655in}{0.547186in}}%
\pgfpathlineto{\pgfqpoint{0.822655in}{0.550135in}}%
\pgfpathlineto{\pgfqpoint{0.827196in}{0.550135in}}%
\pgfpathlineto{\pgfqpoint{0.827196in}{0.547186in}}%
\pgfpathmoveto{\pgfqpoint{0.822655in}{0.550135in}}%
\pgfpathlineto{\pgfqpoint{0.822655in}{0.550135in}}%
\pgfpathlineto{\pgfqpoint{0.822655in}{0.553085in}}%
\pgfpathlineto{\pgfqpoint{0.827196in}{0.553085in}}%
\pgfpathlineto{\pgfqpoint{0.827196in}{0.550135in}}%
\pgfpathmoveto{\pgfqpoint{0.827196in}{0.550135in}}%
\pgfpathlineto{\pgfqpoint{0.827196in}{0.550135in}}%
\pgfpathlineto{\pgfqpoint{0.827196in}{0.553085in}}%
\pgfpathlineto{\pgfqpoint{0.831737in}{0.553085in}}%
\pgfpathlineto{\pgfqpoint{0.831737in}{0.550135in}}%
\pgfpathmoveto{\pgfqpoint{0.822655in}{0.553085in}}%
\pgfpathlineto{\pgfqpoint{0.822655in}{0.553085in}}%
\pgfpathlineto{\pgfqpoint{0.822655in}{0.556034in}}%
\pgfpathlineto{\pgfqpoint{0.827196in}{0.556034in}}%
\pgfpathlineto{\pgfqpoint{0.827196in}{0.553085in}}%
\pgfpathmoveto{\pgfqpoint{0.822655in}{0.556034in}}%
\pgfpathlineto{\pgfqpoint{0.822655in}{0.556034in}}%
\pgfpathlineto{\pgfqpoint{0.822655in}{0.558983in}}%
\pgfpathlineto{\pgfqpoint{0.827196in}{0.558983in}}%
\pgfpathlineto{\pgfqpoint{0.827196in}{0.556034in}}%
\pgfpathmoveto{\pgfqpoint{0.827196in}{0.553085in}}%
\pgfpathlineto{\pgfqpoint{0.827196in}{0.553085in}}%
\pgfpathlineto{\pgfqpoint{0.827196in}{0.556034in}}%
\pgfpathlineto{\pgfqpoint{0.831737in}{0.556034in}}%
\pgfpathlineto{\pgfqpoint{0.831737in}{0.553085in}}%
\pgfpathmoveto{\pgfqpoint{0.827196in}{0.556034in}}%
\pgfpathlineto{\pgfqpoint{0.827196in}{0.556034in}}%
\pgfpathlineto{\pgfqpoint{0.827196in}{0.558983in}}%
\pgfpathlineto{\pgfqpoint{0.831737in}{0.558983in}}%
\pgfpathlineto{\pgfqpoint{0.831737in}{0.556034in}}%
\pgfpathmoveto{\pgfqpoint{0.831737in}{0.553085in}}%
\pgfpathlineto{\pgfqpoint{0.831737in}{0.553085in}}%
\pgfpathlineto{\pgfqpoint{0.831737in}{0.556034in}}%
\pgfpathlineto{\pgfqpoint{0.836278in}{0.556034in}}%
\pgfpathlineto{\pgfqpoint{0.836278in}{0.553085in}}%
\pgfpathmoveto{\pgfqpoint{0.831737in}{0.556034in}}%
\pgfpathlineto{\pgfqpoint{0.831737in}{0.556034in}}%
\pgfpathlineto{\pgfqpoint{0.831737in}{0.558983in}}%
\pgfpathlineto{\pgfqpoint{0.836278in}{0.558983in}}%
\pgfpathlineto{\pgfqpoint{0.836278in}{0.556034in}}%
\pgfpathmoveto{\pgfqpoint{0.836278in}{0.556034in}}%
\pgfpathlineto{\pgfqpoint{0.836278in}{0.556034in}}%
\pgfpathlineto{\pgfqpoint{0.836278in}{0.558983in}}%
\pgfpathlineto{\pgfqpoint{0.840819in}{0.558983in}}%
\pgfpathlineto{\pgfqpoint{0.840819in}{0.556034in}}%
\pgfpathmoveto{\pgfqpoint{0.831737in}{0.558983in}}%
\pgfpathlineto{\pgfqpoint{0.831737in}{0.558983in}}%
\pgfpathlineto{\pgfqpoint{0.831737in}{0.561932in}}%
\pgfpathlineto{\pgfqpoint{0.836278in}{0.561932in}}%
\pgfpathlineto{\pgfqpoint{0.836278in}{0.558983in}}%
\pgfpathmoveto{\pgfqpoint{0.831737in}{0.561932in}}%
\pgfpathlineto{\pgfqpoint{0.831737in}{0.561932in}}%
\pgfpathlineto{\pgfqpoint{0.831737in}{0.564882in}}%
\pgfpathlineto{\pgfqpoint{0.836278in}{0.564882in}}%
\pgfpathlineto{\pgfqpoint{0.836278in}{0.561932in}}%
\pgfpathmoveto{\pgfqpoint{0.836278in}{0.558983in}}%
\pgfpathlineto{\pgfqpoint{0.836278in}{0.558983in}}%
\pgfpathlineto{\pgfqpoint{0.836278in}{0.561932in}}%
\pgfpathlineto{\pgfqpoint{0.840819in}{0.561932in}}%
\pgfpathlineto{\pgfqpoint{0.840819in}{0.558983in}}%
\pgfpathmoveto{\pgfqpoint{0.836278in}{0.561932in}}%
\pgfpathlineto{\pgfqpoint{0.836278in}{0.561932in}}%
\pgfpathlineto{\pgfqpoint{0.836278in}{0.564882in}}%
\pgfpathlineto{\pgfqpoint{0.840819in}{0.564882in}}%
\pgfpathlineto{\pgfqpoint{0.840819in}{0.561932in}}%
\pgfpathmoveto{\pgfqpoint{0.840819in}{0.558983in}}%
\pgfpathlineto{\pgfqpoint{0.840819in}{0.558983in}}%
\pgfpathlineto{\pgfqpoint{0.840819in}{0.561932in}}%
\pgfpathlineto{\pgfqpoint{0.845360in}{0.561932in}}%
\pgfpathlineto{\pgfqpoint{0.845360in}{0.558983in}}%
\pgfpathmoveto{\pgfqpoint{0.840819in}{0.561932in}}%
\pgfpathlineto{\pgfqpoint{0.840819in}{0.561932in}}%
\pgfpathlineto{\pgfqpoint{0.840819in}{0.564882in}}%
\pgfpathlineto{\pgfqpoint{0.845360in}{0.564882in}}%
\pgfpathlineto{\pgfqpoint{0.845360in}{0.561932in}}%
\pgfpathmoveto{\pgfqpoint{0.845360in}{0.561932in}}%
\pgfpathlineto{\pgfqpoint{0.845360in}{0.561932in}}%
\pgfpathlineto{\pgfqpoint{0.845360in}{0.564882in}}%
\pgfpathlineto{\pgfqpoint{0.849901in}{0.564882in}}%
\pgfpathlineto{\pgfqpoint{0.849901in}{0.561932in}}%
\pgfpathmoveto{\pgfqpoint{0.840819in}{0.564882in}}%
\pgfpathlineto{\pgfqpoint{0.840819in}{0.564882in}}%
\pgfpathlineto{\pgfqpoint{0.840819in}{0.567831in}}%
\pgfpathlineto{\pgfqpoint{0.845360in}{0.567831in}}%
\pgfpathlineto{\pgfqpoint{0.845360in}{0.564882in}}%
\pgfpathmoveto{\pgfqpoint{0.840819in}{0.567831in}}%
\pgfpathlineto{\pgfqpoint{0.840819in}{0.567831in}}%
\pgfpathlineto{\pgfqpoint{0.840819in}{0.570780in}}%
\pgfpathlineto{\pgfqpoint{0.845360in}{0.570780in}}%
\pgfpathlineto{\pgfqpoint{0.845360in}{0.567831in}}%
\pgfpathmoveto{\pgfqpoint{0.845360in}{0.564882in}}%
\pgfpathlineto{\pgfqpoint{0.845360in}{0.564882in}}%
\pgfpathlineto{\pgfqpoint{0.845360in}{0.567831in}}%
\pgfpathlineto{\pgfqpoint{0.849901in}{0.567831in}}%
\pgfpathlineto{\pgfqpoint{0.849901in}{0.564882in}}%
\pgfpathmoveto{\pgfqpoint{0.845360in}{0.567831in}}%
\pgfpathlineto{\pgfqpoint{0.845360in}{0.567831in}}%
\pgfpathlineto{\pgfqpoint{0.845360in}{0.570780in}}%
\pgfpathlineto{\pgfqpoint{0.849901in}{0.570780in}}%
\pgfpathlineto{\pgfqpoint{0.849901in}{0.567831in}}%
\pgfpathmoveto{\pgfqpoint{0.849901in}{0.564882in}}%
\pgfpathlineto{\pgfqpoint{0.849901in}{0.564882in}}%
\pgfpathlineto{\pgfqpoint{0.849901in}{0.567831in}}%
\pgfpathlineto{\pgfqpoint{0.854442in}{0.567831in}}%
\pgfpathlineto{\pgfqpoint{0.854442in}{0.564882in}}%
\pgfpathmoveto{\pgfqpoint{0.849901in}{0.567831in}}%
\pgfpathlineto{\pgfqpoint{0.849901in}{0.567831in}}%
\pgfpathlineto{\pgfqpoint{0.849901in}{0.570780in}}%
\pgfpathlineto{\pgfqpoint{0.854442in}{0.570780in}}%
\pgfpathlineto{\pgfqpoint{0.854442in}{0.567831in}}%
\pgfpathmoveto{\pgfqpoint{0.854442in}{0.567831in}}%
\pgfpathlineto{\pgfqpoint{0.854442in}{0.567831in}}%
\pgfpathlineto{\pgfqpoint{0.854442in}{0.570780in}}%
\pgfpathlineto{\pgfqpoint{0.858983in}{0.570780in}}%
\pgfpathlineto{\pgfqpoint{0.858983in}{0.567831in}}%
\pgfpathmoveto{\pgfqpoint{0.849901in}{0.570780in}}%
\pgfpathlineto{\pgfqpoint{0.849901in}{0.570780in}}%
\pgfpathlineto{\pgfqpoint{0.849901in}{0.573729in}}%
\pgfpathlineto{\pgfqpoint{0.854442in}{0.573729in}}%
\pgfpathlineto{\pgfqpoint{0.854442in}{0.570780in}}%
\pgfpathmoveto{\pgfqpoint{0.849901in}{0.573729in}}%
\pgfpathlineto{\pgfqpoint{0.849901in}{0.573729in}}%
\pgfpathlineto{\pgfqpoint{0.849901in}{0.576679in}}%
\pgfpathlineto{\pgfqpoint{0.854442in}{0.576679in}}%
\pgfpathlineto{\pgfqpoint{0.854442in}{0.573729in}}%
\pgfpathmoveto{\pgfqpoint{0.854442in}{0.570780in}}%
\pgfpathlineto{\pgfqpoint{0.854442in}{0.570780in}}%
\pgfpathlineto{\pgfqpoint{0.854442in}{0.573729in}}%
\pgfpathlineto{\pgfqpoint{0.858983in}{0.573729in}}%
\pgfpathlineto{\pgfqpoint{0.858983in}{0.570780in}}%
\pgfpathmoveto{\pgfqpoint{0.854442in}{0.573729in}}%
\pgfpathlineto{\pgfqpoint{0.854442in}{0.573729in}}%
\pgfpathlineto{\pgfqpoint{0.854442in}{0.576679in}}%
\pgfpathlineto{\pgfqpoint{0.858983in}{0.576679in}}%
\pgfpathlineto{\pgfqpoint{0.858983in}{0.573729in}}%
\pgfpathmoveto{\pgfqpoint{0.858983in}{0.570780in}}%
\pgfpathlineto{\pgfqpoint{0.858983in}{0.570780in}}%
\pgfpathlineto{\pgfqpoint{0.858983in}{0.573729in}}%
\pgfpathlineto{\pgfqpoint{0.863524in}{0.573729in}}%
\pgfpathlineto{\pgfqpoint{0.863524in}{0.570780in}}%
\pgfpathmoveto{\pgfqpoint{0.858983in}{0.573729in}}%
\pgfpathlineto{\pgfqpoint{0.858983in}{0.573729in}}%
\pgfpathlineto{\pgfqpoint{0.858983in}{0.576679in}}%
\pgfpathlineto{\pgfqpoint{0.863524in}{0.576679in}}%
\pgfpathlineto{\pgfqpoint{0.863524in}{0.573729in}}%
\pgfpathmoveto{\pgfqpoint{0.863524in}{0.573729in}}%
\pgfpathlineto{\pgfqpoint{0.863524in}{0.573729in}}%
\pgfpathlineto{\pgfqpoint{0.863524in}{0.576679in}}%
\pgfpathlineto{\pgfqpoint{0.868065in}{0.576679in}}%
\pgfpathlineto{\pgfqpoint{0.868065in}{0.573729in}}%
\pgfpathmoveto{\pgfqpoint{0.858983in}{0.576679in}}%
\pgfpathlineto{\pgfqpoint{0.858983in}{0.576679in}}%
\pgfpathlineto{\pgfqpoint{0.858983in}{0.579628in}}%
\pgfpathlineto{\pgfqpoint{0.863524in}{0.579628in}}%
\pgfpathlineto{\pgfqpoint{0.863524in}{0.576679in}}%
\pgfpathmoveto{\pgfqpoint{0.858983in}{0.579628in}}%
\pgfpathlineto{\pgfqpoint{0.858983in}{0.579628in}}%
\pgfpathlineto{\pgfqpoint{0.858983in}{0.582577in}}%
\pgfpathlineto{\pgfqpoint{0.863524in}{0.582577in}}%
\pgfpathlineto{\pgfqpoint{0.863524in}{0.579628in}}%
\pgfpathmoveto{\pgfqpoint{0.863524in}{0.576679in}}%
\pgfpathlineto{\pgfqpoint{0.863524in}{0.576679in}}%
\pgfpathlineto{\pgfqpoint{0.863524in}{0.579628in}}%
\pgfpathlineto{\pgfqpoint{0.868065in}{0.579628in}}%
\pgfpathlineto{\pgfqpoint{0.868065in}{0.576679in}}%
\pgfpathmoveto{\pgfqpoint{0.863524in}{0.579628in}}%
\pgfpathlineto{\pgfqpoint{0.863524in}{0.579628in}}%
\pgfpathlineto{\pgfqpoint{0.863524in}{0.582577in}}%
\pgfpathlineto{\pgfqpoint{0.868065in}{0.582577in}}%
\pgfpathlineto{\pgfqpoint{0.868065in}{0.579628in}}%
\pgfpathmoveto{\pgfqpoint{0.868065in}{0.576679in}}%
\pgfpathlineto{\pgfqpoint{0.868065in}{0.576679in}}%
\pgfpathlineto{\pgfqpoint{0.868065in}{0.579628in}}%
\pgfpathlineto{\pgfqpoint{0.872606in}{0.579628in}}%
\pgfpathlineto{\pgfqpoint{0.872606in}{0.576679in}}%
\pgfpathmoveto{\pgfqpoint{0.868065in}{0.579628in}}%
\pgfpathlineto{\pgfqpoint{0.868065in}{0.579628in}}%
\pgfpathlineto{\pgfqpoint{0.868065in}{0.582577in}}%
\pgfpathlineto{\pgfqpoint{0.872606in}{0.582577in}}%
\pgfpathlineto{\pgfqpoint{0.872606in}{0.579628in}}%
\pgfpathmoveto{\pgfqpoint{0.872606in}{0.579628in}}%
\pgfpathlineto{\pgfqpoint{0.872606in}{0.579628in}}%
\pgfpathlineto{\pgfqpoint{0.872606in}{0.582577in}}%
\pgfpathlineto{\pgfqpoint{0.877147in}{0.582577in}}%
\pgfpathlineto{\pgfqpoint{0.877147in}{0.579628in}}%
\pgfpathmoveto{\pgfqpoint{0.868065in}{0.582577in}}%
\pgfpathlineto{\pgfqpoint{0.868065in}{0.582577in}}%
\pgfpathlineto{\pgfqpoint{0.868065in}{0.585526in}}%
\pgfpathlineto{\pgfqpoint{0.872606in}{0.585526in}}%
\pgfpathlineto{\pgfqpoint{0.872606in}{0.582577in}}%
\pgfpathmoveto{\pgfqpoint{0.868065in}{0.585526in}}%
\pgfpathlineto{\pgfqpoint{0.868065in}{0.585526in}}%
\pgfpathlineto{\pgfqpoint{0.868065in}{0.588476in}}%
\pgfpathlineto{\pgfqpoint{0.872606in}{0.588476in}}%
\pgfpathlineto{\pgfqpoint{0.872606in}{0.585526in}}%
\pgfpathmoveto{\pgfqpoint{0.872606in}{0.582577in}}%
\pgfpathlineto{\pgfqpoint{0.872606in}{0.582577in}}%
\pgfpathlineto{\pgfqpoint{0.872606in}{0.585526in}}%
\pgfpathlineto{\pgfqpoint{0.877147in}{0.585526in}}%
\pgfpathlineto{\pgfqpoint{0.877147in}{0.582577in}}%
\pgfpathmoveto{\pgfqpoint{0.872606in}{0.585526in}}%
\pgfpathlineto{\pgfqpoint{0.872606in}{0.585526in}}%
\pgfpathlineto{\pgfqpoint{0.872606in}{0.588476in}}%
\pgfpathlineto{\pgfqpoint{0.877147in}{0.588476in}}%
\pgfpathlineto{\pgfqpoint{0.877147in}{0.585526in}}%
\pgfpathmoveto{\pgfqpoint{0.877147in}{0.582577in}}%
\pgfpathlineto{\pgfqpoint{0.877147in}{0.582577in}}%
\pgfpathlineto{\pgfqpoint{0.877147in}{0.585526in}}%
\pgfpathlineto{\pgfqpoint{0.881689in}{0.585526in}}%
\pgfpathlineto{\pgfqpoint{0.881689in}{0.582577in}}%
\pgfpathmoveto{\pgfqpoint{0.877147in}{0.585526in}}%
\pgfpathlineto{\pgfqpoint{0.877147in}{0.585526in}}%
\pgfpathlineto{\pgfqpoint{0.877147in}{0.588476in}}%
\pgfpathlineto{\pgfqpoint{0.881689in}{0.588476in}}%
\pgfpathlineto{\pgfqpoint{0.881689in}{0.585526in}}%
\pgfpathmoveto{\pgfqpoint{0.881689in}{0.585526in}}%
\pgfpathlineto{\pgfqpoint{0.881689in}{0.585526in}}%
\pgfpathlineto{\pgfqpoint{0.881689in}{0.588476in}}%
\pgfpathlineto{\pgfqpoint{0.886230in}{0.588476in}}%
\pgfpathlineto{\pgfqpoint{0.886230in}{0.585526in}}%
\pgfpathmoveto{\pgfqpoint{0.877147in}{0.588476in}}%
\pgfpathlineto{\pgfqpoint{0.877147in}{0.588476in}}%
\pgfpathlineto{\pgfqpoint{0.877147in}{0.591425in}}%
\pgfpathlineto{\pgfqpoint{0.881689in}{0.591425in}}%
\pgfpathlineto{\pgfqpoint{0.881689in}{0.588476in}}%
\pgfpathmoveto{\pgfqpoint{0.877147in}{0.591425in}}%
\pgfpathlineto{\pgfqpoint{0.877147in}{0.591425in}}%
\pgfpathlineto{\pgfqpoint{0.877147in}{0.594374in}}%
\pgfpathlineto{\pgfqpoint{0.881689in}{0.594374in}}%
\pgfpathlineto{\pgfqpoint{0.881689in}{0.591425in}}%
\pgfpathmoveto{\pgfqpoint{0.881689in}{0.588476in}}%
\pgfpathlineto{\pgfqpoint{0.881689in}{0.588476in}}%
\pgfpathlineto{\pgfqpoint{0.881689in}{0.591425in}}%
\pgfpathlineto{\pgfqpoint{0.886230in}{0.591425in}}%
\pgfpathlineto{\pgfqpoint{0.886230in}{0.588476in}}%
\pgfpathmoveto{\pgfqpoint{0.881689in}{0.591425in}}%
\pgfpathlineto{\pgfqpoint{0.881689in}{0.591425in}}%
\pgfpathlineto{\pgfqpoint{0.881689in}{0.594374in}}%
\pgfpathlineto{\pgfqpoint{0.886230in}{0.594374in}}%
\pgfpathlineto{\pgfqpoint{0.886230in}{0.591425in}}%
\pgfpathmoveto{\pgfqpoint{0.886230in}{0.588476in}}%
\pgfpathlineto{\pgfqpoint{0.886230in}{0.588476in}}%
\pgfpathlineto{\pgfqpoint{0.886230in}{0.591425in}}%
\pgfpathlineto{\pgfqpoint{0.890771in}{0.591425in}}%
\pgfpathlineto{\pgfqpoint{0.890771in}{0.588476in}}%
\pgfpathmoveto{\pgfqpoint{0.886230in}{0.591425in}}%
\pgfpathlineto{\pgfqpoint{0.886230in}{0.591425in}}%
\pgfpathlineto{\pgfqpoint{0.886230in}{0.594374in}}%
\pgfpathlineto{\pgfqpoint{0.890771in}{0.594374in}}%
\pgfpathlineto{\pgfqpoint{0.890771in}{0.591425in}}%
\pgfpathmoveto{\pgfqpoint{0.890771in}{0.591425in}}%
\pgfpathlineto{\pgfqpoint{0.890771in}{0.591425in}}%
\pgfpathlineto{\pgfqpoint{0.890771in}{0.594374in}}%
\pgfpathlineto{\pgfqpoint{0.895312in}{0.594374in}}%
\pgfpathlineto{\pgfqpoint{0.895312in}{0.591425in}}%
\pgfpathmoveto{\pgfqpoint{0.886230in}{0.594374in}}%
\pgfpathlineto{\pgfqpoint{0.886230in}{0.594374in}}%
\pgfpathlineto{\pgfqpoint{0.886230in}{0.597323in}}%
\pgfpathlineto{\pgfqpoint{0.890771in}{0.597323in}}%
\pgfpathlineto{\pgfqpoint{0.890771in}{0.594374in}}%
\pgfpathmoveto{\pgfqpoint{0.886230in}{0.597323in}}%
\pgfpathlineto{\pgfqpoint{0.886230in}{0.597323in}}%
\pgfpathlineto{\pgfqpoint{0.886230in}{0.600273in}}%
\pgfpathlineto{\pgfqpoint{0.890771in}{0.600273in}}%
\pgfpathlineto{\pgfqpoint{0.890771in}{0.597323in}}%
\pgfpathmoveto{\pgfqpoint{0.890771in}{0.594374in}}%
\pgfpathlineto{\pgfqpoint{0.890771in}{0.594374in}}%
\pgfpathlineto{\pgfqpoint{0.890771in}{0.597323in}}%
\pgfpathlineto{\pgfqpoint{0.895312in}{0.597323in}}%
\pgfpathlineto{\pgfqpoint{0.895312in}{0.594374in}}%
\pgfpathmoveto{\pgfqpoint{0.890771in}{0.597323in}}%
\pgfpathlineto{\pgfqpoint{0.890771in}{0.597323in}}%
\pgfpathlineto{\pgfqpoint{0.890771in}{0.600273in}}%
\pgfpathlineto{\pgfqpoint{0.895312in}{0.600273in}}%
\pgfpathlineto{\pgfqpoint{0.895312in}{0.597323in}}%
\pgfpathmoveto{\pgfqpoint{0.895312in}{0.594374in}}%
\pgfpathlineto{\pgfqpoint{0.895312in}{0.594374in}}%
\pgfpathlineto{\pgfqpoint{0.895312in}{0.597323in}}%
\pgfpathlineto{\pgfqpoint{0.899853in}{0.597323in}}%
\pgfpathlineto{\pgfqpoint{0.899853in}{0.594374in}}%
\pgfpathmoveto{\pgfqpoint{0.895312in}{0.597323in}}%
\pgfpathlineto{\pgfqpoint{0.895312in}{0.597323in}}%
\pgfpathlineto{\pgfqpoint{0.895312in}{0.600273in}}%
\pgfpathlineto{\pgfqpoint{0.899853in}{0.600273in}}%
\pgfpathlineto{\pgfqpoint{0.899853in}{0.597323in}}%
\pgfpathmoveto{\pgfqpoint{0.899853in}{0.597323in}}%
\pgfpathlineto{\pgfqpoint{0.899853in}{0.597323in}}%
\pgfpathlineto{\pgfqpoint{0.899853in}{0.600273in}}%
\pgfpathlineto{\pgfqpoint{0.904394in}{0.600273in}}%
\pgfpathlineto{\pgfqpoint{0.904394in}{0.597323in}}%
\pgfpathmoveto{\pgfqpoint{0.895312in}{0.600273in}}%
\pgfpathlineto{\pgfqpoint{0.895312in}{0.600273in}}%
\pgfpathlineto{\pgfqpoint{0.895312in}{0.603222in}}%
\pgfpathlineto{\pgfqpoint{0.899853in}{0.603222in}}%
\pgfpathlineto{\pgfqpoint{0.899853in}{0.600273in}}%
\pgfpathmoveto{\pgfqpoint{0.895312in}{0.603222in}}%
\pgfpathlineto{\pgfqpoint{0.895312in}{0.603222in}}%
\pgfpathlineto{\pgfqpoint{0.895312in}{0.606171in}}%
\pgfpathlineto{\pgfqpoint{0.899853in}{0.606171in}}%
\pgfpathlineto{\pgfqpoint{0.899853in}{0.603222in}}%
\pgfpathmoveto{\pgfqpoint{0.899853in}{0.600273in}}%
\pgfpathlineto{\pgfqpoint{0.899853in}{0.600273in}}%
\pgfpathlineto{\pgfqpoint{0.899853in}{0.603222in}}%
\pgfpathlineto{\pgfqpoint{0.904394in}{0.603222in}}%
\pgfpathlineto{\pgfqpoint{0.904394in}{0.600273in}}%
\pgfpathmoveto{\pgfqpoint{0.899853in}{0.603222in}}%
\pgfpathlineto{\pgfqpoint{0.899853in}{0.603222in}}%
\pgfpathlineto{\pgfqpoint{0.899853in}{0.606171in}}%
\pgfpathlineto{\pgfqpoint{0.904394in}{0.606171in}}%
\pgfpathlineto{\pgfqpoint{0.904394in}{0.603222in}}%
\pgfpathmoveto{\pgfqpoint{0.904394in}{0.600273in}}%
\pgfpathlineto{\pgfqpoint{0.904394in}{0.600273in}}%
\pgfpathlineto{\pgfqpoint{0.904394in}{0.603222in}}%
\pgfpathlineto{\pgfqpoint{0.908935in}{0.603222in}}%
\pgfpathlineto{\pgfqpoint{0.908935in}{0.600273in}}%
\pgfpathmoveto{\pgfqpoint{0.904394in}{0.603222in}}%
\pgfpathlineto{\pgfqpoint{0.904394in}{0.603222in}}%
\pgfpathlineto{\pgfqpoint{0.904394in}{0.606171in}}%
\pgfpathlineto{\pgfqpoint{0.908935in}{0.606171in}}%
\pgfpathlineto{\pgfqpoint{0.908935in}{0.603222in}}%
\pgfpathmoveto{\pgfqpoint{0.908935in}{0.603222in}}%
\pgfpathlineto{\pgfqpoint{0.908935in}{0.603222in}}%
\pgfpathlineto{\pgfqpoint{0.908935in}{0.606171in}}%
\pgfpathlineto{\pgfqpoint{0.913475in}{0.606171in}}%
\pgfpathlineto{\pgfqpoint{0.913475in}{0.603222in}}%
\pgfpathmoveto{\pgfqpoint{0.904394in}{0.606171in}}%
\pgfpathlineto{\pgfqpoint{0.904394in}{0.606171in}}%
\pgfpathlineto{\pgfqpoint{0.904394in}{0.609120in}}%
\pgfpathlineto{\pgfqpoint{0.908935in}{0.609120in}}%
\pgfpathlineto{\pgfqpoint{0.908935in}{0.606171in}}%
\pgfpathmoveto{\pgfqpoint{0.904394in}{0.609120in}}%
\pgfpathlineto{\pgfqpoint{0.904394in}{0.609120in}}%
\pgfpathlineto{\pgfqpoint{0.904394in}{0.612070in}}%
\pgfpathlineto{\pgfqpoint{0.908935in}{0.612070in}}%
\pgfpathlineto{\pgfqpoint{0.908935in}{0.609120in}}%
\pgfpathmoveto{\pgfqpoint{0.908935in}{0.606171in}}%
\pgfpathlineto{\pgfqpoint{0.908935in}{0.606171in}}%
\pgfpathlineto{\pgfqpoint{0.908935in}{0.609120in}}%
\pgfpathlineto{\pgfqpoint{0.913475in}{0.609120in}}%
\pgfpathlineto{\pgfqpoint{0.913475in}{0.606171in}}%
\pgfpathmoveto{\pgfqpoint{0.908935in}{0.609120in}}%
\pgfpathlineto{\pgfqpoint{0.908935in}{0.609120in}}%
\pgfpathlineto{\pgfqpoint{0.908935in}{0.612070in}}%
\pgfpathlineto{\pgfqpoint{0.913475in}{0.612070in}}%
\pgfpathlineto{\pgfqpoint{0.913475in}{0.609120in}}%
\pgfpathmoveto{\pgfqpoint{0.913475in}{0.606171in}}%
\pgfpathlineto{\pgfqpoint{0.913475in}{0.606171in}}%
\pgfpathlineto{\pgfqpoint{0.913475in}{0.609120in}}%
\pgfpathlineto{\pgfqpoint{0.918016in}{0.609120in}}%
\pgfpathlineto{\pgfqpoint{0.918016in}{0.606171in}}%
\pgfpathmoveto{\pgfqpoint{0.913475in}{0.609120in}}%
\pgfpathlineto{\pgfqpoint{0.913475in}{0.609120in}}%
\pgfpathlineto{\pgfqpoint{0.913475in}{0.612070in}}%
\pgfpathlineto{\pgfqpoint{0.918016in}{0.612070in}}%
\pgfpathlineto{\pgfqpoint{0.918016in}{0.609120in}}%
\pgfpathmoveto{\pgfqpoint{0.918016in}{0.606171in}}%
\pgfpathlineto{\pgfqpoint{0.918016in}{0.606171in}}%
\pgfpathlineto{\pgfqpoint{0.918016in}{0.609120in}}%
\pgfpathlineto{\pgfqpoint{0.922557in}{0.609120in}}%
\pgfpathlineto{\pgfqpoint{0.922557in}{0.606171in}}%
\pgfpathmoveto{\pgfqpoint{0.918016in}{0.609120in}}%
\pgfpathlineto{\pgfqpoint{0.918016in}{0.609120in}}%
\pgfpathlineto{\pgfqpoint{0.918016in}{0.612070in}}%
\pgfpathlineto{\pgfqpoint{0.922557in}{0.612070in}}%
\pgfpathlineto{\pgfqpoint{0.922557in}{0.609120in}}%
\pgfpathmoveto{\pgfqpoint{0.922557in}{0.609120in}}%
\pgfpathlineto{\pgfqpoint{0.922557in}{0.609120in}}%
\pgfpathlineto{\pgfqpoint{0.922557in}{0.612070in}}%
\pgfpathlineto{\pgfqpoint{0.927098in}{0.612070in}}%
\pgfpathlineto{\pgfqpoint{0.927098in}{0.609120in}}%
\pgfpathmoveto{\pgfqpoint{0.922557in}{0.612070in}}%
\pgfpathlineto{\pgfqpoint{0.922557in}{0.612070in}}%
\pgfpathlineto{\pgfqpoint{0.922557in}{0.615019in}}%
\pgfpathlineto{\pgfqpoint{0.927098in}{0.615019in}}%
\pgfpathlineto{\pgfqpoint{0.927098in}{0.612070in}}%
\pgfpathmoveto{\pgfqpoint{0.922557in}{0.615019in}}%
\pgfpathlineto{\pgfqpoint{0.922557in}{0.615019in}}%
\pgfpathlineto{\pgfqpoint{0.922557in}{0.617968in}}%
\pgfpathlineto{\pgfqpoint{0.927098in}{0.617968in}}%
\pgfpathlineto{\pgfqpoint{0.927098in}{0.615019in}}%
\pgfpathmoveto{\pgfqpoint{0.927098in}{0.612070in}}%
\pgfpathlineto{\pgfqpoint{0.927098in}{0.612070in}}%
\pgfpathlineto{\pgfqpoint{0.927098in}{0.615019in}}%
\pgfpathlineto{\pgfqpoint{0.931639in}{0.615019in}}%
\pgfpathlineto{\pgfqpoint{0.931639in}{0.612070in}}%
\pgfpathmoveto{\pgfqpoint{0.927098in}{0.615019in}}%
\pgfpathlineto{\pgfqpoint{0.927098in}{0.615019in}}%
\pgfpathlineto{\pgfqpoint{0.927098in}{0.617968in}}%
\pgfpathlineto{\pgfqpoint{0.931639in}{0.617968in}}%
\pgfpathlineto{\pgfqpoint{0.931639in}{0.615019in}}%
\pgfpathmoveto{\pgfqpoint{0.931639in}{0.615019in}}%
\pgfpathlineto{\pgfqpoint{0.931639in}{0.615019in}}%
\pgfpathlineto{\pgfqpoint{0.931639in}{0.617968in}}%
\pgfpathlineto{\pgfqpoint{0.936180in}{0.617968in}}%
\pgfpathlineto{\pgfqpoint{0.936180in}{0.615019in}}%
\pgfpathmoveto{\pgfqpoint{0.931639in}{0.617968in}}%
\pgfpathlineto{\pgfqpoint{0.931639in}{0.617968in}}%
\pgfpathlineto{\pgfqpoint{0.931639in}{0.620917in}}%
\pgfpathlineto{\pgfqpoint{0.936180in}{0.620917in}}%
\pgfpathlineto{\pgfqpoint{0.936180in}{0.617968in}}%
\pgfpathmoveto{\pgfqpoint{0.931639in}{0.620917in}}%
\pgfpathlineto{\pgfqpoint{0.931639in}{0.620917in}}%
\pgfpathlineto{\pgfqpoint{0.931639in}{0.623867in}}%
\pgfpathlineto{\pgfqpoint{0.936180in}{0.623867in}}%
\pgfpathlineto{\pgfqpoint{0.936180in}{0.620917in}}%
\pgfpathmoveto{\pgfqpoint{0.936180in}{0.617968in}}%
\pgfpathlineto{\pgfqpoint{0.936180in}{0.617968in}}%
\pgfpathlineto{\pgfqpoint{0.936180in}{0.620917in}}%
\pgfpathlineto{\pgfqpoint{0.940721in}{0.620917in}}%
\pgfpathlineto{\pgfqpoint{0.940721in}{0.617968in}}%
\pgfpathmoveto{\pgfqpoint{0.936180in}{0.620917in}}%
\pgfpathlineto{\pgfqpoint{0.936180in}{0.620917in}}%
\pgfpathlineto{\pgfqpoint{0.936180in}{0.623867in}}%
\pgfpathlineto{\pgfqpoint{0.940721in}{0.623867in}}%
\pgfpathlineto{\pgfqpoint{0.940721in}{0.620917in}}%
\pgfpathmoveto{\pgfqpoint{0.940721in}{0.620917in}}%
\pgfpathlineto{\pgfqpoint{0.940721in}{0.620917in}}%
\pgfpathlineto{\pgfqpoint{0.940721in}{0.623867in}}%
\pgfpathlineto{\pgfqpoint{0.945262in}{0.623867in}}%
\pgfpathlineto{\pgfqpoint{0.945262in}{0.620917in}}%
\pgfpathmoveto{\pgfqpoint{0.940721in}{0.623867in}}%
\pgfpathlineto{\pgfqpoint{0.940721in}{0.623867in}}%
\pgfpathlineto{\pgfqpoint{0.940721in}{0.626816in}}%
\pgfpathlineto{\pgfqpoint{0.945262in}{0.626816in}}%
\pgfpathlineto{\pgfqpoint{0.945262in}{0.623867in}}%
\pgfpathmoveto{\pgfqpoint{0.940721in}{0.626816in}}%
\pgfpathlineto{\pgfqpoint{0.940721in}{0.626816in}}%
\pgfpathlineto{\pgfqpoint{0.940721in}{0.629765in}}%
\pgfpathlineto{\pgfqpoint{0.945262in}{0.629765in}}%
\pgfpathlineto{\pgfqpoint{0.945262in}{0.626816in}}%
\pgfpathmoveto{\pgfqpoint{0.945262in}{0.623867in}}%
\pgfpathlineto{\pgfqpoint{0.945262in}{0.623867in}}%
\pgfpathlineto{\pgfqpoint{0.945262in}{0.626816in}}%
\pgfpathlineto{\pgfqpoint{0.949803in}{0.626816in}}%
\pgfpathlineto{\pgfqpoint{0.949803in}{0.623867in}}%
\pgfpathmoveto{\pgfqpoint{0.945262in}{0.626816in}}%
\pgfpathlineto{\pgfqpoint{0.945262in}{0.626816in}}%
\pgfpathlineto{\pgfqpoint{0.945262in}{0.629765in}}%
\pgfpathlineto{\pgfqpoint{0.949803in}{0.629765in}}%
\pgfpathlineto{\pgfqpoint{0.949803in}{0.626816in}}%
\pgfpathmoveto{\pgfqpoint{0.949803in}{0.626816in}}%
\pgfpathlineto{\pgfqpoint{0.949803in}{0.626816in}}%
\pgfpathlineto{\pgfqpoint{0.949803in}{0.629765in}}%
\pgfpathlineto{\pgfqpoint{0.954344in}{0.629765in}}%
\pgfpathlineto{\pgfqpoint{0.954344in}{0.626816in}}%
\pgfpathmoveto{\pgfqpoint{0.949803in}{0.629765in}}%
\pgfpathlineto{\pgfqpoint{0.949803in}{0.629765in}}%
\pgfpathlineto{\pgfqpoint{0.949803in}{0.632714in}}%
\pgfpathlineto{\pgfqpoint{0.954344in}{0.632714in}}%
\pgfpathlineto{\pgfqpoint{0.954344in}{0.629765in}}%
\pgfpathmoveto{\pgfqpoint{0.949803in}{0.632714in}}%
\pgfpathlineto{\pgfqpoint{0.949803in}{0.632714in}}%
\pgfpathlineto{\pgfqpoint{0.949803in}{0.635664in}}%
\pgfpathlineto{\pgfqpoint{0.954344in}{0.635664in}}%
\pgfpathlineto{\pgfqpoint{0.954344in}{0.632714in}}%
\pgfpathmoveto{\pgfqpoint{0.954344in}{0.629765in}}%
\pgfpathlineto{\pgfqpoint{0.954344in}{0.629765in}}%
\pgfpathlineto{\pgfqpoint{0.954344in}{0.632714in}}%
\pgfpathlineto{\pgfqpoint{0.958885in}{0.632714in}}%
\pgfpathlineto{\pgfqpoint{0.958885in}{0.629765in}}%
\pgfpathmoveto{\pgfqpoint{0.954344in}{0.632714in}}%
\pgfpathlineto{\pgfqpoint{0.954344in}{0.632714in}}%
\pgfpathlineto{\pgfqpoint{0.954344in}{0.635664in}}%
\pgfpathlineto{\pgfqpoint{0.958885in}{0.635664in}}%
\pgfpathlineto{\pgfqpoint{0.958885in}{0.632714in}}%
\pgfpathmoveto{\pgfqpoint{0.958885in}{0.632714in}}%
\pgfpathlineto{\pgfqpoint{0.958885in}{0.632714in}}%
\pgfpathlineto{\pgfqpoint{0.958885in}{0.635664in}}%
\pgfpathlineto{\pgfqpoint{0.963426in}{0.635664in}}%
\pgfpathlineto{\pgfqpoint{0.963426in}{0.632714in}}%
\pgfpathmoveto{\pgfqpoint{0.958885in}{0.635664in}}%
\pgfpathlineto{\pgfqpoint{0.958885in}{0.635664in}}%
\pgfpathlineto{\pgfqpoint{0.958885in}{0.638613in}}%
\pgfpathlineto{\pgfqpoint{0.963426in}{0.638613in}}%
\pgfpathlineto{\pgfqpoint{0.963426in}{0.635664in}}%
\pgfpathmoveto{\pgfqpoint{0.958885in}{0.638613in}}%
\pgfpathlineto{\pgfqpoint{0.958885in}{0.638613in}}%
\pgfpathlineto{\pgfqpoint{0.958885in}{0.641562in}}%
\pgfpathlineto{\pgfqpoint{0.963426in}{0.641562in}}%
\pgfpathlineto{\pgfqpoint{0.963426in}{0.638613in}}%
\pgfpathmoveto{\pgfqpoint{0.963426in}{0.635664in}}%
\pgfpathlineto{\pgfqpoint{0.963426in}{0.635664in}}%
\pgfpathlineto{\pgfqpoint{0.963426in}{0.638613in}}%
\pgfpathlineto{\pgfqpoint{0.967967in}{0.638613in}}%
\pgfpathlineto{\pgfqpoint{0.967967in}{0.635664in}}%
\pgfpathmoveto{\pgfqpoint{0.963426in}{0.638613in}}%
\pgfpathlineto{\pgfqpoint{0.963426in}{0.638613in}}%
\pgfpathlineto{\pgfqpoint{0.963426in}{0.641562in}}%
\pgfpathlineto{\pgfqpoint{0.967967in}{0.641562in}}%
\pgfpathlineto{\pgfqpoint{0.967967in}{0.638613in}}%
\pgfpathmoveto{\pgfqpoint{0.967967in}{0.638613in}}%
\pgfpathlineto{\pgfqpoint{0.967967in}{0.638613in}}%
\pgfpathlineto{\pgfqpoint{0.967967in}{0.641562in}}%
\pgfpathlineto{\pgfqpoint{0.972508in}{0.641562in}}%
\pgfpathlineto{\pgfqpoint{0.972508in}{0.638613in}}%
\pgfpathmoveto{\pgfqpoint{0.967967in}{0.641562in}}%
\pgfpathlineto{\pgfqpoint{0.967967in}{0.641562in}}%
\pgfpathlineto{\pgfqpoint{0.967967in}{0.644511in}}%
\pgfpathlineto{\pgfqpoint{0.972508in}{0.644511in}}%
\pgfpathlineto{\pgfqpoint{0.972508in}{0.641562in}}%
\pgfpathmoveto{\pgfqpoint{0.967967in}{0.644511in}}%
\pgfpathlineto{\pgfqpoint{0.967967in}{0.644511in}}%
\pgfpathlineto{\pgfqpoint{0.967967in}{0.647461in}}%
\pgfpathlineto{\pgfqpoint{0.972508in}{0.647461in}}%
\pgfpathlineto{\pgfqpoint{0.972508in}{0.644511in}}%
\pgfpathmoveto{\pgfqpoint{0.972508in}{0.641562in}}%
\pgfpathlineto{\pgfqpoint{0.972508in}{0.641562in}}%
\pgfpathlineto{\pgfqpoint{0.972508in}{0.644511in}}%
\pgfpathlineto{\pgfqpoint{0.977049in}{0.644511in}}%
\pgfpathlineto{\pgfqpoint{0.977049in}{0.641562in}}%
\pgfpathmoveto{\pgfqpoint{0.972508in}{0.644511in}}%
\pgfpathlineto{\pgfqpoint{0.972508in}{0.644511in}}%
\pgfpathlineto{\pgfqpoint{0.972508in}{0.647461in}}%
\pgfpathlineto{\pgfqpoint{0.977049in}{0.647461in}}%
\pgfpathlineto{\pgfqpoint{0.977049in}{0.644511in}}%
\pgfpathmoveto{\pgfqpoint{0.977049in}{0.644511in}}%
\pgfpathlineto{\pgfqpoint{0.977049in}{0.644511in}}%
\pgfpathlineto{\pgfqpoint{0.977049in}{0.647461in}}%
\pgfpathlineto{\pgfqpoint{0.981590in}{0.647461in}}%
\pgfpathlineto{\pgfqpoint{0.981590in}{0.644511in}}%
\pgfpathmoveto{\pgfqpoint{0.977049in}{0.647461in}}%
\pgfpathlineto{\pgfqpoint{0.977049in}{0.647461in}}%
\pgfpathlineto{\pgfqpoint{0.977049in}{0.650410in}}%
\pgfpathlineto{\pgfqpoint{0.981590in}{0.650410in}}%
\pgfpathlineto{\pgfqpoint{0.981590in}{0.647461in}}%
\pgfpathmoveto{\pgfqpoint{0.977049in}{0.650410in}}%
\pgfpathlineto{\pgfqpoint{0.977049in}{0.650410in}}%
\pgfpathlineto{\pgfqpoint{0.977049in}{0.653359in}}%
\pgfpathlineto{\pgfqpoint{0.981590in}{0.653359in}}%
\pgfpathlineto{\pgfqpoint{0.981590in}{0.650410in}}%
\pgfpathmoveto{\pgfqpoint{0.981590in}{0.647461in}}%
\pgfpathlineto{\pgfqpoint{0.981590in}{0.647461in}}%
\pgfpathlineto{\pgfqpoint{0.981590in}{0.650410in}}%
\pgfpathlineto{\pgfqpoint{0.986131in}{0.650410in}}%
\pgfpathlineto{\pgfqpoint{0.986131in}{0.647461in}}%
\pgfpathmoveto{\pgfqpoint{0.981590in}{0.650410in}}%
\pgfpathlineto{\pgfqpoint{0.981590in}{0.650410in}}%
\pgfpathlineto{\pgfqpoint{0.981590in}{0.653359in}}%
\pgfpathlineto{\pgfqpoint{0.986131in}{0.653359in}}%
\pgfpathlineto{\pgfqpoint{0.986131in}{0.650410in}}%
\pgfpathmoveto{\pgfqpoint{0.986131in}{0.650410in}}%
\pgfpathlineto{\pgfqpoint{0.986131in}{0.650410in}}%
\pgfpathlineto{\pgfqpoint{0.986131in}{0.653359in}}%
\pgfpathlineto{\pgfqpoint{0.990672in}{0.653359in}}%
\pgfpathlineto{\pgfqpoint{0.990672in}{0.650410in}}%
\pgfpathmoveto{\pgfqpoint{0.986131in}{0.653359in}}%
\pgfpathlineto{\pgfqpoint{0.986131in}{0.653359in}}%
\pgfpathlineto{\pgfqpoint{0.986131in}{0.656308in}}%
\pgfpathlineto{\pgfqpoint{0.990672in}{0.656308in}}%
\pgfpathlineto{\pgfqpoint{0.990672in}{0.653359in}}%
\pgfpathmoveto{\pgfqpoint{0.986131in}{0.656308in}}%
\pgfpathlineto{\pgfqpoint{0.986131in}{0.656308in}}%
\pgfpathlineto{\pgfqpoint{0.986131in}{0.659258in}}%
\pgfpathlineto{\pgfqpoint{0.990672in}{0.659258in}}%
\pgfpathlineto{\pgfqpoint{0.990672in}{0.656308in}}%
\pgfpathmoveto{\pgfqpoint{0.990672in}{0.653359in}}%
\pgfpathlineto{\pgfqpoint{0.990672in}{0.653359in}}%
\pgfpathlineto{\pgfqpoint{0.990672in}{0.656308in}}%
\pgfpathlineto{\pgfqpoint{0.995213in}{0.656308in}}%
\pgfpathlineto{\pgfqpoint{0.995213in}{0.653359in}}%
\pgfpathmoveto{\pgfqpoint{0.990672in}{0.656308in}}%
\pgfpathlineto{\pgfqpoint{0.990672in}{0.656308in}}%
\pgfpathlineto{\pgfqpoint{0.990672in}{0.659258in}}%
\pgfpathlineto{\pgfqpoint{0.995213in}{0.659258in}}%
\pgfpathlineto{\pgfqpoint{0.995213in}{0.656308in}}%
\pgfpathmoveto{\pgfqpoint{0.995213in}{0.656308in}}%
\pgfpathlineto{\pgfqpoint{0.995213in}{0.656308in}}%
\pgfpathlineto{\pgfqpoint{0.995213in}{0.659258in}}%
\pgfpathlineto{\pgfqpoint{0.999753in}{0.659258in}}%
\pgfpathlineto{\pgfqpoint{0.999753in}{0.656308in}}%
\pgfpathmoveto{\pgfqpoint{0.995213in}{0.659258in}}%
\pgfpathlineto{\pgfqpoint{0.995213in}{0.659258in}}%
\pgfpathlineto{\pgfqpoint{0.995213in}{0.662207in}}%
\pgfpathlineto{\pgfqpoint{0.999753in}{0.662207in}}%
\pgfpathlineto{\pgfqpoint{0.999753in}{0.659258in}}%
\pgfpathmoveto{\pgfqpoint{0.995213in}{0.662207in}}%
\pgfpathlineto{\pgfqpoint{0.995213in}{0.662207in}}%
\pgfpathlineto{\pgfqpoint{0.995213in}{0.665156in}}%
\pgfpathlineto{\pgfqpoint{0.999753in}{0.665156in}}%
\pgfpathlineto{\pgfqpoint{0.999753in}{0.662207in}}%
\pgfpathmoveto{\pgfqpoint{0.999753in}{0.659258in}}%
\pgfpathlineto{\pgfqpoint{0.999753in}{0.659258in}}%
\pgfpathlineto{\pgfqpoint{0.999753in}{0.662207in}}%
\pgfpathlineto{\pgfqpoint{1.004294in}{0.662207in}}%
\pgfpathlineto{\pgfqpoint{1.004294in}{0.659258in}}%
\pgfpathmoveto{\pgfqpoint{0.999753in}{0.662207in}}%
\pgfpathlineto{\pgfqpoint{0.999753in}{0.662207in}}%
\pgfpathlineto{\pgfqpoint{0.999753in}{0.665156in}}%
\pgfpathlineto{\pgfqpoint{1.004294in}{0.665156in}}%
\pgfpathlineto{\pgfqpoint{1.004294in}{0.662207in}}%
\pgfpathmoveto{\pgfqpoint{1.004294in}{0.662207in}}%
\pgfpathlineto{\pgfqpoint{1.004294in}{0.662207in}}%
\pgfpathlineto{\pgfqpoint{1.004294in}{0.665156in}}%
\pgfpathlineto{\pgfqpoint{1.008835in}{0.665156in}}%
\pgfpathlineto{\pgfqpoint{1.008835in}{0.662207in}}%
\pgfpathmoveto{\pgfqpoint{1.004294in}{0.665156in}}%
\pgfpathlineto{\pgfqpoint{1.004294in}{0.665156in}}%
\pgfpathlineto{\pgfqpoint{1.004294in}{0.668105in}}%
\pgfpathlineto{\pgfqpoint{1.008835in}{0.668105in}}%
\pgfpathlineto{\pgfqpoint{1.008835in}{0.665156in}}%
\pgfpathmoveto{\pgfqpoint{1.004294in}{0.668105in}}%
\pgfpathlineto{\pgfqpoint{1.004294in}{0.668105in}}%
\pgfpathlineto{\pgfqpoint{1.004294in}{0.671054in}}%
\pgfpathlineto{\pgfqpoint{1.008835in}{0.671054in}}%
\pgfpathlineto{\pgfqpoint{1.008835in}{0.668105in}}%
\pgfpathmoveto{\pgfqpoint{1.008835in}{0.665156in}}%
\pgfpathlineto{\pgfqpoint{1.008835in}{0.665156in}}%
\pgfpathlineto{\pgfqpoint{1.008835in}{0.668105in}}%
\pgfpathlineto{\pgfqpoint{1.013376in}{0.668105in}}%
\pgfpathlineto{\pgfqpoint{1.013376in}{0.665156in}}%
\pgfpathmoveto{\pgfqpoint{1.008835in}{0.668105in}}%
\pgfpathlineto{\pgfqpoint{1.008835in}{0.668105in}}%
\pgfpathlineto{\pgfqpoint{1.008835in}{0.671054in}}%
\pgfpathlineto{\pgfqpoint{1.013376in}{0.671054in}}%
\pgfpathlineto{\pgfqpoint{1.013376in}{0.668105in}}%
\pgfpathmoveto{\pgfqpoint{1.013376in}{0.668105in}}%
\pgfpathlineto{\pgfqpoint{1.013376in}{0.668105in}}%
\pgfpathlineto{\pgfqpoint{1.013376in}{0.671054in}}%
\pgfpathlineto{\pgfqpoint{1.017917in}{0.671054in}}%
\pgfpathlineto{\pgfqpoint{1.017917in}{0.668105in}}%
\pgfpathmoveto{\pgfqpoint{1.013376in}{0.671054in}}%
\pgfpathlineto{\pgfqpoint{1.013376in}{0.671054in}}%
\pgfpathlineto{\pgfqpoint{1.013376in}{0.674004in}}%
\pgfpathlineto{\pgfqpoint{1.017917in}{0.674004in}}%
\pgfpathlineto{\pgfqpoint{1.017917in}{0.671054in}}%
\pgfpathmoveto{\pgfqpoint{1.013376in}{0.674004in}}%
\pgfpathlineto{\pgfqpoint{1.013376in}{0.674004in}}%
\pgfpathlineto{\pgfqpoint{1.013376in}{0.676953in}}%
\pgfpathlineto{\pgfqpoint{1.017917in}{0.676953in}}%
\pgfpathlineto{\pgfqpoint{1.017917in}{0.674004in}}%
\pgfpathmoveto{\pgfqpoint{1.017917in}{0.671054in}}%
\pgfpathlineto{\pgfqpoint{1.017917in}{0.671054in}}%
\pgfpathlineto{\pgfqpoint{1.017917in}{0.674004in}}%
\pgfpathlineto{\pgfqpoint{1.022458in}{0.674004in}}%
\pgfpathlineto{\pgfqpoint{1.022458in}{0.671054in}}%
\pgfpathmoveto{\pgfqpoint{1.017917in}{0.674004in}}%
\pgfpathlineto{\pgfqpoint{1.017917in}{0.674004in}}%
\pgfpathlineto{\pgfqpoint{1.017917in}{0.676953in}}%
\pgfpathlineto{\pgfqpoint{1.022458in}{0.676953in}}%
\pgfpathlineto{\pgfqpoint{1.022458in}{0.674004in}}%
\pgfpathmoveto{\pgfqpoint{1.022458in}{0.674004in}}%
\pgfpathlineto{\pgfqpoint{1.022458in}{0.674004in}}%
\pgfpathlineto{\pgfqpoint{1.022458in}{0.676953in}}%
\pgfpathlineto{\pgfqpoint{1.026999in}{0.676953in}}%
\pgfpathlineto{\pgfqpoint{1.026999in}{0.674004in}}%
\pgfpathmoveto{\pgfqpoint{1.022458in}{0.676953in}}%
\pgfpathlineto{\pgfqpoint{1.022458in}{0.676953in}}%
\pgfpathlineto{\pgfqpoint{1.022458in}{0.679902in}}%
\pgfpathlineto{\pgfqpoint{1.026999in}{0.679902in}}%
\pgfpathlineto{\pgfqpoint{1.026999in}{0.676953in}}%
\pgfpathmoveto{\pgfqpoint{1.022458in}{0.679902in}}%
\pgfpathlineto{\pgfqpoint{1.022458in}{0.679902in}}%
\pgfpathlineto{\pgfqpoint{1.022458in}{0.682851in}}%
\pgfpathlineto{\pgfqpoint{1.026999in}{0.682851in}}%
\pgfpathlineto{\pgfqpoint{1.026999in}{0.679902in}}%
\pgfpathmoveto{\pgfqpoint{1.026999in}{0.676953in}}%
\pgfpathlineto{\pgfqpoint{1.026999in}{0.676953in}}%
\pgfpathlineto{\pgfqpoint{1.026999in}{0.679902in}}%
\pgfpathlineto{\pgfqpoint{1.031540in}{0.679902in}}%
\pgfpathlineto{\pgfqpoint{1.031540in}{0.676953in}}%
\pgfpathmoveto{\pgfqpoint{1.026999in}{0.679902in}}%
\pgfpathlineto{\pgfqpoint{1.026999in}{0.679902in}}%
\pgfpathlineto{\pgfqpoint{1.026999in}{0.682851in}}%
\pgfpathlineto{\pgfqpoint{1.031540in}{0.682851in}}%
\pgfpathlineto{\pgfqpoint{1.031540in}{0.679902in}}%
\pgfpathmoveto{\pgfqpoint{1.031540in}{0.679902in}}%
\pgfpathlineto{\pgfqpoint{1.031540in}{0.679902in}}%
\pgfpathlineto{\pgfqpoint{1.031540in}{0.682851in}}%
\pgfpathlineto{\pgfqpoint{1.036081in}{0.682851in}}%
\pgfpathlineto{\pgfqpoint{1.036081in}{0.679902in}}%
\pgfpathmoveto{\pgfqpoint{1.031540in}{0.682851in}}%
\pgfpathlineto{\pgfqpoint{1.031540in}{0.682851in}}%
\pgfpathlineto{\pgfqpoint{1.031540in}{0.685801in}}%
\pgfpathlineto{\pgfqpoint{1.036081in}{0.685801in}}%
\pgfpathlineto{\pgfqpoint{1.036081in}{0.682851in}}%
\pgfpathmoveto{\pgfqpoint{1.031540in}{0.685801in}}%
\pgfpathlineto{\pgfqpoint{1.031540in}{0.685801in}}%
\pgfpathlineto{\pgfqpoint{1.031540in}{0.688750in}}%
\pgfpathlineto{\pgfqpoint{1.036081in}{0.688750in}}%
\pgfpathlineto{\pgfqpoint{1.036081in}{0.685801in}}%
\pgfpathmoveto{\pgfqpoint{1.036081in}{0.682851in}}%
\pgfpathlineto{\pgfqpoint{1.036081in}{0.682851in}}%
\pgfpathlineto{\pgfqpoint{1.036081in}{0.685801in}}%
\pgfpathlineto{\pgfqpoint{1.040622in}{0.685801in}}%
\pgfpathlineto{\pgfqpoint{1.040622in}{0.682851in}}%
\pgfpathmoveto{\pgfqpoint{1.036081in}{0.685801in}}%
\pgfpathlineto{\pgfqpoint{1.036081in}{0.685801in}}%
\pgfpathlineto{\pgfqpoint{1.036081in}{0.688750in}}%
\pgfpathlineto{\pgfqpoint{1.040622in}{0.688750in}}%
\pgfpathlineto{\pgfqpoint{1.040622in}{0.685801in}}%
\pgfpathmoveto{\pgfqpoint{1.040622in}{0.685801in}}%
\pgfpathlineto{\pgfqpoint{1.040622in}{0.685801in}}%
\pgfpathlineto{\pgfqpoint{1.040622in}{0.688750in}}%
\pgfpathlineto{\pgfqpoint{1.045163in}{0.688750in}}%
\pgfpathlineto{\pgfqpoint{1.045163in}{0.685801in}}%
\pgfpathmoveto{\pgfqpoint{1.040622in}{0.688750in}}%
\pgfpathlineto{\pgfqpoint{1.040622in}{0.688750in}}%
\pgfpathlineto{\pgfqpoint{1.040622in}{0.691699in}}%
\pgfpathlineto{\pgfqpoint{1.045163in}{0.691699in}}%
\pgfpathlineto{\pgfqpoint{1.045163in}{0.688750in}}%
\pgfpathmoveto{\pgfqpoint{1.040622in}{0.691699in}}%
\pgfpathlineto{\pgfqpoint{1.040622in}{0.691699in}}%
\pgfpathlineto{\pgfqpoint{1.040622in}{0.694648in}}%
\pgfpathlineto{\pgfqpoint{1.045163in}{0.694648in}}%
\pgfpathlineto{\pgfqpoint{1.045163in}{0.691699in}}%
\pgfpathmoveto{\pgfqpoint{1.045163in}{0.688750in}}%
\pgfpathlineto{\pgfqpoint{1.045163in}{0.688750in}}%
\pgfpathlineto{\pgfqpoint{1.045163in}{0.691699in}}%
\pgfpathlineto{\pgfqpoint{1.049704in}{0.691699in}}%
\pgfpathlineto{\pgfqpoint{1.049704in}{0.688750in}}%
\pgfpathmoveto{\pgfqpoint{1.045163in}{0.691699in}}%
\pgfpathlineto{\pgfqpoint{1.045163in}{0.691699in}}%
\pgfpathlineto{\pgfqpoint{1.045163in}{0.694648in}}%
\pgfpathlineto{\pgfqpoint{1.049704in}{0.694648in}}%
\pgfpathlineto{\pgfqpoint{1.049704in}{0.691699in}}%
\pgfpathmoveto{\pgfqpoint{1.049704in}{0.691699in}}%
\pgfpathlineto{\pgfqpoint{1.049704in}{0.691699in}}%
\pgfpathlineto{\pgfqpoint{1.049704in}{0.694648in}}%
\pgfpathlineto{\pgfqpoint{1.054246in}{0.694648in}}%
\pgfpathlineto{\pgfqpoint{1.054246in}{0.691699in}}%
\pgfpathmoveto{\pgfqpoint{1.049704in}{0.694648in}}%
\pgfpathlineto{\pgfqpoint{1.049704in}{0.694648in}}%
\pgfpathlineto{\pgfqpoint{1.049704in}{0.697598in}}%
\pgfpathlineto{\pgfqpoint{1.054246in}{0.697598in}}%
\pgfpathlineto{\pgfqpoint{1.054246in}{0.694648in}}%
\pgfpathmoveto{\pgfqpoint{1.049704in}{0.697598in}}%
\pgfpathlineto{\pgfqpoint{1.049704in}{0.697598in}}%
\pgfpathlineto{\pgfqpoint{1.049704in}{0.700547in}}%
\pgfpathlineto{\pgfqpoint{1.054246in}{0.700547in}}%
\pgfpathlineto{\pgfqpoint{1.054246in}{0.697598in}}%
\pgfpathmoveto{\pgfqpoint{1.054246in}{0.694648in}}%
\pgfpathlineto{\pgfqpoint{1.054246in}{0.694648in}}%
\pgfpathlineto{\pgfqpoint{1.054246in}{0.697598in}}%
\pgfpathlineto{\pgfqpoint{1.058787in}{0.697598in}}%
\pgfpathlineto{\pgfqpoint{1.058787in}{0.694648in}}%
\pgfpathmoveto{\pgfqpoint{1.054246in}{0.697598in}}%
\pgfpathlineto{\pgfqpoint{1.054246in}{0.697598in}}%
\pgfpathlineto{\pgfqpoint{1.054246in}{0.700547in}}%
\pgfpathlineto{\pgfqpoint{1.058787in}{0.700547in}}%
\pgfpathlineto{\pgfqpoint{1.058787in}{0.697598in}}%
\pgfpathmoveto{\pgfqpoint{1.058787in}{0.697598in}}%
\pgfpathlineto{\pgfqpoint{1.058787in}{0.697598in}}%
\pgfpathlineto{\pgfqpoint{1.058787in}{0.700547in}}%
\pgfpathlineto{\pgfqpoint{1.063328in}{0.700547in}}%
\pgfpathlineto{\pgfqpoint{1.063328in}{0.697598in}}%
\pgfpathmoveto{\pgfqpoint{1.058787in}{0.700547in}}%
\pgfpathlineto{\pgfqpoint{1.058787in}{0.700547in}}%
\pgfpathlineto{\pgfqpoint{1.058787in}{0.703496in}}%
\pgfpathlineto{\pgfqpoint{1.063328in}{0.703496in}}%
\pgfpathlineto{\pgfqpoint{1.063328in}{0.700547in}}%
\pgfpathmoveto{\pgfqpoint{1.058787in}{0.703496in}}%
\pgfpathlineto{\pgfqpoint{1.058787in}{0.703496in}}%
\pgfpathlineto{\pgfqpoint{1.058787in}{0.706445in}}%
\pgfpathlineto{\pgfqpoint{1.063328in}{0.706445in}}%
\pgfpathlineto{\pgfqpoint{1.063328in}{0.703496in}}%
\pgfpathmoveto{\pgfqpoint{1.063328in}{0.700547in}}%
\pgfpathlineto{\pgfqpoint{1.063328in}{0.700547in}}%
\pgfpathlineto{\pgfqpoint{1.063328in}{0.703496in}}%
\pgfpathlineto{\pgfqpoint{1.067869in}{0.703496in}}%
\pgfpathlineto{\pgfqpoint{1.067869in}{0.700547in}}%
\pgfpathmoveto{\pgfqpoint{1.063328in}{0.703496in}}%
\pgfpathlineto{\pgfqpoint{1.063328in}{0.703496in}}%
\pgfpathlineto{\pgfqpoint{1.063328in}{0.706445in}}%
\pgfpathlineto{\pgfqpoint{1.067869in}{0.706445in}}%
\pgfpathlineto{\pgfqpoint{1.067869in}{0.703496in}}%
\pgfpathmoveto{\pgfqpoint{1.067869in}{0.703496in}}%
\pgfpathlineto{\pgfqpoint{1.067869in}{0.703496in}}%
\pgfpathlineto{\pgfqpoint{1.067869in}{0.706445in}}%
\pgfpathlineto{\pgfqpoint{1.072411in}{0.706445in}}%
\pgfpathlineto{\pgfqpoint{1.072411in}{0.703496in}}%
\pgfpathmoveto{\pgfqpoint{1.067869in}{0.706445in}}%
\pgfpathlineto{\pgfqpoint{1.067869in}{0.706445in}}%
\pgfpathlineto{\pgfqpoint{1.067869in}{0.709394in}}%
\pgfpathlineto{\pgfqpoint{1.072411in}{0.709394in}}%
\pgfpathlineto{\pgfqpoint{1.072411in}{0.706445in}}%
\pgfpathmoveto{\pgfqpoint{1.067869in}{0.709394in}}%
\pgfpathlineto{\pgfqpoint{1.067869in}{0.709394in}}%
\pgfpathlineto{\pgfqpoint{1.067869in}{0.712343in}}%
\pgfpathlineto{\pgfqpoint{1.072411in}{0.712343in}}%
\pgfpathlineto{\pgfqpoint{1.072411in}{0.709394in}}%
\pgfpathmoveto{\pgfqpoint{1.072411in}{0.706445in}}%
\pgfpathlineto{\pgfqpoint{1.072411in}{0.706445in}}%
\pgfpathlineto{\pgfqpoint{1.072411in}{0.709394in}}%
\pgfpathlineto{\pgfqpoint{1.076952in}{0.709394in}}%
\pgfpathlineto{\pgfqpoint{1.076952in}{0.706445in}}%
\pgfpathmoveto{\pgfqpoint{1.072411in}{0.709394in}}%
\pgfpathlineto{\pgfqpoint{1.072411in}{0.709394in}}%
\pgfpathlineto{\pgfqpoint{1.072411in}{0.712343in}}%
\pgfpathlineto{\pgfqpoint{1.076952in}{0.712343in}}%
\pgfpathlineto{\pgfqpoint{1.076952in}{0.709394in}}%
\pgfpathmoveto{\pgfqpoint{1.076952in}{0.709394in}}%
\pgfpathlineto{\pgfqpoint{1.076952in}{0.709394in}}%
\pgfpathlineto{\pgfqpoint{1.076952in}{0.712343in}}%
\pgfpathlineto{\pgfqpoint{1.081493in}{0.712343in}}%
\pgfpathlineto{\pgfqpoint{1.081493in}{0.709394in}}%
\pgfpathmoveto{\pgfqpoint{1.076952in}{0.712343in}}%
\pgfpathlineto{\pgfqpoint{1.076952in}{0.712343in}}%
\pgfpathlineto{\pgfqpoint{1.076952in}{0.715293in}}%
\pgfpathlineto{\pgfqpoint{1.081493in}{0.715293in}}%
\pgfpathlineto{\pgfqpoint{1.081493in}{0.712343in}}%
\pgfpathmoveto{\pgfqpoint{1.076952in}{0.715293in}}%
\pgfpathlineto{\pgfqpoint{1.076952in}{0.715293in}}%
\pgfpathlineto{\pgfqpoint{1.076952in}{0.718242in}}%
\pgfpathlineto{\pgfqpoint{1.081493in}{0.718242in}}%
\pgfpathlineto{\pgfqpoint{1.081493in}{0.715293in}}%
\pgfpathmoveto{\pgfqpoint{1.081493in}{0.712343in}}%
\pgfpathlineto{\pgfqpoint{1.081493in}{0.712343in}}%
\pgfpathlineto{\pgfqpoint{1.081493in}{0.715293in}}%
\pgfpathlineto{\pgfqpoint{1.086034in}{0.715293in}}%
\pgfpathlineto{\pgfqpoint{1.086034in}{0.712343in}}%
\pgfpathmoveto{\pgfqpoint{1.081493in}{0.715293in}}%
\pgfpathlineto{\pgfqpoint{1.081493in}{0.715293in}}%
\pgfpathlineto{\pgfqpoint{1.081493in}{0.718242in}}%
\pgfpathlineto{\pgfqpoint{1.086034in}{0.718242in}}%
\pgfpathlineto{\pgfqpoint{1.086034in}{0.715293in}}%
\pgfpathmoveto{\pgfqpoint{1.086034in}{0.715293in}}%
\pgfpathlineto{\pgfqpoint{1.086034in}{0.715293in}}%
\pgfpathlineto{\pgfqpoint{1.086034in}{0.718242in}}%
\pgfpathlineto{\pgfqpoint{1.090575in}{0.718242in}}%
\pgfpathlineto{\pgfqpoint{1.090575in}{0.715293in}}%
\pgfpathmoveto{\pgfqpoint{1.086034in}{0.718242in}}%
\pgfpathlineto{\pgfqpoint{1.086034in}{0.718242in}}%
\pgfpathlineto{\pgfqpoint{1.086034in}{0.721191in}}%
\pgfpathlineto{\pgfqpoint{1.090575in}{0.721191in}}%
\pgfpathlineto{\pgfqpoint{1.090575in}{0.718242in}}%
\pgfpathmoveto{\pgfqpoint{1.086034in}{0.721191in}}%
\pgfpathlineto{\pgfqpoint{1.086034in}{0.721191in}}%
\pgfpathlineto{\pgfqpoint{1.086034in}{0.724140in}}%
\pgfpathlineto{\pgfqpoint{1.090575in}{0.724140in}}%
\pgfpathlineto{\pgfqpoint{1.090575in}{0.721191in}}%
\pgfpathmoveto{\pgfqpoint{1.090575in}{0.718242in}}%
\pgfpathlineto{\pgfqpoint{1.090575in}{0.718242in}}%
\pgfpathlineto{\pgfqpoint{1.090575in}{0.721191in}}%
\pgfpathlineto{\pgfqpoint{1.095117in}{0.721191in}}%
\pgfpathlineto{\pgfqpoint{1.095117in}{0.718242in}}%
\pgfpathmoveto{\pgfqpoint{1.090575in}{0.721191in}}%
\pgfpathlineto{\pgfqpoint{1.090575in}{0.721191in}}%
\pgfpathlineto{\pgfqpoint{1.090575in}{0.724140in}}%
\pgfpathlineto{\pgfqpoint{1.095117in}{0.724140in}}%
\pgfpathlineto{\pgfqpoint{1.095117in}{0.721191in}}%
\pgfpathmoveto{\pgfqpoint{1.086034in}{0.724140in}}%
\pgfpathlineto{\pgfqpoint{1.086034in}{0.724140in}}%
\pgfpathlineto{\pgfqpoint{1.086034in}{0.727089in}}%
\pgfpathlineto{\pgfqpoint{1.090575in}{0.727089in}}%
\pgfpathlineto{\pgfqpoint{1.090575in}{0.724140in}}%
\pgfpathmoveto{\pgfqpoint{1.086034in}{0.727089in}}%
\pgfpathlineto{\pgfqpoint{1.086034in}{0.727089in}}%
\pgfpathlineto{\pgfqpoint{1.086034in}{0.730039in}}%
\pgfpathlineto{\pgfqpoint{1.090575in}{0.730039in}}%
\pgfpathlineto{\pgfqpoint{1.090575in}{0.727089in}}%
\pgfpathmoveto{\pgfqpoint{1.090575in}{0.724140in}}%
\pgfpathlineto{\pgfqpoint{1.090575in}{0.724140in}}%
\pgfpathlineto{\pgfqpoint{1.090575in}{0.727089in}}%
\pgfpathlineto{\pgfqpoint{1.095117in}{0.727089in}}%
\pgfpathlineto{\pgfqpoint{1.095117in}{0.724140in}}%
\pgfpathmoveto{\pgfqpoint{1.090575in}{0.727089in}}%
\pgfpathlineto{\pgfqpoint{1.090575in}{0.727089in}}%
\pgfpathlineto{\pgfqpoint{1.090575in}{0.730039in}}%
\pgfpathlineto{\pgfqpoint{1.095117in}{0.730039in}}%
\pgfpathlineto{\pgfqpoint{1.095117in}{0.727089in}}%
\pgfpathmoveto{\pgfqpoint{1.095117in}{0.724140in}}%
\pgfpathlineto{\pgfqpoint{1.095117in}{0.724140in}}%
\pgfpathlineto{\pgfqpoint{1.095117in}{0.727089in}}%
\pgfpathlineto{\pgfqpoint{1.099658in}{0.727089in}}%
\pgfpathlineto{\pgfqpoint{1.099658in}{0.724140in}}%
\pgfpathmoveto{\pgfqpoint{1.095117in}{0.727089in}}%
\pgfpathlineto{\pgfqpoint{1.095117in}{0.727089in}}%
\pgfpathlineto{\pgfqpoint{1.095117in}{0.730039in}}%
\pgfpathlineto{\pgfqpoint{1.099658in}{0.730039in}}%
\pgfpathlineto{\pgfqpoint{1.099658in}{0.727089in}}%
\pgfpathmoveto{\pgfqpoint{1.099658in}{0.727089in}}%
\pgfpathlineto{\pgfqpoint{1.099658in}{0.727089in}}%
\pgfpathlineto{\pgfqpoint{1.099658in}{0.730039in}}%
\pgfpathlineto{\pgfqpoint{1.104199in}{0.730039in}}%
\pgfpathlineto{\pgfqpoint{1.104199in}{0.727089in}}%
\pgfpathmoveto{\pgfqpoint{1.095117in}{0.730039in}}%
\pgfpathlineto{\pgfqpoint{1.095117in}{0.730039in}}%
\pgfpathlineto{\pgfqpoint{1.095117in}{0.732988in}}%
\pgfpathlineto{\pgfqpoint{1.099658in}{0.732988in}}%
\pgfpathlineto{\pgfqpoint{1.099658in}{0.730039in}}%
\pgfpathmoveto{\pgfqpoint{1.095117in}{0.732988in}}%
\pgfpathlineto{\pgfqpoint{1.095117in}{0.732988in}}%
\pgfpathlineto{\pgfqpoint{1.095117in}{0.735937in}}%
\pgfpathlineto{\pgfqpoint{1.099658in}{0.735937in}}%
\pgfpathlineto{\pgfqpoint{1.099658in}{0.732988in}}%
\pgfpathmoveto{\pgfqpoint{1.099658in}{0.730039in}}%
\pgfpathlineto{\pgfqpoint{1.099658in}{0.730039in}}%
\pgfpathlineto{\pgfqpoint{1.099658in}{0.732988in}}%
\pgfpathlineto{\pgfqpoint{1.104199in}{0.732988in}}%
\pgfpathlineto{\pgfqpoint{1.104199in}{0.730039in}}%
\pgfpathmoveto{\pgfqpoint{1.099658in}{0.732988in}}%
\pgfpathlineto{\pgfqpoint{1.099658in}{0.732988in}}%
\pgfpathlineto{\pgfqpoint{1.099658in}{0.735937in}}%
\pgfpathlineto{\pgfqpoint{1.104199in}{0.735937in}}%
\pgfpathlineto{\pgfqpoint{1.104199in}{0.732988in}}%
\pgfpathmoveto{\pgfqpoint{1.104199in}{0.730039in}}%
\pgfpathlineto{\pgfqpoint{1.104199in}{0.730039in}}%
\pgfpathlineto{\pgfqpoint{1.104199in}{0.732988in}}%
\pgfpathlineto{\pgfqpoint{1.108740in}{0.732988in}}%
\pgfpathlineto{\pgfqpoint{1.108740in}{0.730039in}}%
\pgfpathmoveto{\pgfqpoint{1.104199in}{0.732988in}}%
\pgfpathlineto{\pgfqpoint{1.104199in}{0.732988in}}%
\pgfpathlineto{\pgfqpoint{1.104199in}{0.735937in}}%
\pgfpathlineto{\pgfqpoint{1.108740in}{0.735937in}}%
\pgfpathlineto{\pgfqpoint{1.108740in}{0.732988in}}%
\pgfpathmoveto{\pgfqpoint{1.108740in}{0.732988in}}%
\pgfpathlineto{\pgfqpoint{1.108740in}{0.732988in}}%
\pgfpathlineto{\pgfqpoint{1.108740in}{0.735937in}}%
\pgfpathlineto{\pgfqpoint{1.113282in}{0.735937in}}%
\pgfpathlineto{\pgfqpoint{1.113282in}{0.732988in}}%
\pgfpathmoveto{\pgfqpoint{1.104199in}{0.735937in}}%
\pgfpathlineto{\pgfqpoint{1.104199in}{0.735937in}}%
\pgfpathlineto{\pgfqpoint{1.104199in}{0.738886in}}%
\pgfpathlineto{\pgfqpoint{1.108740in}{0.738886in}}%
\pgfpathlineto{\pgfqpoint{1.108740in}{0.735937in}}%
\pgfpathmoveto{\pgfqpoint{1.104199in}{0.738886in}}%
\pgfpathlineto{\pgfqpoint{1.104199in}{0.738886in}}%
\pgfpathlineto{\pgfqpoint{1.104199in}{0.741835in}}%
\pgfpathlineto{\pgfqpoint{1.108740in}{0.741835in}}%
\pgfpathlineto{\pgfqpoint{1.108740in}{0.738886in}}%
\pgfpathmoveto{\pgfqpoint{1.108740in}{0.735937in}}%
\pgfpathlineto{\pgfqpoint{1.108740in}{0.735937in}}%
\pgfpathlineto{\pgfqpoint{1.108740in}{0.738886in}}%
\pgfpathlineto{\pgfqpoint{1.113282in}{0.738886in}}%
\pgfpathlineto{\pgfqpoint{1.113282in}{0.735937in}}%
\pgfpathmoveto{\pgfqpoint{1.108740in}{0.738886in}}%
\pgfpathlineto{\pgfqpoint{1.108740in}{0.738886in}}%
\pgfpathlineto{\pgfqpoint{1.108740in}{0.741835in}}%
\pgfpathlineto{\pgfqpoint{1.113282in}{0.741835in}}%
\pgfpathlineto{\pgfqpoint{1.113282in}{0.738886in}}%
\pgfpathmoveto{\pgfqpoint{1.113282in}{0.735937in}}%
\pgfpathlineto{\pgfqpoint{1.113282in}{0.735937in}}%
\pgfpathlineto{\pgfqpoint{1.113282in}{0.738886in}}%
\pgfpathlineto{\pgfqpoint{1.117823in}{0.738886in}}%
\pgfpathlineto{\pgfqpoint{1.117823in}{0.735937in}}%
\pgfpathmoveto{\pgfqpoint{1.113282in}{0.738886in}}%
\pgfpathlineto{\pgfqpoint{1.113282in}{0.738886in}}%
\pgfpathlineto{\pgfqpoint{1.113282in}{0.741835in}}%
\pgfpathlineto{\pgfqpoint{1.117823in}{0.741835in}}%
\pgfpathlineto{\pgfqpoint{1.117823in}{0.738886in}}%
\pgfpathmoveto{\pgfqpoint{1.117823in}{0.738886in}}%
\pgfpathlineto{\pgfqpoint{1.117823in}{0.738886in}}%
\pgfpathlineto{\pgfqpoint{1.117823in}{0.741835in}}%
\pgfpathlineto{\pgfqpoint{1.122364in}{0.741835in}}%
\pgfpathlineto{\pgfqpoint{1.122364in}{0.738886in}}%
\pgfpathmoveto{\pgfqpoint{1.113282in}{0.741835in}}%
\pgfpathlineto{\pgfqpoint{1.113282in}{0.741835in}}%
\pgfpathlineto{\pgfqpoint{1.113282in}{0.744784in}}%
\pgfpathlineto{\pgfqpoint{1.117823in}{0.744784in}}%
\pgfpathlineto{\pgfqpoint{1.117823in}{0.741835in}}%
\pgfpathmoveto{\pgfqpoint{1.113282in}{0.744784in}}%
\pgfpathlineto{\pgfqpoint{1.113282in}{0.744784in}}%
\pgfpathlineto{\pgfqpoint{1.113282in}{0.747734in}}%
\pgfpathlineto{\pgfqpoint{1.117823in}{0.747734in}}%
\pgfpathlineto{\pgfqpoint{1.117823in}{0.744784in}}%
\pgfpathmoveto{\pgfqpoint{1.117823in}{0.741835in}}%
\pgfpathlineto{\pgfqpoint{1.117823in}{0.741835in}}%
\pgfpathlineto{\pgfqpoint{1.117823in}{0.744784in}}%
\pgfpathlineto{\pgfqpoint{1.122364in}{0.744784in}}%
\pgfpathlineto{\pgfqpoint{1.122364in}{0.741835in}}%
\pgfpathmoveto{\pgfqpoint{1.117823in}{0.744784in}}%
\pgfpathlineto{\pgfqpoint{1.117823in}{0.744784in}}%
\pgfpathlineto{\pgfqpoint{1.117823in}{0.747734in}}%
\pgfpathlineto{\pgfqpoint{1.122364in}{0.747734in}}%
\pgfpathlineto{\pgfqpoint{1.122364in}{0.744784in}}%
\pgfpathmoveto{\pgfqpoint{1.122364in}{0.741835in}}%
\pgfpathlineto{\pgfqpoint{1.122364in}{0.741835in}}%
\pgfpathlineto{\pgfqpoint{1.122364in}{0.744784in}}%
\pgfpathlineto{\pgfqpoint{1.126905in}{0.744784in}}%
\pgfpathlineto{\pgfqpoint{1.126905in}{0.741835in}}%
\pgfpathmoveto{\pgfqpoint{1.122364in}{0.744784in}}%
\pgfpathlineto{\pgfqpoint{1.122364in}{0.744784in}}%
\pgfpathlineto{\pgfqpoint{1.122364in}{0.747734in}}%
\pgfpathlineto{\pgfqpoint{1.126905in}{0.747734in}}%
\pgfpathlineto{\pgfqpoint{1.126905in}{0.744784in}}%
\pgfpathmoveto{\pgfqpoint{1.126905in}{0.744784in}}%
\pgfpathlineto{\pgfqpoint{1.126905in}{0.744784in}}%
\pgfpathlineto{\pgfqpoint{1.126905in}{0.747734in}}%
\pgfpathlineto{\pgfqpoint{1.131447in}{0.747734in}}%
\pgfpathlineto{\pgfqpoint{1.131447in}{0.744784in}}%
\pgfpathmoveto{\pgfqpoint{1.122364in}{0.747734in}}%
\pgfpathlineto{\pgfqpoint{1.122364in}{0.747734in}}%
\pgfpathlineto{\pgfqpoint{1.122364in}{0.750683in}}%
\pgfpathlineto{\pgfqpoint{1.126905in}{0.750683in}}%
\pgfpathlineto{\pgfqpoint{1.126905in}{0.747734in}}%
\pgfpathmoveto{\pgfqpoint{1.122364in}{0.750683in}}%
\pgfpathlineto{\pgfqpoint{1.122364in}{0.750683in}}%
\pgfpathlineto{\pgfqpoint{1.122364in}{0.753632in}}%
\pgfpathlineto{\pgfqpoint{1.126905in}{0.753632in}}%
\pgfpathlineto{\pgfqpoint{1.126905in}{0.750683in}}%
\pgfpathmoveto{\pgfqpoint{1.126905in}{0.747734in}}%
\pgfpathlineto{\pgfqpoint{1.126905in}{0.747734in}}%
\pgfpathlineto{\pgfqpoint{1.126905in}{0.750683in}}%
\pgfpathlineto{\pgfqpoint{1.131447in}{0.750683in}}%
\pgfpathlineto{\pgfqpoint{1.131447in}{0.747734in}}%
\pgfpathmoveto{\pgfqpoint{1.126905in}{0.750683in}}%
\pgfpathlineto{\pgfqpoint{1.126905in}{0.750683in}}%
\pgfpathlineto{\pgfqpoint{1.126905in}{0.753632in}}%
\pgfpathlineto{\pgfqpoint{1.131447in}{0.753632in}}%
\pgfpathlineto{\pgfqpoint{1.131447in}{0.750683in}}%
\pgfpathmoveto{\pgfqpoint{1.131447in}{0.747734in}}%
\pgfpathlineto{\pgfqpoint{1.131447in}{0.747734in}}%
\pgfpathlineto{\pgfqpoint{1.131447in}{0.750683in}}%
\pgfpathlineto{\pgfqpoint{1.135988in}{0.750683in}}%
\pgfpathlineto{\pgfqpoint{1.135988in}{0.747734in}}%
\pgfpathmoveto{\pgfqpoint{1.131447in}{0.750683in}}%
\pgfpathlineto{\pgfqpoint{1.131447in}{0.750683in}}%
\pgfpathlineto{\pgfqpoint{1.131447in}{0.753632in}}%
\pgfpathlineto{\pgfqpoint{1.135988in}{0.753632in}}%
\pgfpathlineto{\pgfqpoint{1.135988in}{0.750683in}}%
\pgfpathmoveto{\pgfqpoint{1.135988in}{0.750683in}}%
\pgfpathlineto{\pgfqpoint{1.135988in}{0.750683in}}%
\pgfpathlineto{\pgfqpoint{1.135988in}{0.753632in}}%
\pgfpathlineto{\pgfqpoint{1.140529in}{0.753632in}}%
\pgfpathlineto{\pgfqpoint{1.140529in}{0.750683in}}%
\pgfpathmoveto{\pgfqpoint{1.131447in}{0.753632in}}%
\pgfpathlineto{\pgfqpoint{1.131447in}{0.753632in}}%
\pgfpathlineto{\pgfqpoint{1.131447in}{0.756581in}}%
\pgfpathlineto{\pgfqpoint{1.135988in}{0.756581in}}%
\pgfpathlineto{\pgfqpoint{1.135988in}{0.753632in}}%
\pgfpathmoveto{\pgfqpoint{1.131447in}{0.756581in}}%
\pgfpathlineto{\pgfqpoint{1.131447in}{0.756581in}}%
\pgfpathlineto{\pgfqpoint{1.131447in}{0.759530in}}%
\pgfpathlineto{\pgfqpoint{1.135988in}{0.759530in}}%
\pgfpathlineto{\pgfqpoint{1.135988in}{0.756581in}}%
\pgfpathmoveto{\pgfqpoint{1.135988in}{0.753632in}}%
\pgfpathlineto{\pgfqpoint{1.135988in}{0.753632in}}%
\pgfpathlineto{\pgfqpoint{1.135988in}{0.756581in}}%
\pgfpathlineto{\pgfqpoint{1.140529in}{0.756581in}}%
\pgfpathlineto{\pgfqpoint{1.140529in}{0.753632in}}%
\pgfpathmoveto{\pgfqpoint{1.135988in}{0.756581in}}%
\pgfpathlineto{\pgfqpoint{1.135988in}{0.756581in}}%
\pgfpathlineto{\pgfqpoint{1.135988in}{0.759530in}}%
\pgfpathlineto{\pgfqpoint{1.140529in}{0.759530in}}%
\pgfpathlineto{\pgfqpoint{1.140529in}{0.756581in}}%
\pgfpathmoveto{\pgfqpoint{1.140529in}{0.753632in}}%
\pgfpathlineto{\pgfqpoint{1.140529in}{0.753632in}}%
\pgfpathlineto{\pgfqpoint{1.140529in}{0.756581in}}%
\pgfpathlineto{\pgfqpoint{1.145070in}{0.756581in}}%
\pgfpathlineto{\pgfqpoint{1.145070in}{0.753632in}}%
\pgfpathmoveto{\pgfqpoint{1.140529in}{0.756581in}}%
\pgfpathlineto{\pgfqpoint{1.140529in}{0.756581in}}%
\pgfpathlineto{\pgfqpoint{1.140529in}{0.759530in}}%
\pgfpathlineto{\pgfqpoint{1.145070in}{0.759530in}}%
\pgfpathlineto{\pgfqpoint{1.145070in}{0.756581in}}%
\pgfpathmoveto{\pgfqpoint{1.145070in}{0.756581in}}%
\pgfpathlineto{\pgfqpoint{1.145070in}{0.756581in}}%
\pgfpathlineto{\pgfqpoint{1.145070in}{0.759530in}}%
\pgfpathlineto{\pgfqpoint{1.149611in}{0.759530in}}%
\pgfpathlineto{\pgfqpoint{1.149611in}{0.756581in}}%
\pgfpathmoveto{\pgfqpoint{1.140529in}{0.759530in}}%
\pgfpathlineto{\pgfqpoint{1.140529in}{0.759530in}}%
\pgfpathlineto{\pgfqpoint{1.140529in}{0.762480in}}%
\pgfpathlineto{\pgfqpoint{1.145070in}{0.762480in}}%
\pgfpathlineto{\pgfqpoint{1.145070in}{0.759530in}}%
\pgfpathmoveto{\pgfqpoint{1.140529in}{0.762480in}}%
\pgfpathlineto{\pgfqpoint{1.140529in}{0.762480in}}%
\pgfpathlineto{\pgfqpoint{1.140529in}{0.765429in}}%
\pgfpathlineto{\pgfqpoint{1.145070in}{0.765429in}}%
\pgfpathlineto{\pgfqpoint{1.145070in}{0.762480in}}%
\pgfpathmoveto{\pgfqpoint{1.145070in}{0.759530in}}%
\pgfpathlineto{\pgfqpoint{1.145070in}{0.759530in}}%
\pgfpathlineto{\pgfqpoint{1.145070in}{0.762480in}}%
\pgfpathlineto{\pgfqpoint{1.149611in}{0.762480in}}%
\pgfpathlineto{\pgfqpoint{1.149611in}{0.759530in}}%
\pgfpathmoveto{\pgfqpoint{1.145070in}{0.762480in}}%
\pgfpathlineto{\pgfqpoint{1.145070in}{0.762480in}}%
\pgfpathlineto{\pgfqpoint{1.145070in}{0.765429in}}%
\pgfpathlineto{\pgfqpoint{1.149611in}{0.765429in}}%
\pgfpathlineto{\pgfqpoint{1.149611in}{0.762480in}}%
\pgfpathmoveto{\pgfqpoint{1.149611in}{0.759530in}}%
\pgfpathlineto{\pgfqpoint{1.149611in}{0.759530in}}%
\pgfpathlineto{\pgfqpoint{1.149611in}{0.762480in}}%
\pgfpathlineto{\pgfqpoint{1.154153in}{0.762480in}}%
\pgfpathlineto{\pgfqpoint{1.154153in}{0.759530in}}%
\pgfpathmoveto{\pgfqpoint{1.149611in}{0.762480in}}%
\pgfpathlineto{\pgfqpoint{1.149611in}{0.762480in}}%
\pgfpathlineto{\pgfqpoint{1.149611in}{0.765429in}}%
\pgfpathlineto{\pgfqpoint{1.154153in}{0.765429in}}%
\pgfpathlineto{\pgfqpoint{1.154153in}{0.762480in}}%
\pgfpathmoveto{\pgfqpoint{1.154153in}{0.762480in}}%
\pgfpathlineto{\pgfqpoint{1.154153in}{0.762480in}}%
\pgfpathlineto{\pgfqpoint{1.154153in}{0.765429in}}%
\pgfpathlineto{\pgfqpoint{1.158694in}{0.765429in}}%
\pgfpathlineto{\pgfqpoint{1.158694in}{0.762480in}}%
\pgfpathmoveto{\pgfqpoint{1.149611in}{0.765429in}}%
\pgfpathlineto{\pgfqpoint{1.149611in}{0.765429in}}%
\pgfpathlineto{\pgfqpoint{1.149611in}{0.768378in}}%
\pgfpathlineto{\pgfqpoint{1.154153in}{0.768378in}}%
\pgfpathlineto{\pgfqpoint{1.154153in}{0.765429in}}%
\pgfpathmoveto{\pgfqpoint{1.149611in}{0.768378in}}%
\pgfpathlineto{\pgfqpoint{1.149611in}{0.768378in}}%
\pgfpathlineto{\pgfqpoint{1.149611in}{0.771327in}}%
\pgfpathlineto{\pgfqpoint{1.154153in}{0.771327in}}%
\pgfpathlineto{\pgfqpoint{1.154153in}{0.768378in}}%
\pgfpathmoveto{\pgfqpoint{1.154153in}{0.765429in}}%
\pgfpathlineto{\pgfqpoint{1.154153in}{0.765429in}}%
\pgfpathlineto{\pgfqpoint{1.154153in}{0.768378in}}%
\pgfpathlineto{\pgfqpoint{1.158694in}{0.768378in}}%
\pgfpathlineto{\pgfqpoint{1.158694in}{0.765429in}}%
\pgfpathmoveto{\pgfqpoint{1.154153in}{0.768378in}}%
\pgfpathlineto{\pgfqpoint{1.154153in}{0.768378in}}%
\pgfpathlineto{\pgfqpoint{1.154153in}{0.771327in}}%
\pgfpathlineto{\pgfqpoint{1.158694in}{0.771327in}}%
\pgfpathlineto{\pgfqpoint{1.158694in}{0.768378in}}%
\pgfpathmoveto{\pgfqpoint{1.158694in}{0.765429in}}%
\pgfpathlineto{\pgfqpoint{1.158694in}{0.765429in}}%
\pgfpathlineto{\pgfqpoint{1.158694in}{0.768378in}}%
\pgfpathlineto{\pgfqpoint{1.163235in}{0.768378in}}%
\pgfpathlineto{\pgfqpoint{1.163235in}{0.765429in}}%
\pgfpathmoveto{\pgfqpoint{1.158694in}{0.768378in}}%
\pgfpathlineto{\pgfqpoint{1.158694in}{0.768378in}}%
\pgfpathlineto{\pgfqpoint{1.158694in}{0.771327in}}%
\pgfpathlineto{\pgfqpoint{1.163235in}{0.771327in}}%
\pgfpathlineto{\pgfqpoint{1.163235in}{0.768378in}}%
\pgfpathmoveto{\pgfqpoint{1.163235in}{0.768378in}}%
\pgfpathlineto{\pgfqpoint{1.163235in}{0.768378in}}%
\pgfpathlineto{\pgfqpoint{1.163235in}{0.771327in}}%
\pgfpathlineto{\pgfqpoint{1.167776in}{0.771327in}}%
\pgfpathlineto{\pgfqpoint{1.167776in}{0.768378in}}%
\pgfpathmoveto{\pgfqpoint{1.158694in}{0.771327in}}%
\pgfpathlineto{\pgfqpoint{1.158694in}{0.771327in}}%
\pgfpathlineto{\pgfqpoint{1.158694in}{0.774276in}}%
\pgfpathlineto{\pgfqpoint{1.163235in}{0.774276in}}%
\pgfpathlineto{\pgfqpoint{1.163235in}{0.771327in}}%
\pgfpathmoveto{\pgfqpoint{1.158694in}{0.774276in}}%
\pgfpathlineto{\pgfqpoint{1.158694in}{0.774276in}}%
\pgfpathlineto{\pgfqpoint{1.158694in}{0.777225in}}%
\pgfpathlineto{\pgfqpoint{1.163235in}{0.777225in}}%
\pgfpathlineto{\pgfqpoint{1.163235in}{0.774276in}}%
\pgfpathmoveto{\pgfqpoint{1.163235in}{0.771327in}}%
\pgfpathlineto{\pgfqpoint{1.163235in}{0.771327in}}%
\pgfpathlineto{\pgfqpoint{1.163235in}{0.774276in}}%
\pgfpathlineto{\pgfqpoint{1.167776in}{0.774276in}}%
\pgfpathlineto{\pgfqpoint{1.167776in}{0.771327in}}%
\pgfpathmoveto{\pgfqpoint{1.163235in}{0.774276in}}%
\pgfpathlineto{\pgfqpoint{1.163235in}{0.774276in}}%
\pgfpathlineto{\pgfqpoint{1.163235in}{0.777225in}}%
\pgfpathlineto{\pgfqpoint{1.167776in}{0.777225in}}%
\pgfpathlineto{\pgfqpoint{1.167776in}{0.774276in}}%
\pgfpathmoveto{\pgfqpoint{1.167776in}{0.771327in}}%
\pgfpathlineto{\pgfqpoint{1.167776in}{0.771327in}}%
\pgfpathlineto{\pgfqpoint{1.167776in}{0.774276in}}%
\pgfpathlineto{\pgfqpoint{1.172318in}{0.774276in}}%
\pgfpathlineto{\pgfqpoint{1.172318in}{0.771327in}}%
\pgfpathmoveto{\pgfqpoint{1.167776in}{0.774276in}}%
\pgfpathlineto{\pgfqpoint{1.167776in}{0.774276in}}%
\pgfpathlineto{\pgfqpoint{1.167776in}{0.777225in}}%
\pgfpathlineto{\pgfqpoint{1.172318in}{0.777225in}}%
\pgfpathlineto{\pgfqpoint{1.172318in}{0.774276in}}%
\pgfpathmoveto{\pgfqpoint{1.172318in}{0.774276in}}%
\pgfpathlineto{\pgfqpoint{1.172318in}{0.774276in}}%
\pgfpathlineto{\pgfqpoint{1.172318in}{0.777225in}}%
\pgfpathlineto{\pgfqpoint{1.176859in}{0.777225in}}%
\pgfpathlineto{\pgfqpoint{1.176859in}{0.774276in}}%
\pgfpathmoveto{\pgfqpoint{1.167776in}{0.777225in}}%
\pgfpathlineto{\pgfqpoint{1.167776in}{0.777225in}}%
\pgfpathlineto{\pgfqpoint{1.167776in}{0.780175in}}%
\pgfpathlineto{\pgfqpoint{1.172318in}{0.780175in}}%
\pgfpathlineto{\pgfqpoint{1.172318in}{0.777225in}}%
\pgfpathmoveto{\pgfqpoint{1.167776in}{0.780175in}}%
\pgfpathlineto{\pgfqpoint{1.167776in}{0.780175in}}%
\pgfpathlineto{\pgfqpoint{1.167776in}{0.783124in}}%
\pgfpathlineto{\pgfqpoint{1.172318in}{0.783124in}}%
\pgfpathlineto{\pgfqpoint{1.172318in}{0.780175in}}%
\pgfpathmoveto{\pgfqpoint{1.172318in}{0.777225in}}%
\pgfpathlineto{\pgfqpoint{1.172318in}{0.777225in}}%
\pgfpathlineto{\pgfqpoint{1.172318in}{0.780175in}}%
\pgfpathlineto{\pgfqpoint{1.176859in}{0.780175in}}%
\pgfpathlineto{\pgfqpoint{1.176859in}{0.777225in}}%
\pgfpathmoveto{\pgfqpoint{1.172318in}{0.780175in}}%
\pgfpathlineto{\pgfqpoint{1.172318in}{0.780175in}}%
\pgfpathlineto{\pgfqpoint{1.172318in}{0.783124in}}%
\pgfpathlineto{\pgfqpoint{1.176859in}{0.783124in}}%
\pgfpathlineto{\pgfqpoint{1.176859in}{0.780175in}}%
\pgfpathmoveto{\pgfqpoint{1.176859in}{0.777225in}}%
\pgfpathlineto{\pgfqpoint{1.176859in}{0.777225in}}%
\pgfpathlineto{\pgfqpoint{1.176859in}{0.780175in}}%
\pgfpathlineto{\pgfqpoint{1.181400in}{0.780175in}}%
\pgfpathlineto{\pgfqpoint{1.181400in}{0.777225in}}%
\pgfpathmoveto{\pgfqpoint{1.176859in}{0.780175in}}%
\pgfpathlineto{\pgfqpoint{1.176859in}{0.780175in}}%
\pgfpathlineto{\pgfqpoint{1.176859in}{0.783124in}}%
\pgfpathlineto{\pgfqpoint{1.181400in}{0.783124in}}%
\pgfpathlineto{\pgfqpoint{1.181400in}{0.780175in}}%
\pgfpathmoveto{\pgfqpoint{1.181400in}{0.780175in}}%
\pgfpathlineto{\pgfqpoint{1.181400in}{0.780175in}}%
\pgfpathlineto{\pgfqpoint{1.181400in}{0.783124in}}%
\pgfpathlineto{\pgfqpoint{1.185941in}{0.783124in}}%
\pgfpathlineto{\pgfqpoint{1.185941in}{0.780175in}}%
\pgfpathmoveto{\pgfqpoint{1.176859in}{0.783124in}}%
\pgfpathlineto{\pgfqpoint{1.176859in}{0.783124in}}%
\pgfpathlineto{\pgfqpoint{1.176859in}{0.786073in}}%
\pgfpathlineto{\pgfqpoint{1.181400in}{0.786073in}}%
\pgfpathlineto{\pgfqpoint{1.181400in}{0.783124in}}%
\pgfpathmoveto{\pgfqpoint{1.176859in}{0.786073in}}%
\pgfpathlineto{\pgfqpoint{1.176859in}{0.786073in}}%
\pgfpathlineto{\pgfqpoint{1.176859in}{0.789022in}}%
\pgfpathlineto{\pgfqpoint{1.181400in}{0.789022in}}%
\pgfpathlineto{\pgfqpoint{1.181400in}{0.786073in}}%
\pgfpathmoveto{\pgfqpoint{1.181400in}{0.783124in}}%
\pgfpathlineto{\pgfqpoint{1.181400in}{0.783124in}}%
\pgfpathlineto{\pgfqpoint{1.181400in}{0.786073in}}%
\pgfpathlineto{\pgfqpoint{1.185941in}{0.786073in}}%
\pgfpathlineto{\pgfqpoint{1.185941in}{0.783124in}}%
\pgfpathmoveto{\pgfqpoint{1.181400in}{0.786073in}}%
\pgfpathlineto{\pgfqpoint{1.181400in}{0.786073in}}%
\pgfpathlineto{\pgfqpoint{1.181400in}{0.789022in}}%
\pgfpathlineto{\pgfqpoint{1.185941in}{0.789022in}}%
\pgfpathlineto{\pgfqpoint{1.185941in}{0.786073in}}%
\pgfpathmoveto{\pgfqpoint{1.185941in}{0.783124in}}%
\pgfpathlineto{\pgfqpoint{1.185941in}{0.783124in}}%
\pgfpathlineto{\pgfqpoint{1.185941in}{0.786073in}}%
\pgfpathlineto{\pgfqpoint{1.190482in}{0.786073in}}%
\pgfpathlineto{\pgfqpoint{1.190482in}{0.783124in}}%
\pgfpathmoveto{\pgfqpoint{1.185941in}{0.786073in}}%
\pgfpathlineto{\pgfqpoint{1.185941in}{0.786073in}}%
\pgfpathlineto{\pgfqpoint{1.185941in}{0.789022in}}%
\pgfpathlineto{\pgfqpoint{1.190482in}{0.789022in}}%
\pgfpathlineto{\pgfqpoint{1.190482in}{0.786073in}}%
\pgfpathmoveto{\pgfqpoint{1.190482in}{0.786073in}}%
\pgfpathlineto{\pgfqpoint{1.190482in}{0.786073in}}%
\pgfpathlineto{\pgfqpoint{1.190482in}{0.789022in}}%
\pgfpathlineto{\pgfqpoint{1.195023in}{0.789022in}}%
\pgfpathlineto{\pgfqpoint{1.195023in}{0.786073in}}%
\pgfpathmoveto{\pgfqpoint{1.185941in}{0.789022in}}%
\pgfpathlineto{\pgfqpoint{1.185941in}{0.789022in}}%
\pgfpathlineto{\pgfqpoint{1.185941in}{0.791972in}}%
\pgfpathlineto{\pgfqpoint{1.190482in}{0.791972in}}%
\pgfpathlineto{\pgfqpoint{1.190482in}{0.789022in}}%
\pgfpathmoveto{\pgfqpoint{1.185941in}{0.791972in}}%
\pgfpathlineto{\pgfqpoint{1.185941in}{0.791972in}}%
\pgfpathlineto{\pgfqpoint{1.185941in}{0.794921in}}%
\pgfpathlineto{\pgfqpoint{1.190482in}{0.794921in}}%
\pgfpathlineto{\pgfqpoint{1.190482in}{0.791972in}}%
\pgfpathmoveto{\pgfqpoint{1.190482in}{0.789022in}}%
\pgfpathlineto{\pgfqpoint{1.190482in}{0.789022in}}%
\pgfpathlineto{\pgfqpoint{1.190482in}{0.791972in}}%
\pgfpathlineto{\pgfqpoint{1.195023in}{0.791972in}}%
\pgfpathlineto{\pgfqpoint{1.195023in}{0.789022in}}%
\pgfpathmoveto{\pgfqpoint{1.190482in}{0.791972in}}%
\pgfpathlineto{\pgfqpoint{1.190482in}{0.791972in}}%
\pgfpathlineto{\pgfqpoint{1.190482in}{0.794921in}}%
\pgfpathlineto{\pgfqpoint{1.195023in}{0.794921in}}%
\pgfpathlineto{\pgfqpoint{1.195023in}{0.791972in}}%
\pgfpathmoveto{\pgfqpoint{1.195023in}{0.789022in}}%
\pgfpathlineto{\pgfqpoint{1.195023in}{0.789022in}}%
\pgfpathlineto{\pgfqpoint{1.195023in}{0.791972in}}%
\pgfpathlineto{\pgfqpoint{1.199564in}{0.791972in}}%
\pgfpathlineto{\pgfqpoint{1.199564in}{0.789022in}}%
\pgfpathmoveto{\pgfqpoint{1.195023in}{0.791972in}}%
\pgfpathlineto{\pgfqpoint{1.195023in}{0.791972in}}%
\pgfpathlineto{\pgfqpoint{1.195023in}{0.794921in}}%
\pgfpathlineto{\pgfqpoint{1.199564in}{0.794921in}}%
\pgfpathlineto{\pgfqpoint{1.199564in}{0.791972in}}%
\pgfpathmoveto{\pgfqpoint{1.199564in}{0.791972in}}%
\pgfpathlineto{\pgfqpoint{1.199564in}{0.791972in}}%
\pgfpathlineto{\pgfqpoint{1.199564in}{0.794921in}}%
\pgfpathlineto{\pgfqpoint{1.204105in}{0.794921in}}%
\pgfpathlineto{\pgfqpoint{1.204105in}{0.791972in}}%
\pgfpathmoveto{\pgfqpoint{1.195023in}{0.794921in}}%
\pgfpathlineto{\pgfqpoint{1.195023in}{0.794921in}}%
\pgfpathlineto{\pgfqpoint{1.195023in}{0.797870in}}%
\pgfpathlineto{\pgfqpoint{1.199564in}{0.797870in}}%
\pgfpathlineto{\pgfqpoint{1.199564in}{0.794921in}}%
\pgfpathmoveto{\pgfqpoint{1.195023in}{0.797870in}}%
\pgfpathlineto{\pgfqpoint{1.195023in}{0.797870in}}%
\pgfpathlineto{\pgfqpoint{1.195023in}{0.800820in}}%
\pgfpathlineto{\pgfqpoint{1.199564in}{0.800820in}}%
\pgfpathlineto{\pgfqpoint{1.199564in}{0.797870in}}%
\pgfpathmoveto{\pgfqpoint{1.199564in}{0.794921in}}%
\pgfpathlineto{\pgfqpoint{1.199564in}{0.794921in}}%
\pgfpathlineto{\pgfqpoint{1.199564in}{0.797870in}}%
\pgfpathlineto{\pgfqpoint{1.204105in}{0.797870in}}%
\pgfpathlineto{\pgfqpoint{1.204105in}{0.794921in}}%
\pgfpathmoveto{\pgfqpoint{1.199564in}{0.797870in}}%
\pgfpathlineto{\pgfqpoint{1.199564in}{0.797870in}}%
\pgfpathlineto{\pgfqpoint{1.199564in}{0.800820in}}%
\pgfpathlineto{\pgfqpoint{1.204105in}{0.800820in}}%
\pgfpathlineto{\pgfqpoint{1.204105in}{0.797870in}}%
\pgfpathmoveto{\pgfqpoint{1.204105in}{0.794921in}}%
\pgfpathlineto{\pgfqpoint{1.204105in}{0.794921in}}%
\pgfpathlineto{\pgfqpoint{1.204105in}{0.797870in}}%
\pgfpathlineto{\pgfqpoint{1.208646in}{0.797870in}}%
\pgfpathlineto{\pgfqpoint{1.208646in}{0.794921in}}%
\pgfpathmoveto{\pgfqpoint{1.204105in}{0.797870in}}%
\pgfpathlineto{\pgfqpoint{1.204105in}{0.797870in}}%
\pgfpathlineto{\pgfqpoint{1.204105in}{0.800820in}}%
\pgfpathlineto{\pgfqpoint{1.208646in}{0.800820in}}%
\pgfpathlineto{\pgfqpoint{1.208646in}{0.797870in}}%
\pgfpathmoveto{\pgfqpoint{1.208646in}{0.797870in}}%
\pgfpathlineto{\pgfqpoint{1.208646in}{0.797870in}}%
\pgfpathlineto{\pgfqpoint{1.208646in}{0.800820in}}%
\pgfpathlineto{\pgfqpoint{1.213187in}{0.800820in}}%
\pgfpathlineto{\pgfqpoint{1.213187in}{0.797870in}}%
\pgfpathmoveto{\pgfqpoint{1.204105in}{0.800820in}}%
\pgfpathlineto{\pgfqpoint{1.204105in}{0.800820in}}%
\pgfpathlineto{\pgfqpoint{1.204105in}{0.803769in}}%
\pgfpathlineto{\pgfqpoint{1.208646in}{0.803769in}}%
\pgfpathlineto{\pgfqpoint{1.208646in}{0.800820in}}%
\pgfpathmoveto{\pgfqpoint{1.204105in}{0.803769in}}%
\pgfpathlineto{\pgfqpoint{1.204105in}{0.803769in}}%
\pgfpathlineto{\pgfqpoint{1.204105in}{0.806718in}}%
\pgfpathlineto{\pgfqpoint{1.208646in}{0.806718in}}%
\pgfpathlineto{\pgfqpoint{1.208646in}{0.803769in}}%
\pgfpathmoveto{\pgfqpoint{1.208646in}{0.800820in}}%
\pgfpathlineto{\pgfqpoint{1.208646in}{0.800820in}}%
\pgfpathlineto{\pgfqpoint{1.208646in}{0.803769in}}%
\pgfpathlineto{\pgfqpoint{1.213187in}{0.803769in}}%
\pgfpathlineto{\pgfqpoint{1.213187in}{0.800820in}}%
\pgfpathmoveto{\pgfqpoint{1.208646in}{0.803769in}}%
\pgfpathlineto{\pgfqpoint{1.208646in}{0.803769in}}%
\pgfpathlineto{\pgfqpoint{1.208646in}{0.806718in}}%
\pgfpathlineto{\pgfqpoint{1.213187in}{0.806718in}}%
\pgfpathlineto{\pgfqpoint{1.213187in}{0.803769in}}%
\pgfpathmoveto{\pgfqpoint{1.213187in}{0.800820in}}%
\pgfpathlineto{\pgfqpoint{1.213187in}{0.800820in}}%
\pgfpathlineto{\pgfqpoint{1.213187in}{0.803769in}}%
\pgfpathlineto{\pgfqpoint{1.217728in}{0.803769in}}%
\pgfpathlineto{\pgfqpoint{1.217728in}{0.800820in}}%
\pgfpathmoveto{\pgfqpoint{1.213187in}{0.803769in}}%
\pgfpathlineto{\pgfqpoint{1.213187in}{0.803769in}}%
\pgfpathlineto{\pgfqpoint{1.213187in}{0.806718in}}%
\pgfpathlineto{\pgfqpoint{1.217728in}{0.806718in}}%
\pgfpathlineto{\pgfqpoint{1.217728in}{0.803769in}}%
\pgfpathmoveto{\pgfqpoint{1.217728in}{0.803769in}}%
\pgfpathlineto{\pgfqpoint{1.217728in}{0.803769in}}%
\pgfpathlineto{\pgfqpoint{1.217728in}{0.806718in}}%
\pgfpathlineto{\pgfqpoint{1.222269in}{0.806718in}}%
\pgfpathlineto{\pgfqpoint{1.222269in}{0.803769in}}%
\pgfpathmoveto{\pgfqpoint{1.213187in}{0.806718in}}%
\pgfpathlineto{\pgfqpoint{1.213187in}{0.806718in}}%
\pgfpathlineto{\pgfqpoint{1.213187in}{0.809667in}}%
\pgfpathlineto{\pgfqpoint{1.217728in}{0.809667in}}%
\pgfpathlineto{\pgfqpoint{1.217728in}{0.806718in}}%
\pgfpathmoveto{\pgfqpoint{1.213187in}{0.809667in}}%
\pgfpathlineto{\pgfqpoint{1.213187in}{0.809667in}}%
\pgfpathlineto{\pgfqpoint{1.213187in}{0.812617in}}%
\pgfpathlineto{\pgfqpoint{1.217728in}{0.812617in}}%
\pgfpathlineto{\pgfqpoint{1.217728in}{0.809667in}}%
\pgfpathmoveto{\pgfqpoint{1.217728in}{0.806718in}}%
\pgfpathlineto{\pgfqpoint{1.217728in}{0.806718in}}%
\pgfpathlineto{\pgfqpoint{1.217728in}{0.809667in}}%
\pgfpathlineto{\pgfqpoint{1.222269in}{0.809667in}}%
\pgfpathlineto{\pgfqpoint{1.222269in}{0.806718in}}%
\pgfpathmoveto{\pgfqpoint{1.217728in}{0.809667in}}%
\pgfpathlineto{\pgfqpoint{1.217728in}{0.809667in}}%
\pgfpathlineto{\pgfqpoint{1.217728in}{0.812617in}}%
\pgfpathlineto{\pgfqpoint{1.222269in}{0.812617in}}%
\pgfpathlineto{\pgfqpoint{1.222269in}{0.809667in}}%
\pgfpathmoveto{\pgfqpoint{1.222269in}{0.806718in}}%
\pgfpathlineto{\pgfqpoint{1.222269in}{0.806718in}}%
\pgfpathlineto{\pgfqpoint{1.222269in}{0.809667in}}%
\pgfpathlineto{\pgfqpoint{1.226810in}{0.809667in}}%
\pgfpathlineto{\pgfqpoint{1.226810in}{0.806718in}}%
\pgfpathmoveto{\pgfqpoint{1.222269in}{0.809667in}}%
\pgfpathlineto{\pgfqpoint{1.222269in}{0.809667in}}%
\pgfpathlineto{\pgfqpoint{1.222269in}{0.812617in}}%
\pgfpathlineto{\pgfqpoint{1.226810in}{0.812617in}}%
\pgfpathlineto{\pgfqpoint{1.226810in}{0.809667in}}%
\pgfpathmoveto{\pgfqpoint{1.226810in}{0.809667in}}%
\pgfpathlineto{\pgfqpoint{1.226810in}{0.809667in}}%
\pgfpathlineto{\pgfqpoint{1.226810in}{0.812617in}}%
\pgfpathlineto{\pgfqpoint{1.231351in}{0.812617in}}%
\pgfpathlineto{\pgfqpoint{1.231351in}{0.809667in}}%
\pgfpathmoveto{\pgfqpoint{1.222269in}{0.812617in}}%
\pgfpathlineto{\pgfqpoint{1.222269in}{0.812617in}}%
\pgfpathlineto{\pgfqpoint{1.222269in}{0.815566in}}%
\pgfpathlineto{\pgfqpoint{1.226810in}{0.815566in}}%
\pgfpathlineto{\pgfqpoint{1.226810in}{0.812617in}}%
\pgfpathmoveto{\pgfqpoint{1.222269in}{0.815566in}}%
\pgfpathlineto{\pgfqpoint{1.222269in}{0.815566in}}%
\pgfpathlineto{\pgfqpoint{1.222269in}{0.818515in}}%
\pgfpathlineto{\pgfqpoint{1.226810in}{0.818515in}}%
\pgfpathlineto{\pgfqpoint{1.226810in}{0.815566in}}%
\pgfpathmoveto{\pgfqpoint{1.226810in}{0.812617in}}%
\pgfpathlineto{\pgfqpoint{1.226810in}{0.812617in}}%
\pgfpathlineto{\pgfqpoint{1.226810in}{0.815566in}}%
\pgfpathlineto{\pgfqpoint{1.231351in}{0.815566in}}%
\pgfpathlineto{\pgfqpoint{1.231351in}{0.812617in}}%
\pgfpathmoveto{\pgfqpoint{1.226810in}{0.815566in}}%
\pgfpathlineto{\pgfqpoint{1.226810in}{0.815566in}}%
\pgfpathlineto{\pgfqpoint{1.226810in}{0.818515in}}%
\pgfpathlineto{\pgfqpoint{1.231351in}{0.818515in}}%
\pgfpathlineto{\pgfqpoint{1.231351in}{0.815566in}}%
\pgfpathmoveto{\pgfqpoint{1.231351in}{0.812617in}}%
\pgfpathlineto{\pgfqpoint{1.231351in}{0.812617in}}%
\pgfpathlineto{\pgfqpoint{1.231351in}{0.815566in}}%
\pgfpathlineto{\pgfqpoint{1.235892in}{0.815566in}}%
\pgfpathlineto{\pgfqpoint{1.235892in}{0.812617in}}%
\pgfpathmoveto{\pgfqpoint{1.231351in}{0.815566in}}%
\pgfpathlineto{\pgfqpoint{1.231351in}{0.815566in}}%
\pgfpathlineto{\pgfqpoint{1.231351in}{0.818515in}}%
\pgfpathlineto{\pgfqpoint{1.235892in}{0.818515in}}%
\pgfpathlineto{\pgfqpoint{1.235892in}{0.815566in}}%
\pgfpathmoveto{\pgfqpoint{1.235892in}{0.815566in}}%
\pgfpathlineto{\pgfqpoint{1.235892in}{0.815566in}}%
\pgfpathlineto{\pgfqpoint{1.235892in}{0.818515in}}%
\pgfpathlineto{\pgfqpoint{1.240433in}{0.818515in}}%
\pgfpathlineto{\pgfqpoint{1.240433in}{0.815566in}}%
\pgfpathmoveto{\pgfqpoint{1.231351in}{0.818515in}}%
\pgfpathlineto{\pgfqpoint{1.231351in}{0.818515in}}%
\pgfpathlineto{\pgfqpoint{1.231351in}{0.821465in}}%
\pgfpathlineto{\pgfqpoint{1.235892in}{0.821465in}}%
\pgfpathlineto{\pgfqpoint{1.235892in}{0.818515in}}%
\pgfpathmoveto{\pgfqpoint{1.231351in}{0.821465in}}%
\pgfpathlineto{\pgfqpoint{1.231351in}{0.821465in}}%
\pgfpathlineto{\pgfqpoint{1.231351in}{0.824414in}}%
\pgfpathlineto{\pgfqpoint{1.235892in}{0.824414in}}%
\pgfpathlineto{\pgfqpoint{1.235892in}{0.821465in}}%
\pgfpathmoveto{\pgfqpoint{1.235892in}{0.818515in}}%
\pgfpathlineto{\pgfqpoint{1.235892in}{0.818515in}}%
\pgfpathlineto{\pgfqpoint{1.235892in}{0.821465in}}%
\pgfpathlineto{\pgfqpoint{1.240433in}{0.821465in}}%
\pgfpathlineto{\pgfqpoint{1.240433in}{0.818515in}}%
\pgfpathmoveto{\pgfqpoint{1.235892in}{0.821465in}}%
\pgfpathlineto{\pgfqpoint{1.235892in}{0.821465in}}%
\pgfpathlineto{\pgfqpoint{1.235892in}{0.824414in}}%
\pgfpathlineto{\pgfqpoint{1.240433in}{0.824414in}}%
\pgfpathlineto{\pgfqpoint{1.240433in}{0.821465in}}%
\pgfpathmoveto{\pgfqpoint{1.240433in}{0.818515in}}%
\pgfpathlineto{\pgfqpoint{1.240433in}{0.818515in}}%
\pgfpathlineto{\pgfqpoint{1.240433in}{0.821465in}}%
\pgfpathlineto{\pgfqpoint{1.244975in}{0.821465in}}%
\pgfpathlineto{\pgfqpoint{1.244975in}{0.818515in}}%
\pgfpathmoveto{\pgfqpoint{1.240433in}{0.821465in}}%
\pgfpathlineto{\pgfqpoint{1.240433in}{0.821465in}}%
\pgfpathlineto{\pgfqpoint{1.240433in}{0.824414in}}%
\pgfpathlineto{\pgfqpoint{1.244975in}{0.824414in}}%
\pgfpathlineto{\pgfqpoint{1.244975in}{0.821465in}}%
\pgfpathmoveto{\pgfqpoint{1.244975in}{0.821465in}}%
\pgfpathlineto{\pgfqpoint{1.244975in}{0.821465in}}%
\pgfpathlineto{\pgfqpoint{1.244975in}{0.824414in}}%
\pgfpathlineto{\pgfqpoint{1.249516in}{0.824414in}}%
\pgfpathlineto{\pgfqpoint{1.249516in}{0.821465in}}%
\pgfpathmoveto{\pgfqpoint{1.240433in}{0.824414in}}%
\pgfpathlineto{\pgfqpoint{1.240433in}{0.824414in}}%
\pgfpathlineto{\pgfqpoint{1.240433in}{0.827363in}}%
\pgfpathlineto{\pgfqpoint{1.244975in}{0.827363in}}%
\pgfpathlineto{\pgfqpoint{1.244975in}{0.824414in}}%
\pgfpathmoveto{\pgfqpoint{1.240433in}{0.827363in}}%
\pgfpathlineto{\pgfqpoint{1.240433in}{0.827363in}}%
\pgfpathlineto{\pgfqpoint{1.240433in}{0.830313in}}%
\pgfpathlineto{\pgfqpoint{1.244975in}{0.830313in}}%
\pgfpathlineto{\pgfqpoint{1.244975in}{0.827363in}}%
\pgfpathmoveto{\pgfqpoint{1.244975in}{0.824414in}}%
\pgfpathlineto{\pgfqpoint{1.244975in}{0.824414in}}%
\pgfpathlineto{\pgfqpoint{1.244975in}{0.827363in}}%
\pgfpathlineto{\pgfqpoint{1.249516in}{0.827363in}}%
\pgfpathlineto{\pgfqpoint{1.249516in}{0.824414in}}%
\pgfpathmoveto{\pgfqpoint{1.244975in}{0.827363in}}%
\pgfpathlineto{\pgfqpoint{1.244975in}{0.827363in}}%
\pgfpathlineto{\pgfqpoint{1.244975in}{0.830313in}}%
\pgfpathlineto{\pgfqpoint{1.249516in}{0.830313in}}%
\pgfpathlineto{\pgfqpoint{1.249516in}{0.827363in}}%
\pgfpathmoveto{\pgfqpoint{1.249516in}{0.824414in}}%
\pgfpathlineto{\pgfqpoint{1.249516in}{0.824414in}}%
\pgfpathlineto{\pgfqpoint{1.249516in}{0.827363in}}%
\pgfpathlineto{\pgfqpoint{1.254057in}{0.827363in}}%
\pgfpathlineto{\pgfqpoint{1.254057in}{0.824414in}}%
\pgfpathmoveto{\pgfqpoint{1.249516in}{0.827363in}}%
\pgfpathlineto{\pgfqpoint{1.249516in}{0.827363in}}%
\pgfpathlineto{\pgfqpoint{1.249516in}{0.830313in}}%
\pgfpathlineto{\pgfqpoint{1.254057in}{0.830313in}}%
\pgfpathlineto{\pgfqpoint{1.254057in}{0.827363in}}%
\pgfpathmoveto{\pgfqpoint{1.254057in}{0.827363in}}%
\pgfpathlineto{\pgfqpoint{1.254057in}{0.827363in}}%
\pgfpathlineto{\pgfqpoint{1.254057in}{0.830313in}}%
\pgfpathlineto{\pgfqpoint{1.258598in}{0.830313in}}%
\pgfpathlineto{\pgfqpoint{1.258598in}{0.827363in}}%
\pgfpathmoveto{\pgfqpoint{1.249516in}{0.830313in}}%
\pgfpathlineto{\pgfqpoint{1.249516in}{0.830313in}}%
\pgfpathlineto{\pgfqpoint{1.249516in}{0.833262in}}%
\pgfpathlineto{\pgfqpoint{1.254057in}{0.833262in}}%
\pgfpathlineto{\pgfqpoint{1.254057in}{0.830313in}}%
\pgfpathmoveto{\pgfqpoint{1.249516in}{0.833262in}}%
\pgfpathlineto{\pgfqpoint{1.249516in}{0.833262in}}%
\pgfpathlineto{\pgfqpoint{1.249516in}{0.836211in}}%
\pgfpathlineto{\pgfqpoint{1.254057in}{0.836211in}}%
\pgfpathlineto{\pgfqpoint{1.254057in}{0.833262in}}%
\pgfpathmoveto{\pgfqpoint{1.254057in}{0.830313in}}%
\pgfpathlineto{\pgfqpoint{1.254057in}{0.830313in}}%
\pgfpathlineto{\pgfqpoint{1.254057in}{0.833262in}}%
\pgfpathlineto{\pgfqpoint{1.258598in}{0.833262in}}%
\pgfpathlineto{\pgfqpoint{1.258598in}{0.830313in}}%
\pgfpathmoveto{\pgfqpoint{1.254057in}{0.833262in}}%
\pgfpathlineto{\pgfqpoint{1.254057in}{0.833262in}}%
\pgfpathlineto{\pgfqpoint{1.254057in}{0.836211in}}%
\pgfpathlineto{\pgfqpoint{1.258598in}{0.836211in}}%
\pgfpathlineto{\pgfqpoint{1.258598in}{0.833262in}}%
\pgfpathmoveto{\pgfqpoint{1.258598in}{0.830313in}}%
\pgfpathlineto{\pgfqpoint{1.258598in}{0.830313in}}%
\pgfpathlineto{\pgfqpoint{1.258598in}{0.833262in}}%
\pgfpathlineto{\pgfqpoint{1.263139in}{0.833262in}}%
\pgfpathlineto{\pgfqpoint{1.263139in}{0.830313in}}%
\pgfpathmoveto{\pgfqpoint{1.258598in}{0.833262in}}%
\pgfpathlineto{\pgfqpoint{1.258598in}{0.833262in}}%
\pgfpathlineto{\pgfqpoint{1.258598in}{0.836211in}}%
\pgfpathlineto{\pgfqpoint{1.263139in}{0.836211in}}%
\pgfpathlineto{\pgfqpoint{1.263139in}{0.833262in}}%
\pgfpathmoveto{\pgfqpoint{1.263139in}{0.833262in}}%
\pgfpathlineto{\pgfqpoint{1.263139in}{0.833262in}}%
\pgfpathlineto{\pgfqpoint{1.263139in}{0.836211in}}%
\pgfpathlineto{\pgfqpoint{1.267680in}{0.836211in}}%
\pgfpathlineto{\pgfqpoint{1.267680in}{0.833262in}}%
\pgfpathmoveto{\pgfqpoint{1.258598in}{0.836211in}}%
\pgfpathlineto{\pgfqpoint{1.258598in}{0.836211in}}%
\pgfpathlineto{\pgfqpoint{1.258598in}{0.839160in}}%
\pgfpathlineto{\pgfqpoint{1.263139in}{0.839160in}}%
\pgfpathlineto{\pgfqpoint{1.263139in}{0.836211in}}%
\pgfpathmoveto{\pgfqpoint{1.258598in}{0.839160in}}%
\pgfpathlineto{\pgfqpoint{1.258598in}{0.839160in}}%
\pgfpathlineto{\pgfqpoint{1.258598in}{0.842110in}}%
\pgfpathlineto{\pgfqpoint{1.263139in}{0.842110in}}%
\pgfpathlineto{\pgfqpoint{1.263139in}{0.839160in}}%
\pgfpathmoveto{\pgfqpoint{1.263139in}{0.836211in}}%
\pgfpathlineto{\pgfqpoint{1.263139in}{0.836211in}}%
\pgfpathlineto{\pgfqpoint{1.263139in}{0.839160in}}%
\pgfpathlineto{\pgfqpoint{1.267680in}{0.839160in}}%
\pgfpathlineto{\pgfqpoint{1.267680in}{0.836211in}}%
\pgfpathmoveto{\pgfqpoint{1.263139in}{0.839160in}}%
\pgfpathlineto{\pgfqpoint{1.263139in}{0.839160in}}%
\pgfpathlineto{\pgfqpoint{1.263139in}{0.842110in}}%
\pgfpathlineto{\pgfqpoint{1.267680in}{0.842110in}}%
\pgfpathlineto{\pgfqpoint{1.267680in}{0.839160in}}%
\pgfpathmoveto{\pgfqpoint{1.267680in}{0.836211in}}%
\pgfpathlineto{\pgfqpoint{1.267680in}{0.836211in}}%
\pgfpathlineto{\pgfqpoint{1.267680in}{0.839160in}}%
\pgfpathlineto{\pgfqpoint{1.272221in}{0.839160in}}%
\pgfpathlineto{\pgfqpoint{1.272221in}{0.836211in}}%
\pgfpathmoveto{\pgfqpoint{1.267680in}{0.839160in}}%
\pgfpathlineto{\pgfqpoint{1.267680in}{0.839160in}}%
\pgfpathlineto{\pgfqpoint{1.267680in}{0.842110in}}%
\pgfpathlineto{\pgfqpoint{1.272221in}{0.842110in}}%
\pgfpathlineto{\pgfqpoint{1.272221in}{0.839160in}}%
\pgfpathmoveto{\pgfqpoint{1.272221in}{0.839160in}}%
\pgfpathlineto{\pgfqpoint{1.272221in}{0.839160in}}%
\pgfpathlineto{\pgfqpoint{1.272221in}{0.842110in}}%
\pgfpathlineto{\pgfqpoint{1.276762in}{0.842110in}}%
\pgfpathlineto{\pgfqpoint{1.276762in}{0.839160in}}%
\pgfpathmoveto{\pgfqpoint{1.267680in}{0.842110in}}%
\pgfpathlineto{\pgfqpoint{1.267680in}{0.842110in}}%
\pgfpathlineto{\pgfqpoint{1.267680in}{0.845059in}}%
\pgfpathlineto{\pgfqpoint{1.272221in}{0.845059in}}%
\pgfpathlineto{\pgfqpoint{1.272221in}{0.842110in}}%
\pgfpathmoveto{\pgfqpoint{1.267680in}{0.845059in}}%
\pgfpathlineto{\pgfqpoint{1.267680in}{0.845059in}}%
\pgfpathlineto{\pgfqpoint{1.267680in}{0.848008in}}%
\pgfpathlineto{\pgfqpoint{1.272221in}{0.848008in}}%
\pgfpathlineto{\pgfqpoint{1.272221in}{0.845059in}}%
\pgfpathmoveto{\pgfqpoint{1.272221in}{0.842110in}}%
\pgfpathlineto{\pgfqpoint{1.272221in}{0.842110in}}%
\pgfpathlineto{\pgfqpoint{1.272221in}{0.845059in}}%
\pgfpathlineto{\pgfqpoint{1.276762in}{0.845059in}}%
\pgfpathlineto{\pgfqpoint{1.276762in}{0.842110in}}%
\pgfpathmoveto{\pgfqpoint{1.272221in}{0.845059in}}%
\pgfpathlineto{\pgfqpoint{1.272221in}{0.845059in}}%
\pgfpathlineto{\pgfqpoint{1.272221in}{0.848008in}}%
\pgfpathlineto{\pgfqpoint{1.276762in}{0.848008in}}%
\pgfpathlineto{\pgfqpoint{1.276762in}{0.845059in}}%
\pgfpathmoveto{\pgfqpoint{1.276762in}{0.842110in}}%
\pgfpathlineto{\pgfqpoint{1.276762in}{0.842110in}}%
\pgfpathlineto{\pgfqpoint{1.276762in}{0.845059in}}%
\pgfpathlineto{\pgfqpoint{1.281303in}{0.845059in}}%
\pgfpathlineto{\pgfqpoint{1.281303in}{0.842110in}}%
\pgfpathmoveto{\pgfqpoint{1.276762in}{0.845059in}}%
\pgfpathlineto{\pgfqpoint{1.276762in}{0.845059in}}%
\pgfpathlineto{\pgfqpoint{1.276762in}{0.848008in}}%
\pgfpathlineto{\pgfqpoint{1.281303in}{0.848008in}}%
\pgfpathlineto{\pgfqpoint{1.281303in}{0.845059in}}%
\pgfpathmoveto{\pgfqpoint{1.281303in}{0.845059in}}%
\pgfpathlineto{\pgfqpoint{1.281303in}{0.845059in}}%
\pgfpathlineto{\pgfqpoint{1.281303in}{0.848008in}}%
\pgfpathlineto{\pgfqpoint{1.285844in}{0.848008in}}%
\pgfpathlineto{\pgfqpoint{1.285844in}{0.845059in}}%
\pgfpathmoveto{\pgfqpoint{1.276762in}{0.848008in}}%
\pgfpathlineto{\pgfqpoint{1.276762in}{0.848008in}}%
\pgfpathlineto{\pgfqpoint{1.276762in}{0.850958in}}%
\pgfpathlineto{\pgfqpoint{1.281303in}{0.850958in}}%
\pgfpathlineto{\pgfqpoint{1.281303in}{0.848008in}}%
\pgfpathmoveto{\pgfqpoint{1.276762in}{0.850958in}}%
\pgfpathlineto{\pgfqpoint{1.276762in}{0.850958in}}%
\pgfpathlineto{\pgfqpoint{1.276762in}{0.853907in}}%
\pgfpathlineto{\pgfqpoint{1.281303in}{0.853907in}}%
\pgfpathlineto{\pgfqpoint{1.281303in}{0.850958in}}%
\pgfpathmoveto{\pgfqpoint{1.281303in}{0.848008in}}%
\pgfpathlineto{\pgfqpoint{1.281303in}{0.848008in}}%
\pgfpathlineto{\pgfqpoint{1.281303in}{0.850958in}}%
\pgfpathlineto{\pgfqpoint{1.285844in}{0.850958in}}%
\pgfpathlineto{\pgfqpoint{1.285844in}{0.848008in}}%
\pgfpathmoveto{\pgfqpoint{1.281303in}{0.850958in}}%
\pgfpathlineto{\pgfqpoint{1.281303in}{0.850958in}}%
\pgfpathlineto{\pgfqpoint{1.281303in}{0.853907in}}%
\pgfpathlineto{\pgfqpoint{1.285844in}{0.853907in}}%
\pgfpathlineto{\pgfqpoint{1.285844in}{0.850958in}}%
\pgfpathmoveto{\pgfqpoint{1.285844in}{0.848008in}}%
\pgfpathlineto{\pgfqpoint{1.285844in}{0.848008in}}%
\pgfpathlineto{\pgfqpoint{1.285844in}{0.850958in}}%
\pgfpathlineto{\pgfqpoint{1.290385in}{0.850958in}}%
\pgfpathlineto{\pgfqpoint{1.290385in}{0.848008in}}%
\pgfpathmoveto{\pgfqpoint{1.285844in}{0.850958in}}%
\pgfpathlineto{\pgfqpoint{1.285844in}{0.850958in}}%
\pgfpathlineto{\pgfqpoint{1.285844in}{0.853907in}}%
\pgfpathlineto{\pgfqpoint{1.290385in}{0.853907in}}%
\pgfpathlineto{\pgfqpoint{1.290385in}{0.850958in}}%
\pgfpathmoveto{\pgfqpoint{1.290385in}{0.850958in}}%
\pgfpathlineto{\pgfqpoint{1.290385in}{0.850958in}}%
\pgfpathlineto{\pgfqpoint{1.290385in}{0.853907in}}%
\pgfpathlineto{\pgfqpoint{1.294926in}{0.853907in}}%
\pgfpathlineto{\pgfqpoint{1.294926in}{0.850958in}}%
\pgfpathmoveto{\pgfqpoint{1.285844in}{0.853907in}}%
\pgfpathlineto{\pgfqpoint{1.285844in}{0.853907in}}%
\pgfpathlineto{\pgfqpoint{1.285844in}{0.856856in}}%
\pgfpathlineto{\pgfqpoint{1.290385in}{0.856856in}}%
\pgfpathlineto{\pgfqpoint{1.290385in}{0.853907in}}%
\pgfpathmoveto{\pgfqpoint{1.285844in}{0.856856in}}%
\pgfpathlineto{\pgfqpoint{1.285844in}{0.856856in}}%
\pgfpathlineto{\pgfqpoint{1.285844in}{0.859806in}}%
\pgfpathlineto{\pgfqpoint{1.290385in}{0.859806in}}%
\pgfpathlineto{\pgfqpoint{1.290385in}{0.856856in}}%
\pgfpathmoveto{\pgfqpoint{1.290385in}{0.853907in}}%
\pgfpathlineto{\pgfqpoint{1.290385in}{0.853907in}}%
\pgfpathlineto{\pgfqpoint{1.290385in}{0.856856in}}%
\pgfpathlineto{\pgfqpoint{1.294926in}{0.856856in}}%
\pgfpathlineto{\pgfqpoint{1.294926in}{0.853907in}}%
\pgfpathmoveto{\pgfqpoint{1.290385in}{0.856856in}}%
\pgfpathlineto{\pgfqpoint{1.290385in}{0.856856in}}%
\pgfpathlineto{\pgfqpoint{1.290385in}{0.859806in}}%
\pgfpathlineto{\pgfqpoint{1.294926in}{0.859806in}}%
\pgfpathlineto{\pgfqpoint{1.294926in}{0.856856in}}%
\pgfpathmoveto{\pgfqpoint{1.294926in}{0.853907in}}%
\pgfpathlineto{\pgfqpoint{1.294926in}{0.853907in}}%
\pgfpathlineto{\pgfqpoint{1.294926in}{0.856856in}}%
\pgfpathlineto{\pgfqpoint{1.299467in}{0.856856in}}%
\pgfpathlineto{\pgfqpoint{1.299467in}{0.853907in}}%
\pgfpathmoveto{\pgfqpoint{1.294926in}{0.856856in}}%
\pgfpathlineto{\pgfqpoint{1.294926in}{0.856856in}}%
\pgfpathlineto{\pgfqpoint{1.294926in}{0.859806in}}%
\pgfpathlineto{\pgfqpoint{1.299467in}{0.859806in}}%
\pgfpathlineto{\pgfqpoint{1.299467in}{0.856856in}}%
\pgfpathmoveto{\pgfqpoint{1.299467in}{0.856856in}}%
\pgfpathlineto{\pgfqpoint{1.299467in}{0.856856in}}%
\pgfpathlineto{\pgfqpoint{1.299467in}{0.859806in}}%
\pgfpathlineto{\pgfqpoint{1.304008in}{0.859806in}}%
\pgfpathlineto{\pgfqpoint{1.304008in}{0.856856in}}%
\pgfpathmoveto{\pgfqpoint{1.294926in}{0.859806in}}%
\pgfpathlineto{\pgfqpoint{1.294926in}{0.859806in}}%
\pgfpathlineto{\pgfqpoint{1.294926in}{0.862755in}}%
\pgfpathlineto{\pgfqpoint{1.299467in}{0.862755in}}%
\pgfpathlineto{\pgfqpoint{1.299467in}{0.859806in}}%
\pgfpathmoveto{\pgfqpoint{1.294926in}{0.862755in}}%
\pgfpathlineto{\pgfqpoint{1.294926in}{0.862755in}}%
\pgfpathlineto{\pgfqpoint{1.294926in}{0.865704in}}%
\pgfpathlineto{\pgfqpoint{1.299467in}{0.865704in}}%
\pgfpathlineto{\pgfqpoint{1.299467in}{0.862755in}}%
\pgfpathmoveto{\pgfqpoint{1.299467in}{0.859806in}}%
\pgfpathlineto{\pgfqpoint{1.299467in}{0.859806in}}%
\pgfpathlineto{\pgfqpoint{1.299467in}{0.862755in}}%
\pgfpathlineto{\pgfqpoint{1.304008in}{0.862755in}}%
\pgfpathlineto{\pgfqpoint{1.304008in}{0.859806in}}%
\pgfpathmoveto{\pgfqpoint{1.299467in}{0.862755in}}%
\pgfpathlineto{\pgfqpoint{1.299467in}{0.862755in}}%
\pgfpathlineto{\pgfqpoint{1.299467in}{0.865704in}}%
\pgfpathlineto{\pgfqpoint{1.304008in}{0.865704in}}%
\pgfpathlineto{\pgfqpoint{1.304008in}{0.862755in}}%
\pgfpathmoveto{\pgfqpoint{1.304008in}{0.859806in}}%
\pgfpathlineto{\pgfqpoint{1.304008in}{0.859806in}}%
\pgfpathlineto{\pgfqpoint{1.304008in}{0.862755in}}%
\pgfpathlineto{\pgfqpoint{1.308549in}{0.862755in}}%
\pgfpathlineto{\pgfqpoint{1.308549in}{0.859806in}}%
\pgfpathmoveto{\pgfqpoint{1.304008in}{0.862755in}}%
\pgfpathlineto{\pgfqpoint{1.304008in}{0.862755in}}%
\pgfpathlineto{\pgfqpoint{1.304008in}{0.865704in}}%
\pgfpathlineto{\pgfqpoint{1.308549in}{0.865704in}}%
\pgfpathlineto{\pgfqpoint{1.308549in}{0.862755in}}%
\pgfpathmoveto{\pgfqpoint{1.308549in}{0.862755in}}%
\pgfpathlineto{\pgfqpoint{1.308549in}{0.862755in}}%
\pgfpathlineto{\pgfqpoint{1.308549in}{0.865704in}}%
\pgfpathlineto{\pgfqpoint{1.313090in}{0.865704in}}%
\pgfpathlineto{\pgfqpoint{1.313090in}{0.862755in}}%
\pgfpathmoveto{\pgfqpoint{1.304008in}{0.865704in}}%
\pgfpathlineto{\pgfqpoint{1.304008in}{0.865704in}}%
\pgfpathlineto{\pgfqpoint{1.304008in}{0.868653in}}%
\pgfpathlineto{\pgfqpoint{1.308549in}{0.868653in}}%
\pgfpathlineto{\pgfqpoint{1.308549in}{0.865704in}}%
\pgfpathmoveto{\pgfqpoint{1.304008in}{0.868653in}}%
\pgfpathlineto{\pgfqpoint{1.304008in}{0.868653in}}%
\pgfpathlineto{\pgfqpoint{1.304008in}{0.871603in}}%
\pgfpathlineto{\pgfqpoint{1.308549in}{0.871603in}}%
\pgfpathlineto{\pgfqpoint{1.308549in}{0.868653in}}%
\pgfpathmoveto{\pgfqpoint{1.308549in}{0.865704in}}%
\pgfpathlineto{\pgfqpoint{1.308549in}{0.865704in}}%
\pgfpathlineto{\pgfqpoint{1.308549in}{0.868653in}}%
\pgfpathlineto{\pgfqpoint{1.313090in}{0.868653in}}%
\pgfpathlineto{\pgfqpoint{1.313090in}{0.865704in}}%
\pgfpathmoveto{\pgfqpoint{1.308549in}{0.868653in}}%
\pgfpathlineto{\pgfqpoint{1.308549in}{0.868653in}}%
\pgfpathlineto{\pgfqpoint{1.308549in}{0.871603in}}%
\pgfpathlineto{\pgfqpoint{1.313090in}{0.871603in}}%
\pgfpathlineto{\pgfqpoint{1.313090in}{0.868653in}}%
\pgfpathmoveto{\pgfqpoint{1.313090in}{0.865704in}}%
\pgfpathlineto{\pgfqpoint{1.313090in}{0.865704in}}%
\pgfpathlineto{\pgfqpoint{1.313090in}{0.868653in}}%
\pgfpathlineto{\pgfqpoint{1.317631in}{0.868653in}}%
\pgfpathlineto{\pgfqpoint{1.317631in}{0.865704in}}%
\pgfpathmoveto{\pgfqpoint{1.313090in}{0.868653in}}%
\pgfpathlineto{\pgfqpoint{1.313090in}{0.868653in}}%
\pgfpathlineto{\pgfqpoint{1.313090in}{0.871603in}}%
\pgfpathlineto{\pgfqpoint{1.317631in}{0.871603in}}%
\pgfpathlineto{\pgfqpoint{1.317631in}{0.868653in}}%
\pgfpathmoveto{\pgfqpoint{1.317631in}{0.868653in}}%
\pgfpathlineto{\pgfqpoint{1.317631in}{0.868653in}}%
\pgfpathlineto{\pgfqpoint{1.317631in}{0.871603in}}%
\pgfpathlineto{\pgfqpoint{1.322172in}{0.871603in}}%
\pgfpathlineto{\pgfqpoint{1.322172in}{0.868653in}}%
\pgfpathmoveto{\pgfqpoint{1.313090in}{0.871603in}}%
\pgfpathlineto{\pgfqpoint{1.313090in}{0.871603in}}%
\pgfpathlineto{\pgfqpoint{1.313090in}{0.874552in}}%
\pgfpathlineto{\pgfqpoint{1.317631in}{0.874552in}}%
\pgfpathlineto{\pgfqpoint{1.317631in}{0.871603in}}%
\pgfpathmoveto{\pgfqpoint{1.313090in}{0.874552in}}%
\pgfpathlineto{\pgfqpoint{1.313090in}{0.874552in}}%
\pgfpathlineto{\pgfqpoint{1.313090in}{0.877501in}}%
\pgfpathlineto{\pgfqpoint{1.317631in}{0.877501in}}%
\pgfpathlineto{\pgfqpoint{1.317631in}{0.874552in}}%
\pgfpathmoveto{\pgfqpoint{1.317631in}{0.871603in}}%
\pgfpathlineto{\pgfqpoint{1.317631in}{0.871603in}}%
\pgfpathlineto{\pgfqpoint{1.317631in}{0.874552in}}%
\pgfpathlineto{\pgfqpoint{1.322172in}{0.874552in}}%
\pgfpathlineto{\pgfqpoint{1.322172in}{0.871603in}}%
\pgfpathmoveto{\pgfqpoint{1.317631in}{0.874552in}}%
\pgfpathlineto{\pgfqpoint{1.317631in}{0.874552in}}%
\pgfpathlineto{\pgfqpoint{1.317631in}{0.877501in}}%
\pgfpathlineto{\pgfqpoint{1.322172in}{0.877501in}}%
\pgfpathlineto{\pgfqpoint{1.322172in}{0.874552in}}%
\pgfpathmoveto{\pgfqpoint{1.322172in}{0.871603in}}%
\pgfpathlineto{\pgfqpoint{1.322172in}{0.871603in}}%
\pgfpathlineto{\pgfqpoint{1.322172in}{0.874552in}}%
\pgfpathlineto{\pgfqpoint{1.326713in}{0.874552in}}%
\pgfpathlineto{\pgfqpoint{1.326713in}{0.871603in}}%
\pgfpathmoveto{\pgfqpoint{1.322172in}{0.874552in}}%
\pgfpathlineto{\pgfqpoint{1.322172in}{0.874552in}}%
\pgfpathlineto{\pgfqpoint{1.322172in}{0.877501in}}%
\pgfpathlineto{\pgfqpoint{1.326713in}{0.877501in}}%
\pgfpathlineto{\pgfqpoint{1.326713in}{0.874552in}}%
\pgfpathmoveto{\pgfqpoint{1.326713in}{0.874552in}}%
\pgfpathlineto{\pgfqpoint{1.326713in}{0.874552in}}%
\pgfpathlineto{\pgfqpoint{1.326713in}{0.877501in}}%
\pgfpathlineto{\pgfqpoint{1.331254in}{0.877501in}}%
\pgfpathlineto{\pgfqpoint{1.331254in}{0.874552in}}%
\pgfpathmoveto{\pgfqpoint{1.322172in}{0.877501in}}%
\pgfpathlineto{\pgfqpoint{1.322172in}{0.877501in}}%
\pgfpathlineto{\pgfqpoint{1.322172in}{0.880450in}}%
\pgfpathlineto{\pgfqpoint{1.326713in}{0.880450in}}%
\pgfpathlineto{\pgfqpoint{1.326713in}{0.877501in}}%
\pgfpathmoveto{\pgfqpoint{1.322172in}{0.880450in}}%
\pgfpathlineto{\pgfqpoint{1.322172in}{0.880450in}}%
\pgfpathlineto{\pgfqpoint{1.322172in}{0.883400in}}%
\pgfpathlineto{\pgfqpoint{1.326713in}{0.883400in}}%
\pgfpathlineto{\pgfqpoint{1.326713in}{0.880450in}}%
\pgfpathmoveto{\pgfqpoint{1.326713in}{0.877501in}}%
\pgfpathlineto{\pgfqpoint{1.326713in}{0.877501in}}%
\pgfpathlineto{\pgfqpoint{1.326713in}{0.880450in}}%
\pgfpathlineto{\pgfqpoint{1.331254in}{0.880450in}}%
\pgfpathlineto{\pgfqpoint{1.331254in}{0.877501in}}%
\pgfpathmoveto{\pgfqpoint{1.326713in}{0.880450in}}%
\pgfpathlineto{\pgfqpoint{1.326713in}{0.880450in}}%
\pgfpathlineto{\pgfqpoint{1.326713in}{0.883400in}}%
\pgfpathlineto{\pgfqpoint{1.331254in}{0.883400in}}%
\pgfpathlineto{\pgfqpoint{1.331254in}{0.880450in}}%
\pgfpathmoveto{\pgfqpoint{1.331254in}{0.877501in}}%
\pgfpathlineto{\pgfqpoint{1.331254in}{0.877501in}}%
\pgfpathlineto{\pgfqpoint{1.331254in}{0.880450in}}%
\pgfpathlineto{\pgfqpoint{1.335795in}{0.880450in}}%
\pgfpathlineto{\pgfqpoint{1.335795in}{0.877501in}}%
\pgfpathmoveto{\pgfqpoint{1.331254in}{0.880450in}}%
\pgfpathlineto{\pgfqpoint{1.331254in}{0.880450in}}%
\pgfpathlineto{\pgfqpoint{1.331254in}{0.883400in}}%
\pgfpathlineto{\pgfqpoint{1.335795in}{0.883400in}}%
\pgfpathlineto{\pgfqpoint{1.335795in}{0.880450in}}%
\pgfpathmoveto{\pgfqpoint{1.335795in}{0.880450in}}%
\pgfpathlineto{\pgfqpoint{1.335795in}{0.880450in}}%
\pgfpathlineto{\pgfqpoint{1.335795in}{0.883400in}}%
\pgfpathlineto{\pgfqpoint{1.340336in}{0.883400in}}%
\pgfpathlineto{\pgfqpoint{1.340336in}{0.880450in}}%
\pgfpathmoveto{\pgfqpoint{1.331254in}{0.883400in}}%
\pgfpathlineto{\pgfqpoint{1.331254in}{0.883400in}}%
\pgfpathlineto{\pgfqpoint{1.331254in}{0.886349in}}%
\pgfpathlineto{\pgfqpoint{1.335795in}{0.886349in}}%
\pgfpathlineto{\pgfqpoint{1.335795in}{0.883400in}}%
\pgfpathmoveto{\pgfqpoint{1.331254in}{0.886349in}}%
\pgfpathlineto{\pgfqpoint{1.331254in}{0.886349in}}%
\pgfpathlineto{\pgfqpoint{1.331254in}{0.889298in}}%
\pgfpathlineto{\pgfqpoint{1.335795in}{0.889298in}}%
\pgfpathlineto{\pgfqpoint{1.335795in}{0.886349in}}%
\pgfpathmoveto{\pgfqpoint{1.335795in}{0.883400in}}%
\pgfpathlineto{\pgfqpoint{1.335795in}{0.883400in}}%
\pgfpathlineto{\pgfqpoint{1.335795in}{0.886349in}}%
\pgfpathlineto{\pgfqpoint{1.340336in}{0.886349in}}%
\pgfpathlineto{\pgfqpoint{1.340336in}{0.883400in}}%
\pgfpathmoveto{\pgfqpoint{1.335795in}{0.886349in}}%
\pgfpathlineto{\pgfqpoint{1.335795in}{0.886349in}}%
\pgfpathlineto{\pgfqpoint{1.335795in}{0.889298in}}%
\pgfpathlineto{\pgfqpoint{1.340336in}{0.889298in}}%
\pgfpathlineto{\pgfqpoint{1.340336in}{0.886349in}}%
\pgfpathmoveto{\pgfqpoint{1.340336in}{0.883400in}}%
\pgfpathlineto{\pgfqpoint{1.340336in}{0.883400in}}%
\pgfpathlineto{\pgfqpoint{1.340336in}{0.886349in}}%
\pgfpathlineto{\pgfqpoint{1.344877in}{0.886349in}}%
\pgfpathlineto{\pgfqpoint{1.344877in}{0.883400in}}%
\pgfpathmoveto{\pgfqpoint{1.340336in}{0.886349in}}%
\pgfpathlineto{\pgfqpoint{1.340336in}{0.886349in}}%
\pgfpathlineto{\pgfqpoint{1.340336in}{0.889298in}}%
\pgfpathlineto{\pgfqpoint{1.344877in}{0.889298in}}%
\pgfpathlineto{\pgfqpoint{1.344877in}{0.886349in}}%
\pgfpathmoveto{\pgfqpoint{1.344877in}{0.886349in}}%
\pgfpathlineto{\pgfqpoint{1.344877in}{0.886349in}}%
\pgfpathlineto{\pgfqpoint{1.344877in}{0.889298in}}%
\pgfpathlineto{\pgfqpoint{1.349418in}{0.889298in}}%
\pgfpathlineto{\pgfqpoint{1.349418in}{0.886349in}}%
\pgfpathmoveto{\pgfqpoint{1.340336in}{0.889298in}}%
\pgfpathlineto{\pgfqpoint{1.340336in}{0.889298in}}%
\pgfpathlineto{\pgfqpoint{1.340336in}{0.892247in}}%
\pgfpathlineto{\pgfqpoint{1.344877in}{0.892247in}}%
\pgfpathlineto{\pgfqpoint{1.344877in}{0.889298in}}%
\pgfpathmoveto{\pgfqpoint{1.340336in}{0.892247in}}%
\pgfpathlineto{\pgfqpoint{1.340336in}{0.892247in}}%
\pgfpathlineto{\pgfqpoint{1.340336in}{0.895196in}}%
\pgfpathlineto{\pgfqpoint{1.344877in}{0.895196in}}%
\pgfpathlineto{\pgfqpoint{1.344877in}{0.892247in}}%
\pgfpathmoveto{\pgfqpoint{1.344877in}{0.889298in}}%
\pgfpathlineto{\pgfqpoint{1.344877in}{0.889298in}}%
\pgfpathlineto{\pgfqpoint{1.344877in}{0.892247in}}%
\pgfpathlineto{\pgfqpoint{1.349418in}{0.892247in}}%
\pgfpathlineto{\pgfqpoint{1.349418in}{0.889298in}}%
\pgfpathmoveto{\pgfqpoint{1.344877in}{0.892247in}}%
\pgfpathlineto{\pgfqpoint{1.344877in}{0.892247in}}%
\pgfpathlineto{\pgfqpoint{1.344877in}{0.895196in}}%
\pgfpathlineto{\pgfqpoint{1.349418in}{0.895196in}}%
\pgfpathlineto{\pgfqpoint{1.349418in}{0.892247in}}%
\pgfpathmoveto{\pgfqpoint{1.349418in}{0.889298in}}%
\pgfpathlineto{\pgfqpoint{1.349418in}{0.889298in}}%
\pgfpathlineto{\pgfqpoint{1.349418in}{0.892247in}}%
\pgfpathlineto{\pgfqpoint{1.353959in}{0.892247in}}%
\pgfpathlineto{\pgfqpoint{1.353959in}{0.889298in}}%
\pgfpathmoveto{\pgfqpoint{1.349418in}{0.892247in}}%
\pgfpathlineto{\pgfqpoint{1.349418in}{0.892247in}}%
\pgfpathlineto{\pgfqpoint{1.349418in}{0.895196in}}%
\pgfpathlineto{\pgfqpoint{1.353959in}{0.895196in}}%
\pgfpathlineto{\pgfqpoint{1.353959in}{0.892247in}}%
\pgfpathmoveto{\pgfqpoint{1.353959in}{0.892247in}}%
\pgfpathlineto{\pgfqpoint{1.353959in}{0.892247in}}%
\pgfpathlineto{\pgfqpoint{1.353959in}{0.895196in}}%
\pgfpathlineto{\pgfqpoint{1.358500in}{0.895196in}}%
\pgfpathlineto{\pgfqpoint{1.358500in}{0.892247in}}%
\pgfpathmoveto{\pgfqpoint{1.349418in}{0.895196in}}%
\pgfpathlineto{\pgfqpoint{1.349418in}{0.895196in}}%
\pgfpathlineto{\pgfqpoint{1.349418in}{0.898145in}}%
\pgfpathlineto{\pgfqpoint{1.353959in}{0.898145in}}%
\pgfpathlineto{\pgfqpoint{1.353959in}{0.895196in}}%
\pgfpathmoveto{\pgfqpoint{1.349418in}{0.898145in}}%
\pgfpathlineto{\pgfqpoint{1.349418in}{0.898145in}}%
\pgfpathlineto{\pgfqpoint{1.349418in}{0.901094in}}%
\pgfpathlineto{\pgfqpoint{1.353959in}{0.901094in}}%
\pgfpathlineto{\pgfqpoint{1.353959in}{0.898145in}}%
\pgfpathmoveto{\pgfqpoint{1.353959in}{0.895196in}}%
\pgfpathlineto{\pgfqpoint{1.353959in}{0.895196in}}%
\pgfpathlineto{\pgfqpoint{1.353959in}{0.898145in}}%
\pgfpathlineto{\pgfqpoint{1.358500in}{0.898145in}}%
\pgfpathlineto{\pgfqpoint{1.358500in}{0.895196in}}%
\pgfpathmoveto{\pgfqpoint{1.353959in}{0.898145in}}%
\pgfpathlineto{\pgfqpoint{1.353959in}{0.898145in}}%
\pgfpathlineto{\pgfqpoint{1.353959in}{0.901094in}}%
\pgfpathlineto{\pgfqpoint{1.358500in}{0.901094in}}%
\pgfpathlineto{\pgfqpoint{1.358500in}{0.898145in}}%
\pgfpathmoveto{\pgfqpoint{1.358500in}{0.895196in}}%
\pgfpathlineto{\pgfqpoint{1.358500in}{0.895196in}}%
\pgfpathlineto{\pgfqpoint{1.358500in}{0.898145in}}%
\pgfpathlineto{\pgfqpoint{1.363041in}{0.898145in}}%
\pgfpathlineto{\pgfqpoint{1.363041in}{0.895196in}}%
\pgfpathmoveto{\pgfqpoint{1.358500in}{0.898145in}}%
\pgfpathlineto{\pgfqpoint{1.358500in}{0.898145in}}%
\pgfpathlineto{\pgfqpoint{1.358500in}{0.901094in}}%
\pgfpathlineto{\pgfqpoint{1.363041in}{0.901094in}}%
\pgfpathlineto{\pgfqpoint{1.363041in}{0.898145in}}%
\pgfpathmoveto{\pgfqpoint{1.363041in}{0.898145in}}%
\pgfpathlineto{\pgfqpoint{1.363041in}{0.898145in}}%
\pgfpathlineto{\pgfqpoint{1.363041in}{0.901094in}}%
\pgfpathlineto{\pgfqpoint{1.367581in}{0.901094in}}%
\pgfpathlineto{\pgfqpoint{1.367581in}{0.898145in}}%
\pgfpathmoveto{\pgfqpoint{1.358500in}{0.901094in}}%
\pgfpathlineto{\pgfqpoint{1.358500in}{0.901094in}}%
\pgfpathlineto{\pgfqpoint{1.358500in}{0.904043in}}%
\pgfpathlineto{\pgfqpoint{1.363041in}{0.904043in}}%
\pgfpathlineto{\pgfqpoint{1.363041in}{0.901094in}}%
\pgfpathmoveto{\pgfqpoint{1.358500in}{0.904043in}}%
\pgfpathlineto{\pgfqpoint{1.358500in}{0.904043in}}%
\pgfpathlineto{\pgfqpoint{1.358500in}{0.906992in}}%
\pgfpathlineto{\pgfqpoint{1.363041in}{0.906992in}}%
\pgfpathlineto{\pgfqpoint{1.363041in}{0.904043in}}%
\pgfpathmoveto{\pgfqpoint{1.363041in}{0.901094in}}%
\pgfpathlineto{\pgfqpoint{1.363041in}{0.901094in}}%
\pgfpathlineto{\pgfqpoint{1.363041in}{0.904043in}}%
\pgfpathlineto{\pgfqpoint{1.367581in}{0.904043in}}%
\pgfpathlineto{\pgfqpoint{1.367581in}{0.901094in}}%
\pgfpathmoveto{\pgfqpoint{1.363041in}{0.904043in}}%
\pgfpathlineto{\pgfqpoint{1.363041in}{0.904043in}}%
\pgfpathlineto{\pgfqpoint{1.363041in}{0.906992in}}%
\pgfpathlineto{\pgfqpoint{1.367581in}{0.906992in}}%
\pgfpathlineto{\pgfqpoint{1.367581in}{0.904043in}}%
\pgfpathmoveto{\pgfqpoint{1.367581in}{0.901094in}}%
\pgfpathlineto{\pgfqpoint{1.367581in}{0.901094in}}%
\pgfpathlineto{\pgfqpoint{1.367581in}{0.904043in}}%
\pgfpathlineto{\pgfqpoint{1.372122in}{0.904043in}}%
\pgfpathlineto{\pgfqpoint{1.372122in}{0.901094in}}%
\pgfpathmoveto{\pgfqpoint{1.367581in}{0.904043in}}%
\pgfpathlineto{\pgfqpoint{1.367581in}{0.904043in}}%
\pgfpathlineto{\pgfqpoint{1.367581in}{0.906992in}}%
\pgfpathlineto{\pgfqpoint{1.372122in}{0.906992in}}%
\pgfpathlineto{\pgfqpoint{1.372122in}{0.904043in}}%
\pgfpathmoveto{\pgfqpoint{1.372122in}{0.904043in}}%
\pgfpathlineto{\pgfqpoint{1.372122in}{0.904043in}}%
\pgfpathlineto{\pgfqpoint{1.372122in}{0.906992in}}%
\pgfpathlineto{\pgfqpoint{1.376663in}{0.906992in}}%
\pgfpathlineto{\pgfqpoint{1.376663in}{0.904043in}}%
\pgfpathmoveto{\pgfqpoint{1.367581in}{0.906992in}}%
\pgfpathlineto{\pgfqpoint{1.367581in}{0.906992in}}%
\pgfpathlineto{\pgfqpoint{1.367581in}{0.909941in}}%
\pgfpathlineto{\pgfqpoint{1.372122in}{0.909941in}}%
\pgfpathlineto{\pgfqpoint{1.372122in}{0.906992in}}%
\pgfpathmoveto{\pgfqpoint{1.367581in}{0.909941in}}%
\pgfpathlineto{\pgfqpoint{1.367581in}{0.909941in}}%
\pgfpathlineto{\pgfqpoint{1.367581in}{0.912890in}}%
\pgfpathlineto{\pgfqpoint{1.372122in}{0.912890in}}%
\pgfpathlineto{\pgfqpoint{1.372122in}{0.909941in}}%
\pgfpathmoveto{\pgfqpoint{1.372122in}{0.906992in}}%
\pgfpathlineto{\pgfqpoint{1.372122in}{0.906992in}}%
\pgfpathlineto{\pgfqpoint{1.372122in}{0.909941in}}%
\pgfpathlineto{\pgfqpoint{1.376663in}{0.909941in}}%
\pgfpathlineto{\pgfqpoint{1.376663in}{0.906992in}}%
\pgfpathmoveto{\pgfqpoint{1.372122in}{0.909941in}}%
\pgfpathlineto{\pgfqpoint{1.372122in}{0.909941in}}%
\pgfpathlineto{\pgfqpoint{1.372122in}{0.912890in}}%
\pgfpathlineto{\pgfqpoint{1.376663in}{0.912890in}}%
\pgfpathlineto{\pgfqpoint{1.376663in}{0.909941in}}%
\pgfpathmoveto{\pgfqpoint{1.376663in}{0.906992in}}%
\pgfpathlineto{\pgfqpoint{1.376663in}{0.906992in}}%
\pgfpathlineto{\pgfqpoint{1.376663in}{0.909941in}}%
\pgfpathlineto{\pgfqpoint{1.381204in}{0.909941in}}%
\pgfpathlineto{\pgfqpoint{1.381204in}{0.906992in}}%
\pgfpathmoveto{\pgfqpoint{1.376663in}{0.909941in}}%
\pgfpathlineto{\pgfqpoint{1.376663in}{0.909941in}}%
\pgfpathlineto{\pgfqpoint{1.376663in}{0.912890in}}%
\pgfpathlineto{\pgfqpoint{1.381204in}{0.912890in}}%
\pgfpathlineto{\pgfqpoint{1.381204in}{0.909941in}}%
\pgfpathmoveto{\pgfqpoint{1.381204in}{0.909941in}}%
\pgfpathlineto{\pgfqpoint{1.381204in}{0.909941in}}%
\pgfpathlineto{\pgfqpoint{1.381204in}{0.912890in}}%
\pgfpathlineto{\pgfqpoint{1.385745in}{0.912890in}}%
\pgfpathlineto{\pgfqpoint{1.385745in}{0.909941in}}%
\pgfpathmoveto{\pgfqpoint{1.376663in}{0.912890in}}%
\pgfpathlineto{\pgfqpoint{1.376663in}{0.912890in}}%
\pgfpathlineto{\pgfqpoint{1.376663in}{0.915840in}}%
\pgfpathlineto{\pgfqpoint{1.381204in}{0.915840in}}%
\pgfpathlineto{\pgfqpoint{1.381204in}{0.912890in}}%
\pgfpathmoveto{\pgfqpoint{1.376663in}{0.915840in}}%
\pgfpathlineto{\pgfqpoint{1.376663in}{0.915840in}}%
\pgfpathlineto{\pgfqpoint{1.376663in}{0.918789in}}%
\pgfpathlineto{\pgfqpoint{1.381204in}{0.918789in}}%
\pgfpathlineto{\pgfqpoint{1.381204in}{0.915840in}}%
\pgfpathmoveto{\pgfqpoint{1.381204in}{0.912890in}}%
\pgfpathlineto{\pgfqpoint{1.381204in}{0.912890in}}%
\pgfpathlineto{\pgfqpoint{1.381204in}{0.915840in}}%
\pgfpathlineto{\pgfqpoint{1.385745in}{0.915840in}}%
\pgfpathlineto{\pgfqpoint{1.385745in}{0.912890in}}%
\pgfpathmoveto{\pgfqpoint{1.381204in}{0.915840in}}%
\pgfpathlineto{\pgfqpoint{1.381204in}{0.915840in}}%
\pgfpathlineto{\pgfqpoint{1.381204in}{0.918789in}}%
\pgfpathlineto{\pgfqpoint{1.385745in}{0.918789in}}%
\pgfpathlineto{\pgfqpoint{1.385745in}{0.915840in}}%
\pgfpathmoveto{\pgfqpoint{1.385745in}{0.912890in}}%
\pgfpathlineto{\pgfqpoint{1.385745in}{0.912890in}}%
\pgfpathlineto{\pgfqpoint{1.385745in}{0.915840in}}%
\pgfpathlineto{\pgfqpoint{1.390286in}{0.915840in}}%
\pgfpathlineto{\pgfqpoint{1.390286in}{0.912890in}}%
\pgfpathmoveto{\pgfqpoint{1.385745in}{0.915840in}}%
\pgfpathlineto{\pgfqpoint{1.385745in}{0.915840in}}%
\pgfpathlineto{\pgfqpoint{1.385745in}{0.918789in}}%
\pgfpathlineto{\pgfqpoint{1.390286in}{0.918789in}}%
\pgfpathlineto{\pgfqpoint{1.390286in}{0.915840in}}%
\pgfpathmoveto{\pgfqpoint{1.390286in}{0.915840in}}%
\pgfpathlineto{\pgfqpoint{1.390286in}{0.915840in}}%
\pgfpathlineto{\pgfqpoint{1.390286in}{0.918789in}}%
\pgfpathlineto{\pgfqpoint{1.394827in}{0.918789in}}%
\pgfpathlineto{\pgfqpoint{1.394827in}{0.915840in}}%
\pgfpathmoveto{\pgfqpoint{1.385745in}{0.918789in}}%
\pgfpathlineto{\pgfqpoint{1.385745in}{0.918789in}}%
\pgfpathlineto{\pgfqpoint{1.385745in}{0.921738in}}%
\pgfpathlineto{\pgfqpoint{1.390286in}{0.921738in}}%
\pgfpathlineto{\pgfqpoint{1.390286in}{0.918789in}}%
\pgfpathmoveto{\pgfqpoint{1.385745in}{0.921738in}}%
\pgfpathlineto{\pgfqpoint{1.385745in}{0.921738in}}%
\pgfpathlineto{\pgfqpoint{1.385745in}{0.924687in}}%
\pgfpathlineto{\pgfqpoint{1.390286in}{0.924687in}}%
\pgfpathlineto{\pgfqpoint{1.390286in}{0.921738in}}%
\pgfpathmoveto{\pgfqpoint{1.390286in}{0.918789in}}%
\pgfpathlineto{\pgfqpoint{1.390286in}{0.918789in}}%
\pgfpathlineto{\pgfqpoint{1.390286in}{0.921738in}}%
\pgfpathlineto{\pgfqpoint{1.394827in}{0.921738in}}%
\pgfpathlineto{\pgfqpoint{1.394827in}{0.918789in}}%
\pgfpathmoveto{\pgfqpoint{1.390286in}{0.921738in}}%
\pgfpathlineto{\pgfqpoint{1.390286in}{0.921738in}}%
\pgfpathlineto{\pgfqpoint{1.390286in}{0.924687in}}%
\pgfpathlineto{\pgfqpoint{1.394827in}{0.924687in}}%
\pgfpathlineto{\pgfqpoint{1.394827in}{0.921738in}}%
\pgfpathmoveto{\pgfqpoint{1.394827in}{0.918789in}}%
\pgfpathlineto{\pgfqpoint{1.394827in}{0.918789in}}%
\pgfpathlineto{\pgfqpoint{1.394827in}{0.921738in}}%
\pgfpathlineto{\pgfqpoint{1.399368in}{0.921738in}}%
\pgfpathlineto{\pgfqpoint{1.399368in}{0.918789in}}%
\pgfpathmoveto{\pgfqpoint{1.394827in}{0.921738in}}%
\pgfpathlineto{\pgfqpoint{1.394827in}{0.921738in}}%
\pgfpathlineto{\pgfqpoint{1.394827in}{0.924687in}}%
\pgfpathlineto{\pgfqpoint{1.399368in}{0.924687in}}%
\pgfpathlineto{\pgfqpoint{1.399368in}{0.921738in}}%
\pgfpathmoveto{\pgfqpoint{1.399368in}{0.921738in}}%
\pgfpathlineto{\pgfqpoint{1.399368in}{0.921738in}}%
\pgfpathlineto{\pgfqpoint{1.399368in}{0.924687in}}%
\pgfpathlineto{\pgfqpoint{1.403909in}{0.924687in}}%
\pgfpathlineto{\pgfqpoint{1.403909in}{0.921738in}}%
\pgfpathmoveto{\pgfqpoint{1.394827in}{0.924687in}}%
\pgfpathlineto{\pgfqpoint{1.394827in}{0.924687in}}%
\pgfpathlineto{\pgfqpoint{1.394827in}{0.927636in}}%
\pgfpathlineto{\pgfqpoint{1.399368in}{0.927636in}}%
\pgfpathlineto{\pgfqpoint{1.399368in}{0.924687in}}%
\pgfpathmoveto{\pgfqpoint{1.394827in}{0.927636in}}%
\pgfpathlineto{\pgfqpoint{1.394827in}{0.927636in}}%
\pgfpathlineto{\pgfqpoint{1.394827in}{0.930585in}}%
\pgfpathlineto{\pgfqpoint{1.399368in}{0.930585in}}%
\pgfpathlineto{\pgfqpoint{1.399368in}{0.927636in}}%
\pgfpathmoveto{\pgfqpoint{1.399368in}{0.924687in}}%
\pgfpathlineto{\pgfqpoint{1.399368in}{0.924687in}}%
\pgfpathlineto{\pgfqpoint{1.399368in}{0.927636in}}%
\pgfpathlineto{\pgfqpoint{1.403909in}{0.927636in}}%
\pgfpathlineto{\pgfqpoint{1.403909in}{0.924687in}}%
\pgfpathmoveto{\pgfqpoint{1.399368in}{0.927636in}}%
\pgfpathlineto{\pgfqpoint{1.399368in}{0.927636in}}%
\pgfpathlineto{\pgfqpoint{1.399368in}{0.930585in}}%
\pgfpathlineto{\pgfqpoint{1.403909in}{0.930585in}}%
\pgfpathlineto{\pgfqpoint{1.403909in}{0.927636in}}%
\pgfpathmoveto{\pgfqpoint{1.403909in}{0.924687in}}%
\pgfpathlineto{\pgfqpoint{1.403909in}{0.924687in}}%
\pgfpathlineto{\pgfqpoint{1.403909in}{0.927636in}}%
\pgfpathlineto{\pgfqpoint{1.408450in}{0.927636in}}%
\pgfpathlineto{\pgfqpoint{1.408450in}{0.924687in}}%
\pgfpathmoveto{\pgfqpoint{1.403909in}{0.927636in}}%
\pgfpathlineto{\pgfqpoint{1.403909in}{0.927636in}}%
\pgfpathlineto{\pgfqpoint{1.403909in}{0.930585in}}%
\pgfpathlineto{\pgfqpoint{1.408450in}{0.930585in}}%
\pgfpathlineto{\pgfqpoint{1.408450in}{0.927636in}}%
\pgfpathmoveto{\pgfqpoint{1.408450in}{0.927636in}}%
\pgfpathlineto{\pgfqpoint{1.408450in}{0.927636in}}%
\pgfpathlineto{\pgfqpoint{1.408450in}{0.930585in}}%
\pgfpathlineto{\pgfqpoint{1.412991in}{0.930585in}}%
\pgfpathlineto{\pgfqpoint{1.412991in}{0.927636in}}%
\pgfpathmoveto{\pgfqpoint{1.403909in}{0.930585in}}%
\pgfpathlineto{\pgfqpoint{1.403909in}{0.930585in}}%
\pgfpathlineto{\pgfqpoint{1.403909in}{0.933534in}}%
\pgfpathlineto{\pgfqpoint{1.408450in}{0.933534in}}%
\pgfpathlineto{\pgfqpoint{1.408450in}{0.930585in}}%
\pgfpathmoveto{\pgfqpoint{1.403909in}{0.933534in}}%
\pgfpathlineto{\pgfqpoint{1.403909in}{0.933534in}}%
\pgfpathlineto{\pgfqpoint{1.403909in}{0.936483in}}%
\pgfpathlineto{\pgfqpoint{1.408450in}{0.936483in}}%
\pgfpathlineto{\pgfqpoint{1.408450in}{0.933534in}}%
\pgfpathmoveto{\pgfqpoint{1.408450in}{0.930585in}}%
\pgfpathlineto{\pgfqpoint{1.408450in}{0.930585in}}%
\pgfpathlineto{\pgfqpoint{1.408450in}{0.933534in}}%
\pgfpathlineto{\pgfqpoint{1.412991in}{0.933534in}}%
\pgfpathlineto{\pgfqpoint{1.412991in}{0.930585in}}%
\pgfpathmoveto{\pgfqpoint{1.408450in}{0.933534in}}%
\pgfpathlineto{\pgfqpoint{1.408450in}{0.933534in}}%
\pgfpathlineto{\pgfqpoint{1.408450in}{0.936483in}}%
\pgfpathlineto{\pgfqpoint{1.412991in}{0.936483in}}%
\pgfpathlineto{\pgfqpoint{1.412991in}{0.933534in}}%
\pgfpathmoveto{\pgfqpoint{1.412991in}{0.930585in}}%
\pgfpathlineto{\pgfqpoint{1.412991in}{0.930585in}}%
\pgfpathlineto{\pgfqpoint{1.412991in}{0.933534in}}%
\pgfpathlineto{\pgfqpoint{1.417532in}{0.933534in}}%
\pgfpathlineto{\pgfqpoint{1.417532in}{0.930585in}}%
\pgfpathmoveto{\pgfqpoint{1.412991in}{0.933534in}}%
\pgfpathlineto{\pgfqpoint{1.412991in}{0.933534in}}%
\pgfpathlineto{\pgfqpoint{1.412991in}{0.936483in}}%
\pgfpathlineto{\pgfqpoint{1.417532in}{0.936483in}}%
\pgfpathlineto{\pgfqpoint{1.417532in}{0.933534in}}%
\pgfpathmoveto{\pgfqpoint{1.417532in}{0.933534in}}%
\pgfpathlineto{\pgfqpoint{1.417532in}{0.933534in}}%
\pgfpathlineto{\pgfqpoint{1.417532in}{0.936483in}}%
\pgfpathlineto{\pgfqpoint{1.422073in}{0.936483in}}%
\pgfpathlineto{\pgfqpoint{1.422073in}{0.933534in}}%
\pgfpathmoveto{\pgfqpoint{1.412991in}{0.936483in}}%
\pgfpathlineto{\pgfqpoint{1.412991in}{0.936483in}}%
\pgfpathlineto{\pgfqpoint{1.412991in}{0.939432in}}%
\pgfpathlineto{\pgfqpoint{1.417532in}{0.939432in}}%
\pgfpathlineto{\pgfqpoint{1.417532in}{0.936483in}}%
\pgfpathmoveto{\pgfqpoint{1.412991in}{0.939432in}}%
\pgfpathlineto{\pgfqpoint{1.412991in}{0.939432in}}%
\pgfpathlineto{\pgfqpoint{1.412991in}{0.942381in}}%
\pgfpathlineto{\pgfqpoint{1.417532in}{0.942381in}}%
\pgfpathlineto{\pgfqpoint{1.417532in}{0.939432in}}%
\pgfpathmoveto{\pgfqpoint{1.417532in}{0.936483in}}%
\pgfpathlineto{\pgfqpoint{1.417532in}{0.936483in}}%
\pgfpathlineto{\pgfqpoint{1.417532in}{0.939432in}}%
\pgfpathlineto{\pgfqpoint{1.422073in}{0.939432in}}%
\pgfpathlineto{\pgfqpoint{1.422073in}{0.936483in}}%
\pgfpathmoveto{\pgfqpoint{1.417532in}{0.939432in}}%
\pgfpathlineto{\pgfqpoint{1.417532in}{0.939432in}}%
\pgfpathlineto{\pgfqpoint{1.417532in}{0.942381in}}%
\pgfpathlineto{\pgfqpoint{1.422073in}{0.942381in}}%
\pgfpathlineto{\pgfqpoint{1.422073in}{0.939432in}}%
\pgfpathmoveto{\pgfqpoint{1.422073in}{0.936483in}}%
\pgfpathlineto{\pgfqpoint{1.422073in}{0.936483in}}%
\pgfpathlineto{\pgfqpoint{1.422073in}{0.939432in}}%
\pgfpathlineto{\pgfqpoint{1.426614in}{0.939432in}}%
\pgfpathlineto{\pgfqpoint{1.426614in}{0.936483in}}%
\pgfpathmoveto{\pgfqpoint{1.422073in}{0.939432in}}%
\pgfpathlineto{\pgfqpoint{1.422073in}{0.939432in}}%
\pgfpathlineto{\pgfqpoint{1.422073in}{0.942381in}}%
\pgfpathlineto{\pgfqpoint{1.426614in}{0.942381in}}%
\pgfpathlineto{\pgfqpoint{1.426614in}{0.939432in}}%
\pgfpathmoveto{\pgfqpoint{1.426614in}{0.939432in}}%
\pgfpathlineto{\pgfqpoint{1.426614in}{0.939432in}}%
\pgfpathlineto{\pgfqpoint{1.426614in}{0.942381in}}%
\pgfpathlineto{\pgfqpoint{1.431155in}{0.942381in}}%
\pgfpathlineto{\pgfqpoint{1.431155in}{0.939432in}}%
\pgfpathmoveto{\pgfqpoint{1.422073in}{0.942381in}}%
\pgfpathlineto{\pgfqpoint{1.422073in}{0.942381in}}%
\pgfpathlineto{\pgfqpoint{1.422073in}{0.945330in}}%
\pgfpathlineto{\pgfqpoint{1.426614in}{0.945330in}}%
\pgfpathlineto{\pgfqpoint{1.426614in}{0.942381in}}%
\pgfpathmoveto{\pgfqpoint{1.422073in}{0.945330in}}%
\pgfpathlineto{\pgfqpoint{1.422073in}{0.945330in}}%
\pgfpathlineto{\pgfqpoint{1.422073in}{0.948280in}}%
\pgfpathlineto{\pgfqpoint{1.426614in}{0.948280in}}%
\pgfpathlineto{\pgfqpoint{1.426614in}{0.945330in}}%
\pgfpathmoveto{\pgfqpoint{1.426614in}{0.942381in}}%
\pgfpathlineto{\pgfqpoint{1.426614in}{0.942381in}}%
\pgfpathlineto{\pgfqpoint{1.426614in}{0.945330in}}%
\pgfpathlineto{\pgfqpoint{1.431155in}{0.945330in}}%
\pgfpathlineto{\pgfqpoint{1.431155in}{0.942381in}}%
\pgfpathmoveto{\pgfqpoint{1.426614in}{0.945330in}}%
\pgfpathlineto{\pgfqpoint{1.426614in}{0.945330in}}%
\pgfpathlineto{\pgfqpoint{1.426614in}{0.948280in}}%
\pgfpathlineto{\pgfqpoint{1.431155in}{0.948280in}}%
\pgfpathlineto{\pgfqpoint{1.431155in}{0.945330in}}%
\pgfpathmoveto{\pgfqpoint{1.431155in}{0.942381in}}%
\pgfpathlineto{\pgfqpoint{1.431155in}{0.942381in}}%
\pgfpathlineto{\pgfqpoint{1.431155in}{0.945330in}}%
\pgfpathlineto{\pgfqpoint{1.435696in}{0.945330in}}%
\pgfpathlineto{\pgfqpoint{1.435696in}{0.942381in}}%
\pgfpathmoveto{\pgfqpoint{1.431155in}{0.945330in}}%
\pgfpathlineto{\pgfqpoint{1.431155in}{0.945330in}}%
\pgfpathlineto{\pgfqpoint{1.431155in}{0.948280in}}%
\pgfpathlineto{\pgfqpoint{1.435696in}{0.948280in}}%
\pgfpathlineto{\pgfqpoint{1.435696in}{0.945330in}}%
\pgfpathmoveto{\pgfqpoint{1.435696in}{0.945330in}}%
\pgfpathlineto{\pgfqpoint{1.435696in}{0.945330in}}%
\pgfpathlineto{\pgfqpoint{1.435696in}{0.948280in}}%
\pgfpathlineto{\pgfqpoint{1.440237in}{0.948280in}}%
\pgfpathlineto{\pgfqpoint{1.440237in}{0.945330in}}%
\pgfpathmoveto{\pgfqpoint{1.431155in}{0.948280in}}%
\pgfpathlineto{\pgfqpoint{1.431155in}{0.948280in}}%
\pgfpathlineto{\pgfqpoint{1.431155in}{0.951229in}}%
\pgfpathlineto{\pgfqpoint{1.435696in}{0.951229in}}%
\pgfpathlineto{\pgfqpoint{1.435696in}{0.948280in}}%
\pgfpathmoveto{\pgfqpoint{1.431155in}{0.951229in}}%
\pgfpathlineto{\pgfqpoint{1.431155in}{0.951229in}}%
\pgfpathlineto{\pgfqpoint{1.431155in}{0.954178in}}%
\pgfpathlineto{\pgfqpoint{1.435696in}{0.954178in}}%
\pgfpathlineto{\pgfqpoint{1.435696in}{0.951229in}}%
\pgfpathmoveto{\pgfqpoint{1.435696in}{0.948280in}}%
\pgfpathlineto{\pgfqpoint{1.435696in}{0.948280in}}%
\pgfpathlineto{\pgfqpoint{1.435696in}{0.951229in}}%
\pgfpathlineto{\pgfqpoint{1.440237in}{0.951229in}}%
\pgfpathlineto{\pgfqpoint{1.440237in}{0.948280in}}%
\pgfpathmoveto{\pgfqpoint{1.435696in}{0.951229in}}%
\pgfpathlineto{\pgfqpoint{1.435696in}{0.951229in}}%
\pgfpathlineto{\pgfqpoint{1.435696in}{0.954178in}}%
\pgfpathlineto{\pgfqpoint{1.440237in}{0.954178in}}%
\pgfpathlineto{\pgfqpoint{1.440237in}{0.951229in}}%
\pgfpathmoveto{\pgfqpoint{1.440237in}{0.948280in}}%
\pgfpathlineto{\pgfqpoint{1.440237in}{0.948280in}}%
\pgfpathlineto{\pgfqpoint{1.440237in}{0.951229in}}%
\pgfpathlineto{\pgfqpoint{1.444777in}{0.951229in}}%
\pgfpathlineto{\pgfqpoint{1.444777in}{0.948280in}}%
\pgfpathmoveto{\pgfqpoint{1.440237in}{0.951229in}}%
\pgfpathlineto{\pgfqpoint{1.440237in}{0.951229in}}%
\pgfpathlineto{\pgfqpoint{1.440237in}{0.954178in}}%
\pgfpathlineto{\pgfqpoint{1.444777in}{0.954178in}}%
\pgfpathlineto{\pgfqpoint{1.444777in}{0.951229in}}%
\pgfpathmoveto{\pgfqpoint{1.444777in}{0.951229in}}%
\pgfpathlineto{\pgfqpoint{1.444777in}{0.951229in}}%
\pgfpathlineto{\pgfqpoint{1.444777in}{0.954178in}}%
\pgfpathlineto{\pgfqpoint{1.449318in}{0.954178in}}%
\pgfpathlineto{\pgfqpoint{1.449318in}{0.951229in}}%
\pgfpathmoveto{\pgfqpoint{1.440237in}{0.954178in}}%
\pgfpathlineto{\pgfqpoint{1.440237in}{0.954178in}}%
\pgfpathlineto{\pgfqpoint{1.440237in}{0.957127in}}%
\pgfpathlineto{\pgfqpoint{1.444777in}{0.957127in}}%
\pgfpathlineto{\pgfqpoint{1.444777in}{0.954178in}}%
\pgfpathmoveto{\pgfqpoint{1.440237in}{0.957127in}}%
\pgfpathlineto{\pgfqpoint{1.440237in}{0.957127in}}%
\pgfpathlineto{\pgfqpoint{1.440237in}{0.960076in}}%
\pgfpathlineto{\pgfqpoint{1.444777in}{0.960076in}}%
\pgfpathlineto{\pgfqpoint{1.444777in}{0.957127in}}%
\pgfpathmoveto{\pgfqpoint{1.444777in}{0.954178in}}%
\pgfpathlineto{\pgfqpoint{1.444777in}{0.954178in}}%
\pgfpathlineto{\pgfqpoint{1.444777in}{0.957127in}}%
\pgfpathlineto{\pgfqpoint{1.449318in}{0.957127in}}%
\pgfpathlineto{\pgfqpoint{1.449318in}{0.954178in}}%
\pgfpathmoveto{\pgfqpoint{1.444777in}{0.957127in}}%
\pgfpathlineto{\pgfqpoint{1.444777in}{0.957127in}}%
\pgfpathlineto{\pgfqpoint{1.444777in}{0.960076in}}%
\pgfpathlineto{\pgfqpoint{1.449318in}{0.960076in}}%
\pgfpathlineto{\pgfqpoint{1.449318in}{0.957127in}}%
\pgfpathmoveto{\pgfqpoint{1.449318in}{0.954178in}}%
\pgfpathlineto{\pgfqpoint{1.449318in}{0.954178in}}%
\pgfpathlineto{\pgfqpoint{1.449318in}{0.957127in}}%
\pgfpathlineto{\pgfqpoint{1.453859in}{0.957127in}}%
\pgfpathlineto{\pgfqpoint{1.453859in}{0.954178in}}%
\pgfpathmoveto{\pgfqpoint{1.449318in}{0.957127in}}%
\pgfpathlineto{\pgfqpoint{1.449318in}{0.957127in}}%
\pgfpathlineto{\pgfqpoint{1.449318in}{0.960076in}}%
\pgfpathlineto{\pgfqpoint{1.453859in}{0.960076in}}%
\pgfpathlineto{\pgfqpoint{1.453859in}{0.957127in}}%
\pgfpathmoveto{\pgfqpoint{1.453859in}{0.957127in}}%
\pgfpathlineto{\pgfqpoint{1.453859in}{0.957127in}}%
\pgfpathlineto{\pgfqpoint{1.453859in}{0.960076in}}%
\pgfpathlineto{\pgfqpoint{1.458400in}{0.960076in}}%
\pgfpathlineto{\pgfqpoint{1.458400in}{0.957127in}}%
\pgfpathmoveto{\pgfqpoint{1.449318in}{0.960076in}}%
\pgfpathlineto{\pgfqpoint{1.449318in}{0.960076in}}%
\pgfpathlineto{\pgfqpoint{1.449318in}{0.963025in}}%
\pgfpathlineto{\pgfqpoint{1.453859in}{0.963025in}}%
\pgfpathlineto{\pgfqpoint{1.453859in}{0.960076in}}%
\pgfpathmoveto{\pgfqpoint{1.449318in}{0.963025in}}%
\pgfpathlineto{\pgfqpoint{1.449318in}{0.963025in}}%
\pgfpathlineto{\pgfqpoint{1.449318in}{0.965974in}}%
\pgfpathlineto{\pgfqpoint{1.453859in}{0.965974in}}%
\pgfpathlineto{\pgfqpoint{1.453859in}{0.963025in}}%
\pgfpathmoveto{\pgfqpoint{1.453859in}{0.960076in}}%
\pgfpathlineto{\pgfqpoint{1.453859in}{0.960076in}}%
\pgfpathlineto{\pgfqpoint{1.453859in}{0.963025in}}%
\pgfpathlineto{\pgfqpoint{1.458400in}{0.963025in}}%
\pgfpathlineto{\pgfqpoint{1.458400in}{0.960076in}}%
\pgfpathmoveto{\pgfqpoint{1.453859in}{0.963025in}}%
\pgfpathlineto{\pgfqpoint{1.453859in}{0.963025in}}%
\pgfpathlineto{\pgfqpoint{1.453859in}{0.965974in}}%
\pgfpathlineto{\pgfqpoint{1.458400in}{0.965974in}}%
\pgfpathlineto{\pgfqpoint{1.458400in}{0.963025in}}%
\pgfpathmoveto{\pgfqpoint{1.458400in}{0.960076in}}%
\pgfpathlineto{\pgfqpoint{1.458400in}{0.960076in}}%
\pgfpathlineto{\pgfqpoint{1.458400in}{0.963025in}}%
\pgfpathlineto{\pgfqpoint{1.462941in}{0.963025in}}%
\pgfpathlineto{\pgfqpoint{1.462941in}{0.960076in}}%
\pgfpathmoveto{\pgfqpoint{1.458400in}{0.963025in}}%
\pgfpathlineto{\pgfqpoint{1.458400in}{0.963025in}}%
\pgfpathlineto{\pgfqpoint{1.458400in}{0.965974in}}%
\pgfpathlineto{\pgfqpoint{1.462941in}{0.965974in}}%
\pgfpathlineto{\pgfqpoint{1.462941in}{0.963025in}}%
\pgfpathmoveto{\pgfqpoint{1.462941in}{0.963025in}}%
\pgfpathlineto{\pgfqpoint{1.462941in}{0.963025in}}%
\pgfpathlineto{\pgfqpoint{1.462941in}{0.965974in}}%
\pgfpathlineto{\pgfqpoint{1.467482in}{0.965974in}}%
\pgfpathlineto{\pgfqpoint{1.467482in}{0.963025in}}%
\pgfpathmoveto{\pgfqpoint{1.458400in}{0.965974in}}%
\pgfpathlineto{\pgfqpoint{1.458400in}{0.965974in}}%
\pgfpathlineto{\pgfqpoint{1.458400in}{0.968923in}}%
\pgfpathlineto{\pgfqpoint{1.462941in}{0.968923in}}%
\pgfpathlineto{\pgfqpoint{1.462941in}{0.965974in}}%
\pgfpathmoveto{\pgfqpoint{1.458400in}{0.968923in}}%
\pgfpathlineto{\pgfqpoint{1.458400in}{0.968923in}}%
\pgfpathlineto{\pgfqpoint{1.458400in}{0.971872in}}%
\pgfpathlineto{\pgfqpoint{1.462941in}{0.971872in}}%
\pgfpathlineto{\pgfqpoint{1.462941in}{0.968923in}}%
\pgfpathmoveto{\pgfqpoint{1.462941in}{0.965974in}}%
\pgfpathlineto{\pgfqpoint{1.462941in}{0.965974in}}%
\pgfpathlineto{\pgfqpoint{1.462941in}{0.968923in}}%
\pgfpathlineto{\pgfqpoint{1.467482in}{0.968923in}}%
\pgfpathlineto{\pgfqpoint{1.467482in}{0.965974in}}%
\pgfpathmoveto{\pgfqpoint{1.462941in}{0.968923in}}%
\pgfpathlineto{\pgfqpoint{1.462941in}{0.968923in}}%
\pgfpathlineto{\pgfqpoint{1.462941in}{0.971872in}}%
\pgfpathlineto{\pgfqpoint{1.467482in}{0.971872in}}%
\pgfpathlineto{\pgfqpoint{1.467482in}{0.968923in}}%
\pgfpathmoveto{\pgfqpoint{1.467482in}{0.965974in}}%
\pgfpathlineto{\pgfqpoint{1.467482in}{0.965974in}}%
\pgfpathlineto{\pgfqpoint{1.467482in}{0.968923in}}%
\pgfpathlineto{\pgfqpoint{1.472023in}{0.968923in}}%
\pgfpathlineto{\pgfqpoint{1.472023in}{0.965974in}}%
\pgfpathmoveto{\pgfqpoint{1.467482in}{0.968923in}}%
\pgfpathlineto{\pgfqpoint{1.467482in}{0.968923in}}%
\pgfpathlineto{\pgfqpoint{1.467482in}{0.971872in}}%
\pgfpathlineto{\pgfqpoint{1.472023in}{0.971872in}}%
\pgfpathlineto{\pgfqpoint{1.472023in}{0.968923in}}%
\pgfpathmoveto{\pgfqpoint{1.472023in}{0.968923in}}%
\pgfpathlineto{\pgfqpoint{1.472023in}{0.968923in}}%
\pgfpathlineto{\pgfqpoint{1.472023in}{0.971872in}}%
\pgfpathlineto{\pgfqpoint{1.476564in}{0.971872in}}%
\pgfpathlineto{\pgfqpoint{1.476564in}{0.968923in}}%
\pgfpathmoveto{\pgfqpoint{1.467482in}{0.971872in}}%
\pgfpathlineto{\pgfqpoint{1.467482in}{0.971872in}}%
\pgfpathlineto{\pgfqpoint{1.467482in}{0.974822in}}%
\pgfpathlineto{\pgfqpoint{1.472023in}{0.974822in}}%
\pgfpathlineto{\pgfqpoint{1.472023in}{0.971872in}}%
\pgfpathmoveto{\pgfqpoint{1.467482in}{0.974822in}}%
\pgfpathlineto{\pgfqpoint{1.467482in}{0.974822in}}%
\pgfpathlineto{\pgfqpoint{1.467482in}{0.977771in}}%
\pgfpathlineto{\pgfqpoint{1.472023in}{0.977771in}}%
\pgfpathlineto{\pgfqpoint{1.472023in}{0.974822in}}%
\pgfpathmoveto{\pgfqpoint{1.472023in}{0.971872in}}%
\pgfpathlineto{\pgfqpoint{1.472023in}{0.971872in}}%
\pgfpathlineto{\pgfqpoint{1.472023in}{0.974822in}}%
\pgfpathlineto{\pgfqpoint{1.476564in}{0.974822in}}%
\pgfpathlineto{\pgfqpoint{1.476564in}{0.971872in}}%
\pgfpathmoveto{\pgfqpoint{1.472023in}{0.974822in}}%
\pgfpathlineto{\pgfqpoint{1.472023in}{0.974822in}}%
\pgfpathlineto{\pgfqpoint{1.472023in}{0.977771in}}%
\pgfpathlineto{\pgfqpoint{1.476564in}{0.977771in}}%
\pgfpathlineto{\pgfqpoint{1.476564in}{0.974822in}}%
\pgfpathmoveto{\pgfqpoint{1.476564in}{0.971872in}}%
\pgfpathlineto{\pgfqpoint{1.476564in}{0.971872in}}%
\pgfpathlineto{\pgfqpoint{1.476564in}{0.974822in}}%
\pgfpathlineto{\pgfqpoint{1.481105in}{0.974822in}}%
\pgfpathlineto{\pgfqpoint{1.481105in}{0.971872in}}%
\pgfpathmoveto{\pgfqpoint{1.476564in}{0.974822in}}%
\pgfpathlineto{\pgfqpoint{1.476564in}{0.974822in}}%
\pgfpathlineto{\pgfqpoint{1.476564in}{0.977771in}}%
\pgfpathlineto{\pgfqpoint{1.481105in}{0.977771in}}%
\pgfpathlineto{\pgfqpoint{1.481105in}{0.974822in}}%
\pgfpathmoveto{\pgfqpoint{1.481105in}{0.974822in}}%
\pgfpathlineto{\pgfqpoint{1.481105in}{0.974822in}}%
\pgfpathlineto{\pgfqpoint{1.481105in}{0.977771in}}%
\pgfpathlineto{\pgfqpoint{1.485646in}{0.977771in}}%
\pgfpathlineto{\pgfqpoint{1.485646in}{0.974822in}}%
\pgfpathmoveto{\pgfqpoint{1.476564in}{0.977771in}}%
\pgfpathlineto{\pgfqpoint{1.476564in}{0.977771in}}%
\pgfpathlineto{\pgfqpoint{1.476564in}{0.980720in}}%
\pgfpathlineto{\pgfqpoint{1.481105in}{0.980720in}}%
\pgfpathlineto{\pgfqpoint{1.481105in}{0.977771in}}%
\pgfpathmoveto{\pgfqpoint{1.476564in}{0.980720in}}%
\pgfpathlineto{\pgfqpoint{1.476564in}{0.980720in}}%
\pgfpathlineto{\pgfqpoint{1.476564in}{0.983669in}}%
\pgfpathlineto{\pgfqpoint{1.481105in}{0.983669in}}%
\pgfpathlineto{\pgfqpoint{1.481105in}{0.980720in}}%
\pgfpathmoveto{\pgfqpoint{1.481105in}{0.977771in}}%
\pgfpathlineto{\pgfqpoint{1.481105in}{0.977771in}}%
\pgfpathlineto{\pgfqpoint{1.481105in}{0.980720in}}%
\pgfpathlineto{\pgfqpoint{1.485646in}{0.980720in}}%
\pgfpathlineto{\pgfqpoint{1.485646in}{0.977771in}}%
\pgfpathmoveto{\pgfqpoint{1.481105in}{0.980720in}}%
\pgfpathlineto{\pgfqpoint{1.481105in}{0.980720in}}%
\pgfpathlineto{\pgfqpoint{1.481105in}{0.983669in}}%
\pgfpathlineto{\pgfqpoint{1.485646in}{0.983669in}}%
\pgfpathlineto{\pgfqpoint{1.485646in}{0.980720in}}%
\pgfpathmoveto{\pgfqpoint{1.485646in}{0.977771in}}%
\pgfpathlineto{\pgfqpoint{1.485646in}{0.977771in}}%
\pgfpathlineto{\pgfqpoint{1.485646in}{0.980720in}}%
\pgfpathlineto{\pgfqpoint{1.490187in}{0.980720in}}%
\pgfpathlineto{\pgfqpoint{1.490187in}{0.977771in}}%
\pgfpathmoveto{\pgfqpoint{1.485646in}{0.980720in}}%
\pgfpathlineto{\pgfqpoint{1.485646in}{0.980720in}}%
\pgfpathlineto{\pgfqpoint{1.485646in}{0.983669in}}%
\pgfpathlineto{\pgfqpoint{1.490187in}{0.983669in}}%
\pgfpathlineto{\pgfqpoint{1.490187in}{0.980720in}}%
\pgfpathmoveto{\pgfqpoint{1.490187in}{0.980720in}}%
\pgfpathlineto{\pgfqpoint{1.490187in}{0.980720in}}%
\pgfpathlineto{\pgfqpoint{1.490187in}{0.983669in}}%
\pgfpathlineto{\pgfqpoint{1.494728in}{0.983669in}}%
\pgfpathlineto{\pgfqpoint{1.494728in}{0.980720in}}%
\pgfpathmoveto{\pgfqpoint{1.485646in}{0.983669in}}%
\pgfpathlineto{\pgfqpoint{1.485646in}{0.983669in}}%
\pgfpathlineto{\pgfqpoint{1.485646in}{0.986619in}}%
\pgfpathlineto{\pgfqpoint{1.490187in}{0.986619in}}%
\pgfpathlineto{\pgfqpoint{1.490187in}{0.983669in}}%
\pgfpathmoveto{\pgfqpoint{1.485646in}{0.986619in}}%
\pgfpathlineto{\pgfqpoint{1.485646in}{0.986619in}}%
\pgfpathlineto{\pgfqpoint{1.485646in}{0.989568in}}%
\pgfpathlineto{\pgfqpoint{1.490187in}{0.989568in}}%
\pgfpathlineto{\pgfqpoint{1.490187in}{0.986619in}}%
\pgfpathmoveto{\pgfqpoint{1.490187in}{0.983669in}}%
\pgfpathlineto{\pgfqpoint{1.490187in}{0.983669in}}%
\pgfpathlineto{\pgfqpoint{1.490187in}{0.986619in}}%
\pgfpathlineto{\pgfqpoint{1.494728in}{0.986619in}}%
\pgfpathlineto{\pgfqpoint{1.494728in}{0.983669in}}%
\pgfpathmoveto{\pgfqpoint{1.490187in}{0.986619in}}%
\pgfpathlineto{\pgfqpoint{1.490187in}{0.986619in}}%
\pgfpathlineto{\pgfqpoint{1.490187in}{0.989568in}}%
\pgfpathlineto{\pgfqpoint{1.494728in}{0.989568in}}%
\pgfpathlineto{\pgfqpoint{1.494728in}{0.986619in}}%
\pgfpathmoveto{\pgfqpoint{1.494728in}{0.983669in}}%
\pgfpathlineto{\pgfqpoint{1.494728in}{0.983669in}}%
\pgfpathlineto{\pgfqpoint{1.494728in}{0.986619in}}%
\pgfpathlineto{\pgfqpoint{1.499269in}{0.986619in}}%
\pgfpathlineto{\pgfqpoint{1.499269in}{0.983669in}}%
\pgfpathmoveto{\pgfqpoint{1.494728in}{0.986619in}}%
\pgfpathlineto{\pgfqpoint{1.494728in}{0.986619in}}%
\pgfpathlineto{\pgfqpoint{1.494728in}{0.989568in}}%
\pgfpathlineto{\pgfqpoint{1.499269in}{0.989568in}}%
\pgfpathlineto{\pgfqpoint{1.499269in}{0.986619in}}%
\pgfpathmoveto{\pgfqpoint{1.499269in}{0.986619in}}%
\pgfpathlineto{\pgfqpoint{1.499269in}{0.986619in}}%
\pgfpathlineto{\pgfqpoint{1.499269in}{0.989568in}}%
\pgfpathlineto{\pgfqpoint{1.503810in}{0.989568in}}%
\pgfpathlineto{\pgfqpoint{1.503810in}{0.986619in}}%
\pgfpathmoveto{\pgfqpoint{1.494728in}{0.989568in}}%
\pgfpathlineto{\pgfqpoint{1.494728in}{0.989568in}}%
\pgfpathlineto{\pgfqpoint{1.494728in}{0.992517in}}%
\pgfpathlineto{\pgfqpoint{1.499269in}{0.992517in}}%
\pgfpathlineto{\pgfqpoint{1.499269in}{0.989568in}}%
\pgfpathmoveto{\pgfqpoint{1.494728in}{0.992517in}}%
\pgfpathlineto{\pgfqpoint{1.494728in}{0.992517in}}%
\pgfpathlineto{\pgfqpoint{1.494728in}{0.995466in}}%
\pgfpathlineto{\pgfqpoint{1.499269in}{0.995466in}}%
\pgfpathlineto{\pgfqpoint{1.499269in}{0.992517in}}%
\pgfpathmoveto{\pgfqpoint{1.499269in}{0.989568in}}%
\pgfpathlineto{\pgfqpoint{1.499269in}{0.989568in}}%
\pgfpathlineto{\pgfqpoint{1.499269in}{0.992517in}}%
\pgfpathlineto{\pgfqpoint{1.503810in}{0.992517in}}%
\pgfpathlineto{\pgfqpoint{1.503810in}{0.989568in}}%
\pgfpathmoveto{\pgfqpoint{1.499269in}{0.992517in}}%
\pgfpathlineto{\pgfqpoint{1.499269in}{0.992517in}}%
\pgfpathlineto{\pgfqpoint{1.499269in}{0.995466in}}%
\pgfpathlineto{\pgfqpoint{1.503810in}{0.995466in}}%
\pgfpathlineto{\pgfqpoint{1.503810in}{0.992517in}}%
\pgfpathmoveto{\pgfqpoint{1.503810in}{0.989568in}}%
\pgfpathlineto{\pgfqpoint{1.503810in}{0.989568in}}%
\pgfpathlineto{\pgfqpoint{1.503810in}{0.992517in}}%
\pgfpathlineto{\pgfqpoint{1.508351in}{0.992517in}}%
\pgfpathlineto{\pgfqpoint{1.508351in}{0.989568in}}%
\pgfpathmoveto{\pgfqpoint{1.503810in}{0.992517in}}%
\pgfpathlineto{\pgfqpoint{1.503810in}{0.992517in}}%
\pgfpathlineto{\pgfqpoint{1.503810in}{0.995466in}}%
\pgfpathlineto{\pgfqpoint{1.508351in}{0.995466in}}%
\pgfpathlineto{\pgfqpoint{1.508351in}{0.992517in}}%
\pgfpathmoveto{\pgfqpoint{1.508351in}{0.992517in}}%
\pgfpathlineto{\pgfqpoint{1.508351in}{0.992517in}}%
\pgfpathlineto{\pgfqpoint{1.508351in}{0.995466in}}%
\pgfpathlineto{\pgfqpoint{1.512892in}{0.995466in}}%
\pgfpathlineto{\pgfqpoint{1.512892in}{0.992517in}}%
\pgfpathmoveto{\pgfqpoint{1.503810in}{0.995466in}}%
\pgfpathlineto{\pgfqpoint{1.503810in}{0.995466in}}%
\pgfpathlineto{\pgfqpoint{1.503810in}{0.998416in}}%
\pgfpathlineto{\pgfqpoint{1.508351in}{0.998416in}}%
\pgfpathlineto{\pgfqpoint{1.508351in}{0.995466in}}%
\pgfpathmoveto{\pgfqpoint{1.503810in}{0.998416in}}%
\pgfpathlineto{\pgfqpoint{1.503810in}{0.998416in}}%
\pgfpathlineto{\pgfqpoint{1.503810in}{1.001365in}}%
\pgfpathlineto{\pgfqpoint{1.508351in}{1.001365in}}%
\pgfpathlineto{\pgfqpoint{1.508351in}{0.998416in}}%
\pgfpathmoveto{\pgfqpoint{1.508351in}{0.995466in}}%
\pgfpathlineto{\pgfqpoint{1.508351in}{0.995466in}}%
\pgfpathlineto{\pgfqpoint{1.508351in}{0.998416in}}%
\pgfpathlineto{\pgfqpoint{1.512892in}{0.998416in}}%
\pgfpathlineto{\pgfqpoint{1.512892in}{0.995466in}}%
\pgfpathmoveto{\pgfqpoint{1.508351in}{0.998416in}}%
\pgfpathlineto{\pgfqpoint{1.508351in}{0.998416in}}%
\pgfpathlineto{\pgfqpoint{1.508351in}{1.001365in}}%
\pgfpathlineto{\pgfqpoint{1.512892in}{1.001365in}}%
\pgfpathlineto{\pgfqpoint{1.512892in}{0.998416in}}%
\pgfpathmoveto{\pgfqpoint{1.512892in}{0.995466in}}%
\pgfpathlineto{\pgfqpoint{1.512892in}{0.995466in}}%
\pgfpathlineto{\pgfqpoint{1.512892in}{0.998416in}}%
\pgfpathlineto{\pgfqpoint{1.517433in}{0.998416in}}%
\pgfpathlineto{\pgfqpoint{1.517433in}{0.995466in}}%
\pgfpathmoveto{\pgfqpoint{1.512892in}{0.998416in}}%
\pgfpathlineto{\pgfqpoint{1.512892in}{0.998416in}}%
\pgfpathlineto{\pgfqpoint{1.512892in}{1.001365in}}%
\pgfpathlineto{\pgfqpoint{1.517433in}{1.001365in}}%
\pgfpathlineto{\pgfqpoint{1.517433in}{0.998416in}}%
\pgfpathmoveto{\pgfqpoint{1.517433in}{0.998416in}}%
\pgfpathlineto{\pgfqpoint{1.517433in}{0.998416in}}%
\pgfpathlineto{\pgfqpoint{1.517433in}{1.001365in}}%
\pgfpathlineto{\pgfqpoint{1.521974in}{1.001365in}}%
\pgfpathlineto{\pgfqpoint{1.521974in}{0.998416in}}%
\pgfpathmoveto{\pgfqpoint{1.512892in}{1.001365in}}%
\pgfpathlineto{\pgfqpoint{1.512892in}{1.001365in}}%
\pgfpathlineto{\pgfqpoint{1.512892in}{1.004314in}}%
\pgfpathlineto{\pgfqpoint{1.517433in}{1.004314in}}%
\pgfpathlineto{\pgfqpoint{1.517433in}{1.001365in}}%
\pgfpathmoveto{\pgfqpoint{1.512892in}{1.004314in}}%
\pgfpathlineto{\pgfqpoint{1.512892in}{1.004314in}}%
\pgfpathlineto{\pgfqpoint{1.512892in}{1.007264in}}%
\pgfpathlineto{\pgfqpoint{1.517433in}{1.007264in}}%
\pgfpathlineto{\pgfqpoint{1.517433in}{1.004314in}}%
\pgfpathmoveto{\pgfqpoint{1.517433in}{1.001365in}}%
\pgfpathlineto{\pgfqpoint{1.517433in}{1.001365in}}%
\pgfpathlineto{\pgfqpoint{1.517433in}{1.004314in}}%
\pgfpathlineto{\pgfqpoint{1.521974in}{1.004314in}}%
\pgfpathlineto{\pgfqpoint{1.521974in}{1.001365in}}%
\pgfpathmoveto{\pgfqpoint{1.517433in}{1.004314in}}%
\pgfpathlineto{\pgfqpoint{1.517433in}{1.004314in}}%
\pgfpathlineto{\pgfqpoint{1.517433in}{1.007264in}}%
\pgfpathlineto{\pgfqpoint{1.521974in}{1.007264in}}%
\pgfpathlineto{\pgfqpoint{1.521974in}{1.004314in}}%
\pgfpathmoveto{\pgfqpoint{1.521974in}{1.001365in}}%
\pgfpathlineto{\pgfqpoint{1.521974in}{1.001365in}}%
\pgfpathlineto{\pgfqpoint{1.521974in}{1.004314in}}%
\pgfpathlineto{\pgfqpoint{1.526515in}{1.004314in}}%
\pgfpathlineto{\pgfqpoint{1.526515in}{1.001365in}}%
\pgfpathmoveto{\pgfqpoint{1.521974in}{1.004314in}}%
\pgfpathlineto{\pgfqpoint{1.521974in}{1.004314in}}%
\pgfpathlineto{\pgfqpoint{1.521974in}{1.007264in}}%
\pgfpathlineto{\pgfqpoint{1.526515in}{1.007264in}}%
\pgfpathlineto{\pgfqpoint{1.526515in}{1.004314in}}%
\pgfpathmoveto{\pgfqpoint{1.526515in}{1.004314in}}%
\pgfpathlineto{\pgfqpoint{1.526515in}{1.004314in}}%
\pgfpathlineto{\pgfqpoint{1.526515in}{1.007264in}}%
\pgfpathlineto{\pgfqpoint{1.531056in}{1.007264in}}%
\pgfpathlineto{\pgfqpoint{1.531056in}{1.004314in}}%
\pgfpathmoveto{\pgfqpoint{1.521974in}{1.007264in}}%
\pgfpathlineto{\pgfqpoint{1.521974in}{1.007264in}}%
\pgfpathlineto{\pgfqpoint{1.521974in}{1.010213in}}%
\pgfpathlineto{\pgfqpoint{1.526515in}{1.010213in}}%
\pgfpathlineto{\pgfqpoint{1.526515in}{1.007264in}}%
\pgfpathmoveto{\pgfqpoint{1.521974in}{1.010213in}}%
\pgfpathlineto{\pgfqpoint{1.521974in}{1.010213in}}%
\pgfpathlineto{\pgfqpoint{1.521974in}{1.013162in}}%
\pgfpathlineto{\pgfqpoint{1.526515in}{1.013162in}}%
\pgfpathlineto{\pgfqpoint{1.526515in}{1.010213in}}%
\pgfpathmoveto{\pgfqpoint{1.526515in}{1.007264in}}%
\pgfpathlineto{\pgfqpoint{1.526515in}{1.007264in}}%
\pgfpathlineto{\pgfqpoint{1.526515in}{1.010213in}}%
\pgfpathlineto{\pgfqpoint{1.531056in}{1.010213in}}%
\pgfpathlineto{\pgfqpoint{1.531056in}{1.007264in}}%
\pgfpathmoveto{\pgfqpoint{1.526515in}{1.010213in}}%
\pgfpathlineto{\pgfqpoint{1.526515in}{1.010213in}}%
\pgfpathlineto{\pgfqpoint{1.526515in}{1.013162in}}%
\pgfpathlineto{\pgfqpoint{1.531056in}{1.013162in}}%
\pgfpathlineto{\pgfqpoint{1.531056in}{1.010213in}}%
\pgfpathmoveto{\pgfqpoint{1.531056in}{1.007264in}}%
\pgfpathlineto{\pgfqpoint{1.531056in}{1.007264in}}%
\pgfpathlineto{\pgfqpoint{1.531056in}{1.010213in}}%
\pgfpathlineto{\pgfqpoint{1.535597in}{1.010213in}}%
\pgfpathlineto{\pgfqpoint{1.535597in}{1.007264in}}%
\pgfpathmoveto{\pgfqpoint{1.531056in}{1.010213in}}%
\pgfpathlineto{\pgfqpoint{1.531056in}{1.010213in}}%
\pgfpathlineto{\pgfqpoint{1.531056in}{1.013162in}}%
\pgfpathlineto{\pgfqpoint{1.535597in}{1.013162in}}%
\pgfpathlineto{\pgfqpoint{1.535597in}{1.010213in}}%
\pgfpathmoveto{\pgfqpoint{1.535597in}{1.010213in}}%
\pgfpathlineto{\pgfqpoint{1.535597in}{1.010213in}}%
\pgfpathlineto{\pgfqpoint{1.535597in}{1.013162in}}%
\pgfpathlineto{\pgfqpoint{1.540138in}{1.013162in}}%
\pgfpathlineto{\pgfqpoint{1.540138in}{1.010213in}}%
\pgfpathmoveto{\pgfqpoint{1.531056in}{1.013162in}}%
\pgfpathlineto{\pgfqpoint{1.531056in}{1.013162in}}%
\pgfpathlineto{\pgfqpoint{1.531056in}{1.016111in}}%
\pgfpathlineto{\pgfqpoint{1.535597in}{1.016111in}}%
\pgfpathlineto{\pgfqpoint{1.535597in}{1.013162in}}%
\pgfpathmoveto{\pgfqpoint{1.531056in}{1.016111in}}%
\pgfpathlineto{\pgfqpoint{1.531056in}{1.016111in}}%
\pgfpathlineto{\pgfqpoint{1.531056in}{1.019061in}}%
\pgfpathlineto{\pgfqpoint{1.535597in}{1.019061in}}%
\pgfpathlineto{\pgfqpoint{1.535597in}{1.016111in}}%
\pgfpathmoveto{\pgfqpoint{1.535597in}{1.013162in}}%
\pgfpathlineto{\pgfqpoint{1.535597in}{1.013162in}}%
\pgfpathlineto{\pgfqpoint{1.535597in}{1.016111in}}%
\pgfpathlineto{\pgfqpoint{1.540138in}{1.016111in}}%
\pgfpathlineto{\pgfqpoint{1.540138in}{1.013162in}}%
\pgfpathmoveto{\pgfqpoint{1.535597in}{1.016111in}}%
\pgfpathlineto{\pgfqpoint{1.535597in}{1.016111in}}%
\pgfpathlineto{\pgfqpoint{1.535597in}{1.019061in}}%
\pgfpathlineto{\pgfqpoint{1.540138in}{1.019061in}}%
\pgfpathlineto{\pgfqpoint{1.540138in}{1.016111in}}%
\pgfpathmoveto{\pgfqpoint{1.540138in}{1.013162in}}%
\pgfpathlineto{\pgfqpoint{1.540138in}{1.013162in}}%
\pgfpathlineto{\pgfqpoint{1.540138in}{1.016111in}}%
\pgfpathlineto{\pgfqpoint{1.544679in}{1.016111in}}%
\pgfpathlineto{\pgfqpoint{1.544679in}{1.013162in}}%
\pgfpathmoveto{\pgfqpoint{1.540138in}{1.016111in}}%
\pgfpathlineto{\pgfqpoint{1.540138in}{1.016111in}}%
\pgfpathlineto{\pgfqpoint{1.540138in}{1.019061in}}%
\pgfpathlineto{\pgfqpoint{1.544679in}{1.019061in}}%
\pgfpathlineto{\pgfqpoint{1.544679in}{1.016111in}}%
\pgfpathmoveto{\pgfqpoint{1.544679in}{1.016111in}}%
\pgfpathlineto{\pgfqpoint{1.544679in}{1.016111in}}%
\pgfpathlineto{\pgfqpoint{1.544679in}{1.019061in}}%
\pgfpathlineto{\pgfqpoint{1.549220in}{1.019061in}}%
\pgfpathlineto{\pgfqpoint{1.549220in}{1.016111in}}%
\pgfpathmoveto{\pgfqpoint{1.540138in}{1.019061in}}%
\pgfpathlineto{\pgfqpoint{1.540138in}{1.019061in}}%
\pgfpathlineto{\pgfqpoint{1.540138in}{1.022010in}}%
\pgfpathlineto{\pgfqpoint{1.544679in}{1.022010in}}%
\pgfpathlineto{\pgfqpoint{1.544679in}{1.019061in}}%
\pgfpathmoveto{\pgfqpoint{1.540138in}{1.022010in}}%
\pgfpathlineto{\pgfqpoint{1.540138in}{1.022010in}}%
\pgfpathlineto{\pgfqpoint{1.540138in}{1.024959in}}%
\pgfpathlineto{\pgfqpoint{1.544679in}{1.024959in}}%
\pgfpathlineto{\pgfqpoint{1.544679in}{1.022010in}}%
\pgfpathmoveto{\pgfqpoint{1.544679in}{1.019061in}}%
\pgfpathlineto{\pgfqpoint{1.544679in}{1.019061in}}%
\pgfpathlineto{\pgfqpoint{1.544679in}{1.022010in}}%
\pgfpathlineto{\pgfqpoint{1.549220in}{1.022010in}}%
\pgfpathlineto{\pgfqpoint{1.549220in}{1.019061in}}%
\pgfpathmoveto{\pgfqpoint{1.544679in}{1.022010in}}%
\pgfpathlineto{\pgfqpoint{1.544679in}{1.022010in}}%
\pgfpathlineto{\pgfqpoint{1.544679in}{1.024959in}}%
\pgfpathlineto{\pgfqpoint{1.549220in}{1.024959in}}%
\pgfpathlineto{\pgfqpoint{1.549220in}{1.022010in}}%
\pgfpathmoveto{\pgfqpoint{1.549220in}{1.019061in}}%
\pgfpathlineto{\pgfqpoint{1.549220in}{1.019061in}}%
\pgfpathlineto{\pgfqpoint{1.549220in}{1.022010in}}%
\pgfpathlineto{\pgfqpoint{1.553761in}{1.022010in}}%
\pgfpathlineto{\pgfqpoint{1.553761in}{1.019061in}}%
\pgfpathmoveto{\pgfqpoint{1.549220in}{1.022010in}}%
\pgfpathlineto{\pgfqpoint{1.549220in}{1.022010in}}%
\pgfpathlineto{\pgfqpoint{1.549220in}{1.024959in}}%
\pgfpathlineto{\pgfqpoint{1.553761in}{1.024959in}}%
\pgfpathlineto{\pgfqpoint{1.553761in}{1.022010in}}%
\pgfpathmoveto{\pgfqpoint{1.553761in}{1.022010in}}%
\pgfpathlineto{\pgfqpoint{1.553761in}{1.022010in}}%
\pgfpathlineto{\pgfqpoint{1.553761in}{1.024959in}}%
\pgfpathlineto{\pgfqpoint{1.558302in}{1.024959in}}%
\pgfpathlineto{\pgfqpoint{1.558302in}{1.022010in}}%
\pgfpathmoveto{\pgfqpoint{1.549220in}{1.024959in}}%
\pgfpathlineto{\pgfqpoint{1.549220in}{1.024959in}}%
\pgfpathlineto{\pgfqpoint{1.549220in}{1.027908in}}%
\pgfpathlineto{\pgfqpoint{1.553761in}{1.027908in}}%
\pgfpathlineto{\pgfqpoint{1.553761in}{1.024959in}}%
\pgfpathmoveto{\pgfqpoint{1.549220in}{1.027908in}}%
\pgfpathlineto{\pgfqpoint{1.549220in}{1.027908in}}%
\pgfpathlineto{\pgfqpoint{1.549220in}{1.030858in}}%
\pgfpathlineto{\pgfqpoint{1.553761in}{1.030858in}}%
\pgfpathlineto{\pgfqpoint{1.553761in}{1.027908in}}%
\pgfpathmoveto{\pgfqpoint{1.553761in}{1.024959in}}%
\pgfpathlineto{\pgfqpoint{1.553761in}{1.024959in}}%
\pgfpathlineto{\pgfqpoint{1.553761in}{1.027908in}}%
\pgfpathlineto{\pgfqpoint{1.558302in}{1.027908in}}%
\pgfpathlineto{\pgfqpoint{1.558302in}{1.024959in}}%
\pgfpathmoveto{\pgfqpoint{1.553761in}{1.027908in}}%
\pgfpathlineto{\pgfqpoint{1.553761in}{1.027908in}}%
\pgfpathlineto{\pgfqpoint{1.553761in}{1.030858in}}%
\pgfpathlineto{\pgfqpoint{1.558302in}{1.030858in}}%
\pgfpathlineto{\pgfqpoint{1.558302in}{1.027908in}}%
\pgfpathmoveto{\pgfqpoint{1.558302in}{1.024959in}}%
\pgfpathlineto{\pgfqpoint{1.558302in}{1.024959in}}%
\pgfpathlineto{\pgfqpoint{1.558302in}{1.027908in}}%
\pgfpathlineto{\pgfqpoint{1.562843in}{1.027908in}}%
\pgfpathlineto{\pgfqpoint{1.562843in}{1.024959in}}%
\pgfpathmoveto{\pgfqpoint{1.558302in}{1.027908in}}%
\pgfpathlineto{\pgfqpoint{1.558302in}{1.027908in}}%
\pgfpathlineto{\pgfqpoint{1.558302in}{1.030858in}}%
\pgfpathlineto{\pgfqpoint{1.562843in}{1.030858in}}%
\pgfpathlineto{\pgfqpoint{1.562843in}{1.027908in}}%
\pgfpathmoveto{\pgfqpoint{1.562843in}{1.027908in}}%
\pgfpathlineto{\pgfqpoint{1.562843in}{1.027908in}}%
\pgfpathlineto{\pgfqpoint{1.562843in}{1.030858in}}%
\pgfpathlineto{\pgfqpoint{1.567384in}{1.030858in}}%
\pgfpathlineto{\pgfqpoint{1.567384in}{1.027908in}}%
\pgfpathmoveto{\pgfqpoint{1.558302in}{1.030858in}}%
\pgfpathlineto{\pgfqpoint{1.558302in}{1.030858in}}%
\pgfpathlineto{\pgfqpoint{1.558302in}{1.033807in}}%
\pgfpathlineto{\pgfqpoint{1.562843in}{1.033807in}}%
\pgfpathlineto{\pgfqpoint{1.562843in}{1.030858in}}%
\pgfpathmoveto{\pgfqpoint{1.558302in}{1.033807in}}%
\pgfpathlineto{\pgfqpoint{1.558302in}{1.033807in}}%
\pgfpathlineto{\pgfqpoint{1.558302in}{1.036756in}}%
\pgfpathlineto{\pgfqpoint{1.562843in}{1.036756in}}%
\pgfpathlineto{\pgfqpoint{1.562843in}{1.033807in}}%
\pgfpathmoveto{\pgfqpoint{1.562843in}{1.030858in}}%
\pgfpathlineto{\pgfqpoint{1.562843in}{1.030858in}}%
\pgfpathlineto{\pgfqpoint{1.562843in}{1.033807in}}%
\pgfpathlineto{\pgfqpoint{1.567384in}{1.033807in}}%
\pgfpathlineto{\pgfqpoint{1.567384in}{1.030858in}}%
\pgfpathmoveto{\pgfqpoint{1.562843in}{1.033807in}}%
\pgfpathlineto{\pgfqpoint{1.562843in}{1.033807in}}%
\pgfpathlineto{\pgfqpoint{1.562843in}{1.036756in}}%
\pgfpathlineto{\pgfqpoint{1.567384in}{1.036756in}}%
\pgfpathlineto{\pgfqpoint{1.567384in}{1.033807in}}%
\pgfpathmoveto{\pgfqpoint{1.567384in}{1.030858in}}%
\pgfpathlineto{\pgfqpoint{1.567384in}{1.030858in}}%
\pgfpathlineto{\pgfqpoint{1.567384in}{1.033807in}}%
\pgfpathlineto{\pgfqpoint{1.571925in}{1.033807in}}%
\pgfpathlineto{\pgfqpoint{1.571925in}{1.030858in}}%
\pgfpathmoveto{\pgfqpoint{1.567384in}{1.033807in}}%
\pgfpathlineto{\pgfqpoint{1.567384in}{1.033807in}}%
\pgfpathlineto{\pgfqpoint{1.567384in}{1.036756in}}%
\pgfpathlineto{\pgfqpoint{1.571925in}{1.036756in}}%
\pgfpathlineto{\pgfqpoint{1.571925in}{1.033807in}}%
\pgfpathmoveto{\pgfqpoint{1.571925in}{1.033807in}}%
\pgfpathlineto{\pgfqpoint{1.571925in}{1.033807in}}%
\pgfpathlineto{\pgfqpoint{1.571925in}{1.036756in}}%
\pgfpathlineto{\pgfqpoint{1.576466in}{1.036756in}}%
\pgfpathlineto{\pgfqpoint{1.576466in}{1.033807in}}%
\pgfpathmoveto{\pgfqpoint{1.567384in}{1.036756in}}%
\pgfpathlineto{\pgfqpoint{1.567384in}{1.036756in}}%
\pgfpathlineto{\pgfqpoint{1.567384in}{1.039705in}}%
\pgfpathlineto{\pgfqpoint{1.571925in}{1.039705in}}%
\pgfpathlineto{\pgfqpoint{1.571925in}{1.036756in}}%
\pgfpathmoveto{\pgfqpoint{1.567384in}{1.039705in}}%
\pgfpathlineto{\pgfqpoint{1.567384in}{1.039705in}}%
\pgfpathlineto{\pgfqpoint{1.567384in}{1.042655in}}%
\pgfpathlineto{\pgfqpoint{1.571925in}{1.042655in}}%
\pgfpathlineto{\pgfqpoint{1.571925in}{1.039705in}}%
\pgfpathmoveto{\pgfqpoint{1.571925in}{1.036756in}}%
\pgfpathlineto{\pgfqpoint{1.571925in}{1.036756in}}%
\pgfpathlineto{\pgfqpoint{1.571925in}{1.039705in}}%
\pgfpathlineto{\pgfqpoint{1.576466in}{1.039705in}}%
\pgfpathlineto{\pgfqpoint{1.576466in}{1.036756in}}%
\pgfpathmoveto{\pgfqpoint{1.571925in}{1.039705in}}%
\pgfpathlineto{\pgfqpoint{1.571925in}{1.039705in}}%
\pgfpathlineto{\pgfqpoint{1.571925in}{1.042655in}}%
\pgfpathlineto{\pgfqpoint{1.576466in}{1.042655in}}%
\pgfpathlineto{\pgfqpoint{1.576466in}{1.039705in}}%
\pgfpathmoveto{\pgfqpoint{1.576466in}{1.036756in}}%
\pgfpathlineto{\pgfqpoint{1.576466in}{1.036756in}}%
\pgfpathlineto{\pgfqpoint{1.576466in}{1.039705in}}%
\pgfpathlineto{\pgfqpoint{1.581007in}{1.039705in}}%
\pgfpathlineto{\pgfqpoint{1.581007in}{1.036756in}}%
\pgfpathmoveto{\pgfqpoint{1.576466in}{1.039705in}}%
\pgfpathlineto{\pgfqpoint{1.576466in}{1.039705in}}%
\pgfpathlineto{\pgfqpoint{1.576466in}{1.042655in}}%
\pgfpathlineto{\pgfqpoint{1.581007in}{1.042655in}}%
\pgfpathlineto{\pgfqpoint{1.581007in}{1.039705in}}%
\pgfpathmoveto{\pgfqpoint{1.581007in}{1.039705in}}%
\pgfpathlineto{\pgfqpoint{1.581007in}{1.039705in}}%
\pgfpathlineto{\pgfqpoint{1.581007in}{1.042655in}}%
\pgfpathlineto{\pgfqpoint{1.585549in}{1.042655in}}%
\pgfpathlineto{\pgfqpoint{1.585549in}{1.039705in}}%
\pgfpathmoveto{\pgfqpoint{1.576466in}{1.042655in}}%
\pgfpathlineto{\pgfqpoint{1.576466in}{1.042655in}}%
\pgfpathlineto{\pgfqpoint{1.576466in}{1.045604in}}%
\pgfpathlineto{\pgfqpoint{1.581007in}{1.045604in}}%
\pgfpathlineto{\pgfqpoint{1.581007in}{1.042655in}}%
\pgfpathmoveto{\pgfqpoint{1.576466in}{1.045604in}}%
\pgfpathlineto{\pgfqpoint{1.576466in}{1.045604in}}%
\pgfpathlineto{\pgfqpoint{1.576466in}{1.048553in}}%
\pgfpathlineto{\pgfqpoint{1.581007in}{1.048553in}}%
\pgfpathlineto{\pgfqpoint{1.581007in}{1.045604in}}%
\pgfpathmoveto{\pgfqpoint{1.581007in}{1.042655in}}%
\pgfpathlineto{\pgfqpoint{1.581007in}{1.042655in}}%
\pgfpathlineto{\pgfqpoint{1.581007in}{1.045604in}}%
\pgfpathlineto{\pgfqpoint{1.585549in}{1.045604in}}%
\pgfpathlineto{\pgfqpoint{1.585549in}{1.042655in}}%
\pgfpathmoveto{\pgfqpoint{1.581007in}{1.045604in}}%
\pgfpathlineto{\pgfqpoint{1.581007in}{1.045604in}}%
\pgfpathlineto{\pgfqpoint{1.581007in}{1.048553in}}%
\pgfpathlineto{\pgfqpoint{1.585549in}{1.048553in}}%
\pgfpathlineto{\pgfqpoint{1.585549in}{1.045604in}}%
\pgfpathmoveto{\pgfqpoint{1.585549in}{1.042655in}}%
\pgfpathlineto{\pgfqpoint{1.585549in}{1.042655in}}%
\pgfpathlineto{\pgfqpoint{1.585549in}{1.045604in}}%
\pgfpathlineto{\pgfqpoint{1.590090in}{1.045604in}}%
\pgfpathlineto{\pgfqpoint{1.590090in}{1.042655in}}%
\pgfpathmoveto{\pgfqpoint{1.585549in}{1.045604in}}%
\pgfpathlineto{\pgfqpoint{1.585549in}{1.045604in}}%
\pgfpathlineto{\pgfqpoint{1.585549in}{1.048553in}}%
\pgfpathlineto{\pgfqpoint{1.590090in}{1.048553in}}%
\pgfpathlineto{\pgfqpoint{1.590090in}{1.045604in}}%
\pgfpathmoveto{\pgfqpoint{1.590090in}{1.045604in}}%
\pgfpathlineto{\pgfqpoint{1.590090in}{1.045604in}}%
\pgfpathlineto{\pgfqpoint{1.590090in}{1.048553in}}%
\pgfpathlineto{\pgfqpoint{1.594631in}{1.048553in}}%
\pgfpathlineto{\pgfqpoint{1.594631in}{1.045604in}}%
\pgfpathmoveto{\pgfqpoint{1.585549in}{1.048553in}}%
\pgfpathlineto{\pgfqpoint{1.585549in}{1.048553in}}%
\pgfpathlineto{\pgfqpoint{1.585549in}{1.051502in}}%
\pgfpathlineto{\pgfqpoint{1.590090in}{1.051502in}}%
\pgfpathlineto{\pgfqpoint{1.590090in}{1.048553in}}%
\pgfpathmoveto{\pgfqpoint{1.585549in}{1.051502in}}%
\pgfpathlineto{\pgfqpoint{1.585549in}{1.051502in}}%
\pgfpathlineto{\pgfqpoint{1.585549in}{1.054452in}}%
\pgfpathlineto{\pgfqpoint{1.590090in}{1.054452in}}%
\pgfpathlineto{\pgfqpoint{1.590090in}{1.051502in}}%
\pgfpathmoveto{\pgfqpoint{1.590090in}{1.048553in}}%
\pgfpathlineto{\pgfqpoint{1.590090in}{1.048553in}}%
\pgfpathlineto{\pgfqpoint{1.590090in}{1.051502in}}%
\pgfpathlineto{\pgfqpoint{1.594631in}{1.051502in}}%
\pgfpathlineto{\pgfqpoint{1.594631in}{1.048553in}}%
\pgfpathmoveto{\pgfqpoint{1.590090in}{1.051502in}}%
\pgfpathlineto{\pgfqpoint{1.590090in}{1.051502in}}%
\pgfpathlineto{\pgfqpoint{1.590090in}{1.054452in}}%
\pgfpathlineto{\pgfqpoint{1.594631in}{1.054452in}}%
\pgfpathlineto{\pgfqpoint{1.594631in}{1.051502in}}%
\pgfpathmoveto{\pgfqpoint{1.594631in}{1.048553in}}%
\pgfpathlineto{\pgfqpoint{1.594631in}{1.048553in}}%
\pgfpathlineto{\pgfqpoint{1.594631in}{1.051502in}}%
\pgfpathlineto{\pgfqpoint{1.599172in}{1.051502in}}%
\pgfpathlineto{\pgfqpoint{1.599172in}{1.048553in}}%
\pgfpathmoveto{\pgfqpoint{1.594631in}{1.051502in}}%
\pgfpathlineto{\pgfqpoint{1.594631in}{1.051502in}}%
\pgfpathlineto{\pgfqpoint{1.594631in}{1.054452in}}%
\pgfpathlineto{\pgfqpoint{1.599172in}{1.054452in}}%
\pgfpathlineto{\pgfqpoint{1.599172in}{1.051502in}}%
\pgfpathmoveto{\pgfqpoint{1.599172in}{1.051502in}}%
\pgfpathlineto{\pgfqpoint{1.599172in}{1.051502in}}%
\pgfpathlineto{\pgfqpoint{1.599172in}{1.054452in}}%
\pgfpathlineto{\pgfqpoint{1.603713in}{1.054452in}}%
\pgfpathlineto{\pgfqpoint{1.603713in}{1.051502in}}%
\pgfpathmoveto{\pgfqpoint{1.594631in}{1.054452in}}%
\pgfpathlineto{\pgfqpoint{1.594631in}{1.054452in}}%
\pgfpathlineto{\pgfqpoint{1.594631in}{1.057401in}}%
\pgfpathlineto{\pgfqpoint{1.599172in}{1.057401in}}%
\pgfpathlineto{\pgfqpoint{1.599172in}{1.054452in}}%
\pgfpathmoveto{\pgfqpoint{1.594631in}{1.057401in}}%
\pgfpathlineto{\pgfqpoint{1.594631in}{1.057401in}}%
\pgfpathlineto{\pgfqpoint{1.594631in}{1.060350in}}%
\pgfpathlineto{\pgfqpoint{1.599172in}{1.060350in}}%
\pgfpathlineto{\pgfqpoint{1.599172in}{1.057401in}}%
\pgfpathmoveto{\pgfqpoint{1.599172in}{1.054452in}}%
\pgfpathlineto{\pgfqpoint{1.599172in}{1.054452in}}%
\pgfpathlineto{\pgfqpoint{1.599172in}{1.057401in}}%
\pgfpathlineto{\pgfqpoint{1.603713in}{1.057401in}}%
\pgfpathlineto{\pgfqpoint{1.603713in}{1.054452in}}%
\pgfpathmoveto{\pgfqpoint{1.599172in}{1.057401in}}%
\pgfpathlineto{\pgfqpoint{1.599172in}{1.057401in}}%
\pgfpathlineto{\pgfqpoint{1.599172in}{1.060350in}}%
\pgfpathlineto{\pgfqpoint{1.603713in}{1.060350in}}%
\pgfpathlineto{\pgfqpoint{1.603713in}{1.057401in}}%
\pgfpathmoveto{\pgfqpoint{1.603713in}{1.054452in}}%
\pgfpathlineto{\pgfqpoint{1.603713in}{1.054452in}}%
\pgfpathlineto{\pgfqpoint{1.603713in}{1.057401in}}%
\pgfpathlineto{\pgfqpoint{1.608254in}{1.057401in}}%
\pgfpathlineto{\pgfqpoint{1.608254in}{1.054452in}}%
\pgfpathmoveto{\pgfqpoint{1.603713in}{1.057401in}}%
\pgfpathlineto{\pgfqpoint{1.603713in}{1.057401in}}%
\pgfpathlineto{\pgfqpoint{1.603713in}{1.060350in}}%
\pgfpathlineto{\pgfqpoint{1.608254in}{1.060350in}}%
\pgfpathlineto{\pgfqpoint{1.608254in}{1.057401in}}%
\pgfpathmoveto{\pgfqpoint{1.608254in}{1.057401in}}%
\pgfpathlineto{\pgfqpoint{1.608254in}{1.057401in}}%
\pgfpathlineto{\pgfqpoint{1.608254in}{1.060350in}}%
\pgfpathlineto{\pgfqpoint{1.612795in}{1.060350in}}%
\pgfpathlineto{\pgfqpoint{1.612795in}{1.057401in}}%
\pgfpathmoveto{\pgfqpoint{1.603713in}{1.060350in}}%
\pgfpathlineto{\pgfqpoint{1.603713in}{1.060350in}}%
\pgfpathlineto{\pgfqpoint{1.603713in}{1.063300in}}%
\pgfpathlineto{\pgfqpoint{1.608254in}{1.063300in}}%
\pgfpathlineto{\pgfqpoint{1.608254in}{1.060350in}}%
\pgfpathmoveto{\pgfqpoint{1.603713in}{1.063300in}}%
\pgfpathlineto{\pgfqpoint{1.603713in}{1.063300in}}%
\pgfpathlineto{\pgfqpoint{1.603713in}{1.066249in}}%
\pgfpathlineto{\pgfqpoint{1.608254in}{1.066249in}}%
\pgfpathlineto{\pgfqpoint{1.608254in}{1.063300in}}%
\pgfpathmoveto{\pgfqpoint{1.608254in}{1.060350in}}%
\pgfpathlineto{\pgfqpoint{1.608254in}{1.060350in}}%
\pgfpathlineto{\pgfqpoint{1.608254in}{1.063300in}}%
\pgfpathlineto{\pgfqpoint{1.612795in}{1.063300in}}%
\pgfpathlineto{\pgfqpoint{1.612795in}{1.060350in}}%
\pgfpathmoveto{\pgfqpoint{1.608254in}{1.063300in}}%
\pgfpathlineto{\pgfqpoint{1.608254in}{1.063300in}}%
\pgfpathlineto{\pgfqpoint{1.608254in}{1.066249in}}%
\pgfpathlineto{\pgfqpoint{1.612795in}{1.066249in}}%
\pgfpathlineto{\pgfqpoint{1.612795in}{1.063300in}}%
\pgfpathmoveto{\pgfqpoint{1.612795in}{1.060350in}}%
\pgfpathlineto{\pgfqpoint{1.612795in}{1.060350in}}%
\pgfpathlineto{\pgfqpoint{1.612795in}{1.063300in}}%
\pgfpathlineto{\pgfqpoint{1.617336in}{1.063300in}}%
\pgfpathlineto{\pgfqpoint{1.617336in}{1.060350in}}%
\pgfpathmoveto{\pgfqpoint{1.612795in}{1.063300in}}%
\pgfpathlineto{\pgfqpoint{1.612795in}{1.063300in}}%
\pgfpathlineto{\pgfqpoint{1.612795in}{1.066249in}}%
\pgfpathlineto{\pgfqpoint{1.617336in}{1.066249in}}%
\pgfpathlineto{\pgfqpoint{1.617336in}{1.063300in}}%
\pgfpathmoveto{\pgfqpoint{1.617336in}{1.063300in}}%
\pgfpathlineto{\pgfqpoint{1.617336in}{1.063300in}}%
\pgfpathlineto{\pgfqpoint{1.617336in}{1.066249in}}%
\pgfpathlineto{\pgfqpoint{1.621877in}{1.066249in}}%
\pgfpathlineto{\pgfqpoint{1.621877in}{1.063300in}}%
\pgfpathmoveto{\pgfqpoint{1.612795in}{1.066249in}}%
\pgfpathlineto{\pgfqpoint{1.612795in}{1.066249in}}%
\pgfpathlineto{\pgfqpoint{1.612795in}{1.069198in}}%
\pgfpathlineto{\pgfqpoint{1.617336in}{1.069198in}}%
\pgfpathlineto{\pgfqpoint{1.617336in}{1.066249in}}%
\pgfpathmoveto{\pgfqpoint{1.612795in}{1.069198in}}%
\pgfpathlineto{\pgfqpoint{1.612795in}{1.069198in}}%
\pgfpathlineto{\pgfqpoint{1.612795in}{1.072147in}}%
\pgfpathlineto{\pgfqpoint{1.617336in}{1.072147in}}%
\pgfpathlineto{\pgfqpoint{1.617336in}{1.069198in}}%
\pgfpathmoveto{\pgfqpoint{1.617336in}{1.066249in}}%
\pgfpathlineto{\pgfqpoint{1.617336in}{1.066249in}}%
\pgfpathlineto{\pgfqpoint{1.617336in}{1.069198in}}%
\pgfpathlineto{\pgfqpoint{1.621877in}{1.069198in}}%
\pgfpathlineto{\pgfqpoint{1.621877in}{1.066249in}}%
\pgfpathmoveto{\pgfqpoint{1.617336in}{1.069198in}}%
\pgfpathlineto{\pgfqpoint{1.617336in}{1.069198in}}%
\pgfpathlineto{\pgfqpoint{1.617336in}{1.072147in}}%
\pgfpathlineto{\pgfqpoint{1.621877in}{1.072147in}}%
\pgfpathlineto{\pgfqpoint{1.621877in}{1.069198in}}%
\pgfpathmoveto{\pgfqpoint{1.621877in}{1.066249in}}%
\pgfpathlineto{\pgfqpoint{1.621877in}{1.066249in}}%
\pgfpathlineto{\pgfqpoint{1.621877in}{1.069198in}}%
\pgfpathlineto{\pgfqpoint{1.626418in}{1.069198in}}%
\pgfpathlineto{\pgfqpoint{1.626418in}{1.066249in}}%
\pgfpathmoveto{\pgfqpoint{1.621877in}{1.069198in}}%
\pgfpathlineto{\pgfqpoint{1.621877in}{1.069198in}}%
\pgfpathlineto{\pgfqpoint{1.621877in}{1.072147in}}%
\pgfpathlineto{\pgfqpoint{1.626418in}{1.072147in}}%
\pgfpathlineto{\pgfqpoint{1.626418in}{1.069198in}}%
\pgfpathmoveto{\pgfqpoint{1.626418in}{1.069198in}}%
\pgfpathlineto{\pgfqpoint{1.626418in}{1.069198in}}%
\pgfpathlineto{\pgfqpoint{1.626418in}{1.072147in}}%
\pgfpathlineto{\pgfqpoint{1.630958in}{1.072147in}}%
\pgfpathlineto{\pgfqpoint{1.630958in}{1.069198in}}%
\pgfpathmoveto{\pgfqpoint{1.621877in}{1.072147in}}%
\pgfpathlineto{\pgfqpoint{1.621877in}{1.072147in}}%
\pgfpathlineto{\pgfqpoint{1.621877in}{1.075097in}}%
\pgfpathlineto{\pgfqpoint{1.626418in}{1.075097in}}%
\pgfpathlineto{\pgfqpoint{1.626418in}{1.072147in}}%
\pgfpathmoveto{\pgfqpoint{1.621877in}{1.075097in}}%
\pgfpathlineto{\pgfqpoint{1.621877in}{1.075097in}}%
\pgfpathlineto{\pgfqpoint{1.621877in}{1.078046in}}%
\pgfpathlineto{\pgfqpoint{1.626418in}{1.078046in}}%
\pgfpathlineto{\pgfqpoint{1.626418in}{1.075097in}}%
\pgfpathmoveto{\pgfqpoint{1.626418in}{1.072147in}}%
\pgfpathlineto{\pgfqpoint{1.626418in}{1.072147in}}%
\pgfpathlineto{\pgfqpoint{1.626418in}{1.075097in}}%
\pgfpathlineto{\pgfqpoint{1.630958in}{1.075097in}}%
\pgfpathlineto{\pgfqpoint{1.630958in}{1.072147in}}%
\pgfpathmoveto{\pgfqpoint{1.626418in}{1.075097in}}%
\pgfpathlineto{\pgfqpoint{1.626418in}{1.075097in}}%
\pgfpathlineto{\pgfqpoint{1.626418in}{1.078046in}}%
\pgfpathlineto{\pgfqpoint{1.630958in}{1.078046in}}%
\pgfpathlineto{\pgfqpoint{1.630958in}{1.075097in}}%
\pgfpathmoveto{\pgfqpoint{1.630958in}{1.072147in}}%
\pgfpathlineto{\pgfqpoint{1.630958in}{1.072147in}}%
\pgfpathlineto{\pgfqpoint{1.630958in}{1.075097in}}%
\pgfpathlineto{\pgfqpoint{1.635499in}{1.075097in}}%
\pgfpathlineto{\pgfqpoint{1.635499in}{1.072147in}}%
\pgfpathmoveto{\pgfqpoint{1.630958in}{1.075097in}}%
\pgfpathlineto{\pgfqpoint{1.630958in}{1.075097in}}%
\pgfpathlineto{\pgfqpoint{1.630958in}{1.078046in}}%
\pgfpathlineto{\pgfqpoint{1.635499in}{1.078046in}}%
\pgfpathlineto{\pgfqpoint{1.635499in}{1.075097in}}%
\pgfpathmoveto{\pgfqpoint{1.635499in}{1.075097in}}%
\pgfpathlineto{\pgfqpoint{1.635499in}{1.075097in}}%
\pgfpathlineto{\pgfqpoint{1.635499in}{1.078046in}}%
\pgfpathlineto{\pgfqpoint{1.640040in}{1.078046in}}%
\pgfpathlineto{\pgfqpoint{1.640040in}{1.075097in}}%
\pgfpathmoveto{\pgfqpoint{1.630958in}{1.078046in}}%
\pgfpathlineto{\pgfqpoint{1.630958in}{1.078046in}}%
\pgfpathlineto{\pgfqpoint{1.630958in}{1.080995in}}%
\pgfpathlineto{\pgfqpoint{1.635499in}{1.080995in}}%
\pgfpathlineto{\pgfqpoint{1.635499in}{1.078046in}}%
\pgfpathmoveto{\pgfqpoint{1.630958in}{1.080995in}}%
\pgfpathlineto{\pgfqpoint{1.630958in}{1.080995in}}%
\pgfpathlineto{\pgfqpoint{1.630958in}{1.083944in}}%
\pgfpathlineto{\pgfqpoint{1.635499in}{1.083944in}}%
\pgfpathlineto{\pgfqpoint{1.635499in}{1.080995in}}%
\pgfpathmoveto{\pgfqpoint{1.635499in}{1.078046in}}%
\pgfpathlineto{\pgfqpoint{1.635499in}{1.078046in}}%
\pgfpathlineto{\pgfqpoint{1.635499in}{1.080995in}}%
\pgfpathlineto{\pgfqpoint{1.640040in}{1.080995in}}%
\pgfpathlineto{\pgfqpoint{1.640040in}{1.078046in}}%
\pgfpathmoveto{\pgfqpoint{1.635499in}{1.080995in}}%
\pgfpathlineto{\pgfqpoint{1.635499in}{1.080995in}}%
\pgfpathlineto{\pgfqpoint{1.635499in}{1.083944in}}%
\pgfpathlineto{\pgfqpoint{1.640040in}{1.083944in}}%
\pgfpathlineto{\pgfqpoint{1.640040in}{1.080995in}}%
\pgfpathmoveto{\pgfqpoint{1.640040in}{1.078046in}}%
\pgfpathlineto{\pgfqpoint{1.640040in}{1.078046in}}%
\pgfpathlineto{\pgfqpoint{1.640040in}{1.080995in}}%
\pgfpathlineto{\pgfqpoint{1.644581in}{1.080995in}}%
\pgfpathlineto{\pgfqpoint{1.644581in}{1.078046in}}%
\pgfpathmoveto{\pgfqpoint{1.640040in}{1.080995in}}%
\pgfpathlineto{\pgfqpoint{1.640040in}{1.080995in}}%
\pgfpathlineto{\pgfqpoint{1.640040in}{1.083944in}}%
\pgfpathlineto{\pgfqpoint{1.644581in}{1.083944in}}%
\pgfpathlineto{\pgfqpoint{1.644581in}{1.080995in}}%
\pgfpathmoveto{\pgfqpoint{1.644581in}{1.080995in}}%
\pgfpathlineto{\pgfqpoint{1.644581in}{1.080995in}}%
\pgfpathlineto{\pgfqpoint{1.644581in}{1.083944in}}%
\pgfpathlineto{\pgfqpoint{1.649122in}{1.083944in}}%
\pgfpathlineto{\pgfqpoint{1.649122in}{1.080995in}}%
\pgfpathmoveto{\pgfqpoint{1.640040in}{1.083944in}}%
\pgfpathlineto{\pgfqpoint{1.640040in}{1.083944in}}%
\pgfpathlineto{\pgfqpoint{1.640040in}{1.086894in}}%
\pgfpathlineto{\pgfqpoint{1.644581in}{1.086894in}}%
\pgfpathlineto{\pgfqpoint{1.644581in}{1.083944in}}%
\pgfpathmoveto{\pgfqpoint{1.640040in}{1.086894in}}%
\pgfpathlineto{\pgfqpoint{1.640040in}{1.086894in}}%
\pgfpathlineto{\pgfqpoint{1.640040in}{1.089843in}}%
\pgfpathlineto{\pgfqpoint{1.644581in}{1.089843in}}%
\pgfpathlineto{\pgfqpoint{1.644581in}{1.086894in}}%
\pgfpathmoveto{\pgfqpoint{1.644581in}{1.083944in}}%
\pgfpathlineto{\pgfqpoint{1.644581in}{1.083944in}}%
\pgfpathlineto{\pgfqpoint{1.644581in}{1.086894in}}%
\pgfpathlineto{\pgfqpoint{1.649122in}{1.086894in}}%
\pgfpathlineto{\pgfqpoint{1.649122in}{1.083944in}}%
\pgfpathmoveto{\pgfqpoint{1.644581in}{1.086894in}}%
\pgfpathlineto{\pgfqpoint{1.644581in}{1.086894in}}%
\pgfpathlineto{\pgfqpoint{1.644581in}{1.089843in}}%
\pgfpathlineto{\pgfqpoint{1.649122in}{1.089843in}}%
\pgfpathlineto{\pgfqpoint{1.649122in}{1.086894in}}%
\pgfpathmoveto{\pgfqpoint{1.649122in}{1.083944in}}%
\pgfpathlineto{\pgfqpoint{1.649122in}{1.083944in}}%
\pgfpathlineto{\pgfqpoint{1.649122in}{1.086894in}}%
\pgfpathlineto{\pgfqpoint{1.653663in}{1.086894in}}%
\pgfpathlineto{\pgfqpoint{1.653663in}{1.083944in}}%
\pgfpathmoveto{\pgfqpoint{1.649122in}{1.086894in}}%
\pgfpathlineto{\pgfqpoint{1.649122in}{1.086894in}}%
\pgfpathlineto{\pgfqpoint{1.649122in}{1.089843in}}%
\pgfpathlineto{\pgfqpoint{1.653663in}{1.089843in}}%
\pgfpathlineto{\pgfqpoint{1.653663in}{1.086894in}}%
\pgfpathmoveto{\pgfqpoint{1.653663in}{1.086894in}}%
\pgfpathlineto{\pgfqpoint{1.653663in}{1.086894in}}%
\pgfpathlineto{\pgfqpoint{1.653663in}{1.089843in}}%
\pgfpathlineto{\pgfqpoint{1.658204in}{1.089843in}}%
\pgfpathlineto{\pgfqpoint{1.658204in}{1.086894in}}%
\pgfpathmoveto{\pgfqpoint{1.649122in}{1.089843in}}%
\pgfpathlineto{\pgfqpoint{1.649122in}{1.089843in}}%
\pgfpathlineto{\pgfqpoint{1.649122in}{1.092792in}}%
\pgfpathlineto{\pgfqpoint{1.653663in}{1.092792in}}%
\pgfpathlineto{\pgfqpoint{1.653663in}{1.089843in}}%
\pgfpathmoveto{\pgfqpoint{1.649122in}{1.092792in}}%
\pgfpathlineto{\pgfqpoint{1.649122in}{1.092792in}}%
\pgfpathlineto{\pgfqpoint{1.649122in}{1.095741in}}%
\pgfpathlineto{\pgfqpoint{1.653663in}{1.095741in}}%
\pgfpathlineto{\pgfqpoint{1.653663in}{1.092792in}}%
\pgfpathmoveto{\pgfqpoint{1.653663in}{1.089843in}}%
\pgfpathlineto{\pgfqpoint{1.653663in}{1.089843in}}%
\pgfpathlineto{\pgfqpoint{1.653663in}{1.092792in}}%
\pgfpathlineto{\pgfqpoint{1.658204in}{1.092792in}}%
\pgfpathlineto{\pgfqpoint{1.658204in}{1.089843in}}%
\pgfpathmoveto{\pgfqpoint{1.653663in}{1.092792in}}%
\pgfpathlineto{\pgfqpoint{1.653663in}{1.092792in}}%
\pgfpathlineto{\pgfqpoint{1.653663in}{1.095741in}}%
\pgfpathlineto{\pgfqpoint{1.658204in}{1.095741in}}%
\pgfpathlineto{\pgfqpoint{1.658204in}{1.092792in}}%
\pgfpathmoveto{\pgfqpoint{1.658204in}{1.089843in}}%
\pgfpathlineto{\pgfqpoint{1.658204in}{1.089843in}}%
\pgfpathlineto{\pgfqpoint{1.658204in}{1.092792in}}%
\pgfpathlineto{\pgfqpoint{1.662745in}{1.092792in}}%
\pgfpathlineto{\pgfqpoint{1.662745in}{1.089843in}}%
\pgfpathmoveto{\pgfqpoint{1.658204in}{1.092792in}}%
\pgfpathlineto{\pgfqpoint{1.658204in}{1.092792in}}%
\pgfpathlineto{\pgfqpoint{1.658204in}{1.095741in}}%
\pgfpathlineto{\pgfqpoint{1.662745in}{1.095741in}}%
\pgfpathlineto{\pgfqpoint{1.662745in}{1.092792in}}%
\pgfpathmoveto{\pgfqpoint{1.662745in}{1.092792in}}%
\pgfpathlineto{\pgfqpoint{1.662745in}{1.092792in}}%
\pgfpathlineto{\pgfqpoint{1.662745in}{1.095741in}}%
\pgfpathlineto{\pgfqpoint{1.667286in}{1.095741in}}%
\pgfpathlineto{\pgfqpoint{1.667286in}{1.092792in}}%
\pgfpathmoveto{\pgfqpoint{1.658204in}{1.095741in}}%
\pgfpathlineto{\pgfqpoint{1.658204in}{1.095741in}}%
\pgfpathlineto{\pgfqpoint{1.658204in}{1.098691in}}%
\pgfpathlineto{\pgfqpoint{1.662745in}{1.098691in}}%
\pgfpathlineto{\pgfqpoint{1.662745in}{1.095741in}}%
\pgfpathmoveto{\pgfqpoint{1.658204in}{1.098691in}}%
\pgfpathlineto{\pgfqpoint{1.658204in}{1.098691in}}%
\pgfpathlineto{\pgfqpoint{1.658204in}{1.101640in}}%
\pgfpathlineto{\pgfqpoint{1.662745in}{1.101640in}}%
\pgfpathlineto{\pgfqpoint{1.662745in}{1.098691in}}%
\pgfpathmoveto{\pgfqpoint{1.662745in}{1.095741in}}%
\pgfpathlineto{\pgfqpoint{1.662745in}{1.095741in}}%
\pgfpathlineto{\pgfqpoint{1.662745in}{1.098691in}}%
\pgfpathlineto{\pgfqpoint{1.667286in}{1.098691in}}%
\pgfpathlineto{\pgfqpoint{1.667286in}{1.095741in}}%
\pgfpathmoveto{\pgfqpoint{1.662745in}{1.098691in}}%
\pgfpathlineto{\pgfqpoint{1.662745in}{1.098691in}}%
\pgfpathlineto{\pgfqpoint{1.662745in}{1.101640in}}%
\pgfpathlineto{\pgfqpoint{1.667286in}{1.101640in}}%
\pgfpathlineto{\pgfqpoint{1.667286in}{1.098691in}}%
\pgfpathmoveto{\pgfqpoint{1.667286in}{1.095741in}}%
\pgfpathlineto{\pgfqpoint{1.667286in}{1.095741in}}%
\pgfpathlineto{\pgfqpoint{1.667286in}{1.098691in}}%
\pgfpathlineto{\pgfqpoint{1.671827in}{1.098691in}}%
\pgfpathlineto{\pgfqpoint{1.671827in}{1.095741in}}%
\pgfpathmoveto{\pgfqpoint{1.667286in}{1.098691in}}%
\pgfpathlineto{\pgfqpoint{1.667286in}{1.098691in}}%
\pgfpathlineto{\pgfqpoint{1.667286in}{1.101640in}}%
\pgfpathlineto{\pgfqpoint{1.671827in}{1.101640in}}%
\pgfpathlineto{\pgfqpoint{1.671827in}{1.098691in}}%
\pgfpathmoveto{\pgfqpoint{1.671827in}{1.098691in}}%
\pgfpathlineto{\pgfqpoint{1.671827in}{1.098691in}}%
\pgfpathlineto{\pgfqpoint{1.671827in}{1.101640in}}%
\pgfpathlineto{\pgfqpoint{1.676368in}{1.101640in}}%
\pgfpathlineto{\pgfqpoint{1.676368in}{1.098691in}}%
\pgfpathmoveto{\pgfqpoint{1.667286in}{1.101640in}}%
\pgfpathlineto{\pgfqpoint{1.667286in}{1.101640in}}%
\pgfpathlineto{\pgfqpoint{1.667286in}{1.104589in}}%
\pgfpathlineto{\pgfqpoint{1.671827in}{1.104589in}}%
\pgfpathlineto{\pgfqpoint{1.671827in}{1.101640in}}%
\pgfpathmoveto{\pgfqpoint{1.667286in}{1.104589in}}%
\pgfpathlineto{\pgfqpoint{1.667286in}{1.104589in}}%
\pgfpathlineto{\pgfqpoint{1.667286in}{1.107538in}}%
\pgfpathlineto{\pgfqpoint{1.671827in}{1.107538in}}%
\pgfpathlineto{\pgfqpoint{1.671827in}{1.104589in}}%
\pgfpathmoveto{\pgfqpoint{1.671827in}{1.101640in}}%
\pgfpathlineto{\pgfqpoint{1.671827in}{1.101640in}}%
\pgfpathlineto{\pgfqpoint{1.671827in}{1.104589in}}%
\pgfpathlineto{\pgfqpoint{1.676368in}{1.104589in}}%
\pgfpathlineto{\pgfqpoint{1.676368in}{1.101640in}}%
\pgfpathmoveto{\pgfqpoint{1.671827in}{1.104589in}}%
\pgfpathlineto{\pgfqpoint{1.671827in}{1.104589in}}%
\pgfpathlineto{\pgfqpoint{1.671827in}{1.107538in}}%
\pgfpathlineto{\pgfqpoint{1.676368in}{1.107538in}}%
\pgfpathlineto{\pgfqpoint{1.676368in}{1.104589in}}%
\pgfpathmoveto{\pgfqpoint{1.676368in}{1.101640in}}%
\pgfpathlineto{\pgfqpoint{1.676368in}{1.101640in}}%
\pgfpathlineto{\pgfqpoint{1.676368in}{1.104589in}}%
\pgfpathlineto{\pgfqpoint{1.680909in}{1.104589in}}%
\pgfpathlineto{\pgfqpoint{1.680909in}{1.101640in}}%
\pgfpathmoveto{\pgfqpoint{1.676368in}{1.104589in}}%
\pgfpathlineto{\pgfqpoint{1.676368in}{1.104589in}}%
\pgfpathlineto{\pgfqpoint{1.676368in}{1.107538in}}%
\pgfpathlineto{\pgfqpoint{1.680909in}{1.107538in}}%
\pgfpathlineto{\pgfqpoint{1.680909in}{1.104589in}}%
\pgfpathmoveto{\pgfqpoint{1.680909in}{1.104589in}}%
\pgfpathlineto{\pgfqpoint{1.680909in}{1.104589in}}%
\pgfpathlineto{\pgfqpoint{1.680909in}{1.107538in}}%
\pgfpathlineto{\pgfqpoint{1.685449in}{1.107538in}}%
\pgfpathlineto{\pgfqpoint{1.685449in}{1.104589in}}%
\pgfpathmoveto{\pgfqpoint{1.676368in}{1.107538in}}%
\pgfpathlineto{\pgfqpoint{1.676368in}{1.107538in}}%
\pgfpathlineto{\pgfqpoint{1.676368in}{1.110488in}}%
\pgfpathlineto{\pgfqpoint{1.680909in}{1.110488in}}%
\pgfpathlineto{\pgfqpoint{1.680909in}{1.107538in}}%
\pgfpathmoveto{\pgfqpoint{1.676368in}{1.110488in}}%
\pgfpathlineto{\pgfqpoint{1.676368in}{1.110488in}}%
\pgfpathlineto{\pgfqpoint{1.676368in}{1.113437in}}%
\pgfpathlineto{\pgfqpoint{1.680909in}{1.113437in}}%
\pgfpathlineto{\pgfqpoint{1.680909in}{1.110488in}}%
\pgfpathmoveto{\pgfqpoint{1.680909in}{1.107538in}}%
\pgfpathlineto{\pgfqpoint{1.680909in}{1.107538in}}%
\pgfpathlineto{\pgfqpoint{1.680909in}{1.110488in}}%
\pgfpathlineto{\pgfqpoint{1.685449in}{1.110488in}}%
\pgfpathlineto{\pgfqpoint{1.685449in}{1.107538in}}%
\pgfpathmoveto{\pgfqpoint{1.680909in}{1.110488in}}%
\pgfpathlineto{\pgfqpoint{1.680909in}{1.110488in}}%
\pgfpathlineto{\pgfqpoint{1.680909in}{1.113437in}}%
\pgfpathlineto{\pgfqpoint{1.685449in}{1.113437in}}%
\pgfpathlineto{\pgfqpoint{1.685449in}{1.110488in}}%
\pgfpathmoveto{\pgfqpoint{1.685449in}{1.107538in}}%
\pgfpathlineto{\pgfqpoint{1.685449in}{1.107538in}}%
\pgfpathlineto{\pgfqpoint{1.685449in}{1.110488in}}%
\pgfpathlineto{\pgfqpoint{1.689990in}{1.110488in}}%
\pgfpathlineto{\pgfqpoint{1.689990in}{1.107538in}}%
\pgfpathmoveto{\pgfqpoint{1.685449in}{1.110488in}}%
\pgfpathlineto{\pgfqpoint{1.685449in}{1.110488in}}%
\pgfpathlineto{\pgfqpoint{1.685449in}{1.113437in}}%
\pgfpathlineto{\pgfqpoint{1.689990in}{1.113437in}}%
\pgfpathlineto{\pgfqpoint{1.689990in}{1.110488in}}%
\pgfpathmoveto{\pgfqpoint{1.689990in}{1.110488in}}%
\pgfpathlineto{\pgfqpoint{1.689990in}{1.110488in}}%
\pgfpathlineto{\pgfqpoint{1.689990in}{1.113437in}}%
\pgfpathlineto{\pgfqpoint{1.694531in}{1.113437in}}%
\pgfpathlineto{\pgfqpoint{1.694531in}{1.110488in}}%
\pgfpathmoveto{\pgfqpoint{1.685449in}{1.113437in}}%
\pgfpathlineto{\pgfqpoint{1.685449in}{1.113437in}}%
\pgfpathlineto{\pgfqpoint{1.685449in}{1.116386in}}%
\pgfpathlineto{\pgfqpoint{1.689990in}{1.116386in}}%
\pgfpathlineto{\pgfqpoint{1.689990in}{1.113437in}}%
\pgfpathmoveto{\pgfqpoint{1.685449in}{1.116386in}}%
\pgfpathlineto{\pgfqpoint{1.685449in}{1.116386in}}%
\pgfpathlineto{\pgfqpoint{1.685449in}{1.119335in}}%
\pgfpathlineto{\pgfqpoint{1.689990in}{1.119335in}}%
\pgfpathlineto{\pgfqpoint{1.689990in}{1.116386in}}%
\pgfpathmoveto{\pgfqpoint{1.689990in}{1.113437in}}%
\pgfpathlineto{\pgfqpoint{1.689990in}{1.113437in}}%
\pgfpathlineto{\pgfqpoint{1.689990in}{1.116386in}}%
\pgfpathlineto{\pgfqpoint{1.694531in}{1.116386in}}%
\pgfpathlineto{\pgfqpoint{1.694531in}{1.113437in}}%
\pgfpathmoveto{\pgfqpoint{1.689990in}{1.116386in}}%
\pgfpathlineto{\pgfqpoint{1.689990in}{1.116386in}}%
\pgfpathlineto{\pgfqpoint{1.689990in}{1.119335in}}%
\pgfpathlineto{\pgfqpoint{1.694531in}{1.119335in}}%
\pgfpathlineto{\pgfqpoint{1.694531in}{1.116386in}}%
\pgfpathmoveto{\pgfqpoint{1.694531in}{1.113437in}}%
\pgfpathlineto{\pgfqpoint{1.694531in}{1.113437in}}%
\pgfpathlineto{\pgfqpoint{1.694531in}{1.116386in}}%
\pgfpathlineto{\pgfqpoint{1.699072in}{1.116386in}}%
\pgfpathlineto{\pgfqpoint{1.699072in}{1.113437in}}%
\pgfpathmoveto{\pgfqpoint{1.694531in}{1.116386in}}%
\pgfpathlineto{\pgfqpoint{1.694531in}{1.116386in}}%
\pgfpathlineto{\pgfqpoint{1.694531in}{1.119335in}}%
\pgfpathlineto{\pgfqpoint{1.699072in}{1.119335in}}%
\pgfpathlineto{\pgfqpoint{1.699072in}{1.116386in}}%
\pgfpathmoveto{\pgfqpoint{1.699072in}{1.116386in}}%
\pgfpathlineto{\pgfqpoint{1.699072in}{1.116386in}}%
\pgfpathlineto{\pgfqpoint{1.699072in}{1.119335in}}%
\pgfpathlineto{\pgfqpoint{1.703613in}{1.119335in}}%
\pgfpathlineto{\pgfqpoint{1.703613in}{1.116386in}}%
\pgfpathmoveto{\pgfqpoint{1.694531in}{1.119335in}}%
\pgfpathlineto{\pgfqpoint{1.694531in}{1.119335in}}%
\pgfpathlineto{\pgfqpoint{1.694531in}{1.122285in}}%
\pgfpathlineto{\pgfqpoint{1.699072in}{1.122285in}}%
\pgfpathlineto{\pgfqpoint{1.699072in}{1.119335in}}%
\pgfpathmoveto{\pgfqpoint{1.694531in}{1.122285in}}%
\pgfpathlineto{\pgfqpoint{1.694531in}{1.122285in}}%
\pgfpathlineto{\pgfqpoint{1.694531in}{1.125234in}}%
\pgfpathlineto{\pgfqpoint{1.699072in}{1.125234in}}%
\pgfpathlineto{\pgfqpoint{1.699072in}{1.122285in}}%
\pgfpathmoveto{\pgfqpoint{1.699072in}{1.119335in}}%
\pgfpathlineto{\pgfqpoint{1.699072in}{1.119335in}}%
\pgfpathlineto{\pgfqpoint{1.699072in}{1.122285in}}%
\pgfpathlineto{\pgfqpoint{1.703613in}{1.122285in}}%
\pgfpathlineto{\pgfqpoint{1.703613in}{1.119335in}}%
\pgfpathmoveto{\pgfqpoint{1.699072in}{1.122285in}}%
\pgfpathlineto{\pgfqpoint{1.699072in}{1.122285in}}%
\pgfpathlineto{\pgfqpoint{1.699072in}{1.125234in}}%
\pgfpathlineto{\pgfqpoint{1.703613in}{1.125234in}}%
\pgfpathlineto{\pgfqpoint{1.703613in}{1.122285in}}%
\pgfpathmoveto{\pgfqpoint{1.703613in}{1.119335in}}%
\pgfpathlineto{\pgfqpoint{1.703613in}{1.119335in}}%
\pgfpathlineto{\pgfqpoint{1.703613in}{1.122285in}}%
\pgfpathlineto{\pgfqpoint{1.708154in}{1.122285in}}%
\pgfpathlineto{\pgfqpoint{1.708154in}{1.119335in}}%
\pgfpathmoveto{\pgfqpoint{1.703613in}{1.122285in}}%
\pgfpathlineto{\pgfqpoint{1.703613in}{1.122285in}}%
\pgfpathlineto{\pgfqpoint{1.703613in}{1.125234in}}%
\pgfpathlineto{\pgfqpoint{1.708154in}{1.125234in}}%
\pgfpathlineto{\pgfqpoint{1.708154in}{1.122285in}}%
\pgfpathmoveto{\pgfqpoint{1.708154in}{1.122285in}}%
\pgfpathlineto{\pgfqpoint{1.708154in}{1.122285in}}%
\pgfpathlineto{\pgfqpoint{1.708154in}{1.125234in}}%
\pgfpathlineto{\pgfqpoint{1.712695in}{1.125234in}}%
\pgfpathlineto{\pgfqpoint{1.712695in}{1.122285in}}%
\pgfpathmoveto{\pgfqpoint{1.703613in}{1.125234in}}%
\pgfpathlineto{\pgfqpoint{1.703613in}{1.125234in}}%
\pgfpathlineto{\pgfqpoint{1.703613in}{1.128183in}}%
\pgfpathlineto{\pgfqpoint{1.708154in}{1.128183in}}%
\pgfpathlineto{\pgfqpoint{1.708154in}{1.125234in}}%
\pgfpathmoveto{\pgfqpoint{1.703613in}{1.128183in}}%
\pgfpathlineto{\pgfqpoint{1.703613in}{1.128183in}}%
\pgfpathlineto{\pgfqpoint{1.703613in}{1.131132in}}%
\pgfpathlineto{\pgfqpoint{1.708154in}{1.131132in}}%
\pgfpathlineto{\pgfqpoint{1.708154in}{1.128183in}}%
\pgfpathmoveto{\pgfqpoint{1.708154in}{1.125234in}}%
\pgfpathlineto{\pgfqpoint{1.708154in}{1.125234in}}%
\pgfpathlineto{\pgfqpoint{1.708154in}{1.128183in}}%
\pgfpathlineto{\pgfqpoint{1.712695in}{1.128183in}}%
\pgfpathlineto{\pgfqpoint{1.712695in}{1.125234in}}%
\pgfpathmoveto{\pgfqpoint{1.708154in}{1.128183in}}%
\pgfpathlineto{\pgfqpoint{1.708154in}{1.128183in}}%
\pgfpathlineto{\pgfqpoint{1.708154in}{1.131132in}}%
\pgfpathlineto{\pgfqpoint{1.712695in}{1.131132in}}%
\pgfpathlineto{\pgfqpoint{1.712695in}{1.128183in}}%
\pgfpathmoveto{\pgfqpoint{1.712695in}{1.125234in}}%
\pgfpathlineto{\pgfqpoint{1.712695in}{1.125234in}}%
\pgfpathlineto{\pgfqpoint{1.712695in}{1.128183in}}%
\pgfpathlineto{\pgfqpoint{1.717236in}{1.128183in}}%
\pgfpathlineto{\pgfqpoint{1.717236in}{1.125234in}}%
\pgfpathmoveto{\pgfqpoint{1.712695in}{1.128183in}}%
\pgfpathlineto{\pgfqpoint{1.712695in}{1.128183in}}%
\pgfpathlineto{\pgfqpoint{1.712695in}{1.131132in}}%
\pgfpathlineto{\pgfqpoint{1.717236in}{1.131132in}}%
\pgfpathlineto{\pgfqpoint{1.717236in}{1.128183in}}%
\pgfpathmoveto{\pgfqpoint{1.717236in}{1.128183in}}%
\pgfpathlineto{\pgfqpoint{1.717236in}{1.128183in}}%
\pgfpathlineto{\pgfqpoint{1.717236in}{1.131132in}}%
\pgfpathlineto{\pgfqpoint{1.721777in}{1.131132in}}%
\pgfpathlineto{\pgfqpoint{1.721777in}{1.128183in}}%
\pgfpathmoveto{\pgfqpoint{1.721777in}{1.128183in}}%
\pgfpathlineto{\pgfqpoint{1.721777in}{1.128183in}}%
\pgfpathlineto{\pgfqpoint{1.721777in}{1.131132in}}%
\pgfpathlineto{\pgfqpoint{1.726318in}{1.131132in}}%
\pgfpathlineto{\pgfqpoint{1.726318in}{1.128183in}}%
\pgfpathmoveto{\pgfqpoint{1.721777in}{1.131132in}}%
\pgfpathlineto{\pgfqpoint{1.721777in}{1.131132in}}%
\pgfpathlineto{\pgfqpoint{1.721777in}{1.134082in}}%
\pgfpathlineto{\pgfqpoint{1.726318in}{1.134082in}}%
\pgfpathlineto{\pgfqpoint{1.726318in}{1.131132in}}%
\pgfpathmoveto{\pgfqpoint{1.721777in}{1.134082in}}%
\pgfpathlineto{\pgfqpoint{1.721777in}{1.134082in}}%
\pgfpathlineto{\pgfqpoint{1.721777in}{1.137031in}}%
\pgfpathlineto{\pgfqpoint{1.726318in}{1.137031in}}%
\pgfpathlineto{\pgfqpoint{1.726318in}{1.134082in}}%
\pgfpathmoveto{\pgfqpoint{1.726318in}{1.131132in}}%
\pgfpathlineto{\pgfqpoint{1.726318in}{1.131132in}}%
\pgfpathlineto{\pgfqpoint{1.726318in}{1.134082in}}%
\pgfpathlineto{\pgfqpoint{1.730859in}{1.134082in}}%
\pgfpathlineto{\pgfqpoint{1.730859in}{1.131132in}}%
\pgfpathmoveto{\pgfqpoint{1.726318in}{1.134082in}}%
\pgfpathlineto{\pgfqpoint{1.726318in}{1.134082in}}%
\pgfpathlineto{\pgfqpoint{1.726318in}{1.137031in}}%
\pgfpathlineto{\pgfqpoint{1.730859in}{1.137031in}}%
\pgfpathlineto{\pgfqpoint{1.730859in}{1.134082in}}%
\pgfpathmoveto{\pgfqpoint{1.730859in}{1.134082in}}%
\pgfpathlineto{\pgfqpoint{1.730859in}{1.134082in}}%
\pgfpathlineto{\pgfqpoint{1.730859in}{1.137031in}}%
\pgfpathlineto{\pgfqpoint{1.735400in}{1.137031in}}%
\pgfpathlineto{\pgfqpoint{1.735400in}{1.134082in}}%
\pgfpathmoveto{\pgfqpoint{1.730859in}{1.137031in}}%
\pgfpathlineto{\pgfqpoint{1.730859in}{1.137031in}}%
\pgfpathlineto{\pgfqpoint{1.730859in}{1.139980in}}%
\pgfpathlineto{\pgfqpoint{1.735400in}{1.139980in}}%
\pgfpathlineto{\pgfqpoint{1.735400in}{1.137031in}}%
\pgfpathmoveto{\pgfqpoint{1.730859in}{1.139980in}}%
\pgfpathlineto{\pgfqpoint{1.730859in}{1.139980in}}%
\pgfpathlineto{\pgfqpoint{1.730859in}{1.142929in}}%
\pgfpathlineto{\pgfqpoint{1.735400in}{1.142929in}}%
\pgfpathlineto{\pgfqpoint{1.735400in}{1.139980in}}%
\pgfpathmoveto{\pgfqpoint{1.735400in}{1.137031in}}%
\pgfpathlineto{\pgfqpoint{1.735400in}{1.137031in}}%
\pgfpathlineto{\pgfqpoint{1.735400in}{1.139980in}}%
\pgfpathlineto{\pgfqpoint{1.739940in}{1.139980in}}%
\pgfpathlineto{\pgfqpoint{1.739940in}{1.137031in}}%
\pgfpathmoveto{\pgfqpoint{1.735400in}{1.139980in}}%
\pgfpathlineto{\pgfqpoint{1.735400in}{1.139980in}}%
\pgfpathlineto{\pgfqpoint{1.735400in}{1.142929in}}%
\pgfpathlineto{\pgfqpoint{1.739940in}{1.142929in}}%
\pgfpathlineto{\pgfqpoint{1.739940in}{1.139980in}}%
\pgfpathmoveto{\pgfqpoint{1.739940in}{1.139980in}}%
\pgfpathlineto{\pgfqpoint{1.739940in}{1.139980in}}%
\pgfpathlineto{\pgfqpoint{1.739940in}{1.142929in}}%
\pgfpathlineto{\pgfqpoint{1.744481in}{1.142929in}}%
\pgfpathlineto{\pgfqpoint{1.744481in}{1.139980in}}%
\pgfpathmoveto{\pgfqpoint{1.739940in}{1.142929in}}%
\pgfpathlineto{\pgfqpoint{1.739940in}{1.142929in}}%
\pgfpathlineto{\pgfqpoint{1.739940in}{1.145879in}}%
\pgfpathlineto{\pgfqpoint{1.744481in}{1.145879in}}%
\pgfpathlineto{\pgfqpoint{1.744481in}{1.142929in}}%
\pgfpathmoveto{\pgfqpoint{1.739940in}{1.145879in}}%
\pgfpathlineto{\pgfqpoint{1.739940in}{1.145879in}}%
\pgfpathlineto{\pgfqpoint{1.739940in}{1.148828in}}%
\pgfpathlineto{\pgfqpoint{1.744481in}{1.148828in}}%
\pgfpathlineto{\pgfqpoint{1.744481in}{1.145879in}}%
\pgfpathmoveto{\pgfqpoint{1.744481in}{1.142929in}}%
\pgfpathlineto{\pgfqpoint{1.744481in}{1.142929in}}%
\pgfpathlineto{\pgfqpoint{1.744481in}{1.145879in}}%
\pgfpathlineto{\pgfqpoint{1.749022in}{1.145879in}}%
\pgfpathlineto{\pgfqpoint{1.749022in}{1.142929in}}%
\pgfpathmoveto{\pgfqpoint{1.744481in}{1.145879in}}%
\pgfpathlineto{\pgfqpoint{1.744481in}{1.145879in}}%
\pgfpathlineto{\pgfqpoint{1.744481in}{1.148828in}}%
\pgfpathlineto{\pgfqpoint{1.749022in}{1.148828in}}%
\pgfpathlineto{\pgfqpoint{1.749022in}{1.145879in}}%
\pgfpathmoveto{\pgfqpoint{1.749022in}{1.145879in}}%
\pgfpathlineto{\pgfqpoint{1.749022in}{1.145879in}}%
\pgfpathlineto{\pgfqpoint{1.749022in}{1.148828in}}%
\pgfpathlineto{\pgfqpoint{1.753563in}{1.148828in}}%
\pgfpathlineto{\pgfqpoint{1.753563in}{1.145879in}}%
\pgfpathmoveto{\pgfqpoint{1.749022in}{1.148828in}}%
\pgfpathlineto{\pgfqpoint{1.749022in}{1.148828in}}%
\pgfpathlineto{\pgfqpoint{1.749022in}{1.151777in}}%
\pgfpathlineto{\pgfqpoint{1.753563in}{1.151777in}}%
\pgfpathlineto{\pgfqpoint{1.753563in}{1.148828in}}%
\pgfpathmoveto{\pgfqpoint{1.749022in}{1.151777in}}%
\pgfpathlineto{\pgfqpoint{1.749022in}{1.151777in}}%
\pgfpathlineto{\pgfqpoint{1.749022in}{1.154727in}}%
\pgfpathlineto{\pgfqpoint{1.753563in}{1.154727in}}%
\pgfpathlineto{\pgfqpoint{1.753563in}{1.151777in}}%
\pgfpathmoveto{\pgfqpoint{1.753563in}{1.148828in}}%
\pgfpathlineto{\pgfqpoint{1.753563in}{1.148828in}}%
\pgfpathlineto{\pgfqpoint{1.753563in}{1.151777in}}%
\pgfpathlineto{\pgfqpoint{1.758104in}{1.151777in}}%
\pgfpathlineto{\pgfqpoint{1.758104in}{1.148828in}}%
\pgfpathmoveto{\pgfqpoint{1.753563in}{1.151777in}}%
\pgfpathlineto{\pgfqpoint{1.753563in}{1.151777in}}%
\pgfpathlineto{\pgfqpoint{1.753563in}{1.154727in}}%
\pgfpathlineto{\pgfqpoint{1.758104in}{1.154727in}}%
\pgfpathlineto{\pgfqpoint{1.758104in}{1.151777in}}%
\pgfpathmoveto{\pgfqpoint{1.758104in}{1.151777in}}%
\pgfpathlineto{\pgfqpoint{1.758104in}{1.151777in}}%
\pgfpathlineto{\pgfqpoint{1.758104in}{1.154727in}}%
\pgfpathlineto{\pgfqpoint{1.762645in}{1.154727in}}%
\pgfpathlineto{\pgfqpoint{1.762645in}{1.151777in}}%
\pgfpathmoveto{\pgfqpoint{1.758104in}{1.154727in}}%
\pgfpathlineto{\pgfqpoint{1.758104in}{1.154727in}}%
\pgfpathlineto{\pgfqpoint{1.758104in}{1.157676in}}%
\pgfpathlineto{\pgfqpoint{1.762645in}{1.157676in}}%
\pgfpathlineto{\pgfqpoint{1.762645in}{1.154727in}}%
\pgfpathmoveto{\pgfqpoint{1.758104in}{1.157676in}}%
\pgfpathlineto{\pgfqpoint{1.758104in}{1.157676in}}%
\pgfpathlineto{\pgfqpoint{1.758104in}{1.160625in}}%
\pgfpathlineto{\pgfqpoint{1.762645in}{1.160625in}}%
\pgfpathlineto{\pgfqpoint{1.762645in}{1.157676in}}%
\pgfpathmoveto{\pgfqpoint{1.762645in}{1.154727in}}%
\pgfpathlineto{\pgfqpoint{1.762645in}{1.154727in}}%
\pgfpathlineto{\pgfqpoint{1.762645in}{1.157676in}}%
\pgfpathlineto{\pgfqpoint{1.767186in}{1.157676in}}%
\pgfpathlineto{\pgfqpoint{1.767186in}{1.154727in}}%
\pgfpathmoveto{\pgfqpoint{1.762645in}{1.157676in}}%
\pgfpathlineto{\pgfqpoint{1.762645in}{1.157676in}}%
\pgfpathlineto{\pgfqpoint{1.762645in}{1.160625in}}%
\pgfpathlineto{\pgfqpoint{1.767186in}{1.160625in}}%
\pgfpathlineto{\pgfqpoint{1.767186in}{1.157676in}}%
\pgfpathmoveto{\pgfqpoint{1.767186in}{1.157676in}}%
\pgfpathlineto{\pgfqpoint{1.767186in}{1.157676in}}%
\pgfpathlineto{\pgfqpoint{1.767186in}{1.160625in}}%
\pgfpathlineto{\pgfqpoint{1.771727in}{1.160625in}}%
\pgfpathlineto{\pgfqpoint{1.771727in}{1.157676in}}%
\pgfpathmoveto{\pgfqpoint{1.767186in}{1.160625in}}%
\pgfpathlineto{\pgfqpoint{1.767186in}{1.160625in}}%
\pgfpathlineto{\pgfqpoint{1.767186in}{1.163574in}}%
\pgfpathlineto{\pgfqpoint{1.771727in}{1.163574in}}%
\pgfpathlineto{\pgfqpoint{1.771727in}{1.160625in}}%
\pgfpathmoveto{\pgfqpoint{1.767186in}{1.163574in}}%
\pgfpathlineto{\pgfqpoint{1.767186in}{1.163574in}}%
\pgfpathlineto{\pgfqpoint{1.767186in}{1.166523in}}%
\pgfpathlineto{\pgfqpoint{1.771727in}{1.166523in}}%
\pgfpathlineto{\pgfqpoint{1.771727in}{1.163574in}}%
\pgfpathmoveto{\pgfqpoint{1.771727in}{1.160625in}}%
\pgfpathlineto{\pgfqpoint{1.771727in}{1.160625in}}%
\pgfpathlineto{\pgfqpoint{1.771727in}{1.163574in}}%
\pgfpathlineto{\pgfqpoint{1.776268in}{1.163574in}}%
\pgfpathlineto{\pgfqpoint{1.776268in}{1.160625in}}%
\pgfpathmoveto{\pgfqpoint{1.771727in}{1.163574in}}%
\pgfpathlineto{\pgfqpoint{1.771727in}{1.163574in}}%
\pgfpathlineto{\pgfqpoint{1.771727in}{1.166523in}}%
\pgfpathlineto{\pgfqpoint{1.776268in}{1.166523in}}%
\pgfpathlineto{\pgfqpoint{1.776268in}{1.163574in}}%
\pgfpathmoveto{\pgfqpoint{1.776268in}{1.163574in}}%
\pgfpathlineto{\pgfqpoint{1.776268in}{1.163574in}}%
\pgfpathlineto{\pgfqpoint{1.776268in}{1.166523in}}%
\pgfpathlineto{\pgfqpoint{1.780809in}{1.166523in}}%
\pgfpathlineto{\pgfqpoint{1.780809in}{1.163574in}}%
\pgfpathmoveto{\pgfqpoint{1.776268in}{1.166523in}}%
\pgfpathlineto{\pgfqpoint{1.776268in}{1.166523in}}%
\pgfpathlineto{\pgfqpoint{1.776268in}{1.169472in}}%
\pgfpathlineto{\pgfqpoint{1.780809in}{1.169472in}}%
\pgfpathlineto{\pgfqpoint{1.780809in}{1.166523in}}%
\pgfpathmoveto{\pgfqpoint{1.776268in}{1.169472in}}%
\pgfpathlineto{\pgfqpoint{1.776268in}{1.169472in}}%
\pgfpathlineto{\pgfqpoint{1.776268in}{1.172422in}}%
\pgfpathlineto{\pgfqpoint{1.780809in}{1.172422in}}%
\pgfpathlineto{\pgfqpoint{1.780809in}{1.169472in}}%
\pgfpathmoveto{\pgfqpoint{1.780809in}{1.166523in}}%
\pgfpathlineto{\pgfqpoint{1.780809in}{1.166523in}}%
\pgfpathlineto{\pgfqpoint{1.780809in}{1.169472in}}%
\pgfpathlineto{\pgfqpoint{1.785350in}{1.169472in}}%
\pgfpathlineto{\pgfqpoint{1.785350in}{1.166523in}}%
\pgfpathmoveto{\pgfqpoint{1.780809in}{1.169472in}}%
\pgfpathlineto{\pgfqpoint{1.780809in}{1.169472in}}%
\pgfpathlineto{\pgfqpoint{1.780809in}{1.172422in}}%
\pgfpathlineto{\pgfqpoint{1.785350in}{1.172422in}}%
\pgfpathlineto{\pgfqpoint{1.785350in}{1.169472in}}%
\pgfpathmoveto{\pgfqpoint{1.785350in}{1.169472in}}%
\pgfpathlineto{\pgfqpoint{1.785350in}{1.169472in}}%
\pgfpathlineto{\pgfqpoint{1.785350in}{1.172422in}}%
\pgfpathlineto{\pgfqpoint{1.789891in}{1.172422in}}%
\pgfpathlineto{\pgfqpoint{1.789891in}{1.169472in}}%
\pgfpathmoveto{\pgfqpoint{1.785350in}{1.172422in}}%
\pgfpathlineto{\pgfqpoint{1.785350in}{1.172422in}}%
\pgfpathlineto{\pgfqpoint{1.785350in}{1.175371in}}%
\pgfpathlineto{\pgfqpoint{1.789891in}{1.175371in}}%
\pgfpathlineto{\pgfqpoint{1.789891in}{1.172422in}}%
\pgfpathmoveto{\pgfqpoint{1.785350in}{1.175371in}}%
\pgfpathlineto{\pgfqpoint{1.785350in}{1.175371in}}%
\pgfpathlineto{\pgfqpoint{1.785350in}{1.178320in}}%
\pgfpathlineto{\pgfqpoint{1.789891in}{1.178320in}}%
\pgfpathlineto{\pgfqpoint{1.789891in}{1.175371in}}%
\pgfpathmoveto{\pgfqpoint{1.789891in}{1.172422in}}%
\pgfpathlineto{\pgfqpoint{1.789891in}{1.172422in}}%
\pgfpathlineto{\pgfqpoint{1.789891in}{1.175371in}}%
\pgfpathlineto{\pgfqpoint{1.794432in}{1.175371in}}%
\pgfpathlineto{\pgfqpoint{1.794432in}{1.172422in}}%
\pgfpathmoveto{\pgfqpoint{1.789891in}{1.175371in}}%
\pgfpathlineto{\pgfqpoint{1.789891in}{1.175371in}}%
\pgfpathlineto{\pgfqpoint{1.789891in}{1.178320in}}%
\pgfpathlineto{\pgfqpoint{1.794432in}{1.178320in}}%
\pgfpathlineto{\pgfqpoint{1.794432in}{1.175371in}}%
\pgfpathmoveto{\pgfqpoint{1.794432in}{1.175371in}}%
\pgfpathlineto{\pgfqpoint{1.794432in}{1.175371in}}%
\pgfpathlineto{\pgfqpoint{1.794432in}{1.178320in}}%
\pgfpathlineto{\pgfqpoint{1.798972in}{1.178320in}}%
\pgfpathlineto{\pgfqpoint{1.798972in}{1.175371in}}%
\pgfpathmoveto{\pgfqpoint{1.794432in}{1.178320in}}%
\pgfpathlineto{\pgfqpoint{1.794432in}{1.178320in}}%
\pgfpathlineto{\pgfqpoint{1.794432in}{1.181269in}}%
\pgfpathlineto{\pgfqpoint{1.798972in}{1.181269in}}%
\pgfpathlineto{\pgfqpoint{1.798972in}{1.178320in}}%
\pgfpathmoveto{\pgfqpoint{1.794432in}{1.181269in}}%
\pgfpathlineto{\pgfqpoint{1.794432in}{1.181269in}}%
\pgfpathlineto{\pgfqpoint{1.794432in}{1.184218in}}%
\pgfpathlineto{\pgfqpoint{1.798972in}{1.184218in}}%
\pgfpathlineto{\pgfqpoint{1.798972in}{1.181269in}}%
\pgfpathmoveto{\pgfqpoint{1.798972in}{1.178320in}}%
\pgfpathlineto{\pgfqpoint{1.798972in}{1.178320in}}%
\pgfpathlineto{\pgfqpoint{1.798972in}{1.181269in}}%
\pgfpathlineto{\pgfqpoint{1.803513in}{1.181269in}}%
\pgfpathlineto{\pgfqpoint{1.803513in}{1.178320in}}%
\pgfpathmoveto{\pgfqpoint{1.798972in}{1.181269in}}%
\pgfpathlineto{\pgfqpoint{1.798972in}{1.181269in}}%
\pgfpathlineto{\pgfqpoint{1.798972in}{1.184218in}}%
\pgfpathlineto{\pgfqpoint{1.803513in}{1.184218in}}%
\pgfpathlineto{\pgfqpoint{1.803513in}{1.181269in}}%
\pgfpathmoveto{\pgfqpoint{1.803513in}{1.181269in}}%
\pgfpathlineto{\pgfqpoint{1.803513in}{1.181269in}}%
\pgfpathlineto{\pgfqpoint{1.803513in}{1.184218in}}%
\pgfpathlineto{\pgfqpoint{1.808054in}{1.184218in}}%
\pgfpathlineto{\pgfqpoint{1.808054in}{1.181269in}}%
\pgfpathmoveto{\pgfqpoint{1.803513in}{1.184218in}}%
\pgfpathlineto{\pgfqpoint{1.803513in}{1.184218in}}%
\pgfpathlineto{\pgfqpoint{1.803513in}{1.187167in}}%
\pgfpathlineto{\pgfqpoint{1.808054in}{1.187167in}}%
\pgfpathlineto{\pgfqpoint{1.808054in}{1.184218in}}%
\pgfpathmoveto{\pgfqpoint{1.803513in}{1.187167in}}%
\pgfpathlineto{\pgfqpoint{1.803513in}{1.187167in}}%
\pgfpathlineto{\pgfqpoint{1.803513in}{1.190117in}}%
\pgfpathlineto{\pgfqpoint{1.808054in}{1.190117in}}%
\pgfpathlineto{\pgfqpoint{1.808054in}{1.187167in}}%
\pgfpathmoveto{\pgfqpoint{1.808054in}{1.184218in}}%
\pgfpathlineto{\pgfqpoint{1.808054in}{1.184218in}}%
\pgfpathlineto{\pgfqpoint{1.808054in}{1.187167in}}%
\pgfpathlineto{\pgfqpoint{1.812595in}{1.187167in}}%
\pgfpathlineto{\pgfqpoint{1.812595in}{1.184218in}}%
\pgfpathmoveto{\pgfqpoint{1.808054in}{1.187167in}}%
\pgfpathlineto{\pgfqpoint{1.808054in}{1.187167in}}%
\pgfpathlineto{\pgfqpoint{1.808054in}{1.190117in}}%
\pgfpathlineto{\pgfqpoint{1.812595in}{1.190117in}}%
\pgfpathlineto{\pgfqpoint{1.812595in}{1.187167in}}%
\pgfpathmoveto{\pgfqpoint{1.812595in}{1.187167in}}%
\pgfpathlineto{\pgfqpoint{1.812595in}{1.187167in}}%
\pgfpathlineto{\pgfqpoint{1.812595in}{1.190117in}}%
\pgfpathlineto{\pgfqpoint{1.817136in}{1.190117in}}%
\pgfpathlineto{\pgfqpoint{1.817136in}{1.187167in}}%
\pgfpathmoveto{\pgfqpoint{1.812595in}{1.190117in}}%
\pgfpathlineto{\pgfqpoint{1.812595in}{1.190117in}}%
\pgfpathlineto{\pgfqpoint{1.812595in}{1.193066in}}%
\pgfpathlineto{\pgfqpoint{1.817136in}{1.193066in}}%
\pgfpathlineto{\pgfqpoint{1.817136in}{1.190117in}}%
\pgfpathmoveto{\pgfqpoint{1.812595in}{1.193066in}}%
\pgfpathlineto{\pgfqpoint{1.812595in}{1.193066in}}%
\pgfpathlineto{\pgfqpoint{1.812595in}{1.196015in}}%
\pgfpathlineto{\pgfqpoint{1.817136in}{1.196015in}}%
\pgfpathlineto{\pgfqpoint{1.817136in}{1.193066in}}%
\pgfpathmoveto{\pgfqpoint{1.817136in}{1.190117in}}%
\pgfpathlineto{\pgfqpoint{1.817136in}{1.190117in}}%
\pgfpathlineto{\pgfqpoint{1.817136in}{1.193066in}}%
\pgfpathlineto{\pgfqpoint{1.821677in}{1.193066in}}%
\pgfpathlineto{\pgfqpoint{1.821677in}{1.190117in}}%
\pgfpathmoveto{\pgfqpoint{1.817136in}{1.193066in}}%
\pgfpathlineto{\pgfqpoint{1.817136in}{1.193066in}}%
\pgfpathlineto{\pgfqpoint{1.817136in}{1.196015in}}%
\pgfpathlineto{\pgfqpoint{1.821677in}{1.196015in}}%
\pgfpathlineto{\pgfqpoint{1.821677in}{1.193066in}}%
\pgfpathmoveto{\pgfqpoint{1.821677in}{1.193066in}}%
\pgfpathlineto{\pgfqpoint{1.821677in}{1.193066in}}%
\pgfpathlineto{\pgfqpoint{1.821677in}{1.196015in}}%
\pgfpathlineto{\pgfqpoint{1.826218in}{1.196015in}}%
\pgfpathlineto{\pgfqpoint{1.826218in}{1.193066in}}%
\pgfpathmoveto{\pgfqpoint{1.821677in}{1.196015in}}%
\pgfpathlineto{\pgfqpoint{1.821677in}{1.196015in}}%
\pgfpathlineto{\pgfqpoint{1.821677in}{1.198964in}}%
\pgfpathlineto{\pgfqpoint{1.826218in}{1.198964in}}%
\pgfpathlineto{\pgfqpoint{1.826218in}{1.196015in}}%
\pgfpathmoveto{\pgfqpoint{1.821677in}{1.198964in}}%
\pgfpathlineto{\pgfqpoint{1.821677in}{1.198964in}}%
\pgfpathlineto{\pgfqpoint{1.821677in}{1.201913in}}%
\pgfpathlineto{\pgfqpoint{1.826218in}{1.201913in}}%
\pgfpathlineto{\pgfqpoint{1.826218in}{1.198964in}}%
\pgfpathmoveto{\pgfqpoint{1.826218in}{1.196015in}}%
\pgfpathlineto{\pgfqpoint{1.826218in}{1.196015in}}%
\pgfpathlineto{\pgfqpoint{1.826218in}{1.198964in}}%
\pgfpathlineto{\pgfqpoint{1.830759in}{1.198964in}}%
\pgfpathlineto{\pgfqpoint{1.830759in}{1.196015in}}%
\pgfpathmoveto{\pgfqpoint{1.826218in}{1.198964in}}%
\pgfpathlineto{\pgfqpoint{1.826218in}{1.198964in}}%
\pgfpathlineto{\pgfqpoint{1.826218in}{1.201913in}}%
\pgfpathlineto{\pgfqpoint{1.830759in}{1.201913in}}%
\pgfpathlineto{\pgfqpoint{1.830759in}{1.198964in}}%
\pgfpathmoveto{\pgfqpoint{1.830759in}{1.198964in}}%
\pgfpathlineto{\pgfqpoint{1.830759in}{1.198964in}}%
\pgfpathlineto{\pgfqpoint{1.830759in}{1.201913in}}%
\pgfpathlineto{\pgfqpoint{1.835300in}{1.201913in}}%
\pgfpathlineto{\pgfqpoint{1.835300in}{1.198964in}}%
\pgfpathmoveto{\pgfqpoint{1.830759in}{1.201913in}}%
\pgfpathlineto{\pgfqpoint{1.830759in}{1.201913in}}%
\pgfpathlineto{\pgfqpoint{1.830759in}{1.204862in}}%
\pgfpathlineto{\pgfqpoint{1.835300in}{1.204862in}}%
\pgfpathlineto{\pgfqpoint{1.835300in}{1.201913in}}%
\pgfpathmoveto{\pgfqpoint{1.830759in}{1.204862in}}%
\pgfpathlineto{\pgfqpoint{1.830759in}{1.204862in}}%
\pgfpathlineto{\pgfqpoint{1.830759in}{1.207812in}}%
\pgfpathlineto{\pgfqpoint{1.835300in}{1.207812in}}%
\pgfpathlineto{\pgfqpoint{1.835300in}{1.204862in}}%
\pgfpathmoveto{\pgfqpoint{1.835300in}{1.201913in}}%
\pgfpathlineto{\pgfqpoint{1.835300in}{1.201913in}}%
\pgfpathlineto{\pgfqpoint{1.835300in}{1.204862in}}%
\pgfpathlineto{\pgfqpoint{1.839841in}{1.204862in}}%
\pgfpathlineto{\pgfqpoint{1.839841in}{1.201913in}}%
\pgfpathmoveto{\pgfqpoint{1.835300in}{1.204862in}}%
\pgfpathlineto{\pgfqpoint{1.835300in}{1.204862in}}%
\pgfpathlineto{\pgfqpoint{1.835300in}{1.207812in}}%
\pgfpathlineto{\pgfqpoint{1.839841in}{1.207812in}}%
\pgfpathlineto{\pgfqpoint{1.839841in}{1.204862in}}%
\pgfpathmoveto{\pgfqpoint{1.839841in}{1.204862in}}%
\pgfpathlineto{\pgfqpoint{1.839841in}{1.204862in}}%
\pgfpathlineto{\pgfqpoint{1.839841in}{1.207812in}}%
\pgfpathlineto{\pgfqpoint{1.844382in}{1.207812in}}%
\pgfpathlineto{\pgfqpoint{1.844382in}{1.204862in}}%
\pgfpathmoveto{\pgfqpoint{1.839841in}{1.207812in}}%
\pgfpathlineto{\pgfqpoint{1.839841in}{1.207812in}}%
\pgfpathlineto{\pgfqpoint{1.839841in}{1.210761in}}%
\pgfpathlineto{\pgfqpoint{1.844382in}{1.210761in}}%
\pgfpathlineto{\pgfqpoint{1.844382in}{1.207812in}}%
\pgfpathmoveto{\pgfqpoint{1.839841in}{1.210761in}}%
\pgfpathlineto{\pgfqpoint{1.839841in}{1.210761in}}%
\pgfpathlineto{\pgfqpoint{1.839841in}{1.213710in}}%
\pgfpathlineto{\pgfqpoint{1.844382in}{1.213710in}}%
\pgfpathlineto{\pgfqpoint{1.844382in}{1.210761in}}%
\pgfpathmoveto{\pgfqpoint{1.844382in}{1.207812in}}%
\pgfpathlineto{\pgfqpoint{1.844382in}{1.207812in}}%
\pgfpathlineto{\pgfqpoint{1.844382in}{1.210761in}}%
\pgfpathlineto{\pgfqpoint{1.848923in}{1.210761in}}%
\pgfpathlineto{\pgfqpoint{1.848923in}{1.207812in}}%
\pgfpathmoveto{\pgfqpoint{1.844382in}{1.210761in}}%
\pgfpathlineto{\pgfqpoint{1.844382in}{1.210761in}}%
\pgfpathlineto{\pgfqpoint{1.844382in}{1.213710in}}%
\pgfpathlineto{\pgfqpoint{1.848923in}{1.213710in}}%
\pgfpathlineto{\pgfqpoint{1.848923in}{1.210761in}}%
\pgfpathmoveto{\pgfqpoint{1.848923in}{1.210761in}}%
\pgfpathlineto{\pgfqpoint{1.848923in}{1.210761in}}%
\pgfpathlineto{\pgfqpoint{1.848923in}{1.213710in}}%
\pgfpathlineto{\pgfqpoint{1.853464in}{1.213710in}}%
\pgfpathlineto{\pgfqpoint{1.853464in}{1.210761in}}%
\pgfpathmoveto{\pgfqpoint{1.848923in}{1.213710in}}%
\pgfpathlineto{\pgfqpoint{1.848923in}{1.213710in}}%
\pgfpathlineto{\pgfqpoint{1.848923in}{1.216659in}}%
\pgfpathlineto{\pgfqpoint{1.853464in}{1.216659in}}%
\pgfpathlineto{\pgfqpoint{1.853464in}{1.213710in}}%
\pgfpathmoveto{\pgfqpoint{1.848923in}{1.216659in}}%
\pgfpathlineto{\pgfqpoint{1.848923in}{1.216659in}}%
\pgfpathlineto{\pgfqpoint{1.848923in}{1.219608in}}%
\pgfpathlineto{\pgfqpoint{1.853464in}{1.219608in}}%
\pgfpathlineto{\pgfqpoint{1.853464in}{1.216659in}}%
\pgfpathmoveto{\pgfqpoint{1.853464in}{1.213710in}}%
\pgfpathlineto{\pgfqpoint{1.853464in}{1.213710in}}%
\pgfpathlineto{\pgfqpoint{1.853464in}{1.216659in}}%
\pgfpathlineto{\pgfqpoint{1.858005in}{1.216659in}}%
\pgfpathlineto{\pgfqpoint{1.858005in}{1.213710in}}%
\pgfpathmoveto{\pgfqpoint{1.853464in}{1.216659in}}%
\pgfpathlineto{\pgfqpoint{1.853464in}{1.216659in}}%
\pgfpathlineto{\pgfqpoint{1.853464in}{1.219608in}}%
\pgfpathlineto{\pgfqpoint{1.858005in}{1.219608in}}%
\pgfpathlineto{\pgfqpoint{1.858005in}{1.216659in}}%
\pgfpathmoveto{\pgfqpoint{1.858005in}{1.216659in}}%
\pgfpathlineto{\pgfqpoint{1.858005in}{1.216659in}}%
\pgfpathlineto{\pgfqpoint{1.858005in}{1.219608in}}%
\pgfpathlineto{\pgfqpoint{1.862546in}{1.219608in}}%
\pgfpathlineto{\pgfqpoint{1.862546in}{1.216659in}}%
\pgfpathmoveto{\pgfqpoint{1.858005in}{1.219608in}}%
\pgfpathlineto{\pgfqpoint{1.858005in}{1.219608in}}%
\pgfpathlineto{\pgfqpoint{1.858005in}{1.222557in}}%
\pgfpathlineto{\pgfqpoint{1.862546in}{1.222557in}}%
\pgfpathlineto{\pgfqpoint{1.862546in}{1.219608in}}%
\pgfpathmoveto{\pgfqpoint{1.858005in}{1.222557in}}%
\pgfpathlineto{\pgfqpoint{1.858005in}{1.222557in}}%
\pgfpathlineto{\pgfqpoint{1.858005in}{1.225506in}}%
\pgfpathlineto{\pgfqpoint{1.862546in}{1.225506in}}%
\pgfpathlineto{\pgfqpoint{1.862546in}{1.222557in}}%
\pgfpathmoveto{\pgfqpoint{1.862546in}{1.219608in}}%
\pgfpathlineto{\pgfqpoint{1.862546in}{1.219608in}}%
\pgfpathlineto{\pgfqpoint{1.862546in}{1.222557in}}%
\pgfpathlineto{\pgfqpoint{1.867087in}{1.222557in}}%
\pgfpathlineto{\pgfqpoint{1.867087in}{1.219608in}}%
\pgfpathmoveto{\pgfqpoint{1.862546in}{1.222557in}}%
\pgfpathlineto{\pgfqpoint{1.862546in}{1.222557in}}%
\pgfpathlineto{\pgfqpoint{1.862546in}{1.225506in}}%
\pgfpathlineto{\pgfqpoint{1.867087in}{1.225506in}}%
\pgfpathlineto{\pgfqpoint{1.867087in}{1.222557in}}%
\pgfpathmoveto{\pgfqpoint{1.867087in}{1.222557in}}%
\pgfpathlineto{\pgfqpoint{1.867087in}{1.222557in}}%
\pgfpathlineto{\pgfqpoint{1.867087in}{1.225506in}}%
\pgfpathlineto{\pgfqpoint{1.871628in}{1.225506in}}%
\pgfpathlineto{\pgfqpoint{1.871628in}{1.222557in}}%
\pgfpathmoveto{\pgfqpoint{1.867087in}{1.225506in}}%
\pgfpathlineto{\pgfqpoint{1.867087in}{1.225506in}}%
\pgfpathlineto{\pgfqpoint{1.867087in}{1.228456in}}%
\pgfpathlineto{\pgfqpoint{1.871628in}{1.228456in}}%
\pgfpathlineto{\pgfqpoint{1.871628in}{1.225506in}}%
\pgfpathmoveto{\pgfqpoint{1.867087in}{1.228456in}}%
\pgfpathlineto{\pgfqpoint{1.867087in}{1.228456in}}%
\pgfpathlineto{\pgfqpoint{1.867087in}{1.231405in}}%
\pgfpathlineto{\pgfqpoint{1.871628in}{1.231405in}}%
\pgfpathlineto{\pgfqpoint{1.871628in}{1.228456in}}%
\pgfpathmoveto{\pgfqpoint{1.871628in}{1.225506in}}%
\pgfpathlineto{\pgfqpoint{1.871628in}{1.225506in}}%
\pgfpathlineto{\pgfqpoint{1.871628in}{1.228456in}}%
\pgfpathlineto{\pgfqpoint{1.876168in}{1.228456in}}%
\pgfpathlineto{\pgfqpoint{1.876168in}{1.225506in}}%
\pgfpathmoveto{\pgfqpoint{1.871628in}{1.228456in}}%
\pgfpathlineto{\pgfqpoint{1.871628in}{1.228456in}}%
\pgfpathlineto{\pgfqpoint{1.871628in}{1.231405in}}%
\pgfpathlineto{\pgfqpoint{1.876168in}{1.231405in}}%
\pgfpathlineto{\pgfqpoint{1.876168in}{1.228456in}}%
\pgfpathmoveto{\pgfqpoint{1.876168in}{1.228456in}}%
\pgfpathlineto{\pgfqpoint{1.876168in}{1.228456in}}%
\pgfpathlineto{\pgfqpoint{1.876168in}{1.231405in}}%
\pgfpathlineto{\pgfqpoint{1.880709in}{1.231405in}}%
\pgfpathlineto{\pgfqpoint{1.880709in}{1.228456in}}%
\pgfpathmoveto{\pgfqpoint{1.876168in}{1.231405in}}%
\pgfpathlineto{\pgfqpoint{1.876168in}{1.231405in}}%
\pgfpathlineto{\pgfqpoint{1.876168in}{1.234354in}}%
\pgfpathlineto{\pgfqpoint{1.880709in}{1.234354in}}%
\pgfpathlineto{\pgfqpoint{1.880709in}{1.231405in}}%
\pgfpathmoveto{\pgfqpoint{1.876168in}{1.234354in}}%
\pgfpathlineto{\pgfqpoint{1.876168in}{1.234354in}}%
\pgfpathlineto{\pgfqpoint{1.876168in}{1.237303in}}%
\pgfpathlineto{\pgfqpoint{1.880709in}{1.237303in}}%
\pgfpathlineto{\pgfqpoint{1.880709in}{1.234354in}}%
\pgfpathmoveto{\pgfqpoint{1.880709in}{1.231405in}}%
\pgfpathlineto{\pgfqpoint{1.880709in}{1.231405in}}%
\pgfpathlineto{\pgfqpoint{1.880709in}{1.234354in}}%
\pgfpathlineto{\pgfqpoint{1.885250in}{1.234354in}}%
\pgfpathlineto{\pgfqpoint{1.885250in}{1.231405in}}%
\pgfpathmoveto{\pgfqpoint{1.880709in}{1.234354in}}%
\pgfpathlineto{\pgfqpoint{1.880709in}{1.234354in}}%
\pgfpathlineto{\pgfqpoint{1.880709in}{1.237303in}}%
\pgfpathlineto{\pgfqpoint{1.885250in}{1.237303in}}%
\pgfpathlineto{\pgfqpoint{1.885250in}{1.234354in}}%
\pgfpathmoveto{\pgfqpoint{1.885250in}{1.234354in}}%
\pgfpathlineto{\pgfqpoint{1.885250in}{1.234354in}}%
\pgfpathlineto{\pgfqpoint{1.885250in}{1.237303in}}%
\pgfpathlineto{\pgfqpoint{1.889791in}{1.237303in}}%
\pgfpathlineto{\pgfqpoint{1.889791in}{1.234354in}}%
\pgfpathmoveto{\pgfqpoint{1.885250in}{1.237303in}}%
\pgfpathlineto{\pgfqpoint{1.885250in}{1.237303in}}%
\pgfpathlineto{\pgfqpoint{1.885250in}{1.240252in}}%
\pgfpathlineto{\pgfqpoint{1.889791in}{1.240252in}}%
\pgfpathlineto{\pgfqpoint{1.889791in}{1.237303in}}%
\pgfpathmoveto{\pgfqpoint{1.885250in}{1.240252in}}%
\pgfpathlineto{\pgfqpoint{1.885250in}{1.240252in}}%
\pgfpathlineto{\pgfqpoint{1.885250in}{1.243201in}}%
\pgfpathlineto{\pgfqpoint{1.889791in}{1.243201in}}%
\pgfpathlineto{\pgfqpoint{1.889791in}{1.240252in}}%
\pgfpathmoveto{\pgfqpoint{1.889791in}{1.237303in}}%
\pgfpathlineto{\pgfqpoint{1.889791in}{1.237303in}}%
\pgfpathlineto{\pgfqpoint{1.889791in}{1.240252in}}%
\pgfpathlineto{\pgfqpoint{1.894332in}{1.240252in}}%
\pgfpathlineto{\pgfqpoint{1.894332in}{1.237303in}}%
\pgfpathmoveto{\pgfqpoint{1.889791in}{1.240252in}}%
\pgfpathlineto{\pgfqpoint{1.889791in}{1.240252in}}%
\pgfpathlineto{\pgfqpoint{1.889791in}{1.243201in}}%
\pgfpathlineto{\pgfqpoint{1.894332in}{1.243201in}}%
\pgfpathlineto{\pgfqpoint{1.894332in}{1.240252in}}%
\pgfpathmoveto{\pgfqpoint{1.894332in}{1.240252in}}%
\pgfpathlineto{\pgfqpoint{1.894332in}{1.240252in}}%
\pgfpathlineto{\pgfqpoint{1.894332in}{1.243201in}}%
\pgfpathlineto{\pgfqpoint{1.898873in}{1.243201in}}%
\pgfpathlineto{\pgfqpoint{1.898873in}{1.240252in}}%
\pgfpathmoveto{\pgfqpoint{1.894332in}{1.243201in}}%
\pgfpathlineto{\pgfqpoint{1.894332in}{1.243201in}}%
\pgfpathlineto{\pgfqpoint{1.894332in}{1.246151in}}%
\pgfpathlineto{\pgfqpoint{1.898873in}{1.246151in}}%
\pgfpathlineto{\pgfqpoint{1.898873in}{1.243201in}}%
\pgfpathmoveto{\pgfqpoint{1.894332in}{1.246151in}}%
\pgfpathlineto{\pgfqpoint{1.894332in}{1.246151in}}%
\pgfpathlineto{\pgfqpoint{1.894332in}{1.249100in}}%
\pgfpathlineto{\pgfqpoint{1.898873in}{1.249100in}}%
\pgfpathlineto{\pgfqpoint{1.898873in}{1.246151in}}%
\pgfpathmoveto{\pgfqpoint{1.898873in}{1.243201in}}%
\pgfpathlineto{\pgfqpoint{1.898873in}{1.243201in}}%
\pgfpathlineto{\pgfqpoint{1.898873in}{1.246151in}}%
\pgfpathlineto{\pgfqpoint{1.903414in}{1.246151in}}%
\pgfpathlineto{\pgfqpoint{1.903414in}{1.243201in}}%
\pgfpathmoveto{\pgfqpoint{1.898873in}{1.246151in}}%
\pgfpathlineto{\pgfqpoint{1.898873in}{1.246151in}}%
\pgfpathlineto{\pgfqpoint{1.898873in}{1.249100in}}%
\pgfpathlineto{\pgfqpoint{1.903414in}{1.249100in}}%
\pgfpathlineto{\pgfqpoint{1.903414in}{1.246151in}}%
\pgfpathmoveto{\pgfqpoint{1.903414in}{1.246151in}}%
\pgfpathlineto{\pgfqpoint{1.903414in}{1.246151in}}%
\pgfpathlineto{\pgfqpoint{1.903414in}{1.249100in}}%
\pgfpathlineto{\pgfqpoint{1.907955in}{1.249100in}}%
\pgfpathlineto{\pgfqpoint{1.907955in}{1.246151in}}%
\pgfpathmoveto{\pgfqpoint{1.903414in}{1.249100in}}%
\pgfpathlineto{\pgfqpoint{1.903414in}{1.249100in}}%
\pgfpathlineto{\pgfqpoint{1.903414in}{1.252049in}}%
\pgfpathlineto{\pgfqpoint{1.907955in}{1.252049in}}%
\pgfpathlineto{\pgfqpoint{1.907955in}{1.249100in}}%
\pgfpathmoveto{\pgfqpoint{1.903414in}{1.252049in}}%
\pgfpathlineto{\pgfqpoint{1.903414in}{1.252049in}}%
\pgfpathlineto{\pgfqpoint{1.903414in}{1.254998in}}%
\pgfpathlineto{\pgfqpoint{1.907955in}{1.254998in}}%
\pgfpathlineto{\pgfqpoint{1.907955in}{1.252049in}}%
\pgfpathmoveto{\pgfqpoint{1.907955in}{1.249100in}}%
\pgfpathlineto{\pgfqpoint{1.907955in}{1.249100in}}%
\pgfpathlineto{\pgfqpoint{1.907955in}{1.252049in}}%
\pgfpathlineto{\pgfqpoint{1.912496in}{1.252049in}}%
\pgfpathlineto{\pgfqpoint{1.912496in}{1.249100in}}%
\pgfpathmoveto{\pgfqpoint{1.907955in}{1.252049in}}%
\pgfpathlineto{\pgfqpoint{1.907955in}{1.252049in}}%
\pgfpathlineto{\pgfqpoint{1.907955in}{1.254998in}}%
\pgfpathlineto{\pgfqpoint{1.912496in}{1.254998in}}%
\pgfpathlineto{\pgfqpoint{1.912496in}{1.252049in}}%
\pgfpathmoveto{\pgfqpoint{1.912496in}{1.252049in}}%
\pgfpathlineto{\pgfqpoint{1.912496in}{1.252049in}}%
\pgfpathlineto{\pgfqpoint{1.912496in}{1.254998in}}%
\pgfpathlineto{\pgfqpoint{1.917037in}{1.254998in}}%
\pgfpathlineto{\pgfqpoint{1.917037in}{1.252049in}}%
\pgfpathmoveto{\pgfqpoint{1.912496in}{1.254998in}}%
\pgfpathlineto{\pgfqpoint{1.912496in}{1.254998in}}%
\pgfpathlineto{\pgfqpoint{1.912496in}{1.257947in}}%
\pgfpathlineto{\pgfqpoint{1.917037in}{1.257947in}}%
\pgfpathlineto{\pgfqpoint{1.917037in}{1.254998in}}%
\pgfpathmoveto{\pgfqpoint{1.912496in}{1.257947in}}%
\pgfpathlineto{\pgfqpoint{1.912496in}{1.257947in}}%
\pgfpathlineto{\pgfqpoint{1.912496in}{1.260897in}}%
\pgfpathlineto{\pgfqpoint{1.917037in}{1.260897in}}%
\pgfpathlineto{\pgfqpoint{1.917037in}{1.257947in}}%
\pgfpathmoveto{\pgfqpoint{1.917037in}{1.254998in}}%
\pgfpathlineto{\pgfqpoint{1.917037in}{1.254998in}}%
\pgfpathlineto{\pgfqpoint{1.917037in}{1.257947in}}%
\pgfpathlineto{\pgfqpoint{1.921578in}{1.257947in}}%
\pgfpathlineto{\pgfqpoint{1.921578in}{1.254998in}}%
\pgfpathmoveto{\pgfqpoint{1.917037in}{1.257947in}}%
\pgfpathlineto{\pgfqpoint{1.917037in}{1.257947in}}%
\pgfpathlineto{\pgfqpoint{1.917037in}{1.260897in}}%
\pgfpathlineto{\pgfqpoint{1.921578in}{1.260897in}}%
\pgfpathlineto{\pgfqpoint{1.921578in}{1.257947in}}%
\pgfpathmoveto{\pgfqpoint{1.921578in}{1.257947in}}%
\pgfpathlineto{\pgfqpoint{1.921578in}{1.257947in}}%
\pgfpathlineto{\pgfqpoint{1.921578in}{1.260897in}}%
\pgfpathlineto{\pgfqpoint{1.926119in}{1.260897in}}%
\pgfpathlineto{\pgfqpoint{1.926119in}{1.257947in}}%
\pgfpathmoveto{\pgfqpoint{1.921578in}{1.260897in}}%
\pgfpathlineto{\pgfqpoint{1.921578in}{1.260897in}}%
\pgfpathlineto{\pgfqpoint{1.921578in}{1.263846in}}%
\pgfpathlineto{\pgfqpoint{1.926119in}{1.263846in}}%
\pgfpathlineto{\pgfqpoint{1.926119in}{1.260897in}}%
\pgfpathmoveto{\pgfqpoint{1.921578in}{1.263846in}}%
\pgfpathlineto{\pgfqpoint{1.921578in}{1.263846in}}%
\pgfpathlineto{\pgfqpoint{1.921578in}{1.266795in}}%
\pgfpathlineto{\pgfqpoint{1.926119in}{1.266795in}}%
\pgfpathlineto{\pgfqpoint{1.926119in}{1.263846in}}%
\pgfpathmoveto{\pgfqpoint{1.926119in}{1.260897in}}%
\pgfpathlineto{\pgfqpoint{1.926119in}{1.260897in}}%
\pgfpathlineto{\pgfqpoint{1.926119in}{1.263846in}}%
\pgfpathlineto{\pgfqpoint{1.930661in}{1.263846in}}%
\pgfpathlineto{\pgfqpoint{1.930661in}{1.260897in}}%
\pgfpathmoveto{\pgfqpoint{1.926119in}{1.263846in}}%
\pgfpathlineto{\pgfqpoint{1.926119in}{1.263846in}}%
\pgfpathlineto{\pgfqpoint{1.926119in}{1.266795in}}%
\pgfpathlineto{\pgfqpoint{1.930661in}{1.266795in}}%
\pgfpathlineto{\pgfqpoint{1.930661in}{1.263846in}}%
\pgfpathmoveto{\pgfqpoint{1.930661in}{1.263846in}}%
\pgfpathlineto{\pgfqpoint{1.930661in}{1.263846in}}%
\pgfpathlineto{\pgfqpoint{1.930661in}{1.266795in}}%
\pgfpathlineto{\pgfqpoint{1.935202in}{1.266795in}}%
\pgfpathlineto{\pgfqpoint{1.935202in}{1.263846in}}%
\pgfpathmoveto{\pgfqpoint{1.930661in}{1.266795in}}%
\pgfpathlineto{\pgfqpoint{1.930661in}{1.266795in}}%
\pgfpathlineto{\pgfqpoint{1.930661in}{1.269745in}}%
\pgfpathlineto{\pgfqpoint{1.935202in}{1.269745in}}%
\pgfpathlineto{\pgfqpoint{1.935202in}{1.266795in}}%
\pgfpathmoveto{\pgfqpoint{1.930661in}{1.269745in}}%
\pgfpathlineto{\pgfqpoint{1.930661in}{1.269745in}}%
\pgfpathlineto{\pgfqpoint{1.930661in}{1.272694in}}%
\pgfpathlineto{\pgfqpoint{1.935202in}{1.272694in}}%
\pgfpathlineto{\pgfqpoint{1.935202in}{1.269745in}}%
\pgfpathmoveto{\pgfqpoint{1.935202in}{1.266795in}}%
\pgfpathlineto{\pgfqpoint{1.935202in}{1.266795in}}%
\pgfpathlineto{\pgfqpoint{1.935202in}{1.269745in}}%
\pgfpathlineto{\pgfqpoint{1.939743in}{1.269745in}}%
\pgfpathlineto{\pgfqpoint{1.939743in}{1.266795in}}%
\pgfpathmoveto{\pgfqpoint{1.935202in}{1.269745in}}%
\pgfpathlineto{\pgfqpoint{1.935202in}{1.269745in}}%
\pgfpathlineto{\pgfqpoint{1.935202in}{1.272694in}}%
\pgfpathlineto{\pgfqpoint{1.939743in}{1.272694in}}%
\pgfpathlineto{\pgfqpoint{1.939743in}{1.269745in}}%
\pgfpathmoveto{\pgfqpoint{1.939743in}{1.269745in}}%
\pgfpathlineto{\pgfqpoint{1.939743in}{1.269745in}}%
\pgfpathlineto{\pgfqpoint{1.939743in}{1.272694in}}%
\pgfpathlineto{\pgfqpoint{1.944284in}{1.272694in}}%
\pgfpathlineto{\pgfqpoint{1.944284in}{1.269745in}}%
\pgfpathmoveto{\pgfqpoint{1.939743in}{1.272694in}}%
\pgfpathlineto{\pgfqpoint{1.939743in}{1.272694in}}%
\pgfpathlineto{\pgfqpoint{1.939743in}{1.275643in}}%
\pgfpathlineto{\pgfqpoint{1.944284in}{1.275643in}}%
\pgfpathlineto{\pgfqpoint{1.944284in}{1.272694in}}%
\pgfpathmoveto{\pgfqpoint{1.939743in}{1.275643in}}%
\pgfpathlineto{\pgfqpoint{1.939743in}{1.275643in}}%
\pgfpathlineto{\pgfqpoint{1.939743in}{1.278592in}}%
\pgfpathlineto{\pgfqpoint{1.944284in}{1.278592in}}%
\pgfpathlineto{\pgfqpoint{1.944284in}{1.275643in}}%
\pgfpathmoveto{\pgfqpoint{1.944284in}{1.272694in}}%
\pgfpathlineto{\pgfqpoint{1.944284in}{1.272694in}}%
\pgfpathlineto{\pgfqpoint{1.944284in}{1.275643in}}%
\pgfpathlineto{\pgfqpoint{1.948825in}{1.275643in}}%
\pgfpathlineto{\pgfqpoint{1.948825in}{1.272694in}}%
\pgfpathmoveto{\pgfqpoint{1.944284in}{1.275643in}}%
\pgfpathlineto{\pgfqpoint{1.944284in}{1.275643in}}%
\pgfpathlineto{\pgfqpoint{1.944284in}{1.278592in}}%
\pgfpathlineto{\pgfqpoint{1.948825in}{1.278592in}}%
\pgfpathlineto{\pgfqpoint{1.948825in}{1.275643in}}%
\pgfpathmoveto{\pgfqpoint{1.948825in}{1.275643in}}%
\pgfpathlineto{\pgfqpoint{1.948825in}{1.275643in}}%
\pgfpathlineto{\pgfqpoint{1.948825in}{1.278592in}}%
\pgfpathlineto{\pgfqpoint{1.953366in}{1.278592in}}%
\pgfpathlineto{\pgfqpoint{1.953366in}{1.275643in}}%
\pgfpathmoveto{\pgfqpoint{1.948825in}{1.278592in}}%
\pgfpathlineto{\pgfqpoint{1.948825in}{1.278592in}}%
\pgfpathlineto{\pgfqpoint{1.948825in}{1.281542in}}%
\pgfpathlineto{\pgfqpoint{1.953366in}{1.281542in}}%
\pgfpathlineto{\pgfqpoint{1.953366in}{1.278592in}}%
\pgfpathmoveto{\pgfqpoint{1.948825in}{1.281542in}}%
\pgfpathlineto{\pgfqpoint{1.948825in}{1.281542in}}%
\pgfpathlineto{\pgfqpoint{1.948825in}{1.284491in}}%
\pgfpathlineto{\pgfqpoint{1.953366in}{1.284491in}}%
\pgfpathlineto{\pgfqpoint{1.953366in}{1.281542in}}%
\pgfpathmoveto{\pgfqpoint{1.953366in}{1.278592in}}%
\pgfpathlineto{\pgfqpoint{1.953366in}{1.278592in}}%
\pgfpathlineto{\pgfqpoint{1.953366in}{1.281542in}}%
\pgfpathlineto{\pgfqpoint{1.957908in}{1.281542in}}%
\pgfpathlineto{\pgfqpoint{1.957908in}{1.278592in}}%
\pgfpathmoveto{\pgfqpoint{1.953366in}{1.281542in}}%
\pgfpathlineto{\pgfqpoint{1.953366in}{1.281542in}}%
\pgfpathlineto{\pgfqpoint{1.953366in}{1.284491in}}%
\pgfpathlineto{\pgfqpoint{1.957908in}{1.284491in}}%
\pgfpathlineto{\pgfqpoint{1.957908in}{1.281542in}}%
\pgfpathmoveto{\pgfqpoint{1.957908in}{1.281542in}}%
\pgfpathlineto{\pgfqpoint{1.957908in}{1.281542in}}%
\pgfpathlineto{\pgfqpoint{1.957908in}{1.284491in}}%
\pgfpathlineto{\pgfqpoint{1.962449in}{1.284491in}}%
\pgfpathlineto{\pgfqpoint{1.962449in}{1.281542in}}%
\pgfpathmoveto{\pgfqpoint{1.957908in}{1.284491in}}%
\pgfpathlineto{\pgfqpoint{1.957908in}{1.284491in}}%
\pgfpathlineto{\pgfqpoint{1.957908in}{1.287440in}}%
\pgfpathlineto{\pgfqpoint{1.962449in}{1.287440in}}%
\pgfpathlineto{\pgfqpoint{1.962449in}{1.284491in}}%
\pgfpathmoveto{\pgfqpoint{1.957908in}{1.287440in}}%
\pgfpathlineto{\pgfqpoint{1.957908in}{1.287440in}}%
\pgfpathlineto{\pgfqpoint{1.957908in}{1.290390in}}%
\pgfpathlineto{\pgfqpoint{1.962449in}{1.290390in}}%
\pgfpathlineto{\pgfqpoint{1.962449in}{1.287440in}}%
\pgfpathmoveto{\pgfqpoint{1.962449in}{1.284491in}}%
\pgfpathlineto{\pgfqpoint{1.962449in}{1.284491in}}%
\pgfpathlineto{\pgfqpoint{1.962449in}{1.287440in}}%
\pgfpathlineto{\pgfqpoint{1.966990in}{1.287440in}}%
\pgfpathlineto{\pgfqpoint{1.966990in}{1.284491in}}%
\pgfpathmoveto{\pgfqpoint{1.962449in}{1.287440in}}%
\pgfpathlineto{\pgfqpoint{1.962449in}{1.287440in}}%
\pgfpathlineto{\pgfqpoint{1.962449in}{1.290390in}}%
\pgfpathlineto{\pgfqpoint{1.966990in}{1.290390in}}%
\pgfpathlineto{\pgfqpoint{1.966990in}{1.287440in}}%
\pgfpathmoveto{\pgfqpoint{1.966990in}{1.287440in}}%
\pgfpathlineto{\pgfqpoint{1.966990in}{1.287440in}}%
\pgfpathlineto{\pgfqpoint{1.966990in}{1.290390in}}%
\pgfpathlineto{\pgfqpoint{1.971531in}{1.290390in}}%
\pgfpathlineto{\pgfqpoint{1.971531in}{1.287440in}}%
\pgfpathmoveto{\pgfqpoint{1.966990in}{1.290390in}}%
\pgfpathlineto{\pgfqpoint{1.966990in}{1.290390in}}%
\pgfpathlineto{\pgfqpoint{1.966990in}{1.293339in}}%
\pgfpathlineto{\pgfqpoint{1.971531in}{1.293339in}}%
\pgfpathlineto{\pgfqpoint{1.971531in}{1.290390in}}%
\pgfpathmoveto{\pgfqpoint{1.966990in}{1.293339in}}%
\pgfpathlineto{\pgfqpoint{1.966990in}{1.293339in}}%
\pgfpathlineto{\pgfqpoint{1.966990in}{1.296288in}}%
\pgfpathlineto{\pgfqpoint{1.971531in}{1.296288in}}%
\pgfpathlineto{\pgfqpoint{1.971531in}{1.293339in}}%
\pgfpathmoveto{\pgfqpoint{1.971531in}{1.290390in}}%
\pgfpathlineto{\pgfqpoint{1.971531in}{1.290390in}}%
\pgfpathlineto{\pgfqpoint{1.971531in}{1.293339in}}%
\pgfpathlineto{\pgfqpoint{1.976072in}{1.293339in}}%
\pgfpathlineto{\pgfqpoint{1.976072in}{1.290390in}}%
\pgfpathmoveto{\pgfqpoint{1.971531in}{1.293339in}}%
\pgfpathlineto{\pgfqpoint{1.971531in}{1.293339in}}%
\pgfpathlineto{\pgfqpoint{1.971531in}{1.296288in}}%
\pgfpathlineto{\pgfqpoint{1.976072in}{1.296288in}}%
\pgfpathlineto{\pgfqpoint{1.976072in}{1.293339in}}%
\pgfpathmoveto{\pgfqpoint{1.976072in}{1.293339in}}%
\pgfpathlineto{\pgfqpoint{1.976072in}{1.293339in}}%
\pgfpathlineto{\pgfqpoint{1.976072in}{1.296288in}}%
\pgfpathlineto{\pgfqpoint{1.980613in}{1.296288in}}%
\pgfpathlineto{\pgfqpoint{1.980613in}{1.293339in}}%
\pgfpathmoveto{\pgfqpoint{1.976072in}{1.296288in}}%
\pgfpathlineto{\pgfqpoint{1.976072in}{1.296288in}}%
\pgfpathlineto{\pgfqpoint{1.976072in}{1.299238in}}%
\pgfpathlineto{\pgfqpoint{1.980613in}{1.299238in}}%
\pgfpathlineto{\pgfqpoint{1.980613in}{1.296288in}}%
\pgfpathmoveto{\pgfqpoint{1.976072in}{1.299238in}}%
\pgfpathlineto{\pgfqpoint{1.976072in}{1.299238in}}%
\pgfpathlineto{\pgfqpoint{1.976072in}{1.302187in}}%
\pgfpathlineto{\pgfqpoint{1.980613in}{1.302187in}}%
\pgfpathlineto{\pgfqpoint{1.980613in}{1.299238in}}%
\pgfpathmoveto{\pgfqpoint{1.980613in}{1.296288in}}%
\pgfpathlineto{\pgfqpoint{1.980613in}{1.296288in}}%
\pgfpathlineto{\pgfqpoint{1.980613in}{1.299238in}}%
\pgfpathlineto{\pgfqpoint{1.985155in}{1.299238in}}%
\pgfpathlineto{\pgfqpoint{1.985155in}{1.296288in}}%
\pgfpathmoveto{\pgfqpoint{1.980613in}{1.299238in}}%
\pgfpathlineto{\pgfqpoint{1.980613in}{1.299238in}}%
\pgfpathlineto{\pgfqpoint{1.980613in}{1.302187in}}%
\pgfpathlineto{\pgfqpoint{1.985155in}{1.302187in}}%
\pgfpathlineto{\pgfqpoint{1.985155in}{1.299238in}}%
\pgfpathmoveto{\pgfqpoint{1.985155in}{1.299238in}}%
\pgfpathlineto{\pgfqpoint{1.985155in}{1.299238in}}%
\pgfpathlineto{\pgfqpoint{1.985155in}{1.302187in}}%
\pgfpathlineto{\pgfqpoint{1.989696in}{1.302187in}}%
\pgfpathlineto{\pgfqpoint{1.989696in}{1.299238in}}%
\pgfpathmoveto{\pgfqpoint{1.985155in}{1.302187in}}%
\pgfpathlineto{\pgfqpoint{1.985155in}{1.302187in}}%
\pgfpathlineto{\pgfqpoint{1.985155in}{1.305136in}}%
\pgfpathlineto{\pgfqpoint{1.989696in}{1.305136in}}%
\pgfpathlineto{\pgfqpoint{1.989696in}{1.302187in}}%
\pgfpathmoveto{\pgfqpoint{1.985155in}{1.305136in}}%
\pgfpathlineto{\pgfqpoint{1.985155in}{1.305136in}}%
\pgfpathlineto{\pgfqpoint{1.985155in}{1.308086in}}%
\pgfpathlineto{\pgfqpoint{1.989696in}{1.308086in}}%
\pgfpathlineto{\pgfqpoint{1.989696in}{1.305136in}}%
\pgfpathmoveto{\pgfqpoint{1.989696in}{1.302187in}}%
\pgfpathlineto{\pgfqpoint{1.989696in}{1.302187in}}%
\pgfpathlineto{\pgfqpoint{1.989696in}{1.305136in}}%
\pgfpathlineto{\pgfqpoint{1.994237in}{1.305136in}}%
\pgfpathlineto{\pgfqpoint{1.994237in}{1.302187in}}%
\pgfpathmoveto{\pgfqpoint{1.989696in}{1.305136in}}%
\pgfpathlineto{\pgfqpoint{1.989696in}{1.305136in}}%
\pgfpathlineto{\pgfqpoint{1.989696in}{1.308086in}}%
\pgfpathlineto{\pgfqpoint{1.994237in}{1.308086in}}%
\pgfpathlineto{\pgfqpoint{1.994237in}{1.305136in}}%
\pgfpathmoveto{\pgfqpoint{1.994237in}{1.305136in}}%
\pgfpathlineto{\pgfqpoint{1.994237in}{1.305136in}}%
\pgfpathlineto{\pgfqpoint{1.994237in}{1.308086in}}%
\pgfpathlineto{\pgfqpoint{1.998778in}{1.308086in}}%
\pgfpathlineto{\pgfqpoint{1.998778in}{1.305136in}}%
\pgfpathmoveto{\pgfqpoint{1.994237in}{1.308086in}}%
\pgfpathlineto{\pgfqpoint{1.994237in}{1.308086in}}%
\pgfpathlineto{\pgfqpoint{1.994237in}{1.311035in}}%
\pgfpathlineto{\pgfqpoint{1.998778in}{1.311035in}}%
\pgfpathlineto{\pgfqpoint{1.998778in}{1.308086in}}%
\pgfpathmoveto{\pgfqpoint{1.994237in}{1.311035in}}%
\pgfpathlineto{\pgfqpoint{1.994237in}{1.311035in}}%
\pgfpathlineto{\pgfqpoint{1.994237in}{1.313984in}}%
\pgfpathlineto{\pgfqpoint{1.998778in}{1.313984in}}%
\pgfpathlineto{\pgfqpoint{1.998778in}{1.311035in}}%
\pgfpathmoveto{\pgfqpoint{1.998778in}{1.308086in}}%
\pgfpathlineto{\pgfqpoint{1.998778in}{1.308086in}}%
\pgfpathlineto{\pgfqpoint{1.998778in}{1.311035in}}%
\pgfpathlineto{\pgfqpoint{2.003319in}{1.311035in}}%
\pgfpathlineto{\pgfqpoint{2.003319in}{1.308086in}}%
\pgfpathmoveto{\pgfqpoint{1.998778in}{1.311035in}}%
\pgfpathlineto{\pgfqpoint{1.998778in}{1.311035in}}%
\pgfpathlineto{\pgfqpoint{1.998778in}{1.313984in}}%
\pgfpathlineto{\pgfqpoint{2.003319in}{1.313984in}}%
\pgfpathlineto{\pgfqpoint{2.003319in}{1.311035in}}%
\pgfpathmoveto{\pgfqpoint{2.003319in}{1.311035in}}%
\pgfpathlineto{\pgfqpoint{2.003319in}{1.311035in}}%
\pgfpathlineto{\pgfqpoint{2.003319in}{1.313984in}}%
\pgfpathlineto{\pgfqpoint{2.007860in}{1.313984in}}%
\pgfpathlineto{\pgfqpoint{2.007860in}{1.311035in}}%
\pgfpathmoveto{\pgfqpoint{2.003319in}{1.313984in}}%
\pgfpathlineto{\pgfqpoint{2.003319in}{1.313984in}}%
\pgfpathlineto{\pgfqpoint{2.003319in}{1.316933in}}%
\pgfpathlineto{\pgfqpoint{2.007860in}{1.316933in}}%
\pgfpathlineto{\pgfqpoint{2.007860in}{1.313984in}}%
\pgfpathmoveto{\pgfqpoint{2.003319in}{1.316933in}}%
\pgfpathlineto{\pgfqpoint{2.003319in}{1.316933in}}%
\pgfpathlineto{\pgfqpoint{2.003319in}{1.319883in}}%
\pgfpathlineto{\pgfqpoint{2.007860in}{1.319883in}}%
\pgfpathlineto{\pgfqpoint{2.007860in}{1.316933in}}%
\pgfpathmoveto{\pgfqpoint{2.007860in}{1.313984in}}%
\pgfpathlineto{\pgfqpoint{2.007860in}{1.313984in}}%
\pgfpathlineto{\pgfqpoint{2.007860in}{1.316933in}}%
\pgfpathlineto{\pgfqpoint{2.012402in}{1.316933in}}%
\pgfpathlineto{\pgfqpoint{2.012402in}{1.313984in}}%
\pgfpathmoveto{\pgfqpoint{2.007860in}{1.316933in}}%
\pgfpathlineto{\pgfqpoint{2.007860in}{1.316933in}}%
\pgfpathlineto{\pgfqpoint{2.007860in}{1.319883in}}%
\pgfpathlineto{\pgfqpoint{2.012402in}{1.319883in}}%
\pgfpathlineto{\pgfqpoint{2.012402in}{1.316933in}}%
\pgfpathmoveto{\pgfqpoint{2.012402in}{1.316933in}}%
\pgfpathlineto{\pgfqpoint{2.012402in}{1.316933in}}%
\pgfpathlineto{\pgfqpoint{2.012402in}{1.319883in}}%
\pgfpathlineto{\pgfqpoint{2.016943in}{1.319883in}}%
\pgfpathlineto{\pgfqpoint{2.016943in}{1.316933in}}%
\pgfpathmoveto{\pgfqpoint{2.012402in}{1.319883in}}%
\pgfpathlineto{\pgfqpoint{2.012402in}{1.319883in}}%
\pgfpathlineto{\pgfqpoint{2.012402in}{1.322832in}}%
\pgfpathlineto{\pgfqpoint{2.016943in}{1.322832in}}%
\pgfpathlineto{\pgfqpoint{2.016943in}{1.319883in}}%
\pgfpathmoveto{\pgfqpoint{2.012402in}{1.322832in}}%
\pgfpathlineto{\pgfqpoint{2.012402in}{1.322832in}}%
\pgfpathlineto{\pgfqpoint{2.012402in}{1.325781in}}%
\pgfpathlineto{\pgfqpoint{2.016943in}{1.325781in}}%
\pgfpathlineto{\pgfqpoint{2.016943in}{1.322832in}}%
\pgfpathmoveto{\pgfqpoint{2.016943in}{1.319883in}}%
\pgfpathlineto{\pgfqpoint{2.016943in}{1.319883in}}%
\pgfpathlineto{\pgfqpoint{2.016943in}{1.322832in}}%
\pgfpathlineto{\pgfqpoint{2.021484in}{1.322832in}}%
\pgfpathlineto{\pgfqpoint{2.021484in}{1.319883in}}%
\pgfpathmoveto{\pgfqpoint{2.016943in}{1.322832in}}%
\pgfpathlineto{\pgfqpoint{2.016943in}{1.322832in}}%
\pgfpathlineto{\pgfqpoint{2.016943in}{1.325781in}}%
\pgfpathlineto{\pgfqpoint{2.021484in}{1.325781in}}%
\pgfpathlineto{\pgfqpoint{2.021484in}{1.322832in}}%
\pgfpathmoveto{\pgfqpoint{2.021484in}{1.322832in}}%
\pgfpathlineto{\pgfqpoint{2.021484in}{1.322832in}}%
\pgfpathlineto{\pgfqpoint{2.021484in}{1.325781in}}%
\pgfpathlineto{\pgfqpoint{2.026025in}{1.325781in}}%
\pgfpathlineto{\pgfqpoint{2.026025in}{1.322832in}}%
\pgfpathmoveto{\pgfqpoint{2.021484in}{1.325781in}}%
\pgfpathlineto{\pgfqpoint{2.021484in}{1.325781in}}%
\pgfpathlineto{\pgfqpoint{2.021484in}{1.328731in}}%
\pgfpathlineto{\pgfqpoint{2.026025in}{1.328731in}}%
\pgfpathlineto{\pgfqpoint{2.026025in}{1.325781in}}%
\pgfpathmoveto{\pgfqpoint{2.021484in}{1.328731in}}%
\pgfpathlineto{\pgfqpoint{2.021484in}{1.328731in}}%
\pgfpathlineto{\pgfqpoint{2.021484in}{1.331680in}}%
\pgfpathlineto{\pgfqpoint{2.026025in}{1.331680in}}%
\pgfpathlineto{\pgfqpoint{2.026025in}{1.328731in}}%
\pgfpathmoveto{\pgfqpoint{2.026025in}{1.325781in}}%
\pgfpathlineto{\pgfqpoint{2.026025in}{1.325781in}}%
\pgfpathlineto{\pgfqpoint{2.026025in}{1.328731in}}%
\pgfpathlineto{\pgfqpoint{2.030566in}{1.328731in}}%
\pgfpathlineto{\pgfqpoint{2.030566in}{1.325781in}}%
\pgfpathmoveto{\pgfqpoint{2.026025in}{1.328731in}}%
\pgfpathlineto{\pgfqpoint{2.026025in}{1.328731in}}%
\pgfpathlineto{\pgfqpoint{2.026025in}{1.331680in}}%
\pgfpathlineto{\pgfqpoint{2.030566in}{1.331680in}}%
\pgfpathlineto{\pgfqpoint{2.030566in}{1.328731in}}%
\pgfpathmoveto{\pgfqpoint{2.030566in}{1.328731in}}%
\pgfpathlineto{\pgfqpoint{2.030566in}{1.328731in}}%
\pgfpathlineto{\pgfqpoint{2.030566in}{1.331680in}}%
\pgfpathlineto{\pgfqpoint{2.035107in}{1.331680in}}%
\pgfpathlineto{\pgfqpoint{2.035107in}{1.328731in}}%
\pgfpathmoveto{\pgfqpoint{2.030566in}{1.331680in}}%
\pgfpathlineto{\pgfqpoint{2.030566in}{1.331680in}}%
\pgfpathlineto{\pgfqpoint{2.030566in}{1.334629in}}%
\pgfpathlineto{\pgfqpoint{2.035107in}{1.334629in}}%
\pgfpathlineto{\pgfqpoint{2.035107in}{1.331680in}}%
\pgfpathmoveto{\pgfqpoint{2.030566in}{1.334629in}}%
\pgfpathlineto{\pgfqpoint{2.030566in}{1.334629in}}%
\pgfpathlineto{\pgfqpoint{2.030566in}{1.337579in}}%
\pgfpathlineto{\pgfqpoint{2.035107in}{1.337579in}}%
\pgfpathlineto{\pgfqpoint{2.035107in}{1.334629in}}%
\pgfpathmoveto{\pgfqpoint{2.035107in}{1.331680in}}%
\pgfpathlineto{\pgfqpoint{2.035107in}{1.331680in}}%
\pgfpathlineto{\pgfqpoint{2.035107in}{1.334629in}}%
\pgfpathlineto{\pgfqpoint{2.039649in}{1.334629in}}%
\pgfpathlineto{\pgfqpoint{2.039649in}{1.331680in}}%
\pgfpathmoveto{\pgfqpoint{2.035107in}{1.334629in}}%
\pgfpathlineto{\pgfqpoint{2.035107in}{1.334629in}}%
\pgfpathlineto{\pgfqpoint{2.035107in}{1.337579in}}%
\pgfpathlineto{\pgfqpoint{2.039649in}{1.337579in}}%
\pgfpathlineto{\pgfqpoint{2.039649in}{1.334629in}}%
\pgfpathmoveto{\pgfqpoint{2.039649in}{1.334629in}}%
\pgfpathlineto{\pgfqpoint{2.039649in}{1.334629in}}%
\pgfpathlineto{\pgfqpoint{2.039649in}{1.337579in}}%
\pgfpathlineto{\pgfqpoint{2.044190in}{1.337579in}}%
\pgfpathlineto{\pgfqpoint{2.044190in}{1.334629in}}%
\pgfpathmoveto{\pgfqpoint{2.039649in}{1.337579in}}%
\pgfpathlineto{\pgfqpoint{2.039649in}{1.337579in}}%
\pgfpathlineto{\pgfqpoint{2.039649in}{1.340528in}}%
\pgfpathlineto{\pgfqpoint{2.044190in}{1.340528in}}%
\pgfpathlineto{\pgfqpoint{2.044190in}{1.337579in}}%
\pgfpathmoveto{\pgfqpoint{2.039649in}{1.340528in}}%
\pgfpathlineto{\pgfqpoint{2.039649in}{1.340528in}}%
\pgfpathlineto{\pgfqpoint{2.039649in}{1.343477in}}%
\pgfpathlineto{\pgfqpoint{2.044190in}{1.343477in}}%
\pgfpathlineto{\pgfqpoint{2.044190in}{1.340528in}}%
\pgfpathmoveto{\pgfqpoint{2.044190in}{1.337579in}}%
\pgfpathlineto{\pgfqpoint{2.044190in}{1.337579in}}%
\pgfpathlineto{\pgfqpoint{2.044190in}{1.340528in}}%
\pgfpathlineto{\pgfqpoint{2.048731in}{1.340528in}}%
\pgfpathlineto{\pgfqpoint{2.048731in}{1.337579in}}%
\pgfpathmoveto{\pgfqpoint{2.044190in}{1.340528in}}%
\pgfpathlineto{\pgfqpoint{2.044190in}{1.340528in}}%
\pgfpathlineto{\pgfqpoint{2.044190in}{1.343477in}}%
\pgfpathlineto{\pgfqpoint{2.048731in}{1.343477in}}%
\pgfpathlineto{\pgfqpoint{2.048731in}{1.340528in}}%
\pgfpathmoveto{\pgfqpoint{2.048731in}{1.340528in}}%
\pgfpathlineto{\pgfqpoint{2.048731in}{1.340528in}}%
\pgfpathlineto{\pgfqpoint{2.048731in}{1.343477in}}%
\pgfpathlineto{\pgfqpoint{2.053272in}{1.343477in}}%
\pgfpathlineto{\pgfqpoint{2.053272in}{1.340528in}}%
\pgfpathmoveto{\pgfqpoint{2.048731in}{1.343477in}}%
\pgfpathlineto{\pgfqpoint{2.048731in}{1.343477in}}%
\pgfpathlineto{\pgfqpoint{2.048731in}{1.346427in}}%
\pgfpathlineto{\pgfqpoint{2.053272in}{1.346427in}}%
\pgfpathlineto{\pgfqpoint{2.053272in}{1.343477in}}%
\pgfpathmoveto{\pgfqpoint{2.048731in}{1.346427in}}%
\pgfpathlineto{\pgfqpoint{2.048731in}{1.346427in}}%
\pgfpathlineto{\pgfqpoint{2.048731in}{1.349376in}}%
\pgfpathlineto{\pgfqpoint{2.053272in}{1.349376in}}%
\pgfpathlineto{\pgfqpoint{2.053272in}{1.346427in}}%
\pgfpathmoveto{\pgfqpoint{2.053272in}{1.343477in}}%
\pgfpathlineto{\pgfqpoint{2.053272in}{1.343477in}}%
\pgfpathlineto{\pgfqpoint{2.053272in}{1.346427in}}%
\pgfpathlineto{\pgfqpoint{2.057813in}{1.346427in}}%
\pgfpathlineto{\pgfqpoint{2.057813in}{1.343477in}}%
\pgfpathmoveto{\pgfqpoint{2.053272in}{1.346427in}}%
\pgfpathlineto{\pgfqpoint{2.053272in}{1.346427in}}%
\pgfpathlineto{\pgfqpoint{2.053272in}{1.349376in}}%
\pgfpathlineto{\pgfqpoint{2.057813in}{1.349376in}}%
\pgfpathlineto{\pgfqpoint{2.057813in}{1.346427in}}%
\pgfpathmoveto{\pgfqpoint{2.057813in}{1.346427in}}%
\pgfpathlineto{\pgfqpoint{2.057813in}{1.346427in}}%
\pgfpathlineto{\pgfqpoint{2.057813in}{1.349376in}}%
\pgfpathlineto{\pgfqpoint{2.062354in}{1.349376in}}%
\pgfpathlineto{\pgfqpoint{2.062354in}{1.346427in}}%
\pgfpathmoveto{\pgfqpoint{2.057813in}{1.349376in}}%
\pgfpathlineto{\pgfqpoint{2.057813in}{1.349376in}}%
\pgfpathlineto{\pgfqpoint{2.057813in}{1.352325in}}%
\pgfpathlineto{\pgfqpoint{2.062354in}{1.352325in}}%
\pgfpathlineto{\pgfqpoint{2.062354in}{1.349376in}}%
\pgfpathmoveto{\pgfqpoint{2.057813in}{1.352325in}}%
\pgfpathlineto{\pgfqpoint{2.057813in}{1.352325in}}%
\pgfpathlineto{\pgfqpoint{2.057813in}{1.355274in}}%
\pgfpathlineto{\pgfqpoint{2.062354in}{1.355274in}}%
\pgfpathlineto{\pgfqpoint{2.062354in}{1.352325in}}%
\pgfpathmoveto{\pgfqpoint{2.062354in}{1.349376in}}%
\pgfpathlineto{\pgfqpoint{2.062354in}{1.349376in}}%
\pgfpathlineto{\pgfqpoint{2.062354in}{1.352325in}}%
\pgfpathlineto{\pgfqpoint{2.066896in}{1.352325in}}%
\pgfpathlineto{\pgfqpoint{2.066896in}{1.349376in}}%
\pgfpathmoveto{\pgfqpoint{2.062354in}{1.352325in}}%
\pgfpathlineto{\pgfqpoint{2.062354in}{1.352325in}}%
\pgfpathlineto{\pgfqpoint{2.062354in}{1.355274in}}%
\pgfpathlineto{\pgfqpoint{2.066896in}{1.355274in}}%
\pgfpathlineto{\pgfqpoint{2.066896in}{1.352325in}}%
\pgfpathmoveto{\pgfqpoint{2.066896in}{1.352325in}}%
\pgfpathlineto{\pgfqpoint{2.066896in}{1.352325in}}%
\pgfpathlineto{\pgfqpoint{2.066896in}{1.355274in}}%
\pgfpathlineto{\pgfqpoint{2.071437in}{1.355274in}}%
\pgfpathlineto{\pgfqpoint{2.071437in}{1.352325in}}%
\pgfpathmoveto{\pgfqpoint{2.066896in}{1.355274in}}%
\pgfpathlineto{\pgfqpoint{2.066896in}{1.355274in}}%
\pgfpathlineto{\pgfqpoint{2.066896in}{1.358223in}}%
\pgfpathlineto{\pgfqpoint{2.071437in}{1.358223in}}%
\pgfpathlineto{\pgfqpoint{2.071437in}{1.355274in}}%
\pgfpathmoveto{\pgfqpoint{2.066896in}{1.358223in}}%
\pgfpathlineto{\pgfqpoint{2.066896in}{1.358223in}}%
\pgfpathlineto{\pgfqpoint{2.066896in}{1.361173in}}%
\pgfpathlineto{\pgfqpoint{2.071437in}{1.361173in}}%
\pgfpathlineto{\pgfqpoint{2.071437in}{1.358223in}}%
\pgfpathmoveto{\pgfqpoint{2.071437in}{1.355274in}}%
\pgfpathlineto{\pgfqpoint{2.071437in}{1.355274in}}%
\pgfpathlineto{\pgfqpoint{2.071437in}{1.358223in}}%
\pgfpathlineto{\pgfqpoint{2.075978in}{1.358223in}}%
\pgfpathlineto{\pgfqpoint{2.075978in}{1.355274in}}%
\pgfpathmoveto{\pgfqpoint{2.071437in}{1.358223in}}%
\pgfpathlineto{\pgfqpoint{2.071437in}{1.358223in}}%
\pgfpathlineto{\pgfqpoint{2.071437in}{1.361173in}}%
\pgfpathlineto{\pgfqpoint{2.075978in}{1.361173in}}%
\pgfpathlineto{\pgfqpoint{2.075978in}{1.358223in}}%
\pgfpathmoveto{\pgfqpoint{2.066896in}{1.361173in}}%
\pgfpathlineto{\pgfqpoint{2.066896in}{1.361173in}}%
\pgfpathlineto{\pgfqpoint{2.066896in}{1.364122in}}%
\pgfpathlineto{\pgfqpoint{2.071437in}{1.364122in}}%
\pgfpathlineto{\pgfqpoint{2.071437in}{1.361173in}}%
\pgfpathmoveto{\pgfqpoint{2.066896in}{1.364122in}}%
\pgfpathlineto{\pgfqpoint{2.066896in}{1.364122in}}%
\pgfpathlineto{\pgfqpoint{2.066896in}{1.367071in}}%
\pgfpathlineto{\pgfqpoint{2.071437in}{1.367071in}}%
\pgfpathlineto{\pgfqpoint{2.071437in}{1.364122in}}%
\pgfpathmoveto{\pgfqpoint{2.071437in}{1.361173in}}%
\pgfpathlineto{\pgfqpoint{2.071437in}{1.361173in}}%
\pgfpathlineto{\pgfqpoint{2.071437in}{1.364122in}}%
\pgfpathlineto{\pgfqpoint{2.075978in}{1.364122in}}%
\pgfpathlineto{\pgfqpoint{2.075978in}{1.361173in}}%
\pgfpathmoveto{\pgfqpoint{2.071437in}{1.364122in}}%
\pgfpathlineto{\pgfqpoint{2.071437in}{1.364122in}}%
\pgfpathlineto{\pgfqpoint{2.071437in}{1.367071in}}%
\pgfpathlineto{\pgfqpoint{2.075978in}{1.367071in}}%
\pgfpathlineto{\pgfqpoint{2.075978in}{1.364122in}}%
\pgfpathmoveto{\pgfqpoint{2.075978in}{1.361173in}}%
\pgfpathlineto{\pgfqpoint{2.075978in}{1.361173in}}%
\pgfpathlineto{\pgfqpoint{2.075978in}{1.364122in}}%
\pgfpathlineto{\pgfqpoint{2.080519in}{1.364122in}}%
\pgfpathlineto{\pgfqpoint{2.080519in}{1.361173in}}%
\pgfpathmoveto{\pgfqpoint{2.075978in}{1.364122in}}%
\pgfpathlineto{\pgfqpoint{2.075978in}{1.364122in}}%
\pgfpathlineto{\pgfqpoint{2.075978in}{1.367071in}}%
\pgfpathlineto{\pgfqpoint{2.080519in}{1.367071in}}%
\pgfpathlineto{\pgfqpoint{2.080519in}{1.364122in}}%
\pgfpathmoveto{\pgfqpoint{2.080519in}{1.364122in}}%
\pgfpathlineto{\pgfqpoint{2.080519in}{1.364122in}}%
\pgfpathlineto{\pgfqpoint{2.080519in}{1.367071in}}%
\pgfpathlineto{\pgfqpoint{2.085060in}{1.367071in}}%
\pgfpathlineto{\pgfqpoint{2.085060in}{1.364122in}}%
\pgfpathmoveto{\pgfqpoint{2.075978in}{1.367071in}}%
\pgfpathlineto{\pgfqpoint{2.075978in}{1.367071in}}%
\pgfpathlineto{\pgfqpoint{2.075978in}{1.370020in}}%
\pgfpathlineto{\pgfqpoint{2.080519in}{1.370020in}}%
\pgfpathlineto{\pgfqpoint{2.080519in}{1.367071in}}%
\pgfpathmoveto{\pgfqpoint{2.075978in}{1.370020in}}%
\pgfpathlineto{\pgfqpoint{2.075978in}{1.370020in}}%
\pgfpathlineto{\pgfqpoint{2.075978in}{1.372969in}}%
\pgfpathlineto{\pgfqpoint{2.080519in}{1.372969in}}%
\pgfpathlineto{\pgfqpoint{2.080519in}{1.370020in}}%
\pgfpathmoveto{\pgfqpoint{2.080519in}{1.367071in}}%
\pgfpathlineto{\pgfqpoint{2.080519in}{1.367071in}}%
\pgfpathlineto{\pgfqpoint{2.080519in}{1.370020in}}%
\pgfpathlineto{\pgfqpoint{2.085060in}{1.370020in}}%
\pgfpathlineto{\pgfqpoint{2.085060in}{1.367071in}}%
\pgfpathmoveto{\pgfqpoint{2.080519in}{1.370020in}}%
\pgfpathlineto{\pgfqpoint{2.080519in}{1.370020in}}%
\pgfpathlineto{\pgfqpoint{2.080519in}{1.372969in}}%
\pgfpathlineto{\pgfqpoint{2.085060in}{1.372969in}}%
\pgfpathlineto{\pgfqpoint{2.085060in}{1.370020in}}%
\pgfpathmoveto{\pgfqpoint{2.085060in}{1.367071in}}%
\pgfpathlineto{\pgfqpoint{2.085060in}{1.367071in}}%
\pgfpathlineto{\pgfqpoint{2.085060in}{1.370020in}}%
\pgfpathlineto{\pgfqpoint{2.089601in}{1.370020in}}%
\pgfpathlineto{\pgfqpoint{2.089601in}{1.367071in}}%
\pgfpathmoveto{\pgfqpoint{2.085060in}{1.370020in}}%
\pgfpathlineto{\pgfqpoint{2.085060in}{1.370020in}}%
\pgfpathlineto{\pgfqpoint{2.085060in}{1.372969in}}%
\pgfpathlineto{\pgfqpoint{2.089601in}{1.372969in}}%
\pgfpathlineto{\pgfqpoint{2.089601in}{1.370020in}}%
\pgfpathmoveto{\pgfqpoint{2.089601in}{1.370020in}}%
\pgfpathlineto{\pgfqpoint{2.089601in}{1.370020in}}%
\pgfpathlineto{\pgfqpoint{2.089601in}{1.372969in}}%
\pgfpathlineto{\pgfqpoint{2.094142in}{1.372969in}}%
\pgfpathlineto{\pgfqpoint{2.094142in}{1.370020in}}%
\pgfpathmoveto{\pgfqpoint{2.085060in}{1.372969in}}%
\pgfpathlineto{\pgfqpoint{2.085060in}{1.372969in}}%
\pgfpathlineto{\pgfqpoint{2.085060in}{1.375919in}}%
\pgfpathlineto{\pgfqpoint{2.089601in}{1.375919in}}%
\pgfpathlineto{\pgfqpoint{2.089601in}{1.372969in}}%
\pgfpathmoveto{\pgfqpoint{2.085060in}{1.375919in}}%
\pgfpathlineto{\pgfqpoint{2.085060in}{1.375919in}}%
\pgfpathlineto{\pgfqpoint{2.085060in}{1.378868in}}%
\pgfpathlineto{\pgfqpoint{2.089601in}{1.378868in}}%
\pgfpathlineto{\pgfqpoint{2.089601in}{1.375919in}}%
\pgfpathmoveto{\pgfqpoint{2.089601in}{1.372969in}}%
\pgfpathlineto{\pgfqpoint{2.089601in}{1.372969in}}%
\pgfpathlineto{\pgfqpoint{2.089601in}{1.375919in}}%
\pgfpathlineto{\pgfqpoint{2.094142in}{1.375919in}}%
\pgfpathlineto{\pgfqpoint{2.094142in}{1.372969in}}%
\pgfpathmoveto{\pgfqpoint{2.089601in}{1.375919in}}%
\pgfpathlineto{\pgfqpoint{2.089601in}{1.375919in}}%
\pgfpathlineto{\pgfqpoint{2.089601in}{1.378868in}}%
\pgfpathlineto{\pgfqpoint{2.094142in}{1.378868in}}%
\pgfpathlineto{\pgfqpoint{2.094142in}{1.375919in}}%
\pgfpathmoveto{\pgfqpoint{2.094142in}{1.372969in}}%
\pgfpathlineto{\pgfqpoint{2.094142in}{1.372969in}}%
\pgfpathlineto{\pgfqpoint{2.094142in}{1.375919in}}%
\pgfpathlineto{\pgfqpoint{2.098683in}{1.375919in}}%
\pgfpathlineto{\pgfqpoint{2.098683in}{1.372969in}}%
\pgfpathmoveto{\pgfqpoint{2.094142in}{1.375919in}}%
\pgfpathlineto{\pgfqpoint{2.094142in}{1.375919in}}%
\pgfpathlineto{\pgfqpoint{2.094142in}{1.378868in}}%
\pgfpathlineto{\pgfqpoint{2.098683in}{1.378868in}}%
\pgfpathlineto{\pgfqpoint{2.098683in}{1.375919in}}%
\pgfpathmoveto{\pgfqpoint{2.098683in}{1.375919in}}%
\pgfpathlineto{\pgfqpoint{2.098683in}{1.375919in}}%
\pgfpathlineto{\pgfqpoint{2.098683in}{1.378868in}}%
\pgfpathlineto{\pgfqpoint{2.103225in}{1.378868in}}%
\pgfpathlineto{\pgfqpoint{2.103225in}{1.375919in}}%
\pgfpathmoveto{\pgfqpoint{2.094142in}{1.378868in}}%
\pgfpathlineto{\pgfqpoint{2.094142in}{1.378868in}}%
\pgfpathlineto{\pgfqpoint{2.094142in}{1.381817in}}%
\pgfpathlineto{\pgfqpoint{2.098683in}{1.381817in}}%
\pgfpathlineto{\pgfqpoint{2.098683in}{1.378868in}}%
\pgfpathmoveto{\pgfqpoint{2.094142in}{1.381817in}}%
\pgfpathlineto{\pgfqpoint{2.094142in}{1.381817in}}%
\pgfpathlineto{\pgfqpoint{2.094142in}{1.384766in}}%
\pgfpathlineto{\pgfqpoint{2.098683in}{1.384766in}}%
\pgfpathlineto{\pgfqpoint{2.098683in}{1.381817in}}%
\pgfpathmoveto{\pgfqpoint{2.098683in}{1.378868in}}%
\pgfpathlineto{\pgfqpoint{2.098683in}{1.378868in}}%
\pgfpathlineto{\pgfqpoint{2.098683in}{1.381817in}}%
\pgfpathlineto{\pgfqpoint{2.103225in}{1.381817in}}%
\pgfpathlineto{\pgfqpoint{2.103225in}{1.378868in}}%
\pgfpathmoveto{\pgfqpoint{2.098683in}{1.381817in}}%
\pgfpathlineto{\pgfqpoint{2.098683in}{1.381817in}}%
\pgfpathlineto{\pgfqpoint{2.098683in}{1.384766in}}%
\pgfpathlineto{\pgfqpoint{2.103225in}{1.384766in}}%
\pgfpathlineto{\pgfqpoint{2.103225in}{1.381817in}}%
\pgfpathmoveto{\pgfqpoint{2.103225in}{1.378868in}}%
\pgfpathlineto{\pgfqpoint{2.103225in}{1.378868in}}%
\pgfpathlineto{\pgfqpoint{2.103225in}{1.381817in}}%
\pgfpathlineto{\pgfqpoint{2.107766in}{1.381817in}}%
\pgfpathlineto{\pgfqpoint{2.107766in}{1.378868in}}%
\pgfpathmoveto{\pgfqpoint{2.103225in}{1.381817in}}%
\pgfpathlineto{\pgfqpoint{2.103225in}{1.381817in}}%
\pgfpathlineto{\pgfqpoint{2.103225in}{1.384766in}}%
\pgfpathlineto{\pgfqpoint{2.107766in}{1.384766in}}%
\pgfpathlineto{\pgfqpoint{2.107766in}{1.381817in}}%
\pgfpathmoveto{\pgfqpoint{2.107766in}{1.381817in}}%
\pgfpathlineto{\pgfqpoint{2.107766in}{1.381817in}}%
\pgfpathlineto{\pgfqpoint{2.107766in}{1.384766in}}%
\pgfpathlineto{\pgfqpoint{2.112307in}{1.384766in}}%
\pgfpathlineto{\pgfqpoint{2.112307in}{1.381817in}}%
\pgfpathmoveto{\pgfqpoint{2.103225in}{1.384766in}}%
\pgfpathlineto{\pgfqpoint{2.103225in}{1.384766in}}%
\pgfpathlineto{\pgfqpoint{2.103225in}{1.387715in}}%
\pgfpathlineto{\pgfqpoint{2.107766in}{1.387715in}}%
\pgfpathlineto{\pgfqpoint{2.107766in}{1.384766in}}%
\pgfpathmoveto{\pgfqpoint{2.103225in}{1.387715in}}%
\pgfpathlineto{\pgfqpoint{2.103225in}{1.387715in}}%
\pgfpathlineto{\pgfqpoint{2.103225in}{1.390665in}}%
\pgfpathlineto{\pgfqpoint{2.107766in}{1.390665in}}%
\pgfpathlineto{\pgfqpoint{2.107766in}{1.387715in}}%
\pgfpathmoveto{\pgfqpoint{2.107766in}{1.384766in}}%
\pgfpathlineto{\pgfqpoint{2.107766in}{1.384766in}}%
\pgfpathlineto{\pgfqpoint{2.107766in}{1.387715in}}%
\pgfpathlineto{\pgfqpoint{2.112307in}{1.387715in}}%
\pgfpathlineto{\pgfqpoint{2.112307in}{1.384766in}}%
\pgfpathmoveto{\pgfqpoint{2.107766in}{1.387715in}}%
\pgfpathlineto{\pgfqpoint{2.107766in}{1.387715in}}%
\pgfpathlineto{\pgfqpoint{2.107766in}{1.390665in}}%
\pgfpathlineto{\pgfqpoint{2.112307in}{1.390665in}}%
\pgfpathlineto{\pgfqpoint{2.112307in}{1.387715in}}%
\pgfpathmoveto{\pgfqpoint{2.112307in}{1.384766in}}%
\pgfpathlineto{\pgfqpoint{2.112307in}{1.384766in}}%
\pgfpathlineto{\pgfqpoint{2.112307in}{1.387715in}}%
\pgfpathlineto{\pgfqpoint{2.116848in}{1.387715in}}%
\pgfpathlineto{\pgfqpoint{2.116848in}{1.384766in}}%
\pgfpathmoveto{\pgfqpoint{2.112307in}{1.387715in}}%
\pgfpathlineto{\pgfqpoint{2.112307in}{1.387715in}}%
\pgfpathlineto{\pgfqpoint{2.112307in}{1.390665in}}%
\pgfpathlineto{\pgfqpoint{2.116848in}{1.390665in}}%
\pgfpathlineto{\pgfqpoint{2.116848in}{1.387715in}}%
\pgfpathmoveto{\pgfqpoint{2.116848in}{1.387715in}}%
\pgfpathlineto{\pgfqpoint{2.116848in}{1.387715in}}%
\pgfpathlineto{\pgfqpoint{2.116848in}{1.390665in}}%
\pgfpathlineto{\pgfqpoint{2.121389in}{1.390665in}}%
\pgfpathlineto{\pgfqpoint{2.121389in}{1.387715in}}%
\pgfpathmoveto{\pgfqpoint{2.112307in}{1.390665in}}%
\pgfpathlineto{\pgfqpoint{2.112307in}{1.390665in}}%
\pgfpathlineto{\pgfqpoint{2.112307in}{1.393614in}}%
\pgfpathlineto{\pgfqpoint{2.116848in}{1.393614in}}%
\pgfpathlineto{\pgfqpoint{2.116848in}{1.390665in}}%
\pgfpathmoveto{\pgfqpoint{2.112307in}{1.393614in}}%
\pgfpathlineto{\pgfqpoint{2.112307in}{1.393614in}}%
\pgfpathlineto{\pgfqpoint{2.112307in}{1.396563in}}%
\pgfpathlineto{\pgfqpoint{2.116848in}{1.396563in}}%
\pgfpathlineto{\pgfqpoint{2.116848in}{1.393614in}}%
\pgfpathmoveto{\pgfqpoint{2.116848in}{1.390665in}}%
\pgfpathlineto{\pgfqpoint{2.116848in}{1.390665in}}%
\pgfpathlineto{\pgfqpoint{2.116848in}{1.393614in}}%
\pgfpathlineto{\pgfqpoint{2.121389in}{1.393614in}}%
\pgfpathlineto{\pgfqpoint{2.121389in}{1.390665in}}%
\pgfpathmoveto{\pgfqpoint{2.116848in}{1.393614in}}%
\pgfpathlineto{\pgfqpoint{2.116848in}{1.393614in}}%
\pgfpathlineto{\pgfqpoint{2.116848in}{1.396563in}}%
\pgfpathlineto{\pgfqpoint{2.121389in}{1.396563in}}%
\pgfpathlineto{\pgfqpoint{2.121389in}{1.393614in}}%
\pgfpathmoveto{\pgfqpoint{2.121389in}{1.390665in}}%
\pgfpathlineto{\pgfqpoint{2.121389in}{1.390665in}}%
\pgfpathlineto{\pgfqpoint{2.121389in}{1.393614in}}%
\pgfpathlineto{\pgfqpoint{2.125930in}{1.393614in}}%
\pgfpathlineto{\pgfqpoint{2.125930in}{1.390665in}}%
\pgfpathmoveto{\pgfqpoint{2.121389in}{1.393614in}}%
\pgfpathlineto{\pgfqpoint{2.121389in}{1.393614in}}%
\pgfpathlineto{\pgfqpoint{2.121389in}{1.396563in}}%
\pgfpathlineto{\pgfqpoint{2.125930in}{1.396563in}}%
\pgfpathlineto{\pgfqpoint{2.125930in}{1.393614in}}%
\pgfpathmoveto{\pgfqpoint{2.125930in}{1.393614in}}%
\pgfpathlineto{\pgfqpoint{2.125930in}{1.393614in}}%
\pgfpathlineto{\pgfqpoint{2.125930in}{1.396563in}}%
\pgfpathlineto{\pgfqpoint{2.130471in}{1.396563in}}%
\pgfpathlineto{\pgfqpoint{2.130471in}{1.393614in}}%
\pgfpathmoveto{\pgfqpoint{2.121389in}{1.396563in}}%
\pgfpathlineto{\pgfqpoint{2.121389in}{1.396563in}}%
\pgfpathlineto{\pgfqpoint{2.121389in}{1.399512in}}%
\pgfpathlineto{\pgfqpoint{2.125930in}{1.399512in}}%
\pgfpathlineto{\pgfqpoint{2.125930in}{1.396563in}}%
\pgfpathmoveto{\pgfqpoint{2.121389in}{1.399512in}}%
\pgfpathlineto{\pgfqpoint{2.121389in}{1.399512in}}%
\pgfpathlineto{\pgfqpoint{2.121389in}{1.402461in}}%
\pgfpathlineto{\pgfqpoint{2.125930in}{1.402461in}}%
\pgfpathlineto{\pgfqpoint{2.125930in}{1.399512in}}%
\pgfpathmoveto{\pgfqpoint{2.125930in}{1.396563in}}%
\pgfpathlineto{\pgfqpoint{2.125930in}{1.396563in}}%
\pgfpathlineto{\pgfqpoint{2.125930in}{1.399512in}}%
\pgfpathlineto{\pgfqpoint{2.130471in}{1.399512in}}%
\pgfpathlineto{\pgfqpoint{2.130471in}{1.396563in}}%
\pgfpathmoveto{\pgfqpoint{2.125930in}{1.399512in}}%
\pgfpathlineto{\pgfqpoint{2.125930in}{1.399512in}}%
\pgfpathlineto{\pgfqpoint{2.125930in}{1.402461in}}%
\pgfpathlineto{\pgfqpoint{2.130471in}{1.402461in}}%
\pgfpathlineto{\pgfqpoint{2.130471in}{1.399512in}}%
\pgfpathmoveto{\pgfqpoint{2.130471in}{1.396563in}}%
\pgfpathlineto{\pgfqpoint{2.130471in}{1.396563in}}%
\pgfpathlineto{\pgfqpoint{2.130471in}{1.399512in}}%
\pgfpathlineto{\pgfqpoint{2.135012in}{1.399512in}}%
\pgfpathlineto{\pgfqpoint{2.135012in}{1.396563in}}%
\pgfpathmoveto{\pgfqpoint{2.130471in}{1.399512in}}%
\pgfpathlineto{\pgfqpoint{2.130471in}{1.399512in}}%
\pgfpathlineto{\pgfqpoint{2.130471in}{1.402461in}}%
\pgfpathlineto{\pgfqpoint{2.135012in}{1.402461in}}%
\pgfpathlineto{\pgfqpoint{2.135012in}{1.399512in}}%
\pgfpathmoveto{\pgfqpoint{2.135012in}{1.399512in}}%
\pgfpathlineto{\pgfqpoint{2.135012in}{1.399512in}}%
\pgfpathlineto{\pgfqpoint{2.135012in}{1.402461in}}%
\pgfpathlineto{\pgfqpoint{2.139554in}{1.402461in}}%
\pgfpathlineto{\pgfqpoint{2.139554in}{1.399512in}}%
\pgfpathmoveto{\pgfqpoint{2.130471in}{1.402461in}}%
\pgfpathlineto{\pgfqpoint{2.130471in}{1.402461in}}%
\pgfpathlineto{\pgfqpoint{2.130471in}{1.405411in}}%
\pgfpathlineto{\pgfqpoint{2.135012in}{1.405411in}}%
\pgfpathlineto{\pgfqpoint{2.135012in}{1.402461in}}%
\pgfpathmoveto{\pgfqpoint{2.130471in}{1.405411in}}%
\pgfpathlineto{\pgfqpoint{2.130471in}{1.405411in}}%
\pgfpathlineto{\pgfqpoint{2.130471in}{1.408360in}}%
\pgfpathlineto{\pgfqpoint{2.135012in}{1.408360in}}%
\pgfpathlineto{\pgfqpoint{2.135012in}{1.405411in}}%
\pgfpathmoveto{\pgfqpoint{2.135012in}{1.402461in}}%
\pgfpathlineto{\pgfqpoint{2.135012in}{1.402461in}}%
\pgfpathlineto{\pgfqpoint{2.135012in}{1.405411in}}%
\pgfpathlineto{\pgfqpoint{2.139554in}{1.405411in}}%
\pgfpathlineto{\pgfqpoint{2.139554in}{1.402461in}}%
\pgfpathmoveto{\pgfqpoint{2.135012in}{1.405411in}}%
\pgfpathlineto{\pgfqpoint{2.135012in}{1.405411in}}%
\pgfpathlineto{\pgfqpoint{2.135012in}{1.408360in}}%
\pgfpathlineto{\pgfqpoint{2.139554in}{1.408360in}}%
\pgfpathlineto{\pgfqpoint{2.139554in}{1.405411in}}%
\pgfpathmoveto{\pgfqpoint{2.139554in}{1.402461in}}%
\pgfpathlineto{\pgfqpoint{2.139554in}{1.402461in}}%
\pgfpathlineto{\pgfqpoint{2.139554in}{1.405411in}}%
\pgfpathlineto{\pgfqpoint{2.144095in}{1.405411in}}%
\pgfpathlineto{\pgfqpoint{2.144095in}{1.402461in}}%
\pgfpathmoveto{\pgfqpoint{2.139554in}{1.405411in}}%
\pgfpathlineto{\pgfqpoint{2.139554in}{1.405411in}}%
\pgfpathlineto{\pgfqpoint{2.139554in}{1.408360in}}%
\pgfpathlineto{\pgfqpoint{2.144095in}{1.408360in}}%
\pgfpathlineto{\pgfqpoint{2.144095in}{1.405411in}}%
\pgfpathmoveto{\pgfqpoint{2.144095in}{1.405411in}}%
\pgfpathlineto{\pgfqpoint{2.144095in}{1.405411in}}%
\pgfpathlineto{\pgfqpoint{2.144095in}{1.408360in}}%
\pgfpathlineto{\pgfqpoint{2.148636in}{1.408360in}}%
\pgfpathlineto{\pgfqpoint{2.148636in}{1.405411in}}%
\pgfpathmoveto{\pgfqpoint{2.139554in}{1.408360in}}%
\pgfpathlineto{\pgfqpoint{2.139554in}{1.408360in}}%
\pgfpathlineto{\pgfqpoint{2.139554in}{1.411309in}}%
\pgfpathlineto{\pgfqpoint{2.144095in}{1.411309in}}%
\pgfpathlineto{\pgfqpoint{2.144095in}{1.408360in}}%
\pgfpathmoveto{\pgfqpoint{2.139554in}{1.411309in}}%
\pgfpathlineto{\pgfqpoint{2.139554in}{1.411309in}}%
\pgfpathlineto{\pgfqpoint{2.139554in}{1.414258in}}%
\pgfpathlineto{\pgfqpoint{2.144095in}{1.414258in}}%
\pgfpathlineto{\pgfqpoint{2.144095in}{1.411309in}}%
\pgfpathmoveto{\pgfqpoint{2.144095in}{1.408360in}}%
\pgfpathlineto{\pgfqpoint{2.144095in}{1.408360in}}%
\pgfpathlineto{\pgfqpoint{2.144095in}{1.411309in}}%
\pgfpathlineto{\pgfqpoint{2.148636in}{1.411309in}}%
\pgfpathlineto{\pgfqpoint{2.148636in}{1.408360in}}%
\pgfpathmoveto{\pgfqpoint{2.144095in}{1.411309in}}%
\pgfpathlineto{\pgfqpoint{2.144095in}{1.411309in}}%
\pgfpathlineto{\pgfqpoint{2.144095in}{1.414258in}}%
\pgfpathlineto{\pgfqpoint{2.148636in}{1.414258in}}%
\pgfpathlineto{\pgfqpoint{2.148636in}{1.411309in}}%
\pgfpathmoveto{\pgfqpoint{2.148636in}{1.408360in}}%
\pgfpathlineto{\pgfqpoint{2.148636in}{1.408360in}}%
\pgfpathlineto{\pgfqpoint{2.148636in}{1.411309in}}%
\pgfpathlineto{\pgfqpoint{2.153177in}{1.411309in}}%
\pgfpathlineto{\pgfqpoint{2.153177in}{1.408360in}}%
\pgfpathmoveto{\pgfqpoint{2.148636in}{1.411309in}}%
\pgfpathlineto{\pgfqpoint{2.148636in}{1.411309in}}%
\pgfpathlineto{\pgfqpoint{2.148636in}{1.414258in}}%
\pgfpathlineto{\pgfqpoint{2.153177in}{1.414258in}}%
\pgfpathlineto{\pgfqpoint{2.153177in}{1.411309in}}%
\pgfpathmoveto{\pgfqpoint{2.153177in}{1.411309in}}%
\pgfpathlineto{\pgfqpoint{2.153177in}{1.411309in}}%
\pgfpathlineto{\pgfqpoint{2.153177in}{1.414258in}}%
\pgfpathlineto{\pgfqpoint{2.157718in}{1.414258in}}%
\pgfpathlineto{\pgfqpoint{2.157718in}{1.411309in}}%
\pgfpathmoveto{\pgfqpoint{2.148636in}{1.414258in}}%
\pgfpathlineto{\pgfqpoint{2.148636in}{1.414258in}}%
\pgfpathlineto{\pgfqpoint{2.148636in}{1.417207in}}%
\pgfpathlineto{\pgfqpoint{2.153177in}{1.417207in}}%
\pgfpathlineto{\pgfqpoint{2.153177in}{1.414258in}}%
\pgfpathmoveto{\pgfqpoint{2.148636in}{1.417207in}}%
\pgfpathlineto{\pgfqpoint{2.148636in}{1.417207in}}%
\pgfpathlineto{\pgfqpoint{2.148636in}{1.420156in}}%
\pgfpathlineto{\pgfqpoint{2.153177in}{1.420156in}}%
\pgfpathlineto{\pgfqpoint{2.153177in}{1.417207in}}%
\pgfpathmoveto{\pgfqpoint{2.153177in}{1.414258in}}%
\pgfpathlineto{\pgfqpoint{2.153177in}{1.414258in}}%
\pgfpathlineto{\pgfqpoint{2.153177in}{1.417207in}}%
\pgfpathlineto{\pgfqpoint{2.157718in}{1.417207in}}%
\pgfpathlineto{\pgfqpoint{2.157718in}{1.414258in}}%
\pgfpathmoveto{\pgfqpoint{2.153177in}{1.417207in}}%
\pgfpathlineto{\pgfqpoint{2.153177in}{1.417207in}}%
\pgfpathlineto{\pgfqpoint{2.153177in}{1.420156in}}%
\pgfpathlineto{\pgfqpoint{2.157718in}{1.420156in}}%
\pgfpathlineto{\pgfqpoint{2.157718in}{1.417207in}}%
\pgfpathmoveto{\pgfqpoint{2.157718in}{1.414258in}}%
\pgfpathlineto{\pgfqpoint{2.157718in}{1.414258in}}%
\pgfpathlineto{\pgfqpoint{2.157718in}{1.417207in}}%
\pgfpathlineto{\pgfqpoint{2.162259in}{1.417207in}}%
\pgfpathlineto{\pgfqpoint{2.162259in}{1.414258in}}%
\pgfpathmoveto{\pgfqpoint{2.157718in}{1.417207in}}%
\pgfpathlineto{\pgfqpoint{2.157718in}{1.417207in}}%
\pgfpathlineto{\pgfqpoint{2.157718in}{1.420156in}}%
\pgfpathlineto{\pgfqpoint{2.162259in}{1.420156in}}%
\pgfpathlineto{\pgfqpoint{2.162259in}{1.417207in}}%
\pgfpathmoveto{\pgfqpoint{2.162259in}{1.417207in}}%
\pgfpathlineto{\pgfqpoint{2.162259in}{1.417207in}}%
\pgfpathlineto{\pgfqpoint{2.162259in}{1.420156in}}%
\pgfpathlineto{\pgfqpoint{2.166800in}{1.420156in}}%
\pgfpathlineto{\pgfqpoint{2.166800in}{1.417207in}}%
\pgfpathmoveto{\pgfqpoint{2.157718in}{1.420156in}}%
\pgfpathlineto{\pgfqpoint{2.157718in}{1.420156in}}%
\pgfpathlineto{\pgfqpoint{2.157718in}{1.423106in}}%
\pgfpathlineto{\pgfqpoint{2.162259in}{1.423106in}}%
\pgfpathlineto{\pgfqpoint{2.162259in}{1.420156in}}%
\pgfpathmoveto{\pgfqpoint{2.157718in}{1.423106in}}%
\pgfpathlineto{\pgfqpoint{2.157718in}{1.423106in}}%
\pgfpathlineto{\pgfqpoint{2.157718in}{1.426055in}}%
\pgfpathlineto{\pgfqpoint{2.162259in}{1.426055in}}%
\pgfpathlineto{\pgfqpoint{2.162259in}{1.423106in}}%
\pgfpathmoveto{\pgfqpoint{2.162259in}{1.420156in}}%
\pgfpathlineto{\pgfqpoint{2.162259in}{1.420156in}}%
\pgfpathlineto{\pgfqpoint{2.162259in}{1.423106in}}%
\pgfpathlineto{\pgfqpoint{2.166800in}{1.423106in}}%
\pgfpathlineto{\pgfqpoint{2.166800in}{1.420156in}}%
\pgfpathmoveto{\pgfqpoint{2.162259in}{1.423106in}}%
\pgfpathlineto{\pgfqpoint{2.162259in}{1.423106in}}%
\pgfpathlineto{\pgfqpoint{2.162259in}{1.426055in}}%
\pgfpathlineto{\pgfqpoint{2.166800in}{1.426055in}}%
\pgfpathlineto{\pgfqpoint{2.166800in}{1.423106in}}%
\pgfpathmoveto{\pgfqpoint{2.166800in}{1.420156in}}%
\pgfpathlineto{\pgfqpoint{2.166800in}{1.420156in}}%
\pgfpathlineto{\pgfqpoint{2.166800in}{1.423106in}}%
\pgfpathlineto{\pgfqpoint{2.171341in}{1.423106in}}%
\pgfpathlineto{\pgfqpoint{2.171341in}{1.420156in}}%
\pgfpathmoveto{\pgfqpoint{2.166800in}{1.423106in}}%
\pgfpathlineto{\pgfqpoint{2.166800in}{1.423106in}}%
\pgfpathlineto{\pgfqpoint{2.166800in}{1.426055in}}%
\pgfpathlineto{\pgfqpoint{2.171341in}{1.426055in}}%
\pgfpathlineto{\pgfqpoint{2.171341in}{1.423106in}}%
\pgfpathmoveto{\pgfqpoint{2.171341in}{1.423106in}}%
\pgfpathlineto{\pgfqpoint{2.171341in}{1.423106in}}%
\pgfpathlineto{\pgfqpoint{2.171341in}{1.426055in}}%
\pgfpathlineto{\pgfqpoint{2.175882in}{1.426055in}}%
\pgfpathlineto{\pgfqpoint{2.175882in}{1.423106in}}%
\pgfpathmoveto{\pgfqpoint{2.166800in}{1.426055in}}%
\pgfpathlineto{\pgfqpoint{2.166800in}{1.426055in}}%
\pgfpathlineto{\pgfqpoint{2.166800in}{1.429004in}}%
\pgfpathlineto{\pgfqpoint{2.171341in}{1.429004in}}%
\pgfpathlineto{\pgfqpoint{2.171341in}{1.426055in}}%
\pgfpathmoveto{\pgfqpoint{2.166800in}{1.429004in}}%
\pgfpathlineto{\pgfqpoint{2.166800in}{1.429004in}}%
\pgfpathlineto{\pgfqpoint{2.166800in}{1.431953in}}%
\pgfpathlineto{\pgfqpoint{2.171341in}{1.431953in}}%
\pgfpathlineto{\pgfqpoint{2.171341in}{1.429004in}}%
\pgfpathmoveto{\pgfqpoint{2.171341in}{1.426055in}}%
\pgfpathlineto{\pgfqpoint{2.171341in}{1.426055in}}%
\pgfpathlineto{\pgfqpoint{2.171341in}{1.429004in}}%
\pgfpathlineto{\pgfqpoint{2.175882in}{1.429004in}}%
\pgfpathlineto{\pgfqpoint{2.175882in}{1.426055in}}%
\pgfpathmoveto{\pgfqpoint{2.171341in}{1.429004in}}%
\pgfpathlineto{\pgfqpoint{2.171341in}{1.429004in}}%
\pgfpathlineto{\pgfqpoint{2.171341in}{1.431953in}}%
\pgfpathlineto{\pgfqpoint{2.175882in}{1.431953in}}%
\pgfpathlineto{\pgfqpoint{2.175882in}{1.429004in}}%
\pgfpathmoveto{\pgfqpoint{2.175882in}{1.426055in}}%
\pgfpathlineto{\pgfqpoint{2.175882in}{1.426055in}}%
\pgfpathlineto{\pgfqpoint{2.175882in}{1.429004in}}%
\pgfpathlineto{\pgfqpoint{2.180424in}{1.429004in}}%
\pgfpathlineto{\pgfqpoint{2.180424in}{1.426055in}}%
\pgfpathmoveto{\pgfqpoint{2.175882in}{1.429004in}}%
\pgfpathlineto{\pgfqpoint{2.175882in}{1.429004in}}%
\pgfpathlineto{\pgfqpoint{2.175882in}{1.431953in}}%
\pgfpathlineto{\pgfqpoint{2.180424in}{1.431953in}}%
\pgfpathlineto{\pgfqpoint{2.180424in}{1.429004in}}%
\pgfpathmoveto{\pgfqpoint{2.180424in}{1.429004in}}%
\pgfpathlineto{\pgfqpoint{2.180424in}{1.429004in}}%
\pgfpathlineto{\pgfqpoint{2.180424in}{1.431953in}}%
\pgfpathlineto{\pgfqpoint{2.184965in}{1.431953in}}%
\pgfpathlineto{\pgfqpoint{2.184965in}{1.429004in}}%
\pgfpathmoveto{\pgfqpoint{2.175882in}{1.431953in}}%
\pgfpathlineto{\pgfqpoint{2.175882in}{1.431953in}}%
\pgfpathlineto{\pgfqpoint{2.175882in}{1.434902in}}%
\pgfpathlineto{\pgfqpoint{2.180424in}{1.434902in}}%
\pgfpathlineto{\pgfqpoint{2.180424in}{1.431953in}}%
\pgfpathmoveto{\pgfqpoint{2.175882in}{1.434902in}}%
\pgfpathlineto{\pgfqpoint{2.175882in}{1.434902in}}%
\pgfpathlineto{\pgfqpoint{2.175882in}{1.437852in}}%
\pgfpathlineto{\pgfqpoint{2.180424in}{1.437852in}}%
\pgfpathlineto{\pgfqpoint{2.180424in}{1.434902in}}%
\pgfpathmoveto{\pgfqpoint{2.180424in}{1.431953in}}%
\pgfpathlineto{\pgfqpoint{2.180424in}{1.431953in}}%
\pgfpathlineto{\pgfqpoint{2.180424in}{1.434902in}}%
\pgfpathlineto{\pgfqpoint{2.184965in}{1.434902in}}%
\pgfpathlineto{\pgfqpoint{2.184965in}{1.431953in}}%
\pgfpathmoveto{\pgfqpoint{2.180424in}{1.434902in}}%
\pgfpathlineto{\pgfqpoint{2.180424in}{1.434902in}}%
\pgfpathlineto{\pgfqpoint{2.180424in}{1.437852in}}%
\pgfpathlineto{\pgfqpoint{2.184965in}{1.437852in}}%
\pgfpathlineto{\pgfqpoint{2.184965in}{1.434902in}}%
\pgfpathmoveto{\pgfqpoint{2.184965in}{1.431953in}}%
\pgfpathlineto{\pgfqpoint{2.184965in}{1.431953in}}%
\pgfpathlineto{\pgfqpoint{2.184965in}{1.434902in}}%
\pgfpathlineto{\pgfqpoint{2.189506in}{1.434902in}}%
\pgfpathlineto{\pgfqpoint{2.189506in}{1.431953in}}%
\pgfpathmoveto{\pgfqpoint{2.184965in}{1.434902in}}%
\pgfpathlineto{\pgfqpoint{2.184965in}{1.434902in}}%
\pgfpathlineto{\pgfqpoint{2.184965in}{1.437852in}}%
\pgfpathlineto{\pgfqpoint{2.189506in}{1.437852in}}%
\pgfpathlineto{\pgfqpoint{2.189506in}{1.434902in}}%
\pgfpathmoveto{\pgfqpoint{2.189506in}{1.434902in}}%
\pgfpathlineto{\pgfqpoint{2.189506in}{1.434902in}}%
\pgfpathlineto{\pgfqpoint{2.189506in}{1.437852in}}%
\pgfpathlineto{\pgfqpoint{2.194047in}{1.437852in}}%
\pgfpathlineto{\pgfqpoint{2.194047in}{1.434902in}}%
\pgfpathmoveto{\pgfqpoint{2.184965in}{1.437852in}}%
\pgfpathlineto{\pgfqpoint{2.184965in}{1.437852in}}%
\pgfpathlineto{\pgfqpoint{2.184965in}{1.440801in}}%
\pgfpathlineto{\pgfqpoint{2.189506in}{1.440801in}}%
\pgfpathlineto{\pgfqpoint{2.189506in}{1.437852in}}%
\pgfpathmoveto{\pgfqpoint{2.184965in}{1.440801in}}%
\pgfpathlineto{\pgfqpoint{2.184965in}{1.440801in}}%
\pgfpathlineto{\pgfqpoint{2.184965in}{1.443750in}}%
\pgfpathlineto{\pgfqpoint{2.189506in}{1.443750in}}%
\pgfpathlineto{\pgfqpoint{2.189506in}{1.440801in}}%
\pgfpathmoveto{\pgfqpoint{2.189506in}{1.437852in}}%
\pgfpathlineto{\pgfqpoint{2.189506in}{1.437852in}}%
\pgfpathlineto{\pgfqpoint{2.189506in}{1.440801in}}%
\pgfpathlineto{\pgfqpoint{2.194047in}{1.440801in}}%
\pgfpathlineto{\pgfqpoint{2.194047in}{1.437852in}}%
\pgfpathmoveto{\pgfqpoint{2.189506in}{1.440801in}}%
\pgfpathlineto{\pgfqpoint{2.189506in}{1.440801in}}%
\pgfpathlineto{\pgfqpoint{2.189506in}{1.443750in}}%
\pgfpathlineto{\pgfqpoint{2.194047in}{1.443750in}}%
\pgfpathlineto{\pgfqpoint{2.194047in}{1.440801in}}%
\pgfpathmoveto{\pgfqpoint{2.194047in}{1.437852in}}%
\pgfpathlineto{\pgfqpoint{2.194047in}{1.437852in}}%
\pgfpathlineto{\pgfqpoint{2.194047in}{1.440801in}}%
\pgfpathlineto{\pgfqpoint{2.198588in}{1.440801in}}%
\pgfpathlineto{\pgfqpoint{2.198588in}{1.437852in}}%
\pgfpathmoveto{\pgfqpoint{2.194047in}{1.440801in}}%
\pgfpathlineto{\pgfqpoint{2.194047in}{1.440801in}}%
\pgfpathlineto{\pgfqpoint{2.194047in}{1.443750in}}%
\pgfpathlineto{\pgfqpoint{2.198588in}{1.443750in}}%
\pgfpathlineto{\pgfqpoint{2.198588in}{1.440801in}}%
\pgfpathmoveto{\pgfqpoint{2.198588in}{1.440801in}}%
\pgfpathlineto{\pgfqpoint{2.198588in}{1.440801in}}%
\pgfpathlineto{\pgfqpoint{2.198588in}{1.443750in}}%
\pgfpathlineto{\pgfqpoint{2.203129in}{1.443750in}}%
\pgfpathlineto{\pgfqpoint{2.203129in}{1.440801in}}%
\pgfpathmoveto{\pgfqpoint{2.194047in}{1.443750in}}%
\pgfpathlineto{\pgfqpoint{2.194047in}{1.443750in}}%
\pgfpathlineto{\pgfqpoint{2.194047in}{1.446699in}}%
\pgfpathlineto{\pgfqpoint{2.198588in}{1.446699in}}%
\pgfpathlineto{\pgfqpoint{2.198588in}{1.443750in}}%
\pgfpathmoveto{\pgfqpoint{2.194047in}{1.446699in}}%
\pgfpathlineto{\pgfqpoint{2.194047in}{1.446699in}}%
\pgfpathlineto{\pgfqpoint{2.194047in}{1.449649in}}%
\pgfpathlineto{\pgfqpoint{2.198588in}{1.449649in}}%
\pgfpathlineto{\pgfqpoint{2.198588in}{1.446699in}}%
\pgfpathmoveto{\pgfqpoint{2.198588in}{1.443750in}}%
\pgfpathlineto{\pgfqpoint{2.198588in}{1.443750in}}%
\pgfpathlineto{\pgfqpoint{2.198588in}{1.446699in}}%
\pgfpathlineto{\pgfqpoint{2.203129in}{1.446699in}}%
\pgfpathlineto{\pgfqpoint{2.203129in}{1.443750in}}%
\pgfpathmoveto{\pgfqpoint{2.198588in}{1.446699in}}%
\pgfpathlineto{\pgfqpoint{2.198588in}{1.446699in}}%
\pgfpathlineto{\pgfqpoint{2.198588in}{1.449649in}}%
\pgfpathlineto{\pgfqpoint{2.203129in}{1.449649in}}%
\pgfpathlineto{\pgfqpoint{2.203129in}{1.446699in}}%
\pgfpathmoveto{\pgfqpoint{2.203129in}{1.443750in}}%
\pgfpathlineto{\pgfqpoint{2.203129in}{1.443750in}}%
\pgfpathlineto{\pgfqpoint{2.203129in}{1.446699in}}%
\pgfpathlineto{\pgfqpoint{2.207670in}{1.446699in}}%
\pgfpathlineto{\pgfqpoint{2.207670in}{1.443750in}}%
\pgfpathmoveto{\pgfqpoint{2.203129in}{1.446699in}}%
\pgfpathlineto{\pgfqpoint{2.203129in}{1.446699in}}%
\pgfpathlineto{\pgfqpoint{2.203129in}{1.449649in}}%
\pgfpathlineto{\pgfqpoint{2.207670in}{1.449649in}}%
\pgfpathlineto{\pgfqpoint{2.207670in}{1.446699in}}%
\pgfpathmoveto{\pgfqpoint{2.207670in}{1.446699in}}%
\pgfpathlineto{\pgfqpoint{2.207670in}{1.446699in}}%
\pgfpathlineto{\pgfqpoint{2.207670in}{1.449649in}}%
\pgfpathlineto{\pgfqpoint{2.212211in}{1.449649in}}%
\pgfpathlineto{\pgfqpoint{2.212211in}{1.446699in}}%
\pgfpathmoveto{\pgfqpoint{2.203129in}{1.449649in}}%
\pgfpathlineto{\pgfqpoint{2.203129in}{1.449649in}}%
\pgfpathlineto{\pgfqpoint{2.203129in}{1.452598in}}%
\pgfpathlineto{\pgfqpoint{2.207670in}{1.452598in}}%
\pgfpathlineto{\pgfqpoint{2.207670in}{1.449649in}}%
\pgfpathmoveto{\pgfqpoint{2.203129in}{1.452598in}}%
\pgfpathlineto{\pgfqpoint{2.203129in}{1.452598in}}%
\pgfpathlineto{\pgfqpoint{2.203129in}{1.455547in}}%
\pgfpathlineto{\pgfqpoint{2.207670in}{1.455547in}}%
\pgfpathlineto{\pgfqpoint{2.207670in}{1.452598in}}%
\pgfpathmoveto{\pgfqpoint{2.207670in}{1.449649in}}%
\pgfpathlineto{\pgfqpoint{2.207670in}{1.449649in}}%
\pgfpathlineto{\pgfqpoint{2.207670in}{1.452598in}}%
\pgfpathlineto{\pgfqpoint{2.212211in}{1.452598in}}%
\pgfpathlineto{\pgfqpoint{2.212211in}{1.449649in}}%
\pgfpathmoveto{\pgfqpoint{2.207670in}{1.452598in}}%
\pgfpathlineto{\pgfqpoint{2.207670in}{1.452598in}}%
\pgfpathlineto{\pgfqpoint{2.207670in}{1.455547in}}%
\pgfpathlineto{\pgfqpoint{2.212211in}{1.455547in}}%
\pgfpathlineto{\pgfqpoint{2.212211in}{1.452598in}}%
\pgfpathmoveto{\pgfqpoint{2.212211in}{1.449649in}}%
\pgfpathlineto{\pgfqpoint{2.212211in}{1.449649in}}%
\pgfpathlineto{\pgfqpoint{2.212211in}{1.452598in}}%
\pgfpathlineto{\pgfqpoint{2.216752in}{1.452598in}}%
\pgfpathlineto{\pgfqpoint{2.216752in}{1.449649in}}%
\pgfpathmoveto{\pgfqpoint{2.212211in}{1.452598in}}%
\pgfpathlineto{\pgfqpoint{2.212211in}{1.452598in}}%
\pgfpathlineto{\pgfqpoint{2.212211in}{1.455547in}}%
\pgfpathlineto{\pgfqpoint{2.216752in}{1.455547in}}%
\pgfpathlineto{\pgfqpoint{2.216752in}{1.452598in}}%
\pgfpathmoveto{\pgfqpoint{2.216752in}{1.452598in}}%
\pgfpathlineto{\pgfqpoint{2.216752in}{1.452598in}}%
\pgfpathlineto{\pgfqpoint{2.216752in}{1.455547in}}%
\pgfpathlineto{\pgfqpoint{2.221293in}{1.455547in}}%
\pgfpathlineto{\pgfqpoint{2.221293in}{1.452598in}}%
\pgfpathmoveto{\pgfqpoint{2.212211in}{1.455547in}}%
\pgfpathlineto{\pgfqpoint{2.212211in}{1.455547in}}%
\pgfpathlineto{\pgfqpoint{2.212211in}{1.458496in}}%
\pgfpathlineto{\pgfqpoint{2.216752in}{1.458496in}}%
\pgfpathlineto{\pgfqpoint{2.216752in}{1.455547in}}%
\pgfpathmoveto{\pgfqpoint{2.212211in}{1.458496in}}%
\pgfpathlineto{\pgfqpoint{2.212211in}{1.458496in}}%
\pgfpathlineto{\pgfqpoint{2.212211in}{1.461446in}}%
\pgfpathlineto{\pgfqpoint{2.216752in}{1.461446in}}%
\pgfpathlineto{\pgfqpoint{2.216752in}{1.458496in}}%
\pgfpathmoveto{\pgfqpoint{2.216752in}{1.455547in}}%
\pgfpathlineto{\pgfqpoint{2.216752in}{1.455547in}}%
\pgfpathlineto{\pgfqpoint{2.216752in}{1.458496in}}%
\pgfpathlineto{\pgfqpoint{2.221293in}{1.458496in}}%
\pgfpathlineto{\pgfqpoint{2.221293in}{1.455547in}}%
\pgfpathmoveto{\pgfqpoint{2.216752in}{1.458496in}}%
\pgfpathlineto{\pgfqpoint{2.216752in}{1.458496in}}%
\pgfpathlineto{\pgfqpoint{2.216752in}{1.461446in}}%
\pgfpathlineto{\pgfqpoint{2.221293in}{1.461446in}}%
\pgfpathlineto{\pgfqpoint{2.221293in}{1.458496in}}%
\pgfpathmoveto{\pgfqpoint{2.221293in}{1.455547in}}%
\pgfpathlineto{\pgfqpoint{2.221293in}{1.455547in}}%
\pgfpathlineto{\pgfqpoint{2.221293in}{1.458496in}}%
\pgfpathlineto{\pgfqpoint{2.225834in}{1.458496in}}%
\pgfpathlineto{\pgfqpoint{2.225834in}{1.455547in}}%
\pgfpathmoveto{\pgfqpoint{2.221293in}{1.458496in}}%
\pgfpathlineto{\pgfqpoint{2.221293in}{1.458496in}}%
\pgfpathlineto{\pgfqpoint{2.221293in}{1.461446in}}%
\pgfpathlineto{\pgfqpoint{2.225834in}{1.461446in}}%
\pgfpathlineto{\pgfqpoint{2.225834in}{1.458496in}}%
\pgfpathmoveto{\pgfqpoint{2.225834in}{1.458496in}}%
\pgfpathlineto{\pgfqpoint{2.225834in}{1.458496in}}%
\pgfpathlineto{\pgfqpoint{2.225834in}{1.461446in}}%
\pgfpathlineto{\pgfqpoint{2.230375in}{1.461446in}}%
\pgfpathlineto{\pgfqpoint{2.230375in}{1.458496in}}%
\pgfpathmoveto{\pgfqpoint{2.221293in}{1.461446in}}%
\pgfpathlineto{\pgfqpoint{2.221293in}{1.461446in}}%
\pgfpathlineto{\pgfqpoint{2.221293in}{1.464395in}}%
\pgfpathlineto{\pgfqpoint{2.225834in}{1.464395in}}%
\pgfpathlineto{\pgfqpoint{2.225834in}{1.461446in}}%
\pgfpathmoveto{\pgfqpoint{2.221293in}{1.464395in}}%
\pgfpathlineto{\pgfqpoint{2.221293in}{1.464395in}}%
\pgfpathlineto{\pgfqpoint{2.221293in}{1.467344in}}%
\pgfpathlineto{\pgfqpoint{2.225834in}{1.467344in}}%
\pgfpathlineto{\pgfqpoint{2.225834in}{1.464395in}}%
\pgfpathmoveto{\pgfqpoint{2.225834in}{1.461446in}}%
\pgfpathlineto{\pgfqpoint{2.225834in}{1.461446in}}%
\pgfpathlineto{\pgfqpoint{2.225834in}{1.464395in}}%
\pgfpathlineto{\pgfqpoint{2.230375in}{1.464395in}}%
\pgfpathlineto{\pgfqpoint{2.230375in}{1.461446in}}%
\pgfpathmoveto{\pgfqpoint{2.225834in}{1.464395in}}%
\pgfpathlineto{\pgfqpoint{2.225834in}{1.464395in}}%
\pgfpathlineto{\pgfqpoint{2.225834in}{1.467344in}}%
\pgfpathlineto{\pgfqpoint{2.230375in}{1.467344in}}%
\pgfpathlineto{\pgfqpoint{2.230375in}{1.464395in}}%
\pgfpathmoveto{\pgfqpoint{2.230375in}{1.461446in}}%
\pgfpathlineto{\pgfqpoint{2.230375in}{1.461446in}}%
\pgfpathlineto{\pgfqpoint{2.230375in}{1.464395in}}%
\pgfpathlineto{\pgfqpoint{2.234916in}{1.464395in}}%
\pgfpathlineto{\pgfqpoint{2.234916in}{1.461446in}}%
\pgfpathmoveto{\pgfqpoint{2.230375in}{1.464395in}}%
\pgfpathlineto{\pgfqpoint{2.230375in}{1.464395in}}%
\pgfpathlineto{\pgfqpoint{2.230375in}{1.467344in}}%
\pgfpathlineto{\pgfqpoint{2.234916in}{1.467344in}}%
\pgfpathlineto{\pgfqpoint{2.234916in}{1.464395in}}%
\pgfpathmoveto{\pgfqpoint{2.234916in}{1.464395in}}%
\pgfpathlineto{\pgfqpoint{2.234916in}{1.464395in}}%
\pgfpathlineto{\pgfqpoint{2.234916in}{1.467344in}}%
\pgfpathlineto{\pgfqpoint{2.239457in}{1.467344in}}%
\pgfpathlineto{\pgfqpoint{2.239457in}{1.464395in}}%
\pgfpathmoveto{\pgfqpoint{2.230375in}{1.467344in}}%
\pgfpathlineto{\pgfqpoint{2.230375in}{1.467344in}}%
\pgfpathlineto{\pgfqpoint{2.230375in}{1.470293in}}%
\pgfpathlineto{\pgfqpoint{2.234916in}{1.470293in}}%
\pgfpathlineto{\pgfqpoint{2.234916in}{1.467344in}}%
\pgfpathmoveto{\pgfqpoint{2.230375in}{1.470293in}}%
\pgfpathlineto{\pgfqpoint{2.230375in}{1.470293in}}%
\pgfpathlineto{\pgfqpoint{2.230375in}{1.473243in}}%
\pgfpathlineto{\pgfqpoint{2.234916in}{1.473243in}}%
\pgfpathlineto{\pgfqpoint{2.234916in}{1.470293in}}%
\pgfpathmoveto{\pgfqpoint{2.234916in}{1.467344in}}%
\pgfpathlineto{\pgfqpoint{2.234916in}{1.467344in}}%
\pgfpathlineto{\pgfqpoint{2.234916in}{1.470293in}}%
\pgfpathlineto{\pgfqpoint{2.239457in}{1.470293in}}%
\pgfpathlineto{\pgfqpoint{2.239457in}{1.467344in}}%
\pgfpathmoveto{\pgfqpoint{2.234916in}{1.470293in}}%
\pgfpathlineto{\pgfqpoint{2.234916in}{1.470293in}}%
\pgfpathlineto{\pgfqpoint{2.234916in}{1.473243in}}%
\pgfpathlineto{\pgfqpoint{2.239457in}{1.473243in}}%
\pgfpathlineto{\pgfqpoint{2.239457in}{1.470293in}}%
\pgfpathmoveto{\pgfqpoint{2.239457in}{1.467344in}}%
\pgfpathlineto{\pgfqpoint{2.239457in}{1.467344in}}%
\pgfpathlineto{\pgfqpoint{2.239457in}{1.470293in}}%
\pgfpathlineto{\pgfqpoint{2.243998in}{1.470293in}}%
\pgfpathlineto{\pgfqpoint{2.243998in}{1.467344in}}%
\pgfpathmoveto{\pgfqpoint{2.239457in}{1.470293in}}%
\pgfpathlineto{\pgfqpoint{2.239457in}{1.470293in}}%
\pgfpathlineto{\pgfqpoint{2.239457in}{1.473243in}}%
\pgfpathlineto{\pgfqpoint{2.243998in}{1.473243in}}%
\pgfpathlineto{\pgfqpoint{2.243998in}{1.470293in}}%
\pgfpathmoveto{\pgfqpoint{2.243998in}{1.470293in}}%
\pgfpathlineto{\pgfqpoint{2.243998in}{1.470293in}}%
\pgfpathlineto{\pgfqpoint{2.243998in}{1.473243in}}%
\pgfpathlineto{\pgfqpoint{2.248539in}{1.473243in}}%
\pgfpathlineto{\pgfqpoint{2.248539in}{1.470293in}}%
\pgfpathmoveto{\pgfqpoint{2.239457in}{1.473243in}}%
\pgfpathlineto{\pgfqpoint{2.239457in}{1.473243in}}%
\pgfpathlineto{\pgfqpoint{2.239457in}{1.476192in}}%
\pgfpathlineto{\pgfqpoint{2.243998in}{1.476192in}}%
\pgfpathlineto{\pgfqpoint{2.243998in}{1.473243in}}%
\pgfpathmoveto{\pgfqpoint{2.239457in}{1.476192in}}%
\pgfpathlineto{\pgfqpoint{2.239457in}{1.476192in}}%
\pgfpathlineto{\pgfqpoint{2.239457in}{1.479141in}}%
\pgfpathlineto{\pgfqpoint{2.243998in}{1.479141in}}%
\pgfpathlineto{\pgfqpoint{2.243998in}{1.476192in}}%
\pgfpathmoveto{\pgfqpoint{2.243998in}{1.473243in}}%
\pgfpathlineto{\pgfqpoint{2.243998in}{1.473243in}}%
\pgfpathlineto{\pgfqpoint{2.243998in}{1.476192in}}%
\pgfpathlineto{\pgfqpoint{2.248539in}{1.476192in}}%
\pgfpathlineto{\pgfqpoint{2.248539in}{1.473243in}}%
\pgfpathmoveto{\pgfqpoint{2.243998in}{1.476192in}}%
\pgfpathlineto{\pgfqpoint{2.243998in}{1.476192in}}%
\pgfpathlineto{\pgfqpoint{2.243998in}{1.479141in}}%
\pgfpathlineto{\pgfqpoint{2.248539in}{1.479141in}}%
\pgfpathlineto{\pgfqpoint{2.248539in}{1.476192in}}%
\pgfpathmoveto{\pgfqpoint{2.248539in}{1.473243in}}%
\pgfpathlineto{\pgfqpoint{2.248539in}{1.473243in}}%
\pgfpathlineto{\pgfqpoint{2.248539in}{1.476192in}}%
\pgfpathlineto{\pgfqpoint{2.253080in}{1.476192in}}%
\pgfpathlineto{\pgfqpoint{2.253080in}{1.473243in}}%
\pgfpathmoveto{\pgfqpoint{2.248539in}{1.476192in}}%
\pgfpathlineto{\pgfqpoint{2.248539in}{1.476192in}}%
\pgfpathlineto{\pgfqpoint{2.248539in}{1.479141in}}%
\pgfpathlineto{\pgfqpoint{2.253080in}{1.479141in}}%
\pgfpathlineto{\pgfqpoint{2.253080in}{1.476192in}}%
\pgfpathmoveto{\pgfqpoint{2.253080in}{1.476192in}}%
\pgfpathlineto{\pgfqpoint{2.253080in}{1.476192in}}%
\pgfpathlineto{\pgfqpoint{2.253080in}{1.479141in}}%
\pgfpathlineto{\pgfqpoint{2.257621in}{1.479141in}}%
\pgfpathlineto{\pgfqpoint{2.257621in}{1.476192in}}%
\pgfpathmoveto{\pgfqpoint{2.248539in}{1.479141in}}%
\pgfpathlineto{\pgfqpoint{2.248539in}{1.479141in}}%
\pgfpathlineto{\pgfqpoint{2.248539in}{1.482091in}}%
\pgfpathlineto{\pgfqpoint{2.253080in}{1.482091in}}%
\pgfpathlineto{\pgfqpoint{2.253080in}{1.479141in}}%
\pgfpathmoveto{\pgfqpoint{2.248539in}{1.482091in}}%
\pgfpathlineto{\pgfqpoint{2.248539in}{1.482091in}}%
\pgfpathlineto{\pgfqpoint{2.248539in}{1.485040in}}%
\pgfpathlineto{\pgfqpoint{2.253080in}{1.485040in}}%
\pgfpathlineto{\pgfqpoint{2.253080in}{1.482091in}}%
\pgfpathmoveto{\pgfqpoint{2.253080in}{1.479141in}}%
\pgfpathlineto{\pgfqpoint{2.253080in}{1.479141in}}%
\pgfpathlineto{\pgfqpoint{2.253080in}{1.482091in}}%
\pgfpathlineto{\pgfqpoint{2.257621in}{1.482091in}}%
\pgfpathlineto{\pgfqpoint{2.257621in}{1.479141in}}%
\pgfpathmoveto{\pgfqpoint{2.253080in}{1.482091in}}%
\pgfpathlineto{\pgfqpoint{2.253080in}{1.482091in}}%
\pgfpathlineto{\pgfqpoint{2.253080in}{1.485040in}}%
\pgfpathlineto{\pgfqpoint{2.257621in}{1.485040in}}%
\pgfpathlineto{\pgfqpoint{2.257621in}{1.482091in}}%
\pgfpathmoveto{\pgfqpoint{2.257621in}{1.479141in}}%
\pgfpathlineto{\pgfqpoint{2.257621in}{1.479141in}}%
\pgfpathlineto{\pgfqpoint{2.257621in}{1.482091in}}%
\pgfpathlineto{\pgfqpoint{2.262162in}{1.482091in}}%
\pgfpathlineto{\pgfqpoint{2.262162in}{1.479141in}}%
\pgfpathmoveto{\pgfqpoint{2.257621in}{1.482091in}}%
\pgfpathlineto{\pgfqpoint{2.257621in}{1.482091in}}%
\pgfpathlineto{\pgfqpoint{2.257621in}{1.485040in}}%
\pgfpathlineto{\pgfqpoint{2.262162in}{1.485040in}}%
\pgfpathlineto{\pgfqpoint{2.262162in}{1.482091in}}%
\pgfpathmoveto{\pgfqpoint{2.262162in}{1.482091in}}%
\pgfpathlineto{\pgfqpoint{2.262162in}{1.482091in}}%
\pgfpathlineto{\pgfqpoint{2.262162in}{1.485040in}}%
\pgfpathlineto{\pgfqpoint{2.266703in}{1.485040in}}%
\pgfpathlineto{\pgfqpoint{2.266703in}{1.482091in}}%
\pgfpathmoveto{\pgfqpoint{2.257621in}{1.485040in}}%
\pgfpathlineto{\pgfqpoint{2.257621in}{1.485040in}}%
\pgfpathlineto{\pgfqpoint{2.257621in}{1.487989in}}%
\pgfpathlineto{\pgfqpoint{2.262162in}{1.487989in}}%
\pgfpathlineto{\pgfqpoint{2.262162in}{1.485040in}}%
\pgfpathmoveto{\pgfqpoint{2.257621in}{1.487989in}}%
\pgfpathlineto{\pgfqpoint{2.257621in}{1.487989in}}%
\pgfpathlineto{\pgfqpoint{2.257621in}{1.490938in}}%
\pgfpathlineto{\pgfqpoint{2.262162in}{1.490938in}}%
\pgfpathlineto{\pgfqpoint{2.262162in}{1.487989in}}%
\pgfpathmoveto{\pgfqpoint{2.262162in}{1.485040in}}%
\pgfpathlineto{\pgfqpoint{2.262162in}{1.485040in}}%
\pgfpathlineto{\pgfqpoint{2.262162in}{1.487989in}}%
\pgfpathlineto{\pgfqpoint{2.266703in}{1.487989in}}%
\pgfpathlineto{\pgfqpoint{2.266703in}{1.485040in}}%
\pgfpathmoveto{\pgfqpoint{2.262162in}{1.487989in}}%
\pgfpathlineto{\pgfqpoint{2.262162in}{1.487989in}}%
\pgfpathlineto{\pgfqpoint{2.262162in}{1.490938in}}%
\pgfpathlineto{\pgfqpoint{2.266703in}{1.490938in}}%
\pgfpathlineto{\pgfqpoint{2.266703in}{1.487989in}}%
\pgfpathmoveto{\pgfqpoint{2.266703in}{1.485040in}}%
\pgfpathlineto{\pgfqpoint{2.266703in}{1.485040in}}%
\pgfpathlineto{\pgfqpoint{2.266703in}{1.487989in}}%
\pgfpathlineto{\pgfqpoint{2.271244in}{1.487989in}}%
\pgfpathlineto{\pgfqpoint{2.271244in}{1.485040in}}%
\pgfpathmoveto{\pgfqpoint{2.266703in}{1.487989in}}%
\pgfpathlineto{\pgfqpoint{2.266703in}{1.487989in}}%
\pgfpathlineto{\pgfqpoint{2.266703in}{1.490938in}}%
\pgfpathlineto{\pgfqpoint{2.271244in}{1.490938in}}%
\pgfpathlineto{\pgfqpoint{2.271244in}{1.487989in}}%
\pgfpathmoveto{\pgfqpoint{2.271244in}{1.487989in}}%
\pgfpathlineto{\pgfqpoint{2.271244in}{1.487989in}}%
\pgfpathlineto{\pgfqpoint{2.271244in}{1.490938in}}%
\pgfpathlineto{\pgfqpoint{2.275785in}{1.490938in}}%
\pgfpathlineto{\pgfqpoint{2.275785in}{1.487989in}}%
\pgfpathmoveto{\pgfqpoint{2.266703in}{1.490938in}}%
\pgfpathlineto{\pgfqpoint{2.266703in}{1.490938in}}%
\pgfpathlineto{\pgfqpoint{2.266703in}{1.493888in}}%
\pgfpathlineto{\pgfqpoint{2.271244in}{1.493888in}}%
\pgfpathlineto{\pgfqpoint{2.271244in}{1.490938in}}%
\pgfpathmoveto{\pgfqpoint{2.266703in}{1.493888in}}%
\pgfpathlineto{\pgfqpoint{2.266703in}{1.493888in}}%
\pgfpathlineto{\pgfqpoint{2.266703in}{1.496837in}}%
\pgfpathlineto{\pgfqpoint{2.271244in}{1.496837in}}%
\pgfpathlineto{\pgfqpoint{2.271244in}{1.493888in}}%
\pgfpathmoveto{\pgfqpoint{2.271244in}{1.490938in}}%
\pgfpathlineto{\pgfqpoint{2.271244in}{1.490938in}}%
\pgfpathlineto{\pgfqpoint{2.271244in}{1.493888in}}%
\pgfpathlineto{\pgfqpoint{2.275785in}{1.493888in}}%
\pgfpathlineto{\pgfqpoint{2.275785in}{1.490938in}}%
\pgfpathmoveto{\pgfqpoint{2.271244in}{1.493888in}}%
\pgfpathlineto{\pgfqpoint{2.271244in}{1.493888in}}%
\pgfpathlineto{\pgfqpoint{2.271244in}{1.496837in}}%
\pgfpathlineto{\pgfqpoint{2.275785in}{1.496837in}}%
\pgfpathlineto{\pgfqpoint{2.275785in}{1.493888in}}%
\pgfpathmoveto{\pgfqpoint{2.275785in}{1.490938in}}%
\pgfpathlineto{\pgfqpoint{2.275785in}{1.490938in}}%
\pgfpathlineto{\pgfqpoint{2.275785in}{1.493888in}}%
\pgfpathlineto{\pgfqpoint{2.280326in}{1.493888in}}%
\pgfpathlineto{\pgfqpoint{2.280326in}{1.490938in}}%
\pgfpathmoveto{\pgfqpoint{2.275785in}{1.493888in}}%
\pgfpathlineto{\pgfqpoint{2.275785in}{1.493888in}}%
\pgfpathlineto{\pgfqpoint{2.275785in}{1.496837in}}%
\pgfpathlineto{\pgfqpoint{2.280326in}{1.496837in}}%
\pgfpathlineto{\pgfqpoint{2.280326in}{1.493888in}}%
\pgfpathmoveto{\pgfqpoint{2.280326in}{1.493888in}}%
\pgfpathlineto{\pgfqpoint{2.280326in}{1.493888in}}%
\pgfpathlineto{\pgfqpoint{2.280326in}{1.496837in}}%
\pgfpathlineto{\pgfqpoint{2.284867in}{1.496837in}}%
\pgfpathlineto{\pgfqpoint{2.284867in}{1.493888in}}%
\pgfpathmoveto{\pgfqpoint{2.275785in}{1.496837in}}%
\pgfpathlineto{\pgfqpoint{2.275785in}{1.496837in}}%
\pgfpathlineto{\pgfqpoint{2.275785in}{1.499786in}}%
\pgfpathlineto{\pgfqpoint{2.280326in}{1.499786in}}%
\pgfpathlineto{\pgfqpoint{2.280326in}{1.496837in}}%
\pgfpathmoveto{\pgfqpoint{2.275785in}{1.499786in}}%
\pgfpathlineto{\pgfqpoint{2.275785in}{1.499786in}}%
\pgfpathlineto{\pgfqpoint{2.275785in}{1.502735in}}%
\pgfpathlineto{\pgfqpoint{2.280326in}{1.502735in}}%
\pgfpathlineto{\pgfqpoint{2.280326in}{1.499786in}}%
\pgfpathmoveto{\pgfqpoint{2.280326in}{1.496837in}}%
\pgfpathlineto{\pgfqpoint{2.280326in}{1.496837in}}%
\pgfpathlineto{\pgfqpoint{2.280326in}{1.499786in}}%
\pgfpathlineto{\pgfqpoint{2.284867in}{1.499786in}}%
\pgfpathlineto{\pgfqpoint{2.284867in}{1.496837in}}%
\pgfpathmoveto{\pgfqpoint{2.280326in}{1.499786in}}%
\pgfpathlineto{\pgfqpoint{2.280326in}{1.499786in}}%
\pgfpathlineto{\pgfqpoint{2.280326in}{1.502735in}}%
\pgfpathlineto{\pgfqpoint{2.284867in}{1.502735in}}%
\pgfpathlineto{\pgfqpoint{2.284867in}{1.499786in}}%
\pgfpathmoveto{\pgfqpoint{2.284867in}{1.496837in}}%
\pgfpathlineto{\pgfqpoint{2.284867in}{1.496837in}}%
\pgfpathlineto{\pgfqpoint{2.284867in}{1.499786in}}%
\pgfpathlineto{\pgfqpoint{2.289408in}{1.499786in}}%
\pgfpathlineto{\pgfqpoint{2.289408in}{1.496837in}}%
\pgfpathmoveto{\pgfqpoint{2.284867in}{1.499786in}}%
\pgfpathlineto{\pgfqpoint{2.284867in}{1.499786in}}%
\pgfpathlineto{\pgfqpoint{2.284867in}{1.502735in}}%
\pgfpathlineto{\pgfqpoint{2.289408in}{1.502735in}}%
\pgfpathlineto{\pgfqpoint{2.289408in}{1.499786in}}%
\pgfpathmoveto{\pgfqpoint{2.289408in}{1.499786in}}%
\pgfpathlineto{\pgfqpoint{2.289408in}{1.499786in}}%
\pgfpathlineto{\pgfqpoint{2.289408in}{1.502735in}}%
\pgfpathlineto{\pgfqpoint{2.293949in}{1.502735in}}%
\pgfpathlineto{\pgfqpoint{2.293949in}{1.499786in}}%
\pgfpathmoveto{\pgfqpoint{2.284867in}{1.502735in}}%
\pgfpathlineto{\pgfqpoint{2.284867in}{1.502735in}}%
\pgfpathlineto{\pgfqpoint{2.284867in}{1.505685in}}%
\pgfpathlineto{\pgfqpoint{2.289408in}{1.505685in}}%
\pgfpathlineto{\pgfqpoint{2.289408in}{1.502735in}}%
\pgfpathmoveto{\pgfqpoint{2.284867in}{1.505685in}}%
\pgfpathlineto{\pgfqpoint{2.284867in}{1.505685in}}%
\pgfpathlineto{\pgfqpoint{2.284867in}{1.508634in}}%
\pgfpathlineto{\pgfqpoint{2.289408in}{1.508634in}}%
\pgfpathlineto{\pgfqpoint{2.289408in}{1.505685in}}%
\pgfpathmoveto{\pgfqpoint{2.289408in}{1.502735in}}%
\pgfpathlineto{\pgfqpoint{2.289408in}{1.502735in}}%
\pgfpathlineto{\pgfqpoint{2.289408in}{1.505685in}}%
\pgfpathlineto{\pgfqpoint{2.293949in}{1.505685in}}%
\pgfpathlineto{\pgfqpoint{2.293949in}{1.502735in}}%
\pgfpathmoveto{\pgfqpoint{2.289408in}{1.505685in}}%
\pgfpathlineto{\pgfqpoint{2.289408in}{1.505685in}}%
\pgfpathlineto{\pgfqpoint{2.289408in}{1.508634in}}%
\pgfpathlineto{\pgfqpoint{2.293949in}{1.508634in}}%
\pgfpathlineto{\pgfqpoint{2.293949in}{1.505685in}}%
\pgfpathmoveto{\pgfqpoint{2.293949in}{1.502735in}}%
\pgfpathlineto{\pgfqpoint{2.293949in}{1.502735in}}%
\pgfpathlineto{\pgfqpoint{2.293949in}{1.505685in}}%
\pgfpathlineto{\pgfqpoint{2.298490in}{1.505685in}}%
\pgfpathlineto{\pgfqpoint{2.298490in}{1.502735in}}%
\pgfpathmoveto{\pgfqpoint{2.293949in}{1.505685in}}%
\pgfpathlineto{\pgfqpoint{2.293949in}{1.505685in}}%
\pgfpathlineto{\pgfqpoint{2.293949in}{1.508634in}}%
\pgfpathlineto{\pgfqpoint{2.298490in}{1.508634in}}%
\pgfpathlineto{\pgfqpoint{2.298490in}{1.505685in}}%
\pgfpathmoveto{\pgfqpoint{2.298490in}{1.505685in}}%
\pgfpathlineto{\pgfqpoint{2.298490in}{1.505685in}}%
\pgfpathlineto{\pgfqpoint{2.298490in}{1.508634in}}%
\pgfpathlineto{\pgfqpoint{2.303031in}{1.508634in}}%
\pgfpathlineto{\pgfqpoint{2.303031in}{1.505685in}}%
\pgfpathmoveto{\pgfqpoint{2.293949in}{1.508634in}}%
\pgfpathlineto{\pgfqpoint{2.293949in}{1.508634in}}%
\pgfpathlineto{\pgfqpoint{2.293949in}{1.511583in}}%
\pgfpathlineto{\pgfqpoint{2.298490in}{1.511583in}}%
\pgfpathlineto{\pgfqpoint{2.298490in}{1.508634in}}%
\pgfpathmoveto{\pgfqpoint{2.293949in}{1.511583in}}%
\pgfpathlineto{\pgfqpoint{2.293949in}{1.511583in}}%
\pgfpathlineto{\pgfqpoint{2.293949in}{1.514533in}}%
\pgfpathlineto{\pgfqpoint{2.298490in}{1.514533in}}%
\pgfpathlineto{\pgfqpoint{2.298490in}{1.511583in}}%
\pgfpathmoveto{\pgfqpoint{2.298490in}{1.508634in}}%
\pgfpathlineto{\pgfqpoint{2.298490in}{1.508634in}}%
\pgfpathlineto{\pgfqpoint{2.298490in}{1.511583in}}%
\pgfpathlineto{\pgfqpoint{2.303031in}{1.511583in}}%
\pgfpathlineto{\pgfqpoint{2.303031in}{1.508634in}}%
\pgfpathmoveto{\pgfqpoint{2.298490in}{1.511583in}}%
\pgfpathlineto{\pgfqpoint{2.298490in}{1.511583in}}%
\pgfpathlineto{\pgfqpoint{2.298490in}{1.514533in}}%
\pgfpathlineto{\pgfqpoint{2.303031in}{1.514533in}}%
\pgfpathlineto{\pgfqpoint{2.303031in}{1.511583in}}%
\pgfpathmoveto{\pgfqpoint{2.303031in}{1.508634in}}%
\pgfpathlineto{\pgfqpoint{2.303031in}{1.508634in}}%
\pgfpathlineto{\pgfqpoint{2.303031in}{1.511583in}}%
\pgfpathlineto{\pgfqpoint{2.307572in}{1.511583in}}%
\pgfpathlineto{\pgfqpoint{2.307572in}{1.508634in}}%
\pgfpathmoveto{\pgfqpoint{2.303031in}{1.511583in}}%
\pgfpathlineto{\pgfqpoint{2.303031in}{1.511583in}}%
\pgfpathlineto{\pgfqpoint{2.303031in}{1.514533in}}%
\pgfpathlineto{\pgfqpoint{2.307572in}{1.514533in}}%
\pgfpathlineto{\pgfqpoint{2.307572in}{1.511583in}}%
\pgfpathmoveto{\pgfqpoint{2.307572in}{1.511583in}}%
\pgfpathlineto{\pgfqpoint{2.307572in}{1.511583in}}%
\pgfpathlineto{\pgfqpoint{2.307572in}{1.514533in}}%
\pgfpathlineto{\pgfqpoint{2.312113in}{1.514533in}}%
\pgfpathlineto{\pgfqpoint{2.312113in}{1.511583in}}%
\pgfpathmoveto{\pgfqpoint{2.303031in}{1.514533in}}%
\pgfpathlineto{\pgfqpoint{2.303031in}{1.514533in}}%
\pgfpathlineto{\pgfqpoint{2.303031in}{1.517482in}}%
\pgfpathlineto{\pgfqpoint{2.307572in}{1.517482in}}%
\pgfpathlineto{\pgfqpoint{2.307572in}{1.514533in}}%
\pgfpathmoveto{\pgfqpoint{2.303031in}{1.517482in}}%
\pgfpathlineto{\pgfqpoint{2.303031in}{1.517482in}}%
\pgfpathlineto{\pgfqpoint{2.303031in}{1.520431in}}%
\pgfpathlineto{\pgfqpoint{2.307572in}{1.520431in}}%
\pgfpathlineto{\pgfqpoint{2.307572in}{1.517482in}}%
\pgfpathmoveto{\pgfqpoint{2.307572in}{1.514533in}}%
\pgfpathlineto{\pgfqpoint{2.307572in}{1.514533in}}%
\pgfpathlineto{\pgfqpoint{2.307572in}{1.517482in}}%
\pgfpathlineto{\pgfqpoint{2.312113in}{1.517482in}}%
\pgfpathlineto{\pgfqpoint{2.312113in}{1.514533in}}%
\pgfpathmoveto{\pgfqpoint{2.307572in}{1.517482in}}%
\pgfpathlineto{\pgfqpoint{2.307572in}{1.517482in}}%
\pgfpathlineto{\pgfqpoint{2.307572in}{1.520431in}}%
\pgfpathlineto{\pgfqpoint{2.312113in}{1.520431in}}%
\pgfpathlineto{\pgfqpoint{2.312113in}{1.517482in}}%
\pgfpathmoveto{\pgfqpoint{2.312113in}{1.514533in}}%
\pgfpathlineto{\pgfqpoint{2.312113in}{1.514533in}}%
\pgfpathlineto{\pgfqpoint{2.312113in}{1.517482in}}%
\pgfpathlineto{\pgfqpoint{2.316654in}{1.517482in}}%
\pgfpathlineto{\pgfqpoint{2.316654in}{1.514533in}}%
\pgfpathmoveto{\pgfqpoint{2.312113in}{1.517482in}}%
\pgfpathlineto{\pgfqpoint{2.312113in}{1.517482in}}%
\pgfpathlineto{\pgfqpoint{2.312113in}{1.520431in}}%
\pgfpathlineto{\pgfqpoint{2.316654in}{1.520431in}}%
\pgfpathlineto{\pgfqpoint{2.316654in}{1.517482in}}%
\pgfpathmoveto{\pgfqpoint{2.316654in}{1.517482in}}%
\pgfpathlineto{\pgfqpoint{2.316654in}{1.517482in}}%
\pgfpathlineto{\pgfqpoint{2.316654in}{1.520431in}}%
\pgfpathlineto{\pgfqpoint{2.321195in}{1.520431in}}%
\pgfpathlineto{\pgfqpoint{2.321195in}{1.517482in}}%
\pgfpathmoveto{\pgfqpoint{2.312113in}{1.520431in}}%
\pgfpathlineto{\pgfqpoint{2.312113in}{1.520431in}}%
\pgfpathlineto{\pgfqpoint{2.312113in}{1.523380in}}%
\pgfpathlineto{\pgfqpoint{2.316654in}{1.523380in}}%
\pgfpathlineto{\pgfqpoint{2.316654in}{1.520431in}}%
\pgfpathmoveto{\pgfqpoint{2.312113in}{1.523380in}}%
\pgfpathlineto{\pgfqpoint{2.312113in}{1.523380in}}%
\pgfpathlineto{\pgfqpoint{2.312113in}{1.526330in}}%
\pgfpathlineto{\pgfqpoint{2.316654in}{1.526330in}}%
\pgfpathlineto{\pgfqpoint{2.316654in}{1.523380in}}%
\pgfpathmoveto{\pgfqpoint{2.316654in}{1.520431in}}%
\pgfpathlineto{\pgfqpoint{2.316654in}{1.520431in}}%
\pgfpathlineto{\pgfqpoint{2.316654in}{1.523380in}}%
\pgfpathlineto{\pgfqpoint{2.321195in}{1.523380in}}%
\pgfpathlineto{\pgfqpoint{2.321195in}{1.520431in}}%
\pgfpathmoveto{\pgfqpoint{2.316654in}{1.523380in}}%
\pgfpathlineto{\pgfqpoint{2.316654in}{1.523380in}}%
\pgfpathlineto{\pgfqpoint{2.316654in}{1.526330in}}%
\pgfpathlineto{\pgfqpoint{2.321195in}{1.526330in}}%
\pgfpathlineto{\pgfqpoint{2.321195in}{1.523380in}}%
\pgfpathmoveto{\pgfqpoint{2.321195in}{1.520431in}}%
\pgfpathlineto{\pgfqpoint{2.321195in}{1.520431in}}%
\pgfpathlineto{\pgfqpoint{2.321195in}{1.523380in}}%
\pgfpathlineto{\pgfqpoint{2.325736in}{1.523380in}}%
\pgfpathlineto{\pgfqpoint{2.325736in}{1.520431in}}%
\pgfpathmoveto{\pgfqpoint{2.321195in}{1.523380in}}%
\pgfpathlineto{\pgfqpoint{2.321195in}{1.523380in}}%
\pgfpathlineto{\pgfqpoint{2.321195in}{1.526330in}}%
\pgfpathlineto{\pgfqpoint{2.325736in}{1.526330in}}%
\pgfpathlineto{\pgfqpoint{2.325736in}{1.523380in}}%
\pgfpathmoveto{\pgfqpoint{2.325736in}{1.523380in}}%
\pgfpathlineto{\pgfqpoint{2.325736in}{1.523380in}}%
\pgfpathlineto{\pgfqpoint{2.325736in}{1.526330in}}%
\pgfpathlineto{\pgfqpoint{2.330277in}{1.526330in}}%
\pgfpathlineto{\pgfqpoint{2.330277in}{1.523380in}}%
\pgfpathmoveto{\pgfqpoint{2.321195in}{1.526330in}}%
\pgfpathlineto{\pgfqpoint{2.321195in}{1.526330in}}%
\pgfpathlineto{\pgfqpoint{2.321195in}{1.529279in}}%
\pgfpathlineto{\pgfqpoint{2.325736in}{1.529279in}}%
\pgfpathlineto{\pgfqpoint{2.325736in}{1.526330in}}%
\pgfpathmoveto{\pgfqpoint{2.321195in}{1.529279in}}%
\pgfpathlineto{\pgfqpoint{2.321195in}{1.529279in}}%
\pgfpathlineto{\pgfqpoint{2.321195in}{1.532228in}}%
\pgfpathlineto{\pgfqpoint{2.325736in}{1.532228in}}%
\pgfpathlineto{\pgfqpoint{2.325736in}{1.529279in}}%
\pgfpathmoveto{\pgfqpoint{2.325736in}{1.526330in}}%
\pgfpathlineto{\pgfqpoint{2.325736in}{1.526330in}}%
\pgfpathlineto{\pgfqpoint{2.325736in}{1.529279in}}%
\pgfpathlineto{\pgfqpoint{2.330277in}{1.529279in}}%
\pgfpathlineto{\pgfqpoint{2.330277in}{1.526330in}}%
\pgfpathmoveto{\pgfqpoint{2.325736in}{1.529279in}}%
\pgfpathlineto{\pgfqpoint{2.325736in}{1.529279in}}%
\pgfpathlineto{\pgfqpoint{2.325736in}{1.532228in}}%
\pgfpathlineto{\pgfqpoint{2.330277in}{1.532228in}}%
\pgfpathlineto{\pgfqpoint{2.330277in}{1.529279in}}%
\pgfpathmoveto{\pgfqpoint{2.330277in}{1.526330in}}%
\pgfpathlineto{\pgfqpoint{2.330277in}{1.526330in}}%
\pgfpathlineto{\pgfqpoint{2.330277in}{1.529279in}}%
\pgfpathlineto{\pgfqpoint{2.334818in}{1.529279in}}%
\pgfpathlineto{\pgfqpoint{2.334818in}{1.526330in}}%
\pgfpathmoveto{\pgfqpoint{2.330277in}{1.529279in}}%
\pgfpathlineto{\pgfqpoint{2.330277in}{1.529279in}}%
\pgfpathlineto{\pgfqpoint{2.330277in}{1.532228in}}%
\pgfpathlineto{\pgfqpoint{2.334818in}{1.532228in}}%
\pgfpathlineto{\pgfqpoint{2.334818in}{1.529279in}}%
\pgfpathmoveto{\pgfqpoint{2.334818in}{1.529279in}}%
\pgfpathlineto{\pgfqpoint{2.334818in}{1.529279in}}%
\pgfpathlineto{\pgfqpoint{2.334818in}{1.532228in}}%
\pgfpathlineto{\pgfqpoint{2.339359in}{1.532228in}}%
\pgfpathlineto{\pgfqpoint{2.339359in}{1.529279in}}%
\pgfpathmoveto{\pgfqpoint{2.330277in}{1.532228in}}%
\pgfpathlineto{\pgfqpoint{2.330277in}{1.532228in}}%
\pgfpathlineto{\pgfqpoint{2.330277in}{1.535177in}}%
\pgfpathlineto{\pgfqpoint{2.334818in}{1.535177in}}%
\pgfpathlineto{\pgfqpoint{2.334818in}{1.532228in}}%
\pgfpathmoveto{\pgfqpoint{2.330277in}{1.535177in}}%
\pgfpathlineto{\pgfqpoint{2.330277in}{1.535177in}}%
\pgfpathlineto{\pgfqpoint{2.330277in}{1.538127in}}%
\pgfpathlineto{\pgfqpoint{2.334818in}{1.538127in}}%
\pgfpathlineto{\pgfqpoint{2.334818in}{1.535177in}}%
\pgfpathmoveto{\pgfqpoint{2.334818in}{1.532228in}}%
\pgfpathlineto{\pgfqpoint{2.334818in}{1.532228in}}%
\pgfpathlineto{\pgfqpoint{2.334818in}{1.535177in}}%
\pgfpathlineto{\pgfqpoint{2.339359in}{1.535177in}}%
\pgfpathlineto{\pgfqpoint{2.339359in}{1.532228in}}%
\pgfpathmoveto{\pgfqpoint{2.334818in}{1.535177in}}%
\pgfpathlineto{\pgfqpoint{2.334818in}{1.535177in}}%
\pgfpathlineto{\pgfqpoint{2.334818in}{1.538127in}}%
\pgfpathlineto{\pgfqpoint{2.339359in}{1.538127in}}%
\pgfpathlineto{\pgfqpoint{2.339359in}{1.535177in}}%
\pgfpathmoveto{\pgfqpoint{2.339359in}{1.532228in}}%
\pgfpathlineto{\pgfqpoint{2.339359in}{1.532228in}}%
\pgfpathlineto{\pgfqpoint{2.339359in}{1.535177in}}%
\pgfpathlineto{\pgfqpoint{2.343900in}{1.535177in}}%
\pgfpathlineto{\pgfqpoint{2.343900in}{1.532228in}}%
\pgfpathmoveto{\pgfqpoint{2.339359in}{1.535177in}}%
\pgfpathlineto{\pgfqpoint{2.339359in}{1.535177in}}%
\pgfpathlineto{\pgfqpoint{2.339359in}{1.538127in}}%
\pgfpathlineto{\pgfqpoint{2.343900in}{1.538127in}}%
\pgfpathlineto{\pgfqpoint{2.343900in}{1.535177in}}%
\pgfpathmoveto{\pgfqpoint{2.343900in}{1.535177in}}%
\pgfpathlineto{\pgfqpoint{2.343900in}{1.535177in}}%
\pgfpathlineto{\pgfqpoint{2.343900in}{1.538127in}}%
\pgfpathlineto{\pgfqpoint{2.348441in}{1.538127in}}%
\pgfpathlineto{\pgfqpoint{2.348441in}{1.535177in}}%
\pgfpathmoveto{\pgfqpoint{2.339359in}{1.538127in}}%
\pgfpathlineto{\pgfqpoint{2.339359in}{1.538127in}}%
\pgfpathlineto{\pgfqpoint{2.339359in}{1.541076in}}%
\pgfpathlineto{\pgfqpoint{2.343900in}{1.541076in}}%
\pgfpathlineto{\pgfqpoint{2.343900in}{1.538127in}}%
\pgfpathmoveto{\pgfqpoint{2.339359in}{1.541076in}}%
\pgfpathlineto{\pgfqpoint{2.339359in}{1.541076in}}%
\pgfpathlineto{\pgfqpoint{2.339359in}{1.544025in}}%
\pgfpathlineto{\pgfqpoint{2.343900in}{1.544025in}}%
\pgfpathlineto{\pgfqpoint{2.343900in}{1.541076in}}%
\pgfpathmoveto{\pgfqpoint{2.343900in}{1.538127in}}%
\pgfpathlineto{\pgfqpoint{2.343900in}{1.538127in}}%
\pgfpathlineto{\pgfqpoint{2.343900in}{1.541076in}}%
\pgfpathlineto{\pgfqpoint{2.348441in}{1.541076in}}%
\pgfpathlineto{\pgfqpoint{2.348441in}{1.538127in}}%
\pgfpathmoveto{\pgfqpoint{2.343900in}{1.541076in}}%
\pgfpathlineto{\pgfqpoint{2.343900in}{1.541076in}}%
\pgfpathlineto{\pgfqpoint{2.343900in}{1.544025in}}%
\pgfpathlineto{\pgfqpoint{2.348441in}{1.544025in}}%
\pgfpathlineto{\pgfqpoint{2.348441in}{1.541076in}}%
\pgfpathmoveto{\pgfqpoint{2.348441in}{1.538127in}}%
\pgfpathlineto{\pgfqpoint{2.348441in}{1.538127in}}%
\pgfpathlineto{\pgfqpoint{2.348441in}{1.541076in}}%
\pgfpathlineto{\pgfqpoint{2.352981in}{1.541076in}}%
\pgfpathlineto{\pgfqpoint{2.352981in}{1.538127in}}%
\pgfpathmoveto{\pgfqpoint{2.348441in}{1.541076in}}%
\pgfpathlineto{\pgfqpoint{2.348441in}{1.541076in}}%
\pgfpathlineto{\pgfqpoint{2.348441in}{1.544025in}}%
\pgfpathlineto{\pgfqpoint{2.352981in}{1.544025in}}%
\pgfpathlineto{\pgfqpoint{2.352981in}{1.541076in}}%
\pgfpathmoveto{\pgfqpoint{2.352981in}{1.541076in}}%
\pgfpathlineto{\pgfqpoint{2.352981in}{1.541076in}}%
\pgfpathlineto{\pgfqpoint{2.352981in}{1.544025in}}%
\pgfpathlineto{\pgfqpoint{2.357522in}{1.544025in}}%
\pgfpathlineto{\pgfqpoint{2.357522in}{1.541076in}}%
\pgfpathmoveto{\pgfqpoint{2.348441in}{1.544025in}}%
\pgfpathlineto{\pgfqpoint{2.348441in}{1.544025in}}%
\pgfpathlineto{\pgfqpoint{2.348441in}{1.546974in}}%
\pgfpathlineto{\pgfqpoint{2.352981in}{1.546974in}}%
\pgfpathlineto{\pgfqpoint{2.352981in}{1.544025in}}%
\pgfpathmoveto{\pgfqpoint{2.348441in}{1.546974in}}%
\pgfpathlineto{\pgfqpoint{2.348441in}{1.546974in}}%
\pgfpathlineto{\pgfqpoint{2.348441in}{1.549924in}}%
\pgfpathlineto{\pgfqpoint{2.352981in}{1.549924in}}%
\pgfpathlineto{\pgfqpoint{2.352981in}{1.546974in}}%
\pgfpathmoveto{\pgfqpoint{2.352981in}{1.544025in}}%
\pgfpathlineto{\pgfqpoint{2.352981in}{1.544025in}}%
\pgfpathlineto{\pgfqpoint{2.352981in}{1.546974in}}%
\pgfpathlineto{\pgfqpoint{2.357522in}{1.546974in}}%
\pgfpathlineto{\pgfqpoint{2.357522in}{1.544025in}}%
\pgfpathmoveto{\pgfqpoint{2.352981in}{1.546974in}}%
\pgfpathlineto{\pgfqpoint{2.352981in}{1.546974in}}%
\pgfpathlineto{\pgfqpoint{2.352981in}{1.549924in}}%
\pgfpathlineto{\pgfqpoint{2.357522in}{1.549924in}}%
\pgfpathlineto{\pgfqpoint{2.357522in}{1.546974in}}%
\pgfpathmoveto{\pgfqpoint{2.357522in}{1.544025in}}%
\pgfpathlineto{\pgfqpoint{2.357522in}{1.544025in}}%
\pgfpathlineto{\pgfqpoint{2.357522in}{1.546974in}}%
\pgfpathlineto{\pgfqpoint{2.362063in}{1.546974in}}%
\pgfpathlineto{\pgfqpoint{2.362063in}{1.544025in}}%
\pgfpathmoveto{\pgfqpoint{2.357522in}{1.546974in}}%
\pgfpathlineto{\pgfqpoint{2.357522in}{1.546974in}}%
\pgfpathlineto{\pgfqpoint{2.357522in}{1.549924in}}%
\pgfpathlineto{\pgfqpoint{2.362063in}{1.549924in}}%
\pgfpathlineto{\pgfqpoint{2.362063in}{1.546974in}}%
\pgfpathmoveto{\pgfqpoint{2.362063in}{1.546974in}}%
\pgfpathlineto{\pgfqpoint{2.362063in}{1.546974in}}%
\pgfpathlineto{\pgfqpoint{2.362063in}{1.549924in}}%
\pgfpathlineto{\pgfqpoint{2.366604in}{1.549924in}}%
\pgfpathlineto{\pgfqpoint{2.366604in}{1.546974in}}%
\pgfpathmoveto{\pgfqpoint{2.357522in}{1.549924in}}%
\pgfpathlineto{\pgfqpoint{2.357522in}{1.549924in}}%
\pgfpathlineto{\pgfqpoint{2.357522in}{1.552873in}}%
\pgfpathlineto{\pgfqpoint{2.362063in}{1.552873in}}%
\pgfpathlineto{\pgfqpoint{2.362063in}{1.549924in}}%
\pgfpathmoveto{\pgfqpoint{2.357522in}{1.552873in}}%
\pgfpathlineto{\pgfqpoint{2.357522in}{1.552873in}}%
\pgfpathlineto{\pgfqpoint{2.357522in}{1.555822in}}%
\pgfpathlineto{\pgfqpoint{2.362063in}{1.555822in}}%
\pgfpathlineto{\pgfqpoint{2.362063in}{1.552873in}}%
\pgfpathmoveto{\pgfqpoint{2.362063in}{1.549924in}}%
\pgfpathlineto{\pgfqpoint{2.362063in}{1.549924in}}%
\pgfpathlineto{\pgfqpoint{2.362063in}{1.552873in}}%
\pgfpathlineto{\pgfqpoint{2.366604in}{1.552873in}}%
\pgfpathlineto{\pgfqpoint{2.366604in}{1.549924in}}%
\pgfpathmoveto{\pgfqpoint{2.362063in}{1.552873in}}%
\pgfpathlineto{\pgfqpoint{2.362063in}{1.552873in}}%
\pgfpathlineto{\pgfqpoint{2.362063in}{1.555822in}}%
\pgfpathlineto{\pgfqpoint{2.366604in}{1.555822in}}%
\pgfpathlineto{\pgfqpoint{2.366604in}{1.552873in}}%
\pgfpathmoveto{\pgfqpoint{2.366604in}{1.549924in}}%
\pgfpathlineto{\pgfqpoint{2.366604in}{1.549924in}}%
\pgfpathlineto{\pgfqpoint{2.366604in}{1.552873in}}%
\pgfpathlineto{\pgfqpoint{2.371145in}{1.552873in}}%
\pgfpathlineto{\pgfqpoint{2.371145in}{1.549924in}}%
\pgfpathmoveto{\pgfqpoint{2.366604in}{1.552873in}}%
\pgfpathlineto{\pgfqpoint{2.366604in}{1.552873in}}%
\pgfpathlineto{\pgfqpoint{2.366604in}{1.555822in}}%
\pgfpathlineto{\pgfqpoint{2.371145in}{1.555822in}}%
\pgfpathlineto{\pgfqpoint{2.371145in}{1.552873in}}%
\pgfpathmoveto{\pgfqpoint{2.371145in}{1.552873in}}%
\pgfpathlineto{\pgfqpoint{2.371145in}{1.552873in}}%
\pgfpathlineto{\pgfqpoint{2.371145in}{1.555822in}}%
\pgfpathlineto{\pgfqpoint{2.375686in}{1.555822in}}%
\pgfpathlineto{\pgfqpoint{2.375686in}{1.552873in}}%
\pgfpathmoveto{\pgfqpoint{2.375686in}{1.552873in}}%
\pgfpathlineto{\pgfqpoint{2.375686in}{1.552873in}}%
\pgfpathlineto{\pgfqpoint{2.375686in}{1.555822in}}%
\pgfpathlineto{\pgfqpoint{2.380227in}{1.555822in}}%
\pgfpathlineto{\pgfqpoint{2.380227in}{1.552873in}}%
\pgfpathmoveto{\pgfqpoint{2.375686in}{1.555822in}}%
\pgfpathlineto{\pgfqpoint{2.375686in}{1.555822in}}%
\pgfpathlineto{\pgfqpoint{2.375686in}{1.558771in}}%
\pgfpathlineto{\pgfqpoint{2.380227in}{1.558771in}}%
\pgfpathlineto{\pgfqpoint{2.380227in}{1.555822in}}%
\pgfpathmoveto{\pgfqpoint{2.375686in}{1.558771in}}%
\pgfpathlineto{\pgfqpoint{2.375686in}{1.558771in}}%
\pgfpathlineto{\pgfqpoint{2.375686in}{1.561720in}}%
\pgfpathlineto{\pgfqpoint{2.380227in}{1.561720in}}%
\pgfpathlineto{\pgfqpoint{2.380227in}{1.558771in}}%
\pgfpathmoveto{\pgfqpoint{2.380227in}{1.555822in}}%
\pgfpathlineto{\pgfqpoint{2.380227in}{1.555822in}}%
\pgfpathlineto{\pgfqpoint{2.380227in}{1.558771in}}%
\pgfpathlineto{\pgfqpoint{2.384768in}{1.558771in}}%
\pgfpathlineto{\pgfqpoint{2.384768in}{1.555822in}}%
\pgfpathmoveto{\pgfqpoint{2.380227in}{1.558771in}}%
\pgfpathlineto{\pgfqpoint{2.380227in}{1.558771in}}%
\pgfpathlineto{\pgfqpoint{2.380227in}{1.561720in}}%
\pgfpathlineto{\pgfqpoint{2.384768in}{1.561720in}}%
\pgfpathlineto{\pgfqpoint{2.384768in}{1.558771in}}%
\pgfpathmoveto{\pgfqpoint{2.384768in}{1.558771in}}%
\pgfpathlineto{\pgfqpoint{2.384768in}{1.558771in}}%
\pgfpathlineto{\pgfqpoint{2.384768in}{1.561720in}}%
\pgfpathlineto{\pgfqpoint{2.389309in}{1.561720in}}%
\pgfpathlineto{\pgfqpoint{2.389309in}{1.558771in}}%
\pgfpathmoveto{\pgfqpoint{2.384768in}{1.561720in}}%
\pgfpathlineto{\pgfqpoint{2.384768in}{1.561720in}}%
\pgfpathlineto{\pgfqpoint{2.384768in}{1.564670in}}%
\pgfpathlineto{\pgfqpoint{2.389309in}{1.564670in}}%
\pgfpathlineto{\pgfqpoint{2.389309in}{1.561720in}}%
\pgfpathmoveto{\pgfqpoint{2.384768in}{1.564670in}}%
\pgfpathlineto{\pgfqpoint{2.384768in}{1.564670in}}%
\pgfpathlineto{\pgfqpoint{2.384768in}{1.567619in}}%
\pgfpathlineto{\pgfqpoint{2.389309in}{1.567619in}}%
\pgfpathlineto{\pgfqpoint{2.389309in}{1.564670in}}%
\pgfpathmoveto{\pgfqpoint{2.389309in}{1.561720in}}%
\pgfpathlineto{\pgfqpoint{2.389309in}{1.561720in}}%
\pgfpathlineto{\pgfqpoint{2.389309in}{1.564670in}}%
\pgfpathlineto{\pgfqpoint{2.393850in}{1.564670in}}%
\pgfpathlineto{\pgfqpoint{2.393850in}{1.561720in}}%
\pgfpathmoveto{\pgfqpoint{2.389309in}{1.564670in}}%
\pgfpathlineto{\pgfqpoint{2.389309in}{1.564670in}}%
\pgfpathlineto{\pgfqpoint{2.389309in}{1.567619in}}%
\pgfpathlineto{\pgfqpoint{2.393850in}{1.567619in}}%
\pgfpathlineto{\pgfqpoint{2.393850in}{1.564670in}}%
\pgfpathmoveto{\pgfqpoint{2.393850in}{1.564670in}}%
\pgfpathlineto{\pgfqpoint{2.393850in}{1.564670in}}%
\pgfpathlineto{\pgfqpoint{2.393850in}{1.567619in}}%
\pgfpathlineto{\pgfqpoint{2.398391in}{1.567619in}}%
\pgfpathlineto{\pgfqpoint{2.398391in}{1.564670in}}%
\pgfpathmoveto{\pgfqpoint{2.393850in}{1.567619in}}%
\pgfpathlineto{\pgfqpoint{2.393850in}{1.567619in}}%
\pgfpathlineto{\pgfqpoint{2.393850in}{1.570568in}}%
\pgfpathlineto{\pgfqpoint{2.398391in}{1.570568in}}%
\pgfpathlineto{\pgfqpoint{2.398391in}{1.567619in}}%
\pgfpathmoveto{\pgfqpoint{2.393850in}{1.570568in}}%
\pgfpathlineto{\pgfqpoint{2.393850in}{1.570568in}}%
\pgfpathlineto{\pgfqpoint{2.393850in}{1.573517in}}%
\pgfpathlineto{\pgfqpoint{2.398391in}{1.573517in}}%
\pgfpathlineto{\pgfqpoint{2.398391in}{1.570568in}}%
\pgfpathmoveto{\pgfqpoint{2.398391in}{1.567619in}}%
\pgfpathlineto{\pgfqpoint{2.398391in}{1.567619in}}%
\pgfpathlineto{\pgfqpoint{2.398391in}{1.570568in}}%
\pgfpathlineto{\pgfqpoint{2.402931in}{1.570568in}}%
\pgfpathlineto{\pgfqpoint{2.402931in}{1.567619in}}%
\pgfpathmoveto{\pgfqpoint{2.398391in}{1.570568in}}%
\pgfpathlineto{\pgfqpoint{2.398391in}{1.570568in}}%
\pgfpathlineto{\pgfqpoint{2.398391in}{1.573517in}}%
\pgfpathlineto{\pgfqpoint{2.402931in}{1.573517in}}%
\pgfpathlineto{\pgfqpoint{2.402931in}{1.570568in}}%
\pgfpathmoveto{\pgfqpoint{2.402931in}{1.570568in}}%
\pgfpathlineto{\pgfqpoint{2.402931in}{1.570568in}}%
\pgfpathlineto{\pgfqpoint{2.402931in}{1.573517in}}%
\pgfpathlineto{\pgfqpoint{2.407472in}{1.573517in}}%
\pgfpathlineto{\pgfqpoint{2.407472in}{1.570568in}}%
\pgfpathmoveto{\pgfqpoint{2.402931in}{1.573517in}}%
\pgfpathlineto{\pgfqpoint{2.402931in}{1.573517in}}%
\pgfpathlineto{\pgfqpoint{2.402931in}{1.576466in}}%
\pgfpathlineto{\pgfqpoint{2.407472in}{1.576466in}}%
\pgfpathlineto{\pgfqpoint{2.407472in}{1.573517in}}%
\pgfpathmoveto{\pgfqpoint{2.402931in}{1.576466in}}%
\pgfpathlineto{\pgfqpoint{2.402931in}{1.576466in}}%
\pgfpathlineto{\pgfqpoint{2.402931in}{1.579416in}}%
\pgfpathlineto{\pgfqpoint{2.407472in}{1.579416in}}%
\pgfpathlineto{\pgfqpoint{2.407472in}{1.576466in}}%
\pgfpathmoveto{\pgfqpoint{2.407472in}{1.573517in}}%
\pgfpathlineto{\pgfqpoint{2.407472in}{1.573517in}}%
\pgfpathlineto{\pgfqpoint{2.407472in}{1.576466in}}%
\pgfpathlineto{\pgfqpoint{2.412013in}{1.576466in}}%
\pgfpathlineto{\pgfqpoint{2.412013in}{1.573517in}}%
\pgfpathmoveto{\pgfqpoint{2.407472in}{1.576466in}}%
\pgfpathlineto{\pgfqpoint{2.407472in}{1.576466in}}%
\pgfpathlineto{\pgfqpoint{2.407472in}{1.579416in}}%
\pgfpathlineto{\pgfqpoint{2.412013in}{1.579416in}}%
\pgfpathlineto{\pgfqpoint{2.412013in}{1.576466in}}%
\pgfpathmoveto{\pgfqpoint{2.412013in}{1.576466in}}%
\pgfpathlineto{\pgfqpoint{2.412013in}{1.576466in}}%
\pgfpathlineto{\pgfqpoint{2.412013in}{1.579416in}}%
\pgfpathlineto{\pgfqpoint{2.416554in}{1.579416in}}%
\pgfpathlineto{\pgfqpoint{2.416554in}{1.576466in}}%
\pgfpathmoveto{\pgfqpoint{2.412013in}{1.579416in}}%
\pgfpathlineto{\pgfqpoint{2.412013in}{1.579416in}}%
\pgfpathlineto{\pgfqpoint{2.412013in}{1.582365in}}%
\pgfpathlineto{\pgfqpoint{2.416554in}{1.582365in}}%
\pgfpathlineto{\pgfqpoint{2.416554in}{1.579416in}}%
\pgfpathmoveto{\pgfqpoint{2.412013in}{1.582365in}}%
\pgfpathlineto{\pgfqpoint{2.412013in}{1.582365in}}%
\pgfpathlineto{\pgfqpoint{2.412013in}{1.585314in}}%
\pgfpathlineto{\pgfqpoint{2.416554in}{1.585314in}}%
\pgfpathlineto{\pgfqpoint{2.416554in}{1.582365in}}%
\pgfpathmoveto{\pgfqpoint{2.416554in}{1.579416in}}%
\pgfpathlineto{\pgfqpoint{2.416554in}{1.579416in}}%
\pgfpathlineto{\pgfqpoint{2.416554in}{1.582365in}}%
\pgfpathlineto{\pgfqpoint{2.421095in}{1.582365in}}%
\pgfpathlineto{\pgfqpoint{2.421095in}{1.579416in}}%
\pgfpathmoveto{\pgfqpoint{2.416554in}{1.582365in}}%
\pgfpathlineto{\pgfqpoint{2.416554in}{1.582365in}}%
\pgfpathlineto{\pgfqpoint{2.416554in}{1.585314in}}%
\pgfpathlineto{\pgfqpoint{2.421095in}{1.585314in}}%
\pgfpathlineto{\pgfqpoint{2.421095in}{1.582365in}}%
\pgfpathmoveto{\pgfqpoint{2.421095in}{1.582365in}}%
\pgfpathlineto{\pgfqpoint{2.421095in}{1.582365in}}%
\pgfpathlineto{\pgfqpoint{2.421095in}{1.585314in}}%
\pgfpathlineto{\pgfqpoint{2.425636in}{1.585314in}}%
\pgfpathlineto{\pgfqpoint{2.425636in}{1.582365in}}%
\pgfpathmoveto{\pgfqpoint{2.421095in}{1.585314in}}%
\pgfpathlineto{\pgfqpoint{2.421095in}{1.585314in}}%
\pgfpathlineto{\pgfqpoint{2.421095in}{1.588263in}}%
\pgfpathlineto{\pgfqpoint{2.425636in}{1.588263in}}%
\pgfpathlineto{\pgfqpoint{2.425636in}{1.585314in}}%
\pgfpathmoveto{\pgfqpoint{2.421095in}{1.588263in}}%
\pgfpathlineto{\pgfqpoint{2.421095in}{1.588263in}}%
\pgfpathlineto{\pgfqpoint{2.421095in}{1.591212in}}%
\pgfpathlineto{\pgfqpoint{2.425636in}{1.591212in}}%
\pgfpathlineto{\pgfqpoint{2.425636in}{1.588263in}}%
\pgfpathmoveto{\pgfqpoint{2.425636in}{1.585314in}}%
\pgfpathlineto{\pgfqpoint{2.425636in}{1.585314in}}%
\pgfpathlineto{\pgfqpoint{2.425636in}{1.588263in}}%
\pgfpathlineto{\pgfqpoint{2.430177in}{1.588263in}}%
\pgfpathlineto{\pgfqpoint{2.430177in}{1.585314in}}%
\pgfpathmoveto{\pgfqpoint{2.425636in}{1.588263in}}%
\pgfpathlineto{\pgfqpoint{2.425636in}{1.588263in}}%
\pgfpathlineto{\pgfqpoint{2.425636in}{1.591212in}}%
\pgfpathlineto{\pgfqpoint{2.430177in}{1.591212in}}%
\pgfpathlineto{\pgfqpoint{2.430177in}{1.588263in}}%
\pgfpathmoveto{\pgfqpoint{2.430177in}{1.588263in}}%
\pgfpathlineto{\pgfqpoint{2.430177in}{1.588263in}}%
\pgfpathlineto{\pgfqpoint{2.430177in}{1.591212in}}%
\pgfpathlineto{\pgfqpoint{2.434718in}{1.591212in}}%
\pgfpathlineto{\pgfqpoint{2.434718in}{1.588263in}}%
\pgfpathmoveto{\pgfqpoint{2.430177in}{1.591212in}}%
\pgfpathlineto{\pgfqpoint{2.430177in}{1.591212in}}%
\pgfpathlineto{\pgfqpoint{2.430177in}{1.594162in}}%
\pgfpathlineto{\pgfqpoint{2.434718in}{1.594162in}}%
\pgfpathlineto{\pgfqpoint{2.434718in}{1.591212in}}%
\pgfpathmoveto{\pgfqpoint{2.430177in}{1.594162in}}%
\pgfpathlineto{\pgfqpoint{2.430177in}{1.594162in}}%
\pgfpathlineto{\pgfqpoint{2.430177in}{1.597111in}}%
\pgfpathlineto{\pgfqpoint{2.434718in}{1.597111in}}%
\pgfpathlineto{\pgfqpoint{2.434718in}{1.594162in}}%
\pgfpathmoveto{\pgfqpoint{2.434718in}{1.591212in}}%
\pgfpathlineto{\pgfqpoint{2.434718in}{1.591212in}}%
\pgfpathlineto{\pgfqpoint{2.434718in}{1.594162in}}%
\pgfpathlineto{\pgfqpoint{2.439259in}{1.594162in}}%
\pgfpathlineto{\pgfqpoint{2.439259in}{1.591212in}}%
\pgfpathmoveto{\pgfqpoint{2.434718in}{1.594162in}}%
\pgfpathlineto{\pgfqpoint{2.434718in}{1.594162in}}%
\pgfpathlineto{\pgfqpoint{2.434718in}{1.597111in}}%
\pgfpathlineto{\pgfqpoint{2.439259in}{1.597111in}}%
\pgfpathlineto{\pgfqpoint{2.439259in}{1.594162in}}%
\pgfpathmoveto{\pgfqpoint{2.439259in}{1.594162in}}%
\pgfpathlineto{\pgfqpoint{2.439259in}{1.594162in}}%
\pgfpathlineto{\pgfqpoint{2.439259in}{1.597111in}}%
\pgfpathlineto{\pgfqpoint{2.443800in}{1.597111in}}%
\pgfpathlineto{\pgfqpoint{2.443800in}{1.594162in}}%
\pgfpathmoveto{\pgfqpoint{2.439259in}{1.597111in}}%
\pgfpathlineto{\pgfqpoint{2.439259in}{1.597111in}}%
\pgfpathlineto{\pgfqpoint{2.439259in}{1.600060in}}%
\pgfpathlineto{\pgfqpoint{2.443800in}{1.600060in}}%
\pgfpathlineto{\pgfqpoint{2.443800in}{1.597111in}}%
\pgfpathmoveto{\pgfqpoint{2.439259in}{1.600060in}}%
\pgfpathlineto{\pgfqpoint{2.439259in}{1.600060in}}%
\pgfpathlineto{\pgfqpoint{2.439259in}{1.603009in}}%
\pgfpathlineto{\pgfqpoint{2.443800in}{1.603009in}}%
\pgfpathlineto{\pgfqpoint{2.443800in}{1.600060in}}%
\pgfpathmoveto{\pgfqpoint{2.443800in}{1.597111in}}%
\pgfpathlineto{\pgfqpoint{2.443800in}{1.597111in}}%
\pgfpathlineto{\pgfqpoint{2.443800in}{1.600060in}}%
\pgfpathlineto{\pgfqpoint{2.448341in}{1.600060in}}%
\pgfpathlineto{\pgfqpoint{2.448341in}{1.597111in}}%
\pgfpathmoveto{\pgfqpoint{2.443800in}{1.600060in}}%
\pgfpathlineto{\pgfqpoint{2.443800in}{1.600060in}}%
\pgfpathlineto{\pgfqpoint{2.443800in}{1.603009in}}%
\pgfpathlineto{\pgfqpoint{2.448341in}{1.603009in}}%
\pgfpathlineto{\pgfqpoint{2.448341in}{1.600060in}}%
\pgfpathmoveto{\pgfqpoint{2.448341in}{1.600060in}}%
\pgfpathlineto{\pgfqpoint{2.448341in}{1.600060in}}%
\pgfpathlineto{\pgfqpoint{2.448341in}{1.603009in}}%
\pgfpathlineto{\pgfqpoint{2.452882in}{1.603009in}}%
\pgfpathlineto{\pgfqpoint{2.452882in}{1.600060in}}%
\pgfpathmoveto{\pgfqpoint{2.448341in}{1.603009in}}%
\pgfpathlineto{\pgfqpoint{2.448341in}{1.603009in}}%
\pgfpathlineto{\pgfqpoint{2.448341in}{1.605958in}}%
\pgfpathlineto{\pgfqpoint{2.452882in}{1.605958in}}%
\pgfpathlineto{\pgfqpoint{2.452882in}{1.603009in}}%
\pgfpathmoveto{\pgfqpoint{2.448341in}{1.605958in}}%
\pgfpathlineto{\pgfqpoint{2.448341in}{1.605958in}}%
\pgfpathlineto{\pgfqpoint{2.448341in}{1.608908in}}%
\pgfpathlineto{\pgfqpoint{2.452882in}{1.608908in}}%
\pgfpathlineto{\pgfqpoint{2.452882in}{1.605958in}}%
\pgfpathmoveto{\pgfqpoint{2.452882in}{1.603009in}}%
\pgfpathlineto{\pgfqpoint{2.452882in}{1.603009in}}%
\pgfpathlineto{\pgfqpoint{2.452882in}{1.605958in}}%
\pgfpathlineto{\pgfqpoint{2.457422in}{1.605958in}}%
\pgfpathlineto{\pgfqpoint{2.457422in}{1.603009in}}%
\pgfpathmoveto{\pgfqpoint{2.452882in}{1.605958in}}%
\pgfpathlineto{\pgfqpoint{2.452882in}{1.605958in}}%
\pgfpathlineto{\pgfqpoint{2.452882in}{1.608908in}}%
\pgfpathlineto{\pgfqpoint{2.457422in}{1.608908in}}%
\pgfpathlineto{\pgfqpoint{2.457422in}{1.605958in}}%
\pgfpathmoveto{\pgfqpoint{2.457422in}{1.605958in}}%
\pgfpathlineto{\pgfqpoint{2.457422in}{1.605958in}}%
\pgfpathlineto{\pgfqpoint{2.457422in}{1.608908in}}%
\pgfpathlineto{\pgfqpoint{2.461963in}{1.608908in}}%
\pgfpathlineto{\pgfqpoint{2.461963in}{1.605958in}}%
\pgfpathmoveto{\pgfqpoint{2.457422in}{1.608908in}}%
\pgfpathlineto{\pgfqpoint{2.457422in}{1.608908in}}%
\pgfpathlineto{\pgfqpoint{2.457422in}{1.611857in}}%
\pgfpathlineto{\pgfqpoint{2.461963in}{1.611857in}}%
\pgfpathlineto{\pgfqpoint{2.461963in}{1.608908in}}%
\pgfpathmoveto{\pgfqpoint{2.457422in}{1.611857in}}%
\pgfpathlineto{\pgfqpoint{2.457422in}{1.611857in}}%
\pgfpathlineto{\pgfqpoint{2.457422in}{1.614806in}}%
\pgfpathlineto{\pgfqpoint{2.461963in}{1.614806in}}%
\pgfpathlineto{\pgfqpoint{2.461963in}{1.611857in}}%
\pgfpathmoveto{\pgfqpoint{2.461963in}{1.608908in}}%
\pgfpathlineto{\pgfqpoint{2.461963in}{1.608908in}}%
\pgfpathlineto{\pgfqpoint{2.461963in}{1.611857in}}%
\pgfpathlineto{\pgfqpoint{2.466504in}{1.611857in}}%
\pgfpathlineto{\pgfqpoint{2.466504in}{1.608908in}}%
\pgfpathmoveto{\pgfqpoint{2.461963in}{1.611857in}}%
\pgfpathlineto{\pgfqpoint{2.461963in}{1.611857in}}%
\pgfpathlineto{\pgfqpoint{2.461963in}{1.614806in}}%
\pgfpathlineto{\pgfqpoint{2.466504in}{1.614806in}}%
\pgfpathlineto{\pgfqpoint{2.466504in}{1.611857in}}%
\pgfpathmoveto{\pgfqpoint{2.466504in}{1.611857in}}%
\pgfpathlineto{\pgfqpoint{2.466504in}{1.611857in}}%
\pgfpathlineto{\pgfqpoint{2.466504in}{1.614806in}}%
\pgfpathlineto{\pgfqpoint{2.471045in}{1.614806in}}%
\pgfpathlineto{\pgfqpoint{2.471045in}{1.611857in}}%
\pgfpathmoveto{\pgfqpoint{2.466504in}{1.614806in}}%
\pgfpathlineto{\pgfqpoint{2.466504in}{1.614806in}}%
\pgfpathlineto{\pgfqpoint{2.466504in}{1.617755in}}%
\pgfpathlineto{\pgfqpoint{2.471045in}{1.617755in}}%
\pgfpathlineto{\pgfqpoint{2.471045in}{1.614806in}}%
\pgfpathmoveto{\pgfqpoint{2.466504in}{1.617755in}}%
\pgfpathlineto{\pgfqpoint{2.466504in}{1.617755in}}%
\pgfpathlineto{\pgfqpoint{2.466504in}{1.620704in}}%
\pgfpathlineto{\pgfqpoint{2.471045in}{1.620704in}}%
\pgfpathlineto{\pgfqpoint{2.471045in}{1.617755in}}%
\pgfpathmoveto{\pgfqpoint{2.471045in}{1.614806in}}%
\pgfpathlineto{\pgfqpoint{2.471045in}{1.614806in}}%
\pgfpathlineto{\pgfqpoint{2.471045in}{1.617755in}}%
\pgfpathlineto{\pgfqpoint{2.475586in}{1.617755in}}%
\pgfpathlineto{\pgfqpoint{2.475586in}{1.614806in}}%
\pgfpathmoveto{\pgfqpoint{2.471045in}{1.617755in}}%
\pgfpathlineto{\pgfqpoint{2.471045in}{1.617755in}}%
\pgfpathlineto{\pgfqpoint{2.471045in}{1.620704in}}%
\pgfpathlineto{\pgfqpoint{2.475586in}{1.620704in}}%
\pgfpathlineto{\pgfqpoint{2.475586in}{1.617755in}}%
\pgfpathmoveto{\pgfqpoint{2.475586in}{1.617755in}}%
\pgfpathlineto{\pgfqpoint{2.475586in}{1.617755in}}%
\pgfpathlineto{\pgfqpoint{2.475586in}{1.620704in}}%
\pgfpathlineto{\pgfqpoint{2.480127in}{1.620704in}}%
\pgfpathlineto{\pgfqpoint{2.480127in}{1.617755in}}%
\pgfpathmoveto{\pgfqpoint{2.475586in}{1.620704in}}%
\pgfpathlineto{\pgfqpoint{2.475586in}{1.620704in}}%
\pgfpathlineto{\pgfqpoint{2.475586in}{1.623654in}}%
\pgfpathlineto{\pgfqpoint{2.480127in}{1.623654in}}%
\pgfpathlineto{\pgfqpoint{2.480127in}{1.620704in}}%
\pgfpathmoveto{\pgfqpoint{2.475586in}{1.623654in}}%
\pgfpathlineto{\pgfqpoint{2.475586in}{1.623654in}}%
\pgfpathlineto{\pgfqpoint{2.475586in}{1.626603in}}%
\pgfpathlineto{\pgfqpoint{2.480127in}{1.626603in}}%
\pgfpathlineto{\pgfqpoint{2.480127in}{1.623654in}}%
\pgfpathmoveto{\pgfqpoint{2.480127in}{1.620704in}}%
\pgfpathlineto{\pgfqpoint{2.480127in}{1.620704in}}%
\pgfpathlineto{\pgfqpoint{2.480127in}{1.623654in}}%
\pgfpathlineto{\pgfqpoint{2.484668in}{1.623654in}}%
\pgfpathlineto{\pgfqpoint{2.484668in}{1.620704in}}%
\pgfpathmoveto{\pgfqpoint{2.480127in}{1.623654in}}%
\pgfpathlineto{\pgfqpoint{2.480127in}{1.623654in}}%
\pgfpathlineto{\pgfqpoint{2.480127in}{1.626603in}}%
\pgfpathlineto{\pgfqpoint{2.484668in}{1.626603in}}%
\pgfpathlineto{\pgfqpoint{2.484668in}{1.623654in}}%
\pgfpathmoveto{\pgfqpoint{2.484668in}{1.623654in}}%
\pgfpathlineto{\pgfqpoint{2.484668in}{1.623654in}}%
\pgfpathlineto{\pgfqpoint{2.484668in}{1.626603in}}%
\pgfpathlineto{\pgfqpoint{2.489209in}{1.626603in}}%
\pgfpathlineto{\pgfqpoint{2.489209in}{1.623654in}}%
\pgfpathmoveto{\pgfqpoint{2.484668in}{1.626603in}}%
\pgfpathlineto{\pgfqpoint{2.484668in}{1.626603in}}%
\pgfpathlineto{\pgfqpoint{2.484668in}{1.629552in}}%
\pgfpathlineto{\pgfqpoint{2.489209in}{1.629552in}}%
\pgfpathlineto{\pgfqpoint{2.489209in}{1.626603in}}%
\pgfpathmoveto{\pgfqpoint{2.484668in}{1.629552in}}%
\pgfpathlineto{\pgfqpoint{2.484668in}{1.629552in}}%
\pgfpathlineto{\pgfqpoint{2.484668in}{1.632501in}}%
\pgfpathlineto{\pgfqpoint{2.489209in}{1.632501in}}%
\pgfpathlineto{\pgfqpoint{2.489209in}{1.629552in}}%
\pgfpathmoveto{\pgfqpoint{2.489209in}{1.626603in}}%
\pgfpathlineto{\pgfqpoint{2.489209in}{1.626603in}}%
\pgfpathlineto{\pgfqpoint{2.489209in}{1.629552in}}%
\pgfpathlineto{\pgfqpoint{2.493750in}{1.629552in}}%
\pgfpathlineto{\pgfqpoint{2.493750in}{1.626603in}}%
\pgfpathmoveto{\pgfqpoint{2.489209in}{1.629552in}}%
\pgfpathlineto{\pgfqpoint{2.489209in}{1.629552in}}%
\pgfpathlineto{\pgfqpoint{2.489209in}{1.632501in}}%
\pgfpathlineto{\pgfqpoint{2.493750in}{1.632501in}}%
\pgfpathlineto{\pgfqpoint{2.493750in}{1.629552in}}%
\pgfpathmoveto{\pgfqpoint{2.493750in}{1.629552in}}%
\pgfpathlineto{\pgfqpoint{2.493750in}{1.629552in}}%
\pgfpathlineto{\pgfqpoint{2.493750in}{1.632501in}}%
\pgfpathlineto{\pgfqpoint{2.498291in}{1.632501in}}%
\pgfpathlineto{\pgfqpoint{2.498291in}{1.629552in}}%
\pgfpathmoveto{\pgfqpoint{2.493750in}{1.632501in}}%
\pgfpathlineto{\pgfqpoint{2.493750in}{1.632501in}}%
\pgfpathlineto{\pgfqpoint{2.493750in}{1.635450in}}%
\pgfpathlineto{\pgfqpoint{2.498291in}{1.635450in}}%
\pgfpathlineto{\pgfqpoint{2.498291in}{1.632501in}}%
\pgfpathmoveto{\pgfqpoint{2.493750in}{1.635450in}}%
\pgfpathlineto{\pgfqpoint{2.493750in}{1.635450in}}%
\pgfpathlineto{\pgfqpoint{2.493750in}{1.638400in}}%
\pgfpathlineto{\pgfqpoint{2.498291in}{1.638400in}}%
\pgfpathlineto{\pgfqpoint{2.498291in}{1.635450in}}%
\pgfpathmoveto{\pgfqpoint{2.498291in}{1.632501in}}%
\pgfpathlineto{\pgfqpoint{2.498291in}{1.632501in}}%
\pgfpathlineto{\pgfqpoint{2.498291in}{1.635450in}}%
\pgfpathlineto{\pgfqpoint{2.502832in}{1.635450in}}%
\pgfpathlineto{\pgfqpoint{2.502832in}{1.632501in}}%
\pgfpathmoveto{\pgfqpoint{2.498291in}{1.635450in}}%
\pgfpathlineto{\pgfqpoint{2.498291in}{1.635450in}}%
\pgfpathlineto{\pgfqpoint{2.498291in}{1.638400in}}%
\pgfpathlineto{\pgfqpoint{2.502832in}{1.638400in}}%
\pgfpathlineto{\pgfqpoint{2.502832in}{1.635450in}}%
\pgfpathmoveto{\pgfqpoint{2.502832in}{1.635450in}}%
\pgfpathlineto{\pgfqpoint{2.502832in}{1.635450in}}%
\pgfpathlineto{\pgfqpoint{2.502832in}{1.638400in}}%
\pgfpathlineto{\pgfqpoint{2.507373in}{1.638400in}}%
\pgfpathlineto{\pgfqpoint{2.507373in}{1.635450in}}%
\pgfpathmoveto{\pgfqpoint{2.502832in}{1.638400in}}%
\pgfpathlineto{\pgfqpoint{2.502832in}{1.638400in}}%
\pgfpathlineto{\pgfqpoint{2.502832in}{1.641349in}}%
\pgfpathlineto{\pgfqpoint{2.507373in}{1.641349in}}%
\pgfpathlineto{\pgfqpoint{2.507373in}{1.638400in}}%
\pgfpathmoveto{\pgfqpoint{2.502832in}{1.641349in}}%
\pgfpathlineto{\pgfqpoint{2.502832in}{1.641349in}}%
\pgfpathlineto{\pgfqpoint{2.502832in}{1.644298in}}%
\pgfpathlineto{\pgfqpoint{2.507373in}{1.644298in}}%
\pgfpathlineto{\pgfqpoint{2.507373in}{1.641349in}}%
\pgfpathmoveto{\pgfqpoint{2.507373in}{1.638400in}}%
\pgfpathlineto{\pgfqpoint{2.507373in}{1.638400in}}%
\pgfpathlineto{\pgfqpoint{2.507373in}{1.641349in}}%
\pgfpathlineto{\pgfqpoint{2.511914in}{1.641349in}}%
\pgfpathlineto{\pgfqpoint{2.511914in}{1.638400in}}%
\pgfpathmoveto{\pgfqpoint{2.507373in}{1.641349in}}%
\pgfpathlineto{\pgfqpoint{2.507373in}{1.641349in}}%
\pgfpathlineto{\pgfqpoint{2.507373in}{1.644298in}}%
\pgfpathlineto{\pgfqpoint{2.511914in}{1.644298in}}%
\pgfpathlineto{\pgfqpoint{2.511914in}{1.641349in}}%
\pgfpathmoveto{\pgfqpoint{2.511914in}{1.641349in}}%
\pgfpathlineto{\pgfqpoint{2.511914in}{1.641349in}}%
\pgfpathlineto{\pgfqpoint{2.511914in}{1.644298in}}%
\pgfpathlineto{\pgfqpoint{2.516455in}{1.644298in}}%
\pgfpathlineto{\pgfqpoint{2.516455in}{1.641349in}}%
\pgfpathmoveto{\pgfqpoint{2.511914in}{1.644298in}}%
\pgfpathlineto{\pgfqpoint{2.511914in}{1.644298in}}%
\pgfpathlineto{\pgfqpoint{2.511914in}{1.647247in}}%
\pgfpathlineto{\pgfqpoint{2.516455in}{1.647247in}}%
\pgfpathlineto{\pgfqpoint{2.516455in}{1.644298in}}%
\pgfpathmoveto{\pgfqpoint{2.511914in}{1.647247in}}%
\pgfpathlineto{\pgfqpoint{2.511914in}{1.647247in}}%
\pgfpathlineto{\pgfqpoint{2.511914in}{1.650196in}}%
\pgfpathlineto{\pgfqpoint{2.516455in}{1.650196in}}%
\pgfpathlineto{\pgfqpoint{2.516455in}{1.647247in}}%
\pgfpathmoveto{\pgfqpoint{2.516455in}{1.644298in}}%
\pgfpathlineto{\pgfqpoint{2.516455in}{1.644298in}}%
\pgfpathlineto{\pgfqpoint{2.516455in}{1.647247in}}%
\pgfpathlineto{\pgfqpoint{2.520996in}{1.647247in}}%
\pgfpathlineto{\pgfqpoint{2.520996in}{1.644298in}}%
\pgfpathmoveto{\pgfqpoint{2.516455in}{1.647247in}}%
\pgfpathlineto{\pgfqpoint{2.516455in}{1.647247in}}%
\pgfpathlineto{\pgfqpoint{2.516455in}{1.650196in}}%
\pgfpathlineto{\pgfqpoint{2.520996in}{1.650196in}}%
\pgfpathlineto{\pgfqpoint{2.520996in}{1.647247in}}%
\pgfpathmoveto{\pgfqpoint{2.520996in}{1.647247in}}%
\pgfpathlineto{\pgfqpoint{2.520996in}{1.647247in}}%
\pgfpathlineto{\pgfqpoint{2.520996in}{1.650196in}}%
\pgfpathlineto{\pgfqpoint{2.525537in}{1.650196in}}%
\pgfpathlineto{\pgfqpoint{2.525537in}{1.647247in}}%
\pgfpathmoveto{\pgfqpoint{2.520996in}{1.650196in}}%
\pgfpathlineto{\pgfqpoint{2.520996in}{1.650196in}}%
\pgfpathlineto{\pgfqpoint{2.520996in}{1.653146in}}%
\pgfpathlineto{\pgfqpoint{2.525537in}{1.653146in}}%
\pgfpathlineto{\pgfqpoint{2.525537in}{1.650196in}}%
\pgfpathmoveto{\pgfqpoint{2.520996in}{1.653146in}}%
\pgfpathlineto{\pgfqpoint{2.520996in}{1.653146in}}%
\pgfpathlineto{\pgfqpoint{2.520996in}{1.656095in}}%
\pgfpathlineto{\pgfqpoint{2.525537in}{1.656095in}}%
\pgfpathlineto{\pgfqpoint{2.525537in}{1.653146in}}%
\pgfpathmoveto{\pgfqpoint{2.525537in}{1.650196in}}%
\pgfpathlineto{\pgfqpoint{2.525537in}{1.650196in}}%
\pgfpathlineto{\pgfqpoint{2.525537in}{1.653146in}}%
\pgfpathlineto{\pgfqpoint{2.530078in}{1.653146in}}%
\pgfpathlineto{\pgfqpoint{2.530078in}{1.650196in}}%
\pgfpathmoveto{\pgfqpoint{2.525537in}{1.653146in}}%
\pgfpathlineto{\pgfqpoint{2.525537in}{1.653146in}}%
\pgfpathlineto{\pgfqpoint{2.525537in}{1.656095in}}%
\pgfpathlineto{\pgfqpoint{2.530078in}{1.656095in}}%
\pgfpathlineto{\pgfqpoint{2.530078in}{1.653146in}}%
\pgfpathmoveto{\pgfqpoint{2.530078in}{1.653146in}}%
\pgfpathlineto{\pgfqpoint{2.530078in}{1.653146in}}%
\pgfpathlineto{\pgfqpoint{2.530078in}{1.656095in}}%
\pgfpathlineto{\pgfqpoint{2.534619in}{1.656095in}}%
\pgfpathlineto{\pgfqpoint{2.534619in}{1.653146in}}%
\pgfpathmoveto{\pgfqpoint{2.530078in}{1.656095in}}%
\pgfpathlineto{\pgfqpoint{2.530078in}{1.656095in}}%
\pgfpathlineto{\pgfqpoint{2.530078in}{1.659044in}}%
\pgfpathlineto{\pgfqpoint{2.534619in}{1.659044in}}%
\pgfpathlineto{\pgfqpoint{2.534619in}{1.656095in}}%
\pgfpathmoveto{\pgfqpoint{2.530078in}{1.659044in}}%
\pgfpathlineto{\pgfqpoint{2.530078in}{1.659044in}}%
\pgfpathlineto{\pgfqpoint{2.530078in}{1.661993in}}%
\pgfpathlineto{\pgfqpoint{2.534619in}{1.661993in}}%
\pgfpathlineto{\pgfqpoint{2.534619in}{1.659044in}}%
\pgfpathmoveto{\pgfqpoint{2.534619in}{1.656095in}}%
\pgfpathlineto{\pgfqpoint{2.534619in}{1.656095in}}%
\pgfpathlineto{\pgfqpoint{2.534619in}{1.659044in}}%
\pgfpathlineto{\pgfqpoint{2.539160in}{1.659044in}}%
\pgfpathlineto{\pgfqpoint{2.539160in}{1.656095in}}%
\pgfpathmoveto{\pgfqpoint{2.534619in}{1.659044in}}%
\pgfpathlineto{\pgfqpoint{2.534619in}{1.659044in}}%
\pgfpathlineto{\pgfqpoint{2.534619in}{1.661993in}}%
\pgfpathlineto{\pgfqpoint{2.539160in}{1.661993in}}%
\pgfpathlineto{\pgfqpoint{2.539160in}{1.659044in}}%
\pgfpathmoveto{\pgfqpoint{2.539160in}{1.659044in}}%
\pgfpathlineto{\pgfqpoint{2.539160in}{1.659044in}}%
\pgfpathlineto{\pgfqpoint{2.539160in}{1.661993in}}%
\pgfpathlineto{\pgfqpoint{2.543701in}{1.661993in}}%
\pgfpathlineto{\pgfqpoint{2.543701in}{1.659044in}}%
\pgfpathmoveto{\pgfqpoint{2.539160in}{1.661993in}}%
\pgfpathlineto{\pgfqpoint{2.539160in}{1.661993in}}%
\pgfpathlineto{\pgfqpoint{2.539160in}{1.664942in}}%
\pgfpathlineto{\pgfqpoint{2.543701in}{1.664942in}}%
\pgfpathlineto{\pgfqpoint{2.543701in}{1.661993in}}%
\pgfpathmoveto{\pgfqpoint{2.539160in}{1.664942in}}%
\pgfpathlineto{\pgfqpoint{2.539160in}{1.664942in}}%
\pgfpathlineto{\pgfqpoint{2.539160in}{1.667892in}}%
\pgfpathlineto{\pgfqpoint{2.543701in}{1.667892in}}%
\pgfpathlineto{\pgfqpoint{2.543701in}{1.664942in}}%
\pgfpathmoveto{\pgfqpoint{2.543701in}{1.661993in}}%
\pgfpathlineto{\pgfqpoint{2.543701in}{1.661993in}}%
\pgfpathlineto{\pgfqpoint{2.543701in}{1.664942in}}%
\pgfpathlineto{\pgfqpoint{2.548242in}{1.664942in}}%
\pgfpathlineto{\pgfqpoint{2.548242in}{1.661993in}}%
\pgfpathmoveto{\pgfqpoint{2.543701in}{1.664942in}}%
\pgfpathlineto{\pgfqpoint{2.543701in}{1.664942in}}%
\pgfpathlineto{\pgfqpoint{2.543701in}{1.667892in}}%
\pgfpathlineto{\pgfqpoint{2.548242in}{1.667892in}}%
\pgfpathlineto{\pgfqpoint{2.548242in}{1.664942in}}%
\pgfpathmoveto{\pgfqpoint{2.548242in}{1.664942in}}%
\pgfpathlineto{\pgfqpoint{2.548242in}{1.664942in}}%
\pgfpathlineto{\pgfqpoint{2.548242in}{1.667892in}}%
\pgfpathlineto{\pgfqpoint{2.552783in}{1.667892in}}%
\pgfpathlineto{\pgfqpoint{2.552783in}{1.664942in}}%
\pgfpathmoveto{\pgfqpoint{2.548242in}{1.667892in}}%
\pgfpathlineto{\pgfqpoint{2.548242in}{1.667892in}}%
\pgfpathlineto{\pgfqpoint{2.548242in}{1.670841in}}%
\pgfpathlineto{\pgfqpoint{2.552783in}{1.670841in}}%
\pgfpathlineto{\pgfqpoint{2.552783in}{1.667892in}}%
\pgfpathmoveto{\pgfqpoint{2.548242in}{1.670841in}}%
\pgfpathlineto{\pgfqpoint{2.548242in}{1.670841in}}%
\pgfpathlineto{\pgfqpoint{2.548242in}{1.673790in}}%
\pgfpathlineto{\pgfqpoint{2.552783in}{1.673790in}}%
\pgfpathlineto{\pgfqpoint{2.552783in}{1.670841in}}%
\pgfpathmoveto{\pgfqpoint{2.552783in}{1.667892in}}%
\pgfpathlineto{\pgfqpoint{2.552783in}{1.667892in}}%
\pgfpathlineto{\pgfqpoint{2.552783in}{1.670841in}}%
\pgfpathlineto{\pgfqpoint{2.557324in}{1.670841in}}%
\pgfpathlineto{\pgfqpoint{2.557324in}{1.667892in}}%
\pgfpathmoveto{\pgfqpoint{2.552783in}{1.670841in}}%
\pgfpathlineto{\pgfqpoint{2.552783in}{1.670841in}}%
\pgfpathlineto{\pgfqpoint{2.552783in}{1.673790in}}%
\pgfpathlineto{\pgfqpoint{2.557324in}{1.673790in}}%
\pgfpathlineto{\pgfqpoint{2.557324in}{1.670841in}}%
\pgfpathmoveto{\pgfqpoint{2.557324in}{1.670841in}}%
\pgfpathlineto{\pgfqpoint{2.557324in}{1.670841in}}%
\pgfpathlineto{\pgfqpoint{2.557324in}{1.673790in}}%
\pgfpathlineto{\pgfqpoint{2.561865in}{1.673790in}}%
\pgfpathlineto{\pgfqpoint{2.561865in}{1.670841in}}%
\pgfpathmoveto{\pgfqpoint{2.557324in}{1.673790in}}%
\pgfpathlineto{\pgfqpoint{2.557324in}{1.673790in}}%
\pgfpathlineto{\pgfqpoint{2.557324in}{1.676739in}}%
\pgfpathlineto{\pgfqpoint{2.561865in}{1.676739in}}%
\pgfpathlineto{\pgfqpoint{2.561865in}{1.673790in}}%
\pgfpathmoveto{\pgfqpoint{2.557324in}{1.676739in}}%
\pgfpathlineto{\pgfqpoint{2.557324in}{1.676739in}}%
\pgfpathlineto{\pgfqpoint{2.557324in}{1.679689in}}%
\pgfpathlineto{\pgfqpoint{2.561865in}{1.679689in}}%
\pgfpathlineto{\pgfqpoint{2.561865in}{1.676739in}}%
\pgfpathmoveto{\pgfqpoint{2.561865in}{1.673790in}}%
\pgfpathlineto{\pgfqpoint{2.561865in}{1.673790in}}%
\pgfpathlineto{\pgfqpoint{2.561865in}{1.676739in}}%
\pgfpathlineto{\pgfqpoint{2.566406in}{1.676739in}}%
\pgfpathlineto{\pgfqpoint{2.566406in}{1.673790in}}%
\pgfpathmoveto{\pgfqpoint{2.561865in}{1.676739in}}%
\pgfpathlineto{\pgfqpoint{2.561865in}{1.676739in}}%
\pgfpathlineto{\pgfqpoint{2.561865in}{1.679689in}}%
\pgfpathlineto{\pgfqpoint{2.566406in}{1.679689in}}%
\pgfpathlineto{\pgfqpoint{2.566406in}{1.676739in}}%
\pgfpathmoveto{\pgfqpoint{2.566406in}{1.676739in}}%
\pgfpathlineto{\pgfqpoint{2.566406in}{1.676739in}}%
\pgfpathlineto{\pgfqpoint{2.566406in}{1.679689in}}%
\pgfpathlineto{\pgfqpoint{2.570947in}{1.679689in}}%
\pgfpathlineto{\pgfqpoint{2.570947in}{1.676739in}}%
\pgfpathmoveto{\pgfqpoint{2.566406in}{1.679689in}}%
\pgfpathlineto{\pgfqpoint{2.566406in}{1.679689in}}%
\pgfpathlineto{\pgfqpoint{2.566406in}{1.682638in}}%
\pgfpathlineto{\pgfqpoint{2.570947in}{1.682638in}}%
\pgfpathlineto{\pgfqpoint{2.570947in}{1.679689in}}%
\pgfpathmoveto{\pgfqpoint{2.566406in}{1.682638in}}%
\pgfpathlineto{\pgfqpoint{2.566406in}{1.682638in}}%
\pgfpathlineto{\pgfqpoint{2.566406in}{1.685587in}}%
\pgfpathlineto{\pgfqpoint{2.570947in}{1.685587in}}%
\pgfpathlineto{\pgfqpoint{2.570947in}{1.682638in}}%
\pgfpathmoveto{\pgfqpoint{2.570947in}{1.679689in}}%
\pgfpathlineto{\pgfqpoint{2.570947in}{1.679689in}}%
\pgfpathlineto{\pgfqpoint{2.570947in}{1.682638in}}%
\pgfpathlineto{\pgfqpoint{2.575489in}{1.682638in}}%
\pgfpathlineto{\pgfqpoint{2.575489in}{1.679689in}}%
\pgfpathmoveto{\pgfqpoint{2.570947in}{1.682638in}}%
\pgfpathlineto{\pgfqpoint{2.570947in}{1.682638in}}%
\pgfpathlineto{\pgfqpoint{2.570947in}{1.685587in}}%
\pgfpathlineto{\pgfqpoint{2.575489in}{1.685587in}}%
\pgfpathlineto{\pgfqpoint{2.575489in}{1.682638in}}%
\pgfpathmoveto{\pgfqpoint{2.575489in}{1.682638in}}%
\pgfpathlineto{\pgfqpoint{2.575489in}{1.682638in}}%
\pgfpathlineto{\pgfqpoint{2.575489in}{1.685587in}}%
\pgfpathlineto{\pgfqpoint{2.580030in}{1.685587in}}%
\pgfpathlineto{\pgfqpoint{2.580030in}{1.682638in}}%
\pgfpathmoveto{\pgfqpoint{2.575489in}{1.685587in}}%
\pgfpathlineto{\pgfqpoint{2.575489in}{1.685587in}}%
\pgfpathlineto{\pgfqpoint{2.575489in}{1.688536in}}%
\pgfpathlineto{\pgfqpoint{2.580030in}{1.688536in}}%
\pgfpathlineto{\pgfqpoint{2.580030in}{1.685587in}}%
\pgfpathmoveto{\pgfqpoint{2.575489in}{1.688536in}}%
\pgfpathlineto{\pgfqpoint{2.575489in}{1.688536in}}%
\pgfpathlineto{\pgfqpoint{2.575489in}{1.691485in}}%
\pgfpathlineto{\pgfqpoint{2.580030in}{1.691485in}}%
\pgfpathlineto{\pgfqpoint{2.580030in}{1.688536in}}%
\pgfpathmoveto{\pgfqpoint{2.580030in}{1.685587in}}%
\pgfpathlineto{\pgfqpoint{2.580030in}{1.685587in}}%
\pgfpathlineto{\pgfqpoint{2.580030in}{1.688536in}}%
\pgfpathlineto{\pgfqpoint{2.584571in}{1.688536in}}%
\pgfpathlineto{\pgfqpoint{2.584571in}{1.685587in}}%
\pgfpathmoveto{\pgfqpoint{2.580030in}{1.688536in}}%
\pgfpathlineto{\pgfqpoint{2.580030in}{1.688536in}}%
\pgfpathlineto{\pgfqpoint{2.580030in}{1.691485in}}%
\pgfpathlineto{\pgfqpoint{2.584571in}{1.691485in}}%
\pgfpathlineto{\pgfqpoint{2.584571in}{1.688536in}}%
\pgfpathmoveto{\pgfqpoint{2.584571in}{1.688536in}}%
\pgfpathlineto{\pgfqpoint{2.584571in}{1.688536in}}%
\pgfpathlineto{\pgfqpoint{2.584571in}{1.691485in}}%
\pgfpathlineto{\pgfqpoint{2.589112in}{1.691485in}}%
\pgfpathlineto{\pgfqpoint{2.589112in}{1.688536in}}%
\pgfpathmoveto{\pgfqpoint{2.584571in}{1.691485in}}%
\pgfpathlineto{\pgfqpoint{2.584571in}{1.691485in}}%
\pgfpathlineto{\pgfqpoint{2.584571in}{1.694435in}}%
\pgfpathlineto{\pgfqpoint{2.589112in}{1.694435in}}%
\pgfpathlineto{\pgfqpoint{2.589112in}{1.691485in}}%
\pgfpathmoveto{\pgfqpoint{2.584571in}{1.694435in}}%
\pgfpathlineto{\pgfqpoint{2.584571in}{1.694435in}}%
\pgfpathlineto{\pgfqpoint{2.584571in}{1.697384in}}%
\pgfpathlineto{\pgfqpoint{2.589112in}{1.697384in}}%
\pgfpathlineto{\pgfqpoint{2.589112in}{1.694435in}}%
\pgfpathmoveto{\pgfqpoint{2.589112in}{1.691485in}}%
\pgfpathlineto{\pgfqpoint{2.589112in}{1.691485in}}%
\pgfpathlineto{\pgfqpoint{2.589112in}{1.694435in}}%
\pgfpathlineto{\pgfqpoint{2.593653in}{1.694435in}}%
\pgfpathlineto{\pgfqpoint{2.593653in}{1.691485in}}%
\pgfpathmoveto{\pgfqpoint{2.589112in}{1.694435in}}%
\pgfpathlineto{\pgfqpoint{2.589112in}{1.694435in}}%
\pgfpathlineto{\pgfqpoint{2.589112in}{1.697384in}}%
\pgfpathlineto{\pgfqpoint{2.593653in}{1.697384in}}%
\pgfpathlineto{\pgfqpoint{2.593653in}{1.694435in}}%
\pgfpathmoveto{\pgfqpoint{2.593653in}{1.694435in}}%
\pgfpathlineto{\pgfqpoint{2.593653in}{1.694435in}}%
\pgfpathlineto{\pgfqpoint{2.593653in}{1.697384in}}%
\pgfpathlineto{\pgfqpoint{2.598194in}{1.697384in}}%
\pgfpathlineto{\pgfqpoint{2.598194in}{1.694435in}}%
\pgfpathmoveto{\pgfqpoint{2.593653in}{1.697384in}}%
\pgfpathlineto{\pgfqpoint{2.593653in}{1.697384in}}%
\pgfpathlineto{\pgfqpoint{2.593653in}{1.700333in}}%
\pgfpathlineto{\pgfqpoint{2.598194in}{1.700333in}}%
\pgfpathlineto{\pgfqpoint{2.598194in}{1.697384in}}%
\pgfpathmoveto{\pgfqpoint{2.593653in}{1.700333in}}%
\pgfpathlineto{\pgfqpoint{2.593653in}{1.700333in}}%
\pgfpathlineto{\pgfqpoint{2.593653in}{1.703282in}}%
\pgfpathlineto{\pgfqpoint{2.598194in}{1.703282in}}%
\pgfpathlineto{\pgfqpoint{2.598194in}{1.700333in}}%
\pgfpathmoveto{\pgfqpoint{2.598194in}{1.697384in}}%
\pgfpathlineto{\pgfqpoint{2.598194in}{1.697384in}}%
\pgfpathlineto{\pgfqpoint{2.598194in}{1.700333in}}%
\pgfpathlineto{\pgfqpoint{2.602735in}{1.700333in}}%
\pgfpathlineto{\pgfqpoint{2.602735in}{1.697384in}}%
\pgfpathmoveto{\pgfqpoint{2.598194in}{1.700333in}}%
\pgfpathlineto{\pgfqpoint{2.598194in}{1.700333in}}%
\pgfpathlineto{\pgfqpoint{2.598194in}{1.703282in}}%
\pgfpathlineto{\pgfqpoint{2.602735in}{1.703282in}}%
\pgfpathlineto{\pgfqpoint{2.602735in}{1.700333in}}%
\pgfpathmoveto{\pgfqpoint{2.602735in}{1.700333in}}%
\pgfpathlineto{\pgfqpoint{2.602735in}{1.700333in}}%
\pgfpathlineto{\pgfqpoint{2.602735in}{1.703282in}}%
\pgfpathlineto{\pgfqpoint{2.607276in}{1.703282in}}%
\pgfpathlineto{\pgfqpoint{2.607276in}{1.700333in}}%
\pgfpathmoveto{\pgfqpoint{2.602735in}{1.703282in}}%
\pgfpathlineto{\pgfqpoint{2.602735in}{1.703282in}}%
\pgfpathlineto{\pgfqpoint{2.602735in}{1.706231in}}%
\pgfpathlineto{\pgfqpoint{2.607276in}{1.706231in}}%
\pgfpathlineto{\pgfqpoint{2.607276in}{1.703282in}}%
\pgfpathmoveto{\pgfqpoint{2.602735in}{1.706231in}}%
\pgfpathlineto{\pgfqpoint{2.602735in}{1.706231in}}%
\pgfpathlineto{\pgfqpoint{2.602735in}{1.709181in}}%
\pgfpathlineto{\pgfqpoint{2.607276in}{1.709181in}}%
\pgfpathlineto{\pgfqpoint{2.607276in}{1.706231in}}%
\pgfpathmoveto{\pgfqpoint{2.607276in}{1.703282in}}%
\pgfpathlineto{\pgfqpoint{2.607276in}{1.703282in}}%
\pgfpathlineto{\pgfqpoint{2.607276in}{1.706231in}}%
\pgfpathlineto{\pgfqpoint{2.611817in}{1.706231in}}%
\pgfpathlineto{\pgfqpoint{2.611817in}{1.703282in}}%
\pgfpathmoveto{\pgfqpoint{2.607276in}{1.706231in}}%
\pgfpathlineto{\pgfqpoint{2.607276in}{1.706231in}}%
\pgfpathlineto{\pgfqpoint{2.607276in}{1.709181in}}%
\pgfpathlineto{\pgfqpoint{2.611817in}{1.709181in}}%
\pgfpathlineto{\pgfqpoint{2.611817in}{1.706231in}}%
\pgfpathmoveto{\pgfqpoint{2.611817in}{1.706231in}}%
\pgfpathlineto{\pgfqpoint{2.611817in}{1.706231in}}%
\pgfpathlineto{\pgfqpoint{2.611817in}{1.709181in}}%
\pgfpathlineto{\pgfqpoint{2.616358in}{1.709181in}}%
\pgfpathlineto{\pgfqpoint{2.616358in}{1.706231in}}%
\pgfpathmoveto{\pgfqpoint{2.611817in}{1.709181in}}%
\pgfpathlineto{\pgfqpoint{2.611817in}{1.709181in}}%
\pgfpathlineto{\pgfqpoint{2.611817in}{1.712130in}}%
\pgfpathlineto{\pgfqpoint{2.616358in}{1.712130in}}%
\pgfpathlineto{\pgfqpoint{2.616358in}{1.709181in}}%
\pgfpathmoveto{\pgfqpoint{2.611817in}{1.712130in}}%
\pgfpathlineto{\pgfqpoint{2.611817in}{1.712130in}}%
\pgfpathlineto{\pgfqpoint{2.611817in}{1.715079in}}%
\pgfpathlineto{\pgfqpoint{2.616358in}{1.715079in}}%
\pgfpathlineto{\pgfqpoint{2.616358in}{1.712130in}}%
\pgfpathmoveto{\pgfqpoint{2.616358in}{1.709181in}}%
\pgfpathlineto{\pgfqpoint{2.616358in}{1.709181in}}%
\pgfpathlineto{\pgfqpoint{2.616358in}{1.712130in}}%
\pgfpathlineto{\pgfqpoint{2.620899in}{1.712130in}}%
\pgfpathlineto{\pgfqpoint{2.620899in}{1.709181in}}%
\pgfpathmoveto{\pgfqpoint{2.616358in}{1.712130in}}%
\pgfpathlineto{\pgfqpoint{2.616358in}{1.712130in}}%
\pgfpathlineto{\pgfqpoint{2.616358in}{1.715079in}}%
\pgfpathlineto{\pgfqpoint{2.620899in}{1.715079in}}%
\pgfpathlineto{\pgfqpoint{2.620899in}{1.712130in}}%
\pgfpathmoveto{\pgfqpoint{2.620899in}{1.712130in}}%
\pgfpathlineto{\pgfqpoint{2.620899in}{1.712130in}}%
\pgfpathlineto{\pgfqpoint{2.620899in}{1.715079in}}%
\pgfpathlineto{\pgfqpoint{2.625440in}{1.715079in}}%
\pgfpathlineto{\pgfqpoint{2.625440in}{1.712130in}}%
\pgfpathmoveto{\pgfqpoint{2.620899in}{1.715079in}}%
\pgfpathlineto{\pgfqpoint{2.620899in}{1.715079in}}%
\pgfpathlineto{\pgfqpoint{2.620899in}{1.718028in}}%
\pgfpathlineto{\pgfqpoint{2.625440in}{1.718028in}}%
\pgfpathlineto{\pgfqpoint{2.625440in}{1.715079in}}%
\pgfpathmoveto{\pgfqpoint{2.620899in}{1.718028in}}%
\pgfpathlineto{\pgfqpoint{2.620899in}{1.718028in}}%
\pgfpathlineto{\pgfqpoint{2.620899in}{1.720977in}}%
\pgfpathlineto{\pgfqpoint{2.625440in}{1.720977in}}%
\pgfpathlineto{\pgfqpoint{2.625440in}{1.718028in}}%
\pgfpathmoveto{\pgfqpoint{2.625440in}{1.715079in}}%
\pgfpathlineto{\pgfqpoint{2.625440in}{1.715079in}}%
\pgfpathlineto{\pgfqpoint{2.625440in}{1.718028in}}%
\pgfpathlineto{\pgfqpoint{2.629981in}{1.718028in}}%
\pgfpathlineto{\pgfqpoint{2.629981in}{1.715079in}}%
\pgfpathmoveto{\pgfqpoint{2.625440in}{1.718028in}}%
\pgfpathlineto{\pgfqpoint{2.625440in}{1.718028in}}%
\pgfpathlineto{\pgfqpoint{2.625440in}{1.720977in}}%
\pgfpathlineto{\pgfqpoint{2.629981in}{1.720977in}}%
\pgfpathlineto{\pgfqpoint{2.629981in}{1.718028in}}%
\pgfpathmoveto{\pgfqpoint{2.629981in}{1.718028in}}%
\pgfpathlineto{\pgfqpoint{2.629981in}{1.718028in}}%
\pgfpathlineto{\pgfqpoint{2.629981in}{1.720977in}}%
\pgfpathlineto{\pgfqpoint{2.634522in}{1.720977in}}%
\pgfpathlineto{\pgfqpoint{2.634522in}{1.718028in}}%
\pgfpathmoveto{\pgfqpoint{2.629981in}{1.720977in}}%
\pgfpathlineto{\pgfqpoint{2.629981in}{1.720977in}}%
\pgfpathlineto{\pgfqpoint{2.629981in}{1.723927in}}%
\pgfpathlineto{\pgfqpoint{2.634522in}{1.723927in}}%
\pgfpathlineto{\pgfqpoint{2.634522in}{1.720977in}}%
\pgfpathmoveto{\pgfqpoint{2.629981in}{1.723927in}}%
\pgfpathlineto{\pgfqpoint{2.629981in}{1.723927in}}%
\pgfpathlineto{\pgfqpoint{2.629981in}{1.726876in}}%
\pgfpathlineto{\pgfqpoint{2.634522in}{1.726876in}}%
\pgfpathlineto{\pgfqpoint{2.634522in}{1.723927in}}%
\pgfpathmoveto{\pgfqpoint{2.634522in}{1.720977in}}%
\pgfpathlineto{\pgfqpoint{2.634522in}{1.720977in}}%
\pgfpathlineto{\pgfqpoint{2.634522in}{1.723927in}}%
\pgfpathlineto{\pgfqpoint{2.639063in}{1.723927in}}%
\pgfpathlineto{\pgfqpoint{2.639063in}{1.720977in}}%
\pgfpathmoveto{\pgfqpoint{2.634522in}{1.723927in}}%
\pgfpathlineto{\pgfqpoint{2.634522in}{1.723927in}}%
\pgfpathlineto{\pgfqpoint{2.634522in}{1.726876in}}%
\pgfpathlineto{\pgfqpoint{2.639063in}{1.726876in}}%
\pgfpathlineto{\pgfqpoint{2.639063in}{1.723927in}}%
\pgfpathmoveto{\pgfqpoint{2.639063in}{1.723927in}}%
\pgfpathlineto{\pgfqpoint{2.639063in}{1.723927in}}%
\pgfpathlineto{\pgfqpoint{2.639063in}{1.726876in}}%
\pgfpathlineto{\pgfqpoint{2.643604in}{1.726876in}}%
\pgfpathlineto{\pgfqpoint{2.643604in}{1.723927in}}%
\pgfpathmoveto{\pgfqpoint{2.639063in}{1.726876in}}%
\pgfpathlineto{\pgfqpoint{2.639063in}{1.726876in}}%
\pgfpathlineto{\pgfqpoint{2.639063in}{1.729825in}}%
\pgfpathlineto{\pgfqpoint{2.643604in}{1.729825in}}%
\pgfpathlineto{\pgfqpoint{2.643604in}{1.726876in}}%
\pgfpathmoveto{\pgfqpoint{2.639063in}{1.729825in}}%
\pgfpathlineto{\pgfqpoint{2.639063in}{1.729825in}}%
\pgfpathlineto{\pgfqpoint{2.639063in}{1.732774in}}%
\pgfpathlineto{\pgfqpoint{2.643604in}{1.732774in}}%
\pgfpathlineto{\pgfqpoint{2.643604in}{1.729825in}}%
\pgfpathmoveto{\pgfqpoint{2.643604in}{1.726876in}}%
\pgfpathlineto{\pgfqpoint{2.643604in}{1.726876in}}%
\pgfpathlineto{\pgfqpoint{2.643604in}{1.729825in}}%
\pgfpathlineto{\pgfqpoint{2.648145in}{1.729825in}}%
\pgfpathlineto{\pgfqpoint{2.648145in}{1.726876in}}%
\pgfpathmoveto{\pgfqpoint{2.643604in}{1.729825in}}%
\pgfpathlineto{\pgfqpoint{2.643604in}{1.729825in}}%
\pgfpathlineto{\pgfqpoint{2.643604in}{1.732774in}}%
\pgfpathlineto{\pgfqpoint{2.648145in}{1.732774in}}%
\pgfpathlineto{\pgfqpoint{2.648145in}{1.729825in}}%
\pgfpathmoveto{\pgfqpoint{2.648145in}{1.729825in}}%
\pgfpathlineto{\pgfqpoint{2.648145in}{1.729825in}}%
\pgfpathlineto{\pgfqpoint{2.648145in}{1.732774in}}%
\pgfpathlineto{\pgfqpoint{2.652686in}{1.732774in}}%
\pgfpathlineto{\pgfqpoint{2.652686in}{1.729825in}}%
\pgfpathmoveto{\pgfqpoint{2.648145in}{1.732774in}}%
\pgfpathlineto{\pgfqpoint{2.648145in}{1.732774in}}%
\pgfpathlineto{\pgfqpoint{2.648145in}{1.735723in}}%
\pgfpathlineto{\pgfqpoint{2.652686in}{1.735723in}}%
\pgfpathlineto{\pgfqpoint{2.652686in}{1.732774in}}%
\pgfpathmoveto{\pgfqpoint{2.648145in}{1.735723in}}%
\pgfpathlineto{\pgfqpoint{2.648145in}{1.735723in}}%
\pgfpathlineto{\pgfqpoint{2.648145in}{1.738672in}}%
\pgfpathlineto{\pgfqpoint{2.652686in}{1.738672in}}%
\pgfpathlineto{\pgfqpoint{2.652686in}{1.735723in}}%
\pgfpathmoveto{\pgfqpoint{2.652686in}{1.732774in}}%
\pgfpathlineto{\pgfqpoint{2.652686in}{1.732774in}}%
\pgfpathlineto{\pgfqpoint{2.652686in}{1.735723in}}%
\pgfpathlineto{\pgfqpoint{2.657227in}{1.735723in}}%
\pgfpathlineto{\pgfqpoint{2.657227in}{1.732774in}}%
\pgfpathmoveto{\pgfqpoint{2.652686in}{1.735723in}}%
\pgfpathlineto{\pgfqpoint{2.652686in}{1.735723in}}%
\pgfpathlineto{\pgfqpoint{2.652686in}{1.738672in}}%
\pgfpathlineto{\pgfqpoint{2.657227in}{1.738672in}}%
\pgfpathlineto{\pgfqpoint{2.657227in}{1.735723in}}%
\pgfpathmoveto{\pgfqpoint{2.657227in}{1.735723in}}%
\pgfpathlineto{\pgfqpoint{2.657227in}{1.735723in}}%
\pgfpathlineto{\pgfqpoint{2.657227in}{1.738672in}}%
\pgfpathlineto{\pgfqpoint{2.661768in}{1.738672in}}%
\pgfpathlineto{\pgfqpoint{2.661768in}{1.735723in}}%
\pgfpathmoveto{\pgfqpoint{2.657227in}{1.738672in}}%
\pgfpathlineto{\pgfqpoint{2.657227in}{1.738672in}}%
\pgfpathlineto{\pgfqpoint{2.657227in}{1.741622in}}%
\pgfpathlineto{\pgfqpoint{2.661768in}{1.741622in}}%
\pgfpathlineto{\pgfqpoint{2.661768in}{1.738672in}}%
\pgfpathmoveto{\pgfqpoint{2.657227in}{1.741622in}}%
\pgfpathlineto{\pgfqpoint{2.657227in}{1.741622in}}%
\pgfpathlineto{\pgfqpoint{2.657227in}{1.744571in}}%
\pgfpathlineto{\pgfqpoint{2.661768in}{1.744571in}}%
\pgfpathlineto{\pgfqpoint{2.661768in}{1.741622in}}%
\pgfpathmoveto{\pgfqpoint{2.661768in}{1.738672in}}%
\pgfpathlineto{\pgfqpoint{2.661768in}{1.738672in}}%
\pgfpathlineto{\pgfqpoint{2.661768in}{1.741622in}}%
\pgfpathlineto{\pgfqpoint{2.666309in}{1.741622in}}%
\pgfpathlineto{\pgfqpoint{2.666309in}{1.738672in}}%
\pgfpathmoveto{\pgfqpoint{2.661768in}{1.741622in}}%
\pgfpathlineto{\pgfqpoint{2.661768in}{1.741622in}}%
\pgfpathlineto{\pgfqpoint{2.661768in}{1.744571in}}%
\pgfpathlineto{\pgfqpoint{2.666309in}{1.744571in}}%
\pgfpathlineto{\pgfqpoint{2.666309in}{1.741622in}}%
\pgfpathmoveto{\pgfqpoint{2.666309in}{1.741622in}}%
\pgfpathlineto{\pgfqpoint{2.666309in}{1.741622in}}%
\pgfpathlineto{\pgfqpoint{2.666309in}{1.744571in}}%
\pgfpathlineto{\pgfqpoint{2.670850in}{1.744571in}}%
\pgfpathlineto{\pgfqpoint{2.670850in}{1.741622in}}%
\pgfpathmoveto{\pgfqpoint{2.666309in}{1.744571in}}%
\pgfpathlineto{\pgfqpoint{2.666309in}{1.744571in}}%
\pgfpathlineto{\pgfqpoint{2.666309in}{1.747520in}}%
\pgfpathlineto{\pgfqpoint{2.670850in}{1.747520in}}%
\pgfpathlineto{\pgfqpoint{2.670850in}{1.744571in}}%
\pgfpathmoveto{\pgfqpoint{2.666309in}{1.747520in}}%
\pgfpathlineto{\pgfqpoint{2.666309in}{1.747520in}}%
\pgfpathlineto{\pgfqpoint{2.666309in}{1.750469in}}%
\pgfpathlineto{\pgfqpoint{2.670850in}{1.750469in}}%
\pgfpathlineto{\pgfqpoint{2.670850in}{1.747520in}}%
\pgfpathmoveto{\pgfqpoint{2.670850in}{1.744571in}}%
\pgfpathlineto{\pgfqpoint{2.670850in}{1.744571in}}%
\pgfpathlineto{\pgfqpoint{2.670850in}{1.747520in}}%
\pgfpathlineto{\pgfqpoint{2.675391in}{1.747520in}}%
\pgfpathlineto{\pgfqpoint{2.675391in}{1.744571in}}%
\pgfpathmoveto{\pgfqpoint{2.670850in}{1.747520in}}%
\pgfpathlineto{\pgfqpoint{2.670850in}{1.747520in}}%
\pgfpathlineto{\pgfqpoint{2.670850in}{1.750469in}}%
\pgfpathlineto{\pgfqpoint{2.675391in}{1.750469in}}%
\pgfpathlineto{\pgfqpoint{2.675391in}{1.747520in}}%
\pgfpathmoveto{\pgfqpoint{2.666309in}{1.750469in}}%
\pgfpathlineto{\pgfqpoint{2.666309in}{1.750469in}}%
\pgfpathlineto{\pgfqpoint{2.666309in}{1.753418in}}%
\pgfpathlineto{\pgfqpoint{2.670850in}{1.753418in}}%
\pgfpathlineto{\pgfqpoint{2.670850in}{1.750469in}}%
\pgfpathmoveto{\pgfqpoint{2.666309in}{1.753418in}}%
\pgfpathlineto{\pgfqpoint{2.666309in}{1.753418in}}%
\pgfpathlineto{\pgfqpoint{2.666309in}{1.756367in}}%
\pgfpathlineto{\pgfqpoint{2.670850in}{1.756367in}}%
\pgfpathlineto{\pgfqpoint{2.670850in}{1.753418in}}%
\pgfpathmoveto{\pgfqpoint{2.670850in}{1.750469in}}%
\pgfpathlineto{\pgfqpoint{2.670850in}{1.750469in}}%
\pgfpathlineto{\pgfqpoint{2.670850in}{1.753418in}}%
\pgfpathlineto{\pgfqpoint{2.675391in}{1.753418in}}%
\pgfpathlineto{\pgfqpoint{2.675391in}{1.750469in}}%
\pgfpathmoveto{\pgfqpoint{2.670850in}{1.753418in}}%
\pgfpathlineto{\pgfqpoint{2.670850in}{1.753418in}}%
\pgfpathlineto{\pgfqpoint{2.670850in}{1.756367in}}%
\pgfpathlineto{\pgfqpoint{2.675391in}{1.756367in}}%
\pgfpathlineto{\pgfqpoint{2.675391in}{1.753418in}}%
\pgfpathmoveto{\pgfqpoint{2.675391in}{1.750469in}}%
\pgfpathlineto{\pgfqpoint{2.675391in}{1.750469in}}%
\pgfpathlineto{\pgfqpoint{2.675391in}{1.753418in}}%
\pgfpathlineto{\pgfqpoint{2.679932in}{1.753418in}}%
\pgfpathlineto{\pgfqpoint{2.679932in}{1.750469in}}%
\pgfpathmoveto{\pgfqpoint{2.675391in}{1.753418in}}%
\pgfpathlineto{\pgfqpoint{2.675391in}{1.753418in}}%
\pgfpathlineto{\pgfqpoint{2.675391in}{1.756367in}}%
\pgfpathlineto{\pgfqpoint{2.679932in}{1.756367in}}%
\pgfpathlineto{\pgfqpoint{2.679932in}{1.753418in}}%
\pgfpathmoveto{\pgfqpoint{2.679932in}{1.753418in}}%
\pgfpathlineto{\pgfqpoint{2.679932in}{1.753418in}}%
\pgfpathlineto{\pgfqpoint{2.679932in}{1.756367in}}%
\pgfpathlineto{\pgfqpoint{2.684473in}{1.756367in}}%
\pgfpathlineto{\pgfqpoint{2.684473in}{1.753418in}}%
\pgfpathmoveto{\pgfqpoint{2.675391in}{1.756367in}}%
\pgfpathlineto{\pgfqpoint{2.675391in}{1.756367in}}%
\pgfpathlineto{\pgfqpoint{2.675391in}{1.759316in}}%
\pgfpathlineto{\pgfqpoint{2.679932in}{1.759316in}}%
\pgfpathlineto{\pgfqpoint{2.679932in}{1.756367in}}%
\pgfpathmoveto{\pgfqpoint{2.675391in}{1.759316in}}%
\pgfpathlineto{\pgfqpoint{2.675391in}{1.759316in}}%
\pgfpathlineto{\pgfqpoint{2.675391in}{1.762266in}}%
\pgfpathlineto{\pgfqpoint{2.679932in}{1.762266in}}%
\pgfpathlineto{\pgfqpoint{2.679932in}{1.759316in}}%
\pgfpathmoveto{\pgfqpoint{2.679932in}{1.756367in}}%
\pgfpathlineto{\pgfqpoint{2.679932in}{1.756367in}}%
\pgfpathlineto{\pgfqpoint{2.679932in}{1.759316in}}%
\pgfpathlineto{\pgfqpoint{2.684473in}{1.759316in}}%
\pgfpathlineto{\pgfqpoint{2.684473in}{1.756367in}}%
\pgfpathmoveto{\pgfqpoint{2.679932in}{1.759316in}}%
\pgfpathlineto{\pgfqpoint{2.679932in}{1.759316in}}%
\pgfpathlineto{\pgfqpoint{2.679932in}{1.762266in}}%
\pgfpathlineto{\pgfqpoint{2.684473in}{1.762266in}}%
\pgfpathlineto{\pgfqpoint{2.684473in}{1.759316in}}%
\pgfpathmoveto{\pgfqpoint{2.684473in}{1.756367in}}%
\pgfpathlineto{\pgfqpoint{2.684473in}{1.756367in}}%
\pgfpathlineto{\pgfqpoint{2.684473in}{1.759316in}}%
\pgfpathlineto{\pgfqpoint{2.689014in}{1.759316in}}%
\pgfpathlineto{\pgfqpoint{2.689014in}{1.756367in}}%
\pgfpathmoveto{\pgfqpoint{2.684473in}{1.759316in}}%
\pgfpathlineto{\pgfqpoint{2.684473in}{1.759316in}}%
\pgfpathlineto{\pgfqpoint{2.684473in}{1.762266in}}%
\pgfpathlineto{\pgfqpoint{2.689014in}{1.762266in}}%
\pgfpathlineto{\pgfqpoint{2.689014in}{1.759316in}}%
\pgfpathmoveto{\pgfqpoint{2.689014in}{1.759316in}}%
\pgfpathlineto{\pgfqpoint{2.689014in}{1.759316in}}%
\pgfpathlineto{\pgfqpoint{2.689014in}{1.762266in}}%
\pgfpathlineto{\pgfqpoint{2.693555in}{1.762266in}}%
\pgfpathlineto{\pgfqpoint{2.693555in}{1.759316in}}%
\pgfpathmoveto{\pgfqpoint{2.684473in}{1.762266in}}%
\pgfpathlineto{\pgfqpoint{2.684473in}{1.762266in}}%
\pgfpathlineto{\pgfqpoint{2.684473in}{1.765215in}}%
\pgfpathlineto{\pgfqpoint{2.689014in}{1.765215in}}%
\pgfpathlineto{\pgfqpoint{2.689014in}{1.762266in}}%
\pgfpathmoveto{\pgfqpoint{2.684473in}{1.765215in}}%
\pgfpathlineto{\pgfqpoint{2.684473in}{1.765215in}}%
\pgfpathlineto{\pgfqpoint{2.684473in}{1.768164in}}%
\pgfpathlineto{\pgfqpoint{2.689014in}{1.768164in}}%
\pgfpathlineto{\pgfqpoint{2.689014in}{1.765215in}}%
\pgfpathmoveto{\pgfqpoint{2.689014in}{1.762266in}}%
\pgfpathlineto{\pgfqpoint{2.689014in}{1.762266in}}%
\pgfpathlineto{\pgfqpoint{2.689014in}{1.765215in}}%
\pgfpathlineto{\pgfqpoint{2.693555in}{1.765215in}}%
\pgfpathlineto{\pgfqpoint{2.693555in}{1.762266in}}%
\pgfpathmoveto{\pgfqpoint{2.689014in}{1.765215in}}%
\pgfpathlineto{\pgfqpoint{2.689014in}{1.765215in}}%
\pgfpathlineto{\pgfqpoint{2.689014in}{1.768164in}}%
\pgfpathlineto{\pgfqpoint{2.693555in}{1.768164in}}%
\pgfpathlineto{\pgfqpoint{2.693555in}{1.765215in}}%
\pgfpathmoveto{\pgfqpoint{2.693555in}{1.762266in}}%
\pgfpathlineto{\pgfqpoint{2.693555in}{1.762266in}}%
\pgfpathlineto{\pgfqpoint{2.693555in}{1.765215in}}%
\pgfpathlineto{\pgfqpoint{2.698096in}{1.765215in}}%
\pgfpathlineto{\pgfqpoint{2.698096in}{1.762266in}}%
\pgfpathmoveto{\pgfqpoint{2.693555in}{1.765215in}}%
\pgfpathlineto{\pgfqpoint{2.693555in}{1.765215in}}%
\pgfpathlineto{\pgfqpoint{2.693555in}{1.768164in}}%
\pgfpathlineto{\pgfqpoint{2.698096in}{1.768164in}}%
\pgfpathlineto{\pgfqpoint{2.698096in}{1.765215in}}%
\pgfpathmoveto{\pgfqpoint{2.698096in}{1.765215in}}%
\pgfpathlineto{\pgfqpoint{2.698096in}{1.765215in}}%
\pgfpathlineto{\pgfqpoint{2.698096in}{1.768164in}}%
\pgfpathlineto{\pgfqpoint{2.702637in}{1.768164in}}%
\pgfpathlineto{\pgfqpoint{2.702637in}{1.765215in}}%
\pgfpathmoveto{\pgfqpoint{2.693555in}{1.768164in}}%
\pgfpathlineto{\pgfqpoint{2.693555in}{1.768164in}}%
\pgfpathlineto{\pgfqpoint{2.693555in}{1.771113in}}%
\pgfpathlineto{\pgfqpoint{2.698096in}{1.771113in}}%
\pgfpathlineto{\pgfqpoint{2.698096in}{1.768164in}}%
\pgfpathmoveto{\pgfqpoint{2.693555in}{1.771113in}}%
\pgfpathlineto{\pgfqpoint{2.693555in}{1.771113in}}%
\pgfpathlineto{\pgfqpoint{2.693555in}{1.774062in}}%
\pgfpathlineto{\pgfqpoint{2.698096in}{1.774062in}}%
\pgfpathlineto{\pgfqpoint{2.698096in}{1.771113in}}%
\pgfpathmoveto{\pgfqpoint{2.698096in}{1.768164in}}%
\pgfpathlineto{\pgfqpoint{2.698096in}{1.768164in}}%
\pgfpathlineto{\pgfqpoint{2.698096in}{1.771113in}}%
\pgfpathlineto{\pgfqpoint{2.702637in}{1.771113in}}%
\pgfpathlineto{\pgfqpoint{2.702637in}{1.768164in}}%
\pgfpathmoveto{\pgfqpoint{2.698096in}{1.771113in}}%
\pgfpathlineto{\pgfqpoint{2.698096in}{1.771113in}}%
\pgfpathlineto{\pgfqpoint{2.698096in}{1.774062in}}%
\pgfpathlineto{\pgfqpoint{2.702637in}{1.774062in}}%
\pgfpathlineto{\pgfqpoint{2.702637in}{1.771113in}}%
\pgfpathmoveto{\pgfqpoint{2.702637in}{1.768164in}}%
\pgfpathlineto{\pgfqpoint{2.702637in}{1.768164in}}%
\pgfpathlineto{\pgfqpoint{2.702637in}{1.771113in}}%
\pgfpathlineto{\pgfqpoint{2.707178in}{1.771113in}}%
\pgfpathlineto{\pgfqpoint{2.707178in}{1.768164in}}%
\pgfpathmoveto{\pgfqpoint{2.702637in}{1.771113in}}%
\pgfpathlineto{\pgfqpoint{2.702637in}{1.771113in}}%
\pgfpathlineto{\pgfqpoint{2.702637in}{1.774062in}}%
\pgfpathlineto{\pgfqpoint{2.707178in}{1.774062in}}%
\pgfpathlineto{\pgfqpoint{2.707178in}{1.771113in}}%
\pgfpathmoveto{\pgfqpoint{2.707178in}{1.771113in}}%
\pgfpathlineto{\pgfqpoint{2.707178in}{1.771113in}}%
\pgfpathlineto{\pgfqpoint{2.707178in}{1.774062in}}%
\pgfpathlineto{\pgfqpoint{2.711719in}{1.774062in}}%
\pgfpathlineto{\pgfqpoint{2.711719in}{1.771113in}}%
\pgfpathmoveto{\pgfqpoint{2.702637in}{1.774062in}}%
\pgfpathlineto{\pgfqpoint{2.702637in}{1.774062in}}%
\pgfpathlineto{\pgfqpoint{2.702637in}{1.777011in}}%
\pgfpathlineto{\pgfqpoint{2.707178in}{1.777011in}}%
\pgfpathlineto{\pgfqpoint{2.707178in}{1.774062in}}%
\pgfpathmoveto{\pgfqpoint{2.702637in}{1.777011in}}%
\pgfpathlineto{\pgfqpoint{2.702637in}{1.777011in}}%
\pgfpathlineto{\pgfqpoint{2.702637in}{1.779961in}}%
\pgfpathlineto{\pgfqpoint{2.707178in}{1.779961in}}%
\pgfpathlineto{\pgfqpoint{2.707178in}{1.777011in}}%
\pgfpathmoveto{\pgfqpoint{2.707178in}{1.774062in}}%
\pgfpathlineto{\pgfqpoint{2.707178in}{1.774062in}}%
\pgfpathlineto{\pgfqpoint{2.707178in}{1.777011in}}%
\pgfpathlineto{\pgfqpoint{2.711719in}{1.777011in}}%
\pgfpathlineto{\pgfqpoint{2.711719in}{1.774062in}}%
\pgfpathmoveto{\pgfqpoint{2.707178in}{1.777011in}}%
\pgfpathlineto{\pgfqpoint{2.707178in}{1.777011in}}%
\pgfpathlineto{\pgfqpoint{2.707178in}{1.779961in}}%
\pgfpathlineto{\pgfqpoint{2.711719in}{1.779961in}}%
\pgfpathlineto{\pgfqpoint{2.711719in}{1.777011in}}%
\pgfpathmoveto{\pgfqpoint{2.711719in}{1.774062in}}%
\pgfpathlineto{\pgfqpoint{2.711719in}{1.774062in}}%
\pgfpathlineto{\pgfqpoint{2.711719in}{1.777011in}}%
\pgfpathlineto{\pgfqpoint{2.716260in}{1.777011in}}%
\pgfpathlineto{\pgfqpoint{2.716260in}{1.774062in}}%
\pgfpathmoveto{\pgfqpoint{2.711719in}{1.777011in}}%
\pgfpathlineto{\pgfqpoint{2.711719in}{1.777011in}}%
\pgfpathlineto{\pgfqpoint{2.711719in}{1.779961in}}%
\pgfpathlineto{\pgfqpoint{2.716260in}{1.779961in}}%
\pgfpathlineto{\pgfqpoint{2.716260in}{1.777011in}}%
\pgfpathmoveto{\pgfqpoint{2.716260in}{1.777011in}}%
\pgfpathlineto{\pgfqpoint{2.716260in}{1.777011in}}%
\pgfpathlineto{\pgfqpoint{2.716260in}{1.779961in}}%
\pgfpathlineto{\pgfqpoint{2.720801in}{1.779961in}}%
\pgfpathlineto{\pgfqpoint{2.720801in}{1.777011in}}%
\pgfpathmoveto{\pgfqpoint{2.711719in}{1.779961in}}%
\pgfpathlineto{\pgfqpoint{2.711719in}{1.779961in}}%
\pgfpathlineto{\pgfqpoint{2.711719in}{1.782910in}}%
\pgfpathlineto{\pgfqpoint{2.716260in}{1.782910in}}%
\pgfpathlineto{\pgfqpoint{2.716260in}{1.779961in}}%
\pgfpathmoveto{\pgfqpoint{2.711719in}{1.782910in}}%
\pgfpathlineto{\pgfqpoint{2.711719in}{1.782910in}}%
\pgfpathlineto{\pgfqpoint{2.711719in}{1.785859in}}%
\pgfpathlineto{\pgfqpoint{2.716260in}{1.785859in}}%
\pgfpathlineto{\pgfqpoint{2.716260in}{1.782910in}}%
\pgfpathmoveto{\pgfqpoint{2.716260in}{1.779961in}}%
\pgfpathlineto{\pgfqpoint{2.716260in}{1.779961in}}%
\pgfpathlineto{\pgfqpoint{2.716260in}{1.782910in}}%
\pgfpathlineto{\pgfqpoint{2.720801in}{1.782910in}}%
\pgfpathlineto{\pgfqpoint{2.720801in}{1.779961in}}%
\pgfpathmoveto{\pgfqpoint{2.716260in}{1.782910in}}%
\pgfpathlineto{\pgfqpoint{2.716260in}{1.782910in}}%
\pgfpathlineto{\pgfqpoint{2.716260in}{1.785859in}}%
\pgfpathlineto{\pgfqpoint{2.720801in}{1.785859in}}%
\pgfpathlineto{\pgfqpoint{2.720801in}{1.782910in}}%
\pgfpathmoveto{\pgfqpoint{2.720801in}{1.779961in}}%
\pgfpathlineto{\pgfqpoint{2.720801in}{1.779961in}}%
\pgfpathlineto{\pgfqpoint{2.720801in}{1.782910in}}%
\pgfpathlineto{\pgfqpoint{2.725342in}{1.782910in}}%
\pgfpathlineto{\pgfqpoint{2.725342in}{1.779961in}}%
\pgfpathmoveto{\pgfqpoint{2.720801in}{1.782910in}}%
\pgfpathlineto{\pgfqpoint{2.720801in}{1.782910in}}%
\pgfpathlineto{\pgfqpoint{2.720801in}{1.785859in}}%
\pgfpathlineto{\pgfqpoint{2.725342in}{1.785859in}}%
\pgfpathlineto{\pgfqpoint{2.725342in}{1.782910in}}%
\pgfpathmoveto{\pgfqpoint{2.725342in}{1.782910in}}%
\pgfpathlineto{\pgfqpoint{2.725342in}{1.782910in}}%
\pgfpathlineto{\pgfqpoint{2.725342in}{1.785859in}}%
\pgfpathlineto{\pgfqpoint{2.729883in}{1.785859in}}%
\pgfpathlineto{\pgfqpoint{2.729883in}{1.782910in}}%
\pgfpathmoveto{\pgfqpoint{2.720801in}{1.785859in}}%
\pgfpathlineto{\pgfqpoint{2.720801in}{1.785859in}}%
\pgfpathlineto{\pgfqpoint{2.720801in}{1.788808in}}%
\pgfpathlineto{\pgfqpoint{2.725342in}{1.788808in}}%
\pgfpathlineto{\pgfqpoint{2.725342in}{1.785859in}}%
\pgfpathmoveto{\pgfqpoint{2.720801in}{1.788808in}}%
\pgfpathlineto{\pgfqpoint{2.720801in}{1.788808in}}%
\pgfpathlineto{\pgfqpoint{2.720801in}{1.791757in}}%
\pgfpathlineto{\pgfqpoint{2.725342in}{1.791757in}}%
\pgfpathlineto{\pgfqpoint{2.725342in}{1.788808in}}%
\pgfpathmoveto{\pgfqpoint{2.725342in}{1.785859in}}%
\pgfpathlineto{\pgfqpoint{2.725342in}{1.785859in}}%
\pgfpathlineto{\pgfqpoint{2.725342in}{1.788808in}}%
\pgfpathlineto{\pgfqpoint{2.729883in}{1.788808in}}%
\pgfpathlineto{\pgfqpoint{2.729883in}{1.785859in}}%
\pgfpathmoveto{\pgfqpoint{2.725342in}{1.788808in}}%
\pgfpathlineto{\pgfqpoint{2.725342in}{1.788808in}}%
\pgfpathlineto{\pgfqpoint{2.725342in}{1.791757in}}%
\pgfpathlineto{\pgfqpoint{2.729883in}{1.791757in}}%
\pgfpathlineto{\pgfqpoint{2.729883in}{1.788808in}}%
\pgfpathmoveto{\pgfqpoint{2.729883in}{1.785859in}}%
\pgfpathlineto{\pgfqpoint{2.729883in}{1.785859in}}%
\pgfpathlineto{\pgfqpoint{2.729883in}{1.788808in}}%
\pgfpathlineto{\pgfqpoint{2.734424in}{1.788808in}}%
\pgfpathlineto{\pgfqpoint{2.734424in}{1.785859in}}%
\pgfpathmoveto{\pgfqpoint{2.729883in}{1.788808in}}%
\pgfpathlineto{\pgfqpoint{2.729883in}{1.788808in}}%
\pgfpathlineto{\pgfqpoint{2.729883in}{1.791757in}}%
\pgfpathlineto{\pgfqpoint{2.734424in}{1.791757in}}%
\pgfpathlineto{\pgfqpoint{2.734424in}{1.788808in}}%
\pgfpathmoveto{\pgfqpoint{2.734424in}{1.788808in}}%
\pgfpathlineto{\pgfqpoint{2.734424in}{1.788808in}}%
\pgfpathlineto{\pgfqpoint{2.734424in}{1.791757in}}%
\pgfpathlineto{\pgfqpoint{2.738965in}{1.791757in}}%
\pgfpathlineto{\pgfqpoint{2.738965in}{1.788808in}}%
\pgfpathmoveto{\pgfqpoint{2.729883in}{1.791757in}}%
\pgfpathlineto{\pgfqpoint{2.729883in}{1.791757in}}%
\pgfpathlineto{\pgfqpoint{2.729883in}{1.794706in}}%
\pgfpathlineto{\pgfqpoint{2.734424in}{1.794706in}}%
\pgfpathlineto{\pgfqpoint{2.734424in}{1.791757in}}%
\pgfpathmoveto{\pgfqpoint{2.729883in}{1.794706in}}%
\pgfpathlineto{\pgfqpoint{2.729883in}{1.794706in}}%
\pgfpathlineto{\pgfqpoint{2.729883in}{1.797655in}}%
\pgfpathlineto{\pgfqpoint{2.734424in}{1.797655in}}%
\pgfpathlineto{\pgfqpoint{2.734424in}{1.794706in}}%
\pgfpathmoveto{\pgfqpoint{2.734424in}{1.791757in}}%
\pgfpathlineto{\pgfqpoint{2.734424in}{1.791757in}}%
\pgfpathlineto{\pgfqpoint{2.734424in}{1.794706in}}%
\pgfpathlineto{\pgfqpoint{2.738965in}{1.794706in}}%
\pgfpathlineto{\pgfqpoint{2.738965in}{1.791757in}}%
\pgfpathmoveto{\pgfqpoint{2.734424in}{1.794706in}}%
\pgfpathlineto{\pgfqpoint{2.734424in}{1.794706in}}%
\pgfpathlineto{\pgfqpoint{2.734424in}{1.797655in}}%
\pgfpathlineto{\pgfqpoint{2.738965in}{1.797655in}}%
\pgfpathlineto{\pgfqpoint{2.738965in}{1.794706in}}%
\pgfpathmoveto{\pgfqpoint{2.738965in}{1.791757in}}%
\pgfpathlineto{\pgfqpoint{2.738965in}{1.791757in}}%
\pgfpathlineto{\pgfqpoint{2.738965in}{1.794706in}}%
\pgfpathlineto{\pgfqpoint{2.743506in}{1.794706in}}%
\pgfpathlineto{\pgfqpoint{2.743506in}{1.791757in}}%
\pgfpathmoveto{\pgfqpoint{2.738965in}{1.794706in}}%
\pgfpathlineto{\pgfqpoint{2.738965in}{1.794706in}}%
\pgfpathlineto{\pgfqpoint{2.738965in}{1.797655in}}%
\pgfpathlineto{\pgfqpoint{2.743506in}{1.797655in}}%
\pgfpathlineto{\pgfqpoint{2.743506in}{1.794706in}}%
\pgfpathmoveto{\pgfqpoint{2.743506in}{1.794706in}}%
\pgfpathlineto{\pgfqpoint{2.743506in}{1.794706in}}%
\pgfpathlineto{\pgfqpoint{2.743506in}{1.797655in}}%
\pgfpathlineto{\pgfqpoint{2.748047in}{1.797655in}}%
\pgfpathlineto{\pgfqpoint{2.748047in}{1.794706in}}%
\pgfpathmoveto{\pgfqpoint{2.738965in}{1.797655in}}%
\pgfpathlineto{\pgfqpoint{2.738965in}{1.797655in}}%
\pgfpathlineto{\pgfqpoint{2.738965in}{1.800605in}}%
\pgfpathlineto{\pgfqpoint{2.743506in}{1.800605in}}%
\pgfpathlineto{\pgfqpoint{2.743506in}{1.797655in}}%
\pgfpathmoveto{\pgfqpoint{2.738965in}{1.800605in}}%
\pgfpathlineto{\pgfqpoint{2.738965in}{1.800605in}}%
\pgfpathlineto{\pgfqpoint{2.738965in}{1.803554in}}%
\pgfpathlineto{\pgfqpoint{2.743506in}{1.803554in}}%
\pgfpathlineto{\pgfqpoint{2.743506in}{1.800605in}}%
\pgfpathmoveto{\pgfqpoint{2.743506in}{1.797655in}}%
\pgfpathlineto{\pgfqpoint{2.743506in}{1.797655in}}%
\pgfpathlineto{\pgfqpoint{2.743506in}{1.800605in}}%
\pgfpathlineto{\pgfqpoint{2.748047in}{1.800605in}}%
\pgfpathlineto{\pgfqpoint{2.748047in}{1.797655in}}%
\pgfpathmoveto{\pgfqpoint{2.743506in}{1.800605in}}%
\pgfpathlineto{\pgfqpoint{2.743506in}{1.800605in}}%
\pgfpathlineto{\pgfqpoint{2.743506in}{1.803554in}}%
\pgfpathlineto{\pgfqpoint{2.748047in}{1.803554in}}%
\pgfpathlineto{\pgfqpoint{2.748047in}{1.800605in}}%
\pgfpathmoveto{\pgfqpoint{2.748047in}{1.797655in}}%
\pgfpathlineto{\pgfqpoint{2.748047in}{1.797655in}}%
\pgfpathlineto{\pgfqpoint{2.748047in}{1.800605in}}%
\pgfpathlineto{\pgfqpoint{2.752588in}{1.800605in}}%
\pgfpathlineto{\pgfqpoint{2.752588in}{1.797655in}}%
\pgfpathmoveto{\pgfqpoint{2.748047in}{1.800605in}}%
\pgfpathlineto{\pgfqpoint{2.748047in}{1.800605in}}%
\pgfpathlineto{\pgfqpoint{2.748047in}{1.803554in}}%
\pgfpathlineto{\pgfqpoint{2.752588in}{1.803554in}}%
\pgfpathlineto{\pgfqpoint{2.752588in}{1.800605in}}%
\pgfpathmoveto{\pgfqpoint{2.752588in}{1.800605in}}%
\pgfpathlineto{\pgfqpoint{2.752588in}{1.800605in}}%
\pgfpathlineto{\pgfqpoint{2.752588in}{1.803554in}}%
\pgfpathlineto{\pgfqpoint{2.757129in}{1.803554in}}%
\pgfpathlineto{\pgfqpoint{2.757129in}{1.800605in}}%
\pgfpathmoveto{\pgfqpoint{2.748047in}{1.803554in}}%
\pgfpathlineto{\pgfqpoint{2.748047in}{1.803554in}}%
\pgfpathlineto{\pgfqpoint{2.748047in}{1.806503in}}%
\pgfpathlineto{\pgfqpoint{2.752588in}{1.806503in}}%
\pgfpathlineto{\pgfqpoint{2.752588in}{1.803554in}}%
\pgfpathmoveto{\pgfqpoint{2.748047in}{1.806503in}}%
\pgfpathlineto{\pgfqpoint{2.748047in}{1.806503in}}%
\pgfpathlineto{\pgfqpoint{2.748047in}{1.809452in}}%
\pgfpathlineto{\pgfqpoint{2.752588in}{1.809452in}}%
\pgfpathlineto{\pgfqpoint{2.752588in}{1.806503in}}%
\pgfpathmoveto{\pgfqpoint{2.752588in}{1.803554in}}%
\pgfpathlineto{\pgfqpoint{2.752588in}{1.803554in}}%
\pgfpathlineto{\pgfqpoint{2.752588in}{1.806503in}}%
\pgfpathlineto{\pgfqpoint{2.757129in}{1.806503in}}%
\pgfpathlineto{\pgfqpoint{2.757129in}{1.803554in}}%
\pgfpathmoveto{\pgfqpoint{2.752588in}{1.806503in}}%
\pgfpathlineto{\pgfqpoint{2.752588in}{1.806503in}}%
\pgfpathlineto{\pgfqpoint{2.752588in}{1.809452in}}%
\pgfpathlineto{\pgfqpoint{2.757129in}{1.809452in}}%
\pgfpathlineto{\pgfqpoint{2.757129in}{1.806503in}}%
\pgfpathmoveto{\pgfqpoint{2.757129in}{1.803554in}}%
\pgfpathlineto{\pgfqpoint{2.757129in}{1.803554in}}%
\pgfpathlineto{\pgfqpoint{2.757129in}{1.806503in}}%
\pgfpathlineto{\pgfqpoint{2.761670in}{1.806503in}}%
\pgfpathlineto{\pgfqpoint{2.761670in}{1.803554in}}%
\pgfpathmoveto{\pgfqpoint{2.757129in}{1.806503in}}%
\pgfpathlineto{\pgfqpoint{2.757129in}{1.806503in}}%
\pgfpathlineto{\pgfqpoint{2.757129in}{1.809452in}}%
\pgfpathlineto{\pgfqpoint{2.761670in}{1.809452in}}%
\pgfpathlineto{\pgfqpoint{2.761670in}{1.806503in}}%
\pgfpathmoveto{\pgfqpoint{2.761670in}{1.806503in}}%
\pgfpathlineto{\pgfqpoint{2.761670in}{1.806503in}}%
\pgfpathlineto{\pgfqpoint{2.761670in}{1.809452in}}%
\pgfpathlineto{\pgfqpoint{2.766211in}{1.809452in}}%
\pgfpathlineto{\pgfqpoint{2.766211in}{1.806503in}}%
\pgfpathmoveto{\pgfqpoint{2.757129in}{1.809452in}}%
\pgfpathlineto{\pgfqpoint{2.757129in}{1.809452in}}%
\pgfpathlineto{\pgfqpoint{2.757129in}{1.812401in}}%
\pgfpathlineto{\pgfqpoint{2.761670in}{1.812401in}}%
\pgfpathlineto{\pgfqpoint{2.761670in}{1.809452in}}%
\pgfpathmoveto{\pgfqpoint{2.757129in}{1.812401in}}%
\pgfpathlineto{\pgfqpoint{2.757129in}{1.812401in}}%
\pgfpathlineto{\pgfqpoint{2.757129in}{1.815350in}}%
\pgfpathlineto{\pgfqpoint{2.761670in}{1.815350in}}%
\pgfpathlineto{\pgfqpoint{2.761670in}{1.812401in}}%
\pgfpathmoveto{\pgfqpoint{2.761670in}{1.809452in}}%
\pgfpathlineto{\pgfqpoint{2.761670in}{1.809452in}}%
\pgfpathlineto{\pgfqpoint{2.761670in}{1.812401in}}%
\pgfpathlineto{\pgfqpoint{2.766211in}{1.812401in}}%
\pgfpathlineto{\pgfqpoint{2.766211in}{1.809452in}}%
\pgfpathmoveto{\pgfqpoint{2.761670in}{1.812401in}}%
\pgfpathlineto{\pgfqpoint{2.761670in}{1.812401in}}%
\pgfpathlineto{\pgfqpoint{2.761670in}{1.815350in}}%
\pgfpathlineto{\pgfqpoint{2.766211in}{1.815350in}}%
\pgfpathlineto{\pgfqpoint{2.766211in}{1.812401in}}%
\pgfpathmoveto{\pgfqpoint{2.766211in}{1.809452in}}%
\pgfpathlineto{\pgfqpoint{2.766211in}{1.809452in}}%
\pgfpathlineto{\pgfqpoint{2.766211in}{1.812401in}}%
\pgfpathlineto{\pgfqpoint{2.770752in}{1.812401in}}%
\pgfpathlineto{\pgfqpoint{2.770752in}{1.809452in}}%
\pgfpathmoveto{\pgfqpoint{2.766211in}{1.812401in}}%
\pgfpathlineto{\pgfqpoint{2.766211in}{1.812401in}}%
\pgfpathlineto{\pgfqpoint{2.766211in}{1.815350in}}%
\pgfpathlineto{\pgfqpoint{2.770752in}{1.815350in}}%
\pgfpathlineto{\pgfqpoint{2.770752in}{1.812401in}}%
\pgfpathmoveto{\pgfqpoint{2.770752in}{1.812401in}}%
\pgfpathlineto{\pgfqpoint{2.770752in}{1.812401in}}%
\pgfpathlineto{\pgfqpoint{2.770752in}{1.815350in}}%
\pgfpathlineto{\pgfqpoint{2.775293in}{1.815350in}}%
\pgfpathlineto{\pgfqpoint{2.775293in}{1.812401in}}%
\pgfpathmoveto{\pgfqpoint{2.766211in}{1.815350in}}%
\pgfpathlineto{\pgfqpoint{2.766211in}{1.815350in}}%
\pgfpathlineto{\pgfqpoint{2.766211in}{1.818299in}}%
\pgfpathlineto{\pgfqpoint{2.770752in}{1.818299in}}%
\pgfpathlineto{\pgfqpoint{2.770752in}{1.815350in}}%
\pgfpathmoveto{\pgfqpoint{2.766211in}{1.818299in}}%
\pgfpathlineto{\pgfqpoint{2.766211in}{1.818299in}}%
\pgfpathlineto{\pgfqpoint{2.766211in}{1.821249in}}%
\pgfpathlineto{\pgfqpoint{2.770752in}{1.821249in}}%
\pgfpathlineto{\pgfqpoint{2.770752in}{1.818299in}}%
\pgfpathmoveto{\pgfqpoint{2.770752in}{1.815350in}}%
\pgfpathlineto{\pgfqpoint{2.770752in}{1.815350in}}%
\pgfpathlineto{\pgfqpoint{2.770752in}{1.818299in}}%
\pgfpathlineto{\pgfqpoint{2.775293in}{1.818299in}}%
\pgfpathlineto{\pgfqpoint{2.775293in}{1.815350in}}%
\pgfpathmoveto{\pgfqpoint{2.770752in}{1.818299in}}%
\pgfpathlineto{\pgfqpoint{2.770752in}{1.818299in}}%
\pgfpathlineto{\pgfqpoint{2.770752in}{1.821249in}}%
\pgfpathlineto{\pgfqpoint{2.775293in}{1.821249in}}%
\pgfpathlineto{\pgfqpoint{2.775293in}{1.818299in}}%
\pgfpathmoveto{\pgfqpoint{2.775293in}{1.815350in}}%
\pgfpathlineto{\pgfqpoint{2.775293in}{1.815350in}}%
\pgfpathlineto{\pgfqpoint{2.775293in}{1.818299in}}%
\pgfpathlineto{\pgfqpoint{2.779834in}{1.818299in}}%
\pgfpathlineto{\pgfqpoint{2.779834in}{1.815350in}}%
\pgfpathmoveto{\pgfqpoint{2.775293in}{1.818299in}}%
\pgfpathlineto{\pgfqpoint{2.775293in}{1.818299in}}%
\pgfpathlineto{\pgfqpoint{2.775293in}{1.821249in}}%
\pgfpathlineto{\pgfqpoint{2.779834in}{1.821249in}}%
\pgfpathlineto{\pgfqpoint{2.779834in}{1.818299in}}%
\pgfpathmoveto{\pgfqpoint{2.779834in}{1.818299in}}%
\pgfpathlineto{\pgfqpoint{2.779834in}{1.818299in}}%
\pgfpathlineto{\pgfqpoint{2.779834in}{1.821249in}}%
\pgfpathlineto{\pgfqpoint{2.784375in}{1.821249in}}%
\pgfpathlineto{\pgfqpoint{2.784375in}{1.818299in}}%
\pgfpathmoveto{\pgfqpoint{2.775293in}{1.821249in}}%
\pgfpathlineto{\pgfqpoint{2.775293in}{1.821249in}}%
\pgfpathlineto{\pgfqpoint{2.775293in}{1.824198in}}%
\pgfpathlineto{\pgfqpoint{2.779834in}{1.824198in}}%
\pgfpathlineto{\pgfqpoint{2.779834in}{1.821249in}}%
\pgfpathmoveto{\pgfqpoint{2.775293in}{1.824198in}}%
\pgfpathlineto{\pgfqpoint{2.775293in}{1.824198in}}%
\pgfpathlineto{\pgfqpoint{2.775293in}{1.827147in}}%
\pgfpathlineto{\pgfqpoint{2.779834in}{1.827147in}}%
\pgfpathlineto{\pgfqpoint{2.779834in}{1.824198in}}%
\pgfpathmoveto{\pgfqpoint{2.779834in}{1.821249in}}%
\pgfpathlineto{\pgfqpoint{2.779834in}{1.821249in}}%
\pgfpathlineto{\pgfqpoint{2.779834in}{1.824198in}}%
\pgfpathlineto{\pgfqpoint{2.784375in}{1.824198in}}%
\pgfpathlineto{\pgfqpoint{2.784375in}{1.821249in}}%
\pgfpathmoveto{\pgfqpoint{2.779834in}{1.824198in}}%
\pgfpathlineto{\pgfqpoint{2.779834in}{1.824198in}}%
\pgfpathlineto{\pgfqpoint{2.779834in}{1.827147in}}%
\pgfpathlineto{\pgfqpoint{2.784375in}{1.827147in}}%
\pgfpathlineto{\pgfqpoint{2.784375in}{1.824198in}}%
\pgfpathmoveto{\pgfqpoint{2.784375in}{1.821249in}}%
\pgfpathlineto{\pgfqpoint{2.784375in}{1.821249in}}%
\pgfpathlineto{\pgfqpoint{2.784375in}{1.824198in}}%
\pgfpathlineto{\pgfqpoint{2.788916in}{1.824198in}}%
\pgfpathlineto{\pgfqpoint{2.788916in}{1.821249in}}%
\pgfpathmoveto{\pgfqpoint{2.784375in}{1.824198in}}%
\pgfpathlineto{\pgfqpoint{2.784375in}{1.824198in}}%
\pgfpathlineto{\pgfqpoint{2.784375in}{1.827147in}}%
\pgfpathlineto{\pgfqpoint{2.788916in}{1.827147in}}%
\pgfpathlineto{\pgfqpoint{2.788916in}{1.824198in}}%
\pgfpathmoveto{\pgfqpoint{2.788916in}{1.824198in}}%
\pgfpathlineto{\pgfqpoint{2.788916in}{1.824198in}}%
\pgfpathlineto{\pgfqpoint{2.788916in}{1.827147in}}%
\pgfpathlineto{\pgfqpoint{2.793457in}{1.827147in}}%
\pgfpathlineto{\pgfqpoint{2.793457in}{1.824198in}}%
\pgfpathmoveto{\pgfqpoint{2.784375in}{1.827147in}}%
\pgfpathlineto{\pgfqpoint{2.784375in}{1.827147in}}%
\pgfpathlineto{\pgfqpoint{2.784375in}{1.830097in}}%
\pgfpathlineto{\pgfqpoint{2.788916in}{1.830097in}}%
\pgfpathlineto{\pgfqpoint{2.788916in}{1.827147in}}%
\pgfpathmoveto{\pgfqpoint{2.784375in}{1.830097in}}%
\pgfpathlineto{\pgfqpoint{2.784375in}{1.830097in}}%
\pgfpathlineto{\pgfqpoint{2.784375in}{1.833046in}}%
\pgfpathlineto{\pgfqpoint{2.788916in}{1.833046in}}%
\pgfpathlineto{\pgfqpoint{2.788916in}{1.830097in}}%
\pgfpathmoveto{\pgfqpoint{2.788916in}{1.827147in}}%
\pgfpathlineto{\pgfqpoint{2.788916in}{1.827147in}}%
\pgfpathlineto{\pgfqpoint{2.788916in}{1.830097in}}%
\pgfpathlineto{\pgfqpoint{2.793457in}{1.830097in}}%
\pgfpathlineto{\pgfqpoint{2.793457in}{1.827147in}}%
\pgfpathmoveto{\pgfqpoint{2.788916in}{1.830097in}}%
\pgfpathlineto{\pgfqpoint{2.788916in}{1.830097in}}%
\pgfpathlineto{\pgfqpoint{2.788916in}{1.833046in}}%
\pgfpathlineto{\pgfqpoint{2.793457in}{1.833046in}}%
\pgfpathlineto{\pgfqpoint{2.793457in}{1.830097in}}%
\pgfpathmoveto{\pgfqpoint{2.793457in}{1.827147in}}%
\pgfpathlineto{\pgfqpoint{2.793457in}{1.827147in}}%
\pgfpathlineto{\pgfqpoint{2.793457in}{1.830097in}}%
\pgfpathlineto{\pgfqpoint{2.797998in}{1.830097in}}%
\pgfpathlineto{\pgfqpoint{2.797998in}{1.827147in}}%
\pgfpathmoveto{\pgfqpoint{2.793457in}{1.830097in}}%
\pgfpathlineto{\pgfqpoint{2.793457in}{1.830097in}}%
\pgfpathlineto{\pgfqpoint{2.793457in}{1.833046in}}%
\pgfpathlineto{\pgfqpoint{2.797998in}{1.833046in}}%
\pgfpathlineto{\pgfqpoint{2.797998in}{1.830097in}}%
\pgfpathmoveto{\pgfqpoint{2.797998in}{1.830097in}}%
\pgfpathlineto{\pgfqpoint{2.797998in}{1.830097in}}%
\pgfpathlineto{\pgfqpoint{2.797998in}{1.833046in}}%
\pgfpathlineto{\pgfqpoint{2.802539in}{1.833046in}}%
\pgfpathlineto{\pgfqpoint{2.802539in}{1.830097in}}%
\pgfpathmoveto{\pgfqpoint{2.793457in}{1.833046in}}%
\pgfpathlineto{\pgfqpoint{2.793457in}{1.833046in}}%
\pgfpathlineto{\pgfqpoint{2.793457in}{1.835995in}}%
\pgfpathlineto{\pgfqpoint{2.797998in}{1.835995in}}%
\pgfpathlineto{\pgfqpoint{2.797998in}{1.833046in}}%
\pgfpathmoveto{\pgfqpoint{2.793457in}{1.835995in}}%
\pgfpathlineto{\pgfqpoint{2.793457in}{1.835995in}}%
\pgfpathlineto{\pgfqpoint{2.793457in}{1.838944in}}%
\pgfpathlineto{\pgfqpoint{2.797998in}{1.838944in}}%
\pgfpathlineto{\pgfqpoint{2.797998in}{1.835995in}}%
\pgfpathmoveto{\pgfqpoint{2.797998in}{1.833046in}}%
\pgfpathlineto{\pgfqpoint{2.797998in}{1.833046in}}%
\pgfpathlineto{\pgfqpoint{2.797998in}{1.835995in}}%
\pgfpathlineto{\pgfqpoint{2.802539in}{1.835995in}}%
\pgfpathlineto{\pgfqpoint{2.802539in}{1.833046in}}%
\pgfpathmoveto{\pgfqpoint{2.797998in}{1.835995in}}%
\pgfpathlineto{\pgfqpoint{2.797998in}{1.835995in}}%
\pgfpathlineto{\pgfqpoint{2.797998in}{1.838944in}}%
\pgfpathlineto{\pgfqpoint{2.802539in}{1.838944in}}%
\pgfpathlineto{\pgfqpoint{2.802539in}{1.835995in}}%
\pgfpathmoveto{\pgfqpoint{2.802539in}{1.833046in}}%
\pgfpathlineto{\pgfqpoint{2.802539in}{1.833046in}}%
\pgfpathlineto{\pgfqpoint{2.802539in}{1.835995in}}%
\pgfpathlineto{\pgfqpoint{2.807080in}{1.835995in}}%
\pgfpathlineto{\pgfqpoint{2.807080in}{1.833046in}}%
\pgfpathmoveto{\pgfqpoint{2.802539in}{1.835995in}}%
\pgfpathlineto{\pgfqpoint{2.802539in}{1.835995in}}%
\pgfpathlineto{\pgfqpoint{2.802539in}{1.838944in}}%
\pgfpathlineto{\pgfqpoint{2.807080in}{1.838944in}}%
\pgfpathlineto{\pgfqpoint{2.807080in}{1.835995in}}%
\pgfpathmoveto{\pgfqpoint{2.807080in}{1.835995in}}%
\pgfpathlineto{\pgfqpoint{2.807080in}{1.835995in}}%
\pgfpathlineto{\pgfqpoint{2.807080in}{1.838944in}}%
\pgfpathlineto{\pgfqpoint{2.811621in}{1.838944in}}%
\pgfpathlineto{\pgfqpoint{2.811621in}{1.835995in}}%
\pgfpathmoveto{\pgfqpoint{2.802539in}{1.838944in}}%
\pgfpathlineto{\pgfqpoint{2.802539in}{1.838944in}}%
\pgfpathlineto{\pgfqpoint{2.802539in}{1.841894in}}%
\pgfpathlineto{\pgfqpoint{2.807080in}{1.841894in}}%
\pgfpathlineto{\pgfqpoint{2.807080in}{1.838944in}}%
\pgfpathmoveto{\pgfqpoint{2.802539in}{1.841894in}}%
\pgfpathlineto{\pgfqpoint{2.802539in}{1.841894in}}%
\pgfpathlineto{\pgfqpoint{2.802539in}{1.844843in}}%
\pgfpathlineto{\pgfqpoint{2.807080in}{1.844843in}}%
\pgfpathlineto{\pgfqpoint{2.807080in}{1.841894in}}%
\pgfpathmoveto{\pgfqpoint{2.807080in}{1.838944in}}%
\pgfpathlineto{\pgfqpoint{2.807080in}{1.838944in}}%
\pgfpathlineto{\pgfqpoint{2.807080in}{1.841894in}}%
\pgfpathlineto{\pgfqpoint{2.811621in}{1.841894in}}%
\pgfpathlineto{\pgfqpoint{2.811621in}{1.838944in}}%
\pgfpathmoveto{\pgfqpoint{2.807080in}{1.841894in}}%
\pgfpathlineto{\pgfqpoint{2.807080in}{1.841894in}}%
\pgfpathlineto{\pgfqpoint{2.807080in}{1.844843in}}%
\pgfpathlineto{\pgfqpoint{2.811621in}{1.844843in}}%
\pgfpathlineto{\pgfqpoint{2.811621in}{1.841894in}}%
\pgfpathmoveto{\pgfqpoint{2.811621in}{1.838944in}}%
\pgfpathlineto{\pgfqpoint{2.811621in}{1.838944in}}%
\pgfpathlineto{\pgfqpoint{2.811621in}{1.841894in}}%
\pgfpathlineto{\pgfqpoint{2.816162in}{1.841894in}}%
\pgfpathlineto{\pgfqpoint{2.816162in}{1.838944in}}%
\pgfpathmoveto{\pgfqpoint{2.811621in}{1.841894in}}%
\pgfpathlineto{\pgfqpoint{2.811621in}{1.841894in}}%
\pgfpathlineto{\pgfqpoint{2.811621in}{1.844843in}}%
\pgfpathlineto{\pgfqpoint{2.816162in}{1.844843in}}%
\pgfpathlineto{\pgfqpoint{2.816162in}{1.841894in}}%
\pgfpathmoveto{\pgfqpoint{2.816162in}{1.841894in}}%
\pgfpathlineto{\pgfqpoint{2.816162in}{1.841894in}}%
\pgfpathlineto{\pgfqpoint{2.816162in}{1.844843in}}%
\pgfpathlineto{\pgfqpoint{2.820703in}{1.844843in}}%
\pgfpathlineto{\pgfqpoint{2.820703in}{1.841894in}}%
\pgfpathmoveto{\pgfqpoint{2.811621in}{1.844843in}}%
\pgfpathlineto{\pgfqpoint{2.811621in}{1.844843in}}%
\pgfpathlineto{\pgfqpoint{2.811621in}{1.847792in}}%
\pgfpathlineto{\pgfqpoint{2.816162in}{1.847792in}}%
\pgfpathlineto{\pgfqpoint{2.816162in}{1.844843in}}%
\pgfpathmoveto{\pgfqpoint{2.811621in}{1.847792in}}%
\pgfpathlineto{\pgfqpoint{2.811621in}{1.847792in}}%
\pgfpathlineto{\pgfqpoint{2.811621in}{1.850742in}}%
\pgfpathlineto{\pgfqpoint{2.816162in}{1.850742in}}%
\pgfpathlineto{\pgfqpoint{2.816162in}{1.847792in}}%
\pgfpathmoveto{\pgfqpoint{2.816162in}{1.844843in}}%
\pgfpathlineto{\pgfqpoint{2.816162in}{1.844843in}}%
\pgfpathlineto{\pgfqpoint{2.816162in}{1.847792in}}%
\pgfpathlineto{\pgfqpoint{2.820703in}{1.847792in}}%
\pgfpathlineto{\pgfqpoint{2.820703in}{1.844843in}}%
\pgfpathmoveto{\pgfqpoint{2.816162in}{1.847792in}}%
\pgfpathlineto{\pgfqpoint{2.816162in}{1.847792in}}%
\pgfpathlineto{\pgfqpoint{2.816162in}{1.850742in}}%
\pgfpathlineto{\pgfqpoint{2.820703in}{1.850742in}}%
\pgfpathlineto{\pgfqpoint{2.820703in}{1.847792in}}%
\pgfpathmoveto{\pgfqpoint{2.820703in}{1.844843in}}%
\pgfpathlineto{\pgfqpoint{2.820703in}{1.844843in}}%
\pgfpathlineto{\pgfqpoint{2.820703in}{1.847792in}}%
\pgfpathlineto{\pgfqpoint{2.825244in}{1.847792in}}%
\pgfpathlineto{\pgfqpoint{2.825244in}{1.844843in}}%
\pgfpathmoveto{\pgfqpoint{2.820703in}{1.847792in}}%
\pgfpathlineto{\pgfqpoint{2.820703in}{1.847792in}}%
\pgfpathlineto{\pgfqpoint{2.820703in}{1.850742in}}%
\pgfpathlineto{\pgfqpoint{2.825244in}{1.850742in}}%
\pgfpathlineto{\pgfqpoint{2.825244in}{1.847792in}}%
\pgfpathmoveto{\pgfqpoint{2.825244in}{1.847792in}}%
\pgfpathlineto{\pgfqpoint{2.825244in}{1.847792in}}%
\pgfpathlineto{\pgfqpoint{2.825244in}{1.850742in}}%
\pgfpathlineto{\pgfqpoint{2.829785in}{1.850742in}}%
\pgfpathlineto{\pgfqpoint{2.829785in}{1.847792in}}%
\pgfpathmoveto{\pgfqpoint{2.820703in}{1.850742in}}%
\pgfpathlineto{\pgfqpoint{2.820703in}{1.850742in}}%
\pgfpathlineto{\pgfqpoint{2.820703in}{1.853691in}}%
\pgfpathlineto{\pgfqpoint{2.825244in}{1.853691in}}%
\pgfpathlineto{\pgfqpoint{2.825244in}{1.850742in}}%
\pgfpathmoveto{\pgfqpoint{2.820703in}{1.853691in}}%
\pgfpathlineto{\pgfqpoint{2.820703in}{1.853691in}}%
\pgfpathlineto{\pgfqpoint{2.820703in}{1.856640in}}%
\pgfpathlineto{\pgfqpoint{2.825244in}{1.856640in}}%
\pgfpathlineto{\pgfqpoint{2.825244in}{1.853691in}}%
\pgfpathmoveto{\pgfqpoint{2.825244in}{1.850742in}}%
\pgfpathlineto{\pgfqpoint{2.825244in}{1.850742in}}%
\pgfpathlineto{\pgfqpoint{2.825244in}{1.853691in}}%
\pgfpathlineto{\pgfqpoint{2.829785in}{1.853691in}}%
\pgfpathlineto{\pgfqpoint{2.829785in}{1.850742in}}%
\pgfpathmoveto{\pgfqpoint{2.825244in}{1.853691in}}%
\pgfpathlineto{\pgfqpoint{2.825244in}{1.853691in}}%
\pgfpathlineto{\pgfqpoint{2.825244in}{1.856640in}}%
\pgfpathlineto{\pgfqpoint{2.829785in}{1.856640in}}%
\pgfpathlineto{\pgfqpoint{2.829785in}{1.853691in}}%
\pgfpathmoveto{\pgfqpoint{2.829785in}{1.850742in}}%
\pgfpathlineto{\pgfqpoint{2.829785in}{1.850742in}}%
\pgfpathlineto{\pgfqpoint{2.829785in}{1.853691in}}%
\pgfpathlineto{\pgfqpoint{2.834326in}{1.853691in}}%
\pgfpathlineto{\pgfqpoint{2.834326in}{1.850742in}}%
\pgfpathmoveto{\pgfqpoint{2.829785in}{1.853691in}}%
\pgfpathlineto{\pgfqpoint{2.829785in}{1.853691in}}%
\pgfpathlineto{\pgfqpoint{2.829785in}{1.856640in}}%
\pgfpathlineto{\pgfqpoint{2.834326in}{1.856640in}}%
\pgfpathlineto{\pgfqpoint{2.834326in}{1.853691in}}%
\pgfpathmoveto{\pgfqpoint{2.834326in}{1.853691in}}%
\pgfpathlineto{\pgfqpoint{2.834326in}{1.853691in}}%
\pgfpathlineto{\pgfqpoint{2.834326in}{1.856640in}}%
\pgfpathlineto{\pgfqpoint{2.838867in}{1.856640in}}%
\pgfpathlineto{\pgfqpoint{2.838867in}{1.853691in}}%
\pgfpathmoveto{\pgfqpoint{2.829785in}{1.856640in}}%
\pgfpathlineto{\pgfqpoint{2.829785in}{1.856640in}}%
\pgfpathlineto{\pgfqpoint{2.829785in}{1.859590in}}%
\pgfpathlineto{\pgfqpoint{2.834326in}{1.859590in}}%
\pgfpathlineto{\pgfqpoint{2.834326in}{1.856640in}}%
\pgfpathmoveto{\pgfqpoint{2.829785in}{1.859590in}}%
\pgfpathlineto{\pgfqpoint{2.829785in}{1.859590in}}%
\pgfpathlineto{\pgfqpoint{2.829785in}{1.862539in}}%
\pgfpathlineto{\pgfqpoint{2.834326in}{1.862539in}}%
\pgfpathlineto{\pgfqpoint{2.834326in}{1.859590in}}%
\pgfpathmoveto{\pgfqpoint{2.834326in}{1.856640in}}%
\pgfpathlineto{\pgfqpoint{2.834326in}{1.856640in}}%
\pgfpathlineto{\pgfqpoint{2.834326in}{1.859590in}}%
\pgfpathlineto{\pgfqpoint{2.838867in}{1.859590in}}%
\pgfpathlineto{\pgfqpoint{2.838867in}{1.856640in}}%
\pgfpathmoveto{\pgfqpoint{2.834326in}{1.859590in}}%
\pgfpathlineto{\pgfqpoint{2.834326in}{1.859590in}}%
\pgfpathlineto{\pgfqpoint{2.834326in}{1.862539in}}%
\pgfpathlineto{\pgfqpoint{2.838867in}{1.862539in}}%
\pgfpathlineto{\pgfqpoint{2.838867in}{1.859590in}}%
\pgfpathmoveto{\pgfqpoint{2.838867in}{1.856640in}}%
\pgfpathlineto{\pgfqpoint{2.838867in}{1.856640in}}%
\pgfpathlineto{\pgfqpoint{2.838867in}{1.859590in}}%
\pgfpathlineto{\pgfqpoint{2.843408in}{1.859590in}}%
\pgfpathlineto{\pgfqpoint{2.843408in}{1.856640in}}%
\pgfpathmoveto{\pgfqpoint{2.838867in}{1.859590in}}%
\pgfpathlineto{\pgfqpoint{2.838867in}{1.859590in}}%
\pgfpathlineto{\pgfqpoint{2.838867in}{1.862539in}}%
\pgfpathlineto{\pgfqpoint{2.843408in}{1.862539in}}%
\pgfpathlineto{\pgfqpoint{2.843408in}{1.859590in}}%
\pgfpathmoveto{\pgfqpoint{2.843408in}{1.859590in}}%
\pgfpathlineto{\pgfqpoint{2.843408in}{1.859590in}}%
\pgfpathlineto{\pgfqpoint{2.843408in}{1.862539in}}%
\pgfpathlineto{\pgfqpoint{2.847949in}{1.862539in}}%
\pgfpathlineto{\pgfqpoint{2.847949in}{1.859590in}}%
\pgfpathmoveto{\pgfqpoint{2.838867in}{1.862539in}}%
\pgfpathlineto{\pgfqpoint{2.838867in}{1.862539in}}%
\pgfpathlineto{\pgfqpoint{2.838867in}{1.865488in}}%
\pgfpathlineto{\pgfqpoint{2.843408in}{1.865488in}}%
\pgfpathlineto{\pgfqpoint{2.843408in}{1.862539in}}%
\pgfpathmoveto{\pgfqpoint{2.838867in}{1.865488in}}%
\pgfpathlineto{\pgfqpoint{2.838867in}{1.865488in}}%
\pgfpathlineto{\pgfqpoint{2.838867in}{1.868438in}}%
\pgfpathlineto{\pgfqpoint{2.843408in}{1.868438in}}%
\pgfpathlineto{\pgfqpoint{2.843408in}{1.865488in}}%
\pgfpathmoveto{\pgfqpoint{2.843408in}{1.862539in}}%
\pgfpathlineto{\pgfqpoint{2.843408in}{1.862539in}}%
\pgfpathlineto{\pgfqpoint{2.843408in}{1.865488in}}%
\pgfpathlineto{\pgfqpoint{2.847949in}{1.865488in}}%
\pgfpathlineto{\pgfqpoint{2.847949in}{1.862539in}}%
\pgfpathmoveto{\pgfqpoint{2.843408in}{1.865488in}}%
\pgfpathlineto{\pgfqpoint{2.843408in}{1.865488in}}%
\pgfpathlineto{\pgfqpoint{2.843408in}{1.868438in}}%
\pgfpathlineto{\pgfqpoint{2.847949in}{1.868438in}}%
\pgfpathlineto{\pgfqpoint{2.847949in}{1.865488in}}%
\pgfpathmoveto{\pgfqpoint{2.847949in}{1.862539in}}%
\pgfpathlineto{\pgfqpoint{2.847949in}{1.862539in}}%
\pgfpathlineto{\pgfqpoint{2.847949in}{1.865488in}}%
\pgfpathlineto{\pgfqpoint{2.852490in}{1.865488in}}%
\pgfpathlineto{\pgfqpoint{2.852490in}{1.862539in}}%
\pgfpathmoveto{\pgfqpoint{2.847949in}{1.865488in}}%
\pgfpathlineto{\pgfqpoint{2.847949in}{1.865488in}}%
\pgfpathlineto{\pgfqpoint{2.847949in}{1.868438in}}%
\pgfpathlineto{\pgfqpoint{2.852490in}{1.868438in}}%
\pgfpathlineto{\pgfqpoint{2.852490in}{1.865488in}}%
\pgfpathmoveto{\pgfqpoint{2.852490in}{1.865488in}}%
\pgfpathlineto{\pgfqpoint{2.852490in}{1.865488in}}%
\pgfpathlineto{\pgfqpoint{2.852490in}{1.868438in}}%
\pgfpathlineto{\pgfqpoint{2.857031in}{1.868438in}}%
\pgfpathlineto{\pgfqpoint{2.857031in}{1.865488in}}%
\pgfpathmoveto{\pgfqpoint{2.857031in}{1.865488in}}%
\pgfpathlineto{\pgfqpoint{2.857031in}{1.865488in}}%
\pgfpathlineto{\pgfqpoint{2.857031in}{1.868438in}}%
\pgfpathlineto{\pgfqpoint{2.861572in}{1.868438in}}%
\pgfpathlineto{\pgfqpoint{2.861572in}{1.865488in}}%
\pgfpathmoveto{\pgfqpoint{2.857031in}{1.868438in}}%
\pgfpathlineto{\pgfqpoint{2.857031in}{1.868438in}}%
\pgfpathlineto{\pgfqpoint{2.857031in}{1.871387in}}%
\pgfpathlineto{\pgfqpoint{2.861572in}{1.871387in}}%
\pgfpathlineto{\pgfqpoint{2.861572in}{1.868438in}}%
\pgfpathmoveto{\pgfqpoint{2.857031in}{1.871387in}}%
\pgfpathlineto{\pgfqpoint{2.857031in}{1.871387in}}%
\pgfpathlineto{\pgfqpoint{2.857031in}{1.874336in}}%
\pgfpathlineto{\pgfqpoint{2.861572in}{1.874336in}}%
\pgfpathlineto{\pgfqpoint{2.861572in}{1.871387in}}%
\pgfpathmoveto{\pgfqpoint{2.861572in}{1.868438in}}%
\pgfpathlineto{\pgfqpoint{2.861572in}{1.868438in}}%
\pgfpathlineto{\pgfqpoint{2.861572in}{1.871387in}}%
\pgfpathlineto{\pgfqpoint{2.866113in}{1.871387in}}%
\pgfpathlineto{\pgfqpoint{2.866113in}{1.868438in}}%
\pgfpathmoveto{\pgfqpoint{2.861572in}{1.871387in}}%
\pgfpathlineto{\pgfqpoint{2.861572in}{1.871387in}}%
\pgfpathlineto{\pgfqpoint{2.861572in}{1.874336in}}%
\pgfpathlineto{\pgfqpoint{2.866113in}{1.874336in}}%
\pgfpathlineto{\pgfqpoint{2.866113in}{1.871387in}}%
\pgfpathmoveto{\pgfqpoint{2.866113in}{1.871387in}}%
\pgfpathlineto{\pgfqpoint{2.866113in}{1.871387in}}%
\pgfpathlineto{\pgfqpoint{2.866113in}{1.874336in}}%
\pgfpathlineto{\pgfqpoint{2.870655in}{1.874336in}}%
\pgfpathlineto{\pgfqpoint{2.870655in}{1.871387in}}%
\pgfpathmoveto{\pgfqpoint{2.866113in}{1.874336in}}%
\pgfpathlineto{\pgfqpoint{2.866113in}{1.874336in}}%
\pgfpathlineto{\pgfqpoint{2.866113in}{1.877286in}}%
\pgfpathlineto{\pgfqpoint{2.870655in}{1.877286in}}%
\pgfpathlineto{\pgfqpoint{2.870655in}{1.874336in}}%
\pgfpathmoveto{\pgfqpoint{2.866113in}{1.877286in}}%
\pgfpathlineto{\pgfqpoint{2.866113in}{1.877286in}}%
\pgfpathlineto{\pgfqpoint{2.866113in}{1.880235in}}%
\pgfpathlineto{\pgfqpoint{2.870655in}{1.880235in}}%
\pgfpathlineto{\pgfqpoint{2.870655in}{1.877286in}}%
\pgfpathmoveto{\pgfqpoint{2.870655in}{1.874336in}}%
\pgfpathlineto{\pgfqpoint{2.870655in}{1.874336in}}%
\pgfpathlineto{\pgfqpoint{2.870655in}{1.877286in}}%
\pgfpathlineto{\pgfqpoint{2.875196in}{1.877286in}}%
\pgfpathlineto{\pgfqpoint{2.875196in}{1.874336in}}%
\pgfpathmoveto{\pgfqpoint{2.870655in}{1.877286in}}%
\pgfpathlineto{\pgfqpoint{2.870655in}{1.877286in}}%
\pgfpathlineto{\pgfqpoint{2.870655in}{1.880235in}}%
\pgfpathlineto{\pgfqpoint{2.875196in}{1.880235in}}%
\pgfpathlineto{\pgfqpoint{2.875196in}{1.877286in}}%
\pgfpathmoveto{\pgfqpoint{2.875196in}{1.877286in}}%
\pgfpathlineto{\pgfqpoint{2.875196in}{1.877286in}}%
\pgfpathlineto{\pgfqpoint{2.875196in}{1.880235in}}%
\pgfpathlineto{\pgfqpoint{2.879737in}{1.880235in}}%
\pgfpathlineto{\pgfqpoint{2.879737in}{1.877286in}}%
\pgfpathmoveto{\pgfqpoint{2.875196in}{1.880235in}}%
\pgfpathlineto{\pgfqpoint{2.875196in}{1.880235in}}%
\pgfpathlineto{\pgfqpoint{2.875196in}{1.883184in}}%
\pgfpathlineto{\pgfqpoint{2.879737in}{1.883184in}}%
\pgfpathlineto{\pgfqpoint{2.879737in}{1.880235in}}%
\pgfpathmoveto{\pgfqpoint{2.875196in}{1.883184in}}%
\pgfpathlineto{\pgfqpoint{2.875196in}{1.883184in}}%
\pgfpathlineto{\pgfqpoint{2.875196in}{1.886134in}}%
\pgfpathlineto{\pgfqpoint{2.879737in}{1.886134in}}%
\pgfpathlineto{\pgfqpoint{2.879737in}{1.883184in}}%
\pgfpathmoveto{\pgfqpoint{2.879737in}{1.880235in}}%
\pgfpathlineto{\pgfqpoint{2.879737in}{1.880235in}}%
\pgfpathlineto{\pgfqpoint{2.879737in}{1.883184in}}%
\pgfpathlineto{\pgfqpoint{2.884278in}{1.883184in}}%
\pgfpathlineto{\pgfqpoint{2.884278in}{1.880235in}}%
\pgfpathmoveto{\pgfqpoint{2.879737in}{1.883184in}}%
\pgfpathlineto{\pgfqpoint{2.879737in}{1.883184in}}%
\pgfpathlineto{\pgfqpoint{2.879737in}{1.886134in}}%
\pgfpathlineto{\pgfqpoint{2.884278in}{1.886134in}}%
\pgfpathlineto{\pgfqpoint{2.884278in}{1.883184in}}%
\pgfpathmoveto{\pgfqpoint{2.884278in}{1.883184in}}%
\pgfpathlineto{\pgfqpoint{2.884278in}{1.883184in}}%
\pgfpathlineto{\pgfqpoint{2.884278in}{1.886134in}}%
\pgfpathlineto{\pgfqpoint{2.888819in}{1.886134in}}%
\pgfpathlineto{\pgfqpoint{2.888819in}{1.883184in}}%
\pgfpathmoveto{\pgfqpoint{2.884278in}{1.886134in}}%
\pgfpathlineto{\pgfqpoint{2.884278in}{1.886134in}}%
\pgfpathlineto{\pgfqpoint{2.884278in}{1.889083in}}%
\pgfpathlineto{\pgfqpoint{2.888819in}{1.889083in}}%
\pgfpathlineto{\pgfqpoint{2.888819in}{1.886134in}}%
\pgfpathmoveto{\pgfqpoint{2.884278in}{1.889083in}}%
\pgfpathlineto{\pgfqpoint{2.884278in}{1.889083in}}%
\pgfpathlineto{\pgfqpoint{2.884278in}{1.892032in}}%
\pgfpathlineto{\pgfqpoint{2.888819in}{1.892032in}}%
\pgfpathlineto{\pgfqpoint{2.888819in}{1.889083in}}%
\pgfpathmoveto{\pgfqpoint{2.888819in}{1.886134in}}%
\pgfpathlineto{\pgfqpoint{2.888819in}{1.886134in}}%
\pgfpathlineto{\pgfqpoint{2.888819in}{1.889083in}}%
\pgfpathlineto{\pgfqpoint{2.893360in}{1.889083in}}%
\pgfpathlineto{\pgfqpoint{2.893360in}{1.886134in}}%
\pgfpathmoveto{\pgfqpoint{2.888819in}{1.889083in}}%
\pgfpathlineto{\pgfqpoint{2.888819in}{1.889083in}}%
\pgfpathlineto{\pgfqpoint{2.888819in}{1.892032in}}%
\pgfpathlineto{\pgfqpoint{2.893360in}{1.892032in}}%
\pgfpathlineto{\pgfqpoint{2.893360in}{1.889083in}}%
\pgfpathmoveto{\pgfqpoint{2.893360in}{1.889083in}}%
\pgfpathlineto{\pgfqpoint{2.893360in}{1.889083in}}%
\pgfpathlineto{\pgfqpoint{2.893360in}{1.892032in}}%
\pgfpathlineto{\pgfqpoint{2.897901in}{1.892032in}}%
\pgfpathlineto{\pgfqpoint{2.897901in}{1.889083in}}%
\pgfpathmoveto{\pgfqpoint{2.893360in}{1.892032in}}%
\pgfpathlineto{\pgfqpoint{2.893360in}{1.892032in}}%
\pgfpathlineto{\pgfqpoint{2.893360in}{1.894981in}}%
\pgfpathlineto{\pgfqpoint{2.897901in}{1.894981in}}%
\pgfpathlineto{\pgfqpoint{2.897901in}{1.892032in}}%
\pgfpathmoveto{\pgfqpoint{2.893360in}{1.894981in}}%
\pgfpathlineto{\pgfqpoint{2.893360in}{1.894981in}}%
\pgfpathlineto{\pgfqpoint{2.893360in}{1.897931in}}%
\pgfpathlineto{\pgfqpoint{2.897901in}{1.897931in}}%
\pgfpathlineto{\pgfqpoint{2.897901in}{1.894981in}}%
\pgfpathmoveto{\pgfqpoint{2.897901in}{1.892032in}}%
\pgfpathlineto{\pgfqpoint{2.897901in}{1.892032in}}%
\pgfpathlineto{\pgfqpoint{2.897901in}{1.894981in}}%
\pgfpathlineto{\pgfqpoint{2.902442in}{1.894981in}}%
\pgfpathlineto{\pgfqpoint{2.902442in}{1.892032in}}%
\pgfpathmoveto{\pgfqpoint{2.897901in}{1.894981in}}%
\pgfpathlineto{\pgfqpoint{2.897901in}{1.894981in}}%
\pgfpathlineto{\pgfqpoint{2.897901in}{1.897931in}}%
\pgfpathlineto{\pgfqpoint{2.902442in}{1.897931in}}%
\pgfpathlineto{\pgfqpoint{2.902442in}{1.894981in}}%
\pgfpathmoveto{\pgfqpoint{2.902442in}{1.894981in}}%
\pgfpathlineto{\pgfqpoint{2.902442in}{1.894981in}}%
\pgfpathlineto{\pgfqpoint{2.902442in}{1.897931in}}%
\pgfpathlineto{\pgfqpoint{2.906983in}{1.897931in}}%
\pgfpathlineto{\pgfqpoint{2.906983in}{1.894981in}}%
\pgfpathmoveto{\pgfqpoint{2.902442in}{1.897931in}}%
\pgfpathlineto{\pgfqpoint{2.902442in}{1.897931in}}%
\pgfpathlineto{\pgfqpoint{2.902442in}{1.900880in}}%
\pgfpathlineto{\pgfqpoint{2.906983in}{1.900880in}}%
\pgfpathlineto{\pgfqpoint{2.906983in}{1.897931in}}%
\pgfpathmoveto{\pgfqpoint{2.902442in}{1.900880in}}%
\pgfpathlineto{\pgfqpoint{2.902442in}{1.900880in}}%
\pgfpathlineto{\pgfqpoint{2.902442in}{1.903829in}}%
\pgfpathlineto{\pgfqpoint{2.906983in}{1.903829in}}%
\pgfpathlineto{\pgfqpoint{2.906983in}{1.900880in}}%
\pgfpathmoveto{\pgfqpoint{2.906983in}{1.897931in}}%
\pgfpathlineto{\pgfqpoint{2.906983in}{1.897931in}}%
\pgfpathlineto{\pgfqpoint{2.906983in}{1.900880in}}%
\pgfpathlineto{\pgfqpoint{2.911524in}{1.900880in}}%
\pgfpathlineto{\pgfqpoint{2.911524in}{1.897931in}}%
\pgfpathmoveto{\pgfqpoint{2.906983in}{1.900880in}}%
\pgfpathlineto{\pgfqpoint{2.906983in}{1.900880in}}%
\pgfpathlineto{\pgfqpoint{2.906983in}{1.903829in}}%
\pgfpathlineto{\pgfqpoint{2.911524in}{1.903829in}}%
\pgfpathlineto{\pgfqpoint{2.911524in}{1.900880in}}%
\pgfpathmoveto{\pgfqpoint{2.911524in}{1.900880in}}%
\pgfpathlineto{\pgfqpoint{2.911524in}{1.900880in}}%
\pgfpathlineto{\pgfqpoint{2.911524in}{1.903829in}}%
\pgfpathlineto{\pgfqpoint{2.916065in}{1.903829in}}%
\pgfpathlineto{\pgfqpoint{2.916065in}{1.900880in}}%
\pgfpathmoveto{\pgfqpoint{2.911524in}{1.903829in}}%
\pgfpathlineto{\pgfqpoint{2.911524in}{1.903829in}}%
\pgfpathlineto{\pgfqpoint{2.911524in}{1.906779in}}%
\pgfpathlineto{\pgfqpoint{2.916065in}{1.906779in}}%
\pgfpathlineto{\pgfqpoint{2.916065in}{1.903829in}}%
\pgfpathmoveto{\pgfqpoint{2.911524in}{1.906779in}}%
\pgfpathlineto{\pgfqpoint{2.911524in}{1.906779in}}%
\pgfpathlineto{\pgfqpoint{2.911524in}{1.909728in}}%
\pgfpathlineto{\pgfqpoint{2.916065in}{1.909728in}}%
\pgfpathlineto{\pgfqpoint{2.916065in}{1.906779in}}%
\pgfpathmoveto{\pgfqpoint{2.916065in}{1.903829in}}%
\pgfpathlineto{\pgfqpoint{2.916065in}{1.903829in}}%
\pgfpathlineto{\pgfqpoint{2.916065in}{1.906779in}}%
\pgfpathlineto{\pgfqpoint{2.920606in}{1.906779in}}%
\pgfpathlineto{\pgfqpoint{2.920606in}{1.903829in}}%
\pgfpathmoveto{\pgfqpoint{2.916065in}{1.906779in}}%
\pgfpathlineto{\pgfqpoint{2.916065in}{1.906779in}}%
\pgfpathlineto{\pgfqpoint{2.916065in}{1.909728in}}%
\pgfpathlineto{\pgfqpoint{2.920606in}{1.909728in}}%
\pgfpathlineto{\pgfqpoint{2.920606in}{1.906779in}}%
\pgfpathmoveto{\pgfqpoint{2.920606in}{1.906779in}}%
\pgfpathlineto{\pgfqpoint{2.920606in}{1.906779in}}%
\pgfpathlineto{\pgfqpoint{2.920606in}{1.909728in}}%
\pgfpathlineto{\pgfqpoint{2.925147in}{1.909728in}}%
\pgfpathlineto{\pgfqpoint{2.925147in}{1.906779in}}%
\pgfpathmoveto{\pgfqpoint{2.920606in}{1.909728in}}%
\pgfpathlineto{\pgfqpoint{2.920606in}{1.909728in}}%
\pgfpathlineto{\pgfqpoint{2.920606in}{1.912677in}}%
\pgfpathlineto{\pgfqpoint{2.925147in}{1.912677in}}%
\pgfpathlineto{\pgfqpoint{2.925147in}{1.909728in}}%
\pgfpathmoveto{\pgfqpoint{2.920606in}{1.912677in}}%
\pgfpathlineto{\pgfqpoint{2.920606in}{1.912677in}}%
\pgfpathlineto{\pgfqpoint{2.920606in}{1.915627in}}%
\pgfpathlineto{\pgfqpoint{2.925147in}{1.915627in}}%
\pgfpathlineto{\pgfqpoint{2.925147in}{1.912677in}}%
\pgfpathmoveto{\pgfqpoint{2.925147in}{1.909728in}}%
\pgfpathlineto{\pgfqpoint{2.925147in}{1.909728in}}%
\pgfpathlineto{\pgfqpoint{2.925147in}{1.912677in}}%
\pgfpathlineto{\pgfqpoint{2.929688in}{1.912677in}}%
\pgfpathlineto{\pgfqpoint{2.929688in}{1.909728in}}%
\pgfpathmoveto{\pgfqpoint{2.925147in}{1.912677in}}%
\pgfpathlineto{\pgfqpoint{2.925147in}{1.912677in}}%
\pgfpathlineto{\pgfqpoint{2.925147in}{1.915627in}}%
\pgfpathlineto{\pgfqpoint{2.929688in}{1.915627in}}%
\pgfpathlineto{\pgfqpoint{2.929688in}{1.912677in}}%
\pgfpathmoveto{\pgfqpoint{2.920606in}{3.402033in}}%
\pgfpathlineto{\pgfqpoint{2.920606in}{3.402033in}}%
\pgfpathlineto{\pgfqpoint{2.920606in}{3.404982in}}%
\pgfpathlineto{\pgfqpoint{2.925147in}{3.404982in}}%
\pgfpathlineto{\pgfqpoint{2.925147in}{3.402033in}}%
\pgfpathmoveto{\pgfqpoint{2.920606in}{3.404982in}}%
\pgfpathlineto{\pgfqpoint{2.920606in}{3.404982in}}%
\pgfpathlineto{\pgfqpoint{2.920606in}{3.407931in}}%
\pgfpathlineto{\pgfqpoint{2.925147in}{3.407931in}}%
\pgfpathlineto{\pgfqpoint{2.925147in}{3.404982in}}%
\pgfpathmoveto{\pgfqpoint{2.925147in}{3.402033in}}%
\pgfpathlineto{\pgfqpoint{2.925147in}{3.402033in}}%
\pgfpathlineto{\pgfqpoint{2.925147in}{3.404982in}}%
\pgfpathlineto{\pgfqpoint{2.929688in}{3.404982in}}%
\pgfpathlineto{\pgfqpoint{2.929688in}{3.402033in}}%
\pgfpathmoveto{\pgfqpoint{2.925147in}{3.404982in}}%
\pgfpathlineto{\pgfqpoint{2.925147in}{3.404982in}}%
\pgfpathlineto{\pgfqpoint{2.925147in}{3.407931in}}%
\pgfpathlineto{\pgfqpoint{2.929688in}{3.407931in}}%
\pgfpathlineto{\pgfqpoint{2.929688in}{3.404982in}}%
\pgfpathmoveto{\pgfqpoint{2.920606in}{3.407931in}}%
\pgfpathlineto{\pgfqpoint{2.920606in}{3.407931in}}%
\pgfpathlineto{\pgfqpoint{2.920606in}{3.410881in}}%
\pgfpathlineto{\pgfqpoint{2.925147in}{3.410881in}}%
\pgfpathlineto{\pgfqpoint{2.925147in}{3.407931in}}%
\pgfpathmoveto{\pgfqpoint{2.920606in}{3.410881in}}%
\pgfpathlineto{\pgfqpoint{2.920606in}{3.410881in}}%
\pgfpathlineto{\pgfqpoint{2.920606in}{3.413830in}}%
\pgfpathlineto{\pgfqpoint{2.925147in}{3.413830in}}%
\pgfpathlineto{\pgfqpoint{2.925147in}{3.410881in}}%
\pgfpathmoveto{\pgfqpoint{2.925147in}{3.407931in}}%
\pgfpathlineto{\pgfqpoint{2.925147in}{3.407931in}}%
\pgfpathlineto{\pgfqpoint{2.925147in}{3.410881in}}%
\pgfpathlineto{\pgfqpoint{2.929688in}{3.410881in}}%
\pgfpathlineto{\pgfqpoint{2.929688in}{3.407931in}}%
\pgfpathmoveto{\pgfqpoint{2.925147in}{3.410881in}}%
\pgfpathlineto{\pgfqpoint{2.925147in}{3.410881in}}%
\pgfpathlineto{\pgfqpoint{2.925147in}{3.413830in}}%
\pgfpathlineto{\pgfqpoint{2.929688in}{3.413830in}}%
\pgfpathlineto{\pgfqpoint{2.929688in}{3.410881in}}%
\pgfpathmoveto{\pgfqpoint{2.911524in}{3.413830in}}%
\pgfpathlineto{\pgfqpoint{2.911524in}{3.413830in}}%
\pgfpathlineto{\pgfqpoint{2.911524in}{3.416779in}}%
\pgfpathlineto{\pgfqpoint{2.916065in}{3.416779in}}%
\pgfpathlineto{\pgfqpoint{2.916065in}{3.413830in}}%
\pgfpathmoveto{\pgfqpoint{2.911524in}{3.416779in}}%
\pgfpathlineto{\pgfqpoint{2.911524in}{3.416779in}}%
\pgfpathlineto{\pgfqpoint{2.911524in}{3.419728in}}%
\pgfpathlineto{\pgfqpoint{2.916065in}{3.419728in}}%
\pgfpathlineto{\pgfqpoint{2.916065in}{3.416779in}}%
\pgfpathmoveto{\pgfqpoint{2.916065in}{3.413830in}}%
\pgfpathlineto{\pgfqpoint{2.916065in}{3.413830in}}%
\pgfpathlineto{\pgfqpoint{2.916065in}{3.416779in}}%
\pgfpathlineto{\pgfqpoint{2.920606in}{3.416779in}}%
\pgfpathlineto{\pgfqpoint{2.920606in}{3.413830in}}%
\pgfpathmoveto{\pgfqpoint{2.916065in}{3.416779in}}%
\pgfpathlineto{\pgfqpoint{2.916065in}{3.416779in}}%
\pgfpathlineto{\pgfqpoint{2.916065in}{3.419728in}}%
\pgfpathlineto{\pgfqpoint{2.920606in}{3.419728in}}%
\pgfpathlineto{\pgfqpoint{2.920606in}{3.416779in}}%
\pgfpathmoveto{\pgfqpoint{2.911524in}{3.419728in}}%
\pgfpathlineto{\pgfqpoint{2.911524in}{3.419728in}}%
\pgfpathlineto{\pgfqpoint{2.911524in}{3.422677in}}%
\pgfpathlineto{\pgfqpoint{2.916065in}{3.422677in}}%
\pgfpathlineto{\pgfqpoint{2.916065in}{3.419728in}}%
\pgfpathmoveto{\pgfqpoint{2.911524in}{3.422677in}}%
\pgfpathlineto{\pgfqpoint{2.911524in}{3.422677in}}%
\pgfpathlineto{\pgfqpoint{2.911524in}{3.425627in}}%
\pgfpathlineto{\pgfqpoint{2.916065in}{3.425627in}}%
\pgfpathlineto{\pgfqpoint{2.916065in}{3.422677in}}%
\pgfpathmoveto{\pgfqpoint{2.916065in}{3.419728in}}%
\pgfpathlineto{\pgfqpoint{2.916065in}{3.419728in}}%
\pgfpathlineto{\pgfqpoint{2.916065in}{3.422677in}}%
\pgfpathlineto{\pgfqpoint{2.920606in}{3.422677in}}%
\pgfpathlineto{\pgfqpoint{2.920606in}{3.419728in}}%
\pgfpathmoveto{\pgfqpoint{2.916065in}{3.422677in}}%
\pgfpathlineto{\pgfqpoint{2.916065in}{3.422677in}}%
\pgfpathlineto{\pgfqpoint{2.916065in}{3.425627in}}%
\pgfpathlineto{\pgfqpoint{2.920606in}{3.425627in}}%
\pgfpathlineto{\pgfqpoint{2.920606in}{3.422677in}}%
\pgfpathmoveto{\pgfqpoint{2.920606in}{3.413830in}}%
\pgfpathlineto{\pgfqpoint{2.920606in}{3.413830in}}%
\pgfpathlineto{\pgfqpoint{2.920606in}{3.416779in}}%
\pgfpathlineto{\pgfqpoint{2.925147in}{3.416779in}}%
\pgfpathlineto{\pgfqpoint{2.925147in}{3.413830in}}%
\pgfpathmoveto{\pgfqpoint{2.920606in}{3.416779in}}%
\pgfpathlineto{\pgfqpoint{2.920606in}{3.416779in}}%
\pgfpathlineto{\pgfqpoint{2.920606in}{3.419728in}}%
\pgfpathlineto{\pgfqpoint{2.925147in}{3.419728in}}%
\pgfpathlineto{\pgfqpoint{2.925147in}{3.416779in}}%
\pgfpathmoveto{\pgfqpoint{2.847949in}{3.496405in}}%
\pgfpathlineto{\pgfqpoint{2.847949in}{3.496405in}}%
\pgfpathlineto{\pgfqpoint{2.847949in}{3.499354in}}%
\pgfpathlineto{\pgfqpoint{2.852490in}{3.499354in}}%
\pgfpathlineto{\pgfqpoint{2.852490in}{3.496405in}}%
\pgfpathmoveto{\pgfqpoint{2.847949in}{3.499354in}}%
\pgfpathlineto{\pgfqpoint{2.847949in}{3.499354in}}%
\pgfpathlineto{\pgfqpoint{2.847949in}{3.502304in}}%
\pgfpathlineto{\pgfqpoint{2.852490in}{3.502304in}}%
\pgfpathlineto{\pgfqpoint{2.852490in}{3.499354in}}%
\pgfpathmoveto{\pgfqpoint{2.852490in}{3.496405in}}%
\pgfpathlineto{\pgfqpoint{2.852490in}{3.496405in}}%
\pgfpathlineto{\pgfqpoint{2.852490in}{3.499354in}}%
\pgfpathlineto{\pgfqpoint{2.857031in}{3.499354in}}%
\pgfpathlineto{\pgfqpoint{2.857031in}{3.496405in}}%
\pgfpathmoveto{\pgfqpoint{2.852490in}{3.499354in}}%
\pgfpathlineto{\pgfqpoint{2.852490in}{3.499354in}}%
\pgfpathlineto{\pgfqpoint{2.852490in}{3.502304in}}%
\pgfpathlineto{\pgfqpoint{2.857031in}{3.502304in}}%
\pgfpathlineto{\pgfqpoint{2.857031in}{3.499354in}}%
\pgfpathmoveto{\pgfqpoint{2.847949in}{3.502304in}}%
\pgfpathlineto{\pgfqpoint{2.847949in}{3.502304in}}%
\pgfpathlineto{\pgfqpoint{2.847949in}{3.505253in}}%
\pgfpathlineto{\pgfqpoint{2.852490in}{3.505253in}}%
\pgfpathlineto{\pgfqpoint{2.852490in}{3.502304in}}%
\pgfpathmoveto{\pgfqpoint{2.847949in}{3.505253in}}%
\pgfpathlineto{\pgfqpoint{2.847949in}{3.505253in}}%
\pgfpathlineto{\pgfqpoint{2.847949in}{3.508202in}}%
\pgfpathlineto{\pgfqpoint{2.852490in}{3.508202in}}%
\pgfpathlineto{\pgfqpoint{2.852490in}{3.505253in}}%
\pgfpathmoveto{\pgfqpoint{2.852490in}{3.502304in}}%
\pgfpathlineto{\pgfqpoint{2.852490in}{3.502304in}}%
\pgfpathlineto{\pgfqpoint{2.852490in}{3.505253in}}%
\pgfpathlineto{\pgfqpoint{2.857031in}{3.505253in}}%
\pgfpathlineto{\pgfqpoint{2.857031in}{3.502304in}}%
\pgfpathmoveto{\pgfqpoint{2.852490in}{3.505253in}}%
\pgfpathlineto{\pgfqpoint{2.852490in}{3.505253in}}%
\pgfpathlineto{\pgfqpoint{2.852490in}{3.508202in}}%
\pgfpathlineto{\pgfqpoint{2.857031in}{3.508202in}}%
\pgfpathlineto{\pgfqpoint{2.857031in}{3.505253in}}%
\pgfpathmoveto{\pgfqpoint{2.838867in}{3.508202in}}%
\pgfpathlineto{\pgfqpoint{2.838867in}{3.508202in}}%
\pgfpathlineto{\pgfqpoint{2.838867in}{3.511151in}}%
\pgfpathlineto{\pgfqpoint{2.843408in}{3.511151in}}%
\pgfpathlineto{\pgfqpoint{2.843408in}{3.508202in}}%
\pgfpathmoveto{\pgfqpoint{2.838867in}{3.511151in}}%
\pgfpathlineto{\pgfqpoint{2.838867in}{3.511151in}}%
\pgfpathlineto{\pgfqpoint{2.838867in}{3.514100in}}%
\pgfpathlineto{\pgfqpoint{2.843408in}{3.514100in}}%
\pgfpathlineto{\pgfqpoint{2.843408in}{3.511151in}}%
\pgfpathmoveto{\pgfqpoint{2.843408in}{3.508202in}}%
\pgfpathlineto{\pgfqpoint{2.843408in}{3.508202in}}%
\pgfpathlineto{\pgfqpoint{2.843408in}{3.511151in}}%
\pgfpathlineto{\pgfqpoint{2.847949in}{3.511151in}}%
\pgfpathlineto{\pgfqpoint{2.847949in}{3.508202in}}%
\pgfpathmoveto{\pgfqpoint{2.843408in}{3.511151in}}%
\pgfpathlineto{\pgfqpoint{2.843408in}{3.511151in}}%
\pgfpathlineto{\pgfqpoint{2.843408in}{3.514100in}}%
\pgfpathlineto{\pgfqpoint{2.847949in}{3.514100in}}%
\pgfpathlineto{\pgfqpoint{2.847949in}{3.511151in}}%
\pgfpathmoveto{\pgfqpoint{2.838867in}{3.514100in}}%
\pgfpathlineto{\pgfqpoint{2.838867in}{3.514100in}}%
\pgfpathlineto{\pgfqpoint{2.838867in}{3.517049in}}%
\pgfpathlineto{\pgfqpoint{2.843408in}{3.517049in}}%
\pgfpathlineto{\pgfqpoint{2.843408in}{3.514100in}}%
\pgfpathmoveto{\pgfqpoint{2.838867in}{3.517049in}}%
\pgfpathlineto{\pgfqpoint{2.838867in}{3.517049in}}%
\pgfpathlineto{\pgfqpoint{2.838867in}{3.519998in}}%
\pgfpathlineto{\pgfqpoint{2.843408in}{3.519998in}}%
\pgfpathlineto{\pgfqpoint{2.843408in}{3.517049in}}%
\pgfpathmoveto{\pgfqpoint{2.843408in}{3.514100in}}%
\pgfpathlineto{\pgfqpoint{2.843408in}{3.514100in}}%
\pgfpathlineto{\pgfqpoint{2.843408in}{3.517049in}}%
\pgfpathlineto{\pgfqpoint{2.847949in}{3.517049in}}%
\pgfpathlineto{\pgfqpoint{2.847949in}{3.514100in}}%
\pgfpathmoveto{\pgfqpoint{2.843408in}{3.517049in}}%
\pgfpathlineto{\pgfqpoint{2.843408in}{3.517049in}}%
\pgfpathlineto{\pgfqpoint{2.843408in}{3.519998in}}%
\pgfpathlineto{\pgfqpoint{2.847949in}{3.519998in}}%
\pgfpathlineto{\pgfqpoint{2.847949in}{3.517049in}}%
\pgfpathmoveto{\pgfqpoint{2.847949in}{3.508202in}}%
\pgfpathlineto{\pgfqpoint{2.847949in}{3.508202in}}%
\pgfpathlineto{\pgfqpoint{2.847949in}{3.511151in}}%
\pgfpathlineto{\pgfqpoint{2.852490in}{3.511151in}}%
\pgfpathlineto{\pgfqpoint{2.852490in}{3.508202in}}%
\pgfpathmoveto{\pgfqpoint{2.847949in}{3.511151in}}%
\pgfpathlineto{\pgfqpoint{2.847949in}{3.511151in}}%
\pgfpathlineto{\pgfqpoint{2.847949in}{3.514100in}}%
\pgfpathlineto{\pgfqpoint{2.852490in}{3.514100in}}%
\pgfpathlineto{\pgfqpoint{2.852490in}{3.511151in}}%
\pgfpathmoveto{\pgfqpoint{2.884278in}{3.449220in}}%
\pgfpathlineto{\pgfqpoint{2.884278in}{3.449220in}}%
\pgfpathlineto{\pgfqpoint{2.884278in}{3.452169in}}%
\pgfpathlineto{\pgfqpoint{2.888819in}{3.452169in}}%
\pgfpathlineto{\pgfqpoint{2.888819in}{3.449220in}}%
\pgfpathmoveto{\pgfqpoint{2.884278in}{3.452169in}}%
\pgfpathlineto{\pgfqpoint{2.884278in}{3.452169in}}%
\pgfpathlineto{\pgfqpoint{2.884278in}{3.455118in}}%
\pgfpathlineto{\pgfqpoint{2.888819in}{3.455118in}}%
\pgfpathlineto{\pgfqpoint{2.888819in}{3.452169in}}%
\pgfpathmoveto{\pgfqpoint{2.888819in}{3.449220in}}%
\pgfpathlineto{\pgfqpoint{2.888819in}{3.449220in}}%
\pgfpathlineto{\pgfqpoint{2.888819in}{3.452169in}}%
\pgfpathlineto{\pgfqpoint{2.893360in}{3.452169in}}%
\pgfpathlineto{\pgfqpoint{2.893360in}{3.449220in}}%
\pgfpathmoveto{\pgfqpoint{2.888819in}{3.452169in}}%
\pgfpathlineto{\pgfqpoint{2.888819in}{3.452169in}}%
\pgfpathlineto{\pgfqpoint{2.888819in}{3.455118in}}%
\pgfpathlineto{\pgfqpoint{2.893360in}{3.455118in}}%
\pgfpathlineto{\pgfqpoint{2.893360in}{3.452169in}}%
\pgfpathmoveto{\pgfqpoint{2.884278in}{3.455118in}}%
\pgfpathlineto{\pgfqpoint{2.884278in}{3.455118in}}%
\pgfpathlineto{\pgfqpoint{2.884278in}{3.458067in}}%
\pgfpathlineto{\pgfqpoint{2.888819in}{3.458067in}}%
\pgfpathlineto{\pgfqpoint{2.888819in}{3.455118in}}%
\pgfpathmoveto{\pgfqpoint{2.884278in}{3.458067in}}%
\pgfpathlineto{\pgfqpoint{2.884278in}{3.458067in}}%
\pgfpathlineto{\pgfqpoint{2.884278in}{3.461016in}}%
\pgfpathlineto{\pgfqpoint{2.888819in}{3.461016in}}%
\pgfpathlineto{\pgfqpoint{2.888819in}{3.458067in}}%
\pgfpathmoveto{\pgfqpoint{2.888819in}{3.455118in}}%
\pgfpathlineto{\pgfqpoint{2.888819in}{3.455118in}}%
\pgfpathlineto{\pgfqpoint{2.888819in}{3.458067in}}%
\pgfpathlineto{\pgfqpoint{2.893360in}{3.458067in}}%
\pgfpathlineto{\pgfqpoint{2.893360in}{3.455118in}}%
\pgfpathmoveto{\pgfqpoint{2.888819in}{3.458067in}}%
\pgfpathlineto{\pgfqpoint{2.888819in}{3.458067in}}%
\pgfpathlineto{\pgfqpoint{2.888819in}{3.461016in}}%
\pgfpathlineto{\pgfqpoint{2.893360in}{3.461016in}}%
\pgfpathlineto{\pgfqpoint{2.893360in}{3.458067in}}%
\pgfpathmoveto{\pgfqpoint{2.875196in}{3.461016in}}%
\pgfpathlineto{\pgfqpoint{2.875196in}{3.461016in}}%
\pgfpathlineto{\pgfqpoint{2.875196in}{3.463965in}}%
\pgfpathlineto{\pgfqpoint{2.879737in}{3.463965in}}%
\pgfpathlineto{\pgfqpoint{2.879737in}{3.461016in}}%
\pgfpathmoveto{\pgfqpoint{2.875196in}{3.463965in}}%
\pgfpathlineto{\pgfqpoint{2.875196in}{3.463965in}}%
\pgfpathlineto{\pgfqpoint{2.875196in}{3.466914in}}%
\pgfpathlineto{\pgfqpoint{2.879737in}{3.466914in}}%
\pgfpathlineto{\pgfqpoint{2.879737in}{3.463965in}}%
\pgfpathmoveto{\pgfqpoint{2.879737in}{3.461016in}}%
\pgfpathlineto{\pgfqpoint{2.879737in}{3.461016in}}%
\pgfpathlineto{\pgfqpoint{2.879737in}{3.463965in}}%
\pgfpathlineto{\pgfqpoint{2.884278in}{3.463965in}}%
\pgfpathlineto{\pgfqpoint{2.884278in}{3.461016in}}%
\pgfpathmoveto{\pgfqpoint{2.879737in}{3.463965in}}%
\pgfpathlineto{\pgfqpoint{2.879737in}{3.463965in}}%
\pgfpathlineto{\pgfqpoint{2.879737in}{3.466914in}}%
\pgfpathlineto{\pgfqpoint{2.884278in}{3.466914in}}%
\pgfpathlineto{\pgfqpoint{2.884278in}{3.463965in}}%
\pgfpathmoveto{\pgfqpoint{2.875196in}{3.466914in}}%
\pgfpathlineto{\pgfqpoint{2.875196in}{3.466914in}}%
\pgfpathlineto{\pgfqpoint{2.875196in}{3.469863in}}%
\pgfpathlineto{\pgfqpoint{2.879737in}{3.469863in}}%
\pgfpathlineto{\pgfqpoint{2.879737in}{3.466914in}}%
\pgfpathmoveto{\pgfqpoint{2.875196in}{3.469863in}}%
\pgfpathlineto{\pgfqpoint{2.875196in}{3.469863in}}%
\pgfpathlineto{\pgfqpoint{2.875196in}{3.472812in}}%
\pgfpathlineto{\pgfqpoint{2.879737in}{3.472812in}}%
\pgfpathlineto{\pgfqpoint{2.879737in}{3.469863in}}%
\pgfpathmoveto{\pgfqpoint{2.879737in}{3.466914in}}%
\pgfpathlineto{\pgfqpoint{2.879737in}{3.466914in}}%
\pgfpathlineto{\pgfqpoint{2.879737in}{3.469863in}}%
\pgfpathlineto{\pgfqpoint{2.884278in}{3.469863in}}%
\pgfpathlineto{\pgfqpoint{2.884278in}{3.466914in}}%
\pgfpathmoveto{\pgfqpoint{2.879737in}{3.469863in}}%
\pgfpathlineto{\pgfqpoint{2.879737in}{3.469863in}}%
\pgfpathlineto{\pgfqpoint{2.879737in}{3.472812in}}%
\pgfpathlineto{\pgfqpoint{2.884278in}{3.472812in}}%
\pgfpathlineto{\pgfqpoint{2.884278in}{3.469863in}}%
\pgfpathmoveto{\pgfqpoint{2.884278in}{3.461016in}}%
\pgfpathlineto{\pgfqpoint{2.884278in}{3.461016in}}%
\pgfpathlineto{\pgfqpoint{2.884278in}{3.463965in}}%
\pgfpathlineto{\pgfqpoint{2.888819in}{3.463965in}}%
\pgfpathlineto{\pgfqpoint{2.888819in}{3.461016in}}%
\pgfpathmoveto{\pgfqpoint{2.884278in}{3.463965in}}%
\pgfpathlineto{\pgfqpoint{2.884278in}{3.463965in}}%
\pgfpathlineto{\pgfqpoint{2.884278in}{3.466914in}}%
\pgfpathlineto{\pgfqpoint{2.888819in}{3.466914in}}%
\pgfpathlineto{\pgfqpoint{2.888819in}{3.463965in}}%
\pgfpathmoveto{\pgfqpoint{2.902442in}{3.425627in}}%
\pgfpathlineto{\pgfqpoint{2.902442in}{3.425627in}}%
\pgfpathlineto{\pgfqpoint{2.902442in}{3.428576in}}%
\pgfpathlineto{\pgfqpoint{2.906983in}{3.428576in}}%
\pgfpathlineto{\pgfqpoint{2.906983in}{3.425627in}}%
\pgfpathmoveto{\pgfqpoint{2.902442in}{3.428576in}}%
\pgfpathlineto{\pgfqpoint{2.902442in}{3.428576in}}%
\pgfpathlineto{\pgfqpoint{2.902442in}{3.431525in}}%
\pgfpathlineto{\pgfqpoint{2.906983in}{3.431525in}}%
\pgfpathlineto{\pgfqpoint{2.906983in}{3.428576in}}%
\pgfpathmoveto{\pgfqpoint{2.906983in}{3.425627in}}%
\pgfpathlineto{\pgfqpoint{2.906983in}{3.425627in}}%
\pgfpathlineto{\pgfqpoint{2.906983in}{3.428576in}}%
\pgfpathlineto{\pgfqpoint{2.911524in}{3.428576in}}%
\pgfpathlineto{\pgfqpoint{2.911524in}{3.425627in}}%
\pgfpathmoveto{\pgfqpoint{2.906983in}{3.428576in}}%
\pgfpathlineto{\pgfqpoint{2.906983in}{3.428576in}}%
\pgfpathlineto{\pgfqpoint{2.906983in}{3.431525in}}%
\pgfpathlineto{\pgfqpoint{2.911524in}{3.431525in}}%
\pgfpathlineto{\pgfqpoint{2.911524in}{3.428576in}}%
\pgfpathmoveto{\pgfqpoint{2.902442in}{3.431525in}}%
\pgfpathlineto{\pgfqpoint{2.902442in}{3.431525in}}%
\pgfpathlineto{\pgfqpoint{2.902442in}{3.434474in}}%
\pgfpathlineto{\pgfqpoint{2.906983in}{3.434474in}}%
\pgfpathlineto{\pgfqpoint{2.906983in}{3.431525in}}%
\pgfpathmoveto{\pgfqpoint{2.902442in}{3.434474in}}%
\pgfpathlineto{\pgfqpoint{2.902442in}{3.434474in}}%
\pgfpathlineto{\pgfqpoint{2.902442in}{3.437423in}}%
\pgfpathlineto{\pgfqpoint{2.906983in}{3.437423in}}%
\pgfpathlineto{\pgfqpoint{2.906983in}{3.434474in}}%
\pgfpathmoveto{\pgfqpoint{2.906983in}{3.431525in}}%
\pgfpathlineto{\pgfqpoint{2.906983in}{3.431525in}}%
\pgfpathlineto{\pgfqpoint{2.906983in}{3.434474in}}%
\pgfpathlineto{\pgfqpoint{2.911524in}{3.434474in}}%
\pgfpathlineto{\pgfqpoint{2.911524in}{3.431525in}}%
\pgfpathmoveto{\pgfqpoint{2.906983in}{3.434474in}}%
\pgfpathlineto{\pgfqpoint{2.906983in}{3.434474in}}%
\pgfpathlineto{\pgfqpoint{2.906983in}{3.437423in}}%
\pgfpathlineto{\pgfqpoint{2.911524in}{3.437423in}}%
\pgfpathlineto{\pgfqpoint{2.911524in}{3.434474in}}%
\pgfpathmoveto{\pgfqpoint{2.893360in}{3.437423in}}%
\pgfpathlineto{\pgfqpoint{2.893360in}{3.437423in}}%
\pgfpathlineto{\pgfqpoint{2.893360in}{3.440372in}}%
\pgfpathlineto{\pgfqpoint{2.897901in}{3.440372in}}%
\pgfpathlineto{\pgfqpoint{2.897901in}{3.437423in}}%
\pgfpathmoveto{\pgfqpoint{2.893360in}{3.440372in}}%
\pgfpathlineto{\pgfqpoint{2.893360in}{3.440372in}}%
\pgfpathlineto{\pgfqpoint{2.893360in}{3.443321in}}%
\pgfpathlineto{\pgfqpoint{2.897901in}{3.443321in}}%
\pgfpathlineto{\pgfqpoint{2.897901in}{3.440372in}}%
\pgfpathmoveto{\pgfqpoint{2.897901in}{3.437423in}}%
\pgfpathlineto{\pgfqpoint{2.897901in}{3.437423in}}%
\pgfpathlineto{\pgfqpoint{2.897901in}{3.440372in}}%
\pgfpathlineto{\pgfqpoint{2.902442in}{3.440372in}}%
\pgfpathlineto{\pgfqpoint{2.902442in}{3.437423in}}%
\pgfpathmoveto{\pgfqpoint{2.897901in}{3.440372in}}%
\pgfpathlineto{\pgfqpoint{2.897901in}{3.440372in}}%
\pgfpathlineto{\pgfqpoint{2.897901in}{3.443321in}}%
\pgfpathlineto{\pgfqpoint{2.902442in}{3.443321in}}%
\pgfpathlineto{\pgfqpoint{2.902442in}{3.440372in}}%
\pgfpathmoveto{\pgfqpoint{2.893360in}{3.443321in}}%
\pgfpathlineto{\pgfqpoint{2.893360in}{3.443321in}}%
\pgfpathlineto{\pgfqpoint{2.893360in}{3.446270in}}%
\pgfpathlineto{\pgfqpoint{2.897901in}{3.446270in}}%
\pgfpathlineto{\pgfqpoint{2.897901in}{3.443321in}}%
\pgfpathmoveto{\pgfqpoint{2.893360in}{3.446270in}}%
\pgfpathlineto{\pgfqpoint{2.893360in}{3.446270in}}%
\pgfpathlineto{\pgfqpoint{2.893360in}{3.449220in}}%
\pgfpathlineto{\pgfqpoint{2.897901in}{3.449220in}}%
\pgfpathlineto{\pgfqpoint{2.897901in}{3.446270in}}%
\pgfpathmoveto{\pgfqpoint{2.897901in}{3.443321in}}%
\pgfpathlineto{\pgfqpoint{2.897901in}{3.443321in}}%
\pgfpathlineto{\pgfqpoint{2.897901in}{3.446270in}}%
\pgfpathlineto{\pgfqpoint{2.902442in}{3.446270in}}%
\pgfpathlineto{\pgfqpoint{2.902442in}{3.443321in}}%
\pgfpathmoveto{\pgfqpoint{2.897901in}{3.446270in}}%
\pgfpathlineto{\pgfqpoint{2.897901in}{3.446270in}}%
\pgfpathlineto{\pgfqpoint{2.897901in}{3.449220in}}%
\pgfpathlineto{\pgfqpoint{2.902442in}{3.449220in}}%
\pgfpathlineto{\pgfqpoint{2.902442in}{3.446270in}}%
\pgfpathmoveto{\pgfqpoint{2.902442in}{3.437423in}}%
\pgfpathlineto{\pgfqpoint{2.902442in}{3.437423in}}%
\pgfpathlineto{\pgfqpoint{2.902442in}{3.440372in}}%
\pgfpathlineto{\pgfqpoint{2.906983in}{3.440372in}}%
\pgfpathlineto{\pgfqpoint{2.906983in}{3.437423in}}%
\pgfpathmoveto{\pgfqpoint{2.902442in}{3.440372in}}%
\pgfpathlineto{\pgfqpoint{2.902442in}{3.440372in}}%
\pgfpathlineto{\pgfqpoint{2.902442in}{3.443321in}}%
\pgfpathlineto{\pgfqpoint{2.906983in}{3.443321in}}%
\pgfpathlineto{\pgfqpoint{2.906983in}{3.440372in}}%
\pgfpathmoveto{\pgfqpoint{2.911524in}{3.425627in}}%
\pgfpathlineto{\pgfqpoint{2.911524in}{3.425627in}}%
\pgfpathlineto{\pgfqpoint{2.911524in}{3.428576in}}%
\pgfpathlineto{\pgfqpoint{2.916065in}{3.428576in}}%
\pgfpathlineto{\pgfqpoint{2.916065in}{3.425627in}}%
\pgfpathmoveto{\pgfqpoint{2.911524in}{3.428576in}}%
\pgfpathlineto{\pgfqpoint{2.911524in}{3.428576in}}%
\pgfpathlineto{\pgfqpoint{2.911524in}{3.431525in}}%
\pgfpathlineto{\pgfqpoint{2.916065in}{3.431525in}}%
\pgfpathlineto{\pgfqpoint{2.916065in}{3.428576in}}%
\pgfpathmoveto{\pgfqpoint{2.893360in}{3.449220in}}%
\pgfpathlineto{\pgfqpoint{2.893360in}{3.449220in}}%
\pgfpathlineto{\pgfqpoint{2.893360in}{3.452169in}}%
\pgfpathlineto{\pgfqpoint{2.897901in}{3.452169in}}%
\pgfpathlineto{\pgfqpoint{2.897901in}{3.449220in}}%
\pgfpathmoveto{\pgfqpoint{2.893360in}{3.452169in}}%
\pgfpathlineto{\pgfqpoint{2.893360in}{3.452169in}}%
\pgfpathlineto{\pgfqpoint{2.893360in}{3.455118in}}%
\pgfpathlineto{\pgfqpoint{2.897901in}{3.455118in}}%
\pgfpathlineto{\pgfqpoint{2.897901in}{3.452169in}}%
\pgfpathmoveto{\pgfqpoint{2.866113in}{3.472812in}}%
\pgfpathlineto{\pgfqpoint{2.866113in}{3.472812in}}%
\pgfpathlineto{\pgfqpoint{2.866113in}{3.475762in}}%
\pgfpathlineto{\pgfqpoint{2.870655in}{3.475762in}}%
\pgfpathlineto{\pgfqpoint{2.870655in}{3.472812in}}%
\pgfpathmoveto{\pgfqpoint{2.866113in}{3.475762in}}%
\pgfpathlineto{\pgfqpoint{2.866113in}{3.475762in}}%
\pgfpathlineto{\pgfqpoint{2.866113in}{3.478711in}}%
\pgfpathlineto{\pgfqpoint{2.870655in}{3.478711in}}%
\pgfpathlineto{\pgfqpoint{2.870655in}{3.475762in}}%
\pgfpathmoveto{\pgfqpoint{2.870655in}{3.472812in}}%
\pgfpathlineto{\pgfqpoint{2.870655in}{3.472812in}}%
\pgfpathlineto{\pgfqpoint{2.870655in}{3.475762in}}%
\pgfpathlineto{\pgfqpoint{2.875196in}{3.475762in}}%
\pgfpathlineto{\pgfqpoint{2.875196in}{3.472812in}}%
\pgfpathmoveto{\pgfqpoint{2.870655in}{3.475762in}}%
\pgfpathlineto{\pgfqpoint{2.870655in}{3.475762in}}%
\pgfpathlineto{\pgfqpoint{2.870655in}{3.478711in}}%
\pgfpathlineto{\pgfqpoint{2.875196in}{3.478711in}}%
\pgfpathlineto{\pgfqpoint{2.875196in}{3.475762in}}%
\pgfpathmoveto{\pgfqpoint{2.866113in}{3.478711in}}%
\pgfpathlineto{\pgfqpoint{2.866113in}{3.478711in}}%
\pgfpathlineto{\pgfqpoint{2.866113in}{3.481660in}}%
\pgfpathlineto{\pgfqpoint{2.870655in}{3.481660in}}%
\pgfpathlineto{\pgfqpoint{2.870655in}{3.478711in}}%
\pgfpathmoveto{\pgfqpoint{2.866113in}{3.481660in}}%
\pgfpathlineto{\pgfqpoint{2.866113in}{3.481660in}}%
\pgfpathlineto{\pgfqpoint{2.866113in}{3.484609in}}%
\pgfpathlineto{\pgfqpoint{2.870655in}{3.484609in}}%
\pgfpathlineto{\pgfqpoint{2.870655in}{3.481660in}}%
\pgfpathmoveto{\pgfqpoint{2.870655in}{3.478711in}}%
\pgfpathlineto{\pgfqpoint{2.870655in}{3.478711in}}%
\pgfpathlineto{\pgfqpoint{2.870655in}{3.481660in}}%
\pgfpathlineto{\pgfqpoint{2.875196in}{3.481660in}}%
\pgfpathlineto{\pgfqpoint{2.875196in}{3.478711in}}%
\pgfpathmoveto{\pgfqpoint{2.870655in}{3.481660in}}%
\pgfpathlineto{\pgfqpoint{2.870655in}{3.481660in}}%
\pgfpathlineto{\pgfqpoint{2.870655in}{3.484609in}}%
\pgfpathlineto{\pgfqpoint{2.875196in}{3.484609in}}%
\pgfpathlineto{\pgfqpoint{2.875196in}{3.481660in}}%
\pgfpathmoveto{\pgfqpoint{2.857031in}{3.484609in}}%
\pgfpathlineto{\pgfqpoint{2.857031in}{3.484609in}}%
\pgfpathlineto{\pgfqpoint{2.857031in}{3.487558in}}%
\pgfpathlineto{\pgfqpoint{2.861572in}{3.487558in}}%
\pgfpathlineto{\pgfqpoint{2.861572in}{3.484609in}}%
\pgfpathmoveto{\pgfqpoint{2.857031in}{3.487558in}}%
\pgfpathlineto{\pgfqpoint{2.857031in}{3.487558in}}%
\pgfpathlineto{\pgfqpoint{2.857031in}{3.490507in}}%
\pgfpathlineto{\pgfqpoint{2.861572in}{3.490507in}}%
\pgfpathlineto{\pgfqpoint{2.861572in}{3.487558in}}%
\pgfpathmoveto{\pgfqpoint{2.861572in}{3.484609in}}%
\pgfpathlineto{\pgfqpoint{2.861572in}{3.484609in}}%
\pgfpathlineto{\pgfqpoint{2.861572in}{3.487558in}}%
\pgfpathlineto{\pgfqpoint{2.866113in}{3.487558in}}%
\pgfpathlineto{\pgfqpoint{2.866113in}{3.484609in}}%
\pgfpathmoveto{\pgfqpoint{2.861572in}{3.487558in}}%
\pgfpathlineto{\pgfqpoint{2.861572in}{3.487558in}}%
\pgfpathlineto{\pgfqpoint{2.861572in}{3.490507in}}%
\pgfpathlineto{\pgfqpoint{2.866113in}{3.490507in}}%
\pgfpathlineto{\pgfqpoint{2.866113in}{3.487558in}}%
\pgfpathmoveto{\pgfqpoint{2.857031in}{3.490507in}}%
\pgfpathlineto{\pgfqpoint{2.857031in}{3.490507in}}%
\pgfpathlineto{\pgfqpoint{2.857031in}{3.493456in}}%
\pgfpathlineto{\pgfqpoint{2.861572in}{3.493456in}}%
\pgfpathlineto{\pgfqpoint{2.861572in}{3.490507in}}%
\pgfpathmoveto{\pgfqpoint{2.857031in}{3.493456in}}%
\pgfpathlineto{\pgfqpoint{2.857031in}{3.493456in}}%
\pgfpathlineto{\pgfqpoint{2.857031in}{3.496405in}}%
\pgfpathlineto{\pgfqpoint{2.861572in}{3.496405in}}%
\pgfpathlineto{\pgfqpoint{2.861572in}{3.493456in}}%
\pgfpathmoveto{\pgfqpoint{2.861572in}{3.490507in}}%
\pgfpathlineto{\pgfqpoint{2.861572in}{3.490507in}}%
\pgfpathlineto{\pgfqpoint{2.861572in}{3.493456in}}%
\pgfpathlineto{\pgfqpoint{2.866113in}{3.493456in}}%
\pgfpathlineto{\pgfqpoint{2.866113in}{3.490507in}}%
\pgfpathmoveto{\pgfqpoint{2.861572in}{3.493456in}}%
\pgfpathlineto{\pgfqpoint{2.861572in}{3.493456in}}%
\pgfpathlineto{\pgfqpoint{2.861572in}{3.496405in}}%
\pgfpathlineto{\pgfqpoint{2.866113in}{3.496405in}}%
\pgfpathlineto{\pgfqpoint{2.866113in}{3.493456in}}%
\pgfpathmoveto{\pgfqpoint{2.866113in}{3.484609in}}%
\pgfpathlineto{\pgfqpoint{2.866113in}{3.484609in}}%
\pgfpathlineto{\pgfqpoint{2.866113in}{3.487558in}}%
\pgfpathlineto{\pgfqpoint{2.870655in}{3.487558in}}%
\pgfpathlineto{\pgfqpoint{2.870655in}{3.484609in}}%
\pgfpathmoveto{\pgfqpoint{2.866113in}{3.487558in}}%
\pgfpathlineto{\pgfqpoint{2.866113in}{3.487558in}}%
\pgfpathlineto{\pgfqpoint{2.866113in}{3.490507in}}%
\pgfpathlineto{\pgfqpoint{2.870655in}{3.490507in}}%
\pgfpathlineto{\pgfqpoint{2.870655in}{3.487558in}}%
\pgfpathmoveto{\pgfqpoint{2.875196in}{3.472812in}}%
\pgfpathlineto{\pgfqpoint{2.875196in}{3.472812in}}%
\pgfpathlineto{\pgfqpoint{2.875196in}{3.475762in}}%
\pgfpathlineto{\pgfqpoint{2.879737in}{3.475762in}}%
\pgfpathlineto{\pgfqpoint{2.879737in}{3.472812in}}%
\pgfpathmoveto{\pgfqpoint{2.875196in}{3.475762in}}%
\pgfpathlineto{\pgfqpoint{2.875196in}{3.475762in}}%
\pgfpathlineto{\pgfqpoint{2.875196in}{3.478711in}}%
\pgfpathlineto{\pgfqpoint{2.879737in}{3.478711in}}%
\pgfpathlineto{\pgfqpoint{2.879737in}{3.475762in}}%
\pgfpathmoveto{\pgfqpoint{2.857031in}{3.496405in}}%
\pgfpathlineto{\pgfqpoint{2.857031in}{3.496405in}}%
\pgfpathlineto{\pgfqpoint{2.857031in}{3.499354in}}%
\pgfpathlineto{\pgfqpoint{2.861572in}{3.499354in}}%
\pgfpathlineto{\pgfqpoint{2.861572in}{3.496405in}}%
\pgfpathmoveto{\pgfqpoint{2.857031in}{3.499354in}}%
\pgfpathlineto{\pgfqpoint{2.857031in}{3.499354in}}%
\pgfpathlineto{\pgfqpoint{2.857031in}{3.502304in}}%
\pgfpathlineto{\pgfqpoint{2.861572in}{3.502304in}}%
\pgfpathlineto{\pgfqpoint{2.861572in}{3.499354in}}%
\pgfpathmoveto{\pgfqpoint{2.929688in}{1.912677in}}%
\pgfpathlineto{\pgfqpoint{2.929688in}{1.912677in}}%
\pgfpathlineto{\pgfqpoint{2.929688in}{1.915627in}}%
\pgfpathlineto{\pgfqpoint{2.934229in}{1.915627in}}%
\pgfpathlineto{\pgfqpoint{2.934229in}{1.912677in}}%
\pgfpathmoveto{\pgfqpoint{2.929688in}{1.915627in}}%
\pgfpathlineto{\pgfqpoint{2.929688in}{1.915627in}}%
\pgfpathlineto{\pgfqpoint{2.929688in}{1.918576in}}%
\pgfpathlineto{\pgfqpoint{2.934229in}{1.918576in}}%
\pgfpathlineto{\pgfqpoint{2.934229in}{1.915627in}}%
\pgfpathmoveto{\pgfqpoint{2.929688in}{1.918576in}}%
\pgfpathlineto{\pgfqpoint{2.929688in}{1.918576in}}%
\pgfpathlineto{\pgfqpoint{2.929688in}{1.921525in}}%
\pgfpathlineto{\pgfqpoint{2.934229in}{1.921525in}}%
\pgfpathlineto{\pgfqpoint{2.934229in}{1.918576in}}%
\pgfpathmoveto{\pgfqpoint{2.934229in}{1.915627in}}%
\pgfpathlineto{\pgfqpoint{2.934229in}{1.915627in}}%
\pgfpathlineto{\pgfqpoint{2.934229in}{1.918576in}}%
\pgfpathlineto{\pgfqpoint{2.938770in}{1.918576in}}%
\pgfpathlineto{\pgfqpoint{2.938770in}{1.915627in}}%
\pgfpathmoveto{\pgfqpoint{2.934229in}{1.918576in}}%
\pgfpathlineto{\pgfqpoint{2.934229in}{1.918576in}}%
\pgfpathlineto{\pgfqpoint{2.934229in}{1.921525in}}%
\pgfpathlineto{\pgfqpoint{2.938770in}{1.921525in}}%
\pgfpathlineto{\pgfqpoint{2.938770in}{1.918576in}}%
\pgfpathmoveto{\pgfqpoint{2.938770in}{1.918576in}}%
\pgfpathlineto{\pgfqpoint{2.938770in}{1.918576in}}%
\pgfpathlineto{\pgfqpoint{2.938770in}{1.921525in}}%
\pgfpathlineto{\pgfqpoint{2.943311in}{1.921525in}}%
\pgfpathlineto{\pgfqpoint{2.943311in}{1.918576in}}%
\pgfpathmoveto{\pgfqpoint{2.938770in}{1.921525in}}%
\pgfpathlineto{\pgfqpoint{2.938770in}{1.921525in}}%
\pgfpathlineto{\pgfqpoint{2.938770in}{1.924474in}}%
\pgfpathlineto{\pgfqpoint{2.943311in}{1.924474in}}%
\pgfpathlineto{\pgfqpoint{2.943311in}{1.921525in}}%
\pgfpathmoveto{\pgfqpoint{2.938770in}{1.924474in}}%
\pgfpathlineto{\pgfqpoint{2.938770in}{1.924474in}}%
\pgfpathlineto{\pgfqpoint{2.938770in}{1.927423in}}%
\pgfpathlineto{\pgfqpoint{2.943311in}{1.927423in}}%
\pgfpathlineto{\pgfqpoint{2.943311in}{1.924474in}}%
\pgfpathmoveto{\pgfqpoint{2.943311in}{1.921525in}}%
\pgfpathlineto{\pgfqpoint{2.943311in}{1.921525in}}%
\pgfpathlineto{\pgfqpoint{2.943311in}{1.924474in}}%
\pgfpathlineto{\pgfqpoint{2.947852in}{1.924474in}}%
\pgfpathlineto{\pgfqpoint{2.947852in}{1.921525in}}%
\pgfpathmoveto{\pgfqpoint{2.943311in}{1.924474in}}%
\pgfpathlineto{\pgfqpoint{2.943311in}{1.924474in}}%
\pgfpathlineto{\pgfqpoint{2.943311in}{1.927423in}}%
\pgfpathlineto{\pgfqpoint{2.947852in}{1.927423in}}%
\pgfpathlineto{\pgfqpoint{2.947852in}{1.924474in}}%
\pgfpathmoveto{\pgfqpoint{2.947852in}{1.924474in}}%
\pgfpathlineto{\pgfqpoint{2.947852in}{1.924474in}}%
\pgfpathlineto{\pgfqpoint{2.947852in}{1.927423in}}%
\pgfpathlineto{\pgfqpoint{2.952394in}{1.927423in}}%
\pgfpathlineto{\pgfqpoint{2.952394in}{1.924474in}}%
\pgfpathmoveto{\pgfqpoint{2.947852in}{1.927423in}}%
\pgfpathlineto{\pgfqpoint{2.947852in}{1.927423in}}%
\pgfpathlineto{\pgfqpoint{2.947852in}{1.930373in}}%
\pgfpathlineto{\pgfqpoint{2.952394in}{1.930373in}}%
\pgfpathlineto{\pgfqpoint{2.952394in}{1.927423in}}%
\pgfpathmoveto{\pgfqpoint{2.947852in}{1.930373in}}%
\pgfpathlineto{\pgfqpoint{2.947852in}{1.930373in}}%
\pgfpathlineto{\pgfqpoint{2.947852in}{1.933322in}}%
\pgfpathlineto{\pgfqpoint{2.952394in}{1.933322in}}%
\pgfpathlineto{\pgfqpoint{2.952394in}{1.930373in}}%
\pgfpathmoveto{\pgfqpoint{2.952394in}{1.927423in}}%
\pgfpathlineto{\pgfqpoint{2.952394in}{1.927423in}}%
\pgfpathlineto{\pgfqpoint{2.952394in}{1.930373in}}%
\pgfpathlineto{\pgfqpoint{2.956935in}{1.930373in}}%
\pgfpathlineto{\pgfqpoint{2.956935in}{1.927423in}}%
\pgfpathmoveto{\pgfqpoint{2.952394in}{1.930373in}}%
\pgfpathlineto{\pgfqpoint{2.952394in}{1.930373in}}%
\pgfpathlineto{\pgfqpoint{2.952394in}{1.933322in}}%
\pgfpathlineto{\pgfqpoint{2.956935in}{1.933322in}}%
\pgfpathlineto{\pgfqpoint{2.956935in}{1.930373in}}%
\pgfpathmoveto{\pgfqpoint{2.956935in}{1.930373in}}%
\pgfpathlineto{\pgfqpoint{2.956935in}{1.930373in}}%
\pgfpathlineto{\pgfqpoint{2.956935in}{1.933322in}}%
\pgfpathlineto{\pgfqpoint{2.961476in}{1.933322in}}%
\pgfpathlineto{\pgfqpoint{2.961476in}{1.930373in}}%
\pgfpathmoveto{\pgfqpoint{2.956935in}{1.933322in}}%
\pgfpathlineto{\pgfqpoint{2.956935in}{1.933322in}}%
\pgfpathlineto{\pgfqpoint{2.956935in}{1.936271in}}%
\pgfpathlineto{\pgfqpoint{2.961476in}{1.936271in}}%
\pgfpathlineto{\pgfqpoint{2.961476in}{1.933322in}}%
\pgfpathmoveto{\pgfqpoint{2.956935in}{1.936271in}}%
\pgfpathlineto{\pgfqpoint{2.956935in}{1.936271in}}%
\pgfpathlineto{\pgfqpoint{2.956935in}{1.939220in}}%
\pgfpathlineto{\pgfqpoint{2.961476in}{1.939220in}}%
\pgfpathlineto{\pgfqpoint{2.961476in}{1.936271in}}%
\pgfpathmoveto{\pgfqpoint{2.961476in}{1.933322in}}%
\pgfpathlineto{\pgfqpoint{2.961476in}{1.933322in}}%
\pgfpathlineto{\pgfqpoint{2.961476in}{1.936271in}}%
\pgfpathlineto{\pgfqpoint{2.966017in}{1.936271in}}%
\pgfpathlineto{\pgfqpoint{2.966017in}{1.933322in}}%
\pgfpathmoveto{\pgfqpoint{2.961476in}{1.936271in}}%
\pgfpathlineto{\pgfqpoint{2.961476in}{1.936271in}}%
\pgfpathlineto{\pgfqpoint{2.961476in}{1.939220in}}%
\pgfpathlineto{\pgfqpoint{2.966017in}{1.939220in}}%
\pgfpathlineto{\pgfqpoint{2.966017in}{1.936271in}}%
\pgfpathmoveto{\pgfqpoint{2.966017in}{1.936271in}}%
\pgfpathlineto{\pgfqpoint{2.966017in}{1.936271in}}%
\pgfpathlineto{\pgfqpoint{2.966017in}{1.939220in}}%
\pgfpathlineto{\pgfqpoint{2.970558in}{1.939220in}}%
\pgfpathlineto{\pgfqpoint{2.970558in}{1.936271in}}%
\pgfpathmoveto{\pgfqpoint{2.966017in}{1.939220in}}%
\pgfpathlineto{\pgfqpoint{2.966017in}{1.939220in}}%
\pgfpathlineto{\pgfqpoint{2.966017in}{1.942169in}}%
\pgfpathlineto{\pgfqpoint{2.970558in}{1.942169in}}%
\pgfpathlineto{\pgfqpoint{2.970558in}{1.939220in}}%
\pgfpathmoveto{\pgfqpoint{2.966017in}{1.942169in}}%
\pgfpathlineto{\pgfqpoint{2.966017in}{1.942169in}}%
\pgfpathlineto{\pgfqpoint{2.966017in}{1.945118in}}%
\pgfpathlineto{\pgfqpoint{2.970558in}{1.945118in}}%
\pgfpathlineto{\pgfqpoint{2.970558in}{1.942169in}}%
\pgfpathmoveto{\pgfqpoint{2.970558in}{1.939220in}}%
\pgfpathlineto{\pgfqpoint{2.970558in}{1.939220in}}%
\pgfpathlineto{\pgfqpoint{2.970558in}{1.942169in}}%
\pgfpathlineto{\pgfqpoint{2.975099in}{1.942169in}}%
\pgfpathlineto{\pgfqpoint{2.975099in}{1.939220in}}%
\pgfpathmoveto{\pgfqpoint{2.970558in}{1.942169in}}%
\pgfpathlineto{\pgfqpoint{2.970558in}{1.942169in}}%
\pgfpathlineto{\pgfqpoint{2.970558in}{1.945118in}}%
\pgfpathlineto{\pgfqpoint{2.975099in}{1.945118in}}%
\pgfpathlineto{\pgfqpoint{2.975099in}{1.942169in}}%
\pgfpathmoveto{\pgfqpoint{2.975099in}{1.942169in}}%
\pgfpathlineto{\pgfqpoint{2.975099in}{1.942169in}}%
\pgfpathlineto{\pgfqpoint{2.975099in}{1.945118in}}%
\pgfpathlineto{\pgfqpoint{2.979640in}{1.945118in}}%
\pgfpathlineto{\pgfqpoint{2.979640in}{1.942169in}}%
\pgfpathmoveto{\pgfqpoint{2.975099in}{1.945118in}}%
\pgfpathlineto{\pgfqpoint{2.975099in}{1.945118in}}%
\pgfpathlineto{\pgfqpoint{2.975099in}{1.948068in}}%
\pgfpathlineto{\pgfqpoint{2.979640in}{1.948068in}}%
\pgfpathlineto{\pgfqpoint{2.979640in}{1.945118in}}%
\pgfpathmoveto{\pgfqpoint{2.975099in}{1.948068in}}%
\pgfpathlineto{\pgfqpoint{2.975099in}{1.948068in}}%
\pgfpathlineto{\pgfqpoint{2.975099in}{1.951017in}}%
\pgfpathlineto{\pgfqpoint{2.979640in}{1.951017in}}%
\pgfpathlineto{\pgfqpoint{2.979640in}{1.948068in}}%
\pgfpathmoveto{\pgfqpoint{2.979640in}{1.945118in}}%
\pgfpathlineto{\pgfqpoint{2.979640in}{1.945118in}}%
\pgfpathlineto{\pgfqpoint{2.979640in}{1.948068in}}%
\pgfpathlineto{\pgfqpoint{2.984181in}{1.948068in}}%
\pgfpathlineto{\pgfqpoint{2.984181in}{1.945118in}}%
\pgfpathmoveto{\pgfqpoint{2.979640in}{1.948068in}}%
\pgfpathlineto{\pgfqpoint{2.979640in}{1.948068in}}%
\pgfpathlineto{\pgfqpoint{2.979640in}{1.951017in}}%
\pgfpathlineto{\pgfqpoint{2.984181in}{1.951017in}}%
\pgfpathlineto{\pgfqpoint{2.984181in}{1.948068in}}%
\pgfpathmoveto{\pgfqpoint{2.984181in}{1.948068in}}%
\pgfpathlineto{\pgfqpoint{2.984181in}{1.948068in}}%
\pgfpathlineto{\pgfqpoint{2.984181in}{1.951017in}}%
\pgfpathlineto{\pgfqpoint{2.988722in}{1.951017in}}%
\pgfpathlineto{\pgfqpoint{2.988722in}{1.948068in}}%
\pgfpathmoveto{\pgfqpoint{2.984181in}{1.951017in}}%
\pgfpathlineto{\pgfqpoint{2.984181in}{1.951017in}}%
\pgfpathlineto{\pgfqpoint{2.984181in}{1.953966in}}%
\pgfpathlineto{\pgfqpoint{2.988722in}{1.953966in}}%
\pgfpathlineto{\pgfqpoint{2.988722in}{1.951017in}}%
\pgfpathmoveto{\pgfqpoint{2.984181in}{1.953966in}}%
\pgfpathlineto{\pgfqpoint{2.984181in}{1.953966in}}%
\pgfpathlineto{\pgfqpoint{2.984181in}{1.956915in}}%
\pgfpathlineto{\pgfqpoint{2.988722in}{1.956915in}}%
\pgfpathlineto{\pgfqpoint{2.988722in}{1.953966in}}%
\pgfpathmoveto{\pgfqpoint{2.988722in}{1.953966in}}%
\pgfpathlineto{\pgfqpoint{2.988722in}{1.953966in}}%
\pgfpathlineto{\pgfqpoint{2.988722in}{1.956915in}}%
\pgfpathlineto{\pgfqpoint{2.993264in}{1.956915in}}%
\pgfpathlineto{\pgfqpoint{2.993264in}{1.953966in}}%
\pgfpathmoveto{\pgfqpoint{2.984181in}{1.956915in}}%
\pgfpathlineto{\pgfqpoint{2.984181in}{1.956915in}}%
\pgfpathlineto{\pgfqpoint{2.984181in}{1.959864in}}%
\pgfpathlineto{\pgfqpoint{2.988722in}{1.959864in}}%
\pgfpathlineto{\pgfqpoint{2.988722in}{1.956915in}}%
\pgfpathmoveto{\pgfqpoint{2.984181in}{1.959864in}}%
\pgfpathlineto{\pgfqpoint{2.984181in}{1.959864in}}%
\pgfpathlineto{\pgfqpoint{2.984181in}{1.962813in}}%
\pgfpathlineto{\pgfqpoint{2.988722in}{1.962813in}}%
\pgfpathlineto{\pgfqpoint{2.988722in}{1.959864in}}%
\pgfpathmoveto{\pgfqpoint{2.988722in}{1.956915in}}%
\pgfpathlineto{\pgfqpoint{2.988722in}{1.956915in}}%
\pgfpathlineto{\pgfqpoint{2.988722in}{1.959864in}}%
\pgfpathlineto{\pgfqpoint{2.993264in}{1.959864in}}%
\pgfpathlineto{\pgfqpoint{2.993264in}{1.956915in}}%
\pgfpathmoveto{\pgfqpoint{2.988722in}{1.959864in}}%
\pgfpathlineto{\pgfqpoint{2.988722in}{1.959864in}}%
\pgfpathlineto{\pgfqpoint{2.988722in}{1.962813in}}%
\pgfpathlineto{\pgfqpoint{2.993264in}{1.962813in}}%
\pgfpathlineto{\pgfqpoint{2.993264in}{1.959864in}}%
\pgfpathmoveto{\pgfqpoint{2.993264in}{1.956915in}}%
\pgfpathlineto{\pgfqpoint{2.993264in}{1.956915in}}%
\pgfpathlineto{\pgfqpoint{2.993264in}{1.959864in}}%
\pgfpathlineto{\pgfqpoint{2.997805in}{1.959864in}}%
\pgfpathlineto{\pgfqpoint{2.997805in}{1.956915in}}%
\pgfpathmoveto{\pgfqpoint{2.993264in}{1.959864in}}%
\pgfpathlineto{\pgfqpoint{2.993264in}{1.959864in}}%
\pgfpathlineto{\pgfqpoint{2.993264in}{1.962813in}}%
\pgfpathlineto{\pgfqpoint{2.997805in}{1.962813in}}%
\pgfpathlineto{\pgfqpoint{2.997805in}{1.959864in}}%
\pgfpathmoveto{\pgfqpoint{2.997805in}{1.959864in}}%
\pgfpathlineto{\pgfqpoint{2.997805in}{1.959864in}}%
\pgfpathlineto{\pgfqpoint{2.997805in}{1.962813in}}%
\pgfpathlineto{\pgfqpoint{3.002346in}{1.962813in}}%
\pgfpathlineto{\pgfqpoint{3.002346in}{1.959864in}}%
\pgfpathmoveto{\pgfqpoint{2.993264in}{1.962813in}}%
\pgfpathlineto{\pgfqpoint{2.993264in}{1.962813in}}%
\pgfpathlineto{\pgfqpoint{2.993264in}{1.965763in}}%
\pgfpathlineto{\pgfqpoint{2.997805in}{1.965763in}}%
\pgfpathlineto{\pgfqpoint{2.997805in}{1.962813in}}%
\pgfpathmoveto{\pgfqpoint{2.993264in}{1.965763in}}%
\pgfpathlineto{\pgfqpoint{2.993264in}{1.965763in}}%
\pgfpathlineto{\pgfqpoint{2.993264in}{1.968712in}}%
\pgfpathlineto{\pgfqpoint{2.997805in}{1.968712in}}%
\pgfpathlineto{\pgfqpoint{2.997805in}{1.965763in}}%
\pgfpathmoveto{\pgfqpoint{2.997805in}{1.962813in}}%
\pgfpathlineto{\pgfqpoint{2.997805in}{1.962813in}}%
\pgfpathlineto{\pgfqpoint{2.997805in}{1.965763in}}%
\pgfpathlineto{\pgfqpoint{3.002346in}{1.965763in}}%
\pgfpathlineto{\pgfqpoint{3.002346in}{1.962813in}}%
\pgfpathmoveto{\pgfqpoint{2.997805in}{1.965763in}}%
\pgfpathlineto{\pgfqpoint{2.997805in}{1.965763in}}%
\pgfpathlineto{\pgfqpoint{2.997805in}{1.968712in}}%
\pgfpathlineto{\pgfqpoint{3.002346in}{1.968712in}}%
\pgfpathlineto{\pgfqpoint{3.002346in}{1.965763in}}%
\pgfpathmoveto{\pgfqpoint{3.002346in}{1.962813in}}%
\pgfpathlineto{\pgfqpoint{3.002346in}{1.962813in}}%
\pgfpathlineto{\pgfqpoint{3.002346in}{1.965763in}}%
\pgfpathlineto{\pgfqpoint{3.006887in}{1.965763in}}%
\pgfpathlineto{\pgfqpoint{3.006887in}{1.962813in}}%
\pgfpathmoveto{\pgfqpoint{3.002346in}{1.965763in}}%
\pgfpathlineto{\pgfqpoint{3.002346in}{1.965763in}}%
\pgfpathlineto{\pgfqpoint{3.002346in}{1.968712in}}%
\pgfpathlineto{\pgfqpoint{3.006887in}{1.968712in}}%
\pgfpathlineto{\pgfqpoint{3.006887in}{1.965763in}}%
\pgfpathmoveto{\pgfqpoint{3.006887in}{1.965763in}}%
\pgfpathlineto{\pgfqpoint{3.006887in}{1.965763in}}%
\pgfpathlineto{\pgfqpoint{3.006887in}{1.968712in}}%
\pgfpathlineto{\pgfqpoint{3.011428in}{1.968712in}}%
\pgfpathlineto{\pgfqpoint{3.011428in}{1.965763in}}%
\pgfpathmoveto{\pgfqpoint{3.002346in}{1.968712in}}%
\pgfpathlineto{\pgfqpoint{3.002346in}{1.968712in}}%
\pgfpathlineto{\pgfqpoint{3.002346in}{1.971661in}}%
\pgfpathlineto{\pgfqpoint{3.006887in}{1.971661in}}%
\pgfpathlineto{\pgfqpoint{3.006887in}{1.968712in}}%
\pgfpathmoveto{\pgfqpoint{3.002346in}{1.971661in}}%
\pgfpathlineto{\pgfqpoint{3.002346in}{1.971661in}}%
\pgfpathlineto{\pgfqpoint{3.002346in}{1.974610in}}%
\pgfpathlineto{\pgfqpoint{3.006887in}{1.974610in}}%
\pgfpathlineto{\pgfqpoint{3.006887in}{1.971661in}}%
\pgfpathmoveto{\pgfqpoint{3.006887in}{1.968712in}}%
\pgfpathlineto{\pgfqpoint{3.006887in}{1.968712in}}%
\pgfpathlineto{\pgfqpoint{3.006887in}{1.971661in}}%
\pgfpathlineto{\pgfqpoint{3.011428in}{1.971661in}}%
\pgfpathlineto{\pgfqpoint{3.011428in}{1.968712in}}%
\pgfpathmoveto{\pgfqpoint{3.006887in}{1.971661in}}%
\pgfpathlineto{\pgfqpoint{3.006887in}{1.971661in}}%
\pgfpathlineto{\pgfqpoint{3.006887in}{1.974610in}}%
\pgfpathlineto{\pgfqpoint{3.011428in}{1.974610in}}%
\pgfpathlineto{\pgfqpoint{3.011428in}{1.971661in}}%
\pgfpathmoveto{\pgfqpoint{3.011428in}{1.968712in}}%
\pgfpathlineto{\pgfqpoint{3.011428in}{1.968712in}}%
\pgfpathlineto{\pgfqpoint{3.011428in}{1.971661in}}%
\pgfpathlineto{\pgfqpoint{3.015969in}{1.971661in}}%
\pgfpathlineto{\pgfqpoint{3.015969in}{1.968712in}}%
\pgfpathmoveto{\pgfqpoint{3.011428in}{1.971661in}}%
\pgfpathlineto{\pgfqpoint{3.011428in}{1.971661in}}%
\pgfpathlineto{\pgfqpoint{3.011428in}{1.974610in}}%
\pgfpathlineto{\pgfqpoint{3.015969in}{1.974610in}}%
\pgfpathlineto{\pgfqpoint{3.015969in}{1.971661in}}%
\pgfpathmoveto{\pgfqpoint{3.015969in}{1.971661in}}%
\pgfpathlineto{\pgfqpoint{3.015969in}{1.971661in}}%
\pgfpathlineto{\pgfqpoint{3.015969in}{1.974610in}}%
\pgfpathlineto{\pgfqpoint{3.020510in}{1.974610in}}%
\pgfpathlineto{\pgfqpoint{3.020510in}{1.971661in}}%
\pgfpathmoveto{\pgfqpoint{3.011428in}{1.974610in}}%
\pgfpathlineto{\pgfqpoint{3.011428in}{1.974610in}}%
\pgfpathlineto{\pgfqpoint{3.011428in}{1.977559in}}%
\pgfpathlineto{\pgfqpoint{3.015969in}{1.977559in}}%
\pgfpathlineto{\pgfqpoint{3.015969in}{1.974610in}}%
\pgfpathmoveto{\pgfqpoint{3.011428in}{1.977559in}}%
\pgfpathlineto{\pgfqpoint{3.011428in}{1.977559in}}%
\pgfpathlineto{\pgfqpoint{3.011428in}{1.980508in}}%
\pgfpathlineto{\pgfqpoint{3.015969in}{1.980508in}}%
\pgfpathlineto{\pgfqpoint{3.015969in}{1.977559in}}%
\pgfpathmoveto{\pgfqpoint{3.015969in}{1.974610in}}%
\pgfpathlineto{\pgfqpoint{3.015969in}{1.974610in}}%
\pgfpathlineto{\pgfqpoint{3.015969in}{1.977559in}}%
\pgfpathlineto{\pgfqpoint{3.020510in}{1.977559in}}%
\pgfpathlineto{\pgfqpoint{3.020510in}{1.974610in}}%
\pgfpathmoveto{\pgfqpoint{3.015969in}{1.977559in}}%
\pgfpathlineto{\pgfqpoint{3.015969in}{1.977559in}}%
\pgfpathlineto{\pgfqpoint{3.015969in}{1.980508in}}%
\pgfpathlineto{\pgfqpoint{3.020510in}{1.980508in}}%
\pgfpathlineto{\pgfqpoint{3.020510in}{1.977559in}}%
\pgfpathmoveto{\pgfqpoint{3.020510in}{1.974610in}}%
\pgfpathlineto{\pgfqpoint{3.020510in}{1.974610in}}%
\pgfpathlineto{\pgfqpoint{3.020510in}{1.977559in}}%
\pgfpathlineto{\pgfqpoint{3.025051in}{1.977559in}}%
\pgfpathlineto{\pgfqpoint{3.025051in}{1.974610in}}%
\pgfpathmoveto{\pgfqpoint{3.020510in}{1.977559in}}%
\pgfpathlineto{\pgfqpoint{3.020510in}{1.977559in}}%
\pgfpathlineto{\pgfqpoint{3.020510in}{1.980508in}}%
\pgfpathlineto{\pgfqpoint{3.025051in}{1.980508in}}%
\pgfpathlineto{\pgfqpoint{3.025051in}{1.977559in}}%
\pgfpathmoveto{\pgfqpoint{3.025051in}{1.977559in}}%
\pgfpathlineto{\pgfqpoint{3.025051in}{1.977559in}}%
\pgfpathlineto{\pgfqpoint{3.025051in}{1.980508in}}%
\pgfpathlineto{\pgfqpoint{3.029593in}{1.980508in}}%
\pgfpathlineto{\pgfqpoint{3.029593in}{1.977559in}}%
\pgfpathmoveto{\pgfqpoint{3.020510in}{1.980508in}}%
\pgfpathlineto{\pgfqpoint{3.020510in}{1.980508in}}%
\pgfpathlineto{\pgfqpoint{3.020510in}{1.983458in}}%
\pgfpathlineto{\pgfqpoint{3.025051in}{1.983458in}}%
\pgfpathlineto{\pgfqpoint{3.025051in}{1.980508in}}%
\pgfpathmoveto{\pgfqpoint{3.020510in}{1.983458in}}%
\pgfpathlineto{\pgfqpoint{3.020510in}{1.983458in}}%
\pgfpathlineto{\pgfqpoint{3.020510in}{1.986407in}}%
\pgfpathlineto{\pgfqpoint{3.025051in}{1.986407in}}%
\pgfpathlineto{\pgfqpoint{3.025051in}{1.983458in}}%
\pgfpathmoveto{\pgfqpoint{3.025051in}{1.980508in}}%
\pgfpathlineto{\pgfqpoint{3.025051in}{1.980508in}}%
\pgfpathlineto{\pgfqpoint{3.025051in}{1.983458in}}%
\pgfpathlineto{\pgfqpoint{3.029593in}{1.983458in}}%
\pgfpathlineto{\pgfqpoint{3.029593in}{1.980508in}}%
\pgfpathmoveto{\pgfqpoint{3.025051in}{1.983458in}}%
\pgfpathlineto{\pgfqpoint{3.025051in}{1.983458in}}%
\pgfpathlineto{\pgfqpoint{3.025051in}{1.986407in}}%
\pgfpathlineto{\pgfqpoint{3.029593in}{1.986407in}}%
\pgfpathlineto{\pgfqpoint{3.029593in}{1.983458in}}%
\pgfpathmoveto{\pgfqpoint{3.029593in}{1.980508in}}%
\pgfpathlineto{\pgfqpoint{3.029593in}{1.980508in}}%
\pgfpathlineto{\pgfqpoint{3.029593in}{1.983458in}}%
\pgfpathlineto{\pgfqpoint{3.034134in}{1.983458in}}%
\pgfpathlineto{\pgfqpoint{3.034134in}{1.980508in}}%
\pgfpathmoveto{\pgfqpoint{3.029593in}{1.983458in}}%
\pgfpathlineto{\pgfqpoint{3.029593in}{1.983458in}}%
\pgfpathlineto{\pgfqpoint{3.029593in}{1.986407in}}%
\pgfpathlineto{\pgfqpoint{3.034134in}{1.986407in}}%
\pgfpathlineto{\pgfqpoint{3.034134in}{1.983458in}}%
\pgfpathmoveto{\pgfqpoint{3.034134in}{1.983458in}}%
\pgfpathlineto{\pgfqpoint{3.034134in}{1.983458in}}%
\pgfpathlineto{\pgfqpoint{3.034134in}{1.986407in}}%
\pgfpathlineto{\pgfqpoint{3.038675in}{1.986407in}}%
\pgfpathlineto{\pgfqpoint{3.038675in}{1.983458in}}%
\pgfpathmoveto{\pgfqpoint{3.029593in}{1.986407in}}%
\pgfpathlineto{\pgfqpoint{3.029593in}{1.986407in}}%
\pgfpathlineto{\pgfqpoint{3.029593in}{1.989356in}}%
\pgfpathlineto{\pgfqpoint{3.034134in}{1.989356in}}%
\pgfpathlineto{\pgfqpoint{3.034134in}{1.986407in}}%
\pgfpathmoveto{\pgfqpoint{3.029593in}{1.989356in}}%
\pgfpathlineto{\pgfqpoint{3.029593in}{1.989356in}}%
\pgfpathlineto{\pgfqpoint{3.029593in}{1.992305in}}%
\pgfpathlineto{\pgfqpoint{3.034134in}{1.992305in}}%
\pgfpathlineto{\pgfqpoint{3.034134in}{1.989356in}}%
\pgfpathmoveto{\pgfqpoint{3.034134in}{1.986407in}}%
\pgfpathlineto{\pgfqpoint{3.034134in}{1.986407in}}%
\pgfpathlineto{\pgfqpoint{3.034134in}{1.989356in}}%
\pgfpathlineto{\pgfqpoint{3.038675in}{1.989356in}}%
\pgfpathlineto{\pgfqpoint{3.038675in}{1.986407in}}%
\pgfpathmoveto{\pgfqpoint{3.034134in}{1.989356in}}%
\pgfpathlineto{\pgfqpoint{3.034134in}{1.989356in}}%
\pgfpathlineto{\pgfqpoint{3.034134in}{1.992305in}}%
\pgfpathlineto{\pgfqpoint{3.038675in}{1.992305in}}%
\pgfpathlineto{\pgfqpoint{3.038675in}{1.989356in}}%
\pgfpathmoveto{\pgfqpoint{3.038675in}{1.986407in}}%
\pgfpathlineto{\pgfqpoint{3.038675in}{1.986407in}}%
\pgfpathlineto{\pgfqpoint{3.038675in}{1.989356in}}%
\pgfpathlineto{\pgfqpoint{3.043216in}{1.989356in}}%
\pgfpathlineto{\pgfqpoint{3.043216in}{1.986407in}}%
\pgfpathmoveto{\pgfqpoint{3.038675in}{1.989356in}}%
\pgfpathlineto{\pgfqpoint{3.038675in}{1.989356in}}%
\pgfpathlineto{\pgfqpoint{3.038675in}{1.992305in}}%
\pgfpathlineto{\pgfqpoint{3.043216in}{1.992305in}}%
\pgfpathlineto{\pgfqpoint{3.043216in}{1.989356in}}%
\pgfpathmoveto{\pgfqpoint{3.043216in}{1.989356in}}%
\pgfpathlineto{\pgfqpoint{3.043216in}{1.989356in}}%
\pgfpathlineto{\pgfqpoint{3.043216in}{1.992305in}}%
\pgfpathlineto{\pgfqpoint{3.047757in}{1.992305in}}%
\pgfpathlineto{\pgfqpoint{3.047757in}{1.989356in}}%
\pgfpathmoveto{\pgfqpoint{3.038675in}{1.992305in}}%
\pgfpathlineto{\pgfqpoint{3.038675in}{1.992305in}}%
\pgfpathlineto{\pgfqpoint{3.038675in}{1.995254in}}%
\pgfpathlineto{\pgfqpoint{3.043216in}{1.995254in}}%
\pgfpathlineto{\pgfqpoint{3.043216in}{1.992305in}}%
\pgfpathmoveto{\pgfqpoint{3.038675in}{1.995254in}}%
\pgfpathlineto{\pgfqpoint{3.038675in}{1.995254in}}%
\pgfpathlineto{\pgfqpoint{3.038675in}{1.998204in}}%
\pgfpathlineto{\pgfqpoint{3.043216in}{1.998204in}}%
\pgfpathlineto{\pgfqpoint{3.043216in}{1.995254in}}%
\pgfpathmoveto{\pgfqpoint{3.043216in}{1.992305in}}%
\pgfpathlineto{\pgfqpoint{3.043216in}{1.992305in}}%
\pgfpathlineto{\pgfqpoint{3.043216in}{1.995254in}}%
\pgfpathlineto{\pgfqpoint{3.047757in}{1.995254in}}%
\pgfpathlineto{\pgfqpoint{3.047757in}{1.992305in}}%
\pgfpathmoveto{\pgfqpoint{3.043216in}{1.995254in}}%
\pgfpathlineto{\pgfqpoint{3.043216in}{1.995254in}}%
\pgfpathlineto{\pgfqpoint{3.043216in}{1.998204in}}%
\pgfpathlineto{\pgfqpoint{3.047757in}{1.998204in}}%
\pgfpathlineto{\pgfqpoint{3.047757in}{1.995254in}}%
\pgfpathmoveto{\pgfqpoint{3.047757in}{1.992305in}}%
\pgfpathlineto{\pgfqpoint{3.047757in}{1.992305in}}%
\pgfpathlineto{\pgfqpoint{3.047757in}{1.995254in}}%
\pgfpathlineto{\pgfqpoint{3.052298in}{1.995254in}}%
\pgfpathlineto{\pgfqpoint{3.052298in}{1.992305in}}%
\pgfpathmoveto{\pgfqpoint{3.047757in}{1.995254in}}%
\pgfpathlineto{\pgfqpoint{3.047757in}{1.995254in}}%
\pgfpathlineto{\pgfqpoint{3.047757in}{1.998204in}}%
\pgfpathlineto{\pgfqpoint{3.052298in}{1.998204in}}%
\pgfpathlineto{\pgfqpoint{3.052298in}{1.995254in}}%
\pgfpathmoveto{\pgfqpoint{3.052298in}{1.995254in}}%
\pgfpathlineto{\pgfqpoint{3.052298in}{1.995254in}}%
\pgfpathlineto{\pgfqpoint{3.052298in}{1.998204in}}%
\pgfpathlineto{\pgfqpoint{3.056839in}{1.998204in}}%
\pgfpathlineto{\pgfqpoint{3.056839in}{1.995254in}}%
\pgfpathmoveto{\pgfqpoint{3.047757in}{1.998204in}}%
\pgfpathlineto{\pgfqpoint{3.047757in}{1.998204in}}%
\pgfpathlineto{\pgfqpoint{3.047757in}{2.001153in}}%
\pgfpathlineto{\pgfqpoint{3.052298in}{2.001153in}}%
\pgfpathlineto{\pgfqpoint{3.052298in}{1.998204in}}%
\pgfpathmoveto{\pgfqpoint{3.047757in}{2.001153in}}%
\pgfpathlineto{\pgfqpoint{3.047757in}{2.001153in}}%
\pgfpathlineto{\pgfqpoint{3.047757in}{2.004102in}}%
\pgfpathlineto{\pgfqpoint{3.052298in}{2.004102in}}%
\pgfpathlineto{\pgfqpoint{3.052298in}{2.001153in}}%
\pgfpathmoveto{\pgfqpoint{3.052298in}{1.998204in}}%
\pgfpathlineto{\pgfqpoint{3.052298in}{1.998204in}}%
\pgfpathlineto{\pgfqpoint{3.052298in}{2.001153in}}%
\pgfpathlineto{\pgfqpoint{3.056839in}{2.001153in}}%
\pgfpathlineto{\pgfqpoint{3.056839in}{1.998204in}}%
\pgfpathmoveto{\pgfqpoint{3.052298in}{2.001153in}}%
\pgfpathlineto{\pgfqpoint{3.052298in}{2.001153in}}%
\pgfpathlineto{\pgfqpoint{3.052298in}{2.004102in}}%
\pgfpathlineto{\pgfqpoint{3.056839in}{2.004102in}}%
\pgfpathlineto{\pgfqpoint{3.056839in}{2.001153in}}%
\pgfpathmoveto{\pgfqpoint{3.056839in}{1.998204in}}%
\pgfpathlineto{\pgfqpoint{3.056839in}{1.998204in}}%
\pgfpathlineto{\pgfqpoint{3.056839in}{2.001153in}}%
\pgfpathlineto{\pgfqpoint{3.061380in}{2.001153in}}%
\pgfpathlineto{\pgfqpoint{3.061380in}{1.998204in}}%
\pgfpathmoveto{\pgfqpoint{3.056839in}{2.001153in}}%
\pgfpathlineto{\pgfqpoint{3.056839in}{2.001153in}}%
\pgfpathlineto{\pgfqpoint{3.056839in}{2.004102in}}%
\pgfpathlineto{\pgfqpoint{3.061380in}{2.004102in}}%
\pgfpathlineto{\pgfqpoint{3.061380in}{2.001153in}}%
\pgfpathmoveto{\pgfqpoint{3.061380in}{2.001153in}}%
\pgfpathlineto{\pgfqpoint{3.061380in}{2.001153in}}%
\pgfpathlineto{\pgfqpoint{3.061380in}{2.004102in}}%
\pgfpathlineto{\pgfqpoint{3.065922in}{2.004102in}}%
\pgfpathlineto{\pgfqpoint{3.065922in}{2.001153in}}%
\pgfpathmoveto{\pgfqpoint{3.056839in}{2.004102in}}%
\pgfpathlineto{\pgfqpoint{3.056839in}{2.004102in}}%
\pgfpathlineto{\pgfqpoint{3.056839in}{2.007051in}}%
\pgfpathlineto{\pgfqpoint{3.061380in}{2.007051in}}%
\pgfpathlineto{\pgfqpoint{3.061380in}{2.004102in}}%
\pgfpathmoveto{\pgfqpoint{3.056839in}{2.007051in}}%
\pgfpathlineto{\pgfqpoint{3.056839in}{2.007051in}}%
\pgfpathlineto{\pgfqpoint{3.056839in}{2.010000in}}%
\pgfpathlineto{\pgfqpoint{3.061380in}{2.010000in}}%
\pgfpathlineto{\pgfqpoint{3.061380in}{2.007051in}}%
\pgfpathmoveto{\pgfqpoint{3.061380in}{2.004102in}}%
\pgfpathlineto{\pgfqpoint{3.061380in}{2.004102in}}%
\pgfpathlineto{\pgfqpoint{3.061380in}{2.007051in}}%
\pgfpathlineto{\pgfqpoint{3.065922in}{2.007051in}}%
\pgfpathlineto{\pgfqpoint{3.065922in}{2.004102in}}%
\pgfpathmoveto{\pgfqpoint{3.061380in}{2.007051in}}%
\pgfpathlineto{\pgfqpoint{3.061380in}{2.007051in}}%
\pgfpathlineto{\pgfqpoint{3.061380in}{2.010000in}}%
\pgfpathlineto{\pgfqpoint{3.065922in}{2.010000in}}%
\pgfpathlineto{\pgfqpoint{3.065922in}{2.007051in}}%
\pgfpathmoveto{\pgfqpoint{3.065922in}{2.004102in}}%
\pgfpathlineto{\pgfqpoint{3.065922in}{2.004102in}}%
\pgfpathlineto{\pgfqpoint{3.065922in}{2.007051in}}%
\pgfpathlineto{\pgfqpoint{3.070463in}{2.007051in}}%
\pgfpathlineto{\pgfqpoint{3.070463in}{2.004102in}}%
\pgfpathmoveto{\pgfqpoint{3.065922in}{2.007051in}}%
\pgfpathlineto{\pgfqpoint{3.065922in}{2.007051in}}%
\pgfpathlineto{\pgfqpoint{3.065922in}{2.010000in}}%
\pgfpathlineto{\pgfqpoint{3.070463in}{2.010000in}}%
\pgfpathlineto{\pgfqpoint{3.070463in}{2.007051in}}%
\pgfpathmoveto{\pgfqpoint{3.070463in}{2.007051in}}%
\pgfpathlineto{\pgfqpoint{3.070463in}{2.007051in}}%
\pgfpathlineto{\pgfqpoint{3.070463in}{2.010000in}}%
\pgfpathlineto{\pgfqpoint{3.075004in}{2.010000in}}%
\pgfpathlineto{\pgfqpoint{3.075004in}{2.007051in}}%
\pgfpathmoveto{\pgfqpoint{3.065922in}{2.010000in}}%
\pgfpathlineto{\pgfqpoint{3.065922in}{2.010000in}}%
\pgfpathlineto{\pgfqpoint{3.065922in}{2.012949in}}%
\pgfpathlineto{\pgfqpoint{3.070463in}{2.012949in}}%
\pgfpathlineto{\pgfqpoint{3.070463in}{2.010000in}}%
\pgfpathmoveto{\pgfqpoint{3.065922in}{2.012949in}}%
\pgfpathlineto{\pgfqpoint{3.065922in}{2.012949in}}%
\pgfpathlineto{\pgfqpoint{3.065922in}{2.015899in}}%
\pgfpathlineto{\pgfqpoint{3.070463in}{2.015899in}}%
\pgfpathlineto{\pgfqpoint{3.070463in}{2.012949in}}%
\pgfpathmoveto{\pgfqpoint{3.070463in}{2.010000in}}%
\pgfpathlineto{\pgfqpoint{3.070463in}{2.010000in}}%
\pgfpathlineto{\pgfqpoint{3.070463in}{2.012949in}}%
\pgfpathlineto{\pgfqpoint{3.075004in}{2.012949in}}%
\pgfpathlineto{\pgfqpoint{3.075004in}{2.010000in}}%
\pgfpathmoveto{\pgfqpoint{3.070463in}{2.012949in}}%
\pgfpathlineto{\pgfqpoint{3.070463in}{2.012949in}}%
\pgfpathlineto{\pgfqpoint{3.070463in}{2.015899in}}%
\pgfpathlineto{\pgfqpoint{3.075004in}{2.015899in}}%
\pgfpathlineto{\pgfqpoint{3.075004in}{2.012949in}}%
\pgfpathmoveto{\pgfqpoint{3.065922in}{3.213281in}}%
\pgfpathlineto{\pgfqpoint{3.065922in}{3.213281in}}%
\pgfpathlineto{\pgfqpoint{3.065922in}{3.216230in}}%
\pgfpathlineto{\pgfqpoint{3.070463in}{3.216230in}}%
\pgfpathlineto{\pgfqpoint{3.070463in}{3.213281in}}%
\pgfpathmoveto{\pgfqpoint{3.065922in}{3.216230in}}%
\pgfpathlineto{\pgfqpoint{3.065922in}{3.216230in}}%
\pgfpathlineto{\pgfqpoint{3.065922in}{3.219180in}}%
\pgfpathlineto{\pgfqpoint{3.070463in}{3.219180in}}%
\pgfpathlineto{\pgfqpoint{3.070463in}{3.216230in}}%
\pgfpathmoveto{\pgfqpoint{3.070463in}{3.213281in}}%
\pgfpathlineto{\pgfqpoint{3.070463in}{3.213281in}}%
\pgfpathlineto{\pgfqpoint{3.070463in}{3.216230in}}%
\pgfpathlineto{\pgfqpoint{3.075004in}{3.216230in}}%
\pgfpathlineto{\pgfqpoint{3.075004in}{3.213281in}}%
\pgfpathmoveto{\pgfqpoint{3.070463in}{3.216230in}}%
\pgfpathlineto{\pgfqpoint{3.070463in}{3.216230in}}%
\pgfpathlineto{\pgfqpoint{3.070463in}{3.219180in}}%
\pgfpathlineto{\pgfqpoint{3.075004in}{3.219180in}}%
\pgfpathlineto{\pgfqpoint{3.075004in}{3.216230in}}%
\pgfpathmoveto{\pgfqpoint{3.065922in}{3.219180in}}%
\pgfpathlineto{\pgfqpoint{3.065922in}{3.219180in}}%
\pgfpathlineto{\pgfqpoint{3.065922in}{3.222129in}}%
\pgfpathlineto{\pgfqpoint{3.070463in}{3.222129in}}%
\pgfpathlineto{\pgfqpoint{3.070463in}{3.219180in}}%
\pgfpathmoveto{\pgfqpoint{3.065922in}{3.222129in}}%
\pgfpathlineto{\pgfqpoint{3.065922in}{3.222129in}}%
\pgfpathlineto{\pgfqpoint{3.065922in}{3.225078in}}%
\pgfpathlineto{\pgfqpoint{3.070463in}{3.225078in}}%
\pgfpathlineto{\pgfqpoint{3.070463in}{3.222129in}}%
\pgfpathmoveto{\pgfqpoint{3.070463in}{3.219180in}}%
\pgfpathlineto{\pgfqpoint{3.070463in}{3.219180in}}%
\pgfpathlineto{\pgfqpoint{3.070463in}{3.222129in}}%
\pgfpathlineto{\pgfqpoint{3.075004in}{3.222129in}}%
\pgfpathlineto{\pgfqpoint{3.075004in}{3.219180in}}%
\pgfpathmoveto{\pgfqpoint{3.070463in}{3.222129in}}%
\pgfpathlineto{\pgfqpoint{3.070463in}{3.222129in}}%
\pgfpathlineto{\pgfqpoint{3.070463in}{3.225078in}}%
\pgfpathlineto{\pgfqpoint{3.075004in}{3.225078in}}%
\pgfpathlineto{\pgfqpoint{3.075004in}{3.222129in}}%
\pgfpathmoveto{\pgfqpoint{3.056839in}{3.225078in}}%
\pgfpathlineto{\pgfqpoint{3.056839in}{3.225078in}}%
\pgfpathlineto{\pgfqpoint{3.056839in}{3.228027in}}%
\pgfpathlineto{\pgfqpoint{3.061380in}{3.228027in}}%
\pgfpathlineto{\pgfqpoint{3.061380in}{3.225078in}}%
\pgfpathmoveto{\pgfqpoint{3.056839in}{3.228027in}}%
\pgfpathlineto{\pgfqpoint{3.056839in}{3.228027in}}%
\pgfpathlineto{\pgfqpoint{3.056839in}{3.230976in}}%
\pgfpathlineto{\pgfqpoint{3.061380in}{3.230976in}}%
\pgfpathlineto{\pgfqpoint{3.061380in}{3.228027in}}%
\pgfpathmoveto{\pgfqpoint{3.061380in}{3.225078in}}%
\pgfpathlineto{\pgfqpoint{3.061380in}{3.225078in}}%
\pgfpathlineto{\pgfqpoint{3.061380in}{3.228027in}}%
\pgfpathlineto{\pgfqpoint{3.065922in}{3.228027in}}%
\pgfpathlineto{\pgfqpoint{3.065922in}{3.225078in}}%
\pgfpathmoveto{\pgfqpoint{3.061380in}{3.228027in}}%
\pgfpathlineto{\pgfqpoint{3.061380in}{3.228027in}}%
\pgfpathlineto{\pgfqpoint{3.061380in}{3.230976in}}%
\pgfpathlineto{\pgfqpoint{3.065922in}{3.230976in}}%
\pgfpathlineto{\pgfqpoint{3.065922in}{3.228027in}}%
\pgfpathmoveto{\pgfqpoint{3.056839in}{3.230976in}}%
\pgfpathlineto{\pgfqpoint{3.056839in}{3.230976in}}%
\pgfpathlineto{\pgfqpoint{3.056839in}{3.233925in}}%
\pgfpathlineto{\pgfqpoint{3.061380in}{3.233925in}}%
\pgfpathlineto{\pgfqpoint{3.061380in}{3.230976in}}%
\pgfpathmoveto{\pgfqpoint{3.056839in}{3.233925in}}%
\pgfpathlineto{\pgfqpoint{3.056839in}{3.233925in}}%
\pgfpathlineto{\pgfqpoint{3.056839in}{3.236874in}}%
\pgfpathlineto{\pgfqpoint{3.061380in}{3.236874in}}%
\pgfpathlineto{\pgfqpoint{3.061380in}{3.233925in}}%
\pgfpathmoveto{\pgfqpoint{3.061380in}{3.230976in}}%
\pgfpathlineto{\pgfqpoint{3.061380in}{3.230976in}}%
\pgfpathlineto{\pgfqpoint{3.061380in}{3.233925in}}%
\pgfpathlineto{\pgfqpoint{3.065922in}{3.233925in}}%
\pgfpathlineto{\pgfqpoint{3.065922in}{3.230976in}}%
\pgfpathmoveto{\pgfqpoint{3.061380in}{3.233925in}}%
\pgfpathlineto{\pgfqpoint{3.061380in}{3.233925in}}%
\pgfpathlineto{\pgfqpoint{3.061380in}{3.236874in}}%
\pgfpathlineto{\pgfqpoint{3.065922in}{3.236874in}}%
\pgfpathlineto{\pgfqpoint{3.065922in}{3.233925in}}%
\pgfpathmoveto{\pgfqpoint{3.065922in}{3.225078in}}%
\pgfpathlineto{\pgfqpoint{3.065922in}{3.225078in}}%
\pgfpathlineto{\pgfqpoint{3.065922in}{3.228027in}}%
\pgfpathlineto{\pgfqpoint{3.070463in}{3.228027in}}%
\pgfpathlineto{\pgfqpoint{3.070463in}{3.225078in}}%
\pgfpathmoveto{\pgfqpoint{3.065922in}{3.228027in}}%
\pgfpathlineto{\pgfqpoint{3.065922in}{3.228027in}}%
\pgfpathlineto{\pgfqpoint{3.065922in}{3.230976in}}%
\pgfpathlineto{\pgfqpoint{3.070463in}{3.230976in}}%
\pgfpathlineto{\pgfqpoint{3.070463in}{3.228027in}}%
\pgfpathmoveto{\pgfqpoint{2.993264in}{3.307657in}}%
\pgfpathlineto{\pgfqpoint{2.993264in}{3.307657in}}%
\pgfpathlineto{\pgfqpoint{2.993264in}{3.310607in}}%
\pgfpathlineto{\pgfqpoint{2.997805in}{3.310607in}}%
\pgfpathlineto{\pgfqpoint{2.997805in}{3.307657in}}%
\pgfpathmoveto{\pgfqpoint{2.993264in}{3.310607in}}%
\pgfpathlineto{\pgfqpoint{2.993264in}{3.310607in}}%
\pgfpathlineto{\pgfqpoint{2.993264in}{3.313556in}}%
\pgfpathlineto{\pgfqpoint{2.997805in}{3.313556in}}%
\pgfpathlineto{\pgfqpoint{2.997805in}{3.310607in}}%
\pgfpathmoveto{\pgfqpoint{2.997805in}{3.307657in}}%
\pgfpathlineto{\pgfqpoint{2.997805in}{3.307657in}}%
\pgfpathlineto{\pgfqpoint{2.997805in}{3.310607in}}%
\pgfpathlineto{\pgfqpoint{3.002346in}{3.310607in}}%
\pgfpathlineto{\pgfqpoint{3.002346in}{3.307657in}}%
\pgfpathmoveto{\pgfqpoint{2.997805in}{3.310607in}}%
\pgfpathlineto{\pgfqpoint{2.997805in}{3.310607in}}%
\pgfpathlineto{\pgfqpoint{2.997805in}{3.313556in}}%
\pgfpathlineto{\pgfqpoint{3.002346in}{3.313556in}}%
\pgfpathlineto{\pgfqpoint{3.002346in}{3.310607in}}%
\pgfpathmoveto{\pgfqpoint{2.993264in}{3.313556in}}%
\pgfpathlineto{\pgfqpoint{2.993264in}{3.313556in}}%
\pgfpathlineto{\pgfqpoint{2.993264in}{3.316505in}}%
\pgfpathlineto{\pgfqpoint{2.997805in}{3.316505in}}%
\pgfpathlineto{\pgfqpoint{2.997805in}{3.313556in}}%
\pgfpathmoveto{\pgfqpoint{2.993264in}{3.316505in}}%
\pgfpathlineto{\pgfqpoint{2.993264in}{3.316505in}}%
\pgfpathlineto{\pgfqpoint{2.993264in}{3.319455in}}%
\pgfpathlineto{\pgfqpoint{2.997805in}{3.319455in}}%
\pgfpathlineto{\pgfqpoint{2.997805in}{3.316505in}}%
\pgfpathmoveto{\pgfqpoint{2.997805in}{3.313556in}}%
\pgfpathlineto{\pgfqpoint{2.997805in}{3.313556in}}%
\pgfpathlineto{\pgfqpoint{2.997805in}{3.316505in}}%
\pgfpathlineto{\pgfqpoint{3.002346in}{3.316505in}}%
\pgfpathlineto{\pgfqpoint{3.002346in}{3.313556in}}%
\pgfpathmoveto{\pgfqpoint{2.997805in}{3.316505in}}%
\pgfpathlineto{\pgfqpoint{2.997805in}{3.316505in}}%
\pgfpathlineto{\pgfqpoint{2.997805in}{3.319455in}}%
\pgfpathlineto{\pgfqpoint{3.002346in}{3.319455in}}%
\pgfpathlineto{\pgfqpoint{3.002346in}{3.316505in}}%
\pgfpathmoveto{\pgfqpoint{2.984181in}{3.319455in}}%
\pgfpathlineto{\pgfqpoint{2.984181in}{3.319455in}}%
\pgfpathlineto{\pgfqpoint{2.984181in}{3.322404in}}%
\pgfpathlineto{\pgfqpoint{2.988722in}{3.322404in}}%
\pgfpathlineto{\pgfqpoint{2.988722in}{3.319455in}}%
\pgfpathmoveto{\pgfqpoint{2.984181in}{3.322404in}}%
\pgfpathlineto{\pgfqpoint{2.984181in}{3.322404in}}%
\pgfpathlineto{\pgfqpoint{2.984181in}{3.325353in}}%
\pgfpathlineto{\pgfqpoint{2.988722in}{3.325353in}}%
\pgfpathlineto{\pgfqpoint{2.988722in}{3.322404in}}%
\pgfpathmoveto{\pgfqpoint{2.988722in}{3.319455in}}%
\pgfpathlineto{\pgfqpoint{2.988722in}{3.319455in}}%
\pgfpathlineto{\pgfqpoint{2.988722in}{3.322404in}}%
\pgfpathlineto{\pgfqpoint{2.993264in}{3.322404in}}%
\pgfpathlineto{\pgfqpoint{2.993264in}{3.319455in}}%
\pgfpathmoveto{\pgfqpoint{2.988722in}{3.322404in}}%
\pgfpathlineto{\pgfqpoint{2.988722in}{3.322404in}}%
\pgfpathlineto{\pgfqpoint{2.988722in}{3.325353in}}%
\pgfpathlineto{\pgfqpoint{2.993264in}{3.325353in}}%
\pgfpathlineto{\pgfqpoint{2.993264in}{3.322404in}}%
\pgfpathmoveto{\pgfqpoint{2.984181in}{3.325353in}}%
\pgfpathlineto{\pgfqpoint{2.984181in}{3.325353in}}%
\pgfpathlineto{\pgfqpoint{2.984181in}{3.328302in}}%
\pgfpathlineto{\pgfqpoint{2.988722in}{3.328302in}}%
\pgfpathlineto{\pgfqpoint{2.988722in}{3.325353in}}%
\pgfpathmoveto{\pgfqpoint{2.984181in}{3.328302in}}%
\pgfpathlineto{\pgfqpoint{2.984181in}{3.328302in}}%
\pgfpathlineto{\pgfqpoint{2.984181in}{3.331252in}}%
\pgfpathlineto{\pgfqpoint{2.988722in}{3.331252in}}%
\pgfpathlineto{\pgfqpoint{2.988722in}{3.328302in}}%
\pgfpathmoveto{\pgfqpoint{2.988722in}{3.325353in}}%
\pgfpathlineto{\pgfqpoint{2.988722in}{3.325353in}}%
\pgfpathlineto{\pgfqpoint{2.988722in}{3.328302in}}%
\pgfpathlineto{\pgfqpoint{2.993264in}{3.328302in}}%
\pgfpathlineto{\pgfqpoint{2.993264in}{3.325353in}}%
\pgfpathmoveto{\pgfqpoint{2.988722in}{3.328302in}}%
\pgfpathlineto{\pgfqpoint{2.988722in}{3.328302in}}%
\pgfpathlineto{\pgfqpoint{2.988722in}{3.331252in}}%
\pgfpathlineto{\pgfqpoint{2.993264in}{3.331252in}}%
\pgfpathlineto{\pgfqpoint{2.993264in}{3.328302in}}%
\pgfpathmoveto{\pgfqpoint{2.993264in}{3.319455in}}%
\pgfpathlineto{\pgfqpoint{2.993264in}{3.319455in}}%
\pgfpathlineto{\pgfqpoint{2.993264in}{3.322404in}}%
\pgfpathlineto{\pgfqpoint{2.997805in}{3.322404in}}%
\pgfpathlineto{\pgfqpoint{2.997805in}{3.319455in}}%
\pgfpathmoveto{\pgfqpoint{2.993264in}{3.322404in}}%
\pgfpathlineto{\pgfqpoint{2.993264in}{3.322404in}}%
\pgfpathlineto{\pgfqpoint{2.993264in}{3.325353in}}%
\pgfpathlineto{\pgfqpoint{2.997805in}{3.325353in}}%
\pgfpathlineto{\pgfqpoint{2.997805in}{3.322404in}}%
\pgfpathmoveto{\pgfqpoint{3.029593in}{3.260469in}}%
\pgfpathlineto{\pgfqpoint{3.029593in}{3.260469in}}%
\pgfpathlineto{\pgfqpoint{3.029593in}{3.263418in}}%
\pgfpathlineto{\pgfqpoint{3.034134in}{3.263418in}}%
\pgfpathlineto{\pgfqpoint{3.034134in}{3.260469in}}%
\pgfpathmoveto{\pgfqpoint{3.029593in}{3.263418in}}%
\pgfpathlineto{\pgfqpoint{3.029593in}{3.263418in}}%
\pgfpathlineto{\pgfqpoint{3.029593in}{3.266367in}}%
\pgfpathlineto{\pgfqpoint{3.034134in}{3.266367in}}%
\pgfpathlineto{\pgfqpoint{3.034134in}{3.263418in}}%
\pgfpathmoveto{\pgfqpoint{3.034134in}{3.260469in}}%
\pgfpathlineto{\pgfqpoint{3.034134in}{3.260469in}}%
\pgfpathlineto{\pgfqpoint{3.034134in}{3.263418in}}%
\pgfpathlineto{\pgfqpoint{3.038675in}{3.263418in}}%
\pgfpathlineto{\pgfqpoint{3.038675in}{3.260469in}}%
\pgfpathmoveto{\pgfqpoint{3.034134in}{3.263418in}}%
\pgfpathlineto{\pgfqpoint{3.034134in}{3.263418in}}%
\pgfpathlineto{\pgfqpoint{3.034134in}{3.266367in}}%
\pgfpathlineto{\pgfqpoint{3.038675in}{3.266367in}}%
\pgfpathlineto{\pgfqpoint{3.038675in}{3.263418in}}%
\pgfpathmoveto{\pgfqpoint{3.029593in}{3.266367in}}%
\pgfpathlineto{\pgfqpoint{3.029593in}{3.266367in}}%
\pgfpathlineto{\pgfqpoint{3.029593in}{3.269317in}}%
\pgfpathlineto{\pgfqpoint{3.034134in}{3.269317in}}%
\pgfpathlineto{\pgfqpoint{3.034134in}{3.266367in}}%
\pgfpathmoveto{\pgfqpoint{3.029593in}{3.269317in}}%
\pgfpathlineto{\pgfqpoint{3.029593in}{3.269317in}}%
\pgfpathlineto{\pgfqpoint{3.029593in}{3.272266in}}%
\pgfpathlineto{\pgfqpoint{3.034134in}{3.272266in}}%
\pgfpathlineto{\pgfqpoint{3.034134in}{3.269317in}}%
\pgfpathmoveto{\pgfqpoint{3.034134in}{3.266367in}}%
\pgfpathlineto{\pgfqpoint{3.034134in}{3.266367in}}%
\pgfpathlineto{\pgfqpoint{3.034134in}{3.269317in}}%
\pgfpathlineto{\pgfqpoint{3.038675in}{3.269317in}}%
\pgfpathlineto{\pgfqpoint{3.038675in}{3.266367in}}%
\pgfpathmoveto{\pgfqpoint{3.034134in}{3.269317in}}%
\pgfpathlineto{\pgfqpoint{3.034134in}{3.269317in}}%
\pgfpathlineto{\pgfqpoint{3.034134in}{3.272266in}}%
\pgfpathlineto{\pgfqpoint{3.038675in}{3.272266in}}%
\pgfpathlineto{\pgfqpoint{3.038675in}{3.269317in}}%
\pgfpathmoveto{\pgfqpoint{3.020510in}{3.272266in}}%
\pgfpathlineto{\pgfqpoint{3.020510in}{3.272266in}}%
\pgfpathlineto{\pgfqpoint{3.020510in}{3.275215in}}%
\pgfpathlineto{\pgfqpoint{3.025051in}{3.275215in}}%
\pgfpathlineto{\pgfqpoint{3.025051in}{3.272266in}}%
\pgfpathmoveto{\pgfqpoint{3.020510in}{3.275215in}}%
\pgfpathlineto{\pgfqpoint{3.020510in}{3.275215in}}%
\pgfpathlineto{\pgfqpoint{3.020510in}{3.278164in}}%
\pgfpathlineto{\pgfqpoint{3.025051in}{3.278164in}}%
\pgfpathlineto{\pgfqpoint{3.025051in}{3.275215in}}%
\pgfpathmoveto{\pgfqpoint{3.025051in}{3.272266in}}%
\pgfpathlineto{\pgfqpoint{3.025051in}{3.272266in}}%
\pgfpathlineto{\pgfqpoint{3.025051in}{3.275215in}}%
\pgfpathlineto{\pgfqpoint{3.029593in}{3.275215in}}%
\pgfpathlineto{\pgfqpoint{3.029593in}{3.272266in}}%
\pgfpathmoveto{\pgfqpoint{3.025051in}{3.275215in}}%
\pgfpathlineto{\pgfqpoint{3.025051in}{3.275215in}}%
\pgfpathlineto{\pgfqpoint{3.025051in}{3.278164in}}%
\pgfpathlineto{\pgfqpoint{3.029593in}{3.278164in}}%
\pgfpathlineto{\pgfqpoint{3.029593in}{3.275215in}}%
\pgfpathmoveto{\pgfqpoint{3.020510in}{3.278164in}}%
\pgfpathlineto{\pgfqpoint{3.020510in}{3.278164in}}%
\pgfpathlineto{\pgfqpoint{3.020510in}{3.281114in}}%
\pgfpathlineto{\pgfqpoint{3.025051in}{3.281114in}}%
\pgfpathlineto{\pgfqpoint{3.025051in}{3.278164in}}%
\pgfpathmoveto{\pgfqpoint{3.020510in}{3.281114in}}%
\pgfpathlineto{\pgfqpoint{3.020510in}{3.281114in}}%
\pgfpathlineto{\pgfqpoint{3.020510in}{3.284063in}}%
\pgfpathlineto{\pgfqpoint{3.025051in}{3.284063in}}%
\pgfpathlineto{\pgfqpoint{3.025051in}{3.281114in}}%
\pgfpathmoveto{\pgfqpoint{3.025051in}{3.278164in}}%
\pgfpathlineto{\pgfqpoint{3.025051in}{3.278164in}}%
\pgfpathlineto{\pgfqpoint{3.025051in}{3.281114in}}%
\pgfpathlineto{\pgfqpoint{3.029593in}{3.281114in}}%
\pgfpathlineto{\pgfqpoint{3.029593in}{3.278164in}}%
\pgfpathmoveto{\pgfqpoint{3.025051in}{3.281114in}}%
\pgfpathlineto{\pgfqpoint{3.025051in}{3.281114in}}%
\pgfpathlineto{\pgfqpoint{3.025051in}{3.284063in}}%
\pgfpathlineto{\pgfqpoint{3.029593in}{3.284063in}}%
\pgfpathlineto{\pgfqpoint{3.029593in}{3.281114in}}%
\pgfpathmoveto{\pgfqpoint{3.029593in}{3.272266in}}%
\pgfpathlineto{\pgfqpoint{3.029593in}{3.272266in}}%
\pgfpathlineto{\pgfqpoint{3.029593in}{3.275215in}}%
\pgfpathlineto{\pgfqpoint{3.034134in}{3.275215in}}%
\pgfpathlineto{\pgfqpoint{3.034134in}{3.272266in}}%
\pgfpathmoveto{\pgfqpoint{3.029593in}{3.275215in}}%
\pgfpathlineto{\pgfqpoint{3.029593in}{3.275215in}}%
\pgfpathlineto{\pgfqpoint{3.029593in}{3.278164in}}%
\pgfpathlineto{\pgfqpoint{3.034134in}{3.278164in}}%
\pgfpathlineto{\pgfqpoint{3.034134in}{3.275215in}}%
\pgfpathmoveto{\pgfqpoint{3.047757in}{3.236874in}}%
\pgfpathlineto{\pgfqpoint{3.047757in}{3.236874in}}%
\pgfpathlineto{\pgfqpoint{3.047757in}{3.239824in}}%
\pgfpathlineto{\pgfqpoint{3.052298in}{3.239824in}}%
\pgfpathlineto{\pgfqpoint{3.052298in}{3.236874in}}%
\pgfpathmoveto{\pgfqpoint{3.047757in}{3.239824in}}%
\pgfpathlineto{\pgfqpoint{3.047757in}{3.239824in}}%
\pgfpathlineto{\pgfqpoint{3.047757in}{3.242773in}}%
\pgfpathlineto{\pgfqpoint{3.052298in}{3.242773in}}%
\pgfpathlineto{\pgfqpoint{3.052298in}{3.239824in}}%
\pgfpathmoveto{\pgfqpoint{3.052298in}{3.236874in}}%
\pgfpathlineto{\pgfqpoint{3.052298in}{3.236874in}}%
\pgfpathlineto{\pgfqpoint{3.052298in}{3.239824in}}%
\pgfpathlineto{\pgfqpoint{3.056839in}{3.239824in}}%
\pgfpathlineto{\pgfqpoint{3.056839in}{3.236874in}}%
\pgfpathmoveto{\pgfqpoint{3.052298in}{3.239824in}}%
\pgfpathlineto{\pgfqpoint{3.052298in}{3.239824in}}%
\pgfpathlineto{\pgfqpoint{3.052298in}{3.242773in}}%
\pgfpathlineto{\pgfqpoint{3.056839in}{3.242773in}}%
\pgfpathlineto{\pgfqpoint{3.056839in}{3.239824in}}%
\pgfpathmoveto{\pgfqpoint{3.047757in}{3.242773in}}%
\pgfpathlineto{\pgfqpoint{3.047757in}{3.242773in}}%
\pgfpathlineto{\pgfqpoint{3.047757in}{3.245722in}}%
\pgfpathlineto{\pgfqpoint{3.052298in}{3.245722in}}%
\pgfpathlineto{\pgfqpoint{3.052298in}{3.242773in}}%
\pgfpathmoveto{\pgfqpoint{3.047757in}{3.245722in}}%
\pgfpathlineto{\pgfqpoint{3.047757in}{3.245722in}}%
\pgfpathlineto{\pgfqpoint{3.047757in}{3.248672in}}%
\pgfpathlineto{\pgfqpoint{3.052298in}{3.248672in}}%
\pgfpathlineto{\pgfqpoint{3.052298in}{3.245722in}}%
\pgfpathmoveto{\pgfqpoint{3.052298in}{3.242773in}}%
\pgfpathlineto{\pgfqpoint{3.052298in}{3.242773in}}%
\pgfpathlineto{\pgfqpoint{3.052298in}{3.245722in}}%
\pgfpathlineto{\pgfqpoint{3.056839in}{3.245722in}}%
\pgfpathlineto{\pgfqpoint{3.056839in}{3.242773in}}%
\pgfpathmoveto{\pgfqpoint{3.052298in}{3.245722in}}%
\pgfpathlineto{\pgfqpoint{3.052298in}{3.245722in}}%
\pgfpathlineto{\pgfqpoint{3.052298in}{3.248672in}}%
\pgfpathlineto{\pgfqpoint{3.056839in}{3.248672in}}%
\pgfpathlineto{\pgfqpoint{3.056839in}{3.245722in}}%
\pgfpathmoveto{\pgfqpoint{3.038675in}{3.248672in}}%
\pgfpathlineto{\pgfqpoint{3.038675in}{3.248672in}}%
\pgfpathlineto{\pgfqpoint{3.038675in}{3.251621in}}%
\pgfpathlineto{\pgfqpoint{3.043216in}{3.251621in}}%
\pgfpathlineto{\pgfqpoint{3.043216in}{3.248672in}}%
\pgfpathmoveto{\pgfqpoint{3.038675in}{3.251621in}}%
\pgfpathlineto{\pgfqpoint{3.038675in}{3.251621in}}%
\pgfpathlineto{\pgfqpoint{3.038675in}{3.254570in}}%
\pgfpathlineto{\pgfqpoint{3.043216in}{3.254570in}}%
\pgfpathlineto{\pgfqpoint{3.043216in}{3.251621in}}%
\pgfpathmoveto{\pgfqpoint{3.043216in}{3.248672in}}%
\pgfpathlineto{\pgfqpoint{3.043216in}{3.248672in}}%
\pgfpathlineto{\pgfqpoint{3.043216in}{3.251621in}}%
\pgfpathlineto{\pgfqpoint{3.047757in}{3.251621in}}%
\pgfpathlineto{\pgfqpoint{3.047757in}{3.248672in}}%
\pgfpathmoveto{\pgfqpoint{3.043216in}{3.251621in}}%
\pgfpathlineto{\pgfqpoint{3.043216in}{3.251621in}}%
\pgfpathlineto{\pgfqpoint{3.043216in}{3.254570in}}%
\pgfpathlineto{\pgfqpoint{3.047757in}{3.254570in}}%
\pgfpathlineto{\pgfqpoint{3.047757in}{3.251621in}}%
\pgfpathmoveto{\pgfqpoint{3.038675in}{3.254570in}}%
\pgfpathlineto{\pgfqpoint{3.038675in}{3.254570in}}%
\pgfpathlineto{\pgfqpoint{3.038675in}{3.257519in}}%
\pgfpathlineto{\pgfqpoint{3.043216in}{3.257519in}}%
\pgfpathlineto{\pgfqpoint{3.043216in}{3.254570in}}%
\pgfpathmoveto{\pgfqpoint{3.038675in}{3.257519in}}%
\pgfpathlineto{\pgfqpoint{3.038675in}{3.257519in}}%
\pgfpathlineto{\pgfqpoint{3.038675in}{3.260469in}}%
\pgfpathlineto{\pgfqpoint{3.043216in}{3.260469in}}%
\pgfpathlineto{\pgfqpoint{3.043216in}{3.257519in}}%
\pgfpathmoveto{\pgfqpoint{3.043216in}{3.254570in}}%
\pgfpathlineto{\pgfqpoint{3.043216in}{3.254570in}}%
\pgfpathlineto{\pgfqpoint{3.043216in}{3.257519in}}%
\pgfpathlineto{\pgfqpoint{3.047757in}{3.257519in}}%
\pgfpathlineto{\pgfqpoint{3.047757in}{3.254570in}}%
\pgfpathmoveto{\pgfqpoint{3.043216in}{3.257519in}}%
\pgfpathlineto{\pgfqpoint{3.043216in}{3.257519in}}%
\pgfpathlineto{\pgfqpoint{3.043216in}{3.260469in}}%
\pgfpathlineto{\pgfqpoint{3.047757in}{3.260469in}}%
\pgfpathlineto{\pgfqpoint{3.047757in}{3.257519in}}%
\pgfpathmoveto{\pgfqpoint{3.047757in}{3.248672in}}%
\pgfpathlineto{\pgfqpoint{3.047757in}{3.248672in}}%
\pgfpathlineto{\pgfqpoint{3.047757in}{3.251621in}}%
\pgfpathlineto{\pgfqpoint{3.052298in}{3.251621in}}%
\pgfpathlineto{\pgfqpoint{3.052298in}{3.248672in}}%
\pgfpathmoveto{\pgfqpoint{3.047757in}{3.251621in}}%
\pgfpathlineto{\pgfqpoint{3.047757in}{3.251621in}}%
\pgfpathlineto{\pgfqpoint{3.047757in}{3.254570in}}%
\pgfpathlineto{\pgfqpoint{3.052298in}{3.254570in}}%
\pgfpathlineto{\pgfqpoint{3.052298in}{3.251621in}}%
\pgfpathmoveto{\pgfqpoint{3.056839in}{3.236874in}}%
\pgfpathlineto{\pgfqpoint{3.056839in}{3.236874in}}%
\pgfpathlineto{\pgfqpoint{3.056839in}{3.239824in}}%
\pgfpathlineto{\pgfqpoint{3.061380in}{3.239824in}}%
\pgfpathlineto{\pgfqpoint{3.061380in}{3.236874in}}%
\pgfpathmoveto{\pgfqpoint{3.056839in}{3.239824in}}%
\pgfpathlineto{\pgfqpoint{3.056839in}{3.239824in}}%
\pgfpathlineto{\pgfqpoint{3.056839in}{3.242773in}}%
\pgfpathlineto{\pgfqpoint{3.061380in}{3.242773in}}%
\pgfpathlineto{\pgfqpoint{3.061380in}{3.239824in}}%
\pgfpathmoveto{\pgfqpoint{3.038675in}{3.260469in}}%
\pgfpathlineto{\pgfqpoint{3.038675in}{3.260469in}}%
\pgfpathlineto{\pgfqpoint{3.038675in}{3.263418in}}%
\pgfpathlineto{\pgfqpoint{3.043216in}{3.263418in}}%
\pgfpathlineto{\pgfqpoint{3.043216in}{3.260469in}}%
\pgfpathmoveto{\pgfqpoint{3.038675in}{3.263418in}}%
\pgfpathlineto{\pgfqpoint{3.038675in}{3.263418in}}%
\pgfpathlineto{\pgfqpoint{3.038675in}{3.266367in}}%
\pgfpathlineto{\pgfqpoint{3.043216in}{3.266367in}}%
\pgfpathlineto{\pgfqpoint{3.043216in}{3.263418in}}%
\pgfpathmoveto{\pgfqpoint{3.011428in}{3.284063in}}%
\pgfpathlineto{\pgfqpoint{3.011428in}{3.284063in}}%
\pgfpathlineto{\pgfqpoint{3.011428in}{3.287012in}}%
\pgfpathlineto{\pgfqpoint{3.015969in}{3.287012in}}%
\pgfpathlineto{\pgfqpoint{3.015969in}{3.284063in}}%
\pgfpathmoveto{\pgfqpoint{3.011428in}{3.287012in}}%
\pgfpathlineto{\pgfqpoint{3.011428in}{3.287012in}}%
\pgfpathlineto{\pgfqpoint{3.011428in}{3.289962in}}%
\pgfpathlineto{\pgfqpoint{3.015969in}{3.289962in}}%
\pgfpathlineto{\pgfqpoint{3.015969in}{3.287012in}}%
\pgfpathmoveto{\pgfqpoint{3.015969in}{3.284063in}}%
\pgfpathlineto{\pgfqpoint{3.015969in}{3.284063in}}%
\pgfpathlineto{\pgfqpoint{3.015969in}{3.287012in}}%
\pgfpathlineto{\pgfqpoint{3.020510in}{3.287012in}}%
\pgfpathlineto{\pgfqpoint{3.020510in}{3.284063in}}%
\pgfpathmoveto{\pgfqpoint{3.015969in}{3.287012in}}%
\pgfpathlineto{\pgfqpoint{3.015969in}{3.287012in}}%
\pgfpathlineto{\pgfqpoint{3.015969in}{3.289962in}}%
\pgfpathlineto{\pgfqpoint{3.020510in}{3.289962in}}%
\pgfpathlineto{\pgfqpoint{3.020510in}{3.287012in}}%
\pgfpathmoveto{\pgfqpoint{3.011428in}{3.289962in}}%
\pgfpathlineto{\pgfqpoint{3.011428in}{3.289962in}}%
\pgfpathlineto{\pgfqpoint{3.011428in}{3.292911in}}%
\pgfpathlineto{\pgfqpoint{3.015969in}{3.292911in}}%
\pgfpathlineto{\pgfqpoint{3.015969in}{3.289962in}}%
\pgfpathmoveto{\pgfqpoint{3.011428in}{3.292911in}}%
\pgfpathlineto{\pgfqpoint{3.011428in}{3.292911in}}%
\pgfpathlineto{\pgfqpoint{3.011428in}{3.295860in}}%
\pgfpathlineto{\pgfqpoint{3.015969in}{3.295860in}}%
\pgfpathlineto{\pgfqpoint{3.015969in}{3.292911in}}%
\pgfpathmoveto{\pgfqpoint{3.015969in}{3.289962in}}%
\pgfpathlineto{\pgfqpoint{3.015969in}{3.289962in}}%
\pgfpathlineto{\pgfqpoint{3.015969in}{3.292911in}}%
\pgfpathlineto{\pgfqpoint{3.020510in}{3.292911in}}%
\pgfpathlineto{\pgfqpoint{3.020510in}{3.289962in}}%
\pgfpathmoveto{\pgfqpoint{3.015969in}{3.292911in}}%
\pgfpathlineto{\pgfqpoint{3.015969in}{3.292911in}}%
\pgfpathlineto{\pgfqpoint{3.015969in}{3.295860in}}%
\pgfpathlineto{\pgfqpoint{3.020510in}{3.295860in}}%
\pgfpathlineto{\pgfqpoint{3.020510in}{3.292911in}}%
\pgfpathmoveto{\pgfqpoint{3.002346in}{3.295860in}}%
\pgfpathlineto{\pgfqpoint{3.002346in}{3.295860in}}%
\pgfpathlineto{\pgfqpoint{3.002346in}{3.298810in}}%
\pgfpathlineto{\pgfqpoint{3.006887in}{3.298810in}}%
\pgfpathlineto{\pgfqpoint{3.006887in}{3.295860in}}%
\pgfpathmoveto{\pgfqpoint{3.002346in}{3.298810in}}%
\pgfpathlineto{\pgfqpoint{3.002346in}{3.298810in}}%
\pgfpathlineto{\pgfqpoint{3.002346in}{3.301759in}}%
\pgfpathlineto{\pgfqpoint{3.006887in}{3.301759in}}%
\pgfpathlineto{\pgfqpoint{3.006887in}{3.298810in}}%
\pgfpathmoveto{\pgfqpoint{3.006887in}{3.295860in}}%
\pgfpathlineto{\pgfqpoint{3.006887in}{3.295860in}}%
\pgfpathlineto{\pgfqpoint{3.006887in}{3.298810in}}%
\pgfpathlineto{\pgfqpoint{3.011428in}{3.298810in}}%
\pgfpathlineto{\pgfqpoint{3.011428in}{3.295860in}}%
\pgfpathmoveto{\pgfqpoint{3.006887in}{3.298810in}}%
\pgfpathlineto{\pgfqpoint{3.006887in}{3.298810in}}%
\pgfpathlineto{\pgfqpoint{3.006887in}{3.301759in}}%
\pgfpathlineto{\pgfqpoint{3.011428in}{3.301759in}}%
\pgfpathlineto{\pgfqpoint{3.011428in}{3.298810in}}%
\pgfpathmoveto{\pgfqpoint{3.002346in}{3.301759in}}%
\pgfpathlineto{\pgfqpoint{3.002346in}{3.301759in}}%
\pgfpathlineto{\pgfqpoint{3.002346in}{3.304708in}}%
\pgfpathlineto{\pgfqpoint{3.006887in}{3.304708in}}%
\pgfpathlineto{\pgfqpoint{3.006887in}{3.301759in}}%
\pgfpathmoveto{\pgfqpoint{3.002346in}{3.304708in}}%
\pgfpathlineto{\pgfqpoint{3.002346in}{3.304708in}}%
\pgfpathlineto{\pgfqpoint{3.002346in}{3.307657in}}%
\pgfpathlineto{\pgfqpoint{3.006887in}{3.307657in}}%
\pgfpathlineto{\pgfqpoint{3.006887in}{3.304708in}}%
\pgfpathmoveto{\pgfqpoint{3.006887in}{3.301759in}}%
\pgfpathlineto{\pgfqpoint{3.006887in}{3.301759in}}%
\pgfpathlineto{\pgfqpoint{3.006887in}{3.304708in}}%
\pgfpathlineto{\pgfqpoint{3.011428in}{3.304708in}}%
\pgfpathlineto{\pgfqpoint{3.011428in}{3.301759in}}%
\pgfpathmoveto{\pgfqpoint{3.006887in}{3.304708in}}%
\pgfpathlineto{\pgfqpoint{3.006887in}{3.304708in}}%
\pgfpathlineto{\pgfqpoint{3.006887in}{3.307657in}}%
\pgfpathlineto{\pgfqpoint{3.011428in}{3.307657in}}%
\pgfpathlineto{\pgfqpoint{3.011428in}{3.304708in}}%
\pgfpathmoveto{\pgfqpoint{3.011428in}{3.295860in}}%
\pgfpathlineto{\pgfqpoint{3.011428in}{3.295860in}}%
\pgfpathlineto{\pgfqpoint{3.011428in}{3.298810in}}%
\pgfpathlineto{\pgfqpoint{3.015969in}{3.298810in}}%
\pgfpathlineto{\pgfqpoint{3.015969in}{3.295860in}}%
\pgfpathmoveto{\pgfqpoint{3.011428in}{3.298810in}}%
\pgfpathlineto{\pgfqpoint{3.011428in}{3.298810in}}%
\pgfpathlineto{\pgfqpoint{3.011428in}{3.301759in}}%
\pgfpathlineto{\pgfqpoint{3.015969in}{3.301759in}}%
\pgfpathlineto{\pgfqpoint{3.015969in}{3.298810in}}%
\pgfpathmoveto{\pgfqpoint{3.020510in}{3.284063in}}%
\pgfpathlineto{\pgfqpoint{3.020510in}{3.284063in}}%
\pgfpathlineto{\pgfqpoint{3.020510in}{3.287012in}}%
\pgfpathlineto{\pgfqpoint{3.025051in}{3.287012in}}%
\pgfpathlineto{\pgfqpoint{3.025051in}{3.284063in}}%
\pgfpathmoveto{\pgfqpoint{3.020510in}{3.287012in}}%
\pgfpathlineto{\pgfqpoint{3.020510in}{3.287012in}}%
\pgfpathlineto{\pgfqpoint{3.020510in}{3.289962in}}%
\pgfpathlineto{\pgfqpoint{3.025051in}{3.289962in}}%
\pgfpathlineto{\pgfqpoint{3.025051in}{3.287012in}}%
\pgfpathmoveto{\pgfqpoint{3.002346in}{3.307657in}}%
\pgfpathlineto{\pgfqpoint{3.002346in}{3.307657in}}%
\pgfpathlineto{\pgfqpoint{3.002346in}{3.310607in}}%
\pgfpathlineto{\pgfqpoint{3.006887in}{3.310607in}}%
\pgfpathlineto{\pgfqpoint{3.006887in}{3.307657in}}%
\pgfpathmoveto{\pgfqpoint{3.002346in}{3.310607in}}%
\pgfpathlineto{\pgfqpoint{3.002346in}{3.310607in}}%
\pgfpathlineto{\pgfqpoint{3.002346in}{3.313556in}}%
\pgfpathlineto{\pgfqpoint{3.006887in}{3.313556in}}%
\pgfpathlineto{\pgfqpoint{3.006887in}{3.310607in}}%
\pgfpathmoveto{\pgfqpoint{2.956935in}{3.354845in}}%
\pgfpathlineto{\pgfqpoint{2.956935in}{3.354845in}}%
\pgfpathlineto{\pgfqpoint{2.956935in}{3.357795in}}%
\pgfpathlineto{\pgfqpoint{2.961476in}{3.357795in}}%
\pgfpathlineto{\pgfqpoint{2.961476in}{3.354845in}}%
\pgfpathmoveto{\pgfqpoint{2.956935in}{3.357795in}}%
\pgfpathlineto{\pgfqpoint{2.956935in}{3.357795in}}%
\pgfpathlineto{\pgfqpoint{2.956935in}{3.360744in}}%
\pgfpathlineto{\pgfqpoint{2.961476in}{3.360744in}}%
\pgfpathlineto{\pgfqpoint{2.961476in}{3.357795in}}%
\pgfpathmoveto{\pgfqpoint{2.961476in}{3.354845in}}%
\pgfpathlineto{\pgfqpoint{2.961476in}{3.354845in}}%
\pgfpathlineto{\pgfqpoint{2.961476in}{3.357795in}}%
\pgfpathlineto{\pgfqpoint{2.966017in}{3.357795in}}%
\pgfpathlineto{\pgfqpoint{2.966017in}{3.354845in}}%
\pgfpathmoveto{\pgfqpoint{2.961476in}{3.357795in}}%
\pgfpathlineto{\pgfqpoint{2.961476in}{3.357795in}}%
\pgfpathlineto{\pgfqpoint{2.961476in}{3.360744in}}%
\pgfpathlineto{\pgfqpoint{2.966017in}{3.360744in}}%
\pgfpathlineto{\pgfqpoint{2.966017in}{3.357795in}}%
\pgfpathmoveto{\pgfqpoint{2.956935in}{3.360744in}}%
\pgfpathlineto{\pgfqpoint{2.956935in}{3.360744in}}%
\pgfpathlineto{\pgfqpoint{2.956935in}{3.363693in}}%
\pgfpathlineto{\pgfqpoint{2.961476in}{3.363693in}}%
\pgfpathlineto{\pgfqpoint{2.961476in}{3.360744in}}%
\pgfpathmoveto{\pgfqpoint{2.956935in}{3.363693in}}%
\pgfpathlineto{\pgfqpoint{2.956935in}{3.363693in}}%
\pgfpathlineto{\pgfqpoint{2.956935in}{3.366642in}}%
\pgfpathlineto{\pgfqpoint{2.961476in}{3.366642in}}%
\pgfpathlineto{\pgfqpoint{2.961476in}{3.363693in}}%
\pgfpathmoveto{\pgfqpoint{2.961476in}{3.360744in}}%
\pgfpathlineto{\pgfqpoint{2.961476in}{3.360744in}}%
\pgfpathlineto{\pgfqpoint{2.961476in}{3.363693in}}%
\pgfpathlineto{\pgfqpoint{2.966017in}{3.363693in}}%
\pgfpathlineto{\pgfqpoint{2.966017in}{3.360744in}}%
\pgfpathmoveto{\pgfqpoint{2.961476in}{3.363693in}}%
\pgfpathlineto{\pgfqpoint{2.961476in}{3.363693in}}%
\pgfpathlineto{\pgfqpoint{2.961476in}{3.366642in}}%
\pgfpathlineto{\pgfqpoint{2.966017in}{3.366642in}}%
\pgfpathlineto{\pgfqpoint{2.966017in}{3.363693in}}%
\pgfpathmoveto{\pgfqpoint{2.947852in}{3.366642in}}%
\pgfpathlineto{\pgfqpoint{2.947852in}{3.366642in}}%
\pgfpathlineto{\pgfqpoint{2.947852in}{3.369592in}}%
\pgfpathlineto{\pgfqpoint{2.952394in}{3.369592in}}%
\pgfpathlineto{\pgfqpoint{2.952394in}{3.366642in}}%
\pgfpathmoveto{\pgfqpoint{2.947852in}{3.369592in}}%
\pgfpathlineto{\pgfqpoint{2.947852in}{3.369592in}}%
\pgfpathlineto{\pgfqpoint{2.947852in}{3.372541in}}%
\pgfpathlineto{\pgfqpoint{2.952394in}{3.372541in}}%
\pgfpathlineto{\pgfqpoint{2.952394in}{3.369592in}}%
\pgfpathmoveto{\pgfqpoint{2.952394in}{3.366642in}}%
\pgfpathlineto{\pgfqpoint{2.952394in}{3.366642in}}%
\pgfpathlineto{\pgfqpoint{2.952394in}{3.369592in}}%
\pgfpathlineto{\pgfqpoint{2.956935in}{3.369592in}}%
\pgfpathlineto{\pgfqpoint{2.956935in}{3.366642in}}%
\pgfpathmoveto{\pgfqpoint{2.952394in}{3.369592in}}%
\pgfpathlineto{\pgfqpoint{2.952394in}{3.369592in}}%
\pgfpathlineto{\pgfqpoint{2.952394in}{3.372541in}}%
\pgfpathlineto{\pgfqpoint{2.956935in}{3.372541in}}%
\pgfpathlineto{\pgfqpoint{2.956935in}{3.369592in}}%
\pgfpathmoveto{\pgfqpoint{2.947852in}{3.372541in}}%
\pgfpathlineto{\pgfqpoint{2.947852in}{3.372541in}}%
\pgfpathlineto{\pgfqpoint{2.947852in}{3.375490in}}%
\pgfpathlineto{\pgfqpoint{2.952394in}{3.375490in}}%
\pgfpathlineto{\pgfqpoint{2.952394in}{3.372541in}}%
\pgfpathmoveto{\pgfqpoint{2.947852in}{3.375490in}}%
\pgfpathlineto{\pgfqpoint{2.947852in}{3.375490in}}%
\pgfpathlineto{\pgfqpoint{2.947852in}{3.378439in}}%
\pgfpathlineto{\pgfqpoint{2.952394in}{3.378439in}}%
\pgfpathlineto{\pgfqpoint{2.952394in}{3.375490in}}%
\pgfpathmoveto{\pgfqpoint{2.952394in}{3.372541in}}%
\pgfpathlineto{\pgfqpoint{2.952394in}{3.372541in}}%
\pgfpathlineto{\pgfqpoint{2.952394in}{3.375490in}}%
\pgfpathlineto{\pgfqpoint{2.956935in}{3.375490in}}%
\pgfpathlineto{\pgfqpoint{2.956935in}{3.372541in}}%
\pgfpathmoveto{\pgfqpoint{2.952394in}{3.375490in}}%
\pgfpathlineto{\pgfqpoint{2.952394in}{3.375490in}}%
\pgfpathlineto{\pgfqpoint{2.952394in}{3.378439in}}%
\pgfpathlineto{\pgfqpoint{2.956935in}{3.378439in}}%
\pgfpathlineto{\pgfqpoint{2.956935in}{3.375490in}}%
\pgfpathmoveto{\pgfqpoint{2.956935in}{3.366642in}}%
\pgfpathlineto{\pgfqpoint{2.956935in}{3.366642in}}%
\pgfpathlineto{\pgfqpoint{2.956935in}{3.369592in}}%
\pgfpathlineto{\pgfqpoint{2.961476in}{3.369592in}}%
\pgfpathlineto{\pgfqpoint{2.961476in}{3.366642in}}%
\pgfpathmoveto{\pgfqpoint{2.956935in}{3.369592in}}%
\pgfpathlineto{\pgfqpoint{2.956935in}{3.369592in}}%
\pgfpathlineto{\pgfqpoint{2.956935in}{3.372541in}}%
\pgfpathlineto{\pgfqpoint{2.961476in}{3.372541in}}%
\pgfpathlineto{\pgfqpoint{2.961476in}{3.369592in}}%
\pgfpathmoveto{\pgfqpoint{2.975099in}{3.331252in}}%
\pgfpathlineto{\pgfqpoint{2.975099in}{3.331252in}}%
\pgfpathlineto{\pgfqpoint{2.975099in}{3.334201in}}%
\pgfpathlineto{\pgfqpoint{2.979640in}{3.334201in}}%
\pgfpathlineto{\pgfqpoint{2.979640in}{3.331252in}}%
\pgfpathmoveto{\pgfqpoint{2.975099in}{3.334201in}}%
\pgfpathlineto{\pgfqpoint{2.975099in}{3.334201in}}%
\pgfpathlineto{\pgfqpoint{2.975099in}{3.337150in}}%
\pgfpathlineto{\pgfqpoint{2.979640in}{3.337150in}}%
\pgfpathlineto{\pgfqpoint{2.979640in}{3.334201in}}%
\pgfpathmoveto{\pgfqpoint{2.979640in}{3.331252in}}%
\pgfpathlineto{\pgfqpoint{2.979640in}{3.331252in}}%
\pgfpathlineto{\pgfqpoint{2.979640in}{3.334201in}}%
\pgfpathlineto{\pgfqpoint{2.984181in}{3.334201in}}%
\pgfpathlineto{\pgfqpoint{2.984181in}{3.331252in}}%
\pgfpathmoveto{\pgfqpoint{2.979640in}{3.334201in}}%
\pgfpathlineto{\pgfqpoint{2.979640in}{3.334201in}}%
\pgfpathlineto{\pgfqpoint{2.979640in}{3.337150in}}%
\pgfpathlineto{\pgfqpoint{2.984181in}{3.337150in}}%
\pgfpathlineto{\pgfqpoint{2.984181in}{3.334201in}}%
\pgfpathmoveto{\pgfqpoint{2.975099in}{3.337150in}}%
\pgfpathlineto{\pgfqpoint{2.975099in}{3.337150in}}%
\pgfpathlineto{\pgfqpoint{2.975099in}{3.340099in}}%
\pgfpathlineto{\pgfqpoint{2.979640in}{3.340099in}}%
\pgfpathlineto{\pgfqpoint{2.979640in}{3.337150in}}%
\pgfpathmoveto{\pgfqpoint{2.975099in}{3.340099in}}%
\pgfpathlineto{\pgfqpoint{2.975099in}{3.340099in}}%
\pgfpathlineto{\pgfqpoint{2.975099in}{3.343049in}}%
\pgfpathlineto{\pgfqpoint{2.979640in}{3.343049in}}%
\pgfpathlineto{\pgfqpoint{2.979640in}{3.340099in}}%
\pgfpathmoveto{\pgfqpoint{2.979640in}{3.337150in}}%
\pgfpathlineto{\pgfqpoint{2.979640in}{3.337150in}}%
\pgfpathlineto{\pgfqpoint{2.979640in}{3.340099in}}%
\pgfpathlineto{\pgfqpoint{2.984181in}{3.340099in}}%
\pgfpathlineto{\pgfqpoint{2.984181in}{3.337150in}}%
\pgfpathmoveto{\pgfqpoint{2.979640in}{3.340099in}}%
\pgfpathlineto{\pgfqpoint{2.979640in}{3.340099in}}%
\pgfpathlineto{\pgfqpoint{2.979640in}{3.343049in}}%
\pgfpathlineto{\pgfqpoint{2.984181in}{3.343049in}}%
\pgfpathlineto{\pgfqpoint{2.984181in}{3.340099in}}%
\pgfpathmoveto{\pgfqpoint{2.966017in}{3.343049in}}%
\pgfpathlineto{\pgfqpoint{2.966017in}{3.343049in}}%
\pgfpathlineto{\pgfqpoint{2.966017in}{3.345998in}}%
\pgfpathlineto{\pgfqpoint{2.970558in}{3.345998in}}%
\pgfpathlineto{\pgfqpoint{2.970558in}{3.343049in}}%
\pgfpathmoveto{\pgfqpoint{2.966017in}{3.345998in}}%
\pgfpathlineto{\pgfqpoint{2.966017in}{3.345998in}}%
\pgfpathlineto{\pgfqpoint{2.966017in}{3.348947in}}%
\pgfpathlineto{\pgfqpoint{2.970558in}{3.348947in}}%
\pgfpathlineto{\pgfqpoint{2.970558in}{3.345998in}}%
\pgfpathmoveto{\pgfqpoint{2.970558in}{3.343049in}}%
\pgfpathlineto{\pgfqpoint{2.970558in}{3.343049in}}%
\pgfpathlineto{\pgfqpoint{2.970558in}{3.345998in}}%
\pgfpathlineto{\pgfqpoint{2.975099in}{3.345998in}}%
\pgfpathlineto{\pgfqpoint{2.975099in}{3.343049in}}%
\pgfpathmoveto{\pgfqpoint{2.970558in}{3.345998in}}%
\pgfpathlineto{\pgfqpoint{2.970558in}{3.345998in}}%
\pgfpathlineto{\pgfqpoint{2.970558in}{3.348947in}}%
\pgfpathlineto{\pgfqpoint{2.975099in}{3.348947in}}%
\pgfpathlineto{\pgfqpoint{2.975099in}{3.345998in}}%
\pgfpathmoveto{\pgfqpoint{2.966017in}{3.348947in}}%
\pgfpathlineto{\pgfqpoint{2.966017in}{3.348947in}}%
\pgfpathlineto{\pgfqpoint{2.966017in}{3.351896in}}%
\pgfpathlineto{\pgfqpoint{2.970558in}{3.351896in}}%
\pgfpathlineto{\pgfqpoint{2.970558in}{3.348947in}}%
\pgfpathmoveto{\pgfqpoint{2.966017in}{3.351896in}}%
\pgfpathlineto{\pgfqpoint{2.966017in}{3.351896in}}%
\pgfpathlineto{\pgfqpoint{2.966017in}{3.354845in}}%
\pgfpathlineto{\pgfqpoint{2.970558in}{3.354845in}}%
\pgfpathlineto{\pgfqpoint{2.970558in}{3.351896in}}%
\pgfpathmoveto{\pgfqpoint{2.970558in}{3.348947in}}%
\pgfpathlineto{\pgfqpoint{2.970558in}{3.348947in}}%
\pgfpathlineto{\pgfqpoint{2.970558in}{3.351896in}}%
\pgfpathlineto{\pgfqpoint{2.975099in}{3.351896in}}%
\pgfpathlineto{\pgfqpoint{2.975099in}{3.348947in}}%
\pgfpathmoveto{\pgfqpoint{2.970558in}{3.351896in}}%
\pgfpathlineto{\pgfqpoint{2.970558in}{3.351896in}}%
\pgfpathlineto{\pgfqpoint{2.970558in}{3.354845in}}%
\pgfpathlineto{\pgfqpoint{2.975099in}{3.354845in}}%
\pgfpathlineto{\pgfqpoint{2.975099in}{3.351896in}}%
\pgfpathmoveto{\pgfqpoint{2.975099in}{3.343049in}}%
\pgfpathlineto{\pgfqpoint{2.975099in}{3.343049in}}%
\pgfpathlineto{\pgfqpoint{2.975099in}{3.345998in}}%
\pgfpathlineto{\pgfqpoint{2.979640in}{3.345998in}}%
\pgfpathlineto{\pgfqpoint{2.979640in}{3.343049in}}%
\pgfpathmoveto{\pgfqpoint{2.975099in}{3.345998in}}%
\pgfpathlineto{\pgfqpoint{2.975099in}{3.345998in}}%
\pgfpathlineto{\pgfqpoint{2.975099in}{3.348947in}}%
\pgfpathlineto{\pgfqpoint{2.979640in}{3.348947in}}%
\pgfpathlineto{\pgfqpoint{2.979640in}{3.345998in}}%
\pgfpathmoveto{\pgfqpoint{2.984181in}{3.331252in}}%
\pgfpathlineto{\pgfqpoint{2.984181in}{3.331252in}}%
\pgfpathlineto{\pgfqpoint{2.984181in}{3.334201in}}%
\pgfpathlineto{\pgfqpoint{2.988722in}{3.334201in}}%
\pgfpathlineto{\pgfqpoint{2.988722in}{3.331252in}}%
\pgfpathmoveto{\pgfqpoint{2.984181in}{3.334201in}}%
\pgfpathlineto{\pgfqpoint{2.984181in}{3.334201in}}%
\pgfpathlineto{\pgfqpoint{2.984181in}{3.337150in}}%
\pgfpathlineto{\pgfqpoint{2.988722in}{3.337150in}}%
\pgfpathlineto{\pgfqpoint{2.988722in}{3.334201in}}%
\pgfpathmoveto{\pgfqpoint{2.966017in}{3.354845in}}%
\pgfpathlineto{\pgfqpoint{2.966017in}{3.354845in}}%
\pgfpathlineto{\pgfqpoint{2.966017in}{3.357795in}}%
\pgfpathlineto{\pgfqpoint{2.970558in}{3.357795in}}%
\pgfpathlineto{\pgfqpoint{2.970558in}{3.354845in}}%
\pgfpathmoveto{\pgfqpoint{2.966017in}{3.357795in}}%
\pgfpathlineto{\pgfqpoint{2.966017in}{3.357795in}}%
\pgfpathlineto{\pgfqpoint{2.966017in}{3.360744in}}%
\pgfpathlineto{\pgfqpoint{2.970558in}{3.360744in}}%
\pgfpathlineto{\pgfqpoint{2.970558in}{3.357795in}}%
\pgfpathmoveto{\pgfqpoint{2.938770in}{3.378439in}}%
\pgfpathlineto{\pgfqpoint{2.938770in}{3.378439in}}%
\pgfpathlineto{\pgfqpoint{2.938770in}{3.381388in}}%
\pgfpathlineto{\pgfqpoint{2.943311in}{3.381388in}}%
\pgfpathlineto{\pgfqpoint{2.943311in}{3.378439in}}%
\pgfpathmoveto{\pgfqpoint{2.938770in}{3.381388in}}%
\pgfpathlineto{\pgfqpoint{2.938770in}{3.381388in}}%
\pgfpathlineto{\pgfqpoint{2.938770in}{3.384338in}}%
\pgfpathlineto{\pgfqpoint{2.943311in}{3.384338in}}%
\pgfpathlineto{\pgfqpoint{2.943311in}{3.381388in}}%
\pgfpathmoveto{\pgfqpoint{2.943311in}{3.378439in}}%
\pgfpathlineto{\pgfqpoint{2.943311in}{3.378439in}}%
\pgfpathlineto{\pgfqpoint{2.943311in}{3.381388in}}%
\pgfpathlineto{\pgfqpoint{2.947852in}{3.381388in}}%
\pgfpathlineto{\pgfqpoint{2.947852in}{3.378439in}}%
\pgfpathmoveto{\pgfqpoint{2.943311in}{3.381388in}}%
\pgfpathlineto{\pgfqpoint{2.943311in}{3.381388in}}%
\pgfpathlineto{\pgfqpoint{2.943311in}{3.384338in}}%
\pgfpathlineto{\pgfqpoint{2.947852in}{3.384338in}}%
\pgfpathlineto{\pgfqpoint{2.947852in}{3.381388in}}%
\pgfpathmoveto{\pgfqpoint{2.938770in}{3.384338in}}%
\pgfpathlineto{\pgfqpoint{2.938770in}{3.384338in}}%
\pgfpathlineto{\pgfqpoint{2.938770in}{3.387287in}}%
\pgfpathlineto{\pgfqpoint{2.943311in}{3.387287in}}%
\pgfpathlineto{\pgfqpoint{2.943311in}{3.384338in}}%
\pgfpathmoveto{\pgfqpoint{2.938770in}{3.387287in}}%
\pgfpathlineto{\pgfqpoint{2.938770in}{3.387287in}}%
\pgfpathlineto{\pgfqpoint{2.938770in}{3.390236in}}%
\pgfpathlineto{\pgfqpoint{2.943311in}{3.390236in}}%
\pgfpathlineto{\pgfqpoint{2.943311in}{3.387287in}}%
\pgfpathmoveto{\pgfqpoint{2.943311in}{3.384338in}}%
\pgfpathlineto{\pgfqpoint{2.943311in}{3.384338in}}%
\pgfpathlineto{\pgfqpoint{2.943311in}{3.387287in}}%
\pgfpathlineto{\pgfqpoint{2.947852in}{3.387287in}}%
\pgfpathlineto{\pgfqpoint{2.947852in}{3.384338in}}%
\pgfpathmoveto{\pgfqpoint{2.943311in}{3.387287in}}%
\pgfpathlineto{\pgfqpoint{2.943311in}{3.387287in}}%
\pgfpathlineto{\pgfqpoint{2.943311in}{3.390236in}}%
\pgfpathlineto{\pgfqpoint{2.947852in}{3.390236in}}%
\pgfpathlineto{\pgfqpoint{2.947852in}{3.387287in}}%
\pgfpathmoveto{\pgfqpoint{2.929688in}{3.390236in}}%
\pgfpathlineto{\pgfqpoint{2.929688in}{3.390236in}}%
\pgfpathlineto{\pgfqpoint{2.929688in}{3.393185in}}%
\pgfpathlineto{\pgfqpoint{2.934229in}{3.393185in}}%
\pgfpathlineto{\pgfqpoint{2.934229in}{3.390236in}}%
\pgfpathmoveto{\pgfqpoint{2.929688in}{3.393185in}}%
\pgfpathlineto{\pgfqpoint{2.929688in}{3.393185in}}%
\pgfpathlineto{\pgfqpoint{2.929688in}{3.396134in}}%
\pgfpathlineto{\pgfqpoint{2.934229in}{3.396134in}}%
\pgfpathlineto{\pgfqpoint{2.934229in}{3.393185in}}%
\pgfpathmoveto{\pgfqpoint{2.934229in}{3.390236in}}%
\pgfpathlineto{\pgfqpoint{2.934229in}{3.390236in}}%
\pgfpathlineto{\pgfqpoint{2.934229in}{3.393185in}}%
\pgfpathlineto{\pgfqpoint{2.938770in}{3.393185in}}%
\pgfpathlineto{\pgfqpoint{2.938770in}{3.390236in}}%
\pgfpathmoveto{\pgfqpoint{2.934229in}{3.393185in}}%
\pgfpathlineto{\pgfqpoint{2.934229in}{3.393185in}}%
\pgfpathlineto{\pgfqpoint{2.934229in}{3.396134in}}%
\pgfpathlineto{\pgfqpoint{2.938770in}{3.396134in}}%
\pgfpathlineto{\pgfqpoint{2.938770in}{3.393185in}}%
\pgfpathmoveto{\pgfqpoint{2.929688in}{3.396134in}}%
\pgfpathlineto{\pgfqpoint{2.929688in}{3.396134in}}%
\pgfpathlineto{\pgfqpoint{2.929688in}{3.399084in}}%
\pgfpathlineto{\pgfqpoint{2.934229in}{3.399084in}}%
\pgfpathlineto{\pgfqpoint{2.934229in}{3.396134in}}%
\pgfpathmoveto{\pgfqpoint{2.929688in}{3.399084in}}%
\pgfpathlineto{\pgfqpoint{2.929688in}{3.399084in}}%
\pgfpathlineto{\pgfqpoint{2.929688in}{3.402033in}}%
\pgfpathlineto{\pgfqpoint{2.934229in}{3.402033in}}%
\pgfpathlineto{\pgfqpoint{2.934229in}{3.399084in}}%
\pgfpathmoveto{\pgfqpoint{2.934229in}{3.396134in}}%
\pgfpathlineto{\pgfqpoint{2.934229in}{3.396134in}}%
\pgfpathlineto{\pgfqpoint{2.934229in}{3.399084in}}%
\pgfpathlineto{\pgfqpoint{2.938770in}{3.399084in}}%
\pgfpathlineto{\pgfqpoint{2.938770in}{3.396134in}}%
\pgfpathmoveto{\pgfqpoint{2.934229in}{3.399084in}}%
\pgfpathlineto{\pgfqpoint{2.934229in}{3.399084in}}%
\pgfpathlineto{\pgfqpoint{2.934229in}{3.402033in}}%
\pgfpathlineto{\pgfqpoint{2.938770in}{3.402033in}}%
\pgfpathlineto{\pgfqpoint{2.938770in}{3.399084in}}%
\pgfpathmoveto{\pgfqpoint{2.938770in}{3.390236in}}%
\pgfpathlineto{\pgfqpoint{2.938770in}{3.390236in}}%
\pgfpathlineto{\pgfqpoint{2.938770in}{3.393185in}}%
\pgfpathlineto{\pgfqpoint{2.943311in}{3.393185in}}%
\pgfpathlineto{\pgfqpoint{2.943311in}{3.390236in}}%
\pgfpathmoveto{\pgfqpoint{2.938770in}{3.393185in}}%
\pgfpathlineto{\pgfqpoint{2.938770in}{3.393185in}}%
\pgfpathlineto{\pgfqpoint{2.938770in}{3.396134in}}%
\pgfpathlineto{\pgfqpoint{2.943311in}{3.396134in}}%
\pgfpathlineto{\pgfqpoint{2.943311in}{3.393185in}}%
\pgfpathmoveto{\pgfqpoint{2.947852in}{3.378439in}}%
\pgfpathlineto{\pgfqpoint{2.947852in}{3.378439in}}%
\pgfpathlineto{\pgfqpoint{2.947852in}{3.381388in}}%
\pgfpathlineto{\pgfqpoint{2.952394in}{3.381388in}}%
\pgfpathlineto{\pgfqpoint{2.952394in}{3.378439in}}%
\pgfpathmoveto{\pgfqpoint{2.947852in}{3.381388in}}%
\pgfpathlineto{\pgfqpoint{2.947852in}{3.381388in}}%
\pgfpathlineto{\pgfqpoint{2.947852in}{3.384338in}}%
\pgfpathlineto{\pgfqpoint{2.952394in}{3.384338in}}%
\pgfpathlineto{\pgfqpoint{2.952394in}{3.381388in}}%
\pgfpathmoveto{\pgfqpoint{2.929688in}{3.402033in}}%
\pgfpathlineto{\pgfqpoint{2.929688in}{3.402033in}}%
\pgfpathlineto{\pgfqpoint{2.929688in}{3.404982in}}%
\pgfpathlineto{\pgfqpoint{2.934229in}{3.404982in}}%
\pgfpathlineto{\pgfqpoint{2.934229in}{3.402033in}}%
\pgfpathmoveto{\pgfqpoint{2.929688in}{3.404982in}}%
\pgfpathlineto{\pgfqpoint{2.929688in}{3.404982in}}%
\pgfpathlineto{\pgfqpoint{2.929688in}{3.407931in}}%
\pgfpathlineto{\pgfqpoint{2.934229in}{3.407931in}}%
\pgfpathlineto{\pgfqpoint{2.934229in}{3.404982in}}%
\pgfpathmoveto{\pgfqpoint{3.075004in}{2.010000in}}%
\pgfpathlineto{\pgfqpoint{3.075004in}{2.010000in}}%
\pgfpathlineto{\pgfqpoint{3.075004in}{2.012949in}}%
\pgfpathlineto{\pgfqpoint{3.079545in}{2.012949in}}%
\pgfpathlineto{\pgfqpoint{3.079545in}{2.010000in}}%
\pgfpathmoveto{\pgfqpoint{3.075004in}{2.012949in}}%
\pgfpathlineto{\pgfqpoint{3.075004in}{2.012949in}}%
\pgfpathlineto{\pgfqpoint{3.075004in}{2.015899in}}%
\pgfpathlineto{\pgfqpoint{3.079545in}{2.015899in}}%
\pgfpathlineto{\pgfqpoint{3.079545in}{2.012949in}}%
\pgfpathmoveto{\pgfqpoint{3.079545in}{2.012949in}}%
\pgfpathlineto{\pgfqpoint{3.079545in}{2.012949in}}%
\pgfpathlineto{\pgfqpoint{3.079545in}{2.015899in}}%
\pgfpathlineto{\pgfqpoint{3.084085in}{2.015899in}}%
\pgfpathlineto{\pgfqpoint{3.084085in}{2.012949in}}%
\pgfpathmoveto{\pgfqpoint{3.075004in}{2.015899in}}%
\pgfpathlineto{\pgfqpoint{3.075004in}{2.015899in}}%
\pgfpathlineto{\pgfqpoint{3.075004in}{2.018848in}}%
\pgfpathlineto{\pgfqpoint{3.079545in}{2.018848in}}%
\pgfpathlineto{\pgfqpoint{3.079545in}{2.015899in}}%
\pgfpathmoveto{\pgfqpoint{3.075004in}{2.018848in}}%
\pgfpathlineto{\pgfqpoint{3.075004in}{2.018848in}}%
\pgfpathlineto{\pgfqpoint{3.075004in}{2.021797in}}%
\pgfpathlineto{\pgfqpoint{3.079545in}{2.021797in}}%
\pgfpathlineto{\pgfqpoint{3.079545in}{2.018848in}}%
\pgfpathmoveto{\pgfqpoint{3.079545in}{2.015899in}}%
\pgfpathlineto{\pgfqpoint{3.079545in}{2.015899in}}%
\pgfpathlineto{\pgfqpoint{3.079545in}{2.018848in}}%
\pgfpathlineto{\pgfqpoint{3.084085in}{2.018848in}}%
\pgfpathlineto{\pgfqpoint{3.084085in}{2.015899in}}%
\pgfpathmoveto{\pgfqpoint{3.079545in}{2.018848in}}%
\pgfpathlineto{\pgfqpoint{3.079545in}{2.018848in}}%
\pgfpathlineto{\pgfqpoint{3.079545in}{2.021797in}}%
\pgfpathlineto{\pgfqpoint{3.084085in}{2.021797in}}%
\pgfpathlineto{\pgfqpoint{3.084085in}{2.018848in}}%
\pgfpathmoveto{\pgfqpoint{3.084085in}{2.015899in}}%
\pgfpathlineto{\pgfqpoint{3.084085in}{2.015899in}}%
\pgfpathlineto{\pgfqpoint{3.084085in}{2.018848in}}%
\pgfpathlineto{\pgfqpoint{3.088626in}{2.018848in}}%
\pgfpathlineto{\pgfqpoint{3.088626in}{2.015899in}}%
\pgfpathmoveto{\pgfqpoint{3.084085in}{2.018848in}}%
\pgfpathlineto{\pgfqpoint{3.084085in}{2.018848in}}%
\pgfpathlineto{\pgfqpoint{3.084085in}{2.021797in}}%
\pgfpathlineto{\pgfqpoint{3.088626in}{2.021797in}}%
\pgfpathlineto{\pgfqpoint{3.088626in}{2.018848in}}%
\pgfpathmoveto{\pgfqpoint{3.088626in}{2.018848in}}%
\pgfpathlineto{\pgfqpoint{3.088626in}{2.018848in}}%
\pgfpathlineto{\pgfqpoint{3.088626in}{2.021797in}}%
\pgfpathlineto{\pgfqpoint{3.093167in}{2.021797in}}%
\pgfpathlineto{\pgfqpoint{3.093167in}{2.018848in}}%
\pgfpathmoveto{\pgfqpoint{3.084085in}{2.021797in}}%
\pgfpathlineto{\pgfqpoint{3.084085in}{2.021797in}}%
\pgfpathlineto{\pgfqpoint{3.084085in}{2.024747in}}%
\pgfpathlineto{\pgfqpoint{3.088626in}{2.024747in}}%
\pgfpathlineto{\pgfqpoint{3.088626in}{2.021797in}}%
\pgfpathmoveto{\pgfqpoint{3.084085in}{2.024747in}}%
\pgfpathlineto{\pgfqpoint{3.084085in}{2.024747in}}%
\pgfpathlineto{\pgfqpoint{3.084085in}{2.027696in}}%
\pgfpathlineto{\pgfqpoint{3.088626in}{2.027696in}}%
\pgfpathlineto{\pgfqpoint{3.088626in}{2.024747in}}%
\pgfpathmoveto{\pgfqpoint{3.088626in}{2.021797in}}%
\pgfpathlineto{\pgfqpoint{3.088626in}{2.021797in}}%
\pgfpathlineto{\pgfqpoint{3.088626in}{2.024747in}}%
\pgfpathlineto{\pgfqpoint{3.093167in}{2.024747in}}%
\pgfpathlineto{\pgfqpoint{3.093167in}{2.021797in}}%
\pgfpathmoveto{\pgfqpoint{3.088626in}{2.024747in}}%
\pgfpathlineto{\pgfqpoint{3.088626in}{2.024747in}}%
\pgfpathlineto{\pgfqpoint{3.088626in}{2.027696in}}%
\pgfpathlineto{\pgfqpoint{3.093167in}{2.027696in}}%
\pgfpathlineto{\pgfqpoint{3.093167in}{2.024747in}}%
\pgfpathmoveto{\pgfqpoint{3.093167in}{2.021797in}}%
\pgfpathlineto{\pgfqpoint{3.093167in}{2.021797in}}%
\pgfpathlineto{\pgfqpoint{3.093167in}{2.024747in}}%
\pgfpathlineto{\pgfqpoint{3.097708in}{2.024747in}}%
\pgfpathlineto{\pgfqpoint{3.097708in}{2.021797in}}%
\pgfpathmoveto{\pgfqpoint{3.093167in}{2.024747in}}%
\pgfpathlineto{\pgfqpoint{3.093167in}{2.024747in}}%
\pgfpathlineto{\pgfqpoint{3.093167in}{2.027696in}}%
\pgfpathlineto{\pgfqpoint{3.097708in}{2.027696in}}%
\pgfpathlineto{\pgfqpoint{3.097708in}{2.024747in}}%
\pgfpathmoveto{\pgfqpoint{3.097708in}{2.024747in}}%
\pgfpathlineto{\pgfqpoint{3.097708in}{2.024747in}}%
\pgfpathlineto{\pgfqpoint{3.097708in}{2.027696in}}%
\pgfpathlineto{\pgfqpoint{3.102249in}{2.027696in}}%
\pgfpathlineto{\pgfqpoint{3.102249in}{2.024747in}}%
\pgfpathmoveto{\pgfqpoint{3.093167in}{2.027696in}}%
\pgfpathlineto{\pgfqpoint{3.093167in}{2.027696in}}%
\pgfpathlineto{\pgfqpoint{3.093167in}{2.030645in}}%
\pgfpathlineto{\pgfqpoint{3.097708in}{2.030645in}}%
\pgfpathlineto{\pgfqpoint{3.097708in}{2.027696in}}%
\pgfpathmoveto{\pgfqpoint{3.093167in}{2.030645in}}%
\pgfpathlineto{\pgfqpoint{3.093167in}{2.030645in}}%
\pgfpathlineto{\pgfqpoint{3.093167in}{2.033594in}}%
\pgfpathlineto{\pgfqpoint{3.097708in}{2.033594in}}%
\pgfpathlineto{\pgfqpoint{3.097708in}{2.030645in}}%
\pgfpathmoveto{\pgfqpoint{3.097708in}{2.027696in}}%
\pgfpathlineto{\pgfqpoint{3.097708in}{2.027696in}}%
\pgfpathlineto{\pgfqpoint{3.097708in}{2.030645in}}%
\pgfpathlineto{\pgfqpoint{3.102249in}{2.030645in}}%
\pgfpathlineto{\pgfqpoint{3.102249in}{2.027696in}}%
\pgfpathmoveto{\pgfqpoint{3.097708in}{2.030645in}}%
\pgfpathlineto{\pgfqpoint{3.097708in}{2.030645in}}%
\pgfpathlineto{\pgfqpoint{3.097708in}{2.033594in}}%
\pgfpathlineto{\pgfqpoint{3.102249in}{2.033594in}}%
\pgfpathlineto{\pgfqpoint{3.102249in}{2.030645in}}%
\pgfpathmoveto{\pgfqpoint{3.102249in}{2.027696in}}%
\pgfpathlineto{\pgfqpoint{3.102249in}{2.027696in}}%
\pgfpathlineto{\pgfqpoint{3.102249in}{2.030645in}}%
\pgfpathlineto{\pgfqpoint{3.106789in}{2.030645in}}%
\pgfpathlineto{\pgfqpoint{3.106789in}{2.027696in}}%
\pgfpathmoveto{\pgfqpoint{3.102249in}{2.030645in}}%
\pgfpathlineto{\pgfqpoint{3.102249in}{2.030645in}}%
\pgfpathlineto{\pgfqpoint{3.102249in}{2.033594in}}%
\pgfpathlineto{\pgfqpoint{3.106789in}{2.033594in}}%
\pgfpathlineto{\pgfqpoint{3.106789in}{2.030645in}}%
\pgfpathmoveto{\pgfqpoint{3.106789in}{2.030645in}}%
\pgfpathlineto{\pgfqpoint{3.106789in}{2.030645in}}%
\pgfpathlineto{\pgfqpoint{3.106789in}{2.033594in}}%
\pgfpathlineto{\pgfqpoint{3.111330in}{2.033594in}}%
\pgfpathlineto{\pgfqpoint{3.111330in}{2.030645in}}%
\pgfpathmoveto{\pgfqpoint{3.102249in}{2.033594in}}%
\pgfpathlineto{\pgfqpoint{3.102249in}{2.033594in}}%
\pgfpathlineto{\pgfqpoint{3.102249in}{2.036544in}}%
\pgfpathlineto{\pgfqpoint{3.106789in}{2.036544in}}%
\pgfpathlineto{\pgfqpoint{3.106789in}{2.033594in}}%
\pgfpathmoveto{\pgfqpoint{3.102249in}{2.036544in}}%
\pgfpathlineto{\pgfqpoint{3.102249in}{2.036544in}}%
\pgfpathlineto{\pgfqpoint{3.102249in}{2.039493in}}%
\pgfpathlineto{\pgfqpoint{3.106789in}{2.039493in}}%
\pgfpathlineto{\pgfqpoint{3.106789in}{2.036544in}}%
\pgfpathmoveto{\pgfqpoint{3.106789in}{2.033594in}}%
\pgfpathlineto{\pgfqpoint{3.106789in}{2.033594in}}%
\pgfpathlineto{\pgfqpoint{3.106789in}{2.036544in}}%
\pgfpathlineto{\pgfqpoint{3.111330in}{2.036544in}}%
\pgfpathlineto{\pgfqpoint{3.111330in}{2.033594in}}%
\pgfpathmoveto{\pgfqpoint{3.106789in}{2.036544in}}%
\pgfpathlineto{\pgfqpoint{3.106789in}{2.036544in}}%
\pgfpathlineto{\pgfqpoint{3.106789in}{2.039493in}}%
\pgfpathlineto{\pgfqpoint{3.111330in}{2.039493in}}%
\pgfpathlineto{\pgfqpoint{3.111330in}{2.036544in}}%
\pgfpathmoveto{\pgfqpoint{3.111330in}{2.033594in}}%
\pgfpathlineto{\pgfqpoint{3.111330in}{2.033594in}}%
\pgfpathlineto{\pgfqpoint{3.111330in}{2.036544in}}%
\pgfpathlineto{\pgfqpoint{3.115871in}{2.036544in}}%
\pgfpathlineto{\pgfqpoint{3.115871in}{2.033594in}}%
\pgfpathmoveto{\pgfqpoint{3.111330in}{2.036544in}}%
\pgfpathlineto{\pgfqpoint{3.111330in}{2.036544in}}%
\pgfpathlineto{\pgfqpoint{3.111330in}{2.039493in}}%
\pgfpathlineto{\pgfqpoint{3.115871in}{2.039493in}}%
\pgfpathlineto{\pgfqpoint{3.115871in}{2.036544in}}%
\pgfpathmoveto{\pgfqpoint{3.115871in}{2.036544in}}%
\pgfpathlineto{\pgfqpoint{3.115871in}{2.036544in}}%
\pgfpathlineto{\pgfqpoint{3.115871in}{2.039493in}}%
\pgfpathlineto{\pgfqpoint{3.120412in}{2.039493in}}%
\pgfpathlineto{\pgfqpoint{3.120412in}{2.036544in}}%
\pgfpathmoveto{\pgfqpoint{3.111330in}{2.039493in}}%
\pgfpathlineto{\pgfqpoint{3.111330in}{2.039493in}}%
\pgfpathlineto{\pgfqpoint{3.111330in}{2.042442in}}%
\pgfpathlineto{\pgfqpoint{3.115871in}{2.042442in}}%
\pgfpathlineto{\pgfqpoint{3.115871in}{2.039493in}}%
\pgfpathmoveto{\pgfqpoint{3.111330in}{2.042442in}}%
\pgfpathlineto{\pgfqpoint{3.111330in}{2.042442in}}%
\pgfpathlineto{\pgfqpoint{3.111330in}{2.045392in}}%
\pgfpathlineto{\pgfqpoint{3.115871in}{2.045392in}}%
\pgfpathlineto{\pgfqpoint{3.115871in}{2.042442in}}%
\pgfpathmoveto{\pgfqpoint{3.115871in}{2.039493in}}%
\pgfpathlineto{\pgfqpoint{3.115871in}{2.039493in}}%
\pgfpathlineto{\pgfqpoint{3.115871in}{2.042442in}}%
\pgfpathlineto{\pgfqpoint{3.120412in}{2.042442in}}%
\pgfpathlineto{\pgfqpoint{3.120412in}{2.039493in}}%
\pgfpathmoveto{\pgfqpoint{3.115871in}{2.042442in}}%
\pgfpathlineto{\pgfqpoint{3.115871in}{2.042442in}}%
\pgfpathlineto{\pgfqpoint{3.115871in}{2.045392in}}%
\pgfpathlineto{\pgfqpoint{3.120412in}{2.045392in}}%
\pgfpathlineto{\pgfqpoint{3.120412in}{2.042442in}}%
\pgfpathmoveto{\pgfqpoint{3.120412in}{2.039493in}}%
\pgfpathlineto{\pgfqpoint{3.120412in}{2.039493in}}%
\pgfpathlineto{\pgfqpoint{3.120412in}{2.042442in}}%
\pgfpathlineto{\pgfqpoint{3.124953in}{2.042442in}}%
\pgfpathlineto{\pgfqpoint{3.124953in}{2.039493in}}%
\pgfpathmoveto{\pgfqpoint{3.120412in}{2.042442in}}%
\pgfpathlineto{\pgfqpoint{3.120412in}{2.042442in}}%
\pgfpathlineto{\pgfqpoint{3.120412in}{2.045392in}}%
\pgfpathlineto{\pgfqpoint{3.124953in}{2.045392in}}%
\pgfpathlineto{\pgfqpoint{3.124953in}{2.042442in}}%
\pgfpathmoveto{\pgfqpoint{3.124953in}{2.042442in}}%
\pgfpathlineto{\pgfqpoint{3.124953in}{2.042442in}}%
\pgfpathlineto{\pgfqpoint{3.124953in}{2.045392in}}%
\pgfpathlineto{\pgfqpoint{3.129493in}{2.045392in}}%
\pgfpathlineto{\pgfqpoint{3.129493in}{2.042442in}}%
\pgfpathmoveto{\pgfqpoint{3.129493in}{2.042442in}}%
\pgfpathlineto{\pgfqpoint{3.129493in}{2.042442in}}%
\pgfpathlineto{\pgfqpoint{3.129493in}{2.045392in}}%
\pgfpathlineto{\pgfqpoint{3.134034in}{2.045392in}}%
\pgfpathlineto{\pgfqpoint{3.134034in}{2.042442in}}%
\pgfpathmoveto{\pgfqpoint{3.129493in}{2.045392in}}%
\pgfpathlineto{\pgfqpoint{3.129493in}{2.045392in}}%
\pgfpathlineto{\pgfqpoint{3.129493in}{2.048341in}}%
\pgfpathlineto{\pgfqpoint{3.134034in}{2.048341in}}%
\pgfpathlineto{\pgfqpoint{3.134034in}{2.045392in}}%
\pgfpathmoveto{\pgfqpoint{3.129493in}{2.048341in}}%
\pgfpathlineto{\pgfqpoint{3.129493in}{2.048341in}}%
\pgfpathlineto{\pgfqpoint{3.129493in}{2.051290in}}%
\pgfpathlineto{\pgfqpoint{3.134034in}{2.051290in}}%
\pgfpathlineto{\pgfqpoint{3.134034in}{2.048341in}}%
\pgfpathmoveto{\pgfqpoint{3.134034in}{2.045392in}}%
\pgfpathlineto{\pgfqpoint{3.134034in}{2.045392in}}%
\pgfpathlineto{\pgfqpoint{3.134034in}{2.048341in}}%
\pgfpathlineto{\pgfqpoint{3.138575in}{2.048341in}}%
\pgfpathlineto{\pgfqpoint{3.138575in}{2.045392in}}%
\pgfpathmoveto{\pgfqpoint{3.134034in}{2.048341in}}%
\pgfpathlineto{\pgfqpoint{3.134034in}{2.048341in}}%
\pgfpathlineto{\pgfqpoint{3.134034in}{2.051290in}}%
\pgfpathlineto{\pgfqpoint{3.138575in}{2.051290in}}%
\pgfpathlineto{\pgfqpoint{3.138575in}{2.048341in}}%
\pgfpathmoveto{\pgfqpoint{3.138575in}{2.048341in}}%
\pgfpathlineto{\pgfqpoint{3.138575in}{2.048341in}}%
\pgfpathlineto{\pgfqpoint{3.138575in}{2.051290in}}%
\pgfpathlineto{\pgfqpoint{3.143116in}{2.051290in}}%
\pgfpathlineto{\pgfqpoint{3.143116in}{2.048341in}}%
\pgfpathmoveto{\pgfqpoint{3.138575in}{2.051290in}}%
\pgfpathlineto{\pgfqpoint{3.138575in}{2.051290in}}%
\pgfpathlineto{\pgfqpoint{3.138575in}{2.054239in}}%
\pgfpathlineto{\pgfqpoint{3.143116in}{2.054239in}}%
\pgfpathlineto{\pgfqpoint{3.143116in}{2.051290in}}%
\pgfpathmoveto{\pgfqpoint{3.138575in}{2.054239in}}%
\pgfpathlineto{\pgfqpoint{3.138575in}{2.054239in}}%
\pgfpathlineto{\pgfqpoint{3.138575in}{2.057189in}}%
\pgfpathlineto{\pgfqpoint{3.143116in}{2.057189in}}%
\pgfpathlineto{\pgfqpoint{3.143116in}{2.054239in}}%
\pgfpathmoveto{\pgfqpoint{3.143116in}{2.051290in}}%
\pgfpathlineto{\pgfqpoint{3.143116in}{2.051290in}}%
\pgfpathlineto{\pgfqpoint{3.143116in}{2.054239in}}%
\pgfpathlineto{\pgfqpoint{3.147657in}{2.054239in}}%
\pgfpathlineto{\pgfqpoint{3.147657in}{2.051290in}}%
\pgfpathmoveto{\pgfqpoint{3.143116in}{2.054239in}}%
\pgfpathlineto{\pgfqpoint{3.143116in}{2.054239in}}%
\pgfpathlineto{\pgfqpoint{3.143116in}{2.057189in}}%
\pgfpathlineto{\pgfqpoint{3.147657in}{2.057189in}}%
\pgfpathlineto{\pgfqpoint{3.147657in}{2.054239in}}%
\pgfpathmoveto{\pgfqpoint{3.147657in}{2.054239in}}%
\pgfpathlineto{\pgfqpoint{3.147657in}{2.054239in}}%
\pgfpathlineto{\pgfqpoint{3.147657in}{2.057189in}}%
\pgfpathlineto{\pgfqpoint{3.152197in}{2.057189in}}%
\pgfpathlineto{\pgfqpoint{3.152197in}{2.054239in}}%
\pgfpathmoveto{\pgfqpoint{3.147657in}{2.057189in}}%
\pgfpathlineto{\pgfqpoint{3.147657in}{2.057189in}}%
\pgfpathlineto{\pgfqpoint{3.147657in}{2.060138in}}%
\pgfpathlineto{\pgfqpoint{3.152197in}{2.060138in}}%
\pgfpathlineto{\pgfqpoint{3.152197in}{2.057189in}}%
\pgfpathmoveto{\pgfqpoint{3.147657in}{2.060138in}}%
\pgfpathlineto{\pgfqpoint{3.147657in}{2.060138in}}%
\pgfpathlineto{\pgfqpoint{3.147657in}{2.063087in}}%
\pgfpathlineto{\pgfqpoint{3.152197in}{2.063087in}}%
\pgfpathlineto{\pgfqpoint{3.152197in}{2.060138in}}%
\pgfpathmoveto{\pgfqpoint{3.152197in}{2.057189in}}%
\pgfpathlineto{\pgfqpoint{3.152197in}{2.057189in}}%
\pgfpathlineto{\pgfqpoint{3.152197in}{2.060138in}}%
\pgfpathlineto{\pgfqpoint{3.156738in}{2.060138in}}%
\pgfpathlineto{\pgfqpoint{3.156738in}{2.057189in}}%
\pgfpathmoveto{\pgfqpoint{3.152197in}{2.060138in}}%
\pgfpathlineto{\pgfqpoint{3.152197in}{2.060138in}}%
\pgfpathlineto{\pgfqpoint{3.152197in}{2.063087in}}%
\pgfpathlineto{\pgfqpoint{3.156738in}{2.063087in}}%
\pgfpathlineto{\pgfqpoint{3.156738in}{2.060138in}}%
\pgfpathmoveto{\pgfqpoint{3.156738in}{2.060138in}}%
\pgfpathlineto{\pgfqpoint{3.156738in}{2.060138in}}%
\pgfpathlineto{\pgfqpoint{3.156738in}{2.063087in}}%
\pgfpathlineto{\pgfqpoint{3.161279in}{2.063087in}}%
\pgfpathlineto{\pgfqpoint{3.161279in}{2.060138in}}%
\pgfpathmoveto{\pgfqpoint{3.156738in}{2.063087in}}%
\pgfpathlineto{\pgfqpoint{3.156738in}{2.063087in}}%
\pgfpathlineto{\pgfqpoint{3.156738in}{2.066037in}}%
\pgfpathlineto{\pgfqpoint{3.161279in}{2.066037in}}%
\pgfpathlineto{\pgfqpoint{3.161279in}{2.063087in}}%
\pgfpathmoveto{\pgfqpoint{3.156738in}{2.066037in}}%
\pgfpathlineto{\pgfqpoint{3.156738in}{2.066037in}}%
\pgfpathlineto{\pgfqpoint{3.156738in}{2.068986in}}%
\pgfpathlineto{\pgfqpoint{3.161279in}{2.068986in}}%
\pgfpathlineto{\pgfqpoint{3.161279in}{2.066037in}}%
\pgfpathmoveto{\pgfqpoint{3.161279in}{2.063087in}}%
\pgfpathlineto{\pgfqpoint{3.161279in}{2.063087in}}%
\pgfpathlineto{\pgfqpoint{3.161279in}{2.066037in}}%
\pgfpathlineto{\pgfqpoint{3.165820in}{2.066037in}}%
\pgfpathlineto{\pgfqpoint{3.165820in}{2.063087in}}%
\pgfpathmoveto{\pgfqpoint{3.161279in}{2.066037in}}%
\pgfpathlineto{\pgfqpoint{3.161279in}{2.066037in}}%
\pgfpathlineto{\pgfqpoint{3.161279in}{2.068986in}}%
\pgfpathlineto{\pgfqpoint{3.165820in}{2.068986in}}%
\pgfpathlineto{\pgfqpoint{3.165820in}{2.066037in}}%
\pgfpathmoveto{\pgfqpoint{3.165820in}{2.066037in}}%
\pgfpathlineto{\pgfqpoint{3.165820in}{2.066037in}}%
\pgfpathlineto{\pgfqpoint{3.165820in}{2.068986in}}%
\pgfpathlineto{\pgfqpoint{3.170361in}{2.068986in}}%
\pgfpathlineto{\pgfqpoint{3.170361in}{2.066037in}}%
\pgfpathmoveto{\pgfqpoint{3.165820in}{2.068986in}}%
\pgfpathlineto{\pgfqpoint{3.165820in}{2.068986in}}%
\pgfpathlineto{\pgfqpoint{3.165820in}{2.071935in}}%
\pgfpathlineto{\pgfqpoint{3.170361in}{2.071935in}}%
\pgfpathlineto{\pgfqpoint{3.170361in}{2.068986in}}%
\pgfpathmoveto{\pgfqpoint{3.165820in}{2.071935in}}%
\pgfpathlineto{\pgfqpoint{3.165820in}{2.071935in}}%
\pgfpathlineto{\pgfqpoint{3.165820in}{2.074884in}}%
\pgfpathlineto{\pgfqpoint{3.170361in}{2.074884in}}%
\pgfpathlineto{\pgfqpoint{3.170361in}{2.071935in}}%
\pgfpathmoveto{\pgfqpoint{3.170361in}{2.068986in}}%
\pgfpathlineto{\pgfqpoint{3.170361in}{2.068986in}}%
\pgfpathlineto{\pgfqpoint{3.170361in}{2.071935in}}%
\pgfpathlineto{\pgfqpoint{3.174901in}{2.071935in}}%
\pgfpathlineto{\pgfqpoint{3.174901in}{2.068986in}}%
\pgfpathmoveto{\pgfqpoint{3.170361in}{2.071935in}}%
\pgfpathlineto{\pgfqpoint{3.170361in}{2.071935in}}%
\pgfpathlineto{\pgfqpoint{3.170361in}{2.074884in}}%
\pgfpathlineto{\pgfqpoint{3.174901in}{2.074884in}}%
\pgfpathlineto{\pgfqpoint{3.174901in}{2.071935in}}%
\pgfpathmoveto{\pgfqpoint{3.174901in}{2.071935in}}%
\pgfpathlineto{\pgfqpoint{3.174901in}{2.071935in}}%
\pgfpathlineto{\pgfqpoint{3.174901in}{2.074884in}}%
\pgfpathlineto{\pgfqpoint{3.179442in}{2.074884in}}%
\pgfpathlineto{\pgfqpoint{3.179442in}{2.071935in}}%
\pgfpathmoveto{\pgfqpoint{3.174901in}{2.074884in}}%
\pgfpathlineto{\pgfqpoint{3.174901in}{2.074884in}}%
\pgfpathlineto{\pgfqpoint{3.174901in}{2.077834in}}%
\pgfpathlineto{\pgfqpoint{3.179442in}{2.077834in}}%
\pgfpathlineto{\pgfqpoint{3.179442in}{2.074884in}}%
\pgfpathmoveto{\pgfqpoint{3.174901in}{2.077834in}}%
\pgfpathlineto{\pgfqpoint{3.174901in}{2.077834in}}%
\pgfpathlineto{\pgfqpoint{3.174901in}{2.080783in}}%
\pgfpathlineto{\pgfqpoint{3.179442in}{2.080783in}}%
\pgfpathlineto{\pgfqpoint{3.179442in}{2.077834in}}%
\pgfpathmoveto{\pgfqpoint{3.179442in}{2.074884in}}%
\pgfpathlineto{\pgfqpoint{3.179442in}{2.074884in}}%
\pgfpathlineto{\pgfqpoint{3.179442in}{2.077834in}}%
\pgfpathlineto{\pgfqpoint{3.183983in}{2.077834in}}%
\pgfpathlineto{\pgfqpoint{3.183983in}{2.074884in}}%
\pgfpathmoveto{\pgfqpoint{3.179442in}{2.077834in}}%
\pgfpathlineto{\pgfqpoint{3.179442in}{2.077834in}}%
\pgfpathlineto{\pgfqpoint{3.179442in}{2.080783in}}%
\pgfpathlineto{\pgfqpoint{3.183983in}{2.080783in}}%
\pgfpathlineto{\pgfqpoint{3.183983in}{2.077834in}}%
\pgfpathmoveto{\pgfqpoint{3.183983in}{2.077834in}}%
\pgfpathlineto{\pgfqpoint{3.183983in}{2.077834in}}%
\pgfpathlineto{\pgfqpoint{3.183983in}{2.080783in}}%
\pgfpathlineto{\pgfqpoint{3.188524in}{2.080783in}}%
\pgfpathlineto{\pgfqpoint{3.188524in}{2.077834in}}%
\pgfpathmoveto{\pgfqpoint{3.183983in}{2.080783in}}%
\pgfpathlineto{\pgfqpoint{3.183983in}{2.080783in}}%
\pgfpathlineto{\pgfqpoint{3.183983in}{2.083732in}}%
\pgfpathlineto{\pgfqpoint{3.188524in}{2.083732in}}%
\pgfpathlineto{\pgfqpoint{3.188524in}{2.080783in}}%
\pgfpathmoveto{\pgfqpoint{3.183983in}{2.083732in}}%
\pgfpathlineto{\pgfqpoint{3.183983in}{2.083732in}}%
\pgfpathlineto{\pgfqpoint{3.183983in}{2.086682in}}%
\pgfpathlineto{\pgfqpoint{3.188524in}{2.086682in}}%
\pgfpathlineto{\pgfqpoint{3.188524in}{2.083732in}}%
\pgfpathmoveto{\pgfqpoint{3.188524in}{2.080783in}}%
\pgfpathlineto{\pgfqpoint{3.188524in}{2.080783in}}%
\pgfpathlineto{\pgfqpoint{3.188524in}{2.083732in}}%
\pgfpathlineto{\pgfqpoint{3.193065in}{2.083732in}}%
\pgfpathlineto{\pgfqpoint{3.193065in}{2.080783in}}%
\pgfpathmoveto{\pgfqpoint{3.188524in}{2.083732in}}%
\pgfpathlineto{\pgfqpoint{3.188524in}{2.083732in}}%
\pgfpathlineto{\pgfqpoint{3.188524in}{2.086682in}}%
\pgfpathlineto{\pgfqpoint{3.193065in}{2.086682in}}%
\pgfpathlineto{\pgfqpoint{3.193065in}{2.083732in}}%
\pgfpathmoveto{\pgfqpoint{3.193065in}{2.083732in}}%
\pgfpathlineto{\pgfqpoint{3.193065in}{2.083732in}}%
\pgfpathlineto{\pgfqpoint{3.193065in}{2.086682in}}%
\pgfpathlineto{\pgfqpoint{3.197606in}{2.086682in}}%
\pgfpathlineto{\pgfqpoint{3.197606in}{2.083732in}}%
\pgfpathmoveto{\pgfqpoint{3.193065in}{2.086682in}}%
\pgfpathlineto{\pgfqpoint{3.193065in}{2.086682in}}%
\pgfpathlineto{\pgfqpoint{3.193065in}{2.089631in}}%
\pgfpathlineto{\pgfqpoint{3.197606in}{2.089631in}}%
\pgfpathlineto{\pgfqpoint{3.197606in}{2.086682in}}%
\pgfpathmoveto{\pgfqpoint{3.193065in}{2.089631in}}%
\pgfpathlineto{\pgfqpoint{3.193065in}{2.089631in}}%
\pgfpathlineto{\pgfqpoint{3.193065in}{2.092580in}}%
\pgfpathlineto{\pgfqpoint{3.197606in}{2.092580in}}%
\pgfpathlineto{\pgfqpoint{3.197606in}{2.089631in}}%
\pgfpathmoveto{\pgfqpoint{3.197606in}{2.086682in}}%
\pgfpathlineto{\pgfqpoint{3.197606in}{2.086682in}}%
\pgfpathlineto{\pgfqpoint{3.197606in}{2.089631in}}%
\pgfpathlineto{\pgfqpoint{3.202146in}{2.089631in}}%
\pgfpathlineto{\pgfqpoint{3.202146in}{2.086682in}}%
\pgfpathmoveto{\pgfqpoint{3.197606in}{2.089631in}}%
\pgfpathlineto{\pgfqpoint{3.197606in}{2.089631in}}%
\pgfpathlineto{\pgfqpoint{3.197606in}{2.092580in}}%
\pgfpathlineto{\pgfqpoint{3.202146in}{2.092580in}}%
\pgfpathlineto{\pgfqpoint{3.202146in}{2.089631in}}%
\pgfpathmoveto{\pgfqpoint{3.202146in}{2.089631in}}%
\pgfpathlineto{\pgfqpoint{3.202146in}{2.089631in}}%
\pgfpathlineto{\pgfqpoint{3.202146in}{2.092580in}}%
\pgfpathlineto{\pgfqpoint{3.206687in}{2.092580in}}%
\pgfpathlineto{\pgfqpoint{3.206687in}{2.089631in}}%
\pgfpathmoveto{\pgfqpoint{3.202146in}{2.092580in}}%
\pgfpathlineto{\pgfqpoint{3.202146in}{2.092580in}}%
\pgfpathlineto{\pgfqpoint{3.202146in}{2.095529in}}%
\pgfpathlineto{\pgfqpoint{3.206687in}{2.095529in}}%
\pgfpathlineto{\pgfqpoint{3.206687in}{2.092580in}}%
\pgfpathmoveto{\pgfqpoint{3.202146in}{2.095529in}}%
\pgfpathlineto{\pgfqpoint{3.202146in}{2.095529in}}%
\pgfpathlineto{\pgfqpoint{3.202146in}{2.098479in}}%
\pgfpathlineto{\pgfqpoint{3.206687in}{2.098479in}}%
\pgfpathlineto{\pgfqpoint{3.206687in}{2.095529in}}%
\pgfpathmoveto{\pgfqpoint{3.206687in}{2.092580in}}%
\pgfpathlineto{\pgfqpoint{3.206687in}{2.092580in}}%
\pgfpathlineto{\pgfqpoint{3.206687in}{2.095529in}}%
\pgfpathlineto{\pgfqpoint{3.211228in}{2.095529in}}%
\pgfpathlineto{\pgfqpoint{3.211228in}{2.092580in}}%
\pgfpathmoveto{\pgfqpoint{3.206687in}{2.095529in}}%
\pgfpathlineto{\pgfqpoint{3.206687in}{2.095529in}}%
\pgfpathlineto{\pgfqpoint{3.206687in}{2.098479in}}%
\pgfpathlineto{\pgfqpoint{3.211228in}{2.098479in}}%
\pgfpathlineto{\pgfqpoint{3.211228in}{2.095529in}}%
\pgfpathmoveto{\pgfqpoint{3.211228in}{2.095529in}}%
\pgfpathlineto{\pgfqpoint{3.211228in}{2.095529in}}%
\pgfpathlineto{\pgfqpoint{3.211228in}{2.098479in}}%
\pgfpathlineto{\pgfqpoint{3.215769in}{2.098479in}}%
\pgfpathlineto{\pgfqpoint{3.215769in}{2.095529in}}%
\pgfpathmoveto{\pgfqpoint{3.211228in}{2.098479in}}%
\pgfpathlineto{\pgfqpoint{3.211228in}{2.098479in}}%
\pgfpathlineto{\pgfqpoint{3.211228in}{2.101428in}}%
\pgfpathlineto{\pgfqpoint{3.215769in}{2.101428in}}%
\pgfpathlineto{\pgfqpoint{3.215769in}{2.098479in}}%
\pgfpathmoveto{\pgfqpoint{3.211228in}{2.101428in}}%
\pgfpathlineto{\pgfqpoint{3.211228in}{2.101428in}}%
\pgfpathlineto{\pgfqpoint{3.211228in}{2.104377in}}%
\pgfpathlineto{\pgfqpoint{3.215769in}{2.104377in}}%
\pgfpathlineto{\pgfqpoint{3.215769in}{2.101428in}}%
\pgfpathmoveto{\pgfqpoint{3.215769in}{2.098479in}}%
\pgfpathlineto{\pgfqpoint{3.215769in}{2.098479in}}%
\pgfpathlineto{\pgfqpoint{3.215769in}{2.101428in}}%
\pgfpathlineto{\pgfqpoint{3.220310in}{2.101428in}}%
\pgfpathlineto{\pgfqpoint{3.220310in}{2.098479in}}%
\pgfpathmoveto{\pgfqpoint{3.215769in}{2.101428in}}%
\pgfpathlineto{\pgfqpoint{3.215769in}{2.101428in}}%
\pgfpathlineto{\pgfqpoint{3.215769in}{2.104377in}}%
\pgfpathlineto{\pgfqpoint{3.220310in}{2.104377in}}%
\pgfpathlineto{\pgfqpoint{3.220310in}{2.101428in}}%
\pgfpathmoveto{\pgfqpoint{3.211228in}{3.024531in}}%
\pgfpathlineto{\pgfqpoint{3.211228in}{3.024531in}}%
\pgfpathlineto{\pgfqpoint{3.211228in}{3.027480in}}%
\pgfpathlineto{\pgfqpoint{3.215769in}{3.027480in}}%
\pgfpathlineto{\pgfqpoint{3.215769in}{3.024531in}}%
\pgfpathmoveto{\pgfqpoint{3.211228in}{3.027480in}}%
\pgfpathlineto{\pgfqpoint{3.211228in}{3.027480in}}%
\pgfpathlineto{\pgfqpoint{3.211228in}{3.030430in}}%
\pgfpathlineto{\pgfqpoint{3.215769in}{3.030430in}}%
\pgfpathlineto{\pgfqpoint{3.215769in}{3.027480in}}%
\pgfpathmoveto{\pgfqpoint{3.215769in}{3.024531in}}%
\pgfpathlineto{\pgfqpoint{3.215769in}{3.024531in}}%
\pgfpathlineto{\pgfqpoint{3.215769in}{3.027480in}}%
\pgfpathlineto{\pgfqpoint{3.220310in}{3.027480in}}%
\pgfpathlineto{\pgfqpoint{3.220310in}{3.024531in}}%
\pgfpathmoveto{\pgfqpoint{3.215769in}{3.027480in}}%
\pgfpathlineto{\pgfqpoint{3.215769in}{3.027480in}}%
\pgfpathlineto{\pgfqpoint{3.215769in}{3.030430in}}%
\pgfpathlineto{\pgfqpoint{3.220310in}{3.030430in}}%
\pgfpathlineto{\pgfqpoint{3.220310in}{3.027480in}}%
\pgfpathmoveto{\pgfqpoint{3.211228in}{3.030430in}}%
\pgfpathlineto{\pgfqpoint{3.211228in}{3.030430in}}%
\pgfpathlineto{\pgfqpoint{3.211228in}{3.033379in}}%
\pgfpathlineto{\pgfqpoint{3.215769in}{3.033379in}}%
\pgfpathlineto{\pgfqpoint{3.215769in}{3.030430in}}%
\pgfpathmoveto{\pgfqpoint{3.211228in}{3.033379in}}%
\pgfpathlineto{\pgfqpoint{3.211228in}{3.033379in}}%
\pgfpathlineto{\pgfqpoint{3.211228in}{3.036328in}}%
\pgfpathlineto{\pgfqpoint{3.215769in}{3.036328in}}%
\pgfpathlineto{\pgfqpoint{3.215769in}{3.033379in}}%
\pgfpathmoveto{\pgfqpoint{3.215769in}{3.030430in}}%
\pgfpathlineto{\pgfqpoint{3.215769in}{3.030430in}}%
\pgfpathlineto{\pgfqpoint{3.215769in}{3.033379in}}%
\pgfpathlineto{\pgfqpoint{3.220310in}{3.033379in}}%
\pgfpathlineto{\pgfqpoint{3.220310in}{3.030430in}}%
\pgfpathmoveto{\pgfqpoint{3.215769in}{3.033379in}}%
\pgfpathlineto{\pgfqpoint{3.215769in}{3.033379in}}%
\pgfpathlineto{\pgfqpoint{3.215769in}{3.036328in}}%
\pgfpathlineto{\pgfqpoint{3.220310in}{3.036328in}}%
\pgfpathlineto{\pgfqpoint{3.220310in}{3.033379in}}%
\pgfpathmoveto{\pgfqpoint{3.202146in}{3.036328in}}%
\pgfpathlineto{\pgfqpoint{3.202146in}{3.036328in}}%
\pgfpathlineto{\pgfqpoint{3.202146in}{3.039277in}}%
\pgfpathlineto{\pgfqpoint{3.206687in}{3.039277in}}%
\pgfpathlineto{\pgfqpoint{3.206687in}{3.036328in}}%
\pgfpathmoveto{\pgfqpoint{3.202146in}{3.039277in}}%
\pgfpathlineto{\pgfqpoint{3.202146in}{3.039277in}}%
\pgfpathlineto{\pgfqpoint{3.202146in}{3.042226in}}%
\pgfpathlineto{\pgfqpoint{3.206687in}{3.042226in}}%
\pgfpathlineto{\pgfqpoint{3.206687in}{3.039277in}}%
\pgfpathmoveto{\pgfqpoint{3.206687in}{3.036328in}}%
\pgfpathlineto{\pgfqpoint{3.206687in}{3.036328in}}%
\pgfpathlineto{\pgfqpoint{3.206687in}{3.039277in}}%
\pgfpathlineto{\pgfqpoint{3.211228in}{3.039277in}}%
\pgfpathlineto{\pgfqpoint{3.211228in}{3.036328in}}%
\pgfpathmoveto{\pgfqpoint{3.206687in}{3.039277in}}%
\pgfpathlineto{\pgfqpoint{3.206687in}{3.039277in}}%
\pgfpathlineto{\pgfqpoint{3.206687in}{3.042226in}}%
\pgfpathlineto{\pgfqpoint{3.211228in}{3.042226in}}%
\pgfpathlineto{\pgfqpoint{3.211228in}{3.039277in}}%
\pgfpathmoveto{\pgfqpoint{3.202146in}{3.042226in}}%
\pgfpathlineto{\pgfqpoint{3.202146in}{3.042226in}}%
\pgfpathlineto{\pgfqpoint{3.202146in}{3.045175in}}%
\pgfpathlineto{\pgfqpoint{3.206687in}{3.045175in}}%
\pgfpathlineto{\pgfqpoint{3.206687in}{3.042226in}}%
\pgfpathmoveto{\pgfqpoint{3.202146in}{3.045175in}}%
\pgfpathlineto{\pgfqpoint{3.202146in}{3.045175in}}%
\pgfpathlineto{\pgfqpoint{3.202146in}{3.048124in}}%
\pgfpathlineto{\pgfqpoint{3.206687in}{3.048124in}}%
\pgfpathlineto{\pgfqpoint{3.206687in}{3.045175in}}%
\pgfpathmoveto{\pgfqpoint{3.206687in}{3.042226in}}%
\pgfpathlineto{\pgfqpoint{3.206687in}{3.042226in}}%
\pgfpathlineto{\pgfqpoint{3.206687in}{3.045175in}}%
\pgfpathlineto{\pgfqpoint{3.211228in}{3.045175in}}%
\pgfpathlineto{\pgfqpoint{3.211228in}{3.042226in}}%
\pgfpathmoveto{\pgfqpoint{3.206687in}{3.045175in}}%
\pgfpathlineto{\pgfqpoint{3.206687in}{3.045175in}}%
\pgfpathlineto{\pgfqpoint{3.206687in}{3.048124in}}%
\pgfpathlineto{\pgfqpoint{3.211228in}{3.048124in}}%
\pgfpathlineto{\pgfqpoint{3.211228in}{3.045175in}}%
\pgfpathmoveto{\pgfqpoint{3.211228in}{3.036328in}}%
\pgfpathlineto{\pgfqpoint{3.211228in}{3.036328in}}%
\pgfpathlineto{\pgfqpoint{3.211228in}{3.039277in}}%
\pgfpathlineto{\pgfqpoint{3.215769in}{3.039277in}}%
\pgfpathlineto{\pgfqpoint{3.215769in}{3.036328in}}%
\pgfpathmoveto{\pgfqpoint{3.211228in}{3.039277in}}%
\pgfpathlineto{\pgfqpoint{3.211228in}{3.039277in}}%
\pgfpathlineto{\pgfqpoint{3.211228in}{3.042226in}}%
\pgfpathlineto{\pgfqpoint{3.215769in}{3.042226in}}%
\pgfpathlineto{\pgfqpoint{3.215769in}{3.039277in}}%
\pgfpathmoveto{\pgfqpoint{3.138575in}{3.118907in}}%
\pgfpathlineto{\pgfqpoint{3.138575in}{3.118907in}}%
\pgfpathlineto{\pgfqpoint{3.138575in}{3.121857in}}%
\pgfpathlineto{\pgfqpoint{3.143116in}{3.121857in}}%
\pgfpathlineto{\pgfqpoint{3.143116in}{3.118907in}}%
\pgfpathmoveto{\pgfqpoint{3.138575in}{3.121857in}}%
\pgfpathlineto{\pgfqpoint{3.138575in}{3.121857in}}%
\pgfpathlineto{\pgfqpoint{3.138575in}{3.124806in}}%
\pgfpathlineto{\pgfqpoint{3.143116in}{3.124806in}}%
\pgfpathlineto{\pgfqpoint{3.143116in}{3.121857in}}%
\pgfpathmoveto{\pgfqpoint{3.143116in}{3.118907in}}%
\pgfpathlineto{\pgfqpoint{3.143116in}{3.118907in}}%
\pgfpathlineto{\pgfqpoint{3.143116in}{3.121857in}}%
\pgfpathlineto{\pgfqpoint{3.147657in}{3.121857in}}%
\pgfpathlineto{\pgfqpoint{3.147657in}{3.118907in}}%
\pgfpathmoveto{\pgfqpoint{3.143116in}{3.121857in}}%
\pgfpathlineto{\pgfqpoint{3.143116in}{3.121857in}}%
\pgfpathlineto{\pgfqpoint{3.143116in}{3.124806in}}%
\pgfpathlineto{\pgfqpoint{3.147657in}{3.124806in}}%
\pgfpathlineto{\pgfqpoint{3.147657in}{3.121857in}}%
\pgfpathmoveto{\pgfqpoint{3.138575in}{3.124806in}}%
\pgfpathlineto{\pgfqpoint{3.138575in}{3.124806in}}%
\pgfpathlineto{\pgfqpoint{3.138575in}{3.127755in}}%
\pgfpathlineto{\pgfqpoint{3.143116in}{3.127755in}}%
\pgfpathlineto{\pgfqpoint{3.143116in}{3.124806in}}%
\pgfpathmoveto{\pgfqpoint{3.138575in}{3.127755in}}%
\pgfpathlineto{\pgfqpoint{3.138575in}{3.127755in}}%
\pgfpathlineto{\pgfqpoint{3.138575in}{3.130704in}}%
\pgfpathlineto{\pgfqpoint{3.143116in}{3.130704in}}%
\pgfpathlineto{\pgfqpoint{3.143116in}{3.127755in}}%
\pgfpathmoveto{\pgfqpoint{3.143116in}{3.124806in}}%
\pgfpathlineto{\pgfqpoint{3.143116in}{3.124806in}}%
\pgfpathlineto{\pgfqpoint{3.143116in}{3.127755in}}%
\pgfpathlineto{\pgfqpoint{3.147657in}{3.127755in}}%
\pgfpathlineto{\pgfqpoint{3.147657in}{3.124806in}}%
\pgfpathmoveto{\pgfqpoint{3.143116in}{3.127755in}}%
\pgfpathlineto{\pgfqpoint{3.143116in}{3.127755in}}%
\pgfpathlineto{\pgfqpoint{3.143116in}{3.130704in}}%
\pgfpathlineto{\pgfqpoint{3.147657in}{3.130704in}}%
\pgfpathlineto{\pgfqpoint{3.147657in}{3.127755in}}%
\pgfpathmoveto{\pgfqpoint{3.129493in}{3.130704in}}%
\pgfpathlineto{\pgfqpoint{3.129493in}{3.130704in}}%
\pgfpathlineto{\pgfqpoint{3.129493in}{3.133654in}}%
\pgfpathlineto{\pgfqpoint{3.134034in}{3.133654in}}%
\pgfpathlineto{\pgfqpoint{3.134034in}{3.130704in}}%
\pgfpathmoveto{\pgfqpoint{3.129493in}{3.133654in}}%
\pgfpathlineto{\pgfqpoint{3.129493in}{3.133654in}}%
\pgfpathlineto{\pgfqpoint{3.129493in}{3.136603in}}%
\pgfpathlineto{\pgfqpoint{3.134034in}{3.136603in}}%
\pgfpathlineto{\pgfqpoint{3.134034in}{3.133654in}}%
\pgfpathmoveto{\pgfqpoint{3.134034in}{3.130704in}}%
\pgfpathlineto{\pgfqpoint{3.134034in}{3.130704in}}%
\pgfpathlineto{\pgfqpoint{3.134034in}{3.133654in}}%
\pgfpathlineto{\pgfqpoint{3.138575in}{3.133654in}}%
\pgfpathlineto{\pgfqpoint{3.138575in}{3.130704in}}%
\pgfpathmoveto{\pgfqpoint{3.134034in}{3.133654in}}%
\pgfpathlineto{\pgfqpoint{3.134034in}{3.133654in}}%
\pgfpathlineto{\pgfqpoint{3.134034in}{3.136603in}}%
\pgfpathlineto{\pgfqpoint{3.138575in}{3.136603in}}%
\pgfpathlineto{\pgfqpoint{3.138575in}{3.133654in}}%
\pgfpathmoveto{\pgfqpoint{3.129493in}{3.136603in}}%
\pgfpathlineto{\pgfqpoint{3.129493in}{3.136603in}}%
\pgfpathlineto{\pgfqpoint{3.129493in}{3.139552in}}%
\pgfpathlineto{\pgfqpoint{3.134034in}{3.139552in}}%
\pgfpathlineto{\pgfqpoint{3.134034in}{3.136603in}}%
\pgfpathmoveto{\pgfqpoint{3.129493in}{3.139552in}}%
\pgfpathlineto{\pgfqpoint{3.129493in}{3.139552in}}%
\pgfpathlineto{\pgfqpoint{3.129493in}{3.142502in}}%
\pgfpathlineto{\pgfqpoint{3.134034in}{3.142502in}}%
\pgfpathlineto{\pgfqpoint{3.134034in}{3.139552in}}%
\pgfpathmoveto{\pgfqpoint{3.134034in}{3.136603in}}%
\pgfpathlineto{\pgfqpoint{3.134034in}{3.136603in}}%
\pgfpathlineto{\pgfqpoint{3.134034in}{3.139552in}}%
\pgfpathlineto{\pgfqpoint{3.138575in}{3.139552in}}%
\pgfpathlineto{\pgfqpoint{3.138575in}{3.136603in}}%
\pgfpathmoveto{\pgfqpoint{3.134034in}{3.139552in}}%
\pgfpathlineto{\pgfqpoint{3.134034in}{3.139552in}}%
\pgfpathlineto{\pgfqpoint{3.134034in}{3.142502in}}%
\pgfpathlineto{\pgfqpoint{3.138575in}{3.142502in}}%
\pgfpathlineto{\pgfqpoint{3.138575in}{3.139552in}}%
\pgfpathmoveto{\pgfqpoint{3.138575in}{3.130704in}}%
\pgfpathlineto{\pgfqpoint{3.138575in}{3.130704in}}%
\pgfpathlineto{\pgfqpoint{3.138575in}{3.133654in}}%
\pgfpathlineto{\pgfqpoint{3.143116in}{3.133654in}}%
\pgfpathlineto{\pgfqpoint{3.143116in}{3.130704in}}%
\pgfpathmoveto{\pgfqpoint{3.138575in}{3.133654in}}%
\pgfpathlineto{\pgfqpoint{3.138575in}{3.133654in}}%
\pgfpathlineto{\pgfqpoint{3.138575in}{3.136603in}}%
\pgfpathlineto{\pgfqpoint{3.143116in}{3.136603in}}%
\pgfpathlineto{\pgfqpoint{3.143116in}{3.133654in}}%
\pgfpathmoveto{\pgfqpoint{3.174901in}{3.071719in}}%
\pgfpathlineto{\pgfqpoint{3.174901in}{3.071719in}}%
\pgfpathlineto{\pgfqpoint{3.174901in}{3.074668in}}%
\pgfpathlineto{\pgfqpoint{3.179442in}{3.074668in}}%
\pgfpathlineto{\pgfqpoint{3.179442in}{3.071719in}}%
\pgfpathmoveto{\pgfqpoint{3.174901in}{3.074668in}}%
\pgfpathlineto{\pgfqpoint{3.174901in}{3.074668in}}%
\pgfpathlineto{\pgfqpoint{3.174901in}{3.077617in}}%
\pgfpathlineto{\pgfqpoint{3.179442in}{3.077617in}}%
\pgfpathlineto{\pgfqpoint{3.179442in}{3.074668in}}%
\pgfpathmoveto{\pgfqpoint{3.179442in}{3.071719in}}%
\pgfpathlineto{\pgfqpoint{3.179442in}{3.071719in}}%
\pgfpathlineto{\pgfqpoint{3.179442in}{3.074668in}}%
\pgfpathlineto{\pgfqpoint{3.183983in}{3.074668in}}%
\pgfpathlineto{\pgfqpoint{3.183983in}{3.071719in}}%
\pgfpathmoveto{\pgfqpoint{3.179442in}{3.074668in}}%
\pgfpathlineto{\pgfqpoint{3.179442in}{3.074668in}}%
\pgfpathlineto{\pgfqpoint{3.179442in}{3.077617in}}%
\pgfpathlineto{\pgfqpoint{3.183983in}{3.077617in}}%
\pgfpathlineto{\pgfqpoint{3.183983in}{3.074668in}}%
\pgfpathmoveto{\pgfqpoint{3.174901in}{3.077617in}}%
\pgfpathlineto{\pgfqpoint{3.174901in}{3.077617in}}%
\pgfpathlineto{\pgfqpoint{3.174901in}{3.080567in}}%
\pgfpathlineto{\pgfqpoint{3.179442in}{3.080567in}}%
\pgfpathlineto{\pgfqpoint{3.179442in}{3.077617in}}%
\pgfpathmoveto{\pgfqpoint{3.174901in}{3.080567in}}%
\pgfpathlineto{\pgfqpoint{3.174901in}{3.080567in}}%
\pgfpathlineto{\pgfqpoint{3.174901in}{3.083516in}}%
\pgfpathlineto{\pgfqpoint{3.179442in}{3.083516in}}%
\pgfpathlineto{\pgfqpoint{3.179442in}{3.080567in}}%
\pgfpathmoveto{\pgfqpoint{3.179442in}{3.077617in}}%
\pgfpathlineto{\pgfqpoint{3.179442in}{3.077617in}}%
\pgfpathlineto{\pgfqpoint{3.179442in}{3.080567in}}%
\pgfpathlineto{\pgfqpoint{3.183983in}{3.080567in}}%
\pgfpathlineto{\pgfqpoint{3.183983in}{3.077617in}}%
\pgfpathmoveto{\pgfqpoint{3.179442in}{3.080567in}}%
\pgfpathlineto{\pgfqpoint{3.179442in}{3.080567in}}%
\pgfpathlineto{\pgfqpoint{3.179442in}{3.083516in}}%
\pgfpathlineto{\pgfqpoint{3.183983in}{3.083516in}}%
\pgfpathlineto{\pgfqpoint{3.183983in}{3.080567in}}%
\pgfpathmoveto{\pgfqpoint{3.165820in}{3.083516in}}%
\pgfpathlineto{\pgfqpoint{3.165820in}{3.083516in}}%
\pgfpathlineto{\pgfqpoint{3.165820in}{3.086465in}}%
\pgfpathlineto{\pgfqpoint{3.170361in}{3.086465in}}%
\pgfpathlineto{\pgfqpoint{3.170361in}{3.083516in}}%
\pgfpathmoveto{\pgfqpoint{3.165820in}{3.086465in}}%
\pgfpathlineto{\pgfqpoint{3.165820in}{3.086465in}}%
\pgfpathlineto{\pgfqpoint{3.165820in}{3.089414in}}%
\pgfpathlineto{\pgfqpoint{3.170361in}{3.089414in}}%
\pgfpathlineto{\pgfqpoint{3.170361in}{3.086465in}}%
\pgfpathmoveto{\pgfqpoint{3.170361in}{3.083516in}}%
\pgfpathlineto{\pgfqpoint{3.170361in}{3.083516in}}%
\pgfpathlineto{\pgfqpoint{3.170361in}{3.086465in}}%
\pgfpathlineto{\pgfqpoint{3.174901in}{3.086465in}}%
\pgfpathlineto{\pgfqpoint{3.174901in}{3.083516in}}%
\pgfpathmoveto{\pgfqpoint{3.170361in}{3.086465in}}%
\pgfpathlineto{\pgfqpoint{3.170361in}{3.086465in}}%
\pgfpathlineto{\pgfqpoint{3.170361in}{3.089414in}}%
\pgfpathlineto{\pgfqpoint{3.174901in}{3.089414in}}%
\pgfpathlineto{\pgfqpoint{3.174901in}{3.086465in}}%
\pgfpathmoveto{\pgfqpoint{3.165820in}{3.089414in}}%
\pgfpathlineto{\pgfqpoint{3.165820in}{3.089414in}}%
\pgfpathlineto{\pgfqpoint{3.165820in}{3.092364in}}%
\pgfpathlineto{\pgfqpoint{3.170361in}{3.092364in}}%
\pgfpathlineto{\pgfqpoint{3.170361in}{3.089414in}}%
\pgfpathmoveto{\pgfqpoint{3.165820in}{3.092364in}}%
\pgfpathlineto{\pgfqpoint{3.165820in}{3.092364in}}%
\pgfpathlineto{\pgfqpoint{3.165820in}{3.095313in}}%
\pgfpathlineto{\pgfqpoint{3.170361in}{3.095313in}}%
\pgfpathlineto{\pgfqpoint{3.170361in}{3.092364in}}%
\pgfpathmoveto{\pgfqpoint{3.170361in}{3.089414in}}%
\pgfpathlineto{\pgfqpoint{3.170361in}{3.089414in}}%
\pgfpathlineto{\pgfqpoint{3.170361in}{3.092364in}}%
\pgfpathlineto{\pgfqpoint{3.174901in}{3.092364in}}%
\pgfpathlineto{\pgfqpoint{3.174901in}{3.089414in}}%
\pgfpathmoveto{\pgfqpoint{3.170361in}{3.092364in}}%
\pgfpathlineto{\pgfqpoint{3.170361in}{3.092364in}}%
\pgfpathlineto{\pgfqpoint{3.170361in}{3.095313in}}%
\pgfpathlineto{\pgfqpoint{3.174901in}{3.095313in}}%
\pgfpathlineto{\pgfqpoint{3.174901in}{3.092364in}}%
\pgfpathmoveto{\pgfqpoint{3.174901in}{3.083516in}}%
\pgfpathlineto{\pgfqpoint{3.174901in}{3.083516in}}%
\pgfpathlineto{\pgfqpoint{3.174901in}{3.086465in}}%
\pgfpathlineto{\pgfqpoint{3.179442in}{3.086465in}}%
\pgfpathlineto{\pgfqpoint{3.179442in}{3.083516in}}%
\pgfpathmoveto{\pgfqpoint{3.174901in}{3.086465in}}%
\pgfpathlineto{\pgfqpoint{3.174901in}{3.086465in}}%
\pgfpathlineto{\pgfqpoint{3.174901in}{3.089414in}}%
\pgfpathlineto{\pgfqpoint{3.179442in}{3.089414in}}%
\pgfpathlineto{\pgfqpoint{3.179442in}{3.086465in}}%
\pgfpathmoveto{\pgfqpoint{3.193065in}{3.048124in}}%
\pgfpathlineto{\pgfqpoint{3.193065in}{3.048124in}}%
\pgfpathlineto{\pgfqpoint{3.193065in}{3.051074in}}%
\pgfpathlineto{\pgfqpoint{3.197606in}{3.051074in}}%
\pgfpathlineto{\pgfqpoint{3.197606in}{3.048124in}}%
\pgfpathmoveto{\pgfqpoint{3.193065in}{3.051074in}}%
\pgfpathlineto{\pgfqpoint{3.193065in}{3.051074in}}%
\pgfpathlineto{\pgfqpoint{3.193065in}{3.054023in}}%
\pgfpathlineto{\pgfqpoint{3.197606in}{3.054023in}}%
\pgfpathlineto{\pgfqpoint{3.197606in}{3.051074in}}%
\pgfpathmoveto{\pgfqpoint{3.197606in}{3.048124in}}%
\pgfpathlineto{\pgfqpoint{3.197606in}{3.048124in}}%
\pgfpathlineto{\pgfqpoint{3.197606in}{3.051074in}}%
\pgfpathlineto{\pgfqpoint{3.202146in}{3.051074in}}%
\pgfpathlineto{\pgfqpoint{3.202146in}{3.048124in}}%
\pgfpathmoveto{\pgfqpoint{3.197606in}{3.051074in}}%
\pgfpathlineto{\pgfqpoint{3.197606in}{3.051074in}}%
\pgfpathlineto{\pgfqpoint{3.197606in}{3.054023in}}%
\pgfpathlineto{\pgfqpoint{3.202146in}{3.054023in}}%
\pgfpathlineto{\pgfqpoint{3.202146in}{3.051074in}}%
\pgfpathmoveto{\pgfqpoint{3.193065in}{3.054023in}}%
\pgfpathlineto{\pgfqpoint{3.193065in}{3.054023in}}%
\pgfpathlineto{\pgfqpoint{3.193065in}{3.056972in}}%
\pgfpathlineto{\pgfqpoint{3.197606in}{3.056972in}}%
\pgfpathlineto{\pgfqpoint{3.197606in}{3.054023in}}%
\pgfpathmoveto{\pgfqpoint{3.193065in}{3.056972in}}%
\pgfpathlineto{\pgfqpoint{3.193065in}{3.056972in}}%
\pgfpathlineto{\pgfqpoint{3.193065in}{3.059922in}}%
\pgfpathlineto{\pgfqpoint{3.197606in}{3.059922in}}%
\pgfpathlineto{\pgfqpoint{3.197606in}{3.056972in}}%
\pgfpathmoveto{\pgfqpoint{3.197606in}{3.054023in}}%
\pgfpathlineto{\pgfqpoint{3.197606in}{3.054023in}}%
\pgfpathlineto{\pgfqpoint{3.197606in}{3.056972in}}%
\pgfpathlineto{\pgfqpoint{3.202146in}{3.056972in}}%
\pgfpathlineto{\pgfqpoint{3.202146in}{3.054023in}}%
\pgfpathmoveto{\pgfqpoint{3.197606in}{3.056972in}}%
\pgfpathlineto{\pgfqpoint{3.197606in}{3.056972in}}%
\pgfpathlineto{\pgfqpoint{3.197606in}{3.059922in}}%
\pgfpathlineto{\pgfqpoint{3.202146in}{3.059922in}}%
\pgfpathlineto{\pgfqpoint{3.202146in}{3.056972in}}%
\pgfpathmoveto{\pgfqpoint{3.183983in}{3.059922in}}%
\pgfpathlineto{\pgfqpoint{3.183983in}{3.059922in}}%
\pgfpathlineto{\pgfqpoint{3.183983in}{3.062871in}}%
\pgfpathlineto{\pgfqpoint{3.188524in}{3.062871in}}%
\pgfpathlineto{\pgfqpoint{3.188524in}{3.059922in}}%
\pgfpathmoveto{\pgfqpoint{3.183983in}{3.062871in}}%
\pgfpathlineto{\pgfqpoint{3.183983in}{3.062871in}}%
\pgfpathlineto{\pgfqpoint{3.183983in}{3.065820in}}%
\pgfpathlineto{\pgfqpoint{3.188524in}{3.065820in}}%
\pgfpathlineto{\pgfqpoint{3.188524in}{3.062871in}}%
\pgfpathmoveto{\pgfqpoint{3.188524in}{3.059922in}}%
\pgfpathlineto{\pgfqpoint{3.188524in}{3.059922in}}%
\pgfpathlineto{\pgfqpoint{3.188524in}{3.062871in}}%
\pgfpathlineto{\pgfqpoint{3.193065in}{3.062871in}}%
\pgfpathlineto{\pgfqpoint{3.193065in}{3.059922in}}%
\pgfpathmoveto{\pgfqpoint{3.188524in}{3.062871in}}%
\pgfpathlineto{\pgfqpoint{3.188524in}{3.062871in}}%
\pgfpathlineto{\pgfqpoint{3.188524in}{3.065820in}}%
\pgfpathlineto{\pgfqpoint{3.193065in}{3.065820in}}%
\pgfpathlineto{\pgfqpoint{3.193065in}{3.062871in}}%
\pgfpathmoveto{\pgfqpoint{3.183983in}{3.065820in}}%
\pgfpathlineto{\pgfqpoint{3.183983in}{3.065820in}}%
\pgfpathlineto{\pgfqpoint{3.183983in}{3.068769in}}%
\pgfpathlineto{\pgfqpoint{3.188524in}{3.068769in}}%
\pgfpathlineto{\pgfqpoint{3.188524in}{3.065820in}}%
\pgfpathmoveto{\pgfqpoint{3.183983in}{3.068769in}}%
\pgfpathlineto{\pgfqpoint{3.183983in}{3.068769in}}%
\pgfpathlineto{\pgfqpoint{3.183983in}{3.071719in}}%
\pgfpathlineto{\pgfqpoint{3.188524in}{3.071719in}}%
\pgfpathlineto{\pgfqpoint{3.188524in}{3.068769in}}%
\pgfpathmoveto{\pgfqpoint{3.188524in}{3.065820in}}%
\pgfpathlineto{\pgfqpoint{3.188524in}{3.065820in}}%
\pgfpathlineto{\pgfqpoint{3.188524in}{3.068769in}}%
\pgfpathlineto{\pgfqpoint{3.193065in}{3.068769in}}%
\pgfpathlineto{\pgfqpoint{3.193065in}{3.065820in}}%
\pgfpathmoveto{\pgfqpoint{3.188524in}{3.068769in}}%
\pgfpathlineto{\pgfqpoint{3.188524in}{3.068769in}}%
\pgfpathlineto{\pgfqpoint{3.188524in}{3.071719in}}%
\pgfpathlineto{\pgfqpoint{3.193065in}{3.071719in}}%
\pgfpathlineto{\pgfqpoint{3.193065in}{3.068769in}}%
\pgfpathmoveto{\pgfqpoint{3.193065in}{3.059922in}}%
\pgfpathlineto{\pgfqpoint{3.193065in}{3.059922in}}%
\pgfpathlineto{\pgfqpoint{3.193065in}{3.062871in}}%
\pgfpathlineto{\pgfqpoint{3.197606in}{3.062871in}}%
\pgfpathlineto{\pgfqpoint{3.197606in}{3.059922in}}%
\pgfpathmoveto{\pgfqpoint{3.193065in}{3.062871in}}%
\pgfpathlineto{\pgfqpoint{3.193065in}{3.062871in}}%
\pgfpathlineto{\pgfqpoint{3.193065in}{3.065820in}}%
\pgfpathlineto{\pgfqpoint{3.197606in}{3.065820in}}%
\pgfpathlineto{\pgfqpoint{3.197606in}{3.062871in}}%
\pgfpathmoveto{\pgfqpoint{3.202146in}{3.048124in}}%
\pgfpathlineto{\pgfqpoint{3.202146in}{3.048124in}}%
\pgfpathlineto{\pgfqpoint{3.202146in}{3.051074in}}%
\pgfpathlineto{\pgfqpoint{3.206687in}{3.051074in}}%
\pgfpathlineto{\pgfqpoint{3.206687in}{3.048124in}}%
\pgfpathmoveto{\pgfqpoint{3.202146in}{3.051074in}}%
\pgfpathlineto{\pgfqpoint{3.202146in}{3.051074in}}%
\pgfpathlineto{\pgfqpoint{3.202146in}{3.054023in}}%
\pgfpathlineto{\pgfqpoint{3.206687in}{3.054023in}}%
\pgfpathlineto{\pgfqpoint{3.206687in}{3.051074in}}%
\pgfpathmoveto{\pgfqpoint{3.183983in}{3.071719in}}%
\pgfpathlineto{\pgfqpoint{3.183983in}{3.071719in}}%
\pgfpathlineto{\pgfqpoint{3.183983in}{3.074668in}}%
\pgfpathlineto{\pgfqpoint{3.188524in}{3.074668in}}%
\pgfpathlineto{\pgfqpoint{3.188524in}{3.071719in}}%
\pgfpathmoveto{\pgfqpoint{3.183983in}{3.074668in}}%
\pgfpathlineto{\pgfqpoint{3.183983in}{3.074668in}}%
\pgfpathlineto{\pgfqpoint{3.183983in}{3.077617in}}%
\pgfpathlineto{\pgfqpoint{3.188524in}{3.077617in}}%
\pgfpathlineto{\pgfqpoint{3.188524in}{3.074668in}}%
\pgfpathmoveto{\pgfqpoint{3.156738in}{3.095313in}}%
\pgfpathlineto{\pgfqpoint{3.156738in}{3.095313in}}%
\pgfpathlineto{\pgfqpoint{3.156738in}{3.098262in}}%
\pgfpathlineto{\pgfqpoint{3.161279in}{3.098262in}}%
\pgfpathlineto{\pgfqpoint{3.161279in}{3.095313in}}%
\pgfpathmoveto{\pgfqpoint{3.156738in}{3.098262in}}%
\pgfpathlineto{\pgfqpoint{3.156738in}{3.098262in}}%
\pgfpathlineto{\pgfqpoint{3.156738in}{3.101212in}}%
\pgfpathlineto{\pgfqpoint{3.161279in}{3.101212in}}%
\pgfpathlineto{\pgfqpoint{3.161279in}{3.098262in}}%
\pgfpathmoveto{\pgfqpoint{3.161279in}{3.095313in}}%
\pgfpathlineto{\pgfqpoint{3.161279in}{3.095313in}}%
\pgfpathlineto{\pgfqpoint{3.161279in}{3.098262in}}%
\pgfpathlineto{\pgfqpoint{3.165820in}{3.098262in}}%
\pgfpathlineto{\pgfqpoint{3.165820in}{3.095313in}}%
\pgfpathmoveto{\pgfqpoint{3.161279in}{3.098262in}}%
\pgfpathlineto{\pgfqpoint{3.161279in}{3.098262in}}%
\pgfpathlineto{\pgfqpoint{3.161279in}{3.101212in}}%
\pgfpathlineto{\pgfqpoint{3.165820in}{3.101212in}}%
\pgfpathlineto{\pgfqpoint{3.165820in}{3.098262in}}%
\pgfpathmoveto{\pgfqpoint{3.156738in}{3.101212in}}%
\pgfpathlineto{\pgfqpoint{3.156738in}{3.101212in}}%
\pgfpathlineto{\pgfqpoint{3.156738in}{3.104161in}}%
\pgfpathlineto{\pgfqpoint{3.161279in}{3.104161in}}%
\pgfpathlineto{\pgfqpoint{3.161279in}{3.101212in}}%
\pgfpathmoveto{\pgfqpoint{3.156738in}{3.104161in}}%
\pgfpathlineto{\pgfqpoint{3.156738in}{3.104161in}}%
\pgfpathlineto{\pgfqpoint{3.156738in}{3.107110in}}%
\pgfpathlineto{\pgfqpoint{3.161279in}{3.107110in}}%
\pgfpathlineto{\pgfqpoint{3.161279in}{3.104161in}}%
\pgfpathmoveto{\pgfqpoint{3.161279in}{3.101212in}}%
\pgfpathlineto{\pgfqpoint{3.161279in}{3.101212in}}%
\pgfpathlineto{\pgfqpoint{3.161279in}{3.104161in}}%
\pgfpathlineto{\pgfqpoint{3.165820in}{3.104161in}}%
\pgfpathlineto{\pgfqpoint{3.165820in}{3.101212in}}%
\pgfpathmoveto{\pgfqpoint{3.161279in}{3.104161in}}%
\pgfpathlineto{\pgfqpoint{3.161279in}{3.104161in}}%
\pgfpathlineto{\pgfqpoint{3.161279in}{3.107110in}}%
\pgfpathlineto{\pgfqpoint{3.165820in}{3.107110in}}%
\pgfpathlineto{\pgfqpoint{3.165820in}{3.104161in}}%
\pgfpathmoveto{\pgfqpoint{3.147657in}{3.107110in}}%
\pgfpathlineto{\pgfqpoint{3.147657in}{3.107110in}}%
\pgfpathlineto{\pgfqpoint{3.147657in}{3.110059in}}%
\pgfpathlineto{\pgfqpoint{3.152197in}{3.110059in}}%
\pgfpathlineto{\pgfqpoint{3.152197in}{3.107110in}}%
\pgfpathmoveto{\pgfqpoint{3.147657in}{3.110059in}}%
\pgfpathlineto{\pgfqpoint{3.147657in}{3.110059in}}%
\pgfpathlineto{\pgfqpoint{3.147657in}{3.113009in}}%
\pgfpathlineto{\pgfqpoint{3.152197in}{3.113009in}}%
\pgfpathlineto{\pgfqpoint{3.152197in}{3.110059in}}%
\pgfpathmoveto{\pgfqpoint{3.152197in}{3.107110in}}%
\pgfpathlineto{\pgfqpoint{3.152197in}{3.107110in}}%
\pgfpathlineto{\pgfqpoint{3.152197in}{3.110059in}}%
\pgfpathlineto{\pgfqpoint{3.156738in}{3.110059in}}%
\pgfpathlineto{\pgfqpoint{3.156738in}{3.107110in}}%
\pgfpathmoveto{\pgfqpoint{3.152197in}{3.110059in}}%
\pgfpathlineto{\pgfqpoint{3.152197in}{3.110059in}}%
\pgfpathlineto{\pgfqpoint{3.152197in}{3.113009in}}%
\pgfpathlineto{\pgfqpoint{3.156738in}{3.113009in}}%
\pgfpathlineto{\pgfqpoint{3.156738in}{3.110059in}}%
\pgfpathmoveto{\pgfqpoint{3.147657in}{3.113009in}}%
\pgfpathlineto{\pgfqpoint{3.147657in}{3.113009in}}%
\pgfpathlineto{\pgfqpoint{3.147657in}{3.115958in}}%
\pgfpathlineto{\pgfqpoint{3.152197in}{3.115958in}}%
\pgfpathlineto{\pgfqpoint{3.152197in}{3.113009in}}%
\pgfpathmoveto{\pgfqpoint{3.147657in}{3.115958in}}%
\pgfpathlineto{\pgfqpoint{3.147657in}{3.115958in}}%
\pgfpathlineto{\pgfqpoint{3.147657in}{3.118907in}}%
\pgfpathlineto{\pgfqpoint{3.152197in}{3.118907in}}%
\pgfpathlineto{\pgfqpoint{3.152197in}{3.115958in}}%
\pgfpathmoveto{\pgfqpoint{3.152197in}{3.113009in}}%
\pgfpathlineto{\pgfqpoint{3.152197in}{3.113009in}}%
\pgfpathlineto{\pgfqpoint{3.152197in}{3.115958in}}%
\pgfpathlineto{\pgfqpoint{3.156738in}{3.115958in}}%
\pgfpathlineto{\pgfqpoint{3.156738in}{3.113009in}}%
\pgfpathmoveto{\pgfqpoint{3.152197in}{3.115958in}}%
\pgfpathlineto{\pgfqpoint{3.152197in}{3.115958in}}%
\pgfpathlineto{\pgfqpoint{3.152197in}{3.118907in}}%
\pgfpathlineto{\pgfqpoint{3.156738in}{3.118907in}}%
\pgfpathlineto{\pgfqpoint{3.156738in}{3.115958in}}%
\pgfpathmoveto{\pgfqpoint{3.156738in}{3.107110in}}%
\pgfpathlineto{\pgfqpoint{3.156738in}{3.107110in}}%
\pgfpathlineto{\pgfqpoint{3.156738in}{3.110059in}}%
\pgfpathlineto{\pgfqpoint{3.161279in}{3.110059in}}%
\pgfpathlineto{\pgfqpoint{3.161279in}{3.107110in}}%
\pgfpathmoveto{\pgfqpoint{3.156738in}{3.110059in}}%
\pgfpathlineto{\pgfqpoint{3.156738in}{3.110059in}}%
\pgfpathlineto{\pgfqpoint{3.156738in}{3.113009in}}%
\pgfpathlineto{\pgfqpoint{3.161279in}{3.113009in}}%
\pgfpathlineto{\pgfqpoint{3.161279in}{3.110059in}}%
\pgfpathmoveto{\pgfqpoint{3.165820in}{3.095313in}}%
\pgfpathlineto{\pgfqpoint{3.165820in}{3.095313in}}%
\pgfpathlineto{\pgfqpoint{3.165820in}{3.098262in}}%
\pgfpathlineto{\pgfqpoint{3.170361in}{3.098262in}}%
\pgfpathlineto{\pgfqpoint{3.170361in}{3.095313in}}%
\pgfpathmoveto{\pgfqpoint{3.165820in}{3.098262in}}%
\pgfpathlineto{\pgfqpoint{3.165820in}{3.098262in}}%
\pgfpathlineto{\pgfqpoint{3.165820in}{3.101212in}}%
\pgfpathlineto{\pgfqpoint{3.170361in}{3.101212in}}%
\pgfpathlineto{\pgfqpoint{3.170361in}{3.098262in}}%
\pgfpathmoveto{\pgfqpoint{3.147657in}{3.118907in}}%
\pgfpathlineto{\pgfqpoint{3.147657in}{3.118907in}}%
\pgfpathlineto{\pgfqpoint{3.147657in}{3.121857in}}%
\pgfpathlineto{\pgfqpoint{3.152197in}{3.121857in}}%
\pgfpathlineto{\pgfqpoint{3.152197in}{3.118907in}}%
\pgfpathmoveto{\pgfqpoint{3.147657in}{3.121857in}}%
\pgfpathlineto{\pgfqpoint{3.147657in}{3.121857in}}%
\pgfpathlineto{\pgfqpoint{3.147657in}{3.124806in}}%
\pgfpathlineto{\pgfqpoint{3.152197in}{3.124806in}}%
\pgfpathlineto{\pgfqpoint{3.152197in}{3.121857in}}%
\pgfpathmoveto{\pgfqpoint{3.102249in}{3.166095in}}%
\pgfpathlineto{\pgfqpoint{3.102249in}{3.166095in}}%
\pgfpathlineto{\pgfqpoint{3.102249in}{3.169044in}}%
\pgfpathlineto{\pgfqpoint{3.106789in}{3.169044in}}%
\pgfpathlineto{\pgfqpoint{3.106789in}{3.166095in}}%
\pgfpathmoveto{\pgfqpoint{3.102249in}{3.169044in}}%
\pgfpathlineto{\pgfqpoint{3.102249in}{3.169044in}}%
\pgfpathlineto{\pgfqpoint{3.102249in}{3.171993in}}%
\pgfpathlineto{\pgfqpoint{3.106789in}{3.171993in}}%
\pgfpathlineto{\pgfqpoint{3.106789in}{3.169044in}}%
\pgfpathmoveto{\pgfqpoint{3.106789in}{3.166095in}}%
\pgfpathlineto{\pgfqpoint{3.106789in}{3.166095in}}%
\pgfpathlineto{\pgfqpoint{3.106789in}{3.169044in}}%
\pgfpathlineto{\pgfqpoint{3.111330in}{3.169044in}}%
\pgfpathlineto{\pgfqpoint{3.111330in}{3.166095in}}%
\pgfpathmoveto{\pgfqpoint{3.106789in}{3.169044in}}%
\pgfpathlineto{\pgfqpoint{3.106789in}{3.169044in}}%
\pgfpathlineto{\pgfqpoint{3.106789in}{3.171993in}}%
\pgfpathlineto{\pgfqpoint{3.111330in}{3.171993in}}%
\pgfpathlineto{\pgfqpoint{3.111330in}{3.169044in}}%
\pgfpathmoveto{\pgfqpoint{3.102249in}{3.171993in}}%
\pgfpathlineto{\pgfqpoint{3.102249in}{3.171993in}}%
\pgfpathlineto{\pgfqpoint{3.102249in}{3.174942in}}%
\pgfpathlineto{\pgfqpoint{3.106789in}{3.174942in}}%
\pgfpathlineto{\pgfqpoint{3.106789in}{3.171993in}}%
\pgfpathmoveto{\pgfqpoint{3.102249in}{3.174942in}}%
\pgfpathlineto{\pgfqpoint{3.102249in}{3.174942in}}%
\pgfpathlineto{\pgfqpoint{3.102249in}{3.177891in}}%
\pgfpathlineto{\pgfqpoint{3.106789in}{3.177891in}}%
\pgfpathlineto{\pgfqpoint{3.106789in}{3.174942in}}%
\pgfpathmoveto{\pgfqpoint{3.106789in}{3.171993in}}%
\pgfpathlineto{\pgfqpoint{3.106789in}{3.171993in}}%
\pgfpathlineto{\pgfqpoint{3.106789in}{3.174942in}}%
\pgfpathlineto{\pgfqpoint{3.111330in}{3.174942in}}%
\pgfpathlineto{\pgfqpoint{3.111330in}{3.171993in}}%
\pgfpathmoveto{\pgfqpoint{3.106789in}{3.174942in}}%
\pgfpathlineto{\pgfqpoint{3.106789in}{3.174942in}}%
\pgfpathlineto{\pgfqpoint{3.106789in}{3.177891in}}%
\pgfpathlineto{\pgfqpoint{3.111330in}{3.177891in}}%
\pgfpathlineto{\pgfqpoint{3.111330in}{3.174942in}}%
\pgfpathmoveto{\pgfqpoint{3.093167in}{3.177891in}}%
\pgfpathlineto{\pgfqpoint{3.093167in}{3.177891in}}%
\pgfpathlineto{\pgfqpoint{3.093167in}{3.180841in}}%
\pgfpathlineto{\pgfqpoint{3.097708in}{3.180841in}}%
\pgfpathlineto{\pgfqpoint{3.097708in}{3.177891in}}%
\pgfpathmoveto{\pgfqpoint{3.093167in}{3.180841in}}%
\pgfpathlineto{\pgfqpoint{3.093167in}{3.180841in}}%
\pgfpathlineto{\pgfqpoint{3.093167in}{3.183790in}}%
\pgfpathlineto{\pgfqpoint{3.097708in}{3.183790in}}%
\pgfpathlineto{\pgfqpoint{3.097708in}{3.180841in}}%
\pgfpathmoveto{\pgfqpoint{3.097708in}{3.177891in}}%
\pgfpathlineto{\pgfqpoint{3.097708in}{3.177891in}}%
\pgfpathlineto{\pgfqpoint{3.097708in}{3.180841in}}%
\pgfpathlineto{\pgfqpoint{3.102249in}{3.180841in}}%
\pgfpathlineto{\pgfqpoint{3.102249in}{3.177891in}}%
\pgfpathmoveto{\pgfqpoint{3.097708in}{3.180841in}}%
\pgfpathlineto{\pgfqpoint{3.097708in}{3.180841in}}%
\pgfpathlineto{\pgfqpoint{3.097708in}{3.183790in}}%
\pgfpathlineto{\pgfqpoint{3.102249in}{3.183790in}}%
\pgfpathlineto{\pgfqpoint{3.102249in}{3.180841in}}%
\pgfpathmoveto{\pgfqpoint{3.093167in}{3.183790in}}%
\pgfpathlineto{\pgfqpoint{3.093167in}{3.183790in}}%
\pgfpathlineto{\pgfqpoint{3.093167in}{3.186739in}}%
\pgfpathlineto{\pgfqpoint{3.097708in}{3.186739in}}%
\pgfpathlineto{\pgfqpoint{3.097708in}{3.183790in}}%
\pgfpathmoveto{\pgfqpoint{3.093167in}{3.186739in}}%
\pgfpathlineto{\pgfqpoint{3.093167in}{3.186739in}}%
\pgfpathlineto{\pgfqpoint{3.093167in}{3.189688in}}%
\pgfpathlineto{\pgfqpoint{3.097708in}{3.189688in}}%
\pgfpathlineto{\pgfqpoint{3.097708in}{3.186739in}}%
\pgfpathmoveto{\pgfqpoint{3.097708in}{3.183790in}}%
\pgfpathlineto{\pgfqpoint{3.097708in}{3.183790in}}%
\pgfpathlineto{\pgfqpoint{3.097708in}{3.186739in}}%
\pgfpathlineto{\pgfqpoint{3.102249in}{3.186739in}}%
\pgfpathlineto{\pgfqpoint{3.102249in}{3.183790in}}%
\pgfpathmoveto{\pgfqpoint{3.097708in}{3.186739in}}%
\pgfpathlineto{\pgfqpoint{3.097708in}{3.186739in}}%
\pgfpathlineto{\pgfqpoint{3.097708in}{3.189688in}}%
\pgfpathlineto{\pgfqpoint{3.102249in}{3.189688in}}%
\pgfpathlineto{\pgfqpoint{3.102249in}{3.186739in}}%
\pgfpathmoveto{\pgfqpoint{3.102249in}{3.177891in}}%
\pgfpathlineto{\pgfqpoint{3.102249in}{3.177891in}}%
\pgfpathlineto{\pgfqpoint{3.102249in}{3.180841in}}%
\pgfpathlineto{\pgfqpoint{3.106789in}{3.180841in}}%
\pgfpathlineto{\pgfqpoint{3.106789in}{3.177891in}}%
\pgfpathmoveto{\pgfqpoint{3.102249in}{3.180841in}}%
\pgfpathlineto{\pgfqpoint{3.102249in}{3.180841in}}%
\pgfpathlineto{\pgfqpoint{3.102249in}{3.183790in}}%
\pgfpathlineto{\pgfqpoint{3.106789in}{3.183790in}}%
\pgfpathlineto{\pgfqpoint{3.106789in}{3.180841in}}%
\pgfpathmoveto{\pgfqpoint{3.120412in}{3.142502in}}%
\pgfpathlineto{\pgfqpoint{3.120412in}{3.142502in}}%
\pgfpathlineto{\pgfqpoint{3.120412in}{3.145451in}}%
\pgfpathlineto{\pgfqpoint{3.124953in}{3.145451in}}%
\pgfpathlineto{\pgfqpoint{3.124953in}{3.142502in}}%
\pgfpathmoveto{\pgfqpoint{3.120412in}{3.145451in}}%
\pgfpathlineto{\pgfqpoint{3.120412in}{3.145451in}}%
\pgfpathlineto{\pgfqpoint{3.120412in}{3.148400in}}%
\pgfpathlineto{\pgfqpoint{3.124953in}{3.148400in}}%
\pgfpathlineto{\pgfqpoint{3.124953in}{3.145451in}}%
\pgfpathmoveto{\pgfqpoint{3.124953in}{3.142502in}}%
\pgfpathlineto{\pgfqpoint{3.124953in}{3.142502in}}%
\pgfpathlineto{\pgfqpoint{3.124953in}{3.145451in}}%
\pgfpathlineto{\pgfqpoint{3.129493in}{3.145451in}}%
\pgfpathlineto{\pgfqpoint{3.129493in}{3.142502in}}%
\pgfpathmoveto{\pgfqpoint{3.124953in}{3.145451in}}%
\pgfpathlineto{\pgfqpoint{3.124953in}{3.145451in}}%
\pgfpathlineto{\pgfqpoint{3.124953in}{3.148400in}}%
\pgfpathlineto{\pgfqpoint{3.129493in}{3.148400in}}%
\pgfpathlineto{\pgfqpoint{3.129493in}{3.145451in}}%
\pgfpathmoveto{\pgfqpoint{3.120412in}{3.148400in}}%
\pgfpathlineto{\pgfqpoint{3.120412in}{3.148400in}}%
\pgfpathlineto{\pgfqpoint{3.120412in}{3.151349in}}%
\pgfpathlineto{\pgfqpoint{3.124953in}{3.151349in}}%
\pgfpathlineto{\pgfqpoint{3.124953in}{3.148400in}}%
\pgfpathmoveto{\pgfqpoint{3.120412in}{3.151349in}}%
\pgfpathlineto{\pgfqpoint{3.120412in}{3.151349in}}%
\pgfpathlineto{\pgfqpoint{3.120412in}{3.154298in}}%
\pgfpathlineto{\pgfqpoint{3.124953in}{3.154298in}}%
\pgfpathlineto{\pgfqpoint{3.124953in}{3.151349in}}%
\pgfpathmoveto{\pgfqpoint{3.124953in}{3.148400in}}%
\pgfpathlineto{\pgfqpoint{3.124953in}{3.148400in}}%
\pgfpathlineto{\pgfqpoint{3.124953in}{3.151349in}}%
\pgfpathlineto{\pgfqpoint{3.129493in}{3.151349in}}%
\pgfpathlineto{\pgfqpoint{3.129493in}{3.148400in}}%
\pgfpathmoveto{\pgfqpoint{3.124953in}{3.151349in}}%
\pgfpathlineto{\pgfqpoint{3.124953in}{3.151349in}}%
\pgfpathlineto{\pgfqpoint{3.124953in}{3.154298in}}%
\pgfpathlineto{\pgfqpoint{3.129493in}{3.154298in}}%
\pgfpathlineto{\pgfqpoint{3.129493in}{3.151349in}}%
\pgfpathmoveto{\pgfqpoint{3.111330in}{3.154298in}}%
\pgfpathlineto{\pgfqpoint{3.111330in}{3.154298in}}%
\pgfpathlineto{\pgfqpoint{3.111330in}{3.157247in}}%
\pgfpathlineto{\pgfqpoint{3.115871in}{3.157247in}}%
\pgfpathlineto{\pgfqpoint{3.115871in}{3.154298in}}%
\pgfpathmoveto{\pgfqpoint{3.111330in}{3.157247in}}%
\pgfpathlineto{\pgfqpoint{3.111330in}{3.157247in}}%
\pgfpathlineto{\pgfqpoint{3.111330in}{3.160197in}}%
\pgfpathlineto{\pgfqpoint{3.115871in}{3.160197in}}%
\pgfpathlineto{\pgfqpoint{3.115871in}{3.157247in}}%
\pgfpathmoveto{\pgfqpoint{3.115871in}{3.154298in}}%
\pgfpathlineto{\pgfqpoint{3.115871in}{3.154298in}}%
\pgfpathlineto{\pgfqpoint{3.115871in}{3.157247in}}%
\pgfpathlineto{\pgfqpoint{3.120412in}{3.157247in}}%
\pgfpathlineto{\pgfqpoint{3.120412in}{3.154298in}}%
\pgfpathmoveto{\pgfqpoint{3.115871in}{3.157247in}}%
\pgfpathlineto{\pgfqpoint{3.115871in}{3.157247in}}%
\pgfpathlineto{\pgfqpoint{3.115871in}{3.160197in}}%
\pgfpathlineto{\pgfqpoint{3.120412in}{3.160197in}}%
\pgfpathlineto{\pgfqpoint{3.120412in}{3.157247in}}%
\pgfpathmoveto{\pgfqpoint{3.111330in}{3.160197in}}%
\pgfpathlineto{\pgfqpoint{3.111330in}{3.160197in}}%
\pgfpathlineto{\pgfqpoint{3.111330in}{3.163146in}}%
\pgfpathlineto{\pgfqpoint{3.115871in}{3.163146in}}%
\pgfpathlineto{\pgfqpoint{3.115871in}{3.160197in}}%
\pgfpathmoveto{\pgfqpoint{3.111330in}{3.163146in}}%
\pgfpathlineto{\pgfqpoint{3.111330in}{3.163146in}}%
\pgfpathlineto{\pgfqpoint{3.111330in}{3.166095in}}%
\pgfpathlineto{\pgfqpoint{3.115871in}{3.166095in}}%
\pgfpathlineto{\pgfqpoint{3.115871in}{3.163146in}}%
\pgfpathmoveto{\pgfqpoint{3.115871in}{3.160197in}}%
\pgfpathlineto{\pgfqpoint{3.115871in}{3.160197in}}%
\pgfpathlineto{\pgfqpoint{3.115871in}{3.163146in}}%
\pgfpathlineto{\pgfqpoint{3.120412in}{3.163146in}}%
\pgfpathlineto{\pgfqpoint{3.120412in}{3.160197in}}%
\pgfpathmoveto{\pgfqpoint{3.115871in}{3.163146in}}%
\pgfpathlineto{\pgfqpoint{3.115871in}{3.163146in}}%
\pgfpathlineto{\pgfqpoint{3.115871in}{3.166095in}}%
\pgfpathlineto{\pgfqpoint{3.120412in}{3.166095in}}%
\pgfpathlineto{\pgfqpoint{3.120412in}{3.163146in}}%
\pgfpathmoveto{\pgfqpoint{3.120412in}{3.154298in}}%
\pgfpathlineto{\pgfqpoint{3.120412in}{3.154298in}}%
\pgfpathlineto{\pgfqpoint{3.120412in}{3.157247in}}%
\pgfpathlineto{\pgfqpoint{3.124953in}{3.157247in}}%
\pgfpathlineto{\pgfqpoint{3.124953in}{3.154298in}}%
\pgfpathmoveto{\pgfqpoint{3.120412in}{3.157247in}}%
\pgfpathlineto{\pgfqpoint{3.120412in}{3.157247in}}%
\pgfpathlineto{\pgfqpoint{3.120412in}{3.160197in}}%
\pgfpathlineto{\pgfqpoint{3.124953in}{3.160197in}}%
\pgfpathlineto{\pgfqpoint{3.124953in}{3.157247in}}%
\pgfpathmoveto{\pgfqpoint{3.129493in}{3.142502in}}%
\pgfpathlineto{\pgfqpoint{3.129493in}{3.142502in}}%
\pgfpathlineto{\pgfqpoint{3.129493in}{3.145451in}}%
\pgfpathlineto{\pgfqpoint{3.134034in}{3.145451in}}%
\pgfpathlineto{\pgfqpoint{3.134034in}{3.142502in}}%
\pgfpathmoveto{\pgfqpoint{3.129493in}{3.145451in}}%
\pgfpathlineto{\pgfqpoint{3.129493in}{3.145451in}}%
\pgfpathlineto{\pgfqpoint{3.129493in}{3.148400in}}%
\pgfpathlineto{\pgfqpoint{3.134034in}{3.148400in}}%
\pgfpathlineto{\pgfqpoint{3.134034in}{3.145451in}}%
\pgfpathmoveto{\pgfqpoint{3.111330in}{3.166095in}}%
\pgfpathlineto{\pgfqpoint{3.111330in}{3.166095in}}%
\pgfpathlineto{\pgfqpoint{3.111330in}{3.169044in}}%
\pgfpathlineto{\pgfqpoint{3.115871in}{3.169044in}}%
\pgfpathlineto{\pgfqpoint{3.115871in}{3.166095in}}%
\pgfpathmoveto{\pgfqpoint{3.111330in}{3.169044in}}%
\pgfpathlineto{\pgfqpoint{3.111330in}{3.169044in}}%
\pgfpathlineto{\pgfqpoint{3.111330in}{3.171993in}}%
\pgfpathlineto{\pgfqpoint{3.115871in}{3.171993in}}%
\pgfpathlineto{\pgfqpoint{3.115871in}{3.169044in}}%
\pgfpathmoveto{\pgfqpoint{3.084085in}{3.189688in}}%
\pgfpathlineto{\pgfqpoint{3.084085in}{3.189688in}}%
\pgfpathlineto{\pgfqpoint{3.084085in}{3.192637in}}%
\pgfpathlineto{\pgfqpoint{3.088626in}{3.192637in}}%
\pgfpathlineto{\pgfqpoint{3.088626in}{3.189688in}}%
\pgfpathmoveto{\pgfqpoint{3.084085in}{3.192637in}}%
\pgfpathlineto{\pgfqpoint{3.084085in}{3.192637in}}%
\pgfpathlineto{\pgfqpoint{3.084085in}{3.195586in}}%
\pgfpathlineto{\pgfqpoint{3.088626in}{3.195586in}}%
\pgfpathlineto{\pgfqpoint{3.088626in}{3.192637in}}%
\pgfpathmoveto{\pgfqpoint{3.088626in}{3.189688in}}%
\pgfpathlineto{\pgfqpoint{3.088626in}{3.189688in}}%
\pgfpathlineto{\pgfqpoint{3.088626in}{3.192637in}}%
\pgfpathlineto{\pgfqpoint{3.093167in}{3.192637in}}%
\pgfpathlineto{\pgfqpoint{3.093167in}{3.189688in}}%
\pgfpathmoveto{\pgfqpoint{3.088626in}{3.192637in}}%
\pgfpathlineto{\pgfqpoint{3.088626in}{3.192637in}}%
\pgfpathlineto{\pgfqpoint{3.088626in}{3.195586in}}%
\pgfpathlineto{\pgfqpoint{3.093167in}{3.195586in}}%
\pgfpathlineto{\pgfqpoint{3.093167in}{3.192637in}}%
\pgfpathmoveto{\pgfqpoint{3.084085in}{3.195586in}}%
\pgfpathlineto{\pgfqpoint{3.084085in}{3.195586in}}%
\pgfpathlineto{\pgfqpoint{3.084085in}{3.198535in}}%
\pgfpathlineto{\pgfqpoint{3.088626in}{3.198535in}}%
\pgfpathlineto{\pgfqpoint{3.088626in}{3.195586in}}%
\pgfpathmoveto{\pgfqpoint{3.084085in}{3.198535in}}%
\pgfpathlineto{\pgfqpoint{3.084085in}{3.198535in}}%
\pgfpathlineto{\pgfqpoint{3.084085in}{3.201485in}}%
\pgfpathlineto{\pgfqpoint{3.088626in}{3.201485in}}%
\pgfpathlineto{\pgfqpoint{3.088626in}{3.198535in}}%
\pgfpathmoveto{\pgfqpoint{3.088626in}{3.195586in}}%
\pgfpathlineto{\pgfqpoint{3.088626in}{3.195586in}}%
\pgfpathlineto{\pgfqpoint{3.088626in}{3.198535in}}%
\pgfpathlineto{\pgfqpoint{3.093167in}{3.198535in}}%
\pgfpathlineto{\pgfqpoint{3.093167in}{3.195586in}}%
\pgfpathmoveto{\pgfqpoint{3.088626in}{3.198535in}}%
\pgfpathlineto{\pgfqpoint{3.088626in}{3.198535in}}%
\pgfpathlineto{\pgfqpoint{3.088626in}{3.201485in}}%
\pgfpathlineto{\pgfqpoint{3.093167in}{3.201485in}}%
\pgfpathlineto{\pgfqpoint{3.093167in}{3.198535in}}%
\pgfpathmoveto{\pgfqpoint{3.075004in}{3.201485in}}%
\pgfpathlineto{\pgfqpoint{3.075004in}{3.201485in}}%
\pgfpathlineto{\pgfqpoint{3.075004in}{3.204434in}}%
\pgfpathlineto{\pgfqpoint{3.079545in}{3.204434in}}%
\pgfpathlineto{\pgfqpoint{3.079545in}{3.201485in}}%
\pgfpathmoveto{\pgfqpoint{3.075004in}{3.204434in}}%
\pgfpathlineto{\pgfqpoint{3.075004in}{3.204434in}}%
\pgfpathlineto{\pgfqpoint{3.075004in}{3.207383in}}%
\pgfpathlineto{\pgfqpoint{3.079545in}{3.207383in}}%
\pgfpathlineto{\pgfqpoint{3.079545in}{3.204434in}}%
\pgfpathmoveto{\pgfqpoint{3.079545in}{3.201485in}}%
\pgfpathlineto{\pgfqpoint{3.079545in}{3.201485in}}%
\pgfpathlineto{\pgfqpoint{3.079545in}{3.204434in}}%
\pgfpathlineto{\pgfqpoint{3.084085in}{3.204434in}}%
\pgfpathlineto{\pgfqpoint{3.084085in}{3.201485in}}%
\pgfpathmoveto{\pgfqpoint{3.079545in}{3.204434in}}%
\pgfpathlineto{\pgfqpoint{3.079545in}{3.204434in}}%
\pgfpathlineto{\pgfqpoint{3.079545in}{3.207383in}}%
\pgfpathlineto{\pgfqpoint{3.084085in}{3.207383in}}%
\pgfpathlineto{\pgfqpoint{3.084085in}{3.204434in}}%
\pgfpathmoveto{\pgfqpoint{3.075004in}{3.207383in}}%
\pgfpathlineto{\pgfqpoint{3.075004in}{3.207383in}}%
\pgfpathlineto{\pgfqpoint{3.075004in}{3.210332in}}%
\pgfpathlineto{\pgfqpoint{3.079545in}{3.210332in}}%
\pgfpathlineto{\pgfqpoint{3.079545in}{3.207383in}}%
\pgfpathmoveto{\pgfqpoint{3.075004in}{3.210332in}}%
\pgfpathlineto{\pgfqpoint{3.075004in}{3.210332in}}%
\pgfpathlineto{\pgfqpoint{3.075004in}{3.213281in}}%
\pgfpathlineto{\pgfqpoint{3.079545in}{3.213281in}}%
\pgfpathlineto{\pgfqpoint{3.079545in}{3.210332in}}%
\pgfpathmoveto{\pgfqpoint{3.079545in}{3.207383in}}%
\pgfpathlineto{\pgfqpoint{3.079545in}{3.207383in}}%
\pgfpathlineto{\pgfqpoint{3.079545in}{3.210332in}}%
\pgfpathlineto{\pgfqpoint{3.084085in}{3.210332in}}%
\pgfpathlineto{\pgfqpoint{3.084085in}{3.207383in}}%
\pgfpathmoveto{\pgfqpoint{3.079545in}{3.210332in}}%
\pgfpathlineto{\pgfqpoint{3.079545in}{3.210332in}}%
\pgfpathlineto{\pgfqpoint{3.079545in}{3.213281in}}%
\pgfpathlineto{\pgfqpoint{3.084085in}{3.213281in}}%
\pgfpathlineto{\pgfqpoint{3.084085in}{3.210332in}}%
\pgfpathmoveto{\pgfqpoint{3.084085in}{3.201485in}}%
\pgfpathlineto{\pgfqpoint{3.084085in}{3.201485in}}%
\pgfpathlineto{\pgfqpoint{3.084085in}{3.204434in}}%
\pgfpathlineto{\pgfqpoint{3.088626in}{3.204434in}}%
\pgfpathlineto{\pgfqpoint{3.088626in}{3.201485in}}%
\pgfpathmoveto{\pgfqpoint{3.084085in}{3.204434in}}%
\pgfpathlineto{\pgfqpoint{3.084085in}{3.204434in}}%
\pgfpathlineto{\pgfqpoint{3.084085in}{3.207383in}}%
\pgfpathlineto{\pgfqpoint{3.088626in}{3.207383in}}%
\pgfpathlineto{\pgfqpoint{3.088626in}{3.204434in}}%
\pgfpathmoveto{\pgfqpoint{3.093167in}{3.189688in}}%
\pgfpathlineto{\pgfqpoint{3.093167in}{3.189688in}}%
\pgfpathlineto{\pgfqpoint{3.093167in}{3.192637in}}%
\pgfpathlineto{\pgfqpoint{3.097708in}{3.192637in}}%
\pgfpathlineto{\pgfqpoint{3.097708in}{3.189688in}}%
\pgfpathmoveto{\pgfqpoint{3.093167in}{3.192637in}}%
\pgfpathlineto{\pgfqpoint{3.093167in}{3.192637in}}%
\pgfpathlineto{\pgfqpoint{3.093167in}{3.195586in}}%
\pgfpathlineto{\pgfqpoint{3.097708in}{3.195586in}}%
\pgfpathlineto{\pgfqpoint{3.097708in}{3.192637in}}%
\pgfpathmoveto{\pgfqpoint{3.075004in}{3.213281in}}%
\pgfpathlineto{\pgfqpoint{3.075004in}{3.213281in}}%
\pgfpathlineto{\pgfqpoint{3.075004in}{3.216230in}}%
\pgfpathlineto{\pgfqpoint{3.079545in}{3.216230in}}%
\pgfpathlineto{\pgfqpoint{3.079545in}{3.213281in}}%
\pgfpathmoveto{\pgfqpoint{3.075004in}{3.216230in}}%
\pgfpathlineto{\pgfqpoint{3.075004in}{3.216230in}}%
\pgfpathlineto{\pgfqpoint{3.075004in}{3.219180in}}%
\pgfpathlineto{\pgfqpoint{3.079545in}{3.219180in}}%
\pgfpathlineto{\pgfqpoint{3.079545in}{3.216230in}}%
\pgfpathmoveto{\pgfqpoint{3.220310in}{2.101428in}}%
\pgfpathlineto{\pgfqpoint{3.220310in}{2.101428in}}%
\pgfpathlineto{\pgfqpoint{3.220310in}{2.104377in}}%
\pgfpathlineto{\pgfqpoint{3.224851in}{2.104377in}}%
\pgfpathlineto{\pgfqpoint{3.224851in}{2.101428in}}%
\pgfpathmoveto{\pgfqpoint{3.220310in}{2.104377in}}%
\pgfpathlineto{\pgfqpoint{3.220310in}{2.104377in}}%
\pgfpathlineto{\pgfqpoint{3.220310in}{2.107326in}}%
\pgfpathlineto{\pgfqpoint{3.224851in}{2.107326in}}%
\pgfpathlineto{\pgfqpoint{3.224851in}{2.104377in}}%
\pgfpathmoveto{\pgfqpoint{3.220310in}{2.107326in}}%
\pgfpathlineto{\pgfqpoint{3.220310in}{2.107326in}}%
\pgfpathlineto{\pgfqpoint{3.220310in}{2.110275in}}%
\pgfpathlineto{\pgfqpoint{3.224851in}{2.110275in}}%
\pgfpathlineto{\pgfqpoint{3.224851in}{2.107326in}}%
\pgfpathmoveto{\pgfqpoint{3.224851in}{2.104377in}}%
\pgfpathlineto{\pgfqpoint{3.224851in}{2.104377in}}%
\pgfpathlineto{\pgfqpoint{3.224851in}{2.107326in}}%
\pgfpathlineto{\pgfqpoint{3.229392in}{2.107326in}}%
\pgfpathlineto{\pgfqpoint{3.229392in}{2.104377in}}%
\pgfpathmoveto{\pgfqpoint{3.224851in}{2.107326in}}%
\pgfpathlineto{\pgfqpoint{3.224851in}{2.107326in}}%
\pgfpathlineto{\pgfqpoint{3.224851in}{2.110275in}}%
\pgfpathlineto{\pgfqpoint{3.229392in}{2.110275in}}%
\pgfpathlineto{\pgfqpoint{3.229392in}{2.107326in}}%
\pgfpathmoveto{\pgfqpoint{3.229392in}{2.107326in}}%
\pgfpathlineto{\pgfqpoint{3.229392in}{2.107326in}}%
\pgfpathlineto{\pgfqpoint{3.229392in}{2.110275in}}%
\pgfpathlineto{\pgfqpoint{3.233933in}{2.110275in}}%
\pgfpathlineto{\pgfqpoint{3.233933in}{2.107326in}}%
\pgfpathmoveto{\pgfqpoint{3.229392in}{2.110275in}}%
\pgfpathlineto{\pgfqpoint{3.229392in}{2.110275in}}%
\pgfpathlineto{\pgfqpoint{3.229392in}{2.113224in}}%
\pgfpathlineto{\pgfqpoint{3.233933in}{2.113224in}}%
\pgfpathlineto{\pgfqpoint{3.233933in}{2.110275in}}%
\pgfpathmoveto{\pgfqpoint{3.229392in}{2.113224in}}%
\pgfpathlineto{\pgfqpoint{3.229392in}{2.113224in}}%
\pgfpathlineto{\pgfqpoint{3.229392in}{2.116174in}}%
\pgfpathlineto{\pgfqpoint{3.233933in}{2.116174in}}%
\pgfpathlineto{\pgfqpoint{3.233933in}{2.113224in}}%
\pgfpathmoveto{\pgfqpoint{3.233933in}{2.110275in}}%
\pgfpathlineto{\pgfqpoint{3.233933in}{2.110275in}}%
\pgfpathlineto{\pgfqpoint{3.233933in}{2.113224in}}%
\pgfpathlineto{\pgfqpoint{3.238474in}{2.113224in}}%
\pgfpathlineto{\pgfqpoint{3.238474in}{2.110275in}}%
\pgfpathmoveto{\pgfqpoint{3.233933in}{2.113224in}}%
\pgfpathlineto{\pgfqpoint{3.233933in}{2.113224in}}%
\pgfpathlineto{\pgfqpoint{3.233933in}{2.116174in}}%
\pgfpathlineto{\pgfqpoint{3.238474in}{2.116174in}}%
\pgfpathlineto{\pgfqpoint{3.238474in}{2.113224in}}%
\pgfpathmoveto{\pgfqpoint{3.238474in}{2.113224in}}%
\pgfpathlineto{\pgfqpoint{3.238474in}{2.113224in}}%
\pgfpathlineto{\pgfqpoint{3.238474in}{2.116174in}}%
\pgfpathlineto{\pgfqpoint{3.243015in}{2.116174in}}%
\pgfpathlineto{\pgfqpoint{3.243015in}{2.113224in}}%
\pgfpathmoveto{\pgfqpoint{3.238474in}{2.116174in}}%
\pgfpathlineto{\pgfqpoint{3.238474in}{2.116174in}}%
\pgfpathlineto{\pgfqpoint{3.238474in}{2.119123in}}%
\pgfpathlineto{\pgfqpoint{3.243015in}{2.119123in}}%
\pgfpathlineto{\pgfqpoint{3.243015in}{2.116174in}}%
\pgfpathmoveto{\pgfqpoint{3.238474in}{2.119123in}}%
\pgfpathlineto{\pgfqpoint{3.238474in}{2.119123in}}%
\pgfpathlineto{\pgfqpoint{3.238474in}{2.122072in}}%
\pgfpathlineto{\pgfqpoint{3.243015in}{2.122072in}}%
\pgfpathlineto{\pgfqpoint{3.243015in}{2.119123in}}%
\pgfpathmoveto{\pgfqpoint{3.243015in}{2.116174in}}%
\pgfpathlineto{\pgfqpoint{3.243015in}{2.116174in}}%
\pgfpathlineto{\pgfqpoint{3.243015in}{2.119123in}}%
\pgfpathlineto{\pgfqpoint{3.247557in}{2.119123in}}%
\pgfpathlineto{\pgfqpoint{3.247557in}{2.116174in}}%
\pgfpathmoveto{\pgfqpoint{3.243015in}{2.119123in}}%
\pgfpathlineto{\pgfqpoint{3.243015in}{2.119123in}}%
\pgfpathlineto{\pgfqpoint{3.243015in}{2.122072in}}%
\pgfpathlineto{\pgfqpoint{3.247557in}{2.122072in}}%
\pgfpathlineto{\pgfqpoint{3.247557in}{2.119123in}}%
\pgfpathmoveto{\pgfqpoint{3.247557in}{2.119123in}}%
\pgfpathlineto{\pgfqpoint{3.247557in}{2.119123in}}%
\pgfpathlineto{\pgfqpoint{3.247557in}{2.122072in}}%
\pgfpathlineto{\pgfqpoint{3.252098in}{2.122072in}}%
\pgfpathlineto{\pgfqpoint{3.252098in}{2.119123in}}%
\pgfpathmoveto{\pgfqpoint{3.247557in}{2.122072in}}%
\pgfpathlineto{\pgfqpoint{3.247557in}{2.122072in}}%
\pgfpathlineto{\pgfqpoint{3.247557in}{2.125021in}}%
\pgfpathlineto{\pgfqpoint{3.252098in}{2.125021in}}%
\pgfpathlineto{\pgfqpoint{3.252098in}{2.122072in}}%
\pgfpathmoveto{\pgfqpoint{3.247557in}{2.125021in}}%
\pgfpathlineto{\pgfqpoint{3.247557in}{2.125021in}}%
\pgfpathlineto{\pgfqpoint{3.247557in}{2.127970in}}%
\pgfpathlineto{\pgfqpoint{3.252098in}{2.127970in}}%
\pgfpathlineto{\pgfqpoint{3.252098in}{2.125021in}}%
\pgfpathmoveto{\pgfqpoint{3.252098in}{2.122072in}}%
\pgfpathlineto{\pgfqpoint{3.252098in}{2.122072in}}%
\pgfpathlineto{\pgfqpoint{3.252098in}{2.125021in}}%
\pgfpathlineto{\pgfqpoint{3.256639in}{2.125021in}}%
\pgfpathlineto{\pgfqpoint{3.256639in}{2.122072in}}%
\pgfpathmoveto{\pgfqpoint{3.252098in}{2.125021in}}%
\pgfpathlineto{\pgfqpoint{3.252098in}{2.125021in}}%
\pgfpathlineto{\pgfqpoint{3.252098in}{2.127970in}}%
\pgfpathlineto{\pgfqpoint{3.256639in}{2.127970in}}%
\pgfpathlineto{\pgfqpoint{3.256639in}{2.125021in}}%
\pgfpathmoveto{\pgfqpoint{3.256639in}{2.125021in}}%
\pgfpathlineto{\pgfqpoint{3.256639in}{2.125021in}}%
\pgfpathlineto{\pgfqpoint{3.256639in}{2.127970in}}%
\pgfpathlineto{\pgfqpoint{3.261180in}{2.127970in}}%
\pgfpathlineto{\pgfqpoint{3.261180in}{2.125021in}}%
\pgfpathmoveto{\pgfqpoint{3.256639in}{2.127970in}}%
\pgfpathlineto{\pgfqpoint{3.256639in}{2.127970in}}%
\pgfpathlineto{\pgfqpoint{3.256639in}{2.130919in}}%
\pgfpathlineto{\pgfqpoint{3.261180in}{2.130919in}}%
\pgfpathlineto{\pgfqpoint{3.261180in}{2.127970in}}%
\pgfpathmoveto{\pgfqpoint{3.256639in}{2.130919in}}%
\pgfpathlineto{\pgfqpoint{3.256639in}{2.130919in}}%
\pgfpathlineto{\pgfqpoint{3.256639in}{2.133868in}}%
\pgfpathlineto{\pgfqpoint{3.261180in}{2.133868in}}%
\pgfpathlineto{\pgfqpoint{3.261180in}{2.130919in}}%
\pgfpathmoveto{\pgfqpoint{3.261180in}{2.127970in}}%
\pgfpathlineto{\pgfqpoint{3.261180in}{2.127970in}}%
\pgfpathlineto{\pgfqpoint{3.261180in}{2.130919in}}%
\pgfpathlineto{\pgfqpoint{3.265721in}{2.130919in}}%
\pgfpathlineto{\pgfqpoint{3.265721in}{2.127970in}}%
\pgfpathmoveto{\pgfqpoint{3.261180in}{2.130919in}}%
\pgfpathlineto{\pgfqpoint{3.261180in}{2.130919in}}%
\pgfpathlineto{\pgfqpoint{3.261180in}{2.133868in}}%
\pgfpathlineto{\pgfqpoint{3.265721in}{2.133868in}}%
\pgfpathlineto{\pgfqpoint{3.265721in}{2.130919in}}%
\pgfpathmoveto{\pgfqpoint{3.265721in}{2.130919in}}%
\pgfpathlineto{\pgfqpoint{3.265721in}{2.130919in}}%
\pgfpathlineto{\pgfqpoint{3.265721in}{2.133868in}}%
\pgfpathlineto{\pgfqpoint{3.270262in}{2.133868in}}%
\pgfpathlineto{\pgfqpoint{3.270262in}{2.130919in}}%
\pgfpathmoveto{\pgfqpoint{3.265721in}{2.133868in}}%
\pgfpathlineto{\pgfqpoint{3.265721in}{2.133868in}}%
\pgfpathlineto{\pgfqpoint{3.265721in}{2.136817in}}%
\pgfpathlineto{\pgfqpoint{3.270262in}{2.136817in}}%
\pgfpathlineto{\pgfqpoint{3.270262in}{2.133868in}}%
\pgfpathmoveto{\pgfqpoint{3.265721in}{2.136817in}}%
\pgfpathlineto{\pgfqpoint{3.265721in}{2.136817in}}%
\pgfpathlineto{\pgfqpoint{3.265721in}{2.139766in}}%
\pgfpathlineto{\pgfqpoint{3.270262in}{2.139766in}}%
\pgfpathlineto{\pgfqpoint{3.270262in}{2.136817in}}%
\pgfpathmoveto{\pgfqpoint{3.270262in}{2.133868in}}%
\pgfpathlineto{\pgfqpoint{3.270262in}{2.133868in}}%
\pgfpathlineto{\pgfqpoint{3.270262in}{2.136817in}}%
\pgfpathlineto{\pgfqpoint{3.274803in}{2.136817in}}%
\pgfpathlineto{\pgfqpoint{3.274803in}{2.133868in}}%
\pgfpathmoveto{\pgfqpoint{3.270262in}{2.136817in}}%
\pgfpathlineto{\pgfqpoint{3.270262in}{2.136817in}}%
\pgfpathlineto{\pgfqpoint{3.270262in}{2.139766in}}%
\pgfpathlineto{\pgfqpoint{3.274803in}{2.139766in}}%
\pgfpathlineto{\pgfqpoint{3.274803in}{2.136817in}}%
\pgfpathmoveto{\pgfqpoint{3.274803in}{2.136817in}}%
\pgfpathlineto{\pgfqpoint{3.274803in}{2.136817in}}%
\pgfpathlineto{\pgfqpoint{3.274803in}{2.139766in}}%
\pgfpathlineto{\pgfqpoint{3.279345in}{2.139766in}}%
\pgfpathlineto{\pgfqpoint{3.279345in}{2.136817in}}%
\pgfpathmoveto{\pgfqpoint{3.274803in}{2.139766in}}%
\pgfpathlineto{\pgfqpoint{3.274803in}{2.139766in}}%
\pgfpathlineto{\pgfqpoint{3.274803in}{2.142715in}}%
\pgfpathlineto{\pgfqpoint{3.279345in}{2.142715in}}%
\pgfpathlineto{\pgfqpoint{3.279345in}{2.139766in}}%
\pgfpathmoveto{\pgfqpoint{3.274803in}{2.142715in}}%
\pgfpathlineto{\pgfqpoint{3.274803in}{2.142715in}}%
\pgfpathlineto{\pgfqpoint{3.274803in}{2.145665in}}%
\pgfpathlineto{\pgfqpoint{3.279345in}{2.145665in}}%
\pgfpathlineto{\pgfqpoint{3.279345in}{2.142715in}}%
\pgfpathmoveto{\pgfqpoint{3.279345in}{2.139766in}}%
\pgfpathlineto{\pgfqpoint{3.279345in}{2.139766in}}%
\pgfpathlineto{\pgfqpoint{3.279345in}{2.142715in}}%
\pgfpathlineto{\pgfqpoint{3.283886in}{2.142715in}}%
\pgfpathlineto{\pgfqpoint{3.283886in}{2.139766in}}%
\pgfpathmoveto{\pgfqpoint{3.279345in}{2.142715in}}%
\pgfpathlineto{\pgfqpoint{3.279345in}{2.142715in}}%
\pgfpathlineto{\pgfqpoint{3.279345in}{2.145665in}}%
\pgfpathlineto{\pgfqpoint{3.283886in}{2.145665in}}%
\pgfpathlineto{\pgfqpoint{3.283886in}{2.142715in}}%
\pgfpathmoveto{\pgfqpoint{3.283886in}{2.142715in}}%
\pgfpathlineto{\pgfqpoint{3.283886in}{2.142715in}}%
\pgfpathlineto{\pgfqpoint{3.283886in}{2.145665in}}%
\pgfpathlineto{\pgfqpoint{3.288427in}{2.145665in}}%
\pgfpathlineto{\pgfqpoint{3.288427in}{2.142715in}}%
\pgfpathmoveto{\pgfqpoint{3.283886in}{2.145665in}}%
\pgfpathlineto{\pgfqpoint{3.283886in}{2.145665in}}%
\pgfpathlineto{\pgfqpoint{3.283886in}{2.148614in}}%
\pgfpathlineto{\pgfqpoint{3.288427in}{2.148614in}}%
\pgfpathlineto{\pgfqpoint{3.288427in}{2.145665in}}%
\pgfpathmoveto{\pgfqpoint{3.283886in}{2.148614in}}%
\pgfpathlineto{\pgfqpoint{3.283886in}{2.148614in}}%
\pgfpathlineto{\pgfqpoint{3.283886in}{2.151563in}}%
\pgfpathlineto{\pgfqpoint{3.288427in}{2.151563in}}%
\pgfpathlineto{\pgfqpoint{3.288427in}{2.148614in}}%
\pgfpathmoveto{\pgfqpoint{3.288427in}{2.145665in}}%
\pgfpathlineto{\pgfqpoint{3.288427in}{2.145665in}}%
\pgfpathlineto{\pgfqpoint{3.288427in}{2.148614in}}%
\pgfpathlineto{\pgfqpoint{3.292968in}{2.148614in}}%
\pgfpathlineto{\pgfqpoint{3.292968in}{2.145665in}}%
\pgfpathmoveto{\pgfqpoint{3.288427in}{2.148614in}}%
\pgfpathlineto{\pgfqpoint{3.288427in}{2.148614in}}%
\pgfpathlineto{\pgfqpoint{3.288427in}{2.151563in}}%
\pgfpathlineto{\pgfqpoint{3.292968in}{2.151563in}}%
\pgfpathlineto{\pgfqpoint{3.292968in}{2.148614in}}%
\pgfpathmoveto{\pgfqpoint{3.292968in}{2.148614in}}%
\pgfpathlineto{\pgfqpoint{3.292968in}{2.148614in}}%
\pgfpathlineto{\pgfqpoint{3.292968in}{2.151563in}}%
\pgfpathlineto{\pgfqpoint{3.297509in}{2.151563in}}%
\pgfpathlineto{\pgfqpoint{3.297509in}{2.148614in}}%
\pgfpathmoveto{\pgfqpoint{3.292968in}{2.151563in}}%
\pgfpathlineto{\pgfqpoint{3.292968in}{2.151563in}}%
\pgfpathlineto{\pgfqpoint{3.292968in}{2.154512in}}%
\pgfpathlineto{\pgfqpoint{3.297509in}{2.154512in}}%
\pgfpathlineto{\pgfqpoint{3.297509in}{2.151563in}}%
\pgfpathmoveto{\pgfqpoint{3.292968in}{2.154512in}}%
\pgfpathlineto{\pgfqpoint{3.292968in}{2.154512in}}%
\pgfpathlineto{\pgfqpoint{3.292968in}{2.157461in}}%
\pgfpathlineto{\pgfqpoint{3.297509in}{2.157461in}}%
\pgfpathlineto{\pgfqpoint{3.297509in}{2.154512in}}%
\pgfpathmoveto{\pgfqpoint{3.297509in}{2.151563in}}%
\pgfpathlineto{\pgfqpoint{3.297509in}{2.151563in}}%
\pgfpathlineto{\pgfqpoint{3.297509in}{2.154512in}}%
\pgfpathlineto{\pgfqpoint{3.302050in}{2.154512in}}%
\pgfpathlineto{\pgfqpoint{3.302050in}{2.151563in}}%
\pgfpathmoveto{\pgfqpoint{3.297509in}{2.154512in}}%
\pgfpathlineto{\pgfqpoint{3.297509in}{2.154512in}}%
\pgfpathlineto{\pgfqpoint{3.297509in}{2.157461in}}%
\pgfpathlineto{\pgfqpoint{3.302050in}{2.157461in}}%
\pgfpathlineto{\pgfqpoint{3.302050in}{2.154512in}}%
\pgfpathmoveto{\pgfqpoint{3.302050in}{2.154512in}}%
\pgfpathlineto{\pgfqpoint{3.302050in}{2.154512in}}%
\pgfpathlineto{\pgfqpoint{3.302050in}{2.157461in}}%
\pgfpathlineto{\pgfqpoint{3.306592in}{2.157461in}}%
\pgfpathlineto{\pgfqpoint{3.306592in}{2.154512in}}%
\pgfpathmoveto{\pgfqpoint{3.302050in}{2.157461in}}%
\pgfpathlineto{\pgfqpoint{3.302050in}{2.157461in}}%
\pgfpathlineto{\pgfqpoint{3.302050in}{2.160410in}}%
\pgfpathlineto{\pgfqpoint{3.306592in}{2.160410in}}%
\pgfpathlineto{\pgfqpoint{3.306592in}{2.157461in}}%
\pgfpathmoveto{\pgfqpoint{3.302050in}{2.160410in}}%
\pgfpathlineto{\pgfqpoint{3.302050in}{2.160410in}}%
\pgfpathlineto{\pgfqpoint{3.302050in}{2.163359in}}%
\pgfpathlineto{\pgfqpoint{3.306592in}{2.163359in}}%
\pgfpathlineto{\pgfqpoint{3.306592in}{2.160410in}}%
\pgfpathmoveto{\pgfqpoint{3.306592in}{2.160410in}}%
\pgfpathlineto{\pgfqpoint{3.306592in}{2.160410in}}%
\pgfpathlineto{\pgfqpoint{3.306592in}{2.163359in}}%
\pgfpathlineto{\pgfqpoint{3.311133in}{2.163359in}}%
\pgfpathlineto{\pgfqpoint{3.311133in}{2.160410in}}%
\pgfpathmoveto{\pgfqpoint{3.302050in}{2.163359in}}%
\pgfpathlineto{\pgfqpoint{3.302050in}{2.163359in}}%
\pgfpathlineto{\pgfqpoint{3.302050in}{2.166308in}}%
\pgfpathlineto{\pgfqpoint{3.306592in}{2.166308in}}%
\pgfpathlineto{\pgfqpoint{3.306592in}{2.163359in}}%
\pgfpathmoveto{\pgfqpoint{3.302050in}{2.166308in}}%
\pgfpathlineto{\pgfqpoint{3.302050in}{2.166308in}}%
\pgfpathlineto{\pgfqpoint{3.302050in}{2.169257in}}%
\pgfpathlineto{\pgfqpoint{3.306592in}{2.169257in}}%
\pgfpathlineto{\pgfqpoint{3.306592in}{2.166308in}}%
\pgfpathmoveto{\pgfqpoint{3.306592in}{2.163359in}}%
\pgfpathlineto{\pgfqpoint{3.306592in}{2.163359in}}%
\pgfpathlineto{\pgfqpoint{3.306592in}{2.166308in}}%
\pgfpathlineto{\pgfqpoint{3.311133in}{2.166308in}}%
\pgfpathlineto{\pgfqpoint{3.311133in}{2.163359in}}%
\pgfpathmoveto{\pgfqpoint{3.306592in}{2.166308in}}%
\pgfpathlineto{\pgfqpoint{3.306592in}{2.166308in}}%
\pgfpathlineto{\pgfqpoint{3.306592in}{2.169257in}}%
\pgfpathlineto{\pgfqpoint{3.311133in}{2.169257in}}%
\pgfpathlineto{\pgfqpoint{3.311133in}{2.166308in}}%
\pgfpathmoveto{\pgfqpoint{3.311133in}{2.163359in}}%
\pgfpathlineto{\pgfqpoint{3.311133in}{2.163359in}}%
\pgfpathlineto{\pgfqpoint{3.311133in}{2.166308in}}%
\pgfpathlineto{\pgfqpoint{3.315674in}{2.166308in}}%
\pgfpathlineto{\pgfqpoint{3.315674in}{2.163359in}}%
\pgfpathmoveto{\pgfqpoint{3.311133in}{2.166308in}}%
\pgfpathlineto{\pgfqpoint{3.311133in}{2.166308in}}%
\pgfpathlineto{\pgfqpoint{3.311133in}{2.169257in}}%
\pgfpathlineto{\pgfqpoint{3.315674in}{2.169257in}}%
\pgfpathlineto{\pgfqpoint{3.315674in}{2.166308in}}%
\pgfpathmoveto{\pgfqpoint{3.315674in}{2.166308in}}%
\pgfpathlineto{\pgfqpoint{3.315674in}{2.166308in}}%
\pgfpathlineto{\pgfqpoint{3.315674in}{2.169257in}}%
\pgfpathlineto{\pgfqpoint{3.320215in}{2.169257in}}%
\pgfpathlineto{\pgfqpoint{3.320215in}{2.166308in}}%
\pgfpathmoveto{\pgfqpoint{3.311133in}{2.169257in}}%
\pgfpathlineto{\pgfqpoint{3.311133in}{2.169257in}}%
\pgfpathlineto{\pgfqpoint{3.311133in}{2.172206in}}%
\pgfpathlineto{\pgfqpoint{3.315674in}{2.172206in}}%
\pgfpathlineto{\pgfqpoint{3.315674in}{2.169257in}}%
\pgfpathmoveto{\pgfqpoint{3.311133in}{2.172206in}}%
\pgfpathlineto{\pgfqpoint{3.311133in}{2.172206in}}%
\pgfpathlineto{\pgfqpoint{3.311133in}{2.175155in}}%
\pgfpathlineto{\pgfqpoint{3.315674in}{2.175155in}}%
\pgfpathlineto{\pgfqpoint{3.315674in}{2.172206in}}%
\pgfpathmoveto{\pgfqpoint{3.315674in}{2.169257in}}%
\pgfpathlineto{\pgfqpoint{3.315674in}{2.169257in}}%
\pgfpathlineto{\pgfqpoint{3.315674in}{2.172206in}}%
\pgfpathlineto{\pgfqpoint{3.320215in}{2.172206in}}%
\pgfpathlineto{\pgfqpoint{3.320215in}{2.169257in}}%
\pgfpathmoveto{\pgfqpoint{3.315674in}{2.172206in}}%
\pgfpathlineto{\pgfqpoint{3.315674in}{2.172206in}}%
\pgfpathlineto{\pgfqpoint{3.315674in}{2.175155in}}%
\pgfpathlineto{\pgfqpoint{3.320215in}{2.175155in}}%
\pgfpathlineto{\pgfqpoint{3.320215in}{2.172206in}}%
\pgfpathmoveto{\pgfqpoint{3.320215in}{2.169257in}}%
\pgfpathlineto{\pgfqpoint{3.320215in}{2.169257in}}%
\pgfpathlineto{\pgfqpoint{3.320215in}{2.172206in}}%
\pgfpathlineto{\pgfqpoint{3.324756in}{2.172206in}}%
\pgfpathlineto{\pgfqpoint{3.324756in}{2.169257in}}%
\pgfpathmoveto{\pgfqpoint{3.320215in}{2.172206in}}%
\pgfpathlineto{\pgfqpoint{3.320215in}{2.172206in}}%
\pgfpathlineto{\pgfqpoint{3.320215in}{2.175155in}}%
\pgfpathlineto{\pgfqpoint{3.324756in}{2.175155in}}%
\pgfpathlineto{\pgfqpoint{3.324756in}{2.172206in}}%
\pgfpathmoveto{\pgfqpoint{3.324756in}{2.172206in}}%
\pgfpathlineto{\pgfqpoint{3.324756in}{2.172206in}}%
\pgfpathlineto{\pgfqpoint{3.324756in}{2.175155in}}%
\pgfpathlineto{\pgfqpoint{3.329297in}{2.175155in}}%
\pgfpathlineto{\pgfqpoint{3.329297in}{2.172206in}}%
\pgfpathmoveto{\pgfqpoint{3.320215in}{2.175155in}}%
\pgfpathlineto{\pgfqpoint{3.320215in}{2.175155in}}%
\pgfpathlineto{\pgfqpoint{3.320215in}{2.178105in}}%
\pgfpathlineto{\pgfqpoint{3.324756in}{2.178105in}}%
\pgfpathlineto{\pgfqpoint{3.324756in}{2.175155in}}%
\pgfpathmoveto{\pgfqpoint{3.320215in}{2.178105in}}%
\pgfpathlineto{\pgfqpoint{3.320215in}{2.178105in}}%
\pgfpathlineto{\pgfqpoint{3.320215in}{2.181054in}}%
\pgfpathlineto{\pgfqpoint{3.324756in}{2.181054in}}%
\pgfpathlineto{\pgfqpoint{3.324756in}{2.178105in}}%
\pgfpathmoveto{\pgfqpoint{3.324756in}{2.175155in}}%
\pgfpathlineto{\pgfqpoint{3.324756in}{2.175155in}}%
\pgfpathlineto{\pgfqpoint{3.324756in}{2.178105in}}%
\pgfpathlineto{\pgfqpoint{3.329297in}{2.178105in}}%
\pgfpathlineto{\pgfqpoint{3.329297in}{2.175155in}}%
\pgfpathmoveto{\pgfqpoint{3.324756in}{2.178105in}}%
\pgfpathlineto{\pgfqpoint{3.324756in}{2.178105in}}%
\pgfpathlineto{\pgfqpoint{3.324756in}{2.181054in}}%
\pgfpathlineto{\pgfqpoint{3.329297in}{2.181054in}}%
\pgfpathlineto{\pgfqpoint{3.329297in}{2.178105in}}%
\pgfpathmoveto{\pgfqpoint{3.329297in}{2.175155in}}%
\pgfpathlineto{\pgfqpoint{3.329297in}{2.175155in}}%
\pgfpathlineto{\pgfqpoint{3.329297in}{2.178105in}}%
\pgfpathlineto{\pgfqpoint{3.333839in}{2.178105in}}%
\pgfpathlineto{\pgfqpoint{3.333839in}{2.175155in}}%
\pgfpathmoveto{\pgfqpoint{3.329297in}{2.178105in}}%
\pgfpathlineto{\pgfqpoint{3.329297in}{2.178105in}}%
\pgfpathlineto{\pgfqpoint{3.329297in}{2.181054in}}%
\pgfpathlineto{\pgfqpoint{3.333839in}{2.181054in}}%
\pgfpathlineto{\pgfqpoint{3.333839in}{2.178105in}}%
\pgfpathmoveto{\pgfqpoint{3.333839in}{2.178105in}}%
\pgfpathlineto{\pgfqpoint{3.333839in}{2.178105in}}%
\pgfpathlineto{\pgfqpoint{3.333839in}{2.181054in}}%
\pgfpathlineto{\pgfqpoint{3.338380in}{2.181054in}}%
\pgfpathlineto{\pgfqpoint{3.338380in}{2.178105in}}%
\pgfpathmoveto{\pgfqpoint{3.329297in}{2.181054in}}%
\pgfpathlineto{\pgfqpoint{3.329297in}{2.181054in}}%
\pgfpathlineto{\pgfqpoint{3.329297in}{2.184003in}}%
\pgfpathlineto{\pgfqpoint{3.333839in}{2.184003in}}%
\pgfpathlineto{\pgfqpoint{3.333839in}{2.181054in}}%
\pgfpathmoveto{\pgfqpoint{3.329297in}{2.184003in}}%
\pgfpathlineto{\pgfqpoint{3.329297in}{2.184003in}}%
\pgfpathlineto{\pgfqpoint{3.329297in}{2.186952in}}%
\pgfpathlineto{\pgfqpoint{3.333839in}{2.186952in}}%
\pgfpathlineto{\pgfqpoint{3.333839in}{2.184003in}}%
\pgfpathmoveto{\pgfqpoint{3.333839in}{2.181054in}}%
\pgfpathlineto{\pgfqpoint{3.333839in}{2.181054in}}%
\pgfpathlineto{\pgfqpoint{3.333839in}{2.184003in}}%
\pgfpathlineto{\pgfqpoint{3.338380in}{2.184003in}}%
\pgfpathlineto{\pgfqpoint{3.338380in}{2.181054in}}%
\pgfpathmoveto{\pgfqpoint{3.333839in}{2.184003in}}%
\pgfpathlineto{\pgfqpoint{3.333839in}{2.184003in}}%
\pgfpathlineto{\pgfqpoint{3.333839in}{2.186952in}}%
\pgfpathlineto{\pgfqpoint{3.338380in}{2.186952in}}%
\pgfpathlineto{\pgfqpoint{3.338380in}{2.184003in}}%
\pgfpathmoveto{\pgfqpoint{3.338380in}{2.181054in}}%
\pgfpathlineto{\pgfqpoint{3.338380in}{2.181054in}}%
\pgfpathlineto{\pgfqpoint{3.338380in}{2.184003in}}%
\pgfpathlineto{\pgfqpoint{3.342921in}{2.184003in}}%
\pgfpathlineto{\pgfqpoint{3.342921in}{2.181054in}}%
\pgfpathmoveto{\pgfqpoint{3.338380in}{2.184003in}}%
\pgfpathlineto{\pgfqpoint{3.338380in}{2.184003in}}%
\pgfpathlineto{\pgfqpoint{3.338380in}{2.186952in}}%
\pgfpathlineto{\pgfqpoint{3.342921in}{2.186952in}}%
\pgfpathlineto{\pgfqpoint{3.342921in}{2.184003in}}%
\pgfpathmoveto{\pgfqpoint{3.342921in}{2.184003in}}%
\pgfpathlineto{\pgfqpoint{3.342921in}{2.184003in}}%
\pgfpathlineto{\pgfqpoint{3.342921in}{2.186952in}}%
\pgfpathlineto{\pgfqpoint{3.347462in}{2.186952in}}%
\pgfpathlineto{\pgfqpoint{3.347462in}{2.184003in}}%
\pgfpathmoveto{\pgfqpoint{3.338380in}{2.186952in}}%
\pgfpathlineto{\pgfqpoint{3.338380in}{2.186952in}}%
\pgfpathlineto{\pgfqpoint{3.338380in}{2.189901in}}%
\pgfpathlineto{\pgfqpoint{3.342921in}{2.189901in}}%
\pgfpathlineto{\pgfqpoint{3.342921in}{2.186952in}}%
\pgfpathmoveto{\pgfqpoint{3.338380in}{2.189901in}}%
\pgfpathlineto{\pgfqpoint{3.338380in}{2.189901in}}%
\pgfpathlineto{\pgfqpoint{3.338380in}{2.192850in}}%
\pgfpathlineto{\pgfqpoint{3.342921in}{2.192850in}}%
\pgfpathlineto{\pgfqpoint{3.342921in}{2.189901in}}%
\pgfpathmoveto{\pgfqpoint{3.342921in}{2.186952in}}%
\pgfpathlineto{\pgfqpoint{3.342921in}{2.186952in}}%
\pgfpathlineto{\pgfqpoint{3.342921in}{2.189901in}}%
\pgfpathlineto{\pgfqpoint{3.347462in}{2.189901in}}%
\pgfpathlineto{\pgfqpoint{3.347462in}{2.186952in}}%
\pgfpathmoveto{\pgfqpoint{3.342921in}{2.189901in}}%
\pgfpathlineto{\pgfqpoint{3.342921in}{2.189901in}}%
\pgfpathlineto{\pgfqpoint{3.342921in}{2.192850in}}%
\pgfpathlineto{\pgfqpoint{3.347462in}{2.192850in}}%
\pgfpathlineto{\pgfqpoint{3.347462in}{2.189901in}}%
\pgfpathmoveto{\pgfqpoint{3.347462in}{2.186952in}}%
\pgfpathlineto{\pgfqpoint{3.347462in}{2.186952in}}%
\pgfpathlineto{\pgfqpoint{3.347462in}{2.189901in}}%
\pgfpathlineto{\pgfqpoint{3.352003in}{2.189901in}}%
\pgfpathlineto{\pgfqpoint{3.352003in}{2.186952in}}%
\pgfpathmoveto{\pgfqpoint{3.347462in}{2.189901in}}%
\pgfpathlineto{\pgfqpoint{3.347462in}{2.189901in}}%
\pgfpathlineto{\pgfqpoint{3.347462in}{2.192850in}}%
\pgfpathlineto{\pgfqpoint{3.352003in}{2.192850in}}%
\pgfpathlineto{\pgfqpoint{3.352003in}{2.189901in}}%
\pgfpathmoveto{\pgfqpoint{3.352003in}{2.189901in}}%
\pgfpathlineto{\pgfqpoint{3.352003in}{2.189901in}}%
\pgfpathlineto{\pgfqpoint{3.352003in}{2.192850in}}%
\pgfpathlineto{\pgfqpoint{3.356544in}{2.192850in}}%
\pgfpathlineto{\pgfqpoint{3.356544in}{2.189901in}}%
\pgfpathmoveto{\pgfqpoint{3.347462in}{2.192850in}}%
\pgfpathlineto{\pgfqpoint{3.347462in}{2.192850in}}%
\pgfpathlineto{\pgfqpoint{3.347462in}{2.195799in}}%
\pgfpathlineto{\pgfqpoint{3.352003in}{2.195799in}}%
\pgfpathlineto{\pgfqpoint{3.352003in}{2.192850in}}%
\pgfpathmoveto{\pgfqpoint{3.347462in}{2.195799in}}%
\pgfpathlineto{\pgfqpoint{3.347462in}{2.195799in}}%
\pgfpathlineto{\pgfqpoint{3.347462in}{2.198748in}}%
\pgfpathlineto{\pgfqpoint{3.352003in}{2.198748in}}%
\pgfpathlineto{\pgfqpoint{3.352003in}{2.195799in}}%
\pgfpathmoveto{\pgfqpoint{3.352003in}{2.192850in}}%
\pgfpathlineto{\pgfqpoint{3.352003in}{2.192850in}}%
\pgfpathlineto{\pgfqpoint{3.352003in}{2.195799in}}%
\pgfpathlineto{\pgfqpoint{3.356544in}{2.195799in}}%
\pgfpathlineto{\pgfqpoint{3.356544in}{2.192850in}}%
\pgfpathmoveto{\pgfqpoint{3.352003in}{2.195799in}}%
\pgfpathlineto{\pgfqpoint{3.352003in}{2.195799in}}%
\pgfpathlineto{\pgfqpoint{3.352003in}{2.198748in}}%
\pgfpathlineto{\pgfqpoint{3.356544in}{2.198748in}}%
\pgfpathlineto{\pgfqpoint{3.356544in}{2.195799in}}%
\pgfpathmoveto{\pgfqpoint{3.356544in}{2.192850in}}%
\pgfpathlineto{\pgfqpoint{3.356544in}{2.192850in}}%
\pgfpathlineto{\pgfqpoint{3.356544in}{2.195799in}}%
\pgfpathlineto{\pgfqpoint{3.361086in}{2.195799in}}%
\pgfpathlineto{\pgfqpoint{3.361086in}{2.192850in}}%
\pgfpathmoveto{\pgfqpoint{3.356544in}{2.195799in}}%
\pgfpathlineto{\pgfqpoint{3.356544in}{2.195799in}}%
\pgfpathlineto{\pgfqpoint{3.356544in}{2.198748in}}%
\pgfpathlineto{\pgfqpoint{3.361086in}{2.198748in}}%
\pgfpathlineto{\pgfqpoint{3.361086in}{2.195799in}}%
\pgfpathmoveto{\pgfqpoint{3.361086in}{2.195799in}}%
\pgfpathlineto{\pgfqpoint{3.361086in}{2.195799in}}%
\pgfpathlineto{\pgfqpoint{3.361086in}{2.198748in}}%
\pgfpathlineto{\pgfqpoint{3.365627in}{2.198748in}}%
\pgfpathlineto{\pgfqpoint{3.365627in}{2.195799in}}%
\pgfpathmoveto{\pgfqpoint{3.356544in}{2.198748in}}%
\pgfpathlineto{\pgfqpoint{3.356544in}{2.198748in}}%
\pgfpathlineto{\pgfqpoint{3.356544in}{2.201697in}}%
\pgfpathlineto{\pgfqpoint{3.361086in}{2.201697in}}%
\pgfpathlineto{\pgfqpoint{3.361086in}{2.198748in}}%
\pgfpathmoveto{\pgfqpoint{3.356544in}{2.201697in}}%
\pgfpathlineto{\pgfqpoint{3.356544in}{2.201697in}}%
\pgfpathlineto{\pgfqpoint{3.356544in}{2.204647in}}%
\pgfpathlineto{\pgfqpoint{3.361086in}{2.204647in}}%
\pgfpathlineto{\pgfqpoint{3.361086in}{2.201697in}}%
\pgfpathmoveto{\pgfqpoint{3.361086in}{2.198748in}}%
\pgfpathlineto{\pgfqpoint{3.361086in}{2.198748in}}%
\pgfpathlineto{\pgfqpoint{3.361086in}{2.201697in}}%
\pgfpathlineto{\pgfqpoint{3.365627in}{2.201697in}}%
\pgfpathlineto{\pgfqpoint{3.365627in}{2.198748in}}%
\pgfpathmoveto{\pgfqpoint{3.361086in}{2.201697in}}%
\pgfpathlineto{\pgfqpoint{3.361086in}{2.201697in}}%
\pgfpathlineto{\pgfqpoint{3.361086in}{2.204647in}}%
\pgfpathlineto{\pgfqpoint{3.365627in}{2.204647in}}%
\pgfpathlineto{\pgfqpoint{3.365627in}{2.201697in}}%
\pgfpathmoveto{\pgfqpoint{3.356544in}{2.835780in}}%
\pgfpathlineto{\pgfqpoint{3.356544in}{2.835780in}}%
\pgfpathlineto{\pgfqpoint{3.356544in}{2.838730in}}%
\pgfpathlineto{\pgfqpoint{3.361086in}{2.838730in}}%
\pgfpathlineto{\pgfqpoint{3.361086in}{2.835780in}}%
\pgfpathmoveto{\pgfqpoint{3.356544in}{2.838730in}}%
\pgfpathlineto{\pgfqpoint{3.356544in}{2.838730in}}%
\pgfpathlineto{\pgfqpoint{3.356544in}{2.841679in}}%
\pgfpathlineto{\pgfqpoint{3.361086in}{2.841679in}}%
\pgfpathlineto{\pgfqpoint{3.361086in}{2.838730in}}%
\pgfpathmoveto{\pgfqpoint{3.361086in}{2.835780in}}%
\pgfpathlineto{\pgfqpoint{3.361086in}{2.835780in}}%
\pgfpathlineto{\pgfqpoint{3.361086in}{2.838730in}}%
\pgfpathlineto{\pgfqpoint{3.365627in}{2.838730in}}%
\pgfpathlineto{\pgfqpoint{3.365627in}{2.835780in}}%
\pgfpathmoveto{\pgfqpoint{3.361086in}{2.838730in}}%
\pgfpathlineto{\pgfqpoint{3.361086in}{2.838730in}}%
\pgfpathlineto{\pgfqpoint{3.361086in}{2.841679in}}%
\pgfpathlineto{\pgfqpoint{3.365627in}{2.841679in}}%
\pgfpathlineto{\pgfqpoint{3.365627in}{2.838730in}}%
\pgfpathmoveto{\pgfqpoint{3.356544in}{2.841679in}}%
\pgfpathlineto{\pgfqpoint{3.356544in}{2.841679in}}%
\pgfpathlineto{\pgfqpoint{3.356544in}{2.844628in}}%
\pgfpathlineto{\pgfqpoint{3.361086in}{2.844628in}}%
\pgfpathlineto{\pgfqpoint{3.361086in}{2.841679in}}%
\pgfpathmoveto{\pgfqpoint{3.356544in}{2.844628in}}%
\pgfpathlineto{\pgfqpoint{3.356544in}{2.844628in}}%
\pgfpathlineto{\pgfqpoint{3.356544in}{2.847577in}}%
\pgfpathlineto{\pgfqpoint{3.361086in}{2.847577in}}%
\pgfpathlineto{\pgfqpoint{3.361086in}{2.844628in}}%
\pgfpathmoveto{\pgfqpoint{3.361086in}{2.841679in}}%
\pgfpathlineto{\pgfqpoint{3.361086in}{2.841679in}}%
\pgfpathlineto{\pgfqpoint{3.361086in}{2.844628in}}%
\pgfpathlineto{\pgfqpoint{3.365627in}{2.844628in}}%
\pgfpathlineto{\pgfqpoint{3.365627in}{2.841679in}}%
\pgfpathmoveto{\pgfqpoint{3.361086in}{2.844628in}}%
\pgfpathlineto{\pgfqpoint{3.361086in}{2.844628in}}%
\pgfpathlineto{\pgfqpoint{3.361086in}{2.847577in}}%
\pgfpathlineto{\pgfqpoint{3.365627in}{2.847577in}}%
\pgfpathlineto{\pgfqpoint{3.365627in}{2.844628in}}%
\pgfpathmoveto{\pgfqpoint{3.347462in}{2.847577in}}%
\pgfpathlineto{\pgfqpoint{3.347462in}{2.847577in}}%
\pgfpathlineto{\pgfqpoint{3.347462in}{2.850526in}}%
\pgfpathlineto{\pgfqpoint{3.352003in}{2.850526in}}%
\pgfpathlineto{\pgfqpoint{3.352003in}{2.847577in}}%
\pgfpathmoveto{\pgfqpoint{3.347462in}{2.850526in}}%
\pgfpathlineto{\pgfqpoint{3.347462in}{2.850526in}}%
\pgfpathlineto{\pgfqpoint{3.347462in}{2.853476in}}%
\pgfpathlineto{\pgfqpoint{3.352003in}{2.853476in}}%
\pgfpathlineto{\pgfqpoint{3.352003in}{2.850526in}}%
\pgfpathmoveto{\pgfqpoint{3.352003in}{2.847577in}}%
\pgfpathlineto{\pgfqpoint{3.352003in}{2.847577in}}%
\pgfpathlineto{\pgfqpoint{3.352003in}{2.850526in}}%
\pgfpathlineto{\pgfqpoint{3.356544in}{2.850526in}}%
\pgfpathlineto{\pgfqpoint{3.356544in}{2.847577in}}%
\pgfpathmoveto{\pgfqpoint{3.352003in}{2.850526in}}%
\pgfpathlineto{\pgfqpoint{3.352003in}{2.850526in}}%
\pgfpathlineto{\pgfqpoint{3.352003in}{2.853476in}}%
\pgfpathlineto{\pgfqpoint{3.356544in}{2.853476in}}%
\pgfpathlineto{\pgfqpoint{3.356544in}{2.850526in}}%
\pgfpathmoveto{\pgfqpoint{3.347462in}{2.853476in}}%
\pgfpathlineto{\pgfqpoint{3.347462in}{2.853476in}}%
\pgfpathlineto{\pgfqpoint{3.347462in}{2.856425in}}%
\pgfpathlineto{\pgfqpoint{3.352003in}{2.856425in}}%
\pgfpathlineto{\pgfqpoint{3.352003in}{2.853476in}}%
\pgfpathmoveto{\pgfqpoint{3.347462in}{2.856425in}}%
\pgfpathlineto{\pgfqpoint{3.347462in}{2.856425in}}%
\pgfpathlineto{\pgfqpoint{3.347462in}{2.859374in}}%
\pgfpathlineto{\pgfqpoint{3.352003in}{2.859374in}}%
\pgfpathlineto{\pgfqpoint{3.352003in}{2.856425in}}%
\pgfpathmoveto{\pgfqpoint{3.352003in}{2.853476in}}%
\pgfpathlineto{\pgfqpoint{3.352003in}{2.853476in}}%
\pgfpathlineto{\pgfqpoint{3.352003in}{2.856425in}}%
\pgfpathlineto{\pgfqpoint{3.356544in}{2.856425in}}%
\pgfpathlineto{\pgfqpoint{3.356544in}{2.853476in}}%
\pgfpathmoveto{\pgfqpoint{3.352003in}{2.856425in}}%
\pgfpathlineto{\pgfqpoint{3.352003in}{2.856425in}}%
\pgfpathlineto{\pgfqpoint{3.352003in}{2.859374in}}%
\pgfpathlineto{\pgfqpoint{3.356544in}{2.859374in}}%
\pgfpathlineto{\pgfqpoint{3.356544in}{2.856425in}}%
\pgfpathmoveto{\pgfqpoint{3.356544in}{2.847577in}}%
\pgfpathlineto{\pgfqpoint{3.356544in}{2.847577in}}%
\pgfpathlineto{\pgfqpoint{3.356544in}{2.850526in}}%
\pgfpathlineto{\pgfqpoint{3.361086in}{2.850526in}}%
\pgfpathlineto{\pgfqpoint{3.361086in}{2.847577in}}%
\pgfpathmoveto{\pgfqpoint{3.356544in}{2.850526in}}%
\pgfpathlineto{\pgfqpoint{3.356544in}{2.850526in}}%
\pgfpathlineto{\pgfqpoint{3.356544in}{2.853476in}}%
\pgfpathlineto{\pgfqpoint{3.361086in}{2.853476in}}%
\pgfpathlineto{\pgfqpoint{3.361086in}{2.850526in}}%
\pgfpathmoveto{\pgfqpoint{3.283886in}{2.930157in}}%
\pgfpathlineto{\pgfqpoint{3.283886in}{2.930157in}}%
\pgfpathlineto{\pgfqpoint{3.283886in}{2.933106in}}%
\pgfpathlineto{\pgfqpoint{3.288427in}{2.933106in}}%
\pgfpathlineto{\pgfqpoint{3.288427in}{2.930157in}}%
\pgfpathmoveto{\pgfqpoint{3.283886in}{2.933106in}}%
\pgfpathlineto{\pgfqpoint{3.283886in}{2.933106in}}%
\pgfpathlineto{\pgfqpoint{3.283886in}{2.936056in}}%
\pgfpathlineto{\pgfqpoint{3.288427in}{2.936056in}}%
\pgfpathlineto{\pgfqpoint{3.288427in}{2.933106in}}%
\pgfpathmoveto{\pgfqpoint{3.288427in}{2.930157in}}%
\pgfpathlineto{\pgfqpoint{3.288427in}{2.930157in}}%
\pgfpathlineto{\pgfqpoint{3.288427in}{2.933106in}}%
\pgfpathlineto{\pgfqpoint{3.292968in}{2.933106in}}%
\pgfpathlineto{\pgfqpoint{3.292968in}{2.930157in}}%
\pgfpathmoveto{\pgfqpoint{3.288427in}{2.933106in}}%
\pgfpathlineto{\pgfqpoint{3.288427in}{2.933106in}}%
\pgfpathlineto{\pgfqpoint{3.288427in}{2.936056in}}%
\pgfpathlineto{\pgfqpoint{3.292968in}{2.936056in}}%
\pgfpathlineto{\pgfqpoint{3.292968in}{2.933106in}}%
\pgfpathmoveto{\pgfqpoint{3.283886in}{2.936056in}}%
\pgfpathlineto{\pgfqpoint{3.283886in}{2.936056in}}%
\pgfpathlineto{\pgfqpoint{3.283886in}{2.939005in}}%
\pgfpathlineto{\pgfqpoint{3.288427in}{2.939005in}}%
\pgfpathlineto{\pgfqpoint{3.288427in}{2.936056in}}%
\pgfpathmoveto{\pgfqpoint{3.283886in}{2.939005in}}%
\pgfpathlineto{\pgfqpoint{3.283886in}{2.939005in}}%
\pgfpathlineto{\pgfqpoint{3.283886in}{2.941954in}}%
\pgfpathlineto{\pgfqpoint{3.288427in}{2.941954in}}%
\pgfpathlineto{\pgfqpoint{3.288427in}{2.939005in}}%
\pgfpathmoveto{\pgfqpoint{3.288427in}{2.936056in}}%
\pgfpathlineto{\pgfqpoint{3.288427in}{2.936056in}}%
\pgfpathlineto{\pgfqpoint{3.288427in}{2.939005in}}%
\pgfpathlineto{\pgfqpoint{3.292968in}{2.939005in}}%
\pgfpathlineto{\pgfqpoint{3.292968in}{2.936056in}}%
\pgfpathmoveto{\pgfqpoint{3.288427in}{2.939005in}}%
\pgfpathlineto{\pgfqpoint{3.288427in}{2.939005in}}%
\pgfpathlineto{\pgfqpoint{3.288427in}{2.941954in}}%
\pgfpathlineto{\pgfqpoint{3.292968in}{2.941954in}}%
\pgfpathlineto{\pgfqpoint{3.292968in}{2.939005in}}%
\pgfpathmoveto{\pgfqpoint{3.274803in}{2.941954in}}%
\pgfpathlineto{\pgfqpoint{3.274803in}{2.941954in}}%
\pgfpathlineto{\pgfqpoint{3.274803in}{2.944904in}}%
\pgfpathlineto{\pgfqpoint{3.279345in}{2.944904in}}%
\pgfpathlineto{\pgfqpoint{3.279345in}{2.941954in}}%
\pgfpathmoveto{\pgfqpoint{3.274803in}{2.944904in}}%
\pgfpathlineto{\pgfqpoint{3.274803in}{2.944904in}}%
\pgfpathlineto{\pgfqpoint{3.274803in}{2.947853in}}%
\pgfpathlineto{\pgfqpoint{3.279345in}{2.947853in}}%
\pgfpathlineto{\pgfqpoint{3.279345in}{2.944904in}}%
\pgfpathmoveto{\pgfqpoint{3.279345in}{2.941954in}}%
\pgfpathlineto{\pgfqpoint{3.279345in}{2.941954in}}%
\pgfpathlineto{\pgfqpoint{3.279345in}{2.944904in}}%
\pgfpathlineto{\pgfqpoint{3.283886in}{2.944904in}}%
\pgfpathlineto{\pgfqpoint{3.283886in}{2.941954in}}%
\pgfpathmoveto{\pgfqpoint{3.279345in}{2.944904in}}%
\pgfpathlineto{\pgfqpoint{3.279345in}{2.944904in}}%
\pgfpathlineto{\pgfqpoint{3.279345in}{2.947853in}}%
\pgfpathlineto{\pgfqpoint{3.283886in}{2.947853in}}%
\pgfpathlineto{\pgfqpoint{3.283886in}{2.944904in}}%
\pgfpathmoveto{\pgfqpoint{3.274803in}{2.947853in}}%
\pgfpathlineto{\pgfqpoint{3.274803in}{2.947853in}}%
\pgfpathlineto{\pgfqpoint{3.274803in}{2.950802in}}%
\pgfpathlineto{\pgfqpoint{3.279345in}{2.950802in}}%
\pgfpathlineto{\pgfqpoint{3.279345in}{2.947853in}}%
\pgfpathmoveto{\pgfqpoint{3.274803in}{2.950802in}}%
\pgfpathlineto{\pgfqpoint{3.274803in}{2.950802in}}%
\pgfpathlineto{\pgfqpoint{3.274803in}{2.953751in}}%
\pgfpathlineto{\pgfqpoint{3.279345in}{2.953751in}}%
\pgfpathlineto{\pgfqpoint{3.279345in}{2.950802in}}%
\pgfpathmoveto{\pgfqpoint{3.279345in}{2.947853in}}%
\pgfpathlineto{\pgfqpoint{3.279345in}{2.947853in}}%
\pgfpathlineto{\pgfqpoint{3.279345in}{2.950802in}}%
\pgfpathlineto{\pgfqpoint{3.283886in}{2.950802in}}%
\pgfpathlineto{\pgfqpoint{3.283886in}{2.947853in}}%
\pgfpathmoveto{\pgfqpoint{3.279345in}{2.950802in}}%
\pgfpathlineto{\pgfqpoint{3.279345in}{2.950802in}}%
\pgfpathlineto{\pgfqpoint{3.279345in}{2.953751in}}%
\pgfpathlineto{\pgfqpoint{3.283886in}{2.953751in}}%
\pgfpathlineto{\pgfqpoint{3.283886in}{2.950802in}}%
\pgfpathmoveto{\pgfqpoint{3.283886in}{2.941954in}}%
\pgfpathlineto{\pgfqpoint{3.283886in}{2.941954in}}%
\pgfpathlineto{\pgfqpoint{3.283886in}{2.944904in}}%
\pgfpathlineto{\pgfqpoint{3.288427in}{2.944904in}}%
\pgfpathlineto{\pgfqpoint{3.288427in}{2.941954in}}%
\pgfpathmoveto{\pgfqpoint{3.283886in}{2.944904in}}%
\pgfpathlineto{\pgfqpoint{3.283886in}{2.944904in}}%
\pgfpathlineto{\pgfqpoint{3.283886in}{2.947853in}}%
\pgfpathlineto{\pgfqpoint{3.288427in}{2.947853in}}%
\pgfpathlineto{\pgfqpoint{3.288427in}{2.944904in}}%
\pgfpathmoveto{\pgfqpoint{3.320215in}{2.882968in}}%
\pgfpathlineto{\pgfqpoint{3.320215in}{2.882968in}}%
\pgfpathlineto{\pgfqpoint{3.320215in}{2.885918in}}%
\pgfpathlineto{\pgfqpoint{3.324756in}{2.885918in}}%
\pgfpathlineto{\pgfqpoint{3.324756in}{2.882968in}}%
\pgfpathmoveto{\pgfqpoint{3.320215in}{2.885918in}}%
\pgfpathlineto{\pgfqpoint{3.320215in}{2.885918in}}%
\pgfpathlineto{\pgfqpoint{3.320215in}{2.888867in}}%
\pgfpathlineto{\pgfqpoint{3.324756in}{2.888867in}}%
\pgfpathlineto{\pgfqpoint{3.324756in}{2.885918in}}%
\pgfpathmoveto{\pgfqpoint{3.324756in}{2.882968in}}%
\pgfpathlineto{\pgfqpoint{3.324756in}{2.882968in}}%
\pgfpathlineto{\pgfqpoint{3.324756in}{2.885918in}}%
\pgfpathlineto{\pgfqpoint{3.329297in}{2.885918in}}%
\pgfpathlineto{\pgfqpoint{3.329297in}{2.882968in}}%
\pgfpathmoveto{\pgfqpoint{3.324756in}{2.885918in}}%
\pgfpathlineto{\pgfqpoint{3.324756in}{2.885918in}}%
\pgfpathlineto{\pgfqpoint{3.324756in}{2.888867in}}%
\pgfpathlineto{\pgfqpoint{3.329297in}{2.888867in}}%
\pgfpathlineto{\pgfqpoint{3.329297in}{2.885918in}}%
\pgfpathmoveto{\pgfqpoint{3.320215in}{2.888867in}}%
\pgfpathlineto{\pgfqpoint{3.320215in}{2.888867in}}%
\pgfpathlineto{\pgfqpoint{3.320215in}{2.891816in}}%
\pgfpathlineto{\pgfqpoint{3.324756in}{2.891816in}}%
\pgfpathlineto{\pgfqpoint{3.324756in}{2.888867in}}%
\pgfpathmoveto{\pgfqpoint{3.320215in}{2.891816in}}%
\pgfpathlineto{\pgfqpoint{3.320215in}{2.891816in}}%
\pgfpathlineto{\pgfqpoint{3.320215in}{2.894766in}}%
\pgfpathlineto{\pgfqpoint{3.324756in}{2.894766in}}%
\pgfpathlineto{\pgfqpoint{3.324756in}{2.891816in}}%
\pgfpathmoveto{\pgfqpoint{3.324756in}{2.888867in}}%
\pgfpathlineto{\pgfqpoint{3.324756in}{2.888867in}}%
\pgfpathlineto{\pgfqpoint{3.324756in}{2.891816in}}%
\pgfpathlineto{\pgfqpoint{3.329297in}{2.891816in}}%
\pgfpathlineto{\pgfqpoint{3.329297in}{2.888867in}}%
\pgfpathmoveto{\pgfqpoint{3.324756in}{2.891816in}}%
\pgfpathlineto{\pgfqpoint{3.324756in}{2.891816in}}%
\pgfpathlineto{\pgfqpoint{3.324756in}{2.894766in}}%
\pgfpathlineto{\pgfqpoint{3.329297in}{2.894766in}}%
\pgfpathlineto{\pgfqpoint{3.329297in}{2.891816in}}%
\pgfpathmoveto{\pgfqpoint{3.311133in}{2.894766in}}%
\pgfpathlineto{\pgfqpoint{3.311133in}{2.894766in}}%
\pgfpathlineto{\pgfqpoint{3.311133in}{2.897715in}}%
\pgfpathlineto{\pgfqpoint{3.315674in}{2.897715in}}%
\pgfpathlineto{\pgfqpoint{3.315674in}{2.894766in}}%
\pgfpathmoveto{\pgfqpoint{3.311133in}{2.897715in}}%
\pgfpathlineto{\pgfqpoint{3.311133in}{2.897715in}}%
\pgfpathlineto{\pgfqpoint{3.311133in}{2.900664in}}%
\pgfpathlineto{\pgfqpoint{3.315674in}{2.900664in}}%
\pgfpathlineto{\pgfqpoint{3.315674in}{2.897715in}}%
\pgfpathmoveto{\pgfqpoint{3.315674in}{2.894766in}}%
\pgfpathlineto{\pgfqpoint{3.315674in}{2.894766in}}%
\pgfpathlineto{\pgfqpoint{3.315674in}{2.897715in}}%
\pgfpathlineto{\pgfqpoint{3.320215in}{2.897715in}}%
\pgfpathlineto{\pgfqpoint{3.320215in}{2.894766in}}%
\pgfpathmoveto{\pgfqpoint{3.315674in}{2.897715in}}%
\pgfpathlineto{\pgfqpoint{3.315674in}{2.897715in}}%
\pgfpathlineto{\pgfqpoint{3.315674in}{2.900664in}}%
\pgfpathlineto{\pgfqpoint{3.320215in}{2.900664in}}%
\pgfpathlineto{\pgfqpoint{3.320215in}{2.897715in}}%
\pgfpathmoveto{\pgfqpoint{3.311133in}{2.900664in}}%
\pgfpathlineto{\pgfqpoint{3.311133in}{2.900664in}}%
\pgfpathlineto{\pgfqpoint{3.311133in}{2.903613in}}%
\pgfpathlineto{\pgfqpoint{3.315674in}{2.903613in}}%
\pgfpathlineto{\pgfqpoint{3.315674in}{2.900664in}}%
\pgfpathmoveto{\pgfqpoint{3.311133in}{2.903613in}}%
\pgfpathlineto{\pgfqpoint{3.311133in}{2.903613in}}%
\pgfpathlineto{\pgfqpoint{3.311133in}{2.906563in}}%
\pgfpathlineto{\pgfqpoint{3.315674in}{2.906563in}}%
\pgfpathlineto{\pgfqpoint{3.315674in}{2.903613in}}%
\pgfpathmoveto{\pgfqpoint{3.315674in}{2.900664in}}%
\pgfpathlineto{\pgfqpoint{3.315674in}{2.900664in}}%
\pgfpathlineto{\pgfqpoint{3.315674in}{2.903613in}}%
\pgfpathlineto{\pgfqpoint{3.320215in}{2.903613in}}%
\pgfpathlineto{\pgfqpoint{3.320215in}{2.900664in}}%
\pgfpathmoveto{\pgfqpoint{3.315674in}{2.903613in}}%
\pgfpathlineto{\pgfqpoint{3.315674in}{2.903613in}}%
\pgfpathlineto{\pgfqpoint{3.315674in}{2.906563in}}%
\pgfpathlineto{\pgfqpoint{3.320215in}{2.906563in}}%
\pgfpathlineto{\pgfqpoint{3.320215in}{2.903613in}}%
\pgfpathmoveto{\pgfqpoint{3.320215in}{2.894766in}}%
\pgfpathlineto{\pgfqpoint{3.320215in}{2.894766in}}%
\pgfpathlineto{\pgfqpoint{3.320215in}{2.897715in}}%
\pgfpathlineto{\pgfqpoint{3.324756in}{2.897715in}}%
\pgfpathlineto{\pgfqpoint{3.324756in}{2.894766in}}%
\pgfpathmoveto{\pgfqpoint{3.320215in}{2.897715in}}%
\pgfpathlineto{\pgfqpoint{3.320215in}{2.897715in}}%
\pgfpathlineto{\pgfqpoint{3.320215in}{2.900664in}}%
\pgfpathlineto{\pgfqpoint{3.324756in}{2.900664in}}%
\pgfpathlineto{\pgfqpoint{3.324756in}{2.897715in}}%
\pgfpathmoveto{\pgfqpoint{3.338380in}{2.859374in}}%
\pgfpathlineto{\pgfqpoint{3.338380in}{2.859374in}}%
\pgfpathlineto{\pgfqpoint{3.338380in}{2.862323in}}%
\pgfpathlineto{\pgfqpoint{3.342921in}{2.862323in}}%
\pgfpathlineto{\pgfqpoint{3.342921in}{2.859374in}}%
\pgfpathmoveto{\pgfqpoint{3.338380in}{2.862323in}}%
\pgfpathlineto{\pgfqpoint{3.338380in}{2.862323in}}%
\pgfpathlineto{\pgfqpoint{3.338380in}{2.865273in}}%
\pgfpathlineto{\pgfqpoint{3.342921in}{2.865273in}}%
\pgfpathlineto{\pgfqpoint{3.342921in}{2.862323in}}%
\pgfpathmoveto{\pgfqpoint{3.342921in}{2.859374in}}%
\pgfpathlineto{\pgfqpoint{3.342921in}{2.859374in}}%
\pgfpathlineto{\pgfqpoint{3.342921in}{2.862323in}}%
\pgfpathlineto{\pgfqpoint{3.347462in}{2.862323in}}%
\pgfpathlineto{\pgfqpoint{3.347462in}{2.859374in}}%
\pgfpathmoveto{\pgfqpoint{3.342921in}{2.862323in}}%
\pgfpathlineto{\pgfqpoint{3.342921in}{2.862323in}}%
\pgfpathlineto{\pgfqpoint{3.342921in}{2.865273in}}%
\pgfpathlineto{\pgfqpoint{3.347462in}{2.865273in}}%
\pgfpathlineto{\pgfqpoint{3.347462in}{2.862323in}}%
\pgfpathmoveto{\pgfqpoint{3.338380in}{2.865273in}}%
\pgfpathlineto{\pgfqpoint{3.338380in}{2.865273in}}%
\pgfpathlineto{\pgfqpoint{3.338380in}{2.868222in}}%
\pgfpathlineto{\pgfqpoint{3.342921in}{2.868222in}}%
\pgfpathlineto{\pgfqpoint{3.342921in}{2.865273in}}%
\pgfpathmoveto{\pgfqpoint{3.338380in}{2.868222in}}%
\pgfpathlineto{\pgfqpoint{3.338380in}{2.868222in}}%
\pgfpathlineto{\pgfqpoint{3.338380in}{2.871171in}}%
\pgfpathlineto{\pgfqpoint{3.342921in}{2.871171in}}%
\pgfpathlineto{\pgfqpoint{3.342921in}{2.868222in}}%
\pgfpathmoveto{\pgfqpoint{3.342921in}{2.865273in}}%
\pgfpathlineto{\pgfqpoint{3.342921in}{2.865273in}}%
\pgfpathlineto{\pgfqpoint{3.342921in}{2.868222in}}%
\pgfpathlineto{\pgfqpoint{3.347462in}{2.868222in}}%
\pgfpathlineto{\pgfqpoint{3.347462in}{2.865273in}}%
\pgfpathmoveto{\pgfqpoint{3.342921in}{2.868222in}}%
\pgfpathlineto{\pgfqpoint{3.342921in}{2.868222in}}%
\pgfpathlineto{\pgfqpoint{3.342921in}{2.871171in}}%
\pgfpathlineto{\pgfqpoint{3.347462in}{2.871171in}}%
\pgfpathlineto{\pgfqpoint{3.347462in}{2.868222in}}%
\pgfpathmoveto{\pgfqpoint{3.329297in}{2.871171in}}%
\pgfpathlineto{\pgfqpoint{3.329297in}{2.871171in}}%
\pgfpathlineto{\pgfqpoint{3.329297in}{2.874120in}}%
\pgfpathlineto{\pgfqpoint{3.333839in}{2.874120in}}%
\pgfpathlineto{\pgfqpoint{3.333839in}{2.871171in}}%
\pgfpathmoveto{\pgfqpoint{3.329297in}{2.874120in}}%
\pgfpathlineto{\pgfqpoint{3.329297in}{2.874120in}}%
\pgfpathlineto{\pgfqpoint{3.329297in}{2.877070in}}%
\pgfpathlineto{\pgfqpoint{3.333839in}{2.877070in}}%
\pgfpathlineto{\pgfqpoint{3.333839in}{2.874120in}}%
\pgfpathmoveto{\pgfqpoint{3.333839in}{2.871171in}}%
\pgfpathlineto{\pgfqpoint{3.333839in}{2.871171in}}%
\pgfpathlineto{\pgfqpoint{3.333839in}{2.874120in}}%
\pgfpathlineto{\pgfqpoint{3.338380in}{2.874120in}}%
\pgfpathlineto{\pgfqpoint{3.338380in}{2.871171in}}%
\pgfpathmoveto{\pgfqpoint{3.333839in}{2.874120in}}%
\pgfpathlineto{\pgfqpoint{3.333839in}{2.874120in}}%
\pgfpathlineto{\pgfqpoint{3.333839in}{2.877070in}}%
\pgfpathlineto{\pgfqpoint{3.338380in}{2.877070in}}%
\pgfpathlineto{\pgfqpoint{3.338380in}{2.874120in}}%
\pgfpathmoveto{\pgfqpoint{3.329297in}{2.877070in}}%
\pgfpathlineto{\pgfqpoint{3.329297in}{2.877070in}}%
\pgfpathlineto{\pgfqpoint{3.329297in}{2.880019in}}%
\pgfpathlineto{\pgfqpoint{3.333839in}{2.880019in}}%
\pgfpathlineto{\pgfqpoint{3.333839in}{2.877070in}}%
\pgfpathmoveto{\pgfqpoint{3.329297in}{2.880019in}}%
\pgfpathlineto{\pgfqpoint{3.329297in}{2.880019in}}%
\pgfpathlineto{\pgfqpoint{3.329297in}{2.882968in}}%
\pgfpathlineto{\pgfqpoint{3.333839in}{2.882968in}}%
\pgfpathlineto{\pgfqpoint{3.333839in}{2.880019in}}%
\pgfpathmoveto{\pgfqpoint{3.333839in}{2.877070in}}%
\pgfpathlineto{\pgfqpoint{3.333839in}{2.877070in}}%
\pgfpathlineto{\pgfqpoint{3.333839in}{2.880019in}}%
\pgfpathlineto{\pgfqpoint{3.338380in}{2.880019in}}%
\pgfpathlineto{\pgfqpoint{3.338380in}{2.877070in}}%
\pgfpathmoveto{\pgfqpoint{3.333839in}{2.880019in}}%
\pgfpathlineto{\pgfqpoint{3.333839in}{2.880019in}}%
\pgfpathlineto{\pgfqpoint{3.333839in}{2.882968in}}%
\pgfpathlineto{\pgfqpoint{3.338380in}{2.882968in}}%
\pgfpathlineto{\pgfqpoint{3.338380in}{2.880019in}}%
\pgfpathmoveto{\pgfqpoint{3.338380in}{2.871171in}}%
\pgfpathlineto{\pgfqpoint{3.338380in}{2.871171in}}%
\pgfpathlineto{\pgfqpoint{3.338380in}{2.874120in}}%
\pgfpathlineto{\pgfqpoint{3.342921in}{2.874120in}}%
\pgfpathlineto{\pgfqpoint{3.342921in}{2.871171in}}%
\pgfpathmoveto{\pgfqpoint{3.338380in}{2.874120in}}%
\pgfpathlineto{\pgfqpoint{3.338380in}{2.874120in}}%
\pgfpathlineto{\pgfqpoint{3.338380in}{2.877070in}}%
\pgfpathlineto{\pgfqpoint{3.342921in}{2.877070in}}%
\pgfpathlineto{\pgfqpoint{3.342921in}{2.874120in}}%
\pgfpathmoveto{\pgfqpoint{3.347462in}{2.859374in}}%
\pgfpathlineto{\pgfqpoint{3.347462in}{2.859374in}}%
\pgfpathlineto{\pgfqpoint{3.347462in}{2.862323in}}%
\pgfpathlineto{\pgfqpoint{3.352003in}{2.862323in}}%
\pgfpathlineto{\pgfqpoint{3.352003in}{2.859374in}}%
\pgfpathmoveto{\pgfqpoint{3.347462in}{2.862323in}}%
\pgfpathlineto{\pgfqpoint{3.347462in}{2.862323in}}%
\pgfpathlineto{\pgfqpoint{3.347462in}{2.865273in}}%
\pgfpathlineto{\pgfqpoint{3.352003in}{2.865273in}}%
\pgfpathlineto{\pgfqpoint{3.352003in}{2.862323in}}%
\pgfpathmoveto{\pgfqpoint{3.329297in}{2.882968in}}%
\pgfpathlineto{\pgfqpoint{3.329297in}{2.882968in}}%
\pgfpathlineto{\pgfqpoint{3.329297in}{2.885918in}}%
\pgfpathlineto{\pgfqpoint{3.333839in}{2.885918in}}%
\pgfpathlineto{\pgfqpoint{3.333839in}{2.882968in}}%
\pgfpathmoveto{\pgfqpoint{3.329297in}{2.885918in}}%
\pgfpathlineto{\pgfqpoint{3.329297in}{2.885918in}}%
\pgfpathlineto{\pgfqpoint{3.329297in}{2.888867in}}%
\pgfpathlineto{\pgfqpoint{3.333839in}{2.888867in}}%
\pgfpathlineto{\pgfqpoint{3.333839in}{2.885918in}}%
\pgfpathmoveto{\pgfqpoint{3.302050in}{2.906563in}}%
\pgfpathlineto{\pgfqpoint{3.302050in}{2.906563in}}%
\pgfpathlineto{\pgfqpoint{3.302050in}{2.909512in}}%
\pgfpathlineto{\pgfqpoint{3.306592in}{2.909512in}}%
\pgfpathlineto{\pgfqpoint{3.306592in}{2.906563in}}%
\pgfpathmoveto{\pgfqpoint{3.302050in}{2.909512in}}%
\pgfpathlineto{\pgfqpoint{3.302050in}{2.909512in}}%
\pgfpathlineto{\pgfqpoint{3.302050in}{2.912461in}}%
\pgfpathlineto{\pgfqpoint{3.306592in}{2.912461in}}%
\pgfpathlineto{\pgfqpoint{3.306592in}{2.909512in}}%
\pgfpathmoveto{\pgfqpoint{3.306592in}{2.906563in}}%
\pgfpathlineto{\pgfqpoint{3.306592in}{2.906563in}}%
\pgfpathlineto{\pgfqpoint{3.306592in}{2.909512in}}%
\pgfpathlineto{\pgfqpoint{3.311133in}{2.909512in}}%
\pgfpathlineto{\pgfqpoint{3.311133in}{2.906563in}}%
\pgfpathmoveto{\pgfqpoint{3.306592in}{2.909512in}}%
\pgfpathlineto{\pgfqpoint{3.306592in}{2.909512in}}%
\pgfpathlineto{\pgfqpoint{3.306592in}{2.912461in}}%
\pgfpathlineto{\pgfqpoint{3.311133in}{2.912461in}}%
\pgfpathlineto{\pgfqpoint{3.311133in}{2.909512in}}%
\pgfpathmoveto{\pgfqpoint{3.302050in}{2.912461in}}%
\pgfpathlineto{\pgfqpoint{3.302050in}{2.912461in}}%
\pgfpathlineto{\pgfqpoint{3.302050in}{2.915411in}}%
\pgfpathlineto{\pgfqpoint{3.306592in}{2.915411in}}%
\pgfpathlineto{\pgfqpoint{3.306592in}{2.912461in}}%
\pgfpathmoveto{\pgfqpoint{3.302050in}{2.915411in}}%
\pgfpathlineto{\pgfqpoint{3.302050in}{2.915411in}}%
\pgfpathlineto{\pgfqpoint{3.302050in}{2.918360in}}%
\pgfpathlineto{\pgfqpoint{3.306592in}{2.918360in}}%
\pgfpathlineto{\pgfqpoint{3.306592in}{2.915411in}}%
\pgfpathmoveto{\pgfqpoint{3.306592in}{2.912461in}}%
\pgfpathlineto{\pgfqpoint{3.306592in}{2.912461in}}%
\pgfpathlineto{\pgfqpoint{3.306592in}{2.915411in}}%
\pgfpathlineto{\pgfqpoint{3.311133in}{2.915411in}}%
\pgfpathlineto{\pgfqpoint{3.311133in}{2.912461in}}%
\pgfpathmoveto{\pgfqpoint{3.306592in}{2.915411in}}%
\pgfpathlineto{\pgfqpoint{3.306592in}{2.915411in}}%
\pgfpathlineto{\pgfqpoint{3.306592in}{2.918360in}}%
\pgfpathlineto{\pgfqpoint{3.311133in}{2.918360in}}%
\pgfpathlineto{\pgfqpoint{3.311133in}{2.915411in}}%
\pgfpathmoveto{\pgfqpoint{3.292968in}{2.918360in}}%
\pgfpathlineto{\pgfqpoint{3.292968in}{2.918360in}}%
\pgfpathlineto{\pgfqpoint{3.292968in}{2.921309in}}%
\pgfpathlineto{\pgfqpoint{3.297509in}{2.921309in}}%
\pgfpathlineto{\pgfqpoint{3.297509in}{2.918360in}}%
\pgfpathmoveto{\pgfqpoint{3.292968in}{2.921309in}}%
\pgfpathlineto{\pgfqpoint{3.292968in}{2.921309in}}%
\pgfpathlineto{\pgfqpoint{3.292968in}{2.924258in}}%
\pgfpathlineto{\pgfqpoint{3.297509in}{2.924258in}}%
\pgfpathlineto{\pgfqpoint{3.297509in}{2.921309in}}%
\pgfpathmoveto{\pgfqpoint{3.297509in}{2.918360in}}%
\pgfpathlineto{\pgfqpoint{3.297509in}{2.918360in}}%
\pgfpathlineto{\pgfqpoint{3.297509in}{2.921309in}}%
\pgfpathlineto{\pgfqpoint{3.302050in}{2.921309in}}%
\pgfpathlineto{\pgfqpoint{3.302050in}{2.918360in}}%
\pgfpathmoveto{\pgfqpoint{3.297509in}{2.921309in}}%
\pgfpathlineto{\pgfqpoint{3.297509in}{2.921309in}}%
\pgfpathlineto{\pgfqpoint{3.297509in}{2.924258in}}%
\pgfpathlineto{\pgfqpoint{3.302050in}{2.924258in}}%
\pgfpathlineto{\pgfqpoint{3.302050in}{2.921309in}}%
\pgfpathmoveto{\pgfqpoint{3.292968in}{2.924258in}}%
\pgfpathlineto{\pgfqpoint{3.292968in}{2.924258in}}%
\pgfpathlineto{\pgfqpoint{3.292968in}{2.927208in}}%
\pgfpathlineto{\pgfqpoint{3.297509in}{2.927208in}}%
\pgfpathlineto{\pgfqpoint{3.297509in}{2.924258in}}%
\pgfpathmoveto{\pgfqpoint{3.292968in}{2.927208in}}%
\pgfpathlineto{\pgfqpoint{3.292968in}{2.927208in}}%
\pgfpathlineto{\pgfqpoint{3.292968in}{2.930157in}}%
\pgfpathlineto{\pgfqpoint{3.297509in}{2.930157in}}%
\pgfpathlineto{\pgfqpoint{3.297509in}{2.927208in}}%
\pgfpathmoveto{\pgfqpoint{3.297509in}{2.924258in}}%
\pgfpathlineto{\pgfqpoint{3.297509in}{2.924258in}}%
\pgfpathlineto{\pgfqpoint{3.297509in}{2.927208in}}%
\pgfpathlineto{\pgfqpoint{3.302050in}{2.927208in}}%
\pgfpathlineto{\pgfqpoint{3.302050in}{2.924258in}}%
\pgfpathmoveto{\pgfqpoint{3.297509in}{2.927208in}}%
\pgfpathlineto{\pgfqpoint{3.297509in}{2.927208in}}%
\pgfpathlineto{\pgfqpoint{3.297509in}{2.930157in}}%
\pgfpathlineto{\pgfqpoint{3.302050in}{2.930157in}}%
\pgfpathlineto{\pgfqpoint{3.302050in}{2.927208in}}%
\pgfpathmoveto{\pgfqpoint{3.302050in}{2.918360in}}%
\pgfpathlineto{\pgfqpoint{3.302050in}{2.918360in}}%
\pgfpathlineto{\pgfqpoint{3.302050in}{2.921309in}}%
\pgfpathlineto{\pgfqpoint{3.306592in}{2.921309in}}%
\pgfpathlineto{\pgfqpoint{3.306592in}{2.918360in}}%
\pgfpathmoveto{\pgfqpoint{3.302050in}{2.921309in}}%
\pgfpathlineto{\pgfqpoint{3.302050in}{2.921309in}}%
\pgfpathlineto{\pgfqpoint{3.302050in}{2.924258in}}%
\pgfpathlineto{\pgfqpoint{3.306592in}{2.924258in}}%
\pgfpathlineto{\pgfqpoint{3.306592in}{2.921309in}}%
\pgfpathmoveto{\pgfqpoint{3.311133in}{2.906563in}}%
\pgfpathlineto{\pgfqpoint{3.311133in}{2.906563in}}%
\pgfpathlineto{\pgfqpoint{3.311133in}{2.909512in}}%
\pgfpathlineto{\pgfqpoint{3.315674in}{2.909512in}}%
\pgfpathlineto{\pgfqpoint{3.315674in}{2.906563in}}%
\pgfpathmoveto{\pgfqpoint{3.311133in}{2.909512in}}%
\pgfpathlineto{\pgfqpoint{3.311133in}{2.909512in}}%
\pgfpathlineto{\pgfqpoint{3.311133in}{2.912461in}}%
\pgfpathlineto{\pgfqpoint{3.315674in}{2.912461in}}%
\pgfpathlineto{\pgfqpoint{3.315674in}{2.909512in}}%
\pgfpathmoveto{\pgfqpoint{3.292968in}{2.930157in}}%
\pgfpathlineto{\pgfqpoint{3.292968in}{2.930157in}}%
\pgfpathlineto{\pgfqpoint{3.292968in}{2.933106in}}%
\pgfpathlineto{\pgfqpoint{3.297509in}{2.933106in}}%
\pgfpathlineto{\pgfqpoint{3.297509in}{2.930157in}}%
\pgfpathmoveto{\pgfqpoint{3.292968in}{2.933106in}}%
\pgfpathlineto{\pgfqpoint{3.292968in}{2.933106in}}%
\pgfpathlineto{\pgfqpoint{3.292968in}{2.936056in}}%
\pgfpathlineto{\pgfqpoint{3.297509in}{2.936056in}}%
\pgfpathlineto{\pgfqpoint{3.297509in}{2.933106in}}%
\pgfpathmoveto{\pgfqpoint{3.247557in}{2.977345in}}%
\pgfpathlineto{\pgfqpoint{3.247557in}{2.977345in}}%
\pgfpathlineto{\pgfqpoint{3.247557in}{2.980294in}}%
\pgfpathlineto{\pgfqpoint{3.252098in}{2.980294in}}%
\pgfpathlineto{\pgfqpoint{3.252098in}{2.977345in}}%
\pgfpathmoveto{\pgfqpoint{3.247557in}{2.980294in}}%
\pgfpathlineto{\pgfqpoint{3.247557in}{2.980294in}}%
\pgfpathlineto{\pgfqpoint{3.247557in}{2.983243in}}%
\pgfpathlineto{\pgfqpoint{3.252098in}{2.983243in}}%
\pgfpathlineto{\pgfqpoint{3.252098in}{2.980294in}}%
\pgfpathmoveto{\pgfqpoint{3.252098in}{2.977345in}}%
\pgfpathlineto{\pgfqpoint{3.252098in}{2.977345in}}%
\pgfpathlineto{\pgfqpoint{3.252098in}{2.980294in}}%
\pgfpathlineto{\pgfqpoint{3.256639in}{2.980294in}}%
\pgfpathlineto{\pgfqpoint{3.256639in}{2.977345in}}%
\pgfpathmoveto{\pgfqpoint{3.252098in}{2.980294in}}%
\pgfpathlineto{\pgfqpoint{3.252098in}{2.980294in}}%
\pgfpathlineto{\pgfqpoint{3.252098in}{2.983243in}}%
\pgfpathlineto{\pgfqpoint{3.256639in}{2.983243in}}%
\pgfpathlineto{\pgfqpoint{3.256639in}{2.980294in}}%
\pgfpathmoveto{\pgfqpoint{3.247557in}{2.983243in}}%
\pgfpathlineto{\pgfqpoint{3.247557in}{2.983243in}}%
\pgfpathlineto{\pgfqpoint{3.247557in}{2.986192in}}%
\pgfpathlineto{\pgfqpoint{3.252098in}{2.986192in}}%
\pgfpathlineto{\pgfqpoint{3.252098in}{2.983243in}}%
\pgfpathmoveto{\pgfqpoint{3.247557in}{2.986192in}}%
\pgfpathlineto{\pgfqpoint{3.247557in}{2.986192in}}%
\pgfpathlineto{\pgfqpoint{3.247557in}{2.989141in}}%
\pgfpathlineto{\pgfqpoint{3.252098in}{2.989141in}}%
\pgfpathlineto{\pgfqpoint{3.252098in}{2.986192in}}%
\pgfpathmoveto{\pgfqpoint{3.252098in}{2.983243in}}%
\pgfpathlineto{\pgfqpoint{3.252098in}{2.983243in}}%
\pgfpathlineto{\pgfqpoint{3.252098in}{2.986192in}}%
\pgfpathlineto{\pgfqpoint{3.256639in}{2.986192in}}%
\pgfpathlineto{\pgfqpoint{3.256639in}{2.983243in}}%
\pgfpathmoveto{\pgfqpoint{3.252098in}{2.986192in}}%
\pgfpathlineto{\pgfqpoint{3.252098in}{2.986192in}}%
\pgfpathlineto{\pgfqpoint{3.252098in}{2.989141in}}%
\pgfpathlineto{\pgfqpoint{3.256639in}{2.989141in}}%
\pgfpathlineto{\pgfqpoint{3.256639in}{2.986192in}}%
\pgfpathmoveto{\pgfqpoint{3.238474in}{2.989141in}}%
\pgfpathlineto{\pgfqpoint{3.238474in}{2.989141in}}%
\pgfpathlineto{\pgfqpoint{3.238474in}{2.992090in}}%
\pgfpathlineto{\pgfqpoint{3.243015in}{2.992090in}}%
\pgfpathlineto{\pgfqpoint{3.243015in}{2.989141in}}%
\pgfpathmoveto{\pgfqpoint{3.238474in}{2.992090in}}%
\pgfpathlineto{\pgfqpoint{3.238474in}{2.992090in}}%
\pgfpathlineto{\pgfqpoint{3.238474in}{2.995040in}}%
\pgfpathlineto{\pgfqpoint{3.243015in}{2.995040in}}%
\pgfpathlineto{\pgfqpoint{3.243015in}{2.992090in}}%
\pgfpathmoveto{\pgfqpoint{3.243015in}{2.989141in}}%
\pgfpathlineto{\pgfqpoint{3.243015in}{2.989141in}}%
\pgfpathlineto{\pgfqpoint{3.243015in}{2.992090in}}%
\pgfpathlineto{\pgfqpoint{3.247557in}{2.992090in}}%
\pgfpathlineto{\pgfqpoint{3.247557in}{2.989141in}}%
\pgfpathmoveto{\pgfqpoint{3.243015in}{2.992090in}}%
\pgfpathlineto{\pgfqpoint{3.243015in}{2.992090in}}%
\pgfpathlineto{\pgfqpoint{3.243015in}{2.995040in}}%
\pgfpathlineto{\pgfqpoint{3.247557in}{2.995040in}}%
\pgfpathlineto{\pgfqpoint{3.247557in}{2.992090in}}%
\pgfpathmoveto{\pgfqpoint{3.238474in}{2.995040in}}%
\pgfpathlineto{\pgfqpoint{3.238474in}{2.995040in}}%
\pgfpathlineto{\pgfqpoint{3.238474in}{2.997989in}}%
\pgfpathlineto{\pgfqpoint{3.243015in}{2.997989in}}%
\pgfpathlineto{\pgfqpoint{3.243015in}{2.995040in}}%
\pgfpathmoveto{\pgfqpoint{3.238474in}{2.997989in}}%
\pgfpathlineto{\pgfqpoint{3.238474in}{2.997989in}}%
\pgfpathlineto{\pgfqpoint{3.238474in}{3.000938in}}%
\pgfpathlineto{\pgfqpoint{3.243015in}{3.000938in}}%
\pgfpathlineto{\pgfqpoint{3.243015in}{2.997989in}}%
\pgfpathmoveto{\pgfqpoint{3.243015in}{2.995040in}}%
\pgfpathlineto{\pgfqpoint{3.243015in}{2.995040in}}%
\pgfpathlineto{\pgfqpoint{3.243015in}{2.997989in}}%
\pgfpathlineto{\pgfqpoint{3.247557in}{2.997989in}}%
\pgfpathlineto{\pgfqpoint{3.247557in}{2.995040in}}%
\pgfpathmoveto{\pgfqpoint{3.243015in}{2.997989in}}%
\pgfpathlineto{\pgfqpoint{3.243015in}{2.997989in}}%
\pgfpathlineto{\pgfqpoint{3.243015in}{3.000938in}}%
\pgfpathlineto{\pgfqpoint{3.247557in}{3.000938in}}%
\pgfpathlineto{\pgfqpoint{3.247557in}{2.997989in}}%
\pgfpathmoveto{\pgfqpoint{3.247557in}{2.989141in}}%
\pgfpathlineto{\pgfqpoint{3.247557in}{2.989141in}}%
\pgfpathlineto{\pgfqpoint{3.247557in}{2.992090in}}%
\pgfpathlineto{\pgfqpoint{3.252098in}{2.992090in}}%
\pgfpathlineto{\pgfqpoint{3.252098in}{2.989141in}}%
\pgfpathmoveto{\pgfqpoint{3.247557in}{2.992090in}}%
\pgfpathlineto{\pgfqpoint{3.247557in}{2.992090in}}%
\pgfpathlineto{\pgfqpoint{3.247557in}{2.995040in}}%
\pgfpathlineto{\pgfqpoint{3.252098in}{2.995040in}}%
\pgfpathlineto{\pgfqpoint{3.252098in}{2.992090in}}%
\pgfpathmoveto{\pgfqpoint{3.265721in}{2.953751in}}%
\pgfpathlineto{\pgfqpoint{3.265721in}{2.953751in}}%
\pgfpathlineto{\pgfqpoint{3.265721in}{2.956701in}}%
\pgfpathlineto{\pgfqpoint{3.270262in}{2.956701in}}%
\pgfpathlineto{\pgfqpoint{3.270262in}{2.953751in}}%
\pgfpathmoveto{\pgfqpoint{3.265721in}{2.956701in}}%
\pgfpathlineto{\pgfqpoint{3.265721in}{2.956701in}}%
\pgfpathlineto{\pgfqpoint{3.265721in}{2.959650in}}%
\pgfpathlineto{\pgfqpoint{3.270262in}{2.959650in}}%
\pgfpathlineto{\pgfqpoint{3.270262in}{2.956701in}}%
\pgfpathmoveto{\pgfqpoint{3.270262in}{2.953751in}}%
\pgfpathlineto{\pgfqpoint{3.270262in}{2.953751in}}%
\pgfpathlineto{\pgfqpoint{3.270262in}{2.956701in}}%
\pgfpathlineto{\pgfqpoint{3.274803in}{2.956701in}}%
\pgfpathlineto{\pgfqpoint{3.274803in}{2.953751in}}%
\pgfpathmoveto{\pgfqpoint{3.270262in}{2.956701in}}%
\pgfpathlineto{\pgfqpoint{3.270262in}{2.956701in}}%
\pgfpathlineto{\pgfqpoint{3.270262in}{2.959650in}}%
\pgfpathlineto{\pgfqpoint{3.274803in}{2.959650in}}%
\pgfpathlineto{\pgfqpoint{3.274803in}{2.956701in}}%
\pgfpathmoveto{\pgfqpoint{3.265721in}{2.959650in}}%
\pgfpathlineto{\pgfqpoint{3.265721in}{2.959650in}}%
\pgfpathlineto{\pgfqpoint{3.265721in}{2.962599in}}%
\pgfpathlineto{\pgfqpoint{3.270262in}{2.962599in}}%
\pgfpathlineto{\pgfqpoint{3.270262in}{2.959650in}}%
\pgfpathmoveto{\pgfqpoint{3.265721in}{2.962599in}}%
\pgfpathlineto{\pgfqpoint{3.265721in}{2.962599in}}%
\pgfpathlineto{\pgfqpoint{3.265721in}{2.965548in}}%
\pgfpathlineto{\pgfqpoint{3.270262in}{2.965548in}}%
\pgfpathlineto{\pgfqpoint{3.270262in}{2.962599in}}%
\pgfpathmoveto{\pgfqpoint{3.270262in}{2.959650in}}%
\pgfpathlineto{\pgfqpoint{3.270262in}{2.959650in}}%
\pgfpathlineto{\pgfqpoint{3.270262in}{2.962599in}}%
\pgfpathlineto{\pgfqpoint{3.274803in}{2.962599in}}%
\pgfpathlineto{\pgfqpoint{3.274803in}{2.959650in}}%
\pgfpathmoveto{\pgfqpoint{3.270262in}{2.962599in}}%
\pgfpathlineto{\pgfqpoint{3.270262in}{2.962599in}}%
\pgfpathlineto{\pgfqpoint{3.270262in}{2.965548in}}%
\pgfpathlineto{\pgfqpoint{3.274803in}{2.965548in}}%
\pgfpathlineto{\pgfqpoint{3.274803in}{2.962599in}}%
\pgfpathmoveto{\pgfqpoint{3.256639in}{2.965548in}}%
\pgfpathlineto{\pgfqpoint{3.256639in}{2.965548in}}%
\pgfpathlineto{\pgfqpoint{3.256639in}{2.968497in}}%
\pgfpathlineto{\pgfqpoint{3.261180in}{2.968497in}}%
\pgfpathlineto{\pgfqpoint{3.261180in}{2.965548in}}%
\pgfpathmoveto{\pgfqpoint{3.256639in}{2.968497in}}%
\pgfpathlineto{\pgfqpoint{3.256639in}{2.968497in}}%
\pgfpathlineto{\pgfqpoint{3.256639in}{2.971446in}}%
\pgfpathlineto{\pgfqpoint{3.261180in}{2.971446in}}%
\pgfpathlineto{\pgfqpoint{3.261180in}{2.968497in}}%
\pgfpathmoveto{\pgfqpoint{3.261180in}{2.965548in}}%
\pgfpathlineto{\pgfqpoint{3.261180in}{2.965548in}}%
\pgfpathlineto{\pgfqpoint{3.261180in}{2.968497in}}%
\pgfpathlineto{\pgfqpoint{3.265721in}{2.968497in}}%
\pgfpathlineto{\pgfqpoint{3.265721in}{2.965548in}}%
\pgfpathmoveto{\pgfqpoint{3.261180in}{2.968497in}}%
\pgfpathlineto{\pgfqpoint{3.261180in}{2.968497in}}%
\pgfpathlineto{\pgfqpoint{3.261180in}{2.971446in}}%
\pgfpathlineto{\pgfqpoint{3.265721in}{2.971446in}}%
\pgfpathlineto{\pgfqpoint{3.265721in}{2.968497in}}%
\pgfpathmoveto{\pgfqpoint{3.256639in}{2.971446in}}%
\pgfpathlineto{\pgfqpoint{3.256639in}{2.971446in}}%
\pgfpathlineto{\pgfqpoint{3.256639in}{2.974396in}}%
\pgfpathlineto{\pgfqpoint{3.261180in}{2.974396in}}%
\pgfpathlineto{\pgfqpoint{3.261180in}{2.971446in}}%
\pgfpathmoveto{\pgfqpoint{3.256639in}{2.974396in}}%
\pgfpathlineto{\pgfqpoint{3.256639in}{2.974396in}}%
\pgfpathlineto{\pgfqpoint{3.256639in}{2.977345in}}%
\pgfpathlineto{\pgfqpoint{3.261180in}{2.977345in}}%
\pgfpathlineto{\pgfqpoint{3.261180in}{2.974396in}}%
\pgfpathmoveto{\pgfqpoint{3.261180in}{2.971446in}}%
\pgfpathlineto{\pgfqpoint{3.261180in}{2.971446in}}%
\pgfpathlineto{\pgfqpoint{3.261180in}{2.974396in}}%
\pgfpathlineto{\pgfqpoint{3.265721in}{2.974396in}}%
\pgfpathlineto{\pgfqpoint{3.265721in}{2.971446in}}%
\pgfpathmoveto{\pgfqpoint{3.261180in}{2.974396in}}%
\pgfpathlineto{\pgfqpoint{3.261180in}{2.974396in}}%
\pgfpathlineto{\pgfqpoint{3.261180in}{2.977345in}}%
\pgfpathlineto{\pgfqpoint{3.265721in}{2.977345in}}%
\pgfpathlineto{\pgfqpoint{3.265721in}{2.974396in}}%
\pgfpathmoveto{\pgfqpoint{3.265721in}{2.965548in}}%
\pgfpathlineto{\pgfqpoint{3.265721in}{2.965548in}}%
\pgfpathlineto{\pgfqpoint{3.265721in}{2.968497in}}%
\pgfpathlineto{\pgfqpoint{3.270262in}{2.968497in}}%
\pgfpathlineto{\pgfqpoint{3.270262in}{2.965548in}}%
\pgfpathmoveto{\pgfqpoint{3.265721in}{2.968497in}}%
\pgfpathlineto{\pgfqpoint{3.265721in}{2.968497in}}%
\pgfpathlineto{\pgfqpoint{3.265721in}{2.971446in}}%
\pgfpathlineto{\pgfqpoint{3.270262in}{2.971446in}}%
\pgfpathlineto{\pgfqpoint{3.270262in}{2.968497in}}%
\pgfpathmoveto{\pgfqpoint{3.274803in}{2.953751in}}%
\pgfpathlineto{\pgfqpoint{3.274803in}{2.953751in}}%
\pgfpathlineto{\pgfqpoint{3.274803in}{2.956701in}}%
\pgfpathlineto{\pgfqpoint{3.279345in}{2.956701in}}%
\pgfpathlineto{\pgfqpoint{3.279345in}{2.953751in}}%
\pgfpathmoveto{\pgfqpoint{3.274803in}{2.956701in}}%
\pgfpathlineto{\pgfqpoint{3.274803in}{2.956701in}}%
\pgfpathlineto{\pgfqpoint{3.274803in}{2.959650in}}%
\pgfpathlineto{\pgfqpoint{3.279345in}{2.959650in}}%
\pgfpathlineto{\pgfqpoint{3.279345in}{2.956701in}}%
\pgfpathmoveto{\pgfqpoint{3.256639in}{2.977345in}}%
\pgfpathlineto{\pgfqpoint{3.256639in}{2.977345in}}%
\pgfpathlineto{\pgfqpoint{3.256639in}{2.980294in}}%
\pgfpathlineto{\pgfqpoint{3.261180in}{2.980294in}}%
\pgfpathlineto{\pgfqpoint{3.261180in}{2.977345in}}%
\pgfpathmoveto{\pgfqpoint{3.256639in}{2.980294in}}%
\pgfpathlineto{\pgfqpoint{3.256639in}{2.980294in}}%
\pgfpathlineto{\pgfqpoint{3.256639in}{2.983243in}}%
\pgfpathlineto{\pgfqpoint{3.261180in}{2.983243in}}%
\pgfpathlineto{\pgfqpoint{3.261180in}{2.980294in}}%
\pgfpathmoveto{\pgfqpoint{3.229392in}{3.000938in}}%
\pgfpathlineto{\pgfqpoint{3.229392in}{3.000938in}}%
\pgfpathlineto{\pgfqpoint{3.229392in}{3.003887in}}%
\pgfpathlineto{\pgfqpoint{3.233933in}{3.003887in}}%
\pgfpathlineto{\pgfqpoint{3.233933in}{3.000938in}}%
\pgfpathmoveto{\pgfqpoint{3.229392in}{3.003887in}}%
\pgfpathlineto{\pgfqpoint{3.229392in}{3.003887in}}%
\pgfpathlineto{\pgfqpoint{3.229392in}{3.006836in}}%
\pgfpathlineto{\pgfqpoint{3.233933in}{3.006836in}}%
\pgfpathlineto{\pgfqpoint{3.233933in}{3.003887in}}%
\pgfpathmoveto{\pgfqpoint{3.233933in}{3.000938in}}%
\pgfpathlineto{\pgfqpoint{3.233933in}{3.000938in}}%
\pgfpathlineto{\pgfqpoint{3.233933in}{3.003887in}}%
\pgfpathlineto{\pgfqpoint{3.238474in}{3.003887in}}%
\pgfpathlineto{\pgfqpoint{3.238474in}{3.000938in}}%
\pgfpathmoveto{\pgfqpoint{3.233933in}{3.003887in}}%
\pgfpathlineto{\pgfqpoint{3.233933in}{3.003887in}}%
\pgfpathlineto{\pgfqpoint{3.233933in}{3.006836in}}%
\pgfpathlineto{\pgfqpoint{3.238474in}{3.006836in}}%
\pgfpathlineto{\pgfqpoint{3.238474in}{3.003887in}}%
\pgfpathmoveto{\pgfqpoint{3.229392in}{3.006836in}}%
\pgfpathlineto{\pgfqpoint{3.229392in}{3.006836in}}%
\pgfpathlineto{\pgfqpoint{3.229392in}{3.009785in}}%
\pgfpathlineto{\pgfqpoint{3.233933in}{3.009785in}}%
\pgfpathlineto{\pgfqpoint{3.233933in}{3.006836in}}%
\pgfpathmoveto{\pgfqpoint{3.229392in}{3.009785in}}%
\pgfpathlineto{\pgfqpoint{3.229392in}{3.009785in}}%
\pgfpathlineto{\pgfqpoint{3.229392in}{3.012735in}}%
\pgfpathlineto{\pgfqpoint{3.233933in}{3.012735in}}%
\pgfpathlineto{\pgfqpoint{3.233933in}{3.009785in}}%
\pgfpathmoveto{\pgfqpoint{3.233933in}{3.006836in}}%
\pgfpathlineto{\pgfqpoint{3.233933in}{3.006836in}}%
\pgfpathlineto{\pgfqpoint{3.233933in}{3.009785in}}%
\pgfpathlineto{\pgfqpoint{3.238474in}{3.009785in}}%
\pgfpathlineto{\pgfqpoint{3.238474in}{3.006836in}}%
\pgfpathmoveto{\pgfqpoint{3.233933in}{3.009785in}}%
\pgfpathlineto{\pgfqpoint{3.233933in}{3.009785in}}%
\pgfpathlineto{\pgfqpoint{3.233933in}{3.012735in}}%
\pgfpathlineto{\pgfqpoint{3.238474in}{3.012735in}}%
\pgfpathlineto{\pgfqpoint{3.238474in}{3.009785in}}%
\pgfpathmoveto{\pgfqpoint{3.220310in}{3.012735in}}%
\pgfpathlineto{\pgfqpoint{3.220310in}{3.012735in}}%
\pgfpathlineto{\pgfqpoint{3.220310in}{3.015684in}}%
\pgfpathlineto{\pgfqpoint{3.224851in}{3.015684in}}%
\pgfpathlineto{\pgfqpoint{3.224851in}{3.012735in}}%
\pgfpathmoveto{\pgfqpoint{3.220310in}{3.015684in}}%
\pgfpathlineto{\pgfqpoint{3.220310in}{3.015684in}}%
\pgfpathlineto{\pgfqpoint{3.220310in}{3.018633in}}%
\pgfpathlineto{\pgfqpoint{3.224851in}{3.018633in}}%
\pgfpathlineto{\pgfqpoint{3.224851in}{3.015684in}}%
\pgfpathmoveto{\pgfqpoint{3.224851in}{3.012735in}}%
\pgfpathlineto{\pgfqpoint{3.224851in}{3.012735in}}%
\pgfpathlineto{\pgfqpoint{3.224851in}{3.015684in}}%
\pgfpathlineto{\pgfqpoint{3.229392in}{3.015684in}}%
\pgfpathlineto{\pgfqpoint{3.229392in}{3.012735in}}%
\pgfpathmoveto{\pgfqpoint{3.224851in}{3.015684in}}%
\pgfpathlineto{\pgfqpoint{3.224851in}{3.015684in}}%
\pgfpathlineto{\pgfqpoint{3.224851in}{3.018633in}}%
\pgfpathlineto{\pgfqpoint{3.229392in}{3.018633in}}%
\pgfpathlineto{\pgfqpoint{3.229392in}{3.015684in}}%
\pgfpathmoveto{\pgfqpoint{3.220310in}{3.018633in}}%
\pgfpathlineto{\pgfqpoint{3.220310in}{3.018633in}}%
\pgfpathlineto{\pgfqpoint{3.220310in}{3.021582in}}%
\pgfpathlineto{\pgfqpoint{3.224851in}{3.021582in}}%
\pgfpathlineto{\pgfqpoint{3.224851in}{3.018633in}}%
\pgfpathmoveto{\pgfqpoint{3.220310in}{3.021582in}}%
\pgfpathlineto{\pgfqpoint{3.220310in}{3.021582in}}%
\pgfpathlineto{\pgfqpoint{3.220310in}{3.024531in}}%
\pgfpathlineto{\pgfqpoint{3.224851in}{3.024531in}}%
\pgfpathlineto{\pgfqpoint{3.224851in}{3.021582in}}%
\pgfpathmoveto{\pgfqpoint{3.224851in}{3.018633in}}%
\pgfpathlineto{\pgfqpoint{3.224851in}{3.018633in}}%
\pgfpathlineto{\pgfqpoint{3.224851in}{3.021582in}}%
\pgfpathlineto{\pgfqpoint{3.229392in}{3.021582in}}%
\pgfpathlineto{\pgfqpoint{3.229392in}{3.018633in}}%
\pgfpathmoveto{\pgfqpoint{3.224851in}{3.021582in}}%
\pgfpathlineto{\pgfqpoint{3.224851in}{3.021582in}}%
\pgfpathlineto{\pgfqpoint{3.224851in}{3.024531in}}%
\pgfpathlineto{\pgfqpoint{3.229392in}{3.024531in}}%
\pgfpathlineto{\pgfqpoint{3.229392in}{3.021582in}}%
\pgfpathmoveto{\pgfqpoint{3.229392in}{3.012735in}}%
\pgfpathlineto{\pgfqpoint{3.229392in}{3.012735in}}%
\pgfpathlineto{\pgfqpoint{3.229392in}{3.015684in}}%
\pgfpathlineto{\pgfqpoint{3.233933in}{3.015684in}}%
\pgfpathlineto{\pgfqpoint{3.233933in}{3.012735in}}%
\pgfpathmoveto{\pgfqpoint{3.229392in}{3.015684in}}%
\pgfpathlineto{\pgfqpoint{3.229392in}{3.015684in}}%
\pgfpathlineto{\pgfqpoint{3.229392in}{3.018633in}}%
\pgfpathlineto{\pgfqpoint{3.233933in}{3.018633in}}%
\pgfpathlineto{\pgfqpoint{3.233933in}{3.015684in}}%
\pgfpathmoveto{\pgfqpoint{3.238474in}{3.000938in}}%
\pgfpathlineto{\pgfqpoint{3.238474in}{3.000938in}}%
\pgfpathlineto{\pgfqpoint{3.238474in}{3.003887in}}%
\pgfpathlineto{\pgfqpoint{3.243015in}{3.003887in}}%
\pgfpathlineto{\pgfqpoint{3.243015in}{3.000938in}}%
\pgfpathmoveto{\pgfqpoint{3.238474in}{3.003887in}}%
\pgfpathlineto{\pgfqpoint{3.238474in}{3.003887in}}%
\pgfpathlineto{\pgfqpoint{3.238474in}{3.006836in}}%
\pgfpathlineto{\pgfqpoint{3.243015in}{3.006836in}}%
\pgfpathlineto{\pgfqpoint{3.243015in}{3.003887in}}%
\pgfpathmoveto{\pgfqpoint{3.220310in}{3.024531in}}%
\pgfpathlineto{\pgfqpoint{3.220310in}{3.024531in}}%
\pgfpathlineto{\pgfqpoint{3.220310in}{3.027480in}}%
\pgfpathlineto{\pgfqpoint{3.224851in}{3.027480in}}%
\pgfpathlineto{\pgfqpoint{3.224851in}{3.024531in}}%
\pgfpathmoveto{\pgfqpoint{3.220310in}{3.027480in}}%
\pgfpathlineto{\pgfqpoint{3.220310in}{3.027480in}}%
\pgfpathlineto{\pgfqpoint{3.220310in}{3.030430in}}%
\pgfpathlineto{\pgfqpoint{3.224851in}{3.030430in}}%
\pgfpathlineto{\pgfqpoint{3.224851in}{3.027480in}}%
\pgfpathmoveto{\pgfqpoint{3.365627in}{2.198748in}}%
\pgfpathlineto{\pgfqpoint{3.365627in}{2.198748in}}%
\pgfpathlineto{\pgfqpoint{3.365627in}{2.201697in}}%
\pgfpathlineto{\pgfqpoint{3.370168in}{2.201697in}}%
\pgfpathlineto{\pgfqpoint{3.370168in}{2.198748in}}%
\pgfpathmoveto{\pgfqpoint{3.365627in}{2.201697in}}%
\pgfpathlineto{\pgfqpoint{3.365627in}{2.201697in}}%
\pgfpathlineto{\pgfqpoint{3.365627in}{2.204647in}}%
\pgfpathlineto{\pgfqpoint{3.370168in}{2.204647in}}%
\pgfpathlineto{\pgfqpoint{3.370168in}{2.201697in}}%
\pgfpathmoveto{\pgfqpoint{3.370168in}{2.201697in}}%
\pgfpathlineto{\pgfqpoint{3.370168in}{2.201697in}}%
\pgfpathlineto{\pgfqpoint{3.370168in}{2.204647in}}%
\pgfpathlineto{\pgfqpoint{3.374709in}{2.204647in}}%
\pgfpathlineto{\pgfqpoint{3.374709in}{2.201697in}}%
\pgfpathmoveto{\pgfqpoint{3.365627in}{2.204647in}}%
\pgfpathlineto{\pgfqpoint{3.365627in}{2.204647in}}%
\pgfpathlineto{\pgfqpoint{3.365627in}{2.207596in}}%
\pgfpathlineto{\pgfqpoint{3.370168in}{2.207596in}}%
\pgfpathlineto{\pgfqpoint{3.370168in}{2.204647in}}%
\pgfpathmoveto{\pgfqpoint{3.365627in}{2.207596in}}%
\pgfpathlineto{\pgfqpoint{3.365627in}{2.207596in}}%
\pgfpathlineto{\pgfqpoint{3.365627in}{2.210545in}}%
\pgfpathlineto{\pgfqpoint{3.370168in}{2.210545in}}%
\pgfpathlineto{\pgfqpoint{3.370168in}{2.207596in}}%
\pgfpathmoveto{\pgfqpoint{3.370168in}{2.204647in}}%
\pgfpathlineto{\pgfqpoint{3.370168in}{2.204647in}}%
\pgfpathlineto{\pgfqpoint{3.370168in}{2.207596in}}%
\pgfpathlineto{\pgfqpoint{3.374709in}{2.207596in}}%
\pgfpathlineto{\pgfqpoint{3.374709in}{2.204647in}}%
\pgfpathmoveto{\pgfqpoint{3.370168in}{2.207596in}}%
\pgfpathlineto{\pgfqpoint{3.370168in}{2.207596in}}%
\pgfpathlineto{\pgfqpoint{3.370168in}{2.210545in}}%
\pgfpathlineto{\pgfqpoint{3.374709in}{2.210545in}}%
\pgfpathlineto{\pgfqpoint{3.374709in}{2.207596in}}%
\pgfpathmoveto{\pgfqpoint{3.374709in}{2.204647in}}%
\pgfpathlineto{\pgfqpoint{3.374709in}{2.204647in}}%
\pgfpathlineto{\pgfqpoint{3.374709in}{2.207596in}}%
\pgfpathlineto{\pgfqpoint{3.379250in}{2.207596in}}%
\pgfpathlineto{\pgfqpoint{3.379250in}{2.204647in}}%
\pgfpathmoveto{\pgfqpoint{3.374709in}{2.207596in}}%
\pgfpathlineto{\pgfqpoint{3.374709in}{2.207596in}}%
\pgfpathlineto{\pgfqpoint{3.374709in}{2.210545in}}%
\pgfpathlineto{\pgfqpoint{3.379250in}{2.210545in}}%
\pgfpathlineto{\pgfqpoint{3.379250in}{2.207596in}}%
\pgfpathmoveto{\pgfqpoint{3.379250in}{2.207596in}}%
\pgfpathlineto{\pgfqpoint{3.379250in}{2.207596in}}%
\pgfpathlineto{\pgfqpoint{3.379250in}{2.210545in}}%
\pgfpathlineto{\pgfqpoint{3.383791in}{2.210545in}}%
\pgfpathlineto{\pgfqpoint{3.383791in}{2.207596in}}%
\pgfpathmoveto{\pgfqpoint{3.374709in}{2.210545in}}%
\pgfpathlineto{\pgfqpoint{3.374709in}{2.210545in}}%
\pgfpathlineto{\pgfqpoint{3.374709in}{2.213494in}}%
\pgfpathlineto{\pgfqpoint{3.379250in}{2.213494in}}%
\pgfpathlineto{\pgfqpoint{3.379250in}{2.210545in}}%
\pgfpathmoveto{\pgfqpoint{3.374709in}{2.213494in}}%
\pgfpathlineto{\pgfqpoint{3.374709in}{2.213494in}}%
\pgfpathlineto{\pgfqpoint{3.374709in}{2.216443in}}%
\pgfpathlineto{\pgfqpoint{3.379250in}{2.216443in}}%
\pgfpathlineto{\pgfqpoint{3.379250in}{2.213494in}}%
\pgfpathmoveto{\pgfqpoint{3.379250in}{2.210545in}}%
\pgfpathlineto{\pgfqpoint{3.379250in}{2.210545in}}%
\pgfpathlineto{\pgfqpoint{3.379250in}{2.213494in}}%
\pgfpathlineto{\pgfqpoint{3.383791in}{2.213494in}}%
\pgfpathlineto{\pgfqpoint{3.383791in}{2.210545in}}%
\pgfpathmoveto{\pgfqpoint{3.379250in}{2.213494in}}%
\pgfpathlineto{\pgfqpoint{3.379250in}{2.213494in}}%
\pgfpathlineto{\pgfqpoint{3.379250in}{2.216443in}}%
\pgfpathlineto{\pgfqpoint{3.383791in}{2.216443in}}%
\pgfpathlineto{\pgfqpoint{3.383791in}{2.213494in}}%
\pgfpathmoveto{\pgfqpoint{3.383791in}{2.210545in}}%
\pgfpathlineto{\pgfqpoint{3.383791in}{2.210545in}}%
\pgfpathlineto{\pgfqpoint{3.383791in}{2.213494in}}%
\pgfpathlineto{\pgfqpoint{3.388332in}{2.213494in}}%
\pgfpathlineto{\pgfqpoint{3.388332in}{2.210545in}}%
\pgfpathmoveto{\pgfqpoint{3.383791in}{2.213494in}}%
\pgfpathlineto{\pgfqpoint{3.383791in}{2.213494in}}%
\pgfpathlineto{\pgfqpoint{3.383791in}{2.216443in}}%
\pgfpathlineto{\pgfqpoint{3.388332in}{2.216443in}}%
\pgfpathlineto{\pgfqpoint{3.388332in}{2.213494in}}%
\pgfpathmoveto{\pgfqpoint{3.388332in}{2.213494in}}%
\pgfpathlineto{\pgfqpoint{3.388332in}{2.213494in}}%
\pgfpathlineto{\pgfqpoint{3.388332in}{2.216443in}}%
\pgfpathlineto{\pgfqpoint{3.392873in}{2.216443in}}%
\pgfpathlineto{\pgfqpoint{3.392873in}{2.213494in}}%
\pgfpathmoveto{\pgfqpoint{3.383791in}{2.216443in}}%
\pgfpathlineto{\pgfqpoint{3.383791in}{2.216443in}}%
\pgfpathlineto{\pgfqpoint{3.383791in}{2.219393in}}%
\pgfpathlineto{\pgfqpoint{3.388332in}{2.219393in}}%
\pgfpathlineto{\pgfqpoint{3.388332in}{2.216443in}}%
\pgfpathmoveto{\pgfqpoint{3.383791in}{2.219393in}}%
\pgfpathlineto{\pgfqpoint{3.383791in}{2.219393in}}%
\pgfpathlineto{\pgfqpoint{3.383791in}{2.222342in}}%
\pgfpathlineto{\pgfqpoint{3.388332in}{2.222342in}}%
\pgfpathlineto{\pgfqpoint{3.388332in}{2.219393in}}%
\pgfpathmoveto{\pgfqpoint{3.388332in}{2.216443in}}%
\pgfpathlineto{\pgfqpoint{3.388332in}{2.216443in}}%
\pgfpathlineto{\pgfqpoint{3.388332in}{2.219393in}}%
\pgfpathlineto{\pgfqpoint{3.392873in}{2.219393in}}%
\pgfpathlineto{\pgfqpoint{3.392873in}{2.216443in}}%
\pgfpathmoveto{\pgfqpoint{3.388332in}{2.219393in}}%
\pgfpathlineto{\pgfqpoint{3.388332in}{2.219393in}}%
\pgfpathlineto{\pgfqpoint{3.388332in}{2.222342in}}%
\pgfpathlineto{\pgfqpoint{3.392873in}{2.222342in}}%
\pgfpathlineto{\pgfqpoint{3.392873in}{2.219393in}}%
\pgfpathmoveto{\pgfqpoint{3.392873in}{2.216443in}}%
\pgfpathlineto{\pgfqpoint{3.392873in}{2.216443in}}%
\pgfpathlineto{\pgfqpoint{3.392873in}{2.219393in}}%
\pgfpathlineto{\pgfqpoint{3.397414in}{2.219393in}}%
\pgfpathlineto{\pgfqpoint{3.397414in}{2.216443in}}%
\pgfpathmoveto{\pgfqpoint{3.392873in}{2.219393in}}%
\pgfpathlineto{\pgfqpoint{3.392873in}{2.219393in}}%
\pgfpathlineto{\pgfqpoint{3.392873in}{2.222342in}}%
\pgfpathlineto{\pgfqpoint{3.397414in}{2.222342in}}%
\pgfpathlineto{\pgfqpoint{3.397414in}{2.219393in}}%
\pgfpathmoveto{\pgfqpoint{3.397414in}{2.219393in}}%
\pgfpathlineto{\pgfqpoint{3.397414in}{2.219393in}}%
\pgfpathlineto{\pgfqpoint{3.397414in}{2.222342in}}%
\pgfpathlineto{\pgfqpoint{3.401955in}{2.222342in}}%
\pgfpathlineto{\pgfqpoint{3.401955in}{2.219393in}}%
\pgfpathmoveto{\pgfqpoint{3.392873in}{2.222342in}}%
\pgfpathlineto{\pgfqpoint{3.392873in}{2.222342in}}%
\pgfpathlineto{\pgfqpoint{3.392873in}{2.225291in}}%
\pgfpathlineto{\pgfqpoint{3.397414in}{2.225291in}}%
\pgfpathlineto{\pgfqpoint{3.397414in}{2.222342in}}%
\pgfpathmoveto{\pgfqpoint{3.392873in}{2.225291in}}%
\pgfpathlineto{\pgfqpoint{3.392873in}{2.225291in}}%
\pgfpathlineto{\pgfqpoint{3.392873in}{2.228240in}}%
\pgfpathlineto{\pgfqpoint{3.397414in}{2.228240in}}%
\pgfpathlineto{\pgfqpoint{3.397414in}{2.225291in}}%
\pgfpathmoveto{\pgfqpoint{3.397414in}{2.222342in}}%
\pgfpathlineto{\pgfqpoint{3.397414in}{2.222342in}}%
\pgfpathlineto{\pgfqpoint{3.397414in}{2.225291in}}%
\pgfpathlineto{\pgfqpoint{3.401955in}{2.225291in}}%
\pgfpathlineto{\pgfqpoint{3.401955in}{2.222342in}}%
\pgfpathmoveto{\pgfqpoint{3.397414in}{2.225291in}}%
\pgfpathlineto{\pgfqpoint{3.397414in}{2.225291in}}%
\pgfpathlineto{\pgfqpoint{3.397414in}{2.228240in}}%
\pgfpathlineto{\pgfqpoint{3.401955in}{2.228240in}}%
\pgfpathlineto{\pgfqpoint{3.401955in}{2.225291in}}%
\pgfpathmoveto{\pgfqpoint{3.401955in}{2.222342in}}%
\pgfpathlineto{\pgfqpoint{3.401955in}{2.222342in}}%
\pgfpathlineto{\pgfqpoint{3.401955in}{2.225291in}}%
\pgfpathlineto{\pgfqpoint{3.406496in}{2.225291in}}%
\pgfpathlineto{\pgfqpoint{3.406496in}{2.222342in}}%
\pgfpathmoveto{\pgfqpoint{3.401955in}{2.225291in}}%
\pgfpathlineto{\pgfqpoint{3.401955in}{2.225291in}}%
\pgfpathlineto{\pgfqpoint{3.401955in}{2.228240in}}%
\pgfpathlineto{\pgfqpoint{3.406496in}{2.228240in}}%
\pgfpathlineto{\pgfqpoint{3.406496in}{2.225291in}}%
\pgfpathmoveto{\pgfqpoint{3.406496in}{2.225291in}}%
\pgfpathlineto{\pgfqpoint{3.406496in}{2.225291in}}%
\pgfpathlineto{\pgfqpoint{3.406496in}{2.228240in}}%
\pgfpathlineto{\pgfqpoint{3.411037in}{2.228240in}}%
\pgfpathlineto{\pgfqpoint{3.411037in}{2.225291in}}%
\pgfpathmoveto{\pgfqpoint{3.401955in}{2.228240in}}%
\pgfpathlineto{\pgfqpoint{3.401955in}{2.228240in}}%
\pgfpathlineto{\pgfqpoint{3.401955in}{2.231189in}}%
\pgfpathlineto{\pgfqpoint{3.406496in}{2.231189in}}%
\pgfpathlineto{\pgfqpoint{3.406496in}{2.228240in}}%
\pgfpathmoveto{\pgfqpoint{3.401955in}{2.231189in}}%
\pgfpathlineto{\pgfqpoint{3.401955in}{2.231189in}}%
\pgfpathlineto{\pgfqpoint{3.401955in}{2.234139in}}%
\pgfpathlineto{\pgfqpoint{3.406496in}{2.234139in}}%
\pgfpathlineto{\pgfqpoint{3.406496in}{2.231189in}}%
\pgfpathmoveto{\pgfqpoint{3.406496in}{2.228240in}}%
\pgfpathlineto{\pgfqpoint{3.406496in}{2.228240in}}%
\pgfpathlineto{\pgfqpoint{3.406496in}{2.231189in}}%
\pgfpathlineto{\pgfqpoint{3.411037in}{2.231189in}}%
\pgfpathlineto{\pgfqpoint{3.411037in}{2.228240in}}%
\pgfpathmoveto{\pgfqpoint{3.406496in}{2.231189in}}%
\pgfpathlineto{\pgfqpoint{3.406496in}{2.231189in}}%
\pgfpathlineto{\pgfqpoint{3.406496in}{2.234139in}}%
\pgfpathlineto{\pgfqpoint{3.411037in}{2.234139in}}%
\pgfpathlineto{\pgfqpoint{3.411037in}{2.231189in}}%
\pgfpathmoveto{\pgfqpoint{3.411037in}{2.228240in}}%
\pgfpathlineto{\pgfqpoint{3.411037in}{2.228240in}}%
\pgfpathlineto{\pgfqpoint{3.411037in}{2.231189in}}%
\pgfpathlineto{\pgfqpoint{3.415578in}{2.231189in}}%
\pgfpathlineto{\pgfqpoint{3.415578in}{2.228240in}}%
\pgfpathmoveto{\pgfqpoint{3.411037in}{2.231189in}}%
\pgfpathlineto{\pgfqpoint{3.411037in}{2.231189in}}%
\pgfpathlineto{\pgfqpoint{3.411037in}{2.234139in}}%
\pgfpathlineto{\pgfqpoint{3.415578in}{2.234139in}}%
\pgfpathlineto{\pgfqpoint{3.415578in}{2.231189in}}%
\pgfpathmoveto{\pgfqpoint{3.415578in}{2.231189in}}%
\pgfpathlineto{\pgfqpoint{3.415578in}{2.231189in}}%
\pgfpathlineto{\pgfqpoint{3.415578in}{2.234139in}}%
\pgfpathlineto{\pgfqpoint{3.420119in}{2.234139in}}%
\pgfpathlineto{\pgfqpoint{3.420119in}{2.231189in}}%
\pgfpathmoveto{\pgfqpoint{3.411037in}{2.234139in}}%
\pgfpathlineto{\pgfqpoint{3.411037in}{2.234139in}}%
\pgfpathlineto{\pgfqpoint{3.411037in}{2.237088in}}%
\pgfpathlineto{\pgfqpoint{3.415578in}{2.237088in}}%
\pgfpathlineto{\pgfqpoint{3.415578in}{2.234139in}}%
\pgfpathmoveto{\pgfqpoint{3.411037in}{2.237088in}}%
\pgfpathlineto{\pgfqpoint{3.411037in}{2.237088in}}%
\pgfpathlineto{\pgfqpoint{3.411037in}{2.240037in}}%
\pgfpathlineto{\pgfqpoint{3.415578in}{2.240037in}}%
\pgfpathlineto{\pgfqpoint{3.415578in}{2.237088in}}%
\pgfpathmoveto{\pgfqpoint{3.415578in}{2.234139in}}%
\pgfpathlineto{\pgfqpoint{3.415578in}{2.234139in}}%
\pgfpathlineto{\pgfqpoint{3.415578in}{2.237088in}}%
\pgfpathlineto{\pgfqpoint{3.420119in}{2.237088in}}%
\pgfpathlineto{\pgfqpoint{3.420119in}{2.234139in}}%
\pgfpathmoveto{\pgfqpoint{3.415578in}{2.237088in}}%
\pgfpathlineto{\pgfqpoint{3.415578in}{2.237088in}}%
\pgfpathlineto{\pgfqpoint{3.415578in}{2.240037in}}%
\pgfpathlineto{\pgfqpoint{3.420119in}{2.240037in}}%
\pgfpathlineto{\pgfqpoint{3.420119in}{2.237088in}}%
\pgfpathmoveto{\pgfqpoint{3.420119in}{2.234139in}}%
\pgfpathlineto{\pgfqpoint{3.420119in}{2.234139in}}%
\pgfpathlineto{\pgfqpoint{3.420119in}{2.237088in}}%
\pgfpathlineto{\pgfqpoint{3.424660in}{2.237088in}}%
\pgfpathlineto{\pgfqpoint{3.424660in}{2.234139in}}%
\pgfpathmoveto{\pgfqpoint{3.420119in}{2.237088in}}%
\pgfpathlineto{\pgfqpoint{3.420119in}{2.237088in}}%
\pgfpathlineto{\pgfqpoint{3.420119in}{2.240037in}}%
\pgfpathlineto{\pgfqpoint{3.424660in}{2.240037in}}%
\pgfpathlineto{\pgfqpoint{3.424660in}{2.237088in}}%
\pgfpathmoveto{\pgfqpoint{3.424660in}{2.237088in}}%
\pgfpathlineto{\pgfqpoint{3.424660in}{2.237088in}}%
\pgfpathlineto{\pgfqpoint{3.424660in}{2.240037in}}%
\pgfpathlineto{\pgfqpoint{3.429201in}{2.240037in}}%
\pgfpathlineto{\pgfqpoint{3.429201in}{2.237088in}}%
\pgfpathmoveto{\pgfqpoint{3.420119in}{2.240037in}}%
\pgfpathlineto{\pgfqpoint{3.420119in}{2.240037in}}%
\pgfpathlineto{\pgfqpoint{3.420119in}{2.242986in}}%
\pgfpathlineto{\pgfqpoint{3.424660in}{2.242986in}}%
\pgfpathlineto{\pgfqpoint{3.424660in}{2.240037in}}%
\pgfpathmoveto{\pgfqpoint{3.420119in}{2.242986in}}%
\pgfpathlineto{\pgfqpoint{3.420119in}{2.242986in}}%
\pgfpathlineto{\pgfqpoint{3.420119in}{2.245935in}}%
\pgfpathlineto{\pgfqpoint{3.424660in}{2.245935in}}%
\pgfpathlineto{\pgfqpoint{3.424660in}{2.242986in}}%
\pgfpathmoveto{\pgfqpoint{3.424660in}{2.240037in}}%
\pgfpathlineto{\pgfqpoint{3.424660in}{2.240037in}}%
\pgfpathlineto{\pgfqpoint{3.424660in}{2.242986in}}%
\pgfpathlineto{\pgfqpoint{3.429201in}{2.242986in}}%
\pgfpathlineto{\pgfqpoint{3.429201in}{2.240037in}}%
\pgfpathmoveto{\pgfqpoint{3.424660in}{2.242986in}}%
\pgfpathlineto{\pgfqpoint{3.424660in}{2.242986in}}%
\pgfpathlineto{\pgfqpoint{3.424660in}{2.245935in}}%
\pgfpathlineto{\pgfqpoint{3.429201in}{2.245935in}}%
\pgfpathlineto{\pgfqpoint{3.429201in}{2.242986in}}%
\pgfpathmoveto{\pgfqpoint{3.429201in}{2.240037in}}%
\pgfpathlineto{\pgfqpoint{3.429201in}{2.240037in}}%
\pgfpathlineto{\pgfqpoint{3.429201in}{2.242986in}}%
\pgfpathlineto{\pgfqpoint{3.433742in}{2.242986in}}%
\pgfpathlineto{\pgfqpoint{3.433742in}{2.240037in}}%
\pgfpathmoveto{\pgfqpoint{3.429201in}{2.242986in}}%
\pgfpathlineto{\pgfqpoint{3.429201in}{2.242986in}}%
\pgfpathlineto{\pgfqpoint{3.429201in}{2.245935in}}%
\pgfpathlineto{\pgfqpoint{3.433742in}{2.245935in}}%
\pgfpathlineto{\pgfqpoint{3.433742in}{2.242986in}}%
\pgfpathmoveto{\pgfqpoint{3.433742in}{2.242986in}}%
\pgfpathlineto{\pgfqpoint{3.433742in}{2.242986in}}%
\pgfpathlineto{\pgfqpoint{3.433742in}{2.245935in}}%
\pgfpathlineto{\pgfqpoint{3.438283in}{2.245935in}}%
\pgfpathlineto{\pgfqpoint{3.438283in}{2.242986in}}%
\pgfpathmoveto{\pgfqpoint{3.429201in}{2.245935in}}%
\pgfpathlineto{\pgfqpoint{3.429201in}{2.245935in}}%
\pgfpathlineto{\pgfqpoint{3.429201in}{2.248884in}}%
\pgfpathlineto{\pgfqpoint{3.433742in}{2.248884in}}%
\pgfpathlineto{\pgfqpoint{3.433742in}{2.245935in}}%
\pgfpathmoveto{\pgfqpoint{3.429201in}{2.248884in}}%
\pgfpathlineto{\pgfqpoint{3.429201in}{2.248884in}}%
\pgfpathlineto{\pgfqpoint{3.429201in}{2.251834in}}%
\pgfpathlineto{\pgfqpoint{3.433742in}{2.251834in}}%
\pgfpathlineto{\pgfqpoint{3.433742in}{2.248884in}}%
\pgfpathmoveto{\pgfqpoint{3.433742in}{2.245935in}}%
\pgfpathlineto{\pgfqpoint{3.433742in}{2.245935in}}%
\pgfpathlineto{\pgfqpoint{3.433742in}{2.248884in}}%
\pgfpathlineto{\pgfqpoint{3.438283in}{2.248884in}}%
\pgfpathlineto{\pgfqpoint{3.438283in}{2.245935in}}%
\pgfpathmoveto{\pgfqpoint{3.433742in}{2.248884in}}%
\pgfpathlineto{\pgfqpoint{3.433742in}{2.248884in}}%
\pgfpathlineto{\pgfqpoint{3.433742in}{2.251834in}}%
\pgfpathlineto{\pgfqpoint{3.438283in}{2.251834in}}%
\pgfpathlineto{\pgfqpoint{3.438283in}{2.248884in}}%
\pgfpathmoveto{\pgfqpoint{3.438283in}{2.245935in}}%
\pgfpathlineto{\pgfqpoint{3.438283in}{2.245935in}}%
\pgfpathlineto{\pgfqpoint{3.438283in}{2.248884in}}%
\pgfpathlineto{\pgfqpoint{3.442824in}{2.248884in}}%
\pgfpathlineto{\pgfqpoint{3.442824in}{2.245935in}}%
\pgfpathmoveto{\pgfqpoint{3.438283in}{2.248884in}}%
\pgfpathlineto{\pgfqpoint{3.438283in}{2.248884in}}%
\pgfpathlineto{\pgfqpoint{3.438283in}{2.251834in}}%
\pgfpathlineto{\pgfqpoint{3.442824in}{2.251834in}}%
\pgfpathlineto{\pgfqpoint{3.442824in}{2.248884in}}%
\pgfpathmoveto{\pgfqpoint{3.442824in}{2.248884in}}%
\pgfpathlineto{\pgfqpoint{3.442824in}{2.248884in}}%
\pgfpathlineto{\pgfqpoint{3.442824in}{2.251834in}}%
\pgfpathlineto{\pgfqpoint{3.447365in}{2.251834in}}%
\pgfpathlineto{\pgfqpoint{3.447365in}{2.248884in}}%
\pgfpathmoveto{\pgfqpoint{3.438283in}{2.251834in}}%
\pgfpathlineto{\pgfqpoint{3.438283in}{2.251834in}}%
\pgfpathlineto{\pgfqpoint{3.438283in}{2.254783in}}%
\pgfpathlineto{\pgfqpoint{3.442824in}{2.254783in}}%
\pgfpathlineto{\pgfqpoint{3.442824in}{2.251834in}}%
\pgfpathmoveto{\pgfqpoint{3.438283in}{2.254783in}}%
\pgfpathlineto{\pgfqpoint{3.438283in}{2.254783in}}%
\pgfpathlineto{\pgfqpoint{3.438283in}{2.257732in}}%
\pgfpathlineto{\pgfqpoint{3.442824in}{2.257732in}}%
\pgfpathlineto{\pgfqpoint{3.442824in}{2.254783in}}%
\pgfpathmoveto{\pgfqpoint{3.442824in}{2.251834in}}%
\pgfpathlineto{\pgfqpoint{3.442824in}{2.251834in}}%
\pgfpathlineto{\pgfqpoint{3.442824in}{2.254783in}}%
\pgfpathlineto{\pgfqpoint{3.447365in}{2.254783in}}%
\pgfpathlineto{\pgfqpoint{3.447365in}{2.251834in}}%
\pgfpathmoveto{\pgfqpoint{3.442824in}{2.254783in}}%
\pgfpathlineto{\pgfqpoint{3.442824in}{2.254783in}}%
\pgfpathlineto{\pgfqpoint{3.442824in}{2.257732in}}%
\pgfpathlineto{\pgfqpoint{3.447365in}{2.257732in}}%
\pgfpathlineto{\pgfqpoint{3.447365in}{2.254783in}}%
\pgfpathmoveto{\pgfqpoint{3.447365in}{2.251834in}}%
\pgfpathlineto{\pgfqpoint{3.447365in}{2.251834in}}%
\pgfpathlineto{\pgfqpoint{3.447365in}{2.254783in}}%
\pgfpathlineto{\pgfqpoint{3.451906in}{2.254783in}}%
\pgfpathlineto{\pgfqpoint{3.451906in}{2.251834in}}%
\pgfpathmoveto{\pgfqpoint{3.447365in}{2.254783in}}%
\pgfpathlineto{\pgfqpoint{3.447365in}{2.254783in}}%
\pgfpathlineto{\pgfqpoint{3.447365in}{2.257732in}}%
\pgfpathlineto{\pgfqpoint{3.451906in}{2.257732in}}%
\pgfpathlineto{\pgfqpoint{3.451906in}{2.254783in}}%
\pgfpathmoveto{\pgfqpoint{3.451906in}{2.254783in}}%
\pgfpathlineto{\pgfqpoint{3.451906in}{2.254783in}}%
\pgfpathlineto{\pgfqpoint{3.451906in}{2.257732in}}%
\pgfpathlineto{\pgfqpoint{3.456447in}{2.257732in}}%
\pgfpathlineto{\pgfqpoint{3.456447in}{2.254783in}}%
\pgfpathmoveto{\pgfqpoint{3.447365in}{2.257732in}}%
\pgfpathlineto{\pgfqpoint{3.447365in}{2.257732in}}%
\pgfpathlineto{\pgfqpoint{3.447365in}{2.260681in}}%
\pgfpathlineto{\pgfqpoint{3.451906in}{2.260681in}}%
\pgfpathlineto{\pgfqpoint{3.451906in}{2.257732in}}%
\pgfpathmoveto{\pgfqpoint{3.447365in}{2.260681in}}%
\pgfpathlineto{\pgfqpoint{3.447365in}{2.260681in}}%
\pgfpathlineto{\pgfqpoint{3.447365in}{2.263630in}}%
\pgfpathlineto{\pgfqpoint{3.451906in}{2.263630in}}%
\pgfpathlineto{\pgfqpoint{3.451906in}{2.260681in}}%
\pgfpathmoveto{\pgfqpoint{3.451906in}{2.257732in}}%
\pgfpathlineto{\pgfqpoint{3.451906in}{2.257732in}}%
\pgfpathlineto{\pgfqpoint{3.451906in}{2.260681in}}%
\pgfpathlineto{\pgfqpoint{3.456447in}{2.260681in}}%
\pgfpathlineto{\pgfqpoint{3.456447in}{2.257732in}}%
\pgfpathmoveto{\pgfqpoint{3.451906in}{2.260681in}}%
\pgfpathlineto{\pgfqpoint{3.451906in}{2.260681in}}%
\pgfpathlineto{\pgfqpoint{3.451906in}{2.263630in}}%
\pgfpathlineto{\pgfqpoint{3.456447in}{2.263630in}}%
\pgfpathlineto{\pgfqpoint{3.456447in}{2.260681in}}%
\pgfpathmoveto{\pgfqpoint{3.456447in}{2.257732in}}%
\pgfpathlineto{\pgfqpoint{3.456447in}{2.257732in}}%
\pgfpathlineto{\pgfqpoint{3.456447in}{2.260681in}}%
\pgfpathlineto{\pgfqpoint{3.460988in}{2.260681in}}%
\pgfpathlineto{\pgfqpoint{3.460988in}{2.257732in}}%
\pgfpathmoveto{\pgfqpoint{3.456447in}{2.260681in}}%
\pgfpathlineto{\pgfqpoint{3.456447in}{2.260681in}}%
\pgfpathlineto{\pgfqpoint{3.456447in}{2.263630in}}%
\pgfpathlineto{\pgfqpoint{3.460988in}{2.263630in}}%
\pgfpathlineto{\pgfqpoint{3.460988in}{2.260681in}}%
\pgfpathmoveto{\pgfqpoint{3.460988in}{2.260681in}}%
\pgfpathlineto{\pgfqpoint{3.460988in}{2.260681in}}%
\pgfpathlineto{\pgfqpoint{3.460988in}{2.263630in}}%
\pgfpathlineto{\pgfqpoint{3.465529in}{2.263630in}}%
\pgfpathlineto{\pgfqpoint{3.465529in}{2.260681in}}%
\pgfpathmoveto{\pgfqpoint{3.456447in}{2.263630in}}%
\pgfpathlineto{\pgfqpoint{3.456447in}{2.263630in}}%
\pgfpathlineto{\pgfqpoint{3.456447in}{2.266580in}}%
\pgfpathlineto{\pgfqpoint{3.460988in}{2.266580in}}%
\pgfpathlineto{\pgfqpoint{3.460988in}{2.263630in}}%
\pgfpathmoveto{\pgfqpoint{3.456447in}{2.266580in}}%
\pgfpathlineto{\pgfqpoint{3.456447in}{2.266580in}}%
\pgfpathlineto{\pgfqpoint{3.456447in}{2.269529in}}%
\pgfpathlineto{\pgfqpoint{3.460988in}{2.269529in}}%
\pgfpathlineto{\pgfqpoint{3.460988in}{2.266580in}}%
\pgfpathmoveto{\pgfqpoint{3.460988in}{2.263630in}}%
\pgfpathlineto{\pgfqpoint{3.460988in}{2.263630in}}%
\pgfpathlineto{\pgfqpoint{3.460988in}{2.266580in}}%
\pgfpathlineto{\pgfqpoint{3.465529in}{2.266580in}}%
\pgfpathlineto{\pgfqpoint{3.465529in}{2.263630in}}%
\pgfpathmoveto{\pgfqpoint{3.460988in}{2.266580in}}%
\pgfpathlineto{\pgfqpoint{3.460988in}{2.266580in}}%
\pgfpathlineto{\pgfqpoint{3.460988in}{2.269529in}}%
\pgfpathlineto{\pgfqpoint{3.465529in}{2.269529in}}%
\pgfpathlineto{\pgfqpoint{3.465529in}{2.266580in}}%
\pgfpathmoveto{\pgfqpoint{3.465529in}{2.263630in}}%
\pgfpathlineto{\pgfqpoint{3.465529in}{2.263630in}}%
\pgfpathlineto{\pgfqpoint{3.465529in}{2.266580in}}%
\pgfpathlineto{\pgfqpoint{3.470070in}{2.266580in}}%
\pgfpathlineto{\pgfqpoint{3.470070in}{2.263630in}}%
\pgfpathmoveto{\pgfqpoint{3.465529in}{2.266580in}}%
\pgfpathlineto{\pgfqpoint{3.465529in}{2.266580in}}%
\pgfpathlineto{\pgfqpoint{3.465529in}{2.269529in}}%
\pgfpathlineto{\pgfqpoint{3.470070in}{2.269529in}}%
\pgfpathlineto{\pgfqpoint{3.470070in}{2.266580in}}%
\pgfpathmoveto{\pgfqpoint{3.470070in}{2.266580in}}%
\pgfpathlineto{\pgfqpoint{3.470070in}{2.266580in}}%
\pgfpathlineto{\pgfqpoint{3.470070in}{2.269529in}}%
\pgfpathlineto{\pgfqpoint{3.474611in}{2.269529in}}%
\pgfpathlineto{\pgfqpoint{3.474611in}{2.266580in}}%
\pgfpathmoveto{\pgfqpoint{3.465529in}{2.269529in}}%
\pgfpathlineto{\pgfqpoint{3.465529in}{2.269529in}}%
\pgfpathlineto{\pgfqpoint{3.465529in}{2.272478in}}%
\pgfpathlineto{\pgfqpoint{3.470070in}{2.272478in}}%
\pgfpathlineto{\pgfqpoint{3.470070in}{2.269529in}}%
\pgfpathmoveto{\pgfqpoint{3.465529in}{2.272478in}}%
\pgfpathlineto{\pgfqpoint{3.465529in}{2.272478in}}%
\pgfpathlineto{\pgfqpoint{3.465529in}{2.275427in}}%
\pgfpathlineto{\pgfqpoint{3.470070in}{2.275427in}}%
\pgfpathlineto{\pgfqpoint{3.470070in}{2.272478in}}%
\pgfpathmoveto{\pgfqpoint{3.470070in}{2.269529in}}%
\pgfpathlineto{\pgfqpoint{3.470070in}{2.269529in}}%
\pgfpathlineto{\pgfqpoint{3.470070in}{2.272478in}}%
\pgfpathlineto{\pgfqpoint{3.474611in}{2.272478in}}%
\pgfpathlineto{\pgfqpoint{3.474611in}{2.269529in}}%
\pgfpathmoveto{\pgfqpoint{3.470070in}{2.272478in}}%
\pgfpathlineto{\pgfqpoint{3.470070in}{2.272478in}}%
\pgfpathlineto{\pgfqpoint{3.470070in}{2.275427in}}%
\pgfpathlineto{\pgfqpoint{3.474611in}{2.275427in}}%
\pgfpathlineto{\pgfqpoint{3.474611in}{2.272478in}}%
\pgfpathmoveto{\pgfqpoint{3.474611in}{2.269529in}}%
\pgfpathlineto{\pgfqpoint{3.474611in}{2.269529in}}%
\pgfpathlineto{\pgfqpoint{3.474611in}{2.272478in}}%
\pgfpathlineto{\pgfqpoint{3.479152in}{2.272478in}}%
\pgfpathlineto{\pgfqpoint{3.479152in}{2.269529in}}%
\pgfpathmoveto{\pgfqpoint{3.474611in}{2.272478in}}%
\pgfpathlineto{\pgfqpoint{3.474611in}{2.272478in}}%
\pgfpathlineto{\pgfqpoint{3.474611in}{2.275427in}}%
\pgfpathlineto{\pgfqpoint{3.479152in}{2.275427in}}%
\pgfpathlineto{\pgfqpoint{3.479152in}{2.272478in}}%
\pgfpathmoveto{\pgfqpoint{3.479152in}{2.272478in}}%
\pgfpathlineto{\pgfqpoint{3.479152in}{2.272478in}}%
\pgfpathlineto{\pgfqpoint{3.479152in}{2.275427in}}%
\pgfpathlineto{\pgfqpoint{3.483693in}{2.275427in}}%
\pgfpathlineto{\pgfqpoint{3.483693in}{2.272478in}}%
\pgfpathmoveto{\pgfqpoint{3.474611in}{2.275427in}}%
\pgfpathlineto{\pgfqpoint{3.474611in}{2.275427in}}%
\pgfpathlineto{\pgfqpoint{3.474611in}{2.278376in}}%
\pgfpathlineto{\pgfqpoint{3.479152in}{2.278376in}}%
\pgfpathlineto{\pgfqpoint{3.479152in}{2.275427in}}%
\pgfpathmoveto{\pgfqpoint{3.474611in}{2.278376in}}%
\pgfpathlineto{\pgfqpoint{3.474611in}{2.278376in}}%
\pgfpathlineto{\pgfqpoint{3.474611in}{2.281326in}}%
\pgfpathlineto{\pgfqpoint{3.479152in}{2.281326in}}%
\pgfpathlineto{\pgfqpoint{3.479152in}{2.278376in}}%
\pgfpathmoveto{\pgfqpoint{3.479152in}{2.275427in}}%
\pgfpathlineto{\pgfqpoint{3.479152in}{2.275427in}}%
\pgfpathlineto{\pgfqpoint{3.479152in}{2.278376in}}%
\pgfpathlineto{\pgfqpoint{3.483693in}{2.278376in}}%
\pgfpathlineto{\pgfqpoint{3.483693in}{2.275427in}}%
\pgfpathmoveto{\pgfqpoint{3.479152in}{2.278376in}}%
\pgfpathlineto{\pgfqpoint{3.479152in}{2.278376in}}%
\pgfpathlineto{\pgfqpoint{3.479152in}{2.281326in}}%
\pgfpathlineto{\pgfqpoint{3.483693in}{2.281326in}}%
\pgfpathlineto{\pgfqpoint{3.483693in}{2.278376in}}%
\pgfpathmoveto{\pgfqpoint{3.483693in}{2.275427in}}%
\pgfpathlineto{\pgfqpoint{3.483693in}{2.275427in}}%
\pgfpathlineto{\pgfqpoint{3.483693in}{2.278376in}}%
\pgfpathlineto{\pgfqpoint{3.488234in}{2.278376in}}%
\pgfpathlineto{\pgfqpoint{3.488234in}{2.275427in}}%
\pgfpathmoveto{\pgfqpoint{3.483693in}{2.278376in}}%
\pgfpathlineto{\pgfqpoint{3.483693in}{2.278376in}}%
\pgfpathlineto{\pgfqpoint{3.483693in}{2.281326in}}%
\pgfpathlineto{\pgfqpoint{3.488234in}{2.281326in}}%
\pgfpathlineto{\pgfqpoint{3.488234in}{2.278376in}}%
\pgfpathmoveto{\pgfqpoint{3.488234in}{2.278376in}}%
\pgfpathlineto{\pgfqpoint{3.488234in}{2.278376in}}%
\pgfpathlineto{\pgfqpoint{3.488234in}{2.281326in}}%
\pgfpathlineto{\pgfqpoint{3.492775in}{2.281326in}}%
\pgfpathlineto{\pgfqpoint{3.492775in}{2.278376in}}%
\pgfpathmoveto{\pgfqpoint{3.483693in}{2.281326in}}%
\pgfpathlineto{\pgfqpoint{3.483693in}{2.281326in}}%
\pgfpathlineto{\pgfqpoint{3.483693in}{2.284275in}}%
\pgfpathlineto{\pgfqpoint{3.488234in}{2.284275in}}%
\pgfpathlineto{\pgfqpoint{3.488234in}{2.281326in}}%
\pgfpathmoveto{\pgfqpoint{3.483693in}{2.284275in}}%
\pgfpathlineto{\pgfqpoint{3.483693in}{2.284275in}}%
\pgfpathlineto{\pgfqpoint{3.483693in}{2.287224in}}%
\pgfpathlineto{\pgfqpoint{3.488234in}{2.287224in}}%
\pgfpathlineto{\pgfqpoint{3.488234in}{2.284275in}}%
\pgfpathmoveto{\pgfqpoint{3.488234in}{2.281326in}}%
\pgfpathlineto{\pgfqpoint{3.488234in}{2.281326in}}%
\pgfpathlineto{\pgfqpoint{3.488234in}{2.284275in}}%
\pgfpathlineto{\pgfqpoint{3.492775in}{2.284275in}}%
\pgfpathlineto{\pgfqpoint{3.492775in}{2.281326in}}%
\pgfpathmoveto{\pgfqpoint{3.488234in}{2.284275in}}%
\pgfpathlineto{\pgfqpoint{3.488234in}{2.284275in}}%
\pgfpathlineto{\pgfqpoint{3.488234in}{2.287224in}}%
\pgfpathlineto{\pgfqpoint{3.492775in}{2.287224in}}%
\pgfpathlineto{\pgfqpoint{3.492775in}{2.284275in}}%
\pgfpathmoveto{\pgfqpoint{3.492775in}{2.281326in}}%
\pgfpathlineto{\pgfqpoint{3.492775in}{2.281326in}}%
\pgfpathlineto{\pgfqpoint{3.492775in}{2.284275in}}%
\pgfpathlineto{\pgfqpoint{3.497316in}{2.284275in}}%
\pgfpathlineto{\pgfqpoint{3.497316in}{2.281326in}}%
\pgfpathmoveto{\pgfqpoint{3.492775in}{2.284275in}}%
\pgfpathlineto{\pgfqpoint{3.492775in}{2.284275in}}%
\pgfpathlineto{\pgfqpoint{3.492775in}{2.287224in}}%
\pgfpathlineto{\pgfqpoint{3.497316in}{2.287224in}}%
\pgfpathlineto{\pgfqpoint{3.497316in}{2.284275in}}%
\pgfpathmoveto{\pgfqpoint{3.497316in}{2.284275in}}%
\pgfpathlineto{\pgfqpoint{3.497316in}{2.284275in}}%
\pgfpathlineto{\pgfqpoint{3.497316in}{2.287224in}}%
\pgfpathlineto{\pgfqpoint{3.501857in}{2.287224in}}%
\pgfpathlineto{\pgfqpoint{3.501857in}{2.284275in}}%
\pgfpathmoveto{\pgfqpoint{3.492775in}{2.287224in}}%
\pgfpathlineto{\pgfqpoint{3.492775in}{2.287224in}}%
\pgfpathlineto{\pgfqpoint{3.492775in}{2.290173in}}%
\pgfpathlineto{\pgfqpoint{3.497316in}{2.290173in}}%
\pgfpathlineto{\pgfqpoint{3.497316in}{2.287224in}}%
\pgfpathmoveto{\pgfqpoint{3.492775in}{2.290173in}}%
\pgfpathlineto{\pgfqpoint{3.492775in}{2.290173in}}%
\pgfpathlineto{\pgfqpoint{3.492775in}{2.293122in}}%
\pgfpathlineto{\pgfqpoint{3.497316in}{2.293122in}}%
\pgfpathlineto{\pgfqpoint{3.497316in}{2.290173in}}%
\pgfpathmoveto{\pgfqpoint{3.497316in}{2.287224in}}%
\pgfpathlineto{\pgfqpoint{3.497316in}{2.287224in}}%
\pgfpathlineto{\pgfqpoint{3.497316in}{2.290173in}}%
\pgfpathlineto{\pgfqpoint{3.501857in}{2.290173in}}%
\pgfpathlineto{\pgfqpoint{3.501857in}{2.287224in}}%
\pgfpathmoveto{\pgfqpoint{3.497316in}{2.290173in}}%
\pgfpathlineto{\pgfqpoint{3.497316in}{2.290173in}}%
\pgfpathlineto{\pgfqpoint{3.497316in}{2.293122in}}%
\pgfpathlineto{\pgfqpoint{3.501857in}{2.293122in}}%
\pgfpathlineto{\pgfqpoint{3.501857in}{2.290173in}}%
\pgfpathmoveto{\pgfqpoint{3.501857in}{2.287224in}}%
\pgfpathlineto{\pgfqpoint{3.501857in}{2.287224in}}%
\pgfpathlineto{\pgfqpoint{3.501857in}{2.290173in}}%
\pgfpathlineto{\pgfqpoint{3.506398in}{2.290173in}}%
\pgfpathlineto{\pgfqpoint{3.506398in}{2.287224in}}%
\pgfpathmoveto{\pgfqpoint{3.501857in}{2.290173in}}%
\pgfpathlineto{\pgfqpoint{3.501857in}{2.290173in}}%
\pgfpathlineto{\pgfqpoint{3.501857in}{2.293122in}}%
\pgfpathlineto{\pgfqpoint{3.506398in}{2.293122in}}%
\pgfpathlineto{\pgfqpoint{3.506398in}{2.290173in}}%
\pgfpathmoveto{\pgfqpoint{3.506398in}{2.290173in}}%
\pgfpathlineto{\pgfqpoint{3.506398in}{2.290173in}}%
\pgfpathlineto{\pgfqpoint{3.506398in}{2.293122in}}%
\pgfpathlineto{\pgfqpoint{3.510939in}{2.293122in}}%
\pgfpathlineto{\pgfqpoint{3.510939in}{2.290173in}}%
\pgfpathmoveto{\pgfqpoint{3.501857in}{2.293122in}}%
\pgfpathlineto{\pgfqpoint{3.501857in}{2.293122in}}%
\pgfpathlineto{\pgfqpoint{3.501857in}{2.296072in}}%
\pgfpathlineto{\pgfqpoint{3.506398in}{2.296072in}}%
\pgfpathlineto{\pgfqpoint{3.506398in}{2.293122in}}%
\pgfpathmoveto{\pgfqpoint{3.501857in}{2.296072in}}%
\pgfpathlineto{\pgfqpoint{3.501857in}{2.296072in}}%
\pgfpathlineto{\pgfqpoint{3.501857in}{2.299021in}}%
\pgfpathlineto{\pgfqpoint{3.506398in}{2.299021in}}%
\pgfpathlineto{\pgfqpoint{3.506398in}{2.296072in}}%
\pgfpathmoveto{\pgfqpoint{3.506398in}{2.293122in}}%
\pgfpathlineto{\pgfqpoint{3.506398in}{2.293122in}}%
\pgfpathlineto{\pgfqpoint{3.506398in}{2.296072in}}%
\pgfpathlineto{\pgfqpoint{3.510939in}{2.296072in}}%
\pgfpathlineto{\pgfqpoint{3.510939in}{2.293122in}}%
\pgfpathmoveto{\pgfqpoint{3.506398in}{2.296072in}}%
\pgfpathlineto{\pgfqpoint{3.506398in}{2.296072in}}%
\pgfpathlineto{\pgfqpoint{3.506398in}{2.299021in}}%
\pgfpathlineto{\pgfqpoint{3.510939in}{2.299021in}}%
\pgfpathlineto{\pgfqpoint{3.510939in}{2.296072in}}%
\pgfpathmoveto{\pgfqpoint{3.501857in}{2.647031in}}%
\pgfpathlineto{\pgfqpoint{3.501857in}{2.647031in}}%
\pgfpathlineto{\pgfqpoint{3.501857in}{2.649980in}}%
\pgfpathlineto{\pgfqpoint{3.506398in}{2.649980in}}%
\pgfpathlineto{\pgfqpoint{3.506398in}{2.647031in}}%
\pgfpathmoveto{\pgfqpoint{3.501857in}{2.649980in}}%
\pgfpathlineto{\pgfqpoint{3.501857in}{2.649980in}}%
\pgfpathlineto{\pgfqpoint{3.501857in}{2.652929in}}%
\pgfpathlineto{\pgfqpoint{3.506398in}{2.652929in}}%
\pgfpathlineto{\pgfqpoint{3.506398in}{2.649980in}}%
\pgfpathmoveto{\pgfqpoint{3.506398in}{2.647031in}}%
\pgfpathlineto{\pgfqpoint{3.506398in}{2.647031in}}%
\pgfpathlineto{\pgfqpoint{3.506398in}{2.649980in}}%
\pgfpathlineto{\pgfqpoint{3.510939in}{2.649980in}}%
\pgfpathlineto{\pgfqpoint{3.510939in}{2.647031in}}%
\pgfpathmoveto{\pgfqpoint{3.506398in}{2.649980in}}%
\pgfpathlineto{\pgfqpoint{3.506398in}{2.649980in}}%
\pgfpathlineto{\pgfqpoint{3.506398in}{2.652929in}}%
\pgfpathlineto{\pgfqpoint{3.510939in}{2.652929in}}%
\pgfpathlineto{\pgfqpoint{3.510939in}{2.649980in}}%
\pgfpathmoveto{\pgfqpoint{3.501857in}{2.652929in}}%
\pgfpathlineto{\pgfqpoint{3.501857in}{2.652929in}}%
\pgfpathlineto{\pgfqpoint{3.501857in}{2.655878in}}%
\pgfpathlineto{\pgfqpoint{3.506398in}{2.655878in}}%
\pgfpathlineto{\pgfqpoint{3.506398in}{2.652929in}}%
\pgfpathmoveto{\pgfqpoint{3.501857in}{2.655878in}}%
\pgfpathlineto{\pgfqpoint{3.501857in}{2.655878in}}%
\pgfpathlineto{\pgfqpoint{3.501857in}{2.658827in}}%
\pgfpathlineto{\pgfqpoint{3.506398in}{2.658827in}}%
\pgfpathlineto{\pgfqpoint{3.506398in}{2.655878in}}%
\pgfpathmoveto{\pgfqpoint{3.506398in}{2.652929in}}%
\pgfpathlineto{\pgfqpoint{3.506398in}{2.652929in}}%
\pgfpathlineto{\pgfqpoint{3.506398in}{2.655878in}}%
\pgfpathlineto{\pgfqpoint{3.510939in}{2.655878in}}%
\pgfpathlineto{\pgfqpoint{3.510939in}{2.652929in}}%
\pgfpathmoveto{\pgfqpoint{3.506398in}{2.655878in}}%
\pgfpathlineto{\pgfqpoint{3.506398in}{2.655878in}}%
\pgfpathlineto{\pgfqpoint{3.506398in}{2.658827in}}%
\pgfpathlineto{\pgfqpoint{3.510939in}{2.658827in}}%
\pgfpathlineto{\pgfqpoint{3.510939in}{2.655878in}}%
\pgfpathmoveto{\pgfqpoint{3.492775in}{2.658827in}}%
\pgfpathlineto{\pgfqpoint{3.492775in}{2.658827in}}%
\pgfpathlineto{\pgfqpoint{3.492775in}{2.661776in}}%
\pgfpathlineto{\pgfqpoint{3.497316in}{2.661776in}}%
\pgfpathlineto{\pgfqpoint{3.497316in}{2.658827in}}%
\pgfpathmoveto{\pgfqpoint{3.492775in}{2.661776in}}%
\pgfpathlineto{\pgfqpoint{3.492775in}{2.661776in}}%
\pgfpathlineto{\pgfqpoint{3.492775in}{2.664726in}}%
\pgfpathlineto{\pgfqpoint{3.497316in}{2.664726in}}%
\pgfpathlineto{\pgfqpoint{3.497316in}{2.661776in}}%
\pgfpathmoveto{\pgfqpoint{3.497316in}{2.658827in}}%
\pgfpathlineto{\pgfqpoint{3.497316in}{2.658827in}}%
\pgfpathlineto{\pgfqpoint{3.497316in}{2.661776in}}%
\pgfpathlineto{\pgfqpoint{3.501857in}{2.661776in}}%
\pgfpathlineto{\pgfqpoint{3.501857in}{2.658827in}}%
\pgfpathmoveto{\pgfqpoint{3.497316in}{2.661776in}}%
\pgfpathlineto{\pgfqpoint{3.497316in}{2.661776in}}%
\pgfpathlineto{\pgfqpoint{3.497316in}{2.664726in}}%
\pgfpathlineto{\pgfqpoint{3.501857in}{2.664726in}}%
\pgfpathlineto{\pgfqpoint{3.501857in}{2.661776in}}%
\pgfpathmoveto{\pgfqpoint{3.492775in}{2.664726in}}%
\pgfpathlineto{\pgfqpoint{3.492775in}{2.664726in}}%
\pgfpathlineto{\pgfqpoint{3.492775in}{2.667675in}}%
\pgfpathlineto{\pgfqpoint{3.497316in}{2.667675in}}%
\pgfpathlineto{\pgfqpoint{3.497316in}{2.664726in}}%
\pgfpathmoveto{\pgfqpoint{3.492775in}{2.667675in}}%
\pgfpathlineto{\pgfqpoint{3.492775in}{2.667675in}}%
\pgfpathlineto{\pgfqpoint{3.492775in}{2.670624in}}%
\pgfpathlineto{\pgfqpoint{3.497316in}{2.670624in}}%
\pgfpathlineto{\pgfqpoint{3.497316in}{2.667675in}}%
\pgfpathmoveto{\pgfqpoint{3.497316in}{2.664726in}}%
\pgfpathlineto{\pgfqpoint{3.497316in}{2.664726in}}%
\pgfpathlineto{\pgfqpoint{3.497316in}{2.667675in}}%
\pgfpathlineto{\pgfqpoint{3.501857in}{2.667675in}}%
\pgfpathlineto{\pgfqpoint{3.501857in}{2.664726in}}%
\pgfpathmoveto{\pgfqpoint{3.497316in}{2.667675in}}%
\pgfpathlineto{\pgfqpoint{3.497316in}{2.667675in}}%
\pgfpathlineto{\pgfqpoint{3.497316in}{2.670624in}}%
\pgfpathlineto{\pgfqpoint{3.501857in}{2.670624in}}%
\pgfpathlineto{\pgfqpoint{3.501857in}{2.667675in}}%
\pgfpathmoveto{\pgfqpoint{3.501857in}{2.658827in}}%
\pgfpathlineto{\pgfqpoint{3.501857in}{2.658827in}}%
\pgfpathlineto{\pgfqpoint{3.501857in}{2.661776in}}%
\pgfpathlineto{\pgfqpoint{3.506398in}{2.661776in}}%
\pgfpathlineto{\pgfqpoint{3.506398in}{2.658827in}}%
\pgfpathmoveto{\pgfqpoint{3.501857in}{2.661776in}}%
\pgfpathlineto{\pgfqpoint{3.501857in}{2.661776in}}%
\pgfpathlineto{\pgfqpoint{3.501857in}{2.664726in}}%
\pgfpathlineto{\pgfqpoint{3.506398in}{2.664726in}}%
\pgfpathlineto{\pgfqpoint{3.506398in}{2.661776in}}%
\pgfpathmoveto{\pgfqpoint{3.429201in}{2.741405in}}%
\pgfpathlineto{\pgfqpoint{3.429201in}{2.741405in}}%
\pgfpathlineto{\pgfqpoint{3.429201in}{2.744355in}}%
\pgfpathlineto{\pgfqpoint{3.433742in}{2.744355in}}%
\pgfpathlineto{\pgfqpoint{3.433742in}{2.741405in}}%
\pgfpathmoveto{\pgfqpoint{3.429201in}{2.744355in}}%
\pgfpathlineto{\pgfqpoint{3.429201in}{2.744355in}}%
\pgfpathlineto{\pgfqpoint{3.429201in}{2.747304in}}%
\pgfpathlineto{\pgfqpoint{3.433742in}{2.747304in}}%
\pgfpathlineto{\pgfqpoint{3.433742in}{2.744355in}}%
\pgfpathmoveto{\pgfqpoint{3.433742in}{2.741405in}}%
\pgfpathlineto{\pgfqpoint{3.433742in}{2.741405in}}%
\pgfpathlineto{\pgfqpoint{3.433742in}{2.744355in}}%
\pgfpathlineto{\pgfqpoint{3.438283in}{2.744355in}}%
\pgfpathlineto{\pgfqpoint{3.438283in}{2.741405in}}%
\pgfpathmoveto{\pgfqpoint{3.433742in}{2.744355in}}%
\pgfpathlineto{\pgfqpoint{3.433742in}{2.744355in}}%
\pgfpathlineto{\pgfqpoint{3.433742in}{2.747304in}}%
\pgfpathlineto{\pgfqpoint{3.438283in}{2.747304in}}%
\pgfpathlineto{\pgfqpoint{3.438283in}{2.744355in}}%
\pgfpathmoveto{\pgfqpoint{3.429201in}{2.747304in}}%
\pgfpathlineto{\pgfqpoint{3.429201in}{2.747304in}}%
\pgfpathlineto{\pgfqpoint{3.429201in}{2.750253in}}%
\pgfpathlineto{\pgfqpoint{3.433742in}{2.750253in}}%
\pgfpathlineto{\pgfqpoint{3.433742in}{2.747304in}}%
\pgfpathmoveto{\pgfqpoint{3.429201in}{2.750253in}}%
\pgfpathlineto{\pgfqpoint{3.429201in}{2.750253in}}%
\pgfpathlineto{\pgfqpoint{3.429201in}{2.753202in}}%
\pgfpathlineto{\pgfqpoint{3.433742in}{2.753202in}}%
\pgfpathlineto{\pgfqpoint{3.433742in}{2.750253in}}%
\pgfpathmoveto{\pgfqpoint{3.433742in}{2.747304in}}%
\pgfpathlineto{\pgfqpoint{3.433742in}{2.747304in}}%
\pgfpathlineto{\pgfqpoint{3.433742in}{2.750253in}}%
\pgfpathlineto{\pgfqpoint{3.438283in}{2.750253in}}%
\pgfpathlineto{\pgfqpoint{3.438283in}{2.747304in}}%
\pgfpathmoveto{\pgfqpoint{3.433742in}{2.750253in}}%
\pgfpathlineto{\pgfqpoint{3.433742in}{2.750253in}}%
\pgfpathlineto{\pgfqpoint{3.433742in}{2.753202in}}%
\pgfpathlineto{\pgfqpoint{3.438283in}{2.753202in}}%
\pgfpathlineto{\pgfqpoint{3.438283in}{2.750253in}}%
\pgfpathmoveto{\pgfqpoint{3.420119in}{2.753202in}}%
\pgfpathlineto{\pgfqpoint{3.420119in}{2.753202in}}%
\pgfpathlineto{\pgfqpoint{3.420119in}{2.756152in}}%
\pgfpathlineto{\pgfqpoint{3.424660in}{2.756152in}}%
\pgfpathlineto{\pgfqpoint{3.424660in}{2.753202in}}%
\pgfpathmoveto{\pgfqpoint{3.420119in}{2.756152in}}%
\pgfpathlineto{\pgfqpoint{3.420119in}{2.756152in}}%
\pgfpathlineto{\pgfqpoint{3.420119in}{2.759101in}}%
\pgfpathlineto{\pgfqpoint{3.424660in}{2.759101in}}%
\pgfpathlineto{\pgfqpoint{3.424660in}{2.756152in}}%
\pgfpathmoveto{\pgfqpoint{3.424660in}{2.753202in}}%
\pgfpathlineto{\pgfqpoint{3.424660in}{2.753202in}}%
\pgfpathlineto{\pgfqpoint{3.424660in}{2.756152in}}%
\pgfpathlineto{\pgfqpoint{3.429201in}{2.756152in}}%
\pgfpathlineto{\pgfqpoint{3.429201in}{2.753202in}}%
\pgfpathmoveto{\pgfqpoint{3.424660in}{2.756152in}}%
\pgfpathlineto{\pgfqpoint{3.424660in}{2.756152in}}%
\pgfpathlineto{\pgfqpoint{3.424660in}{2.759101in}}%
\pgfpathlineto{\pgfqpoint{3.429201in}{2.759101in}}%
\pgfpathlineto{\pgfqpoint{3.429201in}{2.756152in}}%
\pgfpathmoveto{\pgfqpoint{3.420119in}{2.759101in}}%
\pgfpathlineto{\pgfqpoint{3.420119in}{2.759101in}}%
\pgfpathlineto{\pgfqpoint{3.420119in}{2.762050in}}%
\pgfpathlineto{\pgfqpoint{3.424660in}{2.762050in}}%
\pgfpathlineto{\pgfqpoint{3.424660in}{2.759101in}}%
\pgfpathmoveto{\pgfqpoint{3.420119in}{2.762050in}}%
\pgfpathlineto{\pgfqpoint{3.420119in}{2.762050in}}%
\pgfpathlineto{\pgfqpoint{3.420119in}{2.764999in}}%
\pgfpathlineto{\pgfqpoint{3.424660in}{2.764999in}}%
\pgfpathlineto{\pgfqpoint{3.424660in}{2.762050in}}%
\pgfpathmoveto{\pgfqpoint{3.424660in}{2.759101in}}%
\pgfpathlineto{\pgfqpoint{3.424660in}{2.759101in}}%
\pgfpathlineto{\pgfqpoint{3.424660in}{2.762050in}}%
\pgfpathlineto{\pgfqpoint{3.429201in}{2.762050in}}%
\pgfpathlineto{\pgfqpoint{3.429201in}{2.759101in}}%
\pgfpathmoveto{\pgfqpoint{3.424660in}{2.762050in}}%
\pgfpathlineto{\pgfqpoint{3.424660in}{2.762050in}}%
\pgfpathlineto{\pgfqpoint{3.424660in}{2.764999in}}%
\pgfpathlineto{\pgfqpoint{3.429201in}{2.764999in}}%
\pgfpathlineto{\pgfqpoint{3.429201in}{2.762050in}}%
\pgfpathmoveto{\pgfqpoint{3.429201in}{2.753202in}}%
\pgfpathlineto{\pgfqpoint{3.429201in}{2.753202in}}%
\pgfpathlineto{\pgfqpoint{3.429201in}{2.756152in}}%
\pgfpathlineto{\pgfqpoint{3.433742in}{2.756152in}}%
\pgfpathlineto{\pgfqpoint{3.433742in}{2.753202in}}%
\pgfpathmoveto{\pgfqpoint{3.429201in}{2.756152in}}%
\pgfpathlineto{\pgfqpoint{3.429201in}{2.756152in}}%
\pgfpathlineto{\pgfqpoint{3.429201in}{2.759101in}}%
\pgfpathlineto{\pgfqpoint{3.433742in}{2.759101in}}%
\pgfpathlineto{\pgfqpoint{3.433742in}{2.756152in}}%
\pgfpathmoveto{\pgfqpoint{3.465529in}{2.694218in}}%
\pgfpathlineto{\pgfqpoint{3.465529in}{2.694218in}}%
\pgfpathlineto{\pgfqpoint{3.465529in}{2.697167in}}%
\pgfpathlineto{\pgfqpoint{3.470070in}{2.697167in}}%
\pgfpathlineto{\pgfqpoint{3.470070in}{2.694218in}}%
\pgfpathmoveto{\pgfqpoint{3.465529in}{2.697167in}}%
\pgfpathlineto{\pgfqpoint{3.465529in}{2.697167in}}%
\pgfpathlineto{\pgfqpoint{3.465529in}{2.700116in}}%
\pgfpathlineto{\pgfqpoint{3.470070in}{2.700116in}}%
\pgfpathlineto{\pgfqpoint{3.470070in}{2.697167in}}%
\pgfpathmoveto{\pgfqpoint{3.470070in}{2.694218in}}%
\pgfpathlineto{\pgfqpoint{3.470070in}{2.694218in}}%
\pgfpathlineto{\pgfqpoint{3.470070in}{2.697167in}}%
\pgfpathlineto{\pgfqpoint{3.474611in}{2.697167in}}%
\pgfpathlineto{\pgfqpoint{3.474611in}{2.694218in}}%
\pgfpathmoveto{\pgfqpoint{3.470070in}{2.697167in}}%
\pgfpathlineto{\pgfqpoint{3.470070in}{2.697167in}}%
\pgfpathlineto{\pgfqpoint{3.470070in}{2.700116in}}%
\pgfpathlineto{\pgfqpoint{3.474611in}{2.700116in}}%
\pgfpathlineto{\pgfqpoint{3.474611in}{2.697167in}}%
\pgfpathmoveto{\pgfqpoint{3.465529in}{2.700116in}}%
\pgfpathlineto{\pgfqpoint{3.465529in}{2.700116in}}%
\pgfpathlineto{\pgfqpoint{3.465529in}{2.703065in}}%
\pgfpathlineto{\pgfqpoint{3.470070in}{2.703065in}}%
\pgfpathlineto{\pgfqpoint{3.470070in}{2.700116in}}%
\pgfpathmoveto{\pgfqpoint{3.465529in}{2.703065in}}%
\pgfpathlineto{\pgfqpoint{3.465529in}{2.703065in}}%
\pgfpathlineto{\pgfqpoint{3.465529in}{2.706015in}}%
\pgfpathlineto{\pgfqpoint{3.470070in}{2.706015in}}%
\pgfpathlineto{\pgfqpoint{3.470070in}{2.703065in}}%
\pgfpathmoveto{\pgfqpoint{3.470070in}{2.700116in}}%
\pgfpathlineto{\pgfqpoint{3.470070in}{2.700116in}}%
\pgfpathlineto{\pgfqpoint{3.470070in}{2.703065in}}%
\pgfpathlineto{\pgfqpoint{3.474611in}{2.703065in}}%
\pgfpathlineto{\pgfqpoint{3.474611in}{2.700116in}}%
\pgfpathmoveto{\pgfqpoint{3.470070in}{2.703065in}}%
\pgfpathlineto{\pgfqpoint{3.470070in}{2.703065in}}%
\pgfpathlineto{\pgfqpoint{3.470070in}{2.706015in}}%
\pgfpathlineto{\pgfqpoint{3.474611in}{2.706015in}}%
\pgfpathlineto{\pgfqpoint{3.474611in}{2.703065in}}%
\pgfpathmoveto{\pgfqpoint{3.456447in}{2.706015in}}%
\pgfpathlineto{\pgfqpoint{3.456447in}{2.706015in}}%
\pgfpathlineto{\pgfqpoint{3.456447in}{2.708964in}}%
\pgfpathlineto{\pgfqpoint{3.460988in}{2.708964in}}%
\pgfpathlineto{\pgfqpoint{3.460988in}{2.706015in}}%
\pgfpathmoveto{\pgfqpoint{3.456447in}{2.708964in}}%
\pgfpathlineto{\pgfqpoint{3.456447in}{2.708964in}}%
\pgfpathlineto{\pgfqpoint{3.456447in}{2.711913in}}%
\pgfpathlineto{\pgfqpoint{3.460988in}{2.711913in}}%
\pgfpathlineto{\pgfqpoint{3.460988in}{2.708964in}}%
\pgfpathmoveto{\pgfqpoint{3.460988in}{2.706015in}}%
\pgfpathlineto{\pgfqpoint{3.460988in}{2.706015in}}%
\pgfpathlineto{\pgfqpoint{3.460988in}{2.708964in}}%
\pgfpathlineto{\pgfqpoint{3.465529in}{2.708964in}}%
\pgfpathlineto{\pgfqpoint{3.465529in}{2.706015in}}%
\pgfpathmoveto{\pgfqpoint{3.460988in}{2.708964in}}%
\pgfpathlineto{\pgfqpoint{3.460988in}{2.708964in}}%
\pgfpathlineto{\pgfqpoint{3.460988in}{2.711913in}}%
\pgfpathlineto{\pgfqpoint{3.465529in}{2.711913in}}%
\pgfpathlineto{\pgfqpoint{3.465529in}{2.708964in}}%
\pgfpathmoveto{\pgfqpoint{3.456447in}{2.711913in}}%
\pgfpathlineto{\pgfqpoint{3.456447in}{2.711913in}}%
\pgfpathlineto{\pgfqpoint{3.456447in}{2.714862in}}%
\pgfpathlineto{\pgfqpoint{3.460988in}{2.714862in}}%
\pgfpathlineto{\pgfqpoint{3.460988in}{2.711913in}}%
\pgfpathmoveto{\pgfqpoint{3.456447in}{2.714862in}}%
\pgfpathlineto{\pgfqpoint{3.456447in}{2.714862in}}%
\pgfpathlineto{\pgfqpoint{3.456447in}{2.717812in}}%
\pgfpathlineto{\pgfqpoint{3.460988in}{2.717812in}}%
\pgfpathlineto{\pgfqpoint{3.460988in}{2.714862in}}%
\pgfpathmoveto{\pgfqpoint{3.460988in}{2.711913in}}%
\pgfpathlineto{\pgfqpoint{3.460988in}{2.711913in}}%
\pgfpathlineto{\pgfqpoint{3.460988in}{2.714862in}}%
\pgfpathlineto{\pgfqpoint{3.465529in}{2.714862in}}%
\pgfpathlineto{\pgfqpoint{3.465529in}{2.711913in}}%
\pgfpathmoveto{\pgfqpoint{3.460988in}{2.714862in}}%
\pgfpathlineto{\pgfqpoint{3.460988in}{2.714862in}}%
\pgfpathlineto{\pgfqpoint{3.460988in}{2.717812in}}%
\pgfpathlineto{\pgfqpoint{3.465529in}{2.717812in}}%
\pgfpathlineto{\pgfqpoint{3.465529in}{2.714862in}}%
\pgfpathmoveto{\pgfqpoint{3.465529in}{2.706015in}}%
\pgfpathlineto{\pgfqpoint{3.465529in}{2.706015in}}%
\pgfpathlineto{\pgfqpoint{3.465529in}{2.708964in}}%
\pgfpathlineto{\pgfqpoint{3.470070in}{2.708964in}}%
\pgfpathlineto{\pgfqpoint{3.470070in}{2.706015in}}%
\pgfpathmoveto{\pgfqpoint{3.465529in}{2.708964in}}%
\pgfpathlineto{\pgfqpoint{3.465529in}{2.708964in}}%
\pgfpathlineto{\pgfqpoint{3.465529in}{2.711913in}}%
\pgfpathlineto{\pgfqpoint{3.470070in}{2.711913in}}%
\pgfpathlineto{\pgfqpoint{3.470070in}{2.708964in}}%
\pgfpathmoveto{\pgfqpoint{3.483693in}{2.670624in}}%
\pgfpathlineto{\pgfqpoint{3.483693in}{2.670624in}}%
\pgfpathlineto{\pgfqpoint{3.483693in}{2.673573in}}%
\pgfpathlineto{\pgfqpoint{3.488234in}{2.673573in}}%
\pgfpathlineto{\pgfqpoint{3.488234in}{2.670624in}}%
\pgfpathmoveto{\pgfqpoint{3.483693in}{2.673573in}}%
\pgfpathlineto{\pgfqpoint{3.483693in}{2.673573in}}%
\pgfpathlineto{\pgfqpoint{3.483693in}{2.676522in}}%
\pgfpathlineto{\pgfqpoint{3.488234in}{2.676522in}}%
\pgfpathlineto{\pgfqpoint{3.488234in}{2.673573in}}%
\pgfpathmoveto{\pgfqpoint{3.488234in}{2.670624in}}%
\pgfpathlineto{\pgfqpoint{3.488234in}{2.670624in}}%
\pgfpathlineto{\pgfqpoint{3.488234in}{2.673573in}}%
\pgfpathlineto{\pgfqpoint{3.492775in}{2.673573in}}%
\pgfpathlineto{\pgfqpoint{3.492775in}{2.670624in}}%
\pgfpathmoveto{\pgfqpoint{3.488234in}{2.673573in}}%
\pgfpathlineto{\pgfqpoint{3.488234in}{2.673573in}}%
\pgfpathlineto{\pgfqpoint{3.488234in}{2.676522in}}%
\pgfpathlineto{\pgfqpoint{3.492775in}{2.676522in}}%
\pgfpathlineto{\pgfqpoint{3.492775in}{2.673573in}}%
\pgfpathmoveto{\pgfqpoint{3.483693in}{2.676522in}}%
\pgfpathlineto{\pgfqpoint{3.483693in}{2.676522in}}%
\pgfpathlineto{\pgfqpoint{3.483693in}{2.679472in}}%
\pgfpathlineto{\pgfqpoint{3.488234in}{2.679472in}}%
\pgfpathlineto{\pgfqpoint{3.488234in}{2.676522in}}%
\pgfpathmoveto{\pgfqpoint{3.483693in}{2.679472in}}%
\pgfpathlineto{\pgfqpoint{3.483693in}{2.679472in}}%
\pgfpathlineto{\pgfqpoint{3.483693in}{2.682421in}}%
\pgfpathlineto{\pgfqpoint{3.488234in}{2.682421in}}%
\pgfpathlineto{\pgfqpoint{3.488234in}{2.679472in}}%
\pgfpathmoveto{\pgfqpoint{3.488234in}{2.676522in}}%
\pgfpathlineto{\pgfqpoint{3.488234in}{2.676522in}}%
\pgfpathlineto{\pgfqpoint{3.488234in}{2.679472in}}%
\pgfpathlineto{\pgfqpoint{3.492775in}{2.679472in}}%
\pgfpathlineto{\pgfqpoint{3.492775in}{2.676522in}}%
\pgfpathmoveto{\pgfqpoint{3.488234in}{2.679472in}}%
\pgfpathlineto{\pgfqpoint{3.488234in}{2.679472in}}%
\pgfpathlineto{\pgfqpoint{3.488234in}{2.682421in}}%
\pgfpathlineto{\pgfqpoint{3.492775in}{2.682421in}}%
\pgfpathlineto{\pgfqpoint{3.492775in}{2.679472in}}%
\pgfpathmoveto{\pgfqpoint{3.474611in}{2.682421in}}%
\pgfpathlineto{\pgfqpoint{3.474611in}{2.682421in}}%
\pgfpathlineto{\pgfqpoint{3.474611in}{2.685370in}}%
\pgfpathlineto{\pgfqpoint{3.479152in}{2.685370in}}%
\pgfpathlineto{\pgfqpoint{3.479152in}{2.682421in}}%
\pgfpathmoveto{\pgfqpoint{3.474611in}{2.685370in}}%
\pgfpathlineto{\pgfqpoint{3.474611in}{2.685370in}}%
\pgfpathlineto{\pgfqpoint{3.474611in}{2.688319in}}%
\pgfpathlineto{\pgfqpoint{3.479152in}{2.688319in}}%
\pgfpathlineto{\pgfqpoint{3.479152in}{2.685370in}}%
\pgfpathmoveto{\pgfqpoint{3.479152in}{2.682421in}}%
\pgfpathlineto{\pgfqpoint{3.479152in}{2.682421in}}%
\pgfpathlineto{\pgfqpoint{3.479152in}{2.685370in}}%
\pgfpathlineto{\pgfqpoint{3.483693in}{2.685370in}}%
\pgfpathlineto{\pgfqpoint{3.483693in}{2.682421in}}%
\pgfpathmoveto{\pgfqpoint{3.479152in}{2.685370in}}%
\pgfpathlineto{\pgfqpoint{3.479152in}{2.685370in}}%
\pgfpathlineto{\pgfqpoint{3.479152in}{2.688319in}}%
\pgfpathlineto{\pgfqpoint{3.483693in}{2.688319in}}%
\pgfpathlineto{\pgfqpoint{3.483693in}{2.685370in}}%
\pgfpathmoveto{\pgfqpoint{3.474611in}{2.688319in}}%
\pgfpathlineto{\pgfqpoint{3.474611in}{2.688319in}}%
\pgfpathlineto{\pgfqpoint{3.474611in}{2.691268in}}%
\pgfpathlineto{\pgfqpoint{3.479152in}{2.691268in}}%
\pgfpathlineto{\pgfqpoint{3.479152in}{2.688319in}}%
\pgfpathmoveto{\pgfqpoint{3.474611in}{2.691268in}}%
\pgfpathlineto{\pgfqpoint{3.474611in}{2.691268in}}%
\pgfpathlineto{\pgfqpoint{3.474611in}{2.694218in}}%
\pgfpathlineto{\pgfqpoint{3.479152in}{2.694218in}}%
\pgfpathlineto{\pgfqpoint{3.479152in}{2.691268in}}%
\pgfpathmoveto{\pgfqpoint{3.479152in}{2.688319in}}%
\pgfpathlineto{\pgfqpoint{3.479152in}{2.688319in}}%
\pgfpathlineto{\pgfqpoint{3.479152in}{2.691268in}}%
\pgfpathlineto{\pgfqpoint{3.483693in}{2.691268in}}%
\pgfpathlineto{\pgfqpoint{3.483693in}{2.688319in}}%
\pgfpathmoveto{\pgfqpoint{3.479152in}{2.691268in}}%
\pgfpathlineto{\pgfqpoint{3.479152in}{2.691268in}}%
\pgfpathlineto{\pgfqpoint{3.479152in}{2.694218in}}%
\pgfpathlineto{\pgfqpoint{3.483693in}{2.694218in}}%
\pgfpathlineto{\pgfqpoint{3.483693in}{2.691268in}}%
\pgfpathmoveto{\pgfqpoint{3.483693in}{2.682421in}}%
\pgfpathlineto{\pgfqpoint{3.483693in}{2.682421in}}%
\pgfpathlineto{\pgfqpoint{3.483693in}{2.685370in}}%
\pgfpathlineto{\pgfqpoint{3.488234in}{2.685370in}}%
\pgfpathlineto{\pgfqpoint{3.488234in}{2.682421in}}%
\pgfpathmoveto{\pgfqpoint{3.483693in}{2.685370in}}%
\pgfpathlineto{\pgfqpoint{3.483693in}{2.685370in}}%
\pgfpathlineto{\pgfqpoint{3.483693in}{2.688319in}}%
\pgfpathlineto{\pgfqpoint{3.488234in}{2.688319in}}%
\pgfpathlineto{\pgfqpoint{3.488234in}{2.685370in}}%
\pgfpathmoveto{\pgfqpoint{3.492775in}{2.670624in}}%
\pgfpathlineto{\pgfqpoint{3.492775in}{2.670624in}}%
\pgfpathlineto{\pgfqpoint{3.492775in}{2.673573in}}%
\pgfpathlineto{\pgfqpoint{3.497316in}{2.673573in}}%
\pgfpathlineto{\pgfqpoint{3.497316in}{2.670624in}}%
\pgfpathmoveto{\pgfqpoint{3.492775in}{2.673573in}}%
\pgfpathlineto{\pgfqpoint{3.492775in}{2.673573in}}%
\pgfpathlineto{\pgfqpoint{3.492775in}{2.676522in}}%
\pgfpathlineto{\pgfqpoint{3.497316in}{2.676522in}}%
\pgfpathlineto{\pgfqpoint{3.497316in}{2.673573in}}%
\pgfpathmoveto{\pgfqpoint{3.474611in}{2.694218in}}%
\pgfpathlineto{\pgfqpoint{3.474611in}{2.694218in}}%
\pgfpathlineto{\pgfqpoint{3.474611in}{2.697167in}}%
\pgfpathlineto{\pgfqpoint{3.479152in}{2.697167in}}%
\pgfpathlineto{\pgfqpoint{3.479152in}{2.694218in}}%
\pgfpathmoveto{\pgfqpoint{3.474611in}{2.697167in}}%
\pgfpathlineto{\pgfqpoint{3.474611in}{2.697167in}}%
\pgfpathlineto{\pgfqpoint{3.474611in}{2.700116in}}%
\pgfpathlineto{\pgfqpoint{3.479152in}{2.700116in}}%
\pgfpathlineto{\pgfqpoint{3.479152in}{2.697167in}}%
\pgfpathmoveto{\pgfqpoint{3.447365in}{2.717812in}}%
\pgfpathlineto{\pgfqpoint{3.447365in}{2.717812in}}%
\pgfpathlineto{\pgfqpoint{3.447365in}{2.720761in}}%
\pgfpathlineto{\pgfqpoint{3.451906in}{2.720761in}}%
\pgfpathlineto{\pgfqpoint{3.451906in}{2.717812in}}%
\pgfpathmoveto{\pgfqpoint{3.447365in}{2.720761in}}%
\pgfpathlineto{\pgfqpoint{3.447365in}{2.720761in}}%
\pgfpathlineto{\pgfqpoint{3.447365in}{2.723710in}}%
\pgfpathlineto{\pgfqpoint{3.451906in}{2.723710in}}%
\pgfpathlineto{\pgfqpoint{3.451906in}{2.720761in}}%
\pgfpathmoveto{\pgfqpoint{3.451906in}{2.717812in}}%
\pgfpathlineto{\pgfqpoint{3.451906in}{2.717812in}}%
\pgfpathlineto{\pgfqpoint{3.451906in}{2.720761in}}%
\pgfpathlineto{\pgfqpoint{3.456447in}{2.720761in}}%
\pgfpathlineto{\pgfqpoint{3.456447in}{2.717812in}}%
\pgfpathmoveto{\pgfqpoint{3.451906in}{2.720761in}}%
\pgfpathlineto{\pgfqpoint{3.451906in}{2.720761in}}%
\pgfpathlineto{\pgfqpoint{3.451906in}{2.723710in}}%
\pgfpathlineto{\pgfqpoint{3.456447in}{2.723710in}}%
\pgfpathlineto{\pgfqpoint{3.456447in}{2.720761in}}%
\pgfpathmoveto{\pgfqpoint{3.447365in}{2.723710in}}%
\pgfpathlineto{\pgfqpoint{3.447365in}{2.723710in}}%
\pgfpathlineto{\pgfqpoint{3.447365in}{2.726659in}}%
\pgfpathlineto{\pgfqpoint{3.451906in}{2.726659in}}%
\pgfpathlineto{\pgfqpoint{3.451906in}{2.723710in}}%
\pgfpathmoveto{\pgfqpoint{3.447365in}{2.726659in}}%
\pgfpathlineto{\pgfqpoint{3.447365in}{2.726659in}}%
\pgfpathlineto{\pgfqpoint{3.447365in}{2.729608in}}%
\pgfpathlineto{\pgfqpoint{3.451906in}{2.729608in}}%
\pgfpathlineto{\pgfqpoint{3.451906in}{2.726659in}}%
\pgfpathmoveto{\pgfqpoint{3.451906in}{2.723710in}}%
\pgfpathlineto{\pgfqpoint{3.451906in}{2.723710in}}%
\pgfpathlineto{\pgfqpoint{3.451906in}{2.726659in}}%
\pgfpathlineto{\pgfqpoint{3.456447in}{2.726659in}}%
\pgfpathlineto{\pgfqpoint{3.456447in}{2.723710in}}%
\pgfpathmoveto{\pgfqpoint{3.451906in}{2.726659in}}%
\pgfpathlineto{\pgfqpoint{3.451906in}{2.726659in}}%
\pgfpathlineto{\pgfqpoint{3.451906in}{2.729608in}}%
\pgfpathlineto{\pgfqpoint{3.456447in}{2.729608in}}%
\pgfpathlineto{\pgfqpoint{3.456447in}{2.726659in}}%
\pgfpathmoveto{\pgfqpoint{3.438283in}{2.729608in}}%
\pgfpathlineto{\pgfqpoint{3.438283in}{2.729608in}}%
\pgfpathlineto{\pgfqpoint{3.438283in}{2.732558in}}%
\pgfpathlineto{\pgfqpoint{3.442824in}{2.732558in}}%
\pgfpathlineto{\pgfqpoint{3.442824in}{2.729608in}}%
\pgfpathmoveto{\pgfqpoint{3.438283in}{2.732558in}}%
\pgfpathlineto{\pgfqpoint{3.438283in}{2.732558in}}%
\pgfpathlineto{\pgfqpoint{3.438283in}{2.735507in}}%
\pgfpathlineto{\pgfqpoint{3.442824in}{2.735507in}}%
\pgfpathlineto{\pgfqpoint{3.442824in}{2.732558in}}%
\pgfpathmoveto{\pgfqpoint{3.442824in}{2.729608in}}%
\pgfpathlineto{\pgfqpoint{3.442824in}{2.729608in}}%
\pgfpathlineto{\pgfqpoint{3.442824in}{2.732558in}}%
\pgfpathlineto{\pgfqpoint{3.447365in}{2.732558in}}%
\pgfpathlineto{\pgfqpoint{3.447365in}{2.729608in}}%
\pgfpathmoveto{\pgfqpoint{3.442824in}{2.732558in}}%
\pgfpathlineto{\pgfqpoint{3.442824in}{2.732558in}}%
\pgfpathlineto{\pgfqpoint{3.442824in}{2.735507in}}%
\pgfpathlineto{\pgfqpoint{3.447365in}{2.735507in}}%
\pgfpathlineto{\pgfqpoint{3.447365in}{2.732558in}}%
\pgfpathmoveto{\pgfqpoint{3.438283in}{2.735507in}}%
\pgfpathlineto{\pgfqpoint{3.438283in}{2.735507in}}%
\pgfpathlineto{\pgfqpoint{3.438283in}{2.738456in}}%
\pgfpathlineto{\pgfqpoint{3.442824in}{2.738456in}}%
\pgfpathlineto{\pgfqpoint{3.442824in}{2.735507in}}%
\pgfpathmoveto{\pgfqpoint{3.438283in}{2.738456in}}%
\pgfpathlineto{\pgfqpoint{3.438283in}{2.738456in}}%
\pgfpathlineto{\pgfqpoint{3.438283in}{2.741405in}}%
\pgfpathlineto{\pgfqpoint{3.442824in}{2.741405in}}%
\pgfpathlineto{\pgfqpoint{3.442824in}{2.738456in}}%
\pgfpathmoveto{\pgfqpoint{3.442824in}{2.735507in}}%
\pgfpathlineto{\pgfqpoint{3.442824in}{2.735507in}}%
\pgfpathlineto{\pgfqpoint{3.442824in}{2.738456in}}%
\pgfpathlineto{\pgfqpoint{3.447365in}{2.738456in}}%
\pgfpathlineto{\pgfqpoint{3.447365in}{2.735507in}}%
\pgfpathmoveto{\pgfqpoint{3.442824in}{2.738456in}}%
\pgfpathlineto{\pgfqpoint{3.442824in}{2.738456in}}%
\pgfpathlineto{\pgfqpoint{3.442824in}{2.741405in}}%
\pgfpathlineto{\pgfqpoint{3.447365in}{2.741405in}}%
\pgfpathlineto{\pgfqpoint{3.447365in}{2.738456in}}%
\pgfpathmoveto{\pgfqpoint{3.447365in}{2.729608in}}%
\pgfpathlineto{\pgfqpoint{3.447365in}{2.729608in}}%
\pgfpathlineto{\pgfqpoint{3.447365in}{2.732558in}}%
\pgfpathlineto{\pgfqpoint{3.451906in}{2.732558in}}%
\pgfpathlineto{\pgfqpoint{3.451906in}{2.729608in}}%
\pgfpathmoveto{\pgfqpoint{3.447365in}{2.732558in}}%
\pgfpathlineto{\pgfqpoint{3.447365in}{2.732558in}}%
\pgfpathlineto{\pgfqpoint{3.447365in}{2.735507in}}%
\pgfpathlineto{\pgfqpoint{3.451906in}{2.735507in}}%
\pgfpathlineto{\pgfqpoint{3.451906in}{2.732558in}}%
\pgfpathmoveto{\pgfqpoint{3.456447in}{2.717812in}}%
\pgfpathlineto{\pgfqpoint{3.456447in}{2.717812in}}%
\pgfpathlineto{\pgfqpoint{3.456447in}{2.720761in}}%
\pgfpathlineto{\pgfqpoint{3.460988in}{2.720761in}}%
\pgfpathlineto{\pgfqpoint{3.460988in}{2.717812in}}%
\pgfpathmoveto{\pgfqpoint{3.456447in}{2.720761in}}%
\pgfpathlineto{\pgfqpoint{3.456447in}{2.720761in}}%
\pgfpathlineto{\pgfqpoint{3.456447in}{2.723710in}}%
\pgfpathlineto{\pgfqpoint{3.460988in}{2.723710in}}%
\pgfpathlineto{\pgfqpoint{3.460988in}{2.720761in}}%
\pgfpathmoveto{\pgfqpoint{3.438283in}{2.741405in}}%
\pgfpathlineto{\pgfqpoint{3.438283in}{2.741405in}}%
\pgfpathlineto{\pgfqpoint{3.438283in}{2.744355in}}%
\pgfpathlineto{\pgfqpoint{3.442824in}{2.744355in}}%
\pgfpathlineto{\pgfqpoint{3.442824in}{2.741405in}}%
\pgfpathmoveto{\pgfqpoint{3.438283in}{2.744355in}}%
\pgfpathlineto{\pgfqpoint{3.438283in}{2.744355in}}%
\pgfpathlineto{\pgfqpoint{3.438283in}{2.747304in}}%
\pgfpathlineto{\pgfqpoint{3.442824in}{2.747304in}}%
\pgfpathlineto{\pgfqpoint{3.442824in}{2.744355in}}%
\pgfpathmoveto{\pgfqpoint{3.392873in}{2.788593in}}%
\pgfpathlineto{\pgfqpoint{3.392873in}{2.788593in}}%
\pgfpathlineto{\pgfqpoint{3.392873in}{2.791542in}}%
\pgfpathlineto{\pgfqpoint{3.397414in}{2.791542in}}%
\pgfpathlineto{\pgfqpoint{3.397414in}{2.788593in}}%
\pgfpathmoveto{\pgfqpoint{3.392873in}{2.791542in}}%
\pgfpathlineto{\pgfqpoint{3.392873in}{2.791542in}}%
\pgfpathlineto{\pgfqpoint{3.392873in}{2.794491in}}%
\pgfpathlineto{\pgfqpoint{3.397414in}{2.794491in}}%
\pgfpathlineto{\pgfqpoint{3.397414in}{2.791542in}}%
\pgfpathmoveto{\pgfqpoint{3.397414in}{2.788593in}}%
\pgfpathlineto{\pgfqpoint{3.397414in}{2.788593in}}%
\pgfpathlineto{\pgfqpoint{3.397414in}{2.791542in}}%
\pgfpathlineto{\pgfqpoint{3.401955in}{2.791542in}}%
\pgfpathlineto{\pgfqpoint{3.401955in}{2.788593in}}%
\pgfpathmoveto{\pgfqpoint{3.397414in}{2.791542in}}%
\pgfpathlineto{\pgfqpoint{3.397414in}{2.791542in}}%
\pgfpathlineto{\pgfqpoint{3.397414in}{2.794491in}}%
\pgfpathlineto{\pgfqpoint{3.401955in}{2.794491in}}%
\pgfpathlineto{\pgfqpoint{3.401955in}{2.791542in}}%
\pgfpathmoveto{\pgfqpoint{3.392873in}{2.794491in}}%
\pgfpathlineto{\pgfqpoint{3.392873in}{2.794491in}}%
\pgfpathlineto{\pgfqpoint{3.392873in}{2.797441in}}%
\pgfpathlineto{\pgfqpoint{3.397414in}{2.797441in}}%
\pgfpathlineto{\pgfqpoint{3.397414in}{2.794491in}}%
\pgfpathmoveto{\pgfqpoint{3.392873in}{2.797441in}}%
\pgfpathlineto{\pgfqpoint{3.392873in}{2.797441in}}%
\pgfpathlineto{\pgfqpoint{3.392873in}{2.800390in}}%
\pgfpathlineto{\pgfqpoint{3.397414in}{2.800390in}}%
\pgfpathlineto{\pgfqpoint{3.397414in}{2.797441in}}%
\pgfpathmoveto{\pgfqpoint{3.397414in}{2.794491in}}%
\pgfpathlineto{\pgfqpoint{3.397414in}{2.794491in}}%
\pgfpathlineto{\pgfqpoint{3.397414in}{2.797441in}}%
\pgfpathlineto{\pgfqpoint{3.401955in}{2.797441in}}%
\pgfpathlineto{\pgfqpoint{3.401955in}{2.794491in}}%
\pgfpathmoveto{\pgfqpoint{3.397414in}{2.797441in}}%
\pgfpathlineto{\pgfqpoint{3.397414in}{2.797441in}}%
\pgfpathlineto{\pgfqpoint{3.397414in}{2.800390in}}%
\pgfpathlineto{\pgfqpoint{3.401955in}{2.800390in}}%
\pgfpathlineto{\pgfqpoint{3.401955in}{2.797441in}}%
\pgfpathmoveto{\pgfqpoint{3.383791in}{2.800390in}}%
\pgfpathlineto{\pgfqpoint{3.383791in}{2.800390in}}%
\pgfpathlineto{\pgfqpoint{3.383791in}{2.803339in}}%
\pgfpathlineto{\pgfqpoint{3.388332in}{2.803339in}}%
\pgfpathlineto{\pgfqpoint{3.388332in}{2.800390in}}%
\pgfpathmoveto{\pgfqpoint{3.383791in}{2.803339in}}%
\pgfpathlineto{\pgfqpoint{3.383791in}{2.803339in}}%
\pgfpathlineto{\pgfqpoint{3.383791in}{2.806288in}}%
\pgfpathlineto{\pgfqpoint{3.388332in}{2.806288in}}%
\pgfpathlineto{\pgfqpoint{3.388332in}{2.803339in}}%
\pgfpathmoveto{\pgfqpoint{3.388332in}{2.800390in}}%
\pgfpathlineto{\pgfqpoint{3.388332in}{2.800390in}}%
\pgfpathlineto{\pgfqpoint{3.388332in}{2.803339in}}%
\pgfpathlineto{\pgfqpoint{3.392873in}{2.803339in}}%
\pgfpathlineto{\pgfqpoint{3.392873in}{2.800390in}}%
\pgfpathmoveto{\pgfqpoint{3.388332in}{2.803339in}}%
\pgfpathlineto{\pgfqpoint{3.388332in}{2.803339in}}%
\pgfpathlineto{\pgfqpoint{3.388332in}{2.806288in}}%
\pgfpathlineto{\pgfqpoint{3.392873in}{2.806288in}}%
\pgfpathlineto{\pgfqpoint{3.392873in}{2.803339in}}%
\pgfpathmoveto{\pgfqpoint{3.383791in}{2.806288in}}%
\pgfpathlineto{\pgfqpoint{3.383791in}{2.806288in}}%
\pgfpathlineto{\pgfqpoint{3.383791in}{2.809237in}}%
\pgfpathlineto{\pgfqpoint{3.388332in}{2.809237in}}%
\pgfpathlineto{\pgfqpoint{3.388332in}{2.806288in}}%
\pgfpathmoveto{\pgfqpoint{3.383791in}{2.809237in}}%
\pgfpathlineto{\pgfqpoint{3.383791in}{2.809237in}}%
\pgfpathlineto{\pgfqpoint{3.383791in}{2.812187in}}%
\pgfpathlineto{\pgfqpoint{3.388332in}{2.812187in}}%
\pgfpathlineto{\pgfqpoint{3.388332in}{2.809237in}}%
\pgfpathmoveto{\pgfqpoint{3.388332in}{2.806288in}}%
\pgfpathlineto{\pgfqpoint{3.388332in}{2.806288in}}%
\pgfpathlineto{\pgfqpoint{3.388332in}{2.809237in}}%
\pgfpathlineto{\pgfqpoint{3.392873in}{2.809237in}}%
\pgfpathlineto{\pgfqpoint{3.392873in}{2.806288in}}%
\pgfpathmoveto{\pgfqpoint{3.388332in}{2.809237in}}%
\pgfpathlineto{\pgfqpoint{3.388332in}{2.809237in}}%
\pgfpathlineto{\pgfqpoint{3.388332in}{2.812187in}}%
\pgfpathlineto{\pgfqpoint{3.392873in}{2.812187in}}%
\pgfpathlineto{\pgfqpoint{3.392873in}{2.809237in}}%
\pgfpathmoveto{\pgfqpoint{3.392873in}{2.800390in}}%
\pgfpathlineto{\pgfqpoint{3.392873in}{2.800390in}}%
\pgfpathlineto{\pgfqpoint{3.392873in}{2.803339in}}%
\pgfpathlineto{\pgfqpoint{3.397414in}{2.803339in}}%
\pgfpathlineto{\pgfqpoint{3.397414in}{2.800390in}}%
\pgfpathmoveto{\pgfqpoint{3.392873in}{2.803339in}}%
\pgfpathlineto{\pgfqpoint{3.392873in}{2.803339in}}%
\pgfpathlineto{\pgfqpoint{3.392873in}{2.806288in}}%
\pgfpathlineto{\pgfqpoint{3.397414in}{2.806288in}}%
\pgfpathlineto{\pgfqpoint{3.397414in}{2.803339in}}%
\pgfpathmoveto{\pgfqpoint{3.411037in}{2.764999in}}%
\pgfpathlineto{\pgfqpoint{3.411037in}{2.764999in}}%
\pgfpathlineto{\pgfqpoint{3.411037in}{2.767948in}}%
\pgfpathlineto{\pgfqpoint{3.415578in}{2.767948in}}%
\pgfpathlineto{\pgfqpoint{3.415578in}{2.764999in}}%
\pgfpathmoveto{\pgfqpoint{3.411037in}{2.767948in}}%
\pgfpathlineto{\pgfqpoint{3.411037in}{2.767948in}}%
\pgfpathlineto{\pgfqpoint{3.411037in}{2.770898in}}%
\pgfpathlineto{\pgfqpoint{3.415578in}{2.770898in}}%
\pgfpathlineto{\pgfqpoint{3.415578in}{2.767948in}}%
\pgfpathmoveto{\pgfqpoint{3.415578in}{2.764999in}}%
\pgfpathlineto{\pgfqpoint{3.415578in}{2.764999in}}%
\pgfpathlineto{\pgfqpoint{3.415578in}{2.767948in}}%
\pgfpathlineto{\pgfqpoint{3.420119in}{2.767948in}}%
\pgfpathlineto{\pgfqpoint{3.420119in}{2.764999in}}%
\pgfpathmoveto{\pgfqpoint{3.415578in}{2.767948in}}%
\pgfpathlineto{\pgfqpoint{3.415578in}{2.767948in}}%
\pgfpathlineto{\pgfqpoint{3.415578in}{2.770898in}}%
\pgfpathlineto{\pgfqpoint{3.420119in}{2.770898in}}%
\pgfpathlineto{\pgfqpoint{3.420119in}{2.767948in}}%
\pgfpathmoveto{\pgfqpoint{3.411037in}{2.770898in}}%
\pgfpathlineto{\pgfqpoint{3.411037in}{2.770898in}}%
\pgfpathlineto{\pgfqpoint{3.411037in}{2.773847in}}%
\pgfpathlineto{\pgfqpoint{3.415578in}{2.773847in}}%
\pgfpathlineto{\pgfqpoint{3.415578in}{2.770898in}}%
\pgfpathmoveto{\pgfqpoint{3.411037in}{2.773847in}}%
\pgfpathlineto{\pgfqpoint{3.411037in}{2.773847in}}%
\pgfpathlineto{\pgfqpoint{3.411037in}{2.776796in}}%
\pgfpathlineto{\pgfqpoint{3.415578in}{2.776796in}}%
\pgfpathlineto{\pgfqpoint{3.415578in}{2.773847in}}%
\pgfpathmoveto{\pgfqpoint{3.415578in}{2.770898in}}%
\pgfpathlineto{\pgfqpoint{3.415578in}{2.770898in}}%
\pgfpathlineto{\pgfqpoint{3.415578in}{2.773847in}}%
\pgfpathlineto{\pgfqpoint{3.420119in}{2.773847in}}%
\pgfpathlineto{\pgfqpoint{3.420119in}{2.770898in}}%
\pgfpathmoveto{\pgfqpoint{3.415578in}{2.773847in}}%
\pgfpathlineto{\pgfqpoint{3.415578in}{2.773847in}}%
\pgfpathlineto{\pgfqpoint{3.415578in}{2.776796in}}%
\pgfpathlineto{\pgfqpoint{3.420119in}{2.776796in}}%
\pgfpathlineto{\pgfqpoint{3.420119in}{2.773847in}}%
\pgfpathmoveto{\pgfqpoint{3.401955in}{2.776796in}}%
\pgfpathlineto{\pgfqpoint{3.401955in}{2.776796in}}%
\pgfpathlineto{\pgfqpoint{3.401955in}{2.779745in}}%
\pgfpathlineto{\pgfqpoint{3.406496in}{2.779745in}}%
\pgfpathlineto{\pgfqpoint{3.406496in}{2.776796in}}%
\pgfpathmoveto{\pgfqpoint{3.401955in}{2.779745in}}%
\pgfpathlineto{\pgfqpoint{3.401955in}{2.779745in}}%
\pgfpathlineto{\pgfqpoint{3.401955in}{2.782694in}}%
\pgfpathlineto{\pgfqpoint{3.406496in}{2.782694in}}%
\pgfpathlineto{\pgfqpoint{3.406496in}{2.779745in}}%
\pgfpathmoveto{\pgfqpoint{3.406496in}{2.776796in}}%
\pgfpathlineto{\pgfqpoint{3.406496in}{2.776796in}}%
\pgfpathlineto{\pgfqpoint{3.406496in}{2.779745in}}%
\pgfpathlineto{\pgfqpoint{3.411037in}{2.779745in}}%
\pgfpathlineto{\pgfqpoint{3.411037in}{2.776796in}}%
\pgfpathmoveto{\pgfqpoint{3.406496in}{2.779745in}}%
\pgfpathlineto{\pgfqpoint{3.406496in}{2.779745in}}%
\pgfpathlineto{\pgfqpoint{3.406496in}{2.782694in}}%
\pgfpathlineto{\pgfqpoint{3.411037in}{2.782694in}}%
\pgfpathlineto{\pgfqpoint{3.411037in}{2.779745in}}%
\pgfpathmoveto{\pgfqpoint{3.401955in}{2.782694in}}%
\pgfpathlineto{\pgfqpoint{3.401955in}{2.782694in}}%
\pgfpathlineto{\pgfqpoint{3.401955in}{2.785644in}}%
\pgfpathlineto{\pgfqpoint{3.406496in}{2.785644in}}%
\pgfpathlineto{\pgfqpoint{3.406496in}{2.782694in}}%
\pgfpathmoveto{\pgfqpoint{3.401955in}{2.785644in}}%
\pgfpathlineto{\pgfqpoint{3.401955in}{2.785644in}}%
\pgfpathlineto{\pgfqpoint{3.401955in}{2.788593in}}%
\pgfpathlineto{\pgfqpoint{3.406496in}{2.788593in}}%
\pgfpathlineto{\pgfqpoint{3.406496in}{2.785644in}}%
\pgfpathmoveto{\pgfqpoint{3.406496in}{2.782694in}}%
\pgfpathlineto{\pgfqpoint{3.406496in}{2.782694in}}%
\pgfpathlineto{\pgfqpoint{3.406496in}{2.785644in}}%
\pgfpathlineto{\pgfqpoint{3.411037in}{2.785644in}}%
\pgfpathlineto{\pgfqpoint{3.411037in}{2.782694in}}%
\pgfpathmoveto{\pgfqpoint{3.406496in}{2.785644in}}%
\pgfpathlineto{\pgfqpoint{3.406496in}{2.785644in}}%
\pgfpathlineto{\pgfqpoint{3.406496in}{2.788593in}}%
\pgfpathlineto{\pgfqpoint{3.411037in}{2.788593in}}%
\pgfpathlineto{\pgfqpoint{3.411037in}{2.785644in}}%
\pgfpathmoveto{\pgfqpoint{3.411037in}{2.776796in}}%
\pgfpathlineto{\pgfqpoint{3.411037in}{2.776796in}}%
\pgfpathlineto{\pgfqpoint{3.411037in}{2.779745in}}%
\pgfpathlineto{\pgfqpoint{3.415578in}{2.779745in}}%
\pgfpathlineto{\pgfqpoint{3.415578in}{2.776796in}}%
\pgfpathmoveto{\pgfqpoint{3.411037in}{2.779745in}}%
\pgfpathlineto{\pgfqpoint{3.411037in}{2.779745in}}%
\pgfpathlineto{\pgfqpoint{3.411037in}{2.782694in}}%
\pgfpathlineto{\pgfqpoint{3.415578in}{2.782694in}}%
\pgfpathlineto{\pgfqpoint{3.415578in}{2.779745in}}%
\pgfpathmoveto{\pgfqpoint{3.420119in}{2.764999in}}%
\pgfpathlineto{\pgfqpoint{3.420119in}{2.764999in}}%
\pgfpathlineto{\pgfqpoint{3.420119in}{2.767948in}}%
\pgfpathlineto{\pgfqpoint{3.424660in}{2.767948in}}%
\pgfpathlineto{\pgfqpoint{3.424660in}{2.764999in}}%
\pgfpathmoveto{\pgfqpoint{3.420119in}{2.767948in}}%
\pgfpathlineto{\pgfqpoint{3.420119in}{2.767948in}}%
\pgfpathlineto{\pgfqpoint{3.420119in}{2.770898in}}%
\pgfpathlineto{\pgfqpoint{3.424660in}{2.770898in}}%
\pgfpathlineto{\pgfqpoint{3.424660in}{2.767948in}}%
\pgfpathmoveto{\pgfqpoint{3.401955in}{2.788593in}}%
\pgfpathlineto{\pgfqpoint{3.401955in}{2.788593in}}%
\pgfpathlineto{\pgfqpoint{3.401955in}{2.791542in}}%
\pgfpathlineto{\pgfqpoint{3.406496in}{2.791542in}}%
\pgfpathlineto{\pgfqpoint{3.406496in}{2.788593in}}%
\pgfpathmoveto{\pgfqpoint{3.401955in}{2.791542in}}%
\pgfpathlineto{\pgfqpoint{3.401955in}{2.791542in}}%
\pgfpathlineto{\pgfqpoint{3.401955in}{2.794491in}}%
\pgfpathlineto{\pgfqpoint{3.406496in}{2.794491in}}%
\pgfpathlineto{\pgfqpoint{3.406496in}{2.791542in}}%
\pgfpathmoveto{\pgfqpoint{3.374709in}{2.812187in}}%
\pgfpathlineto{\pgfqpoint{3.374709in}{2.812187in}}%
\pgfpathlineto{\pgfqpoint{3.374709in}{2.815136in}}%
\pgfpathlineto{\pgfqpoint{3.379250in}{2.815136in}}%
\pgfpathlineto{\pgfqpoint{3.379250in}{2.812187in}}%
\pgfpathmoveto{\pgfqpoint{3.374709in}{2.815136in}}%
\pgfpathlineto{\pgfqpoint{3.374709in}{2.815136in}}%
\pgfpathlineto{\pgfqpoint{3.374709in}{2.818085in}}%
\pgfpathlineto{\pgfqpoint{3.379250in}{2.818085in}}%
\pgfpathlineto{\pgfqpoint{3.379250in}{2.815136in}}%
\pgfpathmoveto{\pgfqpoint{3.379250in}{2.812187in}}%
\pgfpathlineto{\pgfqpoint{3.379250in}{2.812187in}}%
\pgfpathlineto{\pgfqpoint{3.379250in}{2.815136in}}%
\pgfpathlineto{\pgfqpoint{3.383791in}{2.815136in}}%
\pgfpathlineto{\pgfqpoint{3.383791in}{2.812187in}}%
\pgfpathmoveto{\pgfqpoint{3.379250in}{2.815136in}}%
\pgfpathlineto{\pgfqpoint{3.379250in}{2.815136in}}%
\pgfpathlineto{\pgfqpoint{3.379250in}{2.818085in}}%
\pgfpathlineto{\pgfqpoint{3.383791in}{2.818085in}}%
\pgfpathlineto{\pgfqpoint{3.383791in}{2.815136in}}%
\pgfpathmoveto{\pgfqpoint{3.374709in}{2.818085in}}%
\pgfpathlineto{\pgfqpoint{3.374709in}{2.818085in}}%
\pgfpathlineto{\pgfqpoint{3.374709in}{2.821034in}}%
\pgfpathlineto{\pgfqpoint{3.379250in}{2.821034in}}%
\pgfpathlineto{\pgfqpoint{3.379250in}{2.818085in}}%
\pgfpathmoveto{\pgfqpoint{3.374709in}{2.821034in}}%
\pgfpathlineto{\pgfqpoint{3.374709in}{2.821034in}}%
\pgfpathlineto{\pgfqpoint{3.374709in}{2.823983in}}%
\pgfpathlineto{\pgfqpoint{3.379250in}{2.823983in}}%
\pgfpathlineto{\pgfqpoint{3.379250in}{2.821034in}}%
\pgfpathmoveto{\pgfqpoint{3.379250in}{2.818085in}}%
\pgfpathlineto{\pgfqpoint{3.379250in}{2.818085in}}%
\pgfpathlineto{\pgfqpoint{3.379250in}{2.821034in}}%
\pgfpathlineto{\pgfqpoint{3.383791in}{2.821034in}}%
\pgfpathlineto{\pgfqpoint{3.383791in}{2.818085in}}%
\pgfpathmoveto{\pgfqpoint{3.379250in}{2.821034in}}%
\pgfpathlineto{\pgfqpoint{3.379250in}{2.821034in}}%
\pgfpathlineto{\pgfqpoint{3.379250in}{2.823983in}}%
\pgfpathlineto{\pgfqpoint{3.383791in}{2.823983in}}%
\pgfpathlineto{\pgfqpoint{3.383791in}{2.821034in}}%
\pgfpathmoveto{\pgfqpoint{3.365627in}{2.823983in}}%
\pgfpathlineto{\pgfqpoint{3.365627in}{2.823983in}}%
\pgfpathlineto{\pgfqpoint{3.365627in}{2.826933in}}%
\pgfpathlineto{\pgfqpoint{3.370168in}{2.826933in}}%
\pgfpathlineto{\pgfqpoint{3.370168in}{2.823983in}}%
\pgfpathmoveto{\pgfqpoint{3.365627in}{2.826933in}}%
\pgfpathlineto{\pgfqpoint{3.365627in}{2.826933in}}%
\pgfpathlineto{\pgfqpoint{3.365627in}{2.829882in}}%
\pgfpathlineto{\pgfqpoint{3.370168in}{2.829882in}}%
\pgfpathlineto{\pgfqpoint{3.370168in}{2.826933in}}%
\pgfpathmoveto{\pgfqpoint{3.370168in}{2.823983in}}%
\pgfpathlineto{\pgfqpoint{3.370168in}{2.823983in}}%
\pgfpathlineto{\pgfqpoint{3.370168in}{2.826933in}}%
\pgfpathlineto{\pgfqpoint{3.374709in}{2.826933in}}%
\pgfpathlineto{\pgfqpoint{3.374709in}{2.823983in}}%
\pgfpathmoveto{\pgfqpoint{3.370168in}{2.826933in}}%
\pgfpathlineto{\pgfqpoint{3.370168in}{2.826933in}}%
\pgfpathlineto{\pgfqpoint{3.370168in}{2.829882in}}%
\pgfpathlineto{\pgfqpoint{3.374709in}{2.829882in}}%
\pgfpathlineto{\pgfqpoint{3.374709in}{2.826933in}}%
\pgfpathmoveto{\pgfqpoint{3.365627in}{2.829882in}}%
\pgfpathlineto{\pgfqpoint{3.365627in}{2.829882in}}%
\pgfpathlineto{\pgfqpoint{3.365627in}{2.832831in}}%
\pgfpathlineto{\pgfqpoint{3.370168in}{2.832831in}}%
\pgfpathlineto{\pgfqpoint{3.370168in}{2.829882in}}%
\pgfpathmoveto{\pgfqpoint{3.365627in}{2.832831in}}%
\pgfpathlineto{\pgfqpoint{3.365627in}{2.832831in}}%
\pgfpathlineto{\pgfqpoint{3.365627in}{2.835780in}}%
\pgfpathlineto{\pgfqpoint{3.370168in}{2.835780in}}%
\pgfpathlineto{\pgfqpoint{3.370168in}{2.832831in}}%
\pgfpathmoveto{\pgfqpoint{3.370168in}{2.829882in}}%
\pgfpathlineto{\pgfqpoint{3.370168in}{2.829882in}}%
\pgfpathlineto{\pgfqpoint{3.370168in}{2.832831in}}%
\pgfpathlineto{\pgfqpoint{3.374709in}{2.832831in}}%
\pgfpathlineto{\pgfqpoint{3.374709in}{2.829882in}}%
\pgfpathmoveto{\pgfqpoint{3.370168in}{2.832831in}}%
\pgfpathlineto{\pgfqpoint{3.370168in}{2.832831in}}%
\pgfpathlineto{\pgfqpoint{3.370168in}{2.835780in}}%
\pgfpathlineto{\pgfqpoint{3.374709in}{2.835780in}}%
\pgfpathlineto{\pgfqpoint{3.374709in}{2.832831in}}%
\pgfpathmoveto{\pgfqpoint{3.374709in}{2.823983in}}%
\pgfpathlineto{\pgfqpoint{3.374709in}{2.823983in}}%
\pgfpathlineto{\pgfqpoint{3.374709in}{2.826933in}}%
\pgfpathlineto{\pgfqpoint{3.379250in}{2.826933in}}%
\pgfpathlineto{\pgfqpoint{3.379250in}{2.823983in}}%
\pgfpathmoveto{\pgfqpoint{3.374709in}{2.826933in}}%
\pgfpathlineto{\pgfqpoint{3.374709in}{2.826933in}}%
\pgfpathlineto{\pgfqpoint{3.374709in}{2.829882in}}%
\pgfpathlineto{\pgfqpoint{3.379250in}{2.829882in}}%
\pgfpathlineto{\pgfqpoint{3.379250in}{2.826933in}}%
\pgfpathmoveto{\pgfqpoint{3.383791in}{2.812187in}}%
\pgfpathlineto{\pgfqpoint{3.383791in}{2.812187in}}%
\pgfpathlineto{\pgfqpoint{3.383791in}{2.815136in}}%
\pgfpathlineto{\pgfqpoint{3.388332in}{2.815136in}}%
\pgfpathlineto{\pgfqpoint{3.388332in}{2.812187in}}%
\pgfpathmoveto{\pgfqpoint{3.383791in}{2.815136in}}%
\pgfpathlineto{\pgfqpoint{3.383791in}{2.815136in}}%
\pgfpathlineto{\pgfqpoint{3.383791in}{2.818085in}}%
\pgfpathlineto{\pgfqpoint{3.388332in}{2.818085in}}%
\pgfpathlineto{\pgfqpoint{3.388332in}{2.815136in}}%
\pgfpathmoveto{\pgfqpoint{3.365627in}{2.835780in}}%
\pgfpathlineto{\pgfqpoint{3.365627in}{2.835780in}}%
\pgfpathlineto{\pgfqpoint{3.365627in}{2.838730in}}%
\pgfpathlineto{\pgfqpoint{3.370168in}{2.838730in}}%
\pgfpathlineto{\pgfqpoint{3.370168in}{2.835780in}}%
\pgfpathmoveto{\pgfqpoint{3.365627in}{2.838730in}}%
\pgfpathlineto{\pgfqpoint{3.365627in}{2.838730in}}%
\pgfpathlineto{\pgfqpoint{3.365627in}{2.841679in}}%
\pgfpathlineto{\pgfqpoint{3.370168in}{2.841679in}}%
\pgfpathlineto{\pgfqpoint{3.370168in}{2.838730in}}%
\pgfpathmoveto{\pgfqpoint{3.510939in}{2.293122in}}%
\pgfpathlineto{\pgfqpoint{3.510939in}{2.293122in}}%
\pgfpathlineto{\pgfqpoint{3.510939in}{2.296072in}}%
\pgfpathlineto{\pgfqpoint{3.515479in}{2.296072in}}%
\pgfpathlineto{\pgfqpoint{3.515479in}{2.293122in}}%
\pgfpathmoveto{\pgfqpoint{3.510939in}{2.296072in}}%
\pgfpathlineto{\pgfqpoint{3.510939in}{2.296072in}}%
\pgfpathlineto{\pgfqpoint{3.510939in}{2.299021in}}%
\pgfpathlineto{\pgfqpoint{3.515479in}{2.299021in}}%
\pgfpathlineto{\pgfqpoint{3.515479in}{2.296072in}}%
\pgfpathmoveto{\pgfqpoint{3.515479in}{2.296072in}}%
\pgfpathlineto{\pgfqpoint{3.515479in}{2.296072in}}%
\pgfpathlineto{\pgfqpoint{3.515479in}{2.299021in}}%
\pgfpathlineto{\pgfqpoint{3.520020in}{2.299021in}}%
\pgfpathlineto{\pgfqpoint{3.520020in}{2.296072in}}%
\pgfpathmoveto{\pgfqpoint{3.510939in}{2.299021in}}%
\pgfpathlineto{\pgfqpoint{3.510939in}{2.299021in}}%
\pgfpathlineto{\pgfqpoint{3.510939in}{2.301970in}}%
\pgfpathlineto{\pgfqpoint{3.515479in}{2.301970in}}%
\pgfpathlineto{\pgfqpoint{3.515479in}{2.299021in}}%
\pgfpathmoveto{\pgfqpoint{3.510939in}{2.301970in}}%
\pgfpathlineto{\pgfqpoint{3.510939in}{2.301970in}}%
\pgfpathlineto{\pgfqpoint{3.510939in}{2.304919in}}%
\pgfpathlineto{\pgfqpoint{3.515479in}{2.304919in}}%
\pgfpathlineto{\pgfqpoint{3.515479in}{2.301970in}}%
\pgfpathmoveto{\pgfqpoint{3.515479in}{2.299021in}}%
\pgfpathlineto{\pgfqpoint{3.515479in}{2.299021in}}%
\pgfpathlineto{\pgfqpoint{3.515479in}{2.301970in}}%
\pgfpathlineto{\pgfqpoint{3.520020in}{2.301970in}}%
\pgfpathlineto{\pgfqpoint{3.520020in}{2.299021in}}%
\pgfpathmoveto{\pgfqpoint{3.515479in}{2.301970in}}%
\pgfpathlineto{\pgfqpoint{3.515479in}{2.301970in}}%
\pgfpathlineto{\pgfqpoint{3.515479in}{2.304919in}}%
\pgfpathlineto{\pgfqpoint{3.520020in}{2.304919in}}%
\pgfpathlineto{\pgfqpoint{3.520020in}{2.301970in}}%
\pgfpathmoveto{\pgfqpoint{3.520020in}{2.299021in}}%
\pgfpathlineto{\pgfqpoint{3.520020in}{2.299021in}}%
\pgfpathlineto{\pgfqpoint{3.520020in}{2.301970in}}%
\pgfpathlineto{\pgfqpoint{3.524561in}{2.301970in}}%
\pgfpathlineto{\pgfqpoint{3.524561in}{2.299021in}}%
\pgfpathmoveto{\pgfqpoint{3.520020in}{2.301970in}}%
\pgfpathlineto{\pgfqpoint{3.520020in}{2.301970in}}%
\pgfpathlineto{\pgfqpoint{3.520020in}{2.304919in}}%
\pgfpathlineto{\pgfqpoint{3.524561in}{2.304919in}}%
\pgfpathlineto{\pgfqpoint{3.524561in}{2.301970in}}%
\pgfpathmoveto{\pgfqpoint{3.524561in}{2.301970in}}%
\pgfpathlineto{\pgfqpoint{3.524561in}{2.301970in}}%
\pgfpathlineto{\pgfqpoint{3.524561in}{2.304919in}}%
\pgfpathlineto{\pgfqpoint{3.529102in}{2.304919in}}%
\pgfpathlineto{\pgfqpoint{3.529102in}{2.301970in}}%
\pgfpathmoveto{\pgfqpoint{3.520020in}{2.304919in}}%
\pgfpathlineto{\pgfqpoint{3.520020in}{2.304919in}}%
\pgfpathlineto{\pgfqpoint{3.520020in}{2.307869in}}%
\pgfpathlineto{\pgfqpoint{3.524561in}{2.307869in}}%
\pgfpathlineto{\pgfqpoint{3.524561in}{2.304919in}}%
\pgfpathmoveto{\pgfqpoint{3.520020in}{2.307869in}}%
\pgfpathlineto{\pgfqpoint{3.520020in}{2.307869in}}%
\pgfpathlineto{\pgfqpoint{3.520020in}{2.310818in}}%
\pgfpathlineto{\pgfqpoint{3.524561in}{2.310818in}}%
\pgfpathlineto{\pgfqpoint{3.524561in}{2.307869in}}%
\pgfpathmoveto{\pgfqpoint{3.524561in}{2.304919in}}%
\pgfpathlineto{\pgfqpoint{3.524561in}{2.304919in}}%
\pgfpathlineto{\pgfqpoint{3.524561in}{2.307869in}}%
\pgfpathlineto{\pgfqpoint{3.529102in}{2.307869in}}%
\pgfpathlineto{\pgfqpoint{3.529102in}{2.304919in}}%
\pgfpathmoveto{\pgfqpoint{3.524561in}{2.307869in}}%
\pgfpathlineto{\pgfqpoint{3.524561in}{2.307869in}}%
\pgfpathlineto{\pgfqpoint{3.524561in}{2.310818in}}%
\pgfpathlineto{\pgfqpoint{3.529102in}{2.310818in}}%
\pgfpathlineto{\pgfqpoint{3.529102in}{2.307869in}}%
\pgfpathmoveto{\pgfqpoint{3.529102in}{2.304919in}}%
\pgfpathlineto{\pgfqpoint{3.529102in}{2.304919in}}%
\pgfpathlineto{\pgfqpoint{3.529102in}{2.307869in}}%
\pgfpathlineto{\pgfqpoint{3.533643in}{2.307869in}}%
\pgfpathlineto{\pgfqpoint{3.533643in}{2.304919in}}%
\pgfpathmoveto{\pgfqpoint{3.529102in}{2.307869in}}%
\pgfpathlineto{\pgfqpoint{3.529102in}{2.307869in}}%
\pgfpathlineto{\pgfqpoint{3.529102in}{2.310818in}}%
\pgfpathlineto{\pgfqpoint{3.533643in}{2.310818in}}%
\pgfpathlineto{\pgfqpoint{3.533643in}{2.307869in}}%
\pgfpathmoveto{\pgfqpoint{3.533643in}{2.307869in}}%
\pgfpathlineto{\pgfqpoint{3.533643in}{2.307869in}}%
\pgfpathlineto{\pgfqpoint{3.533643in}{2.310818in}}%
\pgfpathlineto{\pgfqpoint{3.538184in}{2.310818in}}%
\pgfpathlineto{\pgfqpoint{3.538184in}{2.307869in}}%
\pgfpathmoveto{\pgfqpoint{3.529102in}{2.310818in}}%
\pgfpathlineto{\pgfqpoint{3.529102in}{2.310818in}}%
\pgfpathlineto{\pgfqpoint{3.529102in}{2.313767in}}%
\pgfpathlineto{\pgfqpoint{3.533643in}{2.313767in}}%
\pgfpathlineto{\pgfqpoint{3.533643in}{2.310818in}}%
\pgfpathmoveto{\pgfqpoint{3.529102in}{2.313767in}}%
\pgfpathlineto{\pgfqpoint{3.529102in}{2.313767in}}%
\pgfpathlineto{\pgfqpoint{3.529102in}{2.316716in}}%
\pgfpathlineto{\pgfqpoint{3.533643in}{2.316716in}}%
\pgfpathlineto{\pgfqpoint{3.533643in}{2.313767in}}%
\pgfpathmoveto{\pgfqpoint{3.533643in}{2.310818in}}%
\pgfpathlineto{\pgfqpoint{3.533643in}{2.310818in}}%
\pgfpathlineto{\pgfqpoint{3.533643in}{2.313767in}}%
\pgfpathlineto{\pgfqpoint{3.538184in}{2.313767in}}%
\pgfpathlineto{\pgfqpoint{3.538184in}{2.310818in}}%
\pgfpathmoveto{\pgfqpoint{3.533643in}{2.313767in}}%
\pgfpathlineto{\pgfqpoint{3.533643in}{2.313767in}}%
\pgfpathlineto{\pgfqpoint{3.533643in}{2.316716in}}%
\pgfpathlineto{\pgfqpoint{3.538184in}{2.316716in}}%
\pgfpathlineto{\pgfqpoint{3.538184in}{2.313767in}}%
\pgfpathmoveto{\pgfqpoint{3.538184in}{2.310818in}}%
\pgfpathlineto{\pgfqpoint{3.538184in}{2.310818in}}%
\pgfpathlineto{\pgfqpoint{3.538184in}{2.313767in}}%
\pgfpathlineto{\pgfqpoint{3.542725in}{2.313767in}}%
\pgfpathlineto{\pgfqpoint{3.542725in}{2.310818in}}%
\pgfpathmoveto{\pgfqpoint{3.538184in}{2.313767in}}%
\pgfpathlineto{\pgfqpoint{3.538184in}{2.313767in}}%
\pgfpathlineto{\pgfqpoint{3.538184in}{2.316716in}}%
\pgfpathlineto{\pgfqpoint{3.542725in}{2.316716in}}%
\pgfpathlineto{\pgfqpoint{3.542725in}{2.313767in}}%
\pgfpathmoveto{\pgfqpoint{3.542725in}{2.313767in}}%
\pgfpathlineto{\pgfqpoint{3.542725in}{2.313767in}}%
\pgfpathlineto{\pgfqpoint{3.542725in}{2.316716in}}%
\pgfpathlineto{\pgfqpoint{3.547265in}{2.316716in}}%
\pgfpathlineto{\pgfqpoint{3.547265in}{2.313767in}}%
\pgfpathmoveto{\pgfqpoint{3.538184in}{2.316716in}}%
\pgfpathlineto{\pgfqpoint{3.538184in}{2.316716in}}%
\pgfpathlineto{\pgfqpoint{3.538184in}{2.319666in}}%
\pgfpathlineto{\pgfqpoint{3.542725in}{2.319666in}}%
\pgfpathlineto{\pgfqpoint{3.542725in}{2.316716in}}%
\pgfpathmoveto{\pgfqpoint{3.538184in}{2.319666in}}%
\pgfpathlineto{\pgfqpoint{3.538184in}{2.319666in}}%
\pgfpathlineto{\pgfqpoint{3.538184in}{2.322615in}}%
\pgfpathlineto{\pgfqpoint{3.542725in}{2.322615in}}%
\pgfpathlineto{\pgfqpoint{3.542725in}{2.319666in}}%
\pgfpathmoveto{\pgfqpoint{3.542725in}{2.316716in}}%
\pgfpathlineto{\pgfqpoint{3.542725in}{2.316716in}}%
\pgfpathlineto{\pgfqpoint{3.542725in}{2.319666in}}%
\pgfpathlineto{\pgfqpoint{3.547265in}{2.319666in}}%
\pgfpathlineto{\pgfqpoint{3.547265in}{2.316716in}}%
\pgfpathmoveto{\pgfqpoint{3.542725in}{2.319666in}}%
\pgfpathlineto{\pgfqpoint{3.542725in}{2.319666in}}%
\pgfpathlineto{\pgfqpoint{3.542725in}{2.322615in}}%
\pgfpathlineto{\pgfqpoint{3.547265in}{2.322615in}}%
\pgfpathlineto{\pgfqpoint{3.547265in}{2.319666in}}%
\pgfpathmoveto{\pgfqpoint{3.547265in}{2.316716in}}%
\pgfpathlineto{\pgfqpoint{3.547265in}{2.316716in}}%
\pgfpathlineto{\pgfqpoint{3.547265in}{2.319666in}}%
\pgfpathlineto{\pgfqpoint{3.551806in}{2.319666in}}%
\pgfpathlineto{\pgfqpoint{3.551806in}{2.316716in}}%
\pgfpathmoveto{\pgfqpoint{3.547265in}{2.319666in}}%
\pgfpathlineto{\pgfqpoint{3.547265in}{2.319666in}}%
\pgfpathlineto{\pgfqpoint{3.547265in}{2.322615in}}%
\pgfpathlineto{\pgfqpoint{3.551806in}{2.322615in}}%
\pgfpathlineto{\pgfqpoint{3.551806in}{2.319666in}}%
\pgfpathmoveto{\pgfqpoint{3.551806in}{2.319666in}}%
\pgfpathlineto{\pgfqpoint{3.551806in}{2.319666in}}%
\pgfpathlineto{\pgfqpoint{3.551806in}{2.322615in}}%
\pgfpathlineto{\pgfqpoint{3.556347in}{2.322615in}}%
\pgfpathlineto{\pgfqpoint{3.556347in}{2.319666in}}%
\pgfpathmoveto{\pgfqpoint{3.547265in}{2.322615in}}%
\pgfpathlineto{\pgfqpoint{3.547265in}{2.322615in}}%
\pgfpathlineto{\pgfqpoint{3.547265in}{2.325564in}}%
\pgfpathlineto{\pgfqpoint{3.551806in}{2.325564in}}%
\pgfpathlineto{\pgfqpoint{3.551806in}{2.322615in}}%
\pgfpathmoveto{\pgfqpoint{3.547265in}{2.325564in}}%
\pgfpathlineto{\pgfqpoint{3.547265in}{2.325564in}}%
\pgfpathlineto{\pgfqpoint{3.547265in}{2.328514in}}%
\pgfpathlineto{\pgfqpoint{3.551806in}{2.328514in}}%
\pgfpathlineto{\pgfqpoint{3.551806in}{2.325564in}}%
\pgfpathmoveto{\pgfqpoint{3.551806in}{2.322615in}}%
\pgfpathlineto{\pgfqpoint{3.551806in}{2.322615in}}%
\pgfpathlineto{\pgfqpoint{3.551806in}{2.325564in}}%
\pgfpathlineto{\pgfqpoint{3.556347in}{2.325564in}}%
\pgfpathlineto{\pgfqpoint{3.556347in}{2.322615in}}%
\pgfpathmoveto{\pgfqpoint{3.551806in}{2.325564in}}%
\pgfpathlineto{\pgfqpoint{3.551806in}{2.325564in}}%
\pgfpathlineto{\pgfqpoint{3.551806in}{2.328514in}}%
\pgfpathlineto{\pgfqpoint{3.556347in}{2.328514in}}%
\pgfpathlineto{\pgfqpoint{3.556347in}{2.325564in}}%
\pgfpathmoveto{\pgfqpoint{3.556347in}{2.322615in}}%
\pgfpathlineto{\pgfqpoint{3.556347in}{2.322615in}}%
\pgfpathlineto{\pgfqpoint{3.556347in}{2.325564in}}%
\pgfpathlineto{\pgfqpoint{3.560888in}{2.325564in}}%
\pgfpathlineto{\pgfqpoint{3.560888in}{2.322615in}}%
\pgfpathmoveto{\pgfqpoint{3.556347in}{2.325564in}}%
\pgfpathlineto{\pgfqpoint{3.556347in}{2.325564in}}%
\pgfpathlineto{\pgfqpoint{3.556347in}{2.328514in}}%
\pgfpathlineto{\pgfqpoint{3.560888in}{2.328514in}}%
\pgfpathlineto{\pgfqpoint{3.560888in}{2.325564in}}%
\pgfpathmoveto{\pgfqpoint{3.560888in}{2.325564in}}%
\pgfpathlineto{\pgfqpoint{3.560888in}{2.325564in}}%
\pgfpathlineto{\pgfqpoint{3.560888in}{2.328514in}}%
\pgfpathlineto{\pgfqpoint{3.565429in}{2.328514in}}%
\pgfpathlineto{\pgfqpoint{3.565429in}{2.325564in}}%
\pgfpathmoveto{\pgfqpoint{3.556347in}{2.328514in}}%
\pgfpathlineto{\pgfqpoint{3.556347in}{2.328514in}}%
\pgfpathlineto{\pgfqpoint{3.556347in}{2.331463in}}%
\pgfpathlineto{\pgfqpoint{3.560888in}{2.331463in}}%
\pgfpathlineto{\pgfqpoint{3.560888in}{2.328514in}}%
\pgfpathmoveto{\pgfqpoint{3.556347in}{2.331463in}}%
\pgfpathlineto{\pgfqpoint{3.556347in}{2.331463in}}%
\pgfpathlineto{\pgfqpoint{3.556347in}{2.334412in}}%
\pgfpathlineto{\pgfqpoint{3.560888in}{2.334412in}}%
\pgfpathlineto{\pgfqpoint{3.560888in}{2.331463in}}%
\pgfpathmoveto{\pgfqpoint{3.560888in}{2.328514in}}%
\pgfpathlineto{\pgfqpoint{3.560888in}{2.328514in}}%
\pgfpathlineto{\pgfqpoint{3.560888in}{2.331463in}}%
\pgfpathlineto{\pgfqpoint{3.565429in}{2.331463in}}%
\pgfpathlineto{\pgfqpoint{3.565429in}{2.328514in}}%
\pgfpathmoveto{\pgfqpoint{3.560888in}{2.331463in}}%
\pgfpathlineto{\pgfqpoint{3.560888in}{2.331463in}}%
\pgfpathlineto{\pgfqpoint{3.560888in}{2.334412in}}%
\pgfpathlineto{\pgfqpoint{3.565429in}{2.334412in}}%
\pgfpathlineto{\pgfqpoint{3.565429in}{2.331463in}}%
\pgfpathmoveto{\pgfqpoint{3.565429in}{2.328514in}}%
\pgfpathlineto{\pgfqpoint{3.565429in}{2.328514in}}%
\pgfpathlineto{\pgfqpoint{3.565429in}{2.331463in}}%
\pgfpathlineto{\pgfqpoint{3.569970in}{2.331463in}}%
\pgfpathlineto{\pgfqpoint{3.569970in}{2.328514in}}%
\pgfpathmoveto{\pgfqpoint{3.565429in}{2.331463in}}%
\pgfpathlineto{\pgfqpoint{3.565429in}{2.331463in}}%
\pgfpathlineto{\pgfqpoint{3.565429in}{2.334412in}}%
\pgfpathlineto{\pgfqpoint{3.569970in}{2.334412in}}%
\pgfpathlineto{\pgfqpoint{3.569970in}{2.331463in}}%
\pgfpathmoveto{\pgfqpoint{3.569970in}{2.331463in}}%
\pgfpathlineto{\pgfqpoint{3.569970in}{2.331463in}}%
\pgfpathlineto{\pgfqpoint{3.569970in}{2.334412in}}%
\pgfpathlineto{\pgfqpoint{3.574510in}{2.334412in}}%
\pgfpathlineto{\pgfqpoint{3.574510in}{2.331463in}}%
\pgfpathmoveto{\pgfqpoint{3.565429in}{2.334412in}}%
\pgfpathlineto{\pgfqpoint{3.565429in}{2.334412in}}%
\pgfpathlineto{\pgfqpoint{3.565429in}{2.337361in}}%
\pgfpathlineto{\pgfqpoint{3.569970in}{2.337361in}}%
\pgfpathlineto{\pgfqpoint{3.569970in}{2.334412in}}%
\pgfpathmoveto{\pgfqpoint{3.565429in}{2.337361in}}%
\pgfpathlineto{\pgfqpoint{3.565429in}{2.337361in}}%
\pgfpathlineto{\pgfqpoint{3.565429in}{2.340311in}}%
\pgfpathlineto{\pgfqpoint{3.569970in}{2.340311in}}%
\pgfpathlineto{\pgfqpoint{3.569970in}{2.337361in}}%
\pgfpathmoveto{\pgfqpoint{3.569970in}{2.334412in}}%
\pgfpathlineto{\pgfqpoint{3.569970in}{2.334412in}}%
\pgfpathlineto{\pgfqpoint{3.569970in}{2.337361in}}%
\pgfpathlineto{\pgfqpoint{3.574510in}{2.337361in}}%
\pgfpathlineto{\pgfqpoint{3.574510in}{2.334412in}}%
\pgfpathmoveto{\pgfqpoint{3.569970in}{2.337361in}}%
\pgfpathlineto{\pgfqpoint{3.569970in}{2.337361in}}%
\pgfpathlineto{\pgfqpoint{3.569970in}{2.340311in}}%
\pgfpathlineto{\pgfqpoint{3.574510in}{2.340311in}}%
\pgfpathlineto{\pgfqpoint{3.574510in}{2.337361in}}%
\pgfpathmoveto{\pgfqpoint{3.574510in}{2.334412in}}%
\pgfpathlineto{\pgfqpoint{3.574510in}{2.334412in}}%
\pgfpathlineto{\pgfqpoint{3.574510in}{2.337361in}}%
\pgfpathlineto{\pgfqpoint{3.579051in}{2.337361in}}%
\pgfpathlineto{\pgfqpoint{3.579051in}{2.334412in}}%
\pgfpathmoveto{\pgfqpoint{3.574510in}{2.337361in}}%
\pgfpathlineto{\pgfqpoint{3.574510in}{2.337361in}}%
\pgfpathlineto{\pgfqpoint{3.574510in}{2.340311in}}%
\pgfpathlineto{\pgfqpoint{3.579051in}{2.340311in}}%
\pgfpathlineto{\pgfqpoint{3.579051in}{2.337361in}}%
\pgfpathmoveto{\pgfqpoint{3.579051in}{2.337361in}}%
\pgfpathlineto{\pgfqpoint{3.579051in}{2.337361in}}%
\pgfpathlineto{\pgfqpoint{3.579051in}{2.340311in}}%
\pgfpathlineto{\pgfqpoint{3.583592in}{2.340311in}}%
\pgfpathlineto{\pgfqpoint{3.583592in}{2.337361in}}%
\pgfpathmoveto{\pgfqpoint{3.574510in}{2.340311in}}%
\pgfpathlineto{\pgfqpoint{3.574510in}{2.340311in}}%
\pgfpathlineto{\pgfqpoint{3.574510in}{2.343260in}}%
\pgfpathlineto{\pgfqpoint{3.579051in}{2.343260in}}%
\pgfpathlineto{\pgfqpoint{3.579051in}{2.340311in}}%
\pgfpathmoveto{\pgfqpoint{3.574510in}{2.343260in}}%
\pgfpathlineto{\pgfqpoint{3.574510in}{2.343260in}}%
\pgfpathlineto{\pgfqpoint{3.574510in}{2.346209in}}%
\pgfpathlineto{\pgfqpoint{3.579051in}{2.346209in}}%
\pgfpathlineto{\pgfqpoint{3.579051in}{2.343260in}}%
\pgfpathmoveto{\pgfqpoint{3.579051in}{2.340311in}}%
\pgfpathlineto{\pgfqpoint{3.579051in}{2.340311in}}%
\pgfpathlineto{\pgfqpoint{3.579051in}{2.343260in}}%
\pgfpathlineto{\pgfqpoint{3.583592in}{2.343260in}}%
\pgfpathlineto{\pgfqpoint{3.583592in}{2.340311in}}%
\pgfpathmoveto{\pgfqpoint{3.579051in}{2.343260in}}%
\pgfpathlineto{\pgfqpoint{3.579051in}{2.343260in}}%
\pgfpathlineto{\pgfqpoint{3.579051in}{2.346209in}}%
\pgfpathlineto{\pgfqpoint{3.583592in}{2.346209in}}%
\pgfpathlineto{\pgfqpoint{3.583592in}{2.343260in}}%
\pgfpathmoveto{\pgfqpoint{3.583592in}{2.340311in}}%
\pgfpathlineto{\pgfqpoint{3.583592in}{2.340311in}}%
\pgfpathlineto{\pgfqpoint{3.583592in}{2.343260in}}%
\pgfpathlineto{\pgfqpoint{3.588133in}{2.343260in}}%
\pgfpathlineto{\pgfqpoint{3.588133in}{2.340311in}}%
\pgfpathmoveto{\pgfqpoint{3.583592in}{2.343260in}}%
\pgfpathlineto{\pgfqpoint{3.583592in}{2.343260in}}%
\pgfpathlineto{\pgfqpoint{3.583592in}{2.346209in}}%
\pgfpathlineto{\pgfqpoint{3.588133in}{2.346209in}}%
\pgfpathlineto{\pgfqpoint{3.588133in}{2.343260in}}%
\pgfpathmoveto{\pgfqpoint{3.588133in}{2.343260in}}%
\pgfpathlineto{\pgfqpoint{3.588133in}{2.343260in}}%
\pgfpathlineto{\pgfqpoint{3.588133in}{2.346209in}}%
\pgfpathlineto{\pgfqpoint{3.592674in}{2.346209in}}%
\pgfpathlineto{\pgfqpoint{3.592674in}{2.343260in}}%
\pgfpathmoveto{\pgfqpoint{3.583592in}{2.346209in}}%
\pgfpathlineto{\pgfqpoint{3.583592in}{2.346209in}}%
\pgfpathlineto{\pgfqpoint{3.583592in}{2.349158in}}%
\pgfpathlineto{\pgfqpoint{3.588133in}{2.349158in}}%
\pgfpathlineto{\pgfqpoint{3.588133in}{2.346209in}}%
\pgfpathmoveto{\pgfqpoint{3.583592in}{2.349158in}}%
\pgfpathlineto{\pgfqpoint{3.583592in}{2.349158in}}%
\pgfpathlineto{\pgfqpoint{3.583592in}{2.352108in}}%
\pgfpathlineto{\pgfqpoint{3.588133in}{2.352108in}}%
\pgfpathlineto{\pgfqpoint{3.588133in}{2.349158in}}%
\pgfpathmoveto{\pgfqpoint{3.588133in}{2.346209in}}%
\pgfpathlineto{\pgfqpoint{3.588133in}{2.346209in}}%
\pgfpathlineto{\pgfqpoint{3.588133in}{2.349158in}}%
\pgfpathlineto{\pgfqpoint{3.592674in}{2.349158in}}%
\pgfpathlineto{\pgfqpoint{3.592674in}{2.346209in}}%
\pgfpathmoveto{\pgfqpoint{3.588133in}{2.349158in}}%
\pgfpathlineto{\pgfqpoint{3.588133in}{2.349158in}}%
\pgfpathlineto{\pgfqpoint{3.588133in}{2.352108in}}%
\pgfpathlineto{\pgfqpoint{3.592674in}{2.352108in}}%
\pgfpathlineto{\pgfqpoint{3.592674in}{2.349158in}}%
\pgfpathmoveto{\pgfqpoint{3.592674in}{2.346209in}}%
\pgfpathlineto{\pgfqpoint{3.592674in}{2.346209in}}%
\pgfpathlineto{\pgfqpoint{3.592674in}{2.349158in}}%
\pgfpathlineto{\pgfqpoint{3.597215in}{2.349158in}}%
\pgfpathlineto{\pgfqpoint{3.597215in}{2.346209in}}%
\pgfpathmoveto{\pgfqpoint{3.592674in}{2.349158in}}%
\pgfpathlineto{\pgfqpoint{3.592674in}{2.349158in}}%
\pgfpathlineto{\pgfqpoint{3.592674in}{2.352108in}}%
\pgfpathlineto{\pgfqpoint{3.597215in}{2.352108in}}%
\pgfpathlineto{\pgfqpoint{3.597215in}{2.349158in}}%
\pgfpathmoveto{\pgfqpoint{3.597215in}{2.349158in}}%
\pgfpathlineto{\pgfqpoint{3.597215in}{2.349158in}}%
\pgfpathlineto{\pgfqpoint{3.597215in}{2.352108in}}%
\pgfpathlineto{\pgfqpoint{3.601755in}{2.352108in}}%
\pgfpathlineto{\pgfqpoint{3.601755in}{2.349158in}}%
\pgfpathmoveto{\pgfqpoint{3.592674in}{2.352108in}}%
\pgfpathlineto{\pgfqpoint{3.592674in}{2.352108in}}%
\pgfpathlineto{\pgfqpoint{3.592674in}{2.355057in}}%
\pgfpathlineto{\pgfqpoint{3.597215in}{2.355057in}}%
\pgfpathlineto{\pgfqpoint{3.597215in}{2.352108in}}%
\pgfpathmoveto{\pgfqpoint{3.592674in}{2.355057in}}%
\pgfpathlineto{\pgfqpoint{3.592674in}{2.355057in}}%
\pgfpathlineto{\pgfqpoint{3.592674in}{2.358006in}}%
\pgfpathlineto{\pgfqpoint{3.597215in}{2.358006in}}%
\pgfpathlineto{\pgfqpoint{3.597215in}{2.355057in}}%
\pgfpathmoveto{\pgfqpoint{3.597215in}{2.352108in}}%
\pgfpathlineto{\pgfqpoint{3.597215in}{2.352108in}}%
\pgfpathlineto{\pgfqpoint{3.597215in}{2.355057in}}%
\pgfpathlineto{\pgfqpoint{3.601755in}{2.355057in}}%
\pgfpathlineto{\pgfqpoint{3.601755in}{2.352108in}}%
\pgfpathmoveto{\pgfqpoint{3.597215in}{2.355057in}}%
\pgfpathlineto{\pgfqpoint{3.597215in}{2.355057in}}%
\pgfpathlineto{\pgfqpoint{3.597215in}{2.358006in}}%
\pgfpathlineto{\pgfqpoint{3.601755in}{2.358006in}}%
\pgfpathlineto{\pgfqpoint{3.601755in}{2.355057in}}%
\pgfpathmoveto{\pgfqpoint{3.601755in}{2.352108in}}%
\pgfpathlineto{\pgfqpoint{3.601755in}{2.352108in}}%
\pgfpathlineto{\pgfqpoint{3.601755in}{2.355057in}}%
\pgfpathlineto{\pgfqpoint{3.606296in}{2.355057in}}%
\pgfpathlineto{\pgfqpoint{3.606296in}{2.352108in}}%
\pgfpathmoveto{\pgfqpoint{3.601755in}{2.355057in}}%
\pgfpathlineto{\pgfqpoint{3.601755in}{2.355057in}}%
\pgfpathlineto{\pgfqpoint{3.601755in}{2.358006in}}%
\pgfpathlineto{\pgfqpoint{3.606296in}{2.358006in}}%
\pgfpathlineto{\pgfqpoint{3.606296in}{2.355057in}}%
\pgfpathmoveto{\pgfqpoint{3.606296in}{2.355057in}}%
\pgfpathlineto{\pgfqpoint{3.606296in}{2.355057in}}%
\pgfpathlineto{\pgfqpoint{3.606296in}{2.358006in}}%
\pgfpathlineto{\pgfqpoint{3.610837in}{2.358006in}}%
\pgfpathlineto{\pgfqpoint{3.610837in}{2.355057in}}%
\pgfpathmoveto{\pgfqpoint{3.610837in}{2.355057in}}%
\pgfpathlineto{\pgfqpoint{3.610837in}{2.355057in}}%
\pgfpathlineto{\pgfqpoint{3.610837in}{2.358006in}}%
\pgfpathlineto{\pgfqpoint{3.615378in}{2.358006in}}%
\pgfpathlineto{\pgfqpoint{3.615378in}{2.355057in}}%
\pgfpathmoveto{\pgfqpoint{3.610837in}{2.358006in}}%
\pgfpathlineto{\pgfqpoint{3.610837in}{2.358006in}}%
\pgfpathlineto{\pgfqpoint{3.610837in}{2.360955in}}%
\pgfpathlineto{\pgfqpoint{3.615378in}{2.360955in}}%
\pgfpathlineto{\pgfqpoint{3.615378in}{2.358006in}}%
\pgfpathmoveto{\pgfqpoint{3.610837in}{2.360955in}}%
\pgfpathlineto{\pgfqpoint{3.610837in}{2.360955in}}%
\pgfpathlineto{\pgfqpoint{3.610837in}{2.363905in}}%
\pgfpathlineto{\pgfqpoint{3.615378in}{2.363905in}}%
\pgfpathlineto{\pgfqpoint{3.615378in}{2.360955in}}%
\pgfpathmoveto{\pgfqpoint{3.615378in}{2.358006in}}%
\pgfpathlineto{\pgfqpoint{3.615378in}{2.358006in}}%
\pgfpathlineto{\pgfqpoint{3.615378in}{2.360955in}}%
\pgfpathlineto{\pgfqpoint{3.619919in}{2.360955in}}%
\pgfpathlineto{\pgfqpoint{3.619919in}{2.358006in}}%
\pgfpathmoveto{\pgfqpoint{3.615378in}{2.360955in}}%
\pgfpathlineto{\pgfqpoint{3.615378in}{2.360955in}}%
\pgfpathlineto{\pgfqpoint{3.615378in}{2.363905in}}%
\pgfpathlineto{\pgfqpoint{3.619919in}{2.363905in}}%
\pgfpathlineto{\pgfqpoint{3.619919in}{2.360955in}}%
\pgfpathmoveto{\pgfqpoint{3.619919in}{2.360955in}}%
\pgfpathlineto{\pgfqpoint{3.619919in}{2.360955in}}%
\pgfpathlineto{\pgfqpoint{3.619919in}{2.363905in}}%
\pgfpathlineto{\pgfqpoint{3.624460in}{2.363905in}}%
\pgfpathlineto{\pgfqpoint{3.624460in}{2.360955in}}%
\pgfpathmoveto{\pgfqpoint{3.619919in}{2.363905in}}%
\pgfpathlineto{\pgfqpoint{3.619919in}{2.363905in}}%
\pgfpathlineto{\pgfqpoint{3.619919in}{2.366854in}}%
\pgfpathlineto{\pgfqpoint{3.624460in}{2.366854in}}%
\pgfpathlineto{\pgfqpoint{3.624460in}{2.363905in}}%
\pgfpathmoveto{\pgfqpoint{3.619919in}{2.366854in}}%
\pgfpathlineto{\pgfqpoint{3.619919in}{2.366854in}}%
\pgfpathlineto{\pgfqpoint{3.619919in}{2.369803in}}%
\pgfpathlineto{\pgfqpoint{3.624460in}{2.369803in}}%
\pgfpathlineto{\pgfqpoint{3.624460in}{2.366854in}}%
\pgfpathmoveto{\pgfqpoint{3.624460in}{2.363905in}}%
\pgfpathlineto{\pgfqpoint{3.624460in}{2.363905in}}%
\pgfpathlineto{\pgfqpoint{3.624460in}{2.366854in}}%
\pgfpathlineto{\pgfqpoint{3.629001in}{2.366854in}}%
\pgfpathlineto{\pgfqpoint{3.629001in}{2.363905in}}%
\pgfpathmoveto{\pgfqpoint{3.624460in}{2.366854in}}%
\pgfpathlineto{\pgfqpoint{3.624460in}{2.366854in}}%
\pgfpathlineto{\pgfqpoint{3.624460in}{2.369803in}}%
\pgfpathlineto{\pgfqpoint{3.629001in}{2.369803in}}%
\pgfpathlineto{\pgfqpoint{3.629001in}{2.366854in}}%
\pgfpathmoveto{\pgfqpoint{3.629001in}{2.366854in}}%
\pgfpathlineto{\pgfqpoint{3.629001in}{2.366854in}}%
\pgfpathlineto{\pgfqpoint{3.629001in}{2.369803in}}%
\pgfpathlineto{\pgfqpoint{3.633541in}{2.369803in}}%
\pgfpathlineto{\pgfqpoint{3.633541in}{2.366854in}}%
\pgfpathmoveto{\pgfqpoint{3.629001in}{2.369803in}}%
\pgfpathlineto{\pgfqpoint{3.629001in}{2.369803in}}%
\pgfpathlineto{\pgfqpoint{3.629001in}{2.372752in}}%
\pgfpathlineto{\pgfqpoint{3.633541in}{2.372752in}}%
\pgfpathlineto{\pgfqpoint{3.633541in}{2.369803in}}%
\pgfpathmoveto{\pgfqpoint{3.629001in}{2.372752in}}%
\pgfpathlineto{\pgfqpoint{3.629001in}{2.372752in}}%
\pgfpathlineto{\pgfqpoint{3.629001in}{2.375702in}}%
\pgfpathlineto{\pgfqpoint{3.633541in}{2.375702in}}%
\pgfpathlineto{\pgfqpoint{3.633541in}{2.372752in}}%
\pgfpathmoveto{\pgfqpoint{3.633541in}{2.369803in}}%
\pgfpathlineto{\pgfqpoint{3.633541in}{2.369803in}}%
\pgfpathlineto{\pgfqpoint{3.633541in}{2.372752in}}%
\pgfpathlineto{\pgfqpoint{3.638082in}{2.372752in}}%
\pgfpathlineto{\pgfqpoint{3.638082in}{2.369803in}}%
\pgfpathmoveto{\pgfqpoint{3.633541in}{2.372752in}}%
\pgfpathlineto{\pgfqpoint{3.633541in}{2.372752in}}%
\pgfpathlineto{\pgfqpoint{3.633541in}{2.375702in}}%
\pgfpathlineto{\pgfqpoint{3.638082in}{2.375702in}}%
\pgfpathlineto{\pgfqpoint{3.638082in}{2.372752in}}%
\pgfpathmoveto{\pgfqpoint{3.638082in}{2.372752in}}%
\pgfpathlineto{\pgfqpoint{3.638082in}{2.372752in}}%
\pgfpathlineto{\pgfqpoint{3.638082in}{2.375702in}}%
\pgfpathlineto{\pgfqpoint{3.642623in}{2.375702in}}%
\pgfpathlineto{\pgfqpoint{3.642623in}{2.372752in}}%
\pgfpathmoveto{\pgfqpoint{3.638082in}{2.375702in}}%
\pgfpathlineto{\pgfqpoint{3.638082in}{2.375702in}}%
\pgfpathlineto{\pgfqpoint{3.638082in}{2.378651in}}%
\pgfpathlineto{\pgfqpoint{3.642623in}{2.378651in}}%
\pgfpathlineto{\pgfqpoint{3.642623in}{2.375702in}}%
\pgfpathmoveto{\pgfqpoint{3.638082in}{2.378651in}}%
\pgfpathlineto{\pgfqpoint{3.638082in}{2.378651in}}%
\pgfpathlineto{\pgfqpoint{3.638082in}{2.381600in}}%
\pgfpathlineto{\pgfqpoint{3.642623in}{2.381600in}}%
\pgfpathlineto{\pgfqpoint{3.642623in}{2.378651in}}%
\pgfpathmoveto{\pgfqpoint{3.642623in}{2.375702in}}%
\pgfpathlineto{\pgfqpoint{3.642623in}{2.375702in}}%
\pgfpathlineto{\pgfqpoint{3.642623in}{2.378651in}}%
\pgfpathlineto{\pgfqpoint{3.647164in}{2.378651in}}%
\pgfpathlineto{\pgfqpoint{3.647164in}{2.375702in}}%
\pgfpathmoveto{\pgfqpoint{3.642623in}{2.378651in}}%
\pgfpathlineto{\pgfqpoint{3.642623in}{2.378651in}}%
\pgfpathlineto{\pgfqpoint{3.642623in}{2.381600in}}%
\pgfpathlineto{\pgfqpoint{3.647164in}{2.381600in}}%
\pgfpathlineto{\pgfqpoint{3.647164in}{2.378651in}}%
\pgfpathmoveto{\pgfqpoint{3.647164in}{2.378651in}}%
\pgfpathlineto{\pgfqpoint{3.647164in}{2.378651in}}%
\pgfpathlineto{\pgfqpoint{3.647164in}{2.381600in}}%
\pgfpathlineto{\pgfqpoint{3.651705in}{2.381600in}}%
\pgfpathlineto{\pgfqpoint{3.651705in}{2.378651in}}%
\pgfpathmoveto{\pgfqpoint{3.647164in}{2.381600in}}%
\pgfpathlineto{\pgfqpoint{3.647164in}{2.381600in}}%
\pgfpathlineto{\pgfqpoint{3.647164in}{2.384550in}}%
\pgfpathlineto{\pgfqpoint{3.651705in}{2.384550in}}%
\pgfpathlineto{\pgfqpoint{3.651705in}{2.381600in}}%
\pgfpathmoveto{\pgfqpoint{3.647164in}{2.384550in}}%
\pgfpathlineto{\pgfqpoint{3.647164in}{2.384550in}}%
\pgfpathlineto{\pgfqpoint{3.647164in}{2.387499in}}%
\pgfpathlineto{\pgfqpoint{3.651705in}{2.387499in}}%
\pgfpathlineto{\pgfqpoint{3.651705in}{2.384550in}}%
\pgfpathmoveto{\pgfqpoint{3.651705in}{2.381600in}}%
\pgfpathlineto{\pgfqpoint{3.651705in}{2.381600in}}%
\pgfpathlineto{\pgfqpoint{3.651705in}{2.384550in}}%
\pgfpathlineto{\pgfqpoint{3.656246in}{2.384550in}}%
\pgfpathlineto{\pgfqpoint{3.656246in}{2.381600in}}%
\pgfpathmoveto{\pgfqpoint{3.651705in}{2.384550in}}%
\pgfpathlineto{\pgfqpoint{3.651705in}{2.384550in}}%
\pgfpathlineto{\pgfqpoint{3.651705in}{2.387499in}}%
\pgfpathlineto{\pgfqpoint{3.656246in}{2.387499in}}%
\pgfpathlineto{\pgfqpoint{3.656246in}{2.384550in}}%
\pgfpathmoveto{\pgfqpoint{3.647164in}{2.458283in}}%
\pgfpathlineto{\pgfqpoint{3.647164in}{2.458283in}}%
\pgfpathlineto{\pgfqpoint{3.647164in}{2.461232in}}%
\pgfpathlineto{\pgfqpoint{3.651705in}{2.461232in}}%
\pgfpathlineto{\pgfqpoint{3.651705in}{2.458283in}}%
\pgfpathmoveto{\pgfqpoint{3.647164in}{2.461232in}}%
\pgfpathlineto{\pgfqpoint{3.647164in}{2.461232in}}%
\pgfpathlineto{\pgfqpoint{3.647164in}{2.464182in}}%
\pgfpathlineto{\pgfqpoint{3.651705in}{2.464182in}}%
\pgfpathlineto{\pgfqpoint{3.651705in}{2.461232in}}%
\pgfpathmoveto{\pgfqpoint{3.651705in}{2.458283in}}%
\pgfpathlineto{\pgfqpoint{3.651705in}{2.458283in}}%
\pgfpathlineto{\pgfqpoint{3.651705in}{2.461232in}}%
\pgfpathlineto{\pgfqpoint{3.656246in}{2.461232in}}%
\pgfpathlineto{\pgfqpoint{3.656246in}{2.458283in}}%
\pgfpathmoveto{\pgfqpoint{3.651705in}{2.461232in}}%
\pgfpathlineto{\pgfqpoint{3.651705in}{2.461232in}}%
\pgfpathlineto{\pgfqpoint{3.651705in}{2.464182in}}%
\pgfpathlineto{\pgfqpoint{3.656246in}{2.464182in}}%
\pgfpathlineto{\pgfqpoint{3.656246in}{2.461232in}}%
\pgfpathmoveto{\pgfqpoint{3.647164in}{2.464182in}}%
\pgfpathlineto{\pgfqpoint{3.647164in}{2.464182in}}%
\pgfpathlineto{\pgfqpoint{3.647164in}{2.467131in}}%
\pgfpathlineto{\pgfqpoint{3.651705in}{2.467131in}}%
\pgfpathlineto{\pgfqpoint{3.651705in}{2.464182in}}%
\pgfpathmoveto{\pgfqpoint{3.647164in}{2.467131in}}%
\pgfpathlineto{\pgfqpoint{3.647164in}{2.467131in}}%
\pgfpathlineto{\pgfqpoint{3.647164in}{2.470080in}}%
\pgfpathlineto{\pgfqpoint{3.651705in}{2.470080in}}%
\pgfpathlineto{\pgfqpoint{3.651705in}{2.467131in}}%
\pgfpathmoveto{\pgfqpoint{3.651705in}{2.464182in}}%
\pgfpathlineto{\pgfqpoint{3.651705in}{2.464182in}}%
\pgfpathlineto{\pgfqpoint{3.651705in}{2.467131in}}%
\pgfpathlineto{\pgfqpoint{3.656246in}{2.467131in}}%
\pgfpathlineto{\pgfqpoint{3.656246in}{2.464182in}}%
\pgfpathmoveto{\pgfqpoint{3.651705in}{2.467131in}}%
\pgfpathlineto{\pgfqpoint{3.651705in}{2.467131in}}%
\pgfpathlineto{\pgfqpoint{3.651705in}{2.470080in}}%
\pgfpathlineto{\pgfqpoint{3.656246in}{2.470080in}}%
\pgfpathlineto{\pgfqpoint{3.656246in}{2.467131in}}%
\pgfpathmoveto{\pgfqpoint{3.638082in}{2.470080in}}%
\pgfpathlineto{\pgfqpoint{3.638082in}{2.470080in}}%
\pgfpathlineto{\pgfqpoint{3.638082in}{2.473030in}}%
\pgfpathlineto{\pgfqpoint{3.642623in}{2.473030in}}%
\pgfpathlineto{\pgfqpoint{3.642623in}{2.470080in}}%
\pgfpathmoveto{\pgfqpoint{3.638082in}{2.473030in}}%
\pgfpathlineto{\pgfqpoint{3.638082in}{2.473030in}}%
\pgfpathlineto{\pgfqpoint{3.638082in}{2.475979in}}%
\pgfpathlineto{\pgfqpoint{3.642623in}{2.475979in}}%
\pgfpathlineto{\pgfqpoint{3.642623in}{2.473030in}}%
\pgfpathmoveto{\pgfqpoint{3.642623in}{2.470080in}}%
\pgfpathlineto{\pgfqpoint{3.642623in}{2.470080in}}%
\pgfpathlineto{\pgfqpoint{3.642623in}{2.473030in}}%
\pgfpathlineto{\pgfqpoint{3.647164in}{2.473030in}}%
\pgfpathlineto{\pgfqpoint{3.647164in}{2.470080in}}%
\pgfpathmoveto{\pgfqpoint{3.642623in}{2.473030in}}%
\pgfpathlineto{\pgfqpoint{3.642623in}{2.473030in}}%
\pgfpathlineto{\pgfqpoint{3.642623in}{2.475979in}}%
\pgfpathlineto{\pgfqpoint{3.647164in}{2.475979in}}%
\pgfpathlineto{\pgfqpoint{3.647164in}{2.473030in}}%
\pgfpathmoveto{\pgfqpoint{3.638082in}{2.475979in}}%
\pgfpathlineto{\pgfqpoint{3.638082in}{2.475979in}}%
\pgfpathlineto{\pgfqpoint{3.638082in}{2.478928in}}%
\pgfpathlineto{\pgfqpoint{3.642623in}{2.478928in}}%
\pgfpathlineto{\pgfqpoint{3.642623in}{2.475979in}}%
\pgfpathmoveto{\pgfqpoint{3.638082in}{2.478928in}}%
\pgfpathlineto{\pgfqpoint{3.638082in}{2.478928in}}%
\pgfpathlineto{\pgfqpoint{3.638082in}{2.481878in}}%
\pgfpathlineto{\pgfqpoint{3.642623in}{2.481878in}}%
\pgfpathlineto{\pgfqpoint{3.642623in}{2.478928in}}%
\pgfpathmoveto{\pgfqpoint{3.642623in}{2.475979in}}%
\pgfpathlineto{\pgfqpoint{3.642623in}{2.475979in}}%
\pgfpathlineto{\pgfqpoint{3.642623in}{2.478928in}}%
\pgfpathlineto{\pgfqpoint{3.647164in}{2.478928in}}%
\pgfpathlineto{\pgfqpoint{3.647164in}{2.475979in}}%
\pgfpathmoveto{\pgfqpoint{3.642623in}{2.478928in}}%
\pgfpathlineto{\pgfqpoint{3.642623in}{2.478928in}}%
\pgfpathlineto{\pgfqpoint{3.642623in}{2.481878in}}%
\pgfpathlineto{\pgfqpoint{3.647164in}{2.481878in}}%
\pgfpathlineto{\pgfqpoint{3.647164in}{2.478928in}}%
\pgfpathmoveto{\pgfqpoint{3.647164in}{2.470080in}}%
\pgfpathlineto{\pgfqpoint{3.647164in}{2.470080in}}%
\pgfpathlineto{\pgfqpoint{3.647164in}{2.473030in}}%
\pgfpathlineto{\pgfqpoint{3.651705in}{2.473030in}}%
\pgfpathlineto{\pgfqpoint{3.651705in}{2.470080in}}%
\pgfpathmoveto{\pgfqpoint{3.647164in}{2.473030in}}%
\pgfpathlineto{\pgfqpoint{3.647164in}{2.473030in}}%
\pgfpathlineto{\pgfqpoint{3.647164in}{2.475979in}}%
\pgfpathlineto{\pgfqpoint{3.651705in}{2.475979in}}%
\pgfpathlineto{\pgfqpoint{3.651705in}{2.473030in}}%
\pgfpathmoveto{\pgfqpoint{3.574510in}{2.552658in}}%
\pgfpathlineto{\pgfqpoint{3.574510in}{2.552658in}}%
\pgfpathlineto{\pgfqpoint{3.574510in}{2.555607in}}%
\pgfpathlineto{\pgfqpoint{3.579051in}{2.555607in}}%
\pgfpathlineto{\pgfqpoint{3.579051in}{2.552658in}}%
\pgfpathmoveto{\pgfqpoint{3.574510in}{2.555607in}}%
\pgfpathlineto{\pgfqpoint{3.574510in}{2.555607in}}%
\pgfpathlineto{\pgfqpoint{3.574510in}{2.558556in}}%
\pgfpathlineto{\pgfqpoint{3.579051in}{2.558556in}}%
\pgfpathlineto{\pgfqpoint{3.579051in}{2.555607in}}%
\pgfpathmoveto{\pgfqpoint{3.579051in}{2.552658in}}%
\pgfpathlineto{\pgfqpoint{3.579051in}{2.552658in}}%
\pgfpathlineto{\pgfqpoint{3.579051in}{2.555607in}}%
\pgfpathlineto{\pgfqpoint{3.583592in}{2.555607in}}%
\pgfpathlineto{\pgfqpoint{3.583592in}{2.552658in}}%
\pgfpathmoveto{\pgfqpoint{3.579051in}{2.555607in}}%
\pgfpathlineto{\pgfqpoint{3.579051in}{2.555607in}}%
\pgfpathlineto{\pgfqpoint{3.579051in}{2.558556in}}%
\pgfpathlineto{\pgfqpoint{3.583592in}{2.558556in}}%
\pgfpathlineto{\pgfqpoint{3.583592in}{2.555607in}}%
\pgfpathmoveto{\pgfqpoint{3.574510in}{2.558556in}}%
\pgfpathlineto{\pgfqpoint{3.574510in}{2.558556in}}%
\pgfpathlineto{\pgfqpoint{3.574510in}{2.561505in}}%
\pgfpathlineto{\pgfqpoint{3.579051in}{2.561505in}}%
\pgfpathlineto{\pgfqpoint{3.579051in}{2.558556in}}%
\pgfpathmoveto{\pgfqpoint{3.574510in}{2.561505in}}%
\pgfpathlineto{\pgfqpoint{3.574510in}{2.561505in}}%
\pgfpathlineto{\pgfqpoint{3.574510in}{2.564454in}}%
\pgfpathlineto{\pgfqpoint{3.579051in}{2.564454in}}%
\pgfpathlineto{\pgfqpoint{3.579051in}{2.561505in}}%
\pgfpathmoveto{\pgfqpoint{3.579051in}{2.558556in}}%
\pgfpathlineto{\pgfqpoint{3.579051in}{2.558556in}}%
\pgfpathlineto{\pgfqpoint{3.579051in}{2.561505in}}%
\pgfpathlineto{\pgfqpoint{3.583592in}{2.561505in}}%
\pgfpathlineto{\pgfqpoint{3.583592in}{2.558556in}}%
\pgfpathmoveto{\pgfqpoint{3.579051in}{2.561505in}}%
\pgfpathlineto{\pgfqpoint{3.579051in}{2.561505in}}%
\pgfpathlineto{\pgfqpoint{3.579051in}{2.564454in}}%
\pgfpathlineto{\pgfqpoint{3.583592in}{2.564454in}}%
\pgfpathlineto{\pgfqpoint{3.583592in}{2.561505in}}%
\pgfpathmoveto{\pgfqpoint{3.565429in}{2.564454in}}%
\pgfpathlineto{\pgfqpoint{3.565429in}{2.564454in}}%
\pgfpathlineto{\pgfqpoint{3.565429in}{2.567403in}}%
\pgfpathlineto{\pgfqpoint{3.569970in}{2.567403in}}%
\pgfpathlineto{\pgfqpoint{3.569970in}{2.564454in}}%
\pgfpathmoveto{\pgfqpoint{3.565429in}{2.567403in}}%
\pgfpathlineto{\pgfqpoint{3.565429in}{2.567403in}}%
\pgfpathlineto{\pgfqpoint{3.565429in}{2.570353in}}%
\pgfpathlineto{\pgfqpoint{3.569970in}{2.570353in}}%
\pgfpathlineto{\pgfqpoint{3.569970in}{2.567403in}}%
\pgfpathmoveto{\pgfqpoint{3.569970in}{2.564454in}}%
\pgfpathlineto{\pgfqpoint{3.569970in}{2.564454in}}%
\pgfpathlineto{\pgfqpoint{3.569970in}{2.567403in}}%
\pgfpathlineto{\pgfqpoint{3.574510in}{2.567403in}}%
\pgfpathlineto{\pgfqpoint{3.574510in}{2.564454in}}%
\pgfpathmoveto{\pgfqpoint{3.569970in}{2.567403in}}%
\pgfpathlineto{\pgfqpoint{3.569970in}{2.567403in}}%
\pgfpathlineto{\pgfqpoint{3.569970in}{2.570353in}}%
\pgfpathlineto{\pgfqpoint{3.574510in}{2.570353in}}%
\pgfpathlineto{\pgfqpoint{3.574510in}{2.567403in}}%
\pgfpathmoveto{\pgfqpoint{3.565429in}{2.570353in}}%
\pgfpathlineto{\pgfqpoint{3.565429in}{2.570353in}}%
\pgfpathlineto{\pgfqpoint{3.565429in}{2.573302in}}%
\pgfpathlineto{\pgfqpoint{3.569970in}{2.573302in}}%
\pgfpathlineto{\pgfqpoint{3.569970in}{2.570353in}}%
\pgfpathmoveto{\pgfqpoint{3.565429in}{2.573302in}}%
\pgfpathlineto{\pgfqpoint{3.565429in}{2.573302in}}%
\pgfpathlineto{\pgfqpoint{3.565429in}{2.576251in}}%
\pgfpathlineto{\pgfqpoint{3.569970in}{2.576251in}}%
\pgfpathlineto{\pgfqpoint{3.569970in}{2.573302in}}%
\pgfpathmoveto{\pgfqpoint{3.569970in}{2.570353in}}%
\pgfpathlineto{\pgfqpoint{3.569970in}{2.570353in}}%
\pgfpathlineto{\pgfqpoint{3.569970in}{2.573302in}}%
\pgfpathlineto{\pgfqpoint{3.574510in}{2.573302in}}%
\pgfpathlineto{\pgfqpoint{3.574510in}{2.570353in}}%
\pgfpathmoveto{\pgfqpoint{3.569970in}{2.573302in}}%
\pgfpathlineto{\pgfqpoint{3.569970in}{2.573302in}}%
\pgfpathlineto{\pgfqpoint{3.569970in}{2.576251in}}%
\pgfpathlineto{\pgfqpoint{3.574510in}{2.576251in}}%
\pgfpathlineto{\pgfqpoint{3.574510in}{2.573302in}}%
\pgfpathmoveto{\pgfqpoint{3.574510in}{2.564454in}}%
\pgfpathlineto{\pgfqpoint{3.574510in}{2.564454in}}%
\pgfpathlineto{\pgfqpoint{3.574510in}{2.567403in}}%
\pgfpathlineto{\pgfqpoint{3.579051in}{2.567403in}}%
\pgfpathlineto{\pgfqpoint{3.579051in}{2.564454in}}%
\pgfpathmoveto{\pgfqpoint{3.574510in}{2.567403in}}%
\pgfpathlineto{\pgfqpoint{3.574510in}{2.567403in}}%
\pgfpathlineto{\pgfqpoint{3.574510in}{2.570353in}}%
\pgfpathlineto{\pgfqpoint{3.579051in}{2.570353in}}%
\pgfpathlineto{\pgfqpoint{3.579051in}{2.567403in}}%
\pgfpathmoveto{\pgfqpoint{3.610837in}{2.505471in}}%
\pgfpathlineto{\pgfqpoint{3.610837in}{2.505471in}}%
\pgfpathlineto{\pgfqpoint{3.610837in}{2.508420in}}%
\pgfpathlineto{\pgfqpoint{3.615378in}{2.508420in}}%
\pgfpathlineto{\pgfqpoint{3.615378in}{2.505471in}}%
\pgfpathmoveto{\pgfqpoint{3.610837in}{2.508420in}}%
\pgfpathlineto{\pgfqpoint{3.610837in}{2.508420in}}%
\pgfpathlineto{\pgfqpoint{3.610837in}{2.511369in}}%
\pgfpathlineto{\pgfqpoint{3.615378in}{2.511369in}}%
\pgfpathlineto{\pgfqpoint{3.615378in}{2.508420in}}%
\pgfpathmoveto{\pgfqpoint{3.615378in}{2.505471in}}%
\pgfpathlineto{\pgfqpoint{3.615378in}{2.505471in}}%
\pgfpathlineto{\pgfqpoint{3.615378in}{2.508420in}}%
\pgfpathlineto{\pgfqpoint{3.619919in}{2.508420in}}%
\pgfpathlineto{\pgfqpoint{3.619919in}{2.505471in}}%
\pgfpathmoveto{\pgfqpoint{3.615378in}{2.508420in}}%
\pgfpathlineto{\pgfqpoint{3.615378in}{2.508420in}}%
\pgfpathlineto{\pgfqpoint{3.615378in}{2.511369in}}%
\pgfpathlineto{\pgfqpoint{3.619919in}{2.511369in}}%
\pgfpathlineto{\pgfqpoint{3.619919in}{2.508420in}}%
\pgfpathmoveto{\pgfqpoint{3.610837in}{2.511369in}}%
\pgfpathlineto{\pgfqpoint{3.610837in}{2.511369in}}%
\pgfpathlineto{\pgfqpoint{3.610837in}{2.514318in}}%
\pgfpathlineto{\pgfqpoint{3.615378in}{2.514318in}}%
\pgfpathlineto{\pgfqpoint{3.615378in}{2.511369in}}%
\pgfpathmoveto{\pgfqpoint{3.610837in}{2.514318in}}%
\pgfpathlineto{\pgfqpoint{3.610837in}{2.514318in}}%
\pgfpathlineto{\pgfqpoint{3.610837in}{2.517268in}}%
\pgfpathlineto{\pgfqpoint{3.615378in}{2.517268in}}%
\pgfpathlineto{\pgfqpoint{3.615378in}{2.514318in}}%
\pgfpathmoveto{\pgfqpoint{3.615378in}{2.511369in}}%
\pgfpathlineto{\pgfqpoint{3.615378in}{2.511369in}}%
\pgfpathlineto{\pgfqpoint{3.615378in}{2.514318in}}%
\pgfpathlineto{\pgfqpoint{3.619919in}{2.514318in}}%
\pgfpathlineto{\pgfqpoint{3.619919in}{2.511369in}}%
\pgfpathmoveto{\pgfqpoint{3.615378in}{2.514318in}}%
\pgfpathlineto{\pgfqpoint{3.615378in}{2.514318in}}%
\pgfpathlineto{\pgfqpoint{3.615378in}{2.517268in}}%
\pgfpathlineto{\pgfqpoint{3.619919in}{2.517268in}}%
\pgfpathlineto{\pgfqpoint{3.619919in}{2.514318in}}%
\pgfpathmoveto{\pgfqpoint{3.601755in}{2.517268in}}%
\pgfpathlineto{\pgfqpoint{3.601755in}{2.517268in}}%
\pgfpathlineto{\pgfqpoint{3.601755in}{2.520217in}}%
\pgfpathlineto{\pgfqpoint{3.606296in}{2.520217in}}%
\pgfpathlineto{\pgfqpoint{3.606296in}{2.517268in}}%
\pgfpathmoveto{\pgfqpoint{3.601755in}{2.520217in}}%
\pgfpathlineto{\pgfqpoint{3.601755in}{2.520217in}}%
\pgfpathlineto{\pgfqpoint{3.601755in}{2.523166in}}%
\pgfpathlineto{\pgfqpoint{3.606296in}{2.523166in}}%
\pgfpathlineto{\pgfqpoint{3.606296in}{2.520217in}}%
\pgfpathmoveto{\pgfqpoint{3.606296in}{2.517268in}}%
\pgfpathlineto{\pgfqpoint{3.606296in}{2.517268in}}%
\pgfpathlineto{\pgfqpoint{3.606296in}{2.520217in}}%
\pgfpathlineto{\pgfqpoint{3.610837in}{2.520217in}}%
\pgfpathlineto{\pgfqpoint{3.610837in}{2.517268in}}%
\pgfpathmoveto{\pgfqpoint{3.606296in}{2.520217in}}%
\pgfpathlineto{\pgfqpoint{3.606296in}{2.520217in}}%
\pgfpathlineto{\pgfqpoint{3.606296in}{2.523166in}}%
\pgfpathlineto{\pgfqpoint{3.610837in}{2.523166in}}%
\pgfpathlineto{\pgfqpoint{3.610837in}{2.520217in}}%
\pgfpathmoveto{\pgfqpoint{3.601755in}{2.523166in}}%
\pgfpathlineto{\pgfqpoint{3.601755in}{2.523166in}}%
\pgfpathlineto{\pgfqpoint{3.601755in}{2.526115in}}%
\pgfpathlineto{\pgfqpoint{3.606296in}{2.526115in}}%
\pgfpathlineto{\pgfqpoint{3.606296in}{2.523166in}}%
\pgfpathmoveto{\pgfqpoint{3.601755in}{2.526115in}}%
\pgfpathlineto{\pgfqpoint{3.601755in}{2.526115in}}%
\pgfpathlineto{\pgfqpoint{3.601755in}{2.529064in}}%
\pgfpathlineto{\pgfqpoint{3.606296in}{2.529064in}}%
\pgfpathlineto{\pgfqpoint{3.606296in}{2.526115in}}%
\pgfpathmoveto{\pgfqpoint{3.606296in}{2.523166in}}%
\pgfpathlineto{\pgfqpoint{3.606296in}{2.523166in}}%
\pgfpathlineto{\pgfqpoint{3.606296in}{2.526115in}}%
\pgfpathlineto{\pgfqpoint{3.610837in}{2.526115in}}%
\pgfpathlineto{\pgfqpoint{3.610837in}{2.523166in}}%
\pgfpathmoveto{\pgfqpoint{3.606296in}{2.526115in}}%
\pgfpathlineto{\pgfqpoint{3.606296in}{2.526115in}}%
\pgfpathlineto{\pgfqpoint{3.606296in}{2.529064in}}%
\pgfpathlineto{\pgfqpoint{3.610837in}{2.529064in}}%
\pgfpathlineto{\pgfqpoint{3.610837in}{2.526115in}}%
\pgfpathmoveto{\pgfqpoint{3.610837in}{2.517268in}}%
\pgfpathlineto{\pgfqpoint{3.610837in}{2.517268in}}%
\pgfpathlineto{\pgfqpoint{3.610837in}{2.520217in}}%
\pgfpathlineto{\pgfqpoint{3.615378in}{2.520217in}}%
\pgfpathlineto{\pgfqpoint{3.615378in}{2.517268in}}%
\pgfpathmoveto{\pgfqpoint{3.610837in}{2.520217in}}%
\pgfpathlineto{\pgfqpoint{3.610837in}{2.520217in}}%
\pgfpathlineto{\pgfqpoint{3.610837in}{2.523166in}}%
\pgfpathlineto{\pgfqpoint{3.615378in}{2.523166in}}%
\pgfpathlineto{\pgfqpoint{3.615378in}{2.520217in}}%
\pgfpathmoveto{\pgfqpoint{3.629001in}{2.481878in}}%
\pgfpathlineto{\pgfqpoint{3.629001in}{2.481878in}}%
\pgfpathlineto{\pgfqpoint{3.629001in}{2.484827in}}%
\pgfpathlineto{\pgfqpoint{3.633541in}{2.484827in}}%
\pgfpathlineto{\pgfqpoint{3.633541in}{2.481878in}}%
\pgfpathmoveto{\pgfqpoint{3.629001in}{2.484827in}}%
\pgfpathlineto{\pgfqpoint{3.629001in}{2.484827in}}%
\pgfpathlineto{\pgfqpoint{3.629001in}{2.487776in}}%
\pgfpathlineto{\pgfqpoint{3.633541in}{2.487776in}}%
\pgfpathlineto{\pgfqpoint{3.633541in}{2.484827in}}%
\pgfpathmoveto{\pgfqpoint{3.633541in}{2.481878in}}%
\pgfpathlineto{\pgfqpoint{3.633541in}{2.481878in}}%
\pgfpathlineto{\pgfqpoint{3.633541in}{2.484827in}}%
\pgfpathlineto{\pgfqpoint{3.638082in}{2.484827in}}%
\pgfpathlineto{\pgfqpoint{3.638082in}{2.481878in}}%
\pgfpathmoveto{\pgfqpoint{3.633541in}{2.484827in}}%
\pgfpathlineto{\pgfqpoint{3.633541in}{2.484827in}}%
\pgfpathlineto{\pgfqpoint{3.633541in}{2.487776in}}%
\pgfpathlineto{\pgfqpoint{3.638082in}{2.487776in}}%
\pgfpathlineto{\pgfqpoint{3.638082in}{2.484827in}}%
\pgfpathmoveto{\pgfqpoint{3.629001in}{2.487776in}}%
\pgfpathlineto{\pgfqpoint{3.629001in}{2.487776in}}%
\pgfpathlineto{\pgfqpoint{3.629001in}{2.490725in}}%
\pgfpathlineto{\pgfqpoint{3.633541in}{2.490725in}}%
\pgfpathlineto{\pgfqpoint{3.633541in}{2.487776in}}%
\pgfpathmoveto{\pgfqpoint{3.629001in}{2.490725in}}%
\pgfpathlineto{\pgfqpoint{3.629001in}{2.490725in}}%
\pgfpathlineto{\pgfqpoint{3.629001in}{2.493674in}}%
\pgfpathlineto{\pgfqpoint{3.633541in}{2.493674in}}%
\pgfpathlineto{\pgfqpoint{3.633541in}{2.490725in}}%
\pgfpathmoveto{\pgfqpoint{3.633541in}{2.487776in}}%
\pgfpathlineto{\pgfqpoint{3.633541in}{2.487776in}}%
\pgfpathlineto{\pgfqpoint{3.633541in}{2.490725in}}%
\pgfpathlineto{\pgfqpoint{3.638082in}{2.490725in}}%
\pgfpathlineto{\pgfqpoint{3.638082in}{2.487776in}}%
\pgfpathmoveto{\pgfqpoint{3.633541in}{2.490725in}}%
\pgfpathlineto{\pgfqpoint{3.633541in}{2.490725in}}%
\pgfpathlineto{\pgfqpoint{3.633541in}{2.493674in}}%
\pgfpathlineto{\pgfqpoint{3.638082in}{2.493674in}}%
\pgfpathlineto{\pgfqpoint{3.638082in}{2.490725in}}%
\pgfpathmoveto{\pgfqpoint{3.619919in}{2.493674in}}%
\pgfpathlineto{\pgfqpoint{3.619919in}{2.493674in}}%
\pgfpathlineto{\pgfqpoint{3.619919in}{2.496623in}}%
\pgfpathlineto{\pgfqpoint{3.624460in}{2.496623in}}%
\pgfpathlineto{\pgfqpoint{3.624460in}{2.493674in}}%
\pgfpathmoveto{\pgfqpoint{3.619919in}{2.496623in}}%
\pgfpathlineto{\pgfqpoint{3.619919in}{2.496623in}}%
\pgfpathlineto{\pgfqpoint{3.619919in}{2.499573in}}%
\pgfpathlineto{\pgfqpoint{3.624460in}{2.499573in}}%
\pgfpathlineto{\pgfqpoint{3.624460in}{2.496623in}}%
\pgfpathmoveto{\pgfqpoint{3.624460in}{2.493674in}}%
\pgfpathlineto{\pgfqpoint{3.624460in}{2.493674in}}%
\pgfpathlineto{\pgfqpoint{3.624460in}{2.496623in}}%
\pgfpathlineto{\pgfqpoint{3.629001in}{2.496623in}}%
\pgfpathlineto{\pgfqpoint{3.629001in}{2.493674in}}%
\pgfpathmoveto{\pgfqpoint{3.624460in}{2.496623in}}%
\pgfpathlineto{\pgfqpoint{3.624460in}{2.496623in}}%
\pgfpathlineto{\pgfqpoint{3.624460in}{2.499573in}}%
\pgfpathlineto{\pgfqpoint{3.629001in}{2.499573in}}%
\pgfpathlineto{\pgfqpoint{3.629001in}{2.496623in}}%
\pgfpathmoveto{\pgfqpoint{3.619919in}{2.499573in}}%
\pgfpathlineto{\pgfqpoint{3.619919in}{2.499573in}}%
\pgfpathlineto{\pgfqpoint{3.619919in}{2.502522in}}%
\pgfpathlineto{\pgfqpoint{3.624460in}{2.502522in}}%
\pgfpathlineto{\pgfqpoint{3.624460in}{2.499573in}}%
\pgfpathmoveto{\pgfqpoint{3.619919in}{2.502522in}}%
\pgfpathlineto{\pgfqpoint{3.619919in}{2.502522in}}%
\pgfpathlineto{\pgfqpoint{3.619919in}{2.505471in}}%
\pgfpathlineto{\pgfqpoint{3.624460in}{2.505471in}}%
\pgfpathlineto{\pgfqpoint{3.624460in}{2.502522in}}%
\pgfpathmoveto{\pgfqpoint{3.624460in}{2.499573in}}%
\pgfpathlineto{\pgfqpoint{3.624460in}{2.499573in}}%
\pgfpathlineto{\pgfqpoint{3.624460in}{2.502522in}}%
\pgfpathlineto{\pgfqpoint{3.629001in}{2.502522in}}%
\pgfpathlineto{\pgfqpoint{3.629001in}{2.499573in}}%
\pgfpathmoveto{\pgfqpoint{3.624460in}{2.502522in}}%
\pgfpathlineto{\pgfqpoint{3.624460in}{2.502522in}}%
\pgfpathlineto{\pgfqpoint{3.624460in}{2.505471in}}%
\pgfpathlineto{\pgfqpoint{3.629001in}{2.505471in}}%
\pgfpathlineto{\pgfqpoint{3.629001in}{2.502522in}}%
\pgfpathmoveto{\pgfqpoint{3.629001in}{2.493674in}}%
\pgfpathlineto{\pgfqpoint{3.629001in}{2.493674in}}%
\pgfpathlineto{\pgfqpoint{3.629001in}{2.496623in}}%
\pgfpathlineto{\pgfqpoint{3.633541in}{2.496623in}}%
\pgfpathlineto{\pgfqpoint{3.633541in}{2.493674in}}%
\pgfpathmoveto{\pgfqpoint{3.629001in}{2.496623in}}%
\pgfpathlineto{\pgfqpoint{3.629001in}{2.496623in}}%
\pgfpathlineto{\pgfqpoint{3.629001in}{2.499573in}}%
\pgfpathlineto{\pgfqpoint{3.633541in}{2.499573in}}%
\pgfpathlineto{\pgfqpoint{3.633541in}{2.496623in}}%
\pgfpathmoveto{\pgfqpoint{3.638082in}{2.481878in}}%
\pgfpathlineto{\pgfqpoint{3.638082in}{2.481878in}}%
\pgfpathlineto{\pgfqpoint{3.638082in}{2.484827in}}%
\pgfpathlineto{\pgfqpoint{3.642623in}{2.484827in}}%
\pgfpathlineto{\pgfqpoint{3.642623in}{2.481878in}}%
\pgfpathmoveto{\pgfqpoint{3.638082in}{2.484827in}}%
\pgfpathlineto{\pgfqpoint{3.638082in}{2.484827in}}%
\pgfpathlineto{\pgfqpoint{3.638082in}{2.487776in}}%
\pgfpathlineto{\pgfqpoint{3.642623in}{2.487776in}}%
\pgfpathlineto{\pgfqpoint{3.642623in}{2.484827in}}%
\pgfpathmoveto{\pgfqpoint{3.619919in}{2.505471in}}%
\pgfpathlineto{\pgfqpoint{3.619919in}{2.505471in}}%
\pgfpathlineto{\pgfqpoint{3.619919in}{2.508420in}}%
\pgfpathlineto{\pgfqpoint{3.624460in}{2.508420in}}%
\pgfpathlineto{\pgfqpoint{3.624460in}{2.505471in}}%
\pgfpathmoveto{\pgfqpoint{3.619919in}{2.508420in}}%
\pgfpathlineto{\pgfqpoint{3.619919in}{2.508420in}}%
\pgfpathlineto{\pgfqpoint{3.619919in}{2.511369in}}%
\pgfpathlineto{\pgfqpoint{3.624460in}{2.511369in}}%
\pgfpathlineto{\pgfqpoint{3.624460in}{2.508420in}}%
\pgfpathmoveto{\pgfqpoint{3.592674in}{2.529064in}}%
\pgfpathlineto{\pgfqpoint{3.592674in}{2.529064in}}%
\pgfpathlineto{\pgfqpoint{3.592674in}{2.532013in}}%
\pgfpathlineto{\pgfqpoint{3.597215in}{2.532013in}}%
\pgfpathlineto{\pgfqpoint{3.597215in}{2.529064in}}%
\pgfpathmoveto{\pgfqpoint{3.592674in}{2.532013in}}%
\pgfpathlineto{\pgfqpoint{3.592674in}{2.532013in}}%
\pgfpathlineto{\pgfqpoint{3.592674in}{2.534963in}}%
\pgfpathlineto{\pgfqpoint{3.597215in}{2.534963in}}%
\pgfpathlineto{\pgfqpoint{3.597215in}{2.532013in}}%
\pgfpathmoveto{\pgfqpoint{3.597215in}{2.529064in}}%
\pgfpathlineto{\pgfqpoint{3.597215in}{2.529064in}}%
\pgfpathlineto{\pgfqpoint{3.597215in}{2.532013in}}%
\pgfpathlineto{\pgfqpoint{3.601755in}{2.532013in}}%
\pgfpathlineto{\pgfqpoint{3.601755in}{2.529064in}}%
\pgfpathmoveto{\pgfqpoint{3.597215in}{2.532013in}}%
\pgfpathlineto{\pgfqpoint{3.597215in}{2.532013in}}%
\pgfpathlineto{\pgfqpoint{3.597215in}{2.534963in}}%
\pgfpathlineto{\pgfqpoint{3.601755in}{2.534963in}}%
\pgfpathlineto{\pgfqpoint{3.601755in}{2.532013in}}%
\pgfpathmoveto{\pgfqpoint{3.592674in}{2.534963in}}%
\pgfpathlineto{\pgfqpoint{3.592674in}{2.534963in}}%
\pgfpathlineto{\pgfqpoint{3.592674in}{2.537912in}}%
\pgfpathlineto{\pgfqpoint{3.597215in}{2.537912in}}%
\pgfpathlineto{\pgfqpoint{3.597215in}{2.534963in}}%
\pgfpathmoveto{\pgfqpoint{3.592674in}{2.537912in}}%
\pgfpathlineto{\pgfqpoint{3.592674in}{2.537912in}}%
\pgfpathlineto{\pgfqpoint{3.592674in}{2.540861in}}%
\pgfpathlineto{\pgfqpoint{3.597215in}{2.540861in}}%
\pgfpathlineto{\pgfqpoint{3.597215in}{2.537912in}}%
\pgfpathmoveto{\pgfqpoint{3.597215in}{2.534963in}}%
\pgfpathlineto{\pgfqpoint{3.597215in}{2.534963in}}%
\pgfpathlineto{\pgfqpoint{3.597215in}{2.537912in}}%
\pgfpathlineto{\pgfqpoint{3.601755in}{2.537912in}}%
\pgfpathlineto{\pgfqpoint{3.601755in}{2.534963in}}%
\pgfpathmoveto{\pgfqpoint{3.597215in}{2.537912in}}%
\pgfpathlineto{\pgfqpoint{3.597215in}{2.537912in}}%
\pgfpathlineto{\pgfqpoint{3.597215in}{2.540861in}}%
\pgfpathlineto{\pgfqpoint{3.601755in}{2.540861in}}%
\pgfpathlineto{\pgfqpoint{3.601755in}{2.537912in}}%
\pgfpathmoveto{\pgfqpoint{3.583592in}{2.540861in}}%
\pgfpathlineto{\pgfqpoint{3.583592in}{2.540861in}}%
\pgfpathlineto{\pgfqpoint{3.583592in}{2.543810in}}%
\pgfpathlineto{\pgfqpoint{3.588133in}{2.543810in}}%
\pgfpathlineto{\pgfqpoint{3.588133in}{2.540861in}}%
\pgfpathmoveto{\pgfqpoint{3.583592in}{2.543810in}}%
\pgfpathlineto{\pgfqpoint{3.583592in}{2.543810in}}%
\pgfpathlineto{\pgfqpoint{3.583592in}{2.546759in}}%
\pgfpathlineto{\pgfqpoint{3.588133in}{2.546759in}}%
\pgfpathlineto{\pgfqpoint{3.588133in}{2.543810in}}%
\pgfpathmoveto{\pgfqpoint{3.588133in}{2.540861in}}%
\pgfpathlineto{\pgfqpoint{3.588133in}{2.540861in}}%
\pgfpathlineto{\pgfqpoint{3.588133in}{2.543810in}}%
\pgfpathlineto{\pgfqpoint{3.592674in}{2.543810in}}%
\pgfpathlineto{\pgfqpoint{3.592674in}{2.540861in}}%
\pgfpathmoveto{\pgfqpoint{3.588133in}{2.543810in}}%
\pgfpathlineto{\pgfqpoint{3.588133in}{2.543810in}}%
\pgfpathlineto{\pgfqpoint{3.588133in}{2.546759in}}%
\pgfpathlineto{\pgfqpoint{3.592674in}{2.546759in}}%
\pgfpathlineto{\pgfqpoint{3.592674in}{2.543810in}}%
\pgfpathmoveto{\pgfqpoint{3.583592in}{2.546759in}}%
\pgfpathlineto{\pgfqpoint{3.583592in}{2.546759in}}%
\pgfpathlineto{\pgfqpoint{3.583592in}{2.549708in}}%
\pgfpathlineto{\pgfqpoint{3.588133in}{2.549708in}}%
\pgfpathlineto{\pgfqpoint{3.588133in}{2.546759in}}%
\pgfpathmoveto{\pgfqpoint{3.583592in}{2.549708in}}%
\pgfpathlineto{\pgfqpoint{3.583592in}{2.549708in}}%
\pgfpathlineto{\pgfqpoint{3.583592in}{2.552658in}}%
\pgfpathlineto{\pgfqpoint{3.588133in}{2.552658in}}%
\pgfpathlineto{\pgfqpoint{3.588133in}{2.549708in}}%
\pgfpathmoveto{\pgfqpoint{3.588133in}{2.546759in}}%
\pgfpathlineto{\pgfqpoint{3.588133in}{2.546759in}}%
\pgfpathlineto{\pgfqpoint{3.588133in}{2.549708in}}%
\pgfpathlineto{\pgfqpoint{3.592674in}{2.549708in}}%
\pgfpathlineto{\pgfqpoint{3.592674in}{2.546759in}}%
\pgfpathmoveto{\pgfqpoint{3.588133in}{2.549708in}}%
\pgfpathlineto{\pgfqpoint{3.588133in}{2.549708in}}%
\pgfpathlineto{\pgfqpoint{3.588133in}{2.552658in}}%
\pgfpathlineto{\pgfqpoint{3.592674in}{2.552658in}}%
\pgfpathlineto{\pgfqpoint{3.592674in}{2.549708in}}%
\pgfpathmoveto{\pgfqpoint{3.592674in}{2.540861in}}%
\pgfpathlineto{\pgfqpoint{3.592674in}{2.540861in}}%
\pgfpathlineto{\pgfqpoint{3.592674in}{2.543810in}}%
\pgfpathlineto{\pgfqpoint{3.597215in}{2.543810in}}%
\pgfpathlineto{\pgfqpoint{3.597215in}{2.540861in}}%
\pgfpathmoveto{\pgfqpoint{3.592674in}{2.543810in}}%
\pgfpathlineto{\pgfqpoint{3.592674in}{2.543810in}}%
\pgfpathlineto{\pgfqpoint{3.592674in}{2.546759in}}%
\pgfpathlineto{\pgfqpoint{3.597215in}{2.546759in}}%
\pgfpathlineto{\pgfqpoint{3.597215in}{2.543810in}}%
\pgfpathmoveto{\pgfqpoint{3.601755in}{2.529064in}}%
\pgfpathlineto{\pgfqpoint{3.601755in}{2.529064in}}%
\pgfpathlineto{\pgfqpoint{3.601755in}{2.532013in}}%
\pgfpathlineto{\pgfqpoint{3.606296in}{2.532013in}}%
\pgfpathlineto{\pgfqpoint{3.606296in}{2.529064in}}%
\pgfpathmoveto{\pgfqpoint{3.601755in}{2.532013in}}%
\pgfpathlineto{\pgfqpoint{3.601755in}{2.532013in}}%
\pgfpathlineto{\pgfqpoint{3.601755in}{2.534963in}}%
\pgfpathlineto{\pgfqpoint{3.606296in}{2.534963in}}%
\pgfpathlineto{\pgfqpoint{3.606296in}{2.532013in}}%
\pgfpathmoveto{\pgfqpoint{3.583592in}{2.552658in}}%
\pgfpathlineto{\pgfqpoint{3.583592in}{2.552658in}}%
\pgfpathlineto{\pgfqpoint{3.583592in}{2.555607in}}%
\pgfpathlineto{\pgfqpoint{3.588133in}{2.555607in}}%
\pgfpathlineto{\pgfqpoint{3.588133in}{2.552658in}}%
\pgfpathmoveto{\pgfqpoint{3.583592in}{2.555607in}}%
\pgfpathlineto{\pgfqpoint{3.583592in}{2.555607in}}%
\pgfpathlineto{\pgfqpoint{3.583592in}{2.558556in}}%
\pgfpathlineto{\pgfqpoint{3.588133in}{2.558556in}}%
\pgfpathlineto{\pgfqpoint{3.588133in}{2.555607in}}%
\pgfpathmoveto{\pgfqpoint{3.538184in}{2.599844in}}%
\pgfpathlineto{\pgfqpoint{3.538184in}{2.599844in}}%
\pgfpathlineto{\pgfqpoint{3.538184in}{2.602793in}}%
\pgfpathlineto{\pgfqpoint{3.542725in}{2.602793in}}%
\pgfpathlineto{\pgfqpoint{3.542725in}{2.599844in}}%
\pgfpathmoveto{\pgfqpoint{3.538184in}{2.602793in}}%
\pgfpathlineto{\pgfqpoint{3.538184in}{2.602793in}}%
\pgfpathlineto{\pgfqpoint{3.538184in}{2.605742in}}%
\pgfpathlineto{\pgfqpoint{3.542725in}{2.605742in}}%
\pgfpathlineto{\pgfqpoint{3.542725in}{2.602793in}}%
\pgfpathmoveto{\pgfqpoint{3.542725in}{2.599844in}}%
\pgfpathlineto{\pgfqpoint{3.542725in}{2.599844in}}%
\pgfpathlineto{\pgfqpoint{3.542725in}{2.602793in}}%
\pgfpathlineto{\pgfqpoint{3.547265in}{2.602793in}}%
\pgfpathlineto{\pgfqpoint{3.547265in}{2.599844in}}%
\pgfpathmoveto{\pgfqpoint{3.542725in}{2.602793in}}%
\pgfpathlineto{\pgfqpoint{3.542725in}{2.602793in}}%
\pgfpathlineto{\pgfqpoint{3.542725in}{2.605742in}}%
\pgfpathlineto{\pgfqpoint{3.547265in}{2.605742in}}%
\pgfpathlineto{\pgfqpoint{3.547265in}{2.602793in}}%
\pgfpathmoveto{\pgfqpoint{3.538184in}{2.605742in}}%
\pgfpathlineto{\pgfqpoint{3.538184in}{2.605742in}}%
\pgfpathlineto{\pgfqpoint{3.538184in}{2.608692in}}%
\pgfpathlineto{\pgfqpoint{3.542725in}{2.608692in}}%
\pgfpathlineto{\pgfqpoint{3.542725in}{2.605742in}}%
\pgfpathmoveto{\pgfqpoint{3.538184in}{2.608692in}}%
\pgfpathlineto{\pgfqpoint{3.538184in}{2.608692in}}%
\pgfpathlineto{\pgfqpoint{3.538184in}{2.611641in}}%
\pgfpathlineto{\pgfqpoint{3.542725in}{2.611641in}}%
\pgfpathlineto{\pgfqpoint{3.542725in}{2.608692in}}%
\pgfpathmoveto{\pgfqpoint{3.542725in}{2.605742in}}%
\pgfpathlineto{\pgfqpoint{3.542725in}{2.605742in}}%
\pgfpathlineto{\pgfqpoint{3.542725in}{2.608692in}}%
\pgfpathlineto{\pgfqpoint{3.547265in}{2.608692in}}%
\pgfpathlineto{\pgfqpoint{3.547265in}{2.605742in}}%
\pgfpathmoveto{\pgfqpoint{3.542725in}{2.608692in}}%
\pgfpathlineto{\pgfqpoint{3.542725in}{2.608692in}}%
\pgfpathlineto{\pgfqpoint{3.542725in}{2.611641in}}%
\pgfpathlineto{\pgfqpoint{3.547265in}{2.611641in}}%
\pgfpathlineto{\pgfqpoint{3.547265in}{2.608692in}}%
\pgfpathmoveto{\pgfqpoint{3.529102in}{2.611641in}}%
\pgfpathlineto{\pgfqpoint{3.529102in}{2.611641in}}%
\pgfpathlineto{\pgfqpoint{3.529102in}{2.614590in}}%
\pgfpathlineto{\pgfqpoint{3.533643in}{2.614590in}}%
\pgfpathlineto{\pgfqpoint{3.533643in}{2.611641in}}%
\pgfpathmoveto{\pgfqpoint{3.529102in}{2.614590in}}%
\pgfpathlineto{\pgfqpoint{3.529102in}{2.614590in}}%
\pgfpathlineto{\pgfqpoint{3.529102in}{2.617539in}}%
\pgfpathlineto{\pgfqpoint{3.533643in}{2.617539in}}%
\pgfpathlineto{\pgfqpoint{3.533643in}{2.614590in}}%
\pgfpathmoveto{\pgfqpoint{3.533643in}{2.611641in}}%
\pgfpathlineto{\pgfqpoint{3.533643in}{2.611641in}}%
\pgfpathlineto{\pgfqpoint{3.533643in}{2.614590in}}%
\pgfpathlineto{\pgfqpoint{3.538184in}{2.614590in}}%
\pgfpathlineto{\pgfqpoint{3.538184in}{2.611641in}}%
\pgfpathmoveto{\pgfqpoint{3.533643in}{2.614590in}}%
\pgfpathlineto{\pgfqpoint{3.533643in}{2.614590in}}%
\pgfpathlineto{\pgfqpoint{3.533643in}{2.617539in}}%
\pgfpathlineto{\pgfqpoint{3.538184in}{2.617539in}}%
\pgfpathlineto{\pgfqpoint{3.538184in}{2.614590in}}%
\pgfpathmoveto{\pgfqpoint{3.529102in}{2.617539in}}%
\pgfpathlineto{\pgfqpoint{3.529102in}{2.617539in}}%
\pgfpathlineto{\pgfqpoint{3.529102in}{2.620488in}}%
\pgfpathlineto{\pgfqpoint{3.533643in}{2.620488in}}%
\pgfpathlineto{\pgfqpoint{3.533643in}{2.617539in}}%
\pgfpathmoveto{\pgfqpoint{3.529102in}{2.620488in}}%
\pgfpathlineto{\pgfqpoint{3.529102in}{2.620488in}}%
\pgfpathlineto{\pgfqpoint{3.529102in}{2.623437in}}%
\pgfpathlineto{\pgfqpoint{3.533643in}{2.623437in}}%
\pgfpathlineto{\pgfqpoint{3.533643in}{2.620488in}}%
\pgfpathmoveto{\pgfqpoint{3.533643in}{2.617539in}}%
\pgfpathlineto{\pgfqpoint{3.533643in}{2.617539in}}%
\pgfpathlineto{\pgfqpoint{3.533643in}{2.620488in}}%
\pgfpathlineto{\pgfqpoint{3.538184in}{2.620488in}}%
\pgfpathlineto{\pgfqpoint{3.538184in}{2.617539in}}%
\pgfpathmoveto{\pgfqpoint{3.533643in}{2.620488in}}%
\pgfpathlineto{\pgfqpoint{3.533643in}{2.620488in}}%
\pgfpathlineto{\pgfqpoint{3.533643in}{2.623437in}}%
\pgfpathlineto{\pgfqpoint{3.538184in}{2.623437in}}%
\pgfpathlineto{\pgfqpoint{3.538184in}{2.620488in}}%
\pgfpathmoveto{\pgfqpoint{3.538184in}{2.611641in}}%
\pgfpathlineto{\pgfqpoint{3.538184in}{2.611641in}}%
\pgfpathlineto{\pgfqpoint{3.538184in}{2.614590in}}%
\pgfpathlineto{\pgfqpoint{3.542725in}{2.614590in}}%
\pgfpathlineto{\pgfqpoint{3.542725in}{2.611641in}}%
\pgfpathmoveto{\pgfqpoint{3.538184in}{2.614590in}}%
\pgfpathlineto{\pgfqpoint{3.538184in}{2.614590in}}%
\pgfpathlineto{\pgfqpoint{3.538184in}{2.617539in}}%
\pgfpathlineto{\pgfqpoint{3.542725in}{2.617539in}}%
\pgfpathlineto{\pgfqpoint{3.542725in}{2.614590in}}%
\pgfpathmoveto{\pgfqpoint{3.556347in}{2.576251in}}%
\pgfpathlineto{\pgfqpoint{3.556347in}{2.576251in}}%
\pgfpathlineto{\pgfqpoint{3.556347in}{2.579200in}}%
\pgfpathlineto{\pgfqpoint{3.560888in}{2.579200in}}%
\pgfpathlineto{\pgfqpoint{3.560888in}{2.576251in}}%
\pgfpathmoveto{\pgfqpoint{3.556347in}{2.579200in}}%
\pgfpathlineto{\pgfqpoint{3.556347in}{2.579200in}}%
\pgfpathlineto{\pgfqpoint{3.556347in}{2.582149in}}%
\pgfpathlineto{\pgfqpoint{3.560888in}{2.582149in}}%
\pgfpathlineto{\pgfqpoint{3.560888in}{2.579200in}}%
\pgfpathmoveto{\pgfqpoint{3.560888in}{2.576251in}}%
\pgfpathlineto{\pgfqpoint{3.560888in}{2.576251in}}%
\pgfpathlineto{\pgfqpoint{3.560888in}{2.579200in}}%
\pgfpathlineto{\pgfqpoint{3.565429in}{2.579200in}}%
\pgfpathlineto{\pgfqpoint{3.565429in}{2.576251in}}%
\pgfpathmoveto{\pgfqpoint{3.560888in}{2.579200in}}%
\pgfpathlineto{\pgfqpoint{3.560888in}{2.579200in}}%
\pgfpathlineto{\pgfqpoint{3.560888in}{2.582149in}}%
\pgfpathlineto{\pgfqpoint{3.565429in}{2.582149in}}%
\pgfpathlineto{\pgfqpoint{3.565429in}{2.579200in}}%
\pgfpathmoveto{\pgfqpoint{3.556347in}{2.582149in}}%
\pgfpathlineto{\pgfqpoint{3.556347in}{2.582149in}}%
\pgfpathlineto{\pgfqpoint{3.556347in}{2.585098in}}%
\pgfpathlineto{\pgfqpoint{3.560888in}{2.585098in}}%
\pgfpathlineto{\pgfqpoint{3.560888in}{2.582149in}}%
\pgfpathmoveto{\pgfqpoint{3.556347in}{2.585098in}}%
\pgfpathlineto{\pgfqpoint{3.556347in}{2.585098in}}%
\pgfpathlineto{\pgfqpoint{3.556347in}{2.588047in}}%
\pgfpathlineto{\pgfqpoint{3.560888in}{2.588047in}}%
\pgfpathlineto{\pgfqpoint{3.560888in}{2.585098in}}%
\pgfpathmoveto{\pgfqpoint{3.560888in}{2.582149in}}%
\pgfpathlineto{\pgfqpoint{3.560888in}{2.582149in}}%
\pgfpathlineto{\pgfqpoint{3.560888in}{2.585098in}}%
\pgfpathlineto{\pgfqpoint{3.565429in}{2.585098in}}%
\pgfpathlineto{\pgfqpoint{3.565429in}{2.582149in}}%
\pgfpathmoveto{\pgfqpoint{3.560888in}{2.585098in}}%
\pgfpathlineto{\pgfqpoint{3.560888in}{2.585098in}}%
\pgfpathlineto{\pgfqpoint{3.560888in}{2.588047in}}%
\pgfpathlineto{\pgfqpoint{3.565429in}{2.588047in}}%
\pgfpathlineto{\pgfqpoint{3.565429in}{2.585098in}}%
\pgfpathmoveto{\pgfqpoint{3.547265in}{2.588047in}}%
\pgfpathlineto{\pgfqpoint{3.547265in}{2.588047in}}%
\pgfpathlineto{\pgfqpoint{3.547265in}{2.590997in}}%
\pgfpathlineto{\pgfqpoint{3.551806in}{2.590997in}}%
\pgfpathlineto{\pgfqpoint{3.551806in}{2.588047in}}%
\pgfpathmoveto{\pgfqpoint{3.547265in}{2.590997in}}%
\pgfpathlineto{\pgfqpoint{3.547265in}{2.590997in}}%
\pgfpathlineto{\pgfqpoint{3.547265in}{2.593946in}}%
\pgfpathlineto{\pgfqpoint{3.551806in}{2.593946in}}%
\pgfpathlineto{\pgfqpoint{3.551806in}{2.590997in}}%
\pgfpathmoveto{\pgfqpoint{3.551806in}{2.588047in}}%
\pgfpathlineto{\pgfqpoint{3.551806in}{2.588047in}}%
\pgfpathlineto{\pgfqpoint{3.551806in}{2.590997in}}%
\pgfpathlineto{\pgfqpoint{3.556347in}{2.590997in}}%
\pgfpathlineto{\pgfqpoint{3.556347in}{2.588047in}}%
\pgfpathmoveto{\pgfqpoint{3.551806in}{2.590997in}}%
\pgfpathlineto{\pgfqpoint{3.551806in}{2.590997in}}%
\pgfpathlineto{\pgfqpoint{3.551806in}{2.593946in}}%
\pgfpathlineto{\pgfqpoint{3.556347in}{2.593946in}}%
\pgfpathlineto{\pgfqpoint{3.556347in}{2.590997in}}%
\pgfpathmoveto{\pgfqpoint{3.547265in}{2.593946in}}%
\pgfpathlineto{\pgfqpoint{3.547265in}{2.593946in}}%
\pgfpathlineto{\pgfqpoint{3.547265in}{2.596895in}}%
\pgfpathlineto{\pgfqpoint{3.551806in}{2.596895in}}%
\pgfpathlineto{\pgfqpoint{3.551806in}{2.593946in}}%
\pgfpathmoveto{\pgfqpoint{3.547265in}{2.596895in}}%
\pgfpathlineto{\pgfqpoint{3.547265in}{2.596895in}}%
\pgfpathlineto{\pgfqpoint{3.547265in}{2.599844in}}%
\pgfpathlineto{\pgfqpoint{3.551806in}{2.599844in}}%
\pgfpathlineto{\pgfqpoint{3.551806in}{2.596895in}}%
\pgfpathmoveto{\pgfqpoint{3.551806in}{2.593946in}}%
\pgfpathlineto{\pgfqpoint{3.551806in}{2.593946in}}%
\pgfpathlineto{\pgfqpoint{3.551806in}{2.596895in}}%
\pgfpathlineto{\pgfqpoint{3.556347in}{2.596895in}}%
\pgfpathlineto{\pgfqpoint{3.556347in}{2.593946in}}%
\pgfpathmoveto{\pgfqpoint{3.551806in}{2.596895in}}%
\pgfpathlineto{\pgfqpoint{3.551806in}{2.596895in}}%
\pgfpathlineto{\pgfqpoint{3.551806in}{2.599844in}}%
\pgfpathlineto{\pgfqpoint{3.556347in}{2.599844in}}%
\pgfpathlineto{\pgfqpoint{3.556347in}{2.596895in}}%
\pgfpathmoveto{\pgfqpoint{3.556347in}{2.588047in}}%
\pgfpathlineto{\pgfqpoint{3.556347in}{2.588047in}}%
\pgfpathlineto{\pgfqpoint{3.556347in}{2.590997in}}%
\pgfpathlineto{\pgfqpoint{3.560888in}{2.590997in}}%
\pgfpathlineto{\pgfqpoint{3.560888in}{2.588047in}}%
\pgfpathmoveto{\pgfqpoint{3.556347in}{2.590997in}}%
\pgfpathlineto{\pgfqpoint{3.556347in}{2.590997in}}%
\pgfpathlineto{\pgfqpoint{3.556347in}{2.593946in}}%
\pgfpathlineto{\pgfqpoint{3.560888in}{2.593946in}}%
\pgfpathlineto{\pgfqpoint{3.560888in}{2.590997in}}%
\pgfpathmoveto{\pgfqpoint{3.565429in}{2.576251in}}%
\pgfpathlineto{\pgfqpoint{3.565429in}{2.576251in}}%
\pgfpathlineto{\pgfqpoint{3.565429in}{2.579200in}}%
\pgfpathlineto{\pgfqpoint{3.569970in}{2.579200in}}%
\pgfpathlineto{\pgfqpoint{3.569970in}{2.576251in}}%
\pgfpathmoveto{\pgfqpoint{3.565429in}{2.579200in}}%
\pgfpathlineto{\pgfqpoint{3.565429in}{2.579200in}}%
\pgfpathlineto{\pgfqpoint{3.565429in}{2.582149in}}%
\pgfpathlineto{\pgfqpoint{3.569970in}{2.582149in}}%
\pgfpathlineto{\pgfqpoint{3.569970in}{2.579200in}}%
\pgfpathmoveto{\pgfqpoint{3.547265in}{2.599844in}}%
\pgfpathlineto{\pgfqpoint{3.547265in}{2.599844in}}%
\pgfpathlineto{\pgfqpoint{3.547265in}{2.602793in}}%
\pgfpathlineto{\pgfqpoint{3.551806in}{2.602793in}}%
\pgfpathlineto{\pgfqpoint{3.551806in}{2.599844in}}%
\pgfpathmoveto{\pgfqpoint{3.547265in}{2.602793in}}%
\pgfpathlineto{\pgfqpoint{3.547265in}{2.602793in}}%
\pgfpathlineto{\pgfqpoint{3.547265in}{2.605742in}}%
\pgfpathlineto{\pgfqpoint{3.551806in}{2.605742in}}%
\pgfpathlineto{\pgfqpoint{3.551806in}{2.602793in}}%
\pgfpathmoveto{\pgfqpoint{3.520020in}{2.623437in}}%
\pgfpathlineto{\pgfqpoint{3.520020in}{2.623437in}}%
\pgfpathlineto{\pgfqpoint{3.520020in}{2.626387in}}%
\pgfpathlineto{\pgfqpoint{3.524561in}{2.626387in}}%
\pgfpathlineto{\pgfqpoint{3.524561in}{2.623437in}}%
\pgfpathmoveto{\pgfqpoint{3.520020in}{2.626387in}}%
\pgfpathlineto{\pgfqpoint{3.520020in}{2.626387in}}%
\pgfpathlineto{\pgfqpoint{3.520020in}{2.629336in}}%
\pgfpathlineto{\pgfqpoint{3.524561in}{2.629336in}}%
\pgfpathlineto{\pgfqpoint{3.524561in}{2.626387in}}%
\pgfpathmoveto{\pgfqpoint{3.524561in}{2.623437in}}%
\pgfpathlineto{\pgfqpoint{3.524561in}{2.623437in}}%
\pgfpathlineto{\pgfqpoint{3.524561in}{2.626387in}}%
\pgfpathlineto{\pgfqpoint{3.529102in}{2.626387in}}%
\pgfpathlineto{\pgfqpoint{3.529102in}{2.623437in}}%
\pgfpathmoveto{\pgfqpoint{3.524561in}{2.626387in}}%
\pgfpathlineto{\pgfqpoint{3.524561in}{2.626387in}}%
\pgfpathlineto{\pgfqpoint{3.524561in}{2.629336in}}%
\pgfpathlineto{\pgfqpoint{3.529102in}{2.629336in}}%
\pgfpathlineto{\pgfqpoint{3.529102in}{2.626387in}}%
\pgfpathmoveto{\pgfqpoint{3.520020in}{2.629336in}}%
\pgfpathlineto{\pgfqpoint{3.520020in}{2.629336in}}%
\pgfpathlineto{\pgfqpoint{3.520020in}{2.632285in}}%
\pgfpathlineto{\pgfqpoint{3.524561in}{2.632285in}}%
\pgfpathlineto{\pgfqpoint{3.524561in}{2.629336in}}%
\pgfpathmoveto{\pgfqpoint{3.520020in}{2.632285in}}%
\pgfpathlineto{\pgfqpoint{3.520020in}{2.632285in}}%
\pgfpathlineto{\pgfqpoint{3.520020in}{2.635234in}}%
\pgfpathlineto{\pgfqpoint{3.524561in}{2.635234in}}%
\pgfpathlineto{\pgfqpoint{3.524561in}{2.632285in}}%
\pgfpathmoveto{\pgfqpoint{3.524561in}{2.629336in}}%
\pgfpathlineto{\pgfqpoint{3.524561in}{2.629336in}}%
\pgfpathlineto{\pgfqpoint{3.524561in}{2.632285in}}%
\pgfpathlineto{\pgfqpoint{3.529102in}{2.632285in}}%
\pgfpathlineto{\pgfqpoint{3.529102in}{2.629336in}}%
\pgfpathmoveto{\pgfqpoint{3.524561in}{2.632285in}}%
\pgfpathlineto{\pgfqpoint{3.524561in}{2.632285in}}%
\pgfpathlineto{\pgfqpoint{3.524561in}{2.635234in}}%
\pgfpathlineto{\pgfqpoint{3.529102in}{2.635234in}}%
\pgfpathlineto{\pgfqpoint{3.529102in}{2.632285in}}%
\pgfpathmoveto{\pgfqpoint{3.510939in}{2.635234in}}%
\pgfpathlineto{\pgfqpoint{3.510939in}{2.635234in}}%
\pgfpathlineto{\pgfqpoint{3.510939in}{2.638183in}}%
\pgfpathlineto{\pgfqpoint{3.515479in}{2.638183in}}%
\pgfpathlineto{\pgfqpoint{3.515479in}{2.635234in}}%
\pgfpathmoveto{\pgfqpoint{3.510939in}{2.638183in}}%
\pgfpathlineto{\pgfqpoint{3.510939in}{2.638183in}}%
\pgfpathlineto{\pgfqpoint{3.510939in}{2.641132in}}%
\pgfpathlineto{\pgfqpoint{3.515479in}{2.641132in}}%
\pgfpathlineto{\pgfqpoint{3.515479in}{2.638183in}}%
\pgfpathmoveto{\pgfqpoint{3.515479in}{2.635234in}}%
\pgfpathlineto{\pgfqpoint{3.515479in}{2.635234in}}%
\pgfpathlineto{\pgfqpoint{3.515479in}{2.638183in}}%
\pgfpathlineto{\pgfqpoint{3.520020in}{2.638183in}}%
\pgfpathlineto{\pgfqpoint{3.520020in}{2.635234in}}%
\pgfpathmoveto{\pgfqpoint{3.515479in}{2.638183in}}%
\pgfpathlineto{\pgfqpoint{3.515479in}{2.638183in}}%
\pgfpathlineto{\pgfqpoint{3.515479in}{2.641132in}}%
\pgfpathlineto{\pgfqpoint{3.520020in}{2.641132in}}%
\pgfpathlineto{\pgfqpoint{3.520020in}{2.638183in}}%
\pgfpathmoveto{\pgfqpoint{3.510939in}{2.641132in}}%
\pgfpathlineto{\pgfqpoint{3.510939in}{2.641132in}}%
\pgfpathlineto{\pgfqpoint{3.510939in}{2.644081in}}%
\pgfpathlineto{\pgfqpoint{3.515479in}{2.644081in}}%
\pgfpathlineto{\pgfqpoint{3.515479in}{2.641132in}}%
\pgfpathmoveto{\pgfqpoint{3.510939in}{2.644081in}}%
\pgfpathlineto{\pgfqpoint{3.510939in}{2.644081in}}%
\pgfpathlineto{\pgfqpoint{3.510939in}{2.647031in}}%
\pgfpathlineto{\pgfqpoint{3.515479in}{2.647031in}}%
\pgfpathlineto{\pgfqpoint{3.515479in}{2.644081in}}%
\pgfpathmoveto{\pgfqpoint{3.515479in}{2.641132in}}%
\pgfpathlineto{\pgfqpoint{3.515479in}{2.641132in}}%
\pgfpathlineto{\pgfqpoint{3.515479in}{2.644081in}}%
\pgfpathlineto{\pgfqpoint{3.520020in}{2.644081in}}%
\pgfpathlineto{\pgfqpoint{3.520020in}{2.641132in}}%
\pgfpathmoveto{\pgfqpoint{3.515479in}{2.644081in}}%
\pgfpathlineto{\pgfqpoint{3.515479in}{2.644081in}}%
\pgfpathlineto{\pgfqpoint{3.515479in}{2.647031in}}%
\pgfpathlineto{\pgfqpoint{3.520020in}{2.647031in}}%
\pgfpathlineto{\pgfqpoint{3.520020in}{2.644081in}}%
\pgfpathmoveto{\pgfqpoint{3.520020in}{2.635234in}}%
\pgfpathlineto{\pgfqpoint{3.520020in}{2.635234in}}%
\pgfpathlineto{\pgfqpoint{3.520020in}{2.638183in}}%
\pgfpathlineto{\pgfqpoint{3.524561in}{2.638183in}}%
\pgfpathlineto{\pgfqpoint{3.524561in}{2.635234in}}%
\pgfpathmoveto{\pgfqpoint{3.520020in}{2.638183in}}%
\pgfpathlineto{\pgfqpoint{3.520020in}{2.638183in}}%
\pgfpathlineto{\pgfqpoint{3.520020in}{2.641132in}}%
\pgfpathlineto{\pgfqpoint{3.524561in}{2.641132in}}%
\pgfpathlineto{\pgfqpoint{3.524561in}{2.638183in}}%
\pgfpathmoveto{\pgfqpoint{3.529102in}{2.623437in}}%
\pgfpathlineto{\pgfqpoint{3.529102in}{2.623437in}}%
\pgfpathlineto{\pgfqpoint{3.529102in}{2.626387in}}%
\pgfpathlineto{\pgfqpoint{3.533643in}{2.626387in}}%
\pgfpathlineto{\pgfqpoint{3.533643in}{2.623437in}}%
\pgfpathmoveto{\pgfqpoint{3.529102in}{2.626387in}}%
\pgfpathlineto{\pgfqpoint{3.529102in}{2.626387in}}%
\pgfpathlineto{\pgfqpoint{3.529102in}{2.629336in}}%
\pgfpathlineto{\pgfqpoint{3.533643in}{2.629336in}}%
\pgfpathlineto{\pgfqpoint{3.533643in}{2.626387in}}%
\pgfpathmoveto{\pgfqpoint{3.510939in}{2.647031in}}%
\pgfpathlineto{\pgfqpoint{3.510939in}{2.647031in}}%
\pgfpathlineto{\pgfqpoint{3.510939in}{2.649980in}}%
\pgfpathlineto{\pgfqpoint{3.515479in}{2.649980in}}%
\pgfpathlineto{\pgfqpoint{3.515479in}{2.647031in}}%
\pgfpathmoveto{\pgfqpoint{3.510939in}{2.649980in}}%
\pgfpathlineto{\pgfqpoint{3.510939in}{2.649980in}}%
\pgfpathlineto{\pgfqpoint{3.510939in}{2.652929in}}%
\pgfpathlineto{\pgfqpoint{3.515479in}{2.652929in}}%
\pgfpathlineto{\pgfqpoint{3.515479in}{2.649980in}}%
\pgfpathmoveto{\pgfqpoint{3.656246in}{2.384550in}}%
\pgfpathlineto{\pgfqpoint{3.656246in}{2.384550in}}%
\pgfpathlineto{\pgfqpoint{3.656246in}{2.387499in}}%
\pgfpathlineto{\pgfqpoint{3.660787in}{2.387499in}}%
\pgfpathlineto{\pgfqpoint{3.660787in}{2.384550in}}%
\pgfpathmoveto{\pgfqpoint{3.656246in}{2.387499in}}%
\pgfpathlineto{\pgfqpoint{3.656246in}{2.387499in}}%
\pgfpathlineto{\pgfqpoint{3.656246in}{2.390448in}}%
\pgfpathlineto{\pgfqpoint{3.660787in}{2.390448in}}%
\pgfpathlineto{\pgfqpoint{3.660787in}{2.387499in}}%
\pgfpathmoveto{\pgfqpoint{3.656246in}{2.390448in}}%
\pgfpathlineto{\pgfqpoint{3.656246in}{2.390448in}}%
\pgfpathlineto{\pgfqpoint{3.656246in}{2.393397in}}%
\pgfpathlineto{\pgfqpoint{3.660787in}{2.393397in}}%
\pgfpathlineto{\pgfqpoint{3.660787in}{2.390448in}}%
\pgfpathmoveto{\pgfqpoint{3.660787in}{2.387499in}}%
\pgfpathlineto{\pgfqpoint{3.660787in}{2.387499in}}%
\pgfpathlineto{\pgfqpoint{3.660787in}{2.390448in}}%
\pgfpathlineto{\pgfqpoint{3.665328in}{2.390448in}}%
\pgfpathlineto{\pgfqpoint{3.665328in}{2.387499in}}%
\pgfpathmoveto{\pgfqpoint{3.660787in}{2.390448in}}%
\pgfpathlineto{\pgfqpoint{3.660787in}{2.390448in}}%
\pgfpathlineto{\pgfqpoint{3.660787in}{2.393397in}}%
\pgfpathlineto{\pgfqpoint{3.665328in}{2.393397in}}%
\pgfpathlineto{\pgfqpoint{3.665328in}{2.390448in}}%
\pgfpathmoveto{\pgfqpoint{3.665328in}{2.390448in}}%
\pgfpathlineto{\pgfqpoint{3.665328in}{2.390448in}}%
\pgfpathlineto{\pgfqpoint{3.665328in}{2.393397in}}%
\pgfpathlineto{\pgfqpoint{3.669869in}{2.393397in}}%
\pgfpathlineto{\pgfqpoint{3.669869in}{2.390448in}}%
\pgfpathmoveto{\pgfqpoint{3.665328in}{2.393397in}}%
\pgfpathlineto{\pgfqpoint{3.665328in}{2.393397in}}%
\pgfpathlineto{\pgfqpoint{3.665328in}{2.396347in}}%
\pgfpathlineto{\pgfqpoint{3.669869in}{2.396347in}}%
\pgfpathlineto{\pgfqpoint{3.669869in}{2.393397in}}%
\pgfpathmoveto{\pgfqpoint{3.665328in}{2.396347in}}%
\pgfpathlineto{\pgfqpoint{3.665328in}{2.396347in}}%
\pgfpathlineto{\pgfqpoint{3.665328in}{2.399296in}}%
\pgfpathlineto{\pgfqpoint{3.669869in}{2.399296in}}%
\pgfpathlineto{\pgfqpoint{3.669869in}{2.396347in}}%
\pgfpathmoveto{\pgfqpoint{3.669869in}{2.393397in}}%
\pgfpathlineto{\pgfqpoint{3.669869in}{2.393397in}}%
\pgfpathlineto{\pgfqpoint{3.669869in}{2.396347in}}%
\pgfpathlineto{\pgfqpoint{3.674410in}{2.396347in}}%
\pgfpathlineto{\pgfqpoint{3.674410in}{2.393397in}}%
\pgfpathmoveto{\pgfqpoint{3.669869in}{2.396347in}}%
\pgfpathlineto{\pgfqpoint{3.669869in}{2.396347in}}%
\pgfpathlineto{\pgfqpoint{3.669869in}{2.399296in}}%
\pgfpathlineto{\pgfqpoint{3.674410in}{2.399296in}}%
\pgfpathlineto{\pgfqpoint{3.674410in}{2.396347in}}%
\pgfpathmoveto{\pgfqpoint{3.674410in}{2.396347in}}%
\pgfpathlineto{\pgfqpoint{3.674410in}{2.396347in}}%
\pgfpathlineto{\pgfqpoint{3.674410in}{2.399296in}}%
\pgfpathlineto{\pgfqpoint{3.678951in}{2.399296in}}%
\pgfpathlineto{\pgfqpoint{3.678951in}{2.396347in}}%
\pgfpathmoveto{\pgfqpoint{3.674410in}{2.399296in}}%
\pgfpathlineto{\pgfqpoint{3.674410in}{2.399296in}}%
\pgfpathlineto{\pgfqpoint{3.674410in}{2.402245in}}%
\pgfpathlineto{\pgfqpoint{3.678951in}{2.402245in}}%
\pgfpathlineto{\pgfqpoint{3.678951in}{2.399296in}}%
\pgfpathmoveto{\pgfqpoint{3.674410in}{2.402245in}}%
\pgfpathlineto{\pgfqpoint{3.674410in}{2.402245in}}%
\pgfpathlineto{\pgfqpoint{3.674410in}{2.405195in}}%
\pgfpathlineto{\pgfqpoint{3.678951in}{2.405195in}}%
\pgfpathlineto{\pgfqpoint{3.678951in}{2.402245in}}%
\pgfpathmoveto{\pgfqpoint{3.678951in}{2.399296in}}%
\pgfpathlineto{\pgfqpoint{3.678951in}{2.399296in}}%
\pgfpathlineto{\pgfqpoint{3.678951in}{2.402245in}}%
\pgfpathlineto{\pgfqpoint{3.683492in}{2.402245in}}%
\pgfpathlineto{\pgfqpoint{3.683492in}{2.399296in}}%
\pgfpathmoveto{\pgfqpoint{3.678951in}{2.402245in}}%
\pgfpathlineto{\pgfqpoint{3.678951in}{2.402245in}}%
\pgfpathlineto{\pgfqpoint{3.678951in}{2.405195in}}%
\pgfpathlineto{\pgfqpoint{3.683492in}{2.405195in}}%
\pgfpathlineto{\pgfqpoint{3.683492in}{2.402245in}}%
\pgfpathmoveto{\pgfqpoint{3.683492in}{2.402245in}}%
\pgfpathlineto{\pgfqpoint{3.683492in}{2.402245in}}%
\pgfpathlineto{\pgfqpoint{3.683492in}{2.405195in}}%
\pgfpathlineto{\pgfqpoint{3.688033in}{2.405195in}}%
\pgfpathlineto{\pgfqpoint{3.688033in}{2.402245in}}%
\pgfpathmoveto{\pgfqpoint{3.683492in}{2.405195in}}%
\pgfpathlineto{\pgfqpoint{3.683492in}{2.405195in}}%
\pgfpathlineto{\pgfqpoint{3.683492in}{2.408144in}}%
\pgfpathlineto{\pgfqpoint{3.688033in}{2.408144in}}%
\pgfpathlineto{\pgfqpoint{3.688033in}{2.405195in}}%
\pgfpathmoveto{\pgfqpoint{3.683492in}{2.408144in}}%
\pgfpathlineto{\pgfqpoint{3.683492in}{2.408144in}}%
\pgfpathlineto{\pgfqpoint{3.683492in}{2.411093in}}%
\pgfpathlineto{\pgfqpoint{3.688033in}{2.411093in}}%
\pgfpathlineto{\pgfqpoint{3.688033in}{2.408144in}}%
\pgfpathmoveto{\pgfqpoint{3.688033in}{2.405195in}}%
\pgfpathlineto{\pgfqpoint{3.688033in}{2.405195in}}%
\pgfpathlineto{\pgfqpoint{3.688033in}{2.408144in}}%
\pgfpathlineto{\pgfqpoint{3.692574in}{2.408144in}}%
\pgfpathlineto{\pgfqpoint{3.692574in}{2.405195in}}%
\pgfpathmoveto{\pgfqpoint{3.688033in}{2.408144in}}%
\pgfpathlineto{\pgfqpoint{3.688033in}{2.408144in}}%
\pgfpathlineto{\pgfqpoint{3.688033in}{2.411093in}}%
\pgfpathlineto{\pgfqpoint{3.692574in}{2.411093in}}%
\pgfpathlineto{\pgfqpoint{3.692574in}{2.408144in}}%
\pgfpathmoveto{\pgfqpoint{3.683492in}{2.411093in}}%
\pgfpathlineto{\pgfqpoint{3.683492in}{2.411093in}}%
\pgfpathlineto{\pgfqpoint{3.683492in}{2.414043in}}%
\pgfpathlineto{\pgfqpoint{3.688033in}{2.414043in}}%
\pgfpathlineto{\pgfqpoint{3.688033in}{2.411093in}}%
\pgfpathmoveto{\pgfqpoint{3.683492in}{2.414043in}}%
\pgfpathlineto{\pgfqpoint{3.683492in}{2.414043in}}%
\pgfpathlineto{\pgfqpoint{3.683492in}{2.416992in}}%
\pgfpathlineto{\pgfqpoint{3.688033in}{2.416992in}}%
\pgfpathlineto{\pgfqpoint{3.688033in}{2.414043in}}%
\pgfpathmoveto{\pgfqpoint{3.688033in}{2.411093in}}%
\pgfpathlineto{\pgfqpoint{3.688033in}{2.411093in}}%
\pgfpathlineto{\pgfqpoint{3.688033in}{2.414043in}}%
\pgfpathlineto{\pgfqpoint{3.692574in}{2.414043in}}%
\pgfpathlineto{\pgfqpoint{3.692574in}{2.411093in}}%
\pgfpathmoveto{\pgfqpoint{3.688033in}{2.414043in}}%
\pgfpathlineto{\pgfqpoint{3.688033in}{2.414043in}}%
\pgfpathlineto{\pgfqpoint{3.688033in}{2.416992in}}%
\pgfpathlineto{\pgfqpoint{3.692574in}{2.416992in}}%
\pgfpathlineto{\pgfqpoint{3.692574in}{2.414043in}}%
\pgfpathmoveto{\pgfqpoint{3.683492in}{2.416992in}}%
\pgfpathlineto{\pgfqpoint{3.683492in}{2.416992in}}%
\pgfpathlineto{\pgfqpoint{3.683492in}{2.419941in}}%
\pgfpathlineto{\pgfqpoint{3.688033in}{2.419941in}}%
\pgfpathlineto{\pgfqpoint{3.688033in}{2.416992in}}%
\pgfpathmoveto{\pgfqpoint{3.683492in}{2.419941in}}%
\pgfpathlineto{\pgfqpoint{3.683492in}{2.419941in}}%
\pgfpathlineto{\pgfqpoint{3.683492in}{2.422891in}}%
\pgfpathlineto{\pgfqpoint{3.688033in}{2.422891in}}%
\pgfpathlineto{\pgfqpoint{3.688033in}{2.419941in}}%
\pgfpathmoveto{\pgfqpoint{3.688033in}{2.416992in}}%
\pgfpathlineto{\pgfqpoint{3.688033in}{2.416992in}}%
\pgfpathlineto{\pgfqpoint{3.688033in}{2.419941in}}%
\pgfpathlineto{\pgfqpoint{3.692574in}{2.419941in}}%
\pgfpathlineto{\pgfqpoint{3.692574in}{2.416992in}}%
\pgfpathmoveto{\pgfqpoint{3.688033in}{2.419941in}}%
\pgfpathlineto{\pgfqpoint{3.688033in}{2.419941in}}%
\pgfpathlineto{\pgfqpoint{3.688033in}{2.422891in}}%
\pgfpathlineto{\pgfqpoint{3.692574in}{2.422891in}}%
\pgfpathlineto{\pgfqpoint{3.692574in}{2.419941in}}%
\pgfpathmoveto{\pgfqpoint{3.674410in}{2.422891in}}%
\pgfpathlineto{\pgfqpoint{3.674410in}{2.422891in}}%
\pgfpathlineto{\pgfqpoint{3.674410in}{2.425840in}}%
\pgfpathlineto{\pgfqpoint{3.678951in}{2.425840in}}%
\pgfpathlineto{\pgfqpoint{3.678951in}{2.422891in}}%
\pgfpathmoveto{\pgfqpoint{3.674410in}{2.425840in}}%
\pgfpathlineto{\pgfqpoint{3.674410in}{2.425840in}}%
\pgfpathlineto{\pgfqpoint{3.674410in}{2.428790in}}%
\pgfpathlineto{\pgfqpoint{3.678951in}{2.428790in}}%
\pgfpathlineto{\pgfqpoint{3.678951in}{2.425840in}}%
\pgfpathmoveto{\pgfqpoint{3.678951in}{2.422891in}}%
\pgfpathlineto{\pgfqpoint{3.678951in}{2.422891in}}%
\pgfpathlineto{\pgfqpoint{3.678951in}{2.425840in}}%
\pgfpathlineto{\pgfqpoint{3.683492in}{2.425840in}}%
\pgfpathlineto{\pgfqpoint{3.683492in}{2.422891in}}%
\pgfpathmoveto{\pgfqpoint{3.678951in}{2.425840in}}%
\pgfpathlineto{\pgfqpoint{3.678951in}{2.425840in}}%
\pgfpathlineto{\pgfqpoint{3.678951in}{2.428790in}}%
\pgfpathlineto{\pgfqpoint{3.683492in}{2.428790in}}%
\pgfpathlineto{\pgfqpoint{3.683492in}{2.425840in}}%
\pgfpathmoveto{\pgfqpoint{3.674410in}{2.428790in}}%
\pgfpathlineto{\pgfqpoint{3.674410in}{2.428790in}}%
\pgfpathlineto{\pgfqpoint{3.674410in}{2.431739in}}%
\pgfpathlineto{\pgfqpoint{3.678951in}{2.431739in}}%
\pgfpathlineto{\pgfqpoint{3.678951in}{2.428790in}}%
\pgfpathmoveto{\pgfqpoint{3.674410in}{2.431739in}}%
\pgfpathlineto{\pgfqpoint{3.674410in}{2.431739in}}%
\pgfpathlineto{\pgfqpoint{3.674410in}{2.434688in}}%
\pgfpathlineto{\pgfqpoint{3.678951in}{2.434688in}}%
\pgfpathlineto{\pgfqpoint{3.678951in}{2.431739in}}%
\pgfpathmoveto{\pgfqpoint{3.678951in}{2.428790in}}%
\pgfpathlineto{\pgfqpoint{3.678951in}{2.428790in}}%
\pgfpathlineto{\pgfqpoint{3.678951in}{2.431739in}}%
\pgfpathlineto{\pgfqpoint{3.683492in}{2.431739in}}%
\pgfpathlineto{\pgfqpoint{3.683492in}{2.428790in}}%
\pgfpathmoveto{\pgfqpoint{3.678951in}{2.431739in}}%
\pgfpathlineto{\pgfqpoint{3.678951in}{2.431739in}}%
\pgfpathlineto{\pgfqpoint{3.678951in}{2.434688in}}%
\pgfpathlineto{\pgfqpoint{3.683492in}{2.434688in}}%
\pgfpathlineto{\pgfqpoint{3.683492in}{2.431739in}}%
\pgfpathmoveto{\pgfqpoint{3.683492in}{2.422891in}}%
\pgfpathlineto{\pgfqpoint{3.683492in}{2.422891in}}%
\pgfpathlineto{\pgfqpoint{3.683492in}{2.425840in}}%
\pgfpathlineto{\pgfqpoint{3.688033in}{2.425840in}}%
\pgfpathlineto{\pgfqpoint{3.688033in}{2.422891in}}%
\pgfpathmoveto{\pgfqpoint{3.683492in}{2.425840in}}%
\pgfpathlineto{\pgfqpoint{3.683492in}{2.425840in}}%
\pgfpathlineto{\pgfqpoint{3.683492in}{2.428790in}}%
\pgfpathlineto{\pgfqpoint{3.688033in}{2.428790in}}%
\pgfpathlineto{\pgfqpoint{3.688033in}{2.425840in}}%
\pgfpathmoveto{\pgfqpoint{3.692574in}{2.408144in}}%
\pgfpathlineto{\pgfqpoint{3.692574in}{2.408144in}}%
\pgfpathlineto{\pgfqpoint{3.692574in}{2.411093in}}%
\pgfpathlineto{\pgfqpoint{3.697115in}{2.411093in}}%
\pgfpathlineto{\pgfqpoint{3.697115in}{2.408144in}}%
\pgfpathmoveto{\pgfqpoint{3.692574in}{2.411093in}}%
\pgfpathlineto{\pgfqpoint{3.692574in}{2.411093in}}%
\pgfpathlineto{\pgfqpoint{3.692574in}{2.414043in}}%
\pgfpathlineto{\pgfqpoint{3.697115in}{2.414043in}}%
\pgfpathlineto{\pgfqpoint{3.697115in}{2.411093in}}%
\pgfpathmoveto{\pgfqpoint{3.692574in}{2.414043in}}%
\pgfpathlineto{\pgfqpoint{3.692574in}{2.414043in}}%
\pgfpathlineto{\pgfqpoint{3.692574in}{2.416992in}}%
\pgfpathlineto{\pgfqpoint{3.697115in}{2.416992in}}%
\pgfpathlineto{\pgfqpoint{3.697115in}{2.414043in}}%
\pgfpathmoveto{\pgfqpoint{3.665328in}{2.434688in}}%
\pgfpathlineto{\pgfqpoint{3.665328in}{2.434688in}}%
\pgfpathlineto{\pgfqpoint{3.665328in}{2.437638in}}%
\pgfpathlineto{\pgfqpoint{3.669869in}{2.437638in}}%
\pgfpathlineto{\pgfqpoint{3.669869in}{2.434688in}}%
\pgfpathmoveto{\pgfqpoint{3.665328in}{2.437638in}}%
\pgfpathlineto{\pgfqpoint{3.665328in}{2.437638in}}%
\pgfpathlineto{\pgfqpoint{3.665328in}{2.440587in}}%
\pgfpathlineto{\pgfqpoint{3.669869in}{2.440587in}}%
\pgfpathlineto{\pgfqpoint{3.669869in}{2.437638in}}%
\pgfpathmoveto{\pgfqpoint{3.669869in}{2.434688in}}%
\pgfpathlineto{\pgfqpoint{3.669869in}{2.434688in}}%
\pgfpathlineto{\pgfqpoint{3.669869in}{2.437638in}}%
\pgfpathlineto{\pgfqpoint{3.674410in}{2.437638in}}%
\pgfpathlineto{\pgfqpoint{3.674410in}{2.434688in}}%
\pgfpathmoveto{\pgfqpoint{3.669869in}{2.437638in}}%
\pgfpathlineto{\pgfqpoint{3.669869in}{2.437638in}}%
\pgfpathlineto{\pgfqpoint{3.669869in}{2.440587in}}%
\pgfpathlineto{\pgfqpoint{3.674410in}{2.440587in}}%
\pgfpathlineto{\pgfqpoint{3.674410in}{2.437638in}}%
\pgfpathmoveto{\pgfqpoint{3.665328in}{2.440587in}}%
\pgfpathlineto{\pgfqpoint{3.665328in}{2.440587in}}%
\pgfpathlineto{\pgfqpoint{3.665328in}{2.443536in}}%
\pgfpathlineto{\pgfqpoint{3.669869in}{2.443536in}}%
\pgfpathlineto{\pgfqpoint{3.669869in}{2.440587in}}%
\pgfpathmoveto{\pgfqpoint{3.665328in}{2.443536in}}%
\pgfpathlineto{\pgfqpoint{3.665328in}{2.443536in}}%
\pgfpathlineto{\pgfqpoint{3.665328in}{2.446486in}}%
\pgfpathlineto{\pgfqpoint{3.669869in}{2.446486in}}%
\pgfpathlineto{\pgfqpoint{3.669869in}{2.443536in}}%
\pgfpathmoveto{\pgfqpoint{3.669869in}{2.440587in}}%
\pgfpathlineto{\pgfqpoint{3.669869in}{2.440587in}}%
\pgfpathlineto{\pgfqpoint{3.669869in}{2.443536in}}%
\pgfpathlineto{\pgfqpoint{3.674410in}{2.443536in}}%
\pgfpathlineto{\pgfqpoint{3.674410in}{2.440587in}}%
\pgfpathmoveto{\pgfqpoint{3.669869in}{2.443536in}}%
\pgfpathlineto{\pgfqpoint{3.669869in}{2.443536in}}%
\pgfpathlineto{\pgfqpoint{3.669869in}{2.446486in}}%
\pgfpathlineto{\pgfqpoint{3.674410in}{2.446486in}}%
\pgfpathlineto{\pgfqpoint{3.674410in}{2.443536in}}%
\pgfpathmoveto{\pgfqpoint{3.656246in}{2.446486in}}%
\pgfpathlineto{\pgfqpoint{3.656246in}{2.446486in}}%
\pgfpathlineto{\pgfqpoint{3.656246in}{2.449435in}}%
\pgfpathlineto{\pgfqpoint{3.660787in}{2.449435in}}%
\pgfpathlineto{\pgfqpoint{3.660787in}{2.446486in}}%
\pgfpathmoveto{\pgfqpoint{3.656246in}{2.449435in}}%
\pgfpathlineto{\pgfqpoint{3.656246in}{2.449435in}}%
\pgfpathlineto{\pgfqpoint{3.656246in}{2.452384in}}%
\pgfpathlineto{\pgfqpoint{3.660787in}{2.452384in}}%
\pgfpathlineto{\pgfqpoint{3.660787in}{2.449435in}}%
\pgfpathmoveto{\pgfqpoint{3.660787in}{2.446486in}}%
\pgfpathlineto{\pgfqpoint{3.660787in}{2.446486in}}%
\pgfpathlineto{\pgfqpoint{3.660787in}{2.449435in}}%
\pgfpathlineto{\pgfqpoint{3.665328in}{2.449435in}}%
\pgfpathlineto{\pgfqpoint{3.665328in}{2.446486in}}%
\pgfpathmoveto{\pgfqpoint{3.660787in}{2.449435in}}%
\pgfpathlineto{\pgfqpoint{3.660787in}{2.449435in}}%
\pgfpathlineto{\pgfqpoint{3.660787in}{2.452384in}}%
\pgfpathlineto{\pgfqpoint{3.665328in}{2.452384in}}%
\pgfpathlineto{\pgfqpoint{3.665328in}{2.449435in}}%
\pgfpathmoveto{\pgfqpoint{3.656246in}{2.452384in}}%
\pgfpathlineto{\pgfqpoint{3.656246in}{2.452384in}}%
\pgfpathlineto{\pgfqpoint{3.656246in}{2.455334in}}%
\pgfpathlineto{\pgfqpoint{3.660787in}{2.455334in}}%
\pgfpathlineto{\pgfqpoint{3.660787in}{2.452384in}}%
\pgfpathmoveto{\pgfqpoint{3.656246in}{2.455334in}}%
\pgfpathlineto{\pgfqpoint{3.656246in}{2.455334in}}%
\pgfpathlineto{\pgfqpoint{3.656246in}{2.458283in}}%
\pgfpathlineto{\pgfqpoint{3.660787in}{2.458283in}}%
\pgfpathlineto{\pgfqpoint{3.660787in}{2.455334in}}%
\pgfpathmoveto{\pgfqpoint{3.660787in}{2.452384in}}%
\pgfpathlineto{\pgfqpoint{3.660787in}{2.452384in}}%
\pgfpathlineto{\pgfqpoint{3.660787in}{2.455334in}}%
\pgfpathlineto{\pgfqpoint{3.665328in}{2.455334in}}%
\pgfpathlineto{\pgfqpoint{3.665328in}{2.452384in}}%
\pgfpathmoveto{\pgfqpoint{3.660787in}{2.455334in}}%
\pgfpathlineto{\pgfqpoint{3.660787in}{2.455334in}}%
\pgfpathlineto{\pgfqpoint{3.660787in}{2.458283in}}%
\pgfpathlineto{\pgfqpoint{3.665328in}{2.458283in}}%
\pgfpathlineto{\pgfqpoint{3.665328in}{2.455334in}}%
\pgfpathmoveto{\pgfqpoint{3.665328in}{2.446486in}}%
\pgfpathlineto{\pgfqpoint{3.665328in}{2.446486in}}%
\pgfpathlineto{\pgfqpoint{3.665328in}{2.449435in}}%
\pgfpathlineto{\pgfqpoint{3.669869in}{2.449435in}}%
\pgfpathlineto{\pgfqpoint{3.669869in}{2.446486in}}%
\pgfpathmoveto{\pgfqpoint{3.665328in}{2.449435in}}%
\pgfpathlineto{\pgfqpoint{3.665328in}{2.449435in}}%
\pgfpathlineto{\pgfqpoint{3.665328in}{2.452384in}}%
\pgfpathlineto{\pgfqpoint{3.669869in}{2.452384in}}%
\pgfpathlineto{\pgfqpoint{3.669869in}{2.449435in}}%
\pgfpathmoveto{\pgfqpoint{3.674410in}{2.434688in}}%
\pgfpathlineto{\pgfqpoint{3.674410in}{2.434688in}}%
\pgfpathlineto{\pgfqpoint{3.674410in}{2.437638in}}%
\pgfpathlineto{\pgfqpoint{3.678951in}{2.437638in}}%
\pgfpathlineto{\pgfqpoint{3.678951in}{2.434688in}}%
\pgfpathmoveto{\pgfqpoint{3.674410in}{2.437638in}}%
\pgfpathlineto{\pgfqpoint{3.674410in}{2.437638in}}%
\pgfpathlineto{\pgfqpoint{3.674410in}{2.440587in}}%
\pgfpathlineto{\pgfqpoint{3.678951in}{2.440587in}}%
\pgfpathlineto{\pgfqpoint{3.678951in}{2.437638in}}%
\pgfpathmoveto{\pgfqpoint{3.656246in}{2.458283in}}%
\pgfpathlineto{\pgfqpoint{3.656246in}{2.458283in}}%
\pgfpathlineto{\pgfqpoint{3.656246in}{2.461232in}}%
\pgfpathlineto{\pgfqpoint{3.660787in}{2.461232in}}%
\pgfpathlineto{\pgfqpoint{3.660787in}{2.458283in}}%
\pgfpathmoveto{\pgfqpoint{3.656246in}{2.461232in}}%
\pgfpathlineto{\pgfqpoint{3.656246in}{2.461232in}}%
\pgfpathlineto{\pgfqpoint{3.656246in}{2.464182in}}%
\pgfpathlineto{\pgfqpoint{3.660787in}{2.464182in}}%
\pgfpathlineto{\pgfqpoint{3.660787in}{2.461232in}}%
\pgfpathclose%
\pgfusepath{fill}%
\end{pgfscope}%
\begin{pgfscope}%
\pgfpathrectangle{\pgfqpoint{0.750000in}{0.500000in}}{\pgfqpoint{4.650000in}{3.020000in}}%
\pgfusepath{clip}%
\pgfsetbuttcap%
\pgfsetmiterjoin%
\definecolor{currentfill}{rgb}{1.000000,0.000000,0.000000}%
\pgfsetfillcolor{currentfill}%
\pgfsetlinewidth{0.000000pt}%
\definecolor{currentstroke}{rgb}{0.000000,0.000000,0.000000}%
\pgfsetstrokecolor{currentstroke}%
\pgfsetstrokeopacity{0.000000}%
\pgfsetdash{}{0pt}%
\pgfpathmoveto{\pgfqpoint{2.925145in}{3.404980in}}%
\pgfpathlineto{\pgfqpoint{2.925145in}{3.407930in}}%
\pgfpathlineto{\pgfqpoint{2.929686in}{3.407930in}}%
\pgfpathlineto{\pgfqpoint{2.929686in}{3.404980in}}%
\pgfpathmoveto{\pgfqpoint{2.920604in}{3.410879in}}%
\pgfpathlineto{\pgfqpoint{2.920604in}{3.410879in}}%
\pgfpathlineto{\pgfqpoint{2.920604in}{3.413828in}}%
\pgfpathlineto{\pgfqpoint{2.925145in}{3.413828in}}%
\pgfpathlineto{\pgfqpoint{2.925145in}{3.410879in}}%
\pgfpathmoveto{\pgfqpoint{2.925145in}{3.407930in}}%
\pgfpathlineto{\pgfqpoint{2.925145in}{3.407930in}}%
\pgfpathlineto{\pgfqpoint{2.925145in}{3.410879in}}%
\pgfpathlineto{\pgfqpoint{2.929686in}{3.410879in}}%
\pgfpathlineto{\pgfqpoint{2.929686in}{3.407930in}}%
\pgfpathmoveto{\pgfqpoint{2.925145in}{3.410879in}}%
\pgfpathlineto{\pgfqpoint{2.925145in}{3.410879in}}%
\pgfpathlineto{\pgfqpoint{2.925145in}{3.413828in}}%
\pgfpathlineto{\pgfqpoint{2.929686in}{3.413828in}}%
\pgfpathlineto{\pgfqpoint{2.929686in}{3.410879in}}%
\pgfpathmoveto{\pgfqpoint{2.916064in}{3.416777in}}%
\pgfpathlineto{\pgfqpoint{2.916064in}{3.416777in}}%
\pgfpathlineto{\pgfqpoint{2.916064in}{3.419727in}}%
\pgfpathlineto{\pgfqpoint{2.920604in}{3.419727in}}%
\pgfpathlineto{\pgfqpoint{2.920604in}{3.416777in}}%
\pgfpathmoveto{\pgfqpoint{2.911523in}{3.422676in}}%
\pgfpathlineto{\pgfqpoint{2.911523in}{3.422676in}}%
\pgfpathlineto{\pgfqpoint{2.911523in}{3.425625in}}%
\pgfpathlineto{\pgfqpoint{2.916064in}{3.425625in}}%
\pgfpathlineto{\pgfqpoint{2.916064in}{3.422676in}}%
\pgfpathmoveto{\pgfqpoint{2.916064in}{3.419727in}}%
\pgfpathlineto{\pgfqpoint{2.916064in}{3.419727in}}%
\pgfpathlineto{\pgfqpoint{2.916064in}{3.422676in}}%
\pgfpathlineto{\pgfqpoint{2.920604in}{3.422676in}}%
\pgfpathlineto{\pgfqpoint{2.920604in}{3.419727in}}%
\pgfpathmoveto{\pgfqpoint{2.916064in}{3.422676in}}%
\pgfpathlineto{\pgfqpoint{2.916064in}{3.422676in}}%
\pgfpathlineto{\pgfqpoint{2.916064in}{3.425625in}}%
\pgfpathlineto{\pgfqpoint{2.920604in}{3.425625in}}%
\pgfpathlineto{\pgfqpoint{2.920604in}{3.422676in}}%
\pgfpathmoveto{\pgfqpoint{2.920604in}{3.413828in}}%
\pgfpathlineto{\pgfqpoint{2.920604in}{3.413828in}}%
\pgfpathlineto{\pgfqpoint{2.920604in}{3.416777in}}%
\pgfpathlineto{\pgfqpoint{2.925145in}{3.416777in}}%
\pgfpathlineto{\pgfqpoint{2.925145in}{3.413828in}}%
\pgfpathmoveto{\pgfqpoint{2.920604in}{3.416777in}}%
\pgfpathlineto{\pgfqpoint{2.920604in}{3.416777in}}%
\pgfpathlineto{\pgfqpoint{2.920604in}{3.419727in}}%
\pgfpathlineto{\pgfqpoint{2.925145in}{3.419727in}}%
\pgfpathlineto{\pgfqpoint{2.925145in}{3.416777in}}%
\pgfpathmoveto{\pgfqpoint{2.852492in}{3.499357in}}%
\pgfpathlineto{\pgfqpoint{2.852492in}{3.499357in}}%
\pgfpathlineto{\pgfqpoint{2.852492in}{3.502307in}}%
\pgfpathlineto{\pgfqpoint{2.857033in}{3.502307in}}%
\pgfpathlineto{\pgfqpoint{2.857033in}{3.499357in}}%
\pgfpathmoveto{\pgfqpoint{2.847951in}{3.505256in}}%
\pgfpathlineto{\pgfqpoint{2.847951in}{3.505256in}}%
\pgfpathlineto{\pgfqpoint{2.847951in}{3.508205in}}%
\pgfpathlineto{\pgfqpoint{2.852492in}{3.508205in}}%
\pgfpathlineto{\pgfqpoint{2.852492in}{3.505256in}}%
\pgfpathmoveto{\pgfqpoint{2.852492in}{3.502307in}}%
\pgfpathlineto{\pgfqpoint{2.852492in}{3.502307in}}%
\pgfpathlineto{\pgfqpoint{2.852492in}{3.505256in}}%
\pgfpathlineto{\pgfqpoint{2.857033in}{3.505256in}}%
\pgfpathlineto{\pgfqpoint{2.857033in}{3.502307in}}%
\pgfpathmoveto{\pgfqpoint{2.852492in}{3.505256in}}%
\pgfpathlineto{\pgfqpoint{2.852492in}{3.505256in}}%
\pgfpathlineto{\pgfqpoint{2.852492in}{3.508205in}}%
\pgfpathlineto{\pgfqpoint{2.857033in}{3.508205in}}%
\pgfpathlineto{\pgfqpoint{2.857033in}{3.505256in}}%
\pgfpathmoveto{\pgfqpoint{2.843410in}{3.511155in}}%
\pgfpathlineto{\pgfqpoint{2.843410in}{3.511155in}}%
\pgfpathlineto{\pgfqpoint{2.843410in}{3.514104in}}%
\pgfpathlineto{\pgfqpoint{2.847951in}{3.514104in}}%
\pgfpathlineto{\pgfqpoint{2.847951in}{3.511155in}}%
\pgfpathmoveto{\pgfqpoint{2.838869in}{3.517053in}}%
\pgfpathlineto{\pgfqpoint{2.838869in}{3.517053in}}%
\pgfpathlineto{\pgfqpoint{2.838869in}{3.520002in}}%
\pgfpathlineto{\pgfqpoint{2.843410in}{3.520002in}}%
\pgfpathlineto{\pgfqpoint{2.843410in}{3.517053in}}%
\pgfpathmoveto{\pgfqpoint{2.843410in}{3.514104in}}%
\pgfpathlineto{\pgfqpoint{2.843410in}{3.514104in}}%
\pgfpathlineto{\pgfqpoint{2.843410in}{3.517053in}}%
\pgfpathlineto{\pgfqpoint{2.847951in}{3.517053in}}%
\pgfpathlineto{\pgfqpoint{2.847951in}{3.514104in}}%
\pgfpathmoveto{\pgfqpoint{2.843410in}{3.517053in}}%
\pgfpathlineto{\pgfqpoint{2.843410in}{3.517053in}}%
\pgfpathlineto{\pgfqpoint{2.843410in}{3.520002in}}%
\pgfpathlineto{\pgfqpoint{2.847951in}{3.520002in}}%
\pgfpathlineto{\pgfqpoint{2.847951in}{3.517053in}}%
\pgfpathmoveto{\pgfqpoint{2.847951in}{3.508205in}}%
\pgfpathlineto{\pgfqpoint{2.847951in}{3.508205in}}%
\pgfpathlineto{\pgfqpoint{2.847951in}{3.511155in}}%
\pgfpathlineto{\pgfqpoint{2.852492in}{3.511155in}}%
\pgfpathlineto{\pgfqpoint{2.852492in}{3.508205in}}%
\pgfpathmoveto{\pgfqpoint{2.847951in}{3.511155in}}%
\pgfpathlineto{\pgfqpoint{2.847951in}{3.511155in}}%
\pgfpathlineto{\pgfqpoint{2.847951in}{3.514104in}}%
\pgfpathlineto{\pgfqpoint{2.852492in}{3.514104in}}%
\pgfpathlineto{\pgfqpoint{2.852492in}{3.511155in}}%
\pgfpathmoveto{\pgfqpoint{2.888819in}{3.452169in}}%
\pgfpathlineto{\pgfqpoint{2.888819in}{3.452169in}}%
\pgfpathlineto{\pgfqpoint{2.888819in}{3.455118in}}%
\pgfpathlineto{\pgfqpoint{2.893359in}{3.455118in}}%
\pgfpathlineto{\pgfqpoint{2.893359in}{3.452169in}}%
\pgfpathmoveto{\pgfqpoint{2.884278in}{3.458067in}}%
\pgfpathlineto{\pgfqpoint{2.884278in}{3.458067in}}%
\pgfpathlineto{\pgfqpoint{2.884278in}{3.461017in}}%
\pgfpathlineto{\pgfqpoint{2.888819in}{3.461017in}}%
\pgfpathlineto{\pgfqpoint{2.888819in}{3.458067in}}%
\pgfpathmoveto{\pgfqpoint{2.888819in}{3.455118in}}%
\pgfpathlineto{\pgfqpoint{2.888819in}{3.455118in}}%
\pgfpathlineto{\pgfqpoint{2.888819in}{3.458067in}}%
\pgfpathlineto{\pgfqpoint{2.893359in}{3.458067in}}%
\pgfpathlineto{\pgfqpoint{2.893359in}{3.455118in}}%
\pgfpathmoveto{\pgfqpoint{2.888819in}{3.458067in}}%
\pgfpathlineto{\pgfqpoint{2.888819in}{3.458067in}}%
\pgfpathlineto{\pgfqpoint{2.888819in}{3.461017in}}%
\pgfpathlineto{\pgfqpoint{2.893359in}{3.461017in}}%
\pgfpathlineto{\pgfqpoint{2.893359in}{3.458067in}}%
\pgfpathmoveto{\pgfqpoint{2.879737in}{3.463966in}}%
\pgfpathlineto{\pgfqpoint{2.879737in}{3.463966in}}%
\pgfpathlineto{\pgfqpoint{2.879737in}{3.466915in}}%
\pgfpathlineto{\pgfqpoint{2.884278in}{3.466915in}}%
\pgfpathlineto{\pgfqpoint{2.884278in}{3.463966in}}%
\pgfpathmoveto{\pgfqpoint{2.875196in}{3.469865in}}%
\pgfpathlineto{\pgfqpoint{2.875196in}{3.469865in}}%
\pgfpathlineto{\pgfqpoint{2.875196in}{3.472814in}}%
\pgfpathlineto{\pgfqpoint{2.879737in}{3.472814in}}%
\pgfpathlineto{\pgfqpoint{2.879737in}{3.469865in}}%
\pgfpathmoveto{\pgfqpoint{2.879737in}{3.466915in}}%
\pgfpathlineto{\pgfqpoint{2.879737in}{3.466915in}}%
\pgfpathlineto{\pgfqpoint{2.879737in}{3.469865in}}%
\pgfpathlineto{\pgfqpoint{2.884278in}{3.469865in}}%
\pgfpathlineto{\pgfqpoint{2.884278in}{3.466915in}}%
\pgfpathmoveto{\pgfqpoint{2.879737in}{3.469865in}}%
\pgfpathlineto{\pgfqpoint{2.879737in}{3.469865in}}%
\pgfpathlineto{\pgfqpoint{2.879737in}{3.472814in}}%
\pgfpathlineto{\pgfqpoint{2.884278in}{3.472814in}}%
\pgfpathlineto{\pgfqpoint{2.884278in}{3.469865in}}%
\pgfpathmoveto{\pgfqpoint{2.884278in}{3.461017in}}%
\pgfpathlineto{\pgfqpoint{2.884278in}{3.461017in}}%
\pgfpathlineto{\pgfqpoint{2.884278in}{3.463966in}}%
\pgfpathlineto{\pgfqpoint{2.888819in}{3.463966in}}%
\pgfpathlineto{\pgfqpoint{2.888819in}{3.461017in}}%
\pgfpathmoveto{\pgfqpoint{2.884278in}{3.463966in}}%
\pgfpathlineto{\pgfqpoint{2.884278in}{3.463966in}}%
\pgfpathlineto{\pgfqpoint{2.884278in}{3.466915in}}%
\pgfpathlineto{\pgfqpoint{2.888819in}{3.466915in}}%
\pgfpathlineto{\pgfqpoint{2.888819in}{3.463966in}}%
\pgfpathmoveto{\pgfqpoint{2.906982in}{3.428575in}}%
\pgfpathlineto{\pgfqpoint{2.906982in}{3.428575in}}%
\pgfpathlineto{\pgfqpoint{2.906982in}{3.431524in}}%
\pgfpathlineto{\pgfqpoint{2.911523in}{3.431524in}}%
\pgfpathlineto{\pgfqpoint{2.911523in}{3.428575in}}%
\pgfpathmoveto{\pgfqpoint{2.902441in}{3.434473in}}%
\pgfpathlineto{\pgfqpoint{2.902441in}{3.434473in}}%
\pgfpathlineto{\pgfqpoint{2.902441in}{3.437422in}}%
\pgfpathlineto{\pgfqpoint{2.906982in}{3.437422in}}%
\pgfpathlineto{\pgfqpoint{2.906982in}{3.434473in}}%
\pgfpathmoveto{\pgfqpoint{2.906982in}{3.431524in}}%
\pgfpathlineto{\pgfqpoint{2.906982in}{3.431524in}}%
\pgfpathlineto{\pgfqpoint{2.906982in}{3.434473in}}%
\pgfpathlineto{\pgfqpoint{2.911523in}{3.434473in}}%
\pgfpathlineto{\pgfqpoint{2.911523in}{3.431524in}}%
\pgfpathmoveto{\pgfqpoint{2.906982in}{3.434473in}}%
\pgfpathlineto{\pgfqpoint{2.906982in}{3.434473in}}%
\pgfpathlineto{\pgfqpoint{2.906982in}{3.437422in}}%
\pgfpathlineto{\pgfqpoint{2.911523in}{3.437422in}}%
\pgfpathlineto{\pgfqpoint{2.911523in}{3.434473in}}%
\pgfpathmoveto{\pgfqpoint{2.897900in}{3.440372in}}%
\pgfpathlineto{\pgfqpoint{2.897900in}{3.440372in}}%
\pgfpathlineto{\pgfqpoint{2.897900in}{3.443321in}}%
\pgfpathlineto{\pgfqpoint{2.902441in}{3.443321in}}%
\pgfpathlineto{\pgfqpoint{2.902441in}{3.440372in}}%
\pgfpathmoveto{\pgfqpoint{2.893359in}{3.446270in}}%
\pgfpathlineto{\pgfqpoint{2.893359in}{3.446270in}}%
\pgfpathlineto{\pgfqpoint{2.893359in}{3.449220in}}%
\pgfpathlineto{\pgfqpoint{2.897900in}{3.449220in}}%
\pgfpathlineto{\pgfqpoint{2.897900in}{3.446270in}}%
\pgfpathmoveto{\pgfqpoint{2.897900in}{3.443321in}}%
\pgfpathlineto{\pgfqpoint{2.897900in}{3.443321in}}%
\pgfpathlineto{\pgfqpoint{2.897900in}{3.446270in}}%
\pgfpathlineto{\pgfqpoint{2.902441in}{3.446270in}}%
\pgfpathlineto{\pgfqpoint{2.902441in}{3.443321in}}%
\pgfpathmoveto{\pgfqpoint{2.897900in}{3.446270in}}%
\pgfpathlineto{\pgfqpoint{2.897900in}{3.446270in}}%
\pgfpathlineto{\pgfqpoint{2.897900in}{3.449220in}}%
\pgfpathlineto{\pgfqpoint{2.902441in}{3.449220in}}%
\pgfpathlineto{\pgfqpoint{2.902441in}{3.446270in}}%
\pgfpathmoveto{\pgfqpoint{2.902441in}{3.437422in}}%
\pgfpathlineto{\pgfqpoint{2.902441in}{3.437422in}}%
\pgfpathlineto{\pgfqpoint{2.902441in}{3.440372in}}%
\pgfpathlineto{\pgfqpoint{2.906982in}{3.440372in}}%
\pgfpathlineto{\pgfqpoint{2.906982in}{3.437422in}}%
\pgfpathmoveto{\pgfqpoint{2.902441in}{3.440372in}}%
\pgfpathlineto{\pgfqpoint{2.902441in}{3.440372in}}%
\pgfpathlineto{\pgfqpoint{2.902441in}{3.443321in}}%
\pgfpathlineto{\pgfqpoint{2.906982in}{3.443321in}}%
\pgfpathlineto{\pgfqpoint{2.906982in}{3.440372in}}%
\pgfpathmoveto{\pgfqpoint{2.911523in}{3.425625in}}%
\pgfpathlineto{\pgfqpoint{2.911523in}{3.425625in}}%
\pgfpathlineto{\pgfqpoint{2.911523in}{3.428575in}}%
\pgfpathlineto{\pgfqpoint{2.916064in}{3.428575in}}%
\pgfpathlineto{\pgfqpoint{2.916064in}{3.425625in}}%
\pgfpathmoveto{\pgfqpoint{2.911523in}{3.428575in}}%
\pgfpathlineto{\pgfqpoint{2.911523in}{3.428575in}}%
\pgfpathlineto{\pgfqpoint{2.911523in}{3.431524in}}%
\pgfpathlineto{\pgfqpoint{2.916064in}{3.431524in}}%
\pgfpathlineto{\pgfqpoint{2.916064in}{3.428575in}}%
\pgfpathmoveto{\pgfqpoint{2.893359in}{3.449220in}}%
\pgfpathlineto{\pgfqpoint{2.893359in}{3.449220in}}%
\pgfpathlineto{\pgfqpoint{2.893359in}{3.452169in}}%
\pgfpathlineto{\pgfqpoint{2.897900in}{3.452169in}}%
\pgfpathlineto{\pgfqpoint{2.897900in}{3.449220in}}%
\pgfpathmoveto{\pgfqpoint{2.893359in}{3.452169in}}%
\pgfpathlineto{\pgfqpoint{2.893359in}{3.452169in}}%
\pgfpathlineto{\pgfqpoint{2.893359in}{3.455118in}}%
\pgfpathlineto{\pgfqpoint{2.897900in}{3.455118in}}%
\pgfpathlineto{\pgfqpoint{2.897900in}{3.452169in}}%
\pgfpathmoveto{\pgfqpoint{2.870655in}{3.475763in}}%
\pgfpathlineto{\pgfqpoint{2.870655in}{3.475763in}}%
\pgfpathlineto{\pgfqpoint{2.870655in}{3.478712in}}%
\pgfpathlineto{\pgfqpoint{2.875196in}{3.478712in}}%
\pgfpathlineto{\pgfqpoint{2.875196in}{3.475763in}}%
\pgfpathmoveto{\pgfqpoint{2.866114in}{3.481662in}}%
\pgfpathlineto{\pgfqpoint{2.866114in}{3.481662in}}%
\pgfpathlineto{\pgfqpoint{2.866114in}{3.484611in}}%
\pgfpathlineto{\pgfqpoint{2.870655in}{3.484611in}}%
\pgfpathlineto{\pgfqpoint{2.870655in}{3.481662in}}%
\pgfpathmoveto{\pgfqpoint{2.870655in}{3.478712in}}%
\pgfpathlineto{\pgfqpoint{2.870655in}{3.478712in}}%
\pgfpathlineto{\pgfqpoint{2.870655in}{3.481662in}}%
\pgfpathlineto{\pgfqpoint{2.875196in}{3.481662in}}%
\pgfpathlineto{\pgfqpoint{2.875196in}{3.478712in}}%
\pgfpathmoveto{\pgfqpoint{2.870655in}{3.481662in}}%
\pgfpathlineto{\pgfqpoint{2.870655in}{3.481662in}}%
\pgfpathlineto{\pgfqpoint{2.870655in}{3.484611in}}%
\pgfpathlineto{\pgfqpoint{2.875196in}{3.484611in}}%
\pgfpathlineto{\pgfqpoint{2.875196in}{3.481662in}}%
\pgfpathmoveto{\pgfqpoint{2.861573in}{3.487560in}}%
\pgfpathlineto{\pgfqpoint{2.861573in}{3.487560in}}%
\pgfpathlineto{\pgfqpoint{2.861573in}{3.490510in}}%
\pgfpathlineto{\pgfqpoint{2.866114in}{3.490510in}}%
\pgfpathlineto{\pgfqpoint{2.866114in}{3.487560in}}%
\pgfpathmoveto{\pgfqpoint{2.857033in}{3.493459in}}%
\pgfpathlineto{\pgfqpoint{2.857033in}{3.493459in}}%
\pgfpathlineto{\pgfqpoint{2.857033in}{3.496408in}}%
\pgfpathlineto{\pgfqpoint{2.861573in}{3.496408in}}%
\pgfpathlineto{\pgfqpoint{2.861573in}{3.493459in}}%
\pgfpathmoveto{\pgfqpoint{2.861573in}{3.490510in}}%
\pgfpathlineto{\pgfqpoint{2.861573in}{3.490510in}}%
\pgfpathlineto{\pgfqpoint{2.861573in}{3.493459in}}%
\pgfpathlineto{\pgfqpoint{2.866114in}{3.493459in}}%
\pgfpathlineto{\pgfqpoint{2.866114in}{3.490510in}}%
\pgfpathmoveto{\pgfqpoint{2.861573in}{3.493459in}}%
\pgfpathlineto{\pgfqpoint{2.861573in}{3.493459in}}%
\pgfpathlineto{\pgfqpoint{2.861573in}{3.496408in}}%
\pgfpathlineto{\pgfqpoint{2.866114in}{3.496408in}}%
\pgfpathlineto{\pgfqpoint{2.866114in}{3.493459in}}%
\pgfpathmoveto{\pgfqpoint{2.866114in}{3.484611in}}%
\pgfpathlineto{\pgfqpoint{2.866114in}{3.484611in}}%
\pgfpathlineto{\pgfqpoint{2.866114in}{3.487560in}}%
\pgfpathlineto{\pgfqpoint{2.870655in}{3.487560in}}%
\pgfpathlineto{\pgfqpoint{2.870655in}{3.484611in}}%
\pgfpathmoveto{\pgfqpoint{2.866114in}{3.487560in}}%
\pgfpathlineto{\pgfqpoint{2.866114in}{3.487560in}}%
\pgfpathlineto{\pgfqpoint{2.866114in}{3.490510in}}%
\pgfpathlineto{\pgfqpoint{2.870655in}{3.490510in}}%
\pgfpathlineto{\pgfqpoint{2.870655in}{3.487560in}}%
\pgfpathmoveto{\pgfqpoint{2.875196in}{3.472814in}}%
\pgfpathlineto{\pgfqpoint{2.875196in}{3.472814in}}%
\pgfpathlineto{\pgfqpoint{2.875196in}{3.475763in}}%
\pgfpathlineto{\pgfqpoint{2.879737in}{3.475763in}}%
\pgfpathlineto{\pgfqpoint{2.879737in}{3.472814in}}%
\pgfpathmoveto{\pgfqpoint{2.875196in}{3.475763in}}%
\pgfpathlineto{\pgfqpoint{2.875196in}{3.475763in}}%
\pgfpathlineto{\pgfqpoint{2.875196in}{3.478712in}}%
\pgfpathlineto{\pgfqpoint{2.879737in}{3.478712in}}%
\pgfpathlineto{\pgfqpoint{2.879737in}{3.475763in}}%
\pgfpathmoveto{\pgfqpoint{2.857033in}{3.496408in}}%
\pgfpathlineto{\pgfqpoint{2.857033in}{3.496408in}}%
\pgfpathlineto{\pgfqpoint{2.857033in}{3.499357in}}%
\pgfpathlineto{\pgfqpoint{2.861573in}{3.499357in}}%
\pgfpathlineto{\pgfqpoint{2.861573in}{3.496408in}}%
\pgfpathmoveto{\pgfqpoint{2.857033in}{3.499357in}}%
\pgfpathlineto{\pgfqpoint{2.857033in}{3.499357in}}%
\pgfpathlineto{\pgfqpoint{2.857033in}{3.502307in}}%
\pgfpathlineto{\pgfqpoint{2.861573in}{3.502307in}}%
\pgfpathlineto{\pgfqpoint{2.861573in}{3.499357in}}%
\pgfpathmoveto{\pgfqpoint{3.070461in}{3.216230in}}%
\pgfpathlineto{\pgfqpoint{3.070461in}{3.216230in}}%
\pgfpathlineto{\pgfqpoint{3.070461in}{3.219179in}}%
\pgfpathlineto{\pgfqpoint{3.075002in}{3.219179in}}%
\pgfpathlineto{\pgfqpoint{3.075002in}{3.216230in}}%
\pgfpathmoveto{\pgfqpoint{3.065920in}{3.222128in}}%
\pgfpathlineto{\pgfqpoint{3.065920in}{3.222128in}}%
\pgfpathlineto{\pgfqpoint{3.065920in}{3.225077in}}%
\pgfpathlineto{\pgfqpoint{3.070461in}{3.225077in}}%
\pgfpathlineto{\pgfqpoint{3.070461in}{3.222128in}}%
\pgfpathmoveto{\pgfqpoint{3.070461in}{3.219179in}}%
\pgfpathlineto{\pgfqpoint{3.070461in}{3.219179in}}%
\pgfpathlineto{\pgfqpoint{3.070461in}{3.222128in}}%
\pgfpathlineto{\pgfqpoint{3.075002in}{3.222128in}}%
\pgfpathlineto{\pgfqpoint{3.075002in}{3.219179in}}%
\pgfpathmoveto{\pgfqpoint{3.070461in}{3.222128in}}%
\pgfpathlineto{\pgfqpoint{3.070461in}{3.222128in}}%
\pgfpathlineto{\pgfqpoint{3.070461in}{3.225077in}}%
\pgfpathlineto{\pgfqpoint{3.075002in}{3.225077in}}%
\pgfpathlineto{\pgfqpoint{3.075002in}{3.222128in}}%
\pgfpathmoveto{\pgfqpoint{3.061379in}{3.228027in}}%
\pgfpathlineto{\pgfqpoint{3.061379in}{3.228027in}}%
\pgfpathlineto{\pgfqpoint{3.061379in}{3.230976in}}%
\pgfpathlineto{\pgfqpoint{3.065920in}{3.230976in}}%
\pgfpathlineto{\pgfqpoint{3.065920in}{3.228027in}}%
\pgfpathmoveto{\pgfqpoint{3.056838in}{3.233925in}}%
\pgfpathlineto{\pgfqpoint{3.056838in}{3.233925in}}%
\pgfpathlineto{\pgfqpoint{3.056838in}{3.236874in}}%
\pgfpathlineto{\pgfqpoint{3.061379in}{3.236874in}}%
\pgfpathlineto{\pgfqpoint{3.061379in}{3.233925in}}%
\pgfpathmoveto{\pgfqpoint{3.061379in}{3.230976in}}%
\pgfpathlineto{\pgfqpoint{3.061379in}{3.230976in}}%
\pgfpathlineto{\pgfqpoint{3.061379in}{3.233925in}}%
\pgfpathlineto{\pgfqpoint{3.065920in}{3.233925in}}%
\pgfpathlineto{\pgfqpoint{3.065920in}{3.230976in}}%
\pgfpathmoveto{\pgfqpoint{3.061379in}{3.233925in}}%
\pgfpathlineto{\pgfqpoint{3.061379in}{3.233925in}}%
\pgfpathlineto{\pgfqpoint{3.061379in}{3.236874in}}%
\pgfpathlineto{\pgfqpoint{3.065920in}{3.236874in}}%
\pgfpathlineto{\pgfqpoint{3.065920in}{3.233925in}}%
\pgfpathmoveto{\pgfqpoint{3.065920in}{3.225077in}}%
\pgfpathlineto{\pgfqpoint{3.065920in}{3.225077in}}%
\pgfpathlineto{\pgfqpoint{3.065920in}{3.228027in}}%
\pgfpathlineto{\pgfqpoint{3.070461in}{3.228027in}}%
\pgfpathlineto{\pgfqpoint{3.070461in}{3.225077in}}%
\pgfpathmoveto{\pgfqpoint{3.065920in}{3.228027in}}%
\pgfpathlineto{\pgfqpoint{3.065920in}{3.228027in}}%
\pgfpathlineto{\pgfqpoint{3.065920in}{3.230976in}}%
\pgfpathlineto{\pgfqpoint{3.070461in}{3.230976in}}%
\pgfpathlineto{\pgfqpoint{3.070461in}{3.228027in}}%
\pgfpathmoveto{\pgfqpoint{2.997803in}{3.310604in}}%
\pgfpathlineto{\pgfqpoint{2.997803in}{3.310604in}}%
\pgfpathlineto{\pgfqpoint{2.997803in}{3.313553in}}%
\pgfpathlineto{\pgfqpoint{3.002344in}{3.313553in}}%
\pgfpathlineto{\pgfqpoint{3.002344in}{3.310604in}}%
\pgfpathmoveto{\pgfqpoint{2.993262in}{3.316503in}}%
\pgfpathlineto{\pgfqpoint{2.993262in}{3.316503in}}%
\pgfpathlineto{\pgfqpoint{2.993262in}{3.319452in}}%
\pgfpathlineto{\pgfqpoint{2.997803in}{3.319452in}}%
\pgfpathlineto{\pgfqpoint{2.997803in}{3.316503in}}%
\pgfpathmoveto{\pgfqpoint{2.997803in}{3.313553in}}%
\pgfpathlineto{\pgfqpoint{2.997803in}{3.313553in}}%
\pgfpathlineto{\pgfqpoint{2.997803in}{3.316503in}}%
\pgfpathlineto{\pgfqpoint{3.002344in}{3.316503in}}%
\pgfpathlineto{\pgfqpoint{3.002344in}{3.313553in}}%
\pgfpathmoveto{\pgfqpoint{2.997803in}{3.316503in}}%
\pgfpathlineto{\pgfqpoint{2.997803in}{3.316503in}}%
\pgfpathlineto{\pgfqpoint{2.997803in}{3.319452in}}%
\pgfpathlineto{\pgfqpoint{3.002344in}{3.319452in}}%
\pgfpathlineto{\pgfqpoint{3.002344in}{3.316503in}}%
\pgfpathmoveto{\pgfqpoint{2.988721in}{3.322401in}}%
\pgfpathlineto{\pgfqpoint{2.988721in}{3.322401in}}%
\pgfpathlineto{\pgfqpoint{2.988721in}{3.325350in}}%
\pgfpathlineto{\pgfqpoint{2.993262in}{3.325350in}}%
\pgfpathlineto{\pgfqpoint{2.993262in}{3.322401in}}%
\pgfpathmoveto{\pgfqpoint{2.984180in}{3.328299in}}%
\pgfpathlineto{\pgfqpoint{2.984180in}{3.328299in}}%
\pgfpathlineto{\pgfqpoint{2.984180in}{3.331249in}}%
\pgfpathlineto{\pgfqpoint{2.988721in}{3.331249in}}%
\pgfpathlineto{\pgfqpoint{2.988721in}{3.328299in}}%
\pgfpathmoveto{\pgfqpoint{2.988721in}{3.325350in}}%
\pgfpathlineto{\pgfqpoint{2.988721in}{3.325350in}}%
\pgfpathlineto{\pgfqpoint{2.988721in}{3.328299in}}%
\pgfpathlineto{\pgfqpoint{2.993262in}{3.328299in}}%
\pgfpathlineto{\pgfqpoint{2.993262in}{3.325350in}}%
\pgfpathmoveto{\pgfqpoint{2.988721in}{3.328299in}}%
\pgfpathlineto{\pgfqpoint{2.988721in}{3.328299in}}%
\pgfpathlineto{\pgfqpoint{2.988721in}{3.331249in}}%
\pgfpathlineto{\pgfqpoint{2.993262in}{3.331249in}}%
\pgfpathlineto{\pgfqpoint{2.993262in}{3.328299in}}%
\pgfpathmoveto{\pgfqpoint{2.993262in}{3.319452in}}%
\pgfpathlineto{\pgfqpoint{2.993262in}{3.319452in}}%
\pgfpathlineto{\pgfqpoint{2.993262in}{3.322401in}}%
\pgfpathlineto{\pgfqpoint{2.997803in}{3.322401in}}%
\pgfpathlineto{\pgfqpoint{2.997803in}{3.319452in}}%
\pgfpathmoveto{\pgfqpoint{2.993262in}{3.322401in}}%
\pgfpathlineto{\pgfqpoint{2.993262in}{3.322401in}}%
\pgfpathlineto{\pgfqpoint{2.993262in}{3.325350in}}%
\pgfpathlineto{\pgfqpoint{2.997803in}{3.325350in}}%
\pgfpathlineto{\pgfqpoint{2.997803in}{3.322401in}}%
\pgfpathmoveto{\pgfqpoint{3.034132in}{3.263417in}}%
\pgfpathlineto{\pgfqpoint{3.034132in}{3.263417in}}%
\pgfpathlineto{\pgfqpoint{3.034132in}{3.266366in}}%
\pgfpathlineto{\pgfqpoint{3.038673in}{3.266366in}}%
\pgfpathlineto{\pgfqpoint{3.038673in}{3.263417in}}%
\pgfpathmoveto{\pgfqpoint{3.029591in}{3.269315in}}%
\pgfpathlineto{\pgfqpoint{3.029591in}{3.269315in}}%
\pgfpathlineto{\pgfqpoint{3.029591in}{3.272265in}}%
\pgfpathlineto{\pgfqpoint{3.034132in}{3.272265in}}%
\pgfpathlineto{\pgfqpoint{3.034132in}{3.269315in}}%
\pgfpathmoveto{\pgfqpoint{3.034132in}{3.266366in}}%
\pgfpathlineto{\pgfqpoint{3.034132in}{3.266366in}}%
\pgfpathlineto{\pgfqpoint{3.034132in}{3.269315in}}%
\pgfpathlineto{\pgfqpoint{3.038673in}{3.269315in}}%
\pgfpathlineto{\pgfqpoint{3.038673in}{3.266366in}}%
\pgfpathmoveto{\pgfqpoint{3.034132in}{3.269315in}}%
\pgfpathlineto{\pgfqpoint{3.034132in}{3.269315in}}%
\pgfpathlineto{\pgfqpoint{3.034132in}{3.272265in}}%
\pgfpathlineto{\pgfqpoint{3.038673in}{3.272265in}}%
\pgfpathlineto{\pgfqpoint{3.038673in}{3.269315in}}%
\pgfpathmoveto{\pgfqpoint{3.025050in}{3.275214in}}%
\pgfpathlineto{\pgfqpoint{3.025050in}{3.275214in}}%
\pgfpathlineto{\pgfqpoint{3.025050in}{3.278163in}}%
\pgfpathlineto{\pgfqpoint{3.029591in}{3.278163in}}%
\pgfpathlineto{\pgfqpoint{3.029591in}{3.275214in}}%
\pgfpathmoveto{\pgfqpoint{3.020509in}{3.281112in}}%
\pgfpathlineto{\pgfqpoint{3.020509in}{3.281112in}}%
\pgfpathlineto{\pgfqpoint{3.020509in}{3.284061in}}%
\pgfpathlineto{\pgfqpoint{3.025050in}{3.284061in}}%
\pgfpathlineto{\pgfqpoint{3.025050in}{3.281112in}}%
\pgfpathmoveto{\pgfqpoint{3.025050in}{3.278163in}}%
\pgfpathlineto{\pgfqpoint{3.025050in}{3.278163in}}%
\pgfpathlineto{\pgfqpoint{3.025050in}{3.281112in}}%
\pgfpathlineto{\pgfqpoint{3.029591in}{3.281112in}}%
\pgfpathlineto{\pgfqpoint{3.029591in}{3.278163in}}%
\pgfpathmoveto{\pgfqpoint{3.025050in}{3.281112in}}%
\pgfpathlineto{\pgfqpoint{3.025050in}{3.281112in}}%
\pgfpathlineto{\pgfqpoint{3.025050in}{3.284061in}}%
\pgfpathlineto{\pgfqpoint{3.029591in}{3.284061in}}%
\pgfpathlineto{\pgfqpoint{3.029591in}{3.281112in}}%
\pgfpathmoveto{\pgfqpoint{3.029591in}{3.272265in}}%
\pgfpathlineto{\pgfqpoint{3.029591in}{3.272265in}}%
\pgfpathlineto{\pgfqpoint{3.029591in}{3.275214in}}%
\pgfpathlineto{\pgfqpoint{3.034132in}{3.275214in}}%
\pgfpathlineto{\pgfqpoint{3.034132in}{3.272265in}}%
\pgfpathmoveto{\pgfqpoint{3.029591in}{3.275214in}}%
\pgfpathlineto{\pgfqpoint{3.029591in}{3.275214in}}%
\pgfpathlineto{\pgfqpoint{3.029591in}{3.278163in}}%
\pgfpathlineto{\pgfqpoint{3.034132in}{3.278163in}}%
\pgfpathlineto{\pgfqpoint{3.034132in}{3.275214in}}%
\pgfpathmoveto{\pgfqpoint{3.052297in}{3.239824in}}%
\pgfpathlineto{\pgfqpoint{3.052297in}{3.239824in}}%
\pgfpathlineto{\pgfqpoint{3.052297in}{3.242773in}}%
\pgfpathlineto{\pgfqpoint{3.056838in}{3.242773in}}%
\pgfpathlineto{\pgfqpoint{3.056838in}{3.239824in}}%
\pgfpathmoveto{\pgfqpoint{3.047756in}{3.245722in}}%
\pgfpathlineto{\pgfqpoint{3.047756in}{3.245722in}}%
\pgfpathlineto{\pgfqpoint{3.047756in}{3.248671in}}%
\pgfpathlineto{\pgfqpoint{3.052297in}{3.248671in}}%
\pgfpathlineto{\pgfqpoint{3.052297in}{3.245722in}}%
\pgfpathmoveto{\pgfqpoint{3.052297in}{3.242773in}}%
\pgfpathlineto{\pgfqpoint{3.052297in}{3.242773in}}%
\pgfpathlineto{\pgfqpoint{3.052297in}{3.245722in}}%
\pgfpathlineto{\pgfqpoint{3.056838in}{3.245722in}}%
\pgfpathlineto{\pgfqpoint{3.056838in}{3.242773in}}%
\pgfpathmoveto{\pgfqpoint{3.052297in}{3.245722in}}%
\pgfpathlineto{\pgfqpoint{3.052297in}{3.245722in}}%
\pgfpathlineto{\pgfqpoint{3.052297in}{3.248671in}}%
\pgfpathlineto{\pgfqpoint{3.056838in}{3.248671in}}%
\pgfpathlineto{\pgfqpoint{3.056838in}{3.245722in}}%
\pgfpathmoveto{\pgfqpoint{3.043214in}{3.251620in}}%
\pgfpathlineto{\pgfqpoint{3.043214in}{3.251620in}}%
\pgfpathlineto{\pgfqpoint{3.043214in}{3.254569in}}%
\pgfpathlineto{\pgfqpoint{3.047756in}{3.254569in}}%
\pgfpathlineto{\pgfqpoint{3.047756in}{3.251620in}}%
\pgfpathmoveto{\pgfqpoint{3.038673in}{3.257519in}}%
\pgfpathlineto{\pgfqpoint{3.038673in}{3.257519in}}%
\pgfpathlineto{\pgfqpoint{3.038673in}{3.260468in}}%
\pgfpathlineto{\pgfqpoint{3.043214in}{3.260468in}}%
\pgfpathlineto{\pgfqpoint{3.043214in}{3.257519in}}%
\pgfpathmoveto{\pgfqpoint{3.043214in}{3.254569in}}%
\pgfpathlineto{\pgfqpoint{3.043214in}{3.254569in}}%
\pgfpathlineto{\pgfqpoint{3.043214in}{3.257519in}}%
\pgfpathlineto{\pgfqpoint{3.047756in}{3.257519in}}%
\pgfpathlineto{\pgfqpoint{3.047756in}{3.254569in}}%
\pgfpathmoveto{\pgfqpoint{3.043214in}{3.257519in}}%
\pgfpathlineto{\pgfqpoint{3.043214in}{3.257519in}}%
\pgfpathlineto{\pgfqpoint{3.043214in}{3.260468in}}%
\pgfpathlineto{\pgfqpoint{3.047756in}{3.260468in}}%
\pgfpathlineto{\pgfqpoint{3.047756in}{3.257519in}}%
\pgfpathmoveto{\pgfqpoint{3.047756in}{3.248671in}}%
\pgfpathlineto{\pgfqpoint{3.047756in}{3.248671in}}%
\pgfpathlineto{\pgfqpoint{3.047756in}{3.251620in}}%
\pgfpathlineto{\pgfqpoint{3.052297in}{3.251620in}}%
\pgfpathlineto{\pgfqpoint{3.052297in}{3.248671in}}%
\pgfpathmoveto{\pgfqpoint{3.047756in}{3.251620in}}%
\pgfpathlineto{\pgfqpoint{3.047756in}{3.251620in}}%
\pgfpathlineto{\pgfqpoint{3.047756in}{3.254569in}}%
\pgfpathlineto{\pgfqpoint{3.052297in}{3.254569in}}%
\pgfpathlineto{\pgfqpoint{3.052297in}{3.251620in}}%
\pgfpathmoveto{\pgfqpoint{3.056838in}{3.236874in}}%
\pgfpathlineto{\pgfqpoint{3.056838in}{3.236874in}}%
\pgfpathlineto{\pgfqpoint{3.056838in}{3.239824in}}%
\pgfpathlineto{\pgfqpoint{3.061379in}{3.239824in}}%
\pgfpathlineto{\pgfqpoint{3.061379in}{3.236874in}}%
\pgfpathmoveto{\pgfqpoint{3.056838in}{3.239824in}}%
\pgfpathlineto{\pgfqpoint{3.056838in}{3.239824in}}%
\pgfpathlineto{\pgfqpoint{3.056838in}{3.242773in}}%
\pgfpathlineto{\pgfqpoint{3.061379in}{3.242773in}}%
\pgfpathlineto{\pgfqpoint{3.061379in}{3.239824in}}%
\pgfpathmoveto{\pgfqpoint{3.038673in}{3.260468in}}%
\pgfpathlineto{\pgfqpoint{3.038673in}{3.260468in}}%
\pgfpathlineto{\pgfqpoint{3.038673in}{3.263417in}}%
\pgfpathlineto{\pgfqpoint{3.043214in}{3.263417in}}%
\pgfpathlineto{\pgfqpoint{3.043214in}{3.260468in}}%
\pgfpathmoveto{\pgfqpoint{3.038673in}{3.263417in}}%
\pgfpathlineto{\pgfqpoint{3.038673in}{3.263417in}}%
\pgfpathlineto{\pgfqpoint{3.038673in}{3.266366in}}%
\pgfpathlineto{\pgfqpoint{3.043214in}{3.266366in}}%
\pgfpathlineto{\pgfqpoint{3.043214in}{3.263417in}}%
\pgfpathmoveto{\pgfqpoint{3.015968in}{3.287011in}}%
\pgfpathlineto{\pgfqpoint{3.015968in}{3.287011in}}%
\pgfpathlineto{\pgfqpoint{3.015968in}{3.289960in}}%
\pgfpathlineto{\pgfqpoint{3.020509in}{3.289960in}}%
\pgfpathlineto{\pgfqpoint{3.020509in}{3.287011in}}%
\pgfpathmoveto{\pgfqpoint{3.011427in}{3.292909in}}%
\pgfpathlineto{\pgfqpoint{3.011427in}{3.292909in}}%
\pgfpathlineto{\pgfqpoint{3.011427in}{3.295858in}}%
\pgfpathlineto{\pgfqpoint{3.015968in}{3.295858in}}%
\pgfpathlineto{\pgfqpoint{3.015968in}{3.292909in}}%
\pgfpathmoveto{\pgfqpoint{3.015968in}{3.289960in}}%
\pgfpathlineto{\pgfqpoint{3.015968in}{3.289960in}}%
\pgfpathlineto{\pgfqpoint{3.015968in}{3.292909in}}%
\pgfpathlineto{\pgfqpoint{3.020509in}{3.292909in}}%
\pgfpathlineto{\pgfqpoint{3.020509in}{3.289960in}}%
\pgfpathmoveto{\pgfqpoint{3.015968in}{3.292909in}}%
\pgfpathlineto{\pgfqpoint{3.015968in}{3.292909in}}%
\pgfpathlineto{\pgfqpoint{3.015968in}{3.295858in}}%
\pgfpathlineto{\pgfqpoint{3.020509in}{3.295858in}}%
\pgfpathlineto{\pgfqpoint{3.020509in}{3.292909in}}%
\pgfpathmoveto{\pgfqpoint{3.006885in}{3.298807in}}%
\pgfpathlineto{\pgfqpoint{3.006885in}{3.298807in}}%
\pgfpathlineto{\pgfqpoint{3.006885in}{3.301757in}}%
\pgfpathlineto{\pgfqpoint{3.011427in}{3.301757in}}%
\pgfpathlineto{\pgfqpoint{3.011427in}{3.298807in}}%
\pgfpathmoveto{\pgfqpoint{3.002344in}{3.304706in}}%
\pgfpathlineto{\pgfqpoint{3.002344in}{3.304706in}}%
\pgfpathlineto{\pgfqpoint{3.002344in}{3.307655in}}%
\pgfpathlineto{\pgfqpoint{3.006885in}{3.307655in}}%
\pgfpathlineto{\pgfqpoint{3.006885in}{3.304706in}}%
\pgfpathmoveto{\pgfqpoint{3.006885in}{3.301757in}}%
\pgfpathlineto{\pgfqpoint{3.006885in}{3.301757in}}%
\pgfpathlineto{\pgfqpoint{3.006885in}{3.304706in}}%
\pgfpathlineto{\pgfqpoint{3.011427in}{3.304706in}}%
\pgfpathlineto{\pgfqpoint{3.011427in}{3.301757in}}%
\pgfpathmoveto{\pgfqpoint{3.006885in}{3.304706in}}%
\pgfpathlineto{\pgfqpoint{3.006885in}{3.304706in}}%
\pgfpathlineto{\pgfqpoint{3.006885in}{3.307655in}}%
\pgfpathlineto{\pgfqpoint{3.011427in}{3.307655in}}%
\pgfpathlineto{\pgfqpoint{3.011427in}{3.304706in}}%
\pgfpathmoveto{\pgfqpoint{3.011427in}{3.295858in}}%
\pgfpathlineto{\pgfqpoint{3.011427in}{3.295858in}}%
\pgfpathlineto{\pgfqpoint{3.011427in}{3.298807in}}%
\pgfpathlineto{\pgfqpoint{3.015968in}{3.298807in}}%
\pgfpathlineto{\pgfqpoint{3.015968in}{3.295858in}}%
\pgfpathmoveto{\pgfqpoint{3.011427in}{3.298807in}}%
\pgfpathlineto{\pgfqpoint{3.011427in}{3.298807in}}%
\pgfpathlineto{\pgfqpoint{3.011427in}{3.301757in}}%
\pgfpathlineto{\pgfqpoint{3.015968in}{3.301757in}}%
\pgfpathlineto{\pgfqpoint{3.015968in}{3.298807in}}%
\pgfpathmoveto{\pgfqpoint{3.020509in}{3.284061in}}%
\pgfpathlineto{\pgfqpoint{3.020509in}{3.284061in}}%
\pgfpathlineto{\pgfqpoint{3.020509in}{3.287011in}}%
\pgfpathlineto{\pgfqpoint{3.025050in}{3.287011in}}%
\pgfpathlineto{\pgfqpoint{3.025050in}{3.284061in}}%
\pgfpathmoveto{\pgfqpoint{3.020509in}{3.287011in}}%
\pgfpathlineto{\pgfqpoint{3.020509in}{3.287011in}}%
\pgfpathlineto{\pgfqpoint{3.020509in}{3.289960in}}%
\pgfpathlineto{\pgfqpoint{3.025050in}{3.289960in}}%
\pgfpathlineto{\pgfqpoint{3.025050in}{3.287011in}}%
\pgfpathmoveto{\pgfqpoint{3.002344in}{3.307655in}}%
\pgfpathlineto{\pgfqpoint{3.002344in}{3.307655in}}%
\pgfpathlineto{\pgfqpoint{3.002344in}{3.310604in}}%
\pgfpathlineto{\pgfqpoint{3.006885in}{3.310604in}}%
\pgfpathlineto{\pgfqpoint{3.006885in}{3.307655in}}%
\pgfpathmoveto{\pgfqpoint{3.002344in}{3.310604in}}%
\pgfpathlineto{\pgfqpoint{3.002344in}{3.310604in}}%
\pgfpathlineto{\pgfqpoint{3.002344in}{3.313553in}}%
\pgfpathlineto{\pgfqpoint{3.006885in}{3.313553in}}%
\pgfpathlineto{\pgfqpoint{3.006885in}{3.310604in}}%
\pgfpathmoveto{\pgfqpoint{2.961474in}{3.357792in}}%
\pgfpathlineto{\pgfqpoint{2.961474in}{3.357792in}}%
\pgfpathlineto{\pgfqpoint{2.961474in}{3.360741in}}%
\pgfpathlineto{\pgfqpoint{2.966015in}{3.360741in}}%
\pgfpathlineto{\pgfqpoint{2.966015in}{3.357792in}}%
\pgfpathmoveto{\pgfqpoint{2.956933in}{3.363691in}}%
\pgfpathlineto{\pgfqpoint{2.956933in}{3.363691in}}%
\pgfpathlineto{\pgfqpoint{2.956933in}{3.366640in}}%
\pgfpathlineto{\pgfqpoint{2.961474in}{3.366640in}}%
\pgfpathlineto{\pgfqpoint{2.961474in}{3.363691in}}%
\pgfpathmoveto{\pgfqpoint{2.961474in}{3.360741in}}%
\pgfpathlineto{\pgfqpoint{2.961474in}{3.360741in}}%
\pgfpathlineto{\pgfqpoint{2.961474in}{3.363691in}}%
\pgfpathlineto{\pgfqpoint{2.966015in}{3.363691in}}%
\pgfpathlineto{\pgfqpoint{2.966015in}{3.360741in}}%
\pgfpathmoveto{\pgfqpoint{2.961474in}{3.363691in}}%
\pgfpathlineto{\pgfqpoint{2.961474in}{3.363691in}}%
\pgfpathlineto{\pgfqpoint{2.961474in}{3.366640in}}%
\pgfpathlineto{\pgfqpoint{2.966015in}{3.366640in}}%
\pgfpathlineto{\pgfqpoint{2.966015in}{3.363691in}}%
\pgfpathmoveto{\pgfqpoint{2.952392in}{3.369589in}}%
\pgfpathlineto{\pgfqpoint{2.952392in}{3.369589in}}%
\pgfpathlineto{\pgfqpoint{2.952392in}{3.372538in}}%
\pgfpathlineto{\pgfqpoint{2.956933in}{3.372538in}}%
\pgfpathlineto{\pgfqpoint{2.956933in}{3.369589in}}%
\pgfpathmoveto{\pgfqpoint{2.947851in}{3.375488in}}%
\pgfpathlineto{\pgfqpoint{2.947851in}{3.375488in}}%
\pgfpathlineto{\pgfqpoint{2.947851in}{3.378437in}}%
\pgfpathlineto{\pgfqpoint{2.952392in}{3.378437in}}%
\pgfpathlineto{\pgfqpoint{2.952392in}{3.375488in}}%
\pgfpathmoveto{\pgfqpoint{2.952392in}{3.372538in}}%
\pgfpathlineto{\pgfqpoint{2.952392in}{3.372538in}}%
\pgfpathlineto{\pgfqpoint{2.952392in}{3.375488in}}%
\pgfpathlineto{\pgfqpoint{2.956933in}{3.375488in}}%
\pgfpathlineto{\pgfqpoint{2.956933in}{3.372538in}}%
\pgfpathmoveto{\pgfqpoint{2.952392in}{3.375488in}}%
\pgfpathlineto{\pgfqpoint{2.952392in}{3.375488in}}%
\pgfpathlineto{\pgfqpoint{2.952392in}{3.378437in}}%
\pgfpathlineto{\pgfqpoint{2.956933in}{3.378437in}}%
\pgfpathlineto{\pgfqpoint{2.956933in}{3.375488in}}%
\pgfpathmoveto{\pgfqpoint{2.956933in}{3.366640in}}%
\pgfpathlineto{\pgfqpoint{2.956933in}{3.366640in}}%
\pgfpathlineto{\pgfqpoint{2.956933in}{3.369589in}}%
\pgfpathlineto{\pgfqpoint{2.961474in}{3.369589in}}%
\pgfpathlineto{\pgfqpoint{2.961474in}{3.366640in}}%
\pgfpathmoveto{\pgfqpoint{2.956933in}{3.369589in}}%
\pgfpathlineto{\pgfqpoint{2.956933in}{3.369589in}}%
\pgfpathlineto{\pgfqpoint{2.956933in}{3.372538in}}%
\pgfpathlineto{\pgfqpoint{2.961474in}{3.372538in}}%
\pgfpathlineto{\pgfqpoint{2.961474in}{3.369589in}}%
\pgfpathmoveto{\pgfqpoint{2.979639in}{3.334198in}}%
\pgfpathlineto{\pgfqpoint{2.979639in}{3.334198in}}%
\pgfpathlineto{\pgfqpoint{2.979639in}{3.337147in}}%
\pgfpathlineto{\pgfqpoint{2.984180in}{3.337147in}}%
\pgfpathlineto{\pgfqpoint{2.984180in}{3.334198in}}%
\pgfpathmoveto{\pgfqpoint{2.975098in}{3.340096in}}%
\pgfpathlineto{\pgfqpoint{2.975098in}{3.340096in}}%
\pgfpathlineto{\pgfqpoint{2.975098in}{3.343046in}}%
\pgfpathlineto{\pgfqpoint{2.979639in}{3.343046in}}%
\pgfpathlineto{\pgfqpoint{2.979639in}{3.340096in}}%
\pgfpathmoveto{\pgfqpoint{2.979639in}{3.337147in}}%
\pgfpathlineto{\pgfqpoint{2.979639in}{3.337147in}}%
\pgfpathlineto{\pgfqpoint{2.979639in}{3.340096in}}%
\pgfpathlineto{\pgfqpoint{2.984180in}{3.340096in}}%
\pgfpathlineto{\pgfqpoint{2.984180in}{3.337147in}}%
\pgfpathmoveto{\pgfqpoint{2.979639in}{3.340096in}}%
\pgfpathlineto{\pgfqpoint{2.979639in}{3.340096in}}%
\pgfpathlineto{\pgfqpoint{2.979639in}{3.343046in}}%
\pgfpathlineto{\pgfqpoint{2.984180in}{3.343046in}}%
\pgfpathlineto{\pgfqpoint{2.984180in}{3.340096in}}%
\pgfpathmoveto{\pgfqpoint{2.970556in}{3.345995in}}%
\pgfpathlineto{\pgfqpoint{2.970556in}{3.345995in}}%
\pgfpathlineto{\pgfqpoint{2.970556in}{3.348944in}}%
\pgfpathlineto{\pgfqpoint{2.975098in}{3.348944in}}%
\pgfpathlineto{\pgfqpoint{2.975098in}{3.345995in}}%
\pgfpathmoveto{\pgfqpoint{2.966015in}{3.351894in}}%
\pgfpathlineto{\pgfqpoint{2.966015in}{3.351894in}}%
\pgfpathlineto{\pgfqpoint{2.966015in}{3.354843in}}%
\pgfpathlineto{\pgfqpoint{2.970556in}{3.354843in}}%
\pgfpathlineto{\pgfqpoint{2.970556in}{3.351894in}}%
\pgfpathmoveto{\pgfqpoint{2.970556in}{3.348944in}}%
\pgfpathlineto{\pgfqpoint{2.970556in}{3.348944in}}%
\pgfpathlineto{\pgfqpoint{2.970556in}{3.351894in}}%
\pgfpathlineto{\pgfqpoint{2.975098in}{3.351894in}}%
\pgfpathlineto{\pgfqpoint{2.975098in}{3.348944in}}%
\pgfpathmoveto{\pgfqpoint{2.970556in}{3.351894in}}%
\pgfpathlineto{\pgfqpoint{2.970556in}{3.351894in}}%
\pgfpathlineto{\pgfqpoint{2.970556in}{3.354843in}}%
\pgfpathlineto{\pgfqpoint{2.975098in}{3.354843in}}%
\pgfpathlineto{\pgfqpoint{2.975098in}{3.351894in}}%
\pgfpathmoveto{\pgfqpoint{2.975098in}{3.343046in}}%
\pgfpathlineto{\pgfqpoint{2.975098in}{3.343046in}}%
\pgfpathlineto{\pgfqpoint{2.975098in}{3.345995in}}%
\pgfpathlineto{\pgfqpoint{2.979639in}{3.345995in}}%
\pgfpathlineto{\pgfqpoint{2.979639in}{3.343046in}}%
\pgfpathmoveto{\pgfqpoint{2.975098in}{3.345995in}}%
\pgfpathlineto{\pgfqpoint{2.975098in}{3.345995in}}%
\pgfpathlineto{\pgfqpoint{2.975098in}{3.348944in}}%
\pgfpathlineto{\pgfqpoint{2.979639in}{3.348944in}}%
\pgfpathlineto{\pgfqpoint{2.979639in}{3.345995in}}%
\pgfpathmoveto{\pgfqpoint{2.984180in}{3.331249in}}%
\pgfpathlineto{\pgfqpoint{2.984180in}{3.331249in}}%
\pgfpathlineto{\pgfqpoint{2.984180in}{3.334198in}}%
\pgfpathlineto{\pgfqpoint{2.988721in}{3.334198in}}%
\pgfpathlineto{\pgfqpoint{2.988721in}{3.331249in}}%
\pgfpathmoveto{\pgfqpoint{2.984180in}{3.334198in}}%
\pgfpathlineto{\pgfqpoint{2.984180in}{3.334198in}}%
\pgfpathlineto{\pgfqpoint{2.984180in}{3.337147in}}%
\pgfpathlineto{\pgfqpoint{2.988721in}{3.337147in}}%
\pgfpathlineto{\pgfqpoint{2.988721in}{3.334198in}}%
\pgfpathmoveto{\pgfqpoint{2.966015in}{3.354843in}}%
\pgfpathlineto{\pgfqpoint{2.966015in}{3.354843in}}%
\pgfpathlineto{\pgfqpoint{2.966015in}{3.357792in}}%
\pgfpathlineto{\pgfqpoint{2.970556in}{3.357792in}}%
\pgfpathlineto{\pgfqpoint{2.970556in}{3.354843in}}%
\pgfpathmoveto{\pgfqpoint{2.966015in}{3.357792in}}%
\pgfpathlineto{\pgfqpoint{2.966015in}{3.357792in}}%
\pgfpathlineto{\pgfqpoint{2.966015in}{3.360741in}}%
\pgfpathlineto{\pgfqpoint{2.970556in}{3.360741in}}%
\pgfpathlineto{\pgfqpoint{2.970556in}{3.357792in}}%
\pgfpathmoveto{\pgfqpoint{2.943310in}{3.381386in}}%
\pgfpathlineto{\pgfqpoint{2.943310in}{3.381386in}}%
\pgfpathlineto{\pgfqpoint{2.943310in}{3.384335in}}%
\pgfpathlineto{\pgfqpoint{2.947851in}{3.384335in}}%
\pgfpathlineto{\pgfqpoint{2.947851in}{3.381386in}}%
\pgfpathmoveto{\pgfqpoint{2.938768in}{3.387285in}}%
\pgfpathlineto{\pgfqpoint{2.938768in}{3.387285in}}%
\pgfpathlineto{\pgfqpoint{2.938768in}{3.390234in}}%
\pgfpathlineto{\pgfqpoint{2.943310in}{3.390234in}}%
\pgfpathlineto{\pgfqpoint{2.943310in}{3.387285in}}%
\pgfpathmoveto{\pgfqpoint{2.943310in}{3.384335in}}%
\pgfpathlineto{\pgfqpoint{2.943310in}{3.384335in}}%
\pgfpathlineto{\pgfqpoint{2.943310in}{3.387285in}}%
\pgfpathlineto{\pgfqpoint{2.947851in}{3.387285in}}%
\pgfpathlineto{\pgfqpoint{2.947851in}{3.384335in}}%
\pgfpathmoveto{\pgfqpoint{2.943310in}{3.387285in}}%
\pgfpathlineto{\pgfqpoint{2.943310in}{3.387285in}}%
\pgfpathlineto{\pgfqpoint{2.943310in}{3.390234in}}%
\pgfpathlineto{\pgfqpoint{2.947851in}{3.390234in}}%
\pgfpathlineto{\pgfqpoint{2.947851in}{3.387285in}}%
\pgfpathmoveto{\pgfqpoint{2.934227in}{3.393183in}}%
\pgfpathlineto{\pgfqpoint{2.934227in}{3.393183in}}%
\pgfpathlineto{\pgfqpoint{2.934227in}{3.396133in}}%
\pgfpathlineto{\pgfqpoint{2.938768in}{3.396133in}}%
\pgfpathlineto{\pgfqpoint{2.938768in}{3.393183in}}%
\pgfpathmoveto{\pgfqpoint{2.929686in}{3.399082in}}%
\pgfpathlineto{\pgfqpoint{2.929686in}{3.399082in}}%
\pgfpathlineto{\pgfqpoint{2.929686in}{3.402031in}}%
\pgfpathlineto{\pgfqpoint{2.934227in}{3.402031in}}%
\pgfpathlineto{\pgfqpoint{2.934227in}{3.399082in}}%
\pgfpathmoveto{\pgfqpoint{2.934227in}{3.396133in}}%
\pgfpathlineto{\pgfqpoint{2.934227in}{3.396133in}}%
\pgfpathlineto{\pgfqpoint{2.934227in}{3.399082in}}%
\pgfpathlineto{\pgfqpoint{2.938768in}{3.399082in}}%
\pgfpathlineto{\pgfqpoint{2.938768in}{3.396133in}}%
\pgfpathmoveto{\pgfqpoint{2.934227in}{3.399082in}}%
\pgfpathlineto{\pgfqpoint{2.934227in}{3.399082in}}%
\pgfpathlineto{\pgfqpoint{2.934227in}{3.402031in}}%
\pgfpathlineto{\pgfqpoint{2.938768in}{3.402031in}}%
\pgfpathlineto{\pgfqpoint{2.938768in}{3.399082in}}%
\pgfpathmoveto{\pgfqpoint{2.938768in}{3.390234in}}%
\pgfpathlineto{\pgfqpoint{2.938768in}{3.390234in}}%
\pgfpathlineto{\pgfqpoint{2.938768in}{3.393183in}}%
\pgfpathlineto{\pgfqpoint{2.943310in}{3.393183in}}%
\pgfpathlineto{\pgfqpoint{2.943310in}{3.390234in}}%
\pgfpathmoveto{\pgfqpoint{2.938768in}{3.393183in}}%
\pgfpathlineto{\pgfqpoint{2.938768in}{3.393183in}}%
\pgfpathlineto{\pgfqpoint{2.938768in}{3.396133in}}%
\pgfpathlineto{\pgfqpoint{2.943310in}{3.396133in}}%
\pgfpathlineto{\pgfqpoint{2.943310in}{3.393183in}}%
\pgfpathmoveto{\pgfqpoint{2.947851in}{3.378437in}}%
\pgfpathlineto{\pgfqpoint{2.947851in}{3.378437in}}%
\pgfpathlineto{\pgfqpoint{2.947851in}{3.381386in}}%
\pgfpathlineto{\pgfqpoint{2.952392in}{3.381386in}}%
\pgfpathlineto{\pgfqpoint{2.952392in}{3.378437in}}%
\pgfpathmoveto{\pgfqpoint{2.947851in}{3.381386in}}%
\pgfpathlineto{\pgfqpoint{2.947851in}{3.381386in}}%
\pgfpathlineto{\pgfqpoint{2.947851in}{3.384335in}}%
\pgfpathlineto{\pgfqpoint{2.952392in}{3.384335in}}%
\pgfpathlineto{\pgfqpoint{2.952392in}{3.381386in}}%
\pgfpathmoveto{\pgfqpoint{2.929686in}{3.402031in}}%
\pgfpathlineto{\pgfqpoint{2.929686in}{3.402031in}}%
\pgfpathlineto{\pgfqpoint{2.929686in}{3.404980in}}%
\pgfpathlineto{\pgfqpoint{2.934227in}{3.404980in}}%
\pgfpathlineto{\pgfqpoint{2.934227in}{3.402031in}}%
\pgfpathmoveto{\pgfqpoint{2.929686in}{3.404980in}}%
\pgfpathlineto{\pgfqpoint{2.929686in}{3.404980in}}%
\pgfpathlineto{\pgfqpoint{2.929686in}{3.407930in}}%
\pgfpathlineto{\pgfqpoint{2.934227in}{3.407930in}}%
\pgfpathlineto{\pgfqpoint{2.934227in}{3.404980in}}%
\pgfpathmoveto{\pgfqpoint{3.215776in}{3.027482in}}%
\pgfpathlineto{\pgfqpoint{3.215776in}{3.027482in}}%
\pgfpathlineto{\pgfqpoint{3.215776in}{3.030431in}}%
\pgfpathlineto{\pgfqpoint{3.220317in}{3.030431in}}%
\pgfpathlineto{\pgfqpoint{3.220317in}{3.027482in}}%
\pgfpathmoveto{\pgfqpoint{3.211235in}{3.033381in}}%
\pgfpathlineto{\pgfqpoint{3.211235in}{3.033381in}}%
\pgfpathlineto{\pgfqpoint{3.211235in}{3.036330in}}%
\pgfpathlineto{\pgfqpoint{3.215776in}{3.036330in}}%
\pgfpathlineto{\pgfqpoint{3.215776in}{3.033381in}}%
\pgfpathmoveto{\pgfqpoint{3.215776in}{3.030431in}}%
\pgfpathlineto{\pgfqpoint{3.215776in}{3.030431in}}%
\pgfpathlineto{\pgfqpoint{3.215776in}{3.033381in}}%
\pgfpathlineto{\pgfqpoint{3.220317in}{3.033381in}}%
\pgfpathlineto{\pgfqpoint{3.220317in}{3.030431in}}%
\pgfpathmoveto{\pgfqpoint{3.215776in}{3.033381in}}%
\pgfpathlineto{\pgfqpoint{3.215776in}{3.033381in}}%
\pgfpathlineto{\pgfqpoint{3.215776in}{3.036330in}}%
\pgfpathlineto{\pgfqpoint{3.220317in}{3.036330in}}%
\pgfpathlineto{\pgfqpoint{3.220317in}{3.033381in}}%
\pgfpathmoveto{\pgfqpoint{3.206694in}{3.039279in}}%
\pgfpathlineto{\pgfqpoint{3.206694in}{3.039279in}}%
\pgfpathlineto{\pgfqpoint{3.206694in}{3.042229in}}%
\pgfpathlineto{\pgfqpoint{3.211235in}{3.042229in}}%
\pgfpathlineto{\pgfqpoint{3.211235in}{3.039279in}}%
\pgfpathmoveto{\pgfqpoint{3.202152in}{3.045178in}}%
\pgfpathlineto{\pgfqpoint{3.202152in}{3.045178in}}%
\pgfpathlineto{\pgfqpoint{3.202152in}{3.048127in}}%
\pgfpathlineto{\pgfqpoint{3.206694in}{3.048127in}}%
\pgfpathlineto{\pgfqpoint{3.206694in}{3.045178in}}%
\pgfpathmoveto{\pgfqpoint{3.206694in}{3.042229in}}%
\pgfpathlineto{\pgfqpoint{3.206694in}{3.042229in}}%
\pgfpathlineto{\pgfqpoint{3.206694in}{3.045178in}}%
\pgfpathlineto{\pgfqpoint{3.211235in}{3.045178in}}%
\pgfpathlineto{\pgfqpoint{3.211235in}{3.042229in}}%
\pgfpathmoveto{\pgfqpoint{3.206694in}{3.045178in}}%
\pgfpathlineto{\pgfqpoint{3.206694in}{3.045178in}}%
\pgfpathlineto{\pgfqpoint{3.206694in}{3.048127in}}%
\pgfpathlineto{\pgfqpoint{3.211235in}{3.048127in}}%
\pgfpathlineto{\pgfqpoint{3.211235in}{3.045178in}}%
\pgfpathmoveto{\pgfqpoint{3.211235in}{3.036330in}}%
\pgfpathlineto{\pgfqpoint{3.211235in}{3.036330in}}%
\pgfpathlineto{\pgfqpoint{3.211235in}{3.039279in}}%
\pgfpathlineto{\pgfqpoint{3.215776in}{3.039279in}}%
\pgfpathlineto{\pgfqpoint{3.215776in}{3.036330in}}%
\pgfpathmoveto{\pgfqpoint{3.211235in}{3.039279in}}%
\pgfpathlineto{\pgfqpoint{3.211235in}{3.039279in}}%
\pgfpathlineto{\pgfqpoint{3.211235in}{3.042229in}}%
\pgfpathlineto{\pgfqpoint{3.215776in}{3.042229in}}%
\pgfpathlineto{\pgfqpoint{3.215776in}{3.039279in}}%
\pgfpathmoveto{\pgfqpoint{3.143119in}{3.121855in}}%
\pgfpathlineto{\pgfqpoint{3.143119in}{3.121855in}}%
\pgfpathlineto{\pgfqpoint{3.143119in}{3.124804in}}%
\pgfpathlineto{\pgfqpoint{3.147660in}{3.124804in}}%
\pgfpathlineto{\pgfqpoint{3.147660in}{3.121855in}}%
\pgfpathmoveto{\pgfqpoint{3.138577in}{3.127753in}}%
\pgfpathlineto{\pgfqpoint{3.138577in}{3.127753in}}%
\pgfpathlineto{\pgfqpoint{3.138577in}{3.130702in}}%
\pgfpathlineto{\pgfqpoint{3.143119in}{3.130702in}}%
\pgfpathlineto{\pgfqpoint{3.143119in}{3.127753in}}%
\pgfpathmoveto{\pgfqpoint{3.143119in}{3.124804in}}%
\pgfpathlineto{\pgfqpoint{3.143119in}{3.124804in}}%
\pgfpathlineto{\pgfqpoint{3.143119in}{3.127753in}}%
\pgfpathlineto{\pgfqpoint{3.147660in}{3.127753in}}%
\pgfpathlineto{\pgfqpoint{3.147660in}{3.124804in}}%
\pgfpathmoveto{\pgfqpoint{3.143119in}{3.127753in}}%
\pgfpathlineto{\pgfqpoint{3.143119in}{3.127753in}}%
\pgfpathlineto{\pgfqpoint{3.143119in}{3.130702in}}%
\pgfpathlineto{\pgfqpoint{3.147660in}{3.130702in}}%
\pgfpathlineto{\pgfqpoint{3.147660in}{3.127753in}}%
\pgfpathmoveto{\pgfqpoint{3.134036in}{3.133651in}}%
\pgfpathlineto{\pgfqpoint{3.134036in}{3.133651in}}%
\pgfpathlineto{\pgfqpoint{3.134036in}{3.136600in}}%
\pgfpathlineto{\pgfqpoint{3.138577in}{3.136600in}}%
\pgfpathlineto{\pgfqpoint{3.138577in}{3.133651in}}%
\pgfpathmoveto{\pgfqpoint{3.129495in}{3.139550in}}%
\pgfpathlineto{\pgfqpoint{3.129495in}{3.139550in}}%
\pgfpathlineto{\pgfqpoint{3.129495in}{3.142499in}}%
\pgfpathlineto{\pgfqpoint{3.134036in}{3.142499in}}%
\pgfpathlineto{\pgfqpoint{3.134036in}{3.139550in}}%
\pgfpathmoveto{\pgfqpoint{3.134036in}{3.136600in}}%
\pgfpathlineto{\pgfqpoint{3.134036in}{3.136600in}}%
\pgfpathlineto{\pgfqpoint{3.134036in}{3.139550in}}%
\pgfpathlineto{\pgfqpoint{3.138577in}{3.139550in}}%
\pgfpathlineto{\pgfqpoint{3.138577in}{3.136600in}}%
\pgfpathmoveto{\pgfqpoint{3.134036in}{3.139550in}}%
\pgfpathlineto{\pgfqpoint{3.134036in}{3.139550in}}%
\pgfpathlineto{\pgfqpoint{3.134036in}{3.142499in}}%
\pgfpathlineto{\pgfqpoint{3.138577in}{3.142499in}}%
\pgfpathlineto{\pgfqpoint{3.138577in}{3.139550in}}%
\pgfpathmoveto{\pgfqpoint{3.138577in}{3.130702in}}%
\pgfpathlineto{\pgfqpoint{3.138577in}{3.130702in}}%
\pgfpathlineto{\pgfqpoint{3.138577in}{3.133651in}}%
\pgfpathlineto{\pgfqpoint{3.143119in}{3.133651in}}%
\pgfpathlineto{\pgfqpoint{3.143119in}{3.130702in}}%
\pgfpathmoveto{\pgfqpoint{3.138577in}{3.133651in}}%
\pgfpathlineto{\pgfqpoint{3.138577in}{3.133651in}}%
\pgfpathlineto{\pgfqpoint{3.138577in}{3.136600in}}%
\pgfpathlineto{\pgfqpoint{3.143119in}{3.136600in}}%
\pgfpathlineto{\pgfqpoint{3.143119in}{3.133651in}}%
\pgfpathmoveto{\pgfqpoint{3.179447in}{3.074669in}}%
\pgfpathlineto{\pgfqpoint{3.179447in}{3.074669in}}%
\pgfpathlineto{\pgfqpoint{3.179447in}{3.077618in}}%
\pgfpathlineto{\pgfqpoint{3.183988in}{3.077618in}}%
\pgfpathlineto{\pgfqpoint{3.183988in}{3.074669in}}%
\pgfpathmoveto{\pgfqpoint{3.174906in}{3.080567in}}%
\pgfpathlineto{\pgfqpoint{3.174906in}{3.080567in}}%
\pgfpathlineto{\pgfqpoint{3.174906in}{3.083517in}}%
\pgfpathlineto{\pgfqpoint{3.179447in}{3.083517in}}%
\pgfpathlineto{\pgfqpoint{3.179447in}{3.080567in}}%
\pgfpathmoveto{\pgfqpoint{3.179447in}{3.077618in}}%
\pgfpathlineto{\pgfqpoint{3.179447in}{3.077618in}}%
\pgfpathlineto{\pgfqpoint{3.179447in}{3.080567in}}%
\pgfpathlineto{\pgfqpoint{3.183988in}{3.080567in}}%
\pgfpathlineto{\pgfqpoint{3.183988in}{3.077618in}}%
\pgfpathmoveto{\pgfqpoint{3.179447in}{3.080567in}}%
\pgfpathlineto{\pgfqpoint{3.179447in}{3.080567in}}%
\pgfpathlineto{\pgfqpoint{3.179447in}{3.083517in}}%
\pgfpathlineto{\pgfqpoint{3.183988in}{3.083517in}}%
\pgfpathlineto{\pgfqpoint{3.183988in}{3.080567in}}%
\pgfpathmoveto{\pgfqpoint{3.170365in}{3.086466in}}%
\pgfpathlineto{\pgfqpoint{3.170365in}{3.086466in}}%
\pgfpathlineto{\pgfqpoint{3.170365in}{3.089415in}}%
\pgfpathlineto{\pgfqpoint{3.174906in}{3.089415in}}%
\pgfpathlineto{\pgfqpoint{3.174906in}{3.086466in}}%
\pgfpathmoveto{\pgfqpoint{3.165824in}{3.092364in}}%
\pgfpathlineto{\pgfqpoint{3.165824in}{3.092364in}}%
\pgfpathlineto{\pgfqpoint{3.165824in}{3.095313in}}%
\pgfpathlineto{\pgfqpoint{3.170365in}{3.095313in}}%
\pgfpathlineto{\pgfqpoint{3.170365in}{3.092364in}}%
\pgfpathmoveto{\pgfqpoint{3.170365in}{3.089415in}}%
\pgfpathlineto{\pgfqpoint{3.170365in}{3.089415in}}%
\pgfpathlineto{\pgfqpoint{3.170365in}{3.092364in}}%
\pgfpathlineto{\pgfqpoint{3.174906in}{3.092364in}}%
\pgfpathlineto{\pgfqpoint{3.174906in}{3.089415in}}%
\pgfpathmoveto{\pgfqpoint{3.170365in}{3.092364in}}%
\pgfpathlineto{\pgfqpoint{3.170365in}{3.092364in}}%
\pgfpathlineto{\pgfqpoint{3.170365in}{3.095313in}}%
\pgfpathlineto{\pgfqpoint{3.174906in}{3.095313in}}%
\pgfpathlineto{\pgfqpoint{3.174906in}{3.092364in}}%
\pgfpathmoveto{\pgfqpoint{3.174906in}{3.083517in}}%
\pgfpathlineto{\pgfqpoint{3.174906in}{3.083517in}}%
\pgfpathlineto{\pgfqpoint{3.174906in}{3.086466in}}%
\pgfpathlineto{\pgfqpoint{3.179447in}{3.086466in}}%
\pgfpathlineto{\pgfqpoint{3.179447in}{3.083517in}}%
\pgfpathmoveto{\pgfqpoint{3.174906in}{3.086466in}}%
\pgfpathlineto{\pgfqpoint{3.174906in}{3.086466in}}%
\pgfpathlineto{\pgfqpoint{3.174906in}{3.089415in}}%
\pgfpathlineto{\pgfqpoint{3.179447in}{3.089415in}}%
\pgfpathlineto{\pgfqpoint{3.179447in}{3.086466in}}%
\pgfpathmoveto{\pgfqpoint{3.197611in}{3.051076in}}%
\pgfpathlineto{\pgfqpoint{3.197611in}{3.051076in}}%
\pgfpathlineto{\pgfqpoint{3.197611in}{3.054025in}}%
\pgfpathlineto{\pgfqpoint{3.202152in}{3.054025in}}%
\pgfpathlineto{\pgfqpoint{3.202152in}{3.051076in}}%
\pgfpathmoveto{\pgfqpoint{3.193070in}{3.056975in}}%
\pgfpathlineto{\pgfqpoint{3.193070in}{3.056975in}}%
\pgfpathlineto{\pgfqpoint{3.193070in}{3.059924in}}%
\pgfpathlineto{\pgfqpoint{3.197611in}{3.059924in}}%
\pgfpathlineto{\pgfqpoint{3.197611in}{3.056975in}}%
\pgfpathmoveto{\pgfqpoint{3.197611in}{3.054025in}}%
\pgfpathlineto{\pgfqpoint{3.197611in}{3.054025in}}%
\pgfpathlineto{\pgfqpoint{3.197611in}{3.056975in}}%
\pgfpathlineto{\pgfqpoint{3.202152in}{3.056975in}}%
\pgfpathlineto{\pgfqpoint{3.202152in}{3.054025in}}%
\pgfpathmoveto{\pgfqpoint{3.197611in}{3.056975in}}%
\pgfpathlineto{\pgfqpoint{3.197611in}{3.056975in}}%
\pgfpathlineto{\pgfqpoint{3.197611in}{3.059924in}}%
\pgfpathlineto{\pgfqpoint{3.202152in}{3.059924in}}%
\pgfpathlineto{\pgfqpoint{3.202152in}{3.056975in}}%
\pgfpathmoveto{\pgfqpoint{3.188529in}{3.062873in}}%
\pgfpathlineto{\pgfqpoint{3.188529in}{3.062873in}}%
\pgfpathlineto{\pgfqpoint{3.188529in}{3.065822in}}%
\pgfpathlineto{\pgfqpoint{3.193070in}{3.065822in}}%
\pgfpathlineto{\pgfqpoint{3.193070in}{3.062873in}}%
\pgfpathmoveto{\pgfqpoint{3.183988in}{3.068771in}}%
\pgfpathlineto{\pgfqpoint{3.183988in}{3.068771in}}%
\pgfpathlineto{\pgfqpoint{3.183988in}{3.071720in}}%
\pgfpathlineto{\pgfqpoint{3.188529in}{3.071720in}}%
\pgfpathlineto{\pgfqpoint{3.188529in}{3.068771in}}%
\pgfpathmoveto{\pgfqpoint{3.188529in}{3.065822in}}%
\pgfpathlineto{\pgfqpoint{3.188529in}{3.065822in}}%
\pgfpathlineto{\pgfqpoint{3.188529in}{3.068771in}}%
\pgfpathlineto{\pgfqpoint{3.193070in}{3.068771in}}%
\pgfpathlineto{\pgfqpoint{3.193070in}{3.065822in}}%
\pgfpathmoveto{\pgfqpoint{3.188529in}{3.068771in}}%
\pgfpathlineto{\pgfqpoint{3.188529in}{3.068771in}}%
\pgfpathlineto{\pgfqpoint{3.188529in}{3.071720in}}%
\pgfpathlineto{\pgfqpoint{3.193070in}{3.071720in}}%
\pgfpathlineto{\pgfqpoint{3.193070in}{3.068771in}}%
\pgfpathmoveto{\pgfqpoint{3.193070in}{3.059924in}}%
\pgfpathlineto{\pgfqpoint{3.193070in}{3.059924in}}%
\pgfpathlineto{\pgfqpoint{3.193070in}{3.062873in}}%
\pgfpathlineto{\pgfqpoint{3.197611in}{3.062873in}}%
\pgfpathlineto{\pgfqpoint{3.197611in}{3.059924in}}%
\pgfpathmoveto{\pgfqpoint{3.193070in}{3.062873in}}%
\pgfpathlineto{\pgfqpoint{3.193070in}{3.062873in}}%
\pgfpathlineto{\pgfqpoint{3.193070in}{3.065822in}}%
\pgfpathlineto{\pgfqpoint{3.197611in}{3.065822in}}%
\pgfpathlineto{\pgfqpoint{3.197611in}{3.062873in}}%
\pgfpathmoveto{\pgfqpoint{3.202152in}{3.048127in}}%
\pgfpathlineto{\pgfqpoint{3.202152in}{3.048127in}}%
\pgfpathlineto{\pgfqpoint{3.202152in}{3.051076in}}%
\pgfpathlineto{\pgfqpoint{3.206694in}{3.051076in}}%
\pgfpathlineto{\pgfqpoint{3.206694in}{3.048127in}}%
\pgfpathmoveto{\pgfqpoint{3.202152in}{3.051076in}}%
\pgfpathlineto{\pgfqpoint{3.202152in}{3.051076in}}%
\pgfpathlineto{\pgfqpoint{3.202152in}{3.054025in}}%
\pgfpathlineto{\pgfqpoint{3.206694in}{3.054025in}}%
\pgfpathlineto{\pgfqpoint{3.206694in}{3.051076in}}%
\pgfpathmoveto{\pgfqpoint{3.183988in}{3.071720in}}%
\pgfpathlineto{\pgfqpoint{3.183988in}{3.071720in}}%
\pgfpathlineto{\pgfqpoint{3.183988in}{3.074669in}}%
\pgfpathlineto{\pgfqpoint{3.188529in}{3.074669in}}%
\pgfpathlineto{\pgfqpoint{3.188529in}{3.071720in}}%
\pgfpathmoveto{\pgfqpoint{3.183988in}{3.074669in}}%
\pgfpathlineto{\pgfqpoint{3.183988in}{3.074669in}}%
\pgfpathlineto{\pgfqpoint{3.183988in}{3.077618in}}%
\pgfpathlineto{\pgfqpoint{3.188529in}{3.077618in}}%
\pgfpathlineto{\pgfqpoint{3.188529in}{3.074669in}}%
\pgfpathmoveto{\pgfqpoint{3.161283in}{3.098262in}}%
\pgfpathlineto{\pgfqpoint{3.161283in}{3.098262in}}%
\pgfpathlineto{\pgfqpoint{3.161283in}{3.101211in}}%
\pgfpathlineto{\pgfqpoint{3.165824in}{3.101211in}}%
\pgfpathlineto{\pgfqpoint{3.165824in}{3.098262in}}%
\pgfpathmoveto{\pgfqpoint{3.156742in}{3.104160in}}%
\pgfpathlineto{\pgfqpoint{3.156742in}{3.104160in}}%
\pgfpathlineto{\pgfqpoint{3.156742in}{3.107109in}}%
\pgfpathlineto{\pgfqpoint{3.161283in}{3.107109in}}%
\pgfpathlineto{\pgfqpoint{3.161283in}{3.104160in}}%
\pgfpathmoveto{\pgfqpoint{3.161283in}{3.101211in}}%
\pgfpathlineto{\pgfqpoint{3.161283in}{3.101211in}}%
\pgfpathlineto{\pgfqpoint{3.161283in}{3.104160in}}%
\pgfpathlineto{\pgfqpoint{3.165824in}{3.104160in}}%
\pgfpathlineto{\pgfqpoint{3.165824in}{3.101211in}}%
\pgfpathmoveto{\pgfqpoint{3.161283in}{3.104160in}}%
\pgfpathlineto{\pgfqpoint{3.161283in}{3.104160in}}%
\pgfpathlineto{\pgfqpoint{3.161283in}{3.107109in}}%
\pgfpathlineto{\pgfqpoint{3.165824in}{3.107109in}}%
\pgfpathlineto{\pgfqpoint{3.165824in}{3.104160in}}%
\pgfpathmoveto{\pgfqpoint{3.152201in}{3.110058in}}%
\pgfpathlineto{\pgfqpoint{3.152201in}{3.110058in}}%
\pgfpathlineto{\pgfqpoint{3.152201in}{3.113008in}}%
\pgfpathlineto{\pgfqpoint{3.156742in}{3.113008in}}%
\pgfpathlineto{\pgfqpoint{3.156742in}{3.110058in}}%
\pgfpathmoveto{\pgfqpoint{3.147660in}{3.115957in}}%
\pgfpathlineto{\pgfqpoint{3.147660in}{3.115957in}}%
\pgfpathlineto{\pgfqpoint{3.147660in}{3.118906in}}%
\pgfpathlineto{\pgfqpoint{3.152201in}{3.118906in}}%
\pgfpathlineto{\pgfqpoint{3.152201in}{3.115957in}}%
\pgfpathmoveto{\pgfqpoint{3.152201in}{3.113008in}}%
\pgfpathlineto{\pgfqpoint{3.152201in}{3.113008in}}%
\pgfpathlineto{\pgfqpoint{3.152201in}{3.115957in}}%
\pgfpathlineto{\pgfqpoint{3.156742in}{3.115957in}}%
\pgfpathlineto{\pgfqpoint{3.156742in}{3.113008in}}%
\pgfpathmoveto{\pgfqpoint{3.152201in}{3.115957in}}%
\pgfpathlineto{\pgfqpoint{3.152201in}{3.115957in}}%
\pgfpathlineto{\pgfqpoint{3.152201in}{3.118906in}}%
\pgfpathlineto{\pgfqpoint{3.156742in}{3.118906in}}%
\pgfpathlineto{\pgfqpoint{3.156742in}{3.115957in}}%
\pgfpathmoveto{\pgfqpoint{3.156742in}{3.107109in}}%
\pgfpathlineto{\pgfqpoint{3.156742in}{3.107109in}}%
\pgfpathlineto{\pgfqpoint{3.156742in}{3.110058in}}%
\pgfpathlineto{\pgfqpoint{3.161283in}{3.110058in}}%
\pgfpathlineto{\pgfqpoint{3.161283in}{3.107109in}}%
\pgfpathmoveto{\pgfqpoint{3.156742in}{3.110058in}}%
\pgfpathlineto{\pgfqpoint{3.156742in}{3.110058in}}%
\pgfpathlineto{\pgfqpoint{3.156742in}{3.113008in}}%
\pgfpathlineto{\pgfqpoint{3.161283in}{3.113008in}}%
\pgfpathlineto{\pgfqpoint{3.161283in}{3.110058in}}%
\pgfpathmoveto{\pgfqpoint{3.165824in}{3.095313in}}%
\pgfpathlineto{\pgfqpoint{3.165824in}{3.095313in}}%
\pgfpathlineto{\pgfqpoint{3.165824in}{3.098262in}}%
\pgfpathlineto{\pgfqpoint{3.170365in}{3.098262in}}%
\pgfpathlineto{\pgfqpoint{3.170365in}{3.095313in}}%
\pgfpathmoveto{\pgfqpoint{3.165824in}{3.098262in}}%
\pgfpathlineto{\pgfqpoint{3.165824in}{3.098262in}}%
\pgfpathlineto{\pgfqpoint{3.165824in}{3.101211in}}%
\pgfpathlineto{\pgfqpoint{3.170365in}{3.101211in}}%
\pgfpathlineto{\pgfqpoint{3.170365in}{3.098262in}}%
\pgfpathmoveto{\pgfqpoint{3.147660in}{3.118906in}}%
\pgfpathlineto{\pgfqpoint{3.147660in}{3.118906in}}%
\pgfpathlineto{\pgfqpoint{3.147660in}{3.121855in}}%
\pgfpathlineto{\pgfqpoint{3.152201in}{3.121855in}}%
\pgfpathlineto{\pgfqpoint{3.152201in}{3.118906in}}%
\pgfpathmoveto{\pgfqpoint{3.147660in}{3.121855in}}%
\pgfpathlineto{\pgfqpoint{3.147660in}{3.121855in}}%
\pgfpathlineto{\pgfqpoint{3.147660in}{3.124804in}}%
\pgfpathlineto{\pgfqpoint{3.152201in}{3.124804in}}%
\pgfpathlineto{\pgfqpoint{3.152201in}{3.121855in}}%
\pgfpathmoveto{\pgfqpoint{3.106790in}{3.169042in}}%
\pgfpathlineto{\pgfqpoint{3.106790in}{3.169042in}}%
\pgfpathlineto{\pgfqpoint{3.106790in}{3.171991in}}%
\pgfpathlineto{\pgfqpoint{3.111331in}{3.171991in}}%
\pgfpathlineto{\pgfqpoint{3.111331in}{3.169042in}}%
\pgfpathmoveto{\pgfqpoint{3.102249in}{3.174940in}}%
\pgfpathlineto{\pgfqpoint{3.102249in}{3.174940in}}%
\pgfpathlineto{\pgfqpoint{3.102249in}{3.177890in}}%
\pgfpathlineto{\pgfqpoint{3.106790in}{3.177890in}}%
\pgfpathlineto{\pgfqpoint{3.106790in}{3.174940in}}%
\pgfpathmoveto{\pgfqpoint{3.106790in}{3.171991in}}%
\pgfpathlineto{\pgfqpoint{3.106790in}{3.171991in}}%
\pgfpathlineto{\pgfqpoint{3.106790in}{3.174940in}}%
\pgfpathlineto{\pgfqpoint{3.111331in}{3.174940in}}%
\pgfpathlineto{\pgfqpoint{3.111331in}{3.171991in}}%
\pgfpathmoveto{\pgfqpoint{3.106790in}{3.174940in}}%
\pgfpathlineto{\pgfqpoint{3.106790in}{3.174940in}}%
\pgfpathlineto{\pgfqpoint{3.106790in}{3.177890in}}%
\pgfpathlineto{\pgfqpoint{3.111331in}{3.177890in}}%
\pgfpathlineto{\pgfqpoint{3.111331in}{3.174940in}}%
\pgfpathmoveto{\pgfqpoint{3.097708in}{3.180839in}}%
\pgfpathlineto{\pgfqpoint{3.097708in}{3.180839in}}%
\pgfpathlineto{\pgfqpoint{3.097708in}{3.183788in}}%
\pgfpathlineto{\pgfqpoint{3.102249in}{3.183788in}}%
\pgfpathlineto{\pgfqpoint{3.102249in}{3.180839in}}%
\pgfpathmoveto{\pgfqpoint{3.093167in}{3.186737in}}%
\pgfpathlineto{\pgfqpoint{3.093167in}{3.186737in}}%
\pgfpathlineto{\pgfqpoint{3.093167in}{3.189686in}}%
\pgfpathlineto{\pgfqpoint{3.097708in}{3.189686in}}%
\pgfpathlineto{\pgfqpoint{3.097708in}{3.186737in}}%
\pgfpathmoveto{\pgfqpoint{3.097708in}{3.183788in}}%
\pgfpathlineto{\pgfqpoint{3.097708in}{3.183788in}}%
\pgfpathlineto{\pgfqpoint{3.097708in}{3.186737in}}%
\pgfpathlineto{\pgfqpoint{3.102249in}{3.186737in}}%
\pgfpathlineto{\pgfqpoint{3.102249in}{3.183788in}}%
\pgfpathmoveto{\pgfqpoint{3.097708in}{3.186737in}}%
\pgfpathlineto{\pgfqpoint{3.097708in}{3.186737in}}%
\pgfpathlineto{\pgfqpoint{3.097708in}{3.189686in}}%
\pgfpathlineto{\pgfqpoint{3.102249in}{3.189686in}}%
\pgfpathlineto{\pgfqpoint{3.102249in}{3.186737in}}%
\pgfpathmoveto{\pgfqpoint{3.102249in}{3.177890in}}%
\pgfpathlineto{\pgfqpoint{3.102249in}{3.177890in}}%
\pgfpathlineto{\pgfqpoint{3.102249in}{3.180839in}}%
\pgfpathlineto{\pgfqpoint{3.106790in}{3.180839in}}%
\pgfpathlineto{\pgfqpoint{3.106790in}{3.177890in}}%
\pgfpathmoveto{\pgfqpoint{3.102249in}{3.180839in}}%
\pgfpathlineto{\pgfqpoint{3.102249in}{3.180839in}}%
\pgfpathlineto{\pgfqpoint{3.102249in}{3.183788in}}%
\pgfpathlineto{\pgfqpoint{3.106790in}{3.183788in}}%
\pgfpathlineto{\pgfqpoint{3.106790in}{3.180839in}}%
\pgfpathmoveto{\pgfqpoint{3.124954in}{3.145448in}}%
\pgfpathlineto{\pgfqpoint{3.124954in}{3.145448in}}%
\pgfpathlineto{\pgfqpoint{3.124954in}{3.148397in}}%
\pgfpathlineto{\pgfqpoint{3.129495in}{3.148397in}}%
\pgfpathlineto{\pgfqpoint{3.129495in}{3.145448in}}%
\pgfpathmoveto{\pgfqpoint{3.120413in}{3.151346in}}%
\pgfpathlineto{\pgfqpoint{3.120413in}{3.151346in}}%
\pgfpathlineto{\pgfqpoint{3.120413in}{3.154296in}}%
\pgfpathlineto{\pgfqpoint{3.124954in}{3.154296in}}%
\pgfpathlineto{\pgfqpoint{3.124954in}{3.151346in}}%
\pgfpathmoveto{\pgfqpoint{3.124954in}{3.148397in}}%
\pgfpathlineto{\pgfqpoint{3.124954in}{3.148397in}}%
\pgfpathlineto{\pgfqpoint{3.124954in}{3.151346in}}%
\pgfpathlineto{\pgfqpoint{3.129495in}{3.151346in}}%
\pgfpathlineto{\pgfqpoint{3.129495in}{3.148397in}}%
\pgfpathmoveto{\pgfqpoint{3.124954in}{3.151346in}}%
\pgfpathlineto{\pgfqpoint{3.124954in}{3.151346in}}%
\pgfpathlineto{\pgfqpoint{3.124954in}{3.154296in}}%
\pgfpathlineto{\pgfqpoint{3.129495in}{3.154296in}}%
\pgfpathlineto{\pgfqpoint{3.129495in}{3.151346in}}%
\pgfpathmoveto{\pgfqpoint{3.115872in}{3.157245in}}%
\pgfpathlineto{\pgfqpoint{3.115872in}{3.157245in}}%
\pgfpathlineto{\pgfqpoint{3.115872in}{3.160194in}}%
\pgfpathlineto{\pgfqpoint{3.120413in}{3.160194in}}%
\pgfpathlineto{\pgfqpoint{3.120413in}{3.157245in}}%
\pgfpathmoveto{\pgfqpoint{3.111331in}{3.163143in}}%
\pgfpathlineto{\pgfqpoint{3.111331in}{3.163143in}}%
\pgfpathlineto{\pgfqpoint{3.111331in}{3.166093in}}%
\pgfpathlineto{\pgfqpoint{3.115872in}{3.166093in}}%
\pgfpathlineto{\pgfqpoint{3.115872in}{3.163143in}}%
\pgfpathmoveto{\pgfqpoint{3.115872in}{3.160194in}}%
\pgfpathlineto{\pgfqpoint{3.115872in}{3.160194in}}%
\pgfpathlineto{\pgfqpoint{3.115872in}{3.163143in}}%
\pgfpathlineto{\pgfqpoint{3.120413in}{3.163143in}}%
\pgfpathlineto{\pgfqpoint{3.120413in}{3.160194in}}%
\pgfpathmoveto{\pgfqpoint{3.115872in}{3.163143in}}%
\pgfpathlineto{\pgfqpoint{3.115872in}{3.163143in}}%
\pgfpathlineto{\pgfqpoint{3.115872in}{3.166093in}}%
\pgfpathlineto{\pgfqpoint{3.120413in}{3.166093in}}%
\pgfpathlineto{\pgfqpoint{3.120413in}{3.163143in}}%
\pgfpathmoveto{\pgfqpoint{3.120413in}{3.154296in}}%
\pgfpathlineto{\pgfqpoint{3.120413in}{3.154296in}}%
\pgfpathlineto{\pgfqpoint{3.120413in}{3.157245in}}%
\pgfpathlineto{\pgfqpoint{3.124954in}{3.157245in}}%
\pgfpathlineto{\pgfqpoint{3.124954in}{3.154296in}}%
\pgfpathmoveto{\pgfqpoint{3.120413in}{3.157245in}}%
\pgfpathlineto{\pgfqpoint{3.120413in}{3.157245in}}%
\pgfpathlineto{\pgfqpoint{3.120413in}{3.160194in}}%
\pgfpathlineto{\pgfqpoint{3.124954in}{3.160194in}}%
\pgfpathlineto{\pgfqpoint{3.124954in}{3.157245in}}%
\pgfpathmoveto{\pgfqpoint{3.129495in}{3.142499in}}%
\pgfpathlineto{\pgfqpoint{3.129495in}{3.142499in}}%
\pgfpathlineto{\pgfqpoint{3.129495in}{3.145448in}}%
\pgfpathlineto{\pgfqpoint{3.134036in}{3.145448in}}%
\pgfpathlineto{\pgfqpoint{3.134036in}{3.142499in}}%
\pgfpathmoveto{\pgfqpoint{3.129495in}{3.145448in}}%
\pgfpathlineto{\pgfqpoint{3.129495in}{3.145448in}}%
\pgfpathlineto{\pgfqpoint{3.129495in}{3.148397in}}%
\pgfpathlineto{\pgfqpoint{3.134036in}{3.148397in}}%
\pgfpathlineto{\pgfqpoint{3.134036in}{3.145448in}}%
\pgfpathmoveto{\pgfqpoint{3.111331in}{3.166093in}}%
\pgfpathlineto{\pgfqpoint{3.111331in}{3.166093in}}%
\pgfpathlineto{\pgfqpoint{3.111331in}{3.169042in}}%
\pgfpathlineto{\pgfqpoint{3.115872in}{3.169042in}}%
\pgfpathlineto{\pgfqpoint{3.115872in}{3.166093in}}%
\pgfpathmoveto{\pgfqpoint{3.111331in}{3.169042in}}%
\pgfpathlineto{\pgfqpoint{3.111331in}{3.169042in}}%
\pgfpathlineto{\pgfqpoint{3.111331in}{3.171991in}}%
\pgfpathlineto{\pgfqpoint{3.115872in}{3.171991in}}%
\pgfpathlineto{\pgfqpoint{3.115872in}{3.169042in}}%
\pgfpathmoveto{\pgfqpoint{3.088626in}{3.192636in}}%
\pgfpathlineto{\pgfqpoint{3.088626in}{3.192636in}}%
\pgfpathlineto{\pgfqpoint{3.088626in}{3.195585in}}%
\pgfpathlineto{\pgfqpoint{3.093167in}{3.195585in}}%
\pgfpathlineto{\pgfqpoint{3.093167in}{3.192636in}}%
\pgfpathmoveto{\pgfqpoint{3.084085in}{3.198534in}}%
\pgfpathlineto{\pgfqpoint{3.084085in}{3.198534in}}%
\pgfpathlineto{\pgfqpoint{3.084085in}{3.201483in}}%
\pgfpathlineto{\pgfqpoint{3.088626in}{3.201483in}}%
\pgfpathlineto{\pgfqpoint{3.088626in}{3.198534in}}%
\pgfpathmoveto{\pgfqpoint{3.088626in}{3.195585in}}%
\pgfpathlineto{\pgfqpoint{3.088626in}{3.195585in}}%
\pgfpathlineto{\pgfqpoint{3.088626in}{3.198534in}}%
\pgfpathlineto{\pgfqpoint{3.093167in}{3.198534in}}%
\pgfpathlineto{\pgfqpoint{3.093167in}{3.195585in}}%
\pgfpathmoveto{\pgfqpoint{3.088626in}{3.198534in}}%
\pgfpathlineto{\pgfqpoint{3.088626in}{3.198534in}}%
\pgfpathlineto{\pgfqpoint{3.088626in}{3.201483in}}%
\pgfpathlineto{\pgfqpoint{3.093167in}{3.201483in}}%
\pgfpathlineto{\pgfqpoint{3.093167in}{3.198534in}}%
\pgfpathmoveto{\pgfqpoint{3.079543in}{3.204433in}}%
\pgfpathlineto{\pgfqpoint{3.079543in}{3.204433in}}%
\pgfpathlineto{\pgfqpoint{3.079543in}{3.207382in}}%
\pgfpathlineto{\pgfqpoint{3.084085in}{3.207382in}}%
\pgfpathlineto{\pgfqpoint{3.084085in}{3.204433in}}%
\pgfpathmoveto{\pgfqpoint{3.075002in}{3.210331in}}%
\pgfpathlineto{\pgfqpoint{3.075002in}{3.210331in}}%
\pgfpathlineto{\pgfqpoint{3.075002in}{3.213280in}}%
\pgfpathlineto{\pgfqpoint{3.079543in}{3.213280in}}%
\pgfpathlineto{\pgfqpoint{3.079543in}{3.210331in}}%
\pgfpathmoveto{\pgfqpoint{3.079543in}{3.207382in}}%
\pgfpathlineto{\pgfqpoint{3.079543in}{3.207382in}}%
\pgfpathlineto{\pgfqpoint{3.079543in}{3.210331in}}%
\pgfpathlineto{\pgfqpoint{3.084085in}{3.210331in}}%
\pgfpathlineto{\pgfqpoint{3.084085in}{3.207382in}}%
\pgfpathmoveto{\pgfqpoint{3.079543in}{3.210331in}}%
\pgfpathlineto{\pgfqpoint{3.079543in}{3.210331in}}%
\pgfpathlineto{\pgfqpoint{3.079543in}{3.213280in}}%
\pgfpathlineto{\pgfqpoint{3.084085in}{3.213280in}}%
\pgfpathlineto{\pgfqpoint{3.084085in}{3.210331in}}%
\pgfpathmoveto{\pgfqpoint{3.084085in}{3.201483in}}%
\pgfpathlineto{\pgfqpoint{3.084085in}{3.201483in}}%
\pgfpathlineto{\pgfqpoint{3.084085in}{3.204433in}}%
\pgfpathlineto{\pgfqpoint{3.088626in}{3.204433in}}%
\pgfpathlineto{\pgfqpoint{3.088626in}{3.201483in}}%
\pgfpathmoveto{\pgfqpoint{3.084085in}{3.204433in}}%
\pgfpathlineto{\pgfqpoint{3.084085in}{3.204433in}}%
\pgfpathlineto{\pgfqpoint{3.084085in}{3.207382in}}%
\pgfpathlineto{\pgfqpoint{3.088626in}{3.207382in}}%
\pgfpathlineto{\pgfqpoint{3.088626in}{3.204433in}}%
\pgfpathmoveto{\pgfqpoint{3.093167in}{3.189686in}}%
\pgfpathlineto{\pgfqpoint{3.093167in}{3.189686in}}%
\pgfpathlineto{\pgfqpoint{3.093167in}{3.192636in}}%
\pgfpathlineto{\pgfqpoint{3.097708in}{3.192636in}}%
\pgfpathlineto{\pgfqpoint{3.097708in}{3.189686in}}%
\pgfpathmoveto{\pgfqpoint{3.093167in}{3.192636in}}%
\pgfpathlineto{\pgfqpoint{3.093167in}{3.192636in}}%
\pgfpathlineto{\pgfqpoint{3.093167in}{3.195585in}}%
\pgfpathlineto{\pgfqpoint{3.097708in}{3.195585in}}%
\pgfpathlineto{\pgfqpoint{3.097708in}{3.192636in}}%
\pgfpathmoveto{\pgfqpoint{3.075002in}{3.213280in}}%
\pgfpathlineto{\pgfqpoint{3.075002in}{3.213280in}}%
\pgfpathlineto{\pgfqpoint{3.075002in}{3.216230in}}%
\pgfpathlineto{\pgfqpoint{3.079543in}{3.216230in}}%
\pgfpathlineto{\pgfqpoint{3.079543in}{3.213280in}}%
\pgfpathmoveto{\pgfqpoint{3.075002in}{3.216230in}}%
\pgfpathlineto{\pgfqpoint{3.075002in}{3.216230in}}%
\pgfpathlineto{\pgfqpoint{3.075002in}{3.219179in}}%
\pgfpathlineto{\pgfqpoint{3.079543in}{3.219179in}}%
\pgfpathlineto{\pgfqpoint{3.079543in}{3.216230in}}%
\pgfpathmoveto{\pgfqpoint{3.361088in}{2.838732in}}%
\pgfpathlineto{\pgfqpoint{3.361088in}{2.838732in}}%
\pgfpathlineto{\pgfqpoint{3.361088in}{2.841681in}}%
\pgfpathlineto{\pgfqpoint{3.365629in}{2.841681in}}%
\pgfpathlineto{\pgfqpoint{3.365629in}{2.838732in}}%
\pgfpathmoveto{\pgfqpoint{3.356547in}{2.844630in}}%
\pgfpathlineto{\pgfqpoint{3.356547in}{2.844630in}}%
\pgfpathlineto{\pgfqpoint{3.356547in}{2.847580in}}%
\pgfpathlineto{\pgfqpoint{3.361088in}{2.847580in}}%
\pgfpathlineto{\pgfqpoint{3.361088in}{2.844630in}}%
\pgfpathmoveto{\pgfqpoint{3.361088in}{2.841681in}}%
\pgfpathlineto{\pgfqpoint{3.361088in}{2.841681in}}%
\pgfpathlineto{\pgfqpoint{3.361088in}{2.844630in}}%
\pgfpathlineto{\pgfqpoint{3.365629in}{2.844630in}}%
\pgfpathlineto{\pgfqpoint{3.365629in}{2.841681in}}%
\pgfpathmoveto{\pgfqpoint{3.361088in}{2.844630in}}%
\pgfpathlineto{\pgfqpoint{3.361088in}{2.844630in}}%
\pgfpathlineto{\pgfqpoint{3.361088in}{2.847580in}}%
\pgfpathlineto{\pgfqpoint{3.365629in}{2.847580in}}%
\pgfpathlineto{\pgfqpoint{3.365629in}{2.844630in}}%
\pgfpathmoveto{\pgfqpoint{3.352006in}{2.850529in}}%
\pgfpathlineto{\pgfqpoint{3.352006in}{2.850529in}}%
\pgfpathlineto{\pgfqpoint{3.352006in}{2.853479in}}%
\pgfpathlineto{\pgfqpoint{3.356547in}{2.853479in}}%
\pgfpathlineto{\pgfqpoint{3.356547in}{2.850529in}}%
\pgfpathmoveto{\pgfqpoint{3.347465in}{2.856428in}}%
\pgfpathlineto{\pgfqpoint{3.347465in}{2.856428in}}%
\pgfpathlineto{\pgfqpoint{3.347465in}{2.859377in}}%
\pgfpathlineto{\pgfqpoint{3.352006in}{2.859377in}}%
\pgfpathlineto{\pgfqpoint{3.352006in}{2.856428in}}%
\pgfpathmoveto{\pgfqpoint{3.352006in}{2.853479in}}%
\pgfpathlineto{\pgfqpoint{3.352006in}{2.853479in}}%
\pgfpathlineto{\pgfqpoint{3.352006in}{2.856428in}}%
\pgfpathlineto{\pgfqpoint{3.356547in}{2.856428in}}%
\pgfpathlineto{\pgfqpoint{3.356547in}{2.853479in}}%
\pgfpathmoveto{\pgfqpoint{3.352006in}{2.856428in}}%
\pgfpathlineto{\pgfqpoint{3.352006in}{2.856428in}}%
\pgfpathlineto{\pgfqpoint{3.352006in}{2.859377in}}%
\pgfpathlineto{\pgfqpoint{3.356547in}{2.859377in}}%
\pgfpathlineto{\pgfqpoint{3.356547in}{2.856428in}}%
\pgfpathmoveto{\pgfqpoint{3.356547in}{2.847580in}}%
\pgfpathlineto{\pgfqpoint{3.356547in}{2.847580in}}%
\pgfpathlineto{\pgfqpoint{3.356547in}{2.850529in}}%
\pgfpathlineto{\pgfqpoint{3.361088in}{2.850529in}}%
\pgfpathlineto{\pgfqpoint{3.361088in}{2.847580in}}%
\pgfpathmoveto{\pgfqpoint{3.356547in}{2.850529in}}%
\pgfpathlineto{\pgfqpoint{3.356547in}{2.850529in}}%
\pgfpathlineto{\pgfqpoint{3.356547in}{2.853479in}}%
\pgfpathlineto{\pgfqpoint{3.361088in}{2.853479in}}%
\pgfpathlineto{\pgfqpoint{3.361088in}{2.850529in}}%
\pgfpathmoveto{\pgfqpoint{3.288432in}{2.933105in}}%
\pgfpathlineto{\pgfqpoint{3.288432in}{2.933105in}}%
\pgfpathlineto{\pgfqpoint{3.288432in}{2.936055in}}%
\pgfpathlineto{\pgfqpoint{3.292973in}{2.936055in}}%
\pgfpathlineto{\pgfqpoint{3.292973in}{2.933105in}}%
\pgfpathmoveto{\pgfqpoint{3.283891in}{2.939004in}}%
\pgfpathlineto{\pgfqpoint{3.283891in}{2.939004in}}%
\pgfpathlineto{\pgfqpoint{3.283891in}{2.941953in}}%
\pgfpathlineto{\pgfqpoint{3.288432in}{2.941953in}}%
\pgfpathlineto{\pgfqpoint{3.288432in}{2.939004in}}%
\pgfpathmoveto{\pgfqpoint{3.288432in}{2.936055in}}%
\pgfpathlineto{\pgfqpoint{3.288432in}{2.936055in}}%
\pgfpathlineto{\pgfqpoint{3.288432in}{2.939004in}}%
\pgfpathlineto{\pgfqpoint{3.292973in}{2.939004in}}%
\pgfpathlineto{\pgfqpoint{3.292973in}{2.936055in}}%
\pgfpathmoveto{\pgfqpoint{3.288432in}{2.939004in}}%
\pgfpathlineto{\pgfqpoint{3.288432in}{2.939004in}}%
\pgfpathlineto{\pgfqpoint{3.288432in}{2.941953in}}%
\pgfpathlineto{\pgfqpoint{3.292973in}{2.941953in}}%
\pgfpathlineto{\pgfqpoint{3.292973in}{2.939004in}}%
\pgfpathmoveto{\pgfqpoint{3.279350in}{2.944902in}}%
\pgfpathlineto{\pgfqpoint{3.279350in}{2.944902in}}%
\pgfpathlineto{\pgfqpoint{3.279350in}{2.947851in}}%
\pgfpathlineto{\pgfqpoint{3.283891in}{2.947851in}}%
\pgfpathlineto{\pgfqpoint{3.283891in}{2.944902in}}%
\pgfpathmoveto{\pgfqpoint{3.274809in}{2.950800in}}%
\pgfpathlineto{\pgfqpoint{3.274809in}{2.950800in}}%
\pgfpathlineto{\pgfqpoint{3.274809in}{2.953749in}}%
\pgfpathlineto{\pgfqpoint{3.279350in}{2.953749in}}%
\pgfpathlineto{\pgfqpoint{3.279350in}{2.950800in}}%
\pgfpathmoveto{\pgfqpoint{3.279350in}{2.947851in}}%
\pgfpathlineto{\pgfqpoint{3.279350in}{2.947851in}}%
\pgfpathlineto{\pgfqpoint{3.279350in}{2.950800in}}%
\pgfpathlineto{\pgfqpoint{3.283891in}{2.950800in}}%
\pgfpathlineto{\pgfqpoint{3.283891in}{2.947851in}}%
\pgfpathmoveto{\pgfqpoint{3.279350in}{2.950800in}}%
\pgfpathlineto{\pgfqpoint{3.279350in}{2.950800in}}%
\pgfpathlineto{\pgfqpoint{3.279350in}{2.953749in}}%
\pgfpathlineto{\pgfqpoint{3.283891in}{2.953749in}}%
\pgfpathlineto{\pgfqpoint{3.283891in}{2.950800in}}%
\pgfpathmoveto{\pgfqpoint{3.283891in}{2.941953in}}%
\pgfpathlineto{\pgfqpoint{3.283891in}{2.941953in}}%
\pgfpathlineto{\pgfqpoint{3.283891in}{2.944902in}}%
\pgfpathlineto{\pgfqpoint{3.288432in}{2.944902in}}%
\pgfpathlineto{\pgfqpoint{3.288432in}{2.941953in}}%
\pgfpathmoveto{\pgfqpoint{3.283891in}{2.944902in}}%
\pgfpathlineto{\pgfqpoint{3.283891in}{2.944902in}}%
\pgfpathlineto{\pgfqpoint{3.283891in}{2.947851in}}%
\pgfpathlineto{\pgfqpoint{3.288432in}{2.947851in}}%
\pgfpathlineto{\pgfqpoint{3.288432in}{2.944902in}}%
\pgfpathmoveto{\pgfqpoint{3.324760in}{2.885919in}}%
\pgfpathlineto{\pgfqpoint{3.324760in}{2.885919in}}%
\pgfpathlineto{\pgfqpoint{3.324760in}{2.888869in}}%
\pgfpathlineto{\pgfqpoint{3.329301in}{2.888869in}}%
\pgfpathlineto{\pgfqpoint{3.329301in}{2.885919in}}%
\pgfpathmoveto{\pgfqpoint{3.320219in}{2.891818in}}%
\pgfpathlineto{\pgfqpoint{3.320219in}{2.891818in}}%
\pgfpathlineto{\pgfqpoint{3.320219in}{2.894767in}}%
\pgfpathlineto{\pgfqpoint{3.324760in}{2.894767in}}%
\pgfpathlineto{\pgfqpoint{3.324760in}{2.891818in}}%
\pgfpathmoveto{\pgfqpoint{3.324760in}{2.888869in}}%
\pgfpathlineto{\pgfqpoint{3.324760in}{2.888869in}}%
\pgfpathlineto{\pgfqpoint{3.324760in}{2.891818in}}%
\pgfpathlineto{\pgfqpoint{3.329301in}{2.891818in}}%
\pgfpathlineto{\pgfqpoint{3.329301in}{2.888869in}}%
\pgfpathmoveto{\pgfqpoint{3.324760in}{2.891818in}}%
\pgfpathlineto{\pgfqpoint{3.324760in}{2.891818in}}%
\pgfpathlineto{\pgfqpoint{3.324760in}{2.894767in}}%
\pgfpathlineto{\pgfqpoint{3.329301in}{2.894767in}}%
\pgfpathlineto{\pgfqpoint{3.329301in}{2.891818in}}%
\pgfpathmoveto{\pgfqpoint{3.315678in}{2.897716in}}%
\pgfpathlineto{\pgfqpoint{3.315678in}{2.897716in}}%
\pgfpathlineto{\pgfqpoint{3.315678in}{2.900665in}}%
\pgfpathlineto{\pgfqpoint{3.320219in}{2.900665in}}%
\pgfpathlineto{\pgfqpoint{3.320219in}{2.897716in}}%
\pgfpathmoveto{\pgfqpoint{3.311137in}{2.903614in}}%
\pgfpathlineto{\pgfqpoint{3.311137in}{2.903614in}}%
\pgfpathlineto{\pgfqpoint{3.311137in}{2.906563in}}%
\pgfpathlineto{\pgfqpoint{3.315678in}{2.906563in}}%
\pgfpathlineto{\pgfqpoint{3.315678in}{2.903614in}}%
\pgfpathmoveto{\pgfqpoint{3.315678in}{2.900665in}}%
\pgfpathlineto{\pgfqpoint{3.315678in}{2.900665in}}%
\pgfpathlineto{\pgfqpoint{3.315678in}{2.903614in}}%
\pgfpathlineto{\pgfqpoint{3.320219in}{2.903614in}}%
\pgfpathlineto{\pgfqpoint{3.320219in}{2.900665in}}%
\pgfpathmoveto{\pgfqpoint{3.315678in}{2.903614in}}%
\pgfpathlineto{\pgfqpoint{3.315678in}{2.903614in}}%
\pgfpathlineto{\pgfqpoint{3.315678in}{2.906563in}}%
\pgfpathlineto{\pgfqpoint{3.320219in}{2.906563in}}%
\pgfpathlineto{\pgfqpoint{3.320219in}{2.903614in}}%
\pgfpathmoveto{\pgfqpoint{3.320219in}{2.894767in}}%
\pgfpathlineto{\pgfqpoint{3.320219in}{2.894767in}}%
\pgfpathlineto{\pgfqpoint{3.320219in}{2.897716in}}%
\pgfpathlineto{\pgfqpoint{3.324760in}{2.897716in}}%
\pgfpathlineto{\pgfqpoint{3.324760in}{2.894767in}}%
\pgfpathmoveto{\pgfqpoint{3.320219in}{2.897716in}}%
\pgfpathlineto{\pgfqpoint{3.320219in}{2.897716in}}%
\pgfpathlineto{\pgfqpoint{3.320219in}{2.900665in}}%
\pgfpathlineto{\pgfqpoint{3.324760in}{2.900665in}}%
\pgfpathlineto{\pgfqpoint{3.324760in}{2.897716in}}%
\pgfpathmoveto{\pgfqpoint{3.342924in}{2.862326in}}%
\pgfpathlineto{\pgfqpoint{3.342924in}{2.862326in}}%
\pgfpathlineto{\pgfqpoint{3.342924in}{2.865276in}}%
\pgfpathlineto{\pgfqpoint{3.347465in}{2.865276in}}%
\pgfpathlineto{\pgfqpoint{3.347465in}{2.862326in}}%
\pgfpathmoveto{\pgfqpoint{3.338383in}{2.868225in}}%
\pgfpathlineto{\pgfqpoint{3.338383in}{2.868225in}}%
\pgfpathlineto{\pgfqpoint{3.338383in}{2.871174in}}%
\pgfpathlineto{\pgfqpoint{3.342924in}{2.871174in}}%
\pgfpathlineto{\pgfqpoint{3.342924in}{2.868225in}}%
\pgfpathmoveto{\pgfqpoint{3.342924in}{2.865276in}}%
\pgfpathlineto{\pgfqpoint{3.342924in}{2.865276in}}%
\pgfpathlineto{\pgfqpoint{3.342924in}{2.868225in}}%
\pgfpathlineto{\pgfqpoint{3.347465in}{2.868225in}}%
\pgfpathlineto{\pgfqpoint{3.347465in}{2.865276in}}%
\pgfpathmoveto{\pgfqpoint{3.342924in}{2.868225in}}%
\pgfpathlineto{\pgfqpoint{3.342924in}{2.868225in}}%
\pgfpathlineto{\pgfqpoint{3.342924in}{2.871174in}}%
\pgfpathlineto{\pgfqpoint{3.347465in}{2.871174in}}%
\pgfpathlineto{\pgfqpoint{3.347465in}{2.868225in}}%
\pgfpathmoveto{\pgfqpoint{3.333842in}{2.874123in}}%
\pgfpathlineto{\pgfqpoint{3.333842in}{2.874123in}}%
\pgfpathlineto{\pgfqpoint{3.333842in}{2.877072in}}%
\pgfpathlineto{\pgfqpoint{3.338383in}{2.877072in}}%
\pgfpathlineto{\pgfqpoint{3.338383in}{2.874123in}}%
\pgfpathmoveto{\pgfqpoint{3.329301in}{2.880021in}}%
\pgfpathlineto{\pgfqpoint{3.329301in}{2.880021in}}%
\pgfpathlineto{\pgfqpoint{3.329301in}{2.882970in}}%
\pgfpathlineto{\pgfqpoint{3.333842in}{2.882970in}}%
\pgfpathlineto{\pgfqpoint{3.333842in}{2.880021in}}%
\pgfpathmoveto{\pgfqpoint{3.333842in}{2.877072in}}%
\pgfpathlineto{\pgfqpoint{3.333842in}{2.877072in}}%
\pgfpathlineto{\pgfqpoint{3.333842in}{2.880021in}}%
\pgfpathlineto{\pgfqpoint{3.338383in}{2.880021in}}%
\pgfpathlineto{\pgfqpoint{3.338383in}{2.877072in}}%
\pgfpathmoveto{\pgfqpoint{3.333842in}{2.880021in}}%
\pgfpathlineto{\pgfqpoint{3.333842in}{2.880021in}}%
\pgfpathlineto{\pgfqpoint{3.333842in}{2.882970in}}%
\pgfpathlineto{\pgfqpoint{3.338383in}{2.882970in}}%
\pgfpathlineto{\pgfqpoint{3.338383in}{2.880021in}}%
\pgfpathmoveto{\pgfqpoint{3.338383in}{2.871174in}}%
\pgfpathlineto{\pgfqpoint{3.338383in}{2.871174in}}%
\pgfpathlineto{\pgfqpoint{3.338383in}{2.874123in}}%
\pgfpathlineto{\pgfqpoint{3.342924in}{2.874123in}}%
\pgfpathlineto{\pgfqpoint{3.342924in}{2.871174in}}%
\pgfpathmoveto{\pgfqpoint{3.338383in}{2.874123in}}%
\pgfpathlineto{\pgfqpoint{3.338383in}{2.874123in}}%
\pgfpathlineto{\pgfqpoint{3.338383in}{2.877072in}}%
\pgfpathlineto{\pgfqpoint{3.342924in}{2.877072in}}%
\pgfpathlineto{\pgfqpoint{3.342924in}{2.874123in}}%
\pgfpathmoveto{\pgfqpoint{3.347465in}{2.859377in}}%
\pgfpathlineto{\pgfqpoint{3.347465in}{2.859377in}}%
\pgfpathlineto{\pgfqpoint{3.347465in}{2.862326in}}%
\pgfpathlineto{\pgfqpoint{3.352006in}{2.862326in}}%
\pgfpathlineto{\pgfqpoint{3.352006in}{2.859377in}}%
\pgfpathmoveto{\pgfqpoint{3.347465in}{2.862326in}}%
\pgfpathlineto{\pgfqpoint{3.347465in}{2.862326in}}%
\pgfpathlineto{\pgfqpoint{3.347465in}{2.865276in}}%
\pgfpathlineto{\pgfqpoint{3.352006in}{2.865276in}}%
\pgfpathlineto{\pgfqpoint{3.352006in}{2.862326in}}%
\pgfpathmoveto{\pgfqpoint{3.329301in}{2.882970in}}%
\pgfpathlineto{\pgfqpoint{3.329301in}{2.882970in}}%
\pgfpathlineto{\pgfqpoint{3.329301in}{2.885919in}}%
\pgfpathlineto{\pgfqpoint{3.333842in}{2.885919in}}%
\pgfpathlineto{\pgfqpoint{3.333842in}{2.882970in}}%
\pgfpathmoveto{\pgfqpoint{3.329301in}{2.885919in}}%
\pgfpathlineto{\pgfqpoint{3.329301in}{2.885919in}}%
\pgfpathlineto{\pgfqpoint{3.329301in}{2.888869in}}%
\pgfpathlineto{\pgfqpoint{3.333842in}{2.888869in}}%
\pgfpathlineto{\pgfqpoint{3.333842in}{2.885919in}}%
\pgfpathmoveto{\pgfqpoint{3.306596in}{2.909512in}}%
\pgfpathlineto{\pgfqpoint{3.306596in}{2.909512in}}%
\pgfpathlineto{\pgfqpoint{3.306596in}{2.912462in}}%
\pgfpathlineto{\pgfqpoint{3.311137in}{2.912462in}}%
\pgfpathlineto{\pgfqpoint{3.311137in}{2.909512in}}%
\pgfpathmoveto{\pgfqpoint{3.302055in}{2.915411in}}%
\pgfpathlineto{\pgfqpoint{3.302055in}{2.915411in}}%
\pgfpathlineto{\pgfqpoint{3.302055in}{2.918360in}}%
\pgfpathlineto{\pgfqpoint{3.306596in}{2.918360in}}%
\pgfpathlineto{\pgfqpoint{3.306596in}{2.915411in}}%
\pgfpathmoveto{\pgfqpoint{3.306596in}{2.912462in}}%
\pgfpathlineto{\pgfqpoint{3.306596in}{2.912462in}}%
\pgfpathlineto{\pgfqpoint{3.306596in}{2.915411in}}%
\pgfpathlineto{\pgfqpoint{3.311137in}{2.915411in}}%
\pgfpathlineto{\pgfqpoint{3.311137in}{2.912462in}}%
\pgfpathmoveto{\pgfqpoint{3.306596in}{2.915411in}}%
\pgfpathlineto{\pgfqpoint{3.306596in}{2.915411in}}%
\pgfpathlineto{\pgfqpoint{3.306596in}{2.918360in}}%
\pgfpathlineto{\pgfqpoint{3.311137in}{2.918360in}}%
\pgfpathlineto{\pgfqpoint{3.311137in}{2.915411in}}%
\pgfpathmoveto{\pgfqpoint{3.297514in}{2.921309in}}%
\pgfpathlineto{\pgfqpoint{3.297514in}{2.921309in}}%
\pgfpathlineto{\pgfqpoint{3.297514in}{2.924258in}}%
\pgfpathlineto{\pgfqpoint{3.302055in}{2.924258in}}%
\pgfpathlineto{\pgfqpoint{3.302055in}{2.921309in}}%
\pgfpathmoveto{\pgfqpoint{3.292973in}{2.927207in}}%
\pgfpathlineto{\pgfqpoint{3.292973in}{2.927207in}}%
\pgfpathlineto{\pgfqpoint{3.292973in}{2.930156in}}%
\pgfpathlineto{\pgfqpoint{3.297514in}{2.930156in}}%
\pgfpathlineto{\pgfqpoint{3.297514in}{2.927207in}}%
\pgfpathmoveto{\pgfqpoint{3.297514in}{2.924258in}}%
\pgfpathlineto{\pgfqpoint{3.297514in}{2.924258in}}%
\pgfpathlineto{\pgfqpoint{3.297514in}{2.927207in}}%
\pgfpathlineto{\pgfqpoint{3.302055in}{2.927207in}}%
\pgfpathlineto{\pgfqpoint{3.302055in}{2.924258in}}%
\pgfpathmoveto{\pgfqpoint{3.297514in}{2.927207in}}%
\pgfpathlineto{\pgfqpoint{3.297514in}{2.927207in}}%
\pgfpathlineto{\pgfqpoint{3.297514in}{2.930156in}}%
\pgfpathlineto{\pgfqpoint{3.302055in}{2.930156in}}%
\pgfpathlineto{\pgfqpoint{3.302055in}{2.927207in}}%
\pgfpathmoveto{\pgfqpoint{3.302055in}{2.918360in}}%
\pgfpathlineto{\pgfqpoint{3.302055in}{2.918360in}}%
\pgfpathlineto{\pgfqpoint{3.302055in}{2.921309in}}%
\pgfpathlineto{\pgfqpoint{3.306596in}{2.921309in}}%
\pgfpathlineto{\pgfqpoint{3.306596in}{2.918360in}}%
\pgfpathmoveto{\pgfqpoint{3.302055in}{2.921309in}}%
\pgfpathlineto{\pgfqpoint{3.302055in}{2.921309in}}%
\pgfpathlineto{\pgfqpoint{3.302055in}{2.924258in}}%
\pgfpathlineto{\pgfqpoint{3.306596in}{2.924258in}}%
\pgfpathlineto{\pgfqpoint{3.306596in}{2.921309in}}%
\pgfpathmoveto{\pgfqpoint{3.311137in}{2.906563in}}%
\pgfpathlineto{\pgfqpoint{3.311137in}{2.906563in}}%
\pgfpathlineto{\pgfqpoint{3.311137in}{2.909512in}}%
\pgfpathlineto{\pgfqpoint{3.315678in}{2.909512in}}%
\pgfpathlineto{\pgfqpoint{3.315678in}{2.906563in}}%
\pgfpathmoveto{\pgfqpoint{3.311137in}{2.909512in}}%
\pgfpathlineto{\pgfqpoint{3.311137in}{2.909512in}}%
\pgfpathlineto{\pgfqpoint{3.311137in}{2.912462in}}%
\pgfpathlineto{\pgfqpoint{3.315678in}{2.912462in}}%
\pgfpathlineto{\pgfqpoint{3.315678in}{2.909512in}}%
\pgfpathmoveto{\pgfqpoint{3.292973in}{2.930156in}}%
\pgfpathlineto{\pgfqpoint{3.292973in}{2.930156in}}%
\pgfpathlineto{\pgfqpoint{3.292973in}{2.933105in}}%
\pgfpathlineto{\pgfqpoint{3.297514in}{2.933105in}}%
\pgfpathlineto{\pgfqpoint{3.297514in}{2.930156in}}%
\pgfpathmoveto{\pgfqpoint{3.292973in}{2.933105in}}%
\pgfpathlineto{\pgfqpoint{3.292973in}{2.933105in}}%
\pgfpathlineto{\pgfqpoint{3.292973in}{2.936055in}}%
\pgfpathlineto{\pgfqpoint{3.297514in}{2.936055in}}%
\pgfpathlineto{\pgfqpoint{3.297514in}{2.933105in}}%
\pgfpathmoveto{\pgfqpoint{3.252104in}{2.980293in}}%
\pgfpathlineto{\pgfqpoint{3.252104in}{2.980293in}}%
\pgfpathlineto{\pgfqpoint{3.252104in}{2.983242in}}%
\pgfpathlineto{\pgfqpoint{3.256645in}{2.983242in}}%
\pgfpathlineto{\pgfqpoint{3.256645in}{2.980293in}}%
\pgfpathmoveto{\pgfqpoint{3.247563in}{2.986192in}}%
\pgfpathlineto{\pgfqpoint{3.247563in}{2.986192in}}%
\pgfpathlineto{\pgfqpoint{3.247563in}{2.989141in}}%
\pgfpathlineto{\pgfqpoint{3.252104in}{2.989141in}}%
\pgfpathlineto{\pgfqpoint{3.252104in}{2.986192in}}%
\pgfpathmoveto{\pgfqpoint{3.252104in}{2.983242in}}%
\pgfpathlineto{\pgfqpoint{3.252104in}{2.983242in}}%
\pgfpathlineto{\pgfqpoint{3.252104in}{2.986192in}}%
\pgfpathlineto{\pgfqpoint{3.256645in}{2.986192in}}%
\pgfpathlineto{\pgfqpoint{3.256645in}{2.983242in}}%
\pgfpathmoveto{\pgfqpoint{3.252104in}{2.986192in}}%
\pgfpathlineto{\pgfqpoint{3.252104in}{2.986192in}}%
\pgfpathlineto{\pgfqpoint{3.252104in}{2.989141in}}%
\pgfpathlineto{\pgfqpoint{3.256645in}{2.989141in}}%
\pgfpathlineto{\pgfqpoint{3.256645in}{2.986192in}}%
\pgfpathmoveto{\pgfqpoint{3.243022in}{2.992090in}}%
\pgfpathlineto{\pgfqpoint{3.243022in}{2.992090in}}%
\pgfpathlineto{\pgfqpoint{3.243022in}{2.995040in}}%
\pgfpathlineto{\pgfqpoint{3.247563in}{2.995040in}}%
\pgfpathlineto{\pgfqpoint{3.247563in}{2.992090in}}%
\pgfpathmoveto{\pgfqpoint{3.238481in}{2.997989in}}%
\pgfpathlineto{\pgfqpoint{3.238481in}{2.997989in}}%
\pgfpathlineto{\pgfqpoint{3.238481in}{3.000938in}}%
\pgfpathlineto{\pgfqpoint{3.243022in}{3.000938in}}%
\pgfpathlineto{\pgfqpoint{3.243022in}{2.997989in}}%
\pgfpathmoveto{\pgfqpoint{3.243022in}{2.995040in}}%
\pgfpathlineto{\pgfqpoint{3.243022in}{2.995040in}}%
\pgfpathlineto{\pgfqpoint{3.243022in}{2.997989in}}%
\pgfpathlineto{\pgfqpoint{3.247563in}{2.997989in}}%
\pgfpathlineto{\pgfqpoint{3.247563in}{2.995040in}}%
\pgfpathmoveto{\pgfqpoint{3.243022in}{2.997989in}}%
\pgfpathlineto{\pgfqpoint{3.243022in}{2.997989in}}%
\pgfpathlineto{\pgfqpoint{3.243022in}{3.000938in}}%
\pgfpathlineto{\pgfqpoint{3.247563in}{3.000938in}}%
\pgfpathlineto{\pgfqpoint{3.247563in}{2.997989in}}%
\pgfpathmoveto{\pgfqpoint{3.247563in}{2.989141in}}%
\pgfpathlineto{\pgfqpoint{3.247563in}{2.989141in}}%
\pgfpathlineto{\pgfqpoint{3.247563in}{2.992090in}}%
\pgfpathlineto{\pgfqpoint{3.252104in}{2.992090in}}%
\pgfpathlineto{\pgfqpoint{3.252104in}{2.989141in}}%
\pgfpathmoveto{\pgfqpoint{3.247563in}{2.992090in}}%
\pgfpathlineto{\pgfqpoint{3.247563in}{2.992090in}}%
\pgfpathlineto{\pgfqpoint{3.247563in}{2.995040in}}%
\pgfpathlineto{\pgfqpoint{3.252104in}{2.995040in}}%
\pgfpathlineto{\pgfqpoint{3.252104in}{2.992090in}}%
\pgfpathmoveto{\pgfqpoint{3.270268in}{2.956699in}}%
\pgfpathlineto{\pgfqpoint{3.270268in}{2.956699in}}%
\pgfpathlineto{\pgfqpoint{3.270268in}{2.959648in}}%
\pgfpathlineto{\pgfqpoint{3.274809in}{2.959648in}}%
\pgfpathlineto{\pgfqpoint{3.274809in}{2.956699in}}%
\pgfpathmoveto{\pgfqpoint{3.265727in}{2.962597in}}%
\pgfpathlineto{\pgfqpoint{3.265727in}{2.962597in}}%
\pgfpathlineto{\pgfqpoint{3.265727in}{2.965547in}}%
\pgfpathlineto{\pgfqpoint{3.270268in}{2.965547in}}%
\pgfpathlineto{\pgfqpoint{3.270268in}{2.962597in}}%
\pgfpathmoveto{\pgfqpoint{3.270268in}{2.959648in}}%
\pgfpathlineto{\pgfqpoint{3.270268in}{2.959648in}}%
\pgfpathlineto{\pgfqpoint{3.270268in}{2.962597in}}%
\pgfpathlineto{\pgfqpoint{3.274809in}{2.962597in}}%
\pgfpathlineto{\pgfqpoint{3.274809in}{2.959648in}}%
\pgfpathmoveto{\pgfqpoint{3.270268in}{2.962597in}}%
\pgfpathlineto{\pgfqpoint{3.270268in}{2.962597in}}%
\pgfpathlineto{\pgfqpoint{3.270268in}{2.965547in}}%
\pgfpathlineto{\pgfqpoint{3.274809in}{2.965547in}}%
\pgfpathlineto{\pgfqpoint{3.274809in}{2.962597in}}%
\pgfpathmoveto{\pgfqpoint{3.261186in}{2.968496in}}%
\pgfpathlineto{\pgfqpoint{3.261186in}{2.968496in}}%
\pgfpathlineto{\pgfqpoint{3.261186in}{2.971445in}}%
\pgfpathlineto{\pgfqpoint{3.265727in}{2.971445in}}%
\pgfpathlineto{\pgfqpoint{3.265727in}{2.968496in}}%
\pgfpathmoveto{\pgfqpoint{3.256645in}{2.974394in}}%
\pgfpathlineto{\pgfqpoint{3.256645in}{2.974394in}}%
\pgfpathlineto{\pgfqpoint{3.256645in}{2.977344in}}%
\pgfpathlineto{\pgfqpoint{3.261186in}{2.977344in}}%
\pgfpathlineto{\pgfqpoint{3.261186in}{2.974394in}}%
\pgfpathmoveto{\pgfqpoint{3.261186in}{2.971445in}}%
\pgfpathlineto{\pgfqpoint{3.261186in}{2.971445in}}%
\pgfpathlineto{\pgfqpoint{3.261186in}{2.974394in}}%
\pgfpathlineto{\pgfqpoint{3.265727in}{2.974394in}}%
\pgfpathlineto{\pgfqpoint{3.265727in}{2.971445in}}%
\pgfpathmoveto{\pgfqpoint{3.261186in}{2.974394in}}%
\pgfpathlineto{\pgfqpoint{3.261186in}{2.974394in}}%
\pgfpathlineto{\pgfqpoint{3.261186in}{2.977344in}}%
\pgfpathlineto{\pgfqpoint{3.265727in}{2.977344in}}%
\pgfpathlineto{\pgfqpoint{3.265727in}{2.974394in}}%
\pgfpathmoveto{\pgfqpoint{3.265727in}{2.965547in}}%
\pgfpathlineto{\pgfqpoint{3.265727in}{2.965547in}}%
\pgfpathlineto{\pgfqpoint{3.265727in}{2.968496in}}%
\pgfpathlineto{\pgfqpoint{3.270268in}{2.968496in}}%
\pgfpathlineto{\pgfqpoint{3.270268in}{2.965547in}}%
\pgfpathmoveto{\pgfqpoint{3.265727in}{2.968496in}}%
\pgfpathlineto{\pgfqpoint{3.265727in}{2.968496in}}%
\pgfpathlineto{\pgfqpoint{3.265727in}{2.971445in}}%
\pgfpathlineto{\pgfqpoint{3.270268in}{2.971445in}}%
\pgfpathlineto{\pgfqpoint{3.270268in}{2.968496in}}%
\pgfpathmoveto{\pgfqpoint{3.274809in}{2.953749in}}%
\pgfpathlineto{\pgfqpoint{3.274809in}{2.953749in}}%
\pgfpathlineto{\pgfqpoint{3.274809in}{2.956699in}}%
\pgfpathlineto{\pgfqpoint{3.279350in}{2.956699in}}%
\pgfpathlineto{\pgfqpoint{3.279350in}{2.953749in}}%
\pgfpathmoveto{\pgfqpoint{3.274809in}{2.956699in}}%
\pgfpathlineto{\pgfqpoint{3.274809in}{2.956699in}}%
\pgfpathlineto{\pgfqpoint{3.274809in}{2.959648in}}%
\pgfpathlineto{\pgfqpoint{3.279350in}{2.959648in}}%
\pgfpathlineto{\pgfqpoint{3.279350in}{2.956699in}}%
\pgfpathmoveto{\pgfqpoint{3.256645in}{2.977344in}}%
\pgfpathlineto{\pgfqpoint{3.256645in}{2.977344in}}%
\pgfpathlineto{\pgfqpoint{3.256645in}{2.980293in}}%
\pgfpathlineto{\pgfqpoint{3.261186in}{2.980293in}}%
\pgfpathlineto{\pgfqpoint{3.261186in}{2.977344in}}%
\pgfpathmoveto{\pgfqpoint{3.256645in}{2.980293in}}%
\pgfpathlineto{\pgfqpoint{3.256645in}{2.980293in}}%
\pgfpathlineto{\pgfqpoint{3.256645in}{2.983242in}}%
\pgfpathlineto{\pgfqpoint{3.261186in}{2.983242in}}%
\pgfpathlineto{\pgfqpoint{3.261186in}{2.980293in}}%
\pgfpathmoveto{\pgfqpoint{3.233940in}{3.003888in}}%
\pgfpathlineto{\pgfqpoint{3.233940in}{3.003888in}}%
\pgfpathlineto{\pgfqpoint{3.233940in}{3.006837in}}%
\pgfpathlineto{\pgfqpoint{3.238481in}{3.006837in}}%
\pgfpathlineto{\pgfqpoint{3.238481in}{3.003888in}}%
\pgfpathmoveto{\pgfqpoint{3.229399in}{3.009786in}}%
\pgfpathlineto{\pgfqpoint{3.229399in}{3.009786in}}%
\pgfpathlineto{\pgfqpoint{3.229399in}{3.012735in}}%
\pgfpathlineto{\pgfqpoint{3.233940in}{3.012735in}}%
\pgfpathlineto{\pgfqpoint{3.233940in}{3.009786in}}%
\pgfpathmoveto{\pgfqpoint{3.233940in}{3.006837in}}%
\pgfpathlineto{\pgfqpoint{3.233940in}{3.006837in}}%
\pgfpathlineto{\pgfqpoint{3.233940in}{3.009786in}}%
\pgfpathlineto{\pgfqpoint{3.238481in}{3.009786in}}%
\pgfpathlineto{\pgfqpoint{3.238481in}{3.006837in}}%
\pgfpathmoveto{\pgfqpoint{3.233940in}{3.009786in}}%
\pgfpathlineto{\pgfqpoint{3.233940in}{3.009786in}}%
\pgfpathlineto{\pgfqpoint{3.233940in}{3.012735in}}%
\pgfpathlineto{\pgfqpoint{3.238481in}{3.012735in}}%
\pgfpathlineto{\pgfqpoint{3.238481in}{3.009786in}}%
\pgfpathmoveto{\pgfqpoint{3.224858in}{3.015685in}}%
\pgfpathlineto{\pgfqpoint{3.224858in}{3.015685in}}%
\pgfpathlineto{\pgfqpoint{3.224858in}{3.018634in}}%
\pgfpathlineto{\pgfqpoint{3.229399in}{3.018634in}}%
\pgfpathlineto{\pgfqpoint{3.229399in}{3.015685in}}%
\pgfpathmoveto{\pgfqpoint{3.220317in}{3.021583in}}%
\pgfpathlineto{\pgfqpoint{3.220317in}{3.021583in}}%
\pgfpathlineto{\pgfqpoint{3.220317in}{3.024533in}}%
\pgfpathlineto{\pgfqpoint{3.224858in}{3.024533in}}%
\pgfpathlineto{\pgfqpoint{3.224858in}{3.021583in}}%
\pgfpathmoveto{\pgfqpoint{3.224858in}{3.018634in}}%
\pgfpathlineto{\pgfqpoint{3.224858in}{3.018634in}}%
\pgfpathlineto{\pgfqpoint{3.224858in}{3.021583in}}%
\pgfpathlineto{\pgfqpoint{3.229399in}{3.021583in}}%
\pgfpathlineto{\pgfqpoint{3.229399in}{3.018634in}}%
\pgfpathmoveto{\pgfqpoint{3.224858in}{3.021583in}}%
\pgfpathlineto{\pgfqpoint{3.224858in}{3.021583in}}%
\pgfpathlineto{\pgfqpoint{3.224858in}{3.024533in}}%
\pgfpathlineto{\pgfqpoint{3.229399in}{3.024533in}}%
\pgfpathlineto{\pgfqpoint{3.229399in}{3.021583in}}%
\pgfpathmoveto{\pgfqpoint{3.229399in}{3.012735in}}%
\pgfpathlineto{\pgfqpoint{3.229399in}{3.012735in}}%
\pgfpathlineto{\pgfqpoint{3.229399in}{3.015685in}}%
\pgfpathlineto{\pgfqpoint{3.233940in}{3.015685in}}%
\pgfpathlineto{\pgfqpoint{3.233940in}{3.012735in}}%
\pgfpathmoveto{\pgfqpoint{3.229399in}{3.015685in}}%
\pgfpathlineto{\pgfqpoint{3.229399in}{3.015685in}}%
\pgfpathlineto{\pgfqpoint{3.229399in}{3.018634in}}%
\pgfpathlineto{\pgfqpoint{3.233940in}{3.018634in}}%
\pgfpathlineto{\pgfqpoint{3.233940in}{3.015685in}}%
\pgfpathmoveto{\pgfqpoint{3.238481in}{3.000938in}}%
\pgfpathlineto{\pgfqpoint{3.238481in}{3.000938in}}%
\pgfpathlineto{\pgfqpoint{3.238481in}{3.003888in}}%
\pgfpathlineto{\pgfqpoint{3.243022in}{3.003888in}}%
\pgfpathlineto{\pgfqpoint{3.243022in}{3.000938in}}%
\pgfpathmoveto{\pgfqpoint{3.238481in}{3.003888in}}%
\pgfpathlineto{\pgfqpoint{3.238481in}{3.003888in}}%
\pgfpathlineto{\pgfqpoint{3.238481in}{3.006837in}}%
\pgfpathlineto{\pgfqpoint{3.243022in}{3.006837in}}%
\pgfpathlineto{\pgfqpoint{3.243022in}{3.003888in}}%
\pgfpathmoveto{\pgfqpoint{3.220317in}{3.024533in}}%
\pgfpathlineto{\pgfqpoint{3.220317in}{3.024533in}}%
\pgfpathlineto{\pgfqpoint{3.220317in}{3.027482in}}%
\pgfpathlineto{\pgfqpoint{3.224858in}{3.027482in}}%
\pgfpathlineto{\pgfqpoint{3.224858in}{3.024533in}}%
\pgfpathmoveto{\pgfqpoint{3.220317in}{3.027482in}}%
\pgfpathlineto{\pgfqpoint{3.220317in}{3.027482in}}%
\pgfpathlineto{\pgfqpoint{3.220317in}{3.030431in}}%
\pgfpathlineto{\pgfqpoint{3.224858in}{3.030431in}}%
\pgfpathlineto{\pgfqpoint{3.224858in}{3.027482in}}%
\pgfpathmoveto{\pgfqpoint{3.506398in}{2.649981in}}%
\pgfpathlineto{\pgfqpoint{3.506398in}{2.649981in}}%
\pgfpathlineto{\pgfqpoint{3.506398in}{2.652930in}}%
\pgfpathlineto{\pgfqpoint{3.510939in}{2.652930in}}%
\pgfpathlineto{\pgfqpoint{3.510939in}{2.649981in}}%
\pgfpathmoveto{\pgfqpoint{3.501857in}{2.655879in}}%
\pgfpathlineto{\pgfqpoint{3.501857in}{2.655879in}}%
\pgfpathlineto{\pgfqpoint{3.501857in}{2.658829in}}%
\pgfpathlineto{\pgfqpoint{3.506398in}{2.658829in}}%
\pgfpathlineto{\pgfqpoint{3.506398in}{2.655879in}}%
\pgfpathmoveto{\pgfqpoint{3.506398in}{2.652930in}}%
\pgfpathlineto{\pgfqpoint{3.506398in}{2.652930in}}%
\pgfpathlineto{\pgfqpoint{3.506398in}{2.655879in}}%
\pgfpathlineto{\pgfqpoint{3.510939in}{2.655879in}}%
\pgfpathlineto{\pgfqpoint{3.510939in}{2.652930in}}%
\pgfpathmoveto{\pgfqpoint{3.506398in}{2.655879in}}%
\pgfpathlineto{\pgfqpoint{3.506398in}{2.655879in}}%
\pgfpathlineto{\pgfqpoint{3.506398in}{2.658829in}}%
\pgfpathlineto{\pgfqpoint{3.510939in}{2.658829in}}%
\pgfpathlineto{\pgfqpoint{3.510939in}{2.655879in}}%
\pgfpathmoveto{\pgfqpoint{3.497316in}{2.661778in}}%
\pgfpathlineto{\pgfqpoint{3.497316in}{2.661778in}}%
\pgfpathlineto{\pgfqpoint{3.497316in}{2.664727in}}%
\pgfpathlineto{\pgfqpoint{3.501857in}{2.664727in}}%
\pgfpathlineto{\pgfqpoint{3.501857in}{2.661778in}}%
\pgfpathmoveto{\pgfqpoint{3.492775in}{2.667676in}}%
\pgfpathlineto{\pgfqpoint{3.492775in}{2.667676in}}%
\pgfpathlineto{\pgfqpoint{3.492775in}{2.670625in}}%
\pgfpathlineto{\pgfqpoint{3.497316in}{2.670625in}}%
\pgfpathlineto{\pgfqpoint{3.497316in}{2.667676in}}%
\pgfpathmoveto{\pgfqpoint{3.497316in}{2.664727in}}%
\pgfpathlineto{\pgfqpoint{3.497316in}{2.664727in}}%
\pgfpathlineto{\pgfqpoint{3.497316in}{2.667676in}}%
\pgfpathlineto{\pgfqpoint{3.501857in}{2.667676in}}%
\pgfpathlineto{\pgfqpoint{3.501857in}{2.664727in}}%
\pgfpathmoveto{\pgfqpoint{3.497316in}{2.667676in}}%
\pgfpathlineto{\pgfqpoint{3.497316in}{2.667676in}}%
\pgfpathlineto{\pgfqpoint{3.497316in}{2.670625in}}%
\pgfpathlineto{\pgfqpoint{3.501857in}{2.670625in}}%
\pgfpathlineto{\pgfqpoint{3.501857in}{2.667676in}}%
\pgfpathmoveto{\pgfqpoint{3.501857in}{2.658829in}}%
\pgfpathlineto{\pgfqpoint{3.501857in}{2.658829in}}%
\pgfpathlineto{\pgfqpoint{3.501857in}{2.661778in}}%
\pgfpathlineto{\pgfqpoint{3.506398in}{2.661778in}}%
\pgfpathlineto{\pgfqpoint{3.506398in}{2.658829in}}%
\pgfpathmoveto{\pgfqpoint{3.501857in}{2.661778in}}%
\pgfpathlineto{\pgfqpoint{3.501857in}{2.661778in}}%
\pgfpathlineto{\pgfqpoint{3.501857in}{2.664727in}}%
\pgfpathlineto{\pgfqpoint{3.506398in}{2.664727in}}%
\pgfpathlineto{\pgfqpoint{3.506398in}{2.661778in}}%
\pgfpathmoveto{\pgfqpoint{3.433743in}{2.744354in}}%
\pgfpathlineto{\pgfqpoint{3.433743in}{2.744354in}}%
\pgfpathlineto{\pgfqpoint{3.433743in}{2.747303in}}%
\pgfpathlineto{\pgfqpoint{3.438284in}{2.747303in}}%
\pgfpathlineto{\pgfqpoint{3.438284in}{2.744354in}}%
\pgfpathmoveto{\pgfqpoint{3.429202in}{2.750252in}}%
\pgfpathlineto{\pgfqpoint{3.429202in}{2.750252in}}%
\pgfpathlineto{\pgfqpoint{3.429202in}{2.753201in}}%
\pgfpathlineto{\pgfqpoint{3.433743in}{2.753201in}}%
\pgfpathlineto{\pgfqpoint{3.433743in}{2.750252in}}%
\pgfpathmoveto{\pgfqpoint{3.433743in}{2.747303in}}%
\pgfpathlineto{\pgfqpoint{3.433743in}{2.747303in}}%
\pgfpathlineto{\pgfqpoint{3.433743in}{2.750252in}}%
\pgfpathlineto{\pgfqpoint{3.438284in}{2.750252in}}%
\pgfpathlineto{\pgfqpoint{3.438284in}{2.747303in}}%
\pgfpathmoveto{\pgfqpoint{3.433743in}{2.750252in}}%
\pgfpathlineto{\pgfqpoint{3.433743in}{2.750252in}}%
\pgfpathlineto{\pgfqpoint{3.433743in}{2.753201in}}%
\pgfpathlineto{\pgfqpoint{3.438284in}{2.753201in}}%
\pgfpathlineto{\pgfqpoint{3.438284in}{2.750252in}}%
\pgfpathmoveto{\pgfqpoint{3.424661in}{2.756150in}}%
\pgfpathlineto{\pgfqpoint{3.424661in}{2.756150in}}%
\pgfpathlineto{\pgfqpoint{3.424661in}{2.759099in}}%
\pgfpathlineto{\pgfqpoint{3.429202in}{2.759099in}}%
\pgfpathlineto{\pgfqpoint{3.429202in}{2.756150in}}%
\pgfpathmoveto{\pgfqpoint{3.420120in}{2.762048in}}%
\pgfpathlineto{\pgfqpoint{3.420120in}{2.762048in}}%
\pgfpathlineto{\pgfqpoint{3.420120in}{2.764998in}}%
\pgfpathlineto{\pgfqpoint{3.424661in}{2.764998in}}%
\pgfpathlineto{\pgfqpoint{3.424661in}{2.762048in}}%
\pgfpathmoveto{\pgfqpoint{3.424661in}{2.759099in}}%
\pgfpathlineto{\pgfqpoint{3.424661in}{2.759099in}}%
\pgfpathlineto{\pgfqpoint{3.424661in}{2.762048in}}%
\pgfpathlineto{\pgfqpoint{3.429202in}{2.762048in}}%
\pgfpathlineto{\pgfqpoint{3.429202in}{2.759099in}}%
\pgfpathmoveto{\pgfqpoint{3.424661in}{2.762048in}}%
\pgfpathlineto{\pgfqpoint{3.424661in}{2.762048in}}%
\pgfpathlineto{\pgfqpoint{3.424661in}{2.764998in}}%
\pgfpathlineto{\pgfqpoint{3.429202in}{2.764998in}}%
\pgfpathlineto{\pgfqpoint{3.429202in}{2.762048in}}%
\pgfpathmoveto{\pgfqpoint{3.429202in}{2.753201in}}%
\pgfpathlineto{\pgfqpoint{3.429202in}{2.753201in}}%
\pgfpathlineto{\pgfqpoint{3.429202in}{2.756150in}}%
\pgfpathlineto{\pgfqpoint{3.433743in}{2.756150in}}%
\pgfpathlineto{\pgfqpoint{3.433743in}{2.753201in}}%
\pgfpathmoveto{\pgfqpoint{3.429202in}{2.756150in}}%
\pgfpathlineto{\pgfqpoint{3.429202in}{2.756150in}}%
\pgfpathlineto{\pgfqpoint{3.429202in}{2.759099in}}%
\pgfpathlineto{\pgfqpoint{3.433743in}{2.759099in}}%
\pgfpathlineto{\pgfqpoint{3.433743in}{2.756150in}}%
\pgfpathmoveto{\pgfqpoint{3.470071in}{2.697168in}}%
\pgfpathlineto{\pgfqpoint{3.470071in}{2.697168in}}%
\pgfpathlineto{\pgfqpoint{3.470071in}{2.700117in}}%
\pgfpathlineto{\pgfqpoint{3.474612in}{2.700117in}}%
\pgfpathlineto{\pgfqpoint{3.474612in}{2.697168in}}%
\pgfpathmoveto{\pgfqpoint{3.465530in}{2.703066in}}%
\pgfpathlineto{\pgfqpoint{3.465530in}{2.703066in}}%
\pgfpathlineto{\pgfqpoint{3.465530in}{2.706015in}}%
\pgfpathlineto{\pgfqpoint{3.470071in}{2.706015in}}%
\pgfpathlineto{\pgfqpoint{3.470071in}{2.703066in}}%
\pgfpathmoveto{\pgfqpoint{3.470071in}{2.700117in}}%
\pgfpathlineto{\pgfqpoint{3.470071in}{2.700117in}}%
\pgfpathlineto{\pgfqpoint{3.470071in}{2.703066in}}%
\pgfpathlineto{\pgfqpoint{3.474612in}{2.703066in}}%
\pgfpathlineto{\pgfqpoint{3.474612in}{2.700117in}}%
\pgfpathmoveto{\pgfqpoint{3.470071in}{2.703066in}}%
\pgfpathlineto{\pgfqpoint{3.470071in}{2.703066in}}%
\pgfpathlineto{\pgfqpoint{3.470071in}{2.706015in}}%
\pgfpathlineto{\pgfqpoint{3.474612in}{2.706015in}}%
\pgfpathlineto{\pgfqpoint{3.474612in}{2.703066in}}%
\pgfpathmoveto{\pgfqpoint{3.460989in}{2.708964in}}%
\pgfpathlineto{\pgfqpoint{3.460989in}{2.708964in}}%
\pgfpathlineto{\pgfqpoint{3.460989in}{2.711913in}}%
\pgfpathlineto{\pgfqpoint{3.465530in}{2.711913in}}%
\pgfpathlineto{\pgfqpoint{3.465530in}{2.708964in}}%
\pgfpathmoveto{\pgfqpoint{3.456448in}{2.714862in}}%
\pgfpathlineto{\pgfqpoint{3.456448in}{2.714862in}}%
\pgfpathlineto{\pgfqpoint{3.456448in}{2.717812in}}%
\pgfpathlineto{\pgfqpoint{3.460989in}{2.717812in}}%
\pgfpathlineto{\pgfqpoint{3.460989in}{2.714862in}}%
\pgfpathmoveto{\pgfqpoint{3.460989in}{2.711913in}}%
\pgfpathlineto{\pgfqpoint{3.460989in}{2.711913in}}%
\pgfpathlineto{\pgfqpoint{3.460989in}{2.714862in}}%
\pgfpathlineto{\pgfqpoint{3.465530in}{2.714862in}}%
\pgfpathlineto{\pgfqpoint{3.465530in}{2.711913in}}%
\pgfpathmoveto{\pgfqpoint{3.460989in}{2.714862in}}%
\pgfpathlineto{\pgfqpoint{3.460989in}{2.714862in}}%
\pgfpathlineto{\pgfqpoint{3.460989in}{2.717812in}}%
\pgfpathlineto{\pgfqpoint{3.465530in}{2.717812in}}%
\pgfpathlineto{\pgfqpoint{3.465530in}{2.714862in}}%
\pgfpathmoveto{\pgfqpoint{3.465530in}{2.706015in}}%
\pgfpathlineto{\pgfqpoint{3.465530in}{2.706015in}}%
\pgfpathlineto{\pgfqpoint{3.465530in}{2.708964in}}%
\pgfpathlineto{\pgfqpoint{3.470071in}{2.708964in}}%
\pgfpathlineto{\pgfqpoint{3.470071in}{2.706015in}}%
\pgfpathmoveto{\pgfqpoint{3.465530in}{2.708964in}}%
\pgfpathlineto{\pgfqpoint{3.465530in}{2.708964in}}%
\pgfpathlineto{\pgfqpoint{3.465530in}{2.711913in}}%
\pgfpathlineto{\pgfqpoint{3.470071in}{2.711913in}}%
\pgfpathlineto{\pgfqpoint{3.470071in}{2.708964in}}%
\pgfpathmoveto{\pgfqpoint{3.488234in}{2.673575in}}%
\pgfpathlineto{\pgfqpoint{3.488234in}{2.673575in}}%
\pgfpathlineto{\pgfqpoint{3.488234in}{2.676524in}}%
\pgfpathlineto{\pgfqpoint{3.492775in}{2.676524in}}%
\pgfpathlineto{\pgfqpoint{3.492775in}{2.673575in}}%
\pgfpathmoveto{\pgfqpoint{3.483693in}{2.679473in}}%
\pgfpathlineto{\pgfqpoint{3.483693in}{2.679473in}}%
\pgfpathlineto{\pgfqpoint{3.483693in}{2.682422in}}%
\pgfpathlineto{\pgfqpoint{3.488234in}{2.682422in}}%
\pgfpathlineto{\pgfqpoint{3.488234in}{2.679473in}}%
\pgfpathmoveto{\pgfqpoint{3.488234in}{2.676524in}}%
\pgfpathlineto{\pgfqpoint{3.488234in}{2.676524in}}%
\pgfpathlineto{\pgfqpoint{3.488234in}{2.679473in}}%
\pgfpathlineto{\pgfqpoint{3.492775in}{2.679473in}}%
\pgfpathlineto{\pgfqpoint{3.492775in}{2.676524in}}%
\pgfpathmoveto{\pgfqpoint{3.488234in}{2.679473in}}%
\pgfpathlineto{\pgfqpoint{3.488234in}{2.679473in}}%
\pgfpathlineto{\pgfqpoint{3.488234in}{2.682422in}}%
\pgfpathlineto{\pgfqpoint{3.492775in}{2.682422in}}%
\pgfpathlineto{\pgfqpoint{3.492775in}{2.679473in}}%
\pgfpathmoveto{\pgfqpoint{3.479152in}{2.685371in}}%
\pgfpathlineto{\pgfqpoint{3.479152in}{2.685371in}}%
\pgfpathlineto{\pgfqpoint{3.479152in}{2.688320in}}%
\pgfpathlineto{\pgfqpoint{3.483693in}{2.688320in}}%
\pgfpathlineto{\pgfqpoint{3.483693in}{2.685371in}}%
\pgfpathmoveto{\pgfqpoint{3.474612in}{2.691269in}}%
\pgfpathlineto{\pgfqpoint{3.474612in}{2.691269in}}%
\pgfpathlineto{\pgfqpoint{3.474612in}{2.694218in}}%
\pgfpathlineto{\pgfqpoint{3.479152in}{2.694218in}}%
\pgfpathlineto{\pgfqpoint{3.479152in}{2.691269in}}%
\pgfpathmoveto{\pgfqpoint{3.479152in}{2.688320in}}%
\pgfpathlineto{\pgfqpoint{3.479152in}{2.688320in}}%
\pgfpathlineto{\pgfqpoint{3.479152in}{2.691269in}}%
\pgfpathlineto{\pgfqpoint{3.483693in}{2.691269in}}%
\pgfpathlineto{\pgfqpoint{3.483693in}{2.688320in}}%
\pgfpathmoveto{\pgfqpoint{3.479152in}{2.691269in}}%
\pgfpathlineto{\pgfqpoint{3.479152in}{2.691269in}}%
\pgfpathlineto{\pgfqpoint{3.479152in}{2.694218in}}%
\pgfpathlineto{\pgfqpoint{3.483693in}{2.694218in}}%
\pgfpathlineto{\pgfqpoint{3.483693in}{2.691269in}}%
\pgfpathmoveto{\pgfqpoint{3.483693in}{2.682422in}}%
\pgfpathlineto{\pgfqpoint{3.483693in}{2.682422in}}%
\pgfpathlineto{\pgfqpoint{3.483693in}{2.685371in}}%
\pgfpathlineto{\pgfqpoint{3.488234in}{2.685371in}}%
\pgfpathlineto{\pgfqpoint{3.488234in}{2.682422in}}%
\pgfpathmoveto{\pgfqpoint{3.483693in}{2.685371in}}%
\pgfpathlineto{\pgfqpoint{3.483693in}{2.685371in}}%
\pgfpathlineto{\pgfqpoint{3.483693in}{2.688320in}}%
\pgfpathlineto{\pgfqpoint{3.488234in}{2.688320in}}%
\pgfpathlineto{\pgfqpoint{3.488234in}{2.685371in}}%
\pgfpathmoveto{\pgfqpoint{3.492775in}{2.670625in}}%
\pgfpathlineto{\pgfqpoint{3.492775in}{2.670625in}}%
\pgfpathlineto{\pgfqpoint{3.492775in}{2.673575in}}%
\pgfpathlineto{\pgfqpoint{3.497316in}{2.673575in}}%
\pgfpathlineto{\pgfqpoint{3.497316in}{2.670625in}}%
\pgfpathmoveto{\pgfqpoint{3.492775in}{2.673575in}}%
\pgfpathlineto{\pgfqpoint{3.492775in}{2.673575in}}%
\pgfpathlineto{\pgfqpoint{3.492775in}{2.676524in}}%
\pgfpathlineto{\pgfqpoint{3.497316in}{2.676524in}}%
\pgfpathlineto{\pgfqpoint{3.497316in}{2.673575in}}%
\pgfpathmoveto{\pgfqpoint{3.474612in}{2.694218in}}%
\pgfpathlineto{\pgfqpoint{3.474612in}{2.694218in}}%
\pgfpathlineto{\pgfqpoint{3.474612in}{2.697168in}}%
\pgfpathlineto{\pgfqpoint{3.479152in}{2.697168in}}%
\pgfpathlineto{\pgfqpoint{3.479152in}{2.694218in}}%
\pgfpathmoveto{\pgfqpoint{3.474612in}{2.697168in}}%
\pgfpathlineto{\pgfqpoint{3.474612in}{2.697168in}}%
\pgfpathlineto{\pgfqpoint{3.474612in}{2.700117in}}%
\pgfpathlineto{\pgfqpoint{3.479152in}{2.700117in}}%
\pgfpathlineto{\pgfqpoint{3.479152in}{2.697168in}}%
\pgfpathmoveto{\pgfqpoint{3.451907in}{2.720761in}}%
\pgfpathlineto{\pgfqpoint{3.451907in}{2.720761in}}%
\pgfpathlineto{\pgfqpoint{3.451907in}{2.723710in}}%
\pgfpathlineto{\pgfqpoint{3.456448in}{2.723710in}}%
\pgfpathlineto{\pgfqpoint{3.456448in}{2.720761in}}%
\pgfpathmoveto{\pgfqpoint{3.447366in}{2.726659in}}%
\pgfpathlineto{\pgfqpoint{3.447366in}{2.726659in}}%
\pgfpathlineto{\pgfqpoint{3.447366in}{2.729608in}}%
\pgfpathlineto{\pgfqpoint{3.451907in}{2.729608in}}%
\pgfpathlineto{\pgfqpoint{3.451907in}{2.726659in}}%
\pgfpathmoveto{\pgfqpoint{3.451907in}{2.723710in}}%
\pgfpathlineto{\pgfqpoint{3.451907in}{2.723710in}}%
\pgfpathlineto{\pgfqpoint{3.451907in}{2.726659in}}%
\pgfpathlineto{\pgfqpoint{3.456448in}{2.726659in}}%
\pgfpathlineto{\pgfqpoint{3.456448in}{2.723710in}}%
\pgfpathmoveto{\pgfqpoint{3.451907in}{2.726659in}}%
\pgfpathlineto{\pgfqpoint{3.451907in}{2.726659in}}%
\pgfpathlineto{\pgfqpoint{3.451907in}{2.729608in}}%
\pgfpathlineto{\pgfqpoint{3.456448in}{2.729608in}}%
\pgfpathlineto{\pgfqpoint{3.456448in}{2.726659in}}%
\pgfpathmoveto{\pgfqpoint{3.442825in}{2.732557in}}%
\pgfpathlineto{\pgfqpoint{3.442825in}{2.732557in}}%
\pgfpathlineto{\pgfqpoint{3.442825in}{2.735506in}}%
\pgfpathlineto{\pgfqpoint{3.447366in}{2.735506in}}%
\pgfpathlineto{\pgfqpoint{3.447366in}{2.732557in}}%
\pgfpathmoveto{\pgfqpoint{3.438284in}{2.738455in}}%
\pgfpathlineto{\pgfqpoint{3.438284in}{2.738455in}}%
\pgfpathlineto{\pgfqpoint{3.438284in}{2.741405in}}%
\pgfpathlineto{\pgfqpoint{3.442825in}{2.741405in}}%
\pgfpathlineto{\pgfqpoint{3.442825in}{2.738455in}}%
\pgfpathmoveto{\pgfqpoint{3.442825in}{2.735506in}}%
\pgfpathlineto{\pgfqpoint{3.442825in}{2.735506in}}%
\pgfpathlineto{\pgfqpoint{3.442825in}{2.738455in}}%
\pgfpathlineto{\pgfqpoint{3.447366in}{2.738455in}}%
\pgfpathlineto{\pgfqpoint{3.447366in}{2.735506in}}%
\pgfpathmoveto{\pgfqpoint{3.442825in}{2.738455in}}%
\pgfpathlineto{\pgfqpoint{3.442825in}{2.738455in}}%
\pgfpathlineto{\pgfqpoint{3.442825in}{2.741405in}}%
\pgfpathlineto{\pgfqpoint{3.447366in}{2.741405in}}%
\pgfpathlineto{\pgfqpoint{3.447366in}{2.738455in}}%
\pgfpathmoveto{\pgfqpoint{3.447366in}{2.729608in}}%
\pgfpathlineto{\pgfqpoint{3.447366in}{2.729608in}}%
\pgfpathlineto{\pgfqpoint{3.447366in}{2.732557in}}%
\pgfpathlineto{\pgfqpoint{3.451907in}{2.732557in}}%
\pgfpathlineto{\pgfqpoint{3.451907in}{2.729608in}}%
\pgfpathmoveto{\pgfqpoint{3.447366in}{2.732557in}}%
\pgfpathlineto{\pgfqpoint{3.447366in}{2.732557in}}%
\pgfpathlineto{\pgfqpoint{3.447366in}{2.735506in}}%
\pgfpathlineto{\pgfqpoint{3.451907in}{2.735506in}}%
\pgfpathlineto{\pgfqpoint{3.451907in}{2.732557in}}%
\pgfpathmoveto{\pgfqpoint{3.456448in}{2.717812in}}%
\pgfpathlineto{\pgfqpoint{3.456448in}{2.717812in}}%
\pgfpathlineto{\pgfqpoint{3.456448in}{2.720761in}}%
\pgfpathlineto{\pgfqpoint{3.460989in}{2.720761in}}%
\pgfpathlineto{\pgfqpoint{3.460989in}{2.717812in}}%
\pgfpathmoveto{\pgfqpoint{3.456448in}{2.720761in}}%
\pgfpathlineto{\pgfqpoint{3.456448in}{2.720761in}}%
\pgfpathlineto{\pgfqpoint{3.456448in}{2.723710in}}%
\pgfpathlineto{\pgfqpoint{3.460989in}{2.723710in}}%
\pgfpathlineto{\pgfqpoint{3.460989in}{2.720761in}}%
\pgfpathmoveto{\pgfqpoint{3.438284in}{2.741405in}}%
\pgfpathlineto{\pgfqpoint{3.438284in}{2.741405in}}%
\pgfpathlineto{\pgfqpoint{3.438284in}{2.744354in}}%
\pgfpathlineto{\pgfqpoint{3.442825in}{2.744354in}}%
\pgfpathlineto{\pgfqpoint{3.442825in}{2.741405in}}%
\pgfpathmoveto{\pgfqpoint{3.438284in}{2.744354in}}%
\pgfpathlineto{\pgfqpoint{3.438284in}{2.744354in}}%
\pgfpathlineto{\pgfqpoint{3.438284in}{2.747303in}}%
\pgfpathlineto{\pgfqpoint{3.442825in}{2.747303in}}%
\pgfpathlineto{\pgfqpoint{3.442825in}{2.744354in}}%
\pgfpathmoveto{\pgfqpoint{3.397416in}{2.791542in}}%
\pgfpathlineto{\pgfqpoint{3.397416in}{2.791542in}}%
\pgfpathlineto{\pgfqpoint{3.397416in}{2.794491in}}%
\pgfpathlineto{\pgfqpoint{3.401956in}{2.794491in}}%
\pgfpathlineto{\pgfqpoint{3.401956in}{2.791542in}}%
\pgfpathmoveto{\pgfqpoint{3.392875in}{2.797441in}}%
\pgfpathlineto{\pgfqpoint{3.392875in}{2.797441in}}%
\pgfpathlineto{\pgfqpoint{3.392875in}{2.800390in}}%
\pgfpathlineto{\pgfqpoint{3.397416in}{2.800390in}}%
\pgfpathlineto{\pgfqpoint{3.397416in}{2.797441in}}%
\pgfpathmoveto{\pgfqpoint{3.397416in}{2.794491in}}%
\pgfpathlineto{\pgfqpoint{3.397416in}{2.794491in}}%
\pgfpathlineto{\pgfqpoint{3.397416in}{2.797441in}}%
\pgfpathlineto{\pgfqpoint{3.401956in}{2.797441in}}%
\pgfpathlineto{\pgfqpoint{3.401956in}{2.794491in}}%
\pgfpathmoveto{\pgfqpoint{3.397416in}{2.797441in}}%
\pgfpathlineto{\pgfqpoint{3.397416in}{2.797441in}}%
\pgfpathlineto{\pgfqpoint{3.397416in}{2.800390in}}%
\pgfpathlineto{\pgfqpoint{3.401956in}{2.800390in}}%
\pgfpathlineto{\pgfqpoint{3.401956in}{2.797441in}}%
\pgfpathmoveto{\pgfqpoint{3.388334in}{2.803339in}}%
\pgfpathlineto{\pgfqpoint{3.388334in}{2.803339in}}%
\pgfpathlineto{\pgfqpoint{3.388334in}{2.806289in}}%
\pgfpathlineto{\pgfqpoint{3.392875in}{2.806289in}}%
\pgfpathlineto{\pgfqpoint{3.392875in}{2.803339in}}%
\pgfpathmoveto{\pgfqpoint{3.383793in}{2.809238in}}%
\pgfpathlineto{\pgfqpoint{3.383793in}{2.809238in}}%
\pgfpathlineto{\pgfqpoint{3.383793in}{2.812187in}}%
\pgfpathlineto{\pgfqpoint{3.388334in}{2.812187in}}%
\pgfpathlineto{\pgfqpoint{3.388334in}{2.809238in}}%
\pgfpathmoveto{\pgfqpoint{3.388334in}{2.806289in}}%
\pgfpathlineto{\pgfqpoint{3.388334in}{2.806289in}}%
\pgfpathlineto{\pgfqpoint{3.388334in}{2.809238in}}%
\pgfpathlineto{\pgfqpoint{3.392875in}{2.809238in}}%
\pgfpathlineto{\pgfqpoint{3.392875in}{2.806289in}}%
\pgfpathmoveto{\pgfqpoint{3.388334in}{2.809238in}}%
\pgfpathlineto{\pgfqpoint{3.388334in}{2.809238in}}%
\pgfpathlineto{\pgfqpoint{3.388334in}{2.812187in}}%
\pgfpathlineto{\pgfqpoint{3.392875in}{2.812187in}}%
\pgfpathlineto{\pgfqpoint{3.392875in}{2.809238in}}%
\pgfpathmoveto{\pgfqpoint{3.392875in}{2.800390in}}%
\pgfpathlineto{\pgfqpoint{3.392875in}{2.800390in}}%
\pgfpathlineto{\pgfqpoint{3.392875in}{2.803339in}}%
\pgfpathlineto{\pgfqpoint{3.397416in}{2.803339in}}%
\pgfpathlineto{\pgfqpoint{3.397416in}{2.800390in}}%
\pgfpathmoveto{\pgfqpoint{3.392875in}{2.803339in}}%
\pgfpathlineto{\pgfqpoint{3.392875in}{2.803339in}}%
\pgfpathlineto{\pgfqpoint{3.392875in}{2.806289in}}%
\pgfpathlineto{\pgfqpoint{3.397416in}{2.806289in}}%
\pgfpathlineto{\pgfqpoint{3.397416in}{2.803339in}}%
\pgfpathmoveto{\pgfqpoint{3.415579in}{2.767947in}}%
\pgfpathlineto{\pgfqpoint{3.415579in}{2.767947in}}%
\pgfpathlineto{\pgfqpoint{3.415579in}{2.770896in}}%
\pgfpathlineto{\pgfqpoint{3.420120in}{2.770896in}}%
\pgfpathlineto{\pgfqpoint{3.420120in}{2.767947in}}%
\pgfpathmoveto{\pgfqpoint{3.411038in}{2.773846in}}%
\pgfpathlineto{\pgfqpoint{3.411038in}{2.773846in}}%
\pgfpathlineto{\pgfqpoint{3.411038in}{2.776795in}}%
\pgfpathlineto{\pgfqpoint{3.415579in}{2.776795in}}%
\pgfpathlineto{\pgfqpoint{3.415579in}{2.773846in}}%
\pgfpathmoveto{\pgfqpoint{3.415579in}{2.770896in}}%
\pgfpathlineto{\pgfqpoint{3.415579in}{2.770896in}}%
\pgfpathlineto{\pgfqpoint{3.415579in}{2.773846in}}%
\pgfpathlineto{\pgfqpoint{3.420120in}{2.773846in}}%
\pgfpathlineto{\pgfqpoint{3.420120in}{2.770896in}}%
\pgfpathmoveto{\pgfqpoint{3.415579in}{2.773846in}}%
\pgfpathlineto{\pgfqpoint{3.415579in}{2.773846in}}%
\pgfpathlineto{\pgfqpoint{3.415579in}{2.776795in}}%
\pgfpathlineto{\pgfqpoint{3.420120in}{2.776795in}}%
\pgfpathlineto{\pgfqpoint{3.420120in}{2.773846in}}%
\pgfpathmoveto{\pgfqpoint{3.406497in}{2.779744in}}%
\pgfpathlineto{\pgfqpoint{3.406497in}{2.779744in}}%
\pgfpathlineto{\pgfqpoint{3.406497in}{2.782694in}}%
\pgfpathlineto{\pgfqpoint{3.411038in}{2.782694in}}%
\pgfpathlineto{\pgfqpoint{3.411038in}{2.779744in}}%
\pgfpathmoveto{\pgfqpoint{3.401956in}{2.785643in}}%
\pgfpathlineto{\pgfqpoint{3.401956in}{2.785643in}}%
\pgfpathlineto{\pgfqpoint{3.401956in}{2.788592in}}%
\pgfpathlineto{\pgfqpoint{3.406497in}{2.788592in}}%
\pgfpathlineto{\pgfqpoint{3.406497in}{2.785643in}}%
\pgfpathmoveto{\pgfqpoint{3.406497in}{2.782694in}}%
\pgfpathlineto{\pgfqpoint{3.406497in}{2.782694in}}%
\pgfpathlineto{\pgfqpoint{3.406497in}{2.785643in}}%
\pgfpathlineto{\pgfqpoint{3.411038in}{2.785643in}}%
\pgfpathlineto{\pgfqpoint{3.411038in}{2.782694in}}%
\pgfpathmoveto{\pgfqpoint{3.406497in}{2.785643in}}%
\pgfpathlineto{\pgfqpoint{3.406497in}{2.785643in}}%
\pgfpathlineto{\pgfqpoint{3.406497in}{2.788592in}}%
\pgfpathlineto{\pgfqpoint{3.411038in}{2.788592in}}%
\pgfpathlineto{\pgfqpoint{3.411038in}{2.785643in}}%
\pgfpathmoveto{\pgfqpoint{3.411038in}{2.776795in}}%
\pgfpathlineto{\pgfqpoint{3.411038in}{2.776795in}}%
\pgfpathlineto{\pgfqpoint{3.411038in}{2.779744in}}%
\pgfpathlineto{\pgfqpoint{3.415579in}{2.779744in}}%
\pgfpathlineto{\pgfqpoint{3.415579in}{2.776795in}}%
\pgfpathmoveto{\pgfqpoint{3.411038in}{2.779744in}}%
\pgfpathlineto{\pgfqpoint{3.411038in}{2.779744in}}%
\pgfpathlineto{\pgfqpoint{3.411038in}{2.782694in}}%
\pgfpathlineto{\pgfqpoint{3.415579in}{2.782694in}}%
\pgfpathlineto{\pgfqpoint{3.415579in}{2.779744in}}%
\pgfpathmoveto{\pgfqpoint{3.420120in}{2.764998in}}%
\pgfpathlineto{\pgfqpoint{3.420120in}{2.764998in}}%
\pgfpathlineto{\pgfqpoint{3.420120in}{2.767947in}}%
\pgfpathlineto{\pgfqpoint{3.424661in}{2.767947in}}%
\pgfpathlineto{\pgfqpoint{3.424661in}{2.764998in}}%
\pgfpathmoveto{\pgfqpoint{3.420120in}{2.767947in}}%
\pgfpathlineto{\pgfqpoint{3.420120in}{2.767947in}}%
\pgfpathlineto{\pgfqpoint{3.420120in}{2.770896in}}%
\pgfpathlineto{\pgfqpoint{3.424661in}{2.770896in}}%
\pgfpathlineto{\pgfqpoint{3.424661in}{2.767947in}}%
\pgfpathmoveto{\pgfqpoint{3.401956in}{2.788592in}}%
\pgfpathlineto{\pgfqpoint{3.401956in}{2.788592in}}%
\pgfpathlineto{\pgfqpoint{3.401956in}{2.791542in}}%
\pgfpathlineto{\pgfqpoint{3.406497in}{2.791542in}}%
\pgfpathlineto{\pgfqpoint{3.406497in}{2.788592in}}%
\pgfpathmoveto{\pgfqpoint{3.401956in}{2.791542in}}%
\pgfpathlineto{\pgfqpoint{3.401956in}{2.791542in}}%
\pgfpathlineto{\pgfqpoint{3.401956in}{2.794491in}}%
\pgfpathlineto{\pgfqpoint{3.406497in}{2.794491in}}%
\pgfpathlineto{\pgfqpoint{3.406497in}{2.791542in}}%
\pgfpathmoveto{\pgfqpoint{3.379252in}{2.815137in}}%
\pgfpathlineto{\pgfqpoint{3.379252in}{2.815137in}}%
\pgfpathlineto{\pgfqpoint{3.379252in}{2.818086in}}%
\pgfpathlineto{\pgfqpoint{3.383793in}{2.818086in}}%
\pgfpathlineto{\pgfqpoint{3.383793in}{2.815137in}}%
\pgfpathmoveto{\pgfqpoint{3.374711in}{2.821036in}}%
\pgfpathlineto{\pgfqpoint{3.374711in}{2.821036in}}%
\pgfpathlineto{\pgfqpoint{3.374711in}{2.823985in}}%
\pgfpathlineto{\pgfqpoint{3.379252in}{2.823985in}}%
\pgfpathlineto{\pgfqpoint{3.379252in}{2.821036in}}%
\pgfpathmoveto{\pgfqpoint{3.379252in}{2.818086in}}%
\pgfpathlineto{\pgfqpoint{3.379252in}{2.818086in}}%
\pgfpathlineto{\pgfqpoint{3.379252in}{2.821036in}}%
\pgfpathlineto{\pgfqpoint{3.383793in}{2.821036in}}%
\pgfpathlineto{\pgfqpoint{3.383793in}{2.818086in}}%
\pgfpathmoveto{\pgfqpoint{3.379252in}{2.821036in}}%
\pgfpathlineto{\pgfqpoint{3.379252in}{2.821036in}}%
\pgfpathlineto{\pgfqpoint{3.379252in}{2.823985in}}%
\pgfpathlineto{\pgfqpoint{3.383793in}{2.823985in}}%
\pgfpathlineto{\pgfqpoint{3.383793in}{2.821036in}}%
\pgfpathmoveto{\pgfqpoint{3.370170in}{2.826934in}}%
\pgfpathlineto{\pgfqpoint{3.370170in}{2.826934in}}%
\pgfpathlineto{\pgfqpoint{3.370170in}{2.829884in}}%
\pgfpathlineto{\pgfqpoint{3.374711in}{2.829884in}}%
\pgfpathlineto{\pgfqpoint{3.374711in}{2.826934in}}%
\pgfpathmoveto{\pgfqpoint{3.365629in}{2.832833in}}%
\pgfpathlineto{\pgfqpoint{3.365629in}{2.832833in}}%
\pgfpathlineto{\pgfqpoint{3.365629in}{2.835782in}}%
\pgfpathlineto{\pgfqpoint{3.370170in}{2.835782in}}%
\pgfpathlineto{\pgfqpoint{3.370170in}{2.832833in}}%
\pgfpathmoveto{\pgfqpoint{3.370170in}{2.829884in}}%
\pgfpathlineto{\pgfqpoint{3.370170in}{2.829884in}}%
\pgfpathlineto{\pgfqpoint{3.370170in}{2.832833in}}%
\pgfpathlineto{\pgfqpoint{3.374711in}{2.832833in}}%
\pgfpathlineto{\pgfqpoint{3.374711in}{2.829884in}}%
\pgfpathmoveto{\pgfqpoint{3.370170in}{2.832833in}}%
\pgfpathlineto{\pgfqpoint{3.370170in}{2.832833in}}%
\pgfpathlineto{\pgfqpoint{3.370170in}{2.835782in}}%
\pgfpathlineto{\pgfqpoint{3.374711in}{2.835782in}}%
\pgfpathlineto{\pgfqpoint{3.374711in}{2.832833in}}%
\pgfpathmoveto{\pgfqpoint{3.374711in}{2.823985in}}%
\pgfpathlineto{\pgfqpoint{3.374711in}{2.823985in}}%
\pgfpathlineto{\pgfqpoint{3.374711in}{2.826934in}}%
\pgfpathlineto{\pgfqpoint{3.379252in}{2.826934in}}%
\pgfpathlineto{\pgfqpoint{3.379252in}{2.823985in}}%
\pgfpathmoveto{\pgfqpoint{3.374711in}{2.826934in}}%
\pgfpathlineto{\pgfqpoint{3.374711in}{2.826934in}}%
\pgfpathlineto{\pgfqpoint{3.374711in}{2.829884in}}%
\pgfpathlineto{\pgfqpoint{3.379252in}{2.829884in}}%
\pgfpathlineto{\pgfqpoint{3.379252in}{2.826934in}}%
\pgfpathmoveto{\pgfqpoint{3.383793in}{2.812187in}}%
\pgfpathlineto{\pgfqpoint{3.383793in}{2.812187in}}%
\pgfpathlineto{\pgfqpoint{3.383793in}{2.815137in}}%
\pgfpathlineto{\pgfqpoint{3.388334in}{2.815137in}}%
\pgfpathlineto{\pgfqpoint{3.388334in}{2.812187in}}%
\pgfpathmoveto{\pgfqpoint{3.383793in}{2.815137in}}%
\pgfpathlineto{\pgfqpoint{3.383793in}{2.815137in}}%
\pgfpathlineto{\pgfqpoint{3.383793in}{2.818086in}}%
\pgfpathlineto{\pgfqpoint{3.388334in}{2.818086in}}%
\pgfpathlineto{\pgfqpoint{3.388334in}{2.815137in}}%
\pgfpathmoveto{\pgfqpoint{3.365629in}{2.835782in}}%
\pgfpathlineto{\pgfqpoint{3.365629in}{2.835782in}}%
\pgfpathlineto{\pgfqpoint{3.365629in}{2.838732in}}%
\pgfpathlineto{\pgfqpoint{3.370170in}{2.838732in}}%
\pgfpathlineto{\pgfqpoint{3.370170in}{2.835782in}}%
\pgfpathmoveto{\pgfqpoint{3.365629in}{2.838732in}}%
\pgfpathlineto{\pgfqpoint{3.365629in}{2.838732in}}%
\pgfpathlineto{\pgfqpoint{3.365629in}{2.841681in}}%
\pgfpathlineto{\pgfqpoint{3.370170in}{2.841681in}}%
\pgfpathlineto{\pgfqpoint{3.370170in}{2.838732in}}%
\pgfpathmoveto{\pgfqpoint{3.651709in}{2.461232in}}%
\pgfpathlineto{\pgfqpoint{3.651709in}{2.461232in}}%
\pgfpathlineto{\pgfqpoint{3.651709in}{2.464182in}}%
\pgfpathlineto{\pgfqpoint{3.656250in}{2.464182in}}%
\pgfpathlineto{\pgfqpoint{3.656250in}{2.461232in}}%
\pgfpathmoveto{\pgfqpoint{3.647168in}{2.467131in}}%
\pgfpathlineto{\pgfqpoint{3.647168in}{2.467131in}}%
\pgfpathlineto{\pgfqpoint{3.647168in}{2.470080in}}%
\pgfpathlineto{\pgfqpoint{3.651709in}{2.470080in}}%
\pgfpathlineto{\pgfqpoint{3.651709in}{2.467131in}}%
\pgfpathmoveto{\pgfqpoint{3.651709in}{2.464182in}}%
\pgfpathlineto{\pgfqpoint{3.651709in}{2.464182in}}%
\pgfpathlineto{\pgfqpoint{3.651709in}{2.467131in}}%
\pgfpathlineto{\pgfqpoint{3.656250in}{2.467131in}}%
\pgfpathlineto{\pgfqpoint{3.656250in}{2.464182in}}%
\pgfpathmoveto{\pgfqpoint{3.651709in}{2.467131in}}%
\pgfpathlineto{\pgfqpoint{3.651709in}{2.467131in}}%
\pgfpathlineto{\pgfqpoint{3.651709in}{2.470080in}}%
\pgfpathlineto{\pgfqpoint{3.656250in}{2.470080in}}%
\pgfpathlineto{\pgfqpoint{3.656250in}{2.467131in}}%
\pgfpathmoveto{\pgfqpoint{3.642627in}{2.473030in}}%
\pgfpathlineto{\pgfqpoint{3.642627in}{2.473030in}}%
\pgfpathlineto{\pgfqpoint{3.642627in}{2.475979in}}%
\pgfpathlineto{\pgfqpoint{3.647168in}{2.475979in}}%
\pgfpathlineto{\pgfqpoint{3.647168in}{2.473030in}}%
\pgfpathmoveto{\pgfqpoint{3.638086in}{2.478928in}}%
\pgfpathlineto{\pgfqpoint{3.638086in}{2.478928in}}%
\pgfpathlineto{\pgfqpoint{3.638086in}{2.481878in}}%
\pgfpathlineto{\pgfqpoint{3.642627in}{2.481878in}}%
\pgfpathlineto{\pgfqpoint{3.642627in}{2.478928in}}%
\pgfpathmoveto{\pgfqpoint{3.642627in}{2.475979in}}%
\pgfpathlineto{\pgfqpoint{3.642627in}{2.475979in}}%
\pgfpathlineto{\pgfqpoint{3.642627in}{2.478928in}}%
\pgfpathlineto{\pgfqpoint{3.647168in}{2.478928in}}%
\pgfpathlineto{\pgfqpoint{3.647168in}{2.475979in}}%
\pgfpathmoveto{\pgfqpoint{3.642627in}{2.478928in}}%
\pgfpathlineto{\pgfqpoint{3.642627in}{2.478928in}}%
\pgfpathlineto{\pgfqpoint{3.642627in}{2.481878in}}%
\pgfpathlineto{\pgfqpoint{3.647168in}{2.481878in}}%
\pgfpathlineto{\pgfqpoint{3.647168in}{2.478928in}}%
\pgfpathmoveto{\pgfqpoint{3.647168in}{2.470080in}}%
\pgfpathlineto{\pgfqpoint{3.647168in}{2.470080in}}%
\pgfpathlineto{\pgfqpoint{3.647168in}{2.473030in}}%
\pgfpathlineto{\pgfqpoint{3.651709in}{2.473030in}}%
\pgfpathlineto{\pgfqpoint{3.651709in}{2.470080in}}%
\pgfpathmoveto{\pgfqpoint{3.647168in}{2.473030in}}%
\pgfpathlineto{\pgfqpoint{3.647168in}{2.473030in}}%
\pgfpathlineto{\pgfqpoint{3.647168in}{2.475979in}}%
\pgfpathlineto{\pgfqpoint{3.651709in}{2.475979in}}%
\pgfpathlineto{\pgfqpoint{3.651709in}{2.473030in}}%
\pgfpathmoveto{\pgfqpoint{3.579053in}{2.555607in}}%
\pgfpathlineto{\pgfqpoint{3.579053in}{2.555607in}}%
\pgfpathlineto{\pgfqpoint{3.579053in}{2.558556in}}%
\pgfpathlineto{\pgfqpoint{3.583594in}{2.558556in}}%
\pgfpathlineto{\pgfqpoint{3.583594in}{2.555607in}}%
\pgfpathmoveto{\pgfqpoint{3.574512in}{2.561505in}}%
\pgfpathlineto{\pgfqpoint{3.574512in}{2.561505in}}%
\pgfpathlineto{\pgfqpoint{3.574512in}{2.564454in}}%
\pgfpathlineto{\pgfqpoint{3.579053in}{2.564454in}}%
\pgfpathlineto{\pgfqpoint{3.579053in}{2.561505in}}%
\pgfpathmoveto{\pgfqpoint{3.579053in}{2.558556in}}%
\pgfpathlineto{\pgfqpoint{3.579053in}{2.558556in}}%
\pgfpathlineto{\pgfqpoint{3.579053in}{2.561505in}}%
\pgfpathlineto{\pgfqpoint{3.583594in}{2.561505in}}%
\pgfpathlineto{\pgfqpoint{3.583594in}{2.558556in}}%
\pgfpathmoveto{\pgfqpoint{3.579053in}{2.561505in}}%
\pgfpathlineto{\pgfqpoint{3.579053in}{2.561505in}}%
\pgfpathlineto{\pgfqpoint{3.579053in}{2.564454in}}%
\pgfpathlineto{\pgfqpoint{3.583594in}{2.564454in}}%
\pgfpathlineto{\pgfqpoint{3.583594in}{2.561505in}}%
\pgfpathmoveto{\pgfqpoint{3.569971in}{2.567403in}}%
\pgfpathlineto{\pgfqpoint{3.569971in}{2.567403in}}%
\pgfpathlineto{\pgfqpoint{3.569971in}{2.570352in}}%
\pgfpathlineto{\pgfqpoint{3.574512in}{2.570352in}}%
\pgfpathlineto{\pgfqpoint{3.574512in}{2.567403in}}%
\pgfpathmoveto{\pgfqpoint{3.565431in}{2.573302in}}%
\pgfpathlineto{\pgfqpoint{3.565431in}{2.573302in}}%
\pgfpathlineto{\pgfqpoint{3.565431in}{2.576251in}}%
\pgfpathlineto{\pgfqpoint{3.569971in}{2.576251in}}%
\pgfpathlineto{\pgfqpoint{3.569971in}{2.573302in}}%
\pgfpathmoveto{\pgfqpoint{3.569971in}{2.570352in}}%
\pgfpathlineto{\pgfqpoint{3.569971in}{2.570352in}}%
\pgfpathlineto{\pgfqpoint{3.569971in}{2.573302in}}%
\pgfpathlineto{\pgfqpoint{3.574512in}{2.573302in}}%
\pgfpathlineto{\pgfqpoint{3.574512in}{2.570352in}}%
\pgfpathmoveto{\pgfqpoint{3.569971in}{2.573302in}}%
\pgfpathlineto{\pgfqpoint{3.569971in}{2.573302in}}%
\pgfpathlineto{\pgfqpoint{3.569971in}{2.576251in}}%
\pgfpathlineto{\pgfqpoint{3.574512in}{2.576251in}}%
\pgfpathlineto{\pgfqpoint{3.574512in}{2.573302in}}%
\pgfpathmoveto{\pgfqpoint{3.574512in}{2.564454in}}%
\pgfpathlineto{\pgfqpoint{3.574512in}{2.564454in}}%
\pgfpathlineto{\pgfqpoint{3.574512in}{2.567403in}}%
\pgfpathlineto{\pgfqpoint{3.579053in}{2.567403in}}%
\pgfpathlineto{\pgfqpoint{3.579053in}{2.564454in}}%
\pgfpathmoveto{\pgfqpoint{3.574512in}{2.567403in}}%
\pgfpathlineto{\pgfqpoint{3.574512in}{2.567403in}}%
\pgfpathlineto{\pgfqpoint{3.574512in}{2.570352in}}%
\pgfpathlineto{\pgfqpoint{3.579053in}{2.570352in}}%
\pgfpathlineto{\pgfqpoint{3.579053in}{2.567403in}}%
\pgfpathmoveto{\pgfqpoint{3.615381in}{2.508420in}}%
\pgfpathlineto{\pgfqpoint{3.615381in}{2.508420in}}%
\pgfpathlineto{\pgfqpoint{3.615381in}{2.511369in}}%
\pgfpathlineto{\pgfqpoint{3.619922in}{2.511369in}}%
\pgfpathlineto{\pgfqpoint{3.619922in}{2.508420in}}%
\pgfpathmoveto{\pgfqpoint{3.610840in}{2.514318in}}%
\pgfpathlineto{\pgfqpoint{3.610840in}{2.514318in}}%
\pgfpathlineto{\pgfqpoint{3.610840in}{2.517268in}}%
\pgfpathlineto{\pgfqpoint{3.615381in}{2.517268in}}%
\pgfpathlineto{\pgfqpoint{3.615381in}{2.514318in}}%
\pgfpathmoveto{\pgfqpoint{3.615381in}{2.511369in}}%
\pgfpathlineto{\pgfqpoint{3.615381in}{2.511369in}}%
\pgfpathlineto{\pgfqpoint{3.615381in}{2.514318in}}%
\pgfpathlineto{\pgfqpoint{3.619922in}{2.514318in}}%
\pgfpathlineto{\pgfqpoint{3.619922in}{2.511369in}}%
\pgfpathmoveto{\pgfqpoint{3.615381in}{2.514318in}}%
\pgfpathlineto{\pgfqpoint{3.615381in}{2.514318in}}%
\pgfpathlineto{\pgfqpoint{3.615381in}{2.517268in}}%
\pgfpathlineto{\pgfqpoint{3.619922in}{2.517268in}}%
\pgfpathlineto{\pgfqpoint{3.619922in}{2.514318in}}%
\pgfpathmoveto{\pgfqpoint{3.606299in}{2.520217in}}%
\pgfpathlineto{\pgfqpoint{3.606299in}{2.520217in}}%
\pgfpathlineto{\pgfqpoint{3.606299in}{2.523166in}}%
\pgfpathlineto{\pgfqpoint{3.610840in}{2.523166in}}%
\pgfpathlineto{\pgfqpoint{3.610840in}{2.520217in}}%
\pgfpathmoveto{\pgfqpoint{3.601758in}{2.526115in}}%
\pgfpathlineto{\pgfqpoint{3.601758in}{2.526115in}}%
\pgfpathlineto{\pgfqpoint{3.601758in}{2.529064in}}%
\pgfpathlineto{\pgfqpoint{3.606299in}{2.529064in}}%
\pgfpathlineto{\pgfqpoint{3.606299in}{2.526115in}}%
\pgfpathmoveto{\pgfqpoint{3.606299in}{2.523166in}}%
\pgfpathlineto{\pgfqpoint{3.606299in}{2.523166in}}%
\pgfpathlineto{\pgfqpoint{3.606299in}{2.526115in}}%
\pgfpathlineto{\pgfqpoint{3.610840in}{2.526115in}}%
\pgfpathlineto{\pgfqpoint{3.610840in}{2.523166in}}%
\pgfpathmoveto{\pgfqpoint{3.606299in}{2.526115in}}%
\pgfpathlineto{\pgfqpoint{3.606299in}{2.526115in}}%
\pgfpathlineto{\pgfqpoint{3.606299in}{2.529064in}}%
\pgfpathlineto{\pgfqpoint{3.610840in}{2.529064in}}%
\pgfpathlineto{\pgfqpoint{3.610840in}{2.526115in}}%
\pgfpathmoveto{\pgfqpoint{3.610840in}{2.517268in}}%
\pgfpathlineto{\pgfqpoint{3.610840in}{2.517268in}}%
\pgfpathlineto{\pgfqpoint{3.610840in}{2.520217in}}%
\pgfpathlineto{\pgfqpoint{3.615381in}{2.520217in}}%
\pgfpathlineto{\pgfqpoint{3.615381in}{2.517268in}}%
\pgfpathmoveto{\pgfqpoint{3.610840in}{2.520217in}}%
\pgfpathlineto{\pgfqpoint{3.610840in}{2.520217in}}%
\pgfpathlineto{\pgfqpoint{3.610840in}{2.523166in}}%
\pgfpathlineto{\pgfqpoint{3.615381in}{2.523166in}}%
\pgfpathlineto{\pgfqpoint{3.615381in}{2.520217in}}%
\pgfpathmoveto{\pgfqpoint{3.633545in}{2.484827in}}%
\pgfpathlineto{\pgfqpoint{3.633545in}{2.484827in}}%
\pgfpathlineto{\pgfqpoint{3.633545in}{2.487776in}}%
\pgfpathlineto{\pgfqpoint{3.638086in}{2.487776in}}%
\pgfpathlineto{\pgfqpoint{3.638086in}{2.484827in}}%
\pgfpathmoveto{\pgfqpoint{3.629004in}{2.490725in}}%
\pgfpathlineto{\pgfqpoint{3.629004in}{2.490725in}}%
\pgfpathlineto{\pgfqpoint{3.629004in}{2.493674in}}%
\pgfpathlineto{\pgfqpoint{3.633545in}{2.493674in}}%
\pgfpathlineto{\pgfqpoint{3.633545in}{2.490725in}}%
\pgfpathmoveto{\pgfqpoint{3.633545in}{2.487776in}}%
\pgfpathlineto{\pgfqpoint{3.633545in}{2.487776in}}%
\pgfpathlineto{\pgfqpoint{3.633545in}{2.490725in}}%
\pgfpathlineto{\pgfqpoint{3.638086in}{2.490725in}}%
\pgfpathlineto{\pgfqpoint{3.638086in}{2.487776in}}%
\pgfpathmoveto{\pgfqpoint{3.633545in}{2.490725in}}%
\pgfpathlineto{\pgfqpoint{3.633545in}{2.490725in}}%
\pgfpathlineto{\pgfqpoint{3.633545in}{2.493674in}}%
\pgfpathlineto{\pgfqpoint{3.638086in}{2.493674in}}%
\pgfpathlineto{\pgfqpoint{3.638086in}{2.490725in}}%
\pgfpathmoveto{\pgfqpoint{3.624463in}{2.496624in}}%
\pgfpathlineto{\pgfqpoint{3.624463in}{2.496624in}}%
\pgfpathlineto{\pgfqpoint{3.624463in}{2.499573in}}%
\pgfpathlineto{\pgfqpoint{3.629004in}{2.499573in}}%
\pgfpathlineto{\pgfqpoint{3.629004in}{2.496624in}}%
\pgfpathmoveto{\pgfqpoint{3.619922in}{2.502522in}}%
\pgfpathlineto{\pgfqpoint{3.619922in}{2.502522in}}%
\pgfpathlineto{\pgfqpoint{3.619922in}{2.505471in}}%
\pgfpathlineto{\pgfqpoint{3.624463in}{2.505471in}}%
\pgfpathlineto{\pgfqpoint{3.624463in}{2.502522in}}%
\pgfpathmoveto{\pgfqpoint{3.624463in}{2.499573in}}%
\pgfpathlineto{\pgfqpoint{3.624463in}{2.499573in}}%
\pgfpathlineto{\pgfqpoint{3.624463in}{2.502522in}}%
\pgfpathlineto{\pgfqpoint{3.629004in}{2.502522in}}%
\pgfpathlineto{\pgfqpoint{3.629004in}{2.499573in}}%
\pgfpathmoveto{\pgfqpoint{3.624463in}{2.502522in}}%
\pgfpathlineto{\pgfqpoint{3.624463in}{2.502522in}}%
\pgfpathlineto{\pgfqpoint{3.624463in}{2.505471in}}%
\pgfpathlineto{\pgfqpoint{3.629004in}{2.505471in}}%
\pgfpathlineto{\pgfqpoint{3.629004in}{2.502522in}}%
\pgfpathmoveto{\pgfqpoint{3.629004in}{2.493674in}}%
\pgfpathlineto{\pgfqpoint{3.629004in}{2.493674in}}%
\pgfpathlineto{\pgfqpoint{3.629004in}{2.496624in}}%
\pgfpathlineto{\pgfqpoint{3.633545in}{2.496624in}}%
\pgfpathlineto{\pgfqpoint{3.633545in}{2.493674in}}%
\pgfpathmoveto{\pgfqpoint{3.629004in}{2.496624in}}%
\pgfpathlineto{\pgfqpoint{3.629004in}{2.496624in}}%
\pgfpathlineto{\pgfqpoint{3.629004in}{2.499573in}}%
\pgfpathlineto{\pgfqpoint{3.633545in}{2.499573in}}%
\pgfpathlineto{\pgfqpoint{3.633545in}{2.496624in}}%
\pgfpathmoveto{\pgfqpoint{3.638086in}{2.481878in}}%
\pgfpathlineto{\pgfqpoint{3.638086in}{2.481878in}}%
\pgfpathlineto{\pgfqpoint{3.638086in}{2.484827in}}%
\pgfpathlineto{\pgfqpoint{3.642627in}{2.484827in}}%
\pgfpathlineto{\pgfqpoint{3.642627in}{2.481878in}}%
\pgfpathmoveto{\pgfqpoint{3.638086in}{2.484827in}}%
\pgfpathlineto{\pgfqpoint{3.638086in}{2.484827in}}%
\pgfpathlineto{\pgfqpoint{3.638086in}{2.487776in}}%
\pgfpathlineto{\pgfqpoint{3.642627in}{2.487776in}}%
\pgfpathlineto{\pgfqpoint{3.642627in}{2.484827in}}%
\pgfpathmoveto{\pgfqpoint{3.619922in}{2.505471in}}%
\pgfpathlineto{\pgfqpoint{3.619922in}{2.505471in}}%
\pgfpathlineto{\pgfqpoint{3.619922in}{2.508420in}}%
\pgfpathlineto{\pgfqpoint{3.624463in}{2.508420in}}%
\pgfpathlineto{\pgfqpoint{3.624463in}{2.505471in}}%
\pgfpathmoveto{\pgfqpoint{3.619922in}{2.508420in}}%
\pgfpathlineto{\pgfqpoint{3.619922in}{2.508420in}}%
\pgfpathlineto{\pgfqpoint{3.619922in}{2.511369in}}%
\pgfpathlineto{\pgfqpoint{3.624463in}{2.511369in}}%
\pgfpathlineto{\pgfqpoint{3.624463in}{2.508420in}}%
\pgfpathmoveto{\pgfqpoint{3.597217in}{2.532013in}}%
\pgfpathlineto{\pgfqpoint{3.597217in}{2.532013in}}%
\pgfpathlineto{\pgfqpoint{3.597217in}{2.534963in}}%
\pgfpathlineto{\pgfqpoint{3.601758in}{2.534963in}}%
\pgfpathlineto{\pgfqpoint{3.601758in}{2.532013in}}%
\pgfpathmoveto{\pgfqpoint{3.592676in}{2.537912in}}%
\pgfpathlineto{\pgfqpoint{3.592676in}{2.537912in}}%
\pgfpathlineto{\pgfqpoint{3.592676in}{2.540861in}}%
\pgfpathlineto{\pgfqpoint{3.597217in}{2.540861in}}%
\pgfpathlineto{\pgfqpoint{3.597217in}{2.537912in}}%
\pgfpathmoveto{\pgfqpoint{3.597217in}{2.534963in}}%
\pgfpathlineto{\pgfqpoint{3.597217in}{2.534963in}}%
\pgfpathlineto{\pgfqpoint{3.597217in}{2.537912in}}%
\pgfpathlineto{\pgfqpoint{3.601758in}{2.537912in}}%
\pgfpathlineto{\pgfqpoint{3.601758in}{2.534963in}}%
\pgfpathmoveto{\pgfqpoint{3.597217in}{2.537912in}}%
\pgfpathlineto{\pgfqpoint{3.597217in}{2.537912in}}%
\pgfpathlineto{\pgfqpoint{3.597217in}{2.540861in}}%
\pgfpathlineto{\pgfqpoint{3.601758in}{2.540861in}}%
\pgfpathlineto{\pgfqpoint{3.601758in}{2.537912in}}%
\pgfpathmoveto{\pgfqpoint{3.588135in}{2.543810in}}%
\pgfpathlineto{\pgfqpoint{3.588135in}{2.543810in}}%
\pgfpathlineto{\pgfqpoint{3.588135in}{2.546759in}}%
\pgfpathlineto{\pgfqpoint{3.592676in}{2.546759in}}%
\pgfpathlineto{\pgfqpoint{3.592676in}{2.543810in}}%
\pgfpathmoveto{\pgfqpoint{3.583594in}{2.549708in}}%
\pgfpathlineto{\pgfqpoint{3.583594in}{2.549708in}}%
\pgfpathlineto{\pgfqpoint{3.583594in}{2.552657in}}%
\pgfpathlineto{\pgfqpoint{3.588135in}{2.552657in}}%
\pgfpathlineto{\pgfqpoint{3.588135in}{2.549708in}}%
\pgfpathmoveto{\pgfqpoint{3.588135in}{2.546759in}}%
\pgfpathlineto{\pgfqpoint{3.588135in}{2.546759in}}%
\pgfpathlineto{\pgfqpoint{3.588135in}{2.549708in}}%
\pgfpathlineto{\pgfqpoint{3.592676in}{2.549708in}}%
\pgfpathlineto{\pgfqpoint{3.592676in}{2.546759in}}%
\pgfpathmoveto{\pgfqpoint{3.588135in}{2.549708in}}%
\pgfpathlineto{\pgfqpoint{3.588135in}{2.549708in}}%
\pgfpathlineto{\pgfqpoint{3.588135in}{2.552657in}}%
\pgfpathlineto{\pgfqpoint{3.592676in}{2.552657in}}%
\pgfpathlineto{\pgfqpoint{3.592676in}{2.549708in}}%
\pgfpathmoveto{\pgfqpoint{3.592676in}{2.540861in}}%
\pgfpathlineto{\pgfqpoint{3.592676in}{2.540861in}}%
\pgfpathlineto{\pgfqpoint{3.592676in}{2.543810in}}%
\pgfpathlineto{\pgfqpoint{3.597217in}{2.543810in}}%
\pgfpathlineto{\pgfqpoint{3.597217in}{2.540861in}}%
\pgfpathmoveto{\pgfqpoint{3.592676in}{2.543810in}}%
\pgfpathlineto{\pgfqpoint{3.592676in}{2.543810in}}%
\pgfpathlineto{\pgfqpoint{3.592676in}{2.546759in}}%
\pgfpathlineto{\pgfqpoint{3.597217in}{2.546759in}}%
\pgfpathlineto{\pgfqpoint{3.597217in}{2.543810in}}%
\pgfpathmoveto{\pgfqpoint{3.601758in}{2.529064in}}%
\pgfpathlineto{\pgfqpoint{3.601758in}{2.529064in}}%
\pgfpathlineto{\pgfqpoint{3.601758in}{2.532013in}}%
\pgfpathlineto{\pgfqpoint{3.606299in}{2.532013in}}%
\pgfpathlineto{\pgfqpoint{3.606299in}{2.529064in}}%
\pgfpathmoveto{\pgfqpoint{3.601758in}{2.532013in}}%
\pgfpathlineto{\pgfqpoint{3.601758in}{2.532013in}}%
\pgfpathlineto{\pgfqpoint{3.601758in}{2.534963in}}%
\pgfpathlineto{\pgfqpoint{3.606299in}{2.534963in}}%
\pgfpathlineto{\pgfqpoint{3.606299in}{2.532013in}}%
\pgfpathmoveto{\pgfqpoint{3.583594in}{2.552657in}}%
\pgfpathlineto{\pgfqpoint{3.583594in}{2.552657in}}%
\pgfpathlineto{\pgfqpoint{3.583594in}{2.555607in}}%
\pgfpathlineto{\pgfqpoint{3.588135in}{2.555607in}}%
\pgfpathlineto{\pgfqpoint{3.588135in}{2.552657in}}%
\pgfpathmoveto{\pgfqpoint{3.583594in}{2.555607in}}%
\pgfpathlineto{\pgfqpoint{3.583594in}{2.555607in}}%
\pgfpathlineto{\pgfqpoint{3.583594in}{2.558556in}}%
\pgfpathlineto{\pgfqpoint{3.588135in}{2.558556in}}%
\pgfpathlineto{\pgfqpoint{3.588135in}{2.555607in}}%
\pgfpathmoveto{\pgfqpoint{3.542726in}{2.602794in}}%
\pgfpathlineto{\pgfqpoint{3.542726in}{2.602794in}}%
\pgfpathlineto{\pgfqpoint{3.542726in}{2.605743in}}%
\pgfpathlineto{\pgfqpoint{3.547267in}{2.605743in}}%
\pgfpathlineto{\pgfqpoint{3.547267in}{2.602794in}}%
\pgfpathmoveto{\pgfqpoint{3.538185in}{2.608692in}}%
\pgfpathlineto{\pgfqpoint{3.538185in}{2.608692in}}%
\pgfpathlineto{\pgfqpoint{3.538185in}{2.611641in}}%
\pgfpathlineto{\pgfqpoint{3.542726in}{2.611641in}}%
\pgfpathlineto{\pgfqpoint{3.542726in}{2.608692in}}%
\pgfpathmoveto{\pgfqpoint{3.542726in}{2.605743in}}%
\pgfpathlineto{\pgfqpoint{3.542726in}{2.605743in}}%
\pgfpathlineto{\pgfqpoint{3.542726in}{2.608692in}}%
\pgfpathlineto{\pgfqpoint{3.547267in}{2.608692in}}%
\pgfpathlineto{\pgfqpoint{3.547267in}{2.605743in}}%
\pgfpathmoveto{\pgfqpoint{3.542726in}{2.608692in}}%
\pgfpathlineto{\pgfqpoint{3.542726in}{2.608692in}}%
\pgfpathlineto{\pgfqpoint{3.542726in}{2.611641in}}%
\pgfpathlineto{\pgfqpoint{3.547267in}{2.611641in}}%
\pgfpathlineto{\pgfqpoint{3.547267in}{2.608692in}}%
\pgfpathmoveto{\pgfqpoint{3.533644in}{2.614590in}}%
\pgfpathlineto{\pgfqpoint{3.533644in}{2.614590in}}%
\pgfpathlineto{\pgfqpoint{3.533644in}{2.617540in}}%
\pgfpathlineto{\pgfqpoint{3.538185in}{2.617540in}}%
\pgfpathlineto{\pgfqpoint{3.538185in}{2.614590in}}%
\pgfpathmoveto{\pgfqpoint{3.529103in}{2.620489in}}%
\pgfpathlineto{\pgfqpoint{3.529103in}{2.620489in}}%
\pgfpathlineto{\pgfqpoint{3.529103in}{2.623438in}}%
\pgfpathlineto{\pgfqpoint{3.533644in}{2.623438in}}%
\pgfpathlineto{\pgfqpoint{3.533644in}{2.620489in}}%
\pgfpathmoveto{\pgfqpoint{3.533644in}{2.617540in}}%
\pgfpathlineto{\pgfqpoint{3.533644in}{2.617540in}}%
\pgfpathlineto{\pgfqpoint{3.533644in}{2.620489in}}%
\pgfpathlineto{\pgfqpoint{3.538185in}{2.620489in}}%
\pgfpathlineto{\pgfqpoint{3.538185in}{2.617540in}}%
\pgfpathmoveto{\pgfqpoint{3.533644in}{2.620489in}}%
\pgfpathlineto{\pgfqpoint{3.533644in}{2.620489in}}%
\pgfpathlineto{\pgfqpoint{3.533644in}{2.623438in}}%
\pgfpathlineto{\pgfqpoint{3.538185in}{2.623438in}}%
\pgfpathlineto{\pgfqpoint{3.538185in}{2.620489in}}%
\pgfpathmoveto{\pgfqpoint{3.538185in}{2.611641in}}%
\pgfpathlineto{\pgfqpoint{3.538185in}{2.611641in}}%
\pgfpathlineto{\pgfqpoint{3.538185in}{2.614590in}}%
\pgfpathlineto{\pgfqpoint{3.542726in}{2.614590in}}%
\pgfpathlineto{\pgfqpoint{3.542726in}{2.611641in}}%
\pgfpathmoveto{\pgfqpoint{3.538185in}{2.614590in}}%
\pgfpathlineto{\pgfqpoint{3.538185in}{2.614590in}}%
\pgfpathlineto{\pgfqpoint{3.538185in}{2.617540in}}%
\pgfpathlineto{\pgfqpoint{3.542726in}{2.617540in}}%
\pgfpathlineto{\pgfqpoint{3.542726in}{2.614590in}}%
\pgfpathmoveto{\pgfqpoint{3.560890in}{2.579200in}}%
\pgfpathlineto{\pgfqpoint{3.560890in}{2.579200in}}%
\pgfpathlineto{\pgfqpoint{3.560890in}{2.582149in}}%
\pgfpathlineto{\pgfqpoint{3.565431in}{2.582149in}}%
\pgfpathlineto{\pgfqpoint{3.565431in}{2.579200in}}%
\pgfpathmoveto{\pgfqpoint{3.556349in}{2.585098in}}%
\pgfpathlineto{\pgfqpoint{3.556349in}{2.585098in}}%
\pgfpathlineto{\pgfqpoint{3.556349in}{2.588048in}}%
\pgfpathlineto{\pgfqpoint{3.560890in}{2.588048in}}%
\pgfpathlineto{\pgfqpoint{3.560890in}{2.585098in}}%
\pgfpathmoveto{\pgfqpoint{3.560890in}{2.582149in}}%
\pgfpathlineto{\pgfqpoint{3.560890in}{2.582149in}}%
\pgfpathlineto{\pgfqpoint{3.560890in}{2.585098in}}%
\pgfpathlineto{\pgfqpoint{3.565431in}{2.585098in}}%
\pgfpathlineto{\pgfqpoint{3.565431in}{2.582149in}}%
\pgfpathmoveto{\pgfqpoint{3.560890in}{2.585098in}}%
\pgfpathlineto{\pgfqpoint{3.560890in}{2.585098in}}%
\pgfpathlineto{\pgfqpoint{3.560890in}{2.588048in}}%
\pgfpathlineto{\pgfqpoint{3.565431in}{2.588048in}}%
\pgfpathlineto{\pgfqpoint{3.565431in}{2.585098in}}%
\pgfpathmoveto{\pgfqpoint{3.551808in}{2.590997in}}%
\pgfpathlineto{\pgfqpoint{3.551808in}{2.590997in}}%
\pgfpathlineto{\pgfqpoint{3.551808in}{2.593946in}}%
\pgfpathlineto{\pgfqpoint{3.556349in}{2.593946in}}%
\pgfpathlineto{\pgfqpoint{3.556349in}{2.590997in}}%
\pgfpathmoveto{\pgfqpoint{3.547267in}{2.596895in}}%
\pgfpathlineto{\pgfqpoint{3.547267in}{2.596895in}}%
\pgfpathlineto{\pgfqpoint{3.547267in}{2.599844in}}%
\pgfpathlineto{\pgfqpoint{3.551808in}{2.599844in}}%
\pgfpathlineto{\pgfqpoint{3.551808in}{2.596895in}}%
\pgfpathmoveto{\pgfqpoint{3.551808in}{2.593946in}}%
\pgfpathlineto{\pgfqpoint{3.551808in}{2.593946in}}%
\pgfpathlineto{\pgfqpoint{3.551808in}{2.596895in}}%
\pgfpathlineto{\pgfqpoint{3.556349in}{2.596895in}}%
\pgfpathlineto{\pgfqpoint{3.556349in}{2.593946in}}%
\pgfpathmoveto{\pgfqpoint{3.551808in}{2.596895in}}%
\pgfpathlineto{\pgfqpoint{3.551808in}{2.596895in}}%
\pgfpathlineto{\pgfqpoint{3.551808in}{2.599844in}}%
\pgfpathlineto{\pgfqpoint{3.556349in}{2.599844in}}%
\pgfpathlineto{\pgfqpoint{3.556349in}{2.596895in}}%
\pgfpathmoveto{\pgfqpoint{3.556349in}{2.588048in}}%
\pgfpathlineto{\pgfqpoint{3.556349in}{2.588048in}}%
\pgfpathlineto{\pgfqpoint{3.556349in}{2.590997in}}%
\pgfpathlineto{\pgfqpoint{3.560890in}{2.590997in}}%
\pgfpathlineto{\pgfqpoint{3.560890in}{2.588048in}}%
\pgfpathmoveto{\pgfqpoint{3.556349in}{2.590997in}}%
\pgfpathlineto{\pgfqpoint{3.556349in}{2.590997in}}%
\pgfpathlineto{\pgfqpoint{3.556349in}{2.593946in}}%
\pgfpathlineto{\pgfqpoint{3.560890in}{2.593946in}}%
\pgfpathlineto{\pgfqpoint{3.560890in}{2.590997in}}%
\pgfpathmoveto{\pgfqpoint{3.565431in}{2.576251in}}%
\pgfpathlineto{\pgfqpoint{3.565431in}{2.576251in}}%
\pgfpathlineto{\pgfqpoint{3.565431in}{2.579200in}}%
\pgfpathlineto{\pgfqpoint{3.569971in}{2.579200in}}%
\pgfpathlineto{\pgfqpoint{3.569971in}{2.576251in}}%
\pgfpathmoveto{\pgfqpoint{3.565431in}{2.579200in}}%
\pgfpathlineto{\pgfqpoint{3.565431in}{2.579200in}}%
\pgfpathlineto{\pgfqpoint{3.565431in}{2.582149in}}%
\pgfpathlineto{\pgfqpoint{3.569971in}{2.582149in}}%
\pgfpathlineto{\pgfqpoint{3.569971in}{2.579200in}}%
\pgfpathmoveto{\pgfqpoint{3.547267in}{2.599844in}}%
\pgfpathlineto{\pgfqpoint{3.547267in}{2.599844in}}%
\pgfpathlineto{\pgfqpoint{3.547267in}{2.602794in}}%
\pgfpathlineto{\pgfqpoint{3.551808in}{2.602794in}}%
\pgfpathlineto{\pgfqpoint{3.551808in}{2.599844in}}%
\pgfpathmoveto{\pgfqpoint{3.547267in}{2.602794in}}%
\pgfpathlineto{\pgfqpoint{3.547267in}{2.602794in}}%
\pgfpathlineto{\pgfqpoint{3.547267in}{2.605743in}}%
\pgfpathlineto{\pgfqpoint{3.551808in}{2.605743in}}%
\pgfpathlineto{\pgfqpoint{3.551808in}{2.602794in}}%
\pgfpathmoveto{\pgfqpoint{3.524562in}{2.626387in}}%
\pgfpathlineto{\pgfqpoint{3.524562in}{2.626387in}}%
\pgfpathlineto{\pgfqpoint{3.524562in}{2.629337in}}%
\pgfpathlineto{\pgfqpoint{3.529103in}{2.629337in}}%
\pgfpathlineto{\pgfqpoint{3.529103in}{2.626387in}}%
\pgfpathmoveto{\pgfqpoint{3.520021in}{2.632286in}}%
\pgfpathlineto{\pgfqpoint{3.520021in}{2.632286in}}%
\pgfpathlineto{\pgfqpoint{3.520021in}{2.635235in}}%
\pgfpathlineto{\pgfqpoint{3.524562in}{2.635235in}}%
\pgfpathlineto{\pgfqpoint{3.524562in}{2.632286in}}%
\pgfpathmoveto{\pgfqpoint{3.524562in}{2.629337in}}%
\pgfpathlineto{\pgfqpoint{3.524562in}{2.629337in}}%
\pgfpathlineto{\pgfqpoint{3.524562in}{2.632286in}}%
\pgfpathlineto{\pgfqpoint{3.529103in}{2.632286in}}%
\pgfpathlineto{\pgfqpoint{3.529103in}{2.629337in}}%
\pgfpathmoveto{\pgfqpoint{3.524562in}{2.632286in}}%
\pgfpathlineto{\pgfqpoint{3.524562in}{2.632286in}}%
\pgfpathlineto{\pgfqpoint{3.524562in}{2.635235in}}%
\pgfpathlineto{\pgfqpoint{3.529103in}{2.635235in}}%
\pgfpathlineto{\pgfqpoint{3.529103in}{2.632286in}}%
\pgfpathmoveto{\pgfqpoint{3.515480in}{2.638184in}}%
\pgfpathlineto{\pgfqpoint{3.515480in}{2.638184in}}%
\pgfpathlineto{\pgfqpoint{3.515480in}{2.641133in}}%
\pgfpathlineto{\pgfqpoint{3.520021in}{2.641133in}}%
\pgfpathlineto{\pgfqpoint{3.520021in}{2.638184in}}%
\pgfpathmoveto{\pgfqpoint{3.510939in}{2.644083in}}%
\pgfpathlineto{\pgfqpoint{3.510939in}{2.644083in}}%
\pgfpathlineto{\pgfqpoint{3.510939in}{2.647032in}}%
\pgfpathlineto{\pgfqpoint{3.515480in}{2.647032in}}%
\pgfpathlineto{\pgfqpoint{3.515480in}{2.644083in}}%
\pgfpathmoveto{\pgfqpoint{3.515480in}{2.641133in}}%
\pgfpathlineto{\pgfqpoint{3.515480in}{2.641133in}}%
\pgfpathlineto{\pgfqpoint{3.515480in}{2.644083in}}%
\pgfpathlineto{\pgfqpoint{3.520021in}{2.644083in}}%
\pgfpathlineto{\pgfqpoint{3.520021in}{2.641133in}}%
\pgfpathmoveto{\pgfqpoint{3.515480in}{2.644083in}}%
\pgfpathlineto{\pgfqpoint{3.515480in}{2.644083in}}%
\pgfpathlineto{\pgfqpoint{3.515480in}{2.647032in}}%
\pgfpathlineto{\pgfqpoint{3.520021in}{2.647032in}}%
\pgfpathlineto{\pgfqpoint{3.520021in}{2.644083in}}%
\pgfpathmoveto{\pgfqpoint{3.520021in}{2.635235in}}%
\pgfpathlineto{\pgfqpoint{3.520021in}{2.635235in}}%
\pgfpathlineto{\pgfqpoint{3.520021in}{2.638184in}}%
\pgfpathlineto{\pgfqpoint{3.524562in}{2.638184in}}%
\pgfpathlineto{\pgfqpoint{3.524562in}{2.635235in}}%
\pgfpathmoveto{\pgfqpoint{3.520021in}{2.638184in}}%
\pgfpathlineto{\pgfqpoint{3.520021in}{2.638184in}}%
\pgfpathlineto{\pgfqpoint{3.520021in}{2.641133in}}%
\pgfpathlineto{\pgfqpoint{3.524562in}{2.641133in}}%
\pgfpathlineto{\pgfqpoint{3.524562in}{2.638184in}}%
\pgfpathmoveto{\pgfqpoint{3.529103in}{2.623438in}}%
\pgfpathlineto{\pgfqpoint{3.529103in}{2.623438in}}%
\pgfpathlineto{\pgfqpoint{3.529103in}{2.626387in}}%
\pgfpathlineto{\pgfqpoint{3.533644in}{2.626387in}}%
\pgfpathlineto{\pgfqpoint{3.533644in}{2.623438in}}%
\pgfpathmoveto{\pgfqpoint{3.529103in}{2.626387in}}%
\pgfpathlineto{\pgfqpoint{3.529103in}{2.626387in}}%
\pgfpathlineto{\pgfqpoint{3.529103in}{2.629337in}}%
\pgfpathlineto{\pgfqpoint{3.533644in}{2.629337in}}%
\pgfpathlineto{\pgfqpoint{3.533644in}{2.626387in}}%
\pgfpathmoveto{\pgfqpoint{3.510939in}{2.647032in}}%
\pgfpathlineto{\pgfqpoint{3.510939in}{2.647032in}}%
\pgfpathlineto{\pgfqpoint{3.510939in}{2.649981in}}%
\pgfpathlineto{\pgfqpoint{3.515480in}{2.649981in}}%
\pgfpathlineto{\pgfqpoint{3.515480in}{2.647032in}}%
\pgfpathmoveto{\pgfqpoint{3.510939in}{2.649981in}}%
\pgfpathlineto{\pgfqpoint{3.510939in}{2.649981in}}%
\pgfpathlineto{\pgfqpoint{3.510939in}{2.652930in}}%
\pgfpathlineto{\pgfqpoint{3.515480in}{2.652930in}}%
\pgfpathlineto{\pgfqpoint{3.515480in}{2.649981in}}%
\pgfpathmoveto{\pgfqpoint{3.797025in}{2.272480in}}%
\pgfpathlineto{\pgfqpoint{3.797025in}{2.272480in}}%
\pgfpathlineto{\pgfqpoint{3.797025in}{2.275429in}}%
\pgfpathlineto{\pgfqpoint{3.801566in}{2.275429in}}%
\pgfpathlineto{\pgfqpoint{3.801566in}{2.272480in}}%
\pgfpathmoveto{\pgfqpoint{3.792484in}{2.278378in}}%
\pgfpathlineto{\pgfqpoint{3.792484in}{2.278378in}}%
\pgfpathlineto{\pgfqpoint{3.792484in}{2.281327in}}%
\pgfpathlineto{\pgfqpoint{3.797025in}{2.281327in}}%
\pgfpathlineto{\pgfqpoint{3.797025in}{2.278378in}}%
\pgfpathmoveto{\pgfqpoint{3.797025in}{2.275429in}}%
\pgfpathlineto{\pgfqpoint{3.797025in}{2.275429in}}%
\pgfpathlineto{\pgfqpoint{3.797025in}{2.278378in}}%
\pgfpathlineto{\pgfqpoint{3.801566in}{2.278378in}}%
\pgfpathlineto{\pgfqpoint{3.801566in}{2.275429in}}%
\pgfpathmoveto{\pgfqpoint{3.797025in}{2.278378in}}%
\pgfpathlineto{\pgfqpoint{3.797025in}{2.278378in}}%
\pgfpathlineto{\pgfqpoint{3.797025in}{2.281327in}}%
\pgfpathlineto{\pgfqpoint{3.801566in}{2.281327in}}%
\pgfpathlineto{\pgfqpoint{3.801566in}{2.278378in}}%
\pgfpathmoveto{\pgfqpoint{3.787943in}{2.284276in}}%
\pgfpathlineto{\pgfqpoint{3.787943in}{2.284276in}}%
\pgfpathlineto{\pgfqpoint{3.787943in}{2.287225in}}%
\pgfpathlineto{\pgfqpoint{3.792484in}{2.287225in}}%
\pgfpathlineto{\pgfqpoint{3.792484in}{2.284276in}}%
\pgfpathmoveto{\pgfqpoint{3.783402in}{2.290174in}}%
\pgfpathlineto{\pgfqpoint{3.783402in}{2.290174in}}%
\pgfpathlineto{\pgfqpoint{3.783402in}{2.293123in}}%
\pgfpathlineto{\pgfqpoint{3.787943in}{2.293123in}}%
\pgfpathlineto{\pgfqpoint{3.787943in}{2.290174in}}%
\pgfpathmoveto{\pgfqpoint{3.787943in}{2.287225in}}%
\pgfpathlineto{\pgfqpoint{3.787943in}{2.287225in}}%
\pgfpathlineto{\pgfqpoint{3.787943in}{2.290174in}}%
\pgfpathlineto{\pgfqpoint{3.792484in}{2.290174in}}%
\pgfpathlineto{\pgfqpoint{3.792484in}{2.287225in}}%
\pgfpathmoveto{\pgfqpoint{3.787943in}{2.290174in}}%
\pgfpathlineto{\pgfqpoint{3.787943in}{2.290174in}}%
\pgfpathlineto{\pgfqpoint{3.787943in}{2.293123in}}%
\pgfpathlineto{\pgfqpoint{3.792484in}{2.293123in}}%
\pgfpathlineto{\pgfqpoint{3.792484in}{2.290174in}}%
\pgfpathmoveto{\pgfqpoint{3.792484in}{2.281327in}}%
\pgfpathlineto{\pgfqpoint{3.792484in}{2.281327in}}%
\pgfpathlineto{\pgfqpoint{3.792484in}{2.284276in}}%
\pgfpathlineto{\pgfqpoint{3.797025in}{2.284276in}}%
\pgfpathlineto{\pgfqpoint{3.797025in}{2.281327in}}%
\pgfpathmoveto{\pgfqpoint{3.792484in}{2.284276in}}%
\pgfpathlineto{\pgfqpoint{3.792484in}{2.284276in}}%
\pgfpathlineto{\pgfqpoint{3.792484in}{2.287225in}}%
\pgfpathlineto{\pgfqpoint{3.797025in}{2.287225in}}%
\pgfpathlineto{\pgfqpoint{3.797025in}{2.284276in}}%
\pgfpathmoveto{\pgfqpoint{3.724367in}{2.366854in}}%
\pgfpathlineto{\pgfqpoint{3.724367in}{2.366854in}}%
\pgfpathlineto{\pgfqpoint{3.724367in}{2.369803in}}%
\pgfpathlineto{\pgfqpoint{3.728908in}{2.369803in}}%
\pgfpathlineto{\pgfqpoint{3.728908in}{2.366854in}}%
\pgfpathmoveto{\pgfqpoint{3.719826in}{2.372752in}}%
\pgfpathlineto{\pgfqpoint{3.719826in}{2.372752in}}%
\pgfpathlineto{\pgfqpoint{3.719826in}{2.375702in}}%
\pgfpathlineto{\pgfqpoint{3.724367in}{2.375702in}}%
\pgfpathlineto{\pgfqpoint{3.724367in}{2.372752in}}%
\pgfpathmoveto{\pgfqpoint{3.724367in}{2.369803in}}%
\pgfpathlineto{\pgfqpoint{3.724367in}{2.369803in}}%
\pgfpathlineto{\pgfqpoint{3.724367in}{2.372752in}}%
\pgfpathlineto{\pgfqpoint{3.728908in}{2.372752in}}%
\pgfpathlineto{\pgfqpoint{3.728908in}{2.369803in}}%
\pgfpathmoveto{\pgfqpoint{3.724367in}{2.372752in}}%
\pgfpathlineto{\pgfqpoint{3.724367in}{2.372752in}}%
\pgfpathlineto{\pgfqpoint{3.724367in}{2.375702in}}%
\pgfpathlineto{\pgfqpoint{3.728908in}{2.375702in}}%
\pgfpathlineto{\pgfqpoint{3.728908in}{2.372752in}}%
\pgfpathmoveto{\pgfqpoint{3.715285in}{2.378651in}}%
\pgfpathlineto{\pgfqpoint{3.715285in}{2.378651in}}%
\pgfpathlineto{\pgfqpoint{3.715285in}{2.381600in}}%
\pgfpathlineto{\pgfqpoint{3.719826in}{2.381600in}}%
\pgfpathlineto{\pgfqpoint{3.719826in}{2.378651in}}%
\pgfpathmoveto{\pgfqpoint{3.710743in}{2.384549in}}%
\pgfpathlineto{\pgfqpoint{3.710743in}{2.384549in}}%
\pgfpathlineto{\pgfqpoint{3.710743in}{2.387498in}}%
\pgfpathlineto{\pgfqpoint{3.715285in}{2.387498in}}%
\pgfpathlineto{\pgfqpoint{3.715285in}{2.384549in}}%
\pgfpathmoveto{\pgfqpoint{3.715285in}{2.381600in}}%
\pgfpathlineto{\pgfqpoint{3.715285in}{2.381600in}}%
\pgfpathlineto{\pgfqpoint{3.715285in}{2.384549in}}%
\pgfpathlineto{\pgfqpoint{3.719826in}{2.384549in}}%
\pgfpathlineto{\pgfqpoint{3.719826in}{2.381600in}}%
\pgfpathmoveto{\pgfqpoint{3.715285in}{2.384549in}}%
\pgfpathlineto{\pgfqpoint{3.715285in}{2.384549in}}%
\pgfpathlineto{\pgfqpoint{3.715285in}{2.387498in}}%
\pgfpathlineto{\pgfqpoint{3.719826in}{2.387498in}}%
\pgfpathlineto{\pgfqpoint{3.719826in}{2.384549in}}%
\pgfpathmoveto{\pgfqpoint{3.719826in}{2.375702in}}%
\pgfpathlineto{\pgfqpoint{3.719826in}{2.375702in}}%
\pgfpathlineto{\pgfqpoint{3.719826in}{2.378651in}}%
\pgfpathlineto{\pgfqpoint{3.724367in}{2.378651in}}%
\pgfpathlineto{\pgfqpoint{3.724367in}{2.375702in}}%
\pgfpathmoveto{\pgfqpoint{3.719826in}{2.378651in}}%
\pgfpathlineto{\pgfqpoint{3.719826in}{2.378651in}}%
\pgfpathlineto{\pgfqpoint{3.719826in}{2.381600in}}%
\pgfpathlineto{\pgfqpoint{3.724367in}{2.381600in}}%
\pgfpathlineto{\pgfqpoint{3.724367in}{2.378651in}}%
\pgfpathmoveto{\pgfqpoint{3.760696in}{2.319666in}}%
\pgfpathlineto{\pgfqpoint{3.760696in}{2.319666in}}%
\pgfpathlineto{\pgfqpoint{3.760696in}{2.322615in}}%
\pgfpathlineto{\pgfqpoint{3.765237in}{2.322615in}}%
\pgfpathlineto{\pgfqpoint{3.765237in}{2.319666in}}%
\pgfpathmoveto{\pgfqpoint{3.756155in}{2.325565in}}%
\pgfpathlineto{\pgfqpoint{3.756155in}{2.325565in}}%
\pgfpathlineto{\pgfqpoint{3.756155in}{2.328514in}}%
\pgfpathlineto{\pgfqpoint{3.760696in}{2.328514in}}%
\pgfpathlineto{\pgfqpoint{3.760696in}{2.325565in}}%
\pgfpathmoveto{\pgfqpoint{3.760696in}{2.322615in}}%
\pgfpathlineto{\pgfqpoint{3.760696in}{2.322615in}}%
\pgfpathlineto{\pgfqpoint{3.760696in}{2.325565in}}%
\pgfpathlineto{\pgfqpoint{3.765237in}{2.325565in}}%
\pgfpathlineto{\pgfqpoint{3.765237in}{2.322615in}}%
\pgfpathmoveto{\pgfqpoint{3.760696in}{2.325565in}}%
\pgfpathlineto{\pgfqpoint{3.760696in}{2.325565in}}%
\pgfpathlineto{\pgfqpoint{3.760696in}{2.328514in}}%
\pgfpathlineto{\pgfqpoint{3.765237in}{2.328514in}}%
\pgfpathlineto{\pgfqpoint{3.765237in}{2.325565in}}%
\pgfpathmoveto{\pgfqpoint{3.751614in}{2.331463in}}%
\pgfpathlineto{\pgfqpoint{3.751614in}{2.331463in}}%
\pgfpathlineto{\pgfqpoint{3.751614in}{2.334412in}}%
\pgfpathlineto{\pgfqpoint{3.756155in}{2.334412in}}%
\pgfpathlineto{\pgfqpoint{3.756155in}{2.331463in}}%
\pgfpathmoveto{\pgfqpoint{3.747073in}{2.337362in}}%
\pgfpathlineto{\pgfqpoint{3.747073in}{2.337362in}}%
\pgfpathlineto{\pgfqpoint{3.747073in}{2.340311in}}%
\pgfpathlineto{\pgfqpoint{3.751614in}{2.340311in}}%
\pgfpathlineto{\pgfqpoint{3.751614in}{2.337362in}}%
\pgfpathmoveto{\pgfqpoint{3.751614in}{2.334412in}}%
\pgfpathlineto{\pgfqpoint{3.751614in}{2.334412in}}%
\pgfpathlineto{\pgfqpoint{3.751614in}{2.337362in}}%
\pgfpathlineto{\pgfqpoint{3.756155in}{2.337362in}}%
\pgfpathlineto{\pgfqpoint{3.756155in}{2.334412in}}%
\pgfpathmoveto{\pgfqpoint{3.751614in}{2.337362in}}%
\pgfpathlineto{\pgfqpoint{3.751614in}{2.337362in}}%
\pgfpathlineto{\pgfqpoint{3.751614in}{2.340311in}}%
\pgfpathlineto{\pgfqpoint{3.756155in}{2.340311in}}%
\pgfpathlineto{\pgfqpoint{3.756155in}{2.337362in}}%
\pgfpathmoveto{\pgfqpoint{3.756155in}{2.328514in}}%
\pgfpathlineto{\pgfqpoint{3.756155in}{2.328514in}}%
\pgfpathlineto{\pgfqpoint{3.756155in}{2.331463in}}%
\pgfpathlineto{\pgfqpoint{3.760696in}{2.331463in}}%
\pgfpathlineto{\pgfqpoint{3.760696in}{2.328514in}}%
\pgfpathmoveto{\pgfqpoint{3.756155in}{2.331463in}}%
\pgfpathlineto{\pgfqpoint{3.756155in}{2.331463in}}%
\pgfpathlineto{\pgfqpoint{3.756155in}{2.334412in}}%
\pgfpathlineto{\pgfqpoint{3.760696in}{2.334412in}}%
\pgfpathlineto{\pgfqpoint{3.760696in}{2.331463in}}%
\pgfpathmoveto{\pgfqpoint{3.778861in}{2.296072in}}%
\pgfpathlineto{\pgfqpoint{3.778861in}{2.296072in}}%
\pgfpathlineto{\pgfqpoint{3.778861in}{2.299022in}}%
\pgfpathlineto{\pgfqpoint{3.783402in}{2.299022in}}%
\pgfpathlineto{\pgfqpoint{3.783402in}{2.296072in}}%
\pgfpathmoveto{\pgfqpoint{3.774320in}{2.301971in}}%
\pgfpathlineto{\pgfqpoint{3.774320in}{2.301971in}}%
\pgfpathlineto{\pgfqpoint{3.774320in}{2.304920in}}%
\pgfpathlineto{\pgfqpoint{3.778861in}{2.304920in}}%
\pgfpathlineto{\pgfqpoint{3.778861in}{2.301971in}}%
\pgfpathmoveto{\pgfqpoint{3.778861in}{2.299022in}}%
\pgfpathlineto{\pgfqpoint{3.778861in}{2.299022in}}%
\pgfpathlineto{\pgfqpoint{3.778861in}{2.301971in}}%
\pgfpathlineto{\pgfqpoint{3.783402in}{2.301971in}}%
\pgfpathlineto{\pgfqpoint{3.783402in}{2.299022in}}%
\pgfpathmoveto{\pgfqpoint{3.778861in}{2.301971in}}%
\pgfpathlineto{\pgfqpoint{3.778861in}{2.301971in}}%
\pgfpathlineto{\pgfqpoint{3.778861in}{2.304920in}}%
\pgfpathlineto{\pgfqpoint{3.783402in}{2.304920in}}%
\pgfpathlineto{\pgfqpoint{3.783402in}{2.301971in}}%
\pgfpathmoveto{\pgfqpoint{3.769778in}{2.307869in}}%
\pgfpathlineto{\pgfqpoint{3.769778in}{2.307869in}}%
\pgfpathlineto{\pgfqpoint{3.769778in}{2.310819in}}%
\pgfpathlineto{\pgfqpoint{3.774320in}{2.310819in}}%
\pgfpathlineto{\pgfqpoint{3.774320in}{2.307869in}}%
\pgfpathmoveto{\pgfqpoint{3.765237in}{2.313768in}}%
\pgfpathlineto{\pgfqpoint{3.765237in}{2.313768in}}%
\pgfpathlineto{\pgfqpoint{3.765237in}{2.316717in}}%
\pgfpathlineto{\pgfqpoint{3.769778in}{2.316717in}}%
\pgfpathlineto{\pgfqpoint{3.769778in}{2.313768in}}%
\pgfpathmoveto{\pgfqpoint{3.769778in}{2.310819in}}%
\pgfpathlineto{\pgfqpoint{3.769778in}{2.310819in}}%
\pgfpathlineto{\pgfqpoint{3.769778in}{2.313768in}}%
\pgfpathlineto{\pgfqpoint{3.774320in}{2.313768in}}%
\pgfpathlineto{\pgfqpoint{3.774320in}{2.310819in}}%
\pgfpathmoveto{\pgfqpoint{3.769778in}{2.313768in}}%
\pgfpathlineto{\pgfqpoint{3.769778in}{2.313768in}}%
\pgfpathlineto{\pgfqpoint{3.769778in}{2.316717in}}%
\pgfpathlineto{\pgfqpoint{3.774320in}{2.316717in}}%
\pgfpathlineto{\pgfqpoint{3.774320in}{2.313768in}}%
\pgfpathmoveto{\pgfqpoint{3.774320in}{2.304920in}}%
\pgfpathlineto{\pgfqpoint{3.774320in}{2.304920in}}%
\pgfpathlineto{\pgfqpoint{3.774320in}{2.307869in}}%
\pgfpathlineto{\pgfqpoint{3.778861in}{2.307869in}}%
\pgfpathlineto{\pgfqpoint{3.778861in}{2.304920in}}%
\pgfpathmoveto{\pgfqpoint{3.774320in}{2.307869in}}%
\pgfpathlineto{\pgfqpoint{3.774320in}{2.307869in}}%
\pgfpathlineto{\pgfqpoint{3.774320in}{2.310819in}}%
\pgfpathlineto{\pgfqpoint{3.778861in}{2.310819in}}%
\pgfpathlineto{\pgfqpoint{3.778861in}{2.307869in}}%
\pgfpathmoveto{\pgfqpoint{3.783402in}{2.293123in}}%
\pgfpathlineto{\pgfqpoint{3.783402in}{2.293123in}}%
\pgfpathlineto{\pgfqpoint{3.783402in}{2.296072in}}%
\pgfpathlineto{\pgfqpoint{3.787943in}{2.296072in}}%
\pgfpathlineto{\pgfqpoint{3.787943in}{2.293123in}}%
\pgfpathmoveto{\pgfqpoint{3.783402in}{2.296072in}}%
\pgfpathlineto{\pgfqpoint{3.783402in}{2.296072in}}%
\pgfpathlineto{\pgfqpoint{3.783402in}{2.299022in}}%
\pgfpathlineto{\pgfqpoint{3.787943in}{2.299022in}}%
\pgfpathlineto{\pgfqpoint{3.787943in}{2.296072in}}%
\pgfpathmoveto{\pgfqpoint{3.765237in}{2.316717in}}%
\pgfpathlineto{\pgfqpoint{3.765237in}{2.316717in}}%
\pgfpathlineto{\pgfqpoint{3.765237in}{2.319666in}}%
\pgfpathlineto{\pgfqpoint{3.769778in}{2.319666in}}%
\pgfpathlineto{\pgfqpoint{3.769778in}{2.316717in}}%
\pgfpathmoveto{\pgfqpoint{3.765237in}{2.319666in}}%
\pgfpathlineto{\pgfqpoint{3.765237in}{2.319666in}}%
\pgfpathlineto{\pgfqpoint{3.765237in}{2.322615in}}%
\pgfpathlineto{\pgfqpoint{3.769778in}{2.322615in}}%
\pgfpathlineto{\pgfqpoint{3.769778in}{2.319666in}}%
\pgfpathmoveto{\pgfqpoint{3.742531in}{2.343260in}}%
\pgfpathlineto{\pgfqpoint{3.742531in}{2.343260in}}%
\pgfpathlineto{\pgfqpoint{3.742531in}{2.346209in}}%
\pgfpathlineto{\pgfqpoint{3.747073in}{2.346209in}}%
\pgfpathlineto{\pgfqpoint{3.747073in}{2.343260in}}%
\pgfpathmoveto{\pgfqpoint{3.737990in}{2.349158in}}%
\pgfpathlineto{\pgfqpoint{3.737990in}{2.349158in}}%
\pgfpathlineto{\pgfqpoint{3.737990in}{2.352108in}}%
\pgfpathlineto{\pgfqpoint{3.742531in}{2.352108in}}%
\pgfpathlineto{\pgfqpoint{3.742531in}{2.349158in}}%
\pgfpathmoveto{\pgfqpoint{3.742531in}{2.346209in}}%
\pgfpathlineto{\pgfqpoint{3.742531in}{2.346209in}}%
\pgfpathlineto{\pgfqpoint{3.742531in}{2.349158in}}%
\pgfpathlineto{\pgfqpoint{3.747073in}{2.349158in}}%
\pgfpathlineto{\pgfqpoint{3.747073in}{2.346209in}}%
\pgfpathmoveto{\pgfqpoint{3.742531in}{2.349158in}}%
\pgfpathlineto{\pgfqpoint{3.742531in}{2.349158in}}%
\pgfpathlineto{\pgfqpoint{3.742531in}{2.352108in}}%
\pgfpathlineto{\pgfqpoint{3.747073in}{2.352108in}}%
\pgfpathlineto{\pgfqpoint{3.747073in}{2.349158in}}%
\pgfpathmoveto{\pgfqpoint{3.733449in}{2.355057in}}%
\pgfpathlineto{\pgfqpoint{3.733449in}{2.355057in}}%
\pgfpathlineto{\pgfqpoint{3.733449in}{2.358006in}}%
\pgfpathlineto{\pgfqpoint{3.737990in}{2.358006in}}%
\pgfpathlineto{\pgfqpoint{3.737990in}{2.355057in}}%
\pgfpathmoveto{\pgfqpoint{3.728908in}{2.360955in}}%
\pgfpathlineto{\pgfqpoint{3.728908in}{2.360955in}}%
\pgfpathlineto{\pgfqpoint{3.728908in}{2.363905in}}%
\pgfpathlineto{\pgfqpoint{3.733449in}{2.363905in}}%
\pgfpathlineto{\pgfqpoint{3.733449in}{2.360955in}}%
\pgfpathmoveto{\pgfqpoint{3.733449in}{2.358006in}}%
\pgfpathlineto{\pgfqpoint{3.733449in}{2.358006in}}%
\pgfpathlineto{\pgfqpoint{3.733449in}{2.360955in}}%
\pgfpathlineto{\pgfqpoint{3.737990in}{2.360955in}}%
\pgfpathlineto{\pgfqpoint{3.737990in}{2.358006in}}%
\pgfpathmoveto{\pgfqpoint{3.733449in}{2.360955in}}%
\pgfpathlineto{\pgfqpoint{3.733449in}{2.360955in}}%
\pgfpathlineto{\pgfqpoint{3.733449in}{2.363905in}}%
\pgfpathlineto{\pgfqpoint{3.737990in}{2.363905in}}%
\pgfpathlineto{\pgfqpoint{3.737990in}{2.360955in}}%
\pgfpathmoveto{\pgfqpoint{3.737990in}{2.352108in}}%
\pgfpathlineto{\pgfqpoint{3.737990in}{2.352108in}}%
\pgfpathlineto{\pgfqpoint{3.737990in}{2.355057in}}%
\pgfpathlineto{\pgfqpoint{3.742531in}{2.355057in}}%
\pgfpathlineto{\pgfqpoint{3.742531in}{2.352108in}}%
\pgfpathmoveto{\pgfqpoint{3.737990in}{2.355057in}}%
\pgfpathlineto{\pgfqpoint{3.737990in}{2.355057in}}%
\pgfpathlineto{\pgfqpoint{3.737990in}{2.358006in}}%
\pgfpathlineto{\pgfqpoint{3.742531in}{2.358006in}}%
\pgfpathlineto{\pgfqpoint{3.742531in}{2.355057in}}%
\pgfpathmoveto{\pgfqpoint{3.747073in}{2.340311in}}%
\pgfpathlineto{\pgfqpoint{3.747073in}{2.340311in}}%
\pgfpathlineto{\pgfqpoint{3.747073in}{2.343260in}}%
\pgfpathlineto{\pgfqpoint{3.751614in}{2.343260in}}%
\pgfpathlineto{\pgfqpoint{3.751614in}{2.340311in}}%
\pgfpathmoveto{\pgfqpoint{3.747073in}{2.343260in}}%
\pgfpathlineto{\pgfqpoint{3.747073in}{2.343260in}}%
\pgfpathlineto{\pgfqpoint{3.747073in}{2.346209in}}%
\pgfpathlineto{\pgfqpoint{3.751614in}{2.346209in}}%
\pgfpathlineto{\pgfqpoint{3.751614in}{2.343260in}}%
\pgfpathmoveto{\pgfqpoint{3.728908in}{2.363905in}}%
\pgfpathlineto{\pgfqpoint{3.728908in}{2.363905in}}%
\pgfpathlineto{\pgfqpoint{3.728908in}{2.366854in}}%
\pgfpathlineto{\pgfqpoint{3.733449in}{2.366854in}}%
\pgfpathlineto{\pgfqpoint{3.733449in}{2.363905in}}%
\pgfpathmoveto{\pgfqpoint{3.728908in}{2.366854in}}%
\pgfpathlineto{\pgfqpoint{3.728908in}{2.366854in}}%
\pgfpathlineto{\pgfqpoint{3.728908in}{2.369803in}}%
\pgfpathlineto{\pgfqpoint{3.733449in}{2.369803in}}%
\pgfpathlineto{\pgfqpoint{3.733449in}{2.366854in}}%
\pgfpathmoveto{\pgfqpoint{3.688038in}{2.414043in}}%
\pgfpathlineto{\pgfqpoint{3.688038in}{2.414043in}}%
\pgfpathlineto{\pgfqpoint{3.688038in}{2.416992in}}%
\pgfpathlineto{\pgfqpoint{3.692579in}{2.416992in}}%
\pgfpathlineto{\pgfqpoint{3.692579in}{2.414043in}}%
\pgfpathmoveto{\pgfqpoint{3.683497in}{2.419941in}}%
\pgfpathlineto{\pgfqpoint{3.683497in}{2.419941in}}%
\pgfpathlineto{\pgfqpoint{3.683497in}{2.422891in}}%
\pgfpathlineto{\pgfqpoint{3.688038in}{2.422891in}}%
\pgfpathlineto{\pgfqpoint{3.688038in}{2.419941in}}%
\pgfpathmoveto{\pgfqpoint{3.688038in}{2.416992in}}%
\pgfpathlineto{\pgfqpoint{3.688038in}{2.416992in}}%
\pgfpathlineto{\pgfqpoint{3.688038in}{2.419941in}}%
\pgfpathlineto{\pgfqpoint{3.692579in}{2.419941in}}%
\pgfpathlineto{\pgfqpoint{3.692579in}{2.416992in}}%
\pgfpathmoveto{\pgfqpoint{3.688038in}{2.419941in}}%
\pgfpathlineto{\pgfqpoint{3.688038in}{2.419941in}}%
\pgfpathlineto{\pgfqpoint{3.688038in}{2.422891in}}%
\pgfpathlineto{\pgfqpoint{3.692579in}{2.422891in}}%
\pgfpathlineto{\pgfqpoint{3.692579in}{2.419941in}}%
\pgfpathmoveto{\pgfqpoint{3.678955in}{2.425840in}}%
\pgfpathlineto{\pgfqpoint{3.678955in}{2.425840in}}%
\pgfpathlineto{\pgfqpoint{3.678955in}{2.428789in}}%
\pgfpathlineto{\pgfqpoint{3.683497in}{2.428789in}}%
\pgfpathlineto{\pgfqpoint{3.683497in}{2.425840in}}%
\pgfpathmoveto{\pgfqpoint{3.674414in}{2.431739in}}%
\pgfpathlineto{\pgfqpoint{3.674414in}{2.431739in}}%
\pgfpathlineto{\pgfqpoint{3.674414in}{2.434688in}}%
\pgfpathlineto{\pgfqpoint{3.678955in}{2.434688in}}%
\pgfpathlineto{\pgfqpoint{3.678955in}{2.431739in}}%
\pgfpathmoveto{\pgfqpoint{3.678955in}{2.428789in}}%
\pgfpathlineto{\pgfqpoint{3.678955in}{2.428789in}}%
\pgfpathlineto{\pgfqpoint{3.678955in}{2.431739in}}%
\pgfpathlineto{\pgfqpoint{3.683497in}{2.431739in}}%
\pgfpathlineto{\pgfqpoint{3.683497in}{2.428789in}}%
\pgfpathmoveto{\pgfqpoint{3.678955in}{2.431739in}}%
\pgfpathlineto{\pgfqpoint{3.678955in}{2.431739in}}%
\pgfpathlineto{\pgfqpoint{3.678955in}{2.434688in}}%
\pgfpathlineto{\pgfqpoint{3.683497in}{2.434688in}}%
\pgfpathlineto{\pgfqpoint{3.683497in}{2.431739in}}%
\pgfpathmoveto{\pgfqpoint{3.683497in}{2.422891in}}%
\pgfpathlineto{\pgfqpoint{3.683497in}{2.422891in}}%
\pgfpathlineto{\pgfqpoint{3.683497in}{2.425840in}}%
\pgfpathlineto{\pgfqpoint{3.688038in}{2.425840in}}%
\pgfpathlineto{\pgfqpoint{3.688038in}{2.422891in}}%
\pgfpathmoveto{\pgfqpoint{3.683497in}{2.425840in}}%
\pgfpathlineto{\pgfqpoint{3.683497in}{2.425840in}}%
\pgfpathlineto{\pgfqpoint{3.683497in}{2.428789in}}%
\pgfpathlineto{\pgfqpoint{3.688038in}{2.428789in}}%
\pgfpathlineto{\pgfqpoint{3.688038in}{2.425840in}}%
\pgfpathmoveto{\pgfqpoint{3.706202in}{2.390448in}}%
\pgfpathlineto{\pgfqpoint{3.706202in}{2.390448in}}%
\pgfpathlineto{\pgfqpoint{3.706202in}{2.393397in}}%
\pgfpathlineto{\pgfqpoint{3.710743in}{2.393397in}}%
\pgfpathlineto{\pgfqpoint{3.710743in}{2.390448in}}%
\pgfpathmoveto{\pgfqpoint{3.701661in}{2.396346in}}%
\pgfpathlineto{\pgfqpoint{3.701661in}{2.396346in}}%
\pgfpathlineto{\pgfqpoint{3.701661in}{2.399296in}}%
\pgfpathlineto{\pgfqpoint{3.706202in}{2.399296in}}%
\pgfpathlineto{\pgfqpoint{3.706202in}{2.396346in}}%
\pgfpathmoveto{\pgfqpoint{3.706202in}{2.393397in}}%
\pgfpathlineto{\pgfqpoint{3.706202in}{2.393397in}}%
\pgfpathlineto{\pgfqpoint{3.706202in}{2.396346in}}%
\pgfpathlineto{\pgfqpoint{3.710743in}{2.396346in}}%
\pgfpathlineto{\pgfqpoint{3.710743in}{2.393397in}}%
\pgfpathmoveto{\pgfqpoint{3.706202in}{2.396346in}}%
\pgfpathlineto{\pgfqpoint{3.706202in}{2.396346in}}%
\pgfpathlineto{\pgfqpoint{3.706202in}{2.399296in}}%
\pgfpathlineto{\pgfqpoint{3.710743in}{2.399296in}}%
\pgfpathlineto{\pgfqpoint{3.710743in}{2.396346in}}%
\pgfpathmoveto{\pgfqpoint{3.697120in}{2.402245in}}%
\pgfpathlineto{\pgfqpoint{3.697120in}{2.402245in}}%
\pgfpathlineto{\pgfqpoint{3.697120in}{2.405195in}}%
\pgfpathlineto{\pgfqpoint{3.701661in}{2.405195in}}%
\pgfpathlineto{\pgfqpoint{3.701661in}{2.402245in}}%
\pgfpathmoveto{\pgfqpoint{3.692579in}{2.408144in}}%
\pgfpathlineto{\pgfqpoint{3.692579in}{2.408144in}}%
\pgfpathlineto{\pgfqpoint{3.692579in}{2.411093in}}%
\pgfpathlineto{\pgfqpoint{3.697120in}{2.411093in}}%
\pgfpathlineto{\pgfqpoint{3.697120in}{2.408144in}}%
\pgfpathmoveto{\pgfqpoint{3.697120in}{2.405195in}}%
\pgfpathlineto{\pgfqpoint{3.697120in}{2.405195in}}%
\pgfpathlineto{\pgfqpoint{3.697120in}{2.408144in}}%
\pgfpathlineto{\pgfqpoint{3.701661in}{2.408144in}}%
\pgfpathlineto{\pgfqpoint{3.701661in}{2.405195in}}%
\pgfpathmoveto{\pgfqpoint{3.697120in}{2.408144in}}%
\pgfpathlineto{\pgfqpoint{3.697120in}{2.408144in}}%
\pgfpathlineto{\pgfqpoint{3.697120in}{2.411093in}}%
\pgfpathlineto{\pgfqpoint{3.701661in}{2.411093in}}%
\pgfpathlineto{\pgfqpoint{3.701661in}{2.408144in}}%
\pgfpathmoveto{\pgfqpoint{3.701661in}{2.399296in}}%
\pgfpathlineto{\pgfqpoint{3.701661in}{2.399296in}}%
\pgfpathlineto{\pgfqpoint{3.701661in}{2.402245in}}%
\pgfpathlineto{\pgfqpoint{3.706202in}{2.402245in}}%
\pgfpathlineto{\pgfqpoint{3.706202in}{2.399296in}}%
\pgfpathmoveto{\pgfqpoint{3.701661in}{2.402245in}}%
\pgfpathlineto{\pgfqpoint{3.701661in}{2.402245in}}%
\pgfpathlineto{\pgfqpoint{3.701661in}{2.405195in}}%
\pgfpathlineto{\pgfqpoint{3.706202in}{2.405195in}}%
\pgfpathlineto{\pgfqpoint{3.706202in}{2.402245in}}%
\pgfpathmoveto{\pgfqpoint{3.710743in}{2.387498in}}%
\pgfpathlineto{\pgfqpoint{3.710743in}{2.387498in}}%
\pgfpathlineto{\pgfqpoint{3.710743in}{2.390448in}}%
\pgfpathlineto{\pgfqpoint{3.715285in}{2.390448in}}%
\pgfpathlineto{\pgfqpoint{3.715285in}{2.387498in}}%
\pgfpathmoveto{\pgfqpoint{3.710743in}{2.390448in}}%
\pgfpathlineto{\pgfqpoint{3.710743in}{2.390448in}}%
\pgfpathlineto{\pgfqpoint{3.710743in}{2.393397in}}%
\pgfpathlineto{\pgfqpoint{3.715285in}{2.393397in}}%
\pgfpathlineto{\pgfqpoint{3.715285in}{2.390448in}}%
\pgfpathmoveto{\pgfqpoint{3.692579in}{2.411093in}}%
\pgfpathlineto{\pgfqpoint{3.692579in}{2.411093in}}%
\pgfpathlineto{\pgfqpoint{3.692579in}{2.414043in}}%
\pgfpathlineto{\pgfqpoint{3.697120in}{2.414043in}}%
\pgfpathlineto{\pgfqpoint{3.697120in}{2.411093in}}%
\pgfpathmoveto{\pgfqpoint{3.692579in}{2.414043in}}%
\pgfpathlineto{\pgfqpoint{3.692579in}{2.414043in}}%
\pgfpathlineto{\pgfqpoint{3.692579in}{2.416992in}}%
\pgfpathlineto{\pgfqpoint{3.697120in}{2.416992in}}%
\pgfpathlineto{\pgfqpoint{3.697120in}{2.414043in}}%
\pgfpathmoveto{\pgfqpoint{3.669873in}{2.437637in}}%
\pgfpathlineto{\pgfqpoint{3.669873in}{2.437637in}}%
\pgfpathlineto{\pgfqpoint{3.669873in}{2.440587in}}%
\pgfpathlineto{\pgfqpoint{3.674414in}{2.440587in}}%
\pgfpathlineto{\pgfqpoint{3.674414in}{2.437637in}}%
\pgfpathmoveto{\pgfqpoint{3.665332in}{2.443536in}}%
\pgfpathlineto{\pgfqpoint{3.665332in}{2.443536in}}%
\pgfpathlineto{\pgfqpoint{3.665332in}{2.446486in}}%
\pgfpathlineto{\pgfqpoint{3.669873in}{2.446486in}}%
\pgfpathlineto{\pgfqpoint{3.669873in}{2.443536in}}%
\pgfpathmoveto{\pgfqpoint{3.669873in}{2.440587in}}%
\pgfpathlineto{\pgfqpoint{3.669873in}{2.440587in}}%
\pgfpathlineto{\pgfqpoint{3.669873in}{2.443536in}}%
\pgfpathlineto{\pgfqpoint{3.674414in}{2.443536in}}%
\pgfpathlineto{\pgfqpoint{3.674414in}{2.440587in}}%
\pgfpathmoveto{\pgfqpoint{3.669873in}{2.443536in}}%
\pgfpathlineto{\pgfqpoint{3.669873in}{2.443536in}}%
\pgfpathlineto{\pgfqpoint{3.669873in}{2.446486in}}%
\pgfpathlineto{\pgfqpoint{3.674414in}{2.446486in}}%
\pgfpathlineto{\pgfqpoint{3.674414in}{2.443536in}}%
\pgfpathmoveto{\pgfqpoint{3.660791in}{2.449435in}}%
\pgfpathlineto{\pgfqpoint{3.660791in}{2.449435in}}%
\pgfpathlineto{\pgfqpoint{3.660791in}{2.452384in}}%
\pgfpathlineto{\pgfqpoint{3.665332in}{2.452384in}}%
\pgfpathlineto{\pgfqpoint{3.665332in}{2.449435in}}%
\pgfpathmoveto{\pgfqpoint{3.656250in}{2.455334in}}%
\pgfpathlineto{\pgfqpoint{3.656250in}{2.455334in}}%
\pgfpathlineto{\pgfqpoint{3.656250in}{2.458283in}}%
\pgfpathlineto{\pgfqpoint{3.660791in}{2.458283in}}%
\pgfpathlineto{\pgfqpoint{3.660791in}{2.455334in}}%
\pgfpathmoveto{\pgfqpoint{3.660791in}{2.452384in}}%
\pgfpathlineto{\pgfqpoint{3.660791in}{2.452384in}}%
\pgfpathlineto{\pgfqpoint{3.660791in}{2.455334in}}%
\pgfpathlineto{\pgfqpoint{3.665332in}{2.455334in}}%
\pgfpathlineto{\pgfqpoint{3.665332in}{2.452384in}}%
\pgfpathmoveto{\pgfqpoint{3.660791in}{2.455334in}}%
\pgfpathlineto{\pgfqpoint{3.660791in}{2.455334in}}%
\pgfpathlineto{\pgfqpoint{3.660791in}{2.458283in}}%
\pgfpathlineto{\pgfqpoint{3.665332in}{2.458283in}}%
\pgfpathlineto{\pgfqpoint{3.665332in}{2.455334in}}%
\pgfpathmoveto{\pgfqpoint{3.665332in}{2.446486in}}%
\pgfpathlineto{\pgfqpoint{3.665332in}{2.446486in}}%
\pgfpathlineto{\pgfqpoint{3.665332in}{2.449435in}}%
\pgfpathlineto{\pgfqpoint{3.669873in}{2.449435in}}%
\pgfpathlineto{\pgfqpoint{3.669873in}{2.446486in}}%
\pgfpathmoveto{\pgfqpoint{3.665332in}{2.449435in}}%
\pgfpathlineto{\pgfqpoint{3.665332in}{2.449435in}}%
\pgfpathlineto{\pgfqpoint{3.665332in}{2.452384in}}%
\pgfpathlineto{\pgfqpoint{3.669873in}{2.452384in}}%
\pgfpathlineto{\pgfqpoint{3.669873in}{2.449435in}}%
\pgfpathmoveto{\pgfqpoint{3.674414in}{2.434688in}}%
\pgfpathlineto{\pgfqpoint{3.674414in}{2.434688in}}%
\pgfpathlineto{\pgfqpoint{3.674414in}{2.437637in}}%
\pgfpathlineto{\pgfqpoint{3.678955in}{2.437637in}}%
\pgfpathlineto{\pgfqpoint{3.678955in}{2.434688in}}%
\pgfpathmoveto{\pgfqpoint{3.674414in}{2.437637in}}%
\pgfpathlineto{\pgfqpoint{3.674414in}{2.437637in}}%
\pgfpathlineto{\pgfqpoint{3.674414in}{2.440587in}}%
\pgfpathlineto{\pgfqpoint{3.678955in}{2.440587in}}%
\pgfpathlineto{\pgfqpoint{3.678955in}{2.437637in}}%
\pgfpathmoveto{\pgfqpoint{3.656250in}{2.458283in}}%
\pgfpathlineto{\pgfqpoint{3.656250in}{2.458283in}}%
\pgfpathlineto{\pgfqpoint{3.656250in}{2.461232in}}%
\pgfpathlineto{\pgfqpoint{3.660791in}{2.461232in}}%
\pgfpathlineto{\pgfqpoint{3.660791in}{2.458283in}}%
\pgfpathmoveto{\pgfqpoint{3.656250in}{2.461232in}}%
\pgfpathlineto{\pgfqpoint{3.656250in}{2.461232in}}%
\pgfpathlineto{\pgfqpoint{3.656250in}{2.464182in}}%
\pgfpathlineto{\pgfqpoint{3.660791in}{2.464182in}}%
\pgfpathlineto{\pgfqpoint{3.660791in}{2.461232in}}%
\pgfpathmoveto{\pgfqpoint{3.942337in}{2.083733in}}%
\pgfpathlineto{\pgfqpoint{3.942337in}{2.083733in}}%
\pgfpathlineto{\pgfqpoint{3.942337in}{2.086682in}}%
\pgfpathlineto{\pgfqpoint{3.946878in}{2.086682in}}%
\pgfpathlineto{\pgfqpoint{3.946878in}{2.083733in}}%
\pgfpathmoveto{\pgfqpoint{3.937796in}{2.089631in}}%
\pgfpathlineto{\pgfqpoint{3.937796in}{2.089631in}}%
\pgfpathlineto{\pgfqpoint{3.937796in}{2.092580in}}%
\pgfpathlineto{\pgfqpoint{3.942337in}{2.092580in}}%
\pgfpathlineto{\pgfqpoint{3.942337in}{2.089631in}}%
\pgfpathmoveto{\pgfqpoint{3.942337in}{2.086682in}}%
\pgfpathlineto{\pgfqpoint{3.942337in}{2.086682in}}%
\pgfpathlineto{\pgfqpoint{3.942337in}{2.089631in}}%
\pgfpathlineto{\pgfqpoint{3.946878in}{2.089631in}}%
\pgfpathlineto{\pgfqpoint{3.946878in}{2.086682in}}%
\pgfpathmoveto{\pgfqpoint{3.942337in}{2.089631in}}%
\pgfpathlineto{\pgfqpoint{3.942337in}{2.089631in}}%
\pgfpathlineto{\pgfqpoint{3.942337in}{2.092580in}}%
\pgfpathlineto{\pgfqpoint{3.946878in}{2.092580in}}%
\pgfpathlineto{\pgfqpoint{3.946878in}{2.089631in}}%
\pgfpathmoveto{\pgfqpoint{3.933255in}{2.095530in}}%
\pgfpathlineto{\pgfqpoint{3.933255in}{2.095530in}}%
\pgfpathlineto{\pgfqpoint{3.933255in}{2.098479in}}%
\pgfpathlineto{\pgfqpoint{3.937796in}{2.098479in}}%
\pgfpathlineto{\pgfqpoint{3.937796in}{2.095530in}}%
\pgfpathmoveto{\pgfqpoint{3.928714in}{2.101428in}}%
\pgfpathlineto{\pgfqpoint{3.928714in}{2.101428in}}%
\pgfpathlineto{\pgfqpoint{3.928714in}{2.104377in}}%
\pgfpathlineto{\pgfqpoint{3.933255in}{2.104377in}}%
\pgfpathlineto{\pgfqpoint{3.933255in}{2.101428in}}%
\pgfpathmoveto{\pgfqpoint{3.933255in}{2.098479in}}%
\pgfpathlineto{\pgfqpoint{3.933255in}{2.098479in}}%
\pgfpathlineto{\pgfqpoint{3.933255in}{2.101428in}}%
\pgfpathlineto{\pgfqpoint{3.937796in}{2.101428in}}%
\pgfpathlineto{\pgfqpoint{3.937796in}{2.098479in}}%
\pgfpathmoveto{\pgfqpoint{3.933255in}{2.101428in}}%
\pgfpathlineto{\pgfqpoint{3.933255in}{2.101428in}}%
\pgfpathlineto{\pgfqpoint{3.933255in}{2.104377in}}%
\pgfpathlineto{\pgfqpoint{3.937796in}{2.104377in}}%
\pgfpathlineto{\pgfqpoint{3.937796in}{2.101428in}}%
\pgfpathmoveto{\pgfqpoint{3.937796in}{2.092580in}}%
\pgfpathlineto{\pgfqpoint{3.937796in}{2.092580in}}%
\pgfpathlineto{\pgfqpoint{3.937796in}{2.095530in}}%
\pgfpathlineto{\pgfqpoint{3.942337in}{2.095530in}}%
\pgfpathlineto{\pgfqpoint{3.942337in}{2.092580in}}%
\pgfpathmoveto{\pgfqpoint{3.937796in}{2.095530in}}%
\pgfpathlineto{\pgfqpoint{3.937796in}{2.095530in}}%
\pgfpathlineto{\pgfqpoint{3.937796in}{2.098479in}}%
\pgfpathlineto{\pgfqpoint{3.942337in}{2.098479in}}%
\pgfpathlineto{\pgfqpoint{3.942337in}{2.095530in}}%
\pgfpathmoveto{\pgfqpoint{3.869681in}{2.178108in}}%
\pgfpathlineto{\pgfqpoint{3.869681in}{2.178108in}}%
\pgfpathlineto{\pgfqpoint{3.869681in}{2.181057in}}%
\pgfpathlineto{\pgfqpoint{3.874222in}{2.181057in}}%
\pgfpathlineto{\pgfqpoint{3.874222in}{2.178108in}}%
\pgfpathmoveto{\pgfqpoint{3.865140in}{2.184006in}}%
\pgfpathlineto{\pgfqpoint{3.865140in}{2.184006in}}%
\pgfpathlineto{\pgfqpoint{3.865140in}{2.186955in}}%
\pgfpathlineto{\pgfqpoint{3.869681in}{2.186955in}}%
\pgfpathlineto{\pgfqpoint{3.869681in}{2.184006in}}%
\pgfpathmoveto{\pgfqpoint{3.869681in}{2.181057in}}%
\pgfpathlineto{\pgfqpoint{3.869681in}{2.181057in}}%
\pgfpathlineto{\pgfqpoint{3.869681in}{2.184006in}}%
\pgfpathlineto{\pgfqpoint{3.874222in}{2.184006in}}%
\pgfpathlineto{\pgfqpoint{3.874222in}{2.181057in}}%
\pgfpathmoveto{\pgfqpoint{3.869681in}{2.184006in}}%
\pgfpathlineto{\pgfqpoint{3.869681in}{2.184006in}}%
\pgfpathlineto{\pgfqpoint{3.869681in}{2.186955in}}%
\pgfpathlineto{\pgfqpoint{3.874222in}{2.186955in}}%
\pgfpathlineto{\pgfqpoint{3.874222in}{2.184006in}}%
\pgfpathmoveto{\pgfqpoint{3.860599in}{2.189905in}}%
\pgfpathlineto{\pgfqpoint{3.860599in}{2.189905in}}%
\pgfpathlineto{\pgfqpoint{3.860599in}{2.192854in}}%
\pgfpathlineto{\pgfqpoint{3.865140in}{2.192854in}}%
\pgfpathlineto{\pgfqpoint{3.865140in}{2.189905in}}%
\pgfpathmoveto{\pgfqpoint{3.856058in}{2.195803in}}%
\pgfpathlineto{\pgfqpoint{3.856058in}{2.195803in}}%
\pgfpathlineto{\pgfqpoint{3.856058in}{2.198752in}}%
\pgfpathlineto{\pgfqpoint{3.860599in}{2.198752in}}%
\pgfpathlineto{\pgfqpoint{3.860599in}{2.195803in}}%
\pgfpathmoveto{\pgfqpoint{3.860599in}{2.192854in}}%
\pgfpathlineto{\pgfqpoint{3.860599in}{2.192854in}}%
\pgfpathlineto{\pgfqpoint{3.860599in}{2.195803in}}%
\pgfpathlineto{\pgfqpoint{3.865140in}{2.195803in}}%
\pgfpathlineto{\pgfqpoint{3.865140in}{2.192854in}}%
\pgfpathmoveto{\pgfqpoint{3.860599in}{2.195803in}}%
\pgfpathlineto{\pgfqpoint{3.860599in}{2.195803in}}%
\pgfpathlineto{\pgfqpoint{3.860599in}{2.198752in}}%
\pgfpathlineto{\pgfqpoint{3.865140in}{2.198752in}}%
\pgfpathlineto{\pgfqpoint{3.865140in}{2.195803in}}%
\pgfpathmoveto{\pgfqpoint{3.865140in}{2.186955in}}%
\pgfpathlineto{\pgfqpoint{3.865140in}{2.186955in}}%
\pgfpathlineto{\pgfqpoint{3.865140in}{2.189905in}}%
\pgfpathlineto{\pgfqpoint{3.869681in}{2.189905in}}%
\pgfpathlineto{\pgfqpoint{3.869681in}{2.186955in}}%
\pgfpathmoveto{\pgfqpoint{3.865140in}{2.189905in}}%
\pgfpathlineto{\pgfqpoint{3.865140in}{2.189905in}}%
\pgfpathlineto{\pgfqpoint{3.865140in}{2.192854in}}%
\pgfpathlineto{\pgfqpoint{3.869681in}{2.192854in}}%
\pgfpathlineto{\pgfqpoint{3.869681in}{2.189905in}}%
\pgfpathmoveto{\pgfqpoint{3.906009in}{2.130920in}}%
\pgfpathlineto{\pgfqpoint{3.906009in}{2.130920in}}%
\pgfpathlineto{\pgfqpoint{3.906009in}{2.133870in}}%
\pgfpathlineto{\pgfqpoint{3.910550in}{2.133870in}}%
\pgfpathlineto{\pgfqpoint{3.910550in}{2.130920in}}%
\pgfpathmoveto{\pgfqpoint{3.901468in}{2.136819in}}%
\pgfpathlineto{\pgfqpoint{3.901468in}{2.136819in}}%
\pgfpathlineto{\pgfqpoint{3.901468in}{2.139768in}}%
\pgfpathlineto{\pgfqpoint{3.906009in}{2.139768in}}%
\pgfpathlineto{\pgfqpoint{3.906009in}{2.136819in}}%
\pgfpathmoveto{\pgfqpoint{3.906009in}{2.133870in}}%
\pgfpathlineto{\pgfqpoint{3.906009in}{2.133870in}}%
\pgfpathlineto{\pgfqpoint{3.906009in}{2.136819in}}%
\pgfpathlineto{\pgfqpoint{3.910550in}{2.136819in}}%
\pgfpathlineto{\pgfqpoint{3.910550in}{2.133870in}}%
\pgfpathmoveto{\pgfqpoint{3.906009in}{2.136819in}}%
\pgfpathlineto{\pgfqpoint{3.906009in}{2.136819in}}%
\pgfpathlineto{\pgfqpoint{3.906009in}{2.139768in}}%
\pgfpathlineto{\pgfqpoint{3.910550in}{2.139768in}}%
\pgfpathlineto{\pgfqpoint{3.910550in}{2.136819in}}%
\pgfpathmoveto{\pgfqpoint{3.896927in}{2.142717in}}%
\pgfpathlineto{\pgfqpoint{3.896927in}{2.142717in}}%
\pgfpathlineto{\pgfqpoint{3.896927in}{2.145666in}}%
\pgfpathlineto{\pgfqpoint{3.901468in}{2.145666in}}%
\pgfpathlineto{\pgfqpoint{3.901468in}{2.142717in}}%
\pgfpathmoveto{\pgfqpoint{3.892386in}{2.148616in}}%
\pgfpathlineto{\pgfqpoint{3.892386in}{2.148616in}}%
\pgfpathlineto{\pgfqpoint{3.892386in}{2.151565in}}%
\pgfpathlineto{\pgfqpoint{3.896927in}{2.151565in}}%
\pgfpathlineto{\pgfqpoint{3.896927in}{2.148616in}}%
\pgfpathmoveto{\pgfqpoint{3.896927in}{2.145666in}}%
\pgfpathlineto{\pgfqpoint{3.896927in}{2.145666in}}%
\pgfpathlineto{\pgfqpoint{3.896927in}{2.148616in}}%
\pgfpathlineto{\pgfqpoint{3.901468in}{2.148616in}}%
\pgfpathlineto{\pgfqpoint{3.901468in}{2.145666in}}%
\pgfpathmoveto{\pgfqpoint{3.896927in}{2.148616in}}%
\pgfpathlineto{\pgfqpoint{3.896927in}{2.148616in}}%
\pgfpathlineto{\pgfqpoint{3.896927in}{2.151565in}}%
\pgfpathlineto{\pgfqpoint{3.901468in}{2.151565in}}%
\pgfpathlineto{\pgfqpoint{3.901468in}{2.148616in}}%
\pgfpathmoveto{\pgfqpoint{3.901468in}{2.139768in}}%
\pgfpathlineto{\pgfqpoint{3.901468in}{2.139768in}}%
\pgfpathlineto{\pgfqpoint{3.901468in}{2.142717in}}%
\pgfpathlineto{\pgfqpoint{3.906009in}{2.142717in}}%
\pgfpathlineto{\pgfqpoint{3.906009in}{2.139768in}}%
\pgfpathmoveto{\pgfqpoint{3.901468in}{2.142717in}}%
\pgfpathlineto{\pgfqpoint{3.901468in}{2.142717in}}%
\pgfpathlineto{\pgfqpoint{3.901468in}{2.145666in}}%
\pgfpathlineto{\pgfqpoint{3.906009in}{2.145666in}}%
\pgfpathlineto{\pgfqpoint{3.906009in}{2.142717in}}%
\pgfpathmoveto{\pgfqpoint{3.924173in}{2.107327in}}%
\pgfpathlineto{\pgfqpoint{3.924173in}{2.107327in}}%
\pgfpathlineto{\pgfqpoint{3.924173in}{2.110276in}}%
\pgfpathlineto{\pgfqpoint{3.928714in}{2.110276in}}%
\pgfpathlineto{\pgfqpoint{3.928714in}{2.107327in}}%
\pgfpathmoveto{\pgfqpoint{3.919632in}{2.113225in}}%
\pgfpathlineto{\pgfqpoint{3.919632in}{2.113225in}}%
\pgfpathlineto{\pgfqpoint{3.919632in}{2.116174in}}%
\pgfpathlineto{\pgfqpoint{3.924173in}{2.116174in}}%
\pgfpathlineto{\pgfqpoint{3.924173in}{2.113225in}}%
\pgfpathmoveto{\pgfqpoint{3.924173in}{2.110276in}}%
\pgfpathlineto{\pgfqpoint{3.924173in}{2.110276in}}%
\pgfpathlineto{\pgfqpoint{3.924173in}{2.113225in}}%
\pgfpathlineto{\pgfqpoint{3.928714in}{2.113225in}}%
\pgfpathlineto{\pgfqpoint{3.928714in}{2.110276in}}%
\pgfpathmoveto{\pgfqpoint{3.924173in}{2.113225in}}%
\pgfpathlineto{\pgfqpoint{3.924173in}{2.113225in}}%
\pgfpathlineto{\pgfqpoint{3.924173in}{2.116174in}}%
\pgfpathlineto{\pgfqpoint{3.928714in}{2.116174in}}%
\pgfpathlineto{\pgfqpoint{3.928714in}{2.113225in}}%
\pgfpathmoveto{\pgfqpoint{3.915091in}{2.119123in}}%
\pgfpathlineto{\pgfqpoint{3.915091in}{2.119123in}}%
\pgfpathlineto{\pgfqpoint{3.915091in}{2.122073in}}%
\pgfpathlineto{\pgfqpoint{3.919632in}{2.122073in}}%
\pgfpathlineto{\pgfqpoint{3.919632in}{2.119123in}}%
\pgfpathmoveto{\pgfqpoint{3.910550in}{2.125022in}}%
\pgfpathlineto{\pgfqpoint{3.910550in}{2.125022in}}%
\pgfpathlineto{\pgfqpoint{3.910550in}{2.127971in}}%
\pgfpathlineto{\pgfqpoint{3.915091in}{2.127971in}}%
\pgfpathlineto{\pgfqpoint{3.915091in}{2.125022in}}%
\pgfpathmoveto{\pgfqpoint{3.915091in}{2.122073in}}%
\pgfpathlineto{\pgfqpoint{3.915091in}{2.122073in}}%
\pgfpathlineto{\pgfqpoint{3.915091in}{2.125022in}}%
\pgfpathlineto{\pgfqpoint{3.919632in}{2.125022in}}%
\pgfpathlineto{\pgfqpoint{3.919632in}{2.122073in}}%
\pgfpathmoveto{\pgfqpoint{3.915091in}{2.125022in}}%
\pgfpathlineto{\pgfqpoint{3.915091in}{2.125022in}}%
\pgfpathlineto{\pgfqpoint{3.915091in}{2.127971in}}%
\pgfpathlineto{\pgfqpoint{3.919632in}{2.127971in}}%
\pgfpathlineto{\pgfqpoint{3.919632in}{2.125022in}}%
\pgfpathmoveto{\pgfqpoint{3.919632in}{2.116174in}}%
\pgfpathlineto{\pgfqpoint{3.919632in}{2.116174in}}%
\pgfpathlineto{\pgfqpoint{3.919632in}{2.119123in}}%
\pgfpathlineto{\pgfqpoint{3.924173in}{2.119123in}}%
\pgfpathlineto{\pgfqpoint{3.924173in}{2.116174in}}%
\pgfpathmoveto{\pgfqpoint{3.919632in}{2.119123in}}%
\pgfpathlineto{\pgfqpoint{3.919632in}{2.119123in}}%
\pgfpathlineto{\pgfqpoint{3.919632in}{2.122073in}}%
\pgfpathlineto{\pgfqpoint{3.924173in}{2.122073in}}%
\pgfpathlineto{\pgfqpoint{3.924173in}{2.119123in}}%
\pgfpathmoveto{\pgfqpoint{3.928714in}{2.104377in}}%
\pgfpathlineto{\pgfqpoint{3.928714in}{2.104377in}}%
\pgfpathlineto{\pgfqpoint{3.928714in}{2.107327in}}%
\pgfpathlineto{\pgfqpoint{3.933255in}{2.107327in}}%
\pgfpathlineto{\pgfqpoint{3.933255in}{2.104377in}}%
\pgfpathmoveto{\pgfqpoint{3.928714in}{2.107327in}}%
\pgfpathlineto{\pgfqpoint{3.928714in}{2.107327in}}%
\pgfpathlineto{\pgfqpoint{3.928714in}{2.110276in}}%
\pgfpathlineto{\pgfqpoint{3.933255in}{2.110276in}}%
\pgfpathlineto{\pgfqpoint{3.933255in}{2.107327in}}%
\pgfpathmoveto{\pgfqpoint{3.910550in}{2.127971in}}%
\pgfpathlineto{\pgfqpoint{3.910550in}{2.127971in}}%
\pgfpathlineto{\pgfqpoint{3.910550in}{2.130920in}}%
\pgfpathlineto{\pgfqpoint{3.915091in}{2.130920in}}%
\pgfpathlineto{\pgfqpoint{3.915091in}{2.127971in}}%
\pgfpathmoveto{\pgfqpoint{3.910550in}{2.130920in}}%
\pgfpathlineto{\pgfqpoint{3.910550in}{2.130920in}}%
\pgfpathlineto{\pgfqpoint{3.910550in}{2.133870in}}%
\pgfpathlineto{\pgfqpoint{3.915091in}{2.133870in}}%
\pgfpathlineto{\pgfqpoint{3.915091in}{2.130920in}}%
\pgfpathmoveto{\pgfqpoint{3.887845in}{2.154514in}}%
\pgfpathlineto{\pgfqpoint{3.887845in}{2.154514in}}%
\pgfpathlineto{\pgfqpoint{3.887845in}{2.157463in}}%
\pgfpathlineto{\pgfqpoint{3.892386in}{2.157463in}}%
\pgfpathlineto{\pgfqpoint{3.892386in}{2.154514in}}%
\pgfpathmoveto{\pgfqpoint{3.883304in}{2.160413in}}%
\pgfpathlineto{\pgfqpoint{3.883304in}{2.160413in}}%
\pgfpathlineto{\pgfqpoint{3.883304in}{2.163362in}}%
\pgfpathlineto{\pgfqpoint{3.887845in}{2.163362in}}%
\pgfpathlineto{\pgfqpoint{3.887845in}{2.160413in}}%
\pgfpathmoveto{\pgfqpoint{3.887845in}{2.157463in}}%
\pgfpathlineto{\pgfqpoint{3.887845in}{2.157463in}}%
\pgfpathlineto{\pgfqpoint{3.887845in}{2.160413in}}%
\pgfpathlineto{\pgfqpoint{3.892386in}{2.160413in}}%
\pgfpathlineto{\pgfqpoint{3.892386in}{2.157463in}}%
\pgfpathmoveto{\pgfqpoint{3.887845in}{2.160413in}}%
\pgfpathlineto{\pgfqpoint{3.887845in}{2.160413in}}%
\pgfpathlineto{\pgfqpoint{3.887845in}{2.163362in}}%
\pgfpathlineto{\pgfqpoint{3.892386in}{2.163362in}}%
\pgfpathlineto{\pgfqpoint{3.892386in}{2.160413in}}%
\pgfpathmoveto{\pgfqpoint{3.878763in}{2.166311in}}%
\pgfpathlineto{\pgfqpoint{3.878763in}{2.166311in}}%
\pgfpathlineto{\pgfqpoint{3.878763in}{2.169260in}}%
\pgfpathlineto{\pgfqpoint{3.883304in}{2.169260in}}%
\pgfpathlineto{\pgfqpoint{3.883304in}{2.166311in}}%
\pgfpathmoveto{\pgfqpoint{3.874222in}{2.172209in}}%
\pgfpathlineto{\pgfqpoint{3.874222in}{2.172209in}}%
\pgfpathlineto{\pgfqpoint{3.874222in}{2.175159in}}%
\pgfpathlineto{\pgfqpoint{3.878763in}{2.175159in}}%
\pgfpathlineto{\pgfqpoint{3.878763in}{2.172209in}}%
\pgfpathmoveto{\pgfqpoint{3.878763in}{2.169260in}}%
\pgfpathlineto{\pgfqpoint{3.878763in}{2.169260in}}%
\pgfpathlineto{\pgfqpoint{3.878763in}{2.172209in}}%
\pgfpathlineto{\pgfqpoint{3.883304in}{2.172209in}}%
\pgfpathlineto{\pgfqpoint{3.883304in}{2.169260in}}%
\pgfpathmoveto{\pgfqpoint{3.878763in}{2.172209in}}%
\pgfpathlineto{\pgfqpoint{3.878763in}{2.172209in}}%
\pgfpathlineto{\pgfqpoint{3.878763in}{2.175159in}}%
\pgfpathlineto{\pgfqpoint{3.883304in}{2.175159in}}%
\pgfpathlineto{\pgfqpoint{3.883304in}{2.172209in}}%
\pgfpathmoveto{\pgfqpoint{3.883304in}{2.163362in}}%
\pgfpathlineto{\pgfqpoint{3.883304in}{2.163362in}}%
\pgfpathlineto{\pgfqpoint{3.883304in}{2.166311in}}%
\pgfpathlineto{\pgfqpoint{3.887845in}{2.166311in}}%
\pgfpathlineto{\pgfqpoint{3.887845in}{2.163362in}}%
\pgfpathmoveto{\pgfqpoint{3.883304in}{2.166311in}}%
\pgfpathlineto{\pgfqpoint{3.883304in}{2.166311in}}%
\pgfpathlineto{\pgfqpoint{3.883304in}{2.169260in}}%
\pgfpathlineto{\pgfqpoint{3.887845in}{2.169260in}}%
\pgfpathlineto{\pgfqpoint{3.887845in}{2.166311in}}%
\pgfpathmoveto{\pgfqpoint{3.892386in}{2.151565in}}%
\pgfpathlineto{\pgfqpoint{3.892386in}{2.151565in}}%
\pgfpathlineto{\pgfqpoint{3.892386in}{2.154514in}}%
\pgfpathlineto{\pgfqpoint{3.896927in}{2.154514in}}%
\pgfpathlineto{\pgfqpoint{3.896927in}{2.151565in}}%
\pgfpathmoveto{\pgfqpoint{3.892386in}{2.154514in}}%
\pgfpathlineto{\pgfqpoint{3.892386in}{2.154514in}}%
\pgfpathlineto{\pgfqpoint{3.892386in}{2.157463in}}%
\pgfpathlineto{\pgfqpoint{3.896927in}{2.157463in}}%
\pgfpathlineto{\pgfqpoint{3.896927in}{2.154514in}}%
\pgfpathmoveto{\pgfqpoint{3.874222in}{2.175159in}}%
\pgfpathlineto{\pgfqpoint{3.874222in}{2.175159in}}%
\pgfpathlineto{\pgfqpoint{3.874222in}{2.178108in}}%
\pgfpathlineto{\pgfqpoint{3.878763in}{2.178108in}}%
\pgfpathlineto{\pgfqpoint{3.878763in}{2.175159in}}%
\pgfpathmoveto{\pgfqpoint{3.874222in}{2.178108in}}%
\pgfpathlineto{\pgfqpoint{3.874222in}{2.178108in}}%
\pgfpathlineto{\pgfqpoint{3.874222in}{2.181057in}}%
\pgfpathlineto{\pgfqpoint{3.878763in}{2.181057in}}%
\pgfpathlineto{\pgfqpoint{3.878763in}{2.178108in}}%
\pgfpathmoveto{\pgfqpoint{3.833353in}{2.225294in}}%
\pgfpathlineto{\pgfqpoint{3.833353in}{2.225294in}}%
\pgfpathlineto{\pgfqpoint{3.833353in}{2.228243in}}%
\pgfpathlineto{\pgfqpoint{3.837894in}{2.228243in}}%
\pgfpathlineto{\pgfqpoint{3.837894in}{2.225294in}}%
\pgfpathmoveto{\pgfqpoint{3.828812in}{2.231192in}}%
\pgfpathlineto{\pgfqpoint{3.828812in}{2.231192in}}%
\pgfpathlineto{\pgfqpoint{3.828812in}{2.234141in}}%
\pgfpathlineto{\pgfqpoint{3.833353in}{2.234141in}}%
\pgfpathlineto{\pgfqpoint{3.833353in}{2.231192in}}%
\pgfpathmoveto{\pgfqpoint{3.833353in}{2.228243in}}%
\pgfpathlineto{\pgfqpoint{3.833353in}{2.228243in}}%
\pgfpathlineto{\pgfqpoint{3.833353in}{2.231192in}}%
\pgfpathlineto{\pgfqpoint{3.837894in}{2.231192in}}%
\pgfpathlineto{\pgfqpoint{3.837894in}{2.228243in}}%
\pgfpathmoveto{\pgfqpoint{3.833353in}{2.231192in}}%
\pgfpathlineto{\pgfqpoint{3.833353in}{2.231192in}}%
\pgfpathlineto{\pgfqpoint{3.833353in}{2.234141in}}%
\pgfpathlineto{\pgfqpoint{3.837894in}{2.234141in}}%
\pgfpathlineto{\pgfqpoint{3.837894in}{2.231192in}}%
\pgfpathmoveto{\pgfqpoint{3.824271in}{2.237091in}}%
\pgfpathlineto{\pgfqpoint{3.824271in}{2.237091in}}%
\pgfpathlineto{\pgfqpoint{3.824271in}{2.240040in}}%
\pgfpathlineto{\pgfqpoint{3.828812in}{2.240040in}}%
\pgfpathlineto{\pgfqpoint{3.828812in}{2.237091in}}%
\pgfpathmoveto{\pgfqpoint{3.819730in}{2.242989in}}%
\pgfpathlineto{\pgfqpoint{3.819730in}{2.242989in}}%
\pgfpathlineto{\pgfqpoint{3.819730in}{2.245938in}}%
\pgfpathlineto{\pgfqpoint{3.824271in}{2.245938in}}%
\pgfpathlineto{\pgfqpoint{3.824271in}{2.242989in}}%
\pgfpathmoveto{\pgfqpoint{3.824271in}{2.240040in}}%
\pgfpathlineto{\pgfqpoint{3.824271in}{2.240040in}}%
\pgfpathlineto{\pgfqpoint{3.824271in}{2.242989in}}%
\pgfpathlineto{\pgfqpoint{3.828812in}{2.242989in}}%
\pgfpathlineto{\pgfqpoint{3.828812in}{2.240040in}}%
\pgfpathmoveto{\pgfqpoint{3.824271in}{2.242989in}}%
\pgfpathlineto{\pgfqpoint{3.824271in}{2.242989in}}%
\pgfpathlineto{\pgfqpoint{3.824271in}{2.245938in}}%
\pgfpathlineto{\pgfqpoint{3.828812in}{2.245938in}}%
\pgfpathlineto{\pgfqpoint{3.828812in}{2.242989in}}%
\pgfpathmoveto{\pgfqpoint{3.828812in}{2.234141in}}%
\pgfpathlineto{\pgfqpoint{3.828812in}{2.234141in}}%
\pgfpathlineto{\pgfqpoint{3.828812in}{2.237091in}}%
\pgfpathlineto{\pgfqpoint{3.833353in}{2.237091in}}%
\pgfpathlineto{\pgfqpoint{3.833353in}{2.234141in}}%
\pgfpathmoveto{\pgfqpoint{3.828812in}{2.237091in}}%
\pgfpathlineto{\pgfqpoint{3.828812in}{2.237091in}}%
\pgfpathlineto{\pgfqpoint{3.828812in}{2.240040in}}%
\pgfpathlineto{\pgfqpoint{3.833353in}{2.240040in}}%
\pgfpathlineto{\pgfqpoint{3.833353in}{2.237091in}}%
\pgfpathmoveto{\pgfqpoint{3.851517in}{2.201701in}}%
\pgfpathlineto{\pgfqpoint{3.851517in}{2.201701in}}%
\pgfpathlineto{\pgfqpoint{3.851517in}{2.204651in}}%
\pgfpathlineto{\pgfqpoint{3.856058in}{2.204651in}}%
\pgfpathlineto{\pgfqpoint{3.856058in}{2.201701in}}%
\pgfpathmoveto{\pgfqpoint{3.846976in}{2.207600in}}%
\pgfpathlineto{\pgfqpoint{3.846976in}{2.207600in}}%
\pgfpathlineto{\pgfqpoint{3.846976in}{2.210549in}}%
\pgfpathlineto{\pgfqpoint{3.851517in}{2.210549in}}%
\pgfpathlineto{\pgfqpoint{3.851517in}{2.207600in}}%
\pgfpathmoveto{\pgfqpoint{3.851517in}{2.204651in}}%
\pgfpathlineto{\pgfqpoint{3.851517in}{2.204651in}}%
\pgfpathlineto{\pgfqpoint{3.851517in}{2.207600in}}%
\pgfpathlineto{\pgfqpoint{3.856058in}{2.207600in}}%
\pgfpathlineto{\pgfqpoint{3.856058in}{2.204651in}}%
\pgfpathmoveto{\pgfqpoint{3.851517in}{2.207600in}}%
\pgfpathlineto{\pgfqpoint{3.851517in}{2.207600in}}%
\pgfpathlineto{\pgfqpoint{3.851517in}{2.210549in}}%
\pgfpathlineto{\pgfqpoint{3.856058in}{2.210549in}}%
\pgfpathlineto{\pgfqpoint{3.856058in}{2.207600in}}%
\pgfpathmoveto{\pgfqpoint{3.842435in}{2.213498in}}%
\pgfpathlineto{\pgfqpoint{3.842435in}{2.213498in}}%
\pgfpathlineto{\pgfqpoint{3.842435in}{2.216447in}}%
\pgfpathlineto{\pgfqpoint{3.846976in}{2.216447in}}%
\pgfpathlineto{\pgfqpoint{3.846976in}{2.213498in}}%
\pgfpathmoveto{\pgfqpoint{3.837894in}{2.219396in}}%
\pgfpathlineto{\pgfqpoint{3.837894in}{2.219396in}}%
\pgfpathlineto{\pgfqpoint{3.837894in}{2.222345in}}%
\pgfpathlineto{\pgfqpoint{3.842435in}{2.222345in}}%
\pgfpathlineto{\pgfqpoint{3.842435in}{2.219396in}}%
\pgfpathmoveto{\pgfqpoint{3.842435in}{2.216447in}}%
\pgfpathlineto{\pgfqpoint{3.842435in}{2.216447in}}%
\pgfpathlineto{\pgfqpoint{3.842435in}{2.219396in}}%
\pgfpathlineto{\pgfqpoint{3.846976in}{2.219396in}}%
\pgfpathlineto{\pgfqpoint{3.846976in}{2.216447in}}%
\pgfpathmoveto{\pgfqpoint{3.842435in}{2.219396in}}%
\pgfpathlineto{\pgfqpoint{3.842435in}{2.219396in}}%
\pgfpathlineto{\pgfqpoint{3.842435in}{2.222345in}}%
\pgfpathlineto{\pgfqpoint{3.846976in}{2.222345in}}%
\pgfpathlineto{\pgfqpoint{3.846976in}{2.219396in}}%
\pgfpathmoveto{\pgfqpoint{3.846976in}{2.210549in}}%
\pgfpathlineto{\pgfqpoint{3.846976in}{2.210549in}}%
\pgfpathlineto{\pgfqpoint{3.846976in}{2.213498in}}%
\pgfpathlineto{\pgfqpoint{3.851517in}{2.213498in}}%
\pgfpathlineto{\pgfqpoint{3.851517in}{2.210549in}}%
\pgfpathmoveto{\pgfqpoint{3.846976in}{2.213498in}}%
\pgfpathlineto{\pgfqpoint{3.846976in}{2.213498in}}%
\pgfpathlineto{\pgfqpoint{3.846976in}{2.216447in}}%
\pgfpathlineto{\pgfqpoint{3.851517in}{2.216447in}}%
\pgfpathlineto{\pgfqpoint{3.851517in}{2.213498in}}%
\pgfpathmoveto{\pgfqpoint{3.856058in}{2.198752in}}%
\pgfpathlineto{\pgfqpoint{3.856058in}{2.198752in}}%
\pgfpathlineto{\pgfqpoint{3.856058in}{2.201701in}}%
\pgfpathlineto{\pgfqpoint{3.860599in}{2.201701in}}%
\pgfpathlineto{\pgfqpoint{3.860599in}{2.198752in}}%
\pgfpathmoveto{\pgfqpoint{3.856058in}{2.201701in}}%
\pgfpathlineto{\pgfqpoint{3.856058in}{2.201701in}}%
\pgfpathlineto{\pgfqpoint{3.856058in}{2.204651in}}%
\pgfpathlineto{\pgfqpoint{3.860599in}{2.204651in}}%
\pgfpathlineto{\pgfqpoint{3.860599in}{2.201701in}}%
\pgfpathmoveto{\pgfqpoint{3.837894in}{2.222345in}}%
\pgfpathlineto{\pgfqpoint{3.837894in}{2.222345in}}%
\pgfpathlineto{\pgfqpoint{3.837894in}{2.225294in}}%
\pgfpathlineto{\pgfqpoint{3.842435in}{2.225294in}}%
\pgfpathlineto{\pgfqpoint{3.842435in}{2.222345in}}%
\pgfpathmoveto{\pgfqpoint{3.837894in}{2.225294in}}%
\pgfpathlineto{\pgfqpoint{3.837894in}{2.225294in}}%
\pgfpathlineto{\pgfqpoint{3.837894in}{2.228243in}}%
\pgfpathlineto{\pgfqpoint{3.842435in}{2.228243in}}%
\pgfpathlineto{\pgfqpoint{3.842435in}{2.225294in}}%
\pgfpathmoveto{\pgfqpoint{3.815189in}{2.248887in}}%
\pgfpathlineto{\pgfqpoint{3.815189in}{2.248887in}}%
\pgfpathlineto{\pgfqpoint{3.815189in}{2.251836in}}%
\pgfpathlineto{\pgfqpoint{3.819730in}{2.251836in}}%
\pgfpathlineto{\pgfqpoint{3.819730in}{2.248887in}}%
\pgfpathmoveto{\pgfqpoint{3.810648in}{2.254785in}}%
\pgfpathlineto{\pgfqpoint{3.810648in}{2.254785in}}%
\pgfpathlineto{\pgfqpoint{3.810648in}{2.257734in}}%
\pgfpathlineto{\pgfqpoint{3.815189in}{2.257734in}}%
\pgfpathlineto{\pgfqpoint{3.815189in}{2.254785in}}%
\pgfpathmoveto{\pgfqpoint{3.815189in}{2.251836in}}%
\pgfpathlineto{\pgfqpoint{3.815189in}{2.251836in}}%
\pgfpathlineto{\pgfqpoint{3.815189in}{2.254785in}}%
\pgfpathlineto{\pgfqpoint{3.819730in}{2.254785in}}%
\pgfpathlineto{\pgfqpoint{3.819730in}{2.251836in}}%
\pgfpathmoveto{\pgfqpoint{3.815189in}{2.254785in}}%
\pgfpathlineto{\pgfqpoint{3.815189in}{2.254785in}}%
\pgfpathlineto{\pgfqpoint{3.815189in}{2.257734in}}%
\pgfpathlineto{\pgfqpoint{3.819730in}{2.257734in}}%
\pgfpathlineto{\pgfqpoint{3.819730in}{2.254785in}}%
\pgfpathmoveto{\pgfqpoint{3.806107in}{2.260683in}}%
\pgfpathlineto{\pgfqpoint{3.806107in}{2.260683in}}%
\pgfpathlineto{\pgfqpoint{3.806107in}{2.263632in}}%
\pgfpathlineto{\pgfqpoint{3.810648in}{2.263632in}}%
\pgfpathlineto{\pgfqpoint{3.810648in}{2.260683in}}%
\pgfpathmoveto{\pgfqpoint{3.801566in}{2.266581in}}%
\pgfpathlineto{\pgfqpoint{3.801566in}{2.266581in}}%
\pgfpathlineto{\pgfqpoint{3.801566in}{2.269531in}}%
\pgfpathlineto{\pgfqpoint{3.806107in}{2.269531in}}%
\pgfpathlineto{\pgfqpoint{3.806107in}{2.266581in}}%
\pgfpathmoveto{\pgfqpoint{3.806107in}{2.263632in}}%
\pgfpathlineto{\pgfqpoint{3.806107in}{2.263632in}}%
\pgfpathlineto{\pgfqpoint{3.806107in}{2.266581in}}%
\pgfpathlineto{\pgfqpoint{3.810648in}{2.266581in}}%
\pgfpathlineto{\pgfqpoint{3.810648in}{2.263632in}}%
\pgfpathmoveto{\pgfqpoint{3.806107in}{2.266581in}}%
\pgfpathlineto{\pgfqpoint{3.806107in}{2.266581in}}%
\pgfpathlineto{\pgfqpoint{3.806107in}{2.269531in}}%
\pgfpathlineto{\pgfqpoint{3.810648in}{2.269531in}}%
\pgfpathlineto{\pgfqpoint{3.810648in}{2.266581in}}%
\pgfpathmoveto{\pgfqpoint{3.810648in}{2.257734in}}%
\pgfpathlineto{\pgfqpoint{3.810648in}{2.257734in}}%
\pgfpathlineto{\pgfqpoint{3.810648in}{2.260683in}}%
\pgfpathlineto{\pgfqpoint{3.815189in}{2.260683in}}%
\pgfpathlineto{\pgfqpoint{3.815189in}{2.257734in}}%
\pgfpathmoveto{\pgfqpoint{3.810648in}{2.260683in}}%
\pgfpathlineto{\pgfqpoint{3.810648in}{2.260683in}}%
\pgfpathlineto{\pgfqpoint{3.810648in}{2.263632in}}%
\pgfpathlineto{\pgfqpoint{3.815189in}{2.263632in}}%
\pgfpathlineto{\pgfqpoint{3.815189in}{2.260683in}}%
\pgfpathmoveto{\pgfqpoint{3.819730in}{2.245938in}}%
\pgfpathlineto{\pgfqpoint{3.819730in}{2.245938in}}%
\pgfpathlineto{\pgfqpoint{3.819730in}{2.248887in}}%
\pgfpathlineto{\pgfqpoint{3.824271in}{2.248887in}}%
\pgfpathlineto{\pgfqpoint{3.824271in}{2.245938in}}%
\pgfpathmoveto{\pgfqpoint{3.819730in}{2.248887in}}%
\pgfpathlineto{\pgfqpoint{3.819730in}{2.248887in}}%
\pgfpathlineto{\pgfqpoint{3.819730in}{2.251836in}}%
\pgfpathlineto{\pgfqpoint{3.824271in}{2.251836in}}%
\pgfpathlineto{\pgfqpoint{3.824271in}{2.248887in}}%
\pgfpathmoveto{\pgfqpoint{3.801566in}{2.269531in}}%
\pgfpathlineto{\pgfqpoint{3.801566in}{2.269531in}}%
\pgfpathlineto{\pgfqpoint{3.801566in}{2.272480in}}%
\pgfpathlineto{\pgfqpoint{3.806107in}{2.272480in}}%
\pgfpathlineto{\pgfqpoint{3.806107in}{2.269531in}}%
\pgfpathmoveto{\pgfqpoint{3.801566in}{2.272480in}}%
\pgfpathlineto{\pgfqpoint{3.801566in}{2.272480in}}%
\pgfpathlineto{\pgfqpoint{3.801566in}{2.275429in}}%
\pgfpathlineto{\pgfqpoint{3.806107in}{2.275429in}}%
\pgfpathlineto{\pgfqpoint{3.806107in}{2.272480in}}%
\pgfpathmoveto{\pgfqpoint{4.087644in}{1.894980in}}%
\pgfpathlineto{\pgfqpoint{4.087644in}{1.894980in}}%
\pgfpathlineto{\pgfqpoint{4.087644in}{1.897929in}}%
\pgfpathlineto{\pgfqpoint{4.092185in}{1.897929in}}%
\pgfpathlineto{\pgfqpoint{4.092185in}{1.894980in}}%
\pgfpathmoveto{\pgfqpoint{4.083103in}{1.900878in}}%
\pgfpathlineto{\pgfqpoint{4.083103in}{1.900878in}}%
\pgfpathlineto{\pgfqpoint{4.083103in}{1.903827in}}%
\pgfpathlineto{\pgfqpoint{4.087644in}{1.903827in}}%
\pgfpathlineto{\pgfqpoint{4.087644in}{1.900878in}}%
\pgfpathmoveto{\pgfqpoint{4.087644in}{1.897929in}}%
\pgfpathlineto{\pgfqpoint{4.087644in}{1.897929in}}%
\pgfpathlineto{\pgfqpoint{4.087644in}{1.900878in}}%
\pgfpathlineto{\pgfqpoint{4.092185in}{1.900878in}}%
\pgfpathlineto{\pgfqpoint{4.092185in}{1.897929in}}%
\pgfpathmoveto{\pgfqpoint{4.087644in}{1.900878in}}%
\pgfpathlineto{\pgfqpoint{4.087644in}{1.900878in}}%
\pgfpathlineto{\pgfqpoint{4.087644in}{1.903827in}}%
\pgfpathlineto{\pgfqpoint{4.092185in}{1.903827in}}%
\pgfpathlineto{\pgfqpoint{4.092185in}{1.900878in}}%
\pgfpathmoveto{\pgfqpoint{4.078562in}{1.906776in}}%
\pgfpathlineto{\pgfqpoint{4.078562in}{1.906776in}}%
\pgfpathlineto{\pgfqpoint{4.078562in}{1.909725in}}%
\pgfpathlineto{\pgfqpoint{4.083103in}{1.909725in}}%
\pgfpathlineto{\pgfqpoint{4.083103in}{1.906776in}}%
\pgfpathmoveto{\pgfqpoint{4.074022in}{1.912675in}}%
\pgfpathlineto{\pgfqpoint{4.074022in}{1.912675in}}%
\pgfpathlineto{\pgfqpoint{4.074022in}{1.915624in}}%
\pgfpathlineto{\pgfqpoint{4.078562in}{1.915624in}}%
\pgfpathlineto{\pgfqpoint{4.078562in}{1.912675in}}%
\pgfpathmoveto{\pgfqpoint{4.078562in}{1.909725in}}%
\pgfpathlineto{\pgfqpoint{4.078562in}{1.909725in}}%
\pgfpathlineto{\pgfqpoint{4.078562in}{1.912675in}}%
\pgfpathlineto{\pgfqpoint{4.083103in}{1.912675in}}%
\pgfpathlineto{\pgfqpoint{4.083103in}{1.909725in}}%
\pgfpathmoveto{\pgfqpoint{4.078562in}{1.912675in}}%
\pgfpathlineto{\pgfqpoint{4.078562in}{1.912675in}}%
\pgfpathlineto{\pgfqpoint{4.078562in}{1.915624in}}%
\pgfpathlineto{\pgfqpoint{4.083103in}{1.915624in}}%
\pgfpathlineto{\pgfqpoint{4.083103in}{1.912675in}}%
\pgfpathmoveto{\pgfqpoint{4.083103in}{1.903827in}}%
\pgfpathlineto{\pgfqpoint{4.083103in}{1.903827in}}%
\pgfpathlineto{\pgfqpoint{4.083103in}{1.906776in}}%
\pgfpathlineto{\pgfqpoint{4.087644in}{1.906776in}}%
\pgfpathlineto{\pgfqpoint{4.087644in}{1.903827in}}%
\pgfpathmoveto{\pgfqpoint{4.083103in}{1.906776in}}%
\pgfpathlineto{\pgfqpoint{4.083103in}{1.906776in}}%
\pgfpathlineto{\pgfqpoint{4.083103in}{1.909725in}}%
\pgfpathlineto{\pgfqpoint{4.087644in}{1.909725in}}%
\pgfpathlineto{\pgfqpoint{4.087644in}{1.906776in}}%
\pgfpathmoveto{\pgfqpoint{4.014991in}{1.989357in}}%
\pgfpathlineto{\pgfqpoint{4.014991in}{1.989357in}}%
\pgfpathlineto{\pgfqpoint{4.014991in}{1.992306in}}%
\pgfpathlineto{\pgfqpoint{4.019532in}{1.992306in}}%
\pgfpathlineto{\pgfqpoint{4.019532in}{1.989357in}}%
\pgfpathmoveto{\pgfqpoint{4.010450in}{1.995256in}}%
\pgfpathlineto{\pgfqpoint{4.010450in}{1.995256in}}%
\pgfpathlineto{\pgfqpoint{4.010450in}{1.998205in}}%
\pgfpathlineto{\pgfqpoint{4.014991in}{1.998205in}}%
\pgfpathlineto{\pgfqpoint{4.014991in}{1.995256in}}%
\pgfpathmoveto{\pgfqpoint{4.014991in}{1.992306in}}%
\pgfpathlineto{\pgfqpoint{4.014991in}{1.992306in}}%
\pgfpathlineto{\pgfqpoint{4.014991in}{1.995256in}}%
\pgfpathlineto{\pgfqpoint{4.019532in}{1.995256in}}%
\pgfpathlineto{\pgfqpoint{4.019532in}{1.992306in}}%
\pgfpathmoveto{\pgfqpoint{4.014991in}{1.995256in}}%
\pgfpathlineto{\pgfqpoint{4.014991in}{1.995256in}}%
\pgfpathlineto{\pgfqpoint{4.014991in}{1.998205in}}%
\pgfpathlineto{\pgfqpoint{4.019532in}{1.998205in}}%
\pgfpathlineto{\pgfqpoint{4.019532in}{1.995256in}}%
\pgfpathmoveto{\pgfqpoint{4.005909in}{2.001154in}}%
\pgfpathlineto{\pgfqpoint{4.005909in}{2.001154in}}%
\pgfpathlineto{\pgfqpoint{4.005909in}{2.004104in}}%
\pgfpathlineto{\pgfqpoint{4.010450in}{2.004104in}}%
\pgfpathlineto{\pgfqpoint{4.010450in}{2.001154in}}%
\pgfpathmoveto{\pgfqpoint{4.001368in}{2.007053in}}%
\pgfpathlineto{\pgfqpoint{4.001368in}{2.007053in}}%
\pgfpathlineto{\pgfqpoint{4.001368in}{2.010002in}}%
\pgfpathlineto{\pgfqpoint{4.005909in}{2.010002in}}%
\pgfpathlineto{\pgfqpoint{4.005909in}{2.007053in}}%
\pgfpathmoveto{\pgfqpoint{4.005909in}{2.004104in}}%
\pgfpathlineto{\pgfqpoint{4.005909in}{2.004104in}}%
\pgfpathlineto{\pgfqpoint{4.005909in}{2.007053in}}%
\pgfpathlineto{\pgfqpoint{4.010450in}{2.007053in}}%
\pgfpathlineto{\pgfqpoint{4.010450in}{2.004104in}}%
\pgfpathmoveto{\pgfqpoint{4.005909in}{2.007053in}}%
\pgfpathlineto{\pgfqpoint{4.005909in}{2.007053in}}%
\pgfpathlineto{\pgfqpoint{4.005909in}{2.010002in}}%
\pgfpathlineto{\pgfqpoint{4.010450in}{2.010002in}}%
\pgfpathlineto{\pgfqpoint{4.010450in}{2.007053in}}%
\pgfpathmoveto{\pgfqpoint{4.010450in}{1.998205in}}%
\pgfpathlineto{\pgfqpoint{4.010450in}{1.998205in}}%
\pgfpathlineto{\pgfqpoint{4.010450in}{2.001154in}}%
\pgfpathlineto{\pgfqpoint{4.014991in}{2.001154in}}%
\pgfpathlineto{\pgfqpoint{4.014991in}{1.998205in}}%
\pgfpathmoveto{\pgfqpoint{4.010450in}{2.001154in}}%
\pgfpathlineto{\pgfqpoint{4.010450in}{2.001154in}}%
\pgfpathlineto{\pgfqpoint{4.010450in}{2.004104in}}%
\pgfpathlineto{\pgfqpoint{4.014991in}{2.004104in}}%
\pgfpathlineto{\pgfqpoint{4.014991in}{2.001154in}}%
\pgfpathmoveto{\pgfqpoint{4.051317in}{1.942168in}}%
\pgfpathlineto{\pgfqpoint{4.051317in}{1.942168in}}%
\pgfpathlineto{\pgfqpoint{4.051317in}{1.945117in}}%
\pgfpathlineto{\pgfqpoint{4.055858in}{1.945117in}}%
\pgfpathlineto{\pgfqpoint{4.055858in}{1.942168in}}%
\pgfpathmoveto{\pgfqpoint{4.046777in}{1.948066in}}%
\pgfpathlineto{\pgfqpoint{4.046777in}{1.948066in}}%
\pgfpathlineto{\pgfqpoint{4.046777in}{1.951016in}}%
\pgfpathlineto{\pgfqpoint{4.051317in}{1.951016in}}%
\pgfpathlineto{\pgfqpoint{4.051317in}{1.948066in}}%
\pgfpathmoveto{\pgfqpoint{4.051317in}{1.945117in}}%
\pgfpathlineto{\pgfqpoint{4.051317in}{1.945117in}}%
\pgfpathlineto{\pgfqpoint{4.051317in}{1.948066in}}%
\pgfpathlineto{\pgfqpoint{4.055858in}{1.948066in}}%
\pgfpathlineto{\pgfqpoint{4.055858in}{1.945117in}}%
\pgfpathmoveto{\pgfqpoint{4.051317in}{1.948066in}}%
\pgfpathlineto{\pgfqpoint{4.051317in}{1.948066in}}%
\pgfpathlineto{\pgfqpoint{4.051317in}{1.951016in}}%
\pgfpathlineto{\pgfqpoint{4.055858in}{1.951016in}}%
\pgfpathlineto{\pgfqpoint{4.055858in}{1.948066in}}%
\pgfpathmoveto{\pgfqpoint{4.042236in}{1.953965in}}%
\pgfpathlineto{\pgfqpoint{4.042236in}{1.953965in}}%
\pgfpathlineto{\pgfqpoint{4.042236in}{1.956914in}}%
\pgfpathlineto{\pgfqpoint{4.046777in}{1.956914in}}%
\pgfpathlineto{\pgfqpoint{4.046777in}{1.953965in}}%
\pgfpathmoveto{\pgfqpoint{4.037695in}{1.959864in}}%
\pgfpathlineto{\pgfqpoint{4.037695in}{1.959864in}}%
\pgfpathlineto{\pgfqpoint{4.037695in}{1.962813in}}%
\pgfpathlineto{\pgfqpoint{4.042236in}{1.962813in}}%
\pgfpathlineto{\pgfqpoint{4.042236in}{1.959864in}}%
\pgfpathmoveto{\pgfqpoint{4.042236in}{1.956914in}}%
\pgfpathlineto{\pgfqpoint{4.042236in}{1.956914in}}%
\pgfpathlineto{\pgfqpoint{4.042236in}{1.959864in}}%
\pgfpathlineto{\pgfqpoint{4.046777in}{1.959864in}}%
\pgfpathlineto{\pgfqpoint{4.046777in}{1.956914in}}%
\pgfpathmoveto{\pgfqpoint{4.042236in}{1.959864in}}%
\pgfpathlineto{\pgfqpoint{4.042236in}{1.959864in}}%
\pgfpathlineto{\pgfqpoint{4.042236in}{1.962813in}}%
\pgfpathlineto{\pgfqpoint{4.046777in}{1.962813in}}%
\pgfpathlineto{\pgfqpoint{4.046777in}{1.959864in}}%
\pgfpathmoveto{\pgfqpoint{4.046777in}{1.951016in}}%
\pgfpathlineto{\pgfqpoint{4.046777in}{1.951016in}}%
\pgfpathlineto{\pgfqpoint{4.046777in}{1.953965in}}%
\pgfpathlineto{\pgfqpoint{4.051317in}{1.953965in}}%
\pgfpathlineto{\pgfqpoint{4.051317in}{1.951016in}}%
\pgfpathmoveto{\pgfqpoint{4.046777in}{1.953965in}}%
\pgfpathlineto{\pgfqpoint{4.046777in}{1.953965in}}%
\pgfpathlineto{\pgfqpoint{4.046777in}{1.956914in}}%
\pgfpathlineto{\pgfqpoint{4.051317in}{1.956914in}}%
\pgfpathlineto{\pgfqpoint{4.051317in}{1.953965in}}%
\pgfpathmoveto{\pgfqpoint{4.069481in}{1.918573in}}%
\pgfpathlineto{\pgfqpoint{4.069481in}{1.918573in}}%
\pgfpathlineto{\pgfqpoint{4.069481in}{1.921522in}}%
\pgfpathlineto{\pgfqpoint{4.074022in}{1.921522in}}%
\pgfpathlineto{\pgfqpoint{4.074022in}{1.918573in}}%
\pgfpathmoveto{\pgfqpoint{4.064940in}{1.924472in}}%
\pgfpathlineto{\pgfqpoint{4.064940in}{1.924472in}}%
\pgfpathlineto{\pgfqpoint{4.064940in}{1.927421in}}%
\pgfpathlineto{\pgfqpoint{4.069481in}{1.927421in}}%
\pgfpathlineto{\pgfqpoint{4.069481in}{1.924472in}}%
\pgfpathmoveto{\pgfqpoint{4.069481in}{1.921522in}}%
\pgfpathlineto{\pgfqpoint{4.069481in}{1.921522in}}%
\pgfpathlineto{\pgfqpoint{4.069481in}{1.924472in}}%
\pgfpathlineto{\pgfqpoint{4.074022in}{1.924472in}}%
\pgfpathlineto{\pgfqpoint{4.074022in}{1.921522in}}%
\pgfpathmoveto{\pgfqpoint{4.069481in}{1.924472in}}%
\pgfpathlineto{\pgfqpoint{4.069481in}{1.924472in}}%
\pgfpathlineto{\pgfqpoint{4.069481in}{1.927421in}}%
\pgfpathlineto{\pgfqpoint{4.074022in}{1.927421in}}%
\pgfpathlineto{\pgfqpoint{4.074022in}{1.924472in}}%
\pgfpathmoveto{\pgfqpoint{4.060399in}{1.930370in}}%
\pgfpathlineto{\pgfqpoint{4.060399in}{1.930370in}}%
\pgfpathlineto{\pgfqpoint{4.060399in}{1.933320in}}%
\pgfpathlineto{\pgfqpoint{4.064940in}{1.933320in}}%
\pgfpathlineto{\pgfqpoint{4.064940in}{1.930370in}}%
\pgfpathmoveto{\pgfqpoint{4.055858in}{1.936269in}}%
\pgfpathlineto{\pgfqpoint{4.055858in}{1.936269in}}%
\pgfpathlineto{\pgfqpoint{4.055858in}{1.939218in}}%
\pgfpathlineto{\pgfqpoint{4.060399in}{1.939218in}}%
\pgfpathlineto{\pgfqpoint{4.060399in}{1.936269in}}%
\pgfpathmoveto{\pgfqpoint{4.060399in}{1.933320in}}%
\pgfpathlineto{\pgfqpoint{4.060399in}{1.933320in}}%
\pgfpathlineto{\pgfqpoint{4.060399in}{1.936269in}}%
\pgfpathlineto{\pgfqpoint{4.064940in}{1.936269in}}%
\pgfpathlineto{\pgfqpoint{4.064940in}{1.933320in}}%
\pgfpathmoveto{\pgfqpoint{4.060399in}{1.936269in}}%
\pgfpathlineto{\pgfqpoint{4.060399in}{1.936269in}}%
\pgfpathlineto{\pgfqpoint{4.060399in}{1.939218in}}%
\pgfpathlineto{\pgfqpoint{4.064940in}{1.939218in}}%
\pgfpathlineto{\pgfqpoint{4.064940in}{1.936269in}}%
\pgfpathmoveto{\pgfqpoint{4.064940in}{1.927421in}}%
\pgfpathlineto{\pgfqpoint{4.064940in}{1.927421in}}%
\pgfpathlineto{\pgfqpoint{4.064940in}{1.930370in}}%
\pgfpathlineto{\pgfqpoint{4.069481in}{1.930370in}}%
\pgfpathlineto{\pgfqpoint{4.069481in}{1.927421in}}%
\pgfpathmoveto{\pgfqpoint{4.064940in}{1.930370in}}%
\pgfpathlineto{\pgfqpoint{4.064940in}{1.930370in}}%
\pgfpathlineto{\pgfqpoint{4.064940in}{1.933320in}}%
\pgfpathlineto{\pgfqpoint{4.069481in}{1.933320in}}%
\pgfpathlineto{\pgfqpoint{4.069481in}{1.930370in}}%
\pgfpathmoveto{\pgfqpoint{4.074022in}{1.915624in}}%
\pgfpathlineto{\pgfqpoint{4.074022in}{1.915624in}}%
\pgfpathlineto{\pgfqpoint{4.074022in}{1.918573in}}%
\pgfpathlineto{\pgfqpoint{4.078562in}{1.918573in}}%
\pgfpathlineto{\pgfqpoint{4.078562in}{1.915624in}}%
\pgfpathmoveto{\pgfqpoint{4.074022in}{1.918573in}}%
\pgfpathlineto{\pgfqpoint{4.074022in}{1.918573in}}%
\pgfpathlineto{\pgfqpoint{4.074022in}{1.921522in}}%
\pgfpathlineto{\pgfqpoint{4.078562in}{1.921522in}}%
\pgfpathlineto{\pgfqpoint{4.078562in}{1.918573in}}%
\pgfpathmoveto{\pgfqpoint{4.055858in}{1.939218in}}%
\pgfpathlineto{\pgfqpoint{4.055858in}{1.939218in}}%
\pgfpathlineto{\pgfqpoint{4.055858in}{1.942168in}}%
\pgfpathlineto{\pgfqpoint{4.060399in}{1.942168in}}%
\pgfpathlineto{\pgfqpoint{4.060399in}{1.939218in}}%
\pgfpathmoveto{\pgfqpoint{4.055858in}{1.942168in}}%
\pgfpathlineto{\pgfqpoint{4.055858in}{1.942168in}}%
\pgfpathlineto{\pgfqpoint{4.055858in}{1.945117in}}%
\pgfpathlineto{\pgfqpoint{4.060399in}{1.945117in}}%
\pgfpathlineto{\pgfqpoint{4.060399in}{1.942168in}}%
\pgfpathmoveto{\pgfqpoint{4.033154in}{1.965762in}}%
\pgfpathlineto{\pgfqpoint{4.033154in}{1.965762in}}%
\pgfpathlineto{\pgfqpoint{4.033154in}{1.968712in}}%
\pgfpathlineto{\pgfqpoint{4.037695in}{1.968712in}}%
\pgfpathlineto{\pgfqpoint{4.037695in}{1.965762in}}%
\pgfpathmoveto{\pgfqpoint{4.028613in}{1.971661in}}%
\pgfpathlineto{\pgfqpoint{4.028613in}{1.971661in}}%
\pgfpathlineto{\pgfqpoint{4.028613in}{1.974610in}}%
\pgfpathlineto{\pgfqpoint{4.033154in}{1.974610in}}%
\pgfpathlineto{\pgfqpoint{4.033154in}{1.971661in}}%
\pgfpathmoveto{\pgfqpoint{4.033154in}{1.968712in}}%
\pgfpathlineto{\pgfqpoint{4.033154in}{1.968712in}}%
\pgfpathlineto{\pgfqpoint{4.033154in}{1.971661in}}%
\pgfpathlineto{\pgfqpoint{4.037695in}{1.971661in}}%
\pgfpathlineto{\pgfqpoint{4.037695in}{1.968712in}}%
\pgfpathmoveto{\pgfqpoint{4.033154in}{1.971661in}}%
\pgfpathlineto{\pgfqpoint{4.033154in}{1.971661in}}%
\pgfpathlineto{\pgfqpoint{4.033154in}{1.974610in}}%
\pgfpathlineto{\pgfqpoint{4.037695in}{1.974610in}}%
\pgfpathlineto{\pgfqpoint{4.037695in}{1.971661in}}%
\pgfpathmoveto{\pgfqpoint{4.024072in}{1.977560in}}%
\pgfpathlineto{\pgfqpoint{4.024072in}{1.977560in}}%
\pgfpathlineto{\pgfqpoint{4.024072in}{1.980509in}}%
\pgfpathlineto{\pgfqpoint{4.028613in}{1.980509in}}%
\pgfpathlineto{\pgfqpoint{4.028613in}{1.977560in}}%
\pgfpathmoveto{\pgfqpoint{4.019532in}{1.983458in}}%
\pgfpathlineto{\pgfqpoint{4.019532in}{1.983458in}}%
\pgfpathlineto{\pgfqpoint{4.019532in}{1.986408in}}%
\pgfpathlineto{\pgfqpoint{4.024072in}{1.986408in}}%
\pgfpathlineto{\pgfqpoint{4.024072in}{1.983458in}}%
\pgfpathmoveto{\pgfqpoint{4.024072in}{1.980509in}}%
\pgfpathlineto{\pgfqpoint{4.024072in}{1.980509in}}%
\pgfpathlineto{\pgfqpoint{4.024072in}{1.983458in}}%
\pgfpathlineto{\pgfqpoint{4.028613in}{1.983458in}}%
\pgfpathlineto{\pgfqpoint{4.028613in}{1.980509in}}%
\pgfpathmoveto{\pgfqpoint{4.024072in}{1.983458in}}%
\pgfpathlineto{\pgfqpoint{4.024072in}{1.983458in}}%
\pgfpathlineto{\pgfqpoint{4.024072in}{1.986408in}}%
\pgfpathlineto{\pgfqpoint{4.028613in}{1.986408in}}%
\pgfpathlineto{\pgfqpoint{4.028613in}{1.983458in}}%
\pgfpathmoveto{\pgfqpoint{4.028613in}{1.974610in}}%
\pgfpathlineto{\pgfqpoint{4.028613in}{1.974610in}}%
\pgfpathlineto{\pgfqpoint{4.028613in}{1.977560in}}%
\pgfpathlineto{\pgfqpoint{4.033154in}{1.977560in}}%
\pgfpathlineto{\pgfqpoint{4.033154in}{1.974610in}}%
\pgfpathmoveto{\pgfqpoint{4.028613in}{1.977560in}}%
\pgfpathlineto{\pgfqpoint{4.028613in}{1.977560in}}%
\pgfpathlineto{\pgfqpoint{4.028613in}{1.980509in}}%
\pgfpathlineto{\pgfqpoint{4.033154in}{1.980509in}}%
\pgfpathlineto{\pgfqpoint{4.033154in}{1.977560in}}%
\pgfpathmoveto{\pgfqpoint{4.037695in}{1.962813in}}%
\pgfpathlineto{\pgfqpoint{4.037695in}{1.962813in}}%
\pgfpathlineto{\pgfqpoint{4.037695in}{1.965762in}}%
\pgfpathlineto{\pgfqpoint{4.042236in}{1.965762in}}%
\pgfpathlineto{\pgfqpoint{4.042236in}{1.962813in}}%
\pgfpathmoveto{\pgfqpoint{4.037695in}{1.965762in}}%
\pgfpathlineto{\pgfqpoint{4.037695in}{1.965762in}}%
\pgfpathlineto{\pgfqpoint{4.037695in}{1.968712in}}%
\pgfpathlineto{\pgfqpoint{4.042236in}{1.968712in}}%
\pgfpathlineto{\pgfqpoint{4.042236in}{1.965762in}}%
\pgfpathmoveto{\pgfqpoint{4.019532in}{1.986408in}}%
\pgfpathlineto{\pgfqpoint{4.019532in}{1.986408in}}%
\pgfpathlineto{\pgfqpoint{4.019532in}{1.989357in}}%
\pgfpathlineto{\pgfqpoint{4.024072in}{1.989357in}}%
\pgfpathlineto{\pgfqpoint{4.024072in}{1.986408in}}%
\pgfpathmoveto{\pgfqpoint{4.019532in}{1.989357in}}%
\pgfpathlineto{\pgfqpoint{4.019532in}{1.989357in}}%
\pgfpathlineto{\pgfqpoint{4.019532in}{1.992306in}}%
\pgfpathlineto{\pgfqpoint{4.024072in}{1.992306in}}%
\pgfpathlineto{\pgfqpoint{4.024072in}{1.989357in}}%
\pgfpathmoveto{\pgfqpoint{3.978664in}{2.036545in}}%
\pgfpathlineto{\pgfqpoint{3.978664in}{2.036545in}}%
\pgfpathlineto{\pgfqpoint{3.978664in}{2.039495in}}%
\pgfpathlineto{\pgfqpoint{3.983205in}{2.039495in}}%
\pgfpathlineto{\pgfqpoint{3.983205in}{2.036545in}}%
\pgfpathmoveto{\pgfqpoint{3.974123in}{2.042444in}}%
\pgfpathlineto{\pgfqpoint{3.974123in}{2.042444in}}%
\pgfpathlineto{\pgfqpoint{3.974123in}{2.045393in}}%
\pgfpathlineto{\pgfqpoint{3.978664in}{2.045393in}}%
\pgfpathlineto{\pgfqpoint{3.978664in}{2.042444in}}%
\pgfpathmoveto{\pgfqpoint{3.978664in}{2.039495in}}%
\pgfpathlineto{\pgfqpoint{3.978664in}{2.039495in}}%
\pgfpathlineto{\pgfqpoint{3.978664in}{2.042444in}}%
\pgfpathlineto{\pgfqpoint{3.983205in}{2.042444in}}%
\pgfpathlineto{\pgfqpoint{3.983205in}{2.039495in}}%
\pgfpathmoveto{\pgfqpoint{3.978664in}{2.042444in}}%
\pgfpathlineto{\pgfqpoint{3.978664in}{2.042444in}}%
\pgfpathlineto{\pgfqpoint{3.978664in}{2.045393in}}%
\pgfpathlineto{\pgfqpoint{3.983205in}{2.045393in}}%
\pgfpathlineto{\pgfqpoint{3.983205in}{2.042444in}}%
\pgfpathmoveto{\pgfqpoint{3.969583in}{2.048342in}}%
\pgfpathlineto{\pgfqpoint{3.969583in}{2.048342in}}%
\pgfpathlineto{\pgfqpoint{3.969583in}{2.051291in}}%
\pgfpathlineto{\pgfqpoint{3.974123in}{2.051291in}}%
\pgfpathlineto{\pgfqpoint{3.974123in}{2.048342in}}%
\pgfpathmoveto{\pgfqpoint{3.965042in}{2.054241in}}%
\pgfpathlineto{\pgfqpoint{3.965042in}{2.054241in}}%
\pgfpathlineto{\pgfqpoint{3.965042in}{2.057190in}}%
\pgfpathlineto{\pgfqpoint{3.969583in}{2.057190in}}%
\pgfpathlineto{\pgfqpoint{3.969583in}{2.054241in}}%
\pgfpathmoveto{\pgfqpoint{3.969583in}{2.051291in}}%
\pgfpathlineto{\pgfqpoint{3.969583in}{2.051291in}}%
\pgfpathlineto{\pgfqpoint{3.969583in}{2.054241in}}%
\pgfpathlineto{\pgfqpoint{3.974123in}{2.054241in}}%
\pgfpathlineto{\pgfqpoint{3.974123in}{2.051291in}}%
\pgfpathmoveto{\pgfqpoint{3.969583in}{2.054241in}}%
\pgfpathlineto{\pgfqpoint{3.969583in}{2.054241in}}%
\pgfpathlineto{\pgfqpoint{3.969583in}{2.057190in}}%
\pgfpathlineto{\pgfqpoint{3.974123in}{2.057190in}}%
\pgfpathlineto{\pgfqpoint{3.974123in}{2.054241in}}%
\pgfpathmoveto{\pgfqpoint{3.974123in}{2.045393in}}%
\pgfpathlineto{\pgfqpoint{3.974123in}{2.045393in}}%
\pgfpathlineto{\pgfqpoint{3.974123in}{2.048342in}}%
\pgfpathlineto{\pgfqpoint{3.978664in}{2.048342in}}%
\pgfpathlineto{\pgfqpoint{3.978664in}{2.045393in}}%
\pgfpathmoveto{\pgfqpoint{3.974123in}{2.048342in}}%
\pgfpathlineto{\pgfqpoint{3.974123in}{2.048342in}}%
\pgfpathlineto{\pgfqpoint{3.974123in}{2.051291in}}%
\pgfpathlineto{\pgfqpoint{3.978664in}{2.051291in}}%
\pgfpathlineto{\pgfqpoint{3.978664in}{2.048342in}}%
\pgfpathmoveto{\pgfqpoint{3.996828in}{2.012952in}}%
\pgfpathlineto{\pgfqpoint{3.996828in}{2.012952in}}%
\pgfpathlineto{\pgfqpoint{3.996828in}{2.015901in}}%
\pgfpathlineto{\pgfqpoint{4.001368in}{2.015901in}}%
\pgfpathlineto{\pgfqpoint{4.001368in}{2.012952in}}%
\pgfpathmoveto{\pgfqpoint{3.992287in}{2.018850in}}%
\pgfpathlineto{\pgfqpoint{3.992287in}{2.018850in}}%
\pgfpathlineto{\pgfqpoint{3.992287in}{2.021799in}}%
\pgfpathlineto{\pgfqpoint{3.996828in}{2.021799in}}%
\pgfpathlineto{\pgfqpoint{3.996828in}{2.018850in}}%
\pgfpathmoveto{\pgfqpoint{3.996828in}{2.015901in}}%
\pgfpathlineto{\pgfqpoint{3.996828in}{2.015901in}}%
\pgfpathlineto{\pgfqpoint{3.996828in}{2.018850in}}%
\pgfpathlineto{\pgfqpoint{4.001368in}{2.018850in}}%
\pgfpathlineto{\pgfqpoint{4.001368in}{2.015901in}}%
\pgfpathmoveto{\pgfqpoint{3.996828in}{2.018850in}}%
\pgfpathlineto{\pgfqpoint{3.996828in}{2.018850in}}%
\pgfpathlineto{\pgfqpoint{3.996828in}{2.021799in}}%
\pgfpathlineto{\pgfqpoint{4.001368in}{2.021799in}}%
\pgfpathlineto{\pgfqpoint{4.001368in}{2.018850in}}%
\pgfpathmoveto{\pgfqpoint{3.987746in}{2.024748in}}%
\pgfpathlineto{\pgfqpoint{3.987746in}{2.024748in}}%
\pgfpathlineto{\pgfqpoint{3.987746in}{2.027698in}}%
\pgfpathlineto{\pgfqpoint{3.992287in}{2.027698in}}%
\pgfpathlineto{\pgfqpoint{3.992287in}{2.024748in}}%
\pgfpathmoveto{\pgfqpoint{3.983205in}{2.030647in}}%
\pgfpathlineto{\pgfqpoint{3.983205in}{2.030647in}}%
\pgfpathlineto{\pgfqpoint{3.983205in}{2.033596in}}%
\pgfpathlineto{\pgfqpoint{3.987746in}{2.033596in}}%
\pgfpathlineto{\pgfqpoint{3.987746in}{2.030647in}}%
\pgfpathmoveto{\pgfqpoint{3.987746in}{2.027698in}}%
\pgfpathlineto{\pgfqpoint{3.987746in}{2.027698in}}%
\pgfpathlineto{\pgfqpoint{3.987746in}{2.030647in}}%
\pgfpathlineto{\pgfqpoint{3.992287in}{2.030647in}}%
\pgfpathlineto{\pgfqpoint{3.992287in}{2.027698in}}%
\pgfpathmoveto{\pgfqpoint{3.987746in}{2.030647in}}%
\pgfpathlineto{\pgfqpoint{3.987746in}{2.030647in}}%
\pgfpathlineto{\pgfqpoint{3.987746in}{2.033596in}}%
\pgfpathlineto{\pgfqpoint{3.992287in}{2.033596in}}%
\pgfpathlineto{\pgfqpoint{3.992287in}{2.030647in}}%
\pgfpathmoveto{\pgfqpoint{3.992287in}{2.021799in}}%
\pgfpathlineto{\pgfqpoint{3.992287in}{2.021799in}}%
\pgfpathlineto{\pgfqpoint{3.992287in}{2.024748in}}%
\pgfpathlineto{\pgfqpoint{3.996828in}{2.024748in}}%
\pgfpathlineto{\pgfqpoint{3.996828in}{2.021799in}}%
\pgfpathmoveto{\pgfqpoint{3.992287in}{2.024748in}}%
\pgfpathlineto{\pgfqpoint{3.992287in}{2.024748in}}%
\pgfpathlineto{\pgfqpoint{3.992287in}{2.027698in}}%
\pgfpathlineto{\pgfqpoint{3.996828in}{2.027698in}}%
\pgfpathlineto{\pgfqpoint{3.996828in}{2.024748in}}%
\pgfpathmoveto{\pgfqpoint{4.001368in}{2.010002in}}%
\pgfpathlineto{\pgfqpoint{4.001368in}{2.010002in}}%
\pgfpathlineto{\pgfqpoint{4.001368in}{2.012952in}}%
\pgfpathlineto{\pgfqpoint{4.005909in}{2.012952in}}%
\pgfpathlineto{\pgfqpoint{4.005909in}{2.010002in}}%
\pgfpathmoveto{\pgfqpoint{4.001368in}{2.012952in}}%
\pgfpathlineto{\pgfqpoint{4.001368in}{2.012952in}}%
\pgfpathlineto{\pgfqpoint{4.001368in}{2.015901in}}%
\pgfpathlineto{\pgfqpoint{4.005909in}{2.015901in}}%
\pgfpathlineto{\pgfqpoint{4.005909in}{2.012952in}}%
\pgfpathmoveto{\pgfqpoint{3.983205in}{2.033596in}}%
\pgfpathlineto{\pgfqpoint{3.983205in}{2.033596in}}%
\pgfpathlineto{\pgfqpoint{3.983205in}{2.036545in}}%
\pgfpathlineto{\pgfqpoint{3.987746in}{2.036545in}}%
\pgfpathlineto{\pgfqpoint{3.987746in}{2.033596in}}%
\pgfpathmoveto{\pgfqpoint{3.983205in}{2.036545in}}%
\pgfpathlineto{\pgfqpoint{3.983205in}{2.036545in}}%
\pgfpathlineto{\pgfqpoint{3.983205in}{2.039495in}}%
\pgfpathlineto{\pgfqpoint{3.987746in}{2.039495in}}%
\pgfpathlineto{\pgfqpoint{3.987746in}{2.036545in}}%
\pgfpathmoveto{\pgfqpoint{3.960501in}{2.060139in}}%
\pgfpathlineto{\pgfqpoint{3.960501in}{2.060139in}}%
\pgfpathlineto{\pgfqpoint{3.960501in}{2.063088in}}%
\pgfpathlineto{\pgfqpoint{3.965042in}{2.063088in}}%
\pgfpathlineto{\pgfqpoint{3.965042in}{2.060139in}}%
\pgfpathmoveto{\pgfqpoint{3.955960in}{2.066038in}}%
\pgfpathlineto{\pgfqpoint{3.955960in}{2.066038in}}%
\pgfpathlineto{\pgfqpoint{3.955960in}{2.068987in}}%
\pgfpathlineto{\pgfqpoint{3.960501in}{2.068987in}}%
\pgfpathlineto{\pgfqpoint{3.960501in}{2.066038in}}%
\pgfpathmoveto{\pgfqpoint{3.960501in}{2.063088in}}%
\pgfpathlineto{\pgfqpoint{3.960501in}{2.063088in}}%
\pgfpathlineto{\pgfqpoint{3.960501in}{2.066038in}}%
\pgfpathlineto{\pgfqpoint{3.965042in}{2.066038in}}%
\pgfpathlineto{\pgfqpoint{3.965042in}{2.063088in}}%
\pgfpathmoveto{\pgfqpoint{3.960501in}{2.066038in}}%
\pgfpathlineto{\pgfqpoint{3.960501in}{2.066038in}}%
\pgfpathlineto{\pgfqpoint{3.960501in}{2.068987in}}%
\pgfpathlineto{\pgfqpoint{3.965042in}{2.068987in}}%
\pgfpathlineto{\pgfqpoint{3.965042in}{2.066038in}}%
\pgfpathmoveto{\pgfqpoint{3.951419in}{2.071936in}}%
\pgfpathlineto{\pgfqpoint{3.951419in}{2.071936in}}%
\pgfpathlineto{\pgfqpoint{3.951419in}{2.074885in}}%
\pgfpathlineto{\pgfqpoint{3.955960in}{2.074885in}}%
\pgfpathlineto{\pgfqpoint{3.955960in}{2.071936in}}%
\pgfpathmoveto{\pgfqpoint{3.946878in}{2.077834in}}%
\pgfpathlineto{\pgfqpoint{3.946878in}{2.077834in}}%
\pgfpathlineto{\pgfqpoint{3.946878in}{2.080784in}}%
\pgfpathlineto{\pgfqpoint{3.951419in}{2.080784in}}%
\pgfpathlineto{\pgfqpoint{3.951419in}{2.077834in}}%
\pgfpathmoveto{\pgfqpoint{3.951419in}{2.074885in}}%
\pgfpathlineto{\pgfqpoint{3.951419in}{2.074885in}}%
\pgfpathlineto{\pgfqpoint{3.951419in}{2.077834in}}%
\pgfpathlineto{\pgfqpoint{3.955960in}{2.077834in}}%
\pgfpathlineto{\pgfqpoint{3.955960in}{2.074885in}}%
\pgfpathmoveto{\pgfqpoint{3.951419in}{2.077834in}}%
\pgfpathlineto{\pgfqpoint{3.951419in}{2.077834in}}%
\pgfpathlineto{\pgfqpoint{3.951419in}{2.080784in}}%
\pgfpathlineto{\pgfqpoint{3.955960in}{2.080784in}}%
\pgfpathlineto{\pgfqpoint{3.955960in}{2.077834in}}%
\pgfpathmoveto{\pgfqpoint{3.955960in}{2.068987in}}%
\pgfpathlineto{\pgfqpoint{3.955960in}{2.068987in}}%
\pgfpathlineto{\pgfqpoint{3.955960in}{2.071936in}}%
\pgfpathlineto{\pgfqpoint{3.960501in}{2.071936in}}%
\pgfpathlineto{\pgfqpoint{3.960501in}{2.068987in}}%
\pgfpathmoveto{\pgfqpoint{3.955960in}{2.071936in}}%
\pgfpathlineto{\pgfqpoint{3.955960in}{2.071936in}}%
\pgfpathlineto{\pgfqpoint{3.955960in}{2.074885in}}%
\pgfpathlineto{\pgfqpoint{3.960501in}{2.074885in}}%
\pgfpathlineto{\pgfqpoint{3.960501in}{2.071936in}}%
\pgfpathmoveto{\pgfqpoint{3.965042in}{2.057190in}}%
\pgfpathlineto{\pgfqpoint{3.965042in}{2.057190in}}%
\pgfpathlineto{\pgfqpoint{3.965042in}{2.060139in}}%
\pgfpathlineto{\pgfqpoint{3.969583in}{2.060139in}}%
\pgfpathlineto{\pgfqpoint{3.969583in}{2.057190in}}%
\pgfpathmoveto{\pgfqpoint{3.965042in}{2.060139in}}%
\pgfpathlineto{\pgfqpoint{3.965042in}{2.060139in}}%
\pgfpathlineto{\pgfqpoint{3.965042in}{2.063088in}}%
\pgfpathlineto{\pgfqpoint{3.969583in}{2.063088in}}%
\pgfpathlineto{\pgfqpoint{3.969583in}{2.060139in}}%
\pgfpathmoveto{\pgfqpoint{3.946878in}{2.080784in}}%
\pgfpathlineto{\pgfqpoint{3.946878in}{2.080784in}}%
\pgfpathlineto{\pgfqpoint{3.946878in}{2.083733in}}%
\pgfpathlineto{\pgfqpoint{3.951419in}{2.083733in}}%
\pgfpathlineto{\pgfqpoint{3.951419in}{2.080784in}}%
\pgfpathmoveto{\pgfqpoint{3.946878in}{2.083733in}}%
\pgfpathlineto{\pgfqpoint{3.946878in}{2.083733in}}%
\pgfpathlineto{\pgfqpoint{3.946878in}{2.086682in}}%
\pgfpathlineto{\pgfqpoint{3.951419in}{2.086682in}}%
\pgfpathlineto{\pgfqpoint{3.951419in}{2.083733in}}%
\pgfpathmoveto{\pgfqpoint{4.232960in}{1.706229in}}%
\pgfpathlineto{\pgfqpoint{4.232960in}{1.706229in}}%
\pgfpathlineto{\pgfqpoint{4.232960in}{1.709178in}}%
\pgfpathlineto{\pgfqpoint{4.237501in}{1.709178in}}%
\pgfpathlineto{\pgfqpoint{4.237501in}{1.706229in}}%
\pgfpathmoveto{\pgfqpoint{4.228419in}{1.712127in}}%
\pgfpathlineto{\pgfqpoint{4.228419in}{1.712127in}}%
\pgfpathlineto{\pgfqpoint{4.228419in}{1.715076in}}%
\pgfpathlineto{\pgfqpoint{4.232960in}{1.715076in}}%
\pgfpathlineto{\pgfqpoint{4.232960in}{1.712127in}}%
\pgfpathmoveto{\pgfqpoint{4.232960in}{1.709178in}}%
\pgfpathlineto{\pgfqpoint{4.232960in}{1.709178in}}%
\pgfpathlineto{\pgfqpoint{4.232960in}{1.712127in}}%
\pgfpathlineto{\pgfqpoint{4.237501in}{1.712127in}}%
\pgfpathlineto{\pgfqpoint{4.237501in}{1.709178in}}%
\pgfpathmoveto{\pgfqpoint{4.232960in}{1.712127in}}%
\pgfpathlineto{\pgfqpoint{4.232960in}{1.712127in}}%
\pgfpathlineto{\pgfqpoint{4.232960in}{1.715076in}}%
\pgfpathlineto{\pgfqpoint{4.237501in}{1.715076in}}%
\pgfpathlineto{\pgfqpoint{4.237501in}{1.712127in}}%
\pgfpathmoveto{\pgfqpoint{4.223877in}{1.718025in}}%
\pgfpathlineto{\pgfqpoint{4.223877in}{1.718025in}}%
\pgfpathlineto{\pgfqpoint{4.223877in}{1.720974in}}%
\pgfpathlineto{\pgfqpoint{4.228419in}{1.720974in}}%
\pgfpathlineto{\pgfqpoint{4.228419in}{1.718025in}}%
\pgfpathmoveto{\pgfqpoint{4.219336in}{1.723923in}}%
\pgfpathlineto{\pgfqpoint{4.219336in}{1.723923in}}%
\pgfpathlineto{\pgfqpoint{4.219336in}{1.726872in}}%
\pgfpathlineto{\pgfqpoint{4.223877in}{1.726872in}}%
\pgfpathlineto{\pgfqpoint{4.223877in}{1.723923in}}%
\pgfpathmoveto{\pgfqpoint{4.223877in}{1.720974in}}%
\pgfpathlineto{\pgfqpoint{4.223877in}{1.720974in}}%
\pgfpathlineto{\pgfqpoint{4.223877in}{1.723923in}}%
\pgfpathlineto{\pgfqpoint{4.228419in}{1.723923in}}%
\pgfpathlineto{\pgfqpoint{4.228419in}{1.720974in}}%
\pgfpathmoveto{\pgfqpoint{4.223877in}{1.723923in}}%
\pgfpathlineto{\pgfqpoint{4.223877in}{1.723923in}}%
\pgfpathlineto{\pgfqpoint{4.223877in}{1.726872in}}%
\pgfpathlineto{\pgfqpoint{4.228419in}{1.726872in}}%
\pgfpathlineto{\pgfqpoint{4.228419in}{1.723923in}}%
\pgfpathmoveto{\pgfqpoint{4.228419in}{1.715076in}}%
\pgfpathlineto{\pgfqpoint{4.228419in}{1.715076in}}%
\pgfpathlineto{\pgfqpoint{4.228419in}{1.718025in}}%
\pgfpathlineto{\pgfqpoint{4.232960in}{1.718025in}}%
\pgfpathlineto{\pgfqpoint{4.232960in}{1.715076in}}%
\pgfpathmoveto{\pgfqpoint{4.228419in}{1.718025in}}%
\pgfpathlineto{\pgfqpoint{4.228419in}{1.718025in}}%
\pgfpathlineto{\pgfqpoint{4.228419in}{1.720974in}}%
\pgfpathlineto{\pgfqpoint{4.232960in}{1.720974in}}%
\pgfpathlineto{\pgfqpoint{4.232960in}{1.718025in}}%
\pgfpathmoveto{\pgfqpoint{4.160302in}{1.800606in}}%
\pgfpathlineto{\pgfqpoint{4.160302in}{1.800606in}}%
\pgfpathlineto{\pgfqpoint{4.160302in}{1.803555in}}%
\pgfpathlineto{\pgfqpoint{4.164843in}{1.803555in}}%
\pgfpathlineto{\pgfqpoint{4.164843in}{1.800606in}}%
\pgfpathmoveto{\pgfqpoint{4.155761in}{1.806504in}}%
\pgfpathlineto{\pgfqpoint{4.155761in}{1.806504in}}%
\pgfpathlineto{\pgfqpoint{4.155761in}{1.809454in}}%
\pgfpathlineto{\pgfqpoint{4.160302in}{1.809454in}}%
\pgfpathlineto{\pgfqpoint{4.160302in}{1.806504in}}%
\pgfpathmoveto{\pgfqpoint{4.160302in}{1.803555in}}%
\pgfpathlineto{\pgfqpoint{4.160302in}{1.803555in}}%
\pgfpathlineto{\pgfqpoint{4.160302in}{1.806504in}}%
\pgfpathlineto{\pgfqpoint{4.164843in}{1.806504in}}%
\pgfpathlineto{\pgfqpoint{4.164843in}{1.803555in}}%
\pgfpathmoveto{\pgfqpoint{4.160302in}{1.806504in}}%
\pgfpathlineto{\pgfqpoint{4.160302in}{1.806504in}}%
\pgfpathlineto{\pgfqpoint{4.160302in}{1.809454in}}%
\pgfpathlineto{\pgfqpoint{4.164843in}{1.809454in}}%
\pgfpathlineto{\pgfqpoint{4.164843in}{1.806504in}}%
\pgfpathmoveto{\pgfqpoint{4.151219in}{1.812403in}}%
\pgfpathlineto{\pgfqpoint{4.151219in}{1.812403in}}%
\pgfpathlineto{\pgfqpoint{4.151219in}{1.815352in}}%
\pgfpathlineto{\pgfqpoint{4.155761in}{1.815352in}}%
\pgfpathlineto{\pgfqpoint{4.155761in}{1.812403in}}%
\pgfpathmoveto{\pgfqpoint{4.146678in}{1.818302in}}%
\pgfpathlineto{\pgfqpoint{4.146678in}{1.818302in}}%
\pgfpathlineto{\pgfqpoint{4.146678in}{1.821251in}}%
\pgfpathlineto{\pgfqpoint{4.151219in}{1.821251in}}%
\pgfpathlineto{\pgfqpoint{4.151219in}{1.818302in}}%
\pgfpathmoveto{\pgfqpoint{4.151219in}{1.815352in}}%
\pgfpathlineto{\pgfqpoint{4.151219in}{1.815352in}}%
\pgfpathlineto{\pgfqpoint{4.151219in}{1.818302in}}%
\pgfpathlineto{\pgfqpoint{4.155761in}{1.818302in}}%
\pgfpathlineto{\pgfqpoint{4.155761in}{1.815352in}}%
\pgfpathmoveto{\pgfqpoint{4.151219in}{1.818302in}}%
\pgfpathlineto{\pgfqpoint{4.151219in}{1.818302in}}%
\pgfpathlineto{\pgfqpoint{4.151219in}{1.821251in}}%
\pgfpathlineto{\pgfqpoint{4.155761in}{1.821251in}}%
\pgfpathlineto{\pgfqpoint{4.155761in}{1.818302in}}%
\pgfpathmoveto{\pgfqpoint{4.155761in}{1.809454in}}%
\pgfpathlineto{\pgfqpoint{4.155761in}{1.809454in}}%
\pgfpathlineto{\pgfqpoint{4.155761in}{1.812403in}}%
\pgfpathlineto{\pgfqpoint{4.160302in}{1.812403in}}%
\pgfpathlineto{\pgfqpoint{4.160302in}{1.809454in}}%
\pgfpathmoveto{\pgfqpoint{4.155761in}{1.812403in}}%
\pgfpathlineto{\pgfqpoint{4.155761in}{1.812403in}}%
\pgfpathlineto{\pgfqpoint{4.155761in}{1.815352in}}%
\pgfpathlineto{\pgfqpoint{4.160302in}{1.815352in}}%
\pgfpathlineto{\pgfqpoint{4.160302in}{1.812403in}}%
\pgfpathmoveto{\pgfqpoint{4.196631in}{1.753416in}}%
\pgfpathlineto{\pgfqpoint{4.196631in}{1.753416in}}%
\pgfpathlineto{\pgfqpoint{4.196631in}{1.756366in}}%
\pgfpathlineto{\pgfqpoint{4.201172in}{1.756366in}}%
\pgfpathlineto{\pgfqpoint{4.201172in}{1.753416in}}%
\pgfpathmoveto{\pgfqpoint{4.192090in}{1.759315in}}%
\pgfpathlineto{\pgfqpoint{4.192090in}{1.759315in}}%
\pgfpathlineto{\pgfqpoint{4.192090in}{1.762264in}}%
\pgfpathlineto{\pgfqpoint{4.196631in}{1.762264in}}%
\pgfpathlineto{\pgfqpoint{4.196631in}{1.759315in}}%
\pgfpathmoveto{\pgfqpoint{4.196631in}{1.756366in}}%
\pgfpathlineto{\pgfqpoint{4.196631in}{1.756366in}}%
\pgfpathlineto{\pgfqpoint{4.196631in}{1.759315in}}%
\pgfpathlineto{\pgfqpoint{4.201172in}{1.759315in}}%
\pgfpathlineto{\pgfqpoint{4.201172in}{1.756366in}}%
\pgfpathmoveto{\pgfqpoint{4.196631in}{1.759315in}}%
\pgfpathlineto{\pgfqpoint{4.196631in}{1.759315in}}%
\pgfpathlineto{\pgfqpoint{4.196631in}{1.762264in}}%
\pgfpathlineto{\pgfqpoint{4.201172in}{1.762264in}}%
\pgfpathlineto{\pgfqpoint{4.201172in}{1.759315in}}%
\pgfpathmoveto{\pgfqpoint{4.187548in}{1.765214in}}%
\pgfpathlineto{\pgfqpoint{4.187548in}{1.765214in}}%
\pgfpathlineto{\pgfqpoint{4.187548in}{1.768163in}}%
\pgfpathlineto{\pgfqpoint{4.192090in}{1.768163in}}%
\pgfpathlineto{\pgfqpoint{4.192090in}{1.765214in}}%
\pgfpathmoveto{\pgfqpoint{4.183007in}{1.771112in}}%
\pgfpathlineto{\pgfqpoint{4.183007in}{1.771112in}}%
\pgfpathlineto{\pgfqpoint{4.183007in}{1.774062in}}%
\pgfpathlineto{\pgfqpoint{4.187548in}{1.774062in}}%
\pgfpathlineto{\pgfqpoint{4.187548in}{1.771112in}}%
\pgfpathmoveto{\pgfqpoint{4.187548in}{1.768163in}}%
\pgfpathlineto{\pgfqpoint{4.187548in}{1.768163in}}%
\pgfpathlineto{\pgfqpoint{4.187548in}{1.771112in}}%
\pgfpathlineto{\pgfqpoint{4.192090in}{1.771112in}}%
\pgfpathlineto{\pgfqpoint{4.192090in}{1.768163in}}%
\pgfpathmoveto{\pgfqpoint{4.187548in}{1.771112in}}%
\pgfpathlineto{\pgfqpoint{4.187548in}{1.771112in}}%
\pgfpathlineto{\pgfqpoint{4.187548in}{1.774062in}}%
\pgfpathlineto{\pgfqpoint{4.192090in}{1.774062in}}%
\pgfpathlineto{\pgfqpoint{4.192090in}{1.771112in}}%
\pgfpathmoveto{\pgfqpoint{4.192090in}{1.762264in}}%
\pgfpathlineto{\pgfqpoint{4.192090in}{1.762264in}}%
\pgfpathlineto{\pgfqpoint{4.192090in}{1.765214in}}%
\pgfpathlineto{\pgfqpoint{4.196631in}{1.765214in}}%
\pgfpathlineto{\pgfqpoint{4.196631in}{1.762264in}}%
\pgfpathmoveto{\pgfqpoint{4.192090in}{1.765214in}}%
\pgfpathlineto{\pgfqpoint{4.192090in}{1.765214in}}%
\pgfpathlineto{\pgfqpoint{4.192090in}{1.768163in}}%
\pgfpathlineto{\pgfqpoint{4.196631in}{1.768163in}}%
\pgfpathlineto{\pgfqpoint{4.196631in}{1.765214in}}%
\pgfpathmoveto{\pgfqpoint{4.214795in}{1.729821in}}%
\pgfpathlineto{\pgfqpoint{4.214795in}{1.729821in}}%
\pgfpathlineto{\pgfqpoint{4.214795in}{1.732771in}}%
\pgfpathlineto{\pgfqpoint{4.219336in}{1.732771in}}%
\pgfpathlineto{\pgfqpoint{4.219336in}{1.729821in}}%
\pgfpathmoveto{\pgfqpoint{4.210254in}{1.735720in}}%
\pgfpathlineto{\pgfqpoint{4.210254in}{1.735720in}}%
\pgfpathlineto{\pgfqpoint{4.210254in}{1.738670in}}%
\pgfpathlineto{\pgfqpoint{4.214795in}{1.738670in}}%
\pgfpathlineto{\pgfqpoint{4.214795in}{1.735720in}}%
\pgfpathmoveto{\pgfqpoint{4.214795in}{1.732771in}}%
\pgfpathlineto{\pgfqpoint{4.214795in}{1.732771in}}%
\pgfpathlineto{\pgfqpoint{4.214795in}{1.735720in}}%
\pgfpathlineto{\pgfqpoint{4.219336in}{1.735720in}}%
\pgfpathlineto{\pgfqpoint{4.219336in}{1.732771in}}%
\pgfpathmoveto{\pgfqpoint{4.214795in}{1.735720in}}%
\pgfpathlineto{\pgfqpoint{4.214795in}{1.735720in}}%
\pgfpathlineto{\pgfqpoint{4.214795in}{1.738670in}}%
\pgfpathlineto{\pgfqpoint{4.219336in}{1.738670in}}%
\pgfpathlineto{\pgfqpoint{4.219336in}{1.735720in}}%
\pgfpathmoveto{\pgfqpoint{4.205713in}{1.741619in}}%
\pgfpathlineto{\pgfqpoint{4.205713in}{1.741619in}}%
\pgfpathlineto{\pgfqpoint{4.205713in}{1.744568in}}%
\pgfpathlineto{\pgfqpoint{4.210254in}{1.744568in}}%
\pgfpathlineto{\pgfqpoint{4.210254in}{1.741619in}}%
\pgfpathmoveto{\pgfqpoint{4.201172in}{1.747518in}}%
\pgfpathlineto{\pgfqpoint{4.201172in}{1.747518in}}%
\pgfpathlineto{\pgfqpoint{4.201172in}{1.750467in}}%
\pgfpathlineto{\pgfqpoint{4.205713in}{1.750467in}}%
\pgfpathlineto{\pgfqpoint{4.205713in}{1.747518in}}%
\pgfpathmoveto{\pgfqpoint{4.205713in}{1.744568in}}%
\pgfpathlineto{\pgfqpoint{4.205713in}{1.744568in}}%
\pgfpathlineto{\pgfqpoint{4.205713in}{1.747518in}}%
\pgfpathlineto{\pgfqpoint{4.210254in}{1.747518in}}%
\pgfpathlineto{\pgfqpoint{4.210254in}{1.744568in}}%
\pgfpathmoveto{\pgfqpoint{4.205713in}{1.747518in}}%
\pgfpathlineto{\pgfqpoint{4.205713in}{1.747518in}}%
\pgfpathlineto{\pgfqpoint{4.205713in}{1.750467in}}%
\pgfpathlineto{\pgfqpoint{4.210254in}{1.750467in}}%
\pgfpathlineto{\pgfqpoint{4.210254in}{1.747518in}}%
\pgfpathmoveto{\pgfqpoint{4.210254in}{1.738670in}}%
\pgfpathlineto{\pgfqpoint{4.210254in}{1.738670in}}%
\pgfpathlineto{\pgfqpoint{4.210254in}{1.741619in}}%
\pgfpathlineto{\pgfqpoint{4.214795in}{1.741619in}}%
\pgfpathlineto{\pgfqpoint{4.214795in}{1.738670in}}%
\pgfpathmoveto{\pgfqpoint{4.210254in}{1.741619in}}%
\pgfpathlineto{\pgfqpoint{4.210254in}{1.741619in}}%
\pgfpathlineto{\pgfqpoint{4.210254in}{1.744568in}}%
\pgfpathlineto{\pgfqpoint{4.214795in}{1.744568in}}%
\pgfpathlineto{\pgfqpoint{4.214795in}{1.741619in}}%
\pgfpathmoveto{\pgfqpoint{4.219336in}{1.726872in}}%
\pgfpathlineto{\pgfqpoint{4.219336in}{1.726872in}}%
\pgfpathlineto{\pgfqpoint{4.219336in}{1.729821in}}%
\pgfpathlineto{\pgfqpoint{4.223877in}{1.729821in}}%
\pgfpathlineto{\pgfqpoint{4.223877in}{1.726872in}}%
\pgfpathmoveto{\pgfqpoint{4.219336in}{1.729821in}}%
\pgfpathlineto{\pgfqpoint{4.219336in}{1.729821in}}%
\pgfpathlineto{\pgfqpoint{4.219336in}{1.732771in}}%
\pgfpathlineto{\pgfqpoint{4.223877in}{1.732771in}}%
\pgfpathlineto{\pgfqpoint{4.223877in}{1.729821in}}%
\pgfpathmoveto{\pgfqpoint{4.201172in}{1.750467in}}%
\pgfpathlineto{\pgfqpoint{4.201172in}{1.750467in}}%
\pgfpathlineto{\pgfqpoint{4.201172in}{1.753416in}}%
\pgfpathlineto{\pgfqpoint{4.205713in}{1.753416in}}%
\pgfpathlineto{\pgfqpoint{4.205713in}{1.750467in}}%
\pgfpathmoveto{\pgfqpoint{4.201172in}{1.753416in}}%
\pgfpathlineto{\pgfqpoint{4.201172in}{1.753416in}}%
\pgfpathlineto{\pgfqpoint{4.201172in}{1.756366in}}%
\pgfpathlineto{\pgfqpoint{4.205713in}{1.756366in}}%
\pgfpathlineto{\pgfqpoint{4.205713in}{1.753416in}}%
\pgfpathmoveto{\pgfqpoint{4.178466in}{1.777011in}}%
\pgfpathlineto{\pgfqpoint{4.178466in}{1.777011in}}%
\pgfpathlineto{\pgfqpoint{4.178466in}{1.779960in}}%
\pgfpathlineto{\pgfqpoint{4.183007in}{1.779960in}}%
\pgfpathlineto{\pgfqpoint{4.183007in}{1.777011in}}%
\pgfpathmoveto{\pgfqpoint{4.173925in}{1.782910in}}%
\pgfpathlineto{\pgfqpoint{4.173925in}{1.782910in}}%
\pgfpathlineto{\pgfqpoint{4.173925in}{1.785859in}}%
\pgfpathlineto{\pgfqpoint{4.178466in}{1.785859in}}%
\pgfpathlineto{\pgfqpoint{4.178466in}{1.782910in}}%
\pgfpathmoveto{\pgfqpoint{4.178466in}{1.779960in}}%
\pgfpathlineto{\pgfqpoint{4.178466in}{1.779960in}}%
\pgfpathlineto{\pgfqpoint{4.178466in}{1.782910in}}%
\pgfpathlineto{\pgfqpoint{4.183007in}{1.782910in}}%
\pgfpathlineto{\pgfqpoint{4.183007in}{1.779960in}}%
\pgfpathmoveto{\pgfqpoint{4.178466in}{1.782910in}}%
\pgfpathlineto{\pgfqpoint{4.178466in}{1.782910in}}%
\pgfpathlineto{\pgfqpoint{4.178466in}{1.785859in}}%
\pgfpathlineto{\pgfqpoint{4.183007in}{1.785859in}}%
\pgfpathlineto{\pgfqpoint{4.183007in}{1.782910in}}%
\pgfpathmoveto{\pgfqpoint{4.169384in}{1.788808in}}%
\pgfpathlineto{\pgfqpoint{4.169384in}{1.788808in}}%
\pgfpathlineto{\pgfqpoint{4.169384in}{1.791758in}}%
\pgfpathlineto{\pgfqpoint{4.173925in}{1.791758in}}%
\pgfpathlineto{\pgfqpoint{4.173925in}{1.788808in}}%
\pgfpathmoveto{\pgfqpoint{4.164843in}{1.794707in}}%
\pgfpathlineto{\pgfqpoint{4.164843in}{1.794707in}}%
\pgfpathlineto{\pgfqpoint{4.164843in}{1.797656in}}%
\pgfpathlineto{\pgfqpoint{4.169384in}{1.797656in}}%
\pgfpathlineto{\pgfqpoint{4.169384in}{1.794707in}}%
\pgfpathmoveto{\pgfqpoint{4.169384in}{1.791758in}}%
\pgfpathlineto{\pgfqpoint{4.169384in}{1.791758in}}%
\pgfpathlineto{\pgfqpoint{4.169384in}{1.794707in}}%
\pgfpathlineto{\pgfqpoint{4.173925in}{1.794707in}}%
\pgfpathlineto{\pgfqpoint{4.173925in}{1.791758in}}%
\pgfpathmoveto{\pgfqpoint{4.169384in}{1.794707in}}%
\pgfpathlineto{\pgfqpoint{4.169384in}{1.794707in}}%
\pgfpathlineto{\pgfqpoint{4.169384in}{1.797656in}}%
\pgfpathlineto{\pgfqpoint{4.173925in}{1.797656in}}%
\pgfpathlineto{\pgfqpoint{4.173925in}{1.794707in}}%
\pgfpathmoveto{\pgfqpoint{4.173925in}{1.785859in}}%
\pgfpathlineto{\pgfqpoint{4.173925in}{1.785859in}}%
\pgfpathlineto{\pgfqpoint{4.173925in}{1.788808in}}%
\pgfpathlineto{\pgfqpoint{4.178466in}{1.788808in}}%
\pgfpathlineto{\pgfqpoint{4.178466in}{1.785859in}}%
\pgfpathmoveto{\pgfqpoint{4.173925in}{1.788808in}}%
\pgfpathlineto{\pgfqpoint{4.173925in}{1.788808in}}%
\pgfpathlineto{\pgfqpoint{4.173925in}{1.791758in}}%
\pgfpathlineto{\pgfqpoint{4.178466in}{1.791758in}}%
\pgfpathlineto{\pgfqpoint{4.178466in}{1.788808in}}%
\pgfpathmoveto{\pgfqpoint{4.183007in}{1.774062in}}%
\pgfpathlineto{\pgfqpoint{4.183007in}{1.774062in}}%
\pgfpathlineto{\pgfqpoint{4.183007in}{1.777011in}}%
\pgfpathlineto{\pgfqpoint{4.187548in}{1.777011in}}%
\pgfpathlineto{\pgfqpoint{4.187548in}{1.774062in}}%
\pgfpathmoveto{\pgfqpoint{4.183007in}{1.777011in}}%
\pgfpathlineto{\pgfqpoint{4.183007in}{1.777011in}}%
\pgfpathlineto{\pgfqpoint{4.183007in}{1.779960in}}%
\pgfpathlineto{\pgfqpoint{4.187548in}{1.779960in}}%
\pgfpathlineto{\pgfqpoint{4.187548in}{1.777011in}}%
\pgfpathmoveto{\pgfqpoint{4.164843in}{1.797656in}}%
\pgfpathlineto{\pgfqpoint{4.164843in}{1.797656in}}%
\pgfpathlineto{\pgfqpoint{4.164843in}{1.800606in}}%
\pgfpathlineto{\pgfqpoint{4.169384in}{1.800606in}}%
\pgfpathlineto{\pgfqpoint{4.169384in}{1.797656in}}%
\pgfpathmoveto{\pgfqpoint{4.164843in}{1.800606in}}%
\pgfpathlineto{\pgfqpoint{4.164843in}{1.800606in}}%
\pgfpathlineto{\pgfqpoint{4.164843in}{1.803555in}}%
\pgfpathlineto{\pgfqpoint{4.169384in}{1.803555in}}%
\pgfpathlineto{\pgfqpoint{4.169384in}{1.800606in}}%
\pgfpathmoveto{\pgfqpoint{4.123973in}{1.847793in}}%
\pgfpathlineto{\pgfqpoint{4.123973in}{1.847793in}}%
\pgfpathlineto{\pgfqpoint{4.123973in}{1.850743in}}%
\pgfpathlineto{\pgfqpoint{4.128514in}{1.850743in}}%
\pgfpathlineto{\pgfqpoint{4.128514in}{1.847793in}}%
\pgfpathmoveto{\pgfqpoint{4.119432in}{1.853692in}}%
\pgfpathlineto{\pgfqpoint{4.119432in}{1.853692in}}%
\pgfpathlineto{\pgfqpoint{4.119432in}{1.856641in}}%
\pgfpathlineto{\pgfqpoint{4.123973in}{1.856641in}}%
\pgfpathlineto{\pgfqpoint{4.123973in}{1.853692in}}%
\pgfpathmoveto{\pgfqpoint{4.123973in}{1.850743in}}%
\pgfpathlineto{\pgfqpoint{4.123973in}{1.850743in}}%
\pgfpathlineto{\pgfqpoint{4.123973in}{1.853692in}}%
\pgfpathlineto{\pgfqpoint{4.128514in}{1.853692in}}%
\pgfpathlineto{\pgfqpoint{4.128514in}{1.850743in}}%
\pgfpathmoveto{\pgfqpoint{4.123973in}{1.853692in}}%
\pgfpathlineto{\pgfqpoint{4.123973in}{1.853692in}}%
\pgfpathlineto{\pgfqpoint{4.123973in}{1.856641in}}%
\pgfpathlineto{\pgfqpoint{4.128514in}{1.856641in}}%
\pgfpathlineto{\pgfqpoint{4.128514in}{1.853692in}}%
\pgfpathmoveto{\pgfqpoint{4.114890in}{1.859590in}}%
\pgfpathlineto{\pgfqpoint{4.114890in}{1.859590in}}%
\pgfpathlineto{\pgfqpoint{4.114890in}{1.862539in}}%
\pgfpathlineto{\pgfqpoint{4.119432in}{1.862539in}}%
\pgfpathlineto{\pgfqpoint{4.119432in}{1.859590in}}%
\pgfpathmoveto{\pgfqpoint{4.110349in}{1.865488in}}%
\pgfpathlineto{\pgfqpoint{4.110349in}{1.865488in}}%
\pgfpathlineto{\pgfqpoint{4.110349in}{1.868437in}}%
\pgfpathlineto{\pgfqpoint{4.114890in}{1.868437in}}%
\pgfpathlineto{\pgfqpoint{4.114890in}{1.865488in}}%
\pgfpathmoveto{\pgfqpoint{4.114890in}{1.862539in}}%
\pgfpathlineto{\pgfqpoint{4.114890in}{1.862539in}}%
\pgfpathlineto{\pgfqpoint{4.114890in}{1.865488in}}%
\pgfpathlineto{\pgfqpoint{4.119432in}{1.865488in}}%
\pgfpathlineto{\pgfqpoint{4.119432in}{1.862539in}}%
\pgfpathmoveto{\pgfqpoint{4.114890in}{1.865488in}}%
\pgfpathlineto{\pgfqpoint{4.114890in}{1.865488in}}%
\pgfpathlineto{\pgfqpoint{4.114890in}{1.868437in}}%
\pgfpathlineto{\pgfqpoint{4.119432in}{1.868437in}}%
\pgfpathlineto{\pgfqpoint{4.119432in}{1.865488in}}%
\pgfpathmoveto{\pgfqpoint{4.119432in}{1.856641in}}%
\pgfpathlineto{\pgfqpoint{4.119432in}{1.856641in}}%
\pgfpathlineto{\pgfqpoint{4.119432in}{1.859590in}}%
\pgfpathlineto{\pgfqpoint{4.123973in}{1.859590in}}%
\pgfpathlineto{\pgfqpoint{4.123973in}{1.856641in}}%
\pgfpathmoveto{\pgfqpoint{4.119432in}{1.859590in}}%
\pgfpathlineto{\pgfqpoint{4.119432in}{1.859590in}}%
\pgfpathlineto{\pgfqpoint{4.119432in}{1.862539in}}%
\pgfpathlineto{\pgfqpoint{4.123973in}{1.862539in}}%
\pgfpathlineto{\pgfqpoint{4.123973in}{1.859590in}}%
\pgfpathmoveto{\pgfqpoint{4.142137in}{1.824200in}}%
\pgfpathlineto{\pgfqpoint{4.142137in}{1.824200in}}%
\pgfpathlineto{\pgfqpoint{4.142137in}{1.827149in}}%
\pgfpathlineto{\pgfqpoint{4.146678in}{1.827149in}}%
\pgfpathlineto{\pgfqpoint{4.146678in}{1.824200in}}%
\pgfpathmoveto{\pgfqpoint{4.137596in}{1.830099in}}%
\pgfpathlineto{\pgfqpoint{4.137596in}{1.830099in}}%
\pgfpathlineto{\pgfqpoint{4.137596in}{1.833048in}}%
\pgfpathlineto{\pgfqpoint{4.142137in}{1.833048in}}%
\pgfpathlineto{\pgfqpoint{4.142137in}{1.830099in}}%
\pgfpathmoveto{\pgfqpoint{4.142137in}{1.827149in}}%
\pgfpathlineto{\pgfqpoint{4.142137in}{1.827149in}}%
\pgfpathlineto{\pgfqpoint{4.142137in}{1.830099in}}%
\pgfpathlineto{\pgfqpoint{4.146678in}{1.830099in}}%
\pgfpathlineto{\pgfqpoint{4.146678in}{1.827149in}}%
\pgfpathmoveto{\pgfqpoint{4.142137in}{1.830099in}}%
\pgfpathlineto{\pgfqpoint{4.142137in}{1.830099in}}%
\pgfpathlineto{\pgfqpoint{4.142137in}{1.833048in}}%
\pgfpathlineto{\pgfqpoint{4.146678in}{1.833048in}}%
\pgfpathlineto{\pgfqpoint{4.146678in}{1.830099in}}%
\pgfpathmoveto{\pgfqpoint{4.133055in}{1.835997in}}%
\pgfpathlineto{\pgfqpoint{4.133055in}{1.835997in}}%
\pgfpathlineto{\pgfqpoint{4.133055in}{1.838946in}}%
\pgfpathlineto{\pgfqpoint{4.137596in}{1.838946in}}%
\pgfpathlineto{\pgfqpoint{4.137596in}{1.835997in}}%
\pgfpathmoveto{\pgfqpoint{4.128514in}{1.841895in}}%
\pgfpathlineto{\pgfqpoint{4.128514in}{1.841895in}}%
\pgfpathlineto{\pgfqpoint{4.128514in}{1.844844in}}%
\pgfpathlineto{\pgfqpoint{4.133055in}{1.844844in}}%
\pgfpathlineto{\pgfqpoint{4.133055in}{1.841895in}}%
\pgfpathmoveto{\pgfqpoint{4.133055in}{1.838946in}}%
\pgfpathlineto{\pgfqpoint{4.133055in}{1.838946in}}%
\pgfpathlineto{\pgfqpoint{4.133055in}{1.841895in}}%
\pgfpathlineto{\pgfqpoint{4.137596in}{1.841895in}}%
\pgfpathlineto{\pgfqpoint{4.137596in}{1.838946in}}%
\pgfpathmoveto{\pgfqpoint{4.133055in}{1.841895in}}%
\pgfpathlineto{\pgfqpoint{4.133055in}{1.841895in}}%
\pgfpathlineto{\pgfqpoint{4.133055in}{1.844844in}}%
\pgfpathlineto{\pgfqpoint{4.137596in}{1.844844in}}%
\pgfpathlineto{\pgfqpoint{4.137596in}{1.841895in}}%
\pgfpathmoveto{\pgfqpoint{4.137596in}{1.833048in}}%
\pgfpathlineto{\pgfqpoint{4.137596in}{1.833048in}}%
\pgfpathlineto{\pgfqpoint{4.137596in}{1.835997in}}%
\pgfpathlineto{\pgfqpoint{4.142137in}{1.835997in}}%
\pgfpathlineto{\pgfqpoint{4.142137in}{1.833048in}}%
\pgfpathmoveto{\pgfqpoint{4.137596in}{1.835997in}}%
\pgfpathlineto{\pgfqpoint{4.137596in}{1.835997in}}%
\pgfpathlineto{\pgfqpoint{4.137596in}{1.838946in}}%
\pgfpathlineto{\pgfqpoint{4.142137in}{1.838946in}}%
\pgfpathlineto{\pgfqpoint{4.142137in}{1.835997in}}%
\pgfpathmoveto{\pgfqpoint{4.146678in}{1.821251in}}%
\pgfpathlineto{\pgfqpoint{4.146678in}{1.821251in}}%
\pgfpathlineto{\pgfqpoint{4.146678in}{1.824200in}}%
\pgfpathlineto{\pgfqpoint{4.151219in}{1.824200in}}%
\pgfpathlineto{\pgfqpoint{4.151219in}{1.821251in}}%
\pgfpathmoveto{\pgfqpoint{4.146678in}{1.824200in}}%
\pgfpathlineto{\pgfqpoint{4.146678in}{1.824200in}}%
\pgfpathlineto{\pgfqpoint{4.146678in}{1.827149in}}%
\pgfpathlineto{\pgfqpoint{4.151219in}{1.827149in}}%
\pgfpathlineto{\pgfqpoint{4.151219in}{1.824200in}}%
\pgfpathmoveto{\pgfqpoint{4.128514in}{1.844844in}}%
\pgfpathlineto{\pgfqpoint{4.128514in}{1.844844in}}%
\pgfpathlineto{\pgfqpoint{4.128514in}{1.847793in}}%
\pgfpathlineto{\pgfqpoint{4.133055in}{1.847793in}}%
\pgfpathlineto{\pgfqpoint{4.133055in}{1.844844in}}%
\pgfpathmoveto{\pgfqpoint{4.128514in}{1.847793in}}%
\pgfpathlineto{\pgfqpoint{4.128514in}{1.847793in}}%
\pgfpathlineto{\pgfqpoint{4.128514in}{1.850743in}}%
\pgfpathlineto{\pgfqpoint{4.133055in}{1.850743in}}%
\pgfpathlineto{\pgfqpoint{4.133055in}{1.847793in}}%
\pgfpathmoveto{\pgfqpoint{4.105808in}{1.871387in}}%
\pgfpathlineto{\pgfqpoint{4.105808in}{1.871387in}}%
\pgfpathlineto{\pgfqpoint{4.105808in}{1.874336in}}%
\pgfpathlineto{\pgfqpoint{4.110349in}{1.874336in}}%
\pgfpathlineto{\pgfqpoint{4.110349in}{1.871387in}}%
\pgfpathmoveto{\pgfqpoint{4.101267in}{1.877285in}}%
\pgfpathlineto{\pgfqpoint{4.101267in}{1.877285in}}%
\pgfpathlineto{\pgfqpoint{4.101267in}{1.880234in}}%
\pgfpathlineto{\pgfqpoint{4.105808in}{1.880234in}}%
\pgfpathlineto{\pgfqpoint{4.105808in}{1.877285in}}%
\pgfpathmoveto{\pgfqpoint{4.105808in}{1.874336in}}%
\pgfpathlineto{\pgfqpoint{4.105808in}{1.874336in}}%
\pgfpathlineto{\pgfqpoint{4.105808in}{1.877285in}}%
\pgfpathlineto{\pgfqpoint{4.110349in}{1.877285in}}%
\pgfpathlineto{\pgfqpoint{4.110349in}{1.874336in}}%
\pgfpathmoveto{\pgfqpoint{4.105808in}{1.877285in}}%
\pgfpathlineto{\pgfqpoint{4.105808in}{1.877285in}}%
\pgfpathlineto{\pgfqpoint{4.105808in}{1.880234in}}%
\pgfpathlineto{\pgfqpoint{4.110349in}{1.880234in}}%
\pgfpathlineto{\pgfqpoint{4.110349in}{1.877285in}}%
\pgfpathmoveto{\pgfqpoint{4.096726in}{1.883183in}}%
\pgfpathlineto{\pgfqpoint{4.096726in}{1.883183in}}%
\pgfpathlineto{\pgfqpoint{4.096726in}{1.886132in}}%
\pgfpathlineto{\pgfqpoint{4.101267in}{1.886132in}}%
\pgfpathlineto{\pgfqpoint{4.101267in}{1.883183in}}%
\pgfpathmoveto{\pgfqpoint{4.092185in}{1.889081in}}%
\pgfpathlineto{\pgfqpoint{4.092185in}{1.889081in}}%
\pgfpathlineto{\pgfqpoint{4.092185in}{1.892031in}}%
\pgfpathlineto{\pgfqpoint{4.096726in}{1.892031in}}%
\pgfpathlineto{\pgfqpoint{4.096726in}{1.889081in}}%
\pgfpathmoveto{\pgfqpoint{4.096726in}{1.886132in}}%
\pgfpathlineto{\pgfqpoint{4.096726in}{1.886132in}}%
\pgfpathlineto{\pgfqpoint{4.096726in}{1.889081in}}%
\pgfpathlineto{\pgfqpoint{4.101267in}{1.889081in}}%
\pgfpathlineto{\pgfqpoint{4.101267in}{1.886132in}}%
\pgfpathmoveto{\pgfqpoint{4.096726in}{1.889081in}}%
\pgfpathlineto{\pgfqpoint{4.096726in}{1.889081in}}%
\pgfpathlineto{\pgfqpoint{4.096726in}{1.892031in}}%
\pgfpathlineto{\pgfqpoint{4.101267in}{1.892031in}}%
\pgfpathlineto{\pgfqpoint{4.101267in}{1.889081in}}%
\pgfpathmoveto{\pgfqpoint{4.101267in}{1.880234in}}%
\pgfpathlineto{\pgfqpoint{4.101267in}{1.880234in}}%
\pgfpathlineto{\pgfqpoint{4.101267in}{1.883183in}}%
\pgfpathlineto{\pgfqpoint{4.105808in}{1.883183in}}%
\pgfpathlineto{\pgfqpoint{4.105808in}{1.880234in}}%
\pgfpathmoveto{\pgfqpoint{4.101267in}{1.883183in}}%
\pgfpathlineto{\pgfqpoint{4.101267in}{1.883183in}}%
\pgfpathlineto{\pgfqpoint{4.101267in}{1.886132in}}%
\pgfpathlineto{\pgfqpoint{4.105808in}{1.886132in}}%
\pgfpathlineto{\pgfqpoint{4.105808in}{1.883183in}}%
\pgfpathmoveto{\pgfqpoint{4.110349in}{1.868437in}}%
\pgfpathlineto{\pgfqpoint{4.110349in}{1.868437in}}%
\pgfpathlineto{\pgfqpoint{4.110349in}{1.871387in}}%
\pgfpathlineto{\pgfqpoint{4.114890in}{1.871387in}}%
\pgfpathlineto{\pgfqpoint{4.114890in}{1.868437in}}%
\pgfpathmoveto{\pgfqpoint{4.110349in}{1.871387in}}%
\pgfpathlineto{\pgfqpoint{4.110349in}{1.871387in}}%
\pgfpathlineto{\pgfqpoint{4.110349in}{1.874336in}}%
\pgfpathlineto{\pgfqpoint{4.114890in}{1.874336in}}%
\pgfpathlineto{\pgfqpoint{4.114890in}{1.871387in}}%
\pgfpathmoveto{\pgfqpoint{4.092185in}{1.892031in}}%
\pgfpathlineto{\pgfqpoint{4.092185in}{1.892031in}}%
\pgfpathlineto{\pgfqpoint{4.092185in}{1.894980in}}%
\pgfpathlineto{\pgfqpoint{4.096726in}{1.894980in}}%
\pgfpathlineto{\pgfqpoint{4.096726in}{1.892031in}}%
\pgfpathmoveto{\pgfqpoint{4.092185in}{1.894980in}}%
\pgfpathlineto{\pgfqpoint{4.092185in}{1.894980in}}%
\pgfpathlineto{\pgfqpoint{4.092185in}{1.897929in}}%
\pgfpathlineto{\pgfqpoint{4.096726in}{1.897929in}}%
\pgfpathlineto{\pgfqpoint{4.096726in}{1.894980in}}%
\pgfpathmoveto{\pgfqpoint{4.378274in}{1.517479in}}%
\pgfpathlineto{\pgfqpoint{4.378274in}{1.517479in}}%
\pgfpathlineto{\pgfqpoint{4.378274in}{1.520428in}}%
\pgfpathlineto{\pgfqpoint{4.382815in}{1.520428in}}%
\pgfpathlineto{\pgfqpoint{4.382815in}{1.517479in}}%
\pgfpathmoveto{\pgfqpoint{4.373733in}{1.523377in}}%
\pgfpathlineto{\pgfqpoint{4.373733in}{1.523377in}}%
\pgfpathlineto{\pgfqpoint{4.373733in}{1.526326in}}%
\pgfpathlineto{\pgfqpoint{4.378274in}{1.526326in}}%
\pgfpathlineto{\pgfqpoint{4.378274in}{1.523377in}}%
\pgfpathmoveto{\pgfqpoint{4.378274in}{1.520428in}}%
\pgfpathlineto{\pgfqpoint{4.378274in}{1.520428in}}%
\pgfpathlineto{\pgfqpoint{4.378274in}{1.523377in}}%
\pgfpathlineto{\pgfqpoint{4.382815in}{1.523377in}}%
\pgfpathlineto{\pgfqpoint{4.382815in}{1.520428in}}%
\pgfpathmoveto{\pgfqpoint{4.378274in}{1.523377in}}%
\pgfpathlineto{\pgfqpoint{4.378274in}{1.523377in}}%
\pgfpathlineto{\pgfqpoint{4.378274in}{1.526326in}}%
\pgfpathlineto{\pgfqpoint{4.382815in}{1.526326in}}%
\pgfpathlineto{\pgfqpoint{4.382815in}{1.523377in}}%
\pgfpathmoveto{\pgfqpoint{4.369192in}{1.529276in}}%
\pgfpathlineto{\pgfqpoint{4.369192in}{1.529276in}}%
\pgfpathlineto{\pgfqpoint{4.369192in}{1.532225in}}%
\pgfpathlineto{\pgfqpoint{4.373733in}{1.532225in}}%
\pgfpathlineto{\pgfqpoint{4.373733in}{1.529276in}}%
\pgfpathmoveto{\pgfqpoint{4.364651in}{1.535174in}}%
\pgfpathlineto{\pgfqpoint{4.364651in}{1.535174in}}%
\pgfpathlineto{\pgfqpoint{4.364651in}{1.538123in}}%
\pgfpathlineto{\pgfqpoint{4.369192in}{1.538123in}}%
\pgfpathlineto{\pgfqpoint{4.369192in}{1.535174in}}%
\pgfpathmoveto{\pgfqpoint{4.369192in}{1.532225in}}%
\pgfpathlineto{\pgfqpoint{4.369192in}{1.532225in}}%
\pgfpathlineto{\pgfqpoint{4.369192in}{1.535174in}}%
\pgfpathlineto{\pgfqpoint{4.373733in}{1.535174in}}%
\pgfpathlineto{\pgfqpoint{4.373733in}{1.532225in}}%
\pgfpathmoveto{\pgfqpoint{4.369192in}{1.535174in}}%
\pgfpathlineto{\pgfqpoint{4.369192in}{1.535174in}}%
\pgfpathlineto{\pgfqpoint{4.369192in}{1.538123in}}%
\pgfpathlineto{\pgfqpoint{4.373733in}{1.538123in}}%
\pgfpathlineto{\pgfqpoint{4.373733in}{1.535174in}}%
\pgfpathmoveto{\pgfqpoint{4.373733in}{1.526326in}}%
\pgfpathlineto{\pgfqpoint{4.373733in}{1.526326in}}%
\pgfpathlineto{\pgfqpoint{4.373733in}{1.529276in}}%
\pgfpathlineto{\pgfqpoint{4.378274in}{1.529276in}}%
\pgfpathlineto{\pgfqpoint{4.378274in}{1.526326in}}%
\pgfpathmoveto{\pgfqpoint{4.373733in}{1.529276in}}%
\pgfpathlineto{\pgfqpoint{4.373733in}{1.529276in}}%
\pgfpathlineto{\pgfqpoint{4.373733in}{1.532225in}}%
\pgfpathlineto{\pgfqpoint{4.378274in}{1.532225in}}%
\pgfpathlineto{\pgfqpoint{4.378274in}{1.529276in}}%
\pgfpathmoveto{\pgfqpoint{4.305617in}{1.611857in}}%
\pgfpathlineto{\pgfqpoint{4.305617in}{1.611857in}}%
\pgfpathlineto{\pgfqpoint{4.305617in}{1.614806in}}%
\pgfpathlineto{\pgfqpoint{4.310158in}{1.614806in}}%
\pgfpathlineto{\pgfqpoint{4.310158in}{1.611857in}}%
\pgfpathmoveto{\pgfqpoint{4.301076in}{1.617756in}}%
\pgfpathlineto{\pgfqpoint{4.301076in}{1.617756in}}%
\pgfpathlineto{\pgfqpoint{4.301076in}{1.620705in}}%
\pgfpathlineto{\pgfqpoint{4.305617in}{1.620705in}}%
\pgfpathlineto{\pgfqpoint{4.305617in}{1.617756in}}%
\pgfpathmoveto{\pgfqpoint{4.305617in}{1.614806in}}%
\pgfpathlineto{\pgfqpoint{4.305617in}{1.614806in}}%
\pgfpathlineto{\pgfqpoint{4.305617in}{1.617756in}}%
\pgfpathlineto{\pgfqpoint{4.310158in}{1.617756in}}%
\pgfpathlineto{\pgfqpoint{4.310158in}{1.614806in}}%
\pgfpathmoveto{\pgfqpoint{4.305617in}{1.617756in}}%
\pgfpathlineto{\pgfqpoint{4.305617in}{1.617756in}}%
\pgfpathlineto{\pgfqpoint{4.305617in}{1.620705in}}%
\pgfpathlineto{\pgfqpoint{4.310158in}{1.620705in}}%
\pgfpathlineto{\pgfqpoint{4.310158in}{1.617756in}}%
\pgfpathmoveto{\pgfqpoint{4.296535in}{1.623654in}}%
\pgfpathlineto{\pgfqpoint{4.296535in}{1.623654in}}%
\pgfpathlineto{\pgfqpoint{4.296535in}{1.626604in}}%
\pgfpathlineto{\pgfqpoint{4.301076in}{1.626604in}}%
\pgfpathlineto{\pgfqpoint{4.301076in}{1.623654in}}%
\pgfpathmoveto{\pgfqpoint{4.291994in}{1.629553in}}%
\pgfpathlineto{\pgfqpoint{4.291994in}{1.629553in}}%
\pgfpathlineto{\pgfqpoint{4.291994in}{1.632503in}}%
\pgfpathlineto{\pgfqpoint{4.296535in}{1.632503in}}%
\pgfpathlineto{\pgfqpoint{4.296535in}{1.629553in}}%
\pgfpathmoveto{\pgfqpoint{4.296535in}{1.626604in}}%
\pgfpathlineto{\pgfqpoint{4.296535in}{1.626604in}}%
\pgfpathlineto{\pgfqpoint{4.296535in}{1.629553in}}%
\pgfpathlineto{\pgfqpoint{4.301076in}{1.629553in}}%
\pgfpathlineto{\pgfqpoint{4.301076in}{1.626604in}}%
\pgfpathmoveto{\pgfqpoint{4.296535in}{1.629553in}}%
\pgfpathlineto{\pgfqpoint{4.296535in}{1.629553in}}%
\pgfpathlineto{\pgfqpoint{4.296535in}{1.632503in}}%
\pgfpathlineto{\pgfqpoint{4.301076in}{1.632503in}}%
\pgfpathlineto{\pgfqpoint{4.301076in}{1.629553in}}%
\pgfpathmoveto{\pgfqpoint{4.301076in}{1.620705in}}%
\pgfpathlineto{\pgfqpoint{4.301076in}{1.620705in}}%
\pgfpathlineto{\pgfqpoint{4.301076in}{1.623654in}}%
\pgfpathlineto{\pgfqpoint{4.305617in}{1.623654in}}%
\pgfpathlineto{\pgfqpoint{4.305617in}{1.620705in}}%
\pgfpathmoveto{\pgfqpoint{4.301076in}{1.623654in}}%
\pgfpathlineto{\pgfqpoint{4.301076in}{1.623654in}}%
\pgfpathlineto{\pgfqpoint{4.301076in}{1.626604in}}%
\pgfpathlineto{\pgfqpoint{4.305617in}{1.626604in}}%
\pgfpathlineto{\pgfqpoint{4.305617in}{1.623654in}}%
\pgfpathmoveto{\pgfqpoint{4.341945in}{1.564667in}}%
\pgfpathlineto{\pgfqpoint{4.341945in}{1.564667in}}%
\pgfpathlineto{\pgfqpoint{4.341945in}{1.567617in}}%
\pgfpathlineto{\pgfqpoint{4.346486in}{1.567617in}}%
\pgfpathlineto{\pgfqpoint{4.346486in}{1.564667in}}%
\pgfpathmoveto{\pgfqpoint{4.337404in}{1.570566in}}%
\pgfpathlineto{\pgfqpoint{4.337404in}{1.570566in}}%
\pgfpathlineto{\pgfqpoint{4.337404in}{1.573515in}}%
\pgfpathlineto{\pgfqpoint{4.341945in}{1.573515in}}%
\pgfpathlineto{\pgfqpoint{4.341945in}{1.570566in}}%
\pgfpathmoveto{\pgfqpoint{4.341945in}{1.567617in}}%
\pgfpathlineto{\pgfqpoint{4.341945in}{1.567617in}}%
\pgfpathlineto{\pgfqpoint{4.341945in}{1.570566in}}%
\pgfpathlineto{\pgfqpoint{4.346486in}{1.570566in}}%
\pgfpathlineto{\pgfqpoint{4.346486in}{1.567617in}}%
\pgfpathmoveto{\pgfqpoint{4.341945in}{1.570566in}}%
\pgfpathlineto{\pgfqpoint{4.341945in}{1.570566in}}%
\pgfpathlineto{\pgfqpoint{4.341945in}{1.573515in}}%
\pgfpathlineto{\pgfqpoint{4.346486in}{1.573515in}}%
\pgfpathlineto{\pgfqpoint{4.346486in}{1.570566in}}%
\pgfpathmoveto{\pgfqpoint{4.332863in}{1.576465in}}%
\pgfpathlineto{\pgfqpoint{4.332863in}{1.576465in}}%
\pgfpathlineto{\pgfqpoint{4.332863in}{1.579414in}}%
\pgfpathlineto{\pgfqpoint{4.337404in}{1.579414in}}%
\pgfpathlineto{\pgfqpoint{4.337404in}{1.576465in}}%
\pgfpathmoveto{\pgfqpoint{4.328322in}{1.582363in}}%
\pgfpathlineto{\pgfqpoint{4.328322in}{1.582363in}}%
\pgfpathlineto{\pgfqpoint{4.328322in}{1.585313in}}%
\pgfpathlineto{\pgfqpoint{4.332863in}{1.585313in}}%
\pgfpathlineto{\pgfqpoint{4.332863in}{1.582363in}}%
\pgfpathmoveto{\pgfqpoint{4.332863in}{1.579414in}}%
\pgfpathlineto{\pgfqpoint{4.332863in}{1.579414in}}%
\pgfpathlineto{\pgfqpoint{4.332863in}{1.582363in}}%
\pgfpathlineto{\pgfqpoint{4.337404in}{1.582363in}}%
\pgfpathlineto{\pgfqpoint{4.337404in}{1.579414in}}%
\pgfpathmoveto{\pgfqpoint{4.332863in}{1.582363in}}%
\pgfpathlineto{\pgfqpoint{4.332863in}{1.582363in}}%
\pgfpathlineto{\pgfqpoint{4.332863in}{1.585313in}}%
\pgfpathlineto{\pgfqpoint{4.337404in}{1.585313in}}%
\pgfpathlineto{\pgfqpoint{4.337404in}{1.582363in}}%
\pgfpathmoveto{\pgfqpoint{4.337404in}{1.573515in}}%
\pgfpathlineto{\pgfqpoint{4.337404in}{1.573515in}}%
\pgfpathlineto{\pgfqpoint{4.337404in}{1.576465in}}%
\pgfpathlineto{\pgfqpoint{4.341945in}{1.576465in}}%
\pgfpathlineto{\pgfqpoint{4.341945in}{1.573515in}}%
\pgfpathmoveto{\pgfqpoint{4.337404in}{1.576465in}}%
\pgfpathlineto{\pgfqpoint{4.337404in}{1.576465in}}%
\pgfpathlineto{\pgfqpoint{4.337404in}{1.579414in}}%
\pgfpathlineto{\pgfqpoint{4.341945in}{1.579414in}}%
\pgfpathlineto{\pgfqpoint{4.341945in}{1.576465in}}%
\pgfpathmoveto{\pgfqpoint{4.360110in}{1.541072in}}%
\pgfpathlineto{\pgfqpoint{4.360110in}{1.541072in}}%
\pgfpathlineto{\pgfqpoint{4.360110in}{1.544022in}}%
\pgfpathlineto{\pgfqpoint{4.364651in}{1.544022in}}%
\pgfpathlineto{\pgfqpoint{4.364651in}{1.541072in}}%
\pgfpathmoveto{\pgfqpoint{4.355569in}{1.546971in}}%
\pgfpathlineto{\pgfqpoint{4.355569in}{1.546971in}}%
\pgfpathlineto{\pgfqpoint{4.355569in}{1.549921in}}%
\pgfpathlineto{\pgfqpoint{4.360110in}{1.549921in}}%
\pgfpathlineto{\pgfqpoint{4.360110in}{1.546971in}}%
\pgfpathmoveto{\pgfqpoint{4.360110in}{1.544022in}}%
\pgfpathlineto{\pgfqpoint{4.360110in}{1.544022in}}%
\pgfpathlineto{\pgfqpoint{4.360110in}{1.546971in}}%
\pgfpathlineto{\pgfqpoint{4.364651in}{1.546971in}}%
\pgfpathlineto{\pgfqpoint{4.364651in}{1.544022in}}%
\pgfpathmoveto{\pgfqpoint{4.360110in}{1.546971in}}%
\pgfpathlineto{\pgfqpoint{4.360110in}{1.546971in}}%
\pgfpathlineto{\pgfqpoint{4.360110in}{1.549921in}}%
\pgfpathlineto{\pgfqpoint{4.364651in}{1.549921in}}%
\pgfpathlineto{\pgfqpoint{4.364651in}{1.546971in}}%
\pgfpathmoveto{\pgfqpoint{4.351027in}{1.552870in}}%
\pgfpathlineto{\pgfqpoint{4.351027in}{1.552870in}}%
\pgfpathlineto{\pgfqpoint{4.351027in}{1.555819in}}%
\pgfpathlineto{\pgfqpoint{4.355569in}{1.555819in}}%
\pgfpathlineto{\pgfqpoint{4.355569in}{1.552870in}}%
\pgfpathmoveto{\pgfqpoint{4.346486in}{1.558769in}}%
\pgfpathlineto{\pgfqpoint{4.346486in}{1.558769in}}%
\pgfpathlineto{\pgfqpoint{4.346486in}{1.561718in}}%
\pgfpathlineto{\pgfqpoint{4.351027in}{1.561718in}}%
\pgfpathlineto{\pgfqpoint{4.351027in}{1.558769in}}%
\pgfpathmoveto{\pgfqpoint{4.351027in}{1.555819in}}%
\pgfpathlineto{\pgfqpoint{4.351027in}{1.555819in}}%
\pgfpathlineto{\pgfqpoint{4.351027in}{1.558769in}}%
\pgfpathlineto{\pgfqpoint{4.355569in}{1.558769in}}%
\pgfpathlineto{\pgfqpoint{4.355569in}{1.555819in}}%
\pgfpathmoveto{\pgfqpoint{4.351027in}{1.558769in}}%
\pgfpathlineto{\pgfqpoint{4.351027in}{1.558769in}}%
\pgfpathlineto{\pgfqpoint{4.351027in}{1.561718in}}%
\pgfpathlineto{\pgfqpoint{4.355569in}{1.561718in}}%
\pgfpathlineto{\pgfqpoint{4.355569in}{1.558769in}}%
\pgfpathmoveto{\pgfqpoint{4.355569in}{1.549921in}}%
\pgfpathlineto{\pgfqpoint{4.355569in}{1.549921in}}%
\pgfpathlineto{\pgfqpoint{4.355569in}{1.552870in}}%
\pgfpathlineto{\pgfqpoint{4.360110in}{1.552870in}}%
\pgfpathlineto{\pgfqpoint{4.360110in}{1.549921in}}%
\pgfpathmoveto{\pgfqpoint{4.355569in}{1.552870in}}%
\pgfpathlineto{\pgfqpoint{4.355569in}{1.552870in}}%
\pgfpathlineto{\pgfqpoint{4.355569in}{1.555819in}}%
\pgfpathlineto{\pgfqpoint{4.360110in}{1.555819in}}%
\pgfpathlineto{\pgfqpoint{4.360110in}{1.552870in}}%
\pgfpathmoveto{\pgfqpoint{4.364651in}{1.538123in}}%
\pgfpathlineto{\pgfqpoint{4.364651in}{1.538123in}}%
\pgfpathlineto{\pgfqpoint{4.364651in}{1.541072in}}%
\pgfpathlineto{\pgfqpoint{4.369192in}{1.541072in}}%
\pgfpathlineto{\pgfqpoint{4.369192in}{1.538123in}}%
\pgfpathmoveto{\pgfqpoint{4.364651in}{1.541072in}}%
\pgfpathlineto{\pgfqpoint{4.364651in}{1.541072in}}%
\pgfpathlineto{\pgfqpoint{4.364651in}{1.544022in}}%
\pgfpathlineto{\pgfqpoint{4.369192in}{1.544022in}}%
\pgfpathlineto{\pgfqpoint{4.369192in}{1.541072in}}%
\pgfpathmoveto{\pgfqpoint{4.346486in}{1.561718in}}%
\pgfpathlineto{\pgfqpoint{4.346486in}{1.561718in}}%
\pgfpathlineto{\pgfqpoint{4.346486in}{1.564667in}}%
\pgfpathlineto{\pgfqpoint{4.351027in}{1.564667in}}%
\pgfpathlineto{\pgfqpoint{4.351027in}{1.561718in}}%
\pgfpathmoveto{\pgfqpoint{4.346486in}{1.564667in}}%
\pgfpathlineto{\pgfqpoint{4.346486in}{1.564667in}}%
\pgfpathlineto{\pgfqpoint{4.346486in}{1.567617in}}%
\pgfpathlineto{\pgfqpoint{4.351027in}{1.567617in}}%
\pgfpathlineto{\pgfqpoint{4.351027in}{1.564667in}}%
\pgfpathmoveto{\pgfqpoint{4.323781in}{1.588262in}}%
\pgfpathlineto{\pgfqpoint{4.323781in}{1.588262in}}%
\pgfpathlineto{\pgfqpoint{4.323781in}{1.591212in}}%
\pgfpathlineto{\pgfqpoint{4.328322in}{1.591212in}}%
\pgfpathlineto{\pgfqpoint{4.328322in}{1.588262in}}%
\pgfpathmoveto{\pgfqpoint{4.319240in}{1.594161in}}%
\pgfpathlineto{\pgfqpoint{4.319240in}{1.594161in}}%
\pgfpathlineto{\pgfqpoint{4.319240in}{1.597110in}}%
\pgfpathlineto{\pgfqpoint{4.323781in}{1.597110in}}%
\pgfpathlineto{\pgfqpoint{4.323781in}{1.594161in}}%
\pgfpathmoveto{\pgfqpoint{4.323781in}{1.591212in}}%
\pgfpathlineto{\pgfqpoint{4.323781in}{1.591212in}}%
\pgfpathlineto{\pgfqpoint{4.323781in}{1.594161in}}%
\pgfpathlineto{\pgfqpoint{4.328322in}{1.594161in}}%
\pgfpathlineto{\pgfqpoint{4.328322in}{1.591212in}}%
\pgfpathmoveto{\pgfqpoint{4.323781in}{1.594161in}}%
\pgfpathlineto{\pgfqpoint{4.323781in}{1.594161in}}%
\pgfpathlineto{\pgfqpoint{4.323781in}{1.597110in}}%
\pgfpathlineto{\pgfqpoint{4.328322in}{1.597110in}}%
\pgfpathlineto{\pgfqpoint{4.328322in}{1.594161in}}%
\pgfpathmoveto{\pgfqpoint{4.314699in}{1.600060in}}%
\pgfpathlineto{\pgfqpoint{4.314699in}{1.600060in}}%
\pgfpathlineto{\pgfqpoint{4.314699in}{1.603009in}}%
\pgfpathlineto{\pgfqpoint{4.319240in}{1.603009in}}%
\pgfpathlineto{\pgfqpoint{4.319240in}{1.600060in}}%
\pgfpathmoveto{\pgfqpoint{4.310158in}{1.605958in}}%
\pgfpathlineto{\pgfqpoint{4.310158in}{1.605958in}}%
\pgfpathlineto{\pgfqpoint{4.310158in}{1.608908in}}%
\pgfpathlineto{\pgfqpoint{4.314699in}{1.608908in}}%
\pgfpathlineto{\pgfqpoint{4.314699in}{1.605958in}}%
\pgfpathmoveto{\pgfqpoint{4.314699in}{1.603009in}}%
\pgfpathlineto{\pgfqpoint{4.314699in}{1.603009in}}%
\pgfpathlineto{\pgfqpoint{4.314699in}{1.605958in}}%
\pgfpathlineto{\pgfqpoint{4.319240in}{1.605958in}}%
\pgfpathlineto{\pgfqpoint{4.319240in}{1.603009in}}%
\pgfpathmoveto{\pgfqpoint{4.314699in}{1.605958in}}%
\pgfpathlineto{\pgfqpoint{4.314699in}{1.605958in}}%
\pgfpathlineto{\pgfqpoint{4.314699in}{1.608908in}}%
\pgfpathlineto{\pgfqpoint{4.319240in}{1.608908in}}%
\pgfpathlineto{\pgfqpoint{4.319240in}{1.605958in}}%
\pgfpathmoveto{\pgfqpoint{4.319240in}{1.597110in}}%
\pgfpathlineto{\pgfqpoint{4.319240in}{1.597110in}}%
\pgfpathlineto{\pgfqpoint{4.319240in}{1.600060in}}%
\pgfpathlineto{\pgfqpoint{4.323781in}{1.600060in}}%
\pgfpathlineto{\pgfqpoint{4.323781in}{1.597110in}}%
\pgfpathmoveto{\pgfqpoint{4.319240in}{1.600060in}}%
\pgfpathlineto{\pgfqpoint{4.319240in}{1.600060in}}%
\pgfpathlineto{\pgfqpoint{4.319240in}{1.603009in}}%
\pgfpathlineto{\pgfqpoint{4.323781in}{1.603009in}}%
\pgfpathlineto{\pgfqpoint{4.323781in}{1.600060in}}%
\pgfpathmoveto{\pgfqpoint{4.328322in}{1.585313in}}%
\pgfpathlineto{\pgfqpoint{4.328322in}{1.585313in}}%
\pgfpathlineto{\pgfqpoint{4.328322in}{1.588262in}}%
\pgfpathlineto{\pgfqpoint{4.332863in}{1.588262in}}%
\pgfpathlineto{\pgfqpoint{4.332863in}{1.585313in}}%
\pgfpathmoveto{\pgfqpoint{4.328322in}{1.588262in}}%
\pgfpathlineto{\pgfqpoint{4.328322in}{1.588262in}}%
\pgfpathlineto{\pgfqpoint{4.328322in}{1.591212in}}%
\pgfpathlineto{\pgfqpoint{4.332863in}{1.591212in}}%
\pgfpathlineto{\pgfqpoint{4.332863in}{1.588262in}}%
\pgfpathmoveto{\pgfqpoint{4.310158in}{1.608908in}}%
\pgfpathlineto{\pgfqpoint{4.310158in}{1.608908in}}%
\pgfpathlineto{\pgfqpoint{4.310158in}{1.611857in}}%
\pgfpathlineto{\pgfqpoint{4.314699in}{1.611857in}}%
\pgfpathlineto{\pgfqpoint{4.314699in}{1.608908in}}%
\pgfpathmoveto{\pgfqpoint{4.310158in}{1.611857in}}%
\pgfpathlineto{\pgfqpoint{4.310158in}{1.611857in}}%
\pgfpathlineto{\pgfqpoint{4.310158in}{1.614806in}}%
\pgfpathlineto{\pgfqpoint{4.314699in}{1.614806in}}%
\pgfpathlineto{\pgfqpoint{4.314699in}{1.611857in}}%
\pgfpathmoveto{\pgfqpoint{4.269288in}{1.659044in}}%
\pgfpathlineto{\pgfqpoint{4.269288in}{1.659044in}}%
\pgfpathlineto{\pgfqpoint{4.269288in}{1.661993in}}%
\pgfpathlineto{\pgfqpoint{4.273829in}{1.661993in}}%
\pgfpathlineto{\pgfqpoint{4.273829in}{1.659044in}}%
\pgfpathmoveto{\pgfqpoint{4.264747in}{1.664942in}}%
\pgfpathlineto{\pgfqpoint{4.264747in}{1.664942in}}%
\pgfpathlineto{\pgfqpoint{4.264747in}{1.667891in}}%
\pgfpathlineto{\pgfqpoint{4.269288in}{1.667891in}}%
\pgfpathlineto{\pgfqpoint{4.269288in}{1.664942in}}%
\pgfpathmoveto{\pgfqpoint{4.269288in}{1.661993in}}%
\pgfpathlineto{\pgfqpoint{4.269288in}{1.661993in}}%
\pgfpathlineto{\pgfqpoint{4.269288in}{1.664942in}}%
\pgfpathlineto{\pgfqpoint{4.273829in}{1.664942in}}%
\pgfpathlineto{\pgfqpoint{4.273829in}{1.661993in}}%
\pgfpathmoveto{\pgfqpoint{4.269288in}{1.664942in}}%
\pgfpathlineto{\pgfqpoint{4.269288in}{1.664942in}}%
\pgfpathlineto{\pgfqpoint{4.269288in}{1.667891in}}%
\pgfpathlineto{\pgfqpoint{4.273829in}{1.667891in}}%
\pgfpathlineto{\pgfqpoint{4.273829in}{1.664942in}}%
\pgfpathmoveto{\pgfqpoint{4.260206in}{1.670840in}}%
\pgfpathlineto{\pgfqpoint{4.260206in}{1.670840in}}%
\pgfpathlineto{\pgfqpoint{4.260206in}{1.673789in}}%
\pgfpathlineto{\pgfqpoint{4.264747in}{1.673789in}}%
\pgfpathlineto{\pgfqpoint{4.264747in}{1.670840in}}%
\pgfpathmoveto{\pgfqpoint{4.255665in}{1.676738in}}%
\pgfpathlineto{\pgfqpoint{4.255665in}{1.676738in}}%
\pgfpathlineto{\pgfqpoint{4.255665in}{1.679687in}}%
\pgfpathlineto{\pgfqpoint{4.260206in}{1.679687in}}%
\pgfpathlineto{\pgfqpoint{4.260206in}{1.676738in}}%
\pgfpathmoveto{\pgfqpoint{4.260206in}{1.673789in}}%
\pgfpathlineto{\pgfqpoint{4.260206in}{1.673789in}}%
\pgfpathlineto{\pgfqpoint{4.260206in}{1.676738in}}%
\pgfpathlineto{\pgfqpoint{4.264747in}{1.676738in}}%
\pgfpathlineto{\pgfqpoint{4.264747in}{1.673789in}}%
\pgfpathmoveto{\pgfqpoint{4.260206in}{1.676738in}}%
\pgfpathlineto{\pgfqpoint{4.260206in}{1.676738in}}%
\pgfpathlineto{\pgfqpoint{4.260206in}{1.679687in}}%
\pgfpathlineto{\pgfqpoint{4.264747in}{1.679687in}}%
\pgfpathlineto{\pgfqpoint{4.264747in}{1.676738in}}%
\pgfpathmoveto{\pgfqpoint{4.264747in}{1.667891in}}%
\pgfpathlineto{\pgfqpoint{4.264747in}{1.667891in}}%
\pgfpathlineto{\pgfqpoint{4.264747in}{1.670840in}}%
\pgfpathlineto{\pgfqpoint{4.269288in}{1.670840in}}%
\pgfpathlineto{\pgfqpoint{4.269288in}{1.667891in}}%
\pgfpathmoveto{\pgfqpoint{4.264747in}{1.670840in}}%
\pgfpathlineto{\pgfqpoint{4.264747in}{1.670840in}}%
\pgfpathlineto{\pgfqpoint{4.264747in}{1.673789in}}%
\pgfpathlineto{\pgfqpoint{4.269288in}{1.673789in}}%
\pgfpathlineto{\pgfqpoint{4.269288in}{1.670840in}}%
\pgfpathmoveto{\pgfqpoint{4.287453in}{1.635452in}}%
\pgfpathlineto{\pgfqpoint{4.287453in}{1.635452in}}%
\pgfpathlineto{\pgfqpoint{4.287453in}{1.638401in}}%
\pgfpathlineto{\pgfqpoint{4.291994in}{1.638401in}}%
\pgfpathlineto{\pgfqpoint{4.291994in}{1.635452in}}%
\pgfpathmoveto{\pgfqpoint{4.282912in}{1.641350in}}%
\pgfpathlineto{\pgfqpoint{4.282912in}{1.641350in}}%
\pgfpathlineto{\pgfqpoint{4.282912in}{1.644299in}}%
\pgfpathlineto{\pgfqpoint{4.287453in}{1.644299in}}%
\pgfpathlineto{\pgfqpoint{4.287453in}{1.641350in}}%
\pgfpathmoveto{\pgfqpoint{4.287453in}{1.638401in}}%
\pgfpathlineto{\pgfqpoint{4.287453in}{1.638401in}}%
\pgfpathlineto{\pgfqpoint{4.287453in}{1.641350in}}%
\pgfpathlineto{\pgfqpoint{4.291994in}{1.641350in}}%
\pgfpathlineto{\pgfqpoint{4.291994in}{1.638401in}}%
\pgfpathmoveto{\pgfqpoint{4.287453in}{1.641350in}}%
\pgfpathlineto{\pgfqpoint{4.287453in}{1.641350in}}%
\pgfpathlineto{\pgfqpoint{4.287453in}{1.644299in}}%
\pgfpathlineto{\pgfqpoint{4.291994in}{1.644299in}}%
\pgfpathlineto{\pgfqpoint{4.291994in}{1.641350in}}%
\pgfpathmoveto{\pgfqpoint{4.278370in}{1.647248in}}%
\pgfpathlineto{\pgfqpoint{4.278370in}{1.647248in}}%
\pgfpathlineto{\pgfqpoint{4.278370in}{1.650197in}}%
\pgfpathlineto{\pgfqpoint{4.282912in}{1.650197in}}%
\pgfpathlineto{\pgfqpoint{4.282912in}{1.647248in}}%
\pgfpathmoveto{\pgfqpoint{4.273829in}{1.653146in}}%
\pgfpathlineto{\pgfqpoint{4.273829in}{1.653146in}}%
\pgfpathlineto{\pgfqpoint{4.273829in}{1.656095in}}%
\pgfpathlineto{\pgfqpoint{4.278370in}{1.656095in}}%
\pgfpathlineto{\pgfqpoint{4.278370in}{1.653146in}}%
\pgfpathmoveto{\pgfqpoint{4.278370in}{1.650197in}}%
\pgfpathlineto{\pgfqpoint{4.278370in}{1.650197in}}%
\pgfpathlineto{\pgfqpoint{4.278370in}{1.653146in}}%
\pgfpathlineto{\pgfqpoint{4.282912in}{1.653146in}}%
\pgfpathlineto{\pgfqpoint{4.282912in}{1.650197in}}%
\pgfpathmoveto{\pgfqpoint{4.278370in}{1.653146in}}%
\pgfpathlineto{\pgfqpoint{4.278370in}{1.653146in}}%
\pgfpathlineto{\pgfqpoint{4.278370in}{1.656095in}}%
\pgfpathlineto{\pgfqpoint{4.282912in}{1.656095in}}%
\pgfpathlineto{\pgfqpoint{4.282912in}{1.653146in}}%
\pgfpathmoveto{\pgfqpoint{4.282912in}{1.644299in}}%
\pgfpathlineto{\pgfqpoint{4.282912in}{1.644299in}}%
\pgfpathlineto{\pgfqpoint{4.282912in}{1.647248in}}%
\pgfpathlineto{\pgfqpoint{4.287453in}{1.647248in}}%
\pgfpathlineto{\pgfqpoint{4.287453in}{1.644299in}}%
\pgfpathmoveto{\pgfqpoint{4.282912in}{1.647248in}}%
\pgfpathlineto{\pgfqpoint{4.282912in}{1.647248in}}%
\pgfpathlineto{\pgfqpoint{4.282912in}{1.650197in}}%
\pgfpathlineto{\pgfqpoint{4.287453in}{1.650197in}}%
\pgfpathlineto{\pgfqpoint{4.287453in}{1.647248in}}%
\pgfpathmoveto{\pgfqpoint{4.291994in}{1.632503in}}%
\pgfpathlineto{\pgfqpoint{4.291994in}{1.632503in}}%
\pgfpathlineto{\pgfqpoint{4.291994in}{1.635452in}}%
\pgfpathlineto{\pgfqpoint{4.296535in}{1.635452in}}%
\pgfpathlineto{\pgfqpoint{4.296535in}{1.632503in}}%
\pgfpathmoveto{\pgfqpoint{4.291994in}{1.635452in}}%
\pgfpathlineto{\pgfqpoint{4.291994in}{1.635452in}}%
\pgfpathlineto{\pgfqpoint{4.291994in}{1.638401in}}%
\pgfpathlineto{\pgfqpoint{4.296535in}{1.638401in}}%
\pgfpathlineto{\pgfqpoint{4.296535in}{1.635452in}}%
\pgfpathmoveto{\pgfqpoint{4.273829in}{1.656095in}}%
\pgfpathlineto{\pgfqpoint{4.273829in}{1.656095in}}%
\pgfpathlineto{\pgfqpoint{4.273829in}{1.659044in}}%
\pgfpathlineto{\pgfqpoint{4.278370in}{1.659044in}}%
\pgfpathlineto{\pgfqpoint{4.278370in}{1.656095in}}%
\pgfpathmoveto{\pgfqpoint{4.273829in}{1.659044in}}%
\pgfpathlineto{\pgfqpoint{4.273829in}{1.659044in}}%
\pgfpathlineto{\pgfqpoint{4.273829in}{1.661993in}}%
\pgfpathlineto{\pgfqpoint{4.278370in}{1.661993in}}%
\pgfpathlineto{\pgfqpoint{4.278370in}{1.659044in}}%
\pgfpathmoveto{\pgfqpoint{4.251124in}{1.682636in}}%
\pgfpathlineto{\pgfqpoint{4.251124in}{1.682636in}}%
\pgfpathlineto{\pgfqpoint{4.251124in}{1.685585in}}%
\pgfpathlineto{\pgfqpoint{4.255665in}{1.685585in}}%
\pgfpathlineto{\pgfqpoint{4.255665in}{1.682636in}}%
\pgfpathmoveto{\pgfqpoint{4.246583in}{1.688535in}}%
\pgfpathlineto{\pgfqpoint{4.246583in}{1.688535in}}%
\pgfpathlineto{\pgfqpoint{4.246583in}{1.691484in}}%
\pgfpathlineto{\pgfqpoint{4.251124in}{1.691484in}}%
\pgfpathlineto{\pgfqpoint{4.251124in}{1.688535in}}%
\pgfpathmoveto{\pgfqpoint{4.251124in}{1.685585in}}%
\pgfpathlineto{\pgfqpoint{4.251124in}{1.685585in}}%
\pgfpathlineto{\pgfqpoint{4.251124in}{1.688535in}}%
\pgfpathlineto{\pgfqpoint{4.255665in}{1.688535in}}%
\pgfpathlineto{\pgfqpoint{4.255665in}{1.685585in}}%
\pgfpathmoveto{\pgfqpoint{4.251124in}{1.688535in}}%
\pgfpathlineto{\pgfqpoint{4.251124in}{1.688535in}}%
\pgfpathlineto{\pgfqpoint{4.251124in}{1.691484in}}%
\pgfpathlineto{\pgfqpoint{4.255665in}{1.691484in}}%
\pgfpathlineto{\pgfqpoint{4.255665in}{1.688535in}}%
\pgfpathmoveto{\pgfqpoint{4.242042in}{1.694433in}}%
\pgfpathlineto{\pgfqpoint{4.242042in}{1.694433in}}%
\pgfpathlineto{\pgfqpoint{4.242042in}{1.697382in}}%
\pgfpathlineto{\pgfqpoint{4.246583in}{1.697382in}}%
\pgfpathlineto{\pgfqpoint{4.246583in}{1.694433in}}%
\pgfpathmoveto{\pgfqpoint{4.237501in}{1.700331in}}%
\pgfpathlineto{\pgfqpoint{4.237501in}{1.700331in}}%
\pgfpathlineto{\pgfqpoint{4.237501in}{1.703280in}}%
\pgfpathlineto{\pgfqpoint{4.242042in}{1.703280in}}%
\pgfpathlineto{\pgfqpoint{4.242042in}{1.700331in}}%
\pgfpathmoveto{\pgfqpoint{4.242042in}{1.697382in}}%
\pgfpathlineto{\pgfqpoint{4.242042in}{1.697382in}}%
\pgfpathlineto{\pgfqpoint{4.242042in}{1.700331in}}%
\pgfpathlineto{\pgfqpoint{4.246583in}{1.700331in}}%
\pgfpathlineto{\pgfqpoint{4.246583in}{1.697382in}}%
\pgfpathmoveto{\pgfqpoint{4.242042in}{1.700331in}}%
\pgfpathlineto{\pgfqpoint{4.242042in}{1.700331in}}%
\pgfpathlineto{\pgfqpoint{4.242042in}{1.703280in}}%
\pgfpathlineto{\pgfqpoint{4.246583in}{1.703280in}}%
\pgfpathlineto{\pgfqpoint{4.246583in}{1.700331in}}%
\pgfpathmoveto{\pgfqpoint{4.246583in}{1.691484in}}%
\pgfpathlineto{\pgfqpoint{4.246583in}{1.691484in}}%
\pgfpathlineto{\pgfqpoint{4.246583in}{1.694433in}}%
\pgfpathlineto{\pgfqpoint{4.251124in}{1.694433in}}%
\pgfpathlineto{\pgfqpoint{4.251124in}{1.691484in}}%
\pgfpathmoveto{\pgfqpoint{4.246583in}{1.694433in}}%
\pgfpathlineto{\pgfqpoint{4.246583in}{1.694433in}}%
\pgfpathlineto{\pgfqpoint{4.246583in}{1.697382in}}%
\pgfpathlineto{\pgfqpoint{4.251124in}{1.697382in}}%
\pgfpathlineto{\pgfqpoint{4.251124in}{1.694433in}}%
\pgfpathmoveto{\pgfqpoint{4.255665in}{1.679687in}}%
\pgfpathlineto{\pgfqpoint{4.255665in}{1.679687in}}%
\pgfpathlineto{\pgfqpoint{4.255665in}{1.682636in}}%
\pgfpathlineto{\pgfqpoint{4.260206in}{1.682636in}}%
\pgfpathlineto{\pgfqpoint{4.260206in}{1.679687in}}%
\pgfpathmoveto{\pgfqpoint{4.255665in}{1.682636in}}%
\pgfpathlineto{\pgfqpoint{4.255665in}{1.682636in}}%
\pgfpathlineto{\pgfqpoint{4.255665in}{1.685585in}}%
\pgfpathlineto{\pgfqpoint{4.260206in}{1.685585in}}%
\pgfpathlineto{\pgfqpoint{4.260206in}{1.682636in}}%
\pgfpathmoveto{\pgfqpoint{4.237501in}{1.703280in}}%
\pgfpathlineto{\pgfqpoint{4.237501in}{1.703280in}}%
\pgfpathlineto{\pgfqpoint{4.237501in}{1.706229in}}%
\pgfpathlineto{\pgfqpoint{4.242042in}{1.706229in}}%
\pgfpathlineto{\pgfqpoint{4.242042in}{1.703280in}}%
\pgfpathmoveto{\pgfqpoint{4.237501in}{1.706229in}}%
\pgfpathlineto{\pgfqpoint{4.237501in}{1.706229in}}%
\pgfpathlineto{\pgfqpoint{4.237501in}{1.709178in}}%
\pgfpathlineto{\pgfqpoint{4.242042in}{1.709178in}}%
\pgfpathlineto{\pgfqpoint{4.242042in}{1.706229in}}%
\pgfpathmoveto{\pgfqpoint{4.523582in}{1.328729in}}%
\pgfpathlineto{\pgfqpoint{4.523582in}{1.328729in}}%
\pgfpathlineto{\pgfqpoint{4.523582in}{1.331678in}}%
\pgfpathlineto{\pgfqpoint{4.528123in}{1.331678in}}%
\pgfpathlineto{\pgfqpoint{4.528123in}{1.328729in}}%
\pgfpathmoveto{\pgfqpoint{4.519042in}{1.334627in}}%
\pgfpathlineto{\pgfqpoint{4.519042in}{1.334627in}}%
\pgfpathlineto{\pgfqpoint{4.519042in}{1.337576in}}%
\pgfpathlineto{\pgfqpoint{4.523582in}{1.337576in}}%
\pgfpathlineto{\pgfqpoint{4.523582in}{1.334627in}}%
\pgfpathmoveto{\pgfqpoint{4.523582in}{1.331678in}}%
\pgfpathlineto{\pgfqpoint{4.523582in}{1.331678in}}%
\pgfpathlineto{\pgfqpoint{4.523582in}{1.334627in}}%
\pgfpathlineto{\pgfqpoint{4.528123in}{1.334627in}}%
\pgfpathlineto{\pgfqpoint{4.528123in}{1.331678in}}%
\pgfpathmoveto{\pgfqpoint{4.523582in}{1.334627in}}%
\pgfpathlineto{\pgfqpoint{4.523582in}{1.334627in}}%
\pgfpathlineto{\pgfqpoint{4.523582in}{1.337576in}}%
\pgfpathlineto{\pgfqpoint{4.528123in}{1.337576in}}%
\pgfpathlineto{\pgfqpoint{4.528123in}{1.334627in}}%
\pgfpathmoveto{\pgfqpoint{4.514501in}{1.340525in}}%
\pgfpathlineto{\pgfqpoint{4.514501in}{1.340525in}}%
\pgfpathlineto{\pgfqpoint{4.514501in}{1.343474in}}%
\pgfpathlineto{\pgfqpoint{4.519042in}{1.343474in}}%
\pgfpathlineto{\pgfqpoint{4.519042in}{1.340525in}}%
\pgfpathmoveto{\pgfqpoint{4.509960in}{1.346423in}}%
\pgfpathlineto{\pgfqpoint{4.509960in}{1.346423in}}%
\pgfpathlineto{\pgfqpoint{4.509960in}{1.349372in}}%
\pgfpathlineto{\pgfqpoint{4.514501in}{1.349372in}}%
\pgfpathlineto{\pgfqpoint{4.514501in}{1.346423in}}%
\pgfpathmoveto{\pgfqpoint{4.514501in}{1.343474in}}%
\pgfpathlineto{\pgfqpoint{4.514501in}{1.343474in}}%
\pgfpathlineto{\pgfqpoint{4.514501in}{1.346423in}}%
\pgfpathlineto{\pgfqpoint{4.519042in}{1.346423in}}%
\pgfpathlineto{\pgfqpoint{4.519042in}{1.343474in}}%
\pgfpathmoveto{\pgfqpoint{4.514501in}{1.346423in}}%
\pgfpathlineto{\pgfqpoint{4.514501in}{1.346423in}}%
\pgfpathlineto{\pgfqpoint{4.514501in}{1.349372in}}%
\pgfpathlineto{\pgfqpoint{4.519042in}{1.349372in}}%
\pgfpathlineto{\pgfqpoint{4.519042in}{1.346423in}}%
\pgfpathmoveto{\pgfqpoint{4.519042in}{1.337576in}}%
\pgfpathlineto{\pgfqpoint{4.519042in}{1.337576in}}%
\pgfpathlineto{\pgfqpoint{4.519042in}{1.340525in}}%
\pgfpathlineto{\pgfqpoint{4.523582in}{1.340525in}}%
\pgfpathlineto{\pgfqpoint{4.523582in}{1.337576in}}%
\pgfpathmoveto{\pgfqpoint{4.519042in}{1.340525in}}%
\pgfpathlineto{\pgfqpoint{4.519042in}{1.340525in}}%
\pgfpathlineto{\pgfqpoint{4.519042in}{1.343474in}}%
\pgfpathlineto{\pgfqpoint{4.523582in}{1.343474in}}%
\pgfpathlineto{\pgfqpoint{4.523582in}{1.340525in}}%
\pgfpathmoveto{\pgfqpoint{4.450928in}{1.423104in}}%
\pgfpathlineto{\pgfqpoint{4.450928in}{1.423104in}}%
\pgfpathlineto{\pgfqpoint{4.450928in}{1.426054in}}%
\pgfpathlineto{\pgfqpoint{4.455469in}{1.426054in}}%
\pgfpathlineto{\pgfqpoint{4.455469in}{1.423104in}}%
\pgfpathmoveto{\pgfqpoint{4.446387in}{1.429003in}}%
\pgfpathlineto{\pgfqpoint{4.446387in}{1.429003in}}%
\pgfpathlineto{\pgfqpoint{4.446387in}{1.431952in}}%
\pgfpathlineto{\pgfqpoint{4.450928in}{1.431952in}}%
\pgfpathlineto{\pgfqpoint{4.450928in}{1.429003in}}%
\pgfpathmoveto{\pgfqpoint{4.450928in}{1.426054in}}%
\pgfpathlineto{\pgfqpoint{4.450928in}{1.426054in}}%
\pgfpathlineto{\pgfqpoint{4.450928in}{1.429003in}}%
\pgfpathlineto{\pgfqpoint{4.455469in}{1.429003in}}%
\pgfpathlineto{\pgfqpoint{4.455469in}{1.426054in}}%
\pgfpathmoveto{\pgfqpoint{4.450928in}{1.429003in}}%
\pgfpathlineto{\pgfqpoint{4.450928in}{1.429003in}}%
\pgfpathlineto{\pgfqpoint{4.450928in}{1.431952in}}%
\pgfpathlineto{\pgfqpoint{4.455469in}{1.431952in}}%
\pgfpathlineto{\pgfqpoint{4.455469in}{1.429003in}}%
\pgfpathmoveto{\pgfqpoint{4.441846in}{1.434902in}}%
\pgfpathlineto{\pgfqpoint{4.441846in}{1.434902in}}%
\pgfpathlineto{\pgfqpoint{4.441846in}{1.437851in}}%
\pgfpathlineto{\pgfqpoint{4.446387in}{1.437851in}}%
\pgfpathlineto{\pgfqpoint{4.446387in}{1.434902in}}%
\pgfpathmoveto{\pgfqpoint{4.437306in}{1.440800in}}%
\pgfpathlineto{\pgfqpoint{4.437306in}{1.440800in}}%
\pgfpathlineto{\pgfqpoint{4.437306in}{1.443749in}}%
\pgfpathlineto{\pgfqpoint{4.441846in}{1.443749in}}%
\pgfpathlineto{\pgfqpoint{4.441846in}{1.440800in}}%
\pgfpathmoveto{\pgfqpoint{4.441846in}{1.437851in}}%
\pgfpathlineto{\pgfqpoint{4.441846in}{1.437851in}}%
\pgfpathlineto{\pgfqpoint{4.441846in}{1.440800in}}%
\pgfpathlineto{\pgfqpoint{4.446387in}{1.440800in}}%
\pgfpathlineto{\pgfqpoint{4.446387in}{1.437851in}}%
\pgfpathmoveto{\pgfqpoint{4.441846in}{1.440800in}}%
\pgfpathlineto{\pgfqpoint{4.441846in}{1.440800in}}%
\pgfpathlineto{\pgfqpoint{4.441846in}{1.443749in}}%
\pgfpathlineto{\pgfqpoint{4.446387in}{1.443749in}}%
\pgfpathlineto{\pgfqpoint{4.446387in}{1.440800in}}%
\pgfpathmoveto{\pgfqpoint{4.446387in}{1.431952in}}%
\pgfpathlineto{\pgfqpoint{4.446387in}{1.431952in}}%
\pgfpathlineto{\pgfqpoint{4.446387in}{1.434902in}}%
\pgfpathlineto{\pgfqpoint{4.450928in}{1.434902in}}%
\pgfpathlineto{\pgfqpoint{4.450928in}{1.431952in}}%
\pgfpathmoveto{\pgfqpoint{4.446387in}{1.434902in}}%
\pgfpathlineto{\pgfqpoint{4.446387in}{1.434902in}}%
\pgfpathlineto{\pgfqpoint{4.446387in}{1.437851in}}%
\pgfpathlineto{\pgfqpoint{4.450928in}{1.437851in}}%
\pgfpathlineto{\pgfqpoint{4.450928in}{1.434902in}}%
\pgfpathmoveto{\pgfqpoint{4.487255in}{1.375916in}}%
\pgfpathlineto{\pgfqpoint{4.487255in}{1.375916in}}%
\pgfpathlineto{\pgfqpoint{4.487255in}{1.378865in}}%
\pgfpathlineto{\pgfqpoint{4.491796in}{1.378865in}}%
\pgfpathlineto{\pgfqpoint{4.491796in}{1.375916in}}%
\pgfpathmoveto{\pgfqpoint{4.482714in}{1.381815in}}%
\pgfpathlineto{\pgfqpoint{4.482714in}{1.381815in}}%
\pgfpathlineto{\pgfqpoint{4.482714in}{1.384764in}}%
\pgfpathlineto{\pgfqpoint{4.487255in}{1.384764in}}%
\pgfpathlineto{\pgfqpoint{4.487255in}{1.381815in}}%
\pgfpathmoveto{\pgfqpoint{4.487255in}{1.378865in}}%
\pgfpathlineto{\pgfqpoint{4.487255in}{1.378865in}}%
\pgfpathlineto{\pgfqpoint{4.487255in}{1.381815in}}%
\pgfpathlineto{\pgfqpoint{4.491796in}{1.381815in}}%
\pgfpathlineto{\pgfqpoint{4.491796in}{1.378865in}}%
\pgfpathmoveto{\pgfqpoint{4.487255in}{1.381815in}}%
\pgfpathlineto{\pgfqpoint{4.487255in}{1.381815in}}%
\pgfpathlineto{\pgfqpoint{4.487255in}{1.384764in}}%
\pgfpathlineto{\pgfqpoint{4.491796in}{1.384764in}}%
\pgfpathlineto{\pgfqpoint{4.491796in}{1.381815in}}%
\pgfpathmoveto{\pgfqpoint{4.478174in}{1.387713in}}%
\pgfpathlineto{\pgfqpoint{4.478174in}{1.387713in}}%
\pgfpathlineto{\pgfqpoint{4.478174in}{1.390662in}}%
\pgfpathlineto{\pgfqpoint{4.482714in}{1.390662in}}%
\pgfpathlineto{\pgfqpoint{4.482714in}{1.387713in}}%
\pgfpathmoveto{\pgfqpoint{4.473633in}{1.393612in}}%
\pgfpathlineto{\pgfqpoint{4.473633in}{1.393612in}}%
\pgfpathlineto{\pgfqpoint{4.473633in}{1.396561in}}%
\pgfpathlineto{\pgfqpoint{4.478174in}{1.396561in}}%
\pgfpathlineto{\pgfqpoint{4.478174in}{1.393612in}}%
\pgfpathmoveto{\pgfqpoint{4.478174in}{1.390662in}}%
\pgfpathlineto{\pgfqpoint{4.478174in}{1.390662in}}%
\pgfpathlineto{\pgfqpoint{4.478174in}{1.393612in}}%
\pgfpathlineto{\pgfqpoint{4.482714in}{1.393612in}}%
\pgfpathlineto{\pgfqpoint{4.482714in}{1.390662in}}%
\pgfpathmoveto{\pgfqpoint{4.478174in}{1.393612in}}%
\pgfpathlineto{\pgfqpoint{4.478174in}{1.393612in}}%
\pgfpathlineto{\pgfqpoint{4.478174in}{1.396561in}}%
\pgfpathlineto{\pgfqpoint{4.482714in}{1.396561in}}%
\pgfpathlineto{\pgfqpoint{4.482714in}{1.393612in}}%
\pgfpathmoveto{\pgfqpoint{4.482714in}{1.384764in}}%
\pgfpathlineto{\pgfqpoint{4.482714in}{1.384764in}}%
\pgfpathlineto{\pgfqpoint{4.482714in}{1.387713in}}%
\pgfpathlineto{\pgfqpoint{4.487255in}{1.387713in}}%
\pgfpathlineto{\pgfqpoint{4.487255in}{1.384764in}}%
\pgfpathmoveto{\pgfqpoint{4.482714in}{1.387713in}}%
\pgfpathlineto{\pgfqpoint{4.482714in}{1.387713in}}%
\pgfpathlineto{\pgfqpoint{4.482714in}{1.390662in}}%
\pgfpathlineto{\pgfqpoint{4.487255in}{1.390662in}}%
\pgfpathlineto{\pgfqpoint{4.487255in}{1.387713in}}%
\pgfpathmoveto{\pgfqpoint{4.505419in}{1.352322in}}%
\pgfpathlineto{\pgfqpoint{4.505419in}{1.352322in}}%
\pgfpathlineto{\pgfqpoint{4.505419in}{1.355271in}}%
\pgfpathlineto{\pgfqpoint{4.509960in}{1.355271in}}%
\pgfpathlineto{\pgfqpoint{4.509960in}{1.352322in}}%
\pgfpathmoveto{\pgfqpoint{4.500878in}{1.358220in}}%
\pgfpathlineto{\pgfqpoint{4.500878in}{1.358220in}}%
\pgfpathlineto{\pgfqpoint{4.500878in}{1.361170in}}%
\pgfpathlineto{\pgfqpoint{4.505419in}{1.361170in}}%
\pgfpathlineto{\pgfqpoint{4.505419in}{1.358220in}}%
\pgfpathmoveto{\pgfqpoint{4.505419in}{1.355271in}}%
\pgfpathlineto{\pgfqpoint{4.505419in}{1.355271in}}%
\pgfpathlineto{\pgfqpoint{4.505419in}{1.358220in}}%
\pgfpathlineto{\pgfqpoint{4.509960in}{1.358220in}}%
\pgfpathlineto{\pgfqpoint{4.509960in}{1.355271in}}%
\pgfpathmoveto{\pgfqpoint{4.505419in}{1.358220in}}%
\pgfpathlineto{\pgfqpoint{4.505419in}{1.358220in}}%
\pgfpathlineto{\pgfqpoint{4.505419in}{1.361170in}}%
\pgfpathlineto{\pgfqpoint{4.509960in}{1.361170in}}%
\pgfpathlineto{\pgfqpoint{4.509960in}{1.358220in}}%
\pgfpathmoveto{\pgfqpoint{4.496337in}{1.364119in}}%
\pgfpathlineto{\pgfqpoint{4.496337in}{1.364119in}}%
\pgfpathlineto{\pgfqpoint{4.496337in}{1.367068in}}%
\pgfpathlineto{\pgfqpoint{4.500878in}{1.367068in}}%
\pgfpathlineto{\pgfqpoint{4.500878in}{1.364119in}}%
\pgfpathmoveto{\pgfqpoint{4.491796in}{1.370017in}}%
\pgfpathlineto{\pgfqpoint{4.491796in}{1.370017in}}%
\pgfpathlineto{\pgfqpoint{4.491796in}{1.372967in}}%
\pgfpathlineto{\pgfqpoint{4.496337in}{1.372967in}}%
\pgfpathlineto{\pgfqpoint{4.496337in}{1.370017in}}%
\pgfpathmoveto{\pgfqpoint{4.496337in}{1.367068in}}%
\pgfpathlineto{\pgfqpoint{4.496337in}{1.367068in}}%
\pgfpathlineto{\pgfqpoint{4.496337in}{1.370017in}}%
\pgfpathlineto{\pgfqpoint{4.500878in}{1.370017in}}%
\pgfpathlineto{\pgfqpoint{4.500878in}{1.367068in}}%
\pgfpathmoveto{\pgfqpoint{4.496337in}{1.370017in}}%
\pgfpathlineto{\pgfqpoint{4.496337in}{1.370017in}}%
\pgfpathlineto{\pgfqpoint{4.496337in}{1.372967in}}%
\pgfpathlineto{\pgfqpoint{4.500878in}{1.372967in}}%
\pgfpathlineto{\pgfqpoint{4.500878in}{1.370017in}}%
\pgfpathmoveto{\pgfqpoint{4.500878in}{1.361170in}}%
\pgfpathlineto{\pgfqpoint{4.500878in}{1.361170in}}%
\pgfpathlineto{\pgfqpoint{4.500878in}{1.364119in}}%
\pgfpathlineto{\pgfqpoint{4.505419in}{1.364119in}}%
\pgfpathlineto{\pgfqpoint{4.505419in}{1.361170in}}%
\pgfpathmoveto{\pgfqpoint{4.500878in}{1.364119in}}%
\pgfpathlineto{\pgfqpoint{4.500878in}{1.364119in}}%
\pgfpathlineto{\pgfqpoint{4.500878in}{1.367068in}}%
\pgfpathlineto{\pgfqpoint{4.505419in}{1.367068in}}%
\pgfpathlineto{\pgfqpoint{4.505419in}{1.364119in}}%
\pgfpathmoveto{\pgfqpoint{4.509960in}{1.349372in}}%
\pgfpathlineto{\pgfqpoint{4.509960in}{1.349372in}}%
\pgfpathlineto{\pgfqpoint{4.509960in}{1.352322in}}%
\pgfpathlineto{\pgfqpoint{4.514501in}{1.352322in}}%
\pgfpathlineto{\pgfqpoint{4.514501in}{1.349372in}}%
\pgfpathmoveto{\pgfqpoint{4.509960in}{1.352322in}}%
\pgfpathlineto{\pgfqpoint{4.509960in}{1.352322in}}%
\pgfpathlineto{\pgfqpoint{4.509960in}{1.355271in}}%
\pgfpathlineto{\pgfqpoint{4.514501in}{1.355271in}}%
\pgfpathlineto{\pgfqpoint{4.514501in}{1.352322in}}%
\pgfpathmoveto{\pgfqpoint{4.491796in}{1.372967in}}%
\pgfpathlineto{\pgfqpoint{4.491796in}{1.372967in}}%
\pgfpathlineto{\pgfqpoint{4.491796in}{1.375916in}}%
\pgfpathlineto{\pgfqpoint{4.496337in}{1.375916in}}%
\pgfpathlineto{\pgfqpoint{4.496337in}{1.372967in}}%
\pgfpathmoveto{\pgfqpoint{4.491796in}{1.375916in}}%
\pgfpathlineto{\pgfqpoint{4.491796in}{1.375916in}}%
\pgfpathlineto{\pgfqpoint{4.491796in}{1.378865in}}%
\pgfpathlineto{\pgfqpoint{4.496337in}{1.378865in}}%
\pgfpathlineto{\pgfqpoint{4.496337in}{1.375916in}}%
\pgfpathmoveto{\pgfqpoint{4.469092in}{1.399510in}}%
\pgfpathlineto{\pgfqpoint{4.469092in}{1.399510in}}%
\pgfpathlineto{\pgfqpoint{4.469092in}{1.402459in}}%
\pgfpathlineto{\pgfqpoint{4.473633in}{1.402459in}}%
\pgfpathlineto{\pgfqpoint{4.473633in}{1.399510in}}%
\pgfpathmoveto{\pgfqpoint{4.464551in}{1.405409in}}%
\pgfpathlineto{\pgfqpoint{4.464551in}{1.405409in}}%
\pgfpathlineto{\pgfqpoint{4.464551in}{1.408358in}}%
\pgfpathlineto{\pgfqpoint{4.469092in}{1.408358in}}%
\pgfpathlineto{\pgfqpoint{4.469092in}{1.405409in}}%
\pgfpathmoveto{\pgfqpoint{4.469092in}{1.402459in}}%
\pgfpathlineto{\pgfqpoint{4.469092in}{1.402459in}}%
\pgfpathlineto{\pgfqpoint{4.469092in}{1.405409in}}%
\pgfpathlineto{\pgfqpoint{4.473633in}{1.405409in}}%
\pgfpathlineto{\pgfqpoint{4.473633in}{1.402459in}}%
\pgfpathmoveto{\pgfqpoint{4.469092in}{1.405409in}}%
\pgfpathlineto{\pgfqpoint{4.469092in}{1.405409in}}%
\pgfpathlineto{\pgfqpoint{4.469092in}{1.408358in}}%
\pgfpathlineto{\pgfqpoint{4.473633in}{1.408358in}}%
\pgfpathlineto{\pgfqpoint{4.473633in}{1.405409in}}%
\pgfpathmoveto{\pgfqpoint{4.460010in}{1.411307in}}%
\pgfpathlineto{\pgfqpoint{4.460010in}{1.411307in}}%
\pgfpathlineto{\pgfqpoint{4.460010in}{1.414257in}}%
\pgfpathlineto{\pgfqpoint{4.464551in}{1.414257in}}%
\pgfpathlineto{\pgfqpoint{4.464551in}{1.411307in}}%
\pgfpathmoveto{\pgfqpoint{4.455469in}{1.417206in}}%
\pgfpathlineto{\pgfqpoint{4.455469in}{1.417206in}}%
\pgfpathlineto{\pgfqpoint{4.455469in}{1.420155in}}%
\pgfpathlineto{\pgfqpoint{4.460010in}{1.420155in}}%
\pgfpathlineto{\pgfqpoint{4.460010in}{1.417206in}}%
\pgfpathmoveto{\pgfqpoint{4.460010in}{1.414257in}}%
\pgfpathlineto{\pgfqpoint{4.460010in}{1.414257in}}%
\pgfpathlineto{\pgfqpoint{4.460010in}{1.417206in}}%
\pgfpathlineto{\pgfqpoint{4.464551in}{1.417206in}}%
\pgfpathlineto{\pgfqpoint{4.464551in}{1.414257in}}%
\pgfpathmoveto{\pgfqpoint{4.460010in}{1.417206in}}%
\pgfpathlineto{\pgfqpoint{4.460010in}{1.417206in}}%
\pgfpathlineto{\pgfqpoint{4.460010in}{1.420155in}}%
\pgfpathlineto{\pgfqpoint{4.464551in}{1.420155in}}%
\pgfpathlineto{\pgfqpoint{4.464551in}{1.417206in}}%
\pgfpathmoveto{\pgfqpoint{4.464551in}{1.408358in}}%
\pgfpathlineto{\pgfqpoint{4.464551in}{1.408358in}}%
\pgfpathlineto{\pgfqpoint{4.464551in}{1.411307in}}%
\pgfpathlineto{\pgfqpoint{4.469092in}{1.411307in}}%
\pgfpathlineto{\pgfqpoint{4.469092in}{1.408358in}}%
\pgfpathmoveto{\pgfqpoint{4.464551in}{1.411307in}}%
\pgfpathlineto{\pgfqpoint{4.464551in}{1.411307in}}%
\pgfpathlineto{\pgfqpoint{4.464551in}{1.414257in}}%
\pgfpathlineto{\pgfqpoint{4.469092in}{1.414257in}}%
\pgfpathlineto{\pgfqpoint{4.469092in}{1.411307in}}%
\pgfpathmoveto{\pgfqpoint{4.473633in}{1.396561in}}%
\pgfpathlineto{\pgfqpoint{4.473633in}{1.396561in}}%
\pgfpathlineto{\pgfqpoint{4.473633in}{1.399510in}}%
\pgfpathlineto{\pgfqpoint{4.478174in}{1.399510in}}%
\pgfpathlineto{\pgfqpoint{4.478174in}{1.396561in}}%
\pgfpathmoveto{\pgfqpoint{4.473633in}{1.399510in}}%
\pgfpathlineto{\pgfqpoint{4.473633in}{1.399510in}}%
\pgfpathlineto{\pgfqpoint{4.473633in}{1.402459in}}%
\pgfpathlineto{\pgfqpoint{4.478174in}{1.402459in}}%
\pgfpathlineto{\pgfqpoint{4.478174in}{1.399510in}}%
\pgfpathmoveto{\pgfqpoint{4.455469in}{1.420155in}}%
\pgfpathlineto{\pgfqpoint{4.455469in}{1.420155in}}%
\pgfpathlineto{\pgfqpoint{4.455469in}{1.423104in}}%
\pgfpathlineto{\pgfqpoint{4.460010in}{1.423104in}}%
\pgfpathlineto{\pgfqpoint{4.460010in}{1.420155in}}%
\pgfpathmoveto{\pgfqpoint{4.455469in}{1.423104in}}%
\pgfpathlineto{\pgfqpoint{4.455469in}{1.423104in}}%
\pgfpathlineto{\pgfqpoint{4.455469in}{1.426054in}}%
\pgfpathlineto{\pgfqpoint{4.460010in}{1.426054in}}%
\pgfpathlineto{\pgfqpoint{4.460010in}{1.423104in}}%
\pgfpathmoveto{\pgfqpoint{4.414601in}{1.470292in}}%
\pgfpathlineto{\pgfqpoint{4.414601in}{1.470292in}}%
\pgfpathlineto{\pgfqpoint{4.414601in}{1.473241in}}%
\pgfpathlineto{\pgfqpoint{4.419142in}{1.473241in}}%
\pgfpathlineto{\pgfqpoint{4.419142in}{1.470292in}}%
\pgfpathmoveto{\pgfqpoint{4.410060in}{1.476190in}}%
\pgfpathlineto{\pgfqpoint{4.410060in}{1.476190in}}%
\pgfpathlineto{\pgfqpoint{4.410060in}{1.479140in}}%
\pgfpathlineto{\pgfqpoint{4.414601in}{1.479140in}}%
\pgfpathlineto{\pgfqpoint{4.414601in}{1.476190in}}%
\pgfpathmoveto{\pgfqpoint{4.414601in}{1.473241in}}%
\pgfpathlineto{\pgfqpoint{4.414601in}{1.473241in}}%
\pgfpathlineto{\pgfqpoint{4.414601in}{1.476190in}}%
\pgfpathlineto{\pgfqpoint{4.419142in}{1.476190in}}%
\pgfpathlineto{\pgfqpoint{4.419142in}{1.473241in}}%
\pgfpathmoveto{\pgfqpoint{4.414601in}{1.476190in}}%
\pgfpathlineto{\pgfqpoint{4.414601in}{1.476190in}}%
\pgfpathlineto{\pgfqpoint{4.414601in}{1.479140in}}%
\pgfpathlineto{\pgfqpoint{4.419142in}{1.479140in}}%
\pgfpathlineto{\pgfqpoint{4.419142in}{1.476190in}}%
\pgfpathmoveto{\pgfqpoint{4.405519in}{1.482089in}}%
\pgfpathlineto{\pgfqpoint{4.405519in}{1.482089in}}%
\pgfpathlineto{\pgfqpoint{4.405519in}{1.485038in}}%
\pgfpathlineto{\pgfqpoint{4.410060in}{1.485038in}}%
\pgfpathlineto{\pgfqpoint{4.410060in}{1.482089in}}%
\pgfpathmoveto{\pgfqpoint{4.400978in}{1.487987in}}%
\pgfpathlineto{\pgfqpoint{4.400978in}{1.487987in}}%
\pgfpathlineto{\pgfqpoint{4.400978in}{1.490936in}}%
\pgfpathlineto{\pgfqpoint{4.405519in}{1.490936in}}%
\pgfpathlineto{\pgfqpoint{4.405519in}{1.487987in}}%
\pgfpathmoveto{\pgfqpoint{4.405519in}{1.485038in}}%
\pgfpathlineto{\pgfqpoint{4.405519in}{1.485038in}}%
\pgfpathlineto{\pgfqpoint{4.405519in}{1.487987in}}%
\pgfpathlineto{\pgfqpoint{4.410060in}{1.487987in}}%
\pgfpathlineto{\pgfqpoint{4.410060in}{1.485038in}}%
\pgfpathmoveto{\pgfqpoint{4.405519in}{1.487987in}}%
\pgfpathlineto{\pgfqpoint{4.405519in}{1.487987in}}%
\pgfpathlineto{\pgfqpoint{4.405519in}{1.490936in}}%
\pgfpathlineto{\pgfqpoint{4.410060in}{1.490936in}}%
\pgfpathlineto{\pgfqpoint{4.410060in}{1.487987in}}%
\pgfpathmoveto{\pgfqpoint{4.410060in}{1.479140in}}%
\pgfpathlineto{\pgfqpoint{4.410060in}{1.479140in}}%
\pgfpathlineto{\pgfqpoint{4.410060in}{1.482089in}}%
\pgfpathlineto{\pgfqpoint{4.414601in}{1.482089in}}%
\pgfpathlineto{\pgfqpoint{4.414601in}{1.479140in}}%
\pgfpathmoveto{\pgfqpoint{4.410060in}{1.482089in}}%
\pgfpathlineto{\pgfqpoint{4.410060in}{1.482089in}}%
\pgfpathlineto{\pgfqpoint{4.410060in}{1.485038in}}%
\pgfpathlineto{\pgfqpoint{4.414601in}{1.485038in}}%
\pgfpathlineto{\pgfqpoint{4.414601in}{1.482089in}}%
\pgfpathmoveto{\pgfqpoint{4.432765in}{1.446699in}}%
\pgfpathlineto{\pgfqpoint{4.432765in}{1.446699in}}%
\pgfpathlineto{\pgfqpoint{4.432765in}{1.449648in}}%
\pgfpathlineto{\pgfqpoint{4.437306in}{1.449648in}}%
\pgfpathlineto{\pgfqpoint{4.437306in}{1.446699in}}%
\pgfpathmoveto{\pgfqpoint{4.428224in}{1.452597in}}%
\pgfpathlineto{\pgfqpoint{4.428224in}{1.452597in}}%
\pgfpathlineto{\pgfqpoint{4.428224in}{1.455546in}}%
\pgfpathlineto{\pgfqpoint{4.432765in}{1.455546in}}%
\pgfpathlineto{\pgfqpoint{4.432765in}{1.452597in}}%
\pgfpathmoveto{\pgfqpoint{4.432765in}{1.449648in}}%
\pgfpathlineto{\pgfqpoint{4.432765in}{1.449648in}}%
\pgfpathlineto{\pgfqpoint{4.432765in}{1.452597in}}%
\pgfpathlineto{\pgfqpoint{4.437306in}{1.452597in}}%
\pgfpathlineto{\pgfqpoint{4.437306in}{1.449648in}}%
\pgfpathmoveto{\pgfqpoint{4.432765in}{1.452597in}}%
\pgfpathlineto{\pgfqpoint{4.432765in}{1.452597in}}%
\pgfpathlineto{\pgfqpoint{4.432765in}{1.455546in}}%
\pgfpathlineto{\pgfqpoint{4.437306in}{1.455546in}}%
\pgfpathlineto{\pgfqpoint{4.437306in}{1.452597in}}%
\pgfpathmoveto{\pgfqpoint{4.423683in}{1.458495in}}%
\pgfpathlineto{\pgfqpoint{4.423683in}{1.458495in}}%
\pgfpathlineto{\pgfqpoint{4.423683in}{1.461444in}}%
\pgfpathlineto{\pgfqpoint{4.428224in}{1.461444in}}%
\pgfpathlineto{\pgfqpoint{4.428224in}{1.458495in}}%
\pgfpathmoveto{\pgfqpoint{4.419142in}{1.464394in}}%
\pgfpathlineto{\pgfqpoint{4.419142in}{1.464394in}}%
\pgfpathlineto{\pgfqpoint{4.419142in}{1.467343in}}%
\pgfpathlineto{\pgfqpoint{4.423683in}{1.467343in}}%
\pgfpathlineto{\pgfqpoint{4.423683in}{1.464394in}}%
\pgfpathmoveto{\pgfqpoint{4.423683in}{1.461444in}}%
\pgfpathlineto{\pgfqpoint{4.423683in}{1.461444in}}%
\pgfpathlineto{\pgfqpoint{4.423683in}{1.464394in}}%
\pgfpathlineto{\pgfqpoint{4.428224in}{1.464394in}}%
\pgfpathlineto{\pgfqpoint{4.428224in}{1.461444in}}%
\pgfpathmoveto{\pgfqpoint{4.423683in}{1.464394in}}%
\pgfpathlineto{\pgfqpoint{4.423683in}{1.464394in}}%
\pgfpathlineto{\pgfqpoint{4.423683in}{1.467343in}}%
\pgfpathlineto{\pgfqpoint{4.428224in}{1.467343in}}%
\pgfpathlineto{\pgfqpoint{4.428224in}{1.464394in}}%
\pgfpathmoveto{\pgfqpoint{4.428224in}{1.455546in}}%
\pgfpathlineto{\pgfqpoint{4.428224in}{1.455546in}}%
\pgfpathlineto{\pgfqpoint{4.428224in}{1.458495in}}%
\pgfpathlineto{\pgfqpoint{4.432765in}{1.458495in}}%
\pgfpathlineto{\pgfqpoint{4.432765in}{1.455546in}}%
\pgfpathmoveto{\pgfqpoint{4.428224in}{1.458495in}}%
\pgfpathlineto{\pgfqpoint{4.428224in}{1.458495in}}%
\pgfpathlineto{\pgfqpoint{4.428224in}{1.461444in}}%
\pgfpathlineto{\pgfqpoint{4.432765in}{1.461444in}}%
\pgfpathlineto{\pgfqpoint{4.432765in}{1.458495in}}%
\pgfpathmoveto{\pgfqpoint{4.437306in}{1.443749in}}%
\pgfpathlineto{\pgfqpoint{4.437306in}{1.443749in}}%
\pgfpathlineto{\pgfqpoint{4.437306in}{1.446699in}}%
\pgfpathlineto{\pgfqpoint{4.441846in}{1.446699in}}%
\pgfpathlineto{\pgfqpoint{4.441846in}{1.443749in}}%
\pgfpathmoveto{\pgfqpoint{4.437306in}{1.446699in}}%
\pgfpathlineto{\pgfqpoint{4.437306in}{1.446699in}}%
\pgfpathlineto{\pgfqpoint{4.437306in}{1.449648in}}%
\pgfpathlineto{\pgfqpoint{4.441846in}{1.449648in}}%
\pgfpathlineto{\pgfqpoint{4.441846in}{1.446699in}}%
\pgfpathmoveto{\pgfqpoint{4.419142in}{1.467343in}}%
\pgfpathlineto{\pgfqpoint{4.419142in}{1.467343in}}%
\pgfpathlineto{\pgfqpoint{4.419142in}{1.470292in}}%
\pgfpathlineto{\pgfqpoint{4.423683in}{1.470292in}}%
\pgfpathlineto{\pgfqpoint{4.423683in}{1.467343in}}%
\pgfpathmoveto{\pgfqpoint{4.419142in}{1.470292in}}%
\pgfpathlineto{\pgfqpoint{4.419142in}{1.470292in}}%
\pgfpathlineto{\pgfqpoint{4.419142in}{1.473241in}}%
\pgfpathlineto{\pgfqpoint{4.423683in}{1.473241in}}%
\pgfpathlineto{\pgfqpoint{4.423683in}{1.470292in}}%
\pgfpathmoveto{\pgfqpoint{4.396438in}{1.493885in}}%
\pgfpathlineto{\pgfqpoint{4.396438in}{1.493885in}}%
\pgfpathlineto{\pgfqpoint{4.396438in}{1.496835in}}%
\pgfpathlineto{\pgfqpoint{4.400978in}{1.496835in}}%
\pgfpathlineto{\pgfqpoint{4.400978in}{1.493885in}}%
\pgfpathmoveto{\pgfqpoint{4.391897in}{1.499784in}}%
\pgfpathlineto{\pgfqpoint{4.391897in}{1.499784in}}%
\pgfpathlineto{\pgfqpoint{4.391897in}{1.502733in}}%
\pgfpathlineto{\pgfqpoint{4.396438in}{1.502733in}}%
\pgfpathlineto{\pgfqpoint{4.396438in}{1.499784in}}%
\pgfpathmoveto{\pgfqpoint{4.396438in}{1.496835in}}%
\pgfpathlineto{\pgfqpoint{4.396438in}{1.496835in}}%
\pgfpathlineto{\pgfqpoint{4.396438in}{1.499784in}}%
\pgfpathlineto{\pgfqpoint{4.400978in}{1.499784in}}%
\pgfpathlineto{\pgfqpoint{4.400978in}{1.496835in}}%
\pgfpathmoveto{\pgfqpoint{4.396438in}{1.499784in}}%
\pgfpathlineto{\pgfqpoint{4.396438in}{1.499784in}}%
\pgfpathlineto{\pgfqpoint{4.396438in}{1.502733in}}%
\pgfpathlineto{\pgfqpoint{4.400978in}{1.502733in}}%
\pgfpathlineto{\pgfqpoint{4.400978in}{1.499784in}}%
\pgfpathmoveto{\pgfqpoint{4.387356in}{1.505682in}}%
\pgfpathlineto{\pgfqpoint{4.387356in}{1.505682in}}%
\pgfpathlineto{\pgfqpoint{4.387356in}{1.508631in}}%
\pgfpathlineto{\pgfqpoint{4.391897in}{1.508631in}}%
\pgfpathlineto{\pgfqpoint{4.391897in}{1.505682in}}%
\pgfpathmoveto{\pgfqpoint{4.382815in}{1.511580in}}%
\pgfpathlineto{\pgfqpoint{4.382815in}{1.511580in}}%
\pgfpathlineto{\pgfqpoint{4.382815in}{1.514530in}}%
\pgfpathlineto{\pgfqpoint{4.387356in}{1.514530in}}%
\pgfpathlineto{\pgfqpoint{4.387356in}{1.511580in}}%
\pgfpathmoveto{\pgfqpoint{4.387356in}{1.508631in}}%
\pgfpathlineto{\pgfqpoint{4.387356in}{1.508631in}}%
\pgfpathlineto{\pgfqpoint{4.387356in}{1.511580in}}%
\pgfpathlineto{\pgfqpoint{4.391897in}{1.511580in}}%
\pgfpathlineto{\pgfqpoint{4.391897in}{1.508631in}}%
\pgfpathmoveto{\pgfqpoint{4.387356in}{1.511580in}}%
\pgfpathlineto{\pgfqpoint{4.387356in}{1.511580in}}%
\pgfpathlineto{\pgfqpoint{4.387356in}{1.514530in}}%
\pgfpathlineto{\pgfqpoint{4.391897in}{1.514530in}}%
\pgfpathlineto{\pgfqpoint{4.391897in}{1.511580in}}%
\pgfpathmoveto{\pgfqpoint{4.391897in}{1.502733in}}%
\pgfpathlineto{\pgfqpoint{4.391897in}{1.502733in}}%
\pgfpathlineto{\pgfqpoint{4.391897in}{1.505682in}}%
\pgfpathlineto{\pgfqpoint{4.396438in}{1.505682in}}%
\pgfpathlineto{\pgfqpoint{4.396438in}{1.502733in}}%
\pgfpathmoveto{\pgfqpoint{4.391897in}{1.505682in}}%
\pgfpathlineto{\pgfqpoint{4.391897in}{1.505682in}}%
\pgfpathlineto{\pgfqpoint{4.391897in}{1.508631in}}%
\pgfpathlineto{\pgfqpoint{4.396438in}{1.508631in}}%
\pgfpathlineto{\pgfqpoint{4.396438in}{1.505682in}}%
\pgfpathmoveto{\pgfqpoint{4.400978in}{1.490936in}}%
\pgfpathlineto{\pgfqpoint{4.400978in}{1.490936in}}%
\pgfpathlineto{\pgfqpoint{4.400978in}{1.493885in}}%
\pgfpathlineto{\pgfqpoint{4.405519in}{1.493885in}}%
\pgfpathlineto{\pgfqpoint{4.405519in}{1.490936in}}%
\pgfpathmoveto{\pgfqpoint{4.400978in}{1.493885in}}%
\pgfpathlineto{\pgfqpoint{4.400978in}{1.493885in}}%
\pgfpathlineto{\pgfqpoint{4.400978in}{1.496835in}}%
\pgfpathlineto{\pgfqpoint{4.405519in}{1.496835in}}%
\pgfpathlineto{\pgfqpoint{4.405519in}{1.493885in}}%
\pgfpathmoveto{\pgfqpoint{4.382815in}{1.514530in}}%
\pgfpathlineto{\pgfqpoint{4.382815in}{1.514530in}}%
\pgfpathlineto{\pgfqpoint{4.382815in}{1.517479in}}%
\pgfpathlineto{\pgfqpoint{4.387356in}{1.517479in}}%
\pgfpathlineto{\pgfqpoint{4.387356in}{1.514530in}}%
\pgfpathmoveto{\pgfqpoint{4.382815in}{1.517479in}}%
\pgfpathlineto{\pgfqpoint{4.382815in}{1.517479in}}%
\pgfpathlineto{\pgfqpoint{4.382815in}{1.520428in}}%
\pgfpathlineto{\pgfqpoint{4.387356in}{1.520428in}}%
\pgfpathlineto{\pgfqpoint{4.387356in}{1.517479in}}%
\pgfpathmoveto{\pgfqpoint{4.668896in}{1.139980in}}%
\pgfpathlineto{\pgfqpoint{4.668896in}{1.139980in}}%
\pgfpathlineto{\pgfqpoint{4.668896in}{1.142930in}}%
\pgfpathlineto{\pgfqpoint{4.673437in}{1.142930in}}%
\pgfpathlineto{\pgfqpoint{4.673437in}{1.139980in}}%
\pgfpathmoveto{\pgfqpoint{4.664355in}{1.145879in}}%
\pgfpathlineto{\pgfqpoint{4.664355in}{1.145879in}}%
\pgfpathlineto{\pgfqpoint{4.664355in}{1.148828in}}%
\pgfpathlineto{\pgfqpoint{4.668896in}{1.148828in}}%
\pgfpathlineto{\pgfqpoint{4.668896in}{1.145879in}}%
\pgfpathmoveto{\pgfqpoint{4.668896in}{1.142930in}}%
\pgfpathlineto{\pgfqpoint{4.668896in}{1.142930in}}%
\pgfpathlineto{\pgfqpoint{4.668896in}{1.145879in}}%
\pgfpathlineto{\pgfqpoint{4.673437in}{1.145879in}}%
\pgfpathlineto{\pgfqpoint{4.673437in}{1.142930in}}%
\pgfpathmoveto{\pgfqpoint{4.668896in}{1.145879in}}%
\pgfpathlineto{\pgfqpoint{4.668896in}{1.145879in}}%
\pgfpathlineto{\pgfqpoint{4.668896in}{1.148828in}}%
\pgfpathlineto{\pgfqpoint{4.673437in}{1.148828in}}%
\pgfpathlineto{\pgfqpoint{4.673437in}{1.145879in}}%
\pgfpathmoveto{\pgfqpoint{4.659814in}{1.151777in}}%
\pgfpathlineto{\pgfqpoint{4.659814in}{1.151777in}}%
\pgfpathlineto{\pgfqpoint{4.659814in}{1.154726in}}%
\pgfpathlineto{\pgfqpoint{4.664355in}{1.154726in}}%
\pgfpathlineto{\pgfqpoint{4.664355in}{1.151777in}}%
\pgfpathmoveto{\pgfqpoint{4.655272in}{1.157675in}}%
\pgfpathlineto{\pgfqpoint{4.655272in}{1.157675in}}%
\pgfpathlineto{\pgfqpoint{4.655272in}{1.160625in}}%
\pgfpathlineto{\pgfqpoint{4.659814in}{1.160625in}}%
\pgfpathlineto{\pgfqpoint{4.659814in}{1.157675in}}%
\pgfpathmoveto{\pgfqpoint{4.659814in}{1.154726in}}%
\pgfpathlineto{\pgfqpoint{4.659814in}{1.154726in}}%
\pgfpathlineto{\pgfqpoint{4.659814in}{1.157675in}}%
\pgfpathlineto{\pgfqpoint{4.664355in}{1.157675in}}%
\pgfpathlineto{\pgfqpoint{4.664355in}{1.154726in}}%
\pgfpathmoveto{\pgfqpoint{4.659814in}{1.157675in}}%
\pgfpathlineto{\pgfqpoint{4.659814in}{1.157675in}}%
\pgfpathlineto{\pgfqpoint{4.659814in}{1.160625in}}%
\pgfpathlineto{\pgfqpoint{4.664355in}{1.160625in}}%
\pgfpathlineto{\pgfqpoint{4.664355in}{1.157675in}}%
\pgfpathmoveto{\pgfqpoint{4.664355in}{1.148828in}}%
\pgfpathlineto{\pgfqpoint{4.664355in}{1.148828in}}%
\pgfpathlineto{\pgfqpoint{4.664355in}{1.151777in}}%
\pgfpathlineto{\pgfqpoint{4.668896in}{1.151777in}}%
\pgfpathlineto{\pgfqpoint{4.668896in}{1.148828in}}%
\pgfpathmoveto{\pgfqpoint{4.664355in}{1.151777in}}%
\pgfpathlineto{\pgfqpoint{4.664355in}{1.151777in}}%
\pgfpathlineto{\pgfqpoint{4.664355in}{1.154726in}}%
\pgfpathlineto{\pgfqpoint{4.668896in}{1.154726in}}%
\pgfpathlineto{\pgfqpoint{4.668896in}{1.151777in}}%
\pgfpathmoveto{\pgfqpoint{4.596239in}{1.234357in}}%
\pgfpathlineto{\pgfqpoint{4.596239in}{1.234357in}}%
\pgfpathlineto{\pgfqpoint{4.596239in}{1.237307in}}%
\pgfpathlineto{\pgfqpoint{4.600780in}{1.237307in}}%
\pgfpathlineto{\pgfqpoint{4.600780in}{1.234357in}}%
\pgfpathmoveto{\pgfqpoint{4.591698in}{1.240256in}}%
\pgfpathlineto{\pgfqpoint{4.591698in}{1.240256in}}%
\pgfpathlineto{\pgfqpoint{4.591698in}{1.243205in}}%
\pgfpathlineto{\pgfqpoint{4.596239in}{1.243205in}}%
\pgfpathlineto{\pgfqpoint{4.596239in}{1.240256in}}%
\pgfpathmoveto{\pgfqpoint{4.596239in}{1.237307in}}%
\pgfpathlineto{\pgfqpoint{4.596239in}{1.237307in}}%
\pgfpathlineto{\pgfqpoint{4.596239in}{1.240256in}}%
\pgfpathlineto{\pgfqpoint{4.600780in}{1.240256in}}%
\pgfpathlineto{\pgfqpoint{4.600780in}{1.237307in}}%
\pgfpathmoveto{\pgfqpoint{4.596239in}{1.240256in}}%
\pgfpathlineto{\pgfqpoint{4.596239in}{1.240256in}}%
\pgfpathlineto{\pgfqpoint{4.596239in}{1.243205in}}%
\pgfpathlineto{\pgfqpoint{4.600780in}{1.243205in}}%
\pgfpathlineto{\pgfqpoint{4.600780in}{1.240256in}}%
\pgfpathmoveto{\pgfqpoint{4.587157in}{1.246155in}}%
\pgfpathlineto{\pgfqpoint{4.587157in}{1.246155in}}%
\pgfpathlineto{\pgfqpoint{4.587157in}{1.249104in}}%
\pgfpathlineto{\pgfqpoint{4.591698in}{1.249104in}}%
\pgfpathlineto{\pgfqpoint{4.591698in}{1.246155in}}%
\pgfpathmoveto{\pgfqpoint{4.582616in}{1.252053in}}%
\pgfpathlineto{\pgfqpoint{4.582616in}{1.252053in}}%
\pgfpathlineto{\pgfqpoint{4.582616in}{1.255003in}}%
\pgfpathlineto{\pgfqpoint{4.587157in}{1.255003in}}%
\pgfpathlineto{\pgfqpoint{4.587157in}{1.252053in}}%
\pgfpathmoveto{\pgfqpoint{4.587157in}{1.249104in}}%
\pgfpathlineto{\pgfqpoint{4.587157in}{1.249104in}}%
\pgfpathlineto{\pgfqpoint{4.587157in}{1.252053in}}%
\pgfpathlineto{\pgfqpoint{4.591698in}{1.252053in}}%
\pgfpathlineto{\pgfqpoint{4.591698in}{1.249104in}}%
\pgfpathmoveto{\pgfqpoint{4.587157in}{1.252053in}}%
\pgfpathlineto{\pgfqpoint{4.587157in}{1.252053in}}%
\pgfpathlineto{\pgfqpoint{4.587157in}{1.255003in}}%
\pgfpathlineto{\pgfqpoint{4.591698in}{1.255003in}}%
\pgfpathlineto{\pgfqpoint{4.591698in}{1.252053in}}%
\pgfpathmoveto{\pgfqpoint{4.591698in}{1.243205in}}%
\pgfpathlineto{\pgfqpoint{4.591698in}{1.243205in}}%
\pgfpathlineto{\pgfqpoint{4.591698in}{1.246155in}}%
\pgfpathlineto{\pgfqpoint{4.596239in}{1.246155in}}%
\pgfpathlineto{\pgfqpoint{4.596239in}{1.243205in}}%
\pgfpathmoveto{\pgfqpoint{4.591698in}{1.246155in}}%
\pgfpathlineto{\pgfqpoint{4.591698in}{1.246155in}}%
\pgfpathlineto{\pgfqpoint{4.591698in}{1.249104in}}%
\pgfpathlineto{\pgfqpoint{4.596239in}{1.249104in}}%
\pgfpathlineto{\pgfqpoint{4.596239in}{1.246155in}}%
\pgfpathmoveto{\pgfqpoint{4.632567in}{1.187168in}}%
\pgfpathlineto{\pgfqpoint{4.632567in}{1.187168in}}%
\pgfpathlineto{\pgfqpoint{4.632567in}{1.190118in}}%
\pgfpathlineto{\pgfqpoint{4.637108in}{1.190118in}}%
\pgfpathlineto{\pgfqpoint{4.637108in}{1.187168in}}%
\pgfpathmoveto{\pgfqpoint{4.628026in}{1.193067in}}%
\pgfpathlineto{\pgfqpoint{4.628026in}{1.193067in}}%
\pgfpathlineto{\pgfqpoint{4.628026in}{1.196016in}}%
\pgfpathlineto{\pgfqpoint{4.632567in}{1.196016in}}%
\pgfpathlineto{\pgfqpoint{4.632567in}{1.193067in}}%
\pgfpathmoveto{\pgfqpoint{4.632567in}{1.190118in}}%
\pgfpathlineto{\pgfqpoint{4.632567in}{1.190118in}}%
\pgfpathlineto{\pgfqpoint{4.632567in}{1.193067in}}%
\pgfpathlineto{\pgfqpoint{4.637108in}{1.193067in}}%
\pgfpathlineto{\pgfqpoint{4.637108in}{1.190118in}}%
\pgfpathmoveto{\pgfqpoint{4.632567in}{1.193067in}}%
\pgfpathlineto{\pgfqpoint{4.632567in}{1.193067in}}%
\pgfpathlineto{\pgfqpoint{4.632567in}{1.196016in}}%
\pgfpathlineto{\pgfqpoint{4.637108in}{1.196016in}}%
\pgfpathlineto{\pgfqpoint{4.637108in}{1.193067in}}%
\pgfpathmoveto{\pgfqpoint{4.623485in}{1.198966in}}%
\pgfpathlineto{\pgfqpoint{4.623485in}{1.198966in}}%
\pgfpathlineto{\pgfqpoint{4.623485in}{1.201915in}}%
\pgfpathlineto{\pgfqpoint{4.628026in}{1.201915in}}%
\pgfpathlineto{\pgfqpoint{4.628026in}{1.198966in}}%
\pgfpathmoveto{\pgfqpoint{4.618944in}{1.204864in}}%
\pgfpathlineto{\pgfqpoint{4.618944in}{1.204864in}}%
\pgfpathlineto{\pgfqpoint{4.618944in}{1.207814in}}%
\pgfpathlineto{\pgfqpoint{4.623485in}{1.207814in}}%
\pgfpathlineto{\pgfqpoint{4.623485in}{1.204864in}}%
\pgfpathmoveto{\pgfqpoint{4.623485in}{1.201915in}}%
\pgfpathlineto{\pgfqpoint{4.623485in}{1.201915in}}%
\pgfpathlineto{\pgfqpoint{4.623485in}{1.204864in}}%
\pgfpathlineto{\pgfqpoint{4.628026in}{1.204864in}}%
\pgfpathlineto{\pgfqpoint{4.628026in}{1.201915in}}%
\pgfpathmoveto{\pgfqpoint{4.623485in}{1.204864in}}%
\pgfpathlineto{\pgfqpoint{4.623485in}{1.204864in}}%
\pgfpathlineto{\pgfqpoint{4.623485in}{1.207814in}}%
\pgfpathlineto{\pgfqpoint{4.628026in}{1.207814in}}%
\pgfpathlineto{\pgfqpoint{4.628026in}{1.204864in}}%
\pgfpathmoveto{\pgfqpoint{4.628026in}{1.196016in}}%
\pgfpathlineto{\pgfqpoint{4.628026in}{1.196016in}}%
\pgfpathlineto{\pgfqpoint{4.628026in}{1.198966in}}%
\pgfpathlineto{\pgfqpoint{4.632567in}{1.198966in}}%
\pgfpathlineto{\pgfqpoint{4.632567in}{1.196016in}}%
\pgfpathmoveto{\pgfqpoint{4.628026in}{1.198966in}}%
\pgfpathlineto{\pgfqpoint{4.628026in}{1.198966in}}%
\pgfpathlineto{\pgfqpoint{4.628026in}{1.201915in}}%
\pgfpathlineto{\pgfqpoint{4.632567in}{1.201915in}}%
\pgfpathlineto{\pgfqpoint{4.632567in}{1.198966in}}%
\pgfpathmoveto{\pgfqpoint{4.650731in}{1.163574in}}%
\pgfpathlineto{\pgfqpoint{4.650731in}{1.163574in}}%
\pgfpathlineto{\pgfqpoint{4.650731in}{1.166523in}}%
\pgfpathlineto{\pgfqpoint{4.655272in}{1.166523in}}%
\pgfpathlineto{\pgfqpoint{4.655272in}{1.163574in}}%
\pgfpathmoveto{\pgfqpoint{4.646190in}{1.169472in}}%
\pgfpathlineto{\pgfqpoint{4.646190in}{1.169472in}}%
\pgfpathlineto{\pgfqpoint{4.646190in}{1.172422in}}%
\pgfpathlineto{\pgfqpoint{4.650731in}{1.172422in}}%
\pgfpathlineto{\pgfqpoint{4.650731in}{1.169472in}}%
\pgfpathmoveto{\pgfqpoint{4.650731in}{1.166523in}}%
\pgfpathlineto{\pgfqpoint{4.650731in}{1.166523in}}%
\pgfpathlineto{\pgfqpoint{4.650731in}{1.169472in}}%
\pgfpathlineto{\pgfqpoint{4.655272in}{1.169472in}}%
\pgfpathlineto{\pgfqpoint{4.655272in}{1.166523in}}%
\pgfpathmoveto{\pgfqpoint{4.650731in}{1.169472in}}%
\pgfpathlineto{\pgfqpoint{4.650731in}{1.169472in}}%
\pgfpathlineto{\pgfqpoint{4.650731in}{1.172422in}}%
\pgfpathlineto{\pgfqpoint{4.655272in}{1.172422in}}%
\pgfpathlineto{\pgfqpoint{4.655272in}{1.169472in}}%
\pgfpathmoveto{\pgfqpoint{4.641649in}{1.175371in}}%
\pgfpathlineto{\pgfqpoint{4.641649in}{1.175371in}}%
\pgfpathlineto{\pgfqpoint{4.641649in}{1.178320in}}%
\pgfpathlineto{\pgfqpoint{4.646190in}{1.178320in}}%
\pgfpathlineto{\pgfqpoint{4.646190in}{1.175371in}}%
\pgfpathmoveto{\pgfqpoint{4.637108in}{1.181270in}}%
\pgfpathlineto{\pgfqpoint{4.637108in}{1.181270in}}%
\pgfpathlineto{\pgfqpoint{4.637108in}{1.184219in}}%
\pgfpathlineto{\pgfqpoint{4.641649in}{1.184219in}}%
\pgfpathlineto{\pgfqpoint{4.641649in}{1.181270in}}%
\pgfpathmoveto{\pgfqpoint{4.641649in}{1.178320in}}%
\pgfpathlineto{\pgfqpoint{4.641649in}{1.178320in}}%
\pgfpathlineto{\pgfqpoint{4.641649in}{1.181270in}}%
\pgfpathlineto{\pgfqpoint{4.646190in}{1.181270in}}%
\pgfpathlineto{\pgfqpoint{4.646190in}{1.178320in}}%
\pgfpathmoveto{\pgfqpoint{4.641649in}{1.181270in}}%
\pgfpathlineto{\pgfqpoint{4.641649in}{1.181270in}}%
\pgfpathlineto{\pgfqpoint{4.641649in}{1.184219in}}%
\pgfpathlineto{\pgfqpoint{4.646190in}{1.184219in}}%
\pgfpathlineto{\pgfqpoint{4.646190in}{1.181270in}}%
\pgfpathmoveto{\pgfqpoint{4.646190in}{1.172422in}}%
\pgfpathlineto{\pgfqpoint{4.646190in}{1.172422in}}%
\pgfpathlineto{\pgfqpoint{4.646190in}{1.175371in}}%
\pgfpathlineto{\pgfqpoint{4.650731in}{1.175371in}}%
\pgfpathlineto{\pgfqpoint{4.650731in}{1.172422in}}%
\pgfpathmoveto{\pgfqpoint{4.646190in}{1.175371in}}%
\pgfpathlineto{\pgfqpoint{4.646190in}{1.175371in}}%
\pgfpathlineto{\pgfqpoint{4.646190in}{1.178320in}}%
\pgfpathlineto{\pgfqpoint{4.650731in}{1.178320in}}%
\pgfpathlineto{\pgfqpoint{4.650731in}{1.175371in}}%
\pgfpathmoveto{\pgfqpoint{4.655272in}{1.160625in}}%
\pgfpathlineto{\pgfqpoint{4.655272in}{1.160625in}}%
\pgfpathlineto{\pgfqpoint{4.655272in}{1.163574in}}%
\pgfpathlineto{\pgfqpoint{4.659814in}{1.163574in}}%
\pgfpathlineto{\pgfqpoint{4.659814in}{1.160625in}}%
\pgfpathmoveto{\pgfqpoint{4.655272in}{1.163574in}}%
\pgfpathlineto{\pgfqpoint{4.655272in}{1.163574in}}%
\pgfpathlineto{\pgfqpoint{4.655272in}{1.166523in}}%
\pgfpathlineto{\pgfqpoint{4.659814in}{1.166523in}}%
\pgfpathlineto{\pgfqpoint{4.659814in}{1.163574in}}%
\pgfpathmoveto{\pgfqpoint{4.637108in}{1.184219in}}%
\pgfpathlineto{\pgfqpoint{4.637108in}{1.184219in}}%
\pgfpathlineto{\pgfqpoint{4.637108in}{1.187168in}}%
\pgfpathlineto{\pgfqpoint{4.641649in}{1.187168in}}%
\pgfpathlineto{\pgfqpoint{4.641649in}{1.184219in}}%
\pgfpathmoveto{\pgfqpoint{4.637108in}{1.187168in}}%
\pgfpathlineto{\pgfqpoint{4.637108in}{1.187168in}}%
\pgfpathlineto{\pgfqpoint{4.637108in}{1.190118in}}%
\pgfpathlineto{\pgfqpoint{4.641649in}{1.190118in}}%
\pgfpathlineto{\pgfqpoint{4.641649in}{1.187168in}}%
\pgfpathmoveto{\pgfqpoint{4.614403in}{1.210763in}}%
\pgfpathlineto{\pgfqpoint{4.614403in}{1.210763in}}%
\pgfpathlineto{\pgfqpoint{4.614403in}{1.213712in}}%
\pgfpathlineto{\pgfqpoint{4.618944in}{1.213712in}}%
\pgfpathlineto{\pgfqpoint{4.618944in}{1.210763in}}%
\pgfpathmoveto{\pgfqpoint{4.609862in}{1.216662in}}%
\pgfpathlineto{\pgfqpoint{4.609862in}{1.216662in}}%
\pgfpathlineto{\pgfqpoint{4.609862in}{1.219611in}}%
\pgfpathlineto{\pgfqpoint{4.614403in}{1.219611in}}%
\pgfpathlineto{\pgfqpoint{4.614403in}{1.216662in}}%
\pgfpathmoveto{\pgfqpoint{4.614403in}{1.213712in}}%
\pgfpathlineto{\pgfqpoint{4.614403in}{1.213712in}}%
\pgfpathlineto{\pgfqpoint{4.614403in}{1.216662in}}%
\pgfpathlineto{\pgfqpoint{4.618944in}{1.216662in}}%
\pgfpathlineto{\pgfqpoint{4.618944in}{1.213712in}}%
\pgfpathmoveto{\pgfqpoint{4.614403in}{1.216662in}}%
\pgfpathlineto{\pgfqpoint{4.614403in}{1.216662in}}%
\pgfpathlineto{\pgfqpoint{4.614403in}{1.219611in}}%
\pgfpathlineto{\pgfqpoint{4.618944in}{1.219611in}}%
\pgfpathlineto{\pgfqpoint{4.618944in}{1.216662in}}%
\pgfpathmoveto{\pgfqpoint{4.605321in}{1.222560in}}%
\pgfpathlineto{\pgfqpoint{4.605321in}{1.222560in}}%
\pgfpathlineto{\pgfqpoint{4.605321in}{1.225509in}}%
\pgfpathlineto{\pgfqpoint{4.609862in}{1.225509in}}%
\pgfpathlineto{\pgfqpoint{4.609862in}{1.222560in}}%
\pgfpathmoveto{\pgfqpoint{4.600780in}{1.228459in}}%
\pgfpathlineto{\pgfqpoint{4.600780in}{1.228459in}}%
\pgfpathlineto{\pgfqpoint{4.600780in}{1.231408in}}%
\pgfpathlineto{\pgfqpoint{4.605321in}{1.231408in}}%
\pgfpathlineto{\pgfqpoint{4.605321in}{1.228459in}}%
\pgfpathmoveto{\pgfqpoint{4.605321in}{1.225509in}}%
\pgfpathlineto{\pgfqpoint{4.605321in}{1.225509in}}%
\pgfpathlineto{\pgfqpoint{4.605321in}{1.228459in}}%
\pgfpathlineto{\pgfqpoint{4.609862in}{1.228459in}}%
\pgfpathlineto{\pgfqpoint{4.609862in}{1.225509in}}%
\pgfpathmoveto{\pgfqpoint{4.605321in}{1.228459in}}%
\pgfpathlineto{\pgfqpoint{4.605321in}{1.228459in}}%
\pgfpathlineto{\pgfqpoint{4.605321in}{1.231408in}}%
\pgfpathlineto{\pgfqpoint{4.609862in}{1.231408in}}%
\pgfpathlineto{\pgfqpoint{4.609862in}{1.228459in}}%
\pgfpathmoveto{\pgfqpoint{4.609862in}{1.219611in}}%
\pgfpathlineto{\pgfqpoint{4.609862in}{1.219611in}}%
\pgfpathlineto{\pgfqpoint{4.609862in}{1.222560in}}%
\pgfpathlineto{\pgfqpoint{4.614403in}{1.222560in}}%
\pgfpathlineto{\pgfqpoint{4.614403in}{1.219611in}}%
\pgfpathmoveto{\pgfqpoint{4.609862in}{1.222560in}}%
\pgfpathlineto{\pgfqpoint{4.609862in}{1.222560in}}%
\pgfpathlineto{\pgfqpoint{4.609862in}{1.225509in}}%
\pgfpathlineto{\pgfqpoint{4.614403in}{1.225509in}}%
\pgfpathlineto{\pgfqpoint{4.614403in}{1.222560in}}%
\pgfpathmoveto{\pgfqpoint{4.618944in}{1.207814in}}%
\pgfpathlineto{\pgfqpoint{4.618944in}{1.207814in}}%
\pgfpathlineto{\pgfqpoint{4.618944in}{1.210763in}}%
\pgfpathlineto{\pgfqpoint{4.623485in}{1.210763in}}%
\pgfpathlineto{\pgfqpoint{4.623485in}{1.207814in}}%
\pgfpathmoveto{\pgfqpoint{4.618944in}{1.210763in}}%
\pgfpathlineto{\pgfqpoint{4.618944in}{1.210763in}}%
\pgfpathlineto{\pgfqpoint{4.618944in}{1.213712in}}%
\pgfpathlineto{\pgfqpoint{4.623485in}{1.213712in}}%
\pgfpathlineto{\pgfqpoint{4.623485in}{1.210763in}}%
\pgfpathmoveto{\pgfqpoint{4.600780in}{1.231408in}}%
\pgfpathlineto{\pgfqpoint{4.600780in}{1.231408in}}%
\pgfpathlineto{\pgfqpoint{4.600780in}{1.234357in}}%
\pgfpathlineto{\pgfqpoint{4.605321in}{1.234357in}}%
\pgfpathlineto{\pgfqpoint{4.605321in}{1.231408in}}%
\pgfpathmoveto{\pgfqpoint{4.600780in}{1.234357in}}%
\pgfpathlineto{\pgfqpoint{4.600780in}{1.234357in}}%
\pgfpathlineto{\pgfqpoint{4.600780in}{1.237307in}}%
\pgfpathlineto{\pgfqpoint{4.605321in}{1.237307in}}%
\pgfpathlineto{\pgfqpoint{4.605321in}{1.234357in}}%
\pgfpathmoveto{\pgfqpoint{4.559911in}{1.281544in}}%
\pgfpathlineto{\pgfqpoint{4.559911in}{1.281544in}}%
\pgfpathlineto{\pgfqpoint{4.559911in}{1.284493in}}%
\pgfpathlineto{\pgfqpoint{4.564452in}{1.284493in}}%
\pgfpathlineto{\pgfqpoint{4.564452in}{1.281544in}}%
\pgfpathmoveto{\pgfqpoint{4.555370in}{1.287442in}}%
\pgfpathlineto{\pgfqpoint{4.555370in}{1.287442in}}%
\pgfpathlineto{\pgfqpoint{4.555370in}{1.290391in}}%
\pgfpathlineto{\pgfqpoint{4.559911in}{1.290391in}}%
\pgfpathlineto{\pgfqpoint{4.559911in}{1.287442in}}%
\pgfpathmoveto{\pgfqpoint{4.559911in}{1.284493in}}%
\pgfpathlineto{\pgfqpoint{4.559911in}{1.284493in}}%
\pgfpathlineto{\pgfqpoint{4.559911in}{1.287442in}}%
\pgfpathlineto{\pgfqpoint{4.564452in}{1.287442in}}%
\pgfpathlineto{\pgfqpoint{4.564452in}{1.284493in}}%
\pgfpathmoveto{\pgfqpoint{4.559911in}{1.287442in}}%
\pgfpathlineto{\pgfqpoint{4.559911in}{1.287442in}}%
\pgfpathlineto{\pgfqpoint{4.559911in}{1.290391in}}%
\pgfpathlineto{\pgfqpoint{4.564452in}{1.290391in}}%
\pgfpathlineto{\pgfqpoint{4.564452in}{1.287442in}}%
\pgfpathmoveto{\pgfqpoint{4.550829in}{1.293340in}}%
\pgfpathlineto{\pgfqpoint{4.550829in}{1.293340in}}%
\pgfpathlineto{\pgfqpoint{4.550829in}{1.296289in}}%
\pgfpathlineto{\pgfqpoint{4.555370in}{1.296289in}}%
\pgfpathlineto{\pgfqpoint{4.555370in}{1.293340in}}%
\pgfpathmoveto{\pgfqpoint{4.546288in}{1.299238in}}%
\pgfpathlineto{\pgfqpoint{4.546288in}{1.299238in}}%
\pgfpathlineto{\pgfqpoint{4.546288in}{1.302188in}}%
\pgfpathlineto{\pgfqpoint{4.550829in}{1.302188in}}%
\pgfpathlineto{\pgfqpoint{4.550829in}{1.299238in}}%
\pgfpathmoveto{\pgfqpoint{4.550829in}{1.296289in}}%
\pgfpathlineto{\pgfqpoint{4.550829in}{1.296289in}}%
\pgfpathlineto{\pgfqpoint{4.550829in}{1.299238in}}%
\pgfpathlineto{\pgfqpoint{4.555370in}{1.299238in}}%
\pgfpathlineto{\pgfqpoint{4.555370in}{1.296289in}}%
\pgfpathmoveto{\pgfqpoint{4.550829in}{1.299238in}}%
\pgfpathlineto{\pgfqpoint{4.550829in}{1.299238in}}%
\pgfpathlineto{\pgfqpoint{4.550829in}{1.302188in}}%
\pgfpathlineto{\pgfqpoint{4.555370in}{1.302188in}}%
\pgfpathlineto{\pgfqpoint{4.555370in}{1.299238in}}%
\pgfpathmoveto{\pgfqpoint{4.555370in}{1.290391in}}%
\pgfpathlineto{\pgfqpoint{4.555370in}{1.290391in}}%
\pgfpathlineto{\pgfqpoint{4.555370in}{1.293340in}}%
\pgfpathlineto{\pgfqpoint{4.559911in}{1.293340in}}%
\pgfpathlineto{\pgfqpoint{4.559911in}{1.290391in}}%
\pgfpathmoveto{\pgfqpoint{4.555370in}{1.293340in}}%
\pgfpathlineto{\pgfqpoint{4.555370in}{1.293340in}}%
\pgfpathlineto{\pgfqpoint{4.555370in}{1.296289in}}%
\pgfpathlineto{\pgfqpoint{4.559911in}{1.296289in}}%
\pgfpathlineto{\pgfqpoint{4.559911in}{1.293340in}}%
\pgfpathmoveto{\pgfqpoint{4.578075in}{1.257952in}}%
\pgfpathlineto{\pgfqpoint{4.578075in}{1.257952in}}%
\pgfpathlineto{\pgfqpoint{4.578075in}{1.260901in}}%
\pgfpathlineto{\pgfqpoint{4.582616in}{1.260901in}}%
\pgfpathlineto{\pgfqpoint{4.582616in}{1.257952in}}%
\pgfpathmoveto{\pgfqpoint{4.573534in}{1.263850in}}%
\pgfpathlineto{\pgfqpoint{4.573534in}{1.263850in}}%
\pgfpathlineto{\pgfqpoint{4.573534in}{1.266799in}}%
\pgfpathlineto{\pgfqpoint{4.578075in}{1.266799in}}%
\pgfpathlineto{\pgfqpoint{4.578075in}{1.263850in}}%
\pgfpathmoveto{\pgfqpoint{4.578075in}{1.260901in}}%
\pgfpathlineto{\pgfqpoint{4.578075in}{1.260901in}}%
\pgfpathlineto{\pgfqpoint{4.578075in}{1.263850in}}%
\pgfpathlineto{\pgfqpoint{4.582616in}{1.263850in}}%
\pgfpathlineto{\pgfqpoint{4.582616in}{1.260901in}}%
\pgfpathmoveto{\pgfqpoint{4.578075in}{1.263850in}}%
\pgfpathlineto{\pgfqpoint{4.578075in}{1.263850in}}%
\pgfpathlineto{\pgfqpoint{4.578075in}{1.266799in}}%
\pgfpathlineto{\pgfqpoint{4.582616in}{1.266799in}}%
\pgfpathlineto{\pgfqpoint{4.582616in}{1.263850in}}%
\pgfpathmoveto{\pgfqpoint{4.568993in}{1.269748in}}%
\pgfpathlineto{\pgfqpoint{4.568993in}{1.269748in}}%
\pgfpathlineto{\pgfqpoint{4.568993in}{1.272697in}}%
\pgfpathlineto{\pgfqpoint{4.573534in}{1.272697in}}%
\pgfpathlineto{\pgfqpoint{4.573534in}{1.269748in}}%
\pgfpathmoveto{\pgfqpoint{4.564452in}{1.275646in}}%
\pgfpathlineto{\pgfqpoint{4.564452in}{1.275646in}}%
\pgfpathlineto{\pgfqpoint{4.564452in}{1.278595in}}%
\pgfpathlineto{\pgfqpoint{4.568993in}{1.278595in}}%
\pgfpathlineto{\pgfqpoint{4.568993in}{1.275646in}}%
\pgfpathmoveto{\pgfqpoint{4.568993in}{1.272697in}}%
\pgfpathlineto{\pgfqpoint{4.568993in}{1.272697in}}%
\pgfpathlineto{\pgfqpoint{4.568993in}{1.275646in}}%
\pgfpathlineto{\pgfqpoint{4.573534in}{1.275646in}}%
\pgfpathlineto{\pgfqpoint{4.573534in}{1.272697in}}%
\pgfpathmoveto{\pgfqpoint{4.568993in}{1.275646in}}%
\pgfpathlineto{\pgfqpoint{4.568993in}{1.275646in}}%
\pgfpathlineto{\pgfqpoint{4.568993in}{1.278595in}}%
\pgfpathlineto{\pgfqpoint{4.573534in}{1.278595in}}%
\pgfpathlineto{\pgfqpoint{4.573534in}{1.275646in}}%
\pgfpathmoveto{\pgfqpoint{4.573534in}{1.266799in}}%
\pgfpathlineto{\pgfqpoint{4.573534in}{1.266799in}}%
\pgfpathlineto{\pgfqpoint{4.573534in}{1.269748in}}%
\pgfpathlineto{\pgfqpoint{4.578075in}{1.269748in}}%
\pgfpathlineto{\pgfqpoint{4.578075in}{1.266799in}}%
\pgfpathmoveto{\pgfqpoint{4.573534in}{1.269748in}}%
\pgfpathlineto{\pgfqpoint{4.573534in}{1.269748in}}%
\pgfpathlineto{\pgfqpoint{4.573534in}{1.272697in}}%
\pgfpathlineto{\pgfqpoint{4.578075in}{1.272697in}}%
\pgfpathlineto{\pgfqpoint{4.578075in}{1.269748in}}%
\pgfpathmoveto{\pgfqpoint{4.582616in}{1.255003in}}%
\pgfpathlineto{\pgfqpoint{4.582616in}{1.255003in}}%
\pgfpathlineto{\pgfqpoint{4.582616in}{1.257952in}}%
\pgfpathlineto{\pgfqpoint{4.587157in}{1.257952in}}%
\pgfpathlineto{\pgfqpoint{4.587157in}{1.255003in}}%
\pgfpathmoveto{\pgfqpoint{4.582616in}{1.257952in}}%
\pgfpathlineto{\pgfqpoint{4.582616in}{1.257952in}}%
\pgfpathlineto{\pgfqpoint{4.582616in}{1.260901in}}%
\pgfpathlineto{\pgfqpoint{4.587157in}{1.260901in}}%
\pgfpathlineto{\pgfqpoint{4.587157in}{1.257952in}}%
\pgfpathmoveto{\pgfqpoint{4.564452in}{1.278595in}}%
\pgfpathlineto{\pgfqpoint{4.564452in}{1.278595in}}%
\pgfpathlineto{\pgfqpoint{4.564452in}{1.281544in}}%
\pgfpathlineto{\pgfqpoint{4.568993in}{1.281544in}}%
\pgfpathlineto{\pgfqpoint{4.568993in}{1.278595in}}%
\pgfpathmoveto{\pgfqpoint{4.564452in}{1.281544in}}%
\pgfpathlineto{\pgfqpoint{4.564452in}{1.281544in}}%
\pgfpathlineto{\pgfqpoint{4.564452in}{1.284493in}}%
\pgfpathlineto{\pgfqpoint{4.568993in}{1.284493in}}%
\pgfpathlineto{\pgfqpoint{4.568993in}{1.281544in}}%
\pgfpathmoveto{\pgfqpoint{4.541746in}{1.305137in}}%
\pgfpathlineto{\pgfqpoint{4.541746in}{1.305137in}}%
\pgfpathlineto{\pgfqpoint{4.541746in}{1.308086in}}%
\pgfpathlineto{\pgfqpoint{4.546288in}{1.308086in}}%
\pgfpathlineto{\pgfqpoint{4.546288in}{1.305137in}}%
\pgfpathmoveto{\pgfqpoint{4.537205in}{1.311035in}}%
\pgfpathlineto{\pgfqpoint{4.537205in}{1.311035in}}%
\pgfpathlineto{\pgfqpoint{4.537205in}{1.313984in}}%
\pgfpathlineto{\pgfqpoint{4.541746in}{1.313984in}}%
\pgfpathlineto{\pgfqpoint{4.541746in}{1.311035in}}%
\pgfpathmoveto{\pgfqpoint{4.541746in}{1.308086in}}%
\pgfpathlineto{\pgfqpoint{4.541746in}{1.308086in}}%
\pgfpathlineto{\pgfqpoint{4.541746in}{1.311035in}}%
\pgfpathlineto{\pgfqpoint{4.546288in}{1.311035in}}%
\pgfpathlineto{\pgfqpoint{4.546288in}{1.308086in}}%
\pgfpathmoveto{\pgfqpoint{4.541746in}{1.311035in}}%
\pgfpathlineto{\pgfqpoint{4.541746in}{1.311035in}}%
\pgfpathlineto{\pgfqpoint{4.541746in}{1.313984in}}%
\pgfpathlineto{\pgfqpoint{4.546288in}{1.313984in}}%
\pgfpathlineto{\pgfqpoint{4.546288in}{1.311035in}}%
\pgfpathmoveto{\pgfqpoint{4.532664in}{1.316933in}}%
\pgfpathlineto{\pgfqpoint{4.532664in}{1.316933in}}%
\pgfpathlineto{\pgfqpoint{4.532664in}{1.319882in}}%
\pgfpathlineto{\pgfqpoint{4.537205in}{1.319882in}}%
\pgfpathlineto{\pgfqpoint{4.537205in}{1.316933in}}%
\pgfpathmoveto{\pgfqpoint{4.528123in}{1.322831in}}%
\pgfpathlineto{\pgfqpoint{4.528123in}{1.322831in}}%
\pgfpathlineto{\pgfqpoint{4.528123in}{1.325780in}}%
\pgfpathlineto{\pgfqpoint{4.532664in}{1.325780in}}%
\pgfpathlineto{\pgfqpoint{4.532664in}{1.322831in}}%
\pgfpathmoveto{\pgfqpoint{4.532664in}{1.319882in}}%
\pgfpathlineto{\pgfqpoint{4.532664in}{1.319882in}}%
\pgfpathlineto{\pgfqpoint{4.532664in}{1.322831in}}%
\pgfpathlineto{\pgfqpoint{4.537205in}{1.322831in}}%
\pgfpathlineto{\pgfqpoint{4.537205in}{1.319882in}}%
\pgfpathmoveto{\pgfqpoint{4.532664in}{1.322831in}}%
\pgfpathlineto{\pgfqpoint{4.532664in}{1.322831in}}%
\pgfpathlineto{\pgfqpoint{4.532664in}{1.325780in}}%
\pgfpathlineto{\pgfqpoint{4.537205in}{1.325780in}}%
\pgfpathlineto{\pgfqpoint{4.537205in}{1.322831in}}%
\pgfpathmoveto{\pgfqpoint{4.537205in}{1.313984in}}%
\pgfpathlineto{\pgfqpoint{4.537205in}{1.313984in}}%
\pgfpathlineto{\pgfqpoint{4.537205in}{1.316933in}}%
\pgfpathlineto{\pgfqpoint{4.541746in}{1.316933in}}%
\pgfpathlineto{\pgfqpoint{4.541746in}{1.313984in}}%
\pgfpathmoveto{\pgfqpoint{4.537205in}{1.316933in}}%
\pgfpathlineto{\pgfqpoint{4.537205in}{1.316933in}}%
\pgfpathlineto{\pgfqpoint{4.537205in}{1.319882in}}%
\pgfpathlineto{\pgfqpoint{4.541746in}{1.319882in}}%
\pgfpathlineto{\pgfqpoint{4.541746in}{1.316933in}}%
\pgfpathmoveto{\pgfqpoint{4.546288in}{1.302188in}}%
\pgfpathlineto{\pgfqpoint{4.546288in}{1.302188in}}%
\pgfpathlineto{\pgfqpoint{4.546288in}{1.305137in}}%
\pgfpathlineto{\pgfqpoint{4.550829in}{1.305137in}}%
\pgfpathlineto{\pgfqpoint{4.550829in}{1.302188in}}%
\pgfpathmoveto{\pgfqpoint{4.546288in}{1.305137in}}%
\pgfpathlineto{\pgfqpoint{4.546288in}{1.305137in}}%
\pgfpathlineto{\pgfqpoint{4.546288in}{1.308086in}}%
\pgfpathlineto{\pgfqpoint{4.550829in}{1.308086in}}%
\pgfpathlineto{\pgfqpoint{4.550829in}{1.305137in}}%
\pgfpathmoveto{\pgfqpoint{4.528123in}{1.325780in}}%
\pgfpathlineto{\pgfqpoint{4.528123in}{1.325780in}}%
\pgfpathlineto{\pgfqpoint{4.528123in}{1.328729in}}%
\pgfpathlineto{\pgfqpoint{4.532664in}{1.328729in}}%
\pgfpathlineto{\pgfqpoint{4.532664in}{1.325780in}}%
\pgfpathmoveto{\pgfqpoint{4.528123in}{1.328729in}}%
\pgfpathlineto{\pgfqpoint{4.528123in}{1.328729in}}%
\pgfpathlineto{\pgfqpoint{4.528123in}{1.331678in}}%
\pgfpathlineto{\pgfqpoint{4.532664in}{1.331678in}}%
\pgfpathlineto{\pgfqpoint{4.532664in}{1.328729in}}%
\pgfpathmoveto{\pgfqpoint{4.814209in}{0.951232in}}%
\pgfpathlineto{\pgfqpoint{4.814209in}{0.951232in}}%
\pgfpathlineto{\pgfqpoint{4.814209in}{0.954181in}}%
\pgfpathlineto{\pgfqpoint{4.818750in}{0.954181in}}%
\pgfpathlineto{\pgfqpoint{4.818750in}{0.951232in}}%
\pgfpathmoveto{\pgfqpoint{4.809668in}{0.957130in}}%
\pgfpathlineto{\pgfqpoint{4.809668in}{0.957130in}}%
\pgfpathlineto{\pgfqpoint{4.809668in}{0.960080in}}%
\pgfpathlineto{\pgfqpoint{4.814209in}{0.960080in}}%
\pgfpathlineto{\pgfqpoint{4.814209in}{0.957130in}}%
\pgfpathmoveto{\pgfqpoint{4.814209in}{0.954181in}}%
\pgfpathlineto{\pgfqpoint{4.814209in}{0.954181in}}%
\pgfpathlineto{\pgfqpoint{4.814209in}{0.957130in}}%
\pgfpathlineto{\pgfqpoint{4.818750in}{0.957130in}}%
\pgfpathlineto{\pgfqpoint{4.818750in}{0.954181in}}%
\pgfpathmoveto{\pgfqpoint{4.814209in}{0.957130in}}%
\pgfpathlineto{\pgfqpoint{4.814209in}{0.957130in}}%
\pgfpathlineto{\pgfqpoint{4.814209in}{0.960080in}}%
\pgfpathlineto{\pgfqpoint{4.818750in}{0.960080in}}%
\pgfpathlineto{\pgfqpoint{4.818750in}{0.957130in}}%
\pgfpathmoveto{\pgfqpoint{4.805127in}{0.963029in}}%
\pgfpathlineto{\pgfqpoint{4.805127in}{0.963029in}}%
\pgfpathlineto{\pgfqpoint{4.805127in}{0.965978in}}%
\pgfpathlineto{\pgfqpoint{4.809668in}{0.965978in}}%
\pgfpathlineto{\pgfqpoint{4.809668in}{0.963029in}}%
\pgfpathmoveto{\pgfqpoint{4.800586in}{0.968927in}}%
\pgfpathlineto{\pgfqpoint{4.800586in}{0.968927in}}%
\pgfpathlineto{\pgfqpoint{4.800586in}{0.971876in}}%
\pgfpathlineto{\pgfqpoint{4.805127in}{0.971876in}}%
\pgfpathlineto{\pgfqpoint{4.805127in}{0.968927in}}%
\pgfpathmoveto{\pgfqpoint{4.805127in}{0.965978in}}%
\pgfpathlineto{\pgfqpoint{4.805127in}{0.965978in}}%
\pgfpathlineto{\pgfqpoint{4.805127in}{0.968927in}}%
\pgfpathlineto{\pgfqpoint{4.809668in}{0.968927in}}%
\pgfpathlineto{\pgfqpoint{4.809668in}{0.965978in}}%
\pgfpathmoveto{\pgfqpoint{4.805127in}{0.968927in}}%
\pgfpathlineto{\pgfqpoint{4.805127in}{0.968927in}}%
\pgfpathlineto{\pgfqpoint{4.805127in}{0.971876in}}%
\pgfpathlineto{\pgfqpoint{4.809668in}{0.971876in}}%
\pgfpathlineto{\pgfqpoint{4.809668in}{0.968927in}}%
\pgfpathmoveto{\pgfqpoint{4.809668in}{0.960080in}}%
\pgfpathlineto{\pgfqpoint{4.809668in}{0.960080in}}%
\pgfpathlineto{\pgfqpoint{4.809668in}{0.963029in}}%
\pgfpathlineto{\pgfqpoint{4.814209in}{0.963029in}}%
\pgfpathlineto{\pgfqpoint{4.814209in}{0.960080in}}%
\pgfpathmoveto{\pgfqpoint{4.809668in}{0.963029in}}%
\pgfpathlineto{\pgfqpoint{4.809668in}{0.963029in}}%
\pgfpathlineto{\pgfqpoint{4.809668in}{0.965978in}}%
\pgfpathlineto{\pgfqpoint{4.814209in}{0.965978in}}%
\pgfpathlineto{\pgfqpoint{4.814209in}{0.963029in}}%
\pgfpathmoveto{\pgfqpoint{4.741552in}{1.045607in}}%
\pgfpathlineto{\pgfqpoint{4.741552in}{1.045607in}}%
\pgfpathlineto{\pgfqpoint{4.741552in}{1.048556in}}%
\pgfpathlineto{\pgfqpoint{4.746093in}{1.048556in}}%
\pgfpathlineto{\pgfqpoint{4.746093in}{1.045607in}}%
\pgfpathmoveto{\pgfqpoint{4.737011in}{1.051505in}}%
\pgfpathlineto{\pgfqpoint{4.737011in}{1.051505in}}%
\pgfpathlineto{\pgfqpoint{4.737011in}{1.054454in}}%
\pgfpathlineto{\pgfqpoint{4.741552in}{1.054454in}}%
\pgfpathlineto{\pgfqpoint{4.741552in}{1.051505in}}%
\pgfpathmoveto{\pgfqpoint{4.741552in}{1.048556in}}%
\pgfpathlineto{\pgfqpoint{4.741552in}{1.048556in}}%
\pgfpathlineto{\pgfqpoint{4.741552in}{1.051505in}}%
\pgfpathlineto{\pgfqpoint{4.746093in}{1.051505in}}%
\pgfpathlineto{\pgfqpoint{4.746093in}{1.048556in}}%
\pgfpathmoveto{\pgfqpoint{4.741552in}{1.051505in}}%
\pgfpathlineto{\pgfqpoint{4.741552in}{1.051505in}}%
\pgfpathlineto{\pgfqpoint{4.741552in}{1.054454in}}%
\pgfpathlineto{\pgfqpoint{4.746093in}{1.054454in}}%
\pgfpathlineto{\pgfqpoint{4.746093in}{1.051505in}}%
\pgfpathmoveto{\pgfqpoint{4.732470in}{1.057404in}}%
\pgfpathlineto{\pgfqpoint{4.732470in}{1.057404in}}%
\pgfpathlineto{\pgfqpoint{4.732470in}{1.060353in}}%
\pgfpathlineto{\pgfqpoint{4.737011in}{1.060353in}}%
\pgfpathlineto{\pgfqpoint{4.737011in}{1.057404in}}%
\pgfpathmoveto{\pgfqpoint{4.727929in}{1.063302in}}%
\pgfpathlineto{\pgfqpoint{4.727929in}{1.063302in}}%
\pgfpathlineto{\pgfqpoint{4.727929in}{1.066251in}}%
\pgfpathlineto{\pgfqpoint{4.732470in}{1.066251in}}%
\pgfpathlineto{\pgfqpoint{4.732470in}{1.063302in}}%
\pgfpathmoveto{\pgfqpoint{4.732470in}{1.060353in}}%
\pgfpathlineto{\pgfqpoint{4.732470in}{1.060353in}}%
\pgfpathlineto{\pgfqpoint{4.732470in}{1.063302in}}%
\pgfpathlineto{\pgfqpoint{4.737011in}{1.063302in}}%
\pgfpathlineto{\pgfqpoint{4.737011in}{1.060353in}}%
\pgfpathmoveto{\pgfqpoint{4.732470in}{1.063302in}}%
\pgfpathlineto{\pgfqpoint{4.732470in}{1.063302in}}%
\pgfpathlineto{\pgfqpoint{4.732470in}{1.066251in}}%
\pgfpathlineto{\pgfqpoint{4.737011in}{1.066251in}}%
\pgfpathlineto{\pgfqpoint{4.737011in}{1.063302in}}%
\pgfpathmoveto{\pgfqpoint{4.737011in}{1.054454in}}%
\pgfpathlineto{\pgfqpoint{4.737011in}{1.054454in}}%
\pgfpathlineto{\pgfqpoint{4.737011in}{1.057404in}}%
\pgfpathlineto{\pgfqpoint{4.741552in}{1.057404in}}%
\pgfpathlineto{\pgfqpoint{4.741552in}{1.054454in}}%
\pgfpathmoveto{\pgfqpoint{4.737011in}{1.057404in}}%
\pgfpathlineto{\pgfqpoint{4.737011in}{1.057404in}}%
\pgfpathlineto{\pgfqpoint{4.737011in}{1.060353in}}%
\pgfpathlineto{\pgfqpoint{4.741552in}{1.060353in}}%
\pgfpathlineto{\pgfqpoint{4.741552in}{1.057404in}}%
\pgfpathmoveto{\pgfqpoint{4.777881in}{0.998419in}}%
\pgfpathlineto{\pgfqpoint{4.777881in}{0.998419in}}%
\pgfpathlineto{\pgfqpoint{4.777881in}{1.001369in}}%
\pgfpathlineto{\pgfqpoint{4.782422in}{1.001369in}}%
\pgfpathlineto{\pgfqpoint{4.782422in}{0.998419in}}%
\pgfpathmoveto{\pgfqpoint{4.773339in}{1.004318in}}%
\pgfpathlineto{\pgfqpoint{4.773339in}{1.004318in}}%
\pgfpathlineto{\pgfqpoint{4.773339in}{1.007267in}}%
\pgfpathlineto{\pgfqpoint{4.777881in}{1.007267in}}%
\pgfpathlineto{\pgfqpoint{4.777881in}{1.004318in}}%
\pgfpathmoveto{\pgfqpoint{4.777881in}{1.001369in}}%
\pgfpathlineto{\pgfqpoint{4.777881in}{1.001369in}}%
\pgfpathlineto{\pgfqpoint{4.777881in}{1.004318in}}%
\pgfpathlineto{\pgfqpoint{4.782422in}{1.004318in}}%
\pgfpathlineto{\pgfqpoint{4.782422in}{1.001369in}}%
\pgfpathmoveto{\pgfqpoint{4.777881in}{1.004318in}}%
\pgfpathlineto{\pgfqpoint{4.777881in}{1.004318in}}%
\pgfpathlineto{\pgfqpoint{4.777881in}{1.007267in}}%
\pgfpathlineto{\pgfqpoint{4.782422in}{1.007267in}}%
\pgfpathlineto{\pgfqpoint{4.782422in}{1.004318in}}%
\pgfpathmoveto{\pgfqpoint{4.768798in}{1.010216in}}%
\pgfpathlineto{\pgfqpoint{4.768798in}{1.010216in}}%
\pgfpathlineto{\pgfqpoint{4.768798in}{1.013165in}}%
\pgfpathlineto{\pgfqpoint{4.773339in}{1.013165in}}%
\pgfpathlineto{\pgfqpoint{4.773339in}{1.010216in}}%
\pgfpathmoveto{\pgfqpoint{4.764257in}{1.016115in}}%
\pgfpathlineto{\pgfqpoint{4.764257in}{1.016115in}}%
\pgfpathlineto{\pgfqpoint{4.764257in}{1.019064in}}%
\pgfpathlineto{\pgfqpoint{4.768798in}{1.019064in}}%
\pgfpathlineto{\pgfqpoint{4.768798in}{1.016115in}}%
\pgfpathmoveto{\pgfqpoint{4.768798in}{1.013165in}}%
\pgfpathlineto{\pgfqpoint{4.768798in}{1.013165in}}%
\pgfpathlineto{\pgfqpoint{4.768798in}{1.016115in}}%
\pgfpathlineto{\pgfqpoint{4.773339in}{1.016115in}}%
\pgfpathlineto{\pgfqpoint{4.773339in}{1.013165in}}%
\pgfpathmoveto{\pgfqpoint{4.768798in}{1.016115in}}%
\pgfpathlineto{\pgfqpoint{4.768798in}{1.016115in}}%
\pgfpathlineto{\pgfqpoint{4.768798in}{1.019064in}}%
\pgfpathlineto{\pgfqpoint{4.773339in}{1.019064in}}%
\pgfpathlineto{\pgfqpoint{4.773339in}{1.016115in}}%
\pgfpathmoveto{\pgfqpoint{4.773339in}{1.007267in}}%
\pgfpathlineto{\pgfqpoint{4.773339in}{1.007267in}}%
\pgfpathlineto{\pgfqpoint{4.773339in}{1.010216in}}%
\pgfpathlineto{\pgfqpoint{4.777881in}{1.010216in}}%
\pgfpathlineto{\pgfqpoint{4.777881in}{1.007267in}}%
\pgfpathmoveto{\pgfqpoint{4.773339in}{1.010216in}}%
\pgfpathlineto{\pgfqpoint{4.773339in}{1.010216in}}%
\pgfpathlineto{\pgfqpoint{4.773339in}{1.013165in}}%
\pgfpathlineto{\pgfqpoint{4.777881in}{1.013165in}}%
\pgfpathlineto{\pgfqpoint{4.777881in}{1.010216in}}%
\pgfpathmoveto{\pgfqpoint{4.796045in}{0.974826in}}%
\pgfpathlineto{\pgfqpoint{4.796045in}{0.974826in}}%
\pgfpathlineto{\pgfqpoint{4.796045in}{0.977775in}}%
\pgfpathlineto{\pgfqpoint{4.800586in}{0.977775in}}%
\pgfpathlineto{\pgfqpoint{4.800586in}{0.974826in}}%
\pgfpathmoveto{\pgfqpoint{4.791504in}{0.980724in}}%
\pgfpathlineto{\pgfqpoint{4.791504in}{0.980724in}}%
\pgfpathlineto{\pgfqpoint{4.791504in}{0.983673in}}%
\pgfpathlineto{\pgfqpoint{4.796045in}{0.983673in}}%
\pgfpathlineto{\pgfqpoint{4.796045in}{0.980724in}}%
\pgfpathmoveto{\pgfqpoint{4.796045in}{0.977775in}}%
\pgfpathlineto{\pgfqpoint{4.796045in}{0.977775in}}%
\pgfpathlineto{\pgfqpoint{4.796045in}{0.980724in}}%
\pgfpathlineto{\pgfqpoint{4.800586in}{0.980724in}}%
\pgfpathlineto{\pgfqpoint{4.800586in}{0.977775in}}%
\pgfpathmoveto{\pgfqpoint{4.796045in}{0.980724in}}%
\pgfpathlineto{\pgfqpoint{4.796045in}{0.980724in}}%
\pgfpathlineto{\pgfqpoint{4.796045in}{0.983673in}}%
\pgfpathlineto{\pgfqpoint{4.800586in}{0.983673in}}%
\pgfpathlineto{\pgfqpoint{4.800586in}{0.980724in}}%
\pgfpathmoveto{\pgfqpoint{4.786963in}{0.986622in}}%
\pgfpathlineto{\pgfqpoint{4.786963in}{0.986622in}}%
\pgfpathlineto{\pgfqpoint{4.786963in}{0.989572in}}%
\pgfpathlineto{\pgfqpoint{4.791504in}{0.989572in}}%
\pgfpathlineto{\pgfqpoint{4.791504in}{0.986622in}}%
\pgfpathmoveto{\pgfqpoint{4.782422in}{0.992521in}}%
\pgfpathlineto{\pgfqpoint{4.782422in}{0.992521in}}%
\pgfpathlineto{\pgfqpoint{4.782422in}{0.995470in}}%
\pgfpathlineto{\pgfqpoint{4.786963in}{0.995470in}}%
\pgfpathlineto{\pgfqpoint{4.786963in}{0.992521in}}%
\pgfpathmoveto{\pgfqpoint{4.786963in}{0.989572in}}%
\pgfpathlineto{\pgfqpoint{4.786963in}{0.989572in}}%
\pgfpathlineto{\pgfqpoint{4.786963in}{0.992521in}}%
\pgfpathlineto{\pgfqpoint{4.791504in}{0.992521in}}%
\pgfpathlineto{\pgfqpoint{4.791504in}{0.989572in}}%
\pgfpathmoveto{\pgfqpoint{4.786963in}{0.992521in}}%
\pgfpathlineto{\pgfqpoint{4.786963in}{0.992521in}}%
\pgfpathlineto{\pgfqpoint{4.786963in}{0.995470in}}%
\pgfpathlineto{\pgfqpoint{4.791504in}{0.995470in}}%
\pgfpathlineto{\pgfqpoint{4.791504in}{0.992521in}}%
\pgfpathmoveto{\pgfqpoint{4.791504in}{0.983673in}}%
\pgfpathlineto{\pgfqpoint{4.791504in}{0.983673in}}%
\pgfpathlineto{\pgfqpoint{4.791504in}{0.986622in}}%
\pgfpathlineto{\pgfqpoint{4.796045in}{0.986622in}}%
\pgfpathlineto{\pgfqpoint{4.796045in}{0.983673in}}%
\pgfpathmoveto{\pgfqpoint{4.791504in}{0.986622in}}%
\pgfpathlineto{\pgfqpoint{4.791504in}{0.986622in}}%
\pgfpathlineto{\pgfqpoint{4.791504in}{0.989572in}}%
\pgfpathlineto{\pgfqpoint{4.796045in}{0.989572in}}%
\pgfpathlineto{\pgfqpoint{4.796045in}{0.986622in}}%
\pgfpathmoveto{\pgfqpoint{4.800586in}{0.971876in}}%
\pgfpathlineto{\pgfqpoint{4.800586in}{0.971876in}}%
\pgfpathlineto{\pgfqpoint{4.800586in}{0.974826in}}%
\pgfpathlineto{\pgfqpoint{4.805127in}{0.974826in}}%
\pgfpathlineto{\pgfqpoint{4.805127in}{0.971876in}}%
\pgfpathmoveto{\pgfqpoint{4.800586in}{0.974826in}}%
\pgfpathlineto{\pgfqpoint{4.800586in}{0.974826in}}%
\pgfpathlineto{\pgfqpoint{4.800586in}{0.977775in}}%
\pgfpathlineto{\pgfqpoint{4.805127in}{0.977775in}}%
\pgfpathlineto{\pgfqpoint{4.805127in}{0.974826in}}%
\pgfpathmoveto{\pgfqpoint{4.782422in}{0.995470in}}%
\pgfpathlineto{\pgfqpoint{4.782422in}{0.995470in}}%
\pgfpathlineto{\pgfqpoint{4.782422in}{0.998419in}}%
\pgfpathlineto{\pgfqpoint{4.786963in}{0.998419in}}%
\pgfpathlineto{\pgfqpoint{4.786963in}{0.995470in}}%
\pgfpathmoveto{\pgfqpoint{4.782422in}{0.998419in}}%
\pgfpathlineto{\pgfqpoint{4.782422in}{0.998419in}}%
\pgfpathlineto{\pgfqpoint{4.782422in}{1.001369in}}%
\pgfpathlineto{\pgfqpoint{4.786963in}{1.001369in}}%
\pgfpathlineto{\pgfqpoint{4.786963in}{0.998419in}}%
\pgfpathmoveto{\pgfqpoint{4.759716in}{1.022013in}}%
\pgfpathlineto{\pgfqpoint{4.759716in}{1.022013in}}%
\pgfpathlineto{\pgfqpoint{4.759716in}{1.024962in}}%
\pgfpathlineto{\pgfqpoint{4.764257in}{1.024962in}}%
\pgfpathlineto{\pgfqpoint{4.764257in}{1.022013in}}%
\pgfpathmoveto{\pgfqpoint{4.755175in}{1.027911in}}%
\pgfpathlineto{\pgfqpoint{4.755175in}{1.027911in}}%
\pgfpathlineto{\pgfqpoint{4.755175in}{1.030861in}}%
\pgfpathlineto{\pgfqpoint{4.759716in}{1.030861in}}%
\pgfpathlineto{\pgfqpoint{4.759716in}{1.027911in}}%
\pgfpathmoveto{\pgfqpoint{4.759716in}{1.024962in}}%
\pgfpathlineto{\pgfqpoint{4.759716in}{1.024962in}}%
\pgfpathlineto{\pgfqpoint{4.759716in}{1.027911in}}%
\pgfpathlineto{\pgfqpoint{4.764257in}{1.027911in}}%
\pgfpathlineto{\pgfqpoint{4.764257in}{1.024962in}}%
\pgfpathmoveto{\pgfqpoint{4.759716in}{1.027911in}}%
\pgfpathlineto{\pgfqpoint{4.759716in}{1.027911in}}%
\pgfpathlineto{\pgfqpoint{4.759716in}{1.030861in}}%
\pgfpathlineto{\pgfqpoint{4.764257in}{1.030861in}}%
\pgfpathlineto{\pgfqpoint{4.764257in}{1.027911in}}%
\pgfpathmoveto{\pgfqpoint{4.750634in}{1.033810in}}%
\pgfpathlineto{\pgfqpoint{4.750634in}{1.033810in}}%
\pgfpathlineto{\pgfqpoint{4.750634in}{1.036759in}}%
\pgfpathlineto{\pgfqpoint{4.755175in}{1.036759in}}%
\pgfpathlineto{\pgfqpoint{4.755175in}{1.033810in}}%
\pgfpathmoveto{\pgfqpoint{4.746093in}{1.039708in}}%
\pgfpathlineto{\pgfqpoint{4.746093in}{1.039708in}}%
\pgfpathlineto{\pgfqpoint{4.746093in}{1.042658in}}%
\pgfpathlineto{\pgfqpoint{4.750634in}{1.042658in}}%
\pgfpathlineto{\pgfqpoint{4.750634in}{1.039708in}}%
\pgfpathmoveto{\pgfqpoint{4.750634in}{1.036759in}}%
\pgfpathlineto{\pgfqpoint{4.750634in}{1.036759in}}%
\pgfpathlineto{\pgfqpoint{4.750634in}{1.039708in}}%
\pgfpathlineto{\pgfqpoint{4.755175in}{1.039708in}}%
\pgfpathlineto{\pgfqpoint{4.755175in}{1.036759in}}%
\pgfpathmoveto{\pgfqpoint{4.750634in}{1.039708in}}%
\pgfpathlineto{\pgfqpoint{4.750634in}{1.039708in}}%
\pgfpathlineto{\pgfqpoint{4.750634in}{1.042658in}}%
\pgfpathlineto{\pgfqpoint{4.755175in}{1.042658in}}%
\pgfpathlineto{\pgfqpoint{4.755175in}{1.039708in}}%
\pgfpathmoveto{\pgfqpoint{4.755175in}{1.030861in}}%
\pgfpathlineto{\pgfqpoint{4.755175in}{1.030861in}}%
\pgfpathlineto{\pgfqpoint{4.755175in}{1.033810in}}%
\pgfpathlineto{\pgfqpoint{4.759716in}{1.033810in}}%
\pgfpathlineto{\pgfqpoint{4.759716in}{1.030861in}}%
\pgfpathmoveto{\pgfqpoint{4.755175in}{1.033810in}}%
\pgfpathlineto{\pgfqpoint{4.755175in}{1.033810in}}%
\pgfpathlineto{\pgfqpoint{4.755175in}{1.036759in}}%
\pgfpathlineto{\pgfqpoint{4.759716in}{1.036759in}}%
\pgfpathlineto{\pgfqpoint{4.759716in}{1.033810in}}%
\pgfpathmoveto{\pgfqpoint{4.764257in}{1.019064in}}%
\pgfpathlineto{\pgfqpoint{4.764257in}{1.019064in}}%
\pgfpathlineto{\pgfqpoint{4.764257in}{1.022013in}}%
\pgfpathlineto{\pgfqpoint{4.768798in}{1.022013in}}%
\pgfpathlineto{\pgfqpoint{4.768798in}{1.019064in}}%
\pgfpathmoveto{\pgfqpoint{4.764257in}{1.022013in}}%
\pgfpathlineto{\pgfqpoint{4.764257in}{1.022013in}}%
\pgfpathlineto{\pgfqpoint{4.764257in}{1.024962in}}%
\pgfpathlineto{\pgfqpoint{4.768798in}{1.024962in}}%
\pgfpathlineto{\pgfqpoint{4.768798in}{1.022013in}}%
\pgfpathmoveto{\pgfqpoint{4.746093in}{1.042658in}}%
\pgfpathlineto{\pgfqpoint{4.746093in}{1.042658in}}%
\pgfpathlineto{\pgfqpoint{4.746093in}{1.045607in}}%
\pgfpathlineto{\pgfqpoint{4.750634in}{1.045607in}}%
\pgfpathlineto{\pgfqpoint{4.750634in}{1.042658in}}%
\pgfpathmoveto{\pgfqpoint{4.746093in}{1.045607in}}%
\pgfpathlineto{\pgfqpoint{4.746093in}{1.045607in}}%
\pgfpathlineto{\pgfqpoint{4.746093in}{1.048556in}}%
\pgfpathlineto{\pgfqpoint{4.750634in}{1.048556in}}%
\pgfpathlineto{\pgfqpoint{4.750634in}{1.045607in}}%
\pgfpathmoveto{\pgfqpoint{4.705224in}{1.092794in}}%
\pgfpathlineto{\pgfqpoint{4.705224in}{1.092794in}}%
\pgfpathlineto{\pgfqpoint{4.705224in}{1.095743in}}%
\pgfpathlineto{\pgfqpoint{4.709765in}{1.095743in}}%
\pgfpathlineto{\pgfqpoint{4.709765in}{1.092794in}}%
\pgfpathmoveto{\pgfqpoint{4.700683in}{1.098692in}}%
\pgfpathlineto{\pgfqpoint{4.700683in}{1.098692in}}%
\pgfpathlineto{\pgfqpoint{4.700683in}{1.101641in}}%
\pgfpathlineto{\pgfqpoint{4.705224in}{1.101641in}}%
\pgfpathlineto{\pgfqpoint{4.705224in}{1.098692in}}%
\pgfpathmoveto{\pgfqpoint{4.705224in}{1.095743in}}%
\pgfpathlineto{\pgfqpoint{4.705224in}{1.095743in}}%
\pgfpathlineto{\pgfqpoint{4.705224in}{1.098692in}}%
\pgfpathlineto{\pgfqpoint{4.709765in}{1.098692in}}%
\pgfpathlineto{\pgfqpoint{4.709765in}{1.095743in}}%
\pgfpathmoveto{\pgfqpoint{4.705224in}{1.098692in}}%
\pgfpathlineto{\pgfqpoint{4.705224in}{1.098692in}}%
\pgfpathlineto{\pgfqpoint{4.705224in}{1.101641in}}%
\pgfpathlineto{\pgfqpoint{4.709765in}{1.101641in}}%
\pgfpathlineto{\pgfqpoint{4.709765in}{1.098692in}}%
\pgfpathmoveto{\pgfqpoint{4.696142in}{1.104590in}}%
\pgfpathlineto{\pgfqpoint{4.696142in}{1.104590in}}%
\pgfpathlineto{\pgfqpoint{4.696142in}{1.107540in}}%
\pgfpathlineto{\pgfqpoint{4.700683in}{1.107540in}}%
\pgfpathlineto{\pgfqpoint{4.700683in}{1.104590in}}%
\pgfpathmoveto{\pgfqpoint{4.691601in}{1.110489in}}%
\pgfpathlineto{\pgfqpoint{4.691601in}{1.110489in}}%
\pgfpathlineto{\pgfqpoint{4.691601in}{1.113438in}}%
\pgfpathlineto{\pgfqpoint{4.696142in}{1.113438in}}%
\pgfpathlineto{\pgfqpoint{4.696142in}{1.110489in}}%
\pgfpathmoveto{\pgfqpoint{4.696142in}{1.107540in}}%
\pgfpathlineto{\pgfqpoint{4.696142in}{1.107540in}}%
\pgfpathlineto{\pgfqpoint{4.696142in}{1.110489in}}%
\pgfpathlineto{\pgfqpoint{4.700683in}{1.110489in}}%
\pgfpathlineto{\pgfqpoint{4.700683in}{1.107540in}}%
\pgfpathmoveto{\pgfqpoint{4.696142in}{1.110489in}}%
\pgfpathlineto{\pgfqpoint{4.696142in}{1.110489in}}%
\pgfpathlineto{\pgfqpoint{4.696142in}{1.113438in}}%
\pgfpathlineto{\pgfqpoint{4.700683in}{1.113438in}}%
\pgfpathlineto{\pgfqpoint{4.700683in}{1.110489in}}%
\pgfpathmoveto{\pgfqpoint{4.700683in}{1.101641in}}%
\pgfpathlineto{\pgfqpoint{4.700683in}{1.101641in}}%
\pgfpathlineto{\pgfqpoint{4.700683in}{1.104590in}}%
\pgfpathlineto{\pgfqpoint{4.705224in}{1.104590in}}%
\pgfpathlineto{\pgfqpoint{4.705224in}{1.101641in}}%
\pgfpathmoveto{\pgfqpoint{4.700683in}{1.104590in}}%
\pgfpathlineto{\pgfqpoint{4.700683in}{1.104590in}}%
\pgfpathlineto{\pgfqpoint{4.700683in}{1.107540in}}%
\pgfpathlineto{\pgfqpoint{4.705224in}{1.107540in}}%
\pgfpathlineto{\pgfqpoint{4.705224in}{1.104590in}}%
\pgfpathmoveto{\pgfqpoint{4.723388in}{1.069200in}}%
\pgfpathlineto{\pgfqpoint{4.723388in}{1.069200in}}%
\pgfpathlineto{\pgfqpoint{4.723388in}{1.072150in}}%
\pgfpathlineto{\pgfqpoint{4.727929in}{1.072150in}}%
\pgfpathlineto{\pgfqpoint{4.727929in}{1.069200in}}%
\pgfpathmoveto{\pgfqpoint{4.718847in}{1.075099in}}%
\pgfpathlineto{\pgfqpoint{4.718847in}{1.075099in}}%
\pgfpathlineto{\pgfqpoint{4.718847in}{1.078048in}}%
\pgfpathlineto{\pgfqpoint{4.723388in}{1.078048in}}%
\pgfpathlineto{\pgfqpoint{4.723388in}{1.075099in}}%
\pgfpathmoveto{\pgfqpoint{4.723388in}{1.072150in}}%
\pgfpathlineto{\pgfqpoint{4.723388in}{1.072150in}}%
\pgfpathlineto{\pgfqpoint{4.723388in}{1.075099in}}%
\pgfpathlineto{\pgfqpoint{4.727929in}{1.075099in}}%
\pgfpathlineto{\pgfqpoint{4.727929in}{1.072150in}}%
\pgfpathmoveto{\pgfqpoint{4.723388in}{1.075099in}}%
\pgfpathlineto{\pgfqpoint{4.723388in}{1.075099in}}%
\pgfpathlineto{\pgfqpoint{4.723388in}{1.078048in}}%
\pgfpathlineto{\pgfqpoint{4.727929in}{1.078048in}}%
\pgfpathlineto{\pgfqpoint{4.727929in}{1.075099in}}%
\pgfpathmoveto{\pgfqpoint{4.714306in}{1.080997in}}%
\pgfpathlineto{\pgfqpoint{4.714306in}{1.080997in}}%
\pgfpathlineto{\pgfqpoint{4.714306in}{1.083946in}}%
\pgfpathlineto{\pgfqpoint{4.718847in}{1.083946in}}%
\pgfpathlineto{\pgfqpoint{4.718847in}{1.080997in}}%
\pgfpathmoveto{\pgfqpoint{4.709765in}{1.086895in}}%
\pgfpathlineto{\pgfqpoint{4.709765in}{1.086895in}}%
\pgfpathlineto{\pgfqpoint{4.709765in}{1.089845in}}%
\pgfpathlineto{\pgfqpoint{4.714306in}{1.089845in}}%
\pgfpathlineto{\pgfqpoint{4.714306in}{1.086895in}}%
\pgfpathmoveto{\pgfqpoint{4.714306in}{1.083946in}}%
\pgfpathlineto{\pgfqpoint{4.714306in}{1.083946in}}%
\pgfpathlineto{\pgfqpoint{4.714306in}{1.086895in}}%
\pgfpathlineto{\pgfqpoint{4.718847in}{1.086895in}}%
\pgfpathlineto{\pgfqpoint{4.718847in}{1.083946in}}%
\pgfpathmoveto{\pgfqpoint{4.714306in}{1.086895in}}%
\pgfpathlineto{\pgfqpoint{4.714306in}{1.086895in}}%
\pgfpathlineto{\pgfqpoint{4.714306in}{1.089845in}}%
\pgfpathlineto{\pgfqpoint{4.718847in}{1.089845in}}%
\pgfpathlineto{\pgfqpoint{4.718847in}{1.086895in}}%
\pgfpathmoveto{\pgfqpoint{4.718847in}{1.078048in}}%
\pgfpathlineto{\pgfqpoint{4.718847in}{1.078048in}}%
\pgfpathlineto{\pgfqpoint{4.718847in}{1.080997in}}%
\pgfpathlineto{\pgfqpoint{4.723388in}{1.080997in}}%
\pgfpathlineto{\pgfqpoint{4.723388in}{1.078048in}}%
\pgfpathmoveto{\pgfqpoint{4.718847in}{1.080997in}}%
\pgfpathlineto{\pgfqpoint{4.718847in}{1.080997in}}%
\pgfpathlineto{\pgfqpoint{4.718847in}{1.083946in}}%
\pgfpathlineto{\pgfqpoint{4.723388in}{1.083946in}}%
\pgfpathlineto{\pgfqpoint{4.723388in}{1.080997in}}%
\pgfpathmoveto{\pgfqpoint{4.727929in}{1.066251in}}%
\pgfpathlineto{\pgfqpoint{4.727929in}{1.066251in}}%
\pgfpathlineto{\pgfqpoint{4.727929in}{1.069200in}}%
\pgfpathlineto{\pgfqpoint{4.732470in}{1.069200in}}%
\pgfpathlineto{\pgfqpoint{4.732470in}{1.066251in}}%
\pgfpathmoveto{\pgfqpoint{4.727929in}{1.069200in}}%
\pgfpathlineto{\pgfqpoint{4.727929in}{1.069200in}}%
\pgfpathlineto{\pgfqpoint{4.727929in}{1.072150in}}%
\pgfpathlineto{\pgfqpoint{4.732470in}{1.072150in}}%
\pgfpathlineto{\pgfqpoint{4.732470in}{1.069200in}}%
\pgfpathmoveto{\pgfqpoint{4.709765in}{1.089845in}}%
\pgfpathlineto{\pgfqpoint{4.709765in}{1.089845in}}%
\pgfpathlineto{\pgfqpoint{4.709765in}{1.092794in}}%
\pgfpathlineto{\pgfqpoint{4.714306in}{1.092794in}}%
\pgfpathlineto{\pgfqpoint{4.714306in}{1.089845in}}%
\pgfpathmoveto{\pgfqpoint{4.709765in}{1.092794in}}%
\pgfpathlineto{\pgfqpoint{4.709765in}{1.092794in}}%
\pgfpathlineto{\pgfqpoint{4.709765in}{1.095743in}}%
\pgfpathlineto{\pgfqpoint{4.714306in}{1.095743in}}%
\pgfpathlineto{\pgfqpoint{4.714306in}{1.092794in}}%
\pgfpathmoveto{\pgfqpoint{4.687060in}{1.116387in}}%
\pgfpathlineto{\pgfqpoint{4.687060in}{1.116387in}}%
\pgfpathlineto{\pgfqpoint{4.687060in}{1.119336in}}%
\pgfpathlineto{\pgfqpoint{4.691601in}{1.119336in}}%
\pgfpathlineto{\pgfqpoint{4.691601in}{1.116387in}}%
\pgfpathmoveto{\pgfqpoint{4.682519in}{1.122285in}}%
\pgfpathlineto{\pgfqpoint{4.682519in}{1.122285in}}%
\pgfpathlineto{\pgfqpoint{4.682519in}{1.125235in}}%
\pgfpathlineto{\pgfqpoint{4.687060in}{1.125235in}}%
\pgfpathlineto{\pgfqpoint{4.687060in}{1.122285in}}%
\pgfpathmoveto{\pgfqpoint{4.687060in}{1.119336in}}%
\pgfpathlineto{\pgfqpoint{4.687060in}{1.119336in}}%
\pgfpathlineto{\pgfqpoint{4.687060in}{1.122285in}}%
\pgfpathlineto{\pgfqpoint{4.691601in}{1.122285in}}%
\pgfpathlineto{\pgfqpoint{4.691601in}{1.119336in}}%
\pgfpathmoveto{\pgfqpoint{4.687060in}{1.122285in}}%
\pgfpathlineto{\pgfqpoint{4.687060in}{1.122285in}}%
\pgfpathlineto{\pgfqpoint{4.687060in}{1.125235in}}%
\pgfpathlineto{\pgfqpoint{4.691601in}{1.125235in}}%
\pgfpathlineto{\pgfqpoint{4.691601in}{1.122285in}}%
\pgfpathmoveto{\pgfqpoint{4.677978in}{1.128184in}}%
\pgfpathlineto{\pgfqpoint{4.677978in}{1.128184in}}%
\pgfpathlineto{\pgfqpoint{4.677978in}{1.131133in}}%
\pgfpathlineto{\pgfqpoint{4.682519in}{1.131133in}}%
\pgfpathlineto{\pgfqpoint{4.682519in}{1.128184in}}%
\pgfpathmoveto{\pgfqpoint{4.673437in}{1.134082in}}%
\pgfpathlineto{\pgfqpoint{4.673437in}{1.134082in}}%
\pgfpathlineto{\pgfqpoint{4.673437in}{1.137031in}}%
\pgfpathlineto{\pgfqpoint{4.677978in}{1.137031in}}%
\pgfpathlineto{\pgfqpoint{4.677978in}{1.134082in}}%
\pgfpathmoveto{\pgfqpoint{4.677978in}{1.131133in}}%
\pgfpathlineto{\pgfqpoint{4.677978in}{1.131133in}}%
\pgfpathlineto{\pgfqpoint{4.677978in}{1.134082in}}%
\pgfpathlineto{\pgfqpoint{4.682519in}{1.134082in}}%
\pgfpathlineto{\pgfqpoint{4.682519in}{1.131133in}}%
\pgfpathmoveto{\pgfqpoint{4.677978in}{1.134082in}}%
\pgfpathlineto{\pgfqpoint{4.677978in}{1.134082in}}%
\pgfpathlineto{\pgfqpoint{4.677978in}{1.137031in}}%
\pgfpathlineto{\pgfqpoint{4.682519in}{1.137031in}}%
\pgfpathlineto{\pgfqpoint{4.682519in}{1.134082in}}%
\pgfpathmoveto{\pgfqpoint{4.682519in}{1.125235in}}%
\pgfpathlineto{\pgfqpoint{4.682519in}{1.125235in}}%
\pgfpathlineto{\pgfqpoint{4.682519in}{1.128184in}}%
\pgfpathlineto{\pgfqpoint{4.687060in}{1.128184in}}%
\pgfpathlineto{\pgfqpoint{4.687060in}{1.125235in}}%
\pgfpathmoveto{\pgfqpoint{4.682519in}{1.128184in}}%
\pgfpathlineto{\pgfqpoint{4.682519in}{1.128184in}}%
\pgfpathlineto{\pgfqpoint{4.682519in}{1.131133in}}%
\pgfpathlineto{\pgfqpoint{4.687060in}{1.131133in}}%
\pgfpathlineto{\pgfqpoint{4.687060in}{1.128184in}}%
\pgfpathmoveto{\pgfqpoint{4.691601in}{1.113438in}}%
\pgfpathlineto{\pgfqpoint{4.691601in}{1.113438in}}%
\pgfpathlineto{\pgfqpoint{4.691601in}{1.116387in}}%
\pgfpathlineto{\pgfqpoint{4.696142in}{1.116387in}}%
\pgfpathlineto{\pgfqpoint{4.696142in}{1.113438in}}%
\pgfpathmoveto{\pgfqpoint{4.691601in}{1.116387in}}%
\pgfpathlineto{\pgfqpoint{4.691601in}{1.116387in}}%
\pgfpathlineto{\pgfqpoint{4.691601in}{1.119336in}}%
\pgfpathlineto{\pgfqpoint{4.696142in}{1.119336in}}%
\pgfpathlineto{\pgfqpoint{4.696142in}{1.116387in}}%
\pgfpathmoveto{\pgfqpoint{4.673437in}{1.137031in}}%
\pgfpathlineto{\pgfqpoint{4.673437in}{1.137031in}}%
\pgfpathlineto{\pgfqpoint{4.673437in}{1.139980in}}%
\pgfpathlineto{\pgfqpoint{4.677978in}{1.139980in}}%
\pgfpathlineto{\pgfqpoint{4.677978in}{1.137031in}}%
\pgfpathmoveto{\pgfqpoint{4.673437in}{1.139980in}}%
\pgfpathlineto{\pgfqpoint{4.673437in}{1.139980in}}%
\pgfpathlineto{\pgfqpoint{4.673437in}{1.142930in}}%
\pgfpathlineto{\pgfqpoint{4.677978in}{1.142930in}}%
\pgfpathlineto{\pgfqpoint{4.677978in}{1.139980in}}%
\pgfpathmoveto{\pgfqpoint{4.959517in}{0.762481in}}%
\pgfpathlineto{\pgfqpoint{4.959517in}{0.762481in}}%
\pgfpathlineto{\pgfqpoint{4.959517in}{0.765430in}}%
\pgfpathlineto{\pgfqpoint{4.964058in}{0.765430in}}%
\pgfpathlineto{\pgfqpoint{4.964058in}{0.762481in}}%
\pgfpathmoveto{\pgfqpoint{4.954976in}{0.768379in}}%
\pgfpathlineto{\pgfqpoint{4.954976in}{0.768379in}}%
\pgfpathlineto{\pgfqpoint{4.954976in}{0.771329in}}%
\pgfpathlineto{\pgfqpoint{4.959517in}{0.771329in}}%
\pgfpathlineto{\pgfqpoint{4.959517in}{0.768379in}}%
\pgfpathmoveto{\pgfqpoint{4.959517in}{0.765430in}}%
\pgfpathlineto{\pgfqpoint{4.959517in}{0.765430in}}%
\pgfpathlineto{\pgfqpoint{4.959517in}{0.768379in}}%
\pgfpathlineto{\pgfqpoint{4.964058in}{0.768379in}}%
\pgfpathlineto{\pgfqpoint{4.964058in}{0.765430in}}%
\pgfpathmoveto{\pgfqpoint{4.959517in}{0.768379in}}%
\pgfpathlineto{\pgfqpoint{4.959517in}{0.768379in}}%
\pgfpathlineto{\pgfqpoint{4.959517in}{0.771329in}}%
\pgfpathlineto{\pgfqpoint{4.964058in}{0.771329in}}%
\pgfpathlineto{\pgfqpoint{4.964058in}{0.768379in}}%
\pgfpathmoveto{\pgfqpoint{4.950436in}{0.774278in}}%
\pgfpathlineto{\pgfqpoint{4.950436in}{0.774278in}}%
\pgfpathlineto{\pgfqpoint{4.950436in}{0.777227in}}%
\pgfpathlineto{\pgfqpoint{4.954976in}{0.777227in}}%
\pgfpathlineto{\pgfqpoint{4.954976in}{0.774278in}}%
\pgfpathmoveto{\pgfqpoint{4.945895in}{0.780177in}}%
\pgfpathlineto{\pgfqpoint{4.945895in}{0.780177in}}%
\pgfpathlineto{\pgfqpoint{4.945895in}{0.783126in}}%
\pgfpathlineto{\pgfqpoint{4.950436in}{0.783126in}}%
\pgfpathlineto{\pgfqpoint{4.950436in}{0.780177in}}%
\pgfpathmoveto{\pgfqpoint{4.950436in}{0.777227in}}%
\pgfpathlineto{\pgfqpoint{4.950436in}{0.777227in}}%
\pgfpathlineto{\pgfqpoint{4.950436in}{0.780177in}}%
\pgfpathlineto{\pgfqpoint{4.954976in}{0.780177in}}%
\pgfpathlineto{\pgfqpoint{4.954976in}{0.777227in}}%
\pgfpathmoveto{\pgfqpoint{4.950436in}{0.780177in}}%
\pgfpathlineto{\pgfqpoint{4.950436in}{0.780177in}}%
\pgfpathlineto{\pgfqpoint{4.950436in}{0.783126in}}%
\pgfpathlineto{\pgfqpoint{4.954976in}{0.783126in}}%
\pgfpathlineto{\pgfqpoint{4.954976in}{0.780177in}}%
\pgfpathmoveto{\pgfqpoint{4.954976in}{0.771329in}}%
\pgfpathlineto{\pgfqpoint{4.954976in}{0.771329in}}%
\pgfpathlineto{\pgfqpoint{4.954976in}{0.774278in}}%
\pgfpathlineto{\pgfqpoint{4.959517in}{0.774278in}}%
\pgfpathlineto{\pgfqpoint{4.959517in}{0.771329in}}%
\pgfpathmoveto{\pgfqpoint{4.954976in}{0.774278in}}%
\pgfpathlineto{\pgfqpoint{4.954976in}{0.774278in}}%
\pgfpathlineto{\pgfqpoint{4.954976in}{0.777227in}}%
\pgfpathlineto{\pgfqpoint{4.959517in}{0.777227in}}%
\pgfpathlineto{\pgfqpoint{4.959517in}{0.774278in}}%
\pgfpathmoveto{\pgfqpoint{4.886863in}{0.856857in}}%
\pgfpathlineto{\pgfqpoint{4.886863in}{0.856857in}}%
\pgfpathlineto{\pgfqpoint{4.886863in}{0.859807in}}%
\pgfpathlineto{\pgfqpoint{4.891404in}{0.859807in}}%
\pgfpathlineto{\pgfqpoint{4.891404in}{0.856857in}}%
\pgfpathmoveto{\pgfqpoint{4.882322in}{0.862756in}}%
\pgfpathlineto{\pgfqpoint{4.882322in}{0.862756in}}%
\pgfpathlineto{\pgfqpoint{4.882322in}{0.865705in}}%
\pgfpathlineto{\pgfqpoint{4.886863in}{0.865705in}}%
\pgfpathlineto{\pgfqpoint{4.886863in}{0.862756in}}%
\pgfpathmoveto{\pgfqpoint{4.886863in}{0.859807in}}%
\pgfpathlineto{\pgfqpoint{4.886863in}{0.859807in}}%
\pgfpathlineto{\pgfqpoint{4.886863in}{0.862756in}}%
\pgfpathlineto{\pgfqpoint{4.891404in}{0.862756in}}%
\pgfpathlineto{\pgfqpoint{4.891404in}{0.859807in}}%
\pgfpathmoveto{\pgfqpoint{4.886863in}{0.862756in}}%
\pgfpathlineto{\pgfqpoint{4.886863in}{0.862756in}}%
\pgfpathlineto{\pgfqpoint{4.886863in}{0.865705in}}%
\pgfpathlineto{\pgfqpoint{4.891404in}{0.865705in}}%
\pgfpathlineto{\pgfqpoint{4.891404in}{0.862756in}}%
\pgfpathmoveto{\pgfqpoint{4.877781in}{0.868654in}}%
\pgfpathlineto{\pgfqpoint{4.877781in}{0.868654in}}%
\pgfpathlineto{\pgfqpoint{4.877781in}{0.871604in}}%
\pgfpathlineto{\pgfqpoint{4.882322in}{0.871604in}}%
\pgfpathlineto{\pgfqpoint{4.882322in}{0.868654in}}%
\pgfpathmoveto{\pgfqpoint{4.873240in}{0.874553in}}%
\pgfpathlineto{\pgfqpoint{4.873240in}{0.874553in}}%
\pgfpathlineto{\pgfqpoint{4.873240in}{0.877502in}}%
\pgfpathlineto{\pgfqpoint{4.877781in}{0.877502in}}%
\pgfpathlineto{\pgfqpoint{4.877781in}{0.874553in}}%
\pgfpathmoveto{\pgfqpoint{4.877781in}{0.871604in}}%
\pgfpathlineto{\pgfqpoint{4.877781in}{0.871604in}}%
\pgfpathlineto{\pgfqpoint{4.877781in}{0.874553in}}%
\pgfpathlineto{\pgfqpoint{4.882322in}{0.874553in}}%
\pgfpathlineto{\pgfqpoint{4.882322in}{0.871604in}}%
\pgfpathmoveto{\pgfqpoint{4.877781in}{0.874553in}}%
\pgfpathlineto{\pgfqpoint{4.877781in}{0.874553in}}%
\pgfpathlineto{\pgfqpoint{4.877781in}{0.877502in}}%
\pgfpathlineto{\pgfqpoint{4.882322in}{0.877502in}}%
\pgfpathlineto{\pgfqpoint{4.882322in}{0.874553in}}%
\pgfpathmoveto{\pgfqpoint{4.882322in}{0.865705in}}%
\pgfpathlineto{\pgfqpoint{4.882322in}{0.865705in}}%
\pgfpathlineto{\pgfqpoint{4.882322in}{0.868654in}}%
\pgfpathlineto{\pgfqpoint{4.886863in}{0.868654in}}%
\pgfpathlineto{\pgfqpoint{4.886863in}{0.865705in}}%
\pgfpathmoveto{\pgfqpoint{4.882322in}{0.868654in}}%
\pgfpathlineto{\pgfqpoint{4.882322in}{0.868654in}}%
\pgfpathlineto{\pgfqpoint{4.882322in}{0.871604in}}%
\pgfpathlineto{\pgfqpoint{4.886863in}{0.871604in}}%
\pgfpathlineto{\pgfqpoint{4.886863in}{0.868654in}}%
\pgfpathmoveto{\pgfqpoint{4.923190in}{0.809669in}}%
\pgfpathlineto{\pgfqpoint{4.923190in}{0.809669in}}%
\pgfpathlineto{\pgfqpoint{4.923190in}{0.812618in}}%
\pgfpathlineto{\pgfqpoint{4.927731in}{0.812618in}}%
\pgfpathlineto{\pgfqpoint{4.927731in}{0.809669in}}%
\pgfpathmoveto{\pgfqpoint{4.918649in}{0.815568in}}%
\pgfpathlineto{\pgfqpoint{4.918649in}{0.815568in}}%
\pgfpathlineto{\pgfqpoint{4.918649in}{0.818517in}}%
\pgfpathlineto{\pgfqpoint{4.923190in}{0.818517in}}%
\pgfpathlineto{\pgfqpoint{4.923190in}{0.815568in}}%
\pgfpathmoveto{\pgfqpoint{4.923190in}{0.812618in}}%
\pgfpathlineto{\pgfqpoint{4.923190in}{0.812618in}}%
\pgfpathlineto{\pgfqpoint{4.923190in}{0.815568in}}%
\pgfpathlineto{\pgfqpoint{4.927731in}{0.815568in}}%
\pgfpathlineto{\pgfqpoint{4.927731in}{0.812618in}}%
\pgfpathmoveto{\pgfqpoint{4.923190in}{0.815568in}}%
\pgfpathlineto{\pgfqpoint{4.923190in}{0.815568in}}%
\pgfpathlineto{\pgfqpoint{4.923190in}{0.818517in}}%
\pgfpathlineto{\pgfqpoint{4.927731in}{0.818517in}}%
\pgfpathlineto{\pgfqpoint{4.927731in}{0.815568in}}%
\pgfpathmoveto{\pgfqpoint{4.914108in}{0.821466in}}%
\pgfpathlineto{\pgfqpoint{4.914108in}{0.821466in}}%
\pgfpathlineto{\pgfqpoint{4.914108in}{0.824415in}}%
\pgfpathlineto{\pgfqpoint{4.918649in}{0.824415in}}%
\pgfpathlineto{\pgfqpoint{4.918649in}{0.821466in}}%
\pgfpathmoveto{\pgfqpoint{4.909568in}{0.827365in}}%
\pgfpathlineto{\pgfqpoint{4.909568in}{0.827365in}}%
\pgfpathlineto{\pgfqpoint{4.909568in}{0.830314in}}%
\pgfpathlineto{\pgfqpoint{4.914108in}{0.830314in}}%
\pgfpathlineto{\pgfqpoint{4.914108in}{0.827365in}}%
\pgfpathmoveto{\pgfqpoint{4.914108in}{0.824415in}}%
\pgfpathlineto{\pgfqpoint{4.914108in}{0.824415in}}%
\pgfpathlineto{\pgfqpoint{4.914108in}{0.827365in}}%
\pgfpathlineto{\pgfqpoint{4.918649in}{0.827365in}}%
\pgfpathlineto{\pgfqpoint{4.918649in}{0.824415in}}%
\pgfpathmoveto{\pgfqpoint{4.914108in}{0.827365in}}%
\pgfpathlineto{\pgfqpoint{4.914108in}{0.827365in}}%
\pgfpathlineto{\pgfqpoint{4.914108in}{0.830314in}}%
\pgfpathlineto{\pgfqpoint{4.918649in}{0.830314in}}%
\pgfpathlineto{\pgfqpoint{4.918649in}{0.827365in}}%
\pgfpathmoveto{\pgfqpoint{4.918649in}{0.818517in}}%
\pgfpathlineto{\pgfqpoint{4.918649in}{0.818517in}}%
\pgfpathlineto{\pgfqpoint{4.918649in}{0.821466in}}%
\pgfpathlineto{\pgfqpoint{4.923190in}{0.821466in}}%
\pgfpathlineto{\pgfqpoint{4.923190in}{0.818517in}}%
\pgfpathmoveto{\pgfqpoint{4.918649in}{0.821466in}}%
\pgfpathlineto{\pgfqpoint{4.918649in}{0.821466in}}%
\pgfpathlineto{\pgfqpoint{4.918649in}{0.824415in}}%
\pgfpathlineto{\pgfqpoint{4.923190in}{0.824415in}}%
\pgfpathlineto{\pgfqpoint{4.923190in}{0.821466in}}%
\pgfpathmoveto{\pgfqpoint{4.941354in}{0.786075in}}%
\pgfpathlineto{\pgfqpoint{4.941354in}{0.786075in}}%
\pgfpathlineto{\pgfqpoint{4.941354in}{0.789024in}}%
\pgfpathlineto{\pgfqpoint{4.945895in}{0.789024in}}%
\pgfpathlineto{\pgfqpoint{4.945895in}{0.786075in}}%
\pgfpathmoveto{\pgfqpoint{4.936813in}{0.791974in}}%
\pgfpathlineto{\pgfqpoint{4.936813in}{0.791974in}}%
\pgfpathlineto{\pgfqpoint{4.936813in}{0.794923in}}%
\pgfpathlineto{\pgfqpoint{4.941354in}{0.794923in}}%
\pgfpathlineto{\pgfqpoint{4.941354in}{0.791974in}}%
\pgfpathmoveto{\pgfqpoint{4.941354in}{0.789024in}}%
\pgfpathlineto{\pgfqpoint{4.941354in}{0.789024in}}%
\pgfpathlineto{\pgfqpoint{4.941354in}{0.791974in}}%
\pgfpathlineto{\pgfqpoint{4.945895in}{0.791974in}}%
\pgfpathlineto{\pgfqpoint{4.945895in}{0.789024in}}%
\pgfpathmoveto{\pgfqpoint{4.941354in}{0.791974in}}%
\pgfpathlineto{\pgfqpoint{4.941354in}{0.791974in}}%
\pgfpathlineto{\pgfqpoint{4.941354in}{0.794923in}}%
\pgfpathlineto{\pgfqpoint{4.945895in}{0.794923in}}%
\pgfpathlineto{\pgfqpoint{4.945895in}{0.791974in}}%
\pgfpathmoveto{\pgfqpoint{4.932272in}{0.797872in}}%
\pgfpathlineto{\pgfqpoint{4.932272in}{0.797872in}}%
\pgfpathlineto{\pgfqpoint{4.932272in}{0.800821in}}%
\pgfpathlineto{\pgfqpoint{4.936813in}{0.800821in}}%
\pgfpathlineto{\pgfqpoint{4.936813in}{0.797872in}}%
\pgfpathmoveto{\pgfqpoint{4.927731in}{0.803771in}}%
\pgfpathlineto{\pgfqpoint{4.927731in}{0.803771in}}%
\pgfpathlineto{\pgfqpoint{4.927731in}{0.806720in}}%
\pgfpathlineto{\pgfqpoint{4.932272in}{0.806720in}}%
\pgfpathlineto{\pgfqpoint{4.932272in}{0.803771in}}%
\pgfpathmoveto{\pgfqpoint{4.932272in}{0.800821in}}%
\pgfpathlineto{\pgfqpoint{4.932272in}{0.800821in}}%
\pgfpathlineto{\pgfqpoint{4.932272in}{0.803771in}}%
\pgfpathlineto{\pgfqpoint{4.936813in}{0.803771in}}%
\pgfpathlineto{\pgfqpoint{4.936813in}{0.800821in}}%
\pgfpathmoveto{\pgfqpoint{4.932272in}{0.803771in}}%
\pgfpathlineto{\pgfqpoint{4.932272in}{0.803771in}}%
\pgfpathlineto{\pgfqpoint{4.932272in}{0.806720in}}%
\pgfpathlineto{\pgfqpoint{4.936813in}{0.806720in}}%
\pgfpathlineto{\pgfqpoint{4.936813in}{0.803771in}}%
\pgfpathmoveto{\pgfqpoint{4.936813in}{0.794923in}}%
\pgfpathlineto{\pgfqpoint{4.936813in}{0.794923in}}%
\pgfpathlineto{\pgfqpoint{4.936813in}{0.797872in}}%
\pgfpathlineto{\pgfqpoint{4.941354in}{0.797872in}}%
\pgfpathlineto{\pgfqpoint{4.941354in}{0.794923in}}%
\pgfpathmoveto{\pgfqpoint{4.936813in}{0.797872in}}%
\pgfpathlineto{\pgfqpoint{4.936813in}{0.797872in}}%
\pgfpathlineto{\pgfqpoint{4.936813in}{0.800821in}}%
\pgfpathlineto{\pgfqpoint{4.941354in}{0.800821in}}%
\pgfpathlineto{\pgfqpoint{4.941354in}{0.797872in}}%
\pgfpathmoveto{\pgfqpoint{4.945895in}{0.783126in}}%
\pgfpathlineto{\pgfqpoint{4.945895in}{0.783126in}}%
\pgfpathlineto{\pgfqpoint{4.945895in}{0.786075in}}%
\pgfpathlineto{\pgfqpoint{4.950436in}{0.786075in}}%
\pgfpathlineto{\pgfqpoint{4.950436in}{0.783126in}}%
\pgfpathmoveto{\pgfqpoint{4.945895in}{0.786075in}}%
\pgfpathlineto{\pgfqpoint{4.945895in}{0.786075in}}%
\pgfpathlineto{\pgfqpoint{4.945895in}{0.789024in}}%
\pgfpathlineto{\pgfqpoint{4.950436in}{0.789024in}}%
\pgfpathlineto{\pgfqpoint{4.950436in}{0.786075in}}%
\pgfpathmoveto{\pgfqpoint{4.927731in}{0.806720in}}%
\pgfpathlineto{\pgfqpoint{4.927731in}{0.806720in}}%
\pgfpathlineto{\pgfqpoint{4.927731in}{0.809669in}}%
\pgfpathlineto{\pgfqpoint{4.932272in}{0.809669in}}%
\pgfpathlineto{\pgfqpoint{4.932272in}{0.806720in}}%
\pgfpathmoveto{\pgfqpoint{4.927731in}{0.809669in}}%
\pgfpathlineto{\pgfqpoint{4.927731in}{0.809669in}}%
\pgfpathlineto{\pgfqpoint{4.927731in}{0.812618in}}%
\pgfpathlineto{\pgfqpoint{4.932272in}{0.812618in}}%
\pgfpathlineto{\pgfqpoint{4.932272in}{0.809669in}}%
\pgfpathmoveto{\pgfqpoint{4.905027in}{0.833263in}}%
\pgfpathlineto{\pgfqpoint{4.905027in}{0.833263in}}%
\pgfpathlineto{\pgfqpoint{4.905027in}{0.836212in}}%
\pgfpathlineto{\pgfqpoint{4.909568in}{0.836212in}}%
\pgfpathlineto{\pgfqpoint{4.909568in}{0.833263in}}%
\pgfpathmoveto{\pgfqpoint{4.900486in}{0.839162in}}%
\pgfpathlineto{\pgfqpoint{4.900486in}{0.839162in}}%
\pgfpathlineto{\pgfqpoint{4.900486in}{0.842111in}}%
\pgfpathlineto{\pgfqpoint{4.905027in}{0.842111in}}%
\pgfpathlineto{\pgfqpoint{4.905027in}{0.839162in}}%
\pgfpathmoveto{\pgfqpoint{4.905027in}{0.836212in}}%
\pgfpathlineto{\pgfqpoint{4.905027in}{0.836212in}}%
\pgfpathlineto{\pgfqpoint{4.905027in}{0.839162in}}%
\pgfpathlineto{\pgfqpoint{4.909568in}{0.839162in}}%
\pgfpathlineto{\pgfqpoint{4.909568in}{0.836212in}}%
\pgfpathmoveto{\pgfqpoint{4.905027in}{0.839162in}}%
\pgfpathlineto{\pgfqpoint{4.905027in}{0.839162in}}%
\pgfpathlineto{\pgfqpoint{4.905027in}{0.842111in}}%
\pgfpathlineto{\pgfqpoint{4.909568in}{0.842111in}}%
\pgfpathlineto{\pgfqpoint{4.909568in}{0.839162in}}%
\pgfpathmoveto{\pgfqpoint{4.895945in}{0.845060in}}%
\pgfpathlineto{\pgfqpoint{4.895945in}{0.845060in}}%
\pgfpathlineto{\pgfqpoint{4.895945in}{0.848010in}}%
\pgfpathlineto{\pgfqpoint{4.900486in}{0.848010in}}%
\pgfpathlineto{\pgfqpoint{4.900486in}{0.845060in}}%
\pgfpathmoveto{\pgfqpoint{4.891404in}{0.850959in}}%
\pgfpathlineto{\pgfqpoint{4.891404in}{0.850959in}}%
\pgfpathlineto{\pgfqpoint{4.891404in}{0.853908in}}%
\pgfpathlineto{\pgfqpoint{4.895945in}{0.853908in}}%
\pgfpathlineto{\pgfqpoint{4.895945in}{0.850959in}}%
\pgfpathmoveto{\pgfqpoint{4.895945in}{0.848010in}}%
\pgfpathlineto{\pgfqpoint{4.895945in}{0.848010in}}%
\pgfpathlineto{\pgfqpoint{4.895945in}{0.850959in}}%
\pgfpathlineto{\pgfqpoint{4.900486in}{0.850959in}}%
\pgfpathlineto{\pgfqpoint{4.900486in}{0.848010in}}%
\pgfpathmoveto{\pgfqpoint{4.895945in}{0.850959in}}%
\pgfpathlineto{\pgfqpoint{4.895945in}{0.850959in}}%
\pgfpathlineto{\pgfqpoint{4.895945in}{0.853908in}}%
\pgfpathlineto{\pgfqpoint{4.900486in}{0.853908in}}%
\pgfpathlineto{\pgfqpoint{4.900486in}{0.850959in}}%
\pgfpathmoveto{\pgfqpoint{4.900486in}{0.842111in}}%
\pgfpathlineto{\pgfqpoint{4.900486in}{0.842111in}}%
\pgfpathlineto{\pgfqpoint{4.900486in}{0.845060in}}%
\pgfpathlineto{\pgfqpoint{4.905027in}{0.845060in}}%
\pgfpathlineto{\pgfqpoint{4.905027in}{0.842111in}}%
\pgfpathmoveto{\pgfqpoint{4.900486in}{0.845060in}}%
\pgfpathlineto{\pgfqpoint{4.900486in}{0.845060in}}%
\pgfpathlineto{\pgfqpoint{4.900486in}{0.848010in}}%
\pgfpathlineto{\pgfqpoint{4.905027in}{0.848010in}}%
\pgfpathlineto{\pgfqpoint{4.905027in}{0.845060in}}%
\pgfpathmoveto{\pgfqpoint{4.909568in}{0.830314in}}%
\pgfpathlineto{\pgfqpoint{4.909568in}{0.830314in}}%
\pgfpathlineto{\pgfqpoint{4.909568in}{0.833263in}}%
\pgfpathlineto{\pgfqpoint{4.914108in}{0.833263in}}%
\pgfpathlineto{\pgfqpoint{4.914108in}{0.830314in}}%
\pgfpathmoveto{\pgfqpoint{4.909568in}{0.833263in}}%
\pgfpathlineto{\pgfqpoint{4.909568in}{0.833263in}}%
\pgfpathlineto{\pgfqpoint{4.909568in}{0.836212in}}%
\pgfpathlineto{\pgfqpoint{4.914108in}{0.836212in}}%
\pgfpathlineto{\pgfqpoint{4.914108in}{0.833263in}}%
\pgfpathmoveto{\pgfqpoint{4.891404in}{0.853908in}}%
\pgfpathlineto{\pgfqpoint{4.891404in}{0.853908in}}%
\pgfpathlineto{\pgfqpoint{4.891404in}{0.856857in}}%
\pgfpathlineto{\pgfqpoint{4.895945in}{0.856857in}}%
\pgfpathlineto{\pgfqpoint{4.895945in}{0.853908in}}%
\pgfpathmoveto{\pgfqpoint{4.891404in}{0.856857in}}%
\pgfpathlineto{\pgfqpoint{4.891404in}{0.856857in}}%
\pgfpathlineto{\pgfqpoint{4.891404in}{0.859807in}}%
\pgfpathlineto{\pgfqpoint{4.895945in}{0.859807in}}%
\pgfpathlineto{\pgfqpoint{4.895945in}{0.856857in}}%
\pgfpathmoveto{\pgfqpoint{4.850536in}{0.904045in}}%
\pgfpathlineto{\pgfqpoint{4.850536in}{0.904045in}}%
\pgfpathlineto{\pgfqpoint{4.850536in}{0.906994in}}%
\pgfpathlineto{\pgfqpoint{4.855077in}{0.906994in}}%
\pgfpathlineto{\pgfqpoint{4.855077in}{0.904045in}}%
\pgfpathmoveto{\pgfqpoint{4.845995in}{0.909943in}}%
\pgfpathlineto{\pgfqpoint{4.845995in}{0.909943in}}%
\pgfpathlineto{\pgfqpoint{4.845995in}{0.912892in}}%
\pgfpathlineto{\pgfqpoint{4.850536in}{0.912892in}}%
\pgfpathlineto{\pgfqpoint{4.850536in}{0.909943in}}%
\pgfpathmoveto{\pgfqpoint{4.850536in}{0.906994in}}%
\pgfpathlineto{\pgfqpoint{4.850536in}{0.906994in}}%
\pgfpathlineto{\pgfqpoint{4.850536in}{0.909943in}}%
\pgfpathlineto{\pgfqpoint{4.855077in}{0.909943in}}%
\pgfpathlineto{\pgfqpoint{4.855077in}{0.906994in}}%
\pgfpathmoveto{\pgfqpoint{4.850536in}{0.909943in}}%
\pgfpathlineto{\pgfqpoint{4.850536in}{0.909943in}}%
\pgfpathlineto{\pgfqpoint{4.850536in}{0.912892in}}%
\pgfpathlineto{\pgfqpoint{4.855077in}{0.912892in}}%
\pgfpathlineto{\pgfqpoint{4.855077in}{0.909943in}}%
\pgfpathmoveto{\pgfqpoint{4.841454in}{0.915842in}}%
\pgfpathlineto{\pgfqpoint{4.841454in}{0.915842in}}%
\pgfpathlineto{\pgfqpoint{4.841454in}{0.918791in}}%
\pgfpathlineto{\pgfqpoint{4.845995in}{0.918791in}}%
\pgfpathlineto{\pgfqpoint{4.845995in}{0.915842in}}%
\pgfpathmoveto{\pgfqpoint{4.836913in}{0.921740in}}%
\pgfpathlineto{\pgfqpoint{4.836913in}{0.921740in}}%
\pgfpathlineto{\pgfqpoint{4.836913in}{0.924689in}}%
\pgfpathlineto{\pgfqpoint{4.841454in}{0.924689in}}%
\pgfpathlineto{\pgfqpoint{4.841454in}{0.921740in}}%
\pgfpathmoveto{\pgfqpoint{4.841454in}{0.918791in}}%
\pgfpathlineto{\pgfqpoint{4.841454in}{0.918791in}}%
\pgfpathlineto{\pgfqpoint{4.841454in}{0.921740in}}%
\pgfpathlineto{\pgfqpoint{4.845995in}{0.921740in}}%
\pgfpathlineto{\pgfqpoint{4.845995in}{0.918791in}}%
\pgfpathmoveto{\pgfqpoint{4.841454in}{0.921740in}}%
\pgfpathlineto{\pgfqpoint{4.841454in}{0.921740in}}%
\pgfpathlineto{\pgfqpoint{4.841454in}{0.924689in}}%
\pgfpathlineto{\pgfqpoint{4.845995in}{0.924689in}}%
\pgfpathlineto{\pgfqpoint{4.845995in}{0.921740in}}%
\pgfpathmoveto{\pgfqpoint{4.845995in}{0.912892in}}%
\pgfpathlineto{\pgfqpoint{4.845995in}{0.912892in}}%
\pgfpathlineto{\pgfqpoint{4.845995in}{0.915842in}}%
\pgfpathlineto{\pgfqpoint{4.850536in}{0.915842in}}%
\pgfpathlineto{\pgfqpoint{4.850536in}{0.912892in}}%
\pgfpathmoveto{\pgfqpoint{4.845995in}{0.915842in}}%
\pgfpathlineto{\pgfqpoint{4.845995in}{0.915842in}}%
\pgfpathlineto{\pgfqpoint{4.845995in}{0.918791in}}%
\pgfpathlineto{\pgfqpoint{4.850536in}{0.918791in}}%
\pgfpathlineto{\pgfqpoint{4.850536in}{0.915842in}}%
\pgfpathmoveto{\pgfqpoint{4.868700in}{0.880451in}}%
\pgfpathlineto{\pgfqpoint{4.868700in}{0.880451in}}%
\pgfpathlineto{\pgfqpoint{4.868700in}{0.883400in}}%
\pgfpathlineto{\pgfqpoint{4.873240in}{0.883400in}}%
\pgfpathlineto{\pgfqpoint{4.873240in}{0.880451in}}%
\pgfpathmoveto{\pgfqpoint{4.864159in}{0.886350in}}%
\pgfpathlineto{\pgfqpoint{4.864159in}{0.886350in}}%
\pgfpathlineto{\pgfqpoint{4.864159in}{0.889299in}}%
\pgfpathlineto{\pgfqpoint{4.868700in}{0.889299in}}%
\pgfpathlineto{\pgfqpoint{4.868700in}{0.886350in}}%
\pgfpathmoveto{\pgfqpoint{4.868700in}{0.883400in}}%
\pgfpathlineto{\pgfqpoint{4.868700in}{0.883400in}}%
\pgfpathlineto{\pgfqpoint{4.868700in}{0.886350in}}%
\pgfpathlineto{\pgfqpoint{4.873240in}{0.886350in}}%
\pgfpathlineto{\pgfqpoint{4.873240in}{0.883400in}}%
\pgfpathmoveto{\pgfqpoint{4.868700in}{0.886350in}}%
\pgfpathlineto{\pgfqpoint{4.868700in}{0.886350in}}%
\pgfpathlineto{\pgfqpoint{4.868700in}{0.889299in}}%
\pgfpathlineto{\pgfqpoint{4.873240in}{0.889299in}}%
\pgfpathlineto{\pgfqpoint{4.873240in}{0.886350in}}%
\pgfpathmoveto{\pgfqpoint{4.859618in}{0.892248in}}%
\pgfpathlineto{\pgfqpoint{4.859618in}{0.892248in}}%
\pgfpathlineto{\pgfqpoint{4.859618in}{0.895197in}}%
\pgfpathlineto{\pgfqpoint{4.864159in}{0.895197in}}%
\pgfpathlineto{\pgfqpoint{4.864159in}{0.892248in}}%
\pgfpathmoveto{\pgfqpoint{4.855077in}{0.898146in}}%
\pgfpathlineto{\pgfqpoint{4.855077in}{0.898146in}}%
\pgfpathlineto{\pgfqpoint{4.855077in}{0.901096in}}%
\pgfpathlineto{\pgfqpoint{4.859618in}{0.901096in}}%
\pgfpathlineto{\pgfqpoint{4.859618in}{0.898146in}}%
\pgfpathmoveto{\pgfqpoint{4.859618in}{0.895197in}}%
\pgfpathlineto{\pgfqpoint{4.859618in}{0.895197in}}%
\pgfpathlineto{\pgfqpoint{4.859618in}{0.898146in}}%
\pgfpathlineto{\pgfqpoint{4.864159in}{0.898146in}}%
\pgfpathlineto{\pgfqpoint{4.864159in}{0.895197in}}%
\pgfpathmoveto{\pgfqpoint{4.859618in}{0.898146in}}%
\pgfpathlineto{\pgfqpoint{4.859618in}{0.898146in}}%
\pgfpathlineto{\pgfqpoint{4.859618in}{0.901096in}}%
\pgfpathlineto{\pgfqpoint{4.864159in}{0.901096in}}%
\pgfpathlineto{\pgfqpoint{4.864159in}{0.898146in}}%
\pgfpathmoveto{\pgfqpoint{4.864159in}{0.889299in}}%
\pgfpathlineto{\pgfqpoint{4.864159in}{0.889299in}}%
\pgfpathlineto{\pgfqpoint{4.864159in}{0.892248in}}%
\pgfpathlineto{\pgfqpoint{4.868700in}{0.892248in}}%
\pgfpathlineto{\pgfqpoint{4.868700in}{0.889299in}}%
\pgfpathmoveto{\pgfqpoint{4.864159in}{0.892248in}}%
\pgfpathlineto{\pgfqpoint{4.864159in}{0.892248in}}%
\pgfpathlineto{\pgfqpoint{4.864159in}{0.895197in}}%
\pgfpathlineto{\pgfqpoint{4.868700in}{0.895197in}}%
\pgfpathlineto{\pgfqpoint{4.868700in}{0.892248in}}%
\pgfpathmoveto{\pgfqpoint{4.873240in}{0.877502in}}%
\pgfpathlineto{\pgfqpoint{4.873240in}{0.877502in}}%
\pgfpathlineto{\pgfqpoint{4.873240in}{0.880451in}}%
\pgfpathlineto{\pgfqpoint{4.877781in}{0.880451in}}%
\pgfpathlineto{\pgfqpoint{4.877781in}{0.877502in}}%
\pgfpathmoveto{\pgfqpoint{4.873240in}{0.880451in}}%
\pgfpathlineto{\pgfqpoint{4.873240in}{0.880451in}}%
\pgfpathlineto{\pgfqpoint{4.873240in}{0.883400in}}%
\pgfpathlineto{\pgfqpoint{4.877781in}{0.883400in}}%
\pgfpathlineto{\pgfqpoint{4.877781in}{0.880451in}}%
\pgfpathmoveto{\pgfqpoint{4.855077in}{0.901096in}}%
\pgfpathlineto{\pgfqpoint{4.855077in}{0.901096in}}%
\pgfpathlineto{\pgfqpoint{4.855077in}{0.904045in}}%
\pgfpathlineto{\pgfqpoint{4.859618in}{0.904045in}}%
\pgfpathlineto{\pgfqpoint{4.859618in}{0.901096in}}%
\pgfpathmoveto{\pgfqpoint{4.855077in}{0.904045in}}%
\pgfpathlineto{\pgfqpoint{4.855077in}{0.904045in}}%
\pgfpathlineto{\pgfqpoint{4.855077in}{0.906994in}}%
\pgfpathlineto{\pgfqpoint{4.859618in}{0.906994in}}%
\pgfpathlineto{\pgfqpoint{4.859618in}{0.904045in}}%
\pgfpathmoveto{\pgfqpoint{4.832373in}{0.927638in}}%
\pgfpathlineto{\pgfqpoint{4.832373in}{0.927638in}}%
\pgfpathlineto{\pgfqpoint{4.832373in}{0.930588in}}%
\pgfpathlineto{\pgfqpoint{4.836913in}{0.930588in}}%
\pgfpathlineto{\pgfqpoint{4.836913in}{0.927638in}}%
\pgfpathmoveto{\pgfqpoint{4.827832in}{0.933537in}}%
\pgfpathlineto{\pgfqpoint{4.827832in}{0.933537in}}%
\pgfpathlineto{\pgfqpoint{4.827832in}{0.936486in}}%
\pgfpathlineto{\pgfqpoint{4.832373in}{0.936486in}}%
\pgfpathlineto{\pgfqpoint{4.832373in}{0.933537in}}%
\pgfpathmoveto{\pgfqpoint{4.832373in}{0.930588in}}%
\pgfpathlineto{\pgfqpoint{4.832373in}{0.930588in}}%
\pgfpathlineto{\pgfqpoint{4.832373in}{0.933537in}}%
\pgfpathlineto{\pgfqpoint{4.836913in}{0.933537in}}%
\pgfpathlineto{\pgfqpoint{4.836913in}{0.930588in}}%
\pgfpathmoveto{\pgfqpoint{4.832373in}{0.933537in}}%
\pgfpathlineto{\pgfqpoint{4.832373in}{0.933537in}}%
\pgfpathlineto{\pgfqpoint{4.832373in}{0.936486in}}%
\pgfpathlineto{\pgfqpoint{4.836913in}{0.936486in}}%
\pgfpathlineto{\pgfqpoint{4.836913in}{0.933537in}}%
\pgfpathmoveto{\pgfqpoint{4.823291in}{0.939435in}}%
\pgfpathlineto{\pgfqpoint{4.823291in}{0.939435in}}%
\pgfpathlineto{\pgfqpoint{4.823291in}{0.942384in}}%
\pgfpathlineto{\pgfqpoint{4.827832in}{0.942384in}}%
\pgfpathlineto{\pgfqpoint{4.827832in}{0.939435in}}%
\pgfpathmoveto{\pgfqpoint{4.818750in}{0.945334in}}%
\pgfpathlineto{\pgfqpoint{4.818750in}{0.945334in}}%
\pgfpathlineto{\pgfqpoint{4.818750in}{0.948283in}}%
\pgfpathlineto{\pgfqpoint{4.823291in}{0.948283in}}%
\pgfpathlineto{\pgfqpoint{4.823291in}{0.945334in}}%
\pgfpathmoveto{\pgfqpoint{4.823291in}{0.942384in}}%
\pgfpathlineto{\pgfqpoint{4.823291in}{0.942384in}}%
\pgfpathlineto{\pgfqpoint{4.823291in}{0.945334in}}%
\pgfpathlineto{\pgfqpoint{4.827832in}{0.945334in}}%
\pgfpathlineto{\pgfqpoint{4.827832in}{0.942384in}}%
\pgfpathmoveto{\pgfqpoint{4.823291in}{0.945334in}}%
\pgfpathlineto{\pgfqpoint{4.823291in}{0.945334in}}%
\pgfpathlineto{\pgfqpoint{4.823291in}{0.948283in}}%
\pgfpathlineto{\pgfqpoint{4.827832in}{0.948283in}}%
\pgfpathlineto{\pgfqpoint{4.827832in}{0.945334in}}%
\pgfpathmoveto{\pgfqpoint{4.827832in}{0.936486in}}%
\pgfpathlineto{\pgfqpoint{4.827832in}{0.936486in}}%
\pgfpathlineto{\pgfqpoint{4.827832in}{0.939435in}}%
\pgfpathlineto{\pgfqpoint{4.832373in}{0.939435in}}%
\pgfpathlineto{\pgfqpoint{4.832373in}{0.936486in}}%
\pgfpathmoveto{\pgfqpoint{4.827832in}{0.939435in}}%
\pgfpathlineto{\pgfqpoint{4.827832in}{0.939435in}}%
\pgfpathlineto{\pgfqpoint{4.827832in}{0.942384in}}%
\pgfpathlineto{\pgfqpoint{4.832373in}{0.942384in}}%
\pgfpathlineto{\pgfqpoint{4.832373in}{0.939435in}}%
\pgfpathmoveto{\pgfqpoint{4.836913in}{0.924689in}}%
\pgfpathlineto{\pgfqpoint{4.836913in}{0.924689in}}%
\pgfpathlineto{\pgfqpoint{4.836913in}{0.927638in}}%
\pgfpathlineto{\pgfqpoint{4.841454in}{0.927638in}}%
\pgfpathlineto{\pgfqpoint{4.841454in}{0.924689in}}%
\pgfpathmoveto{\pgfqpoint{4.836913in}{0.927638in}}%
\pgfpathlineto{\pgfqpoint{4.836913in}{0.927638in}}%
\pgfpathlineto{\pgfqpoint{4.836913in}{0.930588in}}%
\pgfpathlineto{\pgfqpoint{4.841454in}{0.930588in}}%
\pgfpathlineto{\pgfqpoint{4.841454in}{0.927638in}}%
\pgfpathmoveto{\pgfqpoint{4.818750in}{0.948283in}}%
\pgfpathlineto{\pgfqpoint{4.818750in}{0.948283in}}%
\pgfpathlineto{\pgfqpoint{4.818750in}{0.951232in}}%
\pgfpathlineto{\pgfqpoint{4.823291in}{0.951232in}}%
\pgfpathlineto{\pgfqpoint{4.823291in}{0.948283in}}%
\pgfpathmoveto{\pgfqpoint{4.818750in}{0.951232in}}%
\pgfpathlineto{\pgfqpoint{4.818750in}{0.951232in}}%
\pgfpathlineto{\pgfqpoint{4.818750in}{0.954181in}}%
\pgfpathlineto{\pgfqpoint{4.823291in}{0.954181in}}%
\pgfpathlineto{\pgfqpoint{4.823291in}{0.951232in}}%
\pgfpathmoveto{\pgfqpoint{5.104834in}{0.573728in}}%
\pgfpathlineto{\pgfqpoint{5.104834in}{0.573728in}}%
\pgfpathlineto{\pgfqpoint{5.104834in}{0.576677in}}%
\pgfpathlineto{\pgfqpoint{5.109375in}{0.576677in}}%
\pgfpathlineto{\pgfqpoint{5.109375in}{0.573728in}}%
\pgfpathmoveto{\pgfqpoint{5.100293in}{0.579626in}}%
\pgfpathlineto{\pgfqpoint{5.100293in}{0.579626in}}%
\pgfpathlineto{\pgfqpoint{5.100293in}{0.582576in}}%
\pgfpathlineto{\pgfqpoint{5.104834in}{0.582576in}}%
\pgfpathlineto{\pgfqpoint{5.104834in}{0.579626in}}%
\pgfpathmoveto{\pgfqpoint{5.104834in}{0.576677in}}%
\pgfpathlineto{\pgfqpoint{5.104834in}{0.576677in}}%
\pgfpathlineto{\pgfqpoint{5.104834in}{0.579626in}}%
\pgfpathlineto{\pgfqpoint{5.109375in}{0.579626in}}%
\pgfpathlineto{\pgfqpoint{5.109375in}{0.576677in}}%
\pgfpathmoveto{\pgfqpoint{5.104834in}{0.579626in}}%
\pgfpathlineto{\pgfqpoint{5.104834in}{0.579626in}}%
\pgfpathlineto{\pgfqpoint{5.104834in}{0.582576in}}%
\pgfpathlineto{\pgfqpoint{5.109375in}{0.582576in}}%
\pgfpathlineto{\pgfqpoint{5.109375in}{0.579626in}}%
\pgfpathmoveto{\pgfqpoint{5.095752in}{0.585525in}}%
\pgfpathlineto{\pgfqpoint{5.095752in}{0.585525in}}%
\pgfpathlineto{\pgfqpoint{5.095752in}{0.588474in}}%
\pgfpathlineto{\pgfqpoint{5.100293in}{0.588474in}}%
\pgfpathlineto{\pgfqpoint{5.100293in}{0.585525in}}%
\pgfpathmoveto{\pgfqpoint{5.091211in}{0.591423in}}%
\pgfpathlineto{\pgfqpoint{5.091211in}{0.591423in}}%
\pgfpathlineto{\pgfqpoint{5.091211in}{0.594372in}}%
\pgfpathlineto{\pgfqpoint{5.095752in}{0.594372in}}%
\pgfpathlineto{\pgfqpoint{5.095752in}{0.591423in}}%
\pgfpathmoveto{\pgfqpoint{5.095752in}{0.588474in}}%
\pgfpathlineto{\pgfqpoint{5.095752in}{0.588474in}}%
\pgfpathlineto{\pgfqpoint{5.095752in}{0.591423in}}%
\pgfpathlineto{\pgfqpoint{5.100293in}{0.591423in}}%
\pgfpathlineto{\pgfqpoint{5.100293in}{0.588474in}}%
\pgfpathmoveto{\pgfqpoint{5.095752in}{0.591423in}}%
\pgfpathlineto{\pgfqpoint{5.095752in}{0.591423in}}%
\pgfpathlineto{\pgfqpoint{5.095752in}{0.594372in}}%
\pgfpathlineto{\pgfqpoint{5.100293in}{0.594372in}}%
\pgfpathlineto{\pgfqpoint{5.100293in}{0.591423in}}%
\pgfpathmoveto{\pgfqpoint{5.100293in}{0.582576in}}%
\pgfpathlineto{\pgfqpoint{5.100293in}{0.582576in}}%
\pgfpathlineto{\pgfqpoint{5.100293in}{0.585525in}}%
\pgfpathlineto{\pgfqpoint{5.104834in}{0.585525in}}%
\pgfpathlineto{\pgfqpoint{5.104834in}{0.582576in}}%
\pgfpathmoveto{\pgfqpoint{5.100293in}{0.585525in}}%
\pgfpathlineto{\pgfqpoint{5.100293in}{0.585525in}}%
\pgfpathlineto{\pgfqpoint{5.100293in}{0.588474in}}%
\pgfpathlineto{\pgfqpoint{5.104834in}{0.588474in}}%
\pgfpathlineto{\pgfqpoint{5.104834in}{0.585525in}}%
\pgfpathmoveto{\pgfqpoint{5.032176in}{0.668104in}}%
\pgfpathlineto{\pgfqpoint{5.032176in}{0.668104in}}%
\pgfpathlineto{\pgfqpoint{5.032176in}{0.671053in}}%
\pgfpathlineto{\pgfqpoint{5.036717in}{0.671053in}}%
\pgfpathlineto{\pgfqpoint{5.036717in}{0.668104in}}%
\pgfpathmoveto{\pgfqpoint{5.027634in}{0.674003in}}%
\pgfpathlineto{\pgfqpoint{5.027634in}{0.674003in}}%
\pgfpathlineto{\pgfqpoint{5.027634in}{0.676952in}}%
\pgfpathlineto{\pgfqpoint{5.032176in}{0.676952in}}%
\pgfpathlineto{\pgfqpoint{5.032176in}{0.674003in}}%
\pgfpathmoveto{\pgfqpoint{5.032176in}{0.671053in}}%
\pgfpathlineto{\pgfqpoint{5.032176in}{0.671053in}}%
\pgfpathlineto{\pgfqpoint{5.032176in}{0.674003in}}%
\pgfpathlineto{\pgfqpoint{5.036717in}{0.674003in}}%
\pgfpathlineto{\pgfqpoint{5.036717in}{0.671053in}}%
\pgfpathmoveto{\pgfqpoint{5.032176in}{0.674003in}}%
\pgfpathlineto{\pgfqpoint{5.032176in}{0.674003in}}%
\pgfpathlineto{\pgfqpoint{5.032176in}{0.676952in}}%
\pgfpathlineto{\pgfqpoint{5.036717in}{0.676952in}}%
\pgfpathlineto{\pgfqpoint{5.036717in}{0.674003in}}%
\pgfpathmoveto{\pgfqpoint{5.023093in}{0.679901in}}%
\pgfpathlineto{\pgfqpoint{5.023093in}{0.679901in}}%
\pgfpathlineto{\pgfqpoint{5.023093in}{0.682850in}}%
\pgfpathlineto{\pgfqpoint{5.027634in}{0.682850in}}%
\pgfpathlineto{\pgfqpoint{5.027634in}{0.679901in}}%
\pgfpathmoveto{\pgfqpoint{5.018552in}{0.685800in}}%
\pgfpathlineto{\pgfqpoint{5.018552in}{0.685800in}}%
\pgfpathlineto{\pgfqpoint{5.018552in}{0.688749in}}%
\pgfpathlineto{\pgfqpoint{5.023093in}{0.688749in}}%
\pgfpathlineto{\pgfqpoint{5.023093in}{0.685800in}}%
\pgfpathmoveto{\pgfqpoint{5.023093in}{0.682850in}}%
\pgfpathlineto{\pgfqpoint{5.023093in}{0.682850in}}%
\pgfpathlineto{\pgfqpoint{5.023093in}{0.685800in}}%
\pgfpathlineto{\pgfqpoint{5.027634in}{0.685800in}}%
\pgfpathlineto{\pgfqpoint{5.027634in}{0.682850in}}%
\pgfpathmoveto{\pgfqpoint{5.023093in}{0.685800in}}%
\pgfpathlineto{\pgfqpoint{5.023093in}{0.685800in}}%
\pgfpathlineto{\pgfqpoint{5.023093in}{0.688749in}}%
\pgfpathlineto{\pgfqpoint{5.027634in}{0.688749in}}%
\pgfpathlineto{\pgfqpoint{5.027634in}{0.685800in}}%
\pgfpathmoveto{\pgfqpoint{5.027634in}{0.676952in}}%
\pgfpathlineto{\pgfqpoint{5.027634in}{0.676952in}}%
\pgfpathlineto{\pgfqpoint{5.027634in}{0.679901in}}%
\pgfpathlineto{\pgfqpoint{5.032176in}{0.679901in}}%
\pgfpathlineto{\pgfqpoint{5.032176in}{0.676952in}}%
\pgfpathmoveto{\pgfqpoint{5.027634in}{0.679901in}}%
\pgfpathlineto{\pgfqpoint{5.027634in}{0.679901in}}%
\pgfpathlineto{\pgfqpoint{5.027634in}{0.682850in}}%
\pgfpathlineto{\pgfqpoint{5.032176in}{0.682850in}}%
\pgfpathlineto{\pgfqpoint{5.032176in}{0.679901in}}%
\pgfpathmoveto{\pgfqpoint{5.068505in}{0.620916in}}%
\pgfpathlineto{\pgfqpoint{5.068505in}{0.620916in}}%
\pgfpathlineto{\pgfqpoint{5.068505in}{0.623865in}}%
\pgfpathlineto{\pgfqpoint{5.073046in}{0.623865in}}%
\pgfpathlineto{\pgfqpoint{5.073046in}{0.620916in}}%
\pgfpathmoveto{\pgfqpoint{5.063964in}{0.626814in}}%
\pgfpathlineto{\pgfqpoint{5.063964in}{0.626814in}}%
\pgfpathlineto{\pgfqpoint{5.063964in}{0.629764in}}%
\pgfpathlineto{\pgfqpoint{5.068505in}{0.629764in}}%
\pgfpathlineto{\pgfqpoint{5.068505in}{0.626814in}}%
\pgfpathmoveto{\pgfqpoint{5.068505in}{0.623865in}}%
\pgfpathlineto{\pgfqpoint{5.068505in}{0.623865in}}%
\pgfpathlineto{\pgfqpoint{5.068505in}{0.626814in}}%
\pgfpathlineto{\pgfqpoint{5.073046in}{0.626814in}}%
\pgfpathlineto{\pgfqpoint{5.073046in}{0.623865in}}%
\pgfpathmoveto{\pgfqpoint{5.068505in}{0.626814in}}%
\pgfpathlineto{\pgfqpoint{5.068505in}{0.626814in}}%
\pgfpathlineto{\pgfqpoint{5.068505in}{0.629764in}}%
\pgfpathlineto{\pgfqpoint{5.073046in}{0.629764in}}%
\pgfpathlineto{\pgfqpoint{5.073046in}{0.626814in}}%
\pgfpathmoveto{\pgfqpoint{5.059423in}{0.632713in}}%
\pgfpathlineto{\pgfqpoint{5.059423in}{0.632713in}}%
\pgfpathlineto{\pgfqpoint{5.059423in}{0.635662in}}%
\pgfpathlineto{\pgfqpoint{5.063964in}{0.635662in}}%
\pgfpathlineto{\pgfqpoint{5.063964in}{0.632713in}}%
\pgfpathmoveto{\pgfqpoint{5.054881in}{0.638611in}}%
\pgfpathlineto{\pgfqpoint{5.054881in}{0.638611in}}%
\pgfpathlineto{\pgfqpoint{5.054881in}{0.641561in}}%
\pgfpathlineto{\pgfqpoint{5.059423in}{0.641561in}}%
\pgfpathlineto{\pgfqpoint{5.059423in}{0.638611in}}%
\pgfpathmoveto{\pgfqpoint{5.059423in}{0.635662in}}%
\pgfpathlineto{\pgfqpoint{5.059423in}{0.635662in}}%
\pgfpathlineto{\pgfqpoint{5.059423in}{0.638611in}}%
\pgfpathlineto{\pgfqpoint{5.063964in}{0.638611in}}%
\pgfpathlineto{\pgfqpoint{5.063964in}{0.635662in}}%
\pgfpathmoveto{\pgfqpoint{5.059423in}{0.638611in}}%
\pgfpathlineto{\pgfqpoint{5.059423in}{0.638611in}}%
\pgfpathlineto{\pgfqpoint{5.059423in}{0.641561in}}%
\pgfpathlineto{\pgfqpoint{5.063964in}{0.641561in}}%
\pgfpathlineto{\pgfqpoint{5.063964in}{0.638611in}}%
\pgfpathmoveto{\pgfqpoint{5.063964in}{0.629764in}}%
\pgfpathlineto{\pgfqpoint{5.063964in}{0.629764in}}%
\pgfpathlineto{\pgfqpoint{5.063964in}{0.632713in}}%
\pgfpathlineto{\pgfqpoint{5.068505in}{0.632713in}}%
\pgfpathlineto{\pgfqpoint{5.068505in}{0.629764in}}%
\pgfpathmoveto{\pgfqpoint{5.063964in}{0.632713in}}%
\pgfpathlineto{\pgfqpoint{5.063964in}{0.632713in}}%
\pgfpathlineto{\pgfqpoint{5.063964in}{0.635662in}}%
\pgfpathlineto{\pgfqpoint{5.068505in}{0.635662in}}%
\pgfpathlineto{\pgfqpoint{5.068505in}{0.632713in}}%
\pgfpathmoveto{\pgfqpoint{5.086670in}{0.597322in}}%
\pgfpathlineto{\pgfqpoint{5.086670in}{0.597322in}}%
\pgfpathlineto{\pgfqpoint{5.086670in}{0.600271in}}%
\pgfpathlineto{\pgfqpoint{5.091211in}{0.600271in}}%
\pgfpathlineto{\pgfqpoint{5.091211in}{0.597322in}}%
\pgfpathmoveto{\pgfqpoint{5.082128in}{0.603220in}}%
\pgfpathlineto{\pgfqpoint{5.082128in}{0.603220in}}%
\pgfpathlineto{\pgfqpoint{5.082128in}{0.606169in}}%
\pgfpathlineto{\pgfqpoint{5.086670in}{0.606169in}}%
\pgfpathlineto{\pgfqpoint{5.086670in}{0.603220in}}%
\pgfpathmoveto{\pgfqpoint{5.086670in}{0.600271in}}%
\pgfpathlineto{\pgfqpoint{5.086670in}{0.600271in}}%
\pgfpathlineto{\pgfqpoint{5.086670in}{0.603220in}}%
\pgfpathlineto{\pgfqpoint{5.091211in}{0.603220in}}%
\pgfpathlineto{\pgfqpoint{5.091211in}{0.600271in}}%
\pgfpathmoveto{\pgfqpoint{5.086670in}{0.603220in}}%
\pgfpathlineto{\pgfqpoint{5.086670in}{0.603220in}}%
\pgfpathlineto{\pgfqpoint{5.086670in}{0.606169in}}%
\pgfpathlineto{\pgfqpoint{5.091211in}{0.606169in}}%
\pgfpathlineto{\pgfqpoint{5.091211in}{0.603220in}}%
\pgfpathmoveto{\pgfqpoint{5.077587in}{0.609119in}}%
\pgfpathlineto{\pgfqpoint{5.077587in}{0.609119in}}%
\pgfpathlineto{\pgfqpoint{5.077587in}{0.612068in}}%
\pgfpathlineto{\pgfqpoint{5.082128in}{0.612068in}}%
\pgfpathlineto{\pgfqpoint{5.082128in}{0.609119in}}%
\pgfpathmoveto{\pgfqpoint{5.073046in}{0.615017in}}%
\pgfpathlineto{\pgfqpoint{5.073046in}{0.615017in}}%
\pgfpathlineto{\pgfqpoint{5.073046in}{0.617966in}}%
\pgfpathlineto{\pgfqpoint{5.077587in}{0.617966in}}%
\pgfpathlineto{\pgfqpoint{5.077587in}{0.615017in}}%
\pgfpathmoveto{\pgfqpoint{5.077587in}{0.612068in}}%
\pgfpathlineto{\pgfqpoint{5.077587in}{0.612068in}}%
\pgfpathlineto{\pgfqpoint{5.077587in}{0.615017in}}%
\pgfpathlineto{\pgfqpoint{5.082128in}{0.615017in}}%
\pgfpathlineto{\pgfqpoint{5.082128in}{0.612068in}}%
\pgfpathmoveto{\pgfqpoint{5.077587in}{0.615017in}}%
\pgfpathlineto{\pgfqpoint{5.077587in}{0.615017in}}%
\pgfpathlineto{\pgfqpoint{5.077587in}{0.617966in}}%
\pgfpathlineto{\pgfqpoint{5.082128in}{0.617966in}}%
\pgfpathlineto{\pgfqpoint{5.082128in}{0.615017in}}%
\pgfpathmoveto{\pgfqpoint{5.082128in}{0.606169in}}%
\pgfpathlineto{\pgfqpoint{5.082128in}{0.606169in}}%
\pgfpathlineto{\pgfqpoint{5.082128in}{0.609119in}}%
\pgfpathlineto{\pgfqpoint{5.086670in}{0.609119in}}%
\pgfpathlineto{\pgfqpoint{5.086670in}{0.606169in}}%
\pgfpathmoveto{\pgfqpoint{5.082128in}{0.609119in}}%
\pgfpathlineto{\pgfqpoint{5.082128in}{0.609119in}}%
\pgfpathlineto{\pgfqpoint{5.082128in}{0.612068in}}%
\pgfpathlineto{\pgfqpoint{5.086670in}{0.612068in}}%
\pgfpathlineto{\pgfqpoint{5.086670in}{0.609119in}}%
\pgfpathmoveto{\pgfqpoint{5.091211in}{0.594372in}}%
\pgfpathlineto{\pgfqpoint{5.091211in}{0.594372in}}%
\pgfpathlineto{\pgfqpoint{5.091211in}{0.597322in}}%
\pgfpathlineto{\pgfqpoint{5.095752in}{0.597322in}}%
\pgfpathlineto{\pgfqpoint{5.095752in}{0.594372in}}%
\pgfpathmoveto{\pgfqpoint{5.091211in}{0.597322in}}%
\pgfpathlineto{\pgfqpoint{5.091211in}{0.597322in}}%
\pgfpathlineto{\pgfqpoint{5.091211in}{0.600271in}}%
\pgfpathlineto{\pgfqpoint{5.095752in}{0.600271in}}%
\pgfpathlineto{\pgfqpoint{5.095752in}{0.597322in}}%
\pgfpathmoveto{\pgfqpoint{5.073046in}{0.617966in}}%
\pgfpathlineto{\pgfqpoint{5.073046in}{0.617966in}}%
\pgfpathlineto{\pgfqpoint{5.073046in}{0.620916in}}%
\pgfpathlineto{\pgfqpoint{5.077587in}{0.620916in}}%
\pgfpathlineto{\pgfqpoint{5.077587in}{0.617966in}}%
\pgfpathmoveto{\pgfqpoint{5.073046in}{0.620916in}}%
\pgfpathlineto{\pgfqpoint{5.073046in}{0.620916in}}%
\pgfpathlineto{\pgfqpoint{5.073046in}{0.623865in}}%
\pgfpathlineto{\pgfqpoint{5.077587in}{0.623865in}}%
\pgfpathlineto{\pgfqpoint{5.077587in}{0.620916in}}%
\pgfpathmoveto{\pgfqpoint{5.050340in}{0.644510in}}%
\pgfpathlineto{\pgfqpoint{5.050340in}{0.644510in}}%
\pgfpathlineto{\pgfqpoint{5.050340in}{0.647459in}}%
\pgfpathlineto{\pgfqpoint{5.054881in}{0.647459in}}%
\pgfpathlineto{\pgfqpoint{5.054881in}{0.644510in}}%
\pgfpathmoveto{\pgfqpoint{5.045799in}{0.650408in}}%
\pgfpathlineto{\pgfqpoint{5.045799in}{0.650408in}}%
\pgfpathlineto{\pgfqpoint{5.045799in}{0.653358in}}%
\pgfpathlineto{\pgfqpoint{5.050340in}{0.653358in}}%
\pgfpathlineto{\pgfqpoint{5.050340in}{0.650408in}}%
\pgfpathmoveto{\pgfqpoint{5.050340in}{0.647459in}}%
\pgfpathlineto{\pgfqpoint{5.050340in}{0.647459in}}%
\pgfpathlineto{\pgfqpoint{5.050340in}{0.650408in}}%
\pgfpathlineto{\pgfqpoint{5.054881in}{0.650408in}}%
\pgfpathlineto{\pgfqpoint{5.054881in}{0.647459in}}%
\pgfpathmoveto{\pgfqpoint{5.050340in}{0.650408in}}%
\pgfpathlineto{\pgfqpoint{5.050340in}{0.650408in}}%
\pgfpathlineto{\pgfqpoint{5.050340in}{0.653358in}}%
\pgfpathlineto{\pgfqpoint{5.054881in}{0.653358in}}%
\pgfpathlineto{\pgfqpoint{5.054881in}{0.650408in}}%
\pgfpathmoveto{\pgfqpoint{5.041258in}{0.656307in}}%
\pgfpathlineto{\pgfqpoint{5.041258in}{0.656307in}}%
\pgfpathlineto{\pgfqpoint{5.041258in}{0.659256in}}%
\pgfpathlineto{\pgfqpoint{5.045799in}{0.659256in}}%
\pgfpathlineto{\pgfqpoint{5.045799in}{0.656307in}}%
\pgfpathmoveto{\pgfqpoint{5.036717in}{0.662205in}}%
\pgfpathlineto{\pgfqpoint{5.036717in}{0.662205in}}%
\pgfpathlineto{\pgfqpoint{5.036717in}{0.665155in}}%
\pgfpathlineto{\pgfqpoint{5.041258in}{0.665155in}}%
\pgfpathlineto{\pgfqpoint{5.041258in}{0.662205in}}%
\pgfpathmoveto{\pgfqpoint{5.041258in}{0.659256in}}%
\pgfpathlineto{\pgfqpoint{5.041258in}{0.659256in}}%
\pgfpathlineto{\pgfqpoint{5.041258in}{0.662205in}}%
\pgfpathlineto{\pgfqpoint{5.045799in}{0.662205in}}%
\pgfpathlineto{\pgfqpoint{5.045799in}{0.659256in}}%
\pgfpathmoveto{\pgfqpoint{5.041258in}{0.662205in}}%
\pgfpathlineto{\pgfqpoint{5.041258in}{0.662205in}}%
\pgfpathlineto{\pgfqpoint{5.041258in}{0.665155in}}%
\pgfpathlineto{\pgfqpoint{5.045799in}{0.665155in}}%
\pgfpathlineto{\pgfqpoint{5.045799in}{0.662205in}}%
\pgfpathmoveto{\pgfqpoint{5.045799in}{0.653358in}}%
\pgfpathlineto{\pgfqpoint{5.045799in}{0.653358in}}%
\pgfpathlineto{\pgfqpoint{5.045799in}{0.656307in}}%
\pgfpathlineto{\pgfqpoint{5.050340in}{0.656307in}}%
\pgfpathlineto{\pgfqpoint{5.050340in}{0.653358in}}%
\pgfpathmoveto{\pgfqpoint{5.045799in}{0.656307in}}%
\pgfpathlineto{\pgfqpoint{5.045799in}{0.656307in}}%
\pgfpathlineto{\pgfqpoint{5.045799in}{0.659256in}}%
\pgfpathlineto{\pgfqpoint{5.050340in}{0.659256in}}%
\pgfpathlineto{\pgfqpoint{5.050340in}{0.656307in}}%
\pgfpathmoveto{\pgfqpoint{5.054881in}{0.641561in}}%
\pgfpathlineto{\pgfqpoint{5.054881in}{0.641561in}}%
\pgfpathlineto{\pgfqpoint{5.054881in}{0.644510in}}%
\pgfpathlineto{\pgfqpoint{5.059423in}{0.644510in}}%
\pgfpathlineto{\pgfqpoint{5.059423in}{0.641561in}}%
\pgfpathmoveto{\pgfqpoint{5.054881in}{0.644510in}}%
\pgfpathlineto{\pgfqpoint{5.054881in}{0.644510in}}%
\pgfpathlineto{\pgfqpoint{5.054881in}{0.647459in}}%
\pgfpathlineto{\pgfqpoint{5.059423in}{0.647459in}}%
\pgfpathlineto{\pgfqpoint{5.059423in}{0.644510in}}%
\pgfpathmoveto{\pgfqpoint{5.036717in}{0.665155in}}%
\pgfpathlineto{\pgfqpoint{5.036717in}{0.665155in}}%
\pgfpathlineto{\pgfqpoint{5.036717in}{0.668104in}}%
\pgfpathlineto{\pgfqpoint{5.041258in}{0.668104in}}%
\pgfpathlineto{\pgfqpoint{5.041258in}{0.665155in}}%
\pgfpathmoveto{\pgfqpoint{5.036717in}{0.668104in}}%
\pgfpathlineto{\pgfqpoint{5.036717in}{0.668104in}}%
\pgfpathlineto{\pgfqpoint{5.036717in}{0.671053in}}%
\pgfpathlineto{\pgfqpoint{5.041258in}{0.671053in}}%
\pgfpathlineto{\pgfqpoint{5.041258in}{0.668104in}}%
\pgfpathmoveto{\pgfqpoint{4.995846in}{0.715292in}}%
\pgfpathlineto{\pgfqpoint{4.995846in}{0.715292in}}%
\pgfpathlineto{\pgfqpoint{4.995846in}{0.718242in}}%
\pgfpathlineto{\pgfqpoint{5.000387in}{0.718242in}}%
\pgfpathlineto{\pgfqpoint{5.000387in}{0.715292in}}%
\pgfpathmoveto{\pgfqpoint{4.991305in}{0.721191in}}%
\pgfpathlineto{\pgfqpoint{4.991305in}{0.721191in}}%
\pgfpathlineto{\pgfqpoint{4.991305in}{0.724140in}}%
\pgfpathlineto{\pgfqpoint{4.995846in}{0.724140in}}%
\pgfpathlineto{\pgfqpoint{4.995846in}{0.721191in}}%
\pgfpathmoveto{\pgfqpoint{4.995846in}{0.718242in}}%
\pgfpathlineto{\pgfqpoint{4.995846in}{0.718242in}}%
\pgfpathlineto{\pgfqpoint{4.995846in}{0.721191in}}%
\pgfpathlineto{\pgfqpoint{5.000387in}{0.721191in}}%
\pgfpathlineto{\pgfqpoint{5.000387in}{0.718242in}}%
\pgfpathmoveto{\pgfqpoint{4.995846in}{0.721191in}}%
\pgfpathlineto{\pgfqpoint{4.995846in}{0.721191in}}%
\pgfpathlineto{\pgfqpoint{4.995846in}{0.724140in}}%
\pgfpathlineto{\pgfqpoint{5.000387in}{0.724140in}}%
\pgfpathlineto{\pgfqpoint{5.000387in}{0.721191in}}%
\pgfpathmoveto{\pgfqpoint{4.986764in}{0.727090in}}%
\pgfpathlineto{\pgfqpoint{4.986764in}{0.727090in}}%
\pgfpathlineto{\pgfqpoint{4.986764in}{0.730039in}}%
\pgfpathlineto{\pgfqpoint{4.991305in}{0.730039in}}%
\pgfpathlineto{\pgfqpoint{4.991305in}{0.727090in}}%
\pgfpathmoveto{\pgfqpoint{4.982223in}{0.732988in}}%
\pgfpathlineto{\pgfqpoint{4.982223in}{0.732988in}}%
\pgfpathlineto{\pgfqpoint{4.982223in}{0.735937in}}%
\pgfpathlineto{\pgfqpoint{4.986764in}{0.735937in}}%
\pgfpathlineto{\pgfqpoint{4.986764in}{0.732988in}}%
\pgfpathmoveto{\pgfqpoint{4.986764in}{0.730039in}}%
\pgfpathlineto{\pgfqpoint{4.986764in}{0.730039in}}%
\pgfpathlineto{\pgfqpoint{4.986764in}{0.732988in}}%
\pgfpathlineto{\pgfqpoint{4.991305in}{0.732988in}}%
\pgfpathlineto{\pgfqpoint{4.991305in}{0.730039in}}%
\pgfpathmoveto{\pgfqpoint{4.986764in}{0.732988in}}%
\pgfpathlineto{\pgfqpoint{4.986764in}{0.732988in}}%
\pgfpathlineto{\pgfqpoint{4.986764in}{0.735937in}}%
\pgfpathlineto{\pgfqpoint{4.991305in}{0.735937in}}%
\pgfpathlineto{\pgfqpoint{4.991305in}{0.732988in}}%
\pgfpathmoveto{\pgfqpoint{4.991305in}{0.724140in}}%
\pgfpathlineto{\pgfqpoint{4.991305in}{0.724140in}}%
\pgfpathlineto{\pgfqpoint{4.991305in}{0.727090in}}%
\pgfpathlineto{\pgfqpoint{4.995846in}{0.727090in}}%
\pgfpathlineto{\pgfqpoint{4.995846in}{0.724140in}}%
\pgfpathmoveto{\pgfqpoint{4.991305in}{0.727090in}}%
\pgfpathlineto{\pgfqpoint{4.991305in}{0.727090in}}%
\pgfpathlineto{\pgfqpoint{4.991305in}{0.730039in}}%
\pgfpathlineto{\pgfqpoint{4.995846in}{0.730039in}}%
\pgfpathlineto{\pgfqpoint{4.995846in}{0.727090in}}%
\pgfpathmoveto{\pgfqpoint{5.014011in}{0.691698in}}%
\pgfpathlineto{\pgfqpoint{5.014011in}{0.691698in}}%
\pgfpathlineto{\pgfqpoint{5.014011in}{0.694647in}}%
\pgfpathlineto{\pgfqpoint{5.018552in}{0.694647in}}%
\pgfpathlineto{\pgfqpoint{5.018552in}{0.691698in}}%
\pgfpathmoveto{\pgfqpoint{5.009470in}{0.697597in}}%
\pgfpathlineto{\pgfqpoint{5.009470in}{0.697597in}}%
\pgfpathlineto{\pgfqpoint{5.009470in}{0.700546in}}%
\pgfpathlineto{\pgfqpoint{5.014011in}{0.700546in}}%
\pgfpathlineto{\pgfqpoint{5.014011in}{0.697597in}}%
\pgfpathmoveto{\pgfqpoint{5.014011in}{0.694647in}}%
\pgfpathlineto{\pgfqpoint{5.014011in}{0.694647in}}%
\pgfpathlineto{\pgfqpoint{5.014011in}{0.697597in}}%
\pgfpathlineto{\pgfqpoint{5.018552in}{0.697597in}}%
\pgfpathlineto{\pgfqpoint{5.018552in}{0.694647in}}%
\pgfpathmoveto{\pgfqpoint{5.014011in}{0.697597in}}%
\pgfpathlineto{\pgfqpoint{5.014011in}{0.697597in}}%
\pgfpathlineto{\pgfqpoint{5.014011in}{0.700546in}}%
\pgfpathlineto{\pgfqpoint{5.018552in}{0.700546in}}%
\pgfpathlineto{\pgfqpoint{5.018552in}{0.697597in}}%
\pgfpathmoveto{\pgfqpoint{5.004929in}{0.703495in}}%
\pgfpathlineto{\pgfqpoint{5.004929in}{0.703495in}}%
\pgfpathlineto{\pgfqpoint{5.004929in}{0.706445in}}%
\pgfpathlineto{\pgfqpoint{5.009470in}{0.706445in}}%
\pgfpathlineto{\pgfqpoint{5.009470in}{0.703495in}}%
\pgfpathmoveto{\pgfqpoint{5.000387in}{0.709394in}}%
\pgfpathlineto{\pgfqpoint{5.000387in}{0.709394in}}%
\pgfpathlineto{\pgfqpoint{5.000387in}{0.712343in}}%
\pgfpathlineto{\pgfqpoint{5.004929in}{0.712343in}}%
\pgfpathlineto{\pgfqpoint{5.004929in}{0.709394in}}%
\pgfpathmoveto{\pgfqpoint{5.004929in}{0.706445in}}%
\pgfpathlineto{\pgfqpoint{5.004929in}{0.706445in}}%
\pgfpathlineto{\pgfqpoint{5.004929in}{0.709394in}}%
\pgfpathlineto{\pgfqpoint{5.009470in}{0.709394in}}%
\pgfpathlineto{\pgfqpoint{5.009470in}{0.706445in}}%
\pgfpathmoveto{\pgfqpoint{5.004929in}{0.709394in}}%
\pgfpathlineto{\pgfqpoint{5.004929in}{0.709394in}}%
\pgfpathlineto{\pgfqpoint{5.004929in}{0.712343in}}%
\pgfpathlineto{\pgfqpoint{5.009470in}{0.712343in}}%
\pgfpathlineto{\pgfqpoint{5.009470in}{0.709394in}}%
\pgfpathmoveto{\pgfqpoint{5.009470in}{0.700546in}}%
\pgfpathlineto{\pgfqpoint{5.009470in}{0.700546in}}%
\pgfpathlineto{\pgfqpoint{5.009470in}{0.703495in}}%
\pgfpathlineto{\pgfqpoint{5.014011in}{0.703495in}}%
\pgfpathlineto{\pgfqpoint{5.014011in}{0.700546in}}%
\pgfpathmoveto{\pgfqpoint{5.009470in}{0.703495in}}%
\pgfpathlineto{\pgfqpoint{5.009470in}{0.703495in}}%
\pgfpathlineto{\pgfqpoint{5.009470in}{0.706445in}}%
\pgfpathlineto{\pgfqpoint{5.014011in}{0.706445in}}%
\pgfpathlineto{\pgfqpoint{5.014011in}{0.703495in}}%
\pgfpathmoveto{\pgfqpoint{5.018552in}{0.688749in}}%
\pgfpathlineto{\pgfqpoint{5.018552in}{0.688749in}}%
\pgfpathlineto{\pgfqpoint{5.018552in}{0.691698in}}%
\pgfpathlineto{\pgfqpoint{5.023093in}{0.691698in}}%
\pgfpathlineto{\pgfqpoint{5.023093in}{0.688749in}}%
\pgfpathmoveto{\pgfqpoint{5.018552in}{0.691698in}}%
\pgfpathlineto{\pgfqpoint{5.018552in}{0.691698in}}%
\pgfpathlineto{\pgfqpoint{5.018552in}{0.694647in}}%
\pgfpathlineto{\pgfqpoint{5.023093in}{0.694647in}}%
\pgfpathlineto{\pgfqpoint{5.023093in}{0.691698in}}%
\pgfpathmoveto{\pgfqpoint{5.000387in}{0.712343in}}%
\pgfpathlineto{\pgfqpoint{5.000387in}{0.712343in}}%
\pgfpathlineto{\pgfqpoint{5.000387in}{0.715292in}}%
\pgfpathlineto{\pgfqpoint{5.004929in}{0.715292in}}%
\pgfpathlineto{\pgfqpoint{5.004929in}{0.712343in}}%
\pgfpathmoveto{\pgfqpoint{5.000387in}{0.715292in}}%
\pgfpathlineto{\pgfqpoint{5.000387in}{0.715292in}}%
\pgfpathlineto{\pgfqpoint{5.000387in}{0.718242in}}%
\pgfpathlineto{\pgfqpoint{5.004929in}{0.718242in}}%
\pgfpathlineto{\pgfqpoint{5.004929in}{0.715292in}}%
\pgfpathmoveto{\pgfqpoint{4.977682in}{0.738887in}}%
\pgfpathlineto{\pgfqpoint{4.977682in}{0.738887in}}%
\pgfpathlineto{\pgfqpoint{4.977682in}{0.741836in}}%
\pgfpathlineto{\pgfqpoint{4.982223in}{0.741836in}}%
\pgfpathlineto{\pgfqpoint{4.982223in}{0.738887in}}%
\pgfpathmoveto{\pgfqpoint{4.973141in}{0.744785in}}%
\pgfpathlineto{\pgfqpoint{4.973141in}{0.744785in}}%
\pgfpathlineto{\pgfqpoint{4.973141in}{0.747735in}}%
\pgfpathlineto{\pgfqpoint{4.977682in}{0.747735in}}%
\pgfpathlineto{\pgfqpoint{4.977682in}{0.744785in}}%
\pgfpathmoveto{\pgfqpoint{4.977682in}{0.741836in}}%
\pgfpathlineto{\pgfqpoint{4.977682in}{0.741836in}}%
\pgfpathlineto{\pgfqpoint{4.977682in}{0.744785in}}%
\pgfpathlineto{\pgfqpoint{4.982223in}{0.744785in}}%
\pgfpathlineto{\pgfqpoint{4.982223in}{0.741836in}}%
\pgfpathmoveto{\pgfqpoint{4.977682in}{0.744785in}}%
\pgfpathlineto{\pgfqpoint{4.977682in}{0.744785in}}%
\pgfpathlineto{\pgfqpoint{4.977682in}{0.747735in}}%
\pgfpathlineto{\pgfqpoint{4.982223in}{0.747735in}}%
\pgfpathlineto{\pgfqpoint{4.982223in}{0.744785in}}%
\pgfpathmoveto{\pgfqpoint{4.968599in}{0.750684in}}%
\pgfpathlineto{\pgfqpoint{4.968599in}{0.750684in}}%
\pgfpathlineto{\pgfqpoint{4.968599in}{0.753633in}}%
\pgfpathlineto{\pgfqpoint{4.973141in}{0.753633in}}%
\pgfpathlineto{\pgfqpoint{4.973141in}{0.750684in}}%
\pgfpathmoveto{\pgfqpoint{4.964058in}{0.756582in}}%
\pgfpathlineto{\pgfqpoint{4.964058in}{0.756582in}}%
\pgfpathlineto{\pgfqpoint{4.964058in}{0.759532in}}%
\pgfpathlineto{\pgfqpoint{4.968599in}{0.759532in}}%
\pgfpathlineto{\pgfqpoint{4.968599in}{0.756582in}}%
\pgfpathmoveto{\pgfqpoint{4.968599in}{0.753633in}}%
\pgfpathlineto{\pgfqpoint{4.968599in}{0.753633in}}%
\pgfpathlineto{\pgfqpoint{4.968599in}{0.756582in}}%
\pgfpathlineto{\pgfqpoint{4.973141in}{0.756582in}}%
\pgfpathlineto{\pgfqpoint{4.973141in}{0.753633in}}%
\pgfpathmoveto{\pgfqpoint{4.968599in}{0.756582in}}%
\pgfpathlineto{\pgfqpoint{4.968599in}{0.756582in}}%
\pgfpathlineto{\pgfqpoint{4.968599in}{0.759532in}}%
\pgfpathlineto{\pgfqpoint{4.973141in}{0.759532in}}%
\pgfpathlineto{\pgfqpoint{4.973141in}{0.756582in}}%
\pgfpathmoveto{\pgfqpoint{4.973141in}{0.747735in}}%
\pgfpathlineto{\pgfqpoint{4.973141in}{0.747735in}}%
\pgfpathlineto{\pgfqpoint{4.973141in}{0.750684in}}%
\pgfpathlineto{\pgfqpoint{4.977682in}{0.750684in}}%
\pgfpathlineto{\pgfqpoint{4.977682in}{0.747735in}}%
\pgfpathmoveto{\pgfqpoint{4.973141in}{0.750684in}}%
\pgfpathlineto{\pgfqpoint{4.973141in}{0.750684in}}%
\pgfpathlineto{\pgfqpoint{4.973141in}{0.753633in}}%
\pgfpathlineto{\pgfqpoint{4.977682in}{0.753633in}}%
\pgfpathlineto{\pgfqpoint{4.977682in}{0.750684in}}%
\pgfpathmoveto{\pgfqpoint{4.982223in}{0.735937in}}%
\pgfpathlineto{\pgfqpoint{4.982223in}{0.735937in}}%
\pgfpathlineto{\pgfqpoint{4.982223in}{0.738887in}}%
\pgfpathlineto{\pgfqpoint{4.986764in}{0.738887in}}%
\pgfpathlineto{\pgfqpoint{4.986764in}{0.735937in}}%
\pgfpathmoveto{\pgfqpoint{4.982223in}{0.738887in}}%
\pgfpathlineto{\pgfqpoint{4.982223in}{0.738887in}}%
\pgfpathlineto{\pgfqpoint{4.982223in}{0.741836in}}%
\pgfpathlineto{\pgfqpoint{4.986764in}{0.741836in}}%
\pgfpathlineto{\pgfqpoint{4.986764in}{0.738887in}}%
\pgfpathmoveto{\pgfqpoint{4.964058in}{0.759532in}}%
\pgfpathlineto{\pgfqpoint{4.964058in}{0.759532in}}%
\pgfpathlineto{\pgfqpoint{4.964058in}{0.762481in}}%
\pgfpathlineto{\pgfqpoint{4.968599in}{0.762481in}}%
\pgfpathlineto{\pgfqpoint{4.968599in}{0.759532in}}%
\pgfpathmoveto{\pgfqpoint{4.964058in}{0.762481in}}%
\pgfpathlineto{\pgfqpoint{4.964058in}{0.762481in}}%
\pgfpathlineto{\pgfqpoint{4.964058in}{0.765430in}}%
\pgfpathlineto{\pgfqpoint{4.968599in}{0.765430in}}%
\pgfpathlineto{\pgfqpoint{4.968599in}{0.762481in}}%
\pgfpathmoveto{\pgfqpoint{5.141163in}{0.526541in}}%
\pgfpathlineto{\pgfqpoint{5.141163in}{0.526541in}}%
\pgfpathlineto{\pgfqpoint{5.141163in}{0.529490in}}%
\pgfpathlineto{\pgfqpoint{5.145704in}{0.529490in}}%
\pgfpathlineto{\pgfqpoint{5.145704in}{0.526541in}}%
\pgfpathmoveto{\pgfqpoint{5.136622in}{0.532439in}}%
\pgfpathlineto{\pgfqpoint{5.136622in}{0.532439in}}%
\pgfpathlineto{\pgfqpoint{5.136622in}{0.535388in}}%
\pgfpathlineto{\pgfqpoint{5.141163in}{0.535388in}}%
\pgfpathlineto{\pgfqpoint{5.141163in}{0.532439in}}%
\pgfpathmoveto{\pgfqpoint{5.141163in}{0.529490in}}%
\pgfpathlineto{\pgfqpoint{5.141163in}{0.529490in}}%
\pgfpathlineto{\pgfqpoint{5.141163in}{0.532439in}}%
\pgfpathlineto{\pgfqpoint{5.145704in}{0.532439in}}%
\pgfpathlineto{\pgfqpoint{5.145704in}{0.529490in}}%
\pgfpathmoveto{\pgfqpoint{5.141163in}{0.532439in}}%
\pgfpathlineto{\pgfqpoint{5.141163in}{0.532439in}}%
\pgfpathlineto{\pgfqpoint{5.141163in}{0.535388in}}%
\pgfpathlineto{\pgfqpoint{5.145704in}{0.535388in}}%
\pgfpathlineto{\pgfqpoint{5.145704in}{0.532439in}}%
\pgfpathmoveto{\pgfqpoint{5.132081in}{0.538338in}}%
\pgfpathlineto{\pgfqpoint{5.132081in}{0.538338in}}%
\pgfpathlineto{\pgfqpoint{5.132081in}{0.541287in}}%
\pgfpathlineto{\pgfqpoint{5.136622in}{0.541287in}}%
\pgfpathlineto{\pgfqpoint{5.136622in}{0.538338in}}%
\pgfpathmoveto{\pgfqpoint{5.127540in}{0.544236in}}%
\pgfpathlineto{\pgfqpoint{5.127540in}{0.544236in}}%
\pgfpathlineto{\pgfqpoint{5.127540in}{0.547185in}}%
\pgfpathlineto{\pgfqpoint{5.132081in}{0.547185in}}%
\pgfpathlineto{\pgfqpoint{5.132081in}{0.544236in}}%
\pgfpathmoveto{\pgfqpoint{5.132081in}{0.541287in}}%
\pgfpathlineto{\pgfqpoint{5.132081in}{0.541287in}}%
\pgfpathlineto{\pgfqpoint{5.132081in}{0.544236in}}%
\pgfpathlineto{\pgfqpoint{5.136622in}{0.544236in}}%
\pgfpathlineto{\pgfqpoint{5.136622in}{0.541287in}}%
\pgfpathmoveto{\pgfqpoint{5.132081in}{0.544236in}}%
\pgfpathlineto{\pgfqpoint{5.132081in}{0.544236in}}%
\pgfpathlineto{\pgfqpoint{5.132081in}{0.547185in}}%
\pgfpathlineto{\pgfqpoint{5.136622in}{0.547185in}}%
\pgfpathlineto{\pgfqpoint{5.136622in}{0.544236in}}%
\pgfpathmoveto{\pgfqpoint{5.136622in}{0.535388in}}%
\pgfpathlineto{\pgfqpoint{5.136622in}{0.535388in}}%
\pgfpathlineto{\pgfqpoint{5.136622in}{0.538338in}}%
\pgfpathlineto{\pgfqpoint{5.141163in}{0.538338in}}%
\pgfpathlineto{\pgfqpoint{5.141163in}{0.535388in}}%
\pgfpathmoveto{\pgfqpoint{5.136622in}{0.538338in}}%
\pgfpathlineto{\pgfqpoint{5.136622in}{0.538338in}}%
\pgfpathlineto{\pgfqpoint{5.136622in}{0.541287in}}%
\pgfpathlineto{\pgfqpoint{5.141163in}{0.541287in}}%
\pgfpathlineto{\pgfqpoint{5.141163in}{0.538338in}}%
\pgfpathmoveto{\pgfqpoint{5.159328in}{0.502947in}}%
\pgfpathlineto{\pgfqpoint{5.159328in}{0.502947in}}%
\pgfpathlineto{\pgfqpoint{5.159328in}{0.505896in}}%
\pgfpathlineto{\pgfqpoint{5.163869in}{0.505896in}}%
\pgfpathlineto{\pgfqpoint{5.163869in}{0.502947in}}%
\pgfpathmoveto{\pgfqpoint{5.154787in}{0.508846in}}%
\pgfpathlineto{\pgfqpoint{5.154787in}{0.508846in}}%
\pgfpathlineto{\pgfqpoint{5.154787in}{0.511795in}}%
\pgfpathlineto{\pgfqpoint{5.159328in}{0.511795in}}%
\pgfpathlineto{\pgfqpoint{5.159328in}{0.508846in}}%
\pgfpathmoveto{\pgfqpoint{5.159328in}{0.505896in}}%
\pgfpathlineto{\pgfqpoint{5.159328in}{0.505896in}}%
\pgfpathlineto{\pgfqpoint{5.159328in}{0.508846in}}%
\pgfpathlineto{\pgfqpoint{5.163869in}{0.508846in}}%
\pgfpathlineto{\pgfqpoint{5.163869in}{0.505896in}}%
\pgfpathmoveto{\pgfqpoint{5.159328in}{0.508846in}}%
\pgfpathlineto{\pgfqpoint{5.159328in}{0.508846in}}%
\pgfpathlineto{\pgfqpoint{5.159328in}{0.511795in}}%
\pgfpathlineto{\pgfqpoint{5.163869in}{0.511795in}}%
\pgfpathlineto{\pgfqpoint{5.163869in}{0.508846in}}%
\pgfpathmoveto{\pgfqpoint{5.150245in}{0.514744in}}%
\pgfpathlineto{\pgfqpoint{5.150245in}{0.514744in}}%
\pgfpathlineto{\pgfqpoint{5.150245in}{0.517693in}}%
\pgfpathlineto{\pgfqpoint{5.154787in}{0.517693in}}%
\pgfpathlineto{\pgfqpoint{5.154787in}{0.514744in}}%
\pgfpathmoveto{\pgfqpoint{5.145704in}{0.520642in}}%
\pgfpathlineto{\pgfqpoint{5.145704in}{0.520642in}}%
\pgfpathlineto{\pgfqpoint{5.145704in}{0.523592in}}%
\pgfpathlineto{\pgfqpoint{5.150245in}{0.523592in}}%
\pgfpathlineto{\pgfqpoint{5.150245in}{0.520642in}}%
\pgfpathmoveto{\pgfqpoint{5.150245in}{0.517693in}}%
\pgfpathlineto{\pgfqpoint{5.150245in}{0.517693in}}%
\pgfpathlineto{\pgfqpoint{5.150245in}{0.520642in}}%
\pgfpathlineto{\pgfqpoint{5.154787in}{0.520642in}}%
\pgfpathlineto{\pgfqpoint{5.154787in}{0.517693in}}%
\pgfpathmoveto{\pgfqpoint{5.150245in}{0.520642in}}%
\pgfpathlineto{\pgfqpoint{5.150245in}{0.520642in}}%
\pgfpathlineto{\pgfqpoint{5.150245in}{0.523592in}}%
\pgfpathlineto{\pgfqpoint{5.154787in}{0.523592in}}%
\pgfpathlineto{\pgfqpoint{5.154787in}{0.520642in}}%
\pgfpathmoveto{\pgfqpoint{5.154787in}{0.511795in}}%
\pgfpathlineto{\pgfqpoint{5.154787in}{0.511795in}}%
\pgfpathlineto{\pgfqpoint{5.154787in}{0.514744in}}%
\pgfpathlineto{\pgfqpoint{5.159328in}{0.514744in}}%
\pgfpathlineto{\pgfqpoint{5.159328in}{0.511795in}}%
\pgfpathmoveto{\pgfqpoint{5.154787in}{0.514744in}}%
\pgfpathlineto{\pgfqpoint{5.154787in}{0.514744in}}%
\pgfpathlineto{\pgfqpoint{5.154787in}{0.517693in}}%
\pgfpathlineto{\pgfqpoint{5.159328in}{0.517693in}}%
\pgfpathlineto{\pgfqpoint{5.159328in}{0.514744in}}%
\pgfpathmoveto{\pgfqpoint{5.163869in}{0.499998in}}%
\pgfpathlineto{\pgfqpoint{5.163869in}{0.499998in}}%
\pgfpathlineto{\pgfqpoint{5.163869in}{0.502947in}}%
\pgfpathlineto{\pgfqpoint{5.168410in}{0.502947in}}%
\pgfpathlineto{\pgfqpoint{5.168410in}{0.499998in}}%
\pgfpathmoveto{\pgfqpoint{5.163869in}{0.502947in}}%
\pgfpathlineto{\pgfqpoint{5.163869in}{0.502947in}}%
\pgfpathlineto{\pgfqpoint{5.163869in}{0.505896in}}%
\pgfpathlineto{\pgfqpoint{5.168410in}{0.505896in}}%
\pgfpathlineto{\pgfqpoint{5.168410in}{0.502947in}}%
\pgfpathmoveto{\pgfqpoint{5.145704in}{0.523592in}}%
\pgfpathlineto{\pgfqpoint{5.145704in}{0.523592in}}%
\pgfpathlineto{\pgfqpoint{5.145704in}{0.526541in}}%
\pgfpathlineto{\pgfqpoint{5.150245in}{0.526541in}}%
\pgfpathlineto{\pgfqpoint{5.150245in}{0.523592in}}%
\pgfpathmoveto{\pgfqpoint{5.145704in}{0.526541in}}%
\pgfpathlineto{\pgfqpoint{5.145704in}{0.526541in}}%
\pgfpathlineto{\pgfqpoint{5.145704in}{0.529490in}}%
\pgfpathlineto{\pgfqpoint{5.150245in}{0.529490in}}%
\pgfpathlineto{\pgfqpoint{5.150245in}{0.526541in}}%
\pgfpathmoveto{\pgfqpoint{5.122999in}{0.550134in}}%
\pgfpathlineto{\pgfqpoint{5.122999in}{0.550134in}}%
\pgfpathlineto{\pgfqpoint{5.122999in}{0.553084in}}%
\pgfpathlineto{\pgfqpoint{5.127540in}{0.553084in}}%
\pgfpathlineto{\pgfqpoint{5.127540in}{0.550134in}}%
\pgfpathmoveto{\pgfqpoint{5.118458in}{0.556033in}}%
\pgfpathlineto{\pgfqpoint{5.118458in}{0.556033in}}%
\pgfpathlineto{\pgfqpoint{5.118458in}{0.558982in}}%
\pgfpathlineto{\pgfqpoint{5.122999in}{0.558982in}}%
\pgfpathlineto{\pgfqpoint{5.122999in}{0.556033in}}%
\pgfpathmoveto{\pgfqpoint{5.122999in}{0.553084in}}%
\pgfpathlineto{\pgfqpoint{5.122999in}{0.553084in}}%
\pgfpathlineto{\pgfqpoint{5.122999in}{0.556033in}}%
\pgfpathlineto{\pgfqpoint{5.127540in}{0.556033in}}%
\pgfpathlineto{\pgfqpoint{5.127540in}{0.553084in}}%
\pgfpathmoveto{\pgfqpoint{5.122999in}{0.556033in}}%
\pgfpathlineto{\pgfqpoint{5.122999in}{0.556033in}}%
\pgfpathlineto{\pgfqpoint{5.122999in}{0.558982in}}%
\pgfpathlineto{\pgfqpoint{5.127540in}{0.558982in}}%
\pgfpathlineto{\pgfqpoint{5.127540in}{0.556033in}}%
\pgfpathmoveto{\pgfqpoint{5.113916in}{0.561931in}}%
\pgfpathlineto{\pgfqpoint{5.113916in}{0.561931in}}%
\pgfpathlineto{\pgfqpoint{5.113916in}{0.564880in}}%
\pgfpathlineto{\pgfqpoint{5.118458in}{0.564880in}}%
\pgfpathlineto{\pgfqpoint{5.118458in}{0.561931in}}%
\pgfpathmoveto{\pgfqpoint{5.109375in}{0.567830in}}%
\pgfpathlineto{\pgfqpoint{5.109375in}{0.567830in}}%
\pgfpathlineto{\pgfqpoint{5.109375in}{0.570779in}}%
\pgfpathlineto{\pgfqpoint{5.113916in}{0.570779in}}%
\pgfpathlineto{\pgfqpoint{5.113916in}{0.567830in}}%
\pgfpathmoveto{\pgfqpoint{5.113916in}{0.564880in}}%
\pgfpathlineto{\pgfqpoint{5.113916in}{0.564880in}}%
\pgfpathlineto{\pgfqpoint{5.113916in}{0.567830in}}%
\pgfpathlineto{\pgfqpoint{5.118458in}{0.567830in}}%
\pgfpathlineto{\pgfqpoint{5.118458in}{0.564880in}}%
\pgfpathmoveto{\pgfqpoint{5.113916in}{0.567830in}}%
\pgfpathlineto{\pgfqpoint{5.113916in}{0.567830in}}%
\pgfpathlineto{\pgfqpoint{5.113916in}{0.570779in}}%
\pgfpathlineto{\pgfqpoint{5.118458in}{0.570779in}}%
\pgfpathlineto{\pgfqpoint{5.118458in}{0.567830in}}%
\pgfpathmoveto{\pgfqpoint{5.118458in}{0.558982in}}%
\pgfpathlineto{\pgfqpoint{5.118458in}{0.558982in}}%
\pgfpathlineto{\pgfqpoint{5.118458in}{0.561931in}}%
\pgfpathlineto{\pgfqpoint{5.122999in}{0.561931in}}%
\pgfpathlineto{\pgfqpoint{5.122999in}{0.558982in}}%
\pgfpathmoveto{\pgfqpoint{5.118458in}{0.561931in}}%
\pgfpathlineto{\pgfqpoint{5.118458in}{0.561931in}}%
\pgfpathlineto{\pgfqpoint{5.118458in}{0.564880in}}%
\pgfpathlineto{\pgfqpoint{5.122999in}{0.564880in}}%
\pgfpathlineto{\pgfqpoint{5.122999in}{0.561931in}}%
\pgfpathmoveto{\pgfqpoint{5.127540in}{0.547185in}}%
\pgfpathlineto{\pgfqpoint{5.127540in}{0.547185in}}%
\pgfpathlineto{\pgfqpoint{5.127540in}{0.550134in}}%
\pgfpathlineto{\pgfqpoint{5.132081in}{0.550134in}}%
\pgfpathlineto{\pgfqpoint{5.132081in}{0.547185in}}%
\pgfpathmoveto{\pgfqpoint{5.127540in}{0.550134in}}%
\pgfpathlineto{\pgfqpoint{5.127540in}{0.550134in}}%
\pgfpathlineto{\pgfqpoint{5.127540in}{0.553084in}}%
\pgfpathlineto{\pgfqpoint{5.132081in}{0.553084in}}%
\pgfpathlineto{\pgfqpoint{5.132081in}{0.550134in}}%
\pgfpathmoveto{\pgfqpoint{5.109375in}{0.570779in}}%
\pgfpathlineto{\pgfqpoint{5.109375in}{0.570779in}}%
\pgfpathlineto{\pgfqpoint{5.109375in}{0.573728in}}%
\pgfpathlineto{\pgfqpoint{5.113916in}{0.573728in}}%
\pgfpathlineto{\pgfqpoint{5.113916in}{0.570779in}}%
\pgfpathmoveto{\pgfqpoint{5.109375in}{0.573728in}}%
\pgfpathlineto{\pgfqpoint{5.109375in}{0.573728in}}%
\pgfpathlineto{\pgfqpoint{5.109375in}{0.576677in}}%
\pgfpathlineto{\pgfqpoint{5.113916in}{0.576677in}}%
\pgfpathlineto{\pgfqpoint{5.113916in}{0.573728in}}%
\pgfpathclose%
\pgfusepath{fill}%
\end{pgfscope}%
\begin{pgfscope}%
\pgfpathrectangle{\pgfqpoint{0.750000in}{0.500000in}}{\pgfqpoint{4.650000in}{3.020000in}}%
\pgfusepath{clip}%
\pgfsetbuttcap%
\pgfsetmiterjoin%
\definecolor{currentfill}{rgb}{1.000000,0.000000,0.000000}%
\pgfsetfillcolor{currentfill}%
\pgfsetlinewidth{0.000000pt}%
\definecolor{currentstroke}{rgb}{0.000000,0.000000,0.000000}%
\pgfsetstrokecolor{currentstroke}%
\pgfsetstrokeopacity{0.000000}%
\pgfsetdash{}{0pt}%
\pgfpathmoveto{\pgfqpoint{0.750004in}{0.500002in}}%
\pgfpathlineto{\pgfqpoint{0.750004in}{0.502951in}}%
\pgfpathlineto{\pgfqpoint{0.754545in}{0.502951in}}%
\pgfpathlineto{\pgfqpoint{0.754545in}{0.500002in}}%
\pgfpathmoveto{\pgfqpoint{0.750004in}{0.502951in}}%
\pgfpathlineto{\pgfqpoint{0.750004in}{0.502951in}}%
\pgfpathlineto{\pgfqpoint{0.750004in}{0.505901in}}%
\pgfpathlineto{\pgfqpoint{0.754545in}{0.505901in}}%
\pgfpathlineto{\pgfqpoint{0.754545in}{0.502951in}}%
\pgfpathmoveto{\pgfqpoint{0.754545in}{0.502951in}}%
\pgfpathlineto{\pgfqpoint{0.754545in}{0.502951in}}%
\pgfpathlineto{\pgfqpoint{0.754545in}{0.505901in}}%
\pgfpathlineto{\pgfqpoint{0.759086in}{0.505901in}}%
\pgfpathlineto{\pgfqpoint{0.759086in}{0.502951in}}%
\pgfpathmoveto{\pgfqpoint{0.754545in}{0.505901in}}%
\pgfpathlineto{\pgfqpoint{0.754545in}{0.505901in}}%
\pgfpathlineto{\pgfqpoint{0.754545in}{0.508850in}}%
\pgfpathlineto{\pgfqpoint{0.759086in}{0.508850in}}%
\pgfpathlineto{\pgfqpoint{0.759086in}{0.505901in}}%
\pgfpathmoveto{\pgfqpoint{0.759086in}{0.505901in}}%
\pgfpathlineto{\pgfqpoint{0.759086in}{0.505901in}}%
\pgfpathlineto{\pgfqpoint{0.759086in}{0.508850in}}%
\pgfpathlineto{\pgfqpoint{0.763627in}{0.508850in}}%
\pgfpathlineto{\pgfqpoint{0.763627in}{0.505901in}}%
\pgfpathmoveto{\pgfqpoint{0.759086in}{0.508850in}}%
\pgfpathlineto{\pgfqpoint{0.759086in}{0.508850in}}%
\pgfpathlineto{\pgfqpoint{0.759086in}{0.511799in}}%
\pgfpathlineto{\pgfqpoint{0.763627in}{0.511799in}}%
\pgfpathlineto{\pgfqpoint{0.763627in}{0.508850in}}%
\pgfpathmoveto{\pgfqpoint{0.763627in}{0.508850in}}%
\pgfpathlineto{\pgfqpoint{0.763627in}{0.508850in}}%
\pgfpathlineto{\pgfqpoint{0.763627in}{0.511799in}}%
\pgfpathlineto{\pgfqpoint{0.768168in}{0.511799in}}%
\pgfpathlineto{\pgfqpoint{0.768168in}{0.508850in}}%
\pgfpathmoveto{\pgfqpoint{0.763627in}{0.511799in}}%
\pgfpathlineto{\pgfqpoint{0.763627in}{0.511799in}}%
\pgfpathlineto{\pgfqpoint{0.763627in}{0.514748in}}%
\pgfpathlineto{\pgfqpoint{0.768168in}{0.514748in}}%
\pgfpathlineto{\pgfqpoint{0.768168in}{0.511799in}}%
\pgfpathmoveto{\pgfqpoint{0.768168in}{0.511799in}}%
\pgfpathlineto{\pgfqpoint{0.768168in}{0.511799in}}%
\pgfpathlineto{\pgfqpoint{0.768168in}{0.514748in}}%
\pgfpathlineto{\pgfqpoint{0.772709in}{0.514748in}}%
\pgfpathlineto{\pgfqpoint{0.772709in}{0.511799in}}%
\pgfpathmoveto{\pgfqpoint{0.768168in}{0.514748in}}%
\pgfpathlineto{\pgfqpoint{0.768168in}{0.514748in}}%
\pgfpathlineto{\pgfqpoint{0.768168in}{0.517697in}}%
\pgfpathlineto{\pgfqpoint{0.772709in}{0.517697in}}%
\pgfpathlineto{\pgfqpoint{0.772709in}{0.514748in}}%
\pgfpathmoveto{\pgfqpoint{0.772709in}{0.514748in}}%
\pgfpathlineto{\pgfqpoint{0.772709in}{0.514748in}}%
\pgfpathlineto{\pgfqpoint{0.772709in}{0.517697in}}%
\pgfpathlineto{\pgfqpoint{0.777250in}{0.517697in}}%
\pgfpathlineto{\pgfqpoint{0.777250in}{0.514748in}}%
\pgfpathmoveto{\pgfqpoint{0.772709in}{0.517697in}}%
\pgfpathlineto{\pgfqpoint{0.772709in}{0.517697in}}%
\pgfpathlineto{\pgfqpoint{0.772709in}{0.520646in}}%
\pgfpathlineto{\pgfqpoint{0.777250in}{0.520646in}}%
\pgfpathlineto{\pgfqpoint{0.777250in}{0.517697in}}%
\pgfpathmoveto{\pgfqpoint{0.777250in}{0.517697in}}%
\pgfpathlineto{\pgfqpoint{0.777250in}{0.517697in}}%
\pgfpathlineto{\pgfqpoint{0.777250in}{0.520646in}}%
\pgfpathlineto{\pgfqpoint{0.781791in}{0.520646in}}%
\pgfpathlineto{\pgfqpoint{0.781791in}{0.517697in}}%
\pgfpathmoveto{\pgfqpoint{0.777250in}{0.520646in}}%
\pgfpathlineto{\pgfqpoint{0.777250in}{0.520646in}}%
\pgfpathlineto{\pgfqpoint{0.777250in}{0.523595in}}%
\pgfpathlineto{\pgfqpoint{0.781791in}{0.523595in}}%
\pgfpathlineto{\pgfqpoint{0.781791in}{0.520646in}}%
\pgfpathmoveto{\pgfqpoint{0.781791in}{0.520646in}}%
\pgfpathlineto{\pgfqpoint{0.781791in}{0.520646in}}%
\pgfpathlineto{\pgfqpoint{0.781791in}{0.523595in}}%
\pgfpathlineto{\pgfqpoint{0.786332in}{0.523595in}}%
\pgfpathlineto{\pgfqpoint{0.786332in}{0.520646in}}%
\pgfpathmoveto{\pgfqpoint{0.781791in}{0.523595in}}%
\pgfpathlineto{\pgfqpoint{0.781791in}{0.523595in}}%
\pgfpathlineto{\pgfqpoint{0.781791in}{0.526544in}}%
\pgfpathlineto{\pgfqpoint{0.786332in}{0.526544in}}%
\pgfpathlineto{\pgfqpoint{0.786332in}{0.523595in}}%
\pgfpathmoveto{\pgfqpoint{0.786332in}{0.523595in}}%
\pgfpathlineto{\pgfqpoint{0.786332in}{0.523595in}}%
\pgfpathlineto{\pgfqpoint{0.786332in}{0.526544in}}%
\pgfpathlineto{\pgfqpoint{0.790873in}{0.526544in}}%
\pgfpathlineto{\pgfqpoint{0.790873in}{0.523595in}}%
\pgfpathmoveto{\pgfqpoint{0.786332in}{0.526544in}}%
\pgfpathlineto{\pgfqpoint{0.786332in}{0.526544in}}%
\pgfpathlineto{\pgfqpoint{0.786332in}{0.529493in}}%
\pgfpathlineto{\pgfqpoint{0.790873in}{0.529493in}}%
\pgfpathlineto{\pgfqpoint{0.790873in}{0.526544in}}%
\pgfpathmoveto{\pgfqpoint{0.790873in}{0.526544in}}%
\pgfpathlineto{\pgfqpoint{0.790873in}{0.526544in}}%
\pgfpathlineto{\pgfqpoint{0.790873in}{0.529493in}}%
\pgfpathlineto{\pgfqpoint{0.795414in}{0.529493in}}%
\pgfpathlineto{\pgfqpoint{0.795414in}{0.526544in}}%
\pgfpathmoveto{\pgfqpoint{0.790873in}{0.529493in}}%
\pgfpathlineto{\pgfqpoint{0.790873in}{0.529493in}}%
\pgfpathlineto{\pgfqpoint{0.790873in}{0.532442in}}%
\pgfpathlineto{\pgfqpoint{0.795414in}{0.532442in}}%
\pgfpathlineto{\pgfqpoint{0.795414in}{0.529493in}}%
\pgfpathmoveto{\pgfqpoint{0.795414in}{0.529493in}}%
\pgfpathlineto{\pgfqpoint{0.795414in}{0.529493in}}%
\pgfpathlineto{\pgfqpoint{0.795414in}{0.532442in}}%
\pgfpathlineto{\pgfqpoint{0.799955in}{0.532442in}}%
\pgfpathlineto{\pgfqpoint{0.799955in}{0.529493in}}%
\pgfpathmoveto{\pgfqpoint{0.795414in}{0.532442in}}%
\pgfpathlineto{\pgfqpoint{0.795414in}{0.532442in}}%
\pgfpathlineto{\pgfqpoint{0.795414in}{0.535392in}}%
\pgfpathlineto{\pgfqpoint{0.799955in}{0.535392in}}%
\pgfpathlineto{\pgfqpoint{0.799955in}{0.532442in}}%
\pgfpathmoveto{\pgfqpoint{0.799955in}{0.532442in}}%
\pgfpathlineto{\pgfqpoint{0.799955in}{0.532442in}}%
\pgfpathlineto{\pgfqpoint{0.799955in}{0.535392in}}%
\pgfpathlineto{\pgfqpoint{0.804496in}{0.535392in}}%
\pgfpathlineto{\pgfqpoint{0.804496in}{0.532442in}}%
\pgfpathmoveto{\pgfqpoint{0.799955in}{0.535392in}}%
\pgfpathlineto{\pgfqpoint{0.799955in}{0.535392in}}%
\pgfpathlineto{\pgfqpoint{0.799955in}{0.538341in}}%
\pgfpathlineto{\pgfqpoint{0.804496in}{0.538341in}}%
\pgfpathlineto{\pgfqpoint{0.804496in}{0.535392in}}%
\pgfpathmoveto{\pgfqpoint{0.804496in}{0.535392in}}%
\pgfpathlineto{\pgfqpoint{0.804496in}{0.535392in}}%
\pgfpathlineto{\pgfqpoint{0.804496in}{0.538341in}}%
\pgfpathlineto{\pgfqpoint{0.809037in}{0.538341in}}%
\pgfpathlineto{\pgfqpoint{0.809037in}{0.535392in}}%
\pgfpathmoveto{\pgfqpoint{0.804496in}{0.538341in}}%
\pgfpathlineto{\pgfqpoint{0.804496in}{0.538341in}}%
\pgfpathlineto{\pgfqpoint{0.804496in}{0.541290in}}%
\pgfpathlineto{\pgfqpoint{0.809037in}{0.541290in}}%
\pgfpathlineto{\pgfqpoint{0.809037in}{0.538341in}}%
\pgfpathmoveto{\pgfqpoint{0.809037in}{0.538341in}}%
\pgfpathlineto{\pgfqpoint{0.809037in}{0.538341in}}%
\pgfpathlineto{\pgfqpoint{0.809037in}{0.541290in}}%
\pgfpathlineto{\pgfqpoint{0.813578in}{0.541290in}}%
\pgfpathlineto{\pgfqpoint{0.813578in}{0.538341in}}%
\pgfpathmoveto{\pgfqpoint{0.809037in}{0.541290in}}%
\pgfpathlineto{\pgfqpoint{0.809037in}{0.541290in}}%
\pgfpathlineto{\pgfqpoint{0.809037in}{0.544239in}}%
\pgfpathlineto{\pgfqpoint{0.813578in}{0.544239in}}%
\pgfpathlineto{\pgfqpoint{0.813578in}{0.541290in}}%
\pgfpathmoveto{\pgfqpoint{0.813578in}{0.541290in}}%
\pgfpathlineto{\pgfqpoint{0.813578in}{0.541290in}}%
\pgfpathlineto{\pgfqpoint{0.813578in}{0.544239in}}%
\pgfpathlineto{\pgfqpoint{0.818120in}{0.544239in}}%
\pgfpathlineto{\pgfqpoint{0.818120in}{0.541290in}}%
\pgfpathmoveto{\pgfqpoint{0.813578in}{0.544239in}}%
\pgfpathlineto{\pgfqpoint{0.813578in}{0.544239in}}%
\pgfpathlineto{\pgfqpoint{0.813578in}{0.547188in}}%
\pgfpathlineto{\pgfqpoint{0.818120in}{0.547188in}}%
\pgfpathlineto{\pgfqpoint{0.818120in}{0.544239in}}%
\pgfpathmoveto{\pgfqpoint{0.818120in}{0.544239in}}%
\pgfpathlineto{\pgfqpoint{0.818120in}{0.544239in}}%
\pgfpathlineto{\pgfqpoint{0.818120in}{0.547188in}}%
\pgfpathlineto{\pgfqpoint{0.822661in}{0.547188in}}%
\pgfpathlineto{\pgfqpoint{0.822661in}{0.544239in}}%
\pgfpathmoveto{\pgfqpoint{0.818120in}{0.547188in}}%
\pgfpathlineto{\pgfqpoint{0.818120in}{0.547188in}}%
\pgfpathlineto{\pgfqpoint{0.818120in}{0.550137in}}%
\pgfpathlineto{\pgfqpoint{0.822661in}{0.550137in}}%
\pgfpathlineto{\pgfqpoint{0.822661in}{0.547188in}}%
\pgfpathmoveto{\pgfqpoint{0.822661in}{0.547188in}}%
\pgfpathlineto{\pgfqpoint{0.822661in}{0.547188in}}%
\pgfpathlineto{\pgfqpoint{0.822661in}{0.550137in}}%
\pgfpathlineto{\pgfqpoint{0.827202in}{0.550137in}}%
\pgfpathlineto{\pgfqpoint{0.827202in}{0.547188in}}%
\pgfpathmoveto{\pgfqpoint{0.822661in}{0.550137in}}%
\pgfpathlineto{\pgfqpoint{0.822661in}{0.550137in}}%
\pgfpathlineto{\pgfqpoint{0.822661in}{0.553086in}}%
\pgfpathlineto{\pgfqpoint{0.827202in}{0.553086in}}%
\pgfpathlineto{\pgfqpoint{0.827202in}{0.550137in}}%
\pgfpathmoveto{\pgfqpoint{0.827202in}{0.550137in}}%
\pgfpathlineto{\pgfqpoint{0.827202in}{0.550137in}}%
\pgfpathlineto{\pgfqpoint{0.827202in}{0.553086in}}%
\pgfpathlineto{\pgfqpoint{0.831743in}{0.553086in}}%
\pgfpathlineto{\pgfqpoint{0.831743in}{0.550137in}}%
\pgfpathmoveto{\pgfqpoint{0.827202in}{0.553086in}}%
\pgfpathlineto{\pgfqpoint{0.827202in}{0.553086in}}%
\pgfpathlineto{\pgfqpoint{0.827202in}{0.556035in}}%
\pgfpathlineto{\pgfqpoint{0.831743in}{0.556035in}}%
\pgfpathlineto{\pgfqpoint{0.831743in}{0.553086in}}%
\pgfpathmoveto{\pgfqpoint{0.831743in}{0.553086in}}%
\pgfpathlineto{\pgfqpoint{0.831743in}{0.553086in}}%
\pgfpathlineto{\pgfqpoint{0.831743in}{0.556035in}}%
\pgfpathlineto{\pgfqpoint{0.836284in}{0.556035in}}%
\pgfpathlineto{\pgfqpoint{0.836284in}{0.553086in}}%
\pgfpathmoveto{\pgfqpoint{0.831743in}{0.556035in}}%
\pgfpathlineto{\pgfqpoint{0.831743in}{0.556035in}}%
\pgfpathlineto{\pgfqpoint{0.831743in}{0.558984in}}%
\pgfpathlineto{\pgfqpoint{0.836284in}{0.558984in}}%
\pgfpathlineto{\pgfqpoint{0.836284in}{0.556035in}}%
\pgfpathmoveto{\pgfqpoint{0.836284in}{0.556035in}}%
\pgfpathlineto{\pgfqpoint{0.836284in}{0.556035in}}%
\pgfpathlineto{\pgfqpoint{0.836284in}{0.558984in}}%
\pgfpathlineto{\pgfqpoint{0.840825in}{0.558984in}}%
\pgfpathlineto{\pgfqpoint{0.840825in}{0.556035in}}%
\pgfpathmoveto{\pgfqpoint{0.836284in}{0.558984in}}%
\pgfpathlineto{\pgfqpoint{0.836284in}{0.558984in}}%
\pgfpathlineto{\pgfqpoint{0.836284in}{0.561933in}}%
\pgfpathlineto{\pgfqpoint{0.840825in}{0.561933in}}%
\pgfpathlineto{\pgfqpoint{0.840825in}{0.558984in}}%
\pgfpathmoveto{\pgfqpoint{0.840825in}{0.558984in}}%
\pgfpathlineto{\pgfqpoint{0.840825in}{0.558984in}}%
\pgfpathlineto{\pgfqpoint{0.840825in}{0.561933in}}%
\pgfpathlineto{\pgfqpoint{0.845366in}{0.561933in}}%
\pgfpathlineto{\pgfqpoint{0.845366in}{0.558984in}}%
\pgfpathmoveto{\pgfqpoint{0.840825in}{0.561933in}}%
\pgfpathlineto{\pgfqpoint{0.840825in}{0.561933in}}%
\pgfpathlineto{\pgfqpoint{0.840825in}{0.564883in}}%
\pgfpathlineto{\pgfqpoint{0.845366in}{0.564883in}}%
\pgfpathlineto{\pgfqpoint{0.845366in}{0.561933in}}%
\pgfpathmoveto{\pgfqpoint{0.845366in}{0.561933in}}%
\pgfpathlineto{\pgfqpoint{0.845366in}{0.561933in}}%
\pgfpathlineto{\pgfqpoint{0.845366in}{0.564883in}}%
\pgfpathlineto{\pgfqpoint{0.849907in}{0.564883in}}%
\pgfpathlineto{\pgfqpoint{0.849907in}{0.561933in}}%
\pgfpathmoveto{\pgfqpoint{0.845366in}{0.564883in}}%
\pgfpathlineto{\pgfqpoint{0.845366in}{0.564883in}}%
\pgfpathlineto{\pgfqpoint{0.845366in}{0.567832in}}%
\pgfpathlineto{\pgfqpoint{0.849907in}{0.567832in}}%
\pgfpathlineto{\pgfqpoint{0.849907in}{0.564883in}}%
\pgfpathmoveto{\pgfqpoint{0.849907in}{0.564883in}}%
\pgfpathlineto{\pgfqpoint{0.849907in}{0.564883in}}%
\pgfpathlineto{\pgfqpoint{0.849907in}{0.567832in}}%
\pgfpathlineto{\pgfqpoint{0.854448in}{0.567832in}}%
\pgfpathlineto{\pgfqpoint{0.854448in}{0.564883in}}%
\pgfpathmoveto{\pgfqpoint{0.849907in}{0.567832in}}%
\pgfpathlineto{\pgfqpoint{0.849907in}{0.567832in}}%
\pgfpathlineto{\pgfqpoint{0.849907in}{0.570781in}}%
\pgfpathlineto{\pgfqpoint{0.854448in}{0.570781in}}%
\pgfpathlineto{\pgfqpoint{0.854448in}{0.567832in}}%
\pgfpathmoveto{\pgfqpoint{0.854448in}{0.567832in}}%
\pgfpathlineto{\pgfqpoint{0.854448in}{0.567832in}}%
\pgfpathlineto{\pgfqpoint{0.854448in}{0.570781in}}%
\pgfpathlineto{\pgfqpoint{0.858989in}{0.570781in}}%
\pgfpathlineto{\pgfqpoint{0.858989in}{0.567832in}}%
\pgfpathmoveto{\pgfqpoint{0.854448in}{0.570781in}}%
\pgfpathlineto{\pgfqpoint{0.854448in}{0.570781in}}%
\pgfpathlineto{\pgfqpoint{0.854448in}{0.573730in}}%
\pgfpathlineto{\pgfqpoint{0.858989in}{0.573730in}}%
\pgfpathlineto{\pgfqpoint{0.858989in}{0.570781in}}%
\pgfpathmoveto{\pgfqpoint{0.858989in}{0.570781in}}%
\pgfpathlineto{\pgfqpoint{0.858989in}{0.570781in}}%
\pgfpathlineto{\pgfqpoint{0.858989in}{0.573730in}}%
\pgfpathlineto{\pgfqpoint{0.863530in}{0.573730in}}%
\pgfpathlineto{\pgfqpoint{0.863530in}{0.570781in}}%
\pgfpathmoveto{\pgfqpoint{0.858989in}{0.573730in}}%
\pgfpathlineto{\pgfqpoint{0.858989in}{0.573730in}}%
\pgfpathlineto{\pgfqpoint{0.858989in}{0.576679in}}%
\pgfpathlineto{\pgfqpoint{0.863530in}{0.576679in}}%
\pgfpathlineto{\pgfqpoint{0.863530in}{0.573730in}}%
\pgfpathmoveto{\pgfqpoint{0.863530in}{0.573730in}}%
\pgfpathlineto{\pgfqpoint{0.863530in}{0.573730in}}%
\pgfpathlineto{\pgfqpoint{0.863530in}{0.576679in}}%
\pgfpathlineto{\pgfqpoint{0.868071in}{0.576679in}}%
\pgfpathlineto{\pgfqpoint{0.868071in}{0.573730in}}%
\pgfpathmoveto{\pgfqpoint{0.863530in}{0.576679in}}%
\pgfpathlineto{\pgfqpoint{0.863530in}{0.576679in}}%
\pgfpathlineto{\pgfqpoint{0.863530in}{0.579628in}}%
\pgfpathlineto{\pgfqpoint{0.868071in}{0.579628in}}%
\pgfpathlineto{\pgfqpoint{0.868071in}{0.576679in}}%
\pgfpathmoveto{\pgfqpoint{0.868071in}{0.576679in}}%
\pgfpathlineto{\pgfqpoint{0.868071in}{0.576679in}}%
\pgfpathlineto{\pgfqpoint{0.868071in}{0.579628in}}%
\pgfpathlineto{\pgfqpoint{0.872612in}{0.579628in}}%
\pgfpathlineto{\pgfqpoint{0.872612in}{0.576679in}}%
\pgfpathmoveto{\pgfqpoint{0.868071in}{0.579628in}}%
\pgfpathlineto{\pgfqpoint{0.868071in}{0.579628in}}%
\pgfpathlineto{\pgfqpoint{0.868071in}{0.582577in}}%
\pgfpathlineto{\pgfqpoint{0.872612in}{0.582577in}}%
\pgfpathlineto{\pgfqpoint{0.872612in}{0.579628in}}%
\pgfpathmoveto{\pgfqpoint{0.872612in}{0.579628in}}%
\pgfpathlineto{\pgfqpoint{0.872612in}{0.579628in}}%
\pgfpathlineto{\pgfqpoint{0.872612in}{0.582577in}}%
\pgfpathlineto{\pgfqpoint{0.877153in}{0.582577in}}%
\pgfpathlineto{\pgfqpoint{0.877153in}{0.579628in}}%
\pgfpathmoveto{\pgfqpoint{0.872612in}{0.582577in}}%
\pgfpathlineto{\pgfqpoint{0.872612in}{0.582577in}}%
\pgfpathlineto{\pgfqpoint{0.872612in}{0.585526in}}%
\pgfpathlineto{\pgfqpoint{0.877153in}{0.585526in}}%
\pgfpathlineto{\pgfqpoint{0.877153in}{0.582577in}}%
\pgfpathmoveto{\pgfqpoint{0.877153in}{0.582577in}}%
\pgfpathlineto{\pgfqpoint{0.877153in}{0.582577in}}%
\pgfpathlineto{\pgfqpoint{0.877153in}{0.585526in}}%
\pgfpathlineto{\pgfqpoint{0.881694in}{0.585526in}}%
\pgfpathlineto{\pgfqpoint{0.881694in}{0.582577in}}%
\pgfpathmoveto{\pgfqpoint{0.877153in}{0.585526in}}%
\pgfpathlineto{\pgfqpoint{0.877153in}{0.585526in}}%
\pgfpathlineto{\pgfqpoint{0.877153in}{0.588475in}}%
\pgfpathlineto{\pgfqpoint{0.881694in}{0.588475in}}%
\pgfpathlineto{\pgfqpoint{0.881694in}{0.585526in}}%
\pgfpathmoveto{\pgfqpoint{0.881694in}{0.585526in}}%
\pgfpathlineto{\pgfqpoint{0.881694in}{0.585526in}}%
\pgfpathlineto{\pgfqpoint{0.881694in}{0.588475in}}%
\pgfpathlineto{\pgfqpoint{0.886235in}{0.588475in}}%
\pgfpathlineto{\pgfqpoint{0.886235in}{0.585526in}}%
\pgfpathmoveto{\pgfqpoint{0.881694in}{0.588475in}}%
\pgfpathlineto{\pgfqpoint{0.881694in}{0.588475in}}%
\pgfpathlineto{\pgfqpoint{0.881694in}{0.591424in}}%
\pgfpathlineto{\pgfqpoint{0.886235in}{0.591424in}}%
\pgfpathlineto{\pgfqpoint{0.886235in}{0.588475in}}%
\pgfpathmoveto{\pgfqpoint{0.886235in}{0.588475in}}%
\pgfpathlineto{\pgfqpoint{0.886235in}{0.588475in}}%
\pgfpathlineto{\pgfqpoint{0.886235in}{0.591424in}}%
\pgfpathlineto{\pgfqpoint{0.890776in}{0.591424in}}%
\pgfpathlineto{\pgfqpoint{0.890776in}{0.588475in}}%
\pgfpathmoveto{\pgfqpoint{0.886235in}{0.591424in}}%
\pgfpathlineto{\pgfqpoint{0.886235in}{0.591424in}}%
\pgfpathlineto{\pgfqpoint{0.886235in}{0.594374in}}%
\pgfpathlineto{\pgfqpoint{0.890776in}{0.594374in}}%
\pgfpathlineto{\pgfqpoint{0.890776in}{0.591424in}}%
\pgfpathmoveto{\pgfqpoint{0.890776in}{0.591424in}}%
\pgfpathlineto{\pgfqpoint{0.890776in}{0.591424in}}%
\pgfpathlineto{\pgfqpoint{0.890776in}{0.594374in}}%
\pgfpathlineto{\pgfqpoint{0.895317in}{0.594374in}}%
\pgfpathlineto{\pgfqpoint{0.895317in}{0.591424in}}%
\pgfpathmoveto{\pgfqpoint{0.890776in}{0.594374in}}%
\pgfpathlineto{\pgfqpoint{0.890776in}{0.594374in}}%
\pgfpathlineto{\pgfqpoint{0.890776in}{0.597323in}}%
\pgfpathlineto{\pgfqpoint{0.895317in}{0.597323in}}%
\pgfpathlineto{\pgfqpoint{0.895317in}{0.594374in}}%
\pgfpathmoveto{\pgfqpoint{0.895317in}{0.594374in}}%
\pgfpathlineto{\pgfqpoint{0.895317in}{0.594374in}}%
\pgfpathlineto{\pgfqpoint{0.895317in}{0.597323in}}%
\pgfpathlineto{\pgfqpoint{0.899858in}{0.597323in}}%
\pgfpathlineto{\pgfqpoint{0.899858in}{0.594374in}}%
\pgfpathmoveto{\pgfqpoint{0.895317in}{0.597323in}}%
\pgfpathlineto{\pgfqpoint{0.895317in}{0.597323in}}%
\pgfpathlineto{\pgfqpoint{0.895317in}{0.600272in}}%
\pgfpathlineto{\pgfqpoint{0.899858in}{0.600272in}}%
\pgfpathlineto{\pgfqpoint{0.899858in}{0.597323in}}%
\pgfpathmoveto{\pgfqpoint{0.899858in}{0.597323in}}%
\pgfpathlineto{\pgfqpoint{0.899858in}{0.597323in}}%
\pgfpathlineto{\pgfqpoint{0.899858in}{0.600272in}}%
\pgfpathlineto{\pgfqpoint{0.904398in}{0.600272in}}%
\pgfpathlineto{\pgfqpoint{0.904398in}{0.597323in}}%
\pgfpathmoveto{\pgfqpoint{0.899858in}{0.600272in}}%
\pgfpathlineto{\pgfqpoint{0.899858in}{0.600272in}}%
\pgfpathlineto{\pgfqpoint{0.899858in}{0.603221in}}%
\pgfpathlineto{\pgfqpoint{0.904398in}{0.603221in}}%
\pgfpathlineto{\pgfqpoint{0.904398in}{0.600272in}}%
\pgfpathmoveto{\pgfqpoint{0.904398in}{0.600272in}}%
\pgfpathlineto{\pgfqpoint{0.904398in}{0.600272in}}%
\pgfpathlineto{\pgfqpoint{0.904398in}{0.603221in}}%
\pgfpathlineto{\pgfqpoint{0.908939in}{0.603221in}}%
\pgfpathlineto{\pgfqpoint{0.908939in}{0.600272in}}%
\pgfpathmoveto{\pgfqpoint{0.904398in}{0.603221in}}%
\pgfpathlineto{\pgfqpoint{0.904398in}{0.603221in}}%
\pgfpathlineto{\pgfqpoint{0.904398in}{0.606171in}}%
\pgfpathlineto{\pgfqpoint{0.908939in}{0.606171in}}%
\pgfpathlineto{\pgfqpoint{0.908939in}{0.603221in}}%
\pgfpathmoveto{\pgfqpoint{0.908939in}{0.603221in}}%
\pgfpathlineto{\pgfqpoint{0.908939in}{0.603221in}}%
\pgfpathlineto{\pgfqpoint{0.908939in}{0.606171in}}%
\pgfpathlineto{\pgfqpoint{0.913480in}{0.606171in}}%
\pgfpathlineto{\pgfqpoint{0.913480in}{0.603221in}}%
\pgfpathmoveto{\pgfqpoint{0.908939in}{0.606171in}}%
\pgfpathlineto{\pgfqpoint{0.908939in}{0.606171in}}%
\pgfpathlineto{\pgfqpoint{0.908939in}{0.609120in}}%
\pgfpathlineto{\pgfqpoint{0.913480in}{0.609120in}}%
\pgfpathlineto{\pgfqpoint{0.913480in}{0.606171in}}%
\pgfpathmoveto{\pgfqpoint{0.913480in}{0.606171in}}%
\pgfpathlineto{\pgfqpoint{0.913480in}{0.606171in}}%
\pgfpathlineto{\pgfqpoint{0.913480in}{0.609120in}}%
\pgfpathlineto{\pgfqpoint{0.918021in}{0.609120in}}%
\pgfpathlineto{\pgfqpoint{0.918021in}{0.606171in}}%
\pgfpathmoveto{\pgfqpoint{0.913480in}{0.609120in}}%
\pgfpathlineto{\pgfqpoint{0.913480in}{0.609120in}}%
\pgfpathlineto{\pgfqpoint{0.913480in}{0.612069in}}%
\pgfpathlineto{\pgfqpoint{0.918021in}{0.612069in}}%
\pgfpathlineto{\pgfqpoint{0.918021in}{0.609120in}}%
\pgfpathmoveto{\pgfqpoint{0.918021in}{0.609120in}}%
\pgfpathlineto{\pgfqpoint{0.918021in}{0.609120in}}%
\pgfpathlineto{\pgfqpoint{0.918021in}{0.612069in}}%
\pgfpathlineto{\pgfqpoint{0.922562in}{0.612069in}}%
\pgfpathlineto{\pgfqpoint{0.922562in}{0.609120in}}%
\pgfpathmoveto{\pgfqpoint{0.918021in}{0.612069in}}%
\pgfpathlineto{\pgfqpoint{0.918021in}{0.612069in}}%
\pgfpathlineto{\pgfqpoint{0.918021in}{0.615019in}}%
\pgfpathlineto{\pgfqpoint{0.922562in}{0.615019in}}%
\pgfpathlineto{\pgfqpoint{0.922562in}{0.612069in}}%
\pgfpathmoveto{\pgfqpoint{0.922562in}{0.612069in}}%
\pgfpathlineto{\pgfqpoint{0.922562in}{0.612069in}}%
\pgfpathlineto{\pgfqpoint{0.922562in}{0.615019in}}%
\pgfpathlineto{\pgfqpoint{0.927102in}{0.615019in}}%
\pgfpathlineto{\pgfqpoint{0.927102in}{0.612069in}}%
\pgfpathmoveto{\pgfqpoint{0.922562in}{0.615019in}}%
\pgfpathlineto{\pgfqpoint{0.922562in}{0.615019in}}%
\pgfpathlineto{\pgfqpoint{0.922562in}{0.617968in}}%
\pgfpathlineto{\pgfqpoint{0.927102in}{0.617968in}}%
\pgfpathlineto{\pgfqpoint{0.927102in}{0.615019in}}%
\pgfpathmoveto{\pgfqpoint{0.927102in}{0.615019in}}%
\pgfpathlineto{\pgfqpoint{0.927102in}{0.615019in}}%
\pgfpathlineto{\pgfqpoint{0.927102in}{0.617968in}}%
\pgfpathlineto{\pgfqpoint{0.931643in}{0.617968in}}%
\pgfpathlineto{\pgfqpoint{0.931643in}{0.615019in}}%
\pgfpathmoveto{\pgfqpoint{0.927102in}{0.617968in}}%
\pgfpathlineto{\pgfqpoint{0.927102in}{0.617968in}}%
\pgfpathlineto{\pgfqpoint{0.927102in}{0.620917in}}%
\pgfpathlineto{\pgfqpoint{0.931643in}{0.620917in}}%
\pgfpathlineto{\pgfqpoint{0.931643in}{0.617968in}}%
\pgfpathmoveto{\pgfqpoint{0.931643in}{0.617968in}}%
\pgfpathlineto{\pgfqpoint{0.931643in}{0.617968in}}%
\pgfpathlineto{\pgfqpoint{0.931643in}{0.620917in}}%
\pgfpathlineto{\pgfqpoint{0.936184in}{0.620917in}}%
\pgfpathlineto{\pgfqpoint{0.936184in}{0.617968in}}%
\pgfpathmoveto{\pgfqpoint{0.931643in}{0.620917in}}%
\pgfpathlineto{\pgfqpoint{0.931643in}{0.620917in}}%
\pgfpathlineto{\pgfqpoint{0.931643in}{0.623867in}}%
\pgfpathlineto{\pgfqpoint{0.936184in}{0.623867in}}%
\pgfpathlineto{\pgfqpoint{0.936184in}{0.620917in}}%
\pgfpathmoveto{\pgfqpoint{0.936184in}{0.620917in}}%
\pgfpathlineto{\pgfqpoint{0.936184in}{0.620917in}}%
\pgfpathlineto{\pgfqpoint{0.936184in}{0.623867in}}%
\pgfpathlineto{\pgfqpoint{0.940725in}{0.623867in}}%
\pgfpathlineto{\pgfqpoint{0.940725in}{0.620917in}}%
\pgfpathmoveto{\pgfqpoint{0.936184in}{0.623867in}}%
\pgfpathlineto{\pgfqpoint{0.936184in}{0.623867in}}%
\pgfpathlineto{\pgfqpoint{0.936184in}{0.626816in}}%
\pgfpathlineto{\pgfqpoint{0.940725in}{0.626816in}}%
\pgfpathlineto{\pgfqpoint{0.940725in}{0.623867in}}%
\pgfpathmoveto{\pgfqpoint{0.940725in}{0.623867in}}%
\pgfpathlineto{\pgfqpoint{0.940725in}{0.623867in}}%
\pgfpathlineto{\pgfqpoint{0.940725in}{0.626816in}}%
\pgfpathlineto{\pgfqpoint{0.945265in}{0.626816in}}%
\pgfpathlineto{\pgfqpoint{0.945265in}{0.623867in}}%
\pgfpathmoveto{\pgfqpoint{0.940725in}{0.626816in}}%
\pgfpathlineto{\pgfqpoint{0.940725in}{0.626816in}}%
\pgfpathlineto{\pgfqpoint{0.940725in}{0.629765in}}%
\pgfpathlineto{\pgfqpoint{0.945265in}{0.629765in}}%
\pgfpathlineto{\pgfqpoint{0.945265in}{0.626816in}}%
\pgfpathmoveto{\pgfqpoint{0.945265in}{0.626816in}}%
\pgfpathlineto{\pgfqpoint{0.945265in}{0.626816in}}%
\pgfpathlineto{\pgfqpoint{0.945265in}{0.629765in}}%
\pgfpathlineto{\pgfqpoint{0.949806in}{0.629765in}}%
\pgfpathlineto{\pgfqpoint{0.949806in}{0.626816in}}%
\pgfpathmoveto{\pgfqpoint{0.945265in}{0.629765in}}%
\pgfpathlineto{\pgfqpoint{0.945265in}{0.629765in}}%
\pgfpathlineto{\pgfqpoint{0.945265in}{0.632714in}}%
\pgfpathlineto{\pgfqpoint{0.949806in}{0.632714in}}%
\pgfpathlineto{\pgfqpoint{0.949806in}{0.629765in}}%
\pgfpathmoveto{\pgfqpoint{0.949806in}{0.629765in}}%
\pgfpathlineto{\pgfqpoint{0.949806in}{0.629765in}}%
\pgfpathlineto{\pgfqpoint{0.949806in}{0.632714in}}%
\pgfpathlineto{\pgfqpoint{0.954347in}{0.632714in}}%
\pgfpathlineto{\pgfqpoint{0.954347in}{0.629765in}}%
\pgfpathmoveto{\pgfqpoint{0.949806in}{0.632714in}}%
\pgfpathlineto{\pgfqpoint{0.949806in}{0.632714in}}%
\pgfpathlineto{\pgfqpoint{0.949806in}{0.635664in}}%
\pgfpathlineto{\pgfqpoint{0.954347in}{0.635664in}}%
\pgfpathlineto{\pgfqpoint{0.954347in}{0.632714in}}%
\pgfpathmoveto{\pgfqpoint{0.954347in}{0.632714in}}%
\pgfpathlineto{\pgfqpoint{0.954347in}{0.632714in}}%
\pgfpathlineto{\pgfqpoint{0.954347in}{0.635664in}}%
\pgfpathlineto{\pgfqpoint{0.958888in}{0.635664in}}%
\pgfpathlineto{\pgfqpoint{0.958888in}{0.632714in}}%
\pgfpathmoveto{\pgfqpoint{0.954347in}{0.635664in}}%
\pgfpathlineto{\pgfqpoint{0.954347in}{0.635664in}}%
\pgfpathlineto{\pgfqpoint{0.954347in}{0.638613in}}%
\pgfpathlineto{\pgfqpoint{0.958888in}{0.638613in}}%
\pgfpathlineto{\pgfqpoint{0.958888in}{0.635664in}}%
\pgfpathmoveto{\pgfqpoint{0.958888in}{0.635664in}}%
\pgfpathlineto{\pgfqpoint{0.958888in}{0.635664in}}%
\pgfpathlineto{\pgfqpoint{0.958888in}{0.638613in}}%
\pgfpathlineto{\pgfqpoint{0.963428in}{0.638613in}}%
\pgfpathlineto{\pgfqpoint{0.963428in}{0.635664in}}%
\pgfpathmoveto{\pgfqpoint{0.958888in}{0.638613in}}%
\pgfpathlineto{\pgfqpoint{0.958888in}{0.638613in}}%
\pgfpathlineto{\pgfqpoint{0.958888in}{0.641562in}}%
\pgfpathlineto{\pgfqpoint{0.963428in}{0.641562in}}%
\pgfpathlineto{\pgfqpoint{0.963428in}{0.638613in}}%
\pgfpathmoveto{\pgfqpoint{0.963428in}{0.638613in}}%
\pgfpathlineto{\pgfqpoint{0.963428in}{0.638613in}}%
\pgfpathlineto{\pgfqpoint{0.963428in}{0.641562in}}%
\pgfpathlineto{\pgfqpoint{0.967969in}{0.641562in}}%
\pgfpathlineto{\pgfqpoint{0.967969in}{0.638613in}}%
\pgfpathmoveto{\pgfqpoint{0.963428in}{0.641562in}}%
\pgfpathlineto{\pgfqpoint{0.963428in}{0.641562in}}%
\pgfpathlineto{\pgfqpoint{0.963428in}{0.644512in}}%
\pgfpathlineto{\pgfqpoint{0.967969in}{0.644512in}}%
\pgfpathlineto{\pgfqpoint{0.967969in}{0.641562in}}%
\pgfpathmoveto{\pgfqpoint{0.967969in}{0.641562in}}%
\pgfpathlineto{\pgfqpoint{0.967969in}{0.641562in}}%
\pgfpathlineto{\pgfqpoint{0.967969in}{0.644512in}}%
\pgfpathlineto{\pgfqpoint{0.972510in}{0.644512in}}%
\pgfpathlineto{\pgfqpoint{0.972510in}{0.641562in}}%
\pgfpathmoveto{\pgfqpoint{0.967969in}{0.644512in}}%
\pgfpathlineto{\pgfqpoint{0.967969in}{0.644512in}}%
\pgfpathlineto{\pgfqpoint{0.967969in}{0.647461in}}%
\pgfpathlineto{\pgfqpoint{0.972510in}{0.647461in}}%
\pgfpathlineto{\pgfqpoint{0.972510in}{0.644512in}}%
\pgfpathmoveto{\pgfqpoint{0.972510in}{0.644512in}}%
\pgfpathlineto{\pgfqpoint{0.972510in}{0.644512in}}%
\pgfpathlineto{\pgfqpoint{0.972510in}{0.647461in}}%
\pgfpathlineto{\pgfqpoint{0.977051in}{0.647461in}}%
\pgfpathlineto{\pgfqpoint{0.977051in}{0.644512in}}%
\pgfpathmoveto{\pgfqpoint{0.972510in}{0.647461in}}%
\pgfpathlineto{\pgfqpoint{0.972510in}{0.647461in}}%
\pgfpathlineto{\pgfqpoint{0.972510in}{0.650410in}}%
\pgfpathlineto{\pgfqpoint{0.977051in}{0.650410in}}%
\pgfpathlineto{\pgfqpoint{0.977051in}{0.647461in}}%
\pgfpathmoveto{\pgfqpoint{0.977051in}{0.647461in}}%
\pgfpathlineto{\pgfqpoint{0.977051in}{0.647461in}}%
\pgfpathlineto{\pgfqpoint{0.977051in}{0.650410in}}%
\pgfpathlineto{\pgfqpoint{0.981592in}{0.650410in}}%
\pgfpathlineto{\pgfqpoint{0.981592in}{0.647461in}}%
\pgfpathmoveto{\pgfqpoint{0.981592in}{0.647461in}}%
\pgfpathlineto{\pgfqpoint{0.981592in}{0.647461in}}%
\pgfpathlineto{\pgfqpoint{0.981592in}{0.650410in}}%
\pgfpathlineto{\pgfqpoint{0.986132in}{0.650410in}}%
\pgfpathlineto{\pgfqpoint{0.986132in}{0.647461in}}%
\pgfpathmoveto{\pgfqpoint{0.981592in}{0.650410in}}%
\pgfpathlineto{\pgfqpoint{0.981592in}{0.650410in}}%
\pgfpathlineto{\pgfqpoint{0.981592in}{0.653359in}}%
\pgfpathlineto{\pgfqpoint{0.986132in}{0.653359in}}%
\pgfpathlineto{\pgfqpoint{0.986132in}{0.650410in}}%
\pgfpathmoveto{\pgfqpoint{0.986132in}{0.650410in}}%
\pgfpathlineto{\pgfqpoint{0.986132in}{0.650410in}}%
\pgfpathlineto{\pgfqpoint{0.986132in}{0.653359in}}%
\pgfpathlineto{\pgfqpoint{0.990673in}{0.653359in}}%
\pgfpathlineto{\pgfqpoint{0.990673in}{0.650410in}}%
\pgfpathmoveto{\pgfqpoint{0.986132in}{0.653359in}}%
\pgfpathlineto{\pgfqpoint{0.986132in}{0.653359in}}%
\pgfpathlineto{\pgfqpoint{0.986132in}{0.656309in}}%
\pgfpathlineto{\pgfqpoint{0.990673in}{0.656309in}}%
\pgfpathlineto{\pgfqpoint{0.990673in}{0.653359in}}%
\pgfpathmoveto{\pgfqpoint{0.990673in}{0.653359in}}%
\pgfpathlineto{\pgfqpoint{0.990673in}{0.653359in}}%
\pgfpathlineto{\pgfqpoint{0.990673in}{0.656309in}}%
\pgfpathlineto{\pgfqpoint{0.995214in}{0.656309in}}%
\pgfpathlineto{\pgfqpoint{0.995214in}{0.653359in}}%
\pgfpathmoveto{\pgfqpoint{0.990673in}{0.656309in}}%
\pgfpathlineto{\pgfqpoint{0.990673in}{0.656309in}}%
\pgfpathlineto{\pgfqpoint{0.990673in}{0.659258in}}%
\pgfpathlineto{\pgfqpoint{0.995214in}{0.659258in}}%
\pgfpathlineto{\pgfqpoint{0.995214in}{0.656309in}}%
\pgfpathmoveto{\pgfqpoint{0.995214in}{0.656309in}}%
\pgfpathlineto{\pgfqpoint{0.995214in}{0.656309in}}%
\pgfpathlineto{\pgfqpoint{0.995214in}{0.659258in}}%
\pgfpathlineto{\pgfqpoint{0.999755in}{0.659258in}}%
\pgfpathlineto{\pgfqpoint{0.999755in}{0.656309in}}%
\pgfpathmoveto{\pgfqpoint{0.995214in}{0.659258in}}%
\pgfpathlineto{\pgfqpoint{0.995214in}{0.659258in}}%
\pgfpathlineto{\pgfqpoint{0.995214in}{0.662207in}}%
\pgfpathlineto{\pgfqpoint{0.999755in}{0.662207in}}%
\pgfpathlineto{\pgfqpoint{0.999755in}{0.659258in}}%
\pgfpathmoveto{\pgfqpoint{0.999755in}{0.659258in}}%
\pgfpathlineto{\pgfqpoint{0.999755in}{0.659258in}}%
\pgfpathlineto{\pgfqpoint{0.999755in}{0.662207in}}%
\pgfpathlineto{\pgfqpoint{1.004295in}{0.662207in}}%
\pgfpathlineto{\pgfqpoint{1.004295in}{0.659258in}}%
\pgfpathmoveto{\pgfqpoint{0.999755in}{0.662207in}}%
\pgfpathlineto{\pgfqpoint{0.999755in}{0.662207in}}%
\pgfpathlineto{\pgfqpoint{0.999755in}{0.665157in}}%
\pgfpathlineto{\pgfqpoint{1.004295in}{0.665157in}}%
\pgfpathlineto{\pgfqpoint{1.004295in}{0.662207in}}%
\pgfpathmoveto{\pgfqpoint{1.004295in}{0.662207in}}%
\pgfpathlineto{\pgfqpoint{1.004295in}{0.662207in}}%
\pgfpathlineto{\pgfqpoint{1.004295in}{0.665157in}}%
\pgfpathlineto{\pgfqpoint{1.008836in}{0.665157in}}%
\pgfpathlineto{\pgfqpoint{1.008836in}{0.662207in}}%
\pgfpathmoveto{\pgfqpoint{1.004295in}{0.665157in}}%
\pgfpathlineto{\pgfqpoint{1.004295in}{0.665157in}}%
\pgfpathlineto{\pgfqpoint{1.004295in}{0.668106in}}%
\pgfpathlineto{\pgfqpoint{1.008836in}{0.668106in}}%
\pgfpathlineto{\pgfqpoint{1.008836in}{0.665157in}}%
\pgfpathmoveto{\pgfqpoint{1.008836in}{0.665157in}}%
\pgfpathlineto{\pgfqpoint{1.008836in}{0.665157in}}%
\pgfpathlineto{\pgfqpoint{1.008836in}{0.668106in}}%
\pgfpathlineto{\pgfqpoint{1.013377in}{0.668106in}}%
\pgfpathlineto{\pgfqpoint{1.013377in}{0.665157in}}%
\pgfpathmoveto{\pgfqpoint{1.008836in}{0.668106in}}%
\pgfpathlineto{\pgfqpoint{1.008836in}{0.668106in}}%
\pgfpathlineto{\pgfqpoint{1.008836in}{0.671055in}}%
\pgfpathlineto{\pgfqpoint{1.013377in}{0.671055in}}%
\pgfpathlineto{\pgfqpoint{1.013377in}{0.668106in}}%
\pgfpathmoveto{\pgfqpoint{1.013377in}{0.668106in}}%
\pgfpathlineto{\pgfqpoint{1.013377in}{0.668106in}}%
\pgfpathlineto{\pgfqpoint{1.013377in}{0.671055in}}%
\pgfpathlineto{\pgfqpoint{1.017918in}{0.671055in}}%
\pgfpathlineto{\pgfqpoint{1.017918in}{0.668106in}}%
\pgfpathmoveto{\pgfqpoint{1.013377in}{0.671055in}}%
\pgfpathlineto{\pgfqpoint{1.013377in}{0.671055in}}%
\pgfpathlineto{\pgfqpoint{1.013377in}{0.674004in}}%
\pgfpathlineto{\pgfqpoint{1.017918in}{0.674004in}}%
\pgfpathlineto{\pgfqpoint{1.017918in}{0.671055in}}%
\pgfpathmoveto{\pgfqpoint{1.017918in}{0.671055in}}%
\pgfpathlineto{\pgfqpoint{1.017918in}{0.671055in}}%
\pgfpathlineto{\pgfqpoint{1.017918in}{0.674004in}}%
\pgfpathlineto{\pgfqpoint{1.022459in}{0.674004in}}%
\pgfpathlineto{\pgfqpoint{1.022459in}{0.671055in}}%
\pgfpathmoveto{\pgfqpoint{1.017918in}{0.674004in}}%
\pgfpathlineto{\pgfqpoint{1.017918in}{0.674004in}}%
\pgfpathlineto{\pgfqpoint{1.017918in}{0.676954in}}%
\pgfpathlineto{\pgfqpoint{1.022459in}{0.676954in}}%
\pgfpathlineto{\pgfqpoint{1.022459in}{0.674004in}}%
\pgfpathmoveto{\pgfqpoint{1.022459in}{0.674004in}}%
\pgfpathlineto{\pgfqpoint{1.022459in}{0.674004in}}%
\pgfpathlineto{\pgfqpoint{1.022459in}{0.676954in}}%
\pgfpathlineto{\pgfqpoint{1.026999in}{0.676954in}}%
\pgfpathlineto{\pgfqpoint{1.026999in}{0.674004in}}%
\pgfpathmoveto{\pgfqpoint{1.022459in}{0.676954in}}%
\pgfpathlineto{\pgfqpoint{1.022459in}{0.676954in}}%
\pgfpathlineto{\pgfqpoint{1.022459in}{0.679903in}}%
\pgfpathlineto{\pgfqpoint{1.026999in}{0.679903in}}%
\pgfpathlineto{\pgfqpoint{1.026999in}{0.676954in}}%
\pgfpathmoveto{\pgfqpoint{1.026999in}{0.676954in}}%
\pgfpathlineto{\pgfqpoint{1.026999in}{0.676954in}}%
\pgfpathlineto{\pgfqpoint{1.026999in}{0.679903in}}%
\pgfpathlineto{\pgfqpoint{1.031540in}{0.679903in}}%
\pgfpathlineto{\pgfqpoint{1.031540in}{0.676954in}}%
\pgfpathmoveto{\pgfqpoint{1.026999in}{0.679903in}}%
\pgfpathlineto{\pgfqpoint{1.026999in}{0.679903in}}%
\pgfpathlineto{\pgfqpoint{1.026999in}{0.682852in}}%
\pgfpathlineto{\pgfqpoint{1.031540in}{0.682852in}}%
\pgfpathlineto{\pgfqpoint{1.031540in}{0.679903in}}%
\pgfpathmoveto{\pgfqpoint{1.031540in}{0.679903in}}%
\pgfpathlineto{\pgfqpoint{1.031540in}{0.679903in}}%
\pgfpathlineto{\pgfqpoint{1.031540in}{0.682852in}}%
\pgfpathlineto{\pgfqpoint{1.036081in}{0.682852in}}%
\pgfpathlineto{\pgfqpoint{1.036081in}{0.679903in}}%
\pgfpathmoveto{\pgfqpoint{1.031540in}{0.682852in}}%
\pgfpathlineto{\pgfqpoint{1.031540in}{0.682852in}}%
\pgfpathlineto{\pgfqpoint{1.031540in}{0.685802in}}%
\pgfpathlineto{\pgfqpoint{1.036081in}{0.685802in}}%
\pgfpathlineto{\pgfqpoint{1.036081in}{0.682852in}}%
\pgfpathmoveto{\pgfqpoint{1.036081in}{0.682852in}}%
\pgfpathlineto{\pgfqpoint{1.036081in}{0.682852in}}%
\pgfpathlineto{\pgfqpoint{1.036081in}{0.685802in}}%
\pgfpathlineto{\pgfqpoint{1.040622in}{0.685802in}}%
\pgfpathlineto{\pgfqpoint{1.040622in}{0.682852in}}%
\pgfpathmoveto{\pgfqpoint{1.036081in}{0.685802in}}%
\pgfpathlineto{\pgfqpoint{1.036081in}{0.685802in}}%
\pgfpathlineto{\pgfqpoint{1.036081in}{0.688751in}}%
\pgfpathlineto{\pgfqpoint{1.040622in}{0.688751in}}%
\pgfpathlineto{\pgfqpoint{1.040622in}{0.685802in}}%
\pgfpathmoveto{\pgfqpoint{1.040622in}{0.685802in}}%
\pgfpathlineto{\pgfqpoint{1.040622in}{0.685802in}}%
\pgfpathlineto{\pgfqpoint{1.040622in}{0.688751in}}%
\pgfpathlineto{\pgfqpoint{1.045163in}{0.688751in}}%
\pgfpathlineto{\pgfqpoint{1.045163in}{0.685802in}}%
\pgfpathmoveto{\pgfqpoint{1.040622in}{0.688751in}}%
\pgfpathlineto{\pgfqpoint{1.040622in}{0.688751in}}%
\pgfpathlineto{\pgfqpoint{1.040622in}{0.691700in}}%
\pgfpathlineto{\pgfqpoint{1.045163in}{0.691700in}}%
\pgfpathlineto{\pgfqpoint{1.045163in}{0.688751in}}%
\pgfpathmoveto{\pgfqpoint{1.045163in}{0.688751in}}%
\pgfpathlineto{\pgfqpoint{1.045163in}{0.688751in}}%
\pgfpathlineto{\pgfqpoint{1.045163in}{0.691700in}}%
\pgfpathlineto{\pgfqpoint{1.049704in}{0.691700in}}%
\pgfpathlineto{\pgfqpoint{1.049704in}{0.688751in}}%
\pgfpathmoveto{\pgfqpoint{1.045163in}{0.691700in}}%
\pgfpathlineto{\pgfqpoint{1.045163in}{0.691700in}}%
\pgfpathlineto{\pgfqpoint{1.045163in}{0.694649in}}%
\pgfpathlineto{\pgfqpoint{1.049704in}{0.694649in}}%
\pgfpathlineto{\pgfqpoint{1.049704in}{0.691700in}}%
\pgfpathmoveto{\pgfqpoint{1.049704in}{0.691700in}}%
\pgfpathlineto{\pgfqpoint{1.049704in}{0.691700in}}%
\pgfpathlineto{\pgfqpoint{1.049704in}{0.694649in}}%
\pgfpathlineto{\pgfqpoint{1.054245in}{0.694649in}}%
\pgfpathlineto{\pgfqpoint{1.054245in}{0.691700in}}%
\pgfpathmoveto{\pgfqpoint{1.049704in}{0.694649in}}%
\pgfpathlineto{\pgfqpoint{1.049704in}{0.694649in}}%
\pgfpathlineto{\pgfqpoint{1.049704in}{0.697599in}}%
\pgfpathlineto{\pgfqpoint{1.054245in}{0.697599in}}%
\pgfpathlineto{\pgfqpoint{1.054245in}{0.694649in}}%
\pgfpathmoveto{\pgfqpoint{1.054245in}{0.694649in}}%
\pgfpathlineto{\pgfqpoint{1.054245in}{0.694649in}}%
\pgfpathlineto{\pgfqpoint{1.054245in}{0.697599in}}%
\pgfpathlineto{\pgfqpoint{1.058786in}{0.697599in}}%
\pgfpathlineto{\pgfqpoint{1.058786in}{0.694649in}}%
\pgfpathmoveto{\pgfqpoint{1.054245in}{0.697599in}}%
\pgfpathlineto{\pgfqpoint{1.054245in}{0.697599in}}%
\pgfpathlineto{\pgfqpoint{1.054245in}{0.700548in}}%
\pgfpathlineto{\pgfqpoint{1.058786in}{0.700548in}}%
\pgfpathlineto{\pgfqpoint{1.058786in}{0.697599in}}%
\pgfpathmoveto{\pgfqpoint{1.058786in}{0.697599in}}%
\pgfpathlineto{\pgfqpoint{1.058786in}{0.697599in}}%
\pgfpathlineto{\pgfqpoint{1.058786in}{0.700548in}}%
\pgfpathlineto{\pgfqpoint{1.063328in}{0.700548in}}%
\pgfpathlineto{\pgfqpoint{1.063328in}{0.697599in}}%
\pgfpathmoveto{\pgfqpoint{1.058786in}{0.700548in}}%
\pgfpathlineto{\pgfqpoint{1.058786in}{0.700548in}}%
\pgfpathlineto{\pgfqpoint{1.058786in}{0.703497in}}%
\pgfpathlineto{\pgfqpoint{1.063328in}{0.703497in}}%
\pgfpathlineto{\pgfqpoint{1.063328in}{0.700548in}}%
\pgfpathmoveto{\pgfqpoint{1.063328in}{0.700548in}}%
\pgfpathlineto{\pgfqpoint{1.063328in}{0.700548in}}%
\pgfpathlineto{\pgfqpoint{1.063328in}{0.703497in}}%
\pgfpathlineto{\pgfqpoint{1.067869in}{0.703497in}}%
\pgfpathlineto{\pgfqpoint{1.067869in}{0.700548in}}%
\pgfpathmoveto{\pgfqpoint{1.063328in}{0.703497in}}%
\pgfpathlineto{\pgfqpoint{1.063328in}{0.703497in}}%
\pgfpathlineto{\pgfqpoint{1.063328in}{0.706446in}}%
\pgfpathlineto{\pgfqpoint{1.067869in}{0.706446in}}%
\pgfpathlineto{\pgfqpoint{1.067869in}{0.703497in}}%
\pgfpathmoveto{\pgfqpoint{1.067869in}{0.703497in}}%
\pgfpathlineto{\pgfqpoint{1.067869in}{0.703497in}}%
\pgfpathlineto{\pgfqpoint{1.067869in}{0.706446in}}%
\pgfpathlineto{\pgfqpoint{1.072410in}{0.706446in}}%
\pgfpathlineto{\pgfqpoint{1.072410in}{0.703497in}}%
\pgfpathmoveto{\pgfqpoint{1.067869in}{0.706446in}}%
\pgfpathlineto{\pgfqpoint{1.067869in}{0.706446in}}%
\pgfpathlineto{\pgfqpoint{1.067869in}{0.709396in}}%
\pgfpathlineto{\pgfqpoint{1.072410in}{0.709396in}}%
\pgfpathlineto{\pgfqpoint{1.072410in}{0.706446in}}%
\pgfpathmoveto{\pgfqpoint{1.072410in}{0.706446in}}%
\pgfpathlineto{\pgfqpoint{1.072410in}{0.706446in}}%
\pgfpathlineto{\pgfqpoint{1.072410in}{0.709396in}}%
\pgfpathlineto{\pgfqpoint{1.076951in}{0.709396in}}%
\pgfpathlineto{\pgfqpoint{1.076951in}{0.706446in}}%
\pgfpathmoveto{\pgfqpoint{1.072410in}{0.709396in}}%
\pgfpathlineto{\pgfqpoint{1.072410in}{0.709396in}}%
\pgfpathlineto{\pgfqpoint{1.072410in}{0.712345in}}%
\pgfpathlineto{\pgfqpoint{1.076951in}{0.712345in}}%
\pgfpathlineto{\pgfqpoint{1.076951in}{0.709396in}}%
\pgfpathmoveto{\pgfqpoint{1.076951in}{0.709396in}}%
\pgfpathlineto{\pgfqpoint{1.076951in}{0.709396in}}%
\pgfpathlineto{\pgfqpoint{1.076951in}{0.712345in}}%
\pgfpathlineto{\pgfqpoint{1.081492in}{0.712345in}}%
\pgfpathlineto{\pgfqpoint{1.081492in}{0.709396in}}%
\pgfpathmoveto{\pgfqpoint{1.076951in}{0.712345in}}%
\pgfpathlineto{\pgfqpoint{1.076951in}{0.712345in}}%
\pgfpathlineto{\pgfqpoint{1.076951in}{0.715294in}}%
\pgfpathlineto{\pgfqpoint{1.081492in}{0.715294in}}%
\pgfpathlineto{\pgfqpoint{1.081492in}{0.712345in}}%
\pgfpathmoveto{\pgfqpoint{1.081492in}{0.712345in}}%
\pgfpathlineto{\pgfqpoint{1.081492in}{0.712345in}}%
\pgfpathlineto{\pgfqpoint{1.081492in}{0.715294in}}%
\pgfpathlineto{\pgfqpoint{1.086033in}{0.715294in}}%
\pgfpathlineto{\pgfqpoint{1.086033in}{0.712345in}}%
\pgfpathmoveto{\pgfqpoint{1.081492in}{0.715294in}}%
\pgfpathlineto{\pgfqpoint{1.081492in}{0.715294in}}%
\pgfpathlineto{\pgfqpoint{1.081492in}{0.718243in}}%
\pgfpathlineto{\pgfqpoint{1.086033in}{0.718243in}}%
\pgfpathlineto{\pgfqpoint{1.086033in}{0.715294in}}%
\pgfpathmoveto{\pgfqpoint{1.086033in}{0.715294in}}%
\pgfpathlineto{\pgfqpoint{1.086033in}{0.715294in}}%
\pgfpathlineto{\pgfqpoint{1.086033in}{0.718243in}}%
\pgfpathlineto{\pgfqpoint{1.090575in}{0.718243in}}%
\pgfpathlineto{\pgfqpoint{1.090575in}{0.715294in}}%
\pgfpathmoveto{\pgfqpoint{1.086033in}{0.718243in}}%
\pgfpathlineto{\pgfqpoint{1.086033in}{0.718243in}}%
\pgfpathlineto{\pgfqpoint{1.086033in}{0.721193in}}%
\pgfpathlineto{\pgfqpoint{1.090575in}{0.721193in}}%
\pgfpathlineto{\pgfqpoint{1.090575in}{0.718243in}}%
\pgfpathmoveto{\pgfqpoint{1.090575in}{0.718243in}}%
\pgfpathlineto{\pgfqpoint{1.090575in}{0.718243in}}%
\pgfpathlineto{\pgfqpoint{1.090575in}{0.721193in}}%
\pgfpathlineto{\pgfqpoint{1.095116in}{0.721193in}}%
\pgfpathlineto{\pgfqpoint{1.095116in}{0.718243in}}%
\pgfpathmoveto{\pgfqpoint{1.090575in}{0.721193in}}%
\pgfpathlineto{\pgfqpoint{1.090575in}{0.721193in}}%
\pgfpathlineto{\pgfqpoint{1.090575in}{0.724142in}}%
\pgfpathlineto{\pgfqpoint{1.095116in}{0.724142in}}%
\pgfpathlineto{\pgfqpoint{1.095116in}{0.721193in}}%
\pgfpathmoveto{\pgfqpoint{1.095116in}{0.721193in}}%
\pgfpathlineto{\pgfqpoint{1.095116in}{0.721193in}}%
\pgfpathlineto{\pgfqpoint{1.095116in}{0.724142in}}%
\pgfpathlineto{\pgfqpoint{1.099657in}{0.724142in}}%
\pgfpathlineto{\pgfqpoint{1.099657in}{0.721193in}}%
\pgfpathmoveto{\pgfqpoint{1.095116in}{0.724142in}}%
\pgfpathlineto{\pgfqpoint{1.095116in}{0.724142in}}%
\pgfpathlineto{\pgfqpoint{1.095116in}{0.727091in}}%
\pgfpathlineto{\pgfqpoint{1.099657in}{0.727091in}}%
\pgfpathlineto{\pgfqpoint{1.099657in}{0.724142in}}%
\pgfpathmoveto{\pgfqpoint{1.099657in}{0.724142in}}%
\pgfpathlineto{\pgfqpoint{1.099657in}{0.724142in}}%
\pgfpathlineto{\pgfqpoint{1.099657in}{0.727091in}}%
\pgfpathlineto{\pgfqpoint{1.104198in}{0.727091in}}%
\pgfpathlineto{\pgfqpoint{1.104198in}{0.724142in}}%
\pgfpathmoveto{\pgfqpoint{1.099657in}{0.727091in}}%
\pgfpathlineto{\pgfqpoint{1.099657in}{0.727091in}}%
\pgfpathlineto{\pgfqpoint{1.099657in}{0.730040in}}%
\pgfpathlineto{\pgfqpoint{1.104198in}{0.730040in}}%
\pgfpathlineto{\pgfqpoint{1.104198in}{0.727091in}}%
\pgfpathmoveto{\pgfqpoint{1.104198in}{0.727091in}}%
\pgfpathlineto{\pgfqpoint{1.104198in}{0.727091in}}%
\pgfpathlineto{\pgfqpoint{1.104198in}{0.730040in}}%
\pgfpathlineto{\pgfqpoint{1.108739in}{0.730040in}}%
\pgfpathlineto{\pgfqpoint{1.108739in}{0.727091in}}%
\pgfpathmoveto{\pgfqpoint{1.104198in}{0.730040in}}%
\pgfpathlineto{\pgfqpoint{1.104198in}{0.730040in}}%
\pgfpathlineto{\pgfqpoint{1.104198in}{0.732990in}}%
\pgfpathlineto{\pgfqpoint{1.108739in}{0.732990in}}%
\pgfpathlineto{\pgfqpoint{1.108739in}{0.730040in}}%
\pgfpathmoveto{\pgfqpoint{1.108739in}{0.730040in}}%
\pgfpathlineto{\pgfqpoint{1.108739in}{0.730040in}}%
\pgfpathlineto{\pgfqpoint{1.108739in}{0.732990in}}%
\pgfpathlineto{\pgfqpoint{1.113280in}{0.732990in}}%
\pgfpathlineto{\pgfqpoint{1.113280in}{0.730040in}}%
\pgfpathmoveto{\pgfqpoint{1.108739in}{0.732990in}}%
\pgfpathlineto{\pgfqpoint{1.108739in}{0.732990in}}%
\pgfpathlineto{\pgfqpoint{1.108739in}{0.735939in}}%
\pgfpathlineto{\pgfqpoint{1.113280in}{0.735939in}}%
\pgfpathlineto{\pgfqpoint{1.113280in}{0.732990in}}%
\pgfpathmoveto{\pgfqpoint{1.113280in}{0.732990in}}%
\pgfpathlineto{\pgfqpoint{1.113280in}{0.732990in}}%
\pgfpathlineto{\pgfqpoint{1.113280in}{0.735939in}}%
\pgfpathlineto{\pgfqpoint{1.117822in}{0.735939in}}%
\pgfpathlineto{\pgfqpoint{1.117822in}{0.732990in}}%
\pgfpathmoveto{\pgfqpoint{1.113280in}{0.735939in}}%
\pgfpathlineto{\pgfqpoint{1.113280in}{0.735939in}}%
\pgfpathlineto{\pgfqpoint{1.113280in}{0.738888in}}%
\pgfpathlineto{\pgfqpoint{1.117822in}{0.738888in}}%
\pgfpathlineto{\pgfqpoint{1.117822in}{0.735939in}}%
\pgfpathmoveto{\pgfqpoint{1.117822in}{0.735939in}}%
\pgfpathlineto{\pgfqpoint{1.117822in}{0.735939in}}%
\pgfpathlineto{\pgfqpoint{1.117822in}{0.738888in}}%
\pgfpathlineto{\pgfqpoint{1.122363in}{0.738888in}}%
\pgfpathlineto{\pgfqpoint{1.122363in}{0.735939in}}%
\pgfpathmoveto{\pgfqpoint{1.117822in}{0.738888in}}%
\pgfpathlineto{\pgfqpoint{1.117822in}{0.738888in}}%
\pgfpathlineto{\pgfqpoint{1.117822in}{0.741837in}}%
\pgfpathlineto{\pgfqpoint{1.122363in}{0.741837in}}%
\pgfpathlineto{\pgfqpoint{1.122363in}{0.738888in}}%
\pgfpathmoveto{\pgfqpoint{1.122363in}{0.738888in}}%
\pgfpathlineto{\pgfqpoint{1.122363in}{0.738888in}}%
\pgfpathlineto{\pgfqpoint{1.122363in}{0.741837in}}%
\pgfpathlineto{\pgfqpoint{1.126904in}{0.741837in}}%
\pgfpathlineto{\pgfqpoint{1.126904in}{0.738888in}}%
\pgfpathmoveto{\pgfqpoint{1.122363in}{0.741837in}}%
\pgfpathlineto{\pgfqpoint{1.122363in}{0.741837in}}%
\pgfpathlineto{\pgfqpoint{1.122363in}{0.744787in}}%
\pgfpathlineto{\pgfqpoint{1.126904in}{0.744787in}}%
\pgfpathlineto{\pgfqpoint{1.126904in}{0.741837in}}%
\pgfpathmoveto{\pgfqpoint{1.126904in}{0.741837in}}%
\pgfpathlineto{\pgfqpoint{1.126904in}{0.741837in}}%
\pgfpathlineto{\pgfqpoint{1.126904in}{0.744787in}}%
\pgfpathlineto{\pgfqpoint{1.131445in}{0.744787in}}%
\pgfpathlineto{\pgfqpoint{1.131445in}{0.741837in}}%
\pgfpathmoveto{\pgfqpoint{1.126904in}{0.744787in}}%
\pgfpathlineto{\pgfqpoint{1.126904in}{0.744787in}}%
\pgfpathlineto{\pgfqpoint{1.126904in}{0.747736in}}%
\pgfpathlineto{\pgfqpoint{1.131445in}{0.747736in}}%
\pgfpathlineto{\pgfqpoint{1.131445in}{0.744787in}}%
\pgfpathmoveto{\pgfqpoint{1.131445in}{0.744787in}}%
\pgfpathlineto{\pgfqpoint{1.131445in}{0.744787in}}%
\pgfpathlineto{\pgfqpoint{1.131445in}{0.747736in}}%
\pgfpathlineto{\pgfqpoint{1.135986in}{0.747736in}}%
\pgfpathlineto{\pgfqpoint{1.135986in}{0.744787in}}%
\pgfpathmoveto{\pgfqpoint{1.131445in}{0.747736in}}%
\pgfpathlineto{\pgfqpoint{1.131445in}{0.747736in}}%
\pgfpathlineto{\pgfqpoint{1.131445in}{0.750685in}}%
\pgfpathlineto{\pgfqpoint{1.135986in}{0.750685in}}%
\pgfpathlineto{\pgfqpoint{1.135986in}{0.747736in}}%
\pgfpathmoveto{\pgfqpoint{1.135986in}{0.747736in}}%
\pgfpathlineto{\pgfqpoint{1.135986in}{0.747736in}}%
\pgfpathlineto{\pgfqpoint{1.135986in}{0.750685in}}%
\pgfpathlineto{\pgfqpoint{1.140528in}{0.750685in}}%
\pgfpathlineto{\pgfqpoint{1.140528in}{0.747736in}}%
\pgfpathmoveto{\pgfqpoint{1.135986in}{0.750685in}}%
\pgfpathlineto{\pgfqpoint{1.135986in}{0.750685in}}%
\pgfpathlineto{\pgfqpoint{1.135986in}{0.753634in}}%
\pgfpathlineto{\pgfqpoint{1.140528in}{0.753634in}}%
\pgfpathlineto{\pgfqpoint{1.140528in}{0.750685in}}%
\pgfpathmoveto{\pgfqpoint{1.140528in}{0.750685in}}%
\pgfpathlineto{\pgfqpoint{1.140528in}{0.750685in}}%
\pgfpathlineto{\pgfqpoint{1.140528in}{0.753634in}}%
\pgfpathlineto{\pgfqpoint{1.145069in}{0.753634in}}%
\pgfpathlineto{\pgfqpoint{1.145069in}{0.750685in}}%
\pgfpathmoveto{\pgfqpoint{1.140528in}{0.753634in}}%
\pgfpathlineto{\pgfqpoint{1.140528in}{0.753634in}}%
\pgfpathlineto{\pgfqpoint{1.140528in}{0.756584in}}%
\pgfpathlineto{\pgfqpoint{1.145069in}{0.756584in}}%
\pgfpathlineto{\pgfqpoint{1.145069in}{0.753634in}}%
\pgfpathmoveto{\pgfqpoint{1.145069in}{0.753634in}}%
\pgfpathlineto{\pgfqpoint{1.145069in}{0.753634in}}%
\pgfpathlineto{\pgfqpoint{1.145069in}{0.756584in}}%
\pgfpathlineto{\pgfqpoint{1.149610in}{0.756584in}}%
\pgfpathlineto{\pgfqpoint{1.149610in}{0.753634in}}%
\pgfpathmoveto{\pgfqpoint{1.145069in}{0.756584in}}%
\pgfpathlineto{\pgfqpoint{1.145069in}{0.756584in}}%
\pgfpathlineto{\pgfqpoint{1.145069in}{0.759533in}}%
\pgfpathlineto{\pgfqpoint{1.149610in}{0.759533in}}%
\pgfpathlineto{\pgfqpoint{1.149610in}{0.756584in}}%
\pgfpathmoveto{\pgfqpoint{1.149610in}{0.756584in}}%
\pgfpathlineto{\pgfqpoint{1.149610in}{0.756584in}}%
\pgfpathlineto{\pgfqpoint{1.149610in}{0.759533in}}%
\pgfpathlineto{\pgfqpoint{1.154151in}{0.759533in}}%
\pgfpathlineto{\pgfqpoint{1.154151in}{0.756584in}}%
\pgfpathmoveto{\pgfqpoint{1.149610in}{0.759533in}}%
\pgfpathlineto{\pgfqpoint{1.149610in}{0.759533in}}%
\pgfpathlineto{\pgfqpoint{1.149610in}{0.762482in}}%
\pgfpathlineto{\pgfqpoint{1.154151in}{0.762482in}}%
\pgfpathlineto{\pgfqpoint{1.154151in}{0.759533in}}%
\pgfpathmoveto{\pgfqpoint{1.154151in}{0.759533in}}%
\pgfpathlineto{\pgfqpoint{1.154151in}{0.759533in}}%
\pgfpathlineto{\pgfqpoint{1.154151in}{0.762482in}}%
\pgfpathlineto{\pgfqpoint{1.158692in}{0.762482in}}%
\pgfpathlineto{\pgfqpoint{1.158692in}{0.759533in}}%
\pgfpathmoveto{\pgfqpoint{1.154151in}{0.762482in}}%
\pgfpathlineto{\pgfqpoint{1.154151in}{0.762482in}}%
\pgfpathlineto{\pgfqpoint{1.154151in}{0.765431in}}%
\pgfpathlineto{\pgfqpoint{1.158692in}{0.765431in}}%
\pgfpathlineto{\pgfqpoint{1.158692in}{0.762482in}}%
\pgfpathmoveto{\pgfqpoint{1.158692in}{0.762482in}}%
\pgfpathlineto{\pgfqpoint{1.158692in}{0.762482in}}%
\pgfpathlineto{\pgfqpoint{1.158692in}{0.765431in}}%
\pgfpathlineto{\pgfqpoint{1.163233in}{0.765431in}}%
\pgfpathlineto{\pgfqpoint{1.163233in}{0.762482in}}%
\pgfpathmoveto{\pgfqpoint{1.158692in}{0.765431in}}%
\pgfpathlineto{\pgfqpoint{1.158692in}{0.765431in}}%
\pgfpathlineto{\pgfqpoint{1.158692in}{0.768381in}}%
\pgfpathlineto{\pgfqpoint{1.163233in}{0.768381in}}%
\pgfpathlineto{\pgfqpoint{1.163233in}{0.765431in}}%
\pgfpathmoveto{\pgfqpoint{1.163233in}{0.765431in}}%
\pgfpathlineto{\pgfqpoint{1.163233in}{0.765431in}}%
\pgfpathlineto{\pgfqpoint{1.163233in}{0.768381in}}%
\pgfpathlineto{\pgfqpoint{1.167775in}{0.768381in}}%
\pgfpathlineto{\pgfqpoint{1.167775in}{0.765431in}}%
\pgfpathmoveto{\pgfqpoint{1.163233in}{0.768381in}}%
\pgfpathlineto{\pgfqpoint{1.163233in}{0.768381in}}%
\pgfpathlineto{\pgfqpoint{1.163233in}{0.771330in}}%
\pgfpathlineto{\pgfqpoint{1.167775in}{0.771330in}}%
\pgfpathlineto{\pgfqpoint{1.167775in}{0.768381in}}%
\pgfpathmoveto{\pgfqpoint{1.167775in}{0.768381in}}%
\pgfpathlineto{\pgfqpoint{1.167775in}{0.768381in}}%
\pgfpathlineto{\pgfqpoint{1.167775in}{0.771330in}}%
\pgfpathlineto{\pgfqpoint{1.172316in}{0.771330in}}%
\pgfpathlineto{\pgfqpoint{1.172316in}{0.768381in}}%
\pgfpathmoveto{\pgfqpoint{1.167775in}{0.771330in}}%
\pgfpathlineto{\pgfqpoint{1.167775in}{0.771330in}}%
\pgfpathlineto{\pgfqpoint{1.167775in}{0.774279in}}%
\pgfpathlineto{\pgfqpoint{1.172316in}{0.774279in}}%
\pgfpathlineto{\pgfqpoint{1.172316in}{0.771330in}}%
\pgfpathmoveto{\pgfqpoint{1.172316in}{0.771330in}}%
\pgfpathlineto{\pgfqpoint{1.172316in}{0.771330in}}%
\pgfpathlineto{\pgfqpoint{1.172316in}{0.774279in}}%
\pgfpathlineto{\pgfqpoint{1.176857in}{0.774279in}}%
\pgfpathlineto{\pgfqpoint{1.176857in}{0.771330in}}%
\pgfpathmoveto{\pgfqpoint{1.172316in}{0.774279in}}%
\pgfpathlineto{\pgfqpoint{1.172316in}{0.774279in}}%
\pgfpathlineto{\pgfqpoint{1.172316in}{0.777228in}}%
\pgfpathlineto{\pgfqpoint{1.176857in}{0.777228in}}%
\pgfpathlineto{\pgfqpoint{1.176857in}{0.774279in}}%
\pgfpathmoveto{\pgfqpoint{1.176857in}{0.774279in}}%
\pgfpathlineto{\pgfqpoint{1.176857in}{0.774279in}}%
\pgfpathlineto{\pgfqpoint{1.176857in}{0.777228in}}%
\pgfpathlineto{\pgfqpoint{1.181398in}{0.777228in}}%
\pgfpathlineto{\pgfqpoint{1.181398in}{0.774279in}}%
\pgfpathmoveto{\pgfqpoint{1.176857in}{0.777228in}}%
\pgfpathlineto{\pgfqpoint{1.176857in}{0.777228in}}%
\pgfpathlineto{\pgfqpoint{1.176857in}{0.780178in}}%
\pgfpathlineto{\pgfqpoint{1.181398in}{0.780178in}}%
\pgfpathlineto{\pgfqpoint{1.181398in}{0.777228in}}%
\pgfpathmoveto{\pgfqpoint{1.181398in}{0.777228in}}%
\pgfpathlineto{\pgfqpoint{1.181398in}{0.777228in}}%
\pgfpathlineto{\pgfqpoint{1.181398in}{0.780178in}}%
\pgfpathlineto{\pgfqpoint{1.185939in}{0.780178in}}%
\pgfpathlineto{\pgfqpoint{1.185939in}{0.777228in}}%
\pgfpathmoveto{\pgfqpoint{1.181398in}{0.780178in}}%
\pgfpathlineto{\pgfqpoint{1.181398in}{0.780178in}}%
\pgfpathlineto{\pgfqpoint{1.181398in}{0.783127in}}%
\pgfpathlineto{\pgfqpoint{1.185939in}{0.783127in}}%
\pgfpathlineto{\pgfqpoint{1.185939in}{0.780178in}}%
\pgfpathmoveto{\pgfqpoint{1.185939in}{0.780178in}}%
\pgfpathlineto{\pgfqpoint{1.185939in}{0.780178in}}%
\pgfpathlineto{\pgfqpoint{1.185939in}{0.783127in}}%
\pgfpathlineto{\pgfqpoint{1.190480in}{0.783127in}}%
\pgfpathlineto{\pgfqpoint{1.190480in}{0.780178in}}%
\pgfpathmoveto{\pgfqpoint{1.185939in}{0.783127in}}%
\pgfpathlineto{\pgfqpoint{1.185939in}{0.783127in}}%
\pgfpathlineto{\pgfqpoint{1.185939in}{0.786076in}}%
\pgfpathlineto{\pgfqpoint{1.190480in}{0.786076in}}%
\pgfpathlineto{\pgfqpoint{1.190480in}{0.783127in}}%
\pgfpathmoveto{\pgfqpoint{1.190480in}{0.783127in}}%
\pgfpathlineto{\pgfqpoint{1.190480in}{0.783127in}}%
\pgfpathlineto{\pgfqpoint{1.190480in}{0.786076in}}%
\pgfpathlineto{\pgfqpoint{1.195021in}{0.786076in}}%
\pgfpathlineto{\pgfqpoint{1.195021in}{0.783127in}}%
\pgfpathmoveto{\pgfqpoint{1.190480in}{0.786076in}}%
\pgfpathlineto{\pgfqpoint{1.190480in}{0.786076in}}%
\pgfpathlineto{\pgfqpoint{1.190480in}{0.789025in}}%
\pgfpathlineto{\pgfqpoint{1.195021in}{0.789025in}}%
\pgfpathlineto{\pgfqpoint{1.195021in}{0.786076in}}%
\pgfpathmoveto{\pgfqpoint{1.195021in}{0.786076in}}%
\pgfpathlineto{\pgfqpoint{1.195021in}{0.786076in}}%
\pgfpathlineto{\pgfqpoint{1.195021in}{0.789025in}}%
\pgfpathlineto{\pgfqpoint{1.199562in}{0.789025in}}%
\pgfpathlineto{\pgfqpoint{1.199562in}{0.786076in}}%
\pgfpathmoveto{\pgfqpoint{1.195021in}{0.789025in}}%
\pgfpathlineto{\pgfqpoint{1.195021in}{0.789025in}}%
\pgfpathlineto{\pgfqpoint{1.195021in}{0.791974in}}%
\pgfpathlineto{\pgfqpoint{1.199562in}{0.791974in}}%
\pgfpathlineto{\pgfqpoint{1.199562in}{0.789025in}}%
\pgfpathmoveto{\pgfqpoint{1.199562in}{0.789025in}}%
\pgfpathlineto{\pgfqpoint{1.199562in}{0.789025in}}%
\pgfpathlineto{\pgfqpoint{1.199562in}{0.791974in}}%
\pgfpathlineto{\pgfqpoint{1.204103in}{0.791974in}}%
\pgfpathlineto{\pgfqpoint{1.204103in}{0.789025in}}%
\pgfpathmoveto{\pgfqpoint{1.199562in}{0.791974in}}%
\pgfpathlineto{\pgfqpoint{1.199562in}{0.791974in}}%
\pgfpathlineto{\pgfqpoint{1.199562in}{0.794924in}}%
\pgfpathlineto{\pgfqpoint{1.204103in}{0.794924in}}%
\pgfpathlineto{\pgfqpoint{1.204103in}{0.791974in}}%
\pgfpathmoveto{\pgfqpoint{1.204103in}{0.791974in}}%
\pgfpathlineto{\pgfqpoint{1.204103in}{0.791974in}}%
\pgfpathlineto{\pgfqpoint{1.204103in}{0.794924in}}%
\pgfpathlineto{\pgfqpoint{1.208644in}{0.794924in}}%
\pgfpathlineto{\pgfqpoint{1.208644in}{0.791974in}}%
\pgfpathmoveto{\pgfqpoint{1.204103in}{0.794924in}}%
\pgfpathlineto{\pgfqpoint{1.204103in}{0.794924in}}%
\pgfpathlineto{\pgfqpoint{1.204103in}{0.797873in}}%
\pgfpathlineto{\pgfqpoint{1.208644in}{0.797873in}}%
\pgfpathlineto{\pgfqpoint{1.208644in}{0.794924in}}%
\pgfpathmoveto{\pgfqpoint{1.208644in}{0.794924in}}%
\pgfpathlineto{\pgfqpoint{1.208644in}{0.794924in}}%
\pgfpathlineto{\pgfqpoint{1.208644in}{0.797873in}}%
\pgfpathlineto{\pgfqpoint{1.213185in}{0.797873in}}%
\pgfpathlineto{\pgfqpoint{1.213185in}{0.794924in}}%
\pgfpathmoveto{\pgfqpoint{1.208644in}{0.797873in}}%
\pgfpathlineto{\pgfqpoint{1.208644in}{0.797873in}}%
\pgfpathlineto{\pgfqpoint{1.208644in}{0.800822in}}%
\pgfpathlineto{\pgfqpoint{1.213185in}{0.800822in}}%
\pgfpathlineto{\pgfqpoint{1.213185in}{0.797873in}}%
\pgfpathmoveto{\pgfqpoint{1.213185in}{0.797873in}}%
\pgfpathlineto{\pgfqpoint{1.213185in}{0.797873in}}%
\pgfpathlineto{\pgfqpoint{1.213185in}{0.800822in}}%
\pgfpathlineto{\pgfqpoint{1.217726in}{0.800822in}}%
\pgfpathlineto{\pgfqpoint{1.217726in}{0.797873in}}%
\pgfpathmoveto{\pgfqpoint{1.213185in}{0.800822in}}%
\pgfpathlineto{\pgfqpoint{1.213185in}{0.800822in}}%
\pgfpathlineto{\pgfqpoint{1.213185in}{0.803771in}}%
\pgfpathlineto{\pgfqpoint{1.217726in}{0.803771in}}%
\pgfpathlineto{\pgfqpoint{1.217726in}{0.800822in}}%
\pgfpathmoveto{\pgfqpoint{1.217726in}{0.800822in}}%
\pgfpathlineto{\pgfqpoint{1.217726in}{0.800822in}}%
\pgfpathlineto{\pgfqpoint{1.217726in}{0.803771in}}%
\pgfpathlineto{\pgfqpoint{1.222267in}{0.803771in}}%
\pgfpathlineto{\pgfqpoint{1.222267in}{0.800822in}}%
\pgfpathmoveto{\pgfqpoint{1.217726in}{0.803771in}}%
\pgfpathlineto{\pgfqpoint{1.217726in}{0.803771in}}%
\pgfpathlineto{\pgfqpoint{1.217726in}{0.806720in}}%
\pgfpathlineto{\pgfqpoint{1.222267in}{0.806720in}}%
\pgfpathlineto{\pgfqpoint{1.222267in}{0.803771in}}%
\pgfpathmoveto{\pgfqpoint{1.222267in}{0.803771in}}%
\pgfpathlineto{\pgfqpoint{1.222267in}{0.803771in}}%
\pgfpathlineto{\pgfqpoint{1.222267in}{0.806720in}}%
\pgfpathlineto{\pgfqpoint{1.226808in}{0.806720in}}%
\pgfpathlineto{\pgfqpoint{1.226808in}{0.803771in}}%
\pgfpathmoveto{\pgfqpoint{1.222267in}{0.806720in}}%
\pgfpathlineto{\pgfqpoint{1.222267in}{0.806720in}}%
\pgfpathlineto{\pgfqpoint{1.222267in}{0.809670in}}%
\pgfpathlineto{\pgfqpoint{1.226808in}{0.809670in}}%
\pgfpathlineto{\pgfqpoint{1.226808in}{0.806720in}}%
\pgfpathmoveto{\pgfqpoint{1.226808in}{0.806720in}}%
\pgfpathlineto{\pgfqpoint{1.226808in}{0.806720in}}%
\pgfpathlineto{\pgfqpoint{1.226808in}{0.809670in}}%
\pgfpathlineto{\pgfqpoint{1.231349in}{0.809670in}}%
\pgfpathlineto{\pgfqpoint{1.231349in}{0.806720in}}%
\pgfpathmoveto{\pgfqpoint{1.226808in}{0.809670in}}%
\pgfpathlineto{\pgfqpoint{1.226808in}{0.809670in}}%
\pgfpathlineto{\pgfqpoint{1.226808in}{0.812619in}}%
\pgfpathlineto{\pgfqpoint{1.231349in}{0.812619in}}%
\pgfpathlineto{\pgfqpoint{1.231349in}{0.809670in}}%
\pgfpathmoveto{\pgfqpoint{1.231349in}{0.809670in}}%
\pgfpathlineto{\pgfqpoint{1.231349in}{0.809670in}}%
\pgfpathlineto{\pgfqpoint{1.231349in}{0.812619in}}%
\pgfpathlineto{\pgfqpoint{1.235890in}{0.812619in}}%
\pgfpathlineto{\pgfqpoint{1.235890in}{0.809670in}}%
\pgfpathmoveto{\pgfqpoint{1.231349in}{0.812619in}}%
\pgfpathlineto{\pgfqpoint{1.231349in}{0.812619in}}%
\pgfpathlineto{\pgfqpoint{1.231349in}{0.815568in}}%
\pgfpathlineto{\pgfqpoint{1.235890in}{0.815568in}}%
\pgfpathlineto{\pgfqpoint{1.235890in}{0.812619in}}%
\pgfpathmoveto{\pgfqpoint{1.235890in}{0.812619in}}%
\pgfpathlineto{\pgfqpoint{1.235890in}{0.812619in}}%
\pgfpathlineto{\pgfqpoint{1.235890in}{0.815568in}}%
\pgfpathlineto{\pgfqpoint{1.240431in}{0.815568in}}%
\pgfpathlineto{\pgfqpoint{1.240431in}{0.812619in}}%
\pgfpathmoveto{\pgfqpoint{1.235890in}{0.815568in}}%
\pgfpathlineto{\pgfqpoint{1.235890in}{0.815568in}}%
\pgfpathlineto{\pgfqpoint{1.235890in}{0.818517in}}%
\pgfpathlineto{\pgfqpoint{1.240431in}{0.818517in}}%
\pgfpathlineto{\pgfqpoint{1.240431in}{0.815568in}}%
\pgfpathmoveto{\pgfqpoint{1.240431in}{0.815568in}}%
\pgfpathlineto{\pgfqpoint{1.240431in}{0.815568in}}%
\pgfpathlineto{\pgfqpoint{1.240431in}{0.818517in}}%
\pgfpathlineto{\pgfqpoint{1.244972in}{0.818517in}}%
\pgfpathlineto{\pgfqpoint{1.244972in}{0.815568in}}%
\pgfpathmoveto{\pgfqpoint{1.240431in}{0.818517in}}%
\pgfpathlineto{\pgfqpoint{1.240431in}{0.818517in}}%
\pgfpathlineto{\pgfqpoint{1.240431in}{0.821466in}}%
\pgfpathlineto{\pgfqpoint{1.244972in}{0.821466in}}%
\pgfpathlineto{\pgfqpoint{1.244972in}{0.818517in}}%
\pgfpathmoveto{\pgfqpoint{1.244972in}{0.818517in}}%
\pgfpathlineto{\pgfqpoint{1.244972in}{0.818517in}}%
\pgfpathlineto{\pgfqpoint{1.244972in}{0.821466in}}%
\pgfpathlineto{\pgfqpoint{1.249513in}{0.821466in}}%
\pgfpathlineto{\pgfqpoint{1.249513in}{0.818517in}}%
\pgfpathmoveto{\pgfqpoint{1.244972in}{0.821466in}}%
\pgfpathlineto{\pgfqpoint{1.244972in}{0.821466in}}%
\pgfpathlineto{\pgfqpoint{1.244972in}{0.824416in}}%
\pgfpathlineto{\pgfqpoint{1.249513in}{0.824416in}}%
\pgfpathlineto{\pgfqpoint{1.249513in}{0.821466in}}%
\pgfpathmoveto{\pgfqpoint{1.249513in}{0.821466in}}%
\pgfpathlineto{\pgfqpoint{1.249513in}{0.821466in}}%
\pgfpathlineto{\pgfqpoint{1.249513in}{0.824416in}}%
\pgfpathlineto{\pgfqpoint{1.254054in}{0.824416in}}%
\pgfpathlineto{\pgfqpoint{1.254054in}{0.821466in}}%
\pgfpathmoveto{\pgfqpoint{1.249513in}{0.824416in}}%
\pgfpathlineto{\pgfqpoint{1.249513in}{0.824416in}}%
\pgfpathlineto{\pgfqpoint{1.249513in}{0.827365in}}%
\pgfpathlineto{\pgfqpoint{1.254054in}{0.827365in}}%
\pgfpathlineto{\pgfqpoint{1.254054in}{0.824416in}}%
\pgfpathmoveto{\pgfqpoint{1.254054in}{0.824416in}}%
\pgfpathlineto{\pgfqpoint{1.254054in}{0.824416in}}%
\pgfpathlineto{\pgfqpoint{1.254054in}{0.827365in}}%
\pgfpathlineto{\pgfqpoint{1.258595in}{0.827365in}}%
\pgfpathlineto{\pgfqpoint{1.258595in}{0.824416in}}%
\pgfpathmoveto{\pgfqpoint{1.254054in}{0.827365in}}%
\pgfpathlineto{\pgfqpoint{1.254054in}{0.827365in}}%
\pgfpathlineto{\pgfqpoint{1.254054in}{0.830314in}}%
\pgfpathlineto{\pgfqpoint{1.258595in}{0.830314in}}%
\pgfpathlineto{\pgfqpoint{1.258595in}{0.827365in}}%
\pgfpathmoveto{\pgfqpoint{1.258595in}{0.827365in}}%
\pgfpathlineto{\pgfqpoint{1.258595in}{0.827365in}}%
\pgfpathlineto{\pgfqpoint{1.258595in}{0.830314in}}%
\pgfpathlineto{\pgfqpoint{1.263136in}{0.830314in}}%
\pgfpathlineto{\pgfqpoint{1.263136in}{0.827365in}}%
\pgfpathmoveto{\pgfqpoint{1.258595in}{0.830314in}}%
\pgfpathlineto{\pgfqpoint{1.258595in}{0.830314in}}%
\pgfpathlineto{\pgfqpoint{1.258595in}{0.833263in}}%
\pgfpathlineto{\pgfqpoint{1.263136in}{0.833263in}}%
\pgfpathlineto{\pgfqpoint{1.263136in}{0.830314in}}%
\pgfpathmoveto{\pgfqpoint{1.263136in}{0.830314in}}%
\pgfpathlineto{\pgfqpoint{1.263136in}{0.830314in}}%
\pgfpathlineto{\pgfqpoint{1.263136in}{0.833263in}}%
\pgfpathlineto{\pgfqpoint{1.267677in}{0.833263in}}%
\pgfpathlineto{\pgfqpoint{1.267677in}{0.830314in}}%
\pgfpathmoveto{\pgfqpoint{1.263136in}{0.833263in}}%
\pgfpathlineto{\pgfqpoint{1.263136in}{0.833263in}}%
\pgfpathlineto{\pgfqpoint{1.263136in}{0.836212in}}%
\pgfpathlineto{\pgfqpoint{1.267677in}{0.836212in}}%
\pgfpathlineto{\pgfqpoint{1.267677in}{0.833263in}}%
\pgfpathmoveto{\pgfqpoint{1.267677in}{0.833263in}}%
\pgfpathlineto{\pgfqpoint{1.267677in}{0.833263in}}%
\pgfpathlineto{\pgfqpoint{1.267677in}{0.836212in}}%
\pgfpathlineto{\pgfqpoint{1.272218in}{0.836212in}}%
\pgfpathlineto{\pgfqpoint{1.272218in}{0.833263in}}%
\pgfpathmoveto{\pgfqpoint{1.267677in}{0.836212in}}%
\pgfpathlineto{\pgfqpoint{1.267677in}{0.836212in}}%
\pgfpathlineto{\pgfqpoint{1.267677in}{0.839162in}}%
\pgfpathlineto{\pgfqpoint{1.272218in}{0.839162in}}%
\pgfpathlineto{\pgfqpoint{1.272218in}{0.836212in}}%
\pgfpathmoveto{\pgfqpoint{1.272218in}{0.836212in}}%
\pgfpathlineto{\pgfqpoint{1.272218in}{0.836212in}}%
\pgfpathlineto{\pgfqpoint{1.272218in}{0.839162in}}%
\pgfpathlineto{\pgfqpoint{1.276759in}{0.839162in}}%
\pgfpathlineto{\pgfqpoint{1.276759in}{0.836212in}}%
\pgfpathmoveto{\pgfqpoint{1.272218in}{0.839162in}}%
\pgfpathlineto{\pgfqpoint{1.272218in}{0.839162in}}%
\pgfpathlineto{\pgfqpoint{1.272218in}{0.842111in}}%
\pgfpathlineto{\pgfqpoint{1.276759in}{0.842111in}}%
\pgfpathlineto{\pgfqpoint{1.276759in}{0.839162in}}%
\pgfpathmoveto{\pgfqpoint{1.276759in}{0.839162in}}%
\pgfpathlineto{\pgfqpoint{1.276759in}{0.839162in}}%
\pgfpathlineto{\pgfqpoint{1.276759in}{0.842111in}}%
\pgfpathlineto{\pgfqpoint{1.281300in}{0.842111in}}%
\pgfpathlineto{\pgfqpoint{1.281300in}{0.839162in}}%
\pgfpathmoveto{\pgfqpoint{1.276759in}{0.842111in}}%
\pgfpathlineto{\pgfqpoint{1.276759in}{0.842111in}}%
\pgfpathlineto{\pgfqpoint{1.276759in}{0.845060in}}%
\pgfpathlineto{\pgfqpoint{1.281300in}{0.845060in}}%
\pgfpathlineto{\pgfqpoint{1.281300in}{0.842111in}}%
\pgfpathmoveto{\pgfqpoint{1.281300in}{0.842111in}}%
\pgfpathlineto{\pgfqpoint{1.281300in}{0.842111in}}%
\pgfpathlineto{\pgfqpoint{1.281300in}{0.845060in}}%
\pgfpathlineto{\pgfqpoint{1.285841in}{0.845060in}}%
\pgfpathlineto{\pgfqpoint{1.285841in}{0.842111in}}%
\pgfpathmoveto{\pgfqpoint{1.281300in}{0.845060in}}%
\pgfpathlineto{\pgfqpoint{1.281300in}{0.845060in}}%
\pgfpathlineto{\pgfqpoint{1.281300in}{0.848009in}}%
\pgfpathlineto{\pgfqpoint{1.285841in}{0.848009in}}%
\pgfpathlineto{\pgfqpoint{1.285841in}{0.845060in}}%
\pgfpathmoveto{\pgfqpoint{1.285841in}{0.845060in}}%
\pgfpathlineto{\pgfqpoint{1.285841in}{0.845060in}}%
\pgfpathlineto{\pgfqpoint{1.285841in}{0.848009in}}%
\pgfpathlineto{\pgfqpoint{1.290382in}{0.848009in}}%
\pgfpathlineto{\pgfqpoint{1.290382in}{0.845060in}}%
\pgfpathmoveto{\pgfqpoint{1.285841in}{0.848009in}}%
\pgfpathlineto{\pgfqpoint{1.285841in}{0.848009in}}%
\pgfpathlineto{\pgfqpoint{1.285841in}{0.850958in}}%
\pgfpathlineto{\pgfqpoint{1.290382in}{0.850958in}}%
\pgfpathlineto{\pgfqpoint{1.290382in}{0.848009in}}%
\pgfpathmoveto{\pgfqpoint{1.290382in}{0.848009in}}%
\pgfpathlineto{\pgfqpoint{1.290382in}{0.848009in}}%
\pgfpathlineto{\pgfqpoint{1.290382in}{0.850958in}}%
\pgfpathlineto{\pgfqpoint{1.294923in}{0.850958in}}%
\pgfpathlineto{\pgfqpoint{1.294923in}{0.848009in}}%
\pgfpathmoveto{\pgfqpoint{1.290382in}{0.850958in}}%
\pgfpathlineto{\pgfqpoint{1.290382in}{0.850958in}}%
\pgfpathlineto{\pgfqpoint{1.290382in}{0.853908in}}%
\pgfpathlineto{\pgfqpoint{1.294923in}{0.853908in}}%
\pgfpathlineto{\pgfqpoint{1.294923in}{0.850958in}}%
\pgfpathmoveto{\pgfqpoint{1.294923in}{0.850958in}}%
\pgfpathlineto{\pgfqpoint{1.294923in}{0.850958in}}%
\pgfpathlineto{\pgfqpoint{1.294923in}{0.853908in}}%
\pgfpathlineto{\pgfqpoint{1.299464in}{0.853908in}}%
\pgfpathlineto{\pgfqpoint{1.299464in}{0.850958in}}%
\pgfpathmoveto{\pgfqpoint{1.294923in}{0.853908in}}%
\pgfpathlineto{\pgfqpoint{1.294923in}{0.853908in}}%
\pgfpathlineto{\pgfqpoint{1.294923in}{0.856857in}}%
\pgfpathlineto{\pgfqpoint{1.299464in}{0.856857in}}%
\pgfpathlineto{\pgfqpoint{1.299464in}{0.853908in}}%
\pgfpathmoveto{\pgfqpoint{1.299464in}{0.853908in}}%
\pgfpathlineto{\pgfqpoint{1.299464in}{0.853908in}}%
\pgfpathlineto{\pgfqpoint{1.299464in}{0.856857in}}%
\pgfpathlineto{\pgfqpoint{1.304005in}{0.856857in}}%
\pgfpathlineto{\pgfqpoint{1.304005in}{0.853908in}}%
\pgfpathmoveto{\pgfqpoint{1.299464in}{0.856857in}}%
\pgfpathlineto{\pgfqpoint{1.299464in}{0.856857in}}%
\pgfpathlineto{\pgfqpoint{1.299464in}{0.859806in}}%
\pgfpathlineto{\pgfqpoint{1.304005in}{0.859806in}}%
\pgfpathlineto{\pgfqpoint{1.304005in}{0.856857in}}%
\pgfpathmoveto{\pgfqpoint{1.304005in}{0.856857in}}%
\pgfpathlineto{\pgfqpoint{1.304005in}{0.856857in}}%
\pgfpathlineto{\pgfqpoint{1.304005in}{0.859806in}}%
\pgfpathlineto{\pgfqpoint{1.308546in}{0.859806in}}%
\pgfpathlineto{\pgfqpoint{1.308546in}{0.856857in}}%
\pgfpathmoveto{\pgfqpoint{1.304005in}{0.859806in}}%
\pgfpathlineto{\pgfqpoint{1.304005in}{0.859806in}}%
\pgfpathlineto{\pgfqpoint{1.304005in}{0.862755in}}%
\pgfpathlineto{\pgfqpoint{1.308546in}{0.862755in}}%
\pgfpathlineto{\pgfqpoint{1.308546in}{0.859806in}}%
\pgfpathmoveto{\pgfqpoint{1.308546in}{0.859806in}}%
\pgfpathlineto{\pgfqpoint{1.308546in}{0.859806in}}%
\pgfpathlineto{\pgfqpoint{1.308546in}{0.862755in}}%
\pgfpathlineto{\pgfqpoint{1.313087in}{0.862755in}}%
\pgfpathlineto{\pgfqpoint{1.313087in}{0.859806in}}%
\pgfpathmoveto{\pgfqpoint{1.308546in}{0.862755in}}%
\pgfpathlineto{\pgfqpoint{1.308546in}{0.862755in}}%
\pgfpathlineto{\pgfqpoint{1.308546in}{0.865704in}}%
\pgfpathlineto{\pgfqpoint{1.313087in}{0.865704in}}%
\pgfpathlineto{\pgfqpoint{1.313087in}{0.862755in}}%
\pgfpathmoveto{\pgfqpoint{1.313087in}{0.862755in}}%
\pgfpathlineto{\pgfqpoint{1.313087in}{0.862755in}}%
\pgfpathlineto{\pgfqpoint{1.313087in}{0.865704in}}%
\pgfpathlineto{\pgfqpoint{1.317628in}{0.865704in}}%
\pgfpathlineto{\pgfqpoint{1.317628in}{0.862755in}}%
\pgfpathmoveto{\pgfqpoint{1.313087in}{0.865704in}}%
\pgfpathlineto{\pgfqpoint{1.313087in}{0.865704in}}%
\pgfpathlineto{\pgfqpoint{1.313087in}{0.868654in}}%
\pgfpathlineto{\pgfqpoint{1.317628in}{0.868654in}}%
\pgfpathlineto{\pgfqpoint{1.317628in}{0.865704in}}%
\pgfpathmoveto{\pgfqpoint{1.317628in}{0.865704in}}%
\pgfpathlineto{\pgfqpoint{1.317628in}{0.865704in}}%
\pgfpathlineto{\pgfqpoint{1.317628in}{0.868654in}}%
\pgfpathlineto{\pgfqpoint{1.322169in}{0.868654in}}%
\pgfpathlineto{\pgfqpoint{1.322169in}{0.865704in}}%
\pgfpathmoveto{\pgfqpoint{1.317628in}{0.868654in}}%
\pgfpathlineto{\pgfqpoint{1.317628in}{0.868654in}}%
\pgfpathlineto{\pgfqpoint{1.317628in}{0.871603in}}%
\pgfpathlineto{\pgfqpoint{1.322169in}{0.871603in}}%
\pgfpathlineto{\pgfqpoint{1.322169in}{0.868654in}}%
\pgfpathmoveto{\pgfqpoint{1.322169in}{0.868654in}}%
\pgfpathlineto{\pgfqpoint{1.322169in}{0.868654in}}%
\pgfpathlineto{\pgfqpoint{1.322169in}{0.871603in}}%
\pgfpathlineto{\pgfqpoint{1.326710in}{0.871603in}}%
\pgfpathlineto{\pgfqpoint{1.326710in}{0.868654in}}%
\pgfpathmoveto{\pgfqpoint{1.322169in}{0.871603in}}%
\pgfpathlineto{\pgfqpoint{1.322169in}{0.871603in}}%
\pgfpathlineto{\pgfqpoint{1.322169in}{0.874552in}}%
\pgfpathlineto{\pgfqpoint{1.326710in}{0.874552in}}%
\pgfpathlineto{\pgfqpoint{1.326710in}{0.871603in}}%
\pgfpathmoveto{\pgfqpoint{1.326710in}{0.871603in}}%
\pgfpathlineto{\pgfqpoint{1.326710in}{0.871603in}}%
\pgfpathlineto{\pgfqpoint{1.326710in}{0.874552in}}%
\pgfpathlineto{\pgfqpoint{1.331251in}{0.874552in}}%
\pgfpathlineto{\pgfqpoint{1.331251in}{0.871603in}}%
\pgfpathmoveto{\pgfqpoint{1.326710in}{0.874552in}}%
\pgfpathlineto{\pgfqpoint{1.326710in}{0.874552in}}%
\pgfpathlineto{\pgfqpoint{1.326710in}{0.877501in}}%
\pgfpathlineto{\pgfqpoint{1.331251in}{0.877501in}}%
\pgfpathlineto{\pgfqpoint{1.331251in}{0.874552in}}%
\pgfpathmoveto{\pgfqpoint{1.331251in}{0.874552in}}%
\pgfpathlineto{\pgfqpoint{1.331251in}{0.874552in}}%
\pgfpathlineto{\pgfqpoint{1.331251in}{0.877501in}}%
\pgfpathlineto{\pgfqpoint{1.335792in}{0.877501in}}%
\pgfpathlineto{\pgfqpoint{1.335792in}{0.874552in}}%
\pgfpathmoveto{\pgfqpoint{1.331251in}{0.877501in}}%
\pgfpathlineto{\pgfqpoint{1.331251in}{0.877501in}}%
\pgfpathlineto{\pgfqpoint{1.331251in}{0.880450in}}%
\pgfpathlineto{\pgfqpoint{1.335792in}{0.880450in}}%
\pgfpathlineto{\pgfqpoint{1.335792in}{0.877501in}}%
\pgfpathmoveto{\pgfqpoint{1.335792in}{0.877501in}}%
\pgfpathlineto{\pgfqpoint{1.335792in}{0.877501in}}%
\pgfpathlineto{\pgfqpoint{1.335792in}{0.880450in}}%
\pgfpathlineto{\pgfqpoint{1.340333in}{0.880450in}}%
\pgfpathlineto{\pgfqpoint{1.340333in}{0.877501in}}%
\pgfpathmoveto{\pgfqpoint{1.335792in}{0.880450in}}%
\pgfpathlineto{\pgfqpoint{1.335792in}{0.880450in}}%
\pgfpathlineto{\pgfqpoint{1.335792in}{0.883399in}}%
\pgfpathlineto{\pgfqpoint{1.340333in}{0.883399in}}%
\pgfpathlineto{\pgfqpoint{1.340333in}{0.880450in}}%
\pgfpathmoveto{\pgfqpoint{1.340333in}{0.880450in}}%
\pgfpathlineto{\pgfqpoint{1.340333in}{0.880450in}}%
\pgfpathlineto{\pgfqpoint{1.340333in}{0.883399in}}%
\pgfpathlineto{\pgfqpoint{1.344873in}{0.883399in}}%
\pgfpathlineto{\pgfqpoint{1.344873in}{0.880450in}}%
\pgfpathmoveto{\pgfqpoint{1.340333in}{0.883399in}}%
\pgfpathlineto{\pgfqpoint{1.340333in}{0.883399in}}%
\pgfpathlineto{\pgfqpoint{1.340333in}{0.886349in}}%
\pgfpathlineto{\pgfqpoint{1.344873in}{0.886349in}}%
\pgfpathlineto{\pgfqpoint{1.344873in}{0.883399in}}%
\pgfpathmoveto{\pgfqpoint{1.344873in}{0.883399in}}%
\pgfpathlineto{\pgfqpoint{1.344873in}{0.883399in}}%
\pgfpathlineto{\pgfqpoint{1.344873in}{0.886349in}}%
\pgfpathlineto{\pgfqpoint{1.349414in}{0.886349in}}%
\pgfpathlineto{\pgfqpoint{1.349414in}{0.883399in}}%
\pgfpathmoveto{\pgfqpoint{1.344873in}{0.886349in}}%
\pgfpathlineto{\pgfqpoint{1.344873in}{0.886349in}}%
\pgfpathlineto{\pgfqpoint{1.344873in}{0.889298in}}%
\pgfpathlineto{\pgfqpoint{1.349414in}{0.889298in}}%
\pgfpathlineto{\pgfqpoint{1.349414in}{0.886349in}}%
\pgfpathmoveto{\pgfqpoint{1.349414in}{0.886349in}}%
\pgfpathlineto{\pgfqpoint{1.349414in}{0.886349in}}%
\pgfpathlineto{\pgfqpoint{1.349414in}{0.889298in}}%
\pgfpathlineto{\pgfqpoint{1.353955in}{0.889298in}}%
\pgfpathlineto{\pgfqpoint{1.353955in}{0.886349in}}%
\pgfpathmoveto{\pgfqpoint{1.349414in}{0.889298in}}%
\pgfpathlineto{\pgfqpoint{1.349414in}{0.889298in}}%
\pgfpathlineto{\pgfqpoint{1.349414in}{0.892247in}}%
\pgfpathlineto{\pgfqpoint{1.353955in}{0.892247in}}%
\pgfpathlineto{\pgfqpoint{1.353955in}{0.889298in}}%
\pgfpathmoveto{\pgfqpoint{1.353955in}{0.889298in}}%
\pgfpathlineto{\pgfqpoint{1.353955in}{0.889298in}}%
\pgfpathlineto{\pgfqpoint{1.353955in}{0.892247in}}%
\pgfpathlineto{\pgfqpoint{1.358496in}{0.892247in}}%
\pgfpathlineto{\pgfqpoint{1.358496in}{0.889298in}}%
\pgfpathmoveto{\pgfqpoint{1.353955in}{0.892247in}}%
\pgfpathlineto{\pgfqpoint{1.353955in}{0.892247in}}%
\pgfpathlineto{\pgfqpoint{1.353955in}{0.895196in}}%
\pgfpathlineto{\pgfqpoint{1.358496in}{0.895196in}}%
\pgfpathlineto{\pgfqpoint{1.358496in}{0.892247in}}%
\pgfpathmoveto{\pgfqpoint{1.358496in}{0.892247in}}%
\pgfpathlineto{\pgfqpoint{1.358496in}{0.892247in}}%
\pgfpathlineto{\pgfqpoint{1.358496in}{0.895196in}}%
\pgfpathlineto{\pgfqpoint{1.363037in}{0.895196in}}%
\pgfpathlineto{\pgfqpoint{1.363037in}{0.892247in}}%
\pgfpathmoveto{\pgfqpoint{1.358496in}{0.895196in}}%
\pgfpathlineto{\pgfqpoint{1.358496in}{0.895196in}}%
\pgfpathlineto{\pgfqpoint{1.358496in}{0.898145in}}%
\pgfpathlineto{\pgfqpoint{1.363037in}{0.898145in}}%
\pgfpathlineto{\pgfqpoint{1.363037in}{0.895196in}}%
\pgfpathmoveto{\pgfqpoint{1.363037in}{0.895196in}}%
\pgfpathlineto{\pgfqpoint{1.363037in}{0.895196in}}%
\pgfpathlineto{\pgfqpoint{1.363037in}{0.898145in}}%
\pgfpathlineto{\pgfqpoint{1.367578in}{0.898145in}}%
\pgfpathlineto{\pgfqpoint{1.367578in}{0.895196in}}%
\pgfpathmoveto{\pgfqpoint{1.363037in}{0.898145in}}%
\pgfpathlineto{\pgfqpoint{1.363037in}{0.898145in}}%
\pgfpathlineto{\pgfqpoint{1.363037in}{0.901094in}}%
\pgfpathlineto{\pgfqpoint{1.367578in}{0.901094in}}%
\pgfpathlineto{\pgfqpoint{1.367578in}{0.898145in}}%
\pgfpathmoveto{\pgfqpoint{1.367578in}{0.898145in}}%
\pgfpathlineto{\pgfqpoint{1.367578in}{0.898145in}}%
\pgfpathlineto{\pgfqpoint{1.367578in}{0.901094in}}%
\pgfpathlineto{\pgfqpoint{1.372119in}{0.901094in}}%
\pgfpathlineto{\pgfqpoint{1.372119in}{0.898145in}}%
\pgfpathmoveto{\pgfqpoint{1.367578in}{0.901094in}}%
\pgfpathlineto{\pgfqpoint{1.367578in}{0.901094in}}%
\pgfpathlineto{\pgfqpoint{1.367578in}{0.904044in}}%
\pgfpathlineto{\pgfqpoint{1.372119in}{0.904044in}}%
\pgfpathlineto{\pgfqpoint{1.372119in}{0.901094in}}%
\pgfpathmoveto{\pgfqpoint{1.372119in}{0.901094in}}%
\pgfpathlineto{\pgfqpoint{1.372119in}{0.901094in}}%
\pgfpathlineto{\pgfqpoint{1.372119in}{0.904044in}}%
\pgfpathlineto{\pgfqpoint{1.376660in}{0.904044in}}%
\pgfpathlineto{\pgfqpoint{1.376660in}{0.901094in}}%
\pgfpathmoveto{\pgfqpoint{1.372119in}{0.904044in}}%
\pgfpathlineto{\pgfqpoint{1.372119in}{0.904044in}}%
\pgfpathlineto{\pgfqpoint{1.372119in}{0.906993in}}%
\pgfpathlineto{\pgfqpoint{1.376660in}{0.906993in}}%
\pgfpathlineto{\pgfqpoint{1.376660in}{0.904044in}}%
\pgfpathmoveto{\pgfqpoint{1.376660in}{0.904044in}}%
\pgfpathlineto{\pgfqpoint{1.376660in}{0.904044in}}%
\pgfpathlineto{\pgfqpoint{1.376660in}{0.906993in}}%
\pgfpathlineto{\pgfqpoint{1.381201in}{0.906993in}}%
\pgfpathlineto{\pgfqpoint{1.381201in}{0.904044in}}%
\pgfpathmoveto{\pgfqpoint{1.376660in}{0.906993in}}%
\pgfpathlineto{\pgfqpoint{1.376660in}{0.906993in}}%
\pgfpathlineto{\pgfqpoint{1.376660in}{0.909942in}}%
\pgfpathlineto{\pgfqpoint{1.381201in}{0.909942in}}%
\pgfpathlineto{\pgfqpoint{1.381201in}{0.906993in}}%
\pgfpathmoveto{\pgfqpoint{1.381201in}{0.906993in}}%
\pgfpathlineto{\pgfqpoint{1.381201in}{0.906993in}}%
\pgfpathlineto{\pgfqpoint{1.381201in}{0.909942in}}%
\pgfpathlineto{\pgfqpoint{1.385742in}{0.909942in}}%
\pgfpathlineto{\pgfqpoint{1.385742in}{0.906993in}}%
\pgfpathmoveto{\pgfqpoint{1.381201in}{0.909942in}}%
\pgfpathlineto{\pgfqpoint{1.381201in}{0.909942in}}%
\pgfpathlineto{\pgfqpoint{1.381201in}{0.912891in}}%
\pgfpathlineto{\pgfqpoint{1.385742in}{0.912891in}}%
\pgfpathlineto{\pgfqpoint{1.385742in}{0.909942in}}%
\pgfpathmoveto{\pgfqpoint{1.385742in}{0.909942in}}%
\pgfpathlineto{\pgfqpoint{1.385742in}{0.909942in}}%
\pgfpathlineto{\pgfqpoint{1.385742in}{0.912891in}}%
\pgfpathlineto{\pgfqpoint{1.390283in}{0.912891in}}%
\pgfpathlineto{\pgfqpoint{1.390283in}{0.909942in}}%
\pgfpathmoveto{\pgfqpoint{1.385742in}{0.912891in}}%
\pgfpathlineto{\pgfqpoint{1.385742in}{0.912891in}}%
\pgfpathlineto{\pgfqpoint{1.385742in}{0.915840in}}%
\pgfpathlineto{\pgfqpoint{1.390283in}{0.915840in}}%
\pgfpathlineto{\pgfqpoint{1.390283in}{0.912891in}}%
\pgfpathmoveto{\pgfqpoint{1.390283in}{0.912891in}}%
\pgfpathlineto{\pgfqpoint{1.390283in}{0.912891in}}%
\pgfpathlineto{\pgfqpoint{1.390283in}{0.915840in}}%
\pgfpathlineto{\pgfqpoint{1.394823in}{0.915840in}}%
\pgfpathlineto{\pgfqpoint{1.394823in}{0.912891in}}%
\pgfpathmoveto{\pgfqpoint{1.390283in}{0.915840in}}%
\pgfpathlineto{\pgfqpoint{1.390283in}{0.915840in}}%
\pgfpathlineto{\pgfqpoint{1.390283in}{0.918790in}}%
\pgfpathlineto{\pgfqpoint{1.394823in}{0.918790in}}%
\pgfpathlineto{\pgfqpoint{1.394823in}{0.915840in}}%
\pgfpathmoveto{\pgfqpoint{1.394823in}{0.915840in}}%
\pgfpathlineto{\pgfqpoint{1.394823in}{0.915840in}}%
\pgfpathlineto{\pgfqpoint{1.394823in}{0.918790in}}%
\pgfpathlineto{\pgfqpoint{1.399364in}{0.918790in}}%
\pgfpathlineto{\pgfqpoint{1.399364in}{0.915840in}}%
\pgfpathmoveto{\pgfqpoint{1.394823in}{0.918790in}}%
\pgfpathlineto{\pgfqpoint{1.394823in}{0.918790in}}%
\pgfpathlineto{\pgfqpoint{1.394823in}{0.921739in}}%
\pgfpathlineto{\pgfqpoint{1.399364in}{0.921739in}}%
\pgfpathlineto{\pgfqpoint{1.399364in}{0.918790in}}%
\pgfpathmoveto{\pgfqpoint{1.399364in}{0.918790in}}%
\pgfpathlineto{\pgfqpoint{1.399364in}{0.918790in}}%
\pgfpathlineto{\pgfqpoint{1.399364in}{0.921739in}}%
\pgfpathlineto{\pgfqpoint{1.403905in}{0.921739in}}%
\pgfpathlineto{\pgfqpoint{1.403905in}{0.918790in}}%
\pgfpathmoveto{\pgfqpoint{1.399364in}{0.921739in}}%
\pgfpathlineto{\pgfqpoint{1.399364in}{0.921739in}}%
\pgfpathlineto{\pgfqpoint{1.399364in}{0.924688in}}%
\pgfpathlineto{\pgfqpoint{1.403905in}{0.924688in}}%
\pgfpathlineto{\pgfqpoint{1.403905in}{0.921739in}}%
\pgfpathmoveto{\pgfqpoint{1.403905in}{0.921739in}}%
\pgfpathlineto{\pgfqpoint{1.403905in}{0.921739in}}%
\pgfpathlineto{\pgfqpoint{1.403905in}{0.924688in}}%
\pgfpathlineto{\pgfqpoint{1.408446in}{0.924688in}}%
\pgfpathlineto{\pgfqpoint{1.408446in}{0.921739in}}%
\pgfpathmoveto{\pgfqpoint{1.403905in}{0.924688in}}%
\pgfpathlineto{\pgfqpoint{1.403905in}{0.924688in}}%
\pgfpathlineto{\pgfqpoint{1.403905in}{0.927637in}}%
\pgfpathlineto{\pgfqpoint{1.408446in}{0.927637in}}%
\pgfpathlineto{\pgfqpoint{1.408446in}{0.924688in}}%
\pgfpathmoveto{\pgfqpoint{1.408446in}{0.924688in}}%
\pgfpathlineto{\pgfqpoint{1.408446in}{0.924688in}}%
\pgfpathlineto{\pgfqpoint{1.408446in}{0.927637in}}%
\pgfpathlineto{\pgfqpoint{1.412987in}{0.927637in}}%
\pgfpathlineto{\pgfqpoint{1.412987in}{0.924688in}}%
\pgfpathmoveto{\pgfqpoint{1.408446in}{0.927637in}}%
\pgfpathlineto{\pgfqpoint{1.408446in}{0.927637in}}%
\pgfpathlineto{\pgfqpoint{1.408446in}{0.930586in}}%
\pgfpathlineto{\pgfqpoint{1.412987in}{0.930586in}}%
\pgfpathlineto{\pgfqpoint{1.412987in}{0.927637in}}%
\pgfpathmoveto{\pgfqpoint{1.412987in}{0.927637in}}%
\pgfpathlineto{\pgfqpoint{1.412987in}{0.927637in}}%
\pgfpathlineto{\pgfqpoint{1.412987in}{0.930586in}}%
\pgfpathlineto{\pgfqpoint{1.417528in}{0.930586in}}%
\pgfpathlineto{\pgfqpoint{1.417528in}{0.927637in}}%
\pgfpathmoveto{\pgfqpoint{1.412987in}{0.930586in}}%
\pgfpathlineto{\pgfqpoint{1.412987in}{0.930586in}}%
\pgfpathlineto{\pgfqpoint{1.412987in}{0.933535in}}%
\pgfpathlineto{\pgfqpoint{1.417528in}{0.933535in}}%
\pgfpathlineto{\pgfqpoint{1.417528in}{0.930586in}}%
\pgfpathmoveto{\pgfqpoint{1.417528in}{0.930586in}}%
\pgfpathlineto{\pgfqpoint{1.417528in}{0.930586in}}%
\pgfpathlineto{\pgfqpoint{1.417528in}{0.933535in}}%
\pgfpathlineto{\pgfqpoint{1.422069in}{0.933535in}}%
\pgfpathlineto{\pgfqpoint{1.422069in}{0.930586in}}%
\pgfpathmoveto{\pgfqpoint{1.417528in}{0.933535in}}%
\pgfpathlineto{\pgfqpoint{1.417528in}{0.933535in}}%
\pgfpathlineto{\pgfqpoint{1.417528in}{0.936485in}}%
\pgfpathlineto{\pgfqpoint{1.422069in}{0.936485in}}%
\pgfpathlineto{\pgfqpoint{1.422069in}{0.933535in}}%
\pgfpathmoveto{\pgfqpoint{1.422069in}{0.933535in}}%
\pgfpathlineto{\pgfqpoint{1.422069in}{0.933535in}}%
\pgfpathlineto{\pgfqpoint{1.422069in}{0.936485in}}%
\pgfpathlineto{\pgfqpoint{1.426610in}{0.936485in}}%
\pgfpathlineto{\pgfqpoint{1.426610in}{0.933535in}}%
\pgfpathmoveto{\pgfqpoint{1.422069in}{0.936485in}}%
\pgfpathlineto{\pgfqpoint{1.422069in}{0.936485in}}%
\pgfpathlineto{\pgfqpoint{1.422069in}{0.939434in}}%
\pgfpathlineto{\pgfqpoint{1.426610in}{0.939434in}}%
\pgfpathlineto{\pgfqpoint{1.426610in}{0.936485in}}%
\pgfpathmoveto{\pgfqpoint{1.426610in}{0.936485in}}%
\pgfpathlineto{\pgfqpoint{1.426610in}{0.936485in}}%
\pgfpathlineto{\pgfqpoint{1.426610in}{0.939434in}}%
\pgfpathlineto{\pgfqpoint{1.431151in}{0.939434in}}%
\pgfpathlineto{\pgfqpoint{1.431151in}{0.936485in}}%
\pgfpathmoveto{\pgfqpoint{1.426610in}{0.939434in}}%
\pgfpathlineto{\pgfqpoint{1.426610in}{0.939434in}}%
\pgfpathlineto{\pgfqpoint{1.426610in}{0.942383in}}%
\pgfpathlineto{\pgfqpoint{1.431151in}{0.942383in}}%
\pgfpathlineto{\pgfqpoint{1.431151in}{0.939434in}}%
\pgfpathmoveto{\pgfqpoint{1.431151in}{0.939434in}}%
\pgfpathlineto{\pgfqpoint{1.431151in}{0.939434in}}%
\pgfpathlineto{\pgfqpoint{1.431151in}{0.942383in}}%
\pgfpathlineto{\pgfqpoint{1.435692in}{0.942383in}}%
\pgfpathlineto{\pgfqpoint{1.435692in}{0.939434in}}%
\pgfpathmoveto{\pgfqpoint{1.431151in}{0.942383in}}%
\pgfpathlineto{\pgfqpoint{1.431151in}{0.942383in}}%
\pgfpathlineto{\pgfqpoint{1.431151in}{0.945332in}}%
\pgfpathlineto{\pgfqpoint{1.435692in}{0.945332in}}%
\pgfpathlineto{\pgfqpoint{1.435692in}{0.942383in}}%
\pgfpathmoveto{\pgfqpoint{1.435692in}{0.942383in}}%
\pgfpathlineto{\pgfqpoint{1.435692in}{0.942383in}}%
\pgfpathlineto{\pgfqpoint{1.435692in}{0.945332in}}%
\pgfpathlineto{\pgfqpoint{1.440233in}{0.945332in}}%
\pgfpathlineto{\pgfqpoint{1.440233in}{0.942383in}}%
\pgfpathmoveto{\pgfqpoint{1.435692in}{0.945332in}}%
\pgfpathlineto{\pgfqpoint{1.435692in}{0.945332in}}%
\pgfpathlineto{\pgfqpoint{1.435692in}{0.948281in}}%
\pgfpathlineto{\pgfqpoint{1.440233in}{0.948281in}}%
\pgfpathlineto{\pgfqpoint{1.440233in}{0.945332in}}%
\pgfpathmoveto{\pgfqpoint{1.440233in}{0.945332in}}%
\pgfpathlineto{\pgfqpoint{1.440233in}{0.945332in}}%
\pgfpathlineto{\pgfqpoint{1.440233in}{0.948281in}}%
\pgfpathlineto{\pgfqpoint{1.444773in}{0.948281in}}%
\pgfpathlineto{\pgfqpoint{1.444773in}{0.945332in}}%
\pgfpathmoveto{\pgfqpoint{1.440233in}{0.948281in}}%
\pgfpathlineto{\pgfqpoint{1.440233in}{0.948281in}}%
\pgfpathlineto{\pgfqpoint{1.440233in}{0.951230in}}%
\pgfpathlineto{\pgfqpoint{1.444773in}{0.951230in}}%
\pgfpathlineto{\pgfqpoint{1.444773in}{0.948281in}}%
\pgfpathmoveto{\pgfqpoint{1.444773in}{0.948281in}}%
\pgfpathlineto{\pgfqpoint{1.444773in}{0.948281in}}%
\pgfpathlineto{\pgfqpoint{1.444773in}{0.951230in}}%
\pgfpathlineto{\pgfqpoint{1.449314in}{0.951230in}}%
\pgfpathlineto{\pgfqpoint{1.449314in}{0.948281in}}%
\pgfpathmoveto{\pgfqpoint{1.444773in}{0.951230in}}%
\pgfpathlineto{\pgfqpoint{1.444773in}{0.951230in}}%
\pgfpathlineto{\pgfqpoint{1.444773in}{0.954180in}}%
\pgfpathlineto{\pgfqpoint{1.449314in}{0.954180in}}%
\pgfpathlineto{\pgfqpoint{1.449314in}{0.951230in}}%
\pgfpathmoveto{\pgfqpoint{1.449314in}{0.951230in}}%
\pgfpathlineto{\pgfqpoint{1.449314in}{0.951230in}}%
\pgfpathlineto{\pgfqpoint{1.449314in}{0.954180in}}%
\pgfpathlineto{\pgfqpoint{1.453855in}{0.954180in}}%
\pgfpathlineto{\pgfqpoint{1.453855in}{0.951230in}}%
\pgfpathmoveto{\pgfqpoint{1.449314in}{0.954180in}}%
\pgfpathlineto{\pgfqpoint{1.449314in}{0.954180in}}%
\pgfpathlineto{\pgfqpoint{1.449314in}{0.957129in}}%
\pgfpathlineto{\pgfqpoint{1.453855in}{0.957129in}}%
\pgfpathlineto{\pgfqpoint{1.453855in}{0.954180in}}%
\pgfpathmoveto{\pgfqpoint{1.453855in}{0.954180in}}%
\pgfpathlineto{\pgfqpoint{1.453855in}{0.954180in}}%
\pgfpathlineto{\pgfqpoint{1.453855in}{0.957129in}}%
\pgfpathlineto{\pgfqpoint{1.458396in}{0.957129in}}%
\pgfpathlineto{\pgfqpoint{1.458396in}{0.954180in}}%
\pgfpathmoveto{\pgfqpoint{1.453855in}{0.957129in}}%
\pgfpathlineto{\pgfqpoint{1.453855in}{0.957129in}}%
\pgfpathlineto{\pgfqpoint{1.453855in}{0.960078in}}%
\pgfpathlineto{\pgfqpoint{1.458396in}{0.960078in}}%
\pgfpathlineto{\pgfqpoint{1.458396in}{0.957129in}}%
\pgfpathmoveto{\pgfqpoint{1.458396in}{0.957129in}}%
\pgfpathlineto{\pgfqpoint{1.458396in}{0.957129in}}%
\pgfpathlineto{\pgfqpoint{1.458396in}{0.960078in}}%
\pgfpathlineto{\pgfqpoint{1.462937in}{0.960078in}}%
\pgfpathlineto{\pgfqpoint{1.462937in}{0.957129in}}%
\pgfpathmoveto{\pgfqpoint{1.458396in}{0.960078in}}%
\pgfpathlineto{\pgfqpoint{1.458396in}{0.960078in}}%
\pgfpathlineto{\pgfqpoint{1.458396in}{0.963027in}}%
\pgfpathlineto{\pgfqpoint{1.462937in}{0.963027in}}%
\pgfpathlineto{\pgfqpoint{1.462937in}{0.960078in}}%
\pgfpathmoveto{\pgfqpoint{1.462937in}{0.960078in}}%
\pgfpathlineto{\pgfqpoint{1.462937in}{0.960078in}}%
\pgfpathlineto{\pgfqpoint{1.462937in}{0.963027in}}%
\pgfpathlineto{\pgfqpoint{1.467478in}{0.963027in}}%
\pgfpathlineto{\pgfqpoint{1.467478in}{0.960078in}}%
\pgfpathmoveto{\pgfqpoint{1.462937in}{0.963027in}}%
\pgfpathlineto{\pgfqpoint{1.462937in}{0.963027in}}%
\pgfpathlineto{\pgfqpoint{1.462937in}{0.965976in}}%
\pgfpathlineto{\pgfqpoint{1.467478in}{0.965976in}}%
\pgfpathlineto{\pgfqpoint{1.467478in}{0.963027in}}%
\pgfpathmoveto{\pgfqpoint{1.467478in}{0.963027in}}%
\pgfpathlineto{\pgfqpoint{1.467478in}{0.963027in}}%
\pgfpathlineto{\pgfqpoint{1.467478in}{0.965976in}}%
\pgfpathlineto{\pgfqpoint{1.472019in}{0.965976in}}%
\pgfpathlineto{\pgfqpoint{1.472019in}{0.963027in}}%
\pgfpathmoveto{\pgfqpoint{1.467478in}{0.965976in}}%
\pgfpathlineto{\pgfqpoint{1.467478in}{0.965976in}}%
\pgfpathlineto{\pgfqpoint{1.467478in}{0.968925in}}%
\pgfpathlineto{\pgfqpoint{1.472019in}{0.968925in}}%
\pgfpathlineto{\pgfqpoint{1.472019in}{0.965976in}}%
\pgfpathmoveto{\pgfqpoint{1.472019in}{0.965976in}}%
\pgfpathlineto{\pgfqpoint{1.472019in}{0.965976in}}%
\pgfpathlineto{\pgfqpoint{1.472019in}{0.968925in}}%
\pgfpathlineto{\pgfqpoint{1.476560in}{0.968925in}}%
\pgfpathlineto{\pgfqpoint{1.476560in}{0.965976in}}%
\pgfpathmoveto{\pgfqpoint{1.472019in}{0.968925in}}%
\pgfpathlineto{\pgfqpoint{1.472019in}{0.968925in}}%
\pgfpathlineto{\pgfqpoint{1.472019in}{0.971875in}}%
\pgfpathlineto{\pgfqpoint{1.476560in}{0.971875in}}%
\pgfpathlineto{\pgfqpoint{1.476560in}{0.968925in}}%
\pgfpathmoveto{\pgfqpoint{1.476560in}{0.968925in}}%
\pgfpathlineto{\pgfqpoint{1.476560in}{0.968925in}}%
\pgfpathlineto{\pgfqpoint{1.476560in}{0.971875in}}%
\pgfpathlineto{\pgfqpoint{1.481101in}{0.971875in}}%
\pgfpathlineto{\pgfqpoint{1.481101in}{0.968925in}}%
\pgfpathmoveto{\pgfqpoint{1.476560in}{0.971875in}}%
\pgfpathlineto{\pgfqpoint{1.476560in}{0.971875in}}%
\pgfpathlineto{\pgfqpoint{1.476560in}{0.974824in}}%
\pgfpathlineto{\pgfqpoint{1.481101in}{0.974824in}}%
\pgfpathlineto{\pgfqpoint{1.481101in}{0.971875in}}%
\pgfpathmoveto{\pgfqpoint{1.481101in}{0.971875in}}%
\pgfpathlineto{\pgfqpoint{1.481101in}{0.971875in}}%
\pgfpathlineto{\pgfqpoint{1.481101in}{0.974824in}}%
\pgfpathlineto{\pgfqpoint{1.485642in}{0.974824in}}%
\pgfpathlineto{\pgfqpoint{1.485642in}{0.971875in}}%
\pgfpathmoveto{\pgfqpoint{1.481101in}{0.974824in}}%
\pgfpathlineto{\pgfqpoint{1.481101in}{0.974824in}}%
\pgfpathlineto{\pgfqpoint{1.481101in}{0.977773in}}%
\pgfpathlineto{\pgfqpoint{1.485642in}{0.977773in}}%
\pgfpathlineto{\pgfqpoint{1.485642in}{0.974824in}}%
\pgfpathmoveto{\pgfqpoint{1.485642in}{0.974824in}}%
\pgfpathlineto{\pgfqpoint{1.485642in}{0.974824in}}%
\pgfpathlineto{\pgfqpoint{1.485642in}{0.977773in}}%
\pgfpathlineto{\pgfqpoint{1.490183in}{0.977773in}}%
\pgfpathlineto{\pgfqpoint{1.490183in}{0.974824in}}%
\pgfpathmoveto{\pgfqpoint{1.485642in}{0.977773in}}%
\pgfpathlineto{\pgfqpoint{1.485642in}{0.977773in}}%
\pgfpathlineto{\pgfqpoint{1.485642in}{0.980722in}}%
\pgfpathlineto{\pgfqpoint{1.490183in}{0.980722in}}%
\pgfpathlineto{\pgfqpoint{1.490183in}{0.977773in}}%
\pgfpathmoveto{\pgfqpoint{1.490183in}{0.977773in}}%
\pgfpathlineto{\pgfqpoint{1.490183in}{0.977773in}}%
\pgfpathlineto{\pgfqpoint{1.490183in}{0.980722in}}%
\pgfpathlineto{\pgfqpoint{1.494724in}{0.980722in}}%
\pgfpathlineto{\pgfqpoint{1.494724in}{0.977773in}}%
\pgfpathmoveto{\pgfqpoint{1.490183in}{0.980722in}}%
\pgfpathlineto{\pgfqpoint{1.490183in}{0.980722in}}%
\pgfpathlineto{\pgfqpoint{1.490183in}{0.983671in}}%
\pgfpathlineto{\pgfqpoint{1.494724in}{0.983671in}}%
\pgfpathlineto{\pgfqpoint{1.494724in}{0.980722in}}%
\pgfpathmoveto{\pgfqpoint{1.494724in}{0.980722in}}%
\pgfpathlineto{\pgfqpoint{1.494724in}{0.980722in}}%
\pgfpathlineto{\pgfqpoint{1.494724in}{0.983671in}}%
\pgfpathlineto{\pgfqpoint{1.499265in}{0.983671in}}%
\pgfpathlineto{\pgfqpoint{1.499265in}{0.980722in}}%
\pgfpathmoveto{\pgfqpoint{1.494724in}{0.983671in}}%
\pgfpathlineto{\pgfqpoint{1.494724in}{0.983671in}}%
\pgfpathlineto{\pgfqpoint{1.494724in}{0.986620in}}%
\pgfpathlineto{\pgfqpoint{1.499265in}{0.986620in}}%
\pgfpathlineto{\pgfqpoint{1.499265in}{0.983671in}}%
\pgfpathmoveto{\pgfqpoint{1.499265in}{0.983671in}}%
\pgfpathlineto{\pgfqpoint{1.499265in}{0.983671in}}%
\pgfpathlineto{\pgfqpoint{1.499265in}{0.986620in}}%
\pgfpathlineto{\pgfqpoint{1.503806in}{0.986620in}}%
\pgfpathlineto{\pgfqpoint{1.503806in}{0.983671in}}%
\pgfpathmoveto{\pgfqpoint{1.499265in}{0.986620in}}%
\pgfpathlineto{\pgfqpoint{1.499265in}{0.986620in}}%
\pgfpathlineto{\pgfqpoint{1.499265in}{0.989570in}}%
\pgfpathlineto{\pgfqpoint{1.503806in}{0.989570in}}%
\pgfpathlineto{\pgfqpoint{1.503806in}{0.986620in}}%
\pgfpathmoveto{\pgfqpoint{1.503806in}{0.986620in}}%
\pgfpathlineto{\pgfqpoint{1.503806in}{0.986620in}}%
\pgfpathlineto{\pgfqpoint{1.503806in}{0.989570in}}%
\pgfpathlineto{\pgfqpoint{1.508347in}{0.989570in}}%
\pgfpathlineto{\pgfqpoint{1.508347in}{0.986620in}}%
\pgfpathmoveto{\pgfqpoint{1.503806in}{0.989570in}}%
\pgfpathlineto{\pgfqpoint{1.503806in}{0.989570in}}%
\pgfpathlineto{\pgfqpoint{1.503806in}{0.992519in}}%
\pgfpathlineto{\pgfqpoint{1.508347in}{0.992519in}}%
\pgfpathlineto{\pgfqpoint{1.508347in}{0.989570in}}%
\pgfpathmoveto{\pgfqpoint{1.508347in}{0.989570in}}%
\pgfpathlineto{\pgfqpoint{1.508347in}{0.989570in}}%
\pgfpathlineto{\pgfqpoint{1.508347in}{0.992519in}}%
\pgfpathlineto{\pgfqpoint{1.512888in}{0.992519in}}%
\pgfpathlineto{\pgfqpoint{1.512888in}{0.989570in}}%
\pgfpathmoveto{\pgfqpoint{1.508347in}{0.992519in}}%
\pgfpathlineto{\pgfqpoint{1.508347in}{0.992519in}}%
\pgfpathlineto{\pgfqpoint{1.508347in}{0.995468in}}%
\pgfpathlineto{\pgfqpoint{1.512888in}{0.995468in}}%
\pgfpathlineto{\pgfqpoint{1.512888in}{0.992519in}}%
\pgfpathmoveto{\pgfqpoint{1.512888in}{0.992519in}}%
\pgfpathlineto{\pgfqpoint{1.512888in}{0.992519in}}%
\pgfpathlineto{\pgfqpoint{1.512888in}{0.995468in}}%
\pgfpathlineto{\pgfqpoint{1.517429in}{0.995468in}}%
\pgfpathlineto{\pgfqpoint{1.517429in}{0.992519in}}%
\pgfpathmoveto{\pgfqpoint{1.512888in}{0.995468in}}%
\pgfpathlineto{\pgfqpoint{1.512888in}{0.995468in}}%
\pgfpathlineto{\pgfqpoint{1.512888in}{0.998417in}}%
\pgfpathlineto{\pgfqpoint{1.517429in}{0.998417in}}%
\pgfpathlineto{\pgfqpoint{1.517429in}{0.995468in}}%
\pgfpathmoveto{\pgfqpoint{1.517429in}{0.995468in}}%
\pgfpathlineto{\pgfqpoint{1.517429in}{0.995468in}}%
\pgfpathlineto{\pgfqpoint{1.517429in}{0.998417in}}%
\pgfpathlineto{\pgfqpoint{1.521970in}{0.998417in}}%
\pgfpathlineto{\pgfqpoint{1.521970in}{0.995468in}}%
\pgfpathmoveto{\pgfqpoint{1.517429in}{0.998417in}}%
\pgfpathlineto{\pgfqpoint{1.517429in}{0.998417in}}%
\pgfpathlineto{\pgfqpoint{1.517429in}{1.001366in}}%
\pgfpathlineto{\pgfqpoint{1.521970in}{1.001366in}}%
\pgfpathlineto{\pgfqpoint{1.521970in}{0.998417in}}%
\pgfpathmoveto{\pgfqpoint{1.521970in}{0.998417in}}%
\pgfpathlineto{\pgfqpoint{1.521970in}{0.998417in}}%
\pgfpathlineto{\pgfqpoint{1.521970in}{1.001366in}}%
\pgfpathlineto{\pgfqpoint{1.526511in}{1.001366in}}%
\pgfpathlineto{\pgfqpoint{1.526511in}{0.998417in}}%
\pgfpathmoveto{\pgfqpoint{1.521970in}{1.001366in}}%
\pgfpathlineto{\pgfqpoint{1.521970in}{1.001366in}}%
\pgfpathlineto{\pgfqpoint{1.521970in}{1.004315in}}%
\pgfpathlineto{\pgfqpoint{1.526511in}{1.004315in}}%
\pgfpathlineto{\pgfqpoint{1.526511in}{1.001366in}}%
\pgfpathmoveto{\pgfqpoint{1.526511in}{1.001366in}}%
\pgfpathlineto{\pgfqpoint{1.526511in}{1.001366in}}%
\pgfpathlineto{\pgfqpoint{1.526511in}{1.004315in}}%
\pgfpathlineto{\pgfqpoint{1.531052in}{1.004315in}}%
\pgfpathlineto{\pgfqpoint{1.531052in}{1.001366in}}%
\pgfpathmoveto{\pgfqpoint{1.526511in}{1.004315in}}%
\pgfpathlineto{\pgfqpoint{1.526511in}{1.004315in}}%
\pgfpathlineto{\pgfqpoint{1.526511in}{1.007265in}}%
\pgfpathlineto{\pgfqpoint{1.531052in}{1.007265in}}%
\pgfpathlineto{\pgfqpoint{1.531052in}{1.004315in}}%
\pgfpathmoveto{\pgfqpoint{1.531052in}{1.004315in}}%
\pgfpathlineto{\pgfqpoint{1.531052in}{1.004315in}}%
\pgfpathlineto{\pgfqpoint{1.531052in}{1.007265in}}%
\pgfpathlineto{\pgfqpoint{1.535593in}{1.007265in}}%
\pgfpathlineto{\pgfqpoint{1.535593in}{1.004315in}}%
\pgfpathmoveto{\pgfqpoint{1.531052in}{1.007265in}}%
\pgfpathlineto{\pgfqpoint{1.531052in}{1.007265in}}%
\pgfpathlineto{\pgfqpoint{1.531052in}{1.010214in}}%
\pgfpathlineto{\pgfqpoint{1.535593in}{1.010214in}}%
\pgfpathlineto{\pgfqpoint{1.535593in}{1.007265in}}%
\pgfpathmoveto{\pgfqpoint{1.535593in}{1.007265in}}%
\pgfpathlineto{\pgfqpoint{1.535593in}{1.007265in}}%
\pgfpathlineto{\pgfqpoint{1.535593in}{1.010214in}}%
\pgfpathlineto{\pgfqpoint{1.540134in}{1.010214in}}%
\pgfpathlineto{\pgfqpoint{1.540134in}{1.007265in}}%
\pgfpathmoveto{\pgfqpoint{1.535593in}{1.010214in}}%
\pgfpathlineto{\pgfqpoint{1.535593in}{1.010214in}}%
\pgfpathlineto{\pgfqpoint{1.535593in}{1.013163in}}%
\pgfpathlineto{\pgfqpoint{1.540134in}{1.013163in}}%
\pgfpathlineto{\pgfqpoint{1.540134in}{1.010214in}}%
\pgfpathmoveto{\pgfqpoint{1.540134in}{1.010214in}}%
\pgfpathlineto{\pgfqpoint{1.540134in}{1.010214in}}%
\pgfpathlineto{\pgfqpoint{1.540134in}{1.013163in}}%
\pgfpathlineto{\pgfqpoint{1.544675in}{1.013163in}}%
\pgfpathlineto{\pgfqpoint{1.544675in}{1.010214in}}%
\pgfpathmoveto{\pgfqpoint{1.540134in}{1.013163in}}%
\pgfpathlineto{\pgfqpoint{1.540134in}{1.013163in}}%
\pgfpathlineto{\pgfqpoint{1.540134in}{1.016112in}}%
\pgfpathlineto{\pgfqpoint{1.544675in}{1.016112in}}%
\pgfpathlineto{\pgfqpoint{1.544675in}{1.013163in}}%
\pgfpathmoveto{\pgfqpoint{1.544675in}{1.013163in}}%
\pgfpathlineto{\pgfqpoint{1.544675in}{1.013163in}}%
\pgfpathlineto{\pgfqpoint{1.544675in}{1.016112in}}%
\pgfpathlineto{\pgfqpoint{1.549216in}{1.016112in}}%
\pgfpathlineto{\pgfqpoint{1.549216in}{1.013163in}}%
\pgfpathmoveto{\pgfqpoint{1.544675in}{1.016112in}}%
\pgfpathlineto{\pgfqpoint{1.544675in}{1.016112in}}%
\pgfpathlineto{\pgfqpoint{1.544675in}{1.019061in}}%
\pgfpathlineto{\pgfqpoint{1.549216in}{1.019061in}}%
\pgfpathlineto{\pgfqpoint{1.549216in}{1.016112in}}%
\pgfpathmoveto{\pgfqpoint{1.549216in}{1.016112in}}%
\pgfpathlineto{\pgfqpoint{1.549216in}{1.016112in}}%
\pgfpathlineto{\pgfqpoint{1.549216in}{1.019061in}}%
\pgfpathlineto{\pgfqpoint{1.553757in}{1.019061in}}%
\pgfpathlineto{\pgfqpoint{1.553757in}{1.016112in}}%
\pgfpathmoveto{\pgfqpoint{1.549216in}{1.019061in}}%
\pgfpathlineto{\pgfqpoint{1.549216in}{1.019061in}}%
\pgfpathlineto{\pgfqpoint{1.549216in}{1.022011in}}%
\pgfpathlineto{\pgfqpoint{1.553757in}{1.022011in}}%
\pgfpathlineto{\pgfqpoint{1.553757in}{1.019061in}}%
\pgfpathmoveto{\pgfqpoint{1.553757in}{1.019061in}}%
\pgfpathlineto{\pgfqpoint{1.553757in}{1.019061in}}%
\pgfpathlineto{\pgfqpoint{1.553757in}{1.022011in}}%
\pgfpathlineto{\pgfqpoint{1.558298in}{1.022011in}}%
\pgfpathlineto{\pgfqpoint{1.558298in}{1.019061in}}%
\pgfpathmoveto{\pgfqpoint{1.553757in}{1.022011in}}%
\pgfpathlineto{\pgfqpoint{1.553757in}{1.022011in}}%
\pgfpathlineto{\pgfqpoint{1.553757in}{1.024960in}}%
\pgfpathlineto{\pgfqpoint{1.558298in}{1.024960in}}%
\pgfpathlineto{\pgfqpoint{1.558298in}{1.022011in}}%
\pgfpathmoveto{\pgfqpoint{1.558298in}{1.022011in}}%
\pgfpathlineto{\pgfqpoint{1.558298in}{1.022011in}}%
\pgfpathlineto{\pgfqpoint{1.558298in}{1.024960in}}%
\pgfpathlineto{\pgfqpoint{1.562839in}{1.024960in}}%
\pgfpathlineto{\pgfqpoint{1.562839in}{1.022011in}}%
\pgfpathmoveto{\pgfqpoint{1.558298in}{1.024960in}}%
\pgfpathlineto{\pgfqpoint{1.558298in}{1.024960in}}%
\pgfpathlineto{\pgfqpoint{1.558298in}{1.027909in}}%
\pgfpathlineto{\pgfqpoint{1.562839in}{1.027909in}}%
\pgfpathlineto{\pgfqpoint{1.562839in}{1.024960in}}%
\pgfpathmoveto{\pgfqpoint{1.562839in}{1.024960in}}%
\pgfpathlineto{\pgfqpoint{1.562839in}{1.024960in}}%
\pgfpathlineto{\pgfqpoint{1.562839in}{1.027909in}}%
\pgfpathlineto{\pgfqpoint{1.567380in}{1.027909in}}%
\pgfpathlineto{\pgfqpoint{1.567380in}{1.024960in}}%
\pgfpathmoveto{\pgfqpoint{1.562839in}{1.027909in}}%
\pgfpathlineto{\pgfqpoint{1.562839in}{1.027909in}}%
\pgfpathlineto{\pgfqpoint{1.562839in}{1.030858in}}%
\pgfpathlineto{\pgfqpoint{1.567380in}{1.030858in}}%
\pgfpathlineto{\pgfqpoint{1.567380in}{1.027909in}}%
\pgfpathmoveto{\pgfqpoint{1.567380in}{1.027909in}}%
\pgfpathlineto{\pgfqpoint{1.567380in}{1.027909in}}%
\pgfpathlineto{\pgfqpoint{1.567380in}{1.030858in}}%
\pgfpathlineto{\pgfqpoint{1.571921in}{1.030858in}}%
\pgfpathlineto{\pgfqpoint{1.571921in}{1.027909in}}%
\pgfpathmoveto{\pgfqpoint{1.567380in}{1.030858in}}%
\pgfpathlineto{\pgfqpoint{1.567380in}{1.030858in}}%
\pgfpathlineto{\pgfqpoint{1.567380in}{1.033807in}}%
\pgfpathlineto{\pgfqpoint{1.571921in}{1.033807in}}%
\pgfpathlineto{\pgfqpoint{1.571921in}{1.030858in}}%
\pgfpathmoveto{\pgfqpoint{1.571921in}{1.030858in}}%
\pgfpathlineto{\pgfqpoint{1.571921in}{1.030858in}}%
\pgfpathlineto{\pgfqpoint{1.571921in}{1.033807in}}%
\pgfpathlineto{\pgfqpoint{1.576462in}{1.033807in}}%
\pgfpathlineto{\pgfqpoint{1.576462in}{1.030858in}}%
\pgfpathmoveto{\pgfqpoint{1.571921in}{1.033807in}}%
\pgfpathlineto{\pgfqpoint{1.571921in}{1.033807in}}%
\pgfpathlineto{\pgfqpoint{1.571921in}{1.036756in}}%
\pgfpathlineto{\pgfqpoint{1.576462in}{1.036756in}}%
\pgfpathlineto{\pgfqpoint{1.576462in}{1.033807in}}%
\pgfpathmoveto{\pgfqpoint{1.576462in}{1.033807in}}%
\pgfpathlineto{\pgfqpoint{1.576462in}{1.033807in}}%
\pgfpathlineto{\pgfqpoint{1.576462in}{1.036756in}}%
\pgfpathlineto{\pgfqpoint{1.581003in}{1.036756in}}%
\pgfpathlineto{\pgfqpoint{1.581003in}{1.033807in}}%
\pgfpathmoveto{\pgfqpoint{1.576462in}{1.036756in}}%
\pgfpathlineto{\pgfqpoint{1.576462in}{1.036756in}}%
\pgfpathlineto{\pgfqpoint{1.576462in}{1.039706in}}%
\pgfpathlineto{\pgfqpoint{1.581003in}{1.039706in}}%
\pgfpathlineto{\pgfqpoint{1.581003in}{1.036756in}}%
\pgfpathmoveto{\pgfqpoint{1.581003in}{1.036756in}}%
\pgfpathlineto{\pgfqpoint{1.581003in}{1.036756in}}%
\pgfpathlineto{\pgfqpoint{1.581003in}{1.039706in}}%
\pgfpathlineto{\pgfqpoint{1.585544in}{1.039706in}}%
\pgfpathlineto{\pgfqpoint{1.585544in}{1.036756in}}%
\pgfpathmoveto{\pgfqpoint{1.581003in}{1.039706in}}%
\pgfpathlineto{\pgfqpoint{1.581003in}{1.039706in}}%
\pgfpathlineto{\pgfqpoint{1.581003in}{1.042655in}}%
\pgfpathlineto{\pgfqpoint{1.585544in}{1.042655in}}%
\pgfpathlineto{\pgfqpoint{1.585544in}{1.039706in}}%
\pgfpathmoveto{\pgfqpoint{1.585544in}{1.039706in}}%
\pgfpathlineto{\pgfqpoint{1.585544in}{1.039706in}}%
\pgfpathlineto{\pgfqpoint{1.585544in}{1.042655in}}%
\pgfpathlineto{\pgfqpoint{1.590085in}{1.042655in}}%
\pgfpathlineto{\pgfqpoint{1.590085in}{1.039706in}}%
\pgfpathmoveto{\pgfqpoint{1.585544in}{1.042655in}}%
\pgfpathlineto{\pgfqpoint{1.585544in}{1.042655in}}%
\pgfpathlineto{\pgfqpoint{1.585544in}{1.045604in}}%
\pgfpathlineto{\pgfqpoint{1.590085in}{1.045604in}}%
\pgfpathlineto{\pgfqpoint{1.590085in}{1.042655in}}%
\pgfpathmoveto{\pgfqpoint{1.590085in}{1.042655in}}%
\pgfpathlineto{\pgfqpoint{1.590085in}{1.042655in}}%
\pgfpathlineto{\pgfqpoint{1.590085in}{1.045604in}}%
\pgfpathlineto{\pgfqpoint{1.594626in}{1.045604in}}%
\pgfpathlineto{\pgfqpoint{1.594626in}{1.042655in}}%
\pgfpathmoveto{\pgfqpoint{1.590085in}{1.045604in}}%
\pgfpathlineto{\pgfqpoint{1.590085in}{1.045604in}}%
\pgfpathlineto{\pgfqpoint{1.590085in}{1.048553in}}%
\pgfpathlineto{\pgfqpoint{1.594626in}{1.048553in}}%
\pgfpathlineto{\pgfqpoint{1.594626in}{1.045604in}}%
\pgfpathmoveto{\pgfqpoint{1.594626in}{1.045604in}}%
\pgfpathlineto{\pgfqpoint{1.594626in}{1.045604in}}%
\pgfpathlineto{\pgfqpoint{1.594626in}{1.048553in}}%
\pgfpathlineto{\pgfqpoint{1.599167in}{1.048553in}}%
\pgfpathlineto{\pgfqpoint{1.599167in}{1.045604in}}%
\pgfpathmoveto{\pgfqpoint{1.594626in}{1.048553in}}%
\pgfpathlineto{\pgfqpoint{1.594626in}{1.048553in}}%
\pgfpathlineto{\pgfqpoint{1.594626in}{1.051502in}}%
\pgfpathlineto{\pgfqpoint{1.599167in}{1.051502in}}%
\pgfpathlineto{\pgfqpoint{1.599167in}{1.048553in}}%
\pgfpathmoveto{\pgfqpoint{1.599167in}{1.048553in}}%
\pgfpathlineto{\pgfqpoint{1.599167in}{1.048553in}}%
\pgfpathlineto{\pgfqpoint{1.599167in}{1.051502in}}%
\pgfpathlineto{\pgfqpoint{1.603708in}{1.051502in}}%
\pgfpathlineto{\pgfqpoint{1.603708in}{1.048553in}}%
\pgfpathmoveto{\pgfqpoint{1.599167in}{1.051502in}}%
\pgfpathlineto{\pgfqpoint{1.599167in}{1.051502in}}%
\pgfpathlineto{\pgfqpoint{1.599167in}{1.054451in}}%
\pgfpathlineto{\pgfqpoint{1.603708in}{1.054451in}}%
\pgfpathlineto{\pgfqpoint{1.603708in}{1.051502in}}%
\pgfpathmoveto{\pgfqpoint{1.603708in}{1.051502in}}%
\pgfpathlineto{\pgfqpoint{1.603708in}{1.051502in}}%
\pgfpathlineto{\pgfqpoint{1.603708in}{1.054451in}}%
\pgfpathlineto{\pgfqpoint{1.608249in}{1.054451in}}%
\pgfpathlineto{\pgfqpoint{1.608249in}{1.051502in}}%
\pgfpathmoveto{\pgfqpoint{1.603708in}{1.054451in}}%
\pgfpathlineto{\pgfqpoint{1.603708in}{1.054451in}}%
\pgfpathlineto{\pgfqpoint{1.603708in}{1.057401in}}%
\pgfpathlineto{\pgfqpoint{1.608249in}{1.057401in}}%
\pgfpathlineto{\pgfqpoint{1.608249in}{1.054451in}}%
\pgfpathmoveto{\pgfqpoint{1.608249in}{1.054451in}}%
\pgfpathlineto{\pgfqpoint{1.608249in}{1.054451in}}%
\pgfpathlineto{\pgfqpoint{1.608249in}{1.057401in}}%
\pgfpathlineto{\pgfqpoint{1.612790in}{1.057401in}}%
\pgfpathlineto{\pgfqpoint{1.612790in}{1.054451in}}%
\pgfpathmoveto{\pgfqpoint{1.608249in}{1.057401in}}%
\pgfpathlineto{\pgfqpoint{1.608249in}{1.057401in}}%
\pgfpathlineto{\pgfqpoint{1.608249in}{1.060350in}}%
\pgfpathlineto{\pgfqpoint{1.612790in}{1.060350in}}%
\pgfpathlineto{\pgfqpoint{1.612790in}{1.057401in}}%
\pgfpathmoveto{\pgfqpoint{1.612790in}{1.057401in}}%
\pgfpathlineto{\pgfqpoint{1.612790in}{1.057401in}}%
\pgfpathlineto{\pgfqpoint{1.612790in}{1.060350in}}%
\pgfpathlineto{\pgfqpoint{1.617331in}{1.060350in}}%
\pgfpathlineto{\pgfqpoint{1.617331in}{1.057401in}}%
\pgfpathmoveto{\pgfqpoint{1.612790in}{1.060350in}}%
\pgfpathlineto{\pgfqpoint{1.612790in}{1.060350in}}%
\pgfpathlineto{\pgfqpoint{1.612790in}{1.063299in}}%
\pgfpathlineto{\pgfqpoint{1.617331in}{1.063299in}}%
\pgfpathlineto{\pgfqpoint{1.617331in}{1.060350in}}%
\pgfpathmoveto{\pgfqpoint{1.617331in}{1.060350in}}%
\pgfpathlineto{\pgfqpoint{1.617331in}{1.060350in}}%
\pgfpathlineto{\pgfqpoint{1.617331in}{1.063299in}}%
\pgfpathlineto{\pgfqpoint{1.621872in}{1.063299in}}%
\pgfpathlineto{\pgfqpoint{1.621872in}{1.060350in}}%
\pgfpathmoveto{\pgfqpoint{1.617331in}{1.063299in}}%
\pgfpathlineto{\pgfqpoint{1.617331in}{1.063299in}}%
\pgfpathlineto{\pgfqpoint{1.617331in}{1.066248in}}%
\pgfpathlineto{\pgfqpoint{1.621872in}{1.066248in}}%
\pgfpathlineto{\pgfqpoint{1.621872in}{1.063299in}}%
\pgfpathmoveto{\pgfqpoint{1.621872in}{1.063299in}}%
\pgfpathlineto{\pgfqpoint{1.621872in}{1.063299in}}%
\pgfpathlineto{\pgfqpoint{1.621872in}{1.066248in}}%
\pgfpathlineto{\pgfqpoint{1.626413in}{1.066248in}}%
\pgfpathlineto{\pgfqpoint{1.626413in}{1.063299in}}%
\pgfpathmoveto{\pgfqpoint{1.621872in}{1.066248in}}%
\pgfpathlineto{\pgfqpoint{1.621872in}{1.066248in}}%
\pgfpathlineto{\pgfqpoint{1.621872in}{1.069197in}}%
\pgfpathlineto{\pgfqpoint{1.626413in}{1.069197in}}%
\pgfpathlineto{\pgfqpoint{1.626413in}{1.066248in}}%
\pgfpathmoveto{\pgfqpoint{1.626413in}{1.066248in}}%
\pgfpathlineto{\pgfqpoint{1.626413in}{1.066248in}}%
\pgfpathlineto{\pgfqpoint{1.626413in}{1.069197in}}%
\pgfpathlineto{\pgfqpoint{1.630954in}{1.069197in}}%
\pgfpathlineto{\pgfqpoint{1.630954in}{1.066248in}}%
\pgfpathmoveto{\pgfqpoint{1.626413in}{1.069197in}}%
\pgfpathlineto{\pgfqpoint{1.626413in}{1.069197in}}%
\pgfpathlineto{\pgfqpoint{1.626413in}{1.072147in}}%
\pgfpathlineto{\pgfqpoint{1.630954in}{1.072147in}}%
\pgfpathlineto{\pgfqpoint{1.630954in}{1.069197in}}%
\pgfpathmoveto{\pgfqpoint{1.630954in}{1.069197in}}%
\pgfpathlineto{\pgfqpoint{1.630954in}{1.069197in}}%
\pgfpathlineto{\pgfqpoint{1.630954in}{1.072147in}}%
\pgfpathlineto{\pgfqpoint{1.635495in}{1.072147in}}%
\pgfpathlineto{\pgfqpoint{1.635495in}{1.069197in}}%
\pgfpathmoveto{\pgfqpoint{1.630954in}{1.072147in}}%
\pgfpathlineto{\pgfqpoint{1.630954in}{1.072147in}}%
\pgfpathlineto{\pgfqpoint{1.630954in}{1.075096in}}%
\pgfpathlineto{\pgfqpoint{1.635495in}{1.075096in}}%
\pgfpathlineto{\pgfqpoint{1.635495in}{1.072147in}}%
\pgfpathmoveto{\pgfqpoint{1.635495in}{1.072147in}}%
\pgfpathlineto{\pgfqpoint{1.635495in}{1.072147in}}%
\pgfpathlineto{\pgfqpoint{1.635495in}{1.075096in}}%
\pgfpathlineto{\pgfqpoint{1.640036in}{1.075096in}}%
\pgfpathlineto{\pgfqpoint{1.640036in}{1.072147in}}%
\pgfpathmoveto{\pgfqpoint{1.635495in}{1.075096in}}%
\pgfpathlineto{\pgfqpoint{1.635495in}{1.075096in}}%
\pgfpathlineto{\pgfqpoint{1.635495in}{1.078045in}}%
\pgfpathlineto{\pgfqpoint{1.640036in}{1.078045in}}%
\pgfpathlineto{\pgfqpoint{1.640036in}{1.075096in}}%
\pgfpathmoveto{\pgfqpoint{1.640036in}{1.075096in}}%
\pgfpathlineto{\pgfqpoint{1.640036in}{1.075096in}}%
\pgfpathlineto{\pgfqpoint{1.640036in}{1.078045in}}%
\pgfpathlineto{\pgfqpoint{1.644577in}{1.078045in}}%
\pgfpathlineto{\pgfqpoint{1.644577in}{1.075096in}}%
\pgfpathmoveto{\pgfqpoint{1.640036in}{1.078045in}}%
\pgfpathlineto{\pgfqpoint{1.640036in}{1.078045in}}%
\pgfpathlineto{\pgfqpoint{1.640036in}{1.080995in}}%
\pgfpathlineto{\pgfqpoint{1.644577in}{1.080995in}}%
\pgfpathlineto{\pgfqpoint{1.644577in}{1.078045in}}%
\pgfpathmoveto{\pgfqpoint{1.644577in}{1.078045in}}%
\pgfpathlineto{\pgfqpoint{1.644577in}{1.078045in}}%
\pgfpathlineto{\pgfqpoint{1.644577in}{1.080995in}}%
\pgfpathlineto{\pgfqpoint{1.649118in}{1.080995in}}%
\pgfpathlineto{\pgfqpoint{1.649118in}{1.078045in}}%
\pgfpathmoveto{\pgfqpoint{1.644577in}{1.080995in}}%
\pgfpathlineto{\pgfqpoint{1.644577in}{1.080995in}}%
\pgfpathlineto{\pgfqpoint{1.644577in}{1.083944in}}%
\pgfpathlineto{\pgfqpoint{1.649118in}{1.083944in}}%
\pgfpathlineto{\pgfqpoint{1.649118in}{1.080995in}}%
\pgfpathmoveto{\pgfqpoint{1.649118in}{1.080995in}}%
\pgfpathlineto{\pgfqpoint{1.649118in}{1.080995in}}%
\pgfpathlineto{\pgfqpoint{1.649118in}{1.083944in}}%
\pgfpathlineto{\pgfqpoint{1.653659in}{1.083944in}}%
\pgfpathlineto{\pgfqpoint{1.653659in}{1.080995in}}%
\pgfpathmoveto{\pgfqpoint{1.649118in}{1.083944in}}%
\pgfpathlineto{\pgfqpoint{1.649118in}{1.083944in}}%
\pgfpathlineto{\pgfqpoint{1.649118in}{1.086893in}}%
\pgfpathlineto{\pgfqpoint{1.653659in}{1.086893in}}%
\pgfpathlineto{\pgfqpoint{1.653659in}{1.083944in}}%
\pgfpathmoveto{\pgfqpoint{1.653659in}{1.083944in}}%
\pgfpathlineto{\pgfqpoint{1.653659in}{1.083944in}}%
\pgfpathlineto{\pgfqpoint{1.653659in}{1.086893in}}%
\pgfpathlineto{\pgfqpoint{1.658200in}{1.086893in}}%
\pgfpathlineto{\pgfqpoint{1.658200in}{1.083944in}}%
\pgfpathmoveto{\pgfqpoint{1.653659in}{1.086893in}}%
\pgfpathlineto{\pgfqpoint{1.653659in}{1.086893in}}%
\pgfpathlineto{\pgfqpoint{1.653659in}{1.089842in}}%
\pgfpathlineto{\pgfqpoint{1.658200in}{1.089842in}}%
\pgfpathlineto{\pgfqpoint{1.658200in}{1.086893in}}%
\pgfpathmoveto{\pgfqpoint{1.658200in}{1.086893in}}%
\pgfpathlineto{\pgfqpoint{1.658200in}{1.086893in}}%
\pgfpathlineto{\pgfqpoint{1.658200in}{1.089842in}}%
\pgfpathlineto{\pgfqpoint{1.662741in}{1.089842in}}%
\pgfpathlineto{\pgfqpoint{1.662741in}{1.086893in}}%
\pgfpathmoveto{\pgfqpoint{1.658200in}{1.089842in}}%
\pgfpathlineto{\pgfqpoint{1.658200in}{1.089842in}}%
\pgfpathlineto{\pgfqpoint{1.658200in}{1.092792in}}%
\pgfpathlineto{\pgfqpoint{1.662741in}{1.092792in}}%
\pgfpathlineto{\pgfqpoint{1.662741in}{1.089842in}}%
\pgfpathmoveto{\pgfqpoint{1.662741in}{1.089842in}}%
\pgfpathlineto{\pgfqpoint{1.662741in}{1.089842in}}%
\pgfpathlineto{\pgfqpoint{1.662741in}{1.092792in}}%
\pgfpathlineto{\pgfqpoint{1.667282in}{1.092792in}}%
\pgfpathlineto{\pgfqpoint{1.667282in}{1.089842in}}%
\pgfpathmoveto{\pgfqpoint{1.662741in}{1.092792in}}%
\pgfpathlineto{\pgfqpoint{1.662741in}{1.092792in}}%
\pgfpathlineto{\pgfqpoint{1.662741in}{1.095741in}}%
\pgfpathlineto{\pgfqpoint{1.667282in}{1.095741in}}%
\pgfpathlineto{\pgfqpoint{1.667282in}{1.092792in}}%
\pgfpathmoveto{\pgfqpoint{1.667282in}{1.092792in}}%
\pgfpathlineto{\pgfqpoint{1.667282in}{1.092792in}}%
\pgfpathlineto{\pgfqpoint{1.667282in}{1.095741in}}%
\pgfpathlineto{\pgfqpoint{1.671823in}{1.095741in}}%
\pgfpathlineto{\pgfqpoint{1.671823in}{1.092792in}}%
\pgfpathmoveto{\pgfqpoint{1.667282in}{1.095741in}}%
\pgfpathlineto{\pgfqpoint{1.667282in}{1.095741in}}%
\pgfpathlineto{\pgfqpoint{1.667282in}{1.098690in}}%
\pgfpathlineto{\pgfqpoint{1.671823in}{1.098690in}}%
\pgfpathlineto{\pgfqpoint{1.671823in}{1.095741in}}%
\pgfpathmoveto{\pgfqpoint{1.671823in}{1.095741in}}%
\pgfpathlineto{\pgfqpoint{1.671823in}{1.095741in}}%
\pgfpathlineto{\pgfqpoint{1.671823in}{1.098690in}}%
\pgfpathlineto{\pgfqpoint{1.676364in}{1.098690in}}%
\pgfpathlineto{\pgfqpoint{1.676364in}{1.095741in}}%
\pgfpathmoveto{\pgfqpoint{1.671823in}{1.098690in}}%
\pgfpathlineto{\pgfqpoint{1.671823in}{1.098690in}}%
\pgfpathlineto{\pgfqpoint{1.671823in}{1.101640in}}%
\pgfpathlineto{\pgfqpoint{1.676364in}{1.101640in}}%
\pgfpathlineto{\pgfqpoint{1.676364in}{1.098690in}}%
\pgfpathmoveto{\pgfqpoint{1.676364in}{1.098690in}}%
\pgfpathlineto{\pgfqpoint{1.676364in}{1.098690in}}%
\pgfpathlineto{\pgfqpoint{1.676364in}{1.101640in}}%
\pgfpathlineto{\pgfqpoint{1.680905in}{1.101640in}}%
\pgfpathlineto{\pgfqpoint{1.680905in}{1.098690in}}%
\pgfpathmoveto{\pgfqpoint{1.676364in}{1.101640in}}%
\pgfpathlineto{\pgfqpoint{1.676364in}{1.101640in}}%
\pgfpathlineto{\pgfqpoint{1.676364in}{1.104589in}}%
\pgfpathlineto{\pgfqpoint{1.680905in}{1.104589in}}%
\pgfpathlineto{\pgfqpoint{1.680905in}{1.101640in}}%
\pgfpathmoveto{\pgfqpoint{1.680905in}{1.101640in}}%
\pgfpathlineto{\pgfqpoint{1.680905in}{1.101640in}}%
\pgfpathlineto{\pgfqpoint{1.680905in}{1.104589in}}%
\pgfpathlineto{\pgfqpoint{1.685446in}{1.104589in}}%
\pgfpathlineto{\pgfqpoint{1.685446in}{1.101640in}}%
\pgfpathmoveto{\pgfqpoint{1.680905in}{1.104589in}}%
\pgfpathlineto{\pgfqpoint{1.680905in}{1.104589in}}%
\pgfpathlineto{\pgfqpoint{1.680905in}{1.107538in}}%
\pgfpathlineto{\pgfqpoint{1.685446in}{1.107538in}}%
\pgfpathlineto{\pgfqpoint{1.685446in}{1.104589in}}%
\pgfpathmoveto{\pgfqpoint{1.685446in}{1.104589in}}%
\pgfpathlineto{\pgfqpoint{1.685446in}{1.104589in}}%
\pgfpathlineto{\pgfqpoint{1.685446in}{1.107538in}}%
\pgfpathlineto{\pgfqpoint{1.689987in}{1.107538in}}%
\pgfpathlineto{\pgfqpoint{1.689987in}{1.104589in}}%
\pgfpathmoveto{\pgfqpoint{1.685446in}{1.107538in}}%
\pgfpathlineto{\pgfqpoint{1.685446in}{1.107538in}}%
\pgfpathlineto{\pgfqpoint{1.685446in}{1.110488in}}%
\pgfpathlineto{\pgfqpoint{1.689987in}{1.110488in}}%
\pgfpathlineto{\pgfqpoint{1.689987in}{1.107538in}}%
\pgfpathmoveto{\pgfqpoint{1.689987in}{1.107538in}}%
\pgfpathlineto{\pgfqpoint{1.689987in}{1.107538in}}%
\pgfpathlineto{\pgfqpoint{1.689987in}{1.110488in}}%
\pgfpathlineto{\pgfqpoint{1.694527in}{1.110488in}}%
\pgfpathlineto{\pgfqpoint{1.694527in}{1.107538in}}%
\pgfpathmoveto{\pgfqpoint{1.689987in}{1.110488in}}%
\pgfpathlineto{\pgfqpoint{1.689987in}{1.110488in}}%
\pgfpathlineto{\pgfqpoint{1.689987in}{1.113437in}}%
\pgfpathlineto{\pgfqpoint{1.694527in}{1.113437in}}%
\pgfpathlineto{\pgfqpoint{1.694527in}{1.110488in}}%
\pgfpathmoveto{\pgfqpoint{1.694527in}{1.110488in}}%
\pgfpathlineto{\pgfqpoint{1.694527in}{1.110488in}}%
\pgfpathlineto{\pgfqpoint{1.694527in}{1.113437in}}%
\pgfpathlineto{\pgfqpoint{1.699068in}{1.113437in}}%
\pgfpathlineto{\pgfqpoint{1.699068in}{1.110488in}}%
\pgfpathmoveto{\pgfqpoint{1.694527in}{1.113437in}}%
\pgfpathlineto{\pgfqpoint{1.694527in}{1.113437in}}%
\pgfpathlineto{\pgfqpoint{1.694527in}{1.116386in}}%
\pgfpathlineto{\pgfqpoint{1.699068in}{1.116386in}}%
\pgfpathlineto{\pgfqpoint{1.699068in}{1.113437in}}%
\pgfpathmoveto{\pgfqpoint{1.699068in}{1.113437in}}%
\pgfpathlineto{\pgfqpoint{1.699068in}{1.113437in}}%
\pgfpathlineto{\pgfqpoint{1.699068in}{1.116386in}}%
\pgfpathlineto{\pgfqpoint{1.703609in}{1.116386in}}%
\pgfpathlineto{\pgfqpoint{1.703609in}{1.113437in}}%
\pgfpathmoveto{\pgfqpoint{1.699068in}{1.116386in}}%
\pgfpathlineto{\pgfqpoint{1.699068in}{1.116386in}}%
\pgfpathlineto{\pgfqpoint{1.699068in}{1.119335in}}%
\pgfpathlineto{\pgfqpoint{1.703609in}{1.119335in}}%
\pgfpathlineto{\pgfqpoint{1.703609in}{1.116386in}}%
\pgfpathmoveto{\pgfqpoint{1.703609in}{1.116386in}}%
\pgfpathlineto{\pgfqpoint{1.703609in}{1.116386in}}%
\pgfpathlineto{\pgfqpoint{1.703609in}{1.119335in}}%
\pgfpathlineto{\pgfqpoint{1.708150in}{1.119335in}}%
\pgfpathlineto{\pgfqpoint{1.708150in}{1.116386in}}%
\pgfpathmoveto{\pgfqpoint{1.703609in}{1.119335in}}%
\pgfpathlineto{\pgfqpoint{1.703609in}{1.119335in}}%
\pgfpathlineto{\pgfqpoint{1.703609in}{1.122285in}}%
\pgfpathlineto{\pgfqpoint{1.708150in}{1.122285in}}%
\pgfpathlineto{\pgfqpoint{1.708150in}{1.119335in}}%
\pgfpathmoveto{\pgfqpoint{1.708150in}{1.119335in}}%
\pgfpathlineto{\pgfqpoint{1.708150in}{1.119335in}}%
\pgfpathlineto{\pgfqpoint{1.708150in}{1.122285in}}%
\pgfpathlineto{\pgfqpoint{1.712691in}{1.122285in}}%
\pgfpathlineto{\pgfqpoint{1.712691in}{1.119335in}}%
\pgfpathmoveto{\pgfqpoint{1.708150in}{1.122285in}}%
\pgfpathlineto{\pgfqpoint{1.708150in}{1.122285in}}%
\pgfpathlineto{\pgfqpoint{1.708150in}{1.125234in}}%
\pgfpathlineto{\pgfqpoint{1.712691in}{1.125234in}}%
\pgfpathlineto{\pgfqpoint{1.712691in}{1.122285in}}%
\pgfpathmoveto{\pgfqpoint{1.712691in}{1.122285in}}%
\pgfpathlineto{\pgfqpoint{1.712691in}{1.122285in}}%
\pgfpathlineto{\pgfqpoint{1.712691in}{1.125234in}}%
\pgfpathlineto{\pgfqpoint{1.717232in}{1.125234in}}%
\pgfpathlineto{\pgfqpoint{1.717232in}{1.122285in}}%
\pgfpathmoveto{\pgfqpoint{1.712691in}{1.125234in}}%
\pgfpathlineto{\pgfqpoint{1.712691in}{1.125234in}}%
\pgfpathlineto{\pgfqpoint{1.712691in}{1.128183in}}%
\pgfpathlineto{\pgfqpoint{1.717232in}{1.128183in}}%
\pgfpathlineto{\pgfqpoint{1.717232in}{1.125234in}}%
\pgfpathmoveto{\pgfqpoint{1.717232in}{1.125234in}}%
\pgfpathlineto{\pgfqpoint{1.717232in}{1.125234in}}%
\pgfpathlineto{\pgfqpoint{1.717232in}{1.128183in}}%
\pgfpathlineto{\pgfqpoint{1.721773in}{1.128183in}}%
\pgfpathlineto{\pgfqpoint{1.721773in}{1.125234in}}%
\pgfpathmoveto{\pgfqpoint{1.717232in}{1.128183in}}%
\pgfpathlineto{\pgfqpoint{1.717232in}{1.128183in}}%
\pgfpathlineto{\pgfqpoint{1.717232in}{1.131133in}}%
\pgfpathlineto{\pgfqpoint{1.721773in}{1.131133in}}%
\pgfpathlineto{\pgfqpoint{1.721773in}{1.128183in}}%
\pgfpathmoveto{\pgfqpoint{1.721773in}{1.128183in}}%
\pgfpathlineto{\pgfqpoint{1.721773in}{1.128183in}}%
\pgfpathlineto{\pgfqpoint{1.721773in}{1.131133in}}%
\pgfpathlineto{\pgfqpoint{1.726314in}{1.131133in}}%
\pgfpathlineto{\pgfqpoint{1.726314in}{1.128183in}}%
\pgfpathmoveto{\pgfqpoint{1.721773in}{1.131133in}}%
\pgfpathlineto{\pgfqpoint{1.721773in}{1.131133in}}%
\pgfpathlineto{\pgfqpoint{1.721773in}{1.134082in}}%
\pgfpathlineto{\pgfqpoint{1.726314in}{1.134082in}}%
\pgfpathlineto{\pgfqpoint{1.726314in}{1.131133in}}%
\pgfpathmoveto{\pgfqpoint{1.726314in}{1.131133in}}%
\pgfpathlineto{\pgfqpoint{1.726314in}{1.131133in}}%
\pgfpathlineto{\pgfqpoint{1.726314in}{1.134082in}}%
\pgfpathlineto{\pgfqpoint{1.730855in}{1.134082in}}%
\pgfpathlineto{\pgfqpoint{1.730855in}{1.131133in}}%
\pgfpathmoveto{\pgfqpoint{1.726314in}{1.134082in}}%
\pgfpathlineto{\pgfqpoint{1.726314in}{1.134082in}}%
\pgfpathlineto{\pgfqpoint{1.726314in}{1.137031in}}%
\pgfpathlineto{\pgfqpoint{1.730855in}{1.137031in}}%
\pgfpathlineto{\pgfqpoint{1.730855in}{1.134082in}}%
\pgfpathmoveto{\pgfqpoint{1.730855in}{1.134082in}}%
\pgfpathlineto{\pgfqpoint{1.730855in}{1.134082in}}%
\pgfpathlineto{\pgfqpoint{1.730855in}{1.137031in}}%
\pgfpathlineto{\pgfqpoint{1.735396in}{1.137031in}}%
\pgfpathlineto{\pgfqpoint{1.735396in}{1.134082in}}%
\pgfpathmoveto{\pgfqpoint{1.730855in}{1.137031in}}%
\pgfpathlineto{\pgfqpoint{1.730855in}{1.137031in}}%
\pgfpathlineto{\pgfqpoint{1.730855in}{1.139981in}}%
\pgfpathlineto{\pgfqpoint{1.735396in}{1.139981in}}%
\pgfpathlineto{\pgfqpoint{1.735396in}{1.137031in}}%
\pgfpathmoveto{\pgfqpoint{1.735396in}{1.137031in}}%
\pgfpathlineto{\pgfqpoint{1.735396in}{1.137031in}}%
\pgfpathlineto{\pgfqpoint{1.735396in}{1.139981in}}%
\pgfpathlineto{\pgfqpoint{1.739937in}{1.139981in}}%
\pgfpathlineto{\pgfqpoint{1.739937in}{1.137031in}}%
\pgfpathmoveto{\pgfqpoint{1.735396in}{1.139981in}}%
\pgfpathlineto{\pgfqpoint{1.735396in}{1.139981in}}%
\pgfpathlineto{\pgfqpoint{1.735396in}{1.142930in}}%
\pgfpathlineto{\pgfqpoint{1.739937in}{1.142930in}}%
\pgfpathlineto{\pgfqpoint{1.739937in}{1.139981in}}%
\pgfpathmoveto{\pgfqpoint{1.739937in}{1.139981in}}%
\pgfpathlineto{\pgfqpoint{1.739937in}{1.139981in}}%
\pgfpathlineto{\pgfqpoint{1.739937in}{1.142930in}}%
\pgfpathlineto{\pgfqpoint{1.744478in}{1.142930in}}%
\pgfpathlineto{\pgfqpoint{1.744478in}{1.139981in}}%
\pgfpathmoveto{\pgfqpoint{1.739937in}{1.142930in}}%
\pgfpathlineto{\pgfqpoint{1.739937in}{1.142930in}}%
\pgfpathlineto{\pgfqpoint{1.739937in}{1.145879in}}%
\pgfpathlineto{\pgfqpoint{1.744478in}{1.145879in}}%
\pgfpathlineto{\pgfqpoint{1.744478in}{1.142930in}}%
\pgfpathmoveto{\pgfqpoint{1.744478in}{1.142930in}}%
\pgfpathlineto{\pgfqpoint{1.744478in}{1.142930in}}%
\pgfpathlineto{\pgfqpoint{1.744478in}{1.145879in}}%
\pgfpathlineto{\pgfqpoint{1.749019in}{1.145879in}}%
\pgfpathlineto{\pgfqpoint{1.749019in}{1.142930in}}%
\pgfpathmoveto{\pgfqpoint{1.744478in}{1.145879in}}%
\pgfpathlineto{\pgfqpoint{1.744478in}{1.145879in}}%
\pgfpathlineto{\pgfqpoint{1.744478in}{1.148828in}}%
\pgfpathlineto{\pgfqpoint{1.749019in}{1.148828in}}%
\pgfpathlineto{\pgfqpoint{1.749019in}{1.145879in}}%
\pgfpathmoveto{\pgfqpoint{1.749019in}{1.145879in}}%
\pgfpathlineto{\pgfqpoint{1.749019in}{1.145879in}}%
\pgfpathlineto{\pgfqpoint{1.749019in}{1.148828in}}%
\pgfpathlineto{\pgfqpoint{1.753560in}{1.148828in}}%
\pgfpathlineto{\pgfqpoint{1.753560in}{1.145879in}}%
\pgfpathmoveto{\pgfqpoint{1.749019in}{1.148828in}}%
\pgfpathlineto{\pgfqpoint{1.749019in}{1.148828in}}%
\pgfpathlineto{\pgfqpoint{1.749019in}{1.151778in}}%
\pgfpathlineto{\pgfqpoint{1.753560in}{1.151778in}}%
\pgfpathlineto{\pgfqpoint{1.753560in}{1.148828in}}%
\pgfpathmoveto{\pgfqpoint{1.753560in}{1.148828in}}%
\pgfpathlineto{\pgfqpoint{1.753560in}{1.148828in}}%
\pgfpathlineto{\pgfqpoint{1.753560in}{1.151778in}}%
\pgfpathlineto{\pgfqpoint{1.758101in}{1.151778in}}%
\pgfpathlineto{\pgfqpoint{1.758101in}{1.148828in}}%
\pgfpathmoveto{\pgfqpoint{1.753560in}{1.151778in}}%
\pgfpathlineto{\pgfqpoint{1.753560in}{1.151778in}}%
\pgfpathlineto{\pgfqpoint{1.753560in}{1.154727in}}%
\pgfpathlineto{\pgfqpoint{1.758101in}{1.154727in}}%
\pgfpathlineto{\pgfqpoint{1.758101in}{1.151778in}}%
\pgfpathmoveto{\pgfqpoint{1.758101in}{1.151778in}}%
\pgfpathlineto{\pgfqpoint{1.758101in}{1.151778in}}%
\pgfpathlineto{\pgfqpoint{1.758101in}{1.154727in}}%
\pgfpathlineto{\pgfqpoint{1.762642in}{1.154727in}}%
\pgfpathlineto{\pgfqpoint{1.762642in}{1.151778in}}%
\pgfpathmoveto{\pgfqpoint{1.758101in}{1.154727in}}%
\pgfpathlineto{\pgfqpoint{1.758101in}{1.154727in}}%
\pgfpathlineto{\pgfqpoint{1.758101in}{1.157676in}}%
\pgfpathlineto{\pgfqpoint{1.762642in}{1.157676in}}%
\pgfpathlineto{\pgfqpoint{1.762642in}{1.154727in}}%
\pgfpathmoveto{\pgfqpoint{1.762642in}{1.154727in}}%
\pgfpathlineto{\pgfqpoint{1.762642in}{1.154727in}}%
\pgfpathlineto{\pgfqpoint{1.762642in}{1.157676in}}%
\pgfpathlineto{\pgfqpoint{1.767183in}{1.157676in}}%
\pgfpathlineto{\pgfqpoint{1.767183in}{1.154727in}}%
\pgfpathmoveto{\pgfqpoint{1.762642in}{1.157676in}}%
\pgfpathlineto{\pgfqpoint{1.762642in}{1.157676in}}%
\pgfpathlineto{\pgfqpoint{1.762642in}{1.160626in}}%
\pgfpathlineto{\pgfqpoint{1.767183in}{1.160626in}}%
\pgfpathlineto{\pgfqpoint{1.767183in}{1.157676in}}%
\pgfpathmoveto{\pgfqpoint{1.767183in}{1.157676in}}%
\pgfpathlineto{\pgfqpoint{1.767183in}{1.157676in}}%
\pgfpathlineto{\pgfqpoint{1.767183in}{1.160626in}}%
\pgfpathlineto{\pgfqpoint{1.771724in}{1.160626in}}%
\pgfpathlineto{\pgfqpoint{1.771724in}{1.157676in}}%
\pgfpathmoveto{\pgfqpoint{1.767183in}{1.160626in}}%
\pgfpathlineto{\pgfqpoint{1.767183in}{1.160626in}}%
\pgfpathlineto{\pgfqpoint{1.767183in}{1.163575in}}%
\pgfpathlineto{\pgfqpoint{1.771724in}{1.163575in}}%
\pgfpathlineto{\pgfqpoint{1.771724in}{1.160626in}}%
\pgfpathmoveto{\pgfqpoint{1.771724in}{1.160626in}}%
\pgfpathlineto{\pgfqpoint{1.771724in}{1.160626in}}%
\pgfpathlineto{\pgfqpoint{1.771724in}{1.163575in}}%
\pgfpathlineto{\pgfqpoint{1.776266in}{1.163575in}}%
\pgfpathlineto{\pgfqpoint{1.776266in}{1.160626in}}%
\pgfpathmoveto{\pgfqpoint{1.771724in}{1.163575in}}%
\pgfpathlineto{\pgfqpoint{1.771724in}{1.163575in}}%
\pgfpathlineto{\pgfqpoint{1.771724in}{1.166524in}}%
\pgfpathlineto{\pgfqpoint{1.776266in}{1.166524in}}%
\pgfpathlineto{\pgfqpoint{1.776266in}{1.163575in}}%
\pgfpathmoveto{\pgfqpoint{1.776266in}{1.163575in}}%
\pgfpathlineto{\pgfqpoint{1.776266in}{1.163575in}}%
\pgfpathlineto{\pgfqpoint{1.776266in}{1.166524in}}%
\pgfpathlineto{\pgfqpoint{1.780807in}{1.166524in}}%
\pgfpathlineto{\pgfqpoint{1.780807in}{1.163575in}}%
\pgfpathmoveto{\pgfqpoint{1.776266in}{1.166524in}}%
\pgfpathlineto{\pgfqpoint{1.776266in}{1.166524in}}%
\pgfpathlineto{\pgfqpoint{1.776266in}{1.169473in}}%
\pgfpathlineto{\pgfqpoint{1.780807in}{1.169473in}}%
\pgfpathlineto{\pgfqpoint{1.780807in}{1.166524in}}%
\pgfpathmoveto{\pgfqpoint{1.780807in}{1.166524in}}%
\pgfpathlineto{\pgfqpoint{1.780807in}{1.166524in}}%
\pgfpathlineto{\pgfqpoint{1.780807in}{1.169473in}}%
\pgfpathlineto{\pgfqpoint{1.785348in}{1.169473in}}%
\pgfpathlineto{\pgfqpoint{1.785348in}{1.166524in}}%
\pgfpathmoveto{\pgfqpoint{1.780807in}{1.169473in}}%
\pgfpathlineto{\pgfqpoint{1.780807in}{1.169473in}}%
\pgfpathlineto{\pgfqpoint{1.780807in}{1.172422in}}%
\pgfpathlineto{\pgfqpoint{1.785348in}{1.172422in}}%
\pgfpathlineto{\pgfqpoint{1.785348in}{1.169473in}}%
\pgfpathmoveto{\pgfqpoint{1.785348in}{1.169473in}}%
\pgfpathlineto{\pgfqpoint{1.785348in}{1.169473in}}%
\pgfpathlineto{\pgfqpoint{1.785348in}{1.172422in}}%
\pgfpathlineto{\pgfqpoint{1.789889in}{1.172422in}}%
\pgfpathlineto{\pgfqpoint{1.789889in}{1.169473in}}%
\pgfpathmoveto{\pgfqpoint{1.785348in}{1.172422in}}%
\pgfpathlineto{\pgfqpoint{1.785348in}{1.172422in}}%
\pgfpathlineto{\pgfqpoint{1.785348in}{1.175371in}}%
\pgfpathlineto{\pgfqpoint{1.789889in}{1.175371in}}%
\pgfpathlineto{\pgfqpoint{1.789889in}{1.172422in}}%
\pgfpathmoveto{\pgfqpoint{1.789889in}{1.172422in}}%
\pgfpathlineto{\pgfqpoint{1.789889in}{1.172422in}}%
\pgfpathlineto{\pgfqpoint{1.789889in}{1.175371in}}%
\pgfpathlineto{\pgfqpoint{1.794430in}{1.175371in}}%
\pgfpathlineto{\pgfqpoint{1.794430in}{1.172422in}}%
\pgfpathmoveto{\pgfqpoint{1.789889in}{1.175371in}}%
\pgfpathlineto{\pgfqpoint{1.789889in}{1.175371in}}%
\pgfpathlineto{\pgfqpoint{1.789889in}{1.178320in}}%
\pgfpathlineto{\pgfqpoint{1.794430in}{1.178320in}}%
\pgfpathlineto{\pgfqpoint{1.794430in}{1.175371in}}%
\pgfpathmoveto{\pgfqpoint{1.794430in}{1.175371in}}%
\pgfpathlineto{\pgfqpoint{1.794430in}{1.175371in}}%
\pgfpathlineto{\pgfqpoint{1.794430in}{1.178320in}}%
\pgfpathlineto{\pgfqpoint{1.798971in}{1.178320in}}%
\pgfpathlineto{\pgfqpoint{1.798971in}{1.175371in}}%
\pgfpathmoveto{\pgfqpoint{1.794430in}{1.178320in}}%
\pgfpathlineto{\pgfqpoint{1.794430in}{1.178320in}}%
\pgfpathlineto{\pgfqpoint{1.794430in}{1.181270in}}%
\pgfpathlineto{\pgfqpoint{1.798971in}{1.181270in}}%
\pgfpathlineto{\pgfqpoint{1.798971in}{1.178320in}}%
\pgfpathmoveto{\pgfqpoint{1.798971in}{1.178320in}}%
\pgfpathlineto{\pgfqpoint{1.798971in}{1.178320in}}%
\pgfpathlineto{\pgfqpoint{1.798971in}{1.181270in}}%
\pgfpathlineto{\pgfqpoint{1.803512in}{1.181270in}}%
\pgfpathlineto{\pgfqpoint{1.803512in}{1.178320in}}%
\pgfpathmoveto{\pgfqpoint{1.798971in}{1.181270in}}%
\pgfpathlineto{\pgfqpoint{1.798971in}{1.181270in}}%
\pgfpathlineto{\pgfqpoint{1.798971in}{1.184219in}}%
\pgfpathlineto{\pgfqpoint{1.803512in}{1.184219in}}%
\pgfpathlineto{\pgfqpoint{1.803512in}{1.181270in}}%
\pgfpathmoveto{\pgfqpoint{1.803512in}{1.181270in}}%
\pgfpathlineto{\pgfqpoint{1.803512in}{1.181270in}}%
\pgfpathlineto{\pgfqpoint{1.803512in}{1.184219in}}%
\pgfpathlineto{\pgfqpoint{1.808053in}{1.184219in}}%
\pgfpathlineto{\pgfqpoint{1.808053in}{1.181270in}}%
\pgfpathmoveto{\pgfqpoint{1.803512in}{1.184219in}}%
\pgfpathlineto{\pgfqpoint{1.803512in}{1.184219in}}%
\pgfpathlineto{\pgfqpoint{1.803512in}{1.187168in}}%
\pgfpathlineto{\pgfqpoint{1.808053in}{1.187168in}}%
\pgfpathlineto{\pgfqpoint{1.808053in}{1.184219in}}%
\pgfpathmoveto{\pgfqpoint{1.808053in}{1.184219in}}%
\pgfpathlineto{\pgfqpoint{1.808053in}{1.184219in}}%
\pgfpathlineto{\pgfqpoint{1.808053in}{1.187168in}}%
\pgfpathlineto{\pgfqpoint{1.812595in}{1.187168in}}%
\pgfpathlineto{\pgfqpoint{1.812595in}{1.184219in}}%
\pgfpathmoveto{\pgfqpoint{1.808053in}{1.187168in}}%
\pgfpathlineto{\pgfqpoint{1.808053in}{1.187168in}}%
\pgfpathlineto{\pgfqpoint{1.808053in}{1.190117in}}%
\pgfpathlineto{\pgfqpoint{1.812595in}{1.190117in}}%
\pgfpathlineto{\pgfqpoint{1.812595in}{1.187168in}}%
\pgfpathmoveto{\pgfqpoint{1.812595in}{1.187168in}}%
\pgfpathlineto{\pgfqpoint{1.812595in}{1.187168in}}%
\pgfpathlineto{\pgfqpoint{1.812595in}{1.190117in}}%
\pgfpathlineto{\pgfqpoint{1.817136in}{1.190117in}}%
\pgfpathlineto{\pgfqpoint{1.817136in}{1.187168in}}%
\pgfpathmoveto{\pgfqpoint{1.812595in}{1.190117in}}%
\pgfpathlineto{\pgfqpoint{1.812595in}{1.190117in}}%
\pgfpathlineto{\pgfqpoint{1.812595in}{1.193066in}}%
\pgfpathlineto{\pgfqpoint{1.817136in}{1.193066in}}%
\pgfpathlineto{\pgfqpoint{1.817136in}{1.190117in}}%
\pgfpathmoveto{\pgfqpoint{1.817136in}{1.190117in}}%
\pgfpathlineto{\pgfqpoint{1.817136in}{1.190117in}}%
\pgfpathlineto{\pgfqpoint{1.817136in}{1.193066in}}%
\pgfpathlineto{\pgfqpoint{1.821677in}{1.193066in}}%
\pgfpathlineto{\pgfqpoint{1.821677in}{1.190117in}}%
\pgfpathmoveto{\pgfqpoint{1.817136in}{1.193066in}}%
\pgfpathlineto{\pgfqpoint{1.817136in}{1.193066in}}%
\pgfpathlineto{\pgfqpoint{1.817136in}{1.196015in}}%
\pgfpathlineto{\pgfqpoint{1.821677in}{1.196015in}}%
\pgfpathlineto{\pgfqpoint{1.821677in}{1.193066in}}%
\pgfpathmoveto{\pgfqpoint{1.821677in}{1.193066in}}%
\pgfpathlineto{\pgfqpoint{1.821677in}{1.193066in}}%
\pgfpathlineto{\pgfqpoint{1.821677in}{1.196015in}}%
\pgfpathlineto{\pgfqpoint{1.826218in}{1.196015in}}%
\pgfpathlineto{\pgfqpoint{1.826218in}{1.193066in}}%
\pgfpathmoveto{\pgfqpoint{1.821677in}{1.196015in}}%
\pgfpathlineto{\pgfqpoint{1.821677in}{1.196015in}}%
\pgfpathlineto{\pgfqpoint{1.821677in}{1.198964in}}%
\pgfpathlineto{\pgfqpoint{1.826218in}{1.198964in}}%
\pgfpathlineto{\pgfqpoint{1.826218in}{1.196015in}}%
\pgfpathmoveto{\pgfqpoint{1.826218in}{1.196015in}}%
\pgfpathlineto{\pgfqpoint{1.826218in}{1.196015in}}%
\pgfpathlineto{\pgfqpoint{1.826218in}{1.198964in}}%
\pgfpathlineto{\pgfqpoint{1.830759in}{1.198964in}}%
\pgfpathlineto{\pgfqpoint{1.830759in}{1.196015in}}%
\pgfpathmoveto{\pgfqpoint{1.826218in}{1.198964in}}%
\pgfpathlineto{\pgfqpoint{1.826218in}{1.198964in}}%
\pgfpathlineto{\pgfqpoint{1.826218in}{1.201913in}}%
\pgfpathlineto{\pgfqpoint{1.830759in}{1.201913in}}%
\pgfpathlineto{\pgfqpoint{1.830759in}{1.198964in}}%
\pgfpathmoveto{\pgfqpoint{1.830759in}{1.198964in}}%
\pgfpathlineto{\pgfqpoint{1.830759in}{1.198964in}}%
\pgfpathlineto{\pgfqpoint{1.830759in}{1.201913in}}%
\pgfpathlineto{\pgfqpoint{1.835300in}{1.201913in}}%
\pgfpathlineto{\pgfqpoint{1.835300in}{1.198964in}}%
\pgfpathmoveto{\pgfqpoint{1.830759in}{1.201913in}}%
\pgfpathlineto{\pgfqpoint{1.830759in}{1.201913in}}%
\pgfpathlineto{\pgfqpoint{1.830759in}{1.204863in}}%
\pgfpathlineto{\pgfqpoint{1.835300in}{1.204863in}}%
\pgfpathlineto{\pgfqpoint{1.835300in}{1.201913in}}%
\pgfpathmoveto{\pgfqpoint{1.835300in}{1.201913in}}%
\pgfpathlineto{\pgfqpoint{1.835300in}{1.201913in}}%
\pgfpathlineto{\pgfqpoint{1.835300in}{1.204863in}}%
\pgfpathlineto{\pgfqpoint{1.839841in}{1.204863in}}%
\pgfpathlineto{\pgfqpoint{1.839841in}{1.201913in}}%
\pgfpathmoveto{\pgfqpoint{1.835300in}{1.204863in}}%
\pgfpathlineto{\pgfqpoint{1.835300in}{1.204863in}}%
\pgfpathlineto{\pgfqpoint{1.835300in}{1.207812in}}%
\pgfpathlineto{\pgfqpoint{1.839841in}{1.207812in}}%
\pgfpathlineto{\pgfqpoint{1.839841in}{1.204863in}}%
\pgfpathmoveto{\pgfqpoint{1.839841in}{1.204863in}}%
\pgfpathlineto{\pgfqpoint{1.839841in}{1.204863in}}%
\pgfpathlineto{\pgfqpoint{1.839841in}{1.207812in}}%
\pgfpathlineto{\pgfqpoint{1.844382in}{1.207812in}}%
\pgfpathlineto{\pgfqpoint{1.844382in}{1.204863in}}%
\pgfpathmoveto{\pgfqpoint{1.839841in}{1.207812in}}%
\pgfpathlineto{\pgfqpoint{1.839841in}{1.207812in}}%
\pgfpathlineto{\pgfqpoint{1.839841in}{1.210761in}}%
\pgfpathlineto{\pgfqpoint{1.844382in}{1.210761in}}%
\pgfpathlineto{\pgfqpoint{1.844382in}{1.207812in}}%
\pgfpathmoveto{\pgfqpoint{1.844382in}{1.207812in}}%
\pgfpathlineto{\pgfqpoint{1.844382in}{1.207812in}}%
\pgfpathlineto{\pgfqpoint{1.844382in}{1.210761in}}%
\pgfpathlineto{\pgfqpoint{1.848924in}{1.210761in}}%
\pgfpathlineto{\pgfqpoint{1.848924in}{1.207812in}}%
\pgfpathmoveto{\pgfqpoint{1.844382in}{1.210761in}}%
\pgfpathlineto{\pgfqpoint{1.844382in}{1.210761in}}%
\pgfpathlineto{\pgfqpoint{1.844382in}{1.213710in}}%
\pgfpathlineto{\pgfqpoint{1.848924in}{1.213710in}}%
\pgfpathlineto{\pgfqpoint{1.848924in}{1.210761in}}%
\pgfpathmoveto{\pgfqpoint{1.848924in}{1.210761in}}%
\pgfpathlineto{\pgfqpoint{1.848924in}{1.210761in}}%
\pgfpathlineto{\pgfqpoint{1.848924in}{1.213710in}}%
\pgfpathlineto{\pgfqpoint{1.853465in}{1.213710in}}%
\pgfpathlineto{\pgfqpoint{1.853465in}{1.210761in}}%
\pgfpathmoveto{\pgfqpoint{1.848924in}{1.213710in}}%
\pgfpathlineto{\pgfqpoint{1.848924in}{1.213710in}}%
\pgfpathlineto{\pgfqpoint{1.848924in}{1.216659in}}%
\pgfpathlineto{\pgfqpoint{1.853465in}{1.216659in}}%
\pgfpathlineto{\pgfqpoint{1.853465in}{1.213710in}}%
\pgfpathmoveto{\pgfqpoint{1.853465in}{1.213710in}}%
\pgfpathlineto{\pgfqpoint{1.853465in}{1.213710in}}%
\pgfpathlineto{\pgfqpoint{1.853465in}{1.216659in}}%
\pgfpathlineto{\pgfqpoint{1.858006in}{1.216659in}}%
\pgfpathlineto{\pgfqpoint{1.858006in}{1.213710in}}%
\pgfpathmoveto{\pgfqpoint{1.853465in}{1.216659in}}%
\pgfpathlineto{\pgfqpoint{1.853465in}{1.216659in}}%
\pgfpathlineto{\pgfqpoint{1.853465in}{1.219608in}}%
\pgfpathlineto{\pgfqpoint{1.858006in}{1.219608in}}%
\pgfpathlineto{\pgfqpoint{1.858006in}{1.216659in}}%
\pgfpathmoveto{\pgfqpoint{1.858006in}{1.216659in}}%
\pgfpathlineto{\pgfqpoint{1.858006in}{1.216659in}}%
\pgfpathlineto{\pgfqpoint{1.858006in}{1.219608in}}%
\pgfpathlineto{\pgfqpoint{1.862547in}{1.219608in}}%
\pgfpathlineto{\pgfqpoint{1.862547in}{1.216659in}}%
\pgfpathmoveto{\pgfqpoint{1.858006in}{1.219608in}}%
\pgfpathlineto{\pgfqpoint{1.858006in}{1.219608in}}%
\pgfpathlineto{\pgfqpoint{1.858006in}{1.222557in}}%
\pgfpathlineto{\pgfqpoint{1.862547in}{1.222557in}}%
\pgfpathlineto{\pgfqpoint{1.862547in}{1.219608in}}%
\pgfpathmoveto{\pgfqpoint{1.858006in}{1.222557in}}%
\pgfpathlineto{\pgfqpoint{1.858006in}{1.222557in}}%
\pgfpathlineto{\pgfqpoint{1.858006in}{1.225506in}}%
\pgfpathlineto{\pgfqpoint{1.862547in}{1.225506in}}%
\pgfpathlineto{\pgfqpoint{1.862547in}{1.222557in}}%
\pgfpathmoveto{\pgfqpoint{1.862547in}{1.222557in}}%
\pgfpathlineto{\pgfqpoint{1.862547in}{1.222557in}}%
\pgfpathlineto{\pgfqpoint{1.862547in}{1.225506in}}%
\pgfpathlineto{\pgfqpoint{1.867088in}{1.225506in}}%
\pgfpathlineto{\pgfqpoint{1.867088in}{1.222557in}}%
\pgfpathmoveto{\pgfqpoint{1.862547in}{1.225506in}}%
\pgfpathlineto{\pgfqpoint{1.862547in}{1.225506in}}%
\pgfpathlineto{\pgfqpoint{1.862547in}{1.228456in}}%
\pgfpathlineto{\pgfqpoint{1.867088in}{1.228456in}}%
\pgfpathlineto{\pgfqpoint{1.867088in}{1.225506in}}%
\pgfpathmoveto{\pgfqpoint{1.867088in}{1.225506in}}%
\pgfpathlineto{\pgfqpoint{1.867088in}{1.225506in}}%
\pgfpathlineto{\pgfqpoint{1.867088in}{1.228456in}}%
\pgfpathlineto{\pgfqpoint{1.871629in}{1.228456in}}%
\pgfpathlineto{\pgfqpoint{1.871629in}{1.225506in}}%
\pgfpathmoveto{\pgfqpoint{1.867088in}{1.228456in}}%
\pgfpathlineto{\pgfqpoint{1.867088in}{1.228456in}}%
\pgfpathlineto{\pgfqpoint{1.867088in}{1.231405in}}%
\pgfpathlineto{\pgfqpoint{1.871629in}{1.231405in}}%
\pgfpathlineto{\pgfqpoint{1.871629in}{1.228456in}}%
\pgfpathmoveto{\pgfqpoint{1.871629in}{1.228456in}}%
\pgfpathlineto{\pgfqpoint{1.871629in}{1.228456in}}%
\pgfpathlineto{\pgfqpoint{1.871629in}{1.231405in}}%
\pgfpathlineto{\pgfqpoint{1.876170in}{1.231405in}}%
\pgfpathlineto{\pgfqpoint{1.876170in}{1.228456in}}%
\pgfpathmoveto{\pgfqpoint{1.871629in}{1.231405in}}%
\pgfpathlineto{\pgfqpoint{1.871629in}{1.231405in}}%
\pgfpathlineto{\pgfqpoint{1.871629in}{1.234354in}}%
\pgfpathlineto{\pgfqpoint{1.876170in}{1.234354in}}%
\pgfpathlineto{\pgfqpoint{1.876170in}{1.231405in}}%
\pgfpathmoveto{\pgfqpoint{1.876170in}{1.231405in}}%
\pgfpathlineto{\pgfqpoint{1.876170in}{1.231405in}}%
\pgfpathlineto{\pgfqpoint{1.876170in}{1.234354in}}%
\pgfpathlineto{\pgfqpoint{1.880712in}{1.234354in}}%
\pgfpathlineto{\pgfqpoint{1.880712in}{1.231405in}}%
\pgfpathmoveto{\pgfqpoint{1.876170in}{1.234354in}}%
\pgfpathlineto{\pgfqpoint{1.876170in}{1.234354in}}%
\pgfpathlineto{\pgfqpoint{1.876170in}{1.237303in}}%
\pgfpathlineto{\pgfqpoint{1.880712in}{1.237303in}}%
\pgfpathlineto{\pgfqpoint{1.880712in}{1.234354in}}%
\pgfpathmoveto{\pgfqpoint{1.880712in}{1.234354in}}%
\pgfpathlineto{\pgfqpoint{1.880712in}{1.234354in}}%
\pgfpathlineto{\pgfqpoint{1.880712in}{1.237303in}}%
\pgfpathlineto{\pgfqpoint{1.885253in}{1.237303in}}%
\pgfpathlineto{\pgfqpoint{1.885253in}{1.234354in}}%
\pgfpathmoveto{\pgfqpoint{1.880712in}{1.237303in}}%
\pgfpathlineto{\pgfqpoint{1.880712in}{1.237303in}}%
\pgfpathlineto{\pgfqpoint{1.880712in}{1.240252in}}%
\pgfpathlineto{\pgfqpoint{1.885253in}{1.240252in}}%
\pgfpathlineto{\pgfqpoint{1.885253in}{1.237303in}}%
\pgfpathmoveto{\pgfqpoint{1.885253in}{1.237303in}}%
\pgfpathlineto{\pgfqpoint{1.885253in}{1.237303in}}%
\pgfpathlineto{\pgfqpoint{1.885253in}{1.240252in}}%
\pgfpathlineto{\pgfqpoint{1.889794in}{1.240252in}}%
\pgfpathlineto{\pgfqpoint{1.889794in}{1.237303in}}%
\pgfpathmoveto{\pgfqpoint{1.885253in}{1.240252in}}%
\pgfpathlineto{\pgfqpoint{1.885253in}{1.240252in}}%
\pgfpathlineto{\pgfqpoint{1.885253in}{1.243201in}}%
\pgfpathlineto{\pgfqpoint{1.889794in}{1.243201in}}%
\pgfpathlineto{\pgfqpoint{1.889794in}{1.240252in}}%
\pgfpathmoveto{\pgfqpoint{1.889794in}{1.240252in}}%
\pgfpathlineto{\pgfqpoint{1.889794in}{1.240252in}}%
\pgfpathlineto{\pgfqpoint{1.889794in}{1.243201in}}%
\pgfpathlineto{\pgfqpoint{1.894335in}{1.243201in}}%
\pgfpathlineto{\pgfqpoint{1.894335in}{1.240252in}}%
\pgfpathmoveto{\pgfqpoint{1.889794in}{1.243201in}}%
\pgfpathlineto{\pgfqpoint{1.889794in}{1.243201in}}%
\pgfpathlineto{\pgfqpoint{1.889794in}{1.246150in}}%
\pgfpathlineto{\pgfqpoint{1.894335in}{1.246150in}}%
\pgfpathlineto{\pgfqpoint{1.894335in}{1.243201in}}%
\pgfpathmoveto{\pgfqpoint{1.894335in}{1.243201in}}%
\pgfpathlineto{\pgfqpoint{1.894335in}{1.243201in}}%
\pgfpathlineto{\pgfqpoint{1.894335in}{1.246150in}}%
\pgfpathlineto{\pgfqpoint{1.898876in}{1.246150in}}%
\pgfpathlineto{\pgfqpoint{1.898876in}{1.243201in}}%
\pgfpathmoveto{\pgfqpoint{1.894335in}{1.246150in}}%
\pgfpathlineto{\pgfqpoint{1.894335in}{1.246150in}}%
\pgfpathlineto{\pgfqpoint{1.894335in}{1.249099in}}%
\pgfpathlineto{\pgfqpoint{1.898876in}{1.249099in}}%
\pgfpathlineto{\pgfqpoint{1.898876in}{1.246150in}}%
\pgfpathmoveto{\pgfqpoint{1.898876in}{1.246150in}}%
\pgfpathlineto{\pgfqpoint{1.898876in}{1.246150in}}%
\pgfpathlineto{\pgfqpoint{1.898876in}{1.249099in}}%
\pgfpathlineto{\pgfqpoint{1.903417in}{1.249099in}}%
\pgfpathlineto{\pgfqpoint{1.903417in}{1.246150in}}%
\pgfpathmoveto{\pgfqpoint{1.898876in}{1.249099in}}%
\pgfpathlineto{\pgfqpoint{1.898876in}{1.249099in}}%
\pgfpathlineto{\pgfqpoint{1.898876in}{1.252049in}}%
\pgfpathlineto{\pgfqpoint{1.903417in}{1.252049in}}%
\pgfpathlineto{\pgfqpoint{1.903417in}{1.249099in}}%
\pgfpathmoveto{\pgfqpoint{1.903417in}{1.249099in}}%
\pgfpathlineto{\pgfqpoint{1.903417in}{1.249099in}}%
\pgfpathlineto{\pgfqpoint{1.903417in}{1.252049in}}%
\pgfpathlineto{\pgfqpoint{1.907958in}{1.252049in}}%
\pgfpathlineto{\pgfqpoint{1.907958in}{1.249099in}}%
\pgfpathmoveto{\pgfqpoint{1.903417in}{1.252049in}}%
\pgfpathlineto{\pgfqpoint{1.903417in}{1.252049in}}%
\pgfpathlineto{\pgfqpoint{1.903417in}{1.254998in}}%
\pgfpathlineto{\pgfqpoint{1.907958in}{1.254998in}}%
\pgfpathlineto{\pgfqpoint{1.907958in}{1.252049in}}%
\pgfpathmoveto{\pgfqpoint{1.907958in}{1.252049in}}%
\pgfpathlineto{\pgfqpoint{1.907958in}{1.252049in}}%
\pgfpathlineto{\pgfqpoint{1.907958in}{1.254998in}}%
\pgfpathlineto{\pgfqpoint{1.912499in}{1.254998in}}%
\pgfpathlineto{\pgfqpoint{1.912499in}{1.252049in}}%
\pgfpathmoveto{\pgfqpoint{1.907958in}{1.254998in}}%
\pgfpathlineto{\pgfqpoint{1.907958in}{1.254998in}}%
\pgfpathlineto{\pgfqpoint{1.907958in}{1.257947in}}%
\pgfpathlineto{\pgfqpoint{1.912499in}{1.257947in}}%
\pgfpathlineto{\pgfqpoint{1.912499in}{1.254998in}}%
\pgfpathmoveto{\pgfqpoint{1.912499in}{1.254998in}}%
\pgfpathlineto{\pgfqpoint{1.912499in}{1.254998in}}%
\pgfpathlineto{\pgfqpoint{1.912499in}{1.257947in}}%
\pgfpathlineto{\pgfqpoint{1.917040in}{1.257947in}}%
\pgfpathlineto{\pgfqpoint{1.917040in}{1.254998in}}%
\pgfpathmoveto{\pgfqpoint{1.912499in}{1.257947in}}%
\pgfpathlineto{\pgfqpoint{1.912499in}{1.257947in}}%
\pgfpathlineto{\pgfqpoint{1.912499in}{1.260896in}}%
\pgfpathlineto{\pgfqpoint{1.917040in}{1.260896in}}%
\pgfpathlineto{\pgfqpoint{1.917040in}{1.257947in}}%
\pgfpathmoveto{\pgfqpoint{1.917040in}{1.257947in}}%
\pgfpathlineto{\pgfqpoint{1.917040in}{1.257947in}}%
\pgfpathlineto{\pgfqpoint{1.917040in}{1.260896in}}%
\pgfpathlineto{\pgfqpoint{1.921581in}{1.260896in}}%
\pgfpathlineto{\pgfqpoint{1.921581in}{1.257947in}}%
\pgfpathmoveto{\pgfqpoint{1.917040in}{1.260896in}}%
\pgfpathlineto{\pgfqpoint{1.917040in}{1.260896in}}%
\pgfpathlineto{\pgfqpoint{1.917040in}{1.263846in}}%
\pgfpathlineto{\pgfqpoint{1.921581in}{1.263846in}}%
\pgfpathlineto{\pgfqpoint{1.921581in}{1.260896in}}%
\pgfpathmoveto{\pgfqpoint{1.921581in}{1.260896in}}%
\pgfpathlineto{\pgfqpoint{1.921581in}{1.260896in}}%
\pgfpathlineto{\pgfqpoint{1.921581in}{1.263846in}}%
\pgfpathlineto{\pgfqpoint{1.926122in}{1.263846in}}%
\pgfpathlineto{\pgfqpoint{1.926122in}{1.260896in}}%
\pgfpathmoveto{\pgfqpoint{1.921581in}{1.263846in}}%
\pgfpathlineto{\pgfqpoint{1.921581in}{1.263846in}}%
\pgfpathlineto{\pgfqpoint{1.921581in}{1.266795in}}%
\pgfpathlineto{\pgfqpoint{1.926122in}{1.266795in}}%
\pgfpathlineto{\pgfqpoint{1.926122in}{1.263846in}}%
\pgfpathmoveto{\pgfqpoint{1.926122in}{1.263846in}}%
\pgfpathlineto{\pgfqpoint{1.926122in}{1.263846in}}%
\pgfpathlineto{\pgfqpoint{1.926122in}{1.266795in}}%
\pgfpathlineto{\pgfqpoint{1.930663in}{1.266795in}}%
\pgfpathlineto{\pgfqpoint{1.930663in}{1.263846in}}%
\pgfpathmoveto{\pgfqpoint{1.926122in}{1.266795in}}%
\pgfpathlineto{\pgfqpoint{1.926122in}{1.266795in}}%
\pgfpathlineto{\pgfqpoint{1.926122in}{1.269744in}}%
\pgfpathlineto{\pgfqpoint{1.930663in}{1.269744in}}%
\pgfpathlineto{\pgfqpoint{1.930663in}{1.266795in}}%
\pgfpathmoveto{\pgfqpoint{1.930663in}{1.266795in}}%
\pgfpathlineto{\pgfqpoint{1.930663in}{1.266795in}}%
\pgfpathlineto{\pgfqpoint{1.930663in}{1.269744in}}%
\pgfpathlineto{\pgfqpoint{1.935204in}{1.269744in}}%
\pgfpathlineto{\pgfqpoint{1.935204in}{1.266795in}}%
\pgfpathmoveto{\pgfqpoint{1.930663in}{1.269744in}}%
\pgfpathlineto{\pgfqpoint{1.930663in}{1.269744in}}%
\pgfpathlineto{\pgfqpoint{1.930663in}{1.272694in}}%
\pgfpathlineto{\pgfqpoint{1.935204in}{1.272694in}}%
\pgfpathlineto{\pgfqpoint{1.935204in}{1.269744in}}%
\pgfpathmoveto{\pgfqpoint{1.935204in}{1.269744in}}%
\pgfpathlineto{\pgfqpoint{1.935204in}{1.269744in}}%
\pgfpathlineto{\pgfqpoint{1.935204in}{1.272694in}}%
\pgfpathlineto{\pgfqpoint{1.939745in}{1.272694in}}%
\pgfpathlineto{\pgfqpoint{1.939745in}{1.269744in}}%
\pgfpathmoveto{\pgfqpoint{1.935204in}{1.272694in}}%
\pgfpathlineto{\pgfqpoint{1.935204in}{1.272694in}}%
\pgfpathlineto{\pgfqpoint{1.935204in}{1.275643in}}%
\pgfpathlineto{\pgfqpoint{1.939745in}{1.275643in}}%
\pgfpathlineto{\pgfqpoint{1.939745in}{1.272694in}}%
\pgfpathmoveto{\pgfqpoint{1.939745in}{1.272694in}}%
\pgfpathlineto{\pgfqpoint{1.939745in}{1.272694in}}%
\pgfpathlineto{\pgfqpoint{1.939745in}{1.275643in}}%
\pgfpathlineto{\pgfqpoint{1.944286in}{1.275643in}}%
\pgfpathlineto{\pgfqpoint{1.944286in}{1.272694in}}%
\pgfpathmoveto{\pgfqpoint{1.939745in}{1.275643in}}%
\pgfpathlineto{\pgfqpoint{1.939745in}{1.275643in}}%
\pgfpathlineto{\pgfqpoint{1.939745in}{1.278592in}}%
\pgfpathlineto{\pgfqpoint{1.944286in}{1.278592in}}%
\pgfpathlineto{\pgfqpoint{1.944286in}{1.275643in}}%
\pgfpathmoveto{\pgfqpoint{1.944286in}{1.275643in}}%
\pgfpathlineto{\pgfqpoint{1.944286in}{1.275643in}}%
\pgfpathlineto{\pgfqpoint{1.944286in}{1.278592in}}%
\pgfpathlineto{\pgfqpoint{1.948827in}{1.278592in}}%
\pgfpathlineto{\pgfqpoint{1.948827in}{1.275643in}}%
\pgfpathmoveto{\pgfqpoint{1.944286in}{1.278592in}}%
\pgfpathlineto{\pgfqpoint{1.944286in}{1.278592in}}%
\pgfpathlineto{\pgfqpoint{1.944286in}{1.281541in}}%
\pgfpathlineto{\pgfqpoint{1.948827in}{1.281541in}}%
\pgfpathlineto{\pgfqpoint{1.948827in}{1.278592in}}%
\pgfpathmoveto{\pgfqpoint{1.948827in}{1.278592in}}%
\pgfpathlineto{\pgfqpoint{1.948827in}{1.278592in}}%
\pgfpathlineto{\pgfqpoint{1.948827in}{1.281541in}}%
\pgfpathlineto{\pgfqpoint{1.953368in}{1.281541in}}%
\pgfpathlineto{\pgfqpoint{1.953368in}{1.278592in}}%
\pgfpathmoveto{\pgfqpoint{1.948827in}{1.281541in}}%
\pgfpathlineto{\pgfqpoint{1.948827in}{1.281541in}}%
\pgfpathlineto{\pgfqpoint{1.948827in}{1.284491in}}%
\pgfpathlineto{\pgfqpoint{1.953368in}{1.284491in}}%
\pgfpathlineto{\pgfqpoint{1.953368in}{1.281541in}}%
\pgfpathmoveto{\pgfqpoint{1.953368in}{1.281541in}}%
\pgfpathlineto{\pgfqpoint{1.953368in}{1.281541in}}%
\pgfpathlineto{\pgfqpoint{1.953368in}{1.284491in}}%
\pgfpathlineto{\pgfqpoint{1.957909in}{1.284491in}}%
\pgfpathlineto{\pgfqpoint{1.957909in}{1.281541in}}%
\pgfpathmoveto{\pgfqpoint{1.953368in}{1.284491in}}%
\pgfpathlineto{\pgfqpoint{1.953368in}{1.284491in}}%
\pgfpathlineto{\pgfqpoint{1.953368in}{1.287440in}}%
\pgfpathlineto{\pgfqpoint{1.957909in}{1.287440in}}%
\pgfpathlineto{\pgfqpoint{1.957909in}{1.284491in}}%
\pgfpathmoveto{\pgfqpoint{1.957909in}{1.284491in}}%
\pgfpathlineto{\pgfqpoint{1.957909in}{1.284491in}}%
\pgfpathlineto{\pgfqpoint{1.957909in}{1.287440in}}%
\pgfpathlineto{\pgfqpoint{1.962450in}{1.287440in}}%
\pgfpathlineto{\pgfqpoint{1.962450in}{1.284491in}}%
\pgfpathmoveto{\pgfqpoint{1.957909in}{1.287440in}}%
\pgfpathlineto{\pgfqpoint{1.957909in}{1.287440in}}%
\pgfpathlineto{\pgfqpoint{1.957909in}{1.290389in}}%
\pgfpathlineto{\pgfqpoint{1.962450in}{1.290389in}}%
\pgfpathlineto{\pgfqpoint{1.962450in}{1.287440in}}%
\pgfpathmoveto{\pgfqpoint{1.962450in}{1.287440in}}%
\pgfpathlineto{\pgfqpoint{1.962450in}{1.287440in}}%
\pgfpathlineto{\pgfqpoint{1.962450in}{1.290389in}}%
\pgfpathlineto{\pgfqpoint{1.966991in}{1.290389in}}%
\pgfpathlineto{\pgfqpoint{1.966991in}{1.287440in}}%
\pgfpathmoveto{\pgfqpoint{1.962450in}{1.290389in}}%
\pgfpathlineto{\pgfqpoint{1.962450in}{1.290389in}}%
\pgfpathlineto{\pgfqpoint{1.962450in}{1.293339in}}%
\pgfpathlineto{\pgfqpoint{1.966991in}{1.293339in}}%
\pgfpathlineto{\pgfqpoint{1.966991in}{1.290389in}}%
\pgfpathmoveto{\pgfqpoint{1.966991in}{1.290389in}}%
\pgfpathlineto{\pgfqpoint{1.966991in}{1.290389in}}%
\pgfpathlineto{\pgfqpoint{1.966991in}{1.293339in}}%
\pgfpathlineto{\pgfqpoint{1.971532in}{1.293339in}}%
\pgfpathlineto{\pgfqpoint{1.971532in}{1.290389in}}%
\pgfpathmoveto{\pgfqpoint{1.966991in}{1.293339in}}%
\pgfpathlineto{\pgfqpoint{1.966991in}{1.293339in}}%
\pgfpathlineto{\pgfqpoint{1.966991in}{1.296288in}}%
\pgfpathlineto{\pgfqpoint{1.971532in}{1.296288in}}%
\pgfpathlineto{\pgfqpoint{1.971532in}{1.293339in}}%
\pgfpathmoveto{\pgfqpoint{1.971532in}{1.293339in}}%
\pgfpathlineto{\pgfqpoint{1.971532in}{1.293339in}}%
\pgfpathlineto{\pgfqpoint{1.971532in}{1.296288in}}%
\pgfpathlineto{\pgfqpoint{1.976073in}{1.296288in}}%
\pgfpathlineto{\pgfqpoint{1.976073in}{1.293339in}}%
\pgfpathmoveto{\pgfqpoint{1.971532in}{1.296288in}}%
\pgfpathlineto{\pgfqpoint{1.971532in}{1.296288in}}%
\pgfpathlineto{\pgfqpoint{1.971532in}{1.299237in}}%
\pgfpathlineto{\pgfqpoint{1.976073in}{1.299237in}}%
\pgfpathlineto{\pgfqpoint{1.976073in}{1.296288in}}%
\pgfpathmoveto{\pgfqpoint{1.976073in}{1.296288in}}%
\pgfpathlineto{\pgfqpoint{1.976073in}{1.296288in}}%
\pgfpathlineto{\pgfqpoint{1.976073in}{1.299237in}}%
\pgfpathlineto{\pgfqpoint{1.980614in}{1.299237in}}%
\pgfpathlineto{\pgfqpoint{1.980614in}{1.296288in}}%
\pgfpathmoveto{\pgfqpoint{1.976073in}{1.299237in}}%
\pgfpathlineto{\pgfqpoint{1.976073in}{1.299237in}}%
\pgfpathlineto{\pgfqpoint{1.976073in}{1.302187in}}%
\pgfpathlineto{\pgfqpoint{1.980614in}{1.302187in}}%
\pgfpathlineto{\pgfqpoint{1.980614in}{1.299237in}}%
\pgfpathmoveto{\pgfqpoint{1.980614in}{1.299237in}}%
\pgfpathlineto{\pgfqpoint{1.980614in}{1.299237in}}%
\pgfpathlineto{\pgfqpoint{1.980614in}{1.302187in}}%
\pgfpathlineto{\pgfqpoint{1.985155in}{1.302187in}}%
\pgfpathlineto{\pgfqpoint{1.985155in}{1.299237in}}%
\pgfpathmoveto{\pgfqpoint{1.980614in}{1.302187in}}%
\pgfpathlineto{\pgfqpoint{1.980614in}{1.302187in}}%
\pgfpathlineto{\pgfqpoint{1.980614in}{1.305136in}}%
\pgfpathlineto{\pgfqpoint{1.985155in}{1.305136in}}%
\pgfpathlineto{\pgfqpoint{1.985155in}{1.302187in}}%
\pgfpathmoveto{\pgfqpoint{1.985155in}{1.302187in}}%
\pgfpathlineto{\pgfqpoint{1.985155in}{1.302187in}}%
\pgfpathlineto{\pgfqpoint{1.985155in}{1.305136in}}%
\pgfpathlineto{\pgfqpoint{1.989696in}{1.305136in}}%
\pgfpathlineto{\pgfqpoint{1.989696in}{1.302187in}}%
\pgfpathmoveto{\pgfqpoint{1.989696in}{1.302187in}}%
\pgfpathlineto{\pgfqpoint{1.989696in}{1.302187in}}%
\pgfpathlineto{\pgfqpoint{1.989696in}{1.305136in}}%
\pgfpathlineto{\pgfqpoint{1.994237in}{1.305136in}}%
\pgfpathlineto{\pgfqpoint{1.994237in}{1.302187in}}%
\pgfpathmoveto{\pgfqpoint{1.989696in}{1.305136in}}%
\pgfpathlineto{\pgfqpoint{1.989696in}{1.305136in}}%
\pgfpathlineto{\pgfqpoint{1.989696in}{1.308085in}}%
\pgfpathlineto{\pgfqpoint{1.994237in}{1.308085in}}%
\pgfpathlineto{\pgfqpoint{1.994237in}{1.305136in}}%
\pgfpathmoveto{\pgfqpoint{1.994237in}{1.305136in}}%
\pgfpathlineto{\pgfqpoint{1.994237in}{1.305136in}}%
\pgfpathlineto{\pgfqpoint{1.994237in}{1.308085in}}%
\pgfpathlineto{\pgfqpoint{1.998778in}{1.308085in}}%
\pgfpathlineto{\pgfqpoint{1.998778in}{1.305136in}}%
\pgfpathmoveto{\pgfqpoint{1.994237in}{1.308085in}}%
\pgfpathlineto{\pgfqpoint{1.994237in}{1.308085in}}%
\pgfpathlineto{\pgfqpoint{1.994237in}{1.311035in}}%
\pgfpathlineto{\pgfqpoint{1.998778in}{1.311035in}}%
\pgfpathlineto{\pgfqpoint{1.998778in}{1.308085in}}%
\pgfpathmoveto{\pgfqpoint{1.998778in}{1.308085in}}%
\pgfpathlineto{\pgfqpoint{1.998778in}{1.308085in}}%
\pgfpathlineto{\pgfqpoint{1.998778in}{1.311035in}}%
\pgfpathlineto{\pgfqpoint{2.003319in}{1.311035in}}%
\pgfpathlineto{\pgfqpoint{2.003319in}{1.308085in}}%
\pgfpathmoveto{\pgfqpoint{1.998778in}{1.311035in}}%
\pgfpathlineto{\pgfqpoint{1.998778in}{1.311035in}}%
\pgfpathlineto{\pgfqpoint{1.998778in}{1.313984in}}%
\pgfpathlineto{\pgfqpoint{2.003319in}{1.313984in}}%
\pgfpathlineto{\pgfqpoint{2.003319in}{1.311035in}}%
\pgfpathmoveto{\pgfqpoint{2.003319in}{1.311035in}}%
\pgfpathlineto{\pgfqpoint{2.003319in}{1.311035in}}%
\pgfpathlineto{\pgfqpoint{2.003319in}{1.313984in}}%
\pgfpathlineto{\pgfqpoint{2.007860in}{1.313984in}}%
\pgfpathlineto{\pgfqpoint{2.007860in}{1.311035in}}%
\pgfpathmoveto{\pgfqpoint{2.003319in}{1.313984in}}%
\pgfpathlineto{\pgfqpoint{2.003319in}{1.313984in}}%
\pgfpathlineto{\pgfqpoint{2.003319in}{1.316933in}}%
\pgfpathlineto{\pgfqpoint{2.007860in}{1.316933in}}%
\pgfpathlineto{\pgfqpoint{2.007860in}{1.313984in}}%
\pgfpathmoveto{\pgfqpoint{2.007860in}{1.313984in}}%
\pgfpathlineto{\pgfqpoint{2.007860in}{1.313984in}}%
\pgfpathlineto{\pgfqpoint{2.007860in}{1.316933in}}%
\pgfpathlineto{\pgfqpoint{2.012401in}{1.316933in}}%
\pgfpathlineto{\pgfqpoint{2.012401in}{1.313984in}}%
\pgfpathmoveto{\pgfqpoint{2.007860in}{1.316933in}}%
\pgfpathlineto{\pgfqpoint{2.007860in}{1.316933in}}%
\pgfpathlineto{\pgfqpoint{2.007860in}{1.319882in}}%
\pgfpathlineto{\pgfqpoint{2.012401in}{1.319882in}}%
\pgfpathlineto{\pgfqpoint{2.012401in}{1.316933in}}%
\pgfpathmoveto{\pgfqpoint{2.012401in}{1.316933in}}%
\pgfpathlineto{\pgfqpoint{2.012401in}{1.316933in}}%
\pgfpathlineto{\pgfqpoint{2.012401in}{1.319882in}}%
\pgfpathlineto{\pgfqpoint{2.016942in}{1.319882in}}%
\pgfpathlineto{\pgfqpoint{2.016942in}{1.316933in}}%
\pgfpathmoveto{\pgfqpoint{2.012401in}{1.319882in}}%
\pgfpathlineto{\pgfqpoint{2.012401in}{1.319882in}}%
\pgfpathlineto{\pgfqpoint{2.012401in}{1.322832in}}%
\pgfpathlineto{\pgfqpoint{2.016942in}{1.322832in}}%
\pgfpathlineto{\pgfqpoint{2.016942in}{1.319882in}}%
\pgfpathmoveto{\pgfqpoint{2.016942in}{1.319882in}}%
\pgfpathlineto{\pgfqpoint{2.016942in}{1.319882in}}%
\pgfpathlineto{\pgfqpoint{2.016942in}{1.322832in}}%
\pgfpathlineto{\pgfqpoint{2.021483in}{1.322832in}}%
\pgfpathlineto{\pgfqpoint{2.021483in}{1.319882in}}%
\pgfpathmoveto{\pgfqpoint{2.016942in}{1.322832in}}%
\pgfpathlineto{\pgfqpoint{2.016942in}{1.322832in}}%
\pgfpathlineto{\pgfqpoint{2.016942in}{1.325781in}}%
\pgfpathlineto{\pgfqpoint{2.021483in}{1.325781in}}%
\pgfpathlineto{\pgfqpoint{2.021483in}{1.322832in}}%
\pgfpathmoveto{\pgfqpoint{2.021483in}{1.322832in}}%
\pgfpathlineto{\pgfqpoint{2.021483in}{1.322832in}}%
\pgfpathlineto{\pgfqpoint{2.021483in}{1.325781in}}%
\pgfpathlineto{\pgfqpoint{2.026024in}{1.325781in}}%
\pgfpathlineto{\pgfqpoint{2.026024in}{1.322832in}}%
\pgfpathmoveto{\pgfqpoint{2.021483in}{1.325781in}}%
\pgfpathlineto{\pgfqpoint{2.021483in}{1.325781in}}%
\pgfpathlineto{\pgfqpoint{2.021483in}{1.328730in}}%
\pgfpathlineto{\pgfqpoint{2.026024in}{1.328730in}}%
\pgfpathlineto{\pgfqpoint{2.026024in}{1.325781in}}%
\pgfpathmoveto{\pgfqpoint{2.026024in}{1.325781in}}%
\pgfpathlineto{\pgfqpoint{2.026024in}{1.325781in}}%
\pgfpathlineto{\pgfqpoint{2.026024in}{1.328730in}}%
\pgfpathlineto{\pgfqpoint{2.030565in}{1.328730in}}%
\pgfpathlineto{\pgfqpoint{2.030565in}{1.325781in}}%
\pgfpathmoveto{\pgfqpoint{2.026024in}{1.328730in}}%
\pgfpathlineto{\pgfqpoint{2.026024in}{1.328730in}}%
\pgfpathlineto{\pgfqpoint{2.026024in}{1.331680in}}%
\pgfpathlineto{\pgfqpoint{2.030565in}{1.331680in}}%
\pgfpathlineto{\pgfqpoint{2.030565in}{1.328730in}}%
\pgfpathmoveto{\pgfqpoint{2.030565in}{1.328730in}}%
\pgfpathlineto{\pgfqpoint{2.030565in}{1.328730in}}%
\pgfpathlineto{\pgfqpoint{2.030565in}{1.331680in}}%
\pgfpathlineto{\pgfqpoint{2.035106in}{1.331680in}}%
\pgfpathlineto{\pgfqpoint{2.035106in}{1.328730in}}%
\pgfpathmoveto{\pgfqpoint{2.030565in}{1.331680in}}%
\pgfpathlineto{\pgfqpoint{2.030565in}{1.331680in}}%
\pgfpathlineto{\pgfqpoint{2.030565in}{1.334629in}}%
\pgfpathlineto{\pgfqpoint{2.035106in}{1.334629in}}%
\pgfpathlineto{\pgfqpoint{2.035106in}{1.331680in}}%
\pgfpathmoveto{\pgfqpoint{2.035106in}{1.331680in}}%
\pgfpathlineto{\pgfqpoint{2.035106in}{1.331680in}}%
\pgfpathlineto{\pgfqpoint{2.035106in}{1.334629in}}%
\pgfpathlineto{\pgfqpoint{2.039647in}{1.334629in}}%
\pgfpathlineto{\pgfqpoint{2.039647in}{1.331680in}}%
\pgfpathmoveto{\pgfqpoint{2.035106in}{1.334629in}}%
\pgfpathlineto{\pgfqpoint{2.035106in}{1.334629in}}%
\pgfpathlineto{\pgfqpoint{2.035106in}{1.337578in}}%
\pgfpathlineto{\pgfqpoint{2.039647in}{1.337578in}}%
\pgfpathlineto{\pgfqpoint{2.039647in}{1.334629in}}%
\pgfpathmoveto{\pgfqpoint{2.039647in}{1.334629in}}%
\pgfpathlineto{\pgfqpoint{2.039647in}{1.334629in}}%
\pgfpathlineto{\pgfqpoint{2.039647in}{1.337578in}}%
\pgfpathlineto{\pgfqpoint{2.044188in}{1.337578in}}%
\pgfpathlineto{\pgfqpoint{2.044188in}{1.334629in}}%
\pgfpathmoveto{\pgfqpoint{2.039647in}{1.337578in}}%
\pgfpathlineto{\pgfqpoint{2.039647in}{1.337578in}}%
\pgfpathlineto{\pgfqpoint{2.039647in}{1.340528in}}%
\pgfpathlineto{\pgfqpoint{2.044188in}{1.340528in}}%
\pgfpathlineto{\pgfqpoint{2.044188in}{1.337578in}}%
\pgfpathmoveto{\pgfqpoint{2.044188in}{1.337578in}}%
\pgfpathlineto{\pgfqpoint{2.044188in}{1.337578in}}%
\pgfpathlineto{\pgfqpoint{2.044188in}{1.340528in}}%
\pgfpathlineto{\pgfqpoint{2.048729in}{1.340528in}}%
\pgfpathlineto{\pgfqpoint{2.048729in}{1.337578in}}%
\pgfpathmoveto{\pgfqpoint{2.044188in}{1.340528in}}%
\pgfpathlineto{\pgfqpoint{2.044188in}{1.340528in}}%
\pgfpathlineto{\pgfqpoint{2.044188in}{1.343477in}}%
\pgfpathlineto{\pgfqpoint{2.048729in}{1.343477in}}%
\pgfpathlineto{\pgfqpoint{2.048729in}{1.340528in}}%
\pgfpathmoveto{\pgfqpoint{2.048729in}{1.340528in}}%
\pgfpathlineto{\pgfqpoint{2.048729in}{1.340528in}}%
\pgfpathlineto{\pgfqpoint{2.048729in}{1.343477in}}%
\pgfpathlineto{\pgfqpoint{2.053270in}{1.343477in}}%
\pgfpathlineto{\pgfqpoint{2.053270in}{1.340528in}}%
\pgfpathmoveto{\pgfqpoint{2.048729in}{1.343477in}}%
\pgfpathlineto{\pgfqpoint{2.048729in}{1.343477in}}%
\pgfpathlineto{\pgfqpoint{2.048729in}{1.346426in}}%
\pgfpathlineto{\pgfqpoint{2.053270in}{1.346426in}}%
\pgfpathlineto{\pgfqpoint{2.053270in}{1.343477in}}%
\pgfpathmoveto{\pgfqpoint{2.053270in}{1.343477in}}%
\pgfpathlineto{\pgfqpoint{2.053270in}{1.343477in}}%
\pgfpathlineto{\pgfqpoint{2.053270in}{1.346426in}}%
\pgfpathlineto{\pgfqpoint{2.057810in}{1.346426in}}%
\pgfpathlineto{\pgfqpoint{2.057810in}{1.343477in}}%
\pgfpathmoveto{\pgfqpoint{2.053270in}{1.346426in}}%
\pgfpathlineto{\pgfqpoint{2.053270in}{1.346426in}}%
\pgfpathlineto{\pgfqpoint{2.053270in}{1.349375in}}%
\pgfpathlineto{\pgfqpoint{2.057810in}{1.349375in}}%
\pgfpathlineto{\pgfqpoint{2.057810in}{1.346426in}}%
\pgfpathmoveto{\pgfqpoint{2.057810in}{1.346426in}}%
\pgfpathlineto{\pgfqpoint{2.057810in}{1.346426in}}%
\pgfpathlineto{\pgfqpoint{2.057810in}{1.349375in}}%
\pgfpathlineto{\pgfqpoint{2.062351in}{1.349375in}}%
\pgfpathlineto{\pgfqpoint{2.062351in}{1.346426in}}%
\pgfpathmoveto{\pgfqpoint{2.057810in}{1.349375in}}%
\pgfpathlineto{\pgfqpoint{2.057810in}{1.349375in}}%
\pgfpathlineto{\pgfqpoint{2.057810in}{1.352325in}}%
\pgfpathlineto{\pgfqpoint{2.062351in}{1.352325in}}%
\pgfpathlineto{\pgfqpoint{2.062351in}{1.349375in}}%
\pgfpathmoveto{\pgfqpoint{2.062351in}{1.349375in}}%
\pgfpathlineto{\pgfqpoint{2.062351in}{1.349375in}}%
\pgfpathlineto{\pgfqpoint{2.062351in}{1.352325in}}%
\pgfpathlineto{\pgfqpoint{2.066893in}{1.352325in}}%
\pgfpathlineto{\pgfqpoint{2.066893in}{1.349375in}}%
\pgfpathmoveto{\pgfqpoint{2.062351in}{1.352325in}}%
\pgfpathlineto{\pgfqpoint{2.062351in}{1.352325in}}%
\pgfpathlineto{\pgfqpoint{2.062351in}{1.355274in}}%
\pgfpathlineto{\pgfqpoint{2.066893in}{1.355274in}}%
\pgfpathlineto{\pgfqpoint{2.066893in}{1.352325in}}%
\pgfpathmoveto{\pgfqpoint{2.066893in}{1.352325in}}%
\pgfpathlineto{\pgfqpoint{2.066893in}{1.352325in}}%
\pgfpathlineto{\pgfqpoint{2.066893in}{1.355274in}}%
\pgfpathlineto{\pgfqpoint{2.071434in}{1.355274in}}%
\pgfpathlineto{\pgfqpoint{2.071434in}{1.352325in}}%
\pgfpathmoveto{\pgfqpoint{2.066893in}{1.355274in}}%
\pgfpathlineto{\pgfqpoint{2.066893in}{1.355274in}}%
\pgfpathlineto{\pgfqpoint{2.066893in}{1.358223in}}%
\pgfpathlineto{\pgfqpoint{2.071434in}{1.358223in}}%
\pgfpathlineto{\pgfqpoint{2.071434in}{1.355274in}}%
\pgfpathmoveto{\pgfqpoint{2.071434in}{1.355274in}}%
\pgfpathlineto{\pgfqpoint{2.071434in}{1.355274in}}%
\pgfpathlineto{\pgfqpoint{2.071434in}{1.358223in}}%
\pgfpathlineto{\pgfqpoint{2.075975in}{1.358223in}}%
\pgfpathlineto{\pgfqpoint{2.075975in}{1.355274in}}%
\pgfpathmoveto{\pgfqpoint{2.071434in}{1.358223in}}%
\pgfpathlineto{\pgfqpoint{2.071434in}{1.358223in}}%
\pgfpathlineto{\pgfqpoint{2.071434in}{1.361172in}}%
\pgfpathlineto{\pgfqpoint{2.075975in}{1.361172in}}%
\pgfpathlineto{\pgfqpoint{2.075975in}{1.358223in}}%
\pgfpathmoveto{\pgfqpoint{2.075975in}{1.358223in}}%
\pgfpathlineto{\pgfqpoint{2.075975in}{1.358223in}}%
\pgfpathlineto{\pgfqpoint{2.075975in}{1.361172in}}%
\pgfpathlineto{\pgfqpoint{2.080516in}{1.361172in}}%
\pgfpathlineto{\pgfqpoint{2.080516in}{1.358223in}}%
\pgfpathmoveto{\pgfqpoint{2.075975in}{1.361172in}}%
\pgfpathlineto{\pgfqpoint{2.075975in}{1.361172in}}%
\pgfpathlineto{\pgfqpoint{2.075975in}{1.364122in}}%
\pgfpathlineto{\pgfqpoint{2.080516in}{1.364122in}}%
\pgfpathlineto{\pgfqpoint{2.080516in}{1.361172in}}%
\pgfpathmoveto{\pgfqpoint{2.080516in}{1.361172in}}%
\pgfpathlineto{\pgfqpoint{2.080516in}{1.361172in}}%
\pgfpathlineto{\pgfqpoint{2.080516in}{1.364122in}}%
\pgfpathlineto{\pgfqpoint{2.085057in}{1.364122in}}%
\pgfpathlineto{\pgfqpoint{2.085057in}{1.361172in}}%
\pgfpathmoveto{\pgfqpoint{2.080516in}{1.364122in}}%
\pgfpathlineto{\pgfqpoint{2.080516in}{1.364122in}}%
\pgfpathlineto{\pgfqpoint{2.080516in}{1.367071in}}%
\pgfpathlineto{\pgfqpoint{2.085057in}{1.367071in}}%
\pgfpathlineto{\pgfqpoint{2.085057in}{1.364122in}}%
\pgfpathmoveto{\pgfqpoint{2.085057in}{1.364122in}}%
\pgfpathlineto{\pgfqpoint{2.085057in}{1.364122in}}%
\pgfpathlineto{\pgfqpoint{2.085057in}{1.367071in}}%
\pgfpathlineto{\pgfqpoint{2.089598in}{1.367071in}}%
\pgfpathlineto{\pgfqpoint{2.089598in}{1.364122in}}%
\pgfpathmoveto{\pgfqpoint{2.085057in}{1.367071in}}%
\pgfpathlineto{\pgfqpoint{2.085057in}{1.367071in}}%
\pgfpathlineto{\pgfqpoint{2.085057in}{1.370020in}}%
\pgfpathlineto{\pgfqpoint{2.089598in}{1.370020in}}%
\pgfpathlineto{\pgfqpoint{2.089598in}{1.367071in}}%
\pgfpathmoveto{\pgfqpoint{2.089598in}{1.367071in}}%
\pgfpathlineto{\pgfqpoint{2.089598in}{1.367071in}}%
\pgfpathlineto{\pgfqpoint{2.089598in}{1.370020in}}%
\pgfpathlineto{\pgfqpoint{2.094139in}{1.370020in}}%
\pgfpathlineto{\pgfqpoint{2.094139in}{1.367071in}}%
\pgfpathmoveto{\pgfqpoint{2.089598in}{1.370020in}}%
\pgfpathlineto{\pgfqpoint{2.089598in}{1.370020in}}%
\pgfpathlineto{\pgfqpoint{2.089598in}{1.372969in}}%
\pgfpathlineto{\pgfqpoint{2.094139in}{1.372969in}}%
\pgfpathlineto{\pgfqpoint{2.094139in}{1.370020in}}%
\pgfpathmoveto{\pgfqpoint{2.094139in}{1.370020in}}%
\pgfpathlineto{\pgfqpoint{2.094139in}{1.370020in}}%
\pgfpathlineto{\pgfqpoint{2.094139in}{1.372969in}}%
\pgfpathlineto{\pgfqpoint{2.098680in}{1.372969in}}%
\pgfpathlineto{\pgfqpoint{2.098680in}{1.370020in}}%
\pgfpathmoveto{\pgfqpoint{2.094139in}{1.372969in}}%
\pgfpathlineto{\pgfqpoint{2.094139in}{1.372969in}}%
\pgfpathlineto{\pgfqpoint{2.094139in}{1.375918in}}%
\pgfpathlineto{\pgfqpoint{2.098680in}{1.375918in}}%
\pgfpathlineto{\pgfqpoint{2.098680in}{1.372969in}}%
\pgfpathmoveto{\pgfqpoint{2.098680in}{1.372969in}}%
\pgfpathlineto{\pgfqpoint{2.098680in}{1.372969in}}%
\pgfpathlineto{\pgfqpoint{2.098680in}{1.375918in}}%
\pgfpathlineto{\pgfqpoint{2.103221in}{1.375918in}}%
\pgfpathlineto{\pgfqpoint{2.103221in}{1.372969in}}%
\pgfpathmoveto{\pgfqpoint{2.098680in}{1.375918in}}%
\pgfpathlineto{\pgfqpoint{2.098680in}{1.375918in}}%
\pgfpathlineto{\pgfqpoint{2.098680in}{1.378868in}}%
\pgfpathlineto{\pgfqpoint{2.103221in}{1.378868in}}%
\pgfpathlineto{\pgfqpoint{2.103221in}{1.375918in}}%
\pgfpathmoveto{\pgfqpoint{2.103221in}{1.375918in}}%
\pgfpathlineto{\pgfqpoint{2.103221in}{1.375918in}}%
\pgfpathlineto{\pgfqpoint{2.103221in}{1.378868in}}%
\pgfpathlineto{\pgfqpoint{2.107762in}{1.378868in}}%
\pgfpathlineto{\pgfqpoint{2.107762in}{1.375918in}}%
\pgfpathmoveto{\pgfqpoint{2.103221in}{1.378868in}}%
\pgfpathlineto{\pgfqpoint{2.103221in}{1.378868in}}%
\pgfpathlineto{\pgfqpoint{2.103221in}{1.381817in}}%
\pgfpathlineto{\pgfqpoint{2.107762in}{1.381817in}}%
\pgfpathlineto{\pgfqpoint{2.107762in}{1.378868in}}%
\pgfpathmoveto{\pgfqpoint{2.107762in}{1.378868in}}%
\pgfpathlineto{\pgfqpoint{2.107762in}{1.378868in}}%
\pgfpathlineto{\pgfqpoint{2.107762in}{1.381817in}}%
\pgfpathlineto{\pgfqpoint{2.112303in}{1.381817in}}%
\pgfpathlineto{\pgfqpoint{2.112303in}{1.378868in}}%
\pgfpathmoveto{\pgfqpoint{2.107762in}{1.381817in}}%
\pgfpathlineto{\pgfqpoint{2.107762in}{1.381817in}}%
\pgfpathlineto{\pgfqpoint{2.107762in}{1.384766in}}%
\pgfpathlineto{\pgfqpoint{2.112303in}{1.384766in}}%
\pgfpathlineto{\pgfqpoint{2.112303in}{1.381817in}}%
\pgfpathmoveto{\pgfqpoint{2.112303in}{1.381817in}}%
\pgfpathlineto{\pgfqpoint{2.112303in}{1.381817in}}%
\pgfpathlineto{\pgfqpoint{2.112303in}{1.384766in}}%
\pgfpathlineto{\pgfqpoint{2.116844in}{1.384766in}}%
\pgfpathlineto{\pgfqpoint{2.116844in}{1.381817in}}%
\pgfpathmoveto{\pgfqpoint{2.112303in}{1.384766in}}%
\pgfpathlineto{\pgfqpoint{2.112303in}{1.384766in}}%
\pgfpathlineto{\pgfqpoint{2.112303in}{1.387715in}}%
\pgfpathlineto{\pgfqpoint{2.116844in}{1.387715in}}%
\pgfpathlineto{\pgfqpoint{2.116844in}{1.384766in}}%
\pgfpathmoveto{\pgfqpoint{2.116844in}{1.384766in}}%
\pgfpathlineto{\pgfqpoint{2.116844in}{1.384766in}}%
\pgfpathlineto{\pgfqpoint{2.116844in}{1.387715in}}%
\pgfpathlineto{\pgfqpoint{2.121385in}{1.387715in}}%
\pgfpathlineto{\pgfqpoint{2.121385in}{1.384766in}}%
\pgfpathmoveto{\pgfqpoint{2.116844in}{1.387715in}}%
\pgfpathlineto{\pgfqpoint{2.116844in}{1.387715in}}%
\pgfpathlineto{\pgfqpoint{2.116844in}{1.390664in}}%
\pgfpathlineto{\pgfqpoint{2.121385in}{1.390664in}}%
\pgfpathlineto{\pgfqpoint{2.121385in}{1.387715in}}%
\pgfpathmoveto{\pgfqpoint{2.121385in}{1.387715in}}%
\pgfpathlineto{\pgfqpoint{2.121385in}{1.387715in}}%
\pgfpathlineto{\pgfqpoint{2.121385in}{1.390664in}}%
\pgfpathlineto{\pgfqpoint{2.125926in}{1.390664in}}%
\pgfpathlineto{\pgfqpoint{2.125926in}{1.387715in}}%
\pgfpathmoveto{\pgfqpoint{2.121385in}{1.390664in}}%
\pgfpathlineto{\pgfqpoint{2.121385in}{1.390664in}}%
\pgfpathlineto{\pgfqpoint{2.121385in}{1.393614in}}%
\pgfpathlineto{\pgfqpoint{2.125926in}{1.393614in}}%
\pgfpathlineto{\pgfqpoint{2.125926in}{1.390664in}}%
\pgfpathmoveto{\pgfqpoint{2.125926in}{1.390664in}}%
\pgfpathlineto{\pgfqpoint{2.125926in}{1.390664in}}%
\pgfpathlineto{\pgfqpoint{2.125926in}{1.393614in}}%
\pgfpathlineto{\pgfqpoint{2.130467in}{1.393614in}}%
\pgfpathlineto{\pgfqpoint{2.130467in}{1.390664in}}%
\pgfpathmoveto{\pgfqpoint{2.125926in}{1.393614in}}%
\pgfpathlineto{\pgfqpoint{2.125926in}{1.393614in}}%
\pgfpathlineto{\pgfqpoint{2.125926in}{1.396563in}}%
\pgfpathlineto{\pgfqpoint{2.130467in}{1.396563in}}%
\pgfpathlineto{\pgfqpoint{2.130467in}{1.393614in}}%
\pgfpathmoveto{\pgfqpoint{2.130467in}{1.393614in}}%
\pgfpathlineto{\pgfqpoint{2.130467in}{1.393614in}}%
\pgfpathlineto{\pgfqpoint{2.130467in}{1.396563in}}%
\pgfpathlineto{\pgfqpoint{2.135008in}{1.396563in}}%
\pgfpathlineto{\pgfqpoint{2.135008in}{1.393614in}}%
\pgfpathmoveto{\pgfqpoint{2.130467in}{1.396563in}}%
\pgfpathlineto{\pgfqpoint{2.130467in}{1.396563in}}%
\pgfpathlineto{\pgfqpoint{2.130467in}{1.399512in}}%
\pgfpathlineto{\pgfqpoint{2.135008in}{1.399512in}}%
\pgfpathlineto{\pgfqpoint{2.135008in}{1.396563in}}%
\pgfpathmoveto{\pgfqpoint{2.135008in}{1.396563in}}%
\pgfpathlineto{\pgfqpoint{2.135008in}{1.396563in}}%
\pgfpathlineto{\pgfqpoint{2.135008in}{1.399512in}}%
\pgfpathlineto{\pgfqpoint{2.139549in}{1.399512in}}%
\pgfpathlineto{\pgfqpoint{2.139549in}{1.396563in}}%
\pgfpathmoveto{\pgfqpoint{2.135008in}{1.399512in}}%
\pgfpathlineto{\pgfqpoint{2.135008in}{1.399512in}}%
\pgfpathlineto{\pgfqpoint{2.135008in}{1.402461in}}%
\pgfpathlineto{\pgfqpoint{2.139549in}{1.402461in}}%
\pgfpathlineto{\pgfqpoint{2.139549in}{1.399512in}}%
\pgfpathmoveto{\pgfqpoint{2.139549in}{1.399512in}}%
\pgfpathlineto{\pgfqpoint{2.139549in}{1.399512in}}%
\pgfpathlineto{\pgfqpoint{2.139549in}{1.402461in}}%
\pgfpathlineto{\pgfqpoint{2.144090in}{1.402461in}}%
\pgfpathlineto{\pgfqpoint{2.144090in}{1.399512in}}%
\pgfpathmoveto{\pgfqpoint{2.139549in}{1.402461in}}%
\pgfpathlineto{\pgfqpoint{2.139549in}{1.402461in}}%
\pgfpathlineto{\pgfqpoint{2.139549in}{1.405411in}}%
\pgfpathlineto{\pgfqpoint{2.144090in}{1.405411in}}%
\pgfpathlineto{\pgfqpoint{2.144090in}{1.402461in}}%
\pgfpathmoveto{\pgfqpoint{2.144090in}{1.402461in}}%
\pgfpathlineto{\pgfqpoint{2.144090in}{1.402461in}}%
\pgfpathlineto{\pgfqpoint{2.144090in}{1.405411in}}%
\pgfpathlineto{\pgfqpoint{2.148631in}{1.405411in}}%
\pgfpathlineto{\pgfqpoint{2.148631in}{1.402461in}}%
\pgfpathmoveto{\pgfqpoint{2.144090in}{1.405411in}}%
\pgfpathlineto{\pgfqpoint{2.144090in}{1.405411in}}%
\pgfpathlineto{\pgfqpoint{2.144090in}{1.408360in}}%
\pgfpathlineto{\pgfqpoint{2.148631in}{1.408360in}}%
\pgfpathlineto{\pgfqpoint{2.148631in}{1.405411in}}%
\pgfpathmoveto{\pgfqpoint{2.148631in}{1.405411in}}%
\pgfpathlineto{\pgfqpoint{2.148631in}{1.405411in}}%
\pgfpathlineto{\pgfqpoint{2.148631in}{1.408360in}}%
\pgfpathlineto{\pgfqpoint{2.153172in}{1.408360in}}%
\pgfpathlineto{\pgfqpoint{2.153172in}{1.405411in}}%
\pgfpathmoveto{\pgfqpoint{2.148631in}{1.408360in}}%
\pgfpathlineto{\pgfqpoint{2.148631in}{1.408360in}}%
\pgfpathlineto{\pgfqpoint{2.148631in}{1.411309in}}%
\pgfpathlineto{\pgfqpoint{2.153172in}{1.411309in}}%
\pgfpathlineto{\pgfqpoint{2.153172in}{1.408360in}}%
\pgfpathmoveto{\pgfqpoint{2.153172in}{1.408360in}}%
\pgfpathlineto{\pgfqpoint{2.153172in}{1.408360in}}%
\pgfpathlineto{\pgfqpoint{2.153172in}{1.411309in}}%
\pgfpathlineto{\pgfqpoint{2.157713in}{1.411309in}}%
\pgfpathlineto{\pgfqpoint{2.157713in}{1.408360in}}%
\pgfpathmoveto{\pgfqpoint{2.153172in}{1.411309in}}%
\pgfpathlineto{\pgfqpoint{2.153172in}{1.411309in}}%
\pgfpathlineto{\pgfqpoint{2.153172in}{1.414258in}}%
\pgfpathlineto{\pgfqpoint{2.157713in}{1.414258in}}%
\pgfpathlineto{\pgfqpoint{2.157713in}{1.411309in}}%
\pgfpathmoveto{\pgfqpoint{2.157713in}{1.411309in}}%
\pgfpathlineto{\pgfqpoint{2.157713in}{1.411309in}}%
\pgfpathlineto{\pgfqpoint{2.157713in}{1.414258in}}%
\pgfpathlineto{\pgfqpoint{2.162254in}{1.414258in}}%
\pgfpathlineto{\pgfqpoint{2.162254in}{1.411309in}}%
\pgfpathmoveto{\pgfqpoint{2.157713in}{1.414258in}}%
\pgfpathlineto{\pgfqpoint{2.157713in}{1.414258in}}%
\pgfpathlineto{\pgfqpoint{2.157713in}{1.417207in}}%
\pgfpathlineto{\pgfqpoint{2.162254in}{1.417207in}}%
\pgfpathlineto{\pgfqpoint{2.162254in}{1.414258in}}%
\pgfpathmoveto{\pgfqpoint{2.162254in}{1.414258in}}%
\pgfpathlineto{\pgfqpoint{2.162254in}{1.414258in}}%
\pgfpathlineto{\pgfqpoint{2.162254in}{1.417207in}}%
\pgfpathlineto{\pgfqpoint{2.166795in}{1.417207in}}%
\pgfpathlineto{\pgfqpoint{2.166795in}{1.414258in}}%
\pgfpathmoveto{\pgfqpoint{2.162254in}{1.417207in}}%
\pgfpathlineto{\pgfqpoint{2.162254in}{1.417207in}}%
\pgfpathlineto{\pgfqpoint{2.162254in}{1.420157in}}%
\pgfpathlineto{\pgfqpoint{2.166795in}{1.420157in}}%
\pgfpathlineto{\pgfqpoint{2.166795in}{1.417207in}}%
\pgfpathmoveto{\pgfqpoint{2.166795in}{1.417207in}}%
\pgfpathlineto{\pgfqpoint{2.166795in}{1.417207in}}%
\pgfpathlineto{\pgfqpoint{2.166795in}{1.420157in}}%
\pgfpathlineto{\pgfqpoint{2.171336in}{1.420157in}}%
\pgfpathlineto{\pgfqpoint{2.171336in}{1.417207in}}%
\pgfpathmoveto{\pgfqpoint{2.166795in}{1.420157in}}%
\pgfpathlineto{\pgfqpoint{2.166795in}{1.420157in}}%
\pgfpathlineto{\pgfqpoint{2.166795in}{1.423106in}}%
\pgfpathlineto{\pgfqpoint{2.171336in}{1.423106in}}%
\pgfpathlineto{\pgfqpoint{2.171336in}{1.420157in}}%
\pgfpathmoveto{\pgfqpoint{2.171336in}{1.420157in}}%
\pgfpathlineto{\pgfqpoint{2.171336in}{1.420157in}}%
\pgfpathlineto{\pgfqpoint{2.171336in}{1.423106in}}%
\pgfpathlineto{\pgfqpoint{2.175877in}{1.423106in}}%
\pgfpathlineto{\pgfqpoint{2.175877in}{1.420157in}}%
\pgfpathmoveto{\pgfqpoint{2.171336in}{1.423106in}}%
\pgfpathlineto{\pgfqpoint{2.171336in}{1.423106in}}%
\pgfpathlineto{\pgfqpoint{2.171336in}{1.426055in}}%
\pgfpathlineto{\pgfqpoint{2.175877in}{1.426055in}}%
\pgfpathlineto{\pgfqpoint{2.175877in}{1.423106in}}%
\pgfpathmoveto{\pgfqpoint{2.175877in}{1.423106in}}%
\pgfpathlineto{\pgfqpoint{2.175877in}{1.423106in}}%
\pgfpathlineto{\pgfqpoint{2.175877in}{1.426055in}}%
\pgfpathlineto{\pgfqpoint{2.180418in}{1.426055in}}%
\pgfpathlineto{\pgfqpoint{2.180418in}{1.423106in}}%
\pgfpathmoveto{\pgfqpoint{2.175877in}{1.426055in}}%
\pgfpathlineto{\pgfqpoint{2.175877in}{1.426055in}}%
\pgfpathlineto{\pgfqpoint{2.175877in}{1.429004in}}%
\pgfpathlineto{\pgfqpoint{2.180418in}{1.429004in}}%
\pgfpathlineto{\pgfqpoint{2.180418in}{1.426055in}}%
\pgfpathmoveto{\pgfqpoint{2.180418in}{1.426055in}}%
\pgfpathlineto{\pgfqpoint{2.180418in}{1.426055in}}%
\pgfpathlineto{\pgfqpoint{2.180418in}{1.429004in}}%
\pgfpathlineto{\pgfqpoint{2.184959in}{1.429004in}}%
\pgfpathlineto{\pgfqpoint{2.184959in}{1.426055in}}%
\pgfpathmoveto{\pgfqpoint{2.180418in}{1.429004in}}%
\pgfpathlineto{\pgfqpoint{2.180418in}{1.429004in}}%
\pgfpathlineto{\pgfqpoint{2.180418in}{1.431954in}}%
\pgfpathlineto{\pgfqpoint{2.184959in}{1.431954in}}%
\pgfpathlineto{\pgfqpoint{2.184959in}{1.429004in}}%
\pgfpathmoveto{\pgfqpoint{2.184959in}{1.429004in}}%
\pgfpathlineto{\pgfqpoint{2.184959in}{1.429004in}}%
\pgfpathlineto{\pgfqpoint{2.184959in}{1.431954in}}%
\pgfpathlineto{\pgfqpoint{2.189500in}{1.431954in}}%
\pgfpathlineto{\pgfqpoint{2.189500in}{1.429004in}}%
\pgfpathmoveto{\pgfqpoint{2.184959in}{1.431954in}}%
\pgfpathlineto{\pgfqpoint{2.184959in}{1.431954in}}%
\pgfpathlineto{\pgfqpoint{2.184959in}{1.434903in}}%
\pgfpathlineto{\pgfqpoint{2.189500in}{1.434903in}}%
\pgfpathlineto{\pgfqpoint{2.189500in}{1.431954in}}%
\pgfpathmoveto{\pgfqpoint{2.189500in}{1.431954in}}%
\pgfpathlineto{\pgfqpoint{2.189500in}{1.431954in}}%
\pgfpathlineto{\pgfqpoint{2.189500in}{1.434903in}}%
\pgfpathlineto{\pgfqpoint{2.194041in}{1.434903in}}%
\pgfpathlineto{\pgfqpoint{2.194041in}{1.431954in}}%
\pgfpathmoveto{\pgfqpoint{2.189500in}{1.434903in}}%
\pgfpathlineto{\pgfqpoint{2.189500in}{1.434903in}}%
\pgfpathlineto{\pgfqpoint{2.189500in}{1.437852in}}%
\pgfpathlineto{\pgfqpoint{2.194041in}{1.437852in}}%
\pgfpathlineto{\pgfqpoint{2.194041in}{1.434903in}}%
\pgfpathmoveto{\pgfqpoint{2.194041in}{1.434903in}}%
\pgfpathlineto{\pgfqpoint{2.194041in}{1.434903in}}%
\pgfpathlineto{\pgfqpoint{2.194041in}{1.437852in}}%
\pgfpathlineto{\pgfqpoint{2.198582in}{1.437852in}}%
\pgfpathlineto{\pgfqpoint{2.198582in}{1.434903in}}%
\pgfpathmoveto{\pgfqpoint{2.194041in}{1.437852in}}%
\pgfpathlineto{\pgfqpoint{2.194041in}{1.437852in}}%
\pgfpathlineto{\pgfqpoint{2.194041in}{1.440801in}}%
\pgfpathlineto{\pgfqpoint{2.198582in}{1.440801in}}%
\pgfpathlineto{\pgfqpoint{2.198582in}{1.437852in}}%
\pgfpathmoveto{\pgfqpoint{2.198582in}{1.437852in}}%
\pgfpathlineto{\pgfqpoint{2.198582in}{1.437852in}}%
\pgfpathlineto{\pgfqpoint{2.198582in}{1.440801in}}%
\pgfpathlineto{\pgfqpoint{2.203123in}{1.440801in}}%
\pgfpathlineto{\pgfqpoint{2.203123in}{1.437852in}}%
\pgfpathmoveto{\pgfqpoint{2.198582in}{1.440801in}}%
\pgfpathlineto{\pgfqpoint{2.198582in}{1.440801in}}%
\pgfpathlineto{\pgfqpoint{2.198582in}{1.443750in}}%
\pgfpathlineto{\pgfqpoint{2.203123in}{1.443750in}}%
\pgfpathlineto{\pgfqpoint{2.203123in}{1.440801in}}%
\pgfpathmoveto{\pgfqpoint{2.203123in}{1.440801in}}%
\pgfpathlineto{\pgfqpoint{2.203123in}{1.440801in}}%
\pgfpathlineto{\pgfqpoint{2.203123in}{1.443750in}}%
\pgfpathlineto{\pgfqpoint{2.207664in}{1.443750in}}%
\pgfpathlineto{\pgfqpoint{2.207664in}{1.440801in}}%
\pgfpathmoveto{\pgfqpoint{2.203123in}{1.443750in}}%
\pgfpathlineto{\pgfqpoint{2.203123in}{1.443750in}}%
\pgfpathlineto{\pgfqpoint{2.203123in}{1.446700in}}%
\pgfpathlineto{\pgfqpoint{2.207664in}{1.446700in}}%
\pgfpathlineto{\pgfqpoint{2.207664in}{1.443750in}}%
\pgfpathmoveto{\pgfqpoint{2.207664in}{1.443750in}}%
\pgfpathlineto{\pgfqpoint{2.207664in}{1.443750in}}%
\pgfpathlineto{\pgfqpoint{2.207664in}{1.446700in}}%
\pgfpathlineto{\pgfqpoint{2.212205in}{1.446700in}}%
\pgfpathlineto{\pgfqpoint{2.212205in}{1.443750in}}%
\pgfpathmoveto{\pgfqpoint{2.207664in}{1.446700in}}%
\pgfpathlineto{\pgfqpoint{2.207664in}{1.446700in}}%
\pgfpathlineto{\pgfqpoint{2.207664in}{1.449649in}}%
\pgfpathlineto{\pgfqpoint{2.212205in}{1.449649in}}%
\pgfpathlineto{\pgfqpoint{2.212205in}{1.446700in}}%
\pgfpathmoveto{\pgfqpoint{2.212205in}{1.446700in}}%
\pgfpathlineto{\pgfqpoint{2.212205in}{1.446700in}}%
\pgfpathlineto{\pgfqpoint{2.212205in}{1.449649in}}%
\pgfpathlineto{\pgfqpoint{2.216746in}{1.449649in}}%
\pgfpathlineto{\pgfqpoint{2.216746in}{1.446700in}}%
\pgfpathmoveto{\pgfqpoint{2.212205in}{1.449649in}}%
\pgfpathlineto{\pgfqpoint{2.212205in}{1.449649in}}%
\pgfpathlineto{\pgfqpoint{2.212205in}{1.452598in}}%
\pgfpathlineto{\pgfqpoint{2.216746in}{1.452598in}}%
\pgfpathlineto{\pgfqpoint{2.216746in}{1.449649in}}%
\pgfpathmoveto{\pgfqpoint{2.216746in}{1.449649in}}%
\pgfpathlineto{\pgfqpoint{2.216746in}{1.449649in}}%
\pgfpathlineto{\pgfqpoint{2.216746in}{1.452598in}}%
\pgfpathlineto{\pgfqpoint{2.221287in}{1.452598in}}%
\pgfpathlineto{\pgfqpoint{2.221287in}{1.449649in}}%
\pgfpathmoveto{\pgfqpoint{2.216746in}{1.452598in}}%
\pgfpathlineto{\pgfqpoint{2.216746in}{1.452598in}}%
\pgfpathlineto{\pgfqpoint{2.216746in}{1.455547in}}%
\pgfpathlineto{\pgfqpoint{2.221287in}{1.455547in}}%
\pgfpathlineto{\pgfqpoint{2.221287in}{1.452598in}}%
\pgfpathmoveto{\pgfqpoint{2.221287in}{1.452598in}}%
\pgfpathlineto{\pgfqpoint{2.221287in}{1.452598in}}%
\pgfpathlineto{\pgfqpoint{2.221287in}{1.455547in}}%
\pgfpathlineto{\pgfqpoint{2.225828in}{1.455547in}}%
\pgfpathlineto{\pgfqpoint{2.225828in}{1.452598in}}%
\pgfpathmoveto{\pgfqpoint{2.221287in}{1.455547in}}%
\pgfpathlineto{\pgfqpoint{2.221287in}{1.455547in}}%
\pgfpathlineto{\pgfqpoint{2.221287in}{1.458497in}}%
\pgfpathlineto{\pgfqpoint{2.225828in}{1.458497in}}%
\pgfpathlineto{\pgfqpoint{2.225828in}{1.455547in}}%
\pgfpathmoveto{\pgfqpoint{2.225828in}{1.455547in}}%
\pgfpathlineto{\pgfqpoint{2.225828in}{1.455547in}}%
\pgfpathlineto{\pgfqpoint{2.225828in}{1.458497in}}%
\pgfpathlineto{\pgfqpoint{2.230369in}{1.458497in}}%
\pgfpathlineto{\pgfqpoint{2.230369in}{1.455547in}}%
\pgfpathmoveto{\pgfqpoint{2.225828in}{1.458497in}}%
\pgfpathlineto{\pgfqpoint{2.225828in}{1.458497in}}%
\pgfpathlineto{\pgfqpoint{2.225828in}{1.461446in}}%
\pgfpathlineto{\pgfqpoint{2.230369in}{1.461446in}}%
\pgfpathlineto{\pgfqpoint{2.230369in}{1.458497in}}%
\pgfpathmoveto{\pgfqpoint{2.230369in}{1.458497in}}%
\pgfpathlineto{\pgfqpoint{2.230369in}{1.458497in}}%
\pgfpathlineto{\pgfqpoint{2.230369in}{1.461446in}}%
\pgfpathlineto{\pgfqpoint{2.234910in}{1.461446in}}%
\pgfpathlineto{\pgfqpoint{2.234910in}{1.458497in}}%
\pgfpathmoveto{\pgfqpoint{2.230369in}{1.461446in}}%
\pgfpathlineto{\pgfqpoint{2.230369in}{1.461446in}}%
\pgfpathlineto{\pgfqpoint{2.230369in}{1.464395in}}%
\pgfpathlineto{\pgfqpoint{2.234910in}{1.464395in}}%
\pgfpathlineto{\pgfqpoint{2.234910in}{1.461446in}}%
\pgfpathmoveto{\pgfqpoint{2.234910in}{1.461446in}}%
\pgfpathlineto{\pgfqpoint{2.234910in}{1.461446in}}%
\pgfpathlineto{\pgfqpoint{2.234910in}{1.464395in}}%
\pgfpathlineto{\pgfqpoint{2.239451in}{1.464395in}}%
\pgfpathlineto{\pgfqpoint{2.239451in}{1.461446in}}%
\pgfpathmoveto{\pgfqpoint{2.234910in}{1.464395in}}%
\pgfpathlineto{\pgfqpoint{2.234910in}{1.464395in}}%
\pgfpathlineto{\pgfqpoint{2.234910in}{1.467344in}}%
\pgfpathlineto{\pgfqpoint{2.239451in}{1.467344in}}%
\pgfpathlineto{\pgfqpoint{2.239451in}{1.464395in}}%
\pgfpathmoveto{\pgfqpoint{2.239451in}{1.464395in}}%
\pgfpathlineto{\pgfqpoint{2.239451in}{1.464395in}}%
\pgfpathlineto{\pgfqpoint{2.239451in}{1.467344in}}%
\pgfpathlineto{\pgfqpoint{2.243992in}{1.467344in}}%
\pgfpathlineto{\pgfqpoint{2.243992in}{1.464395in}}%
\pgfpathmoveto{\pgfqpoint{2.239451in}{1.467344in}}%
\pgfpathlineto{\pgfqpoint{2.239451in}{1.467344in}}%
\pgfpathlineto{\pgfqpoint{2.239451in}{1.470294in}}%
\pgfpathlineto{\pgfqpoint{2.243992in}{1.470294in}}%
\pgfpathlineto{\pgfqpoint{2.243992in}{1.467344in}}%
\pgfpathmoveto{\pgfqpoint{2.243992in}{1.467344in}}%
\pgfpathlineto{\pgfqpoint{2.243992in}{1.467344in}}%
\pgfpathlineto{\pgfqpoint{2.243992in}{1.470294in}}%
\pgfpathlineto{\pgfqpoint{2.248533in}{1.470294in}}%
\pgfpathlineto{\pgfqpoint{2.248533in}{1.467344in}}%
\pgfpathmoveto{\pgfqpoint{2.243992in}{1.470294in}}%
\pgfpathlineto{\pgfqpoint{2.243992in}{1.470294in}}%
\pgfpathlineto{\pgfqpoint{2.243992in}{1.473243in}}%
\pgfpathlineto{\pgfqpoint{2.248533in}{1.473243in}}%
\pgfpathlineto{\pgfqpoint{2.248533in}{1.470294in}}%
\pgfpathmoveto{\pgfqpoint{2.248533in}{1.470294in}}%
\pgfpathlineto{\pgfqpoint{2.248533in}{1.470294in}}%
\pgfpathlineto{\pgfqpoint{2.248533in}{1.473243in}}%
\pgfpathlineto{\pgfqpoint{2.253074in}{1.473243in}}%
\pgfpathlineto{\pgfqpoint{2.253074in}{1.470294in}}%
\pgfpathmoveto{\pgfqpoint{2.248533in}{1.473243in}}%
\pgfpathlineto{\pgfqpoint{2.248533in}{1.473243in}}%
\pgfpathlineto{\pgfqpoint{2.248533in}{1.476192in}}%
\pgfpathlineto{\pgfqpoint{2.253074in}{1.476192in}}%
\pgfpathlineto{\pgfqpoint{2.253074in}{1.473243in}}%
\pgfpathmoveto{\pgfqpoint{2.253074in}{1.473243in}}%
\pgfpathlineto{\pgfqpoint{2.253074in}{1.473243in}}%
\pgfpathlineto{\pgfqpoint{2.253074in}{1.476192in}}%
\pgfpathlineto{\pgfqpoint{2.257615in}{1.476192in}}%
\pgfpathlineto{\pgfqpoint{2.257615in}{1.473243in}}%
\pgfpathmoveto{\pgfqpoint{2.253074in}{1.476192in}}%
\pgfpathlineto{\pgfqpoint{2.253074in}{1.476192in}}%
\pgfpathlineto{\pgfqpoint{2.253074in}{1.479142in}}%
\pgfpathlineto{\pgfqpoint{2.257615in}{1.479142in}}%
\pgfpathlineto{\pgfqpoint{2.257615in}{1.476192in}}%
\pgfpathmoveto{\pgfqpoint{2.257615in}{1.476192in}}%
\pgfpathlineto{\pgfqpoint{2.257615in}{1.476192in}}%
\pgfpathlineto{\pgfqpoint{2.257615in}{1.479142in}}%
\pgfpathlineto{\pgfqpoint{2.262156in}{1.479142in}}%
\pgfpathlineto{\pgfqpoint{2.262156in}{1.476192in}}%
\pgfpathmoveto{\pgfqpoint{2.257615in}{1.479142in}}%
\pgfpathlineto{\pgfqpoint{2.257615in}{1.479142in}}%
\pgfpathlineto{\pgfqpoint{2.257615in}{1.482091in}}%
\pgfpathlineto{\pgfqpoint{2.262156in}{1.482091in}}%
\pgfpathlineto{\pgfqpoint{2.262156in}{1.479142in}}%
\pgfpathmoveto{\pgfqpoint{2.262156in}{1.479142in}}%
\pgfpathlineto{\pgfqpoint{2.262156in}{1.479142in}}%
\pgfpathlineto{\pgfqpoint{2.262156in}{1.482091in}}%
\pgfpathlineto{\pgfqpoint{2.266697in}{1.482091in}}%
\pgfpathlineto{\pgfqpoint{2.266697in}{1.479142in}}%
\pgfpathmoveto{\pgfqpoint{2.262156in}{1.482091in}}%
\pgfpathlineto{\pgfqpoint{2.262156in}{1.482091in}}%
\pgfpathlineto{\pgfqpoint{2.262156in}{1.485040in}}%
\pgfpathlineto{\pgfqpoint{2.266697in}{1.485040in}}%
\pgfpathlineto{\pgfqpoint{2.266697in}{1.482091in}}%
\pgfpathmoveto{\pgfqpoint{2.266697in}{1.482091in}}%
\pgfpathlineto{\pgfqpoint{2.266697in}{1.482091in}}%
\pgfpathlineto{\pgfqpoint{2.266697in}{1.485040in}}%
\pgfpathlineto{\pgfqpoint{2.271238in}{1.485040in}}%
\pgfpathlineto{\pgfqpoint{2.271238in}{1.482091in}}%
\pgfpathmoveto{\pgfqpoint{2.266697in}{1.485040in}}%
\pgfpathlineto{\pgfqpoint{2.266697in}{1.485040in}}%
\pgfpathlineto{\pgfqpoint{2.266697in}{1.487989in}}%
\pgfpathlineto{\pgfqpoint{2.271238in}{1.487989in}}%
\pgfpathlineto{\pgfqpoint{2.271238in}{1.485040in}}%
\pgfpathmoveto{\pgfqpoint{2.271238in}{1.485040in}}%
\pgfpathlineto{\pgfqpoint{2.271238in}{1.485040in}}%
\pgfpathlineto{\pgfqpoint{2.271238in}{1.487989in}}%
\pgfpathlineto{\pgfqpoint{2.275779in}{1.487989in}}%
\pgfpathlineto{\pgfqpoint{2.275779in}{1.485040in}}%
\pgfpathmoveto{\pgfqpoint{2.271238in}{1.487989in}}%
\pgfpathlineto{\pgfqpoint{2.271238in}{1.487989in}}%
\pgfpathlineto{\pgfqpoint{2.271238in}{1.490939in}}%
\pgfpathlineto{\pgfqpoint{2.275779in}{1.490939in}}%
\pgfpathlineto{\pgfqpoint{2.275779in}{1.487989in}}%
\pgfpathmoveto{\pgfqpoint{2.275779in}{1.487989in}}%
\pgfpathlineto{\pgfqpoint{2.275779in}{1.487989in}}%
\pgfpathlineto{\pgfqpoint{2.275779in}{1.490939in}}%
\pgfpathlineto{\pgfqpoint{2.280320in}{1.490939in}}%
\pgfpathlineto{\pgfqpoint{2.280320in}{1.487989in}}%
\pgfpathmoveto{\pgfqpoint{2.275779in}{1.490939in}}%
\pgfpathlineto{\pgfqpoint{2.275779in}{1.490939in}}%
\pgfpathlineto{\pgfqpoint{2.275779in}{1.493888in}}%
\pgfpathlineto{\pgfqpoint{2.280320in}{1.493888in}}%
\pgfpathlineto{\pgfqpoint{2.280320in}{1.490939in}}%
\pgfpathmoveto{\pgfqpoint{2.280320in}{1.490939in}}%
\pgfpathlineto{\pgfqpoint{2.280320in}{1.490939in}}%
\pgfpathlineto{\pgfqpoint{2.280320in}{1.493888in}}%
\pgfpathlineto{\pgfqpoint{2.284861in}{1.493888in}}%
\pgfpathlineto{\pgfqpoint{2.284861in}{1.490939in}}%
\pgfpathmoveto{\pgfqpoint{2.280320in}{1.493888in}}%
\pgfpathlineto{\pgfqpoint{2.280320in}{1.493888in}}%
\pgfpathlineto{\pgfqpoint{2.280320in}{1.496837in}}%
\pgfpathlineto{\pgfqpoint{2.284861in}{1.496837in}}%
\pgfpathlineto{\pgfqpoint{2.284861in}{1.493888in}}%
\pgfpathmoveto{\pgfqpoint{2.284861in}{1.493888in}}%
\pgfpathlineto{\pgfqpoint{2.284861in}{1.493888in}}%
\pgfpathlineto{\pgfqpoint{2.284861in}{1.496837in}}%
\pgfpathlineto{\pgfqpoint{2.289402in}{1.496837in}}%
\pgfpathlineto{\pgfqpoint{2.289402in}{1.493888in}}%
\pgfpathmoveto{\pgfqpoint{2.284861in}{1.496837in}}%
\pgfpathlineto{\pgfqpoint{2.284861in}{1.496837in}}%
\pgfpathlineto{\pgfqpoint{2.284861in}{1.499786in}}%
\pgfpathlineto{\pgfqpoint{2.289402in}{1.499786in}}%
\pgfpathlineto{\pgfqpoint{2.289402in}{1.496837in}}%
\pgfpathmoveto{\pgfqpoint{2.289402in}{1.496837in}}%
\pgfpathlineto{\pgfqpoint{2.289402in}{1.496837in}}%
\pgfpathlineto{\pgfqpoint{2.289402in}{1.499786in}}%
\pgfpathlineto{\pgfqpoint{2.293943in}{1.499786in}}%
\pgfpathlineto{\pgfqpoint{2.293943in}{1.496837in}}%
\pgfpathmoveto{\pgfqpoint{2.289402in}{1.499786in}}%
\pgfpathlineto{\pgfqpoint{2.289402in}{1.499786in}}%
\pgfpathlineto{\pgfqpoint{2.289402in}{1.502736in}}%
\pgfpathlineto{\pgfqpoint{2.293943in}{1.502736in}}%
\pgfpathlineto{\pgfqpoint{2.293943in}{1.499786in}}%
\pgfpathmoveto{\pgfqpoint{2.293943in}{1.499786in}}%
\pgfpathlineto{\pgfqpoint{2.293943in}{1.499786in}}%
\pgfpathlineto{\pgfqpoint{2.293943in}{1.502736in}}%
\pgfpathlineto{\pgfqpoint{2.298484in}{1.502736in}}%
\pgfpathlineto{\pgfqpoint{2.298484in}{1.499786in}}%
\pgfpathmoveto{\pgfqpoint{2.293943in}{1.502736in}}%
\pgfpathlineto{\pgfqpoint{2.293943in}{1.502736in}}%
\pgfpathlineto{\pgfqpoint{2.293943in}{1.505685in}}%
\pgfpathlineto{\pgfqpoint{2.298484in}{1.505685in}}%
\pgfpathlineto{\pgfqpoint{2.298484in}{1.502736in}}%
\pgfpathmoveto{\pgfqpoint{2.298484in}{1.502736in}}%
\pgfpathlineto{\pgfqpoint{2.298484in}{1.502736in}}%
\pgfpathlineto{\pgfqpoint{2.298484in}{1.505685in}}%
\pgfpathlineto{\pgfqpoint{2.303025in}{1.505685in}}%
\pgfpathlineto{\pgfqpoint{2.303025in}{1.502736in}}%
\pgfpathmoveto{\pgfqpoint{2.298484in}{1.505685in}}%
\pgfpathlineto{\pgfqpoint{2.298484in}{1.505685in}}%
\pgfpathlineto{\pgfqpoint{2.298484in}{1.508634in}}%
\pgfpathlineto{\pgfqpoint{2.303025in}{1.508634in}}%
\pgfpathlineto{\pgfqpoint{2.303025in}{1.505685in}}%
\pgfpathmoveto{\pgfqpoint{2.303025in}{1.505685in}}%
\pgfpathlineto{\pgfqpoint{2.303025in}{1.505685in}}%
\pgfpathlineto{\pgfqpoint{2.303025in}{1.508634in}}%
\pgfpathlineto{\pgfqpoint{2.307566in}{1.508634in}}%
\pgfpathlineto{\pgfqpoint{2.307566in}{1.505685in}}%
\pgfpathmoveto{\pgfqpoint{2.303025in}{1.508634in}}%
\pgfpathlineto{\pgfqpoint{2.303025in}{1.508634in}}%
\pgfpathlineto{\pgfqpoint{2.303025in}{1.511583in}}%
\pgfpathlineto{\pgfqpoint{2.307566in}{1.511583in}}%
\pgfpathlineto{\pgfqpoint{2.307566in}{1.508634in}}%
\pgfpathmoveto{\pgfqpoint{2.307566in}{1.508634in}}%
\pgfpathlineto{\pgfqpoint{2.307566in}{1.508634in}}%
\pgfpathlineto{\pgfqpoint{2.307566in}{1.511583in}}%
\pgfpathlineto{\pgfqpoint{2.312107in}{1.511583in}}%
\pgfpathlineto{\pgfqpoint{2.312107in}{1.508634in}}%
\pgfpathmoveto{\pgfqpoint{2.307566in}{1.511583in}}%
\pgfpathlineto{\pgfqpoint{2.307566in}{1.511583in}}%
\pgfpathlineto{\pgfqpoint{2.307566in}{1.514533in}}%
\pgfpathlineto{\pgfqpoint{2.312107in}{1.514533in}}%
\pgfpathlineto{\pgfqpoint{2.312107in}{1.511583in}}%
\pgfpathmoveto{\pgfqpoint{2.312107in}{1.511583in}}%
\pgfpathlineto{\pgfqpoint{2.312107in}{1.511583in}}%
\pgfpathlineto{\pgfqpoint{2.312107in}{1.514533in}}%
\pgfpathlineto{\pgfqpoint{2.316648in}{1.514533in}}%
\pgfpathlineto{\pgfqpoint{2.316648in}{1.511583in}}%
\pgfpathmoveto{\pgfqpoint{2.312107in}{1.514533in}}%
\pgfpathlineto{\pgfqpoint{2.312107in}{1.514533in}}%
\pgfpathlineto{\pgfqpoint{2.312107in}{1.517482in}}%
\pgfpathlineto{\pgfqpoint{2.316648in}{1.517482in}}%
\pgfpathlineto{\pgfqpoint{2.316648in}{1.514533in}}%
\pgfpathmoveto{\pgfqpoint{2.316648in}{1.514533in}}%
\pgfpathlineto{\pgfqpoint{2.316648in}{1.514533in}}%
\pgfpathlineto{\pgfqpoint{2.316648in}{1.517482in}}%
\pgfpathlineto{\pgfqpoint{2.321189in}{1.517482in}}%
\pgfpathlineto{\pgfqpoint{2.321189in}{1.514533in}}%
\pgfpathmoveto{\pgfqpoint{2.316648in}{1.517482in}}%
\pgfpathlineto{\pgfqpoint{2.316648in}{1.517482in}}%
\pgfpathlineto{\pgfqpoint{2.316648in}{1.520431in}}%
\pgfpathlineto{\pgfqpoint{2.321189in}{1.520431in}}%
\pgfpathlineto{\pgfqpoint{2.321189in}{1.517482in}}%
\pgfpathmoveto{\pgfqpoint{2.321189in}{1.517482in}}%
\pgfpathlineto{\pgfqpoint{2.321189in}{1.517482in}}%
\pgfpathlineto{\pgfqpoint{2.321189in}{1.520431in}}%
\pgfpathlineto{\pgfqpoint{2.325731in}{1.520431in}}%
\pgfpathlineto{\pgfqpoint{2.325731in}{1.517482in}}%
\pgfpathmoveto{\pgfqpoint{2.321189in}{1.520431in}}%
\pgfpathlineto{\pgfqpoint{2.321189in}{1.520431in}}%
\pgfpathlineto{\pgfqpoint{2.321189in}{1.523380in}}%
\pgfpathlineto{\pgfqpoint{2.325731in}{1.523380in}}%
\pgfpathlineto{\pgfqpoint{2.325731in}{1.520431in}}%
\pgfpathmoveto{\pgfqpoint{2.325731in}{1.520431in}}%
\pgfpathlineto{\pgfqpoint{2.325731in}{1.520431in}}%
\pgfpathlineto{\pgfqpoint{2.325731in}{1.523380in}}%
\pgfpathlineto{\pgfqpoint{2.330272in}{1.523380in}}%
\pgfpathlineto{\pgfqpoint{2.330272in}{1.520431in}}%
\pgfpathmoveto{\pgfqpoint{2.325731in}{1.523380in}}%
\pgfpathlineto{\pgfqpoint{2.325731in}{1.523380in}}%
\pgfpathlineto{\pgfqpoint{2.325731in}{1.526330in}}%
\pgfpathlineto{\pgfqpoint{2.330272in}{1.526330in}}%
\pgfpathlineto{\pgfqpoint{2.330272in}{1.523380in}}%
\pgfpathmoveto{\pgfqpoint{2.330272in}{1.523380in}}%
\pgfpathlineto{\pgfqpoint{2.330272in}{1.523380in}}%
\pgfpathlineto{\pgfqpoint{2.330272in}{1.526330in}}%
\pgfpathlineto{\pgfqpoint{2.334813in}{1.526330in}}%
\pgfpathlineto{\pgfqpoint{2.334813in}{1.523380in}}%
\pgfpathmoveto{\pgfqpoint{2.330272in}{1.526330in}}%
\pgfpathlineto{\pgfqpoint{2.330272in}{1.526330in}}%
\pgfpathlineto{\pgfqpoint{2.330272in}{1.529279in}}%
\pgfpathlineto{\pgfqpoint{2.334813in}{1.529279in}}%
\pgfpathlineto{\pgfqpoint{2.334813in}{1.526330in}}%
\pgfpathmoveto{\pgfqpoint{2.334813in}{1.526330in}}%
\pgfpathlineto{\pgfqpoint{2.334813in}{1.526330in}}%
\pgfpathlineto{\pgfqpoint{2.334813in}{1.529279in}}%
\pgfpathlineto{\pgfqpoint{2.339354in}{1.529279in}}%
\pgfpathlineto{\pgfqpoint{2.339354in}{1.526330in}}%
\pgfpathmoveto{\pgfqpoint{2.334813in}{1.529279in}}%
\pgfpathlineto{\pgfqpoint{2.334813in}{1.529279in}}%
\pgfpathlineto{\pgfqpoint{2.334813in}{1.532228in}}%
\pgfpathlineto{\pgfqpoint{2.339354in}{1.532228in}}%
\pgfpathlineto{\pgfqpoint{2.339354in}{1.529279in}}%
\pgfpathmoveto{\pgfqpoint{2.339354in}{1.529279in}}%
\pgfpathlineto{\pgfqpoint{2.339354in}{1.529279in}}%
\pgfpathlineto{\pgfqpoint{2.339354in}{1.532228in}}%
\pgfpathlineto{\pgfqpoint{2.343895in}{1.532228in}}%
\pgfpathlineto{\pgfqpoint{2.343895in}{1.529279in}}%
\pgfpathmoveto{\pgfqpoint{2.339354in}{1.532228in}}%
\pgfpathlineto{\pgfqpoint{2.339354in}{1.532228in}}%
\pgfpathlineto{\pgfqpoint{2.339354in}{1.535177in}}%
\pgfpathlineto{\pgfqpoint{2.343895in}{1.535177in}}%
\pgfpathlineto{\pgfqpoint{2.343895in}{1.532228in}}%
\pgfpathmoveto{\pgfqpoint{2.343895in}{1.532228in}}%
\pgfpathlineto{\pgfqpoint{2.343895in}{1.532228in}}%
\pgfpathlineto{\pgfqpoint{2.343895in}{1.535177in}}%
\pgfpathlineto{\pgfqpoint{2.348436in}{1.535177in}}%
\pgfpathlineto{\pgfqpoint{2.348436in}{1.532228in}}%
\pgfpathmoveto{\pgfqpoint{2.343895in}{1.535177in}}%
\pgfpathlineto{\pgfqpoint{2.343895in}{1.535177in}}%
\pgfpathlineto{\pgfqpoint{2.343895in}{1.538127in}}%
\pgfpathlineto{\pgfqpoint{2.348436in}{1.538127in}}%
\pgfpathlineto{\pgfqpoint{2.348436in}{1.535177in}}%
\pgfpathmoveto{\pgfqpoint{2.348436in}{1.535177in}}%
\pgfpathlineto{\pgfqpoint{2.348436in}{1.535177in}}%
\pgfpathlineto{\pgfqpoint{2.348436in}{1.538127in}}%
\pgfpathlineto{\pgfqpoint{2.352977in}{1.538127in}}%
\pgfpathlineto{\pgfqpoint{2.352977in}{1.535177in}}%
\pgfpathmoveto{\pgfqpoint{2.348436in}{1.538127in}}%
\pgfpathlineto{\pgfqpoint{2.348436in}{1.538127in}}%
\pgfpathlineto{\pgfqpoint{2.348436in}{1.541076in}}%
\pgfpathlineto{\pgfqpoint{2.352977in}{1.541076in}}%
\pgfpathlineto{\pgfqpoint{2.352977in}{1.538127in}}%
\pgfpathmoveto{\pgfqpoint{2.352977in}{1.538127in}}%
\pgfpathlineto{\pgfqpoint{2.352977in}{1.538127in}}%
\pgfpathlineto{\pgfqpoint{2.352977in}{1.541076in}}%
\pgfpathlineto{\pgfqpoint{2.357518in}{1.541076in}}%
\pgfpathlineto{\pgfqpoint{2.357518in}{1.538127in}}%
\pgfpathmoveto{\pgfqpoint{2.352977in}{1.541076in}}%
\pgfpathlineto{\pgfqpoint{2.352977in}{1.541076in}}%
\pgfpathlineto{\pgfqpoint{2.352977in}{1.544025in}}%
\pgfpathlineto{\pgfqpoint{2.357518in}{1.544025in}}%
\pgfpathlineto{\pgfqpoint{2.357518in}{1.541076in}}%
\pgfpathmoveto{\pgfqpoint{2.357518in}{1.541076in}}%
\pgfpathlineto{\pgfqpoint{2.357518in}{1.541076in}}%
\pgfpathlineto{\pgfqpoint{2.357518in}{1.544025in}}%
\pgfpathlineto{\pgfqpoint{2.362059in}{1.544025in}}%
\pgfpathlineto{\pgfqpoint{2.362059in}{1.541076in}}%
\pgfpathmoveto{\pgfqpoint{2.357518in}{1.544025in}}%
\pgfpathlineto{\pgfqpoint{2.357518in}{1.544025in}}%
\pgfpathlineto{\pgfqpoint{2.357518in}{1.546974in}}%
\pgfpathlineto{\pgfqpoint{2.362059in}{1.546974in}}%
\pgfpathlineto{\pgfqpoint{2.362059in}{1.544025in}}%
\pgfpathmoveto{\pgfqpoint{2.362059in}{1.544025in}}%
\pgfpathlineto{\pgfqpoint{2.362059in}{1.544025in}}%
\pgfpathlineto{\pgfqpoint{2.362059in}{1.546974in}}%
\pgfpathlineto{\pgfqpoint{2.366600in}{1.546974in}}%
\pgfpathlineto{\pgfqpoint{2.366600in}{1.544025in}}%
\pgfpathmoveto{\pgfqpoint{2.362059in}{1.546974in}}%
\pgfpathlineto{\pgfqpoint{2.362059in}{1.546974in}}%
\pgfpathlineto{\pgfqpoint{2.362059in}{1.549923in}}%
\pgfpathlineto{\pgfqpoint{2.366600in}{1.549923in}}%
\pgfpathlineto{\pgfqpoint{2.366600in}{1.546974in}}%
\pgfpathmoveto{\pgfqpoint{2.366600in}{1.546974in}}%
\pgfpathlineto{\pgfqpoint{2.366600in}{1.546974in}}%
\pgfpathlineto{\pgfqpoint{2.366600in}{1.549923in}}%
\pgfpathlineto{\pgfqpoint{2.371141in}{1.549923in}}%
\pgfpathlineto{\pgfqpoint{2.371141in}{1.546974in}}%
\pgfpathmoveto{\pgfqpoint{2.366600in}{1.549923in}}%
\pgfpathlineto{\pgfqpoint{2.366600in}{1.549923in}}%
\pgfpathlineto{\pgfqpoint{2.366600in}{1.552873in}}%
\pgfpathlineto{\pgfqpoint{2.371141in}{1.552873in}}%
\pgfpathlineto{\pgfqpoint{2.371141in}{1.549923in}}%
\pgfpathmoveto{\pgfqpoint{2.371141in}{1.549923in}}%
\pgfpathlineto{\pgfqpoint{2.371141in}{1.549923in}}%
\pgfpathlineto{\pgfqpoint{2.371141in}{1.552873in}}%
\pgfpathlineto{\pgfqpoint{2.375682in}{1.552873in}}%
\pgfpathlineto{\pgfqpoint{2.375682in}{1.549923in}}%
\pgfpathmoveto{\pgfqpoint{2.371141in}{1.552873in}}%
\pgfpathlineto{\pgfqpoint{2.371141in}{1.552873in}}%
\pgfpathlineto{\pgfqpoint{2.371141in}{1.555822in}}%
\pgfpathlineto{\pgfqpoint{2.375682in}{1.555822in}}%
\pgfpathlineto{\pgfqpoint{2.375682in}{1.552873in}}%
\pgfpathmoveto{\pgfqpoint{2.375682in}{1.552873in}}%
\pgfpathlineto{\pgfqpoint{2.375682in}{1.552873in}}%
\pgfpathlineto{\pgfqpoint{2.375682in}{1.555822in}}%
\pgfpathlineto{\pgfqpoint{2.380223in}{1.555822in}}%
\pgfpathlineto{\pgfqpoint{2.380223in}{1.552873in}}%
\pgfpathmoveto{\pgfqpoint{2.375682in}{1.555822in}}%
\pgfpathlineto{\pgfqpoint{2.375682in}{1.555822in}}%
\pgfpathlineto{\pgfqpoint{2.375682in}{1.558771in}}%
\pgfpathlineto{\pgfqpoint{2.380223in}{1.558771in}}%
\pgfpathlineto{\pgfqpoint{2.380223in}{1.555822in}}%
\pgfpathmoveto{\pgfqpoint{2.380223in}{1.555822in}}%
\pgfpathlineto{\pgfqpoint{2.380223in}{1.555822in}}%
\pgfpathlineto{\pgfqpoint{2.380223in}{1.558771in}}%
\pgfpathlineto{\pgfqpoint{2.384764in}{1.558771in}}%
\pgfpathlineto{\pgfqpoint{2.384764in}{1.555822in}}%
\pgfpathmoveto{\pgfqpoint{2.380223in}{1.558771in}}%
\pgfpathlineto{\pgfqpoint{2.380223in}{1.558771in}}%
\pgfpathlineto{\pgfqpoint{2.380223in}{1.561720in}}%
\pgfpathlineto{\pgfqpoint{2.384764in}{1.561720in}}%
\pgfpathlineto{\pgfqpoint{2.384764in}{1.558771in}}%
\pgfpathmoveto{\pgfqpoint{2.384764in}{1.558771in}}%
\pgfpathlineto{\pgfqpoint{2.384764in}{1.558771in}}%
\pgfpathlineto{\pgfqpoint{2.384764in}{1.561720in}}%
\pgfpathlineto{\pgfqpoint{2.389305in}{1.561720in}}%
\pgfpathlineto{\pgfqpoint{2.389305in}{1.558771in}}%
\pgfpathmoveto{\pgfqpoint{2.384764in}{1.561720in}}%
\pgfpathlineto{\pgfqpoint{2.384764in}{1.561720in}}%
\pgfpathlineto{\pgfqpoint{2.384764in}{1.564669in}}%
\pgfpathlineto{\pgfqpoint{2.389305in}{1.564669in}}%
\pgfpathlineto{\pgfqpoint{2.389305in}{1.561720in}}%
\pgfpathmoveto{\pgfqpoint{2.389305in}{1.561720in}}%
\pgfpathlineto{\pgfqpoint{2.389305in}{1.561720in}}%
\pgfpathlineto{\pgfqpoint{2.389305in}{1.564669in}}%
\pgfpathlineto{\pgfqpoint{2.393846in}{1.564669in}}%
\pgfpathlineto{\pgfqpoint{2.393846in}{1.561720in}}%
\pgfpathmoveto{\pgfqpoint{2.389305in}{1.564669in}}%
\pgfpathlineto{\pgfqpoint{2.389305in}{1.564669in}}%
\pgfpathlineto{\pgfqpoint{2.389305in}{1.567619in}}%
\pgfpathlineto{\pgfqpoint{2.393846in}{1.567619in}}%
\pgfpathlineto{\pgfqpoint{2.393846in}{1.564669in}}%
\pgfpathmoveto{\pgfqpoint{2.393846in}{1.564669in}}%
\pgfpathlineto{\pgfqpoint{2.393846in}{1.564669in}}%
\pgfpathlineto{\pgfqpoint{2.393846in}{1.567619in}}%
\pgfpathlineto{\pgfqpoint{2.398387in}{1.567619in}}%
\pgfpathlineto{\pgfqpoint{2.398387in}{1.564669in}}%
\pgfpathmoveto{\pgfqpoint{2.393846in}{1.567619in}}%
\pgfpathlineto{\pgfqpoint{2.393846in}{1.567619in}}%
\pgfpathlineto{\pgfqpoint{2.393846in}{1.570568in}}%
\pgfpathlineto{\pgfqpoint{2.398387in}{1.570568in}}%
\pgfpathlineto{\pgfqpoint{2.398387in}{1.567619in}}%
\pgfpathmoveto{\pgfqpoint{2.398387in}{1.567619in}}%
\pgfpathlineto{\pgfqpoint{2.398387in}{1.567619in}}%
\pgfpathlineto{\pgfqpoint{2.398387in}{1.570568in}}%
\pgfpathlineto{\pgfqpoint{2.402928in}{1.570568in}}%
\pgfpathlineto{\pgfqpoint{2.402928in}{1.567619in}}%
\pgfpathmoveto{\pgfqpoint{2.398387in}{1.570568in}}%
\pgfpathlineto{\pgfqpoint{2.398387in}{1.570568in}}%
\pgfpathlineto{\pgfqpoint{2.398387in}{1.573517in}}%
\pgfpathlineto{\pgfqpoint{2.402928in}{1.573517in}}%
\pgfpathlineto{\pgfqpoint{2.402928in}{1.570568in}}%
\pgfpathmoveto{\pgfqpoint{2.402928in}{1.570568in}}%
\pgfpathlineto{\pgfqpoint{2.402928in}{1.570568in}}%
\pgfpathlineto{\pgfqpoint{2.402928in}{1.573517in}}%
\pgfpathlineto{\pgfqpoint{2.407469in}{1.573517in}}%
\pgfpathlineto{\pgfqpoint{2.407469in}{1.570568in}}%
\pgfpathmoveto{\pgfqpoint{2.402928in}{1.573517in}}%
\pgfpathlineto{\pgfqpoint{2.402928in}{1.573517in}}%
\pgfpathlineto{\pgfqpoint{2.402928in}{1.576466in}}%
\pgfpathlineto{\pgfqpoint{2.407469in}{1.576466in}}%
\pgfpathlineto{\pgfqpoint{2.407469in}{1.573517in}}%
\pgfpathmoveto{\pgfqpoint{2.407469in}{1.573517in}}%
\pgfpathlineto{\pgfqpoint{2.407469in}{1.573517in}}%
\pgfpathlineto{\pgfqpoint{2.407469in}{1.576466in}}%
\pgfpathlineto{\pgfqpoint{2.412010in}{1.576466in}}%
\pgfpathlineto{\pgfqpoint{2.412010in}{1.573517in}}%
\pgfpathmoveto{\pgfqpoint{2.407469in}{1.576466in}}%
\pgfpathlineto{\pgfqpoint{2.407469in}{1.576466in}}%
\pgfpathlineto{\pgfqpoint{2.407469in}{1.579415in}}%
\pgfpathlineto{\pgfqpoint{2.412010in}{1.579415in}}%
\pgfpathlineto{\pgfqpoint{2.412010in}{1.576466in}}%
\pgfpathmoveto{\pgfqpoint{2.412010in}{1.576466in}}%
\pgfpathlineto{\pgfqpoint{2.412010in}{1.576466in}}%
\pgfpathlineto{\pgfqpoint{2.412010in}{1.579415in}}%
\pgfpathlineto{\pgfqpoint{2.416551in}{1.579415in}}%
\pgfpathlineto{\pgfqpoint{2.416551in}{1.576466in}}%
\pgfpathmoveto{\pgfqpoint{2.412010in}{1.579415in}}%
\pgfpathlineto{\pgfqpoint{2.412010in}{1.579415in}}%
\pgfpathlineto{\pgfqpoint{2.412010in}{1.582364in}}%
\pgfpathlineto{\pgfqpoint{2.416551in}{1.582364in}}%
\pgfpathlineto{\pgfqpoint{2.416551in}{1.579415in}}%
\pgfpathmoveto{\pgfqpoint{2.416551in}{1.579415in}}%
\pgfpathlineto{\pgfqpoint{2.416551in}{1.579415in}}%
\pgfpathlineto{\pgfqpoint{2.416551in}{1.582364in}}%
\pgfpathlineto{\pgfqpoint{2.421092in}{1.582364in}}%
\pgfpathlineto{\pgfqpoint{2.421092in}{1.579415in}}%
\pgfpathmoveto{\pgfqpoint{2.416551in}{1.582364in}}%
\pgfpathlineto{\pgfqpoint{2.416551in}{1.582364in}}%
\pgfpathlineto{\pgfqpoint{2.416551in}{1.585314in}}%
\pgfpathlineto{\pgfqpoint{2.421092in}{1.585314in}}%
\pgfpathlineto{\pgfqpoint{2.421092in}{1.582364in}}%
\pgfpathmoveto{\pgfqpoint{2.421092in}{1.582364in}}%
\pgfpathlineto{\pgfqpoint{2.421092in}{1.582364in}}%
\pgfpathlineto{\pgfqpoint{2.421092in}{1.585314in}}%
\pgfpathlineto{\pgfqpoint{2.425633in}{1.585314in}}%
\pgfpathlineto{\pgfqpoint{2.425633in}{1.582364in}}%
\pgfpathmoveto{\pgfqpoint{2.421092in}{1.585314in}}%
\pgfpathlineto{\pgfqpoint{2.421092in}{1.585314in}}%
\pgfpathlineto{\pgfqpoint{2.421092in}{1.588263in}}%
\pgfpathlineto{\pgfqpoint{2.425633in}{1.588263in}}%
\pgfpathlineto{\pgfqpoint{2.425633in}{1.585314in}}%
\pgfpathmoveto{\pgfqpoint{2.425633in}{1.585314in}}%
\pgfpathlineto{\pgfqpoint{2.425633in}{1.585314in}}%
\pgfpathlineto{\pgfqpoint{2.425633in}{1.588263in}}%
\pgfpathlineto{\pgfqpoint{2.430174in}{1.588263in}}%
\pgfpathlineto{\pgfqpoint{2.430174in}{1.585314in}}%
\pgfpathmoveto{\pgfqpoint{2.425633in}{1.588263in}}%
\pgfpathlineto{\pgfqpoint{2.425633in}{1.588263in}}%
\pgfpathlineto{\pgfqpoint{2.425633in}{1.591212in}}%
\pgfpathlineto{\pgfqpoint{2.430174in}{1.591212in}}%
\pgfpathlineto{\pgfqpoint{2.430174in}{1.588263in}}%
\pgfpathmoveto{\pgfqpoint{2.430174in}{1.588263in}}%
\pgfpathlineto{\pgfqpoint{2.430174in}{1.588263in}}%
\pgfpathlineto{\pgfqpoint{2.430174in}{1.591212in}}%
\pgfpathlineto{\pgfqpoint{2.434715in}{1.591212in}}%
\pgfpathlineto{\pgfqpoint{2.434715in}{1.588263in}}%
\pgfpathmoveto{\pgfqpoint{2.430174in}{1.591212in}}%
\pgfpathlineto{\pgfqpoint{2.430174in}{1.591212in}}%
\pgfpathlineto{\pgfqpoint{2.430174in}{1.594161in}}%
\pgfpathlineto{\pgfqpoint{2.434715in}{1.594161in}}%
\pgfpathlineto{\pgfqpoint{2.434715in}{1.591212in}}%
\pgfpathmoveto{\pgfqpoint{2.434715in}{1.591212in}}%
\pgfpathlineto{\pgfqpoint{2.434715in}{1.591212in}}%
\pgfpathlineto{\pgfqpoint{2.434715in}{1.594161in}}%
\pgfpathlineto{\pgfqpoint{2.439256in}{1.594161in}}%
\pgfpathlineto{\pgfqpoint{2.439256in}{1.591212in}}%
\pgfpathmoveto{\pgfqpoint{2.434715in}{1.594161in}}%
\pgfpathlineto{\pgfqpoint{2.434715in}{1.594161in}}%
\pgfpathlineto{\pgfqpoint{2.434715in}{1.597110in}}%
\pgfpathlineto{\pgfqpoint{2.439256in}{1.597110in}}%
\pgfpathlineto{\pgfqpoint{2.439256in}{1.594161in}}%
\pgfpathmoveto{\pgfqpoint{2.439256in}{1.594161in}}%
\pgfpathlineto{\pgfqpoint{2.439256in}{1.594161in}}%
\pgfpathlineto{\pgfqpoint{2.439256in}{1.597110in}}%
\pgfpathlineto{\pgfqpoint{2.443797in}{1.597110in}}%
\pgfpathlineto{\pgfqpoint{2.443797in}{1.594161in}}%
\pgfpathmoveto{\pgfqpoint{2.439256in}{1.597110in}}%
\pgfpathlineto{\pgfqpoint{2.439256in}{1.597110in}}%
\pgfpathlineto{\pgfqpoint{2.439256in}{1.600060in}}%
\pgfpathlineto{\pgfqpoint{2.443797in}{1.600060in}}%
\pgfpathlineto{\pgfqpoint{2.443797in}{1.597110in}}%
\pgfpathmoveto{\pgfqpoint{2.443797in}{1.597110in}}%
\pgfpathlineto{\pgfqpoint{2.443797in}{1.597110in}}%
\pgfpathlineto{\pgfqpoint{2.443797in}{1.600060in}}%
\pgfpathlineto{\pgfqpoint{2.448338in}{1.600060in}}%
\pgfpathlineto{\pgfqpoint{2.448338in}{1.597110in}}%
\pgfpathmoveto{\pgfqpoint{2.443797in}{1.600060in}}%
\pgfpathlineto{\pgfqpoint{2.443797in}{1.600060in}}%
\pgfpathlineto{\pgfqpoint{2.443797in}{1.603009in}}%
\pgfpathlineto{\pgfqpoint{2.448338in}{1.603009in}}%
\pgfpathlineto{\pgfqpoint{2.448338in}{1.600060in}}%
\pgfpathmoveto{\pgfqpoint{2.448338in}{1.600060in}}%
\pgfpathlineto{\pgfqpoint{2.448338in}{1.600060in}}%
\pgfpathlineto{\pgfqpoint{2.448338in}{1.603009in}}%
\pgfpathlineto{\pgfqpoint{2.452879in}{1.603009in}}%
\pgfpathlineto{\pgfqpoint{2.452879in}{1.600060in}}%
\pgfpathmoveto{\pgfqpoint{2.448338in}{1.603009in}}%
\pgfpathlineto{\pgfqpoint{2.448338in}{1.603009in}}%
\pgfpathlineto{\pgfqpoint{2.448338in}{1.605958in}}%
\pgfpathlineto{\pgfqpoint{2.452879in}{1.605958in}}%
\pgfpathlineto{\pgfqpoint{2.452879in}{1.603009in}}%
\pgfpathmoveto{\pgfqpoint{2.452879in}{1.603009in}}%
\pgfpathlineto{\pgfqpoint{2.452879in}{1.603009in}}%
\pgfpathlineto{\pgfqpoint{2.452879in}{1.605958in}}%
\pgfpathlineto{\pgfqpoint{2.457420in}{1.605958in}}%
\pgfpathlineto{\pgfqpoint{2.457420in}{1.603009in}}%
\pgfpathmoveto{\pgfqpoint{2.452879in}{1.605958in}}%
\pgfpathlineto{\pgfqpoint{2.452879in}{1.605958in}}%
\pgfpathlineto{\pgfqpoint{2.452879in}{1.608907in}}%
\pgfpathlineto{\pgfqpoint{2.457420in}{1.608907in}}%
\pgfpathlineto{\pgfqpoint{2.457420in}{1.605958in}}%
\pgfpathmoveto{\pgfqpoint{2.457420in}{1.605958in}}%
\pgfpathlineto{\pgfqpoint{2.457420in}{1.605958in}}%
\pgfpathlineto{\pgfqpoint{2.457420in}{1.608907in}}%
\pgfpathlineto{\pgfqpoint{2.461961in}{1.608907in}}%
\pgfpathlineto{\pgfqpoint{2.461961in}{1.605958in}}%
\pgfpathmoveto{\pgfqpoint{2.457420in}{1.608907in}}%
\pgfpathlineto{\pgfqpoint{2.457420in}{1.608907in}}%
\pgfpathlineto{\pgfqpoint{2.457420in}{1.611856in}}%
\pgfpathlineto{\pgfqpoint{2.461961in}{1.611856in}}%
\pgfpathlineto{\pgfqpoint{2.461961in}{1.608907in}}%
\pgfpathmoveto{\pgfqpoint{2.461961in}{1.608907in}}%
\pgfpathlineto{\pgfqpoint{2.461961in}{1.608907in}}%
\pgfpathlineto{\pgfqpoint{2.461961in}{1.611856in}}%
\pgfpathlineto{\pgfqpoint{2.466502in}{1.611856in}}%
\pgfpathlineto{\pgfqpoint{2.466502in}{1.608907in}}%
\pgfpathmoveto{\pgfqpoint{2.461961in}{1.611856in}}%
\pgfpathlineto{\pgfqpoint{2.461961in}{1.611856in}}%
\pgfpathlineto{\pgfqpoint{2.461961in}{1.614805in}}%
\pgfpathlineto{\pgfqpoint{2.466502in}{1.614805in}}%
\pgfpathlineto{\pgfqpoint{2.466502in}{1.611856in}}%
\pgfpathmoveto{\pgfqpoint{2.466502in}{1.611856in}}%
\pgfpathlineto{\pgfqpoint{2.466502in}{1.611856in}}%
\pgfpathlineto{\pgfqpoint{2.466502in}{1.614805in}}%
\pgfpathlineto{\pgfqpoint{2.471043in}{1.614805in}}%
\pgfpathlineto{\pgfqpoint{2.471043in}{1.611856in}}%
\pgfpathmoveto{\pgfqpoint{2.466502in}{1.614805in}}%
\pgfpathlineto{\pgfqpoint{2.466502in}{1.614805in}}%
\pgfpathlineto{\pgfqpoint{2.466502in}{1.617755in}}%
\pgfpathlineto{\pgfqpoint{2.471043in}{1.617755in}}%
\pgfpathlineto{\pgfqpoint{2.471043in}{1.614805in}}%
\pgfpathmoveto{\pgfqpoint{2.471043in}{1.614805in}}%
\pgfpathlineto{\pgfqpoint{2.471043in}{1.614805in}}%
\pgfpathlineto{\pgfqpoint{2.471043in}{1.617755in}}%
\pgfpathlineto{\pgfqpoint{2.475584in}{1.617755in}}%
\pgfpathlineto{\pgfqpoint{2.475584in}{1.614805in}}%
\pgfpathmoveto{\pgfqpoint{2.471043in}{1.617755in}}%
\pgfpathlineto{\pgfqpoint{2.471043in}{1.617755in}}%
\pgfpathlineto{\pgfqpoint{2.471043in}{1.620704in}}%
\pgfpathlineto{\pgfqpoint{2.475584in}{1.620704in}}%
\pgfpathlineto{\pgfqpoint{2.475584in}{1.617755in}}%
\pgfpathmoveto{\pgfqpoint{2.475584in}{1.617755in}}%
\pgfpathlineto{\pgfqpoint{2.475584in}{1.617755in}}%
\pgfpathlineto{\pgfqpoint{2.475584in}{1.620704in}}%
\pgfpathlineto{\pgfqpoint{2.480125in}{1.620704in}}%
\pgfpathlineto{\pgfqpoint{2.480125in}{1.617755in}}%
\pgfpathmoveto{\pgfqpoint{2.475584in}{1.620704in}}%
\pgfpathlineto{\pgfqpoint{2.475584in}{1.620704in}}%
\pgfpathlineto{\pgfqpoint{2.475584in}{1.623653in}}%
\pgfpathlineto{\pgfqpoint{2.480125in}{1.623653in}}%
\pgfpathlineto{\pgfqpoint{2.480125in}{1.620704in}}%
\pgfpathmoveto{\pgfqpoint{2.480125in}{1.620704in}}%
\pgfpathlineto{\pgfqpoint{2.480125in}{1.620704in}}%
\pgfpathlineto{\pgfqpoint{2.480125in}{1.623653in}}%
\pgfpathlineto{\pgfqpoint{2.484666in}{1.623653in}}%
\pgfpathlineto{\pgfqpoint{2.484666in}{1.620704in}}%
\pgfpathmoveto{\pgfqpoint{2.480125in}{1.623653in}}%
\pgfpathlineto{\pgfqpoint{2.480125in}{1.623653in}}%
\pgfpathlineto{\pgfqpoint{2.480125in}{1.626602in}}%
\pgfpathlineto{\pgfqpoint{2.484666in}{1.626602in}}%
\pgfpathlineto{\pgfqpoint{2.484666in}{1.623653in}}%
\pgfpathmoveto{\pgfqpoint{2.484666in}{1.623653in}}%
\pgfpathlineto{\pgfqpoint{2.484666in}{1.623653in}}%
\pgfpathlineto{\pgfqpoint{2.484666in}{1.626602in}}%
\pgfpathlineto{\pgfqpoint{2.489207in}{1.626602in}}%
\pgfpathlineto{\pgfqpoint{2.489207in}{1.623653in}}%
\pgfpathmoveto{\pgfqpoint{2.484666in}{1.626602in}}%
\pgfpathlineto{\pgfqpoint{2.484666in}{1.626602in}}%
\pgfpathlineto{\pgfqpoint{2.484666in}{1.629551in}}%
\pgfpathlineto{\pgfqpoint{2.489207in}{1.629551in}}%
\pgfpathlineto{\pgfqpoint{2.489207in}{1.626602in}}%
\pgfpathmoveto{\pgfqpoint{2.489207in}{1.626602in}}%
\pgfpathlineto{\pgfqpoint{2.489207in}{1.626602in}}%
\pgfpathlineto{\pgfqpoint{2.489207in}{1.629551in}}%
\pgfpathlineto{\pgfqpoint{2.493748in}{1.629551in}}%
\pgfpathlineto{\pgfqpoint{2.493748in}{1.626602in}}%
\pgfpathmoveto{\pgfqpoint{2.489207in}{1.629551in}}%
\pgfpathlineto{\pgfqpoint{2.489207in}{1.629551in}}%
\pgfpathlineto{\pgfqpoint{2.489207in}{1.632501in}}%
\pgfpathlineto{\pgfqpoint{2.493748in}{1.632501in}}%
\pgfpathlineto{\pgfqpoint{2.493748in}{1.629551in}}%
\pgfpathmoveto{\pgfqpoint{2.493748in}{1.629551in}}%
\pgfpathlineto{\pgfqpoint{2.493748in}{1.629551in}}%
\pgfpathlineto{\pgfqpoint{2.493748in}{1.632501in}}%
\pgfpathlineto{\pgfqpoint{2.498289in}{1.632501in}}%
\pgfpathlineto{\pgfqpoint{2.498289in}{1.629551in}}%
\pgfpathmoveto{\pgfqpoint{2.493748in}{1.632501in}}%
\pgfpathlineto{\pgfqpoint{2.493748in}{1.632501in}}%
\pgfpathlineto{\pgfqpoint{2.493748in}{1.635450in}}%
\pgfpathlineto{\pgfqpoint{2.498289in}{1.635450in}}%
\pgfpathlineto{\pgfqpoint{2.498289in}{1.632501in}}%
\pgfpathmoveto{\pgfqpoint{2.498289in}{1.632501in}}%
\pgfpathlineto{\pgfqpoint{2.498289in}{1.632501in}}%
\pgfpathlineto{\pgfqpoint{2.498289in}{1.635450in}}%
\pgfpathlineto{\pgfqpoint{2.502830in}{1.635450in}}%
\pgfpathlineto{\pgfqpoint{2.502830in}{1.632501in}}%
\pgfpathmoveto{\pgfqpoint{2.498289in}{1.635450in}}%
\pgfpathlineto{\pgfqpoint{2.498289in}{1.635450in}}%
\pgfpathlineto{\pgfqpoint{2.498289in}{1.638399in}}%
\pgfpathlineto{\pgfqpoint{2.502830in}{1.638399in}}%
\pgfpathlineto{\pgfqpoint{2.502830in}{1.635450in}}%
\pgfpathmoveto{\pgfqpoint{2.502830in}{1.635450in}}%
\pgfpathlineto{\pgfqpoint{2.502830in}{1.635450in}}%
\pgfpathlineto{\pgfqpoint{2.502830in}{1.638399in}}%
\pgfpathlineto{\pgfqpoint{2.507371in}{1.638399in}}%
\pgfpathlineto{\pgfqpoint{2.507371in}{1.635450in}}%
\pgfpathmoveto{\pgfqpoint{2.502830in}{1.638399in}}%
\pgfpathlineto{\pgfqpoint{2.502830in}{1.638399in}}%
\pgfpathlineto{\pgfqpoint{2.502830in}{1.641348in}}%
\pgfpathlineto{\pgfqpoint{2.507371in}{1.641348in}}%
\pgfpathlineto{\pgfqpoint{2.507371in}{1.638399in}}%
\pgfpathmoveto{\pgfqpoint{2.507371in}{1.638399in}}%
\pgfpathlineto{\pgfqpoint{2.507371in}{1.638399in}}%
\pgfpathlineto{\pgfqpoint{2.507371in}{1.641348in}}%
\pgfpathlineto{\pgfqpoint{2.511912in}{1.641348in}}%
\pgfpathlineto{\pgfqpoint{2.511912in}{1.638399in}}%
\pgfpathmoveto{\pgfqpoint{2.507371in}{1.641348in}}%
\pgfpathlineto{\pgfqpoint{2.507371in}{1.641348in}}%
\pgfpathlineto{\pgfqpoint{2.507371in}{1.644298in}}%
\pgfpathlineto{\pgfqpoint{2.511912in}{1.644298in}}%
\pgfpathlineto{\pgfqpoint{2.511912in}{1.641348in}}%
\pgfpathmoveto{\pgfqpoint{2.511912in}{1.641348in}}%
\pgfpathlineto{\pgfqpoint{2.511912in}{1.641348in}}%
\pgfpathlineto{\pgfqpoint{2.511912in}{1.644298in}}%
\pgfpathlineto{\pgfqpoint{2.516453in}{1.644298in}}%
\pgfpathlineto{\pgfqpoint{2.516453in}{1.641348in}}%
\pgfpathmoveto{\pgfqpoint{2.511912in}{1.644298in}}%
\pgfpathlineto{\pgfqpoint{2.511912in}{1.644298in}}%
\pgfpathlineto{\pgfqpoint{2.511912in}{1.647247in}}%
\pgfpathlineto{\pgfqpoint{2.516453in}{1.647247in}}%
\pgfpathlineto{\pgfqpoint{2.516453in}{1.644298in}}%
\pgfpathmoveto{\pgfqpoint{2.516453in}{1.644298in}}%
\pgfpathlineto{\pgfqpoint{2.516453in}{1.644298in}}%
\pgfpathlineto{\pgfqpoint{2.516453in}{1.647247in}}%
\pgfpathlineto{\pgfqpoint{2.520994in}{1.647247in}}%
\pgfpathlineto{\pgfqpoint{2.520994in}{1.644298in}}%
\pgfpathmoveto{\pgfqpoint{2.516453in}{1.647247in}}%
\pgfpathlineto{\pgfqpoint{2.516453in}{1.647247in}}%
\pgfpathlineto{\pgfqpoint{2.516453in}{1.650196in}}%
\pgfpathlineto{\pgfqpoint{2.520994in}{1.650196in}}%
\pgfpathlineto{\pgfqpoint{2.520994in}{1.647247in}}%
\pgfpathmoveto{\pgfqpoint{2.520994in}{1.647247in}}%
\pgfpathlineto{\pgfqpoint{2.520994in}{1.647247in}}%
\pgfpathlineto{\pgfqpoint{2.520994in}{1.650196in}}%
\pgfpathlineto{\pgfqpoint{2.525535in}{1.650196in}}%
\pgfpathlineto{\pgfqpoint{2.525535in}{1.647247in}}%
\pgfpathmoveto{\pgfqpoint{2.520994in}{1.650196in}}%
\pgfpathlineto{\pgfqpoint{2.520994in}{1.650196in}}%
\pgfpathlineto{\pgfqpoint{2.520994in}{1.653145in}}%
\pgfpathlineto{\pgfqpoint{2.525535in}{1.653145in}}%
\pgfpathlineto{\pgfqpoint{2.525535in}{1.650196in}}%
\pgfpathmoveto{\pgfqpoint{2.525535in}{1.650196in}}%
\pgfpathlineto{\pgfqpoint{2.525535in}{1.650196in}}%
\pgfpathlineto{\pgfqpoint{2.525535in}{1.653145in}}%
\pgfpathlineto{\pgfqpoint{2.530076in}{1.653145in}}%
\pgfpathlineto{\pgfqpoint{2.530076in}{1.650196in}}%
\pgfpathmoveto{\pgfqpoint{2.525535in}{1.653145in}}%
\pgfpathlineto{\pgfqpoint{2.525535in}{1.653145in}}%
\pgfpathlineto{\pgfqpoint{2.525535in}{1.656095in}}%
\pgfpathlineto{\pgfqpoint{2.530076in}{1.656095in}}%
\pgfpathlineto{\pgfqpoint{2.530076in}{1.653145in}}%
\pgfpathmoveto{\pgfqpoint{2.530076in}{1.653145in}}%
\pgfpathlineto{\pgfqpoint{2.530076in}{1.653145in}}%
\pgfpathlineto{\pgfqpoint{2.530076in}{1.656095in}}%
\pgfpathlineto{\pgfqpoint{2.534616in}{1.656095in}}%
\pgfpathlineto{\pgfqpoint{2.534616in}{1.653145in}}%
\pgfpathmoveto{\pgfqpoint{2.530076in}{1.656095in}}%
\pgfpathlineto{\pgfqpoint{2.530076in}{1.656095in}}%
\pgfpathlineto{\pgfqpoint{2.530076in}{1.659044in}}%
\pgfpathlineto{\pgfqpoint{2.534616in}{1.659044in}}%
\pgfpathlineto{\pgfqpoint{2.534616in}{1.656095in}}%
\pgfpathmoveto{\pgfqpoint{2.534616in}{1.656095in}}%
\pgfpathlineto{\pgfqpoint{2.534616in}{1.656095in}}%
\pgfpathlineto{\pgfqpoint{2.534616in}{1.659044in}}%
\pgfpathlineto{\pgfqpoint{2.539157in}{1.659044in}}%
\pgfpathlineto{\pgfqpoint{2.539157in}{1.656095in}}%
\pgfpathmoveto{\pgfqpoint{2.534616in}{1.659044in}}%
\pgfpathlineto{\pgfqpoint{2.534616in}{1.659044in}}%
\pgfpathlineto{\pgfqpoint{2.534616in}{1.661993in}}%
\pgfpathlineto{\pgfqpoint{2.539157in}{1.661993in}}%
\pgfpathlineto{\pgfqpoint{2.539157in}{1.659044in}}%
\pgfpathmoveto{\pgfqpoint{2.539157in}{1.659044in}}%
\pgfpathlineto{\pgfqpoint{2.539157in}{1.659044in}}%
\pgfpathlineto{\pgfqpoint{2.539157in}{1.661993in}}%
\pgfpathlineto{\pgfqpoint{2.543698in}{1.661993in}}%
\pgfpathlineto{\pgfqpoint{2.543698in}{1.659044in}}%
\pgfpathmoveto{\pgfqpoint{2.539157in}{1.661993in}}%
\pgfpathlineto{\pgfqpoint{2.539157in}{1.661993in}}%
\pgfpathlineto{\pgfqpoint{2.539157in}{1.664942in}}%
\pgfpathlineto{\pgfqpoint{2.543698in}{1.664942in}}%
\pgfpathlineto{\pgfqpoint{2.543698in}{1.661993in}}%
\pgfpathmoveto{\pgfqpoint{2.543698in}{1.661993in}}%
\pgfpathlineto{\pgfqpoint{2.543698in}{1.661993in}}%
\pgfpathlineto{\pgfqpoint{2.543698in}{1.664942in}}%
\pgfpathlineto{\pgfqpoint{2.548239in}{1.664942in}}%
\pgfpathlineto{\pgfqpoint{2.548239in}{1.661993in}}%
\pgfpathmoveto{\pgfqpoint{2.543698in}{1.664942in}}%
\pgfpathlineto{\pgfqpoint{2.543698in}{1.664942in}}%
\pgfpathlineto{\pgfqpoint{2.543698in}{1.667892in}}%
\pgfpathlineto{\pgfqpoint{2.548239in}{1.667892in}}%
\pgfpathlineto{\pgfqpoint{2.548239in}{1.664942in}}%
\pgfpathmoveto{\pgfqpoint{2.548239in}{1.664942in}}%
\pgfpathlineto{\pgfqpoint{2.548239in}{1.664942in}}%
\pgfpathlineto{\pgfqpoint{2.548239in}{1.667892in}}%
\pgfpathlineto{\pgfqpoint{2.552780in}{1.667892in}}%
\pgfpathlineto{\pgfqpoint{2.552780in}{1.664942in}}%
\pgfpathmoveto{\pgfqpoint{2.548239in}{1.667892in}}%
\pgfpathlineto{\pgfqpoint{2.548239in}{1.667892in}}%
\pgfpathlineto{\pgfqpoint{2.548239in}{1.670841in}}%
\pgfpathlineto{\pgfqpoint{2.552780in}{1.670841in}}%
\pgfpathlineto{\pgfqpoint{2.552780in}{1.667892in}}%
\pgfpathmoveto{\pgfqpoint{2.552780in}{1.667892in}}%
\pgfpathlineto{\pgfqpoint{2.552780in}{1.667892in}}%
\pgfpathlineto{\pgfqpoint{2.552780in}{1.670841in}}%
\pgfpathlineto{\pgfqpoint{2.557321in}{1.670841in}}%
\pgfpathlineto{\pgfqpoint{2.557321in}{1.667892in}}%
\pgfpathmoveto{\pgfqpoint{2.552780in}{1.670841in}}%
\pgfpathlineto{\pgfqpoint{2.552780in}{1.670841in}}%
\pgfpathlineto{\pgfqpoint{2.552780in}{1.673790in}}%
\pgfpathlineto{\pgfqpoint{2.557321in}{1.673790in}}%
\pgfpathlineto{\pgfqpoint{2.557321in}{1.670841in}}%
\pgfpathmoveto{\pgfqpoint{2.557321in}{1.670841in}}%
\pgfpathlineto{\pgfqpoint{2.557321in}{1.670841in}}%
\pgfpathlineto{\pgfqpoint{2.557321in}{1.673790in}}%
\pgfpathlineto{\pgfqpoint{2.561862in}{1.673790in}}%
\pgfpathlineto{\pgfqpoint{2.561862in}{1.670841in}}%
\pgfpathmoveto{\pgfqpoint{2.557321in}{1.673790in}}%
\pgfpathlineto{\pgfqpoint{2.557321in}{1.673790in}}%
\pgfpathlineto{\pgfqpoint{2.557321in}{1.676739in}}%
\pgfpathlineto{\pgfqpoint{2.561862in}{1.676739in}}%
\pgfpathlineto{\pgfqpoint{2.561862in}{1.673790in}}%
\pgfpathmoveto{\pgfqpoint{2.561862in}{1.673790in}}%
\pgfpathlineto{\pgfqpoint{2.561862in}{1.673790in}}%
\pgfpathlineto{\pgfqpoint{2.561862in}{1.676739in}}%
\pgfpathlineto{\pgfqpoint{2.566403in}{1.676739in}}%
\pgfpathlineto{\pgfqpoint{2.566403in}{1.673790in}}%
\pgfpathmoveto{\pgfqpoint{2.561862in}{1.676739in}}%
\pgfpathlineto{\pgfqpoint{2.561862in}{1.676739in}}%
\pgfpathlineto{\pgfqpoint{2.561862in}{1.679689in}}%
\pgfpathlineto{\pgfqpoint{2.566403in}{1.679689in}}%
\pgfpathlineto{\pgfqpoint{2.566403in}{1.676739in}}%
\pgfpathmoveto{\pgfqpoint{2.566403in}{1.676739in}}%
\pgfpathlineto{\pgfqpoint{2.566403in}{1.676739in}}%
\pgfpathlineto{\pgfqpoint{2.566403in}{1.679689in}}%
\pgfpathlineto{\pgfqpoint{2.570944in}{1.679689in}}%
\pgfpathlineto{\pgfqpoint{2.570944in}{1.676739in}}%
\pgfpathmoveto{\pgfqpoint{2.566403in}{1.679689in}}%
\pgfpathlineto{\pgfqpoint{2.566403in}{1.679689in}}%
\pgfpathlineto{\pgfqpoint{2.566403in}{1.682638in}}%
\pgfpathlineto{\pgfqpoint{2.570944in}{1.682638in}}%
\pgfpathlineto{\pgfqpoint{2.570944in}{1.679689in}}%
\pgfpathmoveto{\pgfqpoint{2.570944in}{1.679689in}}%
\pgfpathlineto{\pgfqpoint{2.570944in}{1.679689in}}%
\pgfpathlineto{\pgfqpoint{2.570944in}{1.682638in}}%
\pgfpathlineto{\pgfqpoint{2.575485in}{1.682638in}}%
\pgfpathlineto{\pgfqpoint{2.575485in}{1.679689in}}%
\pgfpathmoveto{\pgfqpoint{2.570944in}{1.682638in}}%
\pgfpathlineto{\pgfqpoint{2.570944in}{1.682638in}}%
\pgfpathlineto{\pgfqpoint{2.570944in}{1.685587in}}%
\pgfpathlineto{\pgfqpoint{2.575485in}{1.685587in}}%
\pgfpathlineto{\pgfqpoint{2.575485in}{1.682638in}}%
\pgfpathmoveto{\pgfqpoint{2.575485in}{1.682638in}}%
\pgfpathlineto{\pgfqpoint{2.575485in}{1.682638in}}%
\pgfpathlineto{\pgfqpoint{2.575485in}{1.685587in}}%
\pgfpathlineto{\pgfqpoint{2.580026in}{1.685587in}}%
\pgfpathlineto{\pgfqpoint{2.580026in}{1.682638in}}%
\pgfpathmoveto{\pgfqpoint{2.575485in}{1.685587in}}%
\pgfpathlineto{\pgfqpoint{2.575485in}{1.685587in}}%
\pgfpathlineto{\pgfqpoint{2.575485in}{1.688536in}}%
\pgfpathlineto{\pgfqpoint{2.580026in}{1.688536in}}%
\pgfpathlineto{\pgfqpoint{2.580026in}{1.685587in}}%
\pgfpathmoveto{\pgfqpoint{2.580026in}{1.685587in}}%
\pgfpathlineto{\pgfqpoint{2.580026in}{1.685587in}}%
\pgfpathlineto{\pgfqpoint{2.580026in}{1.688536in}}%
\pgfpathlineto{\pgfqpoint{2.584567in}{1.688536in}}%
\pgfpathlineto{\pgfqpoint{2.584567in}{1.685587in}}%
\pgfpathmoveto{\pgfqpoint{2.580026in}{1.688536in}}%
\pgfpathlineto{\pgfqpoint{2.580026in}{1.688536in}}%
\pgfpathlineto{\pgfqpoint{2.580026in}{1.691485in}}%
\pgfpathlineto{\pgfqpoint{2.584567in}{1.691485in}}%
\pgfpathlineto{\pgfqpoint{2.584567in}{1.688536in}}%
\pgfpathmoveto{\pgfqpoint{2.584567in}{1.688536in}}%
\pgfpathlineto{\pgfqpoint{2.584567in}{1.688536in}}%
\pgfpathlineto{\pgfqpoint{2.584567in}{1.691485in}}%
\pgfpathlineto{\pgfqpoint{2.589108in}{1.691485in}}%
\pgfpathlineto{\pgfqpoint{2.589108in}{1.688536in}}%
\pgfpathmoveto{\pgfqpoint{2.584567in}{1.691485in}}%
\pgfpathlineto{\pgfqpoint{2.584567in}{1.691485in}}%
\pgfpathlineto{\pgfqpoint{2.584567in}{1.694435in}}%
\pgfpathlineto{\pgfqpoint{2.589108in}{1.694435in}}%
\pgfpathlineto{\pgfqpoint{2.589108in}{1.691485in}}%
\pgfpathmoveto{\pgfqpoint{2.589108in}{1.691485in}}%
\pgfpathlineto{\pgfqpoint{2.589108in}{1.691485in}}%
\pgfpathlineto{\pgfqpoint{2.589108in}{1.694435in}}%
\pgfpathlineto{\pgfqpoint{2.593649in}{1.694435in}}%
\pgfpathlineto{\pgfqpoint{2.593649in}{1.691485in}}%
\pgfpathmoveto{\pgfqpoint{2.589108in}{1.694435in}}%
\pgfpathlineto{\pgfqpoint{2.589108in}{1.694435in}}%
\pgfpathlineto{\pgfqpoint{2.589108in}{1.697384in}}%
\pgfpathlineto{\pgfqpoint{2.593649in}{1.697384in}}%
\pgfpathlineto{\pgfqpoint{2.593649in}{1.694435in}}%
\pgfpathmoveto{\pgfqpoint{2.593649in}{1.694435in}}%
\pgfpathlineto{\pgfqpoint{2.593649in}{1.694435in}}%
\pgfpathlineto{\pgfqpoint{2.593649in}{1.697384in}}%
\pgfpathlineto{\pgfqpoint{2.598190in}{1.697384in}}%
\pgfpathlineto{\pgfqpoint{2.598190in}{1.694435in}}%
\pgfpathmoveto{\pgfqpoint{2.593649in}{1.697384in}}%
\pgfpathlineto{\pgfqpoint{2.593649in}{1.697384in}}%
\pgfpathlineto{\pgfqpoint{2.593649in}{1.700333in}}%
\pgfpathlineto{\pgfqpoint{2.598190in}{1.700333in}}%
\pgfpathlineto{\pgfqpoint{2.598190in}{1.697384in}}%
\pgfpathmoveto{\pgfqpoint{2.598190in}{1.697384in}}%
\pgfpathlineto{\pgfqpoint{2.598190in}{1.697384in}}%
\pgfpathlineto{\pgfqpoint{2.598190in}{1.700333in}}%
\pgfpathlineto{\pgfqpoint{2.602731in}{1.700333in}}%
\pgfpathlineto{\pgfqpoint{2.602731in}{1.697384in}}%
\pgfpathmoveto{\pgfqpoint{2.598190in}{1.700333in}}%
\pgfpathlineto{\pgfqpoint{2.598190in}{1.700333in}}%
\pgfpathlineto{\pgfqpoint{2.598190in}{1.703282in}}%
\pgfpathlineto{\pgfqpoint{2.602731in}{1.703282in}}%
\pgfpathlineto{\pgfqpoint{2.602731in}{1.700333in}}%
\pgfpathmoveto{\pgfqpoint{2.602731in}{1.700333in}}%
\pgfpathlineto{\pgfqpoint{2.602731in}{1.700333in}}%
\pgfpathlineto{\pgfqpoint{2.602731in}{1.703282in}}%
\pgfpathlineto{\pgfqpoint{2.607272in}{1.703282in}}%
\pgfpathlineto{\pgfqpoint{2.607272in}{1.700333in}}%
\pgfpathmoveto{\pgfqpoint{2.602731in}{1.703282in}}%
\pgfpathlineto{\pgfqpoint{2.602731in}{1.703282in}}%
\pgfpathlineto{\pgfqpoint{2.602731in}{1.706232in}}%
\pgfpathlineto{\pgfqpoint{2.607272in}{1.706232in}}%
\pgfpathlineto{\pgfqpoint{2.607272in}{1.703282in}}%
\pgfpathmoveto{\pgfqpoint{2.607272in}{1.703282in}}%
\pgfpathlineto{\pgfqpoint{2.607272in}{1.703282in}}%
\pgfpathlineto{\pgfqpoint{2.607272in}{1.706232in}}%
\pgfpathlineto{\pgfqpoint{2.611813in}{1.706232in}}%
\pgfpathlineto{\pgfqpoint{2.611813in}{1.703282in}}%
\pgfpathmoveto{\pgfqpoint{2.607272in}{1.706232in}}%
\pgfpathlineto{\pgfqpoint{2.607272in}{1.706232in}}%
\pgfpathlineto{\pgfqpoint{2.607272in}{1.709181in}}%
\pgfpathlineto{\pgfqpoint{2.611813in}{1.709181in}}%
\pgfpathlineto{\pgfqpoint{2.611813in}{1.706232in}}%
\pgfpathmoveto{\pgfqpoint{2.611813in}{1.706232in}}%
\pgfpathlineto{\pgfqpoint{2.611813in}{1.706232in}}%
\pgfpathlineto{\pgfqpoint{2.611813in}{1.709181in}}%
\pgfpathlineto{\pgfqpoint{2.616354in}{1.709181in}}%
\pgfpathlineto{\pgfqpoint{2.616354in}{1.706232in}}%
\pgfpathmoveto{\pgfqpoint{2.611813in}{1.709181in}}%
\pgfpathlineto{\pgfqpoint{2.611813in}{1.709181in}}%
\pgfpathlineto{\pgfqpoint{2.611813in}{1.712130in}}%
\pgfpathlineto{\pgfqpoint{2.616354in}{1.712130in}}%
\pgfpathlineto{\pgfqpoint{2.616354in}{1.709181in}}%
\pgfpathmoveto{\pgfqpoint{2.616354in}{1.709181in}}%
\pgfpathlineto{\pgfqpoint{2.616354in}{1.709181in}}%
\pgfpathlineto{\pgfqpoint{2.616354in}{1.712130in}}%
\pgfpathlineto{\pgfqpoint{2.620895in}{1.712130in}}%
\pgfpathlineto{\pgfqpoint{2.620895in}{1.709181in}}%
\pgfpathmoveto{\pgfqpoint{2.616354in}{1.712130in}}%
\pgfpathlineto{\pgfqpoint{2.616354in}{1.712130in}}%
\pgfpathlineto{\pgfqpoint{2.616354in}{1.715079in}}%
\pgfpathlineto{\pgfqpoint{2.620895in}{1.715079in}}%
\pgfpathlineto{\pgfqpoint{2.620895in}{1.712130in}}%
\pgfpathmoveto{\pgfqpoint{2.620895in}{1.712130in}}%
\pgfpathlineto{\pgfqpoint{2.620895in}{1.712130in}}%
\pgfpathlineto{\pgfqpoint{2.620895in}{1.715079in}}%
\pgfpathlineto{\pgfqpoint{2.625436in}{1.715079in}}%
\pgfpathlineto{\pgfqpoint{2.625436in}{1.712130in}}%
\pgfpathmoveto{\pgfqpoint{2.620895in}{1.715079in}}%
\pgfpathlineto{\pgfqpoint{2.620895in}{1.715079in}}%
\pgfpathlineto{\pgfqpoint{2.620895in}{1.718029in}}%
\pgfpathlineto{\pgfqpoint{2.625436in}{1.718029in}}%
\pgfpathlineto{\pgfqpoint{2.625436in}{1.715079in}}%
\pgfpathmoveto{\pgfqpoint{2.625436in}{1.715079in}}%
\pgfpathlineto{\pgfqpoint{2.625436in}{1.715079in}}%
\pgfpathlineto{\pgfqpoint{2.625436in}{1.718029in}}%
\pgfpathlineto{\pgfqpoint{2.629977in}{1.718029in}}%
\pgfpathlineto{\pgfqpoint{2.629977in}{1.715079in}}%
\pgfpathmoveto{\pgfqpoint{2.625436in}{1.718029in}}%
\pgfpathlineto{\pgfqpoint{2.625436in}{1.718029in}}%
\pgfpathlineto{\pgfqpoint{2.625436in}{1.720978in}}%
\pgfpathlineto{\pgfqpoint{2.629977in}{1.720978in}}%
\pgfpathlineto{\pgfqpoint{2.629977in}{1.718029in}}%
\pgfpathmoveto{\pgfqpoint{2.629977in}{1.718029in}}%
\pgfpathlineto{\pgfqpoint{2.629977in}{1.718029in}}%
\pgfpathlineto{\pgfqpoint{2.629977in}{1.720978in}}%
\pgfpathlineto{\pgfqpoint{2.634518in}{1.720978in}}%
\pgfpathlineto{\pgfqpoint{2.634518in}{1.718029in}}%
\pgfpathmoveto{\pgfqpoint{2.629977in}{1.720978in}}%
\pgfpathlineto{\pgfqpoint{2.629977in}{1.720978in}}%
\pgfpathlineto{\pgfqpoint{2.629977in}{1.723927in}}%
\pgfpathlineto{\pgfqpoint{2.634518in}{1.723927in}}%
\pgfpathlineto{\pgfqpoint{2.634518in}{1.720978in}}%
\pgfpathmoveto{\pgfqpoint{2.634518in}{1.720978in}}%
\pgfpathlineto{\pgfqpoint{2.634518in}{1.720978in}}%
\pgfpathlineto{\pgfqpoint{2.634518in}{1.723927in}}%
\pgfpathlineto{\pgfqpoint{2.639059in}{1.723927in}}%
\pgfpathlineto{\pgfqpoint{2.639059in}{1.720978in}}%
\pgfpathmoveto{\pgfqpoint{2.634518in}{1.723927in}}%
\pgfpathlineto{\pgfqpoint{2.634518in}{1.723927in}}%
\pgfpathlineto{\pgfqpoint{2.634518in}{1.726876in}}%
\pgfpathlineto{\pgfqpoint{2.639059in}{1.726876in}}%
\pgfpathlineto{\pgfqpoint{2.639059in}{1.723927in}}%
\pgfpathmoveto{\pgfqpoint{2.639059in}{1.723927in}}%
\pgfpathlineto{\pgfqpoint{2.639059in}{1.723927in}}%
\pgfpathlineto{\pgfqpoint{2.639059in}{1.726876in}}%
\pgfpathlineto{\pgfqpoint{2.643600in}{1.726876in}}%
\pgfpathlineto{\pgfqpoint{2.643600in}{1.723927in}}%
\pgfpathmoveto{\pgfqpoint{2.639059in}{1.726876in}}%
\pgfpathlineto{\pgfqpoint{2.639059in}{1.726876in}}%
\pgfpathlineto{\pgfqpoint{2.639059in}{1.729826in}}%
\pgfpathlineto{\pgfqpoint{2.643600in}{1.729826in}}%
\pgfpathlineto{\pgfqpoint{2.643600in}{1.726876in}}%
\pgfpathmoveto{\pgfqpoint{2.643600in}{1.726876in}}%
\pgfpathlineto{\pgfqpoint{2.643600in}{1.726876in}}%
\pgfpathlineto{\pgfqpoint{2.643600in}{1.729826in}}%
\pgfpathlineto{\pgfqpoint{2.648141in}{1.729826in}}%
\pgfpathlineto{\pgfqpoint{2.648141in}{1.726876in}}%
\pgfpathmoveto{\pgfqpoint{2.643600in}{1.729826in}}%
\pgfpathlineto{\pgfqpoint{2.643600in}{1.729826in}}%
\pgfpathlineto{\pgfqpoint{2.643600in}{1.732775in}}%
\pgfpathlineto{\pgfqpoint{2.648141in}{1.732775in}}%
\pgfpathlineto{\pgfqpoint{2.648141in}{1.729826in}}%
\pgfpathmoveto{\pgfqpoint{2.648141in}{1.729826in}}%
\pgfpathlineto{\pgfqpoint{2.648141in}{1.729826in}}%
\pgfpathlineto{\pgfqpoint{2.648141in}{1.732775in}}%
\pgfpathlineto{\pgfqpoint{2.652682in}{1.732775in}}%
\pgfpathlineto{\pgfqpoint{2.652682in}{1.729826in}}%
\pgfpathmoveto{\pgfqpoint{2.648141in}{1.732775in}}%
\pgfpathlineto{\pgfqpoint{2.648141in}{1.732775in}}%
\pgfpathlineto{\pgfqpoint{2.648141in}{1.735724in}}%
\pgfpathlineto{\pgfqpoint{2.652682in}{1.735724in}}%
\pgfpathlineto{\pgfqpoint{2.652682in}{1.732775in}}%
\pgfpathmoveto{\pgfqpoint{2.652682in}{1.732775in}}%
\pgfpathlineto{\pgfqpoint{2.652682in}{1.732775in}}%
\pgfpathlineto{\pgfqpoint{2.652682in}{1.735724in}}%
\pgfpathlineto{\pgfqpoint{2.657224in}{1.735724in}}%
\pgfpathlineto{\pgfqpoint{2.657224in}{1.732775in}}%
\pgfpathmoveto{\pgfqpoint{2.652682in}{1.735724in}}%
\pgfpathlineto{\pgfqpoint{2.652682in}{1.735724in}}%
\pgfpathlineto{\pgfqpoint{2.652682in}{1.738673in}}%
\pgfpathlineto{\pgfqpoint{2.657224in}{1.738673in}}%
\pgfpathlineto{\pgfqpoint{2.657224in}{1.735724in}}%
\pgfpathmoveto{\pgfqpoint{2.657224in}{1.735724in}}%
\pgfpathlineto{\pgfqpoint{2.657224in}{1.735724in}}%
\pgfpathlineto{\pgfqpoint{2.657224in}{1.738673in}}%
\pgfpathlineto{\pgfqpoint{2.661765in}{1.738673in}}%
\pgfpathlineto{\pgfqpoint{2.661765in}{1.735724in}}%
\pgfpathmoveto{\pgfqpoint{2.657224in}{1.738673in}}%
\pgfpathlineto{\pgfqpoint{2.657224in}{1.738673in}}%
\pgfpathlineto{\pgfqpoint{2.657224in}{1.741622in}}%
\pgfpathlineto{\pgfqpoint{2.661765in}{1.741622in}}%
\pgfpathlineto{\pgfqpoint{2.661765in}{1.738673in}}%
\pgfpathmoveto{\pgfqpoint{2.661765in}{1.738673in}}%
\pgfpathlineto{\pgfqpoint{2.661765in}{1.738673in}}%
\pgfpathlineto{\pgfqpoint{2.661765in}{1.741622in}}%
\pgfpathlineto{\pgfqpoint{2.666306in}{1.741622in}}%
\pgfpathlineto{\pgfqpoint{2.666306in}{1.738673in}}%
\pgfpathmoveto{\pgfqpoint{2.661765in}{1.741622in}}%
\pgfpathlineto{\pgfqpoint{2.661765in}{1.741622in}}%
\pgfpathlineto{\pgfqpoint{2.661765in}{1.744572in}}%
\pgfpathlineto{\pgfqpoint{2.666306in}{1.744572in}}%
\pgfpathlineto{\pgfqpoint{2.666306in}{1.741622in}}%
\pgfpathmoveto{\pgfqpoint{2.666306in}{1.741622in}}%
\pgfpathlineto{\pgfqpoint{2.666306in}{1.741622in}}%
\pgfpathlineto{\pgfqpoint{2.666306in}{1.744572in}}%
\pgfpathlineto{\pgfqpoint{2.670847in}{1.744572in}}%
\pgfpathlineto{\pgfqpoint{2.670847in}{1.741622in}}%
\pgfpathmoveto{\pgfqpoint{2.666306in}{1.744572in}}%
\pgfpathlineto{\pgfqpoint{2.666306in}{1.744572in}}%
\pgfpathlineto{\pgfqpoint{2.666306in}{1.747521in}}%
\pgfpathlineto{\pgfqpoint{2.670847in}{1.747521in}}%
\pgfpathlineto{\pgfqpoint{2.670847in}{1.744572in}}%
\pgfpathmoveto{\pgfqpoint{2.670847in}{1.744572in}}%
\pgfpathlineto{\pgfqpoint{2.670847in}{1.744572in}}%
\pgfpathlineto{\pgfqpoint{2.670847in}{1.747521in}}%
\pgfpathlineto{\pgfqpoint{2.675388in}{1.747521in}}%
\pgfpathlineto{\pgfqpoint{2.675388in}{1.744572in}}%
\pgfpathmoveto{\pgfqpoint{2.670847in}{1.747521in}}%
\pgfpathlineto{\pgfqpoint{2.670847in}{1.747521in}}%
\pgfpathlineto{\pgfqpoint{2.670847in}{1.750470in}}%
\pgfpathlineto{\pgfqpoint{2.675388in}{1.750470in}}%
\pgfpathlineto{\pgfqpoint{2.675388in}{1.747521in}}%
\pgfpathmoveto{\pgfqpoint{2.675388in}{1.747521in}}%
\pgfpathlineto{\pgfqpoint{2.675388in}{1.747521in}}%
\pgfpathlineto{\pgfqpoint{2.675388in}{1.750470in}}%
\pgfpathlineto{\pgfqpoint{2.679929in}{1.750470in}}%
\pgfpathlineto{\pgfqpoint{2.679929in}{1.747521in}}%
\pgfpathmoveto{\pgfqpoint{2.675388in}{1.750470in}}%
\pgfpathlineto{\pgfqpoint{2.675388in}{1.750470in}}%
\pgfpathlineto{\pgfqpoint{2.675388in}{1.753419in}}%
\pgfpathlineto{\pgfqpoint{2.679929in}{1.753419in}}%
\pgfpathlineto{\pgfqpoint{2.679929in}{1.750470in}}%
\pgfpathmoveto{\pgfqpoint{2.679929in}{1.750470in}}%
\pgfpathlineto{\pgfqpoint{2.679929in}{1.750470in}}%
\pgfpathlineto{\pgfqpoint{2.679929in}{1.753419in}}%
\pgfpathlineto{\pgfqpoint{2.684470in}{1.753419in}}%
\pgfpathlineto{\pgfqpoint{2.684470in}{1.750470in}}%
\pgfpathmoveto{\pgfqpoint{2.679929in}{1.753419in}}%
\pgfpathlineto{\pgfqpoint{2.679929in}{1.753419in}}%
\pgfpathlineto{\pgfqpoint{2.679929in}{1.756368in}}%
\pgfpathlineto{\pgfqpoint{2.684470in}{1.756368in}}%
\pgfpathlineto{\pgfqpoint{2.684470in}{1.753419in}}%
\pgfpathmoveto{\pgfqpoint{2.684470in}{1.753419in}}%
\pgfpathlineto{\pgfqpoint{2.684470in}{1.753419in}}%
\pgfpathlineto{\pgfqpoint{2.684470in}{1.756368in}}%
\pgfpathlineto{\pgfqpoint{2.689011in}{1.756368in}}%
\pgfpathlineto{\pgfqpoint{2.689011in}{1.753419in}}%
\pgfpathmoveto{\pgfqpoint{2.684470in}{1.756368in}}%
\pgfpathlineto{\pgfqpoint{2.684470in}{1.756368in}}%
\pgfpathlineto{\pgfqpoint{2.684470in}{1.759318in}}%
\pgfpathlineto{\pgfqpoint{2.689011in}{1.759318in}}%
\pgfpathlineto{\pgfqpoint{2.689011in}{1.756368in}}%
\pgfpathmoveto{\pgfqpoint{2.689011in}{1.756368in}}%
\pgfpathlineto{\pgfqpoint{2.689011in}{1.756368in}}%
\pgfpathlineto{\pgfqpoint{2.689011in}{1.759318in}}%
\pgfpathlineto{\pgfqpoint{2.693552in}{1.759318in}}%
\pgfpathlineto{\pgfqpoint{2.693552in}{1.756368in}}%
\pgfpathmoveto{\pgfqpoint{2.689011in}{1.759318in}}%
\pgfpathlineto{\pgfqpoint{2.689011in}{1.759318in}}%
\pgfpathlineto{\pgfqpoint{2.689011in}{1.762267in}}%
\pgfpathlineto{\pgfqpoint{2.693552in}{1.762267in}}%
\pgfpathlineto{\pgfqpoint{2.693552in}{1.759318in}}%
\pgfpathmoveto{\pgfqpoint{2.693552in}{1.759318in}}%
\pgfpathlineto{\pgfqpoint{2.693552in}{1.759318in}}%
\pgfpathlineto{\pgfqpoint{2.693552in}{1.762267in}}%
\pgfpathlineto{\pgfqpoint{2.698093in}{1.762267in}}%
\pgfpathlineto{\pgfqpoint{2.698093in}{1.759318in}}%
\pgfpathmoveto{\pgfqpoint{2.693552in}{1.762267in}}%
\pgfpathlineto{\pgfqpoint{2.693552in}{1.762267in}}%
\pgfpathlineto{\pgfqpoint{2.693552in}{1.765216in}}%
\pgfpathlineto{\pgfqpoint{2.698093in}{1.765216in}}%
\pgfpathlineto{\pgfqpoint{2.698093in}{1.762267in}}%
\pgfpathmoveto{\pgfqpoint{2.698093in}{1.762267in}}%
\pgfpathlineto{\pgfqpoint{2.698093in}{1.762267in}}%
\pgfpathlineto{\pgfqpoint{2.698093in}{1.765216in}}%
\pgfpathlineto{\pgfqpoint{2.702634in}{1.765216in}}%
\pgfpathlineto{\pgfqpoint{2.702634in}{1.762267in}}%
\pgfpathmoveto{\pgfqpoint{2.698093in}{1.765216in}}%
\pgfpathlineto{\pgfqpoint{2.698093in}{1.765216in}}%
\pgfpathlineto{\pgfqpoint{2.698093in}{1.768165in}}%
\pgfpathlineto{\pgfqpoint{2.702634in}{1.768165in}}%
\pgfpathlineto{\pgfqpoint{2.702634in}{1.765216in}}%
\pgfpathmoveto{\pgfqpoint{2.702634in}{1.765216in}}%
\pgfpathlineto{\pgfqpoint{2.702634in}{1.765216in}}%
\pgfpathlineto{\pgfqpoint{2.702634in}{1.768165in}}%
\pgfpathlineto{\pgfqpoint{2.707175in}{1.768165in}}%
\pgfpathlineto{\pgfqpoint{2.707175in}{1.765216in}}%
\pgfpathmoveto{\pgfqpoint{2.702634in}{1.768165in}}%
\pgfpathlineto{\pgfqpoint{2.702634in}{1.768165in}}%
\pgfpathlineto{\pgfqpoint{2.702634in}{1.771114in}}%
\pgfpathlineto{\pgfqpoint{2.707175in}{1.771114in}}%
\pgfpathlineto{\pgfqpoint{2.707175in}{1.768165in}}%
\pgfpathmoveto{\pgfqpoint{2.707175in}{1.768165in}}%
\pgfpathlineto{\pgfqpoint{2.707175in}{1.768165in}}%
\pgfpathlineto{\pgfqpoint{2.707175in}{1.771114in}}%
\pgfpathlineto{\pgfqpoint{2.711716in}{1.771114in}}%
\pgfpathlineto{\pgfqpoint{2.711716in}{1.768165in}}%
\pgfpathmoveto{\pgfqpoint{2.707175in}{1.771114in}}%
\pgfpathlineto{\pgfqpoint{2.707175in}{1.771114in}}%
\pgfpathlineto{\pgfqpoint{2.707175in}{1.774063in}}%
\pgfpathlineto{\pgfqpoint{2.711716in}{1.774063in}}%
\pgfpathlineto{\pgfqpoint{2.711716in}{1.771114in}}%
\pgfpathmoveto{\pgfqpoint{2.711716in}{1.771114in}}%
\pgfpathlineto{\pgfqpoint{2.711716in}{1.771114in}}%
\pgfpathlineto{\pgfqpoint{2.711716in}{1.774063in}}%
\pgfpathlineto{\pgfqpoint{2.716258in}{1.774063in}}%
\pgfpathlineto{\pgfqpoint{2.716258in}{1.771114in}}%
\pgfpathmoveto{\pgfqpoint{2.711716in}{1.774063in}}%
\pgfpathlineto{\pgfqpoint{2.711716in}{1.774063in}}%
\pgfpathlineto{\pgfqpoint{2.711716in}{1.777013in}}%
\pgfpathlineto{\pgfqpoint{2.716258in}{1.777013in}}%
\pgfpathlineto{\pgfqpoint{2.716258in}{1.774063in}}%
\pgfpathmoveto{\pgfqpoint{2.716258in}{1.774063in}}%
\pgfpathlineto{\pgfqpoint{2.716258in}{1.774063in}}%
\pgfpathlineto{\pgfqpoint{2.716258in}{1.777013in}}%
\pgfpathlineto{\pgfqpoint{2.720799in}{1.777013in}}%
\pgfpathlineto{\pgfqpoint{2.720799in}{1.774063in}}%
\pgfpathmoveto{\pgfqpoint{2.716258in}{1.777013in}}%
\pgfpathlineto{\pgfqpoint{2.716258in}{1.777013in}}%
\pgfpathlineto{\pgfqpoint{2.716258in}{1.779962in}}%
\pgfpathlineto{\pgfqpoint{2.720799in}{1.779962in}}%
\pgfpathlineto{\pgfqpoint{2.720799in}{1.777013in}}%
\pgfpathmoveto{\pgfqpoint{2.720799in}{1.777013in}}%
\pgfpathlineto{\pgfqpoint{2.720799in}{1.777013in}}%
\pgfpathlineto{\pgfqpoint{2.720799in}{1.779962in}}%
\pgfpathlineto{\pgfqpoint{2.725340in}{1.779962in}}%
\pgfpathlineto{\pgfqpoint{2.725340in}{1.777013in}}%
\pgfpathmoveto{\pgfqpoint{2.720799in}{1.779962in}}%
\pgfpathlineto{\pgfqpoint{2.720799in}{1.779962in}}%
\pgfpathlineto{\pgfqpoint{2.720799in}{1.782911in}}%
\pgfpathlineto{\pgfqpoint{2.725340in}{1.782911in}}%
\pgfpathlineto{\pgfqpoint{2.725340in}{1.779962in}}%
\pgfpathmoveto{\pgfqpoint{2.725340in}{1.779962in}}%
\pgfpathlineto{\pgfqpoint{2.725340in}{1.779962in}}%
\pgfpathlineto{\pgfqpoint{2.725340in}{1.782911in}}%
\pgfpathlineto{\pgfqpoint{2.729881in}{1.782911in}}%
\pgfpathlineto{\pgfqpoint{2.729881in}{1.779962in}}%
\pgfpathmoveto{\pgfqpoint{2.725340in}{1.782911in}}%
\pgfpathlineto{\pgfqpoint{2.725340in}{1.782911in}}%
\pgfpathlineto{\pgfqpoint{2.725340in}{1.785860in}}%
\pgfpathlineto{\pgfqpoint{2.729881in}{1.785860in}}%
\pgfpathlineto{\pgfqpoint{2.729881in}{1.782911in}}%
\pgfpathmoveto{\pgfqpoint{2.729881in}{1.782911in}}%
\pgfpathlineto{\pgfqpoint{2.729881in}{1.782911in}}%
\pgfpathlineto{\pgfqpoint{2.729881in}{1.785860in}}%
\pgfpathlineto{\pgfqpoint{2.734422in}{1.785860in}}%
\pgfpathlineto{\pgfqpoint{2.734422in}{1.782911in}}%
\pgfpathmoveto{\pgfqpoint{2.729881in}{1.785860in}}%
\pgfpathlineto{\pgfqpoint{2.729881in}{1.785860in}}%
\pgfpathlineto{\pgfqpoint{2.729881in}{1.788809in}}%
\pgfpathlineto{\pgfqpoint{2.734422in}{1.788809in}}%
\pgfpathlineto{\pgfqpoint{2.734422in}{1.785860in}}%
\pgfpathmoveto{\pgfqpoint{2.734422in}{1.785860in}}%
\pgfpathlineto{\pgfqpoint{2.734422in}{1.785860in}}%
\pgfpathlineto{\pgfqpoint{2.734422in}{1.788809in}}%
\pgfpathlineto{\pgfqpoint{2.738963in}{1.788809in}}%
\pgfpathlineto{\pgfqpoint{2.738963in}{1.785860in}}%
\pgfpathmoveto{\pgfqpoint{2.734422in}{1.788809in}}%
\pgfpathlineto{\pgfqpoint{2.734422in}{1.788809in}}%
\pgfpathlineto{\pgfqpoint{2.734422in}{1.791759in}}%
\pgfpathlineto{\pgfqpoint{2.738963in}{1.791759in}}%
\pgfpathlineto{\pgfqpoint{2.738963in}{1.788809in}}%
\pgfpathmoveto{\pgfqpoint{2.738963in}{1.788809in}}%
\pgfpathlineto{\pgfqpoint{2.738963in}{1.788809in}}%
\pgfpathlineto{\pgfqpoint{2.738963in}{1.791759in}}%
\pgfpathlineto{\pgfqpoint{2.743504in}{1.791759in}}%
\pgfpathlineto{\pgfqpoint{2.743504in}{1.788809in}}%
\pgfpathmoveto{\pgfqpoint{2.738963in}{1.791759in}}%
\pgfpathlineto{\pgfqpoint{2.738963in}{1.791759in}}%
\pgfpathlineto{\pgfqpoint{2.738963in}{1.794708in}}%
\pgfpathlineto{\pgfqpoint{2.743504in}{1.794708in}}%
\pgfpathlineto{\pgfqpoint{2.743504in}{1.791759in}}%
\pgfpathmoveto{\pgfqpoint{2.743504in}{1.791759in}}%
\pgfpathlineto{\pgfqpoint{2.743504in}{1.791759in}}%
\pgfpathlineto{\pgfqpoint{2.743504in}{1.794708in}}%
\pgfpathlineto{\pgfqpoint{2.748045in}{1.794708in}}%
\pgfpathlineto{\pgfqpoint{2.748045in}{1.791759in}}%
\pgfpathmoveto{\pgfqpoint{2.743504in}{1.794708in}}%
\pgfpathlineto{\pgfqpoint{2.743504in}{1.794708in}}%
\pgfpathlineto{\pgfqpoint{2.743504in}{1.797657in}}%
\pgfpathlineto{\pgfqpoint{2.748045in}{1.797657in}}%
\pgfpathlineto{\pgfqpoint{2.748045in}{1.794708in}}%
\pgfpathmoveto{\pgfqpoint{2.748045in}{1.794708in}}%
\pgfpathlineto{\pgfqpoint{2.748045in}{1.794708in}}%
\pgfpathlineto{\pgfqpoint{2.748045in}{1.797657in}}%
\pgfpathlineto{\pgfqpoint{2.752586in}{1.797657in}}%
\pgfpathlineto{\pgfqpoint{2.752586in}{1.794708in}}%
\pgfpathmoveto{\pgfqpoint{2.748045in}{1.797657in}}%
\pgfpathlineto{\pgfqpoint{2.748045in}{1.797657in}}%
\pgfpathlineto{\pgfqpoint{2.748045in}{1.800606in}}%
\pgfpathlineto{\pgfqpoint{2.752586in}{1.800606in}}%
\pgfpathlineto{\pgfqpoint{2.752586in}{1.797657in}}%
\pgfpathmoveto{\pgfqpoint{2.752586in}{1.797657in}}%
\pgfpathlineto{\pgfqpoint{2.752586in}{1.797657in}}%
\pgfpathlineto{\pgfqpoint{2.752586in}{1.800606in}}%
\pgfpathlineto{\pgfqpoint{2.757127in}{1.800606in}}%
\pgfpathlineto{\pgfqpoint{2.757127in}{1.797657in}}%
\pgfpathmoveto{\pgfqpoint{2.752586in}{1.800606in}}%
\pgfpathlineto{\pgfqpoint{2.752586in}{1.800606in}}%
\pgfpathlineto{\pgfqpoint{2.752586in}{1.803555in}}%
\pgfpathlineto{\pgfqpoint{2.757127in}{1.803555in}}%
\pgfpathlineto{\pgfqpoint{2.757127in}{1.800606in}}%
\pgfpathmoveto{\pgfqpoint{2.757127in}{1.800606in}}%
\pgfpathlineto{\pgfqpoint{2.757127in}{1.800606in}}%
\pgfpathlineto{\pgfqpoint{2.757127in}{1.803555in}}%
\pgfpathlineto{\pgfqpoint{2.761668in}{1.803555in}}%
\pgfpathlineto{\pgfqpoint{2.761668in}{1.800606in}}%
\pgfpathmoveto{\pgfqpoint{2.757127in}{1.803555in}}%
\pgfpathlineto{\pgfqpoint{2.757127in}{1.803555in}}%
\pgfpathlineto{\pgfqpoint{2.757127in}{1.806505in}}%
\pgfpathlineto{\pgfqpoint{2.761668in}{1.806505in}}%
\pgfpathlineto{\pgfqpoint{2.761668in}{1.803555in}}%
\pgfpathmoveto{\pgfqpoint{2.761668in}{1.803555in}}%
\pgfpathlineto{\pgfqpoint{2.761668in}{1.803555in}}%
\pgfpathlineto{\pgfqpoint{2.761668in}{1.806505in}}%
\pgfpathlineto{\pgfqpoint{2.766209in}{1.806505in}}%
\pgfpathlineto{\pgfqpoint{2.766209in}{1.803555in}}%
\pgfpathmoveto{\pgfqpoint{2.761668in}{1.806505in}}%
\pgfpathlineto{\pgfqpoint{2.761668in}{1.806505in}}%
\pgfpathlineto{\pgfqpoint{2.761668in}{1.809454in}}%
\pgfpathlineto{\pgfqpoint{2.766209in}{1.809454in}}%
\pgfpathlineto{\pgfqpoint{2.766209in}{1.806505in}}%
\pgfpathmoveto{\pgfqpoint{2.766209in}{1.806505in}}%
\pgfpathlineto{\pgfqpoint{2.766209in}{1.806505in}}%
\pgfpathlineto{\pgfqpoint{2.766209in}{1.809454in}}%
\pgfpathlineto{\pgfqpoint{2.770750in}{1.809454in}}%
\pgfpathlineto{\pgfqpoint{2.770750in}{1.806505in}}%
\pgfpathmoveto{\pgfqpoint{2.766209in}{1.809454in}}%
\pgfpathlineto{\pgfqpoint{2.766209in}{1.809454in}}%
\pgfpathlineto{\pgfqpoint{2.766209in}{1.812403in}}%
\pgfpathlineto{\pgfqpoint{2.770750in}{1.812403in}}%
\pgfpathlineto{\pgfqpoint{2.770750in}{1.809454in}}%
\pgfpathmoveto{\pgfqpoint{2.770750in}{1.809454in}}%
\pgfpathlineto{\pgfqpoint{2.770750in}{1.809454in}}%
\pgfpathlineto{\pgfqpoint{2.770750in}{1.812403in}}%
\pgfpathlineto{\pgfqpoint{2.775291in}{1.812403in}}%
\pgfpathlineto{\pgfqpoint{2.775291in}{1.809454in}}%
\pgfpathmoveto{\pgfqpoint{2.770750in}{1.812403in}}%
\pgfpathlineto{\pgfqpoint{2.770750in}{1.812403in}}%
\pgfpathlineto{\pgfqpoint{2.770750in}{1.815352in}}%
\pgfpathlineto{\pgfqpoint{2.775291in}{1.815352in}}%
\pgfpathlineto{\pgfqpoint{2.775291in}{1.812403in}}%
\pgfpathmoveto{\pgfqpoint{2.775291in}{1.812403in}}%
\pgfpathlineto{\pgfqpoint{2.775291in}{1.812403in}}%
\pgfpathlineto{\pgfqpoint{2.775291in}{1.815352in}}%
\pgfpathlineto{\pgfqpoint{2.779833in}{1.815352in}}%
\pgfpathlineto{\pgfqpoint{2.779833in}{1.812403in}}%
\pgfpathmoveto{\pgfqpoint{2.775291in}{1.815352in}}%
\pgfpathlineto{\pgfqpoint{2.775291in}{1.815352in}}%
\pgfpathlineto{\pgfqpoint{2.775291in}{1.818301in}}%
\pgfpathlineto{\pgfqpoint{2.779833in}{1.818301in}}%
\pgfpathlineto{\pgfqpoint{2.779833in}{1.815352in}}%
\pgfpathmoveto{\pgfqpoint{2.779833in}{1.815352in}}%
\pgfpathlineto{\pgfqpoint{2.779833in}{1.815352in}}%
\pgfpathlineto{\pgfqpoint{2.779833in}{1.818301in}}%
\pgfpathlineto{\pgfqpoint{2.784374in}{1.818301in}}%
\pgfpathlineto{\pgfqpoint{2.784374in}{1.815352in}}%
\pgfpathmoveto{\pgfqpoint{2.779833in}{1.818301in}}%
\pgfpathlineto{\pgfqpoint{2.779833in}{1.818301in}}%
\pgfpathlineto{\pgfqpoint{2.779833in}{1.821250in}}%
\pgfpathlineto{\pgfqpoint{2.784374in}{1.821250in}}%
\pgfpathlineto{\pgfqpoint{2.784374in}{1.818301in}}%
\pgfpathmoveto{\pgfqpoint{2.784374in}{1.818301in}}%
\pgfpathlineto{\pgfqpoint{2.784374in}{1.818301in}}%
\pgfpathlineto{\pgfqpoint{2.784374in}{1.821250in}}%
\pgfpathlineto{\pgfqpoint{2.788915in}{1.821250in}}%
\pgfpathlineto{\pgfqpoint{2.788915in}{1.818301in}}%
\pgfpathmoveto{\pgfqpoint{2.784374in}{1.821250in}}%
\pgfpathlineto{\pgfqpoint{2.784374in}{1.821250in}}%
\pgfpathlineto{\pgfqpoint{2.784374in}{1.824200in}}%
\pgfpathlineto{\pgfqpoint{2.788915in}{1.824200in}}%
\pgfpathlineto{\pgfqpoint{2.788915in}{1.821250in}}%
\pgfpathmoveto{\pgfqpoint{2.788915in}{1.821250in}}%
\pgfpathlineto{\pgfqpoint{2.788915in}{1.821250in}}%
\pgfpathlineto{\pgfqpoint{2.788915in}{1.824200in}}%
\pgfpathlineto{\pgfqpoint{2.793456in}{1.824200in}}%
\pgfpathlineto{\pgfqpoint{2.793456in}{1.821250in}}%
\pgfpathmoveto{\pgfqpoint{2.788915in}{1.824200in}}%
\pgfpathlineto{\pgfqpoint{2.788915in}{1.824200in}}%
\pgfpathlineto{\pgfqpoint{2.788915in}{1.827149in}}%
\pgfpathlineto{\pgfqpoint{2.793456in}{1.827149in}}%
\pgfpathlineto{\pgfqpoint{2.793456in}{1.824200in}}%
\pgfpathmoveto{\pgfqpoint{2.793456in}{1.824200in}}%
\pgfpathlineto{\pgfqpoint{2.793456in}{1.824200in}}%
\pgfpathlineto{\pgfqpoint{2.793456in}{1.827149in}}%
\pgfpathlineto{\pgfqpoint{2.797997in}{1.827149in}}%
\pgfpathlineto{\pgfqpoint{2.797997in}{1.824200in}}%
\pgfpathmoveto{\pgfqpoint{2.793456in}{1.827149in}}%
\pgfpathlineto{\pgfqpoint{2.793456in}{1.827149in}}%
\pgfpathlineto{\pgfqpoint{2.793456in}{1.830098in}}%
\pgfpathlineto{\pgfqpoint{2.797997in}{1.830098in}}%
\pgfpathlineto{\pgfqpoint{2.797997in}{1.827149in}}%
\pgfpathmoveto{\pgfqpoint{2.797997in}{1.827149in}}%
\pgfpathlineto{\pgfqpoint{2.797997in}{1.827149in}}%
\pgfpathlineto{\pgfqpoint{2.797997in}{1.830098in}}%
\pgfpathlineto{\pgfqpoint{2.802538in}{1.830098in}}%
\pgfpathlineto{\pgfqpoint{2.802538in}{1.827149in}}%
\pgfpathmoveto{\pgfqpoint{2.797997in}{1.830098in}}%
\pgfpathlineto{\pgfqpoint{2.797997in}{1.830098in}}%
\pgfpathlineto{\pgfqpoint{2.797997in}{1.833048in}}%
\pgfpathlineto{\pgfqpoint{2.802538in}{1.833048in}}%
\pgfpathlineto{\pgfqpoint{2.802538in}{1.830098in}}%
\pgfpathmoveto{\pgfqpoint{2.802538in}{1.830098in}}%
\pgfpathlineto{\pgfqpoint{2.802538in}{1.830098in}}%
\pgfpathlineto{\pgfqpoint{2.802538in}{1.833048in}}%
\pgfpathlineto{\pgfqpoint{2.807079in}{1.833048in}}%
\pgfpathlineto{\pgfqpoint{2.807079in}{1.830098in}}%
\pgfpathmoveto{\pgfqpoint{2.802538in}{1.833048in}}%
\pgfpathlineto{\pgfqpoint{2.802538in}{1.833048in}}%
\pgfpathlineto{\pgfqpoint{2.802538in}{1.835997in}}%
\pgfpathlineto{\pgfqpoint{2.807079in}{1.835997in}}%
\pgfpathlineto{\pgfqpoint{2.807079in}{1.833048in}}%
\pgfpathmoveto{\pgfqpoint{2.807079in}{1.833048in}}%
\pgfpathlineto{\pgfqpoint{2.807079in}{1.833048in}}%
\pgfpathlineto{\pgfqpoint{2.807079in}{1.835997in}}%
\pgfpathlineto{\pgfqpoint{2.811620in}{1.835997in}}%
\pgfpathlineto{\pgfqpoint{2.811620in}{1.833048in}}%
\pgfpathmoveto{\pgfqpoint{2.807079in}{1.835997in}}%
\pgfpathlineto{\pgfqpoint{2.807079in}{1.835997in}}%
\pgfpathlineto{\pgfqpoint{2.807079in}{1.838946in}}%
\pgfpathlineto{\pgfqpoint{2.811620in}{1.838946in}}%
\pgfpathlineto{\pgfqpoint{2.811620in}{1.835997in}}%
\pgfpathmoveto{\pgfqpoint{2.811620in}{1.835997in}}%
\pgfpathlineto{\pgfqpoint{2.811620in}{1.835997in}}%
\pgfpathlineto{\pgfqpoint{2.811620in}{1.838946in}}%
\pgfpathlineto{\pgfqpoint{2.816161in}{1.838946in}}%
\pgfpathlineto{\pgfqpoint{2.816161in}{1.835997in}}%
\pgfpathmoveto{\pgfqpoint{2.811620in}{1.838946in}}%
\pgfpathlineto{\pgfqpoint{2.811620in}{1.838946in}}%
\pgfpathlineto{\pgfqpoint{2.811620in}{1.841895in}}%
\pgfpathlineto{\pgfqpoint{2.816161in}{1.841895in}}%
\pgfpathlineto{\pgfqpoint{2.816161in}{1.838946in}}%
\pgfpathmoveto{\pgfqpoint{2.816161in}{1.838946in}}%
\pgfpathlineto{\pgfqpoint{2.816161in}{1.838946in}}%
\pgfpathlineto{\pgfqpoint{2.816161in}{1.841895in}}%
\pgfpathlineto{\pgfqpoint{2.820702in}{1.841895in}}%
\pgfpathlineto{\pgfqpoint{2.820702in}{1.838946in}}%
\pgfpathmoveto{\pgfqpoint{2.816161in}{1.841895in}}%
\pgfpathlineto{\pgfqpoint{2.816161in}{1.841895in}}%
\pgfpathlineto{\pgfqpoint{2.816161in}{1.844845in}}%
\pgfpathlineto{\pgfqpoint{2.820702in}{1.844845in}}%
\pgfpathlineto{\pgfqpoint{2.820702in}{1.841895in}}%
\pgfpathmoveto{\pgfqpoint{2.820702in}{1.841895in}}%
\pgfpathlineto{\pgfqpoint{2.820702in}{1.841895in}}%
\pgfpathlineto{\pgfqpoint{2.820702in}{1.844845in}}%
\pgfpathlineto{\pgfqpoint{2.825243in}{1.844845in}}%
\pgfpathlineto{\pgfqpoint{2.825243in}{1.841895in}}%
\pgfpathmoveto{\pgfqpoint{2.820702in}{1.844845in}}%
\pgfpathlineto{\pgfqpoint{2.820702in}{1.844845in}}%
\pgfpathlineto{\pgfqpoint{2.820702in}{1.847794in}}%
\pgfpathlineto{\pgfqpoint{2.825243in}{1.847794in}}%
\pgfpathlineto{\pgfqpoint{2.825243in}{1.844845in}}%
\pgfpathmoveto{\pgfqpoint{2.825243in}{1.844845in}}%
\pgfpathlineto{\pgfqpoint{2.825243in}{1.844845in}}%
\pgfpathlineto{\pgfqpoint{2.825243in}{1.847794in}}%
\pgfpathlineto{\pgfqpoint{2.829784in}{1.847794in}}%
\pgfpathlineto{\pgfqpoint{2.829784in}{1.844845in}}%
\pgfpathmoveto{\pgfqpoint{2.825243in}{1.847794in}}%
\pgfpathlineto{\pgfqpoint{2.825243in}{1.847794in}}%
\pgfpathlineto{\pgfqpoint{2.825243in}{1.850743in}}%
\pgfpathlineto{\pgfqpoint{2.829784in}{1.850743in}}%
\pgfpathlineto{\pgfqpoint{2.829784in}{1.847794in}}%
\pgfpathmoveto{\pgfqpoint{2.829784in}{1.847794in}}%
\pgfpathlineto{\pgfqpoint{2.829784in}{1.847794in}}%
\pgfpathlineto{\pgfqpoint{2.829784in}{1.850743in}}%
\pgfpathlineto{\pgfqpoint{2.834325in}{1.850743in}}%
\pgfpathlineto{\pgfqpoint{2.834325in}{1.847794in}}%
\pgfpathmoveto{\pgfqpoint{2.829784in}{1.850743in}}%
\pgfpathlineto{\pgfqpoint{2.829784in}{1.850743in}}%
\pgfpathlineto{\pgfqpoint{2.829784in}{1.853692in}}%
\pgfpathlineto{\pgfqpoint{2.834325in}{1.853692in}}%
\pgfpathlineto{\pgfqpoint{2.834325in}{1.850743in}}%
\pgfpathmoveto{\pgfqpoint{2.834325in}{1.850743in}}%
\pgfpathlineto{\pgfqpoint{2.834325in}{1.850743in}}%
\pgfpathlineto{\pgfqpoint{2.834325in}{1.853692in}}%
\pgfpathlineto{\pgfqpoint{2.838866in}{1.853692in}}%
\pgfpathlineto{\pgfqpoint{2.838866in}{1.850743in}}%
\pgfpathmoveto{\pgfqpoint{2.834325in}{1.853692in}}%
\pgfpathlineto{\pgfqpoint{2.834325in}{1.853692in}}%
\pgfpathlineto{\pgfqpoint{2.834325in}{1.856642in}}%
\pgfpathlineto{\pgfqpoint{2.838866in}{1.856642in}}%
\pgfpathlineto{\pgfqpoint{2.838866in}{1.853692in}}%
\pgfpathmoveto{\pgfqpoint{2.838866in}{1.853692in}}%
\pgfpathlineto{\pgfqpoint{2.838866in}{1.853692in}}%
\pgfpathlineto{\pgfqpoint{2.838866in}{1.856642in}}%
\pgfpathlineto{\pgfqpoint{2.843407in}{1.856642in}}%
\pgfpathlineto{\pgfqpoint{2.843407in}{1.853692in}}%
\pgfpathmoveto{\pgfqpoint{2.838866in}{1.856642in}}%
\pgfpathlineto{\pgfqpoint{2.838866in}{1.856642in}}%
\pgfpathlineto{\pgfqpoint{2.838866in}{1.859591in}}%
\pgfpathlineto{\pgfqpoint{2.843407in}{1.859591in}}%
\pgfpathlineto{\pgfqpoint{2.843407in}{1.856642in}}%
\pgfpathmoveto{\pgfqpoint{2.843407in}{1.856642in}}%
\pgfpathlineto{\pgfqpoint{2.843407in}{1.856642in}}%
\pgfpathlineto{\pgfqpoint{2.843407in}{1.859591in}}%
\pgfpathlineto{\pgfqpoint{2.847948in}{1.859591in}}%
\pgfpathlineto{\pgfqpoint{2.847948in}{1.856642in}}%
\pgfpathmoveto{\pgfqpoint{2.843407in}{1.859591in}}%
\pgfpathlineto{\pgfqpoint{2.843407in}{1.859591in}}%
\pgfpathlineto{\pgfqpoint{2.843407in}{1.862540in}}%
\pgfpathlineto{\pgfqpoint{2.847948in}{1.862540in}}%
\pgfpathlineto{\pgfqpoint{2.847948in}{1.859591in}}%
\pgfpathmoveto{\pgfqpoint{2.847948in}{1.859591in}}%
\pgfpathlineto{\pgfqpoint{2.847948in}{1.859591in}}%
\pgfpathlineto{\pgfqpoint{2.847948in}{1.862540in}}%
\pgfpathlineto{\pgfqpoint{2.852489in}{1.862540in}}%
\pgfpathlineto{\pgfqpoint{2.852489in}{1.859591in}}%
\pgfpathmoveto{\pgfqpoint{2.847948in}{1.862540in}}%
\pgfpathlineto{\pgfqpoint{2.847948in}{1.862540in}}%
\pgfpathlineto{\pgfqpoint{2.847948in}{1.865489in}}%
\pgfpathlineto{\pgfqpoint{2.852489in}{1.865489in}}%
\pgfpathlineto{\pgfqpoint{2.852489in}{1.862540in}}%
\pgfpathmoveto{\pgfqpoint{2.852489in}{1.862540in}}%
\pgfpathlineto{\pgfqpoint{2.852489in}{1.862540in}}%
\pgfpathlineto{\pgfqpoint{2.852489in}{1.865489in}}%
\pgfpathlineto{\pgfqpoint{2.857030in}{1.865489in}}%
\pgfpathlineto{\pgfqpoint{2.857030in}{1.862540in}}%
\pgfpathmoveto{\pgfqpoint{2.852489in}{1.865489in}}%
\pgfpathlineto{\pgfqpoint{2.852489in}{1.865489in}}%
\pgfpathlineto{\pgfqpoint{2.852489in}{1.868439in}}%
\pgfpathlineto{\pgfqpoint{2.857030in}{1.868439in}}%
\pgfpathlineto{\pgfqpoint{2.857030in}{1.865489in}}%
\pgfpathmoveto{\pgfqpoint{2.857030in}{1.865489in}}%
\pgfpathlineto{\pgfqpoint{2.857030in}{1.865489in}}%
\pgfpathlineto{\pgfqpoint{2.857030in}{1.868439in}}%
\pgfpathlineto{\pgfqpoint{2.861571in}{1.868439in}}%
\pgfpathlineto{\pgfqpoint{2.861571in}{1.865489in}}%
\pgfpathmoveto{\pgfqpoint{2.857030in}{1.868439in}}%
\pgfpathlineto{\pgfqpoint{2.857030in}{1.868439in}}%
\pgfpathlineto{\pgfqpoint{2.857030in}{1.871388in}}%
\pgfpathlineto{\pgfqpoint{2.861571in}{1.871388in}}%
\pgfpathlineto{\pgfqpoint{2.861571in}{1.868439in}}%
\pgfpathmoveto{\pgfqpoint{2.861571in}{1.868439in}}%
\pgfpathlineto{\pgfqpoint{2.861571in}{1.868439in}}%
\pgfpathlineto{\pgfqpoint{2.861571in}{1.871388in}}%
\pgfpathlineto{\pgfqpoint{2.866112in}{1.871388in}}%
\pgfpathlineto{\pgfqpoint{2.866112in}{1.868439in}}%
\pgfpathmoveto{\pgfqpoint{2.861571in}{1.871388in}}%
\pgfpathlineto{\pgfqpoint{2.861571in}{1.871388in}}%
\pgfpathlineto{\pgfqpoint{2.861571in}{1.874337in}}%
\pgfpathlineto{\pgfqpoint{2.866112in}{1.874337in}}%
\pgfpathlineto{\pgfqpoint{2.866112in}{1.871388in}}%
\pgfpathmoveto{\pgfqpoint{2.866112in}{1.871388in}}%
\pgfpathlineto{\pgfqpoint{2.866112in}{1.871388in}}%
\pgfpathlineto{\pgfqpoint{2.866112in}{1.874337in}}%
\pgfpathlineto{\pgfqpoint{2.870653in}{1.874337in}}%
\pgfpathlineto{\pgfqpoint{2.870653in}{1.871388in}}%
\pgfpathmoveto{\pgfqpoint{2.866112in}{1.874337in}}%
\pgfpathlineto{\pgfqpoint{2.866112in}{1.874337in}}%
\pgfpathlineto{\pgfqpoint{2.866112in}{1.877286in}}%
\pgfpathlineto{\pgfqpoint{2.870653in}{1.877286in}}%
\pgfpathlineto{\pgfqpoint{2.870653in}{1.874337in}}%
\pgfpathmoveto{\pgfqpoint{2.870653in}{1.874337in}}%
\pgfpathlineto{\pgfqpoint{2.870653in}{1.874337in}}%
\pgfpathlineto{\pgfqpoint{2.870653in}{1.877286in}}%
\pgfpathlineto{\pgfqpoint{2.875194in}{1.877286in}}%
\pgfpathlineto{\pgfqpoint{2.875194in}{1.874337in}}%
\pgfpathmoveto{\pgfqpoint{2.870653in}{1.877286in}}%
\pgfpathlineto{\pgfqpoint{2.870653in}{1.877286in}}%
\pgfpathlineto{\pgfqpoint{2.870653in}{1.880236in}}%
\pgfpathlineto{\pgfqpoint{2.875194in}{1.880236in}}%
\pgfpathlineto{\pgfqpoint{2.875194in}{1.877286in}}%
\pgfpathmoveto{\pgfqpoint{2.875194in}{1.877286in}}%
\pgfpathlineto{\pgfqpoint{2.875194in}{1.877286in}}%
\pgfpathlineto{\pgfqpoint{2.875194in}{1.880236in}}%
\pgfpathlineto{\pgfqpoint{2.879735in}{1.880236in}}%
\pgfpathlineto{\pgfqpoint{2.879735in}{1.877286in}}%
\pgfpathmoveto{\pgfqpoint{2.875194in}{1.880236in}}%
\pgfpathlineto{\pgfqpoint{2.875194in}{1.880236in}}%
\pgfpathlineto{\pgfqpoint{2.875194in}{1.883185in}}%
\pgfpathlineto{\pgfqpoint{2.879735in}{1.883185in}}%
\pgfpathlineto{\pgfqpoint{2.879735in}{1.880236in}}%
\pgfpathmoveto{\pgfqpoint{2.879735in}{1.880236in}}%
\pgfpathlineto{\pgfqpoint{2.879735in}{1.880236in}}%
\pgfpathlineto{\pgfqpoint{2.879735in}{1.883185in}}%
\pgfpathlineto{\pgfqpoint{2.884276in}{1.883185in}}%
\pgfpathlineto{\pgfqpoint{2.884276in}{1.880236in}}%
\pgfpathmoveto{\pgfqpoint{2.879735in}{1.883185in}}%
\pgfpathlineto{\pgfqpoint{2.879735in}{1.883185in}}%
\pgfpathlineto{\pgfqpoint{2.879735in}{1.886134in}}%
\pgfpathlineto{\pgfqpoint{2.884276in}{1.886134in}}%
\pgfpathlineto{\pgfqpoint{2.884276in}{1.883185in}}%
\pgfpathmoveto{\pgfqpoint{2.884276in}{1.883185in}}%
\pgfpathlineto{\pgfqpoint{2.884276in}{1.883185in}}%
\pgfpathlineto{\pgfqpoint{2.884276in}{1.886134in}}%
\pgfpathlineto{\pgfqpoint{2.888817in}{1.886134in}}%
\pgfpathlineto{\pgfqpoint{2.888817in}{1.883185in}}%
\pgfpathmoveto{\pgfqpoint{2.884276in}{1.886134in}}%
\pgfpathlineto{\pgfqpoint{2.884276in}{1.886134in}}%
\pgfpathlineto{\pgfqpoint{2.884276in}{1.889083in}}%
\pgfpathlineto{\pgfqpoint{2.888817in}{1.889083in}}%
\pgfpathlineto{\pgfqpoint{2.888817in}{1.886134in}}%
\pgfpathmoveto{\pgfqpoint{2.888817in}{1.886134in}}%
\pgfpathlineto{\pgfqpoint{2.888817in}{1.886134in}}%
\pgfpathlineto{\pgfqpoint{2.888817in}{1.889083in}}%
\pgfpathlineto{\pgfqpoint{2.893358in}{1.889083in}}%
\pgfpathlineto{\pgfqpoint{2.893358in}{1.886134in}}%
\pgfpathmoveto{\pgfqpoint{2.888817in}{1.889083in}}%
\pgfpathlineto{\pgfqpoint{2.888817in}{1.889083in}}%
\pgfpathlineto{\pgfqpoint{2.888817in}{1.892033in}}%
\pgfpathlineto{\pgfqpoint{2.893358in}{1.892033in}}%
\pgfpathlineto{\pgfqpoint{2.893358in}{1.889083in}}%
\pgfpathmoveto{\pgfqpoint{2.893358in}{1.889083in}}%
\pgfpathlineto{\pgfqpoint{2.893358in}{1.889083in}}%
\pgfpathlineto{\pgfqpoint{2.893358in}{1.892033in}}%
\pgfpathlineto{\pgfqpoint{2.897899in}{1.892033in}}%
\pgfpathlineto{\pgfqpoint{2.897899in}{1.889083in}}%
\pgfpathmoveto{\pgfqpoint{2.893358in}{1.892033in}}%
\pgfpathlineto{\pgfqpoint{2.893358in}{1.892033in}}%
\pgfpathlineto{\pgfqpoint{2.893358in}{1.894982in}}%
\pgfpathlineto{\pgfqpoint{2.897899in}{1.894982in}}%
\pgfpathlineto{\pgfqpoint{2.897899in}{1.892033in}}%
\pgfpathmoveto{\pgfqpoint{2.897899in}{1.892033in}}%
\pgfpathlineto{\pgfqpoint{2.897899in}{1.892033in}}%
\pgfpathlineto{\pgfqpoint{2.897899in}{1.894982in}}%
\pgfpathlineto{\pgfqpoint{2.902440in}{1.894982in}}%
\pgfpathlineto{\pgfqpoint{2.902440in}{1.892033in}}%
\pgfpathmoveto{\pgfqpoint{2.897899in}{1.894982in}}%
\pgfpathlineto{\pgfqpoint{2.897899in}{1.894982in}}%
\pgfpathlineto{\pgfqpoint{2.897899in}{1.897931in}}%
\pgfpathlineto{\pgfqpoint{2.902440in}{1.897931in}}%
\pgfpathlineto{\pgfqpoint{2.902440in}{1.894982in}}%
\pgfpathmoveto{\pgfqpoint{2.902440in}{1.894982in}}%
\pgfpathlineto{\pgfqpoint{2.902440in}{1.894982in}}%
\pgfpathlineto{\pgfqpoint{2.902440in}{1.897931in}}%
\pgfpathlineto{\pgfqpoint{2.906981in}{1.897931in}}%
\pgfpathlineto{\pgfqpoint{2.906981in}{1.894982in}}%
\pgfpathmoveto{\pgfqpoint{2.902440in}{1.897931in}}%
\pgfpathlineto{\pgfqpoint{2.902440in}{1.897931in}}%
\pgfpathlineto{\pgfqpoint{2.902440in}{1.900880in}}%
\pgfpathlineto{\pgfqpoint{2.906981in}{1.900880in}}%
\pgfpathlineto{\pgfqpoint{2.906981in}{1.897931in}}%
\pgfpathmoveto{\pgfqpoint{2.906981in}{1.897931in}}%
\pgfpathlineto{\pgfqpoint{2.906981in}{1.897931in}}%
\pgfpathlineto{\pgfqpoint{2.906981in}{1.900880in}}%
\pgfpathlineto{\pgfqpoint{2.911522in}{1.900880in}}%
\pgfpathlineto{\pgfqpoint{2.911522in}{1.897931in}}%
\pgfpathmoveto{\pgfqpoint{2.906981in}{1.900880in}}%
\pgfpathlineto{\pgfqpoint{2.906981in}{1.900880in}}%
\pgfpathlineto{\pgfqpoint{2.906981in}{1.903830in}}%
\pgfpathlineto{\pgfqpoint{2.911522in}{1.903830in}}%
\pgfpathlineto{\pgfqpoint{2.911522in}{1.900880in}}%
\pgfpathmoveto{\pgfqpoint{2.911522in}{1.900880in}}%
\pgfpathlineto{\pgfqpoint{2.911522in}{1.900880in}}%
\pgfpathlineto{\pgfqpoint{2.911522in}{1.903830in}}%
\pgfpathlineto{\pgfqpoint{2.916063in}{1.903830in}}%
\pgfpathlineto{\pgfqpoint{2.916063in}{1.900880in}}%
\pgfpathmoveto{\pgfqpoint{2.911522in}{1.903830in}}%
\pgfpathlineto{\pgfqpoint{2.911522in}{1.903830in}}%
\pgfpathlineto{\pgfqpoint{2.911522in}{1.906779in}}%
\pgfpathlineto{\pgfqpoint{2.916063in}{1.906779in}}%
\pgfpathlineto{\pgfqpoint{2.916063in}{1.903830in}}%
\pgfpathmoveto{\pgfqpoint{2.916063in}{1.903830in}}%
\pgfpathlineto{\pgfqpoint{2.916063in}{1.903830in}}%
\pgfpathlineto{\pgfqpoint{2.916063in}{1.906779in}}%
\pgfpathlineto{\pgfqpoint{2.920604in}{1.906779in}}%
\pgfpathlineto{\pgfqpoint{2.920604in}{1.903830in}}%
\pgfpathmoveto{\pgfqpoint{2.916063in}{1.906779in}}%
\pgfpathlineto{\pgfqpoint{2.916063in}{1.906779in}}%
\pgfpathlineto{\pgfqpoint{2.916063in}{1.909728in}}%
\pgfpathlineto{\pgfqpoint{2.920604in}{1.909728in}}%
\pgfpathlineto{\pgfqpoint{2.920604in}{1.906779in}}%
\pgfpathmoveto{\pgfqpoint{2.920604in}{1.906779in}}%
\pgfpathlineto{\pgfqpoint{2.920604in}{1.906779in}}%
\pgfpathlineto{\pgfqpoint{2.920604in}{1.909728in}}%
\pgfpathlineto{\pgfqpoint{2.925145in}{1.909728in}}%
\pgfpathlineto{\pgfqpoint{2.925145in}{1.906779in}}%
\pgfpathmoveto{\pgfqpoint{2.920604in}{1.909728in}}%
\pgfpathlineto{\pgfqpoint{2.920604in}{1.909728in}}%
\pgfpathlineto{\pgfqpoint{2.920604in}{1.912678in}}%
\pgfpathlineto{\pgfqpoint{2.925145in}{1.912678in}}%
\pgfpathlineto{\pgfqpoint{2.925145in}{1.909728in}}%
\pgfpathmoveto{\pgfqpoint{2.925145in}{1.909728in}}%
\pgfpathlineto{\pgfqpoint{2.925145in}{1.909728in}}%
\pgfpathlineto{\pgfqpoint{2.925145in}{1.912678in}}%
\pgfpathlineto{\pgfqpoint{2.929686in}{1.912678in}}%
\pgfpathlineto{\pgfqpoint{2.929686in}{1.909728in}}%
\pgfpathmoveto{\pgfqpoint{2.925145in}{1.912678in}}%
\pgfpathlineto{\pgfqpoint{2.925145in}{1.912678in}}%
\pgfpathlineto{\pgfqpoint{2.925145in}{1.915627in}}%
\pgfpathlineto{\pgfqpoint{2.929686in}{1.915627in}}%
\pgfpathlineto{\pgfqpoint{2.929686in}{1.912678in}}%
\pgfpathmoveto{\pgfqpoint{2.929686in}{1.912678in}}%
\pgfpathlineto{\pgfqpoint{2.929686in}{1.912678in}}%
\pgfpathlineto{\pgfqpoint{2.929686in}{1.915627in}}%
\pgfpathlineto{\pgfqpoint{2.934227in}{1.915627in}}%
\pgfpathlineto{\pgfqpoint{2.934227in}{1.912678in}}%
\pgfpathmoveto{\pgfqpoint{2.929686in}{1.915627in}}%
\pgfpathlineto{\pgfqpoint{2.929686in}{1.915627in}}%
\pgfpathlineto{\pgfqpoint{2.929686in}{1.918576in}}%
\pgfpathlineto{\pgfqpoint{2.934227in}{1.918576in}}%
\pgfpathlineto{\pgfqpoint{2.934227in}{1.915627in}}%
\pgfpathmoveto{\pgfqpoint{2.934227in}{1.915627in}}%
\pgfpathlineto{\pgfqpoint{2.934227in}{1.915627in}}%
\pgfpathlineto{\pgfqpoint{2.934227in}{1.918576in}}%
\pgfpathlineto{\pgfqpoint{2.938768in}{1.918576in}}%
\pgfpathlineto{\pgfqpoint{2.938768in}{1.915627in}}%
\pgfpathmoveto{\pgfqpoint{2.934227in}{1.918576in}}%
\pgfpathlineto{\pgfqpoint{2.934227in}{1.918576in}}%
\pgfpathlineto{\pgfqpoint{2.934227in}{1.921525in}}%
\pgfpathlineto{\pgfqpoint{2.938768in}{1.921525in}}%
\pgfpathlineto{\pgfqpoint{2.938768in}{1.918576in}}%
\pgfpathmoveto{\pgfqpoint{2.938768in}{1.918576in}}%
\pgfpathlineto{\pgfqpoint{2.938768in}{1.918576in}}%
\pgfpathlineto{\pgfqpoint{2.938768in}{1.921525in}}%
\pgfpathlineto{\pgfqpoint{2.943309in}{1.921525in}}%
\pgfpathlineto{\pgfqpoint{2.943309in}{1.918576in}}%
\pgfpathmoveto{\pgfqpoint{2.938768in}{1.921525in}}%
\pgfpathlineto{\pgfqpoint{2.938768in}{1.921525in}}%
\pgfpathlineto{\pgfqpoint{2.938768in}{1.924474in}}%
\pgfpathlineto{\pgfqpoint{2.943309in}{1.924474in}}%
\pgfpathlineto{\pgfqpoint{2.943309in}{1.921525in}}%
\pgfpathmoveto{\pgfqpoint{2.943309in}{1.921525in}}%
\pgfpathlineto{\pgfqpoint{2.943309in}{1.921525in}}%
\pgfpathlineto{\pgfqpoint{2.943309in}{1.924474in}}%
\pgfpathlineto{\pgfqpoint{2.947850in}{1.924474in}}%
\pgfpathlineto{\pgfqpoint{2.947850in}{1.921525in}}%
\pgfpathmoveto{\pgfqpoint{2.943309in}{1.924474in}}%
\pgfpathlineto{\pgfqpoint{2.943309in}{1.924474in}}%
\pgfpathlineto{\pgfqpoint{2.943309in}{1.927423in}}%
\pgfpathlineto{\pgfqpoint{2.947850in}{1.927423in}}%
\pgfpathlineto{\pgfqpoint{2.947850in}{1.924474in}}%
\pgfpathmoveto{\pgfqpoint{2.947850in}{1.924474in}}%
\pgfpathlineto{\pgfqpoint{2.947850in}{1.924474in}}%
\pgfpathlineto{\pgfqpoint{2.947850in}{1.927423in}}%
\pgfpathlineto{\pgfqpoint{2.952391in}{1.927423in}}%
\pgfpathlineto{\pgfqpoint{2.952391in}{1.924474in}}%
\pgfpathmoveto{\pgfqpoint{2.947850in}{1.927423in}}%
\pgfpathlineto{\pgfqpoint{2.947850in}{1.927423in}}%
\pgfpathlineto{\pgfqpoint{2.947850in}{1.930373in}}%
\pgfpathlineto{\pgfqpoint{2.952391in}{1.930373in}}%
\pgfpathlineto{\pgfqpoint{2.952391in}{1.927423in}}%
\pgfpathmoveto{\pgfqpoint{2.952391in}{1.927423in}}%
\pgfpathlineto{\pgfqpoint{2.952391in}{1.927423in}}%
\pgfpathlineto{\pgfqpoint{2.952391in}{1.930373in}}%
\pgfpathlineto{\pgfqpoint{2.956932in}{1.930373in}}%
\pgfpathlineto{\pgfqpoint{2.956932in}{1.927423in}}%
\pgfpathmoveto{\pgfqpoint{2.952391in}{1.930373in}}%
\pgfpathlineto{\pgfqpoint{2.952391in}{1.930373in}}%
\pgfpathlineto{\pgfqpoint{2.952391in}{1.933322in}}%
\pgfpathlineto{\pgfqpoint{2.956932in}{1.933322in}}%
\pgfpathlineto{\pgfqpoint{2.956932in}{1.930373in}}%
\pgfpathmoveto{\pgfqpoint{2.956932in}{1.930373in}}%
\pgfpathlineto{\pgfqpoint{2.956932in}{1.930373in}}%
\pgfpathlineto{\pgfqpoint{2.956932in}{1.933322in}}%
\pgfpathlineto{\pgfqpoint{2.961473in}{1.933322in}}%
\pgfpathlineto{\pgfqpoint{2.961473in}{1.930373in}}%
\pgfpathmoveto{\pgfqpoint{2.956932in}{1.933322in}}%
\pgfpathlineto{\pgfqpoint{2.956932in}{1.933322in}}%
\pgfpathlineto{\pgfqpoint{2.956932in}{1.936271in}}%
\pgfpathlineto{\pgfqpoint{2.961473in}{1.936271in}}%
\pgfpathlineto{\pgfqpoint{2.961473in}{1.933322in}}%
\pgfpathmoveto{\pgfqpoint{2.961473in}{1.933322in}}%
\pgfpathlineto{\pgfqpoint{2.961473in}{1.933322in}}%
\pgfpathlineto{\pgfqpoint{2.961473in}{1.936271in}}%
\pgfpathlineto{\pgfqpoint{2.966014in}{1.936271in}}%
\pgfpathlineto{\pgfqpoint{2.966014in}{1.933322in}}%
\pgfpathmoveto{\pgfqpoint{2.961473in}{1.936271in}}%
\pgfpathlineto{\pgfqpoint{2.961473in}{1.936271in}}%
\pgfpathlineto{\pgfqpoint{2.961473in}{1.939220in}}%
\pgfpathlineto{\pgfqpoint{2.966014in}{1.939220in}}%
\pgfpathlineto{\pgfqpoint{2.966014in}{1.936271in}}%
\pgfpathmoveto{\pgfqpoint{2.966014in}{1.936271in}}%
\pgfpathlineto{\pgfqpoint{2.966014in}{1.936271in}}%
\pgfpathlineto{\pgfqpoint{2.966014in}{1.939220in}}%
\pgfpathlineto{\pgfqpoint{2.970555in}{1.939220in}}%
\pgfpathlineto{\pgfqpoint{2.970555in}{1.936271in}}%
\pgfpathmoveto{\pgfqpoint{2.966014in}{1.939220in}}%
\pgfpathlineto{\pgfqpoint{2.966014in}{1.939220in}}%
\pgfpathlineto{\pgfqpoint{2.966014in}{1.942169in}}%
\pgfpathlineto{\pgfqpoint{2.970555in}{1.942169in}}%
\pgfpathlineto{\pgfqpoint{2.970555in}{1.939220in}}%
\pgfpathmoveto{\pgfqpoint{2.970555in}{1.939220in}}%
\pgfpathlineto{\pgfqpoint{2.970555in}{1.939220in}}%
\pgfpathlineto{\pgfqpoint{2.970555in}{1.942169in}}%
\pgfpathlineto{\pgfqpoint{2.975096in}{1.942169in}}%
\pgfpathlineto{\pgfqpoint{2.975096in}{1.939220in}}%
\pgfpathmoveto{\pgfqpoint{2.970555in}{1.942169in}}%
\pgfpathlineto{\pgfqpoint{2.970555in}{1.942169in}}%
\pgfpathlineto{\pgfqpoint{2.970555in}{1.945119in}}%
\pgfpathlineto{\pgfqpoint{2.975096in}{1.945119in}}%
\pgfpathlineto{\pgfqpoint{2.975096in}{1.942169in}}%
\pgfpathmoveto{\pgfqpoint{2.975096in}{1.942169in}}%
\pgfpathlineto{\pgfqpoint{2.975096in}{1.942169in}}%
\pgfpathlineto{\pgfqpoint{2.975096in}{1.945119in}}%
\pgfpathlineto{\pgfqpoint{2.979637in}{1.945119in}}%
\pgfpathlineto{\pgfqpoint{2.979637in}{1.942169in}}%
\pgfpathmoveto{\pgfqpoint{2.975096in}{1.945119in}}%
\pgfpathlineto{\pgfqpoint{2.975096in}{1.945119in}}%
\pgfpathlineto{\pgfqpoint{2.975096in}{1.948068in}}%
\pgfpathlineto{\pgfqpoint{2.979637in}{1.948068in}}%
\pgfpathlineto{\pgfqpoint{2.979637in}{1.945119in}}%
\pgfpathmoveto{\pgfqpoint{2.979637in}{1.945119in}}%
\pgfpathlineto{\pgfqpoint{2.979637in}{1.945119in}}%
\pgfpathlineto{\pgfqpoint{2.979637in}{1.948068in}}%
\pgfpathlineto{\pgfqpoint{2.984178in}{1.948068in}}%
\pgfpathlineto{\pgfqpoint{2.984178in}{1.945119in}}%
\pgfpathmoveto{\pgfqpoint{2.979637in}{1.948068in}}%
\pgfpathlineto{\pgfqpoint{2.979637in}{1.948068in}}%
\pgfpathlineto{\pgfqpoint{2.979637in}{1.951017in}}%
\pgfpathlineto{\pgfqpoint{2.984178in}{1.951017in}}%
\pgfpathlineto{\pgfqpoint{2.984178in}{1.948068in}}%
\pgfpathmoveto{\pgfqpoint{2.984178in}{1.948068in}}%
\pgfpathlineto{\pgfqpoint{2.984178in}{1.948068in}}%
\pgfpathlineto{\pgfqpoint{2.984178in}{1.951017in}}%
\pgfpathlineto{\pgfqpoint{2.988719in}{1.951017in}}%
\pgfpathlineto{\pgfqpoint{2.988719in}{1.948068in}}%
\pgfpathmoveto{\pgfqpoint{2.984178in}{1.951017in}}%
\pgfpathlineto{\pgfqpoint{2.984178in}{1.951017in}}%
\pgfpathlineto{\pgfqpoint{2.984178in}{1.953966in}}%
\pgfpathlineto{\pgfqpoint{2.988719in}{1.953966in}}%
\pgfpathlineto{\pgfqpoint{2.988719in}{1.951017in}}%
\pgfpathmoveto{\pgfqpoint{2.988719in}{1.951017in}}%
\pgfpathlineto{\pgfqpoint{2.988719in}{1.951017in}}%
\pgfpathlineto{\pgfqpoint{2.988719in}{1.953966in}}%
\pgfpathlineto{\pgfqpoint{2.993260in}{1.953966in}}%
\pgfpathlineto{\pgfqpoint{2.993260in}{1.951017in}}%
\pgfpathmoveto{\pgfqpoint{2.988719in}{1.953966in}}%
\pgfpathlineto{\pgfqpoint{2.988719in}{1.953966in}}%
\pgfpathlineto{\pgfqpoint{2.988719in}{1.956915in}}%
\pgfpathlineto{\pgfqpoint{2.993260in}{1.956915in}}%
\pgfpathlineto{\pgfqpoint{2.993260in}{1.953966in}}%
\pgfpathmoveto{\pgfqpoint{2.993260in}{1.953966in}}%
\pgfpathlineto{\pgfqpoint{2.993260in}{1.953966in}}%
\pgfpathlineto{\pgfqpoint{2.993260in}{1.956915in}}%
\pgfpathlineto{\pgfqpoint{2.997801in}{1.956915in}}%
\pgfpathlineto{\pgfqpoint{2.997801in}{1.953966in}}%
\pgfpathmoveto{\pgfqpoint{2.993260in}{1.956915in}}%
\pgfpathlineto{\pgfqpoint{2.993260in}{1.956915in}}%
\pgfpathlineto{\pgfqpoint{2.993260in}{1.959864in}}%
\pgfpathlineto{\pgfqpoint{2.997801in}{1.959864in}}%
\pgfpathlineto{\pgfqpoint{2.997801in}{1.956915in}}%
\pgfpathmoveto{\pgfqpoint{2.997801in}{1.956915in}}%
\pgfpathlineto{\pgfqpoint{2.997801in}{1.956915in}}%
\pgfpathlineto{\pgfqpoint{2.997801in}{1.959864in}}%
\pgfpathlineto{\pgfqpoint{3.002342in}{1.959864in}}%
\pgfpathlineto{\pgfqpoint{3.002342in}{1.956915in}}%
\pgfpathmoveto{\pgfqpoint{2.997801in}{1.959864in}}%
\pgfpathlineto{\pgfqpoint{2.997801in}{1.959864in}}%
\pgfpathlineto{\pgfqpoint{2.997801in}{1.962814in}}%
\pgfpathlineto{\pgfqpoint{3.002342in}{1.962814in}}%
\pgfpathlineto{\pgfqpoint{3.002342in}{1.959864in}}%
\pgfpathmoveto{\pgfqpoint{3.002342in}{1.959864in}}%
\pgfpathlineto{\pgfqpoint{3.002342in}{1.959864in}}%
\pgfpathlineto{\pgfqpoint{3.002342in}{1.962814in}}%
\pgfpathlineto{\pgfqpoint{3.006882in}{1.962814in}}%
\pgfpathlineto{\pgfqpoint{3.006882in}{1.959864in}}%
\pgfpathmoveto{\pgfqpoint{3.002342in}{1.962814in}}%
\pgfpathlineto{\pgfqpoint{3.002342in}{1.962814in}}%
\pgfpathlineto{\pgfqpoint{3.002342in}{1.965763in}}%
\pgfpathlineto{\pgfqpoint{3.006882in}{1.965763in}}%
\pgfpathlineto{\pgfqpoint{3.006882in}{1.962814in}}%
\pgfpathmoveto{\pgfqpoint{3.006882in}{1.962814in}}%
\pgfpathlineto{\pgfqpoint{3.006882in}{1.962814in}}%
\pgfpathlineto{\pgfqpoint{3.006882in}{1.965763in}}%
\pgfpathlineto{\pgfqpoint{3.011423in}{1.965763in}}%
\pgfpathlineto{\pgfqpoint{3.011423in}{1.962814in}}%
\pgfpathmoveto{\pgfqpoint{3.006882in}{1.965763in}}%
\pgfpathlineto{\pgfqpoint{3.006882in}{1.965763in}}%
\pgfpathlineto{\pgfqpoint{3.006882in}{1.968712in}}%
\pgfpathlineto{\pgfqpoint{3.011423in}{1.968712in}}%
\pgfpathlineto{\pgfqpoint{3.011423in}{1.965763in}}%
\pgfpathmoveto{\pgfqpoint{3.011423in}{1.965763in}}%
\pgfpathlineto{\pgfqpoint{3.011423in}{1.965763in}}%
\pgfpathlineto{\pgfqpoint{3.011423in}{1.968712in}}%
\pgfpathlineto{\pgfqpoint{3.015964in}{1.968712in}}%
\pgfpathlineto{\pgfqpoint{3.015964in}{1.965763in}}%
\pgfpathmoveto{\pgfqpoint{3.011423in}{1.968712in}}%
\pgfpathlineto{\pgfqpoint{3.011423in}{1.968712in}}%
\pgfpathlineto{\pgfqpoint{3.011423in}{1.971661in}}%
\pgfpathlineto{\pgfqpoint{3.015964in}{1.971661in}}%
\pgfpathlineto{\pgfqpoint{3.015964in}{1.968712in}}%
\pgfpathmoveto{\pgfqpoint{3.015964in}{1.968712in}}%
\pgfpathlineto{\pgfqpoint{3.015964in}{1.968712in}}%
\pgfpathlineto{\pgfqpoint{3.015964in}{1.971661in}}%
\pgfpathlineto{\pgfqpoint{3.020505in}{1.971661in}}%
\pgfpathlineto{\pgfqpoint{3.020505in}{1.968712in}}%
\pgfpathmoveto{\pgfqpoint{3.015964in}{1.971661in}}%
\pgfpathlineto{\pgfqpoint{3.015964in}{1.971661in}}%
\pgfpathlineto{\pgfqpoint{3.015964in}{1.974610in}}%
\pgfpathlineto{\pgfqpoint{3.020505in}{1.974610in}}%
\pgfpathlineto{\pgfqpoint{3.020505in}{1.971661in}}%
\pgfpathmoveto{\pgfqpoint{3.020505in}{1.971661in}}%
\pgfpathlineto{\pgfqpoint{3.020505in}{1.971661in}}%
\pgfpathlineto{\pgfqpoint{3.020505in}{1.974610in}}%
\pgfpathlineto{\pgfqpoint{3.025046in}{1.974610in}}%
\pgfpathlineto{\pgfqpoint{3.025046in}{1.971661in}}%
\pgfpathmoveto{\pgfqpoint{3.020505in}{1.974610in}}%
\pgfpathlineto{\pgfqpoint{3.020505in}{1.974610in}}%
\pgfpathlineto{\pgfqpoint{3.020505in}{1.977560in}}%
\pgfpathlineto{\pgfqpoint{3.025046in}{1.977560in}}%
\pgfpathlineto{\pgfqpoint{3.025046in}{1.974610in}}%
\pgfpathmoveto{\pgfqpoint{3.025046in}{1.974610in}}%
\pgfpathlineto{\pgfqpoint{3.025046in}{1.974610in}}%
\pgfpathlineto{\pgfqpoint{3.025046in}{1.977560in}}%
\pgfpathlineto{\pgfqpoint{3.029587in}{1.977560in}}%
\pgfpathlineto{\pgfqpoint{3.029587in}{1.974610in}}%
\pgfpathmoveto{\pgfqpoint{3.025046in}{1.977560in}}%
\pgfpathlineto{\pgfqpoint{3.025046in}{1.977560in}}%
\pgfpathlineto{\pgfqpoint{3.025046in}{1.980509in}}%
\pgfpathlineto{\pgfqpoint{3.029587in}{1.980509in}}%
\pgfpathlineto{\pgfqpoint{3.029587in}{1.977560in}}%
\pgfpathmoveto{\pgfqpoint{3.029587in}{1.977560in}}%
\pgfpathlineto{\pgfqpoint{3.029587in}{1.977560in}}%
\pgfpathlineto{\pgfqpoint{3.029587in}{1.980509in}}%
\pgfpathlineto{\pgfqpoint{3.034128in}{1.980509in}}%
\pgfpathlineto{\pgfqpoint{3.034128in}{1.977560in}}%
\pgfpathmoveto{\pgfqpoint{3.029587in}{1.980509in}}%
\pgfpathlineto{\pgfqpoint{3.029587in}{1.980509in}}%
\pgfpathlineto{\pgfqpoint{3.029587in}{1.983458in}}%
\pgfpathlineto{\pgfqpoint{3.034128in}{1.983458in}}%
\pgfpathlineto{\pgfqpoint{3.034128in}{1.980509in}}%
\pgfpathmoveto{\pgfqpoint{3.034128in}{1.980509in}}%
\pgfpathlineto{\pgfqpoint{3.034128in}{1.980509in}}%
\pgfpathlineto{\pgfqpoint{3.034128in}{1.983458in}}%
\pgfpathlineto{\pgfqpoint{3.038669in}{1.983458in}}%
\pgfpathlineto{\pgfqpoint{3.038669in}{1.980509in}}%
\pgfpathmoveto{\pgfqpoint{3.034128in}{1.983458in}}%
\pgfpathlineto{\pgfqpoint{3.034128in}{1.983458in}}%
\pgfpathlineto{\pgfqpoint{3.034128in}{1.986407in}}%
\pgfpathlineto{\pgfqpoint{3.038669in}{1.986407in}}%
\pgfpathlineto{\pgfqpoint{3.038669in}{1.983458in}}%
\pgfpathmoveto{\pgfqpoint{3.038669in}{1.983458in}}%
\pgfpathlineto{\pgfqpoint{3.038669in}{1.983458in}}%
\pgfpathlineto{\pgfqpoint{3.038669in}{1.986407in}}%
\pgfpathlineto{\pgfqpoint{3.043210in}{1.986407in}}%
\pgfpathlineto{\pgfqpoint{3.043210in}{1.983458in}}%
\pgfpathmoveto{\pgfqpoint{3.038669in}{1.986407in}}%
\pgfpathlineto{\pgfqpoint{3.038669in}{1.986407in}}%
\pgfpathlineto{\pgfqpoint{3.038669in}{1.989356in}}%
\pgfpathlineto{\pgfqpoint{3.043210in}{1.989356in}}%
\pgfpathlineto{\pgfqpoint{3.043210in}{1.986407in}}%
\pgfpathmoveto{\pgfqpoint{3.043210in}{1.986407in}}%
\pgfpathlineto{\pgfqpoint{3.043210in}{1.986407in}}%
\pgfpathlineto{\pgfqpoint{3.043210in}{1.989356in}}%
\pgfpathlineto{\pgfqpoint{3.047751in}{1.989356in}}%
\pgfpathlineto{\pgfqpoint{3.047751in}{1.986407in}}%
\pgfpathmoveto{\pgfqpoint{3.043210in}{1.989356in}}%
\pgfpathlineto{\pgfqpoint{3.043210in}{1.989356in}}%
\pgfpathlineto{\pgfqpoint{3.043210in}{1.992305in}}%
\pgfpathlineto{\pgfqpoint{3.047751in}{1.992305in}}%
\pgfpathlineto{\pgfqpoint{3.047751in}{1.989356in}}%
\pgfpathmoveto{\pgfqpoint{3.047751in}{1.989356in}}%
\pgfpathlineto{\pgfqpoint{3.047751in}{1.989356in}}%
\pgfpathlineto{\pgfqpoint{3.047751in}{1.992305in}}%
\pgfpathlineto{\pgfqpoint{3.052292in}{1.992305in}}%
\pgfpathlineto{\pgfqpoint{3.052292in}{1.989356in}}%
\pgfpathmoveto{\pgfqpoint{3.047751in}{1.992305in}}%
\pgfpathlineto{\pgfqpoint{3.047751in}{1.992305in}}%
\pgfpathlineto{\pgfqpoint{3.047751in}{1.995255in}}%
\pgfpathlineto{\pgfqpoint{3.052292in}{1.995255in}}%
\pgfpathlineto{\pgfqpoint{3.052292in}{1.992305in}}%
\pgfpathmoveto{\pgfqpoint{3.052292in}{1.992305in}}%
\pgfpathlineto{\pgfqpoint{3.052292in}{1.992305in}}%
\pgfpathlineto{\pgfqpoint{3.052292in}{1.995255in}}%
\pgfpathlineto{\pgfqpoint{3.056833in}{1.995255in}}%
\pgfpathlineto{\pgfqpoint{3.056833in}{1.992305in}}%
\pgfpathmoveto{\pgfqpoint{3.052292in}{1.995255in}}%
\pgfpathlineto{\pgfqpoint{3.052292in}{1.995255in}}%
\pgfpathlineto{\pgfqpoint{3.052292in}{1.998204in}}%
\pgfpathlineto{\pgfqpoint{3.056833in}{1.998204in}}%
\pgfpathlineto{\pgfqpoint{3.056833in}{1.995255in}}%
\pgfpathmoveto{\pgfqpoint{3.056833in}{1.995255in}}%
\pgfpathlineto{\pgfqpoint{3.056833in}{1.995255in}}%
\pgfpathlineto{\pgfqpoint{3.056833in}{1.998204in}}%
\pgfpathlineto{\pgfqpoint{3.061374in}{1.998204in}}%
\pgfpathlineto{\pgfqpoint{3.061374in}{1.995255in}}%
\pgfpathmoveto{\pgfqpoint{3.056833in}{1.998204in}}%
\pgfpathlineto{\pgfqpoint{3.056833in}{1.998204in}}%
\pgfpathlineto{\pgfqpoint{3.056833in}{2.001153in}}%
\pgfpathlineto{\pgfqpoint{3.061374in}{2.001153in}}%
\pgfpathlineto{\pgfqpoint{3.061374in}{1.998204in}}%
\pgfpathmoveto{\pgfqpoint{3.061374in}{1.998204in}}%
\pgfpathlineto{\pgfqpoint{3.061374in}{1.998204in}}%
\pgfpathlineto{\pgfqpoint{3.061374in}{2.001153in}}%
\pgfpathlineto{\pgfqpoint{3.065915in}{2.001153in}}%
\pgfpathlineto{\pgfqpoint{3.065915in}{1.998204in}}%
\pgfpathmoveto{\pgfqpoint{3.061374in}{2.001153in}}%
\pgfpathlineto{\pgfqpoint{3.061374in}{2.001153in}}%
\pgfpathlineto{\pgfqpoint{3.061374in}{2.004102in}}%
\pgfpathlineto{\pgfqpoint{3.065915in}{2.004102in}}%
\pgfpathlineto{\pgfqpoint{3.065915in}{2.001153in}}%
\pgfpathmoveto{\pgfqpoint{3.065915in}{2.001153in}}%
\pgfpathlineto{\pgfqpoint{3.065915in}{2.001153in}}%
\pgfpathlineto{\pgfqpoint{3.065915in}{2.004102in}}%
\pgfpathlineto{\pgfqpoint{3.070456in}{2.004102in}}%
\pgfpathlineto{\pgfqpoint{3.070456in}{2.001153in}}%
\pgfpathmoveto{\pgfqpoint{3.065915in}{2.004102in}}%
\pgfpathlineto{\pgfqpoint{3.065915in}{2.004102in}}%
\pgfpathlineto{\pgfqpoint{3.065915in}{2.007051in}}%
\pgfpathlineto{\pgfqpoint{3.070456in}{2.007051in}}%
\pgfpathlineto{\pgfqpoint{3.070456in}{2.004102in}}%
\pgfpathmoveto{\pgfqpoint{3.070456in}{2.004102in}}%
\pgfpathlineto{\pgfqpoint{3.070456in}{2.004102in}}%
\pgfpathlineto{\pgfqpoint{3.070456in}{2.007051in}}%
\pgfpathlineto{\pgfqpoint{3.074997in}{2.007051in}}%
\pgfpathlineto{\pgfqpoint{3.074997in}{2.004102in}}%
\pgfpathmoveto{\pgfqpoint{3.070456in}{2.007051in}}%
\pgfpathlineto{\pgfqpoint{3.070456in}{2.007051in}}%
\pgfpathlineto{\pgfqpoint{3.070456in}{2.010000in}}%
\pgfpathlineto{\pgfqpoint{3.074997in}{2.010000in}}%
\pgfpathlineto{\pgfqpoint{3.074997in}{2.007051in}}%
\pgfpathmoveto{\pgfqpoint{3.074997in}{2.007051in}}%
\pgfpathlineto{\pgfqpoint{3.074997in}{2.007051in}}%
\pgfpathlineto{\pgfqpoint{3.074997in}{2.010000in}}%
\pgfpathlineto{\pgfqpoint{3.079538in}{2.010000in}}%
\pgfpathlineto{\pgfqpoint{3.079538in}{2.007051in}}%
\pgfpathmoveto{\pgfqpoint{3.074997in}{2.010000in}}%
\pgfpathlineto{\pgfqpoint{3.074997in}{2.010000in}}%
\pgfpathlineto{\pgfqpoint{3.074997in}{2.012950in}}%
\pgfpathlineto{\pgfqpoint{3.079538in}{2.012950in}}%
\pgfpathlineto{\pgfqpoint{3.079538in}{2.010000in}}%
\pgfpathmoveto{\pgfqpoint{3.079538in}{2.010000in}}%
\pgfpathlineto{\pgfqpoint{3.079538in}{2.010000in}}%
\pgfpathlineto{\pgfqpoint{3.079538in}{2.012950in}}%
\pgfpathlineto{\pgfqpoint{3.084079in}{2.012950in}}%
\pgfpathlineto{\pgfqpoint{3.084079in}{2.010000in}}%
\pgfpathmoveto{\pgfqpoint{3.079538in}{2.012950in}}%
\pgfpathlineto{\pgfqpoint{3.079538in}{2.012950in}}%
\pgfpathlineto{\pgfqpoint{3.079538in}{2.015899in}}%
\pgfpathlineto{\pgfqpoint{3.084079in}{2.015899in}}%
\pgfpathlineto{\pgfqpoint{3.084079in}{2.012950in}}%
\pgfpathmoveto{\pgfqpoint{3.084079in}{2.012950in}}%
\pgfpathlineto{\pgfqpoint{3.084079in}{2.012950in}}%
\pgfpathlineto{\pgfqpoint{3.084079in}{2.015899in}}%
\pgfpathlineto{\pgfqpoint{3.088620in}{2.015899in}}%
\pgfpathlineto{\pgfqpoint{3.088620in}{2.012950in}}%
\pgfpathmoveto{\pgfqpoint{3.084079in}{2.015899in}}%
\pgfpathlineto{\pgfqpoint{3.084079in}{2.015899in}}%
\pgfpathlineto{\pgfqpoint{3.084079in}{2.018848in}}%
\pgfpathlineto{\pgfqpoint{3.088620in}{2.018848in}}%
\pgfpathlineto{\pgfqpoint{3.088620in}{2.015899in}}%
\pgfpathmoveto{\pgfqpoint{3.088620in}{2.015899in}}%
\pgfpathlineto{\pgfqpoint{3.088620in}{2.015899in}}%
\pgfpathlineto{\pgfqpoint{3.088620in}{2.018848in}}%
\pgfpathlineto{\pgfqpoint{3.093161in}{2.018848in}}%
\pgfpathlineto{\pgfqpoint{3.093161in}{2.015899in}}%
\pgfpathmoveto{\pgfqpoint{3.088620in}{2.018848in}}%
\pgfpathlineto{\pgfqpoint{3.088620in}{2.018848in}}%
\pgfpathlineto{\pgfqpoint{3.088620in}{2.021797in}}%
\pgfpathlineto{\pgfqpoint{3.093161in}{2.021797in}}%
\pgfpathlineto{\pgfqpoint{3.093161in}{2.018848in}}%
\pgfpathmoveto{\pgfqpoint{3.093161in}{2.018848in}}%
\pgfpathlineto{\pgfqpoint{3.093161in}{2.018848in}}%
\pgfpathlineto{\pgfqpoint{3.093161in}{2.021797in}}%
\pgfpathlineto{\pgfqpoint{3.097702in}{2.021797in}}%
\pgfpathlineto{\pgfqpoint{3.097702in}{2.018848in}}%
\pgfpathmoveto{\pgfqpoint{3.093161in}{2.021797in}}%
\pgfpathlineto{\pgfqpoint{3.093161in}{2.021797in}}%
\pgfpathlineto{\pgfqpoint{3.093161in}{2.024746in}}%
\pgfpathlineto{\pgfqpoint{3.097702in}{2.024746in}}%
\pgfpathlineto{\pgfqpoint{3.097702in}{2.021797in}}%
\pgfpathmoveto{\pgfqpoint{3.097702in}{2.021797in}}%
\pgfpathlineto{\pgfqpoint{3.097702in}{2.021797in}}%
\pgfpathlineto{\pgfqpoint{3.097702in}{2.024746in}}%
\pgfpathlineto{\pgfqpoint{3.102244in}{2.024746in}}%
\pgfpathlineto{\pgfqpoint{3.102244in}{2.021797in}}%
\pgfpathmoveto{\pgfqpoint{3.097702in}{2.024746in}}%
\pgfpathlineto{\pgfqpoint{3.097702in}{2.024746in}}%
\pgfpathlineto{\pgfqpoint{3.097702in}{2.027696in}}%
\pgfpathlineto{\pgfqpoint{3.102244in}{2.027696in}}%
\pgfpathlineto{\pgfqpoint{3.102244in}{2.024746in}}%
\pgfpathmoveto{\pgfqpoint{3.102244in}{2.024746in}}%
\pgfpathlineto{\pgfqpoint{3.102244in}{2.024746in}}%
\pgfpathlineto{\pgfqpoint{3.102244in}{2.027696in}}%
\pgfpathlineto{\pgfqpoint{3.106785in}{2.027696in}}%
\pgfpathlineto{\pgfqpoint{3.106785in}{2.024746in}}%
\pgfpathmoveto{\pgfqpoint{3.102244in}{2.027696in}}%
\pgfpathlineto{\pgfqpoint{3.102244in}{2.027696in}}%
\pgfpathlineto{\pgfqpoint{3.102244in}{2.030645in}}%
\pgfpathlineto{\pgfqpoint{3.106785in}{2.030645in}}%
\pgfpathlineto{\pgfqpoint{3.106785in}{2.027696in}}%
\pgfpathmoveto{\pgfqpoint{3.106785in}{2.027696in}}%
\pgfpathlineto{\pgfqpoint{3.106785in}{2.027696in}}%
\pgfpathlineto{\pgfqpoint{3.106785in}{2.030645in}}%
\pgfpathlineto{\pgfqpoint{3.111326in}{2.030645in}}%
\pgfpathlineto{\pgfqpoint{3.111326in}{2.027696in}}%
\pgfpathmoveto{\pgfqpoint{3.106785in}{2.030645in}}%
\pgfpathlineto{\pgfqpoint{3.106785in}{2.030645in}}%
\pgfpathlineto{\pgfqpoint{3.106785in}{2.033594in}}%
\pgfpathlineto{\pgfqpoint{3.111326in}{2.033594in}}%
\pgfpathlineto{\pgfqpoint{3.111326in}{2.030645in}}%
\pgfpathmoveto{\pgfqpoint{3.111326in}{2.030645in}}%
\pgfpathlineto{\pgfqpoint{3.111326in}{2.030645in}}%
\pgfpathlineto{\pgfqpoint{3.111326in}{2.033594in}}%
\pgfpathlineto{\pgfqpoint{3.115867in}{2.033594in}}%
\pgfpathlineto{\pgfqpoint{3.115867in}{2.030645in}}%
\pgfpathmoveto{\pgfqpoint{3.111326in}{2.033594in}}%
\pgfpathlineto{\pgfqpoint{3.111326in}{2.033594in}}%
\pgfpathlineto{\pgfqpoint{3.111326in}{2.036543in}}%
\pgfpathlineto{\pgfqpoint{3.115867in}{2.036543in}}%
\pgfpathlineto{\pgfqpoint{3.115867in}{2.033594in}}%
\pgfpathmoveto{\pgfqpoint{3.115867in}{2.033594in}}%
\pgfpathlineto{\pgfqpoint{3.115867in}{2.033594in}}%
\pgfpathlineto{\pgfqpoint{3.115867in}{2.036543in}}%
\pgfpathlineto{\pgfqpoint{3.120408in}{2.036543in}}%
\pgfpathlineto{\pgfqpoint{3.120408in}{2.033594in}}%
\pgfpathmoveto{\pgfqpoint{3.115867in}{2.036543in}}%
\pgfpathlineto{\pgfqpoint{3.115867in}{2.036543in}}%
\pgfpathlineto{\pgfqpoint{3.115867in}{2.039492in}}%
\pgfpathlineto{\pgfqpoint{3.120408in}{2.039492in}}%
\pgfpathlineto{\pgfqpoint{3.120408in}{2.036543in}}%
\pgfpathmoveto{\pgfqpoint{3.120408in}{2.036543in}}%
\pgfpathlineto{\pgfqpoint{3.120408in}{2.036543in}}%
\pgfpathlineto{\pgfqpoint{3.120408in}{2.039492in}}%
\pgfpathlineto{\pgfqpoint{3.124949in}{2.039492in}}%
\pgfpathlineto{\pgfqpoint{3.124949in}{2.036543in}}%
\pgfpathmoveto{\pgfqpoint{3.120408in}{2.039492in}}%
\pgfpathlineto{\pgfqpoint{3.120408in}{2.039492in}}%
\pgfpathlineto{\pgfqpoint{3.120408in}{2.042442in}}%
\pgfpathlineto{\pgfqpoint{3.124949in}{2.042442in}}%
\pgfpathlineto{\pgfqpoint{3.124949in}{2.039492in}}%
\pgfpathmoveto{\pgfqpoint{3.124949in}{2.039492in}}%
\pgfpathlineto{\pgfqpoint{3.124949in}{2.039492in}}%
\pgfpathlineto{\pgfqpoint{3.124949in}{2.042442in}}%
\pgfpathlineto{\pgfqpoint{3.129490in}{2.042442in}}%
\pgfpathlineto{\pgfqpoint{3.129490in}{2.039492in}}%
\pgfpathmoveto{\pgfqpoint{3.124949in}{2.042442in}}%
\pgfpathlineto{\pgfqpoint{3.124949in}{2.042442in}}%
\pgfpathlineto{\pgfqpoint{3.124949in}{2.045391in}}%
\pgfpathlineto{\pgfqpoint{3.129490in}{2.045391in}}%
\pgfpathlineto{\pgfqpoint{3.129490in}{2.042442in}}%
\pgfpathmoveto{\pgfqpoint{3.129490in}{2.042442in}}%
\pgfpathlineto{\pgfqpoint{3.129490in}{2.042442in}}%
\pgfpathlineto{\pgfqpoint{3.129490in}{2.045391in}}%
\pgfpathlineto{\pgfqpoint{3.134031in}{2.045391in}}%
\pgfpathlineto{\pgfqpoint{3.134031in}{2.042442in}}%
\pgfpathmoveto{\pgfqpoint{3.129490in}{2.045391in}}%
\pgfpathlineto{\pgfqpoint{3.129490in}{2.045391in}}%
\pgfpathlineto{\pgfqpoint{3.129490in}{2.048340in}}%
\pgfpathlineto{\pgfqpoint{3.134031in}{2.048340in}}%
\pgfpathlineto{\pgfqpoint{3.134031in}{2.045391in}}%
\pgfpathmoveto{\pgfqpoint{3.134031in}{2.045391in}}%
\pgfpathlineto{\pgfqpoint{3.134031in}{2.045391in}}%
\pgfpathlineto{\pgfqpoint{3.134031in}{2.048340in}}%
\pgfpathlineto{\pgfqpoint{3.138572in}{2.048340in}}%
\pgfpathlineto{\pgfqpoint{3.138572in}{2.045391in}}%
\pgfpathmoveto{\pgfqpoint{3.134031in}{2.048340in}}%
\pgfpathlineto{\pgfqpoint{3.134031in}{2.048340in}}%
\pgfpathlineto{\pgfqpoint{3.134031in}{2.051289in}}%
\pgfpathlineto{\pgfqpoint{3.138572in}{2.051289in}}%
\pgfpathlineto{\pgfqpoint{3.138572in}{2.048340in}}%
\pgfpathmoveto{\pgfqpoint{3.138572in}{2.048340in}}%
\pgfpathlineto{\pgfqpoint{3.138572in}{2.048340in}}%
\pgfpathlineto{\pgfqpoint{3.138572in}{2.051289in}}%
\pgfpathlineto{\pgfqpoint{3.143113in}{2.051289in}}%
\pgfpathlineto{\pgfqpoint{3.143113in}{2.048340in}}%
\pgfpathmoveto{\pgfqpoint{3.138572in}{2.051289in}}%
\pgfpathlineto{\pgfqpoint{3.138572in}{2.051289in}}%
\pgfpathlineto{\pgfqpoint{3.138572in}{2.054238in}}%
\pgfpathlineto{\pgfqpoint{3.143113in}{2.054238in}}%
\pgfpathlineto{\pgfqpoint{3.143113in}{2.051289in}}%
\pgfpathmoveto{\pgfqpoint{3.143113in}{2.051289in}}%
\pgfpathlineto{\pgfqpoint{3.143113in}{2.051289in}}%
\pgfpathlineto{\pgfqpoint{3.143113in}{2.054238in}}%
\pgfpathlineto{\pgfqpoint{3.147654in}{2.054238in}}%
\pgfpathlineto{\pgfqpoint{3.147654in}{2.051289in}}%
\pgfpathmoveto{\pgfqpoint{3.143113in}{2.054238in}}%
\pgfpathlineto{\pgfqpoint{3.143113in}{2.054238in}}%
\pgfpathlineto{\pgfqpoint{3.143113in}{2.057188in}}%
\pgfpathlineto{\pgfqpoint{3.147654in}{2.057188in}}%
\pgfpathlineto{\pgfqpoint{3.147654in}{2.054238in}}%
\pgfpathmoveto{\pgfqpoint{3.147654in}{2.054238in}}%
\pgfpathlineto{\pgfqpoint{3.147654in}{2.054238in}}%
\pgfpathlineto{\pgfqpoint{3.147654in}{2.057188in}}%
\pgfpathlineto{\pgfqpoint{3.152195in}{2.057188in}}%
\pgfpathlineto{\pgfqpoint{3.152195in}{2.054238in}}%
\pgfpathmoveto{\pgfqpoint{3.147654in}{2.057188in}}%
\pgfpathlineto{\pgfqpoint{3.147654in}{2.057188in}}%
\pgfpathlineto{\pgfqpoint{3.147654in}{2.060137in}}%
\pgfpathlineto{\pgfqpoint{3.152195in}{2.060137in}}%
\pgfpathlineto{\pgfqpoint{3.152195in}{2.057188in}}%
\pgfpathmoveto{\pgfqpoint{3.152195in}{2.057188in}}%
\pgfpathlineto{\pgfqpoint{3.152195in}{2.057188in}}%
\pgfpathlineto{\pgfqpoint{3.152195in}{2.060137in}}%
\pgfpathlineto{\pgfqpoint{3.156736in}{2.060137in}}%
\pgfpathlineto{\pgfqpoint{3.156736in}{2.057188in}}%
\pgfpathmoveto{\pgfqpoint{3.152195in}{2.060137in}}%
\pgfpathlineto{\pgfqpoint{3.152195in}{2.060137in}}%
\pgfpathlineto{\pgfqpoint{3.152195in}{2.063086in}}%
\pgfpathlineto{\pgfqpoint{3.156736in}{2.063086in}}%
\pgfpathlineto{\pgfqpoint{3.156736in}{2.060137in}}%
\pgfpathmoveto{\pgfqpoint{3.156736in}{2.060137in}}%
\pgfpathlineto{\pgfqpoint{3.156736in}{2.060137in}}%
\pgfpathlineto{\pgfqpoint{3.156736in}{2.063086in}}%
\pgfpathlineto{\pgfqpoint{3.161277in}{2.063086in}}%
\pgfpathlineto{\pgfqpoint{3.161277in}{2.060137in}}%
\pgfpathmoveto{\pgfqpoint{3.156736in}{2.063086in}}%
\pgfpathlineto{\pgfqpoint{3.156736in}{2.063086in}}%
\pgfpathlineto{\pgfqpoint{3.156736in}{2.066035in}}%
\pgfpathlineto{\pgfqpoint{3.161277in}{2.066035in}}%
\pgfpathlineto{\pgfqpoint{3.161277in}{2.063086in}}%
\pgfpathmoveto{\pgfqpoint{3.161277in}{2.063086in}}%
\pgfpathlineto{\pgfqpoint{3.161277in}{2.063086in}}%
\pgfpathlineto{\pgfqpoint{3.161277in}{2.066035in}}%
\pgfpathlineto{\pgfqpoint{3.165818in}{2.066035in}}%
\pgfpathlineto{\pgfqpoint{3.165818in}{2.063086in}}%
\pgfpathmoveto{\pgfqpoint{3.161277in}{2.066035in}}%
\pgfpathlineto{\pgfqpoint{3.161277in}{2.066035in}}%
\pgfpathlineto{\pgfqpoint{3.161277in}{2.068984in}}%
\pgfpathlineto{\pgfqpoint{3.165818in}{2.068984in}}%
\pgfpathlineto{\pgfqpoint{3.165818in}{2.066035in}}%
\pgfpathmoveto{\pgfqpoint{3.165818in}{2.066035in}}%
\pgfpathlineto{\pgfqpoint{3.165818in}{2.066035in}}%
\pgfpathlineto{\pgfqpoint{3.165818in}{2.068984in}}%
\pgfpathlineto{\pgfqpoint{3.170359in}{2.068984in}}%
\pgfpathlineto{\pgfqpoint{3.170359in}{2.066035in}}%
\pgfpathmoveto{\pgfqpoint{3.165818in}{2.068984in}}%
\pgfpathlineto{\pgfqpoint{3.165818in}{2.068984in}}%
\pgfpathlineto{\pgfqpoint{3.165818in}{2.071934in}}%
\pgfpathlineto{\pgfqpoint{3.170359in}{2.071934in}}%
\pgfpathlineto{\pgfqpoint{3.170359in}{2.068984in}}%
\pgfpathmoveto{\pgfqpoint{3.170359in}{2.068984in}}%
\pgfpathlineto{\pgfqpoint{3.170359in}{2.068984in}}%
\pgfpathlineto{\pgfqpoint{3.170359in}{2.071934in}}%
\pgfpathlineto{\pgfqpoint{3.174900in}{2.071934in}}%
\pgfpathlineto{\pgfqpoint{3.174900in}{2.068984in}}%
\pgfpathmoveto{\pgfqpoint{3.170359in}{2.071934in}}%
\pgfpathlineto{\pgfqpoint{3.170359in}{2.071934in}}%
\pgfpathlineto{\pgfqpoint{3.170359in}{2.074883in}}%
\pgfpathlineto{\pgfqpoint{3.174900in}{2.074883in}}%
\pgfpathlineto{\pgfqpoint{3.174900in}{2.071934in}}%
\pgfpathmoveto{\pgfqpoint{3.174900in}{2.071934in}}%
\pgfpathlineto{\pgfqpoint{3.174900in}{2.071934in}}%
\pgfpathlineto{\pgfqpoint{3.174900in}{2.074883in}}%
\pgfpathlineto{\pgfqpoint{3.179441in}{2.074883in}}%
\pgfpathlineto{\pgfqpoint{3.179441in}{2.071934in}}%
\pgfpathmoveto{\pgfqpoint{3.174900in}{2.074883in}}%
\pgfpathlineto{\pgfqpoint{3.174900in}{2.074883in}}%
\pgfpathlineto{\pgfqpoint{3.174900in}{2.077832in}}%
\pgfpathlineto{\pgfqpoint{3.179441in}{2.077832in}}%
\pgfpathlineto{\pgfqpoint{3.179441in}{2.074883in}}%
\pgfpathmoveto{\pgfqpoint{3.179441in}{2.074883in}}%
\pgfpathlineto{\pgfqpoint{3.179441in}{2.074883in}}%
\pgfpathlineto{\pgfqpoint{3.179441in}{2.077832in}}%
\pgfpathlineto{\pgfqpoint{3.183982in}{2.077832in}}%
\pgfpathlineto{\pgfqpoint{3.183982in}{2.074883in}}%
\pgfpathmoveto{\pgfqpoint{3.179441in}{2.077832in}}%
\pgfpathlineto{\pgfqpoint{3.179441in}{2.077832in}}%
\pgfpathlineto{\pgfqpoint{3.179441in}{2.080781in}}%
\pgfpathlineto{\pgfqpoint{3.183982in}{2.080781in}}%
\pgfpathlineto{\pgfqpoint{3.183982in}{2.077832in}}%
\pgfpathmoveto{\pgfqpoint{3.183982in}{2.077832in}}%
\pgfpathlineto{\pgfqpoint{3.183982in}{2.077832in}}%
\pgfpathlineto{\pgfqpoint{3.183982in}{2.080781in}}%
\pgfpathlineto{\pgfqpoint{3.188523in}{2.080781in}}%
\pgfpathlineto{\pgfqpoint{3.188523in}{2.077832in}}%
\pgfpathmoveto{\pgfqpoint{3.183982in}{2.080781in}}%
\pgfpathlineto{\pgfqpoint{3.183982in}{2.080781in}}%
\pgfpathlineto{\pgfqpoint{3.183982in}{2.083730in}}%
\pgfpathlineto{\pgfqpoint{3.188523in}{2.083730in}}%
\pgfpathlineto{\pgfqpoint{3.188523in}{2.080781in}}%
\pgfpathmoveto{\pgfqpoint{3.188523in}{2.080781in}}%
\pgfpathlineto{\pgfqpoint{3.188523in}{2.080781in}}%
\pgfpathlineto{\pgfqpoint{3.188523in}{2.083730in}}%
\pgfpathlineto{\pgfqpoint{3.193064in}{2.083730in}}%
\pgfpathlineto{\pgfqpoint{3.193064in}{2.080781in}}%
\pgfpathmoveto{\pgfqpoint{3.188523in}{2.083730in}}%
\pgfpathlineto{\pgfqpoint{3.188523in}{2.083730in}}%
\pgfpathlineto{\pgfqpoint{3.188523in}{2.086680in}}%
\pgfpathlineto{\pgfqpoint{3.193064in}{2.086680in}}%
\pgfpathlineto{\pgfqpoint{3.193064in}{2.083730in}}%
\pgfpathmoveto{\pgfqpoint{3.193064in}{2.083730in}}%
\pgfpathlineto{\pgfqpoint{3.193064in}{2.083730in}}%
\pgfpathlineto{\pgfqpoint{3.193064in}{2.086680in}}%
\pgfpathlineto{\pgfqpoint{3.197605in}{2.086680in}}%
\pgfpathlineto{\pgfqpoint{3.197605in}{2.083730in}}%
\pgfpathmoveto{\pgfqpoint{3.193064in}{2.086680in}}%
\pgfpathlineto{\pgfqpoint{3.193064in}{2.086680in}}%
\pgfpathlineto{\pgfqpoint{3.193064in}{2.089629in}}%
\pgfpathlineto{\pgfqpoint{3.197605in}{2.089629in}}%
\pgfpathlineto{\pgfqpoint{3.197605in}{2.086680in}}%
\pgfpathmoveto{\pgfqpoint{3.197605in}{2.086680in}}%
\pgfpathlineto{\pgfqpoint{3.197605in}{2.086680in}}%
\pgfpathlineto{\pgfqpoint{3.197605in}{2.089629in}}%
\pgfpathlineto{\pgfqpoint{3.202147in}{2.089629in}}%
\pgfpathlineto{\pgfqpoint{3.202147in}{2.086680in}}%
\pgfpathmoveto{\pgfqpoint{3.197605in}{2.089629in}}%
\pgfpathlineto{\pgfqpoint{3.197605in}{2.089629in}}%
\pgfpathlineto{\pgfqpoint{3.197605in}{2.092578in}}%
\pgfpathlineto{\pgfqpoint{3.202147in}{2.092578in}}%
\pgfpathlineto{\pgfqpoint{3.202147in}{2.089629in}}%
\pgfpathmoveto{\pgfqpoint{3.202147in}{2.089629in}}%
\pgfpathlineto{\pgfqpoint{3.202147in}{2.089629in}}%
\pgfpathlineto{\pgfqpoint{3.202147in}{2.092578in}}%
\pgfpathlineto{\pgfqpoint{3.206688in}{2.092578in}}%
\pgfpathlineto{\pgfqpoint{3.206688in}{2.089629in}}%
\pgfpathmoveto{\pgfqpoint{3.202147in}{2.092578in}}%
\pgfpathlineto{\pgfqpoint{3.202147in}{2.092578in}}%
\pgfpathlineto{\pgfqpoint{3.202147in}{2.095527in}}%
\pgfpathlineto{\pgfqpoint{3.206688in}{2.095527in}}%
\pgfpathlineto{\pgfqpoint{3.206688in}{2.092578in}}%
\pgfpathmoveto{\pgfqpoint{3.206688in}{2.092578in}}%
\pgfpathlineto{\pgfqpoint{3.206688in}{2.092578in}}%
\pgfpathlineto{\pgfqpoint{3.206688in}{2.095527in}}%
\pgfpathlineto{\pgfqpoint{3.211229in}{2.095527in}}%
\pgfpathlineto{\pgfqpoint{3.211229in}{2.092578in}}%
\pgfpathmoveto{\pgfqpoint{3.206688in}{2.095527in}}%
\pgfpathlineto{\pgfqpoint{3.206688in}{2.095527in}}%
\pgfpathlineto{\pgfqpoint{3.206688in}{2.098477in}}%
\pgfpathlineto{\pgfqpoint{3.211229in}{2.098477in}}%
\pgfpathlineto{\pgfqpoint{3.211229in}{2.095527in}}%
\pgfpathmoveto{\pgfqpoint{3.211229in}{2.095527in}}%
\pgfpathlineto{\pgfqpoint{3.211229in}{2.095527in}}%
\pgfpathlineto{\pgfqpoint{3.211229in}{2.098477in}}%
\pgfpathlineto{\pgfqpoint{3.215770in}{2.098477in}}%
\pgfpathlineto{\pgfqpoint{3.215770in}{2.095527in}}%
\pgfpathmoveto{\pgfqpoint{3.211229in}{2.098477in}}%
\pgfpathlineto{\pgfqpoint{3.211229in}{2.098477in}}%
\pgfpathlineto{\pgfqpoint{3.211229in}{2.101426in}}%
\pgfpathlineto{\pgfqpoint{3.215770in}{2.101426in}}%
\pgfpathlineto{\pgfqpoint{3.215770in}{2.098477in}}%
\pgfpathmoveto{\pgfqpoint{3.215770in}{2.098477in}}%
\pgfpathlineto{\pgfqpoint{3.215770in}{2.098477in}}%
\pgfpathlineto{\pgfqpoint{3.215770in}{2.101426in}}%
\pgfpathlineto{\pgfqpoint{3.220311in}{2.101426in}}%
\pgfpathlineto{\pgfqpoint{3.220311in}{2.098477in}}%
\pgfpathmoveto{\pgfqpoint{3.215770in}{2.101426in}}%
\pgfpathlineto{\pgfqpoint{3.215770in}{2.101426in}}%
\pgfpathlineto{\pgfqpoint{3.215770in}{2.104375in}}%
\pgfpathlineto{\pgfqpoint{3.220311in}{2.104375in}}%
\pgfpathlineto{\pgfqpoint{3.220311in}{2.101426in}}%
\pgfpathmoveto{\pgfqpoint{3.220311in}{2.101426in}}%
\pgfpathlineto{\pgfqpoint{3.220311in}{2.101426in}}%
\pgfpathlineto{\pgfqpoint{3.220311in}{2.104375in}}%
\pgfpathlineto{\pgfqpoint{3.224852in}{2.104375in}}%
\pgfpathlineto{\pgfqpoint{3.224852in}{2.101426in}}%
\pgfpathmoveto{\pgfqpoint{3.220311in}{2.104375in}}%
\pgfpathlineto{\pgfqpoint{3.220311in}{2.104375in}}%
\pgfpathlineto{\pgfqpoint{3.220311in}{2.107324in}}%
\pgfpathlineto{\pgfqpoint{3.224852in}{2.107324in}}%
\pgfpathlineto{\pgfqpoint{3.224852in}{2.104375in}}%
\pgfpathmoveto{\pgfqpoint{3.224852in}{2.104375in}}%
\pgfpathlineto{\pgfqpoint{3.224852in}{2.104375in}}%
\pgfpathlineto{\pgfqpoint{3.224852in}{2.107324in}}%
\pgfpathlineto{\pgfqpoint{3.229393in}{2.107324in}}%
\pgfpathlineto{\pgfqpoint{3.229393in}{2.104375in}}%
\pgfpathmoveto{\pgfqpoint{3.224852in}{2.107324in}}%
\pgfpathlineto{\pgfqpoint{3.224852in}{2.107324in}}%
\pgfpathlineto{\pgfqpoint{3.224852in}{2.110273in}}%
\pgfpathlineto{\pgfqpoint{3.229393in}{2.110273in}}%
\pgfpathlineto{\pgfqpoint{3.229393in}{2.107324in}}%
\pgfpathmoveto{\pgfqpoint{3.229393in}{2.107324in}}%
\pgfpathlineto{\pgfqpoint{3.229393in}{2.107324in}}%
\pgfpathlineto{\pgfqpoint{3.229393in}{2.110273in}}%
\pgfpathlineto{\pgfqpoint{3.233934in}{2.110273in}}%
\pgfpathlineto{\pgfqpoint{3.233934in}{2.107324in}}%
\pgfpathmoveto{\pgfqpoint{3.229393in}{2.110273in}}%
\pgfpathlineto{\pgfqpoint{3.229393in}{2.110273in}}%
\pgfpathlineto{\pgfqpoint{3.229393in}{2.113223in}}%
\pgfpathlineto{\pgfqpoint{3.233934in}{2.113223in}}%
\pgfpathlineto{\pgfqpoint{3.233934in}{2.110273in}}%
\pgfpathmoveto{\pgfqpoint{3.233934in}{2.110273in}}%
\pgfpathlineto{\pgfqpoint{3.233934in}{2.110273in}}%
\pgfpathlineto{\pgfqpoint{3.233934in}{2.113223in}}%
\pgfpathlineto{\pgfqpoint{3.238475in}{2.113223in}}%
\pgfpathlineto{\pgfqpoint{3.238475in}{2.110273in}}%
\pgfpathmoveto{\pgfqpoint{3.233934in}{2.113223in}}%
\pgfpathlineto{\pgfqpoint{3.233934in}{2.113223in}}%
\pgfpathlineto{\pgfqpoint{3.233934in}{2.116172in}}%
\pgfpathlineto{\pgfqpoint{3.238475in}{2.116172in}}%
\pgfpathlineto{\pgfqpoint{3.238475in}{2.113223in}}%
\pgfpathmoveto{\pgfqpoint{3.238475in}{2.113223in}}%
\pgfpathlineto{\pgfqpoint{3.238475in}{2.113223in}}%
\pgfpathlineto{\pgfqpoint{3.238475in}{2.116172in}}%
\pgfpathlineto{\pgfqpoint{3.243016in}{2.116172in}}%
\pgfpathlineto{\pgfqpoint{3.243016in}{2.113223in}}%
\pgfpathmoveto{\pgfqpoint{3.238475in}{2.116172in}}%
\pgfpathlineto{\pgfqpoint{3.238475in}{2.116172in}}%
\pgfpathlineto{\pgfqpoint{3.238475in}{2.119121in}}%
\pgfpathlineto{\pgfqpoint{3.243016in}{2.119121in}}%
\pgfpathlineto{\pgfqpoint{3.243016in}{2.116172in}}%
\pgfpathmoveto{\pgfqpoint{3.243016in}{2.116172in}}%
\pgfpathlineto{\pgfqpoint{3.243016in}{2.116172in}}%
\pgfpathlineto{\pgfqpoint{3.243016in}{2.119121in}}%
\pgfpathlineto{\pgfqpoint{3.247557in}{2.119121in}}%
\pgfpathlineto{\pgfqpoint{3.247557in}{2.116172in}}%
\pgfpathmoveto{\pgfqpoint{3.243016in}{2.119121in}}%
\pgfpathlineto{\pgfqpoint{3.243016in}{2.119121in}}%
\pgfpathlineto{\pgfqpoint{3.243016in}{2.122070in}}%
\pgfpathlineto{\pgfqpoint{3.247557in}{2.122070in}}%
\pgfpathlineto{\pgfqpoint{3.247557in}{2.119121in}}%
\pgfpathmoveto{\pgfqpoint{3.247557in}{2.119121in}}%
\pgfpathlineto{\pgfqpoint{3.247557in}{2.119121in}}%
\pgfpathlineto{\pgfqpoint{3.247557in}{2.122070in}}%
\pgfpathlineto{\pgfqpoint{3.252098in}{2.122070in}}%
\pgfpathlineto{\pgfqpoint{3.252098in}{2.119121in}}%
\pgfpathmoveto{\pgfqpoint{3.247557in}{2.122070in}}%
\pgfpathlineto{\pgfqpoint{3.247557in}{2.122070in}}%
\pgfpathlineto{\pgfqpoint{3.247557in}{2.125020in}}%
\pgfpathlineto{\pgfqpoint{3.252098in}{2.125020in}}%
\pgfpathlineto{\pgfqpoint{3.252098in}{2.122070in}}%
\pgfpathmoveto{\pgfqpoint{3.252098in}{2.122070in}}%
\pgfpathlineto{\pgfqpoint{3.252098in}{2.122070in}}%
\pgfpathlineto{\pgfqpoint{3.252098in}{2.125020in}}%
\pgfpathlineto{\pgfqpoint{3.256639in}{2.125020in}}%
\pgfpathlineto{\pgfqpoint{3.256639in}{2.122070in}}%
\pgfpathmoveto{\pgfqpoint{3.252098in}{2.125020in}}%
\pgfpathlineto{\pgfqpoint{3.252098in}{2.125020in}}%
\pgfpathlineto{\pgfqpoint{3.252098in}{2.127969in}}%
\pgfpathlineto{\pgfqpoint{3.256639in}{2.127969in}}%
\pgfpathlineto{\pgfqpoint{3.256639in}{2.125020in}}%
\pgfpathmoveto{\pgfqpoint{3.256639in}{2.125020in}}%
\pgfpathlineto{\pgfqpoint{3.256639in}{2.125020in}}%
\pgfpathlineto{\pgfqpoint{3.256639in}{2.127969in}}%
\pgfpathlineto{\pgfqpoint{3.261180in}{2.127969in}}%
\pgfpathlineto{\pgfqpoint{3.261180in}{2.125020in}}%
\pgfpathmoveto{\pgfqpoint{3.256639in}{2.127969in}}%
\pgfpathlineto{\pgfqpoint{3.256639in}{2.127969in}}%
\pgfpathlineto{\pgfqpoint{3.256639in}{2.130918in}}%
\pgfpathlineto{\pgfqpoint{3.261180in}{2.130918in}}%
\pgfpathlineto{\pgfqpoint{3.261180in}{2.127969in}}%
\pgfpathmoveto{\pgfqpoint{3.261180in}{2.127969in}}%
\pgfpathlineto{\pgfqpoint{3.261180in}{2.127969in}}%
\pgfpathlineto{\pgfqpoint{3.261180in}{2.130918in}}%
\pgfpathlineto{\pgfqpoint{3.265721in}{2.130918in}}%
\pgfpathlineto{\pgfqpoint{3.265721in}{2.127969in}}%
\pgfpathmoveto{\pgfqpoint{3.261180in}{2.130918in}}%
\pgfpathlineto{\pgfqpoint{3.261180in}{2.130918in}}%
\pgfpathlineto{\pgfqpoint{3.261180in}{2.133867in}}%
\pgfpathlineto{\pgfqpoint{3.265721in}{2.133867in}}%
\pgfpathlineto{\pgfqpoint{3.265721in}{2.130918in}}%
\pgfpathmoveto{\pgfqpoint{3.265721in}{2.130918in}}%
\pgfpathlineto{\pgfqpoint{3.265721in}{2.130918in}}%
\pgfpathlineto{\pgfqpoint{3.265721in}{2.133867in}}%
\pgfpathlineto{\pgfqpoint{3.270262in}{2.133867in}}%
\pgfpathlineto{\pgfqpoint{3.270262in}{2.130918in}}%
\pgfpathmoveto{\pgfqpoint{3.265721in}{2.133867in}}%
\pgfpathlineto{\pgfqpoint{3.265721in}{2.133867in}}%
\pgfpathlineto{\pgfqpoint{3.265721in}{2.136817in}}%
\pgfpathlineto{\pgfqpoint{3.270262in}{2.136817in}}%
\pgfpathlineto{\pgfqpoint{3.270262in}{2.133867in}}%
\pgfpathmoveto{\pgfqpoint{3.270262in}{2.133867in}}%
\pgfpathlineto{\pgfqpoint{3.270262in}{2.133867in}}%
\pgfpathlineto{\pgfqpoint{3.270262in}{2.136817in}}%
\pgfpathlineto{\pgfqpoint{3.274803in}{2.136817in}}%
\pgfpathlineto{\pgfqpoint{3.274803in}{2.133867in}}%
\pgfpathmoveto{\pgfqpoint{3.270262in}{2.136817in}}%
\pgfpathlineto{\pgfqpoint{3.270262in}{2.136817in}}%
\pgfpathlineto{\pgfqpoint{3.270262in}{2.139766in}}%
\pgfpathlineto{\pgfqpoint{3.274803in}{2.139766in}}%
\pgfpathlineto{\pgfqpoint{3.274803in}{2.136817in}}%
\pgfpathmoveto{\pgfqpoint{3.274803in}{2.136817in}}%
\pgfpathlineto{\pgfqpoint{3.274803in}{2.136817in}}%
\pgfpathlineto{\pgfqpoint{3.274803in}{2.139766in}}%
\pgfpathlineto{\pgfqpoint{3.279344in}{2.139766in}}%
\pgfpathlineto{\pgfqpoint{3.279344in}{2.136817in}}%
\pgfpathmoveto{\pgfqpoint{3.274803in}{2.139766in}}%
\pgfpathlineto{\pgfqpoint{3.274803in}{2.139766in}}%
\pgfpathlineto{\pgfqpoint{3.274803in}{2.142715in}}%
\pgfpathlineto{\pgfqpoint{3.279344in}{2.142715in}}%
\pgfpathlineto{\pgfqpoint{3.279344in}{2.139766in}}%
\pgfpathmoveto{\pgfqpoint{3.279344in}{2.139766in}}%
\pgfpathlineto{\pgfqpoint{3.279344in}{2.139766in}}%
\pgfpathlineto{\pgfqpoint{3.279344in}{2.142715in}}%
\pgfpathlineto{\pgfqpoint{3.283885in}{2.142715in}}%
\pgfpathlineto{\pgfqpoint{3.283885in}{2.139766in}}%
\pgfpathmoveto{\pgfqpoint{3.279344in}{2.142715in}}%
\pgfpathlineto{\pgfqpoint{3.279344in}{2.142715in}}%
\pgfpathlineto{\pgfqpoint{3.279344in}{2.145664in}}%
\pgfpathlineto{\pgfqpoint{3.283885in}{2.145664in}}%
\pgfpathlineto{\pgfqpoint{3.283885in}{2.142715in}}%
\pgfpathmoveto{\pgfqpoint{3.283885in}{2.142715in}}%
\pgfpathlineto{\pgfqpoint{3.283885in}{2.142715in}}%
\pgfpathlineto{\pgfqpoint{3.283885in}{2.145664in}}%
\pgfpathlineto{\pgfqpoint{3.288426in}{2.145664in}}%
\pgfpathlineto{\pgfqpoint{3.288426in}{2.142715in}}%
\pgfpathmoveto{\pgfqpoint{3.283885in}{2.145664in}}%
\pgfpathlineto{\pgfqpoint{3.283885in}{2.145664in}}%
\pgfpathlineto{\pgfqpoint{3.283885in}{2.148614in}}%
\pgfpathlineto{\pgfqpoint{3.288426in}{2.148614in}}%
\pgfpathlineto{\pgfqpoint{3.288426in}{2.145664in}}%
\pgfpathmoveto{\pgfqpoint{3.288426in}{2.145664in}}%
\pgfpathlineto{\pgfqpoint{3.288426in}{2.145664in}}%
\pgfpathlineto{\pgfqpoint{3.288426in}{2.148614in}}%
\pgfpathlineto{\pgfqpoint{3.292967in}{2.148614in}}%
\pgfpathlineto{\pgfqpoint{3.292967in}{2.145664in}}%
\pgfpathmoveto{\pgfqpoint{3.288426in}{2.148614in}}%
\pgfpathlineto{\pgfqpoint{3.288426in}{2.148614in}}%
\pgfpathlineto{\pgfqpoint{3.288426in}{2.151563in}}%
\pgfpathlineto{\pgfqpoint{3.292967in}{2.151563in}}%
\pgfpathlineto{\pgfqpoint{3.292967in}{2.148614in}}%
\pgfpathmoveto{\pgfqpoint{3.292967in}{2.148614in}}%
\pgfpathlineto{\pgfqpoint{3.292967in}{2.148614in}}%
\pgfpathlineto{\pgfqpoint{3.292967in}{2.151563in}}%
\pgfpathlineto{\pgfqpoint{3.297509in}{2.151563in}}%
\pgfpathlineto{\pgfqpoint{3.297509in}{2.148614in}}%
\pgfpathmoveto{\pgfqpoint{3.292967in}{2.151563in}}%
\pgfpathlineto{\pgfqpoint{3.292967in}{2.151563in}}%
\pgfpathlineto{\pgfqpoint{3.292967in}{2.154512in}}%
\pgfpathlineto{\pgfqpoint{3.297509in}{2.154512in}}%
\pgfpathlineto{\pgfqpoint{3.297509in}{2.151563in}}%
\pgfpathmoveto{\pgfqpoint{3.297509in}{2.151563in}}%
\pgfpathlineto{\pgfqpoint{3.297509in}{2.151563in}}%
\pgfpathlineto{\pgfqpoint{3.297509in}{2.154512in}}%
\pgfpathlineto{\pgfqpoint{3.302050in}{2.154512in}}%
\pgfpathlineto{\pgfqpoint{3.302050in}{2.151563in}}%
\pgfpathmoveto{\pgfqpoint{3.297509in}{2.154512in}}%
\pgfpathlineto{\pgfqpoint{3.297509in}{2.154512in}}%
\pgfpathlineto{\pgfqpoint{3.297509in}{2.157461in}}%
\pgfpathlineto{\pgfqpoint{3.302050in}{2.157461in}}%
\pgfpathlineto{\pgfqpoint{3.302050in}{2.154512in}}%
\pgfpathmoveto{\pgfqpoint{3.302050in}{2.154512in}}%
\pgfpathlineto{\pgfqpoint{3.302050in}{2.154512in}}%
\pgfpathlineto{\pgfqpoint{3.302050in}{2.157461in}}%
\pgfpathlineto{\pgfqpoint{3.306591in}{2.157461in}}%
\pgfpathlineto{\pgfqpoint{3.306591in}{2.154512in}}%
\pgfpathmoveto{\pgfqpoint{3.302050in}{2.157461in}}%
\pgfpathlineto{\pgfqpoint{3.302050in}{2.157461in}}%
\pgfpathlineto{\pgfqpoint{3.302050in}{2.160411in}}%
\pgfpathlineto{\pgfqpoint{3.306591in}{2.160411in}}%
\pgfpathlineto{\pgfqpoint{3.306591in}{2.157461in}}%
\pgfpathmoveto{\pgfqpoint{3.306591in}{2.157461in}}%
\pgfpathlineto{\pgfqpoint{3.306591in}{2.157461in}}%
\pgfpathlineto{\pgfqpoint{3.306591in}{2.160411in}}%
\pgfpathlineto{\pgfqpoint{3.311132in}{2.160411in}}%
\pgfpathlineto{\pgfqpoint{3.311132in}{2.157461in}}%
\pgfpathmoveto{\pgfqpoint{3.306591in}{2.160411in}}%
\pgfpathlineto{\pgfqpoint{3.306591in}{2.160411in}}%
\pgfpathlineto{\pgfqpoint{3.306591in}{2.163360in}}%
\pgfpathlineto{\pgfqpoint{3.311132in}{2.163360in}}%
\pgfpathlineto{\pgfqpoint{3.311132in}{2.160411in}}%
\pgfpathmoveto{\pgfqpoint{3.311132in}{2.160411in}}%
\pgfpathlineto{\pgfqpoint{3.311132in}{2.160411in}}%
\pgfpathlineto{\pgfqpoint{3.311132in}{2.163360in}}%
\pgfpathlineto{\pgfqpoint{3.315673in}{2.163360in}}%
\pgfpathlineto{\pgfqpoint{3.315673in}{2.160411in}}%
\pgfpathmoveto{\pgfqpoint{3.311132in}{2.163360in}}%
\pgfpathlineto{\pgfqpoint{3.311132in}{2.163360in}}%
\pgfpathlineto{\pgfqpoint{3.311132in}{2.166309in}}%
\pgfpathlineto{\pgfqpoint{3.315673in}{2.166309in}}%
\pgfpathlineto{\pgfqpoint{3.315673in}{2.163360in}}%
\pgfpathmoveto{\pgfqpoint{3.315673in}{2.163360in}}%
\pgfpathlineto{\pgfqpoint{3.315673in}{2.163360in}}%
\pgfpathlineto{\pgfqpoint{3.315673in}{2.166309in}}%
\pgfpathlineto{\pgfqpoint{3.320214in}{2.166309in}}%
\pgfpathlineto{\pgfqpoint{3.320214in}{2.163360in}}%
\pgfpathmoveto{\pgfqpoint{3.315673in}{2.166309in}}%
\pgfpathlineto{\pgfqpoint{3.315673in}{2.166309in}}%
\pgfpathlineto{\pgfqpoint{3.315673in}{2.169258in}}%
\pgfpathlineto{\pgfqpoint{3.320214in}{2.169258in}}%
\pgfpathlineto{\pgfqpoint{3.320214in}{2.166309in}}%
\pgfpathmoveto{\pgfqpoint{3.320214in}{2.166309in}}%
\pgfpathlineto{\pgfqpoint{3.320214in}{2.166309in}}%
\pgfpathlineto{\pgfqpoint{3.320214in}{2.169258in}}%
\pgfpathlineto{\pgfqpoint{3.324755in}{2.169258in}}%
\pgfpathlineto{\pgfqpoint{3.324755in}{2.166309in}}%
\pgfpathmoveto{\pgfqpoint{3.320214in}{2.169258in}}%
\pgfpathlineto{\pgfqpoint{3.320214in}{2.169258in}}%
\pgfpathlineto{\pgfqpoint{3.320214in}{2.172208in}}%
\pgfpathlineto{\pgfqpoint{3.324755in}{2.172208in}}%
\pgfpathlineto{\pgfqpoint{3.324755in}{2.169258in}}%
\pgfpathmoveto{\pgfqpoint{3.324755in}{2.169258in}}%
\pgfpathlineto{\pgfqpoint{3.324755in}{2.169258in}}%
\pgfpathlineto{\pgfqpoint{3.324755in}{2.172208in}}%
\pgfpathlineto{\pgfqpoint{3.329296in}{2.172208in}}%
\pgfpathlineto{\pgfqpoint{3.329296in}{2.169258in}}%
\pgfpathmoveto{\pgfqpoint{3.324755in}{2.172208in}}%
\pgfpathlineto{\pgfqpoint{3.324755in}{2.172208in}}%
\pgfpathlineto{\pgfqpoint{3.324755in}{2.175157in}}%
\pgfpathlineto{\pgfqpoint{3.329296in}{2.175157in}}%
\pgfpathlineto{\pgfqpoint{3.329296in}{2.172208in}}%
\pgfpathmoveto{\pgfqpoint{3.329296in}{2.172208in}}%
\pgfpathlineto{\pgfqpoint{3.329296in}{2.172208in}}%
\pgfpathlineto{\pgfqpoint{3.329296in}{2.175157in}}%
\pgfpathlineto{\pgfqpoint{3.333837in}{2.175157in}}%
\pgfpathlineto{\pgfqpoint{3.333837in}{2.172208in}}%
\pgfpathmoveto{\pgfqpoint{3.329296in}{2.175157in}}%
\pgfpathlineto{\pgfqpoint{3.329296in}{2.175157in}}%
\pgfpathlineto{\pgfqpoint{3.329296in}{2.178106in}}%
\pgfpathlineto{\pgfqpoint{3.333837in}{2.178106in}}%
\pgfpathlineto{\pgfqpoint{3.333837in}{2.175157in}}%
\pgfpathmoveto{\pgfqpoint{3.333837in}{2.175157in}}%
\pgfpathlineto{\pgfqpoint{3.333837in}{2.175157in}}%
\pgfpathlineto{\pgfqpoint{3.333837in}{2.178106in}}%
\pgfpathlineto{\pgfqpoint{3.338378in}{2.178106in}}%
\pgfpathlineto{\pgfqpoint{3.338378in}{2.175157in}}%
\pgfpathmoveto{\pgfqpoint{3.333837in}{2.178106in}}%
\pgfpathlineto{\pgfqpoint{3.333837in}{2.178106in}}%
\pgfpathlineto{\pgfqpoint{3.333837in}{2.181055in}}%
\pgfpathlineto{\pgfqpoint{3.338378in}{2.181055in}}%
\pgfpathlineto{\pgfqpoint{3.338378in}{2.178106in}}%
\pgfpathmoveto{\pgfqpoint{3.338378in}{2.178106in}}%
\pgfpathlineto{\pgfqpoint{3.338378in}{2.178106in}}%
\pgfpathlineto{\pgfqpoint{3.338378in}{2.181055in}}%
\pgfpathlineto{\pgfqpoint{3.342919in}{2.181055in}}%
\pgfpathlineto{\pgfqpoint{3.342919in}{2.178106in}}%
\pgfpathmoveto{\pgfqpoint{3.338378in}{2.181055in}}%
\pgfpathlineto{\pgfqpoint{3.338378in}{2.181055in}}%
\pgfpathlineto{\pgfqpoint{3.338378in}{2.184005in}}%
\pgfpathlineto{\pgfqpoint{3.342919in}{2.184005in}}%
\pgfpathlineto{\pgfqpoint{3.342919in}{2.181055in}}%
\pgfpathmoveto{\pgfqpoint{3.342919in}{2.181055in}}%
\pgfpathlineto{\pgfqpoint{3.342919in}{2.181055in}}%
\pgfpathlineto{\pgfqpoint{3.342919in}{2.184005in}}%
\pgfpathlineto{\pgfqpoint{3.347460in}{2.184005in}}%
\pgfpathlineto{\pgfqpoint{3.347460in}{2.181055in}}%
\pgfpathmoveto{\pgfqpoint{3.342919in}{2.184005in}}%
\pgfpathlineto{\pgfqpoint{3.342919in}{2.184005in}}%
\pgfpathlineto{\pgfqpoint{3.342919in}{2.186954in}}%
\pgfpathlineto{\pgfqpoint{3.347460in}{2.186954in}}%
\pgfpathlineto{\pgfqpoint{3.347460in}{2.184005in}}%
\pgfpathmoveto{\pgfqpoint{3.347460in}{2.184005in}}%
\pgfpathlineto{\pgfqpoint{3.347460in}{2.184005in}}%
\pgfpathlineto{\pgfqpoint{3.347460in}{2.186954in}}%
\pgfpathlineto{\pgfqpoint{3.352001in}{2.186954in}}%
\pgfpathlineto{\pgfqpoint{3.352001in}{2.184005in}}%
\pgfpathmoveto{\pgfqpoint{3.347460in}{2.186954in}}%
\pgfpathlineto{\pgfqpoint{3.347460in}{2.186954in}}%
\pgfpathlineto{\pgfqpoint{3.347460in}{2.189903in}}%
\pgfpathlineto{\pgfqpoint{3.352001in}{2.189903in}}%
\pgfpathlineto{\pgfqpoint{3.352001in}{2.186954in}}%
\pgfpathmoveto{\pgfqpoint{3.352001in}{2.186954in}}%
\pgfpathlineto{\pgfqpoint{3.352001in}{2.186954in}}%
\pgfpathlineto{\pgfqpoint{3.352001in}{2.189903in}}%
\pgfpathlineto{\pgfqpoint{3.356542in}{2.189903in}}%
\pgfpathlineto{\pgfqpoint{3.356542in}{2.186954in}}%
\pgfpathmoveto{\pgfqpoint{3.352001in}{2.189903in}}%
\pgfpathlineto{\pgfqpoint{3.352001in}{2.189903in}}%
\pgfpathlineto{\pgfqpoint{3.352001in}{2.192852in}}%
\pgfpathlineto{\pgfqpoint{3.356542in}{2.192852in}}%
\pgfpathlineto{\pgfqpoint{3.356542in}{2.189903in}}%
\pgfpathmoveto{\pgfqpoint{3.356542in}{2.189903in}}%
\pgfpathlineto{\pgfqpoint{3.356542in}{2.189903in}}%
\pgfpathlineto{\pgfqpoint{3.356542in}{2.192852in}}%
\pgfpathlineto{\pgfqpoint{3.361083in}{2.192852in}}%
\pgfpathlineto{\pgfqpoint{3.361083in}{2.189903in}}%
\pgfpathmoveto{\pgfqpoint{3.356542in}{2.192852in}}%
\pgfpathlineto{\pgfqpoint{3.356542in}{2.192852in}}%
\pgfpathlineto{\pgfqpoint{3.356542in}{2.195802in}}%
\pgfpathlineto{\pgfqpoint{3.361083in}{2.195802in}}%
\pgfpathlineto{\pgfqpoint{3.361083in}{2.192852in}}%
\pgfpathmoveto{\pgfqpoint{3.361083in}{2.192852in}}%
\pgfpathlineto{\pgfqpoint{3.361083in}{2.192852in}}%
\pgfpathlineto{\pgfqpoint{3.361083in}{2.195802in}}%
\pgfpathlineto{\pgfqpoint{3.365624in}{2.195802in}}%
\pgfpathlineto{\pgfqpoint{3.365624in}{2.192852in}}%
\pgfpathmoveto{\pgfqpoint{3.361083in}{2.195802in}}%
\pgfpathlineto{\pgfqpoint{3.361083in}{2.195802in}}%
\pgfpathlineto{\pgfqpoint{3.361083in}{2.198751in}}%
\pgfpathlineto{\pgfqpoint{3.365624in}{2.198751in}}%
\pgfpathlineto{\pgfqpoint{3.365624in}{2.195802in}}%
\pgfpathmoveto{\pgfqpoint{3.365624in}{2.195802in}}%
\pgfpathlineto{\pgfqpoint{3.365624in}{2.195802in}}%
\pgfpathlineto{\pgfqpoint{3.365624in}{2.198751in}}%
\pgfpathlineto{\pgfqpoint{3.370165in}{2.198751in}}%
\pgfpathlineto{\pgfqpoint{3.370165in}{2.195802in}}%
\pgfpathmoveto{\pgfqpoint{3.365624in}{2.198751in}}%
\pgfpathlineto{\pgfqpoint{3.365624in}{2.198751in}}%
\pgfpathlineto{\pgfqpoint{3.365624in}{2.201700in}}%
\pgfpathlineto{\pgfqpoint{3.370165in}{2.201700in}}%
\pgfpathlineto{\pgfqpoint{3.370165in}{2.198751in}}%
\pgfpathmoveto{\pgfqpoint{3.370165in}{2.198751in}}%
\pgfpathlineto{\pgfqpoint{3.370165in}{2.198751in}}%
\pgfpathlineto{\pgfqpoint{3.370165in}{2.201700in}}%
\pgfpathlineto{\pgfqpoint{3.374706in}{2.201700in}}%
\pgfpathlineto{\pgfqpoint{3.374706in}{2.198751in}}%
\pgfpathmoveto{\pgfqpoint{3.370165in}{2.201700in}}%
\pgfpathlineto{\pgfqpoint{3.370165in}{2.201700in}}%
\pgfpathlineto{\pgfqpoint{3.370165in}{2.204649in}}%
\pgfpathlineto{\pgfqpoint{3.374706in}{2.204649in}}%
\pgfpathlineto{\pgfqpoint{3.374706in}{2.201700in}}%
\pgfpathmoveto{\pgfqpoint{3.374706in}{2.201700in}}%
\pgfpathlineto{\pgfqpoint{3.374706in}{2.201700in}}%
\pgfpathlineto{\pgfqpoint{3.374706in}{2.204649in}}%
\pgfpathlineto{\pgfqpoint{3.379247in}{2.204649in}}%
\pgfpathlineto{\pgfqpoint{3.379247in}{2.201700in}}%
\pgfpathmoveto{\pgfqpoint{3.374706in}{2.204649in}}%
\pgfpathlineto{\pgfqpoint{3.374706in}{2.204649in}}%
\pgfpathlineto{\pgfqpoint{3.374706in}{2.207598in}}%
\pgfpathlineto{\pgfqpoint{3.379247in}{2.207598in}}%
\pgfpathlineto{\pgfqpoint{3.379247in}{2.204649in}}%
\pgfpathmoveto{\pgfqpoint{3.379247in}{2.204649in}}%
\pgfpathlineto{\pgfqpoint{3.379247in}{2.204649in}}%
\pgfpathlineto{\pgfqpoint{3.379247in}{2.207598in}}%
\pgfpathlineto{\pgfqpoint{3.383788in}{2.207598in}}%
\pgfpathlineto{\pgfqpoint{3.383788in}{2.204649in}}%
\pgfpathmoveto{\pgfqpoint{3.379247in}{2.207598in}}%
\pgfpathlineto{\pgfqpoint{3.379247in}{2.207598in}}%
\pgfpathlineto{\pgfqpoint{3.379247in}{2.210547in}}%
\pgfpathlineto{\pgfqpoint{3.383788in}{2.210547in}}%
\pgfpathlineto{\pgfqpoint{3.383788in}{2.207598in}}%
\pgfpathmoveto{\pgfqpoint{3.383788in}{2.207598in}}%
\pgfpathlineto{\pgfqpoint{3.383788in}{2.207598in}}%
\pgfpathlineto{\pgfqpoint{3.383788in}{2.210547in}}%
\pgfpathlineto{\pgfqpoint{3.388329in}{2.210547in}}%
\pgfpathlineto{\pgfqpoint{3.388329in}{2.207598in}}%
\pgfpathmoveto{\pgfqpoint{3.383788in}{2.210547in}}%
\pgfpathlineto{\pgfqpoint{3.383788in}{2.210547in}}%
\pgfpathlineto{\pgfqpoint{3.383788in}{2.213497in}}%
\pgfpathlineto{\pgfqpoint{3.388329in}{2.213497in}}%
\pgfpathlineto{\pgfqpoint{3.388329in}{2.210547in}}%
\pgfpathmoveto{\pgfqpoint{3.388329in}{2.210547in}}%
\pgfpathlineto{\pgfqpoint{3.388329in}{2.210547in}}%
\pgfpathlineto{\pgfqpoint{3.388329in}{2.213497in}}%
\pgfpathlineto{\pgfqpoint{3.392870in}{2.213497in}}%
\pgfpathlineto{\pgfqpoint{3.392870in}{2.210547in}}%
\pgfpathmoveto{\pgfqpoint{3.388329in}{2.213497in}}%
\pgfpathlineto{\pgfqpoint{3.388329in}{2.213497in}}%
\pgfpathlineto{\pgfqpoint{3.388329in}{2.216446in}}%
\pgfpathlineto{\pgfqpoint{3.392870in}{2.216446in}}%
\pgfpathlineto{\pgfqpoint{3.392870in}{2.213497in}}%
\pgfpathmoveto{\pgfqpoint{3.392870in}{2.213497in}}%
\pgfpathlineto{\pgfqpoint{3.392870in}{2.213497in}}%
\pgfpathlineto{\pgfqpoint{3.392870in}{2.216446in}}%
\pgfpathlineto{\pgfqpoint{3.397411in}{2.216446in}}%
\pgfpathlineto{\pgfqpoint{3.397411in}{2.213497in}}%
\pgfpathmoveto{\pgfqpoint{3.392870in}{2.216446in}}%
\pgfpathlineto{\pgfqpoint{3.392870in}{2.216446in}}%
\pgfpathlineto{\pgfqpoint{3.392870in}{2.219395in}}%
\pgfpathlineto{\pgfqpoint{3.397411in}{2.219395in}}%
\pgfpathlineto{\pgfqpoint{3.397411in}{2.216446in}}%
\pgfpathmoveto{\pgfqpoint{3.397411in}{2.216446in}}%
\pgfpathlineto{\pgfqpoint{3.397411in}{2.216446in}}%
\pgfpathlineto{\pgfqpoint{3.397411in}{2.219395in}}%
\pgfpathlineto{\pgfqpoint{3.401952in}{2.219395in}}%
\pgfpathlineto{\pgfqpoint{3.401952in}{2.216446in}}%
\pgfpathmoveto{\pgfqpoint{3.397411in}{2.219395in}}%
\pgfpathlineto{\pgfqpoint{3.397411in}{2.219395in}}%
\pgfpathlineto{\pgfqpoint{3.397411in}{2.222344in}}%
\pgfpathlineto{\pgfqpoint{3.401952in}{2.222344in}}%
\pgfpathlineto{\pgfqpoint{3.401952in}{2.219395in}}%
\pgfpathmoveto{\pgfqpoint{3.401952in}{2.219395in}}%
\pgfpathlineto{\pgfqpoint{3.401952in}{2.219395in}}%
\pgfpathlineto{\pgfqpoint{3.401952in}{2.222344in}}%
\pgfpathlineto{\pgfqpoint{3.406493in}{2.222344in}}%
\pgfpathlineto{\pgfqpoint{3.406493in}{2.219395in}}%
\pgfpathmoveto{\pgfqpoint{3.401952in}{2.222344in}}%
\pgfpathlineto{\pgfqpoint{3.401952in}{2.222344in}}%
\pgfpathlineto{\pgfqpoint{3.401952in}{2.225293in}}%
\pgfpathlineto{\pgfqpoint{3.406493in}{2.225293in}}%
\pgfpathlineto{\pgfqpoint{3.406493in}{2.222344in}}%
\pgfpathmoveto{\pgfqpoint{3.406493in}{2.222344in}}%
\pgfpathlineto{\pgfqpoint{3.406493in}{2.222344in}}%
\pgfpathlineto{\pgfqpoint{3.406493in}{2.225293in}}%
\pgfpathlineto{\pgfqpoint{3.411034in}{2.225293in}}%
\pgfpathlineto{\pgfqpoint{3.411034in}{2.222344in}}%
\pgfpathmoveto{\pgfqpoint{3.406493in}{2.225293in}}%
\pgfpathlineto{\pgfqpoint{3.406493in}{2.225293in}}%
\pgfpathlineto{\pgfqpoint{3.406493in}{2.228242in}}%
\pgfpathlineto{\pgfqpoint{3.411034in}{2.228242in}}%
\pgfpathlineto{\pgfqpoint{3.411034in}{2.225293in}}%
\pgfpathmoveto{\pgfqpoint{3.411034in}{2.225293in}}%
\pgfpathlineto{\pgfqpoint{3.411034in}{2.225293in}}%
\pgfpathlineto{\pgfqpoint{3.411034in}{2.228242in}}%
\pgfpathlineto{\pgfqpoint{3.415575in}{2.228242in}}%
\pgfpathlineto{\pgfqpoint{3.415575in}{2.225293in}}%
\pgfpathmoveto{\pgfqpoint{3.411034in}{2.228242in}}%
\pgfpathlineto{\pgfqpoint{3.411034in}{2.228242in}}%
\pgfpathlineto{\pgfqpoint{3.411034in}{2.231191in}}%
\pgfpathlineto{\pgfqpoint{3.415575in}{2.231191in}}%
\pgfpathlineto{\pgfqpoint{3.415575in}{2.228242in}}%
\pgfpathmoveto{\pgfqpoint{3.415575in}{2.228242in}}%
\pgfpathlineto{\pgfqpoint{3.415575in}{2.228242in}}%
\pgfpathlineto{\pgfqpoint{3.415575in}{2.231191in}}%
\pgfpathlineto{\pgfqpoint{3.420116in}{2.231191in}}%
\pgfpathlineto{\pgfqpoint{3.420116in}{2.228242in}}%
\pgfpathmoveto{\pgfqpoint{3.415575in}{2.231191in}}%
\pgfpathlineto{\pgfqpoint{3.415575in}{2.231191in}}%
\pgfpathlineto{\pgfqpoint{3.415575in}{2.234141in}}%
\pgfpathlineto{\pgfqpoint{3.420116in}{2.234141in}}%
\pgfpathlineto{\pgfqpoint{3.420116in}{2.231191in}}%
\pgfpathmoveto{\pgfqpoint{3.420116in}{2.231191in}}%
\pgfpathlineto{\pgfqpoint{3.420116in}{2.231191in}}%
\pgfpathlineto{\pgfqpoint{3.420116in}{2.234141in}}%
\pgfpathlineto{\pgfqpoint{3.424657in}{2.234141in}}%
\pgfpathlineto{\pgfqpoint{3.424657in}{2.231191in}}%
\pgfpathmoveto{\pgfqpoint{3.420116in}{2.234141in}}%
\pgfpathlineto{\pgfqpoint{3.420116in}{2.234141in}}%
\pgfpathlineto{\pgfqpoint{3.420116in}{2.237090in}}%
\pgfpathlineto{\pgfqpoint{3.424657in}{2.237090in}}%
\pgfpathlineto{\pgfqpoint{3.424657in}{2.234141in}}%
\pgfpathmoveto{\pgfqpoint{3.424657in}{2.234141in}}%
\pgfpathlineto{\pgfqpoint{3.424657in}{2.234141in}}%
\pgfpathlineto{\pgfqpoint{3.424657in}{2.237090in}}%
\pgfpathlineto{\pgfqpoint{3.429198in}{2.237090in}}%
\pgfpathlineto{\pgfqpoint{3.429198in}{2.234141in}}%
\pgfpathmoveto{\pgfqpoint{3.424657in}{2.237090in}}%
\pgfpathlineto{\pgfqpoint{3.424657in}{2.237090in}}%
\pgfpathlineto{\pgfqpoint{3.424657in}{2.240039in}}%
\pgfpathlineto{\pgfqpoint{3.429198in}{2.240039in}}%
\pgfpathlineto{\pgfqpoint{3.429198in}{2.237090in}}%
\pgfpathmoveto{\pgfqpoint{3.429198in}{2.237090in}}%
\pgfpathlineto{\pgfqpoint{3.429198in}{2.237090in}}%
\pgfpathlineto{\pgfqpoint{3.429198in}{2.240039in}}%
\pgfpathlineto{\pgfqpoint{3.433739in}{2.240039in}}%
\pgfpathlineto{\pgfqpoint{3.433739in}{2.237090in}}%
\pgfpathmoveto{\pgfqpoint{3.429198in}{2.240039in}}%
\pgfpathlineto{\pgfqpoint{3.429198in}{2.240039in}}%
\pgfpathlineto{\pgfqpoint{3.429198in}{2.242988in}}%
\pgfpathlineto{\pgfqpoint{3.433739in}{2.242988in}}%
\pgfpathlineto{\pgfqpoint{3.433739in}{2.240039in}}%
\pgfpathmoveto{\pgfqpoint{3.433739in}{2.240039in}}%
\pgfpathlineto{\pgfqpoint{3.433739in}{2.240039in}}%
\pgfpathlineto{\pgfqpoint{3.433739in}{2.242988in}}%
\pgfpathlineto{\pgfqpoint{3.438280in}{2.242988in}}%
\pgfpathlineto{\pgfqpoint{3.438280in}{2.240039in}}%
\pgfpathmoveto{\pgfqpoint{3.433739in}{2.242988in}}%
\pgfpathlineto{\pgfqpoint{3.433739in}{2.242988in}}%
\pgfpathlineto{\pgfqpoint{3.433739in}{2.245937in}}%
\pgfpathlineto{\pgfqpoint{3.438280in}{2.245937in}}%
\pgfpathlineto{\pgfqpoint{3.438280in}{2.242988in}}%
\pgfpathmoveto{\pgfqpoint{3.438280in}{2.242988in}}%
\pgfpathlineto{\pgfqpoint{3.438280in}{2.242988in}}%
\pgfpathlineto{\pgfqpoint{3.438280in}{2.245937in}}%
\pgfpathlineto{\pgfqpoint{3.442821in}{2.245937in}}%
\pgfpathlineto{\pgfqpoint{3.442821in}{2.242988in}}%
\pgfpathmoveto{\pgfqpoint{3.438280in}{2.245937in}}%
\pgfpathlineto{\pgfqpoint{3.438280in}{2.245937in}}%
\pgfpathlineto{\pgfqpoint{3.438280in}{2.248886in}}%
\pgfpathlineto{\pgfqpoint{3.442821in}{2.248886in}}%
\pgfpathlineto{\pgfqpoint{3.442821in}{2.245937in}}%
\pgfpathmoveto{\pgfqpoint{3.442821in}{2.245937in}}%
\pgfpathlineto{\pgfqpoint{3.442821in}{2.245937in}}%
\pgfpathlineto{\pgfqpoint{3.442821in}{2.248886in}}%
\pgfpathlineto{\pgfqpoint{3.447362in}{2.248886in}}%
\pgfpathlineto{\pgfqpoint{3.447362in}{2.245937in}}%
\pgfpathmoveto{\pgfqpoint{3.442821in}{2.248886in}}%
\pgfpathlineto{\pgfqpoint{3.442821in}{2.248886in}}%
\pgfpathlineto{\pgfqpoint{3.442821in}{2.251835in}}%
\pgfpathlineto{\pgfqpoint{3.447362in}{2.251835in}}%
\pgfpathlineto{\pgfqpoint{3.447362in}{2.248886in}}%
\pgfpathmoveto{\pgfqpoint{3.447362in}{2.248886in}}%
\pgfpathlineto{\pgfqpoint{3.447362in}{2.248886in}}%
\pgfpathlineto{\pgfqpoint{3.447362in}{2.251835in}}%
\pgfpathlineto{\pgfqpoint{3.451903in}{2.251835in}}%
\pgfpathlineto{\pgfqpoint{3.451903in}{2.248886in}}%
\pgfpathmoveto{\pgfqpoint{3.447362in}{2.251835in}}%
\pgfpathlineto{\pgfqpoint{3.447362in}{2.251835in}}%
\pgfpathlineto{\pgfqpoint{3.447362in}{2.254784in}}%
\pgfpathlineto{\pgfqpoint{3.451903in}{2.254784in}}%
\pgfpathlineto{\pgfqpoint{3.451903in}{2.251835in}}%
\pgfpathmoveto{\pgfqpoint{3.451903in}{2.251835in}}%
\pgfpathlineto{\pgfqpoint{3.451903in}{2.251835in}}%
\pgfpathlineto{\pgfqpoint{3.451903in}{2.254784in}}%
\pgfpathlineto{\pgfqpoint{3.456444in}{2.254784in}}%
\pgfpathlineto{\pgfqpoint{3.456444in}{2.251835in}}%
\pgfpathmoveto{\pgfqpoint{3.451903in}{2.254784in}}%
\pgfpathlineto{\pgfqpoint{3.451903in}{2.254784in}}%
\pgfpathlineto{\pgfqpoint{3.451903in}{2.257734in}}%
\pgfpathlineto{\pgfqpoint{3.456444in}{2.257734in}}%
\pgfpathlineto{\pgfqpoint{3.456444in}{2.254784in}}%
\pgfpathmoveto{\pgfqpoint{3.456444in}{2.254784in}}%
\pgfpathlineto{\pgfqpoint{3.456444in}{2.254784in}}%
\pgfpathlineto{\pgfqpoint{3.456444in}{2.257734in}}%
\pgfpathlineto{\pgfqpoint{3.460985in}{2.257734in}}%
\pgfpathlineto{\pgfqpoint{3.460985in}{2.254784in}}%
\pgfpathmoveto{\pgfqpoint{3.456444in}{2.257734in}}%
\pgfpathlineto{\pgfqpoint{3.456444in}{2.257734in}}%
\pgfpathlineto{\pgfqpoint{3.456444in}{2.260683in}}%
\pgfpathlineto{\pgfqpoint{3.460985in}{2.260683in}}%
\pgfpathlineto{\pgfqpoint{3.460985in}{2.257734in}}%
\pgfpathmoveto{\pgfqpoint{3.460985in}{2.257734in}}%
\pgfpathlineto{\pgfqpoint{3.460985in}{2.257734in}}%
\pgfpathlineto{\pgfqpoint{3.460985in}{2.260683in}}%
\pgfpathlineto{\pgfqpoint{3.465526in}{2.260683in}}%
\pgfpathlineto{\pgfqpoint{3.465526in}{2.257734in}}%
\pgfpathmoveto{\pgfqpoint{3.460985in}{2.260683in}}%
\pgfpathlineto{\pgfqpoint{3.460985in}{2.260683in}}%
\pgfpathlineto{\pgfqpoint{3.460985in}{2.263632in}}%
\pgfpathlineto{\pgfqpoint{3.465526in}{2.263632in}}%
\pgfpathlineto{\pgfqpoint{3.465526in}{2.260683in}}%
\pgfpathmoveto{\pgfqpoint{3.465526in}{2.260683in}}%
\pgfpathlineto{\pgfqpoint{3.465526in}{2.260683in}}%
\pgfpathlineto{\pgfqpoint{3.465526in}{2.263632in}}%
\pgfpathlineto{\pgfqpoint{3.470066in}{2.263632in}}%
\pgfpathlineto{\pgfqpoint{3.470066in}{2.260683in}}%
\pgfpathmoveto{\pgfqpoint{3.465526in}{2.263632in}}%
\pgfpathlineto{\pgfqpoint{3.465526in}{2.263632in}}%
\pgfpathlineto{\pgfqpoint{3.465526in}{2.266581in}}%
\pgfpathlineto{\pgfqpoint{3.470066in}{2.266581in}}%
\pgfpathlineto{\pgfqpoint{3.470066in}{2.263632in}}%
\pgfpathmoveto{\pgfqpoint{3.470066in}{2.263632in}}%
\pgfpathlineto{\pgfqpoint{3.470066in}{2.263632in}}%
\pgfpathlineto{\pgfqpoint{3.470066in}{2.266581in}}%
\pgfpathlineto{\pgfqpoint{3.474607in}{2.266581in}}%
\pgfpathlineto{\pgfqpoint{3.474607in}{2.263632in}}%
\pgfpathmoveto{\pgfqpoint{3.470066in}{2.266581in}}%
\pgfpathlineto{\pgfqpoint{3.470066in}{2.266581in}}%
\pgfpathlineto{\pgfqpoint{3.470066in}{2.269530in}}%
\pgfpathlineto{\pgfqpoint{3.474607in}{2.269530in}}%
\pgfpathlineto{\pgfqpoint{3.474607in}{2.266581in}}%
\pgfpathmoveto{\pgfqpoint{3.474607in}{2.266581in}}%
\pgfpathlineto{\pgfqpoint{3.474607in}{2.266581in}}%
\pgfpathlineto{\pgfqpoint{3.474607in}{2.269530in}}%
\pgfpathlineto{\pgfqpoint{3.479148in}{2.269530in}}%
\pgfpathlineto{\pgfqpoint{3.479148in}{2.266581in}}%
\pgfpathmoveto{\pgfqpoint{3.474607in}{2.269530in}}%
\pgfpathlineto{\pgfqpoint{3.474607in}{2.269530in}}%
\pgfpathlineto{\pgfqpoint{3.474607in}{2.272479in}}%
\pgfpathlineto{\pgfqpoint{3.479148in}{2.272479in}}%
\pgfpathlineto{\pgfqpoint{3.479148in}{2.269530in}}%
\pgfpathmoveto{\pgfqpoint{3.479148in}{2.269530in}}%
\pgfpathlineto{\pgfqpoint{3.479148in}{2.269530in}}%
\pgfpathlineto{\pgfqpoint{3.479148in}{2.272479in}}%
\pgfpathlineto{\pgfqpoint{3.483689in}{2.272479in}}%
\pgfpathlineto{\pgfqpoint{3.483689in}{2.269530in}}%
\pgfpathmoveto{\pgfqpoint{3.479148in}{2.272479in}}%
\pgfpathlineto{\pgfqpoint{3.479148in}{2.272479in}}%
\pgfpathlineto{\pgfqpoint{3.479148in}{2.275428in}}%
\pgfpathlineto{\pgfqpoint{3.483689in}{2.275428in}}%
\pgfpathlineto{\pgfqpoint{3.483689in}{2.272479in}}%
\pgfpathmoveto{\pgfqpoint{3.483689in}{2.272479in}}%
\pgfpathlineto{\pgfqpoint{3.483689in}{2.272479in}}%
\pgfpathlineto{\pgfqpoint{3.483689in}{2.275428in}}%
\pgfpathlineto{\pgfqpoint{3.488230in}{2.275428in}}%
\pgfpathlineto{\pgfqpoint{3.488230in}{2.272479in}}%
\pgfpathmoveto{\pgfqpoint{3.483689in}{2.275428in}}%
\pgfpathlineto{\pgfqpoint{3.483689in}{2.275428in}}%
\pgfpathlineto{\pgfqpoint{3.483689in}{2.278378in}}%
\pgfpathlineto{\pgfqpoint{3.488230in}{2.278378in}}%
\pgfpathlineto{\pgfqpoint{3.488230in}{2.275428in}}%
\pgfpathmoveto{\pgfqpoint{3.483689in}{2.278378in}}%
\pgfpathlineto{\pgfqpoint{3.483689in}{2.278378in}}%
\pgfpathlineto{\pgfqpoint{3.483689in}{2.281327in}}%
\pgfpathlineto{\pgfqpoint{3.488230in}{2.281327in}}%
\pgfpathlineto{\pgfqpoint{3.488230in}{2.278378in}}%
\pgfpathmoveto{\pgfqpoint{3.488230in}{2.278378in}}%
\pgfpathlineto{\pgfqpoint{3.488230in}{2.278378in}}%
\pgfpathlineto{\pgfqpoint{3.488230in}{2.281327in}}%
\pgfpathlineto{\pgfqpoint{3.492771in}{2.281327in}}%
\pgfpathlineto{\pgfqpoint{3.492771in}{2.278378in}}%
\pgfpathmoveto{\pgfqpoint{3.488230in}{2.281327in}}%
\pgfpathlineto{\pgfqpoint{3.488230in}{2.281327in}}%
\pgfpathlineto{\pgfqpoint{3.488230in}{2.284276in}}%
\pgfpathlineto{\pgfqpoint{3.492771in}{2.284276in}}%
\pgfpathlineto{\pgfqpoint{3.492771in}{2.281327in}}%
\pgfpathmoveto{\pgfqpoint{3.492771in}{2.281327in}}%
\pgfpathlineto{\pgfqpoint{3.492771in}{2.281327in}}%
\pgfpathlineto{\pgfqpoint{3.492771in}{2.284276in}}%
\pgfpathlineto{\pgfqpoint{3.497312in}{2.284276in}}%
\pgfpathlineto{\pgfqpoint{3.497312in}{2.281327in}}%
\pgfpathmoveto{\pgfqpoint{3.492771in}{2.284276in}}%
\pgfpathlineto{\pgfqpoint{3.492771in}{2.284276in}}%
\pgfpathlineto{\pgfqpoint{3.492771in}{2.287225in}}%
\pgfpathlineto{\pgfqpoint{3.497312in}{2.287225in}}%
\pgfpathlineto{\pgfqpoint{3.497312in}{2.284276in}}%
\pgfpathmoveto{\pgfqpoint{3.497312in}{2.284276in}}%
\pgfpathlineto{\pgfqpoint{3.497312in}{2.284276in}}%
\pgfpathlineto{\pgfqpoint{3.497312in}{2.287225in}}%
\pgfpathlineto{\pgfqpoint{3.501853in}{2.287225in}}%
\pgfpathlineto{\pgfqpoint{3.501853in}{2.284276in}}%
\pgfpathmoveto{\pgfqpoint{3.497312in}{2.287225in}}%
\pgfpathlineto{\pgfqpoint{3.497312in}{2.287225in}}%
\pgfpathlineto{\pgfqpoint{3.497312in}{2.290174in}}%
\pgfpathlineto{\pgfqpoint{3.501853in}{2.290174in}}%
\pgfpathlineto{\pgfqpoint{3.501853in}{2.287225in}}%
\pgfpathmoveto{\pgfqpoint{3.501853in}{2.287225in}}%
\pgfpathlineto{\pgfqpoint{3.501853in}{2.287225in}}%
\pgfpathlineto{\pgfqpoint{3.501853in}{2.290174in}}%
\pgfpathlineto{\pgfqpoint{3.506394in}{2.290174in}}%
\pgfpathlineto{\pgfqpoint{3.506394in}{2.287225in}}%
\pgfpathmoveto{\pgfqpoint{3.501853in}{2.290174in}}%
\pgfpathlineto{\pgfqpoint{3.501853in}{2.290174in}}%
\pgfpathlineto{\pgfqpoint{3.501853in}{2.293123in}}%
\pgfpathlineto{\pgfqpoint{3.506394in}{2.293123in}}%
\pgfpathlineto{\pgfqpoint{3.506394in}{2.290174in}}%
\pgfpathmoveto{\pgfqpoint{3.506394in}{2.290174in}}%
\pgfpathlineto{\pgfqpoint{3.506394in}{2.290174in}}%
\pgfpathlineto{\pgfqpoint{3.506394in}{2.293123in}}%
\pgfpathlineto{\pgfqpoint{3.510935in}{2.293123in}}%
\pgfpathlineto{\pgfqpoint{3.510935in}{2.290174in}}%
\pgfpathmoveto{\pgfqpoint{3.506394in}{2.293123in}}%
\pgfpathlineto{\pgfqpoint{3.506394in}{2.293123in}}%
\pgfpathlineto{\pgfqpoint{3.506394in}{2.296073in}}%
\pgfpathlineto{\pgfqpoint{3.510935in}{2.296073in}}%
\pgfpathlineto{\pgfqpoint{3.510935in}{2.293123in}}%
\pgfpathmoveto{\pgfqpoint{3.510935in}{2.293123in}}%
\pgfpathlineto{\pgfqpoint{3.510935in}{2.293123in}}%
\pgfpathlineto{\pgfqpoint{3.510935in}{2.296073in}}%
\pgfpathlineto{\pgfqpoint{3.515476in}{2.296073in}}%
\pgfpathlineto{\pgfqpoint{3.515476in}{2.293123in}}%
\pgfpathmoveto{\pgfqpoint{3.510935in}{2.296073in}}%
\pgfpathlineto{\pgfqpoint{3.510935in}{2.296073in}}%
\pgfpathlineto{\pgfqpoint{3.510935in}{2.299022in}}%
\pgfpathlineto{\pgfqpoint{3.515476in}{2.299022in}}%
\pgfpathlineto{\pgfqpoint{3.515476in}{2.296073in}}%
\pgfpathmoveto{\pgfqpoint{3.515476in}{2.296073in}}%
\pgfpathlineto{\pgfqpoint{3.515476in}{2.296073in}}%
\pgfpathlineto{\pgfqpoint{3.515476in}{2.299022in}}%
\pgfpathlineto{\pgfqpoint{3.520017in}{2.299022in}}%
\pgfpathlineto{\pgfqpoint{3.520017in}{2.296073in}}%
\pgfpathmoveto{\pgfqpoint{3.515476in}{2.299022in}}%
\pgfpathlineto{\pgfqpoint{3.515476in}{2.299022in}}%
\pgfpathlineto{\pgfqpoint{3.515476in}{2.301971in}}%
\pgfpathlineto{\pgfqpoint{3.520017in}{2.301971in}}%
\pgfpathlineto{\pgfqpoint{3.520017in}{2.299022in}}%
\pgfpathmoveto{\pgfqpoint{3.520017in}{2.299022in}}%
\pgfpathlineto{\pgfqpoint{3.520017in}{2.299022in}}%
\pgfpathlineto{\pgfqpoint{3.520017in}{2.301971in}}%
\pgfpathlineto{\pgfqpoint{3.524558in}{2.301971in}}%
\pgfpathlineto{\pgfqpoint{3.524558in}{2.299022in}}%
\pgfpathmoveto{\pgfqpoint{3.524558in}{2.299022in}}%
\pgfpathlineto{\pgfqpoint{3.524558in}{2.299022in}}%
\pgfpathlineto{\pgfqpoint{3.524558in}{2.301971in}}%
\pgfpathlineto{\pgfqpoint{3.529099in}{2.301971in}}%
\pgfpathlineto{\pgfqpoint{3.529099in}{2.299022in}}%
\pgfpathmoveto{\pgfqpoint{3.524558in}{2.301971in}}%
\pgfpathlineto{\pgfqpoint{3.524558in}{2.301971in}}%
\pgfpathlineto{\pgfqpoint{3.524558in}{2.304921in}}%
\pgfpathlineto{\pgfqpoint{3.529099in}{2.304921in}}%
\pgfpathlineto{\pgfqpoint{3.529099in}{2.301971in}}%
\pgfpathmoveto{\pgfqpoint{3.529099in}{2.301971in}}%
\pgfpathlineto{\pgfqpoint{3.529099in}{2.301971in}}%
\pgfpathlineto{\pgfqpoint{3.529099in}{2.304921in}}%
\pgfpathlineto{\pgfqpoint{3.533640in}{2.304921in}}%
\pgfpathlineto{\pgfqpoint{3.533640in}{2.301971in}}%
\pgfpathmoveto{\pgfqpoint{3.529099in}{2.304921in}}%
\pgfpathlineto{\pgfqpoint{3.529099in}{2.304921in}}%
\pgfpathlineto{\pgfqpoint{3.529099in}{2.307870in}}%
\pgfpathlineto{\pgfqpoint{3.533640in}{2.307870in}}%
\pgfpathlineto{\pgfqpoint{3.533640in}{2.304921in}}%
\pgfpathmoveto{\pgfqpoint{3.533640in}{2.304921in}}%
\pgfpathlineto{\pgfqpoint{3.533640in}{2.304921in}}%
\pgfpathlineto{\pgfqpoint{3.533640in}{2.307870in}}%
\pgfpathlineto{\pgfqpoint{3.538181in}{2.307870in}}%
\pgfpathlineto{\pgfqpoint{3.538181in}{2.304921in}}%
\pgfpathmoveto{\pgfqpoint{3.533640in}{2.307870in}}%
\pgfpathlineto{\pgfqpoint{3.533640in}{2.307870in}}%
\pgfpathlineto{\pgfqpoint{3.533640in}{2.310819in}}%
\pgfpathlineto{\pgfqpoint{3.538181in}{2.310819in}}%
\pgfpathlineto{\pgfqpoint{3.538181in}{2.307870in}}%
\pgfpathmoveto{\pgfqpoint{3.538181in}{2.307870in}}%
\pgfpathlineto{\pgfqpoint{3.538181in}{2.307870in}}%
\pgfpathlineto{\pgfqpoint{3.538181in}{2.310819in}}%
\pgfpathlineto{\pgfqpoint{3.542722in}{2.310819in}}%
\pgfpathlineto{\pgfqpoint{3.542722in}{2.307870in}}%
\pgfpathmoveto{\pgfqpoint{3.538181in}{2.310819in}}%
\pgfpathlineto{\pgfqpoint{3.538181in}{2.310819in}}%
\pgfpathlineto{\pgfqpoint{3.538181in}{2.313769in}}%
\pgfpathlineto{\pgfqpoint{3.542722in}{2.313769in}}%
\pgfpathlineto{\pgfqpoint{3.542722in}{2.310819in}}%
\pgfpathmoveto{\pgfqpoint{3.542722in}{2.310819in}}%
\pgfpathlineto{\pgfqpoint{3.542722in}{2.310819in}}%
\pgfpathlineto{\pgfqpoint{3.542722in}{2.313769in}}%
\pgfpathlineto{\pgfqpoint{3.547264in}{2.313769in}}%
\pgfpathlineto{\pgfqpoint{3.547264in}{2.310819in}}%
\pgfpathmoveto{\pgfqpoint{3.542722in}{2.313769in}}%
\pgfpathlineto{\pgfqpoint{3.542722in}{2.313769in}}%
\pgfpathlineto{\pgfqpoint{3.542722in}{2.316718in}}%
\pgfpathlineto{\pgfqpoint{3.547264in}{2.316718in}}%
\pgfpathlineto{\pgfqpoint{3.547264in}{2.313769in}}%
\pgfpathmoveto{\pgfqpoint{3.547264in}{2.313769in}}%
\pgfpathlineto{\pgfqpoint{3.547264in}{2.313769in}}%
\pgfpathlineto{\pgfqpoint{3.547264in}{2.316718in}}%
\pgfpathlineto{\pgfqpoint{3.551805in}{2.316718in}}%
\pgfpathlineto{\pgfqpoint{3.551805in}{2.313769in}}%
\pgfpathmoveto{\pgfqpoint{3.547264in}{2.316718in}}%
\pgfpathlineto{\pgfqpoint{3.547264in}{2.316718in}}%
\pgfpathlineto{\pgfqpoint{3.547264in}{2.319667in}}%
\pgfpathlineto{\pgfqpoint{3.551805in}{2.319667in}}%
\pgfpathlineto{\pgfqpoint{3.551805in}{2.316718in}}%
\pgfpathmoveto{\pgfqpoint{3.551805in}{2.316718in}}%
\pgfpathlineto{\pgfqpoint{3.551805in}{2.316718in}}%
\pgfpathlineto{\pgfqpoint{3.551805in}{2.319667in}}%
\pgfpathlineto{\pgfqpoint{3.556346in}{2.319667in}}%
\pgfpathlineto{\pgfqpoint{3.556346in}{2.316718in}}%
\pgfpathmoveto{\pgfqpoint{3.551805in}{2.319667in}}%
\pgfpathlineto{\pgfqpoint{3.551805in}{2.319667in}}%
\pgfpathlineto{\pgfqpoint{3.551805in}{2.322617in}}%
\pgfpathlineto{\pgfqpoint{3.556346in}{2.322617in}}%
\pgfpathlineto{\pgfqpoint{3.556346in}{2.319667in}}%
\pgfpathmoveto{\pgfqpoint{3.556346in}{2.319667in}}%
\pgfpathlineto{\pgfqpoint{3.556346in}{2.319667in}}%
\pgfpathlineto{\pgfqpoint{3.556346in}{2.322617in}}%
\pgfpathlineto{\pgfqpoint{3.560887in}{2.322617in}}%
\pgfpathlineto{\pgfqpoint{3.560887in}{2.319667in}}%
\pgfpathmoveto{\pgfqpoint{3.556346in}{2.322617in}}%
\pgfpathlineto{\pgfqpoint{3.556346in}{2.322617in}}%
\pgfpathlineto{\pgfqpoint{3.556346in}{2.325566in}}%
\pgfpathlineto{\pgfqpoint{3.560887in}{2.325566in}}%
\pgfpathlineto{\pgfqpoint{3.560887in}{2.322617in}}%
\pgfpathmoveto{\pgfqpoint{3.560887in}{2.322617in}}%
\pgfpathlineto{\pgfqpoint{3.560887in}{2.322617in}}%
\pgfpathlineto{\pgfqpoint{3.560887in}{2.325566in}}%
\pgfpathlineto{\pgfqpoint{3.565428in}{2.325566in}}%
\pgfpathlineto{\pgfqpoint{3.565428in}{2.322617in}}%
\pgfpathmoveto{\pgfqpoint{3.560887in}{2.325566in}}%
\pgfpathlineto{\pgfqpoint{3.560887in}{2.325566in}}%
\pgfpathlineto{\pgfqpoint{3.560887in}{2.328515in}}%
\pgfpathlineto{\pgfqpoint{3.565428in}{2.328515in}}%
\pgfpathlineto{\pgfqpoint{3.565428in}{2.325566in}}%
\pgfpathmoveto{\pgfqpoint{3.565428in}{2.325566in}}%
\pgfpathlineto{\pgfqpoint{3.565428in}{2.325566in}}%
\pgfpathlineto{\pgfqpoint{3.565428in}{2.328515in}}%
\pgfpathlineto{\pgfqpoint{3.569969in}{2.328515in}}%
\pgfpathlineto{\pgfqpoint{3.569969in}{2.325566in}}%
\pgfpathmoveto{\pgfqpoint{3.565428in}{2.328515in}}%
\pgfpathlineto{\pgfqpoint{3.565428in}{2.328515in}}%
\pgfpathlineto{\pgfqpoint{3.565428in}{2.331465in}}%
\pgfpathlineto{\pgfqpoint{3.569969in}{2.331465in}}%
\pgfpathlineto{\pgfqpoint{3.569969in}{2.328515in}}%
\pgfpathmoveto{\pgfqpoint{3.569969in}{2.328515in}}%
\pgfpathlineto{\pgfqpoint{3.569969in}{2.328515in}}%
\pgfpathlineto{\pgfqpoint{3.569969in}{2.331465in}}%
\pgfpathlineto{\pgfqpoint{3.574510in}{2.331465in}}%
\pgfpathlineto{\pgfqpoint{3.574510in}{2.328515in}}%
\pgfpathmoveto{\pgfqpoint{3.569969in}{2.331465in}}%
\pgfpathlineto{\pgfqpoint{3.569969in}{2.331465in}}%
\pgfpathlineto{\pgfqpoint{3.569969in}{2.334414in}}%
\pgfpathlineto{\pgfqpoint{3.574510in}{2.334414in}}%
\pgfpathlineto{\pgfqpoint{3.574510in}{2.331465in}}%
\pgfpathmoveto{\pgfqpoint{3.574510in}{2.331465in}}%
\pgfpathlineto{\pgfqpoint{3.574510in}{2.331465in}}%
\pgfpathlineto{\pgfqpoint{3.574510in}{2.334414in}}%
\pgfpathlineto{\pgfqpoint{3.579051in}{2.334414in}}%
\pgfpathlineto{\pgfqpoint{3.579051in}{2.331465in}}%
\pgfpathmoveto{\pgfqpoint{3.574510in}{2.334414in}}%
\pgfpathlineto{\pgfqpoint{3.574510in}{2.334414in}}%
\pgfpathlineto{\pgfqpoint{3.574510in}{2.337363in}}%
\pgfpathlineto{\pgfqpoint{3.579051in}{2.337363in}}%
\pgfpathlineto{\pgfqpoint{3.579051in}{2.334414in}}%
\pgfpathmoveto{\pgfqpoint{3.579051in}{2.334414in}}%
\pgfpathlineto{\pgfqpoint{3.579051in}{2.334414in}}%
\pgfpathlineto{\pgfqpoint{3.579051in}{2.337363in}}%
\pgfpathlineto{\pgfqpoint{3.583592in}{2.337363in}}%
\pgfpathlineto{\pgfqpoint{3.583592in}{2.334414in}}%
\pgfpathmoveto{\pgfqpoint{3.579051in}{2.337363in}}%
\pgfpathlineto{\pgfqpoint{3.579051in}{2.337363in}}%
\pgfpathlineto{\pgfqpoint{3.579051in}{2.340313in}}%
\pgfpathlineto{\pgfqpoint{3.583592in}{2.340313in}}%
\pgfpathlineto{\pgfqpoint{3.583592in}{2.337363in}}%
\pgfpathmoveto{\pgfqpoint{3.583592in}{2.337363in}}%
\pgfpathlineto{\pgfqpoint{3.583592in}{2.337363in}}%
\pgfpathlineto{\pgfqpoint{3.583592in}{2.340313in}}%
\pgfpathlineto{\pgfqpoint{3.588133in}{2.340313in}}%
\pgfpathlineto{\pgfqpoint{3.588133in}{2.337363in}}%
\pgfpathmoveto{\pgfqpoint{3.583592in}{2.340313in}}%
\pgfpathlineto{\pgfqpoint{3.583592in}{2.340313in}}%
\pgfpathlineto{\pgfqpoint{3.583592in}{2.343262in}}%
\pgfpathlineto{\pgfqpoint{3.588133in}{2.343262in}}%
\pgfpathlineto{\pgfqpoint{3.588133in}{2.340313in}}%
\pgfpathmoveto{\pgfqpoint{3.588133in}{2.340313in}}%
\pgfpathlineto{\pgfqpoint{3.588133in}{2.340313in}}%
\pgfpathlineto{\pgfqpoint{3.588133in}{2.343262in}}%
\pgfpathlineto{\pgfqpoint{3.592674in}{2.343262in}}%
\pgfpathlineto{\pgfqpoint{3.592674in}{2.340313in}}%
\pgfpathmoveto{\pgfqpoint{3.588133in}{2.343262in}}%
\pgfpathlineto{\pgfqpoint{3.588133in}{2.343262in}}%
\pgfpathlineto{\pgfqpoint{3.588133in}{2.346211in}}%
\pgfpathlineto{\pgfqpoint{3.592674in}{2.346211in}}%
\pgfpathlineto{\pgfqpoint{3.592674in}{2.343262in}}%
\pgfpathmoveto{\pgfqpoint{3.592674in}{2.343262in}}%
\pgfpathlineto{\pgfqpoint{3.592674in}{2.343262in}}%
\pgfpathlineto{\pgfqpoint{3.592674in}{2.346211in}}%
\pgfpathlineto{\pgfqpoint{3.597215in}{2.346211in}}%
\pgfpathlineto{\pgfqpoint{3.597215in}{2.343262in}}%
\pgfpathmoveto{\pgfqpoint{3.592674in}{2.346211in}}%
\pgfpathlineto{\pgfqpoint{3.592674in}{2.346211in}}%
\pgfpathlineto{\pgfqpoint{3.592674in}{2.349161in}}%
\pgfpathlineto{\pgfqpoint{3.597215in}{2.349161in}}%
\pgfpathlineto{\pgfqpoint{3.597215in}{2.346211in}}%
\pgfpathmoveto{\pgfqpoint{3.597215in}{2.346211in}}%
\pgfpathlineto{\pgfqpoint{3.597215in}{2.346211in}}%
\pgfpathlineto{\pgfqpoint{3.597215in}{2.349161in}}%
\pgfpathlineto{\pgfqpoint{3.601756in}{2.349161in}}%
\pgfpathlineto{\pgfqpoint{3.601756in}{2.346211in}}%
\pgfpathmoveto{\pgfqpoint{3.597215in}{2.349161in}}%
\pgfpathlineto{\pgfqpoint{3.597215in}{2.349161in}}%
\pgfpathlineto{\pgfqpoint{3.597215in}{2.352110in}}%
\pgfpathlineto{\pgfqpoint{3.601756in}{2.352110in}}%
\pgfpathlineto{\pgfqpoint{3.601756in}{2.349161in}}%
\pgfpathmoveto{\pgfqpoint{3.601756in}{2.349161in}}%
\pgfpathlineto{\pgfqpoint{3.601756in}{2.349161in}}%
\pgfpathlineto{\pgfqpoint{3.601756in}{2.352110in}}%
\pgfpathlineto{\pgfqpoint{3.606297in}{2.352110in}}%
\pgfpathlineto{\pgfqpoint{3.606297in}{2.349161in}}%
\pgfpathmoveto{\pgfqpoint{3.601756in}{2.352110in}}%
\pgfpathlineto{\pgfqpoint{3.601756in}{2.352110in}}%
\pgfpathlineto{\pgfqpoint{3.601756in}{2.355059in}}%
\pgfpathlineto{\pgfqpoint{3.606297in}{2.355059in}}%
\pgfpathlineto{\pgfqpoint{3.606297in}{2.352110in}}%
\pgfpathmoveto{\pgfqpoint{3.606297in}{2.352110in}}%
\pgfpathlineto{\pgfqpoint{3.606297in}{2.352110in}}%
\pgfpathlineto{\pgfqpoint{3.606297in}{2.355059in}}%
\pgfpathlineto{\pgfqpoint{3.610838in}{2.355059in}}%
\pgfpathlineto{\pgfqpoint{3.610838in}{2.352110in}}%
\pgfpathmoveto{\pgfqpoint{3.606297in}{2.355059in}}%
\pgfpathlineto{\pgfqpoint{3.606297in}{2.355059in}}%
\pgfpathlineto{\pgfqpoint{3.606297in}{2.358009in}}%
\pgfpathlineto{\pgfqpoint{3.610838in}{2.358009in}}%
\pgfpathlineto{\pgfqpoint{3.610838in}{2.355059in}}%
\pgfpathmoveto{\pgfqpoint{3.610838in}{2.355059in}}%
\pgfpathlineto{\pgfqpoint{3.610838in}{2.355059in}}%
\pgfpathlineto{\pgfqpoint{3.610838in}{2.358009in}}%
\pgfpathlineto{\pgfqpoint{3.615379in}{2.358009in}}%
\pgfpathlineto{\pgfqpoint{3.615379in}{2.355059in}}%
\pgfpathmoveto{\pgfqpoint{3.610838in}{2.358009in}}%
\pgfpathlineto{\pgfqpoint{3.610838in}{2.358009in}}%
\pgfpathlineto{\pgfqpoint{3.610838in}{2.360958in}}%
\pgfpathlineto{\pgfqpoint{3.615379in}{2.360958in}}%
\pgfpathlineto{\pgfqpoint{3.615379in}{2.358009in}}%
\pgfpathmoveto{\pgfqpoint{3.615379in}{2.358009in}}%
\pgfpathlineto{\pgfqpoint{3.615379in}{2.358009in}}%
\pgfpathlineto{\pgfqpoint{3.615379in}{2.360958in}}%
\pgfpathlineto{\pgfqpoint{3.619920in}{2.360958in}}%
\pgfpathlineto{\pgfqpoint{3.619920in}{2.358009in}}%
\pgfpathmoveto{\pgfqpoint{3.615379in}{2.360958in}}%
\pgfpathlineto{\pgfqpoint{3.615379in}{2.360958in}}%
\pgfpathlineto{\pgfqpoint{3.615379in}{2.363907in}}%
\pgfpathlineto{\pgfqpoint{3.619920in}{2.363907in}}%
\pgfpathlineto{\pgfqpoint{3.619920in}{2.360958in}}%
\pgfpathmoveto{\pgfqpoint{3.619920in}{2.360958in}}%
\pgfpathlineto{\pgfqpoint{3.619920in}{2.360958in}}%
\pgfpathlineto{\pgfqpoint{3.619920in}{2.363907in}}%
\pgfpathlineto{\pgfqpoint{3.624461in}{2.363907in}}%
\pgfpathlineto{\pgfqpoint{3.624461in}{2.360958in}}%
\pgfpathmoveto{\pgfqpoint{3.619920in}{2.363907in}}%
\pgfpathlineto{\pgfqpoint{3.619920in}{2.363907in}}%
\pgfpathlineto{\pgfqpoint{3.619920in}{2.366857in}}%
\pgfpathlineto{\pgfqpoint{3.624461in}{2.366857in}}%
\pgfpathlineto{\pgfqpoint{3.624461in}{2.363907in}}%
\pgfpathmoveto{\pgfqpoint{3.624461in}{2.363907in}}%
\pgfpathlineto{\pgfqpoint{3.624461in}{2.363907in}}%
\pgfpathlineto{\pgfqpoint{3.624461in}{2.366857in}}%
\pgfpathlineto{\pgfqpoint{3.629002in}{2.366857in}}%
\pgfpathlineto{\pgfqpoint{3.629002in}{2.363907in}}%
\pgfpathmoveto{\pgfqpoint{3.624461in}{2.366857in}}%
\pgfpathlineto{\pgfqpoint{3.624461in}{2.366857in}}%
\pgfpathlineto{\pgfqpoint{3.624461in}{2.369806in}}%
\pgfpathlineto{\pgfqpoint{3.629002in}{2.369806in}}%
\pgfpathlineto{\pgfqpoint{3.629002in}{2.366857in}}%
\pgfpathmoveto{\pgfqpoint{3.629002in}{2.366857in}}%
\pgfpathlineto{\pgfqpoint{3.629002in}{2.366857in}}%
\pgfpathlineto{\pgfqpoint{3.629002in}{2.369806in}}%
\pgfpathlineto{\pgfqpoint{3.633543in}{2.369806in}}%
\pgfpathlineto{\pgfqpoint{3.633543in}{2.366857in}}%
\pgfpathmoveto{\pgfqpoint{3.629002in}{2.369806in}}%
\pgfpathlineto{\pgfqpoint{3.629002in}{2.369806in}}%
\pgfpathlineto{\pgfqpoint{3.629002in}{2.372755in}}%
\pgfpathlineto{\pgfqpoint{3.633543in}{2.372755in}}%
\pgfpathlineto{\pgfqpoint{3.633543in}{2.369806in}}%
\pgfpathmoveto{\pgfqpoint{3.633543in}{2.369806in}}%
\pgfpathlineto{\pgfqpoint{3.633543in}{2.369806in}}%
\pgfpathlineto{\pgfqpoint{3.633543in}{2.372755in}}%
\pgfpathlineto{\pgfqpoint{3.638084in}{2.372755in}}%
\pgfpathlineto{\pgfqpoint{3.638084in}{2.369806in}}%
\pgfpathmoveto{\pgfqpoint{3.633543in}{2.372755in}}%
\pgfpathlineto{\pgfqpoint{3.633543in}{2.372755in}}%
\pgfpathlineto{\pgfqpoint{3.633543in}{2.375704in}}%
\pgfpathlineto{\pgfqpoint{3.638084in}{2.375704in}}%
\pgfpathlineto{\pgfqpoint{3.638084in}{2.372755in}}%
\pgfpathmoveto{\pgfqpoint{3.638084in}{2.372755in}}%
\pgfpathlineto{\pgfqpoint{3.638084in}{2.372755in}}%
\pgfpathlineto{\pgfqpoint{3.638084in}{2.375704in}}%
\pgfpathlineto{\pgfqpoint{3.642625in}{2.375704in}}%
\pgfpathlineto{\pgfqpoint{3.642625in}{2.372755in}}%
\pgfpathmoveto{\pgfqpoint{3.638084in}{2.375704in}}%
\pgfpathlineto{\pgfqpoint{3.638084in}{2.375704in}}%
\pgfpathlineto{\pgfqpoint{3.638084in}{2.378654in}}%
\pgfpathlineto{\pgfqpoint{3.642625in}{2.378654in}}%
\pgfpathlineto{\pgfqpoint{3.642625in}{2.375704in}}%
\pgfpathmoveto{\pgfqpoint{3.642625in}{2.375704in}}%
\pgfpathlineto{\pgfqpoint{3.642625in}{2.375704in}}%
\pgfpathlineto{\pgfqpoint{3.642625in}{2.378654in}}%
\pgfpathlineto{\pgfqpoint{3.647167in}{2.378654in}}%
\pgfpathlineto{\pgfqpoint{3.647167in}{2.375704in}}%
\pgfpathmoveto{\pgfqpoint{3.642625in}{2.378654in}}%
\pgfpathlineto{\pgfqpoint{3.642625in}{2.378654in}}%
\pgfpathlineto{\pgfqpoint{3.642625in}{2.381603in}}%
\pgfpathlineto{\pgfqpoint{3.647167in}{2.381603in}}%
\pgfpathlineto{\pgfqpoint{3.647167in}{2.378654in}}%
\pgfpathmoveto{\pgfqpoint{3.647167in}{2.378654in}}%
\pgfpathlineto{\pgfqpoint{3.647167in}{2.378654in}}%
\pgfpathlineto{\pgfqpoint{3.647167in}{2.381603in}}%
\pgfpathlineto{\pgfqpoint{3.651708in}{2.381603in}}%
\pgfpathlineto{\pgfqpoint{3.651708in}{2.378654in}}%
\pgfpathmoveto{\pgfqpoint{3.647167in}{2.381603in}}%
\pgfpathlineto{\pgfqpoint{3.647167in}{2.381603in}}%
\pgfpathlineto{\pgfqpoint{3.647167in}{2.384552in}}%
\pgfpathlineto{\pgfqpoint{3.651708in}{2.384552in}}%
\pgfpathlineto{\pgfqpoint{3.651708in}{2.381603in}}%
\pgfpathmoveto{\pgfqpoint{3.651708in}{2.381603in}}%
\pgfpathlineto{\pgfqpoint{3.651708in}{2.381603in}}%
\pgfpathlineto{\pgfqpoint{3.651708in}{2.384552in}}%
\pgfpathlineto{\pgfqpoint{3.656249in}{2.384552in}}%
\pgfpathlineto{\pgfqpoint{3.656249in}{2.381603in}}%
\pgfpathmoveto{\pgfqpoint{3.651708in}{2.384552in}}%
\pgfpathlineto{\pgfqpoint{3.651708in}{2.384552in}}%
\pgfpathlineto{\pgfqpoint{3.651708in}{2.387502in}}%
\pgfpathlineto{\pgfqpoint{3.656249in}{2.387502in}}%
\pgfpathlineto{\pgfqpoint{3.656249in}{2.384552in}}%
\pgfpathmoveto{\pgfqpoint{3.656249in}{2.384552in}}%
\pgfpathlineto{\pgfqpoint{3.656249in}{2.384552in}}%
\pgfpathlineto{\pgfqpoint{3.656249in}{2.387502in}}%
\pgfpathlineto{\pgfqpoint{3.660790in}{2.387502in}}%
\pgfpathlineto{\pgfqpoint{3.660790in}{2.384552in}}%
\pgfpathmoveto{\pgfqpoint{3.656249in}{2.387502in}}%
\pgfpathlineto{\pgfqpoint{3.656249in}{2.387502in}}%
\pgfpathlineto{\pgfqpoint{3.656249in}{2.390451in}}%
\pgfpathlineto{\pgfqpoint{3.660790in}{2.390451in}}%
\pgfpathlineto{\pgfqpoint{3.660790in}{2.387502in}}%
\pgfpathmoveto{\pgfqpoint{3.660790in}{2.387502in}}%
\pgfpathlineto{\pgfqpoint{3.660790in}{2.387502in}}%
\pgfpathlineto{\pgfqpoint{3.660790in}{2.390451in}}%
\pgfpathlineto{\pgfqpoint{3.665331in}{2.390451in}}%
\pgfpathlineto{\pgfqpoint{3.665331in}{2.387502in}}%
\pgfpathmoveto{\pgfqpoint{3.660790in}{2.390451in}}%
\pgfpathlineto{\pgfqpoint{3.660790in}{2.390451in}}%
\pgfpathlineto{\pgfqpoint{3.660790in}{2.393400in}}%
\pgfpathlineto{\pgfqpoint{3.665331in}{2.393400in}}%
\pgfpathlineto{\pgfqpoint{3.665331in}{2.390451in}}%
\pgfpathmoveto{\pgfqpoint{3.665331in}{2.390451in}}%
\pgfpathlineto{\pgfqpoint{3.665331in}{2.390451in}}%
\pgfpathlineto{\pgfqpoint{3.665331in}{2.393400in}}%
\pgfpathlineto{\pgfqpoint{3.669872in}{2.393400in}}%
\pgfpathlineto{\pgfqpoint{3.669872in}{2.390451in}}%
\pgfpathmoveto{\pgfqpoint{3.665331in}{2.393400in}}%
\pgfpathlineto{\pgfqpoint{3.665331in}{2.393400in}}%
\pgfpathlineto{\pgfqpoint{3.665331in}{2.396349in}}%
\pgfpathlineto{\pgfqpoint{3.669872in}{2.396349in}}%
\pgfpathlineto{\pgfqpoint{3.669872in}{2.393400in}}%
\pgfpathmoveto{\pgfqpoint{3.669872in}{2.393400in}}%
\pgfpathlineto{\pgfqpoint{3.669872in}{2.393400in}}%
\pgfpathlineto{\pgfqpoint{3.669872in}{2.396349in}}%
\pgfpathlineto{\pgfqpoint{3.674413in}{2.396349in}}%
\pgfpathlineto{\pgfqpoint{3.674413in}{2.393400in}}%
\pgfpathmoveto{\pgfqpoint{3.669872in}{2.396349in}}%
\pgfpathlineto{\pgfqpoint{3.669872in}{2.396349in}}%
\pgfpathlineto{\pgfqpoint{3.669872in}{2.399298in}}%
\pgfpathlineto{\pgfqpoint{3.674413in}{2.399298in}}%
\pgfpathlineto{\pgfqpoint{3.674413in}{2.396349in}}%
\pgfpathmoveto{\pgfqpoint{3.674413in}{2.396349in}}%
\pgfpathlineto{\pgfqpoint{3.674413in}{2.396349in}}%
\pgfpathlineto{\pgfqpoint{3.674413in}{2.399298in}}%
\pgfpathlineto{\pgfqpoint{3.678954in}{2.399298in}}%
\pgfpathlineto{\pgfqpoint{3.678954in}{2.396349in}}%
\pgfpathmoveto{\pgfqpoint{3.674413in}{2.399298in}}%
\pgfpathlineto{\pgfqpoint{3.674413in}{2.399298in}}%
\pgfpathlineto{\pgfqpoint{3.674413in}{2.402247in}}%
\pgfpathlineto{\pgfqpoint{3.678954in}{2.402247in}}%
\pgfpathlineto{\pgfqpoint{3.678954in}{2.399298in}}%
\pgfpathmoveto{\pgfqpoint{3.678954in}{2.399298in}}%
\pgfpathlineto{\pgfqpoint{3.678954in}{2.399298in}}%
\pgfpathlineto{\pgfqpoint{3.678954in}{2.402247in}}%
\pgfpathlineto{\pgfqpoint{3.683495in}{2.402247in}}%
\pgfpathlineto{\pgfqpoint{3.683495in}{2.399298in}}%
\pgfpathmoveto{\pgfqpoint{3.678954in}{2.402247in}}%
\pgfpathlineto{\pgfqpoint{3.678954in}{2.402247in}}%
\pgfpathlineto{\pgfqpoint{3.678954in}{2.405196in}}%
\pgfpathlineto{\pgfqpoint{3.683495in}{2.405196in}}%
\pgfpathlineto{\pgfqpoint{3.683495in}{2.402247in}}%
\pgfpathmoveto{\pgfqpoint{3.683495in}{2.402247in}}%
\pgfpathlineto{\pgfqpoint{3.683495in}{2.402247in}}%
\pgfpathlineto{\pgfqpoint{3.683495in}{2.405196in}}%
\pgfpathlineto{\pgfqpoint{3.688036in}{2.405196in}}%
\pgfpathlineto{\pgfqpoint{3.688036in}{2.402247in}}%
\pgfpathmoveto{\pgfqpoint{3.683495in}{2.405196in}}%
\pgfpathlineto{\pgfqpoint{3.683495in}{2.405196in}}%
\pgfpathlineto{\pgfqpoint{3.683495in}{2.408146in}}%
\pgfpathlineto{\pgfqpoint{3.688036in}{2.408146in}}%
\pgfpathlineto{\pgfqpoint{3.688036in}{2.405196in}}%
\pgfpathmoveto{\pgfqpoint{3.688036in}{2.405196in}}%
\pgfpathlineto{\pgfqpoint{3.688036in}{2.405196in}}%
\pgfpathlineto{\pgfqpoint{3.688036in}{2.408146in}}%
\pgfpathlineto{\pgfqpoint{3.692577in}{2.408146in}}%
\pgfpathlineto{\pgfqpoint{3.692577in}{2.405196in}}%
\pgfpathmoveto{\pgfqpoint{3.688036in}{2.408146in}}%
\pgfpathlineto{\pgfqpoint{3.688036in}{2.408146in}}%
\pgfpathlineto{\pgfqpoint{3.688036in}{2.411095in}}%
\pgfpathlineto{\pgfqpoint{3.692577in}{2.411095in}}%
\pgfpathlineto{\pgfqpoint{3.692577in}{2.408146in}}%
\pgfpathmoveto{\pgfqpoint{3.692577in}{2.408146in}}%
\pgfpathlineto{\pgfqpoint{3.692577in}{2.408146in}}%
\pgfpathlineto{\pgfqpoint{3.692577in}{2.411095in}}%
\pgfpathlineto{\pgfqpoint{3.697118in}{2.411095in}}%
\pgfpathlineto{\pgfqpoint{3.697118in}{2.408146in}}%
\pgfpathmoveto{\pgfqpoint{3.692577in}{2.411095in}}%
\pgfpathlineto{\pgfqpoint{3.692577in}{2.411095in}}%
\pgfpathlineto{\pgfqpoint{3.692577in}{2.414044in}}%
\pgfpathlineto{\pgfqpoint{3.697118in}{2.414044in}}%
\pgfpathlineto{\pgfqpoint{3.697118in}{2.411095in}}%
\pgfpathmoveto{\pgfqpoint{3.697118in}{2.411095in}}%
\pgfpathlineto{\pgfqpoint{3.697118in}{2.411095in}}%
\pgfpathlineto{\pgfqpoint{3.697118in}{2.414044in}}%
\pgfpathlineto{\pgfqpoint{3.701659in}{2.414044in}}%
\pgfpathlineto{\pgfqpoint{3.701659in}{2.411095in}}%
\pgfpathmoveto{\pgfqpoint{3.697118in}{2.414044in}}%
\pgfpathlineto{\pgfqpoint{3.697118in}{2.414044in}}%
\pgfpathlineto{\pgfqpoint{3.697118in}{2.416993in}}%
\pgfpathlineto{\pgfqpoint{3.701659in}{2.416993in}}%
\pgfpathlineto{\pgfqpoint{3.701659in}{2.414044in}}%
\pgfpathmoveto{\pgfqpoint{3.701659in}{2.414044in}}%
\pgfpathlineto{\pgfqpoint{3.701659in}{2.414044in}}%
\pgfpathlineto{\pgfqpoint{3.701659in}{2.416993in}}%
\pgfpathlineto{\pgfqpoint{3.706201in}{2.416993in}}%
\pgfpathlineto{\pgfqpoint{3.706201in}{2.414044in}}%
\pgfpathmoveto{\pgfqpoint{3.701659in}{2.416993in}}%
\pgfpathlineto{\pgfqpoint{3.701659in}{2.416993in}}%
\pgfpathlineto{\pgfqpoint{3.701659in}{2.419942in}}%
\pgfpathlineto{\pgfqpoint{3.706201in}{2.419942in}}%
\pgfpathlineto{\pgfqpoint{3.706201in}{2.416993in}}%
\pgfpathmoveto{\pgfqpoint{3.706201in}{2.416993in}}%
\pgfpathlineto{\pgfqpoint{3.706201in}{2.416993in}}%
\pgfpathlineto{\pgfqpoint{3.706201in}{2.419942in}}%
\pgfpathlineto{\pgfqpoint{3.710742in}{2.419942in}}%
\pgfpathlineto{\pgfqpoint{3.710742in}{2.416993in}}%
\pgfpathmoveto{\pgfqpoint{3.706201in}{2.419942in}}%
\pgfpathlineto{\pgfqpoint{3.706201in}{2.419942in}}%
\pgfpathlineto{\pgfqpoint{3.706201in}{2.422891in}}%
\pgfpathlineto{\pgfqpoint{3.710742in}{2.422891in}}%
\pgfpathlineto{\pgfqpoint{3.710742in}{2.419942in}}%
\pgfpathmoveto{\pgfqpoint{3.710742in}{2.419942in}}%
\pgfpathlineto{\pgfqpoint{3.710742in}{2.419942in}}%
\pgfpathlineto{\pgfqpoint{3.710742in}{2.422891in}}%
\pgfpathlineto{\pgfqpoint{3.715283in}{2.422891in}}%
\pgfpathlineto{\pgfqpoint{3.715283in}{2.419942in}}%
\pgfpathmoveto{\pgfqpoint{3.710742in}{2.422891in}}%
\pgfpathlineto{\pgfqpoint{3.710742in}{2.422891in}}%
\pgfpathlineto{\pgfqpoint{3.710742in}{2.425840in}}%
\pgfpathlineto{\pgfqpoint{3.715283in}{2.425840in}}%
\pgfpathlineto{\pgfqpoint{3.715283in}{2.422891in}}%
\pgfpathmoveto{\pgfqpoint{3.715283in}{2.422891in}}%
\pgfpathlineto{\pgfqpoint{3.715283in}{2.422891in}}%
\pgfpathlineto{\pgfqpoint{3.715283in}{2.425840in}}%
\pgfpathlineto{\pgfqpoint{3.719824in}{2.425840in}}%
\pgfpathlineto{\pgfqpoint{3.719824in}{2.422891in}}%
\pgfpathmoveto{\pgfqpoint{3.715283in}{2.425840in}}%
\pgfpathlineto{\pgfqpoint{3.715283in}{2.425840in}}%
\pgfpathlineto{\pgfqpoint{3.715283in}{2.428789in}}%
\pgfpathlineto{\pgfqpoint{3.719824in}{2.428789in}}%
\pgfpathlineto{\pgfqpoint{3.719824in}{2.425840in}}%
\pgfpathmoveto{\pgfqpoint{3.719824in}{2.425840in}}%
\pgfpathlineto{\pgfqpoint{3.719824in}{2.425840in}}%
\pgfpathlineto{\pgfqpoint{3.719824in}{2.428789in}}%
\pgfpathlineto{\pgfqpoint{3.724365in}{2.428789in}}%
\pgfpathlineto{\pgfqpoint{3.724365in}{2.425840in}}%
\pgfpathmoveto{\pgfqpoint{3.719824in}{2.428789in}}%
\pgfpathlineto{\pgfqpoint{3.719824in}{2.428789in}}%
\pgfpathlineto{\pgfqpoint{3.719824in}{2.431738in}}%
\pgfpathlineto{\pgfqpoint{3.724365in}{2.431738in}}%
\pgfpathlineto{\pgfqpoint{3.724365in}{2.428789in}}%
\pgfpathmoveto{\pgfqpoint{3.724365in}{2.428789in}}%
\pgfpathlineto{\pgfqpoint{3.724365in}{2.428789in}}%
\pgfpathlineto{\pgfqpoint{3.724365in}{2.431738in}}%
\pgfpathlineto{\pgfqpoint{3.728906in}{2.431738in}}%
\pgfpathlineto{\pgfqpoint{3.728906in}{2.428789in}}%
\pgfpathmoveto{\pgfqpoint{3.724365in}{2.431738in}}%
\pgfpathlineto{\pgfqpoint{3.724365in}{2.431738in}}%
\pgfpathlineto{\pgfqpoint{3.724365in}{2.434688in}}%
\pgfpathlineto{\pgfqpoint{3.728906in}{2.434688in}}%
\pgfpathlineto{\pgfqpoint{3.728906in}{2.431738in}}%
\pgfpathmoveto{\pgfqpoint{3.728906in}{2.431738in}}%
\pgfpathlineto{\pgfqpoint{3.728906in}{2.431738in}}%
\pgfpathlineto{\pgfqpoint{3.728906in}{2.434688in}}%
\pgfpathlineto{\pgfqpoint{3.733447in}{2.434688in}}%
\pgfpathlineto{\pgfqpoint{3.733447in}{2.431738in}}%
\pgfpathmoveto{\pgfqpoint{3.728906in}{2.434688in}}%
\pgfpathlineto{\pgfqpoint{3.728906in}{2.434688in}}%
\pgfpathlineto{\pgfqpoint{3.728906in}{2.437637in}}%
\pgfpathlineto{\pgfqpoint{3.733447in}{2.437637in}}%
\pgfpathlineto{\pgfqpoint{3.733447in}{2.434688in}}%
\pgfpathmoveto{\pgfqpoint{3.733447in}{2.434688in}}%
\pgfpathlineto{\pgfqpoint{3.733447in}{2.434688in}}%
\pgfpathlineto{\pgfqpoint{3.733447in}{2.437637in}}%
\pgfpathlineto{\pgfqpoint{3.737988in}{2.437637in}}%
\pgfpathlineto{\pgfqpoint{3.737988in}{2.434688in}}%
\pgfpathmoveto{\pgfqpoint{3.733447in}{2.437637in}}%
\pgfpathlineto{\pgfqpoint{3.733447in}{2.437637in}}%
\pgfpathlineto{\pgfqpoint{3.733447in}{2.440586in}}%
\pgfpathlineto{\pgfqpoint{3.737988in}{2.440586in}}%
\pgfpathlineto{\pgfqpoint{3.737988in}{2.437637in}}%
\pgfpathmoveto{\pgfqpoint{3.733447in}{2.440586in}}%
\pgfpathlineto{\pgfqpoint{3.733447in}{2.440586in}}%
\pgfpathlineto{\pgfqpoint{3.733447in}{2.443535in}}%
\pgfpathlineto{\pgfqpoint{3.737988in}{2.443535in}}%
\pgfpathlineto{\pgfqpoint{3.737988in}{2.440586in}}%
\pgfpathmoveto{\pgfqpoint{3.737988in}{2.440586in}}%
\pgfpathlineto{\pgfqpoint{3.737988in}{2.440586in}}%
\pgfpathlineto{\pgfqpoint{3.737988in}{2.443535in}}%
\pgfpathlineto{\pgfqpoint{3.742529in}{2.443535in}}%
\pgfpathlineto{\pgfqpoint{3.742529in}{2.440586in}}%
\pgfpathmoveto{\pgfqpoint{3.737988in}{2.443535in}}%
\pgfpathlineto{\pgfqpoint{3.737988in}{2.443535in}}%
\pgfpathlineto{\pgfqpoint{3.737988in}{2.446484in}}%
\pgfpathlineto{\pgfqpoint{3.742529in}{2.446484in}}%
\pgfpathlineto{\pgfqpoint{3.742529in}{2.443535in}}%
\pgfpathmoveto{\pgfqpoint{3.742529in}{2.443535in}}%
\pgfpathlineto{\pgfqpoint{3.742529in}{2.443535in}}%
\pgfpathlineto{\pgfqpoint{3.742529in}{2.446484in}}%
\pgfpathlineto{\pgfqpoint{3.747070in}{2.446484in}}%
\pgfpathlineto{\pgfqpoint{3.747070in}{2.443535in}}%
\pgfpathmoveto{\pgfqpoint{3.742529in}{2.446484in}}%
\pgfpathlineto{\pgfqpoint{3.742529in}{2.446484in}}%
\pgfpathlineto{\pgfqpoint{3.742529in}{2.449433in}}%
\pgfpathlineto{\pgfqpoint{3.747070in}{2.449433in}}%
\pgfpathlineto{\pgfqpoint{3.747070in}{2.446484in}}%
\pgfpathmoveto{\pgfqpoint{3.747070in}{2.446484in}}%
\pgfpathlineto{\pgfqpoint{3.747070in}{2.446484in}}%
\pgfpathlineto{\pgfqpoint{3.747070in}{2.449433in}}%
\pgfpathlineto{\pgfqpoint{3.751611in}{2.449433in}}%
\pgfpathlineto{\pgfqpoint{3.751611in}{2.446484in}}%
\pgfpathmoveto{\pgfqpoint{3.747070in}{2.449433in}}%
\pgfpathlineto{\pgfqpoint{3.747070in}{2.449433in}}%
\pgfpathlineto{\pgfqpoint{3.747070in}{2.452382in}}%
\pgfpathlineto{\pgfqpoint{3.751611in}{2.452382in}}%
\pgfpathlineto{\pgfqpoint{3.751611in}{2.449433in}}%
\pgfpathmoveto{\pgfqpoint{3.751611in}{2.449433in}}%
\pgfpathlineto{\pgfqpoint{3.751611in}{2.449433in}}%
\pgfpathlineto{\pgfqpoint{3.751611in}{2.452382in}}%
\pgfpathlineto{\pgfqpoint{3.756152in}{2.452382in}}%
\pgfpathlineto{\pgfqpoint{3.756152in}{2.449433in}}%
\pgfpathmoveto{\pgfqpoint{3.751611in}{2.452382in}}%
\pgfpathlineto{\pgfqpoint{3.751611in}{2.452382in}}%
\pgfpathlineto{\pgfqpoint{3.751611in}{2.455331in}}%
\pgfpathlineto{\pgfqpoint{3.756152in}{2.455331in}}%
\pgfpathlineto{\pgfqpoint{3.756152in}{2.452382in}}%
\pgfpathmoveto{\pgfqpoint{3.756152in}{2.452382in}}%
\pgfpathlineto{\pgfqpoint{3.756152in}{2.452382in}}%
\pgfpathlineto{\pgfqpoint{3.756152in}{2.455331in}}%
\pgfpathlineto{\pgfqpoint{3.760693in}{2.455331in}}%
\pgfpathlineto{\pgfqpoint{3.760693in}{2.452382in}}%
\pgfpathmoveto{\pgfqpoint{3.756152in}{2.455331in}}%
\pgfpathlineto{\pgfqpoint{3.756152in}{2.455331in}}%
\pgfpathlineto{\pgfqpoint{3.756152in}{2.458280in}}%
\pgfpathlineto{\pgfqpoint{3.760693in}{2.458280in}}%
\pgfpathlineto{\pgfqpoint{3.760693in}{2.455331in}}%
\pgfpathmoveto{\pgfqpoint{3.760693in}{2.455331in}}%
\pgfpathlineto{\pgfqpoint{3.760693in}{2.455331in}}%
\pgfpathlineto{\pgfqpoint{3.760693in}{2.458280in}}%
\pgfpathlineto{\pgfqpoint{3.765235in}{2.458280in}}%
\pgfpathlineto{\pgfqpoint{3.765235in}{2.455331in}}%
\pgfpathmoveto{\pgfqpoint{3.760693in}{2.458280in}}%
\pgfpathlineto{\pgfqpoint{3.760693in}{2.458280in}}%
\pgfpathlineto{\pgfqpoint{3.760693in}{2.461230in}}%
\pgfpathlineto{\pgfqpoint{3.765235in}{2.461230in}}%
\pgfpathlineto{\pgfqpoint{3.765235in}{2.458280in}}%
\pgfpathmoveto{\pgfqpoint{3.765235in}{2.458280in}}%
\pgfpathlineto{\pgfqpoint{3.765235in}{2.458280in}}%
\pgfpathlineto{\pgfqpoint{3.765235in}{2.461230in}}%
\pgfpathlineto{\pgfqpoint{3.769776in}{2.461230in}}%
\pgfpathlineto{\pgfqpoint{3.769776in}{2.458280in}}%
\pgfpathmoveto{\pgfqpoint{3.765235in}{2.461230in}}%
\pgfpathlineto{\pgfqpoint{3.765235in}{2.461230in}}%
\pgfpathlineto{\pgfqpoint{3.765235in}{2.464179in}}%
\pgfpathlineto{\pgfqpoint{3.769776in}{2.464179in}}%
\pgfpathlineto{\pgfqpoint{3.769776in}{2.461230in}}%
\pgfpathmoveto{\pgfqpoint{3.769776in}{2.461230in}}%
\pgfpathlineto{\pgfqpoint{3.769776in}{2.461230in}}%
\pgfpathlineto{\pgfqpoint{3.769776in}{2.464179in}}%
\pgfpathlineto{\pgfqpoint{3.774317in}{2.464179in}}%
\pgfpathlineto{\pgfqpoint{3.774317in}{2.461230in}}%
\pgfpathmoveto{\pgfqpoint{3.769776in}{2.464179in}}%
\pgfpathlineto{\pgfqpoint{3.769776in}{2.464179in}}%
\pgfpathlineto{\pgfqpoint{3.769776in}{2.467128in}}%
\pgfpathlineto{\pgfqpoint{3.774317in}{2.467128in}}%
\pgfpathlineto{\pgfqpoint{3.774317in}{2.464179in}}%
\pgfpathmoveto{\pgfqpoint{3.774317in}{2.464179in}}%
\pgfpathlineto{\pgfqpoint{3.774317in}{2.464179in}}%
\pgfpathlineto{\pgfqpoint{3.774317in}{2.467128in}}%
\pgfpathlineto{\pgfqpoint{3.778858in}{2.467128in}}%
\pgfpathlineto{\pgfqpoint{3.778858in}{2.464179in}}%
\pgfpathmoveto{\pgfqpoint{3.774317in}{2.467128in}}%
\pgfpathlineto{\pgfqpoint{3.774317in}{2.467128in}}%
\pgfpathlineto{\pgfqpoint{3.774317in}{2.470077in}}%
\pgfpathlineto{\pgfqpoint{3.778858in}{2.470077in}}%
\pgfpathlineto{\pgfqpoint{3.778858in}{2.467128in}}%
\pgfpathmoveto{\pgfqpoint{3.778858in}{2.467128in}}%
\pgfpathlineto{\pgfqpoint{3.778858in}{2.467128in}}%
\pgfpathlineto{\pgfqpoint{3.778858in}{2.470077in}}%
\pgfpathlineto{\pgfqpoint{3.783399in}{2.470077in}}%
\pgfpathlineto{\pgfqpoint{3.783399in}{2.467128in}}%
\pgfpathmoveto{\pgfqpoint{3.778858in}{2.470077in}}%
\pgfpathlineto{\pgfqpoint{3.778858in}{2.470077in}}%
\pgfpathlineto{\pgfqpoint{3.778858in}{2.473026in}}%
\pgfpathlineto{\pgfqpoint{3.783399in}{2.473026in}}%
\pgfpathlineto{\pgfqpoint{3.783399in}{2.470077in}}%
\pgfpathmoveto{\pgfqpoint{3.783399in}{2.470077in}}%
\pgfpathlineto{\pgfqpoint{3.783399in}{2.470077in}}%
\pgfpathlineto{\pgfqpoint{3.783399in}{2.473026in}}%
\pgfpathlineto{\pgfqpoint{3.787940in}{2.473026in}}%
\pgfpathlineto{\pgfqpoint{3.787940in}{2.470077in}}%
\pgfpathmoveto{\pgfqpoint{3.783399in}{2.473026in}}%
\pgfpathlineto{\pgfqpoint{3.783399in}{2.473026in}}%
\pgfpathlineto{\pgfqpoint{3.783399in}{2.475975in}}%
\pgfpathlineto{\pgfqpoint{3.787940in}{2.475975in}}%
\pgfpathlineto{\pgfqpoint{3.787940in}{2.473026in}}%
\pgfpathmoveto{\pgfqpoint{3.787940in}{2.473026in}}%
\pgfpathlineto{\pgfqpoint{3.787940in}{2.473026in}}%
\pgfpathlineto{\pgfqpoint{3.787940in}{2.475975in}}%
\pgfpathlineto{\pgfqpoint{3.792481in}{2.475975in}}%
\pgfpathlineto{\pgfqpoint{3.792481in}{2.473026in}}%
\pgfpathmoveto{\pgfqpoint{3.787940in}{2.475975in}}%
\pgfpathlineto{\pgfqpoint{3.787940in}{2.475975in}}%
\pgfpathlineto{\pgfqpoint{3.787940in}{2.478924in}}%
\pgfpathlineto{\pgfqpoint{3.792481in}{2.478924in}}%
\pgfpathlineto{\pgfqpoint{3.792481in}{2.475975in}}%
\pgfpathmoveto{\pgfqpoint{3.792481in}{2.475975in}}%
\pgfpathlineto{\pgfqpoint{3.792481in}{2.475975in}}%
\pgfpathlineto{\pgfqpoint{3.792481in}{2.478924in}}%
\pgfpathlineto{\pgfqpoint{3.797022in}{2.478924in}}%
\pgfpathlineto{\pgfqpoint{3.797022in}{2.475975in}}%
\pgfpathmoveto{\pgfqpoint{3.792481in}{2.478924in}}%
\pgfpathlineto{\pgfqpoint{3.792481in}{2.478924in}}%
\pgfpathlineto{\pgfqpoint{3.792481in}{2.481873in}}%
\pgfpathlineto{\pgfqpoint{3.797022in}{2.481873in}}%
\pgfpathlineto{\pgfqpoint{3.797022in}{2.478924in}}%
\pgfpathmoveto{\pgfqpoint{3.797022in}{2.478924in}}%
\pgfpathlineto{\pgfqpoint{3.797022in}{2.478924in}}%
\pgfpathlineto{\pgfqpoint{3.797022in}{2.481873in}}%
\pgfpathlineto{\pgfqpoint{3.801563in}{2.481873in}}%
\pgfpathlineto{\pgfqpoint{3.801563in}{2.478924in}}%
\pgfpathmoveto{\pgfqpoint{3.797022in}{2.481873in}}%
\pgfpathlineto{\pgfqpoint{3.797022in}{2.481873in}}%
\pgfpathlineto{\pgfqpoint{3.797022in}{2.484823in}}%
\pgfpathlineto{\pgfqpoint{3.801563in}{2.484823in}}%
\pgfpathlineto{\pgfqpoint{3.801563in}{2.481873in}}%
\pgfpathmoveto{\pgfqpoint{3.801563in}{2.481873in}}%
\pgfpathlineto{\pgfqpoint{3.801563in}{2.481873in}}%
\pgfpathlineto{\pgfqpoint{3.801563in}{2.484823in}}%
\pgfpathlineto{\pgfqpoint{3.806104in}{2.484823in}}%
\pgfpathlineto{\pgfqpoint{3.806104in}{2.481873in}}%
\pgfpathmoveto{\pgfqpoint{3.801563in}{2.484823in}}%
\pgfpathlineto{\pgfqpoint{3.801563in}{2.484823in}}%
\pgfpathlineto{\pgfqpoint{3.801563in}{2.487772in}}%
\pgfpathlineto{\pgfqpoint{3.806104in}{2.487772in}}%
\pgfpathlineto{\pgfqpoint{3.806104in}{2.484823in}}%
\pgfpathmoveto{\pgfqpoint{3.806104in}{2.484823in}}%
\pgfpathlineto{\pgfqpoint{3.806104in}{2.484823in}}%
\pgfpathlineto{\pgfqpoint{3.806104in}{2.487772in}}%
\pgfpathlineto{\pgfqpoint{3.810645in}{2.487772in}}%
\pgfpathlineto{\pgfqpoint{3.810645in}{2.484823in}}%
\pgfpathmoveto{\pgfqpoint{3.806104in}{2.487772in}}%
\pgfpathlineto{\pgfqpoint{3.806104in}{2.487772in}}%
\pgfpathlineto{\pgfqpoint{3.806104in}{2.490721in}}%
\pgfpathlineto{\pgfqpoint{3.810645in}{2.490721in}}%
\pgfpathlineto{\pgfqpoint{3.810645in}{2.487772in}}%
\pgfpathmoveto{\pgfqpoint{3.810645in}{2.487772in}}%
\pgfpathlineto{\pgfqpoint{3.810645in}{2.487772in}}%
\pgfpathlineto{\pgfqpoint{3.810645in}{2.490721in}}%
\pgfpathlineto{\pgfqpoint{3.815186in}{2.490721in}}%
\pgfpathlineto{\pgfqpoint{3.815186in}{2.487772in}}%
\pgfpathmoveto{\pgfqpoint{3.810645in}{2.490721in}}%
\pgfpathlineto{\pgfqpoint{3.810645in}{2.490721in}}%
\pgfpathlineto{\pgfqpoint{3.810645in}{2.493671in}}%
\pgfpathlineto{\pgfqpoint{3.815186in}{2.493671in}}%
\pgfpathlineto{\pgfqpoint{3.815186in}{2.490721in}}%
\pgfpathmoveto{\pgfqpoint{3.815186in}{2.490721in}}%
\pgfpathlineto{\pgfqpoint{3.815186in}{2.490721in}}%
\pgfpathlineto{\pgfqpoint{3.815186in}{2.493671in}}%
\pgfpathlineto{\pgfqpoint{3.819727in}{2.493671in}}%
\pgfpathlineto{\pgfqpoint{3.819727in}{2.490721in}}%
\pgfpathmoveto{\pgfqpoint{3.815186in}{2.493671in}}%
\pgfpathlineto{\pgfqpoint{3.815186in}{2.493671in}}%
\pgfpathlineto{\pgfqpoint{3.815186in}{2.496620in}}%
\pgfpathlineto{\pgfqpoint{3.819727in}{2.496620in}}%
\pgfpathlineto{\pgfqpoint{3.819727in}{2.493671in}}%
\pgfpathmoveto{\pgfqpoint{3.819727in}{2.493671in}}%
\pgfpathlineto{\pgfqpoint{3.819727in}{2.493671in}}%
\pgfpathlineto{\pgfqpoint{3.819727in}{2.496620in}}%
\pgfpathlineto{\pgfqpoint{3.824268in}{2.496620in}}%
\pgfpathlineto{\pgfqpoint{3.824268in}{2.493671in}}%
\pgfpathmoveto{\pgfqpoint{3.819727in}{2.496620in}}%
\pgfpathlineto{\pgfqpoint{3.819727in}{2.496620in}}%
\pgfpathlineto{\pgfqpoint{3.819727in}{2.499569in}}%
\pgfpathlineto{\pgfqpoint{3.824268in}{2.499569in}}%
\pgfpathlineto{\pgfqpoint{3.824268in}{2.496620in}}%
\pgfpathmoveto{\pgfqpoint{3.824268in}{2.496620in}}%
\pgfpathlineto{\pgfqpoint{3.824268in}{2.496620in}}%
\pgfpathlineto{\pgfqpoint{3.824268in}{2.499569in}}%
\pgfpathlineto{\pgfqpoint{3.828809in}{2.499569in}}%
\pgfpathlineto{\pgfqpoint{3.828809in}{2.496620in}}%
\pgfpathmoveto{\pgfqpoint{3.824268in}{2.499569in}}%
\pgfpathlineto{\pgfqpoint{3.824268in}{2.499569in}}%
\pgfpathlineto{\pgfqpoint{3.824268in}{2.502518in}}%
\pgfpathlineto{\pgfqpoint{3.828809in}{2.502518in}}%
\pgfpathlineto{\pgfqpoint{3.828809in}{2.499569in}}%
\pgfpathmoveto{\pgfqpoint{3.828809in}{2.499569in}}%
\pgfpathlineto{\pgfqpoint{3.828809in}{2.499569in}}%
\pgfpathlineto{\pgfqpoint{3.828809in}{2.502518in}}%
\pgfpathlineto{\pgfqpoint{3.833350in}{2.502518in}}%
\pgfpathlineto{\pgfqpoint{3.833350in}{2.499569in}}%
\pgfpathmoveto{\pgfqpoint{3.828809in}{2.502518in}}%
\pgfpathlineto{\pgfqpoint{3.828809in}{2.502518in}}%
\pgfpathlineto{\pgfqpoint{3.828809in}{2.505468in}}%
\pgfpathlineto{\pgfqpoint{3.833350in}{2.505468in}}%
\pgfpathlineto{\pgfqpoint{3.833350in}{2.502518in}}%
\pgfpathmoveto{\pgfqpoint{3.833350in}{2.502518in}}%
\pgfpathlineto{\pgfqpoint{3.833350in}{2.502518in}}%
\pgfpathlineto{\pgfqpoint{3.833350in}{2.505468in}}%
\pgfpathlineto{\pgfqpoint{3.837890in}{2.505468in}}%
\pgfpathlineto{\pgfqpoint{3.837890in}{2.502518in}}%
\pgfpathmoveto{\pgfqpoint{3.833350in}{2.505468in}}%
\pgfpathlineto{\pgfqpoint{3.833350in}{2.505468in}}%
\pgfpathlineto{\pgfqpoint{3.833350in}{2.508417in}}%
\pgfpathlineto{\pgfqpoint{3.837890in}{2.508417in}}%
\pgfpathlineto{\pgfqpoint{3.837890in}{2.505468in}}%
\pgfpathmoveto{\pgfqpoint{3.837890in}{2.505468in}}%
\pgfpathlineto{\pgfqpoint{3.837890in}{2.505468in}}%
\pgfpathlineto{\pgfqpoint{3.837890in}{2.508417in}}%
\pgfpathlineto{\pgfqpoint{3.842431in}{2.508417in}}%
\pgfpathlineto{\pgfqpoint{3.842431in}{2.505468in}}%
\pgfpathmoveto{\pgfqpoint{3.837890in}{2.508417in}}%
\pgfpathlineto{\pgfqpoint{3.837890in}{2.508417in}}%
\pgfpathlineto{\pgfqpoint{3.837890in}{2.511366in}}%
\pgfpathlineto{\pgfqpoint{3.842431in}{2.511366in}}%
\pgfpathlineto{\pgfqpoint{3.842431in}{2.508417in}}%
\pgfpathmoveto{\pgfqpoint{3.842431in}{2.508417in}}%
\pgfpathlineto{\pgfqpoint{3.842431in}{2.508417in}}%
\pgfpathlineto{\pgfqpoint{3.842431in}{2.511366in}}%
\pgfpathlineto{\pgfqpoint{3.846972in}{2.511366in}}%
\pgfpathlineto{\pgfqpoint{3.846972in}{2.508417in}}%
\pgfpathmoveto{\pgfqpoint{3.842431in}{2.511366in}}%
\pgfpathlineto{\pgfqpoint{3.842431in}{2.511366in}}%
\pgfpathlineto{\pgfqpoint{3.842431in}{2.514316in}}%
\pgfpathlineto{\pgfqpoint{3.846972in}{2.514316in}}%
\pgfpathlineto{\pgfqpoint{3.846972in}{2.511366in}}%
\pgfpathmoveto{\pgfqpoint{3.846972in}{2.511366in}}%
\pgfpathlineto{\pgfqpoint{3.846972in}{2.511366in}}%
\pgfpathlineto{\pgfqpoint{3.846972in}{2.514316in}}%
\pgfpathlineto{\pgfqpoint{3.851513in}{2.514316in}}%
\pgfpathlineto{\pgfqpoint{3.851513in}{2.511366in}}%
\pgfpathmoveto{\pgfqpoint{3.846972in}{2.514316in}}%
\pgfpathlineto{\pgfqpoint{3.846972in}{2.514316in}}%
\pgfpathlineto{\pgfqpoint{3.846972in}{2.517265in}}%
\pgfpathlineto{\pgfqpoint{3.851513in}{2.517265in}}%
\pgfpathlineto{\pgfqpoint{3.851513in}{2.514316in}}%
\pgfpathmoveto{\pgfqpoint{3.851513in}{2.514316in}}%
\pgfpathlineto{\pgfqpoint{3.851513in}{2.514316in}}%
\pgfpathlineto{\pgfqpoint{3.851513in}{2.517265in}}%
\pgfpathlineto{\pgfqpoint{3.856054in}{2.517265in}}%
\pgfpathlineto{\pgfqpoint{3.856054in}{2.514316in}}%
\pgfpathmoveto{\pgfqpoint{3.851513in}{2.517265in}}%
\pgfpathlineto{\pgfqpoint{3.851513in}{2.517265in}}%
\pgfpathlineto{\pgfqpoint{3.851513in}{2.520214in}}%
\pgfpathlineto{\pgfqpoint{3.856054in}{2.520214in}}%
\pgfpathlineto{\pgfqpoint{3.856054in}{2.517265in}}%
\pgfpathmoveto{\pgfqpoint{3.856054in}{2.517265in}}%
\pgfpathlineto{\pgfqpoint{3.856054in}{2.517265in}}%
\pgfpathlineto{\pgfqpoint{3.856054in}{2.520214in}}%
\pgfpathlineto{\pgfqpoint{3.860595in}{2.520214in}}%
\pgfpathlineto{\pgfqpoint{3.860595in}{2.517265in}}%
\pgfpathmoveto{\pgfqpoint{3.856054in}{2.520214in}}%
\pgfpathlineto{\pgfqpoint{3.856054in}{2.520214in}}%
\pgfpathlineto{\pgfqpoint{3.856054in}{2.523163in}}%
\pgfpathlineto{\pgfqpoint{3.860595in}{2.523163in}}%
\pgfpathlineto{\pgfqpoint{3.860595in}{2.520214in}}%
\pgfpathmoveto{\pgfqpoint{3.860595in}{2.520214in}}%
\pgfpathlineto{\pgfqpoint{3.860595in}{2.520214in}}%
\pgfpathlineto{\pgfqpoint{3.860595in}{2.523163in}}%
\pgfpathlineto{\pgfqpoint{3.865136in}{2.523163in}}%
\pgfpathlineto{\pgfqpoint{3.865136in}{2.520214in}}%
\pgfpathmoveto{\pgfqpoint{3.860595in}{2.523163in}}%
\pgfpathlineto{\pgfqpoint{3.860595in}{2.523163in}}%
\pgfpathlineto{\pgfqpoint{3.860595in}{2.526113in}}%
\pgfpathlineto{\pgfqpoint{3.865136in}{2.526113in}}%
\pgfpathlineto{\pgfqpoint{3.865136in}{2.523163in}}%
\pgfpathmoveto{\pgfqpoint{3.865136in}{2.523163in}}%
\pgfpathlineto{\pgfqpoint{3.865136in}{2.523163in}}%
\pgfpathlineto{\pgfqpoint{3.865136in}{2.526113in}}%
\pgfpathlineto{\pgfqpoint{3.869677in}{2.526113in}}%
\pgfpathlineto{\pgfqpoint{3.869677in}{2.523163in}}%
\pgfpathmoveto{\pgfqpoint{3.869677in}{2.523163in}}%
\pgfpathlineto{\pgfqpoint{3.869677in}{2.523163in}}%
\pgfpathlineto{\pgfqpoint{3.869677in}{2.526113in}}%
\pgfpathlineto{\pgfqpoint{3.874218in}{2.526113in}}%
\pgfpathlineto{\pgfqpoint{3.874218in}{2.523163in}}%
\pgfpathmoveto{\pgfqpoint{3.869677in}{2.526113in}}%
\pgfpathlineto{\pgfqpoint{3.869677in}{2.526113in}}%
\pgfpathlineto{\pgfqpoint{3.869677in}{2.529062in}}%
\pgfpathlineto{\pgfqpoint{3.874218in}{2.529062in}}%
\pgfpathlineto{\pgfqpoint{3.874218in}{2.526113in}}%
\pgfpathmoveto{\pgfqpoint{3.874218in}{2.526113in}}%
\pgfpathlineto{\pgfqpoint{3.874218in}{2.526113in}}%
\pgfpathlineto{\pgfqpoint{3.874218in}{2.529062in}}%
\pgfpathlineto{\pgfqpoint{3.878759in}{2.529062in}}%
\pgfpathlineto{\pgfqpoint{3.878759in}{2.526113in}}%
\pgfpathmoveto{\pgfqpoint{3.874218in}{2.529062in}}%
\pgfpathlineto{\pgfqpoint{3.874218in}{2.529062in}}%
\pgfpathlineto{\pgfqpoint{3.874218in}{2.532011in}}%
\pgfpathlineto{\pgfqpoint{3.878759in}{2.532011in}}%
\pgfpathlineto{\pgfqpoint{3.878759in}{2.529062in}}%
\pgfpathmoveto{\pgfqpoint{3.878759in}{2.529062in}}%
\pgfpathlineto{\pgfqpoint{3.878759in}{2.529062in}}%
\pgfpathlineto{\pgfqpoint{3.878759in}{2.532011in}}%
\pgfpathlineto{\pgfqpoint{3.883299in}{2.532011in}}%
\pgfpathlineto{\pgfqpoint{3.883299in}{2.529062in}}%
\pgfpathmoveto{\pgfqpoint{3.878759in}{2.532011in}}%
\pgfpathlineto{\pgfqpoint{3.878759in}{2.532011in}}%
\pgfpathlineto{\pgfqpoint{3.878759in}{2.534961in}}%
\pgfpathlineto{\pgfqpoint{3.883299in}{2.534961in}}%
\pgfpathlineto{\pgfqpoint{3.883299in}{2.532011in}}%
\pgfpathmoveto{\pgfqpoint{3.883299in}{2.532011in}}%
\pgfpathlineto{\pgfqpoint{3.883299in}{2.532011in}}%
\pgfpathlineto{\pgfqpoint{3.883299in}{2.534961in}}%
\pgfpathlineto{\pgfqpoint{3.887840in}{2.534961in}}%
\pgfpathlineto{\pgfqpoint{3.887840in}{2.532011in}}%
\pgfpathmoveto{\pgfqpoint{3.883299in}{2.534961in}}%
\pgfpathlineto{\pgfqpoint{3.883299in}{2.534961in}}%
\pgfpathlineto{\pgfqpoint{3.883299in}{2.537910in}}%
\pgfpathlineto{\pgfqpoint{3.887840in}{2.537910in}}%
\pgfpathlineto{\pgfqpoint{3.887840in}{2.534961in}}%
\pgfpathmoveto{\pgfqpoint{3.887840in}{2.534961in}}%
\pgfpathlineto{\pgfqpoint{3.887840in}{2.534961in}}%
\pgfpathlineto{\pgfqpoint{3.887840in}{2.537910in}}%
\pgfpathlineto{\pgfqpoint{3.892381in}{2.537910in}}%
\pgfpathlineto{\pgfqpoint{3.892381in}{2.534961in}}%
\pgfpathmoveto{\pgfqpoint{3.887840in}{2.537910in}}%
\pgfpathlineto{\pgfqpoint{3.887840in}{2.537910in}}%
\pgfpathlineto{\pgfqpoint{3.887840in}{2.540859in}}%
\pgfpathlineto{\pgfqpoint{3.892381in}{2.540859in}}%
\pgfpathlineto{\pgfqpoint{3.892381in}{2.537910in}}%
\pgfpathmoveto{\pgfqpoint{3.892381in}{2.537910in}}%
\pgfpathlineto{\pgfqpoint{3.892381in}{2.537910in}}%
\pgfpathlineto{\pgfqpoint{3.892381in}{2.540859in}}%
\pgfpathlineto{\pgfqpoint{3.896922in}{2.540859in}}%
\pgfpathlineto{\pgfqpoint{3.896922in}{2.537910in}}%
\pgfpathmoveto{\pgfqpoint{3.892381in}{2.540859in}}%
\pgfpathlineto{\pgfqpoint{3.892381in}{2.540859in}}%
\pgfpathlineto{\pgfqpoint{3.892381in}{2.543808in}}%
\pgfpathlineto{\pgfqpoint{3.896922in}{2.543808in}}%
\pgfpathlineto{\pgfqpoint{3.896922in}{2.540859in}}%
\pgfpathmoveto{\pgfqpoint{3.896922in}{2.540859in}}%
\pgfpathlineto{\pgfqpoint{3.896922in}{2.540859in}}%
\pgfpathlineto{\pgfqpoint{3.896922in}{2.543808in}}%
\pgfpathlineto{\pgfqpoint{3.901463in}{2.543808in}}%
\pgfpathlineto{\pgfqpoint{3.901463in}{2.540859in}}%
\pgfpathmoveto{\pgfqpoint{3.896922in}{2.543808in}}%
\pgfpathlineto{\pgfqpoint{3.896922in}{2.543808in}}%
\pgfpathlineto{\pgfqpoint{3.896922in}{2.546758in}}%
\pgfpathlineto{\pgfqpoint{3.901463in}{2.546758in}}%
\pgfpathlineto{\pgfqpoint{3.901463in}{2.543808in}}%
\pgfpathmoveto{\pgfqpoint{3.901463in}{2.543808in}}%
\pgfpathlineto{\pgfqpoint{3.901463in}{2.543808in}}%
\pgfpathlineto{\pgfqpoint{3.901463in}{2.546758in}}%
\pgfpathlineto{\pgfqpoint{3.906004in}{2.546758in}}%
\pgfpathlineto{\pgfqpoint{3.906004in}{2.543808in}}%
\pgfpathmoveto{\pgfqpoint{3.901463in}{2.546758in}}%
\pgfpathlineto{\pgfqpoint{3.901463in}{2.546758in}}%
\pgfpathlineto{\pgfqpoint{3.901463in}{2.549707in}}%
\pgfpathlineto{\pgfqpoint{3.906004in}{2.549707in}}%
\pgfpathlineto{\pgfqpoint{3.906004in}{2.546758in}}%
\pgfpathmoveto{\pgfqpoint{3.906004in}{2.546758in}}%
\pgfpathlineto{\pgfqpoint{3.906004in}{2.546758in}}%
\pgfpathlineto{\pgfqpoint{3.906004in}{2.549707in}}%
\pgfpathlineto{\pgfqpoint{3.910545in}{2.549707in}}%
\pgfpathlineto{\pgfqpoint{3.910545in}{2.546758in}}%
\pgfpathmoveto{\pgfqpoint{3.906004in}{2.549707in}}%
\pgfpathlineto{\pgfqpoint{3.906004in}{2.549707in}}%
\pgfpathlineto{\pgfqpoint{3.906004in}{2.552656in}}%
\pgfpathlineto{\pgfqpoint{3.910545in}{2.552656in}}%
\pgfpathlineto{\pgfqpoint{3.910545in}{2.549707in}}%
\pgfpathmoveto{\pgfqpoint{3.910545in}{2.549707in}}%
\pgfpathlineto{\pgfqpoint{3.910545in}{2.549707in}}%
\pgfpathlineto{\pgfqpoint{3.910545in}{2.552656in}}%
\pgfpathlineto{\pgfqpoint{3.915086in}{2.552656in}}%
\pgfpathlineto{\pgfqpoint{3.915086in}{2.549707in}}%
\pgfpathmoveto{\pgfqpoint{3.910545in}{2.552656in}}%
\pgfpathlineto{\pgfqpoint{3.910545in}{2.552656in}}%
\pgfpathlineto{\pgfqpoint{3.910545in}{2.555606in}}%
\pgfpathlineto{\pgfqpoint{3.915086in}{2.555606in}}%
\pgfpathlineto{\pgfqpoint{3.915086in}{2.552656in}}%
\pgfpathmoveto{\pgfqpoint{3.915086in}{2.552656in}}%
\pgfpathlineto{\pgfqpoint{3.915086in}{2.552656in}}%
\pgfpathlineto{\pgfqpoint{3.915086in}{2.555606in}}%
\pgfpathlineto{\pgfqpoint{3.919627in}{2.555606in}}%
\pgfpathlineto{\pgfqpoint{3.919627in}{2.552656in}}%
\pgfpathmoveto{\pgfqpoint{3.915086in}{2.555606in}}%
\pgfpathlineto{\pgfqpoint{3.915086in}{2.555606in}}%
\pgfpathlineto{\pgfqpoint{3.915086in}{2.558555in}}%
\pgfpathlineto{\pgfqpoint{3.919627in}{2.558555in}}%
\pgfpathlineto{\pgfqpoint{3.919627in}{2.555606in}}%
\pgfpathmoveto{\pgfqpoint{3.919627in}{2.555606in}}%
\pgfpathlineto{\pgfqpoint{3.919627in}{2.555606in}}%
\pgfpathlineto{\pgfqpoint{3.919627in}{2.558555in}}%
\pgfpathlineto{\pgfqpoint{3.924168in}{2.558555in}}%
\pgfpathlineto{\pgfqpoint{3.924168in}{2.555606in}}%
\pgfpathmoveto{\pgfqpoint{3.919627in}{2.558555in}}%
\pgfpathlineto{\pgfqpoint{3.919627in}{2.558555in}}%
\pgfpathlineto{\pgfqpoint{3.919627in}{2.561504in}}%
\pgfpathlineto{\pgfqpoint{3.924168in}{2.561504in}}%
\pgfpathlineto{\pgfqpoint{3.924168in}{2.558555in}}%
\pgfpathmoveto{\pgfqpoint{3.924168in}{2.558555in}}%
\pgfpathlineto{\pgfqpoint{3.924168in}{2.558555in}}%
\pgfpathlineto{\pgfqpoint{3.924168in}{2.561504in}}%
\pgfpathlineto{\pgfqpoint{3.928709in}{2.561504in}}%
\pgfpathlineto{\pgfqpoint{3.928709in}{2.558555in}}%
\pgfpathmoveto{\pgfqpoint{3.924168in}{2.561504in}}%
\pgfpathlineto{\pgfqpoint{3.924168in}{2.561504in}}%
\pgfpathlineto{\pgfqpoint{3.924168in}{2.564454in}}%
\pgfpathlineto{\pgfqpoint{3.928709in}{2.564454in}}%
\pgfpathlineto{\pgfqpoint{3.928709in}{2.561504in}}%
\pgfpathmoveto{\pgfqpoint{3.928709in}{2.561504in}}%
\pgfpathlineto{\pgfqpoint{3.928709in}{2.561504in}}%
\pgfpathlineto{\pgfqpoint{3.928709in}{2.564454in}}%
\pgfpathlineto{\pgfqpoint{3.933249in}{2.564454in}}%
\pgfpathlineto{\pgfqpoint{3.933249in}{2.561504in}}%
\pgfpathmoveto{\pgfqpoint{3.928709in}{2.564454in}}%
\pgfpathlineto{\pgfqpoint{3.928709in}{2.564454in}}%
\pgfpathlineto{\pgfqpoint{3.928709in}{2.567403in}}%
\pgfpathlineto{\pgfqpoint{3.933249in}{2.567403in}}%
\pgfpathlineto{\pgfqpoint{3.933249in}{2.564454in}}%
\pgfpathmoveto{\pgfqpoint{3.933249in}{2.564454in}}%
\pgfpathlineto{\pgfqpoint{3.933249in}{2.564454in}}%
\pgfpathlineto{\pgfqpoint{3.933249in}{2.567403in}}%
\pgfpathlineto{\pgfqpoint{3.937790in}{2.567403in}}%
\pgfpathlineto{\pgfqpoint{3.937790in}{2.564454in}}%
\pgfpathmoveto{\pgfqpoint{3.933249in}{2.567403in}}%
\pgfpathlineto{\pgfqpoint{3.933249in}{2.567403in}}%
\pgfpathlineto{\pgfqpoint{3.933249in}{2.570352in}}%
\pgfpathlineto{\pgfqpoint{3.937790in}{2.570352in}}%
\pgfpathlineto{\pgfqpoint{3.937790in}{2.567403in}}%
\pgfpathmoveto{\pgfqpoint{3.937790in}{2.567403in}}%
\pgfpathlineto{\pgfqpoint{3.937790in}{2.567403in}}%
\pgfpathlineto{\pgfqpoint{3.937790in}{2.570352in}}%
\pgfpathlineto{\pgfqpoint{3.942331in}{2.570352in}}%
\pgfpathlineto{\pgfqpoint{3.942331in}{2.567403in}}%
\pgfpathmoveto{\pgfqpoint{3.937790in}{2.570352in}}%
\pgfpathlineto{\pgfqpoint{3.937790in}{2.570352in}}%
\pgfpathlineto{\pgfqpoint{3.937790in}{2.573301in}}%
\pgfpathlineto{\pgfqpoint{3.942331in}{2.573301in}}%
\pgfpathlineto{\pgfqpoint{3.942331in}{2.570352in}}%
\pgfpathmoveto{\pgfqpoint{3.942331in}{2.570352in}}%
\pgfpathlineto{\pgfqpoint{3.942331in}{2.570352in}}%
\pgfpathlineto{\pgfqpoint{3.942331in}{2.573301in}}%
\pgfpathlineto{\pgfqpoint{3.946872in}{2.573301in}}%
\pgfpathlineto{\pgfqpoint{3.946872in}{2.570352in}}%
\pgfpathmoveto{\pgfqpoint{3.942331in}{2.573301in}}%
\pgfpathlineto{\pgfqpoint{3.942331in}{2.573301in}}%
\pgfpathlineto{\pgfqpoint{3.942331in}{2.576251in}}%
\pgfpathlineto{\pgfqpoint{3.946872in}{2.576251in}}%
\pgfpathlineto{\pgfqpoint{3.946872in}{2.573301in}}%
\pgfpathmoveto{\pgfqpoint{3.946872in}{2.573301in}}%
\pgfpathlineto{\pgfqpoint{3.946872in}{2.573301in}}%
\pgfpathlineto{\pgfqpoint{3.946872in}{2.576251in}}%
\pgfpathlineto{\pgfqpoint{3.951413in}{2.576251in}}%
\pgfpathlineto{\pgfqpoint{3.951413in}{2.573301in}}%
\pgfpathmoveto{\pgfqpoint{3.946872in}{2.576251in}}%
\pgfpathlineto{\pgfqpoint{3.946872in}{2.576251in}}%
\pgfpathlineto{\pgfqpoint{3.946872in}{2.579200in}}%
\pgfpathlineto{\pgfqpoint{3.951413in}{2.579200in}}%
\pgfpathlineto{\pgfqpoint{3.951413in}{2.576251in}}%
\pgfpathmoveto{\pgfqpoint{3.951413in}{2.576251in}}%
\pgfpathlineto{\pgfqpoint{3.951413in}{2.576251in}}%
\pgfpathlineto{\pgfqpoint{3.951413in}{2.579200in}}%
\pgfpathlineto{\pgfqpoint{3.955954in}{2.579200in}}%
\pgfpathlineto{\pgfqpoint{3.955954in}{2.576251in}}%
\pgfpathmoveto{\pgfqpoint{3.951413in}{2.579200in}}%
\pgfpathlineto{\pgfqpoint{3.951413in}{2.579200in}}%
\pgfpathlineto{\pgfqpoint{3.951413in}{2.582149in}}%
\pgfpathlineto{\pgfqpoint{3.955954in}{2.582149in}}%
\pgfpathlineto{\pgfqpoint{3.955954in}{2.579200in}}%
\pgfpathmoveto{\pgfqpoint{3.955954in}{2.579200in}}%
\pgfpathlineto{\pgfqpoint{3.955954in}{2.579200in}}%
\pgfpathlineto{\pgfqpoint{3.955954in}{2.582149in}}%
\pgfpathlineto{\pgfqpoint{3.960495in}{2.582149in}}%
\pgfpathlineto{\pgfqpoint{3.960495in}{2.579200in}}%
\pgfpathmoveto{\pgfqpoint{3.955954in}{2.582149in}}%
\pgfpathlineto{\pgfqpoint{3.955954in}{2.582149in}}%
\pgfpathlineto{\pgfqpoint{3.955954in}{2.585098in}}%
\pgfpathlineto{\pgfqpoint{3.960495in}{2.585098in}}%
\pgfpathlineto{\pgfqpoint{3.960495in}{2.582149in}}%
\pgfpathmoveto{\pgfqpoint{3.960495in}{2.582149in}}%
\pgfpathlineto{\pgfqpoint{3.960495in}{2.582149in}}%
\pgfpathlineto{\pgfqpoint{3.960495in}{2.585098in}}%
\pgfpathlineto{\pgfqpoint{3.965036in}{2.585098in}}%
\pgfpathlineto{\pgfqpoint{3.965036in}{2.582149in}}%
\pgfpathmoveto{\pgfqpoint{3.960495in}{2.585098in}}%
\pgfpathlineto{\pgfqpoint{3.960495in}{2.585098in}}%
\pgfpathlineto{\pgfqpoint{3.960495in}{2.588047in}}%
\pgfpathlineto{\pgfqpoint{3.965036in}{2.588047in}}%
\pgfpathlineto{\pgfqpoint{3.965036in}{2.585098in}}%
\pgfpathmoveto{\pgfqpoint{3.965036in}{2.585098in}}%
\pgfpathlineto{\pgfqpoint{3.965036in}{2.585098in}}%
\pgfpathlineto{\pgfqpoint{3.965036in}{2.588047in}}%
\pgfpathlineto{\pgfqpoint{3.969577in}{2.588047in}}%
\pgfpathlineto{\pgfqpoint{3.969577in}{2.585098in}}%
\pgfpathmoveto{\pgfqpoint{3.965036in}{2.588047in}}%
\pgfpathlineto{\pgfqpoint{3.965036in}{2.588047in}}%
\pgfpathlineto{\pgfqpoint{3.965036in}{2.590996in}}%
\pgfpathlineto{\pgfqpoint{3.969577in}{2.590996in}}%
\pgfpathlineto{\pgfqpoint{3.969577in}{2.588047in}}%
\pgfpathmoveto{\pgfqpoint{3.969577in}{2.588047in}}%
\pgfpathlineto{\pgfqpoint{3.969577in}{2.588047in}}%
\pgfpathlineto{\pgfqpoint{3.969577in}{2.590996in}}%
\pgfpathlineto{\pgfqpoint{3.974119in}{2.590996in}}%
\pgfpathlineto{\pgfqpoint{3.974119in}{2.588047in}}%
\pgfpathmoveto{\pgfqpoint{3.969577in}{2.590996in}}%
\pgfpathlineto{\pgfqpoint{3.969577in}{2.590996in}}%
\pgfpathlineto{\pgfqpoint{3.969577in}{2.593946in}}%
\pgfpathlineto{\pgfqpoint{3.974119in}{2.593946in}}%
\pgfpathlineto{\pgfqpoint{3.974119in}{2.590996in}}%
\pgfpathmoveto{\pgfqpoint{3.974119in}{2.590996in}}%
\pgfpathlineto{\pgfqpoint{3.974119in}{2.590996in}}%
\pgfpathlineto{\pgfqpoint{3.974119in}{2.593946in}}%
\pgfpathlineto{\pgfqpoint{3.978660in}{2.593946in}}%
\pgfpathlineto{\pgfqpoint{3.978660in}{2.590996in}}%
\pgfpathmoveto{\pgfqpoint{3.974119in}{2.593946in}}%
\pgfpathlineto{\pgfqpoint{3.974119in}{2.593946in}}%
\pgfpathlineto{\pgfqpoint{3.974119in}{2.596895in}}%
\pgfpathlineto{\pgfqpoint{3.978660in}{2.596895in}}%
\pgfpathlineto{\pgfqpoint{3.978660in}{2.593946in}}%
\pgfpathmoveto{\pgfqpoint{3.978660in}{2.593946in}}%
\pgfpathlineto{\pgfqpoint{3.978660in}{2.593946in}}%
\pgfpathlineto{\pgfqpoint{3.978660in}{2.596895in}}%
\pgfpathlineto{\pgfqpoint{3.983201in}{2.596895in}}%
\pgfpathlineto{\pgfqpoint{3.983201in}{2.593946in}}%
\pgfpathmoveto{\pgfqpoint{3.978660in}{2.596895in}}%
\pgfpathlineto{\pgfqpoint{3.978660in}{2.596895in}}%
\pgfpathlineto{\pgfqpoint{3.978660in}{2.599844in}}%
\pgfpathlineto{\pgfqpoint{3.983201in}{2.599844in}}%
\pgfpathlineto{\pgfqpoint{3.983201in}{2.596895in}}%
\pgfpathmoveto{\pgfqpoint{3.983201in}{2.596895in}}%
\pgfpathlineto{\pgfqpoint{3.983201in}{2.596895in}}%
\pgfpathlineto{\pgfqpoint{3.983201in}{2.599844in}}%
\pgfpathlineto{\pgfqpoint{3.987742in}{2.599844in}}%
\pgfpathlineto{\pgfqpoint{3.987742in}{2.596895in}}%
\pgfpathmoveto{\pgfqpoint{3.983201in}{2.599844in}}%
\pgfpathlineto{\pgfqpoint{3.983201in}{2.599844in}}%
\pgfpathlineto{\pgfqpoint{3.983201in}{2.602793in}}%
\pgfpathlineto{\pgfqpoint{3.987742in}{2.602793in}}%
\pgfpathlineto{\pgfqpoint{3.987742in}{2.599844in}}%
\pgfpathmoveto{\pgfqpoint{3.987742in}{2.599844in}}%
\pgfpathlineto{\pgfqpoint{3.987742in}{2.599844in}}%
\pgfpathlineto{\pgfqpoint{3.987742in}{2.602793in}}%
\pgfpathlineto{\pgfqpoint{3.992283in}{2.602793in}}%
\pgfpathlineto{\pgfqpoint{3.992283in}{2.599844in}}%
\pgfpathmoveto{\pgfqpoint{3.987742in}{2.602793in}}%
\pgfpathlineto{\pgfqpoint{3.987742in}{2.602793in}}%
\pgfpathlineto{\pgfqpoint{3.987742in}{2.605742in}}%
\pgfpathlineto{\pgfqpoint{3.992283in}{2.605742in}}%
\pgfpathlineto{\pgfqpoint{3.992283in}{2.602793in}}%
\pgfpathmoveto{\pgfqpoint{3.992283in}{2.602793in}}%
\pgfpathlineto{\pgfqpoint{3.992283in}{2.602793in}}%
\pgfpathlineto{\pgfqpoint{3.992283in}{2.605742in}}%
\pgfpathlineto{\pgfqpoint{3.996824in}{2.605742in}}%
\pgfpathlineto{\pgfqpoint{3.996824in}{2.602793in}}%
\pgfpathmoveto{\pgfqpoint{3.992283in}{2.605742in}}%
\pgfpathlineto{\pgfqpoint{3.992283in}{2.605742in}}%
\pgfpathlineto{\pgfqpoint{3.992283in}{2.608691in}}%
\pgfpathlineto{\pgfqpoint{3.996824in}{2.608691in}}%
\pgfpathlineto{\pgfqpoint{3.996824in}{2.605742in}}%
\pgfpathmoveto{\pgfqpoint{3.996824in}{2.605742in}}%
\pgfpathlineto{\pgfqpoint{3.996824in}{2.605742in}}%
\pgfpathlineto{\pgfqpoint{3.996824in}{2.608691in}}%
\pgfpathlineto{\pgfqpoint{4.001365in}{2.608691in}}%
\pgfpathlineto{\pgfqpoint{4.001365in}{2.605742in}}%
\pgfpathmoveto{\pgfqpoint{3.996824in}{2.608691in}}%
\pgfpathlineto{\pgfqpoint{3.996824in}{2.608691in}}%
\pgfpathlineto{\pgfqpoint{3.996824in}{2.611641in}}%
\pgfpathlineto{\pgfqpoint{4.001365in}{2.611641in}}%
\pgfpathlineto{\pgfqpoint{4.001365in}{2.608691in}}%
\pgfpathmoveto{\pgfqpoint{4.001365in}{2.608691in}}%
\pgfpathlineto{\pgfqpoint{4.001365in}{2.608691in}}%
\pgfpathlineto{\pgfqpoint{4.001365in}{2.611641in}}%
\pgfpathlineto{\pgfqpoint{4.005906in}{2.611641in}}%
\pgfpathlineto{\pgfqpoint{4.005906in}{2.608691in}}%
\pgfpathmoveto{\pgfqpoint{4.001365in}{2.611641in}}%
\pgfpathlineto{\pgfqpoint{4.001365in}{2.611641in}}%
\pgfpathlineto{\pgfqpoint{4.001365in}{2.614590in}}%
\pgfpathlineto{\pgfqpoint{4.005906in}{2.614590in}}%
\pgfpathlineto{\pgfqpoint{4.005906in}{2.611641in}}%
\pgfpathmoveto{\pgfqpoint{4.005906in}{2.611641in}}%
\pgfpathlineto{\pgfqpoint{4.005906in}{2.611641in}}%
\pgfpathlineto{\pgfqpoint{4.005906in}{2.614590in}}%
\pgfpathlineto{\pgfqpoint{4.010447in}{2.614590in}}%
\pgfpathlineto{\pgfqpoint{4.010447in}{2.611641in}}%
\pgfpathmoveto{\pgfqpoint{4.005906in}{2.614590in}}%
\pgfpathlineto{\pgfqpoint{4.005906in}{2.614590in}}%
\pgfpathlineto{\pgfqpoint{4.005906in}{2.617539in}}%
\pgfpathlineto{\pgfqpoint{4.010447in}{2.617539in}}%
\pgfpathlineto{\pgfqpoint{4.010447in}{2.614590in}}%
\pgfpathmoveto{\pgfqpoint{4.010447in}{2.614590in}}%
\pgfpathlineto{\pgfqpoint{4.010447in}{2.614590in}}%
\pgfpathlineto{\pgfqpoint{4.010447in}{2.617539in}}%
\pgfpathlineto{\pgfqpoint{4.014988in}{2.617539in}}%
\pgfpathlineto{\pgfqpoint{4.014988in}{2.614590in}}%
\pgfpathmoveto{\pgfqpoint{4.010447in}{2.617539in}}%
\pgfpathlineto{\pgfqpoint{4.010447in}{2.617539in}}%
\pgfpathlineto{\pgfqpoint{4.010447in}{2.620488in}}%
\pgfpathlineto{\pgfqpoint{4.014988in}{2.620488in}}%
\pgfpathlineto{\pgfqpoint{4.014988in}{2.617539in}}%
\pgfpathmoveto{\pgfqpoint{4.014988in}{2.617539in}}%
\pgfpathlineto{\pgfqpoint{4.014988in}{2.617539in}}%
\pgfpathlineto{\pgfqpoint{4.014988in}{2.620488in}}%
\pgfpathlineto{\pgfqpoint{4.019529in}{2.620488in}}%
\pgfpathlineto{\pgfqpoint{4.019529in}{2.617539in}}%
\pgfpathmoveto{\pgfqpoint{4.014988in}{2.620488in}}%
\pgfpathlineto{\pgfqpoint{4.014988in}{2.620488in}}%
\pgfpathlineto{\pgfqpoint{4.014988in}{2.623437in}}%
\pgfpathlineto{\pgfqpoint{4.019529in}{2.623437in}}%
\pgfpathlineto{\pgfqpoint{4.019529in}{2.620488in}}%
\pgfpathmoveto{\pgfqpoint{4.019529in}{2.620488in}}%
\pgfpathlineto{\pgfqpoint{4.019529in}{2.620488in}}%
\pgfpathlineto{\pgfqpoint{4.019529in}{2.623437in}}%
\pgfpathlineto{\pgfqpoint{4.024070in}{2.623437in}}%
\pgfpathlineto{\pgfqpoint{4.024070in}{2.620488in}}%
\pgfpathmoveto{\pgfqpoint{4.019529in}{2.623437in}}%
\pgfpathlineto{\pgfqpoint{4.019529in}{2.623437in}}%
\pgfpathlineto{\pgfqpoint{4.019529in}{2.626386in}}%
\pgfpathlineto{\pgfqpoint{4.024070in}{2.626386in}}%
\pgfpathlineto{\pgfqpoint{4.024070in}{2.623437in}}%
\pgfpathmoveto{\pgfqpoint{4.024070in}{2.623437in}}%
\pgfpathlineto{\pgfqpoint{4.024070in}{2.623437in}}%
\pgfpathlineto{\pgfqpoint{4.024070in}{2.626386in}}%
\pgfpathlineto{\pgfqpoint{4.028611in}{2.626386in}}%
\pgfpathlineto{\pgfqpoint{4.028611in}{2.623437in}}%
\pgfpathmoveto{\pgfqpoint{4.024070in}{2.626386in}}%
\pgfpathlineto{\pgfqpoint{4.024070in}{2.626386in}}%
\pgfpathlineto{\pgfqpoint{4.024070in}{2.629335in}}%
\pgfpathlineto{\pgfqpoint{4.028611in}{2.629335in}}%
\pgfpathlineto{\pgfqpoint{4.028611in}{2.626386in}}%
\pgfpathmoveto{\pgfqpoint{4.028611in}{2.626386in}}%
\pgfpathlineto{\pgfqpoint{4.028611in}{2.626386in}}%
\pgfpathlineto{\pgfqpoint{4.028611in}{2.629335in}}%
\pgfpathlineto{\pgfqpoint{4.033152in}{2.629335in}}%
\pgfpathlineto{\pgfqpoint{4.033152in}{2.626386in}}%
\pgfpathmoveto{\pgfqpoint{4.028611in}{2.629335in}}%
\pgfpathlineto{\pgfqpoint{4.028611in}{2.629335in}}%
\pgfpathlineto{\pgfqpoint{4.028611in}{2.632285in}}%
\pgfpathlineto{\pgfqpoint{4.033152in}{2.632285in}}%
\pgfpathlineto{\pgfqpoint{4.033152in}{2.629335in}}%
\pgfpathmoveto{\pgfqpoint{4.033152in}{2.629335in}}%
\pgfpathlineto{\pgfqpoint{4.033152in}{2.629335in}}%
\pgfpathlineto{\pgfqpoint{4.033152in}{2.632285in}}%
\pgfpathlineto{\pgfqpoint{4.037694in}{2.632285in}}%
\pgfpathlineto{\pgfqpoint{4.037694in}{2.629335in}}%
\pgfpathmoveto{\pgfqpoint{4.033152in}{2.632285in}}%
\pgfpathlineto{\pgfqpoint{4.033152in}{2.632285in}}%
\pgfpathlineto{\pgfqpoint{4.033152in}{2.635234in}}%
\pgfpathlineto{\pgfqpoint{4.037694in}{2.635234in}}%
\pgfpathlineto{\pgfqpoint{4.037694in}{2.632285in}}%
\pgfpathmoveto{\pgfqpoint{4.037694in}{2.632285in}}%
\pgfpathlineto{\pgfqpoint{4.037694in}{2.632285in}}%
\pgfpathlineto{\pgfqpoint{4.037694in}{2.635234in}}%
\pgfpathlineto{\pgfqpoint{4.042235in}{2.635234in}}%
\pgfpathlineto{\pgfqpoint{4.042235in}{2.632285in}}%
\pgfpathmoveto{\pgfqpoint{4.037694in}{2.635234in}}%
\pgfpathlineto{\pgfqpoint{4.037694in}{2.635234in}}%
\pgfpathlineto{\pgfqpoint{4.037694in}{2.638183in}}%
\pgfpathlineto{\pgfqpoint{4.042235in}{2.638183in}}%
\pgfpathlineto{\pgfqpoint{4.042235in}{2.635234in}}%
\pgfpathmoveto{\pgfqpoint{4.042235in}{2.635234in}}%
\pgfpathlineto{\pgfqpoint{4.042235in}{2.635234in}}%
\pgfpathlineto{\pgfqpoint{4.042235in}{2.638183in}}%
\pgfpathlineto{\pgfqpoint{4.046776in}{2.638183in}}%
\pgfpathlineto{\pgfqpoint{4.046776in}{2.635234in}}%
\pgfpathmoveto{\pgfqpoint{4.042235in}{2.638183in}}%
\pgfpathlineto{\pgfqpoint{4.042235in}{2.638183in}}%
\pgfpathlineto{\pgfqpoint{4.042235in}{2.641132in}}%
\pgfpathlineto{\pgfqpoint{4.046776in}{2.641132in}}%
\pgfpathlineto{\pgfqpoint{4.046776in}{2.638183in}}%
\pgfpathmoveto{\pgfqpoint{4.046776in}{2.638183in}}%
\pgfpathlineto{\pgfqpoint{4.046776in}{2.638183in}}%
\pgfpathlineto{\pgfqpoint{4.046776in}{2.641132in}}%
\pgfpathlineto{\pgfqpoint{4.051317in}{2.641132in}}%
\pgfpathlineto{\pgfqpoint{4.051317in}{2.638183in}}%
\pgfpathmoveto{\pgfqpoint{4.046776in}{2.641132in}}%
\pgfpathlineto{\pgfqpoint{4.046776in}{2.641132in}}%
\pgfpathlineto{\pgfqpoint{4.046776in}{2.644081in}}%
\pgfpathlineto{\pgfqpoint{4.051317in}{2.644081in}}%
\pgfpathlineto{\pgfqpoint{4.051317in}{2.641132in}}%
\pgfpathmoveto{\pgfqpoint{4.051317in}{2.641132in}}%
\pgfpathlineto{\pgfqpoint{4.051317in}{2.641132in}}%
\pgfpathlineto{\pgfqpoint{4.051317in}{2.644081in}}%
\pgfpathlineto{\pgfqpoint{4.055858in}{2.644081in}}%
\pgfpathlineto{\pgfqpoint{4.055858in}{2.641132in}}%
\pgfpathmoveto{\pgfqpoint{4.051317in}{2.644081in}}%
\pgfpathlineto{\pgfqpoint{4.051317in}{2.644081in}}%
\pgfpathlineto{\pgfqpoint{4.051317in}{2.647030in}}%
\pgfpathlineto{\pgfqpoint{4.055858in}{2.647030in}}%
\pgfpathlineto{\pgfqpoint{4.055858in}{2.644081in}}%
\pgfpathmoveto{\pgfqpoint{4.055858in}{2.644081in}}%
\pgfpathlineto{\pgfqpoint{4.055858in}{2.644081in}}%
\pgfpathlineto{\pgfqpoint{4.055858in}{2.647030in}}%
\pgfpathlineto{\pgfqpoint{4.060399in}{2.647030in}}%
\pgfpathlineto{\pgfqpoint{4.060399in}{2.644081in}}%
\pgfpathmoveto{\pgfqpoint{4.055858in}{2.647030in}}%
\pgfpathlineto{\pgfqpoint{4.055858in}{2.647030in}}%
\pgfpathlineto{\pgfqpoint{4.055858in}{2.649979in}}%
\pgfpathlineto{\pgfqpoint{4.060399in}{2.649979in}}%
\pgfpathlineto{\pgfqpoint{4.060399in}{2.647030in}}%
\pgfpathmoveto{\pgfqpoint{4.055858in}{2.649979in}}%
\pgfpathlineto{\pgfqpoint{4.055858in}{2.649979in}}%
\pgfpathlineto{\pgfqpoint{4.055858in}{2.652929in}}%
\pgfpathlineto{\pgfqpoint{4.060399in}{2.652929in}}%
\pgfpathlineto{\pgfqpoint{4.060399in}{2.649979in}}%
\pgfpathmoveto{\pgfqpoint{4.060399in}{2.649979in}}%
\pgfpathlineto{\pgfqpoint{4.060399in}{2.649979in}}%
\pgfpathlineto{\pgfqpoint{4.060399in}{2.652929in}}%
\pgfpathlineto{\pgfqpoint{4.064940in}{2.652929in}}%
\pgfpathlineto{\pgfqpoint{4.064940in}{2.649979in}}%
\pgfpathmoveto{\pgfqpoint{4.060399in}{2.652929in}}%
\pgfpathlineto{\pgfqpoint{4.060399in}{2.652929in}}%
\pgfpathlineto{\pgfqpoint{4.060399in}{2.655878in}}%
\pgfpathlineto{\pgfqpoint{4.064940in}{2.655878in}}%
\pgfpathlineto{\pgfqpoint{4.064940in}{2.652929in}}%
\pgfpathmoveto{\pgfqpoint{4.064940in}{2.652929in}}%
\pgfpathlineto{\pgfqpoint{4.064940in}{2.652929in}}%
\pgfpathlineto{\pgfqpoint{4.064940in}{2.655878in}}%
\pgfpathlineto{\pgfqpoint{4.069481in}{2.655878in}}%
\pgfpathlineto{\pgfqpoint{4.069481in}{2.652929in}}%
\pgfpathmoveto{\pgfqpoint{4.064940in}{2.655878in}}%
\pgfpathlineto{\pgfqpoint{4.064940in}{2.655878in}}%
\pgfpathlineto{\pgfqpoint{4.064940in}{2.658827in}}%
\pgfpathlineto{\pgfqpoint{4.069481in}{2.658827in}}%
\pgfpathlineto{\pgfqpoint{4.069481in}{2.655878in}}%
\pgfpathmoveto{\pgfqpoint{4.069481in}{2.655878in}}%
\pgfpathlineto{\pgfqpoint{4.069481in}{2.655878in}}%
\pgfpathlineto{\pgfqpoint{4.069481in}{2.658827in}}%
\pgfpathlineto{\pgfqpoint{4.074022in}{2.658827in}}%
\pgfpathlineto{\pgfqpoint{4.074022in}{2.655878in}}%
\pgfpathmoveto{\pgfqpoint{4.069481in}{2.658827in}}%
\pgfpathlineto{\pgfqpoint{4.069481in}{2.658827in}}%
\pgfpathlineto{\pgfqpoint{4.069481in}{2.661776in}}%
\pgfpathlineto{\pgfqpoint{4.074022in}{2.661776in}}%
\pgfpathlineto{\pgfqpoint{4.074022in}{2.658827in}}%
\pgfpathmoveto{\pgfqpoint{4.074022in}{2.658827in}}%
\pgfpathlineto{\pgfqpoint{4.074022in}{2.658827in}}%
\pgfpathlineto{\pgfqpoint{4.074022in}{2.661776in}}%
\pgfpathlineto{\pgfqpoint{4.078563in}{2.661776in}}%
\pgfpathlineto{\pgfqpoint{4.078563in}{2.658827in}}%
\pgfpathmoveto{\pgfqpoint{4.074022in}{2.661776in}}%
\pgfpathlineto{\pgfqpoint{4.074022in}{2.661776in}}%
\pgfpathlineto{\pgfqpoint{4.074022in}{2.664725in}}%
\pgfpathlineto{\pgfqpoint{4.078563in}{2.664725in}}%
\pgfpathlineto{\pgfqpoint{4.078563in}{2.661776in}}%
\pgfpathmoveto{\pgfqpoint{4.078563in}{2.661776in}}%
\pgfpathlineto{\pgfqpoint{4.078563in}{2.661776in}}%
\pgfpathlineto{\pgfqpoint{4.078563in}{2.664725in}}%
\pgfpathlineto{\pgfqpoint{4.083104in}{2.664725in}}%
\pgfpathlineto{\pgfqpoint{4.083104in}{2.661776in}}%
\pgfpathmoveto{\pgfqpoint{4.078563in}{2.664725in}}%
\pgfpathlineto{\pgfqpoint{4.078563in}{2.664725in}}%
\pgfpathlineto{\pgfqpoint{4.078563in}{2.667674in}}%
\pgfpathlineto{\pgfqpoint{4.083104in}{2.667674in}}%
\pgfpathlineto{\pgfqpoint{4.083104in}{2.664725in}}%
\pgfpathmoveto{\pgfqpoint{4.083104in}{2.664725in}}%
\pgfpathlineto{\pgfqpoint{4.083104in}{2.664725in}}%
\pgfpathlineto{\pgfqpoint{4.083104in}{2.667674in}}%
\pgfpathlineto{\pgfqpoint{4.087645in}{2.667674in}}%
\pgfpathlineto{\pgfqpoint{4.087645in}{2.664725in}}%
\pgfpathmoveto{\pgfqpoint{4.083104in}{2.667674in}}%
\pgfpathlineto{\pgfqpoint{4.083104in}{2.667674in}}%
\pgfpathlineto{\pgfqpoint{4.083104in}{2.670623in}}%
\pgfpathlineto{\pgfqpoint{4.087645in}{2.670623in}}%
\pgfpathlineto{\pgfqpoint{4.087645in}{2.667674in}}%
\pgfpathmoveto{\pgfqpoint{4.087645in}{2.667674in}}%
\pgfpathlineto{\pgfqpoint{4.087645in}{2.667674in}}%
\pgfpathlineto{\pgfqpoint{4.087645in}{2.670623in}}%
\pgfpathlineto{\pgfqpoint{4.092186in}{2.670623in}}%
\pgfpathlineto{\pgfqpoint{4.092186in}{2.667674in}}%
\pgfpathmoveto{\pgfqpoint{4.087645in}{2.670623in}}%
\pgfpathlineto{\pgfqpoint{4.087645in}{2.670623in}}%
\pgfpathlineto{\pgfqpoint{4.087645in}{2.673573in}}%
\pgfpathlineto{\pgfqpoint{4.092186in}{2.673573in}}%
\pgfpathlineto{\pgfqpoint{4.092186in}{2.670623in}}%
\pgfpathmoveto{\pgfqpoint{4.092186in}{2.670623in}}%
\pgfpathlineto{\pgfqpoint{4.092186in}{2.670623in}}%
\pgfpathlineto{\pgfqpoint{4.092186in}{2.673573in}}%
\pgfpathlineto{\pgfqpoint{4.096728in}{2.673573in}}%
\pgfpathlineto{\pgfqpoint{4.096728in}{2.670623in}}%
\pgfpathmoveto{\pgfqpoint{4.092186in}{2.673573in}}%
\pgfpathlineto{\pgfqpoint{4.092186in}{2.673573in}}%
\pgfpathlineto{\pgfqpoint{4.092186in}{2.676522in}}%
\pgfpathlineto{\pgfqpoint{4.096728in}{2.676522in}}%
\pgfpathlineto{\pgfqpoint{4.096728in}{2.673573in}}%
\pgfpathmoveto{\pgfqpoint{4.096728in}{2.673573in}}%
\pgfpathlineto{\pgfqpoint{4.096728in}{2.673573in}}%
\pgfpathlineto{\pgfqpoint{4.096728in}{2.676522in}}%
\pgfpathlineto{\pgfqpoint{4.101269in}{2.676522in}}%
\pgfpathlineto{\pgfqpoint{4.101269in}{2.673573in}}%
\pgfpathmoveto{\pgfqpoint{4.096728in}{2.676522in}}%
\pgfpathlineto{\pgfqpoint{4.096728in}{2.676522in}}%
\pgfpathlineto{\pgfqpoint{4.096728in}{2.679471in}}%
\pgfpathlineto{\pgfqpoint{4.101269in}{2.679471in}}%
\pgfpathlineto{\pgfqpoint{4.101269in}{2.676522in}}%
\pgfpathmoveto{\pgfqpoint{4.101269in}{2.676522in}}%
\pgfpathlineto{\pgfqpoint{4.101269in}{2.676522in}}%
\pgfpathlineto{\pgfqpoint{4.101269in}{2.679471in}}%
\pgfpathlineto{\pgfqpoint{4.105810in}{2.679471in}}%
\pgfpathlineto{\pgfqpoint{4.105810in}{2.676522in}}%
\pgfpathmoveto{\pgfqpoint{4.101269in}{2.679471in}}%
\pgfpathlineto{\pgfqpoint{4.101269in}{2.679471in}}%
\pgfpathlineto{\pgfqpoint{4.101269in}{2.682420in}}%
\pgfpathlineto{\pgfqpoint{4.105810in}{2.682420in}}%
\pgfpathlineto{\pgfqpoint{4.105810in}{2.679471in}}%
\pgfpathmoveto{\pgfqpoint{4.105810in}{2.679471in}}%
\pgfpathlineto{\pgfqpoint{4.105810in}{2.679471in}}%
\pgfpathlineto{\pgfqpoint{4.105810in}{2.682420in}}%
\pgfpathlineto{\pgfqpoint{4.110351in}{2.682420in}}%
\pgfpathlineto{\pgfqpoint{4.110351in}{2.679471in}}%
\pgfpathmoveto{\pgfqpoint{4.105810in}{2.682420in}}%
\pgfpathlineto{\pgfqpoint{4.105810in}{2.682420in}}%
\pgfpathlineto{\pgfqpoint{4.105810in}{2.685370in}}%
\pgfpathlineto{\pgfqpoint{4.110351in}{2.685370in}}%
\pgfpathlineto{\pgfqpoint{4.110351in}{2.682420in}}%
\pgfpathmoveto{\pgfqpoint{4.110351in}{2.682420in}}%
\pgfpathlineto{\pgfqpoint{4.110351in}{2.682420in}}%
\pgfpathlineto{\pgfqpoint{4.110351in}{2.685370in}}%
\pgfpathlineto{\pgfqpoint{4.114892in}{2.685370in}}%
\pgfpathlineto{\pgfqpoint{4.114892in}{2.682420in}}%
\pgfpathmoveto{\pgfqpoint{4.110351in}{2.685370in}}%
\pgfpathlineto{\pgfqpoint{4.110351in}{2.685370in}}%
\pgfpathlineto{\pgfqpoint{4.110351in}{2.688319in}}%
\pgfpathlineto{\pgfqpoint{4.114892in}{2.688319in}}%
\pgfpathlineto{\pgfqpoint{4.114892in}{2.685370in}}%
\pgfpathmoveto{\pgfqpoint{4.114892in}{2.685370in}}%
\pgfpathlineto{\pgfqpoint{4.114892in}{2.685370in}}%
\pgfpathlineto{\pgfqpoint{4.114892in}{2.688319in}}%
\pgfpathlineto{\pgfqpoint{4.119433in}{2.688319in}}%
\pgfpathlineto{\pgfqpoint{4.119433in}{2.685370in}}%
\pgfpathmoveto{\pgfqpoint{4.114892in}{2.688319in}}%
\pgfpathlineto{\pgfqpoint{4.114892in}{2.688319in}}%
\pgfpathlineto{\pgfqpoint{4.114892in}{2.691268in}}%
\pgfpathlineto{\pgfqpoint{4.119433in}{2.691268in}}%
\pgfpathlineto{\pgfqpoint{4.119433in}{2.688319in}}%
\pgfpathmoveto{\pgfqpoint{4.119433in}{2.688319in}}%
\pgfpathlineto{\pgfqpoint{4.119433in}{2.688319in}}%
\pgfpathlineto{\pgfqpoint{4.119433in}{2.691268in}}%
\pgfpathlineto{\pgfqpoint{4.123974in}{2.691268in}}%
\pgfpathlineto{\pgfqpoint{4.123974in}{2.688319in}}%
\pgfpathmoveto{\pgfqpoint{4.119433in}{2.691268in}}%
\pgfpathlineto{\pgfqpoint{4.119433in}{2.691268in}}%
\pgfpathlineto{\pgfqpoint{4.119433in}{2.694217in}}%
\pgfpathlineto{\pgfqpoint{4.123974in}{2.694217in}}%
\pgfpathlineto{\pgfqpoint{4.123974in}{2.691268in}}%
\pgfpathmoveto{\pgfqpoint{4.123974in}{2.691268in}}%
\pgfpathlineto{\pgfqpoint{4.123974in}{2.691268in}}%
\pgfpathlineto{\pgfqpoint{4.123974in}{2.694217in}}%
\pgfpathlineto{\pgfqpoint{4.128516in}{2.694217in}}%
\pgfpathlineto{\pgfqpoint{4.128516in}{2.691268in}}%
\pgfpathmoveto{\pgfqpoint{4.123974in}{2.694217in}}%
\pgfpathlineto{\pgfqpoint{4.123974in}{2.694217in}}%
\pgfpathlineto{\pgfqpoint{4.123974in}{2.697166in}}%
\pgfpathlineto{\pgfqpoint{4.128516in}{2.697166in}}%
\pgfpathlineto{\pgfqpoint{4.128516in}{2.694217in}}%
\pgfpathmoveto{\pgfqpoint{4.128516in}{2.694217in}}%
\pgfpathlineto{\pgfqpoint{4.128516in}{2.694217in}}%
\pgfpathlineto{\pgfqpoint{4.128516in}{2.697166in}}%
\pgfpathlineto{\pgfqpoint{4.133057in}{2.697166in}}%
\pgfpathlineto{\pgfqpoint{4.133057in}{2.694217in}}%
\pgfpathmoveto{\pgfqpoint{4.128516in}{2.697166in}}%
\pgfpathlineto{\pgfqpoint{4.128516in}{2.697166in}}%
\pgfpathlineto{\pgfqpoint{4.128516in}{2.700116in}}%
\pgfpathlineto{\pgfqpoint{4.133057in}{2.700116in}}%
\pgfpathlineto{\pgfqpoint{4.133057in}{2.697166in}}%
\pgfpathmoveto{\pgfqpoint{4.133057in}{2.697166in}}%
\pgfpathlineto{\pgfqpoint{4.133057in}{2.697166in}}%
\pgfpathlineto{\pgfqpoint{4.133057in}{2.700116in}}%
\pgfpathlineto{\pgfqpoint{4.137598in}{2.700116in}}%
\pgfpathlineto{\pgfqpoint{4.137598in}{2.697166in}}%
\pgfpathmoveto{\pgfqpoint{4.133057in}{2.700116in}}%
\pgfpathlineto{\pgfqpoint{4.133057in}{2.700116in}}%
\pgfpathlineto{\pgfqpoint{4.133057in}{2.703065in}}%
\pgfpathlineto{\pgfqpoint{4.137598in}{2.703065in}}%
\pgfpathlineto{\pgfqpoint{4.137598in}{2.700116in}}%
\pgfpathmoveto{\pgfqpoint{4.137598in}{2.700116in}}%
\pgfpathlineto{\pgfqpoint{4.137598in}{2.700116in}}%
\pgfpathlineto{\pgfqpoint{4.137598in}{2.703065in}}%
\pgfpathlineto{\pgfqpoint{4.142139in}{2.703065in}}%
\pgfpathlineto{\pgfqpoint{4.142139in}{2.700116in}}%
\pgfpathmoveto{\pgfqpoint{4.137598in}{2.703065in}}%
\pgfpathlineto{\pgfqpoint{4.137598in}{2.703065in}}%
\pgfpathlineto{\pgfqpoint{4.137598in}{2.706014in}}%
\pgfpathlineto{\pgfqpoint{4.142139in}{2.706014in}}%
\pgfpathlineto{\pgfqpoint{4.142139in}{2.703065in}}%
\pgfpathmoveto{\pgfqpoint{4.142139in}{2.703065in}}%
\pgfpathlineto{\pgfqpoint{4.142139in}{2.703065in}}%
\pgfpathlineto{\pgfqpoint{4.142139in}{2.706014in}}%
\pgfpathlineto{\pgfqpoint{4.146680in}{2.706014in}}%
\pgfpathlineto{\pgfqpoint{4.146680in}{2.703065in}}%
\pgfpathmoveto{\pgfqpoint{4.142139in}{2.706014in}}%
\pgfpathlineto{\pgfqpoint{4.142139in}{2.706014in}}%
\pgfpathlineto{\pgfqpoint{4.142139in}{2.708963in}}%
\pgfpathlineto{\pgfqpoint{4.146680in}{2.708963in}}%
\pgfpathlineto{\pgfqpoint{4.146680in}{2.706014in}}%
\pgfpathmoveto{\pgfqpoint{4.146680in}{2.706014in}}%
\pgfpathlineto{\pgfqpoint{4.146680in}{2.706014in}}%
\pgfpathlineto{\pgfqpoint{4.146680in}{2.708963in}}%
\pgfpathlineto{\pgfqpoint{4.151221in}{2.708963in}}%
\pgfpathlineto{\pgfqpoint{4.151221in}{2.706014in}}%
\pgfpathmoveto{\pgfqpoint{4.146680in}{2.708963in}}%
\pgfpathlineto{\pgfqpoint{4.146680in}{2.708963in}}%
\pgfpathlineto{\pgfqpoint{4.146680in}{2.711912in}}%
\pgfpathlineto{\pgfqpoint{4.151221in}{2.711912in}}%
\pgfpathlineto{\pgfqpoint{4.151221in}{2.708963in}}%
\pgfpathmoveto{\pgfqpoint{4.151221in}{2.708963in}}%
\pgfpathlineto{\pgfqpoint{4.151221in}{2.708963in}}%
\pgfpathlineto{\pgfqpoint{4.151221in}{2.711912in}}%
\pgfpathlineto{\pgfqpoint{4.155762in}{2.711912in}}%
\pgfpathlineto{\pgfqpoint{4.155762in}{2.708963in}}%
\pgfpathmoveto{\pgfqpoint{4.151221in}{2.711912in}}%
\pgfpathlineto{\pgfqpoint{4.151221in}{2.711912in}}%
\pgfpathlineto{\pgfqpoint{4.151221in}{2.714862in}}%
\pgfpathlineto{\pgfqpoint{4.155762in}{2.714862in}}%
\pgfpathlineto{\pgfqpoint{4.155762in}{2.711912in}}%
\pgfpathmoveto{\pgfqpoint{4.155762in}{2.711912in}}%
\pgfpathlineto{\pgfqpoint{4.155762in}{2.711912in}}%
\pgfpathlineto{\pgfqpoint{4.155762in}{2.714862in}}%
\pgfpathlineto{\pgfqpoint{4.160304in}{2.714862in}}%
\pgfpathlineto{\pgfqpoint{4.160304in}{2.711912in}}%
\pgfpathmoveto{\pgfqpoint{4.155762in}{2.714862in}}%
\pgfpathlineto{\pgfqpoint{4.155762in}{2.714862in}}%
\pgfpathlineto{\pgfqpoint{4.155762in}{2.717811in}}%
\pgfpathlineto{\pgfqpoint{4.160304in}{2.717811in}}%
\pgfpathlineto{\pgfqpoint{4.160304in}{2.714862in}}%
\pgfpathmoveto{\pgfqpoint{4.160304in}{2.714862in}}%
\pgfpathlineto{\pgfqpoint{4.160304in}{2.714862in}}%
\pgfpathlineto{\pgfqpoint{4.160304in}{2.717811in}}%
\pgfpathlineto{\pgfqpoint{4.164845in}{2.717811in}}%
\pgfpathlineto{\pgfqpoint{4.164845in}{2.714862in}}%
\pgfpathmoveto{\pgfqpoint{4.160304in}{2.717811in}}%
\pgfpathlineto{\pgfqpoint{4.160304in}{2.717811in}}%
\pgfpathlineto{\pgfqpoint{4.160304in}{2.720760in}}%
\pgfpathlineto{\pgfqpoint{4.164845in}{2.720760in}}%
\pgfpathlineto{\pgfqpoint{4.164845in}{2.717811in}}%
\pgfpathmoveto{\pgfqpoint{4.164845in}{2.717811in}}%
\pgfpathlineto{\pgfqpoint{4.164845in}{2.717811in}}%
\pgfpathlineto{\pgfqpoint{4.164845in}{2.720760in}}%
\pgfpathlineto{\pgfqpoint{4.169386in}{2.720760in}}%
\pgfpathlineto{\pgfqpoint{4.169386in}{2.717811in}}%
\pgfpathmoveto{\pgfqpoint{4.164845in}{2.720760in}}%
\pgfpathlineto{\pgfqpoint{4.164845in}{2.720760in}}%
\pgfpathlineto{\pgfqpoint{4.164845in}{2.723709in}}%
\pgfpathlineto{\pgfqpoint{4.169386in}{2.723709in}}%
\pgfpathlineto{\pgfqpoint{4.169386in}{2.720760in}}%
\pgfpathmoveto{\pgfqpoint{4.169386in}{2.720760in}}%
\pgfpathlineto{\pgfqpoint{4.169386in}{2.720760in}}%
\pgfpathlineto{\pgfqpoint{4.169386in}{2.723709in}}%
\pgfpathlineto{\pgfqpoint{4.173927in}{2.723709in}}%
\pgfpathlineto{\pgfqpoint{4.173927in}{2.720760in}}%
\pgfpathmoveto{\pgfqpoint{4.169386in}{2.723709in}}%
\pgfpathlineto{\pgfqpoint{4.169386in}{2.723709in}}%
\pgfpathlineto{\pgfqpoint{4.169386in}{2.726659in}}%
\pgfpathlineto{\pgfqpoint{4.173927in}{2.726659in}}%
\pgfpathlineto{\pgfqpoint{4.173927in}{2.723709in}}%
\pgfpathmoveto{\pgfqpoint{4.173927in}{2.723709in}}%
\pgfpathlineto{\pgfqpoint{4.173927in}{2.723709in}}%
\pgfpathlineto{\pgfqpoint{4.173927in}{2.726659in}}%
\pgfpathlineto{\pgfqpoint{4.178468in}{2.726659in}}%
\pgfpathlineto{\pgfqpoint{4.178468in}{2.723709in}}%
\pgfpathmoveto{\pgfqpoint{4.173927in}{2.726659in}}%
\pgfpathlineto{\pgfqpoint{4.173927in}{2.726659in}}%
\pgfpathlineto{\pgfqpoint{4.173927in}{2.729608in}}%
\pgfpathlineto{\pgfqpoint{4.178468in}{2.729608in}}%
\pgfpathlineto{\pgfqpoint{4.178468in}{2.726659in}}%
\pgfpathmoveto{\pgfqpoint{4.178468in}{2.726659in}}%
\pgfpathlineto{\pgfqpoint{4.178468in}{2.726659in}}%
\pgfpathlineto{\pgfqpoint{4.178468in}{2.729608in}}%
\pgfpathlineto{\pgfqpoint{4.183009in}{2.729608in}}%
\pgfpathlineto{\pgfqpoint{4.183009in}{2.726659in}}%
\pgfpathmoveto{\pgfqpoint{4.178468in}{2.729608in}}%
\pgfpathlineto{\pgfqpoint{4.178468in}{2.729608in}}%
\pgfpathlineto{\pgfqpoint{4.178468in}{2.732557in}}%
\pgfpathlineto{\pgfqpoint{4.183009in}{2.732557in}}%
\pgfpathlineto{\pgfqpoint{4.183009in}{2.729608in}}%
\pgfpathmoveto{\pgfqpoint{4.183009in}{2.729608in}}%
\pgfpathlineto{\pgfqpoint{4.183009in}{2.729608in}}%
\pgfpathlineto{\pgfqpoint{4.183009in}{2.732557in}}%
\pgfpathlineto{\pgfqpoint{4.187550in}{2.732557in}}%
\pgfpathlineto{\pgfqpoint{4.187550in}{2.729608in}}%
\pgfpathmoveto{\pgfqpoint{4.183009in}{2.732557in}}%
\pgfpathlineto{\pgfqpoint{4.183009in}{2.732557in}}%
\pgfpathlineto{\pgfqpoint{4.183009in}{2.735506in}}%
\pgfpathlineto{\pgfqpoint{4.187550in}{2.735506in}}%
\pgfpathlineto{\pgfqpoint{4.187550in}{2.732557in}}%
\pgfpathmoveto{\pgfqpoint{4.187550in}{2.732557in}}%
\pgfpathlineto{\pgfqpoint{4.187550in}{2.732557in}}%
\pgfpathlineto{\pgfqpoint{4.187550in}{2.735506in}}%
\pgfpathlineto{\pgfqpoint{4.192092in}{2.735506in}}%
\pgfpathlineto{\pgfqpoint{4.192092in}{2.732557in}}%
\pgfpathmoveto{\pgfqpoint{4.187550in}{2.735506in}}%
\pgfpathlineto{\pgfqpoint{4.187550in}{2.735506in}}%
\pgfpathlineto{\pgfqpoint{4.187550in}{2.738455in}}%
\pgfpathlineto{\pgfqpoint{4.192092in}{2.738455in}}%
\pgfpathlineto{\pgfqpoint{4.192092in}{2.735506in}}%
\pgfpathmoveto{\pgfqpoint{4.192092in}{2.735506in}}%
\pgfpathlineto{\pgfqpoint{4.192092in}{2.735506in}}%
\pgfpathlineto{\pgfqpoint{4.192092in}{2.738455in}}%
\pgfpathlineto{\pgfqpoint{4.196633in}{2.738455in}}%
\pgfpathlineto{\pgfqpoint{4.196633in}{2.735506in}}%
\pgfpathmoveto{\pgfqpoint{4.192092in}{2.738455in}}%
\pgfpathlineto{\pgfqpoint{4.192092in}{2.738455in}}%
\pgfpathlineto{\pgfqpoint{4.192092in}{2.741405in}}%
\pgfpathlineto{\pgfqpoint{4.196633in}{2.741405in}}%
\pgfpathlineto{\pgfqpoint{4.196633in}{2.738455in}}%
\pgfpathmoveto{\pgfqpoint{4.196633in}{2.738455in}}%
\pgfpathlineto{\pgfqpoint{4.196633in}{2.738455in}}%
\pgfpathlineto{\pgfqpoint{4.196633in}{2.741405in}}%
\pgfpathlineto{\pgfqpoint{4.201174in}{2.741405in}}%
\pgfpathlineto{\pgfqpoint{4.201174in}{2.738455in}}%
\pgfpathmoveto{\pgfqpoint{4.196633in}{2.741405in}}%
\pgfpathlineto{\pgfqpoint{4.196633in}{2.741405in}}%
\pgfpathlineto{\pgfqpoint{4.196633in}{2.744354in}}%
\pgfpathlineto{\pgfqpoint{4.201174in}{2.744354in}}%
\pgfpathlineto{\pgfqpoint{4.201174in}{2.741405in}}%
\pgfpathmoveto{\pgfqpoint{4.201174in}{2.741405in}}%
\pgfpathlineto{\pgfqpoint{4.201174in}{2.741405in}}%
\pgfpathlineto{\pgfqpoint{4.201174in}{2.744354in}}%
\pgfpathlineto{\pgfqpoint{4.205715in}{2.744354in}}%
\pgfpathlineto{\pgfqpoint{4.205715in}{2.741405in}}%
\pgfpathmoveto{\pgfqpoint{4.201174in}{2.744354in}}%
\pgfpathlineto{\pgfqpoint{4.201174in}{2.744354in}}%
\pgfpathlineto{\pgfqpoint{4.201174in}{2.747303in}}%
\pgfpathlineto{\pgfqpoint{4.205715in}{2.747303in}}%
\pgfpathlineto{\pgfqpoint{4.205715in}{2.744354in}}%
\pgfpathmoveto{\pgfqpoint{4.205715in}{2.744354in}}%
\pgfpathlineto{\pgfqpoint{4.205715in}{2.744354in}}%
\pgfpathlineto{\pgfqpoint{4.205715in}{2.747303in}}%
\pgfpathlineto{\pgfqpoint{4.210256in}{2.747303in}}%
\pgfpathlineto{\pgfqpoint{4.210256in}{2.744354in}}%
\pgfpathmoveto{\pgfqpoint{4.205715in}{2.747303in}}%
\pgfpathlineto{\pgfqpoint{4.205715in}{2.747303in}}%
\pgfpathlineto{\pgfqpoint{4.205715in}{2.750252in}}%
\pgfpathlineto{\pgfqpoint{4.210256in}{2.750252in}}%
\pgfpathlineto{\pgfqpoint{4.210256in}{2.747303in}}%
\pgfpathmoveto{\pgfqpoint{4.210256in}{2.747303in}}%
\pgfpathlineto{\pgfqpoint{4.210256in}{2.747303in}}%
\pgfpathlineto{\pgfqpoint{4.210256in}{2.750252in}}%
\pgfpathlineto{\pgfqpoint{4.214797in}{2.750252in}}%
\pgfpathlineto{\pgfqpoint{4.214797in}{2.747303in}}%
\pgfpathmoveto{\pgfqpoint{4.210256in}{2.750252in}}%
\pgfpathlineto{\pgfqpoint{4.210256in}{2.750252in}}%
\pgfpathlineto{\pgfqpoint{4.210256in}{2.753202in}}%
\pgfpathlineto{\pgfqpoint{4.214797in}{2.753202in}}%
\pgfpathlineto{\pgfqpoint{4.214797in}{2.750252in}}%
\pgfpathmoveto{\pgfqpoint{4.214797in}{2.750252in}}%
\pgfpathlineto{\pgfqpoint{4.214797in}{2.750252in}}%
\pgfpathlineto{\pgfqpoint{4.214797in}{2.753202in}}%
\pgfpathlineto{\pgfqpoint{4.219338in}{2.753202in}}%
\pgfpathlineto{\pgfqpoint{4.219338in}{2.750252in}}%
\pgfpathmoveto{\pgfqpoint{4.214797in}{2.753202in}}%
\pgfpathlineto{\pgfqpoint{4.214797in}{2.753202in}}%
\pgfpathlineto{\pgfqpoint{4.214797in}{2.756151in}}%
\pgfpathlineto{\pgfqpoint{4.219338in}{2.756151in}}%
\pgfpathlineto{\pgfqpoint{4.219338in}{2.753202in}}%
\pgfpathmoveto{\pgfqpoint{4.219338in}{2.753202in}}%
\pgfpathlineto{\pgfqpoint{4.219338in}{2.753202in}}%
\pgfpathlineto{\pgfqpoint{4.219338in}{2.756151in}}%
\pgfpathlineto{\pgfqpoint{4.223880in}{2.756151in}}%
\pgfpathlineto{\pgfqpoint{4.223880in}{2.753202in}}%
\pgfpathmoveto{\pgfqpoint{4.219338in}{2.756151in}}%
\pgfpathlineto{\pgfqpoint{4.219338in}{2.756151in}}%
\pgfpathlineto{\pgfqpoint{4.219338in}{2.759100in}}%
\pgfpathlineto{\pgfqpoint{4.223880in}{2.759100in}}%
\pgfpathlineto{\pgfqpoint{4.223880in}{2.756151in}}%
\pgfpathmoveto{\pgfqpoint{4.223880in}{2.756151in}}%
\pgfpathlineto{\pgfqpoint{4.223880in}{2.756151in}}%
\pgfpathlineto{\pgfqpoint{4.223880in}{2.759100in}}%
\pgfpathlineto{\pgfqpoint{4.228421in}{2.759100in}}%
\pgfpathlineto{\pgfqpoint{4.228421in}{2.756151in}}%
\pgfpathmoveto{\pgfqpoint{4.223880in}{2.759100in}}%
\pgfpathlineto{\pgfqpoint{4.223880in}{2.759100in}}%
\pgfpathlineto{\pgfqpoint{4.223880in}{2.762049in}}%
\pgfpathlineto{\pgfqpoint{4.228421in}{2.762049in}}%
\pgfpathlineto{\pgfqpoint{4.228421in}{2.759100in}}%
\pgfpathmoveto{\pgfqpoint{4.228421in}{2.759100in}}%
\pgfpathlineto{\pgfqpoint{4.228421in}{2.759100in}}%
\pgfpathlineto{\pgfqpoint{4.228421in}{2.762049in}}%
\pgfpathlineto{\pgfqpoint{4.232962in}{2.762049in}}%
\pgfpathlineto{\pgfqpoint{4.232962in}{2.759100in}}%
\pgfpathmoveto{\pgfqpoint{4.228421in}{2.762049in}}%
\pgfpathlineto{\pgfqpoint{4.228421in}{2.762049in}}%
\pgfpathlineto{\pgfqpoint{4.228421in}{2.764998in}}%
\pgfpathlineto{\pgfqpoint{4.232962in}{2.764998in}}%
\pgfpathlineto{\pgfqpoint{4.232962in}{2.762049in}}%
\pgfpathmoveto{\pgfqpoint{4.232962in}{2.762049in}}%
\pgfpathlineto{\pgfqpoint{4.232962in}{2.762049in}}%
\pgfpathlineto{\pgfqpoint{4.232962in}{2.764998in}}%
\pgfpathlineto{\pgfqpoint{4.237503in}{2.764998in}}%
\pgfpathlineto{\pgfqpoint{4.237503in}{2.762049in}}%
\pgfpathmoveto{\pgfqpoint{4.232962in}{2.764998in}}%
\pgfpathlineto{\pgfqpoint{4.232962in}{2.764998in}}%
\pgfpathlineto{\pgfqpoint{4.232962in}{2.767948in}}%
\pgfpathlineto{\pgfqpoint{4.237503in}{2.767948in}}%
\pgfpathlineto{\pgfqpoint{4.237503in}{2.764998in}}%
\pgfpathmoveto{\pgfqpoint{4.237503in}{2.764998in}}%
\pgfpathlineto{\pgfqpoint{4.237503in}{2.764998in}}%
\pgfpathlineto{\pgfqpoint{4.237503in}{2.767948in}}%
\pgfpathlineto{\pgfqpoint{4.242044in}{2.767948in}}%
\pgfpathlineto{\pgfqpoint{4.242044in}{2.764998in}}%
\pgfpathmoveto{\pgfqpoint{4.237503in}{2.767948in}}%
\pgfpathlineto{\pgfqpoint{4.237503in}{2.767948in}}%
\pgfpathlineto{\pgfqpoint{4.237503in}{2.770897in}}%
\pgfpathlineto{\pgfqpoint{4.242044in}{2.770897in}}%
\pgfpathlineto{\pgfqpoint{4.242044in}{2.767948in}}%
\pgfpathmoveto{\pgfqpoint{4.242044in}{2.767948in}}%
\pgfpathlineto{\pgfqpoint{4.242044in}{2.767948in}}%
\pgfpathlineto{\pgfqpoint{4.242044in}{2.770897in}}%
\pgfpathlineto{\pgfqpoint{4.246585in}{2.770897in}}%
\pgfpathlineto{\pgfqpoint{4.246585in}{2.767948in}}%
\pgfpathmoveto{\pgfqpoint{4.242044in}{2.770897in}}%
\pgfpathlineto{\pgfqpoint{4.242044in}{2.770897in}}%
\pgfpathlineto{\pgfqpoint{4.242044in}{2.773846in}}%
\pgfpathlineto{\pgfqpoint{4.246585in}{2.773846in}}%
\pgfpathlineto{\pgfqpoint{4.246585in}{2.770897in}}%
\pgfpathmoveto{\pgfqpoint{4.246585in}{2.770897in}}%
\pgfpathlineto{\pgfqpoint{4.246585in}{2.770897in}}%
\pgfpathlineto{\pgfqpoint{4.246585in}{2.773846in}}%
\pgfpathlineto{\pgfqpoint{4.251125in}{2.773846in}}%
\pgfpathlineto{\pgfqpoint{4.251125in}{2.770897in}}%
\pgfpathmoveto{\pgfqpoint{4.246585in}{2.773846in}}%
\pgfpathlineto{\pgfqpoint{4.246585in}{2.773846in}}%
\pgfpathlineto{\pgfqpoint{4.246585in}{2.776795in}}%
\pgfpathlineto{\pgfqpoint{4.251125in}{2.776795in}}%
\pgfpathlineto{\pgfqpoint{4.251125in}{2.773846in}}%
\pgfpathmoveto{\pgfqpoint{4.251125in}{2.773846in}}%
\pgfpathlineto{\pgfqpoint{4.251125in}{2.773846in}}%
\pgfpathlineto{\pgfqpoint{4.251125in}{2.776795in}}%
\pgfpathlineto{\pgfqpoint{4.255666in}{2.776795in}}%
\pgfpathlineto{\pgfqpoint{4.255666in}{2.773846in}}%
\pgfpathmoveto{\pgfqpoint{4.251125in}{2.776795in}}%
\pgfpathlineto{\pgfqpoint{4.251125in}{2.776795in}}%
\pgfpathlineto{\pgfqpoint{4.251125in}{2.779745in}}%
\pgfpathlineto{\pgfqpoint{4.255666in}{2.779745in}}%
\pgfpathlineto{\pgfqpoint{4.255666in}{2.776795in}}%
\pgfpathmoveto{\pgfqpoint{4.255666in}{2.776795in}}%
\pgfpathlineto{\pgfqpoint{4.255666in}{2.776795in}}%
\pgfpathlineto{\pgfqpoint{4.255666in}{2.779745in}}%
\pgfpathlineto{\pgfqpoint{4.260207in}{2.779745in}}%
\pgfpathlineto{\pgfqpoint{4.260207in}{2.776795in}}%
\pgfpathmoveto{\pgfqpoint{4.255666in}{2.779745in}}%
\pgfpathlineto{\pgfqpoint{4.255666in}{2.779745in}}%
\pgfpathlineto{\pgfqpoint{4.255666in}{2.782694in}}%
\pgfpathlineto{\pgfqpoint{4.260207in}{2.782694in}}%
\pgfpathlineto{\pgfqpoint{4.260207in}{2.779745in}}%
\pgfpathmoveto{\pgfqpoint{4.260207in}{2.779745in}}%
\pgfpathlineto{\pgfqpoint{4.260207in}{2.779745in}}%
\pgfpathlineto{\pgfqpoint{4.260207in}{2.782694in}}%
\pgfpathlineto{\pgfqpoint{4.264748in}{2.782694in}}%
\pgfpathlineto{\pgfqpoint{4.264748in}{2.779745in}}%
\pgfpathmoveto{\pgfqpoint{4.260207in}{2.782694in}}%
\pgfpathlineto{\pgfqpoint{4.260207in}{2.782694in}}%
\pgfpathlineto{\pgfqpoint{4.260207in}{2.785643in}}%
\pgfpathlineto{\pgfqpoint{4.264748in}{2.785643in}}%
\pgfpathlineto{\pgfqpoint{4.264748in}{2.782694in}}%
\pgfpathmoveto{\pgfqpoint{4.264748in}{2.782694in}}%
\pgfpathlineto{\pgfqpoint{4.264748in}{2.782694in}}%
\pgfpathlineto{\pgfqpoint{4.264748in}{2.785643in}}%
\pgfpathlineto{\pgfqpoint{4.269289in}{2.785643in}}%
\pgfpathlineto{\pgfqpoint{4.269289in}{2.782694in}}%
\pgfpathmoveto{\pgfqpoint{4.264748in}{2.785643in}}%
\pgfpathlineto{\pgfqpoint{4.264748in}{2.785643in}}%
\pgfpathlineto{\pgfqpoint{4.264748in}{2.788592in}}%
\pgfpathlineto{\pgfqpoint{4.269289in}{2.788592in}}%
\pgfpathlineto{\pgfqpoint{4.269289in}{2.785643in}}%
\pgfpathmoveto{\pgfqpoint{4.269289in}{2.785643in}}%
\pgfpathlineto{\pgfqpoint{4.269289in}{2.785643in}}%
\pgfpathlineto{\pgfqpoint{4.269289in}{2.788592in}}%
\pgfpathlineto{\pgfqpoint{4.273829in}{2.788592in}}%
\pgfpathlineto{\pgfqpoint{4.273829in}{2.785643in}}%
\pgfpathmoveto{\pgfqpoint{4.269289in}{2.788592in}}%
\pgfpathlineto{\pgfqpoint{4.269289in}{2.788592in}}%
\pgfpathlineto{\pgfqpoint{4.269289in}{2.791542in}}%
\pgfpathlineto{\pgfqpoint{4.273829in}{2.791542in}}%
\pgfpathlineto{\pgfqpoint{4.273829in}{2.788592in}}%
\pgfpathmoveto{\pgfqpoint{4.273829in}{2.788592in}}%
\pgfpathlineto{\pgfqpoint{4.273829in}{2.788592in}}%
\pgfpathlineto{\pgfqpoint{4.273829in}{2.791542in}}%
\pgfpathlineto{\pgfqpoint{4.278370in}{2.791542in}}%
\pgfpathlineto{\pgfqpoint{4.278370in}{2.788592in}}%
\pgfpathmoveto{\pgfqpoint{4.273829in}{2.791542in}}%
\pgfpathlineto{\pgfqpoint{4.273829in}{2.791542in}}%
\pgfpathlineto{\pgfqpoint{4.273829in}{2.794491in}}%
\pgfpathlineto{\pgfqpoint{4.278370in}{2.794491in}}%
\pgfpathlineto{\pgfqpoint{4.278370in}{2.791542in}}%
\pgfpathmoveto{\pgfqpoint{4.278370in}{2.791542in}}%
\pgfpathlineto{\pgfqpoint{4.278370in}{2.791542in}}%
\pgfpathlineto{\pgfqpoint{4.278370in}{2.794491in}}%
\pgfpathlineto{\pgfqpoint{4.282911in}{2.794491in}}%
\pgfpathlineto{\pgfqpoint{4.282911in}{2.791542in}}%
\pgfpathmoveto{\pgfqpoint{4.278370in}{2.794491in}}%
\pgfpathlineto{\pgfqpoint{4.278370in}{2.794491in}}%
\pgfpathlineto{\pgfqpoint{4.278370in}{2.797440in}}%
\pgfpathlineto{\pgfqpoint{4.282911in}{2.797440in}}%
\pgfpathlineto{\pgfqpoint{4.282911in}{2.794491in}}%
\pgfpathmoveto{\pgfqpoint{4.282911in}{2.794491in}}%
\pgfpathlineto{\pgfqpoint{4.282911in}{2.794491in}}%
\pgfpathlineto{\pgfqpoint{4.282911in}{2.797440in}}%
\pgfpathlineto{\pgfqpoint{4.287452in}{2.797440in}}%
\pgfpathlineto{\pgfqpoint{4.287452in}{2.794491in}}%
\pgfpathmoveto{\pgfqpoint{4.282911in}{2.797440in}}%
\pgfpathlineto{\pgfqpoint{4.282911in}{2.797440in}}%
\pgfpathlineto{\pgfqpoint{4.282911in}{2.800390in}}%
\pgfpathlineto{\pgfqpoint{4.287452in}{2.800390in}}%
\pgfpathlineto{\pgfqpoint{4.287452in}{2.797440in}}%
\pgfpathmoveto{\pgfqpoint{4.287452in}{2.797440in}}%
\pgfpathlineto{\pgfqpoint{4.287452in}{2.797440in}}%
\pgfpathlineto{\pgfqpoint{4.287452in}{2.800390in}}%
\pgfpathlineto{\pgfqpoint{4.291993in}{2.800390in}}%
\pgfpathlineto{\pgfqpoint{4.291993in}{2.797440in}}%
\pgfpathmoveto{\pgfqpoint{4.287452in}{2.800390in}}%
\pgfpathlineto{\pgfqpoint{4.287452in}{2.800390in}}%
\pgfpathlineto{\pgfqpoint{4.287452in}{2.803339in}}%
\pgfpathlineto{\pgfqpoint{4.291993in}{2.803339in}}%
\pgfpathlineto{\pgfqpoint{4.291993in}{2.800390in}}%
\pgfpathmoveto{\pgfqpoint{4.291993in}{2.800390in}}%
\pgfpathlineto{\pgfqpoint{4.291993in}{2.800390in}}%
\pgfpathlineto{\pgfqpoint{4.291993in}{2.803339in}}%
\pgfpathlineto{\pgfqpoint{4.296533in}{2.803339in}}%
\pgfpathlineto{\pgfqpoint{4.296533in}{2.800390in}}%
\pgfpathmoveto{\pgfqpoint{4.291993in}{2.803339in}}%
\pgfpathlineto{\pgfqpoint{4.291993in}{2.803339in}}%
\pgfpathlineto{\pgfqpoint{4.291993in}{2.806288in}}%
\pgfpathlineto{\pgfqpoint{4.296533in}{2.806288in}}%
\pgfpathlineto{\pgfqpoint{4.296533in}{2.803339in}}%
\pgfpathmoveto{\pgfqpoint{4.296533in}{2.803339in}}%
\pgfpathlineto{\pgfqpoint{4.296533in}{2.803339in}}%
\pgfpathlineto{\pgfqpoint{4.296533in}{2.806288in}}%
\pgfpathlineto{\pgfqpoint{4.301074in}{2.806288in}}%
\pgfpathlineto{\pgfqpoint{4.301074in}{2.803339in}}%
\pgfpathmoveto{\pgfqpoint{4.296533in}{2.806288in}}%
\pgfpathlineto{\pgfqpoint{4.296533in}{2.806288in}}%
\pgfpathlineto{\pgfqpoint{4.296533in}{2.809237in}}%
\pgfpathlineto{\pgfqpoint{4.301074in}{2.809237in}}%
\pgfpathlineto{\pgfqpoint{4.301074in}{2.806288in}}%
\pgfpathmoveto{\pgfqpoint{4.301074in}{2.806288in}}%
\pgfpathlineto{\pgfqpoint{4.301074in}{2.806288in}}%
\pgfpathlineto{\pgfqpoint{4.301074in}{2.809237in}}%
\pgfpathlineto{\pgfqpoint{4.305615in}{2.809237in}}%
\pgfpathlineto{\pgfqpoint{4.305615in}{2.806288in}}%
\pgfpathmoveto{\pgfqpoint{4.301074in}{2.809237in}}%
\pgfpathlineto{\pgfqpoint{4.301074in}{2.809237in}}%
\pgfpathlineto{\pgfqpoint{4.301074in}{2.812187in}}%
\pgfpathlineto{\pgfqpoint{4.305615in}{2.812187in}}%
\pgfpathlineto{\pgfqpoint{4.305615in}{2.809237in}}%
\pgfpathmoveto{\pgfqpoint{4.305615in}{2.809237in}}%
\pgfpathlineto{\pgfqpoint{4.305615in}{2.809237in}}%
\pgfpathlineto{\pgfqpoint{4.305615in}{2.812187in}}%
\pgfpathlineto{\pgfqpoint{4.310156in}{2.812187in}}%
\pgfpathlineto{\pgfqpoint{4.310156in}{2.809237in}}%
\pgfpathmoveto{\pgfqpoint{4.305615in}{2.812187in}}%
\pgfpathlineto{\pgfqpoint{4.305615in}{2.812187in}}%
\pgfpathlineto{\pgfqpoint{4.305615in}{2.815136in}}%
\pgfpathlineto{\pgfqpoint{4.310156in}{2.815136in}}%
\pgfpathlineto{\pgfqpoint{4.310156in}{2.812187in}}%
\pgfpathmoveto{\pgfqpoint{4.310156in}{2.812187in}}%
\pgfpathlineto{\pgfqpoint{4.310156in}{2.812187in}}%
\pgfpathlineto{\pgfqpoint{4.310156in}{2.815136in}}%
\pgfpathlineto{\pgfqpoint{4.314697in}{2.815136in}}%
\pgfpathlineto{\pgfqpoint{4.314697in}{2.812187in}}%
\pgfpathmoveto{\pgfqpoint{4.310156in}{2.815136in}}%
\pgfpathlineto{\pgfqpoint{4.310156in}{2.815136in}}%
\pgfpathlineto{\pgfqpoint{4.310156in}{2.818085in}}%
\pgfpathlineto{\pgfqpoint{4.314697in}{2.818085in}}%
\pgfpathlineto{\pgfqpoint{4.314697in}{2.815136in}}%
\pgfpathmoveto{\pgfqpoint{4.314697in}{2.815136in}}%
\pgfpathlineto{\pgfqpoint{4.314697in}{2.815136in}}%
\pgfpathlineto{\pgfqpoint{4.314697in}{2.818085in}}%
\pgfpathlineto{\pgfqpoint{4.319237in}{2.818085in}}%
\pgfpathlineto{\pgfqpoint{4.319237in}{2.815136in}}%
\pgfpathmoveto{\pgfqpoint{4.314697in}{2.818085in}}%
\pgfpathlineto{\pgfqpoint{4.314697in}{2.818085in}}%
\pgfpathlineto{\pgfqpoint{4.314697in}{2.821034in}}%
\pgfpathlineto{\pgfqpoint{4.319237in}{2.821034in}}%
\pgfpathlineto{\pgfqpoint{4.319237in}{2.818085in}}%
\pgfpathmoveto{\pgfqpoint{4.319237in}{2.818085in}}%
\pgfpathlineto{\pgfqpoint{4.319237in}{2.818085in}}%
\pgfpathlineto{\pgfqpoint{4.319237in}{2.821034in}}%
\pgfpathlineto{\pgfqpoint{4.323778in}{2.821034in}}%
\pgfpathlineto{\pgfqpoint{4.323778in}{2.818085in}}%
\pgfpathmoveto{\pgfqpoint{4.323778in}{2.818085in}}%
\pgfpathlineto{\pgfqpoint{4.323778in}{2.818085in}}%
\pgfpathlineto{\pgfqpoint{4.323778in}{2.821034in}}%
\pgfpathlineto{\pgfqpoint{4.328319in}{2.821034in}}%
\pgfpathlineto{\pgfqpoint{4.328319in}{2.818085in}}%
\pgfpathmoveto{\pgfqpoint{4.323778in}{2.821034in}}%
\pgfpathlineto{\pgfqpoint{4.323778in}{2.821034in}}%
\pgfpathlineto{\pgfqpoint{4.323778in}{2.823984in}}%
\pgfpathlineto{\pgfqpoint{4.328319in}{2.823984in}}%
\pgfpathlineto{\pgfqpoint{4.328319in}{2.821034in}}%
\pgfpathmoveto{\pgfqpoint{4.328319in}{2.821034in}}%
\pgfpathlineto{\pgfqpoint{4.328319in}{2.821034in}}%
\pgfpathlineto{\pgfqpoint{4.328319in}{2.823984in}}%
\pgfpathlineto{\pgfqpoint{4.332860in}{2.823984in}}%
\pgfpathlineto{\pgfqpoint{4.332860in}{2.821034in}}%
\pgfpathmoveto{\pgfqpoint{4.328319in}{2.823984in}}%
\pgfpathlineto{\pgfqpoint{4.328319in}{2.823984in}}%
\pgfpathlineto{\pgfqpoint{4.328319in}{2.826933in}}%
\pgfpathlineto{\pgfqpoint{4.332860in}{2.826933in}}%
\pgfpathlineto{\pgfqpoint{4.332860in}{2.823984in}}%
\pgfpathmoveto{\pgfqpoint{4.332860in}{2.823984in}}%
\pgfpathlineto{\pgfqpoint{4.332860in}{2.823984in}}%
\pgfpathlineto{\pgfqpoint{4.332860in}{2.826933in}}%
\pgfpathlineto{\pgfqpoint{4.337400in}{2.826933in}}%
\pgfpathlineto{\pgfqpoint{4.337400in}{2.823984in}}%
\pgfpathmoveto{\pgfqpoint{4.332860in}{2.826933in}}%
\pgfpathlineto{\pgfqpoint{4.332860in}{2.826933in}}%
\pgfpathlineto{\pgfqpoint{4.332860in}{2.829882in}}%
\pgfpathlineto{\pgfqpoint{4.337400in}{2.829882in}}%
\pgfpathlineto{\pgfqpoint{4.337400in}{2.826933in}}%
\pgfpathmoveto{\pgfqpoint{4.337400in}{2.826933in}}%
\pgfpathlineto{\pgfqpoint{4.337400in}{2.826933in}}%
\pgfpathlineto{\pgfqpoint{4.337400in}{2.829882in}}%
\pgfpathlineto{\pgfqpoint{4.341941in}{2.829882in}}%
\pgfpathlineto{\pgfqpoint{4.341941in}{2.826933in}}%
\pgfpathmoveto{\pgfqpoint{4.337400in}{2.829882in}}%
\pgfpathlineto{\pgfqpoint{4.337400in}{2.829882in}}%
\pgfpathlineto{\pgfqpoint{4.337400in}{2.832831in}}%
\pgfpathlineto{\pgfqpoint{4.341941in}{2.832831in}}%
\pgfpathlineto{\pgfqpoint{4.341941in}{2.829882in}}%
\pgfpathmoveto{\pgfqpoint{4.341941in}{2.829882in}}%
\pgfpathlineto{\pgfqpoint{4.341941in}{2.829882in}}%
\pgfpathlineto{\pgfqpoint{4.341941in}{2.832831in}}%
\pgfpathlineto{\pgfqpoint{4.346482in}{2.832831in}}%
\pgfpathlineto{\pgfqpoint{4.346482in}{2.829882in}}%
\pgfpathmoveto{\pgfqpoint{4.341941in}{2.832831in}}%
\pgfpathlineto{\pgfqpoint{4.341941in}{2.832831in}}%
\pgfpathlineto{\pgfqpoint{4.341941in}{2.835781in}}%
\pgfpathlineto{\pgfqpoint{4.346482in}{2.835781in}}%
\pgfpathlineto{\pgfqpoint{4.346482in}{2.832831in}}%
\pgfpathmoveto{\pgfqpoint{4.346482in}{2.832831in}}%
\pgfpathlineto{\pgfqpoint{4.346482in}{2.832831in}}%
\pgfpathlineto{\pgfqpoint{4.346482in}{2.835781in}}%
\pgfpathlineto{\pgfqpoint{4.351023in}{2.835781in}}%
\pgfpathlineto{\pgfqpoint{4.351023in}{2.832831in}}%
\pgfpathmoveto{\pgfqpoint{4.346482in}{2.835781in}}%
\pgfpathlineto{\pgfqpoint{4.346482in}{2.835781in}}%
\pgfpathlineto{\pgfqpoint{4.346482in}{2.838730in}}%
\pgfpathlineto{\pgfqpoint{4.351023in}{2.838730in}}%
\pgfpathlineto{\pgfqpoint{4.351023in}{2.835781in}}%
\pgfpathmoveto{\pgfqpoint{4.351023in}{2.835781in}}%
\pgfpathlineto{\pgfqpoint{4.351023in}{2.835781in}}%
\pgfpathlineto{\pgfqpoint{4.351023in}{2.838730in}}%
\pgfpathlineto{\pgfqpoint{4.355564in}{2.838730in}}%
\pgfpathlineto{\pgfqpoint{4.355564in}{2.835781in}}%
\pgfpathmoveto{\pgfqpoint{4.351023in}{2.838730in}}%
\pgfpathlineto{\pgfqpoint{4.351023in}{2.838730in}}%
\pgfpathlineto{\pgfqpoint{4.351023in}{2.841679in}}%
\pgfpathlineto{\pgfqpoint{4.355564in}{2.841679in}}%
\pgfpathlineto{\pgfqpoint{4.355564in}{2.838730in}}%
\pgfpathmoveto{\pgfqpoint{4.355564in}{2.838730in}}%
\pgfpathlineto{\pgfqpoint{4.355564in}{2.838730in}}%
\pgfpathlineto{\pgfqpoint{4.355564in}{2.841679in}}%
\pgfpathlineto{\pgfqpoint{4.360104in}{2.841679in}}%
\pgfpathlineto{\pgfqpoint{4.360104in}{2.838730in}}%
\pgfpathmoveto{\pgfqpoint{4.355564in}{2.841679in}}%
\pgfpathlineto{\pgfqpoint{4.355564in}{2.841679in}}%
\pgfpathlineto{\pgfqpoint{4.355564in}{2.844629in}}%
\pgfpathlineto{\pgfqpoint{4.360104in}{2.844629in}}%
\pgfpathlineto{\pgfqpoint{4.360104in}{2.841679in}}%
\pgfpathmoveto{\pgfqpoint{4.360104in}{2.841679in}}%
\pgfpathlineto{\pgfqpoint{4.360104in}{2.841679in}}%
\pgfpathlineto{\pgfqpoint{4.360104in}{2.844629in}}%
\pgfpathlineto{\pgfqpoint{4.364645in}{2.844629in}}%
\pgfpathlineto{\pgfqpoint{4.364645in}{2.841679in}}%
\pgfpathmoveto{\pgfqpoint{4.360104in}{2.844629in}}%
\pgfpathlineto{\pgfqpoint{4.360104in}{2.844629in}}%
\pgfpathlineto{\pgfqpoint{4.360104in}{2.847578in}}%
\pgfpathlineto{\pgfqpoint{4.364645in}{2.847578in}}%
\pgfpathlineto{\pgfqpoint{4.364645in}{2.844629in}}%
\pgfpathmoveto{\pgfqpoint{4.364645in}{2.844629in}}%
\pgfpathlineto{\pgfqpoint{4.364645in}{2.844629in}}%
\pgfpathlineto{\pgfqpoint{4.364645in}{2.847578in}}%
\pgfpathlineto{\pgfqpoint{4.369186in}{2.847578in}}%
\pgfpathlineto{\pgfqpoint{4.369186in}{2.844629in}}%
\pgfpathmoveto{\pgfqpoint{4.364645in}{2.847578in}}%
\pgfpathlineto{\pgfqpoint{4.364645in}{2.847578in}}%
\pgfpathlineto{\pgfqpoint{4.364645in}{2.850527in}}%
\pgfpathlineto{\pgfqpoint{4.369186in}{2.850527in}}%
\pgfpathlineto{\pgfqpoint{4.369186in}{2.847578in}}%
\pgfpathmoveto{\pgfqpoint{4.369186in}{2.847578in}}%
\pgfpathlineto{\pgfqpoint{4.369186in}{2.847578in}}%
\pgfpathlineto{\pgfqpoint{4.369186in}{2.850527in}}%
\pgfpathlineto{\pgfqpoint{4.373727in}{2.850527in}}%
\pgfpathlineto{\pgfqpoint{4.373727in}{2.847578in}}%
\pgfpathmoveto{\pgfqpoint{4.369186in}{2.850527in}}%
\pgfpathlineto{\pgfqpoint{4.369186in}{2.850527in}}%
\pgfpathlineto{\pgfqpoint{4.369186in}{2.853476in}}%
\pgfpathlineto{\pgfqpoint{4.373727in}{2.853476in}}%
\pgfpathlineto{\pgfqpoint{4.373727in}{2.850527in}}%
\pgfpathmoveto{\pgfqpoint{4.373727in}{2.850527in}}%
\pgfpathlineto{\pgfqpoint{4.373727in}{2.850527in}}%
\pgfpathlineto{\pgfqpoint{4.373727in}{2.853476in}}%
\pgfpathlineto{\pgfqpoint{4.378268in}{2.853476in}}%
\pgfpathlineto{\pgfqpoint{4.378268in}{2.850527in}}%
\pgfpathmoveto{\pgfqpoint{4.373727in}{2.853476in}}%
\pgfpathlineto{\pgfqpoint{4.373727in}{2.853476in}}%
\pgfpathlineto{\pgfqpoint{4.373727in}{2.856426in}}%
\pgfpathlineto{\pgfqpoint{4.378268in}{2.856426in}}%
\pgfpathlineto{\pgfqpoint{4.378268in}{2.853476in}}%
\pgfpathmoveto{\pgfqpoint{4.378268in}{2.853476in}}%
\pgfpathlineto{\pgfqpoint{4.378268in}{2.853476in}}%
\pgfpathlineto{\pgfqpoint{4.378268in}{2.856426in}}%
\pgfpathlineto{\pgfqpoint{4.382808in}{2.856426in}}%
\pgfpathlineto{\pgfqpoint{4.382808in}{2.853476in}}%
\pgfpathmoveto{\pgfqpoint{4.378268in}{2.856426in}}%
\pgfpathlineto{\pgfqpoint{4.378268in}{2.856426in}}%
\pgfpathlineto{\pgfqpoint{4.378268in}{2.859375in}}%
\pgfpathlineto{\pgfqpoint{4.382808in}{2.859375in}}%
\pgfpathlineto{\pgfqpoint{4.382808in}{2.856426in}}%
\pgfpathmoveto{\pgfqpoint{4.382808in}{2.856426in}}%
\pgfpathlineto{\pgfqpoint{4.382808in}{2.856426in}}%
\pgfpathlineto{\pgfqpoint{4.382808in}{2.859375in}}%
\pgfpathlineto{\pgfqpoint{4.387350in}{2.859375in}}%
\pgfpathlineto{\pgfqpoint{4.387350in}{2.856426in}}%
\pgfpathmoveto{\pgfqpoint{4.382808in}{2.859375in}}%
\pgfpathlineto{\pgfqpoint{4.382808in}{2.859375in}}%
\pgfpathlineto{\pgfqpoint{4.382808in}{2.862324in}}%
\pgfpathlineto{\pgfqpoint{4.387350in}{2.862324in}}%
\pgfpathlineto{\pgfqpoint{4.387350in}{2.859375in}}%
\pgfpathmoveto{\pgfqpoint{4.387350in}{2.859375in}}%
\pgfpathlineto{\pgfqpoint{4.387350in}{2.859375in}}%
\pgfpathlineto{\pgfqpoint{4.387350in}{2.862324in}}%
\pgfpathlineto{\pgfqpoint{4.391891in}{2.862324in}}%
\pgfpathlineto{\pgfqpoint{4.391891in}{2.859375in}}%
\pgfpathmoveto{\pgfqpoint{4.387350in}{2.862324in}}%
\pgfpathlineto{\pgfqpoint{4.387350in}{2.862324in}}%
\pgfpathlineto{\pgfqpoint{4.387350in}{2.865273in}}%
\pgfpathlineto{\pgfqpoint{4.391891in}{2.865273in}}%
\pgfpathlineto{\pgfqpoint{4.391891in}{2.862324in}}%
\pgfpathmoveto{\pgfqpoint{4.391891in}{2.862324in}}%
\pgfpathlineto{\pgfqpoint{4.391891in}{2.862324in}}%
\pgfpathlineto{\pgfqpoint{4.391891in}{2.865273in}}%
\pgfpathlineto{\pgfqpoint{4.396432in}{2.865273in}}%
\pgfpathlineto{\pgfqpoint{4.396432in}{2.862324in}}%
\pgfpathmoveto{\pgfqpoint{4.391891in}{2.865273in}}%
\pgfpathlineto{\pgfqpoint{4.391891in}{2.865273in}}%
\pgfpathlineto{\pgfqpoint{4.391891in}{2.868222in}}%
\pgfpathlineto{\pgfqpoint{4.396432in}{2.868222in}}%
\pgfpathlineto{\pgfqpoint{4.396432in}{2.865273in}}%
\pgfpathmoveto{\pgfqpoint{4.396432in}{2.865273in}}%
\pgfpathlineto{\pgfqpoint{4.396432in}{2.865273in}}%
\pgfpathlineto{\pgfqpoint{4.396432in}{2.868222in}}%
\pgfpathlineto{\pgfqpoint{4.400973in}{2.868222in}}%
\pgfpathlineto{\pgfqpoint{4.400973in}{2.865273in}}%
\pgfpathmoveto{\pgfqpoint{4.396432in}{2.868222in}}%
\pgfpathlineto{\pgfqpoint{4.396432in}{2.868222in}}%
\pgfpathlineto{\pgfqpoint{4.396432in}{2.871172in}}%
\pgfpathlineto{\pgfqpoint{4.400973in}{2.871172in}}%
\pgfpathlineto{\pgfqpoint{4.400973in}{2.868222in}}%
\pgfpathmoveto{\pgfqpoint{4.400973in}{2.868222in}}%
\pgfpathlineto{\pgfqpoint{4.400973in}{2.868222in}}%
\pgfpathlineto{\pgfqpoint{4.400973in}{2.871172in}}%
\pgfpathlineto{\pgfqpoint{4.405514in}{2.871172in}}%
\pgfpathlineto{\pgfqpoint{4.405514in}{2.868222in}}%
\pgfpathmoveto{\pgfqpoint{4.400973in}{2.871172in}}%
\pgfpathlineto{\pgfqpoint{4.400973in}{2.871172in}}%
\pgfpathlineto{\pgfqpoint{4.400973in}{2.874121in}}%
\pgfpathlineto{\pgfqpoint{4.405514in}{2.874121in}}%
\pgfpathlineto{\pgfqpoint{4.405514in}{2.871172in}}%
\pgfpathmoveto{\pgfqpoint{4.405514in}{2.871172in}}%
\pgfpathlineto{\pgfqpoint{4.405514in}{2.871172in}}%
\pgfpathlineto{\pgfqpoint{4.405514in}{2.874121in}}%
\pgfpathlineto{\pgfqpoint{4.410055in}{2.874121in}}%
\pgfpathlineto{\pgfqpoint{4.410055in}{2.871172in}}%
\pgfpathmoveto{\pgfqpoint{4.405514in}{2.874121in}}%
\pgfpathlineto{\pgfqpoint{4.405514in}{2.874121in}}%
\pgfpathlineto{\pgfqpoint{4.405514in}{2.877070in}}%
\pgfpathlineto{\pgfqpoint{4.410055in}{2.877070in}}%
\pgfpathlineto{\pgfqpoint{4.410055in}{2.874121in}}%
\pgfpathmoveto{\pgfqpoint{4.410055in}{2.874121in}}%
\pgfpathlineto{\pgfqpoint{4.410055in}{2.874121in}}%
\pgfpathlineto{\pgfqpoint{4.410055in}{2.877070in}}%
\pgfpathlineto{\pgfqpoint{4.414597in}{2.877070in}}%
\pgfpathlineto{\pgfqpoint{4.414597in}{2.874121in}}%
\pgfpathmoveto{\pgfqpoint{4.410055in}{2.877070in}}%
\pgfpathlineto{\pgfqpoint{4.410055in}{2.877070in}}%
\pgfpathlineto{\pgfqpoint{4.410055in}{2.880019in}}%
\pgfpathlineto{\pgfqpoint{4.414597in}{2.880019in}}%
\pgfpathlineto{\pgfqpoint{4.414597in}{2.877070in}}%
\pgfpathmoveto{\pgfqpoint{4.414597in}{2.877070in}}%
\pgfpathlineto{\pgfqpoint{4.414597in}{2.877070in}}%
\pgfpathlineto{\pgfqpoint{4.414597in}{2.880019in}}%
\pgfpathlineto{\pgfqpoint{4.419138in}{2.880019in}}%
\pgfpathlineto{\pgfqpoint{4.419138in}{2.877070in}}%
\pgfpathmoveto{\pgfqpoint{4.414597in}{2.880019in}}%
\pgfpathlineto{\pgfqpoint{4.414597in}{2.880019in}}%
\pgfpathlineto{\pgfqpoint{4.414597in}{2.882968in}}%
\pgfpathlineto{\pgfqpoint{4.419138in}{2.882968in}}%
\pgfpathlineto{\pgfqpoint{4.419138in}{2.880019in}}%
\pgfpathmoveto{\pgfqpoint{4.419138in}{2.880019in}}%
\pgfpathlineto{\pgfqpoint{4.419138in}{2.880019in}}%
\pgfpathlineto{\pgfqpoint{4.419138in}{2.882968in}}%
\pgfpathlineto{\pgfqpoint{4.423679in}{2.882968in}}%
\pgfpathlineto{\pgfqpoint{4.423679in}{2.880019in}}%
\pgfpathmoveto{\pgfqpoint{4.419138in}{2.882968in}}%
\pgfpathlineto{\pgfqpoint{4.419138in}{2.882968in}}%
\pgfpathlineto{\pgfqpoint{4.419138in}{2.885918in}}%
\pgfpathlineto{\pgfqpoint{4.423679in}{2.885918in}}%
\pgfpathlineto{\pgfqpoint{4.423679in}{2.882968in}}%
\pgfpathmoveto{\pgfqpoint{4.423679in}{2.882968in}}%
\pgfpathlineto{\pgfqpoint{4.423679in}{2.882968in}}%
\pgfpathlineto{\pgfqpoint{4.423679in}{2.885918in}}%
\pgfpathlineto{\pgfqpoint{4.428220in}{2.885918in}}%
\pgfpathlineto{\pgfqpoint{4.428220in}{2.882968in}}%
\pgfpathmoveto{\pgfqpoint{4.423679in}{2.885918in}}%
\pgfpathlineto{\pgfqpoint{4.423679in}{2.885918in}}%
\pgfpathlineto{\pgfqpoint{4.423679in}{2.888867in}}%
\pgfpathlineto{\pgfqpoint{4.428220in}{2.888867in}}%
\pgfpathlineto{\pgfqpoint{4.428220in}{2.885918in}}%
\pgfpathmoveto{\pgfqpoint{4.428220in}{2.885918in}}%
\pgfpathlineto{\pgfqpoint{4.428220in}{2.885918in}}%
\pgfpathlineto{\pgfqpoint{4.428220in}{2.888867in}}%
\pgfpathlineto{\pgfqpoint{4.432761in}{2.888867in}}%
\pgfpathlineto{\pgfqpoint{4.432761in}{2.885918in}}%
\pgfpathmoveto{\pgfqpoint{4.428220in}{2.888867in}}%
\pgfpathlineto{\pgfqpoint{4.428220in}{2.888867in}}%
\pgfpathlineto{\pgfqpoint{4.428220in}{2.891816in}}%
\pgfpathlineto{\pgfqpoint{4.432761in}{2.891816in}}%
\pgfpathlineto{\pgfqpoint{4.432761in}{2.888867in}}%
\pgfpathmoveto{\pgfqpoint{4.432761in}{2.888867in}}%
\pgfpathlineto{\pgfqpoint{4.432761in}{2.888867in}}%
\pgfpathlineto{\pgfqpoint{4.432761in}{2.891816in}}%
\pgfpathlineto{\pgfqpoint{4.437303in}{2.891816in}}%
\pgfpathlineto{\pgfqpoint{4.437303in}{2.888867in}}%
\pgfpathmoveto{\pgfqpoint{4.432761in}{2.891816in}}%
\pgfpathlineto{\pgfqpoint{4.432761in}{2.891816in}}%
\pgfpathlineto{\pgfqpoint{4.432761in}{2.894765in}}%
\pgfpathlineto{\pgfqpoint{4.437303in}{2.894765in}}%
\pgfpathlineto{\pgfqpoint{4.437303in}{2.891816in}}%
\pgfpathmoveto{\pgfqpoint{4.437303in}{2.891816in}}%
\pgfpathlineto{\pgfqpoint{4.437303in}{2.891816in}}%
\pgfpathlineto{\pgfqpoint{4.437303in}{2.894765in}}%
\pgfpathlineto{\pgfqpoint{4.441844in}{2.894765in}}%
\pgfpathlineto{\pgfqpoint{4.441844in}{2.891816in}}%
\pgfpathmoveto{\pgfqpoint{4.437303in}{2.894765in}}%
\pgfpathlineto{\pgfqpoint{4.437303in}{2.894765in}}%
\pgfpathlineto{\pgfqpoint{4.437303in}{2.897714in}}%
\pgfpathlineto{\pgfqpoint{4.441844in}{2.897714in}}%
\pgfpathlineto{\pgfqpoint{4.441844in}{2.894765in}}%
\pgfpathmoveto{\pgfqpoint{4.441844in}{2.894765in}}%
\pgfpathlineto{\pgfqpoint{4.441844in}{2.894765in}}%
\pgfpathlineto{\pgfqpoint{4.441844in}{2.897714in}}%
\pgfpathlineto{\pgfqpoint{4.446385in}{2.897714in}}%
\pgfpathlineto{\pgfqpoint{4.446385in}{2.894765in}}%
\pgfpathmoveto{\pgfqpoint{4.441844in}{2.897714in}}%
\pgfpathlineto{\pgfqpoint{4.441844in}{2.897714in}}%
\pgfpathlineto{\pgfqpoint{4.441844in}{2.900664in}}%
\pgfpathlineto{\pgfqpoint{4.446385in}{2.900664in}}%
\pgfpathlineto{\pgfqpoint{4.446385in}{2.897714in}}%
\pgfpathmoveto{\pgfqpoint{4.446385in}{2.897714in}}%
\pgfpathlineto{\pgfqpoint{4.446385in}{2.897714in}}%
\pgfpathlineto{\pgfqpoint{4.446385in}{2.900664in}}%
\pgfpathlineto{\pgfqpoint{4.450926in}{2.900664in}}%
\pgfpathlineto{\pgfqpoint{4.450926in}{2.897714in}}%
\pgfpathmoveto{\pgfqpoint{4.446385in}{2.900664in}}%
\pgfpathlineto{\pgfqpoint{4.446385in}{2.900664in}}%
\pgfpathlineto{\pgfqpoint{4.446385in}{2.903613in}}%
\pgfpathlineto{\pgfqpoint{4.450926in}{2.903613in}}%
\pgfpathlineto{\pgfqpoint{4.450926in}{2.900664in}}%
\pgfpathmoveto{\pgfqpoint{4.450926in}{2.900664in}}%
\pgfpathlineto{\pgfqpoint{4.450926in}{2.900664in}}%
\pgfpathlineto{\pgfqpoint{4.450926in}{2.903613in}}%
\pgfpathlineto{\pgfqpoint{4.455467in}{2.903613in}}%
\pgfpathlineto{\pgfqpoint{4.455467in}{2.900664in}}%
\pgfpathmoveto{\pgfqpoint{4.450926in}{2.903613in}}%
\pgfpathlineto{\pgfqpoint{4.450926in}{2.903613in}}%
\pgfpathlineto{\pgfqpoint{4.450926in}{2.906562in}}%
\pgfpathlineto{\pgfqpoint{4.455467in}{2.906562in}}%
\pgfpathlineto{\pgfqpoint{4.455467in}{2.903613in}}%
\pgfpathmoveto{\pgfqpoint{4.455467in}{2.903613in}}%
\pgfpathlineto{\pgfqpoint{4.455467in}{2.903613in}}%
\pgfpathlineto{\pgfqpoint{4.455467in}{2.906562in}}%
\pgfpathlineto{\pgfqpoint{4.460008in}{2.906562in}}%
\pgfpathlineto{\pgfqpoint{4.460008in}{2.903613in}}%
\pgfpathmoveto{\pgfqpoint{4.455467in}{2.906562in}}%
\pgfpathlineto{\pgfqpoint{4.455467in}{2.906562in}}%
\pgfpathlineto{\pgfqpoint{4.455467in}{2.909511in}}%
\pgfpathlineto{\pgfqpoint{4.460008in}{2.909511in}}%
\pgfpathlineto{\pgfqpoint{4.460008in}{2.906562in}}%
\pgfpathmoveto{\pgfqpoint{4.460008in}{2.906562in}}%
\pgfpathlineto{\pgfqpoint{4.460008in}{2.906562in}}%
\pgfpathlineto{\pgfqpoint{4.460008in}{2.909511in}}%
\pgfpathlineto{\pgfqpoint{4.464550in}{2.909511in}}%
\pgfpathlineto{\pgfqpoint{4.464550in}{2.906562in}}%
\pgfpathmoveto{\pgfqpoint{4.460008in}{2.909511in}}%
\pgfpathlineto{\pgfqpoint{4.460008in}{2.909511in}}%
\pgfpathlineto{\pgfqpoint{4.460008in}{2.912461in}}%
\pgfpathlineto{\pgfqpoint{4.464550in}{2.912461in}}%
\pgfpathlineto{\pgfqpoint{4.464550in}{2.909511in}}%
\pgfpathmoveto{\pgfqpoint{4.464550in}{2.909511in}}%
\pgfpathlineto{\pgfqpoint{4.464550in}{2.909511in}}%
\pgfpathlineto{\pgfqpoint{4.464550in}{2.912461in}}%
\pgfpathlineto{\pgfqpoint{4.469091in}{2.912461in}}%
\pgfpathlineto{\pgfqpoint{4.469091in}{2.909511in}}%
\pgfpathmoveto{\pgfqpoint{4.464550in}{2.912461in}}%
\pgfpathlineto{\pgfqpoint{4.464550in}{2.912461in}}%
\pgfpathlineto{\pgfqpoint{4.464550in}{2.915410in}}%
\pgfpathlineto{\pgfqpoint{4.469091in}{2.915410in}}%
\pgfpathlineto{\pgfqpoint{4.469091in}{2.912461in}}%
\pgfpathmoveto{\pgfqpoint{4.469091in}{2.912461in}}%
\pgfpathlineto{\pgfqpoint{4.469091in}{2.912461in}}%
\pgfpathlineto{\pgfqpoint{4.469091in}{2.915410in}}%
\pgfpathlineto{\pgfqpoint{4.473632in}{2.915410in}}%
\pgfpathlineto{\pgfqpoint{4.473632in}{2.912461in}}%
\pgfpathmoveto{\pgfqpoint{4.469091in}{2.915410in}}%
\pgfpathlineto{\pgfqpoint{4.469091in}{2.915410in}}%
\pgfpathlineto{\pgfqpoint{4.469091in}{2.918359in}}%
\pgfpathlineto{\pgfqpoint{4.473632in}{2.918359in}}%
\pgfpathlineto{\pgfqpoint{4.473632in}{2.915410in}}%
\pgfpathmoveto{\pgfqpoint{4.473632in}{2.915410in}}%
\pgfpathlineto{\pgfqpoint{4.473632in}{2.915410in}}%
\pgfpathlineto{\pgfqpoint{4.473632in}{2.918359in}}%
\pgfpathlineto{\pgfqpoint{4.478173in}{2.918359in}}%
\pgfpathlineto{\pgfqpoint{4.478173in}{2.915410in}}%
\pgfpathmoveto{\pgfqpoint{4.473632in}{2.918359in}}%
\pgfpathlineto{\pgfqpoint{4.473632in}{2.918359in}}%
\pgfpathlineto{\pgfqpoint{4.473632in}{2.921308in}}%
\pgfpathlineto{\pgfqpoint{4.478173in}{2.921308in}}%
\pgfpathlineto{\pgfqpoint{4.478173in}{2.918359in}}%
\pgfpathmoveto{\pgfqpoint{4.473632in}{2.921308in}}%
\pgfpathlineto{\pgfqpoint{4.473632in}{2.921308in}}%
\pgfpathlineto{\pgfqpoint{4.473632in}{2.924257in}}%
\pgfpathlineto{\pgfqpoint{4.478173in}{2.924257in}}%
\pgfpathlineto{\pgfqpoint{4.478173in}{2.921308in}}%
\pgfpathmoveto{\pgfqpoint{4.478173in}{2.921308in}}%
\pgfpathlineto{\pgfqpoint{4.478173in}{2.921308in}}%
\pgfpathlineto{\pgfqpoint{4.478173in}{2.924257in}}%
\pgfpathlineto{\pgfqpoint{4.482714in}{2.924257in}}%
\pgfpathlineto{\pgfqpoint{4.482714in}{2.921308in}}%
\pgfpathmoveto{\pgfqpoint{4.478173in}{2.924257in}}%
\pgfpathlineto{\pgfqpoint{4.478173in}{2.924257in}}%
\pgfpathlineto{\pgfqpoint{4.478173in}{2.927207in}}%
\pgfpathlineto{\pgfqpoint{4.482714in}{2.927207in}}%
\pgfpathlineto{\pgfqpoint{4.482714in}{2.924257in}}%
\pgfpathmoveto{\pgfqpoint{4.482714in}{2.924257in}}%
\pgfpathlineto{\pgfqpoint{4.482714in}{2.924257in}}%
\pgfpathlineto{\pgfqpoint{4.482714in}{2.927207in}}%
\pgfpathlineto{\pgfqpoint{4.487255in}{2.927207in}}%
\pgfpathlineto{\pgfqpoint{4.487255in}{2.924257in}}%
\pgfpathmoveto{\pgfqpoint{4.482714in}{2.927207in}}%
\pgfpathlineto{\pgfqpoint{4.482714in}{2.927207in}}%
\pgfpathlineto{\pgfqpoint{4.482714in}{2.930156in}}%
\pgfpathlineto{\pgfqpoint{4.487255in}{2.930156in}}%
\pgfpathlineto{\pgfqpoint{4.487255in}{2.927207in}}%
\pgfpathmoveto{\pgfqpoint{4.487255in}{2.927207in}}%
\pgfpathlineto{\pgfqpoint{4.487255in}{2.927207in}}%
\pgfpathlineto{\pgfqpoint{4.487255in}{2.930156in}}%
\pgfpathlineto{\pgfqpoint{4.491797in}{2.930156in}}%
\pgfpathlineto{\pgfqpoint{4.491797in}{2.927207in}}%
\pgfpathmoveto{\pgfqpoint{4.487255in}{2.930156in}}%
\pgfpathlineto{\pgfqpoint{4.487255in}{2.930156in}}%
\pgfpathlineto{\pgfqpoint{4.487255in}{2.933105in}}%
\pgfpathlineto{\pgfqpoint{4.491797in}{2.933105in}}%
\pgfpathlineto{\pgfqpoint{4.491797in}{2.930156in}}%
\pgfpathmoveto{\pgfqpoint{4.491797in}{2.930156in}}%
\pgfpathlineto{\pgfqpoint{4.491797in}{2.930156in}}%
\pgfpathlineto{\pgfqpoint{4.491797in}{2.933105in}}%
\pgfpathlineto{\pgfqpoint{4.496338in}{2.933105in}}%
\pgfpathlineto{\pgfqpoint{4.496338in}{2.930156in}}%
\pgfpathmoveto{\pgfqpoint{4.491797in}{2.933105in}}%
\pgfpathlineto{\pgfqpoint{4.491797in}{2.933105in}}%
\pgfpathlineto{\pgfqpoint{4.491797in}{2.936054in}}%
\pgfpathlineto{\pgfqpoint{4.496338in}{2.936054in}}%
\pgfpathlineto{\pgfqpoint{4.496338in}{2.933105in}}%
\pgfpathmoveto{\pgfqpoint{4.496338in}{2.933105in}}%
\pgfpathlineto{\pgfqpoint{4.496338in}{2.933105in}}%
\pgfpathlineto{\pgfqpoint{4.496338in}{2.936054in}}%
\pgfpathlineto{\pgfqpoint{4.500879in}{2.936054in}}%
\pgfpathlineto{\pgfqpoint{4.500879in}{2.933105in}}%
\pgfpathmoveto{\pgfqpoint{4.496338in}{2.936054in}}%
\pgfpathlineto{\pgfqpoint{4.496338in}{2.936054in}}%
\pgfpathlineto{\pgfqpoint{4.496338in}{2.939003in}}%
\pgfpathlineto{\pgfqpoint{4.500879in}{2.939003in}}%
\pgfpathlineto{\pgfqpoint{4.500879in}{2.936054in}}%
\pgfpathmoveto{\pgfqpoint{4.500879in}{2.936054in}}%
\pgfpathlineto{\pgfqpoint{4.500879in}{2.936054in}}%
\pgfpathlineto{\pgfqpoint{4.500879in}{2.939003in}}%
\pgfpathlineto{\pgfqpoint{4.505420in}{2.939003in}}%
\pgfpathlineto{\pgfqpoint{4.505420in}{2.936054in}}%
\pgfpathmoveto{\pgfqpoint{4.500879in}{2.939003in}}%
\pgfpathlineto{\pgfqpoint{4.500879in}{2.939003in}}%
\pgfpathlineto{\pgfqpoint{4.500879in}{2.941953in}}%
\pgfpathlineto{\pgfqpoint{4.505420in}{2.941953in}}%
\pgfpathlineto{\pgfqpoint{4.505420in}{2.939003in}}%
\pgfpathmoveto{\pgfqpoint{4.505420in}{2.939003in}}%
\pgfpathlineto{\pgfqpoint{4.505420in}{2.939003in}}%
\pgfpathlineto{\pgfqpoint{4.505420in}{2.941953in}}%
\pgfpathlineto{\pgfqpoint{4.509961in}{2.941953in}}%
\pgfpathlineto{\pgfqpoint{4.509961in}{2.939003in}}%
\pgfpathmoveto{\pgfqpoint{4.505420in}{2.941953in}}%
\pgfpathlineto{\pgfqpoint{4.505420in}{2.941953in}}%
\pgfpathlineto{\pgfqpoint{4.505420in}{2.944902in}}%
\pgfpathlineto{\pgfqpoint{4.509961in}{2.944902in}}%
\pgfpathlineto{\pgfqpoint{4.509961in}{2.941953in}}%
\pgfpathmoveto{\pgfqpoint{4.509961in}{2.941953in}}%
\pgfpathlineto{\pgfqpoint{4.509961in}{2.941953in}}%
\pgfpathlineto{\pgfqpoint{4.509961in}{2.944902in}}%
\pgfpathlineto{\pgfqpoint{4.514503in}{2.944902in}}%
\pgfpathlineto{\pgfqpoint{4.514503in}{2.941953in}}%
\pgfpathmoveto{\pgfqpoint{4.509961in}{2.944902in}}%
\pgfpathlineto{\pgfqpoint{4.509961in}{2.944902in}}%
\pgfpathlineto{\pgfqpoint{4.509961in}{2.947851in}}%
\pgfpathlineto{\pgfqpoint{4.514503in}{2.947851in}}%
\pgfpathlineto{\pgfqpoint{4.514503in}{2.944902in}}%
\pgfpathmoveto{\pgfqpoint{4.514503in}{2.944902in}}%
\pgfpathlineto{\pgfqpoint{4.514503in}{2.944902in}}%
\pgfpathlineto{\pgfqpoint{4.514503in}{2.947851in}}%
\pgfpathlineto{\pgfqpoint{4.519044in}{2.947851in}}%
\pgfpathlineto{\pgfqpoint{4.519044in}{2.944902in}}%
\pgfpathmoveto{\pgfqpoint{4.514503in}{2.947851in}}%
\pgfpathlineto{\pgfqpoint{4.514503in}{2.947851in}}%
\pgfpathlineto{\pgfqpoint{4.514503in}{2.950800in}}%
\pgfpathlineto{\pgfqpoint{4.519044in}{2.950800in}}%
\pgfpathlineto{\pgfqpoint{4.519044in}{2.947851in}}%
\pgfpathmoveto{\pgfqpoint{4.519044in}{2.947851in}}%
\pgfpathlineto{\pgfqpoint{4.519044in}{2.947851in}}%
\pgfpathlineto{\pgfqpoint{4.519044in}{2.950800in}}%
\pgfpathlineto{\pgfqpoint{4.523585in}{2.950800in}}%
\pgfpathlineto{\pgfqpoint{4.523585in}{2.947851in}}%
\pgfpathmoveto{\pgfqpoint{4.519044in}{2.950800in}}%
\pgfpathlineto{\pgfqpoint{4.519044in}{2.950800in}}%
\pgfpathlineto{\pgfqpoint{4.519044in}{2.953749in}}%
\pgfpathlineto{\pgfqpoint{4.523585in}{2.953749in}}%
\pgfpathlineto{\pgfqpoint{4.523585in}{2.950800in}}%
\pgfpathmoveto{\pgfqpoint{4.523585in}{2.950800in}}%
\pgfpathlineto{\pgfqpoint{4.523585in}{2.950800in}}%
\pgfpathlineto{\pgfqpoint{4.523585in}{2.953749in}}%
\pgfpathlineto{\pgfqpoint{4.528126in}{2.953749in}}%
\pgfpathlineto{\pgfqpoint{4.528126in}{2.950800in}}%
\pgfpathmoveto{\pgfqpoint{4.523585in}{2.953749in}}%
\pgfpathlineto{\pgfqpoint{4.523585in}{2.953749in}}%
\pgfpathlineto{\pgfqpoint{4.523585in}{2.956699in}}%
\pgfpathlineto{\pgfqpoint{4.528126in}{2.956699in}}%
\pgfpathlineto{\pgfqpoint{4.528126in}{2.953749in}}%
\pgfpathmoveto{\pgfqpoint{4.528126in}{2.953749in}}%
\pgfpathlineto{\pgfqpoint{4.528126in}{2.953749in}}%
\pgfpathlineto{\pgfqpoint{4.528126in}{2.956699in}}%
\pgfpathlineto{\pgfqpoint{4.532667in}{2.956699in}}%
\pgfpathlineto{\pgfqpoint{4.532667in}{2.953749in}}%
\pgfpathmoveto{\pgfqpoint{4.528126in}{2.956699in}}%
\pgfpathlineto{\pgfqpoint{4.528126in}{2.956699in}}%
\pgfpathlineto{\pgfqpoint{4.528126in}{2.959648in}}%
\pgfpathlineto{\pgfqpoint{4.532667in}{2.959648in}}%
\pgfpathlineto{\pgfqpoint{4.532667in}{2.956699in}}%
\pgfpathmoveto{\pgfqpoint{4.532667in}{2.956699in}}%
\pgfpathlineto{\pgfqpoint{4.532667in}{2.956699in}}%
\pgfpathlineto{\pgfqpoint{4.532667in}{2.959648in}}%
\pgfpathlineto{\pgfqpoint{4.537208in}{2.959648in}}%
\pgfpathlineto{\pgfqpoint{4.537208in}{2.956699in}}%
\pgfpathmoveto{\pgfqpoint{4.532667in}{2.959648in}}%
\pgfpathlineto{\pgfqpoint{4.532667in}{2.959648in}}%
\pgfpathlineto{\pgfqpoint{4.532667in}{2.962597in}}%
\pgfpathlineto{\pgfqpoint{4.537208in}{2.962597in}}%
\pgfpathlineto{\pgfqpoint{4.537208in}{2.959648in}}%
\pgfpathmoveto{\pgfqpoint{4.537208in}{2.959648in}}%
\pgfpathlineto{\pgfqpoint{4.537208in}{2.959648in}}%
\pgfpathlineto{\pgfqpoint{4.537208in}{2.962597in}}%
\pgfpathlineto{\pgfqpoint{4.541749in}{2.962597in}}%
\pgfpathlineto{\pgfqpoint{4.541749in}{2.959648in}}%
\pgfpathmoveto{\pgfqpoint{4.537208in}{2.962597in}}%
\pgfpathlineto{\pgfqpoint{4.537208in}{2.962597in}}%
\pgfpathlineto{\pgfqpoint{4.537208in}{2.965546in}}%
\pgfpathlineto{\pgfqpoint{4.541749in}{2.965546in}}%
\pgfpathlineto{\pgfqpoint{4.541749in}{2.962597in}}%
\pgfpathmoveto{\pgfqpoint{4.541749in}{2.962597in}}%
\pgfpathlineto{\pgfqpoint{4.541749in}{2.962597in}}%
\pgfpathlineto{\pgfqpoint{4.541749in}{2.965546in}}%
\pgfpathlineto{\pgfqpoint{4.546290in}{2.965546in}}%
\pgfpathlineto{\pgfqpoint{4.546290in}{2.962597in}}%
\pgfpathmoveto{\pgfqpoint{4.541749in}{2.965546in}}%
\pgfpathlineto{\pgfqpoint{4.541749in}{2.965546in}}%
\pgfpathlineto{\pgfqpoint{4.541749in}{2.968496in}}%
\pgfpathlineto{\pgfqpoint{4.546290in}{2.968496in}}%
\pgfpathlineto{\pgfqpoint{4.546290in}{2.965546in}}%
\pgfpathmoveto{\pgfqpoint{4.546290in}{2.965546in}}%
\pgfpathlineto{\pgfqpoint{4.546290in}{2.965546in}}%
\pgfpathlineto{\pgfqpoint{4.546290in}{2.968496in}}%
\pgfpathlineto{\pgfqpoint{4.550831in}{2.968496in}}%
\pgfpathlineto{\pgfqpoint{4.550831in}{2.965546in}}%
\pgfpathmoveto{\pgfqpoint{4.546290in}{2.968496in}}%
\pgfpathlineto{\pgfqpoint{4.546290in}{2.968496in}}%
\pgfpathlineto{\pgfqpoint{4.546290in}{2.971445in}}%
\pgfpathlineto{\pgfqpoint{4.550831in}{2.971445in}}%
\pgfpathlineto{\pgfqpoint{4.550831in}{2.968496in}}%
\pgfpathmoveto{\pgfqpoint{4.550831in}{2.968496in}}%
\pgfpathlineto{\pgfqpoint{4.550831in}{2.968496in}}%
\pgfpathlineto{\pgfqpoint{4.550831in}{2.971445in}}%
\pgfpathlineto{\pgfqpoint{4.555372in}{2.971445in}}%
\pgfpathlineto{\pgfqpoint{4.555372in}{2.968496in}}%
\pgfpathmoveto{\pgfqpoint{4.550831in}{2.971445in}}%
\pgfpathlineto{\pgfqpoint{4.550831in}{2.971445in}}%
\pgfpathlineto{\pgfqpoint{4.550831in}{2.974394in}}%
\pgfpathlineto{\pgfqpoint{4.555372in}{2.974394in}}%
\pgfpathlineto{\pgfqpoint{4.555372in}{2.971445in}}%
\pgfpathmoveto{\pgfqpoint{4.555372in}{2.971445in}}%
\pgfpathlineto{\pgfqpoint{4.555372in}{2.971445in}}%
\pgfpathlineto{\pgfqpoint{4.555372in}{2.974394in}}%
\pgfpathlineto{\pgfqpoint{4.559914in}{2.974394in}}%
\pgfpathlineto{\pgfqpoint{4.559914in}{2.971445in}}%
\pgfpathmoveto{\pgfqpoint{4.555372in}{2.974394in}}%
\pgfpathlineto{\pgfqpoint{4.555372in}{2.974394in}}%
\pgfpathlineto{\pgfqpoint{4.555372in}{2.977343in}}%
\pgfpathlineto{\pgfqpoint{4.559914in}{2.977343in}}%
\pgfpathlineto{\pgfqpoint{4.559914in}{2.974394in}}%
\pgfpathmoveto{\pgfqpoint{4.559914in}{2.974394in}}%
\pgfpathlineto{\pgfqpoint{4.559914in}{2.974394in}}%
\pgfpathlineto{\pgfqpoint{4.559914in}{2.977343in}}%
\pgfpathlineto{\pgfqpoint{4.564455in}{2.977343in}}%
\pgfpathlineto{\pgfqpoint{4.564455in}{2.974394in}}%
\pgfpathmoveto{\pgfqpoint{4.559914in}{2.977343in}}%
\pgfpathlineto{\pgfqpoint{4.559914in}{2.977343in}}%
\pgfpathlineto{\pgfqpoint{4.559914in}{2.980293in}}%
\pgfpathlineto{\pgfqpoint{4.564455in}{2.980293in}}%
\pgfpathlineto{\pgfqpoint{4.564455in}{2.977343in}}%
\pgfpathmoveto{\pgfqpoint{4.564455in}{2.977343in}}%
\pgfpathlineto{\pgfqpoint{4.564455in}{2.977343in}}%
\pgfpathlineto{\pgfqpoint{4.564455in}{2.980293in}}%
\pgfpathlineto{\pgfqpoint{4.568996in}{2.980293in}}%
\pgfpathlineto{\pgfqpoint{4.568996in}{2.977343in}}%
\pgfpathmoveto{\pgfqpoint{4.564455in}{2.980293in}}%
\pgfpathlineto{\pgfqpoint{4.564455in}{2.980293in}}%
\pgfpathlineto{\pgfqpoint{4.564455in}{2.983242in}}%
\pgfpathlineto{\pgfqpoint{4.568996in}{2.983242in}}%
\pgfpathlineto{\pgfqpoint{4.568996in}{2.980293in}}%
\pgfpathmoveto{\pgfqpoint{4.568996in}{2.980293in}}%
\pgfpathlineto{\pgfqpoint{4.568996in}{2.980293in}}%
\pgfpathlineto{\pgfqpoint{4.568996in}{2.983242in}}%
\pgfpathlineto{\pgfqpoint{4.573537in}{2.983242in}}%
\pgfpathlineto{\pgfqpoint{4.573537in}{2.980293in}}%
\pgfpathmoveto{\pgfqpoint{4.568996in}{2.983242in}}%
\pgfpathlineto{\pgfqpoint{4.568996in}{2.983242in}}%
\pgfpathlineto{\pgfqpoint{4.568996in}{2.986191in}}%
\pgfpathlineto{\pgfqpoint{4.573537in}{2.986191in}}%
\pgfpathlineto{\pgfqpoint{4.573537in}{2.983242in}}%
\pgfpathmoveto{\pgfqpoint{4.573537in}{2.983242in}}%
\pgfpathlineto{\pgfqpoint{4.573537in}{2.983242in}}%
\pgfpathlineto{\pgfqpoint{4.573537in}{2.986191in}}%
\pgfpathlineto{\pgfqpoint{4.578078in}{2.986191in}}%
\pgfpathlineto{\pgfqpoint{4.578078in}{2.983242in}}%
\pgfpathmoveto{\pgfqpoint{4.573537in}{2.986191in}}%
\pgfpathlineto{\pgfqpoint{4.573537in}{2.986191in}}%
\pgfpathlineto{\pgfqpoint{4.573537in}{2.989140in}}%
\pgfpathlineto{\pgfqpoint{4.578078in}{2.989140in}}%
\pgfpathlineto{\pgfqpoint{4.578078in}{2.986191in}}%
\pgfpathmoveto{\pgfqpoint{4.578078in}{2.986191in}}%
\pgfpathlineto{\pgfqpoint{4.578078in}{2.986191in}}%
\pgfpathlineto{\pgfqpoint{4.578078in}{2.989140in}}%
\pgfpathlineto{\pgfqpoint{4.582619in}{2.989140in}}%
\pgfpathlineto{\pgfqpoint{4.582619in}{2.986191in}}%
\pgfpathmoveto{\pgfqpoint{4.578078in}{2.989140in}}%
\pgfpathlineto{\pgfqpoint{4.578078in}{2.989140in}}%
\pgfpathlineto{\pgfqpoint{4.578078in}{2.992090in}}%
\pgfpathlineto{\pgfqpoint{4.582619in}{2.992090in}}%
\pgfpathlineto{\pgfqpoint{4.582619in}{2.989140in}}%
\pgfpathmoveto{\pgfqpoint{4.582619in}{2.989140in}}%
\pgfpathlineto{\pgfqpoint{4.582619in}{2.989140in}}%
\pgfpathlineto{\pgfqpoint{4.582619in}{2.992090in}}%
\pgfpathlineto{\pgfqpoint{4.587160in}{2.992090in}}%
\pgfpathlineto{\pgfqpoint{4.587160in}{2.989140in}}%
\pgfpathmoveto{\pgfqpoint{4.582619in}{2.992090in}}%
\pgfpathlineto{\pgfqpoint{4.582619in}{2.992090in}}%
\pgfpathlineto{\pgfqpoint{4.582619in}{2.995039in}}%
\pgfpathlineto{\pgfqpoint{4.587160in}{2.995039in}}%
\pgfpathlineto{\pgfqpoint{4.587160in}{2.992090in}}%
\pgfpathmoveto{\pgfqpoint{4.587160in}{2.992090in}}%
\pgfpathlineto{\pgfqpoint{4.587160in}{2.992090in}}%
\pgfpathlineto{\pgfqpoint{4.587160in}{2.995039in}}%
\pgfpathlineto{\pgfqpoint{4.591701in}{2.995039in}}%
\pgfpathlineto{\pgfqpoint{4.591701in}{2.992090in}}%
\pgfpathmoveto{\pgfqpoint{4.587160in}{2.995039in}}%
\pgfpathlineto{\pgfqpoint{4.587160in}{2.995039in}}%
\pgfpathlineto{\pgfqpoint{4.587160in}{2.997988in}}%
\pgfpathlineto{\pgfqpoint{4.591701in}{2.997988in}}%
\pgfpathlineto{\pgfqpoint{4.591701in}{2.995039in}}%
\pgfpathmoveto{\pgfqpoint{4.591701in}{2.995039in}}%
\pgfpathlineto{\pgfqpoint{4.591701in}{2.995039in}}%
\pgfpathlineto{\pgfqpoint{4.591701in}{2.997988in}}%
\pgfpathlineto{\pgfqpoint{4.596242in}{2.997988in}}%
\pgfpathlineto{\pgfqpoint{4.596242in}{2.995039in}}%
\pgfpathmoveto{\pgfqpoint{4.591701in}{2.997988in}}%
\pgfpathlineto{\pgfqpoint{4.591701in}{2.997988in}}%
\pgfpathlineto{\pgfqpoint{4.591701in}{3.000937in}}%
\pgfpathlineto{\pgfqpoint{4.596242in}{3.000937in}}%
\pgfpathlineto{\pgfqpoint{4.596242in}{2.997988in}}%
\pgfpathmoveto{\pgfqpoint{4.596242in}{2.997988in}}%
\pgfpathlineto{\pgfqpoint{4.596242in}{2.997988in}}%
\pgfpathlineto{\pgfqpoint{4.596242in}{3.000937in}}%
\pgfpathlineto{\pgfqpoint{4.600783in}{3.000937in}}%
\pgfpathlineto{\pgfqpoint{4.600783in}{2.997988in}}%
\pgfpathmoveto{\pgfqpoint{4.596242in}{3.000937in}}%
\pgfpathlineto{\pgfqpoint{4.596242in}{3.000937in}}%
\pgfpathlineto{\pgfqpoint{4.596242in}{3.003887in}}%
\pgfpathlineto{\pgfqpoint{4.600783in}{3.003887in}}%
\pgfpathlineto{\pgfqpoint{4.600783in}{3.000937in}}%
\pgfpathmoveto{\pgfqpoint{4.600783in}{3.000937in}}%
\pgfpathlineto{\pgfqpoint{4.600783in}{3.000937in}}%
\pgfpathlineto{\pgfqpoint{4.600783in}{3.003887in}}%
\pgfpathlineto{\pgfqpoint{4.605324in}{3.003887in}}%
\pgfpathlineto{\pgfqpoint{4.605324in}{3.000937in}}%
\pgfpathmoveto{\pgfqpoint{4.600783in}{3.003887in}}%
\pgfpathlineto{\pgfqpoint{4.600783in}{3.003887in}}%
\pgfpathlineto{\pgfqpoint{4.600783in}{3.006836in}}%
\pgfpathlineto{\pgfqpoint{4.605324in}{3.006836in}}%
\pgfpathlineto{\pgfqpoint{4.605324in}{3.003887in}}%
\pgfpathmoveto{\pgfqpoint{4.605324in}{3.003887in}}%
\pgfpathlineto{\pgfqpoint{4.605324in}{3.003887in}}%
\pgfpathlineto{\pgfqpoint{4.605324in}{3.006836in}}%
\pgfpathlineto{\pgfqpoint{4.609865in}{3.006836in}}%
\pgfpathlineto{\pgfqpoint{4.609865in}{3.003887in}}%
\pgfpathmoveto{\pgfqpoint{4.605324in}{3.006836in}}%
\pgfpathlineto{\pgfqpoint{4.605324in}{3.006836in}}%
\pgfpathlineto{\pgfqpoint{4.605324in}{3.009785in}}%
\pgfpathlineto{\pgfqpoint{4.609865in}{3.009785in}}%
\pgfpathlineto{\pgfqpoint{4.609865in}{3.006836in}}%
\pgfpathmoveto{\pgfqpoint{4.609865in}{3.006836in}}%
\pgfpathlineto{\pgfqpoint{4.609865in}{3.006836in}}%
\pgfpathlineto{\pgfqpoint{4.609865in}{3.009785in}}%
\pgfpathlineto{\pgfqpoint{4.614406in}{3.009785in}}%
\pgfpathlineto{\pgfqpoint{4.614406in}{3.006836in}}%
\pgfpathmoveto{\pgfqpoint{4.609865in}{3.009785in}}%
\pgfpathlineto{\pgfqpoint{4.609865in}{3.009785in}}%
\pgfpathlineto{\pgfqpoint{4.609865in}{3.012734in}}%
\pgfpathlineto{\pgfqpoint{4.614406in}{3.012734in}}%
\pgfpathlineto{\pgfqpoint{4.614406in}{3.009785in}}%
\pgfpathmoveto{\pgfqpoint{4.614406in}{3.009785in}}%
\pgfpathlineto{\pgfqpoint{4.614406in}{3.009785in}}%
\pgfpathlineto{\pgfqpoint{4.614406in}{3.012734in}}%
\pgfpathlineto{\pgfqpoint{4.618947in}{3.012734in}}%
\pgfpathlineto{\pgfqpoint{4.618947in}{3.009785in}}%
\pgfpathmoveto{\pgfqpoint{4.614406in}{3.012734in}}%
\pgfpathlineto{\pgfqpoint{4.614406in}{3.012734in}}%
\pgfpathlineto{\pgfqpoint{4.614406in}{3.015684in}}%
\pgfpathlineto{\pgfqpoint{4.618947in}{3.015684in}}%
\pgfpathlineto{\pgfqpoint{4.618947in}{3.012734in}}%
\pgfpathmoveto{\pgfqpoint{4.618947in}{3.012734in}}%
\pgfpathlineto{\pgfqpoint{4.618947in}{3.012734in}}%
\pgfpathlineto{\pgfqpoint{4.618947in}{3.015684in}}%
\pgfpathlineto{\pgfqpoint{4.623488in}{3.015684in}}%
\pgfpathlineto{\pgfqpoint{4.623488in}{3.012734in}}%
\pgfpathmoveto{\pgfqpoint{4.618947in}{3.015684in}}%
\pgfpathlineto{\pgfqpoint{4.618947in}{3.015684in}}%
\pgfpathlineto{\pgfqpoint{4.618947in}{3.018633in}}%
\pgfpathlineto{\pgfqpoint{4.623488in}{3.018633in}}%
\pgfpathlineto{\pgfqpoint{4.623488in}{3.015684in}}%
\pgfpathmoveto{\pgfqpoint{4.623488in}{3.015684in}}%
\pgfpathlineto{\pgfqpoint{4.623488in}{3.015684in}}%
\pgfpathlineto{\pgfqpoint{4.623488in}{3.018633in}}%
\pgfpathlineto{\pgfqpoint{4.628029in}{3.018633in}}%
\pgfpathlineto{\pgfqpoint{4.628029in}{3.015684in}}%
\pgfpathmoveto{\pgfqpoint{4.623488in}{3.018633in}}%
\pgfpathlineto{\pgfqpoint{4.623488in}{3.018633in}}%
\pgfpathlineto{\pgfqpoint{4.623488in}{3.021582in}}%
\pgfpathlineto{\pgfqpoint{4.628029in}{3.021582in}}%
\pgfpathlineto{\pgfqpoint{4.628029in}{3.018633in}}%
\pgfpathmoveto{\pgfqpoint{4.628029in}{3.018633in}}%
\pgfpathlineto{\pgfqpoint{4.628029in}{3.018633in}}%
\pgfpathlineto{\pgfqpoint{4.628029in}{3.021582in}}%
\pgfpathlineto{\pgfqpoint{4.632571in}{3.021582in}}%
\pgfpathlineto{\pgfqpoint{4.632571in}{3.018633in}}%
\pgfpathmoveto{\pgfqpoint{4.628029in}{3.021582in}}%
\pgfpathlineto{\pgfqpoint{4.628029in}{3.021582in}}%
\pgfpathlineto{\pgfqpoint{4.628029in}{3.024532in}}%
\pgfpathlineto{\pgfqpoint{4.632571in}{3.024532in}}%
\pgfpathlineto{\pgfqpoint{4.632571in}{3.021582in}}%
\pgfpathmoveto{\pgfqpoint{4.632571in}{3.021582in}}%
\pgfpathlineto{\pgfqpoint{4.632571in}{3.021582in}}%
\pgfpathlineto{\pgfqpoint{4.632571in}{3.024532in}}%
\pgfpathlineto{\pgfqpoint{4.637112in}{3.024532in}}%
\pgfpathlineto{\pgfqpoint{4.637112in}{3.021582in}}%
\pgfpathmoveto{\pgfqpoint{4.632571in}{3.024532in}}%
\pgfpathlineto{\pgfqpoint{4.632571in}{3.024532in}}%
\pgfpathlineto{\pgfqpoint{4.632571in}{3.027481in}}%
\pgfpathlineto{\pgfqpoint{4.637112in}{3.027481in}}%
\pgfpathlineto{\pgfqpoint{4.637112in}{3.024532in}}%
\pgfpathmoveto{\pgfqpoint{4.637112in}{3.024532in}}%
\pgfpathlineto{\pgfqpoint{4.637112in}{3.024532in}}%
\pgfpathlineto{\pgfqpoint{4.637112in}{3.027481in}}%
\pgfpathlineto{\pgfqpoint{4.641653in}{3.027481in}}%
\pgfpathlineto{\pgfqpoint{4.641653in}{3.024532in}}%
\pgfpathmoveto{\pgfqpoint{4.637112in}{3.027481in}}%
\pgfpathlineto{\pgfqpoint{4.637112in}{3.027481in}}%
\pgfpathlineto{\pgfqpoint{4.637112in}{3.030430in}}%
\pgfpathlineto{\pgfqpoint{4.641653in}{3.030430in}}%
\pgfpathlineto{\pgfqpoint{4.641653in}{3.027481in}}%
\pgfpathmoveto{\pgfqpoint{4.641653in}{3.027481in}}%
\pgfpathlineto{\pgfqpoint{4.641653in}{3.027481in}}%
\pgfpathlineto{\pgfqpoint{4.641653in}{3.030430in}}%
\pgfpathlineto{\pgfqpoint{4.646194in}{3.030430in}}%
\pgfpathlineto{\pgfqpoint{4.646194in}{3.027481in}}%
\pgfpathmoveto{\pgfqpoint{4.641653in}{3.030430in}}%
\pgfpathlineto{\pgfqpoint{4.641653in}{3.030430in}}%
\pgfpathlineto{\pgfqpoint{4.641653in}{3.033379in}}%
\pgfpathlineto{\pgfqpoint{4.646194in}{3.033379in}}%
\pgfpathlineto{\pgfqpoint{4.646194in}{3.030430in}}%
\pgfpathmoveto{\pgfqpoint{4.646194in}{3.030430in}}%
\pgfpathlineto{\pgfqpoint{4.646194in}{3.030430in}}%
\pgfpathlineto{\pgfqpoint{4.646194in}{3.033379in}}%
\pgfpathlineto{\pgfqpoint{4.650735in}{3.033379in}}%
\pgfpathlineto{\pgfqpoint{4.650735in}{3.030430in}}%
\pgfpathmoveto{\pgfqpoint{4.646194in}{3.033379in}}%
\pgfpathlineto{\pgfqpoint{4.646194in}{3.033379in}}%
\pgfpathlineto{\pgfqpoint{4.646194in}{3.036329in}}%
\pgfpathlineto{\pgfqpoint{4.650735in}{3.036329in}}%
\pgfpathlineto{\pgfqpoint{4.650735in}{3.033379in}}%
\pgfpathmoveto{\pgfqpoint{4.650735in}{3.033379in}}%
\pgfpathlineto{\pgfqpoint{4.650735in}{3.033379in}}%
\pgfpathlineto{\pgfqpoint{4.650735in}{3.036329in}}%
\pgfpathlineto{\pgfqpoint{4.655276in}{3.036329in}}%
\pgfpathlineto{\pgfqpoint{4.655276in}{3.033379in}}%
\pgfpathmoveto{\pgfqpoint{4.650735in}{3.036329in}}%
\pgfpathlineto{\pgfqpoint{4.650735in}{3.036329in}}%
\pgfpathlineto{\pgfqpoint{4.650735in}{3.039278in}}%
\pgfpathlineto{\pgfqpoint{4.655276in}{3.039278in}}%
\pgfpathlineto{\pgfqpoint{4.655276in}{3.036329in}}%
\pgfpathmoveto{\pgfqpoint{4.655276in}{3.036329in}}%
\pgfpathlineto{\pgfqpoint{4.655276in}{3.036329in}}%
\pgfpathlineto{\pgfqpoint{4.655276in}{3.039278in}}%
\pgfpathlineto{\pgfqpoint{4.659817in}{3.039278in}}%
\pgfpathlineto{\pgfqpoint{4.659817in}{3.036329in}}%
\pgfpathmoveto{\pgfqpoint{4.655276in}{3.039278in}}%
\pgfpathlineto{\pgfqpoint{4.655276in}{3.039278in}}%
\pgfpathlineto{\pgfqpoint{4.655276in}{3.042227in}}%
\pgfpathlineto{\pgfqpoint{4.659817in}{3.042227in}}%
\pgfpathlineto{\pgfqpoint{4.659817in}{3.039278in}}%
\pgfpathmoveto{\pgfqpoint{4.659817in}{3.039278in}}%
\pgfpathlineto{\pgfqpoint{4.659817in}{3.039278in}}%
\pgfpathlineto{\pgfqpoint{4.659817in}{3.042227in}}%
\pgfpathlineto{\pgfqpoint{4.664358in}{3.042227in}}%
\pgfpathlineto{\pgfqpoint{4.664358in}{3.039278in}}%
\pgfpathmoveto{\pgfqpoint{4.659817in}{3.042227in}}%
\pgfpathlineto{\pgfqpoint{4.659817in}{3.042227in}}%
\pgfpathlineto{\pgfqpoint{4.659817in}{3.045176in}}%
\pgfpathlineto{\pgfqpoint{4.664358in}{3.045176in}}%
\pgfpathlineto{\pgfqpoint{4.664358in}{3.042227in}}%
\pgfpathmoveto{\pgfqpoint{4.664358in}{3.042227in}}%
\pgfpathlineto{\pgfqpoint{4.664358in}{3.042227in}}%
\pgfpathlineto{\pgfqpoint{4.664358in}{3.045176in}}%
\pgfpathlineto{\pgfqpoint{4.668899in}{3.045176in}}%
\pgfpathlineto{\pgfqpoint{4.668899in}{3.042227in}}%
\pgfpathmoveto{\pgfqpoint{4.664358in}{3.045176in}}%
\pgfpathlineto{\pgfqpoint{4.664358in}{3.045176in}}%
\pgfpathlineto{\pgfqpoint{4.664358in}{3.048126in}}%
\pgfpathlineto{\pgfqpoint{4.668899in}{3.048126in}}%
\pgfpathlineto{\pgfqpoint{4.668899in}{3.045176in}}%
\pgfpathmoveto{\pgfqpoint{4.668899in}{3.045176in}}%
\pgfpathlineto{\pgfqpoint{4.668899in}{3.045176in}}%
\pgfpathlineto{\pgfqpoint{4.668899in}{3.048126in}}%
\pgfpathlineto{\pgfqpoint{4.673440in}{3.048126in}}%
\pgfpathlineto{\pgfqpoint{4.673440in}{3.045176in}}%
\pgfpathmoveto{\pgfqpoint{4.668899in}{3.048126in}}%
\pgfpathlineto{\pgfqpoint{4.668899in}{3.048126in}}%
\pgfpathlineto{\pgfqpoint{4.668899in}{3.051075in}}%
\pgfpathlineto{\pgfqpoint{4.673440in}{3.051075in}}%
\pgfpathlineto{\pgfqpoint{4.673440in}{3.048126in}}%
\pgfpathmoveto{\pgfqpoint{4.673440in}{3.048126in}}%
\pgfpathlineto{\pgfqpoint{4.673440in}{3.048126in}}%
\pgfpathlineto{\pgfqpoint{4.673440in}{3.051075in}}%
\pgfpathlineto{\pgfqpoint{4.677981in}{3.051075in}}%
\pgfpathlineto{\pgfqpoint{4.677981in}{3.048126in}}%
\pgfpathmoveto{\pgfqpoint{4.673440in}{3.051075in}}%
\pgfpathlineto{\pgfqpoint{4.673440in}{3.051075in}}%
\pgfpathlineto{\pgfqpoint{4.673440in}{3.054024in}}%
\pgfpathlineto{\pgfqpoint{4.677981in}{3.054024in}}%
\pgfpathlineto{\pgfqpoint{4.677981in}{3.051075in}}%
\pgfpathmoveto{\pgfqpoint{4.677981in}{3.051075in}}%
\pgfpathlineto{\pgfqpoint{4.677981in}{3.051075in}}%
\pgfpathlineto{\pgfqpoint{4.677981in}{3.054024in}}%
\pgfpathlineto{\pgfqpoint{4.682522in}{3.054024in}}%
\pgfpathlineto{\pgfqpoint{4.682522in}{3.051075in}}%
\pgfpathmoveto{\pgfqpoint{4.677981in}{3.054024in}}%
\pgfpathlineto{\pgfqpoint{4.677981in}{3.054024in}}%
\pgfpathlineto{\pgfqpoint{4.677981in}{3.056973in}}%
\pgfpathlineto{\pgfqpoint{4.682522in}{3.056973in}}%
\pgfpathlineto{\pgfqpoint{4.682522in}{3.054024in}}%
\pgfpathmoveto{\pgfqpoint{4.682522in}{3.054024in}}%
\pgfpathlineto{\pgfqpoint{4.682522in}{3.054024in}}%
\pgfpathlineto{\pgfqpoint{4.682522in}{3.056973in}}%
\pgfpathlineto{\pgfqpoint{4.687063in}{3.056973in}}%
\pgfpathlineto{\pgfqpoint{4.687063in}{3.054024in}}%
\pgfpathmoveto{\pgfqpoint{4.682522in}{3.056973in}}%
\pgfpathlineto{\pgfqpoint{4.682522in}{3.056973in}}%
\pgfpathlineto{\pgfqpoint{4.682522in}{3.059922in}}%
\pgfpathlineto{\pgfqpoint{4.687063in}{3.059922in}}%
\pgfpathlineto{\pgfqpoint{4.687063in}{3.056973in}}%
\pgfpathmoveto{\pgfqpoint{4.687063in}{3.056973in}}%
\pgfpathlineto{\pgfqpoint{4.687063in}{3.056973in}}%
\pgfpathlineto{\pgfqpoint{4.687063in}{3.059922in}}%
\pgfpathlineto{\pgfqpoint{4.691604in}{3.059922in}}%
\pgfpathlineto{\pgfqpoint{4.691604in}{3.056973in}}%
\pgfpathmoveto{\pgfqpoint{4.687063in}{3.059922in}}%
\pgfpathlineto{\pgfqpoint{4.687063in}{3.059922in}}%
\pgfpathlineto{\pgfqpoint{4.687063in}{3.062872in}}%
\pgfpathlineto{\pgfqpoint{4.691604in}{3.062872in}}%
\pgfpathlineto{\pgfqpoint{4.691604in}{3.059922in}}%
\pgfpathmoveto{\pgfqpoint{4.691604in}{3.059922in}}%
\pgfpathlineto{\pgfqpoint{4.691604in}{3.059922in}}%
\pgfpathlineto{\pgfqpoint{4.691604in}{3.062872in}}%
\pgfpathlineto{\pgfqpoint{4.696145in}{3.062872in}}%
\pgfpathlineto{\pgfqpoint{4.696145in}{3.059922in}}%
\pgfpathmoveto{\pgfqpoint{4.691604in}{3.062872in}}%
\pgfpathlineto{\pgfqpoint{4.691604in}{3.062872in}}%
\pgfpathlineto{\pgfqpoint{4.691604in}{3.065821in}}%
\pgfpathlineto{\pgfqpoint{4.696145in}{3.065821in}}%
\pgfpathlineto{\pgfqpoint{4.696145in}{3.062872in}}%
\pgfpathmoveto{\pgfqpoint{4.696145in}{3.062872in}}%
\pgfpathlineto{\pgfqpoint{4.696145in}{3.062872in}}%
\pgfpathlineto{\pgfqpoint{4.696145in}{3.065821in}}%
\pgfpathlineto{\pgfqpoint{4.700686in}{3.065821in}}%
\pgfpathlineto{\pgfqpoint{4.700686in}{3.062872in}}%
\pgfpathmoveto{\pgfqpoint{4.696145in}{3.065821in}}%
\pgfpathlineto{\pgfqpoint{4.696145in}{3.065821in}}%
\pgfpathlineto{\pgfqpoint{4.696145in}{3.068770in}}%
\pgfpathlineto{\pgfqpoint{4.700686in}{3.068770in}}%
\pgfpathlineto{\pgfqpoint{4.700686in}{3.065821in}}%
\pgfpathmoveto{\pgfqpoint{4.700686in}{3.065821in}}%
\pgfpathlineto{\pgfqpoint{4.700686in}{3.065821in}}%
\pgfpathlineto{\pgfqpoint{4.700686in}{3.068770in}}%
\pgfpathlineto{\pgfqpoint{4.705227in}{3.068770in}}%
\pgfpathlineto{\pgfqpoint{4.705227in}{3.065821in}}%
\pgfpathmoveto{\pgfqpoint{4.700686in}{3.068770in}}%
\pgfpathlineto{\pgfqpoint{4.700686in}{3.068770in}}%
\pgfpathlineto{\pgfqpoint{4.700686in}{3.071719in}}%
\pgfpathlineto{\pgfqpoint{4.705227in}{3.071719in}}%
\pgfpathlineto{\pgfqpoint{4.705227in}{3.068770in}}%
\pgfpathmoveto{\pgfqpoint{4.705227in}{3.068770in}}%
\pgfpathlineto{\pgfqpoint{4.705227in}{3.068770in}}%
\pgfpathlineto{\pgfqpoint{4.705227in}{3.071719in}}%
\pgfpathlineto{\pgfqpoint{4.709767in}{3.071719in}}%
\pgfpathlineto{\pgfqpoint{4.709767in}{3.068770in}}%
\pgfpathmoveto{\pgfqpoint{4.705227in}{3.071719in}}%
\pgfpathlineto{\pgfqpoint{4.705227in}{3.071719in}}%
\pgfpathlineto{\pgfqpoint{4.705227in}{3.074668in}}%
\pgfpathlineto{\pgfqpoint{4.709767in}{3.074668in}}%
\pgfpathlineto{\pgfqpoint{4.709767in}{3.071719in}}%
\pgfpathmoveto{\pgfqpoint{4.709767in}{3.071719in}}%
\pgfpathlineto{\pgfqpoint{4.709767in}{3.071719in}}%
\pgfpathlineto{\pgfqpoint{4.709767in}{3.074668in}}%
\pgfpathlineto{\pgfqpoint{4.714308in}{3.074668in}}%
\pgfpathlineto{\pgfqpoint{4.714308in}{3.071719in}}%
\pgfpathmoveto{\pgfqpoint{4.709767in}{3.074668in}}%
\pgfpathlineto{\pgfqpoint{4.709767in}{3.074668in}}%
\pgfpathlineto{\pgfqpoint{4.709767in}{3.077618in}}%
\pgfpathlineto{\pgfqpoint{4.714308in}{3.077618in}}%
\pgfpathlineto{\pgfqpoint{4.714308in}{3.074668in}}%
\pgfpathmoveto{\pgfqpoint{4.714308in}{3.074668in}}%
\pgfpathlineto{\pgfqpoint{4.714308in}{3.074668in}}%
\pgfpathlineto{\pgfqpoint{4.714308in}{3.077618in}}%
\pgfpathlineto{\pgfqpoint{4.718849in}{3.077618in}}%
\pgfpathlineto{\pgfqpoint{4.718849in}{3.074668in}}%
\pgfpathmoveto{\pgfqpoint{4.714308in}{3.077618in}}%
\pgfpathlineto{\pgfqpoint{4.714308in}{3.077618in}}%
\pgfpathlineto{\pgfqpoint{4.714308in}{3.080567in}}%
\pgfpathlineto{\pgfqpoint{4.718849in}{3.080567in}}%
\pgfpathlineto{\pgfqpoint{4.718849in}{3.077618in}}%
\pgfpathmoveto{\pgfqpoint{4.718849in}{3.077618in}}%
\pgfpathlineto{\pgfqpoint{4.718849in}{3.077618in}}%
\pgfpathlineto{\pgfqpoint{4.718849in}{3.080567in}}%
\pgfpathlineto{\pgfqpoint{4.723390in}{3.080567in}}%
\pgfpathlineto{\pgfqpoint{4.723390in}{3.077618in}}%
\pgfpathmoveto{\pgfqpoint{4.718849in}{3.080567in}}%
\pgfpathlineto{\pgfqpoint{4.718849in}{3.080567in}}%
\pgfpathlineto{\pgfqpoint{4.718849in}{3.083516in}}%
\pgfpathlineto{\pgfqpoint{4.723390in}{3.083516in}}%
\pgfpathlineto{\pgfqpoint{4.723390in}{3.080567in}}%
\pgfpathmoveto{\pgfqpoint{4.723390in}{3.080567in}}%
\pgfpathlineto{\pgfqpoint{4.723390in}{3.080567in}}%
\pgfpathlineto{\pgfqpoint{4.723390in}{3.083516in}}%
\pgfpathlineto{\pgfqpoint{4.727931in}{3.083516in}}%
\pgfpathlineto{\pgfqpoint{4.727931in}{3.080567in}}%
\pgfpathmoveto{\pgfqpoint{4.723390in}{3.083516in}}%
\pgfpathlineto{\pgfqpoint{4.723390in}{3.083516in}}%
\pgfpathlineto{\pgfqpoint{4.723390in}{3.086465in}}%
\pgfpathlineto{\pgfqpoint{4.727931in}{3.086465in}}%
\pgfpathlineto{\pgfqpoint{4.727931in}{3.083516in}}%
\pgfpathmoveto{\pgfqpoint{4.727931in}{3.083516in}}%
\pgfpathlineto{\pgfqpoint{4.727931in}{3.083516in}}%
\pgfpathlineto{\pgfqpoint{4.727931in}{3.086465in}}%
\pgfpathlineto{\pgfqpoint{4.732472in}{3.086465in}}%
\pgfpathlineto{\pgfqpoint{4.732472in}{3.083516in}}%
\pgfpathmoveto{\pgfqpoint{4.727931in}{3.086465in}}%
\pgfpathlineto{\pgfqpoint{4.727931in}{3.086465in}}%
\pgfpathlineto{\pgfqpoint{4.727931in}{3.089414in}}%
\pgfpathlineto{\pgfqpoint{4.732472in}{3.089414in}}%
\pgfpathlineto{\pgfqpoint{4.732472in}{3.086465in}}%
\pgfpathmoveto{\pgfqpoint{4.732472in}{3.086465in}}%
\pgfpathlineto{\pgfqpoint{4.732472in}{3.086465in}}%
\pgfpathlineto{\pgfqpoint{4.732472in}{3.089414in}}%
\pgfpathlineto{\pgfqpoint{4.737013in}{3.089414in}}%
\pgfpathlineto{\pgfqpoint{4.737013in}{3.086465in}}%
\pgfpathmoveto{\pgfqpoint{4.732472in}{3.089414in}}%
\pgfpathlineto{\pgfqpoint{4.732472in}{3.089414in}}%
\pgfpathlineto{\pgfqpoint{4.732472in}{3.092364in}}%
\pgfpathlineto{\pgfqpoint{4.737013in}{3.092364in}}%
\pgfpathlineto{\pgfqpoint{4.737013in}{3.089414in}}%
\pgfpathmoveto{\pgfqpoint{4.737013in}{3.089414in}}%
\pgfpathlineto{\pgfqpoint{4.737013in}{3.089414in}}%
\pgfpathlineto{\pgfqpoint{4.737013in}{3.092364in}}%
\pgfpathlineto{\pgfqpoint{4.741554in}{3.092364in}}%
\pgfpathlineto{\pgfqpoint{4.741554in}{3.089414in}}%
\pgfpathmoveto{\pgfqpoint{4.737013in}{3.092364in}}%
\pgfpathlineto{\pgfqpoint{4.737013in}{3.092364in}}%
\pgfpathlineto{\pgfqpoint{4.737013in}{3.095313in}}%
\pgfpathlineto{\pgfqpoint{4.741554in}{3.095313in}}%
\pgfpathlineto{\pgfqpoint{4.741554in}{3.092364in}}%
\pgfpathmoveto{\pgfqpoint{4.741554in}{3.092364in}}%
\pgfpathlineto{\pgfqpoint{4.741554in}{3.092364in}}%
\pgfpathlineto{\pgfqpoint{4.741554in}{3.095313in}}%
\pgfpathlineto{\pgfqpoint{4.746095in}{3.095313in}}%
\pgfpathlineto{\pgfqpoint{4.746095in}{3.092364in}}%
\pgfpathmoveto{\pgfqpoint{4.741554in}{3.095313in}}%
\pgfpathlineto{\pgfqpoint{4.741554in}{3.095313in}}%
\pgfpathlineto{\pgfqpoint{4.741554in}{3.098262in}}%
\pgfpathlineto{\pgfqpoint{4.746095in}{3.098262in}}%
\pgfpathlineto{\pgfqpoint{4.746095in}{3.095313in}}%
\pgfpathmoveto{\pgfqpoint{4.746095in}{3.095313in}}%
\pgfpathlineto{\pgfqpoint{4.746095in}{3.095313in}}%
\pgfpathlineto{\pgfqpoint{4.746095in}{3.098262in}}%
\pgfpathlineto{\pgfqpoint{4.750636in}{3.098262in}}%
\pgfpathlineto{\pgfqpoint{4.750636in}{3.095313in}}%
\pgfpathmoveto{\pgfqpoint{4.746095in}{3.098262in}}%
\pgfpathlineto{\pgfqpoint{4.746095in}{3.098262in}}%
\pgfpathlineto{\pgfqpoint{4.746095in}{3.101211in}}%
\pgfpathlineto{\pgfqpoint{4.750636in}{3.101211in}}%
\pgfpathlineto{\pgfqpoint{4.750636in}{3.098262in}}%
\pgfpathmoveto{\pgfqpoint{4.750636in}{3.098262in}}%
\pgfpathlineto{\pgfqpoint{4.750636in}{3.098262in}}%
\pgfpathlineto{\pgfqpoint{4.750636in}{3.101211in}}%
\pgfpathlineto{\pgfqpoint{4.755177in}{3.101211in}}%
\pgfpathlineto{\pgfqpoint{4.755177in}{3.098262in}}%
\pgfpathmoveto{\pgfqpoint{4.750636in}{3.101211in}}%
\pgfpathlineto{\pgfqpoint{4.750636in}{3.101211in}}%
\pgfpathlineto{\pgfqpoint{4.750636in}{3.104160in}}%
\pgfpathlineto{\pgfqpoint{4.755177in}{3.104160in}}%
\pgfpathlineto{\pgfqpoint{4.755177in}{3.101211in}}%
\pgfpathmoveto{\pgfqpoint{4.755177in}{3.101211in}}%
\pgfpathlineto{\pgfqpoint{4.755177in}{3.101211in}}%
\pgfpathlineto{\pgfqpoint{4.755177in}{3.104160in}}%
\pgfpathlineto{\pgfqpoint{4.759718in}{3.104160in}}%
\pgfpathlineto{\pgfqpoint{4.759718in}{3.101211in}}%
\pgfpathmoveto{\pgfqpoint{4.755177in}{3.104160in}}%
\pgfpathlineto{\pgfqpoint{4.755177in}{3.104160in}}%
\pgfpathlineto{\pgfqpoint{4.755177in}{3.107110in}}%
\pgfpathlineto{\pgfqpoint{4.759718in}{3.107110in}}%
\pgfpathlineto{\pgfqpoint{4.759718in}{3.104160in}}%
\pgfpathmoveto{\pgfqpoint{4.759718in}{3.104160in}}%
\pgfpathlineto{\pgfqpoint{4.759718in}{3.104160in}}%
\pgfpathlineto{\pgfqpoint{4.759718in}{3.107110in}}%
\pgfpathlineto{\pgfqpoint{4.764258in}{3.107110in}}%
\pgfpathlineto{\pgfqpoint{4.764258in}{3.104160in}}%
\pgfpathmoveto{\pgfqpoint{4.759718in}{3.107110in}}%
\pgfpathlineto{\pgfqpoint{4.759718in}{3.107110in}}%
\pgfpathlineto{\pgfqpoint{4.759718in}{3.110059in}}%
\pgfpathlineto{\pgfqpoint{4.764258in}{3.110059in}}%
\pgfpathlineto{\pgfqpoint{4.764258in}{3.107110in}}%
\pgfpathmoveto{\pgfqpoint{4.764258in}{3.107110in}}%
\pgfpathlineto{\pgfqpoint{4.764258in}{3.107110in}}%
\pgfpathlineto{\pgfqpoint{4.764258in}{3.110059in}}%
\pgfpathlineto{\pgfqpoint{4.768799in}{3.110059in}}%
\pgfpathlineto{\pgfqpoint{4.768799in}{3.107110in}}%
\pgfpathmoveto{\pgfqpoint{4.764258in}{3.110059in}}%
\pgfpathlineto{\pgfqpoint{4.764258in}{3.110059in}}%
\pgfpathlineto{\pgfqpoint{4.764258in}{3.113008in}}%
\pgfpathlineto{\pgfqpoint{4.768799in}{3.113008in}}%
\pgfpathlineto{\pgfqpoint{4.768799in}{3.110059in}}%
\pgfpathmoveto{\pgfqpoint{4.768799in}{3.110059in}}%
\pgfpathlineto{\pgfqpoint{4.768799in}{3.110059in}}%
\pgfpathlineto{\pgfqpoint{4.768799in}{3.113008in}}%
\pgfpathlineto{\pgfqpoint{4.773340in}{3.113008in}}%
\pgfpathlineto{\pgfqpoint{4.773340in}{3.110059in}}%
\pgfpathmoveto{\pgfqpoint{4.773340in}{3.110059in}}%
\pgfpathlineto{\pgfqpoint{4.773340in}{3.110059in}}%
\pgfpathlineto{\pgfqpoint{4.773340in}{3.113008in}}%
\pgfpathlineto{\pgfqpoint{4.777881in}{3.113008in}}%
\pgfpathlineto{\pgfqpoint{4.777881in}{3.110059in}}%
\pgfpathmoveto{\pgfqpoint{4.773340in}{3.113008in}}%
\pgfpathlineto{\pgfqpoint{4.773340in}{3.113008in}}%
\pgfpathlineto{\pgfqpoint{4.773340in}{3.115957in}}%
\pgfpathlineto{\pgfqpoint{4.777881in}{3.115957in}}%
\pgfpathlineto{\pgfqpoint{4.777881in}{3.113008in}}%
\pgfpathmoveto{\pgfqpoint{4.777881in}{3.113008in}}%
\pgfpathlineto{\pgfqpoint{4.777881in}{3.113008in}}%
\pgfpathlineto{\pgfqpoint{4.777881in}{3.115957in}}%
\pgfpathlineto{\pgfqpoint{4.782422in}{3.115957in}}%
\pgfpathlineto{\pgfqpoint{4.782422in}{3.113008in}}%
\pgfpathmoveto{\pgfqpoint{4.777881in}{3.115957in}}%
\pgfpathlineto{\pgfqpoint{4.777881in}{3.115957in}}%
\pgfpathlineto{\pgfqpoint{4.777881in}{3.118907in}}%
\pgfpathlineto{\pgfqpoint{4.782422in}{3.118907in}}%
\pgfpathlineto{\pgfqpoint{4.782422in}{3.115957in}}%
\pgfpathmoveto{\pgfqpoint{4.782422in}{3.115957in}}%
\pgfpathlineto{\pgfqpoint{4.782422in}{3.115957in}}%
\pgfpathlineto{\pgfqpoint{4.782422in}{3.118907in}}%
\pgfpathlineto{\pgfqpoint{4.786963in}{3.118907in}}%
\pgfpathlineto{\pgfqpoint{4.786963in}{3.115957in}}%
\pgfpathmoveto{\pgfqpoint{4.782422in}{3.118907in}}%
\pgfpathlineto{\pgfqpoint{4.782422in}{3.118907in}}%
\pgfpathlineto{\pgfqpoint{4.782422in}{3.121856in}}%
\pgfpathlineto{\pgfqpoint{4.786963in}{3.121856in}}%
\pgfpathlineto{\pgfqpoint{4.786963in}{3.118907in}}%
\pgfpathmoveto{\pgfqpoint{4.786963in}{3.118907in}}%
\pgfpathlineto{\pgfqpoint{4.786963in}{3.118907in}}%
\pgfpathlineto{\pgfqpoint{4.786963in}{3.121856in}}%
\pgfpathlineto{\pgfqpoint{4.791504in}{3.121856in}}%
\pgfpathlineto{\pgfqpoint{4.791504in}{3.118907in}}%
\pgfpathmoveto{\pgfqpoint{4.786963in}{3.121856in}}%
\pgfpathlineto{\pgfqpoint{4.786963in}{3.121856in}}%
\pgfpathlineto{\pgfqpoint{4.786963in}{3.124805in}}%
\pgfpathlineto{\pgfqpoint{4.791504in}{3.124805in}}%
\pgfpathlineto{\pgfqpoint{4.791504in}{3.121856in}}%
\pgfpathmoveto{\pgfqpoint{4.791504in}{3.121856in}}%
\pgfpathlineto{\pgfqpoint{4.791504in}{3.121856in}}%
\pgfpathlineto{\pgfqpoint{4.791504in}{3.124805in}}%
\pgfpathlineto{\pgfqpoint{4.796045in}{3.124805in}}%
\pgfpathlineto{\pgfqpoint{4.796045in}{3.121856in}}%
\pgfpathmoveto{\pgfqpoint{4.791504in}{3.124805in}}%
\pgfpathlineto{\pgfqpoint{4.791504in}{3.124805in}}%
\pgfpathlineto{\pgfqpoint{4.791504in}{3.127754in}}%
\pgfpathlineto{\pgfqpoint{4.796045in}{3.127754in}}%
\pgfpathlineto{\pgfqpoint{4.796045in}{3.124805in}}%
\pgfpathmoveto{\pgfqpoint{4.796045in}{3.124805in}}%
\pgfpathlineto{\pgfqpoint{4.796045in}{3.124805in}}%
\pgfpathlineto{\pgfqpoint{4.796045in}{3.127754in}}%
\pgfpathlineto{\pgfqpoint{4.800586in}{3.127754in}}%
\pgfpathlineto{\pgfqpoint{4.800586in}{3.124805in}}%
\pgfpathmoveto{\pgfqpoint{4.796045in}{3.127754in}}%
\pgfpathlineto{\pgfqpoint{4.796045in}{3.127754in}}%
\pgfpathlineto{\pgfqpoint{4.796045in}{3.130703in}}%
\pgfpathlineto{\pgfqpoint{4.800586in}{3.130703in}}%
\pgfpathlineto{\pgfqpoint{4.800586in}{3.127754in}}%
\pgfpathmoveto{\pgfqpoint{4.800586in}{3.127754in}}%
\pgfpathlineto{\pgfqpoint{4.800586in}{3.127754in}}%
\pgfpathlineto{\pgfqpoint{4.800586in}{3.130703in}}%
\pgfpathlineto{\pgfqpoint{4.805127in}{3.130703in}}%
\pgfpathlineto{\pgfqpoint{4.805127in}{3.127754in}}%
\pgfpathmoveto{\pgfqpoint{4.800586in}{3.130703in}}%
\pgfpathlineto{\pgfqpoint{4.800586in}{3.130703in}}%
\pgfpathlineto{\pgfqpoint{4.800586in}{3.133653in}}%
\pgfpathlineto{\pgfqpoint{4.805127in}{3.133653in}}%
\pgfpathlineto{\pgfqpoint{4.805127in}{3.130703in}}%
\pgfpathmoveto{\pgfqpoint{4.805127in}{3.130703in}}%
\pgfpathlineto{\pgfqpoint{4.805127in}{3.130703in}}%
\pgfpathlineto{\pgfqpoint{4.805127in}{3.133653in}}%
\pgfpathlineto{\pgfqpoint{4.809668in}{3.133653in}}%
\pgfpathlineto{\pgfqpoint{4.809668in}{3.130703in}}%
\pgfpathmoveto{\pgfqpoint{4.805127in}{3.133653in}}%
\pgfpathlineto{\pgfqpoint{4.805127in}{3.133653in}}%
\pgfpathlineto{\pgfqpoint{4.805127in}{3.136602in}}%
\pgfpathlineto{\pgfqpoint{4.809668in}{3.136602in}}%
\pgfpathlineto{\pgfqpoint{4.809668in}{3.133653in}}%
\pgfpathmoveto{\pgfqpoint{4.809668in}{3.133653in}}%
\pgfpathlineto{\pgfqpoint{4.809668in}{3.133653in}}%
\pgfpathlineto{\pgfqpoint{4.809668in}{3.136602in}}%
\pgfpathlineto{\pgfqpoint{4.814209in}{3.136602in}}%
\pgfpathlineto{\pgfqpoint{4.814209in}{3.133653in}}%
\pgfpathmoveto{\pgfqpoint{4.809668in}{3.136602in}}%
\pgfpathlineto{\pgfqpoint{4.809668in}{3.136602in}}%
\pgfpathlineto{\pgfqpoint{4.809668in}{3.139551in}}%
\pgfpathlineto{\pgfqpoint{4.814209in}{3.139551in}}%
\pgfpathlineto{\pgfqpoint{4.814209in}{3.136602in}}%
\pgfpathmoveto{\pgfqpoint{4.814209in}{3.136602in}}%
\pgfpathlineto{\pgfqpoint{4.814209in}{3.136602in}}%
\pgfpathlineto{\pgfqpoint{4.814209in}{3.139551in}}%
\pgfpathlineto{\pgfqpoint{4.818749in}{3.139551in}}%
\pgfpathlineto{\pgfqpoint{4.818749in}{3.136602in}}%
\pgfpathmoveto{\pgfqpoint{4.814209in}{3.139551in}}%
\pgfpathlineto{\pgfqpoint{4.814209in}{3.139551in}}%
\pgfpathlineto{\pgfqpoint{4.814209in}{3.142500in}}%
\pgfpathlineto{\pgfqpoint{4.818749in}{3.142500in}}%
\pgfpathlineto{\pgfqpoint{4.818749in}{3.139551in}}%
\pgfpathmoveto{\pgfqpoint{4.818749in}{3.139551in}}%
\pgfpathlineto{\pgfqpoint{4.818749in}{3.139551in}}%
\pgfpathlineto{\pgfqpoint{4.818749in}{3.142500in}}%
\pgfpathlineto{\pgfqpoint{4.823290in}{3.142500in}}%
\pgfpathlineto{\pgfqpoint{4.823290in}{3.139551in}}%
\pgfpathmoveto{\pgfqpoint{4.818749in}{3.142500in}}%
\pgfpathlineto{\pgfqpoint{4.818749in}{3.142500in}}%
\pgfpathlineto{\pgfqpoint{4.818749in}{3.145449in}}%
\pgfpathlineto{\pgfqpoint{4.823290in}{3.145449in}}%
\pgfpathlineto{\pgfqpoint{4.823290in}{3.142500in}}%
\pgfpathmoveto{\pgfqpoint{4.823290in}{3.142500in}}%
\pgfpathlineto{\pgfqpoint{4.823290in}{3.142500in}}%
\pgfpathlineto{\pgfqpoint{4.823290in}{3.145449in}}%
\pgfpathlineto{\pgfqpoint{4.827831in}{3.145449in}}%
\pgfpathlineto{\pgfqpoint{4.827831in}{3.142500in}}%
\pgfpathmoveto{\pgfqpoint{4.823290in}{3.145449in}}%
\pgfpathlineto{\pgfqpoint{4.823290in}{3.145449in}}%
\pgfpathlineto{\pgfqpoint{4.823290in}{3.148398in}}%
\pgfpathlineto{\pgfqpoint{4.827831in}{3.148398in}}%
\pgfpathlineto{\pgfqpoint{4.827831in}{3.145449in}}%
\pgfpathmoveto{\pgfqpoint{4.827831in}{3.145449in}}%
\pgfpathlineto{\pgfqpoint{4.827831in}{3.145449in}}%
\pgfpathlineto{\pgfqpoint{4.827831in}{3.148398in}}%
\pgfpathlineto{\pgfqpoint{4.832372in}{3.148398in}}%
\pgfpathlineto{\pgfqpoint{4.832372in}{3.145449in}}%
\pgfpathmoveto{\pgfqpoint{4.827831in}{3.148398in}}%
\pgfpathlineto{\pgfqpoint{4.827831in}{3.148398in}}%
\pgfpathlineto{\pgfqpoint{4.827831in}{3.151348in}}%
\pgfpathlineto{\pgfqpoint{4.832372in}{3.151348in}}%
\pgfpathlineto{\pgfqpoint{4.832372in}{3.148398in}}%
\pgfpathmoveto{\pgfqpoint{4.832372in}{3.148398in}}%
\pgfpathlineto{\pgfqpoint{4.832372in}{3.148398in}}%
\pgfpathlineto{\pgfqpoint{4.832372in}{3.151348in}}%
\pgfpathlineto{\pgfqpoint{4.836913in}{3.151348in}}%
\pgfpathlineto{\pgfqpoint{4.836913in}{3.148398in}}%
\pgfpathmoveto{\pgfqpoint{4.832372in}{3.151348in}}%
\pgfpathlineto{\pgfqpoint{4.832372in}{3.151348in}}%
\pgfpathlineto{\pgfqpoint{4.832372in}{3.154297in}}%
\pgfpathlineto{\pgfqpoint{4.836913in}{3.154297in}}%
\pgfpathlineto{\pgfqpoint{4.836913in}{3.151348in}}%
\pgfpathmoveto{\pgfqpoint{4.836913in}{3.151348in}}%
\pgfpathlineto{\pgfqpoint{4.836913in}{3.151348in}}%
\pgfpathlineto{\pgfqpoint{4.836913in}{3.154297in}}%
\pgfpathlineto{\pgfqpoint{4.841454in}{3.154297in}}%
\pgfpathlineto{\pgfqpoint{4.841454in}{3.151348in}}%
\pgfpathmoveto{\pgfqpoint{4.836913in}{3.154297in}}%
\pgfpathlineto{\pgfqpoint{4.836913in}{3.154297in}}%
\pgfpathlineto{\pgfqpoint{4.836913in}{3.157246in}}%
\pgfpathlineto{\pgfqpoint{4.841454in}{3.157246in}}%
\pgfpathlineto{\pgfqpoint{4.841454in}{3.154297in}}%
\pgfpathmoveto{\pgfqpoint{4.841454in}{3.154297in}}%
\pgfpathlineto{\pgfqpoint{4.841454in}{3.154297in}}%
\pgfpathlineto{\pgfqpoint{4.841454in}{3.157246in}}%
\pgfpathlineto{\pgfqpoint{4.845995in}{3.157246in}}%
\pgfpathlineto{\pgfqpoint{4.845995in}{3.154297in}}%
\pgfpathmoveto{\pgfqpoint{4.841454in}{3.157246in}}%
\pgfpathlineto{\pgfqpoint{4.841454in}{3.157246in}}%
\pgfpathlineto{\pgfqpoint{4.841454in}{3.160195in}}%
\pgfpathlineto{\pgfqpoint{4.845995in}{3.160195in}}%
\pgfpathlineto{\pgfqpoint{4.845995in}{3.157246in}}%
\pgfpathmoveto{\pgfqpoint{4.845995in}{3.157246in}}%
\pgfpathlineto{\pgfqpoint{4.845995in}{3.157246in}}%
\pgfpathlineto{\pgfqpoint{4.845995in}{3.160195in}}%
\pgfpathlineto{\pgfqpoint{4.850536in}{3.160195in}}%
\pgfpathlineto{\pgfqpoint{4.850536in}{3.157246in}}%
\pgfpathmoveto{\pgfqpoint{4.845995in}{3.160195in}}%
\pgfpathlineto{\pgfqpoint{4.845995in}{3.160195in}}%
\pgfpathlineto{\pgfqpoint{4.845995in}{3.163144in}}%
\pgfpathlineto{\pgfqpoint{4.850536in}{3.163144in}}%
\pgfpathlineto{\pgfqpoint{4.850536in}{3.160195in}}%
\pgfpathmoveto{\pgfqpoint{4.850536in}{3.160195in}}%
\pgfpathlineto{\pgfqpoint{4.850536in}{3.160195in}}%
\pgfpathlineto{\pgfqpoint{4.850536in}{3.163144in}}%
\pgfpathlineto{\pgfqpoint{4.855077in}{3.163144in}}%
\pgfpathlineto{\pgfqpoint{4.855077in}{3.160195in}}%
\pgfpathmoveto{\pgfqpoint{4.850536in}{3.163144in}}%
\pgfpathlineto{\pgfqpoint{4.850536in}{3.163144in}}%
\pgfpathlineto{\pgfqpoint{4.850536in}{3.166093in}}%
\pgfpathlineto{\pgfqpoint{4.855077in}{3.166093in}}%
\pgfpathlineto{\pgfqpoint{4.855077in}{3.163144in}}%
\pgfpathmoveto{\pgfqpoint{4.855077in}{3.163144in}}%
\pgfpathlineto{\pgfqpoint{4.855077in}{3.163144in}}%
\pgfpathlineto{\pgfqpoint{4.855077in}{3.166093in}}%
\pgfpathlineto{\pgfqpoint{4.859618in}{3.166093in}}%
\pgfpathlineto{\pgfqpoint{4.859618in}{3.163144in}}%
\pgfpathmoveto{\pgfqpoint{4.855077in}{3.166093in}}%
\pgfpathlineto{\pgfqpoint{4.855077in}{3.166093in}}%
\pgfpathlineto{\pgfqpoint{4.855077in}{3.169042in}}%
\pgfpathlineto{\pgfqpoint{4.859618in}{3.169042in}}%
\pgfpathlineto{\pgfqpoint{4.859618in}{3.166093in}}%
\pgfpathmoveto{\pgfqpoint{4.859618in}{3.166093in}}%
\pgfpathlineto{\pgfqpoint{4.859618in}{3.166093in}}%
\pgfpathlineto{\pgfqpoint{4.859618in}{3.169042in}}%
\pgfpathlineto{\pgfqpoint{4.864158in}{3.169042in}}%
\pgfpathlineto{\pgfqpoint{4.864158in}{3.166093in}}%
\pgfpathmoveto{\pgfqpoint{4.859618in}{3.169042in}}%
\pgfpathlineto{\pgfqpoint{4.859618in}{3.169042in}}%
\pgfpathlineto{\pgfqpoint{4.859618in}{3.171992in}}%
\pgfpathlineto{\pgfqpoint{4.864158in}{3.171992in}}%
\pgfpathlineto{\pgfqpoint{4.864158in}{3.169042in}}%
\pgfpathmoveto{\pgfqpoint{4.864158in}{3.169042in}}%
\pgfpathlineto{\pgfqpoint{4.864158in}{3.169042in}}%
\pgfpathlineto{\pgfqpoint{4.864158in}{3.171992in}}%
\pgfpathlineto{\pgfqpoint{4.868699in}{3.171992in}}%
\pgfpathlineto{\pgfqpoint{4.868699in}{3.169042in}}%
\pgfpathmoveto{\pgfqpoint{4.864158in}{3.171992in}}%
\pgfpathlineto{\pgfqpoint{4.864158in}{3.171992in}}%
\pgfpathlineto{\pgfqpoint{4.864158in}{3.174941in}}%
\pgfpathlineto{\pgfqpoint{4.868699in}{3.174941in}}%
\pgfpathlineto{\pgfqpoint{4.868699in}{3.171992in}}%
\pgfpathmoveto{\pgfqpoint{4.868699in}{3.171992in}}%
\pgfpathlineto{\pgfqpoint{4.868699in}{3.171992in}}%
\pgfpathlineto{\pgfqpoint{4.868699in}{3.174941in}}%
\pgfpathlineto{\pgfqpoint{4.873240in}{3.174941in}}%
\pgfpathlineto{\pgfqpoint{4.873240in}{3.171992in}}%
\pgfpathmoveto{\pgfqpoint{4.868699in}{3.174941in}}%
\pgfpathlineto{\pgfqpoint{4.868699in}{3.174941in}}%
\pgfpathlineto{\pgfqpoint{4.868699in}{3.177890in}}%
\pgfpathlineto{\pgfqpoint{4.873240in}{3.177890in}}%
\pgfpathlineto{\pgfqpoint{4.873240in}{3.174941in}}%
\pgfpathmoveto{\pgfqpoint{4.873240in}{3.174941in}}%
\pgfpathlineto{\pgfqpoint{4.873240in}{3.174941in}}%
\pgfpathlineto{\pgfqpoint{4.873240in}{3.177890in}}%
\pgfpathlineto{\pgfqpoint{4.877781in}{3.177890in}}%
\pgfpathlineto{\pgfqpoint{4.877781in}{3.174941in}}%
\pgfpathmoveto{\pgfqpoint{4.873240in}{3.177890in}}%
\pgfpathlineto{\pgfqpoint{4.873240in}{3.177890in}}%
\pgfpathlineto{\pgfqpoint{4.873240in}{3.180839in}}%
\pgfpathlineto{\pgfqpoint{4.877781in}{3.180839in}}%
\pgfpathlineto{\pgfqpoint{4.877781in}{3.177890in}}%
\pgfpathmoveto{\pgfqpoint{4.877781in}{3.177890in}}%
\pgfpathlineto{\pgfqpoint{4.877781in}{3.177890in}}%
\pgfpathlineto{\pgfqpoint{4.877781in}{3.180839in}}%
\pgfpathlineto{\pgfqpoint{4.882322in}{3.180839in}}%
\pgfpathlineto{\pgfqpoint{4.882322in}{3.177890in}}%
\pgfpathmoveto{\pgfqpoint{4.877781in}{3.180839in}}%
\pgfpathlineto{\pgfqpoint{4.877781in}{3.180839in}}%
\pgfpathlineto{\pgfqpoint{4.877781in}{3.183788in}}%
\pgfpathlineto{\pgfqpoint{4.882322in}{3.183788in}}%
\pgfpathlineto{\pgfqpoint{4.882322in}{3.180839in}}%
\pgfpathmoveto{\pgfqpoint{4.882322in}{3.180839in}}%
\pgfpathlineto{\pgfqpoint{4.882322in}{3.180839in}}%
\pgfpathlineto{\pgfqpoint{4.882322in}{3.183788in}}%
\pgfpathlineto{\pgfqpoint{4.886863in}{3.183788in}}%
\pgfpathlineto{\pgfqpoint{4.886863in}{3.180839in}}%
\pgfpathmoveto{\pgfqpoint{4.882322in}{3.183788in}}%
\pgfpathlineto{\pgfqpoint{4.882322in}{3.183788in}}%
\pgfpathlineto{\pgfqpoint{4.882322in}{3.186737in}}%
\pgfpathlineto{\pgfqpoint{4.886863in}{3.186737in}}%
\pgfpathlineto{\pgfqpoint{4.886863in}{3.183788in}}%
\pgfpathmoveto{\pgfqpoint{4.886863in}{3.183788in}}%
\pgfpathlineto{\pgfqpoint{4.886863in}{3.183788in}}%
\pgfpathlineto{\pgfqpoint{4.886863in}{3.186737in}}%
\pgfpathlineto{\pgfqpoint{4.891404in}{3.186737in}}%
\pgfpathlineto{\pgfqpoint{4.891404in}{3.183788in}}%
\pgfpathmoveto{\pgfqpoint{4.886863in}{3.186737in}}%
\pgfpathlineto{\pgfqpoint{4.886863in}{3.186737in}}%
\pgfpathlineto{\pgfqpoint{4.886863in}{3.189686in}}%
\pgfpathlineto{\pgfqpoint{4.891404in}{3.189686in}}%
\pgfpathlineto{\pgfqpoint{4.891404in}{3.186737in}}%
\pgfpathmoveto{\pgfqpoint{4.891404in}{3.186737in}}%
\pgfpathlineto{\pgfqpoint{4.891404in}{3.186737in}}%
\pgfpathlineto{\pgfqpoint{4.891404in}{3.189686in}}%
\pgfpathlineto{\pgfqpoint{4.895945in}{3.189686in}}%
\pgfpathlineto{\pgfqpoint{4.895945in}{3.186737in}}%
\pgfpathmoveto{\pgfqpoint{4.891404in}{3.189686in}}%
\pgfpathlineto{\pgfqpoint{4.891404in}{3.189686in}}%
\pgfpathlineto{\pgfqpoint{4.891404in}{3.192636in}}%
\pgfpathlineto{\pgfqpoint{4.895945in}{3.192636in}}%
\pgfpathlineto{\pgfqpoint{4.895945in}{3.189686in}}%
\pgfpathmoveto{\pgfqpoint{4.895945in}{3.189686in}}%
\pgfpathlineto{\pgfqpoint{4.895945in}{3.189686in}}%
\pgfpathlineto{\pgfqpoint{4.895945in}{3.192636in}}%
\pgfpathlineto{\pgfqpoint{4.900486in}{3.192636in}}%
\pgfpathlineto{\pgfqpoint{4.900486in}{3.189686in}}%
\pgfpathmoveto{\pgfqpoint{4.895945in}{3.192636in}}%
\pgfpathlineto{\pgfqpoint{4.895945in}{3.192636in}}%
\pgfpathlineto{\pgfqpoint{4.895945in}{3.195585in}}%
\pgfpathlineto{\pgfqpoint{4.900486in}{3.195585in}}%
\pgfpathlineto{\pgfqpoint{4.900486in}{3.192636in}}%
\pgfpathmoveto{\pgfqpoint{4.900486in}{3.192636in}}%
\pgfpathlineto{\pgfqpoint{4.900486in}{3.192636in}}%
\pgfpathlineto{\pgfqpoint{4.900486in}{3.195585in}}%
\pgfpathlineto{\pgfqpoint{4.905027in}{3.195585in}}%
\pgfpathlineto{\pgfqpoint{4.905027in}{3.192636in}}%
\pgfpathmoveto{\pgfqpoint{4.900486in}{3.195585in}}%
\pgfpathlineto{\pgfqpoint{4.900486in}{3.195585in}}%
\pgfpathlineto{\pgfqpoint{4.900486in}{3.198534in}}%
\pgfpathlineto{\pgfqpoint{4.905027in}{3.198534in}}%
\pgfpathlineto{\pgfqpoint{4.905027in}{3.195585in}}%
\pgfpathmoveto{\pgfqpoint{4.905027in}{3.195585in}}%
\pgfpathlineto{\pgfqpoint{4.905027in}{3.195585in}}%
\pgfpathlineto{\pgfqpoint{4.905027in}{3.198534in}}%
\pgfpathlineto{\pgfqpoint{4.909567in}{3.198534in}}%
\pgfpathlineto{\pgfqpoint{4.909567in}{3.195585in}}%
\pgfpathmoveto{\pgfqpoint{4.905027in}{3.198534in}}%
\pgfpathlineto{\pgfqpoint{4.905027in}{3.198534in}}%
\pgfpathlineto{\pgfqpoint{4.905027in}{3.201483in}}%
\pgfpathlineto{\pgfqpoint{4.909567in}{3.201483in}}%
\pgfpathlineto{\pgfqpoint{4.909567in}{3.198534in}}%
\pgfpathmoveto{\pgfqpoint{4.909567in}{3.198534in}}%
\pgfpathlineto{\pgfqpoint{4.909567in}{3.198534in}}%
\pgfpathlineto{\pgfqpoint{4.909567in}{3.201483in}}%
\pgfpathlineto{\pgfqpoint{4.914108in}{3.201483in}}%
\pgfpathlineto{\pgfqpoint{4.914108in}{3.198534in}}%
\pgfpathmoveto{\pgfqpoint{4.909567in}{3.201483in}}%
\pgfpathlineto{\pgfqpoint{4.909567in}{3.201483in}}%
\pgfpathlineto{\pgfqpoint{4.909567in}{3.204432in}}%
\pgfpathlineto{\pgfqpoint{4.914108in}{3.204432in}}%
\pgfpathlineto{\pgfqpoint{4.914108in}{3.201483in}}%
\pgfpathmoveto{\pgfqpoint{4.914108in}{3.201483in}}%
\pgfpathlineto{\pgfqpoint{4.914108in}{3.201483in}}%
\pgfpathlineto{\pgfqpoint{4.914108in}{3.204432in}}%
\pgfpathlineto{\pgfqpoint{4.918649in}{3.204432in}}%
\pgfpathlineto{\pgfqpoint{4.918649in}{3.201483in}}%
\pgfpathmoveto{\pgfqpoint{4.914108in}{3.204432in}}%
\pgfpathlineto{\pgfqpoint{4.914108in}{3.204432in}}%
\pgfpathlineto{\pgfqpoint{4.914108in}{3.207381in}}%
\pgfpathlineto{\pgfqpoint{4.918649in}{3.207381in}}%
\pgfpathlineto{\pgfqpoint{4.918649in}{3.204432in}}%
\pgfpathmoveto{\pgfqpoint{4.918649in}{3.204432in}}%
\pgfpathlineto{\pgfqpoint{4.918649in}{3.204432in}}%
\pgfpathlineto{\pgfqpoint{4.918649in}{3.207381in}}%
\pgfpathlineto{\pgfqpoint{4.923190in}{3.207381in}}%
\pgfpathlineto{\pgfqpoint{4.923190in}{3.204432in}}%
\pgfpathmoveto{\pgfqpoint{4.918649in}{3.207381in}}%
\pgfpathlineto{\pgfqpoint{4.918649in}{3.207381in}}%
\pgfpathlineto{\pgfqpoint{4.918649in}{3.210330in}}%
\pgfpathlineto{\pgfqpoint{4.923190in}{3.210330in}}%
\pgfpathlineto{\pgfqpoint{4.923190in}{3.207381in}}%
\pgfpathmoveto{\pgfqpoint{4.923190in}{3.207381in}}%
\pgfpathlineto{\pgfqpoint{4.923190in}{3.207381in}}%
\pgfpathlineto{\pgfqpoint{4.923190in}{3.210330in}}%
\pgfpathlineto{\pgfqpoint{4.927731in}{3.210330in}}%
\pgfpathlineto{\pgfqpoint{4.927731in}{3.207381in}}%
\pgfpathmoveto{\pgfqpoint{4.923190in}{3.210330in}}%
\pgfpathlineto{\pgfqpoint{4.923190in}{3.210330in}}%
\pgfpathlineto{\pgfqpoint{4.923190in}{3.213280in}}%
\pgfpathlineto{\pgfqpoint{4.927731in}{3.213280in}}%
\pgfpathlineto{\pgfqpoint{4.927731in}{3.210330in}}%
\pgfpathmoveto{\pgfqpoint{4.927731in}{3.210330in}}%
\pgfpathlineto{\pgfqpoint{4.927731in}{3.210330in}}%
\pgfpathlineto{\pgfqpoint{4.927731in}{3.213280in}}%
\pgfpathlineto{\pgfqpoint{4.932272in}{3.213280in}}%
\pgfpathlineto{\pgfqpoint{4.932272in}{3.210330in}}%
\pgfpathmoveto{\pgfqpoint{4.927731in}{3.213280in}}%
\pgfpathlineto{\pgfqpoint{4.927731in}{3.213280in}}%
\pgfpathlineto{\pgfqpoint{4.927731in}{3.216229in}}%
\pgfpathlineto{\pgfqpoint{4.932272in}{3.216229in}}%
\pgfpathlineto{\pgfqpoint{4.932272in}{3.213280in}}%
\pgfpathmoveto{\pgfqpoint{4.932272in}{3.213280in}}%
\pgfpathlineto{\pgfqpoint{4.932272in}{3.213280in}}%
\pgfpathlineto{\pgfqpoint{4.932272in}{3.216229in}}%
\pgfpathlineto{\pgfqpoint{4.936813in}{3.216229in}}%
\pgfpathlineto{\pgfqpoint{4.936813in}{3.213280in}}%
\pgfpathmoveto{\pgfqpoint{4.932272in}{3.216229in}}%
\pgfpathlineto{\pgfqpoint{4.932272in}{3.216229in}}%
\pgfpathlineto{\pgfqpoint{4.932272in}{3.219178in}}%
\pgfpathlineto{\pgfqpoint{4.936813in}{3.219178in}}%
\pgfpathlineto{\pgfqpoint{4.936813in}{3.216229in}}%
\pgfpathmoveto{\pgfqpoint{4.936813in}{3.216229in}}%
\pgfpathlineto{\pgfqpoint{4.936813in}{3.216229in}}%
\pgfpathlineto{\pgfqpoint{4.936813in}{3.219178in}}%
\pgfpathlineto{\pgfqpoint{4.941354in}{3.219178in}}%
\pgfpathlineto{\pgfqpoint{4.941354in}{3.216229in}}%
\pgfpathmoveto{\pgfqpoint{4.936813in}{3.219178in}}%
\pgfpathlineto{\pgfqpoint{4.936813in}{3.219178in}}%
\pgfpathlineto{\pgfqpoint{4.936813in}{3.222127in}}%
\pgfpathlineto{\pgfqpoint{4.941354in}{3.222127in}}%
\pgfpathlineto{\pgfqpoint{4.941354in}{3.219178in}}%
\pgfpathmoveto{\pgfqpoint{4.941354in}{3.219178in}}%
\pgfpathlineto{\pgfqpoint{4.941354in}{3.219178in}}%
\pgfpathlineto{\pgfqpoint{4.941354in}{3.222127in}}%
\pgfpathlineto{\pgfqpoint{4.945895in}{3.222127in}}%
\pgfpathlineto{\pgfqpoint{4.945895in}{3.219178in}}%
\pgfpathmoveto{\pgfqpoint{4.941354in}{3.222127in}}%
\pgfpathlineto{\pgfqpoint{4.941354in}{3.222127in}}%
\pgfpathlineto{\pgfqpoint{4.941354in}{3.225076in}}%
\pgfpathlineto{\pgfqpoint{4.945895in}{3.225076in}}%
\pgfpathlineto{\pgfqpoint{4.945895in}{3.222127in}}%
\pgfpathmoveto{\pgfqpoint{4.945895in}{3.222127in}}%
\pgfpathlineto{\pgfqpoint{4.945895in}{3.222127in}}%
\pgfpathlineto{\pgfqpoint{4.945895in}{3.225076in}}%
\pgfpathlineto{\pgfqpoint{4.950436in}{3.225076in}}%
\pgfpathlineto{\pgfqpoint{4.950436in}{3.222127in}}%
\pgfpathmoveto{\pgfqpoint{4.945895in}{3.225076in}}%
\pgfpathlineto{\pgfqpoint{4.945895in}{3.225076in}}%
\pgfpathlineto{\pgfqpoint{4.945895in}{3.228025in}}%
\pgfpathlineto{\pgfqpoint{4.950436in}{3.228025in}}%
\pgfpathlineto{\pgfqpoint{4.950436in}{3.225076in}}%
\pgfpathmoveto{\pgfqpoint{4.950436in}{3.225076in}}%
\pgfpathlineto{\pgfqpoint{4.950436in}{3.225076in}}%
\pgfpathlineto{\pgfqpoint{4.950436in}{3.228025in}}%
\pgfpathlineto{\pgfqpoint{4.954977in}{3.228025in}}%
\pgfpathlineto{\pgfqpoint{4.954977in}{3.225076in}}%
\pgfpathmoveto{\pgfqpoint{4.950436in}{3.228025in}}%
\pgfpathlineto{\pgfqpoint{4.950436in}{3.228025in}}%
\pgfpathlineto{\pgfqpoint{4.950436in}{3.230974in}}%
\pgfpathlineto{\pgfqpoint{4.954977in}{3.230974in}}%
\pgfpathlineto{\pgfqpoint{4.954977in}{3.228025in}}%
\pgfpathmoveto{\pgfqpoint{4.954977in}{3.228025in}}%
\pgfpathlineto{\pgfqpoint{4.954977in}{3.228025in}}%
\pgfpathlineto{\pgfqpoint{4.954977in}{3.230974in}}%
\pgfpathlineto{\pgfqpoint{4.959517in}{3.230974in}}%
\pgfpathlineto{\pgfqpoint{4.959517in}{3.228025in}}%
\pgfpathmoveto{\pgfqpoint{4.954977in}{3.230974in}}%
\pgfpathlineto{\pgfqpoint{4.954977in}{3.230974in}}%
\pgfpathlineto{\pgfqpoint{4.954977in}{3.233924in}}%
\pgfpathlineto{\pgfqpoint{4.959517in}{3.233924in}}%
\pgfpathlineto{\pgfqpoint{4.959517in}{3.230974in}}%
\pgfpathmoveto{\pgfqpoint{4.959517in}{3.230974in}}%
\pgfpathlineto{\pgfqpoint{4.959517in}{3.230974in}}%
\pgfpathlineto{\pgfqpoint{4.959517in}{3.233924in}}%
\pgfpathlineto{\pgfqpoint{4.964058in}{3.233924in}}%
\pgfpathlineto{\pgfqpoint{4.964058in}{3.230974in}}%
\pgfpathmoveto{\pgfqpoint{4.959517in}{3.233924in}}%
\pgfpathlineto{\pgfqpoint{4.959517in}{3.233924in}}%
\pgfpathlineto{\pgfqpoint{4.959517in}{3.236873in}}%
\pgfpathlineto{\pgfqpoint{4.964058in}{3.236873in}}%
\pgfpathlineto{\pgfqpoint{4.964058in}{3.233924in}}%
\pgfpathmoveto{\pgfqpoint{4.964058in}{3.233924in}}%
\pgfpathlineto{\pgfqpoint{4.964058in}{3.233924in}}%
\pgfpathlineto{\pgfqpoint{4.964058in}{3.236873in}}%
\pgfpathlineto{\pgfqpoint{4.968599in}{3.236873in}}%
\pgfpathlineto{\pgfqpoint{4.968599in}{3.233924in}}%
\pgfpathmoveto{\pgfqpoint{4.964058in}{3.236873in}}%
\pgfpathlineto{\pgfqpoint{4.964058in}{3.236873in}}%
\pgfpathlineto{\pgfqpoint{4.964058in}{3.239822in}}%
\pgfpathlineto{\pgfqpoint{4.968599in}{3.239822in}}%
\pgfpathlineto{\pgfqpoint{4.968599in}{3.236873in}}%
\pgfpathmoveto{\pgfqpoint{4.968599in}{3.236873in}}%
\pgfpathlineto{\pgfqpoint{4.968599in}{3.236873in}}%
\pgfpathlineto{\pgfqpoint{4.968599in}{3.239822in}}%
\pgfpathlineto{\pgfqpoint{4.973140in}{3.239822in}}%
\pgfpathlineto{\pgfqpoint{4.973140in}{3.236873in}}%
\pgfpathmoveto{\pgfqpoint{4.968599in}{3.239822in}}%
\pgfpathlineto{\pgfqpoint{4.968599in}{3.239822in}}%
\pgfpathlineto{\pgfqpoint{4.968599in}{3.242771in}}%
\pgfpathlineto{\pgfqpoint{4.973140in}{3.242771in}}%
\pgfpathlineto{\pgfqpoint{4.973140in}{3.239822in}}%
\pgfpathmoveto{\pgfqpoint{4.973140in}{3.239822in}}%
\pgfpathlineto{\pgfqpoint{4.973140in}{3.239822in}}%
\pgfpathlineto{\pgfqpoint{4.973140in}{3.242771in}}%
\pgfpathlineto{\pgfqpoint{4.977681in}{3.242771in}}%
\pgfpathlineto{\pgfqpoint{4.977681in}{3.239822in}}%
\pgfpathmoveto{\pgfqpoint{4.973140in}{3.242771in}}%
\pgfpathlineto{\pgfqpoint{4.973140in}{3.242771in}}%
\pgfpathlineto{\pgfqpoint{4.973140in}{3.245721in}}%
\pgfpathlineto{\pgfqpoint{4.977681in}{3.245721in}}%
\pgfpathlineto{\pgfqpoint{4.977681in}{3.242771in}}%
\pgfpathmoveto{\pgfqpoint{4.977681in}{3.242771in}}%
\pgfpathlineto{\pgfqpoint{4.977681in}{3.242771in}}%
\pgfpathlineto{\pgfqpoint{4.977681in}{3.245721in}}%
\pgfpathlineto{\pgfqpoint{4.982223in}{3.245721in}}%
\pgfpathlineto{\pgfqpoint{4.982223in}{3.242771in}}%
\pgfpathmoveto{\pgfqpoint{4.977681in}{3.245721in}}%
\pgfpathlineto{\pgfqpoint{4.977681in}{3.245721in}}%
\pgfpathlineto{\pgfqpoint{4.977681in}{3.248670in}}%
\pgfpathlineto{\pgfqpoint{4.982223in}{3.248670in}}%
\pgfpathlineto{\pgfqpoint{4.982223in}{3.245721in}}%
\pgfpathmoveto{\pgfqpoint{4.982223in}{3.245721in}}%
\pgfpathlineto{\pgfqpoint{4.982223in}{3.245721in}}%
\pgfpathlineto{\pgfqpoint{4.982223in}{3.248670in}}%
\pgfpathlineto{\pgfqpoint{4.986764in}{3.248670in}}%
\pgfpathlineto{\pgfqpoint{4.986764in}{3.245721in}}%
\pgfpathmoveto{\pgfqpoint{4.982223in}{3.248670in}}%
\pgfpathlineto{\pgfqpoint{4.982223in}{3.248670in}}%
\pgfpathlineto{\pgfqpoint{4.982223in}{3.251619in}}%
\pgfpathlineto{\pgfqpoint{4.986764in}{3.251619in}}%
\pgfpathlineto{\pgfqpoint{4.986764in}{3.248670in}}%
\pgfpathmoveto{\pgfqpoint{4.986764in}{3.248670in}}%
\pgfpathlineto{\pgfqpoint{4.986764in}{3.248670in}}%
\pgfpathlineto{\pgfqpoint{4.986764in}{3.251619in}}%
\pgfpathlineto{\pgfqpoint{4.991305in}{3.251619in}}%
\pgfpathlineto{\pgfqpoint{4.991305in}{3.248670in}}%
\pgfpathmoveto{\pgfqpoint{4.986764in}{3.251619in}}%
\pgfpathlineto{\pgfqpoint{4.986764in}{3.251619in}}%
\pgfpathlineto{\pgfqpoint{4.986764in}{3.254569in}}%
\pgfpathlineto{\pgfqpoint{4.991305in}{3.254569in}}%
\pgfpathlineto{\pgfqpoint{4.991305in}{3.251619in}}%
\pgfpathmoveto{\pgfqpoint{4.991305in}{3.251619in}}%
\pgfpathlineto{\pgfqpoint{4.991305in}{3.251619in}}%
\pgfpathlineto{\pgfqpoint{4.991305in}{3.254569in}}%
\pgfpathlineto{\pgfqpoint{4.995846in}{3.254569in}}%
\pgfpathlineto{\pgfqpoint{4.995846in}{3.251619in}}%
\pgfpathmoveto{\pgfqpoint{4.991305in}{3.254569in}}%
\pgfpathlineto{\pgfqpoint{4.991305in}{3.254569in}}%
\pgfpathlineto{\pgfqpoint{4.991305in}{3.257518in}}%
\pgfpathlineto{\pgfqpoint{4.995846in}{3.257518in}}%
\pgfpathlineto{\pgfqpoint{4.995846in}{3.254569in}}%
\pgfpathmoveto{\pgfqpoint{4.995846in}{3.254569in}}%
\pgfpathlineto{\pgfqpoint{4.995846in}{3.254569in}}%
\pgfpathlineto{\pgfqpoint{4.995846in}{3.257518in}}%
\pgfpathlineto{\pgfqpoint{5.000387in}{3.257518in}}%
\pgfpathlineto{\pgfqpoint{5.000387in}{3.254569in}}%
\pgfpathmoveto{\pgfqpoint{4.995846in}{3.257518in}}%
\pgfpathlineto{\pgfqpoint{4.995846in}{3.257518in}}%
\pgfpathlineto{\pgfqpoint{4.995846in}{3.260467in}}%
\pgfpathlineto{\pgfqpoint{5.000387in}{3.260467in}}%
\pgfpathlineto{\pgfqpoint{5.000387in}{3.257518in}}%
\pgfpathmoveto{\pgfqpoint{5.000387in}{3.257518in}}%
\pgfpathlineto{\pgfqpoint{5.000387in}{3.257518in}}%
\pgfpathlineto{\pgfqpoint{5.000387in}{3.260467in}}%
\pgfpathlineto{\pgfqpoint{5.004928in}{3.260467in}}%
\pgfpathlineto{\pgfqpoint{5.004928in}{3.257518in}}%
\pgfpathmoveto{\pgfqpoint{5.000387in}{3.260467in}}%
\pgfpathlineto{\pgfqpoint{5.000387in}{3.260467in}}%
\pgfpathlineto{\pgfqpoint{5.000387in}{3.263417in}}%
\pgfpathlineto{\pgfqpoint{5.004928in}{3.263417in}}%
\pgfpathlineto{\pgfqpoint{5.004928in}{3.260467in}}%
\pgfpathmoveto{\pgfqpoint{5.004928in}{3.260467in}}%
\pgfpathlineto{\pgfqpoint{5.004928in}{3.260467in}}%
\pgfpathlineto{\pgfqpoint{5.004928in}{3.263417in}}%
\pgfpathlineto{\pgfqpoint{5.009469in}{3.263417in}}%
\pgfpathlineto{\pgfqpoint{5.009469in}{3.260467in}}%
\pgfpathmoveto{\pgfqpoint{5.004928in}{3.263417in}}%
\pgfpathlineto{\pgfqpoint{5.004928in}{3.263417in}}%
\pgfpathlineto{\pgfqpoint{5.004928in}{3.266366in}}%
\pgfpathlineto{\pgfqpoint{5.009469in}{3.266366in}}%
\pgfpathlineto{\pgfqpoint{5.009469in}{3.263417in}}%
\pgfpathmoveto{\pgfqpoint{5.009469in}{3.263417in}}%
\pgfpathlineto{\pgfqpoint{5.009469in}{3.263417in}}%
\pgfpathlineto{\pgfqpoint{5.009469in}{3.266366in}}%
\pgfpathlineto{\pgfqpoint{5.014010in}{3.266366in}}%
\pgfpathlineto{\pgfqpoint{5.014010in}{3.263417in}}%
\pgfpathmoveto{\pgfqpoint{5.009469in}{3.266366in}}%
\pgfpathlineto{\pgfqpoint{5.009469in}{3.266366in}}%
\pgfpathlineto{\pgfqpoint{5.009469in}{3.269315in}}%
\pgfpathlineto{\pgfqpoint{5.014010in}{3.269315in}}%
\pgfpathlineto{\pgfqpoint{5.014010in}{3.266366in}}%
\pgfpathmoveto{\pgfqpoint{5.014010in}{3.266366in}}%
\pgfpathlineto{\pgfqpoint{5.014010in}{3.266366in}}%
\pgfpathlineto{\pgfqpoint{5.014010in}{3.269315in}}%
\pgfpathlineto{\pgfqpoint{5.018551in}{3.269315in}}%
\pgfpathlineto{\pgfqpoint{5.018551in}{3.266366in}}%
\pgfpathmoveto{\pgfqpoint{5.014010in}{3.269315in}}%
\pgfpathlineto{\pgfqpoint{5.014010in}{3.269315in}}%
\pgfpathlineto{\pgfqpoint{5.014010in}{3.272265in}}%
\pgfpathlineto{\pgfqpoint{5.018551in}{3.272265in}}%
\pgfpathlineto{\pgfqpoint{5.018551in}{3.269315in}}%
\pgfpathmoveto{\pgfqpoint{5.018551in}{3.269315in}}%
\pgfpathlineto{\pgfqpoint{5.018551in}{3.269315in}}%
\pgfpathlineto{\pgfqpoint{5.018551in}{3.272265in}}%
\pgfpathlineto{\pgfqpoint{5.023092in}{3.272265in}}%
\pgfpathlineto{\pgfqpoint{5.023092in}{3.269315in}}%
\pgfpathmoveto{\pgfqpoint{5.018551in}{3.272265in}}%
\pgfpathlineto{\pgfqpoint{5.018551in}{3.272265in}}%
\pgfpathlineto{\pgfqpoint{5.018551in}{3.275214in}}%
\pgfpathlineto{\pgfqpoint{5.023092in}{3.275214in}}%
\pgfpathlineto{\pgfqpoint{5.023092in}{3.272265in}}%
\pgfpathmoveto{\pgfqpoint{5.023092in}{3.272265in}}%
\pgfpathlineto{\pgfqpoint{5.023092in}{3.272265in}}%
\pgfpathlineto{\pgfqpoint{5.023092in}{3.275214in}}%
\pgfpathlineto{\pgfqpoint{5.027633in}{3.275214in}}%
\pgfpathlineto{\pgfqpoint{5.027633in}{3.272265in}}%
\pgfpathmoveto{\pgfqpoint{5.023092in}{3.275214in}}%
\pgfpathlineto{\pgfqpoint{5.023092in}{3.275214in}}%
\pgfpathlineto{\pgfqpoint{5.023092in}{3.278163in}}%
\pgfpathlineto{\pgfqpoint{5.027633in}{3.278163in}}%
\pgfpathlineto{\pgfqpoint{5.027633in}{3.275214in}}%
\pgfpathmoveto{\pgfqpoint{5.027633in}{3.275214in}}%
\pgfpathlineto{\pgfqpoint{5.027633in}{3.275214in}}%
\pgfpathlineto{\pgfqpoint{5.027633in}{3.278163in}}%
\pgfpathlineto{\pgfqpoint{5.032174in}{3.278163in}}%
\pgfpathlineto{\pgfqpoint{5.032174in}{3.275214in}}%
\pgfpathmoveto{\pgfqpoint{5.027633in}{3.278163in}}%
\pgfpathlineto{\pgfqpoint{5.027633in}{3.278163in}}%
\pgfpathlineto{\pgfqpoint{5.027633in}{3.281113in}}%
\pgfpathlineto{\pgfqpoint{5.032174in}{3.281113in}}%
\pgfpathlineto{\pgfqpoint{5.032174in}{3.278163in}}%
\pgfpathmoveto{\pgfqpoint{5.032174in}{3.278163in}}%
\pgfpathlineto{\pgfqpoint{5.032174in}{3.278163in}}%
\pgfpathlineto{\pgfqpoint{5.032174in}{3.281113in}}%
\pgfpathlineto{\pgfqpoint{5.036715in}{3.281113in}}%
\pgfpathlineto{\pgfqpoint{5.036715in}{3.278163in}}%
\pgfpathmoveto{\pgfqpoint{5.032174in}{3.281113in}}%
\pgfpathlineto{\pgfqpoint{5.032174in}{3.281113in}}%
\pgfpathlineto{\pgfqpoint{5.032174in}{3.284062in}}%
\pgfpathlineto{\pgfqpoint{5.036715in}{3.284062in}}%
\pgfpathlineto{\pgfqpoint{5.036715in}{3.281113in}}%
\pgfpathmoveto{\pgfqpoint{5.036715in}{3.281113in}}%
\pgfpathlineto{\pgfqpoint{5.036715in}{3.281113in}}%
\pgfpathlineto{\pgfqpoint{5.036715in}{3.284062in}}%
\pgfpathlineto{\pgfqpoint{5.041256in}{3.284062in}}%
\pgfpathlineto{\pgfqpoint{5.041256in}{3.281113in}}%
\pgfpathmoveto{\pgfqpoint{5.036715in}{3.284062in}}%
\pgfpathlineto{\pgfqpoint{5.036715in}{3.284062in}}%
\pgfpathlineto{\pgfqpoint{5.036715in}{3.287011in}}%
\pgfpathlineto{\pgfqpoint{5.041256in}{3.287011in}}%
\pgfpathlineto{\pgfqpoint{5.041256in}{3.284062in}}%
\pgfpathmoveto{\pgfqpoint{5.041256in}{3.284062in}}%
\pgfpathlineto{\pgfqpoint{5.041256in}{3.284062in}}%
\pgfpathlineto{\pgfqpoint{5.041256in}{3.287011in}}%
\pgfpathlineto{\pgfqpoint{5.045797in}{3.287011in}}%
\pgfpathlineto{\pgfqpoint{5.045797in}{3.284062in}}%
\pgfpathmoveto{\pgfqpoint{5.041256in}{3.287011in}}%
\pgfpathlineto{\pgfqpoint{5.041256in}{3.287011in}}%
\pgfpathlineto{\pgfqpoint{5.041256in}{3.289961in}}%
\pgfpathlineto{\pgfqpoint{5.045797in}{3.289961in}}%
\pgfpathlineto{\pgfqpoint{5.045797in}{3.287011in}}%
\pgfpathmoveto{\pgfqpoint{5.045797in}{3.287011in}}%
\pgfpathlineto{\pgfqpoint{5.045797in}{3.287011in}}%
\pgfpathlineto{\pgfqpoint{5.045797in}{3.289961in}}%
\pgfpathlineto{\pgfqpoint{5.050338in}{3.289961in}}%
\pgfpathlineto{\pgfqpoint{5.050338in}{3.287011in}}%
\pgfpathmoveto{\pgfqpoint{5.045797in}{3.289961in}}%
\pgfpathlineto{\pgfqpoint{5.045797in}{3.289961in}}%
\pgfpathlineto{\pgfqpoint{5.045797in}{3.292910in}}%
\pgfpathlineto{\pgfqpoint{5.050338in}{3.292910in}}%
\pgfpathlineto{\pgfqpoint{5.050338in}{3.289961in}}%
\pgfpathmoveto{\pgfqpoint{5.050338in}{3.289961in}}%
\pgfpathlineto{\pgfqpoint{5.050338in}{3.289961in}}%
\pgfpathlineto{\pgfqpoint{5.050338in}{3.292910in}}%
\pgfpathlineto{\pgfqpoint{5.054879in}{3.292910in}}%
\pgfpathlineto{\pgfqpoint{5.054879in}{3.289961in}}%
\pgfpathmoveto{\pgfqpoint{5.050338in}{3.292910in}}%
\pgfpathlineto{\pgfqpoint{5.050338in}{3.292910in}}%
\pgfpathlineto{\pgfqpoint{5.050338in}{3.295859in}}%
\pgfpathlineto{\pgfqpoint{5.054879in}{3.295859in}}%
\pgfpathlineto{\pgfqpoint{5.054879in}{3.292910in}}%
\pgfpathmoveto{\pgfqpoint{5.054879in}{3.292910in}}%
\pgfpathlineto{\pgfqpoint{5.054879in}{3.292910in}}%
\pgfpathlineto{\pgfqpoint{5.054879in}{3.295859in}}%
\pgfpathlineto{\pgfqpoint{5.059420in}{3.295859in}}%
\pgfpathlineto{\pgfqpoint{5.059420in}{3.292910in}}%
\pgfpathmoveto{\pgfqpoint{5.054879in}{3.295859in}}%
\pgfpathlineto{\pgfqpoint{5.054879in}{3.295859in}}%
\pgfpathlineto{\pgfqpoint{5.054879in}{3.298809in}}%
\pgfpathlineto{\pgfqpoint{5.059420in}{3.298809in}}%
\pgfpathlineto{\pgfqpoint{5.059420in}{3.295859in}}%
\pgfpathmoveto{\pgfqpoint{5.059420in}{3.295859in}}%
\pgfpathlineto{\pgfqpoint{5.059420in}{3.295859in}}%
\pgfpathlineto{\pgfqpoint{5.059420in}{3.298809in}}%
\pgfpathlineto{\pgfqpoint{5.063961in}{3.298809in}}%
\pgfpathlineto{\pgfqpoint{5.063961in}{3.295859in}}%
\pgfpathmoveto{\pgfqpoint{5.059420in}{3.298809in}}%
\pgfpathlineto{\pgfqpoint{5.059420in}{3.298809in}}%
\pgfpathlineto{\pgfqpoint{5.059420in}{3.301758in}}%
\pgfpathlineto{\pgfqpoint{5.063961in}{3.301758in}}%
\pgfpathlineto{\pgfqpoint{5.063961in}{3.298809in}}%
\pgfpathmoveto{\pgfqpoint{5.063961in}{3.298809in}}%
\pgfpathlineto{\pgfqpoint{5.063961in}{3.298809in}}%
\pgfpathlineto{\pgfqpoint{5.063961in}{3.301758in}}%
\pgfpathlineto{\pgfqpoint{5.068502in}{3.301758in}}%
\pgfpathlineto{\pgfqpoint{5.068502in}{3.298809in}}%
\pgfpathmoveto{\pgfqpoint{5.063961in}{3.301758in}}%
\pgfpathlineto{\pgfqpoint{5.063961in}{3.301758in}}%
\pgfpathlineto{\pgfqpoint{5.063961in}{3.304707in}}%
\pgfpathlineto{\pgfqpoint{5.068502in}{3.304707in}}%
\pgfpathlineto{\pgfqpoint{5.068502in}{3.301758in}}%
\pgfpathmoveto{\pgfqpoint{5.068502in}{3.301758in}}%
\pgfpathlineto{\pgfqpoint{5.068502in}{3.301758in}}%
\pgfpathlineto{\pgfqpoint{5.068502in}{3.304707in}}%
\pgfpathlineto{\pgfqpoint{5.073044in}{3.304707in}}%
\pgfpathlineto{\pgfqpoint{5.073044in}{3.301758in}}%
\pgfpathmoveto{\pgfqpoint{5.068502in}{3.304707in}}%
\pgfpathlineto{\pgfqpoint{5.068502in}{3.304707in}}%
\pgfpathlineto{\pgfqpoint{5.068502in}{3.307657in}}%
\pgfpathlineto{\pgfqpoint{5.073044in}{3.307657in}}%
\pgfpathlineto{\pgfqpoint{5.073044in}{3.304707in}}%
\pgfpathmoveto{\pgfqpoint{5.073044in}{3.304707in}}%
\pgfpathlineto{\pgfqpoint{5.073044in}{3.304707in}}%
\pgfpathlineto{\pgfqpoint{5.073044in}{3.307657in}}%
\pgfpathlineto{\pgfqpoint{5.077585in}{3.307657in}}%
\pgfpathlineto{\pgfqpoint{5.077585in}{3.304707in}}%
\pgfpathmoveto{\pgfqpoint{5.073044in}{3.307657in}}%
\pgfpathlineto{\pgfqpoint{5.073044in}{3.307657in}}%
\pgfpathlineto{\pgfqpoint{5.073044in}{3.310606in}}%
\pgfpathlineto{\pgfqpoint{5.077585in}{3.310606in}}%
\pgfpathlineto{\pgfqpoint{5.077585in}{3.307657in}}%
\pgfpathmoveto{\pgfqpoint{5.077585in}{3.307657in}}%
\pgfpathlineto{\pgfqpoint{5.077585in}{3.307657in}}%
\pgfpathlineto{\pgfqpoint{5.077585in}{3.310606in}}%
\pgfpathlineto{\pgfqpoint{5.082126in}{3.310606in}}%
\pgfpathlineto{\pgfqpoint{5.082126in}{3.307657in}}%
\pgfpathmoveto{\pgfqpoint{5.077585in}{3.310606in}}%
\pgfpathlineto{\pgfqpoint{5.077585in}{3.310606in}}%
\pgfpathlineto{\pgfqpoint{5.077585in}{3.313555in}}%
\pgfpathlineto{\pgfqpoint{5.082126in}{3.313555in}}%
\pgfpathlineto{\pgfqpoint{5.082126in}{3.310606in}}%
\pgfpathmoveto{\pgfqpoint{5.082126in}{3.310606in}}%
\pgfpathlineto{\pgfqpoint{5.082126in}{3.310606in}}%
\pgfpathlineto{\pgfqpoint{5.082126in}{3.313555in}}%
\pgfpathlineto{\pgfqpoint{5.086667in}{3.313555in}}%
\pgfpathlineto{\pgfqpoint{5.086667in}{3.310606in}}%
\pgfpathmoveto{\pgfqpoint{5.082126in}{3.313555in}}%
\pgfpathlineto{\pgfqpoint{5.082126in}{3.313555in}}%
\pgfpathlineto{\pgfqpoint{5.082126in}{3.316505in}}%
\pgfpathlineto{\pgfqpoint{5.086667in}{3.316505in}}%
\pgfpathlineto{\pgfqpoint{5.086667in}{3.313555in}}%
\pgfpathmoveto{\pgfqpoint{5.086667in}{3.313555in}}%
\pgfpathlineto{\pgfqpoint{5.086667in}{3.313555in}}%
\pgfpathlineto{\pgfqpoint{5.086667in}{3.316505in}}%
\pgfpathlineto{\pgfqpoint{5.091208in}{3.316505in}}%
\pgfpathlineto{\pgfqpoint{5.091208in}{3.313555in}}%
\pgfpathmoveto{\pgfqpoint{5.086667in}{3.316505in}}%
\pgfpathlineto{\pgfqpoint{5.086667in}{3.316505in}}%
\pgfpathlineto{\pgfqpoint{5.086667in}{3.319454in}}%
\pgfpathlineto{\pgfqpoint{5.091208in}{3.319454in}}%
\pgfpathlineto{\pgfqpoint{5.091208in}{3.316505in}}%
\pgfpathmoveto{\pgfqpoint{5.091208in}{3.316505in}}%
\pgfpathlineto{\pgfqpoint{5.091208in}{3.316505in}}%
\pgfpathlineto{\pgfqpoint{5.091208in}{3.319454in}}%
\pgfpathlineto{\pgfqpoint{5.095749in}{3.319454in}}%
\pgfpathlineto{\pgfqpoint{5.095749in}{3.316505in}}%
\pgfpathmoveto{\pgfqpoint{5.091208in}{3.319454in}}%
\pgfpathlineto{\pgfqpoint{5.091208in}{3.319454in}}%
\pgfpathlineto{\pgfqpoint{5.091208in}{3.322403in}}%
\pgfpathlineto{\pgfqpoint{5.095749in}{3.322403in}}%
\pgfpathlineto{\pgfqpoint{5.095749in}{3.319454in}}%
\pgfpathmoveto{\pgfqpoint{5.095749in}{3.319454in}}%
\pgfpathlineto{\pgfqpoint{5.095749in}{3.319454in}}%
\pgfpathlineto{\pgfqpoint{5.095749in}{3.322403in}}%
\pgfpathlineto{\pgfqpoint{5.100290in}{3.322403in}}%
\pgfpathlineto{\pgfqpoint{5.100290in}{3.319454in}}%
\pgfpathmoveto{\pgfqpoint{5.095749in}{3.322403in}}%
\pgfpathlineto{\pgfqpoint{5.095749in}{3.322403in}}%
\pgfpathlineto{\pgfqpoint{5.095749in}{3.325353in}}%
\pgfpathlineto{\pgfqpoint{5.100290in}{3.325353in}}%
\pgfpathlineto{\pgfqpoint{5.100290in}{3.322403in}}%
\pgfpathmoveto{\pgfqpoint{5.100290in}{3.322403in}}%
\pgfpathlineto{\pgfqpoint{5.100290in}{3.322403in}}%
\pgfpathlineto{\pgfqpoint{5.100290in}{3.325353in}}%
\pgfpathlineto{\pgfqpoint{5.104831in}{3.325353in}}%
\pgfpathlineto{\pgfqpoint{5.104831in}{3.322403in}}%
\pgfpathmoveto{\pgfqpoint{5.100290in}{3.325353in}}%
\pgfpathlineto{\pgfqpoint{5.100290in}{3.325353in}}%
\pgfpathlineto{\pgfqpoint{5.100290in}{3.328302in}}%
\pgfpathlineto{\pgfqpoint{5.104831in}{3.328302in}}%
\pgfpathlineto{\pgfqpoint{5.104831in}{3.325353in}}%
\pgfpathmoveto{\pgfqpoint{5.104831in}{3.325353in}}%
\pgfpathlineto{\pgfqpoint{5.104831in}{3.325353in}}%
\pgfpathlineto{\pgfqpoint{5.104831in}{3.328302in}}%
\pgfpathlineto{\pgfqpoint{5.109372in}{3.328302in}}%
\pgfpathlineto{\pgfqpoint{5.109372in}{3.325353in}}%
\pgfpathmoveto{\pgfqpoint{5.104831in}{3.328302in}}%
\pgfpathlineto{\pgfqpoint{5.104831in}{3.328302in}}%
\pgfpathlineto{\pgfqpoint{5.104831in}{3.331251in}}%
\pgfpathlineto{\pgfqpoint{5.109372in}{3.331251in}}%
\pgfpathlineto{\pgfqpoint{5.109372in}{3.328302in}}%
\pgfpathmoveto{\pgfqpoint{5.109372in}{3.328302in}}%
\pgfpathlineto{\pgfqpoint{5.109372in}{3.328302in}}%
\pgfpathlineto{\pgfqpoint{5.109372in}{3.331251in}}%
\pgfpathlineto{\pgfqpoint{5.113913in}{3.331251in}}%
\pgfpathlineto{\pgfqpoint{5.113913in}{3.328302in}}%
\pgfpathmoveto{\pgfqpoint{5.109372in}{3.331251in}}%
\pgfpathlineto{\pgfqpoint{5.109372in}{3.331251in}}%
\pgfpathlineto{\pgfqpoint{5.109372in}{3.334200in}}%
\pgfpathlineto{\pgfqpoint{5.113913in}{3.334200in}}%
\pgfpathlineto{\pgfqpoint{5.113913in}{3.331251in}}%
\pgfpathmoveto{\pgfqpoint{5.113913in}{3.331251in}}%
\pgfpathlineto{\pgfqpoint{5.113913in}{3.331251in}}%
\pgfpathlineto{\pgfqpoint{5.113913in}{3.334200in}}%
\pgfpathlineto{\pgfqpoint{5.118454in}{3.334200in}}%
\pgfpathlineto{\pgfqpoint{5.118454in}{3.331251in}}%
\pgfpathmoveto{\pgfqpoint{5.113913in}{3.334200in}}%
\pgfpathlineto{\pgfqpoint{5.113913in}{3.334200in}}%
\pgfpathlineto{\pgfqpoint{5.113913in}{3.337150in}}%
\pgfpathlineto{\pgfqpoint{5.118454in}{3.337150in}}%
\pgfpathlineto{\pgfqpoint{5.118454in}{3.334200in}}%
\pgfpathmoveto{\pgfqpoint{5.118454in}{3.334200in}}%
\pgfpathlineto{\pgfqpoint{5.118454in}{3.334200in}}%
\pgfpathlineto{\pgfqpoint{5.118454in}{3.337150in}}%
\pgfpathlineto{\pgfqpoint{5.122995in}{3.337150in}}%
\pgfpathlineto{\pgfqpoint{5.122995in}{3.334200in}}%
\pgfpathmoveto{\pgfqpoint{5.118454in}{3.337150in}}%
\pgfpathlineto{\pgfqpoint{5.118454in}{3.337150in}}%
\pgfpathlineto{\pgfqpoint{5.118454in}{3.340099in}}%
\pgfpathlineto{\pgfqpoint{5.122995in}{3.340099in}}%
\pgfpathlineto{\pgfqpoint{5.122995in}{3.337150in}}%
\pgfpathmoveto{\pgfqpoint{5.122995in}{3.337150in}}%
\pgfpathlineto{\pgfqpoint{5.122995in}{3.337150in}}%
\pgfpathlineto{\pgfqpoint{5.122995in}{3.340099in}}%
\pgfpathlineto{\pgfqpoint{5.127536in}{3.340099in}}%
\pgfpathlineto{\pgfqpoint{5.127536in}{3.337150in}}%
\pgfpathmoveto{\pgfqpoint{5.122995in}{3.340099in}}%
\pgfpathlineto{\pgfqpoint{5.122995in}{3.340099in}}%
\pgfpathlineto{\pgfqpoint{5.122995in}{3.343048in}}%
\pgfpathlineto{\pgfqpoint{5.127536in}{3.343048in}}%
\pgfpathlineto{\pgfqpoint{5.127536in}{3.340099in}}%
\pgfpathmoveto{\pgfqpoint{5.127536in}{3.340099in}}%
\pgfpathlineto{\pgfqpoint{5.127536in}{3.340099in}}%
\pgfpathlineto{\pgfqpoint{5.127536in}{3.343048in}}%
\pgfpathlineto{\pgfqpoint{5.132077in}{3.343048in}}%
\pgfpathlineto{\pgfqpoint{5.132077in}{3.340099in}}%
\pgfpathmoveto{\pgfqpoint{5.127536in}{3.343048in}}%
\pgfpathlineto{\pgfqpoint{5.127536in}{3.343048in}}%
\pgfpathlineto{\pgfqpoint{5.127536in}{3.345997in}}%
\pgfpathlineto{\pgfqpoint{5.132077in}{3.345997in}}%
\pgfpathlineto{\pgfqpoint{5.132077in}{3.343048in}}%
\pgfpathmoveto{\pgfqpoint{5.132077in}{3.343048in}}%
\pgfpathlineto{\pgfqpoint{5.132077in}{3.343048in}}%
\pgfpathlineto{\pgfqpoint{5.132077in}{3.345997in}}%
\pgfpathlineto{\pgfqpoint{5.136619in}{3.345997in}}%
\pgfpathlineto{\pgfqpoint{5.136619in}{3.343048in}}%
\pgfpathmoveto{\pgfqpoint{5.132077in}{3.345997in}}%
\pgfpathlineto{\pgfqpoint{5.132077in}{3.345997in}}%
\pgfpathlineto{\pgfqpoint{5.132077in}{3.348946in}}%
\pgfpathlineto{\pgfqpoint{5.136619in}{3.348946in}}%
\pgfpathlineto{\pgfqpoint{5.136619in}{3.345997in}}%
\pgfpathmoveto{\pgfqpoint{5.136619in}{3.345997in}}%
\pgfpathlineto{\pgfqpoint{5.136619in}{3.345997in}}%
\pgfpathlineto{\pgfqpoint{5.136619in}{3.348946in}}%
\pgfpathlineto{\pgfqpoint{5.141160in}{3.348946in}}%
\pgfpathlineto{\pgfqpoint{5.141160in}{3.345997in}}%
\pgfpathmoveto{\pgfqpoint{5.136619in}{3.348946in}}%
\pgfpathlineto{\pgfqpoint{5.136619in}{3.348946in}}%
\pgfpathlineto{\pgfqpoint{5.136619in}{3.351895in}}%
\pgfpathlineto{\pgfqpoint{5.141160in}{3.351895in}}%
\pgfpathlineto{\pgfqpoint{5.141160in}{3.348946in}}%
\pgfpathmoveto{\pgfqpoint{5.141160in}{3.348946in}}%
\pgfpathlineto{\pgfqpoint{5.141160in}{3.348946in}}%
\pgfpathlineto{\pgfqpoint{5.141160in}{3.351895in}}%
\pgfpathlineto{\pgfqpoint{5.145701in}{3.351895in}}%
\pgfpathlineto{\pgfqpoint{5.145701in}{3.348946in}}%
\pgfpathmoveto{\pgfqpoint{5.141160in}{3.351895in}}%
\pgfpathlineto{\pgfqpoint{5.141160in}{3.351895in}}%
\pgfpathlineto{\pgfqpoint{5.141160in}{3.354845in}}%
\pgfpathlineto{\pgfqpoint{5.145701in}{3.354845in}}%
\pgfpathlineto{\pgfqpoint{5.145701in}{3.351895in}}%
\pgfpathmoveto{\pgfqpoint{5.145701in}{3.351895in}}%
\pgfpathlineto{\pgfqpoint{5.145701in}{3.351895in}}%
\pgfpathlineto{\pgfqpoint{5.145701in}{3.354845in}}%
\pgfpathlineto{\pgfqpoint{5.150242in}{3.354845in}}%
\pgfpathlineto{\pgfqpoint{5.150242in}{3.351895in}}%
\pgfpathmoveto{\pgfqpoint{5.145701in}{3.354845in}}%
\pgfpathlineto{\pgfqpoint{5.145701in}{3.354845in}}%
\pgfpathlineto{\pgfqpoint{5.145701in}{3.357794in}}%
\pgfpathlineto{\pgfqpoint{5.150242in}{3.357794in}}%
\pgfpathlineto{\pgfqpoint{5.150242in}{3.354845in}}%
\pgfpathmoveto{\pgfqpoint{5.150242in}{3.354845in}}%
\pgfpathlineto{\pgfqpoint{5.150242in}{3.354845in}}%
\pgfpathlineto{\pgfqpoint{5.150242in}{3.357794in}}%
\pgfpathlineto{\pgfqpoint{5.154783in}{3.357794in}}%
\pgfpathlineto{\pgfqpoint{5.154783in}{3.354845in}}%
\pgfpathmoveto{\pgfqpoint{5.150242in}{3.357794in}}%
\pgfpathlineto{\pgfqpoint{5.150242in}{3.357794in}}%
\pgfpathlineto{\pgfqpoint{5.150242in}{3.360743in}}%
\pgfpathlineto{\pgfqpoint{5.154783in}{3.360743in}}%
\pgfpathlineto{\pgfqpoint{5.154783in}{3.357794in}}%
\pgfpathmoveto{\pgfqpoint{5.154783in}{3.357794in}}%
\pgfpathlineto{\pgfqpoint{5.154783in}{3.357794in}}%
\pgfpathlineto{\pgfqpoint{5.154783in}{3.360743in}}%
\pgfpathlineto{\pgfqpoint{5.159324in}{3.360743in}}%
\pgfpathlineto{\pgfqpoint{5.159324in}{3.357794in}}%
\pgfpathmoveto{\pgfqpoint{5.154783in}{3.360743in}}%
\pgfpathlineto{\pgfqpoint{5.154783in}{3.360743in}}%
\pgfpathlineto{\pgfqpoint{5.154783in}{3.363692in}}%
\pgfpathlineto{\pgfqpoint{5.159324in}{3.363692in}}%
\pgfpathlineto{\pgfqpoint{5.159324in}{3.360743in}}%
\pgfpathmoveto{\pgfqpoint{5.159324in}{3.360743in}}%
\pgfpathlineto{\pgfqpoint{5.159324in}{3.360743in}}%
\pgfpathlineto{\pgfqpoint{5.159324in}{3.363692in}}%
\pgfpathlineto{\pgfqpoint{5.163865in}{3.363692in}}%
\pgfpathlineto{\pgfqpoint{5.163865in}{3.360743in}}%
\pgfpathmoveto{\pgfqpoint{5.159324in}{3.363692in}}%
\pgfpathlineto{\pgfqpoint{5.159324in}{3.363692in}}%
\pgfpathlineto{\pgfqpoint{5.159324in}{3.366641in}}%
\pgfpathlineto{\pgfqpoint{5.163865in}{3.366641in}}%
\pgfpathlineto{\pgfqpoint{5.163865in}{3.363692in}}%
\pgfpathmoveto{\pgfqpoint{5.163865in}{3.363692in}}%
\pgfpathlineto{\pgfqpoint{5.163865in}{3.363692in}}%
\pgfpathlineto{\pgfqpoint{5.163865in}{3.366641in}}%
\pgfpathlineto{\pgfqpoint{5.168406in}{3.366641in}}%
\pgfpathlineto{\pgfqpoint{5.168406in}{3.363692in}}%
\pgfpathmoveto{\pgfqpoint{5.163865in}{3.366641in}}%
\pgfpathlineto{\pgfqpoint{5.163865in}{3.366641in}}%
\pgfpathlineto{\pgfqpoint{5.163865in}{3.369590in}}%
\pgfpathlineto{\pgfqpoint{5.168406in}{3.369590in}}%
\pgfpathlineto{\pgfqpoint{5.168406in}{3.366641in}}%
\pgfpathmoveto{\pgfqpoint{5.168406in}{3.366641in}}%
\pgfpathlineto{\pgfqpoint{5.168406in}{3.366641in}}%
\pgfpathlineto{\pgfqpoint{5.168406in}{3.369590in}}%
\pgfpathlineto{\pgfqpoint{5.172948in}{3.369590in}}%
\pgfpathlineto{\pgfqpoint{5.172948in}{3.366641in}}%
\pgfpathmoveto{\pgfqpoint{5.168406in}{3.369590in}}%
\pgfpathlineto{\pgfqpoint{5.168406in}{3.369590in}}%
\pgfpathlineto{\pgfqpoint{5.168406in}{3.372540in}}%
\pgfpathlineto{\pgfqpoint{5.172948in}{3.372540in}}%
\pgfpathlineto{\pgfqpoint{5.172948in}{3.369590in}}%
\pgfpathmoveto{\pgfqpoint{5.172948in}{3.369590in}}%
\pgfpathlineto{\pgfqpoint{5.172948in}{3.369590in}}%
\pgfpathlineto{\pgfqpoint{5.172948in}{3.372540in}}%
\pgfpathlineto{\pgfqpoint{5.177489in}{3.372540in}}%
\pgfpathlineto{\pgfqpoint{5.177489in}{3.369590in}}%
\pgfpathmoveto{\pgfqpoint{5.172948in}{3.372540in}}%
\pgfpathlineto{\pgfqpoint{5.172948in}{3.372540in}}%
\pgfpathlineto{\pgfqpoint{5.172948in}{3.375489in}}%
\pgfpathlineto{\pgfqpoint{5.177489in}{3.375489in}}%
\pgfpathlineto{\pgfqpoint{5.177489in}{3.372540in}}%
\pgfpathmoveto{\pgfqpoint{5.177489in}{3.372540in}}%
\pgfpathlineto{\pgfqpoint{5.177489in}{3.372540in}}%
\pgfpathlineto{\pgfqpoint{5.177489in}{3.375489in}}%
\pgfpathlineto{\pgfqpoint{5.182030in}{3.375489in}}%
\pgfpathlineto{\pgfqpoint{5.182030in}{3.372540in}}%
\pgfpathmoveto{\pgfqpoint{5.177489in}{3.375489in}}%
\pgfpathlineto{\pgfqpoint{5.177489in}{3.375489in}}%
\pgfpathlineto{\pgfqpoint{5.177489in}{3.378438in}}%
\pgfpathlineto{\pgfqpoint{5.182030in}{3.378438in}}%
\pgfpathlineto{\pgfqpoint{5.182030in}{3.375489in}}%
\pgfpathmoveto{\pgfqpoint{5.182030in}{3.375489in}}%
\pgfpathlineto{\pgfqpoint{5.182030in}{3.375489in}}%
\pgfpathlineto{\pgfqpoint{5.182030in}{3.378438in}}%
\pgfpathlineto{\pgfqpoint{5.186571in}{3.378438in}}%
\pgfpathlineto{\pgfqpoint{5.186571in}{3.375489in}}%
\pgfpathmoveto{\pgfqpoint{5.182030in}{3.378438in}}%
\pgfpathlineto{\pgfqpoint{5.182030in}{3.378438in}}%
\pgfpathlineto{\pgfqpoint{5.182030in}{3.381387in}}%
\pgfpathlineto{\pgfqpoint{5.186571in}{3.381387in}}%
\pgfpathlineto{\pgfqpoint{5.186571in}{3.378438in}}%
\pgfpathmoveto{\pgfqpoint{5.186571in}{3.378438in}}%
\pgfpathlineto{\pgfqpoint{5.186571in}{3.378438in}}%
\pgfpathlineto{\pgfqpoint{5.186571in}{3.381387in}}%
\pgfpathlineto{\pgfqpoint{5.191112in}{3.381387in}}%
\pgfpathlineto{\pgfqpoint{5.191112in}{3.378438in}}%
\pgfpathmoveto{\pgfqpoint{5.186571in}{3.381387in}}%
\pgfpathlineto{\pgfqpoint{5.186571in}{3.381387in}}%
\pgfpathlineto{\pgfqpoint{5.186571in}{3.384336in}}%
\pgfpathlineto{\pgfqpoint{5.191112in}{3.384336in}}%
\pgfpathlineto{\pgfqpoint{5.191112in}{3.381387in}}%
\pgfpathmoveto{\pgfqpoint{5.191112in}{3.381387in}}%
\pgfpathlineto{\pgfqpoint{5.191112in}{3.381387in}}%
\pgfpathlineto{\pgfqpoint{5.191112in}{3.384336in}}%
\pgfpathlineto{\pgfqpoint{5.195653in}{3.384336in}}%
\pgfpathlineto{\pgfqpoint{5.195653in}{3.381387in}}%
\pgfpathmoveto{\pgfqpoint{5.191112in}{3.384336in}}%
\pgfpathlineto{\pgfqpoint{5.191112in}{3.384336in}}%
\pgfpathlineto{\pgfqpoint{5.191112in}{3.387285in}}%
\pgfpathlineto{\pgfqpoint{5.195653in}{3.387285in}}%
\pgfpathlineto{\pgfqpoint{5.195653in}{3.384336in}}%
\pgfpathmoveto{\pgfqpoint{5.195653in}{3.384336in}}%
\pgfpathlineto{\pgfqpoint{5.195653in}{3.384336in}}%
\pgfpathlineto{\pgfqpoint{5.195653in}{3.387285in}}%
\pgfpathlineto{\pgfqpoint{5.200194in}{3.387285in}}%
\pgfpathlineto{\pgfqpoint{5.200194in}{3.384336in}}%
\pgfpathmoveto{\pgfqpoint{5.195653in}{3.387285in}}%
\pgfpathlineto{\pgfqpoint{5.195653in}{3.387285in}}%
\pgfpathlineto{\pgfqpoint{5.195653in}{3.390235in}}%
\pgfpathlineto{\pgfqpoint{5.200194in}{3.390235in}}%
\pgfpathlineto{\pgfqpoint{5.200194in}{3.387285in}}%
\pgfpathmoveto{\pgfqpoint{5.200194in}{3.387285in}}%
\pgfpathlineto{\pgfqpoint{5.200194in}{3.387285in}}%
\pgfpathlineto{\pgfqpoint{5.200194in}{3.390235in}}%
\pgfpathlineto{\pgfqpoint{5.204735in}{3.390235in}}%
\pgfpathlineto{\pgfqpoint{5.204735in}{3.387285in}}%
\pgfpathmoveto{\pgfqpoint{5.200194in}{3.390235in}}%
\pgfpathlineto{\pgfqpoint{5.200194in}{3.390235in}}%
\pgfpathlineto{\pgfqpoint{5.200194in}{3.393184in}}%
\pgfpathlineto{\pgfqpoint{5.204735in}{3.393184in}}%
\pgfpathlineto{\pgfqpoint{5.204735in}{3.390235in}}%
\pgfpathmoveto{\pgfqpoint{5.204735in}{3.390235in}}%
\pgfpathlineto{\pgfqpoint{5.204735in}{3.390235in}}%
\pgfpathlineto{\pgfqpoint{5.204735in}{3.393184in}}%
\pgfpathlineto{\pgfqpoint{5.209277in}{3.393184in}}%
\pgfpathlineto{\pgfqpoint{5.209277in}{3.390235in}}%
\pgfpathmoveto{\pgfqpoint{5.204735in}{3.393184in}}%
\pgfpathlineto{\pgfqpoint{5.204735in}{3.393184in}}%
\pgfpathlineto{\pgfqpoint{5.204735in}{3.396133in}}%
\pgfpathlineto{\pgfqpoint{5.209277in}{3.396133in}}%
\pgfpathlineto{\pgfqpoint{5.209277in}{3.393184in}}%
\pgfpathmoveto{\pgfqpoint{5.209277in}{3.393184in}}%
\pgfpathlineto{\pgfqpoint{5.209277in}{3.393184in}}%
\pgfpathlineto{\pgfqpoint{5.209277in}{3.396133in}}%
\pgfpathlineto{\pgfqpoint{5.213818in}{3.396133in}}%
\pgfpathlineto{\pgfqpoint{5.213818in}{3.393184in}}%
\pgfpathmoveto{\pgfqpoint{5.209277in}{3.396133in}}%
\pgfpathlineto{\pgfqpoint{5.209277in}{3.396133in}}%
\pgfpathlineto{\pgfqpoint{5.209277in}{3.399082in}}%
\pgfpathlineto{\pgfqpoint{5.213818in}{3.399082in}}%
\pgfpathlineto{\pgfqpoint{5.213818in}{3.396133in}}%
\pgfpathmoveto{\pgfqpoint{5.213818in}{3.396133in}}%
\pgfpathlineto{\pgfqpoint{5.213818in}{3.396133in}}%
\pgfpathlineto{\pgfqpoint{5.213818in}{3.399082in}}%
\pgfpathlineto{\pgfqpoint{5.218359in}{3.399082in}}%
\pgfpathlineto{\pgfqpoint{5.218359in}{3.396133in}}%
\pgfpathmoveto{\pgfqpoint{5.213818in}{3.399082in}}%
\pgfpathlineto{\pgfqpoint{5.213818in}{3.399082in}}%
\pgfpathlineto{\pgfqpoint{5.213818in}{3.402031in}}%
\pgfpathlineto{\pgfqpoint{5.218359in}{3.402031in}}%
\pgfpathlineto{\pgfqpoint{5.218359in}{3.399082in}}%
\pgfpathmoveto{\pgfqpoint{5.218359in}{3.399082in}}%
\pgfpathlineto{\pgfqpoint{5.218359in}{3.399082in}}%
\pgfpathlineto{\pgfqpoint{5.218359in}{3.402031in}}%
\pgfpathlineto{\pgfqpoint{5.222900in}{3.402031in}}%
\pgfpathlineto{\pgfqpoint{5.222900in}{3.399082in}}%
\pgfpathmoveto{\pgfqpoint{5.218359in}{3.402031in}}%
\pgfpathlineto{\pgfqpoint{5.218359in}{3.402031in}}%
\pgfpathlineto{\pgfqpoint{5.218359in}{3.404980in}}%
\pgfpathlineto{\pgfqpoint{5.222900in}{3.404980in}}%
\pgfpathlineto{\pgfqpoint{5.222900in}{3.402031in}}%
\pgfpathmoveto{\pgfqpoint{5.222900in}{3.402031in}}%
\pgfpathlineto{\pgfqpoint{5.222900in}{3.402031in}}%
\pgfpathlineto{\pgfqpoint{5.222900in}{3.404980in}}%
\pgfpathlineto{\pgfqpoint{5.227441in}{3.404980in}}%
\pgfpathlineto{\pgfqpoint{5.227441in}{3.402031in}}%
\pgfpathmoveto{\pgfqpoint{5.222900in}{3.404980in}}%
\pgfpathlineto{\pgfqpoint{5.222900in}{3.404980in}}%
\pgfpathlineto{\pgfqpoint{5.222900in}{3.407929in}}%
\pgfpathlineto{\pgfqpoint{5.227441in}{3.407929in}}%
\pgfpathlineto{\pgfqpoint{5.227441in}{3.404980in}}%
\pgfpathmoveto{\pgfqpoint{5.227441in}{3.404980in}}%
\pgfpathlineto{\pgfqpoint{5.227441in}{3.404980in}}%
\pgfpathlineto{\pgfqpoint{5.227441in}{3.407929in}}%
\pgfpathlineto{\pgfqpoint{5.231982in}{3.407929in}}%
\pgfpathlineto{\pgfqpoint{5.231982in}{3.404980in}}%
\pgfpathmoveto{\pgfqpoint{5.227441in}{3.407929in}}%
\pgfpathlineto{\pgfqpoint{5.227441in}{3.407929in}}%
\pgfpathlineto{\pgfqpoint{5.227441in}{3.410879in}}%
\pgfpathlineto{\pgfqpoint{5.231982in}{3.410879in}}%
\pgfpathlineto{\pgfqpoint{5.231982in}{3.407929in}}%
\pgfpathmoveto{\pgfqpoint{5.227441in}{3.410879in}}%
\pgfpathlineto{\pgfqpoint{5.227441in}{3.410879in}}%
\pgfpathlineto{\pgfqpoint{5.227441in}{3.413828in}}%
\pgfpathlineto{\pgfqpoint{5.231982in}{3.413828in}}%
\pgfpathlineto{\pgfqpoint{5.231982in}{3.410879in}}%
\pgfpathmoveto{\pgfqpoint{5.231982in}{3.410879in}}%
\pgfpathlineto{\pgfqpoint{5.231982in}{3.410879in}}%
\pgfpathlineto{\pgfqpoint{5.231982in}{3.413828in}}%
\pgfpathlineto{\pgfqpoint{5.236523in}{3.413828in}}%
\pgfpathlineto{\pgfqpoint{5.236523in}{3.410879in}}%
\pgfpathmoveto{\pgfqpoint{5.231982in}{3.413828in}}%
\pgfpathlineto{\pgfqpoint{5.231982in}{3.413828in}}%
\pgfpathlineto{\pgfqpoint{5.231982in}{3.416777in}}%
\pgfpathlineto{\pgfqpoint{5.236523in}{3.416777in}}%
\pgfpathlineto{\pgfqpoint{5.236523in}{3.413828in}}%
\pgfpathmoveto{\pgfqpoint{5.236523in}{3.413828in}}%
\pgfpathlineto{\pgfqpoint{5.236523in}{3.413828in}}%
\pgfpathlineto{\pgfqpoint{5.236523in}{3.416777in}}%
\pgfpathlineto{\pgfqpoint{5.241064in}{3.416777in}}%
\pgfpathlineto{\pgfqpoint{5.241064in}{3.413828in}}%
\pgfpathmoveto{\pgfqpoint{5.236523in}{3.416777in}}%
\pgfpathlineto{\pgfqpoint{5.236523in}{3.416777in}}%
\pgfpathlineto{\pgfqpoint{5.236523in}{3.419726in}}%
\pgfpathlineto{\pgfqpoint{5.241064in}{3.419726in}}%
\pgfpathlineto{\pgfqpoint{5.241064in}{3.416777in}}%
\pgfpathmoveto{\pgfqpoint{5.241064in}{3.416777in}}%
\pgfpathlineto{\pgfqpoint{5.241064in}{3.416777in}}%
\pgfpathlineto{\pgfqpoint{5.241064in}{3.419726in}}%
\pgfpathlineto{\pgfqpoint{5.245605in}{3.419726in}}%
\pgfpathlineto{\pgfqpoint{5.245605in}{3.416777in}}%
\pgfpathmoveto{\pgfqpoint{5.241064in}{3.419726in}}%
\pgfpathlineto{\pgfqpoint{5.241064in}{3.419726in}}%
\pgfpathlineto{\pgfqpoint{5.241064in}{3.422675in}}%
\pgfpathlineto{\pgfqpoint{5.245605in}{3.422675in}}%
\pgfpathlineto{\pgfqpoint{5.245605in}{3.419726in}}%
\pgfpathmoveto{\pgfqpoint{5.245605in}{3.419726in}}%
\pgfpathlineto{\pgfqpoint{5.245605in}{3.419726in}}%
\pgfpathlineto{\pgfqpoint{5.245605in}{3.422675in}}%
\pgfpathlineto{\pgfqpoint{5.250147in}{3.422675in}}%
\pgfpathlineto{\pgfqpoint{5.250147in}{3.419726in}}%
\pgfpathmoveto{\pgfqpoint{5.245605in}{3.422675in}}%
\pgfpathlineto{\pgfqpoint{5.245605in}{3.422675in}}%
\pgfpathlineto{\pgfqpoint{5.245605in}{3.425624in}}%
\pgfpathlineto{\pgfqpoint{5.250147in}{3.425624in}}%
\pgfpathlineto{\pgfqpoint{5.250147in}{3.422675in}}%
\pgfpathmoveto{\pgfqpoint{5.250147in}{3.422675in}}%
\pgfpathlineto{\pgfqpoint{5.250147in}{3.422675in}}%
\pgfpathlineto{\pgfqpoint{5.250147in}{3.425624in}}%
\pgfpathlineto{\pgfqpoint{5.254688in}{3.425624in}}%
\pgfpathlineto{\pgfqpoint{5.254688in}{3.422675in}}%
\pgfpathmoveto{\pgfqpoint{5.250147in}{3.425624in}}%
\pgfpathlineto{\pgfqpoint{5.250147in}{3.425624in}}%
\pgfpathlineto{\pgfqpoint{5.250147in}{3.428574in}}%
\pgfpathlineto{\pgfqpoint{5.254688in}{3.428574in}}%
\pgfpathlineto{\pgfqpoint{5.254688in}{3.425624in}}%
\pgfpathmoveto{\pgfqpoint{5.254688in}{3.425624in}}%
\pgfpathlineto{\pgfqpoint{5.254688in}{3.425624in}}%
\pgfpathlineto{\pgfqpoint{5.254688in}{3.428574in}}%
\pgfpathlineto{\pgfqpoint{5.259229in}{3.428574in}}%
\pgfpathlineto{\pgfqpoint{5.259229in}{3.425624in}}%
\pgfpathmoveto{\pgfqpoint{5.254688in}{3.428574in}}%
\pgfpathlineto{\pgfqpoint{5.254688in}{3.428574in}}%
\pgfpathlineto{\pgfqpoint{5.254688in}{3.431523in}}%
\pgfpathlineto{\pgfqpoint{5.259229in}{3.431523in}}%
\pgfpathlineto{\pgfqpoint{5.259229in}{3.428574in}}%
\pgfpathmoveto{\pgfqpoint{5.259229in}{3.428574in}}%
\pgfpathlineto{\pgfqpoint{5.259229in}{3.428574in}}%
\pgfpathlineto{\pgfqpoint{5.259229in}{3.431523in}}%
\pgfpathlineto{\pgfqpoint{5.263770in}{3.431523in}}%
\pgfpathlineto{\pgfqpoint{5.263770in}{3.428574in}}%
\pgfpathmoveto{\pgfqpoint{5.259229in}{3.431523in}}%
\pgfpathlineto{\pgfqpoint{5.259229in}{3.431523in}}%
\pgfpathlineto{\pgfqpoint{5.259229in}{3.434472in}}%
\pgfpathlineto{\pgfqpoint{5.263770in}{3.434472in}}%
\pgfpathlineto{\pgfqpoint{5.263770in}{3.431523in}}%
\pgfpathmoveto{\pgfqpoint{5.263770in}{3.431523in}}%
\pgfpathlineto{\pgfqpoint{5.263770in}{3.431523in}}%
\pgfpathlineto{\pgfqpoint{5.263770in}{3.434472in}}%
\pgfpathlineto{\pgfqpoint{5.268311in}{3.434472in}}%
\pgfpathlineto{\pgfqpoint{5.268311in}{3.431523in}}%
\pgfpathmoveto{\pgfqpoint{5.263770in}{3.434472in}}%
\pgfpathlineto{\pgfqpoint{5.263770in}{3.434472in}}%
\pgfpathlineto{\pgfqpoint{5.263770in}{3.437421in}}%
\pgfpathlineto{\pgfqpoint{5.268311in}{3.437421in}}%
\pgfpathlineto{\pgfqpoint{5.268311in}{3.434472in}}%
\pgfpathmoveto{\pgfqpoint{5.268311in}{3.434472in}}%
\pgfpathlineto{\pgfqpoint{5.268311in}{3.434472in}}%
\pgfpathlineto{\pgfqpoint{5.268311in}{3.437421in}}%
\pgfpathlineto{\pgfqpoint{5.272852in}{3.437421in}}%
\pgfpathlineto{\pgfqpoint{5.272852in}{3.434472in}}%
\pgfpathmoveto{\pgfqpoint{5.268311in}{3.437421in}}%
\pgfpathlineto{\pgfqpoint{5.268311in}{3.437421in}}%
\pgfpathlineto{\pgfqpoint{5.268311in}{3.440370in}}%
\pgfpathlineto{\pgfqpoint{5.272852in}{3.440370in}}%
\pgfpathlineto{\pgfqpoint{5.272852in}{3.437421in}}%
\pgfpathmoveto{\pgfqpoint{5.272852in}{3.437421in}}%
\pgfpathlineto{\pgfqpoint{5.272852in}{3.437421in}}%
\pgfpathlineto{\pgfqpoint{5.272852in}{3.440370in}}%
\pgfpathlineto{\pgfqpoint{5.277393in}{3.440370in}}%
\pgfpathlineto{\pgfqpoint{5.277393in}{3.437421in}}%
\pgfpathmoveto{\pgfqpoint{5.272852in}{3.440370in}}%
\pgfpathlineto{\pgfqpoint{5.272852in}{3.440370in}}%
\pgfpathlineto{\pgfqpoint{5.272852in}{3.443319in}}%
\pgfpathlineto{\pgfqpoint{5.277393in}{3.443319in}}%
\pgfpathlineto{\pgfqpoint{5.277393in}{3.440370in}}%
\pgfpathmoveto{\pgfqpoint{5.277393in}{3.440370in}}%
\pgfpathlineto{\pgfqpoint{5.277393in}{3.440370in}}%
\pgfpathlineto{\pgfqpoint{5.277393in}{3.443319in}}%
\pgfpathlineto{\pgfqpoint{5.281934in}{3.443319in}}%
\pgfpathlineto{\pgfqpoint{5.281934in}{3.440370in}}%
\pgfpathmoveto{\pgfqpoint{5.277393in}{3.443319in}}%
\pgfpathlineto{\pgfqpoint{5.277393in}{3.443319in}}%
\pgfpathlineto{\pgfqpoint{5.277393in}{3.446269in}}%
\pgfpathlineto{\pgfqpoint{5.281934in}{3.446269in}}%
\pgfpathlineto{\pgfqpoint{5.281934in}{3.443319in}}%
\pgfpathmoveto{\pgfqpoint{5.281934in}{3.443319in}}%
\pgfpathlineto{\pgfqpoint{5.281934in}{3.443319in}}%
\pgfpathlineto{\pgfqpoint{5.281934in}{3.446269in}}%
\pgfpathlineto{\pgfqpoint{5.286475in}{3.446269in}}%
\pgfpathlineto{\pgfqpoint{5.286475in}{3.443319in}}%
\pgfpathmoveto{\pgfqpoint{5.281934in}{3.446269in}}%
\pgfpathlineto{\pgfqpoint{5.281934in}{3.446269in}}%
\pgfpathlineto{\pgfqpoint{5.281934in}{3.449218in}}%
\pgfpathlineto{\pgfqpoint{5.286475in}{3.449218in}}%
\pgfpathlineto{\pgfqpoint{5.286475in}{3.446269in}}%
\pgfpathmoveto{\pgfqpoint{5.286475in}{3.446269in}}%
\pgfpathlineto{\pgfqpoint{5.286475in}{3.446269in}}%
\pgfpathlineto{\pgfqpoint{5.286475in}{3.449218in}}%
\pgfpathlineto{\pgfqpoint{5.291016in}{3.449218in}}%
\pgfpathlineto{\pgfqpoint{5.291016in}{3.446269in}}%
\pgfpathmoveto{\pgfqpoint{5.286475in}{3.449218in}}%
\pgfpathlineto{\pgfqpoint{5.286475in}{3.449218in}}%
\pgfpathlineto{\pgfqpoint{5.286475in}{3.452167in}}%
\pgfpathlineto{\pgfqpoint{5.291016in}{3.452167in}}%
\pgfpathlineto{\pgfqpoint{5.291016in}{3.449218in}}%
\pgfpathmoveto{\pgfqpoint{5.291016in}{3.449218in}}%
\pgfpathlineto{\pgfqpoint{5.291016in}{3.449218in}}%
\pgfpathlineto{\pgfqpoint{5.291016in}{3.452167in}}%
\pgfpathlineto{\pgfqpoint{5.295557in}{3.452167in}}%
\pgfpathlineto{\pgfqpoint{5.295557in}{3.449218in}}%
\pgfpathmoveto{\pgfqpoint{5.291016in}{3.452167in}}%
\pgfpathlineto{\pgfqpoint{5.291016in}{3.452167in}}%
\pgfpathlineto{\pgfqpoint{5.291016in}{3.455116in}}%
\pgfpathlineto{\pgfqpoint{5.295557in}{3.455116in}}%
\pgfpathlineto{\pgfqpoint{5.295557in}{3.452167in}}%
\pgfpathmoveto{\pgfqpoint{5.295557in}{3.452167in}}%
\pgfpathlineto{\pgfqpoint{5.295557in}{3.452167in}}%
\pgfpathlineto{\pgfqpoint{5.295557in}{3.455116in}}%
\pgfpathlineto{\pgfqpoint{5.300098in}{3.455116in}}%
\pgfpathlineto{\pgfqpoint{5.300098in}{3.452167in}}%
\pgfpathmoveto{\pgfqpoint{5.295557in}{3.455116in}}%
\pgfpathlineto{\pgfqpoint{5.295557in}{3.455116in}}%
\pgfpathlineto{\pgfqpoint{5.295557in}{3.458065in}}%
\pgfpathlineto{\pgfqpoint{5.300098in}{3.458065in}}%
\pgfpathlineto{\pgfqpoint{5.300098in}{3.455116in}}%
\pgfpathmoveto{\pgfqpoint{5.300098in}{3.455116in}}%
\pgfpathlineto{\pgfqpoint{5.300098in}{3.455116in}}%
\pgfpathlineto{\pgfqpoint{5.300098in}{3.458065in}}%
\pgfpathlineto{\pgfqpoint{5.304639in}{3.458065in}}%
\pgfpathlineto{\pgfqpoint{5.304639in}{3.455116in}}%
\pgfpathmoveto{\pgfqpoint{5.300098in}{3.458065in}}%
\pgfpathlineto{\pgfqpoint{5.300098in}{3.458065in}}%
\pgfpathlineto{\pgfqpoint{5.300098in}{3.461014in}}%
\pgfpathlineto{\pgfqpoint{5.304639in}{3.461014in}}%
\pgfpathlineto{\pgfqpoint{5.304639in}{3.458065in}}%
\pgfpathmoveto{\pgfqpoint{5.304639in}{3.458065in}}%
\pgfpathlineto{\pgfqpoint{5.304639in}{3.458065in}}%
\pgfpathlineto{\pgfqpoint{5.304639in}{3.461014in}}%
\pgfpathlineto{\pgfqpoint{5.309180in}{3.461014in}}%
\pgfpathlineto{\pgfqpoint{5.309180in}{3.458065in}}%
\pgfpathmoveto{\pgfqpoint{5.304639in}{3.461014in}}%
\pgfpathlineto{\pgfqpoint{5.304639in}{3.461014in}}%
\pgfpathlineto{\pgfqpoint{5.304639in}{3.463964in}}%
\pgfpathlineto{\pgfqpoint{5.309180in}{3.463964in}}%
\pgfpathlineto{\pgfqpoint{5.309180in}{3.461014in}}%
\pgfpathmoveto{\pgfqpoint{5.309180in}{3.461014in}}%
\pgfpathlineto{\pgfqpoint{5.309180in}{3.461014in}}%
\pgfpathlineto{\pgfqpoint{5.309180in}{3.463964in}}%
\pgfpathlineto{\pgfqpoint{5.313721in}{3.463964in}}%
\pgfpathlineto{\pgfqpoint{5.313721in}{3.461014in}}%
\pgfpathmoveto{\pgfqpoint{5.309180in}{3.463964in}}%
\pgfpathlineto{\pgfqpoint{5.309180in}{3.463964in}}%
\pgfpathlineto{\pgfqpoint{5.309180in}{3.466913in}}%
\pgfpathlineto{\pgfqpoint{5.313721in}{3.466913in}}%
\pgfpathlineto{\pgfqpoint{5.313721in}{3.463964in}}%
\pgfpathmoveto{\pgfqpoint{5.313721in}{3.463964in}}%
\pgfpathlineto{\pgfqpoint{5.313721in}{3.463964in}}%
\pgfpathlineto{\pgfqpoint{5.313721in}{3.466913in}}%
\pgfpathlineto{\pgfqpoint{5.318262in}{3.466913in}}%
\pgfpathlineto{\pgfqpoint{5.318262in}{3.463964in}}%
\pgfpathmoveto{\pgfqpoint{5.313721in}{3.466913in}}%
\pgfpathlineto{\pgfqpoint{5.313721in}{3.466913in}}%
\pgfpathlineto{\pgfqpoint{5.313721in}{3.469862in}}%
\pgfpathlineto{\pgfqpoint{5.318262in}{3.469862in}}%
\pgfpathlineto{\pgfqpoint{5.318262in}{3.466913in}}%
\pgfpathmoveto{\pgfqpoint{5.318262in}{3.466913in}}%
\pgfpathlineto{\pgfqpoint{5.318262in}{3.466913in}}%
\pgfpathlineto{\pgfqpoint{5.318262in}{3.469862in}}%
\pgfpathlineto{\pgfqpoint{5.322803in}{3.469862in}}%
\pgfpathlineto{\pgfqpoint{5.322803in}{3.466913in}}%
\pgfpathmoveto{\pgfqpoint{5.318262in}{3.469862in}}%
\pgfpathlineto{\pgfqpoint{5.318262in}{3.469862in}}%
\pgfpathlineto{\pgfqpoint{5.318262in}{3.472811in}}%
\pgfpathlineto{\pgfqpoint{5.322803in}{3.472811in}}%
\pgfpathlineto{\pgfqpoint{5.322803in}{3.469862in}}%
\pgfpathmoveto{\pgfqpoint{5.322803in}{3.469862in}}%
\pgfpathlineto{\pgfqpoint{5.322803in}{3.469862in}}%
\pgfpathlineto{\pgfqpoint{5.322803in}{3.472811in}}%
\pgfpathlineto{\pgfqpoint{5.327344in}{3.472811in}}%
\pgfpathlineto{\pgfqpoint{5.327344in}{3.469862in}}%
\pgfpathmoveto{\pgfqpoint{5.322803in}{3.472811in}}%
\pgfpathlineto{\pgfqpoint{5.322803in}{3.472811in}}%
\pgfpathlineto{\pgfqpoint{5.322803in}{3.475760in}}%
\pgfpathlineto{\pgfqpoint{5.327344in}{3.475760in}}%
\pgfpathlineto{\pgfqpoint{5.327344in}{3.472811in}}%
\pgfpathmoveto{\pgfqpoint{5.327344in}{3.472811in}}%
\pgfpathlineto{\pgfqpoint{5.327344in}{3.472811in}}%
\pgfpathlineto{\pgfqpoint{5.327344in}{3.475760in}}%
\pgfpathlineto{\pgfqpoint{5.331885in}{3.475760in}}%
\pgfpathlineto{\pgfqpoint{5.331885in}{3.472811in}}%
\pgfpathmoveto{\pgfqpoint{5.327344in}{3.475760in}}%
\pgfpathlineto{\pgfqpoint{5.327344in}{3.475760in}}%
\pgfpathlineto{\pgfqpoint{5.327344in}{3.478709in}}%
\pgfpathlineto{\pgfqpoint{5.331885in}{3.478709in}}%
\pgfpathlineto{\pgfqpoint{5.331885in}{3.475760in}}%
\pgfpathmoveto{\pgfqpoint{5.331885in}{3.475760in}}%
\pgfpathlineto{\pgfqpoint{5.331885in}{3.475760in}}%
\pgfpathlineto{\pgfqpoint{5.331885in}{3.478709in}}%
\pgfpathlineto{\pgfqpoint{5.336426in}{3.478709in}}%
\pgfpathlineto{\pgfqpoint{5.336426in}{3.475760in}}%
\pgfpathmoveto{\pgfqpoint{5.331885in}{3.478709in}}%
\pgfpathlineto{\pgfqpoint{5.331885in}{3.478709in}}%
\pgfpathlineto{\pgfqpoint{5.331885in}{3.481659in}}%
\pgfpathlineto{\pgfqpoint{5.336426in}{3.481659in}}%
\pgfpathlineto{\pgfqpoint{5.336426in}{3.478709in}}%
\pgfpathmoveto{\pgfqpoint{5.336426in}{3.478709in}}%
\pgfpathlineto{\pgfqpoint{5.336426in}{3.478709in}}%
\pgfpathlineto{\pgfqpoint{5.336426in}{3.481659in}}%
\pgfpathlineto{\pgfqpoint{5.340967in}{3.481659in}}%
\pgfpathlineto{\pgfqpoint{5.340967in}{3.478709in}}%
\pgfpathmoveto{\pgfqpoint{5.336426in}{3.481659in}}%
\pgfpathlineto{\pgfqpoint{5.336426in}{3.481659in}}%
\pgfpathlineto{\pgfqpoint{5.336426in}{3.484608in}}%
\pgfpathlineto{\pgfqpoint{5.340967in}{3.484608in}}%
\pgfpathlineto{\pgfqpoint{5.340967in}{3.481659in}}%
\pgfpathmoveto{\pgfqpoint{5.340967in}{3.481659in}}%
\pgfpathlineto{\pgfqpoint{5.340967in}{3.481659in}}%
\pgfpathlineto{\pgfqpoint{5.340967in}{3.484608in}}%
\pgfpathlineto{\pgfqpoint{5.345508in}{3.484608in}}%
\pgfpathlineto{\pgfqpoint{5.345508in}{3.481659in}}%
\pgfpathmoveto{\pgfqpoint{5.340967in}{3.484608in}}%
\pgfpathlineto{\pgfqpoint{5.340967in}{3.484608in}}%
\pgfpathlineto{\pgfqpoint{5.340967in}{3.487557in}}%
\pgfpathlineto{\pgfqpoint{5.345508in}{3.487557in}}%
\pgfpathlineto{\pgfqpoint{5.345508in}{3.484608in}}%
\pgfpathmoveto{\pgfqpoint{5.345508in}{3.484608in}}%
\pgfpathlineto{\pgfqpoint{5.345508in}{3.484608in}}%
\pgfpathlineto{\pgfqpoint{5.345508in}{3.487557in}}%
\pgfpathlineto{\pgfqpoint{5.350049in}{3.487557in}}%
\pgfpathlineto{\pgfqpoint{5.350049in}{3.484608in}}%
\pgfpathmoveto{\pgfqpoint{5.345508in}{3.487557in}}%
\pgfpathlineto{\pgfqpoint{5.345508in}{3.487557in}}%
\pgfpathlineto{\pgfqpoint{5.345508in}{3.490506in}}%
\pgfpathlineto{\pgfqpoint{5.350049in}{3.490506in}}%
\pgfpathlineto{\pgfqpoint{5.350049in}{3.487557in}}%
\pgfpathmoveto{\pgfqpoint{5.350049in}{3.487557in}}%
\pgfpathlineto{\pgfqpoint{5.350049in}{3.487557in}}%
\pgfpathlineto{\pgfqpoint{5.350049in}{3.490506in}}%
\pgfpathlineto{\pgfqpoint{5.354590in}{3.490506in}}%
\pgfpathlineto{\pgfqpoint{5.354590in}{3.487557in}}%
\pgfpathmoveto{\pgfqpoint{5.350049in}{3.490506in}}%
\pgfpathlineto{\pgfqpoint{5.350049in}{3.490506in}}%
\pgfpathlineto{\pgfqpoint{5.350049in}{3.493455in}}%
\pgfpathlineto{\pgfqpoint{5.354590in}{3.493455in}}%
\pgfpathlineto{\pgfqpoint{5.354590in}{3.490506in}}%
\pgfpathmoveto{\pgfqpoint{5.354590in}{3.490506in}}%
\pgfpathlineto{\pgfqpoint{5.354590in}{3.490506in}}%
\pgfpathlineto{\pgfqpoint{5.354590in}{3.493455in}}%
\pgfpathlineto{\pgfqpoint{5.359131in}{3.493455in}}%
\pgfpathlineto{\pgfqpoint{5.359131in}{3.490506in}}%
\pgfpathmoveto{\pgfqpoint{5.354590in}{3.493455in}}%
\pgfpathlineto{\pgfqpoint{5.354590in}{3.493455in}}%
\pgfpathlineto{\pgfqpoint{5.354590in}{3.496404in}}%
\pgfpathlineto{\pgfqpoint{5.359131in}{3.496404in}}%
\pgfpathlineto{\pgfqpoint{5.359131in}{3.493455in}}%
\pgfpathmoveto{\pgfqpoint{5.359131in}{3.493455in}}%
\pgfpathlineto{\pgfqpoint{5.359131in}{3.493455in}}%
\pgfpathlineto{\pgfqpoint{5.359131in}{3.496404in}}%
\pgfpathlineto{\pgfqpoint{5.363673in}{3.496404in}}%
\pgfpathlineto{\pgfqpoint{5.363673in}{3.493455in}}%
\pgfpathmoveto{\pgfqpoint{5.359131in}{3.496404in}}%
\pgfpathlineto{\pgfqpoint{5.359131in}{3.496404in}}%
\pgfpathlineto{\pgfqpoint{5.359131in}{3.499353in}}%
\pgfpathlineto{\pgfqpoint{5.363673in}{3.499353in}}%
\pgfpathlineto{\pgfqpoint{5.363673in}{3.496404in}}%
\pgfpathmoveto{\pgfqpoint{5.363673in}{3.496404in}}%
\pgfpathlineto{\pgfqpoint{5.363673in}{3.496404in}}%
\pgfpathlineto{\pgfqpoint{5.363673in}{3.499353in}}%
\pgfpathlineto{\pgfqpoint{5.368214in}{3.499353in}}%
\pgfpathlineto{\pgfqpoint{5.368214in}{3.496404in}}%
\pgfpathmoveto{\pgfqpoint{5.363673in}{3.499353in}}%
\pgfpathlineto{\pgfqpoint{5.363673in}{3.499353in}}%
\pgfpathlineto{\pgfqpoint{5.363673in}{3.502303in}}%
\pgfpathlineto{\pgfqpoint{5.368214in}{3.502303in}}%
\pgfpathlineto{\pgfqpoint{5.368214in}{3.499353in}}%
\pgfpathmoveto{\pgfqpoint{5.368214in}{3.499353in}}%
\pgfpathlineto{\pgfqpoint{5.368214in}{3.499353in}}%
\pgfpathlineto{\pgfqpoint{5.368214in}{3.502303in}}%
\pgfpathlineto{\pgfqpoint{5.372755in}{3.502303in}}%
\pgfpathlineto{\pgfqpoint{5.372755in}{3.499353in}}%
\pgfpathmoveto{\pgfqpoint{5.368214in}{3.502303in}}%
\pgfpathlineto{\pgfqpoint{5.368214in}{3.502303in}}%
\pgfpathlineto{\pgfqpoint{5.368214in}{3.505252in}}%
\pgfpathlineto{\pgfqpoint{5.372755in}{3.505252in}}%
\pgfpathlineto{\pgfqpoint{5.372755in}{3.502303in}}%
\pgfpathmoveto{\pgfqpoint{5.372755in}{3.502303in}}%
\pgfpathlineto{\pgfqpoint{5.372755in}{3.502303in}}%
\pgfpathlineto{\pgfqpoint{5.372755in}{3.505252in}}%
\pgfpathlineto{\pgfqpoint{5.377296in}{3.505252in}}%
\pgfpathlineto{\pgfqpoint{5.377296in}{3.502303in}}%
\pgfpathmoveto{\pgfqpoint{5.372755in}{3.505252in}}%
\pgfpathlineto{\pgfqpoint{5.372755in}{3.505252in}}%
\pgfpathlineto{\pgfqpoint{5.372755in}{3.508201in}}%
\pgfpathlineto{\pgfqpoint{5.377296in}{3.508201in}}%
\pgfpathlineto{\pgfqpoint{5.377296in}{3.505252in}}%
\pgfpathmoveto{\pgfqpoint{5.377296in}{3.505252in}}%
\pgfpathlineto{\pgfqpoint{5.377296in}{3.505252in}}%
\pgfpathlineto{\pgfqpoint{5.377296in}{3.508201in}}%
\pgfpathlineto{\pgfqpoint{5.381837in}{3.508201in}}%
\pgfpathlineto{\pgfqpoint{5.381837in}{3.505252in}}%
\pgfpathmoveto{\pgfqpoint{5.377296in}{3.508201in}}%
\pgfpathlineto{\pgfqpoint{5.377296in}{3.508201in}}%
\pgfpathlineto{\pgfqpoint{5.377296in}{3.511150in}}%
\pgfpathlineto{\pgfqpoint{5.381837in}{3.511150in}}%
\pgfpathlineto{\pgfqpoint{5.381837in}{3.508201in}}%
\pgfpathmoveto{\pgfqpoint{5.381837in}{3.508201in}}%
\pgfpathlineto{\pgfqpoint{5.381837in}{3.508201in}}%
\pgfpathlineto{\pgfqpoint{5.381837in}{3.511150in}}%
\pgfpathlineto{\pgfqpoint{5.386378in}{3.511150in}}%
\pgfpathlineto{\pgfqpoint{5.386378in}{3.508201in}}%
\pgfpathmoveto{\pgfqpoint{5.381837in}{3.511150in}}%
\pgfpathlineto{\pgfqpoint{5.381837in}{3.511150in}}%
\pgfpathlineto{\pgfqpoint{5.381837in}{3.514099in}}%
\pgfpathlineto{\pgfqpoint{5.386378in}{3.514099in}}%
\pgfpathlineto{\pgfqpoint{5.386378in}{3.511150in}}%
\pgfpathmoveto{\pgfqpoint{5.386378in}{3.511150in}}%
\pgfpathlineto{\pgfqpoint{5.386378in}{3.511150in}}%
\pgfpathlineto{\pgfqpoint{5.386378in}{3.514099in}}%
\pgfpathlineto{\pgfqpoint{5.390919in}{3.514099in}}%
\pgfpathlineto{\pgfqpoint{5.390919in}{3.511150in}}%
\pgfpathmoveto{\pgfqpoint{5.386378in}{3.514099in}}%
\pgfpathlineto{\pgfqpoint{5.386378in}{3.514099in}}%
\pgfpathlineto{\pgfqpoint{5.386378in}{3.517048in}}%
\pgfpathlineto{\pgfqpoint{5.390919in}{3.517048in}}%
\pgfpathlineto{\pgfqpoint{5.390919in}{3.514099in}}%
\pgfpathmoveto{\pgfqpoint{5.390919in}{3.514099in}}%
\pgfpathlineto{\pgfqpoint{5.390919in}{3.514099in}}%
\pgfpathlineto{\pgfqpoint{5.390919in}{3.517048in}}%
\pgfpathlineto{\pgfqpoint{5.395460in}{3.517048in}}%
\pgfpathlineto{\pgfqpoint{5.395460in}{3.514099in}}%
\pgfpathmoveto{\pgfqpoint{5.390919in}{3.517048in}}%
\pgfpathlineto{\pgfqpoint{5.390919in}{3.517048in}}%
\pgfpathlineto{\pgfqpoint{5.390919in}{3.519998in}}%
\pgfpathlineto{\pgfqpoint{5.395460in}{3.519998in}}%
\pgfpathlineto{\pgfqpoint{5.395460in}{3.517048in}}%
\pgfpathmoveto{\pgfqpoint{5.395460in}{3.517048in}}%
\pgfpathlineto{\pgfqpoint{5.395460in}{3.517048in}}%
\pgfpathlineto{\pgfqpoint{5.395460in}{3.519998in}}%
\pgfpathlineto{\pgfqpoint{5.400001in}{3.519998in}}%
\pgfpathlineto{\pgfqpoint{5.400001in}{3.517048in}}%
\pgfpathclose%
\pgfusepath{fill}%
\end{pgfscope}%
\begin{pgfscope}%
\pgfsetbuttcap%
\pgfsetroundjoin%
\definecolor{currentfill}{rgb}{0.000000,0.000000,0.000000}%
\pgfsetfillcolor{currentfill}%
\pgfsetlinewidth{0.803000pt}%
\definecolor{currentstroke}{rgb}{0.000000,0.000000,0.000000}%
\pgfsetstrokecolor{currentstroke}%
\pgfsetdash{}{0pt}%
\pgfsys@defobject{currentmarker}{\pgfqpoint{0.000000in}{-0.048611in}}{\pgfqpoint{0.000000in}{0.000000in}}{%
\pgfpathmoveto{\pgfqpoint{0.000000in}{0.000000in}}%
\pgfpathlineto{\pgfqpoint{0.000000in}{-0.048611in}}%
\pgfusepath{stroke,fill}%
}%
\begin{pgfscope}%
\pgfsys@transformshift{1.215000in}{2.010000in}%
\pgfsys@useobject{currentmarker}{}%
\end{pgfscope}%
\end{pgfscope}%
\begin{pgfscope}%
\definecolor{textcolor}{rgb}{0.000000,0.000000,0.000000}%
\pgfsetstrokecolor{textcolor}%
\pgfsetfillcolor{textcolor}%
\pgftext[x=1.215000in,y=1.912778in,,top]{\color{textcolor}\sffamily\fontsize{10.000000}{12.000000}\selectfont −4}%
\end{pgfscope}%
\begin{pgfscope}%
\pgfsetbuttcap%
\pgfsetroundjoin%
\definecolor{currentfill}{rgb}{0.000000,0.000000,0.000000}%
\pgfsetfillcolor{currentfill}%
\pgfsetlinewidth{0.803000pt}%
\definecolor{currentstroke}{rgb}{0.000000,0.000000,0.000000}%
\pgfsetstrokecolor{currentstroke}%
\pgfsetdash{}{0pt}%
\pgfsys@defobject{currentmarker}{\pgfqpoint{0.000000in}{-0.048611in}}{\pgfqpoint{0.000000in}{0.000000in}}{%
\pgfpathmoveto{\pgfqpoint{0.000000in}{0.000000in}}%
\pgfpathlineto{\pgfqpoint{0.000000in}{-0.048611in}}%
\pgfusepath{stroke,fill}%
}%
\begin{pgfscope}%
\pgfsys@transformshift{2.145000in}{2.010000in}%
\pgfsys@useobject{currentmarker}{}%
\end{pgfscope}%
\end{pgfscope}%
\begin{pgfscope}%
\definecolor{textcolor}{rgb}{0.000000,0.000000,0.000000}%
\pgfsetstrokecolor{textcolor}%
\pgfsetfillcolor{textcolor}%
\pgftext[x=2.145000in,y=1.912778in,,top]{\color{textcolor}\sffamily\fontsize{10.000000}{12.000000}\selectfont −2}%
\end{pgfscope}%
\begin{pgfscope}%
\pgfsetbuttcap%
\pgfsetroundjoin%
\definecolor{currentfill}{rgb}{0.000000,0.000000,0.000000}%
\pgfsetfillcolor{currentfill}%
\pgfsetlinewidth{0.803000pt}%
\definecolor{currentstroke}{rgb}{0.000000,0.000000,0.000000}%
\pgfsetstrokecolor{currentstroke}%
\pgfsetdash{}{0pt}%
\pgfsys@defobject{currentmarker}{\pgfqpoint{0.000000in}{-0.048611in}}{\pgfqpoint{0.000000in}{0.000000in}}{%
\pgfpathmoveto{\pgfqpoint{0.000000in}{0.000000in}}%
\pgfpathlineto{\pgfqpoint{0.000000in}{-0.048611in}}%
\pgfusepath{stroke,fill}%
}%
\begin{pgfscope}%
\pgfsys@transformshift{3.075000in}{2.010000in}%
\pgfsys@useobject{currentmarker}{}%
\end{pgfscope}%
\end{pgfscope}%
\begin{pgfscope}%
\definecolor{textcolor}{rgb}{0.000000,0.000000,0.000000}%
\pgfsetstrokecolor{textcolor}%
\pgfsetfillcolor{textcolor}%
\pgftext[x=3.075000in,y=1.912778in,,top]{\color{textcolor}\sffamily\fontsize{10.000000}{12.000000}\selectfont 0}%
\end{pgfscope}%
\begin{pgfscope}%
\pgfsetbuttcap%
\pgfsetroundjoin%
\definecolor{currentfill}{rgb}{0.000000,0.000000,0.000000}%
\pgfsetfillcolor{currentfill}%
\pgfsetlinewidth{0.803000pt}%
\definecolor{currentstroke}{rgb}{0.000000,0.000000,0.000000}%
\pgfsetstrokecolor{currentstroke}%
\pgfsetdash{}{0pt}%
\pgfsys@defobject{currentmarker}{\pgfqpoint{0.000000in}{-0.048611in}}{\pgfqpoint{0.000000in}{0.000000in}}{%
\pgfpathmoveto{\pgfqpoint{0.000000in}{0.000000in}}%
\pgfpathlineto{\pgfqpoint{0.000000in}{-0.048611in}}%
\pgfusepath{stroke,fill}%
}%
\begin{pgfscope}%
\pgfsys@transformshift{4.005000in}{2.010000in}%
\pgfsys@useobject{currentmarker}{}%
\end{pgfscope}%
\end{pgfscope}%
\begin{pgfscope}%
\definecolor{textcolor}{rgb}{0.000000,0.000000,0.000000}%
\pgfsetstrokecolor{textcolor}%
\pgfsetfillcolor{textcolor}%
\pgftext[x=4.005000in,y=1.912778in,,top]{\color{textcolor}\sffamily\fontsize{10.000000}{12.000000}\selectfont 2}%
\end{pgfscope}%
\begin{pgfscope}%
\pgfsetbuttcap%
\pgfsetroundjoin%
\definecolor{currentfill}{rgb}{0.000000,0.000000,0.000000}%
\pgfsetfillcolor{currentfill}%
\pgfsetlinewidth{0.803000pt}%
\definecolor{currentstroke}{rgb}{0.000000,0.000000,0.000000}%
\pgfsetstrokecolor{currentstroke}%
\pgfsetdash{}{0pt}%
\pgfsys@defobject{currentmarker}{\pgfqpoint{0.000000in}{-0.048611in}}{\pgfqpoint{0.000000in}{0.000000in}}{%
\pgfpathmoveto{\pgfqpoint{0.000000in}{0.000000in}}%
\pgfpathlineto{\pgfqpoint{0.000000in}{-0.048611in}}%
\pgfusepath{stroke,fill}%
}%
\begin{pgfscope}%
\pgfsys@transformshift{4.935000in}{2.010000in}%
\pgfsys@useobject{currentmarker}{}%
\end{pgfscope}%
\end{pgfscope}%
\begin{pgfscope}%
\definecolor{textcolor}{rgb}{0.000000,0.000000,0.000000}%
\pgfsetstrokecolor{textcolor}%
\pgfsetfillcolor{textcolor}%
\pgftext[x=4.935000in,y=1.912778in,,top]{\color{textcolor}\sffamily\fontsize{10.000000}{12.000000}\selectfont 4}%
\end{pgfscope}%
\begin{pgfscope}%
\definecolor{textcolor}{rgb}{0.000000,0.000000,0.000000}%
\pgfsetstrokecolor{textcolor}%
\pgfsetfillcolor{textcolor}%
\pgftext[x=5.400000in,y=1.722809in,,top]{\color{textcolor}\sffamily\fontsize{10.000000}{12.000000}\selectfont x}%
\end{pgfscope}%
\begin{pgfscope}%
\pgfsetbuttcap%
\pgfsetroundjoin%
\definecolor{currentfill}{rgb}{0.000000,0.000000,0.000000}%
\pgfsetfillcolor{currentfill}%
\pgfsetlinewidth{0.803000pt}%
\definecolor{currentstroke}{rgb}{0.000000,0.000000,0.000000}%
\pgfsetstrokecolor{currentstroke}%
\pgfsetdash{}{0pt}%
\pgfsys@defobject{currentmarker}{\pgfqpoint{-0.048611in}{0.000000in}}{\pgfqpoint{-0.000000in}{0.000000in}}{%
\pgfpathmoveto{\pgfqpoint{-0.000000in}{0.000000in}}%
\pgfpathlineto{\pgfqpoint{-0.048611in}{0.000000in}}%
\pgfusepath{stroke,fill}%
}%
\begin{pgfscope}%
\pgfsys@transformshift{3.075000in}{0.802000in}%
\pgfsys@useobject{currentmarker}{}%
\end{pgfscope}%
\end{pgfscope}%
\begin{pgfscope}%
\definecolor{textcolor}{rgb}{0.000000,0.000000,0.000000}%
\pgfsetstrokecolor{textcolor}%
\pgfsetfillcolor{textcolor}%
\pgftext[x=2.773039in, y=0.749238in, left, base]{\color{textcolor}\sffamily\fontsize{10.000000}{12.000000}\selectfont −4}%
\end{pgfscope}%
\begin{pgfscope}%
\pgfsetbuttcap%
\pgfsetroundjoin%
\definecolor{currentfill}{rgb}{0.000000,0.000000,0.000000}%
\pgfsetfillcolor{currentfill}%
\pgfsetlinewidth{0.803000pt}%
\definecolor{currentstroke}{rgb}{0.000000,0.000000,0.000000}%
\pgfsetstrokecolor{currentstroke}%
\pgfsetdash{}{0pt}%
\pgfsys@defobject{currentmarker}{\pgfqpoint{-0.048611in}{0.000000in}}{\pgfqpoint{-0.000000in}{0.000000in}}{%
\pgfpathmoveto{\pgfqpoint{-0.000000in}{0.000000in}}%
\pgfpathlineto{\pgfqpoint{-0.048611in}{0.000000in}}%
\pgfusepath{stroke,fill}%
}%
\begin{pgfscope}%
\pgfsys@transformshift{3.075000in}{1.406000in}%
\pgfsys@useobject{currentmarker}{}%
\end{pgfscope}%
\end{pgfscope}%
\begin{pgfscope}%
\definecolor{textcolor}{rgb}{0.000000,0.000000,0.000000}%
\pgfsetstrokecolor{textcolor}%
\pgfsetfillcolor{textcolor}%
\pgftext[x=2.773039in, y=1.353238in, left, base]{\color{textcolor}\sffamily\fontsize{10.000000}{12.000000}\selectfont −2}%
\end{pgfscope}%
\begin{pgfscope}%
\pgfsetbuttcap%
\pgfsetroundjoin%
\definecolor{currentfill}{rgb}{0.000000,0.000000,0.000000}%
\pgfsetfillcolor{currentfill}%
\pgfsetlinewidth{0.803000pt}%
\definecolor{currentstroke}{rgb}{0.000000,0.000000,0.000000}%
\pgfsetstrokecolor{currentstroke}%
\pgfsetdash{}{0pt}%
\pgfsys@defobject{currentmarker}{\pgfqpoint{-0.048611in}{0.000000in}}{\pgfqpoint{-0.000000in}{0.000000in}}{%
\pgfpathmoveto{\pgfqpoint{-0.000000in}{0.000000in}}%
\pgfpathlineto{\pgfqpoint{-0.048611in}{0.000000in}}%
\pgfusepath{stroke,fill}%
}%
\begin{pgfscope}%
\pgfsys@transformshift{3.075000in}{2.010000in}%
\pgfsys@useobject{currentmarker}{}%
\end{pgfscope}%
\end{pgfscope}%
\begin{pgfscope}%
\definecolor{textcolor}{rgb}{0.000000,0.000000,0.000000}%
\pgfsetstrokecolor{textcolor}%
\pgfsetfillcolor{textcolor}%
\pgftext[x=2.889413in, y=1.957238in, left, base]{\color{textcolor}\sffamily\fontsize{10.000000}{12.000000}\selectfont 0}%
\end{pgfscope}%
\begin{pgfscope}%
\pgfsetbuttcap%
\pgfsetroundjoin%
\definecolor{currentfill}{rgb}{0.000000,0.000000,0.000000}%
\pgfsetfillcolor{currentfill}%
\pgfsetlinewidth{0.803000pt}%
\definecolor{currentstroke}{rgb}{0.000000,0.000000,0.000000}%
\pgfsetstrokecolor{currentstroke}%
\pgfsetdash{}{0pt}%
\pgfsys@defobject{currentmarker}{\pgfqpoint{-0.048611in}{0.000000in}}{\pgfqpoint{-0.000000in}{0.000000in}}{%
\pgfpathmoveto{\pgfqpoint{-0.000000in}{0.000000in}}%
\pgfpathlineto{\pgfqpoint{-0.048611in}{0.000000in}}%
\pgfusepath{stroke,fill}%
}%
\begin{pgfscope}%
\pgfsys@transformshift{3.075000in}{2.614000in}%
\pgfsys@useobject{currentmarker}{}%
\end{pgfscope}%
\end{pgfscope}%
\begin{pgfscope}%
\definecolor{textcolor}{rgb}{0.000000,0.000000,0.000000}%
\pgfsetstrokecolor{textcolor}%
\pgfsetfillcolor{textcolor}%
\pgftext[x=2.889413in, y=2.561238in, left, base]{\color{textcolor}\sffamily\fontsize{10.000000}{12.000000}\selectfont 2}%
\end{pgfscope}%
\begin{pgfscope}%
\pgfsetbuttcap%
\pgfsetroundjoin%
\definecolor{currentfill}{rgb}{0.000000,0.000000,0.000000}%
\pgfsetfillcolor{currentfill}%
\pgfsetlinewidth{0.803000pt}%
\definecolor{currentstroke}{rgb}{0.000000,0.000000,0.000000}%
\pgfsetstrokecolor{currentstroke}%
\pgfsetdash{}{0pt}%
\pgfsys@defobject{currentmarker}{\pgfqpoint{-0.048611in}{0.000000in}}{\pgfqpoint{-0.000000in}{0.000000in}}{%
\pgfpathmoveto{\pgfqpoint{-0.000000in}{0.000000in}}%
\pgfpathlineto{\pgfqpoint{-0.048611in}{0.000000in}}%
\pgfusepath{stroke,fill}%
}%
\begin{pgfscope}%
\pgfsys@transformshift{3.075000in}{3.218000in}%
\pgfsys@useobject{currentmarker}{}%
\end{pgfscope}%
\end{pgfscope}%
\begin{pgfscope}%
\definecolor{textcolor}{rgb}{0.000000,0.000000,0.000000}%
\pgfsetstrokecolor{textcolor}%
\pgfsetfillcolor{textcolor}%
\pgftext[x=2.889413in, y=3.165238in, left, base]{\color{textcolor}\sffamily\fontsize{10.000000}{12.000000}\selectfont 4}%
\end{pgfscope}%
\begin{pgfscope}%
\definecolor{textcolor}{rgb}{0.000000,0.000000,0.000000}%
\pgfsetstrokecolor{textcolor}%
\pgfsetfillcolor{textcolor}%
\pgftext[x=2.717483in,y=3.520000in,,bottom,rotate=90.000000]{\color{textcolor}\sffamily\fontsize{10.000000}{12.000000}\selectfont y}%
\end{pgfscope}%
\begin{pgfscope}%
\pgfsetrectcap%
\pgfsetmiterjoin%
\pgfsetlinewidth{0.803000pt}%
\definecolor{currentstroke}{rgb}{0.000000,0.000000,0.000000}%
\pgfsetstrokecolor{currentstroke}%
\pgfsetdash{}{0pt}%
\pgfpathmoveto{\pgfqpoint{3.075000in}{0.500000in}}%
\pgfpathlineto{\pgfqpoint{3.075000in}{3.520000in}}%
\pgfusepath{stroke}%
\end{pgfscope}%
\begin{pgfscope}%
\pgfsetrectcap%
\pgfsetmiterjoin%
\pgfsetlinewidth{0.000000pt}%
\definecolor{currentstroke}{rgb}{0.000000,0.000000,0.000000}%
\pgfsetstrokecolor{currentstroke}%
\pgfsetstrokeopacity{0.000000}%
\pgfsetdash{}{0pt}%
\pgfpathmoveto{\pgfqpoint{5.400000in}{0.500000in}}%
\pgfpathlineto{\pgfqpoint{5.400000in}{3.520000in}}%
\pgfusepath{}%
\end{pgfscope}%
\begin{pgfscope}%
\pgfsetrectcap%
\pgfsetmiterjoin%
\pgfsetlinewidth{0.803000pt}%
\definecolor{currentstroke}{rgb}{0.000000,0.000000,0.000000}%
\pgfsetstrokecolor{currentstroke}%
\pgfsetdash{}{0pt}%
\pgfpathmoveto{\pgfqpoint{0.750000in}{2.010000in}}%
\pgfpathlineto{\pgfqpoint{5.400000in}{2.010000in}}%
\pgfusepath{stroke}%
\end{pgfscope}%
\begin{pgfscope}%
\pgfsetrectcap%
\pgfsetmiterjoin%
\pgfsetlinewidth{0.000000pt}%
\definecolor{currentstroke}{rgb}{0.000000,0.000000,0.000000}%
\pgfsetstrokecolor{currentstroke}%
\pgfsetstrokeopacity{0.000000}%
\pgfsetdash{}{0pt}%
\pgfpathmoveto{\pgfqpoint{0.750000in}{3.520000in}}%
\pgfpathlineto{\pgfqpoint{5.400000in}{3.520000in}}%
\pgfusepath{}%
\end{pgfscope}%
\end{pgfpicture}%
\makeatother%
\endgroup%
 \end{solution}
\part[]  $\left\{ {\begin{matrix}
   {y \leqslant 2}  \\ 
   {x + y \leqslant 3}  \\ 
   {x \geqslant 0}  \\ 
   {y \geqslant 0}  \\ 

 \end{matrix} } \right.$
\part[] $\left\{ {\begin{matrix}
   {6x - 5y \leqslant 30}  \\ 
   {4x + 3y \leqslant 0}  \\ 

 \end{matrix} } \right.$
\begin{solution} \end{solution}
\part[] $\left\{ {\begin{matrix}
   {x + y \leqslant 5}  \\ 
   {x + 3y \geqslant 9}  \\ 
   {y \geqslant 0}  \\ 
   {x \geqslant 0}  \\ 
   {y \geqslant 0}  \\ 

 \end{matrix} } \right.$
\begin{solution} \end{solution}

\begin{solution} \end{solution}
\end{parts}
\end{multicols}

\question Resuelve los siguientes problemas
\begin{multicols}{2}
\begin{parts}
\part[] La  base  de  un  rectángulo  excede  en 3 cm a  la altura. Si  elevamos  al  cuadrado  su  diagonal  el  resultado  es menor que 29. Determina las posibles dimensiones del rectángulo. 
\begin{solution} $\left\{ {\begin{matrix}
   {x=y+3}  \\ 
   {x^2+y^2<29}  \\ 
 \end{matrix} } \right. \to x\in \left(-2, 5\right) \land y=x-3 \to y \in \left(0,2\right) \land x=y+3$\end{solution}
\part[]  Una fábrica A paga a sus viajantes 10 euros por artículo vendido más una cantidad fija de 400 euros. Otra fábrica B paga  15  euros  por  artículo  vendido  más  una  cantidad  fija  de  300  euros.  ¿Cuántos artículos  debe vender el viajante de fábrica B para ganar más dinero que el viajante de la fábrica A?
\begin{solution}$300+15x>400+10x\to\left(20, +\infty\right)\to$ más de 20  \end{solution}
\part[]  Se ha medido un campo rectangular con un error menor que un metro. Los valores de los lados, $x$ e $y$, cumplen $70 < x < 71$ y $47 < y < 48$. Determina entre qué valores está el perímetro y el área.
\begin{solution}%234<p<236 y 3290<a<3408 
$ 234<p<238$  y  $3290<a<3408$\end{solution}
\part[]  Un alumno ha realizado dos exámenes de Matemáticas obteniendo calificaciones respectivas de 4,5 y 5,7 puntos. ¿Cuánto ha de sacar como mínimo en el tercero para aprobar, si la nota final es la media aritmética de las tres notas? ¿Y si el primer examen cuenta un 15\%, el segundo un 35\% y el tercero un 50\%?
\begin{solution}%latex(reduce_rational_inequalities([[(4.5+5.7+x)/3>5]], x,relational=0)) y latex(reduce_rational_inequalities([[4.5*0.15+5.7*0.35+0.5*x>5]], x,relational=0))
$\frac{4,5+5,7+x}{3}>\to\left(4.8, \infty\right)$ y  $\left(4.66, \infty\right) $
\end{solution}
\part[]  Entre los triángulos isósceles de lado desigual igual a 20 cm., ¿cuáles tienen perímetro inferior a 200 cm.?
\begin{solution} $ 20+2x<200 \to x<90 \to$ lados iguales menores que 90\end{solution}
\part[]  Un padre tiene 33 años más que su hijo y el abuelo 33 años más que el padre. Hace tres años, sus edades sumaban menos de un siglo. ¿Qué edad puede tener cada uno?
\begin{solution} $x+\left(x+33\right)+\left(x+66\right)-3\cdot3<100 \to x<\frac{10}{3} \to$ Menos de 3 años y 4 meses el hijo\end{solution}

\end{parts}
\end{multicols}

\begin{comment}
\question 
\begin{multicols}{3}
\begin{parts}
\part[]  
\begin{solution} \end{solution}
\end{parts}
\end{multicols}
\end{comment}



\end{questions}

\end{document}


